\documentclass[11pt,a4paper]{ltjsarticle}
\usepackage[margin=24mm]{geometry}
\usepackage{amsmath,amssymb,amsthm,mathtools,booktabs,longtable,array}
\usepackage[haranoaji]{luatexja-preset}
\usepackage[unicode,hidelinks]{hyperref}
\usepackage{xurl}
\newtheorem{theorem}{定理}[section]
\newtheorem{proposition}[theorem]{命題}
\newtheorem{lemma}[theorem]{補題}
\newtheorem{corollary}[theorem]{系}
\theoremstyle{definition}
\newtheorem{definition}[theorem]{定義}
\newtheorem{remark}[theorem]{注意}
\newtheorem{conjecture}[theorem]{予想}
\renewcommand{\proofname}{証明}
\newcommand{\ZZ}{\mathbb Z}
\newcommand{\QQ}{\mathbb Q}
\newcommand{\GL}{\operatorname{GL}}
\newcommand{\PGL}{\operatorname{PGL}}
\newcommand{\Gr}{\operatorname{Gr}}
\newcommand{\Aut}{\operatorname{Aut}}
\newcommand{\Per}{\operatorname{Per}}
\newcommand{\qgr}{\operatorname{qgr}}
\newcommand{\rk}{\operatorname{rank}}
\newcommand{\Spec}{\operatorname{Spec}}
\newcommand{\Sym}{\operatorname{Sym}}
\newcommand{\HH}{\operatorname{HH}}
\newcommand{\NF}{\operatorname{NF}}
\title{標準14641型AS正則$\mathbb Z$-代数の完全分類に向けて\\
\large 有限表示・構造分岐・周期構造と代数的変形}
\author{}
\date{2026年11月27日\quad 研究原稿 v516\\
\small 全分類およびClifford点の高次障害の決定は未完成}
\begin{document}
\maketitle
\begin{abstract}
標準$(1,4,6,4,1)$型の四次元二次Koszul AS正則$\ZZ$-代数について、
添字保持同型による分類の有限表示と、graded代数には由来しない代数的族を研究する。
Frobenius traceにより標準4周期構造を構成し、幾何学的同型類・周期判定を四窓へ還元する。
連続する二つの関係空間から標準延長を復元し、全次数のKoszul性を56次元切断代数の3-representation infinite条件に帰着する。
具体的な正則族として、三生成Clifford部分を動かす九変数の楕円Ore族を構成する。
全ファイバーのAS正則性、同型類の像の次元9、一般点の最小周期2を証明する。
非1周期性は楕円曲線上の12点因子と非零2等分点による平行移動を用いて示す。

有限四辺の点スキームが三次元接点を持つ枝では、三生成基底を内在的に復元し、
滑らかな楕円・節点・横断二成分・三角形の各二型について、特殊値を含む正則性、添字保持同型、周期と平滑化を決定する。

局所変形では、量子多項式族のgraded同型類6次元からindexed同型類3次元へのtwist忘却と、
Clifford点における接空間9次元から15次元への増大を区別する。
後者の全二次障害は、六つの交代変数上で非退化な二次式一つとなる。
九変数族の内在的な2-isogenyから次数4の必要同型不変量を得る。
さらに別のOre表示を持つ境界族を構成し、七・九変数族の有限関係空間における
軌道閉包がClifford点で特異となることを示す。

有限線型代数と多項式恒等式には再現可能な厳密計算を付す。
三因子対称化可能な枝では、基点の存在から非線型Torを強制する左右反転対合を構成し、
その枝の全AS対象が基点自由な完全対称完全交叉チャートに属することを証明する。
中央三次数の加群の剛性から、対称完全交叉チャート全体が三つの二対二階数を6以下とする表示空間の滑らかな開部分であることを示す。
一つの可逆な交代化証明書を持つ有限Frobenius候補についても、Klein二次曲面と六次元Clifford代数で三正規形に還元し、
可換多項式のindexed類だけが有限Frobenius条件を満たすことを示す。
対称成分の境界には、四位置を三対一または二対二に分ける七つの68次元因子を構成する。
さらに、特異な同時移送写像と四周期二次商の指数成長を用いて、
対称閉包内のAS点が対称完全交叉チャートに尽きることを証明する。
一般四窓でも、連続する純粋二次関係から指数成長を導き、標準AS候補を排除できる。
九次元の族の像の成分としての最大性、完全な族内同型判定、
Clifford点の高次障害、全AS正則locusの分類および新規性の最終確認は未完成である。
\end{abstract}
\tableofcontents

\section{対象・比較する同値関係・先行研究}
本稿を通じて$k$は標数零の代数閉体とする。
正連結な局所有限$\ZZ$-代数$C=(C_{ij})$とは
\[
 C_{ii}=k,\qquad C_{ij}=0\ (i>j),\qquad C_{ij}C_{j\ell}\subset C_{i\ell}
\]
を満たすものをいう。$C(s)_{ij}=C_{i+s,j+s}$とし、
\[
 \Per(C)=\{s\in\ZZ:C\simeq C(s)\text{ が添字保持同型で成立する}\}
\]
を周期群とする。通常のconnected graded algebra $A$からは
$\check A_{ij}=A_{j-i}$を作る。

中心対象は\emph{標準14641型}、すなわち次数一生成、二次関係を持ち、
各右単純加群$S_i$に有限最小線型解消
\begin{equation}\label{eq:standard-resolution}
0\longrightarrow P_{i+4}\longrightarrow P_{i+3}^{4}
\longrightarrow P_{i+2}^{6}\longrightarrow P_{i+1}^{4}
\longrightarrow P_i\longrightarrow S_i\longrightarrow0
\end{equation}
があり、左側の対応する解消とAS双対性を持つものとする。
ここで$P_i=e_iC$である。行Hilbert級数は$(1-t)^{-4}$である。
domain、noetherian、PBW、Calabi--Yau性を全対象の追加仮定としない。
一般の四次元quadratic AS正則性からこの数値型が自動的に従うかは別問題とする。

\paragraph{二つの分類問題。}
$\mathcal M_{\rm gr}$ではgraded algebraのgraded同型を、
$\mathcal M_{\ZZ}$では添字保持$\ZZ$-代数同型を使う。
忘却$F:A\mapsto\check A$の像は1周期的部分であるが、
$F$は一般に同型類上で単射でない。
Zhang twistとgraded module圏、$\ZZ$-代数同型との関係は
\cite{Zhang,Sierra,VdB}を参照する。
$qgr$同値、導来同値、添字平行移動をさらに割る同値は、本稿の同値関係ではない。
また本稿でいう\emph{同型類の次元}は幾何学的商・有理商の次元であり、
自己同型群を含めて差し引くstackの次元とは区別する。

\paragraph{既存の周期性と本稿の有限表示。}
Koszul--Frobenius双対性とhelixの方法\cite{BP}、
正則$\ZZ$-空間の圏的構成\cite{MoriNyman}を出発点とする。
標準14641型の一般4周期性は本稿でもtraceにより説明する。
最小周期2の平坦PBW族と最小周期4の例は私的研究原稿\cite{as4period}にある。
その全証明の今回の独立監査、第三者査読、新規性の確定は主張しない。
同原稿に依存する箇所は明記する。

\paragraph{局所モジュライの既知研究との関係。}
Bhatoy--Ingalls--LaRoche--Nookala\cite{BILN}は四次元Koszul AS正則graded代数の
Hochschild cohomologyとKodaira--Spencer写像を調べている。
同論文\S3.1の量子多項式族とTable 2のgraded Clifford族は本稿の比較基準になる。
本稿の追加計算は\emph{周期的な添字依存のある}関係式を許し、
$\zeta=1,-1,\sqrt{-1},-\sqrt{-1}$の四つのFourier成分を分ける点にある。
graded接空間の次元だけで$\ZZ$-代数の変形を決定したとはしない。
新規性の最終確認は未了である。

\paragraph{三次元三次の場合との違い。}
Van den Bergh\cite[Theorem 5.5.10, \S5.6]{VdB}の非可換quadricは、
楕円曲線と次数2の三つの直線束$(E,L_0,L_1,L_2)$で記述される。
滑らかな楕円曲線の一般部分でのパラメータ数は$1+3-1=3$である。
一般の$j$について1周期性は$L_0\otimes L_2\simeq L_1^{\otimes2}$を課し、
2次元の部分になる。ただし特殊な$j$までこの等式だけを必要条件とするのは正しくない。
正確には\cite[\S5.6]{VdB}の自己同型$\sigma$を用いた条件を使う。
この三次元三次の結果を、四次元二次の場合の成分次元にそのまま移してはならない。

本稿の代数的構成（第\ref{sec:algebraic-family}節）は、中心二次式と通常の導分を
偶奇に応じて組み合わせたOre型の構成であり、基底の存在にはDiamond lemmaを使う。
PBW性からKoszul性への含意はPriddyの理論\cite{Priddy}による。
本稿で個別に証明するのは七変数および九変数の具体的関係、全ファイバーのFrobenius双対性、
隣接する点射影の配置による一般点の非1周期性、および同型による重複を除いた次元である。
九変数への拡張では、二次完全交叉のKoszul双対とClifford基底、Ore拡大を使う。
正則性の標準的な道具を新規結果とはせず、楕円曲線の対合が射影変換から来ないことと
12点配置を組み合わせた非1周期性の証明を第\ref{sec:elliptic-ore}節に置く。
これらの具体的族と既存の非可換射影三空間の族との完全な文献上の同定・新規性確認は未了である。

\paragraph{Clifford/Ore構成と幾何的不変量の先行研究。}
Vancliff--Van Rompay--Willaert\cite[\S3]{VVRW}はgraded Clifford代数の
iterated Ore変形と点加群・直線加群を調べている。
Stephenson--Vancliff\cite{SV}も三生成Clifford部分と平面三次曲線を用いる
四次元graded Clifford代数の構成を扱う。
従ってClifford部分・Ore拡大・点配置という方法自体を本稿の新規性とはしない。
本稿で検討すべき追加の寄与は、添字に依存する交代変形、
一般点での非1周期性とindexed同型類像の次元、および境界の具体的な幾何である。
両文献の構成と本稿の全特殊化の完全な同定は未了である。
またBhatoy--Ingalls\cite{BI}は既知の四次元graded族のsuperpotentialから離散不変量を構成する。
これらはgraded側の族との比較に有用だが、本稿のindexed同型関係や全成分数へそのまま適用しない。

対称Frobeniusテンソル部分では、二次完全交叉の既知のClifford手法により正則性を一括認証する。
対称テンソルの独立同値を合同へ直す一般原理は、Belitskii--Sergeichuk
\cite[Theorem 2(a), Lemma 4(a)]{BS2006}の対称同値・行列根の理論と関係する。
本稿は標数零の代数閉体上の必要な四階の場合を直接証明し、
対称チャート内のindexed同型を四二次式の通常の線型同値へ還元する。
この合同原理自体の新規性は主張しない。

対称チャートの最上traceは完全交叉のMacaulay逆系、すなわちassociated formでもある。
その安定性と射の幾何についてはFedorchuk--Isaev\cite{FIstability,FIassociated}を参照する。
ただし本稿の四因子独立作用に対する閉包を、対角群の結果だけから同定してはいけない。
命題\ref{prop:symmetric-regular-specialization}は認証済みチャート内部でも個別軌道が閉じない例を与える。

\subsection*{主結果の位置と読み方}
証明の依存順は、有限四窓、Frobeniusテンソルと被覆分岐、局所変形、代数的族、同型不変量、境界である。
結果の種類と未解決部分を次のように区別する。
\begin{longtable}{p{.24\linewidth}p{.69\linewidth}}
\toprule
内容 & 主な結果と射程\\
\midrule
\endhead
有限表示 &
定理\ref{thm:canonical-four}、\ref{thm:finite-period-test}。
AS正則な標準型対象の同型・周期を有限表示へ還元する。
一般四窓のAS正則性まで判定したものではない。\\
対称閉包のAS点 &
定理\ref{thm:symmetric-closure-as-exact}により
$\mathrm{AS}\cap Z_{\rm sym}=S$。
この閉包外の低階数・高階数領域の被覆は未完成。\\
二つの次元比較 &
定理\ref{thm:quantum-orbits}、\ref{thm:quantum-local}の$6\to3$は同型類の局所成分。
定理\ref{thm:clifford-tangent}の$9\to15$は接空間。
二次障害は定理\ref{thm:clifford-quadratic}、全局所模型は予想\ref{conj:clifford-model}。\\
認証済みの代数的族 &
定理\ref{thm:algebraic-seven}、\ref{thm:algebraic-nine}。
全ファイバー正則性、同型類像の次元、一般最小周期2を別々に証明する。
七変数族の証明も保持する。\\
同型への手掛かり &
定理\ref{thm:ore-derivations}は指定Ore構成の導分分類。
補題\ref{lem:intrinsic-elliptic-tau}と定理\ref{thm:explicit-two-isogeny}は
内在的データと必要不変量を与える。完全なindexed同型分類は未完成。\\
境界の幾何 &
定理\ref{thm:nonore-in-seven-closure}の閉包包含、
定理\ref{thm:seven-closure-singular}、\ref{thm:nine-closure-singular}の特異性、
命題\ref{prop:nine-closure-obstruction}の同時積分障害。
対象は位置0の関係空間の軌道閉包である。\\
\bottomrule
\end{longtable}

有限Frobenius性の十分性については、定理\ref{thm:nonkoszul-symmetric-boundary}で
四次Koszul検査を全て通る具体的な非Koszul境界を構成する。
その五次ホモロジーは一始点で7次元となり、一般四窓の正則性判定に残る問題を具体化する。

\section{周期群と4周期の有限表示}
\begin{lemma}\label{lem:period-subgroup}
$\Per(C)$は$\ZZ$の部分群である。$C$が4周期的なら最小正周期は$1,2,4$のいずれかである。
\end{lemma}
\begin{proof}
$s$周期同型と$t$周期同型は、後者を$s$だけ平行移動して合成できる。
逆同型を$-s$だけ移動すれば$-s$も周期となる。$0$は自明に周期である。
従って$\Per(C)=p\ZZ$であり、$4\in p\ZZ$なら$p\mid4$である。
\end{proof}

4周期同型付き対象$(C,\tau)$とは、積を保つ同型
$\tau_{ij}:C_{ij}\to C_{i+4,j+4}$を選んだ組をいう。
対象の同型は$\tau$と可換する添字保持同型とする。
四つの4次元空間$V_i$を固定し、添字を法4で読むと、
\[
 \mathfrak P=\prod_{i=0}^3\Gr(6,V_i\otimes V_{i+1}),\qquad
 G=\prod_{i=0}^3\GL(V_i)
\]
が有限表示の空間と基底変更群になる。

\begin{theorem}[4周期表示の有限四窓atlas]\label{thm:four-window}
各一次空間が4次元、各二次関係空間が6次元の次数一生成・二次正連結
$\ZZ$-代数$(C,\tau)$の同型群oidは、$\mathfrak P(k)$に対する$G$の作用群oidと同値である。
Koszul性やAS正則性を仮定しなくても成立し、それらを満たす全部分群oidへ制限できる。
\end{theorem}
\begin{proof}
$i=0,1,2,3$で$C_{i,i+1}$の基底を選び、$\tau$の反復で全整数へ延長する。
\[
 R_i=\ker(C_{i,i+1}\otimes C_{i+1,i+2}\longrightarrow C_{i,i+2})
\]
は$R_{i+4}=R_i$を満たすので$\mathfrak P$の点になる。
逆に四窓$(R_i)$から、整数頂点列上の自由経路$\ZZ$-代数を、各位置の$R_i$が生成する
二次両側イデアルで割る。二次成分の核は指定した$R_i$に等しく、厳密4移動が周期同型を与える。
基底変更は$G$の作用である。
$\tau$と可換する同型は四つの可逆写像$g_i:V_i\to V'_i$で決まり、
二次関係の保存は$(g_i\otimes g_{i+1})R_i=R'_i$と同値である。
次数一生成・二次性により、この条件は商代数の同型の十分条件でもある。
合成も一致するので群oid同値を得る。
\end{proof}

\begin{remark}\label{rem:forget4}
一般の4周期同型を選んだ$(C,\tau)$と、裸の$C$は異なる分類問題である。
$\mathfrak P$全体がAS正則とも、$[\mathfrak P/G]$が裸のAS正則代数のモジュライとも主張しない。
$\dim\mathfrak P=240$、$\dim G=64$という数だけからモジュライ次元を計算してはならない。
\end{remark}

\begin{proposition}\label{prop:period-relations}
二次$\ZZ$-代数$C$について$s\in\Per(C)$となることは、全$i\in\ZZ$に可逆写像
$g_i:V_i\to V_{i+s}$があり、$(g_i\otimes g_{i+1})R_i=R_{i+s}$を満たすことと同値である。
\end{proposition}
\begin{proof}
同型の次数一への制限から必要性を得る。逆に指定した可逆写像は自由経路代数の同型を定め、
二次イデアルの生成空間を対応する生成空間へ送るため商代数へ降りる。
\end{proof}
\begin{remark}
この$g_i$は、任意に選んだ4周期同型と可換する必要がない。
初めから$g_{i+4}=g_i$だけを調べると周期同型を見落とし得る。
次節のFrobenius標準化がこの問題を解決する。
\end{remark}

\section{Koszul--Frobenius標準4周期構造と忘却}\label{sec:canonical-four}
前向き二次双対を$E=C^!$とし、$E_{i,i+1}=C_{i,i+1}^*$の規約を使う。
標準14641型の左右AS双対性とKoszul性から、
\[
 \dim E_{i,i+p}={4\choose p}\ (0\le p\le4),\quad E_{i,i+p}=0\ (p>4)
\]
であり、相補次数の積$E_{i,i+p}\otimes E_{i+p,i+4}\to E_{i,i+4}$は完全対合である。
これは有限Koszul解消のAS双対性をYoneda積に移したものである\cite{BP,as4period}。

非零trace $\lambda_i:E_{i,i+4}\to k$を選ぶと、一意の線型同型
$\nu^\lambda_{ij}:E_{ij}\to E_{i+4,j+4}$が
\begin{equation}\label{eq:nakayama-pair}
 \lambda_i(ab)=\lambda_j(b\nu^\lambda_{ij}(a)),\qquad b\in E_{j,i+4}
\end{equation}
で定まる。$a\in E_{ij},c\in E_{j\ell},b\in E_{\ell,i+4}$に対して
\[
 \lambda_i((ac)b)=\lambda_j((cb)\nu(a))
 =\lambda_\ell(b\nu(a)\nu(c))
\]
だから$\nu(ac)=\nu(a)\nu(c)$である。
次数一部分の逆双対$(\nu^{-1}_{i,i+1})^*$は元の二次関係を保存し、
$\tau^\lambda:C\simeq C(4)$を与える。これをFrobenius標準4周期同型と呼ぶ。
特に標準14641型は4周期的である。

\begin{theorem}[標準4周期同型の選択独立性]\label{thm:canonical-four}
任意の二つのtrace族から得る$\tau^\lambda,\tau^{\lambda'}$は、添字保持自己同型$h$により
$\tau^{\lambda'}=h(4)\tau^\lambda h^{-1}$と書ける。
裸の$C,C'$が同型なら、それぞれのFrobenius標準4周期同型付きの組も同型である。
従って幾何学的点の同型類では、標準4周期構造を選ぶことは余分な分類情報を持たない。
\end{theorem}
\begin{proof}
$\lambda'_i=t_i\lambda_i$と置くと
\begin{equation}\label{eq:trace-change}
 \nu^{\lambda'}_{ij}=(t_i/t_j)\nu^\lambda_{ij}.
\end{equation}
四つの整数剰余類ごとに初期値を選び、$s_{i+4}/s_i=t_i$を正負両方向へ解く。
頂点スカラー自己同型$h^E_{ij}(a)=(s_i/s_j)a$に対して
\[
 h^E(4)\nu^\lambda(h^E)^{-1}(a)
 =\frac{s_{i+4}s_j}{s_{j+4}s_i}\nu^\lambda(a)
 =\frac{t_i}{t_j}\nu^\lambda(a)=\nu^{\lambda'}(a).
\]
逆双対で$C$上の共役性を得る。
次に$f:C\simeq C'$があればtraceを双対同型で移送する。
式\eqref{eq:nakayama-pair}の特徴付けにより、移送したtraceの周期構造は$f$と両立する。
別のtraceが指定されても上のスカラー共役で補正できる。
これは同型類の集合上の全単射であり、群oidの完全忠実性を主張するものではない。
\end{proof}

\begin{proposition}[有限四窓での標準化]\label{prop:canonical-matrix}
厳密4周期表示で局所Frobeniusな双対$E$を考える。
四つのtraceを厳密4移動で延長し、$E_{i,i+1}$の基底$u_a$、
$E_{i+1,i+4}$の基底$v_b$と、移動した基底$u'_c$に対して
\[
 P_i=(\lambda_i(u_a v_b))_{ab},\quad
 Q_i=(\lambda_{i+1}(v_bu'_c))_{bc},\quad N_i=Q_i^{-1}P_i^{\mathsf T}
\]
と置く。$N_i$は局所Nakayama同型の行列である。
選んだ厳密4移動がFrobenius標準であることは、traceを選んで
\begin{equation}\label{eq:canonical-matrix-test}
 P_i^{\mathsf T}=Q_i\qquad(i\bmod4)
\end{equation}
を満たせることと同値である。
\end{proposition}
\begin{proof}
$N_i$の列を$\nu(u_a)$の係数とすると、特徴付けの式は
$P_{i,ab}=\sum_cQ_{i,bc}(N_i)_{ca}$、すなわち$P_i^{\mathsf T}=Q_iN_i$である。
対合の完全性から$Q_i$は可逆である。$E$は次数一生成なので、
次数一での$N_i=I$は全積での厳密4移動との一致を意味する。
\end{proof}

\begin{proposition}[有限次数のFrobenius検査]\label{prop:finite-frobenius}
任意の四窓の二次双対$E$について、各始点で次数$0,\ldots,5$の次元が
$(1,4,6,4,1,0)$であり、相補次数の積対合が完全なら、全次数の局所Frobenius性が成立する。
これらは有限回の線型代数で検査できる。ただしこれだけから元代数のKoszul性は従わない。
\end{proposition}
\begin{proof}
次数一生成性により長さ5の全語が零なら全高次も零である。
相補次数の対合の完全性は積行列の行列式の非消滅に等しい。
必要な商空間と積行列は長さ5以下のテンソル空間と二次関係から構成でき、始点も4通りで足りる。
Koszul複体の正確性は、この有限検査とは別の条件である。
\end{proof}

\begin{corollary}\label{cor:canonical-slice}
裸の標準14641型$\ZZ$-代数の幾何学的同型類は、
$\mathfrak P$のうちKoszul AS正則で式\eqref{eq:canonical-matrix-test}を満たすtraceを持つ点の
$G$-軌道と全単射になる。
\end{corollary}
\begin{proof}
任意の$C$には標準4周期構造があり、定理\ref{thm:four-window}で厳密表示に直せる。
命題\ref{prop:canonical-matrix}が標準化条件を与え、
定理\ref{thm:canonical-four}が裸の同型から標準構造を保つ同型への補正を与える。
最後に四窓群oidの同型判定を使う。
\end{proof}

\begin{theorem}[裸の周期性の有限判定]\label{thm:finite-period-test}
Frobenius標準四窓$(R_0,R_1,R_2,R_3)$に対して$C\simeq C(p)$であることは、
四つの可逆写像$g_i:V_i\to V_{i+p}$が
\begin{equation}\label{eq:finite-period-test}
 (g_i\otimes g_{i+1})R_i=R_{i+p}\quad(i\bmod4)
\end{equation}
を満たすことと同値である。
\end{theorem}
\begin{proof}
式があれば$g_{i+4}=g_i$と延長し、命題\ref{prop:period-relations}を適用する。
逆に$f:C\simeq C(p)$と標準$\tau$を取る。
$\tau(p)$と$f(4)\tau f^{-1}$は$C(p)$の二つの標準4周期構造である。
定理\ref{thm:canonical-four}によって頂点スカラー自己同型$h$を選び
$\tau(p)=h(4)[f(4)\tau f^{-1}]h^{-1}$とできる。
$f'=hf$は$\tau(p)f'=f'(4)\tau$を満たすので、その次数一行列は4周期的である。
二次関係の保存が所要の式を与える。
\end{proof}

\begin{proposition}[多項式系]\label{prop:period-polynomial-system}
$R_i$の列基底を階数6の$16\times6$行列$M_i$で表す。
$p=1,2$の周期の存在は
\begin{equation}\label{eq:period-polynomial-system}
 (G_i\otimes G_{i+1})M_i=M_{i+p}H_i,\qquad u_i\det G_i=1
\end{equation}
の可解性に等しい。未知数は$G_i$の64成分、$H_i$の144成分、$u_i$の4個、合計212個、
方程式は384成分等式と4可逆性等式、合計388個である。
\end{proposition}
\begin{proof}
$G_i$が可逆で$M_i$の階数が6なら、等式を満たす$H_i$も自動的に可逆である。
従ってこれは関係空間の一致を行列で書いたものであり、前定理を適用できる。
\end{proof}

\begin{corollary}\label{cor:finite-isom-test}
既にAS正則性が保証された標準型の$C,C'$のFrobenius標準関係行列を$M_i,M'_i$とする。
添字保持同型の存在は
$(G_i\otimes G_{i+1})M_i=M'_iH_i$, $u_i\det G_i=1$の可解性に等しい。
\end{corollary}
\begin{proof}
系\ref{cor:canonical-slice}の同一$G$-軌道という条件を行列表示する。
$H_i$の可逆性は前命題と同様に自動である。
\end{proof}

\begin{corollary}\label{cor:period-orbits}
四窓の添字を一つ回転する操作を$r$とする。
最小周期1は$r(R)\in G\cdot R$、最小周期2は$r^2(R)\in G\cdot R$かつ$r(R)\notin G\cdot R$、
最小周期4は$r^2(R)\notin G\cdot R$で判定できる。
\end{corollary}
\begin{proof}
定理\ref{thm:finite-period-test}と補題\ref{lem:period-subgroup}による。
\end{proof}
\begin{remark}
$G$の頂点ごとのスカラー$(k^\times)^4$は全ての関係空間を固定し、有効作用は$\PGL_4^4$を通る。
しかしstackの自己同型としては残る。$240-60=180$をAS正則モジュライの次元と呼ぶのも誤りである。
有限判定式の解を一つ検証することと、解の非存在を示すことも異なる。
\path{computations/verify_period_witness.py}は正の証明書だけを検査する。
\end{remark}

\section{1周期構造、graded表示、既存の2周期族}
\begin{proposition}\label{prop:graded-period}
二次正連結$C$が1周期的であることと、ある次数一生成二次graded algebra $A$について
$C\simeq\check A$となることは同値である。1周期同型を選んだ組の分類はgraded同型分類に等しい。
\end{proposition}
\begin{proof}
$\check A$には厳密1移動がある。
逆に$\sigma:C\to C(1)$を選び、$A_n=C_{0n}$、
$a*b=a\,\sigma^m_{0n}(b)$（$a\in A_m,b\in A_n$）と定義する。
$\sigma$が積を保つので結合律が成立する。次数一生成・二次性も引き継がれる。
$\sigma$による各成分の同定から$\check A\simeq C$を得る。同型もこの対応を保つ。
\end{proof}

\begin{proposition}\label{prop:compatible-periodic-windows}
$(C,\tau_4)$に対し$p=1,2$とする。
$p$周期同型$\sigma_p$の$(4/p)$乗が$\tau_4$であるなら、四窓の基底を選んで$R_i=R_{i+p}$とできる。
逆にこの等式は、指定した$\tau_4$を$(4/p)$乗として与える$p$周期同型を定める。
\end{proposition}
\begin{proof}
$p$個の剰余類代表の基底を$\sigma_p$で全整数へ運ぶ。
その$(4/p)$乗が$\tau_4$なので、この基底は指定された4周期座標とも両立し、関係が$p$周期になる。
逆に関係が$p$周期なら同一座標の移動が代数同型へ降り、その$(4/p)$乗は厳密4移動である。
\end{proof}
\begin{remark}
これは整合する周期構造を選んだ結果である。
裸の1周期性から、任意の4周期構造が1周期構造の4乗になるとは従わない。
Nakayama自己同型やgraded CY性も、裸の$\ZZ$-代数不変量として無条件に採用しない。
\end{remark}

既存のPBW族\cite{as4period}は、$\varepsilon_i=(-1)^i$として
\begin{equation}\label{eq:period-two-family}
\begin{aligned}
 XY+YX&=Z^2+W^2,& XZ+ZX&=\varepsilon_i\lambda WY,\\
 XW+WX&=Z^2,&YZ+ZY&=W^2,\\
 YW+WY&=0,&ZW+WZ&=0
\end{aligned}
\end{equation}
で与えられる。同原稿では全閉値でのKoszul性・左右AS正則性、14641型、
$\lambda=0$での1周期性と、非空開集合での最小周期2が証明されている。
最小周期4の別の例もある。本稿はその全コードを再実行したとの主張はしない。

\begin{proposition}\label{prop:period-two-sign}
表示上$A(\lambda)(1)=A(-\lambda)$である。
$A(\lambda)$が最小周期2なら、$A(\lambda)$と$A(-\lambda)$は添字保持同型では非同型である。
特に既存の非1周期性の証明を使えば$\lambda=1$に適用できる。
\end{proposition}
\begin{proof}
唯一の添字依存係数$(-1)^i\lambda$は、一つ移動すると$(-1)^i(-\lambda)$になる。
もし両者が添字保持同型なら$A(\lambda)\simeq A(\lambda)(1)$となり、最小周期2に矛盾する。
\end{proof}
\begin{remark}
添字移動をさらに同一視する別のモジュライと混同しない。
また、非1周期的な点を含む族があることだけから、その族の成分次元や1周期的部分の余次元は決まらない。
\end{remark}


\section{最上次数テンソルによるFrobenius候補の復元}
\label{sec:frobenius-tensor}
この節ではKoszul性を仮定する前の有限Frobenius候補を整理する。
最上次数のテンソルを用い、正則性の未決定と、全対象を覆う不変な分岐を分ける。

\begin{proposition}[Frobenius検査の四次への短縮]\label{prop:frobenius-degree-four}
四窓の二次双対$E$について、各始点の次数$0,\ldots,4$の次元が
$(1,4,6,4,1)$であり、相補次数の積対合が完全なら、$E_{i,i+5}=0$は自動的に成立する。
従って命題\ref{prop:finite-frobenius}の五次元数検査はこの仮定の下で省略できる。
これは元の代数のKoszul性の判定を短縮する主張ではない。
\end{proposition}
\begin{proof}
前節の交差空間$W_{i,p}=E_{i,i+p}^*$を用いる。
非零traceを引き戻したテンソル$w_i\in T_{i,4}$は$W_{i,4}$を生成する。
$E_{i,i+1}\otimes E_{i+1,i+4}$の完全対合から、
$w_i$の最初の一因子と残りの三因子の間の行列階数は4である。
一方、交差空間の定義により
\[
 W_{i,5}=(W_{i,4}\otimes V_{i+4})\cap(V_i\otimes W_{i+1,4}).
\]
非零元があれば、ある非零$u\in V_i$, $v\in V_{i+4}$により
$w_i\otimes v=u\otimes w_{i+1}$と書ける。
$v$を1へ送る線型汎関数で末尾を縮約すると$w_i=u\otimes z$となり、
最初の一因子での行列階数が1となって矛盾する。
よって$W_{i,5}=0$であり、次数一生成性から全高次も零となる。
\end{proof}

$0\ne w\in V_0\otimes V_1\otimes V_2\otimes V_3$を取り、
その巡回回転を$w_i\in V_i\otimes\cdots\otimes V_{i+3}$と書く。
回転には符号を入れず、添字は法4で読む。
テンソルを二組の因子に分けた線型写像の階数を、その分け方での行列階数と呼ぶ。
$w_i$の後ろ二因子を縮約して得る像を
\[
 R_i(w)=\operatorname{im}\bigl((V_{i+2}\otimes V_{i+3})^*
                  \xrightarrow{\ w_i\ } V_i\otimes V_{i+1}\bigr)
\]
と定め、これらから交差空間$W_{i,p}(w)$を作る。

\begin{theorem}[標準Frobenius四窓のテンソル復元]\label{thm:frobenius-tensor-reconstruction}
射影テンソル$[w]$に次の有限条件を課す。
\begin{enumerate}
\item 各一因子と残り三因子の間の行列階数は4である。
\item $01\mid23$と$12\mid30$の二つの行列階数は6である。
\item 全ての始点で$\dim W_{i,3}(w)=4$, $\dim W_{i,4}(w)=1$である。
\end{enumerate}
これらの$[w]$と、有限Frobenius検査を通り厳密4移動がFrobenius標準である四窓とは、
$[w]\mapsto(R_i(w))$によって$G$同変な全単射となる。
従って、さらに元代数のKoszul AS正則性を課した部分では、
$G$軌道が裸の標準14641型の添字保持同型類を与える。
上記三条件だけでKoszul性を結論することはしない。
\end{theorem}
\begin{proof}
まず三条件を満たす$w$から出発する。
二つの二対二階数条件により各$R_i(w)$は6次元である。
テンソル$w_i$は、隣接する三つの位置の関係条件全てを満たすので$W_{i,4}(w)$に入り、
第三条件によりその生成元になる。
$w_i$の最初または最後の一因子を縮約した像は、対応する三次交差空間に含まれる。
第一条件と第三条件により、これは4次元の三次交差空間全体に一致する。
$E$をこの関係四窓の二次双対とし、$w_i$を引き戻しとして持つtrace
$\lambda_i:E_{i,i+4}\to k$を選ぶ。
一対三の積対合は第一条件から、二対二の積対合は第二条件から完全であり、
零対四と四対零も完全である。
前命題から五次以上は零なので、有限Frobenius条件が全て成立する。
回転の定義そのものにより
\[
 \lambda_i(\alpha\beta\gamma\delta)
   =\lambda_{i+1}(\beta\gamma\delta\alpha)
\]
が成立する。ここで各文字は対応する次数一双対空間の元である。
次数三の積が$E_{i+1,i+4}$を生成するから、この等式は局所Nakayama写像が
厳密4移動であることを意味する。

逆に標準Frobenius四窓から出発し、標準性を与えるtraceを選ぶ。
その引き戻し$w_i$は$W_{i,4}$を生成し、Nakayamaの特徴付けから互いに巡回回転である。
一対三の完全対合により第一条件が成立する。
二対二の完全対合から、後ろ二因子の縮約像は$E_{i,i+2}^*=R_i$全体になり、
第二条件と関係空間の復元を得る。第三条件は双対の次数元数である。
$W_{0,4}$が一次元なので$[w]$は一意である。
全構成は基底変更と可換し、最後の同型類の主張は系\ref{cor:canonical-slice}による。
\end{proof}

\begin{proposition}[第三の二対二階数による全被覆分岐]\label{prop:opposite-rank-split}
標準14641型の対象について、上の最上次数テンソルから
\[
 \rho(C)=\rk\bigl((V_0\otimes V_2)^*\xrightarrow{\ w\ }V_1\otimes V_3\bigr)
\]
を定めると、これはtrace・基底・標準四窓の選択によらない添字保持同型不変量であり、
$2\leq\rho(C)\leq16$を満たす。
従って$\rho\leq6$と$\rho\geq7$は、全対象を互いに素な二つの部分へ分ける。
前者はテンソル候補空間内で相対的に閉、後者は相対的に開である。
各部分の正規形・AS正則locus・完全同型判定はなお未決定である。
\end{proposition}
\begin{proof}
traceの変更は$w$の非零スカラー倍、基底変更は四因子での可逆変換を与えるため階数を保つ。
裸の同型を標準構造を保つ同型へ補正できることは定理\ref{thm:canonical-four}で示した。
よって標準四窓の選択にもよらない。
非零$16\times16$行列なので$1\leq\rho\leq16$である。
もし$\rho=1$なら$w=a_{02}\otimes b_{13}$と書ける。
四つの一因子階数が4だから、$a_{02},b_{13}$を$4\times4$行列と見た階数はともに4。
すると$01\mid23$での階数はそのKronecker積の階数16となり、6という条件に矛盾する。
二分岐の被覆と非交差性は整数階数の二者択一であり、
閉性は7次小行列式の消滅による。
\end{proof}
\begin{remark}
これは正則な全対象に対する被覆分岐であって、各子の分類が完成したという主張ではない。
同じ$\rho$でも同型とは限らず、異なる$\rho$なら非同型である。
定理のテンソル候補族の中の特殊化では階数は増加しない。
\path{computations/verify_frobenius_tensor.py}は復元と既存族の階数を有理数上で検査する。
テンソル表示の一般原理そのものの新規性は主張しない。
\end{remark}

\subsection{量子多項式上の全階数層と最小階数層の同型}
\begin{proposition}[量子族の第三階数]\label{prop:quantum-opposite-rank}
量子多項式代数
\[
 A(q)=k\langle x_0,x_1,x_2,x_3\rangle/
 (x_bx_a-q_{ab}x_ax_b:\ 0\leq a<b\leq3),\qquad q_{ab}\in k^\times
\]
に対して
\[
 F_1=q_{02}q_{13}-q_{03}q_{12},\quad
 F_2=q_{01}q_{12}q_{23}-q_{03},\quad
 F_3=q_{01}q_{13}-q_{02}q_{23}
\]
と置く。このとき
\[
 \rho(\check A(q))=6+2\#\{j\in\{1,2,3\}:F_j\ne0\}.
\]
従って量子族で現れる階数は正確に$6,8,10,12$であり、一般点では12である。
$\rho=6$の量子対象は、添字保持同型でちょうど二種類、
可換多項式代数と全ての$q_{ab}=-1$である反可換量子多項式代数に帰着する。
これは量子族に限定した主張であり、全対象の$\rho=6$層の分類ではない。
\end{proposition}
\begin{proof}
量子書換えは全重なりを解消し、順序付き単項式が基底になる。
二次双対は量子外積代数で完全対合を持つため、全ての$q_{ab}\ne0$で
標準型Koszul AS正則である。
四次交差空間の生成元は、$\pi\in S_4$に対し
\[
 w=\sum_{\pi\in S_4}c_\pi\,
 x_{\pi(0)}\otimes x_{\pi(1)}\otimes x_{\pi(2)}\otimes x_{\pi(3)},\qquad
 c_\pi=\prod_{a<b}
 \begin{cases}q_{ab},&\pi^{-1}(a)<\pi^{-1}(b),\\-1,&\pi^{-1}(a)>\pi^{-1}(b)\end{cases}
\]
で与えられる。隣接する$a,b$を交換した二係数の比が$-q_{ab}$になるので、
各隣接位置の関係条件を満たすことが直接分かる。
このテンソルは標準化前の最初の四一次空間だけで定まる。
それらの基底を保ったまま標準4周期を延長しても$w$は変わらないので、$\rho$の計算に使える。

$02\mid13$行列の非零行は異なる二文字を持つ12行である。
行の二文字集合$S\subset\{0,1,2,3\}$を固定すると、列の二文字集合はその補集合に限られる。
従って六つの$2\times2$ブロックに分解する。
全ブロックの各成分は非零単項式であり、行集合$S$を順に
$01,02,03,12,13,23$と取った行列式は、可逆単項式を除いてそれぞれ
$F_1,F_2,F_3,F_3,F_2,F_1$である。
よって各ブロックの階数は行列式零で1、非零で2となり公式を得る。
これらの恒等式は\path{computations/verify_quantum_tensor_rank.py}でも記号的に検査する。
例えば$q_{01}=q_{12}=q_{23}=1$と固定し、
$(q_{02},q_{13},q_{03})=(1,1,1),(2,2,1),(2,3,1),(2,3,5)$と取れば、
それぞれ階数$6,8,10,12$を得る。非零条件は開であり一般点の主張も従う。

階数6なら$F_1=F_2=F_3=0$なので
\[
 q_{03}=q_{01}q_{12}q_{23},\qquad
 q_{13}=q_{02}q_{23}/q_{01},\qquad
 q_{02}^2=q_{01}^2q_{12}^2.
\]
標数零より、ある$s\in\{1,-1\}$に対して
$q_{02}=s q_{01}q_{12}$, $q_{13}=s q_{12}q_{23}$である。
$d_0=1$, $d_1=q_{01}/s$, $d_2=q_{01}q_{12}$,
$d_3=q_{01}q_{12}q_{23}/s$と置くと
$q'_{ab}=q_{ab}d_a/d_b=s$が全ての対で成立する。
位置$i$の基底変換$g_i(x_a)=d_a^i x_a$は、
$q'$の関係を$q$の関係へ送る。従ってこれは全整数の添字を保つ同型であり、
graded同型を主張するものではない。

最後に二つの標準形が添字保持同型でないことを示す。
最初の二次関係の左点対応の射影支持は、可換型では全$\mathbb P^3$である。
反可換型では、$x_ay_b+x_by_a=0$を満たす非零$y$が存在するためには、
$x$の非零座標は高々二つでなければならない。
実際、$x_a\ne0$なら$y_a=0$は全$y=0$を強制し、三つの非零座標があれば
三対の式が$y_b/x_b=-y_a/x_a$を互いに矛盾させる。
逆に非零座標が高々二つならそのような$y$を選べる。
従って支持は六本の座標直線の和である。
添字保持同型は二一次空間の射影変換によりこの支持を保つので、両型は非同型である。
\end{proof}

% !TeX root = ../papers/nc-p3.tex
\section{量子部分の全特殊値を含むindexed同型判定}
\label{sec:quantum-indexed-all}

$q_{ba}=q_{ab}^{-1}$, $q_{aa}=1$と拡張し、異なる三添字に対して
$T_{abc}(q)=q_{ab}q_{bc}q_{ca}$と置く。
\begin{theorem}[量子多項式の全パラメータindexed同型]\label{thm:quantum-indexed-all}
全てのパラメータが非零である量子多項式代数$A(q),A(q')$について、次は同値である。
\begin{enumerate}
\item $\check A(q)$と$\check A(q')$が添字保持同型である。
\item ある置換$\sigma\in S_4$と$d_0,\ldots,d_3\in k^\times$が存在し、
全ての異なる$a,b$について
$q'_{\sigma(a)\sigma(b)}=q_{ab}d_a/d_b$となる。
\item ある$\sigma\in S_4$について、全ての三添字で
$T_{abc}(q)=T_{\sigma(a)\sigma(b)\sigma(c)}(q')$となる。
\end{enumerate}
この判定は根が1となる特殊値やパラメータの重複も含み、一般性条件を要しない。
これは量子多項式部分を尽くす同型条件であり、全AS四窓の分類ではない。
\end{theorem}
\begin{proof}
(2)なら位置$i$で$g_i(x_a)=d_a^{-i}x'_{\sigma(a)}$と置く。
二次関係の両項の係数比を取れば、全整数$i$で関係空間が対応し、(1)を得る。
(2)から(3)は三つの比の相殺による。
逆に(3)があれば、置換後の比$r_{ab}=q'_{\sigma(a)\sigma(b)}/q_{ab}$に対し
$r_{ab}r_{bc}=r_{ac}$である。$d_0=1$, $d_b=r_{0b}^{-1}$と取れば$r_{ab}=d_a/d_b$となる。

残る(1)から(2)を示す。
最初の二次関係から得る点対応は
\[
 x_by_a=q_{ab}x_ay_b\qquad(a<b)
\]
である。$x$の非零座標集合を$I$とすると、非零$y$が存在するための必要十分条件は
$I$内の全ての三添字について$T_{abc}=1$となることである。
実際、$y$の非零座標も$I$と一致し、$y_b/x_b=q_{ab}^{-1}(y_a/x_a)$の整合性がこの条件になる。
その場合の$y$は射影的に一意であり、特に各座標点$p_a$の対応先は$p_a$自身である。
従って左射影支持は、全三添字条件を満たす$I$の座標線型空間$\mathbb P(I)$の和となる。

四つの三添字について
$T_{012}T_{023}=T_{013}T_{123}$なので、$T=1$となる三角形の数は$0,1,2,4$のいずれかである。
左支持は順に、六本の座標直線、一平面と外の一点から出る三本の直線、
二平面とそれぞれの外の頂点を結ぶ一本の直線、全$\mathbb P^3$となる。
添字保持同型の最初の二つの行列を$g_0,g_1$とすると、これらは点対応の両因子の射影変換を与える。

三角形数が0の場合、六直線の三重点として四座標点が内在的に復元される。
数が1の場合も、唯一の平面上の三つの付着点と、外の三直線の共通交点として復元される。
従ってこれらの場合$g_0$は単項行列である。
各座標点の右対応が一意なので$g_1$も同じ置換を持つ単項行列となる。
対角係数を$a_j,b_j$と書くと、関係保存から
$q'_{\sigma(a)\sigma(b)}=q_{ab}(a_a/b_a)/(a_b/b_b)$を得る。

三角形数が2の場合は、添字を付け直して共通の辺を$01$、外の二頂点を$2,3$とする。
二平面の交線と、残る直線の二つの付着点が内在的であるから、
$g_0$は$\langle x_0,x_1\rangle$上の$\GL_2$ブロックと外の二頂点のスカラーからなる。
まず(2)から(1)という既に示した基底変換を使い、双方のパラメータで$q_{01}=1$に正規化してよい。
$T_{012}=T_{013}=1$から$q_{02}=q_{12}$, $q_{03}=q_{13}$となる。
従って$\langle x_0,x_1\rangle$内の任意の$\GL_2$変換は量子代数のgraded自己同型である。
これを相手側で合成して$g_0$の$\GL_2$ブロックを除けば、前と同じ単項行列の場合に帰着する。
この正規化と自己同型の除去は、結論(2)の成立を変えない。

三角形数が4なら、双方の全ての$T$が1であるから、(3)から(2)という既証部分を直接適用できる。
以上で全特殊値を尽くした。
\end{proof}
\begin{remark}
三角形数0の一般量子対象で得られる従来の3パラメータの忘却を、
上の証明は平面成分とその交線に$\GL_2$自由度が現れる特殊値まで拡張する。
$\rho$の一致はこの三添字判定より弱い。
同じ$\rho=8$でも、$q_{01}=q_{12}=q_{23}=1$で
$(q_{02},q_{13},q_{03})=(2,2,1)$と$(3,3,1)$の二対象は非同型である。
\path{computations/verify_quantum_indexed.py}は24置換の判定と対角証明書を検査する。
\end{remark}
\begin{corollary}[量子部分の三変数正規形と有限残余作用]\label{cor:quantum-torus-normal-form}
量子多項式の全indexed同型類は、
$q_{01}=q_{02}=q_{03}=1$, $(q_{12},q_{13},q_{23})=(p,q,r)\in(k^\times)^3$
という三変数表示で被覆される。
二つの表示が同型であることは、次の三つの変換が生成する$S_4$作用で同じ軌道に属することと同値である。
\[
 s_{01}(p,q,r)=(p^{-1},q^{-1},pr/q),\qquad
 s_{12}(p,q,r)=(p^{-1},r,q),\qquad
 s_{23}(p,q,r)=(q,p,r^{-1}).
\]
従って量子部分の残余同型は高々24個の明示的候補で判定でき、特殊値もこの式で扱える。
これは粗い同型類の記述であり、自己同型群の次元が特殊値でも一定という主張ではない。
\end{corollary}
\begin{proof}
$d_0=1$, $d_b=q_{0b}$という前定理の変換で最初の三係数を1へ直せる。
残る三係数は$(T_{012},T_{013},T_{023})$であり、第四の三添字値は$T_{123}=pr/q$である。
置換$(01),(12),(23)$の作用をこの三添字値に適用すれば上の三式になる。
前定理が残余同型の必要十分性を与える。
\end{proof}

\section{可換三生成部分上のOreチャートと第三階数}
\label{sec:polynomial-ore}
低階数残余を調べるため、三生成部分が可換な場合を一括して扱う。
以下は既知のOre構成を用いる共通の認証チャートであり、全対象の被覆ではない。

$B=k[x_0,x_1,x_2]$とし、任意の二次式
\[
 f_i=\sum_{0\leq a\leq b\leq2}c_{i,ab}x_ax_b\quad(i=0,1,2)
\]
を取る。$\delta(x_i)=f_i$で定める導分に対し
\[
 A_f=B[X;\delta],\qquad Xx_i-x_iX=f_i
\]
と置く。係数空間は$\mathbb A^{18}$である。

\begin{theorem}[全ファイバー正則な多項式Oreチャート]\label{thm:polynomial-ore-regular}
全ての係数で$A_f$は標準14641型quadratic Koszul AS正則であり、
$\check A_f$は1周期的である。基底は$x_0^{a_0}x_1^{a_1}x_2^{a_2}X^m$である。
これは18次元の同型類像や全モジュライ成分を意味しない。
\end{theorem}
\begin{proof}
任意の三つの$f_i$は多項式環の導分を定める。
$x_0<x_1<x_2<X$の書換えで$Xx_i$を$x_iX+f_i$に直すと、
$X$の個数を第一重み、反転数を第二順序として停止する。
三次重なりは$Xx_jx_i$の三つと$x_2x_1x_0$の一つで、
$\delta(x_jx_i-x_ix_j)=0$により全て解消する。
従って二次Gr\"obner基底であり、主張の基底とKoszul性を得る。

双対生成元を$e_0,e_1,e_2,z$とすると、二次双対は
\begin{align*}
 z^2&=0,& ze_i+e_iz&=0,\\
 e_i^2&=\sum_a c_{a,ii}e_az,&
 e_ie_j+e_je_i&=\sum_a c_{a,ij}e_az\quad(i<j)
\end{align*}
である。$e_a z$は中心的であり、二次補正をもう一度使うと$z^2$が現れる。
$e_i$に重み1、$z$に重み0を与えた書換えは、平方自由の昇順語へ停止する。
20個の三次重なりは恒等的に解消し、随伴次数代数は四生成外積代数となる。
この20恒等式は\path{computations/verify_polynomial_ore.py}で18変数のまま検査する。
従って双対は次数元数14641を持ち、最上積$e_0e_1e_2z$の係数をtraceとする相補積対合は
外積代数の対合のフィルター付き持上げとして完全である。
Koszul性とこの有限双対Frobenius性から、長さ4の線型解消と左右AS条件を得る。
\end{proof}

$w_f$をこの最上traceの引き戻しとし、$z$を第3生成元と読む。
最初の四一次空間の基底を保って標準4周期表示へ直せるので、
この$w_f$の第三階数は内在的な$\rho$に一致する。
少なくとも二つの$z$を含む四次語のtraceは零であり、ちょうど一つの場合は外積の符号である。
$z$を含まない四次語では二次補正が一度だけ寄与するので、そのtraceは$c$に線型である。
$\mathcal P=(00,01,02,11,12,22)$の順に、$U=\langle x_0,x_1,x_2\rangle$の
対称二次テンソルを並べる。各順序対$ac$を$e_a\otimes e_c$と同一視し、
非対角対には両順序を一度ずつ足す規約で
\[
 Q_{ac,bd}(f)=\sum_{(a',c')\in\operatorname{ord}(ac)}
              \sum_{(b',d')\in\operatorname{ord}(bd)}
              (w_f)_{a'b'c'd'}
\]
を定める。これは$6\times6$行列である。

\begin{theorem}[第三階数の交代行列判定]\label{thm:polynomial-ore-rank}
$Q(f)$は交代行列で、$f\mapsto Q(f)$は18次元係数空間から交代$6\times6$行列空間への全射線型写像である。
その核はちょうど$f_i=x_i\ell$（$\ell$は任意の一次式）となる3次元部分空間であり、
全パラメータで
\[
 \rho(\check A_f)=6+\rk Q(f)
\]
が成立する。従ってこのチャートの第三階数は正確に$6,8,10,12$である。
$\rho\leq6$の全ての対象は可換多項式代数のindexed類に属し、
このチャートは\texttt{LOW-UNCLASSIFIED}に対象を与えない。
\end{theorem}
\begin{proof}
第三の二対二行列で、両因子が$U$に属する9成分と、一方が第3生成元である6成分を分けると、
零行・零列を除いて
\[
 M_f=\begin{pmatrix}D(f)&B_0\\ C_0&0\end{pmatrix}
\]
となる。$B_0,C_0$は外積の符号だけで決まり、双方の階数は3である。
$\operatorname{im}B_0=\Lambda^2U$、$\ker C_0=\Sym^2U$なので、
$D(f)$の$\ker C_0\to\operatorname{coker}B_0$への誘導行列が、定義した$Q(f)$である。
従ってブロック消去により$\rk M_f=3+3+\rk Q(f)$となる。

$\mathcal P$を$0,\ldots,5$と番号付けすると、上三角成分は次の通りである。
\[
\begin{array}{c|l@{\qquad}c|l}
 ij&Q_{ij}&ij&Q_{ij}\\\hline
 01&2c_{2,00}&12&2c_{0,00}-c_{1,01}-c_{2,02}\\
 02&-2c_{1,00}&13&2c_{2,11}\\
 03&c_{2,01}&14&c_{0,01}-2c_{1,11}+c_{2,12}\\
 04&-c_{1,01}+c_{2,02}&15&c_{0,02}-c_{1,12}\\
 05&-c_{1,02}&23&-c_{0,01}+c_{2,12}\\
 24&-c_{0,02}-c_{1,12}+2c_{2,22}&25&-2c_{1,22}\\
 34&2c_{0,11}&35&c_{0,12}\\
 45&2c_{0,22}&&
\end{array}
\]
対角成分は零、下三角成分は上三角の負である。
双対の書換えからこの表を直接計算でき、再現コードは四次の全256語で同じ表を検査する。
$Q=0$は、まず$i\notin\{a,b\}$の九係数が零であることを強制し、
残る九係数には
\[
 c_{0,00}=c_{1,01}=c_{2,02},\qquad
 c_{0,01}=c_{1,11}=c_{2,12},\qquad
 c_{0,02}=c_{1,12}=c_{2,22}
\]
だけを課す。従って核は$f_i=x_i\ell$で、次元は3である。
線型写像の階数は15となり、交代行列空間全体への全射性が従う。
交代行列の階数は$0,2,4,6$なので、それぞれの非空性も含めて階数の主張を得る。

最後に$f_i=x_i\ell$なら、位置$j$の一次写像
$x_i\mapsto x_i$, $X\mapsto X-j\ell$を取る。
隣接位置の積に適用すると
\[
 (X-j\ell)x_i-x_i(X-(j+1)\ell)-\ell x_i=0
\]
が可換多項式環内で成立するため、これは全整数で定義されたindexed同型である。
逆に$\rho\leq6$なら階数式から$Q=0$となるので、この場合で尽きる。
\end{proof}
\begin{remark}
$Q$の階数で与える条件はこの係数空間のPfaffian階数条件であり、
indexed同型の完全不変量でも、全モジュライの成分方程式でもない。
Ore構成・PBW・有限双対の手法自体について新規性を主張しない。
\end{remark}

\section{反可換三生成部分の位置依存Ore正規化}
\label{sec:anticommutative-ore}
$B_-=k\langle x_0,x_1,x_2\rangle/(x_bx_a+x_ax_b:a<b)$とする。
四次元indexed表示の残る関係が
\[
 Xx_a+x_aX=f_{n,a}\in(B_-)_2\qquad(a=0,1,2)
\]
である場合を考える。各位置の三二次式は最初は任意で、PBW性を仮定しない。

\begin{theorem}[反可換Ore表示の三次判定と正規形]\label{thm:indexed-anticommutative-ore}
この表示のindexed代数について、標準14641型Koszul AS正則性と
全始点の$\dim C_{n,n+3}=20$は同値である。
この条件が成立することは、一次写像$x_a\mapsto x_a$, $X\mapsto X+u_n$により
\begin{equation}\label{eq:anti-ore-normal}
 Xx_a+x_aX=\sum_b H_{ab}x_b^2+(-1)^n\gamma_a x_bx_c,
 \qquad \{b,c\}=\{0,1,2\}\setminus\{a\},\quad b<c
\end{equation}
へ直せることと同値である。$H$の9成分と$\gamma$の3成分は任意で、全ファイバーが正則である。
元の表示が$p$周期なら同型行列も$p$周期に選べ、$p$が奇数なら$\gamma=0$である。
従ってこの表示を持つ正則対象の最小周期は1または2に限られる。
\end{theorem}
\begin{proof}
反転した文字対を直すと三次の20個の順序付き語が生成系になる。
重なり$Xx_bx_a$の差は$B_-$内の
\begin{equation}\label{eq:anti-ore-overlap}
 f_{n,b}x_a+f_{n,a}x_b-x_bf_{n+1,a}-x_af_{n+1,b}
\end{equation}
である。標準三次元数20なら、この差は$B_-$の順序付き三次基底において零でなければならない。
その係数を比較すると、各$f_{n,a}$は一意に
\[
 f_{n,a}=\sum_b h_{n,ab}x_b^2+
          \sum_{b\ne a}v_{n,b}x_bx_a+
          \gamma_{n,a}x_bx_c\quad(b<c,\ b,c\ne a)
\]
と書け、条件は正確に
\begin{align*}
 h_{n+1,ab}&=h_{n,ab}\quad(a\ne b),&
 \gamma_{n+1,a}&=-\gamma_{n,a},\\
 h_{n+1,aa}-h_{n,aa}&=-v_{n,a}-v_{n+1,a}
\end{align*}
となる。ここで$x_bx_a$は表示通りの順序であり、順序付き基底へ直すときは反可換符号を入れる。
三対の式\eqref{eq:anti-ore-overlap}を全て使うことが重要で、
一対の比較だけでは二成分に現れる同じ$v_{n,b}$の一致をまだ強制しない。
この係数比較は、30行36列の隣接窓行列の階数18と、上記パラメータ化の同じ階数でも検算できる。

$u_n-u_{n+1}=v_n$を解くと、非対角の$H_{ab}=h_{n,ab}$は一定であり、
$H_{aa}=h_{n,aa}-u_{n,a}-u_{n+1,a}$も最後の条件から一定となる。
また$\gamma_n=(-1)^n\gamma_0$である。
$X\mapsto X+u_n$を代入すれば、正確に\eqref{eq:anti-ore-normal}を得る。
共通の$u_n$への定数加算で$H_{aa}=0$とすることもできるが、以下では対称性のため9成分を保持する。

逆に正規形では、$x_b^2$は$B_-$の中心にあるので$H$の部分の重なり差は零である。
$\gamma_a x_bx_c$は、その補集合の各生成元$x_b,x_c$と反可換するため、
隣接位置で符号が変わる$\gamma$部分も重なり差を消す。
全二次重なりは解消し、順序付き基底とKoszul性を得る。
双対を$e_0,e_1,e_2,z$で書くと
\[
 z^2=0,\quad ze_a=e_az,\quad
 e_a^2=-\sum_jH_{ja}e_jz,\quad
 e_be_a-e_ae_b=(-1)^n\gamma_c e_cz\quad(a<b,\ c\ne a,b)
\]
である。補正は全て$z$を一つ増やし、二度使うと零になる。
各$e_a$に重み1、$z$に重み0を与えると随伴次数代数は四生成の可換平方零完全交叉である。
両始点での40個の双対三次重なりは恒等的に解消する。
従って双対は14641型で完全対合を持ち、Koszul性と合わせて左右AS正則性を得る。
再現コード\path{computations/verify_indexed_anticommutative_ore.py}は、
原代数の8個と双対の40個の重なりを全12変数のまま検査する。

$p$周期の場合、対角条件を一周足すと$2\sum_n v_n=0$なので、$u_n$を$p$周期に選べる。
$\gamma$の交代則は奇数周期なら$\gamma=0$を強制する。
正規形は高々2周期なので最小周期の主張も従う。
\end{proof}

\begin{proposition}[反可換Ore四窓の係数空間と第三階数]\label{prop:anticommutative-ore-rank}
四窓72係数での正則条件は階数51の線型条件であり、正則係数空間は21次元である。
正規形\eqref{eq:anti-ore-normal}における第三階数は
\[
 \rho=\begin{cases}6,&\gamma=0,\\8,&\gamma\ne0.\end{cases}
\]
従ってこの表示を持つ低階数対象は全て対称Frobeniusチャートに属し、
\texttt{LOW-UNCLASSIFIED}には入らない。
21は表示係数空間の次元であり、同型類像の次元ではない。
\end{proposition}
\begin{proof}
四窓の係数は$H$の9成分、$\gamma$の3成分、四つの$u_n$の12成分から得られる。
この24変数のパラメータ化の核は、全$u_n$への同じ一次式の加算と$H$の対角成分の補正という3次元である。
従って像の次元は21であり、前定理の必要十分条件から72係数に対する独立条件数は51となる。

最上双対traceを$e_0e_1e_2z$の係数で正規化する。
前節と同様、$z$を二つ含む語は零、一つ含む語のtraceは平方自由の可換積であり、
$z$を含まないtraceは$H,\gamma$に線型である。
第三の二対二行列の$z$を一つ含む二つのブロックは双方階数3であり、
今度は対称な非対角二次テンソルの3次元を像とする。
残る商を$00,11,22,01-10,02-20,12-21$の順に取ると、誘導行列は
\[
 \begin{pmatrix}
 0&-\gamma_2&-\gamma_1\\
 \gamma_2&0&-\gamma_0\\
 \gamma_1&\gamma_0&0
 \end{pmatrix}\oplus0_3.
\]
よって全階数は6にこの行列の階数を加えたものであり、主張を得る。
$\gamma=0$なら正規形の全関係は対称二次テンソルで、双対は可換である。
その可換双対の四二次式は
$z^2$および$e_b^2+\sum_aH_{ab}e_az$で、共通非零零点を持たない。
従って定理\ref{thm:symmetric-frobenius-chart}のチャートに属する。
\end{proof}
\begin{remark}
$\gamma\ne0$で最小周期が常に2とは主張しない。
また$\gamma=0$への帰着を仮定すると、消せない三つの交代方向を誤って失う。
ここでは全方向を保持してから、その非消滅が低階数から外れることを証明した。
\end{remark}

\begin{corollary}[非零交代方向でも1周期となる部分]\label{cor:anti-ore-pure-outer-period}
$H=0$なら、$\gamma$の値にかかわらず正規形は1周期的である。
従って$\gamma\ne0$だけを非1周期性の判定に用いることはできない。
\end{corollary}
\begin{proof}
$x_a$を固定し$X\mapsto-X$とする一次行列は、位置$n$の関係を位置$n+1$の関係へ送る。
あるいは位置$n$で新生成元を$Y=(-1)^nX$と置けば、
$Yx_a-x_aY=\gamma_a x_bx_c$という一定のgraded Ore表示になる。
\end{proof}

% !TeX root = ../papers/nc-p3.tex
\section{対角量子Oreチャートの共通AS判定と低階数の除外}
\label{sec:quantum-ore}
六つの非零係数$q_{ab}$（$0\leq a<b\leq3$）を取り、
三生成量子多項式環$B_q$とその対角自己同型$\sigma(x_a)=q_{a3}x_a$を考える。
$B_q$の順序付き二次基底を用いた任意の$f_a$について、表示
\[
 A(q,f)=k\langle x_0,x_1,x_2,X\rangle/
 (x_bx_a-q_{ab}x_ax_b,\ Xx_a-q_{a3}x_aX-f_a)
\]
を考える。ここで最初の関係は$a<b\leq2$、後者は$a\leq2$である。

\begin{theorem}[全対角量子背景での有限AS判定]\label{thm:quantum-ore-as}
この表示が標準14641型quadratic Koszul AS正則であるための必要十分条件は、
$\delta(x_a)=f_a$が$B_q$上の$\sigma$-導分となることである。
具体的には、三つの式
\begin{equation}\label{eq:quantum-ore-derivation}
 f_bx_a+q_{b3}x_bf_a-q_{ab}f_ax_b-q_{ab}q_{a3}x_af_b=0\quad(a<b\leq2)
\end{equation}
を$B_q$の10個の三次基底で検査すればよい。
これは18個の$f$係数に対する30本の線型方程式であり、$q$の一般性条件を要しない。
特に各固定量子背景で、核の全ファイバーが正則である。
\end{theorem}
\begin{proof}
全ての二次書換えは量子順序付き単項式へ停止する。
重なりは三つの$Xx_bx_a$と一つの$x_2x_1x_0$である。
後者は量子多項式の恒等式、前者の差は正確に\eqref{eq:quantum-ore-derivation}である。
従って標準三次元数20ならこの差は零であり、導分条件が必要である。
逆に導分条件が成立すれば二次Gr\"obner基底で、Hilbert級数$(1-t)^{-4}$とKoszul性を得る。
$X$の重みによる随伴次数代数は四生成量子多項式代数である。
双対側には逆重みを入れると量子外積代数からの全射があり、
Koszul Hilbert逆数公式による14641元数から同型となる。
従って相補積の完全性も量子外積代数から持ち上がり、左右AS正則性を得る。
\end{proof}

\begin{lemma}[量子背景への特殊化と階数下界]\label{lem:quantum-ore-background-rank}
前定理の正則対象について、$A(q,0)$を量子背景とすると
\[
 \rho(\check A(q,0))\leq\rho(\check A(q,f))
\]
である。
\end{lemma}
\begin{proof}
$f_a$を$t f_a$へ置き換える。導分条件は線型なので全$t\in\mathbb A^1$で成立し、
二次書換えから成分ごとに自由な全ファイバー正則族となる。
双対の逆重みの随伴次数は同じ量子外積代数であり、最上traceはその基底で正規化して代数的に選べる。
従って第三階数は一つの代数的行列の階数で、特殊化で増加しない。
$t\ne0$では$X\mapsto tX$が$A(q,t f)\simeq A(q,f)$を与えるため、一般階数は元の階数に等しい。
\end{proof}

\begin{theorem}[対角量子Ore表示から低階数残余は生じない]\label{thm:quantum-ore-low-exclusion}
前定理の正則対象は常に$\rho\geq6$を満たす。
さらに$\rho\leq6$なら、そのindexed類は可換多項式類または対称Frobeniusチャートに属する。
従って\texttt{LOW-UNCLASSIFIED}の対象は、添字保持同型を許してもこの対角量子Ore表示を持たない。
\end{theorem}
\begin{proof}
量子背景の第三階数は命題\ref{prop:quantum-opposite-rank}より6以上である。
前補題により元の階数も6以上で、低階数なら背景・元対象とも階数6である。
同命題の証明から、置換を使わない頂点対角変換により、背景の全六係数を
同じ$s\in\{1,-1\}$へ直せる。
具体的に$q_{ab}d_a/d_b=s$とし、位置$n$の写像を$x_a\mapsto d_a^{-n}x_a$と取る。
補正$f_a$は位置依存になり得るが、三生成部分以外の文字を含まず、
先頭係数は一様な$s$となる。これは全整数での可逆な一次変換であり、正則性を保つ。

$s=1$なら定理\ref{thm:indexed-polynomial-ore}で一定の多項式導分Ore表示へ直せる。
その第三階数が6なら定理\ref{thm:polynomial-ore-rank}から径方向の場合であり、可換多項式のindexed類である。
$s=-1$なら定理\ref{thm:indexed-anticommutative-ore}で$H+(-1)^n\gamma$の正規形へ直せる。
階数6なら命題\ref{prop:anticommutative-ore-rank}から$\gamma=0$なので、対称チャートに属する。
両所属条件と第三階数はindexed同型不変だから、最後の排除も従う。
\end{proof}
\begin{remark}
低階数部分を分類する際に、元対象へPBW・domainを追加仮定していない。
この表示を持つ標準対象なら、三次元数が導分条件とPBW性を強制する。
ただし全正則対象が量子三生成部分や対角Ore表示を持つという構造定理は未証明である。
Ore手法自体の新規性は主張せず、ここでの寄与は共通判定と全低階数部分の既知枝への排除に置く。
\end{remark}

\section{Artin環上の標準化と周期的変形複体}\label{sec:deformation-complex}
以下、変形とは局所Artin $k$-代数$B$上の変形であり、特殊ファイバーとの同定を保持する。
各$C_{ij}$は有限自由$B$-加群とし、同型は特殊ファイバー上で恒等とする。
有限自己同型群による粗モジュライの商と、この基点付き変形問題は区別する。

\begin{lemma}[相対Koszul性と相対4周期性]\label{lem:relative-four}
標準14641型$C_0$の成分ごとに平坦な二次変形$C_B$は相対Koszulであり、
双対の各成分は自由で特殊ファイバーと同じ階数を持ち、相補次数積は完全である。
従って$C_B$はFrobenius標準4周期構造を持つ。
trace変更の頂点スカラー共役と、同型によるtrace移送も$B$上で成立する。
\end{lemma}
\begin{proof}
内部次数$n$を固定した相対bar複体は、有界な有限自由$B$-加群の複体である。
剰余体上ではKoszul性によりホモロジーはホモロジー次数$n$にのみ存在する。
剰余体上の可縮な二項対を、単元の行列成分を使うGaussian消去で$B$上へ持ち上げて除く。
残る複体は一つの次数にのみ自由加群を持つため、ホモロジーは自由で、他の次数では零である。
これを各内部次数で行うと相対Koszul性と、対角Tor、従って二次双対の階数が得られる。
特殊ファイバーで完全な積対合は、その行列式が$B$の単元なので完全である。
最高次数成分は自由階数1だからtraceを選べる。
前節の$\nu$の構成、積の保存、逆双対化は完全対合だけを用い、$B$上でも同じである。
traceの比は$B^\times$に属し、$s_{i+4}/s_i=t_i$は四つの剰余類で解ける。
従って共役性と同型による移送も成立する。
\end{proof}

\paragraph{PBW座標。}
生成元を$x_0<x_1<x_2<x_3$とし、基点の収束二次書換え則を
\[
 x_bx_a\longmapsto F_{ba}(x_0,x_1,x_2,x_3)\quad(a<b)
\]
と書く。右辺は順序付き二次単項式10個の線型結合とする。
位置$i$での変形は$F_{ba}+\alpha_{i,ba}$で表し、$\alpha_i$には$6\cdot10=60$成分がある。
基点での書換えは単項式順序を減らし、変形部分はArtin環の極大イデアルの次数を上げる。
従ってこの二つを組み合わせた有限の還元手順により、Artin環上の正規形を定義できる。
唯一の重なりは$x_cx_bx_a$（$a<b<c$）の4個である。
各重なりの差を順序付き三次単項式20個に展開し、全$4\cdot20=80$成分を零とすることが、
相対Diamond lemmaによる平坦性の必要十分条件になる。
これは通常の数値的Hilbert一致だけをAS正則性の証明に使う議論ではない。

\paragraph{一次複体。}
基底の一次変更を$f_i\in\operatorname{End}(k^4)$とする。
関係の変化と重なりの線型項は
\begin{equation}\label{eq:linear-complex}
 \alpha_i=N_0f_i+N_1f_{i+1},\qquad
 M_0\alpha_i+M_1\alpha_{i+1}=0
\end{equation}
である。$N_0,N_1$は$60\times16$、$M_0,M_1$は$80\times60$行列である。
定義は次の通りで、同梱\path{computations/pbw_complex.py}にもそのまま実装する。
関係$r_{ba}=x_bx_a-F_{ba}$の第$j$因子を$f$で置換し基点で正規化したものが$N_jf$である。
三次重なりを左先行・右先行の二通りで還元し、位置0または1で一度だけ
$F$を$\alpha$に置換した項の差が$M_0\alpha$または$M_1\alpha$である。
全ての語・対・三つ組・単項式は辞書式順序に固定する。

周期$m\mid4$の変形をFourier分解して$\alpha_i=\zeta^i\alpha$と書くと
\begin{equation}\label{eq:fourier-complex}
 k^{16}\xrightarrow{D_1(\zeta)=N_0+\zeta N_1}
 k^{60}\xrightarrow{D_2(\zeta)=M_0+\zeta M_1}k^{80},\qquad\zeta^m=1
\end{equation}
を得る。$D_2D_1=0$は基底変更による変形が結合的であることから従う。
コードでは$M_0N_0=M_1N_1=M_0N_1+M_1N_0=0$も個別に検査する。
一次変形の同値類は$H^2_\zeta=\ker D_2(\zeta)/\operatorname{im}D_1(\zeta)$である。

\begin{lemma}[滑らかな族を用いる局所完全性]\label{lem:smooth-source}
有限型の表示空間の点$x$について、滑らかなパラメータ空間からの射が$x$で接空間を全射するなら、
$x$の表示空間は滑らかであり、その族は形式的局所的に全ての変形を含む。
同じ結論は有限変数・有限形式方程式で与えた完備局所環にも成立する。
\end{lemma}
\begin{proof}
$d=\dim T_x$とする。滑らかな源の完備局所座標から、接写像が同型となる$d$変数の形式的切片を選ぶ。
標的の完備局所環を$k[[u_1,\ldots,u_d]]/I$と書く。
切片への写像における$u_i$の像の線型部は独立なので、形式的逆写像定理により
$k[[u_1,\ldots,u_d]]\to k[[t_1,\ldots,t_d]]$は同型である。
この写像が$I$を零へ送るため$I=0$であり、結論を得る。
\end{proof}
\begin{remark}
\emph{源の滑らかさ}を省いて接写像の全射性だけを使ってはいけない。
例えば尖点曲線$y^2=x^3$の$\mathbb A^2$への包含は原点で接空間を全射するが、
曲線は平面の既約成分ではない。
本稿では\cite[Proposition 5.1]{BILN}をこの一般性では用いず、上の補題を直接使う。
\end{remark}

\section{量子多項式成分：同型類次元6と3}\label{sec:quantum}
$q_{ab}\in k^\times$（$a<b$）に対し
\[
 A_q=k\langle x_0,x_1,x_2,x_3\rangle/
 (x_bx_a-q_{ab}x_ax_b\mid a<b)
\]
と置く。$q_{ba}=q_{ab}^{-1}$と約束する。
全ての$q$で順序付き単項式が基底となり、双対量子外積代数の相補次数積が完全なので、
$A_q$は標準14641型である。次の開集合を考える。
\[
 U=\{q:c_{abc}:=q_{ab}q_{bc}/q_{ac}\ne1\quad(a<b<c)\}.
\]

\begin{theorem}[量子多項式族の二つの同型分類]\label{thm:quantum-orbits}
$U$上で次が成立する。
\begin{enumerate}
\item graded同型は、四生成元の同じ置換による$q$の変換に尽きる。
従って同型類の有理商の次元は6である。
\item 添字保持同型$\check A_q\simeq\check A_{q'}$は、上の置換に加えて
\begin{equation}\label{eq:quantum-twist}
 q'_{ab}=q_{ab}\,s_a/s_b\qquad(s_a\in k^\times)
\end{equation}
を許したものに尽きる。この同型類の次元は3である。
\item 置換で割る前の商は、例えば
$(c_{012},c_{013},c_{023})$で記述できる。
$c_{123}=c_{012}c_{023}/c_{013}$であり、独立な不変量は三つである。
\end{enumerate}
\end{theorem}
\begin{proof}
多重線型化は$p_bq_a-q_{ab}p_aq_b=0$である。
$p$の非零成分が三つ以上あると、対応する三点で$c_{abc}=1$が必要となる。
また$p$の支持と$q$の支持は一致する。従って$U$上の点スキームは六つの座標直線の和であり、
その特異点は四つの座標頂点である。関係対応は各座標頂点を同じ頂点へ送る。

この配置を保つ射影線型写像は四頂点を置換するので単項行列である。
graded同型に対しては両テンソル因子に同じ単項行列が働くため、
対角成分は$q$を変えず、置換だけが残る。

$\ZZ$-代数同型の次数一部分を$g_i$とする。
各$g_i$は単項行列であり、頂点対応の保存により隣接位置の置換は同じである。
共通の置換を除き、$g_i(x_a)=t_{i,a}x'_a$と書くと
\[
 \frac{q'_{ab}}{q_{ab}}
 =\frac{t_{i,a}t_{i+1,b}}{t_{i,b}t_{i+1,a}}.
\]
$i=0$の式だけから$q'/q$は$s_a/s_b$の形になる。
逆に式\eqref{eq:quantum-twist}があれば$t_{i,a}=s_a^{-i}$とすることで全位置に同型を作れる。
有効な対角twist群は$(k^\times)^4/k^\times$、次元3であり、自由に作用する。
その商の次元は$6-3=3$である。
三角積はこの作用で不変であり、三つの辺$q_{01},q_{02},q_{03}$を1に規格化すれば、
残りの三つの辺はこれらの三角積で復元される。
最後の$S_4$という有限群の商は次元を変えない。
\end{proof}

\paragraph{次数零自己同型とstack次元。}
一般点でgraded自己同型群は対角$(k^\times)^4$である。
従ってgraded族のstack次元と同型類次元6を混同してはならない。
$\ZZ$-代数ではさらに各頂点のスカラーがあり、ここでの3もstack次元の主張ではない。

\begin{theorem}[一般量子多項式点の局所成分]\label{thm:quantum-local}
量子多項式族のある非空開部分では、標準AS正則$\ZZ$-代数の全Artin変形は、
同型を除いて量子多項式族から得られる。
従ってこの成分の一般点で、graded同型類の局所次元は6、
裸の添字保持同型類の局所次元は3である。
この点の近傍の次元差はtwist方向の忘却によるものであり、
非1周期的対象を加えたことによる増大ではない。
この主張の行列階数部分は計算補助証明である。
\end{theorem}
\begin{proof}
次の有理数点を使う。順序は$(01,02,03,12,13,23)$である。
\begin{equation}\label{eq:quantum-sample}
 q_*=(2,3,1/6,5,2/5,15).
\end{equation}
この点では各行の積$\prod_{b\ne a}q_{ab}=1$であり、全ての三角積が1と異なる。
双対量子外積代数のNakayama作用は次数一で$-I$、最高次traceを$(-1)^i$倍すると
厳密4移動がFrobenius標準構造になる。

式\eqref{eq:fourier-complex}の厳密階数は
\[
\begin{array}{c|rrr}
\zeta&\rk D_1&\rk D_2&\dim H^2_\zeta\\\hline
1&12&42&6\\
-1&15&45&0\\
\sqrt{-1}&15&45&0\\
-\sqrt{-1}&15&45&0
\end{array}
\]
である。さらに六つの$q$の微分を$K$と書くと
$D_2(1)K=0$、$\rk(D_1(1)\mid K)=18$である。
これらは\path{computations/verify_moduli_jets.py}で$\QQ$および$\QQ(\sqrt{-1})$上で検証する。
行列の定義は前節の有限還元手順であり、入力は式\eqref{eq:quantum-sample}だけである。

四窓表示の接空間は$6+57=63$次元である。
滑らかな族$\GL_4^4\times(k^\times)^6$の接写像は、基底変更の57方向と$q$の6方向により
全63次元を張る。補題\ref{lem:smooth-source}により四周期PBW表示の形式的局所空間は滑らかで、
この族が全変形を含む。
相対Diamond lemmaにより、ここで使う形式方程式は平坦二次変形を記述する。
補題\ref{lem:relative-four}によって、裸の標準AS正則変形も標準4周期座標に直せるため、
この議論から漏れる無周期の形式的方向はない。

従って裸の同型類は量子多項式族における同型類の商で尽くされる。
前定理の三次元twist群と有限置換群による分類を適用して局所次元3を得る。
graded側には同じ議論を$\zeta=1$で適用して局所次元6を得る。
非零小行列式とKodaira--Spencer写像の階数条件は開条件なので、結論は非空開部分で成立する。
滑らかな形式的近傍は既約であるから、これは族内の次元だけでなく、対応する一般点の局所成分の次元である。
\end{proof}

\begin{remark}
この定理は全ての成分が量子多項式成分であるとの主張ではない。
特に次節のClifford点は異なる変形挙動を示す。
また一般の4周期表示の接空間の6を、裸の$\ZZ$-代数の同型類次元と読み替えてはいけない。
標準化しない周期構造と、twistによる同一視の余分な自由度が残っているためである。
\end{remark}

\section{Clifford特殊点の一次変形と二次障害}\label{sec:clifford}
式\eqref{eq:period-two-family}の$\lambda=0$を$C_0=\check A_0$とする。
ここでは$x_0=W,x_1=Y,x_2=Z,x_3=X$の順序を使う。書換えは
\begin{equation}\label{eq:clifford-rules}
\begin{array}{lll}
YW\mapsto-WY,& ZW\mapsto-WZ,& ZY\mapsto-YZ+W^2,\\
XW\mapsto-WX+Z^2,& XY\mapsto-YX+Z^2+W^2,& XZ\mapsto-ZX.
\end{array}
\end{equation}
この六規則の四つの重なりは解消する。
双対を$x,y,z,w$（それぞれ$X,Y,Z,W$に双対）で書くと
\[
 E_0=k[x,y,z,w]/(x^2,y^2,z^2+xw+xy,w^2+yz+xy)
\]
である。四つの二次式の共通零点は原点だけなので正則列であり、
$E_0$はHilbert級数$(1+t)^4$のFrobenius代数である。
元代数の二次PBW性と併せて$C_0$の標準14641型を確認できる。
$E_0$は可換なので一定traceのNakayama作用は恒等であり、厳密4移動がFrobenius標準である。

\begin{theorem}[Clifford点の変形接空間]\label{thm:clifford-tangent}
式\eqref{eq:clifford-rules}に対するFourier複体の厳密階数は
\begin{equation}\label{eq:clifford-ranks}
\begin{array}{c|rrr}
\zeta&\rk D_1&\rk D_2&\dim H^2_\zeta\\\hline
1&15&36&9\\
-1&15&39&6\\
\sqrt{-1}&15&45&0\\
-\sqrt{-1}&15&45&0
\end{array}
\end{equation}
である。選んだgraded表示$A_0$から裸の$C_0$への忘却の接写像は単射で、
その像は9次元の偶成分、余核は6次元の交代成分である。
$C_0$の基点付き標準AS正則$\ZZ$-代数変形の接空間は15次元である。
\end{theorem}
\begin{proof}
階数は式\eqref{eq:clifford-rules}から前節の算法で作った行列を有理数上で消去して得る。
同梱コードは$D_2D_1=0$、各階数、核と像の商の基底を独立に検査する。
四周期変形は四つのFourier成分の直和だから、その商接空間は$9+6=15$であり、
graded変形の像は$\zeta=1$成分である。

裸の変形への移行でこの次元を失わないことを確認する。
$C_0$の任意の無限小自己同型は、標準4周期構造の自然性により、
頂点スカラーを除けば4周期的である。
各Fourier成分で$D_1$の核はスカラー行列の一次元だけなので、
射影的な無限小自己同型は零である。
Artin変形上で恒等に還元される自己同型も、極大イデアルのべきに沿う帰納法により頂点スカラーに限る。
従って厳密4周期構造と相対標準4周期構造との差は、頂点スカラー共役で取り除ける。
二つの裸の同型の比較にも同じ議論が適用できる。
補題\ref{lem:relative-four}により全ての裸の変形を4周期表示にできるので、
ここでは4周期構造を忘れても、スカラーを除いた基点付き変形接空間は変わらない。
\end{proof}

\paragraph{座標の固定。}
各モードで$\ker D_2$と$\operatorname{im}D_1$を通常の有理数rrefで計算し、
像の列基底の後に核の列基底を並べて、ピボットとして新しく残る核の列を商基底とする。
この決定的な手順で偶基底$e_0,\ldots,e_8$、交代基底$o_0,\ldots,o_5$を固定する。
全係数は\path{research/moduli_certificate.json}に記録する。
特に$o_0$は$YW$の右辺へ$ZX$を、$o_5$は$XZ$の右辺へ$WY$を加える基底ベクトルである。
位置$i$での交代変化には$(-1)^i$を掛ける。

\begin{theorem}[全15方向の二次障害]\label{thm:clifford-quadratic}
上の基底で、一次変形$e+o$が二次まで持ち上がるための必要十分条件は
\begin{equation}\label{eq:clifford-Q}
 Q(o)=o_0o_5+o_1o_2-o_2o_3+o_3o_4=0
\end{equation}
である。偶方向同士と偶・交代の混合の二次障害は零である。
$Q$の階数は6である。従って基点付き変形のminiversal環を$R$と書くと、適当な形式座標で
\[
 R/\mathfrak m^3\simeq
 k[e_0,\ldots,e_8,o_0,\ldots,o_5]/(Q,(e,o)^3).
\]
この命題は\emph{二次まで}の主張であり、完備局所環$R$全体を決定するものではない。
\end{theorem}
\begin{proof}[計算補助証明]
書換え則を$F+t\alpha+t^2\beta$とし、三次重なりの$t^2$係数を取る。
その式は$D_2\beta+\operatorname{Ob}_2(\alpha)=0$である。
$\alpha_i=e+(-1)^i o$なので、二次残差は偶・交代成分だけを持つ。
$L_+$を$D_2(1)^{\mathsf T}$の核の転置、$L_-$を$D_2(-1)^{\mathsf T}$の核の転置とする。
それぞれ$44\times80$、$41\times80$であり、
残差が$D_2$の像に入ることは$L_+\operatorname{Ob}_{2,+}=0$と
$L_-\operatorname{Ob}_{2,-}=0$に同値である。

15個の基底ベクトルとその異なる二つの和で残差を評価する。
二次式であるため、この$15+\binom{15}{2}=120$回の厳密計算で全係数が決まる。
得られた85成分の二次式は全て$0,Q,-Q$のいずれかであり、少なくとも一つが$Q$または$-Q$である。
これを\path{computations/verify_moduli_jets.py}が有理数のみで検査する。
$\pm\sqrt{-1}$成分には残差がないので、四周期の二次補正を許しても条件は変わらない。
前定理の相対標準化により裸の変形でも同じ条件である。

この検査は単に正の持上げ候補を見つけるものではない。
$L_\pm$が左核全体の基底なので、像への所属の必要性と十分性を同時に証明している。
また一般の二次式の全係数を検査しており、特定のパラメータ値だけの数値観察ではない。
$rref$によって線型方程式と基底変更方向を除いた局所方程式の二次部分も同じ残差であるため、
表示した$R/\mathfrak m^3$を得る。
最後に
\[
 z_1=o_0,\ z_2=o_5,\ z_3=o_1-o_3,\ z_4=o_2,\ z_5=o_3,\ z_6=o_4
\]
は可逆な線型変換であり、$Q=z_1z_2+z_3z_4+z_5z_6$となる。従って階数は6である。
\end{proof}

\begin{corollary}[Pl\"ucker型の最初の障害]\label{cor:plucker-obstruction}
六つの交代方向に対する二次障害の零点錐は、
Pl\"ucker埋込み$\Gr(2,4)\subset\mathbb P^5$上のアフィン錐と線型同型である。
これは\emph{二次障害が許す候補錐}の記述であり、実際の接錐や成分全体の同定ではない。
\end{corollary}
\begin{proof}
$\Gr(2,4)$のPl\"ucker式$p_{01}p_{23}-p_{02}p_{13}+p_{03}p_{12}=0$に
$(p_{01},p_{23},p_{02},p_{13},p_{03},p_{12})=(z_1,z_2,z_3,-z_4,z_5,z_6)$を代入する。
\end{proof}

\begin{proposition}[二つの積分可能な軸と混合障害]\label{prop:two-axes}
式\eqref{eq:clifford-rules}のうち、それぞれ一方だけを
\[
 YW\mapsto-WY+(-1)^i aZX,\qquad
 XZ\mapsto-ZX+(-1)^i bWY
\]
に置き換えると、各軸上で全パラメータの二次PBW族を得る。
原点での形式変形は相対AS正則である。
しかし両軸の接方向を同時に取った一次変形には二次障害$ab$があり、
$a=b=1$という接方向は、任意の二次補正を許しても持ち上がらない。
\end{proposition}
\begin{proof}
各軸について二つの始点と四つの三次重なり、合計8個の重なり差は$\QQ[a]$または$\QQ[b]$上で恒等的に零である。
同梱コードは記号多項式としてこれを確認する。
$b$軸では通常の次数付き辞書式順序を使える。
$a$軸では重み$(W,Y,Z,X)=(1,4,1,1)$を先に比較し、同じ重みでは$X>Z>Y>W$の辞書式順序を使うと
全規則が減少する。Diamond lemmaから順序付き単項式基底と平坦性を得る。
形式的なAS正則性は補題\ref{lem:relative-four}による。
新しい$a$軸の\emph{全閉ファイバー}のAS正則性は、この議論だけでは主張しない。
接方向は$o_0,o_5$なので、式\eqref{eq:clifford-Q}を制限すると$ab$になる。
\end{proof}

\subsection{まだ結論できないこと}
上の計算から、Clifford点における成分次元が14であるとは結論できない。
三次以上の独立な関係式が残る可能性がある。
また$9\to15$は\emph{選んだgraded表示}$A_0$での忘却写像の微分である。
同じ裸の$C_0$に別の1周期構造を付けたときのgraded表示は別途調べる必要がある。
全ての1周期的部分の局所余次元を6または5と断定してはならない。

\begin{conjecture}[検証対象の局所模型]\label{conj:clifford-model}
Clifford点の基点付き変形環は
\[
 k[[e_0,\ldots,e_8,z_1,\ldots,z_6]]/
 (z_1z_2+z_3z_4+z_5z_6)
\]
と同型である。
\end{conjecture}
この予想を裏付ける確定事項は定理\ref{thm:clifford-quadratic}の二次切断までである。
この模型が正しければ局所次元14となるが、その十分性・高次持上げ・全周期1構造の分類は未解決である。


\section{Clifford点における低階数切断の局所被覆}
\label{sec:clifford-rank-slice}
この節の基点は式\eqref{eq:clifford-rules}の$C_0$であり、
全Cliffordチャートや全低階数領域を一度に分類するものではない。

\begin{lemma}[有限双対のArtin平坦性から元代数の平坦性]\label{lem:dual-artin-flatness}
局所Artin環$B$上の二次四窓について、双対$E_B$が各次数で有限自由であり、
特殊ファイバー$E_0$がKoszulなら、$E_B$は相対Koszulである。
その二次双対の各成分も自由で、特殊ファイバーと同じHilbert元数を持つ。
特に$C_0$を特殊ファイバーとする有限Frobenius候補のArtin変形は、
元代数の成分ごとに平坦な変形を与える。
\end{lemma}
\begin{proof}
補題\ref{lem:relative-four}と同じGaussian消去を、今回は$E_B$のbar複体に適用する。
内部次数を固定すれば有限自由複体であり、剰余体上の可縮対を単元で消去すると、
ホモロジー次数と内部次数が等しい一つの次数だけに自由加群が残る。
従って対角Torは自由で、それを双対にした二次双対の各成分も自由となる。
有限Frobenius候補の定数階数条件は次数$0,\ldots,4$での自由性を保証する。
次数5の特殊ファイバーは零だからNakayama補題で次数5自体が零となり、全高次も零である。
従ってこの場合も全次数で前半の議論を適用できる。
\end{proof}

$w_0$を$C_0$の最上traceテンソルとし、$0123$の係数を1に正規化する。
第三の二対二行列を$M_0$とする。行番号は$4a+c$、列番号は$4b+d$である。
$M_0$は対称な階数6の$16\times16$行列であり、通常の有理数nullspace算法で
$K\in\operatorname{Mat}_{16,10}$を$\ker M_0$の列基底として選ぶ。
一次変形のtrace変化を$\dot w$、対応する行列を$\dot M$とすると、
7次小行列式の一次条件は$K^{\mathsf T}\dot M K=0$である。

\begin{proposition}[低階数条件の一次切断と障害のPfaffian]\label{prop:clifford-rank-derivative}
既定の偶基底$e_0,\ldots,e_8$と交代基底$o_0,\ldots,o_5$に対し、
写像$\dot w\mapsto K^{\mathsf T}\dot M K$は全偶方向と全基底変更方向で零となり、
交代6次元上では単射である。
従って四周期平坦表示の75次元接空間は、低階数条件によって69次元に切られる。
交代方向の六つの$10\times10$行列は共通の6次元核を持ち、商の基底を適切に取ると
\[
 S(o)=\begin{pmatrix}
 0&o_5&o_2-o_4-o_5&-o_4\\
 -o_5&0&o_1&o_1-o_3\\
 -o_2+o_4+o_5&-o_1&0&-o_0-o_1+o_3\\
 o_4&-o_1+o_3&o_0+o_1-o_3&0
 \end{pmatrix}
\]
となる。このPfaffianは既知の二次障害$Q$の負である。
\end{proposition}
\begin{proof}
双対関係を順序付き二次語$m$で番号付けすると、その基底は
$m+\sum_r F_r(m)\,\operatorname{desc}(r)$である。
各一次代表$\alpha$について、この基底の変化を三つの位置へ挿入した行列を$D_1$、
基点の挿入行列を$D_0$とすれば、trace変化は
\[
 D_0^{\mathsf T}\dot w=-D_1^{\mathsf T}w_0,\qquad \dot w_{0123}=0
\]
で一意に決まる。ここで$D_0$は$256\times480$行列である。
480本の式と正規化を有理数上で解き、各解を元の全式へ代入して確認する。
\path{computations/verify_clifford_rank_slice.py}はこの481本の式を15代表について一括で検証する。

同コードで偶9方向の像は零、交代6方向の像は階数6となる。
100成分を行に並べた交代微分行列では、0始まりの行番号
$12,15,19,25,29,59$の小行列式が$-1$である。
64個の初等基底変更方向についても像が零であることを直接確認する。
定理\ref{thm:clifford-tangent}の複体より表示の接次元は$60+9+6=75$なので、切断後は69となる。

共通核を除くと、有理数rrefで選んだ$K$座標の番号$1,2,5,9$が商の基底を与え、
表示した$S(o)$を得る。最後に
\[
 \operatorname{Pf}S=S_{01}S_{23}-S_{02}S_{13}+S_{03}S_{12}
 =-o_0o_5-o_1o_2+o_2o_3-o_3o_4=-Q
\]
を展開すればよい。これは既知の障害の行列的解釈であり、
高次障害や全変形環の新しい決定を主張するものではない。
\end{proof}

\begin{theorem}[Clifford点の低階数近傍は対称チャートで被覆される]\label{thm:clifford-low-rank-local}
テンソル復元定理の有限Frobenius候補空間に、第三の7次小行列式の消滅を課した
局所有限型スキームを$\mathfrak T_{\leq6}$とする。
このスキームは$[w_0]$で滑らかな次元69の表示空間であり、
$[w_0]$を含むある$G$不変Zariski開部分は対称Frobeniusチャートで被覆される。
従ってこの開部分の候補は全てKoszul AS正則であり、
\texttt{LOW-UNCLASSIFIED}はこの近傍に現れない。

より具体的に、既定の9個の偶一次代表をそのまま基点の関係へ加えた9変数族は、
基点近傍で全ての基点付き低階数Artin変形を、添字保持同型を除いて含む。
これは低階数切断内の局所主張であり、無制限のClifford変形環や全低階数領域の被覆ではない。
\end{theorem}
\begin{proof}
定数階数の行列条件は、選んだ非零小行列式の近傍でkernelとcokernelを有限自由にする。
従ってテンソル復元条件は局所有限型の階数層として定義でき、
四窓への写像と最上交差空間の生成線への逆写像はこの階数層上で正則である。
Artin環上でも最上交差空間は自由階数1なので、この写像は無限小にも単射である。
補題\ref{lem:dual-artin-flatness}により、その全Artin接方向は平坦PBW四窓の接方向に含まれる。
前命題から$\dim T_{[w_0]}\mathfrak T_{\leq6}\leq69$となる。

一方、偶基底の9代表は全て二次平方への係数変更だけである。
従ってこの9変数族の関係空間は対称テンソル内にとどまり、双対は可換な四二次式の商である。
基点で完全交叉なので、その非空開近傍$V\subset\mathbb A^9$は
定理\ref{thm:symmetric-frobenius-chart}のチャートへ入る。
特に$V$全体が平坦かつ階数6の族である。
滑らかな源$G\times V$から$\mathfrak T_{\leq6}$への射を考えると、
基底変更の60方向と独立な偶9方向により接写像の階数は69である。
上の上界と合わせて接空間を全射する。
補題\ref{lem:smooth-source}より、標的はこの点で滑らかな次元69となる。

両空間が滑らかで接写像が全射なので、射は対応点で滑らかである。
滑らかな射の像は開であり、その像は全て対称チャートの像に含まれる。
この開近傍を$G$で飽和すれば主張する不変開部分を得る。
形式的な局所完全性も補題\ref{lem:smooth-source}による。
無限小の裸の同型と標準四周期表示の比較は定理\ref{thm:clifford-tangent}と
補題\ref{lem:relative-four}で処理済みであり、全局所対象の漏れはない。
\end{proof}
\begin{remark}
69は基底付き表示空間の局所次元で、stack次元や全Clifford成分の次元ではない。
基点付き同型を除いた接方向は9であるが、有限自己同型で割る粗モジュライの滑らかさは主張しない。
また、低階数の閉条件を課さない全Clifford点には交代変形と障害$Q$が残る。
\end{remark}

\section{四次のKoszul検査を通る非KoszulなFrobenius境界}
\label{sec:nonkoszul-boundary}
有限Frobenius性と元代数のKoszul性を分ける必要は形式上だけではない。
以下では対称チャートの閉包内で、この差を五次に初めて検出する具体例を与える。

\begin{theorem}[有限Frobenius領域内の非対称化境界]\label{thm:nonkoszul-symmetric-boundary}
$Z_{\mathrm{sym}}$には有限Frobenius候補$[w_0]$であって対称化も交代化もできない点が存在する。
その第三階数は6であり、元の四窓代数$C$は全始点で四次Koszul検査を通るが、Koszulではない。
特に対称チャートの飽和像$S$は、有限Frobenius低階数候補領域の中で相対的に閉ではない。
ただし$C$は本研究の標準型Koszul AS正則対象ではなく、未分類AS枝の存在証明にはならない。
\end{theorem}
\begin{proof}
可換二次式の四次元部分空間
\[
 L_t=\langle e_0^2-e_0e_3,\ e_1^2-e_1e_3,\ e_2^2-e_2e_3,
                  \ e_3(e_0+e_1+e_2+t e_3)\rangle
\]
を取る。最初の三式の射影共通零点は
$[\epsilon_0:\epsilon_1:\epsilon_2:1]$, $\epsilon_i\in\{0,1\}$の八点である。
最後の式の値は$\epsilon_0+\epsilon_1+\epsilon_2+t$なので、
$t\notin\{0,-1,-2,-3\}$なら$L_t$は基点自由な完全交叉である。
$t=0$では共通零点$[0:0:0:1]$が残る。

$B_t=\operatorname{Sym}(V^*)/(L_t)$と置く。
対称三テンソル$b=(b_{abc})$が$(B_t)_3^*$に属する条件は
\[
 b_{iia}-b_{i3a}=0\quad(0\leq i\leq2,\ 0\leq a\leq3),\qquad
 \sum_{i=0}^2 b_{i3a}+t b_{33a}=0\quad(0\leq a\leq3).
\]
20個の対称座標に対するこの16本の式に、
$012,133,233,333$の四座標を指定する式を加えた20次行列の行列式は$-2$である
（行順序は再現コードの通り）。
従って全ての$t$で$(B_t)_3$は4次元であり、その双対基底$b_0(t),\ldots,b_3(t)$を
上記四座標が標準基底となるよう一意に多項式で選べる。
因子$0,1,3$に三テンソルを、因子2に対応する基底を置いて
\[
 (w_t)_{abjc}=(b_j(t))_{abc}
\]
と定める。これは全$t$で非零な多項式テンソル族である。
完全交叉となる$t$では、その対称最上traceテンソルの一因子縮約は
$(B_t)_3^*$の基底を与えるので、因子2の可逆変換で$w_t$へ移る。
従って一般の$w_t$は$S$に属し、$w_0\in Z_{\mathrm{sym}}$である。

$t=0$の有理数線型計算は次の通りである：
\[
\begin{array}{c|c}
検査 & 値\\\hline
一因子階数 & (4,4,4,4)\\
三つの二対二階数 & (6,6,6)\\
\dim W_{i,3},\ \dim W_{i,4} & (4,1),(4,1),(4,1),(4,1)\\
\dim C_{i,i+4} & (35,35,35,35)\\
\dim C_{i,i+5} & (63,56,56,56)
\end{array}
\]
ここで三・四次交叉は復元した二窓の直交関係を挿入した行列の核で計算し、
最後の二行は元の六関係を挿入した行列の余核で計算する。
とくに五次挿入行列の階数は$(961,968,968,968)$である。
全係数・行列・Gauss消去の再現は
\path{computations/verify_nonkoszul_boundary.py}に置く。
テンソル復元定理により$w_0$は有限Frobenius候補であり、相補次数の対合も完全である。
さらに線型対称化空間は
\[
 S_1(w_0)=S_3(w_0)=kI,\qquad S_2(w_0)=kE_{33}
\]
となるので対称化できない。符号付き三空間の可逆性判定も失敗するので交代化もできない。
四次の値35は全ての$\delta_i$が零であることを意味する。
しかしKoszulなら、双対のHilbert列$(1,4,6,4,1)$から
$h_C(z)=(1-z)^{-4}$となり、全ての五次元数は56でなければならない。
始点0の値63がこれに反する。より正確には五次の$d_3$階数は全始点で36であり、
命題\ref{prop:quintic-koszul-obstruction}により$H_3=0$、$H_2$の次元は$(7,0,0,0)$である。
\end{proof}

\begin{remark}
この例は有限Frobenius性と四次Koszul切断の十分性をともに否定する。
五次での二階数条件の十分性や、一般の有限次数判定は主張しない。
また$S$の有限Frobenius領域内の非閉性と、Koszul AS領域内での非閉性を区別する。
後者、および対称チャート外の低階数AS対象の存在はなお未解決である。
\end{remark}

% !TeX root = ../papers/nc-p3.tex
\section{三因子対称境界と有限切断代数}
\label{sec:truncated-quadric-boundary}
三次商の双対基底から作る前節のテンソルは、有限切断代数のtrivial extensionとして理解できる。
この記述は新しいAS正則性の十分条件ではないが、個々の大行列に依存せず境界問題を整理する。

\begin{lemma}[四層二次代数のKoszul性]\label{lem:four-layer-koszul}
頂点$0,1,2,3$を持ち、矢が$i$から$i+1$へだけ向かう箙の二次商代数はKoszulであり、
大域次元は3以下である。
\end{lemma}
\begin{proof}
経路の長さは始点と終点の差に等しく、最大3である。
単純加群の極小射影分解では、$p$番目の項の内部次数は$p$以上であり、
始点から終点へ向かう正次数成分の積として長さ3以下でなければならない。
従って$p>3$の項は零である。
第1項は次数1生成、第2項は二次関係による次数2生成である。
第3項は内部次数3以外を持ち得ないので、分解は線型である。
\end{proof}

$L\subset\operatorname{Sym}^2(V^*)$を四次元とし、
$B=\operatorname{Sym}(V^*)/(L)$が$\dim B_1=4$, $\dim B_2=6$, $\dim B_3=4$を満たすとする。
三次双対基底$b_0,\ldots,b_3$から、因子$0,1,3$が対称なテンソル
$w_{abjc}=(b_j)_{abc}$を作る。
有限有向代数$A_L$を、順序づけた頂点$3,0,1,2$に対して
\[
 e_iA_Le_j=B_{\operatorname{pos}(j)-\operatorname{pos}(i)}
 \quad(\operatorname{pos}(i)\leq\operatorname{pos}(j))
\]
とし、逆向き成分は零、積は$B$の積で定める。
ここで$\operatorname{pos}(3)=0$, $\operatorname{pos}(0)=1$,
$\operatorname{pos}(1)=2$, $\operatorname{pos}(2)=3$である。
$A_L$は四層二次代数なのでKoszulであり、$\dim A_L=32$である。
$D A_L=\operatorname{Hom}_k(A_L,k)$の$B_r$双対成分に次数$4-r$を与え、
通常のtrivial extension
\[
 T_L=A_L\ltimes D A_L,
 \qquad (a,f)(b,g)=(ab,a\cdot g+f\cdot b)
\]
にこの次数を入れる。

\begin{proposition}[三因子対称Frobenius候補のtrivial extension表示]
\label{prop:partial-symmetric-trivial-extension}
上記$w$がテンソル復元定理の有限Frobenius条件を満たすなら、
復元した二次双対を四頂点で折り畳んだ代数$E(w)$は、次数付き代数として$T_L$と同型である。
従ってこの部分のKoszul判定は、$T_L$のKoszul判定と一致する。
この主張に$B$の完全交叉性は仮定しない。
\end{proposition}
\begin{proof}
$T_L$の次数一成分は、三本の前向き辺上の$B_1$と戻り辺$2\to3$上の$B_3^*$である。
trivial extensionのtraceは、$D A_L$成分を$1\in A_L$で評価する写像であり、
相補次数の積を完全に対合させる。
三つの$B_1$と一つの$B_3^*$の積のtraceは
\[
 (f,a,b,c)\longmapsto f(abc)
\]
である。選んだ基底でのこの四階テンソルが、まさに$w$とその巡回回転である。
各始点の$T_L$の次数別次元は$(1,4,6,4,1)$となる。

$w$の一因子階数4と隣接二対二階数6により、次数一成分の積は
$T_L$の次数2,3,4を全て生成する。実際、その積写像の階数は、完全trace対合を使うと
対応する一因子・二因子縮約の階数に等しい。
従って$T_L$は次数一生成であり、その二次関係は$R_i(w)^\perp$そのものである。
よって二次表示から全射$E(w)\twoheadrightarrow T_L$を得る。
有限Frobenius仮定から両者の全次元は64で等しいので、これは同型である。
\end{proof}

\begin{remark}[既存のhigher zigzag構成との関係]
Grant\cite[Definitions 2.3, 2.5]{GrantZigzag}は、有限大域次元のKoszul代数$\Lambda$に対し、
その二次双対の符号捩れtrivial extensionをhigher zigzag algebraとして定義している。
ここでは$\Lambda=A_L^!$と取れる。
四層の長さ符号は頂点上の$\operatorname{diag}(1,-1,1,-1)$による内部自己同型なので、
その符号捩れは次数零の可逆元で除ける。
従って上の$T_L$は、この既知構成の四層の場合としても記述できる。
trivial extensionという構成自体の新規性や、それだけからのKoszul性は主張しない。
本研究で必要なのは、基点境界で二次表示とKoszul性がどの段階で破れるかの判定である。
\end{remark}

\begin{lemma}[一般点の有限次数障害による族全体の排除]
\label{lem:generic-quintic-exclusion}
既約なパラメータ空間$X$上の正則な射影テンソル族で、
全ての隣接二対二階数が6以下であるものを取り、
有限Frobenius条件と全四次Koszul検査を通る部分$U\subset X$が空でないとする。
$U$は開部分である。
一般点で五次Koszul微分の階数が
$\operatorname{rk}d_2=a$, $\operatorname{rk}d_3=b$なら、
全ての$U$上の点で$\dim H_2\geq120-a-b$である。
特に$a+b<120$なら、この族は標準型Koszul AS対象を一つも含まない。
\end{lemma}
\begin{proof}
隣接階数の上界により、有限Frobenius条件は$X$内で開であり、四次検査の開性と合わせて$U$は開である。
$U$上では必要な次数2,3,4の商空間は一定階数であり、
$K_2,K_3,K_1$はそれぞれ階数120,40,140のベクトル束になる。
二つの微分はベクトル束の射なので、その階数は一般点の値を超えない。
従って
$\dim H_2=120-\operatorname{rk}d_2-\operatorname{rk}d_3\geq120-a-b$である。
標準型Koszul AS点は必ず$U$に入るので、最後の主張も従う。
\end{proof}

\begin{theorem}[係数を動かした基点境界族の一様なAS排除]
\label{thm:generic-one-point-uniform-exclusion}
$c_0c_1c_2\ne0$とし、
\[
 L_{c,t}=\langle e_i^2-e_ie_3\ (i=0,1,2),\
                    e_3(c_0e_0+c_1e_1+c_2e_2+t e_3)\rangle
\]
から、四座標$012,133,233,333$を正規化した三次双対基底によるテンソル$w_{c,t}$を作る。
$t=0$の係数族は、全て$Z_{\mathrm{sym}}$に属するが、標準型Koszul AS対象を一つも含まない。
より正確には、有限Frobenius条件と全四次検査を通る任意の係数値で、
始点0の五次Koszulホモロジーの次元は7以上である。

一般点では有限Frobenius条件と全四次検査が成立し、
五次の$H_2$次元は$(7,0,0,0)$、$H_3$は全始点で零である。
始点0の七次元空間のうち六次元は、前節と同じ平方型関係と左乗法の核から来る。
ただし係数全域の有限Frobenius層や、一般の四二次式の基点境界を分類したとは主張しない。
\end{theorem}
\begin{proof}
最後の式を$c_0$で割り、$u=c_1/c_0$, $v=c_2/c_0$と正規化する。
$t=0$のテンソル族は$X=(\mathbb G_m)^2$上の正則族になる。
三次の16本の消去式と四本の正規化式を合わせた行列の行列式は$-2uv$である。
従って指定三次基底は$X$全体で定まり、特にその第3ベクトルは$333$のcubeである。
他の$t$も含む行列式は$-2c_0^2c_1c_2$なので、$t$方向にも同じフレームが定まる。
各$c$について、最初の三二次式の八点上で最後の式が消える$t$は有限個である。
それ以外の$t$では基点自由完全交叉になるから、$w_{c,0}\in Z_{\mathrm{sym}}$である。

三因子が対称であり、その二因子に$L_{c,0}$が作用して零になるので、
全ての二対二階数は6以下である。
$u=v=1$では前節の例に一致し、有限Frobenius条件と全四次検査を通る。
従ってそれらを通る部分$U\subset X$は空でない開部分であり、一般点を含む。
同じ標本では、始点1,2,3の五次複体も正確である。
その正確性も一定次元の複体の開条件なので、一般点では保たれる。

始点0は、二つの不定元を持つ体$\mathbb Q(u,v)$上で計算する。
復元した六関係を挿入した行列の直交核から商写像と自由単項式sectionを取ると、
$C_{2,5}$は20次元、$C_{1,5}$は35次元であり、
\[
 \operatorname{rk}_{\mathbb Q(u,v)}d_2=77,\qquad
 \operatorname{rk}_{\mathbb Q(u,v)}d_3=36,\qquad d_2d_3=0
\]
を得る。これは有限個の有理数標本からの推論ではない。
二変数多項式の正確なGCDと割り算による有理関数消去の再現は
\path{computations/verify_generic_one_point_boundary.py}に置く。
各算術演算は分母を約分し、階数消去は非零な有理関数だけをpivotに使う。
従って一般点の$H_2$は$120-77-36=7$次元である。

さらに左乗法$m_{x_{1,3}}:C_{2,5}\to C_{1,5}$の一般階数は14で、その核は6次元である。
その核から平方型関係$x_{0,3}\otimes x_{1,3}$を用いて作ったcycle列を$Z$とすると、
$d_2Z=0$、$\operatorname{rk}(d_3\mid Z)=42$が同じ有理関数体上で成立する。
よって一般点でも六つの独立な平方型障害類があり、その商は一次元である。
一般の混合第七類の短い閉じた式は、ここでは与えない。

最後に補題\ref{lem:generic-quintic-exclusion}を適用すれば、全ての$U$上で
$\dim H_2\geq120-77-36=7$となる。
$X\setminus U$では有限Frobenius条件か四次検査を既に失っている。
従って例外的な係数値にも標準型Koszul AS点は存在しない。
全ての非空部分集合$I\subset\{0,1,2\}$について$\sum_{i\in I}c_i\ne0$とすれば
共通基点は$[0:0:0:1]$だけになるが、一様排除にはこの追加条件を必要としない。
\end{proof}

\section{戻り辺一回のバー複体と双加群Koszul障害}
\label{sec:trivial-extension-bitor}
trivial extensionの構造は、四窓の四次商を直接消去しないKoszul必要条件を与える。
以下のTorは次数付きTorであり、頂点の半単純環上で相対化する。

\begin{proposition}[平方零拡大の一回部分]\label{prop:one-return-bar-tor}
$A_0=S$が半単純である正次数付き代数$A$と、正次数に置かれた次数付き$A$双加群$M$を取り、
$T=A\ltimes M$とする。$T$の相対被約バー複体を$M$因子の個数で分けると、
その一回部分について
\[
 \operatorname{Tor}^{T}_n(S,S)_{\text{one }M,d}
 \simeq\operatorname{Tor}^{A^e}_{n-1}(S\otimes S^{\mathrm{op}},M)_d,
 \qquad A^e=A\otimes A^{\mathrm{op}}.
\]
従って$T$がKoszulなら、右辺は$d=n$以外で零である。
$A$もKoszulなら、右辺は左右のKoszul分解の全複体で計算できる。
\end{proposition}
\begin{proof}
バー微分で$A_+$因子同士を掛ける場合、または$A_+$因子を$M$へ作用させる場合、
$M$の個数は変わらない。隣接する二つの$M$の積は零である。
従って各個数部分は直和因子の複体である。一回部分の次数$n$の項は
\[
 \bigoplus_{p+q=n-1} A_+^{\otimes_Sp}\otimes_S M\otimes_S A_+^{\otimes_Sq}.
\]
これは$S$の右$A$バー分解と左$A$バー分解で$M$を挟んだ二重複体の全複体であり、
全ホモロジー次数が1だけずれている。左右のバー分解の外積は$A^e$上の
$S\otimes S^{\mathrm{op}}$の射影分解なので、主張のTorを計算する。
符号は一方の複体次数に応じた基底符号変更で通常の全微分と一致する。
Koszul代数のTorは内部次数とホモロジー次数が一致するため、必要条件を得る。
$A$がKoszulなら、両側のバー分解を線型Koszul分解に置き換えられる。
\end{proof}

前節の$A_L$について、頂点を位置$0,1,2,3$で表し、
$M=D A_L$には既述の次数$4-r$を入れる。
$Q_0=k$, $Q_1=V^*$,
$Q_2=\bigwedge^2(V^*)+L\subset(V^*)^{\otimes2}$と置く。
$\dim Q_2=10$である。位置1から位置2へ向かう内部次数5の一回部分の双加群Tor複体は
\[
 0\longrightarrow C_4\xrightarrow{\partial_4}C_3
 \xrightarrow{\partial_3}C_2\xrightarrow{\partial_2}C_1\longrightarrow0,
 \qquad
 C_h=\bigoplus_{\substack{p+q=h\\0\leq p,q\leq2}}
 Q_p\otimes B_{h-1}^*\otimes Q_q.
\]
その次元は順に
\[
 (\dim C_1,\dim C_2,\dim C_3,\dim C_4)=(8,144,480,400).
\]
左微分は$Q_p\to Q_{p-1}\otimes V^*$の最後の一次式を$B_{h-1}^*$へ縮約し、
右微分は$Q_q\to V^*\otimes Q_{q-1}$の最初の一次式を縮約する。
全微分を$\partial=\partial_{\mathrm L}+(-1)^p\partial_{\mathrm R}$とする。
これは$B$の可換乗法から作る複体であり、四次商空間の基底を必要としない。

\begin{proposition}[小複体による必要条件]\label{prop:one-return-small-obstruction}
上の$\partial_2$は全射である。
$w$が有限Frobenius候補で、その元の四窓がKoszulなら、
\[
 H_h(C_\bullet)=0\quad(h=1,2,3),\qquad \dim H_4(C_\bullet)=56.
\]
特に$\operatorname{rank}\partial_4<344$ならKoszulではない。
\end{proposition}
\begin{proof}
$Q_2$の交代テンソル$a\otimes b-b\otimes a$を、$b$で1、$a$で0となる一次双対元で縮約すると$a$を得る。
$(p,q)=(2,0)$と$(0,2)$の部分を使えば$C_1$の二つの四次元直和因子へそれぞれ全射となる。
従って$\partial_2$の階数は8であり、$\partial_3$の階数は136以下である。
前命題から$H_h(C_\bullet)$は$T_L$の$\operatorname{Tor}_{h+1,5}$の直和因子である。
$T_L$は復元双対$E(w)$に等しく、Koszul性は二次双対間で同値なので、
Koszulなら$h+1=5$の場合以外は零となる。
Euler標数$400-480+144-8=56$から最後の次元を得る。
さらに$H_3$の次元は$480-\rk\partial_3-\rk\partial_4$なので、最後の主張が従う。
\end{proof}

\begin{proposition}[既知の二変数境界における一回Tor]\label{prop:one-return-two-variable}
定理\ref{thm:generic-one-point-uniform-exclusion}の$\mathbb Q(u,v)$上の一般点では
\[
 (\rk\partial_2,\rk\partial_3,\rk\partial_4)=(8,136,337),\qquad
 (\dim H_1,\dim H_2,\dim H_3,\dim H_4)=(0,0,7,63).
\]
特に双対側の非線型$\operatorname{Tor}_{4,5}$に七次元の直和因子がある。
\end{proposition}
\begin{proof}
$Q_2$の基底を六つの交換子と四つの二次式とし、$B_2^*$は六つの自由二次座標、
$B_3^*$は既存の正規化三次フレームで表示する。
縮約微分の全構成は\path{computations/trivial_extension_bitor.py}、
厳密有理関数体上の零合成と三階数の検算は
\path{computations/verify_trivial_extension_bitor.py --field two}に置く。
既知の有理数点$u=v=1$でも同じ階数を独立に計算した。
完全交叉の対照点$t=1$でも同じコードの\texttt{--field ci}により、
三階数$(8,136,344)$とホモロジー$(0,0,0,56)$を確認した。
八変数チャート全体の階数はこの検算からは結論しない。
\end{proof}

\begin{proposition}[一般点の一回Torから基点枝全体への移行]
\label{prop:one-return-universal-transfer}
定理\ref{thm:product-net-dominance}の八変数チャートで、一般点の$\partial_4$階数を$r$とする。
$r<344$が証明されれば、基点を持つ全ての三因子対称有限Frobenius候補で
$\dim\operatorname{Tor}^{E(w)}_{4,5}\geq344-r>0$となり、Koszul AS候補は存在しない。
この移行自体は四次Koszul検査の正確性を必要としない。
八変数の階数$r$は未確定であり、この排除はここでは条件付きである。
\end{proposition}
\begin{proof}
基点を固定した$Y_3$とその三次フレーム束を
命題\ref{prop:product-net-universal-obstruction-transfer}の通り取る。
$B_1,B_2,B_3$の次元は一定で、$Q_2=\bigwedge^2(V^*)+L$も階数一定の束である。
従って$C_h$とその微分は全フレーム束上で定義されたベクトル束複体となる。
八変数チャートと基底変更による射は支配的なので、その一般階数$r$は
全フレーム束の一般階数であり、他のどの点の階数も$r$を超えない。
また$\partial_2$は全射なので、全ての点で$\rk\partial_3\leq136$である。
従って全点で$\dim H_3\geq480-136-r$となる。
有限Frobenius部分では$E(w)=T_L$であり、前命題の直和因子としてこの$H_3$を持つ。
基点の移動は共通基底変更で処理できる。
三因子対称有限Frobenius候補がこの$Y_3$表示に入ることは、
命題\ref{prop:two-positive-global-split}の三重交差の議論による。
\end{proof}

\section{一般四窓の切断と七次までの共通Tor検査}
\label{sec:general-cut-trivial-extension}
三因子対称性を使わずにも、戻り辺の切断は有限Frobenius候補全体に適用できる。
ここでは有限次数での必要条件を整理し、全Koszul性の有限判定とは区別する。

\begin{proposition}[任意の一辺切断によるtrivial extension]
\label{prop:general-cut-trivial-extension}
標準Frobenius四窓を折り畳んだ代数$E(w)$で、次数一の四辺のうち一辺を選ぶ。
その辺を通らない経路の部分代数を$A$とすると、$A$は32次元の四層二次Koszul代数で、
大域次元は3以下である。また
\[
 E(w)\simeq A\ltimes DA
\]
が成り立つ。$A$の長さ$r$成分の双対には次数$4-r$を入れる。
四つの切断のいずれでもこの表示を取れる。
これは一般対象に可換性・domain・noetherian性を加える主張ではない。
\end{proposition}
\begin{proof}
切断後の頂点を$0,1,2,3$の順に置く。
$A$の次数別次元は$4,12,12,4$であり、全次元は32である。
端点と長さを固定した経路は切断辺を通る回数が一定なので、
二次関係を課しても「切断辺を通らない部分」と「一回通る部分」の直和は保たれる。
前者は四層の二次表示を持つ部分代数であり、補題\ref{lem:four-layer-koszul}を適用できる。
後者を$M$と書くと$E=A\oplus M$で、$M$は双加群イデアルである。
切断辺を二回通る経路の長さは5以上であり、$E$の五次以上は零なので$M^2=0$である。

標準traceの巡回恒等式から、四頂点のtraceの和$\lambda$は対称なFrobenius形式を与える。
$\lambda(AA)=0$、$\lambda(MM)=0$なので、完全対合は$A$と$M$を完全に対合させる。
写像$m\mapsto(a\mapsto\lambda(am))$は$M\simeq DA$を与え、
traceの巡回性から左右の$A$作用を保つ。次数は$\lambda$の次数4から決まる。
従って平方零拡大の積は通常のtrivial extensionの積そのものである。
\end{proof}

切断後の位置$i,j\in\{0,1,2,3\}$を固定し、$d=4+j-i$と置く。
左右のKoszul分解を使って命題\ref{prop:one-return-bar-tor}の一回部分を計算する。
$h$番目の項$K_h^{i,j}$では、左右のKoszul次数を$p,q$として
\[
 p+q=h,\qquad 0\leq p\leq3-i,\quad0\leq q\leq j,
 \qquad r=i-j+h
\]
となり、中央係数は長さ$r$の$A$成分の双対である。
$h_r=(1,4,6,4)_r$とし、範囲外では$h_r=0$と定めれば
\[
 \dim K_h^{i,j}
 =h_{i-j+h}\!!\sum_{\substack{p+q=h\\p\leq3-i,\ q\leq j}}
       \binom{p+3}{3}\binom{q+3}{3}.
\]
ここでKoszul空間の次元$1,4,10,20$は、$A$のHilbert行列の逆行列から従う。
異なる窓の空間を同一視することは仮定していない。

\begin{proposition}[七次切断の必要十分な共通検査]
\label{prop:cut-tor-through-seven}
四つの切断それぞれで、全16組$(i,j)$の上の複体が最上次数$d-1$より下で正確であることは、
$E(w)$の内部次数$1\leq d\leq7$で$\operatorname{Tor}_{n,d}^{E(w)}=0$ ($n\ne d$)となることと同値である。
正確な場合の最上ホモロジー次元は$\binom{d+3}{3}$である。
特に$d=7$の一回部分は、零でない項の次元が
\[
 (\dim K_3,\dim K_4,\dim K_5,\dim K_6)=(120,1040,2400,1600)
\]
となる小複体で検査できる。
この有限検査だけから全Koszul性やAS正則性は結論しない。
\end{proposition}
\begin{proof}
一回部分の内部次数は$4+j-i$、最上Koszul次数は$3+j-i=d-1$である。
そのホモロジーは$E$の$\operatorname{Tor}_{h+1,d}$の対応する直和因子となる。
長さ7以下の有向四周期経路では、少なくとも一つの辺をちょうど一回通る。
その辺で切れば、端点と内部次数を固定した被約バー複体全体が戻り辺一回部分になる。
従って四切断の検査は内部次数7以下の全Tor成分を覆い、述べた同値を得る。
最上次数の符号を正に取った上の次元公式のEuler標数は
$\sum_h(-1)^{d-1-h}\dim K_h^{i,j}=\binom{d+3}{3}$である。
これは$0\leq i,j\leq3$の有限和で直接計算でき、$d=7$の表示も同じ公式による。

長さ8では全ての辺を二回通る経路があり、一回部分では覆えない。
また、より高次の非線型Torの消滅はここでは証明していないので、全Koszul性の十分条件とはしない。
\end{proof}

\path{computations/verify_cut_tor_dimensions.py}は全16組の項次元とEuler標数、
長さ7までの全始点の経路に一回辺があることを整数で検算する。
これは行列の階数や実際の複体の正確性を検算するものではない。

\paragraph{二つの有理点での全微分。}
\path{computations/general_cut_tor.py}では二次関係を挿入した行列の核から
各$Q_p$と中央の双対成分の基底を取る。左作用は左因子と中央因子の末尾の文字同士を縮約する。
右作用は両者の先頭を縮約し、符号$(-1)^p$を付ける。
窓の位置をこの二方向で別々に追うので、中央に可換性を仮定しない。
有理数上で四切断各16組の全ての連続微分の合成と階数を計算した。
符号付きOre例$H=0$, $G=(1,0,0)$では全始点、内部次数7以下で非対角成分は零であり、
特に内部次数7の三階数は小さい$h$から$120,920,1480$となる。
一方、節\ref{sec:trivial-extension-bitor}のBoolean境界$t=0$で非対角成分が出る位置は次の通りである。
表は$\dim\operatorname{Tor}^{E}_{d-1,d}$であり、四つの列は元の始点に対応する。
\[
\begin{array}{c|rrrr}
d&0&1&2&3\\ \hline
5&7&0&0&0\\
6&21&0&0&21\\
7&41&0&41&63
\end{array}
\]
この範囲で他の非対角成分は零、各始点の対角成分の次元は表の値に$\binom{d+3}{3}$を足した数になる。
同じ始点・次数を覆う切断は全て同じホモロジーを与えた。計算入力、再現手順と階数は
\path{computations/verify_general_cut_tor.py}および\path{research/general_cut_tor_certificate.json}に置く。

\section{連続二窓への還元と全次数のSerre判定}
\label{sec:cut-serre}
四窓のうち、連続する二つの関係空間を固定する。以下では、残りの二窓の復元と、
どの入力がAS正則な四窓へ延びるかを同じ切断から扱う。
判定には有限次元代数を使うが、調べるSerre関手の反復回数には上限を置かない。

\begin{theorem}[戻り回数別のTorとSerre反復]\label{thm:cut-serre}
$A$を頂点$0,1,2,3$、矢は位置を一つ増やす四層二次代数とし、
$E=A\ltimes DA$で$(A_r)^*$を内部次数$4-r$に置く。
$S_i$を右単純、$D S_j$を左単純とし、$D$は体上の双対とする。
$\mathcal D=D^b(\mathrm{mod}\text{-}A)$で$\nu_A=-\otimes_A^{\mathbf L}DA$を取る。
$m\geq0$, $d=4m+j-i$に対して
\[
 D\operatorname{Tor}^E_{n,d}(S_i,D S_j)
 \simeq\operatorname{Hom}_{\mathcal D}(\nu_A^m S_i,S_j[n-m]).
\]
各内部次数の非零成分はこの形で尽くされる。
$T=\bigoplus_{i=0}^3 S_i[i]$, $B=\operatorname{End}_{\mathcal D}(T)$と置く。
通常の合成を$B$の積とし、$B$でも右加群を用いる。次は同値である。
\begin{enumerate}
\item $E$が経路長次数についてKoszulである。
\item 全ての$m\geq0$で$(\nu_B[-3])^{-m}B$のコホモロジーが次数零に集中する。
\end{enumerate}
すなわち、切断の$B$が右加群の規約で3-representation infiniteであることが必要十分である。
\end{theorem}
\begin{proof}
被約バー複体の各項を$DA$因子の個数$m$で分ける。
$(DA)^2=0$なのでこれは複体の直和分解である。
$A$の各因子の内部次数は終点位置と始点位置の差であり、$DA$因子ではその差に4を加える。
従って端点$i,j$を固定すると内部次数は$4m+j-i$になり、$m$は一意である。
$m$個の$DA$の両端と間にある$m+1$個の$A_+$の列を、それぞれ正規化された相対両側バーと読む。
隣接する$A_+$同士の積と各列の端での作用がそのバー微分である。
$DA$を一因子として数えた次数$m$を除けば、これは
\[
 S_i\otimes_A^{\mathbf L}(DA)^{\otimes_A^{\mathbf L}m}
       \otimes_A^{\mathbf L}D S_j
\]
のモデルである（$m=0$の中央は$A$）。実際、各接合に
$A\otimes_{A_0}(A_+)^{\otimes_{A_0}p}\otimes_{A_0}A$を挿入する両側バー分解は、
$A_0$が半単純なので導来テンソル積を計算する。全複体の符号は、固定した$DA$因子を
一つずつホモロジー次数1に置く通常のテンソル積の符号である。
従ってバー次数$n$の双対ホモロジーは表示した導来Homとなる。
有限大域次元なので$\nu_A$はSerre自己同値である。

補題\ref{lem:four-layer-koszul}より$A$はKoszulである。全経路の長さが頂点位置差なので
$\operatorname{Ext}^p_A(S_i,S_j)$は$p=j-i$以外で零である。
従って$\operatorname{Hom}(T,T[k])=0$ ($k\ne0$)であり、単純加群のシフトは
$\mathcal D$を生成する。よって$T$はtilting対象で、導来tilting同値
$\mathbf R\operatorname{Hom}_A(T,-)$は$T$を右正則$B$加群へ、Serre関手をSerre関手へ送る。
$B$の非恒等射は位置を真に増やす方向にのみあり、対角成分は$k$なので
（右射影の順序は逆になるが）極小分解の長さは3以下である。
$T_i=S_i[i]$と置けば、先の式の右辺は
\[
 \operatorname{Hom}_{\mathcal D}\bigl((\nu_A[-3])^m T_i,
                                      T_j[n-d]\bigr)
 =\operatorname{Hom}_{\mathcal D}\bigl(T_i,(\nu_A[-3])^{-m}T_j[n-d]\bigr).
\]
固定した$m,i,j$について$n<m$の側の消滅も、上のバーによる導来テンソル積が
ホモロジー負次数を持たないことから従う。
全ての成分で$n\ne d$の消滅を要求することは、tilting同値を通して
$H^k((\nu_B[-3])^{-m}B)=0$ ($k\ne0$)と同値である。
これが3-representation infiniteの条件である\cite{HIOrep,GrantIyama}.
\end{proof}

\begin{corollary}[二窓からの延長と添字保持同型]\label{cor:two-window-extension}
四次元空間$V_0,V_1,V_2$と六次元部分空間
$R_0\subset V_0\otimes V_1$, $R_1\subset V_1\otimes V_2$を与え、
\[
 U=(R_0\otimes V_2)\cap(V_0\otimes R_1),\qquad \dim U=4
\]
とする。矢$V_0^*,V_1^*,V_2^*$、関係$R_0^\perp,R_1^\perp$の四層代数を$A$とする。
この二窓が標準14641型Koszul AS正則対象へ延びる必要十分条件は、
定理\ref{thm:cut-serre}の$B$が3-representation infiniteとなることである。
延長があれば添字保持同型を除き一意であり、二つの延長が添字保持同型となる必要十分条件は
三つの基底変更$(g_0,g_1,g_2)$が二つの関係空間を同時に対応させることである。
ここで始点の添字を固定し、巡回移動で商を取らない。
\end{corollary}
\begin{proof}
$A$の長さ3成分の双対は$U$である。条件$\dim U=4$より、
$A$の非零な端点別次元は長さ順に$1,4,6,4$である。
従って$E=A\ltimes DA$は対称Frobeniusで、各始点の次数別次元は$(1,4,6,4,1)$になる。
定理の条件が成立すると$E$はKoszul、特に次数一生成かつ二次である。
その二次双対を整数頂点列に展開すれば、最初の二窓は指定した$R_0,R_1$となる。
左右のKoszul単純分解の項は相補次数の完全trace対合で互いに双対になり、
双対複体は反対側の単純分解を4だけ移したものになる。従って大域次元4と
左右のAS--Gorenstein条件を得る。Hilbert行列のKoszul逆公式から
各始点の長さ$d$の次元は$\binom{d+3}{3}$であり、成長条件も満たす。
逆に標準AS対象の双対を選んだ辺で切れば命題\ref{prop:general-cut-trivial-extension}から
この$E$そのものであり、Koszul性より定理の条件を得る。

明示的には第四の一次空間を$V_3=U^*$とする。
標準テンソルは包含$U\hookrightarrow V_0\otimes V_1\otimes V_2$を表す
$w=\sum_a u_a\otimes u_a^*$であり、これは基底によらない。
実際、切断における三つの矢の積と$DA$の戻り矢とのtraceはこの評価である。
残る二つの関係を定理\ref{thm:frobenius-tensor-reconstruction}の縮約で復元できる。
三つの基底変更が指定の二窓を対応させれば$U$と$U^*$にも同型を誘導するので、
全四窓が同型になる。逆は添字保持同型の次数一、二への制限による。
\end{proof}

$B$は明示的にも小さい。Koszul分解から
$\operatorname{Ext}_A^p(S_i,S_{i+p})$は長さ$p$の二次双対の経路商であり、
Yoneda積はテンソルの連接を与える（分解のdeconcatenationの双対）。
従って$B^{\mathrm{op}}$は矢$V_0,V_1,V_2$と関係$R_0,R_1$を持つ四層代数である。
長さ別の総次元は$4,12,20,20$である。最後の20は
$64-24-24+\dim U$から得る。従って一般のAS問題を56次元代数上で扱える。
Grant--Iyama\cite[Theorems B, C]{GrantIyama}は$d$-hereditaryなどの仮定のもと、
高次前射影代数の正則性と二次双対からtrivial extensionへの写像を扱う。
上の証明は、候補が最初から3-hereditaryであることを要求せず、各非対角Torの位置を保つ。

% !TeX root = ../papers/nc-p3.tex
\section{二窓のExt微分と通常の経路空間による反復}
\label{sec:two-window-ext}
以下では前節の$B$自体の矢が$0\to1\to2\to3$となるよう番号を振る。
従って元の二窓からは、順序を逆にし各関係の二文字を反転して入力する。
$X_a,L_a$をこの向きの矢空間と関係空間とし、$S=B_0=k^4$、
$K_{a,p}$を長さ$p$のKoszul交叉とする。$p=0,1,2,3$で
それぞれ$k,X_a,L_a,U$（存在する端点のみ）であり、$U$はこの向きの三次交叉である。
$K_p=\bigoplus_a K_{a,p}$と置く。

\begin{proposition}[四項の計算と新しい辺の表示]\label{prop:two-window-ext-complex}
任意の右$B$加群$M$と頂点$i$について、$\operatorname{Ext}_B^p(DB,M)e_i$は、
次の$p=0,1,2,3$の複体で計算できる。
\[
 D_i^p(M)=\bigoplus_{0\le a\le\min(i,3-p)}
 (e_aBe_i)\otimes K_{a,p}^*\otimes M e_{a+p}.
\]
$T=\operatorname{Ext}_B^3(DB,B)$とすると、消滅を仮定せず任意の$M$について
$\operatorname{Ext}_B^3(DB,M)=M\otimes_B T$である。
$T_B(T)$は$B$に新しい辺$X_3=U^*$を加えた四周期の二次経路代数であり、
新しい辺をまたぐ関係は二つの包含
\[
 U\hookrightarrow L_0\otimes X_2,\qquad U\hookrightarrow X_0\otimes L_1
\]
の係数、すなわち$L_0^*\to X_2\otimes U^*$と
$L_1^*\to U^*\otimes X_0$の像である。
\end{proposition}
\begin{proof}
四層二次代数はKoszulである（補題\ref{lem:four-layer-koszul}）。
二側Koszul分解の項は$B\otimes_S K_p\otimes_S B$で、微分は先頭を左に掛ける項と
末尾を右に掛ける項の$(-1)^p$倍の和である。
この分解の正確性は右端の$B$の次数の降順フィルトレーションでも分かる。
随伴次数では末尾を掛ける微分が消え、一側Koszul複体の直和だけが残る。
増大写像を含めた分解は左$B$加群として分裂するため、左から$DB\otimes_B-$を施すと
右射影分解を得る。$e_i$を左から取って$\operatorname{Hom}_B(-,M)$を施し、
$e_i DB e_a=D(e_aBe_i)$を使うと表示式となる。
微分$D^p\to D^{p+1}$の第一項は$a$を一つ減らし、経路と$K^*$に同じ文字とその双対を
左から掛けて和を取る。第二項は$a$を保ち、$K^*$と$M$の右に対応する文字を掛け、
$(-1)^{p+1}$を付ける。従って写像には矢の積と関係の係数だけが必要である。

$M=B$の複体を全$i$で合わせると両側加群の複体となり、
その項は$B\otimes_S K_p^*\otimes_S B$である（双対は端点を反転する）。
任意の$M$の複体は左から$M\otimes_B-$を施したものなので、最終余核の式が従う。
最終項は端点$3\to0$の$U^*$が生成する自由な相対両側加群である。
長さ2の二項から来る微分ではそれぞれ片側にしか文字を付加できないので、
二つの係数写像の像が全関係を生成する。これをテンソル代数に入れれば主張の表示を得る。
\end{proof}

新辺の出現回数を$m$とすれば、$M_m^{(s)}=e_s T^{\otimes_B m}$の位置$i$成分は
上の四周期経路代数の始点$s$、長さ$4m+i-s$の経路空間そのものである。
負の長さは零とする。これは反復の途中で低次Extが出る場合にも\emph{通常の}テンソル積について成立する。
前の全段が純粋なときに初めて、この複体が次の導来反復の全コホモロジーを計算する。
従って両側射影複体を段ごとに作り直す必要はない。
先頭の左端を固定した経路を一文字ずつ延長し、最後の二次関係の像を割れば$M_m^{(s)}$の
次元と矢作用が得られる。最初の段が全端点で純粋なら、上の両側複体を次数3で終わるように
シフトしたものは$T$の左射影分解でもある。
このとき後続の低次Ext消滅は
$\operatorname{Tor}_{j}^{B}(T^{\otimes_B m},T)=0$ ($1\le j\le3$, $m\ge1$)と同値である。

\paragraph{有限の伝播条件を持つ枝。}
復元した四窓が有限Frobenius条件を満たすとする。
各矢に順序付き基底を選び、同じ長さの語を左辞書順で比較する。
位置$i$の六関係を掃き出して得る六つの先頭二文字の集合を$J_i$とする。
各$i$で、先頭二文字が$J_i$に入り末尾二文字が$J_{i+1}$に入る三文字が
ちょうど四つなら、四窓は全次数でKoszul AS正則である。
実際、次数3の正規語の数は$64-24-24+4=20$であり、
実際の経路商次元も三次交叉から20である。従って正規三文字は独立で、
全ての重なりの差は簡約して零になる。辞書順は固定端点・長さで停止し、左右への連接と両立する。
二次関係間の非自明な重なりは長さ3に尽きるため、ダイヤモンド補題から全正規語が基底となる。
二次PBW代数のKoszul性と有限Frobenius条件がAS性を与える。
逆にこの基底順で二次PBWなら、同じ次数3の次元等式より重なり数は四つである。

この枝では、任意の$s$と$m\ge0$について、各位置に長さ$4m+i-s$の正規語を置き、
矢を付けて簡約するという一つの規則で加群のクラスを記述できる。
$m=0$で全射影を含み、上のテンソル表示で$\operatorname{Ext}^3$により$m$が一つ増える。
PBWから全次数Koszul性が従うので、定理\ref{thm:cut-serre}と帰納法により、
これら全加群の$\operatorname{Ext}^{0,1,2}$は消える。異なる順序や他の枝の被覆は別の問題である。

\paragraph{二回戻る次数8の検算。}
\path{computations/two_window_ext.py}に四項の微分を有理数で実装した。
符号付きOre対照点では四つの始点の一段目が全て純粋で、
二段目の同じ端点でのExt次元は全て$(0,0,0,165)$である。
Boolean境界の始点1,2,3では一段目が純粋であるが、同じ二段目の
$\operatorname{Ext}^2$の次元は順に74,122,122である。
始点0では一段目の純粋性がすでに破れるため、同じ手順の二段目を導来反復の値とはしない。
\path{research/two_window_ext_certificate.json}に端点別の値を置く。

% !TeX root = ../papers/nc-p3.tex
\section{四周期二次PBW表示で可能な先頭語グラフの一意性}
\label{sec:periodic-pbw-graph}
各矢空間に順序付き四文字を取り、基底と順序は四位置ごとに繰り返す。
固定した始点・長さで左辞書順を使う。ここでのPBW表示は、窓の六関係を
この順序で掃き出したものが二次Gröbner基底になるという意味である。
\begin{theorem}\label{thm:periodic-pbw-graph}
このようなPBW表示を持つ標準14641型AS窓の二次先頭語代数は、
位置ごとの文字の付け替えを除き一種類である。
具体的には各位置の文字に $1,2,3,4$ を付けると、位置 $p$ の $i$ と
位置 $p+1$ の $j$ を連ねた二文字が正規であることと $i\le j$ が同値になる。
この番号は許容遷移のグラフから一意に回復できる。
これは文字の番号であり、最初の辞書順における文字の大小との一致を要しない。
\end{theorem}
\begin{proof}
頂点を（位置、文字）の十六組とし、正規な二文字に対応する有向辺だけを取る。
各位置から次の位置へ辺が十本ある。正規語はこのグラフの歩道に等しく、
各始点で長さ $n\ge1$ の正規語数は $\binom{n+3}{3}$ である。

閉路を含む強連結成分は単純な有向サイクルでなければならない。
そうでないとある頂点から成分内へ出る辺が少なくとも二本あり、
その二辺から始まり元の頂点へ戻る二つの閉歩道がある。
反復して同じ長さにすれば、これらの任意の連接が相異なる歩道を与え、指数増大に反する。
成分を縮約したグラフは非巡回である。一つの歩道が通るサイクル成分数の最大を $s$ とする。
通る成分列、成分間の辺、サイクルへの入口と出口（端点を含む）を固定する方法は有限である。
残る自由度は各サイクルを余分に何周するかだけで、指定した全長について最後の周回数は
他の周回数から高々一通りに決まる。従って $s\ge1$ のとき全歩道数は $O(n^{s-1})$ であり、
$s=0$ なら長い歩道はない。三次増大により一列に少なくとも四サイクルが必要である。

辺は必ず一位置だけ進むので各サイクルの長さは4の正の倍数である。
十六頂点しかないから、先ほどの列は長さ4の四サイクルで全頂点を使い切る。
この列の到達順に番号を付ける。各サイクルはどの位置にも一つの文字を持ち、
成分をまたぐ辺は小さい番号から大きい番号へのみ進む。
各位置で同番号の四辺以外に可能なのは六対だけである。
辺が十本あるから六対が全て現れる。成分間の到達順は全順序なので番号も一意である。
\end{proof}

\begin{corollary}
標準有限双対Frobenius条件を満たし、全始点の三次経路空間の次元が20の四窓を取る。
四周期の順序付き基底による二次PBW表示が存在することは、基底と各位置の番号付けを選んで、
六関係の先頭二文字をちょうど $i>j$ の六対にできることと同値である。
この枝の全入力は標準Koszul左右AS正則である。
\end{corollary}
\begin{proof}
二次PBWならKoszulであり、有限双対Frobenius条件と標準型からAS性と標準Hilbert級数を得る。
従って定理が必要条件を与える。逆に、番号が非減少な三文字は重複組合せにより
$\binom{6}{3}=20$ 個である。停止する二次簡約でこれらは三次商を張り、次元から独立である。
従って長さ3の全ての重なりが解消する。第\ref{sec:two-window-ext}節の
ダイヤモンド補題の議論で全次数のPBW性、従って正則性が従う。

基底の存在は有限の方程式と非消滅条件にできる。四つの基底変換の行列式を反転し、
各位置で辞書順の文字とサイクル番号の対応を四文字の順列から選ぶ。
関係の $6\times16$ 係数行列の列を辞書順の降順に並べると、最初の $k$ 列の階数は、
その $k$ 列に含まれる指定先頭対の数でなければならず、これで主元集合も決定する。
この有限個の小行列条件を基底変換について射影すればよい。
ここで決めたのは先頭語代数とこのPBW枝への所属である。
同じ先頭対を持つ窓の下位係数には区別すべき添字保持同型類が残り、
PBW表示を持たない窓も別途扱う必要がある。
\end{proof}

\begin{corollary}\label{cor:pbw-nonreduced-four-cycle}
標準14641型窓において四辺の双線型零点スキーム
$E_p\subset\mathbb P(V_p^*)\times\mathbb P(V_{p+1}^*)$ が全て有限であるとする。
四周期の左辞書順で二次PBW表示があれば、四つの点 $[u_p]\in\mathbb P(V_p^*)$
が存在し、$(u_p,u_{p+1})$ は各 $E_p$ の\emph{非被約点}である（$p$ は法4）。
特に有限四辺の少なくとも一辺が被約なら、そのようなPBW表示は存在しない。
\end{corollary}
\begin{proof}
$R_p\subset V_p\otimes V_{p+1}$ を六関係とし、$u_p$ を辞書順で最大の文字の双対とする。
$u_p$ における第一因子の縮約 $L:R_p\to V_{p+1}$ の階数は、最大文字から始まる
先頭対の数である。実際、降順に並べた係数行列の最初の四列の階数に等しい。
この文字のサイクル番号が $i$ なら定理\ref{thm:periodic-pbw-graph}より階数は $i-1$。
一方、階数が2以下なら $u_p$ の上の零点対応は少なくとも射影直線を含むので
有限性に反する。従って四位置の最大文字の番号は全て4で、$L$ の階数は3である。
両位置の最大文字を並べた対は正規で、しかも係数行列の先頭列に当たる。
この列は零でなければならず、$(u_p,u_{p+1})\in E_p$ を得る。

右縮約を $D:R_p\to V_p$, $r\mapsto(1\otimes u_{p+1})(r)$ と書く。
$\ker L$ の元は最大文字から始まる全係数が零である。このとき次に大きい列は
「二番目の文字、次位置の最大文字」の対である。この対も第二の番号が4なので正規であり、
$\ker L$ のいずれの行でもこの列は零となる。さもなければその行自体が正規対を先頭語に持つ。
従って $D(\ker L)$ は残る二文字の張る空間に入り、
$\dim\ker L=3$ から $\ker L\cap\ker D\ne0$ を得る。

この共通核の非零関係は、点 $(u_p,u_{p+1})$ で両因子方向への微分がともに零である。
六本の定義方程式のヤコビ行列の階数は高々5であり、滑らかな六次元の積空間内で
$E_p$ の接空間は少なくとも一次元となる。有限スキームの幾何点が被約ならその局所環は体で
接空間は零である。従って得られた四点は全て非被約である。
\end{proof}

例えば第\ref{sec:finite-point-chart}節の $v_i^2-v_j^2$ による対称AS窓は
被約な二十点を持つため、このPBW枝の外にある。
\path{computations/verify_finite_square_chart.py} は二十組の評価から六関係を回復し、
各関係行列 $R$ の微分 $((R y)^{\mathsf t},x^{\mathsf t} R)$ を用いて全二十点で階数6を直接確認する。

\begin{corollary}\label{cor:pbw-tangent-flags}
有限四辺の場合、候補四周期列を一つ固定し、各辺の点における接空間次元を $t_p$ とする。
辞書順で最初の二つの双対文字が張る旗の二平面の候補は、各位置で
$\mathbb P^{t_p-1}$ に絞られる。特に全て $t_p=1$ なら、有限個の候補点列ごとに
高々 $(\mathbb P^1)^4$ 上で有限個の行列階数条件を調べれば、この表示の存在が決まる。
\end{corollary}
\begin{proof}
前系の $u=u_p$, $v=u_{p+1}$, $L$, $D$ を使う。積の射影接空間での接方程式は
接ベクトルを $\dot u\bmod ku$, $\dot v\bmod kv$ で表して
\[
             \dot u(D r)+\dot v(L r)=0\qquad(r\in R_p)
\]
である。$L$ の像は三次元の $\ker v$ なので、$\dot v\bmod kv$ は $\dot u$ から
一意に決まり、その存在条件は $\dot u(D\ker L)=0$ である。従って接空間の第一射影は
$(D\ker L)^\perp/ku$ に同型である。前系の第二列に関する議論から、
二番目の双対文字はまさにこの消滅条件を満たす必要がある。
$u$ を含む二平面の候補はこの商の直線なので主張の射影空間を得る。

$t_p=1$ ならこの二平面は一意である。残る三平面はその二平面を含むもの全体、
すなわち二次元商の直線からなる。旗を固定すれば具体的な適合基底の選び直しは不要である。
降順に並べた双対基底の変更は各文字にその前の文字の線型結合を加える操作と
非零スカラー倍で尽きる。次位置とテンソル積を取った辞書順の列でも
各先頭区間の張る空間が保たれるので、先の小行列階数は旗だけで決まる。
最大文字の番号は4に固定し、残る番号の対応を各位置で六通りから選べば十分である。
有限辺には有限個の幾何点しかないから、四周期列も有限集合である。
この判定には非被約点の接線を持つことの十分性は仮定しない。
\end{proof}

特に $t_p=t_{p+1}=1$ で、隣接する二位置の第二双対文字の候補を
$u'_p,u'_{p+1}$ と書く。第三の旗を探索する前に
\[
 W_p=(u'_p\otimes1)(\ker L)\subset V_{p+1},\qquad d_p=\dim W_p
\]
を求められる。$u'_p$ に $u_p$ の倍を加えても $W_p$ は変わらない。
$W_p$ の主元が「二番目の文字」で始まる先頭対を与えるので、
その文字の番号は $d_p+1$ に強制される。従って $d_p>2$ は直ちに排除できる。
さらに $W_p\subset\ker u_{p+1}$ であり、次の双対文字が $W_p$ 上で非零となることは、
次位置の第二文字を右に付けた列が主元となることと同値である。
次位置でも $d_{p+1}$ が定まるとき、
\[
         u'_{p+1}(W_p)\ne0 \quad\Longleftrightarrow\quad d_p>d_{p+1}
\]
が必要である。これは各辺の接方向と二次関係だけから四周期列を除外する検査となる。

\begin{corollary}\label{cor:pbw-no-double-cycle}
有限四辺の標準窓が周期的な左辞書順二次PBW表示を持つなら、
上の四周期列に沿う局所環の長さは少なくとも3である。
従って全ての局所因子が長さ1または2である有限辺を一つでも持つ場合もこの表示は存在しない
（他の三辺も有限とする）。
\end{corollary}
\begin{proof}
系\ref{cor:finite-edge-point-sequences}により、共有文字への隣接射影は同一の部分スキームへの同型である。
従って選んだ列に沿う四つの局所環は同型である。長さ1はすでに排除した。
長さ2なら各点は曲線状二重点であり、両射影は接線を同型に運ぶ。
辺 $p$ の接方向の両端には、非零スカラー倍を調整して
上の第二双対文字 $u'_p,u'_{p+1}$ を使える。

この辺の六関係の中で一次微分が消えるものは一次元であり、その生成元を $r_0$ とする。
$r_0(u'_p,u'_{p+1})$ は非零である。
実際、接方向を $\epsilon^3=0$ まで延ばすための方程式は、二次の補正を $b,b'$ として
\[
         b(D r)+b'(L r)=-r(u'_p,u'_{p+1})\qquad(r\in R_p)
\]
となる。左辺の像は $R_p^*$ の余次元1の空間であり、$r_0$ の消滅空間に一致する。
右辺が $r_0$ で零なら補正が存在する。しかし局所環 $k[t]/t^2$ から
$k[\epsilon]/\epsilon^3$ へ、$t$ の像の一次係数が非零となる写像は存在しない。
これは接方向の非零性に反する。$r_0\in\ker L$ なので $u'_{p+1}(W_p)\ne0$ である。
上の主元条件から四位置の全てで $d_p>d_{p+1}$ となり、周期性に反する。
\end{proof}

\paragraph{二重点を持つ正則窓への適用。}
標数零の代数閉体上、対称チャートで四辺の各々に被約点16個と長さ2の点2個を持つAS窓が存在する。
実際、四変数二次式の対称行列の空間 $\mathbb P^9$ で次の直線を取る：
\[
 M(s,t)=\begin{pmatrix}s I_2&t I_2\\ t I_2&0\end{pmatrix}.
\]
$t\ne0$ では可逆で、$P=[M(1,0)]$ だけが階数2である。
近傍で左上ブロックを $A$、右上を $B$、右下を $D$ とすれば、
階数2以下の軌跡の三つの局所式は $D-B^{\mathsf t}A^{-1}B=0$ である。
直線を含む一般の射影三平面を取り、追加する二方向を $N_1,N_2$ とする。
それらの右下対称ブロックは独立で、張る平面は $I_2$ を含まないようにできる。
局所変数を $t,u_1,u_2$ とすると三式の線型項は追加二方向だけに階数2を持つ。
その二成分を解けば $u_i=O(t^2)$ となり、残る成分の二次係数は
$-I_2$ の法線の商における像で非零である。従って $P$ の局所環の長さは2である。

直線を含む超平面系は直線外で基点自由である。その一般の六元で切ると、
Bertiniの定理により階数2の滑らかな六次元軌跡では $P$ 以外の交点は横断的であり、
直線と交わらない三次元の階数1軌跡は避ける。
また直線に属する二次式二本は共通因子を持たず、$\mathbb P^3$ での共通零点は曲線である。
追加の一般二次式二本はその曲線上に共通零点を持たない。
以上の非空開条件を同時に取れば、四式は定理\ref{thm:symmetric-frobenius-chart}の正則窓を与える。

辺スキームは二つの一次式の順序付き積がこの三平面に入る条件である。
積の射は階数2の上で次数2のエタール被覆である。実際、二因子は相異なり、
交換による商から階数2軌跡への写像は一意分解により有限双有理で、像は滑らかなので同型となる。
ゆえに $P$ の上に二つの二重点があり、他の全点は被約である。六つの双一次式の
有限完全交叉の長さは $\binom63=20$ なので被約点は16個である。
この例に系\ref{cor:pbw-no-double-cycle}を適用すると、辺が非被約でも
四周期の左辞書順二次PBW表示が存在しないことが分かる。

上の二方向を有理数で選んだ具体例は
\path{computations/verify_tangent_pencil_window.py}にある。四二次式の基点の不在を確かめ、
三次小行列式で切ったアフィン代数の辞書順基底から、分離座標の十次式とその導関数の最大公約元が
二重点を表す一次式となることを有理数上で検算する。無限遠に別の交点はない。

\begin{proposition}\label{prop:pbw-symmetric-square-flag}
標数零の代数閉体上で対称完全交叉チャートの四次元二次式系を $L$ とする。
全位置で共通の順序基底 $e_1<e_2<e_3<e_4$ を用い、番号もこの順序に合わせて
先頭対を $e_j e_i$ ($j>i$) にできるための必要十分条件は、次の二段の平方条件である。
$L$ がある非零な平方 $u_4^2$ を含み、$u_4=0$ に制限した三次元二次式系も非零な平方を含む。
\end{proposition}
\begin{proof}
$u_i$ を $e_i$ の双対変数とする。六関係は対称である。
指定した先頭対で被約行階段形を取ると、一行に他の下降対は現れず、
対称性により非対角項は指定対とその反対の和だけとなる。
加えられる平方は辞書順から $e_k^2$ ($k<j$) に限る。
従ってその直交する対称二次式系は、各 $k$ で
$u_k^2$ に $u_i u_j$ ($i<j$, $j>k$) の線型結合を加えた四式を基底とする。
逆にこの形なら、各下降対の対称和に下位の平方だけを加えて直交関係を作れる。
三角形の先頭語と三次元数20から前の系を適用できる。

この基底の形からまず $u_4^2$、次に制限後の $u_3^2$ を得る。
逆に、基点のない四式系に平方があれば、対応する超平面上への制限は基点のない三式系である。
実際、次元は高々3で、射影平面上の高々二つの二次式は共通零点を持つ。
次の平方を選んで同様に制限すると、二変数の基点のない鉛筆が残る。
その二元行列束の行列式は二次式で、代数閉体上では特異な非零要素を持ち、これも平方である。
最後の一変数に残るのは非零な平方だけである。逆順に持ち上げると
$q_k\equiv u_k^2\bmod(u_{k+1},\ldots,u_4)$、さらに後の基底式で高い番号の平方項を消せる。
これが求める形である。任意の位置別基底を許す場合にはこの必要性を主張しない。
\end{proof}

この表示と辺の有限性は両立する。同じ検算コードの
\texttt{triangular\_window} は上記三角形の有理四式を与える。
最後の式は $w^2$ である。先頭式の係数を1にしたチャートで
三次小行列式の辞書順基底の先頭項は $b,c,d^5$ となり、
最後の五次式は平方自由である。従って階数2以下の切断はこのチャートで被約五点である。
階数1なら全ての三次小行列式の微分が消えるから、この五点は全て階数2である。
残りのチャートのうち次またはその次の係数が先頭非零となるものは単位イデアルを与え、
残る支持は $w^2$ だけである。順序付きの辺には被約点十個と $(w,w)$ の一点があり、
有限完全交叉の長さから後者の局所長は10である。

\section{Boolean完全交叉上の全一基点二次式のフレーム}
\label{sec:boolean-quadric-boundary}
前節では最後の二次式を$e_3$で割り切れるものに限った。
ここでは交差項も許し、最初の三式を固定した場合の一基点二次式を全て扱うフレームを与える。
これ自体からAS正則性や非Koszul性は結論しない。

\begin{proposition}[全一基点開集合での三次評価フレーム]
\label{prop:boolean-global-cubic-frame}
$Q=(e_0^2-e_0e_3,e_1^2-e_1e_3,e_2^2-e_2e_3)$とし、
$\epsilon\in\{0,1\}^3$に対応する列ベクトルを$p_\epsilon=(\epsilon_0,\epsilon_1,\epsilon_2,1)$とする。
二次式$q$が$p_{000}$でのみ零となり、他の七点で非零であるとする。
$Q$の線型結合を加え、非零スカラー倍で正規化すれば、$q$は一意に
\[
 q=e_0e_3+u e_1e_3+v e_2e_3
 +(a-1-u)e_0e_1+(b-1-v)e_0e_2+(c-u-v)e_1e_2
\]
と書ける。係数域は
\[
 X=\{uvabc(a+b+c-1-u-v)\ne0\}\subset\mathbb A^5
\]
であり、七点での値は順に$1,u,v,a,b,c,a+b+c-1-u-v$である。
$B=\operatorname{Sym}(V^*)/(Q,q)$に対し、
\[
 b_i=\sum_{\epsilon_i=1}
       \frac{(-1)^{|\epsilon|-1}}{q(p_\epsilon)}p_\epsilon^{\otimes3}
       \quad(i=0,1,2),\qquad b_3=p_{000}^{\otimes3}
\]
は$X$の全域で正則な$B_3^*$の基底を与える。
さらに
\[
 (\dim B_n)_{n\geq0}=(1,4,6,4,1,1,1,\ldots)
\]
が全ての$X$上で成立する。
\end{proposition}
\begin{proof}
$Q$を使って$e_i^2$を$e_ie_3$へ置換すると、二次式の標準形は
$e_3^2,e_ie_3,e_ie_j$ $(i<j<3)$の七項である。
$p_{000}$での零条件が$e_3^2$係数を消し、$q(p_{100})\ne0$が$e_0e_3$係数の正規化を許す。
残りは上記の五係数と七点での評価から一意に決まる。

$R=\operatorname{Sym}(V^*)/Q$の次数3以上の成分は、八点の評価によって$k^8$と同型になる。
実際、次数3の多項式
$\prod_{i:\epsilon_i=1}e_i\prod_{i:\epsilon_i=0}(e_3-e_i)$
が各点のindicatorを与える。
従って八つの三次評価テンソル$p_\epsilon^{\otimes3}$は独立である。
$R_1$の七非零点への評価は階数4なので、$qR_1\subset R_3$も4次元である。
よって$\dim B_3=8-4=4$となる。

各$b_i$は$Q$を消す。$q$と任意の一次式との積の評価では分母が消え、
$\epsilon_i=1$の面の四頂点の符号付き和が零になるので、$qR_1$も消す。
$b_3$も$q(p_{000})=0$からこれを消す。
四つの評価重み列を$p_{100},p_{010},p_{001},p_{000}$の行で見ると、
行列は$\operatorname{diag}(1,u^{-1},v^{-1},1)$となる。
従って四テンソルは全ての$X$上で独立であり、基底をなす。

最後に$R_2$は7次元で、七非零点への評価は同型である。
もし次数2の元がその七点で消えるなら、その八点上の関数は$p_{000}$のindicatorの定数倍である。
しかしその一意なmultilinear表示は次数3の項を持つので、次数2では非零な倍を表せない。
従って$\dim B_4=8-7=1$となる。
$n\geq5$では$R_{n-2}$の全八点評価が同型であり、$q$の乗法の階数は7なので、$\dim B_n=1$である。
低次数の値は四つの独立な二次式から従う。
\end{proof}

\begin{corollary}[五係数域全体での対称成分の閉包への所属]
\label{cor:boolean-family-in-symmetric-closure}
前命題の基底から$(w_q)_{abjc}=(b_j)_{abc}$と定めたテンソルは、全て$Z_{\mathrm{sym}}$に属する。
\end{corollary}
\begin{proof}
$q_t=q+t e_3^2$とする。
$i=0,1,2$の基底式では分母を$q(p_\epsilon)+t$へ置き換える。
第3基底は
\[
 b_3(t)=p_{000}^{\otimes3}
 -\frac{t}{u+t}p_{010}^{\otimes3}
 -\frac{t}{v+t}p_{001}^{\otimes3}
 +\frac{t}{c+t}p_{011}^{\otimes3}
\]
とする。これは$\epsilon_0=0$の面の符号付き評価関係を$t$倍したものであり、
全て$q_tR_1$を消す。全ての分母は$t=0$で可逆であり、上の四基底へ特殊化する。
一般の$t$では$q_t$は全八点で非零なので、四二次式は基点自由完全交叉となる。
従って一般テンソルは対称チャートに属し、特殊値$w_q$はその飽和閉包に属する。
\end{proof}

\begin{remark}
このフレームは、$012,133,233,333$を指定する局所正規化の追加分母が零になっても定まる。
その局所正規化を一般点の計算に使う場合も、最後の特殊化議論は上の大域フレーム上で行う必要がある。
最初の三二次式を任意に動かした一般の一基点四二次式系まで被覆したとは主張しない。
\end{remark}

\begin{corollary}[Frobenius境界上の非対称化性]
\label{cor:boolean-frobenius-not-symmetric}
上の$X$で$w_q$が有限Frobenius条件を通るなら、$w_q\notin S$である。
従って、この係数域からKoszul AS点が得られれば、対称チャート外の低階数対象となる。
その存在をここで結論しているわけではない。
\end{corollary}
\begin{proof}
因子$0,1,3$は既に対称なので、$S_1,S_3$は$I$を含む。
対称化可能なら線型判定により$S_2$にも可逆元があり、それだけを因子2へ作用させれば完全対称になる。
この作用は因子$0,1$の縮約像$R_0$を変えない。
$R_0$は6次元であり、その直交二次関係は$\bigwedge^2(V^*)+(Q,q)$である。
従って得られた完全対称Frobenius対象の可換二次双対は$B$そのものとなる。
対称Frobeniusチャート定理はこれが基点自由完全交叉であることを要求するが、
$(Q,q)$には$p_{000}$が共通基点として残る。これは矛盾である。
\end{proof}

\begin{theorem}[固定Boolean完全交叉上の全一基点係数域のAS排除]
\label{thm:boolean-one-point-uniform-exclusion}
命題\ref{prop:boolean-global-cubic-frame}の全係数域$X$は、標準型Koszul AS点を一つも含まない。
有限Frobenius条件と全四次検査を通る任意の$X$上の点で、始点0の五次$H_2$は7次元以上である。
一般点では五次の$H_2$次元は$(7,0,0,0)$であり、$H_3$は全始点で零である。
この結論は最初の三二次式を固定した場合の全域を覆うが、任意の相異なる八点完全交叉への一様排除は主張しない。
\end{theorem}
\begin{proof}
各面の基底$b_i$を、その面上の四つの$q$評価値の積でスカラー倍すると、
$X$上の単元による基底変更で分母のない多項式テンソルを得る。
これを一般点の計算に使う。三因子対称性と四二次式の消去から、全二対二階数は6以下である。
$(u,v,a,b,c)=(1,1,2,2,2)$は前節の検算済み標本の三次基底変更であり、
その変更行列の行列式は$1/12$である。
従って有限Frobenius条件と全四次検査を通る部分$U\subset X$は空でない開部分である。

体$\mathbb Q(u,v,a,b,c)$上で商写像と自由単項式sectionを作り、
始点0の五次微分$d_2$を構成すると、正確な階数は77となる。
行列の大きさは$140\times120$である。
再現コード\path{computations/verify_boolean_quadric_homology.py}は、
浮動小数や有限個の評価点を用いず、五変数の有理関数の消去でこの値を検証する。
計算量を抑えるためpivotの順序を多項式の項数等で選ぶが、全ての操作は体上の可逆な消去である。

$U$上では$\operatorname{rk}d_3\leq36$であり、検算済みの点で36を達成する。
従って一般点での$d_3$階数は36であり、一般$H_2$次元は$120-77-36=7$である。
始点1,2,3の五次正確性と全始点の$H_3=0$は、同じ標本から開性によって一般点へ移る。
補題\ref{lem:generic-quintic-exclusion}を$X$と$U$へ適用すると、
全ての$U$上で$\dim H_2\geq7$となる。
$X\setminus U$では既に必要なFrobenius条件か四次検査を失っているので、そこにもAS点はない。
大域評価フレームを使っているため、途中の局所正規化の追加分母が消える値もこの議論に含まれる。
\end{proof}

\section{相異なる八点の完全交叉から作る一基点境界}
\label{sec:reduced-quadric-boundary}
Boolean配置を固定しないため、三二次式の完全交叉そのものを点の評価で記述する。
この節ではKoszul障害の普遍性を仮定せず、フレームとその適用範囲を証明する。

\begin{lemma}[七点への二次評価と非零重み]\label{lem:reduced-octad-seven-evaluation}
三二次式の完全交叉$Z\subset\mathbb P^3$が八つの相異なる点$p_0,\ldots,p_7$からなるとする。
全点で非零な一次式$h$を取り、各代表を$h(p_i)=1$と正規化する。
次数2の評価像$F_2\subset k^8$は7次元であり、任意の七点への制限は同型である。
従って$F_2$を消す一次元の汎関数の重み$\lambda_0,\ldots,\lambda_7$は全て非零であり、
\[
 P\operatorname{diag}(\lambda_0,\ldots,\lambda_7)P^{\mathsf T}=0,
 \qquad P=(p_0\mid\cdots\mid p_7)
\]
が成立する。
\end{lemma}
\begin{proof}
$R=k[e_0,e_1,e_2,e_3]/(f_1,f_2,f_3)$とする。
完全交叉のHilbert関数は$(1,4,7,8,8,\ldots)$であり、$h$は$R$の非零因子である。
次数$n$の評価像$F_n$は$R_n$と同型で、$h=1$の正規化により$F_n\subset F_{n+1}$となる。
$F_3=k^8$であり、このフィルトレーションの随伴次数代数は$R/(h)$と同型である。
実際、$F_n/F_{n-1}$の核は$R_n$内で$hR_{n-1}$となる。
$R/(h)$は三変数の三二次式Artin完全交叉なので、Hilbert列は$(1,3,3,1)$、socleは次数3だけにある。

もしある一点のindicator $\delta_i$が$F_2$に入れば、最小の$d\leq2$について
$\delta_i\in F_d\setminus F_{d-1}$となる。
全ての$f\in F_1$に対して$\delta_if=f(p_i)\delta_i\in F_d$なので、
$\delta_i$の初期類は随伴次数代数の正次数元で消される。
これは次数$d\leq2$に非零socleを与え、矛盾する。
従って$\delta_i\notin F_2$である。
任意の七点へ制限した評価の核は、欠けた点のindicatorの倍だけなので零であり、次元7から同型となる。
$F_2$の唯一の消去汎関数が一点で重み零なら、そのindicatorが$F_2$に入るため、全重みは非零である。
最後の行列等式は全ての二次単項式への評価をまとめたものである。
\end{proof}

\begin{proposition}[一般の相異なる八点に対する一基点フレーム]
\label{prop:reduced-octad-one-point-frame}
前補題の完全交叉に二次式$q$を加え、$q(p_0)=0$、$q(p_s)\ne0$ $(s=1,\ldots,7)$とする。
$B=R/(q)$と置くと
\[
 (\dim B_n)_{n\geq0}=(1,4,6,4,1,1,1,\ldots).
\]
$P'=(p_1\mid\cdots\mid p_7)$の核の任意の三基底$k_0,k_1,k_2\in k^7$に対し、
\[
 b_j=\sum_{s=1}^7\frac{(k_j)_s}{q(p_s)}p_s^{\otimes3}\quad(j=0,1,2),
 \qquad b_3=p_0^{\otimes3}
\]
は$B_3^*$の基底である。
この基底から三因子対称テンソル$w$を作れば、$w\in Z_{\mathrm{sym}}$である。
さらに$w$が有限Frobenius条件を通るときは$w\notin S$である。
\end{proposition}
\begin{proof}
七点への一次評価も単射なので$\operatorname{rk}P'=4$であり、核は3次元である。
八点の三次評価テンソルは独立である。
上の三つの重みベクトルは$q$の値を掛けた後に$P'$の核になるので、$qR_1$を消す。
基点のcubeもこれを消す。
前三ベクトルは七非基点の評価上で独立であり、最後のcubeは基点成分だけを持つので、四つは独立である。
$q:R_1\to R_3$は単射であり、$\dim B_3=8-4=4$からこれらは基底をなす。
前補題により$q:R_2\to R_4$の階数は7なので$\dim B_4=1$である。
$n\geq5$では$R_{n-2}$の全八点評価が同型となり、$q$乗法の階数は7である。
低次数の値も合わせて所要のHilbert関数を得る。

閉包所属を示すため$q_t=q+t h^2$とし、前三基底では分母を$q(p_s)+t$へ変える。
$P'$は全射なので、ある$l_1,\ldots,l_7$について
$p_0+\sum_{s=1}^7l_sp_s=0$となる。
第3基底を
\[
 b_3(t)=p_0^{\otimes3}
       +t\sum_{s=1}^7\frac{l_s}{q(p_s)+t}p_s^{\otimes3}
\]
とすれば、全ての$q_tR_1$を消し、$t=0$で上の基底へ特殊化する。
一般の$t$では四二次式が基点自由完全交叉なので、対応するテンソルは$S$に入る。
よって$w\in Z_{\mathrm{sym}}$である。
有限Frobenius点での非対称化性は、因子2だけの可逆対称化が最初の二因子の関係空間を変えず、
基点を持つ$B$を完全対称Frobenius双対にしてしまうという、系\ref{cor:boolean-frobenius-not-symmetric}と同じ矛盾による。
\end{proof}

\begin{proposition}[八点の直交行列表示]\label{prop:reduced-octad-orthogonal-form}
上記の相異なる八点完全交叉は、点の順序、代表のスカラー、および共通射影基底を選べば
\[
 (I_4\mid M),\qquad MM^{\mathsf T}=I_4
\]
という列表示を持つ。
この表示は点配置の表示であり、indexed AS対象の完全な同型分類ではない。
また任意の直交行列から相異なる八点完全交叉が得られるとは主張しない。
\end{proposition}
\begin{proof}
四つの独立な点を最初の四列に選ぶ。
前補題の全て非零な重みについて、最初の四列を$\sqrt{\lambda_i}$倍、
後ろ四列を$\sqrt{-\lambda_i}$倍する。標数零の代数閉体なので選べる。
得た列行列を$(A\mid C)$とすれば$AA^{\mathsf T}-CC^{\mathsf T}=0$である。
$A$は可逆だから、左から$A^{-1}$を掛けると
$M=A^{-1}C$は$MM^{\mathsf T}=I_4$を満たす。
\end{proof}

\begin{remark}[古典的な自己付随八点との関係]
これはCayley octadの古典的なGale自己付随性に対応する表示である。
Plaumann--Sturmfels--Vinzant\cite[Theorem 3.10]{PSVQuartics}は
一般位置の八点について完全交叉と自己付随性の同値を述べている。
ここでは完全交叉から表示へ向かう方向を、評価フィルトレーションで証明した。
逆方向を退化した直交行列全体へ拡張せず、この古典的正規形そのものの新規性は主張しない。
\end{remark}

\begin{corollary}[最後の二次式の全係数域]\label{cor:octad-last-quadric-values}
前補題の八点を固定すれば、最初の三二次式を法とした第四二次式は
\[
 \{(z_0,\ldots,z_7)\in k^8:\ \sum_s\lambda_sz_s=0\}
\]
上の評価値と一意に対応する。
基点を$p_0$だけとする条件は$z_0=0$、$z_1\cdots z_7\ne0$であり、
非零スカラー倍を除けば五次元射影空間の開部分である。
これは二次式の係数域の次元であり、indexed同型類の次元を意味しない。
\end{corollary}
\begin{proof}
二次評価像は7次元の$\ker\lambda$に等しく、その核は最初の三二次式の空間である。
一点零条件はその中の一つの独立な線型条件なので、スカラー倍を除いた次元は5である。
残る七点での非消滅は開条件である。
\end{proof}

\begin{proposition}[残る対称化空間と四次双対]\label{prop:partial-symmetrizer-quartics}
$B=\operatorname{Sym}(V^*)/(L)$、$\dim B_3=4$とし、その三次双対基底から
因子$0,1,3$が対称なテンソル$w$を作る。このとき
\[
 S_2(w)\longrightarrow B_4^*,\qquad H\longmapsto H^{(2)}w
\]
は線型同型である。右辺は四階対称テンソルとして解釈する。
さらに$H$の階数は対応する四次テンソルの一因子縮約の階数に等しい。
\end{proposition}
\begin{proof}
$H^{(2)}w$は常に因子$0,1,3$について対称である。
$S_2$の方程式はさらに交換$(02)$についての対称性を要求し、合わせて完全対称性となる。
その三次縮約は全て$B_3^*$に属するので、完全対称な場合は$L\cdot\operatorname{Sym}^2(V^*)$を消し、
$B_4^*$に属する。逆に$F\in B_4^*$の各一次縮約は$B_3^*$に属する。
選んだ三次双対基底により各縮約を一意に表示すれば、$F=H^{(2)}w$となる行列$H$を一意に得る。
$w$の因子2から$B_3^*$への縮約写像は同型なので、最後の階数の等式も従う。
\end{proof}
\begin{corollary}[基点と階数1対称化写像]\label{cor:basepoint-rank-one-symmetrizer}
前命題の仮定の下で、$L$の共通零点$[p]\in\mathbb P(V)$と、
$S_2(w)$の非零階数1行列の射影類は一対一に対応する。
対応は$H^{(2)}w=p^{\otimes4}$で与えられる。
特に基点がないことは$S_2(w)$に非零階数1元がないことと同値である。
これは閉点の対応であり、非被約基点スキームの同定までは主張しない。
\end{corollary}
\begin{proof}
$[p]$が基点なら$p^{\otimes4}$は$L\cdot\operatorname{Sym}^2(V^*)$を消し、
$B_4^*$に属する。逆に非零対称四テンソルの一因子縮約の階数が1なら、
それを$p\otimes T$と書け、対称性から$T$も$p^{\otimes3}$のスカラー倍となる。
従ってそのテンソルは$c p^{\otimes4}$である。
これが$B_4^*$に属する条件は、全ての$q\in L$と二次式$r$について
$q(p)r(p)=0$となることであり、$p\ne0$なので全ての$q(p)=0$と同値である。
前命題の同型と階数の等式を適用すれば結論を得る。
\end{proof}
\begin{corollary}[三因子対称Frobenius候補の尽くされた二分岐]
\label{cor:partial-symmetric-exhaustive-split}
三因子対称な有限Frobenius候補を三次双対基底で表示する。
このとき次の二枝は互いに排反で、全候補を尽くす。
\begin{enumerate}
\item $S_2(w)$に非零階数1元がない。このとき$L$は基点を持たず、
四二次式完全交叉から来る完全対称化可能な枝に属する。
$S_2(w)$は可逆行列で生成される一次元空間である。
\item $S_2(w)$に非零階数1元がある。このとき$L$は基点を持ち、
$w$は完全対称化できない。
\end{enumerate}
第二枝のKoszul/AS正則性はこの二分岐だけからは決まらない。
\end{corollary}
\begin{proof}
前系により二条件は基点の有無そのものである。
基点がなければ、四変数多項式環の四二次式は高さ4の正則列となる。
その完全交叉の四次双対は一次元で、完全な乗法双対性により一因子縮約階数は4である。
三次双対基底をその四次生成元の縮約基底へ変えれば$w$は完全対称となる。
逆に基点を持つ有限Frobenius候補が完全対称化できれば、線型対称化判定から
$S_2$に可逆元があり、対応する四次テンソルが同じ$L$の完全交叉表示を与える。
これは基点に反する（定理\ref{thm:symmetric-frobenius-chart}）。
\end{proof}
\begin{proposition}[二つの対称化証明書を持つ全候補の分岐]
\label{prop:two-positive-global-split}
有限Frobenius候補$w$について、相異なる$j,\ell\in\{1,2,3\}$で
$S_j(w),S_\ell(w)$に可逆元があるとする。残る位置を$m$とすると、次のいずれか一方が成り立つ。
\begin{enumerate}
\item $w$は完全対称化可能であり、$S_m(w)$は可逆行列で生成される一次元空間である。
\item $S_m(w)$の全行列は特異であり、その中に非零階数1行列が存在する。
\end{enumerate}
この分岐は四因子独立の基底変更で不変である。
第二枝を排除するAS定理はここではまだ得られていない。
\end{proposition}
\begin{proof}
二つの証明書で三因子を対称化し、残る因子が位置2となるよう四窓を巡回移動する。
三つの対称因子のいずれを共通因子と見ても、未対称因子への交換方程式は同じ完全対称性を表す。
隣接縮約像の直交空間は$\bigwedge^2(V^*)+L$と書け、
対称な三因子での三重交差は$B_3^*$である。
有限Frobenius条件からその次元は4であり、一因子階数4によって$w$はこの空間の基底テンソルとなる。
従って系\ref{cor:partial-symmetric-exhaustive-split}を適用できる。
元の表示に戻る際には
$S_i((g_0\otimes\cdots\otimes g_3)w)=g_0 S_i(w)g_i^{-1}$
により、次元・行列階数・可逆元の有無は保たれる。
巡回移動で共通因子が変わるときも、三対称因子間の同一視で可逆な合成を施すだけなので、
残る対称化空間の階数条件は同じである。
\end{proof}

\begin{corollary}\label{cor:one-basepoint-symmetrizer}
一点基点の縮約族で$B_4^*=k p_0^{\otimes4}$となる場合、$S_2(w)$は一次元で、
その非零元は全て階数1である。従ってこのチャートの対称化失敗は行列式の消滅として直接検出できる。
\end{corollary}
\begin{proof}
前命題を適用する。$p_0^{\otimes4}$の一因子縮約の階数は1である。
\end{proof}

\begin{lemma}[一点基点の場合の低次socle]\label{lem:one-basepoint-cubic-socle}
$\dim B_3=4$かつ$B_4^*=k p^{\otimes4}$、$p\ne0$なら、
$B_3$内のsocleは三次評価$B_3\to k$の核であり、次元3である。
\end{lemma}
\begin{proof}
$B_4\simeq k$を$p^{\otimes4}$で同定すると、$\ell\in B_1$、$b\in B_3$の積は
$\ell(p)b(p)$に対応する。$\ell(p)\ne0$となる一次式があるので、
全ての一次式に殺されることと$b(p)=0$は同値である。
三次評価は非零であり、$B$は次数1で生成されるから、これがsocleそのものである。
\end{proof}

\begin{proposition}[三因子対称チャートでの四因子同値判定]
\label{prop:partial-symmetric-orbit-congruence}
四次元の二次式空間$L,L'$から三次双対基底で作る二つの三因子対称テンソル$w,w'$について、
$\dim B_3=\dim B'_3=4$、隣接二対二階数6、
\[
 S_1(w)=S_3(w)=kI,\qquad S_1(w')=S_3(w')=kI
\]
を仮定する。
このとき$w,w'$が四因子独立の$G$作用で同値であることと、$L,L'$が通常の$\operatorname{GL}(V^*)$同値であることは同値である。
従ってAS正則部分に対象を限れば、これは当該チャート内の添字保持同型の必要十分条件でもある。
一般の非正則候補については、ここでは四窓の表示同値を述べている。
\end{proposition}
\begin{proof}
$w'=(g_0\otimes g_1\otimes g_2\otimes g_3)w$とする。
線型symmetrizerの定義から
\[
 S_j(w')=g_0S_j(w)g_j^{-1}
\]
が従う。$j=1,3$の一次元条件により、$g_1,g_3$はそれぞれ$g_0$のスカラー倍である。
因子$0,1$の縮約像は、従って共通変換$g_0\otimes g_0$で移る。
その直交関係$\bigwedge^2(V^*)+L$を取れば$L,L'$の通常の同値を得る。

逆に共通の基底変換で$L$を$L'$へ移せば、三つの対称因子上の三次双対空間も移る。
双方は同じ四次元空間の基底なので、残る因子2の独立な可逆変換で基底を合わせられる。
従って$G$同値である。AS正則部分での最後の解釈は、標準四窓の同型判定による。
\end{proof}

\begin{proposition}[一般位置八点でも生じる非Koszul標本]
\label{prop:general-position-octad-nonkoszul}
全ての四点部分集合が独立である八点完全交叉からも、
有限Frobenius条件と全四次検査を通り、五次で七次元のKoszul障害を持つ境界が得られる。
これは当該標本の存在主張であり、一般点の障害次元の認定ではない。
\end{proposition}
\begin{proof}
反対称$4\times4$行列$A$の上三角成分を、順序$01,02,03,12,13,23$で
$(1,2,3,5,7,11)$とし、
\[
 M=(I-A)(I+A)^{-1},\qquad P=(I\mid M)
\]
とする。$MM^{\mathsf T}=I$であり、$P$の70個の四次小行列式は全て非零である。
八列の二次評価行列は階数7なので、その核は三つの二次式を与える。
$h=e_0+e_1+e_2+e_3$で非斉次化した三式の商の次元は8であり、
$h=0$上の斉次三式の商も8次元のArtin環となる。
従って無限遠に追加点はなく、三式の完全交叉は上の相異なる八点そのものである。
これらは有理数上のGroebner基底で検算できる。

列を$h=1$で正規化し、$q=e_1^2+2e_2^2+3e_3^2$とする。
これは$p_0=(1,0,0,0)$だけで零となる。
命題\ref{prop:reduced-octad-one-point-frame}から作ったテンソルについて、厳密消去は
\[
\begin{gathered}
\text{一因子階数 }(4,4,4,4),\quad\text{二対二階数 }(6,6,6),\\
(\dim W_{i,3},\dim W_{i,4})=(4,1)\quad\text{全始点},\qquad
\dim C_{i,i+4}=35\quad\text{全始点},\\
\operatorname{rk}d_2(0,5)=77,\quad\operatorname{rk}d_3(0,5)=36
\end{gathered}
\]
を与える。従って始点0の五次$H_2$は7次元である。
$p_0$の平方型関係による六次元cycle部分も単射であり、その商は一次元となる。
また$\dim S_1=\dim S_2=\dim S_3=1$、中央のsymmetrizerの階数は1である。
再現可能な入力、幾何の検査と同じ五次行列の検算は
\path{computations/verify_reduced_octad_frame.py}に置く。
\end{proof}

\section{長さ二の基点による四次障害}
\label{sec:basepoint-length-two}
一点基点の一般階数問題を残したままでも、複数点と非被約基点は一回Torの部分複体から排除できる。
以下では$\dim B_1=4$, $\dim B_2=6$, $\dim B_3=4$を仮定し、
基点スキームを$Z=\operatorname{Proj}B$とする。

\begin{theorem}[長さ二の基点の排除]\label{thm:length-two-basepoint-obstruction}
$Z$が長さ二の部分スキームを含むとする。
内部次数4、位置1から位置1への一回Tor複体は
\[
 0\longrightarrow D_3\longrightarrow D_2\longrightarrow D_1
 \longrightarrow D_0\longrightarrow0,
 \qquad
 D_h=\bigoplus_{\substack{p+q=h\\0\leq p\leq2,\ 0\leq q\leq1}}
 Q_p\otimes B_h^*\otimes Q_q
\]
であり、次元は$(1,32,156,160)$である。その最上ホモロジーは42次元以上である。
従って$T_L$はKoszulではない。
さらに$w$が有限Frobenius候補なら、その元四窓について位置1に対応する始点で
$\dim C_{i,i+4}\geq42$、従って四次障害$\delta_i\geq7$となる。
特に三因子対称な標準Koszul AS候補の基点スキームは、空か、ただ一つの縮約点である。
\end{theorem}
\begin{proof}
位置と内部次数を命題\ref{prop:one-return-bar-tor}へ代入すると、上の複体を得る。
両側の最大Koszul次数は2と1であり、$M=D A_L$の成分は$B_h^*$となる。

まず一点$[v]\in Z$を取る。$F_h=k v^{\otimes h}\subset B_h^*$は縮約で保たれ、
一点評価の部分複体を与える。
その左側は$Q_2\to Q_1\to Q_0$、右側は$Q_1\to Q_0$のテンソル積である。
一次式の微分は$v$での評価、二次関係の微分は一因子の$v$縮約である。
$L$の二次式は$v$で零なので、$Q_2\to Q_1$の像は$\ker(v:V^*\to k)$に含まれる。
交換子の縮約がこの三次元空間を尽くすので、左側のホモロジーは次数2の七次元だけ、
右側のホモロジーは次数1の三次元だけである。
K\"unnethの公式により一点部分複体のホモロジーは次数3の$7\cdot3=21$次元だけとなる。

長さ二のスキームが相異なる二点なら、両点の三次cubeは独立である。
従ってその二つの21次元核は$D_3$内で交わらず、42次元の核を与える。
ここには$D_4$がないので、そのまま最上ホモロジーに入る。

長さ二のスキームが一点に支持されるなら、標数零でそれは接ベクトル$v+\epsilon u$、
$\epsilon^2=0$、$u\notin kv$で表せる。
次数$h\geq1$で二次元のジェット評価双対$F_h\subset B_h^*$を取り、$F_0=k$とする。
この部分複体は一点部分複体を含む。その商は次数0の項が零で、正次数では一次元である。
実際、$v^{\otimes h}$を除いたジェットの縮約は、商では一次式の$v$評価に等しい。
従って商複体は一点複体の次数0の項だけを除いたものであり、
次数3のホモロジーは21次元、次数2のホモロジーは零である。
部分・全体・商のホモロジー長完全列と、一点部分の次数2ホモロジーの消滅から、
ジェット部分複体の次数3ホモロジーは$21+21=42$次元となる。
これも$D_3$内の核として全複体へ単射する。

命題\ref{prop:one-return-bar-tor}により、この核は
$\operatorname{Tor}^{T_L}_{4,4}(S,S)$の指定頂点成分の直和因子である。
$T_L$の各始点のHilbert列は$(1,4,6,4,1)$なので、Koszulならその成分の次元は
$(1-t)^{-4}$の四次係数35でなければならず、矛盾する。
有限Frobeniusの場合には$E(w)=T_L$であり、二次双対の四次元数は
$\operatorname{Tor}^{E(w)}_{4,4}$の次元だから、述べた四次障害を得る。

最後に、代数閉体上の射影スキームが空でも一縮約点でもなければ、相異なる二閉点、
または一閉点の非零接方向からなる長さ二の部分スキームを含む。
従って最後の結論も従う。
\end{proof}

\begin{corollary}[残る一点枝の縮約八点表示]\label{cor:simple-basepoint-reduced-net}
三因子対称な有限Frobenius候補が全四次Koszul検査を通り、基点を持つなら、その四二次式空間$L$は
相異なる八点完全交叉を定める三二次式と、そのうち一点だけで消える第四式を基底として持つ。
従って\textup{命題\ref{prop:reduced-onepoint-commutative-betti}}の可換最小分解は、
この残余AS候補枝全体に適用できる。ただしその分解だけではAS排除を結論しない。
\end{corollary}
\begin{proof}
前定理から基点スキームは一縮約点$p$である。
$L$のイデアルの高さは3なので、一般の三式は正則列となる。
$p$では四式の一次微分が三次元余接空間を生成するため、一般の三式はそこで横断的に交わる。
$U=\mathbb P^3\setminus\{p\}$上では$L$の評価は各点で全射である。
従って三式とその共通零点のincidence
$\{(x,f_1,f_2,f_3):x\in U,\ f_i(x)=0\}\to U$は階数9のベクトル束であり、滑らかである。
パラメータ空間$L^3$の一般点で、三式は正則列かつ$p$で横断的なので、
その完全交叉の残る長さは$8-1=7$である。従ってincidenceの$L^3$への射は支配的である。
その一般ファイバーは滑らかなincidenceの局所化として正則であり、
パラメータの函数体は標数零なので、この零次元ファイバーは幾何学的に縮約である。
滑らかさの開性から一般の三式の交叉は$U$上でも縮約となる。
従って次数$2^3=8$の相異なる八点完全交叉を得る。
これは本状況でのBertiniの議論をincidenceで書いたものである。
$L$の基底をなす第四式は、他の七点で消えればそこも$L$全体の基点となるので、
$p$でのみ消える。可換最小分解の命題を適用する。
\end{proof}

\begin{remark}
\path{computations/verify_basepoint_length_two.py}は、基点自由な対照点、二点を持つ例、
接ベクトルを含む非被約例について、三次商次元4、微分の零合成と全階数を有理数で検算する。
四次の最上ホモロジーはそれぞれ35,46,84であり、普遍的下界42と整合する。
後二例が有限Frobeniusであるとは主張せず、定理の証明もこの標本の外挿ではない。
\end{remark}

\section{左右反転対合による三因子対称枝のAS決定}
\label{sec:partial-symmetric-as}
ここでは一般点の大行列階数を計算せず、戻り辺一回の複体の対称性から基点枝全体を排除する。

\begin{theorem}[基点を持つ三因子対称切断の普遍的非Koszul性]
\label{thm:partial-symmetric-basepoint-exclusion}
$B=\operatorname{Sym}(V^*)/(L)$について$\dim B_1=4$, $\dim B_2=6$, $\dim B_3=4$とし、
$L$が射影基点を持つとする。このとき$T_L=A_L\ltimes D A_L$はKoszulではない。
より正確には、位置1から位置2の内部次数5成分で
\[
 \dim\operatorname{Tor}^{T_L}_{4,5}(S,S)\geq5.
\]
従って三次双対フレームから作る有限Frobenius候補に、標準Koszul AS対象は存在しない。
基点の縮約性や個数、四次Koszul検査の正確性を仮定しない。
\end{theorem}
\begin{proof}
基点を$p$とし、補題\ref{lem:one-return-point-quotient}の商複体$\overline C_\bullet$を取る。
その係数空間は
\[
 N_r=B_r^*/k p^{\otimes r},\qquad
 (\dim N_1,\dim N_2,\dim N_3)=(3,5,3),
\]
であり、次数1の項は零である。
$Q_a$上の語順反転を$\tau_a$とする。
$Q_0=k$, $Q_1=V^*$では恒等写像であり、
$Q_2=\bigwedge^2(V^*)+L$は語順反転で保たれるので、$\tau_2^2=1$である。

$\overline C_h$の$(a,b)$直和因子上で、
\[
 \sigma_h(u\otimes f\otimes v)
   =(-1)^{ab}\tau_b(v)\otimes f\otimes\tau_a(u)
 \qquad(a+b=h)
\]
と定める。$B$の乗法は可換なので、$B_r^*$への左右の一次式縮約は同一である。
語順反転は$Q_a$の最後の文字による左微分を、$Q_a$の最初の文字による右微分へ移す。
全微分の右成分の符号$(-1)^a$と表示の$(-1)^{ab}$を合わせると、
$\sigma\partial=\partial\sigma$となる。また$\sigma^2=1$である。
点評価部分も保たれるので、これは商複体上の対合として定まっている。

標数零なので、正負の固有部分へ複体を直和分解できる。
負固有部分の次数2,3,4の次元を計算すると
\[
 (\dim\overline C_2^-,\dim\overline C_3^-,\dim\overline C_4^-)
 =(60,200,135)
\]
である。実際、次数2の$(2,0),(0,2)$対から$10\cdot3=30$、
$(1,1)$因子では符号付き交換が$-\mathrm{flip}$なので
$\operatorname{Sym}^2(V^*)\otimes N_1$から$10\cdot3=30$を得る。
次数3は二つの交換される因子から$10\cdot4\cdot5=200$となる。
次数4の$\tau_2$付き交換は、一方の$Q_2$へ$\tau_2$を作用させる共役によって通常の交換へ直るので、
負固有空間は$\bigwedge^2 Q_2\otimes N_3$と同次元の$45\cdot3=135$である。

従って負固有部分のEuler標数は
\[
 135-200+60=-5.
\]
次数1の項がないから
\[
 \dim H_3(\overline C_\bullet^-)
 =5+\dim H_2(\overline C_\bullet^-)+\dim H_4(\overline C_\bullet^-)
 \geq5.
\]
一点部分の$H_2,H_3$は零なので、前記補題から
$H_3(C_\bullet)\simeq H_3(\overline C_\bullet)$となる。
命題\ref{prop:one-return-bar-tor}により、これは$\operatorname{Tor}^{T_L}_{4,5}$の直和因子である。
Koszul性に必要な次数の一致に反する。
有限Frobenius候補では$E(w)=T_L$であり、二次双対間のKoszul性の同値から最後の結論を得る。
\end{proof}

\begin{corollary}[八変数階数に必要な上界]\label{cor:product-net-rank-bound}
八係数チャートの一般点では、戻り辺一回の小複体について
\[
 \rk\partial_3=136,\quad \rk\partial_4\leq339<344,
 \qquad
 \rk\overline\partial_3=107,\quad \rk\overline\partial_4\leq288<293.
\]
従って従来の階数移行命題の前提となる不等式も証明された。
正確な一般階数が337（商では286）であるとの主張はしない。
\end{corollary}
\begin{proof}
$\partial_2$の全射性から$\rk\partial_3\leq136$であり、
チャート内の既知の有理数点で136を達成するので一般点では等しい。
前定理の$\dim H_3\geq5$から
$\rk\partial_4\leq480-136-5=339$となる。
商の二式は補題\ref{lem:one-return-point-quotient}の階数公式による。
\end{proof}

\begin{corollary}[三因子対称枝の完全なAS条件]\label{cor:partial-symmetric-as-complete}
三因子対称な有限Frobenius候補が標準Koszul AS対象を定めることは、
その可換四二次式空間$L$が基点を持たないことと同値である。
従ってこの枝の全AS対象は、既知の完全対称四二次式完全交叉チャートに属する。
これは全四窓の完全分類ではなく、三因子対称化可能な枝の被覆・正則性を決定するものである。
\end{corollary}
\begin{proof}
基点があれば前定理で排除される。なければ四二次式は正則列であり、
定理\ref{thm:symmetric-frobenius-chart}を適用できる。
\end{proof}

\begin{corollary}[AS対象の対称認識には二証明書で十分]
\label{cor:two-positive-as-recognition}
標準Koszul AS対象の標準四窓テンソル$w$について、次は同値である。
\begin{enumerate}
\item 三つの$S_j(w)$のうち少なくとも二つが可逆元を持つ。
\item 三つ全てが可逆元を持つ。
\item $w$は四因子独立の変換で完全対称にできる。
\end{enumerate}
従って完全対称枝・可換多項式の交代枝の双方を除いた未分類AS残余では、
正符号・負符号の可逆証明書はそれぞれ高々一位置にしか存在しない。
両符号が存在するなら、系\ref{cor:mixed-symmetry-witness-exclusion}により同じ位置に限られる。
\end{corollary}
\begin{proof}
二つの正符号証明書があれば命題\ref{prop:two-positive-global-split}の三因子対称表示へ移れる。
基点枝は前定理で排除されるから、完全対称枝に属する。
他の含意は既存の線型対称化判定による。
負符号二証明書の場合は系\ref{cor:two-alternating-witnesses}を適用する。
\end{proof}

\begin{remark}[検算と未確定の精密値]
\path{computations/verify_one_return_involution.py}は有理数境界点で
$\sigma^2=1$、$\sigma\partial=\partial\sigma$、両固有複体の次元と階数を直接検算する。
その点で正負の$H_3$次元は1と6であり、既知の合計7に一致する。
上の普遍的下界5はこの標本から外挿したものではなく、対合とEuler標数のみから従う。
八変数の一般点で正確な障害次元が7であるかは未確定だが、AS排除にはその精密値を必要としない。
\end{remark}

% !TeX root = ../papers/nc-p3.tex
\section{一つの隣接対称性から四平方の場合を認識する条件}
\label{sec:adjacent-four-squares}
指定した隣接正証明書で最初の二因子を同一視し、
$w\in\operatorname{Sym}^2 V\otimes W\otimes U$とする。
第三切断の階数6の分解を$P(x):F\to W$, $Q(x):U^*\to F$,
$\dim F=6$と書く。ここで$x\in V^*$であり、
$P(x)Q(y)=P(y)Q(x)$は$w$の最初の二因子を$x,y$で縮約した写像である。
階数分解より$\bigcap_x\ker P(x)=0$、$\sum_x\operatorname{im}Q(x)=F$である。
$w$の対称二次消去子$L$が独立な四つの平方$x_0^2,\ldots,x_3^2$で張られる場合を考える。
このとき$[x_i]$は$\ell^2\in L$を満たす射影点の全体である。実際、混合項の全消滅より
$\ell$の非零座標は一つである。$P_i=P(x_i)$, $Q_i=Q(x_i)$,
$M_{ij}=P_iQ_j$と書くと、$M_{ii}=0$, $M_{ij}=M_{ji}$である。


\begin{proposition}[四本の核直線による対称認識]\label{prop:square-kernel-recognition}
一因子階数を全て4とする。各$i$で$\rk P_i=\rk Q_i=3$、
四本の直線$\ker Q_i$が$U^*$を張るとする。このとき独立な因子の基底変更により
\[
 w_{abcd}=\begin{cases}1&\{a,b,c,d\}=\{0,1,2,3\},\\0&\text{その他}\end{cases}
\]
となる。特に六つの対全てで$\rk M_{ij}=2$ ($i\ne j$)である。
これは四平方の可換完全交叉の対称チャートの一点である。この認識にAS性は仮定しない。
\end{proposition}
\begin{proof}
$u_i\in\ker Q_i$を基底として選び、$K_i=\operatorname{im}Q_i=\ker P_i$、
$I_{ij}=K_i\cap K_j$と置く。$Q_i u_j\in I_{ij}$である。
実際$P_jQ_i u_j=P_iQ_j u_j=0$である。
各$i$で三つの$Q_i u_j$ ($j\ne i$)は$K_i$の基底である。

各無順序対から一本ずつ$v_{ij}=v_{ji}\in I_{ij}$を選ぶ。
各頂点$i$で三つの$v_{ij}$が$K_i$の基底になるという条件は
$\prod_{i<j}I_{ij}$の非空開条件である。固定した$i$については
$v_{ij}=Q_i u_j$と選べるため非空性が従う。
パラメータ空間は既約なので四つの条件を同時に満たせる。
六つの$v_{ij}$は全ての$K_i$を張り、階数分解の全射性よりその和は六次元の$F$である。
よって六つは基底である。従って$I_{ij}=k v_{ij}$であり、
$\rk M_{ij}=\dim K_j-\dim I_{ij}=2$を得る。

非零の$Q_i u_j$と$Q_j u_i$は同じ直線にあるから、
$Q_i u_j=c_{ij}Q_j u_i$、$c_{ji}=c_{ij}^{-1}$と書ける。
相異なる$i,j,k$について
\[
 c_{ik}M_{jk}u_i=M_{ij}u_k=c_{jk}c_{ij}M_{jk}u_i.
\]
$\ker M_{jk}=\langle u_j,u_k\rangle$なので$M_{jk}u_i\ne0$であり、
$c_{ik}=c_{ij}c_{jk}$を得る。各$u_i$をスカラー調整して全ての$c_{ij}=1$にできる。
従って$M_{ij}u_k$は三添字に対称であり、重複添字では零となる。
相異なる三添字から得られる四つのベクトルは一因子階数4より$W$を張る。
欠けた添字で番号付けし、これらを$W$の基底に取れば表示した全係数を得る。
\end{proof}

四平方の対称テンソルでは、$F$を無順序二添字の六次元空間に取り、
$Q_i u_j$を$i\ne j$でその対の基底、$Q_i u_i=0$とすれば、全ての仮定を満たす。
$P_iQ_i=0$の下で両階数が3以上という開条件は両方3と同値であり、
その上で核直線の独立性も開条件である。階数分解の変更は$F$の基底変更であり、
四平方の方向の並べ替えとスカラーは判定を変えない。
従って指定した隣接対称性と四平方の消去子を持つ候補内で、
この非空な相対開部分はただ一つの添字保持同型類を与える。
双対化して$(P_i,Q_i)$を$(Q_i^{\mathsf T},P_i^{\mathsf T})$に替えても命題を適用できる。
よってこの四平方の枝で標準対称類の外に残るには、
どこかの軸で階数対が$(3,3)$と異なるか、全軸が$(3,3)$なら
入力側の四核直線も出力側の四余核の双対直線も、ともに従属でなければならない。

% !TeX root = ../papers/nc-p3.tex
\section{対称完全交叉の全点での三階数切断の滑らかさ}
\label{sec:ci-global-tangent}
$\mathfrak X_6$は三つの二対二分割について全ての7次小行列式を課した射影スキームである。
定理\ref{thm:symmetric-presentation-product}の対称チャート$S$が、そのどの点でも
$\mathfrak X_6$の局所構造全体を与えることを示す。

\begin{lemma}[中央三次数の次数零剛性]\label{lem:ci-middle-rigidity}
$B=\operatorname{Sym}(V^*)/(L)$を四つの二次式のArtin完全交叉とし、
$\mathfrak m=B_{>0}$, $N=\mathfrak m/B_4$を元の次数$1,2,3$に置く。
このとき$\operatorname{End}_B(N)_0=k$、
$\operatorname{Ext}^1_B(N,N)_0=0$である。
\end{lemma}
\begin{proof}
全てのHomとExtで内部次数零を取る。$D_0=B/B_4$と置く。
$B$はFrobeniusで自己入射的、$B_4=k(-4)$であり、
二次完全交叉はKoszulなので
$\operatorname{Ext}^p_B(k,k(-4))_0=0$ ($p\ne4$)である。
$0\to k(-4)\to B\to D_0\to0$から
$\operatorname{Ext}^p_B(k,D_0)_0=0$ ($p=1,2$)を得る。
また$\operatorname{Hom}_B(k,D_0)_0=0$である。
$0\to N\to D_0\to k\to0$と$\operatorname{Ext}^p_B(k,k)_0=0$ ($p>0$)より
$\operatorname{Ext}^1_B(k,N)_0=k$、$\operatorname{Ext}^2_B(k,N)_0=0$となる。
次に$0\to\mathfrak m\to B\to k\to0$で$N_0=0$を使うと
$\operatorname{Hom}_B(\mathfrak m,N)_0=k$、
$\operatorname{Ext}^1_B(\mathfrak m,N)_0=0$である。
最後に$0\to k(-4)\to\mathfrak m\to N\to0$へ$\operatorname{Hom}_B(-,N)_0$を適用する。
$N_4=0$なので$\operatorname{Hom}_B(k(-4),N)_0=0$となり、両方の主張が従う。
\end{proof}

\begin{theorem}[対称チャートは三階数切断の開部分]\label{thm:symmetric-open-all}
$\mathfrak X_6$は$S$の全点で滑らかであり、$S$は$\mathfrak X_6$の69次元の開部分である。
特に$Z_{\mathrm{sym}}$と異なる既約成分の交わりは$S$を避ける。
\end{theorem}
\begin{proof}
四因子の基底変更により$[w]=[w_L]$を完全対称なCIのtraceとしてよい。
アフィン錐の非零点$w$での接空間をそれぞれ$T_X,T_S$とする。
まず$\dot w\in T_X$が因子$0,1$の交換で不変とする。
$D=k[\epsilon]/(\epsilon^2)$上の$w'=w+\epsilon\dot w$はこの二因子に対称である。
各分割で階数6の小行列式の一つが単元で、全7次小行列式が消えるため、
行列の像・核は直和因子である（単元小行列で掃き出す）。
$01|23$の縮約の$\operatorname{Sym}^2 V^*$内の核は$L$を持ち上げる階数4の$L'$となる。
$B'=\operatorname{Sym}_D(V^*\otimes D)/(L')$は完全交叉族の引戻しであり、
各次数は$D$上自由である。これは、特殊ファイバーで正則列を解消するKoszul複体を
次数ごとに単元掃き出しすることからも従う。

$02|13$の階数分解を$P'(x)Q'(y)$と書く。第0・第1因子の対称性より
$P'(x)Q'(y)=P'(y)Q'(x)$である。従って順に$Q',P'$を一次式の作用とする
階数$(4,6,4)$の三層加群を作れる。$L'$の全二次式が両作用の積を消し、
その加群は$B'$上である。特殊ファイバーは$N$に同型である。
実際、標準の階数分解は$B_1\to B_2\to B_3$の二つの乗法であり、
$B_3$を$B_1$の双対と同定するtraceは最上テンソルを与える。
全階数分解は中間空間の基底変更でこれに移る。

同じ$B'$に対して標準加群$N'_{\rm st}=B'_{>0}/B'_4$がある。
三次数ごとに基底を合わせ、二つの加群の特殊ファイバーを同一にする。
その二つの一次作用の差は$\epsilon$の係数として$B$上の次数付き自己拡大を定める。
具体的には、$B'$の元$b'$の二つの作用の差を$\epsilon$で割れば、
特殊化$b$だけに依存する写像$\delta(b)$が得られる。
$\delta(bc)=b\delta(c)+\delta(b)c$なので、対角に元の作用、
上三角成分に$\delta$を入れた$N\oplus N$が拡大加群である。
補題\ref{lem:ci-middle-rigidity}から次数零で分裂し、差は次数別基底変更
$1+\epsilon f$で消える。従って$w'$は$B'$の対称traceテンソルから、
入力・出力の基底変更とtraceの単元倍によって得られる。中間次数の基底変更は
$P'Q'$の積で相殺する。ゆえに$\dot w\in T_S$である。

$w$は完全対称、両方の錐は四因子の置換で不変なので$T_X,T_S$は
$\mathfrak S_4$表現である。上の議論を共役すれば、各互換の固定部分は$T_S$に入る。
従って商$T_X/T_S$では全互換が$-1$で作用し、交代化射影が恒等になる。
一方、$T_X$に完全交代テンソルは存在しない。実際、$01|23$行列の左右核は
ともに$\bigwedge^2 V^*$を含む。接条件は左右核上で$\dot w$の双線型形式が零となることである。
完全交代テンソルを二組の交代二形式に評価すると、その四文字係数の4倍を得るので、
この条件は全係数を消す。従って交代化射影は$T_X$上零であり、$T_X=T_S$である。

射影化後の接次元は定理\ref{thm:symmetric-presentation-product}より69である。
$S$を含むので局所次元は少なくとも69であり、接次元との一致から
$\mathfrak X_6$はここで正則かつ滑らかである。通過する既約成分は
$Z_{\mathrm{sym}}$一つで、$S$はその中で開である。他の有限個の成分を除けば
各点の近傍を$S$内に取れるため、全体としても$S$は開である。
\end{proof}

この定理は全ての対称正則点を扱うが、決定式条件を外したAS表示空間の滑らかさや、
$\mathfrak X_6$の他の成分の不存在は与えない。
$S$の閉包には既知の有限Frobenius非Koszul境界がある。

\begin{corollary}[未分類低階数部分の特殊化]\label{cor:closed-low-residual}
有限Frobenius標準候補のうち、$\rho\leq6$で、対称CIと可換多項式軌道のいずれにも
入らない部分は、全有限Frobenius標準候補の中で相対閉である。
標準AS対象に制限しても同じ相対閉性が成り立つ。
\end{corollary}
\begin{proof}
二つの既知部分は、定理\ref{thm:symmetric-open-all}と
定理\ref{thm:pure-anti-and-alternating-local}により$\mathfrak X_6$の開部分である。
その補集合は$\mathfrak X_6$で閉であり、$\mathfrak X_6$自体が射影テンソル空間で閉なので、
候補部分・AS部分との交叉を取ればよい。未分類AS部分の非空性はこの議論からは従わない。
\end{proof}

\paragraph{平方を含まない完全交叉での検算。}
巡回添字で$f_i=x_i^2+c_ix_{i+1}x_{i+2}$、$(c_0,c_1,c_2,c_3)=(1,2,3,5)$と取る。
共通零点に一つでも非零座標があれば全座標が非零になり、四式の積から$\prod c_i=1$が必要になる。
また$\ell^2\in L$なら、非零座標$\ell_i$は混合項によって$\ell_{i+2}\ne0$を強制するが、
$x_ix_{i+2}$はどの$f_j$にも現れない。従って射影基点も非零平方もない。
\path{computations/verify_ci_middle_rigidity.py}で四次traceを求め、
固定した$B$の二作用を線型化すると192変数で階数125を得る。
核次元67は三次数の基底変更群の次元からスカラー安定化群を除いた数に等しい。
三分割全部のアフィン接次元は70、その$01$交換固定部分は55で、射影接次元は69である。

% !TeX root = ../papers/nc-p3.tex
\section{対称成分の四つの境界因子}
\label{sec:symmetric-boundary-divisors}
$Z_{\mathrm{sym}}=\overline S$で、添字$e$以外の三因子を同時対称化できる
有限Frobenius候補のうち$S$に属さないものの閉包を$D_e$とする。
添字は固定し、四つの位置を同一視しない。

\begin{theorem}[三因子対称境界の余次元と局所被覆]\label{thm:four-boundary-divisors}
$D_0,D_1,D_2,D_3$は$Z_{\mathrm{sym}}\setminus S$の相異なる68次元既約成分である。
各因子の一般点は有限Frobeniusで非Koszulである。
各因子のある稠密開部分では、$\mathfrak X_6$は滑らかであり、その点の
$\mathfrak X_6$内の開近傍にある標準AS対象は全て$S$に属する。
この開部分には定理\ref{thm:nonkoszul-symmetric-boundary}のBoolean点の対応する巡回回転が入る。
\end{theorem}
\begin{proof}
$e=2$とし、$X=V^*$を対称な三因子の共通双対空間とする。
$\Gr(4,\Sym^2X)$の基点を持つ部分を$\Delta$とする。
基点$[p]$を指定したincidenceは$\mathbb P^3$上の$\Gr(4,9)$束で、既約かつ次元23である。
その射影像が$\Delta$である。Boolean点の基点はただ一つなので一般ファイバーは零次元であり、
$\Delta$は24次元Grassmann多様体の既約因子となる。

$B=\Sym(X)/(L)$で$\dim B_3=4$となる開部分上、$B_3^*$の射影フレームを取り、
対称な位置1と3に独立な$\PGL_4$変換を施す。位置0の座標は固定する。
このパラメータ空間は基底上の相対次元45の滑らかな束である。
一般の$L$はCIなので、一般の像は$S$に入る。従って全ての像は$Z_{\mathrm{sym}}$にある。
$\Delta$上でも$\dim B_3=4$となる開部分はBoolean点を含み、既約である。
この上の有限Frobenius条件は非空開である。実際、縮約の二対二階数は$6$以下であり
（三因子のうち二因子の積が$B_2$を経由する）、残る次元等号は
定理\ref{thm:frobenius-tensor-reconstruction}の開条件になる。
逆に三因子を対称化した有限Frobenius候補では、その連続三因子の交叉が$B_3^*$で
次元4である。系\ref{cor:partial-symmetric-exhaustive-split}から、$S$の外なら
$L\in\Delta$であり、こうした候補の全てがこの像に入る。

一般点では$S_1,S_3$は各々一次元で可逆行列を含む。
正規化した表示では両方が恒等行列を含み、Boolean点で次元1を達成するからである。
この二直線から位置1と3の射影基底変更を一意に戻せる。
三因子を同一視した後で、最初の二因子の縮約の消去子から$L$、残る因子の三次縮約から
$B_3^*$の射影フレームを回収できる。いずれも定数階数の核・像を取る正則な操作である。
従って$D_2$は次元$23+45=68$の既約閉部分である。
$S$は69次元既約な$Z_{\mathrm{sym}}$の開部分なので、$D_2$は境界の既約成分となる。
巡回回転で他の三位置についても同じ結論を得る。
異なる二つの像の一般部分が重なると、共通の因子を基準に取ったとき残る全因子との
交換に可逆な正証明書があり、完全対称化ができてしまう。これは基点枝の二分岐に反する。
同じ既約閉包を持つ稠密な構成可能集合は共通の開部分を持つため、四つの閉包は異なる。
一般部分の非Koszul性は定理\ref{thm:partial-symmetric-basepoint-exclusion}による。

最後にBoolean点で三分割の行列の左右核に対する接条件を全て課すと、
300行256列の有理行列の核次元は70である
（\path{computations/verify_boolean_boundary_divisor.py}）。
$\mathfrak X_6$は69次元成分$Z_{\mathrm{sym}}$を含み、射影接次元も69なので、
Boolean点で滑らか、通過する成分は$Z_{\mathrm{sym}}$だけである。
従ってこの点の近傍で$\mathfrak X_6$を$Z_{\mathrm{sym}}$に制限できる。
$Z_{\mathrm{sym}}$上では対称チャートからの特殊化により全ての正symmetrizer空間が非零である。
Boolean点で$S_1,S_3$は各々可逆な生成元を持つ一次元空間だから、
その性質と有限Frobenius性はこの点の開近傍で続く。
その近傍のAS点は二つの可逆証明書により全て対称CIとなる。
$D_2$との交叉は非空開、従って稠密である。
\end{proof}

\begin{proposition}[境界の横断座標と行列式の重複度]\label{prop:boolean-transverse-order}
Boolean点において$Z_{\mathrm{sym}}\setminus S$の台は滑らかな因子である。
位置2を三対称因子の一つへ同一視する正symmetrizer直線の局所生成行列$H$について、
$\det H$はこの因子に沿ってちょうど三重に消える。
一般点での四つの基本因子の付値は$(0,1,1,1)$である。
\end{proposition}
\begin{proof}
前証明の逆写像により、Boolean点の近傍は$L$、射影三次フレーム、二つの射影基底変更で
局所的に表示される。その中で$S$の補集合は$\Delta$の引戻しである。
Boolean点では最初の三二次式は八つの相異なる点を切り、第四式が零となる点は
$p=(0,0,0,1)$だけである。三式の勾配は$p$のアフィンチャートで独立で、
形式近傍では三式から$p$を一意に持ち上げられる。他の七点では第四式が単元である。
従って残る局所方程式は持ち上げた$p$での第四式の値一つであり、
第四式へ$t e_3^2$を加える曲線はこれを一次で横断する。

$S_2$もBoolean点で一次元なので局所的に直線をなす。
因子の一般点は単一被約基点$p'$を持ち、系\ref{cor:basepoint-rank-one-symmetrizer}から
その直線は階数1となる。従って単元成分で掃き出した後の$3\times3$ブロックの各成分は
滑らかな因子の局所方程式で割り切れ、行列式は少なくとも三重に消える。
正確な次数を上の曲線で計算する。八点を$v=(a,b,c,1)$、$s_v=a+b+c$と書く。
四次traceの正規化として
\[
 F_t=\sum_{v}\frac{(-1)^{s_v}t}{t+s_v}\,v^{\otimes4}
\]
を取れる（$s_v=0$の係数は1）。最初の三二次式は全八点で消え、第四式で縮約した値は
次数2以下の任意の単項式に対する符号付き立方体和で零となる。
$t=0$で$F_0=p^{\otimes4}$である。
三次フレームの正規化座標$012,133,233,333$を$m_j$とすると、
命題\ref{prop:partial-symmetrizer-quartics}の生成行列は$H_{ij}=F_t(i,m_j)$である。
八項の和の行列式を展開すると
\[
 \det H(t)=\frac{-2t^3}{(t+1)^3(t+2)^3(t+3)} .
\]
従って因子の局所方程式の三乗を除いた因子はBoolean点で単元である。
一般点の離散付値環では階数1への落ちと付値の和3から基本因子の主張を得る。
逆向きの射影同一視は逆行列のスカラー倍なので、一般点での極限階数は3となる。
計算コードは四次traceが三次フレームで表示されること、Boolean点の全12個の向き付き
symmetrizer空間が一次元で、通常の三位置相互は階数4、位置2からは1、位置2へは3であることも検算する。
\end{proof}

この四因子は三因子対称な有限Frobenius非正則境界の閉包を尽くす。
他の境界因子の有無や、これらの因子の全特殊点でのAS性は未決定である。

\section{三次縮約からの指数成長障害}
\label{sec:cubic-slice-growth}
有限次数のTorとは別に、商の経路成長から全ての候補に共通の排除条件を得る。
ここで$C$は四窓の元の二次代数とし、$E$はその二次双対である。

\begin{proposition}[一つの矢空間を一次元へ落とす成長検査]\label{prop:cut-arrow-growth}
切る位置を3とし、$0\ne\ell\in V_3^*$について
\[
 X=V_0/(\ell\otimes1)R_3,\qquad Z=V_2/(1\otimes\ell)R_2
\]
と置く。$R_0,R_1$の像を$\bar R_0\subset X\otimes V_1$,
$\bar R_1\subset V_1\otimes Z$とし、
\[
 d_\ell=\dim \frac{X\otimes V_1\otimes Z}
              {\bar R_0\otimes Z+X\otimes\bar R_1}
\]
とする。各$m\ge1$で$\dim C_{0,4m}\ge d_\ell^m$である。
特に標準AS対象では四位置の全ての$\ell$で$d_\ell\le1$である。
\end{proposition}
\begin{proof}
各周期の一次空間を順に$X,V_1,Z,k$へ写し、最後の商写像に$\ell$を使う。
二次関係の像で商を取れば$C$の商代数を得る。位置2,3の二次関係の像は零であり、
残る二つは一つの三文字ブロック内だけに作用する。従って長さ$4m$の経路商は
表示した三文字空間の$m$個のテンソル積である。
特に$d_\ell\ge2$では長さ64の下界$2^{16}=65536$が$\binom{67}{3}=47905$を超える。
標準AS対象の長さ$4m$の次元は$\binom{4m+3}{3}$なので結論を得る。
\end{proof}

\begin{theorem}[連続する二つの純粋関係の排除]\label{thm:no-consecutive-pure-relations}
標準有限Frobenius候補で、非零ベクトル$a,b,c$が
$a\otimes b\in R_0$, $b\otimes c\in R_1$を満たすなら、
$\dim C_{0,4m}\ge6^m$である。特に、標準AS対象のどの連続三窓でも
$W_{i,3}$は非零の可分解三テンソルを含まない。
\end{theorem}
\begin{proof}
$a\otimes b\otimes c$は三次交叉に属する。完全trace対合により
$V_3^*\to W_{0,3}$は同型なので、このテンソルはある$\ell$で$w$を縮約したものになる。
$R_3$を$\ell$で縮約した像はその三テンソルの先頭の支持$ka$、
$R_2$での像は末尾の支持$kc$である。これは$R_2,R_3$が最上テンソルの隣接縮約像だからである。
従って$\dim X=\dim Z=3$で、指定した純粋関係は各々$\bar R_0,\bar R_1$への射影で消える。
両像の次元は5以下であり、$d_\ell\ge36-5\cdot3-3\cdot5=6$となる。
逆に可分解三テンソルが交叉に入れば両端で縮約して指定の二関係を得る。
純粋関係のある始点では長さ16ですでに$6^4=1296>\binom{19}{3}=969$である。
どの位置でも命題\ref{prop:cut-arrow-growth}を適用できる。
\end{proof}

\begin{corollary}[正証明書のある対への階数1写像]\label{cor:no-rank-one-into-positive-pair}
標準AS対象で位置$i,j$を可逆な正証明書により同一視できるとする。
別の位置$b$から$i$への正symmetrizer空間は、階数1の写像を含まない。
ここで指定した二位置が隣接することは要らない。
\end{corollary}
\begin{proof}
$i,j$を同一視してその交換で$w$を不変とする。
階数1の$H=u\ell$が$b$から$i$への正証明書なら、
$H$を位置$b$に作用させたテンソルは位置$i,b$に対称であり、変換した側は$ku$に支持される。
従って$\ell$による元の位置$b$の縮約は位置$i$で$ku$に支持される。
$i,j$の対称性から位置$j$でも同じ直線に支持され、縮約は非零の可分解三テンソルとなる。
これは定理\ref{thm:no-consecutive-pure-relations}に反する。
\end{proof}

\begin{corollary}[三因子対称な基点枝の全次数成長]\label{cor:basepoint-exponential-growth}
三因子対称な標準有限Frobenius候補に基点$[p]$があるとき、
対称な連続三位置から始まる長さ$4m$の次元は少なくとも$9^m$である。
基点で二次式の射影Jacobian階数が3なら、下界は$10^m$になる。
Boolean境界$t=0$では、構成する商の次元はちょうど$10^m$である。
\end{corollary}
\begin{proof}
三次双対空間に$p^{\otimes3}$があり、対称な三因子の両端の商はいずれも$H=V/kp$である。
二つの内側の関係は同じ$R=B_2^*\subset\operatorname{Sym}^2 V$であり、
$\dim R=6$, $p^2\in R$である。片側の射影の核は$kp^2$である
（$p\otimes v$が対称なら$v\in kp$）。
$H\otimes V\otimes H$の両端反転を$\tau$とすると、挿入関係は
$J+\tau J$、$J=\bar R\otimes H$、$\dim J=15$である。
反転の正固有部分は周囲で$6\cdot4=24$次元、関係側で
$(1+\tau)J$の次元は15以下なので、商次元は9以上となる。

Jacobian階数3なら、補空間を選び$V=H\oplus kp$と書いたとき、
$R\to\operatorname{Sym}^2 H$の核も$kp^2$である。
実際、$p\odot v$の形の関係は全$q\in L$について$q(p,v)=0$を満たすので、
$v$の$H$成分が零になる。像を$K$とすると$\dim K=5$であり、
$\bar R\subset H\otimes V$は$K$から$H\otimes kp$への写像のグラフになる。
三文字の中央も$H$へ射影すると、
$J,\tau J$はそれぞれ$K\otimes H$, $H\otimes K$の上のグラフである。
両底空間の交叉は$\operatorname{Sym}^3 H$の中で、
$K^\perp$の生成二次式による縮約が零となる部分である。
その次元は$10-3=7$である。実際、双対写像は非零二次式による三変数一次式の乗法で単射である。
この交叉から二つのグラフへ持ち上げると、中央の$p$成分の差は両端の反転で符号が変わる。
その値域は$\bigwedge^2 H$の三次元空間なので、$\dim(J\cap\tau J)\ge7-3=4$である。
従って$d_\ell=36-30+\dim(J\cap\tau J)\ge10$である。
Boolean点の射影挿入行列では階数が26となることを有理数で直接検算できる。
ここで求めた等号は商に関するものである。
\end{proof}

\paragraph{射影的な有限検査。}
標準有限Frobenius候補の元代数は$\dim C_{0,2}=\dim C_{1,3}=10$、
$\dim C_{0,3}=64-24-24+4=20$を満たす。
命題\ref{prop:cut-arrow-growth}の$d_\ell$は次の写像の余核次元でもある。
\[
 R_3\otimes C_{1,3}\ \oplus\ C_{0,2}\otimes R_2
 \ \xrightarrow{\ \Phi(\ell)\ }\ C_{0,3},\qquad
 (r\otimes z,y\otimes r')\longmapsto
 ((\ell\otimes1)r)z+y((1\otimes\ell)r').
\]
実際、二つの内側の関係は$C_{0,3}$の段階ですでに割られ、残りの商は両端の線型関係による。
したがって$\Phi$は$20\times120$の一次行列である。ある位置で19次小行列式全てが
$\mathbb P(V_3^*)$に共通零点を持つ候補は指数成長を持つ。
この条件は添字保持同型で不変であり、候補族上で相対閉である。
実際、定数階数の経路商束の上の階数条件を射影束から射影すれば閉となり、位置は四つだけである。
これは非正則性の十分条件であり、全ての非Koszul候補の検出は示していない。
\path{computations/verify_cubic_slice_growth.py}に、Boolean点および正則対照点での商の検算を置く。

\begin{corollary}[対称閉包内の二対二の可逆証明書]\label{cor:two-pairs-transfer-ranks}
$Z_{\rm sym}$内の標準AS候補で、四位置が二つの可逆正証明書の対に分かれ、
異なる対の間には可逆な正証明書がないとする。異なる対から一つずつ位置を選ぶと、
その二位置間に両向きの正証明書$H,K$が存在し、
\[
 \rk H=\rk K=2,\qquad \operatorname{im}H=\ker K,\quad
 \operatorname{im}K=\ker H
\]
となる。
\end{corollary}
\begin{proof}
$S$上では全位置を対称化する二方向の同一視は射影的に互いに逆である。
それらの二つの射影行列を記録したグラフの閉包は、射影行列空間の積からの射影の固有性により
$Z_{\rm sym}$へ全射する。従って所与の点には、互換の正証明書方程式と
$HK,KH$がスカラーであるという閉条件を満たす非零の極限$H,K$が存在する。
対をまたぐため両者は特異で、その積のスカラーは零になる。
両方の標的位置がそれぞれ可逆対の中にあるので、系\ref{cor:no-rank-one-into-positive-pair}が
いずれの階数1も排除する。二階数の和は4以下であり、全ての等式を得る。
この結論は当該配置のAS点の存在を仮定する条件付きの制約である。
\end{proof}

% !TeX root = ../papers/nc-p3.tex
\section{正次元ファイバーを許す四移動と反復可能な分岐の排除}
\label{sec:recurrent-point-correspondences}
位置 $i$ を固定し、標準四移動で $V_{i+4}=V_i$ とする。
点は $\mathbb P(V_i^*)$ の任意の幾何点でよい。
$p,q$ に対し
\[
 d(p,q)=\dim\frac{A_{i+1,i+4}}
 {R_i(p) A_{i+2,i+4}+A_{i+1,i+3}R_{i+3}(q)}
\]
を定める。前者の評価は関係の先頭、後者は末尾で取る。
$A_{i+1,i+4}$ は有限Frobenius候補で次元20である。この余核の正次元条件は
両点の対の上で二群の列がそれぞれ $p,q$ に線型な $20\times120$ 行列の階数条件である。
その被約閉対応を $D$ とする。実際の四辺の点列が $p,q$ を結べば必ず $(p,q)\in D$ である。
命題\ref{prop:cut-arrow-growth}は $d(p,p)$ を使う場合に当たる。

\begin{theorem}\label{thm:recurrent-point-correspondences}
標準四窓では以下の三条件が全ての位置で成り立つ。
\begin{enumerate}
\item 任意の有限個の相異なる点を頂点に、$p$ から $q$ に $d(p,q)$ 本の有向辺を置く。
各強連結成分のうち巡回路を持つものは、全ての辺の多重度が1の単純有向巡回路である。
一つの有向道に沿って互いに異なるそのような成分を通過できる個数は4以下である。
点を一般位置に置く仮定はない。
\item 閉既約部分多様体の周期列 $Y_0,\ldots,Y_l=Y_0$ と
$D\cap(Y_j\times Y_{j+1})$ 内の被約既約対応 $B_j$ があり、
両射影がそれぞれの $Y$ に支配的なら、各 $B_j$ の両射影は双有理である。
また $d=1$ が各 $B_j$ の一般点で成立する。
特にこのような周期列では正次元の一般ファイバーも二枚以上の分離被覆も生じない。
\item 同じ既約 $Y$ から $Y$ に戻る二つのこのような対応列があるとする。
長さを反復で揃えると、その二列は一般の始点からの各中間点まで一致する。
従って同じ $Y$ 上の相異なる支配的自己対応を二つ同時に使う反復分岐もない。
\end{enumerate}
\end{theorem}
\begin{proof}
上の余核の双対空間を $K(p,q)$ と書く。三文字テンソルがこの空間に属する条件は、
両内側の関係の零化と、その左右に $p,q$ を付けた二つの境界関係の零化である。
従って $p_0,\ldots,p_m$ の辺列について
\[
 p_0\otimes K(p_0,p_1)\otimes p_1\otimes\cdots
        \otimes K(p_{m-1},p_m)\otimes p_m
 \ \subset\ (A_{i,i+4m+1})^*
\]
となる。点の代表の倍率は独立性に影響しない。
一つの始点から長さ $m$ の二分木を作り、分かれる二辺の終点が違うか、
同じ終点でも $K$ の選んだテンソルが線型独立なら、その全ての葉の積テンソルは独立である。
実際、各辺を三文字テンソルと終点の計四文字で読むと、最初の分岐の二つの四文字は独立である。
それぞれの片方だけを消す線型形式で縮約し、残りの木に帰納する。
各頂点で一分岐だけの部分と、長さが等しい二本の異なる道を一単位として読む場合にも適用できる。
後者では中間点が違えばその点を置いた因子で区別できる。
全頂点を同時に一次独立に取る必要がない。

有限強連結成分が単純な巡回路でなければ、ある頂点に異なる二つの出辺がある。
それぞれに同じ頂点への帰路を付け、反復して正の長さを揃える。
二つの道の四文字ブロックの積テンソルは最初の出辺で区別できる。
任意の順に繰り返せば長さに対して指数個の独立な元ができ、三次のHilbert増大に反する。
次に有向道が $s$ 個の巡回成分を通るとする。相異なる成分間の固定した連結路を選び、
各巡回成分では退出頂点でその巡回路を回ってから次へ出る。
各巡回長の公倍数を $b$ とし、巡回部にさらに $b n_j$ 辺ずつを挿入する。
$\sum n_j=N$ で得た $\binom{N+s-1}{s-1}$ 個の道は同じ長さであり、二分岐の木をなす。
分かれる最初の頂点で退出辺と留まる辺が異なる。従って対応する積テンソルは独立であり、
$N$ に関する三次増大から $s\le4$ である。

支配的対応列では、所定の次側稠密開集合に入る点を各段で選べる。
もし一つの出射影の一般ファイバーが正次元または幾何的に二点以上なら、
どの有限反復回数 $m$ についても、一般始点からその段で毎周二点に分かれる木が作れる。
詳しくは帰納的に終点側の稠密開を引き戻し、対応の稠密開の像で始点側の稠密開を縮める。
選択肢が二つ必要な段では対応の自己ファイバー積から対角を除く。
支配性と一般ファイバーの条件によって、その二つの点をともに指定の稠密開に取れる。
各 $m$ で始点が変わっても、得られる $2^m$ 個の独立な元は同じ次数の経路空間に入る。
従って一般ファイバーは一点であり、標数零より双有理である。
逆向きにも同じ成長制限を使える。
一般の辺で $d\ge2$ の場合は中間点を動かす代わりに $K$ の二独立元を選べば同じ矛盾となる。

最後に二つの列を同じ長さにしたとき、初めて異なる中間点への二つの双有理写像は
稠密開集合上で異なる値を与える。両方の終点の写像はなお支配的なので、
一方か他方を反復ごとに選ぶ上記の二分木が構成できる。
\end{proof}

\paragraph{有限障害としての利用。}
例えば $p\ne q$ で $d(p,p),d(p,q),d(q,p)>0$ が同時に成り立つ点対は、
二つの戻り道による除外を与える。任意の固定した有限有向多重グラフを指定しても、
全頂点の相異条件と必要な次数の小行列式だけでその出現を射影消去で判定できる。
また相異なる点が一次独立な範囲なら、木を二分岐に限らず各 $K$ の基底を同時に使え、
$M=(d(p_a,p_b))$ として $\dim A_{i,i+4m+1}\ge\sum_{a,b}(M^m)_{ab}$ となる。
成長次数だけに基づく巡回成分数の上限4は改善できない。
四文字の非可換単項式代数で降順の隣接二文字を零とすれば、基底は弱増加語である。
座標四点の $M_{ab}$ は $a\le b$ で $\binom{b-a+3}{3}$、ほかは零となり、
四個のループを順に通過でき、上の下界は実際の三次Hilbert関数に一致する。
この単項式代数を標準AS例とする主張は用いていない。

% !TeX root = ../papers/nc-p3.tex
\section{大きな点ファイバーでの端支持と四文字ブロックの増大}
\label{sec:endpoint-plane-growth}
有限Frobenius窓の一次文字を $V_0,\ldots,V_3$ とし、$p\in V_3^*$ の三テンソル $w(p)$ の
端の支持を $V_0,V_2$ の部分空間とする。これらはそれぞれ始点・終点での隣接関係の評価像でもある。
辺の有限性、点列への延長可能性、または反対位置の対称性を仮定しない。
\begin{theorem}\label{thm:all-endpoint-ranks-two}
標準Koszul左右AS正則な四窓では、任意の非零一文字縮約の両端支持はともに次元2以上である。
従って全辺の全射影ファイバーは空、点、または射影直線であり、射影平面は現れない。
さらに二次元の端支持を与える点から、その支持平面への写像は各幾何学的ファイバーに高々一点を持つ。
\end{theorem}
\begin{lemma}\label{lem:rectangular-four-block}
$D\subset V_0$ の次元を $d$ とし、
\[
 k_D=\dim(R_0\cap(D\otimes V_1)),\qquad
 T=\{p\in V_3^*:w(p)\in D\otimes V_1\otimes V_2\},\quad l=\dim T
\]
とする。四位置を順に $V_0/D,V_1,V_2,T^*$ へ落とした商で、
長さ $4n$ の次元は $b^n$ であり
\[
 b\ \ge\ l(20-10d+4k_D-l)-k_D(10-4d+k_D).
\]
\end{lemma}
\begin{proof}
$R_3$ の像は零である。実際、$T$ での全評価が $D$ に入る。
従って四文字ごとの境界に関係はなく、一ブロックの次元を $b$ とすればよい。
第一の二文字の商は $4(4-d)-(6-k_D)=10-4d+k_D$ 次元である。
最初の三文字の周囲は $16(4-d)$ 次元、二つの挿入関係の次元は $4(6-k_D),6(4-d)$ である。
その交叉は元の三次重なりの射影を含み、これは $V_3^*/T$ に等しく $4-l$ 次元である。
よって三文字の商次元は $20-10d+4k_D-l$ 以上である。
$R_2$ の $V_2\otimes T^*$ への像は次元 $k_D$ 以下である。
実際、$w$ の第四文字を $T$ に制限すると、第一・第二文字側の全縮約は
$R_0\cap(D\otimes V_1)$ に入る。対応する隣接切断の他側の階数も同じである。
三文字商と $T^*$ の積から最後の関係像と最初の二文字商の積の次元上界を引けば結論を得る。
\end{proof}

\begin{proof}[一次元端支持の排除]
左端が $ku$ なら $w(p)=u\otimes M$ と書く。$M$ の階数を $r$、その左・右支持を $S,W$ とする。
$R_0$ は $u\otimes S$ を、$R_1$ は $M$ を含む。$r=1$ は連続する純関係により指数成長する。
$k=\dim(R_0\cap(ku\otimes V_1))$ とおく。$k=4$ なら $u$ が右単純のコピーを自由加群内に生成し、
ASの次数零Ext消滅に反する。$k=3$ なら $u\otimes S'\subset R_0$ となる三平面 $S'$ がある。
$R_1$ の $V_1/S'\otimes V_2$ への射影核は次元2以上であり、この核の各元に $u$ を左から付ければ
三次重なりとなる。従って補題の $d=1,k_D=3$ で $l\ge2$ である。
$2\le l\le4$ では $l(22-l)-27\ge13$ なので指数成長となる。

残りは $k=r=2$ としてよい。四位置を
$X=V_0/ku,V_1,Z=V_2/W,k_p$ に落とす。両端の $p$ に接する関係は零になる。
三文字 $X,V_1,Z$ の商の次元は、二つの関係像がそれぞれ4以下、5以下なので
$24-4\cdot2-3\cdot5=1$ に、それらの挿入の交叉次元を足した値以上である。
交叉は非零である。さもなくば全ての三次重なりが両端を落とした後で零となる。
すると任意の $0\ne z\in W^\perp\subset V_2^*$ で $w(z)$ の中央（位置0）の支持は $ku$ に入る。
非零の三テンソルで中央支持が一次元なら両側から同じ中央文字の純関係が得られ、
系\ref{cor:opposite-singular-rank-two}に反する。従って三文字商は二次元以上で、
$k_p$ により互いに独立なブロックを繰り返せば指数成長となる。
右端については語順を逆にする。
\end{proof}

\begin{proof}[二次元支持の一意性]
$d=2$ として $T\ne0$ とする。$w(p)$ の右端支持は次元2以上なので $k_D\ge2$ である。
一方、任意の $0\ne q\in D^\perp$ で $R_0(q)$ は次元2以上で、
この評価は $R_0/(R_0\cap(D\otimes V_1))$ を経由する。従って $k_D\le4$ である。
$2\le l\le4$, $2\le k_D\le4$ なら補題の下界
$l(4k_D-l)-k_D(2+k_D)$ は少なくとも4であり、標準成長に反する。
ゆえに $l=1$ である。
\end{proof}

\begin{corollary}\label{cor:endpoint-line-line-cut}
両端支持が二次元の $w(p)$ を持つ正則窓で、その支持を $D\subset V_0,W\subset V_2$ とすると、
$R_0\cap(D\otimes V_1)$ と $R_1\cap(V_1\otimes W)$ はともに二次元で、
ちょうど $w(p)$ の二端からの縮約像である。
両端を $D,W$ で割った全三次重なりの射影は高々一次元である。
\end{corollary}
\begin{proof}
二つの交叉次元を $k,h$ とすると各々2以上である。
$p$ を一点にし、他の位置を $V_0/D,V_1,V_2/W$ にしたとき、繰り返せる三文字商の次元は
$16-2(6-k)-2(6-h)+j=2(k+h)-8+j$ である。
$j$ は挿入した二つの関係の交叉次元で、全三次重なりの射影を含む。
この値は1以下なので $k=h=2$, $j\le1$ となる。
\end{proof}

% !TeX root = ../papers/nc-p3.tex
\section{直線ファイバーの中心は曲線以下である}
\label{sec:branching-curve-bound}
\begin{theorem}\label{thm:branching-centers-curve}
標準正則四窓の各辺、各射影について、直線ファイバーを持つ点の軌跡は次元1以下である。
従って一つの辺がある曲面の全点の上に直線を持つことはない。
純関係が一つもないという仮定は不要である。
\end{theorem}
\begin{proof}[一般平面における六関係の独立性]
$R\subset V\otimes W$ を辺の六関係とする。
どの $a\ne0$ に対しても $ka\otimes W$ に入る関係は高々二次元である。
三次元なら定理\ref{thm:all-endpoint-ranks-two}の証明で $k=3$ とした議論がそのまま指数成長を与え、
四次元なら右socleがある。逆の $V\otimes kb$ でも同じである。

一般の $a\in V$ については $R\cap(ka\otimes W)=0$ である。
反対なら $\mathbb P(R)$ 内の純関係の既約族の一つが $\mathbb P V$ に支配的となる。
第二の点 $b$ が一定なら $V\otimes kb\subset R$ で矛盾する。従ってその族の関数体で
$b$ の二座標の比 $f$ は非定数である。一座標を自明化した $a$ の四一次式の空間を $A$ とすると、
$\dim(A+fA)\ge7$ となる。実際、$f v=u$ と $f v'=u'$ から $u v'=u' v$ が
$\mathbb P V$ の斉次二次式の恒等式となる。$f$ が非定数なら $u,v$ は独立で互いに素なので
二対は比例する。よって $A\cap fA$ の次元は高々1である。
これら七つ以上の独立関数は $a\otimes b$ の座標の包に入り、六次元性に反する。
従って一般超平面 $a^\perp$ に制限すれば六関係は一次行列の列として独立のままである。

\end{proof}
\begin{proof}[二次・三次の曲面]
$r(p)=\dim R(p)$ と書く。$r(p)=2$ の被約軌跡上の像は階数2の部分束
$S\subset W\otimes\mathcal O$ であり、$S(1)$ は六関係で大域生成される。
$S^*$ の行列式の与えるPlücker写像は閉点上で単射である（定理\ref{thm:all-endpoint-ranks-two}）。
まず全三次小行列式が恒等零の場合は下の平面の場合に含める。
そうでなければ曲面成分は三次以下である。既約な二次または三次の曲面の一般平面切断を $C$ とすれば、
$C$ は整な平面曲線で、二次なら $\mathbb P^1$、三次なら特異な場合も算術種数1で $\omega_C=\mathcal O_C$ である。
六つの生成切断は $H^0(C,S(1))$ で独立である。$C$ は切断平面を張るから、
関係の一次行列が $C$ で消えればその平面で消え、前段が適用できる。
大域生成された階数2束の一般切断は曲線上で零点を持たない
（各点で二条件の入射対応の次元を数える）。従って
$h^0(S(1))\le1+h^0(\det S(1))$ である。$D=\det S^*$ とおけば
$\det(S(1))=\mathcal O_C(2)D^{-1}$ である。
二次なら $\deg D\ge1$ なので右辺は高々5となる。
整な平面三次なら $D$ は正次数であり、Riemann--Rochと双対性から $h^0(D)=\deg D$ である。
$\deg D=1$ は非定数性に反し、$\deg D=2$ は直線への二枚被覆で標数零の単射性に反する。
よって $\deg D\ge3$、$\mathcal O_C(2)D^{-1}$ の大域切断は高々3である。
いずれも六つの独立切断に足りない。
\end{proof}

\begin{proof}[平面の排除]
$r(p)=2$ が平面 $P=\mathbb P(a^\perp)$ 全体で成立すると仮定する。
$\det S^*=\mathcal O_P(e)$ とおく。$\det S(1)$ も大域生成され、単射なGrassmann写像より
$1\le e\le2$ である。$e=2$ では $S(1)$ は大域生成で行列式自明なので自明束である
（一点で独立な二切断の行列式は消えない定数）。
従って六関係のうち少なくとも四次元が $ka\otimes W$ に入るため不可能である。

$e=1$ ならPlücker写像は線型平面を $\operatorname{Gr}(2,W)$ に埋め込む。
そのような平面は全ての二空間が一定の三空間に入るものか、一定の直線を含むものである。
実際、分解可能な二ベクトルの一次結合が分解可能なら二直線は射影空間で交わる。
二本を固定すると、それらと交わる残る直線は同一平面に入るか、共通点を通る。既約性から二型を得る。
前者なら $S$ は $3\mathcal O\to\mathcal O(1)$ の核で、$h^0(S(1))=3$ である
（三個の独立な一次式の一次関係を数える）。
少なくとも三関係が $ka\otimes W$ に入り先の上界に反する。

後者の共通直線を $kb$ とする。$S$ は $\mathcal O$ と $\mathcal O(-1)$ の拡大で
$H^1(\mathcal O(1))=0$ から分裂し、$h^0(S(1))=4$ である。
$R$ の制限は四次元全部でなければ $ka\otimes W$ の関係が三つ以上になる。
従って $R$ 自体は非零の $a\otimes v$ を含み、二文字を
$V/ka,W/kb$ へ落とした後の $R$ の像は非退化な一つの双一次式 $\phi$ で張られる。
実際、$S/kb$ の与える $P\to\mathbb P(W/kb)$ は射影線型同型で、
対応する一次行列以外は全て $b$ に値を持つ。位置を順に0,1とすれば、両文字を
このように落とした四テンソルは $\phi\otimes C$、$0\ne C\in R_2$ である。
他の隣接切断ではこの制限の階数は $3\operatorname{rank} C\le6$ であり、$C$ の左支持 $U$ は高々二次元である。
任意の $y\in b^\perp\setminus0$, $z\in U^\perp$ で $w(y,z)\in R_3$ は
$V_0$ を $ka$ で割ると零なので $V_3\otimes ka$ に入る。
しかし $R_0$ の純関係 $a\otimes v$ と連続する純関係はない。従って $w(y,z)=0$ である。
ゆえに全ての $y\in\mathbb P(b^\perp)$ において $R_1(y)\subset U$ である。
その次元は少なくとも2なので $U$ に等しく、相異なる点は同じ端支持平面を持たないという定理に反する。
これで平面も全射影空間の場合も排除した。
\end{proof}

\begin{corollary}\label{cor:branching-curves-successor-surface}
直線ファイバーの中心が既約曲線をなす場合、それらの直線の合併の閉包は既約曲面である。
曲線の一般点では逆方向にも少なくとも一点が続く。前点を持たず直線の後続を持つ点は各位置に有限個しかない。
\end{corollary}
\begin{proof}
直線は中心から一意に決まり相異なるので、その射影直線束の像は二次元である。
曲線の一般点で逆方向の評価階数が4なら、辺の隣接切断の直交不等式により全ての
後続で順方向の評価階数は2以下となる。すると上の曲面全体が新しい直線ファイバーの中心に含まれて矛盾する。
逆方向の階数4は開条件で、各曲線に稠密にはなれないので、そのような点は有限個である。
\end{proof}

% !TeX root = ../papers/nc-p3.tex
\section{正次元の分岐中心は直線に限る}
\label{sec:branching-lines-only}
\begin{theorem}\label{thm:branching-lines-only}
標準正則四窓の各射影で、直線ファイバーを持つ点の被約軌跡は有限個の点と有限本の直線の和である。
各直線が次の文字に掃く曲面は平面、滑らかな二次曲面、または二次錐である。
滑らかな二次曲面を掃く場合には、その直線上の全点は同じ一つの前点を持つ。
またいずれかの辺の射影像が射影三空間全体なら、全窓は添字付き可換多項式環の一類で尽きる。
\end{theorem}
\begin{proof}[三次元像と曲面次数の上界]
一般点で関係の評価階数が3なら、その辺から後続点への有理写像がある。
一般ファイバーが正次元なら逆方向の直線ファイバーの中心は曲面以上を占める
（全ファイバーは高々直線）。前節に反するので像は三次元、逆も一般一点で双有理である。
座標を共通因子のない四つの独立な同次数式 $f_i$ で表す。
六関係より $f_i$ に全四一次式を掛けた包の次元は10以下である。
このとき $f_i$ は一次式である。実際、基底を消去して異なる四つの先頭単項式を取る。
その各四倍のうち、二単項式間の一致は高々一つである。従って相異なる積は $16-6=10$ 以上で、
等号の場合全対に一致があり三重の一致はない。四単項式は同じ単項式 $h$ に四文字を掛けたものである。
最初の二つを $h x_i,h x_j$ とすると、両方と一文字倍で会う第三のものは
$h$ の倍か $h x_i x_j$ の約数である（指数の差が一文字の交換であることによる）。
後者は三重一致になるからである。各対のS多項式は一つ上の次数で零に簡約できる。
この次数の先頭単項式は既に十個の倍で尽きるためである。
従って四式はGröbner基底であり、$h$ が非定数なら初期イデアルの高さは1である。
元のイデアルも高さ1なら四式に共通因子があり矛盾する。この証明は四次検査による延長を使わない。
従って双有理写像は射影線型で、そのグラフを消す六関係は交代関係全体に移される。
一般の零点対では隣の評価も階数3以上であり、直交不等式から3になるので、隣接辺にも同じ議論を適用できる。
全整数位置で順に基底を運ぶと全関係が可換関係となる。

以後この場合を除く。各辺で一般階数は4である。
四次小行列式の最大公約因子は次数2以下である。仮に次数4なら四行六列の評価行列の核二平面は一般に一定で、
非零関係が全ての点で消えてしまう。次数3なら、小行列式を割ったPlücker写像の像の閉包は線型空間である。
全核に共通の関係はないため定点や直線ではなく、前節と同じ交わる直線族の分類から
核は全て $R$ の一定の三平面 $T$ に入る。すると $T$ の評価は全点で階数1以下である。
一次行列の階数1以下の線型空間は共通像直線または共通余像直線を持つので、
三つの独立な関係が同じ末尾文字を持つか、非零関係が恒等零となる。どちらも不可能である。
従って各射影で全曲面成分の次数の和は2以下である。
\end{proof}

\begin{proof}[曲線と二項表示]
分岐中心の既約曲線を $C$ とする。その後続直線は相異なり既約曲面を掃き、
一般の曲面点には前点が一つしかない。曲面が平面なら、双対平面で後続直線を表した曲線の次数は1である。
実際、その次数が一般の平面点を通る直線の本数である。よって直線束の中心点を共有する。
二次錐の母線も頂点を共有する。共通点があれば、その前点は一つの直線に含まれるため $C$ 自身が直線である。
残りでは後続は滑らかな二次曲面の一方の母線族全部である。$C$ からそのパラメータ直線への射は
単射で、標数零なので有限双有理、標的が正規なので同型である。
すなわち $C=\mathbb P^1$ と書け、後続点は二項座標で $s\otimes t$、$s\in C$、$t$ は自由である。

非線型な $C$ を仮定する。平面に入れば滑らかな有理平面曲線なので二次曲線である。
非平面の場合に $C$ の文字の制限を $A\subset S^d$、$S^j=H^0(\mathbb P^1,\mathcal O(j))$ とする。
$A$ は基点自由な四空間である。この辺の六関係は、$t$ の二座標ごとに
$A\otimes S^1\to S^{d+1}$ の乗法核に入る。従ってこの核は少なくとも三次元であり、
$\dim S^1 A\le5$ である。任意の非零二項形式空間 $B$ で $\dim S^1B\ge\dim B+1$ である
（等次元なら二一次式の非定数比が有限次元空間を保ってしまう）。
等号なら $B=h S^{\dim B-1}$ である。実際、$x b=y b'$ の解を
$b=y c,b'=x c$ と書くと、その $c$ の空間は $\dim B-1$ 次元で、その一次式倍は全て $B$ に入り
同じ最小増加を持つ。次元で帰納すればよい。よって $A=S^3$ で $C$ は有理正規三次である。

この二つを $d=2,3$ とし、射影平面または三空間への文字制限は完全系 $S^d$ である。
位置を $0=C,1=(s,t),2,3$ とする。位置0を $C$ の線型包へ制限した関係 $R'_0$ は
\[
        [x,y]S^{d-1}(x)\otimes S^1(t)
\]
に等しい。$x,y$ は位置0と位置1の母線パラメータの二項変数であり、括弧は二行列式である。
$d=3$ では次元6で等号、$d=2$ ではこの包の次元は4で、平面を消す一文字 $a$ に始まる関係は高々二つなので
やはり等号である。後者では $a$ に始まる純関係があり $a$ に終わる純関係は禁じられる。
従って $R_3$ の位置0への同じ射影 $R'_3$ は六次元のままである。$d=3$ でも当然六次元である。
\end{proof}

\begin{proof}[前後に直線を持つ曲線]
$C$ の一般点の逆方向の評価階数は2か3である。2の場合は $C$ 全体で2となり、
逆順にも同じ曲面の議論が適用できる。$C$ を非線型としているので、前点側も滑らかな二次曲面の
一方の母線族である。その母線パラメータも $C$ からの射影同型で同定し、変数を $v$、
自由なもう一方の二項変数を $u$ とする。
制限された四テンソルは $R'_0,R'_3$ に入るので
\[
              [x,y][x,v]h(x,t,z,u),\qquad \deg_x h=d-2
\]
と書ける（$z$ は位置2の四文字）。二つの括弧は互いに素である。
$d=2$ なら $y,t,z\mid v,u,x$ の階数は高々 $2\cdot2=4$ である。
実際、二つの括弧のこの分割での階数は高々2、残る $h$ の $u$ 側も二次元である。
しかし同じ切断の像は六次元である。位置0だけを制限しており、$R_3$ の六次元像が保たれるからである。
$d=3$ では $s\in C$ を固定すると、その両端支持で文字を割ることは $y=v=s$ での評価である。
系\ref{cor:endpoint-line-line-cut}により残る $[x,s]^2h(x,t,z,u)$ の位置0対他の階数は高々1である。
従って $h$ は $x$ の固定一次式の倍であり、$S^3(x)$ の全支持に反する。
\end{proof}

\begin{proof}[一般前点が唯一の場合]
前点の評価階数が一般に3とする。制限四テンソルを
$[x,y]f(x,t,z,u)$ と書き、今の $u$ は位置3の四文字とする。
$u$ に値を持つ次数 $d-1$ の二項形式の空間 $A'$ を、$t,z$ による $f$ の全縮約の包とする。
$R'_3=S^1 A'$ なので $\dim S^1 A'=6$ である。一般の $s$ での $A'(s)$ の次元は
端支持の階数3である。$s$ の消滅一次式を掛けた $A'$ は評価核に入るので
$6\ge\dim A'+3$ から $\dim A'=3$ である。そこで
\[
               f=\sum_{j=1}^3 K_j(t,z) A_j(x,u)
\]
と分解すると、$K_j$ は独立で、$R_1=S^1(y)\otimes\langle K_j\rangle$ である。
実際、制限後の他側切断の像はこの六次元空間に入り、$R'_3$ の階数により等号である。
制限の前の $R_1$ も六次元だから変わらない。

位置0を $S^d$ に割った商で、位置1の二項二文字 $y,t$ を一旦二つの連続する位置と読む。
$R'_0$ は $x,y$ だけへの関係と自由な $t$ の積、$R_1$ は自由な $y$ と $t,z$ だけへの関係の積である。
従って $y$ と次の $t$ の間には関係がなく、一期を $t,z,u,x,y$ としたブロックが独立に繰り返される。
厳密には元の経路を $z$ から始め $x$ で終え、内部に必要回数の完全ブロックを入れればよい。
端の部分ブロックは完全ブロックに非零商があれば非零である。$d=3$ ではこの分離は元の代数そのもので、
完全ブロックの次元が0,1なら有界、2以上なら指数成長となり標準Hilbert関数に反する。

$d=2$ の商では次元が2以上であることを示す。まず $t,z$ の二文字の商は $8-3=5$ 次元である。
ブロックの右端 $x,y$ の商は乗法で $S^3$ になり、$u,x,y$ の商は
$(V_3\otimes S^3)/(S^2 A')$ なので $16-9=7$ 次元である。
実際、三つの $A_j$ は一般点で独立なため $S^2$ 係数の間にも関係はない。
$R_2$ の六つのうち四つは $A_j=x_0 A_{0j}+x_1 A_{1j}$、
$K_j=t_0 K_{0j}+t_1 K_{1j}$ として
$H_{bi}=\sum_j K_{bj}(z) A_{ij}(u)$ ($b,i=0,1$) である。
これらは制限四テンソルの位置0,1からの縮約で、$R'_0$ の階数4から独立である。

$5\cdot7$ 次元の積へ残る中央関係を挿入する。
$H_{bi}$ の全挿入の像は高々 $3\cdot5=15$ 次元である。
左側で $t_c K_{bj}$ の四つの添字組には全 $j$ で
$\sum_b t_b K_{bj}=0$ がある。右側で $A_{ij}\otimes p_i$、$p_i\in S^3$ の二組には
$(p_0,p_1)=(x_0g,x_1g)$, $g\in S^2$ という三次元の核が全 $j$ に共通である。
従ってこの挿入写像は左右で次元 $4-1$ と $8-3$ の商の積を経由する。
追加二関係の挿入は $2\cdot2\cdot4=16$ 次元以下で、完全ブロックは $35-15-16=4$ 次元以上である。
この商の指数増大も元の標準成長に反する。これで全非線型曲線を除いた。

最後に $C$ が直線で後続が滑らかな二次曲面の母線のとき、位置0の二文字への射影では
$R'_0=[x,y]S^1(t)$ である。包含はグラフでの消滅、等号は各点の端支持2による。
四テンソルは $[x,y]f(t,z,u)$ に制限されるので、位置0の任意の点で逆端の支持は
$f$ の $u$ 支持に等しく一定である。その次元は2以上、一般点で3以下であり、
2なら異なる点が同じ逆支持平面を持つ。従って3で、同じ唯一の前点が全直線に対して決まる。
\end{proof}

% !TeX root = ../papers/nc-p3.tex
\section{分岐中心の各直線が掃く曲面は平面である}
\label{sec:plane-branch-pencils}
\begin{theorem}\label{thm:plane-branch-pencils}
標準正則窓で、正方向の直線ファイバーの中心が一本の直線 $L$ を含むとする。
その後続は一つの平面内の直線束であり、束の中心 $q$ での逆方向ファイバーは $L$ 全体である。
特にどの曲線族も二次曲面の母線を掃くことはない。
さらに $L$ の線型包を点座標に残し、関係六平面を制限すると三つの独立な関係が残る。
点座標で $[x_0:x_1]\in L$、後続平面を $b=0$、束の中心を $b=v=w=0$ とすれば
その関係は $x_0b,x_1b,x_0v+x_1w$ の包である。
\end{theorem}
\begin{proof}[滑らかな母線族の反対側の分離]
位置0を $L$、後続位置1を滑らかな二次曲面 $\mathbb P Y^*\times\mathbb P T^*$ とし
$V_1=Y\otimes T$ とする。前定理の最後の計算より、位置0を $L$ に射影した四テンソルは
$[x,y]f(t,z,u)$、$f$ の $u$ 支持次元は3である。$z,u$ は位置2,3を表す。
従って $R_1=Y\otimes K$ である。実際、$K\subset T\otimes V_2$ は $u$ 縮約による三空間で、
他の隣接切断の階数は $2\cdot3=6$ なので全六関係が得られる。

一般の $t\in\mathbb P T^*$ で $K(t)$ は階数3である。
さもなくば滑らかな二次曲面上で全ての $s\otimes t$ が順方向の直線ファイバーを持ってしまう。
よってその核点 $z(t)$ は $s$ に依存せず、逆方向に $\mathbb P Y^*\times\{t\}$ 全体を持つ。
全ファイバーは高々直線なので、異なる $t$ で同じ $z(t)$ は生じない。
その像の閉包は逆方向の分岐中心の直線 $L'$ である。その母線族は先と同じ二次曲面の他方の族となる。
同じ二項切断を逆向きに適用すれば、$R_0=J\otimes T$ である。
ここで逆向きの二項基底を改めて選んでも、二つの直線族による $V_1=Y\otimes T$ の分解は
各二空間の基底変更を除いて一定である。実際、滑らかな二次曲面の順序付き両直線族のパラメータが
それぞれの射影直線を決め、超平面束はその二つの一次束の外積である。

これで位置1の $Y,T$ の間で分離が起きる。前の関係は $V_0,Y$ だけに関わり、
次の関係は $T,V_2$ だけに関わる。四移動ごとに得られる完全な経路ブロック
$T,V_2,V_3,V_0,Y$ の次元を $b_0$ とすれば、その反復の次元は $b_0$ の冪と
次数に依らない両端ブロックの次元の積である。これは標準の三次Hilbert増大にならない。
\end{proof}

\begin{proof}[二次錐の母線族]
錐の頂点を $q$ とし、$\ker q\subset V_1$ を二項二次形式 $S^2(y)$ に同定する。
$L$ から母線のパラメータ直線への写像は同型である。
実際、単射かつ全射で、標数零の射影直線間のそのような射は次数1である。
位置0の二変数を $x$ として、制限関係 $R'_0$ は $[x,y]S^1(y)$ に等しい。
全ての関係は $q$ で消え、底二次曲線の $y=x$ で消えるのでこの二空間に入り、
各 $x$ で評価の階数が2なので等号である。
従って制限テンソルは $[x,y]f(y,z,u)$ となる。$f$ は $y$ の一次であり、
各 $x$ 評価では非零一次式の乗法だけが加わるので、$u$ 支持は一定である。
前定理の滑らかな場合の末尾と同じくその次元は3である。

$u$ で縮約した $f$ の三次元像を $B\subset S^1(y)\otimes V_2$ とする。
$S^1(y)B\subset R_1$ である（$x,u$ での全縮約を取る）。
これは $\ker q\otimes V_2$ に入るため六次元未満である。
もし六次元なら $R_1$ 全体がその部分に入り、位置1の全支持に反する。
従って $S^1\otimes B\to S^2\otimes V_2$ の核に非零元があり、
一次形式の互いに素な二元を $y_0,y_1$ として
$y_0 b_0+y_1 b_1=0$ となる。各係数で割ればある $0\ne d\in V_2$ に対して
$y_0d,y_1d\in B$ である。すると $S^2(y)\otimes kd\subset R_1$ は
同じ末尾文字の三つの独立な純関係で、定理\ref{thm:branching-centers-curve}の初段に反する。
\end{proof}

\begin{proof}[平面内の表示]
残る平面では、一般点の前点一意性により束となることは定理\ref{thm:branching-lines-only}で見た。
$L$ から束への写像も射影同型なので、後続平面上の入射は一つの非退化な二変数対の式で切れる。
$L$ への制限関係は従って表示した三空間に入る。二次元しかなければ、その二関係の
全点評価は階数2であるから、像束を $S$ として $2\mathcal O\simeq S(1)$ となる。
しかし $S$ は定数直線 $kb$ と母線パラメータの $\mathcal O(-1)$ の拡大であり、
$\det S(1)$ は $\mathcal O(1)$ である。よって三次元全部を得る。
\end{proof}

\begin{corollary}\label{cor:rank-three-plane-branch-slice}
上の平面の場合、制限四テンソルは独立な $C_0,C_1,C_2\in R_2$ により
\[
           b(x_0 C_0+x_1 C_1)+(x_0 v+x_1 w)C_2
\]
である。$C_2$ は純関係である。$D=L^\perp\subset V_0$ について
$\dim(R_0\cap(D\otimes V_1))=3$ であり、
$\dim(R_3\cap(V_3\otimes D))\ge2$ である。
特に純関係を持たない正則窓では、全射影の直線ファイバーの中心は有限個である。
\end{corollary}
\begin{proof}
$R'_0$ の次元から表示と三独立性、交叉次元3が従う。他の隣接切断は $x_0,x_1$ の両成分に
$C_2$ の $V_3$ 支持を含むので階数は $2\operatorname{rank} C_2$ 以上である。
一方、$V_1$ 側の支持は $\langle b,v,w\rangle$ に入り、六次元全部ではなく階数は5以下である。
これは $R_3$ を位置0で射影した階数でもある。
実際には $q$ の逆ファイバーが $L$ 全体なので $w(q)$ の位置0の支持は $D$ である。
位置2からの縮約は少なくとも二次元であり、$R_3\cap(V_3\otimes D)$ に入る。
従って射影階数は4以下である。$C_2$ の階数が2なら、その射影像は
$\langle x_0,x_1\rangle\otimes \operatorname{supp}_{V_3} C_2$ 全体に等しく、
異なる $L$ の点の逆端支持が同じ二平面となって不可能である。
\end{proof}

% !TeX root = ../papers/nc-p3.tex
\section{二次曲面間の双有理対応と孤立ファイバーを許す曲面被覆}
\label{sec:birational-quadric-chains}
\begin{theorem}\label{thm:birational-quadric-chains}
標準正則窓の点対応に、全整数位置で整な二次曲面から次の整な二次曲面への支配的双有理グラフが含まれるとする。
するとこれらの双有理写像は全て、周囲の射影三空間の同型の制限である。
余分な点対や正次元の特殊ファイバーを初めから除く必要はない。
従ってこの全域の窓は既出の滑らかな二次曲面商または二次錐商のリストに入る。
\end{theorem}

\begin{proof}[一般超平面切断と母線束]
辺を $Q\dashrightarrow Q'$ とする。$Q$ の一般超平面切断は滑らかな二次曲線であり、
基点を避け、像曲線に双有理に写る。
像が平面に入れば次数は2以下である（整二次曲面の平面切断に入る）。
非平面ならこの二次曲線上の点座標を二項二次の完全系 $S^2$、相手を基点を除いた
四空間 $W\subset S^d$ とする。定理\ref{thm:branching-centers-curve}の一般平面の補題により
六つの双一次式は線型行列として独立のままであり、全て乗法 $S^2\otimes W\to S^{d+2}$ の核に入る。
従って $\dim S^2 W\le6$。二回の二項増加不等式はどちらも等号であり、
定理\ref{thm:branching-lines-only}の最小増加の議論から $W=S^3$、$d=3$ である。

有理写像の四切断から固定因子を除いた可動系を取る。
滑らかな $Q$ ではその型 $(a,b)$ は $a,b\ge1$, $a+b\le3$ である。
実際、その一般超平面切断との交叉数が先の像曲線の次数である。
錐では解消 $\mathbb F_2$ の例外切断 $E$ と母線 $F$ を使い、超平面を $E+2F$、
可動系を $aE+bF$ と書く。支配性で $a\ge1$、$E$ が固定因子でないことから $b\ge2a$、
また $b=(E+2F)(aE+bF)\le3$ である。
どちらの場合も射影同型でない写像は、少なくとも一方の母線族を相手の母線族へ射影同型で運ぶ。
実際、$Q'$ が滑らかなら二つの射影直線への可動鉛筆の和がこの可動系である。
これは二つの基点自由な有理座標比の積の四座標を取り、各比の固定因子を除いても
積系には新たな固定因子がないことによる。
適切な源母線への次数の和が1なので、一方の鉛筆はその母線で一定である。
その基底写像の次数は1である。一般点での一対一性と母線から相手母線への非定数性による。
$Q'$ が錐なら、頂点からの投影を通じる二項鉛筆を $P$ と書くと、相手の底三座標の系は
可動部分 $2P$ と有効固定部分の和になる。従って源の同じ母線で $2P$ の次数は高々1、
この鉛筆は母線上で一定である。錐がどちらかにある場合、その側の族は唯一である。
滑らかな両側で両族とも対応すれば各基底直線の同型の積で既に射影同型である。

対応する一つの母線の基底直線を同定する。各位置で母線上の一次切断を基底直線上の束 $E_i$ にまとめる。
滑らかな場合 $E_i=\mathcal O(1)^2$、錐では $E_i=\mathcal O(2)\oplus\mathcal O$ である。
どちらも全切断は四文字 $V_i$、評価 $V_i\otimes\mathcal O\to E_i$ の核 $N_i$ は $\mathcal O(-1)^2$ である。
後者は完全二次系の二つの一次関係からも分かる。
相手母線への写像は基礎体 $k(s)$ 上の射影直線の同型なので、
そのグラフを消す双一次式は非退化な一本である。その直線を
$E_i\otimes E_{i+1}$ で飽和して $J_i\simeq\mathcal O(j)$ と書く。
六関係を二位置で同じ基底 $s$ に沿って評価すれば $H^0(J_i)$ に入る。
全二文字の制限は $H^0(E_i\otimes E_{i+1})$ へ全射で、この空間は次元12である
（分裂の各対で二項形式の乗法が全射）。従って $j\ge1$ である。
一方、$J_i\subset E_i\otimes E_{i+1}$ の行列式は次数 $4-2j$ の非零切断なので $j\le2$。
$j=2$ ならこれは定数非零で、$E_i^*\otimes\mathcal O(2)\simeq E_{i+1}$ を与える。
階数2なのでこれは両 $E$ の同型でもある（行列式を固定すれば $E_i^*(2)\simeq E_i$）。
全切断と各母線上の核直線の対応により、曲面間の写像は射影線型同型に延びる。
$j=1$ なら、四次元の制限核全体が $R_i$ に入り、$R_i$ の像は $H^0(J_i)$ 全部である。
\end{proof}

\begin{proof}[母線を交換した二辺の三次元数]
射影同型でない二辺が一つの滑らかな中間曲面で異なる母線を用いることはない。
そのような二辺を $0,1$ とし、中間四文字を $V_1=B\otimes T$, $\dim B=\dim T=2$ とする。
第一辺は $B$ の基底直線、次の辺は $T$ の基底直線で $j=1$ の表示を持つ。
まず双方の四次元制限核だけを三文字で割ると
\[
          H^0(E_0(1))\otimes H^0(E_2(1))
\]
を得て次元は36である。実際、第一の核は $V_0\otimes B\to H^0(E_0(1))$ の核に
自由な $T$ を掛けたもの、第二も $T,V_2$ について同じで、この二商は別因子で取れる。
$J_0\hookrightarrow E_0\otimes\mathcal O(1)\otimes T$ の係数は
$P\in H^0(E_0)\otimes T$、$J_1$ の係数は $Q\in B\otimes H^0(E_2)$ である。
残る二関係の挿入は、第一辺では $P$ を $B\otimes V_2$ の任意の切断と掛ける像、
第二辺では $Q$ を $V_0\otimes T$ の任意の切断と掛ける像となる。
従って二つの像の次元は各々8以下で、共通に積 $P Q$ を含む。
この積は零でない。一般の二つの基点を選び、それぞれの一次形式因子を評価した
$P,Q$ の値は両方非零であり、残る二つの束因子の普通のテンソル積となるからである。
従って三文字商は $36-8-8+1=21$ 次元以上となり、標準値20に反する。
\end{proof}

\begin{proof}[同じ母線を使う隣接辺]
いま射影同型でない辺 $i$ を一つ取る。両隣が射影同型ならその母線をそのまま移せる。
非射影的な隣辺も、上の三次元数と錐での母線の一意性により、同じ母線を使う必要がある。
従って四つの連続位置 $i-1,i,i+1,i+2$ の基底直線を同定でき、両隣辺の一般母線間の写像は同型である。
標準四テンソルの中央二位置をこの共通母線上で制限すると、$J_i$ の包含と
\[
         H\in H^0\big(\mathcal O(1)\otimes V_{i-1}\otimes V_{i+2}\big)
\]
の積になる。実際 $w$ の中央二位置の縮約像は $R_i$ であり、その制限が $H^0(J_i)$ である。
一般の基底点 $s$ を固定し、右隣の点を同じ母線上で任意に与える。
その前点を位置 $i+1$ に取り、位置 $i$ の母線上では中央の非退化な双一次式を非零にできる。
隣接点対で $w$ は消えるため、$H(s)$ は指定した右隣点を消す。同様に左隣の母線全体を消す。
従って $H$ は $\mathcal O(1)\otimes N_{i-1}\otimes N_{i+2}$ の大域切断である
（一般点での包含は部分束への包含なので全点に延びる）。
ところがこの束は $\mathcal O(-1)^4$ で非零切断がない。
これは $R_i$ の制限が二次元であることに反する。
ゆえに全辺は射影同型であり、独立な基底変更で曲面座標環の反復への全射が得られる。
定理\ref{thm:smooth-quadric-quotients}または定理\ref{thm:quadric-cone-quotients}を適用する。
\end{proof}

\begin{corollary}\label{cor:pure-free-surface-complete}
どの関係空間にも純テンソルがない標準正則窓について、任意の辺の射影像が曲面以上なら、
全窓は可換多項式窓、正則正規一次直線の平面商、滑らかな二次曲面商、または整二次錐商の
既出の添字保持分類で尽きる。点対応の全射影が点ファイバーという仮定は不要である。
\end{corollary}
\begin{proof}
三次元像は定理\ref{thm:branching-lines-only}で済む。
純関係がないので全辺で直線ファイバーの中心は有限である。
面上の一般点は隣接辺で曲線に縮約できない。縮約する曲線の一般点には直線ファイバーが生じるからである。
従って与えた面の一般点の相手は再び面の一般点であり、同じ辺を戻れば一意である。
この二点で隣辺の評価階数は一般に3以上で、隣接切断の直交不等式によりともに3となる。
よってこの対応は前後に延び、面の間の双有理写像の列を得る。
面は既に平面または整二次曲面に限られる。

平面が一つあれば、その次も前も平面で写像は射影同型である。
実際、平面への辺の六関係の制限は純関係の不在で独立である。
双有理写像の像が全三空間を張れば、相手の四座標に源の三一次式を掛けた包は次元6以下になる。
しかし四つの異なる先頭単項式の三文字倍の異なるものは7個以上である。
$12-6=6$ 個しかなければ全対が異なる一文字倍で会うので、四単項式は同一の単項式と
四つの異なる文字の積となる。三文字しかなく不可能である。
相手も平面なら、さらに相手を三文字へ射影して失う関係は三次元以下で、独立な線型関係が三つ以上残る。
共通因子を除いた三座標を $f$ とする。独立な二つの一次関係は有理関数体上も独立である。
従属なら共通の定数関係の一次式倍となるが、像は平面を張る。二関係の外積は $f$ の多項式倍なので
$\deg f\le2$。二次なら各外積は $f$ の定数倍で、クラメル則により第三の関係も最初の二つの定数結合となり矛盾する。
従って写像は一次である。
従って平面対上の制限関係はちょうど三つの可換関係となる。

平面方程式 $a_i$ を取ると、この射影の三次元核は
$a_i\otimes V_{i+1}+V_i\otimes a_{i+1}$ に入る。
両混合三空間 $a_i\otimes(V_{i+1}/ka_{i+1})$ と $(V_i/ka_i)\otimes a_{i+1}$ のいずれへも同型である。
実際、一方への射影で零なら元の関係は純となる。
従って $a_i V_{i+1}=V_i a_{i+1}$、全位置で正規となり、商は三変数多項式窓である。
元と商のHilbert関数から両側の乗法は単射であり、定理\ref{thm:plane-quotient-unmarked}の分類に入る。
平面がなければ全位置で整二次曲面となり、前定理を適用できる。
\end{proof}

% !TeX root = ../papers/nc-p3.tex
\section{孤立した直線ファイバーを許す全曲面枝}
\label{sec:isolated-surface-fibers}
\begin{lemma}\label{lem:plane-transition-rigidity}
標準正則窓で、各位置の面の一般点を前後に対応させる双有理な整曲面の列があるとする。
一面が平面なら、全てが平面で全写像は射影同型である。別の成分や特殊ファイバーに仮定はない。
\end{lemma}
\begin{proof}
三次元射影は既に可換なので除く。各面は平面か整二次曲面である。
位置0の平面の方程式を $a$ とする。$a$ に始まる純関係がないときは $R_0$ の平面制限が六次元である。
前節の三文字倍と一次関係の証明により、次の面は平面で、対応は一次になる。
そうでない辺があれば、従って $a$ に始まる純関係があり、その直前の $R_{-1}$ に
$a$ に終わる純関係はない。今の議論を逆順に適用すれば、一つ手前は平面で平面0への対応は射影同型である。

前の平面と今の平面の一次式を同じ三空間 $X$ と同定する。
四テンソル $w$ のこの二位置を両平面に制限したものを $w'$ とする。残る二位置を $Y=V_1,Z=V_2$ と書く。
平面グラフ上で消えるので
\[
 w'\in\Lambda^2 X\otimes Y\otimes Z .
\]
しかも $(X,X)\mid(Y,Z)$ の階数は3である。実際、前の辺をまず平面0に射影した段階では
六関係は独立で、次に前の平面へ制限して失う次元は高々 $\dim X=3$ である。
またこの切断の $Y,Z$ 側の像は $R_1$ に入る。

体積形を一つ固定すると、各 $z\in Z^*$ の縮約は線型写像
$T_z:Y^*\to X^*$ を与える。次の面への双有理写像を $f$ とすると、
一般 $x\in\mathbb P(X^*)$ に対して $T_z f(x)$ は $x$ に比例する（零を含む）。
実際 $w'$ は $x,f(x)$ を代入すれば前位置の任意の値に対して消え、三次の交代行列は
その対応するベクトルを核に持つ。像面が平面のとき、ある $T_z$ がそこに非零に制限されれば、
逆写像はこの線型写像であり、平面間の射影同型である。全ての $T_z$ が像平面上で零なら、
$w'$ の $Y$ 支持は平面方程式一本で、切断の階数3により同じ最初の文字の純関係三つを得て矛盾する。

像面が整二次曲面 $Q$ なら、非零の $T_z$ は全て階数3で、その射影化は $Q$ 上で同じ逆写像を与える。
このような線型写像 $T$ は定数倍を除き一つしかない。実際、一つを座標 $y=(u,t)$ で $T y=u$ とし、
他を $T'y=A u+t b$ とする。$u\times(Au+t b)$ の各座標は $Q$ の二次方程式の定数倍である。
それらと $u$ の内積が恒等零であるため、その定数ベクトルは零である。
$t$ の係数から $b=0$、全ての $u$ が $A$ の固有ベクトルなので $A$ はスカラーである。
従って $w'$ の $Z$ 支持が一本となり、今度は同じ末尾文字の三純関係を生じる。
これは定理\ref{thm:branching-centers-curve}に反する。以上を各位置と逆向きにも適用する。
\end{proof}

\begin{lemma}\label{lem:plane-koszul-isolated-exact}
平面の射影同型の四辺続きに沿った標準Koszul複体が、平面の有限個の点を除いて完全なら、
その平面上の層複体は全点で局所分裂完全である。
\end{lemma}
\begin{proof}
各平面の座標を同一に取ると、複体は
\[
 0\longrightarrow\mathcal O(-4)\longrightarrow 4\mathcal O(-3)
 \longrightarrow 6\mathcal O(-2)\longrightarrow 4\mathcal O(-1)\longrightarrow\mathcal O\longrightarrow0
\]
であり、両端の写像は点の非零評価から部分束への単射、全射となる。
最初の余核を $U$、最後の核を $W$ として $U\to6\mathcal O(-2)\to W$ を得る。
後半の核 $K$ は局所自由である。各閉点の二次元正則局所環で後半の像は自由加群の部分加群なので
深さ1以上であり、深さ補題から核は深さ2、従って自由である。
$U\to K$ は同階数で有限点を除いて同型である。その行列式の零点があれば因子を含むから、全点で同型である。
残る唯一のホモロジーは右の有限長余核 $D$ である。一方で捩り後のEuler標数の交代和は
平面の二次Hilbert多項式の四階差分で零なので $\operatorname{length}D=0$ である。
これで完全であり、各核が束であるためファイバーでも完全となる。
\end{proof}

\begin{theorem}\label{thm:isolated-surface-fibers}
標準正則窓の全射影で、直線ファイバーの中心が有限個であるとする。
いずれかの射影像が二次元以上なら、窓は可換多項式、正規一次直線列による平面商、
滑らかな二次曲面商、整二次錐商の既出の添字保持分類に入る。純関係の有無は仮定しない。
\end{theorem}
\begin{proof}
三次元像を除き、系\ref{cor:pure-free-surface-complete}の冒頭と同様に、
面の一般点の相手は再び面の一般点で、一点対の直交不等式から前後に延長できる。
ここで必要なのは両方向の直線ファイバーの中心が有限であることだけである。
平面がなければ定理\ref{thm:birational-quadric-chains}を使う。
平面があれば補題\ref{lem:plane-transition-rigidity}により、全位置の平面商の座標を三変数多項式の反復に取れる。
この商をテンソルした単純加群の線型解消を平面上に層化する。中央二写像は隣接関係の点評価と相補関係の
点評価の転置であり、各平面で有限点以外は両方階数3である。従って補題\ref{lem:plane-koszul-isolated-exact}
により平面上で完全である。定理\ref{thm:plane-surface-normal}の有限長ホモロジーと両側極小解消の議論を
そのまま適用すると、一次核は全位置で正則正規となり、平面商の分類を得る。
\end{proof}

% !TeX root = ../papers/nc-p3.tex
\section{可動平面束の排除と全曲面射影の被覆}
\label{sec:exclude-moving-plane-pencils}
\begin{theorem}\label{thm:no-moving-plane-pencils}
標準正則窓では各辺の両射影で直線ファイバーの中心は有限個である。
従って一辺でも射影像が二次元以上となる全窓は定理\ref{thm:isolated-surface-fibers}の分類で尽きる。
\end{theorem}
\begin{proof}[二つの平面束]
可動中心の直線 $L$ を位置0に取り、位置1の後続平面の方程式を $b$ とする。
系\ref{cor:rank-three-plane-branch-slice}の表示を使い、反対辺 $R_2$ の純関係を $c\otimes d$ とする。
$R_1$ は $c$ に終わる純関係を持たない。従って $V_2/kc$ へ射影した六関係は独立で、
さらに $V_1/kb$ へ射影しても、失う次元は高々3である。
一方 $w$ を $L$ に制限した後の $R_1$ 側の縮約像は、表示式より
$b\otimes V_2+ V_1\otimes c$ に入るので全て失われる。
ゆえにこの縮約像、すなわち逆辺 $R_3$ の $L$ への制限の階数は高々3である。
ここには $d\otimes H^0(\mathcal O_L(1))$ が入る（独立な二つの動く係数を縮約する）。
各 $x\in L$ で評価した空間は次元2以上であり、従ってちょうど2である。
逆辺にも動く直線束があり、定理\ref{thm:all-endpoint-ranks-two}による支持の一意性と
平面束の分類から、その平面方程式は共通の文字 $d$ である。

$x=(x_0,x_1)$ として順方向の支持を $\langle b,l(x)\rangle$、逆方向を
$\langle d,m(x)\rangle$ とする。$l,m$ の二係数はそれぞれ平面方程式を法として独立である。
厳密四周期表示で $w$ の $L$ への制限は、位置1,2,3の順に
\[
 b\otimes A(x)\otimes d+
 l(x)\otimes c\otimes d+
 b\otimes e\otimes m(x)
\]
と書ける。実際、各支持束は定数直線と $\mathcal O(-1)$ の直和で、
二つの動く因子を同時に使う項は $\mathcal O(-2)(1)$ の切断なので零である。
$c,e$ は非零である（各点の両端階数は2）。$R_1$ は $b\otimes e$ を含むため
$R_0$ に $b$ に終わる純関係はなく、$R_2$ に $e$ に始まる純関係もない。

$R_2$ で $d$ に終わる関係は正確に一本である。実際 $V_2/ke$ を先に取り次に $V_3/kd$ を取ると、
初めは六次元で次も三次元以上である。逆に $d$ を先に消すと、今の表示を位置0,1で縮約して
得る $e\otimes m(x)$ の二係数が射影像にあり、次に $e$ を消すと少なくとも二次元を失う。
従って $d$ だけを消した像は五次元以上で、$c\otimes d$ を失うから正確に五次元である。
\end{proof}

\begin{proof}[二つの独立な束パラメータを持つ三文字]
位置0の相続く二回の訪問で $x,x'\in L$ を別々に指定する。
その間の位置1と3の残る文字はそれぞれ
\[
 E_x=V_1/\langle b,l(x)\rangle,\qquad F_{x'}=V_3/\langle d,m(x')\rangle .
\]
束として $E,F\simeq\mathcal O\oplus\mathcal O(1)$ である。
$R_1$ の $E\otimes V_2$ への像束を $I$、$R_2$ の $V_2\otimes F$ への像束を $J$ とする。
系\ref{cor:endpoint-line-line-cut}により両方の核は各点で二次元であり、
表示式から第一の核は $b\otimes e$ と $b\otimes A(x)+l(x)\otimes c$ で張られる。
第二の核も同様に定数一本と一次の一本である。よって $I,J$ は階数4、次数1である。
$L\times L$ 上で記号の左右の束を各射影で引き戻す。挿入写像
\[
 D:(I\boxtimes F)\oplus(E\boxtimes J)
 \longrightarrow E\boxtimes F\otimes V_2
\]
では三文字を元の順に並べる。両側の階数は16、行列式は
$\mathcal O(2,2)$ の切断である。例えば第一因子での次数は標的8、始域 $2+4=6$ である。
余核は端点を指定した三文字の商であり、対角上の各点でその次元は高々1である
（同じ点を四文字ごとに繰り返したときの独立ブロックの増大による）。

対角上の一般点で余核は非零である。反対に一般点で全射なら、そこでは二つの挿入の交叉がないため、
全ての三次重なりは両端を $E_x,F_x$ に落とすと零になる。
$Y=V_1/kb,\ U=V_3/kd$ と書く。$Y\otimes U\to E_x\otimes F_x$ の全 $x$ に共通の核は
$l,m$ のそれぞれ二係数の張る二空間どうしに入る。実際、この二空間の補直線を取り、
補直線どうしの成分と片側だけ補直線の成分は商の対応する直和成分を全 $x$ で評価すると検出できる。
従って $w$ を $Y,U$ に射影したものの $Y$ 支持は高々二次元である。
さらにその位置0支持は $L^\perp$ に入る（先の三項の表示は $b,d$ を消せば零）。
すると位置0,1対位置2,3の階数は高々4である。
しかし先に $b$ を消すと $R_0$ 側に核はなく階数6であり、続いて $d$ を消すと
$R_2$ 側に失うのは一本だけなので階数5である。これは矛盾する。
従って対角全点で $D$ の階数は15である。
\end{proof}

\begin{proof}[二重対応に沿った指数増大]
行列式が恒等零なら $L$ の異なる二点の任意の順序対で三文字商が非零である。
各対ごとにその非零な一次元商を取り、四文字ごとにこの二点を任意に並べる。
各境界で次に選んだ二点は一次式で分離できるので、長さ $4n+1$ で少なくとも $2^n$ の次元を得る
（以下の経路評価の被約二点版でもある）。

行列式が恒等零でなければ、零因子は対角と型 $(1,1)$ の残余因子の和である。
第一射影の垂直成分を取り除いた有効因子を $C\subset L\times L$ とする。
第一射影は有限平坦次数2である。各ファイバーは射影直線上の非零な二次式の零スキームだからである。
$C$ は二重対角、対角と自己射影同型のグラフ、または対角と定数グラフの和である。
最後の二型では共通成分がないものとする。行列式の単純成分の一般点で $D$ は階数15である。
階数14以下なら行列式はその成分の局所方程式の二乗で割れるからである。
また二重対角を含め、対角上では既に階数15である。

各 $n$ で一般初期点 $x_0$ を一つ固定し、$(x_{j-1},x_j)\in C$ となる
長さ $n$ のスキームを $T_n$ とする。有限平坦性を逐次引き戻すと長さは $2^n$ である。
全ての経路の全辺で $D$ は階数15の開集合に入るようにできる。
自己同型の場合には有限回の合成で有限例外集合を避ければよく、
定数グラフで一度定数点に入った後はその対角点しか通らない。
その開集合上では行列式が唯一の残る関係で、$D$ の余核の $C$ への引戻しは可逆層である。
これは可逆な15次小行列を使う行・列消去から、非被約の場合にも成立する。

各辺で三文字をこの可逆商へ評価し、境界の位置0の文字を $x_j$ で $\mathcal O(1)$ へ評価する。
得られる積評価は元の全二次関係を消す。内部は $I,J$ を割ってあり、
境界の $R_0,R_3$ はそれぞれ $E_{x_j},F_{x_j}$ の定義で零である。
各内部三文字の線型結合を、有限個の支持点全部で非零になるように選べる
（余核は三文字で生成され、体は無限）。これらと初期点の非零評価を固定すると可逆因子を掛けるだけとなる。
残る $n$ 個の境界の二一次式の積は $T_n$ 上の対応する可逆層の全切断を張る。
実際、各 $T_j\to T_{j-1}$ の各ファイバーは射影直線の長さ2の部分スキームで、
二一次式の制限はその $\mathcal O(1)$ の切断への同型である。
有限平坦性と中山の補題により直像束はこの二切断を基底に持ち、射影公式で帰納できる。
従って長さ $4n+1$ の経路空間の次元は少なくとも $2^n$ である。
この下界は二重点でも二つの異なる幾何学的後続を要しない。標準Hilbert関数の三次増大と矛盾する。
これで可動直線束は全て排除された。
\end{proof}

% !TeX root = ../papers/nc-p3.tex
\section{孤立分岐点の二側切断と後続の有理正規曲線}
\label{sec:isolated-branch-blocks}
\begin{proposition}\label{prop:isolated-branch-blocks}
位置0の点 $p$ で順方向支持 $S\subset V_1$ が二次元とする。逆支持を $T\subset V_3$、$t=\dim T$ とする。
$t=2,3$ に限られ、$R_1\cap(S\otimes V_2)$ は $w(p)$ の末尾縮約像に等しく次元 $t$ である。
後続直線 $L=\mathbb P(S^\perp)$ の一般点からさらに点が続く場合、それは一定点、
または $L$ と射影同型にパラメータ付けられる直線、滑らかな二次曲線、有理正規三次曲線である。
$t=3$ なら一般点に後続があり、$L$ 上で順方向階数が落ちるスキーム
（三次小行列式の共通零点）の長さは $3-d$ である。ここで $d$ はこのパラメータ表示の次数である。
$t=2$ で一般点が後続を持たないときは、後続を持つ点が二項四次の零因子をなす。
\end{proposition}
\begin{proof}
$t=4$ なら辺 $(p,q)$ の直交不等式により全 $q\in L$ で順方向支持が二次元となり
定理\ref{thm:no-moving-plane-pencils}に反する。
交叉次元は $t$ 以上である。一方 $L$ 上で残る二次式は一般点で階数3以上なので、交叉次元は3以下である。
$t=2$ のときは系\ref{cor:endpoint-line-line-cut}でちょうど2である。

一般に階数3のときは、$L=\mathbb P^1$ 上の小行列式が定める核点への有理写像を延長し、
像の文字の張る空間を $W\subset H^0(\mathcal O_L(d))$ とする。共通因子を除いてあるので基点はない。
$W$ の評価核を $N$ とする。位置2の全四文字から $W$ に落としても、
二次関係は $N(1)$ の一般ファイバーを張る。もとの点評価が後続点の全消滅超平面だからである。
$N$ の分裂に自明または正次数の因子はない（全切断の線型独立性による）。
また一般点での生成性により $-2$ 以下の因子もない。従って
$N=\mathcal O(-1)^{\dim W-1}$ で $d=\dim W-1$、$W$ は完全系である。
$t=3$ では $L$ に制限した三式の評価が一般に階数3であり、核座標は三次小行列式だけで書ける。
その共通因子の次数は $3-d$ である。$t=2$ で一般階数4の場合、制限した四式の行列式は非零な四次である。
\end{proof}

\begin{proposition}\label{prop:two-sided-center-self-block}
$t=2$ の点では $d(p,p)=1$（定理\ref{thm:recurrent-point-correspondences}の三文字商）であり、
$X=V_1/S$, $F=V_3/T$ に制限した四テンソルは $a\otimes h$、
$0\ne a\in\ker p$, $0\ne h\in X\otimes V_2\otimes F$ となる。
前後どちらの隣接直線にも、その辺で逆向きのファイバーが直線になる点が存在する。
より正確には、$H=a^\perp\subset V_0^*$ とすると、順側で $H$ 内に $p$ 以外の前点を持つ軌跡は
非零二項二次式の零因子の台であり、逆側でも同様である。
\end{proposition}
\begin{proof}
制限した $w$ の位置0対他の階数は高々1である。零なら任意の $u\in F^*$ について
$w(u)$ は位置1で $S$ に支持される。従って $V_0$ からの縮約像は
$R_1\cap(S\otimes V_2)$ に入り、高々二次元である。すると全 $\mathbb P(F^*)$ で
位置0の端支持が二次元となり可動中心が生じる。よって階数は1で、表示の $a$ は $w(p)$ の制限が零なので $\ker p$ に入る。
位置0からのこの非零縮約は二つの内側関係の挿入の交叉に入る。
両挿入は各八次元、周囲は十六次元なので交叉と余核は同じ次元であり、
反復ブロックの多項式増大で1以下となる。従って両者とも一次元である。

$q\in\mathbb P(X^*)$ について $H$ を位置0、$q$ を位置1で縮約すれば $F$ 側への像が零となる。
従って値は $R_2\cap(V_2\otimes T)$ の二空間に入り、常に $p$ が核にある。
これは $H/kp$ からこの二空間への $q$ に線型な二次正方行列である。
一般の $q$ では前点は $p$ 一つなので一般に可逆であり、行列式は非零二次である。
その根では $H$ 内の前点が少なくとも直線をなし、射影平面ファイバーはないので直線になる。
語順を反転すれば逆側での主張を得る。
\end{proof}

% !TeX root = ../papers/nc-p3.tex
\section{孤立した両側直線ファイバーの一次変形による排除}
\label{sec:exclude-two-sided-centers}
\begin{theorem}\label{thm:no-two-sided-centers}
標準正則窓で前点も後続点も直線をなす点は存在しない。
従って直線ファイバーを持つ点には逆側にただ一つの点があり、その逆側評価階数は3である。
\end{theorem}
\begin{proof}[共通の階数二の一次変形]
前節の両側点 $p$ を仮定し、$X,F,S,T,a,H$ を同じ記号とする。
$p$ を含む射影平面 $\mathbb P(H)$ 内での一次近傍を $P$ とする。その台は一点、
局所環は $k[\xi,\eta]/(\xi,\eta)^2$ である。$p$ の任意のこの一次変形について、
順・逆支持は同時に階数2の直和因子に持ち上がる。

実際 $v\in H$ について $w(v)$ の $X,F$ 両側への射影は $v(a)h=0$ である。
$X$ 側だけへの射影は $X\otimes(R_2\cap(V_2\otimes T))$ に入る。
括弧内は $w(p)$ の位置1からの縮約で張られる二空間なので、この射影は
一意な $B_v:S\to X$ を $w(p)$ の位置1に作用させたものである。
同様に $F$ 側だけの射影には $C_v:T\to F$ がある。補空間を固定して持ち上げれば、
$p+\epsilon v$ での順支持は $s\mapsto s+\epsilon B_v(s)$ のグラフ、逆支持も
$t\mapsto t+\epsilon C_v(t)$ のグラフである。
包含は表示式から、等号は閉点での階数と中山の補題から従う。
二方向の一次パラメータの全積が零なので、これらは $P$ 上の支持部分束 $S_P,T_P$ を与える。

商束を $E=V_1/S_P,F'=V_3/T_P$ とする。
$R_1\to E\otimes V_2$, $R_2\to V_2\otimes F'$ の像 $I,J$ は階数4の直和因子である。
例えば第一の核には、変形した $w$ を最後から縮約した二つの独立元があり、
閉点でこれらは核全体を張る。階数4の可逆小行列とこの核を使うと写像を
四次の単位行列と零の和に消去できる。第二も同じである。
\end{proof}

\begin{proof}[始点と終点を固定した十二文字]
$P\times P$ 上の独立な両境界について、三文字挿入の正方行列
\[
 (I\boxtimes F')\oplus(E\boxtimes J)\longrightarrow E\boxtimes F'\otimes V_2
\]
を取る。前節により閉点で階数15である。可逆小行列による消去で、その余核は一つの関数 $f$
の剰余環として表示でき、$f(p,p)=0$ である。特に双対数上で
$(p+\epsilon u,p+\epsilon v)$ に沿って可逆な一文字分の商が得られる条件は
$L(u)+M(v)=0$ の一一次式である。$L,M$ は二次元接空間上の線型形式で、零でもよい。

今、連続する四つの境界で $p, p+\epsilon u,p+\epsilon v,p$ と置く。
三つの三文字商を全て可逆とする条件は
\[
           M(u)=0,\qquad L(u)+M(v)=0,\qquad L(v)=0 .
\]
未知数四個に対する三本の斉次式なので $(u,v)\ne(0,0)$ を取れる。
各三文字では可逆な余核への像が単元となる一元を選ぶ。
境界の一次式をその点で評価し、これらを順に掛けると、
両端を元の $p$ に固定した経路の評価が得られる。内部関係は余核の定義で、境界関係は
$S_P,T_P$ の定義で零である。非零方向のある中間境界では二つの文字の評価が
$k[\epsilon]/\epsilon^2$ の全体を張る（定数と非零一次項を独立に取れる）。
他の境界では単元評価を固定すればよい。

この十二文字の区間を同じ固定点でつなぎ、各区間に独立な双対数 $\epsilon_j$ を与える。
元の窓の $12n+1$ 文字を $\bigotimes_{j=1}^n k[\epsilon_j]/\epsilon_j^2$ に評価する写像は
全関係を消し、上の二つの選択の積により全射である。従って次元は $2^n$ 以上となり、
標準Hilbert関数に反する。逆側が空でないことは命題\ref{prop:isolated-branch-blocks}による。
\end{proof}

% !TeX root = ../papers/nc-p3.tex
\section{片側分岐を結ぶ有限な有理曲線列と次数公式}
\label{sec:finite-branch-corridors}
\begin{theorem}\label{thm:finite-branch-corridors}
全射影像が高々一次元の標準正則窓を考える。後続直線を持つ点 $p$ を一つ固定する。
その直線のどの点でも、その辺における前点は $p$ 一つである。
また $p$ から出発する有理曲線列は、有理的な一般点の後続を一意に取ることにより
\[
 x_0=p,\quad x_1,\ldots,x_m:\mathbb P^1\longrightarrow\mathbb P^3,
 \quad x_{m+1}=z
\]
となる。中間の射は全て像への双有理射で、$x_1,x_m$ の像は直線、$m\ge2$ である。
$z$ は前点の直線を持ち、反対向きの曲線列はこの列を逆にたどる。
従って四周期で同定した順側分岐中心と逆側分岐中心はこの列で一対一に対応する。
\end{theorem}
\begin{proof}
前節により $p$ の逆支持は三次元である。後続直線上の二項変数を $x$ とすると、
$R_1$ のその直線への射影は三次元である。その基底を $K_j(x)\in V_2$ として
$w(x)=\sum_{j=1}^3 K_j(x)\otimes A_j$ と書く（位置を巡回して読む）。
一般点での $K_j(x)$ は独立で、$A_j$ の位置0支持の和は $\ker p$ である。
これはその点の逆支持が、可動中心の不在により一般には三次元だからである。
従って $K_j(x)$ が独立な全ての点で逆支持は三次元である。
従属な点の順支持は二次元であり、定理\ref{thm:no-two-sided-centers}によってやはり逆支持は三次元である。

一般点に先行する点の逆支持が3のとき、直交不等式から一般点の順支持は3以下である。
動く曲線の全点が階数2にはならないので一般に3で、一意な後続への有理写像を得る。
これは射影直線上で射に延びる。像が動けば、その一般点も一意な前点を持つので
標数零より曲線間の写像は双有理である。像が一定なら直前の曲線はその前点の直線全体となる。
その正規化パラメータはその直線との同型である。
いずれの動く像も全射影像の既約成分であり、位置を法4で見れば有限個しかない。
成分が二度現れると、一意な一般前点を順に復元して最初の直線を二度目にも得るが、
初回には前点が一定、二度目には変動するので矛盾する。従って有限回で一定像に達する。
逆の構成も一意なので最後の対応は全単射である。$m=1$ の排除は以下の次数式で示す。
\end{proof}

\begin{proposition}\label{prop:branch-corridor-degree}
$d_i=\deg x_i^*\mathcal O(1)$ とする。$x_{-1},x_{m+2}$ を両端の外の一意な定点とし、
動かない位置の次数を0とする。$1\le i\le m$ で逆・順の階数3評価の三次小行列式の
共通零因子の次数をそれぞれ $b_i,f_i$ とする。$b_0=f_{m+1}=0$ と置く。このとき
\[
 d_{i+2}-3d_{i+1}+3d_i-d_{i-1}= b_i-f_{i+1}\qquad(0\le i\le m).
\]
$b_1=f_m=0$ であり、
\[
 \sum b_i=\sum f_i\ge2,\qquad
 \sum_{i=1}^m\big(i b_i-(i-1)f_i\big)=m+2 .
\]
特に順側分岐中心が一つでもあれば、順側分岐中心を通る実際の周期点列がある。
\end{proposition}
\begin{proof}
辺 $x_i,x_{i+1}$ に沿い、$w$ による隣接六空間の完全対合を使う。
$x_i$ を評価して得る逆側の三空間は、列として $V_{i-1}^*$ から定数六空間へ入る
$d_i$ 次の行列の像であり、有理的な核直線は $x_{i-1}$ である。
従って飽和した像束の双対行列式の次数は $3d_i-d_{i-1}-b_i$ である。
実際 $V_{i-1}^*\otimes\mathcal O(-d_i)$ からの写像の核は
$\mathcal O(-d_i-d_{i-1})$、飽和で増える次数は局所単因子の和、すなわち三次小行列式の共通零因子の長さである。
隣の $x_{i+1}$ の順側評価でも次数は $3d_{i+1}-d_{i+2}-f_{i+1}$ である。
一般点で二つは直交補なので、その飽和三空間へのGrassmann写像は双対同型で対応し、Plücker次数が等しい。
最初・最後の定点で用いた側は階数3一定なので同じ計算が適用できる。
最初の直線上の逆支持は全点で3であるから $b_1=0$、逆に $f_m=0$ である。
$m=1$ なら最初の次数式から $f_1=3$ となり矛盾する。

次数式を加えれば、二端で二階差分がともに1なので総和は零である。
$i$ 倍して加えると、端の傾きが $+1,-1$ であることから表示した $m+2$ となる。
共通総和が1以下なら左辺は高々 $m$ であり矛盾する。
従って各列の内部で少なくとも一つ順側分岐中心に達する。実際のパラメータ値でその点まで進んでから
その中心の後続直線で同じことを繰り返せる。四位置合わせて中心は有限個なので、
ある中心を再び訪れ、間の有限経路を繰り返せば周期点列を得る。
\end{proof}

% !TeX root = ../papers/nc-p3.tex
\section{分岐を通る短い周期点加群と無限生成イデアル}
\label{sec:periodic-branch-obstruction}
\begin{theorem}\label{thm:periodic-branch-obstruction}
曲面射影のない標準正則窓に片側直線ファイバーがあれば、四周期を法として
順側分岐中心は高々20個であり、そのいずれかを通る長さ40以下の周期点列がある。
逆側でも同じである。その右周期点加群と逆向きの左周期点加群の表示イデアルは
それぞれ無限生成である。従ってこの残余枝が実在すれば右・左の両方でnoetherian性に失敗する。
\end{theorem}
\begin{proof}
各辺で六双一次式の零点集合の次元は高々1である。実際、曲線上で動く直線ファイバーはなく、
各孤立中心上のファイバーも高々直線である。この線型系から一般の五式を取れば
$\mathbb P^3\times\mathbb P^3$ で曲線の完全交叉となる。各段で正次元の余分な成分を次の式で切れ、
基底集合が次元1以下なので余次元5まで正則列になるからである。
完全交叉の類 $(H_1+H_2)^5$ により二つの射影方向の次数は10ずつである。
全ての点対応の既約曲線成分はこの交叉の成分なので、両次数の和の総計は一辺で20以下である。
定理\ref{thm:finite-branch-corridors}の各列は、辺ごとに異なる曲線成分を用いる。
一般前点を逆にたどって最初の順側中心を復元できるので、異なる列の間でも重複しない。
列の辺の両次数の和は $2\sum_{j=1}^m d_j$ であるから、
四位置の全列について $\sum m\le\sum d_j\le40$ である。各列は $m\ge2$ なので中心数の上界を得る。
各中心からその列の内部で到達できる順側中心を一つ選ぶ。この有限有向グラフの単純巡回路に沿って
元の点経路を連結すると、その長さは巡回路上の各列の $m$ の和以下である。
この長さは四位置を元に戻す。左向きの結論は語順の反転による。

周期点列 $(p_j)$ の一つの始点 $i$ からの右点加群を $M$ とし、
$M_{i+n}=k$ に次の文字を $p_{i+n}$ で評価して作用させる。
終点 $j=i+n$ の左単純のKoszul解消を $M$ とテンソルする。
$n\ge2$ の第一ホモロジーは
\[
 \operatorname{Tor}^A_1(M,k_j)
   =\ker(p_{j-1}:V_{j-1}\to k)\,/\,R_{j-2}(p_{j-2}) .
\]
実際これは $R_{j-2}\to V_{j-1}\to k$ の二微分そのものである。
従ってその次元は順側中心を訪れた二文字後で1、それ以外は0である。
$n=1$ では一次核が三次元である。極小生成元と第一Torの同定はイデアルの有限生成性を要しない。
各周期に一回以上新たな生成元が必要であるから、$e_i A\to M$ の核は無限生成である。
逆周期列を使う左側も同様である。
\end{proof}

\begin{proposition}\label{prop:terminal-branch-contact}
順側中心を頂点とし、各中心 $p$ の曲線列で $x_j(t)$ が順側中心になる全ての場合に、
$p$ から $x_j(t)$ への辺を置く。この有限有向グラフの終端強連結成分内の各 $p$ について、
その列の順側階数落ちはただ一つのパラメータ $t$ に集中し、到達する各 $x_j$ はそこで不分岐である。
特に長さ2の曲線列がこの成分から出発するとき、次の順側中心の階数二スキームは非被約である。
\end{proposition}
\begin{proof}
異なる二つのパラメータから成分内の中心に着けば、それぞれに $p$ への帰路を付けられる。
二つの閉経路は最初の直線で異なる点を選んでいる。反復で長さを揃えた後、二つを任意順に選べば
指数個の独立な経路評価が生じる（定理\ref{thm:recurrent-point-correspondences}の二分岐の証明）。
また $x_j$ の微分が零なら、そのパラメータの双対数近傍は $p$ と $x_j(t)$ を固定しながら
最初の直線の点を非自明に動かす。定数の帰路を付け、独立な双対数を毎回使えば
定理\ref{thm:no-two-sided-centers}の最後と同じく指数個の独立評価を得る。
$m=2$ なら $d_2=1$、従って $f_1=3-d_2=2$ である。一つの点で二次の接触をするため、
その順側階数二スキームは $x_1$ の非零接ベクトルを含む。
中心は孤立しているのでこのスキームはそこで非被約である。
\end{proof}

% !TeX root = ../papers/nc-p3.tex
\section{反対階数六・七での片側分岐の排除}
\label{sec:low-rank-branch-exclusion}
\begin{theorem}\label{thm:low-rank-branch-exclusion}
反対切断の階数 $\rho\le7$ の標準正則窓で射影像が高々一次元なら、全ての辺の全射影の
ファイバーは空か一点である。従ってこの階数範囲の全窓は、曲面射影の分類と
系\ref{cor:low-opposite-point-graph-as}の点グラフ判定で被覆される。
\end{theorem}
\begin{proof}[直交補と第四の辺]
残る分岐があれば前節の有理曲線列を取る。反対切断を $\rho$ 次元空間とその双対の完全対合に分解し、
点 $x$ における係数行列の行空間を $W_x$ とする。その次元は $w(x)$ の中央支持次元である。
隣接点対の条件は二つの行空間の直交である。
列の $x_i,x_{i+1}$ ($0\le i\le m$) で一般に次元和が $\rho$ なら、直交補が一致するので
$x_{i+2}$ と $x_{i-1}$ も隣接点対になる。従ってこの四点を循環できる。
$i=0,1$ では $x_{-1}$ または $x_0$ が固定されており、動く $x_1$ を積に含む
四文字閉経路の二つを反復できて矛盾する。$i=m-1,m$ でも逆順に同じである。
内部なら一般点に一意な前点を取って $x_{i+2}=x_{i-2}$ を得る。
動く既約成分の四位置内反復または定点への一致になり、曲線列の構成に反する。
従って全辺で次元和は $\rho$ 未満である。
動く一般点の両端支持は3なので中央支持も3以上である（補題\ref{lem:middle-binary-pencil}）。
$m\ge2$ のため $\rho\le6$ はこれで済む。$\rho=7$ では全ての動く一般点で中央支持は3、
両定点でも中央支持は高々3である。
\end{proof}

\begin{proof}[各順側中心に対応する逆側中心と純関係]
順側中心 $p$ を位置0に置き、$S\subset V_1$ を二次元順支持とする。
最初の直線 $L=\mathbb P(S^\perp)$ の一般点 $q$ で
$w(q)=\sum K_j(q)\otimes A_j$ と書け、三つの $K_j(q)\in V_2$ は独立である。
従って $V_3$ の中央支持は $A_j$ の支持の和で一定であり、三次元である。
それを消す一点を $y\in\mathbb P(V_3^*)$ とする。全 $q\in L$ で $w(q,y)=0$ なので
$w(y)$ の中央支持は $S$ に含まれ、一次元支持は不可能なので $S$ そのものである。
$w(q,y)=0$ を隣接二切断の両方で読むと、$y$ の両端評価と $q$ の対応する両端評価の次元和は
それぞれ6以下である。一般の $q$ の両端階数は3なので $y$ の両端も3以下である。
中央二次元の束の補題により $w(y)$ は $2,2,3$ 型またはその逆、中央鉛筆はKroneckerの $L_2$ 型である。
特にその二つの二文字支持における三次元の縮約像は、二項二次の対称平面と同型である。
この平面は基底を取れば $s_0m_0,s_0m_1+s_1m_0,s_1m_1$ で張られ、
全ての $[\lambda:\mu]$ について $(\lambda s_0+\mu s_1)(\lambda m_0+\mu m_1)$ を含む。

この平面が $R_0$ にあれば $S$ の全文字に終わる純関係がある。
しかし $w(p)$ の末尾からの三次元縮約像は $S$ と高々三次元の中央支持の積にあり、純関係を含む。
実際 $2\times3$ 行列の射影二平面は三次元Segre多様体と $\mathbb P^5$ で交わる。
これは連続純関係の禁止に反する。従って平面は $R_1$ にあり、
$y$ は逆側中心、その前点の直線を $N=\mathbb P(M^\perp)$, $M\subset V_2$ と書く。
$R_1\cap(S\otimes V_2)$ の次元は3なのでこの交叉は上の $S\otimes M$ 内の平面ちょうどである。
よって $w(p)$ も中央支持が $M$ で $L_2$ 型である。
$N$ の一般点の中央支持も三次元でそれを消す反対点は $p$ であるから、逆順にも同じ構成をすると元に戻る。
すなわち位置 $i$ の順側中心と位置 $i-1$ の逆側中心は一対一である。
\end{proof}

\begin{proof}[純関係による列長の制限]
一組 $p,y$ から得た $R_1$ の全点対は「第一点が $L$ 上」または「第二点が $N$ 上」を満たす。
実際、純関係の二つの一次評価の積がパラメータ直線上で恒等零だからである。
位置1に別の順側中心 $p'$ があり、その順支持を $S'\subset V_2$ とすると $S'\cap M=0$ である。
$R_1$ には $M$ の全方向に終わる純関係があり、$R_2$ には $S'$ の全方向に始まる純関係があるからである。
従って $p'$ の後続直線は $N$ と交わらず、$p'$ 自身は $L$ にある必要がある。
これらの結論は位置を回転しても使える。

長さ $m=2$ の曲線列はない。その最後の直線は位置3の逆側中心の前点全体なので、
対をなす位置0の順側中心（元の $p$ とは限らない）から得た上の $N$ の一つである。
$d_2=1$ ゆえ $f_1=2$ で、ある位置1の順側中心 $p'$ からこの $N$ に後続があることになり矛盾する。
一位置に二つの順側中心もない。第一の中心の後続直線は第二のものとは異なるので、
その次の動く像は第二の $N$ に含まれ、次に一定となり $m=2$ になってしまう。

連続する二位置0,1に順側中心があるとする。位置1の中心は $L$ にあるので $f_1>0$、
$d_2=3-f_1\le2$ である。$x_2$ の像は位置1の中心の後続直線にはなれない。
そこでは一般前点が一定であるのに $x_1$ は変動するからである。
位置1の対から来る純関係により $x_3$ はその対の逆側直線 $N'$ 上である。
$m=2$ は排除済みなので $x_3$ は動き、$m=3$, $d_3=1$ である。
$0\le f_2=3d_2-3-d_3$ より $d_2=2$, $f_2=2$ となる。しかし位置2の順側中心からは
$N'$ に後続を持てない。従って順側中心は隣接位置に存在できない。逆側も同じである。
\end{proof}

\begin{proof}[三次曲線の次位置]
そこで位置0の順側中心から出る最初の直線上に新しい順側中心はなく、$f_1=0$, $d_2=3$ である。
$x_2$ は有理正規三次で、$V_2$ を二項三次の完全系とする。位置1,2を独立な二項変数 $x,v$ で読むと、
$L$ に制限した $w$ は $[x,v]A(v)$ と書ける。$[x,v]$ は対角の行列式で、$A(v)$ は
$U=\ker p$, $Y=\ker y$ の二つの三空間に値を持つ二次の行列である。
これは $x=v$ での消滅と、$L$ 上の一般点における位置0,3支持による。
$A$ は定数側の二因子のいずれでも全三空間を張る。

$Y^*$ による $A(v)$ の縮約の三空間を $D\subset S^2(v)\otimes U$ とする。
$\dim S^1 D\le5$ である。実際この積は制限テンソルの反対切断で $x,Y$ を縮約した像であり、
全 $x\in L$ について $W_x\subset W_p^\perp$、$\dim W_p=2$ なので五次元を経由する。
従って二つの独立な $u_0,u_1\in Y^*$ があり、
$A(v)(v_0u_0+v_1u_1)=0$ である。一次関係の二係数が比例すれば $D$ が三次元であることに反する。
一般の $A(v)$ は正確に階数2である。階数1なら、二項多項式環の一意分解で
$A(v)=h(v)a(v)b(v)^t$ と斉次多項式ベクトルの積にできる（例えば原始な一列を取り各列をその倍にする）。
両支持が3なので $a,b$ は各々二次以上となり、総次数2に反する。

$x_3(v)$ は $A(v)$ を位置3から消す。実際、$x_2(v),x_3(v)$ で $w$ は零となり、
$[x,v]$ は独立な $x$ で非零にできる。従って $x_3$ の $Y$ への制限は一般に
$v_0u_0+v_1u_1$ の倍で、$x_3$ の像は平面に入る。
さらに $T\subset V_3$ を $p$ の三次元逆支持とすると、$T$ への射影は
$w(p,x_2(v))$ の $T^*$ 側の核点である。$x_2$ は全空間を張り、$M$ 上の二座標比は非定数なので、
一般にこれは $2\times3$ の $L_2$ 鉛筆の核点である。
具体的に非零の二項パラメータ $(s,t)$ で二行は $(s,t,0),(0,s,t)$ となり、
核は $[t^2:-st:s^2]$ である。従って $x_3$ はこの滑らかな二次曲線上の錐に入る
（錐頂点への定数写像は $m\ge3$ に反する）。整二次錐は平面を含まないので
平面との交叉に入る動く $x_3$ の次数は高々2である。各 $x_j$ が像に双有理であることを使った。

いま $f_2=3d_2-3-d_3=6-d_3>0$ なので位置2にも順側中心がある。
$m=3$ なら終点の位置0に逆側中心があり、その対の順側中心は位置1となって連続禁止に反する。
従って $m\ge4$。位置2の中心の後続直線には $x_3$ の像は一致しない（一般前点は $x_2$ で動く）。
その純関係により $x_4$ は逆側直線に入り、$m=4$ となる。
他方、逆向きにも隣接する中心がないため最終直線の次に逆にたどる曲線の次数は3である。
これは $d_3=3$ を意味し、先の上界に矛盾する。これで階数7でも全ての分岐を排除した。
\end{proof}

\begin{corollary}\label{cor:pure-free-branch-rank-nine}
曲面のない標準正則窓に純関係がないとき、片側分岐があれば $\rho\ge9$ である。
さらに分岐中心 $p$ の $2\times4\times3$ 縮約の二項鉛筆は $L_3$ 型である。
\end{corollary}
\begin{proof}
$p$ の中央支持が3以下なら、三次元の末尾縮約像が $2\times3$ 以下の行列の空間となり、
上のSegreとの交叉で純関係を生じる。従って中央支持は4である。
順支持二次元で作る $4\times3$ 鉛筆は全非零パラメータで階数3でなければならない。
さもなくば末尾の核ベクトルによる縮約が順支持一次元の非零な純関係になる。
全支持性からKronecker分解には零行も正方部分もなく、$L_3$ 型である。
また $L$ の動く一般点で中央支持は3以上である。3なら上の一定反対点 $y$ が得られ、
その $L_2$ 型からいずれかの隣辺に純関係ができる。従って最初の二点の中央支持はともに4であり、
それらの次元和が $\rho$ より小さいことから結論を得る。
\end{proof}

% !TeX root = ../papers/nc-p3.tex
\section{二本の直線を通る分岐と純関係のない分岐の三次曲線}
\label{sec:two-line-branch-blocks}
\begin{proposition}\label{prop:two-line-branch-block}
曲面射影のない標準正則窓で、分岐曲線列が
$p,x_1,x_2,z$ と二つの動く位置を持つとする。
二つの直線に制限した四テンソルは $[x,y]h$ であり、$h\in\ker z\otimes\ker p$ は階数3である。
$p,z$ の中央支持次元を $a,b$ とし、$r=\operatorname{rank}w(p,z)$ とすると
\[
                        \rho\ge 6+a+b-r,\qquad r\le2 .
\]
特に純関係がなければ $\rho\ge12$ である。
$\rho=8$ の場合には $a=b=r=2$ となり、二つの直線の消滅二空間どうしに
二項対称二乗と同型な三次元の純関係の線型包がある。
\end{proposition}
\begin{proof}
両直線のパラメータを移動で揃える。$R_1$ が対角で消えるので制限表示が得られ、
両定点との隣接消滅から $h$ の支持は表示通りである。
中間辺の一般直交補はここでは全点で相補的である。実際 $x_1$ の逆支持は常に $\ker p$、
$x_2$ の順支持は常に $\ker z$ である。従って $R_2\simeq R_0^*$ の六空間を使う
定理\ref{thm:finite-branch-corridors}の直交束の列は
\[
          0\longrightarrow3\mathcal O(-1)\longrightarrow6\mathcal O
           \longrightarrow3\mathcal O(1)\longrightarrow0
\]
となる。左右を双対にしても同様であり、$H^0,H^1$ で $\mathcal O(-1)$ が消えるので
$R_0$ の位置1の直線への制限も、$R_2$ の位置2の直線への制限も六次元のままである。
これらは互いに異なる側の切断なので、同時に制限しても隣接切断の階数6は保たれる。
$[x,y]h$ のその階数は $2\operatorname{rank}h$ であり結論を得る。

反対切断でもこの表示は六次の可逆部分行列を与える。
位置0は $kp$ の補空間、位置2は直線上の二点座標の包を使う六行であり、
位置1の二点座標と位置3の $kz$ の補空間が六列である。
追加で $p\otimes V_2^*$ の四行、$V_1^*\otimes z$ の四列を取る。
追加行は先の六列で零（最初の直線は $p$ の後続）、追加列も先の六行で零である。
従って可逆部分を消去しても追加行・列の全階数はそれぞれ $a,b$、その交叉部分行列は $w(p,z)$ のままである。
一般の行列の階数は指定行群の階数と指定列群の階数の和から両群の交叉部分の階数を引いた値以上である。
これは指定行群から独立行を取り、残りの行の指定列への射影を数えれば従う。
$r\le2$ は両端の直線支持である。
純関係がなければ両側から系\ref{cor:pure-free-branch-rank-nine}を使える。
$\rho=8$ なら中央支持は少なくとも2なので全ての等号が必要である。
$p$ の $2\times2\times3$ 縮約の三次元関係像は $L_2$ の二項対称二乗であり、
$r=2$ によってその中央二空間は第二直線の消滅二空間に一致する。
\end{proof}

\begin{proposition}\label{prop:pure-free-branch-cubic}
純関係のない曲面射影のない正則窓で、任意の順側中心 $p$ の直後の辺の点対は、
第一点がその後続直線にあるか、第二点が $p$ から一意に決まる有理正規三次曲線にある。
一位置の順側中心は高々二点で、二点あればその位置の各後続直線の次の曲線は
相手に付随する三次曲線であり、次位置に順側中心はない。
反対向きも同様である。曲線列の最初の二曲線がどちらも直線なら、動く位置はその二つで尽きる。
\end{proposition}
\begin{proof}
系\ref{cor:pure-free-branch-rank-nine}の $L_3$ 表示で、位置2の文字は二項三次となり、
その点 $z$ との縮約は二行三列のHankel行列である。
その階数1以下軌跡はVeronese三次曲線である。外では行支持が $p$ の順支持全体を張るから、
$z$ の全前点は $p$ の後続直線上にある。
次位置に別の順側中心があれば、その後続直線はこの整三次曲線に入らないので
その中心自体が $p$ の後続直線に属する。

同位置の相異なる二中心の直線は相異なる。一方の直線の次の像は一定ではなく、
上の条件より相手の三次曲線に一致する。その二つの三次曲線も相異なる。
同じなら一般点が二つの異なる前点を持つか、元の相異なる二直線に同じ一般前点が
同時に属することになって不可能である。第三の中心があれば第二と第三の三次曲線が
第一の直線の次像として一致して矛盾する。二中心のとき各直線上で $d_2=3$ ゆえ
$f_1=0$ であり、次位置の中心はそこにない。

$d_2=1$ なら $f_1=2$ なので次位置にも順側中心がある。$m\ge3$ とすると、
$x_2$ の像はその新中心の後続直線ではない（一般前点は動く）。
よって $x_3$ はその三次曲線となる。一方 $f_2=3d_2-3-d_3$ は負になってしまう。
これで $m=2$ である。
\end{proof}

\begin{corollary}\label{cor:two-line-secant-roots}
純関係のない二直線列の両側で、最初の直線は終中心 $z$ の反対三次曲線の割線または接線であり、
その長さ二の交叉が最初の直線の順側分岐因子 $f_1$ をスキームとして与える。
最後の直線でも $p$ の三次曲線について逆向きに同じである。
\end{corollary}
\begin{proof}
直線の消滅空間を $S\subset V_1,M\subset V_2$ とする。$p,z$ をつなぐ
$w(p,z)$ は非零である。零なら全ての $p,x_1,x_2,z$ が四文字閉経路となるからである。
これは $L_3$ 鉛筆の末尾からの縮約なので階数2である（非零な二項二次式の一次式倍は独立）。
$w(z)$ の $V_1$ を二項三次 $S^3 B$ と同定すると
$R_1\cap(V_1\otimes M)$ は二項二次式 $g$ に対し
$(s g,t g)$ を対応させる三次元行列空間である。$w(p,z)$ に対応する二次式を $g_0$ とすれば、
その第一支持が $S$ なので $S=g_0 B$ となる。従って $L$ と二項三次Veroneseの交叉は
$g_0$ の二根（重複込み）である。

さらに $R_1\cap(S\otimes M)$ はこの一本に等しい。
二次正方行列の二空間なら非零の可分解元があるからである。
従って $R_1$ は $R_1\cap(S\otimes V_2)$ と $R_1\cap(V_1\otimes M)$ の和五次元に、
両直線上の括弧を持ち上げる一本を足したものである。$q\in L$ で第一の空間は消え、
括弧の評価は $V_2/M$ で常に非零である。順側階数の低下は、第二空間の
$g_0$ 以外の二元を評価した $M$ 内の二行列式の零因子になる。
その行列は、$S^3B/(g_0 B)$ の双対から $g\mapsto (q(sg),q(tg))$ を取ったものなので、
その根は $q$ がVeronese上にある点ちょうどである（Hankel行列の階数1条件）。
両因子はともに長さ2で、二つの相異なる根なら各単根、一つなら二重根となる。
\end{proof}

% !TeX root = ../papers/nc-p3.tex
\section{二直線列の両端中央支持の全制限}
\label{sec:two-line-full-middle-restriction}
\begin{theorem}\label{thm:two-line-full-middle-restriction}
曲面射影のない標準正則窓の二直線列を $p,L,C,z$ とする。
$p$ の中央支持の $C$ への制限と $z$ の中央支持の $L$ への制限はいずれも階数2である。
$L,C$ 上の全点の中央支持は4である。
両端の中央支持次元を $a,b$、$r=\operatorname{rank}w(p,z)$ とすると
\[
             1\le r\le\min(2,a-2,b-2),\qquad \rho\ge6+a+b-r .
\]
従って二直線列には $\rho\ge11$ が必要で、等号なら $a=b=3$, $r=1$ である。
特に反対階数9,10の分岐では純関係を許しても全ての列が少なくとも三つの動く位置を持つ。
\end{theorem}
\begin{proof}
$L,C$ の消去二空間を $S,M$ とする。$P=R_1\cap(S\otimes V_2)$ と
$Q=R_1\cap(V_1\otimes M)$ はそれぞれ三次元で、$w(p),w(z)$ の末尾、先頭の縮約像である。
$w(p)$ の中央支持を $D\subset V_2$ とする。$R_2$ の $C$ への制限は
命題\ref{prop:two-line-branch-block}の束の列によって単射なので $D$ の $C$ への制限は非零である。
さらに $w(p)$ を $C$ に制限しても、位置1の $S$ 支持は二次元のままである。
実際、位置1での縮約の二空間が $R_2$ に入り、この単射で二次元を保つ。

もし $D$ の制限が一次元なら、制限した $w(p)$ はその固定中央文字と両端の行列の積で、
この行列は左支持が二次元なので階数2である。従って $P$ の $C$ への制限の階数は2で、
$P\cap Q$ の次元は高々1。ところが $D$ の制限を消す $y\in C$ があり、
$P+Q$ の少なくとも五関係が $y$ で消えてしまう。
六次元の $R_1$ の逆評価が階数1以下になり、定理\ref{thm:all-endpoint-ranks-two}に反する。
これで制限階数2が従い、逆向きにも同じである。

任意の $x\in L$ に対し、$w(-,x,C,-)$ は $[x,y]h$ の表示により
位置3で $\ker z$ の全三空間を張る。さらに $w(x,z)\ne0$ である。
これは $z$ の中央支持の $L$ への制限の階数が2だからである。
従って $w(x)$ の中央支持は4であり、逆側も同様である。
$w(p,z)$ の位置2支持は $D\cap M$ に含まれその上界は $a-2$ となる。
位置1側の上界も同じである。$w(p,z)=0$ なら $z\to p$ も辺になり、異なる
パラメータの $p,L,C,z$ で二つの四閉路を自由に反復でき、標準成長に反する。
あとは命題\ref{prop:two-line-branch-block}の切断不等式から従う。
\end{proof}

% !TeX root = ../papers/nc-p3.tex
\section{二直線列のない純関係自由な分岐の位置と長さ}
\label{sec:branch-centers-opposite-phases}
\begin{theorem}\label{thm:branch-centers-opposite-phases}
曲面射影と純関係のない標準正則窓で、片側分岐があり、二直線列がないとする。
特に $9\le\rho\le11$ の残余分岐はこの範囲にある。
四位置のうち相対する二位置だけに順中心が一つずつ、同じ二位置に逆中心が一つずつある。
二つの曲線列の動く位置数の組は順序を除いて
\[
                           (3,3),\quad(3,7),\quad(5,5)
\]
に限る。これは必要条件であり、この三配置の存在を主張しない。
\end{theorem}
\begin{proof}[位置の制約と次数十の接空間]
位置を法4で読む。$F_i,B_i$ を位置 $i$ の順・逆中心集合とする。
$F_i$ の一つの中心 $p$ の後続直線を $L_p$、その反対三次曲線を $C_p$ とする。
命題\ref{prop:pure-free-branch-cubic}より辺 $i+1\to i+2$ の動く二成分間の移動は、
$L_p$ からのもの、または $C_p$ に着くものに限る。
一般前点・後点の一意性によりこの相異なる移動は高々二本である。
特に $F_i$ が二点なら両移動は二つの初回の移動で既に尽きる。
二中心の後続直線は交わらない。どの点でも前点は元の中心一つだからである。
$B_{i-1}$ が存在すれば、その直前の動く二成分の移動はこの同じ辺で第二成分が直線なので、
第一成分は $L_p$ に等しい。その一般前点は $p$ で、そこから二直線列になってしまう。
従って $F_i$ があれば $B_{i-1}$ は空である。また $B_{i+2}$ の任意の中心の反対三次曲線に
$F_i$ の全中心が入る。逆向きの二者択一でこれに入らない順中心の後続直線は、
当該逆中心の前点直線に等しくなり、動く位置が一つだけの列を与えてしまうからである。

ある位置で相異なる動く曲線の次数和が10に達するとき、そこにある順中心を通る
それらの曲線の解析的な枝が二本以下で全て滑らかであることは不可能である。
実際、順辺の一般五式の完全交叉は次元1で、その射影は五行四列評価行列の最大小行列式で切られる。
次数を数える曲線を含み、台の次元1なので高さ2のHilbert--Burch解消
$0\to4\mathcal O(-5)\to5\mathcal O(-4)\to\mathcal I\to0$ を持ち、
埋込点のない次数10の曲線である。次数を使い切った各成分は被約で、全スキームがその被約和となる。
一方、順中心では評価階数2なので四次小行列式の一次微分は全て零で、接次元は3である。
滑らかな一枝の接次元は1、滑らかな二枝の和も高々2である。
後者の完備局所環は二つの $k[[t]]$ を共通の交叉商 $k[[t]]/(t^e)$ 上で貼った環で、
パラメータを合わせれば $(t,t),(0,t^e)$ で生成できる。
正規化で不分岐な射影直線の枝は完備局所環に一様化変数を含むので滑らかである。

順中心の位置が一つ $i$ だけと仮定する。各列で内部に順中心へ戻る必要があるため
動く位置は少なくとも五つである。最初の三次数は次数公式で $1,3,6$ となる
（次の二位置の順側欠損が零、$b_1=0$）。中心二つなら相異なる二曲線が同じ位置で
ともに次数6となり上限10に反する。一つなら第四次数は $10+b_2\ge10$ で位置 $i$ の上限を使い切る。
列に同じ成分は再登場しないので内部で順中心を訪れるのはこの成分だけである。
命題\ref{prop:terminal-branch-contact}からそのパラメータは一つで不分岐であり、
先の接空間に反する。逆向きにも占有位置が一つだけの場合はない。
$F_i$ と $B_{i-1}$ の排他性より、順・逆とも正確に二位置ずつを占める。

\end{proof}
\begin{proof}[隣り合う位置の排除]
$F_0,F_1$ が存在すれば $B_1,B_2$ が存在する。$F_0$ は一つである（二つなら $F_1$ は空）。
$F_1$ の全中心は $L_{p_0}$ 上にある。二つならその上の $f_1\ge2$ から第二次数は1以下で、
二直線列になってしまう。逆側も一つずつである。二列の最初の次数は
$p_0$ の側で $1,2,3$、$p_1$ の側で $1,3,6$ である。
前者の第二次数は正の $f_1$ と二直線の排除による。第三の像は $F_1$ の二者択一により三次曲線である
（第二の円錐曲線は $F_1$ の後続直線でない）。第三像が定点なら終中心が位置3に来てしまう。
後者は次の二位置に順中心がないことによる。いずれも少なくとも四つの動く位置がある。
$p_0$ の第四次数は $4+b_2$、同位置0の $p_1$ の第三次数は6なので $b_2=0$ で次数10を使い切る。
$p_1$ の第三の曲線上に順中心 $p_0$ がなければ、第四次数は $10+b_2\ge10$ となり、
位置1にある $p_0$ の最初の直線と両立しない。
従って両中心は互いに到達し、唯一の終端強連結成分をなす。
位置0で $p_0$ を通る二成分の各々は、多くとも一つのパラメータでそこに至り不分岐である。
これも次数10の接空間に反する。

\end{proof}
\begin{proof}[相対する二位置]
添字を命題の表示のために選び、順・逆とも位置0と2を占めるとする。
列長は奇数で3以上である。$F_0$ が二つなら辺 $1\to2$ の移動枠が既に埋まるので、
$F_2$ の全列は長さ3、$F_0$ の列は長さ3または5である。
長さ5のものがなければ $B_0$ は二つで、$F_0$ の列の中央にある二本の三次曲線は
逆向きの二中心の二者択一からちょうど $B_0$ の反対三次曲線二本である。
上で示したことから $F_2$ の中心は両方に属する。しかしその点の前点は一意であり、
交わらない $F_0$ の二つの後続直線の両方には入れない。
長さ5のものがあれば $B_0$ が空でないため残る一つは長さ3で、
$B_2$ の上限二つにより $F_2$ は一つである。二つの $B_2$ の反対三次曲線の一つは
$F_2$ の長さ3の列の中央曲線であり、そこには $F_0$ の二点がともに属する。
一方、$F_0$ の長さ3の列では初めの $1,3,1$ から $f_2=5>0$。
長さ5の列でも最後から第二の次数が3である（逆中心の隣の位置に逆中心がない）。
第三の次数を $d$ とすれば第四次数 $3=3d-8+b_2$ より $d\le3$、
ゆえに $f_2=6-d>0$。両方の $F_0$ が $F_2$ に戻り、全三中心が終端強連結成分となるが、
$F_2$ が相異なる二点へ二パラメータで進むのは不可能である。
これで二重の順中心を排除し、逆側も同じである。

最後に各枝の長さは9以下である。自分の最初の移動位相を再訪して動く曲線を渡れるのは
残る一枠だけであり、三回目には進めない。さらに両方の列が長さ5以上なら、
他の列も第一の移動位相を一回使うので、自身の長さは5を超せない
（奇数長7以上なら自分の第五曲線から第六曲線へ渡る）。
両終端は別々の位置を占めるので二つの長さはともに $3\bmod4$ またはともに $1\bmod4$ となる。
これらの条件から表示の三組に限られる。
\end{proof}

% !TeX root = ../papers/nc-p3.tex
\section{順中心と反対の逆中心による分岐階数十一の下界}
\label{sec:branch-rank-eleven-floor}
\begin{theorem}\label{thm:branch-rank-eleven-floor}
曲面射影のない標準正則窓に純関係がなく片側分岐があれば、反対階数は
$\rho\ge11$ である。$\rho=11$ なら前節の相対二位置で、
順中心 $p$ とその反対位置の逆中心 $d$ を取る。$p$ の後続直線と $d$ の前点直線は
ただ一点 $e$ で交わる。反対切断の両側の行空間について
\[
 \dim W_p=\dim W_d=\dim W_e=4,\quad
 W_p\cap W_d=k w(p,d),\quad W_e=(W_p+W_d)^\perp .
\]
共通元は階数1である。とくに反対階数9と10では純関係のない全ての曲線射影が
一点ファイバーのグラフ判定へ戻る。
\end{theorem}
\begin{proof}
まずどの純関係自由な順中心 $p$ でも後続直線 $L$ の全文字で中央支持は4である。
系\ref{cor:pure-free-branch-rank-nine}で一般点について成立し、
定理\ref{thm:finite-branch-corridors}の $w(x)=\sum K_j(x)\otimes A_j$ の表示で
三つの $K_j$ が独立な点は同じ中央支持を持つ。従属点は順中心であり、その縮約自身が
$L_3$ 型なので中央支持4である。逆向きにも同じである。

二直線列なら命題\ref{prop:two-line-branch-block}から $\rho\ge12$ である。
その他では前定理により順中心 $p$ に反対位置の逆中心 $d$ がある。
$p$ を位置0、$d$ を位置2とし、それぞれの位置1での支持を $S,T$ とする。
どちらも二次元である。$S=T$ なら $p$ の後続直線と $d$ の前点直線が一致し、
一般点はその二定点を結ぶので一つしか動く位置のない列になり矛盾する。
一方 $w(p,d)$ は $w(p)$ の四次元中央支持によって非零であり、
位置1では $S\cap T$ に支持される。よって $S\cap T$ は直線で、
二つの点ファイバーはただ一点 $e$ で交わる。

$W_p=w(p,-)$ と $W_d=w(-,d)$ はどちらも四次元で、それぞれ位置1,3の二次行列の空間である。
共通元は位置1で $S\cap T$ に支持される。$W_p$ は二行三列のHankel四平面であり、
固定した像直線に入る行列は一次元である。例えば非零な左方向で縮約すれば、
二項三次の四係数への隣り合う係数の三本の条件は独立である。
従って共通空間は非零な $w(p,d)$ の直線に一致する。
$W_e$ は $L$ 上なので次元4で、隣接する両中心の行空間を消す。
従って $\rho\ge(4+4-1)+4=11$、等号で直交補の主張も従う。
\end{proof}

\begin{proposition}\label{prop:branch-first-pencil-splitting}
純関係自由な順中心の後続直線では、$W_p^\perp$ の定数束に入る中支持束は
$4\mathcal O(-1)$ であり、余核束の次数は4、階数は $\rho-8$ である。
第二の動く像が有理正規三次なら、余核の各直和因子 $\mathcal O(a)$ は $0\le a\le3$ を満たす。
特に $\rho=11$ では分裂次数は順序を除いて $(1,1,2)$ である。
\end{proposition}
\begin{proof}
位置3の全文字による $w(x,-)$ は直線の各点で単射なので主張の部分束となる。
余核は大域生成され行列式の次数は4である。
第二曲線が三次なら、そこからの偶位置の行像を $W_p$ を法として読むと、
隣接消失により余核束の双対の3捻りに入る切断となる。
これらの切断の全値は定数空間 $W/W_p$ を張る。実際三次曲線の評価点が位置2の全空間を張り、
位置0では全文字を自由に取るためである。
もし $a>3$ の因子があれば総次数4より $a=4$、残りは全て次数0となる。
次数0の双対から定数束への埋込みは定数なので、切断の値はその固定真部分空間にしか入らない。

$\rho=11$ で上の逆中心 $d$ を用いる。直線 $L$ の $e$ における四平面を $G=W_e$ とする。
$W_p^\perp/G$ は三次元で、$W_d$ との対合で識別される。
$e$ の外の $x\in L$ では $x$ の $T$ への制限は固定された非零方向の倍であり、
$d$ の $L_3$ 鉛筆から $w(d,x)$ は階数3である。
従って四列の $G$ を法とした射影は、$e$ で消える一次式 $t$ と定数階数3行列の積となる。
その核の定数元を $s$ とし、$e$ で非零の補座標を $u$ とする。$s$ に対応する列は
$u e_0+t e_1$ と書ける。全文字で中支持が4なのでこの列には零点がなく、$e_0,e_1\in G$ は独立である。
$e$ での評価が $G$ への同型であることを使い、$s$ の三補元を $e_1,e_2,e_3$ に写るよう選ぶ。
残り三列の $t$ 係数 $h_1,h_2,h_3$ は $G$ を法として独立である。
こうして残りの列は $u e_j+t h_j$ ($j=1,2,3$)。
最初の二列は三行二列のKronecker束、残り二列は各々二行一列の束をなし、
余核は $\mathcal O(2)\oplus2\mathcal O(1)$ である。
\end{proof}

% !TeX root = ../papers/nc-p3.tex
\section{反対三次曲線との接触による長い分岐列の排除}
\label{sec:branch-cubic-contact}
\begin{theorem}\label{thm:branch-cubic-contact}
定理\ref{thm:branch-centers-opposite-phases}の三配置のうち、残る可能性は $(3,3)$ だけである。
このとき両方の列の次数は $1,3,1$ で、偶位置における順・逆中心はともにそこでの中央三次曲線に属する。
\end{theorem}
\begin{proof}
順中心 $p$ の最初の直線を $L$、その次の実際の三次曲線を $A$ とする。
直線の位置には順・逆いずれの中心もないので、
命題\ref{prop:isolated-branch-blocks}とその逆により最初から二番目、最後から二番目の像は
完全三次系で埋め込まれる。$S$ を $L$ の消滅空間とする。
純関係自由な $L_3$ 鉛筆の三関係は $S\otimes V$ にあり、偶位置の適切な座標 $y$ で、
その評価は二行三列のHankel行列 $(y_{i+j})_{0\le i\le1,\,0\le j\le2}$ である。
その二次小行列式の零点スキームが反対三次曲線 $K$ である。
$L$ の外で後点があれば、それは $L$ を中心とする直線束への射影に三次埋込みを合成した値である。

同じ奇位置で別の列または同じ列の内部の曲線 $D$ を取り、両隣が埋め込まれた三次曲線とする。
その一般前点は動くので $D\ne L$。後点への写像が双有理なので、正規化からの
二つの射影座標の共通零因子は $\deg D-1$ である。$D$ 上で $L$ に入る点の前点は必ず $p$ で、
前の三次埋込みによりそのようなパラメータは高々一つである。さらにその重複度は高々1である。
実際、評価階数3の開集合では、三本の独立な列を選んで前点の座標比を周囲の開集合の正則関数で書ける。
この関数は $L$ 上で一定なので、$p$ で消える各座標の引戻しは $L$ のイデアルに入る。
前の写像が不分岐なため、正規化上でそのうち一つは位数1で消える。従って $\deg D\le2$。

長さ5なら第三曲線の両隣は上記の三次曲線である。長さ7なら第三曲線はその奇位置の初回直線でなく、
第四曲線はHankelの二者択一によりやはり三次曲線である。従って長い列では $d_3\le2$ で、
$d_1=1,d_2=d_4=3$、奇位置に順中心がないことから命題\ref{prop:branch-corridor-degree}は
$b_2=11-3d_3\ge5$ を与える。逆向きにも最後から二番目の三次曲線の順欠損は少なくとも5である。

$A\ne K$ なら $A$ 上の逆欠損因子はスキームとして $A\cap K$ に含まれる。
実際 $L\to A$ は同型で、各 $y\in A$ に対する前点ファイバーは $L$ と一点で交わる。
よって評価を $S$ で割った行列の階数は各点でちょうど1である。$S$ 外の単元を一つ選ぶと、
Hankelの二列とこの列から作る三次小行列式がその単元とHankel二次小行列式の積になる。
逆向きでも同じ包含が成立する。

$(3,7)$ なら長さ7の列の第二・第六曲線、$(5,5)$ なら一つの列の第二曲線と他方の第四曲線を取る。
いずれも同じ偶位置の相異なる三次曲線である。成分の一致は一般前点をたどって列の初めを復元すると矛盾する。
後者の三次曲線に来る奇曲線は前者の初回直線でないので、それは前者に対する $K$ である。
逆方向に読むと前者も後者に対する反対三次曲線である。従って両者の交叉は、
逆中心と順中心という相異なる二点で各々長さ5以上となる。
一方、相異なる有理正規三次曲線二本の交叉の長さは高々5である。
長さ6以上なら第一曲線の三本の独立な二次方程式は、第二曲線の上では
同じ長さ6の因子で消える六次式なので高々一次元を張る。よって独立な二式が両曲線を含む。
非平面的既約曲線を含む二次式は既約で、二式の完全交叉は次数4に過ぎない。これで長い列は消えた。
$(3,3)$ では次数公式と逆公式から第二曲線で $f_2=b_2=5$ を得る。
\end{proof}

\begin{proposition}\label{prop:branch-three-three-frame}
残った配置で偶位置の三次曲線を $C(a,b)=(a^3,a^2b,ab^2,b^3)$、
順・逆中心をそれぞれ $C(1,0),C(0,1)$ とする基底が取れる。
各奇位置では初回の直線と終端の直線が
$L(a,b)=(a,b,0,0)$ と $M(a,b)=(0,a,b,0)$ となり、動く二辺はパラメータを保って
$L\to C\to M$ をたどる。二偶位置の独立二項変数を $x,t$、間の奇位置を $v,v'$ とし、
\[
 h(x,t;v)=x_1t_1v_0-x_1t_0v_1+x_0t_0v_2
\]
と置く。窓の全係数を偶位置で三次曲線に制限して書いた多項式は
\[
 hh'B+v_3 h'P+v'_3 hQ+v_3v'_3D,\qquad
 h=h(x,t;v),\quad h'=h(t,x;v'),
\]
となる。$B,P,Q,D$ の二重次数は順に $(1,1),(2,2),(2,2),(3,3)$ であり、表示は一意である。
特にこの38次元線型空間内で $B$ は非退化でなければならない。
\end{proposition}
\begin{proof}
両中心は異なるので偶基底を選べる。定理\ref{thm:branch-rank-eleven-floor}の証明により
二つの奇直線は一点で交わる。この点の後点は逆中心、前点は順中心である。
従って各直線で隣の三次曲線から移したパラメータは表示の共通点で整合し、
交点の代表ベクトルを合わせて同時に四次元の基底へ延ばせる。
三次埋込みの座標は三次式の基底なので制限でテンソルの情報は失われない。
$v$ の最初の三係数は $M(x)$ と $L(t)$ をともに消す。従って分数体上でその外積の倍数であり、
三係数に共通の既約因子がないので多項式倍である。反対位置でも同様で、独立な文字 $v_3,v'_3$ で
分ければ表示と一意性を得る。逆に各項は四つの動く辺と両端の定点・直線の条件を満たす。

各六次元関係空間から隣接奇位置の第四文字を除く射影は単射である。核の元は全て純テンソルだからである。
従って両奇位置で第四文字を除いた $hh'B$ は、隣接二切断でいずれも階数6である。
各切断の行・列はそれぞれ固定された六空間に入る。実際、十二次元の座標に動く辺の五係数と
定点の二式を課すと一式だけ重なり六次元になる。$B=\sum B_{ij}x_0^{1-i}x_1^it_0^{1-j}t_1^j$
として両切断を展開すると、この六空間の座標小行列式は符号を除き $(\det B)^3$ である。
添付計算は両方の六座標を指定し多項式恒等式を検算する。
\end{proof}

% !TeX root = ../papers/nc-p3.tex
\section{二つの切断の補完と共通四次根}
\label{sec:branch-pencil-resultant}
\begin{theorem}\label{thm:branch-pencil-resultant}
曲面射影のない標準正則窓で純関係がなく、片側分岐があれば、
二本の動く直線だけからなる列が存在する。とくに反対階数 $\rho\le11$ の純関係自由な窓には
曲線射影の片側分岐はない。
\end{theorem}
\begin{proof}
二直線列がないと仮定すると命題\ref{prop:branch-three-three-frame}が適用できる。
二項双一次形式 $B=x^{\mathsf t}Gt$ の行列式を $\delta\ne0$ とし、
$J=\left(\begin{smallmatrix}0&-1\\1&0\end{smallmatrix}\right)$、
$T=JG^{\mathsf t},\ U=JG$ と置く。$Tx,Ut$ はそれぞれ $B(x,-), B(-,t)$ の零点を表し、
$UT=TU=-\delta I$ である。二重二次式 $E$ に対し
$A_E(x)=E(x,Tx)$、$Z_E(t)=E(Ut,t)$ は二項四次式である。

まず $(x,v)\mid(t,v')$ 切断だけを階数6にする $D=D_0$ は任意の $P,Q$ に対し一意に存在する。
実際、第四文字を除いた両側の対の空間で許される関係は各六次元であり、
中央の可逆六次小行列がその両方を張る。片側だけ第四文字を許したブロックも反対側でその同じ空間に入る。
従って二つの長方形ブロックの間にこの六次行列の逆を挟んだものが唯一の補完となる。
これは $B$ を固定して $P,Q$ の双線型式である。この補完は
\begin{equation}\label{eq:branch-complete-product}
             PQ-BD_0=\delta^{-2}A_P(x) Z_Q(t)
\end{equation}
を満たす。実際、まず一般の $P,Q$ を取る。このとき $A_P,Z_Q$ はそれぞれ単純な四根を持ち、
固定した有限個の点を避ける。グラフ $B=0$ に制限する写像は全四次系へ全射だからである
（$T,U$ が同型なので二次式の積が全四次式を張る）。
$A_P$ の根 $x$ では $P(x,t)=B(x,t)k(t)$ と一次式で書ける。そこで
$h(x,t;v)$ の $v_0,v_1,v_2$ から全ての一次式が得られるので、$k$ を与える $v$ を選ぶ。
$x$ を固定した $v_3$ 行とこの $v$ 行は、第四文字を除く全ての反対列で一致する
（ともに $h'P$）。階数6はその列だけで達しているので、第四文字の列でも一致する。
従って $D_0(x,t)=k(t)Q(x,t)$ である。$Z_Q$ の根で列からの同じ議論を使う。
左辺の次数は $(4,4)$ なので $A_P Z_Q$ の定数倍となる。
$t=Tx$ と置けば $Z_Q(Tx)=\delta^2 A_Q(x)$ によりその定数は $\delta^{-2}$ である。
多項式恒等式として特殊な $P,Q$ でも従う。

もう一つの切断 $(v,t)\mid(v',x)$ は、$x,t$ の側に取る奇文字を交換するため、
右辺が $\delta^{-2} A_Q(x) Z_P(t)$ になる。両切断の階数が6なら
$A_P(x)Z_Q(t)=A_Q(x)Z_P(t)$ である。両四次式 $A_P,A_Q$ が非零なら比例し共通根を持つ。
一方が恒等零のときも他方の根を取り、両方が零なら任意の点を取る。
この共通根を $x$ とすると $P(x,t)=B(x,t)k(t),\ Q(x,t)=B(x,t)l(t)$ で、
式\eqref{eq:branch-complete-product}から $D(x,t)=B(x,t)k(t)l(t)$ である。
従ってこの偶点での全縮約は
\[
       B(x,t)\big(h(x,t;v)+v_3 k(t)\big)
               \big(h(t,x;v')+v'_3 l(t)\big)
\]
である。各括弧は $t$ に一次なので、両奇文字での支持次元はともに2以下である。
三次曲線での評価は四文字の完全座標なので、これは実際の両側支持である。
定理\ref{thm:no-two-sided-centers}に反する。反対階数の系は命題\ref{prop:two-line-branch-block}による。
\end{proof}

% !TeX root = ../papers/nc-p3.tex
\section{分岐を含む射影の次数七上界と二直線列の二重接触}
\label{sec:branch-degree-seven}
\begin{theorem}\label{thm:branch-degree-seven}
曲面射影のない標準正則窓の各辺・各方向で、評価の四次小行列式スキームの一次元部分の
総次数は重複込みで7以下である。従って定理\ref{thm:periodic-branch-obstruction}の
中心数と分岐を通る周期点列の長さの上界は、それぞれ14と28にできる。
\end{theorem}
\begin{proof}
一般平面上で評価行列の転置の余核を $E$ とし、
$0\to4\mathcal O(-1)\to6\mathcal O\to E\to0$ を取る。
一般点で階数4なので左は単射である。平面は孤立した階数2中心を避け、有限個の階数3点以外では
$E$ は束である。局所自由解消から閉点で深さ1以上なので $E$ は無捩れで、
$H=E^{**}$ は階数2の束、$H/E$ は有限長である。
階数3点では単元な三列を消去すると一本の三成分列が残り、
$\mathcal E xt^1(E,\mathcal O)$ はその三成分で割った局所環である。
これは四次小行列式の商環で、一般平面上の長さの和 $d$ が求める次数である。
$0\to E\to H\to H/E\to0$ の双対と局所双対性により、その長さは $H/E$ の長さにも等しい。
従って $c_1(H)=4,c_2(H)=10-d$ である。

$d\ge7$ とすると平面のRiemann--Roch公式で $\chi(H(-3))=d-6>0$、
また $h^2(H(-3))=h^0(H(-4))\le h^0(H(-3))$ なので $H(-3)$ は切断を持つ。
$H(-b)$ が非零切断を持つ最大の $b$ は3以上で、
飽和した切断から $0\to\mathcal O(b)\to H\to I_Z(4-b)\to0$ を得る。$Z$ は有限である。
$H$ は有限個の点の外で六つの定数切断で生成されるので $b\le4$。
$b=4$ なら $I_Z$ に非零切断があり $Z$ は空で、$6\mathcal O\to H\to\mathcal O$ が
評価行列の非零定数関係を与える。定理\ref{thm:branching-centers-curve}の証明で
一般平面への制限後にも六関係が独立であることに反する。
よって $b=3$、$10-d=3+\operatorname{length}Z\ge3$ である。
等号なら平面上で $H=\mathcal O(3)\oplus\mathcal O(1)$ となる。
全列の動く曲線は同じ位置では重複しないので、四位置で次数の和が28以下となる。
先の周期列の証明はこの和と各列の動く位置数が2以上であることだけを用いる。
\end{proof}

\begin{theorem}\label{thm:two-line-double-tangency}
曲面射影と純関係のない標準正則窓で分岐があれば、各位置の順中心と逆中心はそれぞれ高々一つである。
二直線列の初回直線・終端直線は両方とも系\ref{cor:two-line-secant-roots}の三次曲線の接線であり、
いずれの接触因子も一点の二倍である。
\end{theorem}
\begin{proof}
二直線列の始点を位置0とする。初回直線上の順欠損の二根が異なると仮定する。
位置1の二順中心から始まる両列の $d_2$ は3である（命題\ref{prop:pure-free-branch-cubic}）。
二列とも第三曲線が動き、その次数は $6-f_2$ である。
位置3に順中心がなければその二曲線は位置0にある相異なる六次曲線となり、次数上界に反する。
位置3の順中心が存在すれば、その後続直線も位置0にある別の曲線である。
先の二曲線とは一般前点が定点か動点かで区別される。従って一方でも次数6なら
総次数は少なくとも $6+1+1=8$ となる。よっていずれの位置1の列も位置3の順中心に到達する。
そこから位置0の元の中心に進める。実際、順中心の次位置にある全順中心はその後続直線上にある
（同命題のHankelの二者択一）。これで元の中心から異なる二つの初回点を通って戻る二経路ができ、
繰返すと指数増大になり矛盾する。逆向きにも同じで、二根の場合は双方とも排除された。

定理\ref{thm:branch-pencil-resultant}により二直線列は少なくとも一つ存在する。
その位置0は一中心である（二中心なら $d_2=3$）。位置1の順中心は直線上の今の二重根一つで尽きる。
一般に次位置の順中心が二つなら初回直線上の欠損が2以上で $d_2=1$ となり、
二つの異なる根を持つ二直線列を作ってしまう。従って位置2の順中心は高々一つである。
位置3に二つあればその次の位置0には順中心がないという同命題に反する。
従って全位置の主張が得られ、逆向きも同じである。
\end{proof}

% !TeX root = ../papers/nc-p3.tex
\section{純関係のない全分岐の二つの次数配置}
\label{sec:pure-branch-degree-patterns}
\begin{theorem}\label{thm:pure-branch-degree-patterns}
曲面射影と純関係のない標準正則窓で片側分岐があれば、全ての分岐曲線列は
巡回的な記述の起点を適切に取ると、次のいずれかの配置になる。
\begin{enumerate}
\item 順中心は位置 $0,1,2$、逆中心は $0,2,3$ にあり、
順に位置 $0,1,2$ から始まる列の次数は $(1,1),(1,1),(1,3,1)$ である。
\item 四位置の全てに順・逆中心が一つずつあり、四列は全て $(1,1)$ である。
\end{enumerate}
添字移動を同型として割る主張ではない。各場合の動く像は同位置で重複しない。
\end{theorem}
\begin{proof}
定理\ref{thm:branch-pencil-resultant}により二直線列を位置0から始める。
各方向の中心は位置につき一つ以下であり、その占有位置集合 $F,B$ は
$\{0,1\}\subset F,\{2,3\}\subset B$ を満たす。列末尾は $F$ と $B$ の全単射を与える。
初回から終回の動く曲線の間の辺だけを移動として数える。位置 $h$ からの移動は、
$h-1\in F$ または $h+2\in B$ なら高々二つである。前者なら初回の移動以外は
必ずその順中心の反対三次曲線に着き、相異なる移動が一般後点を共有しないからである。
後者なら逆に終回への移動以外は反対三次曲線から出る。
特に12回連続して移動すれば位置1から三回動くため、各列の動く位置数 $m$ は12以下である。

以下は小さい整数の有限列挙である。再現のため全条件を記す。
各 $i\in F$ で $2\le m_i\le12$ を選び、$m_0=2$、終点 $i+m_i+1$ の集合は
重複なしに $B$ とする。位置別移動数に上の条件を課す。
各列の $d_1=d_m=1$、全 $d_j\ge1$、同じ位置の次数の総和は7以下とする。
$m=2$ のとき $i+1\in F,i+2\in B$ が必要である。$m>2$ では $d_2$ は
$i+1\in F$ のとき2、そうでなければ3である。これは $d_2=3-f_1$、次位置の全順中心が
初回直線上にあること、$d_2=1$ なら二直線列となることによる。最後から二番目には逆の条件を課す。
前段の初回・終回以外の移動には、それぞれ到着・出発の次数3を課す。
最後に $d_{-1}=d_0=d_{m+1}=d_{m+2}=0$ とし、$0\le j\le m$ に対する
$d_{j+2}-3d_{j+1}+3d_j-d_{j-1}$ が正なら $j\ge2,\ i+j\in B$、
負なら $j+1<m,\ i+j+1\in F$ とする。これは命題\ref{prop:branch-corridor-degree}による。

$T=(1,1), A=(1,3,1), H=(1,3,3,3,3,3,1)$ と略記すると、
全ループの結果は次の五行である。添付コードは各列の位置ごとの許容次数の直積を取り上記の条件だけで検算する。
\[
\begin{array}{c|c|c}
F&B&\text{開始位置順の列}\\\hline
012&023&(T,T,A)\ \text{または}\ (T,T,H)\\
013&123&(T,A,T)\ \text{または}\ (T,H,T)\\
0123&0123&(T,T,T,T)
\end{array}
\]
$H$ の第三位置の直前・直後は埋め込まれた三次曲線である。直前は初回直線の次の完全三次系、
直後はその位相にある別の初回直線を避ける移動の反対三次曲線である。
実際 $H$ が位置2から始まる場合、第三位置は1で順中心0の初回直線があり、
位置1から始まる場合、第三位置は0で順中心3の初回直線がある。
この直線を中心とする射影を第三曲線に制限すると次数1であるが、
引戻しの固定因子は直前の三次埋込みと一意な前点によって長さ1以下である
（定理\ref{thm:branch-cubic-contact}の最初の射影の議論。初回直線上の逆階数は全点で3である）。
従って第三次数は2以下になり、$H$ を排除する。残る二つの三中心行は巡回表記で一致する。
\end{proof}

% !TeX root = ../papers/nc-p3.tex
\section{純テンソル関係のない曲線射影での全分岐の排除}
\label{sec:final-pure-branch-frames}
\begin{theorem}\label{thm:final-pure-branch-frames}
標準正則窓で全ての射影像が高々一次元なら、片側分岐があるためには少なくとも一辺に
純テンソル関係が必要である。従って純関係自由な曲線射影の全窓では、反対階数によらず、
全辺の両側のファイバーは空か一点であり、点グラフの場合の判定が適用できる。
\end{theorem}
\begin{proof}[共通の基底]
前定理の二配置だけを扱う。位置 $j$ において初回の直線を $L_j$、終回の直線を $M_j$ と書く。
これらがともにあるときは一点 $e_j$ で交わる。
実際、それらを与える隣の順・逆中心の縮約は非零である（順中心の中央支持が4）。
その位置で両直線の消滅空間の交叉に支持されるので直線は交わり、
一致すれば間に一つだけ動く位置を持つ列となってしまう。
$L_j$ の全点は前点が固定した一つなので逆中心を含まず、逆に $M_j$ は順中心を含まない。
二直線列の初回・終回にある順・逆中心は系\ref{cor:two-line-secant-roots}で与えられる。

四中心の場合、各位置の二直線を $L_j(a,b)=(a,b,0,0), M_j(a,b)=(0,a,b,0)$、
順・逆中心を座標点 $0,2$ とすることができる。
ここで座標点の番号は $0,1,2,3$ とし、$e_j$ は1である。
実際、移動 $L_j\to M_{j+1}$ は $F_j$ を $e_{j+1}$ に、$e_j$ を $B_{j+1}$ に送るから、
各移動でそれ以外のパラメータも揃えられる。交わる二直線の基底は交点の代表の倍率を合わせて同時に延ばせる。
三中心の場合には位置0の三次曲線を $C(a,b)=(a^3,a^2b,ab^2,b^3)$ とし、
そこでの順中心を $C(1,0)$、逆中心を $C(0,1)$ とする。両中心がこの曲線にあることは列の次数公式による。
位置 $1,2,3$ の二直線は上と同じ形で、$L_3\to C\to M_1$ のパラメータを揃える。
$e_3$ は逆中心 $C(0,1)$ に移り、$e_1$ は順中心 $C(1,0)$ から来るので、
残り二移動 $L_1\to M_2, L_2\to M_3$ と整合する。

各文字で点座標の最初の三座標だけを使う制限を三座標制限と呼ぶ。
$L_j\to M_{j+1}$ の辺では、$F_j\times L_{j+1}, M_j\times B_{j+1}$ を合わせると
三座標制限後の関係は全て次の形に入る。
\begin{equation}\label{eq:pure-branch-core-four}
         \begin{pmatrix}0&0&c\\ d&-c&0\\ h&f&0\end{pmatrix}.
\end{equation}
また、二直線列の最初の中心 $F_j$ に対して $R_j$ の第二点を $L_{j+1}$ に制限した射影は
$\ker F_j\otimes H^0(L_{j+1},\mathcal O(1))$ への同型である。
逆に終回の直線による反対向きの制限も同型である（命題\ref{prop:two-line-branch-block}の束の証明）。
\end{proof}

\begin{proof}[四中心の排除]
全位置を三座標に制限したテンソルを $W$ とする。その係数は各隣接対で式\eqref{eq:pure-branch-core-four}
を満たし、解空間は一次元である。$W\ne0$ である。実際、一つの二直線列上で窓は括弧と階数3の完全対合の積であり、
その残り二文字を三座標に制限する際に対合の両側で各一次元しか失わず、階数は少なくとも1である。
倍率を定め、行・列の順序を対 $10,11,20,21$ とすると、各切断での四次小行列は
\[
 K=\begin{pmatrix}0&0&0&1\\0&-1&0&0\\0&0&-1&0\\1&0&0&0\end{pmatrix}.
\]
一位置だけ第四座標3、他を三座標とした許容係数の空間は次元4である。位置0をその位置とすると
$W$ の位置0に三座標のベクトルを代入した三行と、単一語 $3210$ の係数だけでこれを張る。
これらの小さい線型計算は各対 $00,01,12,22$ の係数を零、$02$ と $11$ の係数の和を零として得る。
添付コードは有理数上の核と基底を確認する。
従って各位置で第四の点座標ベクトルを三座標分だけ平行移動し、
一位置だけ3の係数を $3210$ の巡回語（位置ごとに独立な倍率）のみにできる。これらの変更は同時に行える。

位置0,2を3に固定した三行三列を $N_{ab}=w_{3a3b}$ とする。$(0,1)\mid(2,3)$ 切断の
$K$ の四行四列を消去してもこの部分は $N$ のままである。
$3a$ 行では $a=2$ のときだけ $10$ 列が残り、$3b$ 列でも $10$ 行しか使わず、
$(K^{-1})_{10,10}=0$ だからである。他の $(1,2)\mid(3,0)$ 切断でも $a3$ 行、$b3$ 列で
修正に使う対はともに $21$ で、再び逆行列の対角係数が零である。
第一の切断における $N$ の $a,b\in\{0,1\}$ 部分は可逆である。
実際 $K$ に対 $30,31$ を両側で追加した六行六列は
$R_0,R_2$ をいずれも初回直線へ制限する完全六空間の座標だからである。
ゆえに $N$ の列空間は座標 $0,1$ への射影が同型な二空間で、座標ベクトル2を含まない。

一方、第二の切断で $K$ を消した残りの全列は、$R_1$ の四座標 $10,11,20,21$ が零である
二次元核を張る。この核の元の $a3$ 成分の全体は $N$ の列空間そのものである
（残りの全階数が2で、既に $N$ が階数2だからである）。
$F_0$ の中心縮約から、$R_1$ には $L_1$ の消滅空間、すなわち行 $2,3$ だけに支持される三関係がある。
それらの三座標制限は式\eqref{eq:pure-branch-core-four}の $20,21$ の二座標だけなので、
非零な一関係がこの核に入る。その $a3$ 成分は座標ベクトル2の倍で、零なら一行3だけに支持されて純関係になる。
従って $N$ の列空間が座標ベクトル2を含み矛盾する。
\end{proof}

\begin{proof}[三中心の排除]
位置0だけは全座標を保ち、他の三位置で三座標制限を行う。
このテンソルを $C(a,b)$ で評価すると全係数が一つの二項一次式で割れる。
実際、二つの内側の対には式\eqref{eq:pure-branch-core-four}を課し、
位置 $0,1$ では $F_0\times L_1$ と $C\to M_1$、位置 $3,0$ では $M_3\times B_0$ と
$L_3\to C$ の全係数を消す。得られる解の基底は、同じ二次式テンソルに $a,b$ をそれぞれ掛けた二本である。
添付計算では四つの三次係数を一座標ずらすことでこれを確認する。零テンソルの場合を含め、
三座標制限が全て消える点 $(a,b)$ を取れる。まず $a\ne0$ とし $u=b/a$ とする。

$U_1,U_2$ を $R_1,R_2$ の両文字の三座標制限の核とする。両核は二次元である。
像次元は式\eqref{eq:pure-branch-core-four}により高々4で、$R_1$ を $L_2$ に、
$R_2$ を $M_2$ に制限した完全六座標のうちそれぞれ四つは三座標だけで読めるからである。
各核から、第四行の最初の三成分を読む写像と、第四列の最初の三成分を読む写像はいずれも単射である。
消えれば第四列だけ、または第四行だけで純テンソルになるためである。
とくに $U_1$ の第四行像は座標ベクトル2を含まない。この向きの $L_2$ への制限が単射だからである。

$w(C(1,u),-,-,-)$ を $v_3$（位置3の文字）の最初の三座標で縮約すると全て $U_1$ に入る。
これらは少なくとも一次元を張る。全て消えれば位置3の支持が高々1となるからである。
その第四行を集めた位置 $2,3$ の三行三列 $H$ は従って非零である。
$H$ は $R_2$ の三座標制限なので式\eqref{eq:pure-branch-core-four}の形であり、
$C(1,u)$ の前点 $L_3(1,u)$ を消すので $d=cu, h=-fu$。
$c=0$ なら非零列は座標ベクトル2の倍に限られ、先の $U_1$ の第四行像に反する。
従って $c\ne0$ で $H$ は階数2である。

同じ $U_1$ の縮約の第四列を集め、位置2を3に固定した位置 $1,3$ の三行三列を $Z$ とする。
核からの二つの読み方が単射なので $Z$ も階数2である。逆に位置1の三座標で縮約すれば全て $U_2$ に入る。
従って $Z$ の行空間は $U_2$ の第四行像の二空間そのものである。
この像は座標ベクトル0を含む。実際、$B_0$ の中心縮約は $R_2$ に $M_3$ を消す列 $0,3$
だけの三関係を与え、三座標制限で $10,20$ の二座標しか持たず、一関係は核に入る。
その第四行の最初の三成分は非零な座標0の倍である。一方 $Z$ の全行は前点
$L_3(1,u)=(1,u,0,0)$ を消すので座標0を含み得ない。これで有限の根は排除された。
$a=0$ の場合は語順を反転し位置 $1,3$ を交換する。この配置は再び同じ三中心配置で、
位置0の順・逆中心が交換される。二項座標を交換すればその根は有限になる。
三座標部分は二直線の線型包そのもので語順の反転によって保たれるから、同じ矛盾を得る。
\end{proof}

\begin{corollary}\label{cor:positive-dimensional-pure-free-decision}
純関係のない標準有限Frobenius候補で、一辺でも点対応の射影が正次元なら、
全次数の標準Koszul左右AS正則性と、正則な候補間の添字保持同型は有限に判定できる。
\end{corollary}
\begin{proof}
二次元以上なら定理\ref{thm:no-moving-plane-pencils}, \ref{thm:isolated-surface-fibers}の被覆判定を使う。
全射影が高々一次元なら、二点以上を持つ射影ファイバーの有無を評価の小行列式で調べる。
一つでもあれば本定理により非正則、なければ系\ref{cor:all-one-dimensional-point-graphs}に帰着する。
\end{proof}

% !TeX root = ../papers/nc-p3.tex
\section{分岐中心に現れる純関係とその反対曲線}
\label{sec:branch-local-pure-pencils}
\begin{proposition}\label{prop:branch-local-pure-pencil}
曲面射影のない標準正則窓で分岐があれば、順・逆の少なくとも一つの中心の
三次元縮約像が純関係を含む。すなわち、順側なら順中心 $p$ の後続直線を消す二空間 $S$ に対して
$P=R_{i+1}\cap(S\otimes V_{i+2})$ に純関係があり、逆側なら語順を反転した条件が成り立つ。
\end{proposition}
\begin{proof}
全てのこの三空間が純関係を含まないと仮定する。
系\ref{cor:pure-free-branch-rank-nine}の鉛筆の議論は各三空間だけを使い、全中心が $L_3$ 型になる。
したがって各中心の後のHankel三次曲線、次中心の初回直線上への包含、
同位相二中心の制約は命題\ref{prop:pure-free-branch-cubic}のままである。
定理\ref{thm:branch-centers-opposite-phases}, \ref{thm:branch-cubic-contact}の
中心数・次数・交叉の証明はこれらと一般の次数公式しか使わない。
二直線列のない場合の最後の $(3,3)$ で、命題\ref{prop:branch-three-three-frame}が使う
第四の奇文字の切除による単射も保たれる。その核の元は第四の奇文字に支持される純関係であり、
その文字は初回・終回両直線を消すので、向きに従い順・逆の中心の三空間自体に入るからである。
従って定理\ref{thm:branch-pencil-resultant}が二直線列を強制する。

二直線列で $w(p,z)$ はどちらの端からも $L_3$ の非零縮約なので階数2である。
系\ref{cor:two-line-secant-roots}で使う $P\cap Q$ の上界1も、$P$ に純関係がないことだけを要する。
これで定理\ref{thm:two-line-double-tangency}, \ref{thm:pure-branch-degree-patterns}の全次数配置を得る。
最後の定理\ref{thm:final-pure-branch-frames}の四中心の場合に排除する一行支持の関係は
初回直線を消す中心の三空間に属する。三中心の場合は $U_1,U_2$ のうち第四行だけ、または
第四列だけに支持される関係を排除すればよい。位置 $1,2,3$ の第四文字はその初回・終回直線の両方を消すので、
それぞれ該当する順・逆中心の三空間に属する。従ってこれらも排除され、同じ二矛盾を得る。
\end{proof}

\begin{proposition}\label{prop:finite-pure-pencil-opposite}
全純関係スキームが有限で、$p$ が順中心であるとする。中央支持次元 $a$ は3か4である。
辺 $i+1$ の第一点 $x$ が初回直線 $L$ の外にあるとき、$P$ の三関係だけから得られる
後点候補は次の通りである。
\begin{enumerate}
\item $a=4$ なら $P$ の鉛筆は $L_e\oplus J_{3-e}$、$1\le e\le3$。
ここで $J$ は正方正則鉛筆で各固有値に対しJordanブロックは一つだけである。
一般の後点候補は $L$ からの射影の値に対応する有理正規 $e$ 次曲線の点であり、
例外の各方向ではその点を通る一直線に増える。
\item $a=3$ なら三行三列の正則鉛筆で、各固有値の幾何学的重複度は1である。
一般の後点候補は中央支持を消す定点、例外の各方向ではこの点を通る一直線である。
\end{enumerate}
三空間内の純関係スキームの長さは $a=4$ のとき $3-e$、$a=3$ のとき3である。
各位相に順中心は高々二つで、二つなら両方とも $a=4$、一方の列の第二次数は他方の $e$ である。
同じ位相で異なる中心の純関係は互いに素である。全て逆向きにも成り立つ。
\end{proposition}
\begin{proof}
$a=2$ のとき三次元の $2\times2$ 行列空間は正次元でSegreと交わる。
$S^*$ 方向 $[s:t]$ で $P$ を評価する $a\times3$ 鉛筆には定数の左右核がない
（両者とも実際の全支持を取ったからである）。
一般階数が3未満なら、各方向で少なくとも一つ非零の純関係が生じるため有限性に反する。
よって三行では全体が正方正則であり、四行ではKronecker標準形に左の特異ブロックが一つあり、
その最小指数は正である。各例外方向の核が二次元以上なら、同じ第一因子で純関係の直線ができる。
従って正方部の各固有値に一つのブロックしかない。
純関係に対しその第一因子を消す $[s:t]$ はスキーム上でも一意である
（純テンソルSegreの第一射影による）。その方向の三列の核線が純関係を復元する。
単元な二次小行列で核線を消去すれば、残る条件は三次小行列で、
局所Smith形によりその長さは正方部の対応する単因子の位数である。

後点は評価鉛筆の左核である。$L_e$ の左核は完全 $e$ 次系の点、三行正則なら
四文字内で使わなかった一方向が常に残る。特異方向ごとに正方部から一次元が加わる。
同位置の別の順中心の初回直線は $L$ と交わらず、そこからの $L$ を中心とする射影は同型なので、
その第二曲線は一般方向の後点軌跡に等しい。定点にはなれず、また異なる二つの初回移動が
同じ第二曲線を共有しない。よって中心数の上界と第二次数が従う。互いに交わらない初回直線の
消滅二空間の交叉は零なので三関係の空間も交わらない。
\end{proof}

\begin{proposition}\label{prop:excess-contact-two-line-pure}
二直線列 $p,L,M,z$ の始点にある別の動く曲線枝を $x(t)$ とする。
$p$ の点イデアルの引戻しの位数を $h$、その枝上の順評価の三次小行列式の
共通位数を $f$ とし、一般順階数は3とする。$f\ge h$ であり、$f>h$ なら
$R_{i+1}$ の $2\times 2$ 部分空間 $L^\perp\otimes M^\perp$ に独立な二関係がある。
従ってその辺の純関係スキームは少なくとも長さ2の部分スキームを含む。
逆向きにも同じである。
\end{proposition}
\begin{proof}
局所パラメータを $t$ とし、$R_i$ の六関係に対する $x(t)$ の評価の四行を考える。
$L$ の二座標で読む最初の二行はともに $t^h$ で割れる。
それらを $t^h$ で割って剰余体に落とした二行は独立である。
実際 $R_i$ の $L$ への制限は $\ker p\otimes H^0(L,\mathcal O(1))$ への同型で、
二行は $x(t)$ の最初の非定数方向と独立な二座標を対にしたものである。
他の二行は $p$ の順階数が2なのでそのままで剰余体上独立である。
三次小行列は必ず最初の行群を使うので $f\ge h$。$f>h$ なら最初の二行を割った
各々と後の二行で作った三次小行列が剰余体で消え、両二行の張る空間は等しい。

六次元の相補切断によりこれらの行を $R_{i+2}$ の二文字行列と同定する。
$w$ の $L\times M$ への制限は二項括弧と、$\ker p,\ker z$ 上の非退化な三次正方行列の積である。
従って割った二行を $M$ に制限すれば、$M$ の独立な二文字と、非定数方向により
$\ker z$ に定まる同一の非零文字との積となる。この二行が $p$ の評価の二行でもあるため、
$w(p)$ の末尾からの三次元縮約像を $M$ に制限すると次元1になる。
その二次元核は $L^\perp\otimes M^\perp$ に入る。この二次元行列空間の射影直線を
二次の行列式で切れば、全直線または長さ2の零因子になる。
\end{proof}

% !TeX root = ../papers/nc-p3.tex
\section{全四辺でただ一つの単純な純関係の排除}
\label{sec:single-pure-branch-exclusion}
\begin{theorem}\label{thm:single-pure-branch-exclusion}
標準Koszul左右AS正則窓において、四つの $\mathbb P(R_i)$ とそれぞれのSegre多様体との
交叉スキームの非交和が長さ1になることはない。
\end{theorem}
\begin{proof}[中心の型と短い列の制限]
四辺が有限なら定理\ref{thm:finite-frobenius-low-opposite-as}による。
曲面以上の射影では定理\ref{thm:isolated-surface-fibers}を用いる。
非平面的な整曲面間の支配的対応では純関係を評価すると、片方の面が一つの超平面に入ることになり矛盾する。
平面列では純関係の一方の文字がその平面の法線であるが、その法線は正則正規である。
一次元の両側一点グラフは定理\ref{thm:curve-graph-pure-exclusion}による。
従って残るのは曲面射影のない分岐である。

唯一の純関係を $R_1$ に置く。命題\ref{prop:branch-local-pure-pencil}, \ref{prop:finite-pure-pencil-opposite}により、
普通の中心は全て $e=3$ で、例外の $e=2$ は順なら位置0、逆なら位置3にそれぞれ高々一つである。
少なくとも一方向に例外がある。例外中心は反対二次曲線に加えて、
ただ一つの例外方向でその曲線の一点を通る線を持つ。
両方向とも例外ならこの二中心は同じ二直線列の始終である。
実際、それらの三空間 $P,Q\subset R_1$ は同じ純関係を含み、その交叉はちょうど一次元である。
二次元なら、初回・終回直線 $L,M$ を消す $L^\perp\otimes M^\perp$ 内の
二つの行列により長さ2以上の純関係スキームを得る。
ゆえに $P+Q$ が五次元で $L\times M$ に残る双一次式は高々一本である。
零なら動く中心があり、階数1なら $L$ の一般点に同じ後点があり、動く位置が一つだけの列となる。
従ってその式は階数2で、初回直線から終回直線への移動を与える。

さらに、連続する位置 $i,i+1$ から二直線列が出て位置 $i+1$ の順中心が一つだけなら不可能である。
位置 $i$ からの初回直線上の順欠損はその一点に長さ2で集中する一方、直線の接触位数は1である。
命題\ref{prop:excess-contact-two-line-pure}から次の列の中央辺で長さ2以上の純関係が生じる。
また同じ位置から二本の $(1,3,1)$ 列が出て、反対位置の順中心が一つだけということもない。
両三次曲線は順欠損5を持ちその中心を通るが、そこでの一意な前点は
元の二中心の相異なる交わらない初回直線の両方には属せないからである。
\end{proof}

\begin{proof}[有限列挙と最後の二三次曲線]
各位置の順・逆中心数をそれぞれ $0,1,2$ から選ぶ。一方向の占有位相が一つだけなら、
そこからの列の初めの次数が $1,3,6$、第四次数が $10+b_2$ となり上界7に反する。
よって各方向に少なくとも二位相あり、普通の順中心が一つ以上ある。その中心の初回移動の位相では
初回以外は同じ反対三次曲線に着くから移動は二回以下であり、各列長は $m\le12$ である。

添付コードの有限列挙で使う条件は以下の通りである。
\begin{enumerate}
\item 順・逆中心の総数は等しい。各順中心から、適合する位相の相異なる逆中心まで
$2\le m\le12$ の列を割り当てる。両方向に例外があれば既述の二直線列を課す。
\item 全次数は正、初回・終回は1、同じ位相の次数和は7以下である。第二次数は $3-f_1$。
次位相に $n$ 中心があれば普通中心では全てが初回直線上にあり、例外中心でも
その外には高々一つしかない。外にある中心の後続直線は唯一の例外直線であって、二中心はこれを共有しない。
次位相が空なら $f_1=0$ である。自位相が二中心なら第二次数は相手の $e$ に等しい。
$m>2$ では第二次数は2以上である（1なら次数公式より第三次数は0以下）。
反対側でも同じ条件を課す。
\item 各移動がある中心の初回移動と同じ位相にありその初回自体でない場合、
行先成分を普通中心なら次数3、例外なら次数1か2に指定し、各中心につき同次数の行先を二度使わない。
この条件は反対曲線の各被約成分がその次数ごとに一本しかないことによる。逆向きも同じである。
\item 定理\ref{thm:pure-branch-degree-patterns}と同じ三階差の符号条件を課し、
前段の連続二直線列と同位相の二三次列の禁止を使う。
\end{enumerate}
例外位置に二中心があれば一番目を例外としてよい。整数計算で残るのは次の二行だけである。
\[
\begin{array}{c|c|c|c}
\text{順中心位相}&\text{逆中心位相}&\text{例外}&\text{二列の次数}\\\hline
0,2&0,2&\text{順 }0 &(1,3,1),(1,3,1)\\
1,3&1,3&\text{逆 }3 &(1,3,1),(1,3,1)
\end{array}
\]

最後は純関係の不在に依存せず定理\ref{thm:branch-pencil-resultant}の補完式を使える。
実際、二直線列がないので定理\ref{thm:ternary-center-cubic-scroll}が適用でき、
各初回直線と第二三次曲線への制限は $[s,t]A(t)$ で $A$ の一般階数は3である。
従って始点の関係空間を初回直線に制限した像は六次元である。異なる二つの一般 $t$ の
各々で $[-,t]$ と始点を消す全三空間の積を含むからである。逆方向でも同様である。
二つの直線が同じ位相に来るとき、それらは交わる。中央三次曲線が逆欠損5により逆中心を通り、
そのパラメータでの初回直線の点は逆中心の前点直線にも属する。両直線の一致は一動位置列を与えてしまう。
よって命題\ref{prop:branch-three-three-frame}の座標と $hh'B$ の表示が使え、
各奇位置の三座標制限は初回または終回直線を含むので前記の完全六空間を保つ。
両切断で $\det B\ne0$ となり、二つの補完式の共通四次根が両側中心を与える矛盾は同じである。
\end{proof}

\begin{corollary}\label{cor:four-two-line-impossible-with-pure}
純関係の個数や有限性を仮定しなくても、各位相一つずつの順・逆中心から四本の
二直線列ができる配置は標準正則窓にはない。
\end{corollary}
\begin{proof}
初回直線 $L_i$ 上にある順中心 $F_i$ は唯一の順欠損の台、終回直線 $M_i$ 上の $B_i$ も同様である。
$L_i$ から $M_{i+1}$ への同型で $B_{i+1}$ を引き戻せば $L_i\cap M_i$ の点となり、
両直線は一致しない。従って定理\ref{thm:final-pure-branch-frames}の四中心の基底、非零な $W$、
両切断の共通部分 $N$、その列空間が座標ベクトル2を含まないことは純関係があっても成立する。
同証明の最後で得た、初回直線を消す三関係のうち一つは、
したがって第四行だけに支持される純関係でなければならない。
巡回して全ての $R_i$ が $z_i$ に始まる純関係を持つ。ここで $z_i$ は
$L_i,M_i$ の張る平面を消す唯一の文字である。語順を逆にした同じ議論では、
全ての $R_{i-1}$ が同じ $z_i$ に終わる純関係を持つ。連続する純関係の禁止に反する。
\end{proof}

% !TeX root = ../papers/nc-p3.tex
\section{全正則窓の純関係スキームの有限性}
\label{sec:finite-pure-relation-locus}
\begin{theorem}\label{thm:finite-pure-relation-locus}
標準Koszul左右AS正則窓の各辺において、純テンソル関係のスキームは零次元以下である。
曲線射影に分岐があれば全ての順・逆中心の中央支持は3か4であり、
命題\ref{prop:finite-pure-pencil-opposite}の鉛筆と位相ごとの中心数上界が無条件に適用できる。
\end{theorem}
\begin{proof}[一文字を共有する純関係の排除]
曲面以上、有限辺、一次元の両側一点グラフでは、前定理の冒頭の理由で純関係はない。
曲面射影のない分岐の場合、一辺 $R_i$ が $u\otimes v,u\otimes v'$ を含み $v,v'$ は独立と仮定する。
平面 $H=\mathbb P(u^\perp)$ に制限して残る四式の四次行列式を $G$ とする。
$G$ は恒等零ではない。零なら全平面が射影像に入るからである。
$G=0$ の各既約成分の一般点で、重複度を保った平面上の局所環 $D=\mathcal O_{H,\eta}/(G)$ を取る。
順評価の三次小行列の一つは単元である。実際、階数2以下の点は有限個である。
ゆえに $D$ 上で行列の核は直和因子の直線で、一意な後点 $y$ が得られる。
その剰余体での値が動けば $y$ の順評価階数は3以上である。定点なら、
元の成分の全点がその前点なので逆中心であり、その順評価階数は3である。
したがって $D$ 上でも $y$ の順評価には階数3の直和因子がある。
隣接切断の完全対合と辺の消滅により元の点での逆評価はこの直交補に入る。
従って逆方向の全四次小行列式は $D$ 上で消える。この議論は $G$ の因子の全重複度を数えている。

平面上で逆四次小行列は全て $G$ の定数倍となり、少なくとも一つが非零である
（全て零なら逆曲面射影となる）。よって平面の一般点における
$R_{i-1}\to V_{i-1}$ の二次元核はPlücker座標が一定で、同じ二関係が平面全体で消える。
これらは $V_{i-1}\otimes ku$ に属し、$u\otimes v$ と連続する純関係を作って矛盾する。
逆向きも同じで、純関係の母線はどちらの向きにも存在しない。
\end{proof}

\begin{proof}[二次曲線と二項中心]
定理\ref{thm:pure-relation-curve-supports}により残る正次元の可能性は滑らかな二次曲線である。
その三関係は二空間 $A,B$ について $A\otimes B$ にあり、$A\otimes V$ に支持される
関係はこの三関係で尽きる。$L=\mathbb P(A^\perp)$ の一般点 $x$ で残り三関係を評価した
三列 $K_j(x)$ は独立である（可動中心の不在）。
$w(x)=\sum K_j(x)\otimes T_j$ と書けばその逆支持は $T_j$ の固定支持の和である。
後点が動けばその順階数は一般に3以上、後点が一定なら逆中心で順階数3なので、
直交により $x$ の逆階数は高々3である。一般に3であるから前点は一定 $p$ で、
$L$ がその後続直線である。命題\ref{prop:isolated-branch-blocks}による $p$ の三次元縮約像は
上記の三関係そのもので、中央支持は2である。

ここからは任意の中央支持2の順中心 $p$ を位置0に取って矛盾を示す。
定理\ref{thm:branch-binary-center-pairing}の対の逆中心を位置3の $z$、
その前点直線を位置2の $N$ とする。$p$ の後続直線を $L$ とする。
$p$ からの列の第二次数は1ではない。1なら次数公式から二直線列となり、
定理\ref{thm:two-line-full-middle-restriction}はその両端中央支持を3以上とするからである。

第二次数が2なら $L$ 上には唯一の順中心 $q$ があり、定理\ref{thm:rank-ten-conic-after-center}の
最初の標準二列の計算を用いる。この計算自体は反対階数を使わない。
独立パラメータでの制限は $[s,t]A(t)$、$A(t)$ は $q$ の鉛筆を、三行と中央支持の列で
書いた一次行列であり、第三点は $A(t)$ を消す。
$q$ は $L$ にあるので $w(q,z)=0$、その中央支持は3以下である。
中央支持2なら $A(t)$ は三行二列の $L_2$ で、左の核行は最小次数2である。
従って $[-,t]\otimes A(t)$ の列の全値を消す一次行はなく、これらは六空間全部を張る。
$R_0$ を $L$ に制限する写像が単射となるため、$A(t)$ の右核直線の全点が第二曲線の各一般点の後点となり、
可動中心になってしまう。中央支持3で常に特異なら、定数の左右核を除いた三次正方鉛筆は
$L_1\oplus L_1^{\mathsf t}$ に限る。一行二列のブロックが固定因子を共有する二純関係を与え、
前段に反する。正則三次正方鉛筆なら一般 $t$ の右核は四文字内で使わない定点一つであり、
二次曲線が終回になってしまう。従って第二次数は3である。
逆に $z$ からも第二次数は3となる。

$L$ 上の順欠損はこれで零である。$L$ の外の点の後点は二項対称二乗の三関係により
必ず $N$ に入る。よって位置1に順中心があればその直線全体が $N$ となり、次は $z$ で
動く位置が一つだけになってしまう。従って位置1には順中心がなく、逆に位置2には逆中心がない。
また二項中心の第二三次曲線の後の第三次数は4以上である。
定理\ref{thm:binary-centers-quadric-cones}の最初の錐の議論がここでも使える。
実際、三次から $N$ を中心とする射影の次数を $k>0$ とすれば第三曲線の頂点からの射影は
その二次写像で、$d_3\ge2k$。$d_3\le3$ なら $k=1$、三次曲線は $N$ と長さ2で会う。
第三曲線は二次または錐上の非平面三次で埋め込まれるが、その長さ2が $N$ の一定の後点 $z$ に写り矛盾する。
この議論は他の列に二直線列があるかによらない。逆向きにも同じ下界がある。

位置3に順中心があればその列の第二曲線は位置1にあり、一般前点が動くので $L$ と異なる。
二直線列なら位置2に逆中心を作り矛盾する。それ以外なら第三曲線は $N$ で、列はそこで動き終わり
$z$ に着く。これを $z$ から逆にたどる第三曲線は直線で、先の下界4に反する。
従って位置3にも順中心がない。残る順中心は偶位置だけで、次位相に順中心がないため
各列の第二次数は3である。これでは窓のどこにも二直線列がなく、
系\ref{cor:no-binary-centers-without-two-lines}が中央支持2自体を排除する。
逆中心も同様である。これで二次曲線もなくなり、有限性を得る。
\end{proof}

% !TeX root = ../papers/nc-p3.tex
\section{純関係を持つ辺の分布と長さ二以上の障害}
\label{sec:sparse-pure-edge-exclusion}
\begin{theorem}\label{thm:sparse-pure-edge-exclusion}
標準正則四窓で純関係が一つでもあれば、純関係を持つ辺は少なくとも二つあり、
そのうち一辺の純関係スキームは長さ2以上である。従って四辺合計の長さ2以下は
空の場合に限り、重複点だけを持つ一辺も任意の長さで排除される。
\end{theorem}
\begin{proof}
定理\ref{thm:finite-pure-relation-locus}により有限性は既知で、
曲面・有限辺・一点グラフの排除は定理\ref{thm:single-pure-branch-exclusion}の冒頭による。
同定理の列挙を、全辺の純関係スキームが長さ1以下の場合と、一辺だけ任意の長さを許す場合に拡張する。
中心の例外固有方向の重複込み個数を $l=3-e$ と置く。
一中心では $0\le l\le3$、同位相二中心では各 $l\le2$ である。
その和は当該辺の純関係スキームの長さを超えない。一辺にだけ純関係を許すときは、
その辺に対する上界4で全ての中心型を列挙できる。これは辺自体の長さの上界を仮定するものではない。
全辺長さ1以下の場合、記述の位相と向きを選んで位置0の最初の順中心を例外にできる。

次数列の整数条件の変更点を記す。$l>0$ の中心に対し、同位相の初回以外の移動は
$e>0$ のときには主曲線（次数 $e$）に一回以下、または例外直線（それぞれ次数1）に合計 $l$ 回以下である。
$e=1$ なら主曲線と例外直線を合わせて高々三回の直線への移動を数える。
$e=0$ では非定数の主曲線がない。初回直線の外にある次位相の順中心も高々 $l$ 個であり、
同位相が二中心のときは引き続き第二次数が相手の $e$ に等しい。
普通中心がどちらかの向きにあれば $m\le12$、全てが例外なら
単一辺だけの設定では二つ以上の占有位相を得られず、全辺単純の設定では
一位相の移動が高々三回なので $m\le16$ である。

非初回の移動の始点曲線を $D$、その一つ前の曲線が埋込みであるとする。
移動先が主曲線なら初回直線 $L$ からの射影が双有理であり、
定理\ref{thm:branch-cubic-contact}の引戻しの議論で $\deg D\le2$ となる。
例外方向にあると分かる場合は $D$ が $L$ を含む平面に入るから、
$D\ne L$ より $D$ 上の $L$ の固定因子の長さ自体が $\deg D$ となり、同じ議論で $\deg D\le1$ である。
主曲線も直線なら緩い方の上界2を使えばよい。
ここで埋込みと数えるのは次数1,2の全て、初回から二番目と終回から二番目、
および反対三次曲線の像である。その他の三次曲線に埋込みを仮定しない。

以上と前定理の終点対応、次数和、三階差の符号を使う。単純な一純関係を両向きの例外中心が共有すれば
その二中心を二直線列で結ぶ。連続二直線列の排除は過剰接触を生じる中央辺の長さが1以下の場合に
両向きで適用する。同位相の二つの $(1,3,1)$ が一つの反対中心を通れない条件も両向きで使う。
添付コードは各中心に許される $l$ を全て選び、上の上限内の全列を相異なる終点に割り当てる。
結果は両向きとも相対二位相に一中心ずつで、次の候補に限られる。
一辺だけを許した場合は両方 $(1,3,1)$ であり、前定理の最後の証明により不可能である。
全辺長さ1以下の場合もこれと同じか、一方 $(1,3,1)$、他方
$(1,3,a,m,c,3,1)$ である。後者では $1\le a,m,c\le2$、
長い列の始点の順鉛筆と終点の逆鉛筆はともに $L_3$ 型である。

長い列の第二・第六三次曲線は同じ位相で相異なり、互いに相手を初回に持つ向きの反対三次曲線となる。
第三曲線と第五曲線は相手側の初回直線ではないので、これは二つの $L_3$ の二者択一に従う。
定理\ref{thm:branch-cubic-contact}の単元を用いた小行列式の包含により、
両者の交叉はそれぞれ逆・順の中心に欠損の長さで含まれる部分因子を持つ。
その長さの和は少なくとも $(m-3a+8)+(m-3c+8)\ge6$ である。
二中心は異なるので、相異なるねじれ三次曲線の交叉長の上界5に反する。
\end{proof}

% !TeX root = ../papers/nc-p3.tex
\section{二次曲線の平面の方程式が与える連続純関係}
\label{sec:conic-corridor-excess}
\begin{proposition}\label{prop:conic-corridor-excess}
曲面射影のない標準正則窓で、順中心 $p$ からの第一直線を $L$、第二曲線を平面二次曲線 $C$、
その平面の方程式を $h_C$ とする。$R_i$ の $L$ への制限は
$\ker p\otimes H^0(L,\mathcal O(1))$ への同型である。
$p$ を通る曲線枝において、順評価の三次小行列式の共通位数が点イデアルの位数を超えるなら、
ある非零 $s\in L^\perp$ に対して $s\otimes h_C\in R_{i+1}$ である。
列の第三次数も2なら、その第三曲線 $C'$ に対して $R_{i+2}$ は第二因子 $h_{C'}$ の純関係を含む。
後者には位数の仮定は不要である。逆方向にも成立する。
\end{proposition}
\begin{proof}
第一・第二点を独立に動かした縮約は $[s,t]A(t)$ と書ける。ここで $A(t)$ は
$L$ 上の唯一の順中心 $q$ の鉛筆で、三行四列の一次行列である
（定理\ref{thm:rank-ten-conic-after-center}の標準二列による同定）。
定理\ref{thm:finite-pure-relation-locus}と命題\ref{prop:finite-pure-pencil-opposite}より
一般階数は3である。中央支持3なら一般右核は固定で、第三点が動かない矛盾となるので
中央支持は4である。一般の $t$ ごとに $A(t)$ の像は $\ker p$ 全体で、
二つの異なる一般 $t$ で $[-,t]$ は独立である。従って制限の像は六次元であり同型を得る。
また第三曲線は $A(t)$ の $L_e$ ブロックの核曲線なので、その次数は $e$ である。
$e=2$ なら残る一行一列が与える一関係はその二項一次式と、残る第四文字との積である。
核曲線は最初の三座標の平面を張るから、第四文字がちょうど $h_{C'}$ である。

命題\ref{prop:excess-contact-two-line-pure}の分割した四行を使う。
制限が同型なので、枝の先頭接方向 $\xi\in(\ker p)^*$ と $L$ の二座標の対が
独立な除算後の二行を与える。過剰な位数があれば、これらは $p$ の残る二行と同じ二空間を張る。
二行を相補切断により $R_{i+2}$ の元として $C$ に制限すると、
符号と基底を除いて $t_0\xi A(t),t_1\xi A(t)$ である。
この二行の末尾の文字の支持は、一次行 $\xi A(t)$ の二係数が張る高々二次元空間に入る。
従って末尾から $w(p)$ を縮約した三空間 $P\subset L^\perp\otimes V_{i+2}$ を
$C$ に制限した像も高々二次元である。核には非零の元があり、
$C$ の線型方程式は $h_C$ の倍だけなので、これが求める純関係である。
\end{proof}

\begin{corollary}\label{cor:conic-corridor-consecutive}
$(1,2,1)$ の列の始点・終点の両方に、それぞれ順・逆の過剰な位数を持つ枝があることはない。
$(1,2,2,1)$ の列では終点に逆の過剰な位数を持つ枝があることはなく、始点についても同様である。
\end{corollary}
\begin{proof}
前者では中央二次曲線の方程式を中間文字として二つの純関係が連続する。
後者では第三曲線の平面の方程式について、順向きの命題の後半と逆向きの前半で同じ矛盾を得る。
連続純関係の禁止は標準正則性の必要条件である。
\end{proof}

% !TeX root = ../papers/nc-p3.tex
\section{二つの二直線列と一つの三次列の五重接触}
\label{sec:cubic-corridor-jets}
\begin{theorem}\label{thm:cubic-corridor-jets}
純関係の有無によらず、順中心を位置 $0,1,2$ に一つずつ持ち、そこから始まる列が順に
$(1,1),(1,1),(1,3,1)$ である標準正則窓は存在しない。
\end{theorem}
\begin{proof}
位置0のねじれ三次曲線上の順・逆欠損は次数公式により各5で、それぞれ相異なる唯一の点
$F_0,B_0$ に集中する。位置 $j=1,2,3$ の初回・終回直線を $L_j,M_j$ と書く。
これらはそれぞれ一点で交わる。実際 $L_1$ には $F_1$、$M_2$ には $B_2$ があり、
$F_1$ の後続直線 $L_2$ は $L_1$ からの対応点でも $M_2$ と会う。
他の二組は三次曲線上の $F_0,B_0$ の後点・前点で会う。
初回と終回が一致すれば移動点一つの列となるので異なる。
定理\ref{thm:final-pure-branch-frames}の三中心の座標を使い、
三次を $C(1,x)=(1,x,x^2,x^3)$、$F_0,B_0$ を $0,\infty$ とする。
$C$ の後点は $M_1(1,x)=(0,1,x,0)$、前点は $L_3(1,x)=(1,x,0,0)$ である。
以下 $v_j$ を位置1の四つの座標関数、$y_j$ を位置3の四つの座標関数とする。

$R_0$ の $L_1$ への制限は完全な六次元である。よって三次上でその関係はちょうど
\[
 p_1h+p_2k+v_3D(h,k),\qquad
 (p_1,p_2)=(xv_0,xv_1-v_2),\qquad \deg h,\deg k\le2
\]
と表せる。実際 $v_0,v_1$ の係数は $F_0$ で消える任意の三次式であり、$v_2$ の係数は
後点から決まる。$D$ の値は次数3以下である。四つの評価行の三次小行列の共通因子は、
$h,k,D$ の評価三行の三次小行列の共通因子に原点でちょうど $x$ を掛けたものである。
$k$ の行は $v_2$ に単元で現れ、$h$ の行は $v_0$ に $x$ 倍で現れるからである。
従って $D(h,k)=[H(x)(h,k)^{\mathsf t}]_{<4}$ となる。
ここで $H$ は二成分の行、括弧は四次以上を捨てる操作である。実際 $h,k$ の定数基底列に
単位小行列があり、それらの列から読む $H$ で全列について $x^4$ を法として一致する。
原点で評価階数2なので $H(0)$ の第一成分は非零である。第四座標方向の倍率と
最初の三方向による平行移動で
\[
 H=H_0+x^2 H_2+x^3 H_3,\qquad H_0=(1,0)
\]
とできる。消した三係数は第一成分の一次係数、第二成分の定数と一次係数である。
$U=(H_2^{\mathsf t},H_3^{\mathsf t})^{\mathsf t}$（行は $H_2,H_3$）は非零である。
零なら三次上で一般評価階数が2に落ちる。
逆に $t=1/x$ で $R_3$ を読むと、座標 $0,2$ を交換した同じ表示になり、
列 $J(t)=J_0+t^2J_2+t^3J_3$、$J_0=(1,0)^{\mathsf t}$、$V=(J_2,J_3)\ne0$ を得る。

両辺の三文字重なりの元を $C(1,x)$ で評価する。$v_3,y_3$ を切除した部分は一意に
\[
 (p_1,p_2)(E_0+xE_1)q,\qquad q=(y_2,y_1-xy_0)^{\mathsf t}
\]
と書ける。実際 $h,k$ の $y_0$ の係数は、前点の消滅から $y_1$ の係数の $-x$ 倍なので、
$y_1$ の係数は一次以下である。$y_2$ の係数も一次以下である。
$B_0$ で $M_3$ を消すから $y_2$ の三次係数が零で、$p_1,p_2$ の一次係数は独立だからである。
この切除は重なりに単射である。一方の辺の三座標への制限は初回直線への単射を含み、
他辺も終回直線への単射を含むので、順に第四座標の行・列を復元できる。

$v_3$ の行を $A(x)q$ とすると、$A$ は二次以下でなければならない（前点の消滅と
$B_0$ での $y_2$ の三次係数の消滅）。そのための条件と係数は
\[
 H_2E_1+H_3E_0=0,\qquad (A_0,A_1,A_2)=(H_0E_0,H_0E_1,H_2E_0).
\]
逆側で $y_3$ の列を $(p_1,p_2)D'(x)$ とすると
\[
 E_0J_2+E_1J_3=0,\qquad (D'_0,D'_1,D'_2)=(E_1J_2,E_0J_0,E_1J_0).
\]
最後の $v_3y_3$ の係数は $[HD']_{<4}$ と $x^3[(t^2 A(1/t))J(t)]_{<4}|_{t=1/x}$ の
両方に等しい。その四係数を比較すると
\[
 H_2E_0J_3=H_2E_0J_2=H_2E_1J_2=H_3E_1J_2=0 .
\]

$U,V$ ともに可逆なら $X_j=UE_jV$ と置く。$X_0$ の第一行、$X_1$ の第一列が零で、
二つの線型条件が残りの行・列も零にする。$U$ が可逆で $V$ が階数1なら
$J_2=v\beta,J_3=v'\beta$ とする。$v\ne0$ なら最後の二条件から $E_1\beta=0$、
次いで $E_0\beta=0$。$v=0$ なら列条件から $E_1\beta=0$、混合条件と行条件から
$H_2E_0\beta=H_3E_0\beta=0$ で同じ結論となる。両行列が $\beta$ を消す四次元空間で
行条件は非自明なので重なり次元は高々3である。逆の場合も同じである。
いずれも位置2による独立な四縮約が重なりに入ることに反する。

残る階数1同士で $H_2=u\alpha,H_3=u'\alpha$ と書く。$uv\ne0$ なら混合条件から
$\alpha E_0\beta=\alpha E_1\beta=0$。$u=0$ なら行条件が $\alpha E_0=0$ を与え、
$v\ne0$ なら最後の混合条件、$v=0$ なら列条件から $\alpha E_1\beta=0$ を得る。
$v=0,u\ne0$ でも列条件と最初の混合条件を用いる。
従って $\alpha E_0,\alpha E_1$ はともに $\beta$ を消す一次元空間に属し、
$u\alpha E_1+u'\alpha E_0=0$ なので、その組の像次元は高々1である。
位置2の四縮約のうち三次元以上の部分空間で $E_0,E_1$ の像は固定直線
$k\lambda=\ker\alpha$ に入る。この部分では全ての係数が
\[
 x(\lambda_1v_0+\lambda_2v_1)+\lambda_1v_3-\lambda_2v_2
\]
で割れる。実際 $y_3$ も含め、対応する二次以下の列は $\lambda$ の倍であり、
$H_2\lambda=H_3\lambda=0$ なので第四行の復元に切捨てがない。
この一次式の $x$ の両係数を消す非零の位置1の点を取る。三次曲線は位置0全体を張るから、
この点で縮約した窓は位置2の三次元以上の点空間を消し、隣接支持が高々1になる。
標準正則窓では隣接評価階数は常に2以上であり、矛盾する。
\end{proof}

% !TeX root = ../papers/nc-p3.tex
\section{全ての中心鉛筆を許した曲線分岐の排除}
\label{sec:unrestricted-corridor-exclusion}
\begin{theorem}\label{thm:unrestricted-corridor-exclusion}
全射影の次元が1以下の標準正則窓には順・逆いずれの分岐中心もない。
\end{theorem}
\begin{proof}[次数・中心型の有限な範囲]
分岐があるとして、各位相の順・逆中心数を $F_i,B_i$ とする。
定理\ref{thm:finite-pure-relation-locus}と命題\ref{prop:finite-pure-pencil-opposite}より
各々高々2で、各中心に主曲線次数 $e=0,1,2,3$ と例外方向高々 $3-e$ がある。$e=0$ は主像が定点の場合である。
同位相同方向に二中心があれば両方 $e>0$。両向きで少なくとも二位相が占有される
（定理\ref{thm:single-pure-branch-exclusion}の第四次数の下界）。
順中心と逆中心を実際の列により全単射で対応させる。中心から $m$ 個の動く曲線を経て終点に至るので
終点位相は始点位相 $i$ に $m+1$ を足したもの。位相ごとの全次数和は7以下である。
ある初回移動の位相では、それ以外の全移動は主曲線一つ以下または例外直線に着く。
各成分を二度使えず、これらの総数は高々3である（$e=0$ なら例外だけ）。
よって同位相の移動は高々4回、$2\le m\le20$ である。

各列の $d_1=d_m=1$、$d_2=3-f_1$ と逆側の同式、および
命題\ref{prop:branch-corridor-degree}の三階差を用いる。$m>2$ なら $d_2,d_{m-1}\ge2$。
初回以外の移動を各該当中心の主像または例外直線に割り当てる。別の中心の初回なら主像である。
動く曲線からの移動と別に、次位相の初回直線外の中心は各々一つの例外直線を占める。
その数は少なくとも $F_{i+1}-f_1$ であり、これらの直線は一般前点が固定なので移動先と重複しない。
$d_2=2$ なら命題\ref{prop:conic-corridor-excess}の $A(t)$ により初回直線上の中心の $e$ は $d_3$ である。
非初回の移動でその前の写像が埋込みなら、始点曲線の次数は主像へ行く場合は2以下、
例外直線へ行く場合は1以下である（定理\ref{thm:sparse-pure-edge-exclusion}の射影の証明）。
ここで埋込みと数えるのは次数1,2、初回と終回から二番目、または主像の有理正規三次曲線である。

始点の中心 $p$ の初回直線 $L$ と、位相 $i+2$ の逆中心 $z$ の前点直線 $M$ を比較する。
交わらなければ $M$ からの $L$ 中心の射影は全単射である。$M$ では後点が常に $z$ なので
$e>0$ の主曲線と両立せず、$e=0$ でも $z$ はただ一つの固定主像に限られる。
交点があればその後点 $z$ は $p$ の第二曲線にあり、逆欠損に少なくとも1寄与する。
よって $b_2\ge B_{i+2}-\mathbf1_{e=0}$ を得る。逆向きにも同じである。
各順欠損の上界は $3d_j$ で、順中心のない位相では0である。
埋込み曲線上で、その位相の全順中心が二直線列の始点で各列のいずれかの端の $e$ が2以上なら、
上界を $F_{i+j}$ に下げられる。実際過剰接触から命題\ref{prop:excess-contact-two-line-pure}で得る
二関係は両端の縮約三空間に属し、両者に長さ2以上の純関係を与えて $3-e\le1$ に反する。
逆欠損についても同様である。
\end{proof}

\begin{proof}[配置の列挙と矛盾]
同位相の初回直線は互いに交わらない。従って $f_1>0$ の列の本数は次位相の順中心数以下である。
$f_2>0$ の列も同じ二位相先の順中心を共有できない。共有すればその一意な前点が
二本の初回直線に属すからである。$f_2=3d_2-3-d_3$ であり、二直線列の場合は $d_3=0$ とする。
同位相に二直線列が二本あれば各始点で $e>0$、かつその終回直線には逆欠損2がある。
上の直線交叉の議論から、この位相の全逆中心が両方の終回直線に属してしまう。
終回直線も互いに交わらないので不可能である。これらを逆向きにも使う。

初回二番目の三次曲線と、同位相に来る終回から二番目の別の三次曲線も比較する。
両始終端の $e$ が3、両列の長さが4以上なら、それぞれを定める二つの移動元は
他方の初回直線と異なるので、互いに相手の主三次曲線である。
同じ列なら二つの成分の一致は列内の再訪となる。別列でも一般前点の一意性により成分は異なる。
定理\ref{thm:branch-cubic-contact}の単元による小行列式の包含はそのまま使える。
前者の逆欠損、後者の順欠損は互いに素な台で両曲線の交叉に含まれるので、その和は5以下である。
両側の次数列を $d,r$（後者は終点から逆順）と書くと、その和の下界は
$(d_4-3d_3+8)+(r_4-3r_3+8)$ である。

以上の整数条件の全列挙は付属コードで行う。まず総数の等しい二組の位相数を全て取り、
各始終点・長さについて許容次数を小さい添字から追加し、確定した三階差の符号を照合する。
次に相異なる終点への完全マッチングを作り、位相次数和と各中心の移動先の本数を加算する。
各 $e$ を全て割り当て、三階差が欠損の上下界の差に入るかを検査し、上記の交叉を検査する。
全列が $(1,3,1)$ の場合は相対二位相各一中心か各二中心に限られ、前者は
定理\ref{thm:single-pure-branch-exclusion}末尾の二補完式、後者は全中心 $e=3$ なので
命題\ref{prop:branch-local-pure-pencil}で消える。全列が二直線列の場合は各位相一中心で
系\ref{cor:four-two-line-impossible-with-pure}が適用される。残る全行を示す。
位相番号を揃えるためだけに巡回しており、添字移動で対象を同一視していない。
\[
\begin{array}{c|l}
\text{順中心の位相（列順）}&\text{列の次数列}\\\hline
1,2,3&11,11,131\\
1,2,3&121,12321,131\\
0,1,2,3&121,121,121,121\\
0,1,2,3&1221,11,1221,11
\end{array}
\]
第一行は定理\ref{thm:cubic-corridor-jets}に反する。
第二行は位相1の順・逆中心で始終する $121$ を持つ。位相3の $131$ の中央三次曲線上の
順・逆欠損は各5で、位相1の各向きの唯一の中心に集中する。第三行も各 $121$ の中央二次曲線上の
両欠損が各2である。いずれも埋込みなので両端で過剰となる $121$ があって
系\ref{cor:conic-corridor-consecutive}に反する。第四行は位相0の $1221$ の始点に、
位相3の二直線列の初回直線から順欠損2が集中し、同じ系に反する。
\end{proof}

\begin{corollary}\label{cor:all-positive-dimensional-window-decision}
標準正則窓の四辺は全て純テンソル関係を含まない。また少なくとも一射影が正次元の
標準有限Frobenius候補で、全次数の標準Koszul左右AS正則性と、正則な窓間の添字保持同型は有限に判定できる。
\end{corollary}
\begin{proof}
純関係の不在は本定理と定理\ref{thm:single-pure-branch-exclusion}冒頭の曲面・一点グラフ・有限辺の議論による。
従って純関係をSegre交叉で検査して非空なら排除でき、空なら系\ref{cor:positive-dimensional-pure-free-decision}を使う。
\end{proof}

% !TeX root = ../papers/nc-p3.tex
\section{有限点窓の三文字評価の階数差と点列のイデアル}
\label{sec:finite-three-slot-evaluations}
\begin{theorem}\label{thm:finite-three-slot-evaluations}
標準有限Frobenius候補の四辺が有限とする。共通点環 $D$ における連続三位置の座標を
$U,V,W$ とし、単元倍の違いを許す。反対位置 $U,W$ が射影同定されれば
$\dim UVW=16$ である。同定されないときは19または20であり、
階数19での三積の唯一の零化汎関数の支持スキームの長さは9以上である。
従って長さ3の経路から点列への評価の核は、順に次元4、1、0で尽きる。
\end{theorem}
\begin{proof}
$W=U$ に規格化できる場合は辺の六式の完全交叉のKoszul分解を $(2,1)$ でひねる。
二つの斉次次数がともに非負なのは初めの二項のみで、最後まで高次コホモロジーから
大域切断の核へ入る項はない（両次数とも $-4$ 以下となるのは第六項だけで全次数0）。
従って像の次元は $10\cdot4-6\cdot4=16$ である。

他の場合は定理\ref{thm:finite-opposite-product-gap}から $\dim UW=16$ である。
$D$ の完全積対合の跡を一つ選び、$E=(UW)^\perp$ とすると $\dim E=4$。
元 $a\ne0$ の支持は $Y_a=\operatorname{Spec}(D/\operatorname{ann}(a))$ とし、
その長さを $m_a$ とする。各座標空間は単元を含み、任意の長さ $m$ の部分スキームへの
制限次元には、$m\ge1,2,4,8$ に応じてそれぞれ $1,2,3,4$ の下界がある。
これは各位置で点・直線・平面の切断長がそれぞれ1,3,7以下であるためである。
$a\in E$ なら $Y_a$ 上の $U,W$ が $a$ による完全跡で直交する。
上の下界を足すと $m_a\ge6$ が必要である。

三積の零化元を $g$ とすると条件は $gV\subset E$ である。$0<m_g\le7$ と仮定し、
$r=\dim(V|_{Y_g})$ とする。$Y_g$ の長さ $r-1$ の部分スキーム上で零で、$Y_g$ 上では非零な
$v\in V$ が存在する（Artin環のイデアルの組成列を使う）。すると $gv\ne0$、
$m_{gv}\le m_g-r+1\le5$ となり $E$ の下界に反する。$m_g=8$ も不可能である。
そのとき $Y_g$ 上の $UV$ の次元は補題\ref{lem:finite-adjacent-strict-half}で5以上、
$W$ は4であり直交できない。従って $m_g\ge9$、特に $gV=E$ となる。

さらに $g'\ne0$ も三積を消せば、両者が生成する $D$ のイデアルは
$DE=(g)=(g')$ である。ゆえに $Y=Y_g$ 上で $g'=ug$ かつ $u$ は単元で、
$u(V|_Y)=V|_Y$。このような $u$ は定数しかない。実際非定数なら剰余値の一つを引き、
$V|_Y$ を保つ非零・非単元 $a$ を取る。$r=\dim a(V|_Y)$ は、その支持への $V$ の制限次元に等しい。
また $Y$ 上の $a=0$ への制限次元は高々 $4-r$ である。
支持と零点スキームの長さの和は $\operatorname{length}Y$。どちらも空でないため
$r=1,2,3$ であり、先の部分空間切断長の上界からその和は
$1+7,3+3,7+1$ 以下である。これは $m_g\ge9$ に反する。
従って零化空間は高々一次元である。最後に候補の三次重なりは次元4なので
経路空間の次元は $4^3-2\cdot 4\cdot6+4=20$ である。
\end{proof}

\begin{corollary}\label{cor:finite-asymmetric-point-ideal}
両方の反対対で射影同定がない有限窓とし、$A$ から共通点列の切断の積への評価の核を $I$ とする。
$I$ の各始点での右イデアルは次数3,4の部分で生成され、左についても同様である。
各始点で三次生成元は高々一本である。全次数のAS性を仮定しない。
\end{corollary}
\begin{proof}
始点を固定した像を $S_n\subset D$ と書く。定理\ref{thm:finite-four-slot-rank}により $S_4=D$、
以後も各座標に単元があるので $S_n=D$ である。$I$ の長さ $n>1$ の右極小生成元の空間は、
$A$ が二次表示なので、点列の切断に関する次の三項の中央ホモロジーである。
\[
 S_{n-2}\otimes R_j\longrightarrow S_{n-1}\otimes V_{j+1}\longrightarrow S_n ,
 \qquad j=i+n-2 .
\]
実際これは $S$ と終点単純の左自由表示をテンソルした第一Torである。
全項の係数空間が $D$ なら点グラフ上の関係の左評価は階数3の直和因子であり、
次の点での評価の核を張るので完全である。これで $n\ge6$ は済む。

$n=5$ で $S_3$ が超平面の場合だけ残る。完全跡で $S_3^\perp=kg$ とし、$Y=Y_g$ とする。
$D$ 上の全体の写像 $\alpha:D\otimes R_j\to D\otimes V_{j+1}$ の像に、
$S_3\otimes R_j$ の像が届かないとすると、両者の余核の双対を比較して、
$0\ne g\otimes c\in\operatorname{im}\alpha^*$、$c\in R_j^*$ が存在する。
実際この制限で消える汎関数の全体は $kg\otimes R_j^*$ である。
相補切断により $R_j^*\simeq R_{j+2}$ とする。$\alpha^*$ の像は
$w(x_j,-,-,-)$ の次文字による縮約で張られるので、どの元も最後の点 $x_{j+3}$
で評価すると零である。これは四移動が点スキーム上で恒等で、辺 $x_{j+3}\to x_j$ が成立するからである。
従って $c$ の末尾支持の文字は $Y$ 上で全て零になる。ところが前定理より
$Y$ の長さは9以上で、一平面に入らないので全て零、矛盾となる。逆向きも同じである。
\end{proof}

\begin{corollary}\label{cor:finite-full-triple-regular-presentation}
十五次元混合成分の幾何学的一般点では全ての三文字評価が同型である。
より一般に全三文字評価が同型の有限窓では、各始点の点列イデアルは四次生成であり、
その極小生成元の間の右関係は全て五次で生成される。左でも同様である。
この関係の本数は各位置で高々24。標準正則なら四次元の窓ごとに
生成元は15個、関係は24個、これらの間の第二関係は六次の10個で尽き、それより高い関係はない。
\end{corollary}
\begin{proof}
定理\ref{thm:marked-twenty-rational-invariants}で構成した二十印の持上げにおいて、
各位置の三文字の全単項式の評価行列を取る。付属検算は素数1009の整数環を十四桁で切った
環上で単元ピボットを消去する。全四位置でピボットは20あり、二位置では全て単元、
残る二位置でも付値は16個が0で4個が1である。従って行列式小行列は標数零の一般側で非零となる。
階数20は開条件なので同成分の一般点に成立する。

各 $S_n$ は $n\ge3$ で $D$ である。三次核が零なので先の系で四次生成である。
終点単純の左Koszul候補複体 $K$ と $S$ をテンソルする。$K$ は二次表示と全三次重なりを使うため、
$\operatorname{Tor}_2(S,k)$ はこの第二ホモロジーの商である。実際、真の第三自由項には
全ての三次重なりが入り、それ以外の生成元があれば境界が増えるだけである。
$n\ge6$ では項 $1,2,3$ の係数が全て $D$ なので、点列での階数 $1,3,3,1$ から
第二ホモロジーは零である。四次生成の極小関係の次数は5以上なので五次だけが残る。

$n=5$ では $S\otimes K$ の項の次元が右端から $20,80,120,40,4$。
最後の写像は単射で、第三ホモロジーも零であることを示す。
全係数を $D$ に広げると第三項の微分の核は、前の点の四座標を並べた列の $D$ 倍である。
これは点列での分裂した完全性と、第四微分がその座標列であることによる。
元の第三項では各係数が $UV$（初めの二位置）に入るので、核の倍率 $a$ の条件は $aV\subset UV$。
$(2,2)$ の完全跡に関し $UV=(UV)^\perp$、また $(1,2)$ の評価次元は上の完全交叉から16である。
従ってこの倍率空間は $(UV^2)^\perp=U$ であり、正確に元の第四項の像である。
第一ホモロジーも前の系で零なので、第二はEuler式により24となる。

正則なら $K$ は単純の解消そのもので、$I$ の一次関係はちょうどこの24個である。
さらに $n\ge7$ なら全項の係数が $D$ で全て完全である。
最小次数と、標準Hilbert関数の35に対する $S_4$ の次元20の差で四次生成は15個。
第三Torの唯一の可能次数6での次元はHilbert分子
$(1-z)^4(1+4z+10z^2+20z^3/(1-z))=1-15z^4+24z^5-10z^6$ で10となる。
第四Torは生成可能な次数が7以上なので存在しない。
\end{proof}

% !TeX root = ../papers/nc-p3.tex
\section{四頂点切断における点加群と分離三次曲面の特異点}
\label{sec:finite-cut-point-pairings}
\begin{theorem}\label{thm:finite-cut-point-pairings}
四辺が有限な標準有限Frobenius候補を取り、連続四頂点 $i,\ldots,i+3$ の経路代数を $B$ とする。
点グラフの点 $p$ は、各頂点一次元、三辺の作用が $p_i,p_{i+1},p_{i+2}$ である右 $B$ 加群 $M(p)$ を定める。
元の無限代数のAS性によらず、$\nu_B M(p)\simeq M(p)[3]$ である。
異なる点 $p,q$ について
\[
 \dim\operatorname{Ext}^1_B(M(p),M(q))
 =\dim\operatorname{Ext}^2_B(M(p),M(q))
 =2-\operatorname{rank} w(p_{i+3},-,q_{i+1},-)
\]
であり、他の次数は零である。この階数は $p,q$ の交換にも切断の始点にも依存しない。
自己Extの次数0,3は一次元、次数1,2はその局所点の埋込み次元である。
\end{theorem}
\begin{proof}
$H_j=\ker p_j$、$P_j=e_jB$ とし、$h=w(p_{i+3},-,-,-)$ は添字を並べ直して
$H_i\otimes V_{i+1}\otimes H_{i+2}$ と読む。次は右加群の線型射影解消である。
\[
 0\to P_{i+3}\to H_{i+2}^*\otimes P_{i+2}
 \to H_i\otimes P_{i+1}\to P_i\to M(p)\to0 .
\]
両端の微分は指定した超平面の三文字で、中央は $h$ による。
第二自由項の係数は $R_i\cap(H_i\otimes V_{i+1})$ に等しい。
実際 $p_i$ での評価が階数3で、$h$ の末尾縮約がこの核内の独立な三列を与える。
最初の余核は点加群である。先頭の三文字を消した後は関係の評価が
次の点の消滅三平面全体を与え、残りの各段でも同じことを繰り返せる。
長さ2まではこの核の記述で完全である。長さ3での第二微分の核は
三次重なりのうち第一位置が $H_i$ に入る元である。全三次重なりは $w(s)$ であり、
条件は $w(s,p_i)=0$ なので $s$ は唯一の前点 $p_{i+3}$ の倍である。
長さ3の各項次元は左から $1,12,30,20,1$ であり、これで完全性が従う。
左加群についても同じ解消があり、右の表示に $\operatorname{Hom}_B(-,B)$ を施したものは
$h$ の中央行列を逆順に読んだその左解消である。従ってコホモロジーは次数3の左点加群だけで、
体上の双対を取ればSerre関手の式を得る。

右解消に $\operatorname{Hom}_B(-,M(q))$ を施すと次元 $1,3,3,1$ の列になり、
中央の微分が表示の反対切断である。異なる点なら両端は階数1で、同じ点なら両端は零である。
自己第一Extの次元は点グラフの接次元である。実際、双対数上への加群の持上げは各頂点で
自由階数1なので基底を選んで三つの非零な矢評価を持ち上げることと同じであり、
頂点基底変更の有効な三倍率は各評価の射影化と一致する。二辺の零点条件で三文字の列は
共通有限スキームと同型である。自己第二Extは同じ次元となる。
交換の不変性はSerre双対とこの二次数の一致による。

位相の不変性には相補切断の完全対合を使う。例えば位相0と2で、
$p_0,-$ と $q_2,-$ により得る二つの三空間の対合の階数が表示の階数である。
それぞれの直交補は辺 $p_3\to p_0, q_1\to q_2$ を使い $-,p_3$ と $-,q_1$ で得る三空間である。
六空間間の完全対合における三空間の対合と直交補の対合は同階数なので
位相3と1の反対切断も同じ階数である。巡回して結論を得る。
\end{proof}

\begin{corollary}\label{cor:finite-separator-singularity-graph}
全ての中央支持が4で二十点が被約なら、いずれの位置の二十点の配置からでも同じ無向グラフが得られる。
各点 $p$ に対し他の十九点で消える唯一の三次曲面を取り、
別の点 $q$ がその特異点となるとき $p,q$ を結ぶ。各頂点の次数は4以下で、
辺は切断上で点加群間の第一Extが一次元となることと同値である。
十五次元混合成分の一般点では辺はなく、二十の点加群は互いに全次数で直交する
三次元球面的対象である。さらにその直和は四頂点切断の忠実表現である。
\end{corollary}
\begin{proof}
分離三次曲面は定理\ref{thm:finite-cubic-separators}の行列式である。異点 $q$ における行列は
階数1または2である。0は全中央支持4に反する。階数1なら余因子は全て零である。
階数2なら、その三行三列での左右核は二つの隣り合う位相における
$q$ の座標を $p$ から射影したものである。従って余因子による微分は、倍率を除き
$w(p_0,q_1,-,q_3)$（曲面の位置を2とした）である。$q_1,q_3$ の反対切断は
$q$ が単純なので階数3で左核が $q_0$ のみである。$p_0\ne q_0$ なので微分は非零である。
これで前定理の交換・位相の不変性とExtの式がグラフに移る。

一つの四次元切断の階数一スキームが有限なら長さは4以下
（定理\ref{thm:finite-marked-secant-pencil}のSegreの完全交叉の計算）。
正次元なら一つの被約直線である（定理\ref{thm:finite-marked-ruling-pencil}の
左右評価階数の証明）。両議論のこの部分は相手の四次元切断との一致を使わない。
後者では元の点配置の直線切断長が3以下なので次数も3以下である。

混合成分の一般点の全二十印と全位置で三次曲面が滑らかなことは
定理\ref{thm:marked-cubic-surface-pairs}による。従って自己Extは次数0,3のみ各一次元、
Serreによる3移動と合わせて球面性が従う。異点間のExtは全て零である。
点加群の直和の表現核は各隅ごとに点列評価の核である。二次まで単射であり、
三次も系\ref{cor:finite-full-triple-regular-presentation}により一般には全位置で単射である。
\end{proof}

% !TeX root = ../papers/nc-p3.tex
\section{有限片側対称窓の三段加群と電荷四のモナド}
\label{sec:finite-symmetric-instanton-monad}
\begin{theorem}\label{thm:finite-symmetric-instanton-monad}
標準有限Frobenius窓の全辺が有限で、一方の反対交換が対称、反対階数 $\rho>6$ とする。
文字を $U,V,U,Z$ と並べ、$R=R_0\subset U\otimes V$、$K=R^\perp$、
$S=\operatorname{Sym}U^*$ と置く。可換次数加群
\[
 N=(S\otimes V^*)/(S K)
\]
はHilbert列 $(4,6,4)$ の有限長加群であり、$N_2=Z$ という同定を除き、
$w$ は二回の乗法 $S_2\otimes N_0\to N_2$ そのものである。
$N$ の極小自由分解は
\[
 0\to S(-6)^4\to S(-5)^{10}\to
 S(-4)^4\oplus S(-3)^4\oplus S(-2)^4\to S(-1)^{10}\to S^4\to N\to0 .
\]
その最初の線型部分から $\mathbb P(U)$ 上で
\[
 L\otimes\mathcal O(-1)\xrightarrow{a}K\otimes\mathcal O
 \xrightarrow{b}V^*\otimes\mathcal O(1),\qquad
 L=(\textstyle\bigwedge^2 U^*\otimes V^*)\cap(U^*\otimes K),\quad \dim L=4
\]
を得る。$b$ は全射である。$a$ が全点で階数4の開部分ではコホモロジー層は
電荷4の安定な階数2インスタントン束 $\mathcal E$ であり、
$H^1(\mathcal E(n))=N_{n+1}$ は全ての整数で乗法と両立する。
この開部分の窓の添字保持同型は、軸 $U$ を位置0に取ったこれらの束の射影同値と必要十分である。
\end{theorem}
\begin{proof}
$N_0=V^*, N_1=R^*$ であり、$N_2^*$ は三次重なりのうち両端対称な部分である。
$R_1=R^{\rm op}$、全三次重なりは $w(Z^*)$ で既に対称なので $N_2=Z$ を得る。
高次項を調べる。$0\ne t\in N_3^*\subset\operatorname{Sym}^3 U\otimes V$ があれば、
その最小 $U$ 支持を $H$ とする。各 $0\ne\phi\in H^*$ で一回の縮約は非零の $w(z^*)$ である。
この両端の $U$ 支持は有限辺の評価階数から3以上なので $\dim H\ge3$。
$\dim H=3$ なら二回の縮約は全て $P=R\cap(H\otimes V)$ に入り、
$H^\perp$ の点での評価階数が3以上なので $\dim P\le3$。
各 $\phi$ を固定した第二縮約 $H^*\to P$ は階数3で、全非零 $\phi$ で可逆となる。
これは射影平面上の三次行列式の零点で矛盾する。従って $H=U$ で、一回の縮約は
$U^*\simeq N_2^*$ を与える。これで $Z$ を同定すると $w$ は三位置で対称になり、
三つの二対二切断はすべて階数6になってしまう。ゆえに $N_3=0$、以後も零である。

反対の辺 $U,Z$ から作る同じ三段加群は $N$ の次数を反転した体上の双対に等しい。
実際、次数0と2は $Z^*,V$ であり、二乗の作用は同じ $w$ の転置である。
中間の六空間は一段目の作用の核を割って得られ、隣接切断は完全対合なので一致する。
従って両側とも四生成・十本の一次関係のみを持つ。四変数のKoszul分解による
有限長加群とその双対のBetti対称性（外次数 $j$ と $4-j$ の完全対合）により
自由分解の最初と最後の二項が決まる。残る中央項は
$(4+6z+4z^2)(1-z)^4=4-10z+4z^2+4z^3+4z^4-10z^5+4z^6$
から表示の通りとなる。

写像 $b$ は $K$ の十元の一次関係であり、その二次シジジーが $L$ である。
有限長なので層化した $b$ は全射である。$a$ が部分束なら
$\mathcal E=\ker b/\operatorname{im}a$ は階数2で、
$c(\mathcal E)=(1-h)^{-4}(1+h)^{-4}=1+4h^2$ である（$h$ は超平面類）。
十本が独立なので $H^0(\mathcal E)=0$、従って安定である。
実際傾きが非負の階数1部分層があればその反射包 $\mathcal O(d)$ からの非零写像が延び、
$d\ge0$ なので大域切断を与える。二つの短完全列と直線束の中間コホモロジーの消滅により、
$H^1(\mathcal E(n))$ は $K\otimes S_n\to V^*\otimes S_{n+1}$ の余核である。
とくに $H^1(\mathcal E(-2))=0$。これでインスタントンの条件が揃う。
逆にこのコホモロジー加群から $N$、それから二回乗法と $w$ を取り戻せる。
また $N$ からは極小表示の線型部分が基底変更まで一意なので束も戻る。
射影空間の座標変更は $U$ の基底変更、初端と終端の加群の基底変更は $V,Z$ の基底変更である。
軸の同定のスカラー自由度はこれらに含まれる。
有限グラフでの反対同定は射影的には一意である。実際全点を固定する射影変換は、
長さ20の座標四空間を保つ単元倍となり、定理\ref{thm:finite-three-slot-evaluations}の
長さ9以上の安定化単元の証明により定数である。従って余分な印は付かない。
\end{proof}

\begin{corollary}\label{cor:instanton-window-reconstruction}
上の束の集合は次の逆構成で尽くせる。
線型モナド $4\mathcal O(-1)\to10\mathcal O\to4\mathcal O(1)$ のコホモロジー束について
$H^1_*(\mathcal E)$ を取り、$N_0,N_1,N_2$ が指定次元であること、その二回乗法を四窓として
隣接切断・有限Frobenius・有限辺・$\rho>6$ の各有限条件で検査する。
ここではモナドの両端が束として完全、定数列の一次関係は独立とする。
この範囲には全次数正則な窓が実在する。
\end{corollary}
\begin{proof}
モナドから $N$ は既に四生成十一次関係の表示を持っている。$N_1$ が六次元で
二回乗法の隣接階数が6ならその一次関係は $R^\perp$ である。
生成次数2の四シジジーが全ての線型シジジーを張ることは上の自由分解と $a$ の単射から従う。
こうして再構成したモナドも元に戻る。定理\ref{thm:twenty-point-distinct-companion}の
反対階数10の正則窓で、付属計算は $a$ の四次小行列式が四変数の全四次式を張ることを
素数1009で検証する。従って標数零でも共通射影零点はなく、この開部分は空でない。
\end{proof}

% !TeX root = ../papers/nc-p3.tex
\section{モナドの退化集合の余次元と反射包}
\label{sec:finite-instanton-sheaf-boundary}
\begin{theorem}\label{thm:finite-instanton-sheaf-boundary}
定理\ref{thm:finite-symmetric-instanton-monad}の全窓で、$a$ の階数降下集合は高々一次元である。
従って同じモナドは常に捩れのない階数2のインスタントン層 $\mathcal E$ を与え、
その $H^1_*$ から窓が復元される。反射包 $F=\mathcal E^{**}$ は局所自由であり、
$Q=F/\mathcal E$ は零か純一次元で $H^*(Q(-2))=0$ である。
$d=\deg Q$ とすると $c_2(F)=4-d$、$d\le4$ で、$d=4$ なら $F\simeq\mathcal O^2$ である。
\end{theorem}
\begin{proof}[退化するシジジーの根基]
$0\ne t\in L\subset\bigwedge^2 U^*\otimes V^*$ を取り、$t$ を $V$ で縮約した二形式たちを考える。
$[s]\in\mathbb P(U)$ で $a(s)t=0$ となる条件は全二形式の根基が $s$ を含むことである。
このとき $t$ は $\bigwedge^2(s^\perp)\otimes V^*$ にあり、二形式の包は三次元全部である。
実際、その包が次元2以下なら全ての二形式が共通の一次因子 $h$ を持つ。
次元2以上でそのような因子があれば $h^\perp$ で $t$ を縮約して得る
$kh\otimes V^*$ 内の像は少なくとも二次元である。次元1なら一個の非零二形式を縮約して
$U^*\otimes kv$ 内に二次元以上を得る。いずれも $K$ に純テンソルの直線を与えて有限辺に反する。
この議論は共通の一次因子を持つ全ての場合に適用できる。
ゆえに $s$ が存在すれば一意で、三次元の二形式の包から正則に復元できる。

このような $[t]$ の集合 $T\subset\mathbb P(L)$ は、全ての二形式同士の外積を係数とする
二次式の共通零点集合である。実際この消滅条件では線型包がKlein二次曲面
$\mathrm{Gr}(2,U^*)\subset\mathbb P(\bigwedge^2 U^*)$ に入る。
そのような線型空間は共通一次因子を持つか、一つの三空間の外積二乗に入る。
これは相交わる二つの二平面を固定し、両方と交わる残りの二平面が
その交点を含むかその和に入るかを見ることで従う。前者は既に排除した。
従って $T$ 上には根基射 $s:T_{\rm red}\to\mathbb P(U)$ があり、像が階数降下集合である。
そのファイバーは $a(s)$ の核の射影空間である。

$T$ 全体が射影三空間、またはその中に射影平面 $D$ を含む場合を考える。
この空間上で $s^*\mathcal O(1)=\mathcal O(m)$、$H=\ker(U^*\otimes\mathcal O\to\mathcal O(m))$ とする。
全射 $V\otimes\mathcal O\to(\bigwedge^2 H)(1)$ があるので、右辺の行列式次数 $3-2m$ は非負である。
$m=0$ なら固定三行四列行列の空間内で次元3以上の線型系が全点で階数3になるが、
階数2以下の行列は余次元2なので不可能である。従って $m=1$ で、根基射は線型埋込みとなる。
像のベクトル空間を $B\subset U$ とする。$\dim B\ge3$ で、$B$ 上の線型族は
$t(s)(v)=\iota_s\omega_v$、$\omega_v\in\bigwedge^3U^*$ と書ける。
実際 $\iota_s t(s)=0$ を偏極すると一つの引数が $B$ に入る三重線型値は完全交代であり、
補空間は高々一次元なので全三形式に延びる。体積形式で $\omega_v$ を $U$ の元と読むと
$V\to U$ は同型である。像が次元2以下、または三次元で $s$ がその像と $B$ の交叉にあるとき、
その $s$ で二形式の包が次元3に達しないからである。
従って全ての $s\in B, u\in U$ の縮約が $K$ に与える空間は二文字の全交代空間となる。
実際 $B\wedge U=\bigwedge^2U$ である。隣接対も対称になり、既知の反対対称性と合わせて
三位置対称、従って $\rho=6$ となる矛盾を得る。これで $T$ は全空間でもなく平面を含まない。

二次式の基点なので残る曲面成分は既約二次曲面のみである。その成分を $D$ とし、像が曲面だと仮定する。
滑らかな場合には $s^*\mathcal O(1)=\mathcal O(m,n)$ として同じ行列式を両母線で読み $m,n\le1$。
曲面像には $(m,n)=(1,1)$ が必要である。錐の場合も超平面束となる。
実際、頂点を解消した $\mathbb F_2$ の負切断 $C_0$ とファイバー $f$ に関し、
引戻しの因子は $C_0$ 上の次数0から $m(C_0+2f)$ であり、母線上で $m\le1$ である。
よってどちらも根基射は周囲の線型写像の制限で、曲面像ならその階数は4である。
階数3では基点を避けた点から平面への二次射影で、標数零の一般点は二つの逆像を持つ。
それらを含む核の射影直線も $T$ に入るが、一般ファイバーは零次元である
（$T$ は高々二次元で、曲面像を持つ）。これが矛盾である。
この線型同定で $s$ を変数にすると $\iota_s t(s)$ は $D$ の方程式の定数ベクトル倍である。
さらに $s$ で縮約すると恒等零なのでこの定数ベクトルも零である。すると $T$ が全空間となってしまう。
従って退化集合は高々一次元である。
\end{proof}

\begin{proof}[特異商の次数]
$a$ は層として単射であり、$\mathcal E$ は局所的に自由分解長1以下である。
余次元1で局所自由、余次元2以上では深さが少なくとも1なので捩れがない。
行列式は自明なため、$F\simeq\mathcal E^*$ である（局所自由な開集合での外積対合を反射延長する）。
$Q$ の台は高々一次元。層の双対複体から $F$ は
$J=\operatorname{coker}(b^\vee)$ の部分層、商は $4\mathcal O(1)$ の部分層である。
$0\to4\mathcal O(-1)\to10\mathcal O\to J\to0$ により、
$H^0(F(-1))=0$、$H^1(F(t))=0$ ($t\le-2$) を得る。
$H^2(\mathcal E(-2))=0$ は元の二つの短完全列から従うので $H^2(F(-2))=0$ でもある。
またSerre双対により $H^3(F(-2))=0$。Riemann--Rochより $c_3(F)=2\chi(F(-2))=0$。
これは反射階数2層の点特異性がないことを意味する。実際その局所射影次元は高々1で、
双対化した自由分解のChern指標は $[F^*]-[\mathcal Ext^1(F,\mathcal O)]$ のものである。
$F^*\simeq F$ と第三指標の符号を比較すると点支持の後者の長さは $c_3(F)=0$ となる。
局所極小表示により $F$ は局所自由である。
$Q\ne0$ なら閉点で $\mathcal E$ の深さは2以上、$F$ は自由なので $Q$ の深さは1以上、
従って純一次元である。$\mathcal E(-2),F(-2)$ の全コホモロジーは零なので $Q(-2)$ も同じである。
$d=c_2(\mathcal E)-c_2(F)$ である。$F^*=F$ と既に示した消滅から
$h^1(F(-1))=c_2(F)$ なので $d\le4$。等号なら
$H^1(F(-1))=H^2(F(-2))=H^3(F(-3))=0$ より $F$ は0正則で大域生成される。
自明な行列式と階数2より $F$ は自明である。
\end{proof}

% !TeX root = ../papers/nc-p3.tex
\section{片側対称な階数十での全ての非局所自由境界}
\label{sec:finite-rank-ten-instanton-boundary}
\begin{theorem}\label{thm:finite-rank-ten-instanton-boundary}
有限点グラフを持つ片側対称な標準有限Frobenius窓で反対階数が10なら、
定理\ref{thm:finite-instanton-sheaf-boundary}の層が局所自由でない全ての窓は
二元三次式・二次式の乗法への辺商を持ち、非正則である。
従ってこの枝の全正則窓の添字保持同型類は、電荷4のインスタントン\emph{束}の
射影同値類に単射する。反対階数7から9の非局所自由窓には既存の全次数判定が適用される。
\end{theorem}
\begin{proof}[特異商上の二つの可換な一次鉛筆]
前定理の非零商 $Q$ を取り、一般の二つの一次式 $x,y$ による台の有限射影を
$f:\operatorname{Supp}_{\rm sch}Q\to\mathbb P^1$ とする。
$G=Q(-1)$ は純なので $f_*G$ は階数 $d$ のベクトル束であり、
$H^0(G(-1))=H^1(G(-1))=0$ により $f_*G=\mathcal O^d$ である。
残り二つの座標はこの自明束からその一次捩りへの可換な二行列
$A=xA_0+yA_1,B=xB_0+yB_1$ で作用する。
$H^0(G)\hookrightarrow H^1(\mathcal E(-1))=N_0$ は $H^0(F(-1))=0$ により単射で、
この埋込みは $H^0(G(1))\to N_1$ と乗法で両立する。
従って四行列の共通不変部分空間 $W\subset k^d$ への制限で成り立つ一次関係は、
全て $K\cap(U^*\otimes W)$ に入る。

次元が最小の非零共通不変部分空間を $W$、その次元を $n$ とする。
$n=1$ なら一つの $v$ に対し四座標の積が二次元以下に入るので、
二独立な一次式が $v$ を消し有限点グラフに反する。
$n\ge2$ の $W$ を同一の非零一次式が消すことも不可能である。いずれも $K$ に純関係でなく
点グラフの直線を与えるからである。
$n=2$ なら二行列の跡零部分は関数体上で比例する（非スカラー二次行列の中心化環は二次元）。
一方が零なら $W$ は一つの固定一次式で消える。ともに非零なら、成分の共通因子を除いた
原始行列が一次のときは比例係数が定数なので同じ矛盾になる。
原始行列が定数ならその固有直線が全係数の不変直線である。従って $n=3,4$ だけが残る。

$W$ が定める部分層のスキーム論的な台を $Y$ とする。
$f_*\mathcal O_Y$ は $\mathcal End(W\otimes\mathcal O)$ の部分代数なので、
直線束分解の次数は全て非正である。また大域切断は $A_0,A_1,B_0,B_1$ と可換な
定数行列であり、$W$ の既約性によりスカラーに限る。その階数 $e$ は $n$ 以下である。
実際二つの可換 $n$ 次行列の生成代数は次元 $n$ 以下である。
関数体の代数閉包に移り、$A$ のJordanブロックごとに異なる固有値を選んで、
$A$ と可換で巡回的な行列 $C$ を作る。一般の $B+tC$ は巡回的で $A$ と可換なので、
二行列の語はその $n$ 次元の中心化環に入る。$t=0$ で語の張る次元は増えない。
よって $f_*\mathcal O_Y=\mathcal O\oplus\bigoplus_{j=2}^e\mathcal O(-a_j)$、$a_j\ge1$ である。
\end{proof}

\begin{proof}[三次元部分と四次元部分]
$n=3$ なら四つの独立な一次座標の存在より $4\le h^0(\mathcal O_Y(1))\le2+(e-1)\le4$。
従って $e=3$、二つの $a_j$ は1である。
$Y$ は滑らかなねじれ三次曲線であることを示す。まず二つの非零一次式の積が
$Y$ を消すことはない。そうなら関数体上の二行列の積が零で、どちらも非零なので
その一つは階数1である。一次行列が関数体上で階数1なら、原始列・行に分けると
次数の和が高々1なので、像直線か核平面が定数である。
他の座標行列と可換だからこの定数部分空間は全係数で不変となり矛盾する。

$x,y$ と他の二座標は $f_*\mathcal O_Y(1)$ の全切断を張るので、
斉次座標環は $k[x,y]$ 上で次数 $0,1,1$ の自由基底を持つ。
実際他の二座標が二つの次数1の基底を与え、以後も二変数を掛けて全切断を張る。
そのイデアルは三二次式で生成される（他の二座標の三つの積を自由基底に戻す）。
高さ2のCohen--Macaulay環なのでHilbert--Burch分解とHilbert分子
$1-3z^2+2z^3$ により、この三式は二行三列の一次行列の小行列式である。
いずれかの非零の行・列結合の交点成分が零なら、その成分を左上に持ってきて非零小行列式に
二一次式の積を得る。三小行列式は独立なのでこれは先の禁止に反する。
よって全ての行結合は三独立な一次式を与え、その零点は一点である。
行方向の射影直線との入射は従って射影直線そのもので、次数3の小行列座標で写る。
三形式の共通の点は全て一意な行方向を持つ。行列全体がある点で零なら
全ての行方向の唯一の零点がそこに集中して曲線を覆えないからである。
単元成分を使って方向を局所的に解けるので入射から $Y$ への射はスキーム同型である。

$W$ に対応する $G$ の部分層はこの滑らかな三次上の階数1の捩れのない層で、
$\chi=3$ なので $\mathcal O_{\mathbb P^1}(2)$ である。
その三つの切断と四座標との積は $P_2\otimes P_3\to P_5$ の通常の乗法に等しい。
よって元の辺をこの三文字へ圧縮すると $R$ の像はこの乗法の双対像に入り、
定理\ref{thm:binary-linear-complex}の非正則な二元商となる。
ここで $G(1)$ から $N_1$ への写像が単射である必要はない。

$n=4$ なら $d=4$ かつ $W=k^4$ で、$N_0=H^0(G)$ である。
したがって二回乗法の作用は $H^0(\mathcal O_Y(2))$ を経由し、
その次元は高々 $3+2(e-1)\le9$ である。これは $\operatorname{Sym}^2 U^*\to
\operatorname{Hom}(N_0,N_2)$ の階数、すなわち反対階数が10であることに反する。
これで全ての非局所自由の場合が尽きる。
\end{proof}

% !TeX root = ../papers/nc-p3.tex
\section{四次Grassmann写像と空間楕円八次曲線による復元}
\label{sec:finite-quartic-grassmann-windows}
\begin{theorem}\label{thm:finite-quartic-grassmann-windows}
反対階数10の有限片側対称窓で、二つの奇位置から得る層を $E,E'$ とする。
両方が局所自由なら $E(2)$ はちょうど四つの独立な大域切断で生成され、
その完全評価列は
\[
 0\to E'^*(-2)\to H\otimes\mathcal O\to E(2)\to0,\qquad \dim H=4
\]
である。従って全ての正則な対象は、二平面核の写像
$f:\mathbb P(U)\to\mathrm{Gr}(2,H)$、$f^*\mathcal O(1)=\mathcal O(4)$ から復元できる。
二窓の添字保持同型はこの写像の始域の射影変更と $H$ の線型変更による同値と必要十分である。
二移動は $f(x)\mapsto f(x)^\perp\subset H^*$ に対応し、正則窓の最小周期は
この双対写像と同値なら2、そうでなければ4である（対称軸の位相を固定する）。
\end{theorem}
\begin{proof}
定理\ref{thm:finite-symmetric-instanton-monad}の極小分解を層化して1捩りする。
最後の二項の余核は、反対の三段加群が次数を反転した双対であることより
$(\ker b')^*(-4)$ と同定できる。中央項への包含から $4\mathcal O(-3)$ への射影は、
反対の線型部分 $a'$ を双対化し $-4$ 捩りした写像である。$a'$ が部分束なので全射で、核は $E'^*(-4)$ である。
中央項の残り $4\mathcal O(-2)\oplus4\mathcal O(-1)$ から $\ker b$ への写像は従って全射で、
核が今の $E'^*(-4)$ となる。さらに $4\mathcal O(-1)$ の写像は $a$ であり、
これを割って2捩りすると表示の列を得る。$H^0(E'^*(-2))=H^1(E'^*(-2))=0$、
かつ行列式は自明なので、これは大域切断全体の評価である。核の双対の評価は他方の完全評価列となる。
$f$ から普遍商の引戻しと2捩り戻しで $E$ が得られ、元の窓を復元する。
$\det E(2)=\mathcal O(4)$ は豊富なので $f$ はその像への有限射である。
二移動は二層を交換する。正則な両対称窓は定理\ref{thm:all-bisymmetric-as-closure}で
反対階数8以下なので、奇数周期で他方へ対称軸が移ることはない。
\end{proof}

\begin{theorem}\label{thm:finite-elliptic-octic-windows}
同じ両局所自由部分は、滑らかな楕円曲線 $C$、次数8の線束 $L$、
$\mathbb P^3$ への閉埋込みを与える $D\subset H^0(L)$（次元4）からの次の開条件付き構成で尽くされる。
\[
 N_j= H^0(L^{j+1})/D^{j+1}\quad (0\le j\le2)
\]
を $D$ の乗法で結び、次元 $4,6,4$、二回乗法に対する有限Frobenius検査・有限辺・反対階数10、
両方のモナドの部分束条件を課す。さらに四次の評価の全射性と、次数1の関係による表示が $N_2$ を与える条件を課す。ここで $D^m$ は積の像である。
同じ窓の楕円八次曲線は射影変更を除いて $E(2)$ の滑らかな零曲線を持つ切断の
$\mathbb P^3$ の開部分で尽きる（束を保つ射影座標変更も作用する）。
この有限候補部分の添字保持同型類は既約で次元14であり、その幾何学的一般点は全次数で正則である。
特にこの位相の反対階数10の有限片側対称\emph{正則窓全体の閉包}は既約である。
\end{theorem}
\begin{proof}
大域生成された $E(2)$ の一般切断の零点 $C$ は滑らかで余次元2である。
実際 $c_2(E(2))=8$ で空ではなく、普遍零点は評価核の束なので標数零の一般滑らか性を使える。
随伴公式は $\omega_C=\mathcal O_C$ を与える。また
\[
 0\to\mathcal O\to E(2)\to I_C(4)\to0
\]
を $-4$ 捩りした列から $H^1(I_C)=H^1(E(-2))=0$、従って $C$ は連結な楕円八次曲線である。
同じ列と直線束のコホモロジーから全次数で $H^1(E(j-1))=H^1(I_C(j+1))$ を得る。
後者は記載の切断の商そのもので、乗法と両立する。

逆に曲線データで三つの次元条件を満たせば $H^0(I_C(1)),H^0(I_C(2)),H^0(I_C(3))$ は零である。
さらに曲線のSerre構成で元の切断を復元できることを直接確認する。
$\operatorname{Ext}^1(I_C(4),\mathcal O)\simeq
H^0(\mathcal Ext^2(\mathcal O_C(4),\mathcal O))=H^0(\omega_C)=k$
である。実際 $H^1,H^2$ の中間直線束コホモロジーは零で、局所完全交叉の他の局所Extは零である。
非零類は各点で二方程式のKoszul表示の単元類なので、拡大の中項は局所自由である。
その $-2$ 捩りを $E_C$ とすれば、行列式は自明で $H^0(E_C)=0$ なので安定である。
$H^1(E_C(-2))=H^1(I_C)=0$。ただし三つの隣合う次数だけから全Rao加群との一致を
仮定しないため、逆構成の開条件には $H^1(I_C(4))=0$ と
$D\otimes N_0\to N_1$ の核による一次表示が $N_2$ を与えることも含める。
順方向ではこれらは既に成立する。
$I_C$ は5正則である（$H^2(I_C(3))=H^1(L^3)=0, H^3(I_C(2))=0$）。
これで高次のRaoは零、負次数と次数0も連結性と負次数切断の不在で零となる。
したがって $H^1_*(E_C)$ は復元した三段加群そのものである。

この加群の線型モナドが与える束 $E$ と $E_C$ の一致は次の初等的な構成でも分かる。
$H^1(E_C(-1))$ の基底で普遍拡大 $0\to E_C\to J\to4\mathcal O(1)\to0$ を取ると、
全次数で乗法がこの四元で生成するので $H^1_*(J)=0$。
双対側でも $E_C^*=E_C$ であるから、$H^1_*(J^*)=H^1_*(E_C)$ を消す四元の
普遍拡大を双対化して $0\to4\mathcal O(-1)\to B\to J\to0$ を得る。
$H^1_*(B)=H^2_*(B)=0$ なのでHorrocksの分裂判定で $B$ は直線束の直和である。
$B\to4\mathcal O(1)$ の切断上の像は三段加群の関係を生成する。
その極小生成は次数1に十本で、$B$ の階数も10なので、十本全てが $\mathcal O$ から来る。
複体の左端も四つの線型シジジーを張るので $E_C=E$ である。
ここで $J^*\to E_C^*$ の第一コホモロジーの同型には $4\mathcal O(-1)$ の中間消滅を用いた。
逆に同じ窓ではこの束が決まり、拡大の一意性により曲線の選択は表示の切断だけである。

基礎の楕円曲線と次数8の線束の同型類は一次元既約である（平行移動で任意の二つの次数8の束が移る）。
その上の $\mathrm{Gr}(4,8)$ の開部分なので、曲線表示の同型類は既約で次元17である。
原点を選んで $L=\mathcal O(8o)$ としても残余自己同型は有限で、同じ次元となる。
開条件は空でない。定理\ref{thm:twenty-point-distinct-companion}の例の両辺は
$B_-=-B_+(J\otimes J)$ で移り、前節のモナドの部分束検算が両側で通る。
その窓の二十点は被約で、その位置を全て固定する射影変換は自明なので一般安定化群は有限である。
よって一般ファイバーの次元3を引いて14となる。
標準複体の各固定内部次数の完全性は開であり、この正則例で全て成立するため
既約なパラメータの幾何学的一般点では各次数で完全である。個別の閉点での正則性はここでは従わない。
\end{proof}

% !TeX root = ../papers/nc-p3.tex
\section{四次K3曲面上の切断と二つの楕円鉛筆}
\label{sec:finite-quartic-k3-cut}
\begin{theorem}\label{thm:finite-quartic-k3-cut}
定理\ref{thm:finite-quartic-grassmann-windows}の両局所自由部分で、
一般の二平面 $P\subset H$ に対し $T=\{x:f(x)\cap P\ne0\}$ は滑らかな四次K3曲面である。
$h$ を超平面類とする。$T$ 上には二つの基点のない楕円鉛筆 $|C|,|F|$ があり、
それぞれの一般ファイバーから元の窓とその二移動を復元できる。
\[
 C+F=4h,\quad hC=hF=8,\quad C^2=F^2=0,\quad CF=32 .
\]
$D=2h-C=F-2h$ と置く。導来圏 $D^b(T)$ の四対象
\[
 G=(\mathcal O,\mathcal O(h),\mathcal O(D)[1],\mathcal O(h+D)[1])
\]
の次数零自己準同型代数は、積の向きを経路順に読んだとき元の小さい四層切断 $A$ である。
次数つき自己準同型代数全体は二次の双対を加えた $A\ltimes DA[-2]$ で、
各双対辺に補経路長を与えれば元の有限Frobenius代数を得る。
\end{theorem}
\begin{proof}
一般切断 $s\in H$ は滑らかな楕円八次曲線を零点に持つ。
$P=\langle s,t\rangle$ とすれば $T$ の方程式は $s\wedge t$ である。
生成性により $s$ の零曲線の外でこの線型系は自由であり、そこで一般の成員は滑らかである。
零曲線上で $t$ も零となる条件の余次元は固定点で2なので、一般の $t$ はその全点で非零である。
$s$ の微分は零曲線上で束へ全射だから $s\wedge t$ の微分もそこで消えない。
こうして得た $T$ の全点で $f(x)\cap P$ は直線であり、これと $f(x)$ の $H/P$ における像直線が
二つの射 $T\to\mathbb P^1$ を与える。前者の繊維は $P$ の指定切断の零曲線である。
後者も $P^\perp$ の指定切断の零曲線であり、双対評価列に同じ議論を適用できる。
一般の二平面を取っているので両側とも一般繊維は滑らかで連結である。
核二平面束の行列式を二直線から計算すれば $C+F=4h$。二鉛筆の自己交叉と曲線の次数から残りの数値が従う。

$k=1,2,3$ に対し $H^1(T,kh-C)=H^1(\mathbb P^3,I_C(k))$ が乗法と両立する。
これは超曲面のイデアル列と射影三空間の直線束の中間消滅による。
また $H^0(T,kh-C)=0$ は $C$ が三次以下の曲面に含まれないことによる。
$C-kh=(4-k)h-F$ に同じ事実を用いれば $H^2(T,kh-C)=0$ である。
従って二つの隣接する無移動直線束の間の順方向射は $H^0(h)$、
シフトをまたぐ順方向射は上の三つのRao空間で、全て導来次数零にある。
合成は二端での $H^0(h)$ の作用、すなわち三段加群の可換乗法である。
一次関係と二回乗法の復元によりこれが $A$ を与える。
逆方向はK3のSerre双対で対応する順方向の双対が次数2にあるだけである。
対角には恒等射とその次数2の双対がある。trace対合の合成との両立性と次数4の消滅から
全積が上の自明拡大となる。この議論に高次積の消滅は必要ない。
\end{proof}

\begin{proposition}\label{prop:quartic-k3-two-pencil-twists}
$T_M$ を球面対象 $M$ に沿うtwist、$L_a$ を $\mathcal O_T(a)$ とのテンソル積とする。
積の右側が先に作用する。$A_0=T_{\mathcal O}T_{\mathcal O(h)}$ とすると
\[
 \Phi=T_{G_0}T_{G_1}T_{G_2}T_{G_3},\qquad
 \Phi[-2]= A_0 L_{-C} A_0^{-1} L_{-F}.
\]
数値的には $v=(0,F,-4)$ を $\Phi$ が固定し、$A_0^{-1}$ はこれを $(0,-C,-12)$ に送る。
\end{proposition}
\begin{proof}
$B_0=T_{\mathcal O(-2h)}T_{\mathcal O(-h)}$ とする。射影三空間の連続した四直線束を四次K3へ制限すると
その順序のtwistの積は $L_{-4h}[2]$ となることを確認する。
各twistの核は $M^*\boxtimes M\to\mathcal O_\Delta$ の錐である。
積を展開すると、恒等項の手前は四直線束の真に増加する列から作る評価複体である。
長さ1から3の差の順方向の導来Homは、射影三空間から制限して同型である
（超曲面列の負の捩りは $-3,-2,-1$）。よって評価複体は射影三空間での同じ複体の導来引戻しである。
射影三空間ではこの複体が対角層を与える。実際、四直線束は充満例外列なので
順方向の射を用いた正規化バー分解の評価が恒等関手の核を解消する
（同じことは右端から各例外成分を順に除く射影の積が零になることからも分かる）。
従って曲面での積は引戻した対角層から曲面自身の対角層への増大の錐である。
四次超曲面の一項のKoszul分解から、この増大は次数零で同型、
唯一の負次数コホモロジーは次数 $-1$ の対角上の $\mathcal O_T(-4h)$ である。
よって錐は $L_{-4h}[2]$ を与え、特に $B_0 A_0=L_{-4h}[2]$。
twistは球面対象のシフトに依存せず、テンソル積による共役と両立する。
$D=F-2h$ より $\Phi=A_0 L_F B_0 L_{-F}$、従って $F-4h=-C$ を代入して表示式を得る。
Mukai対合を $(r,c,s)\cdot(r',c',s')=cc'-rs'-r's$ とすれば
twistとその逆はともに $x\mapsto x+(x\cdot v(M))v(M)$ と作用する。
$v$ にまず $\mathcal O$、次に $\mathcal O(h)$ の反射を施せば $(4,F,0)$ を経て $(0,-C,-12)$ となる。
$F^2=C^2=0$ と分解式で固定性が従う。
\end{proof}

% !TeX root = ../papers/nc-p3.tex
\section{二つの球面twistによる楕円繊維間の移送}
\label{sec:finite-elliptic-fiber-transfer}
\begin{lemma}\label{lem:octic-linear-determinant}
前節のいずれかの滑らかな楕円八次繊維を $F$、$L=\mathcal O_F(h)$、
$D_0=H^0(\mathbb P^3,\mathcal O(1))\subset H^0(L)$ とする。
全ての次数4の線束 $J$ で $D_0\otimes H^0(J)\to H^0(LJ)$ は全射、核 $K_J$ は次元4である。
一般の $x\in T$ と一般の $J$ で、この核を $x$ で評価する写像
$K_J\to H^0(J)\otimes\mathcal O(h)|_x$ は同型となる。
\end{lemma}
\begin{proof}
$I_F(4)$ は大域生成されるので $F$ と長さ5以上で交わる直線はない。
実際その直線を全ての四次切断が含めば、基底軌跡に含まれてしまう。
$M=\ker(D_0\otimes\mathcal O_F\to L)$ とする。
次数 $d\le3$ の線束 $J_d$ による消える積
$t\in D_0\otimes H^0(J_d)$ の非零元はテンソル階数3以上である。
階数1は整性に反する。階数2なら係数の二切断の固定因子を除いた鉛筆は次数 $e\le d$ であり、
対応する二つの超平面切断はこの鉛筆に共通の切断を掛けたものになる。
従って超平面の交線は $F$ と長さ $8-e\ge5$ で交わる。
特に $M^*$ は次数2以下の商直線束を持たない。実際その商への
$D_0^*\otimes\mathcal O_F$ の写像は上の形の関係であり、次数2以下では切断が高々二つである。

$M^*$ には次数4以上の部分直線束もない。飽和してそのような束を割ると階数2・次数4以下の商 $Q$ が残る。
任意のそのような $Q$ は次数2以下の商直線束を持つ。さもなくば全ての次数2の $B$ で
$H^0(Q^* B)=0$ であるが、点 $p$ を一つ選べばRiemann--Rochにより
$h^0(Q^*B(p))\ge2$、これを各点 $q$ で評価した核は $H^0(Q^*B(p-q))=0$。
従って切断の評価は全点で同型となるが行列式の次数が正で矛盾する。
さらに次数3の部分直線束は高々一つである。二本が異なる像を持てばそれらの和の飽和を割って
次数2以下の商を得るからである。同じ像に対する写像は定数倍しかない。
従ってSerre双対より全ての次数4で $H^1(MJ)=0$、これが全射性と核次元を示す。
次数3については $h^0(MJ_3)=1$ で、ただ一つあり得る例外でも高々2である。

$x\notin F$ で $D_x\subset D_0$ を $x$ で消える三切断とし、
$M_x=\ker(D_x\otimes\mathcal O_F\to L)$ とする。この評価は基点を持たない。
$M_x^*$ が半安定でなければ次数3以下の商直線束、従って $D_x$ に支えられる上記の関係がある。
その階数は3、従って次数3で、関係から $D_x$ は一意に決まる。
次数3のPicard曲線の各点で関係の射影空間の次元は零、一つの例外でも高々1であるから、
このような $x$ の族は次元1以下である。これは点の個数からの推測ではなく、
Picard曲線上で切断の積の核を射影化した有限型の入射対応の次元評価である。
従って一般の $x\in T$ で $M_x^*$ は半安定である。
階数2・傾き4の半安定束に入る次数4の線束の同型類は高々二つである。
一本を入れると商も同次数の線束で、他のものはこの部分または商に同型だからである。
よって一般の $J$ に対して $H^0(M_x J)=0$、すなわち $K_J$ の $x$ 評価が単射である。
\end{proof}

\begin{theorem}\label{thm:octic-fiber-twist-transfer}
滑らかな $F\in|F|$ に次数 $-4$ の一般の線束 $J^{-1}$ を置き、その直像を $X$ とする。
$A_0=T_{\mathcal O}T_{\mathcal O(h)}$ に対し、$A_0^{-1}X=Y[-1]$ である。
$Y$ は $|C|$ の一繊維にスキームとして支えられる純一次元層で、Mukaiベクトルは $(0,C,12)$ である。
従って $\Phi[-2](X)\simeq X$ である。$|C|$ の全繊維が整なら $Y$ はその繊維上で階数1の無捩れ層である。
\end{theorem}
\begin{proof}
$i:F\hookrightarrow T$ とする。$J$ は四切断で生成され、評価の核
$W=\ker(H^0(J)\otimes\mathcal O_T\to i_*J)$ は局所自由である。
実際滑らかなCartier因子上の線束は局所射影次元1である。
$\mathcal O_F(F)$ は自明なので、双対列は
$0\to H^0(J)^*\otimes\mathcal O_T\to W^*\to X\to0$ である。
これは $T_{\mathcal O}^{-1}X$ を与える。逆twistの評価は
$\operatorname{Ext}^1_T(X,\mathcal O_T)=H^0(J)$ の恒等テンソルによる普遍拡大であり、
上の双対列は評価そのものを双対化しているからである。
次に $H^*(T,W(h))$ は前補題の積の核 $K_J$ が次数零にあるだけである。
従って $T_{\mathcal O(h)}^{-1}W^*$ は $W^*\to K_J^*\otimes\mathcal O(h)$ の錐の $-1$ 移動である。
$F$ の外ではこれは前補題の四行四列評価の転置で、一般の $J$ で一般点同型である。
束間の単射なので余核 $Y$ は純一次元である（局所射影次元1）。
行列式は $4h-F=C$ の非零切断であり、その随伴行列によって行列式の方程式が余核を消す。
鉛筆は連結繊維なのでこの零因子は実際に一つの繊維で、その上で $\mathcal O_T(C)$ は自明である。
ベクトルは命題\ref{prop:quartic-k3-two-pencil-twists}の計算で決まる。
整な場合の階数は一次の重複度が1であることから従う。
$X$ 上での $L_{-F}$ の自明性と $Y$ 上での $L_{-C}$ の自明性を分解式へ代入すれば最後の同型を得る。
\end{proof}

% !TeX root = ../papers/nc-p3.tex
\section{二つの相対Picard曲面上での四移動}
\label{sec:finite-residual-picard-surfaces}
\begin{theorem}\label{thm:residual-picard-surface-transport}
定理\ref{thm:finite-quartic-k3-cut}の曲面で $\mathrm{NS}(T)=\mathbb Zh\oplus\mathbb ZF$ とする。
$M_F,M_C$ を二つの鉛筆の、それぞれ次数 $-4,12$ の相対コンパクト化Picard曲面とする。
これらは同型なK3曲面である。共通の曲面 $M$ 上に二つの繊維類 $a,b$ があり
$a^2=b^2=0,\ ab=2$。$a+b$ は豊富かつ大域生成である。
ねじれた普遍層による同値 $U_F:D^b(M,\alpha)\to D^b(T)$ の下で
$\Phi^{-1}[2]$ は $\mathcal O_M(a+b)$ とのテンソル積に対応する。
\end{theorem}
\begin{proof}
$(uh+vF)^2=4u(u+4v)$ なので $T$ に自己交叉 $-2$ の曲線はない。
楕円鉛筆の繊維の各成分は繊維と直交し、Hodge指数定理と随伴により繊維の有理数倍である。
一般繊維が連結で滑らかな完全鉛筆の繊維類は原始的である
（非原始なら原始的なnef等方因子はRiemann--Rochから鉛筆を持ち、倍数系の次元は少なくともその倍数）。
両方の全繊維は従って整である。

ここで用いる相対Picard構成を説明する。整な算術種数一のGorenstein繊維について、次数 $d$ の
階数一無捩れ層は全て安定であり、そのコンパクト化Picardは射影的である
（例えば偏極に関する安定層の射影モジュライのこの繊維部分を取れる）。
基底のエタール開集合上で滑らかな切断 $s$ を取れば、点 $x$ に
$I_{x/\text{繊維}}((d+1)s)$ を対応させる同型で元の曲面に局所同定される。
実際、捩り戻した次数 $-1$ の層は、双対の唯一の切断で繊維の構造層に入り、余核の長さが1である。
この構成は相対双対と切断の基底変換により族で逆となる
（Gorenstein曲線の無捩れ層は反射的であり、次数1の双対の第一コホモロジーは消える）。
よって全空間は滑らかである。座標変更は滑らかな一般繊維の平行移動なので
$R^1\pi_*\mathcal O$ と相対双対線束の基底上の遷移は元の鉛筆と同じである。
いずれも $R^1\pi_*\mathcal O=\mathcal O_{\mathbb P^1}(-2)$、
$\omega_{M/\mathbb P^1}=\pi^*\mathcal O(2)$ となり、全空間はK3である。

普遍層は局所普遍層の基底からの線束のずれをBrauer類に入れて取る。
対応するねじれた導来圏からの変換 $U$ が同値であることは局所でも確認できる。
上の座標では、テンソルとシフトを除き、核は相対対角のイデアルである。
各繊維上で、その核による変換の1移動は
$\mathrm{Cone}(\mathbf R\Gamma(-)\otimes\mathcal O\to\mathrm{Id})$ である。
整な種数一Gorenstein曲線の構造層は完全な1球面対象で、逆は
$\mathrm{Cone}(\mathrm{Id}\to\mathbf R\mathrm{Hom}(-,\mathcal O)^*\otimes\mathcal O)[-1]$ となる。
これは特異繊維でも $H^*(\mathcal O)=k\oplus k[-1]$ と双対化層 $\mathcal O$ のtrace対合により、
二つの評価錐を合成して恒等因子を相殺する球面twistの証明がそのまま適用できる。
相対的にも評価には $\pi^*\mathbf R\pi_*$ を用い、その随伴と相対双対を用いるとこの逆評価が得られる。
平坦な曲線族の完全複体に対する導来基底変換により随伴の単位・余単位は各繊維で同型であり、
導来中山により全空間でも同型である（全空間は正則、相対対角イデアルは両因子に平坦）。
この局所検査は普遍層の線束のずれで変わらない。
また繊維積上の核なので基底の $\mathcal O(1)$ とのテンソル積を保つ。

$V=U_C^{-1}A_0^{-1}U_F$ を取る。定理\ref{thm:octic-fiber-twist-transfer}より
一般点の点層を他方の点層の $-1$ 移動に送る。実際整繊維上で得た階数一無捩れ層はちょうど上の対象である。
この対応は稠密開集合の間の同型である。核の一般平坦性から点層を与える開集合で射となり、
同値ゆえ単射、標数零の同次元曲面なので双有理である。

全点へ延ばすため、両Picard曲面上で、正次元の階数を持たない非零ねじれ層 $Q$ について
$\dim\operatorname{Ext}^1(Q,Q)\ge2$ を示す。$U$ によるMukaiベクトル $x=(r,uh+vF,s)$ の像の記法を
$T$ 側で取ると、$Q$ は一般点層と直交するので $x$ は $(0,F,-4)$ の直交格子に入る。
従って $8u+4r=0$、$x^2=4u(u+4v+s)$ は4の倍数である。
他方 $Q$ の自己Euler標数の負はその一次元サイクルの自乗である。
この値が負なら構成成分の少なくとも一本が負の自乗を持ち、K3の随伴でその曲線は滑らかな有理 $(-2)$ 曲線である。
しかしその曲線上でBrauer類は自明で、その直線束の直像は平方 $-2$ の類を与え、4の倍数性に矛盾する。
ゆえに $x^2\ge0$、Serre双対から所望の不等式が従う。
$C$ 側では $F$ を $C$ に替えて格子が同形であり、直交条件 $8u-12r=0$ は $u=3t,r=2t$ を与え、
再び平方は4の倍数である。この議論は逆Brauer類にも双対化で適用できる。

任意の点の $V$ 像は、対応の稠密開集合のうちその点以外の全点と直交するので、各コホモロジー層は真の閉部分に支えられる。
曲面上の自己Extのスペクトル系列では
$\bigoplus_i\operatorname{Ext}^1(H^i,H^i)$ に入る・出る高次微分はなく、
この次元の和は対象の自己 $\operatorname{Ext}^1$ の次元以下である。
点層でこの上限は2なので、上の不等式からコホモロジー層は一つである。
点の数値類は全点で一定で、一般点と同じ符号付き長さ1の点類だから、それは点層の移動である。
平坦性の開性により移動量は連結な曲面上で一定であり、逆同値にも同じ議論を適用できる。
従って $V$ は曲面の同型のグラフ上のねじれ線束と移動による変換である。
これは各繊維が長さ1の平坦族の台が射のグラフであることからも直接分かる。
従って共役によって相手の鉛筆のテンソル積は引き戻した鉛筆のテンソル積に移る。

二つの繊維類の交叉も符号に依存せず計算できる。$u_0=(0,0,1)$ とすると
$A_0 u_0=(2,-h,1)$。$U_F^{-1}$ は $T$ の一般点層を $M_F$ の滑らかな繊維上の線束の移動へ送る。
これは相対核の随伴を一般点で制限すると通常の楕円曲線上の普遍線束の制限になるためである。
同様に $A_0u_0$ は、今得た同型で $M_C$ の繊維上の線束の移動に対応する。
階数零同士のMukai対合の絶対値が二繊維の交叉数なので $ab=|u_0\cdot A_0u_0|=2$ である。
$a+b$ は基点を持たず正の自乗で、零交叉の曲線があればHodge指数により $(-2)$ 曲線が生じるので豊富である。
最後に命題\ref{prop:quartic-k3-two-pencil-twists}の逆を取り、二つのテンソルに上の同値を施せば結論を得る。
\end{proof}

% !TeX root = ../papers/nc-p3.tex
\section{最小Picard格子の四次切断による全次数正則性}
\label{sec:finite-minimal-picard-regularity}
\begin{theorem}\label{thm:finite-minimal-picard-regularity}
反対階数10の両局所自由な有限片側対称窓が
$\mathrm{NS}(T)=\mathbb Zh\oplus\mathbb ZF$ の滑らかなSchubert切断を一つでも持てば、
その窓は全次数でKoszul左右AS正則である。これは個別の窓に適用される十分条件である。
\end{theorem}
\begin{proof}
小切断の四つの射影を順方向の射が添字を増す順に $e_0,\ldots,e_3$ と書く。
これらを $G_i$ に送る完全関手 $J$ を取れる。実際、$D^b(T)$ のdg増強の中で
対角に恒等だけを残し順方向の射複体を全て取った有向圏は、非正次数の良い切断を経て
次数零の小切断と擬同値である。これを完全包へ延ばす。右随伴を $R$、小切断のSerre関手を $S$ とすると
\[
 \mathrm{Cone}(\mathrm{Id}\to RJ)=S[-2],\qquad
 \mathrm{Cone}(JR\to\mathrm{Id})=\Phi,\qquad
 \Phi J=JS[-1].
\]
第一式は単位の錐の両側加群が次数2にのみコホモロジーを持ち、そこでは
定理\ref{thm:finite-quartic-k3-cut}のSerre対合で双対正則両側加群になるためである。
第二式では正規化バー評価の錐を展開すると、四対象の真に増加する列を通る項が一つずつ現れ、
四twistの積の評価錐と同じ微分になる。第三式は $J\to JRJ\to J$ が恒等であり、
その分裂全射の錐が分裂単射の錐の1移動であることから従う。

有向列に左変異を施して得る単純を $s_i$ とし、$p_i=s_i[-i]$ と置く。
実際 $e_i$ からそれ以前の対象を右端から順に左変異で除けば、各 $e_j$ からの導来Homは
$j=i$ の次数零でのみ一次元となり、これが単純を特徴付ける。
定理\ref{thm:cut-serre}の振り方と添字を逆にしたので、これらの直和が大切断のtilting対象であり、
$\mathrm{Hom}(p_i,p_j)$ は $i\ge j$ の向きにある。
$P_i=Jp_i$ は順に次の球面対象である。$M_h=\ker(H^0(h)\otimes\mathcal O_T\to\mathcal O_T(h))$ とすると
\[
 P_0=\mathcal O_T,\quad P_1=M_h,\quad
 P_2=A_0(\mathcal O_T(D))[-1],\quad P_3=A_0(M_h(D))[-1].
\]
実際、順方向の変異は $J$ の下で球面twistの評価そのものであり、
$T_{\mathcal O}\mathcal O(h)=M_h[1]$ と前節の $G_i$ の表示を代入すればよい。

$U_F^{-1}P_i$ は共通の移動を除いて階数4のねじれベクトル束 $V_i$ であることを示す。
一般の点層の $U_F$ 像 $X=i_*J_4^{-1}$ に対し、$J_4$ は $F$ 上の一般の次数4線束である。
$\mathbf R\mathrm{Hom}(P_0,X)$ は次数1のみで次元4、$P_1$ についても
補題\ref{lem:octic-linear-determinant}の証明で $M_h^*|_F$ が次数4の部分線束を持たないので同じである。
$P_2,P_3$ のとき $A_0^{-1}X$ は相手の一般点に対応し、滑らかな $C$ 上の次数12の線束の $-1$ 移動である。
その移動を相殺し $D|_C$（次数16）を引けば、全く同じ次数 $-4$ の計算になる。
従って四対象は一般点で同一の移動を持つ階数4の束である。
各 $P_i$ は球面的なので、その自己 $\operatorname{Ext}^1$ は零である。
前節のスペクトル系列の次元不等式により各コホモロジー層も剛であり、
前節のねじれ階数零層の不等式により、一般点で見える一層以外は存在しない。
導来双対（逆Brauer類）も球面的なので同じことが言える。
層もその導来双対も局所的に次数零の層であれば、正則局所環で極小自由分解を取ることで局所自由性が従う。

どのねじれ層の階数も4の倍数である。実際その一般点層とのEuler対合は
$T$ 上の $(0,F,-4)$ との対合であり、$hF=8$ と格子の仮定から4で割り切れる。
よって各 $V_i$ は二つの鉛筆のいずれの一般繊維上でも幾何学的半安定である。
幾何学的一般繊維で不安定ならその最大不安定部分は一意性により基底の関数体上へ降下し、
曲面の開部分で階数1から3のねじれ部分層を与える。これを延長すれば階数の整除性に反する。
この降下はエタール局所表示のHarder--Narasimhan部分が遷移写像で保たれることからも分かる。
滑らかな一般閉繊維へは半安定性の開性を用いる。

二つの鉛筆上の傾きを計算する。$u_0=(0,0,1)$ とし、$A_0u_0=(2,-h,1)$ とすると
\[
 \big(v(P_i)\cdot u_0\big)_{i=0}^3=(-1,-3,-5,-7),\qquad
 \big(v(P_i)\cdot A_0u_0\big)_{i=0}^3=(-3,-1,1,3).
\]
これは左変異の式からの直接計算である。二つの列はそれぞれ、一般の $a,b$ 繊維上の
ねじれ線束の直像の移動との対合である。従ってそれぞれの繊維でその線束によりねじれを自明化すれば、
制限次数は符号と4の倍数の共通加数を除きこの二列である。
特に次数は奇数、階数は4なので、半安定な一般繊維制限は安定である。
$i>j$ では大切断の順方向射 $p_i\to p_j$ があり、その射は $J$ の下でも消えない
（上の単位の錐とtiltingより次数零の射は同型）。
$V_i\to V_j$ の非零写像は階数整除性から一般点同型で、行列式は有効因子を与える。
二つのnef繊維との次数差は上の計算で絶対値 $2(i-j)$、従って両方で正である。
結局、どちらの一般繊維でも傾きの差は $\mu(V_j)-\mu(V_i)=(i-j)/2$ である。

$L=a+b$ とし、大切断のSerre純粋性の成分を
$X_m^{ij}=\mathbf R\mathrm{Hom}(p_i,S^{-m}p_j[3m])$ と書く。
第一の随伴三角と $\Phi J=JS[-1]$、前節のテンソル表示から、$m\ge1$ で三角
\[
 X_m^{ij}\longrightarrow
 Y_m^{ij}=\mathbf R\mathrm{Hom}_M(V_i,V_j(mL))\longrightarrow X_{m-1}^{ij}[1]
 \longrightarrow X_m^{ij}[1]
\]
を得る。$X_0^{ij}$ はtiltingにより純粋、すなわち次数零に集中する。
$m=0$ の束の間では $H^1(Y_0^{ij})=0$、$H^2(Y_0^{ij})$ は逆方向の大切断の射空間の双対である。
$i<j$ の $X_1^{ij}$ は既に純粋である。実際定理\ref{thm:cut-serre}のバー次数の内部次数は
添字を逆にして $4+i-j\in\{1,2,3\}$ であり、次数一生成二次代数のTorはこの範囲で対角に集中する
（生成元と最小関係の次数、およびバー次数が内部次数以下であることを使う）。
このとき $Y_1^{ij}$ は三角から正次数消滅、束の間なので負次数も消える。

$i>j$ では $Y_0^{ij}$ が既に次数零に集中し、$a$ の一般繊維を一度加え、次に $b$ の一般繊維を一度加える。
各短完全列の制限項では、半安定性と正の傾き差により $H^1=0$ である。
実際楕円曲線のSerre双対で、より大きな傾きからより小さな傾きへの射となる。
よって $Y_1^{ij}$ は純粋である。$i=j$ では最初の $a$ を加える際、
制限自己準同型の $H^1$ から曲面の自己準同型の $H^2$ への連結写像は、
恒等射の制限の双対なので同型である（一般繊維では安定であり法線束は自明）。
従ってこの一段で両方の高次コホモロジーが消える。その次の $b$ の段は $ab=2$ により正の傾き差があり、
ここでも $Y_1^{ii}$ は純粋になる。
$m\ge1$ ではどの対でも $mL$ を既に加えた時点で一般繊維の傾き差は
$2m+(i-j)/2>0$ である。再び $a,b$ を一つずつ加えれば $Y_m^{ij}$ の純粋性は全 $m\ge1$ に伝播する。
上の三角の長完全列により $X_m^{ij}$ の純粋性も全 $m$ に帰納する。
定理\ref{thm:cut-serre}がKoszul性を与え、有限Frobenius条件が左右のAS正則性を与える。
\end{proof}

% !TeX root = ../papers/nc-p3.tex
\section{一般の旗を用いた反対階数10の有限片側対称枝の完結}
\label{sec:finite-rank-ten-full-regularity}
\begin{theorem}\label{thm:finite-rank-ten-full-regularity}
標準有限Frobenius窓の全辺が有限で、一方の反対交換が対称かつ反対階数が10とする。
左右Koszul AS正則性の必要十分条件は、二つの奇位置のインスタントンモナドがともに束を与えることである。
従って定理\ref{thm:finite-elliptic-octic-windows}の開条件を満たす全データが正則である。
四次Grassmann写像による添字保持同型・周期の判定はこの枝の特殊値にも適用できる。
\end{theorem}
\begin{proof}
必要性は前の境界排除で得た。十分性では窓を一つ固定する。
以下で一般に選ぶのはこの固定された写像の旗であり、窓そのものを一般化しない。

\paragraph{次数二の対応の全点での定義。}
片側の一般切断の零曲線 $F$ を固定し、$L=\mathcal O_F(h)$ とする。
補題\ref{lem:octic-linear-determinant}の $J\in\mathrm{Pic}^4(F)$ に対する四次行列式を $q_J$ と書く。
これらは $\mathbb P H^0(I_F(4))=\mathbb P^2$ の中の直線をちょうど二重に覆い、恒等的に零になる $q_J$ はない。
実際、幾何学的なPicard曲線上で正規化した普遍線束を使い、
$H^0(J)$ と $H^0(LJ)$ の束を $B_4,B_{12}$ と書く。階数は添字通りで、次数はともに $-1$。
後者の次数公式は楕円曲線の積の正規化Poincaré束の第一Chern類の平方の繊維積分が
点類の $-2$ 倍であることからRiemann--Rochで従う
（第一Chern類は対角から座標二繊維を引いた類である）。
従って $K_J$ の束の次数は $4(-1)-(-1)=-3$ で、$q_J(x)$ は $J$ の関数として
固定した次数2の線束の切断となる。もし係数が一次元しか張らなければ行列式は
$q(x)e(J)$ に分かれる。非零四次式 $q$ の零曲面上で全 $J$ の評価が失敗するが、
同補題の証明では、このような $x\in\mathbb P^3$ は $F$ と高々一次元の平面射影例外集合に限られた。
従って係数は二次元全部を張り、次数2の完全鉛筆は基点を持たない。
また $x\in F$ で評価した関係は $J_x$ で零になる三次元部分に入るので $q_J$ は $F$ を含む。
これで直線とその二重被覆を得る。

$F$ を与える切断を含む二平面を動かすと対応する $T$ は上の網の全てを動く。
実際、切断と第二切断の行列式を取る写像はSerreの評価列より
切断の四空間を第一切断で割った空間から $H^0(I_F(4))$ への同型である。
一般の旗では $T$ は行列式の直線の外にある。よって全 $J$ で $q_J|_T$ は恒等的に零ではなく、
その $F$ を除いた因子は $|C|$ の一繊維である。
得られた写像 $g_F:\mathrm{Pic}^{-4}(F)\to |C|$ は逆線束を取ってから上の直線を $[T]$ から射影したもので次数2。
ここで残余因子が定理\ref{thm:octic-fiber-twist-transfer}の余核の行列式に等しいことは、
$W\subset H^0(J)\otimes\mathcal O_T$ の行列式が $\mathcal O_T(-F)$ であることからも従う。
二側を交換してもこの旗の開条件は成立する。よって一般の $T$ とその両鉛筆の一般の $F,C$ で同時に使える。

\paragraph{二つの繊維制限。}
$U$ を滑らかな楕円繊維とそのPicard曲線の普遍層による変換とする。
始域にねじれを入れれば基点を選ばず定義でき、幾何学的には通常の楕円曲線の同値である
（前節の相対対角イデアルの構成を一繊維に適用してもよい）。
$F$ で次数 $-4$、$C$ で次数12を使い、
\[
 I_F=i_{F*}U,\quad r_F=U^{-1}\mathbf L i_F^*,\qquad
 I_C=A_0 i_{C*}U,\quad r_C=U^{-1}\mathbf L i_C^* A_0^{-1}
\]
とする。$r$ はそれぞれ $I$ の左随伴である。
$D_F=L_F,\ D_C=A_0 L_C A_0^{-1}$ と書けば $\Psi=\Phi^{-1}[2]=D_F D_C$ である。
$D_F$ は $r_F$ 上で恒等、$r_C$ 上で次数2の線束のテンソル積として働く。
$D_C$ は $r_C$ 上で恒等、$r_F$ 上で次数2の線束のテンソル積として働く。

この最後の主張は各点の数値作用だけからは従わないので族で確認する。
$A_0^{-1}$ を $I_F$ の普遍核に適用すると、前段により全パラメータで
同じ移動の行列式余核になる。各行列は束間の単射で、余核はパラメータに平坦である。
随伴行列によりその余核は $T\times_{|C|}\mathrm{Pic}^{-4}(F)$ にスキームとして支えられる。
従って $L_C$ はこの核に $g_F^*\mathcal O(1)$ を掛けるのと同じであり、
$D_C I_F=I_F(-\otimes g_F^*\mathcal O(1))$、左随伴により所望の $r_F D_C$ の式となる。
逆方向には命題\ref{prop:quartic-k3-two-pencil-twists}の $B_0A_0=L_{-4h}[2]$ より
$A_0=L_{-2h}A_0^{-1}L_{-2h}[2]$ であることを使う。次数12の線束を $-2h$ 捩れば次数 $-4$ なので、
$C,F$ を交換した同じ行列式で、$A_0 I$ の相手繊維への台の写像も次数2となる。
同じ側の作用は繊維の法線束が自明なことから従う。
この証明は一般の印のPicard曲線にわたる全点の族を使い、相手の特異繊維での安定性を必要としない。

\paragraph{制限の階数、次数、半安定性。}
前節の $P_i$ を用いる。$r_F P_i$ は同じ移動を持つ階数4の束である。
$i=0,1$ の全 $J$ でのコホモロジーは前節と補題\ref{lem:octic-linear-determinant}の通りである。
$i=2,3$ には
\[
 \Psi P_2=M_h^*(D+h)[1],\qquad \Psi P_3=\mathcal O(D+h)[1]
\]
を使う。実際 $T_{G_2}^{-1}G_2=G_2[1]$ と双対評価列を左変異表示に代入すればよい。
$\Psi$ は既に $r_F$ 上で線束の捩りなので、右辺で共通の移動を検査すれば足りる。
$D+h$ の $F$ 制限は $-h$ であり、次数 $-4$ の線束に $h$ を足すと次数4になる。
その $M_h$ との積の $H^1$ は同補題で零、$H^0$ は四次元である。
従って四対象とも $\mathbf R\mathrm{Hom}(P_i,I_F k(z))$ は全点で次数1に四次元だけを持つ。
一方 $A_0^{-1}P_i$ は順に
$M_h^*(h),\mathcal O(h),\mathcal O(D)[-1],M_h(D)[-1]$。
$D|_C=2h|_C$ なので、次数12の線束に対して同じ計算が全点で次数零の四次元を与える。
コホモロジーと基底変換により $r_C P_i$ も共通の移動の階数4束となる。

各側の一般の旗の関数体 $K$ で、印を付けた楕円繊維の全ての $K$ 線束の次数は8の倍数である。
実際、旗は切断の直線 $s$ とそれを含む二平面 $P$ からなる。
普遍零繊維の全空間は $x\in\mathbb P^3$、$s\subset f(x)$、$s\subset P\subset H$ を順に選ぶ射影束塔である。
従ってPicard群は $h$ と二つの相対同語反復線束で生成され、後二者は $s,P/s$ による旗底からの引戻しである。
滑らかな一般繊維の線束は有理切断の因子の閉包をこの正則全空間で取って延長でき、繊維に残るのは $h$ のみである。
反対側は双対Grassmann写像で同じである。

従って $K$ 上の対応するねじれPicard曲線の任意の層の階数は4の倍数である。
実際 $U$ による像を楕円繊維上で $(r,d)$ と書けば $d$ は8の倍数で、幾何学的点層とのEuler対合は
$-4r-d$ または $12r-d$ である。この対合がねじれ層の階数である。
各束の幾何学的最大不安定部分は一意であり、エタール局所の表示で降下して $K$ 上のねじれ部分層になる。
階数4未満の正階数を排除したので四束とも幾何学的半安定である。
この性質は一般の旗へ開に広がる。幾何学的Picard曲線で普遍線束を正規化したとき、
次数 $e$ の変換 $U$ の数値行列は階数・次数順で
$\left(\begin{smallmatrix}0&1\\-1&e\end{smallmatrix}\right)$ である
（点は次数 $e$ の線束、構造層は符号反転した長さ1の点類へ移る）。
従って共通移動を除いた $r_F P_i$ の次数は $(-1,-3,-5,-7)$、$r_C P_i$ の次数は $(3,1,-1,-3)$ である。
実際 $P_i|_F$ の階数は $(1,3,5,7)$、次数は $(0,-8,-16,-24)$、
$(A_0^{-1}P_i)|_C$ の階数は $(3,1,-1,-3)$、次数は $(32,8,-16,-40)$ である。
自明化の変更は次数に共通の4の倍数を加えるだけである。
奇数次数なのでこれらは安定であり、両側で傾きの差は $(i-j)/2$ である。

\paragraph{全次数への伝播。}
$Z_m^{ij}=\mathbf R\mathrm{Hom}_T(P_i,\Psi^mP_j)$ とする。
$Q\to D_C Q\to I_C r_C Q$ と $Q\to D_F Q\to I_F r_F Q$ は
それぞれ一般の $C,F$ の繊維方程式による三角である。
$P_i$ から射を取った第三項は楕円Picard曲線上の $r_C$ または $r_F$ の間の導来Homそのものである。
従って $Z_m^{ij}$ から $Z_{m+1}^{ij}$ への二段で加わる項は、
同じ移動の二束間の導来Homであり、その傾き差は順に
\[
 (i-j)/2+2m,\qquad (i-j)/2+2m+2.
\]
ここで先に $D_C$ を掛ける。二関手の可換性は使っていない。
$Z_0^{ij}$ は前節の随伴三角から負次数と次数1で消える。
従って二段を経ても負次数のコホモロジーは出ない。
最初の周で $i<j$ のときは前節の内部次数 $4+i-j\le3$ の議論と随伴三角から $Z_1^{ij}$ の正次数も消える。
$i>j$ のときは初期の次数2も消え、二段とも正の傾き差なので楕円曲線上の $H^1$ は消える。
$i=j$ では第一段の連結写像は、Serre双対の下で恒等射を安定な制限束の恒等射へ送る制限の双対であり同型である。
これは $A_0^{-1}$ で繊維制限の三角へ戻し、繊維の法線束の自明性を使えばよい。
よって次数1と2が消え、次の段では正の傾き差がある。
以後 $m\ge1$ ならいずれの傾き差も正である。よって全 $m\ge1$ で $Z_m^{ij}$ は純粋となる。
前節の $X_m^{ij}\to Z_m^{ij}\to X_{m-1}^{ij}[1]$ は随伴によるので格子の仮定なしに使え、
帰納的に全Serre純粋性を得る。これで十分性と枝の判定が従う。
\end{proof}

% !TeX root = ../papers/nc-p3.tex
\section{非対称な四窓の滑らかな三文字切断と二つの鉛筆}
\label{sec:finite-determinantal-marked-slice}
\begin{theorem}\label{thm:finite-determinantal-marked-slice}
有限な四辺を持つ標準有限Frobenius窓に $x\in\mathbb P(V_0^*)$ を印し、
$t=w(x)\in V_1\otimes V_2\otimes V_3$ の三つの四次行列式曲面が滑らかとする。
中央の曲面 $T\subset\mathbb P(V_2^*)$ はK3である。順序付き三偏極を
$l,h,r$ と書けば
\[
 l^2=h^2=r^2=4,\quad lh=hr=6,\quad l+r=3h .
\]
$A=l+h, B=h+r$ とすると、元の長さ20の点グラフは $T$ 上で
$|A|$ と $|B|$ のそれぞれの鉛筆の基底スキームとなる。
一方の鉛筆を指定すれば他方は一意である。

逆に三文字 $t$ を固定すると、これを指定切断とする全ての窓を次の有限方程式で復元できる。
$\Lambda\in\mathrm{Gr}(2,H^0(T,A))=\mathrm{Gr}(2,12)$ の基底スキーム $X$ が零次元で、
$H^0(I_{X/T}(B))$ が二次元かつその基底スキームも零次元となることを課す。
二鉛筆をそれぞれ双線型式へ引き戻して得た六空間を $R_1,R_2$ とする。
三次重なりが四次元であることと通常の有限Frobenius・有限辺検査を課し、
その重なりに $t$ を印す。これでちょうど指定した候補を得る。
印付き添字保持同型は順序を保つ $t$ の三基底変更と鉛筆の対応と同値である。
\end{theorem}
\begin{proof}
$t$ を中央点の四行四列線型行列と読む。滑らかな行列式上では階数は全点で3であり、左右核点が定まる。
他の二位置の行列式でも同じため、この二つの核写像は四次曲面どうしの同型である。
隣り合う二射影のグラフは四つの $(1,1)$ 式の完全交叉である。
実際単元三次小行列で核を解けばその方程式と残った一つの行列式とに局所分解する。
グラフの積空間内の類は $(l+h)^4$ であり、これから $lh=6$、他側も同様である。
余因子三次式は同時に消えず、左右核のSegre座標を与えるので $l+r=3h$。
完全交叉のKoszul列を $(1,1)$ 捩れば、十六双線型切断から $H^0(l+h)$ への制限は全射で、
その核は $t$ の四縮約である。反対も同様である。
実際負の中間項には $\mathcal O(-j,-j)$、$1\le j\le3$ しか出ず非零コホモロジーを持たない。
六関係中の残り二式が鉛筆を与え、元の六式の全零点はまさにその基底スキームである。
中央射影が点グラフの両側を同じスキームに同定するので二鉛筆の基底は一致する。
$A^2=B^2=20, AB=30$ である。

逆に両方の基底が有限ならそれぞれ長さ20であり、後者が前者を含むので等しい。
有限基底 $X$ の完全交叉列 $0\to\mathcal O(-B)\to\mathcal O^2\to I_{X/T}(B)\to0$
から $H^0(I_{X/T}(B))$ は鉛筆そのものである（$B$ は豊富で $H^1(-B)=0$）。
元の二窓の六式は各鉛筆の四核つき逆像に等しい。
$t$ は三次重なりに入り、系\ref{cor:two-window-extension}の復元によりこれがちょうど印付き四窓を与える。
最後の有限検査は正則性を仮定せずに復元テンソルに施せる。基底変更と逆操作が両立するので同型条件も従う。
\end{proof}

\begin{proposition}\label{prop:finite-marked-eight-determinant}
$\mathrm{NS}(T)=\mathbb Z l\oplus\mathbb Z h$ なら、有限基底を持つ上記の鉛筆で
反対鉛筆の候補は八次元空間間の乗法だけから決まる。
\[
 H^0(I_{X/T}(B))=\ker\big[H^1(B-2A)\longrightarrow
                 \Lambda^*\otimes H^1(B-A)\big]
\]
と体積因子を除き同定でき、両 $H^1$ は八次元である。
従って必要な第一の退化条件は八列・十六行の行列が階数6であることとなる。
\end{proposition}
\begin{proof}
$B-A=3h-2l, B-2A=2h-3l$ はともに平方 $-20$。
格子の平方は4の倍数なので曲面には $(-2)$ 曲線がない。
従ってこれらの因子とその負は有効でなく、Riemann--Rochで $h^1=8$ を得る。
二切断のKoszul列を $B$ 捩ると主張の核になる。
行列の各ブロックは対応する $A$ の切断による乗法である。
\end{proof}

\begin{theorem}\label{thm:dominant-four-by-four-by-four-slice}
十五次元混合成分に一般の $x$ を印し三文字へ縮約すると、
\emph{一般の全ての}四行四列四層テンソルの射影同値類に達し、
その印付き拡張はこの成分で生成的有限である。
従って前定理の滑らかなK3切断は十五次元成分の稠密開集合を被覆する。
\end{theorem}
\begin{proof}
定理\ref{thm:full-fifteen-component}の有理点 $t_0$ で $x=(1,2,3,5)$ を固定する。
基底つき二切断階数6軌跡はそこで滑らかで次元76である。
各切断行列の左右核から $u^t\dot w v=0$ の百式を作り、二組を同じ256座標に戻す。
素数1009でこの連立核は次元76、$\dot w\mapsto\dot w(x)$ の核への制限は階数64である。
元の二切断の核は局所整数環で直和因子であり、連立核次元が標数零と一致するので接基底が持ち上がる。
従って標数零でこの縮約は支配的である。
同じ点で三方向の行列式四次式の四偏微分を全六次単項式倍すると、
三方向とも九次全空間の次元220を張る。従って三つの射影曲面は滑らかである。
正確な行列・座標順と剰余検算は \path{computations/verify_determinantal_marked_slice.py} にある。
一般のファイバーの基底付き次元は $76-64=12$ であり、
$x$ を線型的に固定する第一空間の基底変更が既に次元12を占める。
実際四テンソルの一文字支持は充満なので、それだけを動かしてテンソルを固定する元は恒等である。
固定した $x$ を任意の印へ移すこと、および三文字の基底変更を両側で割ることから結論が従う。
\end{proof}

% !TeX root = ../papers/nc-p3.tex
\section{順次中心になるパラメータに対する一様な有限判定}
\label{sec:successively-central-cutoff}
\begin{theorem}\label{thm:successively-central-cutoff}
標準有限Frobenius候補の折り畳み環 $B$ を取り、正次数 $d_i$ の斉次な頂点対角元
$g_1,\ldots,g_4$ が $B_{i-1}=B/(g_1,\ldots,g_{i-1})$ で順に中心になるとする。
$B_4$ が有限次元なら次が左右Koszul AS正則性の必要十分条件である。
$L=\sum(d_i-1)$、$b=\max(1,L,d_1,\ldots,d_4)$ とする。
\begin{enumerate}
\item 各始点で $B_4$ の次数 $0,\ldots,L+1$ のHilbert関数が
$\prod_i(1+z+\cdots+z^{d_i-1})$ の係数に等しい。
\item $B$ の各始点の経路次元は $3b$ 次まで標準で、左右の一方の標準Koszul複体は
内部次数 $L+4$ まで完全である。
\end{enumerate}
第一条件があれば有限性も同時に証明される。成立すれば四元は順に両側非零因子である。
四元とも四次なら経路次元は36次、複体は16次までで足りる。
各元は前の元を割った環でだけ中心でよく、元の環で四元が生成する部分代数の表示は要らない。
\end{theorem}
\begin{proof}
まず必要性を示す。$B_i$ はいずれも次数付き左右Noether環である。
実際 $R/(g)$ が次数付きNoetherで $g$ が中心なら、$(g)$ の随伴次数環は
$(R/(g))[t]$ の商であり、斉次イデアルの初項を有限生成して持ち上げられる。
内部次数を固定すれば正次数 $g$ の冪のフィルトレーションは有限なので、持上げは停止する。
有限な $B_4$ から逆向きに適用すればよい。
$B$ がAS正則なら、任意の有限生成次数付き $B_i$ 加群 $N$ について
$\underline{\operatorname{Ext}}^j_B(N,B)=0$ ($j<i$) が成り立つ。
$i=4$ では有限長とAS条件による。$i+1$ から $i$ へは
$N(-d)\xrightarrow{g_{i+1}}N$ の核・余核が有限生成 $B_{i+1}$ 加群であることを使う。
像で二つの短完全列に分ければ $j<i$ でこの写像はExtに同型を誘導する。
Extは有限生成自由解消の各項から内部次数が下に有界で、非零の次数移動と同型なら零となる。
以上は左でも右でも通る。$g_1,\ldots,g_{i-1}$ が既に正則列なら中心非零因子の二項分解を繰り返して
$\operatorname{Hom}_{B_{i-1}}(B_i,B_{i-1})$ は捩りを除いて $\operatorname{Ext}_B^{i-1}(B_i,B)$ に等しい。
これが零なので $g_i$ も非零因子である。四段の短完全列で第一条件が従い、第二条件は正則性による。

十分性では $B_4$ の頂点・次数別の語基底の正次数部分を $v$ と書き、各 $v$ を同じ経路で持ち上げる。
各 $B_i$ を順に割ってから持ち上げると、$B$ は順序つき記号
\[
                  g_1^{a_1}g_2^{a_2}g_3^{a_3}g_4^{a_4}v
\]
で張られる（末尾には頂点単位も許す）。実際、各段の核はその段での中心性により
$g_i B_{i-1}$ であり、正次数ゆえ次数帰納が使える。
端点に応じて対角元の対応成分だけを読む。この順序記号の個数は第一条件から各次数で標準値である。
第二条件により $3b$ 次までは既に独立である。

これを全次数へ延ばす。記号 $g_i,v$ の積に、$vv'$、$vg_i$、および $g_j g_i$ ($j>i$) を
それぞれ順序表示へ直す規則を入れる。$vg_i$ は $g_i v$ と、
より早い $g_s$ ($s<i$) が最初に立つ項の和とする。$g_jg_i$ も同様である。
この形は順次中心性と先ほどの持上げで常に選べる。$vv'$ は任意の順序表示とする。
内部次数固定のもとで、各書換えは陽に現れる $g_1,\ldots,g_4$ の個数の列を辞書順に増やすか、
同数なら末尾基底文字の個数を減らすか、どちらも同じなら逆転を一つ減らす。
$vv'$ から中心記号が出ない場合は一つの末尾基底文字だけになる。よって書換えは停止する。
規則は全て二記号に対するもので、重なりは三記号、重みは $3b$ 以下である。
この範囲の独立性により両側の簡約は同じ係数になり、ダイヤモンド補題で全順序記号が
これらの規則で定義した環の基底となる。
これが $B$ と同じ環であることも必要である。
元の矢を次数1の順序表示に送ると、二次関係は $2\le3b$ の独立性で零に簡約する。
また各 $g_i,v$ の元の経路式も重み $b$ 以下の独立性でその記号自身に戻る。
従って元の表示と規則による表示は互いに逆である。

最初の $g_1$ による左乗法は指数を一つ増やす。
その商でも次の $g_2$ について同じであり、順に中心なので両側の正則列を得る。
最後に各標準片側複体を $B_i$ に係数変更したものを $K_i$ とする。
正則列の短完全列により $K_i(-d_{i+1})\to K_i\to K_{i+1}$ は短完全列を与える。
$N=L+4=\sum d_i$ までの仮定から、$K_i$ の非増大複体のホモロジーはこの範囲で
内部次数が高々 $\sum_{j\le i}d_j$ でしか現れない。これは上の短完全列と $K_0$ の既知の消滅による帰納である。
$K_4$ 自体は $N$ より上に項を持たない。
$i=3,2,1,0$ の順に戻ると、$n>N$ では $H(K_i)_{n-d_{i+1}}\to H(K_i)_n$ は全射である。
最初に参照する $N$ 以下の次数は $\sum_{j\le i}d_j$ より大きいので零、以後は内部次数帰納で零となる。
よって $K_0$ は全次数で完全で、有限Frobenius対合が反対側とAS条件を与える。
\end{proof}

% !TeX root = ../papers/nc-p3.tex
\section{非対称成分の滑らかな種数三切断による復元}
\label{sec:finite-genus-three-curve-cut}
\begin{theorem}\label{thm:finite-genus-three-curve-cut}
定理\ref{thm:quartic-center-order}の十五次元成分の開集合にある任意の窓は、
滑らかな平面四次曲線 $C$ と四つの半安定な階数4の束 $E_0,\ldots,E_3$ を用いて復元できる。
$E_{i+4}=E_i\otimes\omega_C$ と置けば $\deg E_{i+1}-\deg E_i=4$ であり、
一般の二本の中心四次式 $g_1,g_2$ を割った全経路代数は
\[
 D_{i,i+n}=H^0(C,E_i^*\otimes E_{i+n})\quad(n\ge0)
\]
の合成代数である。特に二つの隣接窓は $n=2$ の積の核として定まる。
三次元射影中心が二十の既知節点以外で滑らかであるという仮定を置く必要はない。
\end{theorem}
\begin{proof}
五変数の射影四次中心は整かつ正規なので、その特異集合は次元1以下である。
一般平面はそれを避け、Bertiniにより滑らかな整平面四次曲線を切る。
Azumaya軌跡の命題\ref{prop:smooth-azumaya}により全曲線上にAzumaya代数を持つ。
一般の二式は四つのNoether正規化の一次部分の中の最初の二元にできる。
実際、四つの一次式で割ると有限になる条件は開で非空であり、一般の二式に残り二式を足せばよい。
中心と折り畳み環は次元4の極大Cohen--Macaulay加群であるから一般の四式上で自由である。
あるいは一組の既知の多項式自由表示から深さ4、従って任意の斉次パラメータ系の正則性を得る。
二式を割った隅は残りの二変数上の自由加群で、四移動ごとにまとめても深さ2を持つ。
よってこれらは射影曲線上の対応する層の全切断加群に等しい
（無関係イデアルの零次・一次局所コホモロジーが消える）。

一般平面上でこの層表示を具体化する。経路長を $4m+j-i$ にまとめた $ij$ 隅を次数 $m$ として層化する。
中心次数1が $\mathcal O_C(1)=\omega_C$ の捩りに対応する。
全隅を合わせると一般中心点で四冪等元を持つ次数16のAzumaya代数である。
曲線の各点で可逆な中心一次式により剰余次数を自明化して見ると、
これは元の次数環をその値1で割った局所表示である。四射影の行列階数は4ずつで一定となる。
幾何学的曲線上でBrauer群は零なので、全体を一つの束の自己準同型として分裂できる。
積の向きを合わせて必要なら反対代数を取り、その四成分を $E_i$ とすれば上式を得る。
二式の正則性より各始点のHilbert級数は $(1-z^4)^2/(1-z)^4$。
大きな $m$ でRiemann--Rochの $4\,(\deg E_j-\deg E_i)+64m-32$
をこれに比較すると $\deg E_j-\deg E_i=4(j-i)$ となる。

半安定性は固定窓の一般平面で示せる。平面の関数体上の曲線を取ると、
その関数体は射影中心の関数体の純超越拡大である。実際、点とそれを含む平面の入射の射影中心上の
一般ファイバーは点の商空間における二平面のGrassmann多様体で有理である。
従って系\ref{cor:quartic-division-corners}の次数4可除性は保たれる。
四つの束をこの体上でねじれ束として読むと、ねじれた一般階数は4の倍数に限られる。
幾何学的不安定部分は降下して階数1から3の部分を与えてしまうため存在しない。
一般の閉平面での半安定性は開性で従い、Brauer分裂の共通線束の自由度にも依存しない。
\end{proof}

\begin{proposition}\label{prop:finite-genus-three-curve-recognition}
逆に滑らかな種数3の $C$ と上記の次数・半安定性の四束を与える。
$n=0,1,2,3$ の切断次元が各始点で $1,4,10,20$ であるとし、
四つの一次空間の積から四つの二次核を取り、それらで定義した折り畳み環を $B$ とする。
次の範囲の条件だけで、$B$ が標準左右AS正則であり、二次元の中心四次鉛筆 $G$ を割ると
ちょうどこの切断代数になるような表示を判定できる。
\begin{enumerate}
\item 標準有限Frobenius条件と、曲線への像が零である独立な四次中心 $g_1,g_2$ を検査する。
\item $B/(g_1,g_2)$ から曲線の切断代数への評価が経路長13まで全端点で同型であることを課す。
\item 元の経路次元を36次、標準複体の完全性を16次まで調べる。
\end{enumerate}
元の窓が直前の定理の開集合に属することを逆判定の仮定に含めない。
$G$ を指定した折り畳み環（次数・頂点を固定）の同型は、曲線の同型と四束の共通線束捩りまでの同型と必要十分である。
鉛筆を忘れた添字保持同型は、復元した二窓の三基底変更で判定できる。
\end{proposition}
\begin{proof}
半安定性から $n<0$ の切断は零で、$n>4$ ではSerre双対により第一コホモロジーが零となる。
$n=4$ では第一コホモロジーは自己準同型の双対で一次元である。
Riemann--Rochと仮定から全切断代数のHilbert級数は先ほどの二式の商に一致する。
標準束に共通零点のない二切断 $s,t$ を選ぶ。各隅でこれらは切断加群の正則列である。
実際、束の無捩れ性により $s$ は非零因子で、$tf=sf'$ なら $f/s$ は $s$ の零点でも正則である。
負の長さの切断は零なので切り詰めによる例外も生じない。よって二式を割ると
各始点で $(1+z+z^2+z^3)^4$ が残り13次以後は零である。
この二切断の恒等倍を全頂点で取り、第二検査により $B/(g_1,g_2)$ へ持ち上げる。
五次の評価が単射なので全矢とそこで可換である。これで前定理の順次中心四元が揃う。
十三次までの同型より最終商の必要次元を得、第三検査と定理\ref{thm:successively-central-cutoff}が
全次数正則性と四元の順次正則性を与える。
切断代数は $s,t$ と十二次以下の切断で生成されるので評価は全次数で全射であり、
二中心四次の正則性からHilbert級数も一致し同型となる。
逆にこの表示があって正則なら四条件の次元・有限検査は自動的に満たされる。

最後に $G$ を割った零番隅の四Veronese環の中心は $\bigoplus_m H^0(m\omega_C)$ である。
実際高い捩りの全切断が各点で自己準同型束を生成するので、中心元は一般点でスカラー、
そのスカラー切断は全点で正則である。この中心のProjと偏極から曲線と標準束を取り戻す。
隅を層化すると $\mathcal End(E_0)$ が得られ、その分裂束は線束捩りまで一意である。
次に他の三隅の $\mathcal Hom(E_0,E_i)$ を加群として読めば同じ共通捩りで $E_i$ が定まる。
逆向きにはこの曲線データから積と二次関係を復元できる。
鉛筆も評価の四次核のうちの頂点整合中心として一意である。実際 $B/(G)$ で消える四次元は
各頂点で二本の独立な $g_1,g_2$ の成分であり、隣り合う頂点の係数差が作る一次関係は
二元が正則列で全次数商のHilbert級数が上の通りであることから零となる。
具体的には二元の最初の係数シジジーは全次数8で、次数5にはない。
従って全四頂点で同じ二係数を持たねばならず、消える中心は正確に $G$ である。
\end{proof}

% !TeX root = ../papers/nc-p3.tex
\section{十五次元成分の一般中心の孤立特異性と四次K3切断}
\label{sec:finite-central-k3-general-section}
\begin{theorem}\label{thm:finite-central-k3-general-section}
十五次元混合成分のある稠密開集合では、射影中心の一般超平面切断は滑らかな四次K3曲面であり、
射影中心の特異点は有限である。一般窓の既知の二十節点以外の、次元ベクトル $(2,2,2,2)$
を持つ単純も、次数による倍率を除いて高々有限個である。
一般超平面上の次数4のAzumaya隅、四つのねじれ階数4束と偏極から、
元の窓の四層切断を全ての積つきで取り戻せる。
\end{theorem}
\begin{proof}
まず十四次元片側対称枝の一般の窓に一般の二平面を印す。
この十八次元の族では、幾何学的一般の $(T,h,F)$ のPicard格子は $h,F$ で生成される。
実際 $(T,h,F)$ から、一般の繊維でインスタントン束を復元し、鉛筆から切断二平面を復元できる。
一繊維を零点とする切断の直線はSerre拡大の一意性と束の単純性により一意である。
よって偏極K3のモジュライでの像次元は18である。一つの偏極K3上のこの数値の繊維類は
Hodge指数定理で高々有限個である。複素数上のK3の局所周期写像で三独立な因子を保つ変形は
次元17以下なので、可算個のそうした軌跡は幾何学的一般点を含まない。
これは独立な三つの第一Chern類が二形式の周期に独立な三線型条件を課すことによる。
基礎体の変更は幾何学的一般点で行える。生成された格子の判別式の絶対値は64で、
真の飽和があれば指数に2を含む。$F/2$ は繊維の原始性で存在せず、
$(h+jF)/2$ の平方は奇数でK3の偶格子に入らない。これで格子そのものも確定する。

対称枝のこの幾何学的一般点は定理\ref{thm:symmetric-center-cover}の開近傍にも属する。
ここでは四次中心が五次元で、中心のHilbert級数は換算次数で
$(1+t+t^2+t^3)/(1-t)^4$ である。隅の商は可除で、合成可能な非零元の積は零でない。
定理\ref{thm:finite-minimal-picard-regularity}の構成を使い、同所の大切断の四射影に
対応するねじれ束を $V_i$ とする。共通移動を除き、
$\Theta=S^{-1}[3]$ の作用は制限関手 $J$ の下で $L=a+b$ による捩りとなる。
大切断の射を反復して得る軌道代数の各隅を、順序を合わせて
$\mathrm{Hom}(p_i,\Theta^m p_j)$ と書く。これは元の折り畳み経路代数である。
実際定理\ref{thm:cut-serre}と命題\ref{prop:two-window-ext-complex}により
戻り辺一回の両側加群の導来テンソル冪は純粋で通常のテンソル冪と一致し、
合成もその連接である。

随伴三角の連結射 $S[-3]\to\mathrm{Id}$ を $\Theta p_i$ で評価すると
長さ4の四ループ $g_i:p_i\to\Theta p_i$ を得る。
その和を $g$ とする。$J$ と $J\Theta=LJ$ で誘導される評価は乗法と両立し、
各隅でその像は $\mathrm{Hom}(V_i,V_j(mL))$、核は $g$ を始点に付けた短い経路である。
実際 $m\ge1$ は前節の短完全となった随伴三角であり、負次数の束の間のExtはない。
同じ事実は $m=0$ の次数零の射にも成立する。単射は $m=1,i=j$ で $g_i\ne0$ も示す。
連結射の自然性により次数零の大切断の全矢と $g$ は可換である。
また終点に $g$ を付けた像も評価核に入る。この端からの積も非零経路の積の性質で単射で、
次元が等しいので $gB=Bg$。従って $g$ は正則正規元である。
対応自己同型は既に三辺の文字を固定する。四番目は三次重なりの双対であり、
四次socleテンソルの直線を保つので、そこでは一つのスカラー倍である。
正規自己同型は $g$ 自身を固定し、各 $g_i$ は戻り辺をちょうど一つ含むので、このスカラーも1である。
従って $g$ は中心元である。

零番隅の四Veronese商はよって $\bigoplus_m H^0(M,\mathcal End(V_0)(mL))$ である。
その中心は $\bigoplus_m H^0(M,mL)$ である。高い捩りの自己準同型束は全点で生成されるため、
中心元はスカラーであり、スカラー係数は局所行列表示から全点で正則である。
元の全中心 $Z$ の零番隅への写像は単射である（一般商で既に単射）。
さらに $Z/gZ$ もこの商隅へ単射である。実際 $z$ の零番成分が $g$ 倍なら、
評価した中心元の他の成分も、零番から各頂点への非零の束写像との可換性で消える。
各成分は中心のため一般点でスカラー（高い捩りの各隅の全射を使う）であり、
非零の隅間写像は一般点で非零なので同じスカラーを持つからである。
全体で $z=gb$ と書ければ $g$ の中心性と正則性から $b$ は中心である。
K3上で $L^2=4$、$L$ は豊富なので $h^0(mL)=2+2m^2$ ($m>0$)。
これは先の $Z$ のHilbert級数を $(1-t)$ 倍した全係数と一致する。
従って $\operatorname{Proj}(Z/gZ)=M$ は滑らかな曲面である。

中心は定理\ref{thm:symmetric-center-cover}の近傍で重み付き二次・四次完全交叉の平坦族に入り、
四次中心は階数5の束である。超平面切断が滑らかな曲面となる条件は固有性と滑らかさの開性で開である。
上で一つの幾何点と超平面が得られたので、十五次元成分の一般点とその一般超平面で成立する。
対称因子の外の一般点では中心は普通の四次超曲面である。もし射影特異集合が正次元なら、
超平面と会い、その点で切断の接次元は少なくとも3なので滑らかでない。従って特異集合は有限である。
各特異指標のファイバーは有限次元代数で単純は有限個である。
最後に滑らかな超平面を割った隅は、残る三中心変数上で自由で深さ3を持つ。
従って次数別に全切断加群であり、第一コホモロジーも全捩りで消える。
Azumaya層のねじれた分裂束から四冪等元の部分 $V_i$ を取れば、
$n<4$ の積は中心の四次式を割っても変わらず、元の大切断を復元する。
\end{proof}

\begin{proposition}\label{prop:extra-central-singular-characters}
一般の十五次元成分には、既知の二十節点のほかに特異中心指標が実在する。
特に次元ベクトル $(2,2,2,2)$ の単純が存在する。
全ての残りの特異点も通常二重点で、それぞれ非同型単純を二種類持つ場合には、
その個数はちょうど10である。
\end{proposition}
\begin{proof}
複素数上で計算する。使う条件と一般の結論は有限型のデータから標数零の幾何学的一般点に移る。
射影中心 $X$ 上で、四隅の全てを前定理の次数規約で層化した代数を $\mathcal B$ とする。
その完全圏の周期巡回ホモロジーの偶次元数から奇次元数を引いた値は4である。
理由を詳しく述べる。$\operatorname{qgr} B$ は $\mathcal B$ の連接加群圏に同値である。
実際、中心の四次Veronese次数で隅をまとめた有限次数加群の層化と、
高い捩りの全切断を取る逆操作は、通常の有限生成中心に対するSerre対応そのものである。
四つの頂点成分と固定した経路剰余次数を取れば、ここでは整数頂点の一つの連結な圏を使っている。
四つの連続射影の像には高次の順方向Extがなく、逆方向のExtは全て零である。
実際それらの間の層Homは対応する隅の次数付き極大Cohen--Macaulay加群であり、
局所的な射影性と深さ4から中間コホモロジーが消える。
最高コホモロジーは中心のGorenstein双対で計算できる。
射影四次超曲面の双対化層は $\mathcal O_X(-1)$、また命題\ref{prop:smooth-azumaya}のtraceで
隅の双対は反対隅になる。従って最高次の双対は逆順でさらに経路長4を引いた切断であり、
四つの連続頂点の間では常に負の長さで零となる。
次数零の順方向射は元の経路、対角は定数のみである。
四つの単純の線型Koszul分解を層化すると両端が一つずつの射影の完全列になり、
四連続の射影は両向きの全移動を生成する。有限生成加群は有限自由射影分解を持つため
四対象は充満例外列をなす。そのため局所化・加法性から周期巡回のEuler数は4となる。

もし特異点が標準の節点模型
$\operatorname{End}_R(R^a\oplus I^b)$ ($a,b>0$) を持つ $N$ 点のみならば、
同じEuler数は $-56+2N$ である。この計算は大域的な小解消の射影性に依存しない。
滑らかな開集合上はAzumayaなので周期巡回ホモロジーは通常のde Rhamコホモロジーの偶奇化と一致する
（行列traceのエタール降下と滑らかな場合のHochschild--Kostant--Rosenberg同型）。
閉台での局所化列により、これに各点の台つき圏のEuler数を加えればよい。
台つき有界圏は完備化で変わらない。各コホモロジーは有限長なのでこれは有限長加群への平坦基底変換と
有限項の射影解消によるExt比較から従う。局所大域次元は有限である
（元の次数付き正則環を可逆中心元で局所化して次数方向を除き、さらに完備化する）。
従って台つき完全圏を計算している。

標準節点の台つきEuler数は2である。通常のアフィン節点の小解消
$Y=\operatorname{Tot}_{\mathbb P^1}(\mathcal O(-1)^{\oplus2})$ 上で
$\mathcal O\oplus\mathcal O(1)$ の自己準同型環は $\operatorname{End}(R\oplus I)$ のアフィン版である。
実際零切断の外では二冪等元と階数一直線束の表示で一致し、
大域切断はその余次元2の補集合から延長する。高次自己Extは束射影で
$\mathcal O(n+j)$、$n\ge0,\ j=0,\pm1$ の第一コホモロジーに帰着して零で、
基底の例外対の引戻しはアフィン射ゆえ生成する。従ってこのtilting同値で節点台は
$\mathbb P^1$ 台となる。滑らかな $Y$ の開閉局所化によりそのEuler数は
$\chi(Y)-\chi(Y\setminus\mathbb P^1)=2$。
用いている加法性はdg完全圏の局所化に対する混合複体の完全三角を周期化したものであり、
ここでの各偶奇群は有限次元である。非固有な滑らかな複素多様体でもde RhamのEuler数は
Poincaré双対でコンパクト台Euler数と等しいため開閉差が使える。

滑らかな四次三次元多様体のEuler数は $-56$ である。
接束の全Chern類 $(1+h)^5/(1+4h)$ の三次項を $\int h^3=4$ で積分すればよい。
通常二重点のMilnor繊維は三次元球面と同じホモトピー型でEuler数が零なので、
$N$ 節点の補集合のEuler数も $-56$ である。
これで $4=-56+2N$、ゆえに $N=30$ となる。
既知の二十点は全て上の標準模型を持つため、それだけでは矛盾する。
追加点での単純次元は定理\ref{thm:twenty-central-nodes}の直後の系で決まる。
追加点が節点かつ二単純ならその系で非自由な階数一極大Cohen--Macaulay加群 $I$ の模型を持つ。
その最小正方表示は行列式の重複度2より二行二列で、節点の非退化二次主要部により
四成分の一次部は独立である。従ってこれは標準節点模型で、上の計数を適用できる。
\end{proof}

% !TeX root = ../papers/nc-p3.tex
\section{有限窓の頂点次元二の単純と二次曲面の伝播}
\label{sec:finite-two-dimensional-quadrics}
\begin{theorem}\label{thm:finite-two-dimensional-quadrics}
全四辺が有限な標準有限Frobenius窓の、各頂点次元2の折り畳み単純加群を考える。
同型と共通次数倍率を除いた分類は、各文字空間 $V_i$ 上の階数3または4の射影二次式 $q_i$ と、
階数4のときの一方の罫線族の指定による次の有限方程式の解に一致する。
二行二列線型表示 $A_i:V_i\to\mathrm{Hom}(U_i,W_i)$、$\det A_i=q_i$ をその指定に従って一つ取る。
各辺で $C_i:W_i\to U_{i+1}$ を未知とする線型写像
\[
 C_i\longmapsto\left(\sum_{u,v} r_{uv} A_{i+1}(v) C_i A_i(u)\right)_{r\in R_i}
\]
が非零核を持つことが条件である。核があれば次元1で可逆行列からなり（零元を除く）、
四箇所の間に追加の接続条件はない。この判定は未認証の窓と非被約な点グラフにも適用できる。
\end{theorem}
\begin{proof}
まず任意の単純で全頂点が非零なら、四矢に共通する核も、全像の張る真部分空間もない。
前者は頂点単純部分、後者は頂点単純商を与えるためである。
あるベクトルの次位置での像が直線なら矢の作用はその上で非零共変元 $p_i$ を通る。
六関係をこのベクトルに作用させると、次の直線は $R_i(p_i,-)$ で消える。
有限辺でこの評価の階数は少なくとも3なので、消す三平面が定まり、次の像も直線であり、
隣接点は有限グラフの一意な後続 $p_{i+1}$ になる。
四移動はスキーム上で恒等である。各位置で $\ker p_i$ の全矢に消されるベクトル空間を $K_i$ とすると、
$p_i$ で値1の矢は $K_i\to K_{i+1}$ に単射である。共通核の不在と関係評価から従う。
従って四段で同型となり、合成の固有直線を選べば四頂点一次元の部分加群を得る。
全頂点次元2の単純では矛盾する。ゆえに各 $0\ne u\in U_i$ に対し $A_i(V_i)u=W_i$ である。
逆にこの条件が全辺で成立すれば、任意の非零部分加群は次位置を全て含み、巡回して全体に一致する。

二行二列行列空間でこの条件は、像の直交空間が階数1の元を含まないことと同値である
（$w^*A_i(V_i)u=0$ を考える）。行列式二次式のため、直交空間は零か可逆元の張る直線だけである。
後者の像は左右基底変更で跡零超平面にでき、行列式はそこに非退化に制限される。
従って $\det A_i$ の階数はそれぞれ4または3である。
階数4では滑らかな二次曲面の二つの罫線を区別したSegre表示が、左右二次元空間の基底変更まで
行列表示を一意に復元する。階数3では根基を割って滑らかな平面二次曲線を取り、跡零三空間と同定する。
その二次曲線の全射影自己同型は二次元空間の共役作用で誘導されるので表示は一通りである。
射影倍率はいずれも左右変更に吸収できる。

接続方程式の核に階数1の $a b^*$ があれば、$u\in U_i$ と $e^*\in W_{i+1}^*$ を動かした
$b^* A_i(-)u, e^* A_{i+1}(-)a$ の全組が辺の点対応に入る。
両方とも二次元の変数から文字双対への単射である。
従って直線二本の積を含み、有限性に矛盾する。
核の二平面は行列式二次式の零点を持つので、可逆な解直線が唯一となる。
四接続を選べば関係は全て満たされ、上の像条件で単純である。
接続の四倍率は頂点基底の倍率で積だけになり、体が代数閉なので残りは次数倍率で吸収できる。
\end{proof}

\begin{proposition}\label{prop:two-dimensional-quadric-successors}
一つの向き付き $q_i$ だけを与えて次の二次曲面を列挙できる。
$A=A_i$ とし、
\[
 K=\ker\left[V_{i+1}^*\otimes W_i^*\longrightarrow R_i^*\otimes U_i^*\right],
 \qquad b\longmapsto\left(\sum r_{uv} b(v)A(u)\right)_r .
\]
$\dim K\le3$ であり、次位置が存在するのは $\dim K=2,3$ のときだけである。
次位置は $K$ の二平面を選び、それを行とする行列の行列式と罫線指定を取ったものに限り、全てがこの一辺を満たす。
次元2なら一意である。次元3なら次の二次曲面は、$\mathbb P(V_{i+1})$ 内の一つの滑らかな空間三次曲線を
通る二次曲面の網全体を一度ずつ動く。この手順を四辺で閉じることも単純の分類となる。
\end{proposition}
\begin{proof}
$K$ に非零可分解元 $p\otimes b^*$ があれば、前証明と同じく
$(b^*A(-)u,p)$ の直線が有限辺に入ってしまう。
$\mathbb P(W_i^*\otimes V_{i+1}^*)$ のSegreは余次元3なので $\dim K\le3$。
選んだ二行の全非零線型結合は行列階数2である。従ってこれらの行が作る次の矢の像は
全非零ベクトル上で全射で、前定理の階数3または4の二次曲面になる。
逆に後続の行は接続で $W_i$ 上に移せば必ずこの形である。

三行全部を束ねた $B:V_{i+1}\to\mathrm{Hom}(W_i,K^*)$ を考える。
どの $0\ne w\in W_i$ でも $B(-)w$ は三次元へ全射で、核は四次元文字空間の一直線である。
$B$ が単射でなければその共通核を除く三次元からの全射となり、$w$ に線型な三行三列の行列式が
射影直線上で消えないことになり矛盾する。よって核直線の射 $\mathbb P(W_i)\to\mathbb P(V_{i+1})$
の異なる二点は異なる像を持つ。微分も単射である。零なら一次変分でも同じ文字が
独立な二つの $w$ を消し、$B$ の単射性に反する。三次小行列は基点のない三次表示を与え、
従って像は滑らかな空間三次曲線である。$B(v)$ の三つの二次小行列はこの曲線のイデアルを生成する
（階数一条件が曲線をスキームとして定めることは各非零成分で核を解いてもよく、あるいは余次元2の
Hilbert--Burch表示を用いる）。これらは三独立で、二行の外積が三小行列の全一次結合を一意に与える。
\end{proof}

% !TeX root = ../papers/nc-p3.tex
\section{二次元単純の固有な極限と有限窓の分岐}
\label{sec:finite-two-dimensional-specialization}
\begin{theorem}\label{thm:finite-two-dimensional-specialization}
十五次元混合成分の閉包の標準有限Frobenius窓で四辺が有限とする。
各頂点二次元の単純が存在しないなら、長さ二十の点グラフに次のいずれかがある。
\begin{enumerate}
\item 非被約な局所点。
\item 異なる二点 $p,q$ で $\operatorname{rank} w(p_{i+3},-,q_{i+1},-)<2$。
\end{enumerate}
とくに全点が被約でこの異点階数が常に2なら、その閉包の\emph{全ての}窓が
定理\ref{thm:finite-two-dimensional-quadrics}の二次曲面表示を少なくとも一つ持つ。
閉点のAS性や中心の有限性は仮定しない。
より一般に、この結論は一般点に二次元単純がある任意の既約な有限窓の族に通る。
\end{theorem}
\begin{proof}
連続四頂点の大切断 $B$ の表現を、矢が $0,1,2,3$ の順に作用する向きで読む。
$\delta=(1,1,1,1)$ とし、重み $\theta=(1,3,9,-13)$ について非零真部分表現の重みが
負となることを安定、零以下を半安定とする。$0\le d\le2\delta$ で
$\theta d=0$ の可能性は $0,\delta,2\delta$ だけである。三進表示で最初の三成分が決まるためである。
四辺で前節の全射条件を満たす $2\delta$ 表現の切断は安定である。
非零部分が途中から始まっても一回作用させれば次位置は全体となる。
最後の位置の次元が2なら前三位置の真の欠損で重みは負となり、
最後の次元1の部分は最後にのみ支えられるからやはり負である。
安定な $\delta$ 表現は薄い点 $M(p)$ に限る。途中の矢が零なら正重みの初期区間が部分表現となり、
全て非零なら有限な点列である。

逆に切断で安定な $2\delta$ 表現の三辺では、各非零ベクトル上の矢の像が二次元となる。
正重みの各頂点には共通核がない。第一辺で一つの像が直線なら、有限辺の評価階数3以上により
その後も直線に入り、初期から末尾までの $\delta$ 部分または正重みの初期部分を生じる。
双対を逆に読んでも初期区間の重み和は真区間で正（元の終期区間の符号反転）なので、
最終辺でも双対の像は全非零ベクトル上で二次元である。
二行二列の場合この条件は元の像条件と同値である。すなわち共変元とベクトルの非零対が
全矢を同時に消すことを排除するだけである。第二辺で直線が生じれば、
有限辺の関係により最終辺でも直線以下になり矛盾する。
これで安定な切断の三辺のいずれも像条件を満たす。

第四辺を付ける全ての係数に対し、残る二窓の関係は斉次線型である。
安定な切断で非零解が一つでもあれば、それが与える折り畳み表現は単純である。
実際、第四辺であるベクトル上の像が直線なら第一辺へ伝播して矛盾する。
像が零のベクトルがあれば、第四辺の全矢は共通の一次元商を通る。
第四辺は非零なので双対側ではある共変元上の像が直線になり、逆向きに最終辺へ伝播して矛盾する。
従って四辺で全射条件が成立する。

一般窓からの離散付値環上の特殊化を取る。非巡回な四頂点箙に関係を課した
$2\delta$ の $\theta$ 半安定表現の相対GIT商は底に射影的である。
これは半不変式環のProjであり、次数零の不変式は有向閉路がないため底の関数だけである。
半安定縮約により有限な体拡大後、一般の安定表現は特殊繊維が半安定な表現族を持つ
（箙のGITの半安定置換、またはQuot表示による半安定格子の選択）。
関係式の閉条件は保存される。第四辺の解空間が非零という条件も行列階数の閉条件なので極限に残る。
従って特殊な切断が安定なら既に結論を得る。

安定でなければ安定因子は $M(p),M(q)$ の二つである。
$p=q$ を含め、両局所点が単純で異点の反対縮約階数が2なら、
定理\ref{thm:finite-cut-point-pairings}から二因子間・自己の第一Extは全て零である。
従って特殊表現は直和 $N=M(p)\oplus M(q)$ で $\operatorname{Ext}^1_B(N,N)=0$。
長さ二十の有限完全交叉はこの族で平坦で、単純な点はエタール局所で切断へ延びる。
直和の変形と選んだ半安定格子の変形は、完備化後、各次数で基底変更を持ち上げて同型になる。
実際隣り合う冪零厚みで二つの持上げの差は自己第一Extで測られ、零なら
前段の同型を修正して延長できる。行列の完備極限で同型を得るので一般点も直和となり安定性に反する。
一般点の存在は前節までの中心Euler数の結果による。
任意の閉包点への移行には稠密な存在開集合を一般点とする曲線の付値環を使えばよい。
\end{proof}

\begin{corollary}\label{cor:two-dimensional-extension-space}
上の安定な切断 $N$ に対する第四辺の線型解空間は自然に
$D\operatorname{Ext}^3_B(N,N)$ と同型である。従って追加の単純の存在と、
この安定表現が最高自己Extを持つことは同値である。
\end{corollary}
\begin{proof}
戻り辺を付ける表示は $T_B(\operatorname{Ext}^3_B(DB,B))$ なので、
解空間は $\operatorname{Hom}_B(N\otimes_B \operatorname{Ext}^3_B(DB,B),N)$ である。
大域次元3以下、命題\ref{prop:two-window-ext-complex}とSerre関手 $S$ の式より、
左の引数は $H^0(S^{-1}N[3])$ で、負次数以外に他のコホモロジーはない。
次数零加群への導来射はこの $H^0$ からの射と一致する。Serre双対で
$\operatorname{Hom}(S^{-1}N[3],N)=D\operatorname{Hom}(N,N[3])$ となる。
\end{proof}

% !TeX root = ../papers/nc-p3.tex
\section{十五次元成分の三十節点と二十個の二次元表現}
\label{sec:thirty-central-nodes}
\begin{theorem}\label{thm:thirty-central-nodes}
十五次元成分の一般の射影中心はちょうど三十個の通常二重点を持つ四次三次元多様体である。
二十点での単純の頂点次元は1と3、残り十点ではそれぞれ2と2である。
従って二次曲面を用いた定理\ref{thm:finite-two-dimensional-quadrics}の四辺の解は、
この一般窓につきちょうど二十個である。追加点での完備局所模型は
$R=k[[a,b,c,d]]/(ad-bc)$ と
$\operatorname{End}_R(R^{\oplus8}\oplus I^{\oplus8})$ であり、
$I$ は一般二行二列行列の余核である。
\end{theorem}
\begin{proof}
まず一般四次式の特異点でHessianの階数が4であることを、単元小行列を使う局所座標で示す。
$p=1009$ とし、節\ref{sec:unrestricted-doublet}の有理点 $t$ で、二つの隣接切断に
可逆六次小行列を取る。合計200個のSchur補成分の微分は法 $p$ で階数180を持つ。
独立な180行と180列を選ぶと、その180式は残り76変数についての陰関数を与える。
その陰関数上で残る二十式も恒等的に零である。
実際 $t$ の標数零のテンソル成分は滑らかな次元76であり、選んだ式で定まる
滑らかな同次元の局所空間の中にあるため完備局所イデアルが一致する。
$t$ のまわりで原点を移した形式陰関数は微分小行列が単元なので
$\mathbb Z_p$ 係数を持ち、自由変数の変化が $p\mathbb Z_p$ に入れば収束する。
このため有限精度の計算は、標数零の同じ成分上の実際の点の還元になる。
中心と経路空間の基底も同様である。単元小行列で取った基底についての
余分な方程式は標数零の開近傍で恒等的に成立し、積と核を与える消去も整係数で行える。

\path{computations/verify_thirty_central_nodes.py} に用いた入力と消去を全て記す。
固定した剰余ベクトルを先の微分の核へ射影し、$p$ 倍して $t$ に足した後、
選んだ180変数だけをHensel補正し全200式を法 $p^5$ で確認する。
経路長5まで順に単元階数部分を除き、四次中心の可換式の核に
元の五基底を持ち上げる。最初の四元の値を $(1,b,c,d)$ とし、
$b,c,d\ge0$, $b+c+d\le4$ の35組で第五元の左正則作用 $L_z$ を求める。
フィルター経路空間が安定した零番隅は階数64である。
四次の中心関係を $F$ とすれば特性多項式は $F^{16}$ である
（定理\ref{thm:twenty-central-nodes}のノルム恒等式）。
従って $\operatorname{tr}(L_z^j)/16$、$1\le j\le4$ とNewton公式は $F$ の係数を与える。
全次数4以下の三変数補間の行列式は単元であり、斉次化して五変数四次式を得る。

$S=k[a,b,c,d,z]$ と書き、五偏微分のイデアルを $J$、Hessianの全四次小行列の
イデアルを $H$ とする。次数9のMacaulay行列の単元ピボット数を順に示す。
各段で単元ブロックを消し、残りを $p$ で割って精度を一つ下げる。
\[
\begin{array}{c|rrrrr|r}
 &p^0&p^1&p^2&p^3&p^4&\dim S_9\\\hline
 J&637&0&48&0&0&715\\
 J+H&643&0&56&0&16&715
\end{array}
\]
これは元の行列の標数零での階数の下界を与える。各段の単元小行列は前段と独立な
非零ピボットを与えるからで、最後に精度を使い尽くした残りの階数を仮定しない。
従って $(J+H)_9=S_9$ はこの成分の稠密開集合で成立する。
特異点上ではEuler恒等式から動径がHessianの核に入るため階数はちょうど4であり、
アフィンチャートの二次主要部が非退化となる。ゆえに全特異点は通常二重点である。

次に節点 $R$ の閉点を除いた $U$ 上で任意のAzumaya類が零であることを示す。
節点の一方の小解消 $Y\to\operatorname{Spec}R$ は正則で、
$Y\setminus U=\mathbb P^1$ は余次元2である。これは二行二列行列の核直線を付ける
射影的な解消を完備局所環へ引き戻したもので、標準二チャートで正則性も分かる。
Brauer群の純性により $U$ 上の類は $Y$ に延びる。
$R$ は狭義Hensel局所環なので、固有基底変換で
$H^2_{\mathrm{et}}(Y,\mu_n)=H^2_{\mathrm{et}}(\mathbb P^1,\mu_n)=\mathbb Z/n$
となる。最後の群は核直線の射影直線の $\mathcal O(1)$ の類で生成される。
従ってKummer列から $\operatorname{Br}(Y)[n]=0$ が全 $n>0$ で成立する。
以上より完備中心の分数体 $K$ 上の一般行列環は分裂する。

その次数は16である。一般の列加群 $K^{16}$ 内に完備代数で安定な有限格子を取る
（基底の $R$ 張りを代数で生成すればよい）。局所大域次元の有限性により有限生成
射影分解を持つ。分数体へ移して列加群の重複度を数えれば、1は原始射影の
一般行列階数の整数線型結合である。従って追加点で単純が次元8の一種類しかない
可能性はない。その場合には全原始射影の一般行列階数が2だからである。
定理\ref{thm:twenty-central-nodes}の系により二単純となり、
非自由な階数1の極大Cohen--Macaulay加群による局所模型を得る。
その二行二列表示の行列式の二次主要部が非退化なので四成分はパラメータ系である。
以上で命題\ref{prop:extra-central-singular-characters}の計数の仮定が全て成立し、
$4=-56+2N$ より $N=30$、追加の単純は $2(30-20)=20$ 個となる。
\end{proof}

\begin{corollary}\label{cor:quartic-center-unramified-brauer}
一般の射影四次中心の関数体には指数4の非自明な不分岐Brauer類がある。
従ってこの三次元多様体は安定有理でない。
\end{corollary}
\begin{proof}
次数方向を除く中心の関数体で、滑らかな部分の零番隅のAzumaya類を取る。
四次中心の一次数元を可逆化すると次数方向は一つの独立な不定元になる。
その拡大先で零番隅の一般商は既知の次数4の可除環なので、拡大前も指数4である。
節点全部を吹き上げた滑らかな射影解消上で、例外因子の一般点での剰余も零となる。
実際、各節点で局所環を完備化した後の同じ吹上げでは、先の分裂により零である。
吹上げは平坦基底変換と両立し、例外因子はどちらでも同じ滑らかな接錐二次曲面で、
そのイデアルは極大イデアルの引き戻しである。よって該当する付値の分岐指数は1で
剰余体も同じであり、剰余の関手性から元の剰余は零である。
滑らかな解消上の純性から不分岐性を得る。代数閉体上の滑らかな射影多様体の
Brauer群は安定双有理不変で、射影空間では零である。
\end{proof}

% !TeX root = ../papers/nc-p3.tex
\section{反対階数八の分岐を中央二次元の中心で検出する}
\label{sec:branch-middle-defect-bound}
\begin{theorem}\label{thm:branch-middle-defect-bound}
曲面射影のない標準正則窓に片側分岐があるとする。反対階数8ならどちらかの向きに
中央支持2の分岐中心があり、その三テンソルは中央鉛筆 $L_2$ 型である。
反対階数9でこのような中心がないなら、全順・逆中心の中央支持は3であり、
各後続直線（逆向きでは前点直線）の一般点の中央支持は4である。
ここでは純関係を許す。
\end{theorem}
\begin{proof}
中心 $p$ の中央支持を $m$ とし、順向きとして位置0に置く。
後続直線 $L$ の一般点 $x$ の両端支持は3なので中央支持 $h$ は3または4である。
$h=3$ なら低階数分岐の証明で用いた反対点の構成を階数によらず使えることを確認する。
直線上で $w(x)=\sum_{j=1}^3 K_j(x)\otimes A_j$ と書けば、
一般の $K_j(x)$ は独立なので中央支持は一定であり、その消去点 $y$ は全 $x\in L$ と
$w(x,y)=0$ を満たす。よって $w(y)$ の中央支持は $L$ の消去二空間に含まれる。
中央1は連続純関係を与えるからちょうど2である。一般 $x$ の両端階数3と
二つの隣接六切断の直交を使うと $y$ の両端階数はどちらも3以下である。
補題\ref{lem:middle-binary-pencil}と端の下界2よりそれらは $2,3$ または $3,2$ であり、
$y$ は所望の中心である。

そこで $h=4$ とする。分岐列の各辺では中央支持の和が $\rho$ 未満である
（定理\ref{thm:low-rank-branch-exclusion}の最初の直交補の議論）。
以下 $\rho=m+5$ のときも、$L$ 上に中央支持2の中心があることを示す。
反対切断の $W_p^\perp$ はこのとき五次元であり、
直線 $x(t)$ の各中支持を、位置3の四文字を入力とする $5\times4$ の一次行列 $H(t)$ で表す。
一般階数は4である。四次小行列式の共通零因子の次数を $\delta$ とすると、
その行像の消去直線を表す原始な五つの同次式の次数は $4-\delta$。
分岐列の次の曲線 $x_2(t)$ の次数は $d_2=3-f_1$ である。
$w(-,x_2(t))$ の全ての行を $W_p$ を法として読むと、この消去直線に値を持つ
$d_2$ 次の多項式列である。全列が零なら全 $x\in L$ と全 $x_2(t)$ が隣接し、
曲線の積が辺に含まれて可動中心の排除に反する。従って原始列の次数は $d_2$ 以下で
\[
                           \delta\ge f_1+1 .
\]

他方で $H$ の階数落ちの長さを局所的に上から評価する。
先の $K(t)$ は三列の $4\times3$ 行列で、三次小行列式の局所共通零次数を $n$ とすると
全点での和は $f_1$ である。$n>0$ は順中心で、そこでの階数は2なので
局所DVRでの単因子は $1,1,t^n$（今は $t$ を一様化変数とする）。
$H(t)$ は、$A_j$ が与える定数の四次元部分空間に $K(t)\otimes 1$ を制限したものである。
実際 $A_j$ による写像は一般の中央支持4より単射で、
$W_p^\perp$ への定数包含は局所単因子を変えない。
この制限の各単因子の指数は $n$ 以下である。飽和像の元に $t^n$ を掛ければ
$K(t)\otimes1$ の全像に戻り、その逆像は一般体上で元の定数部分空間に属するので
DVR上でもその部分空間に戻る。ゆえに $\delta$ の各点での寄与は
$(4-\operatorname{rank}H(t))n$ 以下となる。
全ての順中心が中央支持3以上なら和は $f_1$ 以下となり、先の下界に反する。
従って $L$ 上に中央支持2の順中心がある。

中央支持2の中心がないと仮定すると、$m\ge3,h=4$ かつ $m+4<\rho$。
階数8なら $m=3$、階数9で $m=4$ の場合も上の $\rho=m+5$ に入り矛盾する。
残る階数9の全中心は $m=3$ であり、逆向きにも同じである。
\end{proof}

% !TeX root = ../papers/nc-p3.tex
\section{中央二次元の中心の対と曲線列の内部での中央支持}
\label{sec:branch-binary-center-pairing}
\begin{theorem}\label{thm:branch-binary-center-pairing}
曲面射影のない標準正則窓で、中央支持2の順中心 $p$（位置0）には
位置3の中央支持2の逆中心 $z$ が一意に対応し、逆も成り立つ。
$p$ の順・中央支持を $S\subset V_1,M\subset V_2$ とすると、$z$ の中央・逆支持も
$S,M$ である。$p$ の後続直線と $z$ の前点直線の一般点は中央支持3である。
逆に分岐中心間の有限曲線列で中央支持が一般に3の動く成分は、
このような対の最初または最後の直線に限る。その他の動く一般点は全て中央支持4である。
反対階数の制限はない。
\end{theorem}
\begin{proof}[二次元の中心からの逆中心の構成]
$w(p)$ の三支持は順に $S,M,C$、次元は $2,2,3$ である。
最後の縮約像 $U\subset R_1\cap(S\otimes M)$ は三次元で、
補題\ref{lem:middle-binary-pencil}の $L_2$ により二項対称二乗の平面となる。
位置2の $x\in N=\mathbf P(M^\perp)$ を一般に取る。$R_1$ を末尾で $x$ に評価すると
$U$ が消えるので逆支持は3以下、動く中心がないので3である。従って一意な前点 $v$ がある。
$v$ が動けば逆支持は3以上で、一定なら $N$ を後続に持ち逆支持3である。
隣接六空間での直交により $x$ の順支持は3以下、再び一般点でちょうど3である。
$x$ の一意な次点を $z$ とする。$z$ が動けば同じ直交により順支持3、
一定なら逆中心なので順支持3である。
$w(z)$ の末尾は $\ker x$ に支持される。その先頭で縮約した三次元像は
$R_1\cap(V_1\otimes\ker x)$ に入る。この評価核は次元3なのでちょうど $U$ に等しい。
従って $w(z)$ の中央・逆支持は $S,M$、特に逆支持2である。
逆中心は有限個なので $z$ は一定、$N$ がその前点直線全体となる。
順逆を入れ替えても全く同じ構成で、$S^\perp$ の射影直線の一意な前点が元の $p$ である。
$w(x,z)=0$ が $p$ の後続直線 $x$ 上で成り立つ（$z$ の中央支持は $S$）。
従って一般点の中央支持は3以下、両端3のため3である。$N$ 上も同じである。
\end{proof}

\begin{proof}[二つの動く反対点]
分岐列の動く一般点 $x$ を位置0とし、その中央支持が3なら唯一の反対消去点を
$y$（位置2）とする。$w(x,y)=0$ であり、両隣の六空間の直交と
$x$ の両端支持3から $y$ の両端支持は3以下である。
$y$ が一定ならその中央支持は動く $x$ の包を消すので2以下となる。中央1は連続純関係を生じ、
中央2では両端は異なる次元で $2,3$ か $3,2$ である。従って $y$ は上の対の中心、
$x$ はその中央支持の消去直線上にあり、一般点の一意な前点または次点が一定となる。

$y$ が動くと仮定する。その両端は一般に3で、$w(x,y)=0$ と中央二項鉛筆の補題から中央も3となる。
$x^-,y^-$ を $x,y$ のそれぞれの前点とする。$w(x)$ の先頭を縮約した三次元関係像は
$R_2\cap(\ker y\otimes V_3)$ であり、その二支持は $\ker y,\ker x^-$ である。
実際包含は三支持から従い、$y$ の順評価の階数3で等号となる。
$w(y^-)$ の先頭も $\ker y$ に入る。その末尾支持は3である。$y^-$ が動けば
逆支持は3以上で、$y$ の順支持3との直交により3以下。定点なら順中心で逆支持3である。
従ってその末尾縮約像が上と同じ三次元関係空間を張る。
ゆえに $y^-$ の順支持は3で、$y^-$ は動く（一定なら順中心となる）。
その中央支持は $\ker x^-$、従って $w(y^-,x^-)=0$ である。
$x,y$ を交換すれば $x^-$ も動き、両者の全三支持が3となる。
この操作を反復すると最初の $x$ の前点を何回取っても動く。
しかし分岐列を逆にたどれば元の順中心という定点に達する。従ってこの場合は起こらない。
\end{proof}

\begin{corollary}\label{cor:rank-nine-nonlinear-corridors}
反対階数9の曲線分岐の各中心の中央支持は2または3である。
どの列でも非直線の連続する成分の一区間内に、射影三空間全体を張る成分は高々一本である。
平面に入る各整成分の次数は3以下である。
\end{corollary}
\begin{proof}
中央支持が3以上の中心では最初の直線の一般中央支持は前定理より4である。
定理\ref{thm:branch-rank-one-residual}より中心の中央支持は $9-6=3$ 以下となる。
連続する非直線成分 $C_i,C_{i+1}$ の一般点の中央支持はともに4である。
同じパラメータの直前・直後の点 $x_{i-1},x_{i+2}$ で縮約した関係は階数1以下である。
実際 $W_{i-1}\subset W_i^\perp$, $W_{i+2}\subset W_{i+1}^\perp$ の間の
対合は直交する $W_i,W_{i+1}$ の一方をそれぞれ割った余次元 $9-4-4=1$ の対合を通る。
この関係は一般に非零である。零なら $x_{i+2}$ から $x_{i-1}$ へ辺ができ、四点を巡回できる。
両成分が非直線なので列の内部にあり、$x_{i+2}$ の一般点での一意な後続をたどると
同じ動く成分を再訪するか、最終直線の一定後続と矛盾する。
従って純関係がある。その一つを固定し、全パラメータで評価すれば、二因子の少なくとも一つは
対応する成分全体を消す。非直線を消せるときはその平面の式に一意である。
連続する純関係の両方に同じ文字は使えないので、非直線成分 $k$ 本の区間の
$k-1$ 辺には異なる $k-1$ 本の平面成分が必要である。

最後に次数4以上の整平面成分を考える。その平面上で順側の四次小行列式は全て
曲線の式で割れる。恒等零なら平面全体が射影像に入り、次数4の場合に限られる。
この場合全小行列式は一つの四次式の定数倍で、一般の評価での
六関係の二次元核は一定である。従って平面上で恒等的に消える非零関係があり、
その平面の式を始文字とする純関係となる。逆側からも同じ平面の式を末文字とする
純関係ができ、連続純関係の禁止に反する。ここは反対階数の制限を使わない。
\end{proof}

% !TeX root = ../papers/nc-p3.tex
\section{純関係を許す反対階数九の曲線分岐の排除}
\label{sec:rank-nine-branch-exclusion}
\begin{lemma}\label{lem:rank-nine-ternary-center-conic}
反対階数9の曲面射影のない標準正則窓で、中央支持3の順中心 $p$ の列の
第二曲線は必ず平面二次曲線である。その平面の式 $b$ を末文字とする純関係がある。
位置を逆にすれば、逆中心の第二曲線の平面の式は純関係の始文字である。
\end{lemma}
\begin{proof}
列を $p,x_1,x_2,x_3,\ldots$ とする。二直線列の下界11により $x_3$ は動く。
$x_2$ の次数は $3-f_1$ で、次数1なら次の次数が $-f_2$ となるから2または3である。
次数2では滑らかな平面二次、次数3では最初の直線上に順側欠損がなく、
$3\mathcal O(-1)\subset4\mathcal O$ の商から完全三次系となる。
その結果 $x_1,x_2$ の一般中央支持はともに4である
（定理\ref{thm:branch-binary-center-pairing}）。反対切断の一余次元の対合により
$w(p,-,-,x_3(t))$ は階数1以下である。一般に零ではない。実際 $w(p)$ の末尾は
三次元支持なのでこれを消す $x_3$ はただ一点である。

$w(p)$ を順に $S\otimes B\otimes C$ の $2\times3\times3$ テンソルとし、
末尾の縮約による三空間を $J\subset S\otimes B$ とする。$S$ 方向での
三行三列鉛筆の非零パラメータに階数1はない。もし値が $b\otimes c$ であれば、
他方の値を $B/kb$ に射影して消す $C^*$ の非零元を取れる。これで末尾を縮約すると
$S\otimes kb$ に非零の元を得る（末尾支持は全 $C$）。二つの連続純関係の中央文字が
$b$ となるので不可能である。階数0は先頭の完全支持に反する。

$J$ 中の上の動く階数1行列は、射影的に一定でないとしてもその $B$ 因子は一定である。
実際 $S$ 因子を固定して $B$ 因子が動けば、$S$ 因子を消す鉛筆の値は
独立な二つの $C^*$ 方向を消すので先の階数下界に反する。両因子が動く場合、
$B$ 因子が二空間を張れば、両因子の積座標は少なくとも三次元を張るため
$J$ 全体がその二空間に支持され矛盾する。$B$ 因子が三空間を張れば積座標は
四次元以上で次元に反する。ここで関数空間の積の次元下界は、一つの滑らかな点で
各空間の基底を異なる位数に取り、二つの位数集合の和集合（対ごとの和）を数えればよい。
従って $w(p,-,-,x_3(t))=a(t)\otimes b$ と書ける。
点対 $x_2(t),x_3(t)$ の条件から $b(x_2(t))=0$ で、動く第二像は非退化三次ではなく
二次曲線に一致する。一般の $t$ を固定すれば主張の純関係が得られる。
\end{proof}

\begin{theorem}\label{thm:rank-nine-branch-exclusion}
反対階数9の標準正則窓で全射影像が高々一次元なら、どの辺のどちらの射影にも
正次元のファイバーはない。従って反対階数9以下の全窓で、全次数の正則性と
添字保持同型が既出の有限判定で決まる。
\end{theorem}
\begin{proof}[二次元中心の一つ前の位置]
順中心をその中央支持によって二次元または三次元と呼ぶ。
系\ref{cor:rank-nine-nonlinear-corridors}によりこの二種類で尽きる。
二次元中心 $q$ の次位置・反対位置に、後続直線 $L_q$ と中央支持の消去直線 $N_q$ を取る。
後者は定理\ref{thm:branch-binary-center-pairing}の対をなす逆中心の前点直線である。
二項対称二乗の三関係により、$L_q$ に始まらない次の辺の点対は $N_q$ に終わる。
同じ位置の別の順中心の後続直線は $L_q$ と交わらず、その第二曲線は $N_q$ に入る。
これは二直線列の排除に反する。よって二次元中心がある位置の順中心はその一点だけである。

一つ前の位置に順中心 $p$ があれば $q\in L_p$ である。$p$ が三次元なら前補題の
$f_1=1$ と $q$ の位置での一意性から従う。$p$ が二次元なら、$q\notin L_p$ のとき
$L_q=N_p$ が必要になる。しかし $N_p$ の各点はその対をなす逆中心一つのみを後続とするので
この直線を初回とする列は動く位置一つで終わり矛盾する。
従って $p$ の第二曲線は二次曲線である（$f_1>0$ かつ二直線列はない）。
この曲線の一般点は $L_q$ の外にあり、第三曲線は $N_q$ となる。次数公式より
$f_2=3\cdot2-3-1=2$ で、二次曲線上のある順中心 $r$ の後続直線が $N_q$ と交わる。

そのような交点はない。$r$ が二次元なら $q$ の中央二空間の全文字に終わる純関係と
$r$ の順二空間の全文字に始まる純関係が連続するため、両二空間は横断し
$L_r\cap N_q=\varnothing$ である。$r$ が三次元なら、$L_r$ の一般点の中央支持は4であり、
順中心でない全点でも4である。実際直線上の表示 $w(x)=\sum K_j(x)\otimes A_j$ の
三列はそこで独立で、固定テンソルの中央支持を保つ。一方 $N_q$ のどの点も後続はただ一つで
順中心ではなく、反対の $q$ に消されるので中央支持は3以下である。
これで二次元順中心の一つ前の位置には順中心がないことが示された。
\end{proof}

\begin{proof}[三次元中心の排除]
三次元順中心が一つあれば前補題より次位置にも順中心があり、前段よりこれは三次元である。
四位置を一周すれば全位置に三次元順中心があり、二次元中心の位置での一意性から
二次元順中心は存在しない。順逆の二次元中心の対より逆中心も全て三次元である。
逆中心は少なくとも一つ存在し、逆の向きにも同じ二次曲線の議論で全四位置に存在する。

同じ位置0にある順中心 $p$ と逆中心 $z$ を取る。$p$ の第二曲線 $C_2$ の平面の式 $b$
は $R_1$ の純関係の末文字である。$z$ の三次縮約を逆の順で
$S'\otimes B'\otimes C'$（位置3,2,1）に置く。$S'$ の二次元鉛筆が正則鉛筆なら、
その全非零値は階数2以上である。前点直線 $L_z$ の外の位置3の点でのこの鉛筆は
一般には可逆で、そのとき対応する位置2の点は $B'$ の消去点のみである。
特異なパラメータは有限通りで、そこでも位置2の点はそれぞれ一本の固定直線に入る。
この制約は辺 $2\to3$ に含まれる $z$ の三縮約関係から従う。
従って $C_2$ からの第三曲線は $L_z$ に一致する。
その直線は $z$ に終わる列の末端であり、逆にたどれば $C_2$ は $z$ の第二曲線である。
逆向きの前補題より $R_2$ に同じ文字 $b$ に始まる純関係が生じる。

鉛筆が恒等的に特異なら、両支持が3で定数左右核がなく、一般階数は2である。
Kronecker分解の左右の特異ブロックの最小指数はともに正であり、
三行三列に収まるのは各一つの指数1だけである。従って逆向きの基底で
\[
                \begin{pmatrix}u&v&0\\0&0&u\\0&0&v\end{pmatrix}
\]
とできる。末尾での縮約の三関係には二次元の $S'\otimes k b'$ が入り、
元の向きではいずれも文字 $b'$ に始まる $R_2$ の純関係である。
第三曲線が $L_z$ の外を動けば、これらから $C_2$ は $b'=0$ の平面に入り
$b'$ は $b$ に比例する。$L_z$ 自身なら先と同じ逆向きの前補題を使う。
いずれも同じ $b$ を介する連続純関係になり、三次元中心は排除される。
\end{proof}

\begin{proof}[二次元中心だけの場合]
全順中心が二次元なら全逆中心も二次元である。実際逆に三次元中心が一つあれば
前段を逆に適用して矛盾する。順中心 $p$ を位置0に置き、一つ前の位置だけでなく
次位置にも順中心がない（次の中心が二次元ならその一つ前に $p$ がある）。
従って $d_2=3$ で非退化三次となる。位置2に順中心がなければ $d_3=6$ である。
平面成分の次数は3以下なので非平面的な非直線が二本連続し、
系\ref{cor:rank-nine-nonlinear-corridors}に反する。
よって順中心は位置0,2にだけ存在し、逆中心は二次元中心の対により位置1,3にだけ存在する。
位置2の順中心を $a$ とする。第三曲線は $L_a$ ではない。一般前点が動くためである。
第四位置には逆中心がなく、従って第四曲線は動いて $N_a$ に一致し、次がその終中心となる。
逆向きでも終中心の隣の位相には逆中心がないから、逆向きの第二曲線、すなわち元の第三曲線は
非退化三次である。第二曲線と連続して同じ矛盾を与える。

分岐が消えたので、系\ref{cor:all-windows-through-eight}の被覆に
系\ref{cor:all-one-dimensional-point-graphs}と定理\ref{thm:finite-rank-nine-all-marks}を適用して
階数9全体の有限判定を得る。
\end{proof}

% !TeX root = ../papers/nc-p3.tex
\section{階数十の四次元分岐中心に続くねじれ三次曲線}
\label{sec:rank-ten-cubic-after-center}
\begin{theorem}\label{thm:rank-ten-cubic-after-center}
反対階数10の標準正則窓で曲面射影がなく、中央支持4の順中心 $p=x_0$ に始まる
列の第二曲線 $x_2$ がねじれ三次曲線とする。第三曲線の次数は4以下であり、第二曲線上に順中心がある。
第三曲線が非平面的なら第四曲線も動き、$p$ を通る平面に入り、
\[
             2\le d_3-1\le d_4\le3
\]
である。さらに第三曲線は滑らかな二次曲面の二つの母線族に対する射影次数が
$1,d_3-1$ の有理曲線である。逆方向でも成立する。
特に中央支持4の中心をどちらかの端に持つ次数列 $(1,3,3,1)$ で、
内側二本がともにねじれ三次曲線であるものはない。
\end{theorem}
\begin{proof}[二次の二項行列と一次の関係]
第一直線と第二曲線のパラメータを $s,t$ とする。$t=(t_0,t_1)$ とし
\[
              w(-,x_1(s),x_2(t),-)=[s,t]A(t).
\]
位置0の行支持は $U=\ker p$、位置3の列の変数空間は四次元 $V_3^*$ とする。
$A(t)$ の全列が張る $D\subset U\otimes k[t_0,t_1]_2$ は四次元で、係数の $U$ 支持は全三次元である。
実際、第二曲線が位置2を張るから、第一直線の一般点における中央支持と逆支持の全てがここに現れる。
逆支持は3、中央支持は4である。後者は定理\ref{thm:branch-binary-center-pairing}より、
この直線の中央支持が3ならその一意な前点の中心が中央支持2になることによる。
$w$ の反対切断で位置1を第一直線に制限するとその像は $k[t_0,t_1]_1D$ である。
これは反対側の $p$ での四次元像を消すので次元6以下である。

$S_j=k[t_0,t_1]_j$ とし $D_-=\{L\in U\otimes S_1:S_1L\subset D\}$、
$D_{--}=\{u\in U:S_1u\subset D_-\}$ と置く。
$t_0 f=t_1 g$ の関係を割ると $\dim S_1D=8-\dim D_-$、
$\dim S_1D_-=2\dim D_- -\dim D_{--}$ を得る。
前者から $m=\dim D_-\ge2$、後者から $2m-\dim D_{--}\le4$ である。
$2\dim D_{--}\le m$ もあるので $m=2$ である。
$D_{--}=0$ なら二つの一次列 $L_1,L_2$ に対し $D=S_1L_1+S_1L_2$、
$D_{--}=ku$ なら $D=S_2 u+kB(t)$ である。
両場合とも一般の値 $P(t)=A(t)V_3^*$ は二次元である。
最初の場合、二つの一次列が常に比例するなら、独立性からともに固定ベクトルの一次式倍となり
$D_{--}=0$ に反する。第二の場合は $U$ の完全支持による。

$R_0$ の第二位置を第一直線に制限する写像を考える。像は
$([-,t]\otimes A(t)z)_{t,z}$ が張る空間である。実際 $R_0$ は反対辺の二変数の
全縮約の像であり、$x_2(t)$ が全空間を張る。制限写像には非零核がある。
さもなくば $z\in\ker A(t)$ の全てが $w(-,-,x_2(t),z)=0$ となり、
動く一般点 $x_2(t)$ が順中心になってしまう。従って六次元の $S_1\otimes U$ における
像の消去式があり、一次行 $\ell(t)$ が全 $P(t)$ を消す。
$U$ の完全支持により定数の消去式はないので $\ell$ の二係数は独立である。
一般の $P(t)$ は $\ker\ell(t)$ に等しく、その全てが固定直線 $ku$ を含む。
一定基底を取り $P(t)=\langle u,v(t)\rangle$、$v(t)=t_0 v_0+t_1 v_1$ と書けば
$u,v_0,v_1$ は独立である。

これで次の二つの形で尽きる。$D_{--}=0$ なら、$L_i$ の $ku$ を法とする像は
それぞれ $v(t)$ の定数倍である。基底を取り直し一方を $n(t)u$、他方を $v(t)$ とできる
（後者の $u$ 方向は $v_0,v_1$ に吸収する）。$D_{--}\ne0$ の場合は $S_2u$ の
$u$ がまさに全 $P(t)$ に共通する直線を張り、$B(t)$ の剰余は $n(t)v(t)$、
$u$ 方向は他の三列で消せる。$n$ は両方で非零一次式である。位置3の列基底の選択により
\[
 \begin{array}{ll}
 \mathrm{I}: & A(t)=(t_0 n(t)u,\ t_1 n(t)u,\ t_0 v(t),\ t_1 v(t)),\\
 \mathrm{II}:& A(t)=(t_0^2 u,\ t_0t_1 u,\ t_1^2u,\ n(t)v(t)).
 \end{array}
\]
一般の $P(t)$ の消去行は上記の原始一次行の多項式倍なので、
第一直線への $R_0$ の制限の余次元はちょうど1、核も一次元である。
$n(t)=0$ では $A(t)$ がちょうど階数1になり、$w(-,-,x_2(t),-)$ の像の制限も一次元。
従って元の像は二次元以下である。二次元の下界より、ここが第二曲線上の順中心となる。
またそこでの $A(t)$ 自身は $R_3$ の純関係である。
\end{proof}
\begin{proof}[後続の曲線の次数]
第三点 $x_3(t)$ は $A(t)$ の核に入る。第三点が一定なら二直線列以下となり、
定理\ref{thm:two-line-full-middle-restriction}に反する。
II型なら $x_3$ は最後の座標の零平面に入り、平面曲線の次数上界3を使える
（系\ref{cor:rank-nine-nonlinear-corridors}の末尾の階数によらない証明）。
I型では、一般の核は二行をそれぞれ $(t_0,t_1)$ が消す二行二列行列であるから
\[
         x_3(t)=(-t_1 f(t),t_0 f(t),-t_1 g(t),t_0 g(t)).
\]
ここで $f,g$ は同次数で基点のない二項式で、次数は $d_3-1$。
射影直線から二次曲面の各因子への射として読むことで、特殊な $t$ にもこの表示が延びる。
第三曲線が平面なら既に十分である。非平面なら $f,g$ は独立であり、
$n=t_0$ とできる。第四点 $x_4(t)$ での三つの評価を $u(x_4),v_0(x_4),v_1(x_4)$ とする。
$x_3(t),x_4(t)$ は辺なので、第二位置には独立なパラメータ $y$ を代入できる。
$A(y)x_3(t)$ を $x_4(t)$ が消す式から $[t,y]$ を除くと
\[
 y_0\big(f(t)u(x_4)+g(t)v_0(x_4)\big)+y_1 g(t)v_1(x_4)=0
\]
を得る。従って $x_4$ は $v_1=0$ の平面に入り、$p$ からの射影は
$[u(x_4):v_0(x_4)]=[g:-f]$ である。
$x_4$ が恒等的に $p$ となることは、順中心 $p$ の前点が一意なので不可能。
射影が非定数なので第四曲線は動き、射影で次数は増えないから所望の不等式となる。
\end{proof}

% !TeX root = ../papers/nc-p3.tex
\section{階数十の四次元中心に続く二次曲線と特異三項鉛筆}
\label{sec:rank-ten-conic-after-center}
\begin{theorem}\label{thm:rank-ten-conic-after-center}
反対階数10、曲面射影のない標準正則窓で、中央支持4の順中心 $p$ の第二曲線が
平面二次曲線であるとする。第一直線の唯一の順中心 $q$ は中央支持3であり、
その二項鉛筆は $L_1\oplus L_1^t$ 型である。第二曲線の次の第三曲線は平面に入り、
その次の位置も $p$ を通る平面に入る。第三曲線が三次なら平面三次の特異点は
$q$ の中央支持の消去点であり、そこに逆中心がある。
逆向きにも成立する。
\end{theorem}
\begin{proof}
前節と同じ $U=\ker p$ を用いる。第一直線と第二曲線の縮約を
$[s,t]A(t)$ とすると今度は $A=t_0A_0+t_1A_1$ が一次である。
第二位置の二次曲線の平面を法線 $h$ で表す。
定理\ref{thm:branch-rank-one-residual}の二次曲線の証明の通り、
第一直線上の三つの関係列を平面へ制限すると二項二次系の二つの標準一次関係と零列になる。
零列の元の値は $m(s)h$ で、$m$ は非零一次式である。
従って第一直線の唯一の順中心は $m=0$ の点 $q$ で、そこで他の二列は平面に制限しても独立である。
$q$ でのテンソルは、これら二つの独立な先頭文字に対応する $A_0,A_1$ の鉛筆に左右同値となる。
実際標準二列を第二曲線に評価すると $[s,t]t_0,[s,t]t_1$ であり、
列に対応する定数係数行列がそのまま $A_0,A_1$ である。
両行列を合わせた $U$ 支持は3（$q$ の前点は $p$ 一つ）、
位置3の支持次元を $k$ とするとこれは $q$ の中央支持である。

$D=A(t)V_3^*\subset U\otimes S_1$ は次元 $k$。
前節と同じ反対切断の制限から $\dim S_1D\le6$ である（今回は第二位置でも平面に制限する）。
もし $k=4$ なら $\dim\{u\in U:S_1u\subset D\}=8-\dim S_1D\ge2$ なので
$D=S_1U'$、$\dim U'=2$ となり $U$ 支持3に反する。
$k=2$ なら中央二項鉛筆の補題から、先頭と中央の二位置の三関係は二項対称二乗の平面である。
すなわち $A(t)$ の有効な三行二列を
$\left(\begin{smallmatrix}t_0&0\\ t_1&t_0\\0&t_1\end{smallmatrix}\right)$ とできる。
その左核行は $(t_1^2,-t_0t_1,t_0^2)$ で、非零の一次関係を持たない。
従って $([-,t]\otimes A(t)z)_{t,z}$ は六次元全体を張り、$R_0$ の第一直線への制限は単射となる。
しかし各 $\ker A(t)$ は二次元以上であり、単射ならその全点が $x_2(t)$ の後続となる。
曲線全体が順中心となって矛盾する。

従って $k=3$。一般の $t$ で $W_{x_2(t)}\subset W_{x_1(t)}^\perp$、
右辺を $W_p$ で割った次元が $10-4-4=2$ である
（第一直線の一般中央支持は前節と同じく4）。
$W_p$ は第一直線で零なので、別の $s$ で評価しても $A(t)$ は階数2以下となる。
有効な三行三列鉛筆は定数の左右核を持たず、従ってKronecker分解は
正の最小指数を持つ左右の特異ブロック各一つ、すなわち $L_1\oplus L_1^t$ である。
適合基底で
\[
 A(t)=u(t_0\alpha_0+t_1\alpha_1)^t+(t_0v_0+t_1v_1)b^t
\]
と書ける。$u,v_0,v_1$ は $U$ の基底、$\alpha_0,\alpha_1,b$ は位置3の独立な文字である。
$A(t)x_3(t)=0$ より $b(x_3)=0$ であり、$x_3$ は平面に入る。
$z$ を三文字の共通消去点とすると、$x_3$ は動くのでその射影は
$[\alpha_0(x_3):\alpha_1(x_3)]=[-t_1:t_0]$ である。
隣接する $x_3(t),x_4(t)$ を固定して全ての $y$ で
$x_4(t)^t A(y)x_3(t)=0$ を使うと $u(x_4(t))=0$ となる。
平面は $p$ を含む。

第三曲線が三次なら $z$ からの射影次数1なので $z$ は重複度2の点である。
そこで逆評価が階数3なら、逆辺の三次小行列の非零な一つを使い、近傍で前点が
$x_3$ の正則関数として定まる。正規化パラメータでの合成は滑らかな二次曲線
$x_2$ への同型である。従って $z$ の二つの逆像は同じ点へ写せず、
一つの逆像であっても正規化の微分が零ではあり得ない。
平面三次の二重点では二枝または一枝で接写像が零なので矛盾する。
逆評価はちょうど階数2である。
\end{proof}

\begin{proposition}\label{prop:singular-ternary-center-projections}
反対階数によらず、中央支持3の順中心 $q$ が $L_1\oplus L_1^t$ 型なら、
その次の位置のさらに次の位置と、その後の位置における二つの平面射影が
一点ずつからの射影として同定される。ただし共通の消去平面に入らない点対での主張である。
特に $q$ の列の第二曲線がねじれ三次曲線なら第三次数は2以上4以下である。
その三次曲線が $q$ の中央支持の消去点を通るときは上界が3になる。
\end{proposition}
\begin{proof}
$q$ の位置を1として上の鉛筆を読む。ここで $u,v_0,v_1$ は $q$ の末尾の三支持であり、
その全消去点は $q$ の一意な前点である。位置3,4の辺の二関係は
$u(x_4)\alpha_j(x_3)+v_j(x_4)b(x_3)=0$ ($j=0,1$)。
従って $b(x_3)\ne0$ で $x_4$ が全消去点でなければ
\[
 [u(x_4):v_0(x_4):v_1(x_4)]
       =[b(x_3):-\alpha_0(x_3):-\alpha_1(x_3)].
\]
$q$ 自身の第二曲線が全空間を張れば右の像は直線に入らない。
第三曲線が動けば左の射影も一般に定義され、元の曲線は直線ではあり得ない。
第二曲線がねじれ三次ならその直後に一定の逆中心へ入ることもない。
このとき右辺の射の引戻し次数は、射影の中心が三次曲線上にあれば2、なければ3である。
左辺の射影の中心 $q^-$ は $q$ の一意な前点で、その第三曲線による逆像長は高々1である。
実際相異なる二点なら、それぞれ $q$ から第三点 $q^-$ までの経路に辺 $q^-\to q$ を加え、
独立に反復できる二つの閉経路ができる。一点に長さ2があれば同じ経路の双対数族があり、
最初の直線で非零の微分を持つ。各反復で独立な双対数を取ると、各ブロックの一次文字の値が
その二次元環を張るので、長さ $4n$ の経路の評価次元が $2^n$ 以上になる。
いずれも標準成長に反する（命題\ref{prop:terminal-branch-contact}の帰路の議論）。
従って左辺の射の引戻し次数に足す基点因子の長さは高々1であり、表示した上下界が従う。
\end{proof}

% !TeX root = ../papers/nc-p3.tex
\section{二項中心の二次錐と第二次数の強制}
\label{sec:binary-centers-quadric-cones}
\begin{theorem}\label{thm:binary-centers-quadric-cones}
曲面射影のない標準正則窓で二直線列がないとする。中央支持2の順中心 $p$ があれば、
その一つ前の位相に順中心はなく、その位相には $p$ しかない。
$p$ の第二曲線は必ず滑らかな平面二次曲線であり、最初の直線上のただ一つの順中心は
中央支持3である。第三曲線は平面二次またはねじれ三次で、第二曲線は
$p$ の中央支持の消去直線に長さ1で会う。逆向きにも成り立つ。
四位相を合わせても中央支持2の順中心は高々一つである。
これらは純関係や反対階数への仮定を要せず、特に反対階数10に適用できる。
\end{theorem}
\begin{proof}[錐と三次曲線の後の下界]
$p$ を位置0とする。定理\ref{thm:branch-binary-center-pairing}の対となる逆中心 $z$ は
位置3にあり、$L=L_p$ は位置1の後続直線、$N$ は位置2の中央支持の消去直線である。
$N$ の各点の後続は $z$ 一つである。$p$ の一意な前点を $p^-$ とする。
二項対称二乗の表示で、中央の二座標を $a,b$、位置3の末尾三座標を $c_0,c_1,c_2$ とすれば、
$R_2$ の二関係は $a c_0+b c_1, a c_1+b c_2$ である。従って位置2,3の点対で
位置2が $N$ の外、位置3が $p^-$ の外なら
\[
             [c_0:c_1:c_2]=[b^2:-ab:a^2].
\]
位置3の動く成分は全て頂点 $p^-$ のこの二次錐上にある。
実際前点が $N$ にあれば後続は $z$ 一点にすぎない。逆向きの対に同じことを使えば
位置0の動く成分も $z$ の一意な次点を頂点とする二次錐にある。
ここで用いる成分は全て列の成分で、前点と次点の存在が保証されている。
一つの位相で同じ二次錐上にある列の異なる成分の次数総和は8以下である。
射影評価の四次小行列式の少なくとも一つはこの錐上で非零であり、
その零点の因子が全成分を含むからである。

第二曲線 $x_2$ がねじれ三次曲線の場合を考える。その $N$ からの射影の次数を
$k=3-n$、$n=\operatorname{length}x_2^{-1}(N)$ とすれば $k>0$ で、
$x_3$ の $p^-$ からの射影は上述の二次写像との合成になる。
第三曲線は動き、その次数は $d_3=2k+m$、$m=\operatorname{length}x_3^{-1}(p^-)\le1$ である。
最後の不等式は $p^-\to p$ の固定帰路を付ける前節の閉経路の証明による。
$d_3\le3$ なら $k=1$、$n=2$ である。$x_3$ が二次なら滑らか、三次なら非平面的である
（二次錐の平面切断には平面三次曲線がない）。従って後者はねじれ三次で、どちらも閉埋込みである。
一方 $x_2^{-1}(N)$ 上で $x_3=z$ がスキームとして成立する。
実際 $N$ 上の順評価は階数3で、$z$ は全 $N$ で解である。局所の単元な三次小行列で解けば
任意の微小部分上でも核直線はこの定点となる。長さ2が一点に写るのは閉埋込みに反する。
従って $d_3\ge4$ であり、特に二項中心側からは $(1,3,1)$ の列は出ない。
\end{proof}
\begin{proof}[一つ前の位相]
$R_1$ の三関係は二項対称二乗なので、位置1が $L$ の外の点対は位置2が $N$ に入る。
従って位置1の順中心は全て $L$ 上である。外にあればその後続直線全体が $N$ で、
ただ一つの動く位置の後に $z$ へ至ってしまう。位置0に他の順中心があればその後続直線は $L$ と交わらず、
第二曲線が $N$ となり二直線列の排除に反する。
位置3に順中心があると仮定する。その第二曲線は非直線（第一直線の後も直線なら
次数式で二直線列となる）なので $L$ ではない。その第三曲線は $N$ であり、
次は $z$ なので次数列は $(1,d,1)$、$d=2,3$ である。
逆に $z$ から始めて前段を使うと $d=3$ は不可能である。

$d=2$ なら次数式により第二曲線での順側欠損は $3\cdot2-3-1=2$。
その上のある順中心 $r$ の後続直線は $N$ に会う。一方 $r\in L$ だから $w(r,z)=0$、
従ってその中央支持は2または3である。2の場合は、$p$ の中央支持の各文字に終わる純関係と、
$r$ の順支持の各文字に始まる純関係により、両二空間は交われない。
従ってその消去直線も交わらない。3の場合、$L_r$ の一般中央支持は4で、
順中心でない全点でも4である。実際そこでは定理\ref{thm:finite-branch-corridors}の
三列が独立で中央支持が落ちない。二項中心の対の定理から中央支持3の最初の直線は二項中心から出るものに限る。
$N$ 上の全点は順中心でなく、反対位置の $p$ に消され中央支持が高々3なので、やはり交われない。
これで位置3には順中心がない。
\end{proof}
\begin{proof}[二次錐の次数と第二曲線]
再び $d_2=3$ と仮定する。位置1の全順中心は $L$ 上にあり、その上の欠損 $f_1=0$ なので存在しない。
位置2の順中心の後続直線は全て位置3の二次錐上に入る。全母線は頂点を通り、
異なる中心の後続直線は交わらないため、この位相の順中心は高々一つである。
一つもなければ $f_2=0$ より $d_3=6$ となり、$f_3=0$ から
$d_4=3d_3-3d_2+1+b_2\ge10$ となる。第三曲線は非直線なので第四曲線も動き、
位置0の二次錐の上界8に反する。

位置2の順中心を $a$ とすれば、その後続直線に順中心はないのでその第二曲線は
位置0を張る次数3の曲線である。$p$ の第四曲線はこれと異なる
（異なる列が同じ動く成分を使えない）。従って $d_4\le8-3=5$。
$d_4=3d_3-8+b_2$ と $d_3\ge4$ から $d_3=4$、$d_4\ge4$ である。
$d_3=2k+m$ より $m=0$、$k=2$、$n=1$。
第三曲線は $p^-$ を避けるため第四曲線には $p$ がなく、位置0の順中心の一意性から $f_4=0$。
一方 $n=1$ から第三曲線は逆中心 $z$ を通り $b_3\ge1$ である。
従って $d_5=3d_4-3d_3+d_2+b_3-f_4\ge4$ となる。
この第五曲線は位置1の $L$ と異なり、非直線なので第六曲線も動いて $N$ に一致する。
次は $z$ であり、逆向きに見ると第五曲線は第二曲線となる。しかし第二曲線の次数は3以下である。
これで $d_2=3$ を排除した。

従って $d_2=2$ で、$f_1=1$ より第一直線に順中心が一つある。
それは $z$ に消され中央支持3以下であり、中央支持2なら前段の一位相前の禁止に反する。
二次曲線から $N$ への交叉長を再び $n$ とする。
$N$ からの射影が一定ならこの曲線の平面が $N$ を含み、$n=2$ である。
この場合第三曲線の射影も一定なので第三曲線は直線（または定点だが二直線列の排除から動く）であり、
先の長さ2の矛盾が成立する。従って射影次数 $k=2-n$ は正である。
$d_3=3-f_2\le3$ と $d_3=2k+m$ より $k=1,n=1$ であり、先と同じ錐の平面切断によって
第三曲線は平面二次またはねじれ三次である。
もう一つ二項順中心があれば相対する位置2に限られるが、その前位相に上記の三項順中心があるため矛盾する。
\end{proof}

\begin{corollary}\label{cor:no-binary-centers-without-two-lines}
曲面射影がなく、どの列も二直線列でない標準正則窓では、順中心・逆中心の中央支持は
全て3または4である。全ての動く成分の一般点の中央支持は4である。
特に反対階数10の曲線分岐に中央支持2の候補はない。
\end{corollary}
\begin{proof}
上の二項順中心 $p$ とその対の逆中心 $z$ があるとする。定理とその逆により
位置3に順中心はなく、位置0に逆中心はない。
位置2に順中心 $a$ があればその最初の直線に順中心はないので、
第二曲線は位置0のねじれ三次曲線である。これは $z$ の後続点 $z^+$ を頂点とする
二次錐に入る。錐上のねじれ三次曲線は必ず頂点を通る。
避けるなら頂点からの射影の三座標は基点のない三次式であり、滑らかな底二次曲線を
経由する射の引戻し次数が奇数となるためである。
ゆえに $z^+$ のある前点が $a$ の後続直線上にある。位置0に逆中心がなく
$z\to z^+$ なのでその前点は $z$ 自身である。しかし最初の後続直線の全点の前点は一意であり、
逆中心 $z$ はその上に乗れない。

これで順中心は位置0の $p$ と位置1の $q$ のみである。後者の一意性は
全て $L$ 上にあり $f_1=1$ であることによる。逆中心は列の対応の全単射により二つで、
$z$ と、上の定理を逆に使って得る位置2の逆中心 $b$ で尽きる。
$p$ の列では $f_2=f_3=0$、従って $d_3=3$, $d_4=4+b_2\ge4$。
非直線の次も動くので第五曲線がある。第五曲線は位置1に戻るが、最初の $L$ を再訪できない。
次が定点なら位置2の逆中心 $b$ であり、逆順の第二次数が $d_4\ge4$ になって矛盾する。
次が動けば $L$ の外から来るので第六曲線は $N$ で、その次が $z$ となる。
この場合には二項中心 $z$ から逆順の第三次数が $d_4\ge4$ となるが、上の定理では3以下である。
よって二項中心は存在せず、動く点の主張は定理\ref{thm:branch-binary-center-pairing}から従う。
\end{proof}

% !TeX root = ../papers/nc-p3.tex
\section{二次行列鉛筆の核直線と次の分岐中心}
\label{sec:rank-ten-flat-kernel-lines}
\begin{theorem}\label{thm:rank-ten-flat-kernel-lines}
定理\ref{thm:rank-ten-cubic-after-center}の四次元順中心 $p$ と第二ねじれ三次曲線の上に、
その後続直線が一般核直線の極限に一致し、自身の第二曲線が平面二次曲線である順中心 $r$ が存在する。
$p$ の第三曲線は直線または平面二次曲線であり、二次曲線なら第四曲線は直線である。
II型の制限では第三曲線は必ず直線である。
従って反対階数10で曲面射影のない窓のどの四次元順中心からも、
高々三辺で中央支持3の順中心に達する。逆向きでも成立する。
\end{theorem}
\begin{proof}[特殊点での核の極限]
同定理の両標準形で $n=t_0$ とし $t_*=(0,1)$ とする。一般核の極限 $K(t)$ は、
I型では二つの二座標組をそれぞれ $(t_0,t_1)$ が消す二平面、
II型では第四座標が零で残る三座標を $(t_0^2,t_0t_1,t_1^2)$ が消す二平面である。
いずれも全 $t$ で部分束を与える。$E(t):V_3^*\to R_0$ を $x_2(t)$ での順評価とする。
第一直線に $R_0$ を制限した後の $E(t)$ は $[-,t]\otimes A(t)$ であり、
制限 $R_0\to R_0|_{L}$ の核は一次元であった。$t\ne t_*$ で順中心があれば
$E(t)$ は階数2なのでその核は $K(t)$ そのものである。

そうでなければ順側欠損 $f_2=6-d_3\ge2$ は全て $t_*$ にある。
この点の局所一様化変数を $v$ とし、局所自由な基底で $E(v)$ を読む。
制限後の二つの有効な行は $L_1(v),vL_2(v)$ と書け、
$L_1,L_2$ は全点で独立な行で $K=\ker(L_1,L_2)$ となる。
実際両標準形の $u,v(t)$ の方向は全点で独立、二行の一方のみに $n$ が掛かる。
したがって単元な主元で消去し、一次元核側の行 $C(v)$ を戻すと、
三次小行列のイデアルは $v$ と、$K$ の局所基底における $C|_K$ の二成分のイデアルとの積である。
これは $L_1,L_2$ を座標行に取れば直ちに従う。位数が2以上のため $C(0)$ は $K(0)$ を消し、
$E(0)$ の階数は下界により2。従ってここでも $\ker E(0)=K(0)$ である。
こうしていずれの場合にも所望の核直線を持つ順中心 $r$ が得られる。
\end{proof}
\begin{proof}[二つの列の平面切断]
I型では $K(t)$ の点を
$z=(-t_1X,t_0X,-t_1Y,t_0Y)$ とでき、$y$ を独立に動かすと
\[
              A(y)z=[t,y]\big(y_0 X u+Y(y_0v_0+y_1v_1)\big).
\]
従って $K(t)$ の一般点のどの後続も $H=\{v_1=0\}$ に入る。
II型では $K(t)$ の第四座標は零で、$A(y)z$ は恒等零でない $u$ の倍となるから
全ての後続が $H=\{u=0\}$ に入る。特に $r$ の第二曲線は $H$ にある。
第二次数は2または3で、3ならねじれ三次となるから平面二次曲線である。

一つの平面内にある異なる動く成分の次数和は3以下である。
系\ref{cor:rank-nine-nonlinear-corridors}の末尾の単一成分の証明で、
その平面内の異なる曲線の式の積を使えばよい。総次数4以上なら順・逆の全四次小行列が
その積で割れ、全て恒等零または定数比例となる。前者は曲面射影、
後者は同じ平面の法線に始まり終わる連続純関係を与える。
$I$ 型で $p$ の第三曲線が非平面的なら、同定理の射影式より第四曲線は $H$ の中の
次数2以上の別の曲線であり、$r$ の二次曲線との次数和に反する。
従って第三曲線は滑らかな二次曲面の平面切断にあり、直線または平面二次である。
二次なら二次曲面の両射影が非定数なので同じ式から第四曲線は $H$ 内で動き、
次数和により直線となる。

II型では $p$ の第三曲線も $r$ の後続直線も位置3の第四座標の零平面にある。
後者と前者の正規化の交叉長は高々1である。
実際 $r$ の後続直線の上では逆評価が全点で階数3で、その前点は定数 $r$ である。
この直線に入る任意の微小因子で $p$ の第二点はよって $r$ に等しく、
第二曲線のパラメータ表示が閉埋込みなのでその因子長は高々1である。
二つの列の成分は異なるから、平面の直線を制限した因子の次数は第三曲線の次数に等しい。
従ってこの次数は1である。

最後に $p$ の第二曲線が最初から二次なら定理\ref{thm:rank-ten-conic-after-center}で
一辺先に三次元順中心がある。ねじれ三次なら $r$ が三次元であるか、
四次元で自身の第二曲線が二次であり、その後続直線に三次元順中心を持つ。
二項中心の排除と二直線列の排除でこれらの場合が尽きる。
\end{proof}

% !TeX root = ../papers/nc-p3.tex
\section{核直線族の非被約因子と四次元中心の三次分岐の排除}
\label{sec:rank-ten-cubic-kernel-divisor}
\begin{theorem}\label{thm:rank-ten-cubic-kernel-divisor}
反対階数10で曲面射影のない標準正則窓の中央支持4の順中心の第二曲線は、
必ず滑らかな平面二次曲線である。第一直線上の唯一の順中心は特異三項鉛筆型である。
第三曲線および次の位置は定理\ref{thm:rank-ten-conic-after-center}の二平面に入る。
逆向きにも成立する。
\end{theorem}
\begin{proof}[核直線上での次辺の必要条件]
第二曲線がねじれ三次であると仮定する。定理\ref{thm:rank-ten-cubic-after-center}の
I, II型のいずれでも、極限を含む核直線を $t=(s,t')$ と $[X:Y]$ の積でパラメータ化する。
混同を避けて以下 $t=(t_0,t_1)$ とし、位置3の座標をそれぞれ
\[
 z_{\rm I}=(-t_1X,t_0X,-t_1Y,t_0Y),\qquad
 z_{\rm II}=(-t_1X,t_0X-t_1Y,t_0Y,0)
\]
とする。いずれも基点のない $(1,1)$ 式であり、$A(t)z=0$。
したがって $E(t)z=w(-,-,x_2(t),z)\in R_0$ は第一直線への制限写像の
固定一次元核に値を取る。その係数は一つの $(4,1)$ 式 $D$ である。
恒等零なら全 $x_2(t)$ が順中心になり不可能である。
$D=0$ は二点 $x_2(t),z$ が実際に辺をなす条件そのもので、各固定 $t$ で $X,Y$ に関して一次である。
従って因子は垂直なファイバー（重複を含む）と、第三点のパラメータを与える重複度1の切断からなる。

この因子の各既約成分の一般点で、その重複度まで含めて $z$ に後続点がある。
実際積曲面の当該素因子での離散付値環を $D$ で割る。水平成分なら一般の $x_2(t)$ は
逆支持3、垂直成分なら $x_2(t)$ は順中心なのでやはり逆支持3である。
よってこの局所商環上でも逆評価の三列が六空間に階数3の直和因子を張る。
$w(-,x_2(t),z,-)=0$（巡回順での二点の辺の式）により、隣接切断の完全対合で
$z$ の順評価像はその直交補の階数3の直和因子に入る。
剰余体上で $z$ は動く一般点であり順支持3なので、単元な三次小行列で解けば
この局所商環で核直線が存在する。これが求める後続 $x_4$ である。
\end{proof}
\begin{proof}[次辺の小行列式の二重次数]
$R_3$ の中に三つの独立な行列 $A(y)$ の係数があり、残る三関係を $B_j$ とする。
I型で $Y$ が単元なら、全 $A(y)z$ との対合が消えるための $x_4$ の条件は、
前節の式より $v_1(x_4)=0$, $Xu(x_4)+Yv_0(x_4)=0$ である。
従って $x_4$ は $p$ と $[u:v_0:v_1]=[Y:-X:0]$ の固定線型持上げ $h(X,Y)$ の包に入る。
この二列はどの $[X:Y]$ でも独立である。残る三関係のこの包への制限行列は
$B_j(z,p)$ と $B_j(z,h)$ を二列に持つので、二次小行列式の二重次数は $(2,3)$ である。
前段の後続の存在により $Y=0$ の外で各小行列式は $D$ で割り切れる
（各素因子の重複度まで消える）。
$D$ の中の $Y$ 因子を全て取り除いても $t$ 次数は4なので、$t$ 次数2の小行列式は
全て恒等零である。こうして全ての一般の $z_{\rm I}$ が後続を持つ。
この $z$ の像は滑らかな二次曲面全体であり、曲面射影の不在に反する。

II型では $A(y)z$ は $u$ の二次式倍で、その全係数の一つは各 $z$ で単元である。
従って $x_4$ は平面 $u=0$ に入る。$B_j$ のその平面への制限を三行三列にすると、
行列式は $(3,3)$ 式である。これも $D$ の全素因子で重複度まで消えるので恒等零となる。
$z_{\rm II}$ の像は第四座標が零の全平面である。実際任意の二項二次式が根 $t$ を持ち、
その $t$ で表示の直線上に乗る。この全平面が後続を持つことになり矛盾する。
以上で第二次数3を排除した。二直線列はなく、残る二次曲線には既出の判定を適用できる。
\end{proof}

% !TeX root = ../papers/nc-p3.tex
\section{三項中心の第二三次曲線と三次スクロールの因子}
\label{sec:ternary-center-cubic-scroll}
\begin{theorem}\label{thm:ternary-center-cubic-scroll}
曲面射影と二直線列のない標準正則窓で、中心 $p$ の第二曲線がねじれ三次曲線なら、
$w(-,x_1(s),x_2(t),-)=[s,t]A(t)$ の二項二次行列は一般階数3である。
反対階数10では、この場合は中央支持3の正則三行三列鉛筆を持つ順中心に限られ、
第三曲線は直線である。第二曲線上の順側欠損は二次列二本と一次列一本の五次行列式で与えられる。
特に特異三項鉛筆の中心の第二曲線は二次である。逆順でも成立する。
\end{theorem}
\begin{proof}[左一次関係のある四空間]
$U=\ker p$ と置く。二項中心の排除により第一直線の一般中央支持は4、
逆支持は3である。従って列の空間 $D=A(t)V_3^*\subset U\otimes S_2$ は四次元で $U$ の全体を使う。
一般階数1なら二項多項式環の一意分解で多項式列と行の積となるが、三次元と四次元を
張る両因子には少なくとも次数2と3が必要である。そこで一般階数2と仮定する。

$R_0$ を第一直線に制限する写像は単射ではない。単射なら一般の $x_2(t)$ の後続に
$\ker A(t)$ の全直線が入る。前節までと同じく制限像は $[-,t]\otimes A(t)z$ の包であり、
一般階数2、完全な $U$ 支持から、その消去式は原始な一次行一つである。
よって制限写像の核は一次元。適合基底を取ると全列は
\[
          A(t)z=uQ_z(t)+v(t)L_z(t),\qquad
          v(t)=t_0v_0+t_1v_1
\]
となり、$u,v_0,v_1$ は独立、$Q_z$ は二次、$L_z$ は一次である。
$D$ は $S_2\oplus S_1$ 内の超平面である。$L$ の像が一次元なら
$D=S_2u+k n(t)v(t)$ で先のII型となるので排除される
（定理\ref{thm:rank-ten-cubic-kernel-divisor}の因子の証明は両標準形の下で反対階数を使わない）。
像が二次元なら超平面の式は $\lambda(Q)+\mu(L)=0$、$\lambda\ne0$ である。
$\lambda$ を $S_1$ 上の対称二形式 $\lambda(lm)$ と見る。
階数1で $\mu$ がこの形式の像に入る場合、$v(t)$ に一次式倍の $u$ を加えて $\mu$ を消せる。
$\ker\lambda$ の二次式は全て固定一次式の倍であり、先のI型となってやはり排除される。
\end{proof}
\begin{proof}[入射曲面上の辺と次の小行列式]
残りは $\lambda$ の階数2、または階数1で $\mu$ がその像に入らない場合である。
$A(t)z=0$ の解を二項変数 $y$ を用いて
\[
      Q_z(y)=[t,y]M(y),\quad L_z(y)=[t,y]b,\qquad M=M_0 y_0+M_1y_1
\]
と書く。$[M_0:M_1:b]\in\mathbb P^2$ と $t\in\mathbb P^1$ の間の
$\lambda([t,y]M)+b\mu([t,y])=0$ は階数2の双一次式である。
実際第一項の $M$ での係数行の写像は $\lambda$ と同階数で、階数1の場合には $b$ の項が独立な一次式を与える。
この曲面 $T$ は一点 $o$ での $\mathbb P^2$ の吹上げである。
二つの独立一次式を消す $o$ の外でその比が $t$ を定め、入射式がこのグラフの閉包を与える。
$F,H$ をそれぞれ $\mathbb P^1,\mathbb P^2$ の超平面類の引戻し、
$J$ を例外曲線とすれば $J=H-F$, $J^2=-1$, $F\cdot J=1$ である。
$z$ の座標は基点のない $H+F$ の切断であり、$b\ne0$ では $L_z$ の根が $t$ を復元するため
その像は曲面である。

前節の $E(t)z$ は同じく位置2と3の辺を与える固定核直線に値を取り、
$T$ 上では類 $H+4F$ の有効因子 $\Delta$ を切る。恒等零は可動中心の禁止に反する。
各ファイバーで一次なので、その水平部分は第三点を与える重複度1の切断である。
各素因子の一般点でその重複度まで次点 $x_4$ が取れることは、前節で使った逆支持3と
直交性による局所商環の議論の通りである。

$p$ からの射影の座標 $(u,v_0,v_1)$ において $h=(b,-M_0,-M_1)$ と置き、位置0への線型持上げも $h$ と書く。
$A(y)z$ との対合は全て零なので $M_j u(x_4)+b v_j(x_4)=0$ ($j=0,1$)。
従って $b$ と $(u(x_4),v_0(x_4),v_1(x_4))\wedge h$ の積は零である。
$R_3$ の $A(y)$ の三係数は独立である。実際 $v_0,v_1$ の方向での係数には
二つの独立な $L_z$ と全ての一次式の積が現れる。残る三式 $B_j$ について
$B_j(z,p),B_j(z,h)$ を二列とする。小行列の類は $3H+2F$ である。
$\Delta$ の素因子の重複度を $n$、$b$ のそこでの位数を $e$ とすれば、その小行列は
少なくとも $n-e$ の位数を持つ。実際離散付値環の商で上記の外積は一様化変数の
$n-e$ 乗を法に消える。$h$ の一成分は単元なので、この商で $x_4$ は $p,h$ の包に入る。
従って二列の間に単元成分を持つ関係がある。
\end{proof}
\begin{proof}[例外曲線での余分な直線の排除]
階数2の $\lambda$ では $b=0$ の直線は $o$ を通らないので、その引戻しは既約な水平曲線である。
この曲線が $\Delta$ の水平切断になることはない。その上で $z$ の像は
$L_z=0$, $Q_z\in\ker\lambda$ の射影直線で、パラメータ $t$ を根に持つ唯一の二次式が $Q_z$ となる。
この直線への射は次数2である。二式の鉛筆 $\ker\lambda$ が共通根を持てば
$\lambda$ はその点の評価の倍で階数1になってしまうからである。
これは全ての動く曲線のパラメータが像へ双有理であることに反する。

従ってこの場合は小行列から $\Delta$ 全体を除けて、残る三成分は類 $2H-2F=2J$ の切断である。
$H^0(\mathcal O(kJ))$ は $k\ge0$ で一次元である（$J^2=-1$ により順に $J$ を除く）。
よって全小行列が零か一定ベクトル倍になる。前者なら一般の $z$ で $p,h$ の包に
後続が存在して曲面射影が生じる。後者なら第一列 $(B_j(z,p))_j$ が固定平面に入る。
しかし $p$ の逆支持は3で $A(y)$ は $p$ を消し、$T$ の $z$ 像は全 $\mathbb P^3$ を張る。
最後の事実は任意の一般の $[t,y](M,b)$ の包、あるいは次の直接計算で分かる。
線型式が $z$ の像で消えればその双一次引戻しは入射式の倍である。
$S_2^*\oplus S_1^*$ から $(t,M,b)$ の双一次式へのこの代入は単射なので
元の式は $D$ の定義式の倍である。従って後者も不可能である。

階数1の残りの場合は $b=0$ が $o$ を通り、その引戻しは被約和 $J+F_e$ となる。
$F_e$ は $\lambda([t,y]M)$ が全 $M$ で零になる唯一のファイバーである。
$J$ と $F_e$ の $z$ 像は同じ直線 $Z=\{L_z=0\}$ であり、どちらの制限も同型である。
実際 $F_e$ 上では $[t,y]$ を固定して $M$ が動き、$J$ 上では $M$ が固定され $[t,y]$ が動く。
$\lambda$ が一つの点での評価なので、どちらもその固定一次因子を持つ二次式全体になる。
$F_e$ が $\Delta$ の成分でなければ、$\Delta$ から除く必要があるのは高々 $J$ 一つで、
小行列の残余類は $2J$ または $3J$。既に排除した定数比例の矛盾となる。
$F_e$ が成分なら、$\Delta$ から $F_e$ を一回と、必要なら $J$ を一回差し引いた因子を除ける。
残余の類は $kJ+F$、$k=2,3$ のいずれかである。$J$ への次数 $1-k<0$ のため
全残余成分が $J$ 上で消える。従って $J$ 上の一般点で $z$ は $p$ と固定した $h(o)$ の結ぶ直線に後続を持つ。
そこでは $b=0$ だが、$p,h(o)$ はいずれも $A(y)z$ を消すのでこの結論に影響しない。
他方、$F_e$ が辺の因子なので $Z$ は順中心 $x_2(e)$ の後続直線そのものである。
一般点の後続は一意なのでこの中心の第二曲線が直線となり、二直線列の禁止に反する。
これで一般階数2を全て排除した。
\end{proof}
\begin{proof}[反対階数十での五次欠損]
中央支持が4なら前定理が第二三次を排除するので、現在は3である。
$\dim S_1D\le10-3=7$ より、四次元 $D$ はある非零一次列 $L$ の
$t_0 L,t_1 L$ を含む（一次関係を割る前節までの計算）。
残る二列を $B,C$ とすれば $A(t)=(t_0L,t_1L,B,C)$。
一般階数3のため核直線は $(-t_1,t_0,0,0)$ であり、第三曲線は直線となる。
この場合 $[-,t]\otimes A(t)z$ は六次元を張るので $R_0$ の第一直線への制限は同型である。
よって $x_2$ の順側欠損は $A(t)$ の三次小行列の共通因子、すなわち非零五次式 $\det(L,B,C)$ が与える。
特異三項鉛筆なら命題\ref{prop:singular-ternary-center-projections}の下界2に反する。
\end{proof}

\begin{corollary}\label{cor:rank-ten-uniform-corridor-types}
反対階数10の曲面射影のない標準正則窓で分岐があれば、全ての列の第二曲線が二次であるか、
全てが次数列 $(1,3,1)$ で両中心とも正則三項鉛筆を持つかのいずれかである。
後者では各位相に順中心・逆中心はそれぞれ高々一つであり、その存在位相は両向きで一致し、
二位置移動で保たれる。前者では順中心も逆中心も全四位相に存在する。
\end{corollary}
\begin{proof}
任意の標準窓の分岐列で、直線からねじれ三次曲線へ写る辺を考える。
その直線の二項座標と次位置の四座標との双一次式で辺上消えるものは、
ねじれ三次の三つの標準一次関係の包に入る。実際直線と三次のパラメータは
同じ正規化で、三次は完全系を使う。この辺の関係の直線への制限は一般点の順評価が
階数3なのでちょうどこの三空間である。直線の一般点 $x$ の縮約を
$\sum K_j(x)\otimes B_j$ と書くと、三つの $K_j(x)$ は独立である。
その逆支持は $B_j$ の固定支持の和に等しく、従って一意な一般前点は一定である。
よってこの直線は列の最初であり、逆にねじれ三次に直線が続けば最後の直線である。
上の定理より第二三次に続くのは直線なので列は $(1,3,1)$ で、逆中心にも同じ定理を用いる。

第二二次の場合は最初の直線の唯一の順中心の鉛筆を読む。その二項行列 $A(t)$ と
この中心の鉛筆の一致は定理\ref{thm:rank-ten-conic-after-center}の最初の標準二列の証明による。
中央支持3でこの三行三列が正則なら $\ker A(t)$ は一般に一定の第四方向だけであるから
第三点が一定になり矛盾する。従ってこの中心は特異三項または四次元であり、既にどちらも
第二曲線が二次と分かっている。これを反復すれば全位相に第二二次の順中心がある。
一つの列で最後から第二曲線が三次なら逆向きの結果からその列自身の次数列が $(1,3,1)$ となる。従って第二二次だけの場合、逆向きでも同じ反復ができる。

一方、$(1,3,1)$ の順中心 $p$ を位置0に取ると同位相に他の順中心はない。
実際その後続直線は $L_p$ と交わらない。$p$ の正則三項鉛筆をその直線上で評価すれば
二項パラメータは全 $\mathbb P^1$ を動くので、一般に三行列は可逆であり、
第二点は $p$ の中央支持の消去点に固定されてしまう。
よって二次の中心鎖とは共存できない。三次の各列の両中心は同位相にあり、
第二三次上での順側欠損は5なので反対位相にも順中心がある。逆向きも同じである。
\end{proof}

% !TeX root = ../papers/nc-p3.tex
\section{五次欠損の排除と階数十の二次曲線列}
\label{sec:rank-ten-quintic-centers}
\begin{theorem}\label{thm:rank-ten-quintic-centers}
反対階数10の曲面射影のない標準正則窓では、次数列 $(1,3,1)$ ばかりの分岐配置は存在しない。
従って分岐があれば全列の第二曲線は二次である。
\end{theorem}
\begin{proof}[二つの反対消去点]
この三次列の始終の中心を位置0の $p,z$ とする。もし $p=z$ なら最初の直線の異なる点から
同じ中心に戻る閉経路を反復できるので $p\ne z$。位置1,3の最初・最後の直線を $L,T$ とする。
$p$ の一意な前点 $p^-$ は $T$ の外にある。実際 $T$ の各点の後続は $z$ 一つである。
同様に $z$ の一意な後続 $z^+$ は $L$ の外にある。
位置2のねじれ三次上の順中心を $a$、逆中心を $b$ とする。前節より各々一つで、
両向きの中心が一致することはない。パラメータの対応する零点の一次式を $l_a,l_b$ とする。

$p$ の中央支持の消去点は $b$ であり、$z$ のそれは $a$ である。これを示す。
$p$ の正則三項鉛筆から、位置3の点を $p^-$ からの平面射影で読み、二つの行を縮約する。
$T$ の平面射影は直線であり、その一般点でこの二行は独立である。
実際正則三行三列鉛筆の可逆元を単位行列にすれば、従属方向はもう一元の左または右固有方向である。
固有空間が二次元なら非零な鉛筆元が階数1以下となってしまう。
これは定理\ref{thm:rank-nine-branch-exclusion}で用いた三項中心の連続純関係による排除で、
その階数下界は反対階数を使わない。従って従属方向は高々三点である。
二つの一次行の外積は $x_2(t)$ を $p$ の中央支持の消去点から射影した三座標に比例する。
この射の引戻し次数は2以下であり、中心を避ける三次曲線の射影ではない。
よって消去点はある $x_2(t)$ に等しい。逆向きの前節の三行行列はこの $t$ で
$p,z$ の両方に消され、階数2以下になる。逆向きでも末端直線への制限は同型なので、
ここが逆中心 $b$ である。順逆を交換して他方も従う。
\end{proof}
\begin{proof}[同じ五次式における階数一の矛盾]
$U=\ker p$ とし $R_0$ を $L$ に制限する。前節よりこの写像は $U\otimes H^0(L,\mathcal O(1))$ への同型である。
その逆写像の後で位置0を $p$ で評価した写像を $\beta$ とし、値域は $p$ の順支持 $S$（次元2）である。
$\beta_t(u)=\beta(u\otimes[-,t])$ と置くと
\[
             w(p,-,x_2(t),y)=\beta_t(A(t)y).
\]
$t_b$ では左辺は恒等零であり、一方 $A(t_b)$ は階数3で $U$ へ全射である。
ゆえに $\beta_t=l_b(t)\beta_0$、$\beta_0:U\to S$ は階数2の固定写像である。
$\ker\beta_0=ku$ と書くと、$p^-\to p$ から $A(t)p^-=h(t)u$ である。
$h$ は非零二次式である。恒等零なら一般の核の点 $x_3(t)$ が一定になる。
その根で $p^-$ が後続となるが、$p^-\notin T$ なので順中心を必要とする。
従って $h$ は $l_a^2$ の非零定数倍である。

$z(u)\ne0$ を確かめる。もし零なら $w(z,L,x_2(t),p^-)=0$ が全 $t$ で成立する。
$x_2$ は全位置2を張るので、$z$ の逆三項鉛筆を $p^-$ で評価した行列の位置1支持は
$L$ の消去二空間と $\ker z^+$ の交叉に入る。$z^+\notin L$ より交叉は一次元である。
しかし $p^-\notin T$ だからこの鉛筆元は非零で階数2以上であり矛盾する。

$T$ を張る二座標、$p^-$、さらに一補元を位置3の基底に取ると
$A(t)=(t_0L(t),t_1L(t),h(t)u,B(t))$ と書ける
（$L(t)$ はここだけ前節の一次列を表す）。五次欠損の式は
$h(t)\det(L(t),u,B(t))$ でその唯一の根は $t_a$、重複度5である。
よって二次因子 $h$ を除いても行列式は $t_a$ で消える。
この点では $w(z,a)=0$ だから $L(t_a),B(t_a)$ はどちらも $z$ に消される。
$z(u)\ne0$ より両列は従属であり、$h(t_a)=0$ なので $A(t_a)$ の階数が1以下になる。
これは順評価階数の下界2に反する。
\end{proof}

\begin{corollary}\label{cor:rank-ten-conic-center-cycles}
反対階数10で曲線分岐が存在するとすれば、各位相に順中心と逆中心が各一つあり、
それぞれが同じ向きの中心だけの四周期点列をなす。どの中心も特異三項または四次元で、
隣り合う二つの同向き中心がともに四次元になることはない。
各列の初めと終わりの三つの動く次数はともに $(1,2,1)$ であり、第二二次曲線は
その位相の順中心と逆中心をともに含む。
\end{corollary}
\begin{proof}
第二二次の最初の直線に乗る唯一の順中心へ写すと、その前点の一意性より中心間の単射ができる。
四位相を一巡すると有限集合の全単射なので、各位相に同数の中心があり、全てが次の中心として現れる。
逆向きも同じで、列の始終の対応で両向きの総数は等しい。
前節の二次曲線の鉛筆の議論より現れ得るのは特異三項または四次元で、四次元の直後は特異三項である。
そこで特異三項の順中心 $p$ を一つ選べる。その後続直線の外から来る次の点は
一つの平面 $b=0$ に入る。これは同鉛筆の末尾縮約に含まれる $S_p\otimes kb$ の純関係による。
$p$ と同位相に他の順中心があれば、その別の直線は $L_p$ と交わらず第二二次曲線はこの平面に入る。
さらにこの位置にある全ての列の最終直線はこの平面に入る。実際その前の曲線は
逆向きの第二二次曲線なので $L_p$ ではない。各位相に二中心以上あればこれら最終直線は二本以上あり、
先の別の二次曲線と合わせた平面内の総次数が4以上になって矛盾する。
従って各々一つである。二つの周期列は同じ位相で二方向中心が一致しないため互いに交わらない。

任意の順中心 $p$ を位置0、次の中心を $q$ とする。
第二二次の行列 $A(t)$ は $q$ の鉛筆である。$q$ が四次元なら $p$ は三次元で、
$D=A(t)V_3^*\subset U\otimes S_1$ は四次元、$\dim S_1D\le10-3=7$。
ゆえに $t_0u,t_1u$ の二列がある。一般階数2なら第一直線への制限像の消去式から
$A(t)$ の全列が $u' S_1+kv(t)$ に入る（前節の左一次関係の議論の一次行列版）ため
四列独立になれない。一般階数1も一次行列の全行または全列が固定方向となり不可能。
従って一般階数3で、第三曲線は二列の一次関係による直線である。

$q$ が特異三項なら $A(t)=u\alpha(t)^t+v(t)b^t$ の全非零値が階数2である。
第三曲線は $b=0$ 内に入り、同じ位置の最終直線も上の純関係によりこの平面に入る。
第三曲線が三次ならこの二成分の次数和に反する。二次なら $f_2=3-d_3=1$ で
第二曲線上の順中心 $r$ があり、その後続直線は階数2の $A(t_r)$ の核全体に等しいので
$b=0$ に入る。最初の後続直線と最終直線は異なる（各点の前点と次点がそれぞれ固定なら
動く位置が一つだけの列になる）。第三二次と二直線が同じ平面に入り再び矛盾する。
よって $d_3=1$ で、逆方向も同様である。
次数式により $f_2=2$、また $d_4=-2+b_2-f_3\ge0$（第三位置が最後なら $d_4=0$）。
従って第二曲線上には同じ位相の順中心と逆中心の双方がある。
\end{proof}

\begin{theorem}\label{thm:rank-ten-branch-exclusion}
反対階数10の標準正則窓で曲面射影がなければ、四辺の両射影の全ファイバーは零次元以下である。
従って同階数の全窓の正則性判定でなお未解決となり得るものは四辺が有限な窓に限られる。
\end{theorem}
\begin{proof}
前の系の分岐配置を仮定し、順側四周期列を $p_i$ とする。四次元の直後は特異三項なので、
巡回して $q=p_0$ を特異三項にできる。その中央支持の消去点を位置2の $d$ とする。
その鉛筆を先と同じく
\[
       A(t)=u(t_0\alpha_0+t_1\alpha_1)^t+(t_0v_0+t_1v_1)b^t
\]
とする。$\alpha_0,\alpha_1,b$ は位置2で $d$ を消す三文字、$u,v_0,v_1$ は位置3で
$q$ の前点 $p_3$ を消す三文字である。

一つ前の順中心 $p_3$ の列で、その第二二次曲線の上の順中心は $p_1$ である。
この二次曲線に沿う行列は、最初の直線上の順中心 $q$ の上記鉛筆であり全非零値が階数2である。
従って $p_1$ の後続直線はある $A(t)$ の核そのもので、$b=0$ の平面に入り $d$ を通る。
位置2における最終直線も $b=0$ に入る。実際その前の成分は二次曲線なので $q$ の後続直線ではなく、
$q$ の鉛筆が与える $S_q\otimes kb$ の純関係を使える。この二直線は異なる。

$q$ 自身の第二二次曲線はしたがって $b=0$ 内にはない。含まれれば別の二直線との次数和が4になる。
命題\ref{prop:singular-ternary-center-projections}を $q$ 自身の列に用いると、
第二二次曲線の $d$ からの射影は、第三直線の $p_3$ からの射影に線型同定される。
二次曲線の点からの射影は非定数であり、第三直線側の引戻し次数は高々1である。
よって $d$ は第二二次曲線上にある。
そのパラメータでの第一点は $p_1$ である。$d$ はこの中心の後続直線上にあり、前点は一意だからである。
一方第二二次曲線は $p_2$ を通り、その一意な前点も順側周期列により $p_1$ である。
第一直線のパラメータ表示は同型なので $d=p_2$ が従う。

ゆえに $w(p_0)$ の末尾縮約の三次元像は
$K=R_1\cap(V_1\otimes\ker p_2)$ に入り、その左支持は二次元の $S_q$ である。
$p_2$ の逆評価は階数3なので $\dim K=6-3=3$、この像は全 $K$ に等しい。
ところが $p_3$ の前点は $p_2$ なので、$w(p_3)$ の位置0での全縮約も $K$ に入る。
従ってその中央位置1での支持は $S_q$ に含まれ次元2以下となる。
これは系\ref{cor:no-binary-centers-without-two-lines}に反する。
こうして二つの分岐配置を全て排除した。
曲面以上なら定理\ref{thm:no-moving-plane-pencils}、正次元の曲線グラフなら
系\ref{cor:all-one-dimensional-point-graphs}が全次数判定を与えるので、有限辺だけが残る。
\end{proof}

% !TeX root = ../papers/nc-p3.tex
\section{直線と第二曲線に制限した反対切断の余次元}
\label{sec:branch-rank-one-residual}
\begin{theorem}\label{thm:branch-rank-one-residual}
曲面射影のない標準正則窓の順中心 $p$ の中央支持次元を $a$、その後続直線 $L$ の
一般点の中央支持次元を $h$ とする。$h=4$ なら $\rho\ge a+6$ である。
二直線列では両中心の中央支持次元 $a,b$ と $r=\operatorname{rank} w(p,z)$ が
$r<\min(a,b)$ を満たす。特に反対階数8で二直線列は存在しない。
全て逆向きにも成り立つ。
\end{theorem}
\begin{proof}
二直線列の第二直線を $C$ とする。命題\ref{prop:two-line-branch-block}の
直交束の列によって $R_2$ の $C$ への制限は単射である。
従って $w(p)$ の中央支持は $C$ の消去空間に含まれない。
実際そうなら $w(p)$ の位置1からの縮約は全てこの単射の核に入る。
$w(p,z)$ の位置2支持は $w(p)$ の中央支持と $C$ の消去空間の交叉に入るので
$r<a$。逆向きに $r<b$。命題の $\rho\ge6+a+b-r$ と $a,b\ge2$ から下界9を得る。

一般の分岐列に戻り、$h=4$, $\rho=a+5$ と仮定する。これより小さい値は
分岐列の隣接中央支持の狭義不等式で除かれる。
第二曲線の次数は $d=3-f_1\in\{1,2,3\}$ である。
二項変数 $s,t$ を独立にして直線 $x(s)$ と第二曲線 $y(t)$ に制限すると
\[
             w(-,x(s),y(t),-)=[s,t]A(t)
\]
であり、$A$ は次数 $d-1$ の位置0と3の間の行列である。対角上の消滅と
$s$ に関する一次性を用いた。各一般 $t$ で $w(-,y(t))$ の全行は
$W_{x(t)}^\perp$ に含まれ、その $W_p$ による商は一次元である。
$W_p$ は全 $x(s)$ 上で消える。従って $s\ne t$ として評価すると
$A(t)$ の階数は高々1である。$d=1$ なら次数式 $d_3=-f_2$ より動く位置は二つだけとなるが、
二直線列の定数 $A$ は命題\ref{prop:two-line-branch-block}により階数3である。

$d=2$ の場合、第二曲線は平面 $\mathbb P(H)$ 内の滑らかな二次曲線である。
$H\subset V_2^*$ の次元は3で、二次の二項式の完全系によるパラメータを取る。
$w(x(s))=\sum_{j=1}^3 K_j(s)\otimes B_j$ とし、一次列 $K_j(s)$ を $H$ 上に制限する。
これらは完全二次系の一次関係であり、その二空間は
$(-t,s,0),(0,-t,s)$ 型の標準二列で張られ、各点で独立である。
三つの制限はこの二空間を張る。実際元の三列は一般点で独立なので、
一次元核 $H^\perp$ を法としても一般階数は2以上である。
そこで定数の列変換により最初の二列が標準二列に制限され、第三列は
$H^\perp$ の固定ベクトルと非零一次式の積となる。その一次式の零点 $q$ では
$K_j(q)$ の包から $H^*$ への射は単射である。
ゆえに $w(q)$ の位置0,3の支持は、その位置2を $H$ に制限しても保たれる。
上の表示の $A(t)$ は階数1以下の一次鉛筆なので、固定した左または右の一次元支持を持つ。
これは二つの可分解行列の和が階数1以下なら一方の因子が比例することから従う。
従って $w(q)$ にも一方の支持が高々1となる。
位置0では $L$ の全点で逆支持3、位置3では中央支持が少なくとも2なので矛盾する。

最後に $d=3$ では三列 $K_j$ は全点で独立である。これらが張る $3\mathcal O(-1)$ の
$4\mathcal O$ における商は $\mathcal O(3)$ で、切断を取れば
$H^0(4\mathcal O)\to H^0(\mathcal O(3))$ は同型である。
従って $y(t)$ は全位置2を張る三次曲線となる。一般の $s$ で $w(x(s))$ の
位置0,3の支持が次元3,4なので、$A(t)$ の左・右支持もそれぞれ少なくとも3,4である。
階数1の二項多項式行列を原始な一列と多項式行に因数分解すると、
左右支持の次元はそれぞれの因子の次数に1を加えた値以下である。
次数2では左因子に次数2以上、右因子に次数3以上を割り振れず矛盾する。
\end{proof}

\begin{corollary}\label{cor:rank-eight-center-types}
反対階数8で曲面射影がなければ、順・逆中心の中央支持は2または4である。
中央支持4の順中心の後続直線は、一つ前の位置の中央支持2の順中心の中央支持の消去直線に等しい。
中央支持2の順中心がある位置に他の順中心はない。特に一位置の順中心は高々一つである。
\end{corollary}
\begin{proof}
$a\ge3$ なら定理により $h=3$ であり、定理\ref{thm:branch-middle-defect-bound}の
反対点 $y$ は一つ前の位置で中央支持2である。$y$ が逆中心ならその対称二乗平面は
$R_1\cap(S\otimes V_2)$ 全体となり、$p$ の中央支持も2になって矛盾する。
ここで $S$ は $p$ の順支持である。$y$ は従って順中心で、その中央支持は $S$。
$R_0$ には $S$ の全方向に終わる純関係がある。$a=3$ なら $p$ の末尾からの
三次元縮約像は $2\times3$ 行列の中でSegreと交わり、連続純関係ができる。
従って $a=4$ である。

中央支持2の順中心を位置 $i$ に置く。その二項対称二乗の三関係により、辺 $i+1\to i+2$ では
第一点が後続直線上にあるか、第二点が中央支持の消去直線上にある。
同位置の別の順中心の後続直線は最初の直線と交わらないので、その第二曲線は直線になり、
次数式から二直線列となる。これは定理に反する。中央支持4の順中心が二つあれば、
それぞれの前位置の中央支持2の順中心は従って同一となり、消去直線も一致してしまう。
二つの順中心の後続直線が交わらないことに反する。
\end{proof}

% !TeX root = ../papers/nc-p3.tex
\section{反対階数八の曲線分岐の排除}
\label{sec:rank-eight-branch-exclusion}
\begin{theorem}\label{thm:rank-eight-branch-exclusion}
反対階数8の標準正則窓で射影像が全て高々一次元なら、順逆いずれの辺の射影にも
正次元のファイバーはない。純関係の有無を問わない。
\end{theorem}
\begin{proof}[中央支持4の中心の排除]
系\ref{cor:rank-eight-center-types}によって各位置の順中心は高々一つである。
中央支持4の順中心 $q$ が位置1にあれば、位置0に中央支持2の順中心 $p$ がある。
$p$ の後続直線を $L$、中央支持の消去直線を $N$ とすると $q$ の後続直線は $N$ である。
$p$ の二項対称二乗平面には $N$ の消去二空間の全方向に終わる純関係がある。
それらの方向は $q$ の順支持全体でもあるので、$w(q)$ の末尾からの縮約像に
純関係はない。従ってこの $2\times4\times3$ テンソルの二項鉛筆は $L_3$ である
（系\ref{cor:pure-free-branch-rank-nine}の鉛筆の証明）。
そこで辺 $2\to3$ では第一点が $N$ 上または第二点が $q$ の反対三次曲線上となり、
位置2にある順中心は必ず $N$ に属する。

まず $q\in L$ は不可能である。このとき $L$ の一般点の中央支持は4で、
$p$ の第二曲線の次数は $d=3-f_1\le2$。二直線列の排除より $d=2$ である。
定理\ref{thm:branch-rank-one-residual}の円錐曲線の場合の記号を使う。
$W_{x(t)}^\perp/W_p$ が今は二次元なので、同じ評価により
$w(-,x(s),y(t),-)=[s,t]A(t)$ の一次鉛筆 $A$ は一般階数が高々2である。
円錐曲線への制限で支持を失わない点は第三列の一次式が零となる順中心で、
唯一の $q$ である。従って $A$ の全左・右支持は少なくとも3,4でなければならない。
このような一次鉛筆はない。二つの一般値の行空間は、次元和が4に達して全四空間を張るので
各々二次元で直和である。するとその二行列の和の像は両列空間の和となり、
階数は少なくとも3になってしまう。

従って $q\notin L$ で、位置1に他の順中心はないから、$p$ の第二曲線は
完全三次系による三次曲線である。これは $N$ とは異なり末端直線でもない。
よって第三曲線は $q$ の反対三次曲線であり、$f_2=6-3=3$ より第二曲線上に順中心がある。
その点は先の二者択一より $N$ 上にある。$N$ の全点の前点は $q$ 一つだが、
$p$ の列では $L$ 上に前点を持つ。これは $q\notin L$ に反する。逆向きにも中央支持4はない。
\end{proof}

\begin{proof}[二組の中央二次元中心と三次曲線]
従って全ての中心の中央支持は2である。各順中心 $p$ の二直線を $L_p,N_p$ と書く。
次位置に順中心 $q$ があれば $L_q\cap N_p=\varnothing$ である。実際、対応する
消去二空間の共通方向は、続く二辺の終文字・始文字として連続純関係を与える。
従って $q\in L_p$、第二曲線の次数は2である。第三曲線は動き、第二曲線が
$L_q$ と異なるので $N_q$ 上にある。ところが $f_2=3\cdot2-3-1=2$ より
第二曲線上のある順中心から $N_q$ に後続があり、同じ不交性に反する。
従って順中心の位置は隣り合わない。逆側も同じである。

順中心が一位置だけの場合は定理\ref{thm:branch-centers-opposite-phases}の
次数十の接空間の議論がそのまま適用される。その部分は純関係の不在を使わない。
詳しくは、その中心の列で内部の順中心に達する必要があるため、
最初の三次数は $1,3,6$、第四次数は $10+b_2$。第四曲線が上限10を使い切り、
この位置の順中心に会うのはその曲線だけである。唯一の順中心は終端強連結成分なので
命題\ref{prop:terminal-branch-contact}により、そこに至るパラメータは一つで不分岐である。
しかし五行評価の次数十の曲線はこの被約曲線に一致し、順中心で接次元3であるから矛盾する。
よって順中心は相対する位置0,2の各一点であり、逆側も相対する二位置を占める。

位置0の順中心 $p$ からの列の第二曲線 $C$ は有理正規三次である。
次の曲線 $D$ は直線にはなれない。実際 $p$ の中央支持 $M$ に $C(t)$ を制限した比は
非定数である。$w(p,C(t))$ が与える $2\times3$ の $L_2$ 鉛筆はこの比を $[u:v]$ として
核点 $[v^2:-uv:u^2]$ を持つ。従って $D$ を $p$ の逆支持の双対平面へ射影すると
この滑らかな二次曲線をたどる。射影中心で一定になることは $D$ が動くため不可能である。
ゆえに列は少なくとも第四曲線まで続く。位置2の順中心を $a$ とすれば $D$ は
$L_a$ とは異なり、第四曲線は $N_a$ に等しい。この直線の次は一定である。
実際、$R_3$ の $N_a$ への制限は、二項対称二乗の三関係が消えるため高々三次元である。
一般点の逆支持は3なので三つの残余関係の評価は独立である。
従ってこの点での三テンソルをそれら三評価と固定した三係数テンソルの積の和で書けば、
順支持は固定係数の支持の和で一定である。列の一般点ではそれが三次元なので後続は定点となる。
この終中心を $z$（位置1）とする。逆側の隣の位置には逆中心がなく、逆順でも
第二曲線は完全三次系なので $D$ も有理正規三次である。
添字を二つ進めて他の順中心から得る列も同様で、その終中心を $b$（位置3）とすれば
$b$ の前点直線は $N_p$ である。

$z$ の中央支持の消去直線は $L_a$ であることにも注意する。
実際 $z$ の逆支持は $N_a$ の消去二空間 $M_a$ であり、
$R_3\cap(V_3\otimes M_a)$ は三次元である（逆中心の直線への制限の階数3）。
これには $a$ の二項対称二乗平面が入り、それで尽くされる。
ゆえにこの平面を張る $w(z)$ の縮約の位置3支持は $L_a$ の消去二空間である。
\end{proof}

\begin{proof}[二つの接触と二次式の二次ジェット]
$p$ の列を $p,L_p,C,D,N_a,z$ と書く。$C,D$ のパラメータを移動で揃える。
$L_2$ の核点表示によれば $C$ の $N_p$ からの射影の次数を $k$ として
$D$ の一点からの平面射影の次数（引戻し直線束の次数）は $2k$ となる。
三次埋込みを一点から射影すると、この次数は3または2である。従って $k=1$ で、
$C\cap N_p$ の長さは2である。$N_p$ の全点は $b$ のみを後続とするので、
この交叉の台は $D$ が $b$ に至る一つのパラメータ $B$ である。
従って $N_p$ は $C$ にそこで接する。
逆の列に $z$ の中央支持の記述を使うと、$D\cap L_a$ も長さ2で、
$C$ が $a$ に至るパラメータ $F$ で接する。$B\ne F$ である。
実際 $D(F)\in L_a$ の逆支持は3であるのに $D(B)=b$ の逆支持は2である。

$C(t),D(u)$ への $R_2$ の制限は両三次曲線が全空間を張るため単射で、
六本の独立な $(3,3)$ 式を与える。全て対角で消えるので括弧 $[t,u]$ で割ると、
二組の二項二次式のテンソル積の中の六空間 $J$ が得られる。双対の消去空間 $E$ は三次元である。
二項二次の評価ベクトルを $v(t)$ と書く。$a\times L_a$ 上の消滅と $D$ の接触から、
$t=F$ として $u$ で微分すると、括弧を割った全式は $(F,F)$ で消える。
同様に $N_p\times b$ 上の消滅と $C$ の接触から $(B,B)$ でも消える。従って
\[
                v(F)\otimes v(F),\quad v(B)\otimes v(B)\ \in E .
\]

一方、次数列 $1,3,3,1$ より $C$ 上の順側欠損は $f_2=3$ で、
位置2の唯一の順中心 $a$ に集中する。$C$ は埋込みであるからパラメータは $F$ 一つである。
$J$ の六行を第一変数で評価した六行三列の行列は、この点で階数2、
三次小行列式の局所共通零次数は3である。実際元の関係評価はこれに $[t,u]$ の乗算で
二項二次式を三次式に入れる写像を合成したもので、この乗算は各 $t$ で分裂単射である。
一様化変数 $l$ で単因子は $1,1,l^3$。従って法 $l^3$ で原始な核ベクトル
$k(l)=k_0+lk_1+l^2k_2$ があり、先の接触より $k_0$ は $v(F)$ の非零倍である。
$v(t(l))=e_0+le_1+l^2e_2$ と書くと、二次完全系により三係数は独立である。
$E$ は $v(t(l))\otimes k(l)$ の最初の三係数を含み、
\[
 e_0\otimes k_0,\qquad e_1\otimes k_0+e_0\otimes k_1,\qquad
 e_2\otimes k_0+e_1\otimes k_1+e_0\otimes k_2
\]
がその基底となる。この三空間の非零な可分解元の右因子は全て $k_0$ に比例する。
実際、非零な最大番号の $e_j$ における行は必ず $k_0$ の非零倍である。
しかし第二の接触は異なる右因子 $v(B)$ を持つ可分解元を与え、矛盾する。
\end{proof}

% !TeX root = ../papers/nc-p3.tex
\section{三次縮約の三支持からの部分代数と低い反対階数の全次数判定}
\label{sec:ternary-point-subalgebras}
有限Frobenius窓の標準四テンソルを $w$ とし、全辺の左右の射影ファイバーが空または一点とする。
点グラフの次元は問わない。一文字を点 $a$ で縮約した三テンソルを巡回順に読む。
\begin{theorem}\label{thm:ternary-point-subalgebras}
ある縮約 $u=w(a)$ の三つの一因子支持が全て三次元であるとする。
各因子の支持に制限した $3\times3\times3$ テンソル $u$ について、次は同値である。
\begin{enumerate}
\item 元の四窓が標準Koszul左右AS正則である。
\item 三方向のいずれの非零一文字縮約にも階数1の行列がない。
\end{enumerate}
従って一つの縮約の三つの射影平面で二次小行列式の共通零点を調べれば足りる。
このとき $A$ は三生成二次AS正則部分代数の上で左右自由であり、
判定に失敗すると指数成長する。
\end{theorem}
\begin{proof}[三文字の列の構成]
最初の三支持を順に $H_i,H_{i+1},H_{i+2}$ とする。添字を四移動で読めば
$H_i=R_{i-1}(a)$、$H_{i+2}=R_{i+2}(a)$ であり、
その直交点 $p_i,p_{i+2}$ は $a$ の次点、前点である。
$u$ の末尾と先頭の縮約像はそれぞれ $U_i\subset H_i\otimes H_{i+1}$,
$U_{i+1}\subset H_{i+1}\otimes H_{i+2}$ という三次元の二次関係空間である。
各 $U_j$ は $R_j$ のうち先頭が $H_j$ に支持される部分および末尾が $H_{j+1}$ に
支持される部分に等しい。実際それぞれの評価核に含まれ、その核の次元はファイバー仮定で高々3である。
特に中支持の直交点 $p_{i+1}$ も共有点グラフに属する。

$p_{i+1}$ の唯一の前点 $b$ で $w$ を縮約すると、位置 $i+1$ から始まる三テンソルの
先頭は $H_{i+1}$ に支持され、末尾の支持も三次元である。その末尾を縮約した像は上の
$U_{i+1}$ と等しいので、中央の支持はちょうど $H_{i+2}$ である。
こうして次の $H_{i+3}$ が定まり、逆向きも同様に延長できる。
得られた全列で $U_j$ は上の左右の評価核であり、三次重なりは既知の三テンソルの直線である。
後者は、元の四文字代数の三次重なりのうち $p_j$ での先頭縮約が零になる部分が一次元であるためである。
実際この部分は $w$ の末尾縮約のうち $p_j$ の前点を使うものだけである。

抽象二次代数 $P=T((H_j))/(U_j)$ の二次双対は従って $(1,3,3,1)$ の有限Frobenius型である。
三次テンソルの各一因子縮約は相補関係への同型、四次重なりは非零なら三次テンソルの片端支持が
一次元になってしまうので零である。次の三テンソルは前のものの巡回回転と、余った一因子の
基底変更で一致する。実際、両者は共通の $U_{i+1}$ への同型を与える。
従って全三位置の縮約階数は最初の $u$ で検査できる。
\end{proof}

\begin{proof}[自由性と有限判定]
各始点の $p_j$ における右点加群を $M_j$ とする。
命題\ref{prop:commuting-three-point-graph}の切断計算と同じく
$M_j=e_jA/H_jA$ であり、$H_j\otimes e_{j+1}A\to e_jA$ の関係は
次数2の高々三個で生成される。この計算に三文字の可換性は必要なく、
元の候補複体上の点列での階数 $1,3,3,1$ を用いる。
二次には独立な $U_j$ が既にあるからそれらが生成する。
$P$ 単純の自由表示の最初の二写像と比較して $\operatorname{Tor}^P_1(k_j,A)=0$。
各終点の有界下加群の極小表示から $A$ は左 $P$ 自由で、各長さに一個の生成元がある。
逆向きの点加群によって右自由性も従う。特に $P$ は $A$ に埋まる。

$P$ の三窓について、補題\ref{lem:ternary-point-graph-growth}と
補題\ref{lem:ternary-finite-graph-test}の証明を添字付きで適用する。
具体的には、三テンソルを順に回転する表示では、一つの点での左評価と右評価は
同じ二文字行列の転置または定数基底変更である。階数1なら、その両端像を割った二文字には
二次元ずつが残り、関係は一本が消えて二本以下である。
三位置を一期として独立に繰り返し、長さ $3m$ で $2^m$ 以上となる。
逆に全て階数2以上なら三位置の両側射影は同じ行列式軌跡をスキームとして共有し、
全て平面または平面三次曲線の同型のグラフになる。平面なら可換多項式型である。
三次の場合はグラフに沿う束 $L_j$ を順次使えば、上記補題と同じ分裂局所完全な
階数 $1,2,1$ の評価複体になる。正の偏極積の $H^1$ は消え、長さ3の左端の
$H^1(\mathcal O)$ だけが残る。よって全切断環を与える商への核は各位置で
左右とも正則な三次の一元で生成され、$P$ のHilbert関数は三変数多項式環のものとなる。
左右のsocle消滅、Euler式および三次対合から $P$ は標準三次元Koszul左右AS正則となる。
これらの計算に同じ束の反復や普通の次数付き表示を仮定しない。

最後に $P$ が正則なら左右自由性と各長さ一個の自由生成元より $H_A=(1-z)^{-4}$。
任意の位置の $H_j$ の二次元部分空間には右乗法が単射の文字が存在する。
平面の場合は明らかであり、三次曲線の場合は独立な二一次式が共に恒等零となる既約曲線成分はないので、
全ての随伴点を避ける一次式を取れる（平面三次は純である）。
その右乗法は切断商で単射であり、先の正則三次元を左側にくくり出す次数降下により $P$ でも単射である。
右 $P$ 自由性から $A$ でも同じである。従って点加群の任意の尾部から $A$ への右Homは零である。
実際、各位置の点加群の元は $H_j$ の二次元以上で消える。
$e_jA_+$ の右線型写像を $e_jP_+$ に制限すれば、三次元AS性と自由性で
左乗法に延びる。それを引いた写像は $H_jA$ を消して点加群の尾部を通るので零である。
これで右単純から $A$ への次数0,1のExtが消え、左側も同様である。
四文字候補複体のFrobenius双対およびHilbert関数のEuler式で正則性が従う。
\end{proof}

\begin{corollary}\label{cor:low-opposite-point-graph-as}
点グラフの四辺の左右射影ファイバーが空または一点なら、反対階数 $\rho\le7$ の
全候補で標準AS正則性が有限に決まる。$\rho<6$ は全て除外される。
$\rho=6,7$ の場合は三支持が三次元の縮約の有無を調べ、なければ除外、
あれば一つの縮約で前定理を適用する。
特に四辺が有限な枝は $\rho\le7$ の全域で、全次数の正則性と印を指定しない添字保持同型・周期が
有限判定となる。
\end{corollary}
\begin{proof}
六双一次式の零点は空でない。実際、Segre埋込みの六次元射影多様体と余次元6の線型空間は交わる。
標準なら任意の零点で三次縮約の両端支持は3であり、中支持は少なくとも3である。
中支持2は補題\ref{lem:middle-binary-pencil}の両端次元の不一致に反し、
中支持1なら同じ中央文字での連続する純粋関係ができる。
反対切断の $\rho$ 次元を経る二つの評価行列は、辺の点対で積が零である。
従って両点での中支持の和は $\rho$ 以下となる。
$\rho\le7$ なら少なくとも一方が3である。有限点の場合の同型と周期には
定理\ref{thm:finite-all-window-keys}を適用する。
\end{proof}

\begin{theorem}\label{thm:finite-frobenius-low-opposite-as}
四辺が有限な標準有限Frobenius候補では、反対階数によらず全ての $R_i$ が
非零の可分解二テンソルを含まない。この枝で $\rho<6$ の候補は存在せず、
$\rho=6,7$ なら全ての候補が標準Koszul左右AS正則である。
\end{theorem}
\begin{proof}
系\ref{cor:finite-edge-point-sequences}の共通長さ20のスキームを $X$ とし、
直線束の自明化後の四文字の切断空間を $V_i\simeq U_i\subset D=H^0(\mathcal O_X)$ と書く。
定理\ref{thm:finite-square-chart}の六双一次式の完全交叉とtraceの議論は、
この共通グラフが成立してから先には元代数のAS性を使わない。
従って各 $\operatorname{Sym}^3 U_i\to D$ は同型、$R_i$ は積の核、
$U_{i+2}^2$ は $U_i^2$ の単元倍である。

もし非零の $a\in U_i,b\in U_{i+1}$ について $ab=0$ なら、
$a U_i^2$ と $b U_{i+1}^2$ は三次の単射性により各十次元である。
$(a)\subset\operatorname{ann}(b)$ かつ $\dim\operatorname{ann}(b)=20-\dim(b)$ なので、
両主イデアルが次元10となる。$Y=\operatorname{Spec}(D/(a))$ を三つ目の四文字 $U_{i+2}$ で見る。
この空間の二次切断の制限像は次元6以下である。実際 $U_i^2$ の制限核は
四次元の $a U_i$ を含み、単元倍で核の次元は変わらない。三次の制限像は全長10である。

この $Y$ は平面に入る。入らないとすると斉次座標環のHilbert関数は
$1,4,h_2,10$ で $h_2\le6$ となる。全ての点で非零の一次式を取って割ると、
非零因子なので差分のHilbert関数を得る。次数2が $h_2-4\le2$ で次数3が $10-h_2\ge4$ となるのは不可能である。
実際、初項イデアルの標準単項式は約数で閉じており、二次の標準単項式高々二つからは
三次のものは高々二つしか作れない。従ってある $0\ne c\in U_{i+2}$ が $Y$ 上で消え、
$c\in(a)$、ゆえに $bc=0$ となる。すると $a\otimes b\otimes c$ が三次重なりに入る。
有限Frobenius対合によりこれは標準四テンソルのある非零一文字縮約である。
その両端支持が一次元ということは、その文字の点で左右の関係評価が階数1ということであり、
有限ファイバーに反する。これで可分解関係は全て排除された。

辺の各点で三次縮約の両端支持は3である。もし中支持が2なら、
中央因子で縮約する $3\times3$ 行列の鉛筆はある非零パラメータで特異になる。
その片側核で元の三テンソルを縮約すれば、残る中央支持が一次元の非零二次関係を与え、
先に得た結論に反する。中支持1も同様に不可能である。
従って上の反対階数の和の不等式はAS性を仮定せず適用でき、$\rho\le7$ なら三支持が3の元が存在する。
その三方向の縮約のうち両端からのものは既に $R_j$ 内にある。
中央からの縮約も階数1になれないことを確認する。もしそうなら、
構成した三文字部分代数の第三の関係空間はその巡回回転を含む。
この空間も延長した $U_{j+2}\subset R_{j+2}$ に可逆基底変更で等しいから矛盾する。
ゆえに前定理により正則である。
\end{proof}

% !TeX root = ../papers/nc-p3.tex
\section{有限点枝の分離三次式とスキーム上の標準四移動}
\label{sec:finite-cubic-separators}
四辺が有限な標準有限Frobenius候補とし、点グラフを $X_i$ の間の同型の列と見る。
標準座標の四移動で $X_{i+4}=X_i$ とする。
\begin{theorem}\label{thm:finite-cubic-separators}
全ての有限Frobenius候補で、標準座標における四移動は点スキーム上で恒等である。
反対階数も元代数の正則性も仮定しない。
中支持3の点を一つでも持つ候補は全次数で正則であり、未認証候補では全ての中支持が4である。
任意の閉点 $p\in X_i$ を固定し、反対位置の変数 $y$ において $w(p,y)$ を
両端の三支持の間の $3\times3$ 行列 $M_p(y)$ と読む。
$\det M_p$ は非零の三次式で、反対位置の長さ20のスキーム上では二移動した対応点の
局所socleを張る。特にその点が被約なら当該点で行列の階数は3である。
\end{theorem}
\begin{proof}
前節により各二次関係空間は可分解元を含まず、中支持の次元は3または4である。
中支持4の三テンソルの両端を基底化する。行列 $M(y)$ の四係数は独立であり、
行または列を非零定数ベクトルで縮約した $3\times4$ 行列はいずれも階数2以上である。
これはそれらが隣接関係の元であり、全一因子支持と純粋関係の不在を使うためである。

$\det M$ が恒等零なら一般階数は2である。一般階数1以下の線型行列空間は
共通な像直線か共通な余像直線を持ち、端支持3に反する。
余因子行列は多項式環上の階数1の二次行列で $c(y)u(y)v(y)^t$ と因数分解できる。
実際、有理関数体上で階数1の積にし、UFD上で二ベクトルの分母と成分の最大公約因子を取り除けばよい。
$M$ は定数左右核を持たないので両ベクトルは正次数であり、従ってともに一次である。
$M(y)u(y)=0$ とする。$u$ の像が一次元なら定数核がある。
二次元なら基底を替えて $u=(y_1,y_2,0)^t$ となり、最初の二列は $y_2 c,-y_1 c$
という定数列 $c$ の倍となる。非零な一列が階数1の切断となり矛盾する。
三次元なら $u=(y_1,y_2,y_3)^t$ としたとき、第四係数行列が三つの基底ベクトルを消すので零である。
これも中支持4の係数の独立性に反する。

中支持3の場合には前節で得た三文字部分代数が標準三次元AS正則であり、
三つの縮約平面に階数1の行列はない。その三行三列の行列式が恒等零なら、
三文字の点対応は射影平面全体の同型のグラフである。
従って元の一辺の六関係のうち三つが二つの三文字空間の間の可換関係に基底変更で一致する。
ところがこの三関係の零点は、両射影三空間の三座標が比例する四次元の閉部分多様体を含む。
実際、共通の射影二座標比が決まった後で双方の補座標がそれぞれ一自由度を持つ。
この四次元多様体と残り三つのSegre超平面の交叉は次元1以上であり、有限性に反する。
従って中支持3の場合も行列式は非零である。

変数 $p$ の属する局所Artin成分を $S$ とする。$p$ から反対位置へ前向き・後向きに
二移動する射を $f,g:S\to X_{i+2}$ とする。
$M_p(y)$ の両端支持は $S$ 上の階数3の部分束となるので、局所基底を選んで
$\delta=\det M_p(y)$ を $S$ 係数の三次式とする。反対文字の全三次式の制限は
$D=H^0(X_{i+2},\mathcal O)$ への同型なので、直線束を自明化して
$\delta\in\mathcal O_S\otimes D$ と読める。閉ファイバーでこれは非零である。

$S\times X_{i+2}$ 上では $y$ 自身の先頭・末尾評価も階数3である。
従って $w(p,y)$ の各外側の支持は、$p$ により定まる超平面と $y$ により定まる超平面の
双方に含まれる。$M_p(y)$ の三列を、$p$ 側の超平面に制限した $y$ 側の点で縮約すると零になる。
余因子行列を掛けると、その点の三つの制限座標は全て $\delta$ を消す。
これら三座標が生成するのは二つの超平面が等しいという対角イデアルの引戻しである
（単元座標で規格化すれば三つの座標差である）。
片端ではこれが $f$ のグラフイデアル $I_f$、他端では $g$ のグラフイデアル $I_g$ となる。
これは隣接辺の射影がどちらも同じ点スキームへの閉埋込みであることによる。

有限Gorenstein環 $D$ の完全traceを用いると
$\operatorname{ann}(I_f)=\operatorname{Hom}_{\mathcal O_S}(\mathcal O_S,\mathcal O_S)$
に可逆層の自明化を除き同定できる。具体的には相対完全積対合のもとで
グラフ商 $\mathcal O_S$ の相対双対に等しく、その形成は閉ファイバーへの制限と可換である。
従って閉ファイバーで非零の $\delta$ は $\operatorname{ann}(I_f)$ を張る。
同じく $\operatorname{ann}(I_g)$ を張り、直交補を二度取れば $I_f=I_g$ を得る。
これで $f=g$ であり、前へ四移動すると $S$ 上で戻る。閉点に制限した消去イデアルは
対応点の極大イデアルの零化、すなわちsocleの直線である。
第一の場合の正則性は定理\ref{thm:ternary-point-subalgebras}と全四位置の純粋関係の不在で従う。
\end{proof}

\begin{corollary}\label{cor:reduced-point-six-mixed-relations}
被約な局所成分を一つ選び、その四周期列に沿って $H_i=\ker p_i$ とし、
$p_i(t_i)=1$ なる補文字を取る。各位置には二つの線型写像 $F_i,G_i$ が一意にあり、
元の六関係は
\[
 t_i v=F_i(v),\qquad u t_{i+1}=G_i(u),\qquad
 u\in H_i,\quad v\in H_{i+1},\quad F_i(v),G_i(u)\in H_i\otimes H_{i+1}
\]
となる。任意の正の長さの経路空間は、全て $t$ の一語と $H$ だけの語の像の直和である。
\end{corollary}
\begin{proof}
$p_i,p_{i+1}$ は辺をなすので $t_i t_{i+1}$ の係数は各関係で零である。
先頭が $H_i$ にある三関係を $p_{i+1}$ で縮約する写像の階数は
反対点 $p_{i-1},p_{i+1}$ での行列の階数に等しい。実際この三関係は
$w(p_{i-1})$ の末尾縮約像である。選んだ点は被約なので階数3であり、
先頭と末尾の双方が $H$ にある関係はない。従って六関係から
$t_i H_{i+1}\oplus H_i t_{i+1}$ への射影は同型である。一つでも $H$ がある語では、
その隣の補文字を順に消せる。全補文字の語の独立性は点列での評価で分かる。
\end{proof}

% !TeX root = ../papers/nc-p3.tex
\section{被約な点を持つ有限窓の行列交換子と三次表示}
\label{sec:marked-point-commutators}
有限Frobeniusの四辺が有限で、長さ1の局所成分を一つ指定する。
前節の四周期列を用いてその評価を各位置で $t_i$ に1、補空間 $H_i$ で0とする。
\begin{theorem}\label{thm:marked-point-commutators}
印付きテンソルは二つの三次元空間 $U,V$ と
\[
 G\in U\otimes\mathfrak{sl}(V),\quad F\in U^*\otimes\mathfrak{sl}(V),\quad
 Q\in V\otimes\mathfrak{sl}(U),\quad P\in V^*\otimes\mathfrak{sl}(U)
\]
によって次の形にできる。位置 $0,1,2,3$ の $H_i$ を $U,V,U^*,V^*$ とし、
各補文字の添字を $*$、残りを $a,b,c,d=1,2,3$ とする。
隣接する二つの $*$ を持つ係数は零で、他の係数は
\[
\begin{array}{ll}
 w_{*b*d}=\delta_{bd},&w_{a*c*}=\delta_{ac},\\
 w_{*bcd}=(F_c)_{bd},& w_{ab*d}=(G_a)_{bd},\\
 w_{a*cd}=(P_d)_{ac},& w_{abc*}=(Q_b)_{ac},\\
 w_{abcd}=(G_a F_c)_{bd}+(Q_b P_d)_{ac}.&
\end{array}
\]
二つの隣接切断の階数6はこの表示ではちょうど
\[
 [G_a,F_c]_{bd}+[Q_b,P_d]_{ac}=0\qquad(1\le a,b,c,d\le3)
\]
に等しい。各三行列組は三次元空間に共通の不変直線も不変平面も持たない。
印を保つ同一視は $A\in\mathrm{GL}(U), B\in\mathrm{GL}(V), s\in k^*$ による
四因子 $(A,B,s A^{-t},s B^{-t})$ の作用と全体の $s^{-1}$ 倍で尽きる。
$*$ 座標はこの作用で固定する。
\end{theorem}
\begin{proof}
隣接する二点での縮約は零で、被約点における反対位置の二つの行列は前節より階数3である。
独立な基底変更とテンソルの規格化で両行列を単位行列にする。
$t_i$ の $H_i$ による持上げ変更は、反対位置の $*$ 一個を持つ三係数行列に
それぞれ独立にスカラー行列を加える。他の三組には影響しないため、
標数零で跡零の持上げが一意に取れる。
$01\mid23$ 切断のうち各側で一つの $*$ を持つ六行六列は置換単位行列である。
従って階数6なら残る九行九列はそのSchur補元を零にする唯一の表示で、
これが上の $w_{abcd}$ である。他の切断で同じことを行うと二積の順序がそれぞれ逆になる。
よって両条件の差が表示の交換子である。各因子の射影スカラーで $*$ の倍率を1に戻せる。
二対合を同じスカラーまで保つことと跡零の一意性で残余群を得る。
逆にこの表で作った係数について、三次・四次重なりと有限点条件は元の有限行列検査で判定できる。
正則性はこの表の二切断条件とは別である。

例えば $F_c$ が共通の不変直線を持つと、$I,F_1,F_2,F_3$ をその直線のベクトルで評価した像は
一次元であり、$w(t_0^*)$ の端からの縮約が可分解な二次関係になる。
双対側の直線にも同じ議論が適用できるので共通不変平面もない。ほかの三組も同様である。
\end{proof}

\begin{corollary}\label{cor:point-commutator-rank}
未認証の場合（全ての中支持が4）、$G,F,Q,P$ は各々三独立行列を与える。
反対階数は、$g=\dim\langle G_a,F_c\rangle$, $q=\dim\langle Q_b,P_d\rangle$ として
\[
 \rho=2+g+q+\operatorname{rank}\bar w .
\]
ここで $\bar w$ は九行九列の上の $w_{abcd}$ を反対切断 $ac\mid bd$ で読み、
それぞれ $\langle I,Q_b,P_d\rangle$, $\langle I,G_a,F_c\rangle$ を法として得るテンソルである。
特に反対階数8では $G,F$ は同じ三平面の二基底、$Q,P$ も同じ三平面の二基底であり $\bar w=0$ である。
\end{corollary}
\begin{proof}
例えば $w(t_0^*)$ の中央評価空間は $I,F_1,F_2,F_3$ の包である。
反対切断で少なくとも一つ $*$ を含む行は $I,G_a,F_c$ を残りの九列にのみ持つ。
列側でも対称的に $I,Q_b,P_d$ を持ち、印のある行と列の間は零である。
行と列を消去して表示式が従う。
\end{proof}

\begin{proposition}\label{prop:point-kernel-cubic-presentation}
指定した四周期点への評価の核は、三文字 $H_i$ だけを使う正次数の代数である。
この代数に各頂点の単位を加えたものを $B$ とすると、関係は三次だけであり、各位置に八個ある。
たとえば $i=0$ では $H_0\otimes H_2=\operatorname{End}(U)$ と読み、
三次関係の空間は
\[
 Z\longmapsto \big([Q_b,Z]_{ac}\big)_{abc}
 \qquad (Z\in\mathfrak{sl}(U))
\]
の像である。他の位置は巡回的に得る。$B$ の三文字点対応は、中央点で評価した行列と、
両端点の与える階数1の行列が交換する条件である。
各正の長さで $A_{i,j}=B_{i,j}\oplus k\,t_i\cdots t_{j-1}$ であり、
これは $B$ の三次表示から全次数を復元する。
\end{proposition}
\begin{proof}
前節の六関係をここでは $t h=L(h),\ h t=N(h)$ と書き、位置は積順で補う。
純 $H$ 語上で $L$ は最初の文字に、$N$ は最後の文字にこの置換を行う作用素とする。
三文字 $a t b$ の二つの書換えの差は $N(a)b-aL(b)$ である。
表から $N_0(a)=-\sum_b Q_b a\otimes e_b$,
$L_1(c)=-\sum_b e_b\otimes Q_b^t c$ となるので、その差の像は表示の交換子空間である。
三行列組の既約性とSchurの補題により共通中心化環はスカラーであり、像の次元は8である。
他の位置も同様である。点対応の条件はこの写像を三点で評価し、
全 $Z$ に対するtrace対合の非退化性を用いれば従う。

二次書換えで三次の経路空間は全 $t$ の一語と純 $H$ の27語で張られ、
有限Frobenius検査からその次元は20である。従って純 $H$ の三次関係は既に得た八個で尽きる。
特に $t h t$ の差 $LN(h)-NL(h)$ もそれらの包にある。
この八関係だけによる自由 $H$ 代数の商で $L,N$ が良定義であることを確かめる。
$N(a)b-aL(b)$ の左端に $L$ を施すと、
$LN(a)b-L(a)L(b)$ となる。新しい語の中央に同じ八関係を使えば
$L(a)L(b)=NL(a)b$ となるため零である。右端での $N$ も逆順で同様である。
関係が端から離れていれば置換はそのまま前後の語に作用する。よって生成イデアルは保たれる。
両置換は長さ2以上で別の文字に作用するので交換し、長さ1でも先の $t h t$ の式で交換する。
さらに任意の非空な純語で $N(a)b=aL(b)$（隣り合う末尾・先頭に八関係を適用）である。

従ってこの三次商の正次数部分と $k[t]$ の反復の直和に、
左 $t$ を $L$、右 $t$ を $N$ とする結合積を入れられる。
結合律は $L(ab)=L(a)b$, $N(ab)=aN(b)$、および直前の二等式である。
この直和が元の六つの二次式を満たし、逆にそれらから全ての指定積が従うので
互いに逆な準同型を得る。
\end{proof}

% !TeX root = ../papers/nc-p3.tex
\section{有限窓の一組の共通三平面と反対対称性}
\label{sec:marked-common-triple-symmetry}
\begin{theorem}\label{thm:marked-common-triple-symmetry}
単純な局所点を持つ有限Frobenius点窓で、既出の中支持3による認証がないとする。
印付き交換子表示で $q=\dim\langle Q,P\rangle=3$ または $g=\dim\langle G,F\rangle=3$ なら、
反対階数の仮定を置かずに少なくとも一方の反対位置が対称であり、$\rho\le10$ である。
$q=3$, $g>3$ なら基底を選んで $P=Q$、$G,F$ はすべて対称とでき、
$g=4,5$ に限る。交叉交換子の次元が2の場合には $g=5$ である。
\end{theorem}
\begin{proof}[交換子の二平面と非対称な基底変換]
巡回して $q=3$ とし、$P_d=Q_{Jd}$ と書く。$J$ は可逆である。
$N=\langle[G_a,F_c]\rangle$ は係数添字 $b,d$ の行列として
$KJ$ ($K$ は交代行列) の包に入る。次元 $r$ は $Q$ の三平面の交換子次元であり、
既約性とSchurの補題により $r=2,3$ である。$r=3$ なら跡零条件
$\operatorname{tr}(KJ)=0$ が全 $K$ に成り立って $J$ は対称である。
また $g=3$ の全場合は定理\ref{thm:rank-nine-finite-symmetry}の証明の共通三平面の場合による。
そこでは混合型も一方が対称、両交代型は点の有限性に反し、反対階数を用いる必要はない。

$g>3,r=2$ として $J$ が非対称なら矛盾することを示す。$D=J-J^t$ と置く。
$N$ を書く二平面の $K$ は全て $J$ と跡直交なので、その交代空間内の直交補は $kD$ である。
一般 $X\in\langle G\rangle$ で $[X,\langle F\rangle]=N$ となる。
そうでなければ $X$ に線型な $\operatorname{Hom}(\langle F\rangle,N)$ の部分空間が全て階数1以下となり、
共通像なら $r=2$ に反し、共通余像なら非零の跡零 $F$ が全 $G$ と交換してしまう。
従って $X,X^2$ は $N$ と跡直交する。多項式性から全 $X$ で同じである。$F$ 側でも同様である。

$JX$ の交代部分は $D$ の倍なので $X'=X-\lambda(X)I$ を $JX'$ が対称になるよう一意に取れる。
$J(X')^2$ の交代部分も $D$ の倍であり、$JX'=X'^tJ^t$ より
\[
                       X'^t D X'\in kD .
\]
$kr_0=\ker D$、$\pi$ をその商への射影とする。右の倍率が一般に非零なら
$\pi X'$ の二次外積と $\pi$ の二次外積が比例し、階数2の両写像の核は同じである。
従って $X'r_0\in kr_0$ が全族に成り立って既約性に反する。
よって $\pi X'$ はいつも階数1以下である。共通像の場合は全ての $X'$ が固定二平面に値を持ち、
元の $X$ もその平面を保ってしまう。残るのは共通余像 $k\alpha$ の場合で、
$X'(\ker\alpha)\subset kr_0$ である。

$\beta=Jr_0$ とすると対称形式 $JX'$ は $\ker\alpha$ を $k\beta$ に送る。
$\alpha,\beta$ が独立ならこうした対称形式は $\alpha^2,\beta^2$ の包に限られる。
実際 $\ker\alpha$ での対称性から制限は $\beta$ に $\beta$ の制限の定数倍を掛けたもので、
それを引くと両因子が $k\alpha$ に入る。$X\mapsto X'$ は跡零空間上で単射なので二次元では足りない。
従って $\alpha$ は $\beta$ に比例し、三次元の対称形式空間は強制的に
$\{\beta\otimes l+l\otimes\beta:l\in(k^3)^*\}$ 全体である。
$F$ 側でも同じ $D,J,r_0,\beta$ を使うため同じ空間となる。
$J^{-1}$ を掛けて跡零部分を取れば二つの元の三平面が等しくなり $g>3$ に反する。
これで $J$ は対称なので合同で単位行列に取れる。
\end{proof}
\begin{proof}[異なる三平面の交叉交換子]
今 $[G_a,F_c]$ は全て交代である。各三平面を跡零対称部分と外積行列に分けて
$G=(A,[u]_\times), F=(B,[v]_\times)$ と書き、
$D_u B=[[u]_\times,B]$ とすると条件は $D_u B=D_v A$ である。
どちらかの交代射影が零なら、相手にも交代部分はない。
非零外積行列と交換する跡零対称行列は一次元しかないからである。
どちらも零でなく片側の交代射影が一次元以下なら、その純対称な核の二元が相手の非零外積と
交換するためやはり不可能である。両側とも交代射影が二次元以上なら、相手の純対称核は零である。
独立な二つの外積と交換する跡零対称行列は零であるためである。
従って $A=A(u), B=B(v)$ と両者を全交代三空間上のグラフに書ける。

和の写像 $H=A+B$ は $D_u H(v)=D_v H(u)$ を満たす。
この十五元の線型系は定理\ref{thm:rank-eight-finite-symmetry}の行列式 $-3456$ の検算から核を持たない。
差側は従って $D_u A(u)=0$ を満たす。一般の $u$ では中心化空間により
$A(u)$ は $u u^t-(u^tu)I/3$ に比例する。この二次行列の成分は互いに全体公約因子を持たない
（既に三つの非対角単項式がそうである）。有理倍率で全成分が一次多項式になることは
零以外になく、$A=B=0$ である。この両交代型なら $g=3$ なので除かれる。
従って $g>3$ では両方とも跡零対称で $g\le5$ である。
星印成分と純成分を交換子表示で比較すれば $b,d$ 交換で $w$ が不変であり、$\rho\le10$ となる。

最後に $r=2$ なら交換子と跡直交する非零外積行列 $K$ を取る。
対称跡零の五空間の交代形式 $\omega(X,Y)=\operatorname{tr}(K[X,Y])$ の根基は一次元である。
実際、対称跡零空間上の跡対合は完全で、根基は $K$ と交換する跡零対称行列である。
二つの三平面は $\omega$ で互いに直交するので根基をともに含み、商の四次元シンプレクティック空間で
互いに直交補の二平面となる。二平面は非退化か全等方であり、後者では二平面が等しく $g=3$ になる。
前者では交叉が零で、元の和は五次元である。
\end{proof}

% !TeX root = ../papers/nc-p3.tex
\section{異なる対称三平面での二交換子鉛筆の復元}
\label{sec:marked-two-commutators}
\begin{theorem}\label{thm:marked-two-commutators}
前節の $q=3$, $g>3$, $r=2$ の場合、$P=Q$ とし、
$Q_s,Q_t$ が可換になるパラメータ二平面を取る。
その平面と補元 $w$ に対して交換子式の既知の右辺を
$C_0=[Q_s,Q_w]$, $C_1=[Q_t,Q_w]$ と読む。この二行列の鉛筆は恒等的に特異にはならない。
従って平面の基底を取り直して $C_0$ は可逆とできる。
$G,F$ とこの基底を固定すると、残る $Q$ の全復元は次の線型演算で決まる。
\begin{enumerate}
\item $D C_0+C_0 D=C_1$ を解く。解 $D$ は一意である。
\item $D$ の最小多項式が三次、$\operatorname{tr}(D^2 C_0)=0$ を要求し、
$[D,Z]=C_0$ の一解 $Z$ を取る。$L_0$ を行列 $L$ の跡零部分と書くと、全候補は
\[
 Q_s=h D_0,\qquad Q_t=h(D^2)_0,\qquad
 Q_w=h^{-1} Z_0+eD_0+f(D^2)_0,\qquad h\ne0 .
\]
\end{enumerate}
ここで $s$ は一般に取って $Q_s$ が巡回となる方向を用いる。
この式と元の有限窓の検査で印付き候補を復元できる。これだけで全次数の正則性は主張しない。
\end{theorem}
\begin{lemma}\label{lem:commuting-linear-symmetric-nets}
跡零対称三行列の二つの線型三列 $G_x,F_x$ がともに独立で既約とし、
$[G_x,F_x]=0$ が全 $x$ に成り立つなら、$F_x=hG_x$ となる。
\end{lemma}
\begin{proof}
$G$ の一般元には相異なる三固有値がある。さもなくば判別式恒等式から、
定理\ref{thm:rank-nine-finite-symmetry}の証明の通り、その中に二次元以上の冪零平面 $N$ があり、
$G$ は $N$ と跡直交する。$N$ が全て階数1以下なら、対称な階数1全体の線型部分空間は
次元1以下なので矛盾する。正則冪零な $X\in N$ を取りJordan基底で $X=E_{12}+E_{23}$ とする。
その基底で計量は反対角単位行列にできる。実際巡回ベクトル $v$ による
$v,Xv,X^2v$ の積対合は $k[t]/t^3$ の完全積跡であり、標準の最上係数との比の単元の平方根で基底を調整できる。
$Y\in N$ への $\operatorname{tr}(X^2Y)=\operatorname{tr}(XY)=0$ は
$Y_{31}=0$, $Y_{21}+Y_{32}=0$ を与え、自己随伴性で両成分は等しいので零である。
ゆえに $N$ は厳密上三角自己随伴の二平面全体であり、$G$ も跡直交と自己随伴性から上三角になって矛盾する。

相異なる三固有値の一般 $G_x$ を直交対角化する。$F_x$ も対角である。
対角固有値をそれぞれ $a_i,b_i$ とし、$k_{ij}=(b_i-b_j)/(a_i-a_j)$ とする。
偏極により全 $y$ で $(F_y)_{ij}=k_{ij}(G_y)_{ij}$ ($i\ne j$) となる。
三つの比が二つでも等しければ全て等しく、跡零性から $F_x$ は $G_x$ の倍である。
一般の $x$ でそうなら線型性により倍率は定数である。

従って三つとも異なる場合を除けばよい。$l_{ij}(y)=(G_y)_{ij}$ とする。
$[G_y,F_y]_{ij}=0$ から $l_{ij}$ のある一次式倍は
$(k_{hj}-k_{ih})l_{ih}l_{hj}$ に等しい（$h$ は第三の添字）。
非零な $l$ の辺グラフは既約性により連結であり、二辺だけなら第三辺のこの等式に反する。
従って全て非零で、各一次式は他の二つの積を割る。三つとも比例し、
$G$ の二次元部分は対角跡零全体になる。そこでは $F$ も対角である。
全非対角成分が非零の一元とこの対角平面を偏極すれば、全ての対角元でも
$F$ の対角差は対応する $G$ の差の $k_{ij}$ 倍である。独立な二差を任意に指定し、
三番目の差が和であることを使うと全 $k_{ij}$ は等しくなって矛盾する。
\end{proof}
\begin{proof}[鉛筆と巡回中心化環]
交換子次元2なので、$\wedge^2\langle Q\rangle$ の核の直線が可換二平面を指定する。
この平面に巡回元がある。なければ、二固有値の半単純元への可換摂動を使う
定理\ref{thm:rank-eight-finite-symmetry}の議論により全元が平方零で、共通像直線または共通核平面を持つ。
そのとき二交換子 $[Q_s,Q_w],[Q_t,Q_w]$ に共通左核または共通右核が生じる。
しかしこれらは $[G_a,F_c]$ の全非零成分を $ac$ 行列として張るため、共通核なら
いずれかの非零な跡零元が相手の既約三平面全体と交換してしまう。

一般の $Q_s$ は巡回で、その中心化環は三次元なので $I,Q_s,Q_t$ の包と一致する。
もし二交換子鉛筆が全て特異なら $\det[Q_s,Q_w]=0$ であり、
定理\ref{thm:rank-eight-finite-symmetry}で示した六元の行列式恒等式により
$Q_s,Q_w$ はある非零対称形式に関して自己随伴である。$Q_t$ も $Q_s$ の多項式であり、
既約性からその形式は可逆である。適切な $U$ 基底で三つの $Q$ を全て対称にする。
この変更は $G,F$ の各三列を線型結合するのでその対称性を保つ。
すると $[Q_b,Q_d]$ は $a,c$ で交代し、交換子式で $[G_x,F_x]=0$ となる。
補題により $g=3$ となって矛盾する。

最後に $C_0$ を可逆に取れば、跡零三行列のどの二固有値も和が零にならない。
したがってSylvester作用素 $L\mapsto LC_0+C_0 L$ は可逆である（上三角化で対角に固有値の和が並ぶ）。
$Q_t=a Q_s^2+bQ_s+cI$ と書き、独立性より $a\ne0$ である。
交換子を取ると一意な $D$ は $a Q_s+bI/2$ であり、$h=1/a$ とすれば表示の二式を得る。
巡回 $D$ に対する $[D,-]$ の像は $I,D,D^2$ の跡直交補である。
$C_0$ は跡零で $2\operatorname{tr}(DC_0)=\operatorname{tr}C_1=0$ だから、残る条件は表示の一つである。
これで $Q_w$ の式も従い、逆に三つの交換子は式から直接指定値に戻る。
\end{proof}

% !TeX root = ../papers/nc-p3.tex
\section{三交換子の判別三次式と巡回行列からの復元}
\label{sec:marked-three-commutators}
\begin{proposition}\label{prop:marked-mixed-discriminant}
単純印を持つ有限窓の $q=3$, $g=4$ の場合に、定理\ref{thm:marked-common-triple-symmetry}の
基底で $P=Q$ とし、跡零対称行列の五空間 $S$ 内の二つの三平面を $G,F$ とする。
$G+F=N^\perp$ と跡対合で書き、$G\cap F=\langle H_1,H_2\rangle$ とする。
交代添字 $s\wedge t$ に対応する三座標を $v=s\times t$ とすれば、固定した非零定数を除いて
\[
          \det[Q_s,Q_t]=
          \operatorname{tr}\!\left([v]_\times[H_1,H_2]\right)\, v^t N v .
\]
従って常に特異交換子の直線があり、恒等零となる必要十分条件は $H_1,H_2$ の可換性である。
重複因子、可約二次式、等方方向を許す。
\end{proposition}
\begin{proof}
$K=[v]_\times$ は外積行列とし、$S$ 上に $\omega_v(X,Y)=\operatorname{tr}(K[X,Y])$ を取る。
交換子式により $[Q_s,Q_t]$ の成分は $G,F$ の順序付き基底での混合行列の成分の一定倍である。
従って $G$ の基底を $H_1,H_2,X$、$F$ の基底を $H_1,H_2,Y$ に替えたとき
行列式は非零定数しか変わらない。混合行列式を展開すると
$\omega_v(H_1,H_2)$ と、$H_1,H_2,X,Y$ 上の $\omega_v$ のPfaffianの積になる。

$\omega_v$ の根基は $K$ と交換する跡零対称行列である。
$0\ne v$ なら $K$ は巡回行列である。実際 $v^tv\ne0$ なら固有値が相異なり、
$v^tv=0$ なら $K^2=vv^t\ne0$, $K^3=0$。従って根基は
$T_v=vv^t-(v^tv)I/3$ の直線である。
体積形式を固定して $\omega_v\wedge\omega_v/2$ を $S$ のベクトルとみなすと
各成分は $v$ の二次式で、根基に入る。$T_v$ の成分に共通因子はない
（非対角の三単項式だけで足りる）。よってそのベクトルは $T_v$ の非零定数倍である。
四平面でのPfaffianは従って $\operatorname{tr}(N T_v)=v^tNv$ の一定倍である。
$N\ne0$ と交代空間での跡対合の非退化性より最後の主張も従う。
\end{proof}

\begin{proposition}\label{prop:marked-three-commutators-reconstruction}
$q=3$, $g>3$ で三つの指定交換子が独立の場合、対称三列 $G,F$ を固定すれば、
適当な有限個の方向を使って、$Q$ の復元を跡零行列 $X$ 一つの二次以下の条件と
線型方程式の整合性で被覆できる。具体的に基底 $s,t,u$ で指定値を
$C=[Q_s,Q_t]$, $D=[Q_s,Q_u]$, $L=[Q_t,Q_u]$ と読む。
\begin{enumerate}
\item $X$ の最小多項式が三次で、$\operatorname{tr}XC=\operatorname{tr}XD
=\operatorname{tr}X^2 C=\operatorname{tr}X^2 D=0$ となる跡零行列を取る。
\item $[X,Y_0]=C$, $[X,Z_0]=D$ の跡零の一解ずつを取る。
$X'=X^2-\operatorname{tr}(X^2)I/3$ とすると、残りの式はただ一つの線型系
\[
 L-[Y_0,Z_0]= aD+b(XD+DX)-cC-d(XC+CX)
\]
である。全ての解から $Q_s=X$, $Q_t=Y_0+aX+bX'$, $Q_u=Z_0+cX+dX'$ を得る。
\end{enumerate}
これは指定した印付き交換子の全解のうち $Q$ が既約なものを覆う。
得られた行列には元の点窓および正則性の条件が別に必要である。
\end{proposition}
\begin{proof}
巡回 $X$ の中心化環は $I,X,X^2$ の張る可換環で、$[X,-]$ の像はその跡直交補である。
$C,D$ の跡は交換子なので既に零である（候補 $G,F$ から始めるなら最初に三つの指定値の
跡を検査する）。これで一解の存在と一般解の表示が従う。
両中心化元の交換子は零、$[X^2,Z_0]=XD+DX$ などを代入すれば表示の線型系となり、
これを解けば三つの指定値を全て復元する。

最後に方向を有限個で選べることを示す。三平面の全元 $X$ が非巡回なら、非零跡零元は
二つの固有値のうち重複固有値の固有空間が二次元以上か、平方零で階数1以下である。
$\operatorname{tr}(X^2)$ が恒等零なら全元が後者となる。それ以外では三平面上の多項式恒等式
$(\operatorname{tr}X^2)^3=6(\operatorname{tr}X^3)^2$ が成立する。
一意分解により非零定数を除いて二次式は一次式の平方、三次式はその立方であり、
重複固有値 $-\operatorname{tr}(X^3)/\operatorname{tr}(X^2)$ は一次式 $\lambda(X)$ に延びる。
従って前者でも $X-\lambda(X)I$ は全三平面で階数1以下である。階数1以下の線型空間は
像直線を共有するか核平面を共有するので、元の三列の既約性に反する。
従って $I,X,X^2$ の独立を検出する三次小行列の一つは $s$ の非零な三次式である。
例えば $s=(1,i,j)$, $0\le i,j\le3$ の十六方向を取ればどれかで非零となる。
各方向は一度だけ基底に補えばよい。
\end{proof}

% !TeX root = ../papers/nc-p3.tex
\section{可換な共通二平面での指定階数一行列}
\label{sec:marked-commuting-intersection}
\begin{lemma}\label{lem:three-bracket-rank-one-compression}
$U$ を三次元、$P=u b^t$ を非零階数一行列とする。
既約な跡零三平面 $\mathcal Q\subset\mathfrak{sl}(U)$ の交換子の包が三次元で、
その全元が $W=\ker b^t$ を $L=ku$ に写すなら
\[
                        P_0=P-\operatorname{tr}(P)I/3\ \in\mathcal Q .
\]
\end{lemma}
\begin{proof}
まず $U=W\oplus L$ の場合、全交換子は $2+1$ ブロックの対角二成分が零である。
$X\in\mathcal Q$ を $\left(\begin{smallmatrix}A_X&p_X\\q_X^t&\alpha_X\end{smallmatrix}\right)$ とする。
最後の対角成分より $(p_X,q_X)$ の像は自然な四次元シンプレクティック空間の等方部分なので次元2以下。
非零の $X$ で $p_X=q_X=0$ とできる。全 $Y$ で $[A_X,A_Y]=0$。
$A_X$ が非スカラーならその二次元中心化環は可換であり、全 $A_Y$ が可換となる。
すると左上の式は全 $Y,Z$ に対して $p_Y q_Z^t=p_Z q_Y^t$ である。
$p,q$ の一方が恒等零なら $W$ または $L$ が共通不変部分空間となる。
それ以外では両方非零となる一元を取り、この等式により $(p_Y,q_Y)$ の像が一次元に限られる。
従って二次元以上のブロック対角部分があり、そこで全元が可換なので交換子の次元3に反する。
よって $A_X$ はスカラーであり、跡零の非零 $X$ が $P_0$ の倍である。

次に $L\subset W$ とし、$P=E_{13}$ に取る。$\mathcal Q$ のパラメータ空間上の一次式を用いて行列を
$\left(\begin{smallmatrix}a&b&c\\d&e&f\\g&h&i\end{smallmatrix}\right)$ と書く。
交換子の左下二行二列（第2,3行、第1,2列）の零は
\[
 d\wedge(a-e)+f\wedge g=0,\quad
 g\wedge(a-i)+h\wedge d=0,\quad
 g\wedge b+h\wedge(e-i)=0,\quad d\wedge b+f\wedge h=0 .
\]
$g\ne0$ なら二番目の式から $d,g,h$ の包の次元は高々2。
次元1なら $d,h$ は $g$ の倍であり、旗を保つ基底変更で $d=h=0$ とする。
実際 $I+tE_{23}$ による共役は $d$ に $tg$ を加え、$I+t'E_{12}$ による共役は
$h$ から $t'g$ を引くのでこの変更ができる。その基底では四式より $f,b,a-i$ も $g$ の倍なので $\ker g$ の二平面は
$\operatorname{diag}(a,e,a)+cE_{13}$ の形で可換となり矛盾する。
従って次元2である。$d,g,h$ を全て消す非零のパラメータ $s$ で四式を縮約すると
$a_s=i_s$, $f_sg=(a_s-e_s)d$, $b_sg=(a_s-e_s)h$ となる。
二次元性から $a_s=e_s$、さらに $f_s=b_s=0$。従って $Q_s$ は $E_{13}$ の非零倍である。
$g=0$ なら $h\wedge d=0$ で、既約性より両方非零だから $h=kd$, $k\ne0$。
残りの式より $a-e,e-i,b-kf$ は $d$ の倍である。$\ker d$ の二平面は
$kE_{12}+E_{23}$ と $E_{13}$ の包となって可換であり、再び交換子の次元に反する。
\end{proof}

\begin{theorem}\label{thm:marked-commuting-intersection}
命題\ref{prop:marked-mixed-discriminant}の三次式が恒等零のとき、
$H=G\cap F$ に対して $W=\{c:F_c\in H\}$、$L=\{a:G_a\in H\}^\perp$ とする。
$L$ を像、$W$ を核とする階数一行列 $P$ は定数倍を除いて指定される。
$\langle Q\rangle$ は $P_0$ を含み、その方向 $[s]$ は次の\emph{既知の射影直線}上にある。
指定交換子写像を $B:\bigwedge^2V\longrightarrow\mathcal C$ とする。
$\operatorname{im}[P,-]$ 上の交代形式を
\[
      \Omega([P,Y],[P,Z])=\operatorname{tr}(P[Y,Z])
\]
とし、$\mathcal C$ への制限の根基直線を $D$ とする。
求める直線は $B^{-1}(D)$ の非零二ベクトルの張る二平面の射影化である。
さらに $s$ は $H$ の共通固有方向を取り得ない。
\end{theorem}
\begin{proof}
$g=4$ なので既出の二交換子の次元判定から $B$ は三次元間の同型である。
$a\in L^\perp,c\in W$ に対し $[G_a,F_c]=0$ なので、全ての指定交換子は $W$ を $L$ に写す。
跡零でこの条件を持つ空間は四次元であり、$[P,-]$ の像と等しい。
実際この像はその中にあり、$P$ の中心化環は $P$ が半単純でも平方零でも五次元である
（それぞれ上の二つの基底で読める）。補題より $Q_s$ は $P_0$ の非零倍に選べる。
$\Omega$ は跡の巡回性により代表によらず、$\operatorname{tr}([P,Y] Z)$ の完全対合から
非退化である。$B(s\wedge V)$ 上では $\operatorname{tr}(P[Q_t,Q_u])=0$ のため
$\Omega=0$（全指定値は $P$ と跡直交する）。四次元シンプレクティック空間内の超平面
$\mathcal C$ の根基は直線であり、その中の等方二平面は必ずこの根基を含む。
$B$ で戻せば $s$ の位置の主張を得る。

最後に $Q_s$ が $W$ 上に作用するスカラーを $\lambda$ とする。
もし $s$ が $H$ の共通固有方向なら、$c\in W$ ごとに $F_c s=\mu(c)s$ で
$\mu$ は一次式である。辺 $12$ の二つの点を $(-\lambda,s)$ と $(-\mu(c),c)$ に取ると
六つの独立な式は $(Q_s-\lambda I)c=0$, $(F_c-\mu(c)I)s=0$ として全て成立する。
これは定理\ref{thm:marked-point-commutators}の一方に星印のある六行の式であり、
$F_c$ が対称なので左右は同じである。$[c]\in\mathbf P(W)$ が動くため辺の有限性に反する。
\end{proof}

\begin{theorem}\label{thm:no-three-bracket-compression}
既約な跡零三平面の三つの交換子が独立なら、その全てが固定二平面を固定直線に送ることはない。
従って前定理の可換な共通二平面の場合は有限窓に存在しない。
さらに $q=3$, $g>3$ で交換子が三独立なら、$g=5$ でも判別三次式は恒等零にならない。
\end{theorem}
\begin{proof}[階数一行列を含んだ後の排除]
補題により $P_0\in\mathcal Q$。まず $W\oplus L$ の場合に、
その補基底 $Y,Z$ の左上二行二列を $A,B$、右上列を $p,r$、左下行を $q^t,t^t$ とする。
$[Y,Z]$ の対角の消失より
\[
               q^t r=t^t p,\qquad [A,B]=r q^t-p t^t .
\]
$[P,Y],[P,Z]$ は独立であり、$(p,q),(r,t)$ は二平面を張る。
$p,r$ が独立なら $q=Sp,t=Sr$ と書け、第一式で $S$ は対称である。
$S=0$ なら $W$ が不変になる。$J=r p^t-p r^t$ は可逆な交代行列で $[A,B]=JS$。
$A,B$ はそれぞれその交換子と跡直交するので、$SA,SB$ は対称である。
$S$ が特異なら $\ker S\subset W$ は共通不変直線となる。
可逆なら $P_0,Y,Z$ は計量 $\operatorname{diag}(S,1)$ に関して自己随伴である。
しかしこの計量に関する交代行列で両対角ブロックが零のものは二次元しかなく、
三交換子の独立性に反する。$q,t$ が独立なら転置して同じである。
どちらの射影も階数1なら、二基底を取り直して
$(p,q)=(p,0)$, $(r,t)=(0,t)$, $p,t\ne0$ とできる。$t^tp=0$ で
$[A,B]=-pt^t$。跡直交により $A,B$ は $\ker t^t=kp$ を保ち、
三行列全てがこの直線を保つ。射影が零の場合も一方のブロックが不変である。

$P=E_{13}$ の場合は補題の四式の記号を使う。既に $g\ne0$ かつ
$d,g,h$ の包は二次元である。$g_Y=1$, $g_Z=0$ として $P_0,Y,Z$ を基底に取る。
旗を保つ先の二つの共役で $d_Y=h_Y=0$ とし、$I+vE_{13}$ による共役で
$a_Y=i_Y$ にする。この最後の変更は $d_Y,h_Y$ を保ち、$a_Y-i_Y$ に $2v$ を加える。
$D=d_Z$, $H=h_Z$ は同時には零でない。四式を $Y,Z$ で評価すると
\[
 a_Z=i_Z,\qquad f_Z=D(e_Y-a_Y),\qquad b_Z=H(e_Y-a_Y),
 \qquad D b_Y=H f_Y .
\]
$D=0$ なら $d=f=0$ で $\langle e_1,e_3\rangle$ が不変、$H=0$ なら $h=b=0$ で
$ke_2$ が不変である。両方非零なら $j=H/D$ として三行列とも
\[
                  M=\begin{pmatrix}0&0&1\\0&j&0\\1&0&0\end{pmatrix}
\]
に関して自己随伴である。実際 $h=jd$, $b=jf$, $i=a$ が必要十分で、$c$ は自由である。
一方、左下二行二列が零の跡零行列で $M$ に関して交代となるものは二次元である。
これで両方の配置で三つの指定値は実現できない。
\end{proof}

\begin{proof}[五次元和を含む判別式]
三独立な指定値の空間 $\mathcal C$ に共通右核 $c$ があれば $F_c$ が全 $G_a$ と交換し、
既約性と跡零性で零になり $F$ の独立性に反する。共通左核もない。
三次式が恒等零なら、このような三次元特異行列空間は固定二平面を固定直線に送るものか、
可逆な左右変換で交代三行列の空間になるものである。ここではこの二択を直接確かめる。

階数2の元を左右変換で $A=\operatorname{diag}(1,1,0)$ とし、補二変数を $x$ とする。
全元が階数1以下の場合は共通核または余核があるため除いた。行列式の $A$ での一次式により
補元は $\left(\begin{smallmatrix}H(x)&p(x)\\q(x)^t&0\end{smallmatrix}\right)$ である。
$\det(tA+M(x))=0$ から $q^tp=q^t\operatorname{adj}(H)p=0$。
$p$ が階数1なら $p=w\alpha$、$q$ は $w$ を消す固定行の倍となる。両者とも恒等零にはなれない。
$w=e_1$ とすれば第二式から $H_{21}=0$ で、全元が $\langle e_1,e_3\rangle$ を $ke_1$ に写す。
$p$ が階数2ならそれを $x$ 自身に取り、$q$ は固定可逆交代行列 $J$ で $Jp$ に書ける。
第二式から $Hp=\lambda(x)p$ であり、二成分が互いに素なので $\lambda$ は一次式である。
$H-\lambda I$ の各行は $p$ を消す一次式の行だから $H=\lambda I+u q^t$、$u$ は定数列である。
補元から $\lambda A$ を引き、最終行の倍を上二行から引けば $H=0$。
上二行に $J$、最終行に $-1$ を掛けると全行列が交代となる。

第一の型は定理前半で生じない。第二の型なら元の空間は $\{KB:K^t=-K\}$、
$B$ 可逆と書ける。全行列の跡が零のため $B$ は対称で、$U$ の基底変更により
$\mathcal C$ を交代空間にできる。この基底変更は $G,F$ の対称行列をそれぞれ線型結合で取り直す。
指定値が $a,c$ で交代であることから $[G_x,F_x]=0$ が恒等的に成立する。
補題\ref{lem:commuting-linear-symmetric-nets}により $F_x$ は $G_x$ の一定倍で $g=3$ に反する。
以上で判別式が恒等零となる場合を全て除いた。
\end{proof}

% !TeX root = ../papers/nc-p3.tex
\section{可逆交換子を持つ鉛筆の復元階数と例外直線}
\label{sec:marked-bracket-pencil-rank}
\begin{theorem}\label{thm:marked-bracket-pencil-rank}
三独立な指定交換子の一般鉛筆 $N=B(s\wedge V)$ を取り、$N$ に可逆元があるとする。
$X=Q_s$ の復元における四係数の線型系
（命題\ref{prop:marked-three-commutators-reconstruction}）の階数は3または4である。
階数3のとき $\{X,K\}=XK+KX\in N$ となる方向 $[K]\in\mathbf P(N)$ はただ一つである。
特異な $K$ は $\det|_N$ の重根に限る。可逆な $K$ が例外方向になる場合には、
全交換子空間 $\mathcal C$ 上の既知の二次式
\[
 \mathcal F(K\wedge T)=(\operatorname{tr} KT)^2-2\operatorname{tr}(KTKT)
\]
の零点が直線 $\mathbf P(K\wedge\mathcal C)$ を含む。
特に $\det|_{\mathcal C}$ が被約三次曲線を定め、$\mathcal F$ が恒等零でなければ、
一般の $s$ で全ての復元解の線型系は階数4になる。
\end{theorem}
\begin{proof}
$C=[X,Y]$ を可逆に選ぶ。$X$ が非巡回ならスカラーを引いて階数1以下であり、
$[X,Y]$ の階数は2以下なので、$X$ は必ず巡回である。
$T_X(C)=XC+CX$ と書くと線型系の像は $N+T_X N$ である。
$C,T_X C,T_X^2 C$ は常に独立である。実際 $T_X C=[X^2,Y]$、
$T_X^2 C=[X^3,Y]+XCX$ で、Cayley--Hamiltonにより $[X^3,Y]$ は $C$ の倍である。
従って従属なら $a C+b(XC+CX)+cXCX=0$ となる。
$c=0$ なら $b\ne0$ であり、$C$ の共役で $X$ が $-X-aI/b$ に移る。
$c\ne0$ なら $T=X+(b/c)I$ として $TCT=\mu C$。
$\mu=0$ では可逆な $C$ により $\operatorname{rank}T\le1$ となって巡回性に反する。
$\mu\ne0$ では $T$ は可逆で、共役で $T$ が $\mu T^{-1}$ に移る。
両場合とも $C$ は可換環 $k[X]$ を保つ位数2以下の自己同型 $\sigma$ を誘導する。
巡回基底で $U=k[X]$ とすれば $C$ は $\sigma$ と単元 $v$ の乗法の合成 $m_v\sigma$ である。
しかし $C=[X,Y]$ は全 $k[X]$ と跡直交するのに
$\operatorname{tr}(m_{v^{-1}}C)=\operatorname{tr}\sigma\ne0$。
最後は三次元で対合の固有値が $\pm1$ だからである。従って $N$ が $T_X$ で閉じることはなく、
階数3ならその欠けた一方向 $K$ は一意である。

$K$ が特異で単根と仮定する。$K=[X,Y']$ と書き、$N$ を補うもう一元を $D=[X,Z]$ とする
（ここで $D$ は $N$ 内で $K$ の補元）。$\operatorname{tr}(\operatorname{adj}K D)\ne0$。
ところが $\operatorname{adj}K$ は $T_X K$ とも $K$ とも跡直交するので、
$T_X K\in N$ は $K$ の倍となる。これより $[X,K^2]=0$ で、跡零のため
$\operatorname{adj}K=K^2-\operatorname{tr}(K^2)I/2$ も $X$ と交換する。
従って $[X,Z]$ と跡直交する矛盾となる。

可逆な $K$ の場合は基底を取り直して $K=[X,Y]$, $D=[X,Z]$, $L=[Y,Z]$ とする。
$T_X K=\alpha K+\beta D$ で $\beta\ne0$。零なら $K,T_X K$ が従属で
上の独立性に反する。$[X,\{X,Y\}]=\{X,K\}$ と中心化環の記述から
$Z$ は $\beta^{-1}(\{X,Y\}-\alpha Y)$ と $X$ の二次以下の多項式の和である。
従って $L$ は $-\beta^{-1}\{Y,K\}$ と $K,D$ の線型結合の和である。
つまり全 $\mathcal C$ が $T_K(\langle I,X,Y\rangle)$ に入る。
$P,Q\in\langle I,X,Y\rangle$ について $\operatorname{tr}(K(PQ+QP))=0$ である。
$X^2,Y^2$ への跡直交は交換子の巡回性から、
$XY+YX$ への直交は $\operatorname{tr}([X,Y](XY+YX))=0$ による。
従って $\operatorname{tr}(T T_K^{-1}T)=0$ が全 $T\in\mathcal C$ で成立する。

可逆な跡零 $K$ と跡零 $T$ については
\[
 2\det K\,T_K^{-1}T=\operatorname{tr}(K^2T)I+
                         \operatorname{tr}(KT)K-2KTK .
\]
相異なる固有値の対角基底で非対角成分は $\lambda_i+\lambda_j=-\lambda_k$、
対角成分はスペクトル射影の二次式から得られ、分母を払えば全元で成立する。
$T_K$ 自体の可逆性は二固有値の和が零にならないためである。
これで $2\det K\operatorname{tr}(T T_K^{-1}T)=\mathcal F(K\wedge T)$。
この表示が外積の二次式であることは
\[
 3\mathcal F(K\wedge T)=(\operatorname{tr} KT)^2-
 \operatorname{tr}(K^2)\operatorname{tr}(T^2)-2\operatorname{tr}([K,T]^2)
\]
から従う。右辺の恒等式は $\operatorname{tr} M^4=(\operatorname{tr}M^2)^2/2$ の偏極である。
非零二次式の零点に含まれる直線は高々二つなので、可逆な例外方向の候補は
鉛筆を動かしても高々二点である。$N$ は $\mathcal C$ の全鉛筆を動くので、
被約三次曲線との横断直線を取り、さらにこれらの候補を避ければ最後の結論を得る。
\end{proof}

\begin{proposition}\label{prop:invertible-defect-bracket-parameters}
上の可逆な例外方向 $K$ を固定する。$K=[X,Y]$, $D=[X,Z]$, $L=[Y,Z]$ を指定値とし
$U=T_K^{-1}D$, $V=T_K^{-1}L$, $u=\operatorname{tr}U/3$, $v=\operatorname{tr}V/3$ と置く。
復元解が存在する必要十分条件は $[U,V]$ が $K$ の非零倍であることである。
$K=-\beta^2[U,V]$ の二つの根と任意の $h\in k$ を取り、跡零部分を下添字0で表すと、
この方向を持つ全解は
\[
 X=\beta U_0,\qquad Y=-\beta V_0+hX,\qquad
 Z=\left(\beta^{-1}\{X,Y\}+2uY-2vX-(h/\beta)X^2\right)_0 .
\]
これは指定した三交換子の段階での必要十分条件である。
\end{proposition}
\begin{proof}
$T_X K=\alpha K+\beta D$ に $T_K^{-1}$ を施して跡零にすると $X=\beta U_0$、
$\alpha=-2\beta u$ である。
前の証明での一般の $Z$ を
$\big(\beta^{-1}(\{X,Y\}-\alpha Y)+c_1X+c_2X^2\big)_0$ とすれば、
第三の交換子は $L=-\beta^{-1}\{Y,K\}-c_1 K-c_2\{X,K\}$。
従って $Y=-\beta V_0+hX$、$c_2=-h/\beta$, $c_1=-2v$ が必要である。
第一交換子を戻せば表示の平方根の条件を得る。逆にこれらの式を順に代入すれば
$[X,Y]=K$, $[X,Z]=D$, $[Y,Z]=L$ であり、指定した $K$ が例外方向となる。
\end{proof}

% !TeX root = ../papers/nc-p3.tex
\section{一般方向にも残る交換子復元の階数降下}
\label{sec:marked-persistent-bracket-defect}
\begin{lemma}\label{lem:nine-symmetric-words}
跡零三行列 $X,Y$ で $\det[X,Y]\ne0$ とする。$S=XY+YX$ と置くと
\[
                I,X,Y,X^2,Y^2,S,XYX,YXY,S^2
\]
は全三行列の基底である。
\end{lemma}
\begin{proof}
行優先で九列を並べた行列式は $-\det[X,Y]^3$ であることを示す。
まず $X$ を相異なる三固有値 $x_i$ の対角行列とする。$\det[X,Y]$ はVandermondeの
非零倍と $f=y_{12}y_{23}y_{31}-y_{13}y_{32}y_{21}$ の積である。
$f$ は既約で、$f=0$ の一般点では全非対角成分が非零となり、
$x_i$ の基底で対角な可逆計量に関して $Y$ も自己随伴になる。
九語は全て反転で不変なので計量に関して自己随伴となり、列階数は6以下である。
従って高さ1の $f$ で局所化して行列を掃き出すと行列式は $f^3$ で割れる。
$Y$ に関する全次数は9なので、商は $Y$ に依存しない。

$Y=E_{12}+E_{23}+E_{31}$ で商を計算する。対角、巡回一段、巡回二段の各三列の係数は、
巡回添字を使い、順に $(1,x_i,x_i^2)$, $(1,x_i+x_{i+1},x_ix_{i+1})$,
$(1,x_{i+1},(x_i+x_{i+1})(x_{i+1}+x_{i+2}))$ である。
跡零性から最後の二群の二次係数はそれぞれ $x_{i+2}^2-\operatorname{tr}(X^2)/2$ と
$x_{i+1}^2-\operatorname{tr}(X^2)/2$ になる。従って三つとも同じVandermondeの符号倍で、
$(x_1,x_2,x_3)=(1,2,-3)$ で順序の符号を決めると表示の恒等式を得る。
共役の全行列空間上の行列式は1であり、対角化可能な稠密開集合から全行列に延びる。
\end{proof}

\begin{theorem}\label{thm:persistent-three-bracket-defect}
既約な跡零三平面 $\mathcal Q$ の交換子写像 $B$ が三独立で、像が可逆元を含むとする。
一般の $X\in\mathcal Q$ でも前節の復元の階数が3に留まる必要十分条件は、
$U=k^2\oplus k$ の適当な基底で
\[
 \mathcal Q=\left\{\begin{pmatrix}aI_2&w\\ w^t S&-2a\end{pmatrix}
               :a\in k,\ w\in k^2\right\},\qquad S\in\mathrm{GL}_2,\quad S\ne S^t
\]
となることである。他の全ての三平面では一般方向の線型復元は階数4である。
表示の場合、判別三次式は非零定数を除いて一つの一次式の三乗であり、
$\mathbf P(kI+\mathcal Q)$ の階数一切断は平方零の平面三重点である。
\end{theorem}
\begin{proof}[二次積と特異方向]
$X$ の鉛筆 $N=[X,\mathcal Q]$ が可逆元を含む開集合で考える。
例外方向 $K=[X,Y]$ と補基底 $Z$ に対し $\{X,K\}=\alpha K+\beta [X,Z]$ とする。
$X$ は巡回であり、$[X,\{X,Y\}]=\{X,K\}$ だから
$\{X,Y\}-\alpha Y-\beta Z$ は $X$ の二次以下の多項式である。
$Y$ に $X$ の適当な倍を足した $Y'$ を使えば $\{X,Y'\}\in kI+\mathcal Q$。
従って $XY'$ は積を対称化する写像
$\mathrm{Sym}^2\mathcal Q\to\operatorname{End}(U)/(kI+\mathcal Q)$ の核に入る可約二次式である。

可逆な $K$ を持つ例外方向が一つでもあれば $\beta\ne0$ であり、
$\mathcal Q=\langle X,Y',\{X,Y'\}_0\rangle$ である。
このとき先の対称積の核はちょうど $k(XY')$ である。
実際補題の九語で $I,X,Y',\{X,Y'\}$ に
$X^2,Y'^2,\{X,\{X,Y'\}\},\{Y',\{X,Y'\}\},\{X,Y'\}^2$ を加えるとやはり基底となる。
$\{X,\{X,Y'\}\}$ は $XY'X$ に低次の $I,X,Y'$ の線型結合を加えたもので、
$X$ と $Y'$ を交換しても同じである（跡零三行列のCayley--Hamiltonの偏極）。
よって例外になり得る方向は二次式の二因子 $[X],[Y']$ だけで、一般方向は階数4である。

一般に階数3が残るとすれば、一般の鉛筆の例外は特異な $K$ である。
前節よりこれは判別三次曲線との重交点である。一般の直線は被約台の滑らかな点で横断するので、
三次式には唯一の重複成分である直線がある。この直線を $\mathbf P([T,\mathcal Q])$ とする
固定方向 $[T]\in\mathbf P(\mathcal Q)$ があり、一般の $X$ で $K=[X,T]$ としてよい。
ここで $\beta\ne0$ なら $\mathcal Q$ は $X,T$ およびその対称積と $X^2$ から生成される。
$\det[X,T]=0$ のため二行列には非零の共通対称形式がある
（定理\ref{thm:rank-eight-finite-symmetry}の六式の行列式恒等式）。
その核は両行列に不変で従って $\mathcal Q$ に不変、よって形式は可逆である。
全 $\mathcal Q$ がその形式に関して自己随伴になり、全交換子が特異になって矛盾する。
ゆえに全 $X$ で多項式恒等式
\[
                    [X^2,T]=\lambda(X)[X,T]
\]
が成立する。倍率が一次式であるのは $[X,T]$ の成分が二つ以上の独立な一次式を張り、
共通因子を持たないからである。$X=T+tY$ として一次項を比べれば
$[Y,T^2-\lambda(T)T]=0$。既約性により $T^2-\lambda(T)T$ はスカラーであり、
$T$ は非スカラー非巡回なので $P_0$ の非零倍、$P$ は階数1である。
\end{proof}

\begin{proof}[階数一方向に沿うブロック]
スカラーを引いても不変部分空間は変わらない。$P=E_{13}$ なら
$X-\lambda(X)I/2$ の平方が $P$ と交換する。
その九成分を前節までと同じ順に $a,b,c,d,e,f,g,h,i$ と書くと
$g(a+i)+hd=0$, $d(a+e)+fg=0$, $h(e+i)+gb=0$ である。
$g=0$ なら $d,h$ の一方が恒等零で、直線または平面が不変となる。
$g\ne0$ なら $g$ が $hd$ を割るので、必要なら反対角転置で双対に移って $d$ が $g$ の倍とできる。
$P$ を保つ旗の基底変更で $d=0$ とすると $f=0$ となり、$\langle e_1,e_3\rangle$ が不変である。

半単純 $P$ では $X=\left(\begin{smallmatrix}A&p\\q^t&b\end{smallmatrix}\right)$ とし、
$l=\lambda(X)-b$ とすると $A p=l p$, $q^t A=l q^t$。
$p,q$ の一方が恒等零なら可約なので、$A-lI$ の全元の階数は1以下である。
この空間が非零で共通像直線を持てば、$q$ は常にこの像を消し、$A$ は像直線を保つため
共通不変直線がある。共通核直線なら $p$ がそこに入り、最後の軸と核直線の張る平面が共通不変である。
よって $A=lI$ が恒等的に成立する。$P_0$ を使って他の二基底は両対角が零にできる。

この二元の交換子 $L$ はブロック対角であり、他の二交換子はブロック非対角である。
$zL+C$ の行列式は非対角部分 $C$ の符号に不変で、$z^3$ と $z$ の項だけを持つ。
$z=0$ は先に得た重複直線なので $\det L\ne0$ である。
従って左上ブロック $p_1 q_2^t-p_2 q_1^t$ も可逆で、$p_1,p_2$ と $q_1,q_2$ はそれぞれ独立。
これで表示の $w^t S$ が得られ、$L$ の最後の成分が非零のため $S$ は非対称である。
逆に表示では $P=E_{33}$ とすると $X^2-(a-2a)X$ はブロック対角で、
$\{X,[X,P]\}$ は $[X,P]$ の倍である。二つの座標方向の交換子は可逆で
（左上は可逆交代行列と $S$ の積、右下は $S_{12}-S_{21}$ の符号倍）、
$P_0$ との二交換子も独立である。表示の三平面は既約である。
実際 $P$ により不変空間は二ブロックの部分空間の和で、二組の非対角成分はいずれも二次元を張る。
よって一般方向に階数3が生じる。上の行列式の議論から三次式は三重直線となる。

最後に $kI+\mathcal Q$ で階数1なら幾何点は $P$ のみである。
右下を1とする近傍では式は $cI_2=w w^t S$ で、対称二行列から
$\mathrm{Mat}_2/kI$ への $M\mapsto MS$ は $S^{-1}$ が非対称のため同型である。
従って $w$ の三つの二次単項式と $c$ が零になり、主張の長さ3局所環を得る。
\end{proof}

\begin{corollary}\label{cor:no-persistent-defect-finite}
単純印を持つ片側対称有限窓で $q=3$, $g>3$、三交換子が独立なら、
復元の線型系は一般方向で必ず階数4である。対称三平面の和の次元4,5の両方を含む。
\end{corollary}
\begin{proof}
前定理の模型とし、$Q_s$ が $P_0$ の倍となる方向を $s$ とする。
$W=k^2$ を上のブロック、$A$ を最後の軸を消す二次元の行空間とする。
$G'=G_A$, $F'=F_W$ はそれぞれ二平面である。
$\det B(v)$ の三重直線を $\ell=0$ とすれば、その上の全行列は非対角ブロックだけを持つ。
ここで $B(v)$ は $v=x\times y$ に対応する指定値 $[Q_x,Q_y]$。
この直線上では列・行の各ブロックが二変数の同型となり、全点で階数2、左核は $A$ 内にある。
ブロックの一成分が単元の近傍で掃き出せばこの核の包含はスキーム上でも成立する。

$Z_G=\{[v]: vv^t\in kI+G\}\subset\mathbf P^2$ とする。
単独の四空間の正次元階数一切断は母線だけで、対称行列の二次Veroneseには母線がないので
$Z_G$ は有限である。二つの二次式の完全交叉ゆえ長さ4である。
$Z_G$ 上で $T_v=vv^t-(v^tv)I/3$ は $G$ の非零元であり、外積行列 $[v]_\times$ と交換する。
従って全 $F_c$ に対し $\operatorname{tr}([v]_\times[T_v,F_c])=0$。
$G$ 基底での $T_v$ の係数が局所的に単元成分を持つ左核を与えるので、
$\det B(v)$、従って $\ell^3$ が $Z_G$ 上で恒等零になる。
直線でさらに切ると左核は $A$ にあるため、$Z_G\cap V(\ell)$ 全体で
$vv^t\in kI+G'$ となる。$\ell$ の冪での三段のフィルトレーションから
$4\le3\operatorname{length}(Z_G\cap V(\ell))$ なので、後の切断は長さ2以上である。

これで $G'$ が可約と分かる。実際その平面 $\mathbf P(kI+G')$ には
Veronese上の長さ2部分の包が入る。相異なる二点なら対応する行列の像は二列の張る二平面にあり、
接線なら $uu^t,u v^t+v u^t$ の像はやはり二平面にある。
どちらもその包には単位行列がないから、これに $I$ を加えると $kI+G'$ 全体となり、
像の二平面を保つ。対称性よりその直交直線 $kn$ も保たれる。

$A,W$ の変数の交叉交換子は模型の左上ブロックから
$[G_a,F_c]=m(a,c)[s]_\times$ と書け、二次元の対合 $m$ は完全である。
実際 $[Q_s,Q_t]$ の左上は全て零、残る二つの非対角生成元の交換子の左上は可逆である。
従って $a$ に $m(a,c)=0$ の方向を対応させる二空間の線型同型 $a\mapsto c(a)$ がある。
一般の $G_a$ は巡回である。全二平面が非巡回なら、
命題\ref{prop:marked-three-commutators-reconstruction}の判別式の議論により
一次スカラーを引いた対称行列が全て階数1以下となる。しかしそのような対称線型空間は一次元以下である。
$F_{c(a)}$ は $G_a$ と交換するためその多項式で、$kn$ を保つ。
稠密性で全 $F'$ が $kn$ を保ち、従って全ての交叉交換子は $n$ を消す。
非零の $[s]_\times$ の核は $ks$ なので $n$ は $s$ に比例する。
$Q_s$ は $W$ 上でスカラーであり、全 $c\in W$ で $F_c s$ は $s$ の倍である。
定理\ref{thm:marked-commuting-intersection}の最後の辺 $12$ の六式を使うと、
$\mathbf P(W)$ の全直線が同じ隣接点の上に入り、有限性に反する。
\end{proof}

% !TeX root = ../papers/nc-p3.tex
\section{指定した判別直線と二次曲線を持つ階数十の正則変形}
\label{sec:marked-prescribed-discriminant-as}
\begin{theorem}\label{thm:marked-prescribed-discriminant-as}
三次元内積空間 $V$ において任意の非零跡零対称行列 $N$ と任意の射影直線 $\ell(v)=0$ を指定する。
単純印付き $q=3$, $g=4$ の四窓に、判別三次式が定数倍を除き
\[
                       \ell(v)\,v^tNv
\]
となる十九次元の既約族があり、その幾何学的一般点は四辺有限、階数10の標準Koszul左右AS正則窓である。
二次曲線の階数、$N$ のJordan型、直線との接触や成分の一致に例外はない。
$[N],[\ell]$ の射影直交群内の同時安定化群の次元を $d$ とすると、
この族の添字保持同型類の像の次元は $9-d$ である。
また $N,\ell$ をともに一般に動かした場合には、同型類の像の次元は12になる。
\end{theorem}
\begin{proof}[混合正則点でのエタールな復元]
$S=\mathrm{Sym}_0(V)$, $D=N^\perp\subset S$ とする。順序付き三列 $G$ を $D$ 内で選び、
三平面 $G'$ の写像 $\bigwedge^2G'\to\mathfrak{so}(V)$ が可逆となる開集合を取る。
実際にここで使える正則な基点は後で全 $N$ について示す。
$H\subset G'$ を、$[H_1,H_2]$ の跡対合による一次式が $\ell$ を定める一意な二平面とする。
これは三次元の二ベクトルが全て可分解であるためで、$G$ を動かすと代数的に動く。
$H\subset F'\subset D$ となる三平面を選ぶ自由度は $\mathbf P^1$ である。
その順序付き基底 $F_i$ に課す式は $\sum_i[G_i,F_i]=0$ である。
これがちょうど指定三交換子の跡零条件である。
$F'=G'$ では写像 $(F_i)\mapsto\sum_i[G_i,F_i]$ が全射なので、
近傍の三式の核は階数6の束となる。その中で基底となることは開条件で
$F=-G$ が使える。得られた滑らかな既約基底の次元は $12+1+6=19$ であり、
$F'\ne G'$ なら $G'\cap F'=H$, $G'+F'=D$ である。

$F=-G$ では三つの指定値は全て交代行列で、その三列は独立である。
ここでいう指定値は $b,d$ を固定し $a,c$ の行列としたものである。
外積行列の二交換子は二列の外積なので、$Q_v=[Lv]_\times$ と取れる。
実際 $\operatorname{cof}(L)=C$ は可逆三行列 $C$ に対し
$L=\sqrt{\det C}\,C^{-t}$ で解ける。$C$ は指定値の外積ベクトルを並べた行列である。
三交換子写像 $\mathfrak{sl}_3^3\to\mathfrak{sl}_3^3$ はこの点でエタールである。
基底を取り直して $Q_i=-[e_i]_\times$ で確かめれば十分で、
微分 $(T_i)\mapsto([T_i,Q_j]+[Q_i,T_j])_{ij=23,31,12}$ の行列式は
本文末の再現計算の基底順で $6912$ である。
この写像を上の十九次元基底に沿って引き戻すと、指定解を通る唯一の局所成分は滑らかで基底に支配的である。
その既約開部分が所要のパラメータ族を与える。一般点での三次式は
命題\ref{prop:marked-mixed-discriminant}による。

次の段で混合点の辺の有限性を確認すれば、定理\ref{thm:mixed-point-frobenius-algebras}により
そこで全次数正則性と階数10が成立する。符号 $F=-G$ は偶二補空間の基底を
互いに逆の平方根で調整して $F=G$ にでき、巡回順を替えれば同定理の交代・対称表示である。
二つの隣接切断の階数6は交換子式による。一文字階数、三・四次交叉、五次双対消滅は
混合点の近傍で標準値を保つ。下界は与えた四テンソルとその縮約、
上界は関係挿入行列の階数による。点辺の有限性も射影族で開であり、
星印は全族で六関係の一次部が独立な単純点を与える。
反対階数は基点で10、近傍で片側対称性により10以下である。
この近傍の有限Frobenius双対の相対bar複体を内部次数ごとに取ると、
非対角ホモロジーの消滅は各次数で基点を含む開条件である。
幾何学的一般点では全次数で消滅し、双対の完全対合により左右AS条件も従う
（定理\ref{thm:reversible-component}の一般点の証明）。
これは十九次元族の全閉点での一様な正則性判定を与える主張ではない。
\end{proof}

\begin{proof}[各跡法線での有限混合点]
半単純な $N$ は直交対角化できる。$G_i=E_{jk}+E_{kj}$（巡回順）を選べば $G_i\in D$ で、
$F=-G,Q_i=-[e_i]_\times$ は定理\ref{thm:twenty-point-distinct-companion}の二十点例を
一位置巡回した星印での表示である。実際同定理の $B_+^t B_-$ の四因子を $(v,x,y,u)$ の順で読むと、
星印一個の係数は順に $D_i,-D_i,-[e_i]_\times,-[e_i]_\times$ となる
（ここで $D_i=E_{jk}+E_{kj}$）。よってこの場合は一般点でも点スキームは被約二十点となる。

非半単純なら、相似と非零スカラー倍でJordan型は平方零階数1、正則冪零、非零重複固有値の二次ブロックの
三通りである。自己随伴行列の相似は計量を保つ相似にできる。標準計量で二つの対称行列を
$T$ で共役すれば $T^tT$ が元の行列と交換するので、その可逆な多項式平方根で $T$ を補正すればよい。
従って反対角単位行列 $J$ を計量にし
\[
 H(a,b,c,d,e)=\begin{pmatrix}a&b&c\\d&-2a&b\\e&d&a\end{pmatrix}
\]
とした次の三代表で尽きる。表では $N$ と $G_i$ を引数の五組で記す。
\[
\begin{array}{c|c|c|c}
N&G_1&G_2&G_3\\ \hline
(0,0,1,0,0)&(-2,-1,2,-2,0)&(-2,-2,1,-2,0)&(2,0,2,-2,0)\\
(0,1,0,0,0)&(2,2,0,0,1)&(-2,-1,-2,0,-1)&(2,-2,1,0,0)\\
(1,0,1,0,0)&(2,12,6,-6,-12)&(0,0,12,-12,0)&(-1,12,0,-6,6)
\end{array}
\]
各々で $\operatorname{tr}(NG_i)=0$ で、$C$ の行列式は順に $-576,18,-1306368$。
ここで $C$ は巡回 $bd$ に対する $(J[G_a,G_c])_{bd}$ を $ac$ 交代行列として外積ベクトルにした三列である。
$L_0=\operatorname{adj}(C)^t$ は $L$ の非零倍であり、星座標の倍率でこの差を吸収できる。
辺の係数行列を点 $(z,v)$ で評価すると、列変数を $(c_*,c)$ として
\[
          \begin{pmatrix}0&zI+[L_0v]_\times\\v&-G_1v\quad -G_2v\quad -G_3v\end{pmatrix}
\]
となる。階数は全点で3以上である。$v\ne0$ なら上の三行の階数は2以上
（非零外積行列は巡回）、下に第一列の非零成分がある。$v=0$ なら上の三行が独立である。
平面 $z=v_1+2v_2+3v_3$ 上、十五個の四次小行列式は三変数四次式の全空間を張る。
実際 $(v_1,v_2,v_3)=(i,j,1)$, $i,j\ge0$, $i+j\le4$ の十五点での値の行列式は
$101$ を法として順に $53,38,82$ である。従って標数零で独立であり、
この平面上に退化点はない。射影的な退化軌跡は有限、各ファイバーは一点以下である。
混合点の逆辺も対称行列の転置と交代行列の符号反転だけが違うので、
星座標の符号変更および二位置の交換で同じ有限性を得る。
\path{computations/verify_prescribed_discriminant_seeds.py} が小行列、微分と次の安定化群の計算を再現する。
\end{proof}

\begin{proof}[同型類の次元]
定理\ref{thm:marked-point-commutators}の残余群では、$Q=P$ の三独立な基底を
保つ条件が $BB^t=sI$ であり、自由な変換群は
$\mathrm{GL}(U)\times\mathrm{GO}(V)$ で次元13である。
$g=4$ なら法線と共通二平面の交換子方向はテンソルと印から固有なので、
両者の射影値を固定した場合の群の次元は $9+1+d$ である。
基点の $G'$ を保つ射影直交群は有限である。半単純な場合は対角平面の安定化群、
残る三例では $J$ 交代な $T$ に対する $[T,G_i]\bmod G'$ の係数行列が階数3である。
従って連結安定化群では $B$ がスカラーで、$G$ の三独立な順序付き基底を保つため $A=I$、
$F$ を保つ条件から $s=1$ となる。近傍でも安定化群は有限である。
連結な変換群は既約成分を保ち、開部分での軌道の次元は上の群次元に等しい。
印を忘れる選択肢も有限であるので $19-(10+d)=9-d$ を得る。
一般に動かす場合は $[N]\in\mathbf P^4$ と $[\ell]\in\mathbf P^2$ の六変数を基底に加えた
同じ束とエタールチャートを用いる。次元 $25-13=12$ が従う。
\end{proof}

% !TeX root = ../papers/nc-p3.tex
\section{半単純な判別法線を指定した族の最小周期}
\label{sec:marked-semisimple-period}
\begin{theorem}\label{thm:marked-semisimple-period}
定理\ref{thm:marked-prescribed-discriminant-as}の半単純な法線 $N$ と任意の直線 $\ell$ に対し、
そこで構成した既約族の幾何学的一般点の最小周期は4である。
印を忘れた場合も各位置への射影作用を取った添字保持自己同型群は自明である。
両因子を動かす十二次元の同型類の族の一般点でも成立する。
\end{theorem}
\begin{proof}
$N$ を対角化し、前節の $G_i=D_i$, $F_i=-D_i$, $Q_i=-[e_i]_\times$ を基点にする。
$B$ を位置1の配置、$A,C$ を位置0,2の配置とする。$b=(r,v)\in B$ の両隣
$a=(s,u)\in A$, $c=(h,w)\in C$ はそれぞれ
\[
 (sI+G_u)v=(rI+Q_v^t)u=0,\qquad
 (hI+F_w)v=(rI+Q_v)w=0
\]
で求まる。基点の $B$ は星点、$r=1$ で三成分のうち一つが零で他の二つが
$\pm\sqrt{-1/2}$ の十二点、$r=0$ で三座標点と $(1,\pm1,\pm1)$ の四点である。
四辺は二十個の被約点のグラフなので、基点の厳密局所化の近くでこれらを別々に標識できる。

一般点の射影同値は、標識のある近傍で、ある一つの置換に対して恒等的に成り立つ。
これを基点で調べる際の行列の退化を除く。非零行列の射影空間内で同値のグラフを閉包し、
各標識点について像と相手の点の二次小行列式を零とする。
基点へ特殊化した行列が特異なら、核に入らない点は像内に写る。
基点のどの平面も高々六点しか含まないので、核と像を合わせても二十点を対応させられない。
よって基点でも可逆である。基点の四つの座標点と全成分非零の一点は射影枠であり、
標識置換を固定すれば近傍で行列はこの五点から正則に復元できる。従って同じ置換が
\emph{全ての}接方向に沿って一次まで持ち上がらなければならない。

以下の小行列計算の入力と全検査は
\path{computations/verify_semisimple_period_obstruction.py}に置く。
分裂素数 $1009$ で還元し、二十点ずつの $A,C$ を上の二組の六式の核で構成する。
両者とも三点が共線になることはなく、六点を超える平面もない。四つの座標点は
四点平面十二枚・六点平面六枚に属し、残りはそれぞれ三枚・九枚に属する。
従って体を拡大しても還元配置間の射影同値は単項行列である。
座標置換と全成分非零の一点の像で列挙すると $A\to A$ と $A\to C$ は各96個である。
この列挙は標数零の候補もすべて検査する。実際、上の核・像の個数による議論は
この素数での特殊化にも適用できる。二十点は還元で分離したままであり、
置換を固定した行列は両配置の枠から単元を分母に持つ式で復元できる。

まず $T$ の行を $(3,2,-1),(1,-3,4),(-1,2,4)$ とし、
$\dot G_i=\sum_jT_{ij}D_j$, $\dot F=-\dot G$ とする。
$G'$ を固定するこの方向は、全ての指定 $N,\ell$ の族で取れる。
$\dot Q$ は前節の可逆微分で復元し、各点の微分は上の六式を二つのアフィンチャートで
微分した六次連立一次方程式で求める。左右から得た $\dot b$ は一致する。
候補 $M$ と対応点 $Mx=\lambda y$ について必要な式は
\[
              (\dot Mx+M\dot x-\lambda\dot y)\bmod ky=0.
\]
行列 $\dot M$ の十六成分への六十式の係数階数は15である。右辺を加えても
階数15となるのは、自己同型では恒等、$A\to C$ では $J=\operatorname{diag}(1,-1,-1,-1)$ だけである。

次に $G$ を固定する。$\dot F$ の対角跡零部分は三列各二成分である。任意の六成分に対して
$\sum[D_i,\dot F_i]=0$ となるように $D_j$ の係数を補える。
$[D_2,D_3],[D_3,D_1],[D_1,D_2]$ が基底であり、補い方の差は
$\dot F_i=\sum K_{ij}D_j$ の対称行列 $K$ である。この差は $J$ の持上げを妨げない。
実際 $F=(-I+tK)G$ の係数行列は非退化対称であり、定理\ref{thm:marked-point-commutators}
の $U$ 側の合同変換により基点で恒等な基底変更で $F=-G$ に戻せる。
その基底で $Q$ は交代解の一意なエタール持上げであり、二組の六式は
$(r,v)\mapsto(r,-v)$ と $(s,u)\mapsto(s,-u)$ で交換される。
従って対称な補正の六列は標数零でも $\dot M$ の像内にある。
これに対し六個の対角成分の単位方向を上の交換子基底で補い、
$J$ の一次方程式の右辺に並べると、還元で拡大階数は21である。
係数側には常にスカラー核があるため、標数零でも余核内の六列は独立である。

指定した $N,\ell$ の族では $G'$ を固定して $H\subset F'\subset N^\perp$ の
射影直線上で $F'=G'$ から離れる接方向がある。前節の階数六の基底束によりこれは
実際に持ち上がる。$N^\perp/G'$ は非零の対角跡零方向なので、この接方向の六成分は非零である。
よって $J$ は持ち上がらず、一般の二つの偶配置は射影同値でない。
二位置移動にはこの同値が必要であり、標準四周期性から最小周期が従う。
位置0で恒等な射影自己同型はグラフに沿って各位置の全点を固定する。
基点の偶配置は射影枠を含む。奇配置の七点を固定すれば無限遠平面上の作用はスカラーであり、
星点と有限の非零点も固定するので全体がスカラーである。この条件は近傍で開である。
反対の奇配置も混合点で符号変更により同じ配置なので、全位置の射影作用が恒等となる。
\end{proof}

% !TeX root = ../papers/nc-p3.tex
\section{混合窓を通る判別層と十五次元成分への特殊化}
\label{sec:marked-mixed-component}
\begin{theorem}\label{thm:marked-mixed-component}
$G,F$ が対称側、$Q,P$ が交代側である有限の混合窓は、印を忘れても、
定理\ref{thm:full-fifteen-component}の十五次元成分に全て属する。
この向きの混合点を通る片側対称な印付き表示では、$g=5,4,3$ の同型類の像を
含む入れ子の既約閉包があり、その次元はそれぞれ $14,12,8$ である。
これらの一般点は全次数で正則で、最小周期は順に $4,4,2$ である。
定理\ref{thm:marked-prescribed-discriminant-as}の指定法線・指定直線の全構成は第二の閉包に属する。
この記述における局所特異点は印と正規化座標を取った表示のものである。
\end{theorem}
\begin{proof}[三層の局所環と閉包]
$S=\mathrm{Sym}_0(V)$ と置き、$G$ は $\bigwedge^2 G'\to\mathfrak{so}(V)$ が可逆な
$S$ の三基底列を動く。ここで $G'=\langle G_i\rangle$ である。
三列 $F$ から $Z\in\mathrm{Hom}(k^3,S/G')$ を取る。交換子の跡条件
$\sum[G_i,F_i]=0$ は、$F_i$ に $G_j$ の倍を加えることで解ける。
補正係数の交代部分に関する微分は三交換子を列に持つ同型である。
よって局所的に十五個の $G$ 座標と任意の $2\times3$ 行列 $Z$、六個の対称な補正係数が自由であり、
$g=3+\operatorname{rank}Z$ となる。$Q=P$ の復元は交代な基点でエタールなので、
完備局所環でもこの積表示が成り立つ。$g\le4$ は滑らかな二十一次元因子と
階数1以下の $2\times3$ 行列錐の積で、特異集合は $Z=0$ である。
実際、階数1なら一つの単元成分で残りを解ける。原点では全小行列式が二次である。
錐は $Z=u v^t$ の座標環で、二次小行列式は同じ行・列多重次数の単項式を
添字交換で全て同一視するため、その関係イデアルを生成する。
この半群は非負で行と列の総次数が等しい全ての次数からなり、飽和しているので錐は正規である。
像直線と核二平面をともに記憶する表示は $\mathbf P^1\times\mathbf P^2$ 上の
$\mathcal O(-1,-1)$ の全空間であり、原点上に零切断を持つ。
像直線からは法線 $N$ が決まり、核二平面の $F$ による像からは共通二平面と
その交換子の直線 $\ell$ が決まる（ランク1で $F$ が独立な開集合）。
従って固定した二因子を動かした前節までの構成もこの同じ局所的な層に入る。

全ての有限混合点を同じ閉包で扱えることを確認する。両基底を合同変換で規格化すれば
$F=-G$ で $Q_i=[Le_i]_\times$ とできる。可逆な $L$ の余因子行列が三交換子から
指定されるため、解は既約な $G$ の開集合の上の次数2のエタール被覆をなす。
この被覆が分離しても二成分を $L\mapsto-L$ が交換する。これは残余群の
$B=-I,A=I,s=1$ で実現され、同型類の像は既約である。
$g\le4$ の局所環は上記の正規な錐により各混合点で幾何学的単枝であり、
制限のない片側対称表示もそこでは滑らかである。
従ってこの混合軌跡の各既約成分に沿って、通過する各層の成分は一つで同じである。
実際、有限個の成分との閉交叉が軌跡を覆い、少なくとも一つが軌跡全体を含む。
有限点条件は開で、有限混合候補は定理\ref{thm:mixed-point-frobenius-algebras}により全て正則である。
このうち分裂二十点例で安定化群は有限であり、印を忘れるファイバーも有限である。
正規化表示の三層の次元 $27,25,21$ から群の次元13を除けば主張の像次元になる。
一般点での正則性は相対bar複体の内部次数ごとの特殊化から得る。
$g=4$ の一般最小周期と射影自己同型の自明性は定理\ref{thm:marked-semisimple-period}より、
配置の番号付けと同じ閉条件を使えば $g=5$ の一般点にも従う。
\end{proof}

\begin{proposition}\label{prop:split-mixed-transverse-curve}
分裂混合点 $G=D$, $F=-D$, $Q=P=K$, $K_i=-[e_i]_\times$ を通り、
\[
             G=D+tK,\quad F=-D+tK,\quad Q=K+tD,\quad P=K-tD
\]
と置く。このアフィン直線上の窓が標準Koszul左右AS正則である必要十分条件は
$(t^2+3)(3t^2+1)\ne0$ である。反対階数は $t=0$ で10、$t^2=-1$ で14、他で16となる。
正則点の周期は2を割る。
\end{proposition}
\begin{proof}
標準印付きの交換子式は全 $t$ で成り立つ。式が高々二次なので $t=0,1,-1$ の
三つで係数ごとに確かめられる。星印の六次単位小行列により二切断は階数6である。
$F_i=-G_i^t$, $P_i=-Q_i^t$ なので二位置移動は各補空間の符号を反転すれば四テンソルを保つ。
反対切断を四元群ラベルの総和ごとに四ブロックにするとその行列式は順に
$-12(1+t^2)^2,16t^2,16t^2,16t^2$ である。零色ブロックの星行と星列には
非対角に1が並び、残る三次正方部の非対角成分は $-2(1+t^2)$、対角は零なので $t^2=-1$ で階数2となる。

$c^2=-2(1+t^2)\ne0$ とする。四位置の星座標を全て $c$ 倍し、さらに四位置で順に
\[
 (i,1,i,1),\quad (i,i,1,1),\quad (-i,1,-i,1),\quad (-i,-i,1,1)
 \qquad (i^2=-1)
\]
の対角倍率を掛けると Walsh 表示になる。すなわち $d=b+e$ とすると
成分 $w_{j,j+b,j+d,j+e}$ は $(-1)^{j_1d_0} m_{be}$ である。
欠色は両方とも0であり、残り三色での符号付き平方行列を $c^2$ で割ると
第一行が $(-2(1+t^2),(1-t)^2,(1+t)^2)$ である巡回行列となる。
行和・列和は零で、全二次余因子は $(t^2+3)(3t^2+1)$ に等しい。
欠色のラベルを移す操作はこの階数を保つ。定理\ref{thm:walsh-normal-quadrics}で判定できる。

$t^2=-1$ では、補文字だけからなる全成分が消え、非零成分は星を一つか二つ持つ。
星を $z$ 倍してテンソル全体を $z$ で割った族において $z=0$ には一つ星のみが残る。
上記の位相倍率を掛けると、これは Walsh 表示で $m_{bb}=0$、
非対角係数は上の $m_{be}/c$ の特殊値である。実際 $b\ne e$ の項は
四ラベルを一度ずつ含み、星倍率が正確に一つである。符号付き平方行列の余因子は
$t^2=-1$ で $-4$ となるのでこの特殊窓は全次数正則である。$z\ne0$ の族は
元の窓と同型であり、標準Frobenius条件および内部次数別の完全性を特殊点から
幾何学的一般点へ移して、元の窓での全次数正則性を得る。
\end{proof}

\begin{proof}[十五次元成分の同定]
上記曲線の $c\ne0$ の全点は定理\ref{thm:walsh-full-component}により既知の
十五次元成分の基底変換による像に属する。分裂混合点で他の成分との交叉の可能性を除くため、
命題\ref{prop:twenty-conormal}を適用する。第一辺の二十点を前節の $(a,b)$ とすれば、
第二切断 $N$ の左右の部分評価は階数3で、核点はそれぞれ $(c,b)$ である
（ここで $c$ は位置2の点）。従って二十個の余法線は
$(c\otimes b)\otimes(a\otimes b)$ である。分裂素数1009でその二十行の階数は20、
二切断の各十次元の左右核の積から得る合計二百行のSchur微分の階数は180となる。
前者による二十個の恒等式があるため後者は標数零で180を超えない。
よって分裂混合点の全テンソル軌跡は滑らかな次元76で、通過成分は既知のもの一つである。
混合点を通る既約な $g=5$ の片側対称族の閉包もこれに入り、三層に共通する成分の同定が従う。
さらに反対行列式は曲線上 $t=0$ に正確に六位の零を持つ。核と余核の六次元間に誘導される
一次写像が可逆であり、この曲線は片側対称な階数十の層から横断方向へ出る。
\path{computations/verify_mixed_component_smoothing.py}が二つの微分、余法線、三値の
交換子、四ブロックとWalshの全係数恒等式を再現する。
\end{proof}

% !TeX root = ../papers/nc-p3.tex
\section{混合点を横断する曲線の全同型と周期}
\label{sec:marked-mixed-curve-isomorphisms}
\begin{theorem}\label{thm:marked-mixed-curve-isomorphisms}
命題\ref{prop:split-mixed-transverse-curve}の標準正則な二窓 $W(t),W(t')$ が
印を忘れて添字保持同型である必要十分条件は $t'^2=t^2$ である。
最小周期は $t^2=1,-1$ で1、他で2である。特に $x=t^2$ を同型類のパラメータとすれば、
除かれる値は $x=-3,-1/3$ の二つである。階数14の $x=-1$ は正則なWalsh窓とも同型にならない。
$x\notin\{0,1,-1\}$ の正則点で四位置に作用する射影自己同型群の位数は48である。
\end{theorem}

\begin{proof}[反対切断の逆から軸を復元する]
反対階数16の四テンソル $w$ に対して、行 $ac$・列 $bd$ の反対切断の逆転置を
同じ添字順に戻した共変テンソルを $w^\vee$ とする。異なる二位置 $i,j$ を残し、
$w\otimes w^\vee$ の他二位置を縮約して
$C_{ij}\in\mathrm{End}(V_i)\otimes\mathrm{End}(V_j)$ を取る。
例えば $C_{01;aa',bb'}=\sum_{c,d}w_{abcd}w^\vee_{a'b'cd}$ である。
両端の跡対合によりこれを写像 $L_{ij}:\mathrm{End}(V_j)\to\mathrm{End}(V_i)$ とみて
$K_{ij}=L_{ij}L_{ij}^*$ を作る。全体倍率は消え、基底変換は各位置で共役になる。
反対切断を転置しても逆テンソルは添字交換で対応するので、この構成は巡回移動とも可換である。

位置 $i$ の二隣 $j=i\pm1$ について $K_{ij}$ の固有値 $16/9$ の同時固有空間を取る。
$t(t^2+1)(t^2-1)\ne0$ なら、$i=0,1$ の少なくとも一方でその空間は
跡零対角空間に等しく、したがってその四共通固有直線が四座標軸を復元する。
これを全ての例外根も含めて示す。各行列単位 $E_{a,a+h}$ を同色 $h$ でまとめると
$C_{ij}$ は四ブロックで、$R_h$ をその $h$ ブロック、$S_h$ を $a\mapsto a+h$ の置換とすると
$K_{ij}$ のブロックは $R_h S_h R_h^t S_h$ である。
$K_{ij}-16I/9$ に $144t^2(1+t^2)^2$ を掛けると、対角ブロックは
全成分が $512t^2(1+t^2)^2$ の行列である。色 $h=1$ では二隣を $+,-$ の順に縦に重ね、
零始まりの行番号 $(0,1,2,3)$ と $(0,1,2,4)$ を取る。
$i=0,1$ の計四小行列式のモニック最大公約数は $(t^2-1)^4(t^2+1)^4$ である。
色 $1,2,3$ の巡回は元の三列を保つため、単元になる小行列のある位置では
三つの非零色の全てが核に寄与しない。この検算は
\path{computations/verify_mixed_curve_axes.py}の有理関数体上の演算である。
$t^2=1$ でも跡零対角空間の包含は保たれる。
従ってこの方法で四軸を復元できる窓との同型があれば、相手の対応位置でも
同じ次元3の空間を保ち、やはり四軸を移す。

$u=1+t$, $v=t-1$ とする。各巡回三組 $(i,j,k)$ で、辺の色 $i$ の列を
$(0i,i0,jk,kj)$ とした二行の関係は
\[
 E(u,v)=\begin{pmatrix}1&0&u&-v\\0&1&v&u\end{pmatrix}
\]
である。位置 $0,1,2,3$ から順にその引数は $(u,v),(v,-u),(-u,-v),(-v,u)$ となる。
どの軸も点グラフで両隣の同名の軸だけに対応する。実際 $0$ を評価すると補三文字が消え、
$i\ne0$ を評価すると $0$ と他の二文字が消える。これは $u,v$ が同時に零でないことによる。
従って一位置から四位置全てに軸を運べて、同型の置換部分は全位置で同じである。
\end{proof}

\begin{proof}[対角同値と特殊値]
Walsh表示の同時平行移動は四軸上で推移的であり、前節の規格化は対角である。
$t^2\ne-1$ の階数16なら、相手側でこの自己同型を合成して $0$ 軸を固定できる。
残る三軸の巡回置換は窓を保つ。奇置換を全位置で施し、
奇位置の補文字をさらに $-1$ 倍すると $W(t)$ は $W(-t)$ になる。
これは $D$ が置換で符号を変えず、外積 $K$ は行列式の符号を取るためで、
階数の制限なしに成り立つ。

よって $s=t$ または $-t$ とし、$E(u_s,v_s)$ を対角倍率だけで移せる場合を調べればよい。
星倍率を各位置で1に戻し、二位置の補文字の倍率を $a_i,b_i$ とする。
行の先頭二列が単位行列なので、$E(u',v')$ に移すには、全ての巡回三組で
\[
 a_j b_k=(u'/u_s)b_i=(v'/v_s)a_i,\qquad
 a_k b_j=(v'/v_s)b_i=(u'/u_s)a_i
\]
を要する。ここでは四係数が非零の場合を扱う。
$a_i/b_i$ が全 $i$ で一定なので左の二つの積は等しく、
$u'/u_s=v'/v_s$ となる。相手が $W(t')$ の位置0ならこれで $t'=s$ である。
相手が $W(t)$ の位置1なら $(u',v')=(v_t,-u_t)$ であり、$s=t$ では
$v_t^2=-u_t^2$ が必要、$s=-t$ では $-1=1$ を要する。
従って対象の階数16で両係数が非零なら最小周期は2である。
いずれかの係数だけが零になる値と両方非零の値の同一視は、一位置の軸の復元後、
上の関係行列の零成分の不一致で排除される。

自己同型ではこの範囲で $s=t$ なので、星を固定する軸置換は巡回三通りである。
置換が恒等なら上式で $a_i=b_i$、$a_ja_k=a_i$ であり、解は積1の三つの符号の四通りである。
隣接位置にも同じ式を適用して全位置の倍率が揃う。これら四変換は実際に全関係を保ち、
星の像の四通りも前記の平行移動で実現する。従って群位数は $4\cdot3\cdot4=48$ である。

$t=1$ では、三補文字の一つの互換を全位置に施し、奇位置の星文字のみ符号を変えると、
上の四つの関係空間は各々一つ先に移る。$t=-1$ はこれと同型である。
$t=0$ は固有の階数10で最小周期2、$t^2=-1$ は固有の階数14で二値は互いに同型である。
後者は一周期の表示も明示できる。位置 $n$ の星文字を $(1+t)t^n z$ に置き、
三補文字を $x_1,x_2,x_3$ と共通に読むと、各位置で
\[
          x_i z=z x_i,\qquad x_jx_k-tx_kx_j=-z x_i
              \quad ((i,j,k)\ \text{巡回})
\]
になる。実際 $v=tu$ で、各辺の第一行の二つの末尾係数は $(ut^n,-ut^{n+1})$ であり、
第二行と第一行の $t$ 倍の差が可換関係を与える。

最後に $t^2=-1$ の四テンソルに、全位置で星を $z^3$ 倍、三補文字を $z^{-1}$ 倍する
特殊線型群の一径数部分群を作用させる。前節の一つ星の極限を得る。その反対切断では
零色のブロックは零、他は四次可逆なので階数12である。従って元のアフィン軌道は閉でない。
Walsh表示の非零テンソルの軌道はノルム最小性により閉である
（定理\ref{thm:walsh-spectral-fibers}の証明）。倍率と任意の基底変更でこの相違は保たれる。
\end{proof}

% !TeX root = ../papers/nc-p3.tex
\section{四つの基点を通る判別三次式と三平面の楕円曲線}
\label{sec:marked-cubic-linear-system}
\begin{theorem}\label{thm:marked-cubic-linear-system}
$G'\subset S=\mathrm{Sym}_0(V)$ を固定した有限混合印の対称三平面とし、
$\bigwedge^2G'\to\mathfrak{so}(V)$ は可逆とする。$G'^\perp$ の基底 $N_1,N_2$ が
定める二次式 $q_i(v)=v^tN_i v$ の完全交叉を $B$ とする。
$G'$ を固定した全ての相手三平面の混合行列式は $H^0(\mathcal I_B(3))$ に入る。
逆にこの六空間の\emph{任意の}非零三次式 $f$ を判別式に指定して、元の混合点を通る
幾何学的一般点で全次数正則・四辺有限・反対階数10の族を取れる。
$f=l_1q_1+l_2q_2$ の一意な表示で $l_1,l_2$ が独立なら、その一般点で $g=5$ である。

特に分裂点 $G_i=D_i$ では $B=\{(1,\pm1,\pm1)\}$ であり、
この四点を通る各 $f$ に対し、同型類の像が六次元、一般最小周期4、一般射影自己同型群が
自明な族となる。従って滑らかな三次曲線の全ての $j$ 不変量を含み、
台が一本の直線より大きいどの平面三次曲線も射影同値までこの構成で実現する。
\end{theorem}
\begin{proof}
二次式に共通因子があれば両方が同じ一次式 $v\cdot n$ を持ち、
跡零性から他の一次因子二つは $n^\perp$ を張る。従って $G'$ は全元が
$kn$ を保つ三平面となり、有限印の既約性に反する。二式は正則列で、飽和イデアルの
三次部分は $q_1,q_2$ の一次式倍の直和である。
$\omega_v(X,Y)=\operatorname{tr}([v]_\times[X,Y])$ の根基を生成する
$T_v=vv^t-(v^tv)I/3$ は $B$ 上スキームとして $G'$ に入る。各点でその一成分は単元なので、
任意の $F$ に対し $\det(\omega_v(G_a,F_c))$ は $B$ 上で零である。

この行列式写像の $F=-G$ での微分を計算する。$F$ の $G'$ 内での変動は微分で零になる。
$Z\in\mathrm{Hom}(k^3,S/G')$ が階数1で $Z_a=\eta_a Y$ なら、
$\ker\eta$ の二基底の $G$ による像を $H_1,H_2$ とする。
命題\ref{prop:marked-mixed-discriminant}のPfaffian展開により、一次係数は
$\omega_v(H_1,H_2)\, q_N(v)$ に比例し非零である。ここで $N$ は $G'$ と $Y$ を消す法線である。
実際行列式において一列だけを $Y$ に替える基底を使えるので係数は一次が全てである。
三交換子の独立性から $[\eta]\mapsto[\omega_v(H_1,H_2)]$ は射影線型同型であり、
$[Y]\mapsto[N]$ は二次元商での消去によりやはり射影線型同型である。
従って微分の像は全ての一次式と任意の鉛筆二次式の積を含む。その線型包は六次元で、
微分は六空間の同型である。特にこの同型は階数1の射影Segreを、
$l_1,l_2$ の従属するSegreに移す。

跡条件 $\sum[G_i,F_i]=0$ は $G'$ 内の係数で任意に補える。
よって固定 $G$ の十二次元の $F$ 座標からこの六空間への写像は混合点で滑らかであり、
指定した直線 $kf$ の逆像は同点で滑らかな七次元で、実際に直線に支配的である。
交代解のエタール復元と内部次数別の相対bar完全性を前節までと同様に適用する。
$l_1,l_2$ が独立なら階数1の変動ではこの直線方向に出られず、階数2の方向があるから
一般点で $g=5$ となる。

分裂点では対角跡零行列が法線鉛筆で、表示した四点を得る。$G'$ を固定して
順序付き基底も動かすと次元は $7+9=16$、$G'$ を保つ射影直交群は有限であり、
残余群の連結成分は次元10である。この群は判別式の射影値を保ち、印は有限なので
像次元は6となる。定理\ref{thm:marked-semisimple-period}の混合点内の基底変動はこの族でも取れ、
$kf$ の非零方向は必ず対角商で非零の変動を持つ。同定理の六列が独立な持上げ障害は
そのまま適用できる。最後に三次曲線の台から射影枠四点を選んで上の四点に移せばよい。
二直線に支持される場合にも各直線から二点を選べる。三重直線だけはこの枠を持たないが、
その射影型は等方一次式とその平方を定理\ref{thm:marked-prescribed-discriminant-as}で指定して実現できる。
この射影型の議論では曲線の座標を任意に取り替えており、指定計量との相対位置は固定していない。
\end{proof}

\begin{theorem}\label{thm:marked-elliptic-complement-planes}
分裂三平面 $G'=\langle D_1,D_2,D_3\rangle$ を固定し、滑らかな $E=V(f)$ が
上の四点を通るとする。$F'\in\mathrm{Gr}(3,S)$ について、混合行列式が $kf$ に入る
全三平面のスキームは、次数5の楕円正規曲線で $\mathrm{Pic}^2(E)$ と同型である。
そのうち行列式が恒等零になるのはちょうど五つの被約点で、
$G'$ と、四点 $v\in B$ について $kv$ を保つ跡零対称三平面である。
他は全て $g=5$ である。これは対称側の相手の\emph{部分空間}を尽くす。
順序付き基底の跡条件、三交換子の持上げとその特殊点のAS条件は別に課す。
\end{theorem}
\begin{proof}
$G'$ と $S$ の基底で $A(v)=(\omega_v(G_a,S_j))$ を三行五列の一次行列とする。
その三次小行列式は前定理により $H^0(\mathcal I_B(3))$ を張る。相手の三平面の
Pl\"ucker埋込み $\mathrm{Gr}(3,5)\subset\mathbf P^9$ において混合行列式は六成分の線型写像である。
中心の $\mathbf P^3$ とGrassmann多様体は表示した五点で横断する。
\path{computations/verify_mixed_cubic_pencil.py}が三つの非対角対称行列と二つの対角跡零行列を
基底に用い、各三平面のアフィンチャートの六座標に関する小行列式微分が
有理数上で階数6であることを検算する。Grassmann多様体の次数は5
（$2\times3$ 長方形の標準盤の数）である。線型切断を順に取ると成分の次数の和は
増えないため、横断的な孤立五点で全切断を尽くす。

$E$ 上の核束 $\mathcal U=\ker(S\otimes\mathcal O_E\xrightarrow A (G')^*\otimes\mathcal O_E(1))$
を取る。$B$ でも像は曲線上で局所自由なので核は部分束である。
$A$ の階数は $T_v\notin G'$ で3、$T_v\in G'$ で2である。
各基点で小行列式は $E$ 上の極大イデアルを生成する。実際六次元の三次式系が
二つの二次生成元の全一次倍を含み、その一方以上は滑らかな $E$ 上で一次に消える。
従ってこの写像の余核の長さは4で、$\deg\mathcal U=-5$、
$\det\mathcal U^*=\mathcal O_E(3)(-B)$ である。
$H^0(\mathcal U)$ の元は線型行列 $A(v)$ により曲線上恒等的に消える定数ベクトルで、
曲線は平面を張るから $G'$ の全元と交換しなければならない。既約性と跡零性からこれは零で、
種数1のRiemann--Rochにより $h^0(\mathcal U^*)=5$ である。

自然写像 $S^*\to H^0(\mathcal U^*)$ は同型である。実際その核に $a\ne0$ があれば、
$a\wedge S^*$ の射影三空間は可分解なPl\"ucker点のみからなり、全て $E$ 上の
$\bigwedge^2\mathcal U^*$ への制限が零である。この制限が零という条件は、
三平面 $F'$ の消去二平面を取り、上の線型写像で行列式が $kf$ に入る条件と一致する。
一般の $v\in E$ で二平面 $\mathcal U_v$ が $F'$ と交わる条件そのものであり、
小行列式は $\det\mathcal U^*$ の切断として基点因子を取り除いたものである。
従って先の射影三空間が $kf$ の逆像の $\mathbf P^4$ に入り、その超平面である中心とは
少なくとも平面で交わり、中心の有限性に反する。

$\mathcal U^*$ は安定である。次数3以上の線部分束があればその切断三つ以上の
二つずつの外積は全て $E$ 上で零となる。三次元切断部分の外積の射影平面は可分解点だけからなり、
上の $\mathbf P^4$ の中で中心と少なくとも直線で交わってしまう。
そこで任意の $M\in\mathrm{Pic}^2(E)$ に対し $\mathcal U^*\otimes M^{-1}$ は次数1の安定階数二束で、
双対に切断がなくRiemann--Rochから切断は一次元である。
得られる $M\hookrightarrow\mathcal U^*$ に零点があれば飽和の次数は3以上になるので部分束である。
$H^0(M)$ は基点自由な二平面で、これを $S^*$ 内に取れば対応する $F'$ が得られる。
逆に $kf$ の逆像のGrassmann点は二つの独立切断を与え、その評価は $E$ 上常に従属する。
像の飽和は独立切断を二つ持つ線束なので次数2以上で、安定性からちょうど2である。
従って元の $M$ が一意に戻る。

この対応はPicard曲線からの射である。普遍線束に対する上の切断空間は次元一定なので
局所的に束となり、二平面はGrassmann多様体へ正則に写る。逆像の台は以上より既約な曲線であり、
その正規化はこの全単射と標数零から $\mathrm{Pic}^2(E)$ である。
五超平面による所期次元の完全交叉はCohen--Macaulayで、$F'=G'$ で滑らかなので被約である。
$K_{\mathrm{Gr}(3,5)}=\mathcal O(-5)$ より随伴公式で算術種数1となり、
正規化の種数と一致するため滑らかである。次数は5であり、五基点のPl\"uckerベクトルは前掲の検算で階数4なので中心を張る。
曲線は中心外にも点を持つので全体の線型包は四空間である。
従って楕円正規五次曲線である。$g=4$ なら判別式は可約なので、
五基点を除いて $g=5$ に限る。
\end{proof}

% !TeX root = ../papers/nc-p3.tex
\section{指定三次式からの相手三平面の全復元}
\label{sec:marked-companion-plane-reconstruction}
\begin{theorem}\label{thm:marked-companion-plane-reconstruction}
定理\ref{thm:marked-cubic-linear-system}の固定三平面 $G'$ を取る。
$B=V(q_1,q_2)$ の非被約性も許す。$0\ne f=a_1q_1+a_2q_2$ に対し、
混合行列式が $f$ の\emph{非零倍}となる相手三平面 $F'$ は以下で全て得られる。
\begin{enumerate}
\item $g=5$ の解には平面の点 $p$ が一意に付く。$p\notin B$ では
$f(p)=0$ かつ $(a_1(p),a_2(p))\ne(0,0)$ の各点にちょうど一解があり、線型方程式で復元される。
$p\in B$ では $a_1(p)=a_2(p)=0$ の場合だけ解があり、各々非零スカラー一つで動く。
\item $g=4$ の解があるのは $f=lq$ で $q\in\langle q_1,q_2\rangle$ のときだけで、
相手はこの場合一本のアフィン直線をなす。
\item 全三次式が零となる中心はスキームとして $\{G'\}\sqcup B$ である。
特に $E=V(f)$ が滑らかなら、中心も含めた $kf$ の逆像は任意のこの $G'$ に対して
$E$ と同型な楕円正規五次曲線となる。
\end{enumerate}
\end{theorem}
\begin{proof}[核の二列と消去]
$A(v)$ を前節の三行五列の行列とし、$X(v)\in G'$ を、$G'$ での三次交代部分の
$23,31,12$ 成分を係数とするベクトルとする。三交換子が独立なので $X:V\to G'$ は同型である。
$A(v)$ は $X(v)$ と $T_v=vv^t-(v^tv)I/3$ を共に消し、$B$ の外でその二列は独立である。
従って $A$ の三次小行列式全体の双対外積は定数倍を除き $X(v)\wedge T_v$ に等しい。
実際両辺は同次数で、いずれも成分の共通零点は有限なので、有理関数の倍率は定数である。
消去二平面を $F'^\perp=\langle\alpha,\beta\rangle\subset S^*$ とすると、混合行列式は
定数倍を除いて
\[
           R(\alpha)H(\beta)-R(\beta)H(\alpha),\qquad
      R(\alpha)(v)=\alpha(X(v)),\quad H(\alpha)(v)=\alpha(T_v)
\]
となる。$R$ は $V^*$ への全射で、$H$ は $S^*$ を跡零二次式の全空間に同定する。
$\ker R$ はこの同定で $\langle q_1,q_2\rangle$ である。

$R$ が消去二平面上で階数2の場合、像を $p^\perp\subset V^*$ と書く
（ここでは点 $p$ を消す線型形式の空間）。$R$ の線型切断 $s$ と体積形式を固定する。
$p$ の局所代表から決まる二ベクトルを $l\wedge m=\iota_p\mathrm{vol}$ と分解する。
基準解 $s(l),s(m)$ の三次式を $b_1(p,v)q_1(v)+b_2(p,v)q_2(v)$ と書くと
各 $b_i$ は $p,v$ に双線型である。$v=p$ では $l,m$ が共に消えるため、
互いに素な二次式の関係から
\[
       (b_1(p,p),b_2(p,p))=\lambda(-q_2(p),q_1(p))
\]
となる。ここで $\lambda$ は定数で、零ではない。もし零なら $b_1,b_2$ は交代形式である。
切断に $q_i$ 倍の線型汎関数を加える操作は $b_i$ に $\mathrm{vol}(p,h_i,v)$ を
加えるので、両交代形式を同時に消せる。そのとき全ての $l,m$ で
$l H(s(m))=m H(s(l))$ である。三つの独立座標で割り切りを用いれば、
ある共通一次式 $L$ により $H(s(l))=lL$ となる。
全ての $l$ で $lL$ が跡零なら計量の非退化性から $L=0$ となり、切断が存在することに反する。

消去二平面を $s(l),s(m)$ のそれぞれへの鉛筆二次式の加算で全て書ける。
四つの自由係数は、各 $b_i(p,-)$ に $p$ を消す任意の一次式を加える四座標に同型に写る。
従って二つの値を $v=p$ に代入した式
\[
       z a_1(p)=-\lambda q_2(p),\qquad z a_2(p)=\lambda q_1(p)
\]
が全ての整合条件である。解 $z$ ごとに残りは一意で、$z\ne0$ の場合が第一項そのものである。
局所代表・基底を替えると $z$ は単元倍される。これは点ごとの解法であると同時に、
四自由係数の消去は各局所チャートで可逆なのでスキームの方程式でもある。
$z=0$ の全解は $B$ 自身に同型である。そこでの三平面を行列で書けば
$kv$ を保つ跡零対称行列の空間となる。$T_v\in G'$ と交換するので行列式は零であり、
$X(v)$ は $G'$ 内の交代行列の根基 $kT_v$ にあるので同じ点 $p=v$ を与える。

階数1では消去平面の基底を $q, s(l)+q'$ と取れる。三次式は $lq$ の定数倍、
$q'\in\langle q_1,q_2\rangle/kq$ が一意な自由座標である。
分解の一意性は $f=a_1q_1+a_2q_2$ の係数の一意性から従う。階数0なら $F'=G'$ だけで、
前節の同型微分により中心の単純な孤立点である。これで中心のスキーム記述も従う。

滑らかな $E$ なら $a_1,a_2$ は独立で、その共通零点 $p_*$ は $B$ の外にある。
もし $B$ 上なら $df$ もそこで消えてしまう。$p_*$ の外では二式から $z$ を正則に
消せて、ちょうど $f=0$ が残る。階数1の解はなく、$p_*$ に対応する補完点は $G'$ だけである。
$G'$ では曲線が滑らかであるから、全逆像は $E$ の滑らかな完備化として $E$ に等しい。
次数と線型正規性については、五超平面でGrassmann多様体を所期次元に切るので次数は5である。
線型包が三空間以下ならその中で中心超平面が切る長さ5の零スキームが平面以下に入る。
より直接に、Grassmann多様体のPl\"ucker環はCohen--Macaulayで、
所期余次元の五つの一次式は正則列となる。商は深さ2で飽和し、残る一次部分の次元は5である。
従って次数五の完全線型系による埋込みである。
\path{computations/verify_companion_plane_reconstruction.py} は四基点が別々の場合と
単一局所環 $k[y]/(y^4)$ の場合に、核の二列と $\lambda\ne0$ の恒等式を検算する。
\end{proof}

% !TeX root = ../papers/nc-p3.tex
\section{計量を固定した滑らかな判別三次曲線の四割線判定}
\label{sec:marked-fixed-metric-elliptic}
\begin{theorem}\label{thm:marked-fixed-metric-elliptic}
滑らかな三次曲線 $E\subset\mathbf P(V)$ と非退化な計量を位置関係も含めて指定する。
$H=\mathcal O_E(1)$, $L=H^2$ とし、$\nu_2(E)\subset\mathbf P(H^0(L)^*)=\mathbf P^5$ を考える。
計量の逆テンソルを $o\in\mathbf P^5$ とする。$M^2=L$ なる四つの次数三線束ごとに、
$H^0(L)=\mathrm{Sym}^2 H^0(M)$ による階数一点のVeronese曲面 $V_M$ を取る。
以下は同値である。
\begin{enumerate}
\item $o$ が $M\ne H$ の三曲面のいずれにも入らない。
\item $E$ が三交換子独立な有限混合印の二次鉛筆の基点を含む。
\item そのような印を、四基点が被約な射影枠となるように選べる。
\end{enumerate}
これらのとき、この計量と $E$ そのものを固定した階数10の族があり、
添字保持同型類の像は八次元、幾何学的一般点は有限な標準Koszul左右AS正則窓で $g=5$ である。
例外では、$E$ 上のどの跡零二次鉛筆の長さ四の基点にも可換二平面を生む長さ三の部分がある。
この障害は混合印からの構成に関するものである。
\end{theorem}
\begin{proof}[六次埋込みの割線]
次数三線束は平面三次埋込みを与え、そこに二次方程式はないから
$\mathrm{Sym}^2H^0(M)\to H^0(L)$ は同型である。$o$ は元の平面に対し階数3なので、
非被約な場合も含め長さ2以下の二次Veronese像の包には入らない。
長さ四の有効因子 $D$ が $o$ をその包に含む条件は、$D$ 上で消える二次式が全て
$o$ と跡直交する条件である。この四点が射影枠ならその二次式はちょうど二次鉛筆で、
平面での共通零スキームが $D$ に等しい。

まず $o$ を包に含む次数三因子を調べる。相異なる二つ $T,T'$ は交わらない。
交われば和集合の長さは5以下で、その包の次元は長さ引く1であり、
二つの包の共通部分は長さ2以下の共通部分の包になる。これは種数1のRiemann--Rochと
$o$ の階数に反する。交わらないなら、六次因子の包が予想より小さい条件は
$\mathcal O_E(T+T')=L$ である。従って三つ以上あれば全て同じ平方根 $M$ の切断である。
$T=(s)_0$ が条件を満たすことは、$o$ が $s H^0(M)$ を消すこと、
すなわち $s$ が $H^0(M)$ 上の対称形式 $o$ の核に入ることに等しい。
よって無限個の場合に限って $o\in V_M$ であり、その全てがこの核の射影直線をなす。
$o\notin V_H$ は階数3より明らかである。

次に $A\in\mathrm{Pic}^4(E)$ とし、$N=L A^{-1}$ と置く。
$D\in|A|$ が $o$ を張る条件は、その切断が積と $o$ で定まる
$H^0(A)\to H^0(N)^*$ の核に入ることである。この写像の階数は常に2である。
実際双対に核 $u$ があれば $o$ は $uH^0(A)$ を消し、長さ2の零因子 $(u)_0$ の包に
入ってしまう。従って全ての $D$ はPicard曲線上の射影直線束を動く。
上記の三因子が有限なら、一般の $A$ でこの鉛筆は基点自由である。
$x$ が基点なら $H^0(A(-x))$ の三次元空間の像が一次元以下となり、
ある非零 $u\in H^0(N)$ について $o$ は $x+(u)_0$ の包に入る。
有限個の三因子とその支持 $x$ から $N=\mathcal O_E((u)_0)$ が決まるため、
該当する $A$ は有限個である。
標数零で一般切断の零点は相異なる。三つが共線になる場合、第四点 $y$ は $A=H(y)$ で一意に定まり、
基点自由な鉛筆の一般元はそれを避ける。有限個の三因子を含むものも避けられる。
従って一般の $D=\{v_1,\ldots,v_4\}$ は射影枠で、$o$ を四つの平方の線型結合にした際の
四係数は全て非零である。

一方 $o\in V_M$ の場合、任意の $A$ に対し $A=M(x)$ となる $x$ は一意である。
上で求めた三因子の鉛筆に固定点 $x$ を加えると $|A|$ 内に条件を満たす射影直線が得られる。
すでに核の次元は2なので、これで全てである。特にどの $D$ も、$o$ を張る長さ三因子 $T$ を含む。
その局所環全体 $R=H^0(\mathcal O_T)$ は三次元である。$T$ 上単元の一次式で自明化すると、
一次式の $R$ への制限は同型である。さもなくば $T$ は一本の直線内にあり、
その平方の包には階数3の $o$ が入らない。$o$ は $T$ の包にあるからある $\tau:R\to k$ で
$\tau(ab)$ の積対合を一次式空間に引き戻したものとして表される。この対合は非退化なので、
$T$ の平方の線型包を計量で自己準同型に戻すと、$R$ の全乗法作用の空間となる。
実際その全双対は $r\mapsto\tau(sr)$ の形に尽きる。これは可換で単位を含む三次元空間である。
$D$ の二次鉛筆の直交補の跡零部分は、従って可換二平面を含み、
三交換子の独立性とは両立しない。これは重なった点の接方向にも適用される。
\end{proof}

\begin{proof}[射影枠から有限混合印へ]
計量を標準形にし、一般の $D$ の四つの代表を列とする $3\times4$ 行列 $P$ を
$PP^t=I$ となるよう取る。これは上の四係数の非零平方根を使えばよい。
$\ker P$ は非退化な直線であり、その生成元を $h^th=1$ に取る。どの三列も独立なので全 $h_i\ne0$。
$K=k^4$ に $\tau(x)=\sum h_i^2 x_i$ を入れる。$U=\ker\tau$ から
$\operatorname{diag}(h_i)$ により $h^\perp$ を経て $P$ で三空間に移ると、
符号の平方根を除いて $-\tau(xy)$ は元の計量になる。
$K$ の原始冪等元による乗法を $U$ に射影した行列は、非零倍率を除いて
$P$ の同名の列の平方である。実際 $h^\perp$ への直交射影は $P^tP$ で、
座標射影を両側から圧縮すると $P e_i e_i^t P^t$ となる。
跡零部分の包は従って $D$ の二次鉛筆に対応する $G'$ である。
定理\ref{thm:mixed-point-frobenius-algebras}の構成で $m$ を変えてもこの包は同じで、
順序付き基底だけが変わる。相手が可逆な外積三列なので $G'$ の三交換子は独立である。

一般の $m$ で辺が有限であることも確認する。上記定理の二周期式を使い、
$L^tL$ に当たる自己随伴同型を $M_0$、$n=L^ty$ と書く。
星点以外なら $x,y$ は共に非零である。$r=0$ のとき $y\parallel Lx$ なので
$B_{M_0x}x\parallel x$ が必要である。ここで $B_zx$ は $K$ の積の $U$ 成分である。
$M_0=I$ では $x^2\in\langle1,x\rangle$、すなわち四座標は相異なる二値を取る。
二つの値の組分けは有限、跡零条件により各組分けで射影点は一意なので、この軌跡は有限である。
射影的な繊維次元の上半連続性により一般の $M_0$ にも有限性が移る。
$r\ne0$ なら外積の式から $y\cdot y=0$ と $x\cdot n=0$ を得る。
従ってもう一方の式は $K$ で $(s'1+n)x=0$ となる。
$\tau(x)=0$ かつ重みが全て非零なので、$n$ の少なくとも二座標が等しい。
また $n$ は $n\cdot M_0^{-1}n=0$ の非退化二次曲線上にある。
一般の $M_0$ ではこれは六本の等値直線の交点を避け、各直線と有限個の点で交わる。
従って重複する一対が決まって $x$ はその二成分に支えられる跡零の一方向である。
残る星座標も式から決まり、符号を逆にした辺も同じ議論で有限となる。

こうして射影直線束の稠密開部分の全ての $D$ で有限混合印を得る。前節の
指定三次式への滑らかさとエタール復元から、各 $D$ で一般点が全次数正則な
十六次元の表示が得られる（$G'$ の九基底座標を含む）。
全体ではこの二次元の基点族に沿う同じ構成を用いる。二つの交代解の符号は残余群が交換し、
他の基底の開条件と混合印を通る局所成分は前節までと同じである。
一般窓の $G'$ は指定印から決まり、$D$ が復元される。計量と $E$ を共に保つ
射影群は有限である（滑らかな平面三次曲線の射影安定化群は有限）。
よって連結な残余群の次元は10、印を忘れる選択肢は有限であり、像次元は $2+16-10=8$ となる。
\end{proof}

\paragraph{例外の具体例。}
$C=V(17xyz-4(x^3+y^3+z^3))$ は滑らかである。
\[
 D(x,y,z)=\begin{pmatrix}x&2z&2y\\2z&y&2x\\2y&2x&z\end{pmatrix},
 \qquad o=[\operatorname{adj}D(1,0,0)]
\]
と置く。余因子の六つの二次式は独立なので $[\operatorname{adj}D]$ の像はVeronese曲面である。
$C$ 上で $D$ は階数2で、核点 $v$ は $D(v)(x,y,z)^t=0$ より再び $C$ 上にある。
この対応は自分自身が逆なので、余因子の曲面は $\nu_2(C)$ を含む。
元のVeroneseとは異なり $o$ は階数3である。従って曲面の線型因子から来る次数三線束は
$H$ と異なる平方根であり、この固定計量では上の例外が実際に生じる。

% !TeX root = ../papers/nc-p3.tex
\section{可換二平面を含む対称三平面の相手と二次曲線}
\label{sec:marked-commuting-plane-fibers}
\begin{theorem}\label{thm:marked-commuting-plane-fibers}
$G'\subset S$ を既約な対称跡零三平面とし、その自己交換子の包の次元を2とする。
調和二次鉛筆の基点 $B$ を含む滑らかな $E=V(f)$ を固定する。
$G'$ からの混合行列式が $f$ の非零倍となる相手三平面の全体は、
Grassmann多様体のPl\"ucker埋込みにおいて一本の滑らかな平面二次曲線の開部分である。
具体的に、$G'$ の可換二平面は $E$ 上の三次因子 $T$ を定め、$B=T+p$ と書ける。
$p\notin T$ なら二次曲線から二点を除き、$p\in T$ なら一点を除く。
全ての解で $g=5$ である。従って前節の固定計量例外も、\emph{対称側の部分空間}
の段階では被約な三角形および重なった三次因子の両方について扱える。
\end{theorem}
\begin{proof}
三交換子写像の核の直線が可換二平面 $H_0$ を一意に決める。
この平面に巡回行列 $Y$ がある。実際、二固有値で二次最小多項式の元があれば、
その二次元固有空間上で独立な可換元の制限はスカラーではなく、
一般結合の最小多項式は三次である。全元が平方零なら全てが対称な階数1以下であり、
二独立な対称階数一行列の和が階数2になることに反する。
よって $kI+H_0=k[Y]$ は三次元の可換環である。巡回ベクトルで $V$ をこの環と同定すれば、
自己随伴性より計量は非退化な積対合となる。この環のスペクトル $T$ は長さ3で、
一次式の制限を全環とする平面埋込みを持つ。すなわち環の単位を含む全切断で線型系を取り、
計量で双対平面を元の点平面に戻す。前節の積対合の議論により二次Veroneseの線型包は
$kI+H_0$ となる。$T$ は平面を張り、$I_T(2)$ は三次元でイデアル層を生成する。
実際一次で制限が全射なので $H^1(\mathcal I_T(1))=0$、また $H^2(\mathcal I_T)=0$ より
$\mathcal I_T$ は2正則である。$T\subset B\subset E$ がスキームとして従い、$B$ の長さは4である。

前節の $X:V\to G'$ の階数は今は2で、像は $H_0$ となる。実際 $H_0$ の基底を
最初の二元に取ると第三の係数は零である。$p=\mathbf P(\ker X)$ と置く。
$B$ の外で $A(v)$ の二次元核は $X(v),T_v$ がともに非零であればその包であり、
両列の外積と三次小行列の間の等式も成り立つ。両者の可能な共通零集合はそれぞれ
$B$ と $B\cup\{p\}$ の有限集合なので倍率は再び定数である。
特に $p\in B$ である。$p$ の外で $T_v\in G'$ なら、$G'$ 上の交代式の根基は
$X(v)$ が生成し、$T_v$ は $H_0$ 内に入る。この議論は単元成分を取って局所環上でも
成立するため、$B$ と $T$ は $p$ の外で等しく、$E$ の因子として $B=T+p$ である。

$R$ の像を $p$ で消える二一次式 $l,m$ で張り、$\ker R$ を二次式で書けば
$J=H^0(\mathcal I_T(2))$ である。消去二平面上 $R$ の階数が1以下では行列式が
零か可約である。階数2の基底では $l,m$ の基準持上げの三次式 $f_0$ に対し、
全ての三次式が一意な二平面座標 $(j_1,j_2)\in J^2$ を用いて
\[
                     f_0+l j_2-m j_1
\]
と書ける。この線型部分の核は $(l h,m h)$ で $l h,m h\in J$ を満たす一次式 $h$ からなる。
$p\notin T$ なら $l,m$ の制限は単位イデアルを生成し、$T$ 上 $h=0$ となるから核は零。
像は六次元の $H^0(\mathcal I_B(3))$ の全体である。従って任意の $zf$、$z\in k$、
ごとに二平面がただ一つあり、非零倍は $z\ne0$ の場合である。

$p\in T$ なら $T$ は滑らかな曲線の因子なので、$l,m$ が $T$ に生成するイデアルの
零化部分は $p$ の局所環のsocleの直線である。一次式の $T$ への制限は同型なので
核はちょうど一次元である。像の三次式は $p$ で特異で、$f$ 自身は滑らかなので
この特異条件が切る五次元部分空間をちょうど張る。$f_0$ は $p$ で滑らかである。
実際 $0\ne T_p\in H_0$ なので $X(v_*)=T_p$ となる $v_*$ を取れ、$v_*$ は $p$ の方向にない。
$df_0$ は $p$ で $m(v_*)l-l(v_*)m\ne0$ である。従って $zf-f_0$ が像に入る
$z$ は一意かつ非零で、二平面の全解は核のアフィン直線となる。

最後にこれらの直線を二列の外積でPl\"ucker座標に戻す。
$p\notin T$ では $z$ の一次係数の二列 $j_1,j_2$ は独立である。
従属なら $f=l j_2-m j_1$ は可約になってしまう。$p\in T$ では核の二列 $lh,mh$ も独立である。
$R$ の像への射影で定数外積は非零、一次係数は像と核の混合部分で非零、
二次係数は核の二次外積で非零となる。この三係数は独立なので閉包は滑らかな二次曲線である。
無限遠では両列とも核に入り、行列式は零となる。$p\notin T$ のときだけ有限の $z=0$ も除く。
$g=4$ では判別式に直線因子があるため滑らかな解の全ては $g=5$ である。
\path{computations/verify_commuting_plane_fibers.py}は剰余点が三角形に重なる前後での二種類の線型階数を再現する。
\end{proof}

% !TeX root = ../papers/nc-p3.tex
\section{滑らかな三交換子の全復元と楕円曲線の負二倍写像}
\label{sec:marked-elliptic-commutator-doubling}
生成添字の三空間を $D$、行列の三空間を $U$ とし、体積形式で
$\bigwedge^2D=D^*$ と読む。$Q:D\to\mathfrak{sl}(U)$ に対し
$C_Q(s\wedge t)=[Q_s,Q_t]$ を考える。滑らかな平面三次曲線 $E\subset\mathbf P(D^*)$
の線型決定式表示 $A(v)$ に対し、その右核点の写像から引き戻した
$\mathcal O(1)$ を $T_A$、平面の線束を $H$、差を $d_A=T_A-H\in\mathrm{Pic}^0(E)$ と書く。
線束の群則は以下では加法で書く。
\begin{theorem}\label{thm:marked-elliptic-commutator-doubling}
$\det C_Q$ が $E$ を定める全ての $Q$ は以下の方法で復元できる。
\begin{enumerate}
\item $d\in\mathrm{Pic}^0(E)\setminus\mathrm{Pic}^0(E)[2]$ を取り、右核線束が
$H+d$ の決定式表示 $A(v)$ を取る。行の線型包の消去空間を $W\subset D^*\otimes U$ とする。
次の写像が $W$ に値を持つ条件は $Q$ に線型で、解空間は六次元である。
\[
 \Gamma_Q:\mathrm{Sym}^2 U^*\longrightarrow D^*\otimes U,\qquad
          B\longmapsto\big(\mathrm{ax}(BQ_i-Q_i^t B)\big)_{i=1}^3 .
\]
ここで $\mathrm{ax}$ は交代形式を体積形式でベクトルにする写像である。
$\Gamma_Q$ が階数6となる非空な開部分の全てが解である。
逆に $d$ が非零二等分点ならこの開部分は空である。
\item 得られた三交換子自身の右核線束は $H-2d$ である。
従って指定した滑らかな三行三列の決定式表示 $C$ を\emph{独立な左右の基底まで}実現する
補助表示の類は、$-2d=d_C$ の相異なる四解ちょうどである。
対称三平面の部分空間を既に固定した場合には、この四選択・六次元線型解の開条件と
最後の基底の線型比較で、順序付き基底と三交換子解を全て得る。
\end{enumerate}
これは印付き二切断方程式の記述であり、各解の四辺有限性および全次数AS性にはなお判定が要る。
\end{theorem}
\begin{proof}[三交換子から補助表示へ]
$\Gamma_Q$ の核に非零 $B$ があれば、$B$ が特異なら $Q_i$ は共通の直線または平面を保ち、
全ての交換子は特異である。$B$ が可逆でも交換子は奇数階の交代形式に対応して特異である。
従って所要の $Q$ では $W=\operatorname{im}\Gamma_Q$ は六次元である。
その消去三平面を三行の双一次形式 $A$ と書く。$v=s\wedge t$ の点で $A(v)u=0$、
$u\ne0$ という条件は、$v\otimes u\in W$、すなわち $Q_s,Q_t$ の両方を
自己随伴にする非零対称形式の存在である。
定理\ref{thm:rank-eight-finite-symmetry}の二行列の対称化行列式により、
これは $\det[Q_s,Q_t]=0$ と同値である。従って二つの三次式は非零定数倍である。
逆向きにも、指定した $A$ の $W$ と $\Gamma_Q$ の像が一致すれば同じ議論が使える。

$\Gamma_Q$ が $W$ に入る線型条件の解次元を示す。変数 $x\in U^*$ について
$Z(x)=\Gamma_Q(xx^t)$ と書く。その各行 $Z_i(x)$ は $x$ との対合が零となる
三成分二次式である。逆にそのような全ての行は一意な跡零 $Q_i$ から来る。
実際 $x\times Q_i^t x$ の符号倍であり、恒等零なら $Q_i$ はスカラー、
二次の三成分から三次式への $x$ との積は全射で次元差は $18-10=8$ である。
従って求める空間は
\[
 W\otimes\mathrm{Sym}^2 U\ \longrightarrow\ D^*\otimes\mathrm{Sym}^3 U
\]
の核である。この積写像は全射である。双対余核の元は、二階偏微分が全て
$W^\perp$ に入る $D$ 値の対称三テンソルである。ある $v_0$ で $W^\perp\to U^*$ の評価が
可逆なので、他の二方向の評価を二つの自己準同型 $B_1,B_2$ で書く。
この対は既約である。共通不変部分空間があれば $\det(v_0'I+v_1'B_1+v_2'B_2)$ が
ブロックで因数分解してしまう。$v_0$ 成分の三テンソルを $P$ とすれば、
他の二成分は $P$ の一因子に $B_i$ を掛けたもので、対称性によりどの因子に掛けても等しい。
$P$ の最小支持空間は両作用で不変となる。さらに二作用を別々の因子に移せるので
$([B_1,B_2]\otimes I\otimes I)P=0$ である。例えば $B_1B_2$ の右の作用を第二因子へ、
左の作用を第三因子へ移してから逆順に第一因子へ戻せる。
非零なら支持が全空間になり二作用が可換となって既約性に反する。
よって余核は零、核次元は $6\cdot6-3\cdot10=6$ である。
\end{proof}

\begin{proof}[決定式表示の座標]
滑らかな決定式では各曲線点で行列の階数は2である。
左核点の線束は $\mathcal O(-1)^3\to\mathcal O^3$ の余核線束 $T'$ であり、
Euler標数より次数3、自由表示を一つ捻れば $H^0(T'-H)=0$ である。
転置により右核点にも同じことが成り立つ。余因子が共通零点のない二次式で
右・左核の積となるので $T_A+T'=2H$ である。
$T\ne H$ なる任意の次数三線束に対し、二つの平面埋込み $(H,T)$ のグラフの
双一次式は三次元をなす。実際 $H$ の基点自由な二切断と $H^0(T)$ の積は、
鉛筆の完全列と $H^0(T-H)=0$ により $H^0(H+T)$ を張る。
この三つの双一次式を行にすれば表示が得られることを次の座標で確かめる。
特に右核線束が等しい二表示は右核の座標を基底変更した後で同じグラフの双一次式を張り、
独立な左右基底まで一致する。

$E[3]$ の二生成平行移動は $H$ の切断へ射影的に作用する。
その二つの持上げの交換倍率は非自明な三乗根である。可換なら共通固有切断の
次数三因子が自由に作用する九元群で不変となり矛盾する。
立方を1に規格化し、一方を対角、他方を三巡回にした基底を取る。
$T$ は $H$ の平行移動による引戻しなので、同じ二作用の基底が取れる。
$E$ の平面式はこの基底でHesse形になる。対角作用の固有文字は自明でなければならない
（座標点はその不動点で $E$ にない）。三巡回の固有文字も同じくその不動点を用いて自明である。
従って三つの立方の係数は等しく、他は積 $xyz$ だけである。
三関係の空間では対角文字を三巡回が推移的に動かすので各文字が一本であり、
三行を同時巡回させて以下のMoore表示となる。
\[
 A_P(x,y,z)=\begin{pmatrix}
 a x&b z&c y\\ c z&a y&b x\\ b y&c x&a z
 \end{pmatrix},\qquad
 \det A_P=(a^3+b^3+c^3)xyz-abc(x^3+y^3+z^3).
\]
$abc=0$ で行列式が恒等零なら、二つの非零係数の三乗和が零で、
行・列の巡回と対角倍率で外積行列にできる。その右核は平面座標の線型同型となり $T=H$ に反する。
行列式が $xyz$ の非零倍でも $E$ は含まれないので、$abc\ne0$、$P=(a:b:c)\in E$ となる。
逆にこのような $P$ は全て滑らかな決定式を与える。
$P\mapsto d(P)=T_{A_P}-H$ は $E$ から $\mathrm{Pic}^0(E)$ への射である。
実際座標の一つが消える九変曲点でも上の外積表示は曲線上常に階数2なので、相対核線束が取れる。
$O=(0:-1:1)$ では核写像は恒等である。この変曲点を原点にすれば、原点を保つ楕円曲線間の射は
群準同型なので $d(-2P)=-2d(P)$ である。
\end{proof}

\begin{proof}[半周期の除外と符号付き二倍]
$d(P)$ が位数2なら左右核線束が一致する。定数基底変更で曲線上の左右核点を
点ごとに揃えると、行列とその転置の差は同じ核を消す。差は交代行列なので、
恒等零でなければその核写像は平面の線型写像である。像が平面を張るから線型同型となり
$T_{A_P}=H$ に反する。よって定数行列を左から掛けて対称にできる。
$L A_P(v)$ の交代成分を消す九線型方程式の行列式は符号を除き
$(a^3-b^3)(a^3-c^3)(b^3-c^3)$ である。

逆に二つの三乗が等しければ、全添字を3で読み、二添字の逆向き巡回で
係数 $(a,b,c)$ を巡回させ $b^3=c^3$ にできる。
非零成分の添字は $i+j+k=0$、係数は $p_{j-k}$（$p_0=a,p_1=b,p_2=c$）である。
二添字を交換すると添字差の符号が変わる。$b/c$ が三乗根なので、一方の添字の
幾何数列状の対角倍率で任意に選んだ二添字について対称にできる。
特に左右核線束は等しく $2d(P)=0$ である。また二変数側を対称にすれば
$W$ は三次交代行列の全空間と、三独立な対称行列 $B_j$ の包の直和になる。
先の $Z(x)$ の対称部分を $\sum h_j(x) B_j$ とする。$Z(x)x=0$ から
$\sum h_j(x)(x^t B_j x)=0$ である。この三つの二次式は基点自由で、
さもなくば $xx^t\in W^\perp$ が階数1となり、その行を最初に選べば決定式に一次因子が出る。
よって三式は正則列で、二次係数の関係は定数の交代三行列によるKoszul関係に尽きる。
従って三係数 $h_j$ は線型従属で $\Gamma_Q$ は $W$ を張れない。

残りの三乗が相異なる場合の非空性と $d_{C_Q}$ を同時に求める。
$Q_0=\operatorname{diag}(d_0,d_1,d_2)$、$Q_1$ の $(i,i+1)$ 成分を $s_i$、
$Q_2$ の $(i,i-1)$ 成分を $r_i$ とし、他の成分は零とする。自由スカラー $h$ について
\[
\begin{split}
 (s_0,s_1,s_2)&=(-3bc(a^3-b^3),3a^2(b^3-c^3),0)
                  +h(a/b,c^2/b^2,c/a),\\
 (m_0,m_1,m_2)&=(c/b,b/a,a/c),\quad r_i=m_i s_i,\\
 k_i&=r_{i+1}-m_i s_{i-1},\quad
 (d_0-d_1,d_1-d_2,d_2-d_0)=(c k_2/b,b k_0/a,a k_1/c).
\end{split}
\]
最後の三成分の和は零で、跡零条件により $d_i$ は決まる。直接代入で
$\operatorname{im}\Gamma_Q\subset W$ である。
$P'=(a(b^3-c^3),b(c^3-a^3),c(a^3-b^3))$ と置くと、$C_Q(v)$ の行 $i$ は
$A_{P'}(v)$ の行 $-i$ に $s_i k_i/P'_{2i}$ を掛けたものである。
各 $k_i$ は $h$ によらず非零、一般 $h$ で $s_i\ne0$ となる。
\path{computations/verify_moore_commutator_doubling.py} が全不定元を残して
六つの対称行列に対する方程式とこの三行の恒等式を検算する。
$P'$ は $P$ の接線と $E$ の第三交点である。曲線式と接線式への代入でこれを確認でき、
変曲点を除く現在の開集合で $P'=-2P$ である。従って行列式は滑らかで
$\Gamma_Q$ は可逆な像を持ち、$d_{C_Q}=d(P')=-2d(P)$ となる。
同じ $A$ の六次元線型空間の階数6の開部分では、$C_Q$ の右核線束も代数的に動く。
有理空間から楕円曲線への有理写像は一定である（一般の直線に制限し完備化する）。
従って得た値は全ての解で同じである。

最後に二つの対称三平面から計算した基底未指定の混合表示を $C$ とする。$d_C\ne0$ なので四つの $-2$ 逆像はいずれも二等分点でない。
各六空間から上の開条件で $Q$ を取ると独立な左右基底変更 $P,R$ で $C_Q=P C R$ にできる。
これは対称三平面の二つの基底の取り直しにほかならず、交換子の跡条件も満たす。
実際の行列を求める際には $C(v_0)$ が可逆な一点で $R$ を消去し、
$C_Q(v)C_Q(v_0)^{-1} P=P C(v)C(v_0)^{-1}$ を二方向について線型に解けばよい。
可逆解は共通スカラーを除き一意で、非スカラーな自己準同型は先の既約性とSchurの補題に反する。
初めから両基底が指定されている場合には同時共役だけを許して $C_Q(v)P=P C(v)$ を用いる。
\end{proof}

% !TeX root = ../papers/nc-p3.tex
\section{基底を固定した滑らかな三交換子の正確な全ファイバー}
\label{sec:marked-fixed-commutator-adjugate}
\begin{theorem}\label{thm:marked-fixed-commutator-adjugate}
標数零の代数閉体上、順序付き跡零三行列 $C_0,C_1,C_2$ を与え、
$\det\sum v_i C_i$ が滑らかな平面三次曲線を定めるとする。
$[Q_1,Q_2]=C_0$, $[Q_2,Q_0]=C_1$, $[Q_0,Q_1]=C_2$ の順序付き跡零解は
必ず存在する。基底も符号も同一視しない解集合は次の三通りで尽くされる。
\begin{enumerate}
\item 一般の場合は相異なる八解である。
\item 一つの余因子条件が階数2となる場合は相異なる六解である。
\item その階数が1となる場合は相異なる六個の孤立解と、これらと交わらない
滑らかなアフィン二次曲面 $uw-v^2=1$ が全解である。
\end{enumerate}
三通りとも実際に起こる。以下の四つの三次行列の階数で区別でき、解の復元は
楕円曲線の二等分、線型演算およびスカラー平方根で行える。
特に三交換子写像 $\mathfrak{sl}_3^3\to\mathfrak{sl}_3^3$ の一般次数は8である。
\end{theorem}
\begin{proof}[六線型座標における余因子]
前節の四つの $d$ を一つずつ扱う。Moore座標の $P=(a:b:c)$ を取り、
$D=(a(b^3-c^3),b(c^3-a^3),c(a^3-b^3))$、$M(v)=J A_D(v)$ とする。
$J$ は行 $1,2$ を交換する置換（零始まり）である。
六次元線型解を対称三行列 $Z$ で座標化する。
前節の疎な三列式で $s=(a/b,c^2/b^2,c/a)$ とした三列を $U_0$ とする
（その対角行列は零）。同式の $h=0$ を $V_0$ とし、
$W_0=-ac U_0/b^2-V_0/(3ab^4)$ と置く。$U_i,W_i$ は、三つの生成添字と
$U$ の三座標を同時に $i$ だけ巡回したものとする。
\[
            Q(Z)=\sum_{i=0}^2\big(Z_{ii}U_i+Z_{i+1,i+2}W_i\big)
\]
は線型同型である。実際各二列は対角行列を持ち得る添字ごとに三つの異なる文字部分に属し、
そこで $U_i$ は対角部分が零、$W_i$ は三乗が相異なるため非零である。

$\Lambda=\{B:\operatorname{tr}(BM(v))=0\ \forall v\}$ とする。線型同型
$\Psi:\mathrm{Sym}_3\to\Lambda$ があって
\[
        C_{Q(Z)}(v)=\Psi(\operatorname{adj} Z) M(v),\qquad
        \Psi(p p^t)=\kappa\,\operatorname{adj}(J A_P(p))
        \quad(\kappa\ne0)
\]
となる。次の不定元計算で恒等式と可逆性を確認する。
$M(e_0)$ は可逆なので左倍率を $C_{Q(Z)}(e_0)M(e_0)^{-1}$ と定める。
$Z$ の二次単項式21個の係数ごとに、他の二スライスが同じ左倍率である。
$Z=p p^t$ では四次単項式15個の係数は全て零になる。対称三次行列の二次式のうち
$pp^t$ で消えるものは六つの余因子で張られる。実際二次単項式21個から
四次15個への代入は全射で、六余因子は独立である。
$Z_{i+1,i+2}^2$ と $Z_{ii} Z_{i+1,i+2}$ の左倍率係数を取れば、それぞれ符号を除き
$\Psi(E_{ii})$, $\Psi(E_{i+1,i+2}+E_{i+2,i+1})$ となる。
$i=0,1,2$ でこの二列の支持は互いに素な三本の巡回対角であり、
各対の最初の二成分の小行列式は順に
\[
       (a^3-c^3)/(a^4b^7),\quad -c(a^3-b^3)/(a^4b^8),
                      \quad -(b^3-c^3)/(a^3b^8)
\]
である。よって六列は独立、像は跡条件の六空間に等しい。
同じ六列を $p_ip_j$ で足すと上記のMoore行列の余因子になる。
\path{computations/verify_fixed_commutator_adjugate.py}はこれらを全係数の有理恒等式として検算する。

入力を $C(v)=L M(v)R$ と書く。前節の線型比較で可逆な $L,R$ は
$(L,R)\mapsto(\lambda L,\lambda^{-1}R)$ を除き一意である。
この補助類の実際の解は
\[
       S=\Psi^{-1}(RL),\qquad \operatorname{adj} Z=S,\qquad
                      Q_i^{\rm actual}=R^{-1} Q_i(Z) R
\]
で全て得られる。跡零条件により $RL\in\Lambda$ である。
十分性は恒等式による。必要性について、補助表示の行基底は $W$ を変えず、
右基底の変更は $Q$ の同時共役により任意に動かせる。
上の積表示から全ての非特異解の右核点は座標ごと $M$ の右核点であるため、
入力に戻すその射影的な共役変換は $R^{-1}$ に限られる。
実際右核点は平面三次曲線を埋め込み、その全点を固定する射影変換は恒等である。
ゆえに独立な左右基底を誤って許すことなく、表示した全解となる。
$S$ は非零で、階数3なら $Z=\pm\sqrt{\det S}\,S^{-1}$ の二解、
階数2なら解なし、階数1なら合同座標で非退化な二次対称ブロックの
行列式を固定した解全体である。後者では $ZS=0$ から他の行と列が零になる。
\end{proof}

\begin{proof}[異なる二選択は同時に退化しない]
入力曲線とその\emph{具体的な右核点写像}を固定し、ある代表 $M_*$ を選んで
入力を $B M_*$ と動かす。可逆な $B$ は $\Lambda_*=\{B:\operatorname{tr}(BM_*)=0\}$ に入る。
この六空間の階数一行列 $u x^t$ は $x^t M_*(v)u=0$ を全 $v$ で満たす。
その射影軌跡 $E'$ は元の曲線 $E$ に同型である。実際 $u$ は $M_*$ の右核点の
平面三次曲線を動き、$M_*(-)u$ の階数は2なので $x$ は一意である。
後者は変数と右列の役割を替えたMoore表示でも行列式が同じ滑らかなHesse三次式に
比例することによる。両核写像は正則である。この同定による列点の線束を
$T_{\rm col}$、Segre埋込みの線束を $L_6$ とする。

四つの $\Psi$ を同じ $\Lambda_*$ に戻す操作は左右の定数行列によるので階数を保つ。
詳しくは $M_*=L_0 M R_0$ なら $B\mapsto R_0 B L_0$ である。
$\Psi(pp^t)$ の公式により、それぞれのVeronese曲面は $E'$ 全体を含む。
その曲面の平面座標から来る $E'$ 上の線束を $N_d$ とすると $2N_d=L_6$ である。
これらは相異なる。実際 $p\in E$ による $E'$ の表示で、行列の列点は
$A_P(p)$ の右核であり、元の $E$ への同定で使う列点は $A_D$ の右核である。
上の左右基底 $R_0,L_0$ はこの比較では打ち消し合う。
Hesse曲線上の $p\mapsto\ker A_P(p)$ は平行移動である。
変数と列を交換しても同じHesse曲線となり、$P$ を変曲点も含めて動かすと
自己同型の代数的な族で、$P=O$ では恒等である。楕円曲線の自己同型の恒等成分は
平行移動群なので主張が従う。$A_D$ にも同じことが成立する。
ゆえに二同定のずれも平行移動で $\mathrm{Pic}^0$ には恒等に作用し、
$T_{\rm col}-N_d=d$ である。四つの $d$ は相異なるので全四平方根を得る。
さらに $E'$ はこれらの曲面上で平面三次曲線であり、$L_6$ の完全系で埋め込まれている。

$S$ が特異である条件は、$B$ の点が対応する $N_d$ のある三次因子の
$L_6$ による包に入ることである。すなわち $B$ が $H^0(N_d)$ 上で定める
対称形式の核に非零切断 $s$ があり、$sH^0(N_d)$ を消すことと同値である。
$B$ は可逆なので、階数一行列の曲線 $E'$ の長さ2以下の部分の包には入らない。
異なる平方根 $N,N'$ の三次因子がともに $B$ を張るなら、
定理\ref{thm:marked-fixed-metric-elliptic}の割線の議論で両因子は交わらず
$N+N'=L_6$ となるが、$2N=L_6$ に反する。
従って四選択のうち特異なのは高々一つであり、定理の個数が従う。

各例外の実在については、$p\notin E$ の $pp^t$ が可逆な $\Psi(pp^t)$ を与えるので
階数1が実現する。そこから対称行列を階数2へ変形しても一般点で倍率は可逆なままである。
四つの固有な三次超曲面の外は非空開であり、一般の順序付き跡零標的では全てが階数3である。
各二解の局所的な復元はこの開集合でエタールである（標数零の二等分と非零平方根、
および単元小行列での線型解）。従って一般次数も8となる。
\end{proof}

% !TeX root = ../papers/nc-p3.tex
\section{滑らかな判別式の全持上げの成分と固定計量例外}
\label{sec:marked-smooth-discriminant-components}
\begin{theorem}\label{thm:marked-smooth-discriminant-components}
$P=Q$、$G,F$ が対称で、判別三次曲線が滑らかな正規化印付き二切断表示の空間は
既約で次元27である。全ての持上げは分裂混合点を通る片側対称族の同じ閉包に入り、
基底変更を許した全窓は既知の十五次元成分に属する。一般点の窓は全次数正則・四辺有限・
反対階数10、最小周期4で射影自己同型群は自明、添字保持同型類の像は十四次元である。

滑らかな $E$ と計量を共に固定すれば、定理\ref{thm:marked-fixed-metric-elliptic}の
例外外では\emph{全て}の持上げの表示空間は既約で次元18である。
従って八次元の一般正則な同型類の族の閉包に、可換二平面を持つ印での解も全て含まれる。
固定計量例外では次元18の既約な正規化表示成分はちょうど四つである。
例外を与える平方根を $M_0$ とすれば、それらは $-2d=H-M_0$ の四つの位数4の点で区別される。
この最後の四成分の個数は座標と印を固定した二切断表示でのものであり、
各成分の一般AS性や、印を忘れたときの個数を前提としない。
\end{theorem}
\begin{proof}[相手の動きに沿う右核線束]
対称三平面 $G',F'$ の順で混合行列を作り、列側を $F'$ とする。
滑らかな行列式があれば両三平面は既約である。共通不変平面は自己随伴性で共通不変直線を与え、
その直線の代表 $n$ を保つ三平面は $n\odot n^\perp$ を消す跡零三平面そのものとなる。
その二次鉛筆は因子 $n\cdot v$ を持つ。この直線上の全 $T_v$ が三平面に入るので
混合行列式にも直線因子が出る。二空間の和の次元も5である。
前節までの $E$ 上の部分束 $\mathcal U$ を用いると
$\det\mathcal U^*=3H-B$ である。実際 $X(v)\wedge T_v$ の三次成分は、
独立交換子の場合には鉛筆基点のイデアルを生成する。
可換な場合も $X$ と $T_v$ が張る飽和部分束の計算を以下で直接行う。

独立交換子の場合、中心が有限なので節\ref{sec:marked-cubic-linear-system}の
安定性・五切断の証明は非被約 $B$ にもそのまま成立する。
$F'$ に対応する次数二線束 $M\subset\mathcal U^*$ の二切断は
$F'^\perp$ の制限である。そこで $\mathcal U\cap(F'\otimes\mathcal O_E)$ の直線は
$(\mathcal U^*/M)^*$ であり、一般点では混合行列の右核そのものである。
両者は飽和した線束なので、右核の射影写像が与える線束は
\[
             T_C=3H-B-M,\qquad d_C=2H-B-M.
\]
従って $F'$ の楕円五次曲線から $\mathrm{Pic}^0(E)$ への $d_C$ は
平行移動を伴う反転同型である。五基点を除いても、$-2$ の逆像の有限エタール被覆は連結である。

可換二平面の場合には $\mathcal U$ 内で $X(v)$ が生成する線を飽和すると
$-H+p$ となる。$X$ の二つの独立座標は $p$ を切り、滑らかな曲線上で共通零は長さ1である。
$T_v$ との外積の共通因子は曲線上でちょうど $B=T+p$ になることを確認する。
$p$ の一次零を除いた $X$ を $\bar X$ とする。$\bar X\wedge T_v$ は $T$ の外では消えない。
$T$ の各点で $T_v\bmod H_0$ の成分が生成するイデアルは $\mathcal I_T\mathcal O_E$ である
（$\ker R$ は $I_T(2)$ 全体）。よって外積の消失次数はそこでの $T$ の長さ以下である。
また $T$ 上で $T_v$ は $H_0$ にあり、$p$ の外では $X$ と比例する。
$p\in T$ の場合にも $B$ 上では $T_v$ が $G'$ にあり、$G'$ 上の交代三行列は $T_v$ を消す。
この条件は $X\wedge T_v=0$ と同値であり $B$ 上で成立する。
従って $p$ の一次零を引いた後も $T$ の全長で外積は消え、主張の共通因子を得る。
ゆえにこの場合も $\det\mathcal U^*=3H-B$ である。
$F'^\perp$ の二切断は飽和した $X$ への制限により $H-p$ の完全で基点自由な二切断に写る。
両者は $\mathcal U^*$ 内で一般に従属するので、評価像は $H-p$ と同型な部分線束となる。
商線束は $3H-B-(H-p)=2H-T$ であり、先と同じ右核の議論から
\[
                         d_C=H-T .
\]
特にこの値は、相手の二次曲線にも剰余点 $p$ にも依存しない。
\end{proof}

\begin{proof}[被覆と表示の既約性]
滑らかな行列式を持つ順序未指定の三平面対は
$\mathrm{Gr}(3,5)^2$ の非空な開部分で既約、次元12である。
楕円曲線の相対 $\mathrm{Pic}^0$ で $-2d=d_C$ を取ると有限エタール次数4の被覆となる。
上で求めた分裂 $G'$ と固定 $E$ の楕円五次曲線の開部分上でこの被覆は連結なので、
元の被覆も連結であり、滑らかな既約基底の上で既約である。
各点で $d$ の具体的な補助三行を行基底まで選ぶ自由度は $\mathrm{PGL}(U)$ の8である。
実際右核埋込みの座標を選ぶことで尽き、その全点を固定する射影変換はない。
さらに六次元線型空間の $\Gamma_Q$ が単射な非空開部分と、
交換子を混合行列に一致させる二つの順序付き基底の共通逆倍率の一変数を取る。
これらはエタール局所的に束の開部分ないしスカラー主束であり、
そこでの復元は線型方程式と定数階数の核による。従って全体は既約で次元 $12+8+6+1=27$ である。
非空な混合点の変形チャートはこの既約族に稠密である。
実際分裂点の構成で滑らかな三次式に出る方向があり、片側対称な正規化チャートも次元27である。
従って全ての持上げの所属成分と一般点の主張が定理\ref{thm:marked-mixed-component}から従う。

$E$ と計量を固定したとき、$G'$ の四基点は定理\ref{thm:marked-fixed-metric-elliptic}の
$\mathrm{Pic}^4$ 上の射影直線束のうち、二次鉛筆の完全交叉となる開部分で尽くされる。
逆にこの開部分では二次鉛筆の跡直交補が $G'$ を一意に復元し、
可換・非可換どちらでも相手三平面が実在する。非空性は例外外では一般の射影枠、
例外では平面を張る一般の三因子と三本の割線を避ける第四点で確かめられる。
この二次元基底上で、$F'$ の空間はGrassmann因子を五つの行列式の一次条件で切り、
行列式が非零の開部分を取ったものである。全繊維は既約一次元で、
全成分の次元は $2+6-5=3$ 以上である。よって縦成分はなく、全体は既約な次元3である
（一般の幾何繊維も前二定理の楕円曲線または二次曲線の開部分）。
例外外では一般の $G'$ が独立交換子となり、一般繊維の $-2$ 被覆も幾何学的既約である。
有限エタール被覆には縦成分が生じないため、全ての可換特殊値を含め同じ既約被覆に入る。
既知の混合変形の開チャートと同じ成分なので一般点は全次数正則である。
例外では全ての $T$ が同じ $M_0$ の切断であり、$d_C=H-M_0$ は固定された非零二等分点となる。
被覆はちょうどその四つの逆像ごとに分離する。どちらの場合も各表示成分の次元は $3+8+6+1=18$ である。
\end{proof}

% !TeX root = ../papers/nc-p3.tex
\section{節点・尖点三次曲線上の順序付き三交換子の全逆像}
\label{sec:marked-singular-integral-commutators}
\begin{theorem}\label{thm:marked-singular-integral-commutators}
標数零の代数閉体上、順序付き跡零三行列の鉛筆 $C(v)$ の行列式が
特異な整三次曲線 $E$ を定めるとする。特異点で $C$ の階数が1なら
$C=C_Q$ を満たす順序付き跡零三行列は存在しない。階数2なら次が全逆像である。
\begin{enumerate}
\item 節点型では四つの孤立解、二つの孤立解、あるいは二つの孤立解と一つの
滑らかなアフィン二次曲面 $uw-v^2=1$ となる。一般には四解である。
\item 尖点型では二つの孤立解、空集合、あるいは上と同じ一つのアフィン二次曲面となる。
一般には二解である。
\end{enumerate}
各型の任意の局所自由な目標表示類で、掲げた三通りが全て起こる。
以下の線型代数と余因子方程式で、基底を固定した入力ごとに区別し、全ての解を与える。
\end{theorem}
\begin{proof}[整三次曲線上の補助表示の被覆]
定理\ref{thm:marked-elliptic-commutator-doubling}の
$W=\operatorname{im}\Gamma_Q$ と六次元線型解の証明は整な行列式にそのまま適用できる。
補助表示 $A$ の左右基底同値類を調べる。転置の余核層 $T$ は $E$ 上の階数一捩れなし層である。
実際、平面上の長さ一の自由分解により各閉点で深さ1となる。
分解から $\chi(T(n))=3n+3$、$h^0(T)=3$、$h^0(T(-1))=0$ である。
正規化 $\pi:\mathbf P^1\to E$ で局所自由でない場合は
$T=\pi_*\mathcal O_{\mathbf P^1}(2)$ しかない。特異点の完備局所環を $R$、正規化を
$\bar R$ とする。分数イデアルとしての局所加群を $I\bar R=\bar R$ に規格化すると、
導手 $\mathfrak c$ は $I$ に入る。$I/\mathfrak c$ は二次元の $\bar R/\mathfrak c$
（節点では $k\oplus k$、尖点では $k[\epsilon]/(\epsilon^2)$）をイデアルとして生成する
部分空間である。一次元なら単元で生成され $I$ は自由、二次元なら $I=\bar R$ である。
後者では正規化上の線束の次数がEuler標数から2になる。
局所自由の場合は正規化で次数3、節点の二値の貼合せは乗法群、
尖点の一次ジェットの貼合せは加法群で指定される。$h^0(T(-1))=0$ により恒等元を除く。

節点を $xyz=y^3+z^3$、正規化を $(r^3+1:r^2:r)$ とすると、右核点は
$(r^3+\lambda:r^2:r)$ であり $\lambda\in k^*\setminus\{1\}$ である。
非自由表示は $\lambda=0$ で共通因子を除いたものとなる。双一次関係の基底は
\[
 N_\lambda=\begin{pmatrix}0&-z&y\\-y&x&(\lambda-1)z\\
 -z&-(\lambda-1)y&\lambda x\end{pmatrix}.
\]
尖点では $xz^2=y^3$、正規化 $(1:r^2:r^3)$ を取り、右核点は
$(1+\lambda r:r^2:r^3)$、$\lambda\ne0$、非自由なら $(1:r:r^2)$ である。
$r$ の倍率による生成添字の基底変更で、非零の $\lambda$ を1にしてよい。二表示は順に
\[
 K=\begin{pmatrix}y&-x&-x\\z&-y&-x\\0&z&-y\end{pmatrix},\qquad
 K_0=\begin{pmatrix}y&0&-x\\z&-y&0\\0&z&-y\end{pmatrix}.
\]
実際各行は正規化上のグラフを消す。自由な三次表示では一次の双一次関係は自由分解により三本で、
その全空間が元の行空間を復元する。明示した各代表の行は独立で、行列式は $E$ の式の非零倍である。
以上は目標 $C$ 自身の左右基底同値類にも適用できる。
\end{proof}

\begin{proof}[負二倍、六座標と余因子の計算]
対称行列 $Z=\left(\begin{smallmatrix}P&A&B\\ A&T&D\\ B&D&U\end{smallmatrix}\right)$ を取る。
六空間 $\Gamma_Q\subset W$ のうち、三行列の次の六成分を指定して $Q(Z)$ を一意に定める。
成分の添字は零始まりである。節点では
\[
 ((Q_1)_{02},(Q_2)_{01},(Q_2)_{02},-\lambda(Q_1)_{21},(Q_2)_{12},(Q_1)_{12})
 =
 (-\lambda((\lambda-1)A+U),-(\lambda-1)B+\lambda T,
 (\lambda-1)^2P+\lambda D,-\lambda A,-\lambda B,(\lambda-1)P).
\]
尖点 $K$ では第四成分を $-(Q_1)_{21}$、第六成分を $(Q_2)_{10}$ とし、
右辺を $(U,-A-B-T,-B-D,A,B,P)$ とする。
これらを代入して六空間の線型方程式を解くと各成分は $Z$ と係数 $\lambda$ の多項式
（定数分母3を許す）である。六成分指定の行列式は符号を除き節点で
$\lambda^5(\lambda-1)$、尖点で1となるので、六次元性により全解が得られる。
節点の参照行列と尖点の参照行列を順に
\[
 M_\lambda=N_{1/\lambda^2}\operatorname{diag}(1,\lambda,\lambda^2),\qquad
 M_K=\begin{pmatrix}-y&x+y&-x+y\\ y+z&-x+y&0\\-y&x+y+z&-x+z\end{pmatrix}
\]
とする。尖点の右核は $((r-1)^2:r^2-r^3:r^3)$、従って $1-2r,r^2,r^3$ の別基底である。
以上の右核の貼合せは節点で $\lambda\mapsto\lambda^{-2}$、尖点で
$\lambda\mapsto-2\lambda$ となる。後者は $r$ の倍率を元に戻した式である。

$\lambda\ne\pm1$ の節点および自由な尖点で、線型同型
$\Psi:\mathrm{Sym}_3\to\{L:\operatorname{tr}(LM(v))=0\ \forall v\}$ は
\[
                   C_{Q(Z)}(v)=\Psi(\operatorname{adj}Z)M(v)
\]
で計算できる。検算の手順を示す。$C_{Q(Z)}(0,1,0)M(0,1,0)^{-1}$ の
$TU,PU,PT,BD,AD,AB$ の係数を順に六列とする。
$\operatorname{adj}Z$ の対角三成分と上三角非対角三成分 $01,02,12$ をこの六列に掛けると、
$Z$ の全二次係数が元の左倍率と一致する。他の二スライスも同じ左倍率と $M$ の積である。
六列の成分を行優先に並べ、節点では行 $0,1,2,3,4,5$、尖点では行 $0,1,2,3,4,7$ を取ると
小行列式はそれぞれ $\lambda^2(\lambda-1)(\lambda+1)$ と $-3$ である。
節点の六列の分母は $\lambda$ の冪だけなので恒等式を $\lambda=-1$ に特殊化できる。
そこで $M_{-1}$ は行列式が恒等零ゆえ補助解は一つも整な目標を与えない。
非自由表示 $N_0,K_0$ の双方では六空間の任意の二基底要素 $q,q'$ について
$[q_1,q'_2]+[q'_1,q_2]=0$ となる。この場合 $C_Q(1,0,0)=0$ で行列式は二変数のみになる。
これらの全係数の厳密検算は \path{computations/verify_singular_integral_commutators.py} にある。
従って非自由な目標に至る補助表示はなく、自由な目標の補助類は上の負二倍の逆像で尽きる。
\end{proof}

\begin{proof}[指定目標と全例外]
自由な目標を各候補について $C=L M R$ と書く。定理\ref{thm:marked-fixed-commutator-adjugate}と同じく
$S=\Psi^{-1}(RL)$ とすれば、解は $\operatorname{adj}Z=S$ の各逆像につき
$R^{-1}Q(Z)R$ である。ここでも右核像の全点を固定する射影変換は恒等であり、他の共役はない。
$S$ の階数3,2,1に対応する解は二点、空、一つのアフィン二次曲面である。

節点の二選択の関係を明示する。$M_{-\lambda}=M_\lambda\operatorname{diag}(1,-1,1)$ ゆえ、
一方の $S$ の座標を $(P,T,U,A,B,D)$ とすれば、他方の $S'$ の座標は
$(P-4D,T,U,2U-A,2T-B,-D)$ となる。先の六列への代入で
\[
          4\det\Psi_\lambda(S)=(\lambda+1)^2\det S-(\lambda-1)^2\det S'
\]
を得る。左辺は零でないので少なくとも一方の $S$ は可逆である。
両行列式の非消失は稠密開条件である。$S=E_{11}$（中央の対角成分、零始まり）
では $\det S'=-4$ なので、階数1の例を得る。階数2の軌跡は階数1を閉包に含むので
$\det S'\ne0$ の階数2の例もある。尖点では同じ $S=E_{11}$ で $\det\Psi(S)=1$、
やはり階数2および3に動かせる。これで固定した各自由表示類における全ての可能性が従う。
\end{proof}

% !TeX root = ../papers/nc-p3.tex
\section{直線と横断二次曲線の三交換子と一次元ファイバー}
\label{sec:marked-transverse-conic-commutators}
\begin{lemma}\label{lem:marked-squarefree-symmetrizer}
$\det A$ が被約な平面三次曲線なら $\Gamma_Q$ が $A$ の行の消去空間に入る
線型解空間は次元 $6+r$ である。$r$ は行空間内の射影的な階数一テンソルの個数であり、
各々は行列式の相異なる直線成分に対応する。実際の三交換子表示で $\det C_Q\ne0$ なら、
そのようなテンソルはちょうど、$Q_n$ がスカラーと階数一行列の和になる方向 $[n]$ に一つずつ対応する。
\end{lemma}
\begin{proof}
定理\ref{thm:marked-elliptic-commutator-doubling}の積写像の双対で、一つの評価を恒等、
残りを $B_1,B_2$ とする。一般直線は被約三次曲線と相異なる三点で横断するので、
一般の $B_1,B_2$ の線型結合は固有値が相異なる。その固有基底で三因子間の作用の移動ができる
対称三テンソルは $\sum c_i e_i^{\otimes3}$ に限る。
両作用とも移動できることから $c_i\ne0$ なら $e_i$ は共通固有ベクトルであり、
逆にそのような立方は全て余核の双対に入る。各共通固有ベクトルは評価を恒等にした
行空間の階数一テンソルを与え、共通固有文字は一般結合上で既に相異なる。
最後に $n\otimes b$ が行空間に入る条件は全ての $h\in U^*$ に対し
$b,h,Q_n^t h$ が従属することである。$h$ と $h+b$ を比較すると $Q_n^t b$ は $kb$ に入る。
商 $U^*/kb$ の全直線を保つので商上でスカラーとなり、$Q_n$ は所要の形である。
スカラーだけなら二つの交換子が消えるから除かれる。階数一補正は一意である。
\end{proof}

\begin{theorem}\label{thm:marked-transverse-conic-commutators}
標数零の代数閉体で $\det C$ が直線とそれに横断する滑らかな二次曲線の積であるとする。
直線を $x=0$ とする座標を選ぶ。順序付き三交換子の全逆像は次の互いに素な二部分から得られる。
\begin{enumerate}
\item $Q_0$ に階数一補正がない部分が存在するのは、目標が両交点で局所自由で
右核の成分次数が $(1,2)$ のときに限り、その条件は十分でもある。
逆像のこの部分は四点、二点、または二点と一つのアフィン二次曲面である。
\item 階数一補正を持つ部分は目標の直線上の鉛筆から高々二つの旗 $(K,kw)$ と以下の
四変数の行列階数判定、線型方程式、スカラー平方根で例外値を含めて得られる。
ここで $K$ は平面で、$w\in K$ は要求しない。両核が直線上で動き、旗が半単純な場合には、
固定倍率ごとにアフィン線型空間上の一つの二次式、または二行三列の階数一以下の条件で尽きる。
\end{enumerate}
局所自由で次数 $(1,2)$ の目標の一般点で、全逆像は四つの孤立解と
$\mathbf G_m$ に同型な二つのアフィン二次曲線である。
\end{theorem}
\begin{proof}[補助表示の被覆]
生成添字の基底変更で $E:x(yz-x^2)=0$ とし、成分を
$(0:s:t)$、$(st:s^2:t^2)$ で正規化する。前節と同様に
$T(-H)$ はEuler標数と $h^0$ が零の階数一捩れなし層である。
局所自由でない交点を分離した部分正規化からの線束の直像となる。
両方で自由なら $T$ の成分次数は $(0,3),(1,2),(2,1)$ のいずれかで、
両端の貼合せの倍率比 $\lambda\ne0$ が残る。$(1,2)$ では $\lambda=1$ を除く。
実際 $T(-H)$ の成分次数和は零、片方が2以上なら他方で零にして両交点で消える切断がある。
一次数 $1,-1$ のときそのような切断はなく、次数 $0,0$ では両定数の貼合せが判定する。
一交点だけ分離すれば次数和が一つ下がり $T$ は $(0,2),(1,1)$、貼合せ倍率は基底変更で消える。
二点を分離すれば $T$ の両次数は $(0,1)$ である。これで全ての可能性を覆う。

直線上の右・左核点の次数の和は、分離した交点の個数を $j$ として $2-j$ である。
余因子の二次式は両核の積で、階数一の交点では直線上にちょうど一次の共通零を持つ。
実際特異点の完備局所環で非自由な余核は二枝の正規化の直和であり、
その平面上の最小自由表示は二局所方程式の対角、残る単位ブロックは核に寄与しない。
従って左核が直線上で定点なら前補題の階数一の場合に入る。
左核が動き右核が定点の場合は、自由な $(0,3)$ または一交点非自由の $(0,2)$ である。
貼合せた切断 $H^0(T)$ の基底と双一次関係を取ると、これらの代表は
\[
 \begin{pmatrix}-x&y&0\\z&-x&0\\\lambda y&z&-x\end{pmatrix},
 \qquad \begin{pmatrix}-x&y&0\\0&-x&y\\0&z&-x\end{pmatrix}
\]
（または交点を交換）である。実際右核はそれぞれ直線で定点、二次曲線で
$(s^2t:st^2:\lambda s^3+t^3)$、$(s^2:st:t^2)$ を与える。
第二表示では直線の定点は $e_0$ であり、二次曲線の二端のうち一端だけこれに一致する。
両表示で六線型解を解くと全 $Q_i$ がそれぞれ $ke_2,ke_0$ を保つ。
従って交換子鉛筆は恒等的に特異である。

残る自由な $(1,2)$ の代表とその交換子の参照行列は
\[
 A_\lambda=\begin{pmatrix}-x&y&0\\z&x(1/\lambda-1)&-y/\lambda\\0&z&-x\end{pmatrix},
 \quad M_\lambda=\begin{pmatrix}-\lambda x&y&0\\z&x(\lambda^2-1)/\lambda&-y\\0&-\lambda z&x\end{pmatrix}.
\]
$A_\lambda$ の核は直線で $(s:0:\lambda t)$、二次曲線で $(s^2:st:t^2)$
（直線上では第二行を用いる）。$M_\lambda$ の核は順に
$(s:0:t)$、$(s^2:\lambda st:\lambda^2t^2)$ であり、目標の貼合せは $\lambda^{-2}$ となる。
前節と同じ $Z$ の座標を使い、$Q$ の六成分
$((Q_1)_{12},(Q_2)_{01},(Q_2)_{02},-(Q_0)_{21}/\lambda,(Q_2)_{12},(Q_2)_{21})$ を
$(-U,-(\lambda-1)D-P,-B,A,D,(\lambda-1)T)$ とする。
方程式 $\operatorname{im}\Gamma_Q\subset W$ のこの解は $Z$ からの線型同型になる。
$C_Q(e_0)M_\lambda(e_0)^{-1}$ から前節と同じ六単項式の係数で作る
$\Psi_\lambda$ に対して、再び $C_Q=\Psi_\lambda(\operatorname{adj}Z)M_\lambda$。
行優先の六行 $0,1,2,3,4,6$ の小行列式は $-\lambda^4$ である。
六列は多項式であり、$\lambda=-1$ では $\det M_{-1}=0$ なので解がない。
その他は各平方根ごとに先の $C=L M R$ の手順が全ての実際の解を与える。
右核の二次曲線の全点を固定する射影変換は恒等なので共役の一意性も成り立つ。
他方の平方根の $S'$ は $(P,T,U,-A,-B,2P-D)$、
\[
      4\det\Psi_\lambda(S)=\lambda(\lambda+1)^2\det S-
                                   \lambda(\lambda-1)^2\det S' .
\]
少なくとも一方は可逆である。$S=E_{00}$ なら $\det S'=-4$ なので
階数一とその近傍の階数二も起こる。
以上の六空間・余因子恒等式・小行列式は
\path{computations/verify_transverse_conic_commutators.py} で検算できる。
\end{proof}

\begin{proof}[階数一補正の全復元]
$N=w b^t$、$K=\ker b^t$ とする。$[N,-]$ の像は
$\{Y:Y K\subset kw,\ \operatorname{tr}Y=0\}$ である。
包含は積から明らかで、両空間は四次元である（半単純型も冪零型も中心化環は
全行列内で五次元）。従って候補 $(K,kw)$ は $C_1K+C_2K\subset kw$ の全解である。

具体的には直線上の一般点の右核は常に $K$ に含まれ、左核は $w$ を消す。
各核の線型包を、余因子の共通因子を除いた多項式ベクトルの係数から求める。
どちらかが平面全体を射影的に張れば候補はない。右核のベクトル包が二次元なら
それが $K$ であり、$C_1K+C_2K$ が一次元のときだけその像を $kw$ とする。
左核の包が二次元のときは転置で同様である。両核が定点 $kr,k\ell$ なら
$U/kr\to\ker\ell^t$ の二行二列の正則鉛筆を得る。可逆な一スライスで割り、他の
スライスの各固有直線を $K/kr$ とすればよい。二スライスが比例なら
$\det C$ の直線に沿う一次係数は一次式の平方の倍になり横断性に反する。
これで旗は高々二つである。

旗ごとに $N$ の代表を一つ固定し、$Q_0=\rho(N-\operatorname{tr}(N)I/3)$、
$X_i=\rho Q_i$ $(i=1,2)$ とする。線型に $[N,H_1]=C_2$、$[N,H_2]=-C_1$ を解き、
跡零中心化環の基底を $Z_j$ $(1\le j\le4)$ とする。
四変数で $X_1=H_1+\sum t_j Z_j$ を自由に動かし、八行五列の拡大前の係数行列と右辺
\[
       D(t)=\big([X_1,Z_1],\ldots,[X_1,Z_4],-C_0\big),
              \qquad k(t)=-[X_1,H_2]
\]
を作る（成分を跡零八座標に並べる）。
$\operatorname{rank}D=\operatorname{rank}(D\mid k)$ を満たす階数層で
$D(v_1,\ldots,v_4,h)^t=k$ を線型に解き、$h\ne0$ の解を全て取る。
$X_2=H_2+\sum v_j Z_j$、$\rho=\pm\sqrt h$ で全逆像となる。
各層の非零小行列式を選ぶと自由変数と有理式で陽に書ける。
入力の任意の特殊化についても四変数だけの小行列式による条件であり、
$h\ne0$ を満たす解の有無は特解と線型核の最終座標で判定できる。
必要十分性は三つの交換子への代入による。

$w\notin K$ なら $N=\operatorname{diag}(1,0,0)$ とする。
$X_i$ を $1+2$ ブロックに分けて上右行を $\alpha_i$、下左列を $u_i$ とすれば
これらは最初の二つの既知の交換子で固定されている。
右下ブロックと左上スカラー倍の差を $D_i$ とすると、跡零条件により
任意の $D_i$ は対角部分を唯一に定める。$h=\rho^2$ は
$h(C_0)_{00}=\alpha_1u_2-\alpha_2u_1$ を満たす非零数を全て取る。残る方程式は
\[
 \alpha_1D_2-\alpha_2D_1=h(C_0)_{01},\quad
 D_1u_2-D_2u_1=h(C_0)_{10},\quad
 [D_1,D_2]=J=h(C_0)_{11}-u_1\alpha_2+u_2\alpha_1 ,
\]
ここで $00,01,10,11$ はブロックを表す。$J$ は跡零である。
$J\ne0$ なら $\operatorname{tr}(D_iJ)=0$ の二つを最初の線型条件に加える。
$\mathfrak{sl}_2$ の跡対合とブラケットが与える交代三形式は体積形式であり、
$J^\perp$ の二元の交換子は $kJ$ に入る。よって最後は $J$ の非零成分一つに
対する二次式だけである。$J=0$ なら $D_1,D_2$ の跡零部分が従属する条件となる。

最後に次数 $(1,2)$ の局所自由目標の空間は既約である。
固定した成分対で貼合せ倍率、右核座標の射影基底、跡零左倍率の六空間の可逆開部分を
取ればよく、横断する成分対も一つの射影軌道をなす。四つの孤立解の条件は既に稠密開である。
半単純な一旗、非零な $h$ の一意性、線型解の次元2と滑らかな二次曲線、相異なる二無限遠点は
いずれも以下の点で確認できる開条件である。$Q_0=\operatorname{diag}(2,-1,-1)/3$ と
\[
 Q_1=\begin{pmatrix}4&1&1\\-1&-2&3\\0&3&-2\end{pmatrix},\qquad
 Q_2=\begin{pmatrix}-1/3&-2&3\\1&2/3&-1\\-2&-3&-1/3\end{pmatrix}
\]
を取ると、目標の二次因子は $38x^2-260xy+308xz+27y^2-30yz-3z^2$、
$h=1$、$J=\left(\begin{smallmatrix}-6&-3\\3&6\end{smallmatrix}\right)$。
$D_2$ の下行を $(g,k)$ として線型消去すれば残る式は
$27g^2-24gk+234g-6k^2-140k+459=0$ になる。両核は直線で次数1であり、全交点で自由である
（$\alpha_i$ 同士と $u_i$ 同士が各々独立）。
従ってこの非空開条件を同時に満たす一般目標では、$\rho$ の二符号ごとに
射影二次曲線から二点を除いたものが得られる。
\end{proof}

% !TeX root = ../papers/nc-p3.tex
\section{接する二成分と三直線を含む被約三次判別式の全持上げ}
\label{sec:marked-reduced-cubic-commutators}
\begin{theorem}\label{thm:marked-reduced-cubic-commutators}
標数零の代数閉体上、判別式が任意の被約な平面三次曲線である順序付き跡零三行列 $C$ の
全三交換子逆像が以下と先の各判別式型の手順で得られる。直線成分に階数一補正を持つ解は各直線の
圧縮対と四変数の階数層から復元する。補正を一つも持たない解について、直線とそれに接する
滑らかな二次曲線、三角形、共点三直線の場合に残る補助表示は、
それぞれ局所自由で多重次数 $(1,2),(1,1,1),(1,1,1)$ の表示だけである。
目標も同じ型でなければこの部分は空である。この型なら三角形で二つ、それ以外で一つの
補助候補があり、各候補の正確な逆像は線型に得る対称行列の階数3,2,1に従って
二点、空、滑らかなアフィン二次曲面である。

直線と接する二次曲線で局所自由な $(1,2)$ 型の一般の目標の全逆像は、
二つの孤立解と、各々 $\mathbf A^1$ に同型な二つのアフィン二次曲線である。
\end{theorem}
\begin{proof}[補助表示の線束と成分次数]
線束の貼合せと双一次関係で決定式の行空間を得る方法は
定理\ref{thm:marked-singular-integral-commutators}と同じである。
直線で左核が一定なら補題\ref{lem:marked-squarefree-symmetrizer}が階数一補正を与える。
全成分が直線で左核がいずれも動くとする。左核を与える余核層の正規化上の線束の
各次数は1以上、その和は3以下である。後者を確かめるには各特異点の完備局所環 $R$ で
加群の分数イデアルを、正規化を生成し単位1を含む $I$ に規格化する。
正規化を生成すれば無限体上で全ての枝で単元な一般元を選べる。
線束として貼り合わせた場合とのEuler標数の差から、正規化の次数和は
$3-\sum\operatorname{length}(I/R)$ となる。従って等号で全て局所自由、
三つの次数は全て1である。余因子の左右核積から右側も同じ次数である。

接する場合は $E:x(xz-y^2)=0$、成分を $(0:s:t),(s^2:st:t^2)$ とする。
接点で $R$ は両枝の値と一次ジェットが一致する部分環で、$\bar R/R$ は
共通ジェット $k[r]/(r^2)$ 上の長さ二の巡回加群である。従って $R\subset I\subset\bar R$
で中間にあるのは値だけを合わせる節点環一つである。
多重次数の候補は、前節の二交点を順に分離した候補と同じである。
実際導手に沿う消失は今度は相異なる二点での消失または単一点での二次消失であり、
$\mathbf P^1$ 上の切断数と二つのジェット評価階数は同じである。
余因子の直線上の共通零の次数は三局所型で順に $0,1,2$。例えば中間の節点環は
二枝の値を合わせるので二つの生成元で表示され、平面局所座標 $u,v$ で
枝を $v=0,v=u^2$ とすれば、基底 $(1,1),(0,u)$ に対する $v$ の作用は
$\left(\begin{smallmatrix}0&0\\u&u^2\end{smallmatrix}\right)$ である。
直線 $v=0$ 上の非単位ブロックの成分の共通零は一次となる。
完全正規化の場合は二枝の対角表示で次数2を得る。

左核が動き右核が直線上で定点の補助解は、自由な $(0,3)$ と中間の $(0,2)$
の場合だけである。順に代表は
\[
 \begin{pmatrix}y&-x&0\\z&-y&0\\0&z-e y&-x\end{pmatrix},\qquad
 \begin{pmatrix}y&-x&0\\0&y&-x\\z&0&-x\end{pmatrix}.
\]
前者の右核は二次曲線上で $(s^3:s^2t:t^3-e st^2)$ であり、後者では $(s^2:st:t^2)$、
直線上はいずれも $e_2$ である。前者のジェット倍率 $e$ は $s/t$ の倍率で $0,1$ にできる。
いずれも六線型解の全行列が $ke_2$ を保ち、非零判別式を生じない。
以上で補正のない解は主張の三型に限られる。
\end{proof}

\begin{proof}[三種類の余因子公式]
接する二成分の自由な $(1,2)$ の右核は、二次曲線上 $(s^2:st:t^2)$、
直線上 $(0:s:t+e s)$ と書ける。$e\ne0$ であり、$e=0$ は平面線束そのものになる。
倍率で $e=1$ とする。以下に補助行列 $A$ と目標の参照行列 $M$ を並べる。
\[
 A=\begin{pmatrix}y&-x&0\\z&0&-x\\0&y+z&-x-y\end{pmatrix},
 \qquad M=\begin{pmatrix}x+y&-x&0\\-x+z&2x&-x\\0&-x+z&x-y\end{pmatrix}.
\]
$M$ の直線上と二次曲線上の右核は $(0:s:t)$ と $(s^2:s(s+t):(s+t)^2)$ であり、
直線の正規化と二次曲線の正規化でのジェットの差は $-2$ となる。

三角形では各頂点で左右二枝の値を貼り合わせる。各成分のスカラーで二頂点の倍率を
規格化すれば、残る一つの倍率が平面線束に対する相対値を分類する。代表は
\[
 A=\begin{pmatrix}x&\lambda z&0\\0&y&x\\y&0&z\end{pmatrix},
 \qquad M=\begin{pmatrix}x&-z&0\\0&\lambda y&-x\\y&0&-\lambda z\end{pmatrix}.
\]
ここで $\lambda\ne0,-1$。三辺で右核の非零二座標を取り、頂点で巡回して比を掛けると
$A$ の相対倍率は $-\lambda$（巡回を逆にすればその逆数）である。
$M$ の行列式の二単項積の係数比は $-1/\lambda^2$ であり、
対角行列による左右の基底変更でこの値を $\lambda$ とする $A$ の形にできる。
従って二候補となる。$\lambda=1$ では $\det M=0$ で補助解はない。

共点三直線を $(r:0:1),(0:r:1),(r:r:1)$ で取る。平面の局所関数は三枝の値が一致し、
第三枝の一次係数が前二枝の和である。二次以上では各枝の係数を独立に指定できる。
自由な多重次数 $(1,1,1)$ の線束は値を一致させた後、第三の一次係数を前二つの和に
$e$ 倍の共通の値を加えたものにする。$e\ne0$ を倍率で1にしてよい。
三右核の埋込みは $(r:0:1),(0:r:1),(r:r:1+r)$ であり、代表は
\[
 A=\begin{pmatrix}y&-x&0\\z&x&-x\\0&x+z&-y\end{pmatrix},
 \qquad M=\begin{pmatrix}y&-x&0\\x+z&-x&-x\\0&-x+y+z&-y\end{pmatrix}.
\]
$M$ の三右核は $(r:0:1+r),(0:r:1+r),(r:r:1)$ なので相対ジェットは $-2$ である。
二つの加法型で倍率を戻せば $e\mapsto-2e$ となり候補は一意である。

三表示とも六空間を $Z$ によって線型座標化すると
$C_{Q(Z)}=\Psi(\operatorname{adj}Z)M$ を得る。再現のため、対角二つと六つの
非対角による跡零の基底を、行優先の行列単位からスカラー部分を引き $E_{22}$ を除いて作る。
$Q_0,Q_1,Q_2$ のこの順に並べた二十四変数で $\operatorname{im}\Gamma_Q\subset W$ を
簡約行階段形にし、自由列の順に変数を置く。前節の $Z=(P,T,U,A',B,D)$ の記号で
（$A'$ は $01$ 成分）、その六自由変数は三型の順に
\[
 (U,-T+D,-D,A'-B,B,P),\quad
 (U,-D,(\lambda-1)A',T,-B,\lambda P),\quad
 (U,-T,-D,P,B,A').
\]
三角形ではまず不定元 $\lambda$ で計算する。得る $Q(Z)$ は多項式で、
$\lambda=1$ でも線型階数6である。他は自由列指定から同型が分かる。
左倍率の六単項式抽出には順に $(1,1,0),(1,1,1),(1,2,1)$ で評価する。
六列の小行列式は順に $-1,-\lambda^5,-1$。各係数の恒等式と特殊化の計算は
\path{computations/verify_reduced_cubic_commutators.py} にある。
$M$ の具体的な右核像は二次曲線または複数の異なる直線を含み全点を固定する
射影変換は恒等なので、従前の固定目標の左右比較と余因子逆像がそのまま使える。
\end{proof}

\begin{proof}[補正を持つ解と接する一般ファイバー]
直線ごとに前節の $K,kw$ を求め、前節の四変数の階数層を用いる。
三直線で複数の補正がある場合は和集合を取り、または直線を順序付けて最初の
補正のみを採用する。核が両方定点の場合の二行二列商鉛筆が比例する場合もここでは起こり得る。
そのとき可逆スライスで割った他の作用がスカラーなので、$K/kr$ は
射影直線の全点を動かす。一般スライスが階数2であることは被約性から従う。
これで交点の階数が零になる場合も含め全候補を得る。

接する二成分の一般ファイバーの型を調べる。前節の半単純な一旗の解法で
$G=hC_0$ とする。一般座標で $\alpha_1=(1,0),\alpha_2=(0,1)$、
$u_1=(0,m)^t,u_2=(n,0)^t$ としてよい。実際二行と二列が各々独立で、
$(\alpha_i u_j)$ の対称部分が非退化なら、その二つの等方直線を添字の基底に取り、
二行を平面の双対基底にできる。$J=\left(\begin{smallmatrix}j&k\\l&-j\end{smallmatrix}\right)$
と置くと $G_{11}=J+\operatorname{diag}(-n,m)$。直線での残余二次曲線の式は
符号・非零倍率と直線座標の変更まで
\[
              -mk s^2+j(m-n)st-nl t^2 .
\]
実際 $(s,t)\operatorname{adj}(G_{11})(nt,ms)^t$ を展開すればよい。
一方解のアフィン平面の無限遠で $c=(D_1)_{10}$、$f=(D_1)_{11}$ とする。
線型条件により $(D_1)_{01}=-kc/l+j(n-m)f/(nl)$、
$(D_2)_{10}=j(n-m)c/(mk)-lf/k$、$(D_2)_{01}=f$ となる。
従って $[D_1,D_2]_{00}=0$ の無限遠式は非零倍率まで
\[
                   nk c^2-j(n-m)cf+ml f^2 .
\]
二式の判別式は一致する。ゆえに接する場合は無限遠点が一つである。
一般条件の実現を $n=2,m=1,j=4,k=1,l=2$ と
$G=\left(\begin{smallmatrix}1&2&-2\\-2&2&1\\-3&2&-3\end{smallmatrix}\right)$ で見る。
残余二次曲線は $-30x^2-15xy-2xz-(y+2z)^2$ で滑らか、$h=1$、
解の平面の式は $(g-v)^2-10g-4v+11=0$ で滑らかである。
被約な接する成分対は一軌道で、自由な $(1,2)$ の目標は貼合せ、右基底、可逆な跡零左倍率から
既約である。よってこの非空な開条件と余因子の階数3が同時に一般点で成立し、
二符号それぞれで射影二次曲線から一点を除いたものを得る。
\end{proof}

% !TeX root = ../papers/nc-p3.tex
\section{非被約三次式を含む非零交換子判別式の全復元}
\label{sec:marked-nonreduced-cubic-commutators}
\begin{theorem}\label{thm:marked-nonreduced-cubic-commutators}
標数零の代数閉体上、三重直線または二重直線と別の単純直線を行列式とする
順序付き跡零三行列 $C$ に対し、全ての $C_Q=C$ は以下で得られる。
各支持直線の法方向の $Q_i$ が階数一行列の跡零部なら圧縮対、零なら一交換子の
二次式で復元する。これらのどちらも起こらない部分は、可逆スライスで割った目標の
残り二作用が既約な場合に限り候補を持つ。候補はただ一つの補助表示になり、
具体的な六座標の線型比較で対称行列 $S$ が得られる。この部分の正確な全逆像は
$S$ の階数3,2,1に応じ二点、空、一つのアフィン二次曲面である。

従って先の各判別式型と合わせ、行列式が恒等零でない全ての順序付き三交換子に
復元手順がある。特に一点を指定した有限窓で $q=3,g>3$ のとき、
正規化した対称側の六行列に対して残る交換子は全判別式型を含め復元できる。
この主張は二切断方程式についての被覆であり、個別の候補の全次数正則性は別に要る。
\end{theorem}
\begin{proof}[補助表示と二つの既約鉛筆]
$\det C_Q\ne0$ なら $\Gamma_Q$ は単射であり、その行消去から得る $A$ は
\emph{重複度も込みで} $\det A$ が $\det C_Q$ の定数倍になる。
実際 $\det A(v)$ は、九空間において $\Gamma_Q$ の六列と $v\otimes U$ の三列を並べた
行列式の定数倍である。$v=s\wedge t$ のときこれは $s,t$ への評価で
$D^*\otimes U$ を $(v\otimes U)$ で割った六次元商への写像の行列式となる。
体積による同定で商の行列式の倍率は一定である
（$s,t$ を基底の最初の二元に拡張し外積の基底変換を取ればよい）。
定理\ref{thm:rank-eight-finite-symmetry}の六変数恒等式からこれは $\det[Q_s,Q_t]$ である。

座標を $\det A$ が $x^3$ または $x^2(x+y)$ の倍になるように取り、
$A=xI+yJ+zD$ と規格化する。二作用が可約なら同時三角化できる。実際不変直線か不変平面の
残る二次元部分でも全行列式は一次式の積である。その線型固有文字の一つを引くと
二スライスの包は階数1以下で、共通核か共通像により同時不変直線がある。
従ってもとの鉛筆に共通左固有ベクトルがあり、行空間中の階数一テンソルを与える。
補題\ref{lem:marked-squarefree-symmetrizer}末尾の階数一テンソルと $Q_n$ の同値の証明は
被約性を用いず、ここでは $Q_n=0$ も許す。逆にこのケースでは必ず可約となる。

既約な場合の左右基底同値類は、固定した $x,y,z$ のままで次の各一変数 $k\ne0$ で尽きる。
\[
\begin{array}{c|c|c}
 &J&D\\ \hline
x^3&E_{12}+E_{23}& k(E_{21}-E_{32})\\
x^2(x+y)&E_{12}+E_{33}&
\left(\begin{smallmatrix}k&-k&0\\0&-k&1\\k^2&-k^2&0\end{smallmatrix}\right)
\end{array}
\]
この表の行列単位だけは $1,2,3$ 始まりである。
三重直線では $J,D$ の全結合は冪零であり、一般元は階数2である（階数1以下の空間は可約）。
まずこの一般元を $J=E_{12}+E_{23}$ として計算する。
$\operatorname{tr}D=\operatorname{tr}(JD)=\operatorname{tr}(J^2D)=0$ から
$D=\left(\begin{smallmatrix}a&b&c\\d&e&f\\0&-d&-a-e\end{smallmatrix}\right)$。
既約性で $d\ne0$、$\operatorname{tr}(JD^2)=d(2a+e)=0$。
$I+sJ$ の共役で対角を消すと $\operatorname{tr}D^2=0$ から $b=f$、
$I+tJ^2$ の共役で両方消え、$\det D=0$ から $c=0$ となる。
得た型ではどの非零の結合も階数2なので、最初から指定した $y$ 方向が使える。
$k$ は $[J,D]$ の二重固有値で一意である。

二重直線でも最初は一般の $y$ 方向（$y$ の係数を1にした $J+tD$）を使う。
二重零固有値が半単純なら階数1、一般にそうなら上と同じく可約である。
従って $J=E_{12}+E_{33}$ にできる。三つの跡条件から
$D=\left(\begin{smallmatrix}a&b&u\\0&-a&v\\w&h&0\end{smallmatrix}\right)$。
$\operatorname{tr}D^2=\operatorname{tr}(JD^2)=0$ から
$uw+vh=-a^2$、$vw=a^2$ を得る。$v=0$ なら第二座標行が共通左固有、
$w=0$ なら第一座標列が共通右固有なので $a,v,w\ne0$。
$J$ の中心化群の上三角剪断とブロック間の倍率で $u=0,v=1$ とし、
$\det D=0$ を用いると上表の $k=a$ となる。残余中心化群は $a$ を変えない。
この型の $J+tD$ は全 $t$ で階数2であり、指定した $y$ に戻せる。
二つの表の型が既約であることも $J$ の左右固有直線を $D$ が保たないことから分かる。
\end{proof}

\begin{proof}[固定基底での逆像]
二つの既約鉛筆の六空間性は定理\ref{thm:marked-elliptic-commutator-doubling}の
双対三テンソルの証明に従う。そこで必要なのは二作用の既約性だけである。
参照目標を $M=xI+yJ'+zD'$ とすると、二型の順に
\[
\begin{split}
 J'&=E_{12}-E_{23}/2,\quad D'=-k E_{21}/2-k E_{32},\\
 J'&=\begin{pmatrix}0&1&-1/k\\0&1&-1/(2k)\\0&0&0\end{pmatrix},\quad
 D'=\begin{pmatrix}0&-k&1\\ k/2&-k&1/2\\0&-k^2&k\end{pmatrix}.
\end{split}
\]
この表示の抽象パラメータは二型とも $-k/2$。
第一の場合は $\operatorname{diag}(1,1,-2)$、
第二では列 $(1,0,0)^t,(0,-1,-2k)^t,(1,1,0)^t$ の基底で $J'$ を表の形に戻すと分かる。
よって既約な目標を規格化すれば $k$ は唯一に求まる。
前節と同じ六座標 $(P,T,U,A,B,D_0)$ と自由列の順で六空間に以下を代入する。
\[
\begin{split}
 &(-U,kP+D_0,B,kA,D_0,kT),\\
 &(-2k^2P+k^2A-kD_0+U,-kP+kT-B+D_0,P-A,-kA+kT-D_0,A-T,k^2T).
\end{split}
\]
これらは $k\ne0$ で同型である。$C_Q(e_0)$ の六余因子係数を $\Psi$ とすると
$C_Q(v)=\Psi(\operatorname{adj}Z)M(v)$ で、六行 $0,1,2,3,4,6$ の小行列式は
$4k^{10},-4k^{10}$ である。\path{computations/verify_nonreduced_commutators.py} で全係数を検算する。
$C=L M R$ の基底比較の自由度は逆倍率のみである。実際 $M(e_0)=I$ なので、
二つの比較の差は $J',D'$ の共通中心化に属し、既約性によりスカラーとなる。
従って $\operatorname{adj}Z=\Psi^{-1}(RL)$ による実際の解とその階数分類が従う。

補正がある各支持直線で法方向を最初の添字に取れば、
$C_1K+C_2K\subset kw$ による全ての圧縮対と定理\ref{thm:marked-transverse-conic-commutators}の
線型復元が使える。直線上の一般階数が1以下でも、全 $K\in\mathbf P(U^*)$ について
$C_1K+C_2K$ の次元を小行列式で判定すればよい。次元1では $kw$ はその像、
次元0では全 $[w]\in\mathbf P(U)$、次元2以上なら解なしである。
$Q_0=0$ なら $C_1=C_2=0$ で三重直線型に限られる。
逆にその入力では $C_0$ が可逆である。$Q_0$ が非スカラーなら、
巡回の場合その中心化環は可換、非巡回の場合その中心化環は
階数一補正の像直線を保つため、いずれもそこにある二元の交換子が可逆になることはない。
ゆえに $Q_0=0$ が必要である。残りは巡回な跡零 $X$ について
$\operatorname{tr}(XC_0)=\operatorname{tr}(X^2 C_0)=0$ の全てを取り、
$Q_1=X$、$[X,Q_2]=C_0$ の跡零解のアフィン二平面を取れば尽くされる。
最後に有限窓では定理\ref{thm:no-three-bracket-compression}とその直前の可換交叉の排除、
および定理\ref{thm:marked-two-commutators}が $q=3,g>3$ の非零判別式または二交換子復元を保証する。
\end{proof}

% !TeX root = ../papers/nc-p3.tex
\section{両側一般の高階数窓の線型復元と八固有直線による全同型}
\label{sec:marked-mixed-bracket-reconstruction}
\begin{theorem}\label{thm:marked-mixed-bracket-reconstruction}
一点を指定した窓で、二つの三平面が同じ空間を張るという仮定を外す。
正規化座標の交叉交換子を
\[
  C_{bd,ac}=[Q_b,P_d]_{ac}=-[G_a,F_c]_{bd},
  \qquad C\in\mathfrak{sl}(V)\otimes\mathfrak{sl}(U)
\]
とする。任意の $C$ に対して印付き二切断方程式の全解は、下記二つの
線型写像の各アフィン逆像と階数一テンソル錐との交叉で決まる。
分裂混合点を通る十五次元成分の一般の窓には、両方の写像が単射となる印がある。
その印では $C$ と $C$ の転置配置のそれぞれから六行列が有理演算だけで逆倍率まで復元できる。
従ってこの開部分の正規化印付き表示は、実際の $C$ の三十二次元軌跡の上で
二つの逆倍率を選ぶものとなる。一般点は $g=q=6$、反対階数16である。

この一般点で $C$ を跡対合により $\mathfrak{sl}(V)\to\mathfrak{sl}(U)$ とすれば
可逆で、$C^*C$ は相異なる八固有値を持つ。この条件と両線型写像の単射性を満たす全ての窓では、
後述の有限個の数表の一致が印を保つ添字保持同型の必要十分条件である。
印を忘れた同型では、相手の有限な印について同じ比較を行えばよい。
\end{theorem}

\begin{proof}[混合交換子からの二つの因子]
まず $C_{ij}$ を一方の三行三列の跡零行列として読む。第一の線型空間を
\[
 L_C=\{X\in\mathfrak{sl}_3^3:
       \operatorname{tr}(X_i C_{kj}+X_k C_{ij})=0\ (i\le k,\ 1\le j\le3)\}
\]
とし、第二を $R_C=L_{(C_{ji})}$ とする。$C_{ij}=[X_i,Y_j]$ なら跡の巡回性で
$X\in L_C,Y\in R_C$ である。係数のブラケットにより
\[
  m_C:L_C\otimes R_C\longrightarrow (\mathfrak{sl}_3)^9,\qquad
                  X\otimes Y\longmapsto ([X_i,Y_j])_{ij}
\]
を作る。この線型写像の $C$ 上の逆像の階数一テンソルは全て、かつそれだけが
求める組である。非零の各テンソルは一つの非零成分を選んで行と列を読むと有理的に分解でき、
自由度は $(X,Y)\mapsto(hX,h^{-1}Y)$ となる。$C=0$ の場合の零テンソルには
一方が零の全組を加える（有限窓の独立な三列ではそもそも不要である）。
他方の組にも $-C_{ij,kl}$ の二添字対を入れ替えて同じ判定を行う。
これで任意の例外次元にも二切断方程式そのものの必要十分条件を得る。
各写像が単射であれば線型逆像は空か一点なので、両方の線型整合性と各一点の
二次小行列式の消失が、実際の $C\ne0$ の明示的条件である。

定理\ref{thm:marked-point-commutators}の残余群 $(A,B,s)$ は $C$ に
$s(\operatorname{Ad}_B\otimes\operatorname{Ad}_A)$ と作用する。
固定した $C$ の二つの逆倍率は $A,B$ 各々のスカラーによって独立に動く。
よって単射な開集合では $[C]$ の $\mathrm{PGL}(V)\times\mathrm{PGL}(U)$ 軌道が、
印を保つ添字保持同型の完全不変量となる。
\end{proof}

\begin{proof}[実在する十五次元成分での単射性と単純固有値]
分裂混合点 $G=D,F=-D,Q=P=K$ で四組の跡零行列の九十六座標を使う。
この点の正規化印付き局所環は滑らかで次元34である。実際、元の二切断テンソル軌跡は
命題\ref{prop:split-mixed-transverse-curve}のSchur微分と二十余法線で滑らかな次元76である。
単純な二十点のうち一つを選ぶ被覆は局所でエタールであり、対応する四点は
各辺の階数3の部分評価から復元する。基底でこれを標準点にする表示へ移る際は
$\mathrm{GL}_4^4$ の枠を用いる。標準点の各直線を保つ枠は四因子で次元52、
二つの反対対合を単位行列にすると18条件、補文字の三つずつの持上げは
跡零化の12条件で一意に決まる。これらは可逆行列の選択・線型方程式で滑らかな
次元22の枠の繊維となる。従って正規化表示の次元は $76+22-64=34$ であり、
九十六変数の方程式の微分階数は62である。
同成分の一般点の $g=q=6$ と反対階数16は既知の横断曲線上の非消失から従う。

素数1009で微分の可逆な六十二次小行列を取り、残りの三十四自由座標を動かす。
この六十二式だけの完備局所環は滑らかで、標数零の全方程式の滑らかな同次元の商と一致する。
従って整数係数の自由座標を $1009$ の倍数だけ動かし、この小行列の単元性で逐次補正すると
実際の標数零の局所族の $1009$ 進点が得られる。より具体的には、形式的陰関数の係数は
$\mathbf Z_{1009}$ に入り、高次係数の収束と残りの全恒等式を保ったまま代入できる。
\path{computations/verify_mixed_bracket_reconstruction.py} は一つの自由方向で72桁まで
九十六変数を求め、全方程式と各段の剰余を検算する。
二つの読み方の各々で左右の十八行の階数はともに18、従って両核とも次元6、
積写像の三十六列は独立である。検算は単元ピボットを選び、残る行が一斉に素数で割れるときは
共通に割って有効桁数を減らす。残りが有効桁まで零となる前に全行にピボットが得られる。
同じ点で、八次の $C$ の行列式の付値は8、$C^*C$ の特性多項式の判別式の付値は62である。
いずれも精度72未満なので真の非消失を証明する。よって各条件はこの印付き成分で稠密開で、
印を忘れた成分の一般点にもそのような印が存在する。二倍率の復元から $C$ の像の
次元は $34-2=32$、射影化して31となる。この検算は個別の閉点の全次数AS判定には使用しない。
\end{proof}

\begin{proof}[有限数表の同型判定]
$C^*C$ は跡対合に関して自己随伴で、八固有直線は互いに直交し各々非退化である。
その順序を取り、最初の固有値が1になるよう $C$ の全体倍率を二通り選ぶ。
各固有直線で $\operatorname{tr}(u_i^2)=1$ となる二つの代表を選び、
$v_i=C u_i$、$\lambda_i=\operatorname{tr}(v_i^2)$ とする。有限表は
\[
    (\lambda_i)_i,\qquad (\operatorname{tr}(u_i u_j u_k))_{ijk},\qquad
                            (\operatorname{tr}(v_i v_j v_k))_{ijk}
\]
である。基底順序、全体倍率、八つの符号を全通り許す。
表の一致から $u_i$ 同士と $v_i$ 同士を送る線型同型は、両方の跡対合を保つ。
さらに単位を単位へ送って延長すると全行列積を保つ。実際、積のスカラー成分は
二元の跡の三分の一、残る成分は三元の順序付き跡を各基底のノルムで割ったものである。
従って全行列環の二つの単位的同型となり、行列単位の像からそれぞれ共役である。
二共役は $C$ の作用を可換にするので上記の完全不変量が一致する。逆は各操作の共変性による。
向きを持つ三元の跡を用いるので転置や添字の交換を余分に同一視しない。
無印の入力には第一の窓で一つの適格な印を固定し、第二の窓の同じ位置の全単純点を調べる。
同型は印と全ての階数条件を運び、逆にどれかでの印付き同型は無印の同型を与える。
\end{proof}

% !TeX root = ../papers/nc-p3.tex
\section{高階数成分の全二十印の分離、周期と有理完全不変量}
\label{sec:marked-twenty-rational-invariants}
\begin{theorem}\label{thm:marked-twenty-rational-invariants}
分裂混合点を通る十五次元成分の非空な稠密開部分では、全二十印で
交叉交換子の両側の線型復元が成立する。
基点を固定した添字保持同型の完全不変量は有理演算だけで求まり、
無印の添字保持同型も一つのモニック二十次式と次数20未満の係数多項式の一致で判定できる。
同じ開部分で射影自己同型群は自明、最小周期は4である。
これらの結論には以下の付値環で全次数正則な実現がある。
素数1009で指定した三十四自由係数の合同近傍の\emph{全て}の持上げが
標準Koszul左右AS正則であり、そこに互いに添字保持非同型な十五パラメータの解析的族を取れる。
代数的数係数でのそのような正則閉点も稠密に存在する。
\end{theorem}

\begin{proof}[二個の行列が与える有理な基底]
前節の $C:\mathfrak{sl}(V)\to\mathfrak{sl}(U)$ を使う。
$T=\operatorname{Tr}(C^*C)\ne0$ とし、両側で $K=C^*C/T,CC^*/T$ を取る。
$\operatorname{Tr}$ は八次元作用の跡であり、以下の $\operatorname{tr}$ は三次行列の跡である。
跡対合の互いに双対な基底 $e_i,e^i$ について、三次行列
\[
 J(H)=\Big(\sum_i H(e_i)e^i+e^i H(e_i)\Big)_0,\quad
 A=J(K),\quad B=J(K^2),\quad m=\operatorname{tr}(A^2),\quad d=\det[A,B]
\]
を作る。添字0は跡零部である。双対基底の縮約だから $J$ は基底の選択に依存せず、
各側の共役と可換である。$K$ は $C$ の全体倍率にも不変である。
入力側を $V$、出力側を $U$ と添字付けし、
$\tau=\operatorname{tr}(A_U C(A_V))$ が非零なら実際の $C$ を $C/\tau$ で規格化する。
次の九語が各側で全行列環の基底である開条件を課す。
\[
           I,A,B,A^2,AB,BA,B^2,A^2B,AB^2 .
\]
各基底での行列積表と、単位を零に送って延長した $C/\tau$ の九語間の線型写像の表を $H_p$
とする。どの表も固定順の基底であり、有理演算だけで得られる。
等しい積表から全行列環の単位的同型が二つ得られ、いずれも内部共役である。
第三の表との一致で $[C]$ を運ぶので、前節の両側単射条件の下でこの表は完全である。
また自己同型では各側の $A,B$ が個別に固定され、九語が基底なので両共役とも恒等である。
残るスカラーは前節の逆倍率の項で尽き、印を固定する射影自己同型はない。
\end{proof}

\begin{proof}[二十個の印の区別と四つの開始位置]
印 $p$ を一つ先に運んだ位置で、元の星係数から読む四組は
$(G',F',Q',P')=(Q^t,P^t,F,G)$ となる。
よって $C_{ij,kl}$ は $C_{kl,ji}$ に移り、二つ先では $C_{ji,lk}$ に移る。
一つ先の八次元写像は $C^*$ の入力に転置を合成したものであり、
$J$ の積を対称化したため行列対 $A,B$ は転置に関しても共変である。
従って四位置の入力側の二数 $(m,d)$ は順に
\[
 (m_V,d_V),\qquad (m_U,-d_U),\qquad (m_V,-d_V),\qquad (m_U,d_U).
\]
全二十印の合計四十個の $m_V,m_U$ が相異なり各 $d$ が非零なら八十組は相異なる。
これで異なる印への自己同型と非四倍の位置移動による同型を同時に排除できる。
四周期自体は標準Frobenius周期から従う。

印を忘れる際に二十回の代数閉体上の比較を除くこともできる。
位置を固定して二十点の有限エタール座標環を $D$ とする。$m=m_V\in D$ の二十値が
相異なるから $D=k[m]$ であり、その特性多項式 $f$ は次数20で分離的である。
$H_p$ の各表成分は $D$ 上で枠に依存せず定まり $h(m)$、$\deg h<20$ と一意に書ける。
局所的に枠を選んで求めた表は共役不変性により貼り合わさるので、分裂体を選ぶ必要はない。
従って $f$ と順序付きのこれらの係数多項式は無印の窓から有理的に求まる。
同型ならこれらは一致し、逆に一致すれば幾何学的な一根を選んだ印付きの積表が等しいので
添字保持同型を得る。このとき開始位置の移動を商にしていない。
\end{proof}

\begin{proof}[稠密開条件と全次数の合同近傍]
\path{computations/verify_twenty_mixed_bracket_invariants.py} は前節の実際の持上げから作る
四テンソルを $1009^{14}$ を法として用いる。剰余の全二十印の四位置の点を、各射影点の
一座標を固定し隣接二評価の零条件で持ち上げる。十二変数のJacobi行列に各々単元小行列があり、
持上げ後は全条件を検査する。各点を枠の第一行にし、二つの反対評価を単位に戻して
$Q,P$ を読む。持上げのスカラー部は交換子に影響しない。
各印の両線型復元の階数は両方で $18,18,36$、九語も各々基底である。
この階数消去で失う精度は最大四桁、$T$ は全て単元、$\tau$ の付値は3、
$d$ の付値は高々9、四十個の $m$ は既に $1009^3$ を法として相異なる。
これらは真の非消失であり、有限エタールな全印に条件を課しても元の十五次元成分に
空でない開集合を与える。剰余の二十点の被約性と軌道の証明は
命題\ref{prop:split-mixed-transverse-curve}の配置を用いる。

同じ三十四自由係数をこの点の値に $1009^{14}$ を法として合同に取れば、
前節の単元小行列で各々唯一の持上げを得る。各解は剰余で分裂混合窓になり、
そのWalsh符号平方行列は倍率を除き第一行 $(-2,1,1)$ の巡回行列で、余因子は $3\ne0$ である。
奇標数のWalsh判定と命題\ref{prop:unrestricted-good-reduction}より、
この近傍の\emph{各}分数体窓が幾何学的に全次数正則である。
反対階数16も前節の横断曲線から一般に成り立つ。必要ならその非消失開集合に近傍を縮められる。
自由係数を有理整数で動かせばエタール射の零次元ファイバーゆえ持上げも代数的で、
この格子は三十四次元の正規化局所成分にZariski稠密である。

最後に逆倍率も含む十九次元残余群の軌道に横断する局所切断を取る。
有理計算した積表の微分の核は正確にその軌道の接空間である。
実際、双対数環上で表が一定なら、二つの九語基底から積と単位と $C/\tau$ を保つ
一次同型が得られる。全行列環の導分は内部であり、二つの単射な積写像は
単元小行列によって双対数環でも単射なので、残りは二つの逆倍率だけである。
ここで $[C]$ の倍率を戻す一自由度も残余群に含まれる。印付き射影安定化群は自明である。
従って三十四個の自由座標のうち十九個をその軌道接空間で独立になるように固定すれば、
残る十五変数上で十五個の表成分の微分が可逆になる。
付値体の逆写像定理で十分小さい開球上ではこの十五成分が単射である。
この球内の異なる印は上の $m$ の剰余で既に区別されるから、
同じ位置で元の印を別の印に移してこれらの窓が同型になることもない。
\end{proof}

% !TeX root = ../papers/nc-p3.tex
\section{五次元交叉交換子の二つの逆像と中心一次枝への分岐}
\label{sec:marked-five-dimensional-bracket-fiber}
\begin{theorem}\label{thm:marked-five-dimensional-bracket-fiber}
同時に二つの座標枠を取って次の形にできる非零交叉交換子テンソルを考える。
$C_{ii}$ は跡零対角空間を張り、$i\ne j$ では $C_{ij}$ は残る添字 $k$ の
$K_k=E_{k+1,k+2}-E_{k+2,k+1}$ の非零倍である（添字は三巡回）。
この形の実現可能な $[C]$ は、独立な二つの共役まで
\[
 C(t)_{ij}=[K_i+tD_i,K_j-tD_j],\qquad
 D_i=E_{i+1,i+2}+E_{i+2,i+1},\qquad t(t^2-1)\ne0
\]
で尽きる。$[C(t')]$ と $[C(t)]$ の軌道が一致する必要十分条件は
$t'^2\in\{t^2,t^{-2}\}$ である。

この軌道上の二切断窓を\emph{全て}復元すると、スカラーと印を忘れた正則な添字保持同型類は以下で尽くされる。
$x=t^2$ とし、$W(t)$ を命題\ref{prop:split-mixed-transverse-curve}の既知の窓とする。
\begin{enumerate}
\item $x\ne-1$ のとき、二つの最小周期1の類 $A(t),A(-1/t)$ が常にある。
$x\ne-3,-1/3$ なら、これらと相異なる最小周期2の二類 $W(t),W(-1/t)$ もある。
$x=-3,-1/3$ では後二候補だけが非正則になる。
\item $x=-1$ なら全候補が一つの正則な周期1の類 $A(t)=W(t)$ になる。
\item 直線の端 $x=1$ で非対角の全 $C_{ij}$ が消えた指定値も含めると、
そこにある正則な添字保持同型類は $W(1)$ ただ一つである。
\end{enumerate}
ここで $A(t)$ は以下の三生成量子多項式環の中心一次拡大である。
この枝全体で $A(t)\simeq A(s)$ と $s=\pm t$ は同値である。
$x=0$ にある交代三空間は五次元像を持たず本定理の軌道条件の外にある。
\end{theorem}
\begin{proof}[軌道の認識と二枚の逆像]
$C_{ij}=c_{ij}K_k$ と書く。$C_{ij}=[X_i,Y_j]$ なら前節の跡の線型条件が成り立つ。
$\operatorname{tr}(X_i C_{ij})=0$ により $X_i$ の交代射影は $a_iK_i$ となる。
また $\operatorname{tr}(X_i C_{jj})=-\operatorname{tr}(X_j C_{ij})=0$ である。
右辺の消失は $i\ne j$ なら $C_{ij}$ が $C_{ji}$ に比例することによる。
従って全 $X_i$ の対角は零である。同様に $Y_i$ の交代射影は $b_iK_i$、対角は零となる。
相異なる $i,j,k$ の跡条件は $a_i c_{kj}+a_k c_{ij}=0$ である。
一つの $a_i$ が零なら全て零となる。全て零の場合、$Y$ の交代部分も零なら
交叉交換子の対角成分が出ない。$b_i$ の一つが非零なら同じ跡条件で全て非零となり、
対称成分を $X_i=\sum h_{li}D_l$ と書いて $[X_i,Y_j]$ の対角成分を $i\ne j$ で取ると
$h_{ji}=0$。残る $X_i=h_{ii}D_i$ に対し非対角の対称成分が
$h_{ii}b_j[D_i,K_j]=0$ となり、やはり解がない。$b_i$ についても同様である。
従って二つの三つ組とも全て非零で、上の三式の積から
$c_{01}c_{12}c_{20}=-c_{10}c_{21}c_{02}$ である。
入力側の対角共役の比の平方を選べば $C_{ji}=-C_{ij}$ を揃えられる。

$C_{ij}=\epsilon_{ijk}\ell_k K_k$ とすれば $a_i=a/\ell_i$、$b_i=b/\ell_i$ となる。
$i\ne j$ の対角比較で $[K_i,D_i]=H_i$ と $H_j$ が独立なので、
$X_i,Y_i$ の対称部分は各々 $h_iD_i,m_iD_i$ に限られる。
$[K_i,D_j]=-[K_j,D_i]\ne0$ より $a m_j\ell_j=-b h_i\ell_i$。
従って $h_i=h/\ell_i,m_i=m/\ell_i$、$am=-bh$ である。
$C_{ii}=(am-bh)H_i/\ell_i^2$ は二次元を張り $\sum C_{ii}=0$ なので、
三つの $\ell_i^2$ は等しい。符号の積が1の変更を対角共役で行い、三つの $\ell_i$ を等しくする。
残る二つの非零倍率の比を $(t^2-1)/t$ の非零定数倍で表せば、全体倍率を除き表示の $C(t)$ になる。

具体的に $C(t)_{ii}=-2t H_i$、$C(t)_{ij}=(1-t^2)[K_i,K_j]$ $(i\ne j)$。
$X_i=aK_i+hD_i$、$Y_i=bK_i+mD_i$ の全解は
\[
            bh=t,\quad am=-t,\quad ab+hm=1-t^2,\qquad (ab-1)(ab+t^2)=0 .
\]
両側の読み方に同じ計算が使える。元の $Q=K+tD,P=K-tD$ の選択を替えると
$Q=K-D/t,P=-t^2K-tD$ になり、元の $G=tK+D,F=tK-D$ の選択を替えると
逆倍率まで $G=K-tD,F=-(K+tD)$ になる。互いに独立な二倍率は残余群で吸収できる。
両方を替えた窓は $W(-1/t)$ である。実際 $C(-1/t)=-C(t)/t^2$ である。
$x=-1$ では各側の二根が一致するので集合として選択は一つである。

入力の形は枠の選択の存在で定義され、各実現上では枠自身が復元できる。
両読みの像五空間は対角跡零と交代行列の和で、その跡直交補は $D_i$ の包である。
この三空間の行列式は座標積の二倍であり、三角形の三頂点の行列の核が元の三座標軸になる。
従って枠を忘れても $i\ne j$ の
$\operatorname{tr}(C_{ij}C_{ji})/\operatorname{tr}(C_{ii}^2)=(x-1)^2/(16x)$ は同じである。
逆に $t\mapsto-t$ は二枠の同時奇置換、$t\mapsto-1/t$ は上記全体倍率で実現する。
\end{proof}

\begin{proof}[一方だけを替えた窓]
一つの側を替えた表示を $G=P=K-tD$、$Q=-F=K+tD$ とする。
先の一位置移動公式で $(G',F',Q',P')=(-P,-Q,-Q,P)$ となるから、
残余の $A=B=-I,s=-1$ で元に戻り最小周期1である。
最初の二位置を同じ四生成に読む際、第二位置の補三文字の符号を反転する。
するとその二次関係は、星文字 $h$ が中心で、巡回的な $(i,j,k)$ ごとに
\[
             (1-t)x_i x_j-(1+t)x_jx_i=-h x_k
\]
となる。実際 $Q_{b,ac}=-G_{a,bc}$ であり、二つの星を用いた六関係を足し引きすればよい。
他方だけの変更は $t\mapsto-1/t$ で得る。$a=1-t,b=1+t$ とし $ab\ne0$。
中心環 $k[h]$ 上で $yx,zy,zx$ を昇順に書き換えると、唯一の重なり $zyx$ は
両順とも $\frac ab xyz+\frac hb(x^2-y^2+z^2)$ になる。
補文字の長さ、次いで辞書式順序を用いれば停止する。従って $h$ は非零因子であり、
その商はパラメータ $q=b/a$ の三生成量子多項式環、中心一次拡大は標準Koszul左右AS正則である。
$q\ne1$。中心文字を規格化し生成元を一斉に倍率変更すれば
$X_i X_j-q X_jX_i=(1-q)h X_k$ と書ける。

$h=0$ の点対応は量子多項式環のもので、三座標直線を含み射影平面を張る。
一方 $h\ne0$ なら中心性により両点は同じであり、$h=1$ で $X_i X_j=X_k$ を解くと
原点と積1の四符号点だけである。従って無印の添字保持同型でも、
正次元の点軌跡の包から無限遠平面が決まり、残る五点の置換は全位置で同じである。
五点は射影枠で、原点は四点の重心という唯一のアフィン従属で識別される。
そのため全位置の射影変換は一致し、四符号点を保つ線型変換は積1の符号と三軸置換の全二十四個である。
巡回の向きを保てば $q'=q$、逆にすれば $q'=q^{-1}$ で、直接関係を送れば両者は十分でもある。
これは $s=\pm t$ と同値である。$W$ 同士の同型と正則性は
定理\ref{thm:marked-mixed-curve-isomorphisms}と命題\ref{prop:split-mixed-transverse-curve}による。
通常の $W$ と $A$ は最小周期で区別され、$x=-1$ での一致は上の二根の合流による。
\end{proof}

\begin{proof}[対角像だけが残る端点]
$t=\pm1$ では $[X_i,Y_j]=0$ $(i\ne j)$、$[X_i,Y_i]=-2t H_i\ne0$ である。
各組とも三独立（他側の各元と交換子を取る）で、全 $X_i,Y_i$ は全 $H_j$ に跡直交し対角が零である。
$X_i$ が非巡回ならスカラーと階数一 $w b^t$ の和であり、
$H_i(\ker b^t)\subset kw$ である。$H_i$ の階数が2なので $\ker b^t$ はその核 $ke_i$ を含み、
$w\in\operatorname{im} H_i$ となる。従って $i$ 番目の成分によりスカラーは零、
残る二行二列でも対角が零なので $X_i$ は他の二軸間の行列単位の倍である。$Y_i$ も同様。
もし $X_i$ が巡回なら $Y_j,Y_k$ は $X_i$ の跡零中心化環の全体を張る。
$Y_j$ が巡回ならこの可換環がその中心化環でもあり、$X_k$ は $Y_j$ と交換して
$Y_k$ とも交換することになり矛盾する。よって $Y_j,Y_k$ はどちらも非巡回である。
上で決めた異なる二辺の行列単位が可換なら行か列が共通となり、その包は階数1以下で
巡回な $X_i$ を含めない。結局全元が上記の行列単位の倍である。

対角交換子の値により $X_i,Y_i$ は逆向きの単位で倍率の積が決まり、
非対角交換子の消失により $X_i$ の三矢は有向三角形またはその全反転となる。
両側の六倍率は枠の対角変更で全て吸収できる。
実際二つの向きを $\epsilon,\delta\in\{1,-1\}$ とし $R$ を巡回する隣接二座標の差の行列とすれば、
倍率トーラス準同型の指数行列の行列式は
$\det(I-\epsilon\delta R^2)=(1+3\epsilon\delta)^2\ne0$ である。
両方向を同時に反転する操作は二枠の同じ奇置換と全体倍率で実現する。
よって元の向きの組合せは $W(1)$、片側反転の組合せは上記の中心表示の $t=1$ 型に尽きる。
後者では $h$ を零にすれば一方向の三積だけが零になる単項式商を得る。
長さ $n\ge1$ の正規語は $3\cdot2^{n-1}$ 個なので元の窓も標準成長を持たない。
$W(1)$ は既知の正則点である。その第一辺は有限なので他の $A(s)$ とも同型にならない。
実際点を $[r:x_1:x_2:x_3],[s:y_1:y_2:y_3]$ とすると
$r y_i+2x_j y_k=s x_i+2x_k y_j=0$ が巡回的に成り立つ。
$r=0$ では座標軸三点だけ、$r\ne0$ では星点または全 $x_i\ne0$ で
$r^3=-8x_1x_2x_3$、$x_1^4=x_2^4=x_3^4$ となり有限である。
\path{computations/verify_five_dimensional_mixed_brackets.py} は二つの選択の恒等式、
十次元の跡条件空間の小行列式、中心表示の係数と端点の倍率指数を検算する。
\end{proof}

% !TeX root = ../papers/nc-p3.tex
\section{四元群次数を持つ混合交換子のトーラスと全閉点の正則性}
\label{sec:marked-klein-mixed-torus}
\begin{theorem}\label{thm:marked-klein-mixed-torus}
二つの三次元空間の座標軸に三つの非自明な四元群文字を対応させ、
四組 $Q_i,P_i,G_i,F_i$ の各第 $i$ 行列が $E_{jk},E_{kj}$
（$i,j,k$ は巡回）の二成分だけを持つとする。
これら二十四成分が全て非零で $[Q_i,P_i]\ne0$ である印付き窓を全て考える。
その二切断方程式の正規形は下記の三変数主枝と一次元の付加枝で尽きる。
主枝の全次数Koszul左右AS正則性は一つの数 $u^6\ne1$ と同値である。
途中でWalsh正規化の対角分母が消える全ての点もこの判定に含まれる。
付加枝は前節の中心一次拡大で常に正則であり、二枝が一致する値も明示できる。
正則点は全て周期が2を割る。印を忘れた全同型と最小周期は、
孤立した単純印の有限比較、軸置換とトーラスの単項式方程式で必要十分に判定できる。
\end{theorem}
\begin{proof}[四組の二成分を全て解く]
$E_{jk}$ の係数を正、$E_{kj}$ の係数を負と呼ぶ。
$\delta_i=Q_i^+P_i^--Q_i^-P_i^+$ とすれば跡条件 $\sum[Q_i,P_i]=0$ より
三つとも同じ $\delta\ne0$ で、対応する $G,F$ の差もこの $\delta$ である。
対角の残余基底と全体倍率を用いて
\[
 Q_i^+=u,\quad Q_i^-=1,\quad P_i^+=v_i r_i,\quad P_i^-=r_i=(u-v_i)^{-1}
\]
とする。実際 $Q_i^-$ は添字の三倍率で揃えられ、$Q_i^+/Q_i^-$ は
三座標の隣接比の平方で変えるため、積を固定して共通値 $u$ にできる。倍率で $\delta=1$ に戻す。
いずれの $u,v_i$ も非零で、$v_i\ne u$ である。

非対角式は二つの階数一行列 $G_i^+ F_j^+$ と $G_i^- F_j^-$ の対角外六成分を指定する。
$G_i^+=c_i,G_i^-=e_i$ とすれば、巡回の順で
\[
 F_i^+=-r_{i-1}/c_{i-1},\quad F_i^-=-r_i/e_{i-1},\qquad
 c_{i+1}/c_{i-1}=u v_i r_i/r_{i-1},\quad
 e_{i+1}/e_{i-1}=u v_{i+1}r_{i+1}/r_i .
\]
この系の必要十分条件は $u^3 v_0 v_1 v_2=1$。$c_i$ は共通倍率を除き一意、
$e_i=h c_{i+1}$ で、$c_0$ は残余のスカラーで1にできる。
$G_i^+ F_i^- -G_i^-F_i^+=1$ はちょうど
\[
                        (hu+1)(h v_i-1)=0\qquad (i=0,1,2)
\]
となる。従って主枝では $h=-1/u$、付加選択は全ての $v_i=v$ が等しく $h=1/v$ のときだけである。
このとき $uv$ は三乗根で、$u,v_i$ を同じ三乗根倍にする残余を用い $uv=1$ にできる。
この規格化では $u^2=-1$ でのみ二根が合流する。
$Q_i,P_i$ の交叉交換子はそこでの規格化 $uv=1$ の全体で、非対角が交代行列となり、
前節の $t=(u+1)/(u-1)$ に当たる（$u=1,-1$ は $v\ne u$ に反する）。
主選択は $G,F$ が $W(t)$ 側、付加選択は一方だけを替えた中心一次側である。
これは $G_i^-/G_{i+1}^+$ の値を代入しても分かる。
\end{proof}

\begin{proof}[Walsh正規化と極限のAS判定]
主枝で四ラベルを $0,1,2,3\in(\mathbf Z/2)^2$ とする（この段では非零ラベルを $1,2,3$ と読む）。
隣接ラベルが等しい成分は零である。残る三十六成分
$w_{i,i+b,i+b+e,i+e}$, $b,e\ne0$ で、四つの相異なるラベルなら元の非零な一星成分である。
二ラベルが反復する成分は $w_{0i0i}=w_{i0i0}=1$、
$w_{jkjk}=w_{kjkj}=s_i$ であり、
\[
       s_i=u(r_j+r_k)-1=(u^2-v_j v_k)/((u-v_j)(u-v_k)) .
\]
まず全 $s_i\ne0$ とする。四位置の四つずつの対角倍率と九個の非零数 $m_{be}$ に関する
単項式系を作り、各 $w$ を $(-1)^{i_1(b+e)_0}m_{be}$ に変える。
指数行列 $M$ の行は、四位置の該当ラベルに1、該当する $m$ に $-1$ を持つ。
三十六行二十五列のSmith標準形の非零対角は $1$ が十四、$2$ が五、$4$ が二である。
$\operatorname{sign}/w$ を右辺とした十五の左核積は全て1なので、代数閉体上で全点に解がある。
ここで検算は独立不定元 $u,a=v_0,b'=v_1$、$v_2=(u^3ab')^{-1}$ における有理恒等式である。
$N$ をWalshの符号付き平方行列とすると、$N$ の符号は行・列別の積なので次の比には残らない。
\[
 \frac{N_{11}N_{22}}{N_{12}N_{21}}-1
       =\frac{(u^6-1)(ab'u^4-1)}{u^4(u-a)(u-b')}.
\]
具体的に $i,b,e$ の辞書式順（$i$ は0、$b,e$ は1始まり）の指数行のうち、
零始まりで $9,10,12,13,18,19,21,22$ を順に $-1,1,1,-1,-1,1,1,-1$ 倍して
足すと、求める平方比の指数行になる。
\path{computations/verify_klein_mixed_torus.py} が全八十一式、整数左核の全積とこの式を検算する。
$ab'u^4-1=u/v_2-1\ne0$。三色ずつあり、行和・列和が零なので全二次余因子は符号まで等しい。
従って定理\ref{thm:walsh-normal-quadrics}により判定は $u^6\ne1$ となる。

$s_i=0$ が一つでもあれば $v_i=u^{-5}\ne u$、従って既に $u^6\ne1$ である。
全位置でラベル $0,i$ を同時に $z$ 倍し、テンソルを $z^2$ で割る。
この操作を欠けた全 $i$ について合成すると四つの異なるラベルの成分は不変であり、
対角色 $i$ の $0,i$ の対はさらに $z^2$ 倍、残る対は既に零である。他の対角色は保たれる。
極限では $s_i=0$ の各色の対角 $m_{ii}$ を丸ごと零とした配置になる。
残った単項式系の整数左核は、消した行に零を補えば元の整数左核に入る。
残る右辺はそこで非零であるから左核の恒等式も特殊化でき、極限をWalsh表示にできる。
各切断は階数6以下で、非零な三つの色ブロックは各々少なくとも二つの非零係数を持つ。
Walshブロックの左右二次行列表現により非零ブロックの階数は少なくとも2
（可逆な外部符号変更の後で二次行列による正則表現である）。
従って三ブロックとも階数2である。$N$ の非対角は全て非零で一つ以上の対角が零なので
階数1にはなれない。極限で全次数正則となり、
命題\ref{prop:unrestricted-good-reduction}を $k[[z]]$ 上で適用すると元の窓も全次数正則である。
\end{proof}

\begin{proof}[全同型と周期の有限判定]
$\delta\ne0$ なので、$\langle Q,P\rangle$ と $\langle G,F\rangle$ は各々
非対角行列の六空間そのものである。その跡直交補は各側の跡零対角Cartan部分代数で、
三つの共通固有直線が軸枠を復元する。印を保つ同型は従って両側の単項行列に限る。
添字と残る二軸の対応から軸置換は同じ置換であり、残る倍率は四組の係数の単項式比較で決まる。
例えば $v_i$ が全て等しい場合以外では主枝しかなく、巡回置換と同じ三乗根による
$(u,v_i)$ の一斉乗法、または奇置換と四数全ての逆数を取って同じ三乗根を掛ける操作が全てである。
実際 $Q_i^+$ と $Q_i^-$ の比に作用する三つの倍率は積1で、
共通の $u$ を共通値に送れば同じ三乗根となる。逆に平方根の対角比を選び
$Q_i^-$ と $\delta$ の規格化を戻せばこの変換は全て実現する。
一致した $v_i$ の二枝の比較は前節の明示解、または四組の単項式比較で行える。

一位置移動の $Q',P'$ は $F,G$ である。これを再規格化すると主枝では
$u'=-1/u,v'_i=-1/v_i$ を取れる。実際 $F_i^+/F_i^-=h c_i r_{i-1}/(c_{i-1}r_i)$
の積は $h^3$、$u'=h$ に揃えた後の $v_i'$ は $1/(hu v_i)$ である。
非定数な三つ組では二移動で元に戻り、定数の場合も前節の周期を用いればよい。
一般に印を忘れるには、基準の単純な星点を固定し相手の同位置の全ての孤立単純点対を調べる。
これらは有限個であり、点対応全体に正次元成分があってもよい。
星点対は局所一次式が六独立なので常に孤立単純である。対応する他の二点は
反対評価の階数3の核から求め、適合しないものを除く。
跡零規格化の後、上記の六空間・Cartanとその固有軸を検査し、それぞれの枠で
全四組の係数を比較する。これは添字保持同型の必要十分条件である。
一移動との同じ有限比較で最小周期1か2かを判定できる。
\end{proof}

% !TeX root = ../papers/nc-p3.tex
\section{非零対角交換子の全ての疎な四元群台}
\label{sec:marked-klein-mixed-supports}
\begin{theorem}\label{thm:marked-klein-mixed-supports}
前節と同じ四元群次数・共通対角差 $\delta\ne0$ を仮定し、二十四係数の消失を全て許す。
全係数非零の枝以外で標準14641型quadratic Koszul左右AS正則になる印付き窓は、
軸置換と位置移動を用いた下表の番号1,8,17で尽きる。各々、許されたパラメータの全点で正則である。
それぞれ点、一次元、二次元の表示で、表に残る二つの自由パラメータの各々は $0,-1$ を避ける。
各正則窓の最小周期は2を割り、1,8では1である。
全係数非零枝との重なりを含む無印の添字保持同型は、前節の有限個の単純印の比較と下記の指標で必要十分に判定される。
\end{theorem}
\begin{proof}[台とトーラス商]
$\delta=1$ にする。三つの軸ごとに $(Q_i^+,Q_i^-,P_i^+,P_i^-)$ の
非零成分に順に $1,2,4,8$ を付け、その和を一個の数とする。$G,F$ 側も同様である。
$\beta_i=Q_i^-P_i^+$, $\gamma_i=G_i^-F_i^+$ と置くと反対の対の積はそれぞれ $1+\beta_i,1+\gamma_i$ である。
$6,14$ では $\beta_i=-1$、$9,11,13$ では $\beta_i=0$、$15$ では自由（$0,-1$ を除く）である。$\gamma$ も同様。
非対角式は $j=i+1$ として
\[
 G_i^-F_j^-=-Q_i^-P_j^-,\quad G_i^+F_j^+=-Q_j^-P_i^-,
 \quad G_j^-F_i^-=-Q_i^+P_j^+,\quad G_j^+F_i^+=-Q_j^+P_i^+ .
\]
まず左右の零・非零だけを一致させる。各側の局所台は七通りなので
両側の $7^3$ 通り同士を上の十二式で突き合わせると293対を得る。
軸の共通置換（奇置換は各 $+,-$ を交換）と、一位置移動
$(Q,P;G,F)\mapsto(F,G;Q^t,P^t)$ で規格化すると十九行となる。
ここで位置移動は場合分けにのみ用い、無印同型で始点を変更しない。

\begin{center}\small
\begin{tabular}{r|l|l|c|l}
番号&$Q,P$&$G,F$&残るトーラス次元&判定\\\hline
1&6,6,6&6,6,6&0&正則\\
2&6,6,6&9,9,9&0&$H_{0,5}=60$\\
3&6,6,9&6,6,9&0&$H_{0,3}=21$\\
4&6,6,9&6,6,13&0&$H_{0,3}=21$\\
5&6,6,9&6,6,15&--&空\\
6&6,6,13&6,6,13&0&$H_{0,3}=21$\\
7&6,6,13&6,6,15&--&空\\
8&6,6,15&6,6,15&0&正則\\
9&6,14,9&6,14,9&0&$H_{0,3}=21$\\
10&6,14,9&6,14,13&0&$H_{0,3}=21$\\
11&6,14,13&6,14,9&0&$H_{0,3}=21$\\
12&6,14,13&6,14,13&1&$H_{0,3}=21$（一般値）\\
13&6,9,15&6,9,15&0&$H_{0,4}=38$（一般値）\\
14&6,13,13&6,13,11&0&$H_{0,3}=21$\\
15&6,13,11&14,13,13&1&$H_{3,3}=21$（一般値）\\
16&6,13,15&6,13,15&0&$H_{0,4}=38$（一般値）\\
17&6,15,15&6,15,15&0&正則\\
18&14,13,15&14,13,15&1&$H_{0,4}=38$（一般値）\\
19&15,15,15&15,15,15&\multicolumn{2}{c}{定理\ref{thm:marked-klein-mixed-torus}}
\end{tabular}
\end{center}

各行の台の全係数をトーラスの座標とする。上の非零な非対角式は比が $-1$、
対角式は二変数の積が $\beta_i,1+\beta_i,\gamma_i,1+\gamma_i$ の指定値になる指標式である。
指数行列を $M$、右辺を $v$ とすれば整数Smith分解 $LMR=D$ の零行で $v^{L_j}=1$ が必要十分である。
行1--18で非零対角は全て1である。自由な積が $0,-1$ を避ける開集合で左核式のイデアルを計算すると、5,7は単位、2は零、
他はちょうど $\gamma_i-\beta_i$ で生成される。よって自由積の間に他の制限はない。
具体的な係数は $R$ の各行の指標を、最初の $\operatorname{rank}M$ 座標を $v^{L_j}$、
残りを自由とした点で評価すれば得られる。対角基底変更の指数行列を $B$ とすると
$Q_i^\pm,P_i^\pm$ の行は第一枠で $\pm(e_j-e_k)$、第二枠でそれぞれ
$e_i,-e_i$、$G,F$ では枠を逆にする。$MB=0$ であり、$R^{-1}B$ の下部をさらにSmith分解して像を除けば
表の次元となる。代数閉体なのでこの二段目の非零対角の冪写像は全射である。
これが余分な連続同一視のない軸付き正規形とその係数復元法である。
\end{proof}

\begin{proof}[全点の判定と周期]
$H_{p,d}$ は始点 $p$ の長さ $d$ の次元である。各位置でテンソルの後二位置を
$(0,c),(c,0)$, $c=1,2,3$ で縮約した六関係が基底となる（対向星の単位成分による）。
四元群の全次数別に関係を全隣接位置に挿入した行列の余核を、有理数または独立なパラメータの有理関数体で求める。
表の一般値は全特殊化で下限となるので、標準値 $20,35,56$ を超える行には正則点が一つもない。

残る1,8,17では $w_{i,i+b,i+b+e,i+e}$ の一般非零台は $i$ に依存しない。
$(b,e)$ の台は順に
\[
 \{11,12,22,23,31,33\},\qquad
 \{11,12,21,22,23,31,33\},\qquad
 \{11,12,13,21,22,23,31,33\}
\]
である（$be$ は順序対）。前節の対角Walsh正規化をこの台に制限すると
右辺の全整数左核積が1となる。1,8ではこの式が全点で定義される。
17では $1+\beta_2+\beta_3$（軸番号は1始まり）が零のときに限り、$w_{2323},w_{3232}$ が消える。
このとき四位置の $0,1$ を $z$ 倍して $z^2$ で割り、残りの $11$ 対を極限で消す。
その台の左核関係は零行補完によって特殊化される。全ての場合で三行三列が非零、各色の係数は二個以上残り、
二切断の各色階数2は前節と同様に従う。$(3,2)$ が常に欠けるので三行三列の符号付き平方行列は階数1ではない。
定理\ref{thm:walsh-normal-quadrics}と命題\ref{prop:unrestricted-good-reduction}が全次数の正則性を与える。
\path{computations/verify_klein_mixed_supports.py} は台を再列挙し、二段のSmith分解、局所化したイデアル、
有理関数体の余核次元、正規化の全左核積を検算する。

二位置移動で各側の四成分は $(p^-,p^+,q^-,q^+)$ の順になり、$6,15$ の台と積 $\beta,\gamma$ は保存される。
正則な三行では表の自由トーラスが零なので元の印付き類に戻る。
一位置移動と共通奇置換を合わせても各側の $6,15$ は元の種類に戻る。
対角差が各操作で符号を変え、合成では元に戻り、各局所積は元の積を置換したものになる。
1,8では等しい二つの局所積の互換を選べるため最小周期は1である。
一般の無印比較は、$\delta\ne0$ により両枠のCartanが内在的に定まる前節の証明に従う。
同位置の印の有限比較で枠を復元し、共通軸置換後の全四組を対角指標で比較する。
位置移動で表に入れた際は元の位置に戻して比較する。一移動とのこの判定は17の全点でも最小周期を決定する。
\end{proof}

% !TeX root = ../papers/nc-p3.tex
\section{対角交換子零の疎な四元群窓と二つの正規平方}
\label{sec:marked-klein-zero-bracket-boundary}
\begin{theorem}\label{thm:marked-klein-zero-bracket-boundary}
定理\ref{thm:marked-klein-mixed-supports}の四元群次数で共通差を $\delta=0$ とし、
少なくとも一成分が零となる場合を考える。標準14641型quadratic Koszul左右AS正則になる台は、
軸の共通置換と四位置の巡回を除き両側とも $(0,15,15)$ に限る。
この台は $x,b\in k^*$ による下記表示で覆われ、その全点が正則である。
特に $x=b$ で二組の平方を結ぶ行列が階数1に落ちても例外は生じない。
\end{theorem}

\begin{proof}[四位置を通る単項式商による一括排除]
各局所台は $2\times2$ 行列の階数が高々1なので空または二つの部分集合の直積であり、十通りある。
前節の十二の非対角式で積の零・非零を一致させると47492対、
共通軸置換と巡回で2038軌道である。
これは台の必要条件による列挙なので、非実現な台を先に除くことは不要である。

全位置で星文字を零にして、残りの関係に現れる二文字も個別に零とした単項式商を取る。
各位置で $jj$ は残る。相異なる $j,k$ と残りの $i$ に対し、$jk$ がその位置で残る必要十分条件は
$(G_j)_{ki}=(Q_k)_{ji}=0$ である。実際、後ろ二位置の $(0,i),(i,0)$ による六縮約が全関係の基底となる。
三頂点有向グラフに、一周期内で少なくとも一度残る全ての $j\to k$（$j\ne k$）を加える。
閉路があれば、各文字で留まる遷移を用いて、閉路を一周して開始位置と開始文字に戻れる。
同じ長さの一定語とこの周回を任意に連接できるため、経路数は指数増大する。
二つの三つ組について十二個の積の真理値を突き合わせ、各軌道の上のグラフの推移閉包を取ると、
閉路を持たないのは両側が $(0,15,15)$ の場合と両側が $(15,15,15)$ の場合だけである。
後者は疎ではない。\path{computations/verify_klein_zero_bracket_boundary.py} が全台を再生成してこの判定を行う。
\end{proof}

\begin{proof}[残る台の二変数表示]
軸を $1,2,3$ とし $Q_1=P_1=G_1=F_1=0$ である。第2,3軸の $Q_i^\pm$ は、
第一枠の二つの独立な隣接対角比の平方、および第二枠の二倍率で全て1にできる。
$P_i$ はそれぞれ $Q_i$ の非零倍である。共通倍率で $P_2=Q_2$, $P_3=xQ_3$ にする。
$G_2,G_3$ の正成分は第一枠の共通倍率および未使用の第二枠第1軸の倍率でともに1にできる。
残る非対角四式は
\[
 (G_2^+,G_2^-)=(1,b),\quad (G_3^+,G_3^-)=(1,x/b),\qquad F_i=-G_i
\]
を与える。実際 $F_i=k_iG_i$ と書けば正成分同士の二式が $k_2=k_3=-1$、
負成分同士の式が $G_2^-G_3^-=x$ である。他の式は恒等的に成り立つ。
\end{proof}

\begin{proof}[正規列と全次数の自由基底]
文字を $0,1,2,3$ で記し、奇数位置でのみ $0,1$ および $2,3$ を交換する。
以後 $ii$ はその位置からの二文字を表す。$i=0,1$ と $j=2,3$ の間の四関係は
$ji=c_{p,ij}ij$、他の二つは $ii=D_{p,i2}\,22+D_{p,i3}\,33$ ($i=0,1$) となる。
$D_{p,2*},D_{p,3*}$ は単位行とする。全 $c$ と $D_{p,0j},D_{p,1j}$ は非零Laurent単項式である。
例えば位置0の二つの行は $(-1,-x/b),(-1,-1)$ である。各位置の係数は上の四組から直接得られる。

$z_2=22,z_3=33$ を次数2の追加文字とする。これらは各一文字を非零スカラー倍で通過する。
実際 $i=0,1$ では二回の交差交換で $i z_j=(c_{p,ij}c_{p+1,ij})^{-1}z_j i$ である。
$i=2,3$ は $i z_i=z_i i$ と、$i$ を $00$ の左から右へ動かす式を使う。
左の平方係数は $D_{p+1,0*}$、右は $c_{p,0i}c_{p+1,0i}D_{p,0*}$。
第 $i$ 単位行と $D_{p+1,0*}$ の行列式は非零単項式なので、二行だけで
通過行列が決まり、検算すると対角で可逆である。これが $x=b$ でも有効な理由である。
$z_3 z_2$ も二回の通過で $z_2 z_3$ の非零倍となる。

書換えは四平方を $z$ の一次結合へ、混合対を $ij$ ($i<2\le j$) へ、$z$ を先頭へ動かして
$z_2,z_3$ の順とする。全て元の関係の帰結であり、逆に $z$ を消去すれば元の表示になる。
平方では総文字数が減り、他の書換えでは $z_2,z_3$, 第一組、第二組という順に対する転倒数が減る。
二文字書換えの重なりは各位置32個で、記載のコードで左右の正規化を Laurent 多項式として照合すると
全て一致する。従って基底は $z_2^a z_3^b$、$0,1$ の交代語、$2,3$ の交代語を順に並べた全語である。
各組の交代語の母関数は $(1+t)/(1-t)$、二つの $z$ は次数2なので
各始点のHilbert級数は $(1-t)^{-4}$ である。両側で $z_2$ およびその商での $z_3$ は正規非零因子である。

最後に定理\ref{thm:klein-rational-period-four}の線型複体の議論を適用する。
折り畳んだ環で全位置の同じ $z_j$ を足せば、通過式と自由基底から正規正則列となる。
従ってどの頂点単純から環への $\operatorname{Hom},\operatorname{Ext}^1$ も両側で消える。
実際最初の正規元による商の二項分解での変換環系列では、一段目のExtは商環へのHomに
正規化自己同型と次数移動を施したものになり、二番目の非零因子がこれを消す。
四テンソルの縮約による $1,4,6,4,1$ 複体は両端で互いの双対であり、
隣接切断の階数6、星成分による一文字縮約の単射性と二次表示から始端の二写像は表示を与える。
上の二消滅で逆端も完全になり、Hilbert級数のEuler等式から中央も完全となる。
左右とも線型な最小分解とその双対の最終単純を得るので、標準Koszul左右AS正則である。
\end{proof}

% !TeX root = ../papers/nc-p3.tex
\section{対角交換子零の全係数非零窓と二つの反対鉛筆}
\label{sec:marked-klein-zero-bracket-torus}
\begin{theorem}\label{thm:marked-klein-zero-bracket-torus}
対角差零の四元群交換子枝で二十四係数が非零とする。全ての窓は二符号 $u,h\in\{1,-1\}$ と
$(r_1,r_2,r_3)=(1,x,y)$, $xy\ne0$ による下記の四表示で覆われる。
$h=u$ の二表示は全点で標準Koszul左右AS正則、最小周期1である。
$h=-u$ の二表示では正則性は $xy+x+y\ne0$ と同値で、正則点の周期は2を割る。
従って前二節と合わせ、これらの四元群次数を持つ印付き交換子表示の全台・全パラメータでAS性が決まる。
\end{theorem}
\begin{proof}[二符号への縮約]
$Q_i^- =1,Q_i^+=u$ と取り、対角差零より $P_i^-=r_i,P_i^+=ur_i$ と書く。
共通倍率で $r_1=1$ とする。全係数非零なので定理\ref{thm:marked-klein-mixed-torus}の
非対角の計算を $v_i=u$ で使える（この段では対角差で割らない）。
二つの対角外三項積の一致は $u^6=1$ を与え、軸の対角比で同じ三乗根を掛けて $u=\pm1$ にできる。
$G_i^+=c_i,G_i^-=hc_j$（以下 $i,j,k$ は巡回）で
$c_j/c_k=r_i/r_k$ となるので共通倍率を選び $c_i=1/r_j$ にする。
$F_i^+=-r_kr_i,F_i^-=-r_ir_j/h$ となり、$G_i^+F_i^-=G_i^-F_i^+$ は $h^2=1$ である。
逆にこれらの係数は全交換子式を満たす。
\end{proof}

\begin{proof}[$h=-u$ の符号付き平方判定]
反復対の純成分は $w_{jkjk}=w_{kjkj}=u(r_j+r_k)$ で、対向星の対は1である。
まず $\prod(r_j+r_k)\ne0$ とする。定理\ref{thm:marked-klein-mixed-torus}と同じ
三十六行のWalsh規格化の整数左核積が全て1となり、同じ八行による余因子の比は
\[
       \frac{N_{11}N_{22}}{N_{12}N_{21}}-1
            =\frac{(x+y)(y+1)}{y^2}-1
            =\frac{xy+x+y}{y^2}.
\]
三行三列は非零で行和・列和が零なので、これが零なら階数1、非零なら階数2である。
各 $r_j+r_k=0$ のときは $r_1r_2+r_2r_3+r_3r_1=-r_j^2\ne0$。
前節までと同じく欠けた色ごとに四位置の $0,i$ を $z$ 倍して $z^2$ で割り、
極限で色対 $(i,i)$ を全て消す。残った整数左核積を特殊化でき、非対角の全成分が非零なので
三色の階数2ブロックと階数2の符号付き平方行列が残る。Walshの定理と良還元によって元の窓も正則である。
二位置移動の周期性は分母が消える点も元の係数で確認できる。
実際移動後の $Q_i^+,Q_i^-$ は $r_i,ur_i$、二対の局所積と $Q_i^+/Q_i^-$ の積は変わらない。
$G$ 側も比の積 $h$ を保つ（$F_i$ は $G_i$ の倍数、転置で比は逆数）。
上の導出で $u,h$ はこの二つの積、$x,y$ は $Q_i^-P_i^+$ の射影比で決まるので元の標準表示に戻る。
\end{proof}

\begin{proof}[$u=h=1$ の両反対核と全評価]
位置2を $\operatorname{diag}(1,-1/(xy),-1/(xy),-1/(xy))$、位置3を
$\operatorname{diag}(1,1,1/x,1/y)$ で変更すると両対称テンソルになる。
反対切断の対称二乗の三つの非自明色の二次行列式はそれぞれ $1/(xy)^2,1/(xy^2),1/(x^2y)$ である。
零色の四次行列は階数2で、その右核と左核の二基底は下の行列の列になる。
\[
 \begin{pmatrix}
 -1+1/x&-1+1/y\\ -1/x&-1/y\\1&0\\0&1
 \end{pmatrix},\qquad
 \begin{pmatrix}
 (1-x)/(xy)&(1-y)/(xy)\\-1&-1\\1&0\\0&1
 \end{pmatrix}.
\]
従って反対階数は全点で8であり、両側とも対角二次式の鉛筆である。
四行のどの二行も独立になる条件は $1,x,y$ が相異なることである。
このとき四つの特異対角二次式は相異なる単根に対応し基底曲線は滑らかな完全交叉である。
実際特異点なら勾配の従属する鉛筆元の核に入るが、階数3の対角元の核は座標軸であり
他の鉛筆元はそこで非零である。重なる場合は二行を同時に消す元が取れて階数2以下の二次式を含む。

第一辺の純関係が存在しないことを示す。非零な純関係の星・星成分は零なので少なくとも一端の星成分が零になる。
星行が零だと星行を含む三関係の係数が零になり、残る三関係の一次結合は列を $a,Q_i a$ とした行列になる。
階数1なら $ka$ は三行列の共通固有直線である。逆端では $G_i$ またはその転置が同じ役割を持つ。
いずれの組でも各第 $i$ 行列の平方は $i$ 軸以外への座標射影の非零倍で、
全軸の間は非零の行列成分で連結するので共通不変直線はない。
滑らかな場合は定理\ref{thm:smooth-apolar-complete}で正則性が従う。

重なる場合のために全評価の階数下界を示す。点 $[t:a]$ での評価の核を $(s,v)$ とすると
\[
       (tI+G_a^t)v=0,\qquad s a+Q_v^t a=0 .
\]
$a=0$ なら核は一次元、$a\ne0$ なら $v$ から $s$ が高々一通りに決まる。
最初の行列が階数2以上ならこれでよい。$t=0$ で $a\ne0$ のとき、
各非零な成分 $a_i$ に対応する二行二列主小行列式が非零なので階数1以下にならない。
$t\ne0$ で階数1なら、対角三成分が全て非零なので全非対角成分も非零、従って $a_1a_2a_3\ne0$。
ところが $v\mapsto Q_v^t a$ の行列式は $2a_1a_2a_3$ なので、これが $ka$ に入る $v$ は高々一次元である。
他端でも同様で、対応する行列式は $2a_1a_2a_3/(xy)$ である。全評価が階数3以上と分かった。
補題\ref{lem:bisymmetric-reducible-quick-test}で重複の場合も全次数正則である。
\end{proof}

\begin{proof}[$u=h=-1$ と両表示の周期]
四元群の二文字積に双指標 $\epsilon(a,b)=(-1)^{a_1 b_0}$ を掛けるコサイクル捻りを行う。
各多重次数のバー複体で語 $(a_1,\ldots,a_m)$ を $\prod_{i<j}\epsilon(a_i,a_j)$ 倍することで
積の微分と捻った微分が一致する。Hilbert級数とKoszul性を保ち、二次双対の
Frobenius対合も各相補次数で非零倍率を掛けるだけなので左右AS性も保つ。
四テンソルでも同じ積を掛け、位置2,3のラベル $a$ を $\epsilon(a,a)$ 倍して星の単位成分を戻す。
得られる $Q_i^+/Q_i^-$ の積と $G_i^+/G_i^-$ の積はともに1、$Q_i^-P_i^+$ は順に $1,x,y$ である。
従って上の規格化により $u=h=1$ の解そのものになり、逆の捻りで正則性を得る。

最後に両符号の周期1をテンソルの軸の混同なしに確認する。位置 $p$ の三つの非星文字に
倍率を選び、全四位置の六関係を同じ六式
\[
       t x_i+A_i x_jx_k+B_i x_kx_j=0,\qquad
       x_i t+u B_i x_jx_k+u A_i x_kx_j=0
\]
に揃えられる。これは十二倍率と六係数を未知とする四十八の指標式であり、
指数行は各非星積の二倍率、星積の非星文字の逆倍率、該当する $A_i,B_i$ の逆数からなる。
全ての右辺は非零Laurent単項式であり、Smith零行の全指標が1であるため全点で解がある。
\path{computations/verify_klein_zero_bracket_torus.py} はこの二符号の指標式、Walshの整数左核、
反対鉛筆、評価の二行列式、コサイクルの符号を検算する。
\end{proof}

なお前節の疎な正則表示も二位置移動で戻る。実際 $Q'_i=P_i^t,P'_i=Q_i^t,G'_i=F_i^t,F'_i=G_i^t$
を前節の表示で取り、第二枠に対角倍率 $(-b/x,1,1/x)$ を掛け、第一枠は単位とする。
$Q_2,Q_3$ が元の1に、$G_2,G_3$ の正成分も1に、負成分も $b,x/b$ に戻る。
従って四元群交換子表示が一つでも存在する標準窓は、対角差と台によらず周期が2を割る。
印を忘れたときの最小周期の残る二択は、同位置の単純印を全て比較することで決まる。
実際星点の局所一次式は六独立であり単純印は有限個、残る枠の同一視は定理\ref{thm:marked-point-commutators}の
二つの可逆三行列と一倍率の式が必要十分である。対角差零では枠を単項行列に限る必要はない。

% !TeX root = ../papers/nc-p3.tex
\section{四つの異なる文字に支えられた窓の全分類と二色欠損への適用}
\label{sec:squarefree-window-coverage}
\begin{theorem}\label{thm:squarefree-window-coverage}
四位置に四ラベルの基底があり、最上テンソルの台が相異なる四文字の置換に含まれるとする。
隣接階数がともに6なら、標準Koszul左右AS正則であるための必要十分条件は
二十四の置換係数が全て非零であることである。零係数がある全ての場合で指数増大する。
正則窓は添字保持同型まで量子多項式窓の全体に一致する。最小周期は1であり、
相互の同型は系\ref{cor:quantum-torus-normal-form}の三変数と $S_4$ で決まる。

特にトーラスが四位置に作用してテンソル直線を保ち、各位置の四つの重みが一斉平行移動を除き
同じアフィン独立な四点になる全ての標準窓は、この量子枝に属する。
孤立単純印や反対切断階数には仮定を課さない。
\end{theorem}
\begin{proof}
各位置の隣接切断は第一対の無順序二文字を固定した六つの二行二列ブロックに分かれる。
一つでも零なら $ij,ji$ は共にその位置の関係の台に現れない。
全位置で $ii,jj$ も関係に現れないため、対角で待つ語を挿入して
$i\to j\to i$ の閉路と一定の閉路を同じ長さで反復でき、指数増大する単項式商を得る。
従って非指数の場合は六ブロックとも非零、切断階数6から各階数1である。

各二行二列の非零階数1行列の台は九通りの非空長方形のいずれかである。
最初の切断の六台を全て選び、次の切断の六台も非空長方形であるものを残すと、
$9^6=531441$ 通りから22785通りとなる。これらについて、ある位置で
$ij$ を行とする二成分がともに零のときに辺 $i\to j$ を引く。
閉路があれば対角語による待機を使った同じ単項式商が指数増大する。
四頂点の推移閉包で対角が一つもないものは二十四成分全て非零の一台だけである。
\path{computations/verify_squarefree_window_coverage.py} は四文字の置換から
十二ブロックの位置を生成し、全ての真理値をこの手順で検査する。

残る台では全位置 $p$ で各無順序対に一つだけ
$x_{p,j}x_{p+1,i}=q_{p,ij}x_{p,i}x_{p+1,j}$ ($i<j$) があり $q_{p,ij}\ne0$。
最上テンソルの係数で隣接二文字を交換する比は同じ関係直線で決まり、残り二文字の順序に依存しない。
三文字の反転を二通りに行い、全係数非零を使って比較すると
\[
   q_{p,ij} q_{p+1,ik} q_{p,jk}
       =q_{p+1,ij}q_{p,ik}q_{p+1,jk}\qquad(i<j<k)
\]
である（関係直線と書換係数の間の負号は両辺で同数）。
よって各 $p$ の $q_{p,ij}/q_{0,ij}$ は四ラベルの倍率の比に書ける。
位置 $p$ の各文字に掛ける倍率と次位置の倍率との比を順に指定し、両方向に延ばせば
全係数を一定の $q_{0,ij}$ にできる。これは添字を動かさない基底変更である。
逆に全ての量子多項式窓は非零の歪交代テンソルをこの台に持つ。
トーラスの場合は重みの平行移動分を除き、アフィン独立性によりテンソルの全項で
各文字の出現回数が一定である。各位置の一次縮約は単射なので各文字が少なくとも一度現れる。
四文字で長さ4のため全て一度ずつであり、本定理に帰着する。
\end{proof}

\begin{proposition}\label{prop:klein-unmarked-matching-reduction}
四元群次数零の四テンソルで両隣接切断の色 $0$ のブロックが零、他の三色が各階数2とする。
色の原点は四位置のラベルの平行移動で選んでよい。四文字が相異なる全二十四成分が非零とする。
相異なる二ラベル $i,j$ について $d_{ij}=w_{ijij}$ と置くと、$d_{ij}=0$ と $d_{ji}=0$ は同値である。
四頂点の無向グラフに $d_{ij}\ne0$ の対を辺として取る。
\begin{enumerate}
\item 次数3の頂点があれば、四元群の印付き交換子表示があり、前三節の全点判定が適用される。
\item 互いに素な二辺がなければ必ず標準Koszul左右AS正則で、量子窓を軌道閉包に持つ。
\item 残るグラフは完全マッチング、四頂点の道、四辺のサイクルの三種類だけである。
道の中央の辺は全置換係数を変えずに対角基底変更の極限で消すことができる。
\end{enumerate}
\end{proposition}
\begin{proof}
二つの欠色と総次数により、置換以外で許されるのはちょうど $ijij$ である。
$\{a,b\},\{c,d\}$ を一つの完全マッチングとして色 $a+b$ のブロックをこの二組の順に並べると、
対角側は $\operatorname{diag}(d_{ab},d_{ba})$ と $\operatorname{diag}(d_{cd},d_{dc})$、
非対角側は全成分非零の二行二列行列 $U,V$ となる。
二つの対角位置 $a,b$ と、残る組で行と列を異ならせた三次小行列式は
$-d_{ab} U_{b d} V_{c b}-d_{ba} U_{a d} V_{c a}$ である。
その消失と両係数の非消失で最初の同値が従う。
次数3の頂点を星とすれば対向星成分の二つの三次対角行列が可逆である。それらを単位に揃えた
Schur補元が印付き交換子表示であり、四元群次数から各行列は残る二軸の間の二成分に限る。

全四位置で同じラベル重み $v_i$ を用い、テンソルを総重み $\sum_i v_i$ で割れば
置換成分は固定される。対 $ab$ の相対重みは $v_a+v_b-v_c-v_d$ である。
三つのマッチングの相対重みは独立で、補対の相対重みはその負である。
従って互いに素な非零辺がなければ全非零辺に正の整数重みを与えられる
（有理解の分母を払う）。極限には全置換成分が残り、各隣接切断には六つの非零ブロックがあり、
元の階数6以下なので全て階数1である。定理\ref{thm:squarefree-window-coverage}と
命題\ref{prop:unrestricted-good-reduction}より元の窓も全次数で正則となる。
最大次数2でマッチングを含む四頂点グラフの三種類は直接列挙できる。
道では二つの端辺の重みを零、中央の辺の重みを正とすればよい。
\end{proof}

% !TeX root = ../papers/nc-p3.tex
\section{無印の欠色窓の二直線配置と三角PBW極限}
\label{sec:unmarked-klein-pbw}
\begin{theorem}\label{thm:unmarked-klein-pbw}
四位置のラベルを $K=(\mathbf Z/2)^2$ とし、総次数零、両隣接切断の色 $0$ は欠け、
三つの他色は全て階数2とする。四位置の全てで、$i\ne j$ の座標点対 $(i,j)$ は
辺の点対応に属さないとする。四文字が相異なる成分の消失は任意に許す。
\begin{enumerate}
\item $w_{ijij}=0$ が相異なるラベルの少なくとも一対で起これば、窓は常に標準14641型quadratic Koszul左右AS正則である。
ある完全マッチングの両辺の反復対が欠ける場合には、元の座標文字を並べ替えるだけで二次PBW表示となる。
それ以外も同じPBW窓への基底変更の一径数極限を持つ。
\item 全ての $w_{ijij}$ が非零なら、どの座標を星に取っても既出の四元群交換子表示が得られ、
定理\ref{thm:marked-klein-mixed-torus}から定理\ref{thm:marked-klein-zero-bracket-torus}までの全点判定が適用できる。
\end{enumerate}
従ってこの無印領域で全ての正則性が決まり、前節に残した完全マッチング、道、四サイクルの
全係数非零の場合はいずれも全点で正則である。
\end{theorem}
\begin{proof}[二次切断の行と列を同じ四つの直線で読む]
$M_b(i,k)=w_{i,i+b,k,k+b}$ と置く。四位置の点対についての仮定は、
$M_b$ と巡回して得られる他切断の各色行列に零行も零列もないことと同値である。
階数分解 $M_b(i,k)=l_i(v_k)$ を二次元空間で取り、非零な各 $l_i$ の核を
$p_i\in\mathbb P^1$ とする。欠色0から $M_b(i,i+b)=0$ であり、
$v_{i+b}$ は非零なので $kv_{i+b}=p_i$ となる。従って
\[
                 M_b(i,k)=0\quad\Longleftrightarrow\quad p_i=p_{k+b}.
\]
特に左辺の零・非零は $(i,k)\mapsto(k+b,i+b)$ で保たれる。二切断で使うと
四位置の反転 $(03)(12)$ と $(01)(23)$ がともに台を保つ。
四文字が相異なる成分 $w_{i,i+b,i+b+e,i+e}$ では、これら二つの反転が生成する四元群は
同じ $(b,e)$ の四つの $i$ に推移的に作用する。従ってこの四項の非零・零は一斉である。
反復対でも $w_{ijij}$ と $w_{jiji}$ の消失は同値である。

$b\ne e$ で四項が零なら $p_i=p_{i+e}$ が全 $i$ に成立する。
残る二組の直線は相異なる。そうでなければ $M_b$ が零になる。
従って全 $p_i\ne p_{i+b}$ であり、この色の反復対 $w_{i,i+b,i,i+b}$ は四つとも非零である。
もう一方の切断でも同様に色 $e$ の反復対が全て非零になる。
ゆえに色 $c$ の反復対に一つでも零があれば $(b,c),(c,b)$ の位置は、
どの $b\ne0,c$ についても四項が全て非零である。
\end{proof}

\begin{proof}[二組の文字の間の四つの主元]
色 $c$ の四反復対を全て零とする。二組のラベル $A=\{0,c\}$, $B=K\setminus A$ を取り、
$A$ の二文字を先に、$B$ の二文字を後に並べる。色 $c$ の二行二列の対角ブロックは
対向反復対と欠色により両方とも零で、非対角の二ブロックは全成分非零である。
各階数1なので $A$ 内に一つ、$B$ 内に一つ、相異なる二文字の両順を非零係数で含む関係がある。

色 $b\ne0,c$ の行列で $B=\{j,k\}$ の二行と $j+b,k+b$ の二列を取る。
対角は欠色で零、非対角は先の $(b,c)$ の全非零条件で非零である。
従って $BA$ で始まる二座標への関係平面の射影は可逆で、二つの $BA$ 語を $AB$ のみで書き換えられる。
他方の切断は $(c,b)$ を使い、反対側の列も同様である。
全四位置で六つの先頭語は降順対の全体となる。停止する二次簡約による三次元数の上界は20、
一文字縮約四本が独立で三次交叉に入ることによる下界も $64-48+4=20$ である。
縮約の単射性は異なるラベルで異なる総次数を持ち、各ラベルが実際に現れることによる。
従って三文字の重なりは全て解消し、Diamond補題で全次数の基底は非減少語となる。
二次PBWからKoszul、双対の次元列は $1,4,6,4,1$ である。一文字縮約の両側単射性と切断階数6から
最上テンソルの対合が次数1対3、2対2でも完全なので左右AS性が従う。

色 $c$ のうち一辺の両向きだけが零なら、その補辺の両ラベルを四位置とも $z$ 倍して $z^2$ で割る。
四文字が異なる項と他色の反復対は不変、補辺の両反復対だけがさらに $z^2$ 倍になる。
極限でも上の主元は残り、色 $c$ の非対角ブロックも非零で、元の階数2以下という上界は保たれる。
従って極限が上記の正則PBW窓になり、命題\ref{prop:unrestricted-good-reduction}で結論を得る。
全反復対が非零の場合は星の正規化が直接可能である。
\path{computations/verify_unmarked_klein_pbw.py} は四直線の全ての一致分割、二切断の反転と主元の色を検算する。
\end{proof}

% !TeX root = ../papers/nc-p3.tex
\section{両切断が欠色を持つ全四元群窓の被覆と正則性判定}
\label{sec:klein-missing-color-coverage}
\begin{theorem}\label{thm:klein-missing-color-coverage}
四位置で四元群の各文字が一回ずつ現れ、最上テンソルは斉次、隣接切断は階数6とする。
両方の切断に少なくとも一つずつ欠色がある場合、印や台への追加仮定なしに次の手順が
標準14641型quadratic Koszul左右AS正則性を必要十分に決定する。
\begin{enumerate}
\item 各切断の色階数が $(0,2,2,2)$ の置換でなければ排除する。
文字原点を選び、総次数および両欠色を零にする。
\item 三つの色ブロックのどれかに零行または零列があれば、二切断のどちらでも指数増大して排除される。
\item それ以外では反復対 $w_{ijij}$ が少なくとも一つ欠ければ常に正則である。
全反復対が非零なら星点で規格化し、四元群印付き交換子の全台判定を用いる。
\end{enumerate}
\end{theorem}
\begin{proof}[Koszul分解の文字別Euler漸化式]
まず一位置のラベルを平行移動して最上テンソルの総次数を零にできる。
反対の二次関係は同じ文字階数を持つ。非自明な指標 $\chi:K\to\{\pm1\}$ について、始点 $p$ の
長さ $n$ の経路の指標値（各文字次元に $\chi$ を掛けた和）を $h_{p,n}$、
関係六空間の値を $s_p$ とする。標準ASなら四項の線型分解はこの多重次数を保ち、一次矢と
三次重なりの各指標値は零、最上の値は1である。$h_{p,0}=1$, $h_{p,1}=0$ とし、
\[
           h_{p,n}=-s_p h_{p+2,n-2}-h_{p+4,n-4}
\]
が成り立つ（負の長さは零）。$s_{p+2}=s_p$ により偶長の母関数は
$1/(1+s_p z+z^2)$ である。各係数の絶対値は標準Hilbert関数以下だから
$|s_p|\le2$ が必要である。実際 $|s_p|>2$ なら相異なる実根の一つが単位円内にあり分子で消えない。
この議論は片方にしか欠色がないときにも各切断に独立に適用できる。

両切断の色を適当に平行移動するには、総次数を零にした後、四位置を
$0,B,B+E,E$ ずらせばよい。欠色を $0$ に取ると、核 $\{0,b\}$ の非自明指標で
$s_p=2r_{p,b}-6$ である。各 $b\ne0$ で $r_{p,b}\ge2$ が必要で、和は6だから全て2となる。
欠色がない側で許される分割は整数の和だけから $(3,1,1,1),(2,2,1,1)$ の二つである。
\end{proof}

\begin{proof}[零行・零列が一方通行だけになることの排除]
自由辺 $i\to j$ を、ある位置で $ij$ がどの関係の台にもない相異なる座標対とする。
欠色0により全ての対角遷移 $ii$ が全位置で使える。従って自由辺に有向閉路があれば
定理\ref{thm:squarefree-window-coverage}の単項式商が指数増大する。
閉路がなければ全自由辺が小から大に向かう共通全順序を選べる。
四頂点の任意の置換は $K$ のアフィン変換なので、共通の付け替えで通常の順序 $0<1<2<3$ にしてよい。

この順序で許される台を二次元直線だけで列挙する。切断色 $b$ の候補は
$w_{i,i+b,i+b+e,i+e}$, $e\ne0$ の十二個の真理値である。行 $i$ と列 $k$ は
それぞれ位置0,2での対 $i(i+b),k(k+b)$ に当たるため、大から小に向かうものは零行・列になれない。
零行と零列を取り除き、残った行・列に二次元空間の直線を一つずつ対応させる。
各行には線型形式の核、各列にはベクトルの像を取る。零成分で結んだ行・列は同じ直線、
非零成分で結んだものは相異なる直線である。零による同値閉包が非零の対を含まず、
行にも列にも二つ以上の同値類があれば、相異なる一般の直線を各類に割り当てて
ちょうどその台の階数2行列を作れる。これで個別の切断での零行・列を含む必要十分条件になる。

十二ビットの4096通りからこの検査で各色171通りを得る。三色を独立に選び、次の切断で
行ラベル・行色・列方向を $(i,b,e)\mapsto(i+b,e,b)$ に置き換える。
この切断でも同じ171通りに入るものは108通りで、どの台も六ブロックの全てに零行も零列も持たない。
\path{computations/verify_klein_missing_color_coverage.py} はこの $171^3$ 通りを二切断のビット位置から直接突き合わせる。
二切断の行列要素の実現を同時に求める必要はなく、必要な台の検査だけで既に零行・列が排除された。
定理\ref{thm:unmarked-klein-pbw}によって最後の段の全次数判定が従う。
\end{proof}

% !TeX root = ../papers/nc-p3.tex
\section{片切断だけが欠色を持つ場合の十個の候補台}
\label{sec:klein-single-missing-supports}
\begin{proposition}\label{prop:klein-single-missing-supports}
前節の四元群次数を仮定し、二切断のうち一つだけに欠色がある場合を考える。
正則窓では欠色側の分割は $(0,2,2,2)$、他側の $(3,1,1,1)$ 型は全て排除される。
残る $(2,2,1,1)$ 型でも、次に定めるラベル変換・巡回・反転で場合分けすると、
可能な台は検算コードに記載した十個だけである。
\end{proposition}
\begin{proof}[一つずつの切断と四層の自由グラフ]
文字の平行移動と共通の線型置換で欠色を $b=0$、奇側の降順階数を $e=0,1,2,3$ の順に取る。
標準正則なら他の分割がないことは前節の文字別分解による。
奇切断 $N_e(j,l)$ では $b=0$ 由来の四成分が零で、各々の台は十二ビットで指定される。
偶切断 $M_b(i,k)$, $b\ne0$ は十六ビットで指定できる。
各単独ブロックの台を次の必要十分条件で検査する。
階数1は非空長方形、階数2は前節の零成分の同値閉包である。
階数3では台の二部グラフに三マッチングがあり、完全マッチング数が1でないことが必要十分である。
実際完全マッチングがなければ一般階数3である。二つ以上なら、あるマッチングに属して別のものには属さない成分を取ればよい。その変数について行列式は
一次係数も定数項も恒等零ではない。他の非零値を一般に固定すれば非零の根で行列式だけを消せ、
その三次余因子は非零なので階数3を得る。

先に三つの奇色 $1,2,3$ を列挙し、各偶色について残る四ビットのうち
階数2可能な補完だけを取り、最後に奇色0の台を検査する。こうして
$(3,1,1,1)$ で15936576通り、$(2,2,1,1)$ で115151588通りが残る。
四層各四頂点のグラフに、零行・列から得る自由な座標遷移を加え、
偶位置では対角の四辺も加える。強連結成分内に二つの出辺を持つ頂点があれば、
二つの異なる閉路の長さを反復で合わせて指数増大商を得る。
これを通過する台の数は順に0、2664である。
\end{proof}

\begin{proof}[連続純関係と一文字商]
$(2,2,1,1)$ 型の台を次の作用でまとめる。文字 $0,1$ を固定する $2,3$ の交換、
四位置への共通平行移動、位置2,3だけへの1の同時加算、二位置巡回と
$(a,b,c,d)\mapsto(b,a,d,c)$ による反転を用いる。これらによる2664台の軌道数は74である。
反転は左右の正則性を同時に扱う場合分けに使っており、反対環を添字保持同型とみなす必要はない。

色階数 $r=1,2$ の行集合 $I$ への射影階数の上界を $u(I)$ とする。零行を除き、
二次元の直線の同値閉包において $I$ が占める類の個数と $r$ の小さい方で上から抑えられる。
階数1では非零行が全て同じ行直線を張るので同じ上界が使える。
一行を除く残りで $u<r$ なら、その行だけに支えられた純関係が存在する。
二つの連続位置でこの関係が繋がれば定理\ref{thm:no-consecutive-pure-relations}で排除する。

次に位置 $p$ を座標点 $i$ に縮約する。命題\ref{prop:cut-arrow-growth}の先頭・末尾の商は
自由な後続ラベル集合 $X$、先行ラベル集合 $Z$ がそのまま基底を与える。
残る三文字の総次数 $g$ ごとに中央文字は両端から決まるので、関係を入れる前の次元は $|X||Z|$。
位置 $p+1,p+2$ の射影階数上界を $u_{p+1,b},u_{p+2,b}$ として
\[
 d_i\ \ge\ \sum_{g\in K}\max\left(0,|X||Z|-
     \sum_{c\in Z}u_{p+1,c+g}(X)-
     \sum_{a\in X}u_{p+2,a+g}(Z+a+g)\right).
\]
各和はもう一端の文字を固定して投影する関係行の上界であり、交叉は引く必要がない。
この下界が2以上でも一文字商の反復が指数増大する。
この二検査を通るものは74軌道中10個である。

具体的には台を非負整数として表し、非零語 $(i,j,j+e,i+e)$ が
$12e+4((i+j)-1)+j$ ビットに対応するとする（$i+j\ne0$ と $e$ を整数ラベルとしてビット位置に使う）。
十代表はこの順に十六進で
\[
\begin{array}{llll}
10330ffffff&660666ff9ff&66066ffffff&110120ffffff\\
110330ffffff&111111ffffff&330330ff0ff3&330330ff3ff3\\
330330ffffff&330cc0ffffff&&
\end{array}
\]
となる。\path{computations/verify_klein_single_missing_supports.py} は付属のC++17検算を一時ディレクトリでビルドし、
全ての単独階数検査、四層の有向グラフの推移閉包、74軌道の二つの成長検査を再現する。
\end{proof}

% !TeX root = ../papers/nc-p3.tex
\section{片側欠色の全正則窓は一つの有理曲線で覆われる}
\label{sec:klein-single-missing-windows}
\begin{theorem}\label{thm:klein-single-missing-windows}
四元群正則表現の四窓で一方の隣接切断だけに欠色があるとする。
標準14641型quadratic Koszul左右AS正則な窓は、前節のラベル変換、巡回、反転、
四位置の基底変更の下、次の一変数表示 $t\ne0$, $1+t^2\ne0$ で全て覆われる。
この表示の全点は正則で、二次PBW基底を持つ。前節の十台のうち実現する正則台は最後の一つだけである。
巡回と反転は被覆における表示操作とする。
\end{theorem}
$L=1+t^2$ と書き、各欄は $w_{j+b,j,j+e,j+e+b}$ を $j=0,1,2,3$ の順に並べた行とする。
$b=0$ は全て零である。
\[
\begin{array}{c|c|c|c}
 e & b=1 & b=2 & b=3\\\hline
 0&(-2t^2/L,-2t^2,1,1)&(t^2,1,-t^2,-1/L)&(1,-t^2,-t^2/L,1)\\
 1&(-2t^3/L,-2t,t,1/t)&(-1,1,-t^2,t^2/L)&(1,1,-t^4/L,-1)\\
 2&(0,0,0,0)&(0,0,-t^2 L,1)&(0,0,-t^2,L)\\
 3&(0,0,0,0)&(1,1,0,0)&(1,1,0,0)
\end{array}
\]
\begin{proof}[トーラス上の小行列式と有限因子]
前節の順に台を $n=0,\ldots,9$ と番号付けする。各台の非零な係数を独立なトーラス座標として
二切断の指定階数より一つ大きい全小行列式を取る。単項式を除いた二項式を先にSmith分解し、
自由トーラス上で四位置の対角基底変更をさらにSmith分解して除く。
座標の構成は定理\ref{thm:marked-klein-mixed-supports}の二段の指標と同じである。
ただし今回は $2,4$ の非自明なSmith因子を全ての根に分ける。
残りの小行列式を代入して非零変数の積を局所化すると、各台で有限因子の選択が一通りだけ残る。
他の根では局所化したイデアルが単位になる。この検査は $\mathbb Q(\sqrt{-1})$ 上で行える。
残る変数を $a,b,c,d,e,f,g$ の必要な個数だけ左から取ると、以下は残る式による
稠密な有理表示である。右辺は上から順に代入する。$H$ は始点0、長さ4、四元群総次数0の次元である。
\[
\begin{array}{c|l|l|r}
 n&\text{自由変数}&\text{他変数}& H\\\hline
0&a,d&b=ad^2+d,\ c=ad/(a-b)&12\\
1&a& &12\\
2&a,b&c=1+1/b&12\\
3&b,c&a=1/b-1/c&12\\
4&c,d,f&a=f/(d(f-1)),\ e=1-ad^2/(cf),\ b=1+ce/a&12\\
5&a,b&c=1+ab&13\\
6&b,c,e&a=e(bc-1)/(b-1),\ d=(1-e)/(c-e)&12\\
7&b,d,e,f&c=def-d+1,\ a=(c-f)/(e-1)&12\\
8&b,c,d,f&e=bf+cf-bcf^2,\ g=(bd^2-e)/(bd^2-1),\ a=bd(c-g)&12\\
9&d&a=-d^2,\ b=1/(a-1),\ c=-2a^3b/d^3&11
\end{array}
\]

分母によって別成分を落としていないことを示す。0--5と9では各式の分母の消失は
元の小行列式と座標の非零性に反する。例えば0で $a=b$ なら $ad=0$、
4で $f=1$ なら $f=0$、9で $a=1$ なら $b(a-1)=1$ に反する。
6の方程式は $(b-1)a=e(bc-1),(c-e)d=1-e$ である。
$b=1$ なら $c=1$、$c=e$ なら $c=e=1$ であり、どちらの境界も次元2以下である。
五変数の二式の各成分は次元3以上なので新しい成分はない。
7で $e=1$ なら $c=f,(d-1)(f-1)=0$ で次元3以下、六変数の二式は次元4以上である。
8では $a,e$ を消去した後の五変数の一式で $bd^2=1$ なら $e=1$ と
$(1-bf)(1-cf)=0$ が必要になり次元3以下である。一般開集合の次元は4であり、これも別成分を持たない。
各行のグラフの残る非零開集合は既約であるから、表の計算は全実現点を含む閉包上の一般点計算になっている。

四次障害を再現するには、総次数零の64語を列に取り、三位置の二次関係の全挿入を行とし、
表の自由変数の有理関数体で余核次元を求める。各色ブロックの実際の階数が指定値に等しいところでは
切断像はベクトル束であり、余核次元は特殊化で下がらない。
定理\ref{thm:klein-missing-color-coverage}の漸化式では、始点0の各非自明指標の長さ4での値が3なので
標準値は $(35+3\cdot3)/4=11$ である。これで初めの九台は例外値を含め排除される。
9の有限因子は $2,4$ の二つであり、残った根で $d=t$ とすれば上の係数表そのものになる。
\path{computations/verify_klein_single_missing_windows.py} は全小行列式、全ての根、
自由変数体のグラフと四次階数を検算する。
\end{proof}

\begin{proof}[二つの二文字窓の線型旗]
位置0,1では第一組 $A=\{0,1\}$、第二組 $B=\{2,3\}$ とし、位置2,3では二組を逆にする。
各辺に第一組内部と第二組内部を結ぶ関係が一つずつあり、ともに可逆な二行二列の双線型形式である。
残る四関係は $A_p B_{p+1}+B_p A_{p+1}$ に入り、$B_p A_{p+1}$ への射影が可逆である。
実際、係数表の各ブロックで辞書順の先頭独立列を取ると、
この混合四行列の行列式は四位置で順に
\[
           L^2,\qquad 4t^6/L^2,\qquad -t^4 L^2,\qquad -16t^8/L^2 .
\]
二つの内部形式の行列式も各位置で非零定数と $t,L$ の冪だけからなる。
これは上記検算の最後の四つの射影で全値について確認できる。

内部形式により、第一組の一文字の射影的評価から次位置の関係を消す相手の評価への
射影直線の同型が定まる。四位置を一巡する合成は射影直線の自己同型なので固定点を持つ。
その点から四位置を一周して、各位置の第一組の最大文字に対応する双対基底ベクトルを選ぶ。
内部関係で最大文字同士の係数が消え、次の列は非零なので先頭対は二文字の下降対になる。
第二組にも同じ選択を独立に行う。各位置で第一組を第二組より小さく置くと、
混合四関係は全ての $BA$ を $AB$ に変える。これで全ての位置で通常の六下降対が得られた。
二十四項全非零をここでは必要とせず、係数表から一文字縮約が全位置で単射である。
三次元数の上下界20による重なり解消、Koszul双対の完全対合は定理\ref{thm:unmarked-klein-pbw}と同じである。
従って全点で左右AS正則である。
\end{proof}

% !TeX root = ../papers/nc-p3.tex
\section{片側欠色曲線の無印同型と二つの周期1点}
\label{sec:single-missing-curve-isomorphism}
\begin{theorem}\label{thm:single-missing-curve-isomorphism}
定理\ref{thm:klein-single-missing-windows}の全ての窓を、四元群作用、二組の文字、巡回表示の印を忘れ、
添字保持同型で分類すると、同型類は $X\in k^*$ と一対一に対応する。
基準の始点では $X=(t+t^{-1})^2$、一位置移動は $X\mapsto16/X$、反転は
前節の基準表示を元の類に戻す。最小周期は $X=4,-4$ の二点で1、その他の全点で2である。
\end{theorem}
\begin{proof}[点対応から二直線を復元する]
前節の二組に対応する点空間を $A,B$ と書き、各々二次元とする。次位置では $A',B'$ とする。
内部関係が定める相手への線型写像を $f:A\to A',g:B\to B'$ とする。
例えば内部形式 $a^t D a'=0$ に対し $f=JD^t$、$J$ は二行の標準交代行列である。
各射影線の外の点 $[a:b]$ について、その相手は $[sfa:t'gb]$ と書ける。
混合関係は $s C(a\otimes b)+t' D'(a\otimes b)=0$ になる。
$C$ は可逆で、$M=C^{-1}D'$ と置く。
ここの $a,b$ の両方が非零のとき $a\otimes b$ は $M$ の固有ベクトルでなければならず、
逆にそのとき $s,t'$ の非零な比が一意に決まる。一方 $a,b$ の片方が零の点は全て相手を持つ。

係数表から四位置で $M^2=kI$, $k\ne0$, $\operatorname{tr} M=0$ となる。
二つの固有空間の射影直線をSegre曲面 $\mathbb P A\times\mathbb P B$ と交わらせる。
第一成分 $a=(1,z)$ が交点の端点になる条件は、$P=a\otimes B$ の二列を取ると
$\det[P,MP]=0$ である。実際 $P\cap MP$ が非零なら $M^2=kI$ により交叉自体が $M$ 不変で、
固有ベクトルを含む。第二成分でも同様である。各位置でこの二つの二元四次式は
$az^4+bz^2+c$ ($ac\ne0$) の形で、$H=t+t^{-1}$ として
\[
       b^2/(4ac)=(H+1)^2,\quad (H-1)^2
\]
をそれぞれ与える。従って固有直線は曲面に含まれない。
各四次式が重根を持つ条件は $(H\pm1)^2=1$ で、$H\ne0$ のため少なくとも一方の四次式は
四つの相異なる端点を与える。他側に重根がある場合にもその重複度は2である。
よって常に相異なる四つの分解可能固有直線、従って射影的には四つの点対 $([a],[b])$ を得る。
第一射影の被約支持は $\mathbb P A,\mathbb P B$ と各点対を結ぶ四本の直線の和である。
二つの基準直線は他の四本と全て交わり、各結合直線の方は二つの基準直線と、
高々もう一つの結合直線と交わる。従って無印の辺から基準直線の無順序対が全点で復元される。
\end{proof}

\begin{proof}[一辺の数値不変量と二窓の比較]
$A\otimes B$ 上に二行二列行列の行列式形式を取り、その随伴を $M^\dagger$ とする。
基準直線対の同型はこの形式の相似変換になる。二つの組を交換すると、
混合式の二項が入れ替わるため、因子の交換を除けば $M$ は $M^{-1}$ になる。
$M^{-1}=M/k$ であり、内部写像のスカラー選択も $M$ をスカラー倍するだけである。
よって
\[
            I=\operatorname{tr}\big((M^\dagger M/k)^2\big)
\]
は無印の辺の不変量である。全四位置で順に
$4-32/X,4-2X,4-32/X,4-2X$ を得る。基準始点の同型では $X$ の一致が必要である。
逆に $X$ が同じならパラメータは $t,-t,t^{-1},-t^{-1}$ のいずれかである。
$-t$ と $t^{-1}$ への同型は四位置のラベル置換と倍率で実現する。
具体的には置換を $0,1,2,3$ の像の列として、$-t$ に対し位置0,1で $2301$、
位置2,3で $3210$、$t^{-1}$ に対し位置0,1で $0132$、位置2,3で $1023$ を取る。
残る非零成分比は四位置各四個の倍率の指標系で、その全整数左核積は1である。
二位置移動は四位置全てで $2310$、四文字の逆順は前二位置 $2310$、後二位置 $3201$ により
同じ指標検査を満たす。従って反転と二移動はパラメータ類を増やさない。

一移動の十分性には混合四行列そのものを比較する。各次位置の $A,B$ を前の位置の $f,g$ で
引き戻し、最初の三辺の $M$ を $M_0,N_1,N_2$ と同じ四空間で読む。
$N_2=-4t^3M_0/(1+t^2)$ である。$i^2=-1$, $s=i(t-i)/(t+i)$, $L=1+t^2$ として
\[
 P=\begin{pmatrix}1&-it/L\\-4it/L&4t^2/L^2\end{pmatrix},
 \qquad Q=\begin{pmatrix}1&-1/t\\-s&-s/t\end{pmatrix},\qquad C=P\otimes Q .
\]
二つの可逆行列は全許容値で
\[
 C M_0(t) C^{-1}\sim N_1(s),\qquad
 C N_1(t) C^{-1}\sim M_0(s)\sim N_2(s)
\]
を満たす。$\sim$ は非零スカラー倍を表す。この比較で二窓の同型が従う理由を述べる。
内部形式を両側の $f,g$ で同定した後、混合関係四空間は二つの混合四空間の間の
可逆写像のグラフであり、その行を左から $C_{\rm left}^{-1}$ で揃えたものが上で用いた $M$ である
（$C_{\rm left}$ は混合式の第一係数行列）。従って $P,Q$ の基底変更でグラフの写像がスカラー倍まで一致する。
二辺のスカラーは、三つの位置で二組の基底に掛ける相対倍率の隣接比で順に吸収できる。
内部関係は各々の直線として保たれる。これで二つの連続窓が同型となり、
系\ref{cor:two-window-extension}で添字保持同型へ延びる。
$(s+1/s)^2=16/X$ なので全一移動が基準曲線に戻り、固定点は $X^2=16$ の二点である。
最後に $X\ne0$ なら $t+t^{-1}=\sqrt X$ は許容根を持つのでパラメータの像は $k^*$ 全体である。
\path{computations/verify_single_missing_isomorphism.py} は八個の四次式、四位置の不変量、
置換後の整数左核および二つの基底変更恒等式を有理関数体で検算する。
\end{proof}

% !TeX root = ../papers/nc-p3.tex
\section{両切断に階数1の色が三つある全窓}
\label{sec:klein-three-lines}
\begin{theorem}\label{thm:klein-three-lines}
正則表現型の四元群窓で隣接する二切断の色階数がともに $(3,1,1,1)$ の置換であるとする。
各切断の階数3の色を位置別の文字平行移動で零に取る。
このとき標準AS窓は、各位置での対角基底変更の後に、次のテンソルとその四元群二コサイクル捻りで尽きる。
\[
 W_{ijkl}={\bf1}_{i+j+k+l=0}
 -\sum_{a\in K}p_a^{-1}{\bf1}_{i=j=k=l=a},
 \qquad \sum_a p_a=1,\quad \prod_a p_a\ne0 .
\]
表示の全点が左右Koszul AS正則で、最小周期は1である。両表示の間に添字保持同型はない。
対称表示の反対切断の階数は6である。捻った表示では
$\sum_a1/p_a\ne4$ のとき16、$(-1,1,1,1)/2$ の置換四点では14、
それ以外の $\sum_a1/p_a=4$ の点では15である。
\end{theorem}
\begin{proof}[零成分を持つ台の検査]
$M_b(i,k)=w_{i,i+b,k,k+b}$ とし、一位置巡回したものを $N_e$ と書く。
階数1の非零行列の台は非空な行集合と列集合の直積であり、225通りある。
$N_1,N_2,N_3$ の三長方形を選び、$M_1,M_2,M_3$ のそれぞれについてこれと
交わる12セルに一致する長方形を全て取る。選択後に未定なのは $w_{iiii}$ の4セルだけである。
$M_0,N_0$ それぞれに大きさ3のマッチングがあり、完全マッチングがちょうど1個でないことを要求すると
39461472台となる。これは階数3の必要条件だけを使う段階である。

零行・列から四層の自由辺グラフを作ると、強連結成分内に分岐を持たない台は5160個である。
このうち16個は六つの長方形が全て四行四列である。残る5144個を同時文字置換、巡回、反転でまとめると
52軌道となる（四文字の任意の置換がアフィン変換である）。
行集合 $I$ への射影階数を、色階数とその行集合上の最大マッチング数の小さい方で抑える。
前節までの連続純関係検査と一文字商の三次元数下界
（\path{sections/klein_single_missing_supports.tex}）をこの上界で実行すると、52軌道は全て排除される。
台のビット番号は $16(i+j)+4i+k$（$i+j$ は四元群の和）である。
\path{computations/klein_three_lines_supports.cpp} と
\path{computations/verify_klein_three_lines.py} が各段階と全代表の検査を再現する。
\end{proof}

\begin{proof}[六長方形の正規化]
残る60個の非純成分は非零である。各階数1ブロックで第零行・列を基準に
$x_{00}x_{ik}=x_{i0}x_{0k}$ ($1\le i,k\le3$) を課す。54行60列の指数行列の
Smith非零対角は $1^{46},2$、四位置の対角基底変更の階数は13である。
従って連結核は基底変更で尽くされ、二つの成分が残る。
一つの代表は全係数1であり、もう一つは
$\prod_{r<s}(-1)^{(i_r)_{\rm hi}(i_s)_{\rm lo}}$ である。
検算では左右の整可逆変換も求め、後者のSmith成分が実際に $-1$ であることを確認する。
捻りは総次数零の四語で巡回不変であり、左右のKoszul複体と双対完全対合を保存する
（\path{sections/klein_second_pairings.tex}）。従って以下は捻りを戻して調べればよい。

対称代表の階数3ブロックは両切断とも $J+\operatorname{diag}(d_a)$、$J$ は全成分1である。
$d_a$ が全て非零なら階数3の条件は $1+\sum 1/d_a=0$ で、核は $p_a=-1/d_a$ で張られる。
零の $d_a$ があるとき核はその場所に支えられ成分和が零のベクトル全体で、階数3にはちょうど二つの零が必要である。
その場合には他の位置の純平方が連続した関係になり、標準型にならない。

全 $p_a\ne0$ の場合、縮約像 $R\subset\operatorname{Sym}^2 V$ は
平方成分で $p$ に直交する超平面と、三つの色ごとの混合対称成分の和の直線からなる。
従って $v^2\in R$ には三通りの $v_0v_i=v_jv_k$ と $\sum p_i v_i^2=0$ が必要十分である。
一成分が零なら前者は座標点のみを許し、全成分が非零なら全 $v_i^2$ は等しい。
$p_i\ne0$ と $\sum p_i=1$ により、どちらにも射影解はない。
従って $R$ の対称二次直交補空間は基点自由な四式を与える。
定理\ref{thm:symmetric-frobenius-chart}の二次完全交叉であり、そのsocleテンソルは
四式で消される非零四次形式 $W$ と比例する。四位置の巡回不変性は捻った後にも成り立つ。
\end{proof}

\begin{proof}[反対階数の特殊化]
対称形では全切断が同階数である。捻った形の反対切断は
$(i,i+b)\mid(j,j+b)$ の各非零色で行列式の絶対値16の符号行列を与える。
色零では $H-\operatorname{diag}(1/p_i)$ となる。$H$ は対称で
$H^2=4I$、対角成分は1、二次主小行列式は0、三次主小行列式は全て $-4$、行列式は $-16$ である。
$P=\operatorname{diag}(p_i)$ と置くと
\[
 \det(I-PH)=1-\sum p_i+4\sum_{i<j<k}p_i p_j p_k-16\prod p_i
 =4\prod p_i\left(\sum_i1/p_i-4\right).
\]
階数2以下なら第 $l$ 行列を除く三次主小行列式も零であり
$p_l+4\prod p_i/p_l=0$ となる。従って全 $p_l$ は一つの $a$ の符号倍である。
負号の個数を $m$ とすれば $(4-2m)a=1$ かつ $1=-4(-1)^m a^2$ なので、
可能なのは上記四点だけである。各点の色零ブロックに対して直接階数2を得る。
従って反対階数は記載の三値で尽くされ、二表示の階数集合は交わらない。
\end{proof}

% !TeX root = ../papers/nc-p3.tex
\section{三重色と二重色が隣接する場合の全排除}
\label{sec:klein-mixed-positive}
\begin{theorem}\label{thm:klein-mixed-positive}
正則表現型四元群窓において、二つの隣接切断の色階数がそれぞれ
$(3,1,1,1)$ と $(2,2,1,1)$ の置換である標準14641型Koszul左右AS正則代数は存在しない。
\end{theorem}
\begin{proof}[台と五つの長方形]
文字を平行移動し共通自己同型を施し、$M_0$ が階数3、$N_0,N_1$ が階数2とする。
$M_1,M_2,M_3$ の非空長方形を選ぶと各 $N_e$ の12セルが決まる。残る4セルの全ての補完を調べ、
階数1では長方形、階数2では零積で定まる行・列直線の同値閉包を使う。
最後に $M_0$ に前節のマッチング検査を施す。台の必要条件を通る数は109034088、
四層の自由辺の指数増大検査を通る数は15000である。
うち256個では五長方形の56係数が全て非零である。残る14744個を
片側欠色窓と同じ文字置換・平行移動・巡回・反転で集めると312軌道となる。
階数2の行射影には直線同値閉包、階数3にはマッチング上界を使い、
連続純関係と一文字商で検査すると、前節と同じビット順で
$\mathtt{0ff00f0f00ffcfff}$ と $\mathtt{0fff0fff0fffcfff}$ の二代表だけが残る。
\end{proof}

\begin{proof}[五長方形が全て四行四列の場合]
階数1の指数行列の非零Smith対角は42個の1であり、対角基底変更は階数13である。
従って56係数の商には一変数 $a\ne0$ が残る。未定の8成分を $x_0,\ldots,x_7$ とすれば、
Smith変換の逆で得る三行列 $N_0,N_1,M_0$ は順に
\[
\begin{pmatrix}x_0&a&a&1\\1/a&x_3&1/a&1\\1/a&a&x_4&1\\1&1&1&x_7\end{pmatrix},\quad
\begin{pmatrix}a&x_1&1&a\\x_2&1/a&1&1/a\\a&1/a&1&x_5\\1&1&x_6&1\end{pmatrix},\quad
\begin{pmatrix}x_0&x_1&1&1\\x_2&x_3&1&1\\1&1&x_4&x_5\\1&1&x_6&x_7\end{pmatrix}.
\]
それぞれの $N$ の四変数を $(u,v,f,h)=(x_0,x_3,x_4,x_7)$ または
$(x_1,x_2,x_5,x_6)$ とする。$a(a^2-1)$ を可逆にすると、三次小行列式のイデアルは
$(u-v,uh-1,f-h,(h-a)(ah-1))$ に等しい。検算は可逆化変数を加えて両側の
Gröbner基底を比較するので、$a^2\ne1$ の特殊値も含む。よって $x_7,x_6$ はそれぞれ
$a,1/a$ の一方である。両者が等しければ $M_0$ は階数1、異なれば行列式 $-(a-1/a)^4$ で階数4になる。

$a=\pm1$ の各々では上記小行列式のイデアルは、異なる二変数 $y,z$ の全ての積
$(y-a)(z-a)$ で生成される。全変数が $a$ のときそれぞれの $N$ と $M_0$ は階数1である。
各 $N$ が階数2ならそれぞれ一セルだけが変わり、$N_0$ 側の変更は対角のある位置 $i$ にある。
$M_0$ が階数3なら二つの変更セルの列ベクトルと元の列直線は一次独立で、その像は $e_i$ を含む。
$N_0$ の像も $e_i$ を含むので、色零の純平方が連続して標準型に反する。
\end{proof}

\begin{proof}[二つの疎な代表]
五長方形の非零行・列の先頭を基準として同じSmith計算を行う。第一代表では30係数に対し
非零対角17個が全て1で自由変数は $a,b$、第二代表では42係数に対し28個が全て1で自由変数は $a$ である。
残る八つの係数 $x_i$ は語の辞書順に取り、ここでは全て非零である。
第一代表では二つの $N$ の三次小行列式から
\[
 x_1=x_2=x_5=1,\quad ax_1=b,\quad x_0=b,\quad ax_3=1,\quad
 bx_4=a,\quad a^3x_0=1
\]
を得る。従って $a=b$, $a^4=1$。第二代表では同じ検査で
$(x_0,x_1,x_2,x_3,x_4,x_5)=(a,a,1,1,a,a)$ かつ $ax_1=1$ となる。
いずれも $a=\pm1$ では $M_0$ の階数が高々2である。

第一代表の $a^2=-1$ では $\det M_0=-4(x_6+x_7)$ なので $x_7=-x_6$ に限られる。
このとき始点1、長さ4、文字次数2の商次元は $k(x_6)$ 上で9である。
実際、総文字次数2の64語を列とし、三位置の二次縮約の全列に外側の二文字を付けて
作る関係行列の階数は55である。特殊化によって階数は増えない。
一方、標準Koszul解消の文字別Euler式はこの次数で8を要求する。
具体的には長さ4で正則表現の畳み込みから来る寄与が相殺し、
長さ2の二関係の寄与 $\sum_b r_b r_{2+b}=8$ が残る（$r=(2,2,1,1)$）。
最上項は総文字次数零にだけ寄与する。従ってこの最後の候補も全点で排除される。

\path{computations/verify_klein_mixed_positive.py} と付属のC++17列挙が、
全台の数、二つの整数正規化、三次小行列式と例外値、最後の文字別階数を再現する。
\end{proof}

% !TeX root = ../papers/nc-p3.tex
\section{二つの厚い色の方向が異なる窓の稠密部分}
\label{sec:klein-transverse-dense}
\begin{theorem}\label{thm:klein-transverse-dense}
二切断にそれぞれ階数2の色が二つ、階数1の色が二つあり、各切断の二つの階数2の色の差が異なるとする。
四つの階数1長方形が四行四列である場合、全正則窓は次の二変数表示とその位置の巡回・反転・文字変更で尽くされる。
和は $K$ の和とし、$b=i+j$, $e=j+k$, $l=i+e$ とする。厚い色を
$b\in\{0,1\}, e\in\{0,2\}$ に取る。総次数零の係数を通常1とし、
両条件を同時に満たす16語だけを
\[
 f=i_{\rm lo}+i_{\rm hi}+b_{\rm lo} e_{\rm hi}\pmod2,\qquad
 w_{ijkl}=\begin{cases}u&f=0,\\v&f=1\end{cases}
\]
に置き換える。正則性の必要十分条件は $u\ne v$ であり、$u=0$ または $v=0$ も許す。
最小周期は常に1である。反対階数は
$(u-1)(v-1)(uv+u+v-3)=0$ で12、その他で16である。
\end{theorem}

\begin{proof}[48成分の被覆]
階数1の全48成分の指数行列のSmith非零対角は35個の1であり、対角基底変更の階数は13である。
よってそれらを同時に1とできる。残る16語を辞書順に $x_0,\ldots,x_{15}$ とし、四つの厚い行列の
全三次小行列式の零点を取る。その極小素イデアルは48個である。
共通平行移動と位置別の $(t,t+h,t+h+k,t+k)$ の加算 ($h\in\{0,1\}$, $k\in\{0,2\}$)、
反転、巡回（奇数回のときは文字1,2も交換）で以下の五軌道になる。
\begin{center}
\begin{tabular}{c|l|c}
個数 & 代表の16係数の順序 & 始点0、四次の零文字次元\\\hline
16 & $(a,b,c,d,-1,\ldots,-1)$ & 11\\
16 & $(a,b,c,d,1,\ldots,1)$ & 11\\
8 & $(a,a,b,b,c,c,d,d,z,z,z,z,z,z,z,z)$ & 11\\
4 & $(a,a,a,a,b,b,b,b,c,c,c,c,d,d,d,d)$ & 13\\
4 & 定理の $(u,v)$ 表示 & 9（一般点）\end{tabular}
\end{center}
第三行には $-bcz+adz-a+b+c-d=0$ の一条件を課す。これは $a$ に関して原始的な一次式であり、
全ての代表は代数閉体への基礎拡大後も既約である。
\path{computations/verify_klein_transverse_dense.py} は極小素分解を五代表から独立に計算し、
各代表の軌道のイデアルとの一致をGröbner基底で検査する。初めの四代表について、
各関数体で文字次数零の関係行列の階数を計算したものが表の次元である。
標準Koszulの文字別Euler式の値は $2^2+2^2+1^2+1^2-1=9$ であり、特殊化で商次元は減らない。
従って最後の行だけが残る。
\end{proof}

\begin{proof}[二次完全交叉への変換]
各位置で $y_{a+2c}=e_a+(-1)^c e_{a+3}$ ($a,c=0,1$、$a+3$ は文字の和) による二文字ずつの
Fourier変換を施す。得たテンソルの文字総次数は零である。新しい文字の二コサイクル
$(-1)^{i_{\rm hi}j_{\rm lo}}$ で捻り、位置1,3で第3成分の符号を反転し、位置2で文字に1を加える。
得られるテンソル $S$ は完全対称である。全体を2で割った $S$ の縮約行列を
$S_b=(S_{i,i+b,k,k+b+1})_{ik}$ と書く。$S_0$ の第0,1列は等しく、第2,3列も等しい。
異なる二列の差と和は
\[
 4(1,1,-1,-1)^t,\qquad 2(u+1,v+1,u+1,v+1)^t
\]
である。$S_1=S_0^t$、$S_2$ の全列は $(u-1,v-1,u-1,v-1)^t$、
$S_3$ の全行は $(u-1,v-1,u-1,v-1)$ に等しい。これらの恒等式は全パラメータで成立する。

$u=v=1,-1$ では縮約階数が6に足りない。その他では $S$ の対称二次直交補空間は四次元である。
その共通零点 $[A:B:C:D]$ の四式は
\[
 A^2-C^2=B^2-D^2,\quad
 (v+1)(A^2+C^2)=(u+1)(B^2+D^2),\quad
 (v-1)AC=(u-1)BD,\quad AD=BC .
\]
最終式で $(A,B,C,D)=(sp,sq,tp,tq)$ と書く。第一式は
$(s^2-t^2)(p^2-q^2)=0$。$s^2=t^2$ なら中央二式の $p^2,q^2$ に関する係数行列の行列式は
符号を除き $2(u-v)$ である。$p^2=q^2$ なら $u\ne v$ のもとで
$s^2+t^2=st=0$ となって同じく解はない。$u=v$ なら $s=t=p=q=1$ が共通解である。
従って基点自由はちょうど $u\ne v$。定理\ref{thm:symmetric-frobenius-chart}と
二コサイクル捻りの左右Koszul複体の対応により、全点の正則性が決まる。
\end{proof}

\begin{proof}[周期と反対階数]
一位置巡回の後に1,2を交換し、四位置で $(0,1,1,0)$ を加えると、16語の条件も $f$ の値も元に戻る。
この四位置の基底変更を繰り返すと一移動と添字保持同型である。
反対切断の四色行列は行・列の置換を除き $J+\operatorname{diag}(u-1,v-1,v-1,u-1)$ であり、
行列式の零点は記載の三曲線である。$u\ne v$ のもとで二つの対角補正が零なら階数3、
全補正が非零なら対角行列に階数1を加えたものなので零行列式でも階数3である。
\end{proof}

% !TeX root = ../papers/nc-p3.tex
\section{両側二重色の疎な台の全被覆}
\label{sec:klein-two-doublet-supports}
\begin{proposition}\label{prop:klein-two-doublet-supports}
二切断の色階数が $(2,2,1,1)$ の置換である正則表現型四元群窓を取る。
少なくとも一つの階数1長方形に零の行または列があるとする。
標準Koszul左右AS窓の台の候補は、厚い色の方向が平行なとき148軌道、異なるとき51軌道の表に尽くされる。
表は \path{computations/klein_doublet_parallel_masks.txt} と
\path{computations/klein_doublet_transverse_masks.txt} である。
\end{proposition}
\begin{proof}
二つの厚い色集合を $\{0,1\},\{0,q\}$ に平行移動し、平行なら $q=1$、
異なれば共通の線型変換で $q=2$ とする。各切断で残る二つの長方形にはそれぞれ $225^2$ 通りの選択がある。
二切断で指定される語の重なりは16セルであり、この上で一致する組を全て結合する。
48セルの指定後に未定なのは16セルである。

同時文字平行移動を含む $(t,t+h,t+h+k,t+k)$ の加算 ($h=0,1$, $k=0,q$)、
巡回、反転を作用させる。$q=2$ の奇数回の巡回に文字1,2の交換を伴わせ、
$q=1$ では文字2,3の同時交換も許す。これらは階数指定を保存する。
まず長方形の組の軌道代表を一つだけ選び、各代表について16セルの全ての補完を調べる。
元の台の個数を数えるときは、長方形の軌道の大きさを掛ける。全台の軌道に直すときは、
代表を保つ群の作用で補完をまとめる。二切断で指定された48成分が全て非零の軌道は、補完の前に除く。
\[
\begin{array}{c|rr}
 & q=1&q=2\\\hline
重なりが一致する長方形の組 & 1033395 & 973921\\
その軌道数（全非零も含む） &4856&8005\\
疎な補完で四つの階数2の台条件を通過&922893330&828689160\\
自由辺の強連結成分に分岐のない台&476008&292056\\
その軌道数&2779&2794\\
連続純関係・一文字商を通過する軌道&148&51
\end{array}
\]
階数2の台条件には零成分に対応する行直線と列の直交直線の一致を用い、
自由辺・純関係・一文字商の検査は片側欠色の場合と同じである。
各軌道は $16(i+j)+4i+k$ のビット順で最小の非負整数として記録する。
\path{computations/verify_klein_two_doublet_supports.py} はC++17列挙で上段の数と正規代表を求め、
下段の全代表について三次の成長検査を行い、二表と項ごとに比較する。
\end{proof}

% !TeX root = ../papers/nc-p3.tex
\section{異方向二重色の疎な窓と二変数族の境界}
\label{sec:klein-sparse-transverse}
\begin{theorem}\label{thm:klein-sparse-transverse}
二切断の厚い色の方向が異なる $(2,2,1,1)$ 型の四元群窓を、係数の零を許して考える。
その全正則窓は相異なる二点 $U,V\in\mathbb P^1$ による一つの族で被覆される。
有限な $U,V$ は定理\ref{thm:klein-transverse-dense}の $u,v$ である。
疎な正則窓には、台を保つ対角基底変更と前節の台の作用のもとで、
一つの孤立表示と一つの $y\in k^*$ による表示があり、
前者は $\{U,V\}=\{\infty,0\}$、後者は $\{\infty,1/y\}$ を与える。
全点で最小周期は1で、二点の順序交換と反転は同じ添字保持同型類を与える。
反対階数12の集合は $U=1$, $V=1$, $(U+1)(V+1)=4$ の射影閉包であり、それ以外は16である。
ここでパラメータの全重複を除いたことは仮定しない。
\end{theorem}
\begin{proof}[51の台から二表示へ]
命題\ref{prop:klein-two-doublet-supports}の異方向の51代表ごとに、各非零語に可逆変数を置く。
薄い行列の二次、厚い行列の三次小行列式から、可逆単項式を除いた二項式の指数行列を取る。
51件の非零Smith対角は全て単元である。符号のある右辺との整合を検査し、さらに対角基底変更を割る。
残りの小行列式を代入し、自由変数の積で飽和して極小素分解を取る。
空なものを除いて合計63成分となる。各素成分は、ある変数に関して一次の式を逐次解いて
残る変数を自由にすることで有理化できる。ここで有理化は関数体の計算であり、
上界は元の全縮約列で書いた関係行列の小行列式の消滅として成分全体に延びる。

始点を $0,1,2,3$ と変えて、長さ4の零文字次元を標準値9と比較する。
61成分ではいずれかの始点ですでに10以上となる。
残る台は十六進数 $\mathtt{0ff066663c3caa55}$ と $\mathtt{0ff066667e7eaff5}$ だけである。
第一代表は全非零係数が1、第二代表はSmith自由変数 $y\in k^*$ 一つで追加方程式はない。
両代表の全係数は \path{computations/verify_klein_sparse_transverse.py} の指数変換から再現される。
同検算は全51台の飽和・極小素分解と全63成分の関数体検査を含む。
\end{proof}

\begin{proof}[境界の完全対称表示]
第二代表で $y=t^2$ と書き、四位置の文字順に $t$ の指数
$(1,-1,0,0),(-1,1,0,0),(1,1,0,0),(1,1,0,0)$ を持つ対角倍率を掛ける。
第一代表には倍率を掛けない。どちらも全位置で文字に3を加える作用に不変となる。
従って定理\ref{thm:klein-transverse-dense}と同じ二文字ずつのFourier変換後の総次数は零である。
同じ二コサイクルで捻り、$z=t^{-2}$（第一代表では $z=1$）として、
位置1,2で $(1,z,1,-z),(1,z,1,z)$ を掛け、位置3では文字に1を加えた後に $(1,1,-1,1)$ を掛ける。
全体を2で割ったものは完全対称で、行列 $S_b=(S_{i,i+b,k,k+b+1})$ は
\[
 S_0=[A,A,C,C],\quad S_1=S_0^t,\quad
 S_2=(1,Q,1,Q)^t{\bf1}^t,\qquad S_3=S_2^t,
\]
ただし $A=(1,P,1,Q)^t$, $C=(1,Q,1,P)^t$、
第一代表で $(P,Q)=(3,-1)$、第二代表で $(P,Q)=((3y+1)/y^2,(1-y)/y^2)$ である。

従って平方成分の二条件は $A_0^2=C_0^2$ と $B_0^2+D_0^2=(P+Q)A_0^2$、
他の二式は $Q A_0 C_0=B_0 D_0$, $A_0D_0=B_0 C_0$ である。
$[A_0:B_0:C_0:D_0]=[sm:sn:rm:rn]$ と書く。$m=0$ なら
$s^2+r^2=sr=0$ となって不可能である。$m\ne0$ なら $s^2=r^2\ne0$ となり、
$2n^2=(P+Q)m^2$, $n^2=Qm^2$ から $P=Q$ が必要である。
実際には全点で $P-Q\ne0$ なので、対称二次完全交叉の判定により全て正則である。
\end{proof}

\begin{proof}[射影パラメータと特殊化]
一般に $S_0$ の二種の列を $(a,b,c,d)^t,(c,d,a,b)^t$、
$S_2$ の列を $(c,d,c,d)^t$、$S_1,S_3$ をそれぞれの転置とする。
これが上記の対称表示をまとめた $S$ である。
全位置に同時に $\operatorname{diag}(s,r,s,r)$ を作用させると、
$(a,c)$ は $s^3r$ 倍、$(b,d)$ は $sr^3$ 倍される。従って二対の独立な非零倍率は基底変更で吸収できる。
有限なパラメータを $[a:c]=[U+3:U-1]$, $[b:d]=[V+3:V-1]$ とし、これを射影直線に拡張する。
疎な第一対は $[1:1]$ で $U=\infty$、第二対を $4/(P-Q)$ 倍すれば本文の $V$ が得られる。
この多項式表示が各境界点への特殊化も与える。

元の無印窓は $S$ の係数に四語の二コサイクル符号を掛けたものと四因子同値である。
実際、途中の文字の平行移動を符号の前後で交換して生ずる差は、双指標性により各位置の対角符号の積である。
$S$ の総次数は一定なので、巡回に伴う符号の差も移動させた一文字だけの関数となる。
反転では指数の差が $(\sum i_{\rm hi})(\sum i_{\rm lo})-\sum i_{\rm hi}i_{\rm lo}$ となり、やはり位置ごとに分離する。
従って全点が一周期で反転は同じ類となる。全位置への1の加算は二つの射影パラメータを交換する。
反対階数は有限部分で前定理に従う。疎な曲線では各反対ブロックの行列式の非単元因子は $y-1,y+1$ だけであり、
両零点で各ブロックが階数3、孤立表示では全ブロックが可逆である。これは上の射影閉包に一致する。
\end{proof}

% !TeX root = ../papers/nc-p3.tex
\section{異方向二重色族の無印同型の完全判定}
\label{sec:transverse-doublet-isomorphism}
\begin{theorem}\label{thm:transverse-doublet-isomorphism}
定理\ref{thm:klein-sparse-transverse}の全窓について、四元群・対称表示・文字・始点での枠を忘れた
添字保持同型類を次のように分類する。二つの射影パラメータを $p=a/c=(U+3)/(U-1)$,
$q=b/d=(V+3)/(V-1)$ と書き、初めから順序を忘れる。常に $p\ne q$ である。
\begin{enumerate}
\item $p,q$ が有限で $p+q\ne0$ の場合は反対階数16であり、唯一の追加同一視は
$\{p,q\}\sim\{-p,-q\}$ である。
\item 一方が $\infty$ の場合は反対階数12であり、他方の有限パラメータ $q$ の二乗が完全不変量である。
\item $\{p,-p\}$ ($p\ne0$) の場合も反対階数12であるが、前項とは非同型であり、$p^4$ が完全不変量である。
\end{enumerate}
全て最小周期1である。従ってこの族の無印像は既約な二次元の族であり、
射影パラメータの二つの階数12境界での追加同一視も上記で尽くされる。
\end{theorem}

\begin{proof}[可逆な反対切断]
$W$ を前節の対称テンソル $S$ に四語の二コサイクル符号を掛けた表示とする。
$V$ を反対切断 $02\mid13$ の逆行列の転置とすると、四位置の独立な基底変更に対して $V$ は四重双対テンソルとして変換する。
従って $L=W_{01\mid23}V_{01\mid23}^{t}$ は相似変換を受ける。
分子・分母を $c=d=1$ に取り $t=p+q$, $e=p-q$ と置く。
$L$ の特性多項式は
\[
 z^{10}(z-4)^2\big((z-2)^2-A/4\big)^2,\qquad A=(t+4/t)^2 .
\]
$T=\operatorname{tr}_{V_1}(L^2)$ を第1位置について部分traceしたものとすると、
$T$ の二つの対角値はそれぞれ二度現れ、その差は $e(t-4/t)/4$ である。
よって $B=16\big(\operatorname{tr}(T^2)-(\operatorname{tr}T)^2/4\big)=e^2(A-16)$ も無印不変量である。

三つ目は三組の $W,V$ を収縮して
\[
 C=\tfrac12\sum W_{abcd}W_{efgh}W_{ijkl}V^{abgh}V^{ejkd}V^{ifcl}
   =-e^2(1+4/t^2)+\tfrac72(t^2+16/t^2)+100
\]
とする。添字の和は四位置ごとに同じ空間とその双対の間で取る。
\path{computations/verify_transverse_rank_sixteen_isomorphism.py} が特性多項式・部分traceを記号計算する。
最後の収縮では $4tV$ と $W$ の各成分が $p,q$ の一次以下であることを確認した上で、
分母を払った次数6以下の恒等式を各変数7点ずつの整数補間で証明する。

$A\ne16$ なら $e^2=B/(A-16)\ne0$ であり、$t^2+16/t^2=A-8$ と $C$ から $1/t^2$ が一意に決まる。
$A=16$ なら $t^2=4$ であり $C$ が直接 $e^2$ を与える。従って全点で二乗 $t^2,e^2$ が復元される。
$t,e$ の符号の選択は $p,q$ の交換と両符号の同時反転に等しい。
後者は全位置で高位ビットが1の文字を $\sqrt{-1}$ 倍し、全体を $-1$ 倍して実現できる。
\end{proof}

\begin{proof}[二つの特異反対切断の区別]
反対切断の左核を四行四列の行列空間と見て、四次元の核の上に行列式を制限する。
この四次曲面の射影同値類は無印不変である。$p=\infty$ では $a=d=1,c=0,b=q$ として、
核のある座標 $z_0,\ldots,z_3$ で行列式は $-z_2^2 z_3^2$ である。
$\{p,-p\}$ では定数倍を除き
\[
 (z_1^2-z_2z_3)
 \left(p^4(z_2+z_3)^2+4p^2 z_0^2-(z_2-z_3)^2\right)
\]
となる。$p\ne0$ なのでいずれの二次因子も階数3であり、二つの重平面の型に一致しない。
\end{proof}

\begin{proof}[四重交代式とその勾配]
四つの $W$ の各位置をそれぞれ四次交代形式一つに縮約し、第一位置の文字順を固定して正規化した
四次相対不変式を $D$ とする。その勾配を $G$ と書く。
四位置の基底変更に対し $D$ には行列式の積が掛かり、$G/D$ は $D\ne0$ で四重双対テンソルとなる。
上の基底では
\[
 D=16(a^2b^2+a^2d^2+2abcd+b^2c^2+7c^2d^2).
\]
$V=G/D$ として前と同じ $L,T$ を作る。$p=\infty$ の場合に $X=q^2$ と書くと
$D=16(X+1)$ であり、$D=0$ はすでに一つの $X$ を決める。
$X\ne-1$ のとき三つの不変量 $\operatorname{tr}(L^2)/2$, $\operatorname{tr}(L^3)/2$,
$\operatorname{tr}(T^2)-(\operatorname{tr}T)^2/4$ は順に
\[
 \frac{16X^2+40X+11}{8(X+1)^2},\quad
 \frac{64X^3+336X^2+252X+31}{32(X+1)^3},\quad
 \frac{(4X-1)^2}{256(X+1)^4}
\]
となる。三関数について $f(X)-f(Y)$ の分子を取り、
$(X+1)(Y+1)(X-Y)$ を可逆化したイデアルのGröbner基底は $(1)$ である。従って $X$ が復元される。
$\infty$ は符号反転で固定され、有限な $q$ の符号反転は同型となる。
実際、前の高位ビットの倍率で $(a,b,c,d)$ は全体スカラーを除き $(-a,-b,c,d)$ になる。
独立な射影倍率で $[-a:0]=[1:0]$ に戻せる。

$\{p,-p\}$ の場合は $X=p^4$, $D=16(X+7)$ である。$X\ne-7$ のとき $L$ の特性多項式は
\[
 z^{10}\left(z-\frac{X+2}{X+7}\right)^4\left(z-\frac{10}{X+7}\right)^2 .
\]
二つの固有値が合流するのは $X=8$ のみであり、そうでなければ後の固有値が重複度で読める。
前の固有値が零の場合にも後者は非零である。$D=0$ は $X=-7$ だけなので全点で $p^4$ は必要である。

十分性には $p\mapsto\sqrt{-1}\,p$ を実現すればよい。高位ビットを固定して二文字の低位ビットをFourier変換する。
この表示で原始八乗根 $\zeta$ の次の指数を四位置にそれぞれ対角作用させる。
\[
 (-7,-7,-6,-10),\quad(3,3,4,0),\quad(3,3,4,0),\quad(3,3,4,0).
\]
係数が $p$ に比例する非零語では指数和は8を法として2、他の非零語では0である。
従って $p$ だけが $\zeta^2$ 倍になる。
\path{computations/verify_transverse_rank_twelve_isomorphism.py} は二つの核の基底からの行列式、
交代式の勾配、三有理関数の単射性と最後の全係数の位相を検算する。
\end{proof}

% !TeX root = ../papers/nc-p3.tex
\section{平行な二重色で片方の文字対に限られる薄い関係}
\label{sec:klein-parallel-flags}
\begin{theorem}\label{thm:klein-parallel-flags}
両切断で厚い色を $0,1$ とし、高位ビットで文字を二対 $\{0,1\},\{2,3\}$ に分ける。
ある位置で、薄い色の関係の行の台が一方の文字対だけに含まれるとする。
平行方向の表を零始まりで番号付けすると、正則な台はちょうど
\[
 50,51,52,54,55,56,57,58,59,62,63,64,65,66
\]
である。これらの台では指定の二切断階数を持つ全係数が二次PBWで左右Koszul AS正則である。
さらにこの行条件を課さない部分でも、表の $76,78,112,115,117$ は実現でき、同じ全点での結論を持つ。
\end{theorem}
\begin{proof}[係数に依らない辞書順]
前の表の148代表のうち、上の行条件は0番から66番にちょうど対応する。
行集合および列集合は四位置を巡るので条件は巡回・反転で保たれる。
位置 $p$ ごとに文字順 $\pi_p$ を選び、語に位置別の左辞書順を入れる。
一色の行 $i$ に対応する二文字語は $i,i+b$ の一個なので、先頭語は行を大きい順に調べて求まる。
階数1なら最初の非零行、階数2なら最初の非零行と、それと独立な後続の最初の行を取る。
零成分が強制する同一直線の類に属する行は飛ばせる。非零成分の台が異なれば二行は必ず独立である。
同じ台だが同一直線と決まらない場合には、その順序では証明を行わない。

各行 $i$ でこの規則が選ぶ相手の集合を $F_i$ とする。$\pi_p$ の順に相手の集合が
大きさ $0,1,2,3$ の入れ子になるなら、足りない最後の文字を加えて次位置の順序 $\pi_{p+1}$ が定まる。
これを四回繰り返して元の順序に戻るものを選ぶ。四位置の $4!$ 通りの初期順序の検査で、
定理に挙げた19台に加え、$77,113,114,116$ の台がこの検査を通る。
各位置で先頭語は番号の降順対の六個であり、長さ3で残る語は20個である。
四つの一文字縮約は異なる文字次数を持ち非零なので、三次関係の交叉が次元4以上で、商の次元は20以上である。
従って全ての三文字重複が解消する。Diamond補題による非減少語の基底と二次双対性から
全次数のKoszul性を得る。$W$ の縮約による $1$ 対 $3$、$2$ 対 $2$ の対合が完全なので、
左右AS性も従う（\path{sections/unmarked_klein_pbw.tex}の辞書順の証明を参照）。
\end{proof}

\begin{proof}[実現と残る53台の全排除]
二切断の必要小行列式で、非零語の単項式を除いて二項式となるものをSmith標準形で解く。
右辺の符号を含め、指数の $1,2,4$ の全ての根を列挙し、連結する対角基底変更を割る。
各代表の係数はこうして得た自由変数のLaurent単項式と定数位相の積で与えられる。
他の小行列式を代入し、変数積を可逆化する。23台でこの解集合が空なのは
$77,113,114,116$ だけである。非空なら上記の行主元が下からも階数を保証するため、
どの解も標準AS窓である。この表示は位相を固定した係数商も再現し、例外値に追加の開条件を要求しない。

残りの0--66番の53台を排除するには、Laurent単項式の比だけで決まる二列基底を使う。
各厚い色で、ある二次小行列式がLaurent単元となる二列を固定する。薄い色では非零列を一つ固定する。
指定階数を持つ任意の実現でこれらは色関係の基底である。残る小行列式をまだ課さない、より大きい
Smithパラメータトーラス上で、固定した列に二文字を付けた長さ4・零文字の関係行列を作る。
その関数体で商次元が9を超えるなら、実現する全特殊化で同じ障害が残る。
5,13,27,36,47番の一部の位相ではこの上界だけを使わず、残る小行列式が可逆化後に単位イデアルとなることを確認する。
これで53台の全位相が尽くされる。
\path{computations/verify_klein_parallel_flags.py} は台と順序の全比較、両切断の係数正規化、
基底の単元小行列式と各位相の障害となる始点、19台の非空性を検算する。
\end{proof}

% !TeX root = ../papers/nc-p3.tex
\section{三つの薄い色を持つ対称族の無印同型}
\label{sec:three-lines-symmetric-isomorphism}
\begin{theorem}\label{thm:three-lines-symmetric-isomorphism}
定理\ref{thm:klein-three-lines}の対称表示の二つの正則窓では、原則として四係数 $p_i$ の多重集合の一致が
添字保持同型の必要十分条件である。例外は二値がそれぞれ二回現れる場合だけである。
この場合の多重集合を $\{a,a,b,b\}$、$r=(a-b)/(a+b)$ と書く。
同じ型の二窓は $r'^2=r^2$ または $r'^2=r^2/(r^2-1)$ のとき、またそのときに限り同型である。
他の型との同型はない。
\end{theorem}
\begin{proof}[四つの二次式の束と自己随伴作用素]
系\ref{cor:symmetric-chart-isomorphism}より二次式の四次元空間 $L$ の合同を判定すればよい。
元の文字順で $p=(a,b,c,d)$ とすると、その行列は
\[
 \begin{pmatrix}
 ta&x&y&z\\ x&tb&-z&-y\\ y&-z&tc&-x\\ z&-y&-x&td
 \end{pmatrix}.
\]
$t=0$ の網を $N$ とする。一文字の符号を変え、四文字の正規直交Fourier基底に移すと、
$N$ は跡零の全対角行列、残る第四行列は対角成分が全て $1/4$ の四元群の巡回行列になる。
非対角成分は色ごとに一定で、4倍すると順序を除いて
$r=a+b-c-d$, $s=a+c-b-d$, $h=a+d-b-c$ となる。

全 $Q\in L$ について $T^tQ=QT$ を満たす作用素の空間を取る。$N$ の条件から $T$ は対角であり、
第四行列の非零な非対角成分の両端で対角値が一致する。従って $(r,s,h)$ の非零成分が
少なくとも二つ、一つ、零の場合に、それぞれ次元1、2、4である。これは合同不変量なので三者を分離する。
次元4なら $p_i$ は全て $1/4$ で一つの類である。
\end{proof}

\begin{proof}[自己随伴空間が一次元の場合の網の復元]
$\mathbb P N$ に階数2の行列が六つあり、跡零対角網の行列式は四直線に分かれる。
$t=1$ の階数2以下の点を調べる。$X=4x,Y=4y,Z=4z$ とする。相補的な二つの余因子の和を取ると
\[
 rYZ=(sh-r)X,\quad sXZ=(rh-s)Y,\quad hXY=(rs-h)Z .
\]
$ r,s,h$ の一つだけが零ならこれらは $X=Y=Z=0$ を強制し、元の対角行列は可逆である。
三つとも非零なら右辺の倍率を左の係数で割り $\alpha=sh/r-1,\beta=rh/s-1,\gamma=rs/h-1$ と置く。
$\alpha=0$ なら $r=sh$ のもとで $4p_i$ は $(1\pm s)(1\pm h)$ の四通りであり、
どれも非零だから $\beta\gamma\ne0$。従って $Y=Z=0$、二つの対角二行ブロックの行列式から
$X$ は高々二通りである。他の倍率が零でも同様である。

$\alpha\beta\gamma\ne0$ のとき、座標零の解は除かれるため
$X^2=\beta\gamma$, $Y^2=\alpha\gamma$, $Z^2=\alpha\beta$ である。
四符号の候補はアフィン独立である。他の平面が六節点（階数2の点）を含むなら、
そのうち三点がアフィン、三点が無限遠でなければならない。無限遠の三点を通る直線は
四直線の一つ $\pm X\pm Y\pm Z=0$ である。実際六点は四直線の二本ずつの交点で、他の三点直線はない。
同じ一次式が四符号点の三点で同じ値になると $X^2=Y^2=Z^2$、
従って $r^2=s^2=h^2$ となる。これは $p_i$ の三つが等しい場合である。
それ以外では六節点の平面は $\mathbb P N$ だけである。

この網の射影安定群はFourier基底の符号付き置換である。四本の行列式直線の各一般点の核が
四座標点を復元し、跡零を保つ対角変換の四平方は等しいからである。
第四行列の対角の跡を揃えた後、非零な相補二辺の成分が等しくなるには四符号の積が $+1$ でなければならない。
すると三つのFourier係数に対する作用は置換と偶数の符号反転であり、これは元の四つの $p_i$ の全置換に正確に対応する。
\end{proof}

\begin{proof}[五つの網が現れる場合]
三つが等しい配置を $(a,b,b,b)$, $a\ne b$ とする。Fourier基底を選び直すと第四行列は $bI+(a-b)J/4$ である。
束は座標四形式とその和の平方、合わせて五つの階数1二次式の張る五次元空間内で
\[
 (1,1,1,1,-16b/(a-b))
\]
を法線とする超平面である。五つの一次形式は射影枠をなし、どの四つも独立である。
法線のどの成分も零でない。五座標のうち二つだけが非零となる十点が階数2であり、これで全てである。
三つか四つだけが非零ならその個数が階数であり、五つ全部が非零なら可逆な対角行列に階数1を加えても階数3以上となる。

法線の各成分で五座標をそれぞれ倍率変更すれば、これらの十点は成分和零の空間の $e_i-e_j$ である。
平面上に入る対は平面の方程式の五係数が等しい対に対応するため、六点を持つのは四係数が等しい五平面だけである。
各平面の網は対応する四つの独立な一次形式の平方で対角化される。
四本の行列式直線の核から双対の四座標形式を復元すると、その平方全体の張る空間も復元される。
これら五つの空間の和が元の五次元空間であり、その階数1点はちょうど五つである。
ゆえにこの射影枠は無印で復元される。五一次形式を和が零になるように規格化すれば、その置換に伴う倍率は
共通スカラーだけであり、法線の係数は置換と共通倍率まで復元される。
四つが1の多重集合は最後の値を一意に決める（五つとも1の場合を含む）。$a+3b=1$ より元の係数も決まる。
\end{proof}

\begin{proof}[二つの二重値]
自己随伴空間が二次元なら、四係数は二値の二対であり、二つの原始冪等元が二平面の無順序対を復元する。
第四行列の対角を1に規格化し、非零な二本の相補辺の成分を $r$ とする。
$r\ne0,\pm1$ である。一般の行列の対角を $d_0,\ldots,d_3$ とすると $t=(\sum d_i)/4$、
二本の非対角成分は $rt$ である。片方の二平面だけに支えられる二次式は各々 $x^2-y^2$ の直線であり、
その等方二直線が決まる。これらの式を割った後の各二平面での座標を、等方二直線での値
$U_\pm=d_0+d_1\pm2rt$, $V_\pm=d_2+d_3\pm2rt$ に取る。
二平面間のグラフ行列は
\[
 \binom{V_+}{V_-}=\frac1r\begin{pmatrix}1&-1-r\\1-r&-1\end{pmatrix}\binom{U_+}{U_-}.
\]
各等方対を保つ変換は二行・二列を別々に置換・倍率変更できる。値に現れる倍率は平方だが代数閉体上で任意に取れる。
逆にこれで商のグラフを合わせれば、その逆像である元の束も一致する。二平面の交換は逆行列を取る。
全成分非零の二行二列行列の行列倍率商は両対角積の比で決まり、
ここでは $1/(1-r^2)$ とその逆だけが同一視される。これが定理の二条件である。
\path{computations/verify_three_lines_symmetric_isomorphism.py} がFourier変換、余因子、五平面と二平面のグラフを検算する。
\end{proof}

% !TeX root = ../papers/nc-p3.tex
\section{三薄色の捻りの無印同型と他枝との交差}
\label{sec:three-lines-twisted-isomorphism}
\begin{theorem}
定理\ref{thm:klein-three-lines}の捻った表示では、四係数 $p_i$ の多重集合の一致が添字保持同型の必要十分条件である。
二対反復の場合にも追加同一視はない。三薄色の二族、定理\ref{thm:transverse-doublet-isomorphism}の異方向族、
および一方の切断だけに欠色のある正則曲線は、無印の同型類として互いに交わらない。
\end{theorem}
\begin{proof}[関係空間の低階数曲線]
捻りを二対の内部と対の間の符号に分け、第 $i$ 行に $(-1)^{i_{\rm hi}i_{\rm lo}}$ を掛ける。
隣接関係空間 $R$ は次の行列の空間になる。
\[
 \begin{pmatrix}A&x&y&z\\x&B&z&y\\y&-z&C&-x\\-z&y&-x&D\end{pmatrix},
 \qquad p_0 A+p_1 B+p_2 C+p_3 D=0 .
\]
以下その対角を $d_i$、混合成分を $\xi_1=x,\xi_2=y,\xi_3=z$、対角部分空間を $D_p$ と書く。
階数2以下の射影閉集合は左右独立基底変更の不変量である。
$i,j$ の補対を $k,l$、$h=i+j\in K$ とすると、随伴行列の $(i,j),(j,i)$ 成分の符号付き和差の半分は、
各々符号を除いて
\[
 \xi_h(x^2+y^2-z^2-d_kd_l),\qquad
 \Big(\prod_{b\ne h,\, b>0}\xi_b\Big)(d_k-d_l)
\]
である。全 $\xi_b\ne0$ なら対角は定数で、$\sum p_i=1$ より零となり、残る点は平面二次曲線上にある。
二成分だけが非零なら、消えた色の二対で対角がそれぞれ一致する。二つの対角値には非自明な一次条件があり、
射影平面内でさらに上の非零な二次式が課される。従っていずれも正種数の成分を持たない。
三成分とも零ならちょうど六つの階数2点である。

一成分 $\xi_h$ だけが非零なら行列は二つの二行ブロックに分かれ、双方の行列式が零になる。
$h=3$ では $\xi_h$ を $\sqrt{-1}$ 倍して、$d_i d_j=d_kd_l=\xi_h^2$ と書ける。
$\xi_h=1$ のアフィン部分は $d_i=s,d_j=1/s,d_k=t,d_l=1/t$ により
\[
 \lambda=p_i s+p_j/s=-(p_k t+p_l/t)
\]
となる。二つの二次被覆の分岐対は $\{\pm 2\sqrt{p_i p_j}\},\{\pm 2\sqrt{p_k p_l}\}$ である。
対が異なれば交わらず、合成体の次数は4、分岐指数2の点が計八つで、正規化の種数は1となる。
この曲線の張る空間は混合成分を一つ許した $\mathbb P^3=P_h$ 全体である。実際 $1,\lambda$ と
二つの相異なる判別式の平方根が定数体上一次独立で、$1,s,s^{-1},t,t^{-1}$ の間の一次関係は表示した一つだけである。
分岐対が等しければ二つの二次体は一致し、成分は有理的である。
以上で階数2以下の集合の次元は高々1、正種数成分の個数は相補対積の異なる分割数と分かった。
\end{proof}

\begin{proof}[楕円曲線が複数ある場合]
二つの相補対積の等式、例えば $p_0p_1=p_2p_3$, $p_0p_2=p_1p_3$ から
$p_1=\pm p_2$, $p_0=\pm p_3$ で符号は共通となる。負号は成分和1に反する。
従って楕円曲線が一本なら二値の二対、零本なら全係数が等しく一つの類だけである。
二本以上ならその張る空間の交わりは $\mathbb P D_p$ であり、色を指定せず復元できる。
内部の六点の像と行空間は両側の四文字から二つを選ぶ六平面である。相異なるそれらの一次元の交わりから
四つの文字直線が決まる。像と行空間で同じ対が組になるため、隣接位置の文字置換は一致する。
四位置に適用すると全位置共通の置換と位置別対角変更だけが可能である。

全ての四文字置換は四元群のアフィン変換である。平行移動の捻り符号への寄与は位置別指標と共通定数に分かれる。
線形変換では交代化が二元体上の唯一の非零交代形式なので、元との比は対称双指標になる。
これは代数閉体上でコバウンダリである（基底の二文字で対角符号の平方根を選び、和での値を乗法則で定める）。
総色零の四語では従って各位置の倍率で吸収できる。よって全ての置換が実現する。

置換を除いた倍率を $f_j(i)$ とする。非純語は全て同じ共通倍率を受ける必要がある。
$(i,i,k,k)$ の全順列を比較すると $f_j(i)/f_j(k)$ は位置に依らない。そのため位置別定数を除き倍率は一つの関数 $f(i)$ である。
全 $i\ne k$ で $f(i)^2 f(k)^2$ が一定だから四平方は等しく、純語も非純語と同じ共通倍率になる。
従って $1-1/p_i$ が一致し、これには純成分が零の値も含まれる。
\end{proof}

\begin{proof}[一本の楕円曲線]
$p=(a,a,b,b)$、$r=(a-b)/(a+b)\ne0,\pm1$ とする。一意な楕円曲線の張る $P_1$ では $y=z=0$ である。
その中の階数2以下の点は楕円曲線の閉包と二つの孤立点からなる。
実際 $x=0$ の六点のうち、二つのブロックから一文字ずつ取る四点が二行列式の共通零点で、上の曲線の無限遠点である
（各被覆の $\lambda=\infty$ の二点は非分岐）。他の二点だけが一つのブロックに支えられる。
孤立点の像と行空間が左右それぞれに二平面の分解を与える。各ブロックへの射影像は対称二行行列の全空間であり、
その左右の同一視はスカラー以外一意である。実際、左右変換がこの空間を保てば全ての階数1平方の像も対称となり、
二つの変換は全ての射影直線上の点で一致する。従って各孤立点 $\operatorname{diag}(1,-1)$ は、合同の意味で固有の等方二直線を与える。

各等方対での評価を $U_\pm=A+B\pm2x$, $V_\pm=C+D\mp2x$ と取る。
$V_++V_-=-(a/b)(U_++U_-)$、$V_+-V_-=-(U_+-U_-)$ である。
このグラフ行列の対角積と反対角積の比は $((a+b)/(a-b))^2$ である。行列の行・列の個別倍率と置換、および
二平面交換による逆行列でこの比は逆数までしか変わらない。従って $r'^2=r^2$ または $r'^2=1/r^2$ が必要である。

後者を調べる。定理\ref{thm:transverse-doublet-isomorphism}の四次交代収縮を $D$、勾配を $G$ とする。
この規格化では $D=64(3r^4-5r^2+6)/(r^2-1)^2$。分子と逆順多項式は互いに素なので、$D=0$ は逆数と同型にならない。
$D\ne0$ のとき $L=W_{01\mid23}(G/D)_{01\mid23}^{t}$ を作る。
$k=2,3$ で $\operatorname{tr}(L^k)(r)-\operatorname{tr}(L^k)(1/r)$ の既約分子の最大公約数は定数倍を除いて $r^4-1$ である。
よって逆数と一致するのは最初から $r^2=1/r^2$ の場合だけである。二対の交換で $r$ の符号は変えられ、
$a+b=1/2$ から多重集合が決まる。
\end{proof}

\begin{proof}[異方向の二重色と一欠色曲線からの分離]
反対階数で三薄色の対称族を分離でき、捻りと異方向族の共通候補は階数16のみである。
その異方向表示の関係空間は、同じ行符号変更を施すと
\[
 \begin{pmatrix}
 ps+t&u&y&z\\u&qs+t&z&y\\y&-z&s+pt&v\\-z&y&v&s+qt
 \end{pmatrix}
\]
となる。$s=t=0$ の射影三空間内の二次曲面 $uv=y^2-z^2$ では全ての三次小行列式が消える。
三薄色の捻りの関係空間にはこの次元の低階数集合がなく、無印でも非同型である。

一欠色曲線が周期1となるのは $X=(t+1/t)^2=\pm4$ のみである。
反対切断の行列式は色0,1で $t^2(t^4+6t^2+1)/(t^2+1)$、色2,3で $-t^2(t^2-1)^2/(t^2+1)$。
$X=4$ は $t=1$ で代表され、左核の一般行列の行列式は基底を選べば
$(z_0^2-z_2^2)(z_1^2+z_3^2)/2$ となる。反対階数は12であり、この四相異平面は
異方向族の階数12の二境界の二重平面対および二つの既約二次錐と異なる。
$X=-4$ では消える二色の行列式が単純零で各階数は3、反対階数は14である。
共通候補は三薄色の捻りの $(-1,1,1,1)/2$ の置換だけである。両者とも $D\ne0$ で
$\operatorname{tr}(L^2)$ は順に $43/16,124/27$ となり分離される。
\path{computations/verify_three_lines_twisted_isomorphism.py} が余因子、例外直線の分子の最大公約数、
二次曲面、および欠色曲線の二つの交点候補を厳密計算する。
\end{proof}

% !TeX root = ../papers/nc-p3.tex
\section{平行二重色の二つの曲面と異方向境界の貼り合わせ}
\label{sec:parallel-symmetric-surface}
\begin{theorem}\label{thm:parallel-symmetric-surface}
平行二重色の台94,95,96の全係数は、位置別文字変更の後、次の対称テンソル $S(p,q)$ とその四元群捻りの二曲面で尽くされる。
順列で同じになる成分を一度ずつ書くと
\[
 S_{0001}=p,\quad S_{0111}=S_{0122}=S_{0023}=S_{2333}=1,\quad S_{2223}=q
\]
であり、他は零とする。$(p,q)\in k^2$ で、両方零、片方零、両方非零が順に台94,95,96を与える。
両曲面の全点が標準Koszul左右AS正則で最小周期1、反対階数は順に6と12である。
各曲面内での無印の添字保持同型は、$(p,q)$ の交換と同時符号反転だけで尽くされる。
異方向二重色族と交わるのは捻りの対角線 $p=q$ だけで、その点は同族の射影対 $\{\infty,p\}$ と一致する。
\end{theorem}

\begin{proof}[全係数の被覆と全点の正則性]
$S$ の総色は1であり、第二文字を1だけ平行移動すると総色零の三つの代表台に正確に一致する。
台ごとに両切断の小行列式の二項式をSmith標準形で解くと、対角倍率を除いた解は順に次元0,1,2のトーラスが二つずつで、
残る小行列式は恒等的に零である。実際、定理\ref{thm:klein-parallel-flags}と同じ指数計算で、
語上の位置別倍率の指数行列の階数は13、非零の $p,q$ の指数列を加えた階数は順に13,14,15になる。
従って各表示のパラメータトーラスは一成分上に全射する。捻りとの相違は以下の反対階数で確かめられる。

$S$ の四文字に対応する線形形式を $X,Y,Z,T$ とする。隣接切断の四ブロックの階数は $2,2,1,1$ である。
対称関係空間中に平方が入る条件は、次の四つの二次式の共通零点である。
\[
 L=\langle X^2-pY^2-T^2,\ Z^2-Y^2-qT^2,\ YT,\ XT-YZ\rangle .
\]
$Y=0$ で $T\ne0$ なら第四式から $X=0$ となり第一式に反する。$T=0$ で $Y\ne0$ でも同様で、
$Y=T=0$ なら $X=Z=0$ が従う。従ってこの四式は基点自由で定理\ref{thm:symmetric-frobenius-chart}を適用できる。
斉次二コサイクル捻りは左右のKoszul性と双対完全性を保つ。捻り符号の一位置巡回による比は
移した一文字と総色1の間の交代双指標、すなわち位置別対角倍率である。従って両方が周期1である。
\end{proof}

\begin{proof}[捻りでの四文字の復元]
反対切断の色0,1の行列式はともに $-4$、色2,3はそれぞれ階数2である。
左核を行列とみなし列を $3,2,1,0$ の順にすると、その一般元は
\[
 \begin{pmatrix}z&0&0&0\\x&-t&0&0\\0&0&z&0\\0&0&y&t\end{pmatrix}.
\]
従って階数1の射影点は $x$ 軸と $y$ 軸の二点だけである。元の文字順でこれは両側の双対空間の文字1,3の二直線を復元する。
右核でも同じ二直線が得られる（次の始点を取ってもよい）。各位置で、その共通零空間 $E=\langle e_0,e_2\rangle$ と、
商空間の二つの文字直線が定まる。

この二段のフィルトレーションの随伴次数で、商文字が三つある成分を見る。各配置で
$0111$ と $2333$ の順列だけがあり、一つの偶位置を除くどの三位置でも同じ商文字を使う。
従って四商での文字交換は全て同じで、そこに伴う偶文字の二直線とその対応も決まる。
二つの文字対の同時交換を除けば、残る変換は対角倍率と、位置 $j$ での
$e_1\mapsto b_j e_1+u_j$, $e_3\mapsto d_j e_3+v_j$ ($u_j,v_j\in E$) である。
商文字を二つだけ残す成分は零である。二つの文字1を残した場合、他の二位置 $j,l$ には
$e_0\otimes u_l$ と $u_j\otimes e_0$ の非零倍率付き和だけが寄与する。
よって全 $u_j$ は $e_0$ 方向で、$e_0$ の位置別倍率 $a_j$ で割ったその係数は
全ての二位置で和が零となる。標数零より全て零である。
文字3を残した場合も同様である。ここで $2333$ の四つの順列の符号は、
唯一の文字2の位置だけで決まり、あらかじめその位置の文字2の倍率に吸収できる。よって $v_j=0$ である。

二文字対の交換は四元群の平行移動で、その捻り符号の変化は位置別倍率で実現する。
これを除くと、位置別対角倍率の比は $0111,0122,0023,2333$ の各語の全順列で同じでなければならない。
同じ語に含まれる二文字を任意の二位置で交換でき、文字間のグラフが連結なので、倍率は位置別定数を除いて四位置共通である。
初めの四つの非零成分を1に保つ文字倍率の指数行列を $M$、$0001,2223$ の指数行列を $N$ とすると
\[
 2NM^{-1}=\begin{pmatrix}1&-1&3&-1\\-1&3&-1&1\end{pmatrix}.
\]
従って $p,q$ の両倍率は共通の符号であり、負号も実現する
（$M$ による写像はトーラス間の全射で、二行とも整数格子に属さず、その和は $NM^{-1}$ で整数になる）。
零係数があってもこの計算は同じである。
\end{proof}

\begin{proof}[対称表示での一般の分解]
対称完全交叉の同型条件は $L$ の合同であり、同値にsocle四次テンソル $S$ の共通四文字変換による射影同値である。
$G$ を交代四次収縮 $D$ の勾配、$B=S_{01\mid23}G_{01\mid23}^{t}$ とする。
$z=pq$ と置くと $D=8(2z+3)$。部分trace $\operatorname{tr}_{1}(B^3)$ は偶文字0,2および奇文字1,3の各組で一つずつの対角値を持ち、
その差は $2048z(z-1)$ である。この作用素は共通倍率と相似を受けるだけで、$D=0$ の場合にも使える。
よって $z\ne0,1$ では二平面の無順序分解が定まる。
奇側三次の二形式は $Y^3,T^3$ であり、その張る空間の純立方はちょうどこの二直線である。
偶側三次は $pX^3+3XZ^2,3X^2Z+qZ^3$。ここに $(uX+vZ)^3$ が入るには
$uv\ne0$, $u^2=pv^2$, $v^2=qu^2$ が必要で、$z\ne1$ に反する。
従って二平面の順序も決まり、奇側の二立方に掛かる偶文字から全四文字が復元される。
残る群は前の指数計算による交換と同時符号だけである。

部分traceがスカラーとなる $z=0,1$ はそれぞれ $\operatorname{tr}(B^2)/D^2=8/3,88/25$ で区別される。
以下この二つの値での合同を直接調べる。
\end{proof}

\begin{proof}[対称表示の二つの例外軌跡]
$E_0=YT$ を指定して合同を取ると、既に同じ交換と同時符号だけが残る。
実際、その二因子と共通核への $L$ の制限の二平方 $X^2,Z^2$ が決まる。
制限の核 $J=\langle YT,XT-YZ\rangle$ により、$X$ は $T$ と、$Z$ は $Y$ と組になる。
従って交換を除くと持ち上げは $X'=aX+uY+vT$, $Z'=cZ+wY+hT$, $Y'=bY,T'=dT$ である。
$X'T'-Y'Z'\in J$ から $ad=bc$, $v=w=0$、第一・第二平方の $XY,ZT$ 係数から $u=h=0$。
二平方の残る係数が $a^2=d^2$, $c^2=b^2$ を強制し、二パラメータの倍率は共通の符号になる。

$E_0$ の像の候補は階数2の式で、その核への $L$ の制限が二次元、制限写像の核の鉛筆には可逆行列がある。
最後の条件は $XT-YZ$ から従う。一般の行列は
\[
 H=\begin{pmatrix}a&0&0&s\\0&-pa-b&-s&t\\0&-s&b&0\\s&t&0&-a-qb\end{pmatrix}.
\]
$a=b=0$ なら階数2の候補は $E_0$ だけである。$a=0,b=1$ なら三次小行列式から $s=0,t^2=q$ となる。
その核への制限は左上が $a$、非対角が $s$、右下が $-(1+pt_0^2)a-2t_0^2 b+2t_0 t$ の二行列である
（$t_0$ は候補点の $t$）。$t_0\ne0$ なら制限は三次元、零なら $a=s=0$ の鉛筆全体が特異で、どちらも除かれる。
$b=0$ の場合も交換で同じである。

$ab\ne0$ ではSchur補行列から $t=0$ と
\[
 s^2=-b(pa+b)=-a(a+qb)
\]
が必要十分である。$a=1$ と置く。$s\ne0$ のとき $b^2+(p-q)b-1$ の根が単純なら、
二つの対角条件の $b,s$ に関する微分は独立である。さらに $t$ の微分を加えると核への制限は三次元である
（階数2点でSchur補行列の微分は核への制限そのもの）。$z=0$ で重根が出るのは $(0,\pm2\sqrt{-1})$ とその交換だけで、
これらはすでに同じ類である。$z=0$ の他の点には $s=0$ の候補もない。
指定した $E_0$ を保存しない同型があればその逆でも保存しないため、他の点からこの例外類への追加同一視もない。

$z=1$ で二次式が重根なら $p+1/p=0$ で $b=-p,s=0$。従って全ての $z=1$ で
$E_0$ 以外の候補は $X^2-pZ^2$ だけである。
その二因子は $X\pm\sqrt p Z$、核上の平方直線は $(\sqrt p Y\pm T)^2$ である。
一方、その制限の核を張るもう一式 $XT-YZ$ は、二因子 $X\pm\sqrt p Z$ に
$T\mp Y/\sqrt p$ を掛け合わせて得る。元の $E_0$ のところでは、
こうして得る二直線が制限の平方直線と一致する。従って新候補に移るには
$\{\pm\sqrt p\}=\{\pm1/\sqrt p\}$、つまり $p=\pm1$ が必要である。
逆同型にも同じ条件が必要で、その二点は $q=p$ で同時符号により同型である。これで全例外を尽くした。
\end{proof}

\begin{proof}[枝間の貼り合わせ]
対称表示は反対階数6で異方向族から分離される。捻りの左反対核の行列式は符号を除き $z^2t^2$ であるため、
定理\ref{thm:transverse-doublet-isomorphism}の二次錐対を持つ階数12枝とは交わらない。
残る $\{\infty,h\}$ の標準形は、全位置の文字に1を加えると、捻りを戻した表示でちょうど $S(h,h)$ になる。
平行移動の捻り符号の変化は位置別倍率なので、証明済みの無印同型により対角線だけが一致する。
従って二つの曲面はいずれも既約な二次元像を持ち、台95と94は各々の座標軸とその交点の特殊化である。
\path{computations/verify_parallel_symmetric_surface.py} が二成分の指数、切断、部分trace、例外の二次式、境界のテンソル一致を検算する。
\end{proof}

% !TeX root = ../papers/nc-p3.tex
\section{平行二重色の残部に対する四次・七次障害}
\label{sec:parallel-excluded-shapes}
\begin{theorem}
平行二重色の代表台のうち $67$--$75$, $86,87$, $91$--$93$, $97$--$111$ の29台には標準AS窓がない。
さらに
\[
 79,83,123,128,130,131,137,141
\]
の八台にもなく、
$80,81,84,119,120,121,122,125,126,127,132,133,135,136,138,140$ の16台には指定階数の係数すらない。
従って定理\ref{thm:klein-parallel-flags}, \ref{thm:parallel-symmetric-surface}と合わせ、
148台のうち、この段階で正則性に未決定部分が残るのは
\[
 82,85,88,89,90,118,124,129,134,139,142,143,144,145,146,147
\]
だけである。
\end{theorem}
\begin{proof}
定理\ref{thm:klein-parallel-flags}のSmith解で、最後の16台では残る小行列式の少なくとも一つが各位相で
Laurent単元となる。29台では同じ単元障害と、自由トーラス上での単元二列小行列式を使った四次障害に加え、
残る係数方程式の解での検査が必要となる。全ての位相を保持して以下の操作を行う。

座標は全て可逆である。可逆と分かった因子を除いた各方程式を既約因子の積とし、ある因子集合が別のものを
含むなら大きい方の方程式を落とす。複数因子の式では最初に零となる因子で分ける。一因子の式 $Ay+B=0$ があれば、
$A=B=0$ の部分と $AB\ne0$, $y=-B/A$ の部分に分ける。以前に除いた因子と代入の分母も追跡する。
よって低次元の特殊化も被覆する。検算は各段での因子分解と、末端で全ての元の式が恒等的に消えることを確かめる。
107番だけは五座標を順に $a,b,c,d,e$ と書き、
\[
 c=d(ab+1)/a,\qquad e^2=(ab+1)(bd+1)/a^2
\]
を先に解く。残る条件は $(a-1)(ab+1)d+a^2=0$ であり、全係数非零だから
$d=-a^2/((a-1)(ab+1))$, $b=(a-1)(1-a^2e^2)/a$ の一つの有理表示で全て覆える。

各末端で実際の六関係の列空間を有理関数体上で取る。始点を固定した零文字の長さ4の関係挿入行列の
余次元が9を超えれば、分母が定義されたその表示のどの点でも同じ下界を持つ。
実際、全ての列を使う元の語基底での挿入行列の小行列式が恒等的に消えることを、関数体で基底を選んで調べただけである。
正則な窓の期待文字次元はKoszul分解から、次数4から7まで順に
$(9,10,8,8),(14,14,14,14),(20,20,22,22),(30,30,30,30)$ となる。
29台は階数不足または四次障害までで全て排除される。

残り八台には同じ全係数被覆の上で、直前の経路基底に一文字を付け、末尾に二次関係を挿入することで
有理関数体上の高次商を作る。各段の文字別の自由列を新しい経路基底とし、全体を分数なしの消去で計算する。
$79,123,128$ は始点0・長さ7・文字0、$83,137,141$ は始点3・長さ7・文字1、
$130$ は始点1・長さ5・文字1、$131$ は始点2・長さ6・文字3で期待値を超える。ここでも元の語への
関係挿入は係数について正則なので、全ての特殊化まで排除される。
\path{computations/verify_parallel_excluded_shapes.py} が全位相と係数分岐に対してこれらの不等式を確認する。
\end{proof}

\begin{proposition}
定理\ref{thm:parallel-symmetric-surface}の二曲面と定理\ref{thm:klein-three-lines}の二族との間に無印の重なりはない。
\end{proposition}
\begin{proof}
反対階数が一致する候補は対称表示同士だけである。三薄色の対称表示では、
定理\ref{thm:three-lines-symmetric-isomorphism}のFourier座標において二次式の束は跡零対角網と第四行列の和になる。
網の射影平面には六つの階数2点があり、第四行列の跡は非零である。
網を跡零の形に規格化すると、その安定群は共通倍率を除き符号付き置換であるため、この跡の零・非零は網を指定した合同の不変量である。

平行曲面の行列 $H$ を前節と同じ記号で書く。階数2点は二平面 $s=0$ と $t=0$ の和に入る。
$s=0$ には $YT$ と、$a=0$ または $b=0$ に各高々二点、さらに $pq=1$ のときだけ $a=1,b=-p,t=0$ がある。
$pq=1$ ならこの平面での行列式は $ab((pa+b)^2/p-t^2)$ で、四直線の相異なる六交点となる。
この場合を含め平面内のどの直線にも三点以下である。$t=0$ のうち $s\ne0$ の点は、
$a=1$, $b^2+(p-q)b-1=0$ の高々二通りの $b$ で、各々一本の直線上に二点以下である。
この二直線は交わり、他の直線は各一本と一点でしか交わらない。従って上の二平面以外の平面には合計五点以下しかない。
$t=0$ で $s=0$ の座標軸点が二つ同時に加わるのは $p=q=0$ のときだけであり、
$pq=1$ の場合 $b=-p$ の根の $s$ は零なので $t=0$ の点は計三点以下である。

よって六点平面の候補は $pq=1$ の $s=0$ と $p=q=0$ の $t=0$ だけである。
前者は $X,\sqrt p Z,\sqrt p Y+T,\sqrt p Y-T$ の座標で跡零対角網になり、第四式 $XT-YZ$ は異なる二組をまたぐため跡零である。
後者の網は $(X,T)$ と $(Z,Y)$ の二ブロックに分かれ、行列式は $(a^2+s^2)(b^2+s^2)$。
四直線に沿って各ブロックを同時対角化し、各座標を倍率変更すれば跡零対角網になる
（四対角値の唯一の線型関係には全四係数が現れる）。第四式 $YT$ は再び組をまたいで跡零である。
従っていずれも三薄色の網には合同にならない。
\end{proof}

% !TeX root = ../papers/nc-p3.tex
\section{可逆な交差ブロックを持つ二つの疎長方形の全補完}
\label{sec:parallel-invertible-rectangles}
\begin{theorem}\label{thm:parallel-invertible-rectangles}
平行二重色の表で76--85または112--147と同じ薄い部分の台を持つ窓を取る。型をそれぞれ76型、112型と呼ぶ。
厚い部分には追加の非零条件を課さない。各型で、位置別対角変更を除いた二切断階数の解は
一つの有理的な主成分と、$t=1,-1$ 上の二つの成分で尽くされる。その係数・交わりは以下の式で与えられる。
主成分の幾何学的一般点は標準Koszul左右AS正則、反対階数16、最小周期4で、
無印の像の次元はそれぞれ3,4である。$t=\pm1$ の二つの特別成分では、後述の例外集合を除く全点が標準正則である。
\end{theorem}

\begin{proof}[薄い部分の規格化]
前節の16台のうち88--90を除いた全てがこの二形のどちらかに含まれる。
薄い色2,3の四長方形で非零成分の二次小行列式を解く。指数行列のSmith非零対角は全て1で、
位置別倍率を除くと一つの可逆な変数 $t$ が残る。検算で用いる順序の規格化は次の通りである。
薄い第一切断の二つの行列を順に、列ベクトルと行ベクトルの外積で書けば
\[
\begin{array}{c|cc}
 & M_2 & M_3\\\hline
76 &(t,0,1,1)^t(1,1,1,0)&(1,1,1,0)^t(1,0,1,1)\\
112 &(t,1,1,1)^t(1,1,1,0)&(t,1,t,0)^t(1,1,1/t,1/t).
\end{array}
\]
薄い第二切断が残る厚いブロックの右上 $U$、左下 $V$ を決定する。
各二行行列を四成分の行順のリストで書き、切断・色を $(0,0),(0,1),(1,0),(1,1)$ の順とすると
\[
\begin{array}{c|l|l}
 &U&V\\\hline
76 &(t,0,0,1);(1,t,1,0);(1,0,0,1);(0,1,1,1)
   &(1,t,1,0);(1,0,0,t);(t,1,1,0);(1,0,0,1)\\
112&(t,0,1,1);(t,t,1,0);(1,0,t,1);(0,1,1,t)
   &(1,1,1,0);(1,1,0,1);(t,1/t,1,0);(1/t,0,1,1).
\end{array}
\]
いずれも可逆であり、各厚い行列が階数2となる条件は左上 $H$ が可逆で右下が $V H^{-1}U$ となることである。
\end{proof}

\begin{proof}[行列式比による全消去]
左上の四行列は順に
\[
 (x_0,x_1,x_6,x_7),\ (x_2,x_3,x_4,x_5),\
 (x_0,x_4,x_3,x_7),\ (x_5,x_1,x_6,x_2)
\]
となる。それらを $H_0,\ldots,H_3$ とし、$\det H_0/\det H_i$ を $i=1,2,3$ の順に
$a,b,c$ と置く。第一切断の右下を $V H^{-1}U$ で定め、第二切断が同じ語に要求する右下と比較する。
共通の $\det H_0$ を払うと、$V\operatorname{adj}(H_i)U$ に順に $1,a,b,c$ を掛けた成分の比較となり、$x_0,\ldots,x_7$ に関する八本の斉次一次式が得られる。
係数行列の行列式は $(at^2-bc)^4$（112型ではさらに $t^{-4}$ 倍）である。
従って $a=bc/t^2$。$A=x_2,B=x_4,C=x_5,D=x_6$ と置く。
両型とも $x_1=bB$ で、残る三つは
\[
\begin{array}{c|ll}
 &76 &112 \\\hline
t x_0 & cA+(t^2-c)B & cA+t(t^2-c)B\\
t x_7 & cC & (t^2-c)B+cC\\
b x_3 &t^2D+(b-1)cC &
 t^2D+(b-1)cA+t(b-1)cC+t(t^2(b-2)-2bc+b+2c)B .
\end{array}
\]
ここまでに割るのはトーラス単元だけである。$\delta=t^2+bc-b-c$ と書く。
残る一次式は $\delta B=0$ と、76型なら $\delta C=0$、112型なら $\delta(A-2tB+tC)=0$ となる。
$\delta\ne0$ のとき前者は $\det H_0=0$ を与える。後者では $B=0,A=-tC$ と代入すると
三つの行列式比の条件は各々単元倍の $C^2(b-c),C^2(b-1),C^2(t^2-c)$ になる。
$C=0$ なら再び行列式が零、$C\ne0$ なら $\delta=0$ に反する。よって必ず
\[
 t^2=b+c-bc .
\]
この曲面上、三つの行列式比の条件は、共通単元と符号を除いて
$(b-c)F,c(b-1)F,b(c-1)F$ である。ただし
\[
 F_{76}=t^2BD+cC(A-B),\qquad
 F_{112}=-t^2BD+cA(B-C)+tb(1-c)B^2+tcBC .
\]
これで $tbc\det H_0\ne0$ の下、全ての係数条件は
$F=0$ または $b=c=1,\ t=\pm1$ となる。逆に三つの行列式比が正しければ他の三行列式も非零で、
右下の二表示が一致し全ての階数を得る。

薄い部分の対角安定群の作用は112型では $A,B,C,D$ の共通倍率であり、
76型では $(A,B)$ と $(C,D)$ の独立倍率である。これは薄い語の位置別指数行列の整数核で直接分かる。
112型のSmith因子2による残余符号も全厚い座標に自明に作用する。
$F=0$ 上で $B=0$ なら $AC=0$ で $\det H_0=0$ となる。従って $B=1$ にでき、$D$ は一次式から一意である。
76型では $C\ne0$ も強制されるので $C=1$ とできる。基底曲面は $b\ne1$ で $c=(t^2-b)/(1-b)$、
$b=1$ なら $t=\pm1$ で、曲面全体は既約である。
従って主成分の次元は3,4で、各特別成分は $AC-BD\ne0$ の開部分を持つ
$\mathbb P^1\times\mathbb P^1$, $\mathbb P^3$ である。それぞれの主成分との交わりは同じ $F=0$ である。
\path{computations/verify_parallel_invertible_rectangles.py} がSmith変換、八行列式、残る一次式と三つの二次式を検算する。
\end{proof}

\begin{proof}[四つの特別成分の全点]
$b=c=1,t=\pm1$ とする。$t=1$ では第四位置の文字に1を加えると偶二位置の交換と奇二位置の交換を個別に保つ。
$t=-1$ でも先に四元群捻りを戻せば同じである。実際、さらに第三位置に符号 $(1,1,-1,1)$、
第四位置の平行移動前に76型で $(1,-1,1,1)$、112型で $(-1,1,-1,1)$ を掛ければよい。
この両対称形の隣接六平面には純テンソルがない。混合色の二つの関係を係数 $u,v$ で取ると、
二つの混合二行ブロックの行列式は $u^2,v^2$ の独立な一次式であり、純関係なら $u=v=0$ となる。
残る対角二ブロックはどちらへの射影も同型であるため、片方が零なら他方も零である。

この両対称表示で、全ての非零片側評価の階数が3以上であることが標準正則性の必要十分条件である。
実際、総色1のため各側の平方四成分には少なくとも二次元の対角反対核がある。
次元3以上なら二つの対角成分を零にでき、可約な二次式を持つ。対角鉛筆が二次元なら、
四列のどれかが零か二列が比例すれば同じく可約な元を持ち、そうでなければ滑らかな四次完全交叉になる。
反対切断の混合色2,3は薄い部分だけで決まり常に階数2である。従って両側で対角核が二次元なら全反対階数は8である。
可約鉛筆では補題\ref{lem:bisymmetric-reducible-quick-test}、両側滑らかなら
定理\ref{thm:smooth-apolar-complete}を適用する。必要性は両対称窓の片側評価障害による。

実際に階数2以下の評価があるのは、$AC-BD\ne0$ の中で次に挙げる各行の条件の和集合に限り、各条件で必ず存在する。
\[
\begin{array}{c|l}
76 &(D=0,\ A=B),\quad(C=D,\ A=0)\\
112 &(B=tD,\ A=D),\quad(C=tD,\ B=tD),\quad(C=tD,\ A=0),\quad(C=0,\ A=D).
\end{array}
\]
これを確かめるには、始点0,1において厚い関係に $\binom{H}{V}$ の二列、
薄い関係にその非零列を一つ取る。$t=-1$ では捻りを戻す行符号を加える。
文字対の一方だけの評価では厚い部分が相手の一対を張り、薄い関係の少なくとも一つが他対に非零となる。
両対で非零なら、階数2以下の像は両対に同型に射影するグラフである。色0関係に限ると両対の係数射影が可逆なので、
評価座標の非零個数が各対の間で一致する。従って調べる台は各対一文字ずつの四通りと全四文字だけである。
各台で最初の座標を1にし、残る非零座標と $AC-BD$ を可逆化する。六行四列の評価行列の三次小行列式の消去根基を求めると、
表示の各一次条件対または単位イデアルを得る。評価点全体は射影的で、階数の下がる射影点の存在は基底上の閉条件である。
したがって各台の像の閉包の和集合を取れば、特殊値も含めて正確な不正則集合になる。
\path{computations/verify_parallel_vertical_regularity.py} は40個の消去根基で両方向のイデアル包含を確認する。
\end{proof}

\begin{proof}[主成分の一般点]
各型の $t=b=c=1$, $(A,B,C)=(2,1,3)$, $F=0$ の点は上の正則判定を通る
（$D$ は順に $-3,-1$）。主成分はこの点を含み、二切断・一文字切断の階数が保たれる。
定理\ref{thm:klein-elliptic-completion}で用いた次数別の開条件を適用する。
すなわち正則点の近くで縮約による三次・四次交叉の下界と上半連続性により有限Frobenius双対が延び、
各内部次数のbar複体の非対角ホモロジーの消滅は一般点へ移る。よって幾何学的一般点が全次数で標準ASとなる。

$t=2,b=3,c=-1/2,(A,B,C)=(7,1,5)$ とし $F=0$ で $D$ を選ぶ。両型で反対四ブロックが可逆、
四位置の無限小線型作用の行列が階数61である。後者の核は自明な倍率だけなので一般の射影安定群は有限である。
有限群への四元群の埋込みは有限個で、各作用で文字が異なるため、印を忘れても像の次元は減らない。
さらに四次交代収縮の勾配 $G$ と値 $D_4$ を使い、各始点 $j$ の隣接切断で $L_j=W_jG_j^t/D_4$ とする。
第一文字を残す部分traceで $\operatorname{tr}((\operatorname{tr}_{j+1} L_j^2)^2)$ はこの有理点の四始点で相異なる。
よって一般点に始点を動かす周期1,2はなく、最小周期は4である。これらの行列計算も前記の補完検算に含まれる。
\end{proof}

% !TeX root = ../papers/nc-p3.tex
\section{第三の疎長方形と三主成分の正則開集合}
\label{sec:parallel-third-rectangle}
\begin{theorem}\label{thm:parallel-third-rectangle}
前節の二形に含まれない未決定の台88--90は一つの有理三変数表示を持つ。
\[
 b\ne0,1,\quad AD\ne0,1,\quad L=AD-1,\quad M=(b-1)L,\quad C=1-(b-1)AD/b
\]
とする。薄い部分の非零成分は全て1とし、厚い部分の16成分を前節の $x_0,\ldots,x_{15}$ の順序で
\[
 \left(A,b,1,1,1,1,C,D,-1/L,-1/L,
 (AD-b)/(bL),D/L,A/L,-b/M,-1/L,-1/L\right)
\]
と取ればよい。指定した文字座標の対角倍率を除いてこの表示は一意である。
この族の一般点は標準左右AS正則、反対階数16、最小周期4で、無印の像の次元は3である。
さらにこの三変数族および前節の二つの主成分のそれぞれに、全閉点が標準ASとなる非空Zariski開集合がある。
そこでは各始点で四中心四次式上の自由階数が256である。
\end{theorem}
\begin{proof}[第一切断の特異な交差ブロック]
薄い四長方形の積条件は連結な一つのトーラスになり、残る係数は全て位置別対角作用で1にできる。
第一切断の厚い行列は
\[
 \begin{pmatrix}x_0&x_1&0&0\\x_6&x_7&1&1\\1&0&x_8&x_9\\1&0&x_{14}&x_{15}\end{pmatrix},
 \qquad
 \begin{pmatrix}x_2&x_3&1&0\\x_4&x_5&1&0\\1&1&x_{10}&x_{11}\\0&0&x_{12}&x_{13}\end{pmatrix}.
\]
第一行列の左上を $P$、第二行列の右下を $Q$ とする。全三台で $x_1 x_{13}\ne0$ である。
各々の指定ブロックを含む三次小行列式から $x_8\det P=-x_1$, $x_2\det Q=x_{13}$ となり両ブロックが可逆である。
従って他方の対角ブロックの四成分は、それぞれ一定値 $u=-x_1/\det P$, $v=x_{13}/\det Q$ である。
薄い部分の安定群は左上の八座標を同じ倍率、右下をその逆で変える。
実際その指数行列の非零Smith因子は $1^{11},2$ で、最後の符号は16座標に自明に作用する。
よって $v=1$ にできる。

$P=(A,b,C,D)$ と置く。第二切断は二つの交差ブロックが両色で可逆なので、左上
$(A,1,1,D)$ と $(1,b,C,1)$ も可逆である。その行列式を $L,M$ と書く。
前者のSchur補式は $u=-1/L$, $x_{12}=A/L$, $x_{11}=D/L$ を与える。
後者は $u=(1-b)/M$, $x_{13}=-b/M$, $x_{10}=(2-b-C)/M$ を与える。
従って $M=(b-1)L$, $C=1-(b-1)AD/b$。逆に両方の第一切断の行列式条件もこれで成立する。
両方の組に渡る残りの零成分候補は $C$ と $x_{10}$ である。
それぞれ $AD=b/(b-1)$ と $AD=b$ で消え、両方なら $b=AD=2$ となる。
どちらか一方だけが零の場合の二つの表示は台表の許した回転・反転・文字変更で同じ代表89となる。
$AD=0$ の境界の二台86,87は既に全係数で排除済みである。
\path{computations/verify_parallel_third_rectangle.py} が全係数表示と残余倍率を検算する。
\end{proof}

\begin{lemma}[中心四次式による有限長の自由性証明]\label{lem:parallel-residue-border}
体上の四頂点二次経路代数を $B$ とする。四つの中心四次式を割った商が、各始点で長さ0から13まで
$(1+z+z^2+z^3)^4$ の係数を持つとする。元の各始点の商次元が長さ14まで $\binom{n+3}{3}$ なら、
各始点は四式の生成する多項式環上で階数256の自由加群である。
一文字階数4と隣接階数6の最上テンソルが与えられていれば、標準Koszul左右AS性が従う。
\end{lemma}
\begin{proof}
中心商は長さ13で零だから、その後も零である。各始点で商の斉次基底の語を長さ12まで持ち上げ、
これらを形式的な基底とする $P=k[z_1,z_2,z_3,z_4]$ 上の自由加群 $F$ を作る（各変数の重み4）。
中心性と次数帰納で自然な写像 $F\to e_i B$ は全射である。次元仮定から長さ14までは同型である。
各基底に右から矢を掛けたものを、長さ13までのこの同型で $F$ に戻すことで、$P$ 線型な矢作用を定義する。
二次関係を基底に作用させると長さ14以下であり、$e_iB$ に写して零だから $F$ 上でも零となる。
従って右作用は $B$ を経由する。空語に基底語を作用させると長さ12以下で対応する形式基底が得られ、
空語に第 $j$ 中心式を作用させると $z_j$ 倍の空語となる。右作用の $P$ 線型性により任意の多項式倍の基底が空語から得られる。
よってこの右作用が先の全射の逆を与える。多項式変数も代数的に独立である。
最後は定理\ref{thm:klein-rational-period-four}の証明と同じである。二つの互いに素な中心変数により
単純から環への $\operatorname{Hom},\operatorname{Ext}^1$ が消え、両側の表示の双対完全性とEuler等式で中央項も完全となる。
\end{proof}

\begin{proof}[三主成分での開集合]
第三の長方形では $(A,D,b)=(1,2,2),(2,3,5)$ を取る。有理数体上で中心四次式を求め、素数1009で還元する。
両点で四中心式の商が四始点で補題の長さ13までの値を持ち、還元後の元の商も長さ14まで標準次元を持つ。
隣接関係の基底、四つの一文字縮約、最上テンソルはこの素数で所定の階数を持つ。
従ってまず剰余体そのもので補題が適用できる。

この証明から標数零へ移るとき、剰余体上の余核次元を標数零での下界として用いる必要はない。
$\mathbb Z_{(1009)}$ 上で単元小行列式により縮約空間の自由基底を選び、テンソルとその縮約による
四項の線型Koszul列を作る。連続微分の積は定義から零である。長さ $n-1$ までの経路加群が自由なら、
内部次数 $n$ のこの列で最後の商以外は自由である。剰余体上の完全性により、始めから単元小行列式で
分裂単射を順に除ける。最後の余核は表示により長さ $n$ の経路加群で、再び自由となる。
初期次数から帰納して全次数で標数零の正則性を得る。左右両向きで同じ議論ができる。
そこで有理中心四式の商の次元上界と正則Hilbert級数を比較すれば標数零でも各始点が自由である。
この二点を含む三変数族の幾何学的一般点は前節の次数別の開条件により正則である。
$(2,3,5)$ で反対切断は可逆、四位置の無限小作用の階数は61、
前節の部分trace二乗の跡は四始点で相異なり、無印次元と最小周期が従う。

前節の二つの主成分では、同節の $t=2,b=3,(A,B,C)=(7,1,5)$ の有理点でそれぞれ四次中心の文字別次元が
$(3,2,0,0)$、四式の商の長さ13までの次元上界が補題の値となる。こちらは既に一般点が正則なので、
半連続性による下界と中心による自由加群の上界から各有理点自体が中心四変数上自由で標準正則となる。

以上の四有理点それぞれで、頂点で零となる内部次数4の導分を六つの二次関係へのLeibniz則で調べると、
224列に対する階数下界は文字順に $188,184,192,192$ である。これは内導分と中心四次式のEuler倍でちょうど尽くされる値である。
定理\ref{thm:full-component-central-open}の持ち上げ論法をそれぞれの既約基底で適用する。
すなわち、正則点での完備局所族の次数別平坦性は上と同じ単元分裂による。
次数4の中心の基底を非零因子で取れば命題\ref{prop:euler-central-lifting}が適用できる。
完備化の忠実平坦性と五次までの定階数性から、中心の階数5の束を一つの開近傍に得る。
長さ13までの中心商の上界は有限個の開条件で保たれ、一般点のHilbert下界から開集合の全点で
四変数上の自由性と標準AS性が従う。これらの主成分の閉包の任意の標準点にも、
系\ref{cor:klein-pi-specializations}と同じ議論で標準恒等式 $s_{512}=0$ が特殊化する。
\path{computations/verify_parallel_residue_border.py} が商経路、中心、導分の計算を再現する。
\end{proof}

% !TeX root = ../papers/nc-p3.tex
\section{第三長方形の二零点条件が定める曲線の完全分類}
\label{sec:parallel-double-zero-curve}
\begin{theorem}\label{thm:parallel-double-zero-curve}
第三長方形で二つの可変成分 $C,x_{10}$ がともに消える台88の全点が標準Koszul左右AS正則で、
各始点は四つの中心四次式上で階数256の自由加群となる。$A=v,D=2/v,b=2$ とすると、無印の完全不変量は
$X=v^2\in k^*$ である。一位置移動した窓はこの指定始点の曲線上のどの窓とも同型でない。
二位置移動は $X\mapsto4/X$ であり、最小周期は $X=\pm2$ のとき2、他では4である。
反対階数は $2v^2\pm5v+4=0$ の四つのテンソル値で15、他は16である。
\end{theorem}
\begin{proof}[無印不変量と文字直線]
四次交代収縮の値は全曲線で $-16$ である。各始点 $p$ での規格化した収縮
$L_p=W_{p,p+1\mid p+2,p+3}(G/D_4)_{p,p+1\mid p+2,p+3}^t$ の二乗を取り、
後の文字だけを縮約した作用素を $T_p$、前の文字だけを縮約した作用素を $U_{p+1}$ とする。
$T_0$ の跡は $-157/16$ で、$T_1$ の跡は $(v^2+4/v^2+316)/16$ である。
後者だけで同型候補は $v'\in\{v,-v,2/v,-2/v\}$ に限られる。前者の二乗跡は
\[
 \operatorname{tr}(T_0(v)^2)-\operatorname{tr}(T_0(2/v)^2)
 =\tfrac{55}{1024}(v^2-4/v^2).
\]
従って元から二乗が等しい場合だけが残る。$v\mapsto-v$ は四位置の文字平行移動量 $(3,2,2,3)$ と対角倍率で実現する。
二位置移動と $v\mapsto2/v$ の合成は全位置で文字0,1を交換した後の対角倍率で実現する。
実際、共通台の語の係数比を取り、位置別倍率の指数行列の左整数核の全基底について指標値が1である。

一位置移動については、$T_p,U_p$ の四組の同時固有値を使える。
いずれの位置でも、異なる二文字での二つの値の差の分子は互いに素なので四同時固有直線を全点で無印復元できる。
元の台の任意の位置別文字置換で一位置移動後の台になるものはない。
これは初めの三位置の $24^3$ 通りを取り、総色零から第四位置でなければならない文字を決める有限検査である。
従って一位置移動でこの曲線に戻ることはない。
反対四ブロックの行列式は $-(2v^2-5v+4)/v,(2v^2+5v+4)/v,-1,-1$ となり、四根は相異なり単純である。
\end{proof}

\begin{proof}[関数体の証明を外れた全根の処理]
前節でこの曲線の点 $v=1$ が正則であるため、曲線の幾何学的一般点も正則で、
各始点・次数のHilbert関数は全特殊点で標準値以上である。有理関数体 $\mathbb Q(v)$ 上で
五次までの経路基底と四次中心を計算する。文字0では三つの自由列で係数 $(1,3,5),(2,7,11)$ の二式、
文字1では二自由列の基底を取り、これらの中心商を四始点で長さ13まで計算すると
$(1+z+z^2+z^3)^4$ の係数を得る。中心四式で割った後は、基底語に一矢を付けて末尾二次関係を課し、
長さ4で始点の中心四式を加える。さらに長い語での中心式は始点に移せるので同じ商である。

この計算が使える開集合を厳密に記録する。経路基底・中心の核・中心商を選ぶ全行列で分母を払い、
行の内容を割る際にはその非定数因子を記録する。主元および後退消去後の各主元の非定数因子も記録する。
一列の消去で二係数の最大公約数を割るのは主元が可逆な局所化での変形である。
従って記録した因子が全て非零なら、全ての基底・次元・中心式がそのまま特殊化する。
標準Hilbert下界と自由加群による上界が一致し、ここでは全次数で正則である。

記録される既約因子は1885個である。ただし正則性は $v\mapsto-v,2/v$ の両方で保たれるため、
その軌道全体が例外集合に入る根だけを調べれば十分である。零を除くと40個の既約多項式、計64根になる。
各多項式 $f$ について先頭係数が可逆な素数 $l$ と非零の一次剰余根を選ぶ。
五次までの経路次元はその剰余根でも標準で、四次中心の文字別次元は $(3,2,0,0)$ である。
剰余体上でその核から四式を取り、全四始点で同じ長さ13までの中心商次元を検算する。

この剰余計算は $f$ の標数零の根での証明になる。実際、数体でその剰余根の上の素点を取り、
経路空間の有限次数の行列を比較する。長さ5まで、曲線の一般点からの下界と剰余体からの上界が等しいため、
各次数は素点で自由である。この基底で四次中心の式を立てると、標数零の核次元も曲線の一般値以上であり、
剰余体の核次元と等しくなる。よって中心核はその素点で直和因子で、指定の剰余四式を中心式として持ち上げられる。
中心商には剰余体からの次元上界が移り、元の全次数のHilbert下界と比較すれば自由性と正則性が従う。
同じ既約多項式の全ての共役根に適用できる。
\path{computations/verify_parallel_double_zero_curve.py} と
\path{computations/verify_parallel_curve_residues.py} が開集合と全残余根を再生成する。
\end{proof}

% !TeX root = ../papers/nc-p3.tex
\section{第三長方形の全点の無印同型と周期}
\label{sec:third-rectangle-orbits}
\begin{theorem}\label{thm:third-rectangle-orbits}
第三長方形全体で $v=A,h=AD$ と置く。指定始点での完全不変量は
\[
 (X,Y,h)=\bigl(v^2,(1-2/b)/v,h\bigr),\qquad Xh(h-1)(1-XY^2)\ne0
\]
であり、この開集合の各点は一つの窓の同型類を与える。
二位置移動は $(X,Y,h)\mapsto(h^2/X,XY/h,h)$ である。
最小周期は $X=h$ または $(X=-h,Y=0)$ のとき2、他は4である。
一零条件の台89の閉包は $(h-2)^2=h^2XY^2$、その二零条件の台88は $h=2,Y=0$ である。
台89は既約で、その閉包は台88に沿って局所エタール位相で二枚が横断する。
\end{theorem}
\begin{proof}[全係数での四座標枠の復元]
$G$ を $W$ の四次交代収縮の勾配とする。$D_4(W)$ では割らずに使う。
$W_{pq}$ を第 $p,q$ 位置を出力とする切断とし、
$T_{pq}=\operatorname{tr}_q((W_{pq}G_{pq}^t)^2)$ とする。
これは基底変更で共役と共通の指標倍を受ける。四元群作用により対角である。
さらにある位置で対角な共変作用素を挿入して $W$ と $G$ を他の位置で縮約できる。
入力側の作用素が一つの場合、位置 $p$ の第 $i$ 対角成分は
$\sum_{k_p=i}W_kG_k T(k_q)$ である。既に復元した位置では、固有直線に沿う射影子も入力に使える。

位置番号を $0,1,2,3$ とする。まず位置1で $T_{10},T_{12},T_{13}$ と、位置0の $T_{01}$ を挿入した縮約を取る。
次に位置3で、位置1の四射影を挿入した四式と $T_{30},T_{31}$ を使う。
最後に位置0,2で、位置1,3の四射影の組を挿入した16式を取る。
各段階、四文字のどの一対でも同時固有値の全差は単位イデアルを生成する。
これは次の小さい有理係数環計算で確認できる。各切断を総色で四ブロックに分け、
$G_k$ を各位置で $k_j$ を除いた三文字の交代和で作る。
$U=vh(h-1)b(b-1)$ とし、$W,G$ の共通因子を割る際には全て $U$ の因子か非零定数であることを確かめる。
各文字対で差の分子を順に加え、既に得たGröbner基底による剰余とその基底自体から $U$ の因子を除く。
この操作で $\mathbb Q[v,h,b,U^{-1}]$ での生成イデアルは変わらない。
位置 $1,3,0,2$ の順にそれぞれ六対とも基底 $\{1\}$ を得る。
従って全幾何点で同時固有直線が復元でき、任意の無印同型が四位置別の単項行列になる。
\path{computations/verify_third_rectangle_frames.py} は差の生成から24回の単位判定までを行う。
\end{proof}

\begin{proof}[台の比較と商]
始点移動後の台と目標の台が同じ大きさのとき、第一から第三位置までの置換を全て選ぶ。
各語の総色零から第四位置の像が強制され、これが置換で実際に目標の台を与える場合を残す。
非零の薄い語は全て薄い語に移る。薄い係数を1に保つ残余対角作用は厚い最初の八成分に一つの倍率、
後の八成分にその逆倍率を掛けるので、$x_2=1$ で再規格化できる。
零条件なしと二零条件では、元の始点で恒等と
\[
 (v,h,b)\longmapsto(-v,h,b/(b-1))
\]
の二つだけを得る。一零条件では恒等が同じ零位置を、後の変換が反対の零位置を与える。
一・三位置移動では台を移す置換がなく、二位置移動では上の結果に $v\mapsto h/v$ を合成する。
\path{computations/verify_third_rectangle_orbits.py} が台ごとの全置換と係数の再規格化を検算する。

$z=1-2/b$ は $z\ne\pm1$ を動き、上の対合は $(v,z)\mapsto(-v,-z)$ となる。
従って $X=v^2,Y=z/v$ が商の座標を与え、二位置移動の固定条件から周期式を得る。
$h=b$ と $h=b/(b-1)$ の二面は交換され、それぞれ既約で $h=b=2$ で横断する。
二重被覆は $v\ne0$ でエタールであり、商に降ろすと定理の方程式になる。
\end{proof}

\begin{corollary}
$r=v+h/v$ と置く。反対階数は16から
$b(r-2)=h-1$ と $b(r+h+1)=h-1$ のうち成立する式の数を引いたものである。
階数14は $h=-3$ かつ $b(v-3)(v+1)=-4v$ の曲線で尽きる。
台88の全点は三変数族全体における標準AS正則開集合に含まれ、特に台89の一般点も正則である。
\end{corollary}
\begin{proof}
反対切断の色2,3の行列式はともに $-1$ である。色0,1には定数単元の三次小行列式があり、
行列式は順に $-[b(r-2)-h+1]/((b-1)(h-1))$ と $[b(r+h+1)-h+1]/(h-1)$ となる。
後二式の同時消滅では $b\ne0$ から $h=-3$ を得る。
台88では前節の全点検算により五次まで標準であり、必要な二文字の四次中心核は次元3,2を持ち、
その中に長さ13で商を零にする四式がある。軌道を移すときもこの二文字は保たれる。
これらの次元は三変数族の一般値である。各点の近傍で五次までの経路と二つの核は局所自由だから、
この四式を局所的に延長できる。中心商の有限個の次元上界と一般点からの全次数の下界に
定理\ref{thm:parallel-third-rectangle}の証明を適用する。
\end{proof}

% !TeX root = ../papers/nc-p3.tex
\section{可逆交差長方形の七つの台での正則性の決定}
\label{sec:parallel-support-fibers}
\begin{theorem}\label{thm:parallel-support-fibers}
平行台の代表82,118,124,129,134,139,142で正則になる窓は全て特別成分 $b=c=1,\ t=\pm1$ に属する。
各成分における正則点は次の表で尽きる。112型で $B=1,a=A/t,y=C,d=D/t$ とし、表の自由な文字には
指定した台の他の係数が非零という条件を課す。
\[
\begin{array}{c|l|l}
 台&\text{特別成分内の等式}&\text{正則点}\\\hline
82&A=B&\text{全点}\\
118&a=y=1,\ d=0&\text{全点}\\
124&a=y=1&\text{全点}\\
129&a=1,\ y=0&d\ne1\\
134&a=1,\ d=0&\text{全点}\\
139&a=1&d\ne1\\
142&d=0,\ a+y=1&\text{全点}
\end{array}
\]
台85,145,147の特別成分上の点も全て正則であり、台143,144,146は特別成分と交わらない。
\end{theorem}
\begin{proof}[全厚い係数の零条件]
特別成分の112型で16成分の零条件は順に
\[
 a,1,a,d,1,y,d,y,\ell,a-1,\ell,y-1,a-1,1,y-1,1,\quad
 \ell=a+y-d-1
\]
の消滅と同値で、$ay-d\ne0$ である。76型でも順に
\[
 A,B,A,D,B,C,D,C,C-D,A-B,C-D,C,A-B,B,C,B
\]
となる。各々、右下八成分の共通分母と単元符号を除いただけである。
従って定理\ref{thm:parallel-invertible-rectangles}の四つの例外直線対を各台に制限すると表の最終列を得る。
例えば112型の対象の台では $a\ne0$ なので障害は
$(a=d=1)$、$(y=d=1)$、$(y=0,a=d)$ の三条件に簡約される。

主成分では $B=1$ とする。76型では $C=1$ で、$x_3$ と $x_9$ は同時に、$x_6$ と $x_{12}$ も同時に消える。
従って $x_9,x_{12}$ だけが零の台82は主成分にない。
112型では $a=A/t,y=C,m=t^2=b+c-bc$ と置く。$n=-c(a-b)(y-b)$ とすれば、
各成分の零条件は次の16多項式の消滅とそれぞれ同値である。
\[
\begin{gathered}
 ca+b(1-c),\ 1,\ a,\ n+b^2,\ 1,\ y,\ c(a-1)(y-1)-m,\ cy+b(1-c),\\
 ay-b,\ a-b,\ n+b(1+b-a-y),\ y-1,\ a-1,\ 1,\ y-b,\ 1 .
\end{gathered}
\]
実際、各分子を $t^2=m$ で簡約すれば表示式との比は単元であり、分母には
$tbc$ と行列式因子 $ay(b+c)-bc(a+y)+b^2(c-1)$ だけを要する。

台118,124,129,134,139は $x_9=x_{12}=0$ を含むため、$a=b=1,t^2=1$ である。
すると $D=t$ で、左上の行列式条件は $y\ne1$ となる。これで118,124,134は主成分上で不可能である。
129では $y=0$ と $cy+1-c=0$ も要るので $c=1$ となる。
残る139については、より強く $a=b=1$ の二つの主成分切片の全点が不正則である。
実際、$t=\pm1$ それぞれで、$\mathbb Q(c,y)$ 上の始点1の長さ7の四文字の次元は
$(34,33,34,33)$、合計134である。標準値は120なので上半連続性によりあらゆる特殊値でも排除される。
計算は各次数の自由語基底に一文字を付加し、末尾での六関係の行空間を割る有理関数体上の掃き出しである。

台142では $x_8=x_{10}=0$ かつ $x_{11}x_{14}\ne0$ より、$ay=b$ および $c=1$ を得る。
実際、$a=b/y$ のとき $x_{10}$ を検出する多項式は $b(y-b)(y-1)(c-1)/y$ である。
$x_3=x_6=0$ よりさらに $a+y=1=b$ となり、特別成分上に戻る。
\end{proof}

\begin{corollary}
主成分で台143は $t=\pm1$, $c=1$, $a=1-y$, $b=y(1-y)$ の二曲線の開部分、
145は $t=\pm1$, $c=1$, $a=b/y$ の二曲面の開部分である。
144は $a+y=1$, $c=b^2/(ay-b+b^2)$ の切片、146は
$c=b(1+b-a-y)/((a-b)(y-b))$ の切片である。いずれも元の台の非零条件と $t^2=b+c-bc$ を課す。
両形の全パラメータで $t$ の符号変更は四元群の二コサイクル捻りと位置別対角変更で実現される。
112型では同時に $A,D$ の符号を変え、76型では $A,B,C,D$ を固定する。
\end{corollary}
\begin{proof}
零条件表で $x_{10}$ と $x_3$ の多項式の差は $-b(a+y-1)$ である。
これと上の $ay=b$ での式を使えば各切片の表示を得る。
捻りについては、符号変更前後の一般16成分と薄い成分の比が全て定数符号になる。
各語でこの比に四元群二コサイクルの符号を掛けると、位置別倍率の重み行列の全左整数核で指標値が1である。
従って一組の定数対角倍率で全パラメータに共通の同値を与え、零成分を持つ場合にも特殊化する。
\path{computations/verify_parallel_support_fibers.py} が係数の単元、台、両符号の長さ7と捻り指標を検算する。
\end{proof}

% !TeX root = ../papers/nc-p3.tex
\section{台143の二曲線の全係数とその周期一点}
\label{sec:parallel-septic-curve}
\begin{theorem}\label{thm:parallel-septic-curve}
台143の全窓と、その112型主成分内の台142への交わりは
\[
 t=\pm1,\quad y\ne0,1,\quad b=y(1-y),\quad c=1,\quad
 (A,B,C,D)=\bigl(t(1-y),1,y,t(y^2-y+1)\bigr)
\]
で尽くされる。各符号で標準Koszul左右AS正則性の必要十分条件は $y^2-y-1\ne0$ である。
指定始点の無印同型類は符号と $y$ で決まり、各符号で $y^2-y+1=0$ の二根だけが互いに同型である。
その一類は最小周期1、$y=1/2$ では2、他の正則点では4である。
二位置移動は $y\mapsto1-y$ を誘導する。
\end{theorem}
\begin{proof}[不正則点と中心式の特殊化]
前節の係数に代入すると $y^2-y+1=0$ でだけ $x_6$ が消えて台142になる。
$t=1$ の係数を使い、$y^2-y-1=0$ では $\mathbb Q[y]/(y^2-y-1)$ 上で四次関係空間を直接計算する。
各総色の64語に、三つの位置の六関係を全て配置して張る行を作る。数体の各係数を基底 $1,y$ での
有理二行行列に置き換えれば、商の四次元数は順に $(10,10,8,8)$ で総計36になる。
従って両根とも正則性を満たさない。捻りでこの障害は両符号に共通である。

他の全点には定理\ref{thm:parallel-double-zero-curve}の局所化と剰余体の論法を使う。
$\mathbb Q(y)$ で正符号の四次中心の文字0,1に三自由列・二自由列を得る。
同定理と同じ係数の四式を選び、四始点で長さ13までの中心商次元は $(1+z+z^2+z^3)^4$ の係数になる。
全ての消去に用いた単元の既約因子は2542個である。下に証明する $y\mapsto1-y$ の移動により
両方の表示で除かれる因子は66個、総次数169である。そのうち $y,y-1,y^2-y-1$ を除く63個それぞれで、
先頭係数を保つ素数の一次剰余根を取り、五次までの標準経路次元、四次中心の次元 $(3,2,0,0)$、
全四始点での中心商を検算する。既約な主成分の一般点が標準なので、そこから特殊化される全次数の
Hilbert下界がある。数体での核の持上げと中心商の上界から全ての共役根で中央四変数上の自由性と正則性が従う。

負符号への捻りは前節で全点共通に構成した。同じHilbert級数とKoszul左右AS性を保つ。
正符号では復元した四中心四次式は文字 $0,0,1,1$ である。二位置移動の文字置換も全位置で0と1の交換なので、
全長4の総色0,1を保つ。後の二式を平方すれば四式とも総色零となり、捻っても中心に残る。
よって負符号でも重み $4,4,8,8$ の多項式環上で各始点は自由階数1024になる。
\end{proof}

\begin{proof}[添字を保つ同型と始点移動]
第三長方形と同じ交代勾配を使う。位置1,3で両側の隣接部分trace二乗の固有値差の分子は、
両符号、全六文字対で互いに素である。これら二位置で得た固有射影を両方挿入して残る位置を縮約すれば、
位置0,2でも各文字対の16個の差の分子は単位イデアルを生成する。
従って全曲線で四つの枠が内在的である。
台143では始点を固定する台置換は恒等だけで、薄いトーラスの規格化 $B=1$ から $y=C$ は一意になる。
奇位置移動では台の置換は存在しない。二位置移動では各位置で0と1を交換するものだけがあり、
$1-y$ の正規係数との比は対角作用の重み行列の全左整数核で指標1となる。
台142の二根では、同じ有限検査を $\mathbb Q[y]/(y^2-y+1)$ で行うと二根を交換する同型と、一位置移動する同型をともに得る。
二つの符号の反対階数は下の通り互いに異なる。これで指定始点の同型と周期式を得る。
\end{proof}

\begin{corollary}
$u=y^2-y$ とする。正符号の反対階数は $u=-1$ で8、$u=1$ で9、他で10である。
負符号の反対階数は $u^3-u^2+3u+1=0$ で15、他で16である。
両符号で、正則な台143の各点は112型主成分全体の正則開集合に属する。
\end{corollary}
\begin{proof}
正符号の反対切断の色0,2,3は全て定階数2である。色1では四次小行列式と三次小行列式の最大公約数が
単元を除き $(u-1)(u+1)^2$ と $u+1$ で、二次小行列式は共通零点を持たない。
負符号では色0,2,3が可逆、色1の行列式分子が上記三次式で、三次小行列式の最大公約数は1である。
開近傍については正符号で、三変数族の二零条件の曲線と同様、経路基底と文字0,1の四次中心核が
一般階数を持つので局所的に四式を延長して中心商の次元上界を保てる。
負符号では112型全体の符号を交換する捻りでこの開近傍を移す。
\path{computations/verify_parallel_septic_curve.py} が二つの数体の計算、共通固有枠、台置換、
因子の局所化と全残余因子での持上げを再現する。
\end{proof}

% !TeX root = ../papers/nc-p3.tex
\section{台145の曲面の商と四次中心の束}
\label{sec:parallel-rectangle-surface}
\begin{theorem}\label{thm:parallel-rectangle-surface}
$t=\pm1,\ c=1,\ b=ay,\ A=ta,\ B=1,\ C=y,\ D=t(a+y-ay)$、
$a,y\ne0,1$ とする。これは台145の水平面と、その中の一零条件・二零条件を含む。
各符号で指定始点の同型は $(a,y)$ と $(1/a,1/y)$ の同一視で尽き、
異なる符号の同型はない。従って $s=(a+1)/(a-1),\ z=(y+1)/(y-1)$ とすると無印商は
\[
 P=s^2,\quad Q=sz,\quad R=z^2,\qquad PR=Q^2,\quad (P-1)(R-1)\ne0 .
\]
二位置移動は $P,R$ の交換である。最小周期2以下は $P=R$、すなわち $a=y$ または $ay=1$ で、
このうち $ay=a+y=1$ の一類だけが周期1である。残る全点の最小周期は4である。
\end{theorem}
\begin{proof}
二つの厚い零成分を固定した後の追加の零は $a+y-1=0$ と $a+y-ay=0$ に限られる。
二面を交換するのが同時逆数変換である。その各曲線と交点は前節のものに一致する。
四枠の復元に定理\ref{thm:third-rectangle-orbits}の部分traceと射影子を同じ順序で使う。
$\mathbb Q[a,y,1/(ay(a-1)(y-1))]$ で、位置1の文字対 $(2,3)$ の差が生成するイデアルは
$(a+1,y+1)$、他の対に用いる差は全て単位イデアルである。従って
$(-1,-1)$ の外で全四枠が復元できる。この一点は位置1で同時固有直線にならない二文字を持つので、
指定始点で他の点と同型にならない。

追加の零がないときは各始点移動で三位置分の置換を列挙し、総色から決まる残りの置換を検査する。
偶位置移動で台を保つ置換が二つ、奇位置移動では零である。各置換について薄い語は薄い語へ移り、
その対角規格化と $B=1$ は一意に厚い係数を定める。両符号で係数比の全トーラス指標を取ると、
元の始点で恒等と同時逆数変換、二移動でこれらに $a,y$ の交換を合成したものとなる。
これらの同型は台が縮む場合にも成り立つ。実際、置換・薄い語の倍率と $B$ の再規格化に使う倍率は
この二曲線上でも単元であり、一零条件上の必要性と二零条件での追加の同型は前節で求めた。
同時逆数変換は $(s,z)\mapsto(-s,-z)$ となり、商の式が従う。

固定点 $(-1,-1)$ でも二位置移動の同型は存在する。一位置移動は、
交代収縮の値が非零であるこの点で、位置0,1の隣接部分trace二乗の固有値を比較すると排除される。
正符号の多重集合は $\{-280,-280,-272,-192\}$ と $\{-456,-104,-232,-232\}$、
負符号は $\{-211,-211,-190,-160\}$ と $\{-313,-109,-175,-175\}$ である。
いずれも同じ共通の交代収縮の二乗で規格化できる。符号間の区別は次の反対階数から従う。
\end{proof}

\begin{proposition}\label{prop:parallel-surface-ranks-centers}
正符号の反対階数は $ay=1$ で8、$ay=-1$ で9、他で10である。
負符号では16から $a+y=0$ および
\[
 b^3-2b^2(a+y)+3b^2+3b-2(a+y)+1=0,\qquad b=ay,
\]
のうち成り立つ条件の数だけ引く。$ay=-1$ の全点が両符号で標準ASから外れる。
その補集合全体では五次までの経路束は標準階数の自由加群であり、正符号の四次中心の
文字別の核は階数 $(3,2,0,0)$ の自由加群となり全点で基底が特殊化する。
\end{proposition}
\begin{proof}
正符号では反対切断の色0,2,3がそれぞれ定階数2で、色1の行列式は
$(b-1)^2(b+1)/(ay(a-1)(y-1))$ である。色1の三次小行列式は共通因子 $b-1$ を除けば単位イデアルを生成し、
二次小行列式も単位イデアルを生成する。負符号では色2,3が可逆で、色0,1の三次小行列式がそれぞれ単位イデアルを生成する。
行列式の分子を取れば表示の条件を得る。正符号で第四文字のラベルを1だけ平行移動すると
奇二位置の交換対称性も全点で直接確認できる。

$y=-1/a$ とした関数体上の長さ4の次元は $(10,10,8,8)$ である。特殊化と捻りでこの障害を保つ。
逆に $ay+1$ も逆にした二変数環上で、五次までの全経路基底と四次中心の全核を求める。
各行列で、各段階の行と列の中から局所化の単元である成分だけを主元に選ぶ。
分母を払う際、行の内容を除く際、後退消去後の正規化の際にも単元以外では割らない。
経路の各総次数の次元は $(1,4,10,20,35,56)$、中心の自由列数は各文字で上記の値になる。
従って関数体の階数計算そのものがこの開集合全体の自由基底を与えている。
\path{computations/verify_parallel_rectangle_surface.py} および
\path{computations/localized_unit_echelon.py} がこれらの検算を行う。
\end{proof}

% !TeX root = ../papers/nc-p3.tex
\section{台145の両曲面の全点での正則性}
\label{sec:parallel-rectangle-surface-regularity}
\begin{theorem}\label{thm:parallel-rectangle-surface-regularity}
前節の両曲面の全幾何点で、標準Koszul左右AS正則性の必要十分条件は $ay\ne-1$ である。
正符号の全正則点は四中心四次式上で各始点が階数256の自由加群になり、
負符号の全正則点では重み $4,4,8,8$ の中心多項式環上で自由階数1024である。
従って前節の二次錐の商から $ay=-1$ の像だけを除けば、ここでの正則代数の無印同型類と全周期層が得られる。
\end{theorem}
\begin{proof}[二変数環上の四つの開集合]
正符号で $ay(a-1)(y-1)(ay+1)$ を逆にする。前節の単元主元だけを用いる経路・中心基底を使い、
文字0の自由列の最初と最後、文字1の両自由列に対応する四中心式を選ぶ。
各始点で四式を割った商を長さ13まで計算すると所定の
$(1+z+z^2+z^3)^4$ の次元になる。ここでは各始点で全ての消去に要した単元以外の既約因子を記録する。
行の内容のうち元の局所化で既に単元である部分だけを除き、主元はまず定数、次に単元、
どちらもない段階で新しい因子を選ぶ。新しい主元で分母や行操作に現れる全ての因子を記録する。
始点を $0,1,2,3$ の順に取ると、開集合を除外する既約因子数は $64,127,106,109$ である。
\path{research/rectangular_surface_cut_factors.json} が四つの完全なリストであり、
\path{computations/verify_rectangular_surface_cut_charts.py} が中心式から再計算する。

始点 $p=0,1$ の各々について、このままの開集合、同時逆数変換で移したもの、
さらに $p+2$ の開集合に $a,y$ の交換を施したものとその同時逆数像を使う。
どれかで証明されればその始点での自由性を得る。$p+2$ も同じ判定を移せる。
全てから外れる点は四つの因子の和集合の共通部分にある。以下では、この共通部分の曲線と孤立点を個別に処理する。
\end{proof}

\begin{proof}[共通因子の全曲線]
$p=0$ の四リストに共通なのは
\[
 a+1,\quad y+1,\quad ay-1,\quad ay-a+1,\quad ay-y+1,\quad ay^2+1
\]
である。$p=1$ ではこの六つに $a-y,ay-2a+1,ay-2y+1,ay^2-1$ を加えた十個になる。
$ay=1$ は定理\ref{thm:parallel-support-fibers}の特別成分に含まれる。この曲線でも四中心四次式上の自由性を得るため、以下で中心商を直接扱う。

これらを二位置移動と同時逆数でまとめ、$v=y$ とすると次の一変数の七行で尽きる。
各行で前記の全曲面と同じ六関係を特殊化し、定理\ref{thm:parallel-double-zero-curve}と同じ有理関数体の中心商を計算する。
表の三数は選んだ開集合の除外因子数、軌道の両表示で除外される非零根の因子数、その総次数である。
\[
\begin{array}{c|c|rrr}
a&v\text{ の移動}& & &\\\hline
-1&1/v&1911&67&136\\
v&1/v&2234&48&126\\
1/v&1/v&1988&47&88\\
1/(1-v)&1-v&2341&68&200\\
1/(2-v)&2-v&2425&58&124\\
1/v^2&1/v&2427&46&124\\
-1/v^2&1/v&2429&39&98
\end{array}
\]
定義域外と既知の $ay=-1$ を除き、各因子の一次剰余根で五次までの経路、中心核、
全四始点の長さ13までの中心商を検算する。一般点からのHilbert下界と剰余体からの上界により全ての共役根に通用する。
\path{computations/verify_rectangular_surface_curve_strata.py} が各行を再現する。
\end{proof}

\begin{proof}[孤立点の消去と全剰余値]
四リストの一つの因子 $f$ を取る。$f$ を含まないリストだけがこの曲線をさらに切る。
縦直線なら変数を交換し、以後 $f(x,Y)$ は $Y$ に正次数とする。
各残りのリスト $T$ について $g\in T$ ごとに $\operatorname{Res}_Y(f,g)$ を取り、平方因子と $x,x-1$ を除く。
異なるリストの多項式を順に最大公約数で合わせる。残る既約な横座標多項式 $A(x)$ と元の $f$ の組は
$p=0$ で721組、$p=1$ で1070組、横座標の次数和はそれぞれ2592,3938である。
これらは射影で得た上界で、余分な組を含んでもよい。最大公約数が1であることを素数で先に検査する際は、
両先頭係数が保たれる場合だけその結論を採用する。

$A$ の根を $x=\alpha$ とする。$\alpha=-1$ は既に扱った。$f(\alpha,Y)$ の次数を求め、定義域と既知の不正則因子から
除く値 $0,1,-1/\alpha$ の重複度を、$Y$ 微分と $A$ による剰余で標数零のまま求める。
$A$ の先頭係数が単元で単純な一次剰余根 $x_0\ne0,1,-1$ があり、
$f(x_0,Y)$ の次数、これら三根の重複度が全て保たれる素点を選ぶ。
そのとき残るどの縦座標もこの素点の上で整であり、禁止された三値以外に還元される。
残った剰余多項式を各 $T$ の $\prod_{g\in T}g(x_0,Y)$ と順に最大公約数で合わせる。
定数になればその素点に還元される例外はない。定数でない場合には、全ての残余根が一次で、
各根の同時逆数・二位置移動を含む表示のいずれかが、五次までの経路・中心と全四始点の中心商の検査に合格する素点を用いる。これらの変換は許容した素点で整なので数体の表示も同じ変換で移せる。
従って全ての $\alpha$ とその上にある全ての幾何点を数体での持上げが覆う。
\path{computations/verify_rectangular_surface_residual_points.py} が終結式から最後の剰余点まで計算する。

各始点で成功する四式の選び方は中心の同じ文字別空間内の非空開集合を与える。二位置移動と同時逆数で四元群の文字が移っても、唯一の次元2の非零文字は1のままである。
有限個の開集合を交差させて四始点に共通の四式を選べる。中心商の上界と元のHilbert下界で
全次数の自由性・正則性が従う。負符号は前節の全曲面の捻りを用いればよい。
\end{proof}

% !TeX root = ../papers/nc-p3.tex
\section{第三長方形の全正則性と無印の四元群}
\label{sec:third-rectangle-regularity}
\begin{theorem}\label{thm:third-rectangle-all-regular}
第三長方形の台88,89,90の全点が標準Koszul左右AS正則である。
各始点は四つの中心四次式上の自由階数256の加群となる。
従って定理\ref{thm:third-rectangle-orbits}の三次元商、その周期層・反対階数層・台の特殊化は
全て正則代数の添字保持同型類を与える。
\end{theorem}
\begin{proof}
$a=A,h=AD,b=x_1$ とし、$ahb(h-1)(b-1)$ を逆にする。この三変数環上で経路を長さ5まで計算し、
文字0,1の四次中心核を求める。局所化の単元主元だけで自由列数3,2を得る。
文字0の最初と最後の自由列、文字1の両自由列の四式を選ぶ。
\begin{enumerate}
\item 始点2では、台の表示をそのまま使うと四式の商の長さ13までの計算にも追加の逆元が一つも要らない。
従って全点で $(1+z+z^2+z^3)^4$ の次元列を得る。二位置移動 $a\mapsto h/a$ により始点0も全点で従う。
\item 始点1では文字を位置順に $(2,2,0,0)$ だけ四元群で平行移動し、同じ自由列選択法で四式を取り直す。
中心商の計算に要る新しい既約因子は
\[
 F=abh+ab-6bh+4h
\]
だけである。指定始点を変えない同型 $(a,b)\mapsto(-a,b/(b-1))$ でこの表示を移すと、
必要な非零条件は単元倍を除いて $F'=-abh-ab-2bh-4h$ になる。
$F+F'=-8bh$ は可逆なので両方が失敗する点はない。始点3も二位置移動で従う。
\end{enumerate}
\path{computations/verify_third_rectangle_regularity.py} が二つの表示の素因子集合を空集合と $\{F\}$ にまで再計算する。
一般点から元の代数には全次数の標準Hilbert下界がある。各始点で成功する中心式を一般に選び直し、
四始点の有限個の開条件を同時に満たす四式を取れる。中心商による自由加群の上界との比較で自由性を得る。
元のテンソルの階数と二つの中心正則元から、補題\ref{lem:parallel-residue-border}と同じ
左右Koszul列の完全性が従う。
\end{proof}

\begin{proposition}\label{prop:rectangular-intrinsic-klein}
四文字空間に射影的に作用して窓の直線を保つ群は、第三長方形の全点、および
台145の両水平面の $a=y=-1$ 以外の全点で、標準の四元群そのものである。
ここには両面の台143の曲線と台142への交わりを含む。
特にこれらの窓に別の四元群の印を選んで異なる文字多重度の枝に移すことはできない。
\end{proposition}
\begin{proof}
既に復元した四枠により各位置で単項行列である。第三長方形で恒等以外の置換は正規表示を $a\mapsto-a$ に移すので、
指定始点で同じ係数を保てない。水平面でも一般の台で追加の置換が係数を保つには同時逆数変換の固定点が必要である。
一零条件の同じ向きの台を保つ置換は恒等だけである。二零条件では $y^2-y+1$ で割った体で、
台の置換のうち係数比の全対角指標が1となるものを数えると各符号とも恒等一つになる。

対角群は台の位置別重み行列から決まる。台88,89,90,142,143,145のいずれも非零Smith因子は
$1^{11},2,2$ である。各位置のスカラーを除けば有限部の位数は4になり、
四元群の二つの独立な文字がこの全てを実現する。これは射影文字空間上の主張であり、
この群に忠実に作用する別の四元群も同じ部分群に一致する。
よってその固有文字線は共通で、印の変更では群の自己同型と各位置の文字平行移動だけが残る。
\path{computations/verify_rectangle_klein_stabilizers.py} が台ごとの有限部と周期1の二境界での全置換を検算する。
\end{proof}

% !TeX root = ../papers/nc-p3.tex
\section{台144の楕円K3曲面と水平三層の無印商}
\label{sec:parallel-horizontal-moduli}
\begin{theorem}\label{thm:parallel-horizontal-moduli}
112型の主成分で、台144,146,147の幾何学的一般点は標準Koszul左右AS正則で、最小周期は4である。
指定始点を固定した無印同型類の関数体は順に次の通りである。
\begin{enumerate}
\item 台144ではPicard数20、Néron--Severi判別式の絶対値8の楕円K3曲面の関数体である。
\item 台146では $k(t,b,a)$、台147では四変数の純超越拡大である。
\end{enumerate}
ここで曲面・三次元多様体は指定した台の層の双有理模型を表す。
\end{theorem}

\begin{proof}[曲面と三次元層]
前節までの $B=1$, $a=A/t$, $y=C$, $m=t^2=b+c-bc$ を使う。
台144では $y=1-a$ かつ $c=b^2/E$, $E=(b-a)(b-1+a)$ である。
$t^2=b a(1-a)/E$ の右辺は $b=0$ で奇数位数を持つから二重被覆は既約である。
$W=tE$ と置くと $W^2=b a(1-a)E$ であり、$b$ の関数体上で
\[
 k_0=-b^2(1-b),\qquad X=k_0/a-k_0,\qquad Z=k_0 W/a^2
\]
により $Z^2=X(X+L)(X+M)$ となる。ただし $L=b(1-b)^2$, $M=b^3$ である。
逆写像は $a=k_0/(X+k_0)$ から得られる。判別式 $16L^2M^2(L-M)^2$ と
$c_4=16(L^2-LM+M^2)$ の位数を調べると
\[
\begin{array}{c|rrrr}
 b&0&1&1/2&\infty\\\hline
 (v(c_4),v(\Delta))&(2,10)&(0,4)&(0,2)&(2,8)\\
 \text{Kodaira型}&I_4^*&I_4&I_2&I_2^*
\end{array}
\]
となる。無限遠では $X/b^4,Z/b^6$ の整モデルを使い、全ての点で $v(\Delta)<12$ なので極小である。
他のファイバーは滑らかで、基本直線束は $\mathcal O_{\mathbf P^1}(2)$。
切断付き楕円面の標準束公式と $R^1f_*\mathcal O=\mathcal O(-2)$ から極小滑らか模型はK3である。
可約ファイバーの根格子 $D_8,A_3,A_1,D_6$ から自明格子の階数は20となり、
標数零のK3のPicard数の上界およびShioda--Tate公式からMordell--Weil階数は0である。
加法的ファイバーの連結な滑らかな部分と還元核には標数零で有限位数点がないため、
捩れ部分群は $I_4^*$ の位数4の成分群に単射する。三つの非自明な有理2捩れ点を既に持つので位数は4であり、
判別式の絶対値は $4\cdot4\cdot2\cdot4/4^2=8$ となる。

台146では $t^2(a-b)(y-b)=b(a-1)(y-1)$ である。$y$ についての係数と定数項は
$k[t,b,a]$ で互いに素であり、全体が既約で、
\[
 y=\frac{b\big(t^2(a-b)-(a-1)\big)}{t^2(a-b)-b(a-1)} .
\]
従って $b=1,t^2=1$ における余分な自由度から別の成分を生じることもない。
元の各台では薄い係数・行列式・指定した厚い係数の非零条件を課す。
\end{proof}

\begin{proof}[文字枠と無印同型]
三層の閉包はいずれも台143上の $t=1,y=2,a=-1,b=-2,c=1$ を含む。
この点は正則で、同じ二切断の階数を持つ。定理\ref{thm:parallel-invertible-rectangles}の次数別の特殊化論証で
有限Frobenius双対および各内部次数でのKoszul性を幾何学的一般点に移せる。
この点では位置1,3の二つの部分traceおよびその固有線を挿入した位置0,2の縮約が四枠を復元する。
各々に要る固有値差の非消滅は開条件なので、三層の一般点でも枠は内在的である。

台144,146で指定始点の台を保つ四置換は恒等のみで、薄い成分の規格化と $B=1$ により全係数が一意である。
奇数移動の台置換はなく、二位置移動の台置換は全位置で $(0\,1)$ の一つだけである。
係数には $(a,y)\mapsto(y,a)$ と作用する。このため台144で二位置移動が $a\mapsto1-a$ になる一方、
$t\mapsto-t$ は別の一般同型類を与える。台146でも二位置移動は一般点を固定しない。

台147は $b\ne1$ で $c=(m-b)/(1-b)$ と書ける。指定始点の全作用は恒等と
\[
 I(t,b,c,a,y)=(t,m/b,m/c,1-c+cy/b,1-c+ca/b)
\]
である。一位置移動の代表は
\[
 J(t,b,c,a,y)=(1/t,c/m,1/b,c(1-a)/m,c(1-y)/m).
\]
$J$ は $I$ と可換で、$J^2$ は $(a,y)$ の交換と $I$ の合成である。
これらを確かめるには台を移す置換を全探索し、薄い係数の $t$ の指数について位置別倍率を消去する。
残る倍率で $B=1$ に戻した後、16個の厚い係数の全てが表示先の補完式に一致する。
一般には $t\ne1/t$ かつ $(a,y)$ の交換が $\{1,I\}$ の同じ軌道に入らず、周期は4である。

$u=(b-t)/(b+t)$ と置くと $I(u)=-u$ なので、$A_0=(a+I(a))/2$,
$A_1=(a-I(a))/(2u)$ は不変である。$b=t(1+u)/(1-u)$, $a=A_0+u A_1$ および
$y=(A_0-u A_1-1+c)b/c$ で全て復元される。従って不変体は $k(t,u^2,A_0,A_1)$ である。
\path{computations/verify_parallel_horizontal_moduli.py} はファイバーの位数、消去式、台の全置換と全係数の変換を検算する。
\end{proof}

% !TeX root = ../papers/nc-p3.tex
\section{第一長方形の特別成分の無印商と共通の四次障害}
\label{sec:first-rectangle-orbits}
\begin{theorem}\label{thm:first-rectangle-orbits}
76型の $b=c=1,t=\pm1$ において、各符号で台85の正則点の添字保持同型類の商は
\[
 \{(z,h)\in\mathbb A^2:z\ne0,\ h\ne(z-1)^2,(z+1)^2\}
\]
である。台82では各符号で $d\in\mathbb A^1\setminus\{0,1\}$ が完全不変量である。
これらの四層の各点の最小周期は2で、両符号の間や両台の間に同型はない。
76型の水平主成分の指定始点の一般無印商は有理三次元で、正規化座標 $(t,b,x)$ は一般同型類を一意に与える。
\end{theorem}
\begin{proof}
台85では $B=C=1$, $x=A$, $d=D$ として $xd(x-1)(d-1)(x-d)\ne0$ となる。
台82では $x=1,d\ne0,1$ である。両方の全点の正則性は定理\ref{thm:parallel-support-fibers}で示した。

まず四文字枠の復元を行う。特別成分で $x+d\ne1$ なら位置1の四線を二次モーメントと挿入付き縮約で分離し、
次に位置3を、最後に位置0,2を復元する。台85の二変数環で各準備段階の差のイデアルは
位置1の文字対 $(0,3),(1,2)$ のみ $(x+d-1)$ で、他はすべて単位イデアルとなる。
台82の一変数環では全て単位イデアルである。

$x+d=1$ のとき、各位置 $p$ に対して異なる位置 $q$ を三通り取り、$p,q$ で切断した
$W_{pq}G_{pq}^t$ の平方の部分traceを取る。$G$ は四次交代収縮の勾配である。
各対角列は $i\leftrightarrow i\mathbin\oplus e_p$, $e=(2,3,2,3)$ で一定で、文字0と1での
三差の分子の最大公約数は定数である。これで各位置の二平面 $E_{p,j}=\langle j,j\oplus e_p\rangle$,
$j=0,1$ を内在的に区別できる。これらの平面への射影で四重テンソルを切り、位置02と13で平坦化する。
$j$ の四つ組が $(0,0,1,0)$ と $(1,1,0,1)$ のブロックは、それぞれ行と列の非零な二座標が
両添字とも異なる位置にあり、その二行二列が可逆である。従って像および双対像の射影直線と
階数一のSegre曲面の交点はそれぞれ座標の二点である。全ての平面の両文字を得る。
ブロックを特定するのに文字名の事前の選択は要らず、復元した平面の全組を取ればよい。
他のブロックも同じ文字対にのみ支持され、階数二のものはいずれも同じ座標線を与える。
同じ論証は水平主成分の $(t,b,x)=(2,3,7)$ でも通る。後述の平行移動対称性により二平面内の
重複固有値は水平主成分上で恒等的なので、この計算は一般点にも適用できる。
\path{computations/verify_first_rectangle_frames.py} が両符号のイデアルと二次元切断の行列式を検算する。

従って全ての同型は単項である。特別成分の台85で指定始点を保つ台置換は二つあり、
係数では恒等と $(x,d)\mapsto(1-d,1-x)$ となる。一位置移動は $(x,d)\mapsto(d,x)$ である。
$z=x-d$, $s=x+d-1$ を使うと指定始点の作用は $(z,s)\mapsto(z,-s)$ であり、$h=s^2$ で
上の商を得る。一位置移動で $z\mapsto-z$ となるので全点の周期は2である。
台82では偶数移動ごとに一つの台置換があり $d$ を保ち、奇数移動の台置換はない。
どちらも $t$ の符号を変える指定始点の台置換はない。

水平主成分では $x=A$, $B=C=1$, $D=c(1-x)/m$, $m=t^2=b+c-bc$ と置く。
台85の指定始点の二置換はともに正規化係数を変えない。一位置移動とその平方は
\[
 J(t,b,c,x)=(1/t,c/m,1/b,c(1-x)/m),\qquad
 J^2(t,b,c,x)=(t,m/b,m/c,1-c+cx/b)
\]
となる。$J^4=1$ である。これらは分母の許す全点で同型を実現する。
枠を復元できる開集合ではこれが全同型である。行列の右下は薄い係数と左上から一意なので、
薄い指数の規格化と左上八係数の比較で十分である。
\path{computations/verify_first_rectangle_orbits.py} が全置換を検算する。
\end{proof}

\begin{proposition}
水平一般点、特別成分の台85の $x+d=1$ 上では、窓の直線の射影安定群は $(\mathbb Z/2)^3$ である。
そのうち四位置の全てで可換な線型持上げを持つ位数4の部分群は標準の四元群一つだけである。
台85の他の特別成分点と台82の特別成分点では射影安定群は標準の四元群である。
\end{proposition}
\begin{proof}
両台とも位置別重み行列の非零Smith因子は $1^{11},2,2$ で、射影対角群 $K_0$ は位数4である。
追加の置換が係数を保つ場合、その平行移動量は各位置で $e=(2,3,2,3)$ である。
その単項持上げ $P$ は文字符号と射影的に可換であり、$P^2\in K_0$ の対角値は各 $e_p$ の対で一定となる。
$2,3$ が四元群を生成するので、全対で一定な文字は自明である。従って射影群は初等可換2群となる。
他の位数4の部分群は $K_0$ と非自明な文字 $\chi$ を共有する。対角外の元とその文字の線型交換子は
位置 $p$ で $\chi(e_p)$ であり、全位置で1にはならない。この交換子は持上げのスカラーを変えても変わらない。
\end{proof}

\begin{proposition}\label{prop:parallel-antidiagonal-obstruction}
76型・112型いずれの主成分でも $b+c=0$ の全点は標準AS正則性を失う。
特に台143における $y^2-y-1=0$ と台145における $ay=-1$ は同じ四次障害の切片である。
\end{proposition}
\begin{proof}
$c=-b$ なら $t=\pm b$ である。各符号で76型は自由な $b,x$ を、
112型は自由な $b,a,y$ を持ち、補完式の定義域は既約である。始点0で一般係数の
長さ4の四文字の経路次元を計算すると、四つの一般点でいずれも $(10,10,8,8)$ である。
その和36は標準値35より大きく、二次関係が一定階数の範囲での上半連続性により特殊点でも排除される。
\path{computations/verify_rectangle_antidiagonal_obstruction.py} が四有理関数体で計算する。
76型の上記の枠復元開集合で $J^2$ が係数を固定するには $b^2=c^2=m$ が必要で、
$c=b$ なら $b=c=1$、$c=-b$ ならこの障害を受ける。従ってこの開集合の正則点は
特別成分との交線では周期2で、その外では周期4である。
\end{proof}

% !TeX root = ../papers/nc-p3.tex
\section{第二長方形の特別成分での無印商と一次元安定群}
\label{sec:parallel-vertical-orbits}
\begin{theorem}\label{thm:parallel-vertical-orbits}
112型の $b=c=1,t=\pm1$ における正規化した八代表台で、全ての正則点の指定始点の同型と最小周期は次表で尽くされる。
$B=1$, $(a,y,d)=(A/t,C,D/t)$ とし、各行で台と行列式の非零条件および
定理\ref{thm:parallel-support-fibers}の正則条件を課す。$\sigma=a+y$, $\pi=ay$ と書く。
\[
\begin{array}{c|c|c|c}
 台&\text{係数の等式}&\text{完全不変量（各符号で）}&\text{最小周期}\\\hline
118&a=y=1,\ d=0&\text{一点}&2\\
124&a=y=1&d&2\\
129&a=1,\ y=0&d&1\\
134&a=1,\ d=0&y&2\\
139&a=1&(y,d)&2\\
142&\sigma=1,\ d=0&\pi&1\\
145&d=\sigma-1&(\sigma,\pi)&2\\
147&\text{開条件のみ}&(\sigma,\pi,d)&1\ (\sigma=1),\ 2\ (\sigma\ne1).
\end{array}
\]
異なる行、異なる符号の間に指定始点の同型はない。台147の商は例えば
$\pi(1-\sigma+\pi)d(\sigma-d-1)(\pi-d)\ne0$ の開集合である。
$\sigma^2=4\pi$ 上では $t=1$ の射影安定群が位数8の二面体群になり、$t=-1$ では
$\mathbb G_m\rtimes\mathbb Z/2$（逆数作用）になる。$a\ne y$ では両符号で標準の四元群である。
\end{theorem}

\begin{proof}[固有空間の分割の有限検査]
各位置で基底に関して対角な内在的作用素の同時固有空間分解を作る。
まず各 $p\ne q$ について $W_{pq}G_{pq}^t$ の平方の部分traceを使う。
さらに他の三位置の既知の固有空間を一つずつ選び、その射影を挿入した $W,G$ の縮約を加える。
位置 $p$ の文字 $i$ での最後の対角値は、残る三文字が指定の三ブロックに入る語での
$W_kG_k$ の和である。固有値の等しい文字を一組のブロックとし、分割が細かくなる限り繰り返す。
各操作はテンソルの基底変更に関して自然で、相対不変量の共通倍率は分割に影響しない。

二位置回転は偶数位置で恒等、奇数位置で四元群の1を加える置換と対角倍率で元の窓に戻る。
従って最終分割は位置0,1の分割 $A,B$ から $(A,B,A,B\oplus1)$ と書ける。
$a\ne y$ では目標は四つの一元集合である。$a=y$ では後述の対称性で偶数位置の $\{0,1\}$、
奇数位置の $\{2,3\}$ は同じブロックに入り、残り二文字を一つずつ分離するのが目標である。
この目標より粗い全候補を取り、そのブロック内の隣り合う二文字について、上の部分traceの差と
全三ブロックへの射影を挿入した差を同時に零とする。これは安定分割であるための必要条件である。
開条件の積、行列式、$a\ne y$ の層では $a-y$ を逆にした環で、これらの差のイデアルは
両符号・全台・全候補で単位イデアルになる。初期の二分割の候補数は非対角で $15^2-1=224$、
対角で $5^2-1=24$ に過ぎない。検算は単元因子を除いてGröbner基底を計算し、
必要な場合は定義域の積で飽和する。\path{computations/verify_rectangle_partition_frames.py} が全条件を生成する。
\end{proof}

\begin{proof}[対角の二文字と二つの安定群]
$a=y$ で四つ組の隣接位置 $0,1$ にそれぞれ二次元固有空間を残し、位置2,3では一元の固有線に射影する。
得られる二行二列行列は $t=1$ なら恒等と交換行列の張る鉛筆で、$t=-1$ ならこれに固定した符号行列を掛けた鉛筆となる。
位置2,3に二次元空間を残しても同様である。いずれも階数一の二点から各二次元空間の和・差の二直線が復元できる。
二次元空間内で和と差を基底とした表示では全同型が単項となり、元の一元固有線の対と復元した二線の対は区別される。

この新表示で指定始点の全台置換は恒等および全位置の和・差の同時交換のみである。
後者は元の四元群の文字符号ですでに実現される。各語の係数を定数と既約因子の積で表示し、
対角倍率の重み行列の左整数核を係数の指数に作用させる。得た格子のHermite標準形は、
正次元の各層で全て単位行列である。例えば台147での因子は
$a,d,a-1,2a-d-1,a^2-d$ であり、台145では $a,a-1,2a-1$ になる。
従って係数を一意に復元できる。台および符号が異なる場合の新表示の台置換は存在しない。

$t=1$ で新表示の重み行列の非零Smith因子は $1^{11},2,2$。対角群の位数4と同時交換から
安定群の位数は8となる。元の表示の標準四元群と二文字交換は非可換で、位数8の非可換群で
四元群を含むものは二面体群である。新表示の対角四元群は元のものと異なり、各位置で一つの固有文字が重なる。
$t=-1$ ではSmith因子は $1^{12}$ で、位置スカラーを除く対角部分は一次元トーラスとなる。
二次元空間の和・差の二線の重みは、各位置のスカラーを除くと $1,-1$ または $-1,1$ で、
元の一元線の重みは両方0である。元の文字符号による同時交換はこのトーラスを反転する。
この場合の全四元部分群は、一つのトーラス内対合と外の対合で生成され、トーラス共役で互いに移る。
\end{proof}

\begin{proof}[指定始点と移動]
$a\ne y$ で元の四枠を復元した場合、台147を保つ二置換は恒等と $I(a,y,d)=(y,a,d)$ である。
薄い係数と $B=1$ の規格化で同型を比較できる。一位置移動の代表は
\[
 J(a,y,d)=(1-a,1-y,1-a-y+d).
\]
$J^2=1$, $IJ=JI$ である。各小さい台での台置換は全て台147の対応する置換の制限である。
異なる台への置換はなく、指定始点で $I$ が残るのは118,124,142,145,147のみである。
二位置回転は常に許され、奇数回転を許すのは129,142,147である。
その場合に同じ類に戻る条件は $\sigma=1$ で、このとき $IJ$ が係数を固定する。
$a=y$ の新表示でも奇数回転が可能な対角の台は142,147だけで、上の変換が実現する同型と全置換数が一致する。
非対角と対角は復元した分割の次元で区別される。
元の枠が四線に分かれる各台の対角Smith因子は $1^{11},2,2$ であり、$a\ne y$ では恒等以外の置換は係数を固定しない。
\path{computations/verify_parallel_vertical_orbits.py} が薄い規格化、全軌道の台置換、和・差の行列鉛筆、両種の格子を検算する。
\end{proof}

% !TeX root = ../papers/nc-p3.tex
\section{第一長方形の水平全候補での文字枠とトーラス座標}
\label{sec:first-rectangle-horizontal-frames}
\begin{theorem}\label{thm:first-rectangle-horizontal-frames}
76型の水平台85で四次の次元が標準値であるための必要十分条件は $b+c\ne0$ である。
この全域において無印の指定始点の同型類は、次のトーラス開集合の各点で一意に表される。
\[
 (u,v,z)\in\mathbb G_m^3,\qquad
 (u+v)(1+uv)(z+v)(1-uz)(z-u)(1+vz)\ne0 .
\]
一位置移動は $J(u,v,z)=(v,1/u,-z)$。周期2の点は $u=v=1$ と $u=v=-1$ の二直線上に限られ、他の全点は周期4である。
全域で射影安定群は $(\mathbb Z/2)^3$、その中で全位置に可換な線型持上げを持つ四元部分群は一意である。
従ってこれらの同型・周期の記述は標準AS点を全て含む。さらに長さ5までの各経路空間は標準次元を持ち、
四次中心の文字別次元は $(3,2,0,0)$ で一定である。
\end{theorem}

\begin{proof}[四平面対を全域で復元する]
$b\ne1$ で $c=(t^2-b)/(1-b)$ と置く。定理\ref{thm:first-rectangle-orbits}の
四位置の対称性 $i\mapsto i\oplus e_p$, $e=(2,3,2,3)$ により、前節の部分traceと
固有空間の射影を挿入した縮約が与える分割は各位置で二対か一ブロックである。
目標より粗い候補は15個ある。少なくとも一位置に二対がある14候補に対し、各ブロックの隣接文字の
対角値の差を集め、台の非零因子および $b+c$ で飽和する。14個とも単位イデアルである。

全位置が一ブロックの場合には12個の部分traceの二次モーメントだけで足りる。
各位置で文字0と1の差の分子から単元を除くと、残りは全て $x=A$ の二次以下である。
係数を $1,x,x^2$ の順に置いた十二行三列行列の全三次小行列式を求める。
$\mathbb Q[b,t]$ で $bt(b-1)(t-b)(t+b)(t^2-b)$ を逆にすると、これらは単位イデアルを生成する。
同時消滅点があれば非零列 $(1,x,x^2)^t$ が核に入ることに矛盾する。
ここでは $b+c=(t^2-b^2)/(1-b)$ を使った。
決まった平面の四つ組 $(0,0,1,0),(1,1,0,1)$ に射影した反対切断には、いずれも
非零な二行二列に定数行列 $\left(\begin{smallmatrix}0&1\\1&0\end{smallmatrix}\right)$ がある。
前節と同じSegre交叉で四文字枠が復元される。
$b=1,c\ne1$ なら $t^2=1$ で、一位置移動により $b'=c\ne1$ に移る。
$b=c=1$ は定理\ref{thm:first-rectangle-orbits}の交線の論証で覆われる。
従って単項性に基づく指定始点の係数の一意性と安定群の結論は、ここでの全候補で成り立つ。
\path{computations/verify_first_rectangle_horizontal_frames.py} が14分割の飽和と十二行の行列の階数を検算する。
\end{proof}

\begin{proof}[基底曲面の一枚の座標]
$u=b/t,v=c/t$ とすると $t(1+uv)=u+v$ である。$b+c\ne0$ なら分母も非零である。
$z=((b+c)x-bc)/(bt)$ と置けば
\[
 t=\frac{u+v}{1+uv},\quad b=ut,\quad c=vt,\quad x=\frac{u(z+v)}{1+uv},\quad
 D=\frac{v(1-uz)}{u+v}.
\]
全非零係数の分子・分母の因子をまとめるとちょうど定理の開集合になる。
一位置移動の係数式に代入して $J$ を得る。$J^2(u,v,z)=(1/u,1/v,z)$ であり、
$J^2$ の固定条件と $u+v\ne0$ から $u=v=\pm1$。$z\ne0$ より奇数移動は固定点を持たない。
この二直線は $b=c=1,t=\pm1$ の特別成分との交線である。

二次関係の各色の像の基底は局所化上の単元小行列式で取れる。
$b\ne1$ の一枚および $b=1,t=\pm1$ の二枚で、それらを経路に順次挿入し、局所化の単元だけで掃き出す。
長さ0から5の各始点の階数は $1,4,10,20,35,56$ となる。同じ基底の四次・五次間の可換方程式の核の階数は
四文字順に $3,2,0,0$。追加の分母は要らず、これらの核と経路の余核は各ファイバーへの特殊化と可換である。
$b+c=0$ の排除は命題\ref{prop:parallel-antidiagonal-obstruction}による。
\end{proof}

\begin{corollary}
112型の水平主成分も $b+c\ne0$ では同じ $u,v$ を持ち、
$a=u(z+v)/(1+uv),\ y=u(w+v)/(1+uv)$ と書ける。左上行列式の非零条件は $zw\ne1$ である。
台147の一般無印商における指定始点の対合および一位置移動は
\[
 I(u,v,z,w)=(1/u,1/v,w,z),\qquad J(u,v,z,w)=(v,1/u,-z,-w).
\]
\end{corollary}
\begin{proof}
左上行列式を検出する $(b+c)ay-bc(a+y)+b^2(c-1)$ は
$u^2(u+v)^2(zw-1)/(1+uv)^3$ となる。残りは定理\ref{thm:parallel-horizontal-moduli}への代入である。
\path{computations/verify_horizontal_torus_coordinates.py} が両形の座標式と76型の全ファイバーに通用する低次の掃き出しを検算する。
\end{proof}

% !TeX root = ../papers/nc-p3.tex
\section{第一長方形の反対階数と階数十の水平面の有限判定}
\label{sec:first-rectangle-rank-strata}
\begin{theorem}\label{thm:first-rectangle-rank-strata}
定理\ref{thm:first-rectangle-horizontal-frames}のトーラス開集合上で、反対階数は $8,10,14,15,16$ を取る。
$C=(u+v)(1+uv)$ と置き
\[
 F=u(v+1)^2z-C,\qquad G=v(u+1)^2z+C
\]
とする。$u=v=1$ で階数8、$u,v$ のちょうど一つが1なら階数10である。
両方とも1と異なるところでは階数は $16-\mathbf1_{F=0}-\mathbf1_{G=0}$ となる。
階数15の閉包は二つの有理曲面で、その共通部分の階数14の曲線は有理かつ滑らかである。
立ち入る定義域を上のトーラス開集合とすると、この曲線の閉包は階数10,8に達しない。
階数10の二面は階数8の直線で交わる。
\end{theorem}
\begin{proof}
反対切断の四色の四行列を取る。色0の行列式から単元を除くと $(v-1)^2F$、色1では $(u-1)^2G$、
色2,3では $(u-1)^2(v-1)^2$ である。三次の全小行列式のイデアルは四つの順に
$(v-1),(u-1),((u-1)(v-1)),((u-1)(v-1))$ で、各場合に一つの小行列式が生成元の単元倍となる。
全ての行列に可逆な二次小行列式がある。$u=1$ では $F=(v+1)^2(z-1)$ が非零であり、
$v=1$ では $G=(u+1)^2(z+1)$ が非零である。これで全ての階数を得る。

$F,G$ の各因子は $z$ について一次でそれぞれ有理グラフを与える。
$F=G=0$ なら $C+4uv=0$, $z=-4v/(v+1)^2$ である。
ここで $u,v\ne-1$ であり、$r=(u+1)/(v+1)$ と置けば
\[
 v=\frac{1-r}{r(r+1)},\qquad u=\frac{r(r-1)}{r+1}.
\]
逆写像も正則で、$r=0,\pm1$ はこの交叉にはない。従って曲線は有理直線の開部分である。
$u+1/u+v+1/v+4=0$ の勾配はこの部分では消えず、二グラフの交叉は横断的である。
$u=1$ または $v=1$ の先の非零条件が階数14の両面への特殊化を排除する。
一方 $F=0$ は $v=1$ の階数10面を $z=(u+1)^2/(4u)$ で切り、
$G=0$ は $u=1$ の面を $z=-(v+1)^2/(4v)$ で切る。
二面の交線 $u=v=1$ ではどちらも非零である。
\path{computations/verify_first_rectangle_rank_strata.py} が全小行列式と交線のパラメータを検算する。
\end{proof}

\begin{theorem}\label{thm:first-rectangle-edge-center}
階数10の両面では、折り畳み環は五つの中心四次式が生成する可換環上有限である。
従って各点について以下は標準Koszul左右AS正則性の必要十分条件である。
四次中心の基底を $s_1,\ldots,s_5$ とする。代数的に独立な $\lambda_1,\ldots,\lambda_4$ を加え、
四式 $s_j-\lambda_j s_5$ で割った各始点の長さ0から13の次元が $(1+T+T^2+T^3)^4$ の係数と一致することを検査する。
$u=-1$ または $v=-1$ の二面でも、符号を戻して同じ有限判定が全点に適用できる。
\end{theorem}
\begin{proof}
$u=1$ なら $t=b=1$, $c=v$, $x=(z+v)/(1+v)$ である。$v=1$ は既に正則と決定した交線であり、
階数10の面の基底環では $v z(v-1)(v+1)(v+z)(z-1)(vz+1)$ を逆にする。
全四次中心は文字0の三式と文字1の二式である。各位置の文字を $(0,2,2,0)$ だけ平行移動した表示で、
五式を全て割った商の経路を長さ9まで計算する。各始点で得られる開集合上の次元列は
\[
 (1,4,10,20,30,36,34,20,1,0)
\]
である。基底の決定と五次までの中心の持上げは単元だけで行え、以後の経路計算で必要になる追加の
既約因子の個数は始点順に $3,16,24,3$ となる。二位置移動は $(v,z)\mapsto(1/v,z)$ であり、
この像の四つの開集合も合わせる。

これら八集合のうち各点で少なくとも一つが使える。実際、始点0と3の因子から一つずつを選んで九つのイデアルを作る。
残る六リストの因子積を順に加え、上記の基底環で飽和すると九つとも単位イデアルになる。
ある始点 $p$ で五式の商の長さ9が零なら、他のどの始点からの経路も高々三歩で $p$ を通るので、
全商の長さ12は零である。従って元の環は五つの中心式上有限である。
\path{computations/verify_first_rectangle_edge_center.py} が八リストと九つの飽和を再計算する。
$v=1$ 側には一位置移動で移せる。

正則なら四次中心が生成する可換環の次元は4で、換算次数1の一般四式は斉次Noether正規化を与える。
定理\ref{thm:central-window-cutoff}の必要性から四式の商の次元列を得る。
逆に表示の列なら、中心商を持ち上げた自由加群の各次数が標準値である。
水平主成分の正則一般点からのHilbert下界と比較すると各始点で自由となり、
補題\ref{lem:parallel-residue-border}の最後の議論により全次数の左右AS正則性が従う。
負の二面は $t$ の符号変更の二コサイクル捻りであり、$u=v=-1$ の既知の交点も含めて判定できる。
\end{proof}

\begin{corollary}
76型の水平正則候補が少なくとも片側対称な表示を持つための必要十分条件は $t=1$ である。
階数10の二面ではそれぞれ一方の偶奇だけが対称化できる。
112型でも $u=1$ で偶位置交換、$v=1$ で奇位置交換を持つ表示が得られる。
\end{corollary}
\begin{proof}
$t=1$ は $u=1$ または $v=1$ と同値である。$b=1$ の補完は位置0,2の交換を保ち、
$c=1$ の補完は第四位置であらかじめ文字に1を加えれば位置1,3の交換を保つ。
これは両形の全係数での恒等式である。逆に片側対称なら反対切断は十次元の対称二乗を経由する。
76型の他の点の階数は14以上である。両側対称なら二位置移動で同型になるので、
周期4の階数10の二面では両側同時の対称化は起きない。
\end{proof}

% !TeX root = ../papers/nc-p3.tex
\section{大きい五つの台に共通する四次障害と階数十の正則曲線}
\label{sec:horizontal-rectangle-low-fibers}
\begin{theorem}\label{thm:horizontal-rectangle-low-fibers}
112型の水平代表台143,144,145,146,147で、各始点の長さ4が標準値35を持つための必要十分条件は
$b+c\ne0$ である。そのとき全始点の長さ5も標準値56を持つ。
\end{theorem}
\begin{proof}
$b+c=0$ は命題\ref{prop:parallel-antidiagonal-obstruction}で排除した。
残りでは前節までの $u,v,z,w$ を使う。指定した五代表で共通に可逆なのは
\[
 M=uv(u+v)(1+uv)(zw-1)\prod_{q\in\{z,w\}}(q+v)(q-u)(uq-1)(vq+1)
\]
である。以下の論証ではこれ以外の厚い係数の非零を使わない。
最初に $b\ne1$ とし $t,b,x=a,y$ を使って $c=(t^2-b)/(1-b)$ を消去する。
各位置各色の関係フレームには単元小行列式がある。五次までの経路を作ったときに必要な追加の因子は
\[
\begin{split}
 L={}&b^2t^2-b^2xy+b^2x+b^2y^2-b^2+b t^2xy-b t^2x-2bt^2y\\
 &{}-by^2+by+t^2y
\end{split}
\]
のみで、この開集合上では全始点で求める次元が得られる。

$J$ の四つの冪と、それぞれへの $I$ の合成を使う。必要な係数変更は、一般テンソルで
全係数に成立する対角倍率と置換であり、薄い係数と規格化 $B$ の倍率が可逆な範囲で厚い零成分にも特殊化できる。
この点を具体的に確認するには、動かした表示の $B$ がいずれの薄いゲージ規格化の後でも
$x_1,x_4,x_{13},x_{15}$ のいずれかの可逆な基準係数またはその対角倍率倍になることを使う。
実際、各変換の逆も同じ表示式で定義され、窓の行列式に対応する $(zw-1)$ は可逆である。
変換先の $b'\ne1$ は $u'^2\ne1$ と同値である。
$u^2,v^2\ne1$ の開集合では八個の $L$ の像を取り、$M(u^2-1)(v^2-1)$ で飽和すると単位イデアルである。
$u=1,-1$ では $v^2\ne1$ の範囲で $b'\ne1$ の四表示だけを使い、両方で再び単位イデアルを得る。
$v=\pm1,u^2\ne1$ は回転でここに移る。残る $u^2=v^2=1$ では $u=v$ であり、
特別成分で $a,y\ne0,1$ の既知の正則点である。
\path{computations/verify_second_rectangle_low_fibers.py} が追加の一因子と三つの飽和を再現する。
\end{proof}

\begin{proposition}
76型の水平台85で $u=1,z=-1$ の全点は標準Koszul左右AS正則である。
その回転、コサイクル捻りにより $(v,z)=(1,1)$, $(u,z)=(-1,1)$, $(v,z)=(-1,-1)$ の各曲線も全点で正則である。
\end{proposition}
\begin{proof}
最初の曲線のパラメータは $v\ne0,\pm1$ である。
各位置を $(0,2,2,0)$ だけ平行移動した座標で、文字0の四次中心核の自由列の最初と最後、
文字1の両自由列を使う。これら四式を割った経路の長さ13までの次元は一般に
$(1+T+T^2+T^3)^4$ である。各始点で掃き出しに要る既約因子は順に $196,183,185,175$ 個である。
$v\mapsto1/v$ で移した二位置先のリストとの共通因子を取ると、始点0では14個、次数和31、
始点1では19個、次数和38となる。和集合は26因子、最大次数4である。
各因子の数体について、定義域を保つ一次剰余根で経路と四次中心の次元、四中心式の商の
長さ13までの全四始点の上界を検算する。中心の持上げと一般点のHilbert下界で標数零の全共役点に結論が従う。
有限個の始点に共通する一般の四式を取ることで自由性と左右正則性が得られる。
\path{computations/verify_first_rectangle_rank_ten_curve.py} は元の一変数体上の掃き出しと全26数体を検算する。
\end{proof}

% !TeX root = ../papers/nc-p3.tex
\section{五中心式上の表示と第一長方形の階数十曲面の全正則性}
\label{sec:first-rectangle-edge-regularity}
\begin{theorem}\label{thm:first-rectangle-edge-regularity}
台85のトーラス開集合で $u=1$ または $v=1$ の二曲面は、交線を含め全点が
標準Koszul左右AS正則である。各始点は一般の四中心四次式上で自由階数256である。
$u=-1$ または $v=-1$ の二曲面も全点が標準正則である。
これら四面の指定始点の無印同型はトーラス座標の一致で決まり、周期は
$u=v=\pm1$ で2、残りで4である。
\end{theorem}

\begin{proof}[五つの中心変数によるHilbert上界]
$u=1$, $v\ne1$ として始点 $i$ を固定する。四次中心の五基底を形式変数とする
$S=k[Z_1,\ldots,Z_5]$, $\deg Z_j=4$ を $M=e_iB$ に作用させる。
以下の開集合では全中心商の基底の長さ別個数が
\[
 (q_0,\ldots,q_9)=(1,4,10,20,30,36,34,20,1,0)
\]
となる。基底語を生成元とする自由加群 $L_0$ から $M$ へ次数帰納で全射がある。
$P(T)=(1-T^4)^5/(1-T)^4$ と置き、$q_n=0$ を高次へ延長し
$\beta_n=q_n-[T^n]P(T)$ とする。$\beta_n$ の非零部分は次数8から16まで順に
$1,20,34,36,30,20,10,4,1$ である。核にこれらの個数の $S$ 上独立な斉次関係を取れば、
生成する自由部分加群 $L_1$ に対し $L_0/L_1$ のHilbert級数は $(1-T)^{-4}$ となる。
主成分の幾何学的一般点から各次数に同じ下界があるので $L_0/L_1\to M$ は同型である。
従って $M$ は $S$ 上長さ1の自由分解を持ち、深さと次元がともに4である。
全始点に共通な四つの一般一次結合を次数4の斉次パラメータとすれば、各 $M$ の正則列となる。
その多項式環上の次数付き自由性と階数256は中心商の基底の持上げと級数比較から従う。
補題\ref{lem:parallel-residue-border}の双対完全性の議論により標準Koszul左右AS性を得る。

独立性は五変数の多項式を展開せずに証明できる。
$L_0$ に各矢の作用の $S$ 線型な持上げを選ぶ。具体的には各長さ・各色ごとに、
一文字延長のうち全中心商で消去する列を選ぶ。末尾の二次関係、および長さ4では
中心基底と $Z_j$ 倍の空語との等式を使う。一段短い語の積を先に代入すると、
最高長の係数行列は全中心商での行列であり、選んだ可逆小行列で解ける。
残す列にはその基底語をそのまま用いる。これは $M$ への評価と可換な持上げである。
$Z=c$ で評価した係数だけで同じ消去ができるため、全関係の多項式表示を保存する必要はない。

各基底語での全二次関係、空語での五つの中心等式から出発する。
それらを持上げた右矢と $Z_j$ 倍で閉じれば、常に $L_0\to M$ の核の斉次元を得る。
$V_n$ を内部次数 $n$ で得たこれらの元の $c$ での評価が張る空間とする。
$c_j=1$ となる $j$ を使えば $V_{n-4}\subset V_n$。長さの4による剰余ごとに、
$n\le16$ で $\dim V_n\ge\sum_{a\ge0}\beta_{n-4a}$ を検算する。ただし負の添字では零とする。
前の次数で選んだ独立列を毎回 $\beta_n$ 個だけ延長でき、異なる剰余は異なる基底語を使う。
評価後の独立小行列式が非零なので、元の斉次列は $S$ 上独立である。
\path{computations/quartic_center_presentation.py} が持上げと全次数の独立数を計算する。
\end{proof}

\begin{proof}[二つの曲面の全点]
$U=vz(v-1)(v+1)(v+z)(z-1)(vz+1)$ を逆にする。
表示Aでは位置別文字平行移動 $(0,2,2,0)$ と $c=(1,0,0,0,0)$、
表示Bでは $(2,0,2,0)$ と $c=(0,1,0,0,0)$ を取る。
各表示と始点ごとに、最高長の小行列・関係の独立性に使う小行列の分母とピボットの非単元既約因子を記録する。
Aの四始点での個数は $5,20,28,5$、Bの始点 $0,1,3$ では $7,19,24$ である。
これらは\path{computations/verify_first_rectangle_edge_presentation.py} で再生成できる。

二位置移動は $v\mapsto1/v$ で $z$ を保つ。その引戻しを上線で表す。
始点0は $A_0,B_0,\overline{A_2}$ で覆われる。実際、最初の二リストの因子の35対ごとに
二因子と $1-TU$ のイデアルに第三リストの因子積を加えると単位イデアルである。
始点1は $v+2z-1\ne0$ で $\overline{A_3},B_1,A_1,\overline{B_3}$ で覆われる。
最初の二リストの95対に対する同じ計算で、今度は $U(v+2z-1)$ を逆にする。
残る直線では $z=(1-v)/2$ を $\overline{A_3}$ の各因子に代入すると、単元を除く零点は
\[
 v^3-4v^2+3v-4,\quad v^2-2,\quad v^2+v+1,\quad 2v^2-3v+2
\]
の四数体に限られる。順に $(p,v)=(1013,412),(1009,570),(1009,634),(1009,274)$ は
単純な剰余根である。各点で四中心式の商を全始点の長さ13まで検算し、前節と同じ持上げで
標数零の全共役点に自由性と左右正則性を得る。
\path{computations/verify_first_rectangle_edge_regularity.py} が35対・95対および四数体を確かめる。
二位置移動で始点2,3も従い、一位置移動で $v=1$ の面へ移る。
交線と負側はそれぞれ既知の垂直成分の定理と四元群二コサイクル捻りによる。
同型と周期は定理\ref{thm:first-rectangle-horizontal-frames}による。
\end{proof}

% !TeX root = ../papers/nc-p3.tex
\section{第二長方形の反対階数と階数十四の特異K3曲面}
\label{sec:second-rectangle-opposite-ranks}
\begin{theorem}\label{thm:second-rectangle-opposite-ranks}
台143から147の水平表示で $b+c\ne0$ とする。$A=u+1/u$, $B=v+1/v$, $P=z+w$, $Q=zw$ と置き
\[
 f=(A+B)(P-2)-(B+2)(Q-1),\quad
 g=(A+B)(P+2)+(A+2)(Q-1)
\]
とする。反対階数は $u=v=1$ で8、ちょうど一方が1なら10、残りでは
$16-\mathbf1_{f=0}-\mathbf1_{g=0}$ である。
二つの階数15因子の閉包は有理三次元多様体で、一位置移動で交換される。
その交叉は階数14の滑らかな既約曲面である。指定した座標でのその極小滑らか射影模型は
Picard数20、Néron--Severi判別式の絶対値16の楕円K3曲面である。
台147でのこの曲面の一般の無印同型類は有理曲面でパラメータ付けられる。
階数14の曲面と両方の階数15因子には、全点が標準Koszul左右AS正則な非空開集合がある。
\end{theorem}

\begin{proof}[階数と分岐]
$C=(u+v)(1+uv)$, $H=u(v+1)^2$, $G=v(u+1)^2$ とする。
四色の反対行列の行列式から定理\ref{thm:horizontal-rectangle-low-fibers}の $M$ の単元を除くと、
順に $(v-1)^2f,(u-1)^2g,(u-1)^2(v-1)^2,(u-1)^2(v-1)^2$ である。
全色に単元の二次小行列式がある。三次小行列式のイデアルはそれぞれ
$(v-1),(u-1),((u-1)(v-1)),((u-1)(v-1))$ となる。
最初の色で共通因子を除いた後には、単元倍を除き $z-1,w-1$ が得られ、第二色では $z+1,w+1$ が得られる。
両方の同時消滅は $Q\ne1$ で排除される。残る二色では共通因子の単元倍が一つの小行列式で得られる。
$u=1$ なら $uv f=-(v+1)^2(z-1)(w-1)$ は可逆、$v=1$ なら
$uv g=(u+1)^2(z+1)(w+1)$ は可逆である。これで全点での階数が従う。
単独の二因子ではそれぞれ $w$ を一次で解け、係数と定数項が互いに素である。
一位置移動は $(A,B,P,Q)\mapsto(B,A,-P,Q)$ と作用する。

$N=A+B+4$ と置くと $f=g=0$ で $N(Q-1)=-4(A+B)$ なので $N\ne0$ であり、
\[
 P=2(A-B)/N,\qquad Q=1-4(A+B)/N .
\]
従って $(z,w)$ はこれらを和と積に持つ順序対である。
$Y=N(z-w)/4$ とすれば
\[
 Y^2=D(A,B)=A^2+AB+B^2+2A+2B-4.
\]
基底のトーラスで分岐曲線の特異点があるとすれば $u=\pm1$ または $v=\pm1$ である。
実際 $D=D_A=D_B=0$ に解はなく、$D(2,B)=(B+2)^2$, $D(-2,B)=B^2-4$ である。
このうち $(u,v)=(1,-1),(-1,1)$ は元のトーラス開集合の外、$(-1,-1)$ は $N=0$ である。
よって階数14の層は滑らかで、二因子はそこで横断的に交わる。
どちらかの座標が1の面上では前記の可逆性により両因子に同時には特殊化しない。
\end{proof}

\begin{proof}[楕円曲面と無印商]
変数 $A$ を底とし $v^2D(A,v+1/v)$ を四次式として取ると
\[
 (vY)^2=v^4+(A+2)v^3+(A^2+2A-2)v^2+(A+2)v+1 .
\]
$v=0$ に有理切断を持つ。二元四次式の不変式 $I,J$ はそれぞれ $A$ の次数4,6、
判別式は $(A-2)^3(A+2)^3(3A+10)^2$ である。
$A=2$ で $I$ は単純零点、$-2,-10/3$ と無限遠の極小表示では単元である。
従って四つのファイバーは $2,-2,-10/3,\infty$ の順に
$\mathrm{III},I_3,I_2,I_4$、基本直線束が次数1の有理楕円面である。
$A=u+1/u$ は $2,-2$ で分岐し、基底変換後のファイバーは
\[
\begin{array}{c|rrrr}
 u&1&-1&-3,-1/3&0,\infty\\\hline
 \text{型}&I_0^*&I_6&I_2,I_2&I_4,I_4
\end{array}
\]
となる。判別式次数24、切断と標準束公式から楕円K3となる。
有理楕円面の自明格子の根の階数は7なのでMordell--Weil階数は1である。
四根を逆数対に分ける分割が有理2捩れ点を与え、加法的ファイバーの成分群により捩れはちょうど位数2である。
有理楕円面の unimodular なNéron--Severi格子から原始生成元の高さは $4/(2\cdot3\cdot2\cdot4)=1/12$ である。

K3の自明格子の階数は19で、引き戻した非捩れ切断によりPicard数は20となる。
四次式の判別式は $u$ の関数体で平方である。四根の三分割に対する作用は、逆数対の分割を固定し
偶置換だけから来るので全ての分割を固定する。従って全2捩れが有理であり、$I_0^*$ の成分群から捩れの位数は4である。
引き戻した切断の高さは $1/6$。ファイバー寄与の分母は全て12を割るため、
これが非捩れ切断の $m$ 倍と捩れの和であれば $1/(6m^2)\in\frac1{12}\mathbb Z$ である。
よって $|m|=1$。格子の判別式の絶対値は
$4\cdot6\cdot2^2\cdot4^2\cdot(1/6)/4^2=16$ となる。

一般点での文字枠の復元は下記の剰余検査から従う。従って指定始点の無印商には
$I:(u,v,Y)\mapsto(1/u,1/v,-Y)$ だけを追加して割ればよい。
コンパクト化上の $u=v=-1,Y=0$ では $D_A=D_B=-4$ であり、局所方程式の二次部が非退化な通常二重点である。
対合の接作用は三座標で全て $-1$ なので解消の例外円錐曲線を点ごとに固定する。
この解消はK3の極小模型にも残る（例えば二重被覆のコンパクト化は分岐の型 $(4,4)$ を持つ。
四隅は二本の異なる接線を持つ節点で、$(1,-1),(-1,1)$ の残る特異性は
分岐曲線の二重接触、二重被覆の $A_3$ であり、全解消はcrepantである）。
固定曲線を持つ対合は正則二形式に $-1$ で作用する。その商は滑らかで、不規則数は0、
二重被覆公式で反標準因子の二倍は非零な有効分岐因子と線型同値である。
従って小平次元が負となり、商は有理曲面である。
\end{proof}

\begin{proposition}\label{prop:second-rectangle-rank-fourteen-period}
上の階数14曲面で $z=w$ の層は滑らかな有理曲線で、一般点は標準正則かつ最小周期2である。
その外の曲面の一般点は最小周期4である。
\end{proposition}
\begin{proof}
$l=(u-1)/(u+1),m=(v-1)/(v+1),r=lm$ と置く。$D=0$ は単元を除き
$1-3l^2m^2+l^2m^2(l^2+m^2)=0$ である。
\[
 S=r(l+m)/(r+1),\qquad E=r(l-m)/(r-1)
\]
と置くと $S^2=2r-1$, $E^2=-2r-1$ となる。逆に滑らかな円錐曲線 $S^2+E^2=-2$ から
$r=(S^2+1)/2$, $l=((r+1)S+(r-1)E)/(2r)$, $m=((r+1)S-(r-1)E)/(2r)$,
$z=w=-SE$ を回収できる。除いた分母の零集合は曲線の成分を含まない。
$J^2$ は指定始点の同型 $I$ を除けば $z,w$ を交換し、一般に奇数移動は固定点を持たない。

階数14曲面とその対角曲線で、それぞれ $u=2$ を取る。前者は $v=3$,
$59z^2+10z-81=0$, $w=-10/59-z$、後者は
$4v^4+18v^3+37v^2+18v+4=0$, $z=w=(A-B)/N$ である。
素数1009の単純根 $z=883$, $v=729$ で台147の全52係数は非零、各位置の隣接二位置に対する
二次モーメント一つだけで異なる四固有値を持つ。
全始点で五次までの経路次元と四次中心の次元 $(3,2,0,0)$ を確認し、中心四式の商は
長さ13まで $(1+T+T^2+T^3)^4$ の係数を持つ。中心核の特殊化については、
主成分一般点からの次元下界と剰余点の上界が一致し、経路も五次まで標準であるため単元小行列で持ち上がる。
前節までの中心商の上界と一般点の全次数下界との比較で、この二数体の点で正則性を得る。
四つの一般中心式の有限判定および枠の四固有値はその各点の近傍で開く。
これにより曲面、含む両因子、対角曲線の各非空開集合での主張が従う。
\path{computations/verify_second_rectangle_opposite_ranks.py} が行列式・曲面・円錐曲線を、
\path{computations/verify_second_rectangle_rank_fourteen_samples.py} が二つの数体の検査を再現する。
\end{proof}

% !TeX root = ../papers/nc-p3.tex
\section{五つの水平台の中心束と第一長方形の階数十四曲線の全正則性}
\label{sec:rectangle-quartic-bundles}
\begin{theorem}\label{thm:rectangle-quartic-bundles}
第二長方形の水平台143から147で $b+c\ne0$ なら、四次中心は四元群の文字順に
階数 $(3,2,0,0)$ のベクトル束をなし、各ファイバーでの中心と可換する。
第一長方形の台85の反対階数14曲線では、全点で四中心四次式上自由階数256の
標準Koszul左右AS正則代数が得られる。全点が最小周期4で、指定始点の無印同型類は
定理\ref{thm:first-rectangle-rank-strata}の $r$ で一意に決まる。
\end{theorem}
\begin{proof}[第二長方形の中心を解く共通表示]
各位置の文字番号に順に $(0,0,2,2)$ を加える。全和が零なので四次の文字は変わらない。
$b\ne1$ の表示 $t,b,x=a,y$ で $c=(t^2-b)/(1-b)$ とし、左上行列式の単元因子を
$E=(b+c)xy-bc(x+y)+b^2(c-1)$ と書く。
$tbc(b-1)(b+c)E$ および $q=x,y$ に対する $q(q-b)(q-1)(cq+b(1-c))$ の分子・分母の因子を逆にする。
これは定理\ref{thm:horizontal-rectangle-low-fibers}の共通の開条件で、五つの台で消えてよい座標を要求しない。
各二次関係のフレームは可逆小行列式を持つ。各始点の長さ5までの経路と、四色の長さ4の経路を一文字ずつ左右から
延長して差を取る可換方程式を、この基底環の単元だけで掃き出せる。経路の次元列は
$1,4,10,20,35,56$、可換方程式の自由列の数は $3,2,0,0$ であり、核・余核は特殊化と可換である。

$b=1,c\ne1$ は一位置移動でこの表示に入る。文字の線型部は厚い二色の差である文字1を保ち、
位置別平行移動の総和は零なのでこの移動でも四次の文字0,1を保つ。
$b=c=1$ では $t=1,-1$ を個別に代入する。今度は $E=2xy-x-y$ と
$x(x-1)y(y-1)$ を逆にした二変数環で、同じ平行移動が両符号で再び単元だけの掃き出しを与える。
\path{computations/verify_second_rectangle_quartic_bundle.py} が三表示を検算する。
\end{proof}

\begin{proof}[第一長方形の曲線を全点で覆う]
$u=r(r-1)/(r+1)$, $v=(1-r)/(r(r+1))$, $z=-4v/(v+1)^2$ を使う。
始点0の五中心式上の表示に定理\ref{thm:first-rectangle-edge-regularity}の独立関係による判定を適用する。
今回は全中心商の基底個数は長さ0から8まで
$(1,4,10,20,30,36,34,20,0)$ となる。従って必要な新関係は次数9から16まで順に
$20,34,36,30,20,10,4,1$ 個である。中心変数の定数評価と位置別文字移動を
$((1,0,0,0,0),(0,2,2,0))$ および $((0,1,0,0,0),(2,0,2,0))$ と取る。
元の曲線で可逆な因子を除けば、中心核・商基底・全ての関係の独立性を同時に保証する二表示の
例外既約因子は120個と141個である。その共通部分は次の八数体で尽くされる。
\[
\begin{gathered}
r^2+6r+1,\quad r^2-3,\quad r^2-6r+1,\quad r^2-6r-1,\\
r^4\pm4r^3+2r^2\mp12r+9,\qquad
r^4\pm4r^3-6r^2\mp12r+1 .
\end{gathered}
\]
それぞれで定義域を保つ次数1の単純剰余根を取り、四中心商を長さ13まで全始点で検算する。
\path{research/first_rectangle_rank_fourteen_residues.json} に素数と根を、
\path{computations/verify_first_rectangle_rank_fourteen_curve.py} に二表示・八数体の再現計算を掲げる。
曲線は反対階数で内在的に定まるため一位置移動で保たれる。従って全点で始点0の表示が得られれば、
全点の全始点で自由性と深さ4が従う。四始点に共通の一般の四中心式を取り左右正則性を得る。
定理\ref{thm:first-rectangle-horizontal-frames}の周期2の二直線とは交わらない。
\end{proof}

% !TeX root = ../papers/nc-p3.tex
\section{三つの長方形の水平成分と両符号の垂直成分を通る無制限変形}
\label{sec:rectangle-unrestricted-components}
\begin{theorem}\label{thm:rectangle-unrestricted-components}
二つの隣接切断がちょうど階数6のアフィンテンソル軌跡を $Z$ とする。
76型と112型の水平主成分、および第三長方形族の全閉包は、$Z$ 内で
定理\ref{thm:full-fifteen-component}の同じ既約成分に属する。
両型の $t=1$ の垂直族は、これと異なる一つの既約成分に属し、その一般点は両側の反対位置を
それぞれ対称化できる。両型の $t=-1$ の垂直族も一つのさらに異なる既約成分に属し、その成分は
定理\ref{thm:klein-other-pairings}の $\lambda=-1$ の表示の全基底変換による閉包である。
三成分のアフィン次元はそれぞれ $76,75,65$、指定始点の一般同型類の次元は $15,14,4$ である。
各々の幾何学的一般点は標準Koszul左右AS正則である。両符号の水平・垂直交叉はそれぞれ異なる次元の成分に属する。
\end{theorem}

\begin{lemma}\label{lem:finite-flat-conormals}
命題\ref{prop:twenty-conormal}の横断的二十点を、長さ20の有限完全交叉 $X$ に置き換えられる。
二つの部分評価には全台点で階数3を要求する。余法線の独立条件は、四配置の斉次座標の積
$a_i b_j c_k d_l$ の $X$ 上への制限が、対応する可逆層の全切断の二十次元空間を張ることである。
この条件下で二つのSchur微分の合計階数が180なら $Z$ はその点で滑らかな76次元である。
\end{lemma}
\begin{proof}
第一切断だけをその滑らかな階数6軌跡上で動かす。第一辺の六双一次式が定める射影族は
この点の近くで有限であり、相対完全交叉の平坦性判定で階数20の有限平坦射になる。
局所基底環上のその代数を $D$ とする。特殊ファイバーは有限個のArtin局所環の積なので、
各射影座標には全台で消えない一次形式が選べる。評価像の補空間 $H$ も同時に選べる。
命題\ref{prop:twenty-conormal}の $X=QJ_a$, $Y=J_b^{\mathsf t}P$ の議論を $D$ 上で行う
（この式の $X$ は行列を表す）。用いる三次小行列は各台点で選べ、核の直線束は
基底を局所化して自明化できる。これは特殊ファイバー上で生成元を選びNakayamaの補題で持ち上げればよい。
第二切断の階数方程式を課す前から、百個のSchur成分に対する同命題の恒等式が $D$ に値を持つ。
$D$ の自由基底でその係数を読めば基底環上の二十本の恒等式となる。
仮定した全切断の生成性は、まさに基点でこの二十行が独立であることを表す。
非主行・列への射影が左右核上の同型であることも各点での可逆小行列から従う。
従って先の命題と同様に八十本の生成系が得られ、微分階数により滑らかである。
\end{proof}

\begin{proof}[水平の三族を通る成分]
76型で $(u,v,z)=(2,3,5)$、112型で $(u,v,z,w)=(2,3,5,4)$、
第三長方形で $(A,D,b)=(2,3,5)$ を取る。素数1009で二切断のSchur微分を四文字に分けて計算すると、
共通接空間の文字別次元は三例とも $(18,18,20,20)$、基底変換からの次元は $(13,16,16,16)$ である。
第一辺を第一射影の四次小行列式で求め、一つのアフィン超平面で全台を覆う。
その座標環を標準単項式の二十次元空間とし、線型形式を可換な乗法行列として扱う。
各隣接評価の三行結合に規格化一次式を加えた四行四列の行列が、この環上で可逆である。
これで $a,b,c,d$ を順に解き、残りの評価式と一周後の一致、$a$ から逆に解いた $d$ も検算する。
したがって第一辺はこの有限スキームのグラフで、六式の有限完全交叉である。
四座標の積の256列は環の単位に作用させると階数20を持つ。
有限平坦な点環と単元小行列を素数で局所化した基底に持ち上げれば、標数零で補題が適用できる。
特にどの点も全 $Z$ の滑らかな点である。

四元群の固定軌跡も標数零では滑らかで、ここでの次元は不変接空間の18である。
これは完備正則局所環の余接空間に同変な持上げを選び、非自明固有文字のパラメータを消せば分かる。
微分階数の剰余検算を使う際、二十恒等式は局所化した整基底上で成り立つので全階数は180を超えず、
文字別の核も特殊化と可換である。
三例とも不変接ベクトルのうち、薄い各行列の零行・零列を一本も恒等的に零に保たないものがある。
このベクトルを固定軌跡で積分すると、一般点で薄い全行・全列が非零となる。
各薄い行列は階数1なので全成分が非零である。
全ての厚い交差ブロックの行列式も、基点で非零またはこの方向の一次項が非零である。
従って一般点では両交差ブロックが可逆で、階数2により対角ブロックも可逆である。
第二切断の色0の右上ブロックでは二つの対角積の和も非零で、薄い規格化の $z=-1$ を排除する。
その行列式の非零が $z=1$ を排除する。これで定理\ref{thm:klein-all-pairings}の全被覆に入る。
この固定成分は対角軌道の13次元を除いても5次元あり、各々4次元以下の他の四表示には入らない。
よって三点はいずれも五次元楕円族の閉包、従って既知の76次元成分に属する。
元の水平族は既約でそれぞれの点で全 $Z$ の成分が一意なので、族の全閉包について従う。
\end{proof}

\begin{proof}[垂直族での二つの別成分]
両型とも $b=c=1$, $t=\pm1$, $A=3t,B=1,D=7t$ とし、$C$ は76型で1、112型で4とする。
不変部分も含む全接空間の文字別次元の上界は正符号で $(17,18,20,20)$、負符号で $(17,16,16,16)$ である。
両側の基底変換の次元は先と同じで、射影安定群は有限である。

正符号では第四位置に文字1を加えると、$\mathrm{Sym}^2 E\otimes\mathrm{Sym}^2 F$ の元である。
この百次元空間の第一切断は $E\otimes F$ 上の対称行列であり、他方の隣接階数も同じになる。
対称な十六行行列に対する階数6以下の余次元は55なので、この固定した二同定の軌跡は各成分で少なくとも45次元である。
二同定を動かす自由度30を加えられる。実際に同じテンソルの同定の二通りから得られる二交換の合成は射影安定群に属し、
その有限性により同定を忘れるファイバーはここで有限である。
次元下界75と接次元上界75は等しく、全 $Z$ は滑らかで、通過する唯一の成分は一般に両対称となる。

同じ論証を四元群を固定して行う。第四位置の変更後の総文字は1であるため、両対称テンソルの空間は
$4\cdot2+2\cdot4+2\cdot2+2\cdot2=24$ 次元である。
二つずつの互いに転置な四色ブロックには階数2と1を一度ずつ課せばよく、余次元の上界は $4+9=13$。
各反対二位置の同定を同変な対角群内で動かす六自由度を加えると、不変軌跡も少なくとも17次元になる。
上の微分上界と一致して滑らかである。従ってこの軌跡の一般点でも、偶二位置の同定は文字を置換しない対角型、
奇二位置の同定は文字1の移動を伴う型である。
水平族と同じ薄い行列の微分検査により、一般点は全薄い成分が非零のチャートに達し、
$z=\pm1$ も除かれる。偶同定が対角型なので規格化では $M_2=N_2$, $M_3=N_3$ とできて $\lambda=1$。
全接次元が75の基点は76次元成分には含まれず、その楕円族側の閉包に入る可能性も消える。
残るのは $\lambda=1,A_1=A_0$ の既約な四変数表示である。
不変成分は対角軌道を除いて4次元なので、両型の正符号の基点はこの同じ表示の稠密開集合からの特殊化となる。
この表示の一般点で全 $Z$ の成分も一意である。実際、その閉包に含まれる上の滑らかな点の近傍は一般点を含む。
両型の正符号の垂直族の全閉包が一つの75次元成分に入る。

負符号は四元群二コサイクル捻りである。固定軌跡上でこの操作は線型な自己同型で、
正の垂直族とそれを含む四変数表示を、負の垂直族と $\lambda=-1$ の表示に移す。
不変軌跡の下界17から、全基底変換を許せば軌道の非自明文字部分の48次元をさらに加えられる。
厳密には正側の17次元不変軌跡が滑らかな点の像でこの48方向は不変接空間と交わらず、
あるいは有限射影安定群内の四元群の選び方が有限であることから次元下界65が従う。
負側の一般垂直点での上界65と一致し、四変数表示の基底変換による閉包そのものが全 $Z$ の既約成分になる。
全ての次元下界は同じ記号の標数零の点で得られ、剰余での接上界の等号もこれによって保証される。

一般同型類の各次元は一般軌道の61次元を引けばよい。垂直成分の両符号とも既知の正則点から
内部次数ごとの完全性を一般点へ持ち上げられる。水平・垂直の交叉点が特異であることは、
各々でこれら二つの異なる成分を通ることからも分かる。
\path{computations/verify_rectangle_component_smoothing.py} が点環の余法線、七つの代表の微分と薄い成分の変形、
および各型・符号の交叉例の文字別接次元 $(19,18,20,20)$ を検算する。
\end{proof}

% !TeX root = ../papers/nc-p3.tex
\section{第二長方形の階数十四・周期二曲線の全点分類}
\label{sec:rank-fourteen-conic-complete}
\begin{theorem}\label{thm:rank-fourteen-conic-complete}
命題\ref{prop:second-rectangle-rank-fourteen-period}の対角曲線の全点は標準Koszul左右AS正則で、
最小周期2、射影安定群は四元群である。指定始点の無印同型類は射影直線から十二点を除いた曲線と一対一に対応する。
座標を $s=h^2$ とすれば、除くのは $0,\infty$ および
\[
 (s^2+8s+4)(s^2-8s+4)(s^2+4)
 \big((s^2+8s+4)^2-16s(s-2)^2\big)=0
\]
の根である。残る点のうち十二点に一つの厚い係数の消滅があり、他では台147の全52係数が非零である。
一位置移動は $s\mapsto-s$ である。
\end{theorem}

\begin{proof}[パラメータと枠]
$i^2=-1$ とし、前節の円錐曲線で $h=S+iE$ と置く。
$S=(h-2/h)/2$, $E=(h+2/h)/(2i)$ である。
$A=h^2+2h-2$, $B=h^2-2h-2$, $C=h^2+2ih+2$, $D=h^2-2ih+2$ とすると
\[
 u=-AC/(BD),\qquad v=-AD/(BC),\qquad z=w=i(h^4-4)/(4h^2).
\]
$r=0,1,-1$ で逆変換が特異となる可能性については、有限な $l,m$ で $r=0$ は
円錐曲線の元の方程式に反し、$r=\pm1$ なら $l^2=m^2=1$ となり $u,v$ はトーラスにない。
$u=-1$ または $v=-1$ は $D(A',B')=0$（$A'=u+1/u,B'=v+1/v$）において他方も $\pm1$ を強制し、
元の階数14層には入らない。$h=0,\infty$ の極限も $u=v=-1$ である。
トーラス開条件の分子・分母の平方自由積は
\[
 h(h^4+8h^2+4)\prod_{e=\pm1}(h^2+2eh-2)(h^2+2eh+2)(h^4+4eh^3+8h^2-8eh+4).
\]
従って有限遠で21点、無限遠を加えて22点が除かれる。

テンソルの交代勾配を使い、各位置で他の三位置に対する三つの二次モーメントの対角成分を作る。
共通分母を払った後、任意の二文字の三つの値の差について $\mathbb Q(i)[h]$ から
$\mathbb Q[h]$ へのノルムの最大公約数を取る。上の境界の全因子を取り除くと四位置の全24対で定数になる。
よって共同固有線は全点で四つの元の文字線であり、印を忘れた同型は単項変換である。

開曲線上で消え得る厚い係数は $x_3,x_6,x_8,x_{10}$ だけである。各々で
境界を除いた零因子は $\mathbb Q(i)$ 上の相異なる二つの三次既約式で、計八式は全て互いに素である。
このため一点で消えるのは一係数だけであり、指定した係数枠で24点が該当する。
これら一欠損台の間の四置換を全探索すると、全て元の52成分の台の置換に延びる。
薄い係数のゲージと $B=1$ の規格化は引き続き可逆なので、
定理\ref{thm:parallel-horizontal-moduli}の $I,J$ がそのまま全点で使える。
対角倍率の非零Smith因子は各台で $1^{11},2,2$ である。
$I$ は $h\mapsto-h$、一位置移動 $J$ は $h\mapsto ih$ となる。
$h\ne0,\infty$ より指定始点で非自明な台置換は点を固定せず、安定群は対角四元群である。
よって指定始点の商は $h^2$ であり、最小周期は2である。
\path{computations/verify_rank_fourteen_conic_frames.py} が境界、係数の零点、共同固有線と全台置換を検算する。
\end{proof}

\begin{proof}[全点の正則性]
始点0で五中心商を計算し、定理\ref{thm:first-rectangle-edge-regularity}の五変数上の表示を使う。
表示Aは位置別文字移動 $(0,0,2,2)$ と第一中心座標の定数1で、全52係数の非零を仮定する。
表示Bは $(2,2,0,0)$ と第二中心座標の定数1で、境界22点だけを除いた基底環を使う。
係数を $(a+ib)/d$, $a,b,d\in\mathbb Q[h]$ で表し、ピボットの分子ノルム $a^2+b^2$ と
分母の既約因子を記録する。前節の単元小行列で特殊化する関係フレームを取り、二次・四次の経路関係を順に消去する。
両表示で全中心商の長さ別基底個数は $(1,4,10,20,30,36,34,20,0)$、
必要な独立関係は次数9から16まで $20,34,36,30,20,10,4,1$ 個である。
評価後の関係の階数条件までに使う境界外の既約因子は各表示で174個と171個である。
最初の表示に追加で仮定した全係数非零の条件も含め、同時に両表示から漏れるのは
次数 $2,2,4,4,6,6,6,6,6,6,8,12,16,16$ の十四個の有理既約式の零点に限られる。

各式について次数1の単純剰余根を持つ素数 $p\equiv1\pmod4$ を選び、
$i$ の二つの剰余を\emph{ともに}使って四中心商を長さ13まで全始点で検算する。
有理既約式の一根で定める数体と $\mathbb Q(i)$ の合成は、直積になる場合でもこの二つの選択で
全ての体因子が覆われる。元の経路の全次数下界と持ち上げた中心商の上界から四中心式上の自由性が従う。
曲線は $J$ で保たれるので、開集合の始点0の結果も全点の全始点に及ぶ。
\path{computations/verify_rank_fourteen_conic_presentation.py} と
\path{research/rank_fourteen_conic_residues.json} に両表示と十四の残余検証を与えた。
\end{proof}

% !TeX root = ../papers/nc-p3.tex
\section{四次元負成分の楕円点グラフと枠を保つ特殊化}
\label{sec:rectangle-component-point-geometry}
\begin{theorem}\label{thm:negative-component-elliptic-graph}
定理\ref{thm:rectangle-unrestricted-components}の次元4の負成分では、
一般の各辺の点対応は滑らかな楕円四次曲線と四つの被約孤立点の同型グラフである。
この成分のどの窓でも各辺の点対応は正次元である。
\end{theorem}
\begin{proof}
$A_0=\left(\begin{smallmatrix}a&b\\c&d\end{smallmatrix}\right)$,
$A_1=-A_0$, $\lambda=-1$, $D=ad-bc$ とする。第一辺の六関係の左評価は
\[
\begin{pmatrix}
 ax_0&cx_1&x_2&-zx_3\\ bx_0&dx_1&-x_2&x_3\\
 -bx_1&-ax_0&x_3&-x_2\\ -dx_1&-cx_0&-zx_3&x_2\\
 x_2&zx_3&x_0&x_1\\ x_3&x_2&x_1&x_0
\end{pmatrix}.
\]
右評価では各双一次式を転置する。各射影に対して二本の二次式を
$Hx_0^2=Lx_2^2+Mx_3^2$, $Hx_1^2=Nx_2^2+Ox_3^2$ と書く。左では
\[
\begin{split}
 H&=(a-d)D,\quad L=-(b+d)(c+d),\quad M=(c+dz)(bz+d),\\
 N&=-(a+b)(a+c),\qquad O=(a+bz)(az+c).
\end{split}
\]
右では
\[
\begin{split}
 H&=(bz-c)D,\quad L=-(c+d)(az+c),\quad M=(a+c)(c+dz),\\
 N&=-(a+b)(bz+d),\qquad O=(a+bz)(b+d).
\end{split}
\]
いずれも一般点で $HLMNO(LO-MN)\ne0$ なので滑らかな四次完全交叉である。
左では $(x_2=x_3=0,\ x_0^2=x_1^2)$ と
$(x_0=x_1=0,\ x_2^2=zx_3^2)$ の各二点を加える。右では二つの係数 $1,z$ を入れ替える。
この和のイデアルで十五個の四次小行列式を割ると、関数体上で余りは全て零である。
$(a,b,c,d,z)=(5,2,3,7,3)$ では左右とも四アフィンチャートで逆包含も成り立ち、
三次小行列式の共通零点はない。\path{computations/verify_negative_component_point_graph.py} に
六関係の元の切断との照合、関数体での除算、有理特殊点の両方向の検算を与えた。

既知の和は標本の近くで平坦である。包含を与える構造層の全射の核は、平坦性により標本上で消え、
台の射影が閉なので近傍でも消える。三次小行列式の無共通零点性も開条件である。
従って両射影は局所的に階数3の核を取るグラフであって、上の和にスキームとして同型である。
他の辺には捻りの前の両対称性を使う。文字を平行移動する際の
双指標の変化は各文字の対角倍率となる。捻りの後の逆辺は転置に加えて
$\mu(i,j)\mu(j,i)$ が掛かるが、$q(i)=\mu(i,i)$ とすればこれは $q(i)q(j)q(i+j)$ である。
関係は色 $i+j$ ごとに斉次なので逆辺も各文字の基底変更で元の辺の転置に戻る。
各辺の射影的点対応のファイバー次元の上半連続性と成分の稠密表示から最後の主張を得る。
\end{proof}

\begin{lemma}\label{lem:stable-twenty-frames}
辺の六式の零スキームが有限なら長さ20の完全交叉で、二射影は閉埋込みである。
一射影の相異なる点で一つの $r$ 次元射影部分空間に入るものは、
$r=0,1,2$ に応じて高々 $1,4,10$ 個である。
離散付値環上の二つのそのような配置で、特殊点の相異なる点の個数を $n_1,n_2$ とする。
$n_1+n_2\ge32$ なら、一般ファイバーの射影同型は特殊点に可逆に延びる。
この主張の一般ファイバーの同型は、スキーム同型を与える射影線型変換をいう。
\end{lemma}
\begin{proof}
有限なら各部分評価は零点上で階数3でなければならない。階数2以下なら核が射影直線を含む。
可逆三次小行列の開集合で核を解くと射影は閉埋込みである。
六因子の交点数は $(H_1+H_2)^6=20H_1^3H_2^3$ である。
$\mathbb P^r\times\mathbb P^3$ に制限して $r+3$ 本の一般線型結合を取る。
共通基点は有限で、各段階で正次元成分を含む切断を避けられるから、その長さは
$\binom{r+3}{r}$ となり点数の上界を得る。

相対完全交叉は有限平坦である。必要なら付値環を有限拡大して一般点の台を分解する。
重複度込みの長さは各特殊点に特殊化する一般点の重複度の和である
（完備化後の有限自由代数を特殊点の冪等元で分解する）。
任意の $r$ 次元空間に載る特殊点の総重複度は上の上界に $20-n_i$ を足した値以下である。
一般ファイバーの射影変換をSmith標準形で最小指数0にすると、
特殊点で階数 $e<4$ なら核の外に特殊化する一般点の像は全て特殊行列の像に特殊化する。
従って総数20は、次元 $3-e$ の核の重複度と次元 $e-1$ の像の重複度の和以下となる。
ところが異なる点の上界の和は $e=1,3$ で11、$e=2$ で8であり、
二つの重複補正の和は8以下である。矛盾により射影変換は可逆に延びる。一意性は分離性による。
\end{proof}

\begin{proposition}\label{prop:stable-component-markings}
四位置の有限配置がそれぞれ十六個以上の相異なる点を持つ軌跡では、
一般点での位置移動を伴う窓同型、四元群の正則文字枠、および一般点での両対称表示は特殊化する。
特に次元14の正成分のこの部分における全ての窓は両対称表示を持ち、
定理\ref{thm:bisymmetric-two-length-test}の経路長3と21で正則性が決まる。
窓の同型では対応先の四配置にも同じ点数を仮定する。
\end{proposition}
\begin{proof}
補題を各位置の射影変換に使う。テンソル直線の一致は閉条件で、
両対称なら二交換を同時に行う写像が二位置移動の同型である。
その各位置の同定を可逆に延ばせば、一交換ずつのテンソルの等式も特殊化する。
成分の閉点には一般表示の稠密開集合から付値曲線を取ればよい。
一般点の射影同型が代数拡大を要する場合はその曲線を有限被覆してよい。

四元群の各生成元は一般点で二乗1の同変な持上げを持つ。
可逆な射影延長と行列式の付値0によりこの持上げ自身が格子を保つ。
文字冪等元の像は直和因子なので四つの階数1は特殊点でも保たれ、テンソルの文字も保たれる。
両対称表示と同様、有限個のテンソル等式で表す文字付き表示にこの手順を使える。
十六点以上を個々に固定する射影変換は恒等写像である。実際非スカラーな行列の固有空間の和集合には、
上の上界で高々十一点しか載らない。従って同型と周期は四配置間の有限個の点の置換を試し、
対応する射影行列の可逆性と四テンソルの直線の一致を調べて判定できる。
\end{proof}

\begin{proposition}\label{prop:elliptic-component-frames}
各位置の射影点スキームの台が滑らかな非退化楕円四次曲線と有限集合の和である
負成分の窓では、一般の四元群文字枠と、逆捻り後の二つの対称化同定が特殊化する。
そのような点での同型は各鉛筆の四つの特異二次曲面の頂点を置換する。
\end{proposition}
\begin{proof}
一般楕円曲線のHilbert点を固有性で延ばす。特殊曲線は台の唯一の一次元成分を含み、
この被約成分だけで同じHilbert多項式 $4n$ を持つので一致する。
従って曲線族は滑らかである。曲線の $4p\sim\mathcal O(1)$ なる十六点は、
相対Picard上の四倍写像から有限エタールに延びる。点・直線・平面内のこの相異なる点の
個数は高々 $1,2,4$（直線上の三点なら、それを含む全平面が三点の残りの
次数1の束を動かすことになる）である。従って補題と同じ核と像の議論で
二配置間の一般射影同型は延びる。四元群については前命題の文字射影を使う。

捻る前の二交換の合成で四文字の前後二組を交換すると、コサイクル符号の比は
$\mu(i+j,t)\mu(t,i+j)$ である。ただし $t$ は四文字の総次数で、比は二文字 $i,j$ の各文字倍率の積となる。
文字の平行移動による補正も一文字の倍率だけである。よって捻り前の各反対同定は、
格子に可逆な文字倍率を除いて捻り後の位置移動の各射影変換に等しく、可逆に延びる。
逆捻りをして二つの交換式を特殊化すればよい。最後に滑らかな四次完全交叉の二次式は鉛筆をなし、
四つの相異なる特異二次曲面を持つ。その頂点は鉛筆を同時対角化する四頂点そのものである。
\end{proof}

\begin{theorem}\label{thm:klein-bisymmetric-irreducible}
$K=(\mathbb Z/2)^2=\{0,1,2,3\}$ とし、偶二位置・奇二位置をそれぞれ同一視した
正則文字基底を固定する。各々の交換に対称、総文字1、第一切断の色順の階数が
$(2,2,1,1)$ のテンソル全体は既約で次元11である。
従ってこの型の任意の文字付き両対称窓とその非自明コサイクル捻りは、それぞれ上の正・負成分に属する。
負成分のうち前命題の楕円曲線が残る部分では、この逆捻り表示の存在が所属の必要十分条件である。
この部分では零係数の有無によらず同じ経路長3と21の有限正則性判定が使える。
\end{theorem}
\begin{proof}
$M_h(i,k)=w(i,i+h,k,k+h+1)$ とする。二重交換から $M_{h+1}=M_h^{\mathsf t}$。
$M_2=ab^{\mathsf t}\ne0$ は階数1であり、$a,b$ は共通逆倍率を除いて決まる。
一交換ずつの対称性を使うと、$M_0=\left(\begin{smallmatrix}X&U\\V&Y\end{smallmatrix}\right)$ の
$X,Y$ は自由な二行二列、交差ブロックは
\[
 U=\begin{pmatrix}a_2b_0&a_0b_3\\a_1b_2&a_3b_1\end{pmatrix},\qquad
 V=\begin{pmatrix}a_0b_2&a_2b_1\\a_3b_0&a_1b_3\end{pmatrix}.
\]
この表示に残る条件は $\operatorname{rank}M_0=2$ だけである。係数の全軌道と
前命題の符号は \path{computations/verify_klein_bisymmetric_closure.py} で検算する。
七次元の階数1行列の上の十五次元空間で階数2条件の余次元は高々4なので、
解の各成分は少なくとも11次元である（階数多様体のグラフとの交叉次元定理）。
$\Delta=\det U=-\det V\ne0$ なら $X$ は可逆で $Y=VX^{-1}U$、従ってこれは既約な十一変数の開部分である。

$\Delta=0$ のファイバーに限った成分がないことを示す。
$M_0=\binom{L}{R}(C\ D')$ を全列・全行階数2の因子に分けると
$LD'=U$, $RC=V$、共通の $\mathrm{GL}_2$ による四自由度を除く。
$U\ne0$ の二行正方分解 $(L,D')$ は高々4次元である。
一方が可逆なら4次元、双方が階数1なら $U=uv^t$ として
$L=ul^t$, $D'=dv^t$, $l^t d=1$ より3次元である。
$U=0$ なら高々5次元であり、非零な二階数1因子の中間直線と各々の二係数で数えられる。
$V$ も同じで、ファイバーの次元は $4+\mathbf1_{U=0}+\mathbf1_{V=0}$ 以下である。
$U,V\ne0$ の判別式上は基底が高々6次元だから合計高々10である。
一方のブロック全体が零なら四成分の各々は相異なる $a_i$ と相異なる $b_j$ の積である。
従って少なくとも四座標が零で、基底は逆倍率を除き高々3次元、合計高々9となる。
これで既約性は疎な全ての台にも及ぶ。

一般の固定枠ではこの既約開部分が定理\ref{thm:rectangle-unrestricted-components}の
正の四変数表示に達する。同変同定を動かす六自由度と固定基底軌道の十三自由度を
考慮すれば $11+6-13=4$ である。捻りは固定色のテンソル空間の可逆線型変換だから
負成分への包含も従う。逆向きに前命題で延びた枠と二同定を取ると、
各色の階数は一般値から増えず、和が6のままなので個々の階数は落ちない。
逆捻りの各隣接切断は色別の左右対角倍率を除いて元の切断に等しい
（交差する二対の符号は色ごとの定数である）。
よってこの型の表示が復元される。捻りは経路の次元と左右Koszul・AS性を保つので
両対称の有限判定が全ての表示に適用できる。
\end{proof}

\begin{corollary}\label{cor:negative-component-period}
負成分の幾何学的一般点の最小周期は2、指定始点の射影安定群は四元群である。
滑らかな楕円四次が各位置の唯一の曲線として残る軌跡では、任意の二点の同型、
各点の周期と安定群を有限個の置換と整数格子で判定できる。
\end{corollary}
\begin{proof}
各鉛筆の四頂点の枠でテンソルを書く。始点の移動を指定し、四位置の置換で台を合わせた後、
対応する非零係数の比を $r_q$ とする。$E$ を台の各語 $q$ を行とし
$E_{q,(p,i)}=[q_p=i]$ とした整数行列とする。対角倍率の存在は
$\prod_q r_q^{v_q}=1$ が $\ker_{\mathbb Z} E^{\mathsf t}$ の基底の全てで成り立つことと
同値である。これは代数閉体上のトーラス準同型の像の文字による表示で、零係数のある台にもそのまま使える。
解は一つあれば同じ対角安定群の剰余類であり、各位置のスカラーを除けば窓の射影安定群を得る。

一般表示では対角鉛筆の四頂点は元の文字線である。非零台は
$i+j+k+l=0$ の64語全体で、これを保つ置換は $i_p\mapsto\phi(i_p)+u_p$、
$\phi\in\mathrm{GL}(2,\mathbb F_2)$, $\sum_p u_p=0$ に限る。実際、各置換の零での値を引き、
二位置に等しい文字を置けば四写像は一致し、三位置に $i,j,i+j$ を置けば共通写像は加法的である。
$E$ のSmith因子は $1^{11},2,2$ である。上の有理標本の剰余1009で各移動につき384通りを調べ、
移動 $0,1,2,3$ の順に残るのは $1,0,1,0$ 通りである。
全ての係数が単元であるこの点で失敗する単項式の等式は、関数体上の恒等式ではない。
恒等置換と既に示した二位置移動は一般点でも成立し、対角安定群の有限部分はSmith因子から位数4である。
\path{computations/verify_negative_component_period.py} が飽和整数核を使って検算する。
\end{proof}

定理\ref{thm:rectangle-unrestricted-components}の三つの水平標本と本節の $A_1=A_0$ の
正標本では、四位置の全てで十六点条件を適用できる。
長さ20の有限自由点環のアフィン四座標を $a_0,\ldots,a_3$ とすると、
$\sum 2^i a_i$ の乗法の特性多項式の平方自由部分の次数は、剰余1009で標本順に $18,19,18,20$ である。
これだけの相異なる値は標数零にも持ち上がる。
\path{computations/verify_rectangle_stable_configurations.py} が全十六配置で検算する。

% !TeX root = ../papers/nc-p3.tex
\section{負交換の射影商、周期一の因子と退化した楕円枠}
\label{sec:negative-pairing-unmarked-quotient}
定理\ref{thm:negative-component-elliptic-graph}の負交換の座標を用い、
$zD(z^2-1)\ne0$, $D=ad-bc$ とする。$A=[a:b:c:d]\in\mathbb P^3$ であり、
正則点とは定理\ref{thm:exceptional-exchange-eight-lines}の八直線を
$(a,-b,-c,d)$ に対して除いた点をいう。以下
\[
 p=a+d,\quad q=a-d,\quad r=zb+c,\quad t=zb-c
\]
とする。
\begin{theorem}\label{thm:negative-pairing-unmarked-quotient}
$qt\ne0$ の全正則点で、指定始点の同型類は次の次数16の群 $H$ の軌道である。
$U,V,S$ はそれぞれ $z$ を $1/z$ とし、射影四座標を
\[
 U:[a:c:b:d],\qquad V:[d:zb:c/z:a],\qquad S:[a:zb:c/z:d]
\]
に移す。第四生成元 $B$ と、一位置移動 $J$ は $z$ を保ち、$(p,r,q,t)$ には
\[
 B:\big(2zp+(z+1)r,-z(z+1)p-2zr,-(z-1)t,-z(z-1)q\big),
 \qquad J:(r,zp,-t,zq)
\]
で作用する。$z$ を固定する部分群は位数8の二面体群である。
この全域で最小周期は $r=-p$ または $r=-zp$ のときちょうど1、その他で2である。
前二平面は商上で一つの既約因子を与える。一般商は有理四次元である。
\end{theorem}

\begin{proof}[全正則点での頂点の復元]
$X_i=x_i^2$, $P=D(X_0-X_1)-(a+b+c+d)X_2+(az+bz^2+c+dz)X_3$,
$Q=aX_0-dX_1-X_2+zX_3$ とする。左の四次完全交叉 $E$ は $P=Q=0$ である。
$Dq\ne0$ により前節の二本の二次式と同じイデアルを張る。
平方座標の直線上で $(X_0,X_1)$ と $(X_2,X_3)$ の各対は独立である。
実際、前節の $LO-MN=(a-d)(bz-c)(z-1)D$ である。
各座標に従う二重被覆の分岐点となる四つの零点は、各対の中では重ならない。
全て異なれば $E$ は滑らかである。重複が一対なら、二文字の平方の比例を因数分解して
二つの滑らかな平面二次曲線を得る。これらは他の二文字が定める座標直線上の二点で横断する。
二対なら二つの独立な平方比例を因数分解し、被約な四直線の四角形を得る。
重複はいずれも $\{0,1\}$ と $\{2,3\}$ から一文字ずつの対である。

全ての四次小行列式は $E$ 上で消える。これは関数体上での割算を
$X_0,X_1$ について行う際の分母が $Dq$ の冪だけであることから特殊化できる。
$E$ 外では点は前節の四孤立点に限ることを示す。行を0から数え、$m_I$ を四次小行列式とする。
$m_{0145},m_{2345}$ の和と差から $x_2x_3 Q=x_0x_1 P=0$。
$Q=0,P\ne0$ なら $m_{0124},m_{0125},m_{0134},m_{0135}$ は
$x_1x_2,x_1x_3,zx_0x_3,x_0x_2$ に $P$ を掛けたものである。
よって片方の座標対は零である。$x_2=x_3=0$ の場合は $x_0x_1=0$ により座標頂点になるが、
$D\ne0$ で両頂点の評価は階数4である。もう一方なら $X_2=zX_3$ となる。

$Q\ne0$ の場合、$x_3=0,x_2\ne0$ なら小行列式から
$(a+b)x_0=(a+c)x_0=(c+d)x_1=(b+d)x_1=0$ となる。
一組がともに零なら正表示の八直線の左側02または12であり、除外済みである。
よって $x_0=x_1=0$ となり $m_{1245}=-Q(b(X_0-X_1)+X_2-X_3)$ に反する。
$x_2=0,x_3\ne0$ でも対応する四係数は $a+bz,az+c,c+dz,bz+d$ であり、
八直線の左側03、13を除くと同じ矛盾を得る。
残りの $x_2=x_3=0$ では $x_0x_1P=0$ と座標頂点の排除から $X_0=X_1$ である。
得た二対の点はいずれも実際の零点で、一方では $Q\ne0$、他方では $P\ne0$ であり、$E$ と交わらない。
\path{computations/verify_negative_smooth_pairing.py} に小行列式と両側の対角鉛筆のPlücker座標を与えた。

滑らかな場合の枠は前節の四頂点である。二次曲線二本なら鉛筆には階数2の曲面一つと
階数3の曲面二つがある。前者の頂点直線は重なった固有値の二頂点を結び、後二者の頂点は
残る固有値の各一頂点である。四つの孤立点の六割線のうち、後二頂点のどちらかを含むのは
座標直線 $01,23$ だけである。この二直線と頂点直線の交点が残りの二頂点である。
四角形なら相対する節点の二割線が重複固有値の二直線である。
孤立点の六割線のうちその両方と交わるのは $01,23$ だけである。実際、異なる組の
孤立点を結ぶ割線では、各組に残る一座標ずつを同時に零にできない。
四つの交点から元の四頂点を取り戻す。右と次辺については以下の $J$ の恒等式が
$qt\ne0$ と正則性を保つので同じである。この恒等式の検証は頂点復元を使わない。
\end{proof}

\begin{proof}[全置換と薄い係数の条件]
非零台には、上位ビットが0の文字と1の文字が二つずつある48語が常に残る。
四文字の各置換は $\mathbb F_2^2$ 上アフィンである。総和零の64語のうち48語で
変換後の総和が零なら、全64語で零である（非恒等なアフィン線型条件の解は高々32語）。
従って置換は四位置で共通の線型部 $\phi$ と総和零の移動量を持つ。
各隣接切断の厚い色が $0,1$ なので、$\phi$ は恒等または $2,3$ の交換で、
四移動量の上位ビットは等しくなければならない。各位置移動につき32候補である。

48語の係数は $\pm z^e$ であり、$z=w_{0022}w_{1133}/(w_{0033}w_{1122})$。
候補の変換後のこの比 $z'$ は $z$ または $1/z$ である。元の表示を $z'$ とした薄い係数と
候補側の薄い係数の比について、文字倍率行列の飽和整数核で符号・指数を比較する。
移動が $k$、各文字移動が $u_i$ のとき $u_0+u_2=k\bmod2$ の16候補では全て恒等式、
他の16候補の差のLaurent多項式の最大公約数は $z-1$ である。
薄い係数を固定した後、左上 $A$ の四成分の倍率は等しい。対応する四語の任意の二つの
文字指数の差は薄い48語の整数行格子に入る。従ってこの主張は $A$ に零成分があっても成り立つ。
共通倍率は全語に対する一体の変更で実現できる。各語で上位ビットの和を $n$ とすると
共通倍率 $v$ に対する因子は $v^{1-n/2}$ である。

指定始点で残る置換は、移動 $(0,1,0,1),(1,0,1,0)$、線型部の交換、および
この交換に移動 $(2,2,2,2)$ を加えたものにより生成される。これが順に $U,V,S,B$ である。
一位置移動にはさらに $(0,0,1,1)$ の文字移動を選ぶと $J$ となる。
これら五変換に対して実際の全64係数の比は、符号と $z,z-1,D$ の整数冪であり、
その全ての文字格子関係が成立する。よって開条件上の各点で対角倍率が存在し、
係数が消える場合にも単元を掛けた恒等式を特殊化できる。
\path{computations/verify_negative_component_quotient.py} が五つの作用と薄い場合分けを検証する。
\end{proof}

\begin{proof}[周期、安定群、有理商]
$US,VS$ は $(q,t)$ の二つの符号を独立に変え $(p,r)$ を保つ。
$B$ の最初の行列を $B_0$、$J$ の最初の行列を $J_0$ とする。
$B_0^2=-z(z-1)^2 I$, $J_0^2=zI$, $B_0J_0=-J_0B_0$ である。
従って $B^2$ は射影的に $(q,t)$ の同時符号変更で、$B$ は二つの単一符号変更を入れ替える。
これで $z$ を固定する八元の二面体群 $H_0$ を得る。$U$ は $z$ を反転し、$H$ の全16元を尽くす。
$J^2$ は二つの符号交換の積を含む $H_0$ の元である
（直接には $(p,r,q,t)\mapsto(p,r,-q,-t)$）。従って常に周期2以下である。

周期1なら $Jx=h x$, $h\in H_0$ が射影的に成立する。
$h$ が符号変更だけの場合、$(p,r)\ne(0,0)$ なら $r^2=zp^2$、さらに $qt\ne0$ から
$t^2=zq^2$ となる。これは $4D=p^2-q^2-(r^2-t^2)/z\ne0$ に反する。
$(p,r)=(0,0)$ は二平面の共通部分である。
残る四元では $J_0(p,r)=cB_0(p,r)$ と書き、後の二座標を比較する。
$qt\ne0$ により $c=\pm1/(z-1)$ で、二つの符号は互いに逆である。
各符号で最初の二式は $r=-zp$ または $r=-p$ と同値で、逆に各平面で必要な二符号を選べる。
二平面は $B$ で入れ替わる。各平面には正則点が稠密である
（除外した八直線と $qtD=0$ はいずれも平面全体を含まない）ので商上の一因子である。

安定群も全点で数えられる。台の文字指数行列は薄い48行と左上一語だけで
Smith因子 $1^{11},2,2$ を持ち、全64語の行格子に等しい。
$D\ne0$ によりそのような左上一語があり、どの左上一語を選んでも同じである。
従って射影安定群の位数は $4|\operatorname{Stab}_{H_0}(x)|$ である。$(p,r)\ne(0,0)$
では後の因子は通常1で、$t^2=-zq^2$ かつ $(p,r)$ が $B_0$ の固有線にあるときだけ2になる。
実際、単一の符号変更では固定できず、$B$ を含む元の二符号はその固有値によって決まる。
$(p,r)=(0,0)$ では同時符号変更も効かず、後の因子は2、ただし $t^2=-zq^2$ なら4である。
$t^2=zq^2$ はここでは $D=0$ で除かれている。群そのものは該当する置換の文字倍率の解を合成すればよい。

最後に関数体を直接求める。二つの独立な符号と $U$ の不変体を
\[
 Z=(z+1)^2/z,\quad x=r/((z+1)p),\quad u=q^2/p^2,\quad v=t^2/(zp^2)
\]
が生成する。実際、先に符号を除いた後は $z$ を戻すだけで元の体に戻る。
$B$ は $x'=-(1+2x)/(2+Zx)$, $\delta=(Z-4)/(2+Zx)^2$ として
$(x,u,v)\mapsto(x',\delta v,\delta u)$ とする。従って
\[
  Z,\quad x+x',\quad (1+\delta)(u+v),\quad
  (1+\delta)(u-v)/(x-x')
\]
は不変である。$x+x'$ は $x$ の真の二次有理式なので、この四元に $x$ を戻せば
$u,v$ まで解け、拡大次数は2である。これで四つの独立な生成元を得た。
\end{proof}

% !TeX root = ../papers/nc-p3.tex
\section{負交換表示の全境界と二次曲面グラフへの特殊化}
\label{sec:negative-pairing-boundary}
前節と同じ $zD(z^2-1)\ne0$ を課し、$[a:b:c:d]$ は八例外直線を避ける。
\begin{theorem}\label{thm:negative-pairing-boundary}
この負交換表示の全正則点で、指定始点の同型の必要十分条件は
定理\ref{thm:negative-pairing-unmarked-quotient}の $H$ の軌道の一致である。
$q,t$ のうちちょうど一方が零なら、最小周期は $r=-p$ または $r=-zp$ で1、それ以外で2である。
両方が零なら、$(a+b)(a+zb)=0$ でちょうど二次曲面グラフとなり周期1、他は四直線のグラフを台に持ち周期2である。
この二次曲面条件も $r=-p$ または $r=-zp$ と同値である。
\end{theorem}
\begin{proof}[一方だけが零の場合の枠]
$B$ で $q,t$ を入れ替えられるので $d=a$, $t\ne0$ とする。前節の二次式 $P,Q$ を用い、
$\alpha=(a+b)(a+c)$, $\beta=(a+bz)(az+c)$ と置く。
$D(X_0-X_1)=(2a+b+c)X_2-(2az+bz^2+c)X_3$ と $\alpha X_2=\beta X_3$ が
$P=Q=0$ と同値である。$D\ne0$, $t\ne0$ により $\alpha,\beta$ は同時には消えない。
$X_0=X_1$ とすれば $Q=0$ が $X_2=zX_3$、$P=0$ が $(z-1)t X_3=0$ を与える。
従って $\alpha\beta\ne0$ なら二つの平面二次曲線が二点
$e_0\pm e_1$ で会い、それ以外なら平面二次曲線の二重化である。後者の被約二次曲線も滑らかである。

二次式はこの境界でも実際の小行列式を消す。具体的に、各小行列式は $P,Q$ の二次式係数結合である。
上位ビット別に混ざる項は前節の表示を使い、平方座標だけの五式は $P$ で $X_0$ を消去して $Q$ で割る。
$d=a$ では五つの商は全て $X_2,X_3$ の線型結合で分母は $D$ の冪だけである。
$Q\ne0$ と $Q=0,P\ne0$ の前節の消去は $q\ne0$ を使わず、
今度は $x_2=x_3=0$, $X_0=X_1$ の二点は曲線上に来る。孤立点は $x_0=x_1=0$, $X_2=zX_3$ の二点だけである。

二重二次曲線の場合にも $e_0\pm e_1$ は固有である。曲線の台で $x_2$ または $x_3$ が非零な点では、
零でない方の位置を使う前節の混合小行列式により、局所的にまず $Q$、続いて $P$ が生成される。
例えば $x_3=0,x_2\ne0$ の点では $(a+b,a+c)$ も $(c+d,b+d)$ も同時零を除いてあり、
$x_0,x_1$ の一方は非零である。この点において $P,Q$ の係数行列の階数が2だから
近傍のイデアルは $(P,Q)$ である。$x_2=0$ の場合は $z$ 付きの四係数で同様である。
二点 $e_0\pm e_1$ では $x_0x_1$ が単元で $P$ が生成される。
$P$ を法として全小行列式は $Q$ の極大イデアル倍であり、平方座標だけの五式は上の除算から分かる。
一方、$P$ による局所商は正則、$Q$ の像は非零である。
従って点スキームから純一次元部分を取る写像の核がこの両点で非零となる。
純部分は $(P,Q)$ そのものである。これは完全交叉で埋込み点がなく、既に二点以外で一致するためである。
よってこの二点は非Cohen--Macaulayな点として復元できる。

二つの二次曲線なら二つの節点と、二つの階数3の鉛筆要素の頂点を取る。
二重二次曲線なら上の二点を取り、二つの孤立点の割線と被約二次曲線の平面の交点を第三点、
その点の孤立二点に関する調和共役を第四点とする。いずれも得られる二組は
$\{e_0+e_1,e_0-e_1\}$ と $\{e_2,e_3\}$ であり、二組の役割は異なる。
奇位置では両組は $\{e_2+e_3,e_2-e_3\}$ と $\{e_0,e_1\}$ である。
実際、$BJ$ は $q=0$ に戻す表示で、位置の上位ビットを交換する。二節点側の元の座標直線では
指数の符号変更が二点を交換する。四位置で示した候補を六式に代入するといずれも評価核を持つので、
この直線と曲線の二交点が上の二点に等しい。孤立二点の座標も
\path{computations/verify_negative_boundary_vertices.py} で元の四辺から検算する。
\end{proof}

\begin{proof}[一方だけが零の場合の商]
復元した点を列とする基底でテンソルを書く。最大台58語、うち42語は
$z$ と非零定数だけで決まる。各位置で二組の各々を交換する $4^4$ 候補のうち、
この42語を最大台へ送るものは16通りだけで、どれも42語自体を保つ。
42語の整数行格子は階数11で、残りの左上係数 $2a,2b,2c$ の任意の一語を加えると、
全58語と同じ階数12の原始格子となる。三語の各対の指数差も42語の格子に入る。
よって $z$ と $[a:b:c]$ が順序付き枠から読め、射影対角安定群は一次元トーラスである。
より詳しくは四組の文字指数の核で42語の係数単項式から $z$ を得て、
三語の各比への補正にも42語の単元だけを用いる。従って射影的な三係数の比較は零成分へも延びる。

16通りのうち8通りは目標で $z'b'=c'$ を要求し、この層には戻らない。
この検査で $b'/c'$ を使った場合、分母が零になる点では分子も同じ元の係数に対応し、
両方とも零である。残りは $U,S$ の生成する四通りの $[a:b:c]$ の変換で、各々が二度現れる。
四通りは前節の実際の同型で、同じパラメータ作用を持つ二置換の差は、
各位置で節点対を交換する元の対角四元群の文字作用で実現する。
\path{computations/verify_negative_pairing_boundary.py} が全候補を照合する。
従って一般の射影安定群の連結成分数は2である。$z$ を保つパラメータ作用は
$[a:b:c]\mapsto[a:c/z:zb]$ だけなので、$a=0,r=0$ なら成分数4、他では2である。
$BJ$ は $[p:r:t]\mapsto[(z+1)p+2r:-2zp-(z+1)r:(z-1)t]$。
固定した $z$ で使える同型は $t$ の符号変更だけなので、最初の二成分の行列の
二固有線 $r=-p,-zp$ がちょうど周期1の条件となる。$(p,r)=(0,0)$ も両式を満たす。
\end{proof}

\begin{proof}[両方が零の場合]
$d=a$, $c=zb$ とする。この面では $D\ne0$, $z\ne0,\pm1$ のもとで八例外直線のどれにも乗らない。
$(a+b)(a+zb)\ne0$ なら $P,Q$ は $X_0-X_1,X_2-zX_3$ を張り、
前節の曲線外の点の計算から射影点スキームの台は四直線の四角形になる。
四つの節点は線型独立で、向かい合う組を入れ替えることができる。
$w^2=z$, $m=a/(wb)\in\mathbb P^1$, $X=(w+1)^2/(w-1)^2$, $Y=(m+1)/(m-1)$ と置く。
$X\ne0,\pm1,\infty$、$Y\ne0,\infty$ である。$b=0$ では $Y=1$ とする。
節点枠は偶位置で $e_0\pm e_1,we_2\pm e_3$、位置1で $we_0\pm e_1,e_2\pm e_3$、
位置3で $e_0\pm we_1,e_2\pm e_3$ である。これは上と同様に元の四辺への評価で確認できる。
この枠の28個の係数は全て非零で、文字倍率行列のSmith因子は $1^{11}$ である。
$b=0$ でも厚い左上を $b$ 倍、右下を $b^{-1}$ 倍して斉次座標に戻すと非零性は保たれる。
トーラス不変な単項式から $X-1,X,Y$ を読み取れる。
四位置で四角形の二面体置換を全て試すと、28語の台に戻るのは8通りで、
$(X,Y)$ には次の四通りが各二回現れる。
\[
 (X,Y),\quad(1/X,1/Y),\quad(X,-1/(XY)),\quad(1/X,-XY).
\]
これらが十分であることも確認する。$X,Y$ が等しい二表示では $w'=w$ または $1/w$ であり、
同じ $m$ である。後者は $U$ の作用で得られる。第二変換は平方根 $w$ の符号変更で、
第三変換は $B$ の作用である。同じ作用を持つ二置換は全位置の両対の入替えで異なり、
元の文字の下位ビットの符号により実現できる。$J$ は $(X,-Y)$ を与えるから
周期1の可能性は $XY^2=1$、すなわち先に除いた $(a+b)(a+zb)=0$ に限る。
この四角形の射影安定群は二次元トーラスを単位成分に持ち、成分数は $XY^2=-1$ で4、その他で2である。

残る二線は $B$ によって交換される。$a=1,b=-1,c=-z,d=1$ を取ると
$Q=x_0^2-x_1^2-x_2^2+zx_3^2$ で、第二位置を
$(x_0,x_1,x_2,x_3)\mapsto(-x_2,x_3,x_0,-x_1)$ により第一位置と同定できる。
六式の対称部分は $Q$ を張り、交代部分にある五平面の零化二ベクトルは
$\Omega=e_0\wedge e_1+e_2\wedge e_3$ である。非退化性により辺はスキームとして
滑らかな $Q$ の対角グラフである。各辺の基底を $B,J$ で戻すと同じ結論が得られる
（$J$ は $m\mapsto1/m$ で二線を交換する）。定理\ref{thm:smooth-quadric-graphs}の二平面の
面積形式の固有値比の平方は $z$ または $1/z$ である。実際、$Q$ の対角行列と
$\Omega$ の交代行列の積の特性多項式は $(h^2-1)(h^2-z)$ である。
従って正確な無印座標は $z+1/z$、周期は1である。ここでも同じ座標を持つ全点は
$B,U$ で同一視できる。この族で滑らかな二次曲面枝の座標 $-2$ 以外を全て実現するため、
その一変数枝の全体が負の四次元成分の閉包に含まれる。
以上の各形は射影点スキームの台の次元と孤立点数 $4,2,0$ で区別でき、層をまたぐ同型も尽きた。
\end{proof}

\begin{corollary}\label{cor:negative-exchange-opposite-ranks}
負交換表示の全域で反対階数は $16,15,12$ のいずれかである。階数15は商上の二因子で、
いずれも最小周期2であり、一位置移動で交換される。これらの接する階数12の全点は周期1で、
その指定始点の商は
\[
 (\mathbb A^1\setminus\{-2,2\})\times(\mathbb P^1\setminus\{2\})
\]
となる。二つの周期1平面は階数12に沿って交わる。
\end{corollary}
\begin{proof}
色0の反対ブロックの列を各二組の和・差に取ると、二行正方ブロック二つの直和になる。
一方は列 $(a,-d)^t,(1,1)^t$、他方は $(1,z)^t,((a+b+c+d)/D,-(az+bz^2+c+dz)/D)^t$ を持つ
（列の定数倍を除く）。色1でも二組の行の和・差で同じ分解があり、
最初の二列 $b,-c$ と $1,z$ の組および右下ブロックから同様に計算できる。
どの二行ブロックも零にならない。色2,3は両方とも行列式 $-4(z+1)^2$ である。
従って正確な階数は
\[
16-[p=0]-[r=0]-[2zp+(z+1)r=0]-[(z+1)p+2r=0].
\]
固定 $z$ で四平面は相異なり、共通の直線 $p=r=0$ を持つ。$B$ の作用は
$p=0$ と $2zp+(z+1)r=0$、$r=0$ と $(z+1)p+2r=0$ をそれぞれ交換し、$J$ は二組を交換する。
$z\ne0,\pm1$ では周期1平面がこれらに重ならない。

$p=r=0$ なら $d=-a,c=-zb$ で、八例外直線は $D\ne0$ の範囲でこの軌跡に会わない。
$h=t^2/(zq^2)\in\mathbb P^1$ とすれば $D\ne0$ は $h\ne1$ と同値である。
符号を除いた商で $B$ は $h\mapsto1/h$、$U$ は $h$ を保ち $z\mapsto1/z$ とする。
平方根の選択も符号変更で尽くされるから、商座標は $z+1/z$ と $h+1/h$ である。
$h=0,\infty$ は一方だけの対角差が消える境界から加わり、第二座標の無限遠の一点を与える。
\end{proof}

% !TeX root = ../papers/nc-p3.tex
\section{全四元群両対称型の評価判定と負交換の端点}
\label{sec:negative-exchange-end-values}
\begin{theorem}\label{thm:all-klein-bisymmetric-evaluation}
定理\ref{thm:klein-bisymmetric-irreducible}の全固定色両対称型において、
全次数の標準Koszul左右AS正則性は次の三条件と同値である。
反対階数が6または8であること、辺の六次元関係空間が非零の純テンソルを含まないこと、
左右の全非零部分評価が階数3以上であることである。捻りの場合は逆捻り後に同じ条件を使う。
この判定には薄い係数、角ブロック、点配置の開条件がない。
\end{theorem}
\begin{proof}
定理\ref{thm:bisymmetric-six-or-eight}、系\ref{cor:one-smooth-apolar-as}の必要条件を使う。
反対切断を $\mathrm{Sym}^2 E$ からの写像として読むと、階数6では核が四次元である。
対称四行行列の階数2以下の射影多様体は次元6なので、核の三次元射影空間がこれと交わる。
従って核に二つの一次式の積が入り、補題\ref{lem:bisymmetric-reducible-quick-test}が使える。
階数8では核は二次元である。四文字の平方の像は総文字1により、相手側の $01,23$ の二積しか使わない。
従って核の全体が四つの平方からなる鉛筆である。四列のいずれか二つを同時に消す鉛筆要素があれば
その要素は階数2以下、なければ四列はいずれも非零で相異なる向きを持ち、基底曲線は滑らかである。
後者には系\ref{cor:one-smooth-apolar-as}、前者には同じ補題を用いる。
\end{proof}

\begin{theorem}\label{thm:negative-exchange-minus-one}
$D\ne0,z=-1$ の負交換表示の指定始点の正則商は重み付き射影空間
$[A:B:C:L]\in\mathbb P(1,1,1,2)$ の次の開集合である。
\[
 A+B-C\ne0,\qquad
 (B,(A-C)^2-4L)\ne(0,0),\qquad
 (A,(B-C)^2-4L)\ne(0,0).
\]
商写像は $A=p^2,B=r^2,C=q^2+t^2,L=q^2t^2$ で、一位置移動は $A,B$ を交換する。
従って最小周期は $A=B$ で1、他で2である。
反対階数は一般に12、$A,B$ の片方だけが零なら10、両方零なら8である。
階数8の全点は周期1で、その商は $L/C^2\in\mathbb A^1\setminus\{1/4\}$ である。
この端点と $z\ne0,\pm1$ の正則交換点には相互に同型がない。
\end{theorem}
\begin{proof}
正の交換に戻して八直線による除外を示す。純テンソルは角二ブロックの行列式から
$z\ne1$ だけで除かれる。反対核の対角鉛筆が滑らかまたは可約元を持つことも以前と同様である。
正確には階数が8に達しなくても、平方の四行二列ブロックの行対が比例すれば可約元が生じ、
どの行対も独立なら他の二ブロックの階数2が $z\ne1$ から従い、全階数は8である。
従って評価階数だけを調べればよい。全四座標が非零の候補を消す行列式の符号の矛盾は
$\det K=-\det L=(z-1)/D\ne0$ だけを使う。全低二座標または全高二座標が零なら
厚い評価の階数2に少なくとも一つの薄い式が加わる。低い側または高い側で一座標だけが零なら、
核が二次元のとき他側にも零座標を要することは、定理\ref{thm:exceptional-exchange-eight-lines}と同じ
可逆な二行列の列比較から従う。従って八つの二座標の場合を調べればよい。
座標を $1,s$ とし $sD$ を可逆化して、六行四列の評価の三次小行列式のイデアルを比較すると、
$z=-1$ では八つ全てが同定理の直線と各々の $s^2$ の式になる。
\path{computations/verify_exchange_end_values.py} は等式を有理数上で両方向に検証する。

前二節の頂点復元・軌道証明は、$z=-1$ でも成り立つ。二次式の行列式に使った非零因子は $D,z,z-1$ と
明示した $q,t$、二重二次曲線の除法には $D$ を用いた。二組の座標の全置換を合わせる48語の検査で
追加解を許す因子は $z-1$ だけである。一方の対角差が零の枠の42定数語も全て非零で、
他の八置換が $z'b'=c'$ を強制する条件も同じである。両方零なら $w^2=-1$ で $w\pm1\ne0$、
四角形の28語の比較も使える。このとき $X=-1$ で二つの有理変換の作用が重なるが、
等しい $X,Y$ を持つ表示が $U$ により同型になる逆向きの証明は同じである。
二次曲面では既知の座標は $-2$ となり、この点も先の負成分上の正則閉点として実現する。
いずれの形でも、前節の $z$ または $1/z$、四角形での $X$ または $1/X$、
二次曲面の座標のいずれかが、従来の交換表示との同型を区別する。

$z=-1$ なら $U$ は $r$ だけの符号変更、$US,VS$ は $q,t$ の符号変更、$B$ はこれらを除けば
$q,t$ の交換である。従って射影作用は四座標の任意の符号変更と $q,t$ の交換の全体である。
不変環の射影商が表示の四座標となる。$4D=A+B-C$ であり、八直線のうち四本は
$r=0$, $p\pm q\pm t=0$、他の四本は $p=0$, $r\pm q\pm t=0$ となるから除外式を得る。
$J$ の固定には $A=B$ のほか射影倍率 $-1$ が考えられるが、そのとき $A+B=C=0$ で除かれる。
反対ブロックの色2,3はこの繊維で階数2となり、色0の二小行列はともに $p$、
色1の二小行列はともに $r$ の定数倍である。これが反対階数の式を与える。
\end{proof}

\begin{proposition}\label{prop:zero-exchange-vertical-identification}
$D\ne0,z=0$ の負交換窓は第二長方形の負垂直族と同じ表示である。
必要十分な正則条件は、$(a',b',c',d')=(a,-b,-c,d)$ が
\[
 (a'=b',c'=b'),\quad(d'=b',c'=b'),\quad
 (a'=b',d'=0),\quad(a'=0,d'=b')
\]
の四直線を避けることである。$z=1,D\ne0$ の交換表示には正則点がない。
\end{proposition}
\begin{proof}
$z=0$ でも正の窓の純関係の排除と二座標への帰着に必要な厚い二行列は可逆である。
四行二列の行そのものが零になった場合も、反対核に平方が入る。
八評価イデアルのうち左02,12と右02,03だけに解があり、順に上の四直線と
$s^2=b'$ が得られる。$sD\ne0$ の局所化での計算は同じスクリプトに含めた。
$z=1$ では正表示の $x=(0,0,1,1)$ の部分評価が階数2で、捻りも非正則である。

$z=0$ の負表示の奇位置の文字に1を加え、目標の四位置で順に
$(-,+,+,-),(+,-,+,+),(+,+,-,+),(-,+,+,+)$ の符号を掛けると、
定理\ref{thm:parallel-vertical-orbits}の $t=-1,b_{\rm rect}=c_{\rm rect}=1$、
$(A,B,C,D_{\rm rect})=(a,c,-d,-b)$ のテンソルと全256係数で一致する。
角ブロックの条件は元の $ad-bc\ne0$ に等しい。この恒等式は $B=0$ でも有効である。
従って $B\ne0$ の正則な八代表台の同型と周期には既存の長方形商を適用できる。
$B=0$ なら $a,d\ne0$ により上の四直線のどれにも乗らず、全点が正則である。
\end{proof}

% !TeX root = ../papers/nc-p3.tex
\section{第二長方形の負垂直族の射影的な全商と二次錐}
\label{sec:zero-exchange-projective-quotient}
\begin{theorem}\label{thm:zero-exchange-projective-quotient}
$z=0,D=ad-bc\ne0$ の前節の負交換表示の全正則点では、
$[a:b:c:d]$ と $[d:b:c:a]$ だけが指定始点で同じ類を与える。
従って射影座標 $[p:b:c:u]$（$p=a+d,u=(a-d)^2$）で、全商は
$\mathbb P(1,1,1,2)$ の
\[
 p^2-u-4bc\ne0,\qquad (c-b,u-(p+2b)^2)\ne(0,0),\qquad
 (p+b,u-p^2)\ne(0,0)
\]
なる開集合である。周期1の軌跡は $p+c=0$ と $c=2b+p=0$ の和、それ以外の最小周期は2である。
反対階数は $16-[p=0]-2[c=0]-[p+2c=0]$ である。
この族の正則点と先の $z\ne0$ の負交換表示の間に指定始点の同型はない。
\end{theorem}
\begin{proof}[二つの射影変換と $c\ne0$]
命題\ref{prop:zero-exchange-vertical-identification}の四直線を
$a,d$ の交換で割ると記載の除外式になる。実際のテンソルで
\[
 I:[d:b:c:a],\qquad J:[-a-c:b+a+c+d:c:-d-c]
\]
がそれぞれ指定始点と一位置移動の変換である。各々で共通の四文字置換をした後
（位置ごとに別の置換を許す）、全非零係数の比は $\pm D^j$, $j=0,\pm1$ となる。
全整数文字関係を検査すると対角倍率が存在し、$D$ が単元なので
厚い係数の消滅を含めて恒等式を特殊化できる。
$c\ne0$ では第二長方形の $B=1$ の座標は $(-a/c,-d/c,b/c)$ である。
定理\ref{thm:parallel-vertical-orbits}の台代表への移動を $I,J$ で記録すれば、
指定始点の比較は $I$ だけで尽きる。代表の被覆は、先の座標を $(v,y,h)$ として
$v,y,1-v,1-y,h,v+y-1-h$ の零条件で確かめられる。
$I$ は $v,y$ を交換し、$J$ は $v,y$ を $1-v,1-y$ に、$h$ を $h+1-v-y$ に移す。
$v,y$ のうち0または1になる個数が0なら、後の二式の消滅数に従い147,145,142へ行く。
一個ならその値を $v=1$ にして139、または $h$ の特別値で134に行く
（$h=v+y-1$ は $vy=h$ により不可）。二個なら124,118,129のいずれかである。
ここで $v=1$ の規格化前に $J$ の適用を要した場合も、添字移動の回数を保存して
同定理の奇・偶両方の置換比較を使う。
\end{proof}

\begin{proof}[無限遠 $c=0$ の点と純曲線]
$ad\ne0$ である。左の点集合を前節の $P,Q$ を用いて調べる。
$a\ne d$ の場合、純一次元部分は
\[
 x_0^2=-\frac{b+d}{a(a-d)}x_2^2,\qquad
 x_1^2=-\frac{a+b}{d(a-d)}x_2^2
\]
の錐で、頂点は $e_3$、残る孤立点は $e_0\pm e_1$ の二点である。
二定数は同時に消えない。両方非零なら相異なる四直線、一方だけ零なら二本の二重直線となる。
実際、混合小行列式で $x_2x_3Q=x_0x_1P=0$ が得られる。
$Q\ne0$ で $x_2,x_3$ の一方が非零の場合を除く計算には、左側四例外直線に乗らないことを使う。
$Q=0,P\ne0$ のとき四つの混合積 $x_ix_jP$ ($i<2\le j$) は全て零である
（$z=0$ でも $m_{1235}$ が $-x_0x_3P$ を与える）。
前節の消去を辿ると、両式の外で可能なのは上の二孤立点だけである。
二式で定まるスキームに全小行列が消え、各直線の一般点では $x_2x_3\ne0$ から $Q$ が、
次に非零な $x_0,x_1$ の一方を含む混合式から $P$ が生成される。
従って埋込み点を除くと上の完全交叉そのものである。

頂点からの射影平面で四つの直線の方向が相異なれば、四方向の完全四辺形の三対角点が
座標点 $e_0,e_1,e_2$ である。孤立二点の像を結ぶ直線は第一・第二点だけを含む。
二重直線なら、方向スキームは二つの長さ2点である。その二接線の交点が消えた方の座標の頂点であり、
二支持の割線と孤立点の割線との交点がもう一方の低位の頂点となる。
二支持の割線上で後者の調和共役を取れば $e_2$ の像を得る。
元の空間で孤立点の割線に戻れば $e_0,e_1$ は交換を除いて決まる。$e_2$ の持上げだけには $e_3$ を加えられる。

$d=a$, $b\ne-a$ なら $P,Q$ は $x_0^2-x_1^2,x_2^2$ を張り、純曲線は二つの二重直線、
孤立点は存在しない。直線の一般点で $x_0x_1\ne0$ より $P$ を得て、$x_0x_3$ を含む式から
$Q$ を得る。台上で頂点以外に埋込みのある点は $e_0\pm e_1$ だけである。
前節の $P$ を法とする五式の除算でもこの両点上で $Q$ は生成されず、
両点と頂点以外では上の同じ二つの単元で $P,Q$ が生成される。
二直線の交点 $e_3$ と二つの埋込み点を枠に取る。残る $e_2$ の像を規定するには
節点枠の二座標 $v_0,v_1$ を使う。純理想は $(v_0v_1,x_2^2)$ であり、
この三点を保つ変換が $v_i$ に $x_2$ の倍を加えると積に二つの独立な項 $v_j x_2$ が生じる。
従って $e_2$ の持上げに許せるのはやはり $e_3$ の加法に限る。
奇位置では前記 $J$ の明示単項同型で同じ結果を使える。
四頂点の順の錐頂点は $(3,1,3,0)$、持上げを変えられる文字は $(2,0,2,1)$ である。
$a=d$ の場合の残りの二線は偶位置で $e_0\pm e_1$、奇位置で $e_2\pm e_3$ である。
直線の方向と頂点座標の零を保つ単項移動、または六式への直接代入で確かめられる。
\end{proof}

\begin{proof}[加法を許す枠の比較]
各位置の表示は、二つの固定文字の交換、文字倍率、持上げ文字に頂点の倍を加える四係数からなる。
四位置の頂点文字を一つ含む語が28個あり、二つ以上含む語はない。
従って持上げの変更は頂点語を固定し、頂点を含まない語に四列の行列の像を加える。
どちらの枠でも、全ての強制非零頂点語を許容台に送る交換は、恒等と四位置同時の交換だけである。
$a\ne d$ の場合には四列の台が元のテンソルの台と交わらず、指定始点の比較は
対角倍率と二つの置換に帰着する。36定数語と $a,b,d$ の各一語による格子差で、
$[a:b:d]$ が $I$ を除いて一意である。加法の四列は階数3で、その核が一次元加法群となる。
$a=d$ の場合は $a=1$ として、まず28語を固定する倍率のトーラスを取る。
そのSmith因子は $1^{11}$ で、残る零頂点語をこのトーラスの文字で分ける。
四列の像と残るテンソルは全て同じ原始文字の19語の空間にある。
四列が張る三次元を割った後、残るテンソルは $v+bw$ の形で $v,w$ は独立である。
従って射影化しても $b$ は一意である。同時交換は元の文字符号によって実現する。
この場合には一次元加法群のほかに一次元の射影対角安定群が残る。
\path{computations/verify_zero_exchange_projective_flags.py} が語の全置換、加法の列とその商を検算する。

$d=a,b=-a,c=0$ はただ一類である。$a=1$ とすれば前節の二次曲面表示を $z=0$ としたもので、
第一辺を対角化する行列はなお可逆、交代五平面の零化二ベクトルもなお非退化である。
従ってスキームとして二次錐 $x_0^2-x_1^2-x_2^2=0$ の対角グラフであり、全辺について $J$ で従う。
定理\ref{thm:quadric-cone-quotients}の余分な部分のない唯一の交代型に等しい。
これで滑らかな二次曲面枝の $z+1/z\to\infty$ における正則な特殊化も決まる。

$c\ne0$ の場合との重複は孤立点数で排除できる。$c\ne0$, $a\ne d$ なら
$P=Q=0$ の外に $e_0\pm e_1$ と $e_3$ の三点、$a=d$ なら $e_3$ だけが来る。
後者でも $\alpha=(a+b)(a+c),\ \beta=ac$ は同時に零でなく、二次式は完全交叉をなす。
小行列による外部の消去は上と同様である。
$z\ne0$ の表示とは三点・一点の孤立点数、または $c=0$ の場合の錐状曲線の交点・多重度や曲面の階数で区別できる。
商上の $J$ は $[p:b:c:u]\mapsto[-p-2c:b+p+c:c:u]$ だから、表示した周期1条件が従う。
反対階数は色0,1の二つずつの二行ブロックと、色2,3の行列式 $-4$ から得られる。
\end{proof}

% !TeX root = ../papers/nc-p3.tex
\section{薄い係数が非零で交差ブロックの階数が一の両対称型}
\label{sec:klein-singular-exchange-corners}
\begin{theorem}\label{thm:klein-singular-exchange-corners}
定理\ref{thm:klein-bisymmetric-irreducible}の固定色で $a_ib_j\ne0$ が全ての $i,j$ に対して成り立ち、
$\det U=0$ とする。正則なものは全て、文字の付け替えを許して
\[
 M_0=\begin{pmatrix}X&X&1&1\\Y&Y&1&1\\1&1&W&W\\1&1&V&V\end{pmatrix},
 \qquad M_1=M_0^{\mathsf t},\qquad M_2=M_3=\boldsymbol1_{4\times4},
 \qquad (X-Y)(W-V)\ne0
\]
で与えられる。この全域で正表示と非自明捻りの双方が標準Koszul左右AS正則で、最小周期は1である。
\end{theorem}
\begin{proof}
偶奇の文字倍率を $e_i,f_i$ とすると、薄い因子 $a_i,b_i$ への作用は
$e_if_{i+2},e_if_{i+3}$ である（文字添字の加法は四元群）。因子の共通逆倍率を選んで
$a_0a_1=b_0b_1$ とでき、$\det U=0$ より $a_2a_3=b_2b_3$ も成り立つ。
すると各組で $f_{i+2}/f_{i+3}=b_i/a_i$ を解いて $a_i=b_i=1$ にできる。
$J=\boldsymbol1_{2\times2}$ とおいて前の階数分解で $LD'=RC=J$ を得る。

$L,C$ が共に階数1なら、$L$ の列と $C$ の行は定数ベクトルの倍数になる。
関係の左評価を $x=(1,1,0,0)$ で取れば、厚い二色は低位側の $(1,1)$ の方向だけを、
薄い二色は高位側の $(1,1)$ の方向だけを与える。定理\ref{thm:all-klein-bisymmetric-evaluation}
に反する。$R,D'$ が共に階数1の場合は高位二文字で同様である。
$L,D'$ の両方または $R,C$ の両方が可逆であることは $J$ の階数から除かれる。
従って $L,R$ が共に可逆、または $C,D'$ が共に可逆である。後者は全ての奇文字を1だけ平行移動し、
$M_0$ と $M_1=M_0^{\mathsf t}$ を交換して前者に帰着する。
前者では $C,D'$ の各二列が一致し、$L,R$ の可逆性が二つの差の非消滅に等しい。

元の総文字1の表示で第一辺の関係空間を $S$ とする。対角色では
$(X,Y,1,1),(1,1,W,V)$、色1では対称積 $01,23$、残る二色では対称積の和
$02+13,03+12$ が基底となる。ここでは異なる二文字の積を対称化して記した。
両対称性から隣の関係は $S^{\mathsf t}=S$ であり、全辺が同じになる。
$tt^{\mathsf t}\in S$ で両半分が非零なら、交差ブロックは対角成分同士、
非対角成分同士が等しい階数1の行列だから $t_0^2=t_1^2,t_2^2=t_3^2$。
対角の二係数がともに零となり矛盾する。片方が零の場合も、対角二平面から各半分への射影が
同型なので非零解はない。従って可換二次双対の四式は共通射影零点を持たず、正則列である。
対称チャートの定理から全次数の結論を得る。
非自明捻りには同じ四文字次数で全辺の関係を捻ればよい。実際テンソルの四位置を平行移動して
総文字を零にした時の符号 $\prod_{i<j}\mu(q_i,q_j)$ の変化は、双指標性により
各位置の文字倍率の積である。従って同じ斉次代数のコサイクル捻りがindexed表示を与え、
周期1も保たれる。
\end{proof}

\begin{theorem}\label{thm:singular-exchange-negative-quotient}
上の捻りの印を忘れた指定始点の商は $\mathbb A^2\times\mathbb G_m$ である。座標は
\[
 u=(X+Y)/(X-Y),\quad v=(W+V)/(W-V),\quad h=(X-Y)(W-V),
 \qquad (s,\ell,H)=(u^2+v^2,uvh,h^2).
\]
反対階数は $16-2[u=0]-2[v=0]$、射影安定群の位数は通常4、$u=\pm v$ で8、
$u=v=0$ で16である。この三次元族の閉包は四次元負成分の周期1点の別の因子を与える。
\end{theorem}
\begin{proof}
$\mu(i,j)=(-1)^{i_{\rm hi}j_{\rm lo}}$ とする。捻りの点対 $(x,y)$ で $x$ の両半分が非零なら、
色1の式より $y=(m x_0,-m x_1,n x_2,n x_3)$。薄い二式は
$(m+n)(x_0x_2+x_1x_3)=(m+n)(x_0x_3+x_1x_2)=0$。
後ろの二因子が両方零なら $x_1=\pm x_0,x_3=\mp x_2$ であり、残りの二式から
$m(X-Y)=n(W-V)=0$ となる。従って $n=-m$ で、点を与える条件は
\[
 Xx_0^2-Yx_1^2=x_2^2-x_3^2,\qquad x_0^2-x_1^2=Wx_2^2-Vx_3^2 .
\]
低位・高位それぞれで二列が独立だから、鉛筆の四列の比例は異なる半分に属する列の間にしか起きない。
従って滑らかな四次曲線、異なる二平面二次曲線、または四直線の四角形を得る。
平方比例式を因数分解すれば後二者も被約である。片半分が零なら色1の式と対角二式で
同じ半分の $y$ は零になり、残る二式から孤立点 $(1,\pm1,0,0),(0,0,1,\pm1)$ を得る。
右側でも両半分の非零な部分は同じ符号変更のグラフで、孤立点は同じ二対である。
定理\ref{thm:negative-pairing-unmarked-quotient}の四頂点の復元が特異四次でも適用される。
すなわち二次曲線二本なら鉛筆の二単頂点を含む孤立点の割線を使い、
四角形なら節点の二対角線の両方と交わる孤立点の割線を使う。

四頂点の枠で総次数零に戻すと、薄い48語の係数は $\pm1$ である。同定できる置換は
前と同じ32候補であり、薄い係数を合わせた後の対角残余は
$(X,Y,W,V)\mapsto(kX,kY,k^{-1}W,k^{-1}V)$ である。
実際、厚い各半分内の二語の文字指数の差と異なる半分から取った二語の指数の和が
薄い48語の整数行格子に入る。逆に上位ビットの個数 $b$ による倍率 $k^{1-b/2}$ は実現する。
四位置の置換補正を薄い文字指数行列のSmith形で解き、残り十六語を比較すると、
八候補は両組内の独立な交換と二組の交換を与え、他の二十四候補には
$X=Y,W=V$ が必要である。この計算は厚い係数で割らずに行える。
一つの低位語を足した指数行列のSmith因子は $1^{11},2,2$ なので、
各実現置換につき射影対角安定群の位数は4である。
\path{computations/verify_exchange_singular_corners.py} は整数行格子、全置換補正および反対切断を検算する。
反対切断の二色は常に可逆で、他の各色は $X+Y,W+V$ の消滅ごとに階数を一つ失う。

共通倍率を除くと $(u,v,h)$ となり、各組内の交換は対応する $u$ または $v$ と $h$ の
符号を同時に変える。これら二交換の不変環は $u^2,v^2,uvh,h^{\pm2}$ で生成され、
$(uvh)^2=u^2v^2h^2$。二組の交換は $u,v$ を交換するので主張の商となる。
安定群の位数もこの作用から従う。負成分への包含は定理\ref{thm:klein-bisymmetric-irreducible}による。
以前の $zD(z^2-1)\ne0$ の表示のうち四孤立点を持つ部分とは、薄い交差比がここでは1、
先の薄い置換では $z^{\pm1}$ であるため一致しない。他の部分には四孤立点がない。
従ってこの既約三次元像は先の周期1因子の稠密開集合と交わらず、閉包も異なる。
\end{proof}

\begin{proposition}\label{prop:singular-corner-positive-forms}
正表示の印を忘れた同型には次の有限な正規形集合が使える。対称関係の直交空間を四次元の
対称行列空間 $L$ とし、階数2の二元を射影的に選んで各々の二つの線型因子を取る。
四因子が独立な全ての組を全ての順序で基底として、$L$ のPlücker座標の台と対角トーラスの
文字不変量を記録する。この有限集合が交わることが同型の必要十分条件である。
\end{proposition}
\begin{proof}
$L$ の行列は
\[
 \begin{pmatrix}A&0&r&t\\0&B&-t&-r\\r&-t&C&0\\t&-r&0&D\end{pmatrix},\qquad
 XA+YB+C+D=A+B+WC+VD=0 .
\]
$r=t=0$ の二次元空間で $(A,B)$ と $(C,D)$ はいずれも座標になる。
三次小行列には符号を除き $rt(A+B),rt(C+D)$ がある。$rt\ne0$ なら
階数2以下では全対角が零、$r=\pm t$ である。この二元はそれぞれ低位と高位の
二文字の和または差を因子とし、合わせて四因子は独立である。
$t=0,r\ne0$ なら $AC=BD=r^2$、$r=0,t\ne0$ なら $AD=BC=t^2$ である。
対角直線上で $AC-BD$ は恒等的に零ではない。恒等的なら独立な一次式 $A,B$
から $C=kB,D=kA$ が従い、第一式より $X=Y$ となる。$AD-BC$ も同様である。
$r=t=0$ の階数2点もどちらかの式に含まれる。従って各場合高々四点、先の二点と合わせて高々十点である。
階数1の元は、非零の非対角成分がある場合は二つのブロックの階数から、対角だけなら
両射影の独立性から除かれる。各枠での文字検査は、台上のPlücker単項式の重み（全体倍率の列も含む）
の整数関係の基底を取り、対応する単項式比を記録すればよい。系\ref{cor:symmetric-chart-isomorphism}
により同型は正確に $L$ の合同であり、合同は因子枠を運ぶので判定は必要十分となる。
同じ検査は任意の他の対称表示との比較にも使え、階数2点の数または独立枠の存在が異なれば直ちに分離する。
\end{proof}

% !TeX root = ../papers/nc-p3.tex
\section{二交換を仮定しない $z=1$ の混合補完}
\label{sec:klein-mixed-exchange-resolution}
節\ref{sec:klein-elliptic-completion}の規格化で $z=1$、
$\lambda\ne0,\pm1$ とし、隣接二切断の色別階数 $(2,2,1,1)$ を課す。
$P=\left(\begin{smallmatrix}0&1\\1&0\end{smallmatrix}\right)$ とする。
$X=X_\lambda,Y=Y_\lambda$ はともに可逆で $Y=\lambda PX$、$XP=PX$ である。
従って最初の二角 $A=A_0,B=A_1$ は必ず可逆であり、$T=YA^{-1}X,T'=YB^{-1}X$ を使える。
\begin{theorem}\label{thm:klein-mixed-exchange-resolution}
この繊維で標準正則となり得る全窓は、以下の表示とその二位置移動・文字置換で覆われる。
\[
 A=\begin{pmatrix}a&b\\b&d\end{pmatrix},\quad
 B=\begin{pmatrix}b&f\\g&b\end{pmatrix},\quad
 x=ad-fg\ne0,\quad t=\frac{a+d-(\lambda+1/\lambda)b}{ad-b^2}
   =\frac{(\lambda+1/\lambda)b-f-g}{b^2-fg}.
\]
この式と $\det A\det B\ne0$ は表示された側の二切断条件に必要十分である。
この側は共通倍率を除き有理既約四次元で、楕円曲線をパラメータに持つ五次元族の閉包に属する。
ここでの幾何学的一般点は標準Koszul左右AS正則、反対階数16、最小周期4、
指定始点の射影安定群は四元群である。印を忘れた像も四次元である。
\end{theorem}
\begin{proof}[全補完の分岐]
第二切断の二つの厚い行列 $N_0,N_1^{\mathsf t}$ の交差ブロックは全て $J=\boldsymbol1_{2\times2}$。
上の角を $C,D$ と書く。$C$ は $A$ の対角と $B$ の非対角、$D$ は $B$ の対角と
$A$ の非対角を使う。下の角はそれぞれ $T,T'$ の対応する成分を同じ順序で使う。
一つの厚い行列を $\binom{L}{R}(K\ S)$ と分解すると $LS=RK=J$ である。
上の角 $LK$ が特異なら $L$ の列が定数、または $K$ の行が定数である。
従ってその二行または二列が一致する。二列だけが一致する場合 $L$ は可逆なので
$S$ の二列も、下の角の二列も一致する。行についても同じである。

$C,D$ がともに特異な場合、片方の二列と他方の二行が一致すれば
$B=A^{\mathsf t}P$ または $B=PA^{\mathsf t}$ である。
それ以外では両方で二列だけ、または両方で二行だけが一致する。
前者で $B=AP,T'=PT$、しかも二つの下の角の両方で二列が一致する。
従って $T_{00}=T_{11},T_{01}=T_{10}$、$T$ と $A$ は $P$ と可換となり、
$C,D$ はともに $J$ の倍数で矛盾する。後者は転置で同じである。

残った二つの特異角表示を経路長4で排除する。
$B=A^{\mathsf t}P$ なら $T'=T^{\mathsf t}P$ であり、
$N_0$ の二つの列対は各々同一、$N_1^{\mathsf t}$ の二つの行対も各々同一である。
従って始点3の厚い二色の関係の係数空間はいずれも $(1,1,0,0),(0,0,1,1)$ の張る平面となる。
薄い二色は定数行である。始点2の長さ4で両端の二次関係を先に割ると、各総文字での次元は
$26,26,24,24$。中央の関係の挿入の関数体上の階数は $13,13,12,12$ で、
残る全次元は50である。具体的には端の各色の商を、厚いブロック $M$ の上の角を $G$、
左下を $H$ として $(-HG^{-1}\ I_2)$ で与え、薄い色では $(I_3\ {-\boldsymbol1})$ を使う。
左右のこの二射影に中央の関係の各項を通した行列を有理係数多項式環上で掃き出す。
\path{computations/verify_exchange_mixed_resolution.py} が独立な五変数で検算する。
階数6を保つ全特殊化でも次元は50以上であり、標準値35に反する。
$B=PA^{\mathsf t}$ は二位置移動で $A,B$ を転置し $\lambda$ を逆数にすれば同じである。

$C$ が可逆なら $N_0$ の下の角は $JC^{-1}J=tJ$。
$T_{00}=T_{11}$ から $A$ の非対角が等しく、$T'_{01}=T'_{10}$ から $B$ の対角が等しい。
これは $\lambda-1/\lambda\ne0$ による。もし $D$ も可逆なら下の角も定数で、
$T,T'$ は $P$ と可換、しかも $T'=PT$。従って $A,B$ も可換で $B=AP$、
$C$ が特異となり矛盾する。$D$ は対角が等しく非対角も等しい特異行列だから、
先の行または列の一致より $D=bJ$ である。$C$ の行列式は $x$ であり、$T_{00}=T'_{01}$ が
記載の $t$ の等式となる。$t(ad-fg)=a+d-f-g$ はその二式の和から従うので $tJ=JC^{-1}J$。
他方の下の角を $Z=\left(\begin{smallmatrix}T'_{00}&T_{01}\\T_{10}&T'_{11}\end{smallmatrix}\right)$ とすると、
残る条件は $b\det Z=Z_{00}+Z_{11}-Z_{01}-Z_{10}$ である。実際に両辺の差は
$T_{00}-T'_{01}$ の倍数で、商の分母は $\lambda\det A\det B$ を割る。
$Z$ も可逆である。さもなくば高位と低位の交換で上の角が二つとも特異となる。
この交換は $A,B$ を $T,T'$ にし $\lambda$ を逆数にするので、先の特異角分岐から
$T'=T^{\mathsf t}P$ または $PT^{\mathsf t}$、従って $B=A^{\mathsf t}P$ または $PA^{\mathsf t}$
となって $C$ が特異になる。ここで$Z$ の可逆性の議論に階数2の下界が必要な場合、
階数1なら $b\ne0,Z=J/b$ であり、二つの下の角が定数なので先と同じ矛盾である。
最後に $D$ が可逆な側は、二位置移動、奇文字への1の加法、全位置の $2,3$ の交換で
$(A,B,\lambda)$ を $(B,A,1/\lambda)$ として得る。これは始点の変更を明示した被覆である。
\end{proof}

\begin{proof}[特殊化と正則点]
$(tf-1)g=f+tb^2-(\lambda+1/\lambda)b$ と書ける。$tf-1$ で割った係数環は
$t$ を逆にした整域であり、その上で右辺は恒等的に零でない。
従って $g$ の一次方程式は原始的かつ既約で、$tf=1$ 上に別成分はない。
共通倍率を除けば四自由度である。

一般点を滑らかにする一変数を $\epsilon=z-1$ とし、節\ref{sec:klein-elliptic-completion}の
$A$ の左下を $c=b+\epsilon k$ とおく。$x=ad-fg$ を固定し、同節の
$f=f_0+f_1x,g=g_0+g_1x$ による $fg+x-ad=0$ を $k$ の式と見る。
全係数は $\epsilon=0$ で正則であり、$f_0=g_0=(a+d)/2$、
$f_1+g_1=-t$、$\partial f_1/\partial k=(1/\lambda-\lambda)/(2\det A)$ がその剰余である。
従って与えた $f,g$ を実現する剰余 $k$ があり、$f\ne g$ の一般点で方程式の
$k$ 微分は $(g-f)x(1/\lambda-\lambda)/(2\det A)\ne0$。形式的陰関数定理が持ち上げを与える。
同節の二次式を $p x^2+m x+n$, $q y^2+j y+o$ と書けば、正しい枝は
$y=-(2px+m+j)/(2q)$。$p,m,j$ は正則で、$e_1,h_1$ の単純極の主要部は逆、
$q=e_1h_1$ の極の位数は2である。従って $y=O(\epsilon^2)$ で $e,h$ は $b,b$ に戻る。
一般の剰余で $o=(b-c)(z^2b-c)/(z+1)^2$ の主要係数は非零なので一般繊維では $y\ne0$。
以上は五次元族からの特殊化であり、既約性により表示全体が閉包に入る。

二つの有理点 $(a,b,d,f,\lambda,g)=(2,3,7,5,2,2/5),(3,5,7,4,3,87/17)$ では
反対切断が可逆で、全始点の次数5までの次元は標準値である。文字別の四次中心次元は
$(3,2,0,0)$。前掲の検算が有理数上で四式を選び、その商の素数1009での次元を全始点で
次数13まで調べると $(1+t+t^2+t^3)^4$ の係数である。
五次元族の一般点から全次数の標準Hilbert関数が下界であるため、
命題\ref{prop:klein-rational-full-opposite}の自由加群・Koszul双対の証明が適用される。
従って両点で全次数の正則性が従い、有限Frobenius条件と各内部次数のbar複体から
本表示の幾何学的一般点も正則である。

第一の有理点を素数1753で還元すると、各位置で被約二十点、三点直線なし、四点平面52枚で
五点平面はない。各点を通る四点平面の数を保存するように枠の像を全列挙した
射影写像数の行は $(4,0,0,4),(0,4,4,0),(0,4,4,0),(4,0,0,4)$ である。
\path{computations/verify_exchange_mixed_configurations.py} は各配置を独立にも消去して照合する。
有限エタールな配置の固有な枠比較よりこの計算は上界を与え、位置0と2は標数零および
一般点で同値にならない。よって最小周期は4である。各配置の射影群は指定四元群を含むため
位数4である。窓を保つ四位置の写像は位置0の写像で決まる。実際二十点間の一意な点対応が
残りの配置の各点の像を決め、各配置は射影枠を含む。従って安定群は四元群であり、印を忘れる
ファイバーも有限である。厚い角の共通倍率以外の連続な対角残余がないことは、薄い
48語の文字指数行列の階数12と、上の角の文字指数差がその整数行格子に入ることから従う。
\end{proof}

% !TeX root = ../papers/nc-p3.tex
\section{二切断の交差階数が同時に一になる繊維の分類}
\label{sec:klein-double-exchange-fourier}
\begin{theorem}\label{thm:full-thin-double-exchange}
節\ref{sec:klein-elliptic-completion}の全薄係数非零の表示で $z=1$, $\lambda=1$ とする。
全ての標準Koszul左右AS正則窓は、位置ごとの文字置換により次のいずれかである。
\begin{enumerate}
\item 定理\ref{thm:klein-singular-exchange-corners}の正表示（最終文字を1だけ移して総文字を零にしたもの）。
\item $g_0,g_1,r_0,r_1\ne0$, $U=r_0+r_1,V=g_0+g_1$, $k=1-UV\ne0$ を取る。
二文字ずつ両半分に属する48語の係数は1とし、全て低位の八語の係数は $U$、
全て高位の八語の係数は $V$ とする。ただし $0000,1111$ の係数にはそれぞれ
$k/g_0,k/g_1$ を、$2323,3232$ にはそれぞれ $k/r_0,k/r_1$ を加える。
\end{enumerate}
記載の全点が正則である。$\lambda=-1$ の全解はこれらの非自明コサイクル捻りで得られる。
四つの族は全点で最小周期1である。
\end{theorem}
\begin{proof}[二位置ずつのFourier変換による被覆]
総文字零を課し、同じ半分の四文字の下位ビットの符号を $e_i\in\{1,-1\}$ とする。
$\prod e_i=1$ なので、低位八係数は一意に
\[
 m+\sum l_i e_i+A e_0e_1+B e_1e_2+C e_0e_2
\]
と書ける。高位は $n,h_i,A',B',C'$ と書く。切断 $01|23$ で $e_0e_1=e_2e_3=s$ とする。
低位角の二行・二列を定数と符号の基底で読んだ係数行列は
\[
 \begin{pmatrix}m+sA&l_2+s l_3\\l_0+s l_1&C+sB\end{pmatrix}.
\]
高位角も同じで、交差ブロックは左上の成分1のみである。切断 $12|30$ では
位置を巡回して $A,B$ を交換する。いずれかの角の右下成分が非零なら、
その成分をピボットに取った残り三行の階数は1以下である。
二つの交差成分が1なので、相手の角は左上の一成分だけとなり、それは非零である。
特に $(C+sB)(C'+sB')=(C+sA)(C'+sA')=0$ が両符号で成り立つ。

一方の半分の全八係数が定数なら正則性を失う。その半分の二文字の和で左右いずれから
評価しても、他方の変数の各半分で係数が一定になる。各奇位置でこの和の一次元を、
各偶位置で二つの半分の差の張る二次元を双対矢空間に取れば、全ての制限関係が零である。
交互にこれらの空間を双対とする経路代数への商で指数増大する。

$C,C'$ がともに非零なら $A/C,B/C=\pm1$、相手の比は逆の符号である。
全位置の下位ビットの平行移動（四移動量の和は零）で $A=B=C$, $A'=B'=-C'$ とする。
低位にピボットを持つ二切断の条件から $(h_0,h_1,h_2,h_3)=(y,-y,y,-y)$、
高位のピボットから $l_i=x$ を得る。もし $C\ne0,C'=0$ なら $A'=B'=0$。
どちらかの切断で低位の両符号にピボットがあれば $h_i=0$ が全てで成立し、高位が定数となる。
従って再び $A=B=C$ にでき、$h=(y,-y,y,-y)$ で $y\ne0$ である。
低位のピボットのない符号では高位角の二非対角成分がともに非零で右下が零だから可逆である。
Schur補完により低位の角は定数であり、やはり $l_i=x$ である。半分を交換すれば逆の場合も従う。
得られた型では低位係数は二つの定数語を除き一定、高位係数は二つの交代語を除き一定となる。

$C=C'=0$ で $A,B,A',B'$ のうち非零があるとする。低位の対係数が一つでも非零なら、
対応する切断で両符号にピボットがあり、高位の全一次係数が零になる。高位の対係数も
非零でなければ定数である。非零なものは $AA'=BB'=0$ により切断の異なる組で、
この場合低位の一次係数も全て零となる。例えば $A,B'\ne0$ なら $B=A'=0$、
二切断条件から $m(n\pm B')=1$ が両方に必要で矛盾する。

残る全対係数零の場合、切断 $ij|kl$ の符号 $s$ に対して
$u=l_i+s l_j,v=l_k+s l_l$, $p=h_i+s h_j,q=h_k+s h_l$ と置く。
三次小行列の条件は $uq=pv=nuv=mpq=0$ で尽きる。
どちらの一次係数も全ては零でない。隣接する $l_i^2,l_j^2$ が異なれば対向する組の
$h_k,h_l$ はともに零である。$h$ の支持が隣接二位置ちょうどなら、その両端と外側との
二つの隣接対で平方が異なるので逆に $l$ 全体が零になる。
一位置だけなら $l$ の支持もその位置だけである。この場合、その位置を含まない連続二辺では
各厚色の係数空間が、各半分で一定なベクトルの二平面になる。
両方の関係空間に $(e_0+e_1)\otimes(e_0+e_1)$ が入り、
定理\ref{thm:no-consecutive-pure-relations}で排除される（ここで $e_i$ は文字基底）。
従って全 $l_i^2$ は同一非零であり、$h_i^2$ も同様である。
平行移動後 $l=(x,x,x,\delta x)$, $\delta=\pm1$ とすると
$h=(y,-\delta y,y,-y)$。$\delta=-1$ は
$(X,Y,W,V)=(m+2x,m-2x,n+2y,n-2y)$ の第一型である。
$\delta=1$ は先の定数語・交代語だけが異なる型に含まれる（$nuv=mpq=0$ から $m=n=0$）。

第二型の四例外語への加算量を順に $p,q,j,t$ と置くと厚い階数の上界は
$Vpq=(1-UV)(p+q),\ Ujt=(1-UV)(j+t)$ である。
これは定数値 $U,V$ を持つ二行対に対する三次小行列式そのものである。
$k\ne0$ なら片方の加算量だけ零にはならず、二つとも零なら定数半分として除かれる。
従って $g_0=k/p,g_1=k/q,r_0=k/j,r_1=k/t$ が記載の表示を与える。
$k=0$ なら $pq=jt=0$ で、両方の低位加算量が零なら第一厚色の階数が1に落ちる。
一方だけ非零ならその二列の差が文字の平方の純関係となる。
テンソルは偶位置の交換にも奇位置の交換にも対称であるため、
補題\ref{lem:bisymmetric-reducible-quick-test}の評価判定の必要条件、または連続純関係の排除で落ちる。
\end{proof}

\begin{proof}[十分性と周期]
第二型の色0関係の係数平面は $(a,b,g_0a+g_1b,g_0a+g_1b)$、
色1は $(r_0c+r_1d,r_0c+r_1d,c,d)$ である。薄い二色は全成分1。
各々を $ij$ の係数行列にまとめると交差ブロックはともに
$\left(\begin{smallmatrix}u&v\\v&u\end{smallmatrix}\right)$。
非零の交差を持つ純関係なら $v=\pm u$ であり、同じ符号の二文字ベクトルが左右各半分に現れる。
従って $a=b$, $c=d=\pm Va$, $a=UVa$ となり矛盾する。交差が零なら片方の角が零で、
四係数 $g_0,g_1,r_0,r_1$ の非消滅により残る角も階数1にはならない。

評価点を $(x_0,x_1,x_2,x_3)$ とし $S=x_2y_2+x_3y_3,T=x_0y_1+x_1y_0$ と置く。
厚い四式は $x_0y_0=-g_0S,x_1y_1=-g_1S,x_2y_3=-r_0T,x_3y_2=-r_1T$。
全座標が非零なら $S$ が $y$ を決める。低位の一方だけが零なら $S=0$ となり、
高位二座標が非零のとき残りの自由度は高々一つである。高位にも一つ零があれば $T=0$ で、
四式と $S=T=0$ が核を零にする。高位だけ一座標が零の場合も $T$ から先に同じ議論をする。
片半分だけが非零なら厚い式が同じ半分の二座標を消し、薄い式が少なくとも一条件を加える。
右評価は $r_0,r_1$ の交換で同じである。反対切断の薄い二色の各核は
二つの交差積の差を含み、その和が二つの一次式の積になる。よって
補題\ref{lem:bisymmetric-reducible-quick-test}により全次数で正則である。

$P$ を高位二文字の交換とすると関係六平面は $R^{\mathsf t}=(P\otimes P)R$ を満たす。
全辺が $R,R^{\mathsf t}$ を反復するので最小周期は1である。
非自明捻りで $P$ は次数の自己同型 $i\mapsto i+(i\mathbin{\mathrm{div}}2)$ を誘導し、
二コサイクルの比は $(-1)^{i_{\rm hi}j_{\rm hi}}$、すなわち
$\beta(i)=(\sqrt{-1})^{i_{\rm hi}}$ の余境界である。よって一文字倍率で補正した $P$ が
捻り後も一位置移動を実現する。さらに全薄い48語上で
$\prod_{i<j}\mu(q_i,q_j)$ に対角文字倍率を掛けると、ちょうど $\lambda$ を $-\lambda$ にする。
低位の厚角 $A,B$ への作用は $DAD,-DBD$ ($D=\operatorname{diag}(1,-1)$)。
\path{computations/verify_double_exchange_fourier.py} が全語でこの符号と切断の式を確かめる。
全次数の正則性は捻りで保存される。
\end{proof}

\begin{theorem}\label{thm:double-exchange-frames}
第二型は正表示で反対階数8、捻りで
\[
 16-[U=0]-[V=0]-[k^2+(4Ug_0g_1+kV)(4Vr_0r_1+kU)=0]
\]
を持つ。両表示とも印を忘れた添字保持同型の必要十分条件を、全パラメータで次の有限な枠比較により得る。
正表示では各反対核にある四つの純テンソルの因子を枠とする。捻りでは各位置の被約二十点から
枠を取り、座標平面を含む全ての候補を調べればよい。四位置の枠を順序付け、テンソルの
非零台とその整数文字関係の単項式比を記録する。正規形としてこの有限集合を取れば、
二集合が交わることと同型であることは同値である。第一型と第二型および各々の正負の間に同型はない。
第二型の正負の全点が前節の混合補完族の特殊化となる。
\end{theorem}
\begin{proof}[反対核と捻りの二十点]
正表示の反対切断は色順に階数 $(4,2,1,1)$ である。色0の行列式は
$-k^3/(g_0g_1r_0r_1)$、色1は二組で行・列が等しく縮約行列の行列式が $UV-1$ となる。
対称二次式の核は $e_0 e_2-e_1 e_3,e_0e_3-e_1e_2$ の鉛筆で、
積となる二元は $(e_0\pm e_1)(e_2\mp e_3)$ である（符号を入れ替えても同じ二組）。
二つの反対位置を別空間としても、純テンソルの対称化はその二因子の積であり、
従って全反対核の純テンソルはこれらの各順序の四点で尽きる。
左右の四因子は独立で、相手側の反対切断にも同じ議論を適用できる。
捻りの薄い反対二色は行列式 $-16$、色1は二組の和と差に分けると階数
$2+[U\ne0]+[V\ne0]$ である。色0の二つの角は可逆、交差は階数1であり、
Schur補完の階数損失は高々1である。行列式を計算すれば上記の式を得る。

捻りの左点を斉次座標 $(a,b,c,d)$ で書く。二次関係から
$S=c y_2-d y_3$, $T=a y_1+b y_0$ とおいて
\[
 ay_0=-g_0 S,\quad by_1=-g_1 S,\quad cy_3=r_0T,\quad dy_2=-r_1T .
\]
薄い二式は $ay_2+by_3+cy_0-dy_1=ay_3+by_2-cy_1+dy_0=0$。
全座標が非零なら $z=b/a,w=d/c$, $G=g_1/z+g_0z$, $H=r_1/w+r_0w$ に対し
$T=-GS,S=-HT$ で $GH=1$。薄い二式は $G,c^2/a^2$ の同次二式となるが係数行列式は
\[
 (r_1/w-r_0z)(g_1/z-g_0w)-(-r_0+r_1z/w)(-g_0+g_1w/z)
 =GH-UV=k\ne0
\]
となり矛盾する。一座標が零で他の三つが非零の場合は次の四組を得る。
\[
\begin{array}{c|cc}
d=0 & a^2=g_0c^2& b^2=-g_1c^2\\
c=0 & a^2=-g_0d^2& b^2=g_1d^2\\
b=0 & c^2=r_0a^2& d^2=-r_1a^2\\
a=0 & c^2=-r_0b^2& d^2=r_1b^2 .
\end{array}
\]
二座標以上が消える場合は $(1,\pm1,0,0),(0,0,1,\pm1)$ のみである。
実際片半分が零なら厚い式が同じ半分の $y$ を零とし、残りの二式の行列式は
二つの平方の差となる。両半分に一零ずつあれば $S=T=0$ から四座標とも消える。
これで相異なる二十点があり、有限な期待余次元の最大小行列式スキームの長さ20から被約である。
他の位置も既に構成した一位置同型で同じとなる。

各座標平面はその上の四つの片零点で張られ、そこに純半分の二点もある。
各位置で配置の独立三点が張り六点以上を含む平面を全て取り、その中の法線が
独立な四平面を選べば、座標四面体を含む固有な有限集合が得られる。
追加の六点平面を持つ特殊値もこの規則に含む。どちらの符号でも残る対角倍率は
非零語の文字指数行列の飽和整数核で正確に比較できる。
正表示二型は反対階数6と8で異なり、負表示はいずれもそれより大きい。
負表示二型は点配置の曲線部分の有無で異なる。
\end{proof}

\begin{proof}[混合補完からの特殊化]
第二型の低位角を $A=\left(\begin{smallmatrix}a&b\\b&d\end{smallmatrix}\right)$ として
$b=U,a=b+k/g_0,d=b+k/g_1$ と取る。$\det A=k/(g_0g_1)$ である。
$\lambda=1+\epsilon,j=r_0-r_1,f=b+\epsilon j$ とし、前節の式で $t,g$ を決める。
$tb-1$ の剰余は $-k$、$g$ の剰余は $b$ となり、$\det B$ は正確に位数2を持つ。
実際その主要係数は $(b-j)(b+j)\det A/((a-b)(b-d))$ である。
$T'$ の非対角の剰余は $V$、対角の剰余は順に $V+2k/(b+j),V+2k/(b-j)$ となる。
従って第二型の四例外語が全て希望の値を持つ。一般繊維では
$x=ad-fg$ の剰余も $\det A\ne0$ であり、前節の混合補完条件に属する。
コサイクルの符号補正はこの補完で $b$ と $\lambda$ の符号を変え、他の $a,d,f,g$ を保つ。
よって負表示も $\lambda=-1$ での同じ族の特殊化である。
\end{proof}

% !TeX root = ../papers/nc-p3.tex
\section{平行な二重色と四つの非欠損長方形の全被覆}
\label{sec:full-thin-parallel-cover}
\begin{theorem}\label{thm:full-thin-parallel-cover}
四元群が四位置で正則表現型に作用し、二つの隣接切断の色階数がともに $(2,2,1,1)$ の
並べ替えで、二つの厚い色の差が一致するとする。四つの薄い階数1長方形の全ての行・列が
非零なら、全ての標準Koszul左右AS窓は、文字の付け替えと始点の移動を含めて、
次の三族の閉包で覆われる。
\begin{enumerate}
\item 定理\ref{thm:klein-elliptic-completion}の五次元族。
\item $\lambda=1,B=A$ の両対称な四次元交換表示。
\item その非自明コサイクル捻り。
\end{enumerate}
これは四つの薄い長方形の係数だけに非消滅を課した被覆であり、
十六個の厚い係数のうちいくつが消えてもよい。
異なる厚色方向では定理\ref{thm:klein-transverse-dense}の二変数表示が適用されるため、
二重色の両方向について、薄い欠損のない場合に正則候補を含む族が尽くされた。
\end{theorem}
\begin{proof}[薄い係数の商と切断の入替え]
節\ref{sec:klein-elliptic-completion}の規格化を使う。薄い48語は二文字ずつが両半分にある語で、
係数は $z,\lambda$ のLaurent単項式である。各語の四位置での文字指数行列を $E$、
その転置の整数核の基底を行とする行列を $L$、薄い係数の二指数の行列を $F$ とすれば、
$LF$ の整数行格子は $\mathbb Z^2$ そのものである。従って対角文字倍率を除いた薄い情報は
この二単元で正確に表される。共通の線型文字置換と位置別平行移動で両厚色集合 $\{0,1\}$ を
保存する32通りを考える。線型部は恒等または2と3の交換、四移動量は総和零で全上位ビットが等しい。
置換・始点移動後の指数行列を $F'$ として $L F'=LF M$ を解くと、
$(t,v)=(\lambda,z\lambda)$ に対する作用は次の通りである。
始点固定では $(t,v)$ の交換と同時逆数が生成する四元群であり、
一位置移動を許せば独立な逆数と交換が生成する位数8の全群となる。
\path{computations/verify_full_thin_exchange_cover.py} は位置移動の偶奇別に全候補を検算する。
これらは実際の対角基底変更を伴う。代数閉体上でトーラス写像の像はその整数文字関係で決まるからである。

$x=t+t^{-1},y=v+v^{-1}$ と書くと、始点を指定した薄い表示の商は
\[
 S=x+y,\quad P=xy,\quad D=(t-t^{-1})(v-v^{-1}),\qquad
 D^2=(P+4)^2-4S^2 .
\]
実際に同時逆数の不変環は $x,y,D$ で生成され、交換は $x,y$ を交換して $D$ を固定する。
一位置移動は $D$ の符号を変える。この商は指定した文字作用についてのものである。
同じ二点が $\{z,z\lambda^2\}$ について順序交換と独立な逆数を受ける。
数 $z$ と $z\lambda^2$ の一方が1なら、切断の向きを選んで $z=1$ にできる。
他方も1なら定理\ref{thm:full-thin-double-exchange}による。第二型は五次元族へ、
第一型の正負は定理\ref{thm:klein-bisymmetric-irreducible}により対応する交換表示の閉包に入る。
他方が1でなければ定理\ref{thm:klein-mixed-exchange-resolution}による。

それ以外は両切断の交差ブロックが全て可逆であり、各厚い角も階数2条件から可逆である。
$z,z\lambda^2$ の一方が $-1$ とも異なれば、そちらを $z$ にして
定理\ref{thm:klein-all-pairings}を適用できる。四つの例外表示の二つの正表示は
$(t,v)$ の交換で相互に移る。具体的に $v=1$ の一般点は $t'=1$ へ移り、
正交換の一般点は両対称なので負交換の一般点とは反対階数により異なる。
五次元主枝は $A,z,\lambda$ の一般値の上の既約な解であり、
いずれの付け替えもこの次元の唯一の枝を保つ。
従って付け替え後にも両対称例外枝の稠密開部は主枝の外にある。
同じことを捻りの後にも使える（薄い符号は $\lambda\mapsto-\lambda$）。
四例外表示は各々既約なので、一般点での帰着はその閉包全体の帰着を与える。
\end{proof}

\begin{proof}[最後の繊維 $z=z\lambda^2=-1$]
コサイクル捻りで $\lambda=1$ とする。
$A=\left(\begin{smallmatrix}a&b\\c&d\end{smallmatrix}\right)$ に対し
$(p,q,r,w)=(a+d,a-d,b+c,b-c)$、$B$ の対応する四数を $(P,Q,R,W)$ と書く。
節\ref{sec:klein-elliptic-completion}の角 $C,D$ を使い、
$u=\det C/\det A$, $v=\det D/\det A$, $t=\det B/\det A=u+v-1$ とする。
いずれも非零である。$X=\left(\begin{smallmatrix}1&1\\1&-1\end{smallmatrix}\right)$,
$Y=\left(\begin{smallmatrix}1&1\\-1&1\end{smallmatrix}\right)$ とすると
\[
 Y\operatorname{adj}(A) X=
 \begin{pmatrix}p-r&-q+w\\q+w&-p-r\end{pmatrix}.
\]
$C$ は $(p,q,R,W)$、$D$ は $(P,Q,r,w)$ で与えられる。
従って二つの切断の条件はちょうど
\[
 R=ur,\quad Q=vq,\quad (u-1)(p,P)=0,\quad (v-1)(w,W)=0,
 \quad (u-1)(v-1)(q,r)=0
\]
である。最後の式は $uv-t=(u-1)(v-1)$ から来る。
$u,v$ がともに1と異なる場合は $A=0$ となるので解がない。

$u=v=1$ なら $R=r,Q=q$, $P=\pm p,W=\pm w$ で尽きる。
この四表示は $B=A,A^{\mathsf t},
\left(\begin{smallmatrix}-d&b\\c&-a\end{smallmatrix}\right),
\left(\begin{smallmatrix}-d&c\\b&-a\end{smallmatrix}\right)$ の順である。
第一は正交換、第三は $b\ne c$ なら節\ref{sec:klein-doublet-boundary}の
第二補完を $z\to-1$ とした極限そのものであり、$b=c$ にも特殊化する。
第二は一位置移動を含む下記の第一変換で第三へ、第四は捻りを含む第二変換で第一へ移る。
前者の第二補完枝は $\lambda=1$ で五次元族に含まれた有理切断である。
実際八角を $(a,b,c,d,e,f,g,h)$ と並べると二変換は
\[
 (g,a,-d,-f,c,e,-h,-b),\qquad (b,a,-d,-c,g,-h,e,-f)
\]
である。いずれも後二文字に1を足す。第一は一位置の巡回を、第二は非自明捻りを伴い、
薄い符号を合わせる文字倍率だけで上式を得る。薄い語上の整数Smith方程式と八角で検算でき、
高位角はSchur補完で一意である。

$u=1,v\ne1$ なら $w=W=0,R=r$ なので二つの非対角はすべて $b$ で一致する。
$A,B$ は対称であり、残りは
\[
 P^2=v p^2+(1-v)r^2+v(v-1)q^2,\qquad Q=vq .
\]
これがこの側の全解である。同じ非対角を持つため偶交換そのものが成立する
（両切断で $C=A,D=B$ が成立する）。相異なる補完は全て既出の第二補完の閉包に入る。
直接には $e=(kd-b^2)/(d-k)$, $h=(ka-b^2)/(a-k)$ と取って $k$ を動かせばよい。
この式は $(e-h)(ad-b^2)=(a-d)(eh-b^2)$ と同値な一般表示であり、
上記の方程式は $v$ に関する二次式の平方根のモデルで既約である。
実際右辺は $p,q,r$ の階数3の二次形式で、$v(v-1)\ne0$ で平方にならない。
よって特殊値も閉包に入る。$v=1,u\ne1$ では $p=P=0,Q=q$ であり、
上の第一変換で対称な二角・共通非対角の側に戻る。
これが最後の繊維の全被覆を与える。
\end{proof}

\begin{corollary}\label{cor:minus-one-second-completion-periods}
第三表示 $B=\left(\begin{smallmatrix}-d&b\\c&-a\end{smallmatrix}\right)$ と
$u=1,v\ne1$ の表示はともに反対階数10となる稠密開集合を持つ。
両者の各々の稠密開集合の全閉点は標準Koszul左右AS正則である。
第三表示の幾何学的一般点の最小周期は4、射影安定群の位数は8である。
$u=1,v\ne1$ 表示の反対階数10の正則点は全て最小周期2である。
\end{corollary}
\begin{proof}
前者で $(a,b,c,d)=(2,3,5,7)$、後者で $(a,b,d,e,h)=(2,3,7,5,1)$ とする。
両点で反対階数10、全始点で五次までの経路次元は標準値、四次中心の色別次元は
$(3,2,0,0)$。有理中心四式で割った商の素数1009での次元は全始点で次数13まで
$(1+t+t^2+t^3)^4$ の係数である。定理\ref{thm:klein-generic-central-parameters}の
第二補完枝の一般中心次元に等しいので、この二点の周囲で中心基底が延び、
有限個の商次元上界と全次数のHilbert下界を比較する自由加群の証明が適用できる。

前者の標本の素数109での四つの配置は各々相異なる二十点で、
射影写像数の四行は $(8,0,8,0),(0,8,0,0),(8,0,8,0),(0,0,0,8)$ である。
三点直線を許して、各点を通る三点直線上の三つ組と大きさ別の多点平面を数え、残る枠の像を全列挙した。
被約二十点の固有な枠比較でこの上界を一般点に使える。位置1と3の非同値性で周期2が除かれる。
位置0の射影安定群が窓の安定群の上界になることは前節の有限グラフの場合と同じである。
四元群に加えて位置別移動 $(1,0,1,0)$ が文字倍率補正により窓を保存する。
実際その八角の変換は共通倍率を除き $(h,-f,-g,e,d,-b,-c,a)$ であり、
第三表示では元の八角の $-1$ 倍になる。よって位数8が達成される。
後者の表示では二位置移動が文字倍率を除き $(A,B)\mapsto(A^{\mathsf t},B^{\mathsf t})$ を与え、
両角が対称なので周期2を持つ。周期1なら偶交換を一位置運んで奇交換も得る。
反対対は互いに素であるため両同定を同時に使えて両対称となるが、
定理\ref{thm:bisymmetric-six-or-eight}に反する。
\end{proof}

\begin{corollary}\label{cor:mixed-ordinary-open-center}
定理\ref{thm:klein-mixed-exchange-resolution}の四次元表示でも、同節の二つの有理点を
含む一つの稠密開集合の全閉点が四中心四次式上自由であり、標準Koszul左右AS正則である。
\end{corollary}
\begin{proof}
五次元楕円族は定理\ref{thm:full-fifteen-component}の一意な十五次元成分に含まれ、
従って今回の混合表示もそこに含まれる。この成分の一般の四次中心は五次元である。
二つの有理点で全始点の長さ4,5の次元は標準値に等しく、経路束はその近くで局所自由である。
中心方程式の核次元はそこで5以上で、二点ではちょうど5である。よって中心はその近くで部分束となり、
認証に使った四元を局所切断として延ばせる。中心商の長さ13までの次元上界は開条件なので、
二点それぞれの近傍の和で自由加群と全次数正則性の証明が通用する。
\end{proof}

% !TeX root = ../papers/nc-p3.tex
\section{等行列式の第二補完における八つの評価障害と無印商}
\label{sec:equal-determinant-second-sheet}
前節の $z=-1,\lambda=1$ の第三表示 $B=\left(\begin{smallmatrix}-d&b\\c&-a\end{smallmatrix}\right)$ を取る。
$D=ad-bc$, $p=a+d,q=a-d,r=b+c,w=b-c$ とし $D\ne0$ を保つ。
$p=0$ では $B=A$、$w=0$ では前節の第四表示に一致する。
従ってこの二平面では既に正交換と負交換の端点で全次数正則性が決定されている。
\begin{theorem}\label{thm:equal-determinant-point-configurations}
$pw\ne0$ とする。全四辺の左右の非零評価が階数3以上であるための必要十分条件は
\[
    (a^2-b^2)(d^2-c^2)(a^2-c^2)(d^2-b^2)\ne0 .
\]
条件が破れれば非標準ASである。条件が成立すれば反対階数は10、四辺はいずれも
長さ20の有限スキームの同型のグラフである。位置1の座標の平方を順に並べると、支持は
$(1,-1,0,0),(0,0,1,1)$ と次の四点の、座標ごとの平方写像による全逆像である。
\[
\begin{array}{c|cccc}
 &Y_0&Y_1&Y_2&Y_3\\ \hline
 &4b&0&(a+b)^2&(a-b)^2\\
 &0&4c&(d+c)^2&(d-c)^2\\
 &(d-c)^2&-(a-b)^2&D(d-c-a+b)&0\\
 &(d+c)^2&-(a+b)^2&0&D(a+b-c-d)
\end{array}
\]
位置3では $b,c$ を交換する（文字倍率の符号を除く）。
各位置の相異なる点の数は $20-4[b=0]-4[c=0]-4[q=w]-4[q=-w]$。
合流した各支持の局所環は $k[t]/(t^3)$、他の支持は被約である。
\end{theorem}
\begin{proof}
奇位置の座標を $y=(y_0,y_1,y_2,y_3)$、位置0の座標を $(u,v,s,t)$ とする。
厚い二色の四式を
\[
 D\binom{uy_0}{vy_1}=\begin{pmatrix}c-d&c+d\\-a+b&-a-b\end{pmatrix}\binom{sy_2}{ty_3},\quad
 D\binom{uy_1}{vy_0}=\begin{pmatrix}a+b&a-b\\c+d&c-d\end{pmatrix}\binom{sy_3}{ty_2}
\]
と取れ、薄い二式は $uy_2+vy_3+sy_0-ty_1=uy_3+vy_2+sy_1+ty_0=0$。
奇位置では全座標の独立な符号変更が点配置を保存する。実際四元群の高低二ビットの文字に加え、
前節の移動 $(1,0,1,0)$ は偶位置で両組内の交換、奇位置1で第1番目だけ、
位置3で第0番目だけの符号変更を与える（番号は0から）。
最大小行列式の各文字成分もスキーム上零となる。$y_0y_1y_2y_3$ の成分の係数に
$\pm2(b-c)D\ne0$ があるため、必ず一座標が零である。

$y_3=0,y_2=1$ なら $u=-sy_0+ty_1,v=-sy_1-ty_0$。
$L=(d-c)/D,M=(a-b)/D$ と置く。残る式は二行二列の回転
$(s,t)\mapsto(-t,s)$ の二固有線ごとに解ける。$t=\alpha s$, $\alpha^2=-1$ の線では
\[
 y_0(y_0-\alpha y_1)=L,\qquad y_1(y_0-\alpha y_1)=-\alpha M .
\]
これは四つの厚い式を $(s,t)$ と $(-t,s)$ の係数で読むことでも従う。
$L,M$ はともには零でない。$L=M$ なら解はなく、異なれば第三行の平方値が得られる。
$L,M$ の一方だけが零なら同じ $y$ で両固有線が核に入り、評価階数が2以下になる。
$y_2=0,y_3=1$ でも $L,M$ を $-(d+c)/D,-(a+b)/D$ とすれば同じ消去で第四行を得る。
この二平面の検査で $a=\pm b$ と $d=\pm c$ を全て排除した。

$y_1=0,y_0=1$, $y_2y_3\ne0$ で、上で除いた平面を避けるとする。
二番目と三番目の厚い式から $t=\alpha s$, $\alpha^2=1$、
$y_2=-\alpha(a+b)y_3/(a-b)$。$s=t=0$ は他の二式により不可である。
残る二つの厚い式より $u/s=2\alpha y_3/(a-b), v=\alpha u$。
薄い式から $4b y_3^2=(a-b)^2$ となる。これで第一行の四点を得る。
$b=0$ ならこのチャートに点はない。$y_0=0$ 側は同じ代入で第二行となる。
片半分だけが非零なら、厚い式により対応する半分の相手座標が零となり、
薄い二式から低位の平方比 $-1$ または高位の平方比 $1$ の二点となる。
いずれも核は直線である。各三座標の点でも記載の比と非零の分母から核は直線である。
二位置移動は二つの角を転置するので、もう一方の奇位置で同じ判定をすると
$a=\pm c,d=\pm b$ が加わる。片側反転の評価下界定理\ref{thm:single-opposite-evaluation-ranks}
より、八平面のいずれも非標準ASである。

残った全点で奇位置の射影は有限で核が直線だから、辺自体も有限である。
従って偶位置の核も高々直線である。点対は六つの双一次式の期待余次元の完全交叉で長さ20であり、
両射影は最大小行列式のスキームへの同型である。
反対切断では色0の行列式が $8p^2/D$、薄い二色は各階数2。
色1の二組の行は組ごとに等しく、残る二行は $(b,c,1,1)$,
$(1,-1,(-q+w)/D,(q+w)/D)$ である。比例すれば $q=r=0$ かつ $a^2=b^2$ で除かれている。
従って階数の和は10である。

点対の族は上の評価開集合で有限平坦である。実際、滑らかな基底上の相対完全交叉で
全ファイバーが同じ零次元を持つ。一般配置からの特殊化で第一・第二行の各四点は
$b=0,c=0$ でそれぞれ純高位二点に二点ずつ合流する。同時の消滅は $w\ne0$ に反する。
第三・第四行も $q=w,q=-w$ で純低位二点に二点ずつ合流し、この二条件も同時には起きない。
他に重なりはない。局所長さ3のGorenstein環は接次元1であるから局所構造も従う。
元の六式、符号成分と後の置換作用は \path{computations/verify_equal_determinant_second_sheet.py} が検算する。
\end{proof}

\begin{theorem}\label{thm:equal-determinant-generic-quotient}
$x=q/p,y=r/p,s=w/p$ とする。射影安定群の位数8の部分では、指定始点の無印同型は
\[
    Q=x^2+y^2,\qquad L=xys,\qquad H=s^2
\]
の一致と同値である。この軌跡は稠密開集合を含み、その商は有理三次元である。
一位置だけの周期は反対階数10の正則点では起きず、二位置移動は $L$ の符号を変える。
従ってこの部分の正則点では $L=0$ で周期2、他で周期4である。
評価開集合の残りにも、支持点の有限枠と文字倍率による無印同型・周期の必要十分判定がある。
\end{theorem}
\begin{proof}
四元群の二指標を $H_0$（高位ビット）と $J$（低位ビット）、前記の追加対称を $S_0$ とする。
射影的には可換な三つの対合である。偶位置の線型持上げでは $J$ と $S_0$ だけが反可換となる。
従ってこの群のうち四位置の文字線を与え得る階数2の群は、交換子形式の根基
$\langle H_0\rangle$ を含まねばならない。$\langle H_0,S_0\rangle$ は奇位置1の
高位二文字が同じ文字を持ち、$\langle H_0,JS_0\rangle$ では低位二文字が同じになる。
よって四位置で正則表現を与える四元群は元のもの一つである。
系\ref{cor:minus-one-second-completion-periods}の二十点標本により一般点の安定群はちょうど位数8である。
有限エタールな点配置が枠を持つ開部分では自己同型候補の数は局所的に有限であり、
標本で失敗する候補を除けば共通の稠密開集合を取れる。

この固有な四元群を運ぶ写像は、各位置の文字基底の置換と対角倍率である。
薄い48語から線型部分は共通で、移動量の上位ビットも共通になる。
32候補から薄い符号を合わせられる指定始点の16個を取り、四角の共通倍率を除いて
$(x,y,s)$ を読むと、作用は $x,y$ の交換と
$(x,y,s)\mapsto(\epsilon x,\eta y,\epsilon\eta s)$ で尽きる。
実際符号補正のSmith方程式を解くと、$[a:b:c:d]$ に
$(a,-b,-c,d),(d,c,b,a),(p+r,q+w,q-w,p-r)$ の作用がある。
各射影変換を二置換ずつが与える。$H\ne0$ だから $x^2y^2=L^2/H$ であり記載の三数が軌道を分ける。
$D\ne0$ は $1-Q+H\ne0$。二位置移動は角の転置により $s$ の符号を変える。
周期1の排除は前節の二交換の議論による。

安定群の拡大点でも、両奇位置では表の支持が、両偶位置ではその評価核の支持が得られる。
各々が独立四点を含む。実際、二組の純半分の点の相手は反対の半分にある相異なる二点である。
全位置で独立四点の順序付き組を列挙し、その各枠で四窓の非零台と飽和整数文字関係の単項式比を取れば
共通の正規形を持つかが判定できる。局所環を点数だけで同一視する必要はなく、比較するのは元の四窓自体である。
各移動について同じ有限検査をすれば周期も決まる。これらの同型判定は正則点では添字保持同型を与える。
\end{proof}

% !TeX root = ../papers/nc-p3.tex
\section{行列式比の円錐曲線族の全評価と三次元無印商}
\label{sec:unequal-determinant-conic}
節\ref{sec:full-thin-parallel-cover}の $u=1,v\ne1,z=-1,\lambda=1$ の表示で、
$A=\left(\begin{smallmatrix}a&b\\b&d\end{smallmatrix}\right)$,
$B=\left(\begin{smallmatrix}e&b\\b&h\end{smallmatrix}\right)$,
$D=ad-b^2$, $\beta=(eh-b^2)/D$ とする。$D\beta(1-\beta)\ne0$ である。
$e-h=\beta(a-d)$ により一意な $k$ があり、
\[
 e=k+\beta(a-k),\qquad h=k+\beta(d-k),\qquad
 \beta(a-k)(d-k)=(k-b)(k+b).
\]
最後の式は行列式比から得られ、逆も成り立つ。以下
$I=(a-b)(e-b),\ J=(a+b)(e+b),\ K=(d-b)(h-b),\ N_0=(d+b)(h+b)$、
$T_\pm=e\pm b+\beta(d\pm b)$ とする。

\begin{theorem}\label{thm:unequal-conic-all-evaluations}
四辺の全ての非零左右評価が階数3以上となる必要十分条件は
$(1+\beta)I J K N_0\ne0$ であり、その外に標準AS窓はない。
この開集合の反対階数は10、四辺は各々長さ20の有限スキームの同型のグラフである。
偶位置0の平方座標の支持は次の六点の全逆像である。
\[
\begin{gathered}
 (1,1,0,0),\quad(0,0,1,-1),\\
 (0,-2(d+h),-N_0,K),\quad(-2(a+e),0,-J,I),\\
 (-N_0,-J,0,D T_+),\quad(K,I,D T_-,0).
\end{gathered}
\]
異なる支持点数は $20-4[d+h=0]-4[a+e=0]-4[T_+=0]-4[T_-=0]$。
合流点は長さ3の曲線状局所環を持つ。正則なら最小周期は2である。
\end{theorem}
\begin{proof}
偶点 $x=(u,v,s,t)$ で奇点に作用する六式の行列は
\[
 E(x)=\begin{pmatrix}
 u&0&ms&-nt\\0&v&fs&gt\\0&u&Nt&Ms\\v&0&-Gt&Fs\\
 s&-t&u&v\\t&s&v&u
 \end{pmatrix},
\]
ここで $(m,n,f,g)=(d-b,d+b,a-b,a+b)/D$,
$(M,N,F,G)=(h-b,h+b,e-b,e+b)/(\beta D)$。
厚い二角の列基底から得られる。反対の偶位置では最後の偶座標の符号を変えれば同じである。
実際二移動の対角倍率は偶で $(-1,-1,-1,1)$、奇で恒等となり、
偶交換から逆辺は転置である。

まず $I J K N_0\ne0$ とする。$k=\pm b$ なら $a$ または $d$ が $k$ に等しくなるので、
ここでは $(a-k)(d-k)(k^2-b^2)\ne0$。行番号を0から数え、四行の最大行列式を $[ijkl]$ とすると
\[
 (m-M)[0123]+Mf[0345]-Fm[1245]=cuvst,
\quad c=\frac{-2(a-b)(a+b)(b-d)(b+d)(a-k)(d-k)(1+\beta)}
 {(b-k)(b+k)^2D^3}.
\]
これは上の行列と円錐の式の展開である。$\beta\ne-1$ なら一座標が消える。
$u=0,vst\ne0$ では核の高位座標を $(nt,ms)$ に取れ、
第三式で $mM s^2+nN t^2=0$、第二・第四式で二低位座標が決まる。
薄い二式は $mM v^2=-2(d+h)t^2/(\beta D^2)$ に還元される。
$v=0,ust\ne0$ では同じく高位を $(gt,-fs)$ に取ると第四表示を得る。
$s=0,uvt\ne0$ では高位座標の比は $y_2/y_3=nv/(Gu)=gu/(Nv)$、
薄い式は $(G+n)v^2=-gG t^2$ となる。
$t=0,uvs\ne0$ では $y_3/y_2=mv/(Fu)=fu/(Mv)$ と $(M+f)u^2=mM s^2$ である。
ここでは $fn+gm=2/D$, $FN+GM=2/(\beta D)$,
$M+f=F+m$, $N+g=G+n$ を使った。各計算の核は一本である。
二座標が零なら、低位と高位に一つずつ非零のときは厚い式が核を零にする。
同じ組に二つ非零のときは表の最初の二表示となり、厚い式は他方の二座標を自由にする。
薄い式でその一本が残る。一座標のみの点では階数4である。
よって左の像は有限で全核は一本、右の点に直線の逆像があれば左の像の有限性に反する。
六つの双一次式はゼロ次元完全交叉で、その長さは $\binom63=20$ である。
反対の二移動と転置により全四辺に通用する。反対切断の四色の階数は $4,2,2,2$。
第一行列式の分子は定数倍を除いて
$((a-k)(d-k)(1-\beta))^2(a-k)(d-k)(1+\beta)$、分母はここで単元であり、
他の三行列の独立二行は直接取れる。周期は系\ref{cor:minus-one-second-completion-periods}による。

$a+e=d+h=0$ は $\beta=1$ を意味し、$T_+=T_-=0$ なら $b=0$ かつ
$e=-\beta d,h=-\beta a$ となり再び $\beta=1$ である。
従って高位と低位の各二点対には一つの四点軌道しか衝突しない。
稠密開部の異なる二十点から一般の付値曲線で特殊化し、平方根を分岐後に取ると、
各衝突点に元の一点と動く二点が来る。特殊側も長さ20なので局所長はちょうど3である
（付値環上の一般ファイバーの閉包は有限平坦）。低位点 $(1,1,0,0)$ では
核を $(0,0,1,-1)$ として一次化すると $dv=0$ と
$(M+f)ds=(G+n)dt=0$ が得られる。少なくとも一係数が非零である。
高位点 $(0,0,1,i)$, $i^2=-1$ では薄い式で $dt=0$ と低位核の比の微分が決まる。
二つの高位核座標の変分について残る二式の係数は
$(i(m+N),n+M)$ と $(-i(f+G),g+F)$ で、$u,v$ の変分はこれらの座標から決まる。
$d+h=0$ なら第一行は $(1-1/\beta)(im,n)$、$a+e=0$ なら第二行が
$(1-1/\beta)(-if,g)$ となる。他の符号でも同様なので接次元は1である。

$\beta=-1$ では $I J K N_0\ne0$ の場合、$Z^2=-mN/(nM)$,
$X^2=m(M+f)/M$ とする。$M+f=(D/k)/D$ は符号を除いて $1/k\ne0$ である
（$D=k(a+d-2k)$）。奇点 $(X,MZX/m,1,Z)$ での評価の階数は2である。
実際 $FN=-fn,GM=-gm$ で後ろの厚い二式は先頭の二式の倍数となり、
薄い二式にも同じ代入をすれば両方がその包に入る。

最後に $I J K N_0=0$ では、両角の交換、両対角の同時交換、$b$ の符号変更を使い
$a=b$ としてよい。これらは薄い標準語を保つ付け替えと倍率である（下の商の計算）。
$D\ne0$ なので $b(d-b)\ne0$。円錐の式により $e=b$ または $h=-b$。
前者では $f=F=0$, $m=M=1/b$。$x=(r,0,1,0)$, $r^2=1/b$ で行列が階数2になる。
後者では $f=N=0$, $m=-G=1/b$ なので、同じ座標の奇点で転置評価が階数2になる。
定理\ref{thm:single-opposite-evaluation-ranks}で全てが排除される。
\end{proof}

\begin{corollary}\label{cor:unequal-conic-unmarked-quotient}
一般の射影安定群は文字の四元群である。この族の一般の無印商は有理三次元である。
$b(a-d)\ne0$ の一般点で
\[
 x=(a+d)/(2b),\quad y=(e+h)/(2b),\quad z=(a-d)/(2b),\quad
 A_0=y/x,\quad B_0=\beta/A_0,\quad C_0=\beta z^2
\]
と置く。$A_0,B_0,C_0$ の各々に $t\mapsto(t-1)/(t+1)$ を施した値を $r,s,t$ とすると、
三つの無印座標は $(r+s)^2,(r-s)^2,(r+s)(r-s)t$ である。
特別な全ての評価適合点にも、独立四点の組と文字倍率による有限の同型比較が適用できる。
\end{corollary}
\begin{proof}
$(a,b,d,e,h)=(5,2,11,7,37)$ の素数761での偶配置・奇配置に対し、
上の表の平方根と一意な核を取る。両方とも相異なる二十点である。
三点の従属と平面上の点数で点を分け、独立四点と全係数非零の第五点の像を列挙すると、
射影自己写像数は偶で8、奇で4である。両方で一平面の最大点数は6であり、
任意の直線上の点数も6以下である。従って一様重みの配置は安定である。
具体的には、付値環上で異なる二十点を分離しておくと、射影写像の行列の原始な特殊化が
階数 $j<4$ なら核の外の全点が $j$ 次元線型空間の射影に入る。
核内と像内の点数の上界はそれぞれ $5(4-j)$ と $5j$ より真に小さいので不可能である。
従って一般の自己同型候補の置換もこの標本に単射で特殊化する。
全窓の自己同型は奇位置の配置に忠実に作用する。実際一点の相手は一意で、
各位置の配置が射影枠を含むからである。文字群が位数4を実現するので一般の群はこれに等しい。

従って一般点の無印同型は必ず文字群を運ぶ。四位置の四固有線の付け替えと
文字倍率だけを考えてよい。$y^2=\beta x^2+(1-\beta)(1-\beta z^2)$。
薄い二色のラベルを保ち規格化した語を移す置換を、節\ref{sec:full-thin-parallel-cover}と
同じ整数文字行列で解くと、$b$ による共通倍率除去の後の作用は16個である。生成元は
$(x,y)\mapsto(-x,-y)$, $z\mapsto-z$ に加えて
\[
 (x,y,z,\beta)\longmapsto(y,x,\beta z,1/\beta),\qquad
 (x,y,z,\beta)\longmapsto(-x/z,-y/(\beta z),1/z,1/\beta).
\]
前二者では書かなかった座標を固定する。
実際、低位と高位を動かさない置換では両角の交換、対角の同時交換と符号を得る。
全位置で高低を交換すればSchur補完が新しい上角となり、最後の式を得る。
再現計算は全位置移動と共通の文字せん断を直接列挙する。
最初の二符号の不変体は $A_0,B_0,C_0$ が生成する。実際 $y/x,\beta,z^2$ から
二次式で $x^2$ が復元できる。残る群は $(A_0,B_0,C_0)$ の初めの二数を同時に逆数にする操作と、
その二数の交換および最後の逆数を同時に取る操作である。Cayley変換の後は
$(r+s,r-s,t)$ の符号積を保つ四通りの変更なので、表示の不変体を得る。
例外値の比較には、全位置で元の配置の二つの固定点対が独立な四点を与える。
核の移送でも二つの固定対は隣の二座標直線上の二対に移る。
従って各位置で支持の独立四点を選び、前節の対角倍率の整数文字判定で四窓全体を比較できる。
これは各閉点の全次数正則性を仮定せずに同型を判定する。
\end{proof}

% !TeX root = ../papers/nc-p3.tex
\section{点と二平面の巡回からの全次数判定}
\label{sec:point-line-flag-cycles}
四線型テンソル $w$ の隣接二対二切断がともに階数6、一文字切断が全階数とし、
各辺の六次元像 $R_i\subset V_i\otimes V_{i+1}$ で定まる二次代数を $A$ とする。
両向きの非零評価が全て階数3以上で、どの $R_i$ も階数1テンソルを含まないとする。
以下の点は関係で結ばれた両無限点列に属するものを使う。
位置 $i$ の点 $p$ の消滅超平面を $H(p)$、前点を $p^-$ と書く。
$U\subset H(p)$ が二平面なら、$U^\perp$ で $p$ から独立な $q$ を一つ取り、
\[
 E(U,p)=\operatorname{supp}_{\rm left}w(p^-,q,-,-)\ \subset V_{i+1}
\]
とする。$w(p^-,p,-,-)=0$ のため選び方に依存しない。

\begin{theorem}\label{thm:recurrent-point-line-flags}
各位置で有限個の二平面を用意し、以下の二条件を満たすとする。有限周期に展開してよい。
\begin{enumerate}
\item 各二平面 $U$ に $U\subset H(p)$ となる指定点があり、$E(U,p)$ は次の位置で
用意した二平面の一つである。
\item この集合から各位置で $U_i$ を選んだ周期列があり、
$R_i\cap(U_i\otimes U_{i+1})\ne0$ が全ての位置で成り立つ。
\end{enumerate}
このとき $A$ は全次数で標準Koszul左右AS正則である。
有限周期の指定点列と二平面の遷移先を固定したこの判定は、上の評価下界と階数1関係の不在を
保つ特殊化の下で閉じている。
\end{theorem}
\begin{proof}
点 $p$ の右点加群を $M(p)$ とする。$H(p)$ の右イデアルはその全関係イデアルで、
$H(p)\otimes e_{i+1}A\to e_i A$ の核は
$K(p)=R_i\cap(H(p)\otimes V_{i+1})$ の三元で生成される。
命題\ref{prop:commuting-three-point-graph}の点上のTor計算を用いるが、
ここで必要な部分は標準有限双対をまだ仮定せずに通る。
実際、終点単純の左極小解消の最初の三項は矢と二次関係で決まり、次の項に
$w$ の一文字縮約の四次元空間を三次の極小重なりとして含められる。
点加群にテンソルしたこれらの射の階数は内部次数3以上で $1,3,3$ である。
二番目は関係の点評価、三番目も反対の隣接切断の完全対合を通した点評価の転置である。
従って第一Torは次数1で三個、第二Torは次数2の三個の商となる。
実際に $K(p)$ の三元が独立な二次関係を与える。

$h=w(p^-,-,-,-)$ の最終文字の縮約は $K(p)$ を張る。
その像は $K(p)$ に含まれ、最終文字側での階数3は $p^-$ への前辺の評価下界と
相補切断の完全対合から従う。$q$ で $K(p)$ を評価した像は $E(U,p)$ である。
$L_i(U)=e_i A/UA$ とし、$H(p)/U$ の生成元で $q$ が1を取る持上げを $v$ とする。
上記の三関係を $q$ で読むと $(UA:v)=E(U,p) A$、従って
\[
 0\longrightarrow L_{i+1}(E(U,p))(-1)\xrightarrow{\ v\ }L_i(U)
 \longrightarrow M(p)\longrightarrow0 .
\]
任意の長さについて条件1を反復すると $\dim L_i(U)_n=n+1$ である。
さらに $UA$ の関係はただ一つの二次関係 $r\in R_i\cap(U\otimes V_{i+1})$ で生成される。
実際 $K(p)$ の $q$ 評価の核は $kr$ であり、その基底を $r,s_1,s_2$ に延ばす。
$h=r\otimes c_0+s_1\otimes c_1+s_2\otimes c_2$ とすると
$r'=q(s_1)\otimes c_1+q(s_2)\otimes c_2$ が次の平面の非零関係である。
非零性は $p^-$ の唯一の後続が $p$ であることによる。
所与の $U$ の関係を $ra_0+s_1a_1+s_2a_2$ と表し、一つ短い長さで帰納仮定を
$E(U,p)$ に使うと $(a_1,a_2)=(c_1b,c_2b)$ となる。
$h\in H(p)\otimes R_{i+1}$ だから元の関係は $r(a_0-c_0b)$ である。

条件2の各 $r_i$ は純関係がないので階数2である。$U_i$ と $kr_i$ による抽象的二文字代数 $P$ は
位置ごとの基底で通常の交換子表示になり、長さ $n$ で次元 $n+1$ を持つ。
$L_i(U_i)$ の基底を持ち上げ左から $P$ の語を掛けると $A$ を張るので、各成分のHilbert上界は
$(1-z)^{-4}$ である。固定終点の経路次元を $a_n$ とすると上の一関係表示から
$a_n-2a_{n-1}+a_{n-2}=n+1+e_n$, $e_n\ge0$ を得る。
$(1-z)^{-2}$ の係数により上界で等号である。従って持上げは左 $P$ 自由基底を与え、
どの $U_i$ の任意の非零文字による左乗法も単射である。

任意の両無限点列について、点加群の右表示を三次の列 $h$ で延ばす。
得られる線形列の階数は $1,3,3,1$ で、既に項1まで完全である。
上の $c_0,c_1,c_2$（一般の点では $K(p)$ の任意の基底についての三係数）は独立である。
この三平面は同じ位置の $U_{i+2}$ と交わるので、三次の列の少なくとも一つの線型結合が
単射な左乗法となり、最後の列は単射である。点加群のHilbert関数と得られた標準Hilbert関数から
中央も完全となる。従って点加群とその次の点加群はともに線形解消を持つ。
後者を一つ移動して前者に入れた商は始点単純なので、錐によるその解消も線形である。
全始点でのKoszul性から双対の次元は $1,4,6,4,1$、$w$ とその縮約は全重なりとなる。
一文字の全階数と隣接切断の階数6で相補次数の積対合は完全なので、左右AS条件が従う。

最後に曲線上の付値環で、一般側の点、全平面、条件2の関係直線を各射影空間・Grassmann多様体に延ばす。
点列と関係の包含は閉条件である。$U\subset H(p)$ で独立な $q$ を特殊側で選び直し、
局所自由な $U^\perp$ に持ち上げる。一般側で $w(p^-,q,-,-)$ の左支持が次の平面に入るので
特殊側でも入る。$q\ne p$ と一意な隣接点から非零で、純関係がないため階数2のままである。
全ての遷移が残り、条件2の非零関係直線の特殊化も階数2を保つ。付値判定または射影像の閉性で従う。
\end{proof}

\begin{corollary}\label{cor:rank-ten-sheets-complete-regularity}
節\ref{sec:equal-determinant-second-sheet}の $Dpw\ne0$ の第二補完では
$(a^2-b^2)(d^2-c^2)(a^2-c^2)(d^2-b^2)\ne0$ が標準Koszul左右AS正則性の必要十分条件である。
節\ref{sec:unequal-determinant-conic}の $D\beta(1-\beta)\ne0$ の円錐曲線族では
$(1+\beta)I J K N_0\ne0$ が必要十分である。ともに三重衝突を含む。
従って $z=z\lambda^2=-1$ の可逆な薄い繊維では、全表示の各閉点の正則性が決定される。
\end{corollary}
\begin{proof}
必要性は両節の評価障害である。両厚色で関係空間は低位と高位への射影がともに同型であり、
各薄い色は全成分が非零である。階数1の関係があれば組をまたぐ二部分行列の行列式はともに零となる。
二薄色の係数を $x,y$ とすると、符号・単元を除いて $x^2-y^2$ と $x^2+y^2$ なので両係数が零である。
残った対角二ブロックの片方は零となり、二厚色の同型射影によりもう片方も零である。
逆辺にも同じ計算が使える。

\path{computations/verify_point_line_flag_cycles.py}は両枝の関数体上で以下の有限証明書を検算する。
等行列式では $b=B^2,c=C^2$ として位置1の表から五点を取る。
点名 $1,4,9,10,15$ は、$a=2,B^2=3,C^2=5,d=7$ の素数109での座標が順に
$(0,0,1,1),(0,1,47,26),(1,0,59,10),(1,0,50,99),(0,1,62,83)$ となる符号を指定するだけである。
他の位置の点は四文字縮約の一意な核へ順に移す。
位置0,2で、前々位置の点4と前位置の点1による縮約の左支持を $U_i$、右支持を $U_{i+1}$ に取る。
位置1,3では同じ五点の収縮のうち両支持が所定の二平面となるものを取れる。
これが条件2である。条件1の指定点を、各 $U_i$ から順に
\[
 (1,9,10,15,4),\quad(15,4,10,9),\quad(1,9,10,15,4),\quad(15,4,9,10)
\]
と選ぶ。五歩列では最後が一歩後の平面に戻り、四歩列では最初に戻る。
従って平面の和集合に必要な遷移が存在する。計算は各 $4\times4$ 縮約の二行の外積と
三行の外積で階数と平面の一致を整数多項式として確認する。

円錐曲線族では共通倍率でまず $b=1$ とし、$l^2=-I/J$ と取れば
$e=(kd-1)/(d-k)$, $a=((e-1)-l^2(e+1))/((e-1)+l^2(e+1))$ として
$k,d,l$ が独立なパラメータになる。位置0で二つの点対を使う。
第一対は $v=0$ 面の $(1,0,j,lj),(1,0,-j,-lj)$、第二対は $u=0$ 面の
$(0,1,v,Lv),(0,1,-v,-Lv)$ とし $L=-(k-1)/((k+1)l)$ である。
$j^2=J/(2(a+e)),v^2=N_0/(2(d+h))$ である。位置別に各対の割線の消滅二平面を作る。
偶では第一対、奇では第二対を取ると、各辺で六関係にこの二平面の積四式を加えた階数は9である。
従って条件2が成立する。どの位置のどちらの割線からも、その二点を順に使う二歩で
二位置後の他方の割線の消滅二平面に移る。これを繰り返して条件1を得る。
この二歩の検査は、使う対ごとに共通倍率で $j=1$ または $v=1$ として行える。
他方の対の二平面は二点の低位部と高位部の各直線で決まり、平方根を要しない。
従ってここでも三独立変数の多項式恒等式で十分である。
二つの既約表示に定理の特殊化を適用すれば評価障害の補集合の全閉点が正則となる。
最後に等行列式の $p w=0$ と正負の交換表示は既出の全点判定を持ち、その他の表示は
節\ref{sec:full-thin-parallel-cover}の移動・捻りで以上に戻る。
\end{proof}

% !TeX root = ../papers/nc-p3.tex
\section{対合が交換する四点と七次元成分の全次数判定}
\label{sec:involution-quartet-component}
四本の四周期点列が各位置で独立とし、その点座標を $0,1,2,3$ の座標点とする。
従って同じラベルが巡回して隣接する係数は零である。残る84係数が非零である開集合を取る。
対合が両対 $I=\{0,1\},J=\{2,3\}$ の各々を交換するとし、各基底の倍率を合わせて
四位置で $(01)(23)$ とする。テンソルを対合の正固有空間に置く。
四つの関係空間で対合の正・負部分の次元は $4,2$ とする。

\begin{lemma}\label{lem:four-track-opposite-minors}
任意の反対位置で、一方に $I$ の点、他方に $J$ の点を入れた二文字縮約が全て階数1とする。
階数1関係がなく、四点列での左右の関係評価が各々階数3なら、窓は標準Koszul左右AS正則である。
\end{lemma}
\begin{proof}
各位置で一対の点の消滅二平面を $U$ とする。その対の一方を $p$、他方を $q$ と書く。
四次元の正関係を $p$ で評価すると後続点の消滅超平面に入る。
従って正関係の一つ $r\ne0$ が $p$ を消し、対合で $q$ も消す。
点 $p^-$ で縮約した三テンソルの末尾縮約は $K(p)$ 全体を張るので
$r=w(p^-,-,-,l)$ と書ける。$w(p^-,q,-,-)$ の階数は2である。
実際、$K(p)$ の $q$ 評価は $r$ を消し、点評価階数3・二文字の完全対合によりこの縮約は非零の関係である。
次位置の他方の対のどちらの点 $b$ に対しても $w(p^-,-,b,-)$ は階数1である。
その $q$ 行は零でない。四点が座標点の表示で許された成分が一つ以上あり、全て非零だからである。
$l$ はこの行を消すので行列全体を消す。よって $r$ の右支持は他方の対の消滅二平面であり、
両対を交互に取ると定理\ref{thm:recurrent-point-line-flags}の条件2になる。

一歩の後続は $U'=\operatorname{supp}_{\rm left}w(p^-,q,-,-)$ で $q$ の後続点を消す。
二位置後の他方の対の点 $c$ では $w(q,-,c,-)$ が階数1で、その $p^-$ での非零列は $U'$ に入る。
従って $U'^\perp$ の任意の元 $s$ は $w(q,s,c,-)=0$ を満たす。
$s$ を $q$ の後続とは独立に取れば、二歩の後続関係は非零で階数2、左支持は二位置後の他方の対の消滅二平面に等しい。
これで条件1も得る。定理の証明で評価下界を使うのは指定点列上での三生成表示、前後の核の一意性、
および最後に単純を解消する一つの点列だけであり、全てこの四点列から取れる。
\end{proof}

\begin{theorem}\label{thm:involution-quartet-component}
反対階数16の点 $t$（命題\ref{prop:klein-rational-full-opposite}）を通る、
対合付きの七次元の局所成分がある。一般点は上の四点と階数1縮約の表示を持つ。
この成分の基底変更による像の閉包では、二隣接階数6、一文字全階数、階数1関係の不在、
全位置の左右評価下界3を満たす全ての閉点が標準Koszul左右AS正則である。
特に五次元楕円補完の既約族全体の閉包にこの含意が成立する。
\end{theorem}
\begin{proof}[四点表示の局所モデル]
同時に1を加える対合で残った42係数を使う。$w_{ijkl}$ は第一ラベルが0または2の方で表す。
左二文字・右二文字をいずれも $01,23,02,03,20,21$ の順の六軌道で取る。
開始位置 $s=0,1$、符号 $\epsilon=\pm1$ ごとに
$w_{ijkl}+\epsilon w_{ij,k+1,l+1}$ の六行列を $M_s^\epsilon$ と書く
（加算はビット排他的論理和、$s=1$ では四位置を巡回する）。
前半二行・列を一組、後ろを二行・列ずつ二組とする。
反対階数1の条件は前後に残る二ラベルの二行列式が零であることなので八つの二項式になる。
次の表の記号を使う。
\[
\begin{array}{c|rrrrrrrrrr}
 &P&Q&A&B&C&D&E&F&G&H\\ \hline
 w&0101&2323&0123&0132&2301&2310&0231&0321&2013&2103\\
\end{array}\qquad
\begin{array}{c|rrrrrrrr}
 &i&j&k&l&m&n&o&r\\ \hline
 w&0102&0103&0121&0131&0201&0301&2010&2101
\end{array}
\]
八つの二項式は別々の成分 $2023,2123,0323,0232,2320,2321,2313,2303$ を消す。
例えば $w_{0102}w_{2023}=w_{0132}w_{2013}$ である。
全体倍率と、各位置の二対の相対倍率で $P=Q=A=E=l=1$ とできる。
実際、全体を $1/P$ 倍し、位置 $2,3,1,0$ の順に相対倍率を決めればよい。

七変数 $x,X,u,U,B,F,k$ を取り、次の値を入れる。
\[
\begin{aligned}
 C&=(1-x)(1-X),&D&=xX/B,&G&=uU,&H&=(1-u)(1-U)/F,\\
 i&=Bu,&j&=(1-u)/k,&m&=x,&n&=F(1-x)/k,\\
 o&=XU/B,&r&=k(1-X)(1-U)/F .
\end{aligned}
\]
これらの変数は表示の係数から復元できる。実際に $x=lm/E$, $X=io/G$,
$u=il/B$, $U=mo/D$ である。
両 $M_s^-$ の先頭二行列を可逆とし、そのSchur補を取る。後ろ四行列の二つの交差ブロックは
この代入で恒等的に零となる。対角二ブロックの二切断合わせて16式は、残った交互配置
（例えば $0202,0213$）の16成分について定数係数の一次式で、係数階数は12である。
辞書式の自由座標を $w_{0312},w_{0313},w_{2130},w_{2131}$ とする。四つの整合条件を $f_j$ と書く。
第一式から単元を除けば $U$ について一次であり、第二式は $k^2$ について一次である。
これらを順に解くと残り二式は恒等的に零となる。従って
$x,X,u,B,F$ は独立であり、$k$ の二次拡大を許せば負切断の全式を満たす。

次に $M_s^+$ の先頭二行列のSchur補は二つの対角二行列からなり、交差ブロックは零である。
一方の対角行列には最初の自由二座標 $v,w$ だけ、他方には残る二座標だけが現れる。
各々の対角は $2v$ に定数を、非対角は $2w$ に定数を加えたものである。
第二切断では両組が交換され、対応する二行列式は二次項、一次項、定数項まで一致する。
従って両符号の必要な階数を得るため、各対の二座標についてちょうど一つずつのアフィン二次曲線を取ればよい
（各二行列が零になる閉集合を避ける）。これが七次元の表示である。
\path{computations/verify_involution_quartet_component.py}は上の二つの一次消去、四ブロックの
Schur補、二次曲線の一致を五変数有理関数体とその二次拡大で検算する。
この消去はこの表示を構成するためのもので、他の根の組合せを全てこの表示に含める必要はない。
\end{proof}

\begin{proof}[元の階数軌跡への支配と特殊化]
点 $t$ の素数2657での二十点のうち、位置0の
\[
 (1,2197,1186,92),\quad(1,2197,1471,2565),\quad
 (1,1667,2322,492),\quad(1,1667,335,2165)
\]
を順に取る。各対は高位二座標の符号を変える対合で交換される。
前節までの計算で点を四位置に移すと元に戻る。
これらを座標枠にすると84係数は全て非零で、八二項式を満たす。
再現計算では全42係数を上記表示から復元し、$U,k^2$ の二主係数と全ての使用分母が単元であることも確認する。

この剰余点の近傍で、まだ八二項式を課さずに二つの隣接切断の階数6条件を微分する。
各切断の左核と右核の間への微分の射影を、42座標について縦に並べた行列の階数は30である。
従って $\mathbb Z_{(2657)}$ 上でこの30個の独立な式だけを課した周囲は剰余点で相対次元12の滑らかな空間である。
上で構成した部分はそこに、五つの倍率を戻した相対次元12で入る。
次元の下界はこの点自身で成立する。五独立座標と二次根は主係数可逆の開部分にあり、
残る四変数には二式しかない。係数から全変数を復元する局所逆もある。
特に付加条件を課した局所モデルの絶対次元は少なくとも13で、上の正則局所環と等しい。
正則局所環は整域なのでこの中に同次元の真の閉部分はなく、階数軌跡と表示は局所的に一致する。
この議論は有限体の標本から標数零の恒等式を推測するものではない。
四点を忘れたとき一般の選択は有限で、枠を固定した後の連続な同一視は上の五倍率だけである。
点配置の安定群も標本で有限なので一般点で有限であり、基底変更を除いた次元は7となる。

五次元族の標本 $t$ は標数零から同じ剰余点に特殊化する。
単純な点配置を局所エタールに番号付けし、上の二点と対合による二つの像を取れる。
四移動は単純な点の番号の置換で、剰余点で恒等であり、点列の四周期性が保たれる。
標数零の一般点の点対応と四周期性は定理\ref{thm:exact-two-configurations}の証明からも従う。
従って五次元族の一般点のこの表示が上の開近傍に入り、既約性でその全閉包が七次元成分の
基底変更像の閉包に含まれる。一般点の評価下界・階数1関係の不在はそれらが成立する $t$ から開性で従う。
補題\ref{lem:four-track-opposite-minors}の巡回旗と定理\ref{thm:recurrent-point-line-flags}の固有性で、
表示が消える境界でも定理に述べた全閉点の正則性を得る。
\end{proof}

\begin{corollary}\label{cor:full-thin-finite-graph-regularity}
全成分非零の薄い二色を持つ平行切断の五次元主閉包では、純関係がなく、
四辺が有限であれば、一文字全階数・隣接階数6の候補は全て標準Koszul左右AS正則である。
特に混合交換表示でも、以前の四中心式で認証した開集合の外までこの結論が通用する。
\end{corollary}
\begin{proof}
片側の評価が階数2以下の点があれば、その点の相手が少なくとも直線をなし有限性に反する。
従って左右評価の下界を満たす。
\end{proof}

% !TeX root = ../papers/nc-p3.tex
\section{三つの水平長方形の特殊値を含む正則性}
\label{sec:horizontal-rectangle-flags}
\begin{theorem}\label{thm:horizontal-rectangle-flags}
標数零で次が成立する。
\begin{enumerate}
\item 76型の水平台85、112型の水平台143--147では、標準Koszul左右AS正則性の
必要十分条件は $b+c\ne0$ である。反対階数の制限はない。
\item 第三長方形の台88--90では、定理\ref{thm:parallel-third-rectangle}の
$b(b-1)AD(AD-1)\ne0$ の全点が標準Koszul左右AS正則である。
\end{enumerate}
従って台85では定理\ref{thm:first-rectangle-horizontal-frames}の三変数の開集合が、
第三長方形では定理\ref{thm:third-rectangle-orbits}の $(X,Y,h)$ の開集合が、
指定始点の正則同型類を例外なく表示する。そこでの最小周期の式も全点に適用される。
\end{theorem}

\begin{proof}[五次元固定成分への包含]
定理\ref{thm:rectangle-unrestricted-components}の水平標本三点の証明を、
四元群の作用を固定した軌跡で用いる。各標本でこの軌跡は滑らかで次元18であり、
薄い零行・零列が全て消え、厚い交差ブロックが全て可逆になる変形がある。
対角軌道の次元は13であるから、唯一の局所固定成分の一般点は
定理\ref{thm:klein-all-pairings}の五次元楕円補完に入る。
元の三つの水平族の各々は既にこの固定軌跡にあり、既約で、該当する標本を通る。
従って三族の全閉包が、基底変更を許した五次元楕円補完の閉包に入る。
定理\ref{thm:involution-quartet-component}により、以後は各辺で純関係がなく、
左右の評価階数が全て3以上であることを調べれば十分である。使う台では隣接切断の色階数は
$(2,2,1,1)$、一文字は全階数である。$b+c=0$ の最初の二族の排除は
命題\ref{prop:parallel-antidiagonal-obstruction}での四次経路次元の障害による。
\end{proof}

\begin{proof}[開集合全体での評価と純関係]
76型の十分性では定理\ref{thm:first-rectangle-horizontal-frames}のトーラス開集合を用いる。
112型では定理\ref{thm:horizontal-rectangle-low-fibers}により、使う五台で共通して次が可逆である。
\[
 uv(u+v)(1+uv)(zw-1)
 \prod_{q\in\{z,w\}}(v+q)(q-u)(uq-1)(vq+1).
\]
以下の計算はこの積だけを逆にするので、その開集合のどの点にも通用する。
両型とも一位置移動 $J$ が該当する開集合を保ち、移動後のテンソルの
全係数を与えるので、位置0の両向きだけを調べればよい。第三長方形では四位置を直接調べる。

各色 $h$ について切断行列の列空間の基底を列に持つ $4\times d_h$ 行列を $B_h$ とする。
指定した基底開集合で単元となる $d_h$ 次小行列を使えば、$B_h$ の対応する行ブロックを単位行列にできる。
全ての切断列をこの基底で復元して等式を検算する。
最初の二型の薄い列を $c,d$ と書くと、純関係の二つの交差ブロックの行列式がともに零であるため
\[
 \begin{pmatrix}c_0c_1&-d_0d_1\\c_2c_3&-d_2d_3\end{pmatrix}
 \binom{\alpha^2}{\beta^2}=0
\]
となる。行列式は単元なので両薄色が消える。残ったブロック対角行列が階数1なら
片側の二行ブロックが零となる。両厚色とも各側二行への射影が同型なので全て零である。
第三長方形では $x\otimes y$ の色成分 $c_h=(x_j y_{j\oplus h})_j$ を作り、
$B_h$ の単位行ブロックの添字を $I$ として $c_h=B_h(c_h)_I$ を課す。
$x,y$ の先頭非零座標をそれぞれ1とした十六開片全てで、基底の非零積による飽和は単位イデアルとなる。

左右評価には、非空の座標台 $S$ ごとに $x_{\min S}=1$ とし、$S$ の他の座標も可逆にする。
一つの関係列 $c$ に対応する評価の行は、一方では $(c_{j\oplus h} x_{j\oplus h})_j$、
他方では $(c_j x_{j\oplus h})_j$ である。六行四列の三次小行列式を課す。
単元の二次小行列があればそれを含む縁の三次小行列式だけで十分である。
各小行列式の分子を因数分解し、可逆な因子を除いた積を零にする最小因子集合ごとに
飽和する。十五台のいずれでも左右とも消滅集合は空である。
\path{computations/verify_horizontal_rectangle_flags.py} は三つのパラメータ環上でこの全検査を行う。
従って巡回旗の特殊化による全次数判定が定理の全域で成立する。
\end{proof}

% !TeX root = ../papers/nc-p3.tex
\section{七次元対合成分の順序付き有理商と一般安定群}
\label{sec:rational-ordered-quartet}
\begin{theorem}\label{thm:rational-ordered-quartet}
定理\ref{thm:involution-quartet-component}の成分で、対合により第一点と第二点、
第三点と第四点が交換される順序付き四点列を、その表示の通り指定する。
この印付き成分は既約で、指定始点の同型類の関数体は七変数の純超越拡大である。
一般の無印窓の射影安定群は指定対合の生成する位数2の群であり、最小周期は4である。
\end{theorem}

\begin{proof}[二次曲線と倍率の消去]
節\ref{sec:involution-quartet-component}の $x,X,u,U,B,F,k$ を用いる。
整合条件の第一式から $U$ を消した第二式を $k^2=N(X)/D(X)$ と書く。
分子と分母はともに $X$ について一次で、その係数行列式は非零である。
従って $x,u,B,F,k$ を五つの独立座標にでき、$X,U$ はこれらの有理式になる。
この置換の逆は標本 $t$ でも定義される。
各自由二座標 $v_i,w_i$ の行列を
\[
 \begin{pmatrix}2v_i+a_i&2w_i+b_i\\2w_i+c_i&2v_i+d_i\end{pmatrix}
\]
と書き、$V_i=4v_i+a_i+d_i$, $W_i=4w_i+b_i+c_i$ とする。
その行列式の消滅は
\[
 (V_i+W_i)(V_i-W_i)=(a_i-d_i+b_i-c_i)(a_i-d_i-b_i+c_i)
\]
である。右辺は標本 $t$ の素数2657で順に $2218,1246$ となりともに単元である。
従って各曲線が開トーラスになり、五つの基底変数を含むこの稠密表示は既約である。
前節の $t$ での局所一致と、そこで全逆変換が存在することから、これが所定の成分を与える。

順序を固定した四点枠を保つ基底変更は位置別の対角行列である。
射影対合の行列を再び $(01)(23)$ にするため、各位置で両対内の倍率比は共通の符号となる。
隣接二位置の符号の積が負なら正4、負2の関係固有空間が交換されるため、四つの符号は全て等しい。
各位置スカラーを除く連結な倍率と全体倍率は $P=Q=A=E=l=1$ で一意に除かれた。
残った同時負符号は、この規格化に戻したとき各語に
$(-1)^{\sum_j (i_j\bmod 2)+\sum_j\lfloor i_j/2\rfloor}$ を掛ける。
従って $x,u,B,F$ を固定し $k$ を反転する。各ブロックの $v_i,a_i,d_i$ は偶、
$w_i,b_i,c_i$ は奇である。
\[
 s_i=\frac{V_i+W_i}{a_i-d_i+b_i-c_i},\qquad
 y_i=\frac{s_i-1}{k(s_i+1)}
\]
と置く。符号は $s_i\mapsto 1/s_i$ と作用するので $y_i$ は不変である。
逆に $k$ を一つ取れば $s_i=(1+ky_i)/(1-ky_i)$ から二座標が戻る。
よって所求の順序付き射影商の関数体は
$k_0(x,u,B,F,k^2,y_1,y_2)$（$k_0$ は基礎体）である。
\path{computations/verify_quartet_rational_chart.py} が一次分数式の次数・行列式、
二つの右辺と全係数の符号を検算する。
\end{proof}

\begin{proof}[印を外した自己同型]
$t$ で二十点が被約で、一般位置の五点から各位置の射影写像を一意に決められる。
ここで一般位置の五点とは、そのどの四点も独立であることをいう。
有限エタールに番号を付ければ、一般側の点配置間の射影写像は $t$ まで特殊化する。
実際、行列を射影空間に延ばし、階数が落ちた場合の核と像に入る点数を比べる。
$t$ では三点直線も五点平面もないので、階数 $r<4$ の行列では、
核の外の全点を相異なる像点に送れない。従って特殊化行列も可逆である。
隣接辺のグラフから一位置での番号の置換は全位置の置換を決定する。
系\ref{cor:klein-four-configurations}により、安定群は一般側で $t$ の四元群の部分群であり、
異なる位置の配置間には同型がない。従って一般周期は4である。

一般側で四元群全体が残ると仮定する。付値環上で二つの射影対合を線型に持ち上げ、
単元の平方根を加えて二乗を1とできる。交換子の符号は特殊側で正なので一般側でも正である。
四文字の多重度は各位置で独立な四つの文字のままであり、固有射影で局所的に文字枠を取れる。
隣接関係の色階数も固有空間の束の階数で一定、薄い全成分と可逆な交差ブロックは
$t$ で非零なので一般側でも非零である。
定理\ref{thm:klein-all-pairings}によればこの四元作用を持つ窓の無印像は次元5以下である。
これは前節の七次元性に反する。指定対合が残るので一般安定群は位数2となる。
\end{proof}

% !TeX root = ../papers/nc-p3.tex
\section{十五次元成分の三次曲面対と残余八点}
\label{sec:marked-cubic-surface-pairs}
\begin{theorem}\label{thm:marked-cubic-surface-pairs}
十五次元成分の稠密開集合で、二十点のいずれを印にしても以下が成立する。
位置 $0,2$ を反対対とし、印を入れた二つの三文字テンソルから、同じ
$\mathbb P^2\times\mathbb P^2$ 内の順序付き二曲面を得る。いずれも滑らかな三次曲面で、
二つの平面への射影を通じて五次の平面Cremona変換のグラフになる。
両曲面は長さ27で交わる。そのうち19点が元の印を除く点列の射影で、各々交点重複度は1である。
残余は長さ8である。二曲面の順序と二つの平面の順序を記録した射影同値類は、
その印を持つ窓の添字保持同型の完全不変量である。もう一つの反対対にも同じ結論が通る。
\end{theorem}

\begin{proof}[二曲面とその逆構成]
定理\ref{thm:marked-point-commutators}の四組のうち $G,F$ を使う。$T=\langle I,G_1,G_2,G_3\rangle$,
$S=\langle I,F_1,F_2,F_3\rangle$ は同じ三行三列の行列空間に入る二つの四平面である。
各 $E=T,S$ について $\det Z=0$ を $\mathbb P E$ 内の三次曲面とし、
$Z$ の左核 $l$ と右核 $r$ を二平面の積に写す。用いる開条件とその非空性は下で検証する。
滑らかな三次曲面には階数1の点がない。$l,r$ をともに核に持つ行列の空間は
高々一次元である。二次元なら、その鉛筆の二行二列の行列式に射影的零点があり階数1以下になるためである。
従って像上で行列を線型方程式の唯一の核として復元でき、像への写像は同型である。

$l$ を固定すると $lZ=0$ は $E$ に三式であり常に解を持つ。
一般ファイバーは零次元なので一点である。右核も同様で、二つの平面射影は双有理である。
両引戻し直線束を $L,R$、三次曲面の超平面類を $H$ と書くと $L^2=R^2=1$。
余因子行列の二次式に基点はなく、左右核のSegre写像を与えるので $L+R=2H$。
$H^2=3$ により $LR=5$、したがって平面間の双有理写像の次数は5である。

逆にこの埋め込まれた曲面から $E$ を元の行列空間の部分空間として取り戻せる。
反標準束は $H$ であるから、それによる埋込みは元の $\mathbb P^3$ 表示を座標変更まで与える。
各点で指定された左右核を持つ一次形式行列をこの四つの射影座標で解く。
元の $M$ に加え $M'$ も解なら、$M'\operatorname{adj}M$ の各成分が曲面上で零になる。
各成分は三次式なので $M'\operatorname{adj}M=C\det M$ で $C$ は定数行列、従って $M'=CM$。
左核の一致条件は $l C$ が $l$ の倍であることを全ての $l$ について与える。従って $C$ はスカラーである。
一次形式の係数の張る四平面は反標準座標の選び方にも依存しない。

二曲面から復元した $T,S$ は可逆な一行列の直線で交わる。これを単位行列として二つの三次元空間を双対に同定すれば、
残る基底変更は共役と倍率である。両四平面の跡零部分を $T_0,S_0$ とする。
混合交換子 $T_0\otimes S_0\to\mathfrak{sl}_3$ の階数は8、核線はもとの表示で
$\sum G_i\otimes F_i$ である。これは定理\ref{thm:marked-point-commutators}の式の跡であり、
非退化なので両三平面の双対基底を共通倍率まで指定する。あとは得られた
$[G_a,F_c]$ に定理\ref{thm:marked-mixed-bracket-reconstruction}の逆写像を用いればよい。
両側単射の開条件は全二十印で成立する（定理\ref{thm:marked-twenty-rational-invariants}）。
この復元の基底変更と倍率は同定理の印付き同型作用に入る。逆方向は構成の自然性から従う。
\end{proof}

\begin{proof}[交点の長さと開条件の算術的検証]
第一平面の点 $l$ を変数とし、$E$ の基底 $M_0,\ldots,M_3$ の四行 $l M_i$ の符号付き三次小行列式を $a_i$ とする。
$Z_E(l)=\sum a_i M_i$ の右核が五次写像 $f_E(l)$ である。実際、$Z_E$ の二行の外積は
除いた行の $l$ 座標倍の同一の五次ベクトルとなる。これは $lZ_E=0$ と座標環の一意分解から従う。
$a_i$ の全てが零のときだけ第一射影のファイバーが正次元となる。

二曲面の第一射影の正次元ファイバーの基点は互いに素である。
さらに第一平面の一般直線上で二つの五次写像の像が一致する点はない。
従って曲面の交叉には第一平面の曲線を覆う曲線も、第一平面の一点上の曲線もない。
各曲面の二つの射影次数と混合次数は $1,5,1$ であるから、積のChow環での交点数は
$1+25+1=27$ である。滑らかな四次元多様体の中の滑らかな余次元2の二曲面の固有交叉なので、
これはスキームの長さである。

他の点列の位置1,3の座標を印から射影すれば、その二点は
$w(p_0,-,x_2,-)$ と $w(x_0,-,p_2,-)$ の双方の左右核である。
印を除く19個の第一射影は相異なる。一般位置での三点共線の不在は
系\ref{cor:klein-four-configurations}の標本とエタール特殊化から従う。
これらの点は両写像の定義域にあり、二つの五次ベクトルの外積の微分は各点で階数2を持つ。
従ってそれぞれ単純交点で、残余の長さは8である。

\path{computations/verify_marked_cubic_surface_pairs.py} は
定理\ref{thm:marked-twenty-rational-invariants}の実在する1009進持上げを14桁で用いる。
全二十印と四位置で三次行列式の偏微分のイデアルが五次斉次成分を張ることを検査し、滑らかさを得る。
二方向の各曲面対で、両方の $a_i$ が生成するイデアルは第一平面の八次斉次成分を張る。
直線 $l=(1,s,2+3s)$ に沿う二つの五次ベクトルの外積から二つの十次斉次二元式を取り、
そのSylvester行列が全階数であることを検査する。無限遠もこの斉次終結式に含まれる。
全ての点を持ち上げて他の19点での両写像の非消滅と二つの接方向の独立性を確認する。
四平面の和の次元7および交換子の階数8も全対で検査する。いずれも単元消去および
残余行列の共通1009因子による除算後に十分な精度で非零な小行列を得る検査である。
従って標数零のこの持上げで開条件が同時に成立し、既約成分の一般点にも成立する。
\end{proof}

\begin{corollary}
この稠密開部分では、定理\ref{thm:recurrent-point-line-flags}の指定点を持つ二平面と、
同じく指定点を持つその後続を同時に取れない。
\end{corollary}
\begin{proof}
出発点列の印を $a$、後続平面の印を $b$ とする。
$U\subset H(a_i)$ なので $U^\perp$ の独立な第二元を $q$ とすると、
$w(a_{i-1},q,b_{i+1},-)=0$ が必要である。$a=b$ なら反対行列は階数3で、
その左核は $a_i$ のみである。異なる場合は $a$ の滑らかな三次曲面の点なので
反対行列の階数は2で、左核は $a_i,b_i$ の張る平面である。
従って後続の左支持を与える行列は非零倍率を除いて $w(a_{i-1},b_i,-,-)$ である。
$b_i$ の印からの射影では上の三次小行列式 $a_j$ が全消滅しないので、この縮約は階数3である。
二平面を後続にすることはできない。
\end{proof}

% !TeX root = ../papers/nc-p3.tex
\section{単純点を持つ一辺の十九点模型と全延長の線型核}
\label{sec:nineteen-point-edge}
二つの三次元空間を用い、$X=\mathbb P(V^*)\times\mathbb P(U)$ の点を行 $l$ と列 $z$ で表す。
$T_1,T_2$ を両平面の接束とする。
\begin{theorem}\label{thm:nineteen-point-edge}
左右四文字・六二次関係の一辺の点対が有限で、単純な一点を指定する。
その印を保つ一辺の基底変更類は、二平面の積内の長さ19のスキーム $Y$ から復元される。
現れる $Y$ は
\[
 E_1=T_1\boxtimes\mathcal O(1),\qquad E_2=\mathcal O(1)\boxtimes T_2
\]
の切断を一つずつ取った同時零スキームで、零次元であるものにちょうど一致する。
非被約な $Y$ を許す。二つの切断は $Y$ からそれぞれ非零倍率まで一意である。
\end{theorem}
\begin{proof}
印をそれぞれの星座標点とする。単純点なので局所アフィン六変数について六双一次式の
一次部行列が可逆である。関係の基底を取り直し、座標を $(\gamma,l),(\alpha,z)$ と書けば六式は
\[
 l(\alpha I+F_z)=0,\qquad (\gamma I+Q_l)z=0
\]
となる。ここで $F_z$ と $Q_l$ は添字に線型な三次行列である。
$\alpha,\gamma$ に三座標の一次式を加えることで各々跡零とする。
Euler列によればこの二つの跡零行列系がちょうど上の二束の大域切断である。
星点以外では $l,z$ とも零でない。片方でも零ならもう片方も零となるからである。
二平面に射影した像は二つの固有ベクトル条件であり、固有値 $-\alpha,-\gamma$ は
各単元座標で割ればスキーム上でも一意に戻る。逆に任意の二切断からこの式を作ると星点は単純である。
星点の補集合が有限なら零スキームは階数4の束の正則切断となる。
二平面の超平面類を $h,k$ とすれば最高Chern類は
$(3h^2+3hk+k^2)(h^2+3hk+3k^2)=19h^2k^2$ である。

より強く $h^0(\mathcal I_Y E_i)=1$ を示す。$E=E_1\oplus E_2$ の正則切断のKoszul解消を $E_1$ でひねる。
$H^0(E_1\otimes E^*)$ は恒等写像だけである。実際平面接束の自己準同型はスカラーで、
異なる束間には $H^0(\Omega^1(1))=0$ を使える。
さらに $j=2,3,4$ で $H^{j-1}(E_1\otimes\bigwedge^j E^*)=0$ である。
$j=2$ の三項は順に
$T(-3)\boxtimes\mathcal O(-1)$、$\operatorname{End}(T)(-1)\boxtimes\Omega^1$、
$T(-2)\boxtimes\mathcal O(-2)$ となる。
$j=3$ では $T(-4)\boxtimes\Omega^1(-1)$ と
$\operatorname{End}(T)(-2)\boxtimes\mathcal O(-3)$、
$j=4$ では $T(-5)\boxtimes\mathcal O(-4)$ である。
Euler列とKünneth公式で所要の消滅が従う。
最後は三個の一次式による $H^0(\mathcal O(1))^3\to H^0(\mathcal O(2))$ の全射性を双対化する。
また $H^0(\operatorname{End}T(-1))=H^0(\operatorname{End}T(-2))=0$ は
Euler列、または接束の安定性から従う。よってイデアルに値を持つ切断は元の切断の倍数だけで、
二束を交換しても同じである。二つの倍率は星座標の独立な倍率で吸収される。
印を保つ射影変更は二平面で射影変更を誘導するので、$Y$ の順序付き二因子の射影同値が一辺の同型の必要十分条件である。
\end{proof}

\begin{theorem}\label{thm:nineteen-point-linear-companion}
この一辺を位置 $1,2$ に置く。両反対の印の縮約が階数3の四周期の印付き二切断補完は、
跡零の $F,Q$ を上の十九点から取り、次の線型核で表示できる。
\[
 K_{F,Q}=\{(G,P):[G_a,F_c]_{bd}+[Q_b,P_d]_{ac}=0\}.
\]
非零な組の共通倍率を除き、同じ $Y$ の安定群からの同一視だけが残る。
この表示では有限な標準Frobenius・Koszul左右ASとなるものへの条件を引き続き課す。
特に四辺が有限な標準Frobenius候補に単純な印があればこの表示で尽くされる。
十五次元成分の一般点の全二十印・四始点では $\dim K_{F,Q}=1$ なので、
十九点の配置だけで印付き窓そのものの添字保持同型類が決まる。
\end{theorem}
\begin{proof}
反対二行列を単位行列とし、星文字の持上げで残る二組も跡零にできる。
定理\ref{thm:marked-point-commutators}の二切断のSchur補の計算は、既に分かった $F,Q$ に
ついて読めば、48未知係数 $(G,P)$ の81本の一次式である。
解ごとに星のない四成分を同定理の積和で置けば、両切断がちょうど階数6のテンソルに一意に戻る。
$G=P=0$ の補完も式上は存在するが、反対の辺では星座標を両方とも零とした二平面の積が点対に入り四辺有限の候補から外れる。
非零解では、$F,Q$ を固定したまま位置0,3の三文字を $h$ 倍し、全体を $h^{-1}$ 倍すれば
$(G,P)\mapsto(hG,hP)$ となる。このほかの同一視は元の印付き一辺から誘導される。
具体的には $Y$ を保つ両平面の変換を持ち上げ、二切断の単位行列と $F,Q$ の規格化へ戻して核に作用させる。
持上げの違いは既に除いた共通倍率に過ぎない。
有限Frobenius候補の単純な印で反対縮約がともに階数3となることは
定理\ref{thm:finite-cubic-separators}である。
\path{computations/verify_nineteen_point_edge.py} は分裂混合窓の素数1009の横断二十点を用い、
全八十通りでこの線型写像の階数が47であることを確認する。上界は実際の非零解で従う。
エタールな点の近傍で小行列を持ち上げて一般点での主張を得る。
\end{proof}

% !TeX root = ../papers/nc-p3.tex
\section{七十一節点の双三次式による十五次元成分の局所分類}
\label{sec:marked-bicubic-nodes}
\begin{theorem}\label{thm:marked-bicubic-nodes}
十五次元成分上の印付き被覆のうち、例 $t$ の素数2657での第一位置の印
$(1,1,0,0)$ に達する既約成分を取る。その一般の窓で
\[
 D=\{\det[Q_l,P_r]=0\}\subset\mathbb P^2_l\times\mathbb P^2_r,\qquad
 D'=\{\det[G_x,F_z]=0\}\subset\mathbb P^2_x\times\mathbb P^2_z
\]
はともに71個の通常二重点だけを特異点とする双次数 $(3,3)$ の超曲面である。
各々36点では行列階数が1である。残る35点は二方向で標準的に対応し、
十九点の窓の点列と、二つの互いに素な八点集合に分かれる。
特に定理\ref{thm:marked-cubic-surface-pairs}の残余八点はこの一般点で被約である。

$D$ の71節点を保つ双三次式の被約局所成分は滑らかで、式の全体倍率を含む次元は32、
両平面の独立な射影変更と倍率を除いた一般次元は15である。
印付き窓からこの成分の商への写像は支配的かつ生成的有限で、一般点で分離的な局所座標を与える。
従って一般の局所変形は、この節点数を保つ双三次式の変形で追跡できる。
\end{theorem}

\begin{proof}[二方向の行列と節点]
四線型形式を $f(l,r,x,z)=x[Q_l,P_r]z=-l[G_x,F_z]r$ と置く。
四つの文字について全ての一階偏微分が零となる四つ組を考える。
$[Q_l,P_r]$ の階数が2ならその左右核 $x,z$ は一意で、
$\det[Q_l,P_r]$ の微分は余因子公式によって $f$ の $l,r$ 微分の非零倍率である。
従って $D$ の階数2の特異点はこのような四つ組に持ち上がる。
通常節点しかないときには、四つ組のいずれの行列も階数1以下にならない。
実際、階数0なら行列式が三重以上になる。
階数1なら単元を消した二行二列のSchur補の行列式が局所式である。
$x,z$ はその左右核の二空間に属し、$x\,d[Q_l,P_r]\,z=0$ が全方向で成立する。
これは四成分の一次部の非自明な線型関係であり、二次主要部は退化して通常節点でない。
もう一方についても同じなので、二つの階数2の特異点集合が左右核によって一対一に対応する。

定理\ref{thm:marked-cubic-surface-pairs}の二つの曲面対の交叉からこの四重臨界集合に点が入る。
例えば第一対で $G_x,F_z$ はいずれも $l,r$ を左右の同一固有値の固有ベクトルとして持つ。
従って $l[G_x,F_z]=[G_x,F_z]r=0$ であり、他方の二式は
$l[G_x,F_{z'}]r=l[G_{x'},F_z]r=0$ から従う。これは交点スキーム上の等式でもある。
どちらの交叉も対応する双三次式の特異スキームに含まれ、通常節点だけなら長さ27の交叉は被約である。
両方向の27点が同じ四つ組になる条件は、$Q_l,P_r$ についても $x,z$ がそれぞれ左右の
同一固有値の組になることである。四つの固有値を星座標に戻せば四辺の点対の全式が成り立つ。
逆も明らかで、重なるのは星点以外の19本の点列に限る。従って階数2の節点は少なくとも35点である。

一般点では交叉交換子は跡零八空間の同型である。
$[Q_l,P_r]$ の射影行列はSegre四次元多様体を可逆行列の点から跡零空間へ射影し、
線型同型を施したものになる。射影の中心は割線多様体の外なのでその像は次数6である。
階数1・跡零の行列の多様体はもう一つのSegreの超平面切断で次元3、次数6である。
$\mathbb P^7$ での交叉が有限なら長さ36、通常節点ではSchur補の四一次部が独立なので単純である。
よって節点の上界71と通常性が成立すれば、二方向に全ての点が上の記述で尽くされる。
\end{proof}

\begin{proof}[上界、節点条件の接空間と支配性]
上の印の剰余点で両行列式を直接計算する。各平面の先頭非零座標を1とした九つの座標片で、
元の四アフィン方向の微分を全て課した後に先行座標を零にする。
二方向とも特異点の長さは71、二次小行列式を加えた長さは36、
四行Hesse行列式も零とした軌跡は空である。
特異点の射影性と非退化Hesse条件から、局所エタールに各点では四偏微分を解いた
唯一の臨界点列があり、その点での式の値の消滅が節点が残る条件である。
従ってどちらも近傍で特異点は高々71、全て通常節点である。
標数零の一般点に前段の下界を適用すれば等号であり、71個全てが残る。
このため $t$ の該当する標数零の持上げにもその全てが特殊化する。

100次元の双三次式のアフィン空間で、節点存続式の微分は各節点での変分三次式の評価である。
$D$ の剰余の全節点環に100個の単項式を写した階数は68である。
従って七十一節点を保つ局所軌跡の接次元は32以下である。
一方、四つの跡零三行列組の96係数について交換子方程式の微分核は34次元である。
この印付き正規化表示が滑らかな34次元であることは、$t$ で滑らかな76次元の
二切断軌跡、単純点のエタールな選択、二つの反対行列と跡の規格化から従う
（定理\ref{thm:marked-mixed-bracket-reconstruction}と同じ枠の計数）。
剰余でも微分核が同次元なので接方向の基底を持ち上げられる。
これらの方向で $D$ の変分を計算し、さらに二平面の独立な基底変更と式の倍率を加える。
軌道方向だけで17、窓の変分だけで24、合わせて32個の一次独立な式が得られる。
従って窓と座標変更の像の次元は少なくとも32である。近傍ではその像が全節点の存続式を満たすので
先の接上界と一致し、局所環は正則、像はこの成分に稠密となる。
射影同値を除いた次元は $32-17=15$、元の印付き商と一致する。
\path{computations/verify_marked_bicubic_nodes.py} は両特異イデアル、評価行列、
および行列式の余因子微分を検算する。変分多項式の独立性の下界には130組の代入値を用いる。
\end{proof}

% !TeX root = ../papers/nc-p3.tex
\section{交差ブロック可逆な主枝の全点での正則性}
\label{sec:elliptic-evaluation-obstructions}
定理\ref{thm:klein-elliptic-completion}の全薄係数非零の規格化を用い、
$z\ne1$, $z\lambda^2\ne1$ とする。以下の同時文字置換、始点移動、左右反転は
判定を実行するために使う。分類する同型に始点移動や反転を加える必要はない。
\begin{theorem}\label{thm:elliptic-evaluation-obstructions}
五次元主族の閉包にあるこの表示の窓で、二隣接階数6、一文字全階数を仮定する。
標準Koszul左右AS正則となる必要十分条件は、四辺のいずれを左右いずれから
非零点で評価しても階数が3以上となることである。下界を破る表示は、
辺・向き・文字を選び直した後、次の二型で尽くされる。
\begin{enumerate}
\item $A=\left(\begin{smallmatrix}1&\lambda\\c&d\end{smallmatrix}\right)$,
$B=\left(\begin{smallmatrix}\lambda&f\\1&h\end{smallmatrix}\right)$
であり、節\ref{sec:klein-elliptic-completion}の二切断条件を課す。
\item $s,m\ne0$, $(1-s^2)(1-zm^2)\ne0$ とし、
\[
 T=\begin{pmatrix}1&zm\\m&1\end{pmatrix}^{-1}
       \begin{pmatrix}1&s\\s&1\end{pmatrix},\qquad
 D_0^{\mathsf t}=\operatorname{diag}(1,s)T\operatorname{diag}(1,m)^{-1},
 \qquad
 D_1^{\mathsf t}=\begin{pmatrix}0&1\\s&0\end{pmatrix}
                  T\begin{pmatrix}0&1\\m&0\end{pmatrix}^{-1}.
\]
$A=D_0^{-1}Y_\lambda$, $B=(D_1^{-1}\lambda X_\lambda^{\mathsf t})^{\mathsf t}$。
四つの厚い角が可逆なところで二切断条件は恒等的に成り立つ。
\end{enumerate}
いずれの型も指数増大する商を持つ。この排除には五次元主族の閉包への所属を要しない。
\end{theorem}
\begin{proof}
関係空間の最初の厚い二色の列基底は $\binom{I}{D_0},\binom{I}{D_1}$ で、
$D_0=Y_\lambda A^{-1}$, $D_1=\lambda X_\lambda^{\mathsf t}B^{-\mathsf t}$。
両者は可逆である。薄い二色の係数ベクトルは $(1,1,1,z)$ と $(1,1,1,1)$。
階数1関係において薄い係数を $\alpha,\beta$ とすれば、二つの交差ブロックの
行列式 $\alpha^2-\beta^2,z\alpha^2-\beta^2$ は零であり、$\alpha=\beta=0$。
残った対角二ブロックの一方は零でなければならず、厚い各色の両半分への射影が
同型なので関係全体が零である。巡回・反転しても交差ブロックの可逆性は保たれ、
薄い規格化を取り直せる（反転は行と列を入れ替える）。
全位置で純関係がないので、十分性は定理\ref{thm:involution-quartet-component}による。

低位二座標・高位二座標への評価像の射影を考える。評価点の片方の半分が非零なら、
その側の射影は厚い二色だけで全射である。実際、各色の係数射影は同型で、
下位ビットを入れ替える色も使える。片方だけが非零なら薄い関係が反対側の
独立方向を与える。従って階数2以下なら両側が非零で、像は二平面間の同型のグラフである。
色0だけの像をこのグラフの両側に射影すれば、両半分の非零座標数は等しい。
二支持の場合は共通のアフィン文字置換（差1を保つ）で支持を $0,2$ にできる。
全支持の場合も含め $p_0=p_2=1$ としてよい。射影化に加え、四位置の高位の
倍率を同じ数にしテンソル全体の倍率を調整すると、薄い48語を保ったまま
低位と高位の相対倍率を自由に取れるためである。

$p=(1,0,1,0)$ では薄い二行がグラフの行列を $I$ に決める。
両 $D_i$ の第一行は $(1,0)$ でなければならず、これが第一型である。
$p=(1,s,1,m)$ では薄い二行が定理の $T$ の満たす行列方程式を与える。
左の行列が特異なら右の行列も特異で、二行の比例係数が一致するから
$m=s$, $s^2=zm^2=1$ となり $z\ne1$ に反する。厚い二色の評価をこの
グラフに入れると表示した $D_i$ を得る。逆にこれらの式で全評価は階数2となる。

最後に、位置0の双対矢を $kp$、位置1を $R_0(p,-)$ の核 $L$、
位置3を $R_3(-,p)$ の核 $H$、位置2を元の四次元全体に制限して周期的な商を取る。
次の計算は八つのSchur等式（第二型）とこの商の関係階数を有理関数体上で確認する。
各辺の色行列の最初の二列（厚色）または一列（薄色）を使う。厚色の最初の二列は
交差ブロックにより独立である。$E$ を両端それぞれの薄い二行の評価行列とすると
$\binom{-E_{\rm low}^{-1}E_{\rm high}}{I}$ が核の基底であり、厚い四行もこの基底を消す。
第一型で $L=\binom{-I}{I}$, $H=\binom{\operatorname{diag}(-1,-1/\lambda)}{I}$。
第二型で消去の分母は $\lambda,1-zm^2$ の単元倍の積である。
これらの基底を同じ記号で表し、中央の辺の関係行列を $r$ とする。
$L^{\mathsf t}r$ と $rH$ を一列に並べ、それぞれ右・左に二次元の単位基底を
テンソルして得る $16\times24$ 行列の階数は、第一型の独立六変数
$c,d,f,h,z,\lambda$ 上で14、第二型の四変数 $s,m,z,\lambda$ 上で12である。
\path{computations/verify_elliptic_evaluation_obstructions.py} はこの行列を生成して厳密に掃き出す。
第一型では第二切断の条件を課す前から選んだ列についての上界が成立し、
実際の窓に制限したときそれらは関係の全基底である。従って例外値でも階数は
この上界を超えない。位置0に接する二辺の制限関係は零なので、
長さ $4n$ の商空間はこの三文字の商の $n$ 回のテンソル積に一次元因子を掛けたものとなる。
次元は少なくとも $2^n$ または $4^n$ で、標準Hilbert関数に反する。
\end{proof}

% !TeX root = ../papers/nc-p3.tex
\section{混合交換境界の全特殊値の正則性}
\label{sec:mixed-exchange-regularity}
\begin{theorem}\label{thm:mixed-exchange-regularity}
定理\ref{thm:klein-mixed-exchange-resolution}の $z=1$, $\lambda\ne0,\pm1$ の
表示で、$x=ad-fg$, $\det A$, $\det B$ を非零とする。
標準Koszul左右AS正則性の必要十分条件は
\[
                 k=(tb-\lambda)(tb-1/\lambda)\ne0
\]
である。$t=0$ の全点を含む。$k=0$ の全点は $4n$ 文字の商空間に
少なくとも $4^n$ の次元を持つ始点を有する。
\end{theorem}
\begin{proof}
規格化前の二つの式から
$(ta-1)(td-1)=(tf-1)(tg-1)=k$ を得る。
$t\ne0$ では $a,b,d,f,g$ の共通倍率で $t=1$ にできる。
$k\ne0$ のとき $(a,d,f,g)=(1+u,1+k/u,1+v,1+k/v)$ であり、
$x=(u-v)(1-k/(uv))$。$t=0$ では $a+d=f+g=(\lambda+1/\lambda)b$ である。

まず $k\ne0$ とする。偶数辺の薄い二関係は全係数1である。
最初の偶数辺で厚色の係数は前節の $D_0,D_1$ のグラフをなす。
$Y_\lambda=\lambda P X_\lambda$ と書く（$P$ は二文字の置換）。
$D_0=Y_\lambda A^{-1}$ の全成分は非零である。
実際、零なら $a,d$ のいずれかが $\lambda b,b/\lambda$ のいずれかに等しく、
$t\det A=a+d-(\lambda+1/\lambda)b$ と $\det A\ne0$ により $k=0$ となる。
階数1関係があるとする。薄い二係数の符号付き和の一方が零でなければ
交差ブロックは可逆であり不可能。薄い二係数がともに零なら厚色のグラフ性に反する。
従って二つの交差ブロックはともに $(1,\sigma)^{\mathsf t}(1,\sigma)$ の
非零倍である（$\sigma=\pm1$）。二つの対角ブロックも同じ行列の非零倍となる。
これは $D_0,D_1$ が $(1,1)^{\mathsf t}$ を同じ倍数で保つことを要する。
$A,B^{\mathsf t}$ がこの直線を保つので $a=d,f=g$、さらにその固有値の一致から
$(1+\lambda)/(a+b)=(1+\lambda)/(b+f)$、従って $x=0$ で矛盾する。

同じ偶数辺の左評価が低階数になるとする。前節と同じ射影で両半分の支持数は等しい。
二支持では、二つの薄い関係が定めるグラフは単項行列である。色0の評価をこの
グラフに入れると $D_0$ の一成分が零になり矛盾する。
全支持では $p=(1,s,1,m)$ とする。前節の薄い二行列が可逆なら、$z=1$ とした
第二型の式から $B=PA^{\mathsf t}$ が従う（$T$ が $P$ と可換であることでも確認できる）。
従って $C$ は特異となる。薄い二行列が特異なら $s=m=\pm1$ である。
厚い二色の式から $D_1=PD_0P$、ゆえに $B=A^{\mathsf t}P$ となり、やはり $x=0$。
反対の偶数辺では $A,B,\lambda$ を $A,B^{\mathsf t},1/\lambda$ に取り替えると
同じグラフになる。右評価には辺の関係行列を転置すればよく、色1のグラフは
$PD_1P$ となる。これは $B$ と $PBP=B^{\mathsf t}$ の交換である。
以上の操作は $x$, $k$ と非消滅条件を保つ。

奇数辺では、厚い二色の係数平面がそれぞれ
\[
 (u_0,u_1,\alpha u_0+\beta u_1,\alpha u_0+\beta u_1),\qquad
 (\gamma v_0+\delta v_1,\gamma v_0+\delta v_1,v_0,v_1)
\]
となり、四つの係数は全て非零である。色0の基底は $\binom{C}{J}$ または
$\binom{C^{\mathsf t}}{J}$、色1では上が $J$、下が可逆な角である。$J$ は全成分1の行列とする。
\path{computations/verify_mixed_exchange_regularity.py} は、上記の $t=0,1$ の
両座標で残る角と各係数を計算し、その分子・分母の既約因子が
$\lambda(\lambda^2-1)kx\det A\det B$ と座標の単元の因子であることを検算する。
奇数辺の薄い色の二つの交差行列式は比例しない（積比は $\lambda^{\pm2}\ne1$）。
よって純関係では薄い係数がともに零となり、一つの対角ブロックが零になる。
残るブロックには一色の一次元しか残らず、その二つの成分はともに非零で階数2である。

奇数辺の評価点で低位半分の両座標が非零なら、色0の像は低位への射影が同型の
二平面で、高位での像は高位座標対が張る直線以下である。高位座標が非零なら
色1の高位像はこの直線から出る（支持二なら全射、支持一なら別の座標）。
高位の両座標が非零の場合も役割を逆にして同じである。片半分だけの支持なら
厚い二色でその半分を張り、薄い色で別の方向を得る。
一座標ずつの支持なら薄い二色だけで両半分間の単項同型のグラフが得られ、
色0がそこに入るには $\alpha\beta=0$ が必要である。
右評価では各色内の添字が入れ替わるだけで結論は同じである。
一文字全階数は薄い非零成分と文字次数分解から従う。
定理\ref{thm:involution-quartet-component}と混合表示の特殊化定理を適用する。

$k=0$ なら $t=1$ にでき、$b=\lambda,1/\lambda$ である。
$a,d$ の第 $i$ 成分、$f,g$ の第 $j$ 成分を1に選ぶ（$i,j=0,1$）。
残りを $u,v$ とし、$b=\lambda,1/\lambda$ の順に $e=0,1$ とする。
$i=j$ なら位置 $1+e$、$i\ne j$ なら位置 $3e$ に
$p=e_i+e_{i+2}$ を置く。両隣の核平面と中央三文字商を前節と同じ方法で取る。
検算では奇数辺の色1だけは最後の二列を基底に使う。核平面の薄い二行の低位小行列式は
単項式である。八表示の全てで二平面が六関係を消し、三独立変数
$u,v,\lambda$ 上で挿入関係の階数は12である。特殊化で上がらないため
三文字商の次元は4以上となり、命題\ref{prop:cut-arrow-growth}で結論を得る。
\end{proof}

\begin{corollary}\label{cor:regular-full-thin-all-pairings}
二隣接切断の色次元が $(2,2,1,1)$ の並べ替えで四薄色の全成分が非零な四元群窓では、
候補の正則性に全パラメータで有限な必要十分判定がある。
\end{corollary}
\begin{proof}
異方向は定理\ref{thm:klein-transverse-dense}を使う。
平行方向の被覆は定理\ref{thm:full-thin-parallel-cover}である。両交換境界の
正負表示は定理\ref{thm:full-thin-double-exchange}、一交換境界は直前の定理で決まる。
残った可逆交差部の主枝には定理\ref{thm:elliptic-evaluation-obstructions}を使う。
二つの交換枝には定理\ref{thm:bisymmetric-two-length-test}を正表示に使い、
負表示は指定コサイクルを外して同じ検査をする。文字次数を保つ捻りはKoszul性、
左右の完全対合、各始点のHilbert関数を保つ。各閉包への所属も、明示された
有限型表示の稠密開部のグラフ閉包を射影的に消去して判定できる。
\end{proof}

% !TeX root = ../papers/nc-p3.tex
\section{正則表現型四元群を持つ全窓の判定と印を忘れる有限不変量}
\label{sec:klein-unmarked-atlas}
$K=(\mathbb Z/2)^2$ とし $G=\prod_{i=0}^3\operatorname{PGL}(V_i)$ とする。
ここで扱う印は $K\hookrightarrow G_{[w]}$ で、各位置で四本の固有直線に分かれ、
線形持上げで四文字が一回ずつ現れるものである。最上テンソルの総文字は各位置の
持上げの選択で零とする。非可換な射影持上げしかない印はこの型に含まれない。
\begin{theorem}\label{thm:klein-unmarked-atlas}
この型の印が一つでもある四窓全体について、標準Koszul左右AS正則性は有限判定できる。
全正則窓の添字保持同型類には、次に定める有限集合 $\mathcal D(w)$ が完全不変量を与える。
異なる色階数、疎密の異なる台、例外値、正次元の安定群を持つ場合の相互比較にも使える。
最小周期は $\mathcal D(w)=\mathcal D(\operatorname{rot}^j w)$ となる最小の $j=1,2,4$ である。
\end{theorem}
\begin{proof}[全切断の判定]
一文字全階数・隣接階数6を確認する。色別Euler式による
定理\ref{thm:klein-missing-color-coverage}の制限により、許される各切断の分割は
$0222,3111,2211$ に限られる。以下の参照は窓を受理・排除するための必要十分条件である。
\begin{center}
\begin{tabular}{c|l}
二切断における条件 & 判定箇所\\\hline
両方に欠色 & 定理\ref{thm:klein-missing-color-coverage}\\
片方に欠色 & 定理\ref{thm:klein-single-missing-windows}\\
両方 $3111$ & 定理\ref{thm:klein-three-lines}\\
$3111$ と $2211$ & 定理\ref{thm:klein-mixed-positive}で全排除\\
両方 $2211$、薄い四長方形が完備 & 系\ref{cor:regular-full-thin-all-pairings}\\
両方 $2211$、疎で異方向 & 定理\ref{thm:klein-sparse-transverse}\\
両方 $2211$、疎で平行 & 下記の148代表の判定
\end{tabular}
\end{center}
最後の二行の被覆は命題\ref{prop:klein-two-doublet-supports}である。平行方向では
定理\ref{thm:klein-parallel-flags}, \ref{thm:parallel-symmetric-surface}と
節\ref{sec:parallel-excluded-shapes}の全排除により残る代表のうち、
$82,118,124,129,134,139,142$ は定理\ref{thm:parallel-support-fibers}で決まる。
$85,88$--$90,143$--$147$ の各補完は同定理と定理\ref{thm:horizontal-rectangle-flags}
で全て決まる。以上は指定した四元群を保つ表示の被覆である。一般の入力でも
後述の印の方程式を先に解いてこの型の印が存在するかを判定できる。
\end{proof}

\begin{proof}[一つの印の文字不変量]
四つの共通固有基底を選び、総文字零の64語のうち非零の語の集合を $S$ とする。
整数行列 $E_S$ の語 $q$ の行の $(i,j)$ 成分を $[q_i=j]$ とする。
$\ker_{\mathbb Z}E_S^{\mathsf t}$ の基底 $m_1,\ldots,m_h$ を、固定した語順と整数行操作で取る。
この印の記録は
\[
 \left(S,\;\big(\prod_{q\in S}w_q^{m_{a,q}}\big)_{1\le a\le h}\right)
\]
の以下の有限な付け替えによる軌道である。
四位置で同じ $L\in\operatorname{GL}_2(\mathbb F_2)$ と、
$\sum_i d_i=0$ となる四つの平行移動を取り、$q_i\mapsto Lq_i+d_i$ とする。
これら $384$ 通りは、共通の群の生成元を選び直し、各線形持上げを文字倍して得られる。
逆に、指定した射影群を自身へ共役する四位置写像は各文字直線をこう置換する。
残る変換は各直線の倍率である。二記録の一致は二組の係数比が対角トーラスの像に入ることに等しい。
これはトーラスの指標による像の記述から従い、非飽和な格子の有理基底で代用しない。
射影全体倍率もこの像に含まれる。従って記録は部分群を指定した同型の完全不変量である。
\path{computations/verify_klein_atlas_codes.py} は付け替えを独立に列挙し、疎な台でも飽和指標による比較を行う。
\end{proof}

\begin{proof}[全印の有限性と復元]
$H=G_{[w]}$ とする。$\operatorname{Hom}(K,H)$ の一点 $\rho$ の接空間は
$Z^1(K,\operatorname{Lie}H_{\operatorname{Ad}\rho})$、
$H^\circ$ 共役軌道の接空間は同じ係数の余境界 $B^1$ である。
実際各群元の一次摂動に積の等式を課せばコサイクル式そのものを得る。
標数零の有限群ではこの二空間が等しい。コサイクル $c$ に対し
$v=|K|^{-1}\sum_{g\in K}c(g)$ とすると $c(g)=v-gv$ である。
軌道は滑らかで局所閉、周囲の接次元がその軌道次元に等しいので周囲もその点で滑らかで
軌道は開である。全軌道が開なのでそれぞれ閉でもあり、有限型性で有限個しかない。
各既約成分は一軌道である。この議論は $H$ の簡約性や有限性を要しない。

具体的には $H$ を $g w\wedge w=0$ と各行列の可逆性で与え、
$H$ の可換な二つの射影対合の方程式を取る（恒等元を含む $g^2=1$ でよい）。
アフィン標での既約分解後、各成分から幾何点を一つ取る。
各位置で二つの行列の二乗を1に規格化し、通常の行列として可換で、
同時固有空間が各一次元の成分のみを残す。この条件は共役不変である。
射影基底の全選択による重複は前段の有限軌道に吸収される。
残った印についての記録の集合を $\mathcal D(w)$ とする。
別の実装では、台と文字不変量の値を座標とし、固有基底による変換の存在を
消去してもよい。その値の像は今示した有限集合であり、固有基底自体の有限性を要しない。

添字保持同型は標準テンソルの四位置射影同値である
（系\ref{cor:two-window-extension}の最後の構成）。
従って同型なら全印の記録が一致する。逆に二集合に一つでも共通の記録があれば、
その二つの固有基底を使って添字を保つ同型を構成できる。
窓には四移動との標準同型があるため周期の判定も従う。
\end{proof}

% !TeX root = ../papers/nc-p3.tex
\section{四元群窓の閉軌道への全退化とトーリック商}
\label{sec:klein-polystable-limits}
$W$ を総文字零の64語の線形空間、$T$ を四位置の対角トーラスの射影像、
$G=\prod\operatorname{PGL}(V_i)$ とする。射影GITには
$\prod\operatorname{SL}(V_i)$ の自然な線形化を用いる。以下の閉軌道は半安定集合内で取る。
\begin{theorem}\label{thm:klein-polystable-limits}
正則表現型四元群を持つ全標準AS窓は半安定であり、その軌道閉包にある唯一の閉軌道も
標準Koszul左右AS正則である。従って同じ印を持つ二隣接階数6の候補が正則であることと、
半安定で閉軌道の窓が正則であることは同値である。閉軌道は次の有限操作で求まる。
\begin{enumerate}
\item 両切断に欠色があれば、四文字が相異なる語は全て残す。
反復対 $ijij,jiji$ は、その補対 $klkl,lklk$ がない場合だけ二つとも消す。
\item 疎な平行二重色の正則台で、番号
$51,54,55,56,57,58,62,63,64,65,76,78,112,115,117$ は指定始点を保つ
文字置換を含めて台50に移り、$59,66$ は台52に移る。
残す語の係数はそのまま使う。
\item 他の全正則表示は既に閉軌道である。
\end{enumerate}
さらに四元群印を仮定しない二隣接階数6の任意の窓も、ここで得た正則な閉軌道のいずれかを
極限に持てば標準Koszul左右AS正則である。
\end{theorem}
\begin{proof}[台の多面体と四位置での閉軌道]
語 $q$ の重みを前節の $E$ の行とし、各位置の四座標が全て $1/4$ の点を $o$ と書く。
台 $S$ の重みの凸包を $P$ とする。$o\notin P$ なら対角一径数群が零へ送る。
$o\in P$ なら $o$ を相対内部に含む最小面 $F$ があり、外の語を消す
四つの跡零整数重みを一つ取れる。残った台には、全語で正かつ各位置の周辺和が全て等しい
有理重みがある。このとき対角トーラスでの軌道は多安定である。
複素数上では実対角倍率によるノルム最小化で全ての一位置周辺ノルムが等しくなる。
実際、対数変数での二乗ノルムは重みの指数関数の正係数和で、台が張る方向では
$o$ が相対内部にあるので無限遠で発散する。
各位置の異なる二文字の切片の台は総文字零により交わらず、Gram行列の非対角成分は元から零である。
従って全 $G$ に対しても最小ノルム条件を満たし、閉軌道を得る。
有限個の係数の体を複素数へ埋め込んでから代数閉拡大に戻せる。

$F$ が真の面なら元の窓はこの閉軌道と同型にならない。実際、一つの $G$ 軌道内で
$K$ が固定する射影点の部分の接空間は接表現の不変部分であり、
$\operatorname{Lie}G$ から軌道接空間への全射を $K$ 不変に取っても全射である。
よって各固定点で中心化群の恒等成分の軌道が開である（固定部分は滑らかで、
または接次元と軌道次元が等しいことを使える）。その恒等成分は $T$ なので各軌道は
固定部分で閉でもある。相異なる台の対角極限が元の $G$ 軌道に入ることはない。
\end{proof}

\begin{proof}[正則な台の全場合]
両欠色では節\ref{sec:unmarked-klein-pbw}の二反転により、全相異なる語は同じ $(b,e)$ の
四項が同時に現れる。この四項は各位置で全文字を一度ずつ使う。
補い合う二反復対の四項も同じ性質を持つ。残る反復対の片割れを消すには、
残った辺の二ラベルに四位置で重み $+1$、その補辺に $-1$ を与える。
全相異なる語、他色の反復対の重みは零であり、残った辺の二項だけが正である。
各欠損色でこれを加える。全相異なる語が全くない場合は零行の不在から全反復対がある。
従って残る台は空でなく、各残存語を正の重みで覆えるのでちょうど最小面である。
一回の片割れ消去は節\ref{sec:unmarked-klein-pbw}の良還元の証明にある階数2の
主元と非零行・列を保つ。その段階で新しく全部欠けた反復色があれば同節のPBW判定が
適用でき、残る片割れを順に消しても同じ主元の議論が使える。

一方だけに欠色の正則台、異方向の疎な二正則台では各台に正の重みを直接与えられる。
$3111$ 型の全非定数60語は各位置の辺際が等しく、その重み行列の階数13は全64語の階数に等しい。
従って定数四語の任意の追加に対し $o$ は相対内部にとどまる。
異方向の完備な薄色48語でも同じ階数13と周辺和の一致が成り立つ。
平行方向の48語は同じ周辺和を持ち、階数12であり、そのアフィン包は高位文字数2の超平面である。
可逆交差部と混合表示には全低位・全高位の語が双方存在する。
二重交換表示でも節\ref{sec:klein-double-exchange-fourier}の第一型には $X\ne Y,W\ne V$、
第二型には $k/g_i,k/r_i\ne0$ があり、両側に語がある。
よって超平面の両側に重みが追加され $o$ は再び相対内部となる。

平行疎型の全正則代表については
\path{computations/verify_klein_polystable_limits.py} が各表示の台を検査する。
定理の17台には消す語で正、残す語で零の整数重みを与え、残存台が実際に50,52の
始点固定の付け替えに属することを確かめる。残りの正則台、50,52自身、上記の片欠色と
異方向の疎台の全てで、正有理ベクトル $c$ により $E_S^{\mathsf t}c=\boldsymbol1$ を検算する。
二つの終着台の各厚色・各切断には非零単項式の二次小行列式があり、全薄色は非零である。
階数上界は特殊化で保存されるので、係数の特殊値を含め定理\ref{thm:klein-parallel-flags}で正則である。
最後に極限が正則なら、元の軌道からその極限を通る曲線を取って
命題\ref{prop:unrestricted-good-reduction}を適用できる。これは元の点の対称性を使わない。
\end{proof}

\begin{theorem}\label{thm:klein-toric-normalization}
$\mathbb P(W)//T$ は次元51の正規射影トーリック多様体で、その前節の384元群による商は
$\mathbb P(W)$ の全基底変更GIT商における像の正規化である。
特に任意の印付き族の多安定な特殊化の像はこの有限写像の下でその閉包を取って得られる。
\end{theorem}
\begin{proof}
四位置の跡零対角群を $\widetilde T$ とする。
$B=k[W]^{\widetilde T}$ は、各位置の四ラベルで周辺次数が等しい非負整数指数の
単項式環である。直交格子と正錐の交わりは飽和しており正規である。
全重みの階数13、$o$ は内部なので $\operatorname{Proj} B$ は次元 $64-13=51$ を持つ。
全 $\prod\operatorname{SL}(V_i)$ の斉次不変式の制限の像を $A\subset B$ とする。
上の最小面の議論は階数条件を使わず、対角群と全基底群の零錐が $W$ 内で一致することを示す。
従って $A_+$ の $B$ での零点は斉次頂点だけである
（アフィン対角商 $W\to\operatorname{Spec}B$ の全射性を使う）。
$B/A_+B$ は有限次元であり、斉次次数の帰納により $B$ は $A$ 上有限である。
これで射影商の像への有限全射が従う。射影群による付け替えは全基底変更商上で自明なので
384元群の商を経由する。

系\ref{cor:klein-four-configurations}の標本は今示した多安定な点で、射影安定群が正確に $K$ である。
有限安定群なので安定点であり、Lunaのエタール切片によりその飽和近傍で各安定群は
この有限群の部分群に共役である。従って $\mathbb P(W)$ の一般点では指定 $K$ が
既に全安定群である。安定点の二つの表示の同値は指定 $K$ を正規化する四位置変換だけから
生じ、前節より対角群と384置換で尽くされる。有限写像はこの商の後で一般に一対一だから
標数零で双有理である。正規環の有限群不変環の正規性で結論が従う。
有限射は閉写像なので、任意の部分集合の像の閉包についての主張も従う。
\end{proof}

% !TeX root = ../papers/nc-p3.tex
\section{任意のトーラス作用に対する四文字のEuler行列式}
\label{sec:torus-euler-characters}
四元群印・点配置・領域性を仮定しない。標準Koszul左右AS四窓を保つトーラスの
作用を各 $V_i$ に線形に持ち上げる。$A_i$ を四矢、$B_i$ を六関係の指標とし、
$h$ を最上四テンソルの一次元指標とする。以下、$h$ は指標環の単項式、
$A^*$ は指数を逆にした指標とする。始点を法4で並べる。
\begin{theorem}\label{thm:torus-euler-characters}
Koszul分解の行列 $Q(x)$ の第 $i$ 行には、巡回順に
\[
 1+x^4h,\qquad -x A_i,\qquad x^2 B_i,\qquad -x^3 h A_{i+3}^*
\]
が入る。四つのトーラス文字 $m_1,\dots,m_4$ が重複を込めて一意に存在し、
\[
 \det Q(x)=\prod_{j=1}^4(1-x^4m_j),\qquad
 \prod_j m_j=h^4,\qquad
 \prod_{i=0}^3\det(V_i)=h^4 .
\]
最後の等式は射影安定群の線形逆像の恒等成分上での指標の等式でもある。
とくに四位置で行列式1の連続安定群は最上テンソルそのものを固定する。
\end{theorem}
\begin{proof}
まず一つの乗法群で文字変数を $t$ とする。$Q^{-1}$ は始点・終点の組ごとの
経路の指標Hilbert行列である。$|t|=1$ なら各係数の絶対値は
$\binom{n+3}{3}$ 以下であり、全成分が $|x|<1$ で正則である。
従って $Q(x)$ の行列式はこの円板で消えない。置換展開の各項の経路長の和は4の倍数なので
$y=x^4$ とおけば行列式 $D(y,t)$ は $y$ の四次式で、定数項1、最高次係数 $h^4$ である。
$|t|=1$ では両端係数の絶対値は1なので四つの逆根 $\lambda_j(t)$ は全て単位円にある。
$t=1$ では $Q=(I-xS)^4$（$S$ は巡回置換）より $D(y,1)=(1-y)^4$。

ここで四次多項式の逆根が実際に単項式になることを示す。
べき和 $p_n(t)=\sum_j\lambda_j(t)^n$ はNewton恒等式で整数Laurent多項式であり、
単位円上の絶対値が4以下なので各Fourier係数の絶対値も4以下である。
素数 $\ell=11$ を使う。Frobenius合同
$p_{n\ell}(t)\equiv p_n(t)^\ell\equiv p_n(t^\ell)\bmod\ell$ と係数の上界から、
$p_{n\ell}(t)=p_n(t^\ell)$ が全 $n$ で成立する。従って $t^\ell$ の逆根の多重集合は
$t$ の逆根の $\ell$ 乗の多重集合である。$t=\exp u$ とおき、$u=0$ でPuiseux級数の
枝を全て取る。全逆根の定数項は1なので形式対数 $g_j(u)$ が取れる。
多重集合 $\{g_j(\ell u)\}$ と $\{\ell g_j(u)\}$ の一致を与える置換を一つ固定し、
その位数 $N$ 回反復すると各枝で $g_j(\ell^N u)=\ell^N g_j(u)$ となる
（分数冪には正の実数の同じ根を使う）。係数比較により $g_j(u)=a_j u$。
$p_1(e^u)$ は有限個の $e^{nu}$, $n\in\mathbb Z$ の線形結合であり、
異なる指数の指数関数は原点の導関数のVandermonde行列により一次独立である。
従って全 $a_j$ は整数で、$p_1(t)$ の係数は四文字の多重度である。
トーラスが高次元なら $p_1$ の台およびその四重までの和と $D$ の全係数の台を
分離する一般の整数一径数を選ぶ。この方向に制限した係数比較は元の指標環にも持ち上がる。

最後に、$p_1=-[y]D$ を展開すると
\[
 p_1=\prod_i A_i-\sum_i B_i A_{i+2}A_{i+3}
       +h\sum_i A_i A_i^*+B_0B_2+B_1B_3-4h .
\]
任意の一径数で $h=t^H$ とし、$M=\sum_i A_i'(1)$ と置く。
$B_{i+2}=hB_i^*$ より $\sum_i B_i'(1)=12H$。
上式を微分すると $p_1'(1)=16M-60H$ である。一方、四文字の積は最高次係数
$h^4$ だから $p_1'(1)=4H$。従って $M=4H$ で、これが所望の行列式指標である。
連結線形群の文字は極大トーラスで決まることから最後の主張を得る。
\end{proof}

\begin{corollary}\label{cor:two-weight-profiles}
各位置の非自明な重みが二つまでで、スカラーを引いて重み $0,1$ となる一径数作用を考える。
スカラーだけの位置の重みは0とし、少なくとも一位置が非自明とする。
$k_i$ を重み1の矢次元、$H$ を最上の重みとする。可能な型は巡回を含め下表に限られる。
$B_i$ は重み $0,1,2$ の多重度順で、$B_{i+2}(t)=t^HB_i(t^{-1})$ である。
\[
\begin{array}{c|c|cc|c}
 (k_i)&H&B_0&B_1&\{\deg m_j\}\\ \hline
 1111&1&330&330&0004\\
 2222&2&141&141&0044\\
 3333&3&033&033&0444\\
 0121&1&420&240&0022\\
 0323&2&240&042&1133\\
 1133&2&321&060&2222\\
 1133&2&420&141&2222
\end{array}
\]
これは必要条件である。特に各位置が $2+2$ なら各辺に低位二文字の関係と
高位二文字の関係がちょうど一つずつ存在し、混合関係はちょうど四つである。
この場合の二重Hilbert級数は全始点で $(1-u)^{-2}(1-v)^{-2}$ である。
\end{corollary}
\begin{proof}
一文字全階数と重み0の存在から $1\le H\le3$。各辺の $B_i$ の重み別次元は
その辺の重み空間と、対辺の相補重み空間の両次元以下である。$p_1$ の上式から
四指数は全て0から4の間にあり、その和は $4H$。
\path{computations/verify_torus_euler_profiles.py} はこの範囲の整数多重度を尽くす。
行・列の経路長の差を取り、行列式には継ぎ目をまたぐ成分だけ周期変数 $y$ を付ければよい。
整数点 $(t,y)=(11,7)$ で全因子候補と比較して不一致を落とし、残った各組では
多項式の両変数の次数上界を超える格子で係数の一致も確かめる。
表のほかには $0103,H=1,B_0=B_1=510$ のみである。この型では位置1の高位文字を $v$ として、
両隣が重み0のみなので $a\otimes v$ と $v\otimes b$ が続けて関係に入り、
定理\ref{thm:no-consecutive-pure-relations}に反する。
$2222$ のHilbert級数は位置によらない上記の矢・関係の多重度をKoszul逆公式に入れると得られる。
\end{proof}

% !TeX root = ../papers/nc-p3.tex
\section{二平面分解の可逆部における全正則性と相補逆行列}
\label{sec:split-inverse-windows}
四位置で $V_i=X_i\oplus Y_i$、両者二次元、重みを $0,1$ とする。
全体重み2、各辺の色階数 $(1,4,1)$ とし、一文字全階数を仮定する。
各辺の純次数関係を $r_i\in X_i\otimes X_{i+1}$,
$s_i\in Y_i\otimes Y_{i+1}$ とし、混合次数の関係四平面を $M_i$ とする。
\begin{theorem}\label{thm:split-inverse-regularity}
全ての $r_i,s_i$ が行列として可逆であることと、
全位置で $M_i\to Y_i\otimes X_{i+1}$ が同型であることは同値である。
この場合 $M_i\to X_i\otimes Y_{i+1}$ も全て同型で、
これらの窓は係数への追加条件なしに全て標準Koszul左右AS正則である。
\end{theorem}
\begin{proof}
全体の二平面別の六成分のうち $XXYY, XYYX, YYXX,YXXY$ は、
色0と2の切断階数が1なので、各々対応する $r_i\otimes s_{i+2}$ の位置置換・非零倍である。
混合対混合の $8\times8$ 切断の二交差ブロックはこれらの二次行列二つの
Kronecker積の置換である。各因子が可逆なら両交差ブロックが可逆で、切断階数が4なので
列空間・行空間の両半分への射影は全て同型である。

逆に一方の射影だけを全位置で仮定する。二次双対では双対文字の $XY$ を $YX$ に
常に書き直せる（元の $M_i$ を $YX$ 上のグラフとして直交補を取る）。
双対の任意の語は従って全 $Y$ の後に全 $X$ が来る語で張られる。
もし $r_i$ が特異なら、その右核に非零の $x\in X_{i+1}^*$ があり、$X_i^*x=0$ である。
位置 $i+1$ の $x$ に終わる四文字で最上形式が評価できるのは $Y$ が二つのときだけである。
その前の三文字を $YYX$ に並べ替えると全て零で、最上テンソルの一文字全階数に反する。
この評価は二次双対の関係を消すので、まだKoszul性を要しない。$s_i$ が特異な場合は
左核の $y$ から始めて後続三文字を $YXX$ に並べれば同じ矛盾を得る。

元の代数で $YX$ は $XY$ に書き直せる。二本の非退化二文字関係列はそれぞれ、
整数位置に沿って基底を順に選ぶと通常の交代関係になる。高位文字を低位文字より大きくした
位置別左辞書順を使えば混合四式と両半分の各一式はすべて降順対を先頭語に持つ。
従って長さ3での正規候補は20語である。最上テンソルの一文字縮約4本は独立な三次重なりを与え、
実際の次元も20以上である。全三文字の重なりが解消し、Diamond補題で全次数の二次PBW基底を得る。
四位置の有限Frobenius完全対合は、Koszul双対次元が $1,4,6,4,1$ となった後に
元の切断階数と最上形式から従う。
\end{proof}

\begin{theorem}\label{thm:split-inverse-transporter}
この可逆部で、順序付き二平面分解を保つ添字同型類は次のデータで全て与えられる。
四つの二次元空間 $U_i$ 上の射影テンソル $P$ で、$01|23$ と $12|30$ の両切断が可逆なものを取る。
それぞれの切断で逆行列の転置を四つの双対空間に戻した射影テンソルを $I_0(P),I_1(P)$ とする。
$D=\prod\operatorname{PGL}(U_i)$ とし、二次元空間の射影線とその双対は交代形式で同一視する。
$\tau=(02)(13)$ と書き、$U_i$ を固定した二次元模型に移す。
\[
  H=\operatorname{Stab}_D(P),\qquad
  \mathcal B(P)=\{v\in \tau D:\ v I_1(P)=I_0(P)\}.
\]
$P$ を $D$ 同値まで取り、$\mathcal B(P)$ が空なら捨てる。各 $P$ の上の
同型類は $\mathcal B(P)$ の $H$ 共役類である。非空なら $\mathcal B(P)=Hv$、
$vHv^{-1}=H$。これは非半単純な $P$ にも成立する。
\end{theorem}
\begin{proof}
$U=(X_0,Y_1,X_2,Y_3)$ とし、各位置の相補空間を $Z$ と書く。
交代二成分を $P\in\bigotimes U_i$, $Q\in\bigotimes Z_i$ とする。
第一切断の交差成分は非退化な二文字形式の積であり、
$\pi_0=(03)(12)$ の四つの対応で射影同型 $L_i^0:U_i^*\to Z_{\pi_0(i)}$ を与える。
第二切断では $\pi_1=(01)(23)$ の四対応 $L_i^1$ となる。
各切断で $P$ のブロックを左上とした行・列順にする。交差ブロックが可逆で全階数が4なので
$P,Q$ の当該ブロックも可逆であり、Schur式 $Q=V P^{-1}U$ が成り立つ（右辺の $U,V$ は二交差ブロック）。
添字を戻すと $[Q]=\pi_j(\bigotimes L_i^j)[I_j(P)]$ である。
\path{computations/verify_split_inverse_windows.py} は両方の添字順を直接検算する。

逆に、二つの射影等式を満たすデータから、六成分の代表倍率を決められる。
$P,Q$ の代表を固定し、各切断で二交差成分の一方を倍率変更してSchur式を厳密に合わせる。
二切断で使う交差成分は別々なので両立する。全解は一つの添字同型類を与える。
実際、高位位置が $i,j$ の六成分の代表の比を $c_{ij}$ とすれば
$c_{01}c_{23}=c_{02}c_{13}=c_{03}c_{12}$ である。これらは $c_{ij}=d_i d_j$ であるための
必要十分条件であり、例えば $d_0^2=c_{01}c_{02}/c_{12}$ から順に取れる。
四位置の高位を $d_i$ 倍すればよい。再構成した窓は純次数の全行列が可逆であり、
各文字はそれらの積で検出されるので一文字全階数も持つ。

第一の四つの $L_i^0$ を相補空間の座標に使うと、残る一組は
$\pi_0^{-1}\pi_1=\tau$ を伴う $v$ となる。
$P$ を変更する残りの $D$ は $v$ に共役で作用する。逆転置写像は射影双対上で
$D$ 同変かつ各切断で対合なので $I_0(P),I_1(P),P$ の安定群は同じ $H$ である。
従って求める移送集合と共役商が得られる。ここで $\tau$ は相補文字の同定から生じ、
窓の始点を移動する操作ではない。
\end{proof}

\begin{corollary}\label{cor:split-finite-unmarking}
一つの窓におけるこの分解の選択は、射影安定群の恒等成分での共役まで有限個である。
$H$ が有限なら順序付き分解を忘れる追加操作は全位置の $X,Y$ の同時交換だけである。
\end{corollary}
\begin{proof}
各分解は重み差1の一径数で決まり、全一径数は安定群の固定した極大トーラスに
共役で入る。そのトーラスの各位置の矢の射影重み差は指標格子の有理空間を張る。
これら有限個の重み差の値が $0,\pm1$ に限られるので、該当する余指標は有限個である。
$H$ が有限なら指定一径数を中心化する群の恒等成分はその一径数だけである。
実際、八つの二次元空間の射影作用は $P$ と四同定から $H$ に入り、残るスカラーで
テンソル直線を保つには六つの高位対の倍率 $d_i d_j$ が全て等しく、全 $d_i$ が等しい。
よって安定群の階数は1で、原始余指標は符号まで一つである。
\end{proof}

% !TeX root = ../papers/nc-p3.tex
\section{相補逆行列の半単純スペクトルと四元群印を持たない孤立型}
\label{sec:split-bell-locus}
前節の $P$ の第一切断を四次元空間の間の可逆行列 $M$ と見る。
各側は二つの二次元空間の積なので行列式二次形式を持つ。その形式による随伴を $M^\dagger$ とする。
ここでいう半単純とは $M^\dagger M$ が対角化可能であることを指す。
\begin{theorem}\label{thm:split-bell-locus}
半単純な $P$ は以下の射影四数 $x=(a,b,c,d)$ で、座標置換と偶数個の符号反転を除いて分類される。
\[
  f_0=abcd,\qquad
  -16 f_1=(a-b-c-d)(a-b+c+d)(a+b-c+d)(a+b+c-d).
\]
$f_0f_1\ne0$ を課す。前節の移送集合が非空となる $P$ は、以下とその付け替えで尽くされる。
\begin{enumerate}
\item 曲面 $ab+cd=0$。
\item 直線 $a=b,c=d$。
\item $a=1$, $(d^2+1)^3=4d^3$,
$2c=d^3+d^2+3d-1$, $2b=d^5-d^4+3d^3-7d^2+6d-2$。
この六点は付け替えの下でちょうど二つの軌道を与える。
\end{enumerate}
曲面の一般点の上には分解付きで四つの添字同型類がある。第三型では二軌道の各々に
ちょうど一つあり、分解を忘れると両者は同型である。この一つの標準AS窓は最小周期1で、
四位置の射影安定群は $\mathbb G_m$ 自身である。従って四元群印を持つ窓のどれとも同型でない。
孤立は本定理の $P$ の分類についての主張である。
\end{theorem}
\begin{proof}[標準形と九十六通りの比較]
$\operatorname{SL}_2\times\operatorname{SL}_2$ の四次元像は行列式形式の特殊直交群である
（核は同時の $\pm I$、像は同次元の連結閉部分群）。
$M^\dagger M$ の異なる固有空間は互いに直交し、各空間は非退化である。
各々の正規直交基底を右側に取り、その $M$ 像を固有値の平方根で割ると左側の
正規直交基底を得る。向きを揃えると、対角の四係数の自由度は順序と偶符号変更および共通倍率だけである。
これは固有値が重なっても同じで、二乗の多重集合と四係数の積が残る向きを決める。
通常の二進基底に戻せば非零になり得る四対は
$0000,0011,0101,0110$ とその全ビット反転で、各対の値は
\[
             (a+b)/2,\ (a-b)/2,\ (c+d)/2,\ (c-d)/2
\]
である。この規約は二側のBell基底 $00\pm11,01\pm10$ を使うもので、
二次形式の符号に合わせた単位因子を両側で逆に取れる。
四成分の一斉符号は射影的に消えるので、有限群 $W$ の位数は96である。

両切断の行列式は $f_0,f_1$ であり、逆転置の射影Bell四数はそれぞれ
$\nabla f_0,\nabla f_1$ である。これは二つの二行列の逆転置を直接取った式で、
二次元双対を同一視する交代形式四つもこの表示を保つ。
二重互換 $\tau$ は四対の値を変えない。従って比較は $g=(\pi,e)\in W$ に対する三条件
\[
  x_i e_i\,\partial_{\pi(i)}f_1=x_0 e_0\,\partial_{\pi(0)}f_1
                    \qquad(i=1,2,3)
\]
に等しい。各射影勾配は $W$ 同変である。実際、四文字の局所基底変換に対し
両切断の逆転置はともに双対に移り、交代形式による射影同定は同変である。
従って共役な $g$ の解は $W$ で移る。

$a=1$ として十共役類の各代表で三式の極小素成分を求め、$f_0f_1=0$ に入る成分を除く。
\path{computations/verify_split_bell_locus.py} は四つの非零座標と四一次因子を個別に調べる。
代表の順に共役類の大きさは $3,1,12,12,32,6,3,3,12,12$、残る有理素成分数は
$0,6,0,2,2,1,0,3,0,3$ である。各素成分で三曲面の積の式または六直線の一つの二式が
消えることを確認し、残る一成分では消去イデアルが第三型の三式と一致する。
逆向きには恒等変換で第二型、$(01)(23)$ で第一型の三式が恒等的に消える。
第五代表の六次式は既約で、三曲面式の積の剰余は $-4$ である。
六点では $\sum x_i^2\ne0$ かつ平方は全て相異なるので、$[x]$ を保つ $W$ は自明である。
第一逆写像は対合だからその像の $W$ 安定群も自明で、六点それぞれの比較 $g$ は一意である。
第五代表の32共役を動かして得る $32\cdot6$ 点は従って相異なり、二つの96点軌道となる。
\end{proof}

\begin{proof}[曲面の一般ファイバーと孤立窓]
平方が相異なり、$\sum x_i^2\ne0$ なら $D$ の安定群 $H$ はちょうど四元群である。
左右の直交行列は共通の四固有線上の偶符号変更となり、全符号を除くと四つ残る。
元の各二次元因子の射影線では、全位置同時のビット反転と全位置同時の
$\operatorname{diag}(1,-1)$ が生成する。$g$ のベル四線の置換が $H$ の三つの非自明元に
誘導する作用は、四線の三つの二対分割への置換である。各偶符号はその正負で四線を二分するためである。
$\tau$ は上の二生成元をそれぞれ保つ。従って曲面の一般点で $g=(01)(23)$ は
$H$ に自明に作用し、前節の $Hv$ の共役類は四個ある。
六次式では第五代表の置換が三サイクルで $H$ の非自明三元を巡回する。
このとき $1-\operatorname{Ad}(v)$ は $H$ 上全単射なので $Hv$ は単一共役類である。

分解を交換すると第一交代成分は旧 $Q$ になり、最初の四対応で座標を揃えると
そのBell四数は $[1/a:1/b:1/c:1/d]$ である。上記検算は六次式を法として
この四数を元に戻す $W$ 元が一つもないことを確認する。交換は対合であり、
曲面・直線の付け替えの和集合は逆数で保たれるので二つの孤立軌道を交換する。
一位置移動の新しい交代成分は $Q$ の一巡回である。同じ六次剰余体で、
そのBell四数を旧 $P$ に戻す $W$ 元の存在も検算する。移送の共役類は一つなので
分解付きでも最小周期1になる。

最後に孤立窓の多安定性を確かめる。複素数上で、各BellベクトルをEuclidノルム1にし、
行列式形式に対する符号が負のものに虚数単位を掛けると、対称形式とHermite形式の
双方について正規直交となる。行列式1の実符号置換はこの基底の実特殊直交群に属し、
二つの二次元因子でユニタリ行列に持ち上がる。実際この実群は
$\operatorname{SU}(2)\times\operatorname{SU}(2)$ の像である
（コンパクトな連結六次元群で、二つの形式を保つ）。右左で必要な置換・偶符号も
この持上げを用いる。第一の四対応と第二の四対応を射影ユニタリに選べるので、
四つの境界成分は最大階数の二文字ユニタリ形式二つの積である。
Bell型の $P,Q$ でも、各一文字Gram行列は偶奇の台と全反転によりスカラーである。
六つの二重次数成分は各一文字の同じ半分で異なる相手台を持つ。
各高位半分の実倍率を動かすと、六つの二要素集合の重みの凸包の相対内部に
均一周辺和の点があるので、全位置で両半分のGram値を一致させられる
（節\ref{sec:klein-polystable-limits}のノルム最小化）。
よって全基底群で多安定で、射影安定群は簡約群となる。

指定トーラス $T$ の中心化群を $T$ で割ったものは前節の構成で $H$ のうち
$v$ と可換な群に入り、ここでは自明である。$T$ の反転は分解の交換に当たり、
二軌道が印付きで相異なるので安定群内で起こらない。従って $N(T)=T$。
$T$ は極大トーラスで、連結簡約群の非自明な根は非自明なWeyl反射を与えるから
安定群の恒等成分は $T$ である。どの連結成分にも $T$ の正規化元があるため全群も $T$ となる。
\end{proof}

% !TeX root = ../papers/nc-p3.tex
\section{半単純なBell直線の全補完・無印同型・最小周期}
\label{sec:split-spectral-line-classification}
\begin{theorem}\label{thm:split-spectral-line-classification}
定理\ref{thm:split-bell-locus}の直線の全点で、交代成分が半単純な可逆二平面分解窓を
添字保持同型で割ると、全ての量子多項式類と以下の四行で尽きる。全て標準Koszul左右AS正則で最小周期1である。
\begin{enumerate}
\item $s\in k^\times$, $s^2\ne1$ の一族、$s\sim-s\sim s^{-1}$。
\item $s^2=-1$ で追加される $a\in k^\times$ の一族、$a\sim-a\sim a^{-1}$。
\item $s^2=1$ で一方が非半単純Jordan、他方が半単純な一族。後者の固有値比
$\lambda\in k^\times$ を $\lambda\sim\lambda^{-1}$ で割り、$\lambda=1$ はスカラーとする。
\item $s^2=1$ で両方とも非半単純Jordanの一点。
\end{enumerate}
この四行は相互に交わらない。量子部分の全同型は定理\ref{thm:quantum-indexed-all}による。
\end{theorem}
\begin{proof}[印付きの全移送]
$P$ のBell四数は $(1,1,s,s)$ に取れ、付け替えによる $s$ の自由度は
$\pm s^{\pm1}$ である。通常の座標では $P_{abab}=K_{ab}$ が全ての非零成分で、
$K=\left(\begin{smallmatrix}1&s\\s&1\end{smallmatrix}\right)$ である。
両隣接切断の逆転置はどちらも同じ表示の $P(1/s)$ であり、
$\pi_0=(03)(12),\pi_1=(01)(23)$ の両方がこれを保つ。
第一の四対応を恒等に取れば、第二の四対応 $L_i$ の全体は
$P(1/s)$ の局所射影安定群 $H$ の元である。
$P(s)$ の安定群も $H$ と同じで、$b=(b_i)\in H$ による同型は
\[
                   L_i\longmapsto b_{i+2}^{-t}L_i b_i^t
\]
である。これは補空間側の基底変更を第一の四対応で消した式である。

$D(a)=\operatorname{diag}(1,a)$, $X=\left(\begin{smallmatrix}0&1\\1&0\end{smallmatrix}\right)$,
$Z=D(-1)$, $T(a,b)=(D(a),D(b),D(a^{-1}),D(b^{-1}))$ とする。
$s^2\ne1$ なら $02\mid13$ 切断の両支持は $00,11$ の張る二平面であり、
その中の可分解直線はこの二本だけである。ゆえに局所作用は、位置0,2で同じビット置換、
位置1,3でも同じビット置換を持つ単項行列である。係数表 $K$ の行列行・列倍率に関する不変量は
交差比 $s^{-2}$ である。片側だけの反転が可能なのは $s^4=1$ のときだけである。
$s^4\ne1$ なら $H$ は全 $T(a,b)$ と一斉反転 $F=(X,X,X,X)$ で生成される。
従って $L=T(a,b)$ は全て $L=1$ に同型、$L=T(a,b)F$ は
$(a,b)\sim(a^{-1},b^{-1})$ だけで割る。実際、前者には対角基底の二乗が作用し、
後者で対角基底は自明に作用する。

$s^2=-1$ では二つの独立な反転の持上げを
\[
           A=(XZ,Z,X,1),\qquad B=(Z,XZ,1,X)
\]
に取る。これらは局所射影群で可換な対合であり、$AB=F$ である。
全 $L$ は $T(a,b)A^e B^f$ ($e,f\in\{0,1\}$) となる。
対角基底で第一座標は $e=0$ のとき、第二は $f=0$ のとき自由に消せる。
残る座標で $A,B$ による作用は、それぞれ
$(a,b)\mapsto(-a^{-1},-b)$ と $(a,b)\mapsto(-a,-b^{-1})$ である。
ここでは消した座標を無視し、$e=f=0$ なら一点となる。
従って片側だけの反転は各側で一つの $a\sim-a\sim a^{-1}$ の族となる。
両側反転 $e=f=1$ では表示した二つの作用で二座標を割る。

$s=1$ に残る場合は $K$ の階数が1で、$P$ は $02,13$ の二つの非退化形式の積である。
このとき $H$ は $(C,D,C^{-t},D^{-t})$ の全体で $\mathrm{PGL}_2^2$ と同型である。
$L$ の最初の二行列 $C,D$ を各々独立な射影合同で割ればよい。
$J=\left(\begin{smallmatrix}0&-1\\1&0\end{smallmatrix}\right)$ とすると
これは $CJ,DJ$ の独立な射影共役に等しい。
\end{proof}
\begin{proof}[量子窓の認識と印の除去]
両側反転 $L=T(a,b)F$ なら各位置2,3の全二文字でビットを反転すると、
四テンソルの24個の非零語はすべて四文字の置換となる。
逆Schur式の倍率は全て非零である。各辺の切断は、互いに補となる無順序二文字組について
六つの二行二列ブロックに分かれ、全階数6、各係数非零だから各ブロックが階数1である。
よって全関係はスカラー付きの二文字交換である。三次PBW次元から相異なる三文字を
逆順から並べる両経路の係数が一致し、各三角積は位置に依存しない。位置別の対角倍率で
交換係数を一定にできる。具体的に $L=T(a,b)F$ での四三角積は、文字を二色の順に並べて
$(-as,-a/s,-1/(bs),-s/b)$ である。\path{computations/verify_split_spectral_line.py} は両切断の再構成から
全位置でこの式を直接求め、例外点の基底変換と一移動も検算する。$s=1$ で $CJ,DJ$ がともに半単純な場合も、各々を対角化すれば
$C,D$ が反対角になり同じ議論が通る。逆に任意の量子窓で四座標を二組に分けると
射影座標二直線が指定できる。その交代四テンソルは位置別のビット変更の後 $abab$ にのみ支持され、
四係数は全て非零である。行・列倍率で対角が1、非対角が共通の $s\ne0$ にできるので
直線スペクトルに入る。

分解を与える一径数群を含む極大トーラスの次元は、共通法線倍率の1に、
$H$ 内の $L$ の上記移送作用による安定群の階数を足したものである。
実際この一径数群の中心化群は分解を保ち、第一の四同定を固定すれば
残る相対倍率の核は六つの二重次数から共通の法線倍率だけになる。
$s^2\ne1$ の反転のない場合は階数1、片側だけの反転なら2である。
$s=1$ で非半単純な行列の射影中心化群の恒等成分は加法群なので、一つだけ非半単純なら階数2、
両方なら1である。量子窓の階数は少なくとも3なのでこれらは量子部分と交わらない。

非量子の場合の二平面分解を調べる。階数1なら分解を与える原始余指標は共役と符号を除いて一つである。
階数2では共通法線倍率のほか、片側の二文字だけを非スカラーに動かすトーラスがある。
これは $e=1,f=0$ なら $T(a,1)$、$s=1$ で一方だけ非半単純なら
半単純側の対角トーラスである。極大トーラス内の余指標が二文字ずつの二重みを持つには、
スカラーで作用される側の二文字が既に同じ重みなので、残る二文字も同じ重みでなければならない。
よってやはり共通法線倍率またはその逆である。
従って無印同型で追加されるのは分解の交換だけである。

交換では第一の四対応は恒等、第二は $L_{\pi_1(i)}^t$ となる。
交代成分 $P(1/s)$ を $P(s)$ に戻すには新しい能動位置0,2で $X$ を作用させればよい。
これで $e,f$ が入れ替わり、反転なしの一点は同じ一点、片側反転の二族は互いに移る。
片側反転を位置0,2に選んで比較すれば、追加の自己比較はなく第二行の同値が残る。
$s=1$ では交換に際して $P$ の変更は不要で、射影合同類の対が入れ替わる。
二次行列とその転置は射影合同である。実際 $C^t J$ は $CJ$ の逆行列と
射影共役になり、二次元射影行列はその逆と共役である。
従って最後の二行を得る。両側の取替えと逆行列は $02\mid13$ の支持階数も保つので
$s^2=1$ と他の非量子行は交わらない。
\end{proof}
\begin{proof}[一移動]
一移動と同時に分解を交換すると、新しい能動四テンソルは $P$ の一巡回で、
$K$ が対称だから元の $P$ に等しい。第一の四対応は旧第二の $L_{i+1}$ であり、
これを恒等に戻した後の第二の四対応は $L_{i-1}^{-1}$ である。
反転なしでは $L=1$ が戻る。片側反転 $L=T(a,1)A$ ではここから再び分解を交換し、
能動位置0,2の $X$ で $P$ を戻すと、第二の四対応は
$X^{[i\ {\rm even}]}L_{\pi_1(i)-1}^{-t}X^{[i\ {\rm even}]}$ である。
偶位置は元と同じで、奇位置の $Z,1$ が入れ替わるだけであり、自由に消せる対角係数で元に戻る。
$s=1$ の一移動と交換では、先頭の二行列が $D^t,C^{-1}$ になる。
$C^{-1}$ は $C^t$ と射影合同なので無印の対は等しい。量子部は一定の次数付き表示を持つ。
これで全点の最小周期が1である。
\end{proof}

% !TeX root = ../papers/nc-p3.tex
\section{非半単純な相補逆行列の全補完と閉軌道}
\label{sec:split-jordan-transport}
前節の二次形式に関する行列の同一視を使い、
$A=M^\dagger M$, $S=\operatorname{tr}A$, $d=\det M$ と置く。
固有値を $x_i^2$ と書き、$\prod x_i=d$ とする。
\begin{theorem}\label{thm:split-jordan-transport}
$\det M\ne0$ とする。可逆二平面分解の窓に $M$ が現れる必要十分条件は、
$x$ が $f_0f_1\ne0$ を満たす前節の曲面・直線・六次式のいずれかの付け替えであることである。
重固有値ごとのJordan分割への追加条件はない。非常に偶な場合（一つ一つの鎖長が全て偶数）
の直交基底の向きの双方でも成立する。全ての補完の分解付き同型は
定理\ref{thm:split-inverse-transporter}の共役商である。
\end{theorem}
\begin{proof}[二つの逆テンソルの単一行列による表示]
標準的な二側の形式を互いに同一視する。$f_0=d$ として
\[
 16 f_1=S^2-2\operatorname{tr}(A^2)+8d,\qquad
 F(z)=Sz-2z^2+2d,\qquad G(z)=F(z)^2/z .
\]
実際、最初の等式はBell基底で前節の四一次因子を展開した等式であり、
双方が特殊直交群二つの不変式なので、単純スペクトルの稠密性により全行列に成立する。
勾配を形式で同一視すると、$\nabla f_0=d M A^{-1}$ と
$4\nabla f_1=M A^{-1}F(A)$ を得る。従って二つの逆テンソルは射影的に
$M_0=M A^{-1},\ M_1=M A^{-1}F(A)$ と書ける。その右Gram行列は
$A^{-1},G(A)$ である。第二式の勾配と二つの切断の逆転置との対応は
\path{computations/verify_split_jordan_transport.py} でも通常基底で確認する。
$\tau$ は行側と列側を交換する。転置後も可逆行列の左右Gramは相似なので、
移送があれば平方固有値の多重集合・行列式は半単純化した二つの逆テンソルでも合う。
これが半単純な $P$ の局所同値の必要十分条件だったので、前節のスペクトル条件は必要である。

自己随伴行列は各一般固有空間の中で、鎖のGram行列が反対角単位行列となる
Jordan鎖の直交和にできる。実際、最大鎖長を $m$ とし、冪零部 $N$ について
$\langle v,N^{m-1}v\rangle\ne0$ と取れば、その巡回空間は非退化である。
$k[N]/N^m$ の単元の平方根で生成元を変更して最後のGram係数だけを1にし、
直交補に帰納する。従って同じJordan型の自己随伴行列は直交共役である。
一つでも奇数長の鎖があればその上だけの符号で共役の行列式を変更できる。

各許容スペクトル上で $M=O L$ と取る。$L$ は上の各鎖で $A$ の単元平方根を取った
可逆自己随伴行列である。奇数長の鎖があれば符号を選んで $\det O=1$ にできる。
偶数長だけなら $\det L=\prod_\nu\nu^{\dim E_\nu/2}$ である（$E_\nu$ は一般固有空間）。
直線上のスペクトルは $(1,1,r,r)$ に $d=r$、曲面上で全多重度が偶数となるのは
これと重なる $(1,1,-1,-1)$ に $d=-1$ である。従ってここでも $O$ は特殊直交である。
座標・共通倍率の変更はこの存在性を変えない。
$M_1=O L A^{-1}F(A)$ の右因子は自己随伴なので、その随伴も二側の特殊直交作用で
$M_1$ に戻る。よって $\tau$ に伴う転置を許した相補逆行列の比較を、
$M_0,M_1$ そのものの二側特殊直交同値として判定できる。

直線上では $\nu\in\{1,r\}$ で $F(\nu)=4r$, $G'(\nu)=-16r\ne0$ である。
従って $G(A)$ は $(4r)^2 A^{-1}$ と同じ固有値・Jordan分割を持つ。
曲面では $x=(1,-uv,u,v)$、$r=u^2,s=v^2,C=(r-1)(s-1)\ne0$、
$d=-rs$ とおく。$\nu=(1,rs,r,s)$ の順に
$F(\nu)/\nu=(-C,-C,C,C)$。$G(\nu)=C^2\nu$ は、$A^{-1}$ の固有値を
$\nu\mapsto rs/\nu$ で対応させた後 $rsC^2$ 倍したものである。
$1$ と $rs$ または $r$ と $s$ が重なる場合だけ非自明な鎖が可能で、
どちらもこの置換がその重固有値を固定する。
$r=s$ の $\nu=r$ で $G'(r)=(r-1)^4$、$rs=1$ の $\nu=1$ で
$G'(1)=-C(r+r^{-1}-2)$ なので鎖の長さは保たれる。
単純固有値の六次式の場合は既に前節で解いた。
上で得たGram倍率を $c^2$ とすれば、直線と曲面の双方で $\det F(A)=c^4$ である。

こうして $M_1$ と $cM_0$ の行列式が等しく、右Gramが直交共役となった。
必要なのは特殊直交共役である。奇数鎖の場合は先の符号を使う。
偶数鎖のみなら、長さ4が一つのときは冪零部の平方の核、
二つの長さ2のときは各一般固有空間内の冪零部の核の直和を取る。
これは双方で同じ極大全等方二平面である。導関数が零でないので同じ核を持ち、
スペクトルの並べ替えもこの合計を変えない。この平面を保つ直交行列の対角ブロックは
$B,B^{-\mathsf t}$ なので行列式は1である。
右の同値行列が特殊直交と分かれば、左の行列もGramの等式と行列式の一致から特殊直交である。
\end{proof}

\begin{theorem}\label{thm:split-polystable-good-limits}
可逆二平面分解の全窓は半安定で、唯一の閉軌道も同じ可逆二平面分解の
標準Koszul左右AS正則窓である。非半単純な $P$ のどの補完も、同じ平方固有値と
行列式を保つ半単純な $P$ の補完を軌道閉包内に持つ。
\end{theorem}
\begin{proof}
標準化した右Jordan鎖に $1-m,3-m,\ldots,m-1$ の重みをつけ、
$N$ が重みを上げる向きとする。反対角Gramよりこれは直交一径数である。
誘導される減少フィルトレーションは鎖の選択に依存しない。実際、中心0の冪零重み
フィルトレーションは、各Jordan鎖で上記の重みになるものとして、核と像の交叉の和で書ける。
具体的に重み $\ge j$ の部分は $\sum_{a\ge0}\ker N^{a+1}\cap
\operatorname{im}N^{\max(0,a+j)}$ である。一本の長さ $m$ の鎖で、先頭から
$h$ 番目（番号0始まり）のベクトルについての条件が $2h-(m-1)\ge j$ になる。
一般固有空間ごとにこれを取り、左側には $O$ により同じ一径数を運ぶ。
$L$ はフィルトレーションを両向きに保つので $M$ も左右のフィルトレーションを対応させる。
次数付きの $M$ は可逆で右Gramの冪零部は消える。

$M_0,M_1$ の各Gram上でも、非自明鎖での導関数の非消滅から冪零部の全ての核・像は
対応する元のものと同じである。従って左右のこのフィルトレーションは両逆テンソルから
同じ規則で再現できる。移送 $v\in\tau D$ はこれらを行と列の交換を伴って保つ。
上の一径数を $\lambda(t)\in D$ とすれば $\lambda(t)v\lambda(t)^{-1}$ は群内に極限を持つ。
実際、フィルトレーションを保つ行列は重みを下げる成分がなく、逆にも同じことが成立する。
二側の特殊直交一径数を射影二文字因子に戻し、必要なら $t$ の有限冪を取って線形に持ち上げる。
第一の四対応を固定した前節の構成で、交代成分と第二の対応がともにこの次数付きに特殊化する。
各局所射影行列は可逆な極限を持ち、特殊線形持上げでは発散する共通倍率はない。
$f_0,f_1$ も非零のままであり、得られた窓は半単純な $P$ の補完である。

さらに、全基底群内で分解を保つ中心化群でのアフィン軌道閉包を考える。
六つの二重次数の成分ごとに、可逆な二対二切断の四次行列式を一つ取る。
高位位置集合 $I$ に対応する行列式は四位置の使用した二平面の行列式の二乗の積で変換する。
従って $I$ と補集合の二行列式の積は、各位置の行列式を1とした中心化群に不変である。
両交代成分では両切断を使ってよい。これらの不変式は元の窓で非零なので、
閉包内でも各境界成分は非退化二形式二つの積のままであり、両交代成分の必要な切断も可逆である。
切断階数の上界は極限へ延び、可逆な交差ブロックで下界も保たれる。
従ってこの閉包内の中心化群の閉軌道は零でなく、同じ正則枝に属する。
その閉軌道でHermiteノルムが最小になるように基底を取れば、各位置の
二半分はトーラスの異なる重みでGram行列の交差ブロックが零であり、
残る全ての跡零方向は中心化群の極小条件に含まれる。
よって全基底群でも閉軌道であり、還元と最小ノルムの議論は標数零の代数閉体に戻る。
\end{proof}

% !TeX root = ../papers/nc-p3.tex
\section{半単純対称性の四位置での一致と三つの完全分類}
\label{sec:sl2-window-profiles}
$L_d=\operatorname{Sym}^{d-1}k^2$ を $\mathrm{SL}_2$ の $d$ 次元既約表現とする。
\begin{theorem}\label{thm:sl2-window-profiles}
標準AS窓の射影安定群への非自明な $\mathrm{SL}_2$ 準同型を取る。四位置の矢表現は全て同じ型であり、
各関係表現も全位置で同じになる。可能な型と、最後の三行で印を忘れた全正則窓は次の通りである。
\[
\begin{array}{c|c|l}
 V_i & R_i & \text{窓}\\\hline
 L_2+2L_1 & 2L_2+2L_1 & \text{下記の二次元係数表示}\\
 L_3+L_1 & 2L_3 & \text{多項式環または }U(\mathfrak{sl}_2)\text{ のPBW斉次化}\\
 2L_2 & L_3+3L_1 & \text{多項式環または下記の非自明な単冪型}\\
 L_4 & L_5+L_1 & \text{多項式環}
\end{array}
\]
三つの同型類の最小周期は全て1で、反対階数は多項式環6、斉次包絡環8、単冪型12である。
\end{theorem}
\begin{proof}[全表現の制限]
単連結性により四射影表現を線形に持ち上げ、最上文字は自明である。
各四次元表現は4の分割、各六次元表現は6の分割であり、相補二位置の関係表現は双対、従って同型である。
$L_i\otimes L_j=L_{i+j-1}+L_{i+j-3}+\cdots+L_{|i-j|+1}$ から各隣接積と対辺の積の
両方に入る六次元部分表現を列挙する。極大トーラスのEuler四文字はWeyl反転で不変だから
$\{m,-m,n,-n\}$ である。$p_1$ の公式で各文字の絶対値は12以下となる。
\path{computations/verify_sl2_window_profiles.py} は整数値を一つの素数で還元した不一致だけで
この表と自明作用以外を全て落とす。$L_4$ の場合 $L_5+L_1=\bigwedge^2 L_4$ は積で一意な部分表現なので、
全辺の関係は四位置を同一視した通常の交換子である。
\end{proof}

\begin{proof}[$L_3+L_1$]
矢をスカラー $t_i$ と標準三次元直交空間の座標 $x$ に分ける。各辺に現れる三次元型の
三つのコピーは $t_i x,x t_{i+1},x\times x$ で、最後は三次元交代形式による包含を表す。
隣接する二スカラーの成分はない。四つとも $x$ の部分は第一切断で交代積二つの内積、
第二切断でも巡回した同じ形の倍数である。独立な三つの計量の対積で展開するとこの二式は比例せず、
四つとも $x$ の部分は零である。残る一スカラー成分を
$a_i t_i\epsilon(x_{i+1},x_{i+2},x_{i+3})$、二つの反対スカラーの成分を順に
$u t_0t_2(x_1,x_3),v t_1t_3(x_2,x_0)$ と書く（添字順のテンソルを表す）。
上述の三コピーで第一切断の重複度行列は
\[
 \begin{pmatrix}u&0&a_0\\0&v&a_1\\a_2&a_3&0\end{pmatrix}.
\]
他の辺では巡回する。必要な階数2から $v a_0a_2+u a_1a_3=0$。
$u=0$ なら一文字全階数から $a_0a_2\ne0$、従って $v=0$ となる。
このとき全 $a_i\ne0$ で、列空間は $x\times x$ と
$a_i t_i x+a_{i+1}x t_{i+1}$ を与える。位置別の $t_i$ の倍率で通常の交換子表示になる。
$uv\ne0$ なら $a_{i+2}=0$ のとき、前辺に $x t_i$、次辺に $t_i x$ が純関係として続くので
正則でない。全 $a_i$ を1にして共通倍率を調整すると $u=1,v=-1$。
$(-1)^i t_i$ を新しい中心文字にすれば残る三式は三次元交代積によるLie括弧の
斉次化である。包絡環のPBW定理または定理\ref{thm:regular-normal-linear}で全次数正則となる。
\end{proof}

\begin{proof}[$2L_2$]
$V_i=U_i\otimes L_2$, $\dim U_i=2$ とする。
四つの標準表現の不変形式は $[01][23],[12][30]$ が張る。ここで $[ij]$ は交代二形式である。
第一切断の $L_3$ の重複度は1なので、第二の不変形式の係数四テンソルは
$01|23$ で階数1である。第二切断から第一の係数も $12|30$ で階数1である。
従って非零二行列 $r_i$ を使い、全窓のテンソルが
\[
   A\, r_1 r_3[01][23]+B\,r_0 r_2[12][30],\qquad AB\ne0
\]
となる。$r_i$ は位置 $i,i+1$ の重複度添字に置く。
$J=\left(\begin{smallmatrix}0&1\\-1&0\end{smallmatrix}\right)$ として
$r_i=JF_i^{\mathsf t}$ と書くと、$F_i$ の基底変更は倍率を除き有向サイクルの写像の基底変更である。
自明表現成分の階数は各切断で3でなければならない。第一で
$2A r_1r_3+B r_0r_2$、第二で $A r_1r_3+2B r_0r_2$ を当該順で切る。
各々一方の項は階数1、他方は対辺の二行列の階数の積なので、二つの対辺がともに特異にはならない。

特異な辺が高々一つなら、巡回と基底変更後に $(F_0,F_1,F_2,F_3)=(I,I,I,M)$ とする。
二つの四次行列式は
\[
 -8A^3\det M(2A\det M+B\operatorname{tr}M),\qquad
 -8B^3(A\operatorname{tr}M+2B).
\]
$M$ が可逆なら $(\operatorname{tr}M)^2=4\det M$。
倍率を吸収して $A=1,B=-1$、$M=I$ または非自明な単冪Jordan行列とできる。
前者のテンソルは四つの積文字の交代体積である。両者とも標準表現の対角トーラスで
二平面ずつに分ければ純次数の二形式は $r_i$ の非零倍で、全て可逆である。
従って定理\ref{thm:split-inverse-regularity}が使える。
$M$ の行列式1の四乗根を全辺に置く表示に移せるので、どちらも一巡回で符号だけが変わり周期1である。
$M$ が特異なら跡は非零なので $M=\operatorname{diag}(1,0)$, $B=-A/2$ とでき、
始点1の四次元数39で除外される。

特異な二辺は隣接するので $F_0=F_1=I,F_2=\operatorname{diag}(1,0)$ とする。
二つの頂点空間で核と相手の像が一致するかどうかに従って
$F_3=E_{00},E_{11},E_{01},E_{10}$ の四場合になる。中央二つは隣接階数が5以下である。
他の二場合では $A=1,B$ 独立の関数体上で始点1の四次元数52となる。
これは実際の関係像の階数が6の全特殊化でも下界である。上記コードは四文字の全語に
二次関係を三通りに挿入し、標準表現のトーラス次数ごとに正確に掃き出す。
反対階数6,8,12も表示自体から直接得られるので、印を忘れても三種は区別される。
\end{proof}

\begin{proposition}\label{prop:one-standard-sl2-slice}
$L_2+2L_1$ の場合、二次元の自明文字を $U_i$、標準文字を $S_i$ とする。
テンソルは $P\in\bigotimes U_i$ と、二つの $S$ に交代形式、その補二位置に二行列を置く六成分で尽くされる。
$U_i,U_{i+1}$ を使用する境界二行列を $b_i$ とする。正則性には全 $b_i\ne0$ が必要である。
全 $b_i$ が可逆なら所定の二隣接表現階数を満たす全係数が標準Koszul左右AS正則である。
\end{proposition}
\begin{proof}
辺に $L_3$ がないので、四標準文字の成分は四つの各隣接対で交代的であり零である。
奇数個の標準文字の成分もない。$U$ を全て消す商では各辺の関係の像は標準二文字の
交代関係の倍数だけである。$b_i=0$ なら対辺のこの関係も零になる。
実際この射影には二つの標準文字を含む当該成分と、既に消えた四標準文字成分だけが寄与する。
四文字ごとに少なくとも一つの関係のない境界があり、間のブロックは少なくとも二つの
対称語を残すので指数増大となる。

全 $b_i$ 可逆の場合、各 $S_i$ を同率で拡大して射影極限を取り、$P$ だけを零にする。
混合対切断の標準表現成分は変わらず重複度2である。自明部分では、
純 $UU$ 対と $SS$ の交代成分との間の非零列・行だけが残り、その和の階数は2となる。
極限は $2+2$ の二平面分解を持ち、純次数関係は $b_i$ と $S_i$ の交代形式である。
全て可逆なので前記定理で正則であり、命題\ref{prop:unrestricted-good-reduction}で元の窓も正則となる。
\end{proof}

% !TeX root = ../papers/nc-p3.tex
\section{標準二文字と自明二文字を持つ全窓}
\label{sec:one-standard-sl2}
$V_i=U_i\oplus S_i$ とし、$S_i$ は同じ $\mathrm{SL}_2$ の標準表現、$U_i$ は二次元自明表現とする。
$J=\left(\begin{smallmatrix}0&1\\-1&0\end{smallmatrix}\right)$ と置く。
四文字テンソルの純 $U$ 成分を $P$ とし、異なる位置 $i,j$ が $U$ である成分を、
残る二位置 $k,l$ での交代形式に $\operatorname{sgn}(i,j,k,l) B_{ij}$ を掛けて表す。
これにより $B_{ji}=-B_{ij}^{\mathsf t}$ となる。
$b_i=B_{i,i+1}$ とし、$p=B_{02},q=B_{13}$ と書く。添字は4を法とする。

\begin{theorem}\label{thm:one-standard-sl2-classification}
この表現型の標準Koszul左右AS正則性は、所定の二隣接切断の表現階数と四つの $b_i$ の可逆性と同値である。
正則窓の添字保持同型類の重複のないパラメータは次の通りで、各類の最小周期は1である。
\begin{enumerate}
\item $p=J+Q$, $Q$ は非零対称行列、$P=0$。
$Q$ の階数1が一類、階数2では $\det Q\in k\setminus\{0,-1\}$ が完全不変量である。
\item $p=J$。追加データは零または射影二元二次形式 $T$ で、$T$ の階数 $0,1,2$ が三類を与える。
階数0は多項式環である。
\item $p=\left(\begin{smallmatrix}0&1\\1&0\end{smallmatrix}\right)$。
追加データは同じく $T=\left(\begin{smallmatrix}a&b\\b&c\end{smallmatrix}\right)$。
$T=0$ は一類である。$b=0$, $T\ne0$ のときは $a,c$ の非零成分数で二類に分かれる。
$b\ne0$ なら $b=1$ として $ac$ と $a,c$ の非零成分数が完全不変量となる。
\end{enumerate}
復元では $b_0=b_1=b_2=J$, $b_3=JM^{\mathsf t}$, $M=-pp^{-\mathsf t}$,
$q=p/\det p$ とする。後二行の $P$ は下記の $T$ に対応するテンソルである。
この表と定理\ref{thm:sl2-window-profiles}の他の表現型が共有する類は多項式環だけである。
\end{theorem}

\begin{proof}[切断と退化境界の排除]
各辺の標準型の関係は二平面 $E_i\subset U_i\oplus U_{i+1}$ と標準表現の積であり、
二つの直和因子は $U_i S_{i+1},S_i U_{i+1}$ に対応する。標準型の切断は、行と列の符号の変更を除き、
\[
 \begin{pmatrix}p&-b_3^{\mathsf t}\\ b_1&q\end{pmatrix},\qquad
 \begin{pmatrix}q&-b_0^{\mathsf t}\\ b_2&-p^{\mathsf t}\end{pmatrix}
\]
であり、各階数は2である。従って各ブロックは相補二辺の $E$ の完全ペアリングを二つの
座標射影で挟んだ行列であり、可逆性はその両射影が全射であることと同値である。
自明型の切断は、$UU$ と交代 $SS$ の順に、右上と左下が非零の境界ベクトル、右下が零、
左上が $P$ の行列である。非零性は命題\ref{prop:one-standard-sl2-slice}による。
階数2の条件は、各 $i=0,1$ で
\[
 P\in kb_i\otimes (U_{i+2}U_{i+3})+(U_i U_{i+1})\otimes kb_{i+2}
 \tag{*}
\]
となる。並びは巡回して読む。

$E_i$ の左・右射影の全射性を $L_i,D_i$ で表す。各頂点で $D_{i-1}$ または $L_i$ が必要である。
両方が欠ければ $u S_i$ と $S_i v$ が連続純関係になる。
各辺でも $L_i$ または $D_i$ が必要である。実際、$b_{i-1},b_i\ne0$ から両射影の階数は少なくとも1である。
両方が1なら $E_i$ はその左右の像直線の直和である。$b_{i-1}$ の右支持はその左像直線に入り、
$b_{i-1}$ 自身が純 $UU$ 関係なので、その終文字は次の $U_i S_{i+1}$ 純関係の始文字にもなる。
いずれも定理\ref{thm:no-consecutive-pure-relations}に反する。

各反対ブロックの可逆性は $L_i L_{i+2}$ と $D_{i-1}D_{i+1}$ のいずれでも表せる（積は論理積）。
$p,q$ が両方特異なら各反対頂点対は、各頂点の必要条件を満たすために、一方で左だけ、他方で右だけを持つ。
全四射影条件を各辺にも一個ずつ配るには各頂点の選択が巡回して一定でなければならず、矛盾する。
両方可逆なら全射影と全境界が可逆となる。一方だけ可逆なら、巡回によりそれを $p$ としてよい。
頂点0,2は両射影を持ち、1,3では反対の側を一個ずつ選ぶ。
$b_i$ の可逆性は $D_{i-1}L_{i+1}$ なので、二つの特異境界は隣接する。

特異辺を0,1に移す。$b_i=JF_i^{\mathsf t}$ で表すと $F_i$ は射影倍率を除き有向四辺の写像として基底変更される。
可逆二辺を恒等に縮めれば、二つの二次元空間を行き来する二本の階数1写像が残る。
各空間で入射像を $ke_0$ にし、出射核を $ke_0$ または $ke_1$ にする。
位置1での一致は連続する純 $b_0,b_1$ を与えるので許されない。
従って $F_1=E_{00}$, $F_2=F_3=I$, $F_0=E_{00}$ または $E_{01}$ としてよい。ここで $E_{ab}=e_a e_b^{\mathsf t}$ である。
両方で $(*)$ の18本の線形条件は階数16で $P=0$ を強制する。
すると各辺の自明型関係の $UU$ 部分は $kb_i$ のみである。
全射でない左射影の像を消す非零 $U_i^*$ でその辺を評価すると、標準型の関係は消え、評価階数が1以下となる。
右も同じで、定理\ref{thm:all-endpoint-ranks-two}により全射影が必要になる。特異境界は全て除かれた。
逆に可逆境界に対する十分性は命題\ref{prop:one-standard-sl2-slice}の全次数の極限論法である。
\end{proof}

\begin{proof}[正規形と周期]
標準文字のスカラーを四つの基底変更に吸収すると、各 $B_{ij}$ には $g_i B_{ij}g_j^{\mathsf t}$、
$P$ には $\bigotimes g_i$ が作用し、さらに窓全体の射影倍率 $\lambda$ がある。
三辺を $J$ に揃え、上記二切断のSchur補を取ると
$q=-Jp^{-1}b_3^{\mathsf t}=-Jp^{-\mathsf t}J=p/\det p$, $M=-pp^{-\mathsf t}$ となる。
残る変更は $g_0=g_2=g$, $g_1=g_3=g/(\lambda\det g)$ である。
$M$ の射影Jordan型に従い $(*)$ を解く。$M=\operatorname{diag}(1,t)$ のとき18行16列の係数行列に
行列式 $-(t-1)^3/t^3$ の小行列があり、非自明単冪型でも階数16となる。
$M$ がスカラーなら核の次元は3で、$\Sym^2 U\otimes\det U$ である。
明示的には $(a,2b,c)$ を、語 $0000,0001,\ldots,1111$ の辞書順で次の三列の係数とする。
各列は表示した語の符号付き和である。
\[
 \begin{aligned}
 &-0001+0010-0100+1000,\\
 &-0101+1010,\\
 &-0111+1011-1101+1110.
 \end{aligned}
\]
$g$ の四重作用をこの三列に掛けると $T\mapsto(\det g)gTg^{\mathsf t}$ を得る。
従って正規化後は $(p,T)\mapsto(\lambda gpg^{\mathsf t},gTg^{\mathsf t}/(\lambda\det g))$ となる。
$g$ のスカラー部分と $\lambda$ により二成分の倍率は独立である。
$M$ がスカラーであることは $p$ が対称または交代であることと同値である。
それ以外は交代部を $J$ にして $p=J+Q$ とすれば、$Q$ には $\mathrm{SL}_2$ の合同が作用する。
二次形式の合同分類と $\det(J+Q)=1+\det Q$ が第一行を与える。
交代の場合は $\PGL_2$ 全体、対称の場合は等方二直線を保つ対角トーラスとその交換が $[T]$ に作用し、表の残りを得る。

一巡回して符号規約の共通符号を吸収すると新しい反対ブロックは $q$ である。
三つ目の境界を $J$ にするには新位置3だけに $M^{-1}$ を作用させればよい。
従って新 $p$ は元の $p$ に比例する。$P\ne0$ のとき $M$ はスカラーであり、上記三列の巡回は全て符号反転である。
独立な倍率を使えば元の記録に戻るので周期は1である。
\end{proof}

\begin{proof}[印を外した同型]
射影安定群のLevi部分群を取り、その半単純導来群の単連結被覆を $D$ とする。
標数零で半単純部分群はLevi内へ共役できる。位置0への射影は $D$ 上有限核である。
実際、正次元の正規核には単純因子の根 $\mathrm{SL}_2$ が含まれ、四位置の型一致に反する。
従って $D$ を位置0の四次元表現で調べられる。
四次元以下に非自明な表現を持つ単純Lie代数はWeyl次元公式により
$\mathfrak{sl}_2,\mathfrak{sl}_3,\mathfrak{sp}_4,\mathfrak{sl}_4$ である。
単純な場合、表の印は一つの共役類になる。$\mathrm{SL}_3$ では三次元標準または双対上で $2+1$、
$\mathrm{SL}_4$ では四次元標準または双対上で $2+1+1$ の表現による。
$\mathrm{Sp}_4$ では標準二平面は自明文字と直交する非退化平面であり、共役で全て移り合う。

複数因子が作用できるのは二つの $\mathrm{SL}_2$ の標準の直和またはテンソル積に限る。
テンソル積ではどの埋込みも表の型にならない。直和では表の印は個別の因子である。
このとき第二因子は四位置とも $U_i$ に標準として作用し、$S_i$ には自明に作用する。
三辺 $J$ により四つの $U_i$ 上の作用は同じになり、$p$ は交代形でなければならない。
$\Sym^2 U$ の不変元は零なので $P=0$ である。このテンソルは四文字の全交代体積で、
多項式環では二つの印は実際に共役である。従ってどの記録も印を外すことで他の記録と一致しない。

斉次包絡環と非自明単冪型について、窓テンソル上の四位置の無限小基底変更の和を零として解くと、
位置0への安定Lie代数の射影次元は順に3,4である（スカラーを除く）。
既に三次元単純因子があり、半単純部分の次元はこれより大きくなれない。
その既知の矢表現はそれぞれ $L_3+L_1,2L_2$ であるから、いずれも本表と重ならない。
全線形計算は \path{computations/verify_one_standard_sl2.py} で再現できる。
\end{proof}

\begin{corollary}\label{cor:semisimple-symmetry-closed-limits}
非自明な連結半単純対称性を持つ全正則窓は半安定であり、軌道閉包の唯一の多重安定類も正則である。
定理\ref{thm:one-standard-sl2-classification}の表では $P$ を零にし、第一行で $\rk Q=1$ なら
さらに $Q$ を零にすればその類を得る。斉次包絡環と非自明単冪型の閉類は多項式環である。
\end{corollary}
\begin{proof}
各位置で $U$ を $t$, $S$ を $t^{-1}$ 倍する特殊線形変換の極限は二標準文字成分を保ち $P$ を消す。
階数1の $Q$ は $\mathrm{SL}_2$ の対角一径数群で零にでき、三つの $J$ と $p$ の交代部を保つ。
残る $p$ は全て非零の二つの反対角成分だけを持つ表示を取れる。非退化対称部の二等方方向を基底にすればよい。
すると各 $b_i,p,q$ も同じ支持で可逆で、四文字テンソルの支持は四文字 $u_0,u_1,s_0,s_1$ の
24通りの順列全体である。全順列の一様平均は四位置全てで均一で、重み多面体の相対内部にある。
従って対角群の軌道はアフィンで閉であり、複素数上では対角変換により各位置の対角モーメントを零にできる。
一位置だけが異なる二語はないから非対角モーメントも零となり、Kempf--Ness判定で全群について閉である。
一般の標数零代数閉体へは有限生成定義体からの基底変更で移る。

斉次包絡環では前節の表示でスカラー文字を $t^3$, 三次元文字を $t^{-1}$ 倍すると
一スカラー成分が残り二スカラー成分が消え、多項式環を得る。
非自明単冪型では重複度の $M$ を行列式1の対角変換で共役して $I$ に落とせる。
非自明半単純群は根部分群から来る $\mathrm{SL}_2$ を含むので二つの定理が全場合を覆う。
述べた極限はすべて非零アフィン閉軌道であり、半安定性と一意性が従う。
\end{proof}

% !TeX root = ../papers/nc-p3.tex
\section{各矢に高々二重みのある全ての乗法群作用}
\label{sec:binary-torus-windows}
\begin{theorem}\label{thm:binary-torus-windows}
標準正則窓の非自明な乗法群が各矢に高々二つの重みを持つとする。各位置のスカラーを除くと
四位置とも同じ正の重み差を持ち、重み重複度も同じである。共通の重み差を1にする。
\begin{enumerate}
\item 重複度 $3+1$（または反転）の全窓は、三生成二次三次元Koszul左右AS正則な通常の
次数付き代数 $S$ に中心文字を加えた $S[z]$ の添字反復であり、最小周期は1である。
重みを保つ添字同型は $S$ の次数付き同型で完全に決まる。
\item 重複度 $2+2$ の全窓は定理\ref{thm:split-inverse-transporter}の可逆チャートで尽きる。
特に二軸の純二次関係のどれか一つでも退化すれば正則でない。
定理\ref{thm:split-jordan-transport}と系\ref{cor:split-finite-unmarking}の
スペクトル条件・残余共役・有限個の印の検査はこの表現型全体に適用できる。
全て二次PBWで、多重安定極限も同じチャートで正則である。
\end{enumerate}
\end{theorem}

\begin{proof}[等間隔の不均一型の排除]
まず全ての非自明な位置で差が1の場合を扱う。系\ref{cor:two-weight-profiles}の $0121$ では
位置1の高位文字の双対点で縮約する。位置0の端支持は $B_0$ の重み1の次元から2以下であり、
位置2の端支持は最上重み1なので二次元の低位空間に入る。
定理\ref{thm:all-endpoint-ranks-two}によって両端が2となり、定理\ref{thm:no-two-sided-centers}に反する。
$0323$ では重みを反転しスカラー位置を正規化すれば同じである。

$1133$ の $B_1=(0,6,0)$ の場合、全ての台で位置1,2の高位文字の個数は1である。
従ってこの二位置だけで高位を拡大する別の乗法群が窓を保つ。それは隣接二位置だけに非自明で、
すでに系\ref{cor:two-weight-profiles}のどの行にもない。
残る $B_0=(4,2,0),B_1=(1,4,1)$ では、位置0,1の高位直線を $kh_0,kh_1$、
位置2,3の低位直線を $kx_2,kx_3$ とし、その補三空間を $W_i$ とする。
$01|23$ の混合切断は二つの三空間ずつの間の階数2行列で、各側の像は二平面である。
$12|30$ の重み0と2が各1なので、$R_1$ に $a_1x_2,h_1a_2$、
$R_3$ に $a_3h_0,x_3a_0$ という純関係がある（$0\ne a_i\in W_i$）。
従って $R_0$ の混合二平面は $h_0W_1,W_0h_1$ の一方だけに属する非零元を持てない。
どちらも上の純関係と連続するためである。同じく $R_2$ の二平面にも $x_2 W_3,W_2x_3$ の単独支持はない。
よって四つの三空間への射影は全て階数2で、混合切断の各三次正方ブロックは階数2となる。
しかし $h_0W_1x_2W_3$ 成分は $12|30$ の重み0のブロックそのもので階数1である。これも矛盾する。
\end{proof}

\begin{proof}[異なる重み差の排除]
非自明な位置数 $n\le4$ とし、高位・低位の選択を立方体 $\{0,1\}^n$ の頂点で表す。
実際に係数のある頂点は、正の差の列 $d$ と最上重み $H$ に対する $d\cdot x=H$ に含まれる。
この台上で定数となる任意の整数係数の一次式も射影窓を保つ一径数を与える。
従って台のアフィン包に直交する非零の $\{-1,0,1\}$ ベクトルで、
四位置へ拡張したとき零成分を持つものはない。あれば負符号の位置では二重みを取り替えてスカラーを引き、
前段で排除したスカラー位置付きの等間隔作用になるからである。

台のアフィン基底として高々 $n$ 頂点を取れる。その差ベクトルの直交を全て調べると、上記条件に適い
正の法線を持つ基底は $n=4$ の四点のみで、原始法線の成分は等しいか、並べ替えて $(1,1,1,2)$ である。
後者の $H$ は2または3で、反転により2とできる。巡回で差2を最後に置き、
$k_i\in\{1,2,3\}$ に対して $\sum d_i k_i=8$ を課す。各隣接対と相補対の重み次元から
六関係の重複度を列挙する。定理\ref{thm:torus-euler-characters}の $p_1$ の式に出る重みは0から5に入り、
四指数の和は8である。全候補が $(t,y)=(11,7)$ でその四一次因子の積と不一致となる。
二つの有限整数計算は \path{computations/verify_binary_torus_windows.py} による。
\end{proof}

\begin{proof}[純粋関係と曲面、動く軌道の終点]
以後 $X_i,Y_i$ を低位・高位空間とする。曲面射影を持つ正則窓には純粋関係 $u v$ はない。
定理\ref{thm:no-moving-plane-pencils}と定理\ref{thm:isolated-surface-fibers}によれば、
整二次曲面商では全空間を張る整なグラフの一般点でこの積は消えない。
平面商では正則中心一次拡大の表示で、中心文字による随伴次数は三変数多項式環の中心多項式拡大であり整域である。
多項式環自身でも明らかである。従って以下で純関係が生じれば、全射影像が高々曲線である。

曲面のない場合に、乗法群が非自明に作用する後続直線を持つ固定点から
定理\ref{thm:finite-branch-corridors}の曲線列をたどる。
一般後続は一意なので全ての中間写像は同変である。
各中間像は非自明軌道の閉包で、二重みしかないため高位軸の一点と低位軸の一点を結ぶ直線である。
最後の直線では全点の後続がただ一つの同じ終点になる
（逆向きに同定理を適用、または命題\ref{prop:branch-corridor-degree}の $f_m=0$）。
特にこの直線の二つの軸端点にも同じことがいえる。
\end{proof}

\begin{proof}[$3+1$ の正規直線]
高位文字を $z_i$ とし、その双対点を $p_i$ とする。$B_i=(3,3,0)$ なので全 $p_i$ は連続して点列となる。
混合三関係は $X_i z_{i+1}\oplus z_i X_{i+1}$ の三平面である。
一方の射影が同型でないと純関係があり、対応する $p_j$ での前または後の評価階数が2以下である。
時間を必要なら逆にして $p_j$ の後続が直線となる場合を扱う。
その直線は $p_{j+1}$ を含み、高位空間は一次元なので群が非自明に作用する。
前段の曲線列の各中間直線の高位端はその位置の $p_i$ である。
最後の高位端には元の全高位点列で高位の後続があるから、列の定終点も高位点である。
従って固定した二つの高位点の間に、最初の直線の点を変えた異なる経路がある。
高位点列で同じ開始位置へ四周期を合わせて戻れば、二経路を独立に反復して指数増大する。
実際異なる内部点での評価は非比例で、両端を固定した各ブロックの独立な評価をテンソル積にできる。
これは標準成長に反する。

よって全ての混合射影は同型で、$X_i z_{i+1}=z_i X_{i+1}$ となる。
$z$ を右へ動かすと、全語は低位語の後に $z$ の列を付けたもので張られる。
重み0の部分代数の各長さの次元はKoszul指標の逆公式で $\binom{n+2}{2}$ である。
これは各位置の重み0のKoszul係数が $(1,3,3,1)$ だからである。
全長 $n$ のこれらの上界を足すとちょうど $\binom{n+3}{3}$ になるため各正規表示は独立である。
逆向きにも同じで $z$ 列は正則正規となる。定理\ref{thm:regular-normal-linear}で $z$ を中心化する
位置移動は低位空間を保つ。低位関係は純粋なので補正 $\alpha,\beta$ は零であり、$S[z]$ を得る。
逆も同定理による。印を保つ同型は中心化と可換で、元の $S$ の次数付き同型を与える。
三生成関係の全次数判定には補題\ref{lem:ternary-finite-graph-test}を使える。
\end{proof}

\begin{proof}[$2+2$ の特異純次数関係の排除]
色階数は $(1,4,1)$ で、各軸に一本ずつの関係がある。どの軸上の点も同じ軸上に少なくとも一つの後続を持つ。
従って非自明に動く後続直線から始まる上記の列は存在できない。
最後の直線の二軸端点が、それぞれ自分と同じ軸に後続を持ち、定終点が互いに素な二つの軸に同時に属するからである。

低位の関係 $r_i$ が特異と仮定する。これは純関係で、曲面はない。$r_i$ を左で消す低位双対点 $p$ は、
低位の順軸直線全体を後続に持つ。一段後のその直線から始めて一般後続を取る際、
次の低位関係 $r_{i+1}$ も特異なら一般点の低位後続は一定である。
一般後続は一意なので動く位置が一個だけとなり、上記の列の $m\ge2$ に反する。
従って $r_{i+1}$ は可逆であり、次の動く像も低位の直線で $d_1=d_2=1$ である。
命題\ref{prop:branch-corridor-degree}の次数式の最初の項から $f_1=3-d_2=2$。
ゆえに最初の低位直線上に順側評価が階数2の点がある。そこでの純低位後続は $r_{i+1}$ の可逆性から一点のみである。
従ってその後続直線は不変で低位・高位をまたぎ、群が非自明に作用する。
これは前段で排除した場合である。高位の関係も重みの反転で同様であり、
定理\ref{thm:split-inverse-regularity}と定理\ref{thm:split-polystable-good-limits}が結論を与える。
\end{proof}

\begin{corollary}\label{cor:binary-torus-unmark}
二つの表現型の内部と相互に、印を外した同型は有限個の重み付き記録の集合の一致で決定される。
\end{corollary}
\begin{proof}
窓の射影安定群の一つの極大トーラス内で、上のいずれかの正規化された射影重みを持つ全ての一径数を探す。
各矢の固有文字の任意の二つの差への値は $0,1,-1$ に限られ、これらの差は極大トーラスの
有理指標空間を張る（射影四位置への作用が忠実）。従って候補は有限である。
全ての印は共役によりここに入り、共役な印は同じ重み付き記録を与える。
逆に一つでも記録の一致があれば、重みを保つ同型が実際の窓の同型である。
位置を巡回した窓との比較にも同じ判定を用いる。
\end{proof}

% !TeX root = ../papers/nc-p3.tex
\section{有限点グラフの連結トーラスと二つの重み付き局所型}
\label{sec:finite-torus-profiles}
\begin{theorem}\label{thm:finite-torus-profiles}
四辺の点対応が有限な標準正則窓では、各辺の六関係を保つ射影群
$\PGL(V_i)\times\PGL(V_{i+1})$ の部分群のトーラス階数は高々1である。
窓全体を保つ非自明な乗法群があれば、スカラーと共通倍率を除き四位置の矢および関係の重みは
全て等しく、可能性は下表に限られる。
\[
\begin{array}{c|c|c|l}
 V & R & h\text{ の指数}&\text{点スキームの長さ}\\\hline
 (0,1,2,3)&(1,2,3,3,4,5)&6&\text{最低・最高の二点に各10}\\
 (0,1,1,2)&(1,1,2,2,3,3)&4&\text{最低・最高に各9、中間の固定直線上に計2}
\end{array}
\]
後者の中間部分は二つの射影直線の積で二本の $(1,1)$ 形式が切る全交叉である。
いずれも必要条件で、値の列だけで正則性は主張しない。
\end{theorem}

\begin{proof}[露出面の次元と二端の重み差]
各辺の六空間 $R$ は純テンソルを含まない（定理\ref{thm:finite-frobenius-low-opposite-as}）。
最小重みの矢部分の次元を二端で $a,b$ とする。最小和の関係空間はこの $a\times b$ ブロックに
そのまま含まれ、Segreを避けるので次元は $(a-1)(b-1)$ 以下である。
二つの固有空間の射影双対の積に制限するとこの関係だけが残る。
有限交叉のためには少なくとも $a+b-2$ 本が必要である。従って $a=b=1$ または $a,b\ge3$ が必要で、特に2は現れない。

辺を保つトーラスは各位置で忠実に作用する。片端の作用が恒等なら、有限グラフのもう片端に
誘導される作用も点スキーム上恒等である。よって $x\wedge gx$ の各二次式はその上で消え、
定理\ref{thm:finite-square-chart}の二次制限の単射性から $g$ はスカラーである。
有効階数が2以上なら重みの凸包は二次元以上である。頂点の重複度が2ならその頂点を露出させれば矛盾する。
それ以外は三角形・四角形・四面体の中に、重複込みでちょうど二つの重みを持つ辺がある
（三角形に残りの一点があっても高々一辺だけに入る）。
この辺を露出させる一径数を使っても同じ矛盾となる。

窓作用に戻る。全位置で最低重みを0にする。最低または最高が三重なら全位置で同じ側が三重であり、
定理\ref{thm:binary-torus-windows}の $S[z]$ になる。これは有限点グラフを持たない。
実際 $z=0$ に三生成商の点対応が残り、三本の $(1,1)$ 形式を二平面の積に課しても次元は1以上である。
従って両極は単純である。各辺の最低どうしは点対となり、そこでの各側の評価は階数3である。
二端の最初の正の重みが違えば、小さい方の重みで可能な関係は片方が最低直線である純ブロックだけに入り、
そこには非零関係がない。しかし最低点で相手の全ての正の重みが評価像に必要である。
従って最初の差 $d$ は等しい。その重みでの二ブロックは各々片側が直線であり、
純関係がないことから両射影は単射、最低点での評価から両射影は全射となる。
ゆえに第二重みの重複度も等しい。反転すれば最後の差 $e$ とその重複度も全位置で共通である。
二重みだけなら定理\ref{thm:binary-torus-windows}で既に尽き、$2+2$ も純軸の点直線を持つので有限でない。
\end{proof}

\begin{proof}[単純四重みと固定点の重複度]
隣接矢の重みを $a=(0,d,d+s,d+s+e)$, $b=(0,d,d+t,d+t+e)$ と置く。
各和 $a_i+b_j$ は $(d,e,s,t)$ の整数一次式である。実際に成立する和の一致を全て保つ
新しい正の四数は、各関係重み空間を細分しないので、同じ辺の別の作用を与える。
整数解空間から射影二位置への写像は単射で、前段の階数上界から等号の行列は階数3である。
正の解を許す18枚の異なる超平面から三枚ずつを選び、符号付き三次小行列式を取ると、
正の原始候補は31本である。

各候補では $c_z=\#\{(i,j):a_i+b_j=z\}$ として重み $z$ の関係数 $0\le r_z<c_z$ を使う
（存在する和だけを数える）。全ブロックを関係にすることは純関係を含むため不可能である。
$r_{a_i+b_j}=0$ なら座標点対 $(i,j)$ は固定点である。逆も、局所指標の次式から従う。
その点対における局所環の指標（点がなければ零）は
\[
 L_{ij}(u)=
 \frac{\prod_z(1-u^{z-a_i-b_j})^{r_z}}
 {\prod_{k\ne i}(1-u^{a_k-a_i})\prod_{l\ne j}(1-u^{b_l-b_j})}.
\]
実際 $(n,m)$ を十分大きくし、六本の完全交叉Koszul分解の指標を取る。
完全対称式 $\operatorname{ch}\Sym^n V=\sum_i u^{n a_i}/\prod_{k\ne i}(1-u^{a_k-a_i})$
を両位置に入れると $u^{na_i+mb_j}$ の係数がこの式となる。
他方、有限スキームの全閉点は固定され、各局所因子で座標自明化により同じ係数は局所環指標となる。
異なる $(i,j)$ の係数は $n,m$ を動かしてVandermondeで分離できる。
これで $r_{a_i+b_j}>0$ の対には点がない。残る点対は部分マッチングをなし、$L_{ij}(1)$ は正整数で総和20である。

\path{computations/verify_finite_torus_profiles.py} がこれらの条件を全31候補の全六関係で比較する。
結果は $(d,e)=(1,1)$ として $(s,t)=(1,1),(1,2),(2,1)$ だけで、いずれも両極の点対が各長さ10である。
最初の関係重みは表の上段、他の二つでは $1,2,3,4,5,6$ である。
全四位置で倍率が共通なのは $d$ が共通だからである。
中央差2の位置数 $j$ は隣接禁止から高々2、矢の行列式文字の等式から $4H=24+2j$ となる。
従って $j=2$ なら交互に現れ $H=7$。この場合のEuler第一べき和は重み2の係数が $-2$ で、
四文字の指標にならない。以上で上段が従う。
\end{proof}

\begin{proof}[中央二重の四重み]
全位置が $(0,d,d,d+e)$ で共通である。$d+e\ne2d$ と仮定すると、重み $2d$ の関係は
二つの中央二平面の積に純粋に入り、その次元はSegreを避けるため高々1である。
両端で最高文字を零にして最低座標を1と置く。重み $d$ の二関係は中央二変数間の可逆な線形対応を与え、
残る式は重み $2d$ の高々一つの斉次二次式のみとなる。二変数上なので正次元の解が生じ、有限性に反する。
従って原始化すれば $d=e=1$。重み1と3の各二関係に加え重み2が二つで、他の成分は零である。
最低点のアフィン座標の重みは両端とも $1,1,2$ で、関係の重みは $1,1,2,2,3,3$ である。
正に次数付けられた六変数上の有限完全交叉なので局所指標は $(1+u+u^2)^2$ となる。
このチャートの閉点は最低点だけであり、その長さが9である。反転すれば最高点も同じとなる。
残りの閉点はグラフの一意性と固定性から中央どうしである。
中央二直線の積には重み2の二関係だけが制限され、有限なので交叉長は2である。
既に全長を使うため、これはスキームとして残余全部である。
\end{proof}

% !TeX root = ../papers/nc-p3.tex
\section{二つの長さ十の点を持つ乗法群窓の二成分}
\label{sec:simple-weight-finite-windows}
\begin{theorem}\label{thm:simple-weight-finite-windows}
有限点グラフに作用する乗法群の矢重みが $0,1,2,3$ である全窓は、添字を保ち印を外すと次の二族になる。
\begin{enumerate}
\item $(u,v)\in(k\setminus\{0,1\})^2$,
$4u^2-4u+4v^2-4v+1\ne0$。異なるパラメータは全て非同型である。
一添字の巡回は向きを一つ選ぶと $(u,v)\mapsto(v,1-u)$ であり、最小周期は
$(1/2,1/2)$ で1、それ以外で4である。
\item $b\in k\setminus\{0,1\}$, $3b^2-3b+1\ne0$、
同型は $b'=b$ または $1-b$ と同値で、全て最小周期1である。
\end{enumerate}
両成分はそれぞれ二次元と一次元で、第一の $(1/2,1/2)$ と第二の $1/2$ だけが一致する。
二つの判別式の零点でも下記の表示は正則で、そこで点対応が正次元になる。
全表示は二次PBWで多重安定である。
\end{theorem}

\begin{proof}[必須係数とPBW]
基底を重み順に $0,1,2,3$ と呼び、四つの数字でテンソル係数を指定する。
前定理の点対は両極の二対だけである。重み1,2,4,5の切断は階数1で、両側の全係数が非零となる。
実際、各和にあるどの座標点対も辺に入らないので、その座標で少なくとも一関係は非零である。
重み0,6では関係がない。重み3の二関係を両端の $0,3$ の直線の包に制限すると、
$03,30$ 上への射影は同型である。階数が1以下なら二つの射影直線の積に高々一本の式しかなく、
グラフが有限でない。従って重み3切断の両側で $03,30$ を選んだ小行列は可逆で、
隣り合う $00,33$ は消えるから $0303,3030$ の非零対角行列である。
片側だけ $12,21$ を選ぶ八係数は、他方向の階数1ブロックに入り全て非零である。

以上から関係の先頭語を、左辞書順で数字の降順二文字六つにできる。重み1,2,4,5で順に
$10,20,31,32$ が得られる。重み3では $30$ を先に掃き出した後の $03$ の関係に
$21$ が非零係数で現れるため $30,21$ の二つになる。これは上の対角小行列と非零の交差係数から従う。
長さ3の標準候補は20個であり、四テンソルの四つの独立な一文字縮約から実次元も20以上である。
従って全ての三次の重なりが解消し、定理\ref{thm:split-inverse-regularity}と同じDiamond補題と
最上テンソルの完全対合で全次数AS正則となる。この議論には点グラフの有限性の代わりに
今得た非零係数と切断階数だけで十分である。
\end{proof}

\begin{proof}[六変数の再現正規形]
和6の語のうち隣り合う和が全て1以上5以下であるものを用いる。
$1212,2121$ の二係数を $x,y$ とし、残る38語の係数は全て非零である。
二切断の重み1,2,4,5の階数1は40本の二項式で与えられ、その指数行列は階数20である。
各語の四位置と数字への接続行列の階数は12である。辞書順にこの行列の独立行を最初から選び、
その12語の係数を1にする。行格子のSmith非零対角成分は全て1で、残余有限被覆はない。
40式にこの12正規化を加えて指数を行簡約すると、次の辞書順の六係数が独立変数となる。
\[
 (w_{3030},w_{3102},w_{3111},w_{3120},w_{3201},w_{3210})=(a,b,c,d,h,f).
\]
残りはこれらのLaurent単項式である。この規則と \path{computations/verify_simple_weight_torus.py} の
\texttt{monomial\_chart} は全係数を復元する。

重み3の両切断の順を $03,30,12,21$ とすると行列は順に
\[
 \begin{pmatrix}1&0&1&d/c\\0&a&1&c\\h/c&1&x&1\\b&dh/(bc)&1&y\end{pmatrix},
 \qquad
 \begin{pmatrix}a&0&1&1\\0&1&dh^2/(bcf)&b/c\\d&1&y&1\\f&1&h/b&x\end{pmatrix}.
\]
左上を掃き出して右下を零とする。最初の対角式は $x=h/c+1/a$, $y=bd/c+dh/(ab)$ を与える。
残りから
\[
 d=a(c-b)/c,\qquad h=bc^2(1-b)/(c-b),\qquad f=1+a(h-b)/c,
 \qquad (ab-c)(a(1-b)-c)=0 .
\]
従って $b,c$ は非零、$b\ne1$, $b\ne c$ である。$a=c/b$ のとき全残余式は恒等的に零。
$a=c/(1-b)$ では $2b=1$ は前枝で、他は $c=1$ が必要十分である。
実際一式の分子は単項式を除き $(2b-1)^2(c-1)(bc+b-c)$ で、最後の因子を零にすると
残る式は $b(2b-1)\ne0$ の下で非零である。
第一族の座標を $u=b$, $v=1-b/c$ とし、第二族は $c=1$ の $b$ を用いる。
この範囲で $a,b,c,d,h,f$ は全て単元であり、$x,y$ の零は許される。
\end{proof}

\begin{proof}[有限グラフと同型]
辺の二点の最低座標が一方だけ非零なら重み1,2,3の式を順に使って他方の全座標が消える。
最高座標も同じである。両方で最低・最高ともに零の場合は重み2,4で $11,22$ が消え、
重み3の交差係数から点対はない。従って最低どうしか最高どうかのアフィンチャートで尽きる。
最低座標を共に1とし第一点の重み1,2の座標を $t,z$ とする。重み1,2の式で第二点の重み1,2を解き、
重み3の二式で両方の最高座標を解く。残る重み4,5の式は
\[
 A z^2+Bzt^2+Ct^4,\qquad t(Dz^2+Ezt^2+Ft^4)
\]
で、$A\ne0$ は重み2,4の必須係数から従う。原点以外では $t\ne0$ で $t=1$ にできる。
全辺と反転した最高チャートで二式の終結式を取ると、単元を除き第一族で
$4u^2-4u+4v^2-4v+1$、第二族で $3b^2-3b+1$ となる。
終結式が零なら非固定軌道の曲線があり、非零なら四つのグラフ全てが有限である。
PBWの証明はこの因子を使わない。

有限グラフ上では前定理の階数1と極大トーラスの共役により、印を外す追加操作は重みの反転だけである。
反転後の12係数を再び1にすると第一族は恒等、第二族は $b\mapsto1-b$ となる。
同じ再正規化で一巡回は第一族で $(u,v)\mapsto(v,1-u)$、第二族で $b\mapsto1-b$ となり、
同型・交叉・周期の結論を得る。上のコードは終結式およびこの単項式座標変換も検算する。

最後に全ての語は一位置だけ違う別の語と共存しない。24の順列は全て現れ、
それらの重みのアフィン包は行和・列和を固定した全単体行列の空間である（置換行列の凸包）。
追加語の数字頻度を一様頻度から引いた差は、符号を揃えて
$A=(1,-1,-1,1),-A,C=(-1,2,-1,0),D=(0,-1,2,-1)$ で尽き、全四種が現れる。
$2A-A+C+D=0$ なので各追加語に十分小さい正の重みを与えて列和を揃えられる。
残る各位置の偏りは、24順列の一様平均が相対内部にあるためその正の係数を少し調整して打ち消せる。
よって対角群について多重安定であり、非対角モーメントが恒等零なので全群でも多重安定となる。
\end{proof}

\begin{proposition}\label{prop:double-middle-finite-pbw}
有限点グラフのもう一方の重み $0,1,1,2$ でも、前定理の関係重複度と端点の評価、
二つの隣接切断、四つの一文字全階数を持つ有限候補は全て二次PBWかつ標準Koszul左右AS正則である。
\end{proposition}
\begin{proof}
最低・最高文字を $x,z$、中央を二平面 $U_i$ とする。重み1の二関係は $xU,Ux$ の間の可逆なグラフ、
重み3の二関係も $zU,Uz$ の間の可逆なグラフである（両極の一意な点評価）。
重み2の $xz,zx$ への射影は同型である。そうでなければ両極文字の二直線の積に曲線が残る。
よって $xz=f_i, zx=g_i$, $f_i,g_i\in U_i U_{i+1}$ と書ける。
両者は二次行列として可逆である。実際、四極文字成分は $xzxz,zxzx$ の非零対角のみであり、
$f_i,g_i$ は相補位置を $xz,zx$ にした二中央文字成分の非零倍である。
その文字並びをもう一方向で切ると、重み1対3の階数2行列の二次正方部分である。
両側で二つの射影が同型なため、各部分が可逆になる。

各 $U_i^*$ に $f_i$ の零点列を一つ取る。可逆性から任意の一点より両方向へ延ばせる。
位置ごとに $u_i,v_i$ のうち大きい方 $v_i$ の双対がその点となる基底を取れば、
$f_i$ に $v_i v_{i+1}$ は現れず $v_i u_{i+1}$ は現れる。
$x<u<v<z$ の左辞書順で、重み1では $ux,vx$、重み3では $zu,zv$、
重み2では $zx,vu$ が先頭語となる。先と同じ三文字計算で全ての重なりが解消する。
\end{proof}

% !TeX root = ../papers/nc-p3.tex
\section{中央二重の有限点窓の唯一の曲線}
\label{sec:double-middle-finite-windows}
\begin{theorem}\label{thm:double-middle-finite-windows}
有限点グラフ上の $0,1,1,2$ 型の全正則窓は
$\alpha\in k\setminus\{0,1,-1\}$ による一族で、添字保持同型は
$\alpha'=\alpha$ または $\alpha'=\alpha^{-1}$ と同値である。全て最小周期1で多重安定、
点スキームは長さ9の二点と二つの単純点からなる。$\alpha=-1$ の同じ表示も正則で点対応は正次元である。
従って定理\ref{thm:simple-weight-finite-windows}と本定理で、有限点グラフと非自明トーラスを同時に持つ
全窓のパラメータ・例外値・印なし同型・最小周期が決まる。いずれの射影安定群の恒等成分も乗法群である。
\end{theorem}
\begin{proof}[中央四文字と二つの鉛筆]
$x,U_i,z$ を重み別に取る。全体重み4のうち両極だけの成分は交互に並べた二語のみで、どちらも非零である。
四位置の $x,z$ の倍率を使って両係数を1にする。中央四文字成分を $P$ とし、$xz,zx$ を位置 $i,i+1$ に持つ成分の
残り二文字の行列をそれぞれ $A_i,B_i$ とする。$x,z$ を位置 $i,i+2$ に持つ成分の行列を
行 $i+1$、列 $i+3$ の順に $C_i$ とする。これらは全て可逆である。実際、重み1対3の切断の両側で
二つの射影が同型であり、各行列はその切断の一隅である。

重み2対2のSchur補は次の二式である。反対位置の二関係から来る中央行列は
$A_{i+2},B_{i+2}$ そのものである。
\[
 P=A_{i+2}\otimes A_i+B_{i+2}\otimes B_i\qquad (i=0,1).
\]
各辺の二行列は二次元鉛筆を張る。そうでなければ中央二直線の積に曲線が残る。
また鉛筆は共通の行核・列核を持たない。各鉛筆と階数1の二次曲面の交叉は、
相異なる二点または重複した一点である。

第一鉛筆が二点型なら左二位置の基底で $00,11$ を張る形となる。
$P=00\otimes H+11\otimes K$ の次切断の階数は $\rk H+\rk K=2$。
両者の行直線と列直線はそれぞれ独立なので $P=0000+1111$ にできる。
重根型なら第一鉛筆は $00,01+10$ を張り、$P=00\otimes H+(01+10)\otimes K$ となる。
次切断は転置や行列並べ替えを除き
$\left(\begin{smallmatrix}H&K\\K&0\end{smallmatrix}\right)$ の階数2を要求する。
従って $K$ は階数1で $H$ はそのSegre接平面にあり、最後の二位置で $K=00,H=01+10$ にできる
（行列鉛筆の正則性で両交差係数が非零である）。よって $P=1000+0100+0010+0001$。
これで全辺の型がそれぞれ同一になった。
\end{proof}

\begin{proof}[混合切断の一巡]
二点型では鉛筆の順序付き基底を $00,11$、重根型では $00,01+10$ とする。
$D_i$ をこの基底での $A_i,B_i$ の二列とすると
$D_{i+2}=J D_i^{-\mathsf t}$ であり、$J$ は二点型で恒等、重根型で二座標の交換である。
さらに、重み1対3の行を $xU,Ux$、列を $zU,Uz$ とすると
\[
 \begin{pmatrix} C_i&B_{i+3}\\ A_{i+1}^{\mathsf t}&C_{i+1}^{\mathsf t}\end{pmatrix}
 \quad\Longrightarrow\quad C_{i+1}=B_{i+3}^{\mathsf t} C_i^{-\mathsf t}A_{i+1}.
 \tag{1}
\]
従って $D_0,D_1,C=C_0$ が四回の式(1)で閉じることだけが残る。

二点型では中央文字の対角倍率と極文字の倍率を使い
\[
 D_0=\begin{pmatrix}1&1\\1&\alpha\end{pmatrix},\qquad
 D_1=\begin{pmatrix}1&1\\1&\beta\end{pmatrix}
\]
にできる。$\alpha,\beta\ne0,1$ である。実際中央二文字の倍率比 $r_i$ は
$\prod r_i=1$ としてよく、$r_2r_3,r_3r_0$ を独立に指定できる。
各列には極文字の倍率積 $x_0z_1,z_0x_1,x_1z_2,z_1x_2$ が独立に作用する
（残りの二つの極倍率は交互極語の正規化から決まる）。
(1)を四回用いると $C=\operatorname{diag}(l,m)C\operatorname{diag}(m,l)$ となる。ただし
\[
 l=(\beta-1)/(\alpha-1),\qquad m=\alpha l/\beta .
\]
可逆 $C$ があるには $lm=1$ または $l^2=m^2=1$ が必要である。
前者は $(\beta-\alpha)(\alpha\beta-1)=0$、後者でも $l=1$ または $m=1$ なら $\beta=\alpha$、
両方 $-1$ は不可能である。従って $\beta=\alpha$ で任意の $C$、または $\beta=1/\alpha$ で
対角 $C$ だけが残る（重複時は前者）。後者では共通倍率および両側からの同じ対角合同で $C=I$ にできる。

重根型の中央基底の残余は、位置別のスカラーを除き
$e_0\mapsto e_0,e_1\mapsto k e_1+m_i e_0$, $\sum m_i=0$ である。
$P$ の線を保つには $k$ が共通であり、これらは全てその線を保つ。
各列の第二座標は可逆性から非零で、極倍率を使い第一・第二鉛筆で1に揃える。
第一座標は $m_2+m_3,m_3+m_0$ の平行移動と共通の $k^{-1}$ 倍を使って
\[
 D_0=\begin{pmatrix}0&-1\\1&1\end{pmatrix},\qquad
 D_1=\begin{pmatrix}0&-\beta\\1&1\end{pmatrix},\qquad \beta\ne0
\]
にできる。$E=E_{01}$ とすれば一巡条件は
$C=\beta^2(I+(\beta-1)E) C(I-(\beta-1)E^{\mathsf t})$。
右の単冪部分の固有値は1なので、$\beta=1$ で任意の $C$ か、
$\beta=-1$ で $EC=CE^{\mathsf t}$ となる。後者は
$C=\left(\begin{smallmatrix}u&v\\v&0\end{smallmatrix}\right)$, $v\ne0$ で、
残余平行移動 $m_0=m_2=-m_1=-m_3$ と共通倍率で交換行列にできる。
\end{proof}

\begin{proof}[二つの二元三次式と同型]
辺の最低文字を双方1にする。一方だけ非零なら混合グラフと $xz$ の式から他方の全座標が零になり、
最高文字でも同じである。中央二座標を $u,v$ とすると重み1より
$v=-C_i^{-\mathsf t}A_{i+1}u$。重み2で両最高座標は $-u^{\mathsf t}B_{i+2}v,-u^{\mathsf t}A_{i+2}v$ となる。
重み3の二関係は従って次の二元三次式の同時消滅である。
\[
 (u^{\mathsf t}B_{i+2}v)C_{i+2}v+(u^{\mathsf t}A_{i+2}v)B_{i+1}u=0.
\]
非零共通根があれば非自明な軌道がある。
二点型 $\beta=\alpha$ と重根型 $\beta=1$ は一般行列 $C$ の関数体上で始点0の二式に共通二次因子がある。
重根型 $\beta=-1$ は共通一次因子がある。終結式が恒等零なので可逆な全特殊値でも有限性は成立しない。
残る二点相反型では、四位置と極反転後の計八終結式は全て $\alpha+1$ の二乗に
$\alpha,\alpha-1$ の単元を掛けたものになる。中央だけの辺は二点型鉛筆の二つの単純点のみである。
正確な二次行列反復・共通因子・Sylvester行列式は \path{computations/verify_double_middle_torus.py} で検算できる。
これで定理のパラメータ領域と復元式が得られる。逆向きの正則性は命題\ref{prop:double-middle-finite-pbw}による。
$\alpha=-1$ でもその証明の可逆性と三文字の基底は保たれる。

二点型の二つのSegre分解項は順序を除き一意であり、各位置で交換するなら四位置全てで同時に交換する。
列倍率と行倍率を除く $D_0$ の交差比が $\alpha$ で、同時交換は $\alpha^{-1}$ を与える。
極文字の交換は $A_i,B_i$ を交換するので同じ反転しか加えない。相反型の対角 $C$ は常に一類であるから十分性も従う。
有限グラフでは極大トーラスの共役によりこれが印なし同型であり、一巡回は $\alpha$ を $\beta=\alpha^{-1}$ に移すため周期1となる。

最後の標準形の支持は二つの交互極語、中央のみで全て同じ文字の二語、そして任意の異なる二位置に $x,z$ を置き
残る二位置で同じ中央文字を使う24語である。それぞれの群での一様平均を合わせると全位置四文字の頻度を均一にでき、
一位置だけの相違はない。従って前節と同じモーメント判定で多重安定となる。
両節の閉軌道の射影安定群は還元的で階数1である。もし半単純因子があれば、閉点の有限集合を
その連結群が各点ごとに固定する。他方その根トーラスの最低点は単純な最低文字で、
非自明な $\mathrm{SL}_2$ 表現の極重み直線であるから群全体には固定されない。これで恒等成分の主張も従う。
\end{proof}

% !TeX root = ../papers/nc-p3.tex
\section{有限トーラス窓に特殊化する対称性を指定しない全窓}
\label{sec:finite-torus-slices}
定理\ref{thm:simple-weight-finite-windows}の曲面で $(b,c)=(2,3)$、曲線で $b=2$、
定理\ref{thm:double-middle-finite-windows}で $\alpha=2$ の閉軌道を $A,B,C$ とする。
ここで $b,c$ は六単元正規形の座標である。
\begin{theorem}\label{thm:finite-torus-slices}
二隣接切断の階数6軌跡は三点で滑らかで、アフィン次元は $76,75,75$、各点を通る既約成分は一意である。
$A$ の成分と $B$ の成分は異なり、両者とも $(u,v)=(1/2,1/2)$ の軌道を含む。
一般の添字保持類の次元はそれぞれ $15,14,14$ で、各幾何学的一般点は標準正則である。

各閉軌道に特殊化する全窓は有限点グラフを持ち、トーラスの存在を仮定せず全て正則である。
閉軌道自身を除くそれらの添字保持類は、集合として次の通りである。
\begin{enumerate}
\item $A$ では $\mathbb P(1,2,3,3,4,4,5)$ 一つで、全て最小周期4である。
\item $B$ では同じ重み付き射影空間の二つの互いに素なコピーである。
\item $C$ では $\mathbb P(1,1,2,2,2,3,3)$ の対合商である。重み $1,2,3$ ごとに
符号を変える座標数はそれぞれ一つで、他は固定する。
\end{enumerate}
同じ計算は各表示族の稠密開部分で成立する。
\end{theorem}
\begin{proof}[非被約点での余法線と接空間]
補題\ref{lem:finite-flat-conormals}の証明では四文字積の評価の階数が $p$ のとき、
有限平坦な点環上で作る恒等式の係数行列から $p$ 本の独立行を取れる。
従って第一切断の滑らかな156次元軌跡上で第二切断のSchur方程式は局所的に $100-p$ 本以下で生成される。
$p=19$ の場合も同じ単元小行列の消去でよい。合わせた微分の階数が $200-p$ なら滑らかである。

極小点のアフィンチャートの局所環は前二節の消去式で、単重みなら重み $1,2$ の二変数上の重み $4,5$ の二式、
中央二重なら二変数上の二本の三次式である。最高点にも反転した式を用いる。
四つの一文字の積を、順に位置 $3,0,1,2$ の座標を使い、これらの有限環で剰余にする。
$C$ では中間二つの単純点の評価も加える。256列の階数は順に $20,19,19$ となる。
一方、各切断の左右核のベクトル $l,r$ に対する $l^{\mathsf t}(\delta w)r$ の全てで一次条件を作り、
位置別 $\mathfrak{gl}_4$ の軌道方向を除く。最上重みを引き、
正負の同じ絶対重みをそれぞれ一回ずつ数えたときの各側の多重度は次表になる。
\[
\begin{array}{c|c|ccccc}
 &0&\pm1&\pm2&\pm3&\pm4&\pm5\\\hline
 A&2&1&1&2&2&1\\
 B&1&1&1&2&2&1\\
 C&1&2&3&2&0&0
\end{array}
\]
これ以外の重みは零である。アフィン軌道は各60次元なので上記の微分階数が得られる。
正確な剰余・階数は \path{computations/verify_finite_torus_slices.py} による。
各族は既約なので、その全体が標本での一意な成分に入る。稠密開部分では滑らかさが保たれ、
トーラス同変接束の重み階数は一定である。
二族の共通パラメータが両成分に入るのは先の正規形による。
一般点の正則性は命題\ref{prop:unrestricted-good-reduction}の次数ごとの持上げによる。

小次元側の二つの表示は両対称化もできる。位置0,1にそれぞれ作用させて反対相手との交換を対称にする
対角行列は、単重み曲線では $\operatorname{diag}(b,1,1,1),I$、中央二重では
$\operatorname{diag}(1,1,\alpha,1),\operatorname{diag}(-1/\alpha,1-1/\alpha,-1+1/\alpha,1)$ である。
これは全語の係数比較で直接得られる。
\end{proof}

\begin{proof}[閉軌道の全ファイバー]
射影基底群 $G$ の半安定軌跡で各閉点のアフィン不変近傍を使う。
$H$ を閉軌道の安定群とする。標数零の滑らかな点でのLunaのエタールスライス定理により、
軌道に垂直な上表の $H$ 表現 $N$ とスライスの間に強エタールな局所模型がある。
商で閉点の剰余体は $k$ なので、商上のその一点のファイバーを取ると、
実際の $G$ ファイバーは $G\times^H (N\mathbin{\times}_{N/\!/H}\{0\})$ と同型である。
すなわち線形近似だけでなく、この一点に特殊化する全軌道の同型類を与える。
ランク6と有限辺は閉軌道を含む不変開条件であり全ファイバーで保たれる。
正則性は閉軌道への付値特殊化と命題\ref{prop:unrestricted-good-reduction}から従う。

$H^\circ=T$ が一次元トーラスなので、線形零ファイバーの台は正・負の重み空間の和集合である。
零重みは零でなければならず、正負がともに非零なら次数を揃えた積の不変式がある。
有限拡大 $H/T$ を加えても零ファイバーの台は変わらない。
$A,B$ で向きを保つ $H$ は $T$ である。実際四固有線は単純で、先の12主元の指標格子は飽和している。
$A$ は反転があり正負を交換し、$B$ では $b\ne1-b$ なので反転がない。
$C$ で向きを保ち二つの中間線を交換するには $\alpha=\alpha^{-1}$ が必要で、これは許されない。
従って対角正規化の残余だけで、$C_0=I$ を保つ条件が比の二乗を1にする。
$T$ のほかに一方の中間文字だけ四位置で符号反転する対合が生じる。
重みを反転する同型は極の交換と中間二線の交換を組み合わせて得られるため正負は一つになる。
符号を変える接方向の多重度は上記コードでテンソル単位をその文字数の偶奇にも分けて得る。
非零方向をトーラスで割れば記載した重み付き射影空間となる。
$A$ の場合に一添字または二添字の移動で一般軌道が自分に戻れば、その閉極限も同じ移動で自分に戻る。
この閉極限は周期4なのでファイバー内のどの軌道も周期4である。
いずれの法表現にも非零重みがあり、零でない一般方向の安定群は有限であるから
一般軌道のアフィン次元は61、従って表示した一般同型類次元となる。
\end{proof}

% !TeX root = ../papers/nc-p3.tex
\section{周期一の共通窓における全成分の局所方程式}
\label{sec:torus-component-crossing}
\begin{theorem}\label{thm:torus-component-crossing}
単重み曲面と単重み曲線の共通点 $w_*$、すなわち $(b,c)=(1/2,1)$ において、
トーラスを仮定しない二切断階数6軌跡の軌道横断完備局所環は
\[
 k[[u,v,z,t_1,\ldots,t_{14}]]/(uv,uz)
\]
である。$u,v,z$ は重み0、残りは正負各側で $1,2,3,3,4,4,5$ の絶対重みを持つ。
アフィンテンソルの完備局所環にはさらに60個の滑らかな軌道変数を加える。
従ってこの点を通る全体成分はちょうど二つで、双方ともここで滑らか、交叉も滑らかである。
二成分のアフィン次元は76と75、共通部分は74であり、交叉のスキームに重複はない。
この閉軌道に特殊化する窓は、閉軌道自身のほか、一つの $\mathbb P(1,2,3,3,4,4,5)$ で分類される。
これらは全て有限点グラフの標準正則窓で両成分の共通部分にある。
\end{theorem}
\begin{proof}[重み0におけるスキーム構造]
前節の三十八単元の規格化を使う。重み0の基底変換に横断な十二主元を1にし、
重み3ブロックのSchur方程式から $d,h,f,x,y$ を解いた残りの三つの単元 $a,b,c$ を取る。
全ての消去の分母は $w_*$ で単元である。新しいエタール座標を
\[
 u=ab-c,\qquad v=a(1-b)-c,\qquad z=c-1
\]
とする。逆に $c=1+z$, $a=u+v+2c$, $b=(u+c)/a$。
重み0のイデアルは $I_0=(uv,uz)$ である。実際一式は単元倍の $uv$ となり、
別の式の $v=0$ での制限は $uz$ の単元倍となる。全式は $u=0$ 上と $v=z=0$ 上でそれぞれ恒等零である。
よって包含 $I_0\subset (u)\cap(v,z)=(uv,uz)$ を得る。
逆向きには後者の式を $u$ で割って $v$ を法とした剰余が $z$ の単元倍であることと $uv\in I_0$ を使う。
\end{proof}

\begin{proof}[全横断環での共通因子]
前節と同じ一次条件で接空間の軌道商は17次元、うち重み0が3次元で正負は上表のままである。
両極の局所環上の四文字積の階数は19である。これらと上記単元の検査は
\path{computations/verify_torus_component_crossing.py} による。
有限平坦な点環の19恒等式により、軌道を除き線形独立な微分を持つ方程式を消去した後では、
17変数正則完備局所環 $L$ に残るイデアル $I$ は高々二本で生成される。
実際アフィンでは独立微分が $256-60-17=179$、局所生成元が高々 $200-19=181$ 本である。
トーラスは線形還元的なので、この消去では独立な微分の同変な持上げを選べる。
固定部分の軌道は対角基底変更の軌道であり、その完備横断環は上の十二主元による規格化と同型である。
従って $L$ の非零重み変数を零にしたときの $I$ の像は $I_0$ である。
$I_0$ に二生成元が必要なので、$I/\mathfrak m_L I$ は二次元で $I_0/\mathfrak m I_0$ へ同型に写る。
特に二つの生成元 $f_1,f_2$ を重み0に選べる。

前節の76次元成分がこの点を通るため、$I$ は $L$ の高さ1の素イデアルに入る。
正則局所環は一意分解環なので $f_1,f_2$ に非単元共通因子 $F$ がある。
重み0の三変数へのその制限は、零でも単元でもなく、$uv,uz$ の最大公約元 $u$ を割る。
従って制限は $u$ の単元倍であり、$F$ は滑らかな助変数である。
別の共通素因子もその二乗もあり得ない。この一意な主イデアルはトーラス安定であり、
$F$ の重み0部分も同じイデアルを生成するので $F$ を不変にできる。
$f_i=F g_i$ と置けば $g_i$ の固定部分での制限は $(v,z)$ を生成する。
従って $F,g_1,g_2$ は独立な一次部を持つ三つの重み0の変数に取れ、残りの変数を
同変に補えば主張した完備局所方程式になる。既に相異なる二つの実際の成分がここを通るため、
二つの完備既約枝はその二成分に一つずつ対応する。
\end{proof}

\begin{proof}[特異点の零ファイバー]
スライスを滑らかな線形表現に局所閉埋込みし、独立な微分を持つ方程式の持上げを
安定群同変に取ると、その零点はスライスを含む滑らかな17次元の周囲を与える。
ここで安定群はトーラスとそれを反転する有限成分からなり線形還元的である。
周囲のスライスの零ファイバーは前節の強エタール性から正・負の各七次元空間の和集合になる。
この全体が実際のテンソルスライスに入る。実際、残るイデアルの極小生成空間は上で示した通りトーラス重み0だけなので、
アフィン近傍でトーラス不変な生成元に持ち上げられる。これらは原点で消え、原点を極限に持つ全点で消える。
同じ議論を各枝のイデアルに用いることもできる。上の完備表示で両枝の極小生成元は重み0であり、
完備化前の各枝の極小生成空間も同じである。従って両枝とも零ファイバー全体を含む。
向きを保つ安定群は対角トーラスで、重み反転が正負を交換するため、一つの重み付き射影空間を得る。
\end{proof}

% !TeX root = ../papers/nc-p3.tex
\section{吸引ファイバーの全周期と重み付き座標での判別}
\label{sec:torus-fiber-periods}
\begin{theorem}\label{thm:torus-fiber-periods}
定理\ref{thm:finite-torus-slices}および定理\ref{thm:torus-component-crossing}の
非閉軌道の全てについて、最小周期は次の通りである。
\begin{enumerate}
\item 単重み曲面の一般点 $A$ に落ちるものは全て4、単重み曲線の一般点 $B$ に落ちるものは全て2である。
\item 共通点 $w_*$ の重み付き射影空間上で、重み3と4の各二座標のうち一つだけを符号反転し、
残りを固定する対合を $F$ とする。$F$ の不動点がちょうど周期1、それ以外は2である。
\item 中央二重型の一般点 $C$ に落ちる空間の重み $1,3$ の各二座標を、符号対合 $D$ の
正負の固有順に取る。$F$ はこの二座標を交換する。重み2の三座標では
$D=\operatorname{diag}(1,1,-1)$, $F=\operatorname{diag}(1,-1,1)$ とする。
$F$ が $D$ 商上で固定する点がちょうど周期1、それ以外は2である。
\end{enumerate}
具体的に重み $j$ のベクトルを $v_j$ とすると、最後の場合はある $t\in k^\times$ と
$e\in\{0,1\}$ が全 $j$ で $Fv_j=t^j D^e v_j$ を満たすこと、二番目なら同じ式で $e=0$ とすることが必要十分である。
従って座標が消える部分や二つの符号成分が同時に現れる場合も含む。
\end{theorem}
\begin{proof}
位置巡回を基底群に加えた有限拡大は還元的である。閉軌道を保つ巡回と基底変更の組にも同変に
前二節のスライスを取れる。特異な共通点でも滑らかな周囲のスライスを同変に取るので、
零ファイバー上は実際に正負の線形重み空間の作用になる。群の有限拡大によって
元の同型関係を変えることはせず、巡回後に元の基底群軌道へ戻るかを計算する。
$A$ は既に閉極限の周期で決まる。

$B$ の一移動は重みを反転する同型でのみ閉点に戻せ、正負の二コピーを交換する。
二移動は正負を保ち、位置0に $\operatorname{diag}(b,1,1,1)$ を作用させた表示では
単に二つの反対位置を同時交換する位数2の作用である。
$b=2$ でその正側の軌道法表現への作用を計算すると全て恒等となる。負側も一移動による共役で同じである。

共通点では向きを保つ一移動がある。四テンソルの座標を一つ巡回し、新位置2の数字0を $1/2$ 倍、
新位置3の数字0を2倍すれば固定される（他は恒等）。
正側での各重みの作用は全て二乗が恒等で、負固有空間の次元が重み順に $0,0,1,1,0$ である。
従って二移動は全点で恒等、一移動の不動条件は記載したものとなる。

$C$ では巡回と二中間文字の交換を組み合わせれば向きを保つ。
$\alpha=2$ の基底で、この単項式変換に対し四位置に必要な対角倍率を作用させた写像を $L$ とすると、
重み $j$ 上で $L^2=(\sqrt{2}/4)^j I$, $LD=(-1)^j DL$ となる。
トーラスの元を掛けて全重みで $F^2=1$ にできる。奇数重みの二座標は
$D$ の固有空間が一つずつで互いに交換されるので、基底を選んで表示通りとなる。
偶数重みでは $D$ の正側で $F$ は正負一つずつ、負側で正一つとなる。
持上げに使ったトーラス倍率と $D$ は既に同型関係に含まれる。

\path{computations/verify_torus_fiber_periods.py} は元の窓上で各変換を検算し、
二切断の微分の核から軌道方向を除いた各空間でこの作用を計算する。
各族の稠密開では既に重み別法束が局所自由であり、シフトの正規化に必要な根をエタールに添加すれば
有限位数作用の固有束の階数は一定である。実際、中央二重型で巡回の四乗は向きと中間二線を保ち、
二移動自体は両対称表示の二交換による位数2の持上げを持つ。
より直接には一移動の四乗はトーラスと $D$ の積であり、トーラスの四乗根で除いた後の位数は高々8である。
これで一般点にも標本と同じ零ファイバー判定が移る。
\end{proof}

% !TeX root = ../papers/nc-p3.tex
\section{有限点対応を保つ加法群と四つの推移変換}
\label{sec:finite-additive-windows}
標数零の代数閉体上、四辺の点対応が有限なFrobenius窓を取る。
有限完全交叉の射影は閉埋込みで、二つの隣接対応は共有位置で同じ長さ20の像を持つ
（定理\ref{thm:finite-frobenius-low-opposite-as}の証明）。
各辺の関係六平面に階数1のテンソルはなく、どの端点での評価階数も3以上である。
各位置の三次式の制限は全切断への同型である（定理\ref{thm:finite-square-chart}の完全交叉計算）。
ここでは窓の射影安定群内の非自明な加法群を考え、四つの文字空間上の冪零生成元を $N_i$ とする。

\begin{proposition}\label{prop:finite-additive-jordan}
各 $N_i$ のJordan型は $(4),(3,1),(2,1,1)$ のいずれかである。
各位置の最大ブロックの末端を検出する点、すなわち
$[p_i]=\mathbb P(\operatorname{im}(N_i^*)^{a_i-1})$（$a_i$ は最大ブロック長）は
四周期点列を与える。型 $(4)$ と型 $(2,1,1)$ は一窓に共存しない。
\end{proposition}
\begin{proof}
一位置が自明ならグラフ性により隣の作用も点スキーム上恒等である。
射影変換 $g$ の固定方程式 $v\wedge gv$ は二次式なのでその位置の全射影空間でも零になる。
従って四位置すべてが自明となる。隣接文字を $A,B$、最大長を $a,b$ とし、
$D=N_A\otimes1+1\otimes N_B$、辺の関係を $R$ とする。
$D^{a+b-2}$ は $N_A^{a-1}\otimes N_B^{b-1}$ の非零倍である。
$A$ の最大ブロックが一個ならこの作用の各像は純関係なので零である。
従って $R(p_A)\subset\ker N_B^{b-1}$、評価階数から $B$ の最大ブロックも一個で等号である。
これが $p_B$ を一意な後続点とする。最大ブロックが二個なら全四位置が $(2,2)$ である。
この場合 $D^2 R$ は二文字の像二平面の積に入り、純テンソルを含まないため次元は高々1である
（射影直線は二次曲面と交わる）。両商 $A/\operatorname{im}N_A,B/\operatorname{im}N_B$ への
関係の射影は両 $N$ による同型を通じて同じ次元を持つ。二つの固定直線の積に高々一式しかなく有限性に反する。

型 $(4)$ があれば全支持はその固定点 $p$ 一個にある。その位置の射影余接空間で誘導される作用は
長さ3の一ブロックである。型 $(2,1,1)$ の正準点の余接空間では作用は零である。
同じ局所環の余接空間は両方の商なので両型があれば次元は1であり、環は $C=k[t]/(t^{20})$ となる。
実は型 $(4)$ の単独支持でこの環は不可能である。$p(z)=1$ となる長い鎖の始文字を取り
$\phi=Nz/z$、$C$ に誘導する微分を $f(t)\partial_t$ と書く。
$\ker(N|_{\Sym^2 V})$ の次元は2である（四次元既約表現の対称平方は次元7と3の和）。
どの不変二次式 $q$ も $z^2$ の係数が零である。$N^6$ でその係数だけを検出できるからである。
非零な $Q=q/z^2$ は従って消失位数 $s\ge1$ を持ち、$\phi Q\ne0$ は三次式の単射性による。
微分方程式 $fQ'=-2\phi Q$ の右辺が零でないので先頭項を比較できる。
$e=\operatorname{ord}\phi$、$d$ をその先頭係数とすると $\operatorname{ord} f=e+1$ で、
$f$ の先頭係数 $c$ は $sc=-2d$ を満たす。すべての非零不変二次式に同じ $s$ を強制するが、
二つの独立な切断の先頭係数を一次結合で消せるため矛盾する。
\end{proof}

\begin{theorem}\label{thm:finite-four-transvections}
四つの $N_i$ がすべて型 $(2,1,1)$ なら、有限点対応を持つFrobenius窓は
標準quadratic Koszul左右AS正則である。この窓は半安定で、唯一の多重安定極限は
定理\ref{thm:one-standard-sl2-classification}の交代型または対称型の $P=0$ である。
加法群の固定平面どうしで切った各辺は長さ6で、元の辺の全ての閉点を台に含む。
\end{theorem}
\begin{proof}
各文字空間を $kz\oplus U\oplus kx$、$Nz=x$、$\dim U=2$ と書く。
重みを順に $-1,0,1$ とし、$z,x$ を標準表現にして四位置同時の $\mathrm{SL}_2$ に拡張する。
辺のテンソル積で重み $h$ 以上を $F^h$ と書く。$D^2R=0$ から $zz$ の係数は常に零である。
$R\cap F^1$ の重み1の成分を $xU$ と $Ux$ のそれぞれへ射影するといずれも単射である。
片側と $xx$ だけなら純テンソルになるからである。
一方 $DR\subset R\cap F^1$ の両射影は全射である。実際、$zU$ と $Uz$ の係数を読む写像は
正準点での評価の像 $U\oplus kx$ によりそれぞれ全射で、$D$ はこれらをそのまま重み1へ送る。
従って $R\cap F^1=DR$ は二次元、$R\cap F^0=\ker(D|_R)$ は四次元である。
後者では $D$ が重み2以上に入って消え、逆に重み $-1$ 上で $D$ は単射である。
残る重み0の二本の初期項は $UU\oplus k(xz-zx)$ にある。

重み $-1$ の二つの代表とその $D$ 像、重み0の不変な二代表を基底にすると、
抽象的な $R$ 自身に標準表現二個と自明表現二個の構造を与えられる。
この構造での辺の埋込みは重みを下げず、次数零の部分も埋込みである。
反対辺間の窓の完全対合 $R_i^*\simeq R_{i+2}$ と逆写像は加法群同変である。
二つの $\mathrm{SL}_2$ 表現の間の加法群同変写像は $\operatorname{Hom}$ の最高重みベクトルの和なので
重みを下げず、次数零の部分は $\mathrm{SL}_2$ 同変である。逆写像との積により次数零の部分は同型である。
従って四位置の対角一径数群による窓の極限 $w_0$ は存在し、
両方の隣接切断で六次元の初期関係空間の間の完全対合を与える。
この極限は $\mathrm{SL}_2$ 不変で、標準型の重複度空間から両座標への射影は上記の二つの全射である。
したがって標準型切断の四つの $b_i$ は可逆である（同定理の射影による分解）。
各文字の支持も全体である。重み $-1$ と $1$ の混合関係の両射影から $U,z,x$ を全て検出できる。
同定理で $w_0$ は正則、命題\ref{prop:unrestricted-good-reduction}で元の窓も正則である。

二つの固定平面 $x=0$ 上で $DR$ の二式が消える。残り四式は
$\mathbb P^2\times\mathbb P^2$ 上で有限、かつ正準点対を含むので、独立な完全交叉で長さは
$\binom42=6$ である。連結群は有限支持の各点を固定する。
重み0の不変な二代表の制限はその $UU$ 成分のみで、これらも独立である。
$w_0$ が $P=0$ なら自明表現部分の切断の $UU$ 射影は $kb_i$ だけになり、矛盾する。
従って $w_0$ は交代型または対称型の $P\ne0$ であり、系\ref{cor:semisimple-symmetry-closed-limits}
から多重安定極限が従う。この必要条件は二つの型の有限例の存在を仮定しない。
\end{proof}

% !TeX root = ../papers/nc-p3.tex
\section{正次元安定群を持つ有限候補の全次数正則性}
\label{sec:finite-connected-regularity}
\begin{theorem}\label{thm:finite-connected-regularity}
四辺の点対応が有限な標準有限Frobenius窓で、射影安定群が正次元なら、
全次数で標準quadratic Koszul左右AS正則であり、群作用を忘れた窓は半安定である。
加法群が作用すればそのJordan型は四位置で等しく、次のいずれかである。
\[
\begin{array}{c|c|l}
V_i\text{ の型}&R_i\text{ の型}&\text{多重安定極限}\\\hline
(4)&(5,1)&\text{多項式環}\\
(3,1)&(3,3)&\text{多項式環}\\
(2,1,1)&(2,2,1,1)&\text{多項式環または対称型の }P=0
\end{array}
\]
これは有限候補を入力とする判定で、各行の有限例の存在を必要としない。
乗法群の場合は定理\ref{thm:simple-weight-finite-windows}と
\ref{thm:double-middle-finite-windows}の既出の二曲線と一曲面に尽きる。
\end{theorem}
\begin{proof}[異なるJordan型の辺の排除]
前節の $D$ を使い、各Jordanブロックを $L_d$ に延長してその対角重みを使う。
各位置の最大ブロックの最低文字を $z$、最高文字を $x$ と書く。
$zz$ 項は必ず零である。$(3,1)$ と $(2,1,1)$ の順に隣接すると、
最低となり得る重み $-2$ の成分は $z\otimes U$（$U$ は第二位置の二自明文字）
だけである。$D^2r$ はそこから来る $x\otimes U$ と重み3の $kx\otimes x$ だけに入る。
純関係がないので $D^2 R=0$ となり、$z$ での評価が $U$ を含めない。
逆順も同じである。$(4)$ と $(3,1)$ の辺では最低 $-5$ が消え、
重み $-3$ の成分は $L_6$ の次最低と二個の $L_4$ の最低の直和にある。
$D^4R$ は最高 $kx\otimes x$ に入り零なので $L_6$ の成分はない。
$z$ での評価は第二位置の二つの重み0をともに含むから、独立な二つの最低成分が必要になる。
$D^j$, $0\le j<4$ によるそれらの像は異なる重み層ごとに独立なので $\dim R\ge8$ となって矛盾する。
$(4)$ と $(2,1,1)$ は命題\ref{prop:finite-additive-jordan}によっても排除済みである。

\end{proof}
\begin{proof}[同じJordan型の初期関係]
型 $(4)$ では $V_i\otimes V_{i+1}=L_7+L_5+L_3+L_1$ で、最低 $-6$ は消える。
$D^5 R\subset kx\otimes x$ は零なので、重み $-4$ の成分は $L_5$ の最低文字だけである。
正準点で次の文字への評価が必要なためこの成分は現れる。
従って $R$ は長さ5の鎖と長さ1の鎖である。四位置の同じ対角一径数群で窓 $w$ を還元する。
この極限は常に存在する。実際、加法群同変な各辺の埋込み
$L_5+L_1\to V_i\otimes V_{i+1}$ は重みを下げず、次数零写像は半単純群同変である
（前節の最高重みの議論）。ここで長さ5の始元は非零の重み $-4$ 成分を持つように取る。
相補辺間の抽象的な完全対合の次数零部分は同型で、特に重複度1の $L_5$ 間で非零である。
従って極限 $w_0$ は非零である。その全辺の関係像は $L_5+L_1=\bigwedge^2 L_4$ に入る。
これは四つの $L_4$ を同一視したとき各隣接交換で交代的ということで、$w_0$ は交代体積となる。
両切断で階数6が回復し、命題\ref{prop:unrestricted-good-reduction}が使える。

型 $(3,1)$ では重み $-4$ が消え、$D^3 R$ は最高の純直線に入り零である。
重み $-2$ 部分で $L_5$ の次最低が消え、残るのは三つの $L_3$ の最低である。
$z$ で評価すると相手の重み0の二文字へ全射なので独立な最低二成分がある。
これらとその一次、二次の $D$ 像だけで六次元を使い尽くす。
従って $R=2L_3$ と読んだ埋込みの次数零部分は単射で、両切断で $w_0$ は階数6を保つ。
その重複度二平面は、三コピー $t_i X,Xt_{i+1},X\times X$ のうち、
最後と最初又は最後と二番目を取る両射影で同型になる。
これは元の辺の正準点で左右それぞれ、$t$ と中間の鎖文字の両方への評価が全射であるためである。
特にこの二平面には最初又は二番目だけのベクトルはなく、各文字の支持も全体である。
定理\ref{thm:sl2-window-profiles}の $L_3+L_1$ の証明で、階数2かつ全一文字支持でも
非正則になり得るのは $uv\ne0$ である $a_i$ が消える場合に限られる。このとき隣辺に純混合コピー
$t_i X$ 又は $X t_i$ が入るが、上の二射影がこれを排除する。
従って $w_0$ は多項式環または $U(\mathfrak{sl}_2)$ の斉次化である。
良い還元と系\ref{cor:semisimple-symmetry-closed-limits}で正則性と閉極限を得る。
型 $(2,1,1)$ は定理\ref{thm:finite-four-transvections}である。
\end{proof}

\begin{proof}[正則性を仮定しない乗法群の重み]
定理\ref{thm:finite-torus-profiles}の露出面、端点評価、辺ごとの局所指標の証明は
有限完全交叉と純関係の排除だけによる。そこで正則窓のEuler文字を用いた箇所を次の二点で置き換える。
第一に、最低が三重なら隣接位置でも三重であり、最低の $3\times3$ 部分の関係は
純テンソルを避ける上界4と有限交叉の下界4から四次元である。
全六関係中、最高関係重みの重複度は高々2となる。反対辺と窓文字への完全対合では
最低と最高の関係重複度が対応するため不可能である。最高が三重の場合も反転すればよい。

第二に、両極単純で中間も単純なとき、辺の指標の列挙は中央差を $1,2$ に制限し、
差2の隣接を禁じる。異なる差の辺では関係重み $1,\ldots,6$ が各一本で、両位置の
両極どうし以外の組は辺の点ではない。一方、差1どうしの辺の中央関係重複度は2である。
よって混合辺が一つでもあれば、その反対辺も混合辺である。四位置で差2は交互に並ぶしかない。
このとき第一切断で「最高、最低、最低、最高」の係数を見る。行の位置対も列の位置対も
グラフにないから、それぞれの関係線のその座標は非零である。
その二重みは補い合う。実際全体重みは $1+6=7$ で、二つの最高文字は重み3と4である。
従ってこの係数は階数1の完全対合により非零となる。他方、次の切断では中央の最低どうしに
関係が一つもなく、この係数は零である。矛盾する。
中央二重の議論は同定理の有限交叉のままである。残る二型が満たす係数条件とPBW性は
前掲の二分類定理の切断計算による。特に中央二重の両極の交互語係数は、二つの極の包の積に
中央重みの関係が二本独立に制限されることと、同じ極の連続を含む語の係数が零であることから非零である。
最後に標数零の正次元線形代数群は乗法群又は加法群を含むので、以上ですべてを覆う。
\end{proof}

% !TeX root = ../papers/nc-p3.tex
\section{有限点グラフに加法群が作用しないこと}
\label{sec:finite-unipotent-exclusion}
\begin{lemma}\label{lem:maximal-tangent-no-additive}
三次元接点を持つ標準正則有限窓の射影安定群は加法群を含まない。
\end{lemma}
\begin{proof}
定理\ref{thm:maximal-tangent-complete}の一意な三次元接点の各位置を $p_i$ とする。
加法群はこの点列を保ち、$H_i=\ker p_i$ も保つ。
定理\ref{thm:maximal-tangent-ore}の内在的な三文字関係 $S_i=R_i\cap(H_i\otimes H_{i+1})$ と
その点対応の射影像も保存される。定理\ref{thm:reduced-ternary-graphs}により
各 $\mathbb P(H_i^*)$ でこの像は滑らかな三次曲線または節点だけを特異点とする被約三次曲線である。
どの許容曲線にも非自明な射影加法群は作用しない。
滑らかな種数1曲線には射影直線からの非定数射がなく、加法群の軌道写像も固有性で射影直線へ延びる。
節点有理三次曲線では正規化上の二点が固定され、横断する直線と二次曲線、
非共点三直線でも各成分上で少なくとも二つの交点が固定される。
射影直線の二点を固定する乗法群には加法群からの非自明な準同型がない。
これらの平面曲線は射影標構を含むため曲線上の恒等作用は平面上でも恒等である。

したがって冪零生成元は $H_i$ を零にし、残る文字 $Z_i$ をある $n_i\in H_i$ に送る。
同定理による各位置の $H_i$ の移送を行い、混合関係を
$Z_i h-h Z_{i+1}-d_i(h)$ と書く。微分すると
$n_i\otimes h-h\otimes n_{i+1}\in S_i$ が全 $h$ について必要である。
三文字の点グラフの $(p,\sigma p)$ で評価すれば
$n_i(p)\sigma p=n_{i+1}(\sigma p)p$ となる（代表の倍率は両辺で同じ）。
補題\ref{lem:reduced-ternary-fourth-shift}により全既約成分の一般点で $p$ と $\sigma p$ は独立である。
ゆえに $n_i$ は被約三次曲線全体で消える一次式であり零である。これで全位置の作用が消える。
\end{proof}

\begin{theorem}\label{thm:finite-no-additive}
四辺の点グラフが有限な標準有限Frobenius窓の射影安定群は加法群を含まない。
従ってその恒等成分は自明または一次元トーラスである。後者の添字保持類は
定理\ref{thm:simple-weight-finite-windows}と\ref{thm:double-middle-finite-windows}の
一曲面と二曲線で尽き、正則性・周期・重複同型も両定理の通りである。
\end{theorem}
\begin{proof}
加法群があれば定理\ref{thm:finite-connected-regularity}により既に全次数正則である。
型 $(2,1,1)$ の正準点対の六方程式を一次化する。二つの負重みの鎖の始元からは高々二本、
その $D$ 像からは零本である。実際 $xU,Ux,xx$ は両点で一次部を持たない。
重み0以上の不変二元は、重み0では $UU\oplus k(xz-zx)$ に入り、
それより高い項にも $z$ がないので、この二元の一次部は合わせて高々一本である。
六次元の周囲でJacobi階数が3以下となり、グラフの射影は閉埋込みなので接次元はちょうど3である。

型 $(3,1)$ では正準点で非零の文字 $z$、同じ鎖の $v=Nz,x=Nv$、自明文字 $t$ を使う。
二つの鎖始元 $r_1,r_2$ の一次部は高々二本である。重み $-2$ の初期成分は
$zv-vz,zt,tz$ の線型結合で、$D$ を掛けた重み0の部分は $zx-xz,vt,tv$ の線型結合となる。
よって $Dr_1,Dr_2$ の一次部は合わせて高々一本である。重み0より大きい項は $z$ を含めない。
$D^2 r_1,D^2 r_2$ はそのような項だけである。ここでも接次元3となる。
両型とも前補題に反する。

型 $(4)$ では同定理の対角還元により六関係の初期空間は通常の交代二形式全体である。
従って文字を重み $0,1,2,3$ の順に取ると、各交換子 $ij-ji$ ($i<j$) に
重みが $i+j$ より真に大きい項を足した六式が辺の基底になる。
両側で文字 $2,3$ を零にする。残る二直線の積に制限されるのは $01-10$ の持上げだけである。
他の交換子は初期部が消え、補正の重みは2より大きい。従ってこの積に曲線が残って有限性に反する。
最後に正次元安定群のトーラスの場合は前定理を適用する。
\end{proof}

\begin{proposition}\label{prop:finite-polynomial-limit-flag}
対称性を仮定せず、有限点グラフの窓が一径数群によって多項式環の四テンソルに還元されるとする。
窓と極限が異なる軌道なら、四位置を交代体積で同一視した共通の矢重みは
スカラーと正倍率を除き $0<d<e<f$（四重みは $0,d,e,f$）であり、
\[
                  e<2d,\qquad d+e<f<2e
\]
を満たす。特に重複した重みや等差四重みでは、この閉軌道の有限グラフを得ない。
\end{proposition}
\begin{proof}
極限を固定する四位置の射影作用は同じである。実際各辺は二位置の対角グラフである。
その共通の重み基底で、関係の初期空間は交換子全体となる。
各関係には対応する交換子よりも真に大きい重みの補正だけがある。
最低重みが二重以上なら、その中の対応する二文字だけに制限すると一本の交換子しか残らない。
最低を0とし重みを $0<d\le e\le f$ とする。$e\ge2d$ なら最初の二文字の積に同じ理由で
高々一式、$f\ge2e$ なら最初の三文字の積に高々三式しかなく、いずれも有限性に反する。

$f\le d+e$ と仮定する。上三文字どうしの交換子の補正は全て $d+e$ より上なので、
最低文字を含まない。この三式は最低点対で一次部を持たず、残る三式の一次部は
相手の上三文字に関して三角可逆である。よってそこは三次元接点であり、
上三文字とその三式は定理\ref{thm:maximal-tangent-ore}の内在的な三生成基底である。
良い還元によって元の窓は標準正則であるから、定理\ref{thm:maximal-tangent-complete}により
この基底の各点曲線は滑らか又は節点のみの被約平面三次曲線となる。
ところがその左評価の三行三列行列式 $F$ には重み $d+e+f$ 以下の項がない。
実際各行の最低重みは当該二文字の和で、行列式では列文字の総重み $d+e+f$ を引く。
等号の項は通常の交換子三式の行列式で、恒等零である。
上三文字を $u,v,y$ とすれば、$u^3,u^2v,u^2y,uv^2,uvy$ の重みはすべて $d+e+f$ 以下である。
従って $F$ が恒等零なら平面全体が残り、零でなければ $[1:0:0]$ は
二次接錐が零または $y^2$ の倍数となる特異点である。いずれも許容曲線でない。
以上で $f>d+e$、既得の $f<2e$ と合わせて $e>d$ も従う。
\end{proof}

% !TeX root = ../papers/nc-p3.tex
\section{有限点グラフの閉極限となる量子窓は二類だけである}
\label{sec:finite-quantum-closed-limits}
\begin{proposition}\label{prop:maximal-tangent-closed-limits}
三次元接点を持つ標準正則有限窓は半安定である。三生成基底 $B$ の点三次曲線が滑らかなら、
唯一の多重安定極限は $B[Z]$ であり、量子窓ではない。特異なら、位数2と4の二型についてそれぞれ、
中心文字 $Z$ と $xy=\zeta yx,yv=\zeta vy,vx=\zeta xv$ を持つ量子窓に退化する。
ここで $\zeta=-1$ 又は $\zeta^2=-1$ で、逆数の選択は同じ添字保持類を与える。
\end{proposition}
\begin{proof}
定理\ref{thm:maximal-tangent-ore}の移送表示で補文字を $Z$ とする。$Z$ より三つの基底文字を
高く取れば混合補正が消えて $B[Z]$ になる。この極限の関係は六本で、
四テンソルは三生成基底の巡回三体積 $w_B$ に $Z$ を交代符号で四通りに挿入したものである。
実際六関係の先頭部分はこれで、各辺の像は直和因子として極限と可換である。
三次体積の巡回性は滑らかな場合はHesse表示、特異な場合は定理\ref{thm:ternary-torsion-smoothing}の三関係表示による。
四テンソルの極限もこれらの辺の交叉の一意な直線に入る。

滑らかな場合の三体積は三文字の巡回置換と $(1,\omega,\omega^2)$ 倍で不変である（$\omega$ は原始三乗根）。
複素数上でこの二つのユニタリ行列は三文字上既約であり、各位置のGram行列の三文字部分はスカラーである。
$Z$ と三文字の成分間は、残る三位置の $Z$ の個数が異なり直交する。体積は四位置で巡回し、
全語に $Z$ がちょうど一個なので各位置でその重みのノルム平方は全体の四分の一である。
従ってGram行列は四文字全体でスカラーとなり、Kempf--Ness判定で閉軌道である。
標数零体上へは係数の有限生成定義体から移る。
この窓の点対応は $Z=0$ で滑らかな平面三次曲線を含み、その外で三文字が非零なら
混合の可換関係が平行移動の固定点を強制する。従って一次元支持は滑らかな三次曲線である。
量子窓の点支持は座標線型空間の和であるから同型にならない。

特異な三型では三生成基底の関係は巡回量子交換子に $0,1,2$ 個の純平方を加えたものである。
平方のある文字を残る文字より高く取ると、その各平方だけが初期関係で消える。
従って $B[Z]$ は記載した量子点へ落ちる。量子点は多重安定
（定理\ref{thm:quantum-face-resonance}の証明）であり、二閉類は面値で区別される。
\end{proof}

\begin{theorem}\label{thm:finite-quantum-closed-limits}
二隣接切断が階数6で四辺の点対応が有限な窓の軌道閉包に量子点があれば、
その量子点の四面値はすべて同じ $\zeta$ で、$\zeta=-1$ 又は $\zeta^2=-1$ である。
これら二つの添字保持類はいずれも実現する。特に多項式環のGITファイバーには
有限点グラフの窓は一つもない。元の窓に群対称性もAS性も仮定しない。
\end{theorem}
\begin{proof}[三つの必要な一次係数]
定理\ref{thm:quantum-universal-chambers}により元の窓は正則で、その三余指標の一つの正の変位で表示できる。
各辺の階数6は還元で保たれる。量子二項式 $r_{ij}$ ($i<j$) の各々を、それより真に高い重みだけで補正し、
六つの順序対 $ij$ での係数が単位行列である関係基底を取る。
背景で同じ射影文字を各位置に使える。余指標 $(0,1,4,6)$ では最初の二文字の両直線の積に一式しか残らず有限でない。
$(0,6,8,9)$ では上三文字どうしの三式は最低文字を含まないので最低点に三次元接点が生じる。
これは命題\ref{prop:finite-polynomial-limit-flag}のJacobi階数の議論と同じで、初期二項式の交換倍率は非零で足りる。
この場合の量子極限は直前の命題で尽きる。

$(0,3,5,9)$ では順に
\[
  r_{02}\longrightarrow 11,\qquad r_{12}\longrightarrow30,\qquad
  r_{03}\longrightarrow22
\]
への補正はすべて重み差1である。この順に係数を $A,B,C$ とする。
$A=0$ なら最初の二文字に $r_{01}$ 以外の制限がない。
$C=0$ なら最初の三文字に $r_{01},r_{02},r_{12}$ 以外の制限がない。
どちらも有限性に反する。$B=0$ なら $r_{12}$ に最低文字がなく、
$r_{13},r_{23}$ にもないので上と同じ三次元接点を得る。
一辺でもこの場合に入れば直前の命題を使える。
残りでは三係数が非零である。余重み1は初期の一次変形に直接現れる。
各補正の負添字対は $02,12,03$ であり、定理\ref{thm:quantum-face-resonance}の証明の
外積複体による必要条件は単位面や三重、四重の面値でも適用できる。
これらは文字変更の一次境界を持たない。従って
$h_0=h_2=h_1$ かつ $h_0=h_3$ である。面値の積は1なので共通値は四乗根である。
\end{proof}

\begin{proof}[共通値1での二次障害]
全て1の背景は通常の交代体積 $w$ である。上記余重みでのテンソルの最初の二項を
$w+V_1+V_2$ とし、$V_j$ の語の重みは $17+j$ とする。
二切断で $w$ の行列を $M$、$V_j$ の行列を $M_j$ と書く。
$P=(01,02,03,12,13,23)$、$I=M_{PP}^{-1}$、$L,R$ をそれぞれ左核の行、右核の列の基底とすれば、
Schur補の最初の二式は
\[
   L M_1 R=0,\qquad L M_2 R=L M_{1,*,P}\, I\,M_{1,P,*} R .
\]
両切断の式を積み重ねる。重み1では14列の線型部分の階数は11で、三自由変数を $a,b,c$ と取ると
上記の $(A,B,C)=(c,-b,-a)$ である。重み2では12列の線型部分の階数は7で、
右辺をその像で割った条件は $4ab=0$ となる。
\path{computations/verify_finite_quantum_closed_limits.py} が体積の24係数から両行列を構成し、
主元による消去と三つの関係係数を有理数上で検算する。
従って三次元接点がない場合の $BC\ne0$ と両立しない。
三次元接点があっても上述の閉極限は多項式環でない。
逆に二つの非自明型の背景は定理\ref{thm:quantum-tangent-finite}の
有限条件を満たす補正を選べば実現する。全ての面値が等しい場合に奇置換は
共通値を反転するので、残る類数は二つである。
\end{proof}

% !TeX root = ../papers/nc-p3.tex
\section{有限量子ファイバーの三つの六次元正規形}
\label{sec:finite-quantum-planar-fibers}
\begin{theorem}\label{thm:finite-quantum-planar-fibers}
四辺が有限で量子閉極限を持ち、三次元接点がない全窓の添字保持類は、次の互いに素な三表で尽きる。
各表で異なる六座標は非同型であり、全点の射影安定群は自明である。
\begin{enumerate}
\item 共通面値 $-1$ では $(a,b,c,d,e,f)\in\mathbb A^6$ の一表である。
\item 原始四乗根の閉類では、二つの向きにそれぞれ
$\mathbb A^6\setminus\{b=c=d=e=f=0\}$ が一表ずつある。
除いた直線の窓も全次数で正則であり、そこでは点対応に曲線が残る。
\end{enumerate}
向きを選んで $\zeta=-1,i,-i$ と書くと、一移動又はその逆は六座標を順に
$\zeta^{2},\zeta^{3},\zeta^{3},1,1,\zeta$ 倍する。最小周期は非零座標に対応する倍率の位数の最小公倍数である。
前節の最大接点の分類と合わせて、二つの量子ファイバーにある有限グラフを全て扱っている。
\end{theorem}
\begin{proof}[余重みと印なし同型]
前定理で最大接点がないとき、余重みは $w=(0,3,5,9)$ で三係数 $A,B,C$ は非零である。
背景の二文字交換倍率は位置0の文字を他と可換として、
残る $12,13,23$ で $\zeta,\zeta^{-1},\zeta$ に取る。
体積の順列係数は符号に各反転対のこの倍率を掛ける（逆の規約は向きを反転する）。
多重次数の差は四添字の和が零の格子に取る。
三つの余重み1の方向を
\[
 (-1,0,2,-1),\quad(-1,2,-1,0),\quad(1,-1,-1,1)
\]
の順に取れば格子の整数基底である。よって三係数は共通対角群で指定した非零値に一意にできる。
正の残る自由方向は次表で、項目は多重次数、余重み、座標の順である。
\[
\begin{array}{c|r|c@{\qquad}c|r|c}
(-1,1,-1,1)&7&a&(-1,-1,2,0)&7&b\\
(0,-1,-1,2)&10&c&(-1,-1,1,1)&11&d\\
(-1,0,-1,2)&13&e&(-1,-1,0,2)&15&f
\end{array}
\]

具体的な復元を述べる。定理\ref{thm:quantum-universal-chambers}のSchur補の二切断を使う。
各重みで語を辞書順に並べ、線型写像の核は既約行階段形の自由変数を順に1とする基底に取る。
非斉次式の特解には、この核と標準Hermite内積で直交する唯一の解を取る。
ここで複素共役は定数 $\mathbb Q(i)$ のみで使い、パラメータには作用させない。
重み1の核の三座標は全て1とする。表の六方向では座標を加え、他の核方向は軌道接なので零とする。
$\zeta=i$ では正規形の $d,e$ に順に $504i/78125,-7911i/12500000$ を加えて使う。
反対の向きでは共役する。余重み19まで行えば全256係数が定まる。

\path{computations/verify_finite_quantum_planar_fibers.py} が各重みで二つの
Schur補の右辺が線型部分の像に入ることを六変数の多項式恒等式として検算する。
実装では主六行列の逆の次数 $j>0$ 成分を $I_j=-I_0\sum_{l=1}^j M_{l,PP} I_{j-l}$ で得る。
主行列式は重み零で一定なのでこの有限再帰で評価できる。核の負添字二個の成分は記載した九線、
他の非零核は全て根指標である。根指標の方向は二つの面値1を要する横断変形がなく
（定理\ref{thm:quantum-face-resonance}の外積計算）、全て軌道接となる。
したがって直交条件は実際のLuna横断面を取っている。
背景の共通トーラスと有限単項式安定群、巡回作用はこの補空間を保つ。
実際コンパクト対角部分と単位絶対値の置換持上げはユニタリで、対角部分は稠密である。
これで多項式再帰のパラメータは飽和スライス内の軌道の代表を与える。

上の三基底を正にする余指標の開錐は、最低を0とすると正確に
$e_2<2e_1$, $e_1+e_2<e_3<2e_2$ である（ここだけ $e_j$ は矢重みを示す）。
この錐の内部では常に $0<e_1<e_2<e_3$ である。表の六方向は三基底の非負整数和で、
再帰の他の成分もそれらの和のみから生じる。従って流入する余指標の開錐は
三非零一次方向だけで既に決まり、どのパラメータでも同じになる。
この錐から文字の順が復元されるため、同じ順を保つ置換は恒等である。
背景安定群の置換は $-1$ では全対称群、原始四乗根では交代群なので、
順の同値類はそれぞれ一つ、二つである。トーラスにも三単元を固定する非自明な元はない。
Lunaのファイバー同値から元の窓でも安定群と同型は記載通りである。
巡回は正規化前の各核線で根付き単位倍率を与え、指標基底の三倍率を共通対角群で戻す。
上記計算は全テンソルに対するこの巡回と六座標の指定倍率を直接比較する。
\end{proof}

\begin{proof}[有限性の全特殊値]
辺の二端を $[z:u:v:H],[z':u':v':H']$ とする。$H=0$ なら $r_{23},r_{13}$ から
$H'=0$ 又は $u=v=0$ となり、後者では最低点の相手も最低点である。逆も同じである。
$H=H'=0$ では $r_{03}$ が $Cvv'=0$、$r_{12}$ が $uv'-qvu'+kvv'=0$ を与える。
片側だけ $v$ が非零なら他側が最低点となり矛盾する。両 $v$ が零なら
$r_{02}$ から $Auu'=0$ となり、両側とも最低点である。
従ってこの点以外は $H=H'=1$ のアフィンチャートで尽きる。

このチャートでは $r_{23},r_{13},r_{12},r_{03}$ の順に $v',u',z',z$ を解く。
最初の二式はアフィン線型で、後二式から $z,z'$ は $u,v$ の二次以下の多項式となる
（$r_{12}$ の $z'$ 係数は $B\ne0$）。残る二式 $F,G$ は $u,v$ の三次式である。
従って共通成分がなければ有限で、アフィンにある全局所環の埋込み次元は高々2である。
最低点では $B\ne0$ により $r_{12}$ の一次部が $r_{01},r_{02},r_{03}$ の一次部から独立で、
残る二式の一次部は零なので接次元2である。

二平面三次式を斉次化し、無限遠で調べる。$[1:0:0]$ は全パラメータで横断交点である。
もう一方のチャート $[t:1:0]$ で $F,G$ の最高部はそれぞれ次数2と1であり、後者の単根を $r$ とする。
$-1$ 型では前者の $r$ での値が各辺で $1/2$ なのでこれ以外の無限遠交点はない。
従って共通曲線はなくアフィン長は $9-1=8$ である。
四乗根ではこの値が零になる。共通成分があれば無限遠で $r$ だけを通り、
$G$ の無限遠への制限のこの根が単純なので、その成分は直線 $u=rv+k$ でなければならない。
$G(rv+k,v)$ の二次係数が $k$ を一意に定める。その $k$ を代入した $F,G$ が
ともに恒等零となる条件の既約行階段消去（又はGröbner基底）は、全四辺で
ちょうど $(b,c,d,e,f)$ となる。前記コードは最高部・横断性・五式のイデアルを記号計算する。
これで例外の全てが決まる。例外上でも同じ階数6の再帰と良い還元から正則性が従う。
\end{proof}

\begin{corollary}\label{cor:finite-quantum-planar-lengths}
各正規形には長さが二十の半分を超える唯一の点がある。$-1$ 型での長さは12である。
原始四乗根では次表の行を上から読み、満たす最初の行がその長さを与える。
\[
\begin{array}{c|c}
\text{条件}&\text{長さ}\\\hline
 b\ne0&13\\
 d\ne0&14\\
 f\ne0&15\\
 ce\ne0&17\\
 e\ne0\text{ 又は }ac\ne0&18\\
 c\ne0&19
\end{array}
\]
\end{corollary}
\begin{proof}
前証明の第二無限遠点の交点数を $m$ とすればアフィン長は $8-m$、
元の窓の最低点の長さは $12+m$ である。斉次化の第三変数を $h$ とし、
$F(t,1,h),G(t,1,h)$ の $t$ に関する終結式の $h$ での位数が $m$ である。
$F(t,1,0)$ の二次の先頭係数は非零定数で $G(t,1,0)$ の零点は $r$ のみだからである。
前記コードでこの終結式の次数 $1,2,3,4,5$ までの逐次係数イデアルを取ると、
順に新しく $b,d,f,0,ce$ が加わる。次数6でさらに $ac^2-e^2,e^3$、
次数7でさらに $c^3$ が加わり、次数8で変わらない。有限点の条件の下で
これらの消失を順に読めば表になる。$-1$ 型はアフィン長8と全長20から直ちに従う。
\end{proof}

% !TeX root = ../papers/nc-p3.tex
\section{全ての有限点グラフでの四文字評価と成分の次元下界}
\label{sec:finite-four-slot-evaluations}
標準有限Frobenius窓とし、四辺の点対応が有限であるとする。
定理\ref{thm:finite-cubic-separators}で四移動が恒等なので、各位置の四文字を
同じ長さ20のスキーム $X$ 上で単元倍の自由度を持つ $U_i\subset D=H^0(\mathcal O_X)$ として読む。
\begin{lemma}\label{lem:finite-adjacent-strict-half}
隣接する積 $A=U_iU_{i+1}$ と、長さ $0<m<20$ の任意の部分スキーム $Y\subset X$ について
\[
                     \dim(A|_Y)>m/2
\]
である。被約性・支持数・局所埋込み次元の仮定は要らない。
\end{lemma}
\begin{proof}
$C=H^0(\mathcal O_Y)$ は $m$ 次元で、各座標空間の像 $U,V$ に単元1を含めてよい。
$U,V$ はそれぞれ $C$ を代数として生成する。実際、各三次式の制限が $D$ に全射である。
$m\ge2,4,8$ なら $\dim U,\dim V$ はそれぞれ $2,3,4$ 以上である。
これは射影が閉埋込みで、点・直線・平面の切断の長さの上限が $1,3,7$ であることによる
（系\ref{cor:finite-no-four-on-line}、定理\ref{thm:finite-plane-section-seven}）。
$\dim UV<m$ なら $\dim UV>\max(\dim U,\dim V)$ である。
等号なら、例えば $UV=U$ となり、$V$ の生成する代数で $U$ が閉じるため $U=C$ となる。
上の三段の下界を使えば $1\le m\le9$ で主張を得る。

$m=10$ で $\dim UV=5$ を仮定する。$U+V\subset UV$ なので
$\dim(U\cap V)\ge3$ である。$U(U\cap V)\subset UV$、かつ $U/(U\cap V)$ の次元は高々1なので
$\dim U^2\le6$ となる。$U$ による閉埋込みのHilbert関数は従って
$1,4,h_2,10$ で $h_2\le6$ である。全台で非零の一次式で割ったArtin還元では
二次の次元が高々2だが三次が4以上になる。単項式初期イデアルに移すと、
高々二つの標準二次単項式から、それらを全ての二次約数に持つ標準三次単項式は高々二つしか作れない。
矛盾により $m=10$ も従う。

$A$ は十次元で、辺の $(2,2)$ 完全跡により $A=A^\perp$ である
（定理\ref{thm:finite-square-chart}の証明）。$Y$ の定義イデアルを $I$、
残余 $Y'$ のイデアルを $I^\perp=\operatorname{ann}(I)$ とすれば、$Y'$ の長さは $20-m$ である。
$A\cap I=(A+I^\perp)^\perp$ の次元を取って
$\dim(A|_Y)=m-10+\dim(A|_{Y'})$ を得る。これが $10<m<20$ を既に示した範囲に戻す。
\end{proof}

\begin{theorem}\label{thm:finite-four-slot-rank}
$P=U_0U_1U_2U_3$ の次元は19または20である。19になる必要十分条件は、
両側で反対二位置の座標が射影線型に同定されること、同値に四テンソルが
両方の反対位置の交換で対称な表示を持つことである。
階数19の場合の唯一の零化汎関数は完全跡である。
従ってどの枝にも $P$ を消す非零の退化積対合はない。
\end{theorem}
\begin{proof}
$P$ を消す汎関数を一つの完全跡 $\tau$ に対し $\tau(g\,\cdot)$ と書く。
$g\ne0$ が非単元なら $C=D/\operatorname{ann}(g)$ の長さは $0<m<20$ で、
$\tau(g\,\cdot)$ は $C$ 上の完全跡となる。隣接二積 $U_0U_1,U_2U_3$ の $C$ での像が
互いに直交するが、前補題で両次元が $m/2$ を超え、矛盾する。
ゆえに非零の零化汎関数は全て完全である。二本独立にあれば、
一つの局所因子の剰余体で対応する二単元を比例させて非零の非単元を作れるので、余次元は高々1である。

完全跡で $P$ が消えれば $U_0U_2$ と $U_1U_3$ が直交する。
定理\ref{thm:finite-opposite-product-gap}で各積の次元は10又は16であるから、両方が10である。
同定理が両側の射影同定と対称表示を与える。二つの交換は異なる位置の組で行うので同時の基底選択でよい。
逆に両側が同定されれば $P$ は単元倍を除き $U_0^2U_1^2$ であり、
辺の $(2,2)$ 完全交叉の跡により正確に超平面である。
\end{proof}

\begin{corollary}\label{cor:finite-component-dimension-bound}
二隣接階数6のアフィンテンソル軌跡で、有限Frobenius点を通る各成分の次元は
両対称点なら75以上、その他の点なら76以上である。
従って次元74以下の成分は有限Frobenius窓を含まない。
次元75の成分の有限Frobenius部分はすべて両対称である。
\end{corollary}
\begin{proof}
有限グラフの各部分評価は各台点で階数3である。補題\ref{lem:finite-flat-conormals}の
有限平坦なSchur恒等式は $\dim P$ 本の独立行を与える。実際、その補題の四文字は
中央の零点対と、グラフで一意な直前・直後の座標である。
第一切断の滑らかな156次元軌跡の上で残る式は高々 $100-\dim P$ 本であり、
高さ定理により各局所成分の次元は少なくとも $56+\dim P$ となる。
\end{proof}

% !TeX root = ../papers/nc-p3.tex
\section{全有限窓の半安定性と厳密半安定枝の全次数正則性}
\label{sec:finite-window-semistability}
\begin{theorem}\label{thm:all-finite-windows-semistable}
四辺の点対応が有限な全ての標準有限Frobenius窓は半安定である。
厳密半安定、すなわち安定でない場合は全て標準quadratic Koszul左右AS正則で、
唯一の多重安定極限も標準正則である。特に有限点グラフにおいて全次数判定がなお必要な候補は、
閉軌道と有限安定群を持つ安定点に限られる。
\end{theorem}

\begin{proof}[Hilbert--Mumfordの全不等式]
四位置の任意の一径数群を対角化し、各位置の文字重みから最低値を引いて
$0,d_i,e_i,f_i$ の昇順とする。窓で比較すべき平均は $M=\sum_i(d_i+e_i+f_i)/4$ である。
辺 $i$ の六関係に低い重みから初期部を取り、その独立な初期基底の重みを
$r_{i,1}\le\cdots\le r_{i,6}$ とする。
最低文字の一点と隣の低い二文字の直線との積に関係が一つもなければ有限性に反する。
最高文字だけを片側に持つ関係は純テンソルになり存在しない。従って
\[
 r_{i,1}\le\min(d_i,d_{i+1}),\qquad
 r_{i,6}\le\min(e_i+f_{i+1},f_i+e_{i+1}).
\]
窓テンソルの最低重み $m$ は各 $k$ で $r_{i,k}+r_{i+2,7-k}$ 以下である。
実際、完全対合の係数行列で最初の $k$ 行と最初の $7-k$ 列の交叉が恒等零になることはない。
適合基底の初期テンソル積は各重み層で独立である。
$k=1$ の四本を平均し、二つの最小値をそれぞれ算術平均で抑えれば $m\le M$ を得る。
これが半安定性である。

等号なら用いた各不等式が等号である。特に $d_i=d$、$f_i-e_i=g$ は共通であり、
$r_{i,1}=d$, $r_{i,6}=e_i+e_{i+1}+g$ となる。
$m=d+e_i+e_{i+1}+g$ が全 $i$ で同じなので $e_i=e_{i+2}$ である。
したがって矢重みは $0,d,d+s,d+s+g$ と $0,d,d+t,d+t+g$ が交互に並ぶ。
$0\le s\le t$ としてよい。以下スカラーでない等号方向を調べる。
\end{proof}

\begin{proof}[すべての切断階数を保つこと]
まず $d=0$ は不可能である。$g>0$ なら最低空間の次元 $l_i$ は
中央差が正なら2、零なら3で、最高の次の空間の次元 $h_i$ はそれぞれ1,3である。
最低関係の本数は低い二固有空間の積への有限切断から $l_i+l_{i+1}-2$ 以上である。
これらは反対辺の最高関係としか対合できない。最高関係は最高文字と次高文字の二通りの積と
最高どうしの積だけに入り、両側の次高部分への射影が純関係の不在から単射である。
よって本数は $\min(h_i,h_{i+1})$ 以下である。$l$ の組 $2,2$、$2,3$、$3,3$ の全てで矛盾する。
$d=g=0$ なら非スカラー位置の重みは二重ずつ、スカラー位置の重複度は4である。
非自明な方向では各辺の最低空間の次元は $(2,2)$ 又は $(2,4)$ で、最高も同じ組である。
有限な最低切断に必要な本数は2又は4だが、最高ブロックでSegreを避ける六関係部分の
次元上界は1又は3である。これも両立しない。

これで $d>0$ であり、最低点どうしが点列をなす。三次元接点がある場合は最後に扱う。
それ以外ではこの点対で六関係の一次部は少なくとも四本独立である。
二端の重みを上の $s,t$ の順に並べると $M=3d+s+t+g$ である。
各辺で添字 $i$ を略し、次の上界が得られる。
\[
\begin{array}{ll}
 r_2\le d+\min(d,s), & r_5\le 2d+t+g,\\
 r_3\le d+s+\min(d,g), & r_4\le\min(d+t+g,\,2d+s+t).
\end{array}
\]
実際、低い $a,b$ 文字の射影空間の積への有限切断には $a+b-2$ 本以上必要である。
$r_2$ には $2,2$ と $3,1$（小さい中央差の側で3）、
$r_3$ には $4,1$ と $3,2$、$r_4$ には $3,3$ を使う。
$r_4$ のもう一つの上界は、$d+t+g$ を超える関係が最低点で一次部を持たないことによる。
$r_5$ の上界を超える項は最高二文字ずつの積だけに入る。ここで純テンソルを避ける空間は高々一次元である。

反対辺も同じ上界を持ち、$r_{i,2}+r_{i+2,5}\ge M$ から $s\le d$、
$r_{i,3}+r_{i+2,4}\ge M$ から $d\le g\le d+s$ を得る。
後者では $\min(d,g)+\min(g,d+s)\ge d+g$ を使った。
これらの範囲では各対の上界の和がちょうど $M$ である。
よって六重みは全辺で正確に
\[
 d,\ d+s,\ 2d+s,\ d+t+g,\ 2d+t+g,\ 2d+s+t+g
\]
となる。相補切断の行列式のどの非零項も六係数の重みの和が $6M$ で、
各非零係数の重みは $M$ 以上である。従って重み $M$ の部分だけでも行列式は非零であり、
極限窓の隣接二切断は階数6である。
\end{proof}

\begin{proof}[中央差の一致]
$s<t$ で $s<g$ とすると、重み $r_5$ 以上の全項のうち、$t$ 側の最高文字を含まないものは
もう一方の最高文字と $t$ 側の次高文字の積一つだけである。
実際次高どうしの重みは $2d+s+t<r_5$ で、最高と重み $d$ の組も $t$ 側が最高でなければ $r_5$ 未満である。
上位二関係をこの一成分へ射影した核には純関係が生じ、矛盾する。
従って $s<t$ なら $s=d=g$ である。単位を $d$ として矢重みは
$0,1,2,3$ と $0,1,1+u,2+u$ ($u>1$)、関係は $1,2,3,2+u,3+u,4+u$ で全て単純である。
文字を昇順に $0,1,2,3$ と呼ぶ。最低関係の $01,10$ は最低点の両評価からともに非零である。
最高関係の $23,32$ の初期係数も、どちらかが欠ければ最高どうしの補正だけが残って純関係となるため、ともに非零である。
極限の第一切断の階数1部分で $0132$ は非零で、次の切断には二文字 $13$（第二位置の重み1、第三位置の重み3）があり、
その和4が関係重みでなければならない。従って $u=2$ である。
このとき $30$（重み3）は最初の三関係のうち第三の初期部に必要である。
第二位置の最低点で評価した階数が3であり、第四以後の関係は第一位置の文字をそこへ供給できないからである。
$03$（重み4）は第四関係の初期部に必要である。第五以後は一次部が零であり、最低点のJacobi階数は4以上だからである。
ゆえに極限の $3003$ の係数がこの二つの単純重み間の完全対合により非零となる。
ところが次の切断には $00$ の関係がない。これで $s=t$ が従う。
\end{proof}

\begin{proof}[等号方向のPBW還元]
$s=t=0$ なら $g=d$ で、矢重みは $0,1,1,2$ にできる。重み1の両混合射影は、
各最低点での評価により二次元の同型である。重み3の両混合射影は
実際の最高二関係における純関係の排除から同型である。重み2では $03,30$ への射影が
最低点のJacobi階数4以上から同型である。命題\ref{prop:double-middle-finite-pbw}の証明で
極限窓のPBW性を得る。その証明で用いる中央の二行列の可逆性は、
重み1対3の完全対合の両混合射影が同型であることから従い、極限の点グラフの有限性は要らない。

$s=t>0$ の場合、四文字は全て単純である。重み $d$ の $01,10$ と最高重みの $23,32$ は
前段と同様に全辺で非零である。この二つの関係線が対辺で完全に対合することを別の切断に読むと、
$02,20$ と $13,31$ の四係数も各々の関係線で非零となる。
実際 $0132,1023$ の中央対と外対を使い、開始位置を全て回せばよい。
六つの関係重みのうち中央二つを除く四つは単純で、各関係の左辞書順先頭は $10,20,31,32$ である。
三つの不等式 $0<s\le d\le g\le d+s$ からこれらの重みで
あり得る余分な項は第二の $11$、第五の $22$ だけで、いずれも先頭を変えない。

$g>d$ なら中央二重みも単純で、$12,21$ と $03,30$ は $0123,1032$ を別切断で読むとともに非零である。
前者の先頭は $21$、後者では $22$ の追加が可能でも先頭は $30$ である。
$g=d$ なら中央二関係の $03,30$ への射影は最低点でのJacobi階数から同型である。
両対の $03,30$ 主小行列は $00,33$ の不在から可逆対角である。
行 $21$、列 $03$ の係数は別切断の $10,32$ により非零なので、
$30$ を除いた関係の先頭は $21$ となる。これで全場合で六つの下降対が先頭である。
各一文字の全支持はこれらの非零係数と単純重みから従う。
四テンソルの独立な四縮約が三次の重なりを全て解消し、定理\ref{thm:simple-weight-finite-windows}の
PBW証明により極限は標準Koszul左右AS正則である。命題\ref{prop:unrestricted-good-reduction}で元の窓にも戻る。

三次元接点があれば定理\ref{thm:maximal-tangent-ore}の純三文字関係が
ある点縮約の中央支持をも三次元にする。実際その縮約の末尾縮約像は
点での左評価核の三空間であり、これは接次元3により右評価核でもある。
定理\ref{thm:finite-cubic-separators}により元の窓は正則で、命題\ref{prop:maximal-tangent-closed-limits}で
多重安定極限も正則である。以上、三次元接点のない場合には任意の非自明な等号方向で上の証明が働く。
厳密半安定点では数値判定と閉軌道への一径数群判定により、必要な等号方向を選べる。
すでに多重安定で安定でない場合は正次元の簡約安定群内の一径数群を使う。
\end{proof}

% !TeX root = ../papers/nc-p3.tex
\section{有限点窓の安定性を二つの内在的な旗だけで判定する}
\label{sec:finite-intrinsic-stability-test}
\begin{theorem}\label{thm:finite-intrinsic-stability-test}
有限Frobenius窓の四辺が有限であるとする。三次元接点があれば厳密半安定である。
それ以外では次の有限検査で安定性が決まる。長さ9以上で埋込み次元2の各局所点 $p$ に対し、
四位置で $p_i$ とその射影接平面を取る。このような点は高々二つである。
\begin{enumerate}
\item $p_i$ で非零な文字 $z_i$、接平面を切る文字 $h_i$、その間の二文字 $U_i$ を
順に重み $0,1,2$ の各部分に対応させる。
$K_i=R_i\cap(\ker p_i\otimes\ker p_{i+1})$ を、
$(\ker p_i/kh_i)\otimes(\ker p_{i+1}/kh_{i+1})$ へ写す。
全四辺で像が零なら厳密半安定である。
\item 点の長さが10以上で、局所環の極小二変数表示の二次初期式空間が
一つの一次式の平方で張られるとする。その接方向直線を各接平面内に戻し、
点、接方向直線、接平面の完全旗を取る。対応する四文字の重みを $0,1,2,3$ とする。
窓の重み6未満の全係数が零なら厳密半安定である。
\end{enumerate}
すべての点で両検査が不成立なら安定である。重み検査は旗に適合する任意の基底で行える。
したがって曲線状の点だけを持つ有限窓は、重複度によらず安定である。
三次元接点も長さ9以上の平面局所点もない場合にも同じ結論となる。
\end{theorem}
\begin{proof}[二重の中央重みと二次関係の検査]
点対の周囲で両 $z$ を1にする。グラフの接空間は両接平面の間の接線型同型のグラフなので、
六方程式の微分の階数は4である。両方で最低文字を用いない関係は正確に $K_i$ で、次元2となる。
$w(p_{i-2},\eta_{i-1})$ を、$\eta_{i-1}$ が点の代表を法として二つの接方向を動くときに取れば、
その像は $K_i$ である。実際、$(p_{i-2},p_{i-1})$ でのグラフ上の微分により
$w(p_{i-2},\eta_{i-1})=-w(\eta_{i-2},p_{i-1})$ である。
一方は末尾が、他方は先頭が点で消える関係に入る。点での評価が階数3で、
その零化点は $p_{i-1}$ だけだから接方向二つの像は独立である。

重み4未満の四語を調べる。最低文字二個が隣り合えば点対条件から零である。
反対位置に最低文字が二個なら、その間の残る少なくとも一文字は中央文字である。
対応する係数は上の $K_i$ のどちらかの端に最低文字を入れたものなので零となる。
最低文字が一個なら残る三文字は全て中央文字である。この全係数の消失は
$K_i$ の両中央文字への射影の消失とちょうど同値である。
これで第一検査が重み4未満の消失と同値であることが分かる。
第二検査と同様に数値判定の等号を与える。各検査は重み以上の部分空間のみによるので、
旗の分裂の選択には依存しない。
\end{proof}

\begin{proof}[一般の等号方向が与える旗の復元]
三次元接点があれば前定理の証明と命題\ref{prop:maximal-tangent-closed-limits}によって
正次元の点グラフを持つ閉窓へ落ち、安定ではない。
残る厳密半安定窓に、定理\ref{thm:all-finite-windows-semistable}のスカラーでない等号方向を取る。
その最低点 $p$ で、上二関係には一次部がなく、最初の四関係の一次部は独立である。
よって埋込み次元は2で、上二関係がちょうど $K_i$ である。
中央二関係の一次部は両最高文字 $h_i,h_{i+1}$ の独立な二式である。
実際、単純文字の場合にその重みは $2d+s,d+s+g$ 以上で、
最低文字を含む一次部に現れる文字は最高文字だけである。
中央二重の場合にも同じである。これで $h_i$ が実際の接平面の式となる。

中央二重の場合、$K_i$ は重み3以上であり、両中央文字への射影は零である。
単純四重み $0,d,d+s,d+s+g$ の場合も、上二関係の重みは $2d+s+g$ 以上であり、
両中央文字のどの積よりも真に大きければ第一検査が成立する。
この不等式が破れるのは $s=g$、従って $s=d=g$ の場合だけである。
ここで上二関係に残り得る中央積は $22$ のみである。
ある辺でこの係数が非零なら、全辺で非零である。
実際その重み4の関係線と反対の重み2の線の完全対合から
$2220,2202$ の係数はともに非零である（相手の $20,02$ はともに非零）。
次の切断で隣接する $22$ を読めば両隣の辺にも伝播する。

各位置の文字を $z,u,v,h$ とする。局所二変数を $u,v$ として最初の四式で他の四変数を形式的に消去する。
一次では $h,h'$ は零となり、重み2の関係の $02,20$ の両係数が非零なので
$v'$ は $v$ の非零倍となる。残る二式の二次部は一方が $vv'$ の非零倍、他方は零である。
したがって二次初期式空間は $v^2$ で張られる。この平方が決める接方向は重み
$0,1,2,3$ の元の旗の点と接平面の間の直線である。
元の等号条件は正確に第二検査の重み6未満の消失となる。
\end{proof}

\begin{proof}[点の長さ]
第一検査なら形式的な二変数への消去後、$K_i$ のどの項も一次の中央文字と二次以上の
$h,h'$ の積か最高文字どうしの積である。従って二式はいずれも位数3以上で、局所交点数は9以上となる。
第二の旗を与えた等号方向では、アフィン六変数の重みは両側で $1,2,3$、
六関係の最低重みは $1,2,3,3,4,5$ である。各変数をその重み乗で置き換えると、
形式局所環の拡大は階数36の有限自由拡大になる。置換後の各式の通常の位数は
対応する最低重み以上である。正則局所環の局所Bézout不等式で、
拡大後の完全交叉長は六つの位数の積以上、従って元の長さは $360/36=10$ 以上となる。
全長20と合わせて、検査する点数と判定の必要十分性が従う。
\end{proof}

% !TeX root = ../papers/nc-p3.tex
\section{有限両対称窓の非正則境界を与える二元十次形式と三割線}
\label{sec:binary-secant-obstructions}
$P_j=k[s,t]_j$ とし、$D\in P_{a+b}^*$ の積対合を
$H_{a,b}(D):P_b\to P_a^*$ と書く。以下の幾何的記述は代数閉包上で行う。
\begin{theorem}\label{thm:binary-secant-obstructions}
点グラフが有限な両対称Frobenius窓のうち非正則なものは、必要なら始点を一つ移し、
以下の構成で全て得られる。次数十の汎関数 $T$ は $G=H_{5,5}(T)$ が可逆なもの、
次数八の汎関数 $Y$ は $T\mathbin{\lrcorner}P_2$ の外にあるものを取る。
$U=P_2\oplus k$, $Q=P_3$ とおき、積を
\[
 B((p,a),q)=pq+aG^{-1}H_{5,3}(Y)q\in P_5
\]
とする。四窓は $w(u,q,u',q')=G(B(u,q),B(u',q'))$ である。
固定した二元座標と $[T]$ に対し、余分な文字の選び方を除いたパラメータは
$[\bar Y]\in\mathbb P(P_8^*/(T\mathbin{\lrcorner}P_2))=\mathbb P^5$ である。
ここで始点の移動は二つの向きを列挙するために用いる。
二元形式への商の印を忘れた同型類の重複は、このパラメータ表示では取り除いていない。

$S_j$ を有理正規八次曲線の長さ $j$ の割線多様体とする。
中心 $\Pi=\mathbb P(T\mathbin{\lrcorner}P_2)$ からの射影
$\pi:S_3\to\mathbb P^5$ は次数20の有限平坦被覆である。
点グラフが有限であるための必要十分条件は
$[\bar Y]\notin\pi(S_2)$ であり、そのときグラフは被覆のスキーム論的ファイバーである。
この開集合では定理\ref{thm:frobenius-tensor-reconstruction}の三次・四次交叉を調べれば
求める非正則Frobenius窓が正確に得られる。一般の $(T,[\bar Y])$ は両交叉検査を通り、
相異なる20点を持つ安定な非正則窓を与える。固定した $T$ での有限平坦被覆の
トレース判別式は $\mathbb P^5$ 上次数90の因子を定める（重複度を含む）。
\end{theorem}
\begin{proof}[二元商からの復元]
有限Frobenius窓では純関係がなく、両評価の階数は3以上である。
非正則なら定理\ref{thm:bisymmetric-two-length-test}により一方に二元形式への商がある。
関係六平面を商に写す核は、一方の文字が固定された純関係だけからなり零。
六次元どうしなので二元乗法の双対像に一致する。双対に移り、関係空間の双対を $P_5$ と同一視し、
指定された三文字の補空間を選ぶと積は $B((p,a),q)=pq+aLq$ である。
両対称性と隣接切断の階数6により $w$ は非退化な対称形式 $G$ によるこの積のGram形式である。
逆に、一方の $q,q'$ の交換不変性だけを課せば、Gram対称性から他方も交換不変となる。

二元単項式を $t$ の指数で番号付けする。三文字どうしの交換条件は
$G_{a+i,b+j}=G_{a+j,b+i}$ ($0\le a,b\le2$, $0\le i,j\le3$) であり、
対称行列 $G$ の解はちょうどHankel行列である。
$J=GL$ に対する条件は $J_{a+i,j}=J_{a+j,i}$ となり、解も全てHankelである。
実際どちらも座標どうしの同一視だけで、第一の21座標は11個、第二の24座標は9個の
反対角線ごとに連結する（添付検算では各等式の整数行列の階数10と15を求める）。
従って $G,J$ は一意に $T,Y$ を与える。
第四文字の平行移動と非零倍は $Y$ の三次元の極性空間による変更とスカラー倍となる。
極性空間に入れば第四文字が他の三つに従属する。外では四文字は独立であり、
$Q$ 側の独立性は二元乗法から従う。$B$ は全射で、両隣接階数は6である。
逆向きには二元商による指数成長が全てのパラメータで正則性を排除する。
\end{proof}
\begin{proof}[有限被覆と点グラフ]
$u=[p:a]$ での関係評価は $G$ を掛けると $H_{5,3}(T\mathbin{\lrcorner}p+aY)$ であり、
その核にある $[q]$ がグラフの他の座標になる。$a=0$ かつ $p\ne0$ なら
これは $q\mapsto G(pq)$ で階数4。従って中心は階数3以下の多様体と交わらない。
$([Z],[q])\in\mathbb P(P_8^*)\times\mathbb P(P_3)$ で $Z(qP_5)=0$ の入射を取る。
非零 $q$ による乗法は単射なので、これは $\mathbb P^3$ 上の階数3の束
$E=\ker(9\mathcal O\to6\mathcal O(1))$ の直線の射影化である。
相異なる三根を持つ一般の $q$ ではその平面は根の三つの評価点の張る平面であり、
一般の $Z$ は階数3でただ一つの核 $[q]$ を持つ。従って入射像は $S_3$ で次元5。
六行四列の最大小行列は期待余次元3を持ち、Eagon--Northcottにより商はCohen--Macaulayである。
階数3上は核線の消去で滑らかな入射と局所同型だから一般被約、従ってこの小行列スキームは被約である。
核に独立な三次式 $q,r$ がある場合、その最大公約式 $d$ に対し $qP_5+rP_5=dP_{8-\deg d}$
（最小公倍式で次元を数える）。$\deg d\le2$ より、これはちょうど $S_2$ の条件となる。
こちらも二次式に対する入射により次元3である。

射影の各ファイバーは中心を超平面に持つ線型空間内の、中心と交わらない射影スキームなので有限である。
像の次元は5で全射となり、Cohen--Macaulay性と標的の正則性から平坦である。
$h$ を三次式の空間の超平面、$H$ を入射の $Z$ 座標の超平面とすると、
$c(E)=(1+h)^{-6}$ により $\int H^5=\binom63=20$, $\int hH^4=\binom62=15$ となる。
階数3上は核線が逆射を与え、グラフはファイバーに一致する。
反対射影でも $q$ を固定した $p\mapsto pq$ の単射性から核は高々一次元で、同じく閉埋込みである。
階数2以下の点があれば、その上の $[q]$ が正次元となる。
\end{proof}
\begin{proof}[判別式と一般パラメータでの非正則Frobenius性]
$\pi(S_2)$ は次元3なので、一般直線はこれと交わらない。
入射上に四つの一般超平面を戻して滑らかな曲線を得る（標数零のBertini）。
入射の標準類は $-3H+2h$ なので、この曲線から直線への射の分岐の次数は
$\int (3H+2h)H^4=90$ である。一般直線上での有限平坦トレース判別式の位数は、
滑らかな曲線の分岐指数引く一の和であり、これで全体の次数を得る。
特に一般ファイバーは被約である。

三次交叉の次元4と四次交叉の次元1は、上述の隣接階数と一文字階数を固定すれば開条件である。
空でないことは単項式番号順に
\[
 T=(2,3,5,7,11,13,17,19,23,29,31),\qquad Y=(3,1,4,1,5,9,2,6,5)
\]
を使って有理数上で確認できる。
各向きで隣接二挿入の連結行列の階数は44、三次縮約を両側から挿入した連結行列の階数は31である。
全パラメータ空間は既約で、二条件の開集合は被約有限グラフの非空開集合と交わる。
定理\ref{thm:finite-intrinsic-stability-test}によりこれらは安定である。
計算は\path{computations/verify_binary_secant_obstructions.py}で再現できる。
\end{proof}

% !TeX root = ../papers/nc-p3.tex
\section{二元商の全例外値、商の一意性と同型}
\label{sec:binary-apolar-rank-test}
\begin{theorem}\label{thm:binary-apolar-rank-test}
前節の $([T],[\bar Y])$ に対し、$I_6=\{f\in P_6:T(fP_4)=0\}$ とおく。
$\bar Y$ は自然に $(I_6\otimes P_2)^*$ の非零元をスカラーまで定める。
この二行三列行列が階数2であることと、構成した窓が有限Frobenius検査を通ることは同値である。
このとき反対階数は常に8である。階数1では三次交叉が既に過大となる。
従って有限点グラフを持つ非正則Frobenius窓の指定の向きでのパラメータは、
階数1のSegre多様体と前定理の $\pi(S_2)$ を除けば、例外なく全て許される。

指定の向きで二元商の核直線は一意で、二元座標は同時の $\operatorname{PGL}_2$ まで一意である。
有限グラフのこの向きの二組の表示が添字保持同型である必要十分条件は、
$([T],[\bar Y])$ が同じ $\operatorname{PGL}_2$ 軌道に属することである。
従ってこの向きの非正則有限境界は既約で、基底変更を除いた次元は12である。
\end{theorem}
\begin{proof}[差分式]
$P_d$ を一変数 $x$ の次数 $d$ 以下の空間として書く。モニック六次式 $f\in I_6$ を選べる。
実際次数5以下の係数に関する五条件の行は可逆Hankel行列 $G$ の初めの五行である。
$e\in P_5$ を $G(e,x^j)=\delta_{j5}$ で定め、斉次化ではこれを次数6にも見る。
$f,e$ は $I_6$ の基底で、共通零点を持たない。
実際 $A=k[x]/f$ で $T$ の次数10以下の積評価は剰余の評価だけで決まる
（$T(fP_4)=0$）。従って $G(e,b)$ は次数5以下の剰余 $b$ の最高次係数である。
もし $a$ が共通根なら $b=f/(x-a)$ に対し $eb=0$ が $A$ で成立し、その最高次係数1と矛盾する。
無限遠では $f$ は零でない。従って $I_6P_2$ は六次元で、その零化空間は
$T\mathbin{\lrcorner}P_2$ に次元から一致する。

$L_j=G^{-1}H_{5,3}(Y)x^j$、$\alpha_j=Y(fx^j),\ \beta_j=Y(ex^j)$ ($0\le j\le2$) とおく。
通常の多項式の等式として
\[
                L_{j+1}-xL_j=\alpha_j e-\beta_j f
\]
が成り立つ。実際 $A$ で左辺は $1,\ldots,x^4$ と直交し、$x^5$ との対合は
$x^6$ の剰余を用いて $\alpha_j$ である。また $L_j$ の五次係数は
$G(e,L_j)=\beta_j$ なので、多項式への持上げも式の通りである。
これから $K_{ij}=G(L_i,L_j)$ について
$K_{i+1,j}-K_{i,j+1}=\alpha_i\beta_j-\beta_i\alpha_j$ を得る。

反対切断を $U$ の対称二乗から $Q$ 上の対称形式への写像として読む。
$P_2\cdot P_2$ と $P_2\cdot z$ の像は六次のHankel空間内でそれぞれ
$T\mathbin{\lrcorner}P_4$ と $Y\mathbin{\lrcorner}P_2$ である。
前者は五次元で、その零化空間が $I_6$ なので、和の次元は $5+\operatorname{rank}(\alpha,\beta)$。
$z^2$ の像 $K$ は上式から階数2の場合に限りHankelでない。従ってこの場合全階数は8となる。
\end{proof}
\begin{proof}[両始点の交叉]
$Q,U,Q$ 順の三次交叉の元を、両側からそれぞれ関係 $B^*(P_5^*)$ に入る三重線型形式として読む。
$U$ を $P_2$ に制限すると、両方のHankel条件から係数は三指数の和だけに依存し、
一意な $D\in P_8^*$ になる。この制限から両側の関係係数は全て戻る
（$P_2P_3=P_5$）。$z$ で両側が一致するための残る条件は
$M_{ij}=G(L_i,L'_j)$ が対称であることであり、$L'_j$ は $D$ から同様に作る。
$D$ の二行を $\alpha',\beta'$ とすると
\[
 M_{i,j+1}-M_{i+1,j}=\beta_i\alpha'_j-\alpha_i\beta'_j.
\]
各反対角線で差を加えて $M_{ij}-M_{ji}$ を取ると、対称性は
$\beta\odot\alpha'=\alpha\odot\beta'$ と同値である。実際長さ1の差で対角三成分、
長さ2で隣の二成分、長さ3で残る一成分が順に得られる。
階数2なら独立な一次式 $\alpha,\beta$ の一意分解から
$(\alpha',\beta')$ は $(\alpha,\beta)$ のスカラー倍だけである。
従って交叉は $\dim(T\mathbin{\lrcorner}P_2)+1=4$。
階数1なら行基底を変えて $\beta=0$ とすれば $\beta'=0$ だけが条件で、交叉は六次元となる。

$U,Q,U$ 順では両端を $P_2$ に制限すると、同じ指数のずらしにより
$D\in P_7^*$ を得る。片方の端だけを $z$ にした値は
$D(x^i L_j)$ ($0\le i\le2$, $0\le j\le3$) に強制される。
これがHankelなら、端が $z$ の側の六つの関係係数も一意に定まり、両端が $z$ の値は
左右でその同じ汎関数を $L_j$ に作用させるので一致する。
従って条件は $D(x^i(\alpha_j e-\beta_j f))=0$ ($i=0,1$, $0\le j\le2$) で尽きる。
階数2なら $D$ は四次元の $eP_1+fP_1$ を零化し、解も四次元となる。

以上より四次交叉の元は、第一文字または末尾文字の自己写像 $S,N$ を $w$ に作用させたものとして
両方に表せる。例えば $w(Su,q,u',q')=w(u,q,u',Nq')$ である。
左辺の $q,q'$ 交換不変性を使えば右辺の $N$ を $q$ へ移せる。
$B$ の全射性と $G$ の非退化性から $B(Su,q)=B(u,Nq)$ となる。
さらに元の等式を使えば $B(u,Nq)$ は $B(u,q)$ に一つの線型写像を作用させたものになる。
従って $B(p,q)=pq$ に制限して $N$ は二元積の全ての一致を保つ。
$p=1,x$ により $N(x^{j+1})=xN(x^j)$ ($j=0,1,2$) となり、全ての値の次数上界3から $N$ はスカラー。
他方の向きでは役割を交換して同じ等式を得る。両四次交叉は一次元で、
復元定理が有限Frobenius性を与える。
\end{proof}
\begin{proof}[商の印の復元]
$Q$ 側の反対核は二次元で、二元三次評価のねじれ三次曲線 $C\subset\mathbb P(Q^*)$
を消す三つの二次式の中の鉛筆である。各非零式は既約である
（非退化整曲線 $C$ は二平面の和に入らない）。
従って二つの二次式は次数4の完全交叉を切り、その中で $C$ は唯一のねじれ三次曲線である。
よって $C$ とその二次Veroneseの七次元線型包であるHankel空間は内在的に定まる。
この空間で反対切断の値を割ると、$P_2\cdot U$ では零で $z^2$ だけが非零。
この階数1の二次形式の根基が元の三平面 $P_2\subset U$ を復元する。
これで商の核直線が決まり、命題\ref{prop:binary-quotient-rank-two-curve}で
商における階数一曲線と両側の二元座標が共通の射影座標変換まで定まる。
第四文字を加える自由度は前定理で既に割っている。

独立な四位置の同型については、点グラフが反転を交互に繰り返すので両偶数位置の誘導写像は
点スキーム上で一致する。点スキームの二次式制限は単射であり、二つの射影線型写像が
そこで一致すれば比例小行列式が恒等零となる。奇数位置についても同じで、
従って二位置の基底変更だけを考えればよい。逆に二元座標の変換は積、極性空間、Gram形式をすべて運ぶ。
既約なパラメータ空間の次元は $10+5=15$、一般の十次形式の射影安定群は有限なので
商の次元を得る。差分式と両階数での交叉検算は
\path{computations/verify_binary_apolar_rank_test.py}にある。
\end{proof}

% !TeX root = ../papers/nc-p3.tex
\section{二辺だけの二元商から対称性を仮定せず四窓を復元する}
\label{sec:binary-two-edge-rigidity}
\begin{theorem}\label{thm:binary-two-edge-rigidity}
標準有限Frobenius四窓の連続二辺に、両端の文字をそれぞれ三次元に割る二元形式の商があるとする。
中央の四文字は両方の商で三次二元形式の双対となる。両辺には純関係がないとする。
このとき中央に現れる二つのねじれ三次曲線は自動的に一致し、窓全体が
定理\ref{thm:binary-apolar-rank-test}の階数2の両対称非正則窓となる。
最初から二辺の二元座標の整合性、四周期の印、あるいは交換対称性を仮定する必要はない。
特に全辺に有限点グラフを持つFrobenius窓では、この隣接二辺の検査だけで、
交互の二元形式による指数成長商の存在が必要十分に判定できる（始点は動かす）。
そのような商を持つ窓は全て両対称で反対階数8、二元座標の一周の移動は恒等である。
\end{theorem}
\begin{proof}[中央の三次曲線が異なる場合の圧縮交叉]
各辺の六関係は端の一文字を割っても単射で、双対の積は端の三平面では二元乗法である。
三次交叉を両端で同時に圧縮しても単射となる。
実際三重形式を二つの三平面と中央文字で知れば、第一辺の関係の単射性から第一端全体へ、
次に第二辺の関係の単射性から最終端全体へ一意に戻る。
従って圧縮された二辺の三次交叉 $S$ は少なくとも四次元である。

中央の二曲線 $C,C'\subset\mathbb P^3$ が相異なると仮定し、交叉長を $m$ とする。
$m\ge6$ なら $C'$ の三つの二次式は $C$ 上で高々一次元の六次式空間に制限される。
従って二つの独立な二次式が両曲線を含む。非退化整曲線を含む二次式は既約なので、
次数4の完全交叉が次数6の和を含むことになり矛盾。よって $m\le5$。
$C$ の座標を $x$、$C'$ の座標を $y=[s:t]$ とする。$C'$ の標準座標で $C$ を
独立な三次形式 $a_0(x),\ldots,a_3(x)$ によって表示する。
$S^*$ は $P_5(x)\otimes P_2(y)$ を
\[
 F_i(x,y)=a_i(x)t-a_{i+1}(x)s\qquad (0\le i\le2)
\]
の $P_2(x)\otimes P_1(y)$ 倍で割ったものになる。これは左辺でまず二元乗法を取り、
右辺の乗法の核を対角式の倍で表示したものである。関係を入れる前後がともに18次元なので、
$\dim S$ は $\sum b_i F_i=0$ を満たす双次数 $(2,1)$ の三係数の空間の次元でもある。

$F_i=u_i s+v_i t$ とし、$C'$ の両座標端点を交叉から外して $u,v$ をともに基点自由な三次式の組にする。
$b_i=A_i s+B_i t$ と置けば $A\cdot u=B\cdot v=0$ で、残る条件は
$A\cdot v+B\cdot u=0$。最初の二つの解空間は各三次元である。
実際、各基点自由空間には互いに素な二つの三次式があり、その二次式倍だけで $P_5$ を張る。
$u\times v$ の成分は $C'$ の三つの二次式の制限であり、共通因子 $d$ はちょうど交叉を切り、次数 $m$。
両方に直交する二次式列は原始列 $(u\times v)/d$ の多項式倍だから、その次元は
$\max(3-n,0)$、ただし $n=6-m$。$A\cdot v,B\cdot u$ はともに $dP_{5-m}$ に入る。
これは $u$ のある成分が各点の局所環で単元で、$d$ を法として $v$ は $u$ の倍数であることから分かる。
従って $m\le2$ なら両写像は単射で最後の一致の空間は高々三次元。
$m\ge3$ なら各像は $n$ 次元の $dP_{5-m}$ に全射し、$\dim S=6-n=m$ となる。

よって残る $m=4,5$ で $S$ は交叉スキーム $D=C\cap C'$ 上の三文字評価の双対像に等しい。
実際この像は両辺のHankel条件を満たし、$P_2(x)P_3(x)=P_5(x)$ は $D$ に全射、
残る二次文字にも $D$ 全体で単元なものを取れるから次元は $m$ である。
\end{proof}
\begin{proof}[非被約交叉を含む二つの追加条件]
第一端の第四文字での積を $L:P_3(x)\to P_5(x)$ と書く。
$D$ 上の一点評価をこの文字まで戻し、第二端を圧縮したまま取ると、中央と第二端には
$L(-)(x)\otimes\operatorname{ev}_{y}^{(2)}$ という値が残る。
これが右辺のHankel関係に入る条件をその零化空間に評価する。
右辺の乗法の核は二組の二元変数の対角式と双次数 $(2,1)$ の全空間の積なので、
得られる制限写像の双対は $D$ 上
\[
                b(x)\otimes\operatorname{ev}^{(1)}_{y}
\]
という評価列を与える。ここで $b$ は中央の二次式に、$y$ で消える一次式を掛けてから
$L(-)(x)$ を評価した三成分である。各台点で $b$ は非零である。
消えれば $L(-)(x)$ が中央の $y$ における評価の倍数となり、
$C$ と $C'$ の同じ点であるため第一辺の関係に純テンソルを与えてしまう。
線束を $D$ 上で自明化し、$b$ の一般の一次結合を取れば全台で単元になる。
従ってこの追加条件の階数は $P_1(y)$ の $D$ への制限階数2以上である。
$D$ は $C'$ 上の長さ4または5の部分スキームだから、全てが一点に支持されても階数2である。
これで持ち上がる三次交叉の次元は $m-2\le3$ となり矛盾する。
\end{proof}
\begin{proof}[一致した印からの全窓の復元]
二つの曲線が一致したので共通の二元座標を選び、両端での第四文字の積を
$L,\widetilde L:P_3\to P_5$ と書く。両端を圧縮した交叉は今度は全指数和だけに依存する $P_7^*$ である。
各持上げ $L$ の三つの差分 $L(x^{j+1})-xL(x^j)$ を六次に斉次化し、
両持上げの合計六元の張る空間を $D_6\subset P_6$ とする。
片端を第四文字に戻した値が他端でHankelである条件から、持ち上がる $D\in P_7^*$ は
$P_1D_6$ 全体を消す。従って $\dim(P_1D_6)\le4$ である。
各持上げの差分は基点自由である。共通零点では $L(-)$ の評価が三次評価の倍数となり、
その点での $P_5$ の評価は元の関係の中で純テンソルになるからである。
$\dim D_6\ge3$ なら一次式倍で少なくとも一次元増えるので、可能なのは次元3から4への極小増加だけ。
このとき $D_6=hP_2$ となる（命題\ref{prop:binary-quotient-rank-two-curve}の
$U+xU$ の証明）。共通因子があるため不可能である。
よって $D_6$ は互いに素な二つの六次式の鉛筆で、両辺の三差分はそれぞれこの鉛筆を張る。

二元多項式環をこの正則列で割るとソクル次数10のArtin完全交叉を得る。
その最上次数の汎関数 $T$ は $D_6P_4$ を消す唯一の射影汎関数で、
五次までは関係がないので $G=H_{5,5}(T)$ は可逆である。
差分を $P_4$ 倍して $T$ に入れれば $GL,G\widetilde L$ はそれぞれHankelとなる。
$P_1D_6$ の零化空間は四次元で $T\mathbin{\lrcorner}P_3$ に一致するので、全ての
$T\mathbin{\lrcorner}q'$ が実際に三次交叉に持ち上がる。
両端を第四文字にして左右から計算すると、一方は $G(Lq,\widetilde Lq')$、
他方は $G(Lq',\widetilde Lq)$ となる。実際片端だけ第四文字にすると
$T(pq' Lq)=T(pq Lq')$ のHankel等式がその側の関係係数を一意に決める。
従って混合行列は対称である。前定理の差分計算から、独立な二行を持つ二つの八次パラメータは
極性空間を法として比例する。係数は非零なので、第四文字を選び直せば両持上げは一致する。
これで三次交叉が前節のGram窓の四つの縮約で張られた。
元の窓も第四文字との対合が完全なので、第四位置の基底変更だけでGram窓に一致する。
全四辺の一意の二元印は同じ二元座標で実現され、一周の移動も恒等となる。
圧縮交叉の有理数標本と長さ5の一点における検算は
\path{computations/verify_binary_two_edge_rigidity.py}で再現できる。
\end{proof}

% !TeX root = ../papers/nc-p3.tex
\section{二元障害の二方向が重なる十四次因子}
\label{sec:binary-reverse-orientation}
\begin{theorem}\label{thm:binary-reverse-orientation}
階数2の二元障害の行列の核を $k\ell\subset P_2$ とし、
$Y\mathbin{\lrcorner}\ell=T\mathbin{\lrcorner}c$ で四次式 $c$ を定める。
$Y$ の代表による $c$ の変更は $\ell P_2$ の加算である。
\begin{enumerate}
\item $\ell\mid c$ なら点グラフは正次元であり、逆も成立する。
\item $\gcd(\ell,c)=1$ なら有限窓の二元商は指定の向きにしか存在しない。
\item $\gcd(\ell,c)=d$ が一次なら有限窓の二元商は両方向に存在する。
固定した $T$ の $\mathbb P^5$ 上、この両向きの軌跡の閉包は被約既約な次数14の超曲面である。
\end{enumerate}
両向きのときは以下の線型変換が逆向きの十次・八次パラメータを計算する。
二組が異なる向きで提示された場合の添字保持同型は、一方をこの変換で戻して
定理\ref{thm:binary-apolar-rank-test}の同じ向きの軌道検査を行うことと同値である。
これら非正則候補の窓の最小周期は2で、両方向が許され逆向きのパラメータが同じ軌道に戻る場合に限り1となる。
両向き因子上でも一般最小周期は2である。
\end{theorem}
\begin{proof}[反対四次曲線と逆向きの三平面]
$T\mathbin{\lrcorner}P_4\subset P_6^*$ の零化空間は $I_6$ で、対応する写像は単射である。
従って行列の核条件がちょうど $c$ の存在と一意性を与える。
$\ell\mid c$ なら $Y$ は極性空間による変更後 $\ell P_6$ を消すので第二割線の射影にある。
逆に第二割線にあれば、動かした代表はある二次式の倍数を消し、その二次式は核の一意性から $\ell$ である。

$U=P_2\oplus kz$ 側の反対核は、三つの $P_2$ 座標の二次錐と、
$z\ell-c$ の二つの式を持つ。ここで $c$ は $\operatorname{Sym}^2(P_2)$ への任意の持上げである。
二式は反対積の式から消え、階数8なので核を尽くす。
錐の頂点以外を、母線の二元座標 $t$ から読むと、共通零点は
$[\ell(t)\operatorname{ev}^{(2)}_t:c(t)]$ の閉包と、共通根の上の母線とである。
$\ell,c$ が互いに素なら上記パラメータは基点を持たず、二次曲線への射影から生成的一対一である。
従って整な四次曲線を切り、逆向きの商に必要なねじれ三次曲線を含まない。

共通因子が一次なら $\ell=dl$, $c=dc_3$, $(l,c_3)=1$ とする。
二つの端を逆にして $dP_2\subset Q$ を新しい三平面、
$J:U\simeq P_3$, $J(p,a)=lp+a c_3$ を新しい全四文字とする。
新しい積の値を元の $P_5$ に移す線型写像 $I$ を
\[
         I(lh)=dh\quad(h\in P_4),\qquad I(c_3 r)=L(dr)\quad(r\in P_2)
\]
で定める。二空間の交わりは $lc_3P_1$ であり、
$Y\mathbin{\lrcorner}(dl)=T\mathbin{\lrcorner}(dc_3)$ から $L(dl v)=dc_3v$ ($v\in P_1$)。
従って定義は一致し全空間に及ぶ。像は $dP_4$ を含む。
もしこれだけなら $d$ の根での評価を $B$ に合成すると $dP_2\subset Q$ が核に入るため純関係となる。
有限点Frobenius窓の純関係の不在に反するので $I$ は可逆。
これで $B(u,dr)=I(J(u)r)$ となり逆の商が得られた。
新たなGram行列は $I^tGI$、第四文字 $q_0\notin dP_2$ の積は
$j\mapsto I^{-1} B(J^{-1}j,q_0)$ であり、その積にGram行列を掛ければ逆向きの八次式が戻る。
前定理の一意性により二方向を比較する同型と周期の判定が従う。
\end{proof}
\begin{proof}[次数と既約性]
$\ell$ は行列の二次小行列で与えられ、$c$ の係数は $Y$ と $\ell$ に双線型である。
従って $\operatorname{Res}(\ell,c)$ は $\bar Y$ の斉次次数 $4\cdot2+2\cdot3=14$ の式である。
代表を変えても終結式は不変である。階数1の集合は余次元2なので因子にはならない。
$[\ell]\in\mathbb P(P_2)$ を固定し $Y(\ell I_6)=0$ とすると、
$[\bar Y]$ の選択は射影三空間である。$c\bmod\ell P_2$ の二係数はその射影化前で全射の線型像である。
実際 $Y\mapsto Y\mathbin{\lrcorner}\ell$ は $P_8^*\to P_6^*$ の全射なので、
希望する $T\mathbin{\lrcorner}c$ を指定できる。階数2の選択はこの三空間内で稠密である
（二行が $P_2/k\ell$ の独立な二汎関数となればよい）。
$\ell$ の二つの相異なる根の一方を選びそこで $c$ が零となることは一つの線型条件で、
他方で零とする条件とは独立である。順序付きの根対を動かしたこの超平面束は既約である。
各ファイバーの階数1条件は二行二列の行列式の整な二次曲面なので、この超平面を含まない。
二重根の底でも $c$ の値が零となるのはなお真の超平面で、別成分を増やさない。
一般点ではちょうど一根の単純な消失であり次数に重複はない。

\path{computations/verify_binary_reverse_orientation.py}では $d=x$, $l=x-1$, $c_3=x^3+2$ の例で
両表示の積を直接照合する。新旧十次形式を通常の二項係数付き表示 $F$ とし、
$A=(F,F)_{10}$, $B=(F,F)_8$ に対する絶対不変量 $(B,B)_4/A^2$ を比較すると値が異なる。
従って既約な両向き軌跡の一般点は一周期でない。
\end{proof}

% !TeX root = ../papers/nc-p3.tex
\section{安定な非正則境界を通る全成分と非対称正則窓}
\label{sec:binary-component-crossing}
\begin{theorem}\label{thm:binary-component-crossing}
定理\ref{thm:binary-secant-obstructions}の有理標本および向きを固定した二元障害の一般点で、
二隣接切断階数6の全軌跡の軌道横断完備局所環は
\[
                  k[[x_1,\ldots,x_{13},h,a,b]]/(ah,bh)
\]
である（幾何点での剰余体を $k$ と書いた）。アフィンテンソルでは滑らかな軌道変数61個を加える。
二つの成分は滑らかで、添字保持基底変更を除いた次元は14と15、交叉は被約滑らかな13次元である。
前者は両対称部分で、後者の一般点はどちらの反対交換でも対称化できず、
有限被約点グラフを持つ最小周期4の標準quadratic Koszul左右AS正則窓である。
後者上の片側対称な二因子の一般点も同じ正則性と周期4を持つ。
一般の安定な二元障害点は、これらの正則な成分が交わるところで非正則に特殊化した点である。
\end{theorem}
\begin{proof}[四指標の法線と二次式]
標本の $\ell,c$ は互いに素で二行は独立なので、有限グラフでFrobeniusかつ非正則である。
定理\ref{thm:all-finite-windows-semistable}から安定で、そのアフィン軌道は61次元。
反対二交換の符号で周囲の線型空間を四指標 $++,+-,-+,--$ に分ける。
第一切断 $w=rGr^t$ で $N$ を $r^t$ の十次元核の基底列とする。
各指標の変分 $X$ に対し、両切断の一次条件は $N^tXN=0$ の一組と等しい。
二つのSchur行列を交換に関して和差に分ければ、各指標の式数は順に
$55,45,45,55$ である。実際一つの交換は二切断を交換し、両交換の積は行列を転置する。
同じ可逆な主六次小行列を両切断で選べばSchur補もこの作用と可換である。
標本での核次元は有理数上で正確に $45,16,16,0$ となる。
軌道接の各指標次元 $31,15,15,0$ を除くと、横断変数は自明指標14と両混合指標各1である。
軌道横断は線型補空間と基底変更の積で形式的に取れる。有限な二交換群に関し
軌道接の補空間も、消去する方程式と変数の組も同変に選べる。

陰関数定理で独立な一次式を解く。自明指標の式は全て消去される。
二つの混合指標の残る式は $af$ と $bg$ 各一本、最後の指標は $ab h_j$ の形である。
ここで $a,b$ が対応する混合変数で、$f,g,h_j$ は符号変換に不変。
$f,g$ の定数項は零で、両方とも自明指標方向での微分が非零となる。
この計算では一般逆行列 $D$ を同じ可逆主小行列の逆で取り、
自明指標接 $X$ と各混合接 $A$ に対して
\[
              N^t(XDA+ADX)N
\]
を対応する一次写像の一次元余核に射影する。両指標で非零であることを
\path{computations/verify_binary_component_crossing.py}が検算する。
全Jacobi階数を有理数と素数1009の両方で照合してから二次式を還元しているため、
剰余体での非零性は標数零の接・余核上での非零性を与える。
これは元のSchur形式の二次項であり、解いた一次変数からの補正は余核で消える。
従って $f,g$ はそれぞれ正則パラメータである。
\end{proof}
\begin{proof}[有限グラフの次元下界による高次式の決定]
近傍で基底変更を許した両対称部分は、横断上では $a=b=0$ である。
実際両対称なら反対位置の点座標を点ごとに移す射影線型写像が各側に存在する。
基点でその線型方程式の核は恒等写像の一次元だけである。
点座標の二次制限が単射なので、点上で比例する二つの線型列は全空間で比例するためである。
有限平坦な点代数上のこれらの方程式は近傍で書け、その核が存在すれば基点で選んだ小行列により
一意に正規化して可逆な同定を戻せる。両同定で基底を合わせれば文字通り交換不変である
（定理\ref{thm:finite-four-slot-rank}）。片側の交換についても同様で、
同定後のテンソルが比例するとき符号は基点の正符号のままである。
従って対称部分の飽和は局所的に固定部分に滑らかな対応軌道を加えたものとなる。
指標ごとに横断を取った上では、各片側の固定条件は対応する混合座標の消失である。

横断の16変数正則完備環において $(b,f)$ は高さ2の素イデアルで $a$ を含まない。
この素点は全方程式を満たし、両対称でない。ここを通る極小成分は
系\ref{cor:finite-component-dimension-bound}により横断次元15以上である。
この下界は形式環でも使える。同系の証明は有限平坦な辺のSchur恒等式を開集合上で消去するもので、
両対称部分の外で同じ高さの上界が働く。あるいは有限型体上の局所環の優秀性と
完備化の等次元性を使い、両対称成分の完備化が $a=b=0$ に含まれることを使ってもよい。
従って全方程式の極小素イデアルで $(b,f)$ に含まれるものは正則環で高さ1以下であり、
$af$ を含み $a$ を含まないため $(f)$ そのものである。
特に $bg$ と全ての $ab h_j$ が $f$ の倍数である。
$f$ は $a,b$ と互いに素で $g$ は正則パラメータなので、$g$ は $f$ と同伴、
全 $h_j$ は $f$ の倍数となる。これで全イデアルは $(af,bf)$ と正確に決まり、
$h=f$ と取れる。次元の異なる二成分と交叉の主張が従う。
\end{proof}
\begin{proof}[二元軌跡の一般点と各成分での全次数]
この交叉が二元障害を持つ一般点に延びることを示す。
任意の二元表示で $K=\ker B$ は、$P_2(x)\otimes P_3(y)$ の乗法核
$K_0=\Delta(x,y)(P_1(x)\otimes P_2(y))$ を含む。$\Delta$ は二変数組の対角式である。
従って $\bigwedge^2 K$ 内に
\[
 \Delta(x,y)\Delta(x',y')\Delta(x',y)\Delta(x,y')\Delta(y,y')
\]
という非零元がある。これは $x,x'$ に対称、$y,y'$ に交代するので、反対の混合指標の
$N^tXN$ の値を全て消す。従って常に少なくとも一方の混合余核が非零である。
上の局所模型の両対称部分では、これが起こるのは交叉 $h=0$ に限られ、そこでは両方とも非零となる。
同じ結論を近傍で代数的に読むには、Jacobi階数の小行列と
共通の二成分および交叉の滑らかさを用いればよい。両対称成分だけの滑らかな点では
追加の混合接は存在しないので、標本近傍の二元軌跡は交叉に含まれる。
階数および混合二次式の非退化性は、この交叉上の開条件である。
二元パラメータの既約性から幾何学的一般点で同じ局所模型を得る。

交叉は横断次元13である。有限両対称Frobenius窓の非正則部分は向きごとに次元12しかないので、
交叉の一般点は全次数で正則となる。従ってこの交叉を含む二成分、および15次元成分内で
$a=0$ または $b=0$ の各片側対称因子の幾何学的一般点も正則である。これは各内部次数の
相対bar複体を正則点から一般点へ移す既出の特殊化判定による。
被約有限グラフは二元パラメータの一般点ですでに成立するから全てのこれらの一般点にも成立する。

一般二元表示の十次形式の射影安定群は自明である（一般位置の十根の集合を保つ射影直線の変換は恒等）。
商の一意性から $Q$ 座標の作用も自明で、$P_2$ の作用をスカラーまで固定した後に
第四文字を $sz+p$ に変える作用は $(s-1)Y+T\mathbin{\lrcorner}p=0$ を満たす必要があり恒等である。
従って添字を保つ射影安定群は自明となる。この点で二移動の同定は標準の反転以外になく、
番号付き二十点には恒等として働く。近傍で二移動がある場合も同じ置換でなければならない。
実際有限エタールな番号付けの後、指定置換の座標変換は射影行列空間で固有に極限を取れる。
特異行列が二組の二十点を対応させることはない。核と像の射影空間に入る点数の和は
階数1または3なら $7+1$ 以下、階数2なら $3+3$ 以下だからである。
従って各極限は可逆で、非恒等置換は特殊点でも同じ置換を与えてしまう。
恒等置換による二移動は反対両配置を点ごとに射影的に同定し、両対称性を強制する。
両対称部分の外にある15次元成分およびその片側対称因子の一般点では二移動も一移動もなく、
四移動により最小周期4である。
\end{proof}

% !TeX root = ../papers/nc-p3.tex
\section{非対称成分の反対行列式と二元障害を切る残余因子}
\label{sec:binary-opposite-determinant}
\begin{theorem}\label{thm:binary-opposite-determinant}
前定理の15次元成分の、同じ有理標本と二元障害の一般点での軌道横断局所環を取る。
二つの片側対称因子の式を $a,b$ とすると、反対切断の行列式は単元倍を除き
\[
                         a^6b^6 D
\]
となる。$a,b,D$ は正則パラメータ系の一部で、$a=b=0$ において
$D=0$ は指定の向きの二元障害の閉包と局所的に一致する。
15次元成分の一般反対階数は16、二つの片側対称因子の一般反対階数はともに10、
残余因子の一般反対階数は15である。残余因子の一般点は有限被約グラフを持ち、
両方の反対交換対称性を欠く最小周期4の安定なFrobenius窓である。
この階数15の因子自身の一般点のAS正則性は、周囲の一般点の正則性からは決まらない。
\end{theorem}
\begin{proof}[因子の抽出と二元錐の一次関係]
局所模型の混合座標をそれぞれ指標 $+-$ と $-+$ の $a,b$ とし、
テンソルの各反対組を対称十成分と交代六成分に分ける。
同変な陰関数消去で戻した反対行列は
\[
           \begin{pmatrix} O&aP\\ bQ&abF\end{pmatrix}
\]
で、四ブロック $O,P,Q,F$ は二交換で不変である。交代六列から $a$、交代六行から $b$ をくくると
残余行列 $M=\left(\begin{smallmatrix}O&P\\Q&F\end{smallmatrix}\right)$ とその行列式 $D$ が出る。

指定の側が $U=P_2\oplus kz$ の二元点では、錐の二次式
$C=\Delta(x,x')^2\in\operatorname{Sym}^2 P_2$ が $O$ の左核に入る。
この式は $P$ の値も消す。実際 $P$ は $+-$ の接変分の対称・交代ブロックであり、
接条件からその変分は $K_0(x,y)\otimes K_0(x',y')$ と片側を交換した空間の双方を消す。
ここで $K_0=\Delta(x,y)(P_1(x)\otimes P_2(y))$ である。
$\bigwedge^2 P_3(y)$ の元は $\Delta(y,y')$ と対称な双次数 $(2,2)$ の式の積である。
二次元のPlücker恒等式
\[
 \Delta(x,x')\Delta(y,y')
 =\Delta(x,y)\Delta(x',y')-\Delta(x,y')\Delta(x',y)
\]
に残る $\Delta(x,x')$ と双次数 $(2,2)$ の式を掛ければ、
$C\otimes\bigwedge^2P_3(y)$ は上記の消える二空間の和に入る。
従って $M$ の左核に $(C,0)$ があり、二元障害の上で $D=0$ となる。
\end{proof}
\begin{proof}[交叉に沿った残余行列式の微分]
計算する成分と持上げを明示する。$A,B$ を二つの混合接の非軌道代表とし、
$X$ を $a=b=0$ の13次元交叉に接する対称方向とする。第一切断の一般逆を $J$、
核列を $N$ とし、$[V,W]=VJW+WJV$ と略記する。
$M$ の定数項では $O,P,Q$ はそれぞれ $w,A,B$ の反対切断の該当ブロック、
$F$ は混合項 $Z$ の交代・交代ブロックである。$Z$ は $--$ の空間で
$N^t ZN=N^t[A,B]N$ の唯一の解である。

次に方向 $X$ でこれらを微分する。$A_1,B_1$ は対応する混合指標で
$N^t A_1 N=N^t[X,A]N$, $N^t B_1N=N^t[X,B]N$ を解く。
$Z_1$ は $--$ の空間で、両辺に $N^t, N$ を掛けたとき次の式の右辺に一致する解である。
\[
 Z_1\ \equiv\ [X,Z]+[A,B_1]+[B,A_1]
          -\sum_{(U,V,W)\in\operatorname{Perm}(X,A,B)} UJVJW .
\]
これはSchur補の三次までの展開で、$X,A,B$ の三方向にそれぞれ一次の係数を取った式である。
$M$ の微分の四ブロックは $X,A_1,B_1,Z_1$ から取れる。
解の存在は前定理の滑らかな15次元枝で $X$ を交叉に接して取ったことによる。
混合の一次補正の解は自由な混合接を加えられるが、これは混合座標の一階の規格化と
文字基底の一階変更で実現できるので任意の一組でよい。

\path{computations/verify_binary_opposite_determinant.py}が有理標本で素数1009により次を検査する。
定数 $O$ は階数8、その左右二次元核に制限した $P,Q$ の階数は順に1と2、
$M$ の階数は15、ある $X$ に対する微分を $M$ の左右核線で縮約した値は非零である。
上の錐の恒等式が $M$ の階数の上界15を与える。
線型系の全Jacobi階数は前定理の有理数検証と一致し、
$X$ に課す二つの混合障害は標数零でも同じ非零の微分を与える。
従って各接と補正をこの良い素数での指定解に還元するよう持ち上げられ、非零微分は標数零でも成り立つ。
これで $D$ が $a=b=0$ 上でも正則パラメータとなる。
その零点は滑らかな12次元で、そこに入る既約な12次元の二元障害の閉包は局所的に全体と一致する。
微分の非消滅と階数の小行列は開条件なので、二元障害の一般点にも同じ結論が延びる。
\end{proof}
\begin{proof}[反対階数]
$a,b,D$ は独立だから三因子以外では反対階数16である。
残余因子の一般点は $ab\ne0$ で $M$ は階数15のままなので、実際の反対階数も15。
被約有限グラフと最小周期については、一般二元点で前定理の番号付けを取り、
その点を通る $ab\ne0$ の部分に同じ置換の特殊化を適用すればよい。
有限点グラフでの安定性も一般二元点から開条件で従う。

片側対称因子上では階数は10以下で、上記の $P,Q$ の非零制限からそれぞれ9以上、10以上である。
中央支持3以下の点の存在は射影的な入射から閉条件で、有理標本は非正則なのでこの条件を満たさない
（定理\ref{thm:finite-cubic-separators}）。従って両因子の一般点でもこの条件はない。
ところが定理\ref{thm:finite-rank-nine-all-marks}により、有限グラフで階数9の正則窓には
中央支持3の点がある。前定理の両因子の一般正則性から階数9が除かれ、いずれも10となる。
\end{proof}

% !TeX root = ../papers/nc-p3.tex
\section{十三次元の交叉を与える楕円曲線の非自明な二捩れ}
\label{sec:elliptic-two-torsion-crossing}
\begin{theorem}\label{thm:elliptic-two-torsion-crossing}
安定な二元障害の一般点を通る14次元成分と15次元成分の共通13次元部分は、
滑らかな楕円鉛筆 $(E,L,M,W)$ で $ML^{-1}$ が位数ちょうど2であるものの閉包と、
その点の近くで一致する。この滑らかな楕円曲線上の族は既約で、全て標準Koszul左右AS正則である。
有限二元障害の各向きの軌跡はその閉包に含まれる。特に共通部分の幾何学的一般点では
二つの反対鉛筆は同じ滑らかな楕円四次曲線を二捩れだけ異なる二束で埋め込んだものである。
\end{theorem}
\begin{proof}[四本の対角と二捩れ条件]
$E$ の原点を固定し、$\operatorname{Pic}^d(E)$ を因子の和で $E$ に同一視する。
$L,M$ に対応する点を $l,m$ とし、$U=H^0(L),Q=H^0(M)$ とする。
$K_0=\ker(U\otimes Q\to H^0(LM))$ は $E^2$ 上の対角で消える切断の空間である。
四変数を $x,x',y,y'$ とし、$\mathcal B=L\boxtimes L\boxtimes M\boxtimes M$ から
四つの交差対角 $x=y,x'=y,x=y',x'=y'$ を引いた線束を取る。
$K_0\otimes K_0$ と片側交換後の同じ空間との交わりは $H^0(\mathcal B)$
（四対角を引いた線束をここでは $\mathcal B$ と書く）である。
実際異なる素因子全てでの消滅はその和での消滅であり、二つの独立な組の対角を引くときは外積の切断空間となる。

この線束を $(x,x',y,y')=(z,z,-z,-z)$ の平行移動に制限すると次数は零である。
そのNéron--Severi形式は、$u=x+x'+y+y'$, $v=x-x'$, $w=y-y'$ として
$u^2+v^2+w^2$ である。実際四座標の各四次から四つの座標差の平方を引けばこの式となる。
従って上の楕円部分群を根基に持ち、平行移動しても制限の同型類は変わらず、
原点を通るものにおける類は $2(l-m)\in\operatorname{Pic}^0(E)$ である。
これが非零なら $H^0(\mathcal B)=0$ となる。

$\eta=m-l\ne0$ が二捩れとする。$b\in E[2]$ の各々に対し、
線束 $\mathcal B$ は次に同型である。
\[
 u^*\mathcal O_E([l-b])\otimes v^*\mathcal O_E([b])
                  \otimes w^*\mathcal O_E([b+\eta]).
\]
上の平方形式で同じNéron--Severi類を持ち、四座標軸への制限の次数2の類は
$l-b+b,l-b-b,l-b+b+\eta,l-b-b-\eta$ で、これは $l,l,m,m$ となる。
積多様体の $\operatorname{Pic}^0$ が各軸の積であることから同型が従う。
各点を切る切断を引き戻した積に四交差対角の切断を掛ける。
$b=0$ のとき $x=x'$ で一位の零、$y=y'$ の一般点で非零なので指標 $-+$ を持つ。
$b=\eta$ のときは逆の $+-$ となる。実際除数は各交換で保たれ、同一束二因子の
自然な交換作用に対する符号は固定対角での位数の偶奇である。
他の二つの $b$ は $++$ の相異なる除数を与える。
これで四つの切断は独立で全空間を張る。対角楕円部分群で割った後の
$(u,v,w)$ は行列式の絶対値2の同種写像で次数4、主偏極の積の引戻しの切断次元は4だからである。

$K=\ker(U\otimes Q\to H^0(LM)/W)$ は $K_0$ を含む。
二つの混合切断は $\bigwedge^2 K$ にあり、反対指標のSchur微分にそれぞれ非零の余核を与える。
これは任意の基点自由二次元鉛筆 $W$ で成り立つ。
非二捩れのときに消えると示したのは $K_0$ による四対角の切断であって、
個別の $W$ に由来する他の余核をこの議論で排除する必要はない。
\end{proof}
\begin{proof}[二元窓の基点自由鉛筆と二角形の平滑化]
一般二元点で $\ell$ は異なる二根 $a,b$ を持ち $c$ と互いに素である。
二本の射影直線を二節点で結んだ二角形 $E_0$ を取り、第一成分の節点を $a,b$ とする。
多重次数 $(4,0)$ の束 $L_0$ の貼合せを、主成分への制限空間が
$\ell P_2+kc\subset P_4$ になるように取る。この空間は $a,b$ での値の比を
$c(a),c(b)$ と同じにする四次式の全体である。多重次数 $(3,1)$ の $M_0$ は
主成分で全 $P_3$、$L_0M_0$ は全 $P_7$ を与える。残る成分での一次式は二端の値で一意に定まる。
写像
\[
 J:P_7\longrightarrow P_5,\qquad J(\ell h)=h,\quad J(cq)=L_Y(q)
\]
を取る。ここで $L_Y=G^{-1}H_{5,3}(Y)$ である。
二つの制限は交叉 $c\ell P_1$ で一致する。これは $Y\mathbin{\lrcorner}\ell=T\mathbin{\lrcorner}c$ による。
$J$ は全射で、積をその二次元核 $W_0$ で割ると元の二元辺になる。
$W_0$ の基点があれば、その点の評価は $J$ を経由する。従って二束の切断の積を
評価した純関係を与える。二束は全ての点で生成されるのでこの関係は非零となり、
有限Frobenius条件に反する。よって $W_0$ は二角形全体で基点自由である。

平滑化を具体的に取る。Legendre族 $y^2=x(x-1)(x-t)$ の原点の二重点を一回解消する。
全空間の二次初期式 $y^2+x^2-tx$ は非退化で、特別ファイバーの二つの接線は相異なる。
従って特別ファイバーは重複度一の主成分と例外成分が横断的に二回交わる上の二角形である。
無限遠の零切断 $O$ は主成分、一般ファイバーで位数2の切断 $P=(0,0)$ は例外成分の滑らかな点に着く。
$\mathcal L=\mathcal O(4O)$ と $\mathcal M=\mathcal L(P-O)$ は所望の多重次数で、
一般ファイバーでの差は非自明な二捩れである。
二角形上の指定された束対はこの特殊対に同型である。各束の貼合せの自由度は一つの単元であり、
両成分の二端を固定する二つの拡大縮小の作用の指数行列は、座標の向きを除き
$\left(\begin{smallmatrix}4&0\\3&1\end{smallmatrix}\right)$ となる。
従って代数閉体上で両単元を同時に合わせられる。

三束 $L_0,M_0,L_0M_0$ の一次コホモロジーは零である。双対束の切断は
正次数成分で消え、次数零の成分にも節点から零が伝わるためである。
基底変更定理で切断束と鉛筆 $W_0$ を持ち上げ、開条件により基点自由な一般鉛筆を得る。
二次辺は特殊点でも階数6を保つので族として延びる。
一般楕円鉛筆の最上テンソルの射影極限は特殊点の四次交叉にあり、その交叉は
一次元だから二元窓自身に一致する。この平滑化が二元パラメータの稠密開集合で働き、
既約性から軌跡全体の閉包含を得る。
\end{proof}
\begin{proof}[交叉成分の同定]
滑らかな楕円曲線と指定した非零二捩れ点の族は既約である。
実際指定点を $(0,0)$ とする $y^2=x(x-r)(x-s)$ の既約なパラメータで全てを覆える。
平行移動で次数4の $L$ を固定すれば $M=L\otimes\eta$ も決まり、その後の自動同型群は有限である。
$H^0(LM)$ の八次元空間の二次元鉛筆は12次元で、曲線の一変数と合わせて13次元となる。
反対鉛筆と積からのデータの復元により像も同じ次元である。
定理\ref{thm:elliptic-pencil-all-as}が全ファイバーの正則性を与える。
この軌跡は二元点の近傍で両対称であり、四対角の切断が混合接の跳躍を強制する。
前定理までの二成分模型では、そのような両対称点は13次元の共通部分に入る。
平滑化による閉包含、次元と既約性から局所的な一致を得る。
\path{computations/verify_elliptic_two_torsion_crossing.py}で四対角の形式、
解消の接錐、元の二元標本における $J$ と鉛筆を検算できる。
\end{proof}

% !TeX root = ../papers/nc-p3.tex
\section{一辺だけからの二元障害の伝播と指数成長}
\label{sec:single-binary-edge-obstruction}
\begin{theorem}\label{thm:single-binary-edge-obstruction}
標準有限Frobenius四窓の一辺に二元形式の商があるとする。
その辺の六関係は商のHankel六空間に同型に写り、三次形式の四文字の側で続く隣の辺には、
その有理正規三次曲線の点を因子とする純関係がないとする（逆向きにも同じ主張が成り立つ）。
このとき元の窓は指数成長し、標準AS正則にはならない。
特に全辺の点グラフが有限なFrobenius窓では、\emph{一辺だけ}の二元商の存在で非正則性が決まる。
長さ $4m+1$ のある始点で経路次元は $4\cdot3^m$ 以上で、長さ29ですでに標準次元を超える。
有限グラフの場合、もとの四窓は三次形式の側の反対二位置で必ず交換対称な表示を持つ。
ここでは残る三辺の商、正則性、交換対称性を事前に仮定しない。
\end{theorem}
\begin{proof}[ベクトル値二元逆系の分裂]
向きを $U,Q,V,W$ と書き、点座標側の $U$ に $P_2$ の三平面を取り、$Q=P_3$ とする。
関係の空間はこれらの双対に入る。指定の商の双対は $P_2\otimes P_3\to P_5$ の乗法である。
$U$ を $P_2$ に制限した四テンソルの値は $p,q$ に関して $pq$ だけに依存する。
続く辺の六関係を $R_1\subset P_3^*\otimes V^*$ とし、その零化空間 $B\subset P_3\otimes V$ を取る。
仮定により次数三の評価 $B\otimes\mathcal O\to V\otimes\mathcal O(3)$ は
$\mathbb P^1$ 上全射である。ある点 $t$ で全射でなければちょうど
$\operatorname{ev}_t\otimes v^*\in R_1$ ($v^*\ne0$) となるためである。
この十列から四行への全射の核束を $\bigoplus_{j=1}^6\mathcal O(-d_j)$ と分裂させる。
十列は独立なので $d_j\ge1$、次数から $\sum d_j=12$ である。

$n=3+j$ とし、$M_n\subset P_n^*\otimes V^*$ を $B P_j$ の零化空間とする。
$P_j$ の因子との全ての縮約が $R_1$ に入ることと同じである。
上の束列のコホモロジーから
\[
       M_{3+j}^*=H^1\!\left(\bigoplus\mathcal O(j-d_i)\right),
       \qquad \dim M_{3+j}=\sum_i\max(d_i-j-1,0).
\]
圧縮した四テンソルの第四文字は $M_5$ に単射で写る。
実際 $P_2$ との全縮約は $R_1$ に入り、最初の二位置の六関係を単射に圧縮したので
反対二位置の六関係も第四文字の全支持も失われない。
従って $\sum(d_i-3)_+\ge4$。六つの $d_i-1$ の和が6であるから、可能なのは
$(d_i)=(1,1,1,1,1,7)$ だけである。
この唯一の長い直和因子とSerre双対性により $M_8$ は一本のベクトル値八次汎関数
$\Phi\in P_8^*\otimes V^*$ で張られ、$3\le n\le8$ について
$P_{8-n}\to M_n$, $h\mapsto\Phi\mathbin{\lrcorner}h$ は同型である。
実際積による遷移は束列の接続写像と可換で、長い直和因子の双対では通常の二元乗法である。

これで第四位置は $M_5\simeq P_3$ と同定され、各 $v\in V$ を $D=\Phi(v)\in P_8^*$ として
圧縮した窓は $D(pqq')$ となる。$V\to P_8^*$ は単射である。核があれば
$R_1=M_3$ の一文字支持が落ちる。像を $S$ とする。
$R_1$ とその次の辺 $R_2$ は、$Q,S$ と $S,Q$ における同じ写像
\[
           (q,D)\longmapsto D\mathbin{\lrcorner}q\in P_5^*
\]
の双対像とその転置である。$R_1$ については $M_3$ の表示、$R_2$ については
対になる最初の辺の圧縮が単射であることによる。
$R_3$ を $P_3^*\otimes P_2^*$ へ圧縮した像も二元乗法の双対像になる。
実際、圧縮テンソルのその側の像は $R_1$ との元の完全対合から圧縮 $R_3$ の像全体で、
圧縮窓の式と $P_5\simeq M_3$ の単射性から六次元のHankel空間に等しい。
\end{proof}
\begin{proof}[切り離せる三文字ブロック]
$U$ 位置の生成文字を三つに割った経路代数を考える。
この位置と両隣の $P_3$ の間の二元乗法の双対関係を消すため、双対の経路空間では
各隣接組の対角式 $\Delta(u,q)$ を掛ける。内部の $u$ の次数2は両側の対角式でちょうど尽き、
各 $q$ では次数2が残る。$S$ 位置を挟む三文字に
$F\in P_2\otimes S\otimes P_2$ を置く。
両隣の対角式を掛けても中央の両辺の関係を消す十分条件は
\[
 \bar\gamma:P_2\otimes S\to P_6^*,\qquad (r,D)\longmapsto D\mathbin{\lrcorner}r
\]
をどちらの端から作用させても零となることである。実際、対角式の係数に現れる
二つの一次式を $P_5$ に掛けたものは $P_6$ を張る。
両端について対称な $F$ の空間は $6\cdot4=24$ 次元で、一方からの収縮先は
$7\cdot3=21$ 次元、他方の条件は対称性から従う。よって少なくとも三つの独立な $F$ がある。

長さ $4m+1$ の $u,q,S,q',u,\ldots,q,S,q',u$ という並びで $m$ 個のブロックを
独立に選び、全対角式を掛ける。両端の $u$ には各々まだ任意の一次式を掛けられる。
異なるブロックは変数を共有せず、固定した対角式の積による乗法も単射なので、
次元は少なくとも $4\cdot3^m$。$m=7$ で $8748>\binom{32}3=4960$ となる。
元の窓はこの経路代数に全射するので同じ下界を持つ。
\end{proof}
\begin{proof}[有限グラフでの対称性]
有限グラフの全関係に純テンソルがないので定理の二つの仮定は自動的である。
$Q,S,Q$ の両辺が互いに転置であることから、その四次元の三次交叉は両端の交換で保たれる。
$P_2\subset U$ の三文字による四次窓の縮約は端の交換で不変で互いに独立。
残りの一文字を選び直せば、全四文字が不変か、第四文字だけ符号が変わるかのいずれかである。
後者なら $Q$ 両端の交換と $U$ でのこの射影反射、$S$ での恒等の組が窓を保つ。
これは四辺の点列を逆順にするが、$S$ 位置への射影は点列の閉埋込みで、そこで各点を
スキームとして固定するため、$U$ 位置の配置も点ごとに固定しなければならない。
配置への二次式制限は単射であり、全ての点で恒等と比例する行列はスカラーなので
この反射は不可能である。従って元の窓自身で $Q$ の二位置は交換対称となる。
\path{computations/verify_single_binary_edge_obstruction.py}は分裂次数、対角式の両係数と成長下界を検算する。
\end{proof}

% !TeX root = ../papers/nc-p3.tex
\section{単一二元辺を持つ全有限窓の直線複体と階数層}
\label{sec:binary-linear-complex}
\begin{theorem}\label{thm:binary-linear-complex}
一辺の二元商に向きと印を付けた有限点Frobenius窓は次の一つの既約なパラメータ空間の開部分で尽くされる。
$L:P_3\to P_5$ と $[\omega]\in\mathbb P(\bigwedge^2 P_3)$ に対し
\[
 F_L(q\wedge r)=qL(r)-rL(q)\in P_8,\qquad F_L(\omega)=0
\]
を課す。$F_L$ の階数は5とし、$S=(\operatorname{im}F_L)^\perp\subset P_8^*$ とする。
窓を $Q,U,Q,S$ 順に、$Q=P_3$, $U=P_2\oplus k$ として
\[
  w(q,(p,a),q',D)=D\big((pq+aL(q))q'\big)
\]
で与える。これに有限点グラフと有限Frobenius検査を課すだけで全許容値となり、全て非正則である。
第四文字を加える自由度 $L\mapsto sL+M_p$ と同時の二元座標変換を除く。
ここで $M_p$ は $p$ による乗法である。

非退化な $\omega$ なら窓の反対階数は10で、この層は稠密である。
$\omega=q\wedge r$ が分解するとき、$e=\deg\gcd(q,r)$ とする。
$e=0$ は有限グラフを持たず、$e=1$ は反対階数9、$e=2$ は反対階数8である。
後二層とも有限Frobenius窓が実在し、$e=2$ はこの向きの両対称二元窓に一致する。
従って非正則な有限窓の中に、階数10で安定・最小周期4・二十点全て被約の
既約な13次元の添字保持同型類の族がある。この族の一般点で二元商の印は一意に戻る。
各有限非両対称点では、商を取る位置は四位置のうちただ一つで、最小周期は常に4である。
\end{theorem}
\begin{proof}[全パラメータと既約性]
前定理で $Q$ の両位置は同定され、第四文字による縮約は $Q$ の交換で不変である。
二元乗法を $L$ で拡張すれば、その不変性は $D$ が $F_L$ の像を消すこととなる。
実際、三次交叉 $Q,U,Q$ は $P_2$ への制限で $P_8^*$ に入り、両辺のHankel条件から
四文字目の値まで一意に決まる。従ってこの交叉はちょうど $(\operatorname{im}F_L)^\perp$ で、
Frobenius条件により階数5である。逆に上の式は両辺で指定の関係を実現し、
開な有限Frobenius条件とグラフの有限性を通常の有限行列・入射検査で決定できる。
$L$ の第四文字の変更と二元座標を除けば、印を保つ同型条件も必要十分である。
これは圧縮辺の階数一五次曲線が二つの座標を共通の変換まで決め、$S$ は三次交叉で内在的に決まるためである。

非退化な $\omega$ では $L\mapsto F_L(\omega)$ は $P_8$ に全射である。
実際非退化な四行交代行列で単項式基底を変換した四つの三次式と任意の五次式の積を足す写像である。
従って $\mathbb P^5$ の開部分上で、入射は階数15のベクトル束である。
分解するときの像は $qP_5+rP_5=dP_{8-e}$ ($d=\gcd(q,r)$) で、核は $15+e$ 次元。
$[d]$ と基点自由な二平面 $\langle q/d,r/d\rangle\subset P_{3-e}$ は計 $4-e$ 次元を動く。
よって $e=0,1,2$ の入射は各19次元である。全入射は29次元の滑らかな空間で九式により定まり、
各成分は20次元以上なので、これらは全て既約な非退化部分の閉包に入る。

$e=0$ では $L(q)=cq,L(r)=cr$ ($c\in P_2$) となり $(-c,1)\in U$ の評価核が二次元以上となる。
$e=2$ なら $\langle q,r\rangle=dP_1$ で、$L(dv)=cv$ ($v\in P_1$, $c\in P_4$)。
従って前定理群で用いた三差分は $d$ の係数による定数関係を持つ。
純関係がないためこれらは基点自由であり、六次の鉛筆を張る。
その完全交叉の十次汎関数を取り $GL$ がHankelになることから
定理\ref{thm:binary-apolar-rank-test}の両対称窓に戻る。
逆に全両対称二元点では、核二次式 $d$ に対する同じ式から $\omega$ はこの分解型である。
$e=1$ の有限例は定理\ref{thm:binary-reverse-orientation}の $d=d_1d_2$, $c=d_2c_3$ の有限例から得る。
指定の二平面を $d_1\langle g_t,h_t\rangle$ として
$(g_0,h_0)=d_2(s,t')$ から互いに素な二次式対へ動かす（$s,t'$ は二元変数）。
$L_t(d_1g_t)=c_3g_t$, $L_t(d_1h_t)=c_3h_t$ とし、補二文字の値も延ばせば所望の $L_0$ に特殊化する。
階数5、有限性、Frobenius性は特殊点に成立する開条件である。
稠密な非退化 $\omega$ の側でもこれらが成り立ち、一般二元点からの開性により点は一般に被約である。
\end{proof}
\begin{proof}[射影直線上の二平面束と反対階数]
$v(t),\lambda(t)\in Q^*$ をそれぞれ $q(t),L(q)(t)$ の評価とする。
これらは純関係の不在により各点で独立で、生成する二平面は
$\mathcal E=\mathcal O(-3)\oplus\mathcal O(-5)\subset Q^*\otimes\mathcal O$ を与える。
そのPlücker座標系 $H=\operatorname{im}F_L$ は $\mathbb P^1$ を $\mathbb P^4$ へ埋め込む。
実際長さ2を分離できなければ $S=H^\perp$ に長さ2以下の因子に支持される八次汎関数がある。
その $H_{5,3}$ の核は二次元以上となり、$Q,S$ 辺の有限性に反する。

反対切断を $\operatorname{Sym}^2 Q$ からの写像とする。その核の二次式 $A$ の
対角制限 $f=A(v,v)\in P_6$ は $fP_2\subset H$ を満たす。
$f\ne0$ ならその長さ6の零因子で $H$ の像は高々二次元で、その埋込み像は直線に含まれる。
Plücker二次式はこの直線全体で零である。Grassmann多様体の直線は固定した平面内の直線の鉛筆だから、
長さ6の零因子に沿い $v(t)$ がその固定平面に入ってしまう。三次正規曲線の平面切断は長さ3で矛盾する。
よって $A(v,v)=0$。第四文字の条件は $A(v,\lambda)\in H$ で、
これは適当な交代行列を $A$ に足した行列 $B$ で $B(v,\lambda)=0$ となることに等しい。
同時に $B(v,v)=0$。$B(v,-)$ は任意の三次係数の四成分を一意に表せるので、
反対核の次元は $h^0(\mathcal E^\perp(3))-1$ である。引いた一次元は交代行列 $\omega$ の核線である。

非退化なら $\omega$ で $\mathcal E^\perp\simeq\mathcal E$ だから反対核は零である。
ここで反対核は対称二乗十次元内で取っている。
分解型なら $\omega(v,-)/d$ は次数 $3-e$ の零点を持たない切断列となる。
従って $\mathcal E^\perp$ は次数 $-(3-e)$ の線部分束を持ち、商の次数は $-(5+e)$。
この拡大は $\mathbb P^1$ 上で分裂し、反対核の次元は $e$ となる。これで全階数を得る。
\end{proof}
\begin{proof}[印を忘れた像と周期]
非退化部分の20次元から第四文字の自由度4と $\operatorname{PGL}_2$ の3を除き、一般の印付き像は
高々13次元である（一般の二元四次成分を持つ $[\omega]$ の射影安定群は有限）。
その閉包は12次元の両対称二元軌跡を含み、稠密部の反対階数は10であるから、
印を忘れた像の次元も13である。一般の印の一意性には命題\ref{prop:binary-quotient-rank-two-curve}と
定理\ref{thm:binary-incidence-codimension}の残余平面六次を使える。
\path{computations/verify_binary_linear_complex.py}の非退化標本ではその六次が滑らかである。
階数一Hankel点 $h(t)$ が残余六次の一般点の核になることはない。
その場合の二次式は $t$ で消え、残る一次因子は追加行と二つの条件を満たす。
純関係がなければその一次因子はあるとしても一意である
（二条件がともに恒等零なら追加行は $t$ の評価の倍数になる）。
従ってこの型の曲線は有理で、滑らかな平面六次を支配できない。
よって残余成分はその六次と双有理で、関係六平面の唯一の有理正規五次曲線が指定の印を復元する。

両対称でない窓では、別の向きで全三次文字を $U$ または $S$ とする商があれば、
前定理によりもう一方の反対交換も対称となってしまう。
同じ $Q$ を全三次文字、$S$ を圧縮側にする商も、既存の $U$ を圧縮する辺と
$Q$ を挟んで隣接するため定理\ref{thm:binary-two-edge-rigidity}に反する。
従って圧縮する位置は一意であり、二移動では保持されない。四移動と合わせて最小周期を得る。
非正則有限窓は定理\ref{thm:all-finite-windows-semistable}から安定である。
\end{proof}

% !TeX root = ../papers/nc-p3.tex
\section{安定な非正則軌跡の既知十五次元成分への大域的同定}
\label{sec:binary-global-component}
\begin{theorem}\label{thm:binary-global-component}
定理\ref{thm:binary-component-crossing}の15次元成分は
定理\ref{thm:full-fifteen-component}の既約成分と同一である。
非自明な二捩れ楕円鉛筆の13次元軌跡と、定理\ref{thm:binary-linear-complex}の
有限非正則窓の全てがこの成分に属する。両軌跡の一般点は異なり、前者は両対称で正則、
後者は片側対称で非正則である。

この成分の一般の単一二元辺を持つ非正則窓では全窓の軌跡は滑らかで、
片側対称因子への横断変数を $a$ とすると反対行列式は単元倍の $a^6$ である。
従って前節までの残余階数15因子はこの片側対称な非正則13次元軌跡を含まない。
一般二元点の近くで二つの軌跡が片側対称因子上に切る13次元部分も互いに異なる。
\end{theorem}
\begin{proof}[一つの印からの成分の照合]
線型複体の非退化部分において、印を選ぶ有限被約開集合は空でない。
以下、滑らかな判別三次を持つ印が一般の窓に存在することを示す。三次形式を
$q=1-x+x^3$ とし、$L(q)=(2+3x+x^2)q$ となる有理標本を
\path{computations/verify_binary_global_component.py}で取る。
指定点で六つの双一次式のアフィンJacobi行列は可逆なので、
この印は辺と共に局所エタールに動き、入射の像はパラメータ空間の開集合を含む。
$D\in S$ 側の座標は $H_{5,3}(D)q=0$ の唯一の核線で、標本での非零小行列により近傍にも戻せる。
一般に有限Frobeniusであることは前定理による。従って印の開条件の検査を、この有理標本で行えばよい。

$U,Q,S,Q$ 順に回転して印を先頭の座標とし、反対評価の二行列を単位に戻す。
この片側対称表示では共通の非退化対称計量で $G,F$ は自己随伴である。
直交基底で $P=Q$ かつ $G,F$ が対称の表示にできるが、判別三次を求めるだけなら
計量の平方根を取る必要はない。$[G_i,F_j]$ の三次元線型包の三座標を取れば、
対応する三行三列一次式行列の行列式が同じ平面曲線を与える。
素数1009で三つの偏微分は共通射影零点を持たない。
従って標数零の印付き開部分にも滑らかな判別三次があり、
定理\ref{thm:marked-smooth-discriminant-components}からその窓は既知の15次元成分に入る。

同じ標本で $Q$ の反対交換の二指標の接次元の上界は $60,16$ である。
一方、一般の単一二元窓は有限かつ非両対称なので全成分のアフィン次元が76以上である。
よって一般点は滑らかな76次元で、成分は一つしかない。
この点は二元境界から得た15次元成分にも入る。実際既約な単一二元辺の軌跡は
一般両対称二元点を含み、そこでの二つの局所成分のうち、両対称でない一般点を受けられるのは
15次元のものだけである。以上で二成分の大域的な一致を得る。
全ての許容有限パラメータも、既約な入射からの連続なテンソルの構成により同成分に入る。
\end{proof}
\begin{proof}[片側対称因子の横断]
一つの交換の正指標の160変数には一切断の100本のSchur式だけが課されるので、
正指標の接次元は少なくとも60であり、一般点ではちょうど60である。
軌道の正負指標次元はそれぞれ46と15である。従って軌道横断の15変数で
固定部分は滑らかな14次元、負指標は一変数 $a$ のみである。
交換による同定は有限グラフ上で唯一なので、前節の局所議論と同じく基底変更による飽和もこの固定条件となる。
反対切断で交代六行の $a$ をくくった行列の、$a=0$ での可逆性を調べる。
正十行は階数10であるから、残る六次元核から交代六行への負接による写像を取ればよい。
\path{computations/verify_binary_linear_complex.py}の標本でこの階数は6である。
負接の一般次元16は上の下界と標本のJacobi階数の双方で確認できるため、
核とその変分を良い素数から標数零へ持ち上げられる。可逆性は既約な族の稠密開集合で成り立つ。
従ってこの点には片側対称因子以外の反対行列式の局所成分は来ない。
\end{proof}

% !TeX root = ../papers/nc-p3.tex
\section{階数15の反対切断の核線と添字保持同型}
\label{sec:rank-fifteen-opposite-kernels}
\begin{theorem}\label{thm:rank-fifteen-opposite-kernels}
定理\ref{thm:binary-opposite-determinant}の残余因子の一般点で、
反対切断の左右の核線はどちらも非退化な四元双線型形式で生成される。
\end{theorem}
\begin{proof}
同定理の有理標本で、滑らかな15次元枝の残余因子内に形式曲線を取る。
二交換のそれぞれで符号を変える座標を $a,b$ とし、その一次方向を同定理の $A,B$ に選ぶ。
$a=b=t$ とおき、残りの不変座標を $t^2$ の級数に取れる。
実際、交叉上でも $D$ の微分は非零なので、一つの不変座標を $D=0$ で消去できる。
曲線の二次までの不変座標は、15次元枝の式と $D=0$ の微分を満たす範囲で選べる。
ここで $a,b$ は文字基底の軌道も含め同変な滑らかな座標を取り直してよい。

$w(t)$ の一次項を $A+B$、二次項を $X+Z$、三次項を $A_3+B_3$ とし、
四次項の交代・交代部分を $Z_4$ と書く。
$X,Z,A_3,B_3$ の指標は順に $++,--,+-,-+$ である。
対向する二組をそれぞれ対称・交代に分けた基底で、実際の反対切断は
$\operatorname{diag}(I_{10},tI_6)M(t)\operatorname{diag}(I_{10},tI_6)$ である。
$M(t)=M_0+t^2M_2+O(t^4)$ の四ブロックは、順に
$(w,A,B,Z)$ と $(X,A_3,B_3,Z_4)$ の反対切断の対応部分から取る。

係数の決め方を記す。第一隣接切断の核列を $N$、一般逆を $J$、
増分を $E$ とし、Schur方程式
\[
 N^t(E-EJE+EJEJE-EJEJEJE+\cdots)N=0
\]
の指定次数を取る。もう一つの隣接切断も課すには、
二変数 $a,b$ の各指標の項を別々に足す。これは一方の交換が移す切断にも
同じ式を課すことに等しい。二次で $X,Z$ を解き、三次の混合障害が消えるよう
$X$ に対称接を加える。定理\ref{thm:binary-opposite-determinant}の滑らかな枝の式により、
二つの三次の障害は同じ一条件である。$A_3,B_3$ の自由接成分は基底と混合座標の
三次補正で選べ、$Z_4$ は第四次の $--$ 方程式から一意に決まる。
最後に $M_0$ の左右の核生成元 $l,r$ に対し $l^tM_2r=0$ を課す。
$X$ に交叉に接する対称方向を加えたときのこの式の微分は前定理の非零微分である。
これで残余因子内の曲線の必要な係数が決まる。以降の係数は上の陰関数消去で延長できる。

\path{computations/verify_rank_fifteen_opposite_kernels.py} はこれらの線型系を素数1009で
解き、全方程式への代入も検査する。$l=(C,0)$ の対称形式 $C$ は第四行・列が零で
左上 $3\times3$ ブロック $C_3$ が可逆であり、$r$ の交代部分は階数4である。
従って実際の右核の生成元を $t$ 倍した極限は可逆な交代形式である。
左側では $l_2^t M_0=-l^tM_2$ を解くと、実際の核形式は
$C+tK+t^2L+O(t^3)$ となる。$K,L$ は $l_2$ の交代・対称部分で、
$K_{44}=0$ である。左上ブロックに関するSchur補の $t^2$ の係数は
\[
 L_{44}-K_{4,1:3}C_3^{-1}K_{1:3,4}\equiv 22\pmod {1009}.
\]
$l_2$ に $l$ を加える自由度はこの係数を変えない。
前定理で標数零とこの素数で隣接Jacobi階数、混合障害の階数および
残余因子の横断性が一致しており、全ての線型系は同じ可逆小行列で持ち上がる。
$M_0$ の階数15の上界・錐の式も標数零で既知である。従って左右の非零小行列は
標数零でも非零である。形式曲線の一般点での非消滅と開性から主張が従う。
\end{proof}

\begin{proposition}\label{prop:opposite-adjoint-equivalence}
任意の反対階数15の四窓テンソルで、核双線型形式が両方とも非退化とする。
四文字空間を順に $A,B,A^*,B^*$ と同定すると、この部分の添字保持同型は
射影線型同型 $[f]\colon\mathfrak{sl}(A)\longrightarrow\mathfrak{sl}(B)$ の
独立な左右の内部共役にちょうど一致する。一添字移動の式は、トレース随伴を
$f^\dagger$、双対空間に値を取る転置を $\tau$ として $[\tau_A f^\dagger]$ であり、
二添字移動は $[\tau_B f\tau_A^{-1}]$ である。
\end{proposition}
\begin{proof}
それぞれの核線を標準評価に取る。その縮約が零だからテンソルは
$\mathfrak{sl}(A)\otimes\mathfrak{sl}(B)$ に属し、非退化なトレース対で
線型写像 $f$ を与える。反対階数15がその可逆性に等しい。
標準評価の線を保つ基底変更は一方の文字の基底変更と反対文字での逆転置のスカラー倍に限られる。
これが $f\mapsto \lambda\operatorname{Ad}_h f\operatorname{Ad}_g^{-1}$ を与え、
逆に全てのそのような変換が実現する。テンソル $\sum_i P_i\otimes Q_i$ は
一添字移動で $\sum_i Q_i\otimes P_i^t$ となるので移動の式が従う。
従って各移動と元の写像の内部共役二重軌道の一致が、その移動が周期となる必要十分条件である。
窓の隣接階数やFrobenius条件には、$f$ から戻した同一テンソルを用いる。
\end{proof}

四文字に体積形式を選ぶと、反対切断の余因子行列の四乗を
各文字の交代体積で縮約して次数60の $\mathrm{SL}_4^4$ 不変式を得る。
反対階数15では余因子が核生成元の単純積 $x\otimes y$ であるから、この値は
非零定数倍を除き $\det(x)\det(y)$ である。従って命題の開部分はグラフの次元によらず半安定である。
半安定商における閉軌道代表でもこの不変式は非零であり、階数15以下は閉条件であるため、
同じ正規化でその代表を調べられる。

% !TeX root = ../papers/nc-p3.tex
\section{両向き二元点での四次接錐と残余階数15因子の全次数正則性}
\label{sec:double-binary-residual-regularity}
\begin{theorem}\label{thm:double-binary-residual-regularity}
定理\ref{thm:binary-opposite-determinant}の残余階数15因子の幾何学的一般点は
標準14641型quadratic Koszul左右AS正則である。向きを一つ移した因子にも同じ結論が成り立つ。
これらの一般窓は安定で、二十個の被約点を持ち、最小周期4である。
さらに、各向きの因子を二つの片側対称因子で切ったうち、以下で指定する側の交叉の
幾何学的一般点も全次数で正則となり、反対階数は10である。
\end{theorem}
\begin{proof}[二つの二元軌跡の横断面]
定理\ref{thm:binary-reverse-orientation}の $d=x$, $l=x-1$, $c_3=x^3+2$ の有理点を取る。
ここでは二行は独立で $\gcd(\ell,c)=d$ なので、グラフは有限で窓は安定Frobeniusである。
両向きを $B_1,B_2$ とする。各指標の隣接Jacobi核次元は有理数上で $45,16,16,0$、
両混合指標の二次障害はともに非零である。従って定理\ref{thm:binary-component-crossing}の
高さによる証明がこの点にも適用できる。滑らかな15次元枝上に $a,b$ を取り、その両対称交叉
$H=(a=b=0)$ は滑らかな13次元である。以下の次元では文字基底の軌道を除く。

各 $B_i$ の被約な芽がこの点で滑らかなことを確かめる。
向きを固定した二元点では、反対切断の一方の核は、ねじれ三次曲線 $C$ を含む二次式の鉛筆である。
その完全交叉は $C$ と残余直線 $s$ の和である。実際二つの二次式は共通因子を持たず、
三次曲線が二重に現れる次数はない。残余次数1で、完全交叉は純次元である。
この中の直線は $s$ 一つだけで、直線の入射はこのファイバーで非分岐である。
$s$ の一般点で完全交叉は滑らかなので、
$N_{s/\mathbb P^3}\to\mathcal O_s(2)^2$ は一般点で単射、従って大域切断の核が零だからである。
階数8が一定な二元軌跡上で鉛筆を束とし、その中の直線をGrassmann多様体上で取る。
入射は固有で、標本の近くで有限である。特殊ファイバーが被約な一点なので、
中山の補題で関数環は底の代数の商となり、被約な二元軌跡を覆うことからその上で同型である。
残余三次曲線も代数的に回復する。直線の二つの一次式に関して鉛筆を二行二列の一次式で書き、
その行列式を鉛筆に加えると残余曲線のイデアルになる（直線との残余完全交叉の式）。
従って滑らかなねじれ三次曲線の族を得る。曲線上の射影座標はエタール局所で選べる。
その二次像の線型包と階数1の形式の根基から三平面を戻す操作は
定理\ref{thm:binary-apolar-rank-test}の証明の定階数の線型演算である。
そこで積の二元座標と第四文字を選べば、Gram行列と八次形式も同じく代数的に戻る。
従って滑らかな二元パラメータと文字枠による表示は、これらの基底選択の自由度を除いて局所逆を持つ。
特に各被約像は滑らかな12次元である。この議論は他方の向きにも使える。

元の表示の八次形式に $(3,1,4,1,5,9,2,6,5)$ を加える一次方向を $u$ とし、
逆表示の八次形式で同じ変更をして元の基底に戻した方向を $v$ とする。
十次・八次形式の微分に二文字の軌道微分を加えた二空間の階数はそれぞれ43、
その和は44である。また第一空間に $v$、第二空間に $u$ を一本加えるといずれも44となる。
小行列は素数1009で非零であり、上界は上の次元と両対称軌道の次元31である。
よって $B_1,B_2$ は $H$ 内で互いに横断し、それぞれに含まれる選んだ二方向が
その共通部分への横断を与える。以下、$u,v$ をこの接基底の座標にも用いる。
この二方向を残した $H$ の滑らかな二次元切片では $B_1$ の接線が $v=0$、
$B_2$ の接線が $u=0$ である。混合方向 $a,b$ を加えた同変な四次元切片を取る。
\end{proof}
\begin{proof}[余階数二の残余行列]
前節と同じく交代六列と六行から $a,b$ をくくって得た行列を $M$ とする。
その四ブロックは二交換で不変である。
\path{computations/verify_double_binary_transversal.py}は $M_0=M(0)$ の階数が
有理数上で14であること、$u,v$ に沿う微分がそれぞれ
$\ker M_0\to\operatorname{coker}M_0$ 上で零になることを有理数上で検算する。
計算式は定理\ref{thm:binary-opposite-determinant}の四ブロック微分である。
$a,b$ に沿う一次微分も、二交換の符号により零である。
従って可逆な14次小行列で消去した二行二列のSchur補は四次元切片上で二次から始まり、
$D=\det M$ の消失次数は4以上である。

良い素数1009で四次斉次項は単元倍を除いて
\[
               (664v^2+546b^2)(691u^2+570a^2)
\]
となる。計算を説明する。$a,b$ 両方の速度を非零に取り、不変方向の速度を
$X=uX_1+vX_2$ とおく。隣接Schur方程式の二次を指標ごとに解き、
対称な二次補正を選んで三次の混合障害を消す。三次の混合解と四次の $--$ 解までを
用いて $M=M_0+tM_1+t^2M_2+O(t^3)$ を計算する。
行列をくくるときには $a,b$ の各速度も割り、定数行列と核基底を全方向で共通に取る。
$L,R$ を左・右核の基底列、$J_0$ を $M_0$ の一般逆として、
Schur補の二次項は $L^t(M_2-M_1J_0M_1)R$ である。
\path{computations/verify_double_binary_residual_quartic.py}はその左上成分が零、
右上・左下成分が表示の両二次式となることを検査する。
二交換で不変な二次式の単項式は $u^2,uv,v^2,a^2,b^2$ の五つだけであり、
コードはこれらを分離する方向で偏極する。符号空間の線型補正は
滑らかな枝の同変な陰関数消去で選べる。混合接の自由項は基底変更と混合座標の選択にあたり、
行列式の最低項には単元倍しか与えない。自由な対称二次補正についても
$H$ の接空間全体で左右核への微分がこの素数で零と検算しているので還元値は変わらない。
全ての消去に単元小行列を用い、標数零での三次以下の消失は上記の有理数計算で確認済みである。
従って標数零の最低次斉次項は四次で、その還元が表示式となる。
\end{proof}
\begin{proof}[両対称正則点への特殊化]
15次元枝上、$D$ の局所既約因子 $F$ の幾何学的一般点が非正則であると仮定する。
有限Frobenius開部分であるから、その因子内に全次数で正則な点が一つでもあれば、
各内部次数でのbar複体の特殊化により一般点も全次数で正則となる。
従って $F$ の $H$ への制限の零点は非正則な点だけからなる。
$D|_H$ は上の四次項から恒等零でない。有限な両対称Frobenius窓の
非正則なものは二方向の二元軌跡だけである。各 $B_i$ の一般点では
定理\ref{thm:binary-opposite-determinant}により $D|_H$ は単純に消える。
従って $F|_H$ は単元倍を除き $B_1,B_2$ の一方の局所式、または両者の積である。
この議論は完備局所環でもよい。各 $B_i$ は滑らかで、全次数の特殊化は
局所環の基底変換後にも各次数の有限複体で行える。

二交換は $H$ を点ごとに固定する。$F|_H$ に現れる $B_i$ に沿う消失が単純だから、
同じ $B_i$ を含む別の局所因子はなく、二交換は $(F)$ を保つ。
生成元は二交換で平均して不変にできる。実際 $H$ 上で元の非零な制限が変わらないため、
各単元倍率は閉点で1である。四次元切片に制限すると、
$B_i$ が一つなら $F$ の最低項は $u$ または $v$ の定数倍である。
二つなら最低項は二次で、不変かつ $a=b=0$ への制限が $uv$ の非零定数倍となる。

これらはいずれも四次項の因子にならない。標数1009の代数閉包で
表示した四次式は相異なる四つの一次式の積となる。$b$ の符号は $v,b$ の二因子を交換し、
$a$ の符号は $u,a$ の二因子を交換する。従って不変な一次因子はなく、
不変な二次因子の $a=b=0$ への制限は $u^2$ または $v^2$ の倍数である。
標数零で仮定した最低項の因数分解は、係数体上で1009を延長した付値を取って、
各因子をGauss付値零に規格化すればこの素数にも還元できる。
四次式全体の $a=b=0$ への制限がこの素数でも非零なので、どの因子の制限も
還元で零にならず、上で求めた $u,v,uv$ のいずれかがそのまま残る。これは矛盾である。
従ってこの両向き点を通る残余因子はいずれも正則な両対称窓に特殊化でき、
幾何学的一般点が全次数で正則である。各向きの一般二元点で一意な残余因子は
その二元軌跡の閉包を含むので、この点にも来る。既知の有限被約性・安定性・周期と合わせて主張を得る。

同様に $a=0$ に制限すると最低項は $u^2(664v^2+546b^2)$ の定数倍である。
$B_1$ を含む因子で非正則性を仮定すれば、不変な最低項の制限は $v$ または $uv$ の倍数を要する。
$b$ の符号を保つ因子でこれは不可能である。従って $B_1$ から来る残余因子のこの片側切片も
一般に正則である。$b=0$ では $B_2$ について同じである。
これらの一般点は両対称でなく、二元点で欠けていた中央支持3の点も閉条件により一般点では欠ける。
有限片側対称で階数9以下の正則窓の判定
（系\ref{cor:all-defective-one-sym-finite}と定理\ref{thm:finite-rank-nine-all-marks}）から
階数10を得る。残余因子上の任意の特殊値の正則性はここでは一般点への特殊化からは従わない。
\end{proof}

% !TeX root = ../papers/nc-p3.tex
\section{点グラフの次元を問わない一辺の二元障害}
\label{sec:binary-edge-basepoints}
\begin{theorem}\label{thm:binary-edge-basepoints}
標準有限Frobenius四窓で、一辺の六関係を一方の一文字を落として
$P_2^*\otimes P_3^*$ に射影した像が、$P_5^*$ の乗法双対像に等しいとする。
この窓は標準Koszul左右AS正則にはならない。点グラフの次元、純関係の有無、
他の三辺の表示は任意でよい。順序を逆にした商についても同じである。
\end{theorem}
\begin{proof}[欠損を持つ評価層]
定理\ref{thm:single-binary-edge-obstruction}の記号を用い、指定辺を $U,Q$、
次の二文字を $V,W$ とする。$Q=P_3$ で、$P_2\subset U$ への圧縮を施す。
指定の辺で六関係の射影が同型なので、圧縮しても対向する辺 $R_2$ の像と階数6は保たれる。
特に圧縮窓でも $V,W$ の各支持は四次元、最初の二位置を合わせた $P_5^*$ も全支持である。

$B=R_1^\perp\subset P_3\otimes V$ の十列の評価を取り、その像層を $E$、余核を $C$ とする。
像は射影直線上の束で、核は $K=\bigoplus_{i=1}^{10-r}\mathcal O(-d_i)$、
$r=\operatorname{rank}E$ と書く。定数列が独立なので $d_i\ge1$、
$\sum_i d_i=\deg E\le3r$ である。最後の上界は $E\hookrightarrow V\otimes\mathcal O(3)$
の最高外冪から従う。$E$ は十列で生成されるので $j\ge0$ では $H^1(E(j))=0$ である。
$M_{3+j}=(BP_j)^\perp$ に対して二つの束列より
\[
 0\longrightarrow H^1(K(j))\longrightarrow M_{3+j}^*
       \longrightarrow H^0(C(j))\longrightarrow0
\]
を得る。右端からの双対包含は評価と可換である。特に $j=2$ のとき圧縮窓は
$W\hookrightarrow M_5\subset P_5^*\otimes V^*$ を定める。

$10-r\le 3r$ から $r\ge3$。$r=3$ では核は七項、$7\le\deg E\le9$ なので全て $d_i\le3$、
$H^1(K(2))=0$ である。$C$ の捩れ長を $t$ とすると自由商は
$\mathcal O(12-\deg E-t)$ である。この次数は $V\otimes\mathcal O(3)$ からの全射により3以上である。
従って
\[
   h^0(C(-3))=t+12-\deg E-t-3+1=10-\deg E\le3.
\]
捩れ層の一次コホモロジーが零なのでこの式は拡大によらない。
$V\to H^0(C(-3))$ に核があり、その文字と $P_5$ の任意の積は
$H^0(C(2))=M_5^*$ で零になる。これは圧縮窓の $V$ が全支持であることに反する。
\end{proof}
\begin{proof}[有限長の欠損]
残る $r=4$ では $C$ は捩れ長 $t=12-\deg E\le6$ であり、
\[
 \dim M_5=t+\sum_{i=1}^6(d_i-3)_+\ \ge4,\qquad
 \sum_{i=1}^6(d_i-1)=6-t .
\]
$t=0$ は全射の場合で、定理\ref{thm:single-binary-edge-obstruction}が指数成長を与える。
$1\le t\le3$ なら上式から唯一の分裂は $(1,1,1,1,1,7-t)$ で、$\dim M_5=4$ である。
従って $W=M_5$。欠損の台の一点 $s$ と非零の剰余体商 $C\to k_s$ を取ると
$H^0(C(2))^*\subset M_5$ に $\operatorname{ev}_s^{(5)}\otimes v_s$ という非零元がある。
これは $P_5\otimes V\to C(2)|_s\to k$ の積評価である。
この元に対応する第四文字で元の四テンソルを縮約する。圧縮後の三テンソルの末尾支持は $kv_s$ である。
圧縮前にも末尾支持は同じである。実際最初の辺の六関係から $P_5^*$ への射影が同型なので、
この三テンソルは $\operatorname{ev}_s^{(5)}$ の一意な持上げと $v_s$ の単純積となる。
定理\ref{thm:all-endpoint-ranks-two}に反する。

$4\le t\le6$ なら $d_i\le3$ だから $M_5=H^0(C(2))^*$ である。
$C$ の零化イデアルが定める有効因子を $Z$ とすれば $\deg Z\le t$ である。
もし $\deg Z\le5$ なら $Z$ 上で消える非零の五次式が全ての $V,W$ の組を消し、
圧縮の $P_5$ 側の階数6に反する。従って $t=\deg Z=6$。
局所の離散付値環で零化指数は基本因子の長さの最大値、層の長さはそれらの和であるから、
各点で基本因子は一つである。従って $C$ は $Z$ 上可逆である。
直線束を $Z$ 上で自明化して $D=H^0(\mathcal O_Z)$ を取ると
$P_5\to D$ は同型である。$D$ は有限有効因子の環なので、
完全積対合を与える汎関数 $\lambda$ があり $M_5\simeq D^*\simeq D$ となる。
$V\to D$ の像を $v\mapsto\bar v$、$W\hookrightarrow D^*\simeq D$ を $w\mapsto\bar w$ とすると、
圧縮窓は
\[
                (p,q,v,w)\longmapsto \lambda(\bar p\bar q\bar v\bar w)
\]
である。$Z$ の閉点 $s$ での評価は、辺 $R_1$ に
$\operatorname{ev}_s^{(3)}\otimes v_s$ を与える。これは $R_1=M_3$ に対する上の右端の双対包含から従う。
また $P_5\simeq D$ と完全対合により、同じ点での積評価
$v_s\otimes w_s$ が $R_2$ に入る。$v_s$ は層の生成性から非零で、
$w_s$ も非零である。零なら $\bar V\bar W$ の積がその極大イデアルに入り、
$P_5$ に対する切断階数が6にならない。こうして中央文字 $v_s$ を共有する
二つの純関係を\emph{元の}二辺に得る。定理\ref{thm:no-consecutive-pure-relations}で指数成長する。
分裂次数の列挙は\path{computations/verify_binary_edge_basepoints.py}で再検算できる。
\end{proof}

% !TeX root = ../papers/nc-p3.tex
\section{二元Hankel空間への非全射な射影と残る二つの捩れ長}
\label{sec:deficient-binary-support}
\begin{theorem}\label{thm:deficient-binary-support}
一文字の直線を落とした一辺の像が $P_5^*\subset P_2^*\otimes P_3^*$ の
真の部分空間 $H$ となる標準正則窓を仮定する。$s=\dim H$ は4か5である。
前節の評価層は一般階数4で、その余核の長さ $t$ は2または3に限られる。
十列の核束は $\mathcal O(-1)^5\oplus\mathcal O(t-7)$、第四文字は $W=M_5$ である。
$H$ を消す各五次式 $f$ について、十列の評価行列 $P$ は二次行列 $A_f$ と
$P A_f=fI_4$ を満たす。五つの一次核列で自明十束を割った束を $J$ とすると、
$H^\perp$ から $\operatorname{Hom}(J,\mathcal O(t-2))$ への単射がある。
特に長さ2の場合、$J$ の自明な直和因子が $6-s$ 個以上必要である。
\end{theorem}
\begin{proof}[圧縮で保持される支持]
落とす文字を $z$ とし、元の辺の射影核を $z\otimes S$ とする。
$\dim S=6-s$ が3以上なら $R_1\cap(S\otimes V^*)\ne0$ であり、その非零元に
$z$ を付けて三次交叉の一次元端支持を得る。従って $s\ge4$。
同じ理由で任意の $s$ について $R_1\cap(S\otimes V^*)=0$ を要する。
圧縮によって $W$ の一文字が消えれば元の対応する三テンソルの $U$ 側の端支持は一次元、
$V$ の一文字が消えれば三テンソルの中央の $U$ 支持が一次元になる。
後者では同じ $z$ を共有する連続純関係ができる。よって $V,W$ とも圧縮後も四次元支持である。
これだけで前節の一般階数3の排除と $\dim M_5\ge4$ の議論が再び使える。

$t=0$ だと $W=M_5$ は一つの八次逆系 $\Phi$ の三次極性空間になり、
$P_5\to M_3$, $f\mapsto\Phi\mathbin{\lrcorner}f$ は同型である。
従って圧縮窓 $\Phi(v)(pqq')$ の $P_5$ 側の階数は6で、真の部分空間ではない。
$1\le t\le3$ では分裂と $W=M_5$ は前節の次数計算で従う。
このとき $f\in H^\perp$ は $fV\subset BP_2$ と同値である。各文字の原像を取って
四本の二次列を並べれば $P A_f=fI_4$ である。
$C$ の零化因子を $Z$ とすると $f$ は当然 $Z$ 上で消える。

核の中の一次直和因子の和 $K_1$ は $K(1)$ の大域評価で内在的に決まる。
$J=\mathcal O^{10}/K_1$ とし、$\mathcal O^{10}$ 上で
$D_f=fI_{10}-A_fP$ を考える。$P D_f=0$ なので像は $K(5)$ に入り、
$K_1$ 上での作用は $fI$ だから商では零である。
従って $D_f$ は $J\to(K/K_1)(5)=\mathcal O(t-2)$ を誘導する。
$K/K_1\hookrightarrow J$ への制限が $f$ の乗法なので、$A_f$ を線型に選んで単射を得る。
$J$ は大域生成だから $t=1$ ではそのような非零写像はない。
$t=2$ の自明因子の下界も束の分裂から従う。
\end{proof}
\begin{proof}[有限長部分の左右同時縮約]
残りの $4\le t\le6$ では $M_5=H^0(C(2))^*$ である。
$C$ の台上の線束を自明化し、零化因子の環を $D$ と書く。
$V$ は $C$ を $D$ 加群として生成し、$W\subset C^*$ は四つの汎関数を与える。
$DW$ が消す $C$ の部分加群を割り、有効な加群 $N$ を取る。
このとき $W$ は $N^*$ を $D$ 加群として生成する。
$N$ の長さは $t$ 以下で、零化因子の次数を $e$ とすれば $e\le6$ である。
五次式はその因子の環へ全射する。圧縮の $P_5$ 側の階数 $s$ は $e$ 以下である。
よって $e\ge4$ となる。

$N$ を各閉点における離散付値環の基本因子に分解する。
各点で最大長の基本因子がすべて二つ以上あれば、長さは最大長の和の二倍、
すなわち $2e$ 以上になり矛盾する。従ってある点 $x$ で最大長 $m$ の因子が一つだけである。
他点で零、$x$ で局所パラメータの $m-1$ 乗となる元を五次式で表す。
この元が $N$ に与える乗法は階数1であり、その像はsocle、その余像は
剰余体商 $N\to k_x$ である。$V$ に制限した余像汎関数 $v_x$ は生成性から非零である。
$W$ の全てがこのsocle上で消えれば $DW$ もそこで消えるので、
その評価による汎関数 $w_x$ も非零である。従ってこの五次式で圧縮窓を縮約すると
$0\ne v_x\otimes w_x\in R_2$ を得る。ここでは圧縮後の $R_2$ の像は元の辺に含まれれば十分である。
一方 $C\twoheadrightarrow N\to k_x$ と $M_3$ への双対包含から
$\operatorname{ev}_x^{(3)}\otimes v_x\in R_1$ である。元の二辺での連続純関係となり排除される。
長さの組合せを\path{computations/verify_deficient_binary_support.py}でも検算する。
\end{proof}

\begin{corollary}\label{cor:no-full-pure-projected-curve}
任意の標準正則な辺で一方の四文字を直線で割ると、射影された関係空間の階数一部分は、
両射影がそれぞれ射影平面と射影三空間を線型に張る整曲線を含まない。
元の辺に純関係があっても同じである。
\end{corollary}
\begin{proof}
その曲線の関数体で両側の座標空間の積は少なくとも $3+4-1=6$ 次元である。
辺は六次元しかないので射影は単射で積次元は最小値になる。
命題\ref{prop:binary-quotient-rank-two-curve}の付値と一次元増大の証明は
この等号だけから両空間を $P_2,P_3$ に同定する。従って前節の定理に反する。
\end{proof}

% !TeX root = ../papers/nc-p3.tex
\section{非全射・非被約を含む一辺の二元射影の全面排除}
\label{sec:deficient-binary-quintic-pencil}
\begin{theorem}\label{thm:deficient-binary-quintic-pencil}
標準Koszul左右AS正則な四窓では、一文字を直線で割ったとき、射影された一辺の関係空間が
$P_5^*\subset P_2^*\otimes P_3^*$ の乗法双対像に含まれる基底は存在しない。
像の次元に条件はなく、辺の向きはどちらでもよい。点グラフの次元・被約性も任意である。
\end{theorem}
\begin{proof}[捩れ評価における二つの端支持]
像が六次元の場合は定理\ref{thm:binary-edge-basepoints}による。
真の部分空間 $H$ の場合は定理\ref{thm:deficient-binary-support}により
$s=\dim H$ は4か5、隣接評価の捩れ長 $t$ は2か3で $W=M_5$ である。
任意の捩れの台の点 $x$ における余核の剰余体商から、第四文字 $l_x\in W$ を得る。
その三テンソルは圧縮後 $\operatorname{ev}_x^{(5)}\otimes v_x$ に等しく、
圧縮前の左端支持は $D_x=\langle z,\operatorname{ev}_x^{(2)}\rangle$ に入る。
正則窓での端支持下界から丁度 $D_x$ である。同一点で二つの独立な剰余体商があれば
同じ二次元端支持を与えるので、定理\ref{thm:all-endpoint-ranks-two}に反する。
従って余核は次数 $t$ の有効因子 $T$ 上可逆である。

$s=5$ とすると、指定の辺の射影核 $z\otimes S$ は一次元である。
$w(l_x)$ は圧縮後の評価の一つの持上げと $v_x$ の積に、
この射影核と一つの末尾ベクトルの積を加えたものである。よって右端支持も高々二次元となる。
両端とも丁度二次元なら $l_x$ が前点・後続点の双方の直線を持ち、
定理\ref{thm:no-two-sided-centers}に反する。これで階数5が除かれる。
\end{proof}
\begin{proof}[二次式による両端の同時切断]
$s=4$ とする。$T$ 上で束を自明化し、$T^*$ を $H^0(\mathcal O_T)^*$ の
各次数の評価逆系と書く。前節の二つの束列から $T^*$ が $M_n$ に入り、
その上の多項式による縮約は $T$ の環の乗法の双対である。
$\dim(W/T^*)=4-t$ である。また $H\subset P_5^*$ は通常のスカラーの評価空間
$H^0(\mathcal O_T(5))^*$ を含む。実際 $V$ の文字で $T$ 上いたるところ単元となる
一つを選べる。各閉点で剰余値が非零の条件は真の超平面の補であり、文字は層を生成するからである。
この文字と $W$ 内の $T^*$ の全てとで圧縮窓を評価すると、$T$ 上の全五次汎関数を得る。
五次式の $T$ への制限は全射なのでこの部分は $t$ 次元である。

$T$ の長さ2の部分因子 $Z$ を選び、それを定める二次式 $p$ を第一文字の点として取る。
$T$ 上で $p$ の乗法の階数は $t-2$ である。各局所因子で零点の位数を数えればよく、
二重点・三重点でも同じである。$R_0$ の $p$ による評価は
$H\to P_3^*$, $h\mapsto p\mathbin{\lrcorner}h$ である。
$T$ 部分からの階数は高々 $t-2$、他の部分は $s-t$ 次元なので合計は $s-2=2$ 以下となる。
他方、$R_3$ の同じ文字での評価階数は三テンソル $w(p)$ の $W$ 側の支持次元で、
$W=M_5\to M_3$, $l\mapsto p\mathbin{\lrcorner}l$ の階数に等しい。
ここでも $T^*$ 部分の階数は高々 $t-2$、それ以外は $4-t$ 次元だから合計は2以下である。
$z$ を落とした側に $p$ を取ったため、どちらも元の窓の評価である。
両方とも2以上でなければならず、またしても両側の点ファイバーが直線となって矛盾する。
以上で全ての階数と捩れ型を排除した。
\end{proof}
\path{computations/verify_deficient_binary_quintic_pencil.py}は長さ2・3の全基本因子型で、
二次式による両評価の階数上界を局所乗法行列で検算する。

% !TeX root = ../papers/nc-p3.tex
\section{任意の正則辺における純関係の正次元成分}
\label{sec:pure-relation-curve-supports}
\begin{theorem}\label{thm:pure-relation-curve-supports}
標準正則な辺 $R\subset U\otimes V$ の純関係の被約軌跡を考える。
正次元なら次の二型に限る。
\begin{enumerate}
\item 一方の因子が固定された直線があり、そのような直線は各向き高々一本である。
正次元部分はこれらの直線の和である。
\item 二つの射影直線間の射影同型のグラフである滑らかな平面二次曲線が一本ある。
その外の純関係は高々三点であり、他の正次元成分はない。
\end{enumerate}
特に純関係の曲面は存在しない。これは元の辺の中の分解可能テンソルについての制限で、
辺の点対応（関係の零点）とは区別する。
\end{theorem}
\begin{proof}[部分空間に支持される関係数]
$D\subset U$ に対して $k_D=\dim(R\cap(D\otimes V))$ とする。
$\dim D=1,2,3$ に対して上界は順に $2,3,4$ である。
一次元の場合は定理\ref{thm:branching-centers-curve}の最初の評価、
三次元の場合は法線での端支持の下界2による。
二次元で $k_D\ge4$ なら $D^\perp$ の全点で評価階数は高々2、従って常に2である。
しかし定理\ref{thm:plane-branch-pencils}によれば、この分岐直線への制限に残る関係は
二次元以下でなく丁度三次元である。矛盾によって上界3を得る。右支持にも同じ上界がある。

純関係の整曲線の両支持の次元を $a,b$、そのテンソルの包の次元を $c$ とする。
関数体の積空間の付値から $c\ge a+b-1$ である。両方が4なら $c\ge7$ は不可能である。
一方の支持が真なら $c\le4$、二次元支持なら $c\le3$ である。
従って両方が動く場合は $a=b=2,c=3$ しかない。固定二平面のSegre二次曲面の平面切断で、
整かつ両因子が動くため滑らかな二次曲線である。片方が固定の場合は $k_D\le2$ により直線である。
二次曲線の三空間 $P$ は、両端の二平面 $A,B$ を同定すれば対称行列全体である。
\end{proof}
\begin{proof}[複数の曲線と孤立した純関係]
同じ向きの異なる二直線の第一因子が同じなら $k_D\le2$ に反し、異なるなら
その二因子を張る二平面に四次元の関係が支持される。従って同じ向きには一本しかない。
二次曲線と純直線の共存も不可能である。固定因子が $A$ 内なら三空間 $P$ の外の関係を
$A$ 内に生じ、外なら $A$ とその固定因子を張る三平面に五次元の関係を生じる。
右側でも同様である。

二つの二次曲線の左支持を $A,A'$ とする。同じなら両者とも同じ三空間 $P$ に入り、
そこにある二次曲線は一つである。交わりが一次元なら二曲線の包の交叉は高々一次元である。
実際、各包のうちその固定第一因子に支持されるのは対応する一点の直線のみである。
よって $A+A'$ の三平面に五次元以上の関係が入り矛盾する。
交わりが零なら $R$ 全体が両三空間の直和となる。$A^\perp$ のどの非零点で評価しても
$A'$ の対称二行列の評価だけが残り、その像は常に第二の曲線の同じ右二平面である。
これは二次元端支持を定める点の一意性に反する。以上で曲線の個数が決まり、
曲面なら無限個の異なる整曲線を含むので除かれる。

二次曲線がある場合に $U$ を $A$ で割った像 $R'$ は三次元である。
二次曲線の外の純点 $x\otimes y$ は $x\notin A,y\notin B$ を満たす。
例えば $y\in B$ なら $B$ に支持された全関係は $P$ に等しいためである。
$\bar x\otimes y\in R'$ はその純点を一意に決める。持上げの差は $A\otimes ky$ と
$P\subset A\otimes B$ の双方に入るからである。
$R'$ の要素を $U/A$ の双対の非零点で評価する四行三列の一次式鉛筆を取る。
全点で階数は2以上である。一般階数3なら三次小行列の共通零点は高々三点で、
階数2の各点にただ一つの核線があり、これが $R'$ の対応する純関係である。
一般階数2なら $A^\perp$ は分岐中心の直線になる。定理\ref{thm:plane-branch-pencils}の
三関係表示で、$R'$ の純関係は全て同じ第二因子 $b$ を持つ。
実際第三関係の二つの係数は $b$ を法として独立である。
従って二つの異なる純点が持ち上がれば $R$ に $b$ を持つ純直線を生じてしまう。
この場合に持ち上がるのは高々一点である。
\end{proof}
\path{computations/verify_pure_relation_curve_supports.py}は両支持の数値条件と三次鉛筆の次数を検算する。

% !TeX root = ../papers/nc-p3.tex
\section{曲線グラフでの長方形切断と四つの純関係の座標枠}
\label{sec:curve-graph-pure-frames}
\begin{theorem}\label{thm:curve-graph-pure-frames}
四辺の全射影ファイバーが空または一点の標準有限Frobenius窓で、共通グラフ $X$ が一次元とする。
一辺の点座標でベクトル次元 $a,b$ の任意の二部分空間へ制限した関係の階数 $r$ は
\[
        r\ge a+b-3+\mathbf1_{a\ne b}\qquad(1\le a,b\le4)
\]
を満たす。各辺の相異なる純関係は高々四点で、どの四点以下についても両側の文字は
それぞれ線型独立である。元代数の全次数の正則性は仮定しない。
\end{theorem}
\begin{proof}[制限した辺の両次数]
$\mathbb P^{a-1}\times\mathbb P^{b-1}$ 上で $r$ 個の双一次式を取る。
交叉は共通グラフの部分スキームなので次元1以下であり、Segre超平面の次元定理から
$r\ge a+b-3$ である。等号で積の次元が1以上なら、純一次元の完全交叉となる。
両射影への次数は $(H_1+H_2)^{a+b-3}$ の交叉数から
$\binom{a+b-3}{a-2},\binom{a+b-3}{a-1}$ である（範囲外は零）。
しかし各整曲線は $X$ の整成分であり、定理\ref{thm:point-graph-curve-degree}により
両次数は等しい。重複度を掛けて加えても等しいので $a=b$ を要する。

一方の文字の $d$ 次元真部分空間に支持される関係数は、
$(a,b)=(4-d,4)$ を用いると高々 $d$ となる。
特に固定因子を共有する二つの純関係はない。
相異なる純関係 $u_i\otimes v_i$ のうち四個以下で $u_i$ が従属とし、極小従属部分を取る。
その包の次元は個数より一つ小さいから、その純関係自身にも従属が必要である。
しかし $u_i$ の唯一の従属は全係数が非零であり、テンソルの従属を各 $v_i$ 側の汎関数で
評価すると、全 $v_i$ が比例することになる。固定因子の上界に反する。右側も同じである。
五つの相異なる純関係があれば、辺の各点対はそれぞれの第一因子または第二因子を消す。
四つを同時に消せないので、五添字を二個と三個に分け、一方を左、残りを右で消す有限個の
割当ての和集合に点対応が入る。各割当ての場所は直線と一点の積またはその逆である。
ここにグラフの整曲線は入れず、一次元性に反する。
\end{proof}

\begin{proposition}\label{prop:curve-graph-four-pure-frame}
四つの純関係がある場合、両文字の基底を取ればそれらは $u_i\otimes v_i$ ($1\le i\le4$) である。
二つの残余関係を対角項なしの二行列で表示する。
一次元の純曲線部分は被約な座標直線だけからなり、
各座標対 $I$ の左の直線を補対 $I^c$ の右の直線に射影同型で対応させる。
その直線が現れる必要十分条件は、両行列の $I\times I^c$ の二行二列ブロックが
一つの非退化行列を張ることである。
現れる左座標対のグラフは、一本の辺、交わる二辺、離れた二辺、三辺の道、または四辺の閉路のいずれかである。
対応するグラフの有限長部分（孤立部分と埋込み部分）の長さは順に $14,10,8,6,4$ となる。
\end{proposition}
\begin{proof}
四つの対角関係が消える場所は、座標集合を左右に分けた直積たちの和である。
各整曲線は両側に非定数なので、一点を含む分割は使えず、二対二の分割に入る。
その正方形は $\mathbb P^1\times\mathbb P^1$ で、残る二式はともに双次数 $(1,1)$。
グラフ成分を切る条件はちょうど表示した階数条件である。二式とも零または
因数分解する曲線全体が含まれることは一点ファイバー条件に反する。
この曲線の一般点では二つずつの非零座標が単元なので、四つの対角式が残り四座標を
スキームとして零にする。従って重複度は1で、有限台の部分を除いた純曲線は被約である。

ある頂点に三辺が入れば、その左座標点の像は三つの相補直線の共通点でなければならず、そのような点はない。
右側でも同様である。$n$ 本の左辺での頂点次数を $d_i$ とすれば相補グラフの頂点次数は $n-d_i$。
全て $d_i,n-d_i\le2$ なので非空グラフは列挙した五型に限られる。
交わる座標直線は異なる接方向を持つ通常二重点である。よって純曲線の算術種数は順に
$0,0,-1,0,1$ であり、定理\ref{thm:pure-curve-defect}の $18-4n+2g$ を代入して長さを得る。
この列挙は正則窓における各型の実在を主張するものではない。
\end{proof}
\path{computations/verify_curve_graph_pure_frames.py}は両次数の等号と相補座標グラフの全列挙を検算する。

% !TeX root = ../papers/nc-p3.tex
\section{曲線グラフにおけるただ一つの単純な純関係と直線への還元}
\label{sec:curve-graph-single-pure}
\begin{theorem}\label{thm:curve-graph-single-pure}
四辺の全射影ファイバーが空または一点で、点グラフが一次元の標準AS正則窓では、
各辺の関係空間とSegre多様体のスキーム論的交叉は空か被約な一点である。
もし $R_0$ に $u\otimes v$ があれば、共通点グラフの全ての一次元被約成分は全位置で直線となり、
位置0では全て $u^\perp$ の平面に、位置1では全て $v^\perp$ の平面に入る。
\end{theorem}
\begin{proof}[長さ三の無限小純関係]
$R_0$ の純関係の長さ3の部分スキームがあれば、その線型包 $P$ はベクトル次元3である。
二次元以下なら長さ3のスキームが乗る射影直線上で全二次小行列が恒等零となり、
純関係の直線が生じて前定理に反する。両文字の射影の包は各々高々3で、
前定理の部分支持上界から各々丁度3である。第一の包を消す文字 $p$ を取ると $R_0(p)$ は
残る三関係の像であり、点ファイバー条件により階数3である。従って $p$ は点グラフ上にある。
その前点を $y$ とする。$w(y)$ の末尾縮約が張る三関係の空間は $R_0$ の
$p$ で消える部分に入る。両者とも三次元で後者は $P$ だから等しい。
従って $w(y)$ は三つの支持がいずれも三次元で、非零の一文字縮約に純テンソルがある。
定理\ref{thm:ternary-point-subalgebras}により非正則となる。以上は任意の辺に使える。
\end{proof}
\begin{proof}[長さ二の純関係と四次評価鉛筆]
長さ2の部分スキームの包を $P$ とする。次元は2、両支持も2である。
$P$ の補空間を $E\subset R_0$ とし、完全対合でテンソルを分解したときのこの補空間の部分を
$B\in E\otimes V_2\otimes V_3$ と書く。$P$ の第一支持の消滅直線 $L$ 上に
四行四列の一次評価 $F(p):E\to V_1$ ができ、全点で階数3以上である。
$B$ は位置2,3とも全支持である。一文字で $B$ が消えれば、残る $P$ の二項だけからなる
三テンソルは少なくとも一方の端支持が2以下となり、一点ファイバー条件に反する。

まず $\det F$ が恒等零とする。各 $p\in L$ に前点 $y(p)$ と次点 $q(p)$ があり、
次数の一致によりどちらの射も次数1である。$\ker F(p)$ の射を $\alpha:L\to\mathbb P E$ とする。
$4\mathcal O(-1)\to4\mathcal O$ の核は $\mathcal O(-1-\deg\alpha)$、
余核は $\mathcal O(\deg q)$ であり、$\deg\alpha+\deg q=3$、従って $\deg\alpha=2$。
隣接点での零関係から $B(y(p))$ は $\alpha(p)\otimes d(p)$ であり、全支持のため非零である。
$y(p)$ からの基点なし線型写像としては次数1だが、このSegre射の次数は2以上であり矛盾する。

$\det F\ne0$ ならその零因子 $T$ は長さ4で、$L$ と点グラフのスキーム論的交叉に等しい。
実際単元な三次小行列の近くで核点を消去すれば残るのは行列式のみである。
核線 $\alpha:T\to\mathbb P E$ は完全な四次元切断空間を使う。
転置した鉛筆の列 $0\to4\mathcal O(-1)\to4\mathcal O\to\mathcal L_T\to0$ で
余核は $T$ 上可逆、$H^0(4\mathcal O)\to H^0(\mathcal L_T)$ が同型だからである。
点グラフから $y:T\to X_3$ を戻し、$D=H^0(\mathcal O_T)$ とおく。
$B(y)=\alpha\otimes d$ はここでも定階数1の因数分解である。
実際 $F(p)B(y)=0$ は点グラフ上の恒等式であり、$\ker F$ は $T$ 上の直線束である。
線束を自明化し、$d$ 上でいたるところ非零となる一文字を選ぶと、
$B(y)$ のその文字による縮約は $\alpha$ の単元倍である。
従って $y$ の線型座標も $D$ 全体を張る。$d$ の線型座標も同様に $D$ 全体を張る。
さもなくば四文字の一つで $B(y)$ が $T$ 上恒等零となり、$y$ の全線型包を用いると
$B$ の位置2の全支持に反するからである。

これら三つの同定の下で、$B$ により $R_2$ に入る四空間は $D$ の全ての乗法行列の空間である。
実際 $E^*$ と位置2の文字を縮約し、位置3の文字を $y$ 上で読むと対応する二関数の積が戻る。
$D$ は直線上の因子のFrobenius環であり、完全積対合を取ればこの四空間は乗法の双対像となる。
その純関係スキームは $T$ 自身の閉埋込みを含む。各試験環への点評価は積を保ち、
商の二因子の評価の単純積を与える。用いる線型系が $D$ 全体なのでこれは閉埋込みである。
この中の長さ3の部分因子に先の排除を適用すれば矛盾する。
これで長さ2もなく、前定理の有限性からスキーム論的な主張を得る。
重根での固有直線と乗法双対の明示検算は\path{computations/verify_curve_graph_pure_pair.py}にある。
\end{proof}
\begin{proof}[一純点が全曲線に課す条件]
$P=k(u\otimes v)$ とし、今度は $E$ を五次元補空間とする。
ある整曲線成分の位置0で $p(u)=0$ と仮定する。その正規化上の前点を $y$、次点を $q$ とし、
各文字束の共通次数を $d>0$ とする。$F(p):E\otimes\mathcal O\to V_1\otimes L_0$ は階数3一定で、
余核は $L_0L_1$、核 $K$ は階数2、次数 $-2d$ である。
$B(y):V_2^*\otimes\mathcal O\to K\otimes L_3$ は全射である。
実際 $w(y)$ の末尾縮約像は三次元で、その $E$ への射影で失う次元は高々1、像は $K$ に入る。
標的は大域生成で行列式の次数が零なので、この商を与えるGrassmann写像は定数となる。
従って $B(y)$ の $V_2$ 側の支持は一つの固定二平面 $D_2$ である。

$D_2$ を消す非零点 $x$ を取ると $B(x)$ の位置3の支持は $y$ の曲線の消滅空間に入る。
$y$ が三次元以上のベクトル空間を張れば、この支持は一次元以下となる。
$P$ の一項を戻しても $w(x)$ の位置3の端支持が2以下となり矛盾する。
従って $y$ の像は直線、$d=1$ で全位置とも直線である。
次に位置2の対応点は $B(y)$ の位置2の全支持を消すので、位置2の直線は
$\mathbb P(D_2^\perp)$ そのものである。
これで位置2,3の二直線の\emph{直積全体}で $B$ は消える。
その直積への $R_2$ の制限は $(1,1)$ のグラフの式一本で、階数2の双線型形式である。
実際グラフは射影同型のグラフであり、全式が零なら一点ファイバーでなくなる。
従って制限四テンソルはその形式と $u\otimes v$ の積である。
位置2または3の各点で相手側の直線に沿う縮約が非零となるので、
位置1の対応点は $v$ を消し、位置0の対応点は $u$ を消す。
最初に $q(v)=0$ から始めても語順の反転で同じ結論になる。
各整曲線上では $p(u)q(v)=0$ なので少なくとも一方から始められる。
\end{proof}

% !TeX root = ../papers/nc-p3.tex
\section{曲線グラフからの純関係の排除と直線を保つ窓の基底}
\label{sec:curve-graph-pure-exclusion}
\begin{theorem}\label{thm:curve-graph-pure-exclusion}
四辺の射影ファイバーが空または一点で、点グラフが一次元の標準有限Frobenius窓とする。
一辺に純関係があれば全ての整曲線成分は全位置で直線となり、この窓はAS正則でない。
従って直線成分を持つ場合、標準Koszul左右AS正則であることと
全四辺の $\mathbb P R_i$ がそれぞれのSegre多様体と交わらないことは同値である。
また、一次元点グラフの全ての標準正則窓は四テンソルとして半安定である。
\end{theorem}
\begin{proof}
純関係 $r=u\otimes v$ と任意の整曲線成分を選ぶ。
前節の最後の束の議論にはAS性も純関係の一意性も使わない。
その議論によれば曲線は直線となり、反対辺の二直線への制限は
そのグラフの式と $r$ の積となる（第二文字側で恒等零なら語順を逆にする）。
ここで前後に対応する成分は実際の点写像で運ぶ。四移動が成分を置換しても、
得られた二直線は同じ一つの辺のグラフで結ばれているから、
定理\ref{thm:line-component-as-test}の当該反対縮約は階数1である。
同定理の必要性から非正則性が従う。逆に純関係が一つもなく、直線成分が存在すれば、
その軌道の各反対縮約は非零で階数2以下なので常に階数2である。同定理の十分性を使う。

半安定性については定理\ref{thm:all-finite-windows-semistable}の最初の重み評価を使う。
一点と直線の積は辺に含まれないので最低初期重みは $\min(d_i,d_{i+1})$ 以下である。
純関係がないので最高初期重みは $\min(e_i+f_{i+1},f_i+e_{i+1})$ 以下である。
相補辺の完全対合で両者を足し四位置で平均すれば、テンソルの最低重みは
各文字重みの平均の和以下となる。これは純関係のない有限Frobenius候補にも成立する。
\end{proof}

\begin{theorem}\label{thm:four-shift-line-pbw}
標準有限Frobenius窓の各位置に射影直線 $L_i$ を取り、標準四移動で
$L_i$ が $L_{i+4}$ に移るとする。各辺は $L_i$ から $L_{i+1}$ への射影同型の
グラフを含み、その各点で当該辺の左右の評価階数がともに3とする。
直線の外の対応には次元・ファイバーの条件を置かない。
$N_i=L_i^\perp$ を二次元文字空間とする。このとき
\[
 R_i\cap (N_i\otimes V_{i+1})=R_i\cap(V_i\otimes N_{i+1})
 =k C_i\subset N_i\otimes N_{i+1}.
\]
窓が標準Koszul左右AS正則となるための必要十分条件は四本の $C_i$ が全て階数2であることである。
その場合、全整数位置に整合する基底で、各辺の六関係は四文字の六つの下降対を先頭とする
二次Gröbner基底になる。
\end{theorem}
\begin{proof}
直線上の同型に沿って二切断を同定する。反対の二直線における制限はグラフの行列式の倍で、
全て零なら評価階数3に反する。従って完全対合により制限四テンソルは
非零二文字関係 $C_i$ とこの行列式の積である。両端で隣接する点を代入すれば
$C_i\in N_i\otimes N_{i+1}$ である。ここでは四移動が指定直線を保つことを使った。

$L_i$ の点 $p$ における評価 $6\mathcal O\to V_{i+1}\otimes\mathcal O(1)$
の階数は3で、余核は次点による $\mathcal O(2)$ である。核 $K$ は階数3、次数 $-2$。
前点で $w$ を縮約しさらに位置 $i+2$ を縮約すれば、$4\mathcal O\to K(1)$ が得られる。
像は隣接関係により $K(1)$ に入り、階数は前点での反対端の辺評価の階数3である。
ゆえに $K(1)$ は大域生成であり、$K$ は自明束の部分束なので、
分解型は $(0,-1,-1)$ である。従って $H^0(K)$ は一次元で第一の交叉を与える。逆順にも同じである。

必要性は単純加群の解消から直接確認できる。直線への商を $D$ とする。
各位置で基底を運べば $D$ は二変数多項式環の反復である。
ASの場合、左単純のKoszul解消に右直線加群 $e_jD$ をテンソルしてTorを求める。
定理\ref{thm:line-component-as-test}の切断計算は指定直線上の評価階数 $1,3,3,1$ だけを使う。
ゆえに $e_jD$ の右極小自由解消は線型で個数は $1,2,1$、その二次関係が $C_j$ である。
右単純の $D$ 上の二変数解消を、各項のこれらの $A$ 上の自由解消で持ち上げる。
具体的には末尾の単射を比較写像で持ち上げて錐を取り、その余核から $e_iD$ への単射を
再び持ち上げて錐を取る。全ての第 $h$ 項が内部次数 $i+h$ で生成されるから得られた解消は極小である。
最終写像の四係数のうち二つは $e_{i+2}D$ の解消の最終写像から来る。
$C_{i+2}$ が階数1ならこの二係数は比例し、四係数の包は三次元以下となる。
一方、最終写像の双対の余核はAS性により単純で、その一次成分を消すには四文字の全支持が必要である。
全ての $i$ に適用する。

十分性のため $N_i$ の補空間 $T_i$ を取り、四整数ごとに基底が戻ることは要求せず、
$N_i$ の基底を順に運んで $C_i$ を二文字の交換子とする。
$T_i$ は直線の切断で基底を運ぶ。$R_i$ の $T_i\otimes T_{i+1}$ への射影は接線文字の交換子一本を張る。
この射影の五次元核から $T_i\otimes N_{i+1}$ への射影は核 $kC_i$ を持ち全射である。
従って四つの混合下降対を、法線文字から始まる対で書き換えられる。
最後の接線文字の降順対は昇順対と接線文字の少ない対で書き換えられる。
同じ長さでは接線文字の個数を先に、次に法線二文字、接線二文字の順の辞書順で比較する。
全ての簡約は小さい語に移る。長さ3の昇順語は20個で、
実際の三次元数も $4^3-2\cdot4\cdot6+4=20$ である。
従って全ての三次曖昧さが解消する。二次先頭対の包含曖昧さはなく、
ダイヤモンド補題により全長で昇順語が基底となる。二次PBW性からKoszul性、
所与の標準双対Frobenius性から左右AS性が従う。
\end{proof}

% !TeX root = ../papers/nc-p3.tex
\section{曲線グラフの厳密半安定窓と二直線への特殊化}
\label{sec:curve-graded-limits}
\begin{theorem}\label{thm:curve-strictly-semistable}
左右点ファイバーで一次元のグラフを持つ標準有限Frobenius窓とする。
純関係がなく、四テンソルが厳密半安定なら、全次数で標準Koszul左右AS正則である。
従ってこのグラフにおける非正則窓は、純関係を持つ直線配置か、純関係のない安定軌道に限られる。
\end{theorem}
\begin{proof}
純関係がないので前定理の重み計算から半安定性が成立する。
直線を含む場合は定理\ref{thm:curve-graph-pure-exclusion}、
純な平面三次部分曲線を含む場合は定理\ref{thm:plane-cubic-mark-normality}で正則である。
残る窓の任意の $a\times b$ 部分への関係の制限は階数 $a+b-2$ 以上である。
相異なる寸法は定理\ref{thm:curve-graph-pure-frames}による。等しい場合は
$1,4$ は直接、$2,3$ は同定理の完全交叉の議論を使う。
$2\times2$ の階数1は二直線の同型のグラフを与える。
$3\times3$ の階数3は両次数3の純な局所完全交叉を与え、
一方の平面への閉埋込みでその像は平面三次因子になる。どちらも除外した。

次にある点対の接次元が3なら、定理\ref{thm:maximal-tangent-ore}の最初のJacobi計算により
左右の三次元評価核が一致する。前点での三次縮約の末尾の縮約像がこの核なので、
中支持は高々3である。2以下なら両端支持が3なので、中央で得る三行三列の特異な行列の核で
端を縮約すると非零純関係が生じる（中央支持1のときは端を直接縮約する）。
従って三支持が3である。定理\ref{thm:ternary-point-subalgebras}の巡回移送により
この三テンソルの三方向の二文字縮約空間はいずれも基底変換で窓の関係部分空間になるため、
どれにも純関係はない。同定理で全次数の正則性を得る。

これで全点で接次元が2以下、辺のJacobi階数は4以上である。
定理\ref{thm:all-finite-windows-semistable}の等号方向の証明はこの階数、
各 $a\times b$ での下界 $a+b-2$、純関係の不在、左右点ファイバーだけを使う。
実際、最低重みの重複の排除では低い部分の長方形下界と最高部分のSegre回避を使い、
次の四つの関係重みの評価では $(2,2),(3,1),(4,1),(3,2),(3,3)$ とJacobi階数を使う。
これで相補六重みが確定し、同定理の中央差の一致と下降対の証明をそのまま適用できる。
非自明な等号方向は厳密半安定性から存在し、その極限はPBW正則である。
良い特殊化で元の窓も正則になる。
\end{proof}

\begin{theorem}\label{thm:line-polystable-regular-limit}
定理\ref{thm:four-shift-line-pbw}の四移動不変直線とその上の評価階数3を仮定し、
四本の法線関係の階数が2とする。四テンソルは厳密半安定で、
その唯一の多重安定極限は二組の相補的な直線グラフを持つ標準Koszul左右AS正則PBW窓である。
直線外の点ファイバーは任意でよい。各位置の二文字ずつの直和を選ぶと、
まず法線文字がちょうど二つの六種類のブロックだけに特殊化できる。
隣接二位置が法線文字の場合の四ブロックは非退化二次形式二本の積で、
法線文字が交互の二ブロックは各隣接二対への切断でいずれも可逆な四行四列行列となる。
多重安定極限でもこれらの性質を保つ。
\end{theorem}
\begin{proof}
各文字空間を $N_i\oplus T_i$ とする。隣接二接線文字への制限は
反対二法線文字の非退化関係と接線行列式の積なので、
法線文字が二個未満の項はない。法線文字を重み $1$、接線文字を $-1$ にする
四位置の一径数群での極限 $w_0$ は法線二個の部分である。
法線二個が隣接する各ブロックは上述の積そのものである。

各辺の切断の階数は極限で6以下である。両組の同種文字の間の二ブロックは各階数1で、
混合対どうしの $8\times8$ ブロックは階数4以下になる。その行列を
$(NT,TN)$ の順に四分割すると、非対角の両ブロックは上述の非退化な積を
別の二対へ切断したものなので各階数4である。従って混合ブロック全体が階数4で、
その行・列空間は各混合四空間へ同型に射影する。対角の二ブロックも可逆である。
これで両隣接切断の階数6と交互型の可逆性を得た。

従って $w_0$ は定理\ref{thm:split-inverse-regularity}の可逆二平面分解そのものである。
混合四関係の両射影は同型、二つの同種関係は非退化なので、両側の座標直線はいずれも
左右評価階数3のグラフを与える。定理\ref{thm:split-polystable-good-limits}により
その唯一の閉軌道も同じ可逆二平面分解の正則PBW窓である。
同定理の中心化群の証明で使う不変量は、四つの隣接型ブロックを各々二コピー取り、
各二文字位置で交代形式を縮約した四つの値の積である。
各値は二行列式の積の非零定数倍であり、閉軌道でも四ブロックの非退化性を保つ。
この閉軌道は元の四テンソルの閉包に入り、非零なので元の窓も半安定である。
最初の非自明な一径数群が重み零の極限を持つので安定でない。
極限の交代四テンソルの許容スペクトルとJordan分割は
定理\ref{thm:split-bell-locus}、\ref{thm:split-jordan-transport}に従い、
分解付き同型は定理\ref{thm:split-inverse-transporter}の移送集合の共役類で決まる。
\end{proof}

\begin{corollary}\label{cor:curve-two-by-two-equality}
純関係のない一次元点グラフで、四移動不変の直線成分を持つことと、
全位置で最低・最高各二文字の重みを持つ一径数群によるHilbert--Mumford等号があることは同値である。
このような等号方向の正の重み差は全位置で等しい。
\end{corollary}
\begin{proof}
直線からの方向は前定理にある。逆に最低二空間の双対の直線どうしへの制限は各辺で非零であり、
最高二空間の積だけに支持される関係空間は、Segreを避けるので高々一次元である。
定理\ref{thm:all-finite-windows-semistable}の最初の等号計算から重み差は位置ごとに
$s,t,s,t$ と交互になる。最低関係は反対の最高関係としか対合できない。
従って直線の積への制限は一次元である。その双一次式は点ファイバーから非退化で、
四位置の旗が四移動不変の直線グラフを指定する。各法線関係の階数は純関係の不在で2である。
最低を零とした重みの平均の和は $s+t$ である。
前定理の交互二ブロックはどちらも非零で重みは $2s,2t$ であるから、
等号には $2s,2t\ge s+t$ が必要であり $s=t$ となる。
\end{proof}

% !TeX root = ../papers/nc-p3.tex
\section{不変直線から相補逆行列窓への全持上げと二次障害の個数}
\label{sec:invariant-line-filtered-lifts}
可逆二平面分解の任意の窓の四テンソルを $w_2$ とする。
二文字ずつの空間を $N_i,T_i$ とし、$N$ の文字数を重みとする。
$w_2$ は重み2であり、$N_i\otimes N_{i+1}$ 内の関係を $c_i$ と書く。
四位置の文字列の長さは常に4であり、以下の添字は標準四移動で読む。

\begin{theorem}\label{thm:invariant-line-filtered-lifts}
$w_2$ に指定した二平面の重みで特殊化する全窓は、重み3の64係数 $u$ と
重み4の16係数 $v$ から $w_2+u+v$ として得られる。許容条件は次の通りである。
\begin{enumerate}
\item $u$ には48行の斉一次方程式があり、その解空間の次元は16以上24以下である。
\item $c_0,c_1,c_2$ を交換子にする基底を順に取り、残る $c_3$ の対称部分を $q$ とする。
$q=0$（四つの射影的な同一視の合成が恒等）のとき、許される $u$ は
この一次解空間上の五成分の斉二次障害の零点で、各 $u$ の上の $v$ は三次元のアフィン空間である。
$q\ne0$ なら二成分の斉二次障害となり、各零点の上の $v$ は一意である。
成分数は標的の次元を表し、二次式の独立性や高さを要求しない。
\item 全解は全次数で標準Koszul左右AS正則で、指定直線上で前後とも評価階数3である。
指定直線を保つ添字同型は下記の線型群とせん断の軌道である。
従ってこの構成と $w_2$ の相補逆行列の正規形で、直線外に分岐を許した
定理\ref{thm:four-shift-line-pbw}の正則枝全体を被覆する。
\end{enumerate}
\end{theorem}
\begin{proof}[二切断の消去]
切断 $i,i+1\mid i+2,i+3$ を $i=0,1$ で取る。各二文字位置を $TT$, 混合, $NN$ の順に書くと行列は
\[
 \begin{pmatrix}0&0&A\\0&B&C\\D&E&F\end{pmatrix}.
\]
$A,D$ が階数1、$B$ が階数4で、これらは $w_2$ の係数である。
$C,E$ は $u$、$F$ は $v$ だけで決まる。$A$ の $NN$ 側の直線と
$D$ の $NN$ 側の直線がちょうど $k c_{i+2},k c_i$ である。
この両直線で割った $C,E,F$ の像を $C',E',F'$ とする。行・列のピボット二つを外すと
残るこの行列の階数が4である条件は、$BJB=B$ なる一つの一般逆 $J$ を用いて
\[
 \operatorname{im} C'\subset\operatorname{im} B,\qquad
 \operatorname{row} E'\subset\operatorname{row} B,\qquad F'=E'JC'
\]
である。前二条件の下で右辺は一般逆の選択に依存しない。
各切断で一次条件は $4\times3$ 個ずつ二組、二次の右辺は $3\times3$ 個である。

一次条件の余核次元は高々8である。実際、接線文字が位置 $j$ だけにある十六空間を $U_j$ とする。
辺 $j,j+1$ の一次条件は $U_j\oplus U_{j+1}$ から十二空間への全射で、各直和因子への制限も全射である。
これは混合関係四平面が両混合半分に同型に射影し、反対辺では $NN$ を $kc_{j+2}$ で割るからである。
四つの条件の線型従属では、どれか一辺の係数で残り三辺の係数が決まる。
実際隣り合う条件の係数は $U_j^*$ 内で互いに打ち消し、各々への双対写像は単射である。
その共通像は $U_j^*$ の中で三法線文字の相続く二辺の $c$ をそれぞれ零化する。
これら二辺を交換子に揃えると、三つの法線文字で完全に対称な三次式と接線文字の積だから次元8である。
これで四十八行の階数は40以上48以下である。

次に $v$ の二つの射影の合計を $L:k^{16}\to k^9\oplus k^9$ とする。
各射影の双対像は、互いに反対の二辺でそれぞれ $c$ を零化する四テンソルの九空間である。
二つの像の交わりは四辺全てで $c$ を零化する。最初の三辺で交換子を零化すれば
通常の $\operatorname{Sym}^4(k^2)^*$ となり、残る縮約は $q$ による二次の縮約である。
$q\ne0$ ならその双対は二次式への $q$ 倍で単射だから、縮約の階数は3である。
従って交わりの次元は $q=0$ なら5、他は2、$L$ の階数はそれぞれ13と16である。
二つの $E'JC'$ の組の $\operatorname{coker}L$ での類が主張の明示的な二次障害となる。
\end{proof}
\begin{proof}[全次数性と同型]
上の式から各切断は階数6である。重みの一径数で $u,v$ にそれぞれ $t,t^2$ を掛けると、
全ての $t\ne0$ で同じ窓、$t=0$ で $w_2$ が得られる。
命題\ref{prop:unrestricted-good-reduction}から非零ファイバーも標準Koszul左右AS正則である。
直線上の評価を調べる。関係空間の重みによる初期空間は $w_2$ の関係空間である。
実際二つの切断はいずれも階数6一定なので、像の極限は特殊点の像である。
従って接線二文字への射影はその行列式であり、混合四関係の持上げの評価は
任意の指定直線上の点で二つの法線文字を張る。左右とも階数3である。
逆に指定不変直線を持つ初めの正則窓では定理\ref{thm:line-polystable-regular-limit}の
最初の特殊化がちょうどこの $w_2$ を与え、係数は以上の全条件を満たす。

$w_2$ の射影類を保つ $\prod_i(\mathrm{GL}(N_i)\times\mathrm{GL}(T_i))$ と
$h_i:T_i\to N_i$ のせん断で全ての指定付き同型が得られる。
先の群では $w_2$ の代表を揃えるため全テンソルの共通倍率を正規化する。
せん断の作用は各 $h_i$ をその位置だけに作用させて
\[
 u\mapsto u+\textstyle\sum_i h_iw_2,\qquad
 v\mapsto v+\textstyle\sum_i h_i u+\sum_{i<j}h_i h_j w_2
\]
である。これは指定直線を保つ全基底群のブロック三角分解から必要十分で、
余分な変換は生じない。印を忘れるときは第二の窓の指定可能な四移動不変直線を動かして比較する。
一次元グラフならこの選択は成分数以下である。一移動・二移動にも同じ比較を適用できる。
\end{proof}

% !TeX root = ../papers/nc-p3.tex
\section{六次式の正則一点が対称性を外した窓全体でも独立な成分になること}
\label{sec:isolated-split-component}
\begin{theorem}\label{thm:isolated-split-component}
定理\ref{thm:split-bell-locus}の六次式から得られる無印の一類を $[w]$ とする。
隣接二切断の階数がともに6以下となる四テンソルの射影スキームを $Y$、
基底群を $G=\prod_{i=0}^3\mathrm{PGL}(V_i)$ とする。
$Y$ は $[w]$ で滑らかで、その唯一の既約成分は $\overline{G[w]}$、次元は59である。
$G[w]$ は $Y$ の開集合で、この成分の半安定部分全体に等しい。
従って半安定商ではこの最小周期1の左右AS正則窓は被約な開かつ閉な一点である。
標準窓の任意の族で $[w]$ に特殊化する幾何学的一般ファイバーはこれと添字保持同型である。
\end{theorem}
\begin{proof}[標数零の代表と良い剰余]
分解の八つの二文字空間を定理\ref{thm:split-inverse-transporter}の順とする。
$P$ のBell四数は定理\ref{thm:split-bell-locus}第三型の $(1,b,c,d)$ で、
$(d^2+1)^3-4d^3=0$ とする。$\iota^2=-1$ とし、第一切断の四対応を恒等、
第二切断の四対応を順に
\[
 \begin{pmatrix}1&1\\-\iota&\iota\end{pmatrix},\quad
 \tfrac12\begin{pmatrix}1&1\\\iota&-\iota\end{pmatrix},\quad
 \begin{pmatrix}1&-1\\\iota&\iota\end{pmatrix},\quad
 \tfrac12\begin{pmatrix}1&-1\\-\iota&-\iota\end{pmatrix}
\]
に取る。両逆転置の余因子をこれらの対応で運び、第二の組と第一の組の射影比例を
六次式を法として全十六座標で検算した。従って六ブロックのSchur式の倍率が決まり、
同定理により標数零の代表となる。$109$ を法とし $d=17,\iota=33$ に取れる。
六次式の根は単純で、両切断行列式、倍率の規格化に用いた座標、四対応の行列式は
全て単元である。この整数環の局所化上で両方の切断を比較できる。
\end{proof}
\begin{proof}[制約なしの接空間]
$16\times16$ の階数6行列 $M$ において七次小行列の一次条件は、
左右の十次元核の基底を行列 $L,R$ として $L\dot M R=0$ である。
二切断それぞれから100行を得て、全256係数を列とする。重みは法線文字の個数に取る。
基点の重みが一定なので核は重みごとに分解する。
前記の良い剰余での列数、制約階数、接次元は順に
\[
\begin{array}{c|rrrrr}
\text{法線文字数}&0&1&2&3&4\\\hline
\text{列数}&16&64&96&64&16\\
\text{階数}&16&48&68&48&16\\
\text{核次元}&0&16&28&16&0 .
\end{array}
\]
全行列の階数も196である。\path{computations/verify_isolated_split_component.py}
は上述の体上の射影比例、剰余での四テンソルの構成、両核からの全行列を再現する。
六次小行列の単元に沿って両核の基底を持ち上げられるので、標数零のアフィン接次元は60以下となる。
一方、定理\ref{thm:split-bell-locus}により射影安定群は一次元トーラス、
射影軌道の次元は $4(4^2-1)-1=59$ である。原点以外の錐でスカラーを戻すと60となる。
従って接次元は局所次元と一致し、標数零で滑らかである。
その点を通る成分はただ一つで、その次元に等しい軌道の閉包がこの成分になる。
軌道はその閉包で開、他成分はそのどの点も含まないので $Y$ でも開である。

この軌道は半安定部分で閉であることも定理\ref{thm:split-bell-locus}で示した。
従って軌道閉包のそれ以外の点は全て不安定である。商でこの軌道と同じ像を持つ
半安定点は、自身の軌道閉包がこの開軌道に会うので初めから同じ軌道に属する。
従って半安定部分の開かつ閉な飽和集合であり、その商は被約一点である。
最後の族に関する主張も軌道の開性と二窓の同型定理による。
\end{proof}

% !TeX root = ../papers/nc-p3.tex
\section{正則格子の比較による射影軌道の全正則特殊化}
\label{sec:regular-orbit-valuations}
\begin{theorem}\label{thm:regular-orbit-valuations}
$D$ を離散付値環、$K$ をその分数体とする。二つの標準四窓の $D$ 上の模型を取り、
一般・特殊ファイバーがいずれも標準quadratic Koszul左右AS正則とする。
分数体で一方を他方に運ぶ四文字の基底変換を $g_i$ とし、
標準四位置の文字格子に関する $\operatorname{ord}\det g_i$ を $v_i$ とする。
原始的な最上テンソルの生成元どうしを比較した倍率の付値を $m$ とする。このとき
\[
 \sum_{i=0}^3 v_i=4m,\qquad
 \sum_{i=0}^3\operatorname{ord}\det(g:A_{i,i+n}\longrightarrow A'_{i,i+n})
       =4m\binom{n+3}{4}\quad(n\ge0).
\]
従って、半安定な標準AS窓の射影軌道閉包に属する標準AS窓はすべて半安定で、
半安定商において元の窓と同じ点に写る。とくに多重安定な標準AS窓の
射影軌道閉包に属する標準AS窓は元の添字保持同型類だけである。
\end{theorem}
\begin{proof}[付値のEuler等式]
各模型で窓の二切断の像を二次関係の格子に取る。階数6は一般・特殊で等しく
単元な六次小行列を持つので、これは直和因子である。生成した二次代数の各長さ $n$ の経路加群は
有限表示で、両ファイバーで同じ次元 $h_n=\binom{n+3}3$ だから自由である。
単純の線型解消に使う二次双対の各次数の双対格子 $W_{i,j}$ も、
両ファイバーの次元一致と単元な小行列から自由で、基底変換と可換である。
最高積の左右対合は剰余で完全なので格子上で完全である。
左右Koszul複体の各長さでの項も自由で、特殊ファイバーの完全性から正の長さでは
格子上で分裂完全である（単元なピボットを端から順に除ける）。

長さ $n$ の経路格子の行列式付値を $a_i(n)$ とする。比較は $K$ 上の同型でよく、
行列式は両側の $D$ 基底で取る。$W_{i,2}$ の比較付値を $r_i$ とし、
最上直線 $W_{i,4}$ の付値は原始テンソルの巡回で全位置で $m$ である。
完全対合から $r_i+r_{i+2}=6m$、$W_{i,3}$ の比較付値は $4m-v_{i+3}$ となる。
分裂完全な二列の間の分数体同型では各項の行列式付値の交代和が零である。
各項 $W_{i,j}\otimes A_{i+j,i+n}$ に適用し、始点四つの和を取る。
$A(x)=\sum_{n\ge0}(\sum_i a_i(n))x^n$, $V=\sum_i v_i$ とすれば
\[
 (1-x)^4 A(x)=(1-x)^{-4}\big(Vx-12m x^2+(16m-V)x^3-4m x^4\big).
\]
実際、テンソル積の付値は二因子の付値に相手の階数を掛けた和である。
次数零は同一の単純の比較で寄与しない。

両模型の一次格子の間で $g_i,g_i^{-1}$ の全係数の極の次数を定数 $c$ 以下に取る。
長さ $n$ の積で各経路格子が生成されるので、$g$ は一方を他方の $\pi^{-cn}$ 倍以内に送り、
逆も同様である。従って各基本因子の付値の絶対値は $cn$ 以下、
$|a_i(n)|\le cn h_n=O(n^4)$ である。上式の右の分子は $x=1$ で少なくとも三重に消える。
さもなくば $(1-x)^{-8}$ を掛けた係数の増大がこの評価を超えるためである。
一次微分は $8m-2V$ なので $V=4m$ を得る。分子は $4mx(1-x)^3$ に等しく、
係数を取ると第二式も得る。
\end{proof}
\begin{proof}[軌道閉包]
射影軌道の境界点を特殊点に持つ曲線をその閉包内で取り、正規化・有限体拡大の後、
一般点を基底群の $K$ 点から来るものとできる。
実際、軌道写像の一般ファイバーには有限拡大上の閉点がある。
有限拡大を追加して各四次行列の行列式の四乗根を取れば、スカラー倍で全位置を
特殊線形基底変換に取れる。特殊点に延びる原始的な四テンソルは、
一定な原始テンソルの像を $\pi^m$ で割ったものに単元を掛けて得る。
特殊点も標準正則なら二つの模型に上式が適用でき、$v_i=0$ なので $m=0$ となる。
ゆえに特殊窓には、元のテンソルの特殊線形群の\emph{アフィン}軌道閉包内の非零代表がある。
元が半安定ならある斉次不変式が非零で、その値はこの代表にも保たれる。
全不変式の値が一致するので射影商でも同じ点である。
元のアフィン軌道が閉ならこの代表はその軌道自身に入り、添字保持同型を得る。
\end{proof}

\begin{corollary}\label{cor:isolated-split-component-all-regular}
定理\ref{thm:isolated-split-component}の59次元成分に属する標準Koszul左右AS正則窓は、
不安定境界も含めて調べると六次式の一類で尽きる。
この一類は標準正則軌跡で開かつ閉であり、他の類への特殊化も他の類からの特殊化もない。
\end{corollary}
\begin{proof}
この成分は一つの多重安定な標準窓の射影軌道閉包である。上の定理を直接適用する。
逆向きは同成分の滑らかな開軌道であることから従う。
\end{proof}

% !TeX root = ../papers/nc-p3.tex
\section{可逆二平面分解の点スキームと孤立成分の五直線}
\label{sec:split-point-pencils}
\begin{theorem}\label{thm:split-point-pencils}
可逆二平面分解を持つ全窓は、全位置で左右の射影ファイバーが空または一点である。
一辺の法線・接線の非退化関係を二行列 $C,D$、混合四関係の $NT,TN$ ブロックをそれぞれ
$R_a,S_a$ ($1\le a\le4$) とする。$J=\left(\begin{smallmatrix}0&-1\\1&0\end{smallmatrix}\right)$ とおく。
$n\in\mathbb P(N^*)$ のまわりで座標を一つ1にすれば、$n$ 側の点グラフは
次の四行四列鉛筆と一つの二次式でスキーム論的に与えられる。
\[
 M(n)=\left[\,n^tR_a\ \ (S_a JC^t n)^t\,\right]_{a=1}^4,\qquad
 M(n)\binom{t'}t=0,\qquad t^t D t'=0 .
\]
$t$ は始点の接線座標、$t'$ は終点の接線座標である。
接線側の直線の近傍は $C,D$ と $R_a,S_a$ を交換した式で与えられ、両者で全グラフを覆う。
\end{theorem}
\begin{proof}
始点を $[n:t]$、終点を $[n':t']$ と書く。$n\ne0$ のとき $n'=0$ なら四混合式は
$n\otimes t'$ を零化し、$R_a$ が全二行二列空間を張るので $t'=0$ となる。
従って $n'\ne0$ で、同種式 $n^t Cn'=0$ から、局所座標では $n'=JC^tn$ と規格化できる。
これを他の式へ代入すれば表示の式を得る。鉛筆の左二列は各 $n$ で階数2なので
$t$ を指定すれば $t'$ は高々一つである。交換した近傍でも同じであり、
語順を逆にした混合射影も同型なので反対側も同じである。
これらの消去には単元な小行列しか使わないため、非被約基底でも有効である。
\end{proof}

\begin{proposition}\label{prop:split-simple-four-roots}
両側の座標直線の鉛筆の行列式がそれぞれ相異なる四根を持つとする。
その根で $M(n)$ の核を張る $(t',t)$ が $t^tDt'=0$ を満たすものの数を $g$ とする。
二側でこの数は等しい。被約グラフの像は二本のねじれ直線と、両者を結ぶ
互いに交わらない $g$ 本の直線である。各結線は異なる根どうしを結び、
両端は通常二重点である。残る $4-g$ 個ずつの根では、各々の直線上に長さ1の
埋込み部分があり、これが全ての非被約構造である。
\end{proposition}
\begin{proof}
単根の局所環で鉛筆の階数は3、三つの単元ピボットを除くと最後の式は $zu=0$
（$z$ は根の局所座標）、残りは $h(z)u^2=0$ である。$h(0)$ はまさに述べた核での値である。
$h(0)\ne0$ なら $u^2=zu=0$ で一次元の埋込みイデアル $ku$ を持つ。
$h(0)=0$ なら第二式は第一式から従い、$z=0$ 上の $u$ は任意である。
両方の二列ブロックが単射なので核で $t,t'$ はともに非零である。
従って始点では $[n:ut]$ が座標二直線を結ぶ射影直線の開部分である。
反対の座標直線の近傍でその無限遠点も同じ単根の形となる。
単根の核は一本だけなので、異なる結線は各側で異なる点に至り、
二つの相補平面の直和表示から結線どうしも交わらない。
行列式が単元の点では $t=t'=0$ となるから他の成分はない。
\end{proof}

\begin{corollary}\label{cor:isolated-split-point-scheme}
六次式の孤立成分の唯一の標準窓の点グラフは、二本のねじれ直線と三本の結線を持つ。
六交点は通常二重点で、これとは別にねじれ二直線の上に長さ1の埋込み部分が一つずつある。
特に被約曲線は次数5、算術種数2で、分解の二本の直線は三つの交点を持つ成分として無印で区別できる。
\end{corollary}
\begin{proof}
定理\ref{thm:isolated-split-component}の標数零の代表を使う。
\path{computations/verify_isolated_split_point_scheme.py} は
$\mathbb Q(\sqrt{-1})[d]/((d^2+1)^3-4d^3)$ 上で二切断の相補式から四テンソルを復元する。
各座標直線で $n=(1,y)$ として得る行列式は次数4で、その微分と互いに素である。
随伴行列の四列それぞれを $(t',t)$ として求めた $t^tDt'$ の四式と
行列式との最大公約因子は、両座標直線で次数3である。単根では随伴行列の階数が1のため、
この検査は核による判定と一致する。無限遠で行列式が単元なので全ての点を数えている。
前命題に $g=3$ を適用する。五本の射影直線を六個の通常二重点で貼ると
算術種数は $6-5+1=2$ である。既に最小周期1であるから全位置で同じ記述になる。
\end{proof}

\begin{corollary}\label{cor:no-branch-on-invariant-line-lifts}
標準正則窓に四移動不変の直線グラフがあり、その上の左右評価が全て階数3なら、
この窓のどの辺のどの点にも分岐はない。さらに全点グラフの次元は、
二平面分解への重み特殊化で得られる点グラフの次元以下である。
\end{corollary}
\begin{proof}
定理\ref{thm:invariant-line-filtered-lifts}により重み特殊化の特殊ファイバーは可逆二平面分解になる。
評価階数2以下の点の入射は、階数6の関係束と各位置の射影三空間との積上で閉である。
射影空間は固有なので、族の基底への像も閉である。特殊ファイバーには前定理によりこの点がない。
非零ファイバーは互いに同型なので、どれにも分岐はない。点グラフの次元も射影族の
上半連続性で特殊ファイバーの次元以下となる。
\end{proof}

% !TeX root = ../papers/nc-p3.tex
\section{相補逆行列の曲面の一般点に至る全窓の六次元族}
\label{sec:split-surface-full-fibers}
\begin{theorem}\label{thm:split-surface-full-fibers}
Bell曲面の一般の $P=(1,-uv,u,v)$ を取り、定理\ref{thm:split-bell-locus}の四つの
指定付き補完の一つを $w_2$ とする。各補完で次が成立する。
\begin{enumerate}
\item 指定法線側の全持上げは、せん断を除くと重み $1,2$ の四座標と三座標のアフィン空間になる。
すべての値で標準Koszul左右AS正則である。原点以外の添字保持同型類は
$\mathbb P(1,1,1,1,2,2,2)$ を残る四元群で割ったものである。
四元群の作用は相補逆行列表示で $P$ を保つ四つの局所基底変換から誘導する。
\item $w_2$ は多重安定で、この六次元族と原点の一類が、
分解も不変直線も指定しない全窓の中で $[w_2]$ に特殊化する同型類を尽くす。
各補完で分解の交換は自己同型により実現する。従って一般の $P$ の四つの補完は
分解を忘れても、$P$ の座標付け替えによる同一視以外には重ならない。
\item 原点の点グラフの純部分は二本のねじれ直線で、各々に四つの長さ一の埋込み部分がある。
非零の持上げの点グラフでは一次元部分は一本の直線で一般被約である。
従ってこの直線は無印で復元できる。曲面基底をも動かす全族の同型類の像は八次元である。
\end{enumerate}
\end{theorem}
\begin{proof}[相補行列、軸交換と閉軌道]
$X=\left(\begin{smallmatrix}0&1\\1&0\end{smallmatrix}\right)$, $Z=\operatorname{diag}(1,-1)$ とする。
第一の四対応を恒等にした表示で、第二は
\[
                  (hZ,h,hZ,h),\qquad h=1,X,Z,XZ
\]
と取れる。実際、両Schur式への代入で逆行列の移送を満たし、
定理\ref{thm:split-bell-locus}の四つの補完を全て与える。
両分解を交換すると第一の四対応は恒等のまま、第二は
$(h,hZ,h,hZ)$ となる。これは対応を表す行列を両端の交換で転置することによる。
新しい交代テンソルのBell四数は逆数である。その四位置に $(1,Z,Z,1)$ を作用させると
旧 $P$ に射影的に戻る。実際 $(01)(23)$ の四数の入替えと第三・第四係数の符号反転で、
$a b=-c d$ を用いる。第一の四対応を保つよう相補側にも対応する逆転置を作用させると、
第二の四対応は、各 $i$ で $k_{i+2}^{-t} L_i k_i^t$ ($k=(1,Z,Z,1)$) となり、
元の $(hZ,h,hZ,h)$ に戻る。各局所行列の共通倍率は相補二切断の復元で吸収できる。
これで全四補完における交換自己同型と、二つの向きの同一視を得る。

複素数上で四数をBell基底に取り、両方の交代テンソルにビットの全反転と全位置の $Z$ が作用すると、
各二文字空間上のGram行列はスカラーとなる。隣接二文字型ではユニタリ行列の積で同じことが従う。
全体の各一文字Gram行列は重みからブロック対角である。六種類の法線二文字の台が全て非零なので、
定理\ref{thm:split-bell-locus}のノルム最小化で二ブロックのスカラーを各位置で一致させられる。
従って $w_2$ は多重安定である。
\end{proof}
\begin{proof}[一般係数体での全障害]
定理\ref{thm:invariant-line-filtered-lifts}の線型条件の核は
$\mathbb Q(u,v)$ 上で20次元、せん断の像は16次元である。
四法線係数の線型写像の階数は13で、その像に、
せん断の余空間の四基底の二次補正の全十係数が既に入る。
\path{computations/verify_split_surface_lifts.py} が四つの $h$ でこの等式を有理関数体上で検算する。
具体的に二つの切断 $M$ で $L,R$ を左右核、$J$ を六次ピボット上の一般逆とし、
三法線係数の基底候補 $U$ を $LUR=0$ で求める。その核にせん断の十六列を入れ、
余空間の四列 $U_b$ を取る。四法線十六列の制約行列に
$L U_b J U_b R$ と $L(U_b J U_c+U_c J U_b)R$ ($b<c$) の両切断の十列を追加しても、
階数は13である。ピボットの選択だけに有理標本を使い、零の判定は一般係数で行う。
これらは有限個の行列の恒等式と非零小行列なので、共通の非空開集合の全点で成立する。

せん断は一次部分で自由に作用するので四座標の切片をちょうど一度通る。
切片上では全ての二次障害が消え、四法線部の一つの二次の持上げに三次元の核を足せる。
残るブロック対角群は共通の法線倍率と、$h$ の移送を保つ $P$ の四元群である。
これは定理\ref{thm:split-inverse-transporter}の安定群計算による。
この群は簡約的なので一次余空間および四法線の線型像の切断を同変に選べる。
共通法線倍率は二群の座標に重み1,2で作用し、表示した商を与える。
正方向のせん断安定群が零であること、軸交換で負方向も零であることから、
射影安定群の恒等成分は上の一次元トーラスだけである。
\end{proof}
\begin{proof}[印を仮定しない被覆]
$w_2$ に特殊化する隣接階数6以下の窓は半安定である。多重安定軌道への
Hilbert--Mumford判定により、一つの基底変更の後に一径数群でこの極限を取れる。
極限を保つ一径数群はその安定群の恒等成分に入るので、
原点自身を除けば上の重みまたはその逆の正の倍である。前者は全持上げで尽くし、
後者は軸交換で前者に戻る。

点グラフについては命題\ref{prop:split-simple-four-roots}の二つの鉛筆を使う。
$u=2,v=3$ では四補完とも各行列式は単純四根、四つの核二次式との共通因子は定数である。
上の検算でこれらの開条件を両座標直線上で確認する。軸以外に曲線はなく、
これらの軸はともに一般被約である。任意の持上げは分岐を持たず、
その点曲線の重複込み次数は2以下となる。実際、非零ファイバーの純曲線を一径数で動かし
射影Hilbertスキームで平坦極限を取れば、極限は基点の点グラフ内に含まれ、
次数は保たれる。二重の指定直線は、その台が一定なので一般被約な同じ直線に二重には特殊化できない。
さらに別の曲線があれば一本の直線であり、指定直線と相補な方の軸に特殊化する。
二直線が交わる条件は閉で、極限の両軸が交わらないので最初から二直線は交わらない。
第二の直線は他に成分がないため四移動不変である。この直線を消す平面を新しい補空間に取ると、
二つの直線のそれぞれについて反対辺の制限が残る二文字平面の行列式関係に入る。
従って法線文字は二個以上かつ二個以下で、テンソルは重み2だけとなる。
補空間の変更はせん断なので、これは切片の原点である。
これで非零点の指定直線の一意性が分かり、放物群による比較が無印でも必要十分となる。
基底の $P$ は閉軌道から有限通りに復元できるので、次元は $2+6=8$ である。
\end{proof}

% !TeX root = ../papers/nc-p3.tex
\section{Bell曲面の全流入窓の周期二部分の四成分}
\label{sec:split-surface-periods}
\begin{theorem}\label{thm:split-surface-periods}
前節の曲面の一般の閉軌道の四補完のいずれでも、閉軌道の最小周期は2である。
閉軌道に至る非自明な全窓のうち、周期2の軌跡は同型を除いて
次元 $4,2,2,2$ の四つの既約成分を持ち、それ以外の点は全て最小周期4である。

具体的に、重み1の四座標を $b$、重み2の三座標を $c$ とし、
$\rho$ を下記の補正二移動、$k$ を $1,Z,X,XZ$ の四つの補正四元群元とする。
$E_\pm(k\rho)$ を該当する座標空間での固有空間と書く。
$b\ne0$ の周期2条件は、ある $k$ で $b\in E_+(k\rho)\cup E_-(k\rho)$ かつ
$c\in E_+(k\rho)$ となることである。$b=0$ なら常に周期2である。
\end{theorem}
\begin{proof}[一般点での二移動]
$P$ の一巡回あるいはその相補逆の一巡回が元のBell四数と局所同値でなければ、
軸交換も許して一移動同型は生じない。$(u,v)=(4,7)$ で
$(1,-28,4,7)$ に対する比較相手は $(-15,-12,20,-9)$ と $(15,12,20,-9)/28$ となり、
どちらも射影的な偶符号変更・置換で元に戻らない。
ゆえに一般点で一移動はなく、閉極限の一意性からその上の全窓でもない。

$\kappa=-(u^2-1)(v^2-1)/(4uv)$ とする。二移動で四テンソルの
六ブロックのうち法線位置が $01$ のものは $\kappa$ 倍、$23$ は $\kappa^{-1}$ 倍、
他の四つは一定である。二移動の後で位置別の法線文字だけに
$(1/s,1/s,s,s)$ を掛けると $w_2$ を固定する（$s^2=\kappa$）。これを $\rho$ とする。
$Z$ は全二文字のビット符号、$X$ は全ビット反転の後に法線だけに
$(\iota,\iota,-\iota,-\iota)$ を掛ける補正とし、$\iota^2=-1$ とする。
これらも $w_2$ を固定する。テンソル全体で $X,Z$ は可換、
$X^2$ は法線数 $n$ の空間で $(-1)^n$、$Z^2=\rho^2=1$ である。
$\rho Z=Z\rho$、$\rho X=(-1)^n X\rho$ なので、四つの $k\rho$ は全て対合となる。
この有限群と法線倍率を含めて前節の同変な二つの切断を選べば、
同じ作用が切片の $(b,c)$ に線型に作用する。

重み1側は三法線の一次核をせん断で割った四空間に等しく、
重み2側は四法線の二つの射影の共通核である。
前者の $Z$ 正負固有空間は各二次元で、各 $k\rho$ の正負固有空間も各二次元である。
後者では $k=hZ$（$h$ は補完の添字）のとき $k\rho$ は恒等、
他の三つでは正固有空間は一次元、負固有空間は二次元となる。
\path{computations/verify_split_surface_periods.py} は四つの切断の倍率を一般体で、
固有空間を $u=2,v=3$, 素数109、$s=\iota=33$ で計算する。
ここで一次核・せん断・四法線核の次元は一般値を保つ。階数が一定の近傍では
これらの格子は局所自由である。2を可逆にすると対合の二射影は直和因子となるので
固有空間の次元が一般体にも戻る。$s$ のもう一方の選択は重み1側で符号を交換するだけである。
\end{proof}
\begin{proof}[全座標での判定と成分数]
二移動後に $k$ と法線倍率を用いて戻ればよい。
$b\ne0$ のとき $k\rho$ の $b$ 上の固有値は $1$ または $-1$ であり、
その平方が $c$ 上の倍率となるため $c$ は正固有空間に入る。
$b=0$ は $k=hZ$ で全て含まれる。$X$ は重み1上で $k\rho$ と反交換し、
重み2上では可換なので二符号は同型になる。
各 $k$ の像は従って一つの既約な重み付き線型部分空間の有限商であり、
その次元は $2+\dim E_+(k\rho)|_c-1$ である。

相異なる $k$ の像の一般の $b$ の方向は異なる。
$Z$ の各固有二平面で $X$ の二固有線は $\rho$ によって交換されるので、
$k\rho$ の各固有平面は $Z$ の二平面それぞれから一本の固有線を取る。
$k$ を $kZ$ に替えると一方だけの符号が反転する。
$k$ を $kX$ または $kXZ$ に替えると、各 $Z$ 固有二平面で元の対合と反交換するため
固有線は互いに異なる。四元群での共役は各 $k\rho$ の符号だけを変える。
従って小次元の三つの像も他の像に含まれず、表示した四成分で尽きる。
\end{proof}

% !TeX root = ../papers/nc-p3.tex
\section{Bell曲面の全流入と十五次元成分およびその滑らかな近傍}
\label{sec:split-surface-component}
\begin{theorem}\label{thm:split-surface-component}
$u v(u^2-1)(v^2-1)\ne0$ における前節の四つの明示的な二平面分解族は
全て定理\ref{thm:full-fifteen-component}の十五次元成分に属する。
各族の一般の閉軌道と、それに至る\emph{全て}の無印の窓で、
隣接二切断の階数スキームはアフィン錐上で滑らかな76次元であり、この成分だけが通る。
さらにそれらの全点を含む不変なザリスキー開集合で、全窓は標準Koszul左右AS正則で、
四次中心は五次元、代数は四つの中心四次式の生成する部分環上有限である。
一般閉軌道での射影切片の接表現のトーラス重み $-2,-1,0,1,2$ の多重度は $3,4,2,4,3$ である。
\end{theorem}
\begin{proof}[五中心の計算と量子特殊化]
四位置で二進表示した四文字の排他的論理和により各関係に四元群次数を与える。
両行列切断の六関係を $\mathbb Q(u,v)$ で取り、各始点で経路長5まで一次関係を挿入して
経路基底を取り、一文字との可換式の両側を射影する。
長さ4の中心切断の次元は四文字次数について $(5,0,0,0)$ である。
\path{computations/verify_split_surface_component.py} は四つの $h$ 全てで、
有理関数体上の経路基底から計224行・44列となる自明指標の可換式および残り三指標を計算する。

$u=\iota$, $v=\iota$ を $h=X,XZ$ に、$u=\iota$, $v=-\iota$ を $h=1,Z$ に使う。
全て $\iota^2=-1$ とする。どちらの四文字表示も正則で、二切断の双方の
相補逆行列の分母は単元である。前者は位置2,3、後者は位置1,3で各二文字のビットを反転する。
非零の24語は全て四文字の置換になる。各辺の二つの文字 $a,b$ の交換係数を $q_{i,ab}$ とすると、
前記の検算で、各向きの相補二文字をどちらの順序にとっても同じ比が得られる。
四つの三角積 $q_{i,ab}q_{i,bc}/q_{i,ac}$ ($a<b<c$) は全 $i$ で同一の
$(\iota,-\iota,\iota,-\iota)$ またはその逆数である。
従って相続く位置の交換係数の比は文字ごとの倍率の比であり、添字付きの対角基底変更で
一定な量子交換係数にできる。この三角積は定理\ref{thm:quantum-component-intersection}の
唯一の位数4の量子類を与える。

二変数族の全点で経路の各次数は標準次元なので、四次中心の可換式は束の射であり、
五次元以上の条件は閉である。四族とも既約で、一般点で五次元なので全体がこの閉条件に入る。
量子点の近傍では定理\ref{thm:quantum-ore-incidence}により、この閉条件は
既知の十五次元成分そのものである。上の特殊化を用いると、各既約族全体が同じ成分に入る。
\end{proof}
\begin{proof}[全接方向、全流入と正則開集合]
$u=2,v=3$ を素数109で還元し、定理\ref{thm:isolated-split-component}の
左右核の全200行・256列で接方程式を取る。法線数 $0,1,2,3,4$ のアフィン接次元は
四補完とも $3,20,30,20,3$ である。各切断の階数6は保たれており、
標数零の幾何学的一般点でもこの和76が接次元の上界である。
既知成分のアフィン次元が76だから等号で、局所環は正則、通る成分は一つだけである。
射影化と軌道方向を除くには、重み $\pm1$ のせん断各16、重み零の28
（ブロック対角群の次元32から三つの自明スカラーと一つのトーラスを除く）を引く。
この28にはアフィンの共通倍率の方向も入る。これで表示した多重度を得る。
滑らかでこの一成分だけに属する条件は基底群不変な開条件なので、
それを満たす点を軌道閉包に持つ全窓にも成立する。

最後に量子点では通常の四乗四元が中心で、これらの商は
各文字の指数が3以下の量子単項式で張られる。標準四移動表示への輸送では中心切断として取る。
量子点の四次中心は五次元であり、この成分上で一般点も五次元なので量子点の近傍で束となる。
その四元を局所切断に持ち上げ、中心商の有限性と標準解消を認証する有限次数検査の
開性（定理\ref{thm:central-window-cutoff}）を適用する。
これにより量子点を含む成分内の開集合の全点で主張が成立する。
その基底群による全ての移動の和も同様で、四族の各一般点は上の特殊化からこの開集合に入る。
さらに各一般点で得た滑らかな一成分だけの近傍に制限すれば、全階数スキーム内の不変開集合が得られる。
軌道閉包にその一般点を持つ全窓もこの開集合に入る。
\end{proof}

% !TeX root = ../papers/nc-p3.tex
\section{十五次元成分の全点列に対する行列恒等式と分岐周期の上界}
\label{sec:pi-periodic-point-paths}
$B$ は標準四頂点で折り畳んだ代数とする。点列の長さは文字の個数で数え、
周期は頂点ラベルも含めて最小に取る。$s_m$ は $m$ 変数の積の交代和である。
\begin{theorem}\label{thm:pi-periodic-point-paths}
十五次元成分の閉包の任意の標準AS正則窓について次が成立する。
\begin{enumerate}
\item 周期点列の最小周期は $4,8,12,16$ のいずれかである。
順または逆の分岐を一度でも訪れる周期列の最小周期は $4,8,12$ に限る。
従って曲面射影を持たない分岐窓では、定理\ref{thm:periodic-branch-obstruction}の
分岐を通る周期列の長さの上界を40から12にできる。
\item 左右点ファイバーなら、グラフの次元を問わず、全閉点が周期点である。
四移動 $\sigma$ の各閉点での周期は $1,2,3,4$ のいずれかであり、
$X_{\mathrm{red}}$ 上では $\sigma^{12}=1$ である。
\item 二つの周期列を、途中で $\ell\ge2$ 個の文字を通って切り替えられるとする。
すなわち最初の列のある状態からその後続文字と異なる文字で出て、
$\ell$ 番目では第二列の到着状態への周期からの入射文字と異なる文字で入る。
全隣接対が点対応に属することを要求する。それぞれの最小周期を $p,q$ とすると
\[
                 p+q+\min(\ell-1,4)\le16 .
\]
とくに順序を除く周期の組は $(4,4)$ または $(4,8)$ に限る。
\end{enumerate}
\end{theorem}
\begin{proof}[全特殊点に残る恒等式]
系\ref{cor:quartic-division-corners}により成分の一般折り畳み代数は次数16の中心単純環に入り、
各隅は次数4の可除代数に入る。分裂体でのAmitsur--Levitzki恒等式から
$B$ は $s_{32}$、各隅は $s_8$ を満たす。標準正則な特殊点には、この成分の一般点から
離散付値環を通じて至ることができる。一般・特殊の全経路次元が標準値なので
各次数の経路格子は自由である。したがって任意の有限個の経路の線型結合に恒等式を代入し、
分数体上での零を格子上の零として還元すれば、両恒等式が特殊点でも成立する。
特殊点の隅が整域であるという条件は使わない。

周期 $p$ の点列から、一周の各状態に一元を置き、文字をその点で評価して次の状態に移る
有限次元表現を作る。その像は $M_p$ 全体である。実際、同じ頂点で始まる異なる状態の
一周の評価テンソルは射影的に異なる（等しければより短い周期になる）。
各相手の一周を消して指定状態の一周で非零となる線型結合を選び、
相手の個数だけ一周の区間を繰り返して掛けると、指定状態だけに作用する非零対角単位が得られる。
隣状態どうしは非零文字で結べるので全行列単位が得られる。
$m$ 個の直交冪等元 $e_1,\ldots,e_m$ と各 $e_jAe_{j+1}$ の元で
全積が非零なものがあれば、$s_{2m-2}$ は消えない。
実際 $e_1,a_1,e_2,a_2,\ldots,e_{m-1},a_{m-1}$ を代入するとこの順の積だけが残る。
従って $p\le16$ である。

周期列の順分岐では、元の文字と比例しない別の後続文字を使い、
その一歩だけを新しい終端の一次元状態に着ける。
追加した長さ二の経路は一つ前の文字とこの後続文字を通るので、これは依然 $B$ の表現である。
行ベクトルの右作用で終端は頂点単純、商は前の周期単純である。この拡大は分裂しない。
実際、周期状態の持上げを終端の元の定数倍だけ変えたとき、
追加の係数を消すにはそれが周期の同じ位置の文字の倍数であることが必要になる。
像の二つの対角表現は非同型単純でそれぞれ全行列環なので、半単純商は $M_p\oplus k$ である。
非分裂性から非対角根基は非零で、単純な矩形行列両側加群なのでその全体となる。
従って像は二ブロック上三角環であり、上の行列単位列を $p+1$ 個まで取れる。
$p+1\le16$ と四の倍数性から $p\le12$ を得る。逆側は反対環で同じである。
\end{proof}
\begin{proof}[決定的な点列の非周期性の排除]
グラフの四移動で非周期な閉点があれば、その軌道の相異なる点は永久に一致しない。
連続する64文字を八文字ずつ八区間に切る。各区間の相対位置を一つ固定すれば、
そこで八区間に現れる点は同じ文字空間内で相異なる。
目標の一区間を残すため、七つの他区間それぞれに異なる相対位置を割り当て、
そこで他区間の点を消し目標の点で非零な文字を取る。余った一位置も目標で非零に取れば、
この長さ8の経路は目標だけで非零である。八区間全てを順に目標にして規格化し、
長さ8の八経路を得る。いずれも同じ四頂点の隅に属する。これらの $s_8$ を
元の64文字の点列で評価すると、八区間の評価行列は単位行列なので行列式は1である。
隅の恒等式に反する。よって周期的で、第一の上界から四移動十二回が全閉点を固定する。
被約有限型スキーム上の射は閉点上で決まる。
\end{proof}
\begin{proof}[二周期を結ぶ場合]
二つの周期の状態を別々に取り、最初の状態 $A_s$ から最後の状態 $B_0$ への
$\ell$ 本の矢と途中の $\ell-1$ 個の一次元状態を追加する。
二次関係は各二歩の係数が点対応で消えるため全て成立する。
各頂点から出る一文字ずつを全ての評価で非零に選んで和を $x$ とする。
各周期の各矢の代表倍率を選び、一方と他方で $x$ の一周の積の倍率を独立に調節できる。
$p,q,4$ の公倍数 $N\ge\ell$ を取り、$x^N$ の両周期上の倍率を相異なる非零値にする。
中間ブロックは冪零である。$x^N$ の三つの一般固有空間への射影 $E_A,E_B,E_T$
はその多項式で表され、元の四頂点冪等元と可換である。
各周期への圧縮環は、射影像の次元と周期商上の行列単位から $M_p,M_q$ 全体である。

最初の追加文字で非零、同じ状態での周期文字で零となる元を第一の矢に、
最後の追加文字で非零、$B_0$ へ入る周期文字で零となる元を最後の矢に使う。
中間の各矢では所定の点で非零な元を使う。最初の $k=\min(\ell-1,4)$ 個の途中状態の直後に、
その頂点冪等元と $E_T$ の積を挿入する。この $k$ 個は位置を法4で異にするから互いに直交する。
この積は $E_A$ と $E_B$ の間で依然非零である。実際 $A_s E_A$ は $A_s$ と
中間・終周期の補正だけからなる。第一矢で始周期に残る項は消え、中間の第一状態への係数は非零となる。
$E_T$ は中間ブロックで恒等であり、中間状態からの補正は終周期にだけ出る。
従って各段で所定の途中状態への係数を保ち、他の項はそれより先の状態か終周期にある。
最後の矢では終周期からの $B_0$ への係数は零なので、追加経路の非零係数が打ち消されない。
終周期では $E_B$ は恒等である。

両端を各圧縮環の行列単位で挟み、非零な一成分を選ぶ。
その前後にそれぞれの全対角単位を通る行列単位列をつなげば、
$p+q+k$ 個の相異なる直交冪等元を通る非零積になる。
17個以上ならその先頭17個で $s_{32}$ が非零となって矛盾する。以上で上界を得る。
\end{proof}

\begin{corollary}\label{cor:periodic-branch-central-basepoints}
任意の標準正則窓で、周期列の右点加群上に正次数の斉次中心元が非零に作用するとする。
その次数は列の最小周期の倍数であり、この列上には順分岐も逆分岐もない。
特に四次中心の一つでも非零に作用する周期列は長さ4で分岐がない。
分岐を訪れる全周期列は全正次数中心元の共通零点に限られる。
\end{corollary}
\begin{proof}
折り畳みの各頂点に成分を持つ中心元 $z$ の次数を $d$ とし、点加群の各次数での作用係数を $c_j$ とする。
矢との可換式は $c_j p_{j+d}=c_{j+1}p_j$ である。どの $p$ も非零なので、
一つの $c_j$ が非零なら全て非零で、列が $d$ 移動で戻る。右点加群は $k[z]$ 上で
最初の $d$ 個の同次元を基底とする自由加群である。従って右 $B$ 加群として有限表示である。
実際この基底各元と有限個の代数生成元の積を基底の $k[z]$ 線型結合で記すと、
これら有限個の関係で任意の語を簡約でき、簡約後の独立性は $k[z]$ 上の基底性から従う。
しかし順中心を繰り返せば定理\ref{thm:periodic-branch-obstruction}の第一Tor計算は
無限個の次数で新しい関係を与える。この計算は曲面の有無を問わない。
逆順で同じ周期の左点加群を取ると、$z$ の経路評価も同じ（開始位置を取り直すだけ）なので
左側の同じ議論で逆分岐も排除できる。
\end{proof}

% !TeX root = ../papers/nc-p3.tex
\section{Bell直線の十五次元成分との全交差と片側反転の既約成分}
\label{sec:split-line-component-incidence}
\begin{theorem}\label{thm:split-line-component-incidence}
半単純なBell直線の全補完のうち、定理\ref{thm:full-fifteen-component}の成分に入る添字保持同型類は
ちょうど二つである。一つは既知の位数四の量子窓、他方は $s^2=-1$ の反転なしの窓である。
片側反転族の基底群による移動の閉包は、それ自身が対称性を仮定しない二切断階数スキームの
既約成分をなし、射影次元59、半安定商の次元1である。この成分の商の関数体は $k(a^2+a^{-2})$ である。
\end{theorem}
\begin{proof}[中心の全跳躍と軸の四移動]
直線スペクトルを前節の通常基底で再構成する。反転なし $L=1$ と片側反転
$L=T(a,1)A$ の六関係は、全位置において二項式で、その係数比はそれぞれ
$s$ または $a$ の整数冪と四乗根の積である。
\path{computations/verify_split_line_center.py} では各始点で長さ4と5の全語を頂点とし、
関係に従う隣接二文字の書換えを辺とした。各連結成分で根の語を一つ選び、
各語から根への倍率を指数対 $\mathbb Z\times\mathbb Z/4$ で記録する。
全閉路で倍率が1であり、成分数は35と56である。従って得た全係数は
$\mathbb Q(\iota)[t,t^{-1}]$ 上の単元であり、$t\ne0$ の全特殊値で同じ基底に使える。

四位置の長さ4の各基底語に未知係数を置き、左・右から同じ一文字を掛けた係数を比較する。
固定した文字の片側乗法で相異なる基底語の像の台は相異なるので、得られる式は一変数を零とするものか、
二変数間の単元倍率の等式である。一つの連結成分に自由度が残るのは、
零を強制されず、閉路の全倍率が1のときだけである。全ての閉路差の最大公約式を、
零を強制された成分を除いて、単元倍と $t$ の冪を除いて並べると
\[
\begin{array}{c|c}
 L & \text{最大公約式とその重複度}\\\hline
 1 & 0\ (3),\quad t^2-1\ (6),\quad t^4-1\ (2)\\
 T(t,1)A & 1\ (14)
\end{array}
\]
である。検算は各辺で割った倍率と両隣接切断を元のテンソルから取り直す。
従って反転なしで $s^2\ne1$ の四次中心は $s^2=-1$ のとき5次元、他で3次元である。
片側反転では例外なく四次中心が零である。十五次元成分の全正則点では、
各長さの経路階数が一定であることと可換式の閉条件から、四次中心は五次元以上である。
残る反転なしの $s^2=-1$ の点は、定理\ref{thm:split-surface-component}に
$u=v=\iota$, $h=1$ としたものに等しい。実際 $(Z,1,Z,1)$ は対角型で、前節の基底変更で全て恒等となる。

$s=1$ で $CJ$ が非自明なJordanなら、それを上三角の標準単位Jordan行列とする。
最初の純色の軸上で、各辺の二文字関係を $r_i$ とすると、点の一移動は列ベクトルに
$Jr_i^t$ を作用させる。上の検算は、他方の $DJ$ が $\operatorname{diag}(t,1)$ または
非自明な単位Jordanの両場合について、四つの積が射影的に
$\left(\begin{smallmatrix}1&0\\-2&1\end{smallmatrix}\right)$ であることを確かめる。
これは標数零で無限位数であり、左右に一価な分解窓に関する定理\ref{thm:pi-periodic-point-paths}により
十五次元成分には入らない。量子部分については定理\ref{thm:quantum-component-intersection}を適用する。
\end{proof}
\begin{proof}[片側反転成分]
各 $a\ne0$ で六ブロックの各々に四つずつ非零語がある。一つの位置だけ異なる二語はない。
同じブロックでは二つの対応が単項であり、交代ブロックでも各ビットが二度現れるためである。
従って対角基底変更の後にも各位置の一文字Hermite Gramの非対角項は零である。
24語に一様な正重みを与えると各位置の四文字の周辺分布は一様になる。
ゆえに対角トーラス重みの凸包の相対内部に原点があり、対角軌道でHermiteノルムが最小になる基底を取れる。
全ての跡零方向の微分がこの基底で消えるので全特殊線形群でもアフィン軌道は閉である。
この議論は係数を含む有限生成部分体から複素数への埋め込みと基礎体の拡大で標数零に通用する。

$a=3$ を $109$ を法として $\iota=33$ に還元すると、アフィン基底軌道の
64個の行列単位による微分の階数は59である。接方程式の全200行を前と同様に使うと、
法線数順で接次元は $(0,16,28,16,0)$ である。非零の小行列式から、標数零の一般点で
軌道次元は59以上、全アフィン接次元は60以下である。前節の二次元トーラスと
自明なスカラー作用から軌道次元は59以下でもある。
$a$ の同値は有限なので、族と基底軌道を合わせた既約閉包の次元はアフィンで60である。
従って一般点で局所次元と接次元が一致し、閉包は全階数スキームの既約成分となる。
これを射影化すれば59次元である。一般の閉軌道の同値作用は $a\mapsto-a,a^{-1}$ で、
その不変体は $k(a^2+a^{-2})$ である。標数零で一般幾何ファイバーの次数はこの四つの
異なる点の個数なので、商の関数体もこの体に等しい。
\end{proof}

% !TeX root = ../papers/nc-p3.tex
\section{PBW経路格子の行列式による動く族の正則境界の排除}
\label{sec:pbw-valuative-boundaries}
\begin{theorem}\label{thm:pbw-valuative-boundaries}
$K$ を離散付値体とし、標準四窓が $K$ 上で標準Koszul左右AS正則とする。
四位置の文字にそれぞれ全順序を選び、各位置の四文字をその順位で標識する。
位置別左辞書順で六つの二次関係の先頭語が全ての降順対であり、弱い昇順語がPBW基底であるとする。
任意の長さ $n$ の語の正規形の各係数の付値が $-Cn^2$ 以上となる定数 $C$ があるとする。

$W_{i,j}$ を始点 $i$、長さ $j$ のKoszul交叉とし、各々を厳密降順語上の係数で座標化する。
この座標で最上テンソル $w_i$ の唯一の厳密降順語の係数を1に取る。
$w_i\in V_i\otimes W_{i+1,3}$ の四行四列の係数行列の行列式付値を $d_i$ とする。
任意の四位置の基底変更と有限体拡大の後に標準AS正則な特殊ファイバーを持ち得るためには
\[
                         \sum_i d_i\le0
\]
が必要である。これは条件を満たす\emph{全ての}文字順序に課せる。
\end{theorem}
\begin{proof}
各 $W_{i,j}$ の任意の非零元の先頭語は、各隣接対で二次関係の先頭語でなければならない。
実際、関係を置いた位置の外の文字列を固定した各係数列について既約な先頭語を取れる。
従って厳密降順語への射影は単射で、次元 $\binom4j$ が等しいので全単射である。
降順語の単位座標に対応するモニックな基底を取る。全長 $n>0$ のKoszul完全列の各項には
この基底と残りの昇順経路基底との積基底を入れる。この基底での完全列の行列式トーションは $\pm1$ である。
実際、全語の辞書順で各項を濾過すると、次数付き微分で残るのは降順頭部の最後の一文字を
後ろの昇順部へ移しても昇順のままのときの係数1だけである。既約化は全て語を下げ、
モニックな交叉基底の他項も語を下げる。固定全語に寄与する項は、その語が厳密降順の後に
弱い昇順となるときだけで、谷の一文字を左右どちらに含めるかの隣接二項が単位写像で結ばれる。
全語順にこのピボットを除けば主張を得る。

標準正則な格子が存在すると仮定する。その各項の整基底で同じ完全列を取ると、
剰余でも完全なのでトーションは単元である。分数体で二つの列を比較し、
参照PBW座標から整座標への $W_{i,j}$ の行列式付値の四始点の和を
$j=1,2,3,4$ に応じて $V,R,U,M$ と書く。長さ $n$ の経路での和を $b(n)$ とする。
両列のトーション付値が等しいので、定理\ref{thm:regular-orbit-valuations}と同じ
テンソル積の行列式の交代和から
\[
       \sum_{n\ge0} b(n)x^n =
       \frac{P(x)}{(1-x)^8},\qquad P(x)=Vx-Rx^2+Ux^3-Mx^4
\]
を得る。最上四基底は各始点で独立にモニックにしたので、その比較付値も合計 $M$ を使う。
正則格子の $1+3$ 対合の行列式は単元であり、
$(g_i\otimes g_{i+1,3})w_i$ と格子の最上元との倍率を行列式で比較すると
\[
                    V+U+\delta=4M,\qquad\delta=\sum_i d_i .
\]

一次格子の基本因子は両向きに一定範囲にあるから、参照文字の\emph{全語}が生成する
長さ $n$ の格子と正則側の経路格子との差は基本因子ごとに $O(n)$ である。
昇順語格子は全語格子に含まれ、逆には付値の仮定から一様化元の $O(n^2)$ 乗を
全語格子に掛ければ昇順語格子に入る。経路の階数が $\binom{n+3}3$ なので、
ある定数 $C'$ で $-C'n^4\le b(n)\le C'n^5$ を得る。
従って $P(1)=P'(1)=0$ でなければならない。この二式と対合の式を解くと
$R=3M-\delta$, $V=M-\delta/2$, $P''(1)/2=-\delta/2$ となる。
ゆえに $b(n)=-\tfrac12\delta\binom{n+5}{5}+O(n^4)$ であり、下界から $\delta\le0$ を得る。
有限分岐で付値は全て同じ正整数倍となる。
\end{proof}

\begin{corollary}\label{cor:quantum-no-polar-limits}
量子窓の動く族が、任意の基底変更の後に標準AS正則な特殊点を持つなら、
全ての量子三角積の付値は零である。従ってその特殊点は半安定であり、
その半安定商での像は三角積の剰余で定まる量子窓の像に等しい。
\end{corollary}
\begin{proof}
有限拡大後、定理\ref{thm:quantum-component-intersection}の対角twistで各行の交換係数の積を
1に取る。このとき通常の四位置反復がFrobenius標準である。任意の四文字順序を選び、
$j<k$ で $x_kx_j=c_{jk}x_jx_k$ とする。書換えは最大 $n(n-1)/2$ 回で終わる。
降順の最上語から文字 $j$ を先頭に持つ残り降順の語への係数は、符号を除き
$\prod_{k>j}c_{jk}$ である。従って $\delta=4\sum_{j<k}\operatorname{ord}c_{jk}$。
順序反転でこの和は符号を変えるため上の必要条件から零である。さらに任意の隣接二文字だけを
交換した順序と比較すれば各 $\operatorname{ord}c_{jk}=0$ である。
全係数が単元になったので、この量子表示自体が正則な格子を与える。
二つの正則格子の最上テンソルを特殊線形比較に直すと、定理\ref{thm:regular-orbit-valuations}で最上倍率は単元となる。
従って全斉次不変式の特殊値は共通の射影倍率で一致する。
量子窓は多重安定である（24置換語の一様分布で一文字周辺は一様、各Gramは対角）から結論を得る。
\end{proof}

\begin{theorem}\label{thm:all-mixed-flip-component-regular}
片側反転の射影59次元成分の\emph{全ての}標準AS正則点は
$a\in k^\times$ の同型類 $a\sim -a\sim a^{-1}$ で尽きる。
従って正則な同型類の集合は一つのアフィン直線 $a^2+a^{-2}$ であり、
商曲線の無限遠および不安定境界には正則点がない。
\end{theorem}
\begin{proof}
成分の任意の正則点を特殊点に取り、稠密な片側反転族からの曲線と基底変更を、
正規化と有限拡大後の離散付値環上で比較する。$a$ の付値が非零なら逆数同値で正としてよい。
元の文字を $0,1,2,3$ とし、前二文字の標識ビットを位置2,3で反転し、
後二文字の標識ビットを奇位置で反転する。得た標識の数値順を使う。
原稿の再構成から六関係は各辺の六降順対を全て異なる昇順対へ書き換える。
より詳しくは、標識0、標識1、残る二文字の三群の順序を先に整える。
異なる群の入替えは群の組を逆順にするだけで、二文字群内のビットだけが変わり得る。
最後にその群内を整えればよく、全体で $2n^2$ 回以内である。全次数の標準次元から昇順語は基底である。
係数は $a$ の有界な整数冪と四乗根の積なので定理の付値条件も満たす。
\path{computations/verify_pbw_valuative_boundary.py} は、各巡回の
最上語をモニックにし、第一文字と残りの厳密降順三文字の係数行列を直接取る。
四行列式の $a$ の指数は $(2,1,0,-1)$、従って $\delta=2\operatorname{ord}a>0$ となり矛盾する。

残る場合は $a$ が単元なので、片側反転の表示のまま正則に延び、その特殊点は前節で示した閉軌道に属する。
任意の正則模型の特殊点は、二正則格子の比較により半安定で同じ閉軌道を持つ。
この閉軌道の射影次元は全ての $a\ne0$ で58である。実際、上の検算が選ぶ
アフィン軌道微分の59次小行列式は $-64a^{13}(a^2+1)$ であり、
$a=\pm\iota$ でも法109で別の小行列式が非零となる。逆向きの上界は二次元トーラスから従う。
異なる半安定軌道がこれを閉包に含めば、その次元は59以上で、この既約成分で稠密となる。
これは成分上の不変有理関数 $a^2+a^{-2}$ が非定数であることに反する。
従って正則な特殊点は全て表示済みの元の閉軌道に属する。
\end{proof}

% !TeX root = ../papers/nc-p3.tex
\section{量子点を閉軌道に持つ全三表の正則特殊化による閉性}
\label{sec:quantum-regular-closure}
\begin{theorem}\label{thm:quantum-regular-closure}
定理\ref{thm:quantum-universal-chambers}の三表の和、すなわち閉包に量子点を持つ窓全体の
射影閉包を、標準Koszul左右AS正則窓に制限すると、元の三表の和に一致する。
より詳しく、離散付値環上で幾何学的一般ファイバーの閉軌道が量子窓で、
特殊ファイバーが標準正則なら、一般の量子三角積は全て単元である。
特殊窓は半安定で、その閉軌道は三角積を還元して得た量子類である。
従って、固定した一つの量子類への全流入も標準正則軌跡で閉じている。
\end{theorem}
\begin{proof}[一般ファイバーの濾過と全順序での行列式]
有限拡大後、一般ファイバーの四位置の基底と一径数群を選び、最下重み部分を量子体積とできる。
実際、三表の構成では閉軌道の共通対角トーラスの余指標で基点に流入する。
各点で必要な有限個の基底係数は代数閉包上で選んだ後に有限拡大上で定義できる。
全行積が1の量子代表を使い、基点を各始点で同じ四文字順列の台を持つ体積とする。
余指標の四文字重みを $l_j$、体積の重みを $L=\sum_j l_j$ と書く。

各長さでは、小さい余指標重みを先頭とし、同じ重みでは任意に選んだ文字順序の左辞書順を使う。
元の関係空間の最下重み付きの随伴次数は量子関係空間である。実際、テンソルを一径数群で
量子体積に特殊化する間に両切断の階数6が保たれ、像六平面もその像に特殊化する。
従って六つの降順対に対し、量子二項関係にこれより高い重みの項を足した関係が取れる。
書換えは同じ重みでは転倒数を減らし、補正項では重みを真に上げる。長さごとに停止する。
昇順語数と標準次元の一致から全順序でこれがPBW基底となる。

補正項を使う回数は一つの書換え経路で $n(\max l_j-\min l_j)$ 以下である。
連続する二文字を他の二文字に替える一回の補正で転倒数の増加は $2n$ 以下であり、
同じ重みの書換えは転倒数を1減らす。従って全経路の書換え回数は一様に $O(n^2)$ である。
有限個の規則の係数から各経路の係数付値が $-Cn^2$ 以上となり、和でも同じ下界が保たれる。
定理\ref{thm:pbw-valuative-boundaries}の証明はこの重み優先辞書順にもそのまま適用できる。
実際、同じ両側の語を付加しても順序が保たれるため、交叉の先頭語は厳密降順であり、
経路と交叉の語の濾過で得る単項式完全列も同じである。

$w_i$ を降順四文字上でモニックにする倍率は背景と同じである。
文字 $j$ を除いた降順三文字が $W_{i+1,3}$ の座標を与えるので、
$1+3$ 対合の行 $a$、列 $j$ の成分の余指標重みは $L+l_a-l_j$ である。
$l_a<l_j$ の成分は零、等しい重みのブロックは量子体積の成分そのものである。
ゆえに行列式は背景の対角項の積に等しい。これは重みが重複しても変わらない。
\path{computations/verify_filtered_quantum_valuations.py} は三表の各重みと重複重みで、
全24順序の行列式と隣接交換による付値の差を検算する。
よって一般の変形係数によらず、前節の量子窓と全く同じ
$4\sum_{j<k}\operatorname{ord}c_{jk}\le0$ が全順序で必要である。
逆順と隣接交換により全ての $c_{jk}$ は単元である。
\end{proof}
\begin{proof}[量子閉軌道の還元と全被覆]
全交換係数が単元になった背景の各体積係数も、降順四文字でモニックにすれば単元である。
残る有限個のテンソル係数は全て重み $L$ を超えるので、余指標に一様化元の十分大きな正冪を
代入してから共通の重み $L$ の倍率を除くと、それらは全て極大イデアルに入る。
これで\emph{同じ一般窓}のもう一つの整模型が得られ、その特殊窓は剰余量子窓で標準正則である。
定理\ref{thm:regular-orbit-valuations}で元の正則模型と比較する。
四つの行列式の四乗根を取って特殊線形比較にすると、二つの原始最上元の倍率は単元である。
従って特殊窓での不変斉次式は共通の射影倍率を除いて同値であり、
元の特殊窓も半安定で剰余量子窓と同じ閉軌道を持つ。三角積の還元に対する同型は
定理\ref{thm:quantum-indexed-all}の文字置換を許す判定に従う。

最後にこの軌跡は構成可能である。量子体積とその基底移動は有限型データであり、
量子に流入する条件は三表それぞれで重みの下の座標が零、同じ重みの座標が背景に等しい
という条件と基底変換の存在で表せる。逆にこの条件で二切断階数6を満たせば
定理\ref{thm:quantum-universal-chambers}の範囲に入る。
その閉包の任意の正則点には、稠密な構成可能部分から曲線を取り正規化して上の付値判定を適用できる。
従ってその点も三表に入る。量子三角積の組を有限置換まで固定した場合も、
同じ構成可能性と還元の一意性から最後の主張を得る。
\end{proof}

% !TeX root = ../papers/nc-p3.tex
\section{異なる二固有値を持つBell直線の全Jordan窓}
\label{sec:split-two-eigenvalue-jordan}
\begin{theorem}\label{thm:split-two-eigenvalue-jordan}
直線型で平方固有値が $1,s^2$, $s\ne0$, $s^2\ne1$、少なくとも一方が非半単純となる
全ての可逆分解窓は、次の表と下記の例外一類で尽きる。全点が標準Koszul左右AS正則で最小周期は1である。
表の正規形の不変量は四文字の量子三角積 $T$ と、指定した無向辺の集合 $E$ である。
\[
\begin{array}{c|c|c}
\text{長さ二の鎖} & a & E\\\hline
1 & k^\times & \{02\}\\
2 & 1,-1 & \{02,03\}\\
2 & s^{-2},-s^{-2} & \{02,12\}
\end{array}
\qquad
 T_{012},T_{013},T_{023},T_{123}= -as,\ -a/s,\ -as,\ -as^3 .
\]
表内・表間の無印同型の必要十分条件は、\emph{一つの}四文字置換が $E,T$ の双方を対応させることである。
三角積は奇置換で逆数にする交代関数として三文字の全順序に延長する。
残る一類は $s^2=-1$ で二鎖のときだけ現れ、閉軌道は片側反転の $a^2=1$ である。
表の閉軌道は全て表示した量子類である。十五次元成分との交差は、表の $s^2=-1$, $a^2=1$、
すなわち一鎖と二鎖の各一類で尽きる。全て半単純なBell直線の無印同型類と相異なる。
\end{theorem}
\begin{proof}[全ての四対応]
第一の交代成分の半単純化を前節の $P(s)$ とし、$E_{abcd}$ を二進語上の単位係数テンソルとする。
一鎖では固有値の順と共通倍率を選び、加える項を $E_{0011}$ とする。
二鎖ではこれに $E_{0110}$ または $E_{1001}$ も加える。それぞれ $+,-$ 型と呼ぶ。
実際、第一切断の二つの双曲平面は $00,11$ と $01,10$ の対である。
その各平面上でこの行列は対角 $1$ または $s$ に一つ以下の非対角行列単位を加えた形で、
右Gramの冪零部はその行列単位の非零倍となる。
自己随伴な二鎖の核は等方直線であり、双曲対の入替えと逆倍率でどちらの向きの非零係数も作れる。
二平面の向きを同時に反転することは特殊直交で、一鎖なら他平面で向きを合わせられる。
二鎖なら相対向き二通りを残す。右Gramの特殊直交同値と第一切断の行列式一致から
左側も特殊直交で同定される（定理\ref{thm:split-jordan-transport}）。

$P=P(s)+E_{0011}+\beta E_q$、$\beta=0,1$ とし、$\bar q$ を全ビット反転とする。
二逆転置の通常の成分は
\[
 I_0=P(1/s)-E_{1100}-\beta s^{-2} E_{\bar q},\qquad
 I_1=P(1/s)-s^{-2}E_{1100}-\beta E_{\bar q}.
\]
第一の四対応を恒等にしたとき、$Q=\pi_0 I_0$ に $\pi_1(\bigotimes_i L_i)I_1$ を射影的に合わせる。
比較は全て単項行列で行えるし、これで全てである。実際、差のある二固有値の一般固有平面の対を
右側および左側において保つか交換する必要がある。行列積の二側において各平面の二本の
可分解線が四つの二進座標線を定める。各逆行列の二つのブロックは対角または三角であり、
その可逆行列の台の完全マッチングは一つだけなので、四つの対角係数の台も保存される。
ここでは $\pi_1^{-1}\pi_0=(02)(13)$ が第一切断の二側を交換することを用いる。
従って必ず半単純部の比較にもなり、前の $T(c,d)$ と反転のリストを使える。

$F$ を伴う場合、第一の付加語を比較すると $cd=s^{-2}$、第二があれば
$+$ 型で $d/c=s^{-2}$、$-$ 型で $c/d=s^{-2}$ を得る。
$L=T(a,s^{-2}/a)F$ の全解が表の三行となる。反転なしでは第一の付加語の位置が合わない。
片側反転は $s^2=-1$ に限り、$s=\iota$, $\iota^2=-1$ を選ぶ。
一鎖ではその付加語が合わず、二鎖では $+$ 型で $L=T(e,\iota e)A$、
$-$ 型で $L=T(\iota e,e)B$, $e=\pm1$ のみが残る。
実際 $+$ 型の追加二式は $cd=\iota$, $c=-\iota d$、$-$ 型では
$cd=\iota$, $c=\iota d$ であり、反対の片側反転では台が合わない。
これで全表示を得た。正則性は定理\ref{thm:split-inverse-regularity}による。
\end{proof}

\begin{proof}[量子部分と追加二語による無印同型]
$F$ の行では両色のビットを位置2,3で反転する。六ブロックの非零語は24置換語と、
各辺 $uv\in E$ に対する $uvuv,vuvu$ だけである。置換語だけを残せば表の $T$ の量子窓となる。
一辺または共通端を持つ二辺なので、一つの共通端の文字だけを高重みにすれば量子窓へ流入する。

一鎖の場合、全位置に同じ対角行列を作用させ、$l_0+l_2=l_1+l_3$ とした二次元トーラスが
各位置で相異なる四指標を持つ。これは極大である。実際、分解トーラスの中心化群の階数は、
共通法線倍率の1に $P$ の安定群の階数を足した数以下である。
$P$ を保つ単項行列の対角部分は $T(c,d)$ で、$E_{0011}$ を保つため $cd=1$、
従って階数は高々1である。よって同型があれば極大トーラスの共役後に単項同型がある。
二鎖ならさらに $c/d=1$ が必要で、$P$ の安定群は有限、分解トーラス自体が極大である。
その共役後、同型は両側の二平面を保つか同時に交換する。後者でも旧 $Q$ は逆転置と
二側の交換から得られ同じ二固有平面と二鎖を持つ。各可分解線と第一の四対応から
全八つの二文字基底が線の順序・倍率まで定まるので、やはり単項同型で比較できる。

一辺の台では各位置の二文字の出現回数が7、他二文字が6で、追加語は7回の二文字だけを使う。
二辺の台では出現回数は $8,7,7,6$ で、追加語は6回の文字を含まない語として認識できる。
従って置換語の部分は保たれる。それを保つ位置別の置換は全位置で同じ置換である。
異なる二位置で相異なる入力文字を同じ出力にすることはできないからである。
これで $E,T$ の必要性を得る。

逆にその置換をした後、同じ $T$ の二つの置換語部分は位置別対角倍率で一致させられる。
これは三角積の等しい交換係数の比が文字倍率比となることと量子最上直線の一意性による。
相補二文字を $b,d$ とし、各辺の追加係数の比を測る次の量は、位置別の対角変更に不変である。
\[
   \frac{w_{uvuv} w_{vubd} w_{bdvu}}{w_{vuvu} w_{uvbd} w_{bduv}}=-1.
\]
相補文字の順序では量子切断内の階数1から値は変わらない。
\path{computations/verify_split_two_eigenvalue_jordan.py} はこの等式、$T$、および一巡回後も同じ値を
取ることを全行で有理関数として確認する。ゆえに各辺に残るのは両追加項共通の倍率だけである。
全位置に同じ対角基底を使うと、これは $d_ud_v/(d_bd_d)$ 倍になる。
二辺ならその二指標は独立なので、代数閉体上で二倍率を同時に所望の値に取れる。
これで十分性と、量子部分の一移動と同型なことを用いて最小周期1を得る。
\end{proof}

\begin{proof}[例外一類と成分]
二鎖の片側反転四表示を上の検算で相互に同定する。具体的には28個の非零座標で
位置別文字の出現回数を保つ置換だけを列挙する。対応後の係数比は全て四乗根である。
位置別の16個の倍率による作用の整数行列は語ごとにその四文字に1を置いた列を持つ。
Smith標準形で整数核の基底を取り、係数比の四乗根指数との積が全て法4で零なら
所望の倍率がある。四表示の各々と、代表の一移動についてこのような同型が得られる。
Jordan項は、$T(c,d)$ の適切な一径数で両項の重みを上げて消せる。
補完の半単純化は、二鎖の重みフィルトレーションによる定理\ref{thm:split-polystable-good-limits}
の操作である。ここでは四対応が単項で鎖フィルトレーションを運ぶので次数付きでも
同じ四対応が残る。より明示的には $+$ 型は $T(t^{-1},1)$ と旧第一対応で決まる相補作用、
$-$ 型は $T(1,t^{-1})$ を使い、両付加語は重み1となる。片側反転の四対応はこの変更で不変である。
この極限は半単純な片側反転の $a=\pm1$ の類であり、量子閉軌道と相異なる。
同じ分解において非半単純な $P,Q$ を持ち、階数1による分解の一意性から、
半単純直線の階数1のどの類とも同型にならない。一鎖の階数2では四指標が相異なり、
半単純直線の非量子階数2に現れる重複指標と区別できる。量子窓自身の階数は3以上である。

量子閉軌道を持つ行では定理\ref{thm:quantum-ore-incidence}により、十五次元成分に入る必要十分条件は
同じ成分の量子類を閉軌道に持つことである。四面の向きに合わせればこの条件は
$T=(\iota,-\iota,\iota,-\iota)$ またはその逆である。表の式に代入して $s^2=-1$, $a^2=1$ を得る。
この位数四の背景では偶置換が向き付けた面値を保ち、奇置換がその逆へ運ぶ。
無向一辺にも共通端を持つ無向二辺にも奇置換の安定化元があるので、各場合ちょうど一類である。
片側反転に流入する例外は、その閉軌道が十五次元成分に入らないため除かれる。
\end{proof}

% !TeX root = ../papers/nc-p3.tex
\section{単一固有値の全Jordan窓とBell直線の完結}
\label{sec:split-single-eigenvalue-jordan}
\begin{theorem}\label{thm:split-single-eigenvalue-jordan}
平方固有値が一つの直線型可逆分解窓のうち、非半単純なものの無印添字保持同型類は次で尽きる。
全て標準Koszul左右AS正則、最小周期1であり、十五次元成分には属さない。
\begin{center}
\begin{tabular}{cl}
鎖分割 & 同型類\\\hline
$4$ & 一点\\
$3+1$ & $b\in k$, $b\sim-b$\\
$2+2$ & $(u,v)\in k\times k^\times$, $(u,v)\sim(-u,v^{-1})$、および別の一点\\
$2+1+1$ & $a\in k^\times$, $a\sim a^{-1}$、別の一点、さらに $p\in k^\times$, $p\sim-p$
\end{tabular}
\end{center}
各行と各行内の別欄は互いに交わらない。全て量子閉軌道へ流入する。
$2+1+1$ の $a$ 欄は前節の一辺の表を $s=1$ に延ばしたものであり、他の量子閉軌道との
照合にも同じ辺と三角積による判定が使える。
\end{theorem}

\begin{proof}[斉一な移送と四つの中心化群]
四次元模型を二つの標準二文字空間の積とし、形式を $B=J\otimes J$ とする。
通常の四テンソルの第一切断行列を右から $B$ 倍して作用素として見る。
左右の局所射影群は二つの $\operatorname{PSO}(B)$、$\tau=(02)(13)$ は随伴を取る。
共通倍率と左特殊直交作用により $P$ の作用素を
$L=(1+N)^{1/2}=1+N/2-N^2/8+N^3/16$ に取れる。ここで $N$ は自己随伴、$N^4=0$ である。
これは定理\ref{thm:split-jordan-transport}の極分解で、単一固有値1と $\det P=1$ に規格化したものである。
$1+N$ の固有値は1なので、$P$ の射影安定群 $H$ は
$N$ を中心化する特殊直交群を $\{\pm1\}$ で割ったものの左右対角埋込みである。
実際、作用素の射影倍率の二乗が1となり、左右Gramの等式から右側は $N$ と可換、
左側も射影的に同じとなる。

双対を各二文字で $J$ により同定する。二逆テンソルの作用素は、それぞれ $L^{-1}$ と
$L^{-1}(1-N^2/2)=(1-N)^{1/2}$ である。長さ $m$ の鎖を $k[z]/z^m$、
$N$ を $z$ 倍、対合を積の $z^{m-1}$ 係数として
\[
 U_t f(z)=(1+tz)^{(m-2)/2} f\big(z/(1+tz)\big)\pmod {z^m}
\]
と置く。$U_tU_h=U_{t+h}$ であり、$U_t$ は直交ユニポテントで
$U_tNU_t^{-1}=N/(1+tN)$ を満たす。直交性は、積の最上係数を留数
$\operatorname{Res} f(z)g(z)\,dz/z^m$ と書き変数変換すれば従う。
よって $U_1(1-N)^{1/2}U_1^{-1}=L^{-1}$ である。
第一の四対応は双対同定後に恒等、第二の四対応は二位置反復で
$g=U_1h$ の二つの射影二文字因子を取ればよい。全ての解はこれと $h\in H$ で得られる。
下記の通り $U_1$ は $H$ と可換なので、印付き同型は $h$ の通常の共役である。

長さ4の一鎖なら、それと可換な行列は $N$ の多項式、直交条件はその二乗が1であるため $H=1$。
$3+1$ では三鎖に $1,z,z^2$、一鎖に $f$ を取り $\langle f,f\rangle=1$ とする。
$H$ は $1\mapsto f$, $f\mapsto-z^2$, $z,z^2\mapsto0$ なる $D$ の
$\exp(bD)$ 全体である。実際、$1$ の像を $c+c_1z+c_2z^2+b f$、
$f$ の像を $e z^2+d f$ とすると、直交性と行列式1より
$c=d=\pm1$, $c_1=0$, $2cc_2=-b^2$, $ce=-bd$ を得る。全符号を除けば主張通りである。
$U_t$ は定数項を保ち、$z^2,f$ を固定するので $D$ と可換である。

以下 $N^2=0$ では $U_t=1$ である。$E$ を二次の上三角行列単位 $E_{12}$ とする。
$2+2$ は $N=E\otimes Z$ とその二因子交換、$Z=\operatorname{diag}(1,-1)$ である。
両者は像の極大全等方平面の二つの族に対応する。$\operatorname{PSO}(B)=\mathrm{PGL}_2^2$ で
$E,Z$ を逆の倍率で保つ必要があり、後者の倍率は $\pm1$ に限る。
従って $H=(\mathbb G_a\times\mathbb G_m)\rtimes C_2$ で、対合は $(u,v)\mapsto(-u,v^{-1})$ である。
非恒等成分は平方写像の全射性から単一共役類となる。
$2+1+1$ は $N=E\otimes E$ である。二つの上三角群の対角倍率が互いに逆になり、
$H=\mathbb G_a^2\rtimes\mathbb G_m$、$t$ の重みは $1,-1$ である。
$t\ne1$ では二つの平行移動を共役で除ける。$t=1$ なら零ベクトル、二つの非零座標軸、
積 $p=uv\ne0$ ごとの一類に分かれる。これで印付きの全解を得た。
\end{proof}

\begin{proof}[一移動と分解の交換]
$S$ を二つの二文字因子を交換する行列とし $K=L-1$ と書く。
分解交換後の第一交代成分は $S L^{-1}S^{-1}$ である。実際 $\pi_0=(03)(12)$ は
第一切断の両側の交換と各側での因子交換であり、双対同定後の第一逆行列は自己随伴である。
一移動と分解交換を同時に行ったときは $S(1-K)S^{-1}$ となる。
実際、自己随伴作用素の空間はスカラーと $\mathfrak{sl}_2\otimes\mathfrak{sl}_2$ の直和で、
四文字の巡回は因子交換と第一因子だけの交代形式による随伴を取る。後者は跡零行列に負号を付ける。
$K$ は跡零自己随伴なので表示式を得る（体積の共通符号は射影的に消える）。

どちらの操作でも第二の四対応は $Sg^{-1}S^{-1}$ の二位置反復となる。
実際、双対同定後の四つの行列を $d_i$ とすると、分解交換では $d_{\pi_1(i)}^{-1}$、
一移動と交換では旧第二対応を新第一の座標に用いて $d_{i-1}^{-1}$ になる。
今は $d_{i+2}=d_i$ なので両者は同じである。
さらに $U_{1/2}(1-K)U_{1/2}^{-1}=L^{-1}$ である。これは $N^4=0$ の上で
$(1+N/(1+N/2))^{1/2}-1=K/(1+K)$ を展開した等式である。
$U_{1/2}$ は $g$ と可換だから、これを $S$ で運んで左右同じ基底変更に使えば、
二窓が一致する。従って無印最小周期は全て1である。
\path{computations/verify_split_single_eigenvalue_jordan.py} は各鎖の反対角形式、両移送、
通常二進座標での巡回・交換および $3+1$ の中心化作用を直接検算する。

指定分解の一径数を含む安定群の階数は、その中心化群の階数であり、$1+\operatorname{rank}C_H(h)$ である。
六ブロックの相対倍率の核が共通法線倍率で、残りは第一の四対応と $h$ の比較から従う。
階数1なら原始余指標は共役を除き符号だけである。$2+2$ の恒等成分内の $h$ では階数2だが、
第二因子の対角トーラスは奇位置の能動二文字と偶位置の相補二文字という一つの色にだけ非自明に作用する。
他の色の二文字は同じ重みなので、二文字ずつの二重みを与える余指標は元の分解とその逆だけである。

交換は $N$ を $S[-N/(1+N)]S^{-1}$ に運ぶ。$-N/(1+N)$ は元と同じ核の旗を持ち、
$\det S=-1$ なので、全鎖長が偶数のとき極大全等方平面の二族が交換される。
従って長さ4の二向きから一点、$2+2$ では片方の向きの印付き分類をそのまま得る。
$3+1$ では $1,z,z^2,f$ に対する $D_0=\operatorname{diag}(1,-1,1,1)$ を使う。
$D_0 U_{-1}S^{-1}$ は特殊直交で新しい $N$ を元に戻す。$D_0$ は $D$ を固定し $U_t$ を
$U_{-t}$ に変えるので、移送 $g$ の新しいパラメータは $-b$ である。
$2+1+1$ の非零平行移動では階数1で、交換は $(u,v)$ をまず $(-v,-u)$ にする。
さらに $-E\otimes E$ を $E\otimes E$ に戻す二対角倍率の積は $-1$ なので、
二座標軸は一類、積は $p\sim-p$ となる。
\end{proof}

\begin{proof}[残りのトーラスと他の行との比較]
$2+1+1$ の対角な $h$（恒等を含む）では、
通常二進座標で $P=P(1)+E_{0011}$ として $L_i^1=T(a,1/a)_i F_i$ を使える。
実際、この $P$ の第一切断は単位行列に一つの上三角行列単位を加えた形である。
$B$ を一側に作用させて随伴作用素に戻せば $1+N/2$ であり、二つの因子のビット変更と
対角倍率で $N=E\otimes E$ にできる。この変更と双対の $J$ はいずれも単項なので、
前節の二逆転置に $s=1$ を代入して得る上の解は、$H$ の中でも単項、従って上三角の中の対角になる。
$a$ の倍率を変えると両因子の対角比は互いに逆の任意の倍率で変わるので、対角な $h$ を尽くす。
全位置で相異なる四指標を持つ極大トーラスがあるから、前節の一辺の無印同型証明もそのまま使える。
この辺は $02$、全三角積は $-a$ である。置換 $(02)$ が全三角積を逆数にするので $a\sim a^{-1}$ となる。

階数1あるいは一つの色にだけ追加トーラスが作用する行では、分解は上で述べた通り一意であり、
$P,Q$ の鎖分割は等しい。相互の異なる鎖分割や半単純な $P$ と混同されない。
四指標が相異なる一辺の行では前節の表とも同じ有限置換判定で比較できる。
具体的に、向き付き四面値を $T_{123},1/T_{023},T_{013},1/T_{012}$ と並べると、
$s=1$ では互いに逆な二値が二個ずつである。前節の一辺でこの形を持つには既に等しい
第二・第四値以外に第一・第三値も等しく、従って $s^4=1$ が必要である。
$s^2=-1$ なら二値の積は $-1$ であって条件を満たさない。置換による全値の同時反転もこの区別を保つ。

最後にJordan鎖の標準重みで $N$ を零にすれば、$4,3+1$ では $g$ の次数付きは恒等、
$2+2$ では第一因子の平行移動が消え第二因子のトーラス元または法線化群の反転が残る。
$2+1+1$ では両因子で $E$ の重みが正なので平行移動は消える。
従って残る二因子はいずれも半単純である。$L=1$ でのこの二因子は、
通常座標 $P(1)$ に戻すと定理\ref{thm:split-spectral-line-classification}の $CJ,DJ$ の共役類である。
実際、第一二位置の $J$ で $B$ を単位行列にし、相補基底で第一の四対応を恒等に戻すと
最初の二行列は各因子の左に $J^{-1}$ を掛けたものとなる。
よって閉軌道は単一平方固有値の量子類であり、定理\ref{thm:split-line-component-incidence}により
十五次元成分とは交わらない。
\end{proof}

% !TeX root = ../papers/nc-p3.tex
\section{Bell曲面の固有値衝突における全正規形・周期・成分交差}
\label{sec:split-repeated-surface}
\begin{theorem}\label{thm:split-repeated-surface}
可逆分解窓の曲面スペクトルで平方固有値に重複があり、直線スペクトルには属さないとする。
$t\ne0$, $t^4\ne1$, $f=(t^2-1)/(2t)$ として、以下の三行が全ての無印添字保持同型を与える。
三行は互いに交わらない。
\begin{enumerate}
\item 半単純・反転なしでは $u=f^2\in k\setminus\{0,-1\}$ が完全不変量である。
最小周期は $u=1$ で1、他で2である。
\item 半単純・反転ありでは $(u,v)=(f^2,fc)$, $c\ne0$ を取り、
$(u,v)\mapsto(u,u/v),(u^{-1},v/u)$ で割る。最小周期は $v^2=-u$ で1、他で2である。
\item 非半単純では $v=\pm f$, $v\ne0$, $v^2\ne-1$ 自身が完全不変量である。
全点の最小周期は2で、閉軌道は第二行の $(v^2,v)$ である。
\end{enumerate}
全点が標準Koszul左右AS正則である。十五次元成分に入るのは第一・第三行の全体と、
第二行の $v^2=u$ だけである。
\end{theorem}
\begin{proof}[全補完]
曲面の平方固有値は $1,r,s,rs$、$r,s\ne0,1$ なので、衝突は $r=s$ または $rs=1$ である。
共通倍率・置換・偶符号でBell四数を $(1,-t^2,t,t)$ に取れる。
実際、重なる二つの平方と残りの二平方の積の関係からこの規格化を得て、
行列式の符号は曲面式が決める。残り二平方の衝突は直線上の $t^2=-1$ である。
$P$ の通常二進成分は $0000,0011,0101$ と補語に順に $(1-t^2)/2,(1+t^2)/2,t$。
非半単純ならさらに $E_{0110}$ を加える。第一切断の奇ビット対平面で一鎖を作り、
残りの奇数長の鎖で向きを調整できるのでこれで十分である。

$P$ の安定群は、半単純なら $T(c,1/c)$ と全反転 $F$ の生成群、非半単純なら
そのうちの $T(\pm1,\pm1)$ だけである。実際、唯一の重固有平面とその二本の可分解線は
左右で単項作用を強制する。残りの二固有線は $00\pm11$ でそれぞれ保たれる。
従って正規化した対角作用は $T(c,1/c)$、反転は全ビット同時だけであり、付加語があれば
その係数と位置の保存から後者が消え $c^2=1$ となる。
$D_0=(Z,1,Z,1)$ として、半単純の全第二対応は $T(c,1/c)F^eD_0$ である。
$D_0$ が移送になるのは定理\ref{thm:split-surface-component}の $h=1$ の表示による。
対角基底は $e=0$ で $c$ を自由に変え、$e=1$ では自明に作用し、$F$ は $c$ を逆にする。
付加語があるときも重固有平面から全局所行列は単項で、両逆テンソルの奇ビット対ブロックは
三角になる。その唯一の完全マッチングを保つため半単純部も移送する。
付加項は $I_0,I_1$ に順に $-t^{-2}E_{1001},-4(t^2-1)^{-2}E_{1001}$ として現れる。
両切断で比較すると $e=1$, $c^2=1$ が必要十分である。
この表示を $T(z,z^{-1})$ の比を用いて付加項だけ消せば同じ $t,c$ の半単純部へ流入する。
半単純部の閉性は、各ブロック内の補語に同じ正重み、各ブロックに等しい全重みを与えれば
一様な周辺を持ち、一文字だけ異なる非零語もないことから従う。
\end{proof}

\begin{proof}[無印同型]
反転なしは対角移送に対する固定群が有限なので全安定群の階数1であり、指定分解は交換を除き一意である。
逆数Bell四数に移る交換は $t\mapsto1/t$ となる。重固有値の規格化と二単純固有線の交換から
許される $t$ は $\pm t^{\pm1}$ であり、これらは全て同型になる。
実際、元の四数を共通倍率、置換、偶符号で移せば、移送側は階数1の唯一の半単純補完である。
よって不変量は $u$ である。一移動と分解交換では旧 $P$ を一巡回するため、
新しいBell四数は $((1-t^2)/2,(1-t^2)/2,(t+1)^2/2,-(t-1)^2/2)$ である。
従って $u$ は $1/u$ に移り第一行の周期を得る。

反転ありでは、位置2,3でビットを反転する。非零語は24置換語に
$0202,2020,1313,3131$ を加えたもの、Jordanの場合はさらに $0303,3030$ を加えたものである。
半単純では全位置共通の対角トーラス $l_0+l_2=l_1+l_3$ が極大で各位置に異なる四指標を持つため、
全同型を単項同型で比較できる。非半単純では階数1で分解が符号まで一意、
重固有平面の二本の可分解線が各局所二文字の座標線を決めるため、交換の場合を含めやはり単項で足りる。
台で隣接対として現れないのは等しい二文字だけなので、位置別の許される置換は四位置共通である。
従って第二行ではマッチング $02,13$ を保つ八置換、第三行ではさらに辺 $03$ を保つ
恒等と $(03)(12)$ の二つに限られる（四頂点における二端の次数も保存する）。

対角同値を具体的に計算する。$w_q$ を語 $q$ の係数として、二つの不変量は
\[
 u=\frac{w_{0213}w_{1302}}{w_{0312}w_{1203}},\qquad
 v=-\frac{w_{2031}w_{0132}w_{1230}}{w_{0231}w_{2130}w_{1032}} .
\]
六ブロックの式でこれらは $f^2,fc$ となる。
\path{computations/verify_split_repeated_surface.py} は位置別文字の28列の整数重み行列の核を取り、
各核指標の値を $t,t-1,t+1,t^2+1,c$ の指数ベクトルに写す。
像の格子の基底は $(-2,2,2,0,0),(-2,0,0,2,0),(-1,1,1,0,1)$ である。
従って指標値は定数を除き $u,1+u,v$ の冪積で、同じ $u,v$ で元の二表示は対角同値である。
八置換を行うと対角同値の範囲で $t$ は $t,(t+1)/(t-1)$ に、$c$ は $c,c^{-1}$ に独立に移る。
上の検算は各比が全整数核上で1であることを符号も含めて確認する。
これが表示した二つの生成変換である。一移動を合わせるとさらに $c$ が $-c$ になる。
従って無印では $v\mapsto-v$ であり、$u=-1$ は除外されているから第二行の周期条件を得る。

Jordanの場合、二置換による第二行の比較では $u,v$ ともに保たれる（$c=\pm1$）。
付加二語の相対倍率について、前節の辺 $03$ の六係数比はやはり $-1$ である。
よって $u,v$ が等しければ、マッチングの対角トーラスの残りの非自明指標で
二項共通の倍率も一致させられる。これで第三行を得る。一移動同型は閉極限の周期から存在せず、
二移動では付加語を含む全整数核上でも係数比が1となる。
最後に第二行の別分解の可能性を確かめる。その極大トーラスの四指標から、
二重みずつの分解は $01\mid23$ と $03\mid12$ だけである。
後者の交代部分の $00,11$ ブロックは反対角で二係数の積 $f^2$、
$01,10$ ブロックの二つの平方固有値は $((t^2+1)/(2t)\pm1)^2$ である。
従ってこちらの印でも直線外の同じ重複型である。第一・第三行では印が符号まで一意で、
逆転置もこの重複型とJordan分割を保つ。これで直線部や単純平方スペクトルとの重複も起きない。
\end{proof}

\begin{proof}[十五次元成分との全交差]
反転なしと $c=\pm1$ の半単純窓は曲面の四族に $u=v=t$ を代入して得られる。
Jordanの二表示も同じ成分に入る。実際、曲面一般の四数 $(1,-xy,x,y)$ において
$(x-y)/2=z^2$、$x,y\to t$ とし、能動側に $T(z,z^{-1})$ を作用させる。
$0110$ の係数だけが1として残り、補語は零になる。$F$ を伴う二表示の四対応は
この対角変更で不変なので所定のJordan窓になる。

逆向きに半単純反転窓では、両色の純軸で純二文字関係を $r_i$ としたとき、
$Jr_i^t$ の四移動の積は射影的に $\operatorname{diag}(1,c^{-2})$ である。
これは $T(c,1/c)FD_0$ の四対を純ブロックで順に取った式である。
従って十五次元成分に入るなら定理\ref{thm:pi-periodic-point-paths}により、
$c^2$ の位数は $1,2,3,4$ である。$c^2\ne1$ の候補の最小多項式は
$c^2+1,c^2+c+1,c^2-c+1,c^4+1$ に限る。
各剰余体上で隣接二切断の接方程式を取り、$t(t^4-1)\ne0$ の全域で
アフィン接次元は65である。具体的に各辺の六降順対をピボットに取ると
残り十座標から核基底が得られる。ピボットの可逆性は純二式が交換式、混合四式が
二方向の間の同型であることから全域で成立する。
二切断のそれぞれについて両核基底の100積を作り、トーラスの二指標別の25ブロックに分ける。
上の検算は表示中の分母が指定開集合で単元であることを調べ、
一変数多項式環上のSmith形の非零不変因子が全てこの開集合で単元となること、
零方向の次元和が65であることを確認する。十五次元成分のアフィン次元76を下回るため、
これらの候補は例外 $t$ を含め全て除かれる。
\end{proof}

% !TeX root = ../papers/nc-p3.tex
\section{Bell曲面の全単純スペクトルと可逆二平面分解の分類}
\label{sec:split-simple-surface}
\begin{theorem}\label{thm:split-simple-surface}
$x=(1,-uv,u,v)$, $f(t)=(t^2-1)/(2t)$, $A=f(u)^2$, $B=f(v)^2$ とする。
平方が相異なる可逆曲面部の基底は $uv\ne0$, $AB\ne0$, $A\ne B$ である。
以下が全ての無印添字保持同型類である。
\begin{enumerate}
\item 通常の四表示を $h=1,X,Z,XZ$ の順に $0,1,2,3$ と標識する。
$(u,v,h)$ に対して以下の三変換とそれらの合成が全ての同型を与える。
右欄は旧番号 $0,1,2,3$ に対する新番号の順である。
\[
\begin{array}{c|c}
(u,v)\mapsto(-u,v)&3,2,1,0\\
(u,v)\mapsto(1/u,v)&1,0,3,2\\
(u,v)\mapsto(v,u)&2,1,0,3
\end{array}
\]
各無順序対 $\{A,B\}$ の上には、どちらも $-1$ と異なれば四類、
一つが $-1$ なら二類がある。全て最小周期2で、十五次元成分に属する。
\item $\{A,B\}=\{-1,b\}$, $b\ne0,-1$ のとき、各 $b$ にさらに一類がある。
これは十五次元成分に属さず、最小周期は $b=1$ で1、他で2である。
この一変数族の基底軌道の閉包は、制約なしの階数スキームの射影60次元既約成分で、
各 $b\ne0,-1$ の軌道で滑らかである。
\end{enumerate}
いずれも標準Koszul左右AS正則かつ多重安定である。
\end{theorem}
\begin{proof}[基底の同型と移送商]
平方の集合は $\{1,u^2,v^2,u^2v^2\}$ であり、可逆性の下では平方の衝突は $A=B$ と同値である。
この点は三つのBell曲面のうちただ一つにある。例えば $ab+cd=ac+bd=0$ なら
非零座標から $b^2=c^2$ を得る。ゆえに曲面内の基底比較は
$ab+cd=0$ を保つBell付け替えに限られる。四座標の二対分割を保つ八置換と射影偶符号四通りで位数32、
$(u,v)$ 上では独立な符号・逆数と交換である。その作用の生成は下で実際の二文字基底でも与える。
従って基底の完全不変量は $\{A,B\}$ である。

左右の直交四次元模型で単純平方スペクトルを持つため $H$ のうちこれらの固有線を個別に保つ部分は
同時の二文字Pauli四元群 $H_0$ である。固有線に非自明な置換が生じるなら
$P$ の射影倍率の四乗は1で、その二乗は $-1$ でなければならない。
平方集合が負号で不変となるのは $u^2=-1$ または $v^2=-1$ のときだけである。
実際 $u^2v^2=-1$ だけで他の二平方も負号で組を作れば $u^4=1$ となる。
従って一般には $H=H_0$、通常の四対応 $(hZ,h,hZ,h)$ が四共役類を尽くす。

例外では $u=\iota$, $\iota^2=-1$ に取れる。$E=(1,X,Z,XZ)$ が $P$ の射影直線を保ち、
$H$ は $H_0$ と $E$ の生成する位数8の可換群である。位数の上界は固有線を個別に保つ核と
可能なただ一つの負号置換による。$D_0=(Z,1,Z,1)$ とすると通常基底の第二対応全体は $H D_0$ である。
$H_0$ は比較作用で自明に働くが、$E$ は四箇所同じ $Z$ を掛ける。
従って $D_0,FD_0,ED_0,FED_0$ が印付きの相異なる四類である（$F$ は全位置の $X$）。
分解交換では第一交代成分が逆数Bell四数に移る。
$(1,Z,Z,1)$ でこれを元に戻すと、通常表示は各々自分に戻る
（定理\ref{thm:split-surface-full-fibers}の基底等式は特殊点でも成立する）。
追加表示は $ED_0\leftrightarrow FED_0$ となる。
$H$ は有限で分解トーラスが階数1の極大トーラスなので、分解の符号と共役を除いて追加同型はない。

通常表示において、三つの基底生成変換を持ち上げる能動基底を順に
$(1,1,XZ,XZ),(1,X,X,1),(D,D^{-1},D^{-1},D)$, $D=\operatorname{diag}(1,\iota)$ とする。
各々を $P$ と四対応に作用させると表を得る。
\path{computations/verify_split_simple_surface.py} はこの三つの射影等式を $\mathbb Q(\iota,u,v)$ で検算する。
付け替え群の作用は一般の $(u,v)$ で自由であり、$u^2=-1$ の固定変換 $u\mapsto-1/u$ は
ちょうど同時 $Z$ による二番号の同一視を与える。これで無印分類を得る。
各 $P,Q$ は補語に等しい係数を持つBell型、四対応はPauli型なので、
各一文字のGramの非対角成分が消え、各ブロック内の両ビットの周辺は等しい。
既出の六ブロックのノルム最小化により多重安定であり、正則性は可逆分解の定理による。
\end{proof}

\begin{proof}[全ての一移動]
$\phi(t)=(t-1)/(t+1)$ とすると、巡回したBell四数を先頭1に直したパラメータは
$(-1/\phi(v),\phi(u))$ となる。一移動と分解交換では第二対応が $L_{i-1}^{-1}$ となり、
通常表示の番号は同時 $Z$ 倍になる。二移動で第一成分は元に戻り、第二成分は通常表示で一致し、
追加表示でも $E$ の比較で消せる同時 $Z$ 倍にしか変わらない。
一移動があり得るのは $\{A,B\}=\{1/A,1/B\}$ の場合だけである。
$AB=1$ かつ $A\ne B$ なら、付け替えで $v=\phi(u)$ に取れる。
このとき一移動後の $(u,v)$ は既に元通りだが、番号が $Z$ 倍で変わり、基底に固定変換はない。
残りの候補は $\{A,B\}=\{-1,1\}$ だけで、$u=\iota,v=1+\sqrt2$ に取る。
一移動は $(-v,\iota)$ と $Z$ 倍へ移り、交換と第二変数の負号で基底を戻すと、
二つの通常番号類 $\{0,2\}$ と $\{1,3\}$ が入れ替わる。
追加類はその基底上で唯一の追加類なので固定される。これで全周期を得る。
\end{proof}

\begin{proof}[追加表示の成分]
通常表示は定理\ref{thm:split-surface-component}による。追加表示では分解の交換が
印付きの\emph{別の}類へ移るので、この窓に指定分解を逆転する自己同型はない。
射影安定群は多重安定性から簡約的であり、もし恒等成分が階数1のトーラスより大きければ
非自明なWeyl反射がこれを逆転する。従って射影軌道次元は $60-1=59$ である。
$b$ は非定数でその上に一類なので族の閉包は射影60次元である。

$u=\iota$, $v(v^4-1)\ne0$ において $ED_0$ の二切断のアフィン接次元は常に61である。
上の検算では二色内の純関係のビットに合わせ、六降順座標の射影で関係を表示し、
二切断の左右の核の全200積を取る。分母は $v(v^4-1)$ の局所化で単元である。
同一方程式に現れる列をつないでブロックを分け、
$\mathbb Q(\iota)[v]$ 上のSmith形で非零不変因子の根が全て許容域外にあることと、
核次元和61を調べる。これは上の軌道族のアフィン次元であるから全点で滑らかで、
通る成分はこの軌道族の閉包一つだけである。特にアフィン76次元の既知成分の点はない。
\end{proof}

\begin{corollary}\label{cor:all-invertible-split-atlas}
可逆二平面分解を持つ全窓の添字保持同型類と最小周期は、
直線の半単純表・二固有値Jordan表・単一固有値Jordan表、曲面の重複表・単純表、
および六次式の一点で尽きる。全ての最小周期は1または2である。
\end{corollary}
\begin{proof}
補完可能な全スペクトルと鎖分割は定理\ref{thm:split-bell-locus}、
\ref{thm:split-jordan-transport}の三領域で尽くされる。
曲面と直線の共通部分は直線の表で処理した。
異なる印があり得るときは、極大トーラスでの分解判定、または四指標が異なるときの単項同型判定を
各表で使った。量子窓では座標四文字を二対に分ける他の印もあるが、
どの可逆分解も、対応する一径数を量子窓の共通対角極大トーラスへ共役すれば二対分割に戻るので、
その交代成分は常に半単純直線型である。表外の重複は生じない。六次式の一点とは第三スペクトルで区別できる。
\end{proof}

% !TeX root = ../papers/nc-p3.tex
\section{片側反転への全流入と二成分の交点の局所環}
\label{sec:mixed-flip-slice}
\begin{theorem}\label{thm:mixed-flip-slice}
半単純直線の片側反転 $s=\iota$, $L=T(a,1)A$ を閉軌道とする二切断階数スキームを考える。
$a^4\ne1$ では、全基底群の半安定商のこのファイバーは一軌道だけである。
$a^2=1$ では閉軌道と、定理\ref{thm:split-two-eigenvalue-jordan}の例外Jordan窓の
ちょうど二軌道で尽きる。双方とも最小周期1である。
この閉点での軌道横断面の完備局所環は
\[
                   k[[z,x,y]]/(zx,zy).
\]
直線が片側反転成分、平面が定理\ref{thm:split-simple-surface}の追加成分である。
二つは軌道方向を除いて一点だけで交わり、埋込み成分はない。
残る $a^2=-1$ の横断接表現は自明指標一つと、下記の相異なる七指標各一つである。
\end{theorem}
\begin{proof}[全パラメータの横断接指標]
片側反転窓で最初の色のビットを位置2,3だけで反転し、この色を二本の文字 $x_0,x_1$ に分ける。
他方の二文字を $Y$ とする。全非零語に両 $x$ が一つずつ現れる。
二つを独立に拡大するトーラスは射影安定群の恒等成分である。
実際、定理\ref{thm:all-mixed-flip-component-regular}の軌道微分の全値での階数から
恒等成分の次元は2である。語の横断指標は $(n_0-1,n_1-1)$、
$n_j$ は $x_j$ の出現数である。

二切断の接方程式をこの指標ごとに $\mathbb Q(\iota)[a,a^{-1}]$ で取る。
\path{computations/verify_mixed_flip_slice.py} は純関係に合わせて第二色のビットも奇位置で反転し、
各辺の六降順対をピボットにした左右の核から200積を作る。分母は $a$ の冪と定数だけである。
Smith形の一般核次元は和60、自明指標で20である。
非単元不変因子の全ての根と、各根で一次元ずつ増える指標は
\[
\begin{array}{c|l}
a^2=1 & (1,-1),(-1,1)\\
a^2=-1 & (-1,0),(0,-1),(-1,2),(2,-1),(0,2),(1,1),(2,0)
\end{array}
\]
である。59次元のアフィン基底軌道のうち自明指標の次元は19である。
実際、各位置の同重みブロックのLie環次元は $1+1+4=6$、
四位置の線形逆像でアフィン安定群の次元は $64-59=5$ で、全て重み零に入る。
他の指標でも軌道次元は一定である。よって横断接の全指標は、一般の自明指標一つと
表の該当箇所の和である。

閉軌道でLuna横断面を取り、安定群の極大トーラスで点へ流入する部分を考える。
量子三表の証明と同じく、強エタールな飽和開近傍での商の基底変換により
この基点のファイバー全体が横断面の零ファイバーで得られる。
接余接の指標座標を同変に持ち上げれば、それらは局所環の完備化を位相的に生成する。
一径数で基点に流れる弧では重み零・負の接座標は全て零であり、使える座標は正のものだけである。
$a^4\ne1$ では何もなく、閉点だけを得る。
$a^2=1$ ではどの余指標にも正となる接座標は高々一つで、二つの符号は互いに反対である。
\end{proof}

\begin{proof}[横断的な実現と全局所式]
同型な片側反転表示に移り、$E=(1,X,Z,XZ)$、$D_0=(Z,1,Z,1)$ とする。
第一対応は恒等、第二対応は $ED_0$ とし、第一交代テンソルを
\[
 P=P(\iota)+x(E_{0011}+\iota E_{1001})+
                  y(E_{1100}+\iota E_{0110})
\]
とする。原点の近傍ではこの全平面で両Schur式が成り立つ。
実際、追加曲面表示 $(1,-\iota v,\iota,v)$ の最初の係数を1に直すと、
上の二括弧の係数が共に $(1+\iota v)/(1-\iota v)$ になる。
$T(1,d)$ がこの係数を $d^{-1},d$ 倍し四対応を保つため、$xy\ne0$ の稠密な点で表示を得る。
原点でも両隣接逆行列の分母は単元で、Schur式は平面上で延びる。
座標軸の非零点は第二対応の片側反転を保つ二鎖Jordanであり、前記の例外一類になる。
一方、$x=y=0$ のまま $L=T(1,a)ED_0$ を動かすと片側反転曲線が得られる。
$z=a-1$ をその局所座標とする。

三接方向が軌道横断で独立で、混合二次障害 $zx,zy$ も独立であることを示す。
上の検算は原点のテンソルを二つのSchur式で再構成し、平面と曲線を各座標で微分する。
法109で $\iota=33$ とし、64個の基底行列単位の列に順に $z,x,y$ の微分を加えた階数は
$59,60,61,62$ である。各切断行列 $M$ の単元ピボット上の一般逆を $K$、
左右核を $L_0,R_0$ とし、二つの一次変分を $S,T$ とする。
Schur式の混合二次係数を消すには、二切断を合わせた接方程式の像に
$L_0(SKT+TKS)R_0$ がなければならない。
全接行列の階数は上で標数零でも194であり、法109で $zx,zy$ の右辺を順に加えた階数は
$194,195,196$ である。よって標数零の余核で二列は独立である。

滑らかな軌道座標を除くには、周囲の射影空間において軌道接の補空間と群の局所座標で
形式同型を取り、上の曲線・平面を横断面に射影する。
接が独立なので、三つの極小形式座標の中でこの二つの滑らかな像を座標直線と座標平面に同時に直せる。
実際、曲線と平面の各パラメータ表示の和から共通原点を引いたものは可逆な一次項を持つ。
従って横断面のイデアルは $(zx,zy)$ に含まれる。他方、接行列の独立な行から
余分な周囲座標を形式的に消去した後、残る方程式の二次項は上記の余核での障害である。
独立な二式が $(zx,zy)/\mathfrak m(zx,zy)$ を張るので、中山の補題からイデアルは全体に等しい。
これで完備環と二成分への所属を得る。

二つの正指標方向の各々には実際に平面の座標軸上のJordan表示がある。
横断面の同じ正接座標を持つ流入点は一意である。実際、トーラスで規格化して値を合わせ、
流入弧に沿って他の全座標を完備環の冪級数として比較すれば二弧は一致する。
非零値どうしはトーラスで移り、二軸は既に同じ無印Jordan類である。
従って全基底群の商ファイバーには他の軌道はない。
\end{proof}

% !TeX root = ../papers/nc-p3.tex
\section{片側反転の最後の例外への全流入}
\label{sec:mixed-flip-fiber}
\begin{theorem}\label{thm:mixed-flip-fiber}
$w$ を $s=a=\iota$ の片側反転窓とする。全基底群の半安定商で $[w]$ の上にある
窓は全て標準Koszul左右AS正則である。同型類は以下の七座標の集合を以下の作用で割ったものに等しい。
座標を $b=(A,C,B,E,F,D,G)$ とし、この順の二次元トーラス指標を
\[
 (-1,0),\ (-1,2),\ (0,-1),\ (0,2),\ (1,1),\ (2,-1),\ (2,0)
\]
とする。許す台は $ABC$, $AC EF$, $CDEFG$ の部分集合、およびそれらの
反射 $A\leftrightarrow B,C\leftrightarrow D,E\leftrightarrow G$ である。
トーラス作用のほかに二作用 $Q,J$ で割る。
$Q$ はこの順に $(-1,-1,-1,1,1,-1,1)$ 倍、
$J$ は $(\iota,\iota,\iota,-\iota,1,\iota,\iota)$ 倍した後に上記の反射を行う。
一移動 $S$ は $(1,-1,-\iota,-1,-\iota,-\iota,1)$ 倍である。
最小周期は $Sb$ が $b$ と同値なら1、他は2である。
同じ五空間内外の同型・この周期条件はいずれも次の有限判定でよい。
$b'$ を $b,Qb,Jb,QJb$ のいずれかと比較し、台が一致し、
台上の指標行列の整数核の基底の各ベクトル $m$ について $\prod_j(b'_j/b_j^{*})^{m_j}=1$
となるか調べる。ここで $b^{*}$ は選んだ比較元である。
固定軌道の半安定な特殊化はこの同値を許し、台の指標が張る錐の任意の面の外の座標を零にしたものに尽きる。
\end{theorem}
\begin{proof}[五空間の構成と被覆]
前節の文字順で実際の四テンソルを与える再帰を記す。基底係数は $\mathbb Q(\iota)$ に置き、
共役転置による直交補をこの体で選んだ後 $k$ へ拡大する。
アフィン基底軌道の接空間を $O$、二切断の全接核を $K$ とする。
周囲の軌道横断面を $w+O^\perp$ に取り、$K\cap O^\perp$ の各一次元指標で
ガウス消去の自由成分を1にしたベクトル $b_j^0$ を取る。$O$ はスカラー方向も含む。
横断成分への直交射影を七座標と重み零の一座標に使う。
この選択は下記の位相付き単項行列と巡回に同変であり、トーラスについては重み空間が直交する。
周囲の作用写像の接は軌道に沿った因子を取れば同型であり、
閉軌道のLuna横断面定理で飽和開近傍に縮めると商からの基底変換として
この横断面を使える。基点へ縮む全軌道はその開近傍に入る。
安定群の恒等成分はトーラスなので、流入点はある余指標で正となる接座標だけを使う。
七つの内積の符号を平面で並べると極大の正台は上記五つで、
余指標は $(-3,-1),(-1,-3),(-1,2),(2,-1),(1,1)$ に取れる。
元の余指標での流入弧では周囲の座標が正の接座標の形式級数で決まり、
各項の指標はその非負和である。従って台を含む極大欄の余指標でも流入する。

各欄の座標を不定元とし、補正テンソル $D$ の単項式係数を総次数順に求める。
一次は $b_j^0$ とし、高次は $K^\perp$ のみに取る。各切断で $M$ を背景行列、
$L,R$ を左右核基底、$g$ を六つの独立行・列の交点上のピボット逆とする。
$MgM=M$ であり、六階数のSchur式は
$L(1+Dg)^{-1}DR=0$ である（ここでは $D$ は該当する切断を表す）。
指数 $m$ に対して $B_m=D_m-\sum_{u+v=m}D_u g B_v$ とすれば、
二切断の $LB_mR=0$ を同時に解けばよい。
最初の一次方程式 $D\mapsto (LDR)$ の像上で、$K^\perp$ への逆は
独立な行を $H$ として $H^*(HH^*)^{-1}$ で与える。
全ての行に戻して残差の消滅を確認する。各 $M,g$ は行列の源と先の
指標分解を保つ（独立行・列を選ぶ各ブロックは斉次）。
正の余重みしか使わないので $Dg$ は冪零であり、必要な単項式は各周囲語の指標に達するものだけで尽きる。
従ってこの再帰が通れば多項式による大域的な正規表示である。

\path{computations/verify_mixed_flip_fiber.py} は五つ全てで上の逆を取って
200行の残差が係数ごとに零になることを $\mathbb Q(\iota)$ 上で検算する。
この多項式表示の接射影は七つの入力そのものである。重なる台で高次補正は同じ再帰で一意である。
逆に、基点への流入弧の横断表示ではこの射影で全ての形式係数が一意に決まり、
従ってこの多項式表示が台を覆う。収縮の原点で六階数なので軌道上でも六階数、
命題\ref{prop:unrestricted-good-reduction}から全次数で正則となる。
\end{proof}
\begin{proof}[全安定群、一移動と軌道境界]
恒等成分の二文字・一文字・一文字の指標分解は内在的である。
従って二平面分解を保ち、二つの一文字を置換できる。
定理\ref{thm:split-spectral-line-classification}の交代成分 $02|13$ の二本の可分解直線と
第一の四対応により、二重文字の側も各位置で単項行列でなければならない。
ゆえに各位置各色の二文字の入替えの全 $2^8$ 通りを試せば十分である。
残る対角倍率は24個の支持語 $q$ と16個の位置文字 $(i,j)$ の行列
$N_{q,ij}=[q_i=j]$ で決まる。$N$ の非零Smith因子は1が十個と2が一個であり、
有限対角核はちょうど $Q$ で生成される（位置スカラーの積1は全空間に自明）。
上記検算では支持の一致後、係数比を四乗根で測り、Smith左変換で像への所属を判定し、
右変換で対角倍率を一組復元する。入替えで残るのは恒等と全位置の第一色のビット反転だけで、
後者の持上げが $J$ である。巡回したテンソルから同じ方法で一つの移送を選ぶと $S$ を得る。
全ての持上げは絶対値1の位相付き置換であり、スカラー規格化した背景を固定し、
横断接射影に上記の通り作用する。これで高次の座標変更も不要である。
$J^2$ はトーラスと $Q$ が生成する対角群に入り、$Q$ は反射不変である。
$S^2$ はトーラスの $(-1,1)$ と $Q$ の積に等しい。
代数閉体上で対角トーラスの像は指標の全整数関係で定まるので有限判定を得る。
横断面の商ファイバーと全基底群の商ファイバーの軌道は対応し、
その特殊化順序も一致する。対角トーラス軌道閉包の面による記述を有限個の移送と合わせれば最後の主張となる。
\end{proof}

% !TeX root = ../papers/nc-p3.tex
\section{二平面分解の正則特殊化と追加単純成分の全点}
\label{sec:split-regular-limits}
\begin{theorem}\label{thm:split-regular-limits}
分数体上の窓が可逆二平面分解を持つか、または可逆二平面分解のある閉軌道へ流入するとする。
任意の基底変更後に標準正則な特殊点を持てば、その特殊点は半安定で、
商で対応する唯一の閉軌道も可逆二平面分解を持ち標準正則である。
従ってこのような閉軌道への全流入の合併は正則軌跡で特殊化に閉じる。
\end{theorem}
\begin{proof}[整模型の体積と四つの純成分]
全体テンソル $w$ の $Y_iY_{i+1}X_{i+2}X_{i+3}$ 成分を
$\kappa_i s_i\otimes r_{i+2}$ と書く。これを色が混合した二組の間で切った四次正方行列を $C_i$ とする。
その行列式は符号を除き $\kappa_i^4(\det s_i\det r_{i+2})^2$ である。
有限拡大後、$X_i$ の各位置の高位文字を選ぶ際、四つの $r_i$ で
高位・高位の係数を同時に零にできる。実際、双対射影直線上の高位係数を読む点を、
各非退化形式の直交で次の点へ送り、四つの射影同型の合成の不動点から始めればよい。
$Y_i$ でも同様に選ぶ。二色を $X<Y$ とし、$r_i,s_i$ の高位・低位の係数を $r_i',s_i'$ と書く。
どちらも非零である。全色順の降順対が先頭語のPBW表示となる。まず色だけを並べ替え、
次に純二文字の降順対を昇順にする。前者は $O(n^2)$ 回であり、後者で $10$ を $01$ に
するたび逆転数が減り、$00$ にするたび高位文字数が減って逆転数の増加は高々 $n$ である。
よって各項の書換え経路が $O(n^2)$ で、各係数の付値に定理\ref{thm:pbw-valuative-boundaries}の二次下界がある。

四つの巡回における最上降順語の係数は $p_i=\kappa_i s_i' r_{i+2}'$。
モニックな最上元での $1+3$ 行列の行列式は符号を除き
\[
 \frac{\kappa_i^2\kappa_{i+1}^2\det(s_i)\det(r_{i+3})(r_{i+2}')^2(s_{i+1}')^2}{p_i^4}.
\]
実際、先頭が $Y$ の二行では残りが $YXX$、先頭が $X$ の二行では $YYX$ であり、
各 $2\times2$ ブロックを上記純成分の降順対で縮約した式である。
四行列式の付値和 $\delta$ は $\prod_i\det(s_i)\det(r_i)/(s_i'r_i')^2$ の付値となる。

一般の正則格子の座標への基底変更 $g$ を有限拡大と四位置のスカラー倍により
各位置特殊線形に取る。$gw$ の最小係数付値を $\ell$ とする。全巡回で同じであり、
正則側の最上格子は原始的なので、参照モニック基底との比較付値和は
$M=4\ell-\sum_i\operatorname{ord}p_i$ である。
定理\ref{thm:pbw-valuative-boundaries}の証明中の $V=M-\delta/2$ に $V=0$ を入れれば
\[
              16\ell=\sum_{i=0}^3\operatorname{ord}\det C_i .
\]
特にこれは純関係代表の倍率や高位文字の選び方に依らない。
\end{proof}
\begin{proof}[閉軌道への流入を一般ファイバーに許す場合]
一般テンソルを $z$ とし、一般体の有限拡大と位置別の基底変更後に、
$\lambda(u)z\to w$ ($u\to0$) とする。ここで $u$ は元の離散付値の一様化元とは独立であり、
$\lambda$ は特殊線形群の余指標である。半安定商の閉軌道への一径数判定によりこの形を使える。
$w$ を保つ $\lambda$ と可逆分解の余指標は、後者の印を $w$ の安定群で共役すれば
同じ極大トーラスに入る。このため $X,Y$ を $\lambda$ 不変に取れる。

二色各々について、純二文字のモニック書換えに負付値係数を生じないよう、
さらに $\lambda$ 固有基底を選べる。実際、非退化純形式の直交が定める射影直線同型に沿って
基点の基底を順に運べば、前三辺の形式はスカラー倍を除き交代形式になり、最後の一辺だけに
四同型の合成 $B\in\mathrm{PGL}_2$ が残る。この輸送は $\lambda$ の射影作用と可換である。
基点で作用が非自明なら $B$ はそのトーラスと可換なので同じ二固有線を持つ。
作用がスカラーなら自由にJordan基底を取れる。半単純な場合、高位と低位を交換して
最後のモニック規則 $10\to c\,01$ で $\operatorname{ord} c\ge0$ に取る。
非半単純なら作用はスカラーで、通常のJordan行列の代表を選んで
$10\to c\,01+d\,00$ の両係数を整にできる。スカラー $B$ も含む。
こうして純書換えの回数は付値下界に影響しない。

$z$ の関係を降順対の六列で行消去する。その小行列式の非消滅は $w$ での
非消滅と $\lambda$ 流入から従う。Grassmannチャートのグラフ係数は固有座標だから、
各降順対の右辺で同じ余重みの部分は $w$ の規則、残りは全て高い余重みを持つ
（行が張る関係部分空間への作用で右辺と左辺の文字重み差を取る）。
低い余重みを優先して同重みでは元の左辞書順にすると全規則は降下する。
$z$ は正則背景への良い還元で標準正則なので次元比較から昇順語は全次数で基底である。
最上元の各巡回の降順語は $w$ に非零係数を持つので、$z$ でも同じ $p_i$ である。
また $1+3$ 行列式では、背景の非零な置換項があるから、各置換項に現れる
四語の余重みの総和は背景重みの四倍である（行文字と列三文字の多重集合が一定）。
背景より高い項を一つでも含めば背景より低い項も要り、後者は存在しない。
従って四つのモニック行列式も全て変わらない。

長さ $n$ の各書換え経路上、余重みが上がる回数は、全語の余重みの幅から $O(n)$ である。
一回の増加で二色の逆転数は高々 $2n$ 増える。同重みで異色を動かせば逆転数が1減り、
同色の背景書換えでは不変である。従って同重みの異色交換も計 $O(n^2)$ 回に限られる。
背景の同色係数は全て整で、それ以外の各一歩の係数付値は一定値以上だから、
全語の正規形係数に $-Cn^2$ の下界を得る。定理\ref{thm:pbw-valuative-boundaries}の
トーションの議論はこの余重み優先順でも成立する。中央の微分の初期形はやはり
降順頭部と昇順尾部の単位写像であり、固定長における順序は有限で連接と両立するからである。
ゆえに $z$ の正則特殊模型への特殊線形比較での最小付値 $\ell$ は
先の体積公式の右辺（$w$ で計算）で決まる。
$z$ と $w$ の斉次特殊線形不変式は最初の余指標極限で等しいため、
次の整従属性の議論はこの同じ $\ell$ で適用できる。
\end{proof}
\begin{proof}[分解を保つ群の商への延長]
$G=\prod_i\mathrm{SL}_4$、各位置の $X,Y$ に重み $-1,1$ を与えるトーラスを $T$、
その中心化群を $C$ とする。周囲の四テンソル表現を $V$ とし、$V^T$ 上の
$C$ 不変式環は $k[V]^G$ の制限上整であることを確認する。
零錐を比較すれば十分である。実際、$C$ 軌道閉包の閉軌道において、複素数上で
最大コンパクト群に関する極小ベクトルを取ると、各位置の一文字Gramは $T$ 不変性で
二色間の非対角成分が零である。中心化群のモーメントが零なので各位置のGram全体がスカラーとなり、
全群でも極小ベクトルで閉軌道である。従って $G$ 零錐の制限は $C$ 零錐に等しい。
この代数的な結論は標数零代数閉体へ基底拡大で移る。
有限生成された二つの次数付き不変式環で、正次数制限のイデアルの根基は
$C$ 商の頂点イデアルだから、その剰余は有限次元である。
斉次な剰余基底を持ち上げ次数帰納すれば、制限環上の加群有限性を得る。

$\pi^{-\ell}w$ の $G$ 不変式は、正則特殊模型における原始 $z$（分解を初めから持つ場合は原始 $w$）の係数から整である。整従属性により全 $C$ 不変式も整であり、
その剰余が $V^T/\!/C$ の一点を定める。そのファイバーの閉軌道代表を $v$ とする。
$\det C_0\det C_2,\det C_1\det C_3$ はいずれも次数8の $C$ 不変式である。
各位置で一方の因子に $X$、他方に $Y$ の行列式が同じ指数2で掛かるからである。
両不変式は整で、その積は体積公式から単元だから $v$ でも四つの $C_i$ は全て可逆である。
全二切断が階数6以下、四純成分が純対どうしの切断で階数1以下という閉条件も $v$ に残る。
実際、これらは $C$ 不変な閉集合でその良い商への像は閉で、像上の閉軌道はこの集合にある。
従って四純成分は全て非退化二次行列の積であり、両隣接切断は色別に少なくとも $(1,4,1)$、
一文字も全階数である。定理\ref{thm:split-inverse-regularity}で $v$ は可逆分解の正則窓となる。
前段の極小ベクトル論法で $v$ は $G$ についても閉軌道である。
正則模型の原始特殊テンソルは $v$ と全 $G$ 不変式が等しく、
$v\ne0$ により半安定で、同じ閉軌道に対応する。
\end{proof}

\begin{theorem}\label{thm:extra-simple-all-regular}
定理\ref{thm:split-simple-surface}の射影60次元の追加単純成分の
標準正則点は、$b\in k\setminus\{0,-1\}$ の表示済みの各一類と、
$b=-1$ における片側反転の閉軌道一類・その例外二鎖Jordan一類で尽きる。
$b=0,\infty$ に正則な特殊化はない。
\end{theorem}
\begin{proof}
追加表示 $(1,-\iota t,\iota,t)$, $L=ED_0$ を用いる。
$\infty$ に向かう $t$ は逆数で $0$ にでき、$t\to\pm1$ も符号と
一添字移動の $t\mapsto(t-1)/(t+1)$ および逆数で $0$ に移る。
従って $t$ の付値 $e>0$ を扱えばよい。四位置の文字の標識を順に
$(0,1,2,3),(1,0,2,3),(0,1,3,2),(1,0,3,2)$ とする。
モニック $1+3$ 行列式の積は1で、四つの降順語係数は $((1-t^2)/(2t),1,1,1)$ となる。
従って正則な特殊点があれば上の特殊線形規格化で $\ell=-e/4$ である。
ところが反対切断 $02|13$ の行列式は次数16の全群不変式で
\[
 \det w_{02|13}=
 \frac{(t^2+1)^4(t^2-2t-1)^2(t^2+2t-1)^2}{256t^8}.
\]
\path{computations/verify_split_extra_boundary.py} は四位置の関係の先頭語、体積と
この恒等式を有理関数体で検算する。原始模型に直した不変式の付値は $-8e-16\ell=-4e<0$ で矛盾する。

既約成分の任意の点を稠密な表示からの曲線の特殊点とし、必要なら分数体を有限拡大する。
残る $t$ は単元で剰余は $\pm1$ でなく、$b=f(t)^2$ が $-1$ なら
定理\ref{thm:mixed-flip-slice}の平面で正則な閉模型に延びる。他は元の単純表示で正則かつ閉軌道に延びる。
二正則格子の比較によりその商ファイバーに限られる。
$b=-1$ の二類は同平面で実現される。他の閉軌道の射影次元は59で、
それを境界に持つ別軌道が成分にあれば次元60で稠密となり、
成分上の非定数不変有理関数 $b$ と矛盾する。
\end{proof}

% !TeX root = ../papers/nc-p3.tex
\section{通常Bell曲面の単純点と反転なし衝突点に至る全窓}
\label{sec:split-surface-all-fibers}
\begin{theorem}\label{thm:split-surface-all-fibers}
通常Bell曲面の全ての単純スペクトルの閉軌道、および重複曲面の半単純・反転なしの
全閉軌道を取る。これらの全流入ファイバーで、原点を除く無印同型類は
$\mathbb P(1^4,2^3)$ の補正Pauli四元群による商である。
全点が標準Koszul左右AS正則で、階数スキームで滑らかであり、十五次元成分だけが通る。
各非零軌道の半安定境界は指定閉軌道のみである。異なる閉軌道上の類は重ならない。
追加単純表示の滑らかな60次元成分では、単純スペクトルの全商ファイバーが閉軌道一つだけになる。

通常表示では、非零類のうち周期が2を割る軌跡の既約成分は全てのファイバーで
同型類次元 $4,2,2,2$ である。定理\ref{thm:split-surface-periods}の固有空間式が
各点での必要十分条件で、それ以外の類は最小周期4である。
閉軌道が最小周期2ならこの四成分上の全点は最小周期2である。
残る反転なし衝突の $f(t)^2=1$ の閉軌道では、
非零類で最小周期1の軌跡は次元 $2,1,1,0,0$ の五つの既約成分を持つ。
\end{theorem}
\begin{proof}[パラメータ全域での階数]
通常四表示は基底パラメータの符号・逆数・交換の持上げで $h=1$ の表示へ移せる。
衝突の反転なしは $u=v=t$, $t^4\ne1$ の $h=1$ を使う。法線文字を0,1、他を2,3とする。
二つの切断の左右核から取る全接行列を、まず法線数 $n$ で分ける。
さらに定理\ref{thm:split-surface-periods}の符号 $Z$ と補正ビット反転 $X$ を使う。
前者の符号は語の総ビットの偶奇 $p$ で定まり、後者の固有値は
$n=2$ なら $\epsilon=\pm1$、$n=3$ なら $\epsilon\iota$ である。
$X$ は各語を全補語へ送り、法線位置で $(\iota,\iota,-\iota,-\iota)$ の積を掛ける。
二つの隣接切断の行・列核基底は、四つの異色対と二つの純関係ピボットを用いて選べる。
それらへの射影が全域で同型であることは可逆分解と純行列の単項形による。

\path{computations/verify_split_surface_all_fibers.py} はこの核の全係数の分母が
$D=uv(u^2-1)(v^2-1)$ の因子だけであることを調べる。
両核を掛けた各接方程式で、各語を一方のビットが0の代表と全補語の固有和に制限する。
$(p,\epsilon)=(0,+),(0,-),(1,+),(1,-)$ の順に、各ブロックの階数の下界は次の通りである。
\[
\begin{array}{c|c|c}
 &\text{各ブロックの列数}&\text{階数下界}\\\hline
 n=2 &24&18,16,16,16\\
 n=3 &16&11,11,11,11\\
 n=3\text{のせん断作用}&4&4,4,4,4
\end{array}
\]
下二行では $D$ だけを逆にする。第一行は
$D(u^2-v^2)(u^2v^2-1)$ と $D(u+v)(uv-1)(u^2+1)$ の各主開集合上で別々に検算する。
後者は上記の全衝突 $u=v=t$ を含む。実際の証明書は、各行の単元分母を払って単元ピボットで消去し、
残る必要階数の全小行列式と $1-zD'$ が単位イデアルを生成することを有理係数で確認するもの
（$D'$ は当該開集合の式）である。せん断行列には各位置の第一文字だけへのせん断を
補語と固有和にして使う。全16本のせん断空間の各指標をこれで尽くす。

$n=4$ の純行列の接核は三次元である。実際、通常表示の法線の純二次関係は
四つのPauli関係でパラメータに依存せず、定理\ref{thm:split-surface-full-fibers}の一般値をそのまま使える。
軸交換の自己同型は全パラメータでの有理式で与えられ、分母は $D$ の因子だけなので
$n=0,1$ にも $n=4,3$ の次元上界が成立する。従ってアフィン接次元は高々
$3+20+30+20+3=76$。全点を通る既約なアフィン76次元成分があるから全ての上界は等号であり、
局所環は正則である。滑らかな一成分点という不変な開条件は流入窓にも移る。
\end{proof}
\begin{proof}[横断表示と周期層]
単純点で分解の向きを保つ安定群は法線倍率と補正Pauli群である。
$u^2=-1$ などで交代成分の安定群が増えても、追加元は通常の第二対応を同時 $Z$ 倍に移し、
指定した第二対応を保たない。反転なし衝突では $T(c,1/c)$ の二乗が第二対応に自明となるには
$c=\pm1$ が必要で、全反転と合わせてやはり同じ四元群である。
両向きのせん断の作用が単射なので、全安定群の恒等成分も指定トーラスだけである。
軸交換で逆方向の流入は順方向と同一視できる。
滑らかな閉軌道横断面で順方向の接指標は重み1が4個、重み2が3個である。
定理\ref{thm:split-surface-full-fibers}の正規表示の一次核とせん断余空間は定階数の束となり、
四法線式の余核で測る二次障害は一般点で恒等零なので全域で恒等零である。
これにより全パラメータで同じ重み付き射影商と全次数正則性を得る。
横断面でのトーラス軌道境界は原点のみである。
追加単純表示では接次元61、アフィン軌道次元60で、一つの横断方向はパラメータを動かす
重み零の方向である（同じ軌道次元の曲線族が全て多重安定）。
従って正の横断座標はなく、全流入が閉軌道だけで尽きる。

二移動の持上げは旧式の $\kappa=-(u^2-1)(v^2-1)/(4uv)$ の平方根を使う。
余空間と核が束になったので、標数零の各補正対合の固有次元は被覆上で局所一定であり、
既出の一般値がそのまま全点の値になる。平方根の符号は重み1側の正負を交換するだけである。
非零正方向を固定する比較は指定トーラスを逆にできないので、可能な追加の基底比較は上記四元群に尽きる。
四成分の相違の証明も対合の交換則と固有次元によるため全域に適用できる。

残る一移動は $t=1+\sqrt2$、$h=1$ に取る。重み1,2の座標を $b,c$ とする。
補正 $X$ を法線数 $n$ でさらに $\iota^{n-2}$ 倍して対合に取り直す。
補正一移動 $S$ は四語を左巡回した後に全文字2,3だけのビットを反転し、
旧語の各位置・文字に順に次表の係数を掛け、全体を $t^{n-2}$ 倍する。
\[
\begin{pmatrix}
 \iota&1&\iota t^2&-t^2\\
 -\iota&-1&\iota&1\\
 -\iota/t^2&-1/t^2&-\iota&1\\
 \iota&1&\iota&1
\end{pmatrix}.
\]
これが背景を保つこと、$S^4=1$ であることを直接検算できる。横断座標の切断は
法線トーラスと $X,Z,S$ を含む群に同変に取れる。この群は分解トーラスの有限拡大で簡約的である。
従って作用は $b,c$ 上の線型作用で決まる。$g=1,Z,X,XZ$ とし、
$b\ne0$ なら $gSb=\lambda b,\ gSc=\lambda^2c$、$b=0$ なら $c$ が $gS$ の固有ベクトル
という条件が一移動の必要十分条件である。
\path{computations/verify_split_collision_shift.py} は両切断の接核、せん断商、
各作用を $\mathbb Q(\sqrt2,\iota)$ で取り、全ての $gS$ が $b$ 上に四乗根各一つを持つことを確認する。
$c$ 上の固有値は $g=1,Z$ で $1,1,-1$、$g=X,XZ$ で $1,\iota,-\iota$ である。
各固有部分空間の組の包含を四元群の全像で比較すると極大なものの代表は
$(g,\lambda)=(1,1),(1,\iota),(X,1),(X,\iota)$ と、
$b=0$, $g=X$ で $c$ が $\iota$ 固有の一点である。
その $b,c$ の許容次元の対は $(1,2),(1,1),(1,1),(1,0),(0,1)$ となり、主張の五成分を得る。
\end{proof}

% !TeX root = ../papers/nc-p3.tex
\section{反転を伴う固有値衝突の全流入ファイバー}
\label{sec:reflected-collision-fibers}
\begin{theorem}\label{thm:reflected-collision-fibers}
定理\ref{thm:split-repeated-surface}第二行の閉軌道を用い、$t^4\ne1$, $t c\ne0$,
$f=(t^2-1)/(2t)$, $u=f^2$, $v=fc$ とする。
$c^2\ne1$ ならその上の全窓の添字保持同型類は、$u\ne1$ でちょうど五つ、
$u=1$ でちょうど三つである。正規表示の台は、二座標軸について各々
正負のどちらかまたは空を選んだものとする。一斉符号反転で割り、$u=1$ ではさらに
平面の四分の一回転で割る。特殊化は台を小さくすることとこの同値で尽きる。
$c^2\ne-1$ なら全て最小周期2である。$c^2=-1$ なら両軸を使って $u\ne1$ のときだけ
最小周期2、それ以外は1である。

$c^2=1$ では次の16個の相異なる指標を各一次元とする。
\[
 H=\{(\pm2,0),(0,\pm2),(\pm2,\pm2),(\pm3,\pm1),(\pm1,\pm3)\}.
\]
許す非零台は $H$ のうちある開半平面に入る任意の部分集合である。
これらの座標空間を指標通りの二次元トーラスと下記の明示的な有限作用で割ったものが
全無印同型類である。原点を含む全窓が十五次元成分に属し、標準Koszul左右AS正則である。
極大の座標空間は八次元が16個で、有限作用を除くと $u\ne1$ で四つ、$u=1$ で二つ、
各々の一般同型類の次元は6である。各軌道の半安定特殊化は有限作用と台の指標錐の面で決まる。
最小周期は下記の符号表と有限同型検査で2または4に決まり、各極大台の一般点では4である。
\end{theorem}
\begin{proof}[全ての横断指標と非例外値]
位置2,3でビットを反転した既出の28語表示を使う。語 $q$ の重みは
$h=(n_0-n_2,n_1-n_3)$ で、$n_j$ は文字 $j$ の個数である。
全基底群のアフィン軌道次元は59、自明指標の軌道次元は11である。
二切断の接核をこの軌道接で割ると、自明指標が常に二つあり、非自明な指標は
$c^2\ne1$ で四軸指標 $(\pm2,0),(0,\pm2)$ 各一つ、$c^2=1$ で上記 $H$ の各一つとなる。
\path{computations/verify_reflected_collision_fibers.py} は各指標で二切断の左右核の全積を行とし、
分数体での階数を計算する。一方 $tc(t^4-1)$ を逆にした全域で軌道行列の必要小行列式の
同時零点がないことを確認する。接行列はさらに $c^2-1$ を逆にした場合と
$c=1$ に制限した場合で同じ検査を行う。一移動で $c=-1$ に移せる。
下界検査は前節と同じ単元ピボット消去および小行列式イデアルの局所化による。
従って全ての階数が各指定範囲で一定である。アフィン安定群は自明スカラー三次元と
上記トーラスで恒等成分を尽くす。

非例外値では任意の正台を含む余指標は $(\pm1,\pm1)$ のいずれかでよい。
全正座標は余重み2で、最大語重みは4なので二次補正だけを要する。
$M$ の一般逆を $g$、左右核を $L,R$ とすると二座標の代表 $A,B$ の補正右辺は
$L(AgB+BgA)R$ である。検算は $\mathbb Q(t,c)$ で各軌道余空間の一本を選び、
両軸の平方および二種類の象限の混合積の全四列が対応する接行列の像に入ることを確かめる。
反対象限は全符号反転で移る。定階数の接核・余核上での二次障害は一般体上で零なので全域で零である。
余重み4に新たな横断接はなく、各象限で二成分とも自由である。トーラスは非零な二値を独立に規格化する。
\end{proof}
\begin{proof}[例外値で全ての正方向が積分できる理由]
全階数スキームの軌道横断面を取り、接射影の座標で滑らかな18変数形式空間に埋め込む。
二つの中性座標は $t,c$ で張れる。実際、28語上の既出の対角不変比 $u,v$ の微分は
全域で独立であり、軌道の重み零部分は対角作用である。
既約なアフィン76次元の十五次元成分は $c=1$ の閉軌道を通るので、
その被約横断面は純次元17である。形式的な軌道との積分解と有限型代数の完備化の
被約性・等次元性による。従って正則形式局所環でのこの横断面のイデアルは高さ1の
素イデアルの交わりであり、一つの非零元で生成される。
トーラス同変な生成元を選べる。実際、各有限jetで一次元の極小生成空間の
固有ベクトルを同じ重みで両立して持ち上げ、中山の補題を使えばよい。
この生成元の重みは零である。さもなくば中性二平面への制限が恒等零となる。
中性平面には上記二変数の固定点族が形式同型に入り、$c^2\ne1$ の一般点は
定理\ref{thm:split-repeated-surface}によりこの成分の外にあるので、これは不可能である。

従って中性座標を原点に固定し、一余指標で正になる座標のみを任意に取れば、
この重み零の方程式は定数項以外を持てず恒等零となる。これで全ての正座標が同時に積分できる。
実際のテンソルは定理\ref{thm:mixed-flip-fiber}のSchur再帰と同様、
軌道接の補および接核の補を同変に選び、正座標の単項式ごとに
$L(D_m-\sum_{a+b=m}D_a g B_b)R=0$ を解いて得る。
正余重みは有界なのでこれは多項式であり、入力の接射影は指定した座標そのものとなる。
Luna横断面での流入弧の一意性と閉軌道への一径数判定で全体を覆い、
重なりでは同じ台の同じ点になる。全次数性は正則な閉点への良い還元から従う。
\end{proof}
\begin{proof}[有限比較と周期]
四相異なる指標は全位置の四文字を区別し、恒等安定群の正規化群は各位置で単項である。
支持にない隣接対は同一文字の対だけなので各位置の置換は同じである。
$u,v$ の比較により、ある固定基点の置換は一般に $-1$（二軸の同時反転）までで、
$c^2=1$ なら軸交換、$u=1$ なら四分の一回転も許す。
これは前の八置換をこの二不変量に作用させた結果である。$c^2\ne1$ での個数が従う。
この場合二移動は置換なしで可能で、一移動は $c^2=-1$ でちょうど一軸の符号反転として実現する。
台の二つの位相はトーラスで自由に変わるので、台だけで判定して表示通りの周期となる。

例外値での有限作用を係数まで指定する。$N_{q,ij}=[q_i=j]$ を28語の行と16文字の列で作ると、
その非零Smith因子は11個全て1である。許した置換 $p$ を実装するには
比 $w_{p q}/w_q$ の列に、$N$ のSmith左変換の整数行を指数とする冪積を取る。
最初の11積を残し、残る対角パラメータを1とし、Smith右変換の整数行の冪積で16個の係数 $d_{i,j}$ に戻す。
この単項作用は語 $q$ を $p q$ へ移し $\prod_i d_{i,q_i}$ 倍する。他の解との差は恒等トーラスだけである。
各 $h$ の座標線の基底は、接行列の核の通常の自由列基底で軌道列の張る空間の外にある最初の列の剰余類
とすればよい。上の単項作用をその列に掛け、行消去で新しい重みでの軌道列を引いた一係数が
座標間の倍率である。横断面と接核への射影を同変に取るので、この剰余作用だけで全座標に作用する。
トーラスについては、同じ台での係数比の冪積を指標行列の整数核基底上で比較すればよい。
従って有限候補と線型演算・整数冪積のみの同型判定であり、二つの基点の表示が異なる場合は
前の $u,v$ の基点同型を一つ作用させて同じ判定を行う。

$H$ の八つの反対対から壁の16個の室を得る。各室にちょうど八指標が正である。
$u\ne1$ の四元群、$u=1$ の二面体群はいずれも室へ自由に作用する。
各室の全台の軌道次元は2なので表示した極大族と次元を得る。
トーラス軌道の閉包の面判定と有限比較から軌道の特殊化も従う。

$c=1$ の二移動は、語を二位置巡回した後、旧位置0,1の文字0,1に $f$、
旧位置2,3の文字0,1に $1/f$ を掛けて補正する。これは背景を保つ対合である。
接座標で負号となる $h$ は
\[
       (-3,1),\ (-2,0),\ (-1,3),\ (1,3),\ (2,0),\ (3,1)
\]
だけで、他は正号である。各指標の接商は全域で直線束なので、対合の符号は一定である。
上記検算は $t=2$ で軌道列と核基底からこの全符号を有理数上で求める。
従って座標にこの六符号を掛け、前段の有限同型検査を通れば最小周期2、他は4である。
$c=-1$ では一移動で $c=1$ へ移して使える。極大台を保つ有限置換は恒等だけで、
二つの軸での符号は $-,+$、それらの和または差の指標では全て $+$ なのでトーラス倍率では一致しない。
これで全台での最小周期と一般周期が決まる。
\end{proof}

% !TeX root = ../papers/nc-p3.tex
\section{反転なしBell直線の全流入と二平面分解を閉極限に持つ枝の被覆}
\label{sec:noflip-line-fibers}
\begin{theorem}\label{thm:noflip-line-fibers}
半単純直線の $s\ne0$, $s^2\ne1$, $L=1$ を閉軌道とする。
全流入窓は標準Koszul左右AS正則で、閉軌道一類と
$\mathbb P(1^4,2^3)$ の次の有限群による商で尽きる。各非零軌道の半安定境界は閉軌道だけである。
最小周期は1または2で、非零類の最小周期1の軌跡の既約成分次元は
\[
 \begin{array}{c|l}
 s^4\ne1 & 4,4,3,3,3,3,2,2\\
 s^2=-1 & 4,4,3,3,2 .
 \end{array}
\]
閉軌道の基点同型は $s\sim-s\sim1/s$ である。基点を揃えた後、重み付き射影座標に関する下記の
有限行列の軌道比較が全ての添字保持同型を決定する。
\end{theorem}
\begin{proof}[全正座標と同型作用]
$P_{abab}=1$ ($a=b$) または $s$ ($a\ne b$) の通常基底で六成分を再構成する。
各成分の四支持語に等しい正重みを与えれば一文字周辺が一様で、一位置だけ異なる二語もないので
この窓は多重安定である。指定した二平面の法線文字を0,1、他を2,3とし、その語中の法線数を $n$ とする。
\path{computations/verify_noflip_line_fibers.py} は $\mathbb Q[s,s^{-1}]$ 上で各切断の
左右核の積を取り、$n=3,4$ の接方程式の核次元が20,3で一定であることをSmith形で確認する。
使う核は単位座標を含み、分母は $s$ の冪だけである。背景の四純関係が単項可逆で混合ブロックも
可逆なので核の次元は特殊化で保たれる。せん断16列にも $s(s^2-1)$ の因子のみを持つ
非零な16次小行列式がある。残る重み1の四方向の各平方・混合積の両Schur右辺を
四法線係数の行列に加えても、一般体で階数は13のままである。
定階数の束の余核への二次写像が一般体で零なので全 $s$ で零であり、
全 $b\in k^4,c\in k^3$ が自由に取れる。定理\ref{thm:invariant-line-filtered-lifts}の
再構成により全てを実際の四テンソルに戻せ、閉点への良い還元で正則性が従う。

分解トーラスは階数1で、定理\ref{thm:split-spectral-line-classification}の軸交換が
指定分解を逆にする自己同型を与える。正負のせん断安定群は零なので恒等安定群はトーラスだけである。
従って無印の全流入が正方向で得られ、その非零点間の比較は向きを保つ群でよい。
この群の有限部分を、法線数 $n$ のテンソル上で以下のように作用させる。
$A,B$ は語中の文字0または2の個数についてそれぞれ $(-1)^{n_0},(-1)^{n_2}$ 倍、
$F$ は四位置の全ビット反転である。通常は $F^i A^j B^k$ の八個で尽きる。
実際、交代成分の $02|13$ 支持の二可分解線により全て単項で、
$L=1$ を保つ対角比較は二つの局所ビット比の二乗が1のものだけである。$s^4\ne1$ では
行列 $\left(\begin{smallmatrix}1&s\\s&1\end{smallmatrix}\right)$ の交差比から
独立な入替えは全反転だけである。$s^2=-1$ なら $s=\iota$ とし、さらに一つの作用 $E$ を加え、
八個の元と $E$ の左積を用いる。$E$ は各位置の文字2,3のビットだけを反転し、
旧語の各位置・文字に次表の倍率を掛ける。
\[
 \begin{pmatrix}
 -\iota&1&-1&-\iota\\
 \iota&1&\iota&1\\
 1&\iota&-\iota&1\\
 \iota&1&\iota&1
 \end{pmatrix}.
\]
これは28語の場合と同じ位置別対角比較を24支持語に適用するか、直接代入で確かめられる。
追加できる二つの独立な反転は $E,F$ で張られる。法線倍率は既に重み付き射影化に含まれる。
作用は三法線の接核をせん断列で割った $b$ 空間および四法線の接核 $c$ 空間にそのまま降ろす。
二次補正の切断をこのトーラスの有限拡大で同変に選ぶことで、この線型作用が全同型を与える。
\end{proof}
\begin{proof}[巡回の作用と全固有空間比較]
$S$ は左一巡回の後に文字2,3だけのビットを反転する。
旧語の位置0が文字2,3なら $s$、旧位置2が文字0,1なら $1/s$ を掛け、
最後に選んだ $r^2=s$ で $r^{n-2}$ 倍する。これは $n=2$ の背景を保ち、$S^4=1$ となる。
重み付き射影点 $[b,c]$ に周期 $m=1,2$ があるための必要十分条件は、前段の有限元 $g$ の一つで
\[
 b\ne0:\quad gS^m b=\lambda b,\quad gS^m c=\lambda^2 c;\qquad
 b=0:\quad c\text{ は }gS^m\text{ の固有ベクトル}
\]
となることである。二移動は全点で成立する。

計算を特殊値から全域へ戻す根拠を記す。$s^4\ne1$ では $r$ を変数とした連結な
開集合上で上の四元 $A,B,F,S$ が生成する作用は一つの有限2群を通る。
実際 $S$ の倍率は色だけによるので、$F$ と可換、$A$ と可換で $B$ を
法線倍率 $(-1)^{n-2}$ 倍に移す。$F$ も $A,B$ を共役で高々この中心符号だけ変える。
各生成元の位数は4以下である。$s=\iota$ では位相は八乗根、各位置内のビット反転と
巡回で作る置換も2群なので、ここでも全生成群が有限2群となる。
固有空間成分表は、二つの表現上のこの群の指標で決定される。
特に二つの固有空間の包含を調べるとき、その共通固有空間は対応する二作用の生成部分群の
一次元指標空間である。指定した二固有値が群関係に反すれば零であり、他は
指標冪等元で次元が決まる。これらは群の位数を逆にし必要な一の根を加えた
定階数格子上で直和因子になる。従って連結なパラメータ域で次元は一定であり、
良い奇素数への還元でも同じ包含表を得ることができる。

上記検算は通常行で $s=4,r=2$、例外行で $s=\iota$ を取り、法257ではそれぞれ
$(s,r)=(4,2),(16,4)$ とする。両接核とせん断商で次元を直接検査し、
各 $gS^m$ の全固有空間を取る。各語置換の位数は8以下、例外の位相は八乗根なので
六十四乗根までで足り、通常の生成関係でも十分である。法257にはこれらの根が全て存在する。
$b$ の固有空間と $c$ の固有値二乗空間の対、および $c$ のみの対から、
群作用の後に他の対に含まれるものを除く。$m=2$ では $k^4,k^3$ の対だけが残る。
$m=1$ では通常行の極大対の次元は $(2,3)$ 二つ、$(2,2)$ 四つ、$(2,1)$ 二つ、
例外行は $(2,3),(3,2),(2,2),(2,2),(2,1)$ である。
各々で射影化の一を引けば主張した成分を得る。
\end{proof}

\begin{corollary}\label{cor:all-split-attraction-atlas}
半安定商の閉軌道が可逆二平面分解を持つ全ての窓について、正則性、添字保持同型および
最小周期を決める有限な表と検査が得られる。この集合は正則特殊化で閉じ、
片側分岐中心を持つ正則窓とは交わらない。
\end{corollary}
\begin{proof}
閉軌道自身の候補は系\ref{cor:all-invertible-split-atlas}で尽くされる。
二固有値・単一固有値の非半単純直線は既出の量子または例外片側反転の閉軌道へ流入し、
$s=1$ の非半単純な射影合同の場合も二つの通常の行列のJordan成分を共役で消せば量子点になる。
重複曲面の非半単純型は反転あり半単純型へ流入する。従って新しい閉軌道を生じない。
量子基点には定理\ref{thm:quantum-universal-chambers}の全三表と安定群比較を用い、
一移動・二移動したテンソルに同じ有限比較を適用する。片側反転は
定理\ref{thm:mixed-flip-slice}, \ref{thm:mixed-flip-fiber}、
反転なし直線は上の定理、単純・重複曲面は定理\ref{thm:split-surface-all-fibers},
\ref{thm:reflected-collision-fibers}、六次式の孤立型は定理\ref{thm:isolated-split-component}である。
閉極限の一意性で異なる基点間の比較も終わる。特殊化閉性は定理\ref{thm:split-regular-limits}。
点対応が直線ファイバーを持つ条件は、点と隣の直線からなる射影パラメータ上で
全関係の評価消滅を課して射影した閉条件である。正則な閉極限では関係平面も階数6で特殊化する。
分解窓には定理\ref{thm:split-point-pencils}で左右ともそのようなファイバーがないため、
流入窓にも分岐はない。
\end{proof}

% !TeX root = ../papers/nc-p3.tex
\section{印付き四平面対の交換子階数と階数一の差}
\label{sec:paired-marked-four-planes}
定理\ref{thm:marked-point-commutators}の印で $g=q=4$ とする。
$D=\langle G,F\rangle$, $D'=\langle Q,P\rangle$ とし、$[D,D]$ は交換子の線型包を表す。
\begin{theorem}\label{thm:paired-marked-four-planes}
各三列が独立なら $r=\dim[D,D]=\dim[D',D']$ は4か5である。
両側で同時に基底を選び
\[
 F_x=G_x+\alpha(x)T,\qquad P_y=Q_y+\beta(y)L
\]
とできる。ここで $\alpha,\beta$ は非零一次形式で、$T,L$ は各三列の包の外にある。
このとき $[G_{\alpha^t},T]=[Q_{\beta^t},L]=0$ である。
$r=5$ では各四平面に可換二平面がただ一つある。
$r=4$ では各四平面の交換子核は二次元で、表示
$\ker(\wedge^2D\to[D,D])=h\wedge E$ を持つ。ここで $E$ は三次元で
$h\in E$ はスカラーを加えれば階数1であり、可換二平面は $k h$ を共有する射影直線をなす。
\end{theorem}
\begin{proof}
$G,F$ の二つの像超平面の交わりを
$G_{a(t)}=F_{c(t)}$ と書く。$a,c$ は二次元の $t$ に関する単射である。
$G_a\wedge F_c$ は $\wedge^2D$ を張る（交叉の二基底と各補元の対を取る）。
従って両側の交換子次元は交換子テンソルの行列階数 $r$ に一致する。
印付き等式から
\[
 \mathcal S=\langle c(t)a(t)^t\rangle\subset [D',D']^\perp
\]
が跡対合に関し成立する。$\mathcal S$ は二次元から二次元への対称二行列全体を
単射と全射で挟んだ三次元空間で、全体で共通核を持つ。単位行列も直交するので $r\le5$。
逆側にも同じ構成がある。

どちらの四平面も非退化対称形式について自己随伴にはならない。
例えば $D'$ を対称にできれば印付き等式より全ての $x$ で $[G_x,F_x]=0$ である。
$D$ の座標を交叉二元と補元二つで選び
$G_x=(u_1,u_2,a,0)$, $F_x=(v_1,v_2,0,b)$ と書く（文字は $x$ の一次式）。
$x\mapsto G_x\wedge F_x$ の座標二次式は
$a\langle v_1,v_2,b\rangle+b\langle u_1,u_2,a\rangle$
を含む。各括弧は全一次式空間なので $a,b$ が独立ならその次元は5となる。
一方この二次写像の像の包は補題\ref{lem:irreducible-four-plane-brackets}により
次元3以下の交換子核に入る。よって $b$ は $a$ の倍、全座標二次式は $a$ で割り切れる。
$\ker a$ 上で $u_1v_2-u_2v_1=0$ となり、この二平面から交叉への二つの同型は比例する。
従って $F_x=dG_x+a(x)Z$、可換性により $Z$ は既約な $G$ と交換する跡零元なので零となり矛盾する。

次に、既約な跡零四平面で $r=3$ なら必ず自己随伴にできることを示す。
前掲補題の証明により、一般 $X\in D$ は三固有値相異なり、
$\dim[X,D]=3$ である（一般で2以下なら五次元の既約Jordan部分空間が生じる矛盾）。
$X$ を対角化すると $[D,D]=[X,D]$ は非対角である。
$Y$ の成分を $l_{ij}(Y)$ とすると、全ての対角交換子成分が零なので
$l_{12}\wedge l_{21}=l_{23}\wedge l_{32}=l_{31}\wedge l_{13}$。
この交代形式が非零ならその台二平面に全非対角成分が属して $[X,D]$ の階数に反する。
よって各対の二成分は比例し、三対の寄与は各一次元、互いに独立である。
三対を別々に支える元を選び、向き $12,23,31$ でその成分を $p_i$、逆成分を $q_i$ とする。
最初の二元の交換子の $13,31$ は $p_1p_2,-q_1q_2$。
$[X,D]$ の同成分対が $(q_3,-p_3)$ の倍なので $p_1p_2p_3=q_1q_2q_3$。
既約性により有向成分グラフは強連結である。一方向だけの辺があればそれを含む有向三角形があり、
この積等式で逆向き三角形も存在して矛盾する。従って全成分は非零で、
同じ積等式から $h_iY_{ij}=h_jY_{ji}$ を満たす非零対角計量が存在する。これで $r=3$ を排除した。

残る交換子核 $K\subset\wedge^2D$ の全元が可分解であることを示す。
$\phi(X,Y,Z)=\operatorname{tr}(X[Y,Z])$ と置く。
$\phi\ne0$ なら四次元の三形式には根基直線 $kh$ があり、その二重縮約の核は $h\wedge D$。
$K$ はこの核に含まれる。$\phi=0$ なら $kI+D\subset[D,D]^\perp$ より $r=4$ で等号となる。
そこで反対側からの三平面 $\mathcal S'=\{C S A:S=S^t\}$ が $kI+D$ に入る。
$C$ は $3\times2$ の単射、$A$ は $2\times3$ の全射とし、$J=AC$ は階数1以上である。
$\mathcal S'$ の三基底の三形式は符号を除き
$(J_{12}-J_{21})\det J$ となる（対称二行列の標準三基底で展開する）。
$J$ が非対称なら階数1であり、合同変換で $E_{12}$ にできる。
$C$ の二列を基底の初めに取り、第三列を選べば $\mathcal S'$ の基底は
$E_{12},E_{13}+E_{22},E_{23}$。その交換子は厳密上三角全体を張るので、
$\phi=0$ から $D$ は上三角に含まれてしまう。従って $J$ は対称である。
このとき $[CSA,CTA]=C(SJT-TJS)A$ の中項は交代二行列で、全交換子は一次元以下である。
$\mathcal S'$ を跡零化した三平面を $E$ とすると $\wedge^2 E\cap K$ は二次元以上、
$\dim K=2$ だから等号である。三次元の全二ベクトルは可分解であり、
二次元の束は共通直線 $kh\subset E$ を持つ。
$\phi\ne0$ で $r=4$ の場合にも $K\subset h\wedge D$ より同じ表示を得る。
いずれでも $h$ は $E$ 全体と交換するので三行の巡回元にはなりえず、
非零跡零であることからあるスカラー行列との差が階数1である。
$r=5$ のときは $K$ が一本の可分解直線である。

最後に $D$ への全射 $\pi:(a,c)\mapsto G_a+F_c$ を考える。
入力六空間に標準交代形式 $a^tc'-a'^tc$ を入れると、逆二ベクトルを商へ移したものは
符号を除き $\Omega=\sum_i G_i\wedge F_i$ である。交換子式の跡から $\Omega\in K$。
商の二ベクトルの核次元は $\dim(\ker\pi\cap(\ker\pi)^\perp)$ である
（商の双対は交代形式で $(\ker\pi)^\perp$ と同一視される）。
よって可分解性から二次元の $\ker\pi$ は等方的である。
$\ker\pi=\{(a(t),-c(t))\}$ なので $a(t)^tc(u)=a(u)^tc(t)$。
対応 $a(t)\mapsto c(t)$ は可逆対称行列へ延長できる。実際、指定二列は独立であり、
その主二行二列が可逆なら最後の対角値を一般に選べばよく、階数1なら最後の列との対合が
その根基で非零であるため既に行列式が非零である。合同変換でこの延長を単位行列にする。
すると $G,F$ は同じ二平面上で等しく、主張の差表示が従う。
反対側の同じ基底選択はこの一致を同時共役するだけなので両側で成立する。
交換子の跡縮約をもう一度用いると最後の二可換条件も従う。
\end{proof}

% !TeX root = ../papers/nc-p3.tex
\section{有限点枝の反対階数八と両対称性}
\label{sec:rank-eight-finite-symmetry}
\begin{theorem}\label{thm:rank-eight-finite-symmetry}
標準有限Frobenius窓の四辺が有限、反対階数が8で、長さ20の点スキームに長さ1の局所成分が一つあるとする。
中支持3の点があれば全次数で正則である。それ以外の窓は全て独立な基底変更で両対称反転窓になる。
従ってこの範囲では全次数の正則性と添字保持同型類は有限に決まる。両対称の場合は最小周期が1または2で、
純関係と大きな射影ファイバーは既に排除されているので、残る正則性の検査は二変数Hankel商の不在だけである。
\end{theorem}
\begin{proof}[二つの交換子空間]
系\ref{cor:point-commutator-rank}の印付き表示を使う。中支持3の場合は定理\ref{thm:finite-cubic-separators}による。
それ以外では $G,F$ は同じ既約な跡零三平面 $\mathcal G\subset\operatorname{End}(V)$ の二基底、
$Q,P$ は同様の $\mathcal Q\subset\operatorname{End}(U)$ の二基底である。
既約とは共通不変部分空間が非自明にはないという意味である。$F_c=G_{Tc}$、$P_d=Q_{Jd}$ と書き、
交換子式を
\[
 [G_a,G_{Tc}]_{bd}=-[Q_b,Q_{Jd}]_{ac}
\]
とする。$T,J$ は可逆である。この式の $ac\mid bd$ 階数 $r$ は両三平面の交換子の線型包の次元に等しい。
$r\le1$ なら $\wedge^2\mathcal G$ の核の二次元部分が共通因子を持ち、
その非零因子が $\mathcal G$ 全体と交換する。Schurの補題と跡零性に反する。
従って $r=2,3$ である。さらに各交換子空間は、三行三列の交代行列空間の片側を可逆変換したものに入る。
実際、右辺は $b,Jd$ について交代し、左辺にも同じ議論が適用できる。
\end{proof}

\begin{proof}[交換子の次元が二の場合]
$\mathcal G$ 内の二平面の基底 $X,Y$ が可換であり、第三の基底を $Z$ とする。
その可換平面に最小多項式が三次の $X$ を取れる。実際、二固有値の半単純元しか三次に届かない
可能性を考えると、その二次元固有空間上で別の可換元がスカラーでなければ、一般線型結合の
最小多項式は三次になる。スカラーなら跡零性から二元が従属となる。
よって三次の元がなければ全元が平方零となる（単一固有値は跡零なので0）。
この二平面は全て階数1以下で、共通像直線または共通核平面を持つ。
前者なら $[X,Z],[Y,Z]$ の像はその直線と $Z$ での像の和に入り、後者なら共通核とその $Z$ 逆像の交わりを消す。
これは独立な二つの交代三行列が全体で三次元を張る像と自明な共通核を持つことに反する。
この二性質は上の片側可逆変換でも保たれる。

巡回行列 $X$ に可換な $Y$ はその多項式なので $(X,Z)$ は既約である。
$\det[X,Z]=0$ は、ある非零対称行列 $H$ が $HX=X^tH,\ HZ=Z^tH$ を満たすことと同値である。
実際、対称行列六次元から交代行列二組の六次元へのこの線型写像の行列式は符号を除き $\det[X,Z]$ である。
一般の $X$ を対角化すると、三つの非対角未知数の係数は固有値差で、対角未知数の残りの
行列式は $Z_{12}Z_{23}Z_{31}-Z_{13}Z_{32}Z_{21}$ の符号倍である。
同時共役で行列式の変換倍率は1となる（対称二乗の行列式の指標は四乗、交代二乗二組も四乗）。
多項式恒等式なので全域で成り立つ。ここでは交換子空間が交代行列の片側変換に入るため行列式は零である。
$H$ の核は不変なので $H$ は可逆で、$Y$ も同じ形式に関して自己随伴である。
$\mathcal Q$ にも適用し、二空間の基底を独立に変えれば四組 $G,F,Q,P$ はいずれも対称行列になる。

このとき交換子式の右辺は行列として $a,c$ の交換で符号を変える。従って全ての $a$ で $[G_a,G_{Ta}]=0$。
可換二平面の外の $G_a$ と交換する $\mathcal G$ の元は $G_a$ の倍しかない
（交換子写像の核はこの平面の外積の直線である）。従って $T$ はスカラーで、$J$ も同様である。
\end{proof}

\begin{proof}[交換子の次元が三の場合]
交換子式を $a,c$ で跡を取れば $\sum_a[G_a,G_{Ta}]=0$。三つの交換子の独立性により $T$ は対称で、
$J$ も対称である。可逆対称形式をそれぞれ合同変換で単位行列に取り、$F=G,P=Q$ とする。
このとき $[G_a,G_c]$ と $[Q_b,Q_d]$ は交代行列である。

三基底 $G_i=S_i+K_i$ を跡零対称部分と交代部分に分けると $[K_i,S_j]=[K_j,S_i]$ である。
$K_i=[v_i]_\times$ は三次元外積作用とする。$v_i$ の張る次元が1なら二つの独立な対称部分が
同じ非零外積行列と交換することになる。しかしその行列は巡回行列で（等方ベクトルでも平方が非零の冪零行列）、
跡零対称な中心化空間は一次元である。次元2なら第三の非零対称部分が独立な二つの外積行列と交換する。
その平方 $vv^t-(v^tv)I$ とも交換するため、二平面内の各 $v$ はその固有ベクトルである。
従ってこの対称部分はスカラーを除くと $cww^t$、$w$ は平面に垂直となる。
$c\ne0$ なら交換関係から全ての $v$ で $v\times w=0$ となり矛盾する。
次元3なら三基底を線型結合して $v_i=e_i$ とする。上の三つの交換子式は $S_i$ の十五成分に非退化である。
具体的に各対称跡零行列に基底
$E_{11}-E_{33},E_{22}-E_{33},E_{12}+E_{21},E_{13}+E_{31},E_{23}+E_{32}$ を取り、
対 $23,31,12$ の順に差の成分 $11,22,12,13,23$ を読むと行列式は $-3456$ である。
従って $\mathcal G$ は全て対称か全て交代であり、$\mathcal Q$ も同様である。

混合型は反対階数条件に反する。$G$ を交代側として $G_i=[Le_i]_\times$ と書き、
定理\ref{thm:marked-opposite-eight}と同じ向きで交換子行列 $D_G$ と
$P_G(Z)_{ac}=\operatorname{tr}(ZG_aG_c)$ を取る。対称跡零 $Z$ に対して
\[
 D_G P_G(Z)D_G^t=-\det(D_G) Z
\]
である。$D_G=-\operatorname{cof}(L)$ と $P_G(Z)=L^t ZL$ は外積の積の式による。
一方、対称な $Q$ では補題\ref{lem:ternary-net-commutator}により、
$I,Q_1,Q_2,Q_3$ に直交する対称行列 $Y$ に対して同じ式の符号が正である。
$D_Q=-D_G^t$ なので残余対合は
\[
 \operatorname{tr}(Y P_G(Z))+\operatorname{tr}(Z P_Q(Y))
 =-2\det(D_G)\operatorname{tr}(Z D_G^{-t}Y D_G^{-1}).
\]
$Y$ の空間は二次元以上であり、その可逆合同変換の像が全てスカラーになることはないから、
右辺は恒等零ではない。ところが系\ref{cor:point-commutator-rank}の $\bar w=0$ はまさにこれを消す。
逆の混合型も位置を交換すればよい。

双方が交代なら $Q_b=[Ne_b]_\times$ と書いて
$\operatorname{cof}(N)=-\operatorname{cof}(L)^t$ より $N$ は $L^t$ の非零倍となる。
最初の辺に点対 $((r,x),(s,y))$ を入れた六式は
$r y+G_x^t y=0,\ s x+Q_y^t x=0$ である。実際、これらは反対側の星印一つを含む六成分であり、
印付き表示の可逆な六次小行列により独立である。$x,y$ をそれぞれ非零等方ベクトルとし、
$y^tLx=0$ を課す。両等方二次曲線の積上のこの零点は次元1以上である。
$(Lx)\times y$ は $y$ の倍であり、$(Ny)\times x$ も $x$ の倍なので
$r,s$ が取れて辺の実際の点対となる。これは有限性に反する。
\end{proof}

\begin{proof}[両対称表示]
残る全ての場合に $G,Q$ が対称で $F=hG,P=kQ$、$h,k\ne0$ である。
位置0,1の星印の係数をそれぞれ $1/h,1/k$ 倍すれば、星印を持つ項は各対向位置の交換で不変となる。
純成分 $h(G_aG_c)_{bd}+k(Q_bQ_d)_{ac}$ も交換子式で各交換に不変である。
よって四テンソルは両対称であり、定理\ref{thm:bisymmetric-two-length-test}により結論を得る。
同型・周期は定理\ref{thm:finite-all-window-keys}で決まる。
\path{computations/verify_rank_eight_finite_symmetry.py}は十五成分の行列式、六未知数の対称化方程式、
および二種の符号のうち交代側を検算する。
\end{proof}

% !TeX root = ../papers/nc-p3.tex
\section{非被約な印から反対位置の射影同一視を延ばす}
\label{sec:finite-marked-projection-symmetry}
四辺が有限な標準有限Frobenius候補で全ての中支持が4とする。
共通の長さ20の点スキームを $X$、その各座標を $x_i$ とし、一つの閉点の座標を $p_i$ とする。
三行三列の行列を、位置1と3で $p_1,p_3$ を消す双線型形式として読む。
\begin{theorem}\label{thm:finite-marked-projection-symmetry}
二つの四次元空間 $w(p_0,-,V_2^*,-)$, $w(V_0^*,-,p_2,-)$ が一致することを仮定し、
共通空間を $E$ とする。反対階数8ならこの一致はどの印でも自動である。
次の二条件を課す。
\begin{enumerate}
\item 印以外の各閉点で、$x_1,x_3$ の $p_1,p_3$ からの射影を $l,r$ とする。
$E$ の中で $l^tZ=Zr=0$ を満たす空間は一次元である。
\item 印の接空間次元は高々2で、各非零接方向を二つの射影座標の微分で
$l,r$ に写したときにも同じ一次元条件が成り立つ。
ここで微分は $p_i$ のアフィン近傍でその点を原点とした三座標について取る。
\end{enumerate}
このとき唯一の線型同定 $T:V_0^*\to V_2^*$,
$w(x,-,p_2,-)=w(p_0,-,Tx,-)$ は、$X$ 上で反対位置への二移動を射影的に与える。
窓はこの位置について対称になる。他方でも同じ条件を満たす印があれば両対称となり、
定理\ref{thm:bisymmetric-two-length-test}によって全次数の必要十分判定が得られる。
二つの印は異なってよい。

特に $E$ の行列式三次曲面が滑らかなら二つの一次元条件が成り立つ。
より一般に、共通の左右核を持つ行列の二次元鉛筆が $E$ 内にないことでも十分である。
\end{theorem}
\begin{proof}
二つの部分切断はどちらも四文字空間から $E$ への同型である。
独立な二つの側の係数を反対切断の行と列に分けて考える。
偶位置の一つが印の行は、奇位置の一つが印の列を全て消し、その逆も成り立つ。
行側の像と列側の像は各々次元4以上で、その和は反対階数以下である
（独立行を選んだ後、それらを消す列を選ぶ）。階数8なら両側で次元4なので一致を得る。

局所的に各点座標の単元倍を選ぶ。$A=w(p_0,-,x_2,-)$ と $B=w(x_0,-,p_2,-)$ は
いずれも $E$ に値を取り、閉点では零にならない。
両方とも、左右から同じ $x_1,x_3$ で縮約すると零である。これは四つの隣接点対の関係そのもので、
$X$ 上の恒等式である。$E$ の全元は $p_1,p_3$ をも消すため、これらの縮約には三座標 $l,r$ を用いられる。
印以外では六本の線型条件の階数は3であり、$A$ がその核の直和因子を張る。
従って $B$ は $A$ の単元倍である。これは局所環が非被約でも成り立つ。

印の局所環を $S$、極大イデアルを $\mathfrak m$ とする。$A,B$ の剰余は単元倍を除いて
$M=w(p_0,-,p_2,-)\ne0$ である。$\lambda(M)=1$ となる $E$ の汎関数を取り、
$Z=B-\lambda(B)A/\lambda(A)$ と置く。$Z$ は三次元の $N=\ker\lambda$ に値を持ち、
三座標 $l,r\in\mathfrak m^{\oplus3}$ による六つの縮約を全て消す。
接方向上の一次の $l,r$ は $M$ を消す（$l^tA=Ar=0$ を一次で読む）。
仮定から $N$ 上の $6\times3$ 縮約行列の一次項は全非零接方向で階数3である。

この条件が $\mathfrak m^3Z=0$ を与えることを示す。
接次元2なら、一次項の転置により接方向の $\mathbb P^1$ 上に全射
$\mathcal O^6\to\mathcal O(1)^3$ がある。核は三つの
$\mathcal O(-a_i)$ の直和で、$a_i\ge0,\ \sum a_i=3$。
従って2だけひねれば核の一次コホモロジーが消え、行の二次式倍は
三個の任意の二変数三次式を生成する。接次元1では非零方向での定数係数の階数3から同じ主張を得る。
接次元0なら自明である。局所環の次数環は対称代数の商なので、
実際の行の $\mathfrak m^2$ 倍は $\mathfrak m^3 S^{\oplus3}$ を $\mathfrak m^4$ を法として張る。
さらに $\mathfrak m$ の冪を掛けて冪零性を使えば $\mathfrak m^3 S^{\oplus3}$ の全体が行加群に入る。
これらの行は全て $Z$ を消すので主張が従う。

最後にこの三乗の誤差を取り除く。$D=H^0(\mathcal O_X)$ とし、
$J$ を印で $\mathfrak m^3$、他の因子で単位イデアルとする。
隣接二位置の座標積 $C=U_0U_1\subset D$ は定理\ref{thm:finite-square-chart}の完全跡対合で
自分自身と直交する。実際 $C^2=U_0^2U_1^2$ をこの跡が消す。
さらに $C\to D/J$ は全射である。各位置で単元座標を1に自明化すると、両位置の
一次座標差は閉埋込みのためそれぞれ $\mathfrak m/\mathfrak m^2$ を張り、
両者の積は $\mathfrak m^2/\mathfrak m^3$ を張る。従って
$C\cap\operatorname{ann}(J)=0$ である。実際この共通部分の元は
$J$ と $C$ の両方に跡直交し、$C+J=D$ である。
第二辺の任意の関係を $(x_1,T x_0)$ で評価した値は $C$ に入り、
印以外で零、印では $\mathfrak m^3$ に消される。後者は $T x_0$ と $x_2$ が
単元倍を除いて $Z$ の線型成分だけ異なることによる。
よってこの評価値は全て零である。第二辺の一意なグラフにより $T$ は $X$ 上で二移動を与える。
六本の関係は有限完全交叉の $(1,1)$ 消去式を全て張るので最初の二辺は反転となり、
命題\ref{prop:finite-flip-obstructions}の交代三次交叉の消失からテンソルの対称性が従う。

二つの縮約を消す元は各閉点では $A$ により、接方向では $M$ により既に一つある。
二元あれば、それらが張る射影直線では基底を取って全行列の最終行と最終列が零になる。
行列式が滑らかな曲面なら階数1以下の点は存在しないが、この直線の二行二列の行列式は
二変数二次式であり零点を持つ。従って滑らかな場合にもこのような直線はない。
\end{proof}

% !TeX root = ../papers/nc-p3.tex
\section{階数一切断の長さと非被約な例外鉛筆}
\label{sec:finite-marked-secant-pencil}
前節の共通行列四平面を $E$ とし、$\Sigma=\mathbb P(E)\cap(\mathbb P^2\times\mathbb P^2)$ を
階数1のスキーム切断とする。接線を数えるためスキームの長さを使う。
\begin{theorem}\label{thm:finite-marked-secant-pencil}
$\Sigma$ が有限なら長さは4以下である。長さ2以下のとき、
印の接次元が0または1なら、他の局所成分への仮定なしに前節の $T$ は反対対称性を与える。
印の接次元が2でも、その非零接方向に前節の一次元条件が成立すれば同じである。
従って両側でこの判定を通れば全次数AS判定ができる。
\end{theorem}
\begin{lemma}\label{lem:finite-adjacent-word-support}
長さ20の有限辺で $C=U_0U_1\subset D$ とし、$0\ne f\in C$ の支持を
$Z=\operatorname{Spec}(D/\operatorname{ann} f)$ とする。
$Z$ 上での二つの座標系の線型包の次元を $a,b$（射影次元に1を加えた値）とすると
\[
 a+b\le d=\operatorname{length} Z\le\binom{a+b-2}{a-1}.
\]
特に $a,b\ge3$, $d\ge6$ で、$d=6$ なら $a=b=3$、$Z$ は二平面の積での
四つの双一次式の完全交叉である。第二位置のどの直線 $L$ に対しても
$d-\operatorname{length}(Z\cap L)\ge b$ が成り立つ。ここで交叉は第二射影で取る。
\end{lemma}
\begin{proof}
$\tau(f\,\cdot)$ は商の完全跡であり、$C$ が $C$ と跡直交するので
二座標空間の制限像はこの対合で互いに直交する。
両射影の包 $\mathbb P^{a-1}\times\mathbb P^{b-1}$ に元の六式を制限すれば
基点スキームは有限で $Z$ を含む。従ってそのうち $a+b-2$ 本の一般線型結合が
固有に交わる（有限基点の外で各段の正次元成分を避ける）。
交点数は表示の二項係数である。$a$ または $b$ が1,2なら二つの不等式は矛盾する。
$d=6$ では全ての不等式が等号となり、完全交叉に含まれる同じ長さの部分は全体である。
$Z\cap L$ は第一射影の包と $L$ の積で同様に長さ $a$ 以下なので最後の主張を得る。
\end{proof}
\begin{proof}[階数一切断と固定鉛筆]
$\mathbb P(E)$ を定める五超平面のうち一般の四線型結合でSegre四次元多様体を切ると、
有限基点仮定により長さ6の完全交叉 $Y$ になる。これは切断の射影四平面を張る。
実際 $(1,1)$ で消える切断は四つの定義式で尽きる（Koszul分解で $j=2,3$ の
$\mathcal O(1-j,1-j)$ は全コホモロジーが零、$j=4$ の非零コホモロジーは次数4だけである）。
$Y$ での $(1,1)$ 切断の像は五次元で、随伴 $\omega_Y=\mathcal O_Y(1,1)$ の完全跡の
核である。これは同じKoszul分解の余核への写像が双対で単位の制限になることによる。
残った超平面の切断は非零で、生成するイデアルの長さが1なら、ある局所因子のsocleを張る。
完全跡はこの元で非零なので矛盾する。よって残余長さは2以上、$\Sigma$ の長さは4以下である。

左右の指定核 $l,r$ を消す二次元鉛筆は二行二列の行列表示を持つ。その射影直線を
階数1で切る二次式は恒等零ではない（$\Sigma$ は有限）。従って長さ2の部分を切り出す。
$\Sigma$ の長さが2以下なら、そのような直線がある場合は長さが2で、直線は
$\Lambda=\langle\Sigma\rangle$ に一意である。二行二列の一般元は階数2なので
共通左右核も一意である。この一意性は局所Artin基底上でも成立する。
実際、二つの元が剰余体で独立で、共通左右核が直和因子の直線なら、
二次式による切断は相対射影直線で有限平坦な次数2の因子となる。
局所的には非根を無限遠に選び、単元を先頭係数とする二次式の商で確認できる。
この因子は $\Sigma$ の基底変換に含まれ、長さ2の自由加群間の全射は同型である。
従って相対鉛筆も相対核も一定である。$\Sigma$ が非被約長さ2でも、その一次式の核が
張る射影直線は一意なので同じ証明でよい。

以下、例外鉛筆があるとし、その左核が印からの射影で指定する位置1の直線を $L$ とする。
印以外の点で前節の $A,B$ の比例からのずれを調べる。
縮約核が二次元以上ならちょうど二次元である。非零の左縮約は階数2以上であることが
有限辺での純関係の不在から従う（定理\ref{thm:finite-cubic-separators}の証明）。
$\lambda(A)$ を単元として三次元の $\ker\lambda$ にある差 $Z$ を取ると、
その六本の縮約条件の階数は剰余体上で2である。二つの主元で消去した後は
$Z=f_0 V$、$V$ の一成分は1、$V$ と $A$ は剰余体で独立となる。
$f_0\ne0$ なら $\operatorname{ann}(f_0)$ を法として $A,V$ は実際に二本の共通核を消す。
相対鉛筆の一意性よりこの商で $x_1$ は $L$ にある。ゆえに $I_L Z=0$。
縮約核が一次元のときは最初から $Z=0$ である。

印が曲線状なら局所環を $k[t]/(t^m)$ とする。印では $A,B$ は比例するので
$Z\ne0$ なら $Z=t^s V$, $s\ge1$ と書け、$V$ は $A$ と剰余体で独立である。
射影三座標 $l,r$ はいずれも $t$ で割れ、それぞれの商は一成分が単元である。
$m-s\ge2$ なら $t^{m-s-1}$ を法として、$A,V$ は $l/t,r/t$ を共通核に持つ。
よって $l/t$ はその商で固定方向に入り、$l$ は $t^{m-s}$ を法としてその方向に入る。
$m-s=1$ のときも $l$ の剰余は零である。結局 $I_L Z=0$ が印でも成立する。
印が被約なら差は零である。従って各第二辺関係を $(x_1,T x_0)$ で評価した $f\in C$ は
全て $I_L f=0$ となる。非零なら補題の $b\ge3$ に反する。

印が接次元2で全接方向の条件が成立する場合、前節の行加群の証明から印で
$\mathfrak m^3 f=0$ である。$f\ne0$ と仮定し、補題の支持 $Z$ を考える。
印以外は $L$ 上にあり、印の局所成分があればGorensteinで極大イデアルの三乗が零である。
従ってその長さは4以下である。実際二乗が非零ならsocleそのもので次元1である。
$Z$ の印の接次元が1以下ならここで $L$ を課して失う長さは高々2であり、
補題の $d-\operatorname{length}(Z\cap L)\ge b\ge3$ に反する。
接次元2なら長さ4で、その接平面は元の $X$ の接平面である。
$L$ がこの平面に含まれるなら $I_L$ の一次部分の像は高々一次元で、
二乗部分も一次元なので同じ矛盾を得る。
$L$ がこの平面に含まれなければ、失う長さは高々3で、補題から $b=3$。
従って第二射影の包はまさに印の接平面である。
$d\ge6$ なので印以外にも支持点があり、それと印を結ぶ $L$ はその平面に入らねばならず矛盾する。
以上で評価は全て零となり、前節と同じグラフの議論で結論を得る。
\end{proof}

% !TeX root = ../papers/nc-p3.tex
\section{例外割線に接する平面局所点と三つの被約な割線}
\label{sec:finite-marked-tangent-secants}
\begin{theorem}\label{thm:finite-marked-tangent-secants}
定理\ref{thm:finite-marked-secant-pencil}で $\operatorname{length}\Sigma\le2$ なら、
印の接次元2でも全接方向の追加条件を除ける。
また $\Sigma$ が相異なる三点なら、接次元0または1の印で全く同じ反対対称性が成り立つ。
印以外の点はどちらの場合も非被約でよい。
\end{theorem}
\begin{proof}[平面局所環での二次の行関係]
長さ2で例外方向がある場合を扱う。印の局所環 $(S,\mathfrak m)$ の二つの生成元を $s,t$ とし、
例外接方向は $s=0$ とする。これは左右の各接写像が単射、例外割線の左右核が一意なので
ただ一つである。三座標を $l=s l_1+t l_0$, $r=s r_1+t r_0$ と書き、
係数ベクトル $l_i,r_i$ は $S$ 上に取る。前節の $A,B,Z=B-\lambda(B)A/\lambda(A)$ と
$N=\ker\lambda$ を使う。独立な形式変数 $v,w$ に関し、$N$ の行列を
$v l_1+w l_0$, $v r_1+w r_0$ で左右から縮約する六行三列を $H(v,w)$ とする。
実際の局所環では $H(s,t)Z=0$ である。

剰余体の射影直線上で $H$ の転置は $\mathcal O^6\to\mathcal O(1)^3$ を与える。
余核は $[0:1]$ にある長さ1の層であることを確かめる。
その点では階数2なので、局所的に余核は最大小行列式のイデアルによる環の商である。
そこで商が長さ2以上なら、その一次近傍の双対数環で $H$ の右核に、剰余が非零の $V\in N$ が取れる
（二つの主元で消去する）。$M=A\bmod\mathfrak m$ は全接方向の二つの縮約を消す。
従って $M,V$ が共通左右核を持つ相対鉛筆になり、前節の一意性からその核方向は一定である。
しかし $\bar l_0,\bar l_1$ は独立なので $v/w$ の一次変化で左方向は実際に動く。
これで長さ1が従う。核束は次数 $-2$、階数3で $\mathcal O^6$ の部分束なので、
分裂の各次数は $0,-1,-2$ のいずれかである。従って1だけひねった核の一次コホモロジーは消え、
像層の一次コホモロジーも $\mathcal O(1)^6$ の商であることから消える。
六行の一次式倍による三成分の二次式の余核は次元1で、例外方向の値で検出される。

$N^*$ の定数元 $e$ をこの方向の余核で非零に取り、上記二次式の余核を今度は $S$ 係数で作る。
剰余体上の結果と中山の補題で $w^2e$ がこの有限 $S$ 加群を生成する。
したがって、ある $a\in\mathfrak m$ について $vwe=a w^2 e$ が行を法として成り立つ
（例外方向で $v=0$）。各 $n\in N^*$ の $w^2 n$ も $c_n w^2 e$ に等しくなる。
基底に $e$ を含めて $c_e=1$ とできる。$(v,w)=(s,t)$ として $Z$ に作用させ、
$g=t^2 e(Z)$ と置くと
\[
 \mathfrak m^2\langle n(Z):n\in N^*\rangle=(g),\qquad
 t^2 Z=g V,\qquad (s-at)g=0.
\]
$V\in N\otimes S$ は基底に対する $c_n$ を成分に持ち、$e(V)=1$ である。
最後の式は $st e(Z)=a t^2e(Z)$ にさらに $t$ を掛けたものである。

$g$ 上では $l,r$ を $t(a l_1+l_0),t(a r_1+r_0)$ に置き換えられる。
$tg\ne0$ なら $\operatorname{ann}(tg)$ を法として、$A,V$ はこの二つの括弧に入る方向を
共通核に持つ。実際 $A$ の縮約は元から零、$V$ の縮約は $gV=t^2 Z$ から $g$ を掛けて零であり、
$t$ を取り出せばよい。両方向は剰余で非零、$A,V$ は剰余で独立である。
よって剰余長さにかかわらず例外割線の固定左核が得られ、$I_L g=0$ となる。
$tg=0$ なら $sg=0$ でもあり同じ結論となる。
前節の第二辺の評価 $f\in C$ は差 $Z$ の成分の局所線型結合なので、印では
$\mathfrak m^2 I_L f=0$ である。

印以外では既に $I_L f=0$。従って補題\ref{lem:finite-adjacent-word-support}の支持に
$L$ を課して失う長さは、印における $I_L f$ の長さに等しく、3以下である。
実際 $S$ は局所Gorensteinで $\operatorname{ann}(\mathfrak m^2)$ の長さは
$\operatorname{length}(S/\mathfrak m^2)=3$ である。
$f\ne0$ なら同補題により $b=3$ で、失う長さもちょうど3となる。
支持の長さは6以上ゆえ $Z(f)=\operatorname{Spec}(D/\operatorname{ann} f)$ の $L$ との交叉は
長さ3以上である。その第二射影は閉埋込みだから $L$ を張り、第二射影の包である平面は $L$ を含む。
この平面で $I_L$ は一本の線型式で生成され、特に印の $I_L f$ は単項である。
局所Gorenstein環の単項イデアルは零でなければその生成元の零化イデアルによる
Gorenstein商と加群同型である。ここではその極大イデアルの二乗が零なので、
socle次元1より長さは2以下である。これは必要な長さ3と矛盾する。
\end{proof}

\begin{proof}[三つの被約な階数一点]
三つの階数一行列の像ベクトルは独立である。
さもないと三行列はある共通左核を持ち、左縮約が階数1以下になる。
ここで三行列は独立である。三つが射影直線に入れば全二次小行列式がその直線で恒等零となり
$\Sigma$ の有限性に反する。右側でも同じ独立性が成り立つ。
従って三割線の左核は双対基底の三方向を与え、位置1では印を通る非共面な三直線 $L_1,L_2,L_3$ を得る。
局所Artin基底の相対鉛筆が階数一で切る次数2の有限平坦族は、三つの定数点の和のうち
二点の和の基底変換に限る。実際その三つの直和因子は各々平坦な環の直和因子で、
剰余で長さ1なら基底環と等しい。従って前節と同じ主元の消去、印での一様化変数の割算により、
評価 $f$ は各局所因子で三直線のどれか一つのイデアルで消される。点ごとに選ばれる直線は異なってよい。

二直線を含む平面の式 $h\in U_1$ を取ると、$hf$ の支持は残る直線に入る。
この支持に制限した第一座標の次元を $a'$ とすると、その長さは $a'$ 以下である
（補題の証明での $\mathbb P^{a'-1}\times\mathbb P^1$ 上の有限切断）。よって第一座標だけで
支持環へ全射する。一方 $\tau(hf U_0)=0$ は $f\in C$ の自己直交性から従う。
$hf$ が非零ならその支持商での完全跡に反する。三つの二直線の組の全てで $hf=0$ となる。
三平面の式は独立で共通零点は印だけなので、$f$ の第二射影の支持はスキームとして一つの被約点に含まれる。
補題の $b\ge3$ に反するため $f=0$ である。
\end{proof}

% !TeX root = ../papers/nc-p3.tex
\section{正次元の階数一切断と核方向の剛性}
\label{sec:finite-marked-ruling-pencil}
\begin{theorem}\label{thm:finite-marked-ruling-pencil}
定理\ref{thm:finite-marked-projection-symmetry}の二つの四次元切断が一致するとする。
$\Sigma$ が正次元なら、その全体がSegreの一方の族の被約直線であり、
どの印でも $T$ はスキーム上の反対位置を同定する。従ってその側のテンソル対称性が従う。
\end{theorem}
\begin{proof}
$E$ の非零の片側ベクトルでの評価はどちらからも階数2以上である。
よって階数一点の整曲線の一方の射影が直線内にあれば、その全体は
一つの非零片側縮約の核に入り、線型包のベクトル次元は2以下である。
両射影とも平面を張るなら、その関数体の中の三次元座標空間二つの積は次元5以上である。
実際、滑らかな一点の位数で各空間に三つの異なる位数の基底を取り、
二つの三元整数集合の和集合（対ごとの和全体）に五つ以上の異なる値があることを使う。
これはSegre像の線型包が $\mathbb P^3$ に入ることに反する。
従って直線があり、左右を必要に応じて交換して、固定像を持つ
$K=\langle u v_1^t,u v_2^t\rangle$ と書く。

補元 $C,D$ を取り、上線を像直線 $ku$ を法とした二行三列への射影とする。
全ての $0\ne l\in u^\perp$ で $l^t C,l^t D$ は独立である。これは任意の
非零な $s\bar C+t\bar D$ の階数が2であることとも同値である。
従って $K$ の外には階数一点がない。$uv^t\in\mathbb P(K)$ においてSegreに接する
摂動は $ku$ を法として像の階数が高々1なので、$E$ 内の全接方向も $K$ に入る。
各閉点で局所次元と接次元が1なので交叉スキームはこの被約直線である。
さらに $\bar C,\bar D$ の共通右核は零である。さもなくば三列を二列に下げた
二行二列鉛筆が全非零元で可逆になり、その行列式の零点に反する。

二次元鉛筆が左右核を持ち、一般元が階数2なら、制限した二行二列の行列式が
$\Sigma$ の長さ2の部分を切る。その線型包は $\mathbb P(K)$ そのものなので、
実際には鉛筆の全元が階数1となり矛盾する。従って鉛筆は $K$ のみである。
共通右核は $kr_0=\ker v_1^t\cap\ker v_2^t$、左方向は任意の $[l]\in\mathbb P(u^\perp)$ である。
この記述が相対的にも成り立つことを示す。二平面と左右の二つの核直線をパラメータとする
Grassmann多様体と射影空間の積で縮約の消失スキームを取る。
$(K,l,r_0)$ で二平面の一次変形を、基底 $uv_i^t$ に $s_i C+t_i D$ を加えるものとする。
左縮約の一次式は
\[
 \alpha v_i^t+s_i l^t C+t_i l^t D=0,\qquad \alpha=\dot l^t u .
\]
$\alpha\ne0$ なら二行 $(s_i,t_i)$ は独立である。右縮約の一次式を $ku$ を法として読むと
$(s_i\bar C+t_i\bar D)r_0=0$ なので $r_0$ が共通右核となって矛盾する。
従って $\alpha=0$、$s_i=t_i=0$、$\dot r$ は射影的に零である。
このパラメータスキームは各閉点で接次元1、台は上の射影直線なので正則であり、
その被約射影直線と等しい。よって任意のArtin基底上でも右方向と鉛筆は固定である。

定理\ref{thm:finite-marked-secant-pencil}の局所消去を右方向で用いる。
印以外の誤差は、位置3で印と $r_0$ を結ぶ固定直線 $L$ のイデアルに消される。
印の接次元が1以下なら同定理の一様化変数の議論もそのまま適用できる。
評価には $(T x_0,x_3)$ における辺 $23$ の関係を用いる。これは $U_0U_3$ に属し、
辺 $30$ の自己直交性と支持補題が位置3の直線について適用できる。
接次元2で例外接方向がなければ前節までの三乗の証明による。
例外方向がある場合、、位置3の微分が単射なので $r_0$ へ行く射影接方向はただ一つである。
定理\ref{thm:finite-marked-tangent-secants}の二変数行行列の余核はここでも単純な一点である。
実際そこの双対数上の二独立元があれば上の相対剛性が右方向を一定にするが、
射影接直線上で右方向の微分は非零である。従って同定理の二次の余核による
$\mathfrak m^2 I_L f=0$ の証明は全て右方向で成り立つ。
支持長の矛盾により評価は零となる。接次元3は生じない。
実際、辺 $01$ の二つの端の評価像 $B,C$ は六次元空間の三平面で、
接次元は $\dim(B\cap C)$ である。辺 $12$ の消失から $w(p_0,-,p_2,-)$ の
行列は $B$ と $C^\perp$ の対合（順序は端の取り方による）であり、その階数は
$3-\dim(B\cap C)$。中支持4では $p_2$ によるこの行列は非零である。
\end{proof}

% !TeX root = ../papers/nc-p3.tex
\section{有限な階数一切断のArtin環と三方向の分離}
\label{sec:finite-marked-artin-secants}
\begin{theorem}\label{thm:finite-marked-artin-secants}
共通切断 $E$ の $\Sigma$ が長さ3のGorensteinスキームなら、
接次元0または1の印でその側の反対対称性が成り立つ。三重の一点も含む。
長さ4なら $B=H^0(\mathcal O_\Sigma)$ は必ずGorensteinであり、
非退化な積跡 $\tau:B\to k$ を選ぶと、左右の三次元基底変更で $E$ は
\[
       z\longmapsto \big((a,b)\longmapsto\tau(zab)\big),
       \qquad z\in B,\quad a,b\in\ker\tau
\]
に一致する。特に $E$ のこの行列空間としての同型型は $B$ で決まる。
この部分切断の対称表示だけでは四窓の対称性はまだ従わない。
\end{theorem}
\begin{proof}[射影と積の表示]
$\Sigma$ のどの長さ3の部分も、どちらの射影でも一本の直線に入らない。
仮に入れば $E$ の非零片側縮約の核に全体が入り、この核は次元2以下である。
その射影直線の全二次小行列式は長さ3で消えるので直線全体で消え、有限性に反する。
直線束を自明化して二つの三座標の像を $U,V\subset B$ とする。
長さ3なら $U=V=B$ で、$\langle\Sigma\rangle$ のベクトル空間は
$B\otimes B\to B$ の積の双対像となる。
長さ4なら $U,V$ は超平面であり、どちらも非零イデアルを含まない。
実際イデアルは長さ1のイデアルを含み、対応する長さ3の商に座標空間の三次元性が必要である。
よって $U,V$ を零化する二つの汎関数はともに完全積跡であり、その比は単元倍である。
一方の自明化を変えて $U=V=\ker\tau$ とする。積 $UV$ は $B$ を張る。
さもなくば $0\ne z\in B$ について $zU\subset U^\perp=k1$ となる。
単元なら次元に反し、非単元なら $zB$ は真のイデアルなので $zU=0$。
このとき $U$ は $\operatorname{ann}(z)$ に含まれるが、$\operatorname{ann}(z)\ne0$ の場合
$U+\operatorname{ann}(z)=B$（$U$ は非零イデアルを含まない）なので矛盾する。
従って $\langle\Sigma\rangle=\mathbb P(E)$、積の双対で主張の表示を得る。
他の完全跡は単元を掛けた跡であり、二つの引数と $z$ に単元を配分すれば同型になる。
逆向きにも行列空間から階数1切断の環が復元される。

\end{proof}
\begin{proof}[長さ3の相対割線]
$B$ の完全跡を固定する。局所Artin基底 $S$ 上で $\Sigma$ の次数2の部分を指定することは、
$B_S=B\otimes S$ に直和因子の直線イデアル $I$ を指定することと同じである。
$B$ は $I$ に指標 $\chi:B\to S$ で作用する。$I$ は $\ker(\chi_S)$ の
完全跡による直交補、すなわちこの核の零化イデアルに一致する。
実際そこに含まれ、どちらも直和因子の直線である。逆に任意の指標がこの直線を与える。
相対割線に対応する二次元行列空間は $I$ の零化汎関数、すなわち完全跡で同定すれば
$\ker(\chi_S)$ の元による乗法である。その共通左核は $I$ である。
従って全ての相対左方向は、$B^*$ 内の指標を完全跡で $B$ に移した
一定の長さ3のスキーム $D_3\subset\mathbb P^2$ を通る。
これは $B$ の完全一次系での埋込みであり、その三座標は自明化後に $B$ そのものを張る。

定理\ref{thm:finite-marked-secant-pencil}の差の局所消去をもう一度用いる。
印以外で差が非零なら $Z=f_0 V$ とし、$\operatorname{ann}(f_0)$ を法として相対割線を得る。
その行列式二次で切る平坦族は上に述べた次数2の部分である。
よって位置1の印から射影した三座標 $l$ に対して、$D_3$ で消える二次式 $Q$ は $Q(l)Z=0$ を満たす。
曲線状の印では $Z=t^sV$ のとき、$l/t$ の方向が $t^{m-s-1}$ を法として $D_3$ を通る。
従って $Q(l)=t^2 Q(l/t)$ は $t^{m-s}$ を法として零である。
$m-s\le1$ でもこの結論は成り立つ。被約の印で差は零である。
ゆえに各第二辺関係の評価 $f\in C=U_0U_1$ に対し $Q(l)f=0$ が全域で成立する。

このような $f$ は零である。$B$ に長さが一つずつ増えるイデアルの列を取り、
各イデアルを $l$ の三つの一次式の中に座標制限の同型で持ち上げる。
前のイデアルの全ての持上げが $f$ を消すとし、次のイデアルの元を $h$ と書く。
長さ1の商の支持方向で消える一次式 $k$ に対し、$kh$ の値は前のイデアルに入る。
自明化が単位 $1\in B$ を与える一次式を $h_0$ とし、その値と同じ前段の持上げを $h'$ とすれば
$kh-h_0 h'$ は $D_3$ 上で消える二次式である。従って $khf=0$。
これは $hf$ の支持が位置1でその方向と印を結ぶ直線に入ることを意味する。
第一座標のみでこの支持環へ全射することは三つの被約割線の場合と同じである。
$h\in U_1$ と $\tau_X(fhU_0)=0$ により $hf=0$ を得る（$\tau_X$ は長さ20の辺の完全跡）。
帰納して全 $l$ が $f$ を消す。補題\ref{lem:finite-adjacent-word-support}から $f=0$ である。
\end{proof}

% !TeX root = ../papers/nc-p3.tex
\section{平方零の三重点と長さ六の平面切断}
\label{sec:finite-marked-planar-triple}
\begin{theorem}\label{thm:finite-marked-planar-triple}
共通の印付き切断で階数1軌跡が非Gorensteinの長さ3なら、その四平面は
\[
 E_D=\left\{\begin{pmatrix}a&b^t\\ b&tD\end{pmatrix}:
       a,t\in k,\ b\in k^2\right\},\qquad D\in\mathrm{GL}_2,\quad D\ne D^t
\]
に左右基底変更で一致する。$D$ の同一視はスカラー倍と合同で尽きる。
印が曲線状または被約なら、この模型でも $T$ が反対位置の対称性を与える。
従って接次元1以下の印では長さ3の全ての階数1切断がこの対称性判定に含まれる。
\end{theorem}
\begin{proof}[模型と二つの平面]
非Gorensteinの長さ3の環は一つの点に支持され、その極大イデアルは二次元で平方零である。
支持行列を $M=E_{11}$ とする。Segreの接行列は右下二行二列が零である。
二つの接ベクトルの左下二座標も右上二座標も各々独立である。
一方で従属なら線型結合で $M$ と共通像または共通余像を持つ接方向があり、
その直線全体が $E$ 中の階数1行列になる。従って基底変更で双方を同じ $b$ とできる。
第四の行列の右下 $D$ は可逆である。さもなければその左核の元は $E$ 全体を階数1以下に縮約する。
この可逆性による行・列加算で第四行列の上段と左列を消せる。接二行列には $M$ の倍を
加えるだけである。$a=1$ で階数1の式は $tD=bb^t$。
$D$ が対称なら、$t$ を消去した二本の二次式は二変数平方零の三独立二次式を張れず矛盾する。
非対称なら $t=0$, $bb^t=0$ を得て、実際に平方零環になる。$a=0$ に他の階数一点はない。
この支持を保つ左右変換はそれぞれ $ke_1$ を保つ。二つの下側二行二列の変換は
共通の接変数 $b$ を対応させるので互いに比例する。ゆえに $D$ の変換はスカラー倍と合同であり逆も明らかである。

中央文字は $E_D$ の座標 $(a,b,t)$、位置1の座標は印の補座標と射影左方向 $l=(l_0,l')$ とする。
$T$ で同定した位置0,2の三平面内の射影平面を $P=\{t=0\}$、位置1の平面を $L=\{l_0=0\}$ とする。
$R$ を辺 $01$ の六関係、$R'$ を辺 $12$ を反転して $T$ で移した六関係とする。
両空間は $l^t E_D$ の三成分を共有する。実際これらは $p_2$ と $p_0$ で縮約した後に位置3で縮約した
もので、共通切断の同定そのものである。$P\times L$ では三式は $(l'^t b,0,0)$ に制限される。
$R$ の全制限の独立数は4以下で、その基点が有限なのでちょうど4である。
そこで $Y=X_{01}\cap(P\times L)$ は四双一次式の完全交叉で長さ6である。
その $(1,1)$ の制限像は次元5、位置1だけの二次式の制限は $H^0(\mathcal O_Y)$ への
自明化後の同型である。前者は消える四切断が定義式で尽きること、後者は四式のKoszul分解を
$(0,2)$ でひねった負の全項のコホモロジーが消えることによる。

\end{proof}
\begin{proof}[誤差と三行三列の余因子]
平方零環の相対割線は全て $P$ にあり、左核は $L$ に対応する方向にある。
実際、階数1の切断に含まれる相対次数2の部分の線型包は $P$ に入る。
その鉛筆は剰余で $M$ と一つの非零接行列に張られる。相対的にも $M$ を含む：
平方零環の基底を $1,x,y$ とすると、相対直線イデアルの生成元
$u+vx+wy$ は $v,w$ の一つが単元であり、$x,y$ の乗法でその直線を保つため $u=0$ である
（例えば $v$ が単元なら $y$ による積 $uy$ が生成元の倍であることからその倍数が零になる）。
従って長さ2の商は定数点を含む。$M$ を消す左核は $l_0=0$ を満たす。

$X=X_{01}$ 上で $R'$ を評価する。印以外の局所因子で差 $Z=f_0 V$ があれば、
$\operatorname{ann}(f_0)$ を法として、前節までの消去で独立な二元 $A,V$ が相対割線を張る。
したがって任意の評価 $f$ について $t(x_0) f=l_0(x_1) f=0$ である。
ここで位置0の $t$ は $T$ を通した座標、$x_0$ は適宜 $T x_0$ と読む。
印の環を $k[z]/(z^m)$ とし、$Z=z^s V$ とする。$z^{m-s-1}$ を法として同じ結論を
$l/z,r/z$ に適用すると $l_0(x_1) f=0$ で、$t(x_0) f$ は少なくともsocleに入る。
$m-s=1$ の場合も $l_0$ が印で零なので同様、被約の印では差は零である。
実は $t(x_0) f=0$ である。これは印以外で零、印でsocleに入り、
$\tau_X(t(x_0) f U_1)=0$ は隣接座標積の自己直交性である。$U_1$ は印で単元を含むので
完全跡により零となる。ゆえに全評価は $C\cap\operatorname{ann}(I_Y)$ に入る。
この空間は次元1である。実際 $C=C^\perp$ よりその次元は $C+I_Y$ の余次元、
すなわち $6-5$ に等しい。

非零評価 $f$ があると仮定する。その支持は補題\ref{lem:finite-adjacent-word-support}より
長さ6以上で、$Y$ に含まれるから $Y$ に一致する。$H=R\cap R'$ は五次元である。
各局所環で二つの偶座標の差の成分が張るイデアルは $(f)$ と等しい。
実際、第二辺の位置1での評価行列は $X$ 上で階数3で、核直線は $x_2$ である。
単元な三次小行列で消去すれば、その六式の $T x_0$ での値が張るイデアルは
規格化した $T x_0,x_2$ の座標差が張るものに等しい。これが共通四平面の $Z$ の成分のイデアルである。
一方で全評価は $f$ の体スカラー倍である。

共有三式の制限が一本なので、$H$ は $P\times L$ 上の三本以下の双一次式を与える。
さらに一本で長さ6の有限完全交叉を切るのでちょうど三本で、その三行三列の係数行列を
$J(x_1)$ とする。これは $L$ の一般点で可逆である。そうでなければ三式の零点が平面 $L$
を覆って次元2以上となり、さらに一本のSegre超平面では有限にならない。
印以外の $Y$ 上では $A,V$ が $P$ に入り $H(x_1)$ の独立な二核元となる。
実際 $H$ は両方のグラフを消し、差の共通生成元で割れば $V$ を $\operatorname{ann} f$ を法として消す。
印では同じことが局所長さを1だけ減らせば成り立つ。上で $z^{m-s-1}$ を法として得た
相対割線全体が $P$ に入るからである。従って $J$ の全二次小行列式は、$Y$ の長さ5以上の部分で消える。
位置1の二次式が $Y$ の長さ6を分離するので、全余因子は同一の二次式の定数倍である。
一般可逆性から $\operatorname{adj}J=qK$ の定数行列 $K$ は可逆である。
随伴行列の積の式で $J$ 自身が一本の一次式の $K^{-1}$ 倍となり、その直線上で階数0に落ちる。
第四の一式を加えてもその上で有限交叉とならず矛盾する。従って全評価は零である。
\end{proof}

% !TeX root = ../papers/nc-p3.tex
\section{有限窓の反対座標積は十次元または十六次元}
\label{sec:finite-opposite-product-gap}
\begin{theorem}\label{thm:finite-opposite-product-gap}
四辺が有限な標準有限Frobenius窓で、反対二位置の座標を $U,W\subset D=H^0(\mathcal O_X)$ とする。
$U\otimes W\to D$ が単射でないこと、二移動での座標射が射影線型に同定されること、
その側の反対交換で四テンソルが不変になる表示を持つことは同値である。
この場合は $\dim UW=10$、それ以外では $\dim UW=16$ である。
点スキームの被約性、反対切断階数、印の存在、全次数の正則性を仮定しない。
\end{theorem}
\begin{proof}[階数一と二の双一次式]
定理\ref{thm:finite-square-chart}の有限完全交叉の証明により
$W^2=hU^2$ となる単元 $h$ があり、両座標の二次・三次の制限は単射である。
これらは有限Frobeniusの四移動定理により現在の候補にも適用できる。
$u b=0$ ($0\ne u\in U,0\ne b\in W$) なら $u(h^{-1}b^2)=0$ が三次多項式の恒等式となり矛盾する。
階数2の式を $u b=v c$ とする。各 $w\in W$ に対し $h^{-1}bw,h^{-1}cw$ を
$U$ 上の一意な二次式として持ち上げる。これらに $u,v$ を掛けたものは三次式として等しい。
一次式 $u,v$ は独立なので、ある $A(w)\in U$ により
\[
     h^{-1}bw=v A(w),\qquad h^{-1}cw=u A(w).
\]
$A$ は線型で、$bw=0$ は二次制限の単射性から $w=0$ を与えるため同型である。
隣の位置の座標と $W$ の間の任意の関係で $w$ を $A(w)$ に置き換えて $X$ 上で評価する。
得た $f$ は隣接座標積に属し、二等式から $uf=vf=0$ である。
その支持は $U$ 側の直線に入るので補題\ref{lem:finite-adjacent-word-support}により $f=0$。
グラフの一意性で $A$ が二移動を与え、有限辺の反転からテンソル対称性を得る。
逆に対称表示から二移動の同定が従う。実際二辺が反転し、有限グラフの両射影が同型である。
二移動の射影同定があれば、線束の自明化の単元を除いて $W=U$ であり積は十次元である。

\end{proof}
\begin{proof}[階数三は不可能]
$W$ の基底について積 $h^{-1}w_i w_j$ の二次持上げを対称行列 $P(x)$ に並べる。
その十成分は $U$ の全二次式の基底であり、$X$ 上では全二次小行列式がスキームとして零である。
混合双一次関係の係数を一次式の列 $m(x)$ とすれば、$P(x)m(x)=0$ が多項式として成り立つ。
実際各成分は元の式に $W$ の一文字を掛けた三次式である。逆も、$W$ に単元を取れることから従う。

階数3の関係があれば基底で $m(x)=(v_1,v_2,v_3,0)^t$ とできる。
$P$ の左上三行三列 $H$ は $H v=0$ により
\[
 H=[v]_\times^t C [v]_\times
\]
と一定の対称行列 $C$ で表せる。ここで $[v]_\times z=v\times z$。
この表示を確認する。第四変数で展開した各係数 $F(v,z)=z^t H_j(v)z$ は
$z\mapsto z+t v$ に不変で $v$ に関する次数が $j\le2$ である。
$F(v+t z,z)=t^{-2}F(v+t z,v)$ なので、$j<2$ の係数は零、$j=2$ では
左辺が $t$ に依存しない。よって $F$ は二つの剪断に不変で、独立な二列の対では
$v\times z$ だけに依存する（同じ外積を持つ基底は $\mathrm{SL}_2$ で移る）。
外積の非零座標ごとにその座標だけを分母に用いて二列を復元できる。
従ってこの関数は三空間の原点外で正則で、座標環の一意分解により原点にも延びる。
次数から外積の二次式となる。これで表示を得た。
$C$ が階数1以下なら $H$ は高々二独立な一次式の列の平方で、その六成分は独立になれない。
従って $q(v)=v^t\operatorname{adj}(C)v$ は非零である。
余因子行列の積の式から $\operatorname{adj}(H)=q(v)vv^t$ なので $X$ 上で $q(v)v_i v_j=0$。
三一次式の共通零点が $X$ にないなら直ちに $q=0$ となる。
あるならその点の極大イデアル $\mathfrak m$ をこの三式が生成し、$q$ はそこで
$\mathfrak m^2$ に消され、他では零である。$q\in U^2$ は隣接座標の平方全体と完全跡で直交する。
その平方は $\mathcal O_{X,p}/\mathfrak m^2$ に全射する（単元と、閉埋込みの一次座標差を用いる）。
ゆえに $q$ は零でなければならず、二次制限の単射性に反する。

\end{proof}
\begin{proof}[階数四の関係の二次行列]
階数4なら $m(x)=x$ として $P(x)x=0$ となる。一般階数3と仮定すると
$\operatorname{adj} P=\lambda(x)xx^t$ で、$\lambda$ は四次式である。
実際関数体での核が $kx$ なのでこの形になり、四変数に共通因子がないため係数は多項式である。
全三次小行列式とその各偏微分は $X$ 上で零である。偏微分は二次小行列式の多項式結合である。
$X$ の各点で単元な $x_i$ を選び $\lambda x_i^2$ とその偏微分を使うと、
全ての $\partial_j\lambda$ が $X$ 上で零となる。三次制限の単射性と標数零により $\lambda=0$ となり矛盾する。
従って $P$ の一般階数は2以下である。

一般階数2のとき $\wedge^2 P$ は対称な階数1の四次多項式行列で、$d(x)b(x)b(x)^t$ と書ける。
これは関数体で階数1の列を取り、多項式環で分母と全成分の共通因子を除き、
原始列への比例係数が多項式になることと対称性を使えばよい。
$b$ は一般点で $P$ の像の二平面を表す二形式で、$P x=0$ から $\iota_x b(x)=0$。
$\deg b=0$ は非零の定数二形式では不可能である。
$\deg b=1$ なら $b(x)=\iota_x\omega$ となる一定の三形式 $\omega$ がある。
実際偏極すると係数の $x$ 添字と縮約される添字も交代する。
四次元の三形式は非零の根基 $kz$ を持つので、一般の $x$ における $P$ の像二平面は
$z^\perp$ に入る。これは全二次式を基底とする $P$ に定数核を与えて矛盾する。
よって $\deg b=2$, $d$ は非零定数である。関数体上で非零な主二次小行列式は定数倍の平方なので、
標数零の代数閉基礎体より、階数2の二次式 $y^t P(x)y$ は関数体上で二一次式に分解する。
係数二次式は共通因子を持たないのでGaussの補題で多項式環上に分解が降りる。
一方の因子が $x$ に関して定数なら値域は固定一次式の倍だけになり、$P$ の全線型包に反する。
従って両因子は双一次で
\[
 y^t P(x)y=(y^t Sx)(y^t Tx).
\]
$P(x)\ne0$ は全射影点で成り立つので $S,T$ は可逆である。
一般の $Sx,Tx$ は独立で $P(x)x=0$ のため $S,T$ はともに交代行列である。
その鉛筆のPfaffianの根を取り $L=T-cS$ を階数2にする。両行列は比例せず、
$S$ は可逆なのでこのような有限な $c$ がある。
$L$ の像を根基に持つ階数2の交代行列 $J$ を取れば
$(Sx)^tJSx=(Tx)^tJSx=0$。従って $P(x)JSx=0$ は階数2の混合関係を与え、最初の証明が適用される。

一般階数1の場合も同じ原始列の議論で $P=d b b^t$ と書け、全成分の次数と独立性から $d$ は定数、
$b$ は可逆な一次列である。この場合 $X$ 上の $W$ 座標点は $[b(x)]$ に等しい。
実際二次Veroneseは閉埋込みで、$P$ は $X$ 上でも両者の同じ二次座標を与える。
これで全階数の場合が尽きる。
\end{proof}

% !TeX root = ../papers/nc-p3.tex
\section{被約四頂点切断の六つの反対座標積}
\label{sec:finite-marked-tetrahedral-products}
\begin{theorem}\label{thm:finite-marked-tetrahedral-products}
共通の印付き切断で $\Sigma$ が四つの被約点とし、印の接次元は1以下とする。
偶位置の座標を四頂点に対応する $a_i,b_i$ ($1\le i\le4$) に選び、
二乗空間を一致させる単元により自明化すると
\[
                      b_i b_j=a_i a_j \quad(i\ne j)
\]
となる。偶位置が対称でなければ六つの $g_{ij}=a_i b_j-a_j b_i$ ($i<j$) は独立で、
その各支持スキームは長さ1または2であり、異なる六つの支持は交わらない。
$d_i=b_i^2-a_i^2$ を $a$ の一意な二次式と見ると、どの三つの異なる添字でも
$d_i a_j a_k=0$ が点環で成り立つ。少なくとも三つの $d_i$ は非零である。
\end{theorem}
\begin{proof}
長さ4の積跡模型を使う。四つの階数一行列を $u_i u_i^t$ とし、
三次元の $u_i$ はどの三個も独立である。
$E$ の二つの評価を $A=\sum b_i u_i u_i^t$, $B=\sum a_i u_i u_i^t$ とする。
局所自明化による各列全体への単元倍は以下の消失条件に影響しない。
位置1の三射影座標を $l$、$l_i=u_i^tl$ と書く。位置3にも同じ座標表示を使う。
例外鉛筆は二頂点の張る直線で、Artin基底上でもそのような鉛筆は一定である。
実際相対的な次数2の階数一因子は定数な四点の直和因子二つに限られる。
鉛筆 $ij$ に対応する位置1の直線は $l_i=l_j=0$ である。

印以外の任意の局所因子で $A,B$ が比例しないとき、前節までの縮約行列の主元消去により、
正規化した差は $z V$ となる。$S/\operatorname{ann}(z)$ 上で $A,V$ は剰余で独立であり、
二点 $i,j$ の張る固定鉛筆の基底となる。従って外積 $A\wedge B$ の $ij$ 以外の成分は零、
$ij$ 成分は $a_k,b_k$ ($k\notin\{i,j\}$) および $l_i,l_j$ に消される。
印でも、$S=k[t]/(t^m)$, $z=t^s$ のとき $l/t$ を用いて
$t^{m-s-1}$ を法として同じ鉛筆を得る。従って全ての異なる $i,j,k$ で $b_k g_{ij}$ は
印に支持されるsocleに入り、$l_i g_{ij}=l_j g_{ij}=0$ が全域で成り立つ。
実際、指定鉛筆以外の外積成分が印で残ってもsocleに属し、$l_i$ は印で零である。
$m-s\le1$ なら初めから全外積成分がsocleに入る。被約の印でも結論は成立する。

単元 $h$ を $W^2=h U^2$ と取り、$P_{ij}(a)$ を $h^{-1} b_i b_j$ の二次持上げとする。
$x$ を印の一つのsocle生成元とし、$C(a)$ をそれを表す一意な三次式とする。
各相異なる $i,j,k$ について
\[
 a_i P_{jk}-a_j P_{ik}=c_{ijk} C
\]
が多項式恒等式である。$k$ を固定し、他の三添字 $i,j,n$ の三つの式に
$a_n,a_i,a_j$ を掛けて巡回和を取ると左辺は零である。
$C\ne0$ と三一次式の独立性から全係数が零である。
従って $P_{ik}=a_i h_k$ ($i\ne k$) と書け、$P$ の対称性から $h_k=c a_k$ が
共通の非零定数 $c$ で成立する。$ch$ の平方根で $b$ の自明化を変えれば最初の式を得る。
これで $a_k g_{ij}=b_k g_{ij}=0$ は印を含む全域で成り立つ。
偶側が非対称なら定理\ref{thm:finite-opposite-product-gap}より六元は独立で全て非零である。

各 $g_{ij}$ の支持は偶位置の座標直線 $ij$、位置1の直線 $l_i=l_j=0$ の積で
六関係を満たす。従って一般の二式の交点数から長さは2以下である。
印以外の局所因子上では既に高々一つの割線に限られる。
印では $A(p)=B(p)$ の階数は $3-\dim T_pX\ge2$ なので、これは高々一本の座標直線に入る
（異なる直線の交点は階数一頂点である）。偶直線条件が全域で成り立つから印でも一つに限られる。
最後に $d_i a_j=-b_i g_{ij}$ ($i\ne j$、外積は添字交換で符号を変える) より四次の消失を得る。
二つの $d_i,d_j$ がともに零なら $b_i,b_j$ は $g_{ij}$ をともに消す。
その支持は $b$ 側でも直線 $ij$ に入るので空となり、独立性に反する。
\end{proof}

% !TeX root = ../papers/nc-p3.tex
\section{被約四頂点切断の左右核座標の反対対称性}
\label{sec:finite-tetrahedral-kernel-symmetry}
\begin{theorem}\label{thm:finite-tetrahedral-kernel-symmetry}
四辺が有限で全中央支持4の有限Frobenius窓に共通印付き切断 $E$ があり、
その階数一切断 $\Sigma$ が被約四点であるとする。印の接次元が高々1なら、$E$ の
行・列に当たる二位置（以下、奇位置）の間に射影線型な二移動があり、その側の反対対称性が成り立つ。
他の局所因子の接次元・被約性を仮定しない。
\end{theorem}
\begin{proof}[外積の支持と三次式への持上げ]
位置1,3の印からの三射影座標をそれぞれ $l,r$ とする。
四頂点表示を $A=\sum_i a_i u_i u_i^t$, $B=\sum_i b_i u_i u_i^t$ と取り、
$\sum u_i=0$ となるよう四方向を規格化する（係数は二つの偶座標に吸収する）。
どの三方向も独立である。$l_i=u_i^tl$ とすると $Al=0$ より
$a_i l_i$ は四つとも等しい。$Al=Ar=Bl=Br=0$ は点環 $D$ 上の等式である。
奇側が非対称と仮定する。$k=l\times r$ の三成分は定理\ref{thm:finite-opposite-product-gap}より独立である。

$A$ の各列 $v$ は $l,r$ の両方を消すから $v\times k=0$ である。
各局所因子で $A$ の単元成分を持つ列を取り、その列に比例して $k=\eta v$ と書ける。
他の全列に同じ恒等式を使うと $A$ の全二次小行列式は $\eta$ に消される。
従って $k$ の成分の生成イデアルは局所主で、その零化イデアルが定める $Z$ は
$A$ により $\Sigma$ に入る。偶座標から $\mathbb P(E)$ への写像が線型同型であるため、
$Z$ は相異なる $n\le4$ 個の被約点である。$B$ にも同じ議論を適用すると、
$Z$ 上で $A,B$ の列は同じ $v$ の方向にあり、二つの偶座標は同じ頂点に対応する。
$Z$ の各点 $j$ での $k$ は零化イデアルが極大イデアルのベクトル、すなわち
$u_j$ と非零socle元の積である。三成分の独立性から $n=3,4$ である。
印での $A$ は階数 $3-\dim T_pX\ge2$ なので $Z$ は印を含まない。

位置1,3の四座標を $U,W$ とし、単元 $h$ を $h^{-1}W^2=U^2$ と取る。
$c_j$ を上のsocle元の $h^{-1}$ 倍を表す位置1の三次式とする。三次制限が同型なので
これは一意である。局所自明化に従い、$b\in W$ の値を $b(j)$ と書く。このとき
\[
                    \sum_{j\in Z} b(j)\,l_j c_j=0
                    \qquad (b\in W)
\]
が多項式恒等式である。実際 $h^{-1}b k$ の三次持上げは一方では
$\sum_j b(j)c_j u_j$、他方では $l$ と、三つの $h^{-1}b r_\nu$ の二次持上げの列との外積である。
後者と $l$ の内積は多項式として零である。
$W$ の $Z$ 上の値の階数は2以上で、$n$ より小さい。全射なら各非零積 $l_j c_j$ が零になるからである。

階数が $n-1$ なら定数な一つの線型従属の係数は全て非零であり、$l_j c_j$ は
いずれも同じ多項式の非零定数倍である。異なる $l_j$ は比例しないので
\[
                    c_j=\gamma_j H\prod_{i\in Z,\ i\ne j}l_i ,
      \qquad \gamma_j\in k^\times,\quad\deg H=4-n .
\]
\end{proof}

\begin{proof}[長さ三と四の排除]
$n=3$ なら $Z$ の添字を $1,2,3$ とし $H$ は非零一次式である。
隣接座標積 $H a_4$ の支持長は高々5であることを示す。
印以外では四つの $l_i$ のうち少なくとも二つが単元である（三方向ずつ独立である）。
三つの $c_j$ を法とすると、$1,2,3$ のうち二つが単元の場合には $H=0$ となる。
残りの場合は $l_4$ と一つが単元で、他の二つの $l_i$ に $H$ を掛けると零になる。
$a_4 l_4=a_i l_i$ より $H a_4=0$ である。従って印以外で $H a_4$ は三つのsocleの包に入る。
印の環を $k[t]/(t^m)$ とし $m>2$ の場合だけを考える。$l/t$ の剰余は非零なので
$l_i/t$ の少なくとも二つが単元である。$c_j$ はここで零であり、$t^2$ で割った式と
$a_i l_i/t$ の一致を法 $t^{m-2}$ で用いると同じ議論で $Ha_4=0\bmod t^{m-2}$。
これで支持長の上界を得る。補題\ref{lem:finite-adjacent-word-support}から $Ha_4=0$ となるが、
これは有限Frobenius窓に純関係がないことに反する。

$n=4$ で値の階数が3なら $H$ は定数である。印以外で二つの $l_i,l_j$ を単元に取ると、
他の二つの $l$ 座標は $c_j$ 全体のイデアルに入る。
従って $a_i,a_j$ もそのイデアルに入る。これは $Z$ 上だけのsocleイデアルである。
$Z$ の各点では $a$ の四座標のうちただ一つだけが単元であるため、相異なる三つの
$a$ 座標の積は印以外で零である。実際少なくとも一因子が上のsocleに入り、
さらに一因子が極大イデアルに入る。印で $m>3$ なら $t^3$ で割った三積の消滅と
$a_i l_i/t$ の一致から、少なくとも二つの $a_i$ は $t^{m-3}$ の倍となる。
従って相異なる三座標の四つの三積は、印の長さ高々3のイデアルに全て入る。
これは偶位置の三次制限の単射性に反する。

残りは $n=4$ で位置3での四点が共線の場合である。有限窓において、ある位置で
四点が共線で、その前点と後点がそれぞれ射影基底をなすことは不可能である。
ここで射影基底とは四次元線型空間の基底の四点を指す。
実際、前・後の点をそれぞれの座標頂点に取り、中央位置の直線を $y(t)$ とする。
制限 $w(-,y(t),-,-)$ の前後の座標 $i,j$ の成分は残る位置に値を持つ一次式で、
前の点との隣接から $t_i$ で、後の点との隣接から $t_j$ で零になる。
$i\ne j$ なら恒等零で、対角成分は $t_i$ で消える一次式と定数ベクトルの積である。
制限全体は零ではない（一文字縮約は非零）。非零の対角成分を前位置で取り出し
一般の $t$ を代入すれば、後位置から始まる関係空間に非零な純関係を得て矛盾する。
いまは $A,B$ による両偶座標が四頂点なのでまさにこの禁止配置である。
全ての場合が尽きたため奇側の反対対称性が従う。
\end{proof}

% !TeX root = ../papers/nc-p3.tex
\section{長さ四の全Artin切断と核側の反対対称性}
\label{sec:finite-artin-kernel-symmetry}
\begin{theorem}\label{thm:finite-artin-kernel-symmetry}
定理\ref{thm:finite-tetrahedral-kernel-symmetry}で $\Sigma$ は長さ4であればよく、
非被約でも奇位置の反対対称性が成立する。従って四辺が有限、反対階数8で、
接次元1以下の局所成分が一つでもある候補は、中央支持3の点から正則性が決まるか、
少なくとも片側の反対対称性を持つ。いずれも全次数の標準正則性が有限に決まる。
\end{theorem}
\begin{proof}[外積の環と線型な分母]
$B=H^0(\mathcal O_\Sigma)$ とし、定理\ref{thm:finite-marked-artin-secants}の完全跡表示を使う。
偶座標を $a\in B$ と書くと、その対称行列は $\tau(a u v)$ ($u,v\in V=\ker\tau$) である。
点環上で奇位置の射影座標 $l,r$ は $V$ に値を持ち、$al=\alpha 1$ があるスカラー $\alpha$ で成り立つ。
$V$ の非零元の乗法階数は2以上である。階数1ならその元の張る直線がイデアルとなるが、
完全跡の核は非零イデアルを含まない。

非対称を仮定し、前定理と同じく $k=l\times r$ の成分イデアルを $(\eta)$ とし、
$S=D/\operatorname{ann}(\eta)$ と置く。局所生成元は直積環上の一元にまとめられる。
前定理の外積の恒等式から $Z=\operatorname{Spec}S\subset\Sigma$ であり、$S=B/I$ でもある。
$Z$ 上で二つの偶行列の全列は同じ直線に入り、対称性と単元成分により両行列は単元倍で一致する。
$\Sigma$ の $\mathbb P(B)$ での普遍点 $a$ は、完全跡により評価指標に双対な直線である。
実際指標 $\chi$ の双対元は $b a=\chi(b)a$ を満たし、制限形式は
$\chi(u)\chi(v)$ となる。この写像が完全な $B$ の座標系による閉埋込みで、
長さ4の切断に含まれることから切断全体に一致する。
従って $Z$ 上で $a$ の全乗法行列 $m_a$ も階数1である。
$k=\eta (u_1,u_2,u_3)$ の $u_i\in S$ は、単元倍を $\eta$ に吸収すれば
$V$ の三基底の $S$ への像である。$\dim S=n$ は3または4、印はその台にない。
$S$ は $D$ の完全跡に $\eta$ を掛けて得る完全積対合によりGorensteinである。

位置3の四座標の $S$ への像を $W'$ とする。$h^{-1}W^2=U^2$ として
$h^{-1}\eta s$ の一意な位置1三次持上げを $\Phi(s)$ と書く。
$\Phi$ は単射で、$S$ の完全跡 $\tau_S$ により $\Phi(s)=\tau_S(Cs)$ と表せる。
ここで $C$ は位置1の三次同次式を係数とする $S$ の元である。
独立変数での $l$ の $S$ への射影を $L$ とする。前定理と同じ外積に $l$ を掛けた恒等式より
$\tau_S(LC b)=0$ ($b\in W'$) である。一般体上で $C$ は単元である。
そうでなければ、ある局所因子の定数socle元を $\Phi$ が消す。$L$ も一般体上で単元である。
よって $W'\ne S$ で、その次元は閉埋込みより2以上である。

まず $W'$ が余次元1なら $LC=\lambda c$、$c$ はその直交補を張る定数単元、$\lambda$ は四次式である。
$S$ が曲線状の局所因子の積なら、$L^{-1}$ の最小共通分母は $N=\det m_L$（定数倍を除く）。
実際各因子 $k[e]/e^d$ で $L=L_0+eL_1+\cdots$ と書くと、$L_0$ に沿う極は正確に $d$ 次である。
$d\ge2$ では方向射影 $\Sigma\to\mathbb P(V^*)$ の閉埋込みにより $L_1$ は $L_0$ に比例せず、
逆数展開の $e^{d-1}$ 係数の分子は法 $L_0$ で $(-L_1)^{d-1}$ となる。
この閉埋込みは異なる二点と全接方向を $V$ が分離することによる
（任意の長さ3の部分が三座標を張る）。異なる局所因子の $L_0$ は比例しない。
従って $C=H c\,N L^{-1}$、$\deg H=4-n$ である。$NL^{-1}$ の乗法行列は
$\operatorname{adj}m_L$ であり、全係数を実際の点環に戻すと
$H\operatorname{adj}m_L$ の成分は $(\eta)$ に入る（$C$ の係数関数は $\Phi(S)$ を張る）。
\end{proof}

\begin{proof}[乗法の階数と印での長さ]
$n=3$ のGorenstein環は曲線状の因子の積である（局所埋込次元2なら長さ3では平方零となる）。
$I$ は長さ1のイデアルで、$B$ は $I$ に指標 $\chi$ で作用する。
印以外で $H\chi(a)\in(\eta)$ を示す。まず $(\eta)$ を法とし、さらに $H$ の零化イデアルで割る。
この局所環上で $\operatorname{adj}m_L=0$、すなわち $S$ 上の乗法の全二次小行列式は零。
一方 $l$ の剰余は非零で、$B$ での乗法には単元な二次小行列式がある。
$B\to S$ の核は一次元なので、$S$ 上での像は階数1の直和因子である。
$B$ 上の $lB$ はこの直線の逆像全体（階数2の直和因子）である。
実際その逆像に入り、剰余で二つの独立列を持つためである。特に $I$ を含む。
$a(lB)=\alpha B$ は二元で生成されるため $\alpha=0$。非零なら $\alpha B$ は基底環の
非零な主イデアルの四コピーで、Nakayamaにより四生成元を要する。
従って $aI=0$、すなわち $\chi(a)=0$ であり、元の環に戻して所望の包含を得る。
印 $k[t]/t^m$ で $m>2$ なら $l/t$ を法 $t^{m-2}$ で同じ議論に用いる。
$H\operatorname{adj}m_L$ の次数2の因子を $t^2$ で割れ、$a(l/t)$ もスカラーである。
よって $H\chi(a)$ の支持長は $3+2$ 以下である。非零な隣接語の下界6と純関係の不在に反する。

$n=4$, $S=B$ が曲線状の因子の積なら $H$ は定数である。印以外で $(\eta)$ を法とすると
$\operatorname{adj}m_l=0$ かつ $m_l$ に単元な二次小行列式がある。従って像は階数2の
直和因子で、先と同じ生成元数から $\alpha=0\bmod(\eta)$ である。
元の局所環でも $lB$ の二列を基底の二元に取れる。$m_a$ はそれらを $\alpha$ の倍に移すので、
三次小行列式は $\alpha$ と $m_a$ の二次小行列式のイデアルとの積に入る。
列の基底変更を行ってその列で展開すればよい。
$Z$ 上で $m_a$ が階数1であったため、この積は零である。
従って印以外で $\operatorname{adj}m_a=0$。印で $m>3$ なら $l/t$ と法 $t^{m-3}$ での
同じ階数2の像を使い、$a$ がそれを消すのでやはり $\operatorname{adj}m_a=0\bmod t^{m-3}$。
しかし $N(a)a^{-1}$ の四座標は独立な三次式である。単元の反転が射影単元開集合を
全体へ移すことから従う。その四式が長さ3以下のイデアルに入るのは不可能である。

長さ4の残る非曲線状Gorenstein環 $B$ は局所で Hilbert 関数 $(1,2,1)$ を持つ。
実際埋込次元3なら平方零である。極大イデアルの二乗はsocle $k\omega$、三乗は零で、
一次の二座標上の乗法は非退化二次形式である。
$V\to B/k\omega$ は同型なので、$l=s+v$（$v$ は根基に値を持つ）に対し
逆数の最小分母は $s^3$ で、分子は $s^2-sv+v^2$。
二次形式 $v^2/\omega$ が法 $s$ で零にならないことを用いた。
従って $C=H c(s^2-sv+v^2)$、$H$ は非零一次式である。
剰余指標を $\chi$ とすると印以外では $H\chi(a)\in(\eta)$ である。
実際法 $(\eta)$ で $s$ が単元なら $H=0$。そうでなければ $v$ の一次二座標の一つが単元で、
$Hsv=0$ を根基の二乗を法として読めば $Hs=0$。
$al=\alpha 1$ の根基の一次部分に $H$ を掛けると $H\chi(a)v=0$ であり結論を得る。
印で $m>2$ なら次数2の分子を $t^2$ で割り同じ議論を行う。
従って隣接語 $H\chi(a)$ の支持長は $4+2$ 以下。6未満なら不可能で、6なら
印以外で長さ4を全て使い、その支持は $Z$ を含む。$Z=\Sigma$ の偶座標は全四空間を張るが、
補題\ref{lem:finite-adjacent-word-support}の長さ6の語はどちらの座標でも平面に入る。これも矛盾する。
\end{proof}

\begin{proof}[長さ四の共線スキーム]
最後に $n=4$, $\dim W'=2$ とする。$Z$ は位置3の直線における長さ4の因子である。
無限遠を外し $S=k[t]/P(t)$ とし、$P$ は四次モニックとする。両隣の偶座標はそれぞれ
$Z$ に全射する完全一次系なので、個別の自明化と基底で両方を
$e(t)=(1,t,t^2,t^3)^t$ と書ける。$t e=T e\bmod P$ と定数行列 $T$ を取る。
直線上の四テンソルを、残る一位置の任意の縮約ごとに $M(t)=M_0+tM_1$ として両偶位置の間で読む。
隣接条件 $e(t)^t M(t)=M(t)e(t)=0\bmod P$ と四基底の独立性より
$M_0=-T^t M_1=-M_1T$。$P$ の根 $d$ で $Te(d)=d e(d)$ となる。
従って $e(d)^t M_1$ は $T$ の固有値 $d$ の左固有ベクトルまたは零で、
この固有空間は巡回行列のため重根でも一次元である。
よって前位置を $e(d)$ に、中央位置を直線の $t\ne d$ に制限すれば、
残る位置の全縮約の後位置での像は固定直線に入る。
この二文字縮約は零ではない。零なら $e(d)$ が点辺で中央直線の全点を後続に持ち、
辺の有限性に反する。従って後位置の関係空間に純関係があり矛盾する。

これで長さ4の全切断が尽きる。反対階数8の系では各側で共通切断が自動であり、
有限長3以下と正次元の切断では先の各割線・母線の定理が切断の偶座標の対称性を与える。
長さ4では今回の奇側の結論を使えばよい。片側対称になった場合の全次数判定は
系\ref{cor:all-defective-one-sym-finite}を適用する。具体的には対称核に階数2以下の元があれば認証され、
それ以外は両対称窓の有限判定で決まる。
\end{proof}

% !TeX root = ../papers/nc-p3.tex
\section{極大イデアルの二乗が単項となる平面局所点}
\label{sec:finite-stretched-planar-marks}
\begin{theorem}\label{thm:finite-stretched-planar-marks}
四辺が有限な反対階数8の標準有限Frobenius窓に、接次元2で極大イデアルの二乗が
単項の局所点が一つあれば、全次数の標準正則性は有限に決まる。
従ってこの階数の有限枝で既出の正則性判定が被覆していない候補は、
閉点が高々三つで、全ての局所環が接次元2、長さ6以上、極大イデアルの二乗が非単項の場合に限る。
\end{theorem}
\begin{proof}[局所環と三重点切断]
三文字認証がないと仮定し、当該点を印とする。局所環 $P$ はGorensteinであり、
パラメータを取れば
\[
 P=k[x,y]/(xy,\ y^2-x^e,\ x^{e+1}),\qquad e\ge2 .
\]
実際 $\mathfrak m^2/\mathfrak m^3$ は一次元なので標数零で一般の $x$ の平方がその生成元となる。
$e>2$ なら $\mathfrak m^3/\mathfrak m^4$ も一次元で、一般の $x$ を $x^3\ne0\bmod\mathfrak m^4$ とも
選べる（三次の偏極による）。このとき $\mathfrak m^3=(x^3)$ で、順に $\mathfrak m^j=(x^j)$ ($j\ge2$)。
実際 $\mathfrak m x^j=x^{j-2}\mathfrak m x^2$ である。$e=2$ でも以下同様である。
第二生成元から $x$ の多項式を引いて $xy=0$ とする。すると $y^2\in\mathfrak m^2$
は $x$ に消されるので $x^e$ の倍で、零ならsocleの次元に反する。倍率を揃えて表示を得る。
どの $f\in\mathfrak m$ に対しても、非零の商 $P/\operatorname{ann}f$ は曲線状である。
実際 $\mathfrak m f$ は一列のイデアル $(x^j)$ からなる $\mathfrak m^2$ の部分加群だからである。
$P/(y)=k[x]/x^e$ とし、各元を $x$ の次数 $e$ 以下の多項式と $y$ の定数倍に分ける。
$\mathfrak m$ の三元以上の積の $x$ 表示は、各 $y$ 項を捨てて通常に掛け法 $x^{e+1}$ で取ったものに等しい。

共通切断の $\Sigma$ が正次元または長さ2以下なら既出の平面点判定を使う。
長さ3のGorenstein型での印の評価誤差を考える。正規化した二行列の差を $Z$ とする。
その全成分は $\mathfrak m$ に入る。定理\ref{thm:finite-marked-artin-secants}で使った三方向の
スキームで消える二次式 $Q$ について、$Q(l) Z=0$ を示せば以後の証明は変わらない。
$y$ を捨てた $Z$ の成分の最小 $x$ 次数を $s\ge1$ とし、$s\ge e-1$ なら三因子の積より直ちに従う。
その他では $k[x]/x^e$ において $Z=x^s V'$ とし $l,r$ を $x$ で割る。
商での $l,r$ の座標は極大イデアルを生成するので、割った列の剰余は非零。
$A,V'$ は法 $x^{e-s-1}$ で独立な相対割線の基底となる（正規化によりその剰余は独立）。
従って $Q(l/x)=0\bmod x^{e-s-1}$。二次の $x^2$ と $x^s$ を戻すと少なくとも
$e+1$ 次となる。三因子の積の表示により元の環でも零である。

非Gorenstein三重点では定理\ref{thm:finite-marked-planar-triple}の $t,l_0$ と評価 $f$ を使う。
印では $f,t,l_0\in\mathfrak m$ である。曲線状商 $P/(y)$ での元の割線の証明は
$t f$ がそのsocleに、$l_0 f$ が零になることを与える。元の環でこれらは $\mathfrak m^2$
に入るので、各々 $(x^{e-1}),(x^e)$ に含まれる（$e=2$ でも長さの上界は同じ）。
印以外で零だからそれぞれ長さ2以下で、対応する隣の一次座標がこの支持環に全射する。
隣接座標積の自己直交性より両方とも零である。これで前定理の長さ6の完全交叉 $Y$ と、
誤差全体を生成する一元 $f$ の議論が適用できる。
印の $P/\operatorname{ann} f$ 自身が曲線状なので、$Z=fV'$ を書いた後、
この環で $l,r$ を一様化変数で割れば、局所長さを一つ減らした上で $A,V'$ の相対割線を得る。
$A$ と $V'$ は五つの共有関係の評価核にあり（後者は差を $f$ で割る）、
ともに割線の平面 $t=0$ に入る。従って $Y$ の長さ5以上で余因子が消えるという矛盾まで同じである。
\end{proof}

\begin{proof}[長さ四の切断での二次の誤差]
定理\ref{thm:finite-artin-kernel-symmetry}と同じ外積 $k$、その成分イデアル $(\eta)$、
支持環 $S=B/I$ を用いる。今は印も $Z=\operatorname{Spec}S$ に入ってよい。
$W'$ が全射の場合と共線四点因子の場合の排除には印の仮定がない。
残る余次元1での分母と、印以外での乗法行列の議論もそのまま成り立つ。

$n=3$ のとき $f=H\chi(a)$ と書く。印で
\[
                          \mathfrak m^2 f\subset(\eta)
\]
を示す。$(\eta)$ を法とし、さらに $H$ の零化イデアルで割る。
この環上で三行の乗法行列 $M=m_L$ は全二次小行列式が零で $m_{\bar a}M=\alpha I_3$。
従って $\alpha M=0$ である。実際 $\alpha M_{ij}$ は $k\ne j$ として
$\sum_h(m_{\bar a})_{kh}(M_{hk}M_{ij}-M_{hj}M_{ik})$ に等しい。
$M$ の成分は $l$ の三座標を張る（$V\to S$ は同型）。
$I$ を先頭の基底とする $m_l$ の下三行は $(0,M)$ である。
それらの二次小行列式が零で、$m_a m_l=\alpha I_4$ の第一行を用いれば
$\chi(a)$ は $m_l$ の二次小行列式を全て消す。実際第一行を含む小行列式に
$\chi(a)$ を掛け、第一行を積の第一行で置き換えると、残りの項は下三行の小行列式、
または $\alpha$ と下三行の成分との積である。

これらの小行列式は印で $\mathfrak m^2$ を生成する。接方向の二次元鉛筆は
$V$ の二平面（$l$ が接ベクトルを分離する）で、どの非零方向でも乗法階数は2以上である。
その二次小行列式が全二項二次式を張ることを確認する。さもなくばそれらを消す
二次汎関数を取り、平方ならある方向の階数が1以下になる。
相異なる二因子なら基底変更後の二行列 $U,V'$ の混合二次小行列式が全て零である。
$U$ を階数 $r\ge2$ の単位ブロックにして微分すると $V'$ はそのブロックの外で零となる
（非対角成分には別の非零対角元を使う）。$r\ge3$ ならブロック内でも非対角は零、
対角はどの二つの和も零なので $V'=0$。$r=2$ なら二行二列の鉛筆となり、
行列式に根がある。いずれも矛盾する。次数環へ写しNakayamaを使えば生成が従う。
先の二つの商から戻れば表示の包含を得る。

$B$ が局所 $(1,2,1)$ で $n=4$ の場合にも $f=H\chi(a)$ に同じ包含が成り立つ。
同じ二つの商で $l=s+v$、$v$ の二つの接座標を $v_1,v_2$ とすると、逆数の分子より
$s^2=s v_1=s v_2=0$。$al=\alpha 1$ の一次部分は $\chi(a)v_i+s a_i=0$ なので
$\chi(a)v_i v_j=0$ でもある。$s,v_1,v_2$ は $l$ の三座標の基底である。

どちらでも $f$ が印で単元で $(\eta)$ が零なら $\mathfrak m^2=0$ で矛盾する。
印で $(\eta)\ne0$ なら、$a(p)$ は $S$ の支持に属する指標の跡双対元である。
$n=3$ の $\chi$ は欠けた長さ1のイデアルの指標で、この双対元に非零となるのは
その指標の因子が単純で同じ指標の場合だけである。その場合当該因子は $S$ に残っていない。
局所 $(1,2,1)$ の場合もsocle元の剰余は零。従って常に $f$ は印で非単元であり、
その局所支持環は曲線状である。よって $(f)/\mathfrak m^2(f)$ の長さは高々2であり、
全体の支持長は $n+2$ 以下となる。$n=3$ は下界6に反し、$n=4$ で等号なら
$(\eta)\subset(f)$ が全点で成り立つため支持が $Z$ を含み、偶座標の平面条件に反する。
\end{proof}

\begin{proof}[曲線状因子の積からなる長さ四の切断]
$n=4$ で $B$ が曲線状因子の積の場合は $\operatorname{adj}m_l$ の各成分が $(\eta)$ に入り、
印以外の $\operatorname{adj}m_a$ の消滅は既に証明した。印では $a(p)$ の乗法階数は1である。
実際 $A(p)$ は接次元2により階数1で $\Sigma$ に属する。
従って $m_a$ の三次小行列式は $\mathfrak m^2$ に入る。$e\le4$ ではこれは長さ3以下である。
$e\ge5$ とする。$\eta$ は外積から $\mathfrak m^2$ に入り、
$\eta\ne0$ なら $(\eta)=(x^s)$、$n_p=e+1-s$ とする。$\eta=0$ なら $s=e+1,n_p=0$ とする。
$Z$ 上で $m_a$ が階数1なので、その二次小行列式の $x$ 部分は法 $x^{n_p}$ で零である。
$s\le3$ ならこれだけで三次小行列式には $x^{e-2}$ 以上の次数がある。
$s>3$ では、$l$ の三因子の積の公式から、$y$ を捨てた $l$ を $x$ で割った $l'$ の
乗法の三次小行列式は法 $x^{s-3}$ で零である。この割算は $s=e+1$ でも、
三因子の積を法 $x^{e+1}$ で読めることから正当である。$l'$ の乗法には単元な二次小行列式がある。
法 $x^{e-1}$ で $al'=\alpha'1$ と書け、さらに法 $x^{s-3}$ では像が階数2なので
前定理の生成元数の議論で $\alpha'=0$。従って法 $x^{e-1}$ では、二つの基底列を $l'B$ から選んで
$m_a$ の三次小行列式を展開すると全て $x^{s-3} I_2(m_a)$ に入る。
$s-3+n_p=e-2$ だから、元の環でも三次小行列式は長さ3の $(x^{e-2})$ に入る
（既に $\mathfrak m^2$ に入るため $y$ の一次項はない）。
独立な四つの三次式をここに入れることはできず、長さ4の場合も核側の対称性を得た。

片側対称性からの正則性判定は系\ref{cor:all-defective-one-sym-finite}による。
中央支持が全て4なら印で $w(p_0,p_2)\ne0$ であり、接次元は3以下の行列の余階数として高々2である。
接次元2の長さ4または5の局所Gorenstein環は、平方零の二次元socleを持てないため
$\dim\mathfrak m^2/\mathfrak m^3=1$ である（長さ5でこれが2なら $\mathfrak m^3=0$）。
長さ20の和を取れば残余の閉点数の上界を得る。
\end{proof}

% !TeX root = ../papers/nc-p3.tex
\section{節点または尖点の曲線上の局所印と長さ六の場合}
\label{sec:finite-nodal-cuspidal-marks}
\begin{theorem}\label{thm:finite-nodal-cuspidal-marks}
有限な共通切断で印の接次元が2であり、その局所環が通常の節点または通常の尖点を持つ
平面曲線上の非零因子による商なら、長さ3の階数一切断では常に偶側の反対対称性が成立する。
反対階数8の有限窓では、長さ6の局所因子が一つでもあれば全次数の正則性が有限に決まる。
従って残余の有限窓の閉点数は高々二つである。
\end{theorem}
\begin{proof}[曲線の導手による印での評価]
印の局所環を $P=C/(g)$ とする。完備な節点は $C=k[[u,v]]/(uv)$ で、
正規化の二枝での $g$ の位数を $n_j$、各枝の極大イデアルの最小位数を $d_j=1$ とする。
尖点は $C=k[[t^2,t^3]]$、$n=\operatorname{ord}_t g$、$d=2$ とする。
いずれも正規化で各位数が $d_j$ 以上の元は $C$ に入る。
従ってある持上げの各位数が $n_j+d_j$ 以上なら、それは $gC$ に入る。
位置1,3の $l,r$ の各列は $P$ の極大イデアルを生成する。$g$ は極大イデアルの二乗に入り、
各枝で両列の最小位数がちょうど $d_j$ である。

$Z$ を印の二行列 $A,B$ の規格化した差とする。各枝で $Z$ の最小位数を $s$、一様化変数を $t$ と書く。
$s<n_j-d_j$ なら、法 $t^{n_j-s-d_j}$ で $A,Z/t^s$ は独立な二元であり、
$l/t^{d_j},r/t^{d_j}$ の剰余で非零な二方向を消す。共通の独立性は $Z$ を固定超平面に正規化したことによる。
従ってこれらは相対割線の基底である。Gorenstein三重点の模型では三方向のスキームで消える二次式 $Q$ に対し、
$Q(l/t^{d_j})=0\bmod t^{n_j-s-d_j}$ である。ゆえに $Q(l)Z$ の各位数は
$2d_j+s+(n_j-s-d_j)=n_j+d_j$ 以上となる。
$s\ge n_j-d_j$ なら同じ位数は二次式の次数だけで足りる。導手により元の $P$ 上で $Q(l)Z=0$ である。
定理\ref{thm:finite-marked-artin-secants}のイデアルの列と完全跡の議論から対称性が従う。

非Gorenstein三重点では定理\ref{thm:finite-marked-planar-triple}の例外平面の式を $T,l_0$ と書く。
上の相対割線では $T(A)$ が零で、正規化した左方向の $l_0$ も零である。
第二辺関係の評価 $f$ は $Z$ の成分の局所線型結合であるから、
$T f,l_0 f$ の枝での位数は少なくとも $n_j-d_j$ である。$s\ge n_j-d_j$ でも同じである。
従って $\mathfrak m^2 T f=\mathfrak m^2 l_0 f=0$ が導手で従う。印以外でこれらは零であり、
各々の支持環はGorensteinで極大イデアルの二乗が零なので長さ2以下である。
隣の一次座標の全射と跡直交で両方とも零になる。よって既出の長さ6の完全交叉 $Y$ 上の誤差議論が適用できる。

そこで得る誤差イデアルの一生成元を再び $f$ とする。
印で $f\ne0$ なら $P/\operatorname{ann}f$ は曲線状であることを一般に確かめる。
$Z=fV'$ とし $V'$ の剰余は $A$ の剰余 $M$ と独立である。
この商で $V'$ は $l,r$ の両方を消す。もし接次元2なら、両者の一次方向はそれぞれ
$M$ の左・右核の二平面全体を張る。$V'$ の剰余もこれら両平面を消し、階数1の $M$ の倍となって矛盾する。
よってこの商で一様化変数を割り、一つだけ局所長さを落として相対割線を得る。
$Y$ の長さ5以上での余因子の消滅は既出の証明と同じ矛盾を与える。
\end{proof}

\begin{proof}[長さ六の平面局所環]
前定理で扱われない長さ6の平面Gorenstein局所環は Hilbert 関数 $(1,2,2,1)$ を持つ。
実際二次部分が三次元ならそれがsocleとなってしまう。
有限辺は六次元滑らかな多様体での六双一次式の完全交叉なので、
完備局所環で独立な一次式四つを消せば $P=k[[x,y]]/(F,G)$ であり、両式の位数は2以上である。
二次初期式は一次元を張る。$F$ に非零の二次初期式を持たせ、$G$ の位数は3以上とする。
$F$ の二次が二つの異なる方向を持てば通常の節点である（異なる一次因子を逐次持ち上げる）。
二次が $y^2$ なら、$F$ に $G$ の適当な定数倍を足して通常の尖点にできる。
実際両式の $x^3$ 係数が零と仮定する。Weierstrassの準備と除法で
\[
                F=y^2+a(x)y+b(x),\qquad G=c(x)+d(x)y
\]
とでき、各位数は $2,4,4,2$ 以上である。長さは自由な二次の $k[[x]]$ 加群での
$G$ の乗法行列の行列式 $c^2-acd+bd^2$ の位数なので8以上となり矛盾する。
非零の $x^3$ 項があれば平方完成と単元の三乗根で通常の尖点となる。
これで長さ3の切断は最初の証明で尽きた。

長さ4の場合は定理\ref{thm:finite-stretched-planar-marks}の二次小行列式の議論を使う。
$n=3$ または局所 $(1,2,1)$ 型の $n=4$ では、非零であるはずの隣接語 $f=H\chi(a)$ は
印以外で $(\eta)$ に入り、印で $\mathfrak m^2 f\subset(\eta)$ である。
印が $Z$ の台から外れれば印の寄与は高々2であり、先の証明の通り矛盾する。
印が台に入るとき $n=4$ のこの場合は $Z$ が印に集中し、$f$ は非単元なので
支持長は6未満である。$n=3$ では $\operatorname{length}(f)\le3+3=6$ であり、
等号ならその支持は $Z$ を含む。長さ6の隣接語の偶側の包は平面なので、
$Z$ を張る平面の式 $\chi(a)$ は $f$ を消す。この式が印で一次の非零方向を持てば
$f$ の局所支持環は曲線状であり、$\mathfrak m^2 f\subset(\eta)$ と合わせ全長は5以下となる。
そうでなければ $\chi(a)\in\mathfrak m^2$ で $f$ の局所寄与は2以下である
（$\mathfrak m^4=0$ ゆえ $f$ の局所支持環の極大イデアルの二乗は零）。
いずれも矛盾する。最後に $n=4$ で曲線状因子の積の場合、四つの独立な偶三次式
$\operatorname{adj}_B(a)$ は印以外で零であり、印では乗法階数1より $\mathfrak m^2$ に入る。
その長さはちょうど3なので矛盾する。共線因子の排除と短い切断の場合は既出の通りである。
長さ20で三局所因子があれば一つは長さ6以下であるから最後の主張も従う。
\end{proof}

\begin{corollary}\label{cor:finite-no-four-on-line}
四辺が有限な標準有限Frobenius窓の各位置の点スキームは、
どの射影直線とも長さ3以下で交わる。反対階数の制限はない。
\end{corollary}
\begin{proof}
直線上に長さ4の因子があれば、その前または後の座標の包を $\mathbb P^{a-1}$ とする。
隣接辺の六式をこの包と直線の積に制限し、一般の $a$ 本を取れば
有限基点なので固有交叉となり、その長さは $a$ である。従って $a\ge4$ であり、両隣の座標は
因子の全長さ4の環をそれぞれ張る。定理\ref{thm:finite-artin-kernel-symmetry}の巡回行列の証明によって
純関係ができ、有限性に反する。
\end{proof}

% !TeX root = ../papers/nc-p3.tex
\section{平面に分解する分離三次式による有限切断の識別}
\label{sec:finite-split-separator-slices}
\begin{theorem}\label{thm:finite-split-separator-slices}
有限窓で全中央支持が4とする。ある局所点の分離三次式に二重平面因子があれば、
その反対切断は三次元の交代行列空間に一つの対称な階数一行列を加えた空間に左右同値である。
三次式は異なる二平面の積 $h s^2$、階数一切断は平方零の長さ3で、唯一の階数一点で $h\ne0$ である。
三重平面は生じない。分離三次式が三つの異なる平面に分かれるなら、その切断の階数一軌跡は
正次元または長さ2以下である。
\end{theorem}
\begin{proof}[二重因子]
写像 $E\to k^3$ の非零の左右評価はどれも階数2以上である。
二重因子の平面をなす三空間の行列 $M(x)$ は一般階数2である。
一般階数1以下なら共通左核または共通右核があり、この評価下界に反する。
同じ下界から一般階数2の $M$ にも定数左右核はない。
$\operatorname{adj}M$ は二次の階数1行列なので原始な列と行に因数分解し、
定数核がないため両ベクトルはともに一次である。
右核が三つの独立な一次式なら、核座標を $x$ として $M(x)x=0$ から
$M(x)=D[x]_\times$ である（各行の係数が交代）。$D$ は定数左核の不在から可逆。
交代空間に第四の行列 $J$ を加えて $\det([x]_\times+sJ)$ を取ると $s$ の一次係数は
$x^t Jx$ である。これが零なら $J$ も交代で四次元性に反する。左側も同様である。

従って両核はそれぞれ二つの独立な一次式を用いる。右核を $(u_2,-u_1,0)$ とすれば
最初の二列は $u_1,u_2$ と共通定数列の積である。その列を第一基底に取ると
下二行の最終列の二式 $v_1,v_2$ が独立である。二重因子条件は
$(v_2,-v_1)K(u_2,-u_1)^t=0$ であり、$K$ は第四行列の下二行・初め二列のブロックである。
$K$ は可逆である。特異なら下二行のある結合は全ての $E$ 上で最終列だけに値を持ち評価下界に反する。
もし二つの一次式空間が異なれば各基底変換で $(u_1,u_2)=(x,y),(v_1,v_2)=(x,z)$
にでき、四積 $x^2,xy,xz,yz$ の独立性から $K=0$ となる。よって同じ二空間である。
行列の基底を揃えると三空間は
$\bigl(\begin{smallmatrix}a&b^t\\ b&0\end{smallmatrix}\bigr)$ となる。残る一成分 $a$ が
$b$ の二座標と独立なのは三空間の単射性による。第四行列の外側成分を行・列加算で消せば
定理\ref{thm:finite-marked-planar-triple}の $E_D$ 型で $D$ は可逆である。
行列式 $a s^2\det D-s b^t\operatorname{adj}(D)b$ が $s^2$ で割れるので $D$ は交代。
第一行または第一列の符号変更で表示した交代空間と階数一行列の和になる。
階数一点とそのスキームも同定理の計算通りである。

\end{proof}
\begin{proof}[三つの異なる平面]
有限な長さ3以上の切断の全模型でこの分解が起きないことを確認する。
非Gorenstein三重点の場合、上の $a s\det D-b^t\operatorname{adj}(D)b$ は後項が非零なら既約二次式である
（$a$ について一次で、係数と定数項は互いに素）。後項が零なら二重因子の場合である。
Gorenstein三重点の三空間は長さ3環 $B$ の全乗法行列で、$E=m_B+kJ$ である。
行列式が一次因子に分かれるなら、単位行列を使った特性多項式の根は
$\chi_i(b)+s\lambda_i$ である。実際 $s=0$ で乗法行列の三指標に制限され、
各一次因子の $s$ 係数は定数である。
$B=k^3$ なら対角部分を除いた $J$ を選ぶ。各 $\lambda_i=0$ は行列式の $s$ 一次項より従う。
二次・三次項から $J_{ij}J_{ji}=0$、$\det J=0$ である。
有向グラフには二閉路がなく三閉路もない（三閉路の逆の積は零）。
同時に上三角化でき、共通不変直線が評価下界に反する。
$B=k[t]/t^3$ では $N=E_{12}+E_{23}$ として、$N,N^2$ の方向はどの固有値も変えない。
平方・三乗の跡の係数から $\operatorname{tr}(N^2J)=\operatorname{tr}(NJ)=
\operatorname{tr}(N^2J^2)=0$。従って $J_{31}=0,J_{21}=-J_{32},J_{21}J_{32}=0$ で共通不変直線がある。
$B=k[t]/t^2\times k$ では同様に $E_{12}$ を使い、$J_{21}=J_{23}J_{31}=0$ なので
共通不変直線または双対の共通不変直線がある。

長さ4の跡模型では、三行の行列式は非零定数倍を除き
$\tau(\det(m_a)a^{-1})$ である。実際跡超平面への形式の制限の行列式は、
周囲の行列式とその逆形式による法線 $\tau$ の自己対合との積であり、単元開集合上でこの式になる。
曲線状の長さ $d_i$ の因子では、剰余座標を $s_i$ としてノルムは $\prod s_i^{d_i}$。
$d_i\ge2$ の因子があればその最高冪の係数座標についてこの三次式は一次で、係数は
非零定数倍の $s_i^{d_i-2}\prod_{j\ne i}s_j^{d_j}$ である。
定数項はこの係数のどの一次因子でも割れない。実際 $s_j=0$ で第 $j$ 因子の
逆数分子の最高冪成分が残り、その最初の接座標の $d_j-1$ 乗の係数は完全跡のため非零である。
単純因子ではこの冪は0乗である。$j=i$ を調べる必要のある $d_i>2$ でも同じである。
ゆえに三次式は既約である。四つの単純因子の場合は
$\sum c_i\prod_{j\ne i}s_j$ ($c_i\ne0$) を $s_1$ の一次式として同じ互いに素の判定を使う。
局所 $(1,2,1)$ の場合は $s$ と、$sw$ 項及び二つの接変数の非退化二次形式を持つ
非退化二次式との積である（$w$ はsocle係数）。以上で三つの異なる一次因子を排除した。
\end{proof}

\begin{corollary}\label{cor:single-support-quartic-slice}
有限な共通切断が長さ4で、点スキームの台が一点なら、任意の接次元の印で
行・列側の反対対称性が成立する。反対階数8で台が二点以下の残余候補でも、
この切断からの非対称な外積の支持環が曲線状の長さ4である場合は排除される。
\end{corollary}
\begin{proof}
接次元2だけを扱えばよい。定理\ref{thm:finite-artin-kernel-symmetry}の外積の
支持環 $S$ と三次持上げを使う。$S$ の一つの局所socleに対応する三次分離式は
曲線状因子の積なら $H\prod_i l_i^{d_i}/l_j$ である。ここで $d_i$ は各局所長、
$l_i$ は各方向での $L$ の剰余一次式、$j$ はそのsocleの因子である。
これは $C=H c N L^{-1}$ を第 $j$ の指標で読むことで従う。
socle元を $\eta$ 倍しても元の点環の非零socle元になるので、
定理\ref{thm:finite-cubic-separators}の当該点の行列式とこの三次式は定数倍で一致する。
局所 $(1,2,1)$ の $S$ では同様に分離式は $H s^2$ となる。

台が一点で曲線状 $S$ の長さが4なら分離式は三重平面となり不可能である。
残る長さ3または局所 $(1,2,1)$ 型では $H$ は一次で、分離式は $H s^2$ である。
上の定理から $H$ は印で非零なので点環全体で単元である。
しかし $H$ と三次持上げの式から、逆数分子の非零な二次座標（例えば $s^2$）は
$Z=\operatorname{Spec}S$ に支持される。隣の偶座標は $S$ に全射するのでその平方も全射し、
平方どうしの完全跡直交によりこの二次座標は零になる。これは二次式の単射性に反する。
全射の場合および共線因子の排除は既出の通りで、奇側の対称性が従う。

台が二点で $S$ が曲線状長さ4かつ局所なら同じ三重平面の矛盾となる。
二因子なら $Z$ は印を含む。$H$ は定数であり、印でのsocle分離式は
印で消える三つの一次式の積なので射影三空間で印における重複度が3となる。
一方、階数一点での行列式曲面の重複度は正確に2である。
その点を $E_{11}$ としたとき二次初期式は右下二行二列の行列式で、恒等零なら
そのブロックの線型空間は共通左核か共通右核を持つ。従って全四空間の一側評価が
高々一次元となり、純関係の不在に反する。これで二因子の場合も除かれる。
\end{proof}

% !TeX root = ../papers/nc-p3.tex
\section{長さ七以下の平面切断と二つの局所点の接平面}
\label{sec:finite-plane-sections-and-tangent-incidence}
\begin{theorem}\label{thm:finite-plane-section-seven}
四辺が有限な標準有限Frobenius窓の各座標において、点スキームの任意の平面切断は
長さ7以下である。長さ7なら、その平面と反対辺でない隣の三空間との積で、
一般の五関係が切る長さ10の完全交叉における残余は、平面へ共線な長さ3のスキームとして埋め込まれる。
\end{theorem}
\begin{proof}
隣接する辺を $\mathbb P^3\times\mathbb P^3$ の六双一次式とし、第一位置の平面 $H$ に制限する。
六式は独立なままである。依存すれば平面の式と第二位置の一次式との純関係が生じる。
基点は有限なので、一般の五式は $H\times\mathbb P^3$ に長さ $\binom52=10$ の完全交叉 $Y$ を与える。
$Y$ の $H$ への射影は閉埋込みである。全六式の評価が平面上の点で階数4なら五式でも少なくとも3、
階数3の有限個の点では一般の五式を選んでも階数3である。従って $Y$ はその退化スキーム上の核直線のグラフである。

残った一式の値 $f$ は $Y$ 上で零ではない。実際 $Y$ 上で消える双一次式は定義五式で尽きる。
Koszul分解を $(1,1)$ でひねると他の負の項は寄与しない（最終項の非零コホモロジーも次数5である）。
随伴により $\omega_Y=\mathcal O_Y(2,1)$ で、その完全跡は周囲から制限した $(2,1)$ 式を消す。
これはKoszul分解で単位写像に双対な最終接続写像がこの跡となるためである。
従って $f$ と $H$ の各一次式との積の跡は零。$Z=\operatorname{Spec}(\mathcal O_Y/\operatorname{ann} f)$ の
完全跡はこの $f$ 倍であり、平面の一次式は $Z$ に全射できない。平面射影が閉埋込みなので
$\operatorname{length}Z\ge3$ で、等号なら共線である。
元の切断長は $10-\operatorname{length} Z$ だから主張を得る。
\end{proof}

\begin{proposition}\label{prop:finite-two-planar-tangent-incidence}
反対階数8の有限窓で全中央支持4とする。二つの異なる平面局所点 $p,q$ について
$q_i\in T_{p_i}X_i$ は反対位置 $i+2$ での同じ包含と同値である。
一つの位置で $q_i\in T_{p_i}X_i$ と $p_i\in T_{q_i}X_i$ は両立しない。
また印 $p$ の両側のいずれかの分離三次式が $h s^2$ のとき、$X$ の台は一点ではあり得ない。
台が $p,q$ の二点で $q$ の局所長が8以上かつ $p$ の極大イデアルの二乗が非単項なら、四位置の全てで $q_i\in T_{p_i}X_i$ である。
\end{proposition}
\begin{proof}
偶側で $M_p=w(p_0,p_2)$ は階数1で、左右の因子は奇位置の接平面の法線である。
実際、印からの三座標に対して $M_p$ の左右核はその一次接方向を含み、各々二次元なので一致する。
反対切断の八空間の中で、$q_1$ による縮約が零の空間は次元4である。
それは $q$ の偶評価の共通四空間であり、$q_3$ での縮約が零である空間にも等しい。
従って $M_p$ の一方の法線が $q$ を消すことは他方でも同じである。添字を回せば偶位置側も従う。
両向きに包含すれば二点を結ぶ直線が両端で長さ2以上の切断を持つので、
系\ref{cor:finite-no-four-on-line}に反する。

二重因子の印付き切断を奇座標を変数として読む。定理\ref{thm:finite-split-separator-slices}により
$A=h vv^t+K$、$K$ は交代、$h(p)\ne0$ である。隣の偶位置の射影座標 $l$ に対して
$l^tA=0$ より $h(v^tl)^2=0$。$h$ が全ての局所因子で単元なら偶側の二次制限の単射性に反する。
従って一点支持は不可能で、二点支持なら $h(q)=0$。
$q$ の長さが8以上なら平面 $h=0$ にこの局所スキーム全体は入らない。
よって $v^tl$ は $q$ で非単元である。これは $p$ の偶接平面の式であり、反対位置でも同じである。
また分離式は $q$ の局所環で零、$h$ はそこで恒等零ではないので $s(q)=0$。
この $s$ の平面が奇変数での $p$ の接平面になることを、$p$ の局所環が非伸長型の場合に確かめる。
$E_D$ 表示の右下ブロックは $sD$、$a=h$ は単元、他の二座標を $b$ とすれば
局所式で $sD r'=a^{-1} b b^t r'$。$r'$ は二つの独立な一次方向を持つので
$s\mathfrak m\subset\mathfrak m^3$。もし $s$ の一次部が非零であれば
$\dim\mathfrak m^2/\mathfrak m^3\le1$ となる。従って非伸長型では $s\in\mathfrak m^2$ で所望の接平面である。
\end{proof}

\begin{corollary}\label{cor:two-support-quartic-discrepancy}
未認証の反対階数8の二点支持窓で、一つの印の階数一切断が長さ4とする。
核側が非対称のままなら、その外積の支持 $Z$ はその印だけに集中し、長さ3の曲線状環または
長さ4の局所 $(1,2,1)$ 環である。後者では $\Sigma$ も同じ一点に支持される。
どちらの場合もその印に重平面因子を持つ分離式が別の切断から得られる。
\end{corollary}
\begin{proof}
共線四点因子の排除と前節の結果から、残る支持環は上の二つまたは局所長 $2,1$ に分かれる長さ3である。
後者では単純因子での三次分離式は $H l_2^2$ である（$l_2$ は二重因子の方向）。
その局所点では $H$ が単元、$l_2$ は零であり、単純因子の方向 $l_1$ もその点で消える。
これは分離行列の階数1と外積支持の $Al=0$ による。従って $H,l_2,l_1$ は互いに比例せず、
二重因子の局所点の分離式 $H l_2 l_1$ は相異なる三平面の積になる。
定理\ref{thm:finite-split-separator-slices}からその印の切断は長さ2以下または正次元となり、既出の平面点判定で認証または排除できる。
最後に $Z$ が印を含まなければ、定理\ref{thm:finite-stretched-planar-marks}の
$\mathfrak m^2 f\subset(\eta)$ は印で $\mathfrak m^2 f=0$ を与える。
この包含自体の証明は伸長型を仮定しない。局所支持長は高々2なので、同定理の長さ下界との矛盾がそのまま成り立つ。
よって残る $Z$ は印そのものに支持され、分離式は $H s^2$ である。
\end{proof}

% !TeX root = ../papers/nc-p3.tex
\section{平面局所環の定義式が四次以上のときの共通八点切断}
\label{sec:finite-high-order-common-octad}
\begin{theorem}\label{thm:finite-high-order-common-octad}
反対階数8で四辺が有限な候補のうち、既出の対称性または三文字の判定に入らない場合を考える。
平面局所点で $\dim\mathfrak m^j/\mathfrak m^{j+1}=j+1$ が $j=0,1,2,3$ に成り立てば、
その点は唯一の台である。各座標の接平面の式を $n_i$ とし、局所的に自明化する。
全ての $\operatorname{ann}(n_i^2)$ が同じイデアルとなり、これが定める長さ8の部分スキームは
各位置で三つの二次式の完全交叉である。この部分上で四つの座標は射影線型に同定される。
隣接する二つの接平面の積による辺の切断は長さ6、各座標の接平面切断は長さ7である。
\end{theorem}
\begin{proof}[二つの接平面と四つの関係]
平面局所環は二式完全交叉で、両式の位数が4以上なので局所交点数は16以上である。
他に未処理の局所因子があれば長さ7以上で全長20に反する。従って台は一点 $p$ である。
各側の共通四空間 $E$ の印での行列を $M$ とし、$M$ が第一対角成分だけを持つように選ぶ。
右下二行二列への射影を $C$ とする。印の二つの接座標を $x$、三座標の右・左核の一次部の
下二成分を $Rx,Lx$ とすると、どちらも二次元の同型である。
$A r=0$ および $l^t A=0$ から、$A$ の右下ブロックの一次部 $C_1(x)$ は
$C_1(x)Rx=(Lx)^tC_1(x)=0$（両式それぞれ零）を二次の多項式として満たす。
ここでは上・下の一次オフ対角成分と核の第一座標との積は三次以上で、二次の定義式がないことを使った。
非零の階数1以下の一次鉛筆なら共通像か共通余像を持つので、両核がともに動くことに反する。
従って $C_1=0$。$A$ が閉埋込みを与えるから、$M$ での階数一切断の接次元は少なくとも2である。
長さ4または正次元の切断は一点支持では既出の判定に入る。
従って両側とも接次元2の長さ3切断であり、定理\ref{thm:finite-marked-planar-triple}の $E_D$ 模型となる。
その $t$ 座標は非零の一次式で印の二次近傍の極大イデアルの二乗に入るので、その位置の接平面の式である。
行・列の先頭の射影座標も $M$ の核により対応する位置の接平面の式となる。

辺の二位置を $[s:x:n],[t:y:m]$ と書く。$x,y$ は各々二文字で、$n,m$ は接平面の式である。
両側の $E_D$ の行または列の評価から来る三関係ずつを取る。最初の組の下二式は
基底を選べば $xm=ny$ である。実際 $b$ の二文字、相手の下二文字、及び可逆な $D$ による式で、
接平面の式を加える二文字のリフトも選べる。他の組の下二式は同じ二空間を張る。
両組とも印の両方の評価で消える二空間に入り、その次元は正確に2だからである。
後者は辺の六式の一次部分が接次元2より階数4であることから従う。
それぞれの上式では相手側の下二文字と $b$ の二文字が非退化に対合する。
異なるリフトからの $xm,ny,nm$ の項を下二式と $s,t$ への一次式の加算で消せば、
四つの共有部分からの関係は
\[
                 xm=ny,\qquad sm=q(x,y),\qquad nt=r(x,y)
\]
となる。$q,r$ は二行二列の可逆双一次形式である。
二つは独立である。実際 $n=m=0$ の平面の積に制限した六式は、有限交叉のため少なくとも四つの独立式を必要とする。
従ってこの積で $q,r$ と残り二式は完全交叉で、長さは $\binom42=6$ である。

\end{proof}
\begin{proof}[八点の連結と二乗イデアル]
印のアフィン座標 $s=t=1$ で $m=q,n=r$ を消去する。辺の各射影は閉埋込みなので
残り二式の $y$ での微分は可逆である。形式的に $y=Tx+O(x^2)$ を解くと $T$ も可逆となる。
$xq-yr$ の三次初期式が恒等零であることから
\[
                     q(x,Tx)x=r(x,Tx)Tx .
\]
どちらの二次式も零なら、独立な $q,r$ が可逆な $T$ を通してともに唯一の交代形式の倍となり矛盾する。
よって一般の $x$ が $T$ の固有ベクトルであり $T=cI$, $c\ne0$。
従って $q=c r+\kappa[x,y]$、$\kappa\ne0$ である。
射影線型同型 $[s:x:n]\mapsto[s:cx:cn]$ のグラフに全六式を制限すると、
上の四式は $sn=c r(x,x)$ の一式となる。残り二式とともに有限な辺の部分を切るので
三二次の完全交叉、長さ8である。その支持も $p$ だけである。

この部分の局所イデアルを同定する。$y=y(x)$ を解いた二変数冪級数環で $d=y-cx$ と置く。
辺の二式は $Md$ で、$M=\kappa x[x,-]-r I$、$\det M=r^2$ である。
線型同型のグラフを課した部分のイデアルは $(d)$ である。
実際 $d=0$ とすれば $m=c n$ も従う。
二つの正則列 $(Md)\subset(d)$ の商の間で、$\det M$ 倍は後者の商を前者の商に単射し、
その像は $(d)$ の零化イデアルである。これはKoszul分解の比較で最上段が $\det M$ 倍で、
二次の局所双対が商の全射を単射に移すためである。
従って八点部分の定義イデアルは $\operatorname{ann}(r^2)=\operatorname{ann}(n^2)$ となる。
また $g=[x,y]$ は $xm=ny$ より $n,m$ の両方に消され、$m-cn=\kappa g$。
従って $m^2=c^2n^2$ がこの自明化で成り立つ。自明化を変えても二乗の主イデアルは変わらないから、
全辺で同じ八点部分を得る。四つの座標はこの部分の上で各辺の線型同型によって対応する。

最後に $n=0$ の切断を同じ二変数環で読むとそのイデアルは $(r,x_1g,x_2g)$ である。
$(r,g)$ の商は $n=m=0$ の長さ6切断であり、$g$ は一次元環 $k[[x_1,x_2]]/(r)$ の非零因子である。
従って $g$ の像を極大イデアル倍まで残した商の長さはちょうど一つ増えて7である。
\end{proof}

% !TeX root = ../papers/nc-p3.tex
\section{共通二次式の網による反対位置の対称性}
\label{sec:finite-common-net-symmetry}
\begin{lemma}\label{lem:common-net-alternating-part}
四位置を同じ四次元空間 $V$ で表し、窓を四重線型形式 $w$ とする。
四辺の関係式を対角に制限した二次式が全て同じ正則列
$Q_1,Q_2,Q_3$ の線型包に入るとする。
すると $w=G+\omega$ と一意に書ける。ここで $G$ は位置 $0,2$ の交換と
位置 $1,3$ の交換の両方で不変、$\omega\in\bigwedge^4 V^*$ である。
\end{lemma}
\begin{proof}
$V$ が生成するテンソル代数を、対称二次部分で三つの $Q_a$ の全てに
直交する元により割る。商では $x^2=\sum_a Q_a(x)z_a$ と書ける。
$[x^2,x]=0$ を $x$ の成分の三次式として比較すると各 $[z_a,x]=0$ である。
実際、正則列をなす三二次式には一次係数の関係がなく、係数空間の任意の線型汎関数を
適用できる。連続する三辺についての仮定は $w$ がこの商の次数4部分の汎関数に降りることをいう。
よって $P_x(y,z)=w(x,x,y,z)$ と置くと
\[
 w(y,x,x,z)=w(y,z,x,x)=P_x(y,z).
\]
巡回移動した窓にも同じ論法を適用して $w(x,z,y,x)=w(y,x,x,z)=P_x(y,z)$ を得る。
従って位置 $0,2$ を交換した形式と元の形式との差は、三つの連続する組のいずれに
同じベクトルを置いても零であり、完全交代形式である。位置 $1,3$ の交換についても、
四つの等式から同じことが従う。完全交代化を取ればどちらの差も
$2\operatorname{Alt}(w)$ である。$\omega=\operatorname{Alt}(w)$ が求める分解を与える。
\end{proof}

\begin{theorem}\label{thm:high-order-octad-bisymmetry}
有限な反対階数8の窓で、平面局所点の二つの定義式がともに四次以上なら、
既出の三文字または対称性判定に含まれる。
具体的に、定理\ref{thm:finite-high-order-common-octad}の共通八点切断を持つ未判定候補には
両対称性が強制され、そこで必要な非対称模型 $E_D$ に反する。
\end{theorem}
\begin{proof}
各辺の線型同定の一周の合成は八点切断を各点上スキームとして固定する。
これを $B$ と書くと $u\wedge Bu$ の成分は八点上消える二次式なので三二次式の網に入る。
$u\wedge(u\wedge Bu)=0$ と網に一次関係がないことから、$u\wedge Bu=0$ が恒等的に成り立つ。
実際 $u\wedge Bu=\sum Q_a(u)K_a$ と書けば各定数二ベクトルに $u\wedge K_a=0$ がすべての $u$ で成立する。
従って $B$ は射影的に恒等であり、四位置を整合的に同定できる。全辺の対角二次式は同じ八点を含み、
補題\ref{lem:common-net-alternating-part}を適用できる。

印を $p$ とし、位置0で $p$、位置2で $u$ を代入した位置1,3の行列を $A(u)$ とする。
この四空間は左右とも $p$ を核に含むので、$V/kp$ の三行三列の行列とみなす。
辺12の前定理の四式を対角に制限するとただ一つの二次式
$Q_{12}=sn-c r(x,x)$ を得る。$r(x,x)$ は零でない（前定理の三次恒等式と $q,r$ の独立性）。
よって $Q_{12}$ は既約で $p$ はその滑らかな点である。
四式には $A(u)$ の左核を与える三式が含まれるので、
$Q_{12}$ 上では $u$ が左核に入る。同じく辺23の二次式 $Q_{23}$ 上では $u$ が右核に入る。
いずれも $\det A(u)$ を割る。$E_D$ の式によりこの三次式には $at^2\det D$ の項があり非零である。
従って二つの既約二次因子は比例し、同じ曲面 $Q$ 上で両方の核条件が成立する。

$A(u)$ の対称部分が張る空間の次元は3以上である。実際 $A(p)$ は両隣の接平面の式の積である。
前定理の各同定が接平面を接平面に移すため、これは同じ一次式の平方の非零倍であり、
対称部分は $p$ で非零である。$Q$ は全空間を張るから、その一般点でも非零である。
$p$ からの $Q$ の射影は $\mathbf P(V/kp)$ に稠密である（一般直線は $p$ の他にもう一度交わる）。
従ってこれらの非零対称行列の射影核が平面で稠密になる。
対称行列の一本の鉛筆では不可能である。階数2の元の射影核は各一点であり、階数1の元は有限個しかない。
後者は対称階数1のVeronese曲面が直線を含まないためで、一次元以下の線型空間の場合も同じ結論となる。

反対切断を対称二乗と外積二乗に分けると、$G$ と $\omega$ は独立な対角ブロックを与える。
前段から $G$ の階数は3以上であり、非零の四形式は外積二乗に階数6の完全対合を与える。
全階数8より $\omega=0$ である。特に全ての $A(u)$ は対称になる。ところが $E_D$ の全行列を左右基底変更で
対称にできるなら、可逆な一つの行列 $J$ で $E_D J$ を対称にできる。
$E_{11}J$ の対称性より $J$ の先頭行は $(a,0,0)$ であり、$a\ne0$。
二つの $b$ 方向の対称性より $J$ の右下ブロックは $aI$、
$t$ 方向の対称性より $aD$ も対称となり、模型の $D\ne D^t$ に矛盾する。
従って未判定の場合自体が起こらない。
\end{proof}

% !TeX root = ../papers/nc-p3.tex
\section{平面曲線の一回のブローアップと有限な反対階数八の全次数判定}
\label{sec:finite-rank-eight-all-marks}
\begin{lemma}\label{lem:finite-carrier-blowup}
有限窓の共通印付き四空間の階数一切断が長さ3とする。
印の平面局所環を $P=C/(g)$、$C$ を重複度 $d$ の平面超曲面完備局所環、
$g$ を極大イデアルの二乗に入る非零因子とする。
$d=2$ なら対称性判定が成立する。切断が非Gorensteinなら $d=3$ でも成立する。
曲線の被約性は仮定しない。
\end{lemma}
\begin{proof}
$C$ の一般の一次パラメータを $x$、もう一つの座標を $y$ とする。
Weierstrass表示を $y^d+\sum_{j=1}^d a_j(x)y^{d-j}=0$ とすれば
$a_j\in(x^j)$ である（準備で除いた因子は単元なので位数を保つ）。
$C'=C[y/x]\subset C[1/x]$ は $1,z,\ldots,z^{d-1}$, $z=y/x$ を
$k[[x]]$ 基底に持ち、$x^{d-1}C'\subset C$、$\mathfrak m_C C'=x C'$。
実際 $x^d$ で割ったモニック方程式を使え、線型独立性は $x$ の逆元を加えて従う。
$h=g/x\in xC'$ で $C'/(h)$ はArtin環である。

共通行列を $A$、規格化したもう一つの行列との差を $Z$ とし、行・列核の
三座標を $l,r$ とする。全てを $C$ に持ち上げ、差は固定三空間 $N$ 内で取る。
$l,r$ の成分はどちらも $\mathfrak m_C$ を生成する。実際 $P$ で極大イデアルを生成し、
$g\in\mathfrak m_C^2$ なので一次部が基底を含む。
従って $L=l/x$, $R=r/x$ は $C'$ 上で各々単位イデアルを生成する。
$A,Z$ は $h$ を法として $L,R$ を共通核に持つ。これは元の縮約が $g$ の倍であるためである。

$C'/(h)$ の各局所因子で、剰余体でのこの二核を消す $E$ の空間は次元2以下である。
実際三次元なら二行二列行列の射影平面に階数1の曲線ができる。
$A$ は固定超平面 $N$ の外の単元座標を持つので、$N$ 上の縮約核は剰余体で高々一次元。
二つの主元で消去すれば、差は零、または $Z=eV$ で $V$ は剰余体で $A$ と独立である。
後者で $\operatorname{ann}e$ を法とすれば $A,V$ は共通核を持つ相対割線を与える。
この商でも $L,R$ は射影座標となり、剰余体で階数1集合は有限なので
行列式が割線上に切る相対次数2の部分は平坦である。

Gorenstein三重点では定理\ref{thm:finite-marked-artin-secants}の三方向上で
消える任意の二次式 $Q$ が $Q(L)Z\in h C'$ を成分ごとに満たす。
従って $Q(l)Z\in xg C'$ であり、$d=2$ なら $Q(l)Z=0$ が $P$ 上で成り立つ。
印の外での同じ消滅と三方向のイデアル列を用いる同定理の証明で対称性が従う。

非Gorenstein三重点では定理\ref{thm:finite-marked-planar-triple}の $E_D$ 座標を用いる。
$L=(L_0,L')$ とすると、相対割線の左核には $L_0=0$ が必要なので $L_0 Z\in h C'$。
$A=\bigl(\begin{smallmatrix}a&b^t\\b&tD\end{smallmatrix}\bigr)$ では
$a$ が単元、$b$ と $Z$ の全成分は印で消え $x C'$ に入る。
各局所因子で $D^t L'$ の少なくとも一成分は単元である。
実際、$A$ の剰余は $E_{11}$ の非零倍であり、$L$ の剰余は非零でこれを消す。
左核の下二式 $tD^t L'+b L_0\in h C'$ に $Z$ の各成分を掛けると
$b L_0 Z\in xh C'$ かつ $hZ\in xh C'$ である。可逆な成分で消去して
\[
                 t Z\in g C',\qquad l_0 Z=x L_0 Z\in g C'
\]
を得る。ここで単元での消去は $C'$ の各極大イデアルで局所化して行える。
$d\le3$ なら $x^2 C'\subset C$ により $\mathfrak m_P^2 t Z=\mathfrak m_P^2 l_0 Z=0$。
第二辺の評価 $f$ は局所的に $Z$ の成分の線型結合であり、印の外では既出の相対割線の
議論により $tf,l_0f$ は零である。印で極大イデアルの二乗が消す各元の支持環は
Gorensteinで長さ2以下、対応する隣の一次座標が全射する。
隣接座標積の完全跡直交から $tf=l_0 f=0$ が全域で成立する。
以後は長さ6の二平面切断の証明を適用できる。
すなわち全評価の空間は高々一次元で、非零ならその一生成元の支持は二平面切断全体である。
定理\ref{thm:finite-nodal-cuspidal-marks}の証明の通り、平面印におけるこの生成元の
支持環は曲線状であり（この段に曲線 $C$ の条件は要らない）、
長さ5以上で消える余因子と長さ6の二次補間が矛盾する。
\end{proof}

\begin{theorem}\label{thm:finite-rank-eight-all-marks}
四辺が有限な反対階数8の標準有限Frobenius窓は、全て次の有限判定に入る。
中央支持3の点があれば全次数で正則である。それ以外では少なくとも一方の反対対称性が成立する。
その対称核が階数2以下の非零元を含めば全次数で正則であり、残りでは
両対称性と定理\ref{thm:bisymmetric-two-length-test}の条件が正則性に必要十分である。
\end{theorem}
\begin{proof}
中央支持3の認証は定理\ref{thm:finite-cubic-separators}による。
曲線状の印・伸長型の印・短い局所環の場合を除けば、残る印は全て平面局所完全交叉で
閉点数は高々2である。両側の共通四空間は階数8で自動的に存在する。
未判定なら定理\ref{thm:high-order-octad-bisymmetry}により、印における
二定義式の最小位数 $d$ は2または3である。位数 $d$ の式で曲線 $C$ を取り、
他方をその上の非零因子 $g$ とする。

一つの印を取る。いずれかの側の階数一切断が正次元または長さ2以下なら既出の対称性を使う。
長さ4では一点支持なら系\ref{cor:single-support-quartic-slice}を使い、
二点支持で核側が非対称なら系\ref{cor:two-support-quartic-discrepancy}が同じ印に
重平面因子を持つ別の分離式を与える。その切断は定理\ref{thm:finite-split-separator-slices}より
非Gorensteinの長さ3である。従って残る場合にはこの印で長さ3の切断を使える。
$d=2$ なら前補題を適用する。$d=3$ では局所定義式に二次の初期式がない。
定理\ref{thm:finite-high-order-common-octad}の最初のブロック計算
$C_1(x)Rx=(Lx)^t C_1(x)=0$ はまさにこの次数までしか使わず、
階数一切断の印での接次元が2以上となる。長さ3なら平方零の非Gorenstein局所環なので、
再び前補題が適用できる。これで少なくとも片側が対称になり、
系\ref{cor:all-defective-one-sym-finite}が述べた判定を与える。
\end{proof}

% !TeX root = ../papers/nc-p3.tex
\section{片側だけの共通切断と有限な反対階数九の判定}
\label{sec:finite-rank-nine-all-marks}
\begin{lemma}\label{lem:quartic-two-marks-incidence}
全中央支持4の有限窓が、二つの非伸長型平面局所点 $p,q$ に支持されるとする。
各点で少なくとも一方に共通四空間があり、その階数一切断が長さ4であるなら、
少なくとも片側の反対対称性が成立する。二点が選ぶ側は異なってよい。
\end{lemma}
\begin{proof}
対称性がないとし、まず $p$ の共通切断の核側の一座標で三次持上げを取る。
定理\ref{thm:finite-artin-kernel-symmetry}の外積の支持 $Z$ は長さ3または4である。
共線の場合は一つの巡回行列による純関係で排除される。
長さ4の曲線状因子の積は系\ref{cor:single-support-quartic-slice}の三重平面または
印で重複度3の行列式による矛盾で除かれる。この計算には反対切断階数も他側の一致も使わない。
残るのは曲線状因子の積の長さ3または局所 $(1,2,1)$ 型である。
$Z$ は $p$ を含む。実際、含まなければ定理\ref{thm:finite-stretched-planar-marks}の
小行列式による $\mathfrak m_p^2 f\subset(\eta)$ の評価で
隣接語 $f=H\chi(a)$ の印での支持長は高々2、外では $(\eta)$ に入る。
長さ3なら全長6の下界に反し、長さ4で全長6に達すれば $f$ の支持が
偶座標で全空間を張る $Z=\Sigma$ を含み、長さ6の平面条件に反する。

三次持上げの一次因子 $H$ は $p$ で非零である。実際印でのsocle分離式は
$H$ と印からの二つの射影一次式の積で、後の二式は印で消える。
行列式の印での重複度はちょうど2である（系\ref{cor:single-support-quartic-slice}の二次初期式の計算）。
一方 $H$ は点環全体の単元ではない。単元なら、逆数分子の零でない二次座標が
$(\eta)$ に入り、その支持環 $Z$ には隣の偶座標が全射する。隣り合う座標の二次空間どうしの
完全跡直交からこの二次式は零になり、二次制限の単射性に反する。
よって $H(q)=0$。$Z$ が局所なら印での分離式は $Hs^2$ である。
$Z$ の長さが $2+1$ に分かれれば、単純因子の分離式が $Hs^2$ となる。
定理\ref{thm:finite-split-separator-slices}よりその因子の閉点で $H$ は非零であるから、
単純因子の閉点は必ず $p$ である。従ってこの場合も $p$ の分離式が $Hs^2$ である。
核の二位置を交換した同じ議論を適用できる。これで二つの核位置のそれぞれに同じ印 $p$ の重平面分離式がある。

長さの大きい方を $q$ と取る。その長さは10以上なので、どの位置でも $H$ は
$q$ の局所環で恒等零ではない（定理\ref{thm:finite-plane-section-seven}）。
各重平面分離式が $q$ の局所環で零だから $s(q)=0$ であり、これが当該位置の
$p$ の接平面である。実際重平面模型 $E_D$ の $D$ は交代で、その $s$ 座標が
極大イデアルの二乗に入ることは命題\ref{prop:finite-two-planar-tangent-incidence}の
$s\mathfrak m\subset\mathfrak m^3$ の計算による。隣の二位置でも、同模型
$A=H vv^t+K$（$K$ は交代）の左右核から $H(v^tl)^2=H(v^tr)^2=0$ となる。
左右の基底はこの模型に従って取り、$v^tl,v^tr$ はその二位置の印の接平面の式である。
よって $q_i\in T_{p_i}X_i$ が四位置全てで成立する。この推論には反対階数8での
二位置の接平面入射の同値を使わない。

今度は $q$ の共通切断で同じ証明を行う。$q$ における各重平面分離式の単因子は
$q$ で非零なので、$p$ の接平面とは異なる。従ってこの一次式は $p$ の局所環で
恒等零とはなり得ない。実際、恒等零の平面は二次元接平面を含み、これに等しくなる。
$p$ の長さに関係なく上の核計算から $p_i\in T_{q_i}X_i$ が成立する。
両端で接する直線は切断長4以上となり系\ref{cor:finite-no-four-on-line}に反する。
\end{proof}

\begin{lemma}\label{lem:high-order-point-rank-nine}
有限窓で全中央支持が4、点スキームが一つの平面局所点に支持され、
その二つの定義式がともに四次以上なら、四位置を揃えると共通八点切断があり、
窓は補題\ref{lem:common-net-alternating-part}の分解を持つ。
四形式が非零なら反対階数は少なくとも12である。
\end{lemma}
\begin{proof}
ここでは反対位置での二つの四空間が一致することを仮定しない。
一つの文字を印で固定し、反対位置の変数 $u$ による四空間 $A(u)$ を一つずつ取る。
隣の二座標はそれぞれ $A(u)$ の核を与える。
定理\ref{thm:finite-high-order-common-octad}の一次ブロック計算は
この一行列だけを使い、印での二つの接方向では右下ブロックが零になる。
そのブロックへの射影の核はちょうど三次元である。全四空間がそこにあれば
下の行からの非零評価の階数が1以下になる。ブロック像の生成元 $D$ も同じ理由で可逆である。
印の階数一切断の接次元は少なくとも2で、正次元切断ではあり得ない。
実際定理\ref{thm:finite-marked-ruling-pencil}の母線だけとなる証明の前半は
四空間単独の左右評価階数の下界から従い、母線では全点で接次元が1である。
従って印以外の二つの接行列の上行への射影も左列への射影も独立であり、
さもなければ階数一行列の直線ができる。これで各四空間は個別に $E_D$ のブロック表示を持つ。
ここでは対称な $D$ も許す。

各辺では、残る二位置の一方を印に固定して得る三関係ずつを取る。
各組の下二式は両位置の印での評価が零の関係空間を張る。これは辺の一次部分の階数4から
二次元であり、二組で共有される。従って各辺に定理\ref{thm:finite-high-order-common-octad}の
四式 $xm=ny,\ sm=q(x,y),\ nt=r(x,y)$ がある。ここまでに同じ位置を変数とする
別の四空間との同一視は要らない。その後の三次初期式の比較、局所完全交叉の連結、
$n_i^2$ の零化イデアルによる四辺の同定はこの四式のみを使う。唯一の支持が印なので
八点切断は全辺で一致し、補題\ref{lem:common-net-alternating-part}の分解を得る。

定理\ref{thm:high-order-octad-bisymmetry}の既約二次曲面上の左右核の証明も、
$A(u)$ だけの行列式とその両隣の辺だけを使う。従って $A(u)$ の対称部分は三次元以上を張る。
これらは奇側の $p$ を消す対称形式である。奇偶を交換して同じ結論を得る。
$G$ の反対切断をその階数だけの空間と双対との完全対合に分解すると、
偶側で $p$ と任意の $u$ を入れた像と、奇側で $p$ と任意の $v$ を入れた像は直交する。
実際 $w(p,p,u,v)=0$ は辺の関係であり、四形式もこの評価では零である。
両方がそれぞれ三次元以上なので $G$ の階数は6以上である。
四形式が非零なら独立な外積ブロックの階数6を加えて下界12を得る。
\end{proof}

\begin{theorem}\label{thm:finite-rank-nine-all-marks}
四辺が有限な反対階数9の標準有限Frobenius候補が標準Koszul左右AS正則となる
必要十分条件は、中央支持3の点が存在することである。印の被約性を仮定しない。
\end{theorem}
\begin{proof}
中央支持3があれば定理\ref{thm:finite-cubic-separators}を使う。残りでは全て4である。
各印で、偶二位置を別々に印に固定した二つの四空間の和と、
奇二位置に対する同様の和は、反対切断の完全対合で相互に直交する。
これは四つの隣接点対全てでの消滅である。和の次元の合計は9以下であり、
各々4以上だから少なくとも一方は4、その側に共通切断がある。

対称性がどちらにもないとする。曲線状の印、伸長型の印、長さ6以下の印での
三重点・四点の各証明は一つの共通切断だけを使う。長さ3ではその変数側、
長さ4では核側で対称性を与える（定理\ref{thm:finite-artin-kernel-symmetry}、
\ref{thm:finite-stretched-planar-marks}、\ref{thm:finite-nodal-cuspidal-marks}の各局所計算）。
長さ2以下と正次元の場合も一つの共通切断で足りる。
従って残る点は非伸長型平面局所点で各長さ7以上、台は高々二点となる。
四次以上で始まる局所完全交叉は長さ16以上なので唯一の台となり、
補題\ref{lem:high-order-point-rank-nine}から両対称である。
それ以外の印では局所曲線の最小重複度が2または3である。
共通切断が長さ3なら、重複度3でその切断は二次初期式の不在から平方零型であり、
いずれも補題\ref{lem:finite-carrier-blowup}で対称になる。
よって各点に長さ4の共通切断がなければならない。一点支持なら
系\ref{cor:single-support-quartic-slice}、二点支持なら
補題\ref{lem:quartic-two-marks-incidence}で対称性が従い、矛盾する。

片側対称性を得た後は系\ref{cor:all-defective-one-sym-finite}を適用する。
対称核に階数2以下の非零元があれば中央支持3の点ができるので今の場合にはない。
他の場合に正則なら両対称であるが、定理\ref{thm:all-bisymmetric-as-closure}より両対称な標準正則窓の反対階数は8以下である。
従って中央支持が全て4の候補は除外される。
\end{proof}

\begin{corollary}\label{cor:all-windows-through-eight}
反対階数8以下では、有限Frobenius条件を通る全窓の全次数正則性と、正則な窓どうしの
添字保持同型を有限に判定できる。階数9でこれらの判定の被覆に残る正則候補は、
曲面射影を持たず、純関係と正次元射影ファイバーの双方を持つ窓に限る。
\end{corollary}
\begin{proof}
曲面以上の射影には定理\ref{thm:no-moving-plane-pencils}の被覆を使う。
曲線射影で分岐があれば、階数7以下は定理\ref{thm:low-rank-branch-exclusion}、
階数8は定理\ref{thm:rank-eight-branch-exclusion}で排除される。
階数9で純関係がない場合も第\ref{sec:branch-rank-eleven-floor}節の下界に反する。
点ファイバーで一次元なら系\ref{cor:all-one-dimensional-point-graphs}の全被覆を使う。
四辺が有限なら定理\ref{thm:finite-frobenius-low-opposite-as}、
\ref{thm:finite-rank-eight-all-marks}及び上の定理を用い、同型判定には
定理\ref{thm:finite-all-window-keys}を用いる。
\end{proof}

% !TeX root = ../papers/nc-p3.tex
\section{高重複度の局所点での半回転と全階数での両対称性}
\label{sec:finite-high-order-half-turn}
\begin{theorem}\label{thm:finite-high-order-half-turn}
四辺が有限で全中央支持4の標準有限Frobenius窓に、平面局所完全交叉の
二つの定義式がともに四次以上となる点 $p$ があるとする。
他の局所成分の有無、反対切断の階数、印付き四空間の一致を仮定せず、
窓は両対称である。特に正則なら反対階数は8以下である。
\end{theorem}
\begin{proof}[他の局所因子を含む共通八点切断]
補題\ref{lem:high-order-point-rank-nine}の最初のブロック計算は、印 $p$ の局所環だけで
各辺の四式を与える。定理\ref{thm:finite-high-order-common-octad}の記号で
\[
      xm=ny,\quad sm=q(x,y),\quad nt=r(x,y),\quad
      q=c r+\kappa[x,y],\quad c\kappa\ne0
\]
である。ここで $c,\kappa$ は辺ごとの定数である。射影同型
$[s:x:n]\mapsto[s:cx:cn]$ のグラフとの辺の交叉 $Z$ は、
対角で残る三二次式が固有に交わるため全体で長さ8である。
印ではその局所イデアルが $\operatorname{ann}(n^2)=\operatorname{ann}(m^2)$ で、
接空間は $X$ の接空間に等しい。後者はグラフの局所方程式 $y(x)-cx$ に一次項がないことによる。

各々の残りの局所因子でもこの零化イデアルの記述を示す。
$n=0$ と $m=0$ は同じ閉点を通る。実際、例えば $n=0,m\ne0$ なら $x=0$、
次に $sm=0$ が $s\ne0$ に反する。両者が単元なら $xm=ny$ と
$sm-cnt=\kappa[x,y]$ からこの局所因子全体は同型のグラフ上にある。
残る閉点は印と異なるので、両射影の閉埋込みにより $x,y$ はいずれも非零である。
同じ一次座標変換を二組に施し、一つの成分を両方で単元として
$x=(1,X),y=(1,Y)$ と自明化する。括弧の非零定数倍を $\kappa$ に吸収する。
すると $m=n$, $n(Y-X)=0$, $n(s-c t)=\kappa(Y-X)$ である。
$m,Y$ を消した後、この滑らかな四次元座標近傍で辺の局所方程式は
\[
     n^2(s-c t),\quad nt-r(x,y),\quad R_5,\quad R_6
\]
となる。後の三式では消した変数を代入する。グラフは最初の式を $s-ct$ に置き換えて切られる。
$s-ct$ が単元なら局所辺で $n^2=0$ となり $Z$ はそこにない。
非単元なら、どちらも有限交叉なので四式は正則列である。二つのKoszul分解を比較すると
最上段は $n^2$ 倍で、局所双対により $Z$ のイデアルは $\operatorname{ann}(n^2)$ となる。
ここでは $n=m$ なので左右の零化イデアルも一致する。
以上で四辺がスキームとして同じ八点部分を持つ。
\end{proof}

\begin{proof}[局所自己同型の位数と完全跡による延長]
八点上で四つの線型同定の合成が射影的に恒等であることは、三二次式の網に
一次係数の関係がないことから定理\ref{thm:high-order-octad-bisymmetry}の通りである。
従って全位置を同じ四空間で表せる。補題\ref{lem:common-net-alternating-part}より
窓は両対称部分と四形式の和であり、どちらも二位置の巡回移動で不変である。
関係六空間は相補二位置での全縮約像なので $R_0=R_2,R_1=R_3$。
共通点グラフ上の二移動は従って自己同型 $\sigma$ をなし、その二乗は
定理\ref{thm:finite-cubic-separators}により点環上の恒等である。
$Z$ 上で隣接位置が幾何学的に一致するため、この二移動は $Z$ 上で恒等である。
印の接空間が $Z$ と全局所環とで一致するので $\sigma$ は印を固定し、接空間で恒等である。
よって局所環上でも恒等である。実際、接空間の双対の任意の基底を持ち上げ、
各持上げ $v$ を $(v+\sigma^*v)/2$ で置き換えれば、固定される元がなお極大イデアルを生成し、
局所Artin環全体を代数として生成する。

印の局所長は少なくとも16で、残る局所因子の全長は4以下である。
反転した第二辺の六式を第一辺の座標で読むと、印での二移動が恒等なのでその全局所環で消える。
各評価は第一辺の隣接語であり、非零なら補題\ref{lem:finite-adjacent-word-support}より
支持長6以上が必要なので他の因子でも消える。各辺で同じことがいえ、有限完全交叉の
$(1,1)$ 消去式の空間がちょうど六次元であることから、全ての隣接辺は互いの反転である。
命題\ref{prop:finite-flip-obstructions}の交代三次交叉の消失で両対称性が従う。
最後の階数上界は定理\ref{thm:all-bisymmetric-as-closure}による。
\end{proof}

% !TeX root = ../papers/nc-p3.tex
\section{被約二十点での五次元印付き空間と六点の割線}
\label{sec:finite-marked-five-space-secants}
\begin{theorem}\label{thm:finite-marked-five-space-secants}
四辺が有限で点スキームが被約二十点、全中央支持4とする。
印 $p$ で二つの異なる四空間 $E_A=w(p_0,-,V_2^*,-)$ と
$E_B=w(V_0^*,-,p_2,-)$ が五空間 $E$ を張れば、
$\mathbf P(E)$ の三行三列の階数一切断 $\Sigma$ は必ず正次元である。
特に反対階数10で両側の和が五次元ずつなら、各側に階数一行列の曲線が必要である。
\end{theorem}
\begin{proof}[十三本以上の割線]
$J=\mathbf P(E_A\cap E_B)$ は平面である。位置2の座標を $A$ により $\mathbf P(E_A)$ と
同定する。平面切断の長さは7以下なので少なくとも13点 $q$ が $J$ の外にある。
印は $J$ に入る。これらの点では $A(q),B(q)$ は独立で、両方とも $p_1,p_3$ から射影した
$q_1,q_3$ を左右核に持つ。よってその割線 $\Lambda_q$ は二行二列行列の空間に入る。
$\Sigma$ が有限と仮定すると行列式がこの直線にちょうど長さ2の部分を切る。
同じ割線は二度使えない。少なくとも $\Lambda_q$ は $\mathbf P(E_B)$ に含まれず、
その交点 $B(q)$ が一点に決まり、位置0の射影は閉埋込みだからである。

有限な $\Sigma$ は $\mathbf P^2\times\mathbf P^2$ の四双一次式の完全交叉で長さ6。
射影四平面を張り、被約なら六つの代表ベクトルの唯一の従属関係には全ての項が現れる。
実際 $(1,1)$ の制限像は五次元の完全跡超平面で、被約因子の跡係数は全て非零である
（定理\ref{thm:finite-marked-secant-pencil}のKoszul計算）。
非被約で全因子が曲線状なら長さ2の部分の台は二点の組または非単純因子の唯一の二重点である。
支持数 $k\le5$ なら総数は $\binom{k}{2}+6-k\le11$。
曲線状でない点があればその接次元は2以下である。
実際Segreの接ベクトルを行と列の二つの変化に分けると、一方の変化が零である
非零接方向は射影四平面内に階数一の直線を作ってしまう。
有限完全交叉なので、接次元2なら二つの局所定義式の位数はともに2以上で長さ4以上となる。
従ってそのような点は一つであり、その点の二重点の割線は全て固定接平面内にある。
この平面内の $A(q)$ は多くとも7点である。相異なる支持間の割線は高々3本、
残る長さ2の因子の接線は高々1本なので13点には達しない。

以下 $\Sigma$ は被約六点とする。もし $k$ 点が $\mathbf P(E_A)$ に入れば $k\le4$ である。
両端が内の割線と両端が外の割線が与える点数はそれぞれ $\binom{k}{2},\binom{6-k}{2}$ 以下。
一端だけが内なら交点 $A(q)$ はその端点自身なので全て合わせて $k$ 点以下となる。
$1\le k\le4$ なら合計は11以下である。よって全六点が外にある。
このとき $J$ の外の各 $A(q)$ に対して $B(q)$ も $J$ の外にあり、
同じ計算で $\mathbf P(E_B)$ も六点を全て避ける。

\end{proof}
\begin{proof}[六点の差の平方関係]
六点のベクトル $z_i$ を $\sum_i z_i=0$ と規格化する。
二超平面の式を $a,b$ とし、$a_i=a(z_i),b_i=b(z_i)$ は全て非零で和は零である。
使用された割線の添字の組を $I$ とする。欠けた組は高々二つである。
$ij\in I$ のとき二つの交点は $z_i/a_i-z_j/a_j$ と $z_i/b_i-z_j/b_j$ により表される。
それぞれの交点の二次Veronese像の線型関係空間を比較する。
定理\ref{thm:finite-square-chart}の反対位置の二次座標像の単元倍での一致により、
二つの関係空間は各組ごとの非零定数倍で同定される。

その関係空間を明示する。六列 $z_i/a_i$ の行列を $U$ とすると核は $k(a_i)$。
差の外積でなく対称平方の関係は $U M U^t=0$、$M\mathbf 1=0$ なる対称行列で表される。
すなわち組 $ij$ の係数は $-M_{ij}$ で、対角成分は行和が零となるように決める。
$M=av^t+va^t$（ここでは $a=(a_i)$ は列）、$\sum v_i=0$ で尽きる。
これは五次元商の対称二乗への射の核が $a$ の対称積全体であるためである。
$v_i=a_i t_i$ とおき、各辺の定数倍を除くと関係空間は
\[
         (t_i+t_j)_{ij\in I},\qquad
         \sum_i a_i t_i=0,\qquad t_i+t_j=0\quad(ij\notin I)
\]
である。$b$ でも同じ表示。全ての辺和を零にすれば $t=0$ なのでパラメータ表示は単射である。
従って $D_0=k^6/\langle e_i+e_j:ij\notin I\rangle$ とすると、
$D_0/k\bar a$ と $D_0/k\bar b$ の間に線型同型があり、
各 $e_i+e_j$ の像を対応する像の非零倍に移す。棒は $D_0$ への像とする。

実は $k\bar a=k\bar b$ である。欠辺がなければ相異なる四添字による
$(e_i+e_j)+(e_k+e_l)=(e_i+e_k)+(e_j+e_l)$ を使う。
外の二添字で $a,b$ は非零なので商でもこの四元の関係は一つだけであり、
同型の四つの倍率は等しい。四添字を動かすと全倍率が等しくなり、
辺和が $D_0$ を張るため主張が従う。
欠辺が $12$ だけ、または $12,13$ なら $e_1=u,e_2=-u$、後者では $e_3=-u$ と置く。
残りの独立基底元を $r,s,\ldots$ と書くと少なくとも三つある。
$(u+r)+(-u+s)=r+s$ の三元は、外の基底元に $\bar a,\bar b$ の非零成分があるので
両商でも唯一の関係を持つ。$r,s$ の全順序対を使えばこれら全ての辺和の倍率が揃う。
それらだけで $D_0$ を張る。
欠辺が $12,34$ なら $D_0$ の基底を $u,v,r,s$ とし $e_1=-e_2=u,e_3=-e_4=v$ とする。
$(\epsilon u+\delta v)+(-\delta v+r)=\epsilon u+r$（$\epsilon,\delta=\pm1$）と
$r$ を $s$ に替えた三元を用いる。こちらも残る $s$ または $r$ の係数が非零で
関係は唯一、全ての符号を動かすと倍率が揃い $D_0$ を張る。
各場合とも、同型がこの生成系の自然な像に共通定数を掛けるので、二つの商の核直線は等しい。

従って欠辺に現れない全添字で $a_i/b_i$ は同じである。そのような添字は少なくとも二つある。
その二つを $i,j$ とすると $ij\in I$ にもかかわらず二つの割線交点は比例し、
$A(q)\notin J$ に反する。これで有限切断を全て除外した。
\end{proof}

\begin{proposition}\label{prop:single-quartic-slice-kernel-symmetry}
四辺が有限で全中央支持4とし、接次元1以下の印において、
一つの四空間 $E_A=w(p_0,-,V_2^*,-)$ の階数一切断が有限長4であるとする。
もう一つの四空間との一致を仮定せず、核側の二位置は反対対称性を持つ。
従って上の被約二十点の五空間で、階数一の整曲線が $\mathbf P(E)$ を張る場合も
核側の反対対称性が成立する。
\end{proposition}
\begin{proof}
$E_A$ だけで長さ4環の積跡による対称表示が得られる。
左右核をその表示で $l,r$ とすれば $A l=A r=0$ が点環全域で成立する。
奇側の非対称を仮定し、定理\ref{thm:finite-artin-kernel-symmetry}の証明で
どの偶座標が必要になるかを確認する。
外積 $l\times r$ の成分は独立で、単元成分を持つ $A$ の列との比例により
局所主イデアル $(\eta)$ を生成する。その支持 $Z$ は $A$ の二次小行列を消すので
$A$ 側の四座標による閉埋込みで長さ4の切断の部分となる。三次持上げに使う指標と
$Z$ 上のベクトルは $A$ の列から決まり、$A$ だけの跡模型で $a l=\alpha 1$ が成立する。
従って同定理の余次元1の三次持上げ、逆数分子の小行列、印での一様化変数による除法は
そのまま適用できる。長さ3および局所 $(1,2,1)$ で使う非零隣接語 $H\chi(a)$、
長さ4の曲線状因子の積での四つの独立な三次式も全て $A$ 側の座標を使う。
同定理のもう一つの偶行列との比較で得ていた反対座標の $Z$ への全射は、
最後の長さ4の共線因子の排除だけに必要であった。現在これは
系\ref{cor:finite-no-four-on-line}が独立に排除する。これで全ての選択肢が矛盾する。

$\mathbf P(E)$ を張る整曲線 $C$ の超平面切断 $C\cap\mathbf P(E_A)$ は
$\mathbf P(E_A)$ を張る。実際切断上消える一次式の $C$ 上の制限は、
切断式で割ると $H^0(\mathcal O_C)=k$ に入る。この切断は従って長さ4以上である。
四空間の階数一軌跡が正次元なら一本の被約母線だけであり、こうした切断は入らない。
有限なら長さ4以下なので等号となり、前半を適用する。
\end{proof}

% !TeX root = ../papers/nc-p3.tex
\section{被約二十点の五空間における階数一曲線と反対対称性}
\label{sec:finite-five-space-rank-one-curves}
\begin{theorem}\label{thm:finite-five-space-rank-one-curves}
定理\ref{thm:finite-marked-five-space-secants}の被約二十点の五空間では、
少なくとも一方の反対位置にテンソル対称性がある。
従って被約な有限窓で全中央支持が4、反対階数10なら片側対称であり、反対階数11は生じない。
\end{theorem}
\begin{proof}[母線の排除]
前節の記号を使う。$E_A,E_B$ はそれぞれ非零の片側評価が全て階数2以上で、
共通の $A(p)=B(p)$ は可逆である。後者は印での行列の階数が
$3-\dim T_pX$ であることによる。左右を交換して $E$ に母線 $ku\otimes T$、
$\dim T=2$ があると仮定する。行列の像を $ku$ で割った二行三列空間
$\bar E$ の次元は2または3である（像が $ku$ に入る行列は高々三次元）。
前節の独立な鉛筆を与える13点以上のうち、左核 $l$ が $l(u)=0$ を
満たす点は印を通る位置1の平面にあり、印を除いて高々6点である。
残る各点の鉛筆は $\bar E$ に単射し、射影後も共通右核 $r$ がある。
$\dim\bar E=2$ なら鉛筆の像は全体で、その中に階数2の行列があるため $r$ は一意。
固定 $r$ は印からの射影の一方向なので位置3で高々2点しか与えない。

$\dim\bar E=3$ とする。射影平面の階数一軌跡が有限なら、それは
$\mathbb P^1\times\mathbb P^2$ の三超平面完全交叉で長さ3である。
各接空間は第一因子への射影で単射である（核に非零方向があれば母線が丸ごと入る）。
従って曲線状で長さ2の部分スキームは高々三つ。
共通右核を持つ鉛筆は二行二列の行列式でそのような部分を切り、その線型包に等しい。
各々の右核が一意なので高々6点となり、残る7点以上に足りない。
正次元なら整曲線の両射影の包のベクトル次元 $a,b$ に対し
Segre像の包の次元は $a+b-1$ 以上である
（定理\ref{thm:finite-marked-ruling-pencil}の付値による積次元評価）。
第一因子が固定された母線はその消滅行による $E$ の評価階数を1以下にする。
平面全体が一つの固定因子を持つことも階数2の元に反する。
両因子とも動くなら $a=b=2$ で曲線が平面を張り、全 $\bar E$ が固定した
二列空間に入る。その消滅列での $E$ の評価像は $ku$ 以下となって矛盾する。
従って第二因子の固定された母線 $k^2\otimes v$ が入る。
$r(v)\ne0$ なら $\bar E$ の $r$ 評価は全射で核は一次元にすぎない。
全ての独立な鉛筆の点が $l(u)=0$ または $r(v)=0$ に入り、
それぞれ印以外に高々6点なので再び13点には足りない。
\end{proof}
\begin{proof}[平面二次曲線の場合]
階数一軌跡には前節より整曲線がある。それが全体を張れば
命題\ref{prop:single-quartic-slice-kernel-symmetry}を適用する。
それ以外の曲線では両因子が動く。実際、固定因子の線型空間と $\mathbf P(E)$ の交叉は線型であり、
曲線を含めば母線を含む。両射影が平面を張れば積次元評価により全体を張る。
一方が直線を張ればその法線での評価核は三次元以下であり、
同じ評価から他方も直線を張る。このとき曲線は二行二列のSegre二次曲面の平面切断で、
母線がないから滑らかな二次曲線である。

この曲線の包を $D$ とし、左右の基底を選んで $D$ を左上二行二列の対称行列全体とする。
$E/D$ の同次座標を $y=(s,t)$ として一つの線型補空間を選び、その行列を
\[
 M(y)=\begin{pmatrix}K(y)&c(y)\\ d(y)^t&h(y)\end{pmatrix}
\]
とする。$y\mapsto (c,h)$ と $y\mapsto(d,h)$ はともに単射である。最後の列または行で
$E$ を評価した階数が2以上だからである。$[N]$ が $D$ の外の階数一点で $y$ がその商座標なら、
$h(y)\ne0$ である。零なら最終列と最終行がともに非零なのに隅成分が零で階数一に反する。
左上ブロックは $c d^t/h$ に決まる。よってこの点の存在条件は二次式
\[
 Q(y)=c_1d_2-c_2d_1-h(K_{12}-K_{21})=0
\]
であり、点も $[y]$ から一意である。

$Q\ne0$ とする。二つまでの根から次のように高々二点を選ぶ。
$h(y)\ne0$ の根からは上の階数一点を選ぶ。$h(y)=0$ の根では
非零な $c(y),d(y)$ が比例するので二次曲線上の $[c(y)c(y)^t]$ を選ぶ。
この選択は曲線上の接線で $D$ の外へ出るものも捉える。
実際 $[vv^t]$ での階数一行列の接方向は、一方の因子だけを動かす二項の和である。
その最終行・列は $h=0$ かつ $c,d\in kv$ を要求するため、商座標が非零なら上の根から得た点である。

各 $q$ で $A(q),B(q)$ が比例すれば、その任意の共通商への像も比例する。
独立なら鉛筆は共通左右核を持ち、母線がないので行列式はその上で長さ2の階数一因子を切る。
その因子に $D$ 外の点があれば鉛筆は選んだ点を通る。相異なる二点が $D$ 内にあれば鉛筆も
$D$ 内にある。一点に集中すれば接線であり、$D$ 外に出る場合も選んだ点を通る。
一方 $D$ 内の独立な鉛筆の共通左右核は二行二列部分の外の固定方向である。
実際二次元対称行列の空間で指定した非零二列を殺す部分は一次元である。
従ってこの例外点は位置1で印を通る固定直線にあり、印自身を除いて高々2点である。

選んだ点の包を含む任意のベクトル二平面 $T\subset E$ を取り $F=E/T$ とする。
$A,B$ の $F$ への像の外積三成分は例外の高々2点だけに支持される。
これらは形式的な位置2と0の双一次式としては三つ独立である。
実際 $F_A,F_B\subset F$ を両写像の像とすれば各次元は2以上で和は三次元。
交代形式が $F_A\times F_B$ 上で消えるとする。両方が平面なら異なる平面で、
共通の基底元を含む適合基底の三対を全て消す。一方が全空間なら根基が二次元以上になり、
やはり零である。よって二十点への双一次評価は単射でなく、
定理\ref{thm:finite-opposite-product-gap}が偶位置の対称性を与える。

$Q=0$ が恒等的で $h\ne0$ の場合、$y$ を動かした
$(c(y),h(y))(d(y),h(y))^t$ はもう一つの滑らかな二次曲線となる。
実際両列写像は単射であり、二次Veroneseを与える。この曲線を張る平面と最初の平面は
ちょうど一点で交わる。二つの行支持および列支持はいずれも異なる二次元空間であり、
共通部分が入る一行一列空間の次元は1、二平面は射影四空間で必ず交わるためである。
この二曲線はその一点で交わる（共通行列は階数一で、第二の曲線については
$h=0$ の値を取る）。その被約和 $C$ は連結で $\mathbf P(E)$ を張る。
$E_A$ はどちらの曲線も含めない（単独の四空間では正次元切断が母線だけである）。
従って超平面切断は非零因子であり、$H^0(C,\mathcal O_C)=k$ からその切断は
$\mathbf P(E_A)$ を張る。長さ4の判定は前節の命題の証明と同じで核側対称性が従う。

最後に $Q=h=0$ がともに恒等式なら $c,d$ は二次元への同型で、$d$ は $c$ の定数倍である。
最後の行の倍率を揃え、さらに最後の行の定数倍を上二行に足せば $K_{12}-K_{21}$ を消せる。
実際 $c$ の二成分は独立であり、この操作で任意の一次の歪成分を加えられる。
こうして $E$ 全体が対称となる。$D$ 外の階数一点はない。
$E_A$ は $D$ を含めないから、その階数一点は二次曲線上の高々二点である。
$A(q)$ が階数2なら射影した左右核は一致し、印では両射影ベクトルは零である。
従って左右三座標の外積も高々二点に支持され、三つの独立な奇位置の双一次式から
同じ積次元定理で核側対称性が従う。

反対階数については、印の両側の和の次元の和は反対階数以下である
（定理\ref{thm:finite-rank-nine-all-marks}の直交性の証明）。
従って階数11以下なら一方の和は4または5である。4なら同定理の曲線状印での
共通切断の議論、5なら本定理を使う。片側対称なら反対切断は四次元の対称二乗を経由するため階数は10以下である。
\end{proof}

% !TeX root = ../papers/nc-p3.tex
\section{一つの二重点を含む十九点支持での五空間判定}
\label{sec:finite-one-double-five-space}
\begin{theorem}\label{thm:finite-one-double-five-space}
有限辺・全中央支持4の標準有限Frobenius窓で、点スキームが十八の単純点と
一つの二重点 $d$ からなるとする。反対階数10なら少なくとも片側対称であり、
反対階数11は不可能である。
\end{theorem}
\begin{proof}[超平面の残余と六点の割線]
単純印 $p$ で四空間 $E_A,E_B$ の和が五空間 $E$ の場合を考える。
記号は定理\ref{thm:finite-marked-five-space-secants}に従う。$E_B$ の式を $h$、
$J=\mathbf P(E_A\cap E_B)$ とする。印以外で両行列は共通左右核 $l,r$ を持つ。
$A$ 側の平面 $J$ の外には少なくとも $19-7=12$ 個の異なる点があり、
各々で独立な鉛筆がある。さらに $Z=V(\operatorname{ann}h(A))$ は長さ13以上で、
単純印を含まない。$Z$ 上で使う鉛筆を局所的に明示する。
$A,B$ が剰余で独立ならそのまま二本を取れる。
比例する場合は $h(A)$ が非零の二重点のsocleを張り、$Z$ 上では長さ1である。
このとき定数倍を揃え $B=\lambda A+\epsilon V$ と書けば
（$\epsilon^2=0$, $\lambda$ は単元）、$V$ の剰余は $A$ の剰余と独立で、
共通の左右核を消す。$h(V)\ne0$ も従う。従ってどの局所因子でも $Z$ 上に独立な鉛筆が得られる。

$E$ の階数一切断が有限なら長さ6である。非被約切断は既出の割線計数で不可能となる。
全局所因子が曲線状なら異なる割線は高々11本、曲線状でない点があれば
その点での接線は固定平面に入り他の割線と合わせ $J$ の外で高々11点しか与えない。
ここでは $J$ の外の12個の異なる点だけで足りる。

被約六点なら相対的にも各鉛筆は一定である。共通左右核を持つ二行二列の鉛筆の
行列式が相対長さ2の因子を切り、六つの定数点の中の二点とその線型包は微小変形できない。
特に $h(A)$ が単元の因子で $B$ は一定の割線と $E_B$ の一点の交点に等しく、
$B$ が閉埋込みのためその因子は単純である。$h(A)$ が非零の冪零元のときも、
剰余で $A,B$ が独立なら二重点全体で割線が固定され、$E_B$ 内の割線では $h(A)=0$ となる。
比例するときの $Z$ 上の鉛筆も $h(V)\ne0$ のため $E_B$ の外に出る。
これで $Z$ は被約で、全ての割線が $E_B$ と一点で交わり、
異なる支持点は異なる割線を使う。よって十三本以上である。

六節点が両超平面の外にあることと、欠辺高々二本の平方の応力計算は同定理と同じである。
実際どちらの超平面についても $k$ 個の内部節点があれば混合割線の交点は
それらの点に限られ、二座標の閉埋込みから数は $\binom{k}{2}+k+\binom{6-k}{2}$ 以下である。
平方の対角同値には $Z$ の被約な行だけを使えばよい。
欠辺に現れない二添字を取ると、対応する使用割線で $A,B$ の値が比例する。
現在これは矛盾ではなく、その点が必ず $d$ であることを示す。
単純な $Z$ の点では $h(A)$ は単元で、二つの値は独立だからである。
従って有限切断のこの側では $A(d)$ と $B(d)$ の剰余値が比例する。
\end{proof}
\begin{proof}[母線での長さ二の計数]
$ku\otimes T$ が $E$ に含まれる場合を定理\ref{thm:finite-five-space-rank-one-curves}と同じ二行三列の
$\bar E$ で調べる。$l(u)=0$ は位置1の平面条件で印以外の長さが6以下である。
$\dim\bar E=2$ ならその外で $r$ は一方向に固定されるので支持数は高々 $6+2$、
二重点による超過を加えても13に足りない。
$\dim\bar E=3$ でその射影平面の階数一切断が有限なら長さ3で曲線状である。
非被約なら長さ2の部分の選び方は高々二通りで、$l(u)\ne0$ の支持は高々4点。
従って超過1を足しても足りない。被約三点なら許される右方向は高々三つである。

ここで平面と三つの固定右方向による長さの計数が二重点にも使えることを確認する。
$l(u)$ が単元なら、射影鉛筆の相対二点は上の三点中の二点に固定され、右方向も局所環上で固定される。
右方向が許される三方向の外なら、$\bar E$ の右評価階数は剰余で2以上である
（核が二平面ならそこに二点の因子を切る）。従って右評価核は局所環上で
高々一つの自由直線に含まれる。共通鉛筆の二本の射影はこの直線に入るので、
単元成分を持つ係数で線型結合して射影を零にできる。
実際二係数の生成イデアルは体または双対数環で単項であり、一方が他方の倍である。
元の独立な二本のその結合は $ku$ に像を持ち単元成分を持つので、$l(u)=0$ が局所環全体で従う。
二つの条件がともに剰余で成立するときは、それぞれ少なくとも長さ1を捉えるため長さ2まで数えられる。
従って $\operatorname{length}Z\le6+3\cdot2=12$ となり矛盾する。

$\bar E$ の階数一切断が正次元なら、同定理の支持の議論より固定列の母線
$k^2\otimes v$ が含まれる。$l(u)$ が単元なら射影鉛筆は右評価の核に入るので、
$r(v)=0$ が局所環上で成立する。実際その評価列は $r(v)e_1,r(v)e_2,Cr$ である。
剰余で $r(v)=0$ のとき $Cr$ は非零でなければならない（元の $E$ の片側評価は階数2以上）。
従って二次小行列の消失が $r(v)=0$ を与える。逆に $r(v)$ が単元なら右評価階数2で、
先と同じ単項イデアルの議論が $l(u)=0$ を与える。両方が剰余で零なら両平面で長さ2を数えられる。
従って二平面が印以外に捉える長さ $6+6$ では足りない。これで母線はない。
\end{proof}
\begin{proof}[二次曲線の孤立点と微小鉛筆]
全 $E$ を張る整曲線があれば単純印での一切断の長さ4判定で核側対称となる。
それ以外の整曲線は同定理により左上の対称二行二列空間 $D$ に入る滑らかな二次曲線である。
商の二変数の二次式 $Q$ が恒等零で隅成分が非零一次式なら、二つの二次曲線の連結和が全体を張り、
やはり長さ4である。両方が恒等零なら $E$ 全体を対称にでき、階数一切断の台はこの二次曲線だけである。
この場合を含め、埋め込まれた接線を以下で数える。

$J$ の外の12点の独立鉛筆を考える。$D$ 内の鉛筆では左右の射影核は固定方向 $l_0,r_0$ であり、
印以外で高々2点である。$D$ 外の孤立階数一点を通る鉛筆の点は、その行因子を消す
位置1の固定平面上にあるので高々6点である。$D$ 外の点を通らず鉛筆自体は $D$ 外に出るには、
二次曲線上の一点で外へ出る接線となる必要がある。
$Q\ne0$ ならその接点は二次式の根の一つであり、孤立点と合わせて高々二つである。
各接点の外向き接線は一つの接平面に入る。実際 $E/D$ でのその一次方向は唯一である
（$h=0$、$c,d$ が接点の固定因子に比例する条件と最終列評価の単射性）。
$D$ 内の一次方向も曲線の接方向だけである。
この固定平面のうち少なくとも一方の超平面での交点は直線以下であるか、
両超平面がこの平面を含んで $J$ の外の点に使えない。
よって各接点は高々3点を与える。一つでも孤立点をこのような接点に替えた場合は
$2+6+3=11$ 以下である。同じ孤立節点の非被約化を使っても節点は一つにしか数えられない。
従って $Q\ne0$ で12点に達するには、二つの異なる根がどちらも $D$ 外の孤立点を与えねばならない。

この場合 $\Sigma$ は二次曲線とこれら二つの被約点とのスキームとしての非交和である。
孤立点では隅成分が単元で $Q=0$ が局所式となり、単根のため被約である。
曲線上には外向きの接線がないから全点で局所次元と接次元はともに1である。
孤立二点のベクトルの包を $T$ とし $E/T$ での $A,B$ の外積三成分を取る。
剰余で二本が独立なら二重点全体で鉛筆上の階数一因子は有限平坦長2である。
剰余で孤立点を含めばその部分は単根として定数点に一意に延びる
（その点の接空間は零）。従って鉛筆が $T$ に会い商の外積は局所環上で零である。
全因子が二次曲線上なら相対的にも $D$ 内であり、二点因子の線型包として鉛筆全体が $D$ に入る。
その一般行列は左上二行二列で可逆なので $l=l_0$ が局所環上で成立する。
剰余で比例する二重点では外積はsocleに入り、$B=\lambda A+\epsilon V$ の剰余鉛筆で
非零のずれだけを調べればよい。そこでも $D$ 内の場合に限り外積に寄与し得るので、
位置1の固定直線のイデアルが外積を消す。単純印では外積自体が零である。
三成分は従って印以外でこの直線上の長さ2以下の環の双対に入る
（点環の完全跡による零化イデアルと商の双対との同定）。
形式的には三つ独立であることは同定理の商三空間の証明によるので偶側対称となる。

最後に $Q=h=0$ の対称表示では、$E_A$ の階数一切断は曲線との高々二つの幾何点を持つ。
$A(q)$ が階数2なら、その単元小行列を使い左右核 $l,r$ は二重点でも比例する。
従って三座標の奇側外積が長さ3以上の環に支持されて独立となるためには、
二つの相異なる階数一点の一つが二重点で、そこに長さ2、他方に長さ1が必要である。
長さ2に寄与するには剰余で外積が非零、すなわち $l,r$ が独立である。
このとき二重点上で独立な二核が対称行列を消すので、その二次小行列は全て恒等零となる。
従って $A$ の階数一切断はこの長さ2のスキームと他の単純点を含む。
実は $E_A$ の切断は長さ4である。二次曲線の二つの交点は相異なるため横断的で、
$E$ の階数一切断は両点で接次元2以上である。各点で $Q=h=0$ の模型を
$\left(\begin{smallmatrix}S&c\\c^t&0\end{smallmatrix}\right)$ として、
曲線の行因子 $v$ に対し $c=v$ の方向がSegreの接方向にあり、曲線の方向と独立になることから分かる。
超平面を課した後も各点の接次元は1以上であり、合計長は4以上となる。
単独四空間に母線はなく切断長は高々4なので長さ4の判定が適用できる。
外積が最初から従属なら奇側の積次元定理を使えばよい。
\end{proof}
\begin{proof}[二重点を印に取り直す]
対称性がないと仮定する。反対階数が10または11なら直交する二つの和の次元から、
各単純印で少なくとも一方が四空間または五空間である。
四空間なら曲線状の印での共通切断の判定が既に対称性を与える。
五空間なら以上より有限六節点の場合だけが残り、各印 $p$ でこの側の
$w(p_0,-,d_2,-)$ と $w(d_0,-,p_2,-)$（または巡回した二式）が比例する。
十八印中少なくとも九印で同じ側を選べる。今度は $d$ を固定すると、
これら九点の像は同時にその両四空間に入る。各座標の平面切断は長さ7以下で、
中支持は4なので九点はそれぞれの四空間を張る。よって二重点 $d$ に共通四空間がある。
これは曲線状の印であり、定理\ref{thm:finite-rank-nine-all-marks}の共通切断の証明のうち
曲線状印を扱う部分は反対階数を使わず対称性を与える。矛盾である。
片側対称なら反対切断は対称二乗を経由するため階数11は不可能となる。
\end{proof}

% !TeX root = ../papers/nc-p3.tex
\section{曲線状有限窓での母線の排除と三重点の取り直し}
\label{sec:finite-curvilinear-pencils}
\begin{proposition}\label{prop:curvilinear-marked-pencils}
全中央支持4の有限窓の全局所環を曲線状とし、単純印の二つの四空間が五空間 $E$ を張るとする。
$\mathbf P(E)$ の階数一切断は母線を含まない。切断が被約六点なら、
印と異なる非単純点の少なくとも一つで $A,B$ の剰余値は比例する。
切断が平面二次曲線と二つの孤立被約点の非交和なら偶側対称である。
\end{proposition}
\begin{proof}[付値で測った鉛筆]
前節と同じく $h$ を $E_B$ の式とする。$h(A)$ の零化イデアルで定まる残余 $Z$ は長さ13以上。
印を除いた局所環 $R=k[t]/t^m$ で、$A,B$ が剰余で独立なら元の鉛筆を $R$ 上で使う。
比例なら、$A$ の単元成分に合わせて単元 $\lambda$ を選び、
$B=\lambda A+t^a V$ とする。差が零なら $h(A)=0$ で外積も零である。
そうでなければ $1\le a<m$ で、$A,V$ は $R'=k[t]/t^{m-a}$ 上で独立かつ
共通左右核を持つ。これは $A$ の選んだ成分で差が零、別の成分でちょうど位数 $a$ となるためである。
$h(A)=-t^a h(V)/\lambda$ より $Z$ の局所因子は $R'$ の商である。
独立の場合も $R'=R$ として、$Z$ 上には常に直和因子である二次元鉛筆がある。

母線 $ku\otimes T$ を仮定し、定理\ref{thm:finite-five-space-rank-one-curves}の
二行三列への射影 $E\to\bar E$ を用いる。$Z$ の一因子 $k[t]/t^s$ での位数は $s$ で打ち切る。
右評価 $\bar E\to k^2$ は剰余で非零である（零なら元の右評価階数は1以下）。
$\dim\bar E=2$ なら単元の軸を掃き出すと核は一本の自由直線に入る。
独立鉛筆の二本の射影を、その直線上で単元成分を持つ係数対で結合して零にできる。
二係数の一方が他方を割るからである。元の結合は単元成分を持ち $ku$ に値を取るので $l(u)=0$。
これだけでは印以外の長さ6までである。

$\dim\bar E=3$ では右評価の二次小行列の最小位数を $b$ とする。
一つの単元の軸の掃き出しと主イデアル性により行列を対角 $1,t^b$ と零列にできる。
鉛筆の二本の射影は第一成分零、第二成分が $(t^{s-b})$ にあり、
第三成分を上と同じく原始的な係数対で打ち消せる。元の結合を $t^{s-b}$ を法として取れば
$ku$ に値を持つ。従って
\[
                \operatorname{ord} l(u)+b\ \ge\ s .
\]
$\bar E$ の平面の階数一切断が有限なら、右評価の二次小行列の零点スキーム $Y$ も有限である。
実際、評価核が二次元となる各右方向は、平面の有限曲線状長3切断の
長さ2部分の割線で一意に決まる（同定理の母線の議論）。
$Y\subset\mathbf P^2$ の長さは3である。高さ2の極大小行列イデアルの
Hilbert--Burch分解は $0\to\mathcal O(-3)^2\to\mathcal O(-2)^3\to\mathcal I_Y\to0$ である。
空の場合も上界だけなら自明である。長さ $e_i$ の各局所因子のイデアルは
その点の極大イデアルの $e_i$ 乗を含む。従って全因子での $b$ の和は、
対応する右方向の逆像長を各 $e_i$ 倍した和以下である。
一定の右方向は印を通る直線の印以外の部分で長さ2以下ゆえ、合計は6以下である。
階数一切断が正次元なら、同じ平面の分類により $\bar E=\langle e_1v^t,e_2v^t,C\rangle$。
評価列は $v(r)e_1,v(r)e_2,Cr$、かつ $v(r)=0$ なら $Cr$ の剰余は非零なので、
小行列イデアルは層として $(v(r))$ である。この場合も印を通る平面で $\sum b\le6$ となる。
$\sum\operatorname{ord} l(u)\le6$ と合わせると $13\le\operatorname{length}Z\le12$ に矛盾する。

被約六点の場合、$R'$ 上の鉛筆の行列式で切る二点因子は定数六点の部分であり、鉛筆はその一定割線である。
$h(A)\ne0$ ならこの割線は $E_B$ 内にない。比例の剰余の場合も、
$A,V$ がともに $R'$ 上で $h$ を消せば $t^a h(V)=0$ に反する。
$B$ の $R'$ 上の像は従って割線と超平面の固定交点となり、閉埋込みより $R'$ は長さ1。
$Z$ の長さ13以上が互いに異なる割線を与え、前節の六点の応力の証明により一つ以上で
両超平面の交点が等しい。その点は単純になれない。

最後に二次曲線と被約二点の非交和では、$R'$ 上の鉛筆の二点因子が孤立点を含めば
その点は単根として一定に持ち上がる（被約孤立点に二重接線はない）。
その場合は二つの孤立点の包で割った商の外積が零となる。
残る鉛筆は相対的にも曲線の平面 $D$ 内にあり、左核は一定方向になる。
比例の剰余なら外積に $t^a$ が掛かるので、$R'$ 上で得たこの固定方向のイデアルによる
零化は元の環にも戻る。印で外積は零である。前節の商の三成分と長さ2の判定で偶側対称となる。
\end{proof}

\begin{theorem}\label{thm:finite-single-triple-five-space}
全中央支持4の有限窓で点スキームが十七単純点と一つの曲線状三重点からなるとき、
反対階数10なら少なくとも片側対称であり、反対階数11は不可能である。
\end{theorem}
\begin{proof}
単純印で五空間が出る場合を調べる。$J$ の外に少なくとも $18-7=11$ 個の異なる支持があり、
各々で元の二本が独立である。以下、対称性がまだ得られないと仮定する。

切断が有限非被約なら、接次元2の点では全接線がその接平面にある。
両超平面がこの平面を含めば $J$ の外には寄与しない。それ以外は片方で直線以下を切り
高々3点しか捉えず、他の割線の高々4本を加えても足りない。
全点曲線状なら11本に達するのは単純四点と一つの二重点の場合だけで、
十本の異なる二点の組と一つの接線が全て使われる。
各超平面内の節点が四つ以上なら、単一四空間の長さ4判定が対称性を与える。
一つから三つならその一つと外の異なる二節点との二割線は同じ点で超平面に交わり、
座標埋込みの単射性に反する。従って両超平面とも全節点を避ける。
$Z$ の別の支持を加える余地はない。実際そこでも先の残余鉛筆の剰余割線は既に使われ、
その $E_B$ との唯一の交点が別の支持と一致してしまう。
従って $Z$ の長さ13を得るには三重点が十一支持の一つであり、そこで全長3を使う。
しかし元の独立鉛筆の行列式因子は長さ3の基底上では存在してもその割線の動きは
高々長さ2の基底を経由する。異なる二節点の場合、各根は単根で一意に分離し、
動ける一根は長さ2の節点への写像だけで決まる。一節点に重なる場合は一定の
長さ2節点全体に入る有限平坦次数2の閉部分であり、同次数ゆえ等しい。
いずれでも $E_B$ との単一交点 $B$ の閉埋込みは長さ3を通せない。

正次元切断で整曲線が全体を張る場合と二つの二次曲線が張る場合は、既出の長さ4判定を使う。
命題で母線はない。平面二次曲線の模型で $Q=h=0$ なら全 $E$ が対称にでき、
$E_A$ の階数一部分は $\mathbf P^2$ の二次Veroneseに二超平面を課したものとなる。
これは有限である（単一四空間の正次元切断なら母線を含む）。よって二つの二次式の完全交叉で長さ4。
$Q\ne0$ の二つの孤立被約点の場合も命題で済む。残る場合、前節の支持上界では11点に達するには
孤立点一つと外向き接点一つが必要である。$D$ 内の鉛筆で2点、孤立点を通る鉛筆で6点、
外向き接線で3点の各上界を同時に達し、三集合は互いに素となる。
最後の三点は両超平面の少なくとも一方と固定接平面との交線にあり、
前の二群もそれぞれ印を通る直線・平面の印外の長さ上限を使い切っている。
$Z$ の他の支持は残余鉛筆によってこの三条件の一つに入らねばならず、加える余地がない。
再び三重点は十一支持の一つで、そこに元の独立鉛筆がある。

この場合の外向き接点での完備局所環を求める。二つの相異なる根の一つは $h=0$ の単根で、
$h$ は非零一次式である。接点の左上を1とし、階数一行列を
\[
 \begin{pmatrix}1&x&z\\ y&yx&yz\\ w&wx&wz\end{pmatrix}
\]
と書く。商座標に $z,h$ を取れる。実際 $h=0$ の方向で $c,d$ の第一成分だけが非零となるよう
曲線の基底を選べるからである。$w=bz+ch$, $yz=\alpha h$, $wx=\beta h$,
$x-y=kz+k'h$, $h=wz$ で $b\ne0$。$Q=h((\beta-\alpha b-k)z+(-\alpha c-k')h)$
より $\beta-\alpha b-k\ne0$ である。$w=z b/(1-cz)$ と代入し、
$z(x-y)$ を二通りに計算すれば $z^2$ と単元の積が零となる。
従って局所環はちょうど $k[[x,z]]/(xz,z^2)$ で、$D$ は $z=0$ である。
孤立点はもう一つの単根なので被約であり、曲線の残りでは接次元1なので滑らかである。

三重点での鉛筆が孤立点を通れば、単根の分離によって全長でその点を通り、六点で飽和した
平面に入ってしまう。外向き接線の場合、行列式の相対長さ2の代数は剰余で $1,z$ が基底なので
全体でもこれが基底である。$x=A_0+A_1z$ と書けば $xz=z^2=0$ より $A_0=0$。
従って長さ2因子は固定接平面にスキームとして入り、その包である鉛筆もそこに入る。
三点で飽和した交線に長さを加えることはできない。
$D$ 内の場合、剰余鉛筆が特別接点に集中した二重点を切れば、それは曲線の接線でやはり固定接平面にあり、
既にその交線上にある別の三点と両立しない。特別接点を避けるなら全長で $D$ 内にある。
特別接点ともう一つの相異なる点を切れば単根二つに分離し、特別接点に延びる根の
$z$ は $z^2=0$ より $k[t]/t^3$ の $(t^2)$ に入る。
従って少なくとも $t^2$ を法として二根とも $D$ 内にあり、鉛筆全体とその固定左核も
その長さ2で $D$ 内の条件を満たす。二つの異なる支持で飽和した印外の直線に更に一重を加える矛盾となる。

以上で各単純印の五空間は、対称性がなければ被約六節点の場合に限られ、
命題よりその割線の一致点は唯一の三重点である。階数10または11では各単純印で一方の和の次元が
5以下で、4の場合は既に対称性が従う。十七印の少なくとも九印で同じ側の一致が生じるから、
前節と同様に三重点に印を移すと共通四空間となる。曲線状印での共通切断判定が対称性を与える。
\end{proof}

% !TeX root = ../papers/nc-p3.tex
\section{二つの二重点における割線の一階平方関係}
\label{sec:finite-two-double-five-space}
\begin{theorem}\label{thm:finite-two-double-five-space}
四辺が有限で全中央支持4の標準有限Frobenius窓において、点スキームの幾何学的支持が
十八点以上なら、反対階数10で少なくとも片側対称性があり、反対階数11は生じない。
\end{theorem}
\begin{proof}[合流する二節点の平方]
長さ20のGorenstein環なので長さ3以下の局所因子は曲線状である（接次元2で長さ3ならsocleが二次元）。
前節までの結果から十六単純点と二つの二重点 $d,e$ を扱えば十分である。
単純印 $p$ の五空間で、対称性がまだないとする。階数一切断が有限非被約なら
前節の十一割線の場合だけが残り、両超平面が全節点を避け、残余 $Z$ の別の支持も加えられない。
よって $J$ の外にある十一点のうち二つが二重点で、双方の一次方向が使われる。
そこにある割線は、非被約節点と単純節点を結ぶ異なる二本に限る。
実際、単純節点同士の割線は動かず、二重点そのものを切る接線も前節の有限平坦次数2の議論で動かない。
固定割線と超平面の交点は一意で、点座標の閉埋込みでは二重点を受け取れない。

四つの単純節点のベクトルを $u_1,\ldots,u_4$、二重節点を $v+\epsilon v'$ とする。
六つの列にただ一つある従属関係において $v'$ と四つの $u_i$ の係数は全て非零である。
これは定理\ref{thm:finite-marked-secant-pencil}の完全跡超平面が各局所socle上で非零であることによる。
代表の調整と $v'$ への $v$ の定数倍の加算により $v'=-\sum_i u_i$ とでき、
$v,u_1,\ldots,u_4$ は独立である。
$E_A$ の式 $a$ を $a(v)=1$ と規格化し $a_i=a(u_i)\ne0$,
$S_a=\sum a_i$, $x_i=v-u_i/a_i$ と置く。$x_i$ は $E_A$ の基底である。
四本の混合割線での超平面交点は $x_i$、六本の単純節点同士の割線では $x_i-x_j$、
二重節点での接線では $W_a=\sum a_i x_i$ となる。
二重点の一次方向を使う混合割線では、交点が非零の速度で
$x_i+\epsilon W_a$ の方向に動く。実際、$v+\epsilon v'$ を $a$ の値で割った微分は
$v'-a(v')v=W_a$ である。$b$ でも同じ表示をダッシュ付きの基底で使う。

十本の平方 $X_i=x_i^2$, $X_{ij}=(x_i-x_j)^2$ は対称二乗の基底である。
定理\ref{thm:finite-square-chart}の平方空間の単元倍での一致を双対化すると、
この十本を対応する十本の非零倍に送る線型同型が得られる。倍率を $s_i,s_{ij}$ とする。
この同型は十一番目の平方 $W_a^2$ を $cW_b^2$ に送り、微分した二本 $2x_i W_a$ を
それぞれ $2c_i x'_iW_b$ に $X'_i$ の倍を足したものに送る。
後者の不定項は各二重点での単元による代表の変更であり、$c,c_i\ne0$。
$W_a^2$ の $X_{jk}$ 係数は $-a_ja_k$ なので
\[
                   s_{jk}=c\frac{b_j b_k}{a_j a_k}.
\]
微分した式の $X_{ik}$ 係数は $-a_k$ ($k\ne i$) だから $c_i=c b_i/a_i$。
同じ式の $X_k$ 係数は $a_k$ ($k\ne i$) であり、$s_k a_k=c_i b_k$ となる。
微分する二添字を $1,2$ とすれば $k=3$ で比較して $b_1/a_1=b_2/a_2$。
従って使用された単純節点間の割線 $12$ において $a,b$ の交点が一致し、
そのグラフ点が $J$ の外にあることに反する。$S_a=0$ の場合も同じである。
\end{proof}
\begin{proof}[有限六節点と埋込み接点の印の取り直し]
切断が被約六点なら、命題\ref{prop:curvilinear-marked-pencils}の十三割線の議論で
両超平面は全節点を避ける。ここでは非単純点の\emph{全て}で二行列が剰余で比例する。
実際、剰余で独立なら局所環全域での鉛筆が一定割線となり、$B$ の値も超平面との交点に
固定される。これは非単純点を閉埋込みで写せない。この議論は任意の曲線状グラフにも成立する。

切断が正次元なら前節の証明により、未処理の場合は平面二次曲線の一つの外向き接点と
一つの孤立点で、$D$ 内の二支持、孤立点を通る六支持、外向き接線の三支持が各上界を使い切る。
$Z$ に別の支持が入らないことも同じである。孤立点の単根を通る相対鉛筆は全長で
その点を通り、外向き接線を持つ相対鉛筆も全長で固定接平面にあることは
前節の局所環 $k[[x,z]]/(xz,z^2)$ の計算による。従ってその六点と三点に二重点は入れられない。
二つの二重点はともに $D$ 内の鉛筆を与える二支持である。
この二つの剰余鉛筆の左核は同じ方向なので、核側の第一の位置で
$p,d,e$ は共線である。$d,e$ を先に固定すれば、この直線に入る単純印は
系\ref{cor:finite-no-four-on-line}により高々一つである。

階数10または11なら各単純印で少なくとも一方の和が五次元以下である。
共通四空間なら既に対称性が得られる。対称性がないとして、十六印のそれぞれについて
五空間の側を一つ選ぶ。埋込み接点のために除く印は側ごとに高々一つであり、
少なくとも十四印が有限六節点を与える。その七印以上で同じ側が選ばれる。
その側で $d$ を印に取り直すと七印の行列の剰余値が両四空間の交叉にある。
さらに $d$ 自身の二つの値は同じであり、それぞれの座標が閉埋込みなので
合計八つの異なる点が交叉に入る。どの平面切断も長さ7以下なので両四空間は一致する。
曲線状印での共通切断判定が対称性を与える矛盾である。対称二乗の次元10により階数11も排除される。
\end{proof}

% !TeX root = ../papers/nc-p3.tex
\section{十七点支持の曲線状有限窓と接平面の接触次数}
\label{sec:finite-seventeen-site-pencils}
\begin{lemma}\label{lem:embedded-conic-contact-budget}
命題\ref{prop:curvilinear-marked-pencils}の五空間の階数一切断が平面二次曲線、
一つの単純な孤立点と一つの単純な埋込み接点からなる場合を考える。
$J$ の外に十以上の支持があるとする。核側の二位置の各々で、単純印は非単純点のみに
支持される長さ2部分スキームの割線上にある。それが一つの局所因子に集中するには、
その因子は長さ4以上を要する。
\end{lemma}
\begin{proof}
$Z$ 上の独立な相対鉛筆を用いる。孤立点を通る剰余鉛筆は全長でその点を通り、
その固定行因子による平面が捉える印外の長さは6以下である。
外向き接線を持つ剰余鉛筆は全長で固定接平面にある。
この接平面は少なくとも一方の超平面で直線以下を切る（さもなければ外の支持は $2+6$ 以下）。
その直線で数えられる長さは3以下である。

二次曲線の平面 $D$ に入る剰余鉛筆では、局所基底環を $R=k[t]/t^s$ とする。
埋込み接点を避ければ全長で $D$ 内にあり、その長さ $f=s$ を固定左核の直線に数えられる。
接点と別の点を切るなら、接点に延びる単根の法線座標 $z$ は $z^2=0$ を満たす。
$f=\operatorname{ord}z$ は $2f\ge s$ を満たし、長さ $f$ で両根は $D$ にあり、再び核が固定される。

最後に曲線の内側の接線として埋込み接点に集中する剰余鉛筆を扱う。局所模型
$k[[x,z]]/(xz,z^2)$ で行列式の二点因子の代数は $R$ 上 $1,x$ を基底とする。
$x^2=A+Bx$, $z=C+Dx$ と書くと $xz=0$ より $DA=0$, $C=-DB$ で、
$z^2=0$ の $x$ 係数より $D^2B=0$ である。従って
\[
    f=\operatorname{ord}D,\qquad e=\min(\operatorname{ord}A,\operatorname{ord}B),
    \qquad s\le 2f+e
\]
となる（位数は $s$ で打ち切る）。長さ $f$ では $z=0$ より鉛筆は $D$ 内にある。
長さ $e$ では $x^2=0$ なので固定接平面にある。実際前節までの局所座標で階数一行列は
定数行列、$x,z$ の二つの接行列、および $x^2$ の一行列の線型結合である。
いずれも二点因子全体が線型部分空間に入ればその相対線型包も入る。

こうして各因子の長さは、孤立点の固定平面、接平面が与える直線、および
$D$ の固定左核の直線での長さを順に重み $1,1,2$ で数えたもの以下である。
最後の直線が印以外に捉える長さは2以下であるから
\[
              13\le \operatorname{length} Z\le 6+3+2\cdot 2=13 .
\]
使った全ての不等式は等号である。固定左核に割り当てた $f$ の和はちょうど2で、
$f>0$ の各因子の長さは $s=2f+e\ge2f$ である（単根分離では $e=0$）。
よってその直線が捉える長さ2は非単純因子だけにあり、一つなら元の因子の長さは4以上である。
右核も同じ計算で扱える。
\end{proof}

\begin{lemma}\label{lem:confluent-three-first-jets}
十七点支持で三つの二重点があるとする。単純印の五空間の有限な階数一切断は、
対称性がないなら六つの被約点に限られる。
\end{lemma}
\begin{proof}
$J$ の外に十点以上ある。接次元2の節点と、曲線状で支持四点以下の切断は
それぞれ上界7と8で除かれる。五節点なら前節の $v,u_1,\ldots,u_4$ の模型である。
どちらの超平面も四節点以上を含めば単独の長さ4判定で対称性が出る。
内部節点が $1\le k\le3$ 個なら各内部節点から外に出る $5-k$ 本の割線の超平面交点は同じで、
外のグラフ点の総数は $11-k(4-k)\le8$ に落ちる。従って両超平面とも節点を避ける。
残余 $Z$ は各割線上の $B$ の剰余値で区別され、支持数は10または11である。
長さ2の因子は全て、相対割線が動く混合辺 $i$ にあり、元の二重点全体を使って $h(A)$ は単元である。

前節の十基底 $X_i,X_{ij}$、接線平方 $W_a^2$ と一次方向 $D_i=2x_i W_a$ を使う。
二本の交点は\emph{同じ}混合割線の変形から来る。すなわち $v+\epsilon v'$ を動かす速度は
両側で同じである。このため平方の同型が $X_i$ を $s_i X'_i$ に送れば、
$D_i$ は $s_iD'_i$ と $X'_i$ の倍の和に移る。これは両側で $a(v)=b(v)=1$ とした
交点式を同じ一様化変数で微分し、二重点上の単元を微分して得る。
$r_i=b_i/a_i$ と書く。$D_i$ の $X_j,X_{ij}$ 係数は $j\ne i$ で $a_j,-a_j$、
$i$ に接しない二添字成分は零である。

十基底が全て使われていれば二つ以上の可動混合辺 $i,j$ があるので、
それらの微分を比較すると $s_j=s_i r_j=s_{ij}=s_j r_i=s_i$。
よって $r_i=1$ であり、混合辺 $i$ の二交点が一致して $h(A)$ が単元であることに反する。
これは接線平方だけが欠ける場合も含む。他の場合は支持十点、三つの可動混合辺が必要であり、
全二重点をそこで使うので $Z$ の全ての支持で $h(A)$ は単元である。
欠けた基底 $X_m$ だけは $\tau X'_m+\sum_{j\ne m}\delta_j X'_j$ に送られ、
$\tau\ne0$ とする（ここで $j$ は十基底の添字）。

$m$ が単一添字 $k$ なら残りの三混合辺が全て可動である。
$D_i$ の $m$ 係数から $\tau a_k=s_i b_k$ が全ての可動 $i$ で成立し、$s_i$ は共通である。
可動二添字 $i,j$ に対して、第三の可動添字の微分の $ij$ 係数は両側で零なので
$\delta_{ij}=0$。従って $s_{ij}=s_i r_j=s_j r_i$ で $r_i=r_j$。
使われた単純節点間の割線 $ij$ で二超平面の交点が一致して矛盾する。

$m=12$ とする。$3,4$ がともに可動なら両者の微分は $m$ 係数が零なので、
先の十基底での二辺比較をそのまま適用できる。残るのは $1,2,3$ が可動の場合である。
$D_2$ の $13$ 係数は零で $m$ 係数は非零だから $\delta_{13}=0$。
従って $D_1,D_3$ の $13$ 係数から $s_{13}=s_1 r_3=s_3 r_1$。
一方 $D_3$ は $m$ 係数が零なので、その $1$ 係数によって $s_1=s_3 r_1$。
再び $r_3=1$ が単元条件に反する。
\end{proof}

\begin{theorem}\label{thm:finite-seventeen-site-pencils}
全中央支持4の有限窓で十七の幾何点を持つ曲線状点スキームについて、
反対階数10なら少なくとも片側対称であり、反対階数11は不可能である。
\end{theorem}
\begin{proof}
長さの分割は一つの四重点、三重点と二重点、または三つの二重点である。
単純印の五空間で対称性がなければ、正次元切断は先の接触次数補題の場合だけが残る。
実際、外向き接点だけの配置は $2+3+3$ 以下、孤立点が一つだけで接点のない場合は
$2+6$ 以下で、$J$ の外の十点に足りない。他の配置は
命題\ref{prop:curvilinear-marked-pencils}と定理\ref{thm:finite-single-triple-five-space}の
長さ4切断の議論による。
有限切断は被約六点または単純四点と二重節点の場合に限り、両超平面とも節点を避ける。
後者の証明は直前の $11-k(4-k)$ の計数をそのまま使える。
前者で十三本の割線による節点回避は曲線状窓すべてに使える。
長さ3以上のグラフ点ではどちらの場合も二行列の剰余値が比例する。
剰余で独立なら相対割線の変形と超平面の交点が閉埋込みでその局所環を受け取る必要があるが、
六節点では動かず、五節点でも混合割線の一階変形までしか動かないからである。

四重点だけなら接触次数補題の長さ2はその接線、三重点と二重点なら両者を結ぶ割線である。
各側においてこの固定直線には高々一つの単純印しか加えられない。
階数10または11なら各印で次元5以下の側を選ぶ。共通四空間なら対称性が得られ、
それ以外でも十五以上の単純印から高々二印を除けば、少なくとも七印が同じ側で有限切断を与える。
長さ3以上の点に印を移せば、その点自身を含め八つの異なる交点があり、共通四空間が強制される。

三つの二重点の場合、十四の単純印のそれぞれで次元5の側を一つ選ぶ。
各側の埋込み接点の例外印は二重点の三対の各割線で高々一印、従って高々三印である。
残る有限切断の印では補題\ref{lem:confluent-three-first-jets}により被約六節点となり、
三つの二重点\emph{全て}で二行列が比例する。
一つの側でこうした有限切断印が四つ以上あるとする。三つの二重点それぞれに印を取り直すと、
その二つの四空間の交叉を一方の変数の射影三空間に引き戻した線型空間には
同じ四つ以上の単純印が入り、さらに取り直した二重点自身が入る。
共通四空間なら対称性が得られるので、この線型空間は平面以下である。
四つの異なる印は共線になれないから平面を一意に張る。
従って三つの二重点は\emph{同じ}平面に全て入る。
この平面の長さ上界7より当該側の有限切断印は全体で四つ以下である。

各側で有限切断が高々四印、例外が高々三印なので十四印を覆うには
両側とも四印と三印で等号でなければならない。一方の側で得た平面は既に
四つの単純印と三つの二重点の剰余を含む。他方の例外印はこの位置で
二つの二重点を通る直線にある（接触次数補題は両核位置に成立する）。
この別の単純印も同じ平面に入り、長さ7を超える。これで階数11を含めて矛盾する。
\end{proof}

% !TeX root = ../papers/nc-p3.tex
\section{平面長さ四の一点を持つ十七支持の窓}
\label{sec:finite-planar-four-seventeen}
\begin{theorem}\label{thm:finite-seventeen-all}
全中央支持4の標準有限Frobenius窓で、四辺が有限で幾何学的支持が十七点以上であれば、
反対階数10で少なくとも片側対称性があり、反対階数11は生じない。
\end{theorem}
\begin{proof}[独立な鉛筆を持つ平面四重点]
前節で残るのは十六単純点と平面四重点 $d$ である。局所環は
$R=k[x,y]/(x^2,y^2)$ と書ける。実際 $\mathfrak m/\mathfrak m^2$ の積は
一次元のsocleへの非退化対称対合であり、二つの等方直線の基底を用いればよい。
$R$ での非零主イデアルの長さは単元なら4、非零一次部分があれば2、socleなら1である。
単純印 $p$ の五空間で $A(d),B(d)$ の剰余値が比例する印をまず数える側に回す。
残りの印では $d$ の環全体に独立な鉛筆があるので、非主イデアルによる差の除法は不要である。
$Z$ は前節までの $h(A)$ の零化イデアルの商とする。$d$ でのその因子は
長さ1,2の曲線状環、$R$、または空である。印以外の全単純点にも剰余二行列が比例すれば
外積は零で、その点は $Z$ に入らない。

有限な階数一切断で $J$ の外の十支持を受け取れるのは被約六点または単純四点と二重節点の型である。
後者では前節の計数で両超平面が節点を避ける。前者で超平面内の節点数が4以上なら
単独四空間の長さ4判定で対称性が得られる。2または3なら、混合割線を内部の同じ節点でまとめた
$\binom{k}{2}+k+\binom{6-k}{2}$ は9以下で十点に足りない。
従ってどちらの型でも超平面は節点間の割線や二重節点の接線を含まない。
$d$ での相対割線は定数か一つの長さ2節点を介するので、その超平面との交点である
$B$ の閉埋込みが長さ4を受け取れない。

曲線がある場合、母線を除くと全体を張る曲線の長さ4切断の議論はそのまま使える。
二次曲線と二つの被約孤立点の場合も $d$ では環全体で独立な鉛筆を使えるため、
定数孤立点の分離と固定核の直線による商外積の零化の証明が成立する。
前節の十支持の計数により、残るのは一つの埋込み接点と一つの孤立点がある場合である。
ここで補題\ref{lem:embedded-conic-contact-budget}の局所上界は任意のArtin基底 $T$ 上で使える。
単根が埋込み点に延びる場合、$z^2=0$ なら長さは $2\dim T/(z)$ 以下である。
二点因子の剰余が曲線の内側に接する場合、同補題の $D^2(A,B)=0$ より
\[
 \dim T=\dim T/(D^2)+\dim (D^2)
          \le 2\dim T/(D)+\dim T/(A,B).
\]
実際 $(D)/(D^2)$ は $T/(D)$ の商、$(D^2)$ は $T/(A,B)$ の商である。
前者の一次商 $T/(D)$ 上では鉛筆は曲線の張る平面に入り、後者 $T/(A,B)$ 上では固定接平面に入る。
こうして全長13の等号まで同じである。単独の非単純因子で固定核の直線に長さ2を使うので
$d$ の全長4を要する。内側に接する場合は $(A,B)$ が剰余で零だから接平面への長さも正となり、
等号は不可能である。残る単根の場合は $z^2=0$, $\dim R/(z)=2$ である。
従って $(z)$ は先の二つの等方方向の一つが定めるイデアルである。
socleは $(z)$ に入るので一次方向の選択で商は一意に決まる。
この場合の $p$ は核側でその長さ2部分の接線にあり、一つの側で高々二つの単純印しか得られない。
\end{proof}

\begin{proof}[母線に落ちる鉛筆とsocle]
行因子が一定の母線 $ku\otimes T_0$ を仮定する。列因子の場合は転置して同様である。
$\bar E$ と右評価を命題\ref{prop:curvilinear-marked-pencils}の通り取る。
$\dim\bar E=3$ のとき右評価の二次小行列のイデアルを $I$、
$\dim\bar E=2$ なら $I=(1)$ とする。$Z$ の各因子 $T$ で $l(u)I=0$ が成立すれば矛盾する。
実際 $(l(u))$ は $T/I$ の商なので $\dim T\le\dim T/(l(u))+\dim T/I$。
第一項の総和は印を通る平面で6以下である。第二項も同命題の評価の退化スキームにより6以下である。
長さ3の退化スキームの場合、各局所環を剰余体で組成列分解し、グラフへの引戻しの右完全性で
その逆像長を各方向の逆像長に組成長を掛けた和で抑えればよい。
これは曲線状でない $T$ でも成立する。曲線状因子では先の付値の不等式が $l(u)I=0$ を与える。

残る $T=R$ 上でこの積が非零になるには、鉛筆全体が剰余で $ku$ に値を取り、
$l(u)$ が非零socleに入る必要があることを示す。右評価には単元の軸があるので
射影した鉛筆の係数は一成分を消去した自由加群に入る。その階数は
$\dim\bar E-1\le2$ で、階数2の場合は残る一行の評価式の二係数が $I$ を生成する。
射影した二本の一方が剰余で非零ならこれを基底元にできる。
第二の射影からその成分を引くと、残る一成分 $b$ は $bI=0$ を満たす
（自由加群が一次元なら $b=0$）。$b$ が単元なら $I=0$。
さもなくば元の行列の結合は剰余で非零かつ $ku$ に値を持ち、
$b$ を法としても $ku$ に値を取るため $l(u)\in(b)$ である。

二本とも剰余射影が零なら、上の消去後の射影係数を二列の行列 $P$ として全成分が
$\mathfrak m$ にある。元の二行列の $u$ 成分は剰余で独立な二行ベクトルであり、
二つの座標を取れば可逆な二次行列 $S$ になる。共通左核の式をこの二座標で読むと
$l(u) S+KP=0$ となる。ここで $K$ は射影した行列に残りの左核成分を掛けて同じ二座標を読む写像である。
$l(u)$ に非零の一次部分があれば、その部分を消さない接ベクトルで微分した等式から
$K$ の剰余は階数2でなければならない。従って $K$ は二行二列可逆であり、
$P=-l(u)K^{-1}S$。残りの右評価行が $P$ を消すので $l(u)I=0$ である。
これで積の非零は主張のsocleの場合に限られる。

この例外で $u$ の定める平面は $d$ の位置1の接平面である。
実際 $p_1$ からの射影での式 $l(u)$ は元のこの平面式の単元倍であり、$d$ で一次部分が零となる。
$d$ で二つの偶位置を固定した階数1の行列を $u_d v_d^t$ とすると $u$ は $u_d$ に比例する。
行列の左右因子が両奇位置の接平面の式であることは、隣接する二点の縮約消失を微分すれば従う。
例えば $w(d_0,-,q_2,q_3)=0$ の $q=d$ での微分では $d_3,d_0$ が隣接するため
$q_2$ を微分した項が零で、$u_d v_d^t$ が位置3の接方向を消す。もう一側も同じである。
\end{proof}

\begin{proof}[四重点における二つの母線]
今 $d$ を印に取り直す。母線の例外印 $p$ では $A^d(p),B^d(p)$ はともに
$u_d$ を行因子とする階数1行列である。それぞれ $A^d(d)=B^d(d)=u_dv_d^t$ と異なる。
よって両四空間に $u_d$ を固定した母線が一本ずつある。
単一四空間の階数一切断はこの場合その一本の直線だけである。従って同じ側で
列因子固定型の例外印を併せて持つこともなく、全ての母線の例外印と $d$ は
一方の変数側で同じ直線に入る。そのような単純印は高々二つである。

この場合 $d$ における二つの四空間の和は、対称性がなければ六次元以上である。
四次元なら伸長型の印における共通切断の判定
（定理\ref{thm:finite-rank-nine-all-marks}の証明）を使う。
五次元と仮定し、$u_d$ による射影を取る。各単一四空間の射影には全非零元が
階数2の二次元鉛筆がある（定理\ref{thm:finite-marked-ruling-pencil}の冒頭の計算）。
$d$ を通るこの側の変数の平面 $J_d$ の外には十以上の単純点があり、
その各点で共通左右核を持つ二つの行列は独立である。$l(u_d)=0$ の点では、
階数2の射影鉛筆の非零元は左核を持たないから各行列は対応する母線上にある。
これで高々二点である。$l(u_d)\ne0$ の点では共通鉛筆が射影後も独立である。
射影空間が二次元なら共通右核がなく不可能である。三次元ならその射影平面の
階数一切断は有限である（正次元なら全非零元が階数2の射影直線と交わってしまう）。
従って共通右核の方向は長さ3の平面切断の高々三本の割線から得られ、
各方向は印以外に高々二単純点である。合計八点では足りない。

最後に対称性がないとし、各単純印で直交性から次元5以下の側を一つ選ぶ。
単純印の共通四空間なら既に矛盾する。各側で $d$ で剰余値が比例する印の集合を取ると、
その集合と $d$ は $d$ で取り直した両四空間の交叉に入る。
交叉が平面以下ならこれだけで単純印は高々五つである。実際 $d$ を通る任意の平面は
$R$ で少なくとも長さ2を切り、全長は7以下である。母線の例外があれば先の和の下界から
交叉は直線以下であり、比例する単純印は高々二つに落ちる。
よって母線例外のない側で選べる印は高々 $5+2=7$、
ある側でも高々 $2+2+2=6$ である。二側合わせても十六単純印を覆えない。
これで残る平面四重点の場合の対称性と反対階数11の排除が示された。
\end{proof}

% !TeX root = ../papers/nc-p3.tex
\section{有限窓の反対階数九における片側対称性}
\label{sec:rank-nine-finite-symmetry}
\begin{theorem}\label{thm:rank-nine-finite-symmetry}
長さ1の局所点を含む有限Frobenius点グラフで、反対階数が9とする。
中支持3の点による認証がないとき、窓は少なくとも一方の対向位置の交換で不変な表示を持つ。
両対称なら既出の有限AS判定を使える。片側対称表示の対称二乗の核は一つの二次形式で、
その階数が2以下なら全次数で正則である。従ってこの範囲の未判定窓では、核の階数は3か4に限られる。
\end{theorem}

\begin{lemma}\label{lem:irreducible-four-plane-brackets}
三行跡零行列の既約な四次元部分空間 $D$ の交換子の線型包は三次元以上である。
\end{lemma}
\begin{proof}
$D$ に相異なる三固有値の行列があることをまず示す。
反対なら各元 $X$ の三次特性式の判別式が零で、$\operatorname{tr}(X^2)^3$ と
$\operatorname{tr}(X^3)^2$ は恒等的に一定の非零倍率で一致する。
ともに恒等零か、$D$ の座標環の一意分解で共通の一次式の二乗・三乗の非零倍である。
従って次元3以上の冪零部分空間 $N$ がある。三行の冪零部分空間で次元3以上なら
厳密上三角全体に同時共役であることをここで確認する。
巡回冪零元がないと全元が階数1以下で、共通像直線または共通核平面および跡零条件から次元は2以下である。
巡回元を $E_{12}+E_{23}$ とすれば、これと $Y\in N$ の全線型結合の平方・三乗の跡が零なので
$Y_{31}=0,Y_{32}=-Y_{21}$、$Y_{21}(Y_{11}-Y_{33})=0$ となる。
$Y_{21}$ が恒等零でなければ $Y_{11}=Y_{33}=a,\ Y_{21}=c$ と書き
$3a^2+c(Y_{12}-Y_{23})=0$。よって $a=kc,\ Y_{23}=Y_{12}+3k^2c$ である。
さらに $\det Y=0$ は $c^2(k^3c-Y_{13})=0$ を与える。自由座標は $Y_{12},c$ の二個で矛盾する。
従って $Y_{21}=0$ で全元が上三角冪零である。
$D$ 全体が冪零であることはこれで排除された。残るとき $N$ は一次式の核であり、
平方の跡の偏極より $\operatorname{tr}(D N)=0$。従って $D$ も上三角で既約性に反する。

交換子の包が二次元以下と仮定する。一般の $X\in D$ の $D$ 内の中心化空間は次元2以上となり、
三固有値が異なるので跡零中心化空間の全体である。特に $X^2-\frac13\operatorname{tr}(X^2)I\in D$。
稠密性と偏極で $k I+D$ はJordan積 $(UV+VU)/2$ で閉じる五次元空間になる。
ある三固有値行列を対角化すればその三つの射影冪等元がこの空間に入る。
各冪等元とのJordan積で対角成分と三つの非対角対 $ij,ji$ に分解できる（可換な対角作用の同時重み分解）。
非対角部分の総次元は2である。既約なので、可能な非零行列成分を辺とする三頂点有向グラフは強連結である。
ゆえに相異なる二つの非対角対で各一次元を使い、両辺は双方向でなければならない。
その二元のJordan積は三番目の対に非零成分を持ち、五次元性に反する。
\end{proof}

\begin{proof}[定理の対称性]
印付きの $g=\dim\langle G,F\rangle,\ q=\dim\langle Q,P\rangle$ はともに3以上で $g+q\le7$。
両方3なら前節の証明で $\bar w=0$ を必要としたのは対称・交代の混合型を排除する段だけである。
混合型でも $G$ を交代、$Q$ を対称とすれば $F=G,P=Q$ である。
$a,c$ 交換では星印を含む項は不変、純成分の差は
$[G_a,G_c]_{bd}+[Q_b,Q_d]_{ac}=0$ である。よって少なくとも一方は対称となる。

残りは巡回して $g=4,q=3$ とする。$P_d=Q_{Jd}$ であり、
$G,F$ の三平面は二平面で交わる。このため $D=\langle G,F\rangle$ 全体の外積は交叉交換子の入力
$G\wedge F$ で張られる。実際、共通二元とそれぞれの補元による六基底外積を使える。
既約な $D$ に前補題を適用すると、$[G,F]$ の包は三次元以上、一方で $[Q,Q]$ の包と同じ次元なので
ちょうど三次元である。交換子式の $b,d$ 跡から $\sum_b[Q_b,Q_{Jb}]=0$、
独立性から $J$ は対称である。合同変換で $J=I$ と取れば $D$ の全交換子は交代行列となる。
ここで $D$ は全て対称である。実際、交代部分の像次元が1,2なら、前節と同じ中心化空間の議論により
その核に必要な二つ以上の跡零対称行列がない。像次元3なら三元の交代部分を外積基底に取れ、
同じ十五次の検査でその三元は純交代となる。残りの対称元は三外積の全てと交換するので零になってしまう。
よって $G,F$ は対称で、$P=Q$ はこの段では任意の跡零三行列組でよい。
$b,d$ を交換しても星印を持つ項は不変、純成分の差も交換子式で零となる。
これで全ての場合の片側対称性を得た。

対称側の二次核は $10-9=1$ 次元である。これが階数2以下なら $x y+ y x$ 型に書け、
対称性で $w(x,-,y,-)=0$ となる（$x=y$ も許す）。
六次元標的の隣接切断 $\beta,\gamma$ では $\beta(x,-)$ と $\gamma(y,-)$ の像が互いに直交する。
有限ファイバーより両評価は階数3以上なので両方3、従って $x$ は点グラフ上にある。
その三テンソルは中央を $y$ で消せるから中支持3となり既知の全次数認証が適用できる。
\end{proof}

% !TeX root = ../papers/nc-p3.tex
\section{階数四の対称核と二項切断の全射性}
\label{sec:smooth-kernel-binary-transfer}
四辺が有限な標準有限Frobenius窓で、一方の対向位置が対称、反対階数9とする。
対称二乗の核の生成元が階数4なら、対称側の点座標を $E=A\otimes B$、$\dim A=\dim B=2$ とし、
核を二行二列の行列式に対応するテンソルに取れる。
従って四テンソルは二つの $A$ 位置の交換と二つの $B$ 位置の交換で別々に不変である。
これは $\operatorname{Sym}^2(A\otimes B)^*=
(\operatorname{Sym}^2 A^*\otimes\operatorname{Sym}^2 B^*)\oplus
(\wedge^2 A^*\otimes\wedge^2 B^*)$ の分解による。
反対側の点座標を $F,G$ とし、隣接分解を $\beta:E\otimes F\to H$、$\gamma:E\otimes G\to H^*$、$\dim H=6$ とする。
\begin{theorem}\label{thm:smooth-kernel-binary-transfer}
いずれかの二項因子、例えば $A$ において、ある $a\in A$ で
$\beta(a,-,-):B\otimes F\to H$ と $\gamma(a,-,-):B\otimes G\to H^*$ がともに全射なら、
窓は両対称反転窓である。長さ1の局所点の存在は不要である。
従ってこの有限開条件では全次数の正則性が既出の有限判定で決まる。
\end{theorem}
\begin{proof}[一つの作用素]
$a$ を第一基底に取り、二写像の二係数を $P_0,P_1$、$Q_0,Q_1$ とする。
$A$ の対称性から $P_0^t Q_1=P_1^t Q_0$。$Q_0$ の全射性により $P_1$ は
$P_0$ を経由し、両全射性から $P_1=T P_0,\ Q_1=T^t Q_0$ となる単一の $T\in\operatorname{End} H$ がある。
$T$ の各左固有汎関数は $R_\beta$ の元を与え、その偶位置の支持は固有値に対応する
$A^*$ の一次式と $B^*$ の積の二平面に入る。純関係がないので階数は常に2である。
同じ固有値の独立な二汎関数があれば、これらの右支持の和は $F^*$ 全体である。
実際、和が二次元以下なら二行二列の鉛筆に階数1がある。
三次元なら二関係の四つの行の依存は一つで、関係の係数と二行の間の二行列として非退化である
（階数1の依存なら非零の関係が階数1以下となる）。
基底を合わせると二関係は $u_1f_1+u_2f_2,\ u_1f_2+u_2f_3$ と書ける。
これらを順に $u_1,u_2$ と右からテンソルして加えたものは、両端対称の非零三次重なりで両端支持が2以下である。
$R_\beta,R_\beta^{\rm op}$ は連続し、全ての三次重なりは四テンソルの一文字縮約なので、
その点での辺の評価階数が2以下となり有限性に反する。$\gamma$ についても同じである。
固有空間の次元が3以上なら、対応する $a'\in\mathbb P A$ の $P(a')$ は階数3以下となり、
$\{a'\}\times\mathbb P B\times\mathbb P F$ の四次元上で三つ以下の双一次式を課すだけで
正次元の辺の点対ができる。従って各固有値のJordanブロックは高々二つである。
\end{proof}

\begin{proof}[巡回性と二次行列]
不変因子を $p$ および $q\mid p$ とし、巡回の場合 $q=1$ とする。
$m=\deg p,\ e=\deg q,\ m+e=6$ で、$H=D_p\oplus D_q$、$D_h=k[t]/(h)$ と書く。
$T$ は $t$ 倍である。各モニック $h$ に対し剰余代表の最高次係数（次数 $\deg h-1$）を取る
$\ell_h$ で積を対合させ、双対も同じ表示に取る。
$P_0,Q_0$ の $D_h$ 成分を $B$ の二基底で $U_{h,0},U_{h,1}$ と $V_{h,0},V_{h,1}$ と書く。
前者は $F^*$、後者は $G^*$ に値を持つ $D_h$ ベクトルである。$W_h=U_{h,0}V_{h,1}^t-U_{h,1}V_{h,0}^t$ とする。
第二の二項因子の対称性から $W=W_p+(p/q)W_q$ の各成分は $D_p$ 上で $1,t,t^2$ に直交する
（$q=1$ なら第二項なし）。実際 $\ell_q(s)=\ell_p((p/q)s)$ であり、二つの $A$ 変数の積の係数を読めばよい。
$m\ge3$ なので直交部分は次数 $m-4$ 以下の通常の多項式全体である。次数上限が負なら零と読む。
以下この代表 $W(t)$ を使う。

$e>0$ なら $m\le5$ で、$W$ の三次小行列式は $p/q$ で割れるが次数は $3(m-4)<m-e$ である
（$m=3$ では $W=0$）。従って全て恒等零となる。
一方、$q$ の根で $p,q$ の重複度を $l_1\ge l_2>0$ とする。
その局所環で両不変因子の成分を規定次数まで持ち上げると、先の固有空間の議論により
四列 $U_{p,0},U_{p,1},U_{q,0},U_{q,1}$ は剰余体で独立であり、$V$ の四列も独立である。
従って $W\bmod (t-\lambda)^{l_1}$ は左右可逆変換で、二次交代ブロックと
$(t-\lambda)^{l_1-l_2}$ 倍の二次交代ブロック（単元倍を許す）の直和である。
ある三次小行列式はこの局所環で非零であり矛盾する。
従って $T$ は巡回であり、$m=6$、$W$ は二次以下になる。

全三次小行列式は六次までで $p$ で割れるので $\operatorname{adj} W=p C$、$C$ は定数行列である。
$W$ が一般に可逆なら $C$ は可逆で $W$ は定数可逆行列のスカラー多項式倍となるが、
各根での階数は2である。一般階数3なら $C$ の非零行・列から定数左右核が生じる。
例えば $b^tW=0$ なら、全ての局所環で $V_{p,0},V_{p,1}$ は独立な直和因子を張るから
$b^t U_{p,0}=b^t U_{p,1}=0$ が $D_p$ で成り立つ。この文字は $P_0$ と $P_1=TP_0$ の双方で消え、
一因子全支持に反する。一般階数は従って2である。この議論は $p$ が重根だけの場合にも成立する。
\end{proof}

\begin{proof}[二つの四文字の同定]
$W(t)$ を二次に斉次化し、両側の一般像を射影直線上で飽和させた階数2の部分束を
$S_F\subset F^*\otimes\mathcal O$, $S_G\subset G^*\otimes\mathcal O$ とする。
$S_F$ の双対の次数を $d_F$ とすると、大域生成と直線上の分裂から $h^0(S_F^*)=2+d_F$。
先の定数核の排除により $F\to H^0(S_F^*)$ は単射で、$d_F\ge2$。
同じく $d_G\ge2$ である。一方、$W$ は $S_G^*\to S_F(2)$ を誘導し、行列式は
次数 $4-d_F-d_G$ の非零切断である。従って $d_F=d_G=2$、誘導写像は全点で同型、
二つの大域切断の単射も同型である。

$\det S_F\simeq\mathcal O(-2)$ を選ぶ。$W$ による同型と階数2の外積対合により
$S_G\simeq S_F$ と同定して $W$ を $S_F$ 上の非退化交代二テンソル（捻り2）にできる。
両者の双対の全切断が四文字そのものなので、この束同型は周囲の定数線型同型を誘導し、
両埋込みも双対評価で対応する。従って定数同定後の $W(t)$ は全域で交代である。
$D_p$ に戻ると、二列ずつの $U,V$ は同じ階数2直和因子を張る。
局所的に $V=U L$ と書けば $U\bigl(\begin{smallmatrix}0&1\\-1&0\end{smallmatrix}\bigr)L^t U^t$ の交代性から
$L$ は単元スカラー行列である。全局所成分を合わせ $V=h U$、$h\in D_p^\times$ と書ける。
従って $\gamma$ は相手四文字を同定した $\beta$ に対し、完全対称積対合
$(s,t)\mapsto\ell_p(hst)$ で与えられる。これは隣接二組の同時交換で四テンソルを保ち、
既知の偶位置の交換と合わせて奇位置の交換も保つ。これで両対称性が従う。
\end{proof}

\begin{theorem}\label{thm:smooth-symmetric-kernel-all}
同じ有限点グラフで対称核の生成元が階数4なら、どちらかの二項因子が前定理の同時全射条件を満たすか、
元代数が指数成長するかのいずれかである。従ってこの滑らかな対称核の全域で正則性は有限に決まる。
特に長さ1の局所成分を持つ反対階数9では、片側対称表示で残る未判定の二次形式は階数3だけである。
\end{theorem}
\begin{proof}[一つの写像の二方向の階数]
まず $A$ を動かす $P(a):B\otimes F\to H$ の一般階数は5以上である。
4以下なら各 $a$ で $\mathbb P B\times\mathbb P F$ の四次元に高々四つの超平面を課すため、
各 $a$ に辺の点対が存在して有限性に反する。$B$ 方向でも、また $\gamma$ でも同様である。

$\beta$ は両方向とも一般階数5にはなれないことを証明する。
そうと仮定し、左核の原始斉次ベクトルをそれぞれ $l(a),n(b)\in H^*$、次数を $k,h$ とする。
五次小行列式で表示できるので $k,h\le5$、原始性により各射影直線で基点はない。
一般の $(a,b)$ で $\beta(a,b,F)$ の階数は4で、その零化空間は $l(a),n(b)$ の二元で張られる。
実際、後二者が一般に比例すれば両方が一定となり $\beta$ の全射性に反する。
$6\times4$ 行列 $\beta(a,b,-)$ の四次小行列式は二次数 $(4,4)$ で最大公約因子はない（退化点は有限）。
従ってこれらの相補Plücker座標は $l(a)\wedge n(b)$ をその成分の最大公約因子で割ったものに比例する。
特に $k,h\ge4$ である。

$P_0,P_1$ の共通核を割った入力次元を $r$ とする。$r\ge6$ なら一般に階数5である六列を独立な入力から取る。
$6\times6$ 鉛筆の随伴行列の非零行は $l(a)$ の多項式倍であり、
次数が5なので、この六列の非零右核ベクトルが次数 $5-k\le1$ で得られる
（随伴行列を共通の左核でくくる）。定数核は既に割ったため、一次の
$a_0u_0+a_1u_1$ である。$P_0u_0=P_1u_1=0$, $v=P_0u_1=-P_1u_0\ne0$ となる。
ところが $P_0^t Q_1=P_1^tQ_0$ にこの二入力を入れると $v$ は $Q_0,Q_1$ の両像を消す。
$\gamma$ の全射性に反するので $r=5$、他方向でも同じである。

従って $B\otimes F$ の三次元共通核 $K$ には階数1元がない。あれば $a$ 全体で辺の零点となる。
この核は基底変更で $B\otimes S^3 B\to S^4 B$ の乗法核である。
実際、$\mathbb P(B^*)$ 上で評価した $K\otimes\mathcal O(-1)\to F\otimes\mathcal O$ は
全点で階数3、余核は $\mathcal O(3)$ で、大域切断で $F$ はその全切断に等しい。
完全三次系の線形関係は二変数の通常の乗法核である（体積形式で $B$ とその双対を同定してよい）。
よって $\beta(a,b,f)=\widetilde P(a)(bf)$、$\widetilde P$ は五列の鉛筆である。他方向にも同じ記述がある。

滑らかな二次曲面 $\mathbb P A\times\mathbb P B$ 上の点グラフの制限長は $20-4k$ である。
実際、$l$ の核束を $K_l=\ker(H\otimes\mathcal O\to\mathcal O(k))$ とし、
$\mathbb P A\times\mathbb P B\times\mathbb P F$ に引き戻すと、辺の零点は
階数5束 $K_l(1,1,1)$ の切断の零点である。有限なので空の場合も含めて最高Chern数が長さとなる。
各超平面類を $x,y,z$ とすると $(x+y+z)^5-kx(x+y+z)^4$ の $xyz^3$ 係数は $20-4k$ である。
射影はグラフ上で閉埋込みなので他方向の計算と比較して $k=h$ となる。
$k=4$ なら $\widetilde P(a)$ の最大小行列式は $l(a)$ の各成分と同じ一次式の積である。
零点上では $bf$ の単元成分があるため、$\widetilde P(a)(bf)=0$ により全最大小行列式が
スキーム上で消える。共通の一次式が消え、グラフの制限はこの一つの被約な母線に入る。
逆方向の一母線にも入るため長さは高々1となり、計算値4に反する。

$k=5$ なら $\widetilde P$ の階数は全点で5、束の余核は
$5\mathcal O(-1)\to6\mathcal O\to\mathcal O(5)$ である。
従って $l$ は完全五次系による有理正規曲線であり、$n$ も同様である。
$l\wedge n$ の最大公約因子は二次数 $(1,1)$ なので両曲線は一致する
（一方の変数を固定した成分は、他側が埋込みなのでありえない）。
二つのパラメータを射影同型で同定し $H=S^5 A$ と書くと、$\beta(a,b,f)$ は
全ての $a,b$ についてそれぞれが指定する二つの一次因子を持つ五次式である。
二つの普遍一次因子は独立な変数を用い互いに素なのでその積で割れる。
$\beta$ 自身は $a,b$ にそれぞれ一次で、商は $f$ のみに依存する。
従って偶位置の $A\otimes B$ は対称二乗を経由し、一文字全支持に反する。
これで $\beta$ の両方向同時の階数5を除外した。$\gamma$ にも同じ証明を適用できる。
\end{proof}

\begin{proof}[交叉する二方向の非全射]
同時全射の方向が一つもなければ、入れ替えを除いて $A$ 方向で $\beta$ の一般階数が5、
$\gamma$ が6、$B$ 方向では逆となる。一般の $a$ で $P(a)v=0$ とすると、
対称性と $Q(a)$ の全射により全 $a'$ で $P(a')v=0$ である。従って共通核の次元は3で、
先ほどの乗法核の同定により $F=S^3B$、$\beta$ は $B\otimes S^3B\to S^4B$ を経由する。
反対側では $G=S^3A$ で、$\gamma$ は $A\otimes S^3A\to S^4A$ を経由する。
隣接関係の四列は $R_\beta,R_\beta^{\rm op},R_\gamma,R_\gamma^{\rm op}$ である。

経路空間の双対を各点文字のテンソルに取り、二項多項式で表す。
各 $E,F$ 辺では、偶文字の $B$ 変数と奇文字の $B$ 変数の角括弧があれば二変数を掛ける写像で零となる。
$G,E$ 辺では $A$ 変数どうしの角括弧で同じことが成り立つ。
全辺の括弧を掛けると、内部の偶文字では $A,B$ の各一次をちょうど使い、内部の奇文字では
三次のうち二次しか使わない。各内部奇位置に任意の一次形式を掛け、端の不足次数を固定した非零形式で埋める。
得たテンソルは全隣接関係を消す。多項式環で固定した非零積の乗法は単射なので、内部奇位置の数を $s$ として
同じ長さに $2^s$ 個の線型独立元がある。これは標準Hilbert増大を排除する。
\end{proof}

% !TeX root = ../papers/nc-p3.tex
\section{階数三の対称核の二項四次部分}
\label{sec:cone-kernel-binary-transfer}
\begin{theorem}\label{thm:cone-kernel-binary-transfer}
四辺が有限な標準有限Frobenius窓が片側対称で、対称二乗に階数3の核元があるとする。
対称な二位置の点座標を $E=k e\oplus S^2 B$, $\dim B=2$ とし、
核元を二項二次形式の判別式に対応させる。
前節と同じ六次元への隣接分解 $\beta,\gamma$ に対し、一変数の二つの鉛筆
\[
 P(a)(b\otimes f)=\beta(ab,f),\qquad Q(a)(b\otimes g)=\gamma(ab,g)
 \quad(a,b\in B)
\]
を作る。双方が一般に全射なら窓は両対称である。
少なくとも一方が一般に非全射なら元代数は指数成長する。
特に長さ1の点を持つ反対階数9の全有限点窓で全次数の正則性は有限判定できる。
\end{theorem}
核元を消すことは、二つの $S^2 B$ への制限で四つの二項変数の全てが対称であることに等しい。
実際 $S^2(S^2 B)\to S^4 B$ の積の核は判別式の一次元である。
また関係に純テンソル $u f$ があれば、隣り合う転置窓で $u f u$ が三次重なりに入り、
その三テンソルに対応する非零点で辺の評価階数が1となり有限性に反する。
従って $S^2 B$ 上の $\beta,\gamma$ の値は各々六次元標的を張る。
実際、張らなければそれを消す汎関数の関係は $e$ 成分だけにあり純関係となる。

\begin{proof}[両方が全射となる方向]
一般の二項基底 $x,y$ で $P(x),Q(x)$ は全射とする。四二項位置の対称性から
前節の二行列分解と同じく $P(y)=TP(x),\ Q(y)=T^tQ(x)$ である。
従って $x^2,xy,y^2$ で評価した写像は、それぞれ
$U,TU,T^2U$ および $V,T^tV,(T^t)^2 V$ である。
ここで $U:F\to H,\ V:G\to H^*$ は単射である。実際、$[x^2]$ は有限点の外に選べる。
さらに $U,TU$ の像は合わせて $H$ を張り、$V,T^tV$ の像は $H^*$ を張る。
$U',V'$ を $e$ での評価とすれば、偶位置対称性により
\[
 (T^jU')^t V=(T^j U)^t V'\qquad(j=0,1,2).
\]
これらの式を使って $S(Uf+TUf')=U'f+TU'f'$ と置く。
左辺の引数が零なら右辺は $V,T^t V$ を消すので零であり、$S$ は良定義である。
また $S^t V=V',\ S^t T^t V=T^t V'$ も同じ式による。
$C=[S,T]$ とすると $CU=0,\ C^t V=0$ である。
各 $x^2,xy,y^2,e$ での評価は、以後 $M U,M^t V$ と書ける。
ここで $M$ は順に $I,T,T^2,S$ である。

$U$ の像の余次元から $\operatorname{rank} C\le2$ だが、実は等しい。
$h\in\ker C^t$ ごとに $M\mapsto M^t h$ は $P(\beta)$ の対称三次重なりを与える。
実際 $h^t[M,N]U=0$ を調べればよく、唯一残る $[S,T^2]=CT+TC$ でも $h^t C=0$, $CU=0$ である
（ここで $h^t$ は行汎関数を表す）。$I$ での評価によりこれらは独立で、
三次重なりは四次元である。従って
\[
 \operatorname{rank} C=2,\qquad U(F)=\ker C,\quad V(G)=\ker C^t .
\]
各固有値 $\lambda$ で $T$ の固有空間と $U(F)$ の交叉は零である。
交叉元 $Uf\ne0$ があれば、$p(\lambda)=0$ を満たす全ての二次以下の多項式で
$p(T)Uf=0$ となり、偶文字の射影直線が辺に射影される。転置についても同様である。
$C=A D$ を $A:k^2\hookrightarrow H,\ D:H\twoheadrightarrow k^2$ と階数分解する。
同じ固有値の左右固有ベクトル $s,r$ には $sCr=s[S,T]r=0$ が成り立つ。
もし固有空間が二次元以上なら、$D$ のその上での単射性と左固有空間上での $A^t$ の単射性により
$sA,Dr$ はそれぞれ二次元全部を走り、矛盾する。従って $T$ は巡回である。

六次の最小多項式 $p$ により $H=k[t]/p$ とし、最高次係数汎関数 $\ell$ の完全積対合を使う。
$C(s)=u_0\ell(v_1 s)-u_1\ell(v_0s)$ と書く。
$\operatorname{tr}(T^j [S,T])=0$ より、全 $j$ に対して $\ell(t^j(u_0v_1-u_1v_0))=0$ であり、
$u_0v_1=u_1v_0$ となる。各根で $u$ の二元がともに零なら、その根での評価が左固有汎関数で
$C$ の像を消してしまう。$v$ についても、二元がともに零ならその局所因子のsocleは
$T$ の右固有空間で $C$ に殺される。いずれも先の交叉零に反する。
従って $v_i=h u_i$ ($i=0,1$) となる単元 $h\in H^\times$ がある。
完全対称対合 $\ell(h a b)$ では $T$ は自己随伴、$C$ は反自己随伴である。
この随伴を $\dagger$ とすると $[S^\dagger,T]=C$、ゆえに $S-S^\dagger$ は巡回な $T$ と交換し、
$T$ の多項式で自己随伴となる。反自己随伴でもあるため $S=S^\dagger$ である。
最後に $\ker C$ と $\ker C^t$ はこの対合で対応し、$U,V$ で二つの奇位置を同定できる。
全 $M$ が自己随伴なので得られる窓はGram表示の両対称窓である。
\end{proof}

\begin{proof}[非全射の鉛筆]
必要なら奇位置を交換し、$P$ が一般に非全射とする。一般階数は5である。
実際4以下なら各 $a$ に対して $\mathbb P B\times\mathbb P F$ 上の高々四式の共通零点があり、
有限個の辺点では全 $a$ を覆えない（二項二次形式の一次因子の個数は有限）。
原始左核 $l(a)$ の次数は五次以下である。$\beta(ab,-)$ の一般階数は4で、
$ab$ の評価の像は $l(a),l(b)$ を消す。この二元が一般に比例すれば $l$ は定数となり
六次元を張ることに反する。最大四次小行列式の相補座標の次数は $(4,4)$ で、
退化が有限なので共通因子を持たない。従ってこれらは $l(a)\wedge l(b)$ の原始化である。
対角 $a=b$ が共通因子なので $\deg l\ge5$、ゆえに等しい。
共通右核を割った列数は5である。もし6以上なら、独立な入力から一般階数5の六列を選び、
随伴行列を五次の $l$ でくくると定数右核が生じて矛盾する。
残る六行五列の鉛筆の最大階数は全点で5である（最大五次小行列式は $l$ の非零定数倍）。
完全列 $5\mathcal O(-1)\to6\mathcal O\to\mathcal O(5)$ により
$l$ は有理正規五次曲線であり、$H=S^5B$ と取ると $l(a)$ の核は五次式のうち $a$ の倍数全体である。
従って $\beta(ab,f)$ は $a,b$ 両方の倍数である。普遍一次式の互いに素な二因子で割って
$\beta(z,f)=zf$ ($z\in S^2 B,\ f\in F=S^3 B$) を得る。
商の $F\to S^3 B$ は単射である。核があれば固定した奇点の相手が射影平面全部を含む。

$D:S^3 B\to S^5 B$ を $\beta(e,-)$ と書き、基底 $x,y$ で
\[
 K=\langle xD(yu)-yD(xu):u\in S^2 B\rangle\subset S^6 B
\]
と置く。$G$ は $S^7 B$ 上で $BK$ を消す全汎関数と同定できる。
実際、四二項位置の対称性と $zf$ の五変数の全対称性から、各 $g$ で
$w(z,f,z',g)=L_g(z f z')$ となる七次の汎関数がある。
これは五変数と四変数の置換が重なった二位置を介して七変数の対称群を生成することによる。
$e$ との対称性は $\gamma(e,g)(zf)=L_g(z D f)$ である。
$S^2 B\otimes S^3 B\to S^5 B$ の乗法核は $xv\otimes yu-yv\otimes xu$ ($v\in B,u\in S^2 B$)
で張られる。各全次数の単項式の隣り合う分解の差を取ればよい。
従って $\gamma(e,g)$ の存在・一意性は $L_g(BK)=0$ と同値である。
逆にその任意の $L$ は同じ式から全偶位置対称の三次重なりを定める（$e,e$ の条件は自動）。
三次重なりの次元は4であり、その中に元の $G$ が単射で入るから主張が従う。
特に $\dim BK=4$、$\dim K\ge2$ である。

一次縮約 $J:B\otimes G\to (S^6 B)^*$ を $J(a,g)(s)=L_g(as)$ とする。
全像は $K$ を消すので階数は $7-\dim K\le5$ である。
両端対称な $B\otimes G\otimes B$ の部分（次元12）において、
左から $J$ で縮約して零となる部分の次元は少なくとも $12-2\operatorname{rank}J\ge2$ である。
対称性から右の縮約でも零となる。
その非零元 $h$ を固定する。
偶位置を $S^2 B$ に制限した経路の双対テンソルを作る。
各内部の $E,G,E$ 三位置には $h$ を一つ置き、各 $E,F$ 辺には
二つの二項変数の角括弧を一つずつ掛ける。
例えば両端を $F$ とする長さ $4s+1$ の列を使えば、$s$ 個の三位置は互いに素で、
偶位置は一次と括弧の一次で丁度二次になる。
各 $F$ 位置は三次まで任意の残りの因子を掛けられる。
全 $F$ 辺で括弧により乗法は零で、$G$ 辺でも $J$ の零条件が、偶位置に
さらに掛けた一次因子によらず $\gamma(z,g)(s')=L_g(zs')$ を消す。
よって得たテンソルは全関係を消す。固定した括弧積と互いに異なる位置の $h$ の積は非零である。
各内部 $F$ の残る一次を独立に選んだテンソルは独立である
（$G$ の係数を評価して各 $h$ を非零の二項多項式にしてから整域で固定積を消去してよい）。
従って次元は少なくとも $2^{s-1}$ であり、指数成長する。
\end{proof}

\begin{corollary}\label{cor:rank-nine-simple-ternary-equivalence}
反対階数9の有限点グラフに長さ1の局所点があるとする。標準有限Frobenius検査を通った窓が
全次数で正則であることは、いずれかの点で中央の支持が3であることと同値である。
特に全正則例は三文字部分代数の既出の表示を持つ。
\end{corollary}
\begin{proof}
中支持3の十分性は定理\ref{thm:finite-cubic-separators}である。全中支持が4なら
定理\ref{thm:rank-nine-finite-symmetry}で片側対称な表示があり、核の階数は3か4である。
この二つの階数では前定理と定理\ref{thm:smooth-symmetric-kernel-all}から
両対称または指数成長である。両対称な標準AS窓の反対階数は6か8なので
（定理\ref{thm:bisymmetric-six-or-eight}）、正則例はない。
\end{proof}

% !TeX root = ../papers/nc-p3.tex
\section{交代・対称な印付き窓と可換な長さ四の環}
\label{sec:mixed-point-frobenius-algebras}
有限点枝の印付き表示で $g=q=3$ のとき、
定理\ref{thm:rank-eight-finite-symmetry}の交換子の証明のうち
反対階数8を使うのは混合型を除外する段のみである。以下ではこの混合型そのものを分類する。
$U,V$ を三次元の標準内積空間とし、$V$ に外積を選ぶ。
記号の向きを固定して $G=F$ を交代側、$Q=P$ を対称側に置く。逆の配置では開始位置を移す。
\begin{theorem}\label{thm:mixed-point-frobenius-algebras}
混合型の印付きデータは、四次元可換単位的代数 $K$ と線型汎関数 $\tau$ で
$\tau(1)=1$, $\tau(ab)$ が完全であるもの、および $U=\ker\tau$ 上の
相似までの非退化対称形式 $m$ に対応する。印を保つ添字保持同型の必要十分条件は
$\tau$ と $[m]$ を保つ単位代数同型の存在である。
$U$ の内積を $x\cdot z=-\tau(xz)$ とし、
$m(x,z)=(Lx)\cdot(Lz)$ となる同型 $L:U\to V$ を選ぶ。
各局所因子のsocleの生成元 $d\in K$ について
\[
      y_d=(L^t)^{-1}\big(d-\tau(d)1\big)\in V
\]
を作る。元の窓が標準Koszul左右AS正則であることと、全ての $d$ について
$y_d\cdot y_d\ne0$ であることは同値である。局所因子は非被約でもよい。
特に四辺が有限な混合型候補は全て正則で、反対階数10、最小周期2である。
\end{theorem}

\begin{proof}[長さ四の環と二周期表示]
$G_a=[La]_\times$ と置く。三独立な交代行列を持つので $L$ は可逆である。
$n_b=L^t b$ と置けば印付きの交換子条件は
\[
 [Q_b,Q_d]=n_b n_d^t-n_d n_b^t
\]
となる。Jacobi恒等式から $Q_b n_d-Q_d n_b=h_b n_d-h_d n_b$ となる線型形式 $h$ が一意にある。
実際、$b,d$ を添字対 $23,31,12$ として左辺を列に持つ行列を $J$、$n_i$ を列に持つ行列を $N$ とする。
上の交換子をJacobi恒等式に代入すれば $JN^t+NJ^t=0$ である。
三行の交代行列 $N^{-1}J$ はちょうど表示した $h$ で一意に書ける。
$B_{n_b}=Q_b-h_b I$ とおくと $B_z x=B_x z$ であり、$K=k1\oplus U$ に積
\[
 (r,x)(s,z)=(rs-x\cdot z,\ rz+sx+B_z x)
\]
を入れられる。各 $z$ の乗法行列
$\bigl(\begin{smallmatrix}0&-z^t\\ z&B_z\end{smallmatrix}\bigr)$ は互いに交換する。
非対角ブロックは上の対称性、右下は交換子式そのものである。
単位的可換な積で左乗法が互いに交換すれば結合律が従う。
スカラー部 $\tau$ の積対合は $1\oplus(-I_U)$ で完全である。
逆にこのような $(K,\tau)$ から $B_z$ は $U$ 上の乗法の $U$ 成分として復元され、
$Q_b$ は $B_{n_b}$ の跡零部分である。従って元の交換子式も成立する。

位置を偶・奇に折り畳むには、二位置先の偶の補文字 $r$ の符号を反転し、奇文字はそのまま用いる。
さらに奇の補座標を $s'=s+h_y$ と書き直す。隣接写像の六成分は、二つの順で
\[
           r y\mp (Lx)\times y,\qquad s'x+B_{n_y} x
\]
である。実際、最初の三成分は印付き四テンソルの $c=*,d=j$、残りは $c=i,d=*$ の列である。
反対向きの窓に移ると $G^t=-G$、$Q^t=Q$ が二つの符号を与える。
従って関係にはこの二組の行空間を取ればよい。
印付き同型公式で $F=G,P=Q$ の各三基底を保つには
$A^tA=B^t B=sI$ が必要十分である（ここだけ $B$ は奇位置の基底変換を表す）。
$A=a O,\ B=a P$、$O,P$ を直交変換とすれば
$L'=a\det(P)P L O^t$, $Q'_y=a O Q_{P^ty}O^t$ である。
従って $U$ 上の $\det(P)O$ が単位を固定して $(K,\tau)$ を対応させ、
$m$ は倍率 $a^2$ まで保たれる。逆にそのような代数同型を $O$ とし、
$L'=a P L O^t$ と書く。平方根 $a$ の符号で $\det P=1$ とできる。
$Q$ は上記の復元で対応するため、同じ印付き変換を得る。
\end{proof}

\begin{proof}[二次双対の構成と必要性]
$M=k\oplus V$ を積
$(r,u)(s,v)=(rs-u\cdot v,rv+su-u\times v)$ の分裂四元数環とする。
標数零の代数閉体上で $M\simeq M_2(k)$ であり、スカラー部は半跡、純成分の行列式は $v\cdot v$ である。
$D=M\oplus K$ で $\ell(m,k)=\operatorname{tr}(m)/2-\tau(k)$,
$H=\ker\ell$, $z=(1,1)$ と置く。点座標の二つの部分空間を
\[
 E=\{(r+Lx,x):r\in k,x\in U\},\qquad
 F=\{(y,s'+n_y):y\in V,s'\in k\}
\]
とする。$\ell(ab)$ は非退化で $F=E^\perp$。$E F=F E=H$（積の線型包）である。
積のスカラー部は両環で $-(Lx)\cdot y$、残る六成分は先の二周期写像である。
六成分の独立性は $r,s'$ の係数から、スカラー部の独立性も同じ二係数とその非零性から従う。
従って $H/kz$ への積が二次関係の評価そのものである。

どの一次元イデアルも一つの局所因子のsocleであり、その生成元 $d$ は $\tau(d)\ne0$ である。
これは有限Frobenius環の完全対合と $dK=kd$ による。$F$ の $K$ 射影は同型で、
$d$ の持上げ $f$ の行列成分は $y_d$ である。$y_d$ が非可逆なら $fD$ は次元3以下で
（行列側が高々2、可換側が1）、$\ell$ は $fD$ で非零である。
ゆえに $fE\subset fD\cap H$ の次元は2以下で、二次評価の階数も2以下となる。
片側対称表示に戻して定理\ref{thm:single-opposite-evaluation-ranks}を適用すれば非正則である。
また幾何学的にはこの評価に射影直線以上のファイバーがある。
\end{proof}

\begin{proof}[非被約な環を含む十分性]
以下全てのsocle持上げが可逆行列成分を持つとする。
$D$ の次数別コピーに $E,F$ を交互に掛けた部分からなる添字付き代数を $C$ とする。
各一次空間は $D$ の単元を稠密に含む。
$E$ の $K$ 成分 $U$ はどの極大イデアルとも異なる。もし一致すれば $\tau$ はそのイデアルを消し、
完全積対合に反する。$F$ では $K$ への射影が同型、行列成分の三平面にも単元がある。
$H$ の右安定化単元はスカラーのみである（完全対合で $H^\perp=k1$）。
従って $H$ に各次位置の一次空間の独立な二単元を掛けた二超平面は異なる。
長さ3の積は $D$ であり、以後も $D$ である。

$K$ のイデアル $I$ と $k^2$ の部分空間 $N$ に対し、$J_{I,N}$ を
$C$ の正次数成分と右イデアル $\{(X,u):\operatorname{im}X\subset N,\ u\in I\}$ の次数別交叉とする。
両始点で同じ記号を使う。二次元イデアル $I$ を任意に取り、$I'=\operatorname{ann}_K I$ とする。
$\dim I'=2$, $\operatorname{ann} I'=I$ は完全対合から従う。$\tau$ はどの非零イデアル上でも非零である。
$E,F$ のうち $K$ 成分を $I$ に制限した部分はともに二次元で、行列成分への射影は単射である。
$E$ 側は単位行列を含み、$F$ 側は純四元数の二平面で非退化三変数二次形式の全等方平面ではない。
従って行列式が零の射影根 $b$ があり、その行列 $X$ は階数1である。
この $b$ の $K$ 成分は $I$ を生成する。実際、生成しなければ積イデアルは零か一次元である。
$E$ 側の零成分は単位行列の倍、一次元イデアルは $U$ と交わらない。
$F$ 側は $K$ 射影が単射で、一次元イデアルでは仮定により行列が可逆である。
この議論は二重根と $I=I'$ も含む。

次位置の $I'$ に制限した行列二平面は、$I$ に制限した元位置の二平面の行列跡直交補である。
実際、$I I'=0$ と $F=E^\perp$、両二平面の次元による。
このうち像が $\ker X$ にある行列は一次元であり、そこに次の根 $c$ がある。
実際、今の二平面で $X$ と独立な可逆 $X'$ を取る。$\operatorname{im}Y\subset\ker X$ の二平面では
$\operatorname{tr}(XY)=0$ は自動、$\operatorname{tr}(X'Y)=0$ は非零の一式である。
従って $b$ の右零化空間は次の一次空間で $kc$ であり、二次以後では $D$ の右零化イデアルとの交叉である。
$b$ と次位置の一次空間の積は三次元、従って $J_{I,\operatorname{im}X}$ の二次成分全部である。
また $bH=bD$ である。その右零化イデアルに $I'$ の元で $\tau$ が非零になるものがあるからである。
一次でもこの行列像条件の解は $kb$ である（同じ二平面に可逆 $X'$ がある）。
ゆえに $bC=J_{I,\operatorname{im}X}$、その右零化イデアルは次位置の $cC$ となる。

ここからの線型イデアルの構成を記す。$I$ に制限した二平面で $X'$ に対応する元 $b'$ を加える。
$bC$ に対する $b'$ の右コロンは $J_{K,N}$, $N=(X')^{-1}\operatorname{im}X$ である。
このイデアルの次元は長さ1,2,3以後で $2,5,6$ である。一次では $E\to M$ は同型、
$F\to M$ は核一次元で像が跡零三平面、$\operatorname{im}Y\subset N$ の行列のうち跡零のものは一次元だからである。
従って短完全列で $(b,b')C=J_{I,k^2}$ が従う。右辺の次元も長さ1,2,3以後で $2,5,6$ であり、
$bC$ とコロン商の次元を足すとこれに達する。

この $J_{K,N}$ も一次生成元の追加で得られる。$X,X'$ の $N$ への制限は比例するので、
像が $N$ に入る行列二平面では二つの跡条件は一つの非零条件である。
従って次位置に $K$ 成分が $I'$ を生成し、行列像が $N$ の根 $c'$ がある。
同じ行列像条件の二平面で、$K$ 成分が単元の独立な $c''$ を取れる。
$F$ 側では $(0,1)$ が使える。$E$ 側では各剰余体への写像は、
反対位置の該当socleの持上げとの行列跡対合の非零定数倍である。
実際、$\ell(EF)=0$ と $\tau(kd)=\tau(d)\bar k$ である。
その行列は可逆なので、像が $N$ に入る行列全体ではこの対合は恒等零でない。
有限個の極大イデアルを避けられることが従う。
$c'C$ に対する $c''$ の右コロンは $J_{I',k^2}$ であり、同じ次元比較で
$(c',c'')C=J_{K,N}$ を得る。

最後に $I$ を二次元から一次元ずつ三次元のイデアル、さらに $K$ へ増やす。
各包含 $I_j\subset I_{j+1}$ に対し、$K$ 成分が $I_{j+1}/I_j$ で非零の一次生成元を追加する。
$U\cap I_j$ は $I_j$ の余次元1であり両始点でこの選択ができる。
右コロンはこの単純商を消す極大イデアル $\mathfrak m$ に対する $J_{\mathfrak m,k^2}$ である。
その商は各次数一次元であり、得られるイデアルは $J_{I_{j+1},k^2}$ に等しい。
全ての三次元イデアルが二次元イデアルを含むことは組成列から従う。
以上のイデアルを両始点で一緒に取れば、零から各イデアルへ一つずつ線型生成元を足す際の
コロンも同じ集合に属する。従って定理\ref{thm:half-period-semisimple}と同じ短完全列の
二重帰納で $C$ はKoszulである。この証明で $K$ の根基を割る必要はない。

$z$ は $C$ の両始点の二次に属する中心非零因子である。
従って $C/(z)$ はKoszul、各始点のHilbert関数は $1,4,6,4,1$ である。
最上次数の汎関数 $D/H\to k$ は $\ell$ で、$E,F$ の完全直交と $H^\perp=kz$ により
全ての相補次数で積対合は完全である。
二次の積はちょうど所与の六成分なので二次双対性で元の窓の標準Koszul左右AS正則性を得る。
\end{proof}

\begin{proof}[有限枝の階数と周期]
混合型の残余対合は定理\ref{thm:rank-eight-finite-symmetry}の式で非零なので階数は少なくとも9、
片側対称なので高々10である。階数9の場合でも全点で中支持4を仮定する必要はない。
この混合型の対称核は階数4であることを直接確かめる。
$D_G=-\operatorname{cof}(L)$ として、残余対合が階数1なら、
$I,Q_b$ と跡直交する対称行列の二平面中に、$D_G^tD_G$ の非零倍 $Y$ がある。
実際、残余対合の核は $D_G^{-t}Y D_G^{-1}$ がスカラーとなる条件そのものである。
従って $Y$ は $(L^tL)^{-1}$ の定数倍であり、後者に規格化する。
偶文字対の純成分に $Y$ を縮約すると $G$ 側の寄与は $-2I$ である
（外積積公式）。$Q$ 側は補題\ref{lem:ternary-net-commutator}から
$D_Q P_Q(Y)D_Q^t=(\det D_Q)Y$ で、$D_Q=-D_G^t=\operatorname{cof}(L)^t$ を入れると
$P_Q(Y)=I$ である。星印一個の成分は $\operatorname{tr}(YQ_b)=0$ で消え、
両星印側の $\operatorname{tr}Y$ も零である。ゆえに対称核の生成元はブロック対角 $1\oplus Y$ で階数4である。
有限枝では定理\ref{thm:smooth-symmetric-kernel-all}と両対称AS階数二分法により階数9は不可能である。
残る階数10で周期1なら偶位置の交換を一文字ずらして奇位置の交換も得るので両対称となり再び矛盾する。
既に二周期表示があるため最小周期は2である。
有限枝で印を忘れた同型は、相手の高々二十個の局所点のうち対応する混合型の単点を順に印として
上の $(K,\tau,[m])$ を比較すれば必要十分である。
\end{proof}

% !TeX root = ../papers/nc-p3.tex
\section{混合型の正規形、点の有限性と半周期成分への特殊化}
\label{sec:mixed-point-orbits}
\begin{corollary}\label{cor:all-defective-one-sym-finite}
片側対称な標準有限Frobenius窓の四辺が有限で、反対階数が9以下なら、
長さ1の点がなくても全次数の正則性は有限に決まる。
対称核に階数2以下の非零元があれば正則である。それ以外は両対称な正則窓に限る。
\end{corollary}
\begin{proof}
低階数元から中支持3を得る定理\ref{thm:rank-nine-finite-symmetry}の末尾の証明は核の次元を使わない。
それ以外では任意の非零核元が階数3または4である。階数3では定理\ref{thm:cone-kernel-binary-transfer}、
階数4では定理\ref{thm:smooth-symmetric-kernel-all}の証明を用いる。
後者で使う核の条件は、選んだ行列式の消滅、すなわち二組の二項位置の各対称性だけである。
いずれも両対称または指数成長となり、前者の判定は定理\ref{thm:bisymmetric-two-length-test}である。
\end{proof}

\begin{proposition}\label{prop:mixed-quartic-normal-forms}
前節の混合型で正則な全パラメータは、点グラフの次元によらず反対階数10、最小周期2である。
印付き正規形の可換環は $4=n_1+\cdots+n_l$ に対する
$\prod_i k[t_i]/t_i^{n_i}$ または $k[u,v]/(u^2,v^2)$ である。
\begin{enumerate}
\item 前者では各局所単位の値 $\omega_i$ の和は1であり、$n_i=1$ のとき $\omega_i\ne0$ である。
$n_i>1$ なら $\tau(t_i^{n_i-1})=1$、その間の冪の値は全て零とする。
残余同型は同じ $(n_i,\omega_i)$ の置換と $t_i\mapsto\zeta_i t_i$, $\zeta_i^{n_i-1}=1$ ($n_i>1$) で尽きる。
\item 後者では $\tau(u)=\tau(v)=0$, $\tau(uv)=1$ とする。
残余群は $(u,v)\mapsto(cu,c^{-1}v)$ と $u,v$ の交換が生成する。
\end{enumerate}
いずれも $\ker\tau$ 上の射影的な可逆対称三行列をこの残余群の合同作用で比較すれば必要十分である。
\end{proposition}
\begin{proof}
非被約な局所Frobenius環で接空間が一次元なら一文字の打切り環である。
接空間が二次元以上ならsocleの次元1と長さ4からHilbert関数は $1,2,1$ で、
中央二次元の積はsocleに値を持つ非退化対称形式となる。
二つの等方方向を取れば第二の環を得る。三次元接空間は平方零の三次元socleとなり不可能である。
一文字の環では最上冪の値は非零である。文字の倍率でそれを1にしてから
$t\mapsto t+c_jt^j$ ($j=2,\ldots,n-1$) で順に $\tau(t^{n-j})=0$ とする。
実際、変化量は $(n-j)c_j$ で既に揃えた高次の値は保たれる。
従ってこの表示の自己同型は表示した根の倍率だけである。
第二の場合は $uv$ の倍率を揃え、$u,v$ へsocle元を加えて各値を零とする。
自己同型は二つの等方方向を保ち、跡の条件がsocleへの追加を禁じる。

階数について、前節の残余対合が階数1となるなら
$Y=(L^t L)^{-1}$ は $I,Q_b$ 全てと跡直交する。
$U$ の直交基底 $e_i$ を使い $a=\sum Y_{ij} e_i e_j\in K$ とする。
$\tau(a)=-\operatorname{tr}Y=0$ であり、$a$ の各 $z\in U$ との積跡は
$-\operatorname{tr}(Y B_z)=0$ である。従って $a=0$ となる。
一つの局所因子の剰余体で評価し、その指標を $\chi$、socleの生成元を $d$ とする。
$e_i\cdot(d-\tau(d)1)=-\tau(d)\chi(e_i)$ なので
$y_d\cdot y_d=\tau(d)^2\chi(a)=0$ である。正則性条件に反する。
残余対合の階数0も既に排除したので全正則点の反対階数は10である。
周期1なら二つの反対位置で対称となり、両対称な標準AS窓の階数制限に反する。
\end{proof}

\begin{proposition}\label{prop:mixed-finite-eigenvectors}
正則な混合型で四辺が有限となる必要十分条件は、$K$ が一文字局所因子の積であり、
写像 $x\mapsto B_{(L^tL)x}x$ の $\mathbb P U$ 上の固有方向が有限となることである。
後者は $x\wedge B_{(L^tL)x}x$ の三つの三次式に非定数共通因子がないことと同値である
（三式全零の場合を除く）。
\end{proposition}
\begin{proof}
二周期の二つの辺は補文字 $r$ の符号変更と順序反転で対応するので、第一の辺の台を調べればよい。
原点対の外では $x,y$ はともに非零である。$r=0$ なら $y$ は $Lx$ の倍なので
$s'x+B_{n_y}x=0$ はちょうど固有方向の条件になる。
$r\ne0$ なら $y\cdot y=(Lx)\cdot y=0$ が外積式から従う。後者により
$x(s'+n_y)=0$ が $K$ の全積として成り立つ。
逆に $k'=s'+n_y$ が $x\in U$ を消し $y\cdot y=0$ なら $(Lx)\cdot y=0$ であり、
三次元の等方直線の性質から $(Lx)\times y$ は $y$ の倍で、一意な $r$ が取れる。
この逆構成で $r=0$ も許す。

$\operatorname{ann}(k')$ は $U$ と非自明に交わるので次元2以上である。
一次元イデアルは $U$ と交わらないからである。
一方、$k'$ が一次元イデアルを生成するなら $y$ は前節のsocle判定で非等方となる。
従ってこの構成では $I=k'K$ はちょうど二次元、$kx=\operatorname{ann}(I)\cap U$ である。
各二次元イデアル $I$ の $F$ への持上げの行列鉛筆は純四元数二平面であり、
等方根は重複込み二個である。各根が実際に $I$ を生成することも前節で示した。
異なる $I$ は異なる奇点を与える。一文字局所環の積に二次元イデアルは有限個しかなく、
$k[u,v]/(u^2,v^2)$ には $\langle a u+b v,uv\rangle$ の射影直線がある。
これらと $r=0$ の全点で被覆されるので結論を得る。
\end{proof}

\begin{theorem}\label{thm:mixed-half-period-specialization}
全ての正則な混合型の窓は、九次元半周期成分の正則な枠の特殊化である。
有限点を持つ混合型の稠密部分は印を忘れて八次元の既約な軌跡をなし、
その閉包はこの成分で余次元1である。
\end{theorem}
\begin{proof}
各 $(K,\tau)$ を被約な長さ4の可換代数へ平滑化する。一文字局所因子は
$t^n$ を $\prod_{j=1}^n(t-ja)$ に、第二の環は $(u^2,v^2)$ を $(u(u-a),v(v-a))$
に変えればよい。パラメータは $a$ で、単項式基底は一定である。
$\tau$ の基底上の値を固定すれば完全対合は零ファイバーで非退化なのでその近傍でも非退化である。
特に一般被約ファイバーでは全ての原始冪等元で値が非零である。
$\ker\tau$ とその内積は束をなし、$L$ は基底上で延ばせる。随伴は延ばした内積で取ればよい。

この近傍上で $D=M_2\oplus K$ の四次元部分空間 $E'$ をGrassmann束全体に動かし、
$F'=(E')^\perp$ とする。二次積 $E'F',F'E'$ がともに $H$ となる条件は開であり、
指定した混合型点で満たす。両積を $H/kz$ に写した六成分を二周期の関係として取れば、
この点の元の関係へ特殊化する。四階テンソル自体も
$\ell(e f e' f')$ で規格化できる。これは $z=(1,1)$ を割った
相補二次積の対合であり、$H^\perp=kz$ から隣接階数6である。

一般点では $K=k^4$ に分裂し、両四平面の両環因子への射影は同型である。
そのとき偶文字を行列 $X$、奇文字を $\operatorname{adj}Y$ により四次元ずつ同定する。
$\omega_i=\tau(e_i)$ を用いスカラー積の四成分をそれぞれ $\omega_i$ 倍して読むと、
偶側の四座標を $u_i$、奇側を $v_i$（この倍率は片方に入れる）として
$\sum u_i\otimes v_i$ は行列式の偏極となり、$z$ は $(\omega_1,\ldots,\omega_4,I_2)$ となる。
従って定理\ref{thm:half-period-frame}の枠表示である。一般枠では各単一座標の行列式と
その双対側の行列式が非零なので、定理\ref{thm:half-period-semisimple}により正則である。
一般の枠が現れることは、$E'$ が二つの四次元の間の任意の同型のグラフを走り、
$F'$ は積跡直交で一意に定まるためである。必要なら $\tau(1)=1$ を保って $\tau$ も動かせる。

混合型の被約 $K$ では三つの重み自由度と $U$ 上の射影二次形式の五つの自由度があり、
残余群は有限である。この八次元既約パラメータの一般点は正則で四辺が有限である。
実際、$K=k^4$ に等しい重みを取り、$(\mathbb Z/2)^2$ の三つの非自明指標に
$\sqrt{-1}$ を掛けた直交基底を $U$ に取って $L=I$ とすれば、$B_yx=\sqrt{-1}D(x,y)$
（ここで $D$ は定理\ref{thm:twenty-point-distinct-companion}の対称積）である。
補座標を定数倍すれば第一辺はその被約二十点の例になる。有限性・横断性と前節の非等方性は開条件である。
印の選び方が有限なので印なしの像も八次元であり、九次元成分内の閉包の余次元を得る。
非被約入力も同じ八次元混合型の閉包に入る。実際、先の平滑化で混合型表示を保って
形式を一般に動かすと一般ファイバーは正則かつ有限である。
\end{proof}

% !TeX root = ../papers/nc-p3.tex
\section{二対に分かれる対称境界の排除}
\label{sec:two-pair-boundary}
二つの互いに素な位置対を可逆な正証明書で同一視し、テンソルを
$\operatorname{Sym}^2 V\otimes\operatorname{Sym}^2 W$の元として書く。
この表示は記述用であり、元の窓の巡回順序は保つ。
対をまたぐ特異な正証明書$H:W\to V$があるとし、$A=\operatorname{im}H$、
$D=\ker H$と置く。このとき
\[
 w\in\operatorname{Sym}^2 A\otimes\operatorname{Sym}^2 W+
        \operatorname{Sym}^2 V\otimes\operatorname{Sym}^2 D.                 \tag{*}
\]
実際、一つの$W$因子に$H$を作用させたテンソルは三つの$V$位置に対称であり、
その三位置の全てが$A$に支持される。従って片方の$W$文字を
$D^\perp=\operatorname{im}H^*$の元で縮約すると両$V$文字が$A$に入る。
もう一つの$W$位置についても対称性を使えば$(*)$を得る。
一つずつ異なる種類の因子へ縮約した像を$F\subset V\otimes W$と書く。
対ごとの対称性から四つの混合対の像は順序を揃えれば同じ$F$であり、
$(*)$から$F\subset A\otimes W+V\otimes D$である。

\begin{proposition}[交互に並ぶ二対]\label{prop:alternating-position-pairs}
窓の巡回順序が$V,W,V,W$なら、$0<r=\rk H<4$であるどの交差証明書も
標準AS性を排除する。
\end{proposition}
\begin{proof}
全ての隣接関係空間は$F$またはその反転である。一次空間を位置ごとに
$V/A,W/D$へ落とすと、全関係が消える。従って長さ$4m$の商の次元は
$\bigl(r(4-r)\bigr)^{2m}$となる。
\end{proof}

\begin{proposition}[隣接する二対の四次障害]\label{prop:adjacent-position-pairs}
窓の巡回順序が$V,V,W,W$で、$w$が標準有限Frobenius候補、
$H$が階数2の交差証明書であるとする。このとき、連続純粋関係があるか、
$\dim C_{0,4}\ge41$である。従ってどちらの場合も標準AS性を満たさない。
\end{proposition}
\begin{proof}
$\dim A=\dim D=2$である。対を分ける切断の階数は6であり、
$(*)$からその上界は$\dim\operatorname{Sym}^2 A+\dim\operatorname{Sym}^2 D=6$となる。
等号より$R_0\supset\operatorname{Sym}^2 A$、$R_2\supset\operatorname{Sym}^2 D$である。
実際、$\operatorname{Sym}^2 A$で像側を割った行列は$(\operatorname{Sym}^2 D)^*$からの階数3写像となり、
元の像が6次元なので核に当たる3次元も全て像に入る。転置にも同じ議論を適用する。
$R_0\subset\operatorname{Sym}^2 V$, $R_2\subset\operatorname{Sym}^2 W$でもある。

以後、連続純粋関係がないとする。
$F=R_1$のうち$F_A=F\cap(A\otimes W)$または$F_D=F\cap(V\otimes D)$に入る
非零純テンソルは存在しない。例えば$a\otimes w'\in F_A$なら$a\otimes a\in R_0$と連続し、
$F_D$でも後続の平方と連続する。$2\times4$行列空間の射影階数1部分は
$\mathbb P^1\times\mathbb P^3\subset\mathbb P^7$の4次元部分なので、
射影次元定理から$\dim F_A,\dim F_D\le3$である。

一つの$V$位置を$A^\perp$、一つの$W$位置を任意の汎関数で縮約する像を$J_D$とする。
$(*)$から$J_D\subset F_D$である。一方、その次元は混合切断の転置を使えば
$F\to(V/A)\otimes W$の像次元$6-\dim F_A$である。よって
$\dim F_A=\dim F_D=3$、$J_D=F_D$である。
$J_D$の$V$側の支持は$V$全体である。もし三次元以下なら、
$3\times2$行列の三次元部分空間には射影次元定理で純テンソルが入り、矛盾する。

ここで$(F_A\otimes W)\cap(A\otimes R_2)=0$を示す。
交叉の元は$W_{1,3}$に入り、有限Frobenius性から、先頭$V$である位置0に
ある$\phi\in V^*$を作用させた$w(\phi)$と書ける。
残る$V$が$A$に支持されるので、そこを$A^\perp$で、さらに$W$の一位置を任意に縮約すると零となる。
従って$\phi$は$J_D$の$V$支持を消し、$\phi=0$である。
ゆえに$A\otimes W\otimes W$から$C_{1,4}$への像次元は
$32-3\cdot4-2\cdot6=8$以下となる。
四次Koszul複体の$d_2$を$\operatorname{Sym}^2 A\otimes C_{2,4}$に制限すると、
その定義域は30次元で、像は$A$と今の8次元以下の像のテンソル積に入る。
従ってこの部分から少なくとも14次元のサイクルが生じる。

このうち境界は8次元以下であることを確認する。
$R'_0=R_0/\operatorname{Sym}^2 A$とし、$d_3:W_{0,3}\otimes W\to R_0\otimes C_{2,4}$を
$R'_0\otimes C_{2,4}$へ射影する。$(*)$の最初の項を除いたテンソルは、
先の階数3から$(\operatorname{Sym}^2 D)^*\simeq R'_0$という同型を与える。
従って$W_{0,3}\simeq W^*$を使うと、射影した$d_3$は
$W^*\twoheadrightarrow D^*$に続いて
\[
 D^*\otimes W\longrightarrow
  (\operatorname{Sym}^2 D)^*\otimes ((D\otimes W)/\operatorname{Sym}^2 D)
  \ \subset R'_0\otimes C_{2,4}
\]
を作用させる写像である。ここで$R_2\cap(D\otimes W)=\operatorname{Sym}^2D$を使った。
写像は二次テンソル$z^2$で評価すると$\lambda\otimes y\mapsto\lambda(z)[z\otimes y]$で与えられる。
$W/D$方向では単射で階数4、$y\in D$の部分では$D^*\otimes D$の元を$T$として
$z\mapsto z\wedge Tz$となり、核はスカラーのみで階数3である。
二つの像は補空間の選択で直和となり、射影の階数は7である。
命題\ref{prop:quartic-koszul-obstruction}より$d_3$自体の階数は15なので、
$\operatorname{Sym}^2 A\otimes C_{2,4}$内に入る境界の次元は$15-7=8$である。
従って$\dim H_2\ge14-8=6$、$\dim C_{0,4}\ge35+6=41$となる。
\end{proof}

\begin{theorem}[対称閉包のAS残余では可逆対は高々一つ]\label{thm:no-two-pair-AS-boundary}
$Z_{\rm sym}\setminus S$内の標準AS候補では、六つの無順序位置対のうち
可逆な正証明書を持つものは高々一対である。
\end{theorem}
\begin{proof}
位置を共有する二対なら三因子を同時対称化でき、基点排除によりASの場合は$S$に入る。
共有しない二対があるなら、それらをまたぐ可逆証明書も同じ理由で存在しない。
系\ref{cor:two-pairs-transfer-ranks}より両方向で階数2の交差証明書を取れる。
二対の巡回順序は上の二命題で尽くされ、定理\ref{thm:no-consecutive-pure-relations}と合わせて
いずれもAS性に反する。よって主張を得る。
\end{proof}
$Z_{\rm sym}$の外では、向かい合う二対に正証明書を持つ符号付きOreの既知AS例と両立する。

\subsection{二対二の三つの境界因子の実現}
ここでは排除条件を満たす境界が実在することと、その余次元を確認する。
$V=A\oplus B$, $W=C\oplus D$と二次元ずつに分け、
$H:C\simeq A$, $K:B\simeq D$を用いて$A,C$を色$r$、
$B,D$を色$s$の共通二次元空間として扱う。
二対内の対称性とこの二つの交差証明書の条件を満たす模型を次で作る。
$V$側に$s$文字があり同時に$W$側に$r$文字がある係数を零とし、
他の係数には同じ色の文字同士で対称な次の二重次数の形式を入れる。
\[
 \alpha_{4,0},\quad\beta_{3,1},\quad\gamma_{2,2},\quad
 \delta_{1,3},\quad\eta_{0,4}.
\]
添字は二色の個数である。係数は対称テンソルの文字座標で取り、多項式の多項係数を掛けない。
$rr,rs,ss$の対の空間の次元を重複を除いて$3,4,3$とする。
$\alpha,\eta$の二対二行列を$A_0,E_0$、$\beta,\delta$の$rr|rs,rs|ss$行列を$B_0,D_0$とする。
$\gamma$には二つの読み方$G_\gamma$ ($rr|ss$, $3\times3$)、
$C_\gamma$ ($rs|rs$, $4\times4$)がある。後者は対称行列全体の中の超平面
\[
 \sigma(C)=C_{00,11}-C_{01,10}=0
\]
の任意の元を一意に表す。混合切断では$sr$の行・列が零で、$rr,ss$内の反対称成分も零である。
重複を除くと混合切断は
\[
 \begin{pmatrix}A_0&B_0&0\\ B_0^{\mathsf T}&C_\gamma&D_0\\
                 0&D_0^{\mathsf T}&E_0\end{pmatrix}.
\]
$A_0,E_0$可逆のとき、階数6となる条件は
$C_\gamma=B_0^{\mathsf T}A_0^{-1}B_0+D_0E_0^{-1}D_0^{\mathsf T}$である。
右辺を$T$と書く。$\sigma(T)=0$が$\alpha,\beta,\delta,\eta$に対する一条件であり、
成立すれば$\gamma$は一意である。

\begin{theorem}[二対二配置による三つの追加境界因子]\label{thm:three-two-pair-divisors}
四位置の二対への三つの分割ごとに、$Z_{\rm sym}\setminus S$には
その二対だけが一般点で可逆正証明書を持つ68次元既約成分が存在する。
これらと定理\ref{thm:four-boundary-divisors}の四成分は全て異なる。
各成分の一般点は有限Frobenius非Koszulであり、$\mathfrak X_6$のある開近傍で
標準AS点は全て$S$に属する。
\end{theorem}
\begin{proof}
まず具体的な非空開チャートを定める。二進添字の和の昇順で
\[
\begin{aligned}
 \alpha&=(1,0,2,1,3), &\eta&=(2,1,3,0,5),\\
 \beta&=((2,3,5,7),(11,13,17,19)), &\delta&=((2,1,0,3),\,\lambda(2,4,1,-1)).
\end{aligned}
\]
を対照点に使う。$\delta$の第二列ベクトルのスカラー$\lambda$に関し$\sigma(T)$は一次である。
実際、$\sigma$の二つの成分はいずれも異なる$r$文字から一つずつ取る。
対照点での傾きは$-54$、根は$\lambda=-23/162$となる。
一般にもこの第二ベクトルの一つの非零座標をスカラーとして分離し、
その方向を射影チャートで変化させれば、傾きが非零な開部分で一意に解ける。
よって$5+8+8+5-1=25$次元の既約な係数チャートを得る。
$\gamma$を加えた五形式全体に一つのスカラーを掛ける自由を除けば24次元である。

このチャートの有限Frobenius部分が$Z_{\rm sym}$に入ることを示す。
完全対称四テンソルの二重次数成分を順に
$\alpha,\beta,\gamma(t),t^{-1}\delta(t),t^{-2}\eta$とする。
その中間縮約の$ss$の行・列に$t$を掛けると、上のブロック表示の二つの零を
$tG_{\gamma(t)}$とその転置で置き換えたものになる。
階数6を保つSchur補の方程式は
\[
 C_{\gamma(t)}=
 (B_0^{\mathsf T}\ \ D_0(t))
 \begin{pmatrix}A_0&tG_{\gamma(t)}\\tG_{\gamma(t)}^{\mathsf T}&E_0\end{pmatrix}^{-1}
 \binom{B_0}{D_0(t)^{\mathsf T}} .
\]
$t=0$での未知数を$\gamma$の9座標と上の$\lambda$とする。
対称$4\times4$行列への線型化は、$\gamma$から$\ker\sigma$へ同型であり、
残る一方向は選んだ非零の傾きとなる。従ってJacobianは可逆で、
次数ごとにこの十元連立一次方程式を解いて形式的な持ち上げを得る。
最初の二位置の$s$文字を$t$倍すれば全係数は$k[[t]]$に入り、特殊化が指定した模型になる。
実際、$s$文字をできるだけ後の二位置に置く成分だけが最小付値を持つ。
一般ファイバーは完全対称テンソルの四因子基底変更で、三分割の階数は6以下である。
有限Frobenius条件はこの三階数切断内で開なので、特殊点がそれを満たせば一般ファイバーも満たす。
完全対称な有限Frobeniusテンソルの認識定理より一般ファイバーは$S$に属する。
従って所与の特殊点は$Z_{\rm sym}$に入る。

この対照点では、一因子階数は全て4、三切断の階数は全て6である。
隣接二対の順序でも、一対を交互に配置する順序でも、四始点全てで
$(\dim W_{i,3},\dim W_{i,4})=(4,1)$である。
全12個の向き付き正symmetrizer空間は一次元であり、対の内部で階数4、対をまたぐと階数2である。
従ってこれらの次元1条件は上の族の非空開部分で保たれる。
三切断の射影接次元の有理数検算は69である。
数値による近似は使わず、\path{computations/verify_two_pair_boundary.py}で関係挿入の核、
行列の階数とSchur方程式の傾きを計算する。

位置のラベルを固定して表示の自由度を数える。二対の各内部で一方の因子の基底を合わせる
射影変換が二つで30次元。共通空間$V,W$内の二平面$A,D$の選択が$4+4$次元、
射影同型$W/D\simeq A$, $V/A\simeq D$の選択が$3+3$次元である。
この14次元の基底データのGrassmann・行列チャート内で補空間と基底を正則に選べる。
係数を合わせて$30+14+24=68$次元の既約な表示族となる。
一般点から、二対内の一次元可逆証明書で二対をそれぞれ正規化し、
両向きの交差証明書の射影類を回収すれば$A,D$と二つの射影同型が決まる。
固定した局所基底では五形式も共通スカラー以外一意である。よって一般ファイバーは零次元である。
この族は$S$と交わらない一般部分を持つので、その68次元閉包は69次元成分の境界因子である。

四位置を置換すると他の二分割の因子も得られる。巡回順序を変えた場合の有限Frobenius性は
上の二順序の検算と開性により確かめられる。一般点で可逆な位置対が指定した二対だけであることから、
それら同士および三因子対称の四つの境界因子との区別がつく。
対照点で$\mathfrak X_6$は69次元成分を含み接次元も69なので滑らかであり、近傍を
$Z_{\rm sym}$のみに取れる。その中で二対の可逆証明書は、次元1と可逆性の開性で残る。
定理\ref{thm:no-two-pair-AS-boundary}より、この近傍にある標準AS点は$S$に属する。
各因子の稠密開部分について主張が従う。七つで全境界を尽くすかは未決定である。
\end{proof}

対照点の隣接配置では$F_A,F_D$の二行または二列の行列束を調べると、
どちらも四つの最大minorの最大公約式が1であり、純テンソルを含まない。
$F\subset A\otimes W+V\otimes D$内の純テンソルはどちらかに入るため、混合関係には純テンソルがない。
このとき四次経路次元は$(41,42,41,42)$で、隣接配置の下界41は始点0で達成される。
同じ係数を交互に配置すると四次経路次元は全始点で35だが、自由な交互商による指数成長がある。

\begin{corollary}[二対二境界の局所方程式]\label{cor:two-pair-transverse-coordinate}
上の各対照点の近くで、対をまたぐ両向きの正証明書直線の局所生成元を$H,K$と取ると、
$HK=KH=zI$のスカラー$z$が滑らかな被約境界の局所方程式である。
両行列式はともにこの因子に沿って二重に消える。
\end{corollary}
\begin{proof}
二対の内部を同一視した局所座標を使う。
$Z_{\rm sym}$の稠密開部分$S$の上で$HK,KH$は同じスカラー行列であり、
この等式は正則な局所生成元に対して近傍全体に延びる。
対内の可逆性も保たれるため、$z\ne0$は四因子全てを対称化できることと同値である。
前定理の形式曲線では最初の二位置の色$s$を$t$倍するので、
$H,K$は単元を除いて$\operatorname{diag}(1,1,t,t)$と
$\operatorname{diag}(t,t,1,1)$である。従って$z$は曲線上で単元倍の$t$となり、
滑らかな周囲で$dz\ne0$を得る。$z=0$で両行列は階数2
（階数2以上は開であり、積が零）なので、単元二次minorによる掃き出しで
両行列式が$z^2$で割り切れる。積は$z^4$であり、両方の消滅次数は正確に2である。
\end{proof}

\begin{corollary}[二対を持つAS点の交差証明書空間]\label{cor:two-pair-cross-space-zero}
$Z_{\rm sym}$への所属を仮定せず、標準AS対象が互いに素な二つの位置対で
可逆な正証明書を持つとする。$S$に属さなければ、異なる対の位置間の
正symmetrizer空間は両向きとも零である。
\end{corollary}
\begin{proof}
非零の交差証明書$H$が可逆なら、二対の同一視と合わせて四因子を対称化できる。
階数1は系\ref{cor:no-rank-one-into-positive-pair}で排除され、
階数2はいずれの巡回配置でも上の二命題で排除される。
階数3では$H$を$\operatorname{diag}(I_3,0)$に取る。対象二因子について、
残る因子の全縮約行列$M$は$MH^{\mathsf T}$が対称であるという条件を満たす。
従って$M$の左下$1\times3$ブロックは零である。
逆向きに$K=\operatorname{diag}(0,0,0,1)$を取れば、$KM$は右下の一要素しか持たず対称である。
よって$K$は逆向きの階数1の正証明書となり、その標的位置も可逆対に入っているので矛盾する。
\end{proof}

\section{対称閉包での隣接正証明書の伝播}
\label{sec:adjacent-witness-boundary}
\begin{theorem}[対称閉包内での一つの隣接正証明書]\label{thm:one-adjacent-AS-in-symmetric-closure}
$Z_{\rm sym}$に属する標準AS対象が一つの隣接可逆正証明書を持てば、その対象は$S$に属する。
\end{theorem}
\begin{proof}
巡回順序を$V,V,W,U$とし、最初の二位置で$w$を対称にする。
$w\notin S$と仮定する。外側の位置とこの対の間に可逆証明書があれば、
三因子対称のAS認識から矛盾する。$S$の上の両向きの同一視のグラフを閉じて、
各$T=W,U$について非零の正証明書$H_T:T\to V$, $K_T:V\to T$を
$H_TK_T=K_TH_T=0$となるように取れる
（系\ref{cor:two-pairs-transfer-ranks}の固有射影の議論）。
どちらの$T$も$V$二位置の一方とは向かい合い、二位置の交換で$w$は不変なので、
系\ref{cor:opposite-singular-rank-two}により四行列全てが階数2である。
\[
 A_T=\operatorname{im} H_T=\ker K_T,\qquad D_T=\ker H_T=\operatorname{im}K_T
\]
と置く。$H_T$を作用させた三位置は同時対称であり、その三因子全てが$A_T$に支持される。
従って$V$と$T$にわたる混合縮約像は$A_T\otimes T+V\otimes D_T$に入る。

まず一次空間を順に$V/A_U,V/A_W,W/D_W,U/D_U$へ落とす。
位置1,3の関係の像は零なので、残る二つの二次経路商の次元を$e_0,e_2$とすれば、
長さ$4m$の経路商次元は$(e_0e_2)^m$になる。
$H_W,H_U$を同時に作用させたテンソルは
$\operatorname{Sym}^4(A_W\cap A_U)$に入る。実際、それぞれの変換は別々の位置に作用し、
変換後の各位置と元の二つの$V$位置との交換対称性が全置換を生成する。
従って$A_W\ne A_U$なら、位置2の投影関係の次元は高々1で$e_2\ge3$である。

各$T$について、$H_T$を作用させたテンソルの二つの元の$V$位置への縮約像は
$R_0\cap\operatorname{Sym}^2 A_T$の中で少なくとも二次元である。
そうでなければ、残る一因子を縮約して得られる全ての対称三テンソルが
一つの二次テンソル$q$に対し$q\otimes a$の形になる。
非零で対称なら$q$は$a^2$の倍であり、全ての$a$の方向も等しい。
これは変換した$T$位置の一因子階数が$\rk H_T=2$であることと矛盾する。
$A_W\ne A_U$の場合は二つの対称二乗の交叉が高々一次元なので、
$R_0$の少なくとも三次元が位置0の射影で消える。よって$e_0\ge1$となる。
これは指数成長を与えるので、$A_W=A_U=A$の場合だけが残る。

$B=V/A$と置く。位置0の投影関係は$\operatorname{Sym}^2 B$に入る。
位置2の投影関係も、$H_W,H_U$で商を$A$と同一視すると
$\operatorname{Sym}^2 A$に入る。従って$e_0,e_2\ge1$であり、
多項式成長には両方が1でなければならない。二つの投影の階数は従って3である。
特に$H_W,H_U$を両方作用させた二変数四次形式は中間縮約の階数が3となり、
$R_0\supset\operatorname{Sym}^2A$を得る。
また二つの$V$位置を$B$へ落としたテンソルは$D_W\otimes D_U$に外側支持を持つ。
$K_W,K_U$の同型$B\simeq D_W,D_U$で読むと、これは$\operatorname{Sym}^4 B$の元である。
実際、二変換を別の$V$位置に施すと元の二位置間の対称性と、
それぞれの変換位置と対応する外側位置との対称性が残る。
こちらも中間縮約の階数が3なので、この同一視の下で
$G=\operatorname{Sym}^2 B\subset D_W\otimes D_U$が$R_2$に含まれる。
$R_2$から$(W/D_W)\otimes(U/D_U)$への射影の核はちょうど三次元だから
$R_2\cap (D_W\otimes U)=G$である。
さらに$(R_0/\operatorname{Sym}^2 A)\simeq G^*$が最上テンソルで誘導される。
実際、$W$または$U$をそれぞれ$D_W^\perp,D_U^\perp$で縮約すると
両$V$位置は$A$に入るので、$\operatorname{Sym}^2 A$で割ったテンソルの外側支持は
$R_2\cap(D_W\otimes D_U)=G$であり、階数は3である。

最後に命題\ref{prop:adjacent-position-pairs}の四次サイクルの計数を、
$W,U$の同一視なしに行う。$F_A=R_1\cap(A\otimes W)$は純テンソルを含まない。
含めば$R_0$内の平方から連続純粋関係になる。
逆順に読んだ$R_3$の$V\otimes D_U$部分にも純テンソルはない。
実際、$G$は$B\simeq D_W,D_U$で同じ文字を二回使う全ての平方を含み、
純テンソルがあれば$R_2,R_3$に連続純粋関係ができる。
先の射影次元定理より両空間の次元は3以下である。
$V$を$A^\perp$、$W$を任意に縮約した像は転置切断によって次元$6-\dim F_A$で、
逆順の$R_3$の$V\otimes D_U$部分に入る。従って両次元は3であり、
この像の$V$支持は4である（$3\times2$内の三次元部分空間は純テンソルを含む）。
これにより$(F_A\otimes U)\cap(A\otimes R_2)=0$である。
実際、その元を$w(\phi)$、$\phi\in V^*$と書くと、$\phi$は今の四次元支持を消す。

$A\otimes W\otimes U$の$C_{1,4}$での像は$32-12-12=8$次元以下であり、
$\operatorname{Sym}^2 A\otimes C_{2,4}$から少なくとも14次元の四次サイクルを得る。
$R_0/\operatorname{Sym}^2 A=G^*$へ射影した$d_3$は、
$D_U^*\otimes U$から$G^*\otimes((D_W\otimes U)/G)$への縮約である。
$B\simeq D_W,D_U$のもとで前命題と同じく、$U/D_U$方向の階数4と
内側の$z\wedge Tz$の階数3で合計7となる。
従って先のサイクル空間内の境界は$15-7=8$次元以下であり、
$\dim C_{0,4}\ge41$となって矛盾する。
この証明は外の二位置の間の対称性を要求しない。
\end{proof}

\begin{corollary}[向かい合う一対が残る場合の相補平面]\label{cor:opposite-witness-complementary-planes}
$Z_{\rm sym}\setminus S$内の標準AS対象に向かい合う可逆正証明書があると仮定し、
巡回順序を$V,W,V,U$とする。$S$上の全同一視を同時に極限に取って得られる
$H_T:T\to V$, $K_T:V\to T$（$T=W,U$）と$L:U\to W$, $M:W\to U$は
いずれも階数2である。さらに
\[
 V=A_W\oplus A_U,\qquad A_T=\operatorname{im}H_T=\ker K_T,\qquad
 \ker H_W=\operatorname{im}L,\quad \ker H_U=\operatorname{im}M .
\]
ここで二つの$V$位置の一方を同一視の基準に選ぶ。
\end{corollary}
\begin{proof}
定理\ref{thm:no-two-pair-AS-boundary}と三因子対称化の認識により、
列挙した六つの写像は非零かつ特異である。$L,M$は向かい合うため階数2で、
対の外から入る$H_W,H_U$は系\ref{cor:no-rank-one-into-positive-pair}により階数2以上である。
$A_W+A_U\ne V$なら、二つの$V$位置をともに$V/(A_W+A_U)$へ、
$W,U$をそれぞれ$W/\ker H_W,U/\ker H_U$へ落とす。
$H_T$を施した三位置が$A_T$に支持されるので、全ての隣接関係が消える。
長さ$4m$の商次元は少なくとも$4^m$となり、従って$A_W+A_U=V$である。

全射影行列のグラフ閉包では、三位置間の合成行列と直接の同一視は線型従属である。
これは$S$上の恒等式から従い、合成が零のときも閉条件である。
例えば$H_W L$が非零なら$H_U$に比例し、$A_U\subset A_W$となる。
$A_W$は真部分空間だから、これは上の和に反する。従って$H_W L=H_U M=0$である。
よって$\ker H_W,\ker H_U$は二次元以上であり、階数の下界からいずれもちょうど二次元である。
従って像$A_W,A_U$は相補的であり、$\ker H_W=\operatorname{im}L$,
$\ker H_U=\operatorname{im}M$を得る。
$K_WH_W=0$かつ$K_W\ne0$なので$K_WH_U\ne0$であり、合成は$L$に比例する。
従って$K_W$の階数も2である。同様に$K_U$の階数も2で、主張を得る。
\end{proof}

\begin{theorem}[対称閉包内での任意の可逆正証明書]\label{thm:one-positive-AS-in-symmetric-closure}
$Z_{\rm sym}$に属する標準AS対象について、六つの位置対の少なくとも一つに
可逆な正証明書があることは、対象が$S$に属することと同値である。
\end{theorem}
\begin{proof}
$S$では全六対に可逆な正証明書がある。逆向きは隣接対について既に示したので、
$Z_{\rm sym}\setminus S$内のAS対象に向かい合う可逆対があると仮定する。
系\ref{cor:opposite-witness-complementary-planes}の記号で
$D_T=\ker H_T$とし、順に一次空間を
\[
 V/A_W,\quad W/D_W,\quad V/A_U,\quad U/D_U
\]
へ射影する。位置0,2の隣接関係が消え、残る二つの独立な辺の関係の像次元を
$s_W,s_U$とすると、位置1から始まる長さ$4m$の商次元は
$((4-s_W)(4-s_U))^m$である。$W/D_W\simeq A_W$、
$V/A_U\simeq A_W$の下で辺1の関係像は$\operatorname{Sym}^2 A_W$に入る。
これは$H_W$を$W$位置に施した三位置の対称性と支持から従う。
辺3でも同様に$\operatorname{Sym}^2 A_U$に入る。
従って$s_W,s_U\le3$であり、多項式成長には$s_W=s_U=3$が必要となる。

今度は両方の$V$を$P=V/A_U$に落とし、$W$は全て残し、$U$を$U/D_U$に落とす。
$U$に接する二辺の関係像は零である。
$V,W$に残る関係の像次元が高々3であることを示す。
元の$V\otimes W$内の関係空間を$F$とすると$\dim F=6$で、
他方の$V$位置と$U$位置を縮約した全像が$F$である。
先の辺3への投影の転置は、$U$位置を$D_U^\perp$、
対応する$V$位置を$A_W^\perp$で縮約する写像である。
$s_U=3$よりその像は三次元であり、$H_U$の支持条件から
$F\cap(A_U\otimes W)$に入る。従って$F$の$P\otimes W$への像次元は高々3である。
二つの$V$位置の対称性から、$W\otimes P$側の関係像にも同じ上界がある。

残る三文字$P\otimes W\otimes P$の商は少なくとも
$2\cdot4\cdot2-3\cdot2-2\cdot3=4$次元である。
これに二次元の$U/D_U$が関係なしに続くので、長さ$4m$の商次元は少なくとも$8^m$となる。
標準AS対象の多項式成長に矛盾する。
\end{proof}

% !TeX root = ../papers/nc-p3.tex
\section{全ての対称化写像が特異な境界}
\label{sec:singular-transfer-planes}
$Z_{\rm sym}$では、$S$上の四位置間の同一視のグラフを射影行列空間の積で閉じると、
任意の点の上で同時に正証明書$H_{ij}:V_j\to V_i$を得る。
以下ではこのように選んだ一組を固定する。各行列は非零で、
$H_{ik}H_{kj}$と$H_{ij}$は線型従属である。逆向きとの積はスカラーである。

\begin{lemma}[一位置からの三つの写像が二次元の核を共有する場合]\label{lem:three-common-binary-quotient}
一因子階数が全て4の四テンソルで、ある位置$j$から他の三位置への
階数2の正証明書が共通の核$D\subset V_j$を持つとする。
二次関係を隣接対の縮約像で定めると、連続純粋関係が存在する。
\end{lemma}
\begin{proof}
$j$を$B=V_j/D$へ落とす。他の位置$i$は
$A_i=\operatorname{im}H_{ij}$に支持される。実際、各$B\simeq A_i$で
射影した$j$を$i$の空間で読むと、二位置の交換で不変となる。
従って全ての$A_i$を$B$と同一視したテンソルは
三つの互換$(j,i)$で不変であり、$q\in\operatorname{Sym}^4 B$である。
$j$位置の一因子階数は2である。従って$q$の二対二切断の像
$E\subset\operatorname{Sym}^2B$は少なくとも二次元である。
実際、階数1なら対称な切断はスカラーを除き$c\otimes c$で、
中央の二文字の交換から二次形式$c$の二次小行列式は全て零となる。
すると$q$は四乗の倍で、一因子階数に反する。

代数閉体上の射影平面$\mathbb P(\operatorname{Sym}^2 B)$では、
正の次元の線型部分空間は平方の二次曲線に交わる。
よってある$0\ne b\in B$で$b^2\in E$である。
位置$j$以外の三文字を巡回順に見ると、その二つの隣接関係像は、
いずれも対応する同一視のもとで$E$を含む。これは$j$を$D^\perp$で縮約して得る部分像である。
従って$b$に対応する三文字を取れば、その前二文字も後二文字も純粋関係となる。
\end{proof}

\begin{theorem}[対称閉包のAS残余に必要な隣接階数1]\label{thm:singular-transfer-rank-one-required}
$Z_{\rm sym}\setminus S$内の標準AS対象を仮定する。どの同時極限の組にも、
隣接位置間の階数1の正証明書が少なくとも一つある。
\end{theorem}
\begin{proof}
定理\ref{thm:one-positive-AS-in-symmetric-closure}により全ての写像が特異で、
向かい合う位置間では系\ref{cor:opposite-singular-rank-two}から階数2である。
まず全ての向きで階数2と仮定する。各位置$k$で
\[
 A_{ki}=\operatorname{im} H_{ki}=\ker H_{ik}\subset V_k\quad(i\ne k)
\]
と書ける。逆向きとの積が零であることと次元から等号を得る。
$i,j,k$が異なるとき、$H_{ik}H_{kj}$が零なら$A_{ki}=A_{kj}$である。
零でなければ$H_{ij}$に比例し、階数2なのでこの二平面は相補的である。
この場合、合成の像と核から
\[
 A_{ij}=A_{ik},\qquad A_{ji}=A_{jk}
\]
も従う。

ある位置では三平面が全て一致することを示す。
どの位置でも全て一致する場合は明らかである。
そうでなければ位置$k$の三平面を等しいものごとに分け、
単独となる位置$i$を一つ選ぶ。残りの二位置$j,\ell$に対し、
$k$での$i$の平面はそれぞれの平面と相補的なので、
上の合成則は$A_{ij}=A_{ik}=A_{i\ell}$を与える。
位置$i$から出る三写像の核が一致するため、
補題\ref{lem:three-common-binary-quotient}と定理\ref{thm:no-consecutive-pure-relations}が
AS性に矛盾する。従って全階数2は不可能である。

残る写像には階数1か3があり、それは向かい合う位置間にはない。
階数3の場合、同時極限の逆向き写像との積が零であり、逆向きは非零なので階数1である。
従っていずれの場合も同じ組に隣接階数1が存在する。
\end{proof}

三平面の等号関係は各位置で五通りであり、四位置で625通りである。
合成から得た含意を満たすものは29通りで、共通核を持つ位置が二、三、四個のものは
それぞれ12、16、1通りとなる（\path{computations/verify_transfer_plane_patterns.py}）。
この個数は等号関係の組合せに関するものである。

\begin{proposition}[隣接階数1から得る二変数の部分テンソル]\label{prop:forced-binary-endpoint-pencil}
上のAS残余を仮定すると、巡回順またはその逆順で位置を$0,1,2,3$と選び、
位置0を適当な二次元商$B$へ落としたテンソルは、位置2の平面を$B$と同一視して
\[
 x^2\otimes u\otimes v_0+
 (x\otimes y+y\otimes x)\otimes u\otimes v_1+
 y^2\otimes u\otimes v_2
 \quad\in\operatorname{Sym}^2 B_{0,2}\otimes V_1\otimes V_3
\]
と書ける。ここで$x,y$は$B$の基底で、$0\ne u\in V_1$、
$v_0,v_1,v_2\in V_3$は一次独立である。
特に一つの隣接対には像直線が共通の階数1正証明書からなる二次元線型空間が必要である。
\end{proposition}
\begin{proof}
階数1の同時極限を$H_{10}=u\ell$とする。
$\ell$が$\ker H_{20}=\operatorname{im} H_{02}$を消せば
$B=V_0/\ker H_{20}$上の汎関数となる。そうでなければ
$H_{10}H_{02}$は階数1で非零であり、$H_{12}$に比例する。
その共変ベクトルは$\ker H_{02}$を消すから、位置0と2を交換し逆順を取れば
同じ状況に帰着する。$H_{20}$で$B$を位置2の平面と同一視すると、
位置0を$B$へ射影したテンソルは位置0,2の交換で対称である。

$\ell$で縮約した三テンソルは$u\otimes M$で
$M\in B\otimes V_3$である。一因子階数から$M\ne0$であり、
階数1なら連続純粋関係ができるので$\rk M=2$である。
従って$R_1$は$u\otimes B$を含む。$\ell(x)=1,\ell(y)=0$と基底を選べば、
二つの$B$位置における$x^2,xy+yx$の係数は$ku\otimes V_3$に入る。
残る$y^2$の係数を$T\in V_1\otimes V_3$とする。
$T$が$ku\otimes V_3$に入らなければ、$u$を消す汎関数と
位置0の$y$の双対基底で縮約し、$R_2$に$y\otimes v\ne0$を得る。
$u\otimes y\in R_1$に続くので矛盾である。これで表示式を得る。

どの非零の$B^*$による縮約も同じ理由で$M$の階数が2でなければならない。
三係数$v_0,v_1,v_2$が高々二次元しか張らないと、二次元に埋め込んだときの
$2\times2$行列束の行列式は二変数の斉次二次式であり、代数閉体上で射影零点を持つ。
従って三係数は独立である。任意の$\lambda\in B^*$に対し
$u\lambda:V_0\to V_1$は、作用後の二位置をともに$ku$に支持させるので正証明書である。
\end{proof}

この表示で$Z=\langle v_0,v_1,v_2\rangle$とすると、もう一つの向かい合う対について
\[
 u\in\ker H_{31}=\operatorname{im}H_{13},\qquad Z=H_{13}^{-1}(ku)
\]
である。実際、$H_{31}$を位置1に作用させると位置1,3は対称になる。
$H_{31}u\ne0$なら位置3の三次元の支持$Z$が変換後にそのまま残るが、
対称な相手の支持は一次元なので矛盾する。逆方向に$H_{13}$を位置3に作用させると、
位置1の支持が$ku$に入るため$H_{13}Z\subset ku$である。
$\ker H_{13}$は二次元で$\dim Z=3$なので$H_{13}Z=ku$であり、主張を得る。

\begin{theorem}[対称閉包のAS点の完全な認識]\label{thm:symmetric-closure-as-exact}
$Z_{\rm sym}$内の標準AS対象の集合は、対称完全交叉チャート$S$に等しい。
\end{theorem}
\begin{proof}
$S$の全点のAS性は定理\ref{thm:symmetric-frobenius-chart}による。
逆に、$Z_{\rm sym}\setminus S$に標準AS対象があると仮定する。
定理\ref{thm:singular-transfer-rank-one-required}と
命題\ref{prop:forced-binary-endpoint-pencil}の表示を用いる。
そこでの巡回順序の変更は証明の始点の選択であり、逆順の場合は反対代数で議論する。
標準AS性と標準Hilbert級数はこの操作で保たれる。
分類の同値関係に添字移動や反転を加えるものではない。

次の二次元空間と三次元空間を記す。
\[
 \begin{gathered}
 D_0=\ker H_{20},\qquad B_0=V_0/D_0,\qquad
 B_2=\operatorname{im}H_{20},\\
 A_1=\operatorname{im}H_{13},\qquad
 Z=\langle v_0,v_1,v_2\rangle=H_{13}^{-1}(ku).
 \end{gathered}
\]
特に$u\in A_1$であり、命題の表示から
\[
 R_1\supset ku\otimes B_2,\qquad
 \operatorname{im}\bigl(R_0\longrightarrow B_0\otimes V_1\bigr)
       =B_0\otimes ku.                                      \tag{1}
\]
第二の等号は、位置0だけを商にしたテンソルを残る位置2,3で全て縮約して得る。
表示中の三係数$v_0,v_1,v_2$の独立性により、その像は二次元である。

$Z$を消す非零の$\ell\in V_3^*$を取り、位置3で縮約した三テンソルを$c=w(\ell)$とする。
位置3の一因子階数が4なので$c\ne0$である。
位置0を$B_0$へ落とした表示の位置3の支持は$Z$に入るため、
$c\in D_0\otimes V_1\otimes V_2$である。
また$\ker H_{13}\subset Z$より、ある$\eta\in V_1^*$で$\ell=\eta H_{13}$と書ける。
位置3に$H_{13}$を作用させたテンソルは位置1,3の交換で不変であり、
位置3の支持が$A_1$なので位置1の支持も$A_1$に入る。
その位置3を$\eta$で縮約すると
\[
                         c\in D_0\otimes A_1\otimes V_2                 \tag{2}
\]
を得る。

$c$の末端の支持を$E\subset V_2$とし、$e=\dim E$と置く。
このとき$1\le e\le2$である。
実際、位置2で$c$を縮約した像
$J\subset D_0\otimes A_1$の次元は$e$である。
$e\ge3$なら、四次元の$D_0\otimes A_1$の中で
$J$は二次元の$D_0\otimes ku$と非自明に交わる。
従ってある$0\ne d\in D_0$について$d\otimes u\in J\subset R_0$となる。
(1)の$u\otimes B_2\subset R_1$と合わせると連続純粋関係が存在し、
定理\ref{thm:no-consecutive-pure-relations}に反する。

ここで\emph{元の二次代数の商}を作る。
$Q=V_2/E$、$z=\dim Q=4-e\ge2$とし、四位置の一次空間を順に
\[
                        B_0,\quad V_1,\quad Q,\quad k_\ell              \tag{3}
\]
へ写す。最後の写像は$\ell:V_3\to k$である。
$R_2$を$\ell$で縮約した像は$c$の末端支持$E$であり、
$R_3$を$\ell$で縮約した像は$c$の先頭支持で、$D_0$に含まれる。
従って(3)では辺2,3の関係空間全体が零になる。
辺0の関係像は(1)より$B_0\otimes ku$である。
辺1の関係像を$F\subset V_1\otimes Q$とし、
\[
 s=\dim(B_2\cap E),\qquad r=\dim F,\qquad
 t=\dim\bigl(F\cap(ku\otimes Q)\bigr)
\]
と置く。$\dim R_1=6$と(1)により
\[
                             r\le6-s,\qquad t\ge2-s.                    \tag{4}
\]
すなわち、$R_1$の射影の核は$ku\otimes(B_2\cap E)$を含み、
像は$ku\otimes\operatorname{im}(B_2\to Q)$を含む。

三文字ブロックの商
\[
 T=\frac{B_0\otimes V_1\otimes Q}
          {B_0\otimes ku\otimes Q+B_0\otimes F}
\]
の次元$d$は、分母の二部分空間の交叉が
$B_0\otimes(F\cap(ku\otimes Q))$であることから
\[
 d=8z-2z-2r+2t
   \ge 6z-2(6-s)+2(2-s)
   =6z-8\ge4.                                                        \tag{5}
\]
これは射影した最上テンソルの縮約空間の次元ではなく、
各辺の\emph{全二次関係を射影した像}による商の次元である。

辺2,3の関係が零なので、(3)で得た四周期二次商代数の長さ$4m$の経路空間は、
一次元の矢を挿入した$T^{\otimes m}$である。
従って元の代数$C$について$\dim C_{0,4m}\ge d^m\ge4^m$となる。
例えば$m=6$で
\[
               \dim C_{0,24}\ge4096>2925=\binom{27}{3}
\]
となり、標準AS対象のHilbert級数$(1-t)^{-4}$に反する。
これで対称閉包の残余全体を排除した。
\end{proof}

% !TeX root = ../papers/nc-p3.tex
\section{反対位置の階数と二次関係の零点}
\label{sec:opposite-evaluation}
四窓の標準テンソルを反対位置の切断で最小分解し、
$w=\sum_{h=1}^{\rho} A_h\otimes B_h$と書く。
ここで$A_h\in V_0\otimes V_2$, $B_h\in V_1\otimes V_3$であり、
それぞれの組は一次独立である。各位置の非零共変ベクトル$x$で、
対応する$A_h$または$B_h$を評価した$\rho$本の列の階数を$r_i(x)$とする。
隣接対の二次関係$R_i$の共通零点スキームを
$\Gamma_i\subset\mathbb P(V_i^*)\times\mathbb P(V_{i+1}^*)$とする。

\begin{theorem}[AS対象の第三階数は4以上]\label{thm:opposite-rank-at-least-four}
標準AS対象では全ての位置で$r_i(x)\ge2$である。
また各$\Gamma_i$は非空であり、その全ての点$(x,y)$で
$r_i(x)+r_{i+1}(y)\le\rho$である。従って$\rho\ge4$である。
\end{theorem}
\begin{proof}
位置$i$を$x$で縮約した三テンソルは一因子階数4から非零であり、
その中央因子（位置$i+2$）の支持の次元は、反対側の行列の一次独立性から$r_i(x)$である。
系\ref{cor:opposite-singular-rank-two}の中央支持の障害により階数1は不可能である。
一方、$\Gamma_i$はSegre像$\mathbb P^3\times\mathbb P^3\subset\mathbb P^{15}$と
$\mathbb P(R_i^\perp)=\mathbb P^9$の交叉なので、射影次元定理から非空である。
例えば位置0,1を縮約した二テンソルは、列を$A_h(x),B_h(y)$とする
$4\times\rho$行列$P(x),Q(y)$を用いて$P(x)Q(y)^{\mathsf T}$となる。
これが零であることは$R_0$の全元が$(x,y)$で消えることと同値であり、
$\operatorname{im} Q(y)^{\mathsf T}\subset\ker P(x)$から不等式が従う。
他の辺でも同様である。
\end{proof}

\begin{theorem}[第三階数4では全ての隣接点対応が正次元]\label{thm:rho-four-positive-incidence}
標準AS対象で$\rho=4$なら、各$\Gamma_i$の次元は1以上である。
その全ての点では両評価階数が2で、上の像と核は同じ平面になる。
\end{theorem}
\begin{proof}
$\Gamma_0$が有限と仮定する。六本の双一次式の零点は滑らかな六次元の積空間で
余次元6なのでKoszul分解を持つ（零点の外では局所的に単元を含む）。
その分解を$\mathcal O(3,0)$でひねると、外次数$j=1,\ldots,6$の
$\mathcal O(3-j,-j)$は全コホモロジーが零である。
$j=1,2,3$は第二因子、$j=4,5,6$は第一因子の次数が$-1,-2,-3$のいずれかになる。
よって三次式の$\Gamma_0$への制限は単射（実際には同型）である。
非被約な零点にも同じ結論が成立する。

ところが$P(x)Q(y)^{\mathsf T}=0$は$\Gamma_0$上のスキームとしての方程式である。
各点の近傍で$Q(y)$の二次小行列式の一つが単元なので、その転置の像は
階数2の局所直和因子を含む。従って$P(x)$の全三次小行列式が$\Gamma_0$上で消える。
単射性によりこれらは$x$の三次式として恒等的に零である。
従って$P(x)$は$\mathbb P^3$の全点でちょうど階数2となる。

この最後の場合を除く。行列束
$4\mathcal O(-1)\to4\mathcal O$の像を階数2のベクトル束$L$、余核を$N$とする。
超平面類を$h$として$c_1(L)=-a h$, $c_2(L)=b h^2$と書く。
$L^*$と$L(1)$は大域生成なので行列式の次数から$0\le a\le2$である。
$c(N)=1/c(L)$と$\operatorname{rk}N=2$より$a^3-2ab=0$であり、整数$b$に対し$a=1$はない。
$a=0$なら$L^*$、$a=2$なら$L(1)$が行列式自明な大域生成束であり、自明束となる
（一点で基底を与える二切断の外積は非零の定数）。
前者では$L\to4\mathcal O$の像が固定平面となって位置2の一因子階数4に反する。
後者では$4\mathcal O(-1)\to L$が定数行列で二次元の共通核を持ち、
四つの$A_h$の一次独立性に反する。従って$\Gamma_0$は有限ではない。
巡回して全辺について従い、評価の等号は前定理による。
\end{proof}

\begin{corollary}[第三階数5、有限な点対応の三次式制約]\label{cor:rho-five-minor-budget}
$\rho=5$とし、位置$i$の$4\times5$評価行列の三次小行列式が張る三次式空間の次元を$c_i$とする。
標準AS対象で$\Gamma_i$が有限なら$c_i+c_{i+1}\le20$である。
\end{corollary}
\begin{proof}
前定理のKoszul分解により、両因子の三次式はそれぞれ$\Gamma_i$上に単射に制限される。
また有限完全交叉$\Gamma_i$の長さは$(h_1+h_2)^6$の$h_1^3 h_2^3$係数20である。
局所環上で$P: \mathcal O^5\to\mathcal O^4$の二次小行列式の単元を使い、
階数2の恒等ブロックを分裂させる。$PQ^{\mathsf T}=0$より$Q^{\mathsf T}$は
残る三次元の中間空間に入る。さらに$Q^{\mathsf T}$の単元小行列式でその中に
階数2の恒等ブロックを分裂させる。対応する二列では$P$が零である。
残る中間空間は一次元なので、$P,Q^{\mathsf T}$の残りは二行一列$u$と一行二列$v$で、
$uv=0$である。この表示で三次小行列式イデアルはそれぞれ$u$と$v$の成分で生成される。
従って二イデアルの積は零である。

有限スキーム上の直線束を自明化し、二つのイデアルを$I,J$とする。
$H^0(\mathcal O_{\Gamma_i})$はArtin局所完全交叉環の積でFrobeniusである。
従って積とtraceの完全対合で$I,J$は直交し、$\dim I+\dim J\le20$となる。
制限の単射性から$\dim I\ge c_i$, $\dim J\ge c_{i+1}$である。
これは非被約な点も長さに含む。
\end{proof}

\begin{proposition}[階数4の平面の照合と定数ファイバーの全型]\label{prop:rho-four-plane-fibers}
$\rho=4$の標準AS対象を仮定する。位置0,1で係数空間を$H=k^4$とし、
$P:H\to V_2$, $Q^{\mathsf T}:V_3^*\to H$と読む。
評価階数2以下の小行列式スキームを$D_0,D_1$とすると
\[
 \Gamma_0=D_0\mathbin{\times}_{\mathrm{Gr}(2,H)}D_1,qquad
 x\longmapsto\ker P(x),\quad y\longmapsto\operatorname{im}Q(y)^{\mathsf T}.
\]
平面$K\subset H$におけるファイバーは二つの射影線型空間の積である。
片側の射影次元を$a-1$とすると、非空なら$a=1,2,3$である。
この部分で残る二つの係数行列を制限した$4\times a$束は、左右定数基底変更で
$L_e^{\mathsf T}$の直和となり、その添字は順序を除き
\[
 a=1:\ (1,0,0),\qquad a=2:\ (2,0)\ \text{または}\ (1,1),\qquad a=3:\ (3)
\]
で尽くされる。$L_0^{\mathsf T}$は入力が零次元の一行を表す。
\end{proposition}
\begin{proof}
非空な点では両行列の階数が2である。二次の単元小行列式を選べば、
積が零のスキーム上で双方の三次小行列式は全て消え、階数2の像と核が一致する。
逆も局所直和因子の包含から従うのでファイバー積はスキームとして成立する。
$K$を固定すると$P(x)|_K=0$と$Q(y)|_{K^\perp}=0$は各変数について線型である。
例えば第一の解空間を$X$とする。$H/K$の基底に対応する二行列を$A,B:X\to V_2$とすれば、
$0\ne x\in X$で$A x,B x$は独立である。従って全ての$(s:t)$で$sA+tB$は単射である。
Kronecker標準形\cite[Theorem 1]{DDLpencils}では右特異ブロックは常に核を持ち、
非空の正方正則ブロックはある射影点で行列式が消える。
よって左特異ブロックしかなく、その個数は$4-a$、添字の和は$a$となる。
この分割を列挙すると表示の全型を得る。他方も同じである。
\end{proof}

特に次の二つの幾何的可能性を調べればよい。正次元の既約成分上で平面が一定なら、
少なくとも一方で$a\ge2$となり、その側で消える係数の二次元束は
全ての行列が階数2以下の直線をなす。平面が動くなら、係数行列式の四次曲面
$\det(\sum z_h A_h)=0$は動く直線$\mathbb P(K)$を含み、もう一方の行列式も対応する
$\mathbb P(K^\perp)$を含む。恒等的に消える行列式もこの条件に含める。
従って両方が滑らかな四次曲面なら標準AS性は不可能である。
実際、滑らかな行列式曲面上の行列は階数3であり、またその上の直線は
自己交叉$-2$で法束に切断がなく、射影的な直線スキームは有限である。

\subsection{定数平面の排除}
\begin{lemma}[中央二次元の三テンソルの束]\label{lem:middle-binary-pencil}
標準AS対象の三次交叉の元で中央の支持が二次元なら、その二つの端の支持の上で、
中央を縮約する行列束の非零パラメータの全てで階数は端の二次元数の小さい方に等しい。
両端の次元は異なる。
\end{lemma}
\begin{proof}
中央の共変ベクトル$\lambda\ne0$で得た行列が両端に非零の核を持つと、
一方の核で元のテンソルを縮約して中央が$\ker\lambda$に支持される二次関係を得る。
端の空間は実際の支持なのでこの縮約は非零である。反対の核でも同じ中央直線の純粋関係が生じ、
連続純粋関係の禁止に反する。正方なら斉次の行列式に射影零点があるので同様に矛盾する。
\end{proof}

\begin{theorem}[各照合平面での零点対の一意性]\label{thm:matched-plane-unique-pair}
第三階数4の標準AS対象では、実現される各平面$K$の上の
命題\ref{prop:rho-four-plane-fibers}のファイバーは一点である。
対応する二つの係数行列束は、それぞれ零の行または列を除くと
$L_3$または$L_3^{\mathsf T}$である。従って各辺の正次元の点対応では平面そのものが動く。
\end{theorem}
\begin{proof}
位置0の非零解$x$で縮約すると、中央支持は二次元であり、
両端の行列束は$B$の係数平面$K^\perp$そのものである。
共通零行・零列を除き、必要なら転置すると、前補題とKronecker標準形により
この束は$L_e^{\mathsf T}$の直和で、今度は零添字もない。
二つの係数行列が独立で両端は高々四次元なので、その添字の組は
$(1),(2),(1,1),(3)$で尽くされる。位置1の解から$A$の係数平面$K$にも同じ結論を得る。

位置0の解空間$X$が次元$a\ge2$と仮定する。ここで
$K$の行列は$X$を消し、残り二行列の$X$への制限を$C_s,C_t$とする。
命題\ref{prop:rho-four-plane-fibers}により、その非零行のブロックは
$(3),(2),(1,1)$のいずれかである。平面$K^\perp$の束を$D_s,D_t$と書くと、
位置0を$X$で縮約した二つの隣接切断の行列は、零行・零列を除き
$C_s\otimes D_s+C_t\otimes D_t$と、$D_s,D_t$を転置した同じ式である。
両束の標準化はパラメータを固定する厳密同値で行える。
行は$C$、列は$D$のブロック添字として、この二つの階数は順に
\[
\begin{array}{c|rrrr}
 & (1)&(2)&(1,1)&(3)\\\hline
 (3)&3,4&6,8&6,8&9,12\\
 (2)&2,3&4,6&4,6&6,8\\
 (1,1)&2,4&4,6&4,8&6,8
\end{array}
\]
である。これは各$L_e^{\mathsf T}$の二係数を$(e+1)\times e$単位埋込みの
上詰め、下詰めとして代入するだけで再計算できる
（\path{computations/verify_matched_pencils.py}）。
元の切断の階数は6であり、制限の出力のいずれか一因子の支持が真部分空間なら
制限階数は5以下である。等号6なら制限の像が元の像全体となり、一因子階数4に反する。
表の$(1)$列以外では8以上があるか、第二の階数が6でその$D$側の出力が二次元となる。
従って$D$は$(1)$に限られる。

先に$a=3$を除く。このとき$K$の二行列は$X$を消すので位置0に共通の一次元支持を持ち、
$(1)$型の$D$の二行列も位置1または3に共通の一次元支持を持つ。
従って隣接する二位置に固定文字$u,v$があり、テンソルの各項はそのどちらかを持つ。
順序をこの二位置から始め、必要なら逆順として、
その辺の関係$R$は$ku\otimes V+V'\otimes kv$に入る。
$\dim R=6$なので$V'\otimes kv$との交叉は少なくとも$6-3=3$次元である。
最初の文字を$u$の零化空間で縮約すると、隣の位置は$kv$に支持される。
一因子階数4なので非零縮約があり、次の辺には$v\otimes z\ne0$が入る。
従って連続純粋関係となる。これでどの辺のどちら側でも$a=3$は不可能である。
$a=2$でも$(1)$型の$D$は一端に三次元の共通核を持つ。この端は位置0と隣接し、
$K$の共通核$X$が非零なので、そこにも同じ平面を照合する零点がある。
既に排除した三次元解空間に帰着する。

こうして照合される両束のいずれの端にも二次元以上の共通核はない。
実際、もう一方の束の共通核が非零な端はその両端と隣接する。
簡約後の両端の次元は従って少なくとも3であり、列挙した型でこれを満たすのは$(3)$だけである。
一点のファイバーは射影線型空間の積なので被約でもある。
\end{proof}

\subsection{動く平面の曲面と第三階数4の排除}
\begin{lemma}[動く照合直線が掃く曲面]\label{lem:matched-ruling-quadric}
上の階数4のAS仮定の下で、辺0の照合平面の動く既約族が
$\mathbb P(H)$に掃く曲面は滑らかな二次曲面である。
その直線族は一方の母線族の全体であり、直線ごとに位置0で固定される核点は
三次の有理正規曲線を描く。
\end{lemma}
\begin{proof}
係数行列を$A(z):V_0^*\to V_2$と読む。各照合直線では全点で階数3、
位置0の核は固定であり、余核の双対直線は$L_3^{\mathsf T}$の三次小行列式により
$\mathbb P(V_2^*)$の有理正規三次曲線を動く。
従って$\det A$は恒等的に零ではない。実際、固定核を$x$、余核の双対を$y(z)$とすると、
行列式の任意の方向$z'$への微分は非零倍を除いて$y(z)A(z')x$である。
全て零なら、$y(z)$が$V_2^*$を張るので全ての$A(z')$が$x$を消し、
一因子階数に反する。

同じ向きの異なる照合直線は交わらない。交点での階数3行列の核$x$は一意で、
$x$の評価の核平面も一意だからである。射影的な正次元の既約族から
既約射影曲線を取り、その正規化を$T$とする。ここではGrassmann多様体内の
曲線を取る。固定核$x$は$T$上の射である（束上の共通核は定数階数）。
その平面束の射影化$\pi:Y\to T$を$\mathbb P(H)$へ写すと、像は
$\det A=0$に入る既約曲面$D$である。相異なる直線の非交叉と標数零から
$Y\to D$は双有理であり、超平面の引戻しを$h$とすると$\int_Y h^2=\deg D$である。
同じ非交叉性から照合直線の既約族の次元自体も高々1である。

$d=\pi^*c_1(x^*\mathcal O(1))$とする。$Y$上で
$4\mathcal O(-h)\to4\mathcal O$の核は第一Chern類$-h-d$の直線束、
余核も直線束でその第一Chern類は$3h-d$である。$d^2=0$を使い、
核と余核の完全列からChern指標の次数2部分を比較すると
\[
 (3h-d)^2=-4h^2+(h+d)^2,\qquad 3\deg D=2\deg_T x^*\mathcal O(1).
\]
従って$\deg D$は2または4である。

次数4なら行列式曲面は既約でその台は$D$である。
位置2でも正次元の照合族があるので、その直線族も同じ$D$を掃く。
こちらは余核側が固定され、$L_3$型により核点は動くから、二族に共通の直線はない。
第二族の任意の直線は第一族の全ての直線に会う。
実際、この直線は$D$の中にあり、その各点は第一族の直線上にある。
ただ一本の第一族の直線には収まらず、会う直線のパラメータ集合は閉なので、曲線全体となる。
第一族の相異なる三直線は互いにねじれの位置にあり、これら全てに会う直線は
滑らかな二次曲面上に入る。これは最初の二直線を相補二平面の射影、第三をその間の
同型のグラフとして読むと得られる。従って第二族の掃く$D$は二次曲面に入り、次数4に反する。
ゆえに$D$は既約二次曲面である。円錐では全直線が頂点を通るため、$D$は滑らかである。
射影的な動く族はその一方の母線族$\mathbb P^1$全体を与える。

$D=\mathbb P^1_p\times\mathbb P^1_q$で第一変数$p$を固定核のパラメータとする。
上の等式から$x^*\mathcal O(1)$の次数は3であり、余核直線束は
$\mathcal O(3,3)\otimes\mathcal O(-3,0)=\mathcal O(0,3)$である。
従って余核点は$q$だけで決まる。$q$を固定する直線上では余核が固定で常に階数3なので、
共通零行を除いた$3\times4$束は$L_3$である。その核点$x(p)$は有理正規三次曲線を描く。
同様に$q$側の余核点も有理正規三次曲線である。
\end{proof}

\begin{theorem}[標準AS対象の第三階数は5以上]\label{thm:as-opposite-rank-five-lower}
標準14641型の全AS対象について$5\le\rho\le16$である。
\end{theorem}
\begin{proof}
第三階数4を仮定する。辺0の正次元の照合から前補題の二次曲面を取り、
$H=P\otimes Q$と書いて分解可能元$z=p\otimes q$のSegre像をその曲面とする。
$A(p,q)$は両変数に線型で、全ての点で階数3である。核点は$p$だけ、
余核点は$q$だけで決まる。双対係数空間$H^*=P^*\otimes Q^*$の対応する直線は
$(kp)^\perp\otimes Q^*$なので、こちらも双対Segre曲面の一方の母線族の全体を掃く。
$B(p',q'):V_1^*\to V_3$もそこで全点階数3、核点は$p'$のみに依存する。
この曲面上でも、先の核・余核の直線束の計算と母線への制限により
二つの点は各変数の有理正規三次曲線を描く。

この条件を満たす$A$の表示を明示する。標準三次曲線を核とする
$3\times4$の線型束$D(p)$と、相手の三次曲線を余核点とする
$4\times3$の線型束$M(q)$をそれぞれ$L_3,L_3^{\mathsf T}$から取る。
$D(p)$は$4\mathcal O(-1,-1)\twoheadrightarrow3\mathcal O(0,-1)$で
$A$と核を共有し、$M(q)$は$3\mathcal O(0,-1)\hookrightarrow4\mathcal O$で
$A$と像を共有する。従って$A(p,q)=M(q)C D(p)$であり、
中央の自己同型$C$は定数可逆$3\times3$行列である。
同様に$B(p',q')=M'(q')C'D'(p')$である。
これらは分解可能元での表示だが、それが係数空間を張るので全係数が決まる。

$P,Q$の基底と双対基底で係数を読み、位置0,1から2,3への切断行列を取ると
\[
 \sum_{i,j} A_{ij}\otimes B^{ij}
 =\Big(\sum_j M_j\otimes M'{}^{j}\Big)(C\otimes C')
   \Big(\sum_i D_i\otimes D'{}^{i}\Big).
\]
左の因子は九次元から十六次元への単射である。
実際、$L_3^{\mathsf T}$の二係数は二変数二次式に各変数を掛ける二つの写像である。
従ってテンソル積の和は、別々の二組の変数について非零な双一次式による乗法になる。
変数の基底変更もこの積の非零性を保つ。右の因子は転置で同じ議論を適用して全射である。
ゆえに切断の階数は9であり、標準の隣接切断階数6に反する。
定理\ref{thm:opposite-rank-at-least-four}と合わせて主張を得る。
\end{proof}

\subsection{有限な三つの辺と階数5の残余}
非零$x\in V_i^*$に対し、$R_i$を$x$で縮約した像の次元を$f_i(x)$、
$R_{i-1}$の縮約像の次元を$b_i(x)$と書く。この二数は三テンソル$w(x)$の
巡回順での先頭・末端の支持次元である。$\Gamma_i$が有限なら
どの$x$でも$f_i(x)\ge3$であり、$(x,y)\in\Gamma_i$なら
$f_i(x)=b_{i+1}(y)=3$である。零点の片方を固定したファイバーが線型空間であるためである。

\begin{proposition}[連続三辺の有限性と第三階数]\label{prop:three-finite-edges}
標準AS対象の連続する三つの$\Gamma$が有限ならば$\rho\ge6$である。
従って第三階数5では少なくとも二辺の点対応が正次元である。
\end{proposition}
\begin{proof}
$(x,y)\in\Gamma_i$では常に
\[
 b_i(x)+f_{i+1}(y)\le6 .
\]
実際、$w(x)$の末端因子を縮約して得る$R_{i+1}$の$b_i(x)$次元部分空間は、
$w(x,y)=0$なので$y$による縮約で消える。
前後の辺がともに有限なら左辺の両数は3以上であり、従って両方3である。
中央の辺$\Gamma_i$も有限なら、$w(x)$、$w(y)$のそれぞれで両端の支持次元は3,3になる。
このとき補題\ref{lem:middle-binary-pencil}より中央の支持次元は2になれず、1も既に排除した。
従って$r_i(x),r_{i+1}(y)\ge3$であり、像と核の不等式が結論を与える。
\end{proof}

\begin{corollary}[四辺が有限な場合の点列]\label{cor:finite-edge-point-sequences}
標準有限Frobenius候補で四つの$\Gamma_i$が全て有限ならば、それぞれ長さ20であり、
共有する文字への射影は$\Gamma_i$と$\Gamma_{i+1}$を同じ部分スキームへ同型に写す。
従って隣接零点対は両方向に一意に延長し、始点を固定した点加群の点列を与える。
\end{corollary}
\begin{proof}
上記の支持次元の不等式と有限性による等号にはAS性を使っていない。
$\Gamma_i$上で$w(x)$の末端縮約の像は各閉点で三次元なので、
局所的に$R_{i+1}$の階数3の直和因子を含む。これは$y$による縮約で零になり、
残りの生成元は三本で足りる。従って$R_{i+1}(y)$の階数4の小行列式は
スキームとして$\Gamma_i$上で消える。各閉点で階数3なので、次の文字$z$は
局所的な線型方程式の核直線から射として一意に決まり、$(y,z)\in\Gamma_{i+1}$となる。
逆順にも同じ操作があり、共有文字$y$から元の文字を復元する一意性で両射は逆である。
同じ核直線の構成により共有文字への射影は階数3以下の小行列式スキームへの閉埋込みである。
長さは六本の双一次式の完全交叉から20である。各文字列の関係は隣接二文字だけなので、
この射の反復が全点列を与える。
\end{proof}

\begin{proposition}[任意の第三階数での階数2評価のファイバー]\label{prop:rank-two-evaluation-fiber}
標準AS対象の任意の位置で、階数2評価の行空間を係数空間の双対Grassmann多様体へ送る写像は
各平面上に高々一つの点を持つ。特に$\rho=5$でも、階数2の点が動けばその平面は動く。
\end{proposition}
\begin{proof}
同じ平面に対応する共変ベクトルが独立に二つあるなら、その線型包$X$で
評価は固定した二次元の係数だけを使い、全ての非零評価の階数が2である。
これを位置0とする。制限した二行列$C_s,C_t:X\to V_2$は先の表で
$(2)$または$(1,1)$の行に属し、出力支持$E$の次元は少なくとも3である。
相手の二係数$D_s,D_t$は補題\ref{lem:middle-binary-pencil}から同じ表の四列で尽くされる。
切断の階数6と一因子階数4を使う表の議論は$\rho$を使わないので、
$(1)$列だけが残る。従って相手には固定文字$u$と二次元の支持$Z$がある。
逆順も許して固定文字を位置1に取る。位置0を$X^*$へ射影した部分は
\[
  u\otimes\bigl(C_s\otimes z_s+C_t\otimes z_t\bigr),
 \qquad Z=\langle z_s,z_t\rangle\subset V_3
\]
である（表示は因子を並べ替えたもの）。

元の四窓の一次空間を$X^*,V_1/ku,V_2,V_3/Z$へ落とすと、
位置0に接する二辺の関係像が消える。
$R_1$は$ku\otimes E$を含むので残る辺1の像次元は高々3である。
$R_2$は位置0を$X$で、位置1を$u$に非零な一つの汎関数で縮約した
二次元部分空間を含み、これは$V_2\otimes Z$に入る。
その次元は全ての非零$x\in X$の評価階数2から従う。辺2の像次元は従って高々4である。
位置1からの三文字の商は少なくとも$3\cdot4\cdot2-3\cdot2-3\cdot4=6$次元で、
位置0の二次元因子が独立に繰り返せる。長さ$4m$の商の次元は少なくとも$12^m$となり矛盾する。
\end{proof}

% !TeX root = ../papers/nc-p3.tex
\section{第三階数5で強制される低階数行列の超曲面}
\label{sec:rank-five-pencils}
本節では標準AS対象で$\rho=5$と仮定する。
$H$を五次元の係数空間、$A:H\to V_0\otimes V_2$と
$B:H^*\to V_1\otimes V_3$を最小分解による単射とする。
位置0の共変ベクトル$x$の評価$P(x):H\to V_2$の階数が2なら、
$W_x=(\ker P(x))^\perp\subset H^*$は二次元で、
$w(x)$の中央を縮約する行列束は$B(W_x)$である。他の位置も同様に読む。
その両端支持の次元がともに3以下となる階数2点の集合を$E_i$とする。
補題\ref{lem:middle-binary-pencil}より、このとき実際の束は転置を除き
$L_1$または$L_2$であり、両端の次元は$(1,2)$または$(2,3)$である。

\begin{lemma}[評価階数3の点は三文字列の中央に来ない]\label{lem:rho-five-rank-three-neighbors}
$r_i(y)=3$なら、各側の隣接零点は高々一つで、両側が同時に非空になることもない。
また各$(x,y)\in\Gamma_j$には$E_j$または$E_{j+1}$に属する端点がある。
\end{lemma}
\begin{proof}
例えば$y$を位置1とする。左の隣接点$x$は$r_0(x)=2$で、
$W_x$は位置1の評価の核（$H^*$内の二次元空間）に等しい。
命題\ref{prop:rank-two-evaluation-fiber}より$x$は一意である。
右の隣接点$z$についても同じで、その行空間$W_z$は同じ平面になる。
両方存在すると仮定すると$b_1(y)=f_1(y)=3$である。
$x,z$でそれぞれ縮約した三テンソルの位置1,3の支持は、いずれも
$B(W_x)=B(W_z)$のその両端の支持$L\subset V_1$, $N\subset V_3$である。
$y$が$L$を消すので$\dim L\le3$であり、
$b_0(x)+f_1(y)\le6$より$\dim N=b_0(x)\le3$である。
中央二次元の束の両支持次元は等しくない。そこで元の一次空間を順に
$k_x,V_1/L,k_z,V_3/N$へ射影すると全ての隣接関係が消え、
長さ$4m$の経路商次元は$((4-\dim L)(4-\dim N))^m\ge2^m$となる。矛盾である。

最後の主張を示す。$\Gamma_j$上で$f_j(x),b_{j+1}(y)\le3$である。
両方の評価階数が2で、どちらの端も$E$に入らなければ
$b_j(x)=f_{j+1}(y)=4$となり、その和は6以下という不等式に反する。
片方、例えば$y$が階数3ならば、$x$は階数2である。
$y$には前の隣接点があるため次の隣接点はなく、$f_{j+1}(y)=4$である。
従って$b_j(x)\le2$となって$x\in E_j$である。逆順も同じである。
\end{proof}

\begin{proposition}[一つの点から動く階数2束を得る]\label{prop:rho-five-propagate-small}
$x\in E_i$なら、その片側の隣接零点に$E$に属する点の正次元族がある。
さらに$x$の評価の三次元核に属する係数行列は全て階数2以下である。
\end{proposition}
\begin{proof}
$i=0$として順序の反転も許す。束が$L_2$型なら$b_0(x)=3$, $f_0(x)=2$とする。
隣接点$y$は射影直線をなす。全ての$y$について$f_1(y)\le6-b_0(x)=3$なので
次の隣接点も存在する。上の補題より$r_1(y)=2$であり、前後の支持がともに3以下なので
この直線全体が$E_1$に入る。

$L_1$型なら$f_0(x)=1$, $b_0(x)=2$とする。唯一の前方支持を$ku\subset V_1$と書くと、
$w(x)=u\otimes M$で$M\in V_2\otimes V_3$は階数2である。
従って$R_1$は$ku$と$M$の位置2の支持とのテンソル積を含む。
$R_1$の$(V_1/ku)\otimes V_2$への像は高々四次元である。
隣接点全体は$\mathbb P(u^\perp)=\mathbb P^2$であり、その像の基底を$y$で評価して得る列が
$V_2$を張らない点では次の隣接点も存在する。列が四本ならその行列式は
平面四次式であり、恒等的に零の場合を含め零点の次元は1以上である。
四本未満なら平面全体が該当する。これらの点は再び$E_1$に属する。

いずれも$y$の二次元の係数平面$W_y$は$K=\ker P(x)\subset H$に含まれる。
異なる$y$には異なる平面が対応する。各直線$\mathbb P(W_y)\subset\mathbb P(K)$上の
$A$行列は$L_1$または$L_2$型なので階数2以下である。
射影平面内の無限個の異なる直線は稠密であり、階数条件は閉であるため
$\mathbb P(K)$全体で階数2以下になる。
\end{proof}

\begin{theorem}[二つの係数空間の三次小行列式に共通因子が必要]\label{thm:rho-five-two-divisors}
標準AS対象で$\rho=5$なら、$A$と$B$の\emph{それぞれ}の射影四次元係数空間で、
階数2以下の部分は射影平面の動く族が掃く少なくとも三次元の閉部分を含む。
従ってどちらの側でも、全三次小行列式が恒等的に零であるか、
その非零なもの全てに共通の非定数既約因子がある。
特に各側の最大階数が4なら、その側で得た因子の次数は1または2である。
\end{theorem}
\begin{proof}
どの辺の$\Gamma$も非空で、少なくとも一端に$E$の点がある。
前命題をその点に適用すると他方の偶奇の位置に$E$の無限族が得られ、
そこから一点を選んで再び適用すると元の偶奇にも無限族が得られる。
位置は有限個であるから各偶奇で一位置を固定できる。
$E_i$は行列の階数条件で構成可能であり、正次元の既約な部分族を選べる。
各点の評価の核は命題\ref{prop:rank-two-evaluation-fiber}により互いに異なる。
階数2以下の閉部分が無限個の異なる射影平面を含むので、その次元は少なくとも3である。
実際、次元2以下の閉集合に含まれる平面はその既約成分に限られる。
射影四空間の既約な余次元1成分は一つの既約因子で定まり、
全ての三次小行列式はそれで割り切れる。
行列式が非零多項式である場合、その微分も三次小行列式の線型結合なので
同じ因子で割り切れる。標数零により四次の行列式はその因子の二乗で割り切れ、次数の上界を得る。
\end{proof}

\subsection{低階数因子の次数と正規形}
\begin{lemma}[階数2以下の係数空間と平面族が掃く超平面の排除]\label{lem:rank-five-hyperplane}
上の仮定の下では係数空間全体が階数2以下になることはない。
また前定理の$\mathbb P(\ker P(x))$の動く族の閉包が超平面になることもない。
\end{lemma}
\begin{proof}
まず一般行列の階数が2の場合を考え、係数行列の一つを$M=\operatorname{diag}(I_2,0)$とする。
$M+tA$のSchur補が恒等的に零なので、全ての係数行列は
$A=\left(\begin{smallmatrix}P&Q\\R&0\end{smallmatrix}\right)$であり、$RQ=0$となる。
$Q$が一般に可逆なら$R$は恒等的に零で一因子階数4に反する。
従って$Q$の線型空間は階数1以下である。階数1以下の線型空間は共通の像直線を持つか、
全元が一つの共変ベクトルとの外積である。実際、二つの純テンソルの和が純テンソルになるには
一方の因子が比例しなければならず、線型空間全体でも共通の側が決まる。
従ってある$0\ne v\in\ker M$で$Qv$の動く空間は一次元以下となり、$r_i(v)\ge2$に反する。
一般階数が1以下の場合も同じ純テンソルの事実が一因子階数4に反する。

次に動く平面の掃く既約超曲面が$\mathbb P(S)$、$S\subset H$が超平面であるとする。
対応する位置を0とする。動く点$x$で$S$の評価像は一次元である。
実際、$\ker P(x)\subset S$であり、全評価の階数は2である。既約な曲線族上で
その像直線を$ku_x$とする。一般の$z\in S$について$A(z)x$は一般の$x$で非零であり、
$u_x$の曲線の線型包は$\operatorname{im}A(z)$に入る。
その線型包が二次元以上ならば$S$上の階数上界2から、全$A(z)$の像が同じ平面に入る。
この平面を消す共変ベクトルで$H$を評価すると階数は高々1となり矛盾する。
従って$ku_x=ku$は固定である。点$x$の線型包を$U$とすれば$\dim U\ge2$で、
$S$は$U$を$ku$に送る。一般の$A(z)$は$U$上で非零なので、
$V_0^*/U\to V_2/ku$の誘導行列は全$S$上で階数1以下である。
共通像直線なら$A(S)$の像全体が一平面に入るため、上と同じ矛盾になる。
よって共通の共変ベクトルがあり、三次元の超平面$U'\supset U$が$S$により$ku$へ送られる。
誘導行列が全て零の場合も共通像平面の矛盾で除かれる。

$S=\ker\ell$とする。$U'$上では$z\in S$の作用を$A(z)v=f(v)(z)u$と書ける。
$f:U'\to S^*$には核がない（核点では全評価が階数1以下になる）。
$f$を$H^*$に線型に持ち上げれば、$0\ne v\in U'$について
$W_v=\langle\ell,f(v)\rangle$である。実際、$S$の補方向の作用が$ku$に入れば評価階数は1になる。
これらの直線は射影三空間$\mathbb P(\langle\ell,f(U')\rangle)$を覆う。
各$B(W_v)$は中央二次元の束なので、全点で$B(\ell)$と同じ階数である。
こうして$B$の超平面上に共通核・共通余核のない定数階数の$4\times4$行列空間を得る。
共通核や共通余核があれば残る一係数だけで評価することになり、階数下界2に反する。

この定数階数は1,2,3のいずれも不可能である。階数1は純テンソルの性質、
階数2は定理\ref{thm:rho-four-positive-incidence}の$\mathbb P^3$上の束の議論による。
階数3なら$4\mathcal O(-1)\to4\mathcal O$の核と余核は直線束となる。
それらの次数を$k,c$とするとChern指標から$c-k=4$, $c^2-k^2=-4$で、
$c+k=-1$となり整数解がない。これで超平面の場合を除いた。
\end{proof}

\begin{proposition}[係数の五次元空間の正規形]\label{prop:rho-five-spin-form}
$\rho=5$の標準AS性を仮定すると、各係数空間は左右および係数の基底変更で
\[
 A(S,n)=\begin{pmatrix}S&nI_2\\0&\operatorname{adj}S\end{pmatrix},
 \qquad S\in\operatorname{Mat}_{2\times2}(k),\quad n\in k
\]
に限られる。
\end{proposition}
\begin{proof}
上の平面族を含む既約超曲面の方程式を$F$とする。前補題から$\deg F\ge2$であり、
全随伴行列が零になる場合も除かれた。一般階数が3なら$\operatorname{adj}A=F L$で、
$L$は定数または一次の行列、一般に階数1である。
この線型空間（定数の場合はその張る直線）には固定像または固定共変ベクトルがある。
$AL=LA=0$より元の全係数にも共通の核または余核が生じ、一因子階数に反する。
従って一般階数は4で、前定理より$\det A=c F^2$、$F$は既約二次式である。
その階数は3または4である。実際、平面を含むので五変数非退化二次式ではなく、既約性から階数は3以上である。

$F(e)=1$とし、左基底変更で$A(e)=I$とする。$b$を$F$の極化で$b(e,e)=1$とする。
一般の$z$について$F(z-te)$は異なる二根を持ち、それぞれで$A(z)-tI$の核次元は2以上である。
従って二つの固有空間が全空間を張り、密度から全ての$z$で
\[
 A(z)^2-2b(e,z)A(z)+F(z)I=0 .
\]
$e^\perp$の非退化な直交方向の像は、適当にスカラーを掛けると二乗が恒等の反可換行列になる。
根基の像はそれらと反可換で二乗零である。

まず二つの非退化方向を取り、その像を$g_1,g_2$とする。
$g_2$が$g_1$の二固有空間を交換するので両方二次元であり、基底を選べば二行列は
$\operatorname{diag}(1,-1)\otimes I_2$と
$\left(\begin{smallmatrix}0&1\\1&0\end{smallmatrix}\right)\otimes I_2$である。
両者と反可換な行列は$g_1g_2(I_2\otimes T)$の形に限られる。
$F$が階数3なら根基は二次元であり、係数の単射性から二次元の$T$の空間の全元が二乗零になる。
しかし二行二列の跡零行列の射影平面で行列式零は滑らかな二次曲線なので直線を含まない。

従って$F$は階数4である。第三の非退化方向$g_3$も上の$T$を使って表され、
$g_3^2=I$より$T$の二乗は非零スカラーである。
根基の非零な像はこれと反可換であるため、この$T$はスカラーではない。
その二固有空間に従って四次元を二つの二次元ブロックに分けると、
初めの二行列の作用は両ブロックで同じ、第三の作用は符号が逆である。
第二ブロックを$g_1g_2$の二次元因子で共役すれば三行列全てが第一ブロックの負になる。
三行列は第一ブロックの跡零空間を張る。従って$ke$と三方向の作用は
$\operatorname{diag}(S,\operatorname{adj}S)$である。
根基は三つ全てと反可換なので対角ブロックは零、非対角ブロックはスカラーとなる
（最初の二行列との反可換行列はそれらの積の倍であり、第三とは可換する）。
二乗零と非零性より片側の非対角ブロックだけが非零であり、
ブロックの交換と根基の基底変更で表示を得る。
\end{proof}

\begin{theorem}[第三階数の鋭い下界]\label{thm:as-opposite-rank-sharp-six}
標準14641型の全AS対象の第三階数は$6\le\rho\le16$を満たす。
下界6は可換多項式環で達成される。
\end{theorem}
\begin{proof}
階数5を仮定して位置0,2の係数を前命題の表示にする。
$S=\left(\begin{smallmatrix}a&b\\c&d\end{smallmatrix}\right)$と記す。
位置0で入力$(v,w)$を固定して評価すると、$w=0$, $v=(s,t)\ne0$のとき階数2で、
その行空間は$\langle sa+tb,sc+td\rangle$である。
$w\ne0$なら$\operatorname{adj}S\,w$が二次元を張り、$n w$が上ブロックにもう一方向を与えるので
階数は3以上である。転置側は上ブロックの共変ベクトルが非零なら同じく階数3以上である。
下ブロックだけの共変ベクトル$(u,v')$では行空間は
$\langle ud-v'c,-ub+v'a\rangle$となる。
従って両側の階数2点は同じ二次元行空間の族を与え、$(s,t)=(v',-u)$で対応する。

命題\ref{prop:rho-five-propagate-small}と定理\ref{thm:rho-five-two-divisors}の証明により、
位置0または2には$E$の点がある。その反対側にも同じ行空間$W$の点がある。
$B(W)$の両端の支持の次元は$(1,2)$または$(2,3)$（順序の交換を許す）。
この二つの支持を除く商を位置1,3に取り、位置0,2を対応する二共変ベクトルで一次元に落とせば、
補題\ref{lem:rho-five-rank-three-neighbors}と同じく四辺の関係は全て零である。
長さ$4m$の商が少なくとも$2^m$次元となり、AS対象の多項式成長に反する。
階数4以下の排除は定理\ref{thm:as-opposite-rank-five-lower}、可換例の階数は
交代四テンソルの二対二切断の階数$\dim\bigwedge^2 k^4=6$である。
\end{proof}

% !TeX root = ../papers/nc-p3.tex
\section{有限な点対応の二窓の復元}
\label{sec:finite-point-chart}
四つの隣接点対応が有限な枝を扱う。系\ref{cor:finite-edge-point-sequences}により、
各$\Gamma_i$は長さ20で、共有する位置の射影はスキームとして同じ像を持つ。

\begin{proposition}[第三階数6の有限点列の四移動]\label{prop:rank-six-finite-four-shift}
標準AS対象で$\rho=6$、四つの隣接点対応が有限なら、四回の移動は点スキーム上の恒等射である。
\end{proposition}
\begin{proof}
隣接零点のどちらの端点でも、三テンソルの両端支持はともに3次元である
（命題\ref{prop:three-finite-edges}の証明）。中央支持は1でも2でもない。
隣接する二つの評価階数の和は6以下なので、全点列で各評価階数が3になる。
これは点列スキーム上でも定数階数である。実際、隣の評価行列の三次小行列式の単元で
階数3の直和因子を分裂させれば、積が零なので相手の四次小行列式はスキームとして消える。
従って二つの係数束の階数3の行空間は互いに零化空間である。
四移動した行空間は最初の行空間に戻り、その次の元の文字と再び隣接零点対をなす。
次の文字から前の文字への射が一意なので、元の文字に一致する。
この議論は直線束の局所自明化で行えるため非被約点の上でも成立する。
\end{proof}

\begin{corollary}[三辺からの周期点列]\label{cor:three-finite-periodic-sequences}
第三階数6で、連続する三辺$\Gamma_0,\Gamma_1,\Gamma_2$だけに有限性を仮定しても、
中央の長さ20のスキームの点列は両端へ一意に延び、この三辺の全零点を含む四周期点列を与える。
\end{corollary}
\begin{proof}
中央の辺で$b_1=f_2=3$となるのは前後の有限性だけから従う。
系\ref{cor:finite-edge-point-sequences}の局所直和因子の議論により、
中央の零点対$(x_1,x_2)$から前の$x_0$と次の$x_3$がスキーム上で一意に定まる。
両端の辺への射は共有文字への閉埋込みを通じて閉埋込みとなり、同じ長さ20なので同型である。
中央の二文字の両端支持は3,3で、評価もともに階数3となる。
従って$W_1,W_2$をその行空間とすれば$W_2=W_1^\perp$である。
端の二文字の行はそれぞれ$W_1^\perp,W_2^\perp$に入るから互いの積が零となり、
$(x_3,x_0)$は第四の辺の関係も満たす。ここでは端の評価の定数階数を必要としない。
四文字を繰り返せば結論を得る。第四の辺全体の有限性はこの証明からは従わない。
\end{proof}

\begin{theorem}[三つの座標空間からの復元と平方の直交]\label{thm:finite-square-chart}
標準AS対象で四辺の点対応が有限とし、移動で最初の三位置の点スキームを$X$に同一視する。
$D=H^0(X,\mathcal O_X)$とし、各座標直線束を自明化して
四次元の座標空間を$U_0,U_1,U_2\subset D$とする。
$UV$は積の線型包、$U^2=UU$と書く。次が成立する。
\begin{enumerate}
\item 隣接する$U,V$について$U^2,UV,V^2$は十次元で、それぞれの三乗$U^3,V^3$は$D$である。
$U^2 V^2$は超平面で、その零化汎関数$\tau$はスカラーを除き一意である。
$\tau(ab)$は完全対合であり、$U^2$と$V^2$は互いの直交補である。
\item 二つの対象が添字保持同型であることは、その$D$の代数同型と三つの独立な単元倍が
順序付き三空間を対応させることと同値である。
\item $U_2$を単元倍すると$U_2^2=U_0^2$とできる。従って第一窓を固定したとき、
次の座標空間は、$\mathbb P(D)$内の十本の二次式
$z^2\bmod U_0^2=0$が定める交叉に入る射影三平面のスキームから取れる。
\end{enumerate}
\end{theorem}
\begin{proof}
隣接対を$\mathbb P^3\times\mathbb P^3$の六本の双一次式の有限完全交叉と読む。
三次式の制限は定理\ref{thm:rho-four-positive-incidence}の証明の通り同型で、
一次・二次の制限も単射となる。$(1,1)$で消える切断は定義六式が張る。
実際、イデアルのKoszul分解を$(1,1)$でひねると、六本の定数切断のほかには
全次数0以下にコホモロジーを与える項がない。従って$\dim UV=16-6=10$であり、
$R_i$は対応する$U\otimes V\to D$の核そのものである。

$(2,2)$でひねった構造層のKoszul分解では、大域切断の列は次元$15,96,100$で、
ほかに外次数6の$H^6(\mathcal O(-4,-4))=k$が全次数0に出る。
他のコホモロジー項は零なので大域切断の制限像は余次元1である。
余核への辺写像は随伴同型$\omega_X=\mathcal O_X(2,2)$でのtraceである。
実際、Koszul分解を$\omega_{\mathbb P^3\times\mathbb P^3}$で双対化すると外次数$j$と$6-j$が対応し、
この写像はSerre双対の下で$k\to H^0(\mathcal O_X)$の単位に対応する。
従って自明化後に得る$\tau$の積対合は有限Gorenstein環$D$の完全対合である。
$\tau(U^2 V^2)=0$と次元から直交補の主張が従う。

添字保持同型は点列を対応させ、座標直線束の自明化の差は単元倍である。
逆に三空間が対応すれば二つの積の核も対応し、系\ref{cor:two-window-extension}から
AS延長の同型が従う。最初の三つの位置を順序付きで用いている。
最後に辺0と辺1のtraceを$\tau_0,\tau_1$とすると$\tau_1=\tau_0(g\,\cdot)$となる単元$g\in D^\times$がある。
従って$g U_2^2=U_0^2$である。各局所Artin因子では剰余体で平方根を取り、
冪零部分は有限二項展開で持ち上がるので$g$に平方根がある。
この単元で$U_2$を拡大縮小すればよい。偏極により三平面の条件は$U_2^2\subset U_0^2$と同値である。
候補から二窓を作った後には重なりと全次数AS条件を課す必要がある。
\end{proof}

\paragraph{孤立した第三の座標空間を持つ例。}
$v_0,\ldots,v_3$を一次式の基底、$v_4=\sum_{i=0}^3v_i$とし、四次元の二次式系
$\langle v_i^2-v_4^2:0\le i<4\rangle$を取る。共通零点があれば四つの符号の和が$\pm1$となるので、
この系は基点自由で定理\ref{thm:symmetric-frobenius-chart}の対称AS枝を与える。
辺の零点対は対称積がこの系に入る分解である。
$\sum_i t_i v_i^2$、$\sum_i t_i=0$で非零係数が三本または四本なら、その数が階数である。
五本なら$4\times5$の基底枠、可逆対角行列、その転置の積なので階数は3以上である。
従って階数2の二次式は十個の$v_i^2-v_j^2$だけで、
零点は$v_i+v_j,v_i-v_j$の順序を両方取る二十点である。完全交叉の長さから全て被約である。

ここでは$D=k^{20}$で、行をこの二十点とする評価行列から全積を成分ごとに計算できる。
$U_2=U_0$の平方を固定した三平面スキームの接空間は零である。
実際、四つの基底列$a$の変分$h(a)\in D$に対し、十組の
$a h(b)+b h(a)\bmod U_0^2$を零とする線型方程式は、80変数で階数64である。
残る16次元は基底変更で尽きる。
\path{computations/verify_finite_square_chart.py}で積・trace・接階数を有理数で検算できる。
従ってこの例で第一窓を固定すると、第三空間のこの正規化された解は被約孤立点である。
この例の他の解も含む全決定は定理\ref{thm:finite-all-square-roots}で行う。

\subsection{一つの乗法行列による平方根の探索}
\begin{proposition}[長さ20の環での比による復元]\label{prop:finite-square-ratio}
上の有限枝の$D$と十次元空間$Q=U_0^2$を固定する。
単元を含み、$L^2=Q$かつ$\operatorname{Sym}^3L\to D$が同型となる
四次元空間$L$を、次の方法で全て求められる。
$[\alpha]\in\mathbb P(D/k)$に対し
\[
 K_\alpha=\ker\bigl(Q\xrightarrow{\ \alpha\cdot\ }D/Q\bigr),\qquad
 J_\alpha=\{h\in D:hQ\subset K_\alpha^2\}
\]
を作る。十行十列の写像の階数が6、$\dim K_\alpha^2=10$で、
$J_\alpha$に単元$h$があれば$L=h^{-1/2}K_\alpha$を取り、上の単元と三乗の条件を検査する。
通過した場合$J_\alpha=kh$であり、一つの$[\alpha]$から得る空間は平方が1の単元倍を除いて一意である。
\end{proposition}
\begin{proof}
まず$L$を所与とし、単元$s\in L$と$s$に独立な$t\in L$を取り$\alpha=t/s$と置く。
$q\in K_\alpha$なら$q'=\alpha q\in Q$である。$S^2L\to Q$の同型で$q,q'$を二次式として読むと、
$sq'=tq$は$\operatorname{Sym}^3 L$での恒等式になる。
独立な二一次式は互いに素なので$q=s a$、$q'=t a$、$a\in L$である。
逆も成立し、$K_\alpha=sL$、$K_\alpha^2=s^2 Q$となる。

$gQ\subset Q$なら$g$は定数であることを示す。単元で$L$を割って$1\in L$としてよい。
$L=\langle 1,u,v,w\rangle$と書けば、$1,u,v,w$の斉次三次単項式の評価が基底であり、
同じことを$u,v,w$の次数3以下の単項式の独立性として読める。
$g\cdot1\in Q$より$g$は次数2以下で、$gu\in Q$より次数1以下、
$gu^2\in Q$より定数である。この議論の最初の二段は
$gL\subset L^2\Rightarrow g\in L$も与える（自明化を戻しても成立する）。

逆に上記の線型計算から単元$h$が得られれば、次元により$hQ=K_\alpha^2$で、
$h^{-1/2}K_\alpha$の平方は$Q$となる。
三乗等の条件を通過すれば$h^{-1}J_\alpha=k$で、単元の平方根の存在は前定理通りである。
さらに$s=h^{1/2}$とおくと$sL=K_\alpha\subset Q$かつ$\alpha sL\subset Q$なので、
先の含意から$s,\alpha s\in L$である。従って実際に二つの一次式の比を取っている。
$\alpha$への定数の加法と非零スカラー倍は$K_\alpha$を変えない。
一つの解に対応する比の族は、$L$の単元直線$ks$とそれを含む二平面$\langle s,t\rangle$の選択で
記述できる。二つの比の一致を交叉乗算して$\operatorname{Sym}^2L$の中で比較すれば
この旗は一意であり、そのパラメータは$3+2=5$次元である。
\end{proof}
この表示の探索変数は射影十八次元に入り、階数は小行列式で検査できる。
正規化した平方根を得た後に、二次関係・三次交叉・正則性を課す点は同じである。

有理数検算では、先の二十点例で比の乗法行列の階数6、$J_\alpha$の余次元19と正規化された平方を検査する。
さらに平方の復元だけの非被約な対照例として$D=k[t]/(t^{20})$で
$L=\langle\exp(at):a=0,1,4,16\rangle$を取る（指数関数は打ち切る）。
三重和二十個が相異なるので三乗はVandermonde行列によって$D$全体になる。
$\alpha=\exp(12t)$では乗法行列の核は$\exp(4t)L$である。
この対照例の検査は平方と比の計算に限る。

\begin{proposition}[比の階数条件に含まれる二つの成分]\label{prop:finite-ratio-components}
上の二十点の被約例で$Q=U_0^2$とする。
$\mathbb P(D/k)$内で比の乗法行列の階数が6以下となる軌跡には、
$U_0$の旗と$U_1$の旗から来る異なる二つの五次元既約成分がある。
後者の一般点では$J_\alpha=0$である。
\end{proposition}
\begin{proof}
前者の旗で階数が6であることは前命題による。後者については$Q$と$U_1^2$が
あるFrobenius積対合で互いに直交補であることを使う。
$Q\to D/Q$と$U_1^2\to D/U_1^2$の乗法行列はこの対合で互いに転置になるので、
$U_1$の比も同じ階数6を与える。いずれの旗の像も既約で五次元である。

各側で$s,t$の係数を順に$(1,3,9,27),(2,-1,4,7)$とする。
$\alpha=t/s$での階数条件の接方程式は
$\beta K_\alpha\subset Q+\alpha Q$である。実際、階数6の可逆小行列を分裂させると、
一次項は核から余核への写像だけである。通常の成分積と成分和による対合を使えば、
二十変数$\beta$への条件の階数は
$K_\alpha\,(Q+\alpha Q)^\perp$の次元であり、両側で13になる。
$k\cdot1+k\alpha$を除いた接次元は$20-13-2=5$である。
従って各試験点は滑らかで、その旗の像の閉包が既約成分となる。
双方で$\dim K_\alpha^2=10$で、$hQ\subset K_\alpha^2$の$h$に対する線型方程式は
前者で階数19、後者で20である。全て上記の検算ファイルで有理数計算できる。
$\dim J_\alpha\ge1$はこれらの定数階数開部分で閉条件であり、前者の旗で常に成立する。
後者の試験点では成立しないので二成分は異なる。
二成分以外の解の有無はこの接計算からは決まらない。
\end{proof}

この例では座標台から来る第三の五次元成分もある。
一般に上の真の比について、任意の$c\in k$で
$\dim D/(\alpha-c)D\le10$である。
実際、$t-cs$は非零一次式で、その$L^2$倍は三乗の独立性により十次元である。
$D$が被約なら二十座標のうち十一個以上が同じ値になる比を初めから除ける。

例の五つの一次式から辞書順の添字対を取り、各対で和と差、次に差と和を置いて
二十行を0から番号付けする。$S=\{0,2,3,7,8,16\}$上に支持される
$\alpha$の像は$\mathbb P(D/k)$の射影五平面で、乗法の階数は高々6である。
$\alpha$を$S$の指示関数とすると階数は6、接方程式の階数は13であり、
この五平面は再び既約成分である。これは台の補で同じ値を持つので
$U_0,U_1$いずれの旗からも稠密には覆えない。この成分の除外は単元の検査でも行える。

% !TeX root = ../papers/nc-p3.tex
\subsection{有限な四辺での同型不変量と最小周期}
\label{sec:curvilinear-window-orbits}
基礎体は標数零の代数閉体とする。定理\ref{thm:finite-square-chart}の三つの
四次元空間を使い、局所環が全て曲線状の場合の同型を数値の一致に帰着する。
重複点の長さと個数には制限を置かない。

\begin{theorem}\label{thm:curvilinear-window-orbits}
四辺が有限な標準AS窓で、点環が
$D=\prod_{a=1}^h k[t_a]/(t_a^{m_a})$、$\sum m_a=20$とする。
以下の手順で得る、長さ順の各標識付けに対する記録の有限集合は
添字保持同型の完全不変量である。記録は長さ、三組の主元行、
零でない座標の位置、非零な体の元の列からなる。
同じ長さの局所因子の置換以外の軌道探索を必要としない。
\end{theorem}
\begin{proof}
局所因子を順に並べ、各因子では $1,t,\ldots,t^{m-1}$ の順を使う。
三位置 $s=0,1,2$ の空間を列行列 $B_s$ で表示する。
最初から行を読んで独立になるたびに採用した四つの行を
$P=(p_1,\ldots,p_4)$ とし、右から主元小行列の逆を掛けて $B_P=I$ とする。
下三角可逆行列による左作用は先頭の全ての区間の階数を保つので $P$ を変えない。

各位置で独立な単元倍と、三位置共通の代数自己同型を作用させる。
因子の標識を固定すると、この群は対角トーラスと下単三角群 $N$ との積である。
$N$ のLie環の生成ベクトルは、各因子の長さを $m$ として
\[
 \begin{array}{ll}
 t^k\text{ の乗法}&(1\le k<m,\ \text{一つの位置に作用}),\\
 t^{k+1}\partial_t &(1\le k<m-1,\ \text{三位置に同時に作用})
 \end{array}
\]
で張られる。実際、定数項1の単元と一次係数1の座標変換の接空間が
この二種類であり、各群は係数を座標とする連結な単冪群である。
トーラスの $(s,a,d)$ 行での重みは $c_{s,a} b_a^d$ である。
ここで $c_{s,a}$ は単元の定数、$b_a$ は $t_a$ の一次係数である。
$m_a=1$ のときは $b_a$ を省いてよい。

各位置内で非主元行 $r$ を昇順に、列 $j$ を降順に取り、$p_j<r$ の座標だけを読む。
$g\in N$ の作用は $B\mapsto (gB)((gB)_P)^{-1}$ であり、
各座標に先行座標の多項式を加える。実際、主元小行列は下単三角で、
その逆の第 $i$ 行の非対角成分は $p_i$ より前の行だけで決まる。
$r$ 行の積で使うのは $p_i<r$ の列に限られ、$j$ より大きい列は先に読んでいる。
従って既読の値を固定する部分群は次の一座標に平行移動で作用する。

現在の安定化Lie環の元 $X$ の平行移動速度は、対応する成分を同じ記号で書けば
\[
                 (XB-B(XB)_P)_{rj}
\]
である。これは線型汎関数で、零なら値を残す。非零なら速度 $v\ne0$ の
$X$ を選び、三位置へ $\exp(-B_{rj}X/v)$ を作用させ、この座標を零にする。
次は速度の核を使う。単冪群では閉部分群も指数・対数でそのLie環と対応する。
ここでは直接、部分群の元 $g$ が既読値を保てば $g^n$ も全て保ち、
$\exp(z\log g)$ の正規化後の各成分が $z$ の多項式なので全ての $z$ で保つ。
逆向きは既読座標への三角な微分方程式の一意性から従う。
従って速度が零の場合にはその値は軌道上で強制される。
帰納的に同一軌道の代表をただ一つ選んでいるので、速度方向の選択は最終値を変えない。
非零速度を使う回数は当初のLie環の次元以下である。

この正規化はトーラスと可換である。トーラスは $N$ を正規化し、
速度が零かどうかと零にした座標を保つので、各段階の代表の一意性を適用できる。
残った非零座標を $z_1,\ldots,z_l$ とする。行 $(a,d)$、主元行 $(a',e)$ の
座標の重みは $c_{s,a}b_a^d/(c_{s,a'}b_{a'}^e)$ である。
これらの重みを行にした整数行列を $W$ とし、
$\ker_{\mathbb Z}(W^{\mathsf t})$ の整数基底 $u$ ごとに $z^u$ を記録する。
主元と台が等しいときこの値の一致はトーラス同値と同値である。
実際、一つのトーラス変数について行を整数の可逆基本変形し、その指数を
一本の行に集める。対応する単項式座標変換も可逆である。
その変数は非零べき乗を通じこの一座標だけに全射に作用するので、
残る行に同じ手順を施せばよい。除いた座標は逆順に合わせられる。
有理核の独立なベクトルを個別に整数化するだけでは整数核の基底を保証しない。
最後に代数同型は同じ長さの局所因子だけを交換できるから、
標識付けの有限集合を比較し、定理\ref{thm:finite-square-chart}を適用する。
\end{proof}

\begin{corollary}\label{cor:curvilinear-window-period}
この枝で最小周期を求めるには、開始位置 $i=0,1,2$ の各々で、実際の辺の
二つの射影を通じて三位置を同じスキームへ移し、その三空間の記録集合を
$K_i$ とする。$K_0=K_1$ なら1、そうでなく $K_0=K_2$ なら2、残りは4である。
\end{corollary}
\begin{proof}
$K_i=K_0$ と $i$ 移動への添字保持同型は定理から同値である。
同型を許す移動量は $\mathbb Z$ の部分群であり、標準四周期性により4を含む。
境界を越える三位置にも辺の射影を用いるので、四移動の点スキーム上の作用を
恒等写像と仮定する必要はない。
\end{proof}

\begin{proposition}\label{prop:finite-point-aut-solvable}
一般の有限四辺でも、各局所点環 $C$ の自己同型群の単位成分は可解である。
従って一つの点環 $D$ の同型型とその表示を固定すれば、三空間の同型判定は
有限個の成分代表と、三角行列に対する上の消去法に帰着する。
\end{proposition}
\begin{proof}
辺は滑らかな六次元多様体内の有限完全交叉で、第一射影は閉埋込みである。
従って $C$ の埋込み次元 $e$ は高々3である。六つの局所方程式の独立な
一次部を形式座標に採用して消去すると、極小表示は
$C=k[[x_1,\ldots,x_e]]/I$、$I\subset (x)^2$ であり、
極小関係空間 $I/(x)I$ の次元は $e$ である（高さも $e$）。

自己同型が余接空間上で恒等なら、極大イデアルの冪による全ての次数商上で
恒等なので作用は単冪である。単位成分が可解でなければ標数零のLevi分解と
半単純Lie環の根部分環から、余接空間に非自明に作用する
$\mathfrak{sl}_2$ が取れる。完全可約性で余接空間を極大イデアル内に
同変に持ち上げ、上の形式変数をその基底に取る。
$W=I/(x)I$ にも同変な関係式代表を取れる。実際、
$(x)^n\subset (x)I$ となる $n$ で切り、$I$ 内の次数 $n$ 以下の多項式から
$W$ への全射を分裂させればよい。線型変数への作用はこの有限次元空間を保つ。
この $e$ 本は中山の補題で $I$ を生成する。

$e\le3$ で非自明な $\mathfrak{sl}_2$ 表現は、二次元標準表現、これと
自明な一次元の和、または最高重み2の三次元既約表現である。
変数の最高重みを $M>0$ とし、最高重みの変数一つだけを残す。
全ての関係がこの直線上で消えれば商が有限にならない。
一次式がないので $W$ に重み $2M$ 以上が必要である。
$e=2$ または $M=2$ なら次元 $e$ の表現にその重みはない。
残る $e=3,M=1$ なら $W$ は三次元既約で、自明な変数の軸への制限は
自明表現への準同型なので零である。やはり有限性に反する。

$D$ に対し、三組の単元群と代数自己同型群の半直積の単位成分は以上により可解である。
各位置の $D$ 上でLie--Kolchinの定理による不変完全旗を取り、極大トーラスの
固有ベクトルで旗に適合する基底を選ぶ。連結可解群は極大トーラスと単冪根基の積で、
後者はこの基底で下単三角となる（旗の順を反転する）。
従って単冪根基のLie環と対角重みを用いれば先の消去証明がそのまま働く。
有限成分群の各代表を先に作用させれば全群の軌道が得られる。
この還元で非曲線状の局所環の同型型そのものを列挙したことにはならない。
\end{proof}

\path{computations/verify_curvilinear_window_orbits.py} は被約二十点の対称AS例と、
局所長さ4の点を含む長さ20の三空間で、基底変更・単元倍・共通の座標置換による
値の不変性を検査する。後者で一位置だけの接係数を変える対照では
標識付きの値が異なる。後者は軌道計算の試験入力であり、AS性は用いない。

% !TeX root = ../papers/nc-p3.tex
\subsection{接空間が三次元の有限点からのOre表示}
\label{sec:maximal-tangent-ore}
有限四辺で最も大きい接空間を持つ点を用いると、三生成部分の存在を仮定せずに
四生成の正則性を三生成の正則性へ還元できる。
以下の周期に関する点の同一視にはFrobenius標準四窓を使う。

\begin{theorem}\label{thm:maximal-tangent-ore}
標準有限Frobenius候補の四辺が全て有限で、ある辺の閉点の接空間が三次元とする。
その点から両方向に延びる点列を $[p_i]\in\mathbb P(V_i^*)$ とする。
\begin{enumerate}
\item $[p_{i+4}]=[p_i]$ であり、このような列は始点を固定すると高々二本である。
$H_i=\ker p_i$ によって内在的に三生成部分が定まる。
\item 各位置で一次基底を選ぶと、ある一定の三生成三関係の二次代数
$B=T(H)/(S)$ と写像 $d_i:H\to B_2$ により、候補の全関係は
\[
         S=0,\qquad X_i h-h X_{i+1}=d_i(h)\quad(h\in H)
\]
になる。ここで隣接する $H$ は選んだ同一視で表している。三次における条件は
\begin{equation}\label{eq:ternary-shift-derivation}
       (d_i\otimes 1+1\otimes d_{i+1})(S)=0\quad\text{in } B_3
\end{equation}
で、この候補では全始点で成立する。各区間の長さ $n$ の空間は
$\bigoplus_{r=0}^n B_{n-r}X^r$ と同型である。特に $B$ は実際に埋め込まれる。
\item 候補が標準14641型Koszul左右AS正則であることと、$B$ がHilbert級数
$(1-t)^{-3}$ の三次元quadratic Koszul左右AS正則代数であることは同値である。
逆構成ではこのような $B$ と任意の列 $d_i$ について
\eqref{eq:ternary-shift-derivation}だけで標準AS正則性を得る。
この逆構成全体には辺が正次元のものも含まれる。
\end{enumerate}
\end{theorem}
\begin{proof}
系\ref{cor:finite-edge-point-sequences}から点列は一意で、全位置で同じ三次元接空間を持つ。
$X_i$ を $p_i(X_i)=1$ と選ぶ。関係の $X_iX_{i+1}$ 成分は零である。
$R_i$ から $X_iH_{i+1}$ と $H_i X_{i+1}$ の二成分への写像はそれぞれ階数3で、
両方を合わせた写像は六つの双一次式の点 $(p_i,p_{i+1})$ における
Jacobi行列の転置である。接空間次元が3なので合わせても階数3である。
従って両核は同じ三次元空間 $S_i=R_i\cap(H_i\otimes H_{i+1})$ で、残る関係は
\[
 X_i h=f_i(h)X_{i+1}+d_i(h),\qquad f_i:H_{i+1}\xrightarrow{\sim}H_i
\]
と書ける。補正は $H_i\otimes H_{i+1}$ の $S_i$ を法とした類を使う。

まず $B^{(0)}$ を $H_i,S_i$ だけで定義したindexed代数とし、
$d=\dim((S_i\otimes H_{i+2})\cap(H_i\otimes S_{i+1}))$ とする。
純粋な三次の商次元は $27-18+d=9+d$ である。
混合語の $X$ を右へ動かせば、四生成の三次全体は
$(9+d)+6+3+1=19+d$ 個で生成される。
四生成の三次交叉 $W_{i,3}$ の次元4から三次商は20次元なので $d\ge1$ である。

一方、純粋な交叉は $W_{i,3}$ のうち第一位置を $p_i$ で縮約して零になる部分に入る。
定理\ref{thm:frobenius-tensor-reconstruction}により $W_{i,3}$ は四テンソルの
第四位置の縮約と同一視される。この写像に第一位置の $p_i$ を続けたものの核は
$R_{i-1}(p_i)$ の零化空間であり、直前の点 $kp_{i-1}$ である。
ここで第四位置は標準四周期性によって位置 $i-1$ と読んだ。
よって $d=1$ で、純粋な交叉の生成元は $w(p_{i-1})$ である。
その末端支持は次元3で、直前の辺の零点対から $\ker p_{i-2}$ に等しい。
純粋性から同じ支持は $H_{i+2}$ に含まれ、従って $[p_{i+2}]=[p_{i-2}]$ となる。
先頭支持も次元3である（命題\ref{prop:three-finite-edges}の証明）。
また各局所点環は埋込み次元3の完全交叉なので、次数2の次数商の次元は
$6-3=3$ 以上である。三乗イデアルが零ならこの商がsocleに入り
Gorenstein性に反する。従って長さは $1+3+3+1=8$ 以上で、全長20から個数を得る。

上述の20個の三次生成元は基底である。$X_i S_{i+1}$ を右へ還元すれば、
一つの $X$ の係数は $(f_i\otimes f_{i+1})S_{i+1}$ である。
還元結果が零なので $S_i$ に含まれ、両方三次元で等しい。
全整数上で $f_i$ に沿って順次 $H_i$ を同一視すれば、$S_i=S$、
$f_i=1$ となる。この基底に四周期性を要求していない。
$X$ を含まない残りの係数がちょうど\eqref{eq:ternary-shift-derivation}である。

この三次条件で高次に新しい障害がないことを直接示す。
$d_i$ を積の規則 $D_i(ab)=D_i(a)b+a D_{i+\deg a}(b)$、
$D_i(1)=0$ により自由代数から $B$ へ延ばす。
全始点で関係 $S$ を殺すので、その両側イデアルも殺し $B$ 上に降りる。
位置ごとに $bZ^r$ の形式的直和を作り、左から $h$ は係数に掛け、
$X_i$ は $bZ^{r+1}+D_i(b)Z^r$ に送る。積の規則によって全関係が作用で零になる。
右へ還元した語を終点の $1$ に作用させるとその係数を復元できるので、
表示した直和が全次数で成立する。$X$ の個数で濾過した随伴次数は
中心変数 $Z$ 一個を加えた $\check{B[Z]}$ となる。

$B$ が定理の三次元正則代数なら、$B[Z]$ は二つの線型解消のテンソル積によって
四次元Koszulで、二次双対は $B^!$ と一生成外積代数の次数付きテンソル積で
14641型の完全対合を持つ。元の候補で各内部次数のbar複体を $X$ の総数で濾過する。
濾過は有限で、その随伴次数は直前の代数のbar複体である。
そのホモロジーは内部次数とホモロジー次数が等しいところだけにあるので
スペクトル列はそこで退化し、元の代数もKoszulである。
barの分割余積が濾過と両立するから、有限次元の対角部分を双対化して得る
Extの随伴次数の積も上記の二次双対の積になる。
従って元の二次双対にも同じ次元と完全対合が持ち上がり、左右AS条件を得る。
Hilbert級数は直和表示から $(1-t)^{-4}$ である。

逆に元の代数 $A$ が標準ASなら直和表示から $H_B(t)=(1-t)^{-3}$ である。
$B$ の右自明加群の極小次数付き自由解消を、始点を固定して
$\check B\to A$ に沿って誘導する。各終点では左 $B$ 側で上の直和を取っているので
誘導は完全であり、微分の正次数性も保たれる。
得られるのは右点加群（各長さで $kX^r$）の極小解消である。
$A$ の大域次元4により $g=\operatorname{pd}_B k\le4$ となる。
以下、無限個の生成元を持つ次数付き加群にも使える形で、$B$ の両側のsocleが零であることを確認する。
左右の $k$ の極小解消の最大長は $\operatorname{Tor}^B(k,k)$ の最終次数で一致する。
従って任意の下に有界な次数付き右加群 $M$ に対して
$\operatorname{Tor}^B_{>g}(M,k)=0$ であり、左加群でも同様である。
正の大域次元 $g$ で同次socle元があれば $k(-a)\hookrightarrow B$ があり、
その余核の $\operatorname{Tor}_{g+1}$ は $\operatorname{Tor}_g(k(-a),k)\ne0$ に同型となる。
これは矛盾である。

$B$ の生成元、三関係、唯一の純粋な三次交叉から右Koszul複体の先頭
$B(-3)\to B(-2)^3\to B(-1)^3\to B\to k$ を作る。
最左写像の三係数は交叉生成元の末端支持が全三次元であるため一次生成元の基底である。
従って左socleが零であることから単射である。
生成関係の定義から $B(-1)^3$ 以降は完全で、残る一箇所のホモロジーは
既知のHilbert級数とEuler等式から零である。逆側も先頭支持を使って同じ結論を得る。
両端の三次元支持は $B^!$ の次数1と2の両向きの積を完全にする。
以上で $B$ は三次元quadratic Koszul左右AS正則である。
\end{proof}

\begin{corollary}\label{cor:maximal-tangent-ore-isomorphism}
上記の点列を指定した二対象の正規表示を $(B,d_i),(B',d'_i)$ とする。
指定列を保つ添字保持同型の存在は、次数付き代数同型 $G:B\to B'$ と
各位置の $v_i\in B'_1$ が存在して
\[
       Gd_i(h)=d'_i(Gh)+v_i(Gh)-(Gh)v_{i+1}\quad(h\in H)
\]
となることと同値である。列を指定しない場合は、始点の指定点の像として
標的の高々二つの三次元接点を試せば十分である。
\end{corollary}
\begin{proof}
指定点を保つ一次同型は $h\mapsto g_i h$, $X_i\mapsto c_i X'_i+u_i$
の形である。$c_i\ne0$ で、混合関係の $X'$ 係数から
$c_i g_{i+1}=c_{i+1}g_i$ となる。
各位置の一次写像を $c_i$ で割っても二次関係の像は同じである。
従って $g_i/c_i$ は一定の $G$ となり、純粋な三関係を対応させる。
残る係数が表示した式である。逆にこれらの式は全関係を対応させ、一次写像が可逆なので同型を与える。
\end{proof}

\path{computations/verify_maximal_tangent_ore.py} では、
四式 $x^2+4yz+2wx,y^2+4xz,z^2+4xy,w^2$ の対称AS窓で辺の有限性を検査する。
平方点での接空間、純粋な三関係、三次交叉の両端支持、四生成の三次元数を有理数で計算する。
一つの位置の混合補正だけを変えると、同じ接空間と純粋関係を保っても三次元数が20未満になる。

\begin{theorem}[一位置の線型パラメータと移動作用]\label{thm:ternary-ore-linear-space}
三次元Koszul左右AS正則な $B=T(H)/(S)$、$H_B(t)=(1-t)^{-3}$ を固定する。
$D=\operatorname{Hom}(H,B_2)$ は18次元で、$M=\operatorname{Hom}(S,B_3)$ に値を取る写像
$\lambda d=(d\otimes1)|_S$, $\rho d=(1\otimes d)|_S$ を積で定める。
$T=\{(d,e):\lambda d+\rho e=0\}$ と置き、
\[
 E_0=D,\qquad E_{n+1}=\operatorname{pr}_1(T\cap E_n^2)
                 \cap\operatorname{pr}_2(T\cap E_n^2)
\]
を安定するまで計算して $E$ とする。$L_v(h)=vh$, $R_v(h)=hv$ と書けば
\[
                 P_B=E/(L_H+R_H)
\]
は線型空間である。$B$ の同一視と点列を指定した上のOre表示の同型類は
$P_B$ と全単射である。$B$ の同一視を外すと $\Aut_{\mathrm{gr}}(B)$ の作用で割る。
点列も外した有限四辺の枝では、標的の高々二つの指定点でこの条件を試す。
$T\cap E^2$ は $P_B$ 上に一つの線型自己同型 $F_B$ のグラフを誘導し、
添字を一つ移すとパラメータは $F_B$ で移る。
\end{theorem}
\begin{proof}
各段階で計算するのは部分空間の交叉と射影であり、真に減る回数は高々18である。
両方向に延びる列の各項は常に $E_n$ に入る。逆に安定後は
$T\cap E^2$ の両射影が $E$ に全射で、任意の点を前後に延長できる。
式\eqref{eq:ternary-shift-derivation}により、これが全ての実現可能な初期値である。
$L_{v_i}-R_{v_{i+1}}$ は任意の列 $v_i\in H$ に対して同式を満たし、
従って $L_H+R_H\subset E$ である。

さらに
\[
                   \ker\rho=R_H,\qquad \ker\lambda=L_H .
\]
実際、$\rho d=0$ は生成元上で $h\mapsto d(h)$ とした次数1の
左 $B$ 加群写像 $B_{\ge1}\to B$ が定まることと同値である。
左イデアルの表示の最初の関係は $S$ が生成するためである。
$0\to B_{\ge1}\to B\to k\to0$ と左右AS条件
$\operatorname{Hom}_B(k,B)=\operatorname{Ext}^1_B(k,B)=0$ により、
この写像は $B\to B$ に一意に延びるので右乗法である。反対側も同様である。

二つの列の初期項が $P_B$ で一致すれば、位置0,1の $v_0,v_1$ で補正して
初期項を一致させる。次の差は $\ker\rho$ に入り $v_2$ だけで消せる。
順方向にこれを繰り返し、逆方向は $\ker\lambda$ と $v_{-1},v_{-2},\ldots$ を使う。
従って一位置の一致で全位置の同型が実現する。必要性は前系による。
この議論は指定した点列と $B$ を保つ場合には辺の有限性を使わない。
同じ初期類からの全ての次項が同じ類を与え、直前の項にも同じことが成立するので、
線型関係の像は両方向に一価、すなわち線型自己同型のグラフになる。

有限四辺の条件を課すには、$E$ 内で次項を三回取って四つの辺の六双一次式の
零点スキームの有限性を調べればよい。実現した対象は標準ASで四周期であり、
初期項が等しい全ての延長は指定点を保って同型なので、この検査は選択によらない。
指定点を保つ最小周期は $j=1,2,4$ の順に
$F_B^j[d]$ が $[d]$ と同じ自己同型軌道にあるかで得られる。
有限四辺では指定点を標準四移動が固定することをすでに示した。
指定点を外す場合は同じ移動先の他の三次元接点も候補に入れる。
\end{proof}

上記有理例の三生成部分では $E_n$ の次元は $18,15,15$、$L_H+R_H$ は六次元である。
九次元の商で $F_B^2=1$ となり、固有値 $+1,-1$ の重複度は順に $3,6$ である。
反可換三生成基底でも同じ次元と重複度を得る。可換三生成基底では
$E=D$、補生成元の取り替えによる部分空間は三次元で、十五次元の商上の移動は恒等である。
ここではすべて $X_i h-hX_{i+1}$ の符号規約を使い、
$B$ の自己同型群による商を取る前の数値を記している。

% !TeX root = ../papers/nc-p3.tex
\subsection{三次元接点の基底が量子多項式環となる有限四辺の分類}
\label{sec:quantum-tangent-windows}
前節の $B$ が量子多項式環
$B_Q=k\langle x_0,x_1,x_2\rangle/(x_jx_k-q_{jk}x_kx_j)$ になる枝を扱う。
$q_{jk}q_{kj}=q_{jj}=1$ とし、$q_{01}=p,q_{12}=q,q_{20}=r$ とする。
この基底は順序付き単項式の基底と量子外積双対を持つ三次元Koszul左右AS正則代数である。

\begin{theorem}\label{thm:quantum-tangent-finite}
有限四辺で三次元接点からこの形の $B_Q$ が得られる全AS対象は、
次の正規表示で尽くされる。添字 $a,b,c$ は $0,1,2$ の巡回順である。
\begin{gather*}
 p=q=r=\zeta,\qquad \zeta=-1\ \text{または}\ \zeta^2=-1,\\
 X_n x_a-x_a X_{n+1}
  =\sum_{s=0}^2 C_{a,s}\,\zeta^{n(s+1)}x_b^s x_c^{2-s}.
\end{gather*}
全 $C\in k^{3\times3}$ が標準AS正則である。四辺が有限である必要十分条件は
\[
 (C_{a,0},C_{a,1},C_{a,2})\ne0\quad(\forall a),\qquad
 (C_{j-1,2},C_{j-2,0})\ne(0,0)\quad(\forall j)
\]
であり、行添字は法3で読む。三次元接点は一意である。
\end{theorem}
\begin{proof}
量子トーラスで $x_j x^\alpha=\chi_j x^\alpha x_j$、
$\chi_j=\prod_k q_{jk}^{\alpha_k}$ と書く。
写像 $d_n$ を多次数で分け、$d_n(x_j)=c_j(n)x^\alpha x_j$ と置く。
$\alpha+e_j$ が非負でない場合は $c_j(n)=0$ とする。
式\eqref{eq:ternary-shift-derivation}は共通の量子単項式の係数を取ると
\[
          c_j(n)-c_k(n)=\chi_k c_j(n+1)-\chi_j c_k(n+1)
\]
となる。$\alpha=e_a$ のとき、$\chi=(\chi_0,\chi_1,\chi_2)$ と
$\mathbf1=(1,1,1)$ が独立なら、
$\mathbf1\wedge k^3$ と $\chi\wedge k^3$ の交叉は
$k(\mathbf1\wedge\chi)$ であり、$c(n)$ は $\langle\mathbf1,\chi\rangle$ に入る。
これは $L_H+R_H$ のこの次数の部分である。
比例する場合 $\chi_a=1$ なので $x_a$ は中心的で、
$k^3/k\mathbf1$ の二次元が残り、移動は恒等である。

他の次数は $\alpha=-e_a+s e_b+(2-s)e_c$、$s=0,1,2$ であり、
ただ一つの成分 $c_a$ が許される。
非零な両側列の存在は $\chi_b=\chi_c$ に同値で、成立すれば
$c_a(n+1)=\chi_b^{-1}c_a(n)$ となる。
条件は $q_{bc}^2=q_{ba}/q_{ca}$ で、$s$ によらない。
成立する軸では三次元が残り、この次数に補生成元の変更は作用しない。
従ってこの計算は任意の非零 $p,q,r$ に対する $P_{B_Q}$ も決めている。

軸 $a$ の条件が破れれば、$x_a=0$ 上に制限した $d_n(x_a)$ は恒等的に零となる。
実際、非負多次数の項は $x_a$ を含み、負の次数の項は排除される。
座標直線 $[u_b:u_c]$ にある基底の点とその次点を
$u=(u_b,u_c)$、$v=(u_b,q_{bc}u_c)$ とする。混合関係は
二つの追加座標 $z,z'$ について $zv-uz'=d_n(u,v)$ である。
$q_{bc}\ne1$ なら稠密開集合で二行が独立になり全点が持ち上がる。
$q_{bc}=1$ なら二行の両立条件は二元三次式一本であり、
射影直線上の少なくとも一点で解が一次元になる。
従って有限四辺には三軸全ての条件が必要である。
その三式は $pqr\cdot q=pqr\cdot r=pqr\cdot p=1$ なので、
$p=q=r=\zeta$, $\zeta^4=1$ に限る。

$\zeta=1$ なら射影平面上の $u=v$ について $d_n(u,u)$ は三つの二次式である。
共通零点があればそこで、なければこれらが定める次数2の自己射の固定点で、
$d_n(u,u)$ は $u$ に比例する。固定点の存在は射影平面の積でグラフと対角の交叉数が
$1+2+4=7$ であることから従う。そこで $z,z'$ に一次元の解があり有限でない。
残る $\zeta$ には中心的な一次生成元がなく、負の次数の三ブロックだけが残る。
$\chi_b^{-1}=\zeta^{s+1}$ となり、単項式の定数を吸収すれば定理の式を得る。
前節によって全係数のAS正則性が従う。

$\zeta^3\ne1$ のため、三つの純粋関係を満たす基底の射影零点対は
座標三角形上のグラフである。三座標が全て非零なら巡回積で
$\zeta^3=1$ を得て矛盾し、各直線では先の $u,v$ の表示になる。
スキームとしても二座標非零の開部分では、二つの式で標的を消去した残る式が
$(1-\zeta^3)u_0u_1u_2=0$ の単元倍である。
基底への射影が一方の位置で零なら混合関係から両方で零となり、指定点だけを得る。
他では軸 $a$ を欠く直線の開トーラス上で $z,z'$ が一意に解ける。残る式は始点0で
\[
 g_a(u_b,u_c)=C_{a,2}u_b^2+\zeta C_{a,1}u_bu_c+\zeta C_{a,0}u_c^2=0
\]
である。頂点 $j$ は他の二つの混合関係に現れる $x_j^2$ の係数がともに零なら
一次元の持上げを持ち、一方が非零なら持上げを持たない。
各始点で非零係数の位置は同じなので、主張した二条件を得る。
指定点で混合三式の微分は階数3、他の局所環は直線上の二次式の根の局所環なので
三次元接点は一意である。
\end{proof}

\begin{corollary}\label{cor:quantum-tangent-invariants}
有限な正規表示の添字保持同型は以下の作用で厳密に決まる。
対角基底変更 $(t_0,t_1,t_2)$ は各 $C_{a,s}$ に
$t_b^s t_c^{2-s}/t_a$ を掛ける。偶置換は三つの行を巡回させる。
奇置換は $\zeta$ を $\zeta'=\zeta^{-1}$ に送り、対応する行で
$s$ 列を $2-s$ 列に移し、係数に $(\zeta')^{-s(2-s)}$ を掛ける。
従って原始四乗根は一方だけ採用すればよく、$-1$ 型と原始四乗根型は非同型である。
各型は既約な六次元の同型類の族で、一般最小周期はそれぞれ2と4である。
\end{corollary}
\begin{proof}
一意な点から $B_Q$ と $P_{B_Q}$ が復元される。
基底の次数付き同型は点スキームの座標三角形を保つので単項行列であり、
関係を代入すると示した置換だけが可能になる。
奇置換の係数は $x_{c'}^s x_{b'}^{2-s}$ を巡回順に戻す交換係数である。
従って指定点を外しても群作用の条件は必要十分である。
任意の零係数を許して、有限個の置換を試した後、非零座標の台を比較し、
重み行列の転置の整数核から来るLaurent単項式を比較すればよい
（定理\ref{thm:curvilinear-window-orbits}のトーラスの証明）。
移動は各列への $\zeta^{s+1}$ 倍なので、この判定を
$j=1,2,4$ 回の移動に順に用いれば特殊値の最小周期も決まる。

有限条件は九次元係数空間の非空開集合である。全係数が非零なら
重みの差により三つの軸の差を張り、各重みの成分和が1なので重み階数は3である。
従って像の次元は $9-3=6$。一般には三つの比
$C_{a,1}^2/(C_{a,0}C_{a,2})$ は互いに異なる。
これらは移動と対角変換で不変で、同じ型の置換で行に沿って移る
（$\zeta=-1$ の逆順の符号も平方で消える）。従って一般点の周期判定では置換は恒等である。
隣り合う列の重み比から、$j$ 移動を実現するには全巡回順で
$t_b/t_c=\zeta^j$ が必要で、三つを掛けると $\zeta^{3j}=1$ を得る。
四乗根なので最小の $j$ はその位数であり、そこで表示そのものが戻る。
\end{proof}

有限な各始点で、$S_a=\{s:C_{a,s}\ne0\}$ と置けば指定点の長さは
\[
                  20-\sum_a(\max S_a-\min S_a)
\]
である。実際、各開トーラスでの根の重複度の和は表示した幅であり、
頂点の上に点はなく、全長は20である。
幅2の二根は $\zeta^2 C_{a,1}^2-4\zeta C_{a,0} C_{a,2}=0$ で二重点になる。
始点を移すと判別式に非零定数が掛かるだけである。
端の係数を零にして有限条件を保つ特殊化では、失われた開トーラス上の長さが
指定点に加わる。全ての行で第0列だけが非零の場合は長さ20の一点となり最小周期は1である。
原始四乗根で三行を $(2,0,5),(7,0,13),(17,0,23)$ とすれば最小周期2の例になる。
\path{computations/verify_quantum_tangent_windows.py} はこれらの同型・周期判定と
$P_B$ の独立な行列計算を検査する。九係数の512通りの台のうち有限なものを
指定点の長さ14から20の順に数えると $8,24,54,50,36,6,2$ 通りである。
この数は標識付き台の数であり、同型類数ではない。

% !TeX root = ../papers/nc-p3.tex
\subsection{三次元接点の三生成基底が滑らかな平面三次曲線を持つ場合}
\label{sec:elliptic-tangent-windows}
三生成基底 $B=T(H)/(S)$ の三つの双一次関係が
滑らかな平面三次曲線 $E\subset\mathbb P(H^*)$ 上の自己同型 $\sigma$ のグラフを
定める枝を扱う。以下 $L=\mathcal O_E(1)$ とする。

\begin{theorem}\label{thm:elliptic-tangent-classification}
この枝で有限四辺を持つ標準AS窓の基底は、平面三次曲線 $E$ と、
その上の位数 $m=2$ または $4$ の平行移動 $\sigma$ から得るものに限り、全て実現する。
各基底に対する全補正は定理\ref{thm:ternary-ore-linear-space}の九次元空間 $P_B$ である。
非零の全ての補正が有限四辺を与え、零のみ辺が正次元である。
非零な射影補正は次数12の有効因子 $D$ として復元できる。
二つの窓の添字保持同型は、$(E,L,\sigma,D)$ を対応させる曲線同型の存在と同値である。
最小周期は $\sigma^jD$ を $D$ と同じ $\Aut(E,L,\sigma)$ 軌道に送る最小の $j>0$ である。
二つの位数の族はそれぞれ既約で同型類の次元は9、一般最小周期は $m$ である。
\end{theorem}

\begin{proof}[滑らかな曲線上での四移動の必要性]
定理\ref{thm:maximal-tangent-ore}によって $B$ は三次元標準二次ASである。
唯一の三次交叉を張るテンソル $w$ の先頭・末端の縮約を
$\ell,r:H^*\xrightarrow{\sim}S$ とする。連続する三点を $p,q=\sigma p,u=\sigma^2p$
とし、混合補正 $d_i$ の $(p,q)$ での評価を $v_i\in H^*$ とする。
$p,q$ が独立なら、二つの第四座標が存在するのは $v_i\in\langle p,q\rangle$ の場合だけである。
$\sigma$ が恒等なら、三次曲線には消える二次式がないため $S$ は交代式の全空間となり、
そのグラフの台は平面全体になって矛盾する。従って一般に $p,q$ は独立である。
有限辺なら一般の $p$ で $v_i\notin\langle p,q\rangle$ である。

$A_u(v)=(s(v,u))_{s\in S}$、$B_p(v)=(s(p,v))_{s\in S}$ と書く。
両写像の核は $kq$ で、像は $S^*$ の平面である。
式\eqref{eq:ternary-shift-derivation}を $(p,q,u)$ で評価すれば
$A_u(v_i)+B_p(v_{i+1})=0$ となる。
$A_u(p)=B_p(u)\ne0$（一般に $u\ne q$）で、$A_u(v_i)$ はこれと独立である。
よって二平面は等しい。その零化直線はそれぞれ $kr(\sigma u)$ と
$k\ell(\sigma^{-1}p)$ である。実際 $w(-,u,\sigma u)=w(\sigma^{-1}p,p,-)=0$ である。
従って全曲線上で $\sigma^4=\mathbb P(r^{-1}\ell)|_E$ である。

もし $\sigma p=p$ なら $r(p)$ と $\ell(p)$ は比例する。
線型行列 $B_p$ の階数は2で、その核は $kp$、余核の双対は $k\ell(p)$ である。
従って行列式の微分は $v\mapsto\ell(p)(v,p)$ の定数倍だが、
これは $r(p)(v,p)=w(v,p,p)=0$ に比例する。行列式は $E$ を定める三次式なので
滑らかさに反する。従って固定点はなく、$\sigma$ は非恒等な平行移動である。
実際、楕円曲線の自己同型の原点保持部分 $\alpha$ が非恒等なら $1-\alpha$ は全射である。
\end{proof}

\begin{proof}[基底の表示と全ファイバーの正則性]
変曲点 $O$ を原点にしたHesse表示
$x^3+y^3+z^3-3hxyz=0$、$O=[1:-1:0]$ を使う。
この表示は標数零の代数閉体上の滑らかな平面三次曲線を覆う。
例えば変曲点を原点にWeierstrass表示にすれば埋込みは $|3O|$ であり、
Hesse族の $j=27h^3(h^3+8)^3/(h^3-1)^3$ は射影直線から全射で
特異パラメータは極だけである\cite[\S7.3]{KLV}。
$\sigma$ の平行移動点 $P=[a:b:c]$ は三等分点でない。
三等分点なら $\sigma$ は平面に延び、二次式の制限の単射性から
$S$ は射影変換の平面グラフの三関係となる。従って $abc\ne0$ である。
ここで座標軸との九交点は全て変曲点であり、接線を代入しても確かめられる。

文字を $x_0,x_1,x_2$ として、三つの関係を巡回順に
\[
          c x_j^2+a x_{j+1}x_{j+2}+b x_{j+2}x_{j+1}
\]
とする。左点 $p=(x,y,z)$ に対する行列の三行は
$(cx,bz,ay),(az,cy,bx),(by,ax,cz)$ である。
行列式は $(a^3+b^3+c^3)xyz-abc(x^3+y^3+z^3)$ で、$E$ 上の全点で階数2である。
逆方向も $a,b$ を交換すれば同じで、グラフを定める。
次点 $q=(v_0,v_1,v_2)$ に対し、最初の二関係の差は
$p,P,(v_1,v_0,v_2)$ の共線条件である。
最初の二座標の交換は $O$ に関する群逆（$q$ と交換像を結ぶ直線は $O$ を通る）なので、
このグラフは $P$ の平行移動を与える。
グラフで消える双一次式の空間は三次元である。
実際、$L\not\simeq\sigma^*L$ なので、$L$ の基点自由な二切断の完全列を
$\sigma^*L$ で捩ると、次数零の非自明な束の $H^1$ の消滅により
二次の積は次数6の全切断に全射となる。従って元の三関係もこの表示になる。
三関係には巡回対称な三次交叉 $w$（各関係に対応する文字を左から掛けて足したもの）があり、
$\ell=r$ である。先の必要条件は $4P=0$ となり、残る位数は2と4である。

逆にこれらの $P$ の全てで三生成基底が正則となることを示す。
$L_j=(\sigma^j)^*L$、$L_{[j,n)}=\bigotimes_{j\le h<n}L_h$ と置く。
$R_n=H^0(L_{[0,n)})$ に積 $f\cdot g=f(\sigma^{\deg f})^*g$ を入れる。
$n\ge3$ で、$B$ の生成元・関係・三次交叉による右複体を曲線上で評価した列
\[
 0\longrightarrow L_{[3,n)}\longrightarrow S\otimes L_{[2,n)}
  \longrightarrow H\otimes L_{[1,n)}\longrightarrow L_{[0,n)}\longrightarrow0
\]
は束として完全である。最後の写像は $p$ での評価、中央写像は第二文字に
$\sigma p$ を代入したもの、最初の写像は $w$ の末尾に $\sigma^2p$ を代入したものである。
各点で合成は零で、階数が順に $1,2,1$ なので完全となる。
$n>3$ では全項の $H^1$ が零で、大域切断の列も完全である。
$n=3$ では左端の $H^1(\mathcal O_E)=k$ だけが残り、
大域切断では $H\otimes H^0(L_{[1,3)})$ の位置に限り一次元のホモロジーがある。
$n=2$ での積の全射性と核は上の三関係の計算である。
従って $R$ は一次生成で $B\to R$ は全射である。

$k$ の右 $B$ 自由解消の最初の二微分は生成元と $S$ で与えられる。
その第三項に入れる全てのsyzygyに先の $w$ も含められる。
従ってこの解消を $R$ でテンソルしたときの第1ホモロジーは今計算した
$k(-3)$ であり、第2ホモロジーは零である。後者は $w$ の像だけですでに核を埋めるためである。
次数付き左 $B$ 加群 $R$ の極小解消から、$B\to R$ の核は
$B(-3)$ と同型で、右からある三次元 $g$ を掛けて埋め込まれる。
よって $H_B(t)=H_R(t)/(1-t^3)=(1-t)^{-3}$ である。
$R$ は整域であるので、$B$ の左socle元は $b'g$ と書け、右からの $g$ の単射性により
$b'$ も左socleに入る。次数降下で零となる。反対側も同じ証明である。
$w$ の両端支持が全三次元であるので、定理\ref{thm:maximal-tangent-ore}の末尾と同様に、
三次までのKoszul複体の最左写像は単射で、Hilbert級数から全列が完全となる。
完全な相補積から三次元左右AS条件も従う。
\end{proof}

\begin{proof}[九次元空間の一様な計算]
$a=1$ とし、対角変数変更と各関係のスカラー倍で三つの平方係数を $1,1,u$、
$u=c^3$ にできる。実際、各平方係数は変数倍率 $\lambda_j$ により
$c\lambda_j^3/(\lambda_0\lambda_1\lambda_2)$ となる。
位数2なら $a=b=1$、$u=t$ と置く。
位数4では $b=t$、$u=t(1+t^2)/(1+t)$ となる。
実際、$2P$ の座標は行列の二行の外積から
$(b(c^3-a^3),a(b^3-c^3),c(a^3-b^3))$ である。
三座標の三乗が全て等しい場合は特異なので除いてある。
$2P$ が $O$ でない二等分点となる条件は前二座標の一致であり、表示式を与える。
各場合 $bu\ne0$、$(1+b^3+u)^3\ne27b^3u$ を課し、位数4では分母 $1+t\ne0$ とする。
$t=1$ や $t^2+t+1=0$ は位数4の滑らかな領域に入らない。

$\omega$ を原始三乗根とする。$x_j\mapsto\omega^j x_j$ の三つの重みに分けて
$\lambda,\rho$ を計算できる。三つの重みの計算は、平方係数を一定に戻した基底での
文字の巡回置換により対応する。実際、二操作の交換子は文字全てへの同じ三乗根倍であり、
次数1の補正空間上で重みを一つ移す。
重み0では各 $x_j$ の二次補正に、$i+k=j\bmod3$, $i\ne k$ の二単項式 $x_ix_k$
を使う。平方は関係で消えるので六変数である。$\lambda,\rho$ は各々 $10\times6$ となる。
有理関数体上の厳密計算は次表となる。
\[
\begin{array}{c|c|c|c}
m&[\lambda\ \rho]\text{ の六次小行列式}
     &\mu&\rk(\lambda+\mu\rho)\\ \hline
2&-t(t-1)&1,-1&4,3\\
4&-t^2(t-1)(t^2+1)(t^2+t+1)/(t+1)^2&-1,\sqrt{-1},-\sqrt{-1}&4,4,4
\end{array}
\]
\path{computations/verify_elliptic_tangent_windows.py} は各三次関係の重みの
九つの三次語から $(S\otimes H)+(H\otimes S)$ の行を零化する汎関数を計算し、
その行に各補正を入れてこの表を再現する。小行列式に使う汎関数の分母の最小公倍式は
それぞれ $1,t(t+1)$ である。従って示した小行列式は全ての許される $t$ で有効かつ非零である。
一方、三次商が定数階数なので最後の欄の一般階数は全特殊値での上界である。

$\rho$ の核は右乗法のこの重みの一方向なので階数5である。
連結行列の階数下界6から、次項を許す初期値の次元は高々5である。
左右乗法の二方向は独立である。比例すれば $x_0 h$ と $h x_0$ が全ての文字で比例するが、
非零の平方係数を持つ三関係を比較すると $x_0x_1$ と $x_1x_0$ は独立である。
ゆえに商は一重み当たり高々3次元である。
$\ker(\lambda+\mu\rho)$ の元は $d_i=\mu^i d$ で両側に延び、移動の $\mu$ 固有類を与える。
補生成元の変更の二方向のうちこの核に入るものは $L_v-\mu R_v$ の一直線だけである。
実際 $\lambda R_v=\rho L_v$ は零でなく（$\ker\rho=R_H$）、$\lambda L_v=\rho R_v=0$ である。
よって表から $m=2$ の固有次元は少なくとも $1,2$、
$m=4$ ではそれぞれ少なくとも1となる。異なる固有値は独立なので全て等号であり、
三重みを合わせて常に九次元となる。特に検算に使った有理点への特殊化だけに依存しない。
\end{proof}

\begin{proof}[有限条件、同型、特殊化と周期]
$N=L\otimes\sigma^*L$ と置く。$B_2=H^0(N)$ なので、補正 $d$ の評価ベクトル
$v(p)$ は $H^*\otimes H^0(N)$ に属する。
$p,\sigma p$ が全点で独立であることから
\[
 0\longrightarrow \sigma^*L\oplus L\longrightarrow H^*\otimes N
 \xrightarrow{\ \det(p,\sigma p,-)\ } N^2\longrightarrow0
\]
は完全である。最初の写像は二つの係数を $p,\sigma p$ に掛けて足す。
$H^1$ の消滅と一次の完全線型系から、大域切断での核はちょうど左右乗法である。
従って $P_B$ は $H^0(N^2)$ 内の九次元部分空間と同一視される。
非零切断の零因子 $D$ は長さ12である。各零点では独立な二行から第四座標二つを解けるため、
持ち上がった局所環は曲線上の因子の局所環と一致する。
一方だけで基底への射影が零なら他方でも零となる。
残る指定点は有限完全交叉の全長20から長さ8で、唯一の三次元接点である。
特殊値で重根が生じれば重複度がそのまま残余局所長になる。
切断零では全曲線が持ち上がる。

一意な点と三生成部分を保つ同型は $(E,L,\sigma)$ の射影同型を誘導し、
逆にこの同型はグラフで消える三関係を対応させる。
同じ因子を与える非零切断の比は定数であり、この定数は文字全てを同じ定数倍する
基底自己同型で実現できる。従って因子を含む表示が全ての特殊値でも必要十分である。
$\Aut(E,L)$ は有限である。例えば九つの変曲点の像を試し、四点の射影枠で変換を定めて
曲線・グラフ・因子を比較できる。隣接辺のスキームが同じ点列を与えるので、
次の残余因子は重複度を込めて $\sigma D$ となる。

各位数の基底パラメータは上の $t$ の非空開部分で尽くされる。
その像は一次元である。Hesseの $j$ は各場合に非定数であり、
固定した楕円曲線では変曲原点と所定の平行移動点はいずれも有限個である。
補正の核と商は上の定数階数計算によりベクトル束をなし、射影ファイバーは八次元、
残る射影群は有限なので像は既約九次元である。
固定基底について、九次元切断空間の線型系の像は $\mathbb P^8$ を張る曲線で次数8以上である。
固定零因子を除いても移動部分の次数は高々12なので、この写像は像へ双有理である。
もし一般因子の $j$ 移動を有限群 $\Aut(E,L,\sigma)$ で戻せれば、
群のある一元が切断の全射影空間でこの移動と一致する
（各候補の一致は射影線型な閉条件）。その元と $\sigma^j$ の差は全因子を保ち、
上記の双有理写像により恒等である。平行移動 $\sigma^j$ が $L$ を保つには
$3jP=0$ が必要十分で、$m=2,4$ では $m\mid j$ である。これが一般最小周期を与える。
位数2と4は滑らかな領域内では互いに特殊化しない。変曲原点を局所的に選んだ後の
四等分点の群スキームは有限エタールで、二等分点はその開かつ閉な部分である。
\end{proof}

定理\ref{thm:algebraic-nine}の九変数族の一般点では、三生成部分の曲線は滑らかで、
移動は固定点のない対合であり、指定点は接空間三次元である。
従ってその九次元の像は上の位数2の九次元族に稠密である。
今回の因子による同型判定は、この共通部分で定理\ref{thm:algebraic-nine}の
次元と一般周期の計算を補う。全モジュライでの既約成分性はこの比較からは出ない。

% !TeX root = ../papers/nc-p3.tex
\subsection{特異三次曲線の三生成基底と楕円枝への平滑化}
\label{sec:singular-tangent-smoothing}
引き続き四辺が有限で三次元接点を持つ枝を扱う。
三生成基底の点対応が既約特異三次曲線のグラフの場合も決定できる。

\begin{theorem}\label{thm:nodal-tangent-windows}
三生成基底の対応が射影平面の射影自己同型のグラフ、または尖点三次曲線の自己同型の
グラフなら、この有限四辺の枝に入らない。既約節点三次曲線のグラフであれば、
正規化の座標を $p(t)=[t:t^2:1+t^3]$ と選んで移動は $t\mapsto\lambda t$ であり、
$\lambda=-1$ または $\lambda^2=-1$ である。基底の三関係は
\[
 xy-\lambda yx,\qquad yz-\lambda zy+\beta x^2,\qquad
 zx-\lambda xz+\beta y^2,\qquad \beta=(\lambda^3-1)/\lambda .
\]
どの基底にも九次元の全補正空間 $P_B$ がある。
有限条件は $v=(0,0,1)$ で $(d_0(x)(v,v),d_0(y)(v,v))\ne(0,0)$ であり、
独立な二本の線型式の共通核を除く。指定点は一意である。
各反転対 $\{\lambda,\lambda^{-1}\}$ の有限な射影補正は、
$t\mapsto ct$ と $t\mapsto c/t$（$c^3=1$）のうち移動を対応させるもので割ると
添字保持同型類をちょうど一度ずつ与える。二型の同型類の族は各八次元で、
一般最小周期は順に2と4である。
\end{theorem}
\begin{proof}
平面グラフ $\sigma=[G]$ では $S$ の評価は外積 $u\wedge G^{-1}v$ で与えられる。
実際、このグラフで消える双一次式の空間は三次元である。
$p,Gp,G^2p$ で前節の二つの縮約平面は
$Gp\wedge H^*$ と $p\wedge H^*$ となる。
$Gp\not\in kp$ なら異なるため、補正の評価は $\langle p,Gp\rangle$ に入り、
そのような全ての点で第四座標を解ける。恒等の場合の固定点による除外は
定理\ref{thm:quantum-tangent-finite}の証明による。

既約三次曲線で移動が恒等なら、対角で消える双一次式は交代式だけなので
やはり平面グラフになる。一般点で非恒等なら前節の四移動の必要条件の証明は
そのまま使え、$\sigma^4$ は平面に延び、滑らかな固定点を持つことは許されない
（ここで固定点は $\sigma$ の固定点をいう）。
尖点三次曲線の正規化で滑らかな部分はアフィン直線なので、固定点のない自己同型は
非零の平行移動である。四乗も非零の平行移動で、平面三次曲線の唯一の滑らかな
変曲点を保たない。変曲点の一意性は尖点三次曲線 $y^2z=x^3$ の接線との交叉でも従う。

節点三次曲線の正規化では節点上の二点を $0,\infty$ と置く。
三次の斉次式の両端係数が所定の非零比を持つという一条件が降下条件なので、
$t$ の倍率と平面の基底を選べば $p(t)$ の表示を得る。
自己同型 $t\mapsto c/t$ は滑らかな固定点を持つので、移動は $t\mapsto\lambda t$ に限る。
$\lambda^3=1$ なら平面に延び、三次曲線には二次の方程式がないので平面グラフとなる。
$\lambda^3\ne1$ とする。$(p(t),p(\lambda t))$ を代入すると表示した三関係は零となる。
九つの評価積は $t,t^2,t^3,t^4,t^5,1+\lambda^3t^6$ を張るので核はこれで尽きる。
三次交叉は三つの関係を対応する文字で左から掛けて足した巡回テンソルであり、
前節の $\ell=r$ が成立する。従って $\lambda^4=1$ が必要となる。

逆に、この基底は量子平面 $xy=\lambda yx$ に $z$ を加えたOre拡大である。
自己同型は $(x,y)\mapsto(\lambda x,\lambda^{-1}y)$、補正は
$x\mapsto-\beta y^2$, $y\mapsto(\beta/\lambda)x^2$ であり、
$xy-\lambda yx$ に捩れた積の規則を適用すると零になる。
従って順序付き基底を持ち、$z$ の個数による随伴次数は三生成量子多項式環である。
定理\ref{thm:maximal-tangent-ore}のbar濾過の証明により標準三次元Koszul左右AS性を得る。
点対応の行列式は上の節点三次曲線を定め、節点でも行列は階数2なので
両方向の核直線によりグラフが得られる。

九次元と節点での二方向の全射性は下の六基底の同時計算で示す。
補正の節点での評価を $w_n=d_n(v,v)$ と書く。その $kv$ を法とする類は補生成元変更に依存しない。
$v$ が固定点なので、類が零なら節点の上に一次元の解がある。非零なら節点の上に解はない。
節点における式\eqref{eq:ternary-shift-derivation}の両縮約行列の核は $kv$ なので、
この非零性は次の位置と同値である。滑らかな部分では四座標の残りの条件は
\[
                 f_n(t)=\det[p(t),p(\lambda t),d_n(p(t),p(\lambda t))]=0
\]
という一変数の式である。節点の $y$ 成分が非零なら $t=0$ の近傍で
$f_n$ の一次係数が非零である。$x$ 成分が非零なら $t=\infty$ の自明化で
同様である。これは二枝の接線がそれぞれ $x,y$ 方向で、移動の微分が
$\lambda,\lambda^{-1}\ne1$ であることによる。従ってこの式は恒等零にならない。
各根の局所環は重複度を長さとする曲線状環、指定点の長さは
$20-(\deg f_n-\operatorname{ord}_0 f_n)$ であり、三次元接点は一意である。

節点を保つ正規化の変換のうち、埋込みの三次切断空間を保つのは表示した六つだけである。
反転は $\lambda$ と $\lambda^{-1}$ を交換する。定理\ref{thm:ternary-ore-linear-space}
から補正にはその変換の誘導作用を使えばよく、共通スカラーは射影化で除かれる。
従って零係数や重根でも有限群による一致が同型に必要十分である。
特殊値の最小周期は同じ射影軌道を与える $F_B^j$ の最小の正べきである。
有限な射影開集合は八次元である。
原始四乗根の型では比較に残る群の位数は3なので、固有値が下表の $F_B$ の
真のべき $j=1,2$ と射影変換として一致する元はない。
$-1$ 型で $F_B$ は全群と可換である。三乗根倍は節点評価の二方向に異なる
重みで作用するので、その位数は射影作用でも3である。
従って二面体群の射影像は忠実で中心が自明であり、非スカラーの $F_B$ に等しい元はない。
各群元との一致は閉条件なので一般最小周期が従う。
\end{proof}

\begin{theorem}\label{thm:ternary-torsion-smoothing}
上の節点基底の全ての有限窓と、定理\ref{thm:quantum-tangent-finite}の全ての有限窓は、
それぞれ同じ位数の滑らかな楕円基底を持つ窓の平坦特殊化である。
さらに二次曲線と横断的な直線の和に対応する下記の一係数基底でも同じことが成立する。
\end{theorem}
\begin{proof}
三関係を $x_{j+1}x_{j+2}+b x_{j+2}x_{j+1}+c_jx_j^2$ と書く。
$b=-\lambda$ を特別値とし、巡回順に $c_0,\ldots,c_{h-1}$ は非零、
残りは零とする（$h=0,1,2$）。対角変数変更で非零係数は1にできる。
$h=0$ は三角形、$h=2$ は前記の節点基底であり、$h=1$ では行列式が
$x_0((1-\lambda^3)x_1x_2+\lambda c_0x_0^2)$ の定数倍となる。
全基底の正則性は節点の場合と同じ量子平面上のOre表示から従う
（補正の一方または両方を零にしても関係上の導分条件は保たれる）。

実際の三次商を使う行列計算は
\path{computations/verify_singular_tangent_smoothing.py}にある。
平方係数を $1$ から $h$ 個置く上記六基底で、$[\lambda_B\ \rho_B]$ の階数は18である。
ここで $\lambda_B,\rho_B$ は定理\ref{thm:ternary-ore-linear-space}の二写像である。
$\rho_B$ は階数15なので、初期補正の次元は高々15、六次元の左右乗法を割れば高々9となる。
\[
\begin{array}{c|c|c|c}
\lambda&\mu&\dim\ker(\lambda_B+\mu\rho_B)&
       \dim (P_B)_{\mu}\\ \hline
-1&1,-1&6,9&3,6\\
\sqrt{-1}&-1,\sqrt{-1},-\sqrt{-1}&6,6,6&3,3,3
\end{array}
\]
各固有核では $L_v-\mu R_v$ の三次元を除く。各核の残りが独立なので上界に達する。
さらに $c_j=0$ の頂点 $v=e_j^*$ ごとに、全固有核から
$d(v,v)\bmod kv$ への写像の階数は2である。コードは二次商の基底を
語空間へ持ち上げてこの二行を直接計算する。

特別な窓の $\mu$ 固有補正の代表を一つずつ取る。
離散付値環 $k[[s]]$ 上で、零だった $c_j$ は $s$、非零だったものはその値に取る。
位数2の場合は $b=1$ を保つ。位数4の場合は $b(1+b^2)/(1+b)=\prod_j c_j$
を $b(0)=-\lambda$ から解く。この点で左辺の導関数は非零なので解が存在する。
一般ファイバーは前節の同じ位数の滑らかな楕円基底である。
各次数の三生成代数は特殊・一般とも同じHilbert級数を持ち、各次数で有限生成の
離散付値環加群なので自由である。従って混合補正の固有核も自由加群の写像の核として
計算できる。前節と上表で一般・特殊の核次元が等しいので、単元小行列を分裂して
全ての特殊固有ベクトルを持ち上げられる。
その持上げを $\mu^n$ 倍して足せば全始点の混合三次条件が成立する。
定理\ref{thm:maximal-tangent-ore}の語の直和の証明はこの係数環上でも働くため、
四生成族は成分ごとに自由である。特殊補正が有限窓を与えれば、
射影的な各辺のファイバー次元の上半連続性から一般窓も四辺とも有限である。
反転で残る原始四乗根も従う。この平滑化は対称性を特別ファイバーに追加していない。
\end{proof}

% !TeX root = ../papers/nc-p3.tex
\subsection{三生成基底の被約な点グラフの全被覆}
\label{sec:reduced-ternary-point-graphs}
ここでは三次元接点を持つ有限四辺で、基底の三関係が被約な平面三次曲線 $E$ の
自己同型 $\sigma$ のグラフを定める場合を扱う。曲線の成分を入れ替える移動も含める。

\begin{lemma}\label{lem:reduced-ternary-fourth-shift}
$\sigma$ はどの既約成分上でも点ごとの恒等にならない。
全 $E$ 上で $\sigma^4=\nu:=\mathbb P(r^{-1}\ell)$ が成立し、$\sigma$ の固定点は特異点に限る。
また $\sigma$ は既約成分の次数を保つ。
\end{lemma}
\begin{proof}
固定される成分があれば、その一般の滑らかな点 $p$ で、前節の二平面
$\operatorname{im} A_p,\operatorname{im} B_p$ は異なる。もし一致すれば
$\ell(p),r(p)$ が比例し、定理\ref{thm:elliptic-tangent-classification}の行列式の
微分計算により $p$ は特異となる。固定成分上の正規化に $L$ を引き戻した次数を $d>0$ とする。
$A_p(d_i(p,p))$ は次数 $3d$ の束の三つの切断で、その方向は一般点で
$\ell(p)\times r(p)$ の方向に入る（恒等的に零なら直ちに全点が持ち上がる）。
二次式の三切断の共通零因子を除くと、この方向への写像は次数高々 $2d$ の束から来る。
従って前者の三切断には次数少なくとも $d$ の共通零因子がある。
その零点では、グラフの仮定から特異点でも $\ker A_p=kp$ であり、
$d_i(p,p)$ が $p$ に比例する。第四座標に一次元の解が生じるので有限性に反する。
こうして各成分の一般点で $p,\sigma p$ は独立となり、前節の二平面の証明が
各成分に適用できる。被約性により $\sigma^4=\nu$ はスキーム上の等式である。
固定点での微分計算も前節通りである。

次数については前節の三次交叉を用いた曲線上の束の完全列を $n=3$ で使う。
この列の完全性には曲線の滑らかさを使わず、グラフでの階数2だけでよい。
行列式の等式は $L\otimes(\sigma^2)^*L=(\sigma^*L)^2$ を与える。
成分次数ベクトルに対する置換の作用を $P$ とすれば $(P-1)^2\deg L=0$ である。
有限位数の置換は標数零で半単純なので $P\deg L=\deg L$ となる。
\end{proof}

\begin{theorem}\label{thm:reduced-ternary-graphs}
被約三次曲線のグラフとなる上記の枝において、許される $E$ は滑らかな三次曲線、
節点三次曲線、横断する直線と滑らかな二次曲線、非共点三直線だけである。
いずれも位数2と4の二型があり、後の三種は最初の対応する型から平坦特殊化する。
横断二成分の場合は定理\ref{thm:ternary-torsion-smoothing}の $h=1$ が全てであり、
非共点三直線では定理\ref{thm:quantum-tangent-finite}の量子基底が全てである。
\end{theorem}
\begin{proof}
既約曲線は前二節による。横断する直線と二次曲線は次数により個別に保たれる。
二つの交点を交換するなら各射影直線の内部に固定点が生じるので、両交点は固定される。
座標を直線 $x=0$ と二次曲線 $yz=x^2$ とし、それぞれ $[0:1:t],[s:1:s^2]$ と表示する。
移動は $t\mapsto At$, $s\mapsto Bs$ である。$C=B^{-1}$ と置けば三関係は
\[
 R_x=yz-Azy+k x^2,\qquad R_y=zx-Cxz,\qquad R_z=xy-Cyx,\qquad k=(A-B^2)/B .
\]
両成分への代入で消え、全ての消える式がこの三本であることは
直線への制限と二次曲線上の $s^0,\ldots,s^4$ の係数から分かる。
実際 $A=B^2$ なら移動が平面に延びて平面グラフになり、他では二成分の
同時評価は階数6である（直線で階数3、その核から $s,s^2,s^3$ が出る）。
三次交叉は
\[
 xR_x+yR_y+(A/C)zR_z=R_xx+(A/C)R_yy+R_zz .
\]
従って $\nu$ は $[x:(C/A)y:(A/C)z]$ で、二成分での四乗条件は
$B^4=A B$、$A^4=(A B)^2$ となる。よって $B^4=1$, $A=C=B^{-1}$ である。
恒等を除くと $k\ne0$ で、対角変換で $k=1$ として前節の $h=1$ を得る。

非共点三直線を座標三角形とし、$\sigma$ の頂点置換を $\pi$ とする。
第0,1,2位置で順に座標を $1,\pi,\pi^2$ の順に揃えると、
各関係空間は異なる二文字の交換二項式三つからなる。
実際、頂点を揃えた各成分の移動は両端を保つ倍率であり、辺の内部で消えるためには
対角係数が零で、各非対角二項の比がこの倍率により決まる。
三次交叉のテンソルは、この座標では互いに異なる三文字の六順列だけに台を持つ。
両側で隣接交換が非零倍を保つので、繰り返し文字を隣接させた係数は零、
異なる三文字の六係数は全て非零となる。
元の座標へ戻すと台は $(i,\pi(j),\pi^2(k))$ である。
頂点 $i$ の先頭縮約 $\ell(e_i^*)$ は第一文字 $\pi(i)$ を欠く関係に、
末尾縮約 $r(e_j^*)$ は第一文字 $\pi^{-2}(j)$ を欠く関係に対応する。
従って $\nu$ の頂点置換は $\pi^3$ であり、$\pi^4=\pi^3$ より $\pi=1$ である。
これで三関係は一つの共通基底で量子交換式となり、量子基底の分類を適用できる。

接する直線と二次曲線は個別に保たれ、唯一の交点以外の各成分はアフィン直線である。
その上の移動は滑らかな固定点を持たず、いずれも非零の平行移動である。
直線 $y=0$、二次曲線 $yz=x^2$ とし、$[1:0:u]$ と $[s:1:s^2]$ で
$u\mapsto u+a,s\mapsto s+b$ と書く。
二次曲線上で $s\mapsto s+4b$ を実現する射影変換は一意で、直線上で $u\mapsto u+8b$ と作用する。
$\sigma^4$ が射影変換なので $a=2b$、従って $\sigma$ 自身が同様に射影変換となる。
被約三次曲線上で消える二次式はないので三関係は平面グラフを与え、仮定に反する。

最後に共点三直線の三重点を $p$ とする。この点は固定され、グラフの両側の評価階数は2である。
一次文字を $H=K\oplus kz$、$K=\ker p$、$p(z)=1$ とする。関係の二つの混合成分の
同時評価階数は、接空間二次元により $4-2=2$ である。
従って関係は二次元 $K$ の文字 $x_1,x_2$ を使い
\[
 R=\sum R_{ab}x_ax_b,\qquad m_j=zx_j-\sum_a F_{aj}x_a z-d_j,\quad d_j\in K^2
\]
と書け、$F$ は可逆である。$K^2$ はここでは二文字のテンソル積を表す。
$\ell(p)$ は $(p,p)$ が零点であることと $\nu p=p$ から両端で $p$ に消されるので、
$R$ を $\ell(p)$ に取れる。その二次行列は階数2である。実際 $w(p,-,-)$ の行列階数は、
末尾縮約 $r$ が同型であることから関係の $p$ での左評価階数である。
接平面での移動の微分は $D=F^{\mathsf t}$ である。零点近傍で混合二式から次点を解くと、
残る式の二次項は $u^{\mathsf t}R D u$ となる。三重点の重複度が3なのでこれは恒等零で、
$J=R D$ は可逆交代行列である。

$w$ の $z$ を一つ持つ部分を、左からは $zR+\sum A_{aj}x_a m_j$、
右からは $cRz+\sum B_{jb}m_j x_b$ の対応部分として比較する。
$z$ の先頭・中央・末尾の係数から $B=R$, $A=-FR$, $FRF^{\mathsf t}=cR$ を得る。
よって縮約を $R$ の係数で正規化して微分すると
$d\nu_p=cR^{-1}A^{\mathsf t}=-cR^{-1}R^{\mathsf t}D$ である。
$R=JD^{-1}$ と二次交代行列の恒等式から $c=\det D$、
$R^{-1}R^{\mathsf t}=-D^2/\det D$ なので $d\nu_p=D^3$ となる。
$d\sigma^4_p=D^4$ と比較して $D=1$ が従う。
特に三つの接方向は個別に保たれ、三成分の各々で移動は平行移動である。
三本の共点直線を個別に保ち、共点での微分が恒等の射影変換は、
共点 $[0:0:1]$ の座標で $[x:y:z]\mapsto[x:y:z+u x+v y]$ と書ける。
$\sigma^4$ のこの式の $u,v$ を4で割ると、各直線上の平行移動量の一意性から
$\sigma$ 自身を実現する。これも平面グラフになり矛盾する。
\end{proof}

\paragraph{横断二成分での特殊値と周期。}
上記 $h=1$ の九次元 $P_B$ のどの点でもAS正則である。辺の有限条件は、
二つの頂点で $d(p,p)\bmod kp$ がともに非零で、各成分の開トーラスで
$\det[p,\sigma p,d(p,\sigma p)]$ が恒等零でないことである。
実際、滑らかな部分では第四座標が二行から一意に解け、各根がその重複度のまま局所長になる。
頂点では節点基底と同じ検査になる。これらは一位置だけで検査できる。
一般の $p$ で $A_{\sigma^2p}(p)=B_p(\sigma^2p)\ne0$ であるから、
三次の混合条件は、$d_i(p,\sigma p)\in\langle p,\sigma p\rangle$ と
$d_{i+1}(\sigma p,\sigma^2p)\in\langle\sigma p,\sigma^2p\rangle$ を同値にする。
頂点でも剰余評価の非零性は移動で保たれる。
指定点の長さは20から両開トーラスの根の重複度を引いたものである。
有限部分は非空開である。両頂点で評価は接平面へ全射であり、
一般の補正では各枝の接線を避ける。各枝で移動の微分は1と異なるため、
その行列式には頂点の各逆像で単純な零があり、どの成分でも恒等零ではない。

同型は正規表示の射影補正を、二次曲線上の $s\mapsto l s$ および
$s\mapsto l/s$（移動の向きを対応させるもの）によって比較すれば必要十分である。
二次曲線と二つの交点を保つ射影変換はこれで尽きる。
元の座標で倍率の作用は対角 $(l,1,l^2)$ であり、補正への作用は
$d\mapsto GdG^{-1}$（値には二次の作用）で計算できる。
指定点以外の全局所環は曲線状なので、ここでも指定を忘れて同じ判定になる。
特定の係数での最小周期は $F_B^j$ に対してこの判定を行う。
倍率のトーラスは節点評価の二方向に異なる重みで作用するので、
射影空間の一般軌道は一次元、同型類の族は既約七次元である。

一般最小周期はやはり $m=2,4$ である。実際、上記の倍率の重み $-1$ の部分には、
$m=2$ で固有値 $1,-1$、$\lambda=\sqrt{-1}$ で $-1,-\sqrt{-1}$ の
$F_B$ 固有直線が各一本ある。重み3の部分には順に $1,-1$ および
$-1,\sqrt{-1}$ がある。これは前節の固有核計算を文字の重み $(1,0,2)$ で
分けて得られ、両重みに補生成元変更はない。
一つの重みに属する二係数の比は倍率で変わらず、移動の固有値比は位数 $m$ を持つ。
$m=2$ で自己同型に残る反転は重み $-1$ と3を交換し、
それらの係数比は一般には独立なので反転でも真の移動は戻せない。
以上の四種類は二つの楕円族とその境界の記述であり、非被約三次曲線をこの被覆には含めていない。

% !TeX root = ../papers/nc-p3.tex
\subsection{三次元接点を持つ有限四辺の完全分類}
\label{sec:maximal-tangent-classification}
三生成部分に点対応がグラフであるという仮定を置かず、残る非被約曲線も判定する。

\begin{lemma}\label{lem:ternary-point-graph-growth}
三生成三関係の標準三次元quadratic Koszul左右AS正則代数 $B=T(H)/(S)$ の
点対応は、平面全体または平面三次曲線上の自己同型のグラフである。三次曲線は非被約でもよい。
\end{lemma}
\begin{proof}
三次交叉を張る $w$ の両端縮約は $H^*\simeq S$ である。
$p\in H^*\setminus0$ での関係の左右評価の階数は等しい。
実際、これらは二次双対の $1$ 次部分から $2$ 次部分への $p$ の左右乗法の階数である。
Frobenius対合で左乗法に随伴なのはNakayama自己同型で移した $p$ の右乗法であり、
この右乗法は元の $p$ の右乗法と一定の基底変更で対応する。
従って左右評価の行列式も同一の $p$ の多項式として定数倍である。

もしこの階数が1なら、右隣の評価像を $ku$、左隣の評価像を $kv$ とし、
三つの位置で順に $H/ku,H/kv,k_p$ を循環する生成空間の商を取る。
$k_p$ へは $p$ で写す。その両側に接する関係は全て零になる。
残る隣接対への三関係の像の次元は高々2である。
実際、$w(p,-,-)$ は $S$ の非零元で、その先頭支持は
残りの末端を縮約して得る $(p\otimes1)S=ku$ だからこの商で消える。
従って三文字の一組で残る商の次元は $2\cdot2-2=2$ 以上となる。
組をつなぐ境界に関係がないので長さ $3n$ で $2^n$ 以上となり、
$H_B(t)=(1-t)^{-3}$ に反する。階数0も両端支持から不可能である。
左右とも各点で階数2以上なので、非零行列式の場合、対応はその三次式の
零点スキーム上で単元二次小行列から復元される唯一の核直線のグラフである。
逆射も右評価から得られる。行列式が恒等零の場合は両射影が平面全体へ同型となる。
\end{proof}

\begin{theorem}\label{thm:maximal-tangent-complete}
標数零の代数閉体上の標準14641型のうち、四辺が有限で、ある点の接空間が三次元のものは
定理\ref{thm:reduced-ternary-graphs}の四種類・各二型で尽きる。
三次元接点はどの対象でも一意である。埋込み次元3という条件は窓の添字保持同型不変量で、
この枝とその補集合の間に同型はない。
\end{theorem}
\begin{proof}
三生成基底は定理\ref{thm:maximal-tangent-ore}で復元され、前補題と
定理\ref{thm:nodal-tangent-windows}により非被約三次曲線の場合だけが残る。
その形は三重直線または二重直線と単直線である。

まず重複直線の台上で $\sigma$ が恒等でないとする。
単直線があればその上でも点ごとの恒等でない。
補題\ref{lem:reduced-ternary-fourth-shift}の固定成分の証明は、その成分の
一般点で曲線が被約でありさえすれば働くからである。
この場合、$\sigma^4=\nu$ は非被約構造上でも成立する。
各成分の一般点の局所環はArtin環であり、$p,\sigma p$ の独立性と
$\det[p,\sigma p,d_i(p,\sigma p)]$ の非零性はその剰余体で成立する。
従って前節までの二平面の証明に使う小行列式は局所環の単元であり、
二つの像は同じ階数2の直和因子となる。両零化直線の一致が
$\sigma^4=\nu$ をこの環上で与える。平面の超曲面には埋込み随伴点がないので、
全一般点の局所環上の等式は全スキーム上の等式である。

二直線なら交点、三重直線なら台の自己同型の固定点を $p$ とする。
$p$ は曲線の重複度3の点で接空間は平面全体である。
定理\ref{thm:reduced-ternary-graphs}の $RD=J$ による微分計算は、
局所方程式の二次項が零であることだけを使うので、この場合も
$d\sigma_p=1$ となる。従って射影変換 $\nu$ は $p$ で恒等微分を持ち、
行列を $I+N$, $N^2=0$ と書ける。$I+tN$ は全 $t$ で曲線を保ち、
制限は $\sigma$ と可換である。実際、整数 $t\ge0$ で既に成立し、
曲線の保存と二つの射の一致は閉条件である。
$\rho=I+N/4$ の曲線への制限を取れば $\tau=\sigma\rho^{-1}$ は位数が4を割り、
$p$ を固定し接空間で恒等である。局所環の各次数商で恒等なので、
各有限次数の近傍で単冪かつ有限位数、従って恒等である。
局所環の分離性から $p$ の局所環で恒等であり、全成分が $p$ を通るので
埋込み随伴点の不在から全曲線上で恒等となる。
これで $\sigma$ は平面に延びるが、非被約な三次超曲面にも消える二次式はなく、
このグラフ上で消える三つの双一次式は平面グラフを定めて矛盾する。

残りは重複直線 $z=0$ の台上で $\sigma=1$ の場合である。
文字を $x,y,z$ とすると、$z=0$ の直線の積への三関係の制限の像は
$k(xy-yx)$ である。像が零なら直線の積全体を含みグラフにならない。
核の二関係の $xz,yz$ 係数は独立である。従属ならある左点 $[x:y:0]$ で
二関係の評価がともに零となり階数2に反する。右側も同様である。
従って $z$ は正規元で、$hz=z\theta(h)$、$\theta(z)=z$ と書ける。
$B/(z)=k[x,y]$ とHilbert級数から $z$ の両側乗法は単射である。
従って $\theta$ は生成元上の可逆一次写像から定まる代数自己同型である。
重複直線の一般点で $z$ は余接の非零な横断方向である。
関係 $h(p)z(\sigma p)=z(p)(\theta h)(\sigma p)$ を $z(p)^2$ を法として読むと、
$z(\sigma p)/z(p)$ は剰余体で非零で、$h=x,y$ に共通である。
ゆえに $\theta$ の $H/kz$ 上の作用は射影直線で恒等、すなわちスカラー $\gamma$ である。
全関係は定数を選んで
\[
 xz=z(\gamma x+a z),\qquad yz=z(\gamma y+b z),\qquad
 xy-yx=z(lx+my+rz)
\]
となる。

$\gamma\ne1$ なら $x,y$ に $z$ の倍を足して $a=b=0$ とする。
第三式の $\theta$ 不変性から $l=m=0$、かつ $(\gamma^2-1)r=0$ を得る。
$r=0$ なら平面グラフなので、倍率変更によって $(\gamma,r)=(-1,1)$ だけが残る。
$\gamma=1$, $a=b=0$ なら $z$ は中心である。$(l,m)\ne(0,0)$ なら
$z\ne0$ 上で対応は対角で、$lx+my+rz=0$ は点ごとに固定される単直線である。
実際、二つの中心関係で次点を消去するとこの一次式がそのまま残る。
これは除外済みであり、残る三重直線型は $l=m=0,r=1$ に取れる
（$r=0$ は平面グラフ）。

最後に $\gamma=1$, $(a,b)\ne(0,0)$ なら $(a,b)=(1,0)$ に直せる。
第三式の $\theta$ 不変性は $l=0$ を与える。
$m=1$ なら $r=0$ は射影移動 $[x:y:z]\mapsto[x-z:y:z]$ の平面グラフであり、
他では $r=1$ にできる。$m\ne1$ なら $y$ への $z$ の倍の加法で $r=0$ とする。
$z\ne0$ での関係からこのとき $y=0$ に単直線があり、移動は $x\mapsto x-z$ である。
$R=yz-zy$, $M_y=xy-yx-mzy$, $M_z=xz-zx-z^2$ に対して三次交叉は
\[
 (x-(m+1)z)R-yM_z+zM_y=R(x+z)+M_yz-M_zy .
\]
従って $\nu[x:y:z]=[x-(m+2)z:y:z]$ となる。この単直線の一般点では
二平面から四移動の等式が得られるので、$m+2=4$、すなわち $m=2$ に限る。

残った四つの基底では行列の係数に自由パラメータはない。
全位置に延びる初期補正の線型空間 $E$ を定理\ref{thm:ternary-ore-linear-space}で計算する。
$P=[u:v:0]$ とし、$g=d(z)(P,P)$、$h=u d(y)(P,P)-v d(x)(P,P)$ と置く。
全ての補正を含む行列を使うと結果は次の通りである。
\[
\begin{array}{c|c|c|c|c|c}
\gamma&a&m&r&\dim E&E\text{ の全域で零}\\ \hline
-1&0&0&1&9&h\\
1&0&0&1&12&g\\
1&1&1&1&9&g\\
1&1&2&0&12&g
\end{array}
\]
再現計算は \path{computations/verify_thick_ternary_exclusion.py} で、
三十行の混合条件から両方向の射影の交叉を反復し、$E$ の全基底を上記の二元斉次式に代入する。
$g$ は二次、$h$ は三次なので一方が恒等零なら必ず射影直線上に共通根がある。
そこで $d(P,P)$ は $P$ に比例し、第四座標は一次元の解を持つ。
これで非被約三次曲線を全て排除した。前節までの各有限表示で指定点が一意であることと、
局所環の接空間が同型で保たれることから残りの主張を得る。
\end{proof}

% !TeX root = ../papers/nc-p3.tex
\subsection{有限四辺の全域での添字保持同型と周期}
\label{sec:planar-point-flags}
三次元接点を持たない場合も、局所点環どうしの同型写像をあらかじめ求めずに、
第一・第二窓の添字保持同型を有限個の数値記録で判定できる。

\begin{lemma}\label{lem:planar-point-wronskian}
長さ高々20、埋込み次元2の局所Artin完全交叉 $C$ に対し、
余接平面の直線からなる非空な有限集合が同型不変に定まり、その個数は高々6である。
\end{lemma}
\begin{proof}
$V=\mathfrak m/\mathfrak m^2$ とし、$\Sym V\to\operatorname{gr}_{\mathfrak m} C$ の
核が初めて非零となる次数を $d$、その次数の核を $J_d$ とする。
極小表示 $C=k[[x,y]]/(f,g)$ を取ると、$f,g$ の位数は $d$ 以上で、
次数 $d$ の初期式で $J_d$ が張られる。特に $1\le\dim J_d\le2$ である。
次数 $2d-2$ までの打切り環の次元は $d(2d-1)$ であり、
その中で両関係の倍数の次元は高々 $2\binom d2=d(d-1)$ である。
従って $\dim C\ge d^2$、すなわち $2\le d\le4$ である。

$J_d$ が一次元ならその生成元 $f_d$ を、二次元なら基底 $f_d,g_d$ の
二元Jacobian $W=(f_d)_x(g_d)_y-(f_d)_y(g_d)_x$ を取る。
後者は独立な同次数二元形式の比が標数零で非定数であるため零でない。
例えばアフィン座標で微分して斉次Euler等式を使えばよい。
基底変更と変数変更に対してこの式の射影類は自然に移る。
次数は $d$ または $2d-2$ なので、$\mathbb P(V^*)$ 上の相異なる根の数は1から6である。
各根の零化直線を $V$ 内に取れば主張を得る。
\end{proof}

\begin{theorem}\label{thm:finite-all-window-keys}
四辺が有限な任意の二つの標準AS対象について、添字保持同型の存在と最小周期は、
連続する三つの四次元座標空間からの有限な記録の比較で決定できる。
埋込み次元2以下の部分では、点環の同型型、三空間の独立な単元倍、共通の
局所パラメータ変更を同じ計算に含める。使うのは高々四次の二元形式または
その次数6以下のJacobianの根、因子の有限置換、線型代数および整数核の単項式である。
三次元接点の部分は定理\ref{thm:maximal-tangent-complete}の表示で決定される。
\end{theorem}
\begin{proof}
点環を $D=\prod C_a$ とし、局所因子の長さ・埋込み次元 $e_a$・
$\mathfrak m_a^{n_a}=0$ となる最小の $n_a$ を記録する。
データが等しい因子に標識を付け、$e_a=2$ の各因子で前補題の直線を一つ選ぶ。
適合する余接基底を $x,y$（指定直線 $ky$）とし、これを極大イデアルに持ち上げる。
$e_a=1$ では一変数を取り、$e_a=0$ では変数は不要である。
単項式を評価して得る全射を
\[
 T=\prod_a k[t_1,\ldots,t_{e_a}]/(t_1,\ldots,t_{e_a})^{n_a}
            \longrightarrow D
\]
とし、核を $I$ とする。$U_s\subset D$ を定理\ref{thm:finite-square-chart}の
座標空間とし、その全逆像 $P_s\subset T$ を取る。比較に使うのは
四つの部分空間 $(I,P_0,P_1,P_2)$ である。
この形なら三つの自明化に共通して同じ点環を用いるという条件が失われない。

各因子の単項式を全次数昇順、同じ次数では $y$ の指数昇順に並べる。
標識と直線を保つ任意の点環同型は、$T$ の因子ごとの変数置換に持ち上がる。
実際、選んだ生成元の像の代表多項式を取れば、一次係数は可逆で
$x\mapsto a x+b y$, $y\mapsto c y$ の形を持ち、打切り環の自己同型となる。
逆に $I$ を $I'$ に運ぶ置換は点環同型を与える。任意の局所単元も単元へ持ち上がるので、
三空間の同値は $T$ のこの共通変数変更と $P_s$ ごとの独立な単元倍に等しい。
定数因子のスカラーと、各 $a,c$ をトーラスに分けると、残る群は下単三角である。
そのLie環は次の操作で張られる。各位置別の単元の冪零部分による乗法、
全四空間に共通の次数2以上の単項式を係数とする偏微分、
および二変数因子の $y\partial_x$ である。
いずれも上の行順で次数を上げるか、同次数で $y$ の指数を上げる。
定数項1の単元と指定した線型部の変数変更はその係数で連結にパラメータ化される。

以下は定理\ref{thm:curvilinear-window-orbits}の三角消去と同じであるが、
全逆像の列数は4には限らない。各空間の独立な列基底を取り、
最初の独立な主元行を単位行列に正規化する。非主元行を昇順、列を降順に読み、
既読値の安定化Lie環の平行移動速度の核を順次取る。
速度非零の座標を指数作用で零にすると、単冪軌道の代表は一意である。
この証明は任意の列数で成立し、核 $I$ の次元が零の場合はその行列を空にしてよい。
行 $x^iy^j$ のトーラス重みは $a^ic^j$、単元倍を許す側では
さらにその位置・因子の定数倍を掛ける。主元行で割った重みの整数核を使って
非零値の不変単項式を取る。主元・台・単項式値を全て記録すれば
標識と指定直線を保つ同値と記録の一致が必要十分になる。

従って根と同じ型の因子の全ての標識付けにわたる記録集合が完全不変量となる。
局所完全交叉であることは有限な辺の完全交叉性から従う。
埋込み次元3を含む場合はその存在自体が不変で、前節までの分類を使う。
最小周期は実際の射影対応に沿って得た開始位置 $0,1,2$ の記録を比較して
系\ref{cor:curvilinear-window-period}と同じく決まる。
任意の三空間から作った二窓がAS延長を持つかどうかは、この同型判定とは別の条件である。
\end{proof}

\path{computations/verify_finite_planar_window_orbits.py} では
局所点環 $k[x,y]/(x^2,y^2)$ と16個の単純点の直積を使い、二元Jacobian、
非線型な共通変数変更・三位置の単元倍・核を足した持上げの変更が記録を保つことを検査する。
一位置だけで接係数を変えると標識付き記録が異なることも検査する。
この三空間は軌道計算の試験入力である。

% !TeX root = ../papers/nc-p3.tex
\subsection{正次元の点対応でも検出できるOre窓と楕円基底の分類}
\label{sec:elliptic-ore-recognition}
\begin{proposition}\label{prop:ore-mark-recognition}
標準有限Frobenius四窓の中で、定理\ref{thm:maximal-tangent-ore}の三生成Ore表示を持つ
標準AS対象とその指定点を、辺の有限性を仮定せずに検出できる。
標準四窓上で四点 $[p_i]\in\mathbb P(V_i^*)$ を取り、
各 $(p_i,p_{i+1})$ が辺の零点で、関係の左右評価が個別に階数3、
二評価を合わせた射影Jacobi行列の階数も3という条件を課す（添字は法4）。
この条件の下で四生成候補の標準AS正則性は、復元した三生成基底の標準二次AS正則性に同値である。
点を指定した同型と移動は $P_B$ と $F_B$ で決まる。
\end{proposition}
\begin{proof}
点を反復すれば全整数上で同じ評価条件を得る。定理\ref{thm:maximal-tangent-ore}で
$S_i$ と可逆な $f_i$ を得る部分はこのJacobi行列だけを使う。
純粋な三次交叉の次元が1となる証明で用いたのも、四生成三次元数20と
直前の点での評価階数3だけである。その交叉の両端支持も連続する指定点での階数3である。
従って高次の直和表示・正則性の両方向・指定同型の証明が全て適用される。
逆に標準ASなOre表示では $H_i$ の零化点の両評価とJacobi行列が上記階数を持つ。
この列を標準四窓へ運ぶと、同定前の列に対する純粋交叉の末端支持の証明から
$[p_{i+2}]=[p_{i-2}]$ が従うので、実際に四点の検査で拾える。
有限性から得た「高々二本」という個数をここでは使っていない。
\end{proof}

\begin{theorem}\label{thm:all-elliptic-ore-windows}
上の検出条件を満たし、三生成部分の三関係が滑らかな平面三次曲線 $E$ の
自己同型 $\sigma$ のグラフを定める標準AS対象を考える。正次元の辺も許す。
指定点は対象から一意に復元できる。$L=\mathcal O_E(1)$ として基底は
\[
        L\ne\sigma^* L,\qquad L\otimes(\sigma^2)^*L\simeq(\sigma^*L)^2
\]
を満たす $(E,L,\sigma)$ で尽くされる。各基底に対する補正は次の通りである。
\begin{enumerate}
\item $\sigma$ が位数2,4の平行移動なら定理\ref{thm:elliptic-tangent-classification}の
九次元空間である。他の平行移動の補正は零だけである。
\item 固定点があるとき、原点保持部分 $\alpha$ の位数は2,3,4のいずれかである。
位数2では補正空間は二次元、位数3,4では零である。
\end{enumerate}
零補正は常に $\check{B[Z]}$ である。有限四辺の部分以外は全て最小周期1である。
添字保持同型類を得るには基底を偏極曲線と移動の同型で比較し、
非零補正はその射影類を有限群 $\Aut(E,L,\sigma)$ の誘導作用で比較すればよい。
\end{theorem}
\begin{proof}[基底と指定点]
二つの束条件の必要性は定理\ref{thm:elliptic-tangent-classification}と
補題\ref{lem:reduced-ternary-fourth-shift}の行列式計算による。
逆にこのデータから $S=\ker(H^0(L)\otimes H^0(L)\to H^0(L\otimes\sigma^*L))$ を取る。
第二因子の切断は $\sigma$ で引き戻す。前節のpencilの完全列から $\dim S=3$ である。
平面の点 $p$ の課す二切断の鉛筆は、$p\notin E$ なら $E$ 上で基点自由であり、
$\sigma^*L\otimes L^{-1}$ に切断がないことから $S$ の $p$ での評価核が零になる。
$p\in E$ なら $L(-p)$ の基点自由な二切断を使い、
$\sigma^*L\otimes L^{-1}(p)$ の切断が一次元なので評価階数は2である。
従って行列式は $E$ の三次方程式であり、対応はスキームとして所定のグラフとなる。

定理\ref{thm:elliptic-tangent-classification}の $n=3$ の束の列で、
中央写像の核直線束は行列式から
$L\otimes(\sigma^2)^*L\otimes(\sigma^*L)^{-2}\simeq\mathcal O_E$ となる。
その生成切断は $S\otimes H$ に入り、像が零なので $H\otimes S$ にも入る。
従って一次元の三次交叉を生成する。片側縮約が全三次元を張ることも確認する。
$p$ に消える核関係の直線 $E\to\mathbb P(S)$ は単射である。
二点で同じ関係の左評価が零なら、その行列は階数1以下である。
しかし可分解な非零関係が既約曲線のグラフ上で消えることはない。
従ってこの写像の像は直線に入らない（楕円曲線から射影直線への単射はない）。
逆側も同じである。これで先の束の完全列からbarの初項、三次の非零因子、
Hilbert級数とKoszul--Frobenius性を得た証明がそのまま適用され、$B$ は標準三次元ASである。

指定点以外の辺零点では両端で三生成部分への射影が定まる。
片方だけが零となることは混合関係が禁じる。
局所的に基底の両射影の自明化を選ぶと $(p,\sigma p)$ の上で増えるのは第四座標二つであり、
混合関係の一行で一つを消せる。従って局所的に $E\times\mathbb A^1$ の部分スキームに入り、
接空間は高々二次元である。三次元の指定点は唯一で、正次元の辺でも指定同型条件が使える。
\end{proof}

\begin{proof}[移動と有理切断による補正]
$\sigma$ が固定点を持たない場合は平行移動である。$\sigma^4\ne1$ のとき
Hesse表示で三次交叉は巡回対称なので $\nu=1$ である。
前節までの二平面の議論を逆に読めば、任意の可延長補正について
$\det[p,\sigma p,d_i(p,\sigma p)]$ は一般点で零である。
よって全点で零となり、定理\ref{thm:elliptic-tangent-classification}の束の完全列から
初期補正は左右乗法で尽きる。位数2,4の結果は同定理による。

固定点があるときも、$\nu=\sigma^4$ が全曲線上で成立すれば固定点で
行列式の微分が消えて滑らかさに反する。従ってやはり上の行列式は恒等零である。
$q=\sigma p$ とし、一般点での二ベクトルの独立性により
\[
    d_i(h)(p,q)=a_i(p)h(q)-h(p)b_{i+1}(q)
\]
と一意に書ける。$a_i,b_{i+1}$ は $L$ の有理切断である。
固定点以外では二行を解けるので正則である。
$\sigma$ の不変微分に対する倍率を $\eta\ne1$ とすれば、各固定点の微分も同じ $\eta$ である。
その近傍では $p$ と $q$ の二次小行列のうち一つは単純零を持つ。
実際、埋込みの非零接ベクトルに対し二つの微分の差は $1-\eta$ 倍である。
従って極は高々単純である。
三次の混合条件を $(p,q,\sigma q)$ で評価すると内側の係数の差
$a_{i+1}(q)-b_{i+1}(q)$ に各 $s(p,\sigma q)$ が掛かる。
一般点ではこれらの全てが零にはならないので $a_{i+1}=b_{i+1}$ である。

固定点の全てで $a_i$ の主要部を取り、その列を $\xi_i$ とする。
局所座標で $h(p)$ と $h(q)$ は同じ定数項を持つため、
$d_i(h)$ の正則性は $\xi_{i+1}=\eta\xi_i$ を与える。
固定点因子を $T$ とすれば $H^1(L)=0$ により
$\xi_0$ は $H^0(L(T))$ の一元 $a$ の主要部となる。
$a_i-\eta^i a$ は正則で、これらの差は補生成元を変更して消せる。
従って全ての類は $d_i=\eta^i d$ の形に表せる。正確には
\[
        P_B=\ker(\lambda_B+\eta\rho_B)/
                 \{L_v-\eta R_v:v\in H\},\qquad F_B=\eta\,1 .
\]
分子の任意の元は実際に全整数列に延びる。
分母の次元は3であり、左右乗法の部分空間の交叉は零である。
実際、$vh=hu$ を全 $h$ についてグラフで評価すると、一般に $p,q$ は独立なので
$v(p)=u(q)=0$、従って $u=v=0$ となる。定理\ref{thm:ternary-ore-linear-space}の
両核の計算で、純粋に補生成元変更から来るこの固有値の元は上記の分母で尽きる。
$X_i$ を $\eta^{-i}$ 倍すれば混合関係は一定になり、最小周期1を得る。
また行列式が恒等零なので一般の $p$ が持ち上がり、辺は正次元である。
\end{proof}

\begin{proof}[固定点を持つ三つの場合の全計算]
変曲原点で $\sigma=t_P\alpha$ と書く。
$L=\mathcal O(3O)$ なので二つの束条件は $3P\ne0$ と $(1-\alpha)3P=0$ である。
原点保持自己同型の位数は $2,3,4,6$ であり、位数6では $1-\alpha$ の核が自明なので除かれる。
$b$ を三等分点として原点を移す射影変換で $P$ を $(1-\alpha)b$ だけ変えられる。
$\alpha$ の位数2,4の場合、$1-\alpha$ は三等分点上可逆であり、
その核上では3倍が可逆である。従って $P$ を $\alpha$ が固定する非零二等分点に取れる。

$\alpha=-1$ ではHesse座標の前二文字の交換 $G$ が逆元を与え、
$P=[1:1:c]$、$c\ne0$, $c^3\notin\{1,-8\}$ を取れる。
従って関係は前の対称三関係に $1\otimes G$ を施したものである。
前節の三次商の零化行列を使って $\lambda_B-\rho_B$ を計算すると、
$\mathbb Q(c)$ 上の階数は13で、分母のない行の十三次小行列式が
\[
                       2 c^{60}(c-1)^{28}(c^2+c+1)^{28}
\]
となる。全ての許容値で非零なので核次元は5、商は二次元である。

位数4なら $y^2=x^3-x$ で $\alpha(x,y)=(-x,\sqrt{-1}\,y)$ またはその逆とする。
唯一の非零固定二等分点 $(0,0)$ の平行移動を合成した像の斉次座標は
$[x^2:x:\sqrt{-1}\,y]$ となる。元の $[1:x:y]$ との九つの積から関係を取り、
$\lambda_B+\sqrt{-1}\rho_B$ の階数は15で、共役も同じである。

位数3では $y^2=x^3+1$、$\alpha(x,y)=(\omega x,y)$、
$\omega^2+\omega+1=0$ と取る。$3P=(0,1)$ または $(0,-1)$ の全18点を考える。
良い素数109ではこれら18点は相異なる有理点へ還元され、$\omega$ は45または63となる。
検算はこの曲線の全108点でグラフを評価して三関係を復元し、
36通り全てで $\lambda_B+\omega\rho_B$ の階数15を得る。
標数零への持上げを明示する。三等分写像は109上エタールであるため
18点は全てこの計算で尽くされる。評価に用いた曲線の各点も良い還元の滑らかさにより持ち上がる。
九列の評価で階数6の単元小行列を取れるので、標数零の三関係の自由基底が持ち上がる。
三次商の余次元は両標数で10である（剰余では直接行列階数を確認）。
従って二次・三次商の基底を局所的に選んで同じ非零小行列式を持ち上げられる。
階数上界15は三次元の補生成元変更から従うので、標数零で商は零である。
これら全ての厳密計算は \path{computations/verify_elliptic_fixed_ore.py} にある。
非零の射影補正への作用は $d(h)\mapsto Gd(G^{-1}h)$ であり、
基底の関係はグラフから一意なので、所定の有限群の比較が零の場合も含め必要十分である。
\end{proof}

% !TeX root = ../papers/nc-p3.tex
\subsection{平面グラフ基底の固定直線と指定点を忘れた軌道}
\label{sec:plane-graph-ore-windows}
命題\ref{prop:ore-mark-recognition}の三生成基底の点対応が射影平面全体の自己同型
$[G]$ のグラフである場合を考える。三生成関係は $p\wedge G^{-1}r$ の三成分で、
$B$ は多項式環の次数付きtwistである。実際、長さ$n$のテンソルを
$(p,Gp,\ldots,G^{n-1}p)$ で評価すると、各位置を順次 $G$ の冪で
基底変更した後の関係は可換関係となり、商は全$n$次式である。

\begin{theorem}\label{thm:plane-ore-fixed-lines}
平面グラフについて、指定列を持つOre補正の空間 $P_B$ と移動 $F_B$ は次で尽くされる。
\begin{enumerate}
\item $[G]=1$ なら二次導分をEuler導分の一次式倍で割った十五次元空間で、移動は恒等である。
\item $[G]\ne1$ が直線を点ごとに固定しなければ $P_B=0$ である。
\item 残りでは $G=1+e\ell$ と取る。ここで $e$ は点ベクトル、$\ell$ は一次形式で、
$t=1+\ell(e)\ne0$ とする。$P_B$ は固定直線 $\ell=0$ 上の二次式の三次元空間、
移動は $t$ 倍である。全ての例の最小周期は1である。
\end{enumerate}
\end{theorem}
\begin{proof}
恒等の場合、三次の条件は $(d_{i+1}-d_i)(h)k=(d_{i+1}-d_i)(k)h$ である。
従って差が一次式倍のEuler導分に限り、定理\ref{thm:ternary-ore-linear-space}から結論を得る。
非恒等の場合、定理\ref{thm:nodal-tangent-windows}の平面における二つの縮約平面の計算により
$d_i$ の評価は一般点で $p,Gp$ の一次結合となる。この二項を
$a_i(p)Gp-b_{i+1}(Gp)p$ と書く。斉次性から係数は $\mathcal O(1)$ の有理切断である。
混合三次関係を $(p,Gp,G^2p)$ で評価すると
$(a_{i+1}-b_{i+1})(Gp)$ と三関係の $(p,G^2p)$ での評価の積が零である。
後者が全て消えるのは $Gp=p$ が射影的に成立するときなので、$a_i=b_i$ である。

二次小行列で二係数を解くと、極因子は固定直線に限る。実際、固定点集合は
固有空間の射影化の和で、平面上の有理切断は余次元一の全点で正則なら正則である。
従って固定直線がなければ補生成元を変えて全補正を消せる。
固定直線があればその上での固有値を1に選ぶと $G=1+e\ell$ となり、
$p\wedge Gp=\ell(p)p\wedge e$ のため一般の極は高々単純である。
よって一意な主要部は $a_i=f_i/\ell$ の二次分子の直線への制限で与えられる。
$\ell(Gp)=t\ell(p)$ より二次補正の正則性は
$f_{i+1}|_{\ell=0}=t f_i|_{\ell=0}$ と同値である。逆に任意の二次式 $f$ に対し
$a_i=t^i f/\ell$ の表示は全位置で正則な二次補正を与える。
有理式での相殺と三次評価の単射性により全混合三次関係を満たす。
正則な一次切断の加算はちょうど補生成元変更である。初項が左右乗法の和なら
一般点での係数の一意性から主要部は零であり、主張の同定を得る。
\end{proof}

\begin{theorem}\label{thm:plane-ore-unmarked}
非恒等の平面グラフ基底を少なくとも一つ持つOre窓どうしの添字保持同型を考える。
固定直線のない型は $[G]$ の射影共役類で分類される。
固定直線のある型は次の表で分類され、表の各行および固定直線のない型は互いに素である。
$r$ は二元二次式 $f|_{\ell=0}$ の階数である。
\[
\begin{array}{c|l}
 t\ne1 & \{t,t^{-1}\},\quad r=0,1,2\\
 t=1 & r=0;\quad r=1,2\ \text{ではさらに}\ f(e)=0\ \text{か否か}
\end{array}
\]
全係数で正則である。各 $t=1$ の類は $t\ne1$ の類の平坦特殊化である。
\end{theorem}
\begin{proof}
固定直線のある型に対し、位置 $i$ の旧水平座標を $G^ip$ で置換する。
旧第四座標を $u+t^i f(G^ip)/\ell(G^ip)-f(p)/\ell(p)$ とする。
後二項の差は一次式である。実際その分子は直線上で消え、$G^i$ は直線上恒等である。
新しい座標 $(p,u),(r,v)$ では六関係は
\[
 p\wedge r=0,\qquad u G r-vp=e f(p,r)
\]
となる。$f(p,r)$ は $f$ の対称偏極である。これは純粋な三関係を移し、
$r=p$ で旧有理補正を代入して $G^{-i}$ を掛ければ確認できる。
対角で消える双一次式は純粋な可換関係に吸収される。
特に一定表示である。$C=k[p_1,p_2,p_3]$ に対する通常のOre拡大としても、
$\gamma(h)(p)=h(G^{-1}p)$、$\delta(h)=h(G^{-1}e)f$ である。
$h-\gamma(h)$ が $h(e)\ell/t$ に等しいため、$\delta$ の捩れたLeibniz則は可換関係を保つ。
従って順序単項式基底があり、また定理\ref{thm:maximal-tangent-ore}で全次数の正則性を得る。

$Q$ を二次曲面 $u\ell(p)-f(p)=0$、$D$ を直線 $p\wedge e=0$ とする。
辺の第一射影はイデアル $\mathcal I_Q\mathcal I_D$ のスキームへ同型である。
$p\ne0$ の開集合では $r=p$ と取れ、混合式の一行で $v$ を消すと
$(u\ell(p)-f(p))p\wedge e=0$ が残る。$p=0$ の近傍では $u=v=1$ と取れ、
$(G-e f(p,-))r=p$ である。左の係数はそこで可逆で、階数一更新の余因子を掛けると
$p\wedge r$ は単元倍を除いて同じ積となる。辺の対応は $Q$ 上恒等である。
$f$ を $\ell$ と独立な二変数で取れば $\operatorname{rk}Q=2+r$ であり、
その被約な曲面部分から $Q$、イデアル商から $D$ を復元できる。
各位置の射影同型は $Q$ 上で隣の位置と等しい。$Q$ を点ごとに固定する射影変換は恒等である。
実際、固定集合は固有空間の和で、平面を二枚含む場合も二固有値は交線上で一致する。
従って一つの射影同型で全位置を比較できる。

$t\ne1$ のとき $e$ を第三方向に選び、$\ell=(t-1)z$、$f$ を前二変数の式に取れる。
$Q\cap D$ は異なる二点で、$D$ 上の対応は $u/z\mapsto t u/z$ である。
従って倍率は反転を除いて不変である。前二変数の基底変更と倍率で同階数の $f$ を移せる。
$u$ と $z$ を交換すると最後の関係 $tuz'-zv=f(p,r)$ は倍率変更で逆倍率の表示となり、
非零な $f$ の定数倍も吸収できる。$t=1$ なら $e\in\ker\ell$ で、
$Q$ は $D$ を含むか、$D$ と一点で二重に交わる。この二つを区別する条件は $f(e)=0$ である。
逆に固定直線上で $[e]$ を保つ二変数の射影変換は全て $[G]$ と可換な変換に延びる。
実際 $e$ の倍率に等しく $\ell$ の商方向を伸縮すればよい。
この群は射影直線の $[e]$ の補集合上のアフィン群であるから、
二元二次式の根の重複度と $[e]$ が根に含まれるかで軌道が定まる。

固定直線がないときは $f=0$、同じ位置変更で $p\wedge r=0,\ u G r=vp$ となる。
辺の台の二次元成分は平面 $u=0$ 一枚だけで、それ以外は固有点と指定点を結ぶ有限本の直線である。
これで二次曲面を持つ型から区別できる。この平面の外では $u=v=1$ で $r=G^{-1}p$、
接空間が三次元の点は $p=0$ だけである。実際 $G^{-1}p=\mu p$ となる非零点では
微分は $p\wedge(G^{-1}-\mu)dp$ で、その恒等消滅は $G^{-1}-\mu$ の階数が高々1となることを強いる。
従って同型は指定点を保ち、指定同型の基底共役条件が必要十分である。

最後に $e=(1,0,0)$、$\ell=z+s x$、$f=f(x,y)$ と取ると $t=1+s$ である。
$s=0$ で所定の任意の移行型を得る。$1+s$ を逆にした係数環上でも上記Ore拡大は
順序単項式基底を持つので、この特殊化は平坦である。固定した $G$ で $f$ の任意の退化も同様である。
記号的な余因子・三次重なりの検算と、Jordan各型の有理標本での独立した
補正空間検算は \path{computations/verify_plane_graph_ore.py} にある。
\end{proof}

% !TeX root = ../papers/nc-p3.tex
\section{有限な第一窓の剛性と三次障害}
\label{sec:finite-window-rigidity}
前節の二十点例で、正規化した第三空間の候補を退化解まで含めて決定する。
これにより第一窓を動かしたときの非延長開集合も得られる。以下では第三階数を指定しない。

\begin{theorem}[二十点例の全平方根]\label{thm:finite-all-square-roots}
前節の被約例の第一座標空間を $U\subset D=k^{20}$、$Q=U^2$とする。
$\Gr(4,D)$で$L^2\subset Q$を課した閉スキームは、ちょうど$2^{19}$個の被約点からなる。
その全点は$L=\epsilon U$で、$\epsilon$は各座標で独立に取った符号、全体符号は同一視する。
従って、この第一窓から四辺の点対応が有限な標準AS対象への延長は添字保持同型を除き一つである。
\end{theorem}
\begin{proof}
二十点を行とする第一座標行列を$X$とする。
四座標関数を$x_1,\ldots,x_4$とし、$x_0=0$と置く。
十個の一次式$a_i$として$x_r-x_s$ ($0\le r<s\le4$)を取り、
その零点の行集合を$P_i$とする。以下の小さい行列計算は
\path{computations/verify_finite_window_rigidity.py}で有理数上再現できる。
\begin{enumerate}
\item 各$P_i$は七行で、$Q$の制限は階数6である。その評価七行の唯一の依存の係数は全て非零である。
$P_i\cup P_j$への制限は$i\ne j$で階数9である。他の九つの一次式は常に$U$を張る。
\item $P_i\cup P_j\cup P_k$の補で$X$が階数4なら三つ組を採用する。
採用されるのは120組中80組である。任意の一対から始め、採用した三つ組の一辺があれば残り二辺も加えると、
どの始点からも全45対に達する。
\item 最初の平面と各他の平面の和集合の補には、どの四行も独立な五行がある。
十平面は共通点を持たない。
\end{enumerate}
$L$の平方が$Q$に入るとする。第一の依存の係数を対角に置けば七次元非退化対称形式を得る。
$L$の七行への制限は全等方なので次元は3以下である。
従って各$i$で$P_i$上零になる$0\ne h_i\in L$を選べる。
$P_i$上零の$Q$の元はちょうど$a_iU$である。これは次元が4で、$U$の二次単項式は独立だからである。
また$P_i\cup P_j$上零の$Q$の元は$ka_i a_j$である。

$h_i h_j\ne0$ならこれは$a_i a_j$の非零スカラー倍である。そのような対が存在する。
全対で零なら十個の$h_i$の非零座標の台が互いに交わらず、四次元空間内で十個が独立になってしまう。
一対$i,j$で積が非零のとき、採用三つ組$i,j,k$を取る。
もし$h_i h_k=0$なら$h_k$は$P_i\cup P_j$の補で零である。
$h_k^2=a_k b$ ($b\in U$)なので、三平面の補で$b$が零になり、階数4から$b=0$となる。
被約性と$h_k\ne0$に反する。同じ議論を$j$にも用いて第二の計算を適用すれば、全対で積が非零となる。

従って各$i$で$h_i L=a_i U$である。右辺の基底を他の$a_j$との積から取れるためである。
$P_i$の補では右辺に基点がなく、$h_i$はそこで消えない。
この等号の与える線型同型から、$L$と$U$のその補上の評価点列は一つの射影同型で対応する。
最初の補との共通部分に射影枠があるので、これらの射影同型は全て等しい。
補が全二十行を覆うため、各行の比例定数をまとめた単元$d$で$L=dU$となる。
$d^2Q\subset Q$なら$d^2$はスカラーであることは命題\ref{prop:finite-square-ratio}の証明による。
また$gU=U$なら$gQ=Q$なので$g$もスカラーである。従って符号の数は表示通りである。
前節で求めた$U$での接次元は0であり、各符号倍はスキームの自己同型なので全点が被約である。
最後に単元倍は第二窓の積の核を変えず、$U$自身は対称完全交叉からAS延長を持つ。
定理\ref{thm:finite-square-chart}と系\ref{cor:two-window-extension}で一意性を得る。
\end{proof}

\begin{theorem}[一般の第一窓に対する非延長]\label{thm:generic-finite-window-obstruction}
四次元の$E,F$を固定する。$\Gr(6,E\otimes F)$の稠密開集合上で、
第一関係空間$R$は四つの点対応が全て有限な標準AS対象に延びない。
より詳しくは、前定理の第一窓を含む開集合で、有限枝の第二窓は基底変更を除き
$R^{\mathrm{op}}\subset F\otimes E$に限られる。ここで$\mathrm{op}$は二文字の交換を表す。
\end{theorem}
\begin{proof}
順序付き二十点の配置を動かし、一次式と二次式の評価階数がそれぞれ4と10の開集合を取る。
各点の斉次座標を局所的に自明化する。平方が$Q$に入る四空間の相対Grassmann閉スキームは基底上固有であり、
各符号倍はその切断となる。前定理のファイバーでは全点の相対接空間が零なので、射はそこで不分岐である。
固有性により不分岐でない部分の像を除ける。不分岐射の切断は開埋込みであり、ここでは分離性から閉でもある。
全切断の補の像も閉で、前定理の点を含まない。従ってその近傍では全解が符号切断で尽くされる。
異なる切断の一致する閉集合も除けばこの有限被覆の次数は$2^{19}$のままである。

関係空間のGrassmann多様体上で、六本の双一次式の共通零点が被約な長さ20の完全交叉となる条件は開である。
実際、例の近傍で普遍零点は射影的で相対次元零、各零点で六式の相対Jacobi行列は可逆なので有限エタールである。
点を順序付ける有限エタール被覆と各行の自明化の上で、第一射影へ先の開条件を引き戻す。
結論は行の順序と単元倍に不変なので第一窓の開近傍を得る。
有限枝の延長があれば定理\ref{thm:finite-square-chart}から第三空間は第一空間の単元倍である。
従って第二の積の核は$R^{\mathrm{op}}$となる。

$A=R\otimes E\subset E\otimes F\otimes E$とし、両端の交換を$\sigma$と書く。
必要な三次交叉は$A\cap\sigma A$であり、標準AS対象では四次元である。
ところが一般の$R$でこの交叉は零となる。検算ファイルの整数係数の六列では、
$A$と$\sigma A$の挿入行列を並べた$64\times48$行列が標数101で階数48を持つので、
同じ小行列式は標数零でも非零である。階数48の開集合と先の非空開近傍は
Grassmann多様体の既約性で稠密開集合に交わる。これが最初の主張を与える。
\end{proof}

\begin{proposition}[反転する二窓の障害と二周期性]\label{prop:finite-flip-obstructions}
第一窓の点対応が有限とし、第二窓が$R^{\mathrm{op}}$であるとする。
三次交叉は両端を交換する対合の正固有空間に入り、その次元は
\[
 R\otimes E\longrightarrow \bigwedge^2 E\otimes F,\qquad
 (e\otimes f)\otimes e'\longmapsto(e\wedge e')\otimes f
\]
（単純テンソルからの制限）の余階数である。この行列は$24\times24$であり、
AS延長があれば階数は20、第三階数は$\rho\le10$となる。

さらに前定理の符号倍による復元が両射影に成立する被約有限枝では、任意のAS延長の周期は2を割る。
この場合、以下の$36\times21$行列の核はちょうど一次元で、
その生成元は六次元空間$R$上の係数行列として可逆である。
\[
 \operatorname{Sym}^2 R\ \longrightarrow\
 \bigwedge^2 E\otimes\bigwedge^2 F .
\]
写像は$R\subset E\otimes F$を用いた二組の交代化で、$r\otimes r$では
$4\times4$行列$r$の二次小行列式を並べる規約を取る。
\end{proposition}
\begin{proof}
両端交代な三次交叉の元は、$R$の基底六式の、第一文字の一次式を係数とする多項式恒等式を与える。
この一次関係は零である。実際、六式の有限完全交叉のKoszul分解を$(2,1)$でひねると、
外次数2から5の項の直線束$\mathcal O(2-j,1-j)$は全コホモロジーが零であり、
外次数6の非零項は$H^6$だけである。従って外次数1の大域切断から
$H^0(\mathcal O(2,1))$への写像は単射である。
交叉は対合で保たれ、正固有部分はちょうど表示した交代化の核である。
AS延長の標準テンソルはこの三次交叉と第四文字のテンソル積に入るので、
反対位置$(0,2)$では$\operatorname{Sym}^2 E$に入り、階数は10以下となる。

両射影の復元を順に使うと最初の四文字の座標空間を
$E,F,E,F$と同定でき、最初の三つの関係は$R,R^{\rm op},R$となる。
ここで第二の適用では最初の反転窓の第一射影が元の第二射影である。
先の交代部分の消滅を続く窓にも使えば標準テンソル$w$は
二つの端点交換$(0,2),(1,3)$のそれぞれで不変である。
従って第四の関係も$R^{\rm op}$、標準四周期表示の関係列は二周期となる。
さらに$w\in R\otimes R$は二組を交換しても不変なので$\operatorname{Sym}^2 R$に入り、
隣接切断階数6よりその係数行列は可逆である。
標数零の分解
\[
 \operatorname{Sym}^2(E\otimes F)=
 (\operatorname{Sym}^2 E\otimes\operatorname{Sym}^2 F)
 \oplus(\bigwedge^2 E\otimes\bigwedge^2 F)
\]
より、表示した写像の核は両方の端点交換で不変なテンソルからなる。

核の任意の元$w'$は第一と第二の関係条件を満たすので、
最後の文字を縮約して得る写像は同じ四次元三次交叉に入る。
$w$の対応する写像は同型だから、最後の$F$にある$g\in\operatorname{End}(F)$を
作用させて$w'$を得る。$w'$の最後の二文字は$R$に入り、$w$の隣接切断は全階数なので
$(1\otimes g)R\subset R$である。
各零点対$(x,y)$について$(x,g^*y)$も全関係で消える。$x$上の核は直線のため
$y\wedge g^*y=0$である。第二射影では二次式の評価が単射なのでこれは多項式恒等式となり、
$g$はスカラーである。従って核は$kw$で尽くされる。
二十点例では二行列の階数はそれぞれ20と20で、後者の核の係数行列は可逆である。
一般の階数条件からAS正則性を結論するにはなお全次数条件が要る。
\end{proof}

% !TeX root = ../papers/nc-p3.tex
\section{被約二十点を通る反転二窓の成分}
\label{sec:reversible-component}
$E,F$を四次元とし、$\delta$を命題\ref{prop:finite-flip-obstructions}の二重交代化とする。
第一関係空間$R\in\Gr(6,E\otimes F)$に対し、$\delta|_{\operatorname{Sym}^2R}$の階数が
20以下となる小行列式スキームを$Z$とする。以下の次元は以前のClifford点での
一次変形の次元を用いず、この$Z$の局所環から計算する。

\begin{theorem}[反転二窓の成分と幾何学的一般点]\label{thm:reversible-component}
被約二十点例の第一窓$R_*$を通る$Z$の成分は一つであり、その次元は44である。
$R_*$での接次元は45で、滑らかな周囲の45変数の形式座標における局所方程式は
階数6の二次式を初項とする一式である。
この成分の幾何学的一般点では、第一窓は四辺が有限な標準Koszul AS正則対象へ一意に延び、
周期は2を割り、$8\le\rho\le10$である。
第一窓の基底変換を除いた成分の有理商の次元は14で、その幾何学的一般点が表す添字保持同型類はこの商に従う。
\end{theorem}
\begin{proof}
例で$\delta$の核は可逆な六行対称行列$H$が張る直線である。
$R_*$の選んだ基底列を$P$ ($16\times6$)とし、$C=PHP^{\mathsf t}$と書く。
階数がちょうど20の近傍では、核直線は小行列式から正則に復元される。
従って$Z$はそこで、対称階数6の射影行列$[C]$で$\delta(C)=0$かつ上記の開条件を満たすものと
スキームとして同型である。逆射は$C$の像を取る。
列が独立な行列$P'$から$[P' H P'^{\mathsf t}]$への射は、相対次元16の滑らかな被覆である。
実際、像空間を局所的に自明化すれば可逆な対称形式を基底変換で$H$へ移す問題となり、
標数零ではエタール局所的に実行できる。ファイバーは直交相似群の捩れで、次元は$15+1$である。

この被覆の96変数を$P'=P+YH^{-1}$と取る。36個の方程式はちょうど
\[
 G(Y)+\delta(YH^{-1}Y^{\mathsf t})=0,\qquad
 G(Y)=\delta(PY^{\mathsf t}+YP^{\mathsf t})
\]
である。$G$の階数は35で、その余核の双対を張る行を$\lambda$とすると、
$\ker G$上の二次式 $q(Y)=\lambda\delta(YH^{-1}Y^{\mathsf t})$の階数は6である。
\path{computations/verify_reversible_component.py}では第一・第二座標の積の核を$P$、
二重交代化の核を$H$としてこの二つの階数を有理数で検算する。
形式陰関数定理で独立な35式を解くと、残る61変数の一式の初項は$q$である。
独立式の解の二次以上の補正に$G$を施しても$\lambda$で消えるためである。
階数6の斉次二次式は既約で、完備局所環の随伴次数環は多項式環をこの二次式で割った整域となる。
従って被覆の局所環は整域で、次元60、接次元61である。
滑らかな16変数を除けば元の局所環は次元44、接次元45で、一つの成分を持つ。
完備局所環への滑らかな変数の添加は最小表示の関係式数や二次初項の階数を変えないので、
元の局所方程式についても主張を得る。

対称行列$C$を$(0,2)\mid(1,3)$で読み替え、対称な添字対それぞれ十個に制限した行列を$B(C)$とする。
例では$B(C)$の階数は6である。上記検算では$Y\in\ker G$で$q(Y)=0$かつ、
$B(PY^{\mathsf t}+YP^{\mathsf t})$が$\ker B(C)$から$\operatorname{coker}B(C)$に
階数2以上を誘導する有理方向を求める。可逆な六行小行列を局所消去すれば、残る四行四列ブロックの
ある二次小行列式の二次初項がこの方向で非零である。
従って初項は$q$の倍数ではなく、局所環でもこの小行列式は非零である。
一般点の第三階数は8以上となり、上界は両端対称性から従う。

例の近傍で四窓の一因子階数は4、両隣接切断の階数は6である。
後者は対称性と$C$の階数6による。三次・四次交叉は上半連続で、
$w=C$とその縮約からそれぞれの下界4と1も成立するので、開近傍で例と同じ次元になる。
定理\ref{thm:frobenius-tensor-reconstruction}を適用すると、二次双対は
各次数の次元が一定の有限Frobenius代数族になる。実際、四頂点に折り返した各成分は
関係を挿入する行列の余核で、五次は零、関係が二次のため以後も零である。
階数一定の各成分と積はこの近傍でベクトル束とその射を定める。
四つの点対応は$R,R^{\rm op}$の反復であり、例での有限エタール性も開条件である。

内部次数$n$を固定した相対被約bar複体は有界なベクトル束複体である。
例ではKoszulなので全ての非対角ホモロジーは零であり、固定した$n$ごとに同じ消滅が
例を含む開集合で成立する。既約成分の幾何学的一般点はこれら全てに属するので、
そこでは全次数でKoszulである。有限Frobenius双対の完全対合から左右のAS条件と標準Hilbert級数が従う。
同じ証明は非可算な代数閉体上の非常に一般の閉点にも適用できる。
単一の開近傍の全閉点で全次数を認証したことにはならない。

両射影で定理\ref{thm:finite-all-square-roots}の剛性が成立する開集合も例を含む
（第二射影は例の二十行を組ごとに交換したもの）。
そこで各AS延長の次の窓は反転窓に限られ、二窓の一意性から第一窓の軌道が添字保持同型を決める。
両射影の点列は射影枠を含み、第一窓の安定化群は二十点への置換に忠実に作用する。
両射影の枠を点ごとに固定する射影変換は恒等だからである。
従って$\PGL(E)\times\PGL(F)$の軌道次元は30で、有理商の次元は$44-30=14$である。
\end{proof}

\begin{corollary}[対称部分に横断的な二次錐]\label{cor:reversible-quadratic-cone}
$Z$の$R_*$における完備局所環は、39個の自由変数と六変数の非退化二次錐の積である。
特異部分の形式近傍は、第一窓を対称化できる完全交叉の部分に一致する。
\end{corollary}
\begin{proof}
$E=F$と同定する。$R\subset\operatorname{Sym}^2 E$の開部分
（直交する四二次式に基点がないもの）は24次元で、対称完全交叉の定理から$Z$に入る。
二つの空間の相対的同定を$\PGL_4$で動かすと、例の近くで滑らかな39次元部分を得る。
実際、$(1\otimes g)R$が対称に入る条件は$g$の十六係数で線型で、$R_*$で階数15である
（前記検算）。階数15の開部分では唯一の対称化写像の直線を小行列式から復元でき、可逆性も開である。
ここで階数15の開部分は対称化可能な軌跡内で取る。
対称な各$R$では$G$の像は、$\bigwedge^2 E\otimes\bigwedge^2 E\to\bigwedge^4 E$
の積で消える。実際、$P$の各列は対称二階テンソルであり、一列と任意の列との積を四文字で
全交代化すると零である。従ってこの39次元部分に沿って微分の階数は35以下で、
近傍では35となる。ゆえに全体が特異軌跡に入る。

一式の二次初項の非退化な六変数を取り、残る39変数をパラメータにする。
六つの偏微分を形式陰関数定理で解いて臨界点を原点へ平行移動し、六変数を含む高次項を
次数順に非退化二次式の偏微分で消せば、方程式は$q_6(z)+g(t)$となる（標数零）。
特異部分は$z=0$に入るが、上述の39次元部分の完備局所環は滑らかで同次元なので、
その部分は$z=0$全体を占める。従って$g=0$であり、特異部分のスキーム構造も主張通りである。
\end{proof}

\begin{corollary}[一般最小周期]\label{cor:reversible-generic-period}
この成分の幾何学的一般点のAS延長の最小周期は2である。
\end{corollary}
\begin{proof}
一移動との添字保持同型があれば、第一窓は基底変換で$R^{\rm op}$に移る。
この条件は例の近傍で閉である。点を順序付けたエタール被覆で、各順列に対し、
二つの射影行列が第一・第二射影の対応する点を入れ替える条件を課す。
行列を射影化した全線型写像にコンパクト化し、行の比例を二次小行列式で表す。
この閉条件で特異な射影行列は現れない。実際、三次評価の独立性からベクトル次元
$r=1,2,3$の空間に属する行は高々$1,4,10$本である。
階数$r$の行列は次元$4-r$の核以外の全行を次元$r$の像に送るが、両方の上界の和は20未満である。
従って固有射の像は求める閉集合である。点対応から双一次式の空間も復元できるので、
全点を対応させることは窓を対応させることと同値である。

例の第一射影の二十点の自己同型群は$S_5$である。
七点以上を含む平面は先の十平面で尽くされることは、独立な三行が張る全平面を列挙して検査できる。
その法線の位置は五頂点の完全グラフの辺であり、三法線が従属なのは三角形の三辺の場合に限る。
従って置換は辺の接続関係を保ち、互いに接続する四辺の星から五頂点を復元できる。
全法線を固定する射影変換は、座標とその差を固定するので恒等である。
逆に$v_4$の符号を変えて五つの枠ベクトルの和を零にすれば、その任意の置換は
各対の和・差の二点とその相手を保つ。従って窓を保つ基底変換は対角な$S_5$であり、
反転と合成して窓を保つものはこれと因子交換との積で尽くされる。

射影スカラーを含めた基底変換軌道の接空間を$O$とする。$C$の対称行列表示ではこれは
$LC+CL^{\mathsf t}$、$L=u\otimes1+1\otimes v$が張る31次元空間である。
$S_5$の七共役類の各代表と因子交換との積を$J$とすると、検算ファイルでは$JCJ^{\mathsf t}=C$となる。
各代表に対し$Y\in\ker G$のある基底列で$D=YP^{\mathsf t}+PY^{\mathsf t}$とおくと、
$JDJ^{\mathsf t}-D$を$O$に加えた階数が標数101で32になる。
軌道次元と合わせ、標数零でも$JDJ^{\mathsf t}-D\in O$は真の線型条件である。
共役で他の全置換にも同じことがいえる。

階数6の既約二次錐$q=0$は$\ker G$を線型に張るので、これら有限個の真の部分空間を
全て避ける滑らかな方向を取れる。これは形式曲線の一階項になる。
実際、先に35式を消去した残りの式に変数$t(v+u)$（$v$は選んだ方向、$u(0)=0$）を代入し、
$t^2$で割った式を$q$の$v$での非零偏微分で解けばよい。
この曲線の一般点で反転窓に同型なら、二十点の順序を固定した後には、ある一定の順列で対応する。
特殊点でも行列が可逆であり、対応する五点の射影枠から変換は形式曲線上に一意に延びる。
従って射影スカラーを含めて一次で比較すると、特殊点の対応$J$に対し
$JDJ^{\mathsf t}-D\in O$となって矛盾する。
反転同型の閉集合は成分全体を含まず、一般点は一周期にならない。二周期性は命題\ref{prop:finite-flip-obstructions}による。
\end{proof}

\begin{proposition}[第三階数6と7の局所イデアル]\label{prop:reversible-low-opposite-ideals}
$Z$の$R_*$での局所環で、対称化可能な部分の被約イデアルを$I$とする。
第三階数$\rho\le6$の小行列式イデアルは$I$、$\rho\le7$のものは$I^2$である。
従ってこの点の開近傍では、第三階数6の点は対称部分に尽き、ちょうど7の点は存在しない。
\end{proposition}
\begin{proof}
$B(C)$の可逆な六行六列部分を消去し、残りを四行四列のSchur補行列$S$とする。
対称完全交叉では三つの二対二切断の階数が同じ6なので、$S$の全成分は$I$に入る。
$R_*$での$B(C)$は対称であり、$K$をその四次元核の基底行列とする。核と余核の基底を合わせると
$S$の微分は$T(Y)=K^{\mathsf t}B(PY^{\mathsf t}+YP^{\mathsf t})K$で計算できる。
被覆の接空間$\ker G$上でこれは交代行列全体に全射である。
具体的には、前記検算で$T+T^{\mathsf t}=0$、上三角六成分の線型階数6、
$\operatorname{Pf}(T)=15q/16$となる（定数は検算の基底による）。
滑らかな被覆の垂直方向は$T$で消えるので、元の接空間でもその階数は6である。

系\ref{cor:reversible-quadratic-cone}の局所表示により、$I$は六つの法線座標で生成され、
その一次部分は独立である。$\mathfrak m$を極大イデアルとすれば、$S$の
$I/\mathfrak m I$での成分は上記の一般交代行列の六座標となる。中山の補題から$S$の成分は$I$を生成する。
一般の四行交代行列の二次小行列式は六座標の全二次式を張る
（係数行列の階数は21）。従って$S$の二次小行列式の像は
$I^2/\mathfrak m I^2$全体を張り、再び中山の補題でそのイデアルは$I^2$になる。
六行を消去した後のこれら二つのイデアルが第三階数の各小行列式イデアルそのものである。
\end{proof}

% !TeX root = ../papers/nc-p3.tex
\section{一点を指定した反転窓の標準形と一般切断階数}
\label{sec:marked-opposite-rank}
\begin{theorem}\label{thm:marked-opposite-eight}
定理\ref{thm:reversible-component}の44次元成分の一般点の第三階数は8である。
従って14次元の有理商の幾何学的一般点では、標準AS正則性、最小周期2、第三階数8が同時に成立する。
\end{theorem}

隣接切断階数6の反転テンソルを、非退化対称対合を持つ六次元空間$H$への双線型写像
$\beta:E^*\times F^*\to H$のGramテンソルで表す。二重交代化の消滅条件は
\[
 \langle\beta(e,f),\beta(e',f')\rangle=\langle\beta(e,f'),\beta(e',f)\rangle .
\]
零点対$([e_0],[f_0])$では$\beta(e_0,f_0)=0$である。
$U=\beta(e_0,F^*)$と$V=\beta(E^*,f_0)$が三次元で、$U$の対合が非退化だとする。
両者は直交し、$H=U\perp V$となる。直交標準基底と一次空間の補空間を取ると、表示は
\[
 \beta((r,x),(s,y))=(r y+A_x y,\ s x+B_y x),\qquad x\in V,\quad y\in U
\]
となる。一文字に$e_0$または$f_0$を入れたGram等式は各$A_x,B_y$の対称性を与える。
補空間を変えると各作用素に恒等写像の線型なスカラー倍を加えられ、跡を零にできる。

三つの三行対称行列$X=(X_1,X_2,X_3)$について、巡回順$(23),(31),(12)$を使い
$D_X$の第$(ij)$列、第$(mn)$行を$[X_i,X_j]_{mn}$とする。
上の$A,B$を基底上の三行列として読むと、残る九つのGram等式はすべて
\[
 D_B=-D_A^{\mathsf t}
\]
にまとめられる。これは第$(ik),(jl)$の式が
$[A_i,A_k]_{jl}+[B_j,B_l]_{ik}=0$であることから従う。

\begin{lemma}[三行の交換子と直交二次式]\label{lem:ternary-net-commutator}
対称行列$S$が$\operatorname{tr}S=\operatorname{tr}(SX_i)=0$を満たすとき、
$P_X(S)_{ij}=\operatorname{tr}(SX_iX_j)$とおけば
\[
 D_X P_X(S) D_X^{\mathsf t}=(\det D_X) S .
\]
\end{lemma}
\begin{proof}
$P_X$は転置と跡の巡回性から対称である。各$X_i$へのスカラー行列の加法は両辺を変えない。
三行列の並べ方を列変換$g$で変えると$D$は右から$(\det g)g^{-\mathsf t}$を掛けられ、
$P$は$g^{\mathsf t}Pg$となるので等式は共変である。三次元空間の向きを保つ直交変換にも共変である。
非零な無跡対称$S$を動かしたとき$\operatorname{tr}(SX_i)=0$の三つ組はベクトル束をなし、
この全空間は既約である。従って相異なる固有値を持つ一般の$S$と一般の三つ組で証明すれば足りる。
$S=\operatorname{diag}(p,q,r)$、$p+q+r=0$とし、
$L=\operatorname{diag}(q-r,r-p,p-q)$とする。基底変換とスカラーの加法により
$X_i$の非対角部分を順に$E_{23}+E_{32},E_{31}+E_{13},E_{12}+E_{21}$、
対角部分を順に$aL,bL,cL$とできる。
$v=(a,b,c)^{\mathsf t}$、$[v]_\times$を$v$との外積の行列とすると直接乗算で
\[
 D=-I+3S[v]_\times,\qquad
 P=-S+\operatorname{tr}(SL^2)vv^{\mathsf t}.
\]
この二式に$r=-p-q$を代入して$DPD^{\mathsf t}-(\det D)S$を展開すると全九成分は零である。
\path{computations/verify_marked_opposite_rank.py}は同じ計算を五つの不定元を残して検査する。
$S=0$での等式も直ちに従う。
\end{proof}

\begin{proof}[定理\ref{thm:marked-opposite-eight}の証明]
まず$D_A$が可逆な表示を考える。対称対の最初の四つを$00,01,02,03$とすれば、
反対切断の十行十列行列は$\left(\begin{smallmatrix}0&L_A\\L_B^{\mathsf t}&W\end{smallmatrix}\right)$である。
上の四行は六次元の対称行列空間で$I,A_1,A_2,A_3$、左の四列は$I,B_1,B_2,B_3$に対応する。
$D_A,D_B$が可逆ならこれらは各々独立である。両側の直交補の行列をそれぞれ$S,T$とすると、
消去後の残り二行二列の対合は
\[
 \operatorname{tr}(T P_A(S))+\operatorname{tr}(S P_B(T))
\]
である。実際、右下の$ik,jl$成分は$(A_iA_k)_{jl}+(B_jB_l)_{ik}$の対称対への制限である。
補題から$P_A(S)=\det(D_A)D_A^{-1} S D_A^{-\mathsf t}$となり、
$D_B=-D_A^{\mathsf t}$と列数3の符号で二つの跡が打ち消し合う。従って第三階数はちょうど8となる。

次にこの開条件が目的の成分の一般点で成立することを確かめる。
二十点例の第一の零点対を指定する。点を指定した有限エタール被覆の完備局所環は
元の点のものと同じである。$P,H,Y,G,q$は定理\ref{thm:reversible-component}のGram表示とする。
選んだ$e_0,f_0$を標準行の先頭へ置き、各々に残り三つの座標行を補う。
十六行の変換後の$P$のうち$(0,j)$の三行を$U_0$、$(i,0)$の三行を$V_0$と書く ($i,j>0$)。
$J=U_0 H U_0^{\mathsf t}$は可逆である。$(i,1),(i,2),(i,3)$の行列を$P_i$とすると、
直交化前には$C_i=P_i H U_0^{\mathsf t}$が対称で、交換子は
$C_iJ^{-1}C_j-C_jJ^{-1}C_i$で計算する。基底を直交化してもその三列の独立性は変わらない。

例ではこの交換子行列は$\operatorname{diag}(-15/32,0,0)$である。
$\ker G$上の変分で点も動かし、可逆な最初の一行を消去した残り二行二列の微分を取る。
具体的には$\dot P=YH^{-1}$に対して六つの零点条件を微分し、
$V_0,U_0$を重ねた六行六列の転置で割って点の六座標の変分を得る。
上記ファイルの有理な接線の直線上では、接錐の方程式と二行二列の微分の行列式は順に
\[
 -(928t^2+704t+232)/15,\qquad (112t^2+184t-37)/256 .
\]
両式は互いに素なので、後者の斉次二次式は$q$の倍数でない。
完備局所環の随伴次数環が$(q)$で割った整域であることから、この小行列式の二次初項は消えず、
交換子行列の行列式は成分上で恒等的に零にならない。従って一般点で使った開条件が成立する。
さらに特殊ファイバーの二十の印は対角$S_5$と両因子交換で推移的に移り、
両因子交換は同じ成分と可逆条件を保つ。よってこの成分の一般点では二十通りの印のどれも使える。
\end{proof}

\begin{proposition}[印付きの一般窓を構成するパラメータ]\label{prop:marked-net-parameters}
上の可逆な開部分の印付きテンソルは、独立な無跡対称三行列$A$と、別の無跡対称三行列空間内の
三次元部分空間$\mathcal B$から構成できる。いずれの交換子行列も可逆なものを取る。
$\mathcal B$の基底$B^0$を取り、$K=D_{B^0}^{-1}(-D_A^{\mathsf t})$とすれば
\[
 B=B^0 g,\qquad g=\pm\sqrt{\det K}\,K^{-\mathsf t}.
\]
印を保つ同型ではこの二符号は同じである。正規化された二つの表示が同型となる条件は
$O_U\in\mathrm O(U), O_V\in\mathrm O(V), c\in k^\times$で
\[
 A'_{O_V x}=c O_U A_x O_U^{-1},\qquad
 B'_{O_U y}=c O_V B_y O_V^{-1}
\]
とできることである。従って上記14次元の有理商は単有理的である。
\end{proposition}
\begin{proof}
三列の基底変換の公式$D_B=D_{B^0}(\det g)g^{-\mathsf t}$を解くと表示式となる。
このとき九つのGram等式を満たし、定理の点付き開部分の全てのデータが得られる。
六次元空間への写像は全射であるから、Gramテンソルの射影同型は$H$での相似変換を誘導する。
指定の二直線を保つ基底変換を$(r,x)\mapsto(a r+u(x),C x)$、
$(s,y)\mapsto(b s+v(y),D y)$と書くと、相似変換の二つの直交因子への制限は
$aD=d O_U$, $bC=d O_V$と書ける。同型条件の跡を取ると、無跡性から$u=v=0$になる。
残る式は$c=ab/d$とした表示条件であり、逆にこの条件から基底変換を戻せる。
$O_U=-I,O_V=I,c=1$で$B$の符号だけを変えられる。
最初の関係空間に戻るときは二重交代化の核がただ一つの直線である開部分を使う。

$A$と$\mathcal B$のパラメータの基底はアフィン15空間と$\Gr(3,5)$の積の開部分で有理的である。
平方根を選ぶ被覆の二点は同じ軌道を与えるので、この基底から有理商への写像は有理関数体上に降下する。
目的の成分の幾何学的一般点はこの構成に含まれ、テンソルの局所成分は例で一つだから、
パラメータからの像の閉包はその成分の商となる（正則性を使うのは幾何学的一般点）。
これにより支配有理写像を得る。特殊な全データのAS性には別の判定が必要である。
\end{proof}

% !TeX root = ../papers/nc-p3.tex
\section{楕円曲線の切断環からの反転窓と全次数認証}
\label{sec:elliptic-reversible-pencil}
滑らかな種数1の射影曲線$C$、順序付きの次数4の直線束$L,M$、$N=LM$と
基点を持たない鉛筆$W\in\Gr(2,H^0(N))$を取る。直線束の積はテンソル積を表す。
本節の曲線は元代数の点加群を並べた曲線ではない。反対切断の左右の核にある二次式から復元する。

\begin{theorem}[楕円鉛筆の全ファイバーの正則性]\label{thm:elliptic-pencil-all-as}
一次空間を$E=H^0(L)^*,F=H^0(M)^*$と置く。第一窓の六次元の関係空間を
\[
 R=\ker\bigl(H^0(L)\otimes H^0(M)\longrightarrow H^0(N)/W\bigr)^\perp
 \subset E\otimes F
\]
とし、$R,R^{\rm op}$を交互に繰り返すと標準14641型の左右Koszul AS正則対象を得る。
$L\not\simeq M$なら第三階数は常に8であり、その両側の二次核はそれぞれ
$C\subset\mathbb P(E), C\subset\mathbb P(F)$を定義する楕円四次曲線の二次式の鉛筆である。
$L\simeq M$の場合は四二次式の対称完全交叉のチャートに入る。
\end{theorem}
\begin{proof}
$L_i$を交互に$L,M$とし、$D_{ij}=L_i\cdots L_{j-1}$、$D_{ii}=\mathcal O_C$とする。
次数$i,j$に$H^0(D_{ij})$を置き、切断の積で結んだ代数$S$がKoszulであることを示す。
各始点$i$で$K_0=\mathcal O_C$とし、
\[
 G_n=K_n L_{i+n},\qquad
 0\longrightarrow K_{n+1}\longrightarrow H^0(G_n)\otimes\mathcal O_C
 \longrightarrow G_n\longrightarrow0
\]
を順次作る。次数2以上の直線束は大域生成されて$H^1=0$である。
次数2の束$D$の評価核は$D^{-1}$である。次数3または4では滑らかな点$p$を取り、
$D(-p)$と$D$の評価列を比較する。切断列は全射であり、商の評価は
$\mathcal O_C\to k(p)$なので、$D$の評価核は$D(-p)$の評価核と$\mathcal O(-p)$の拡大である。
よって次数2,3,4の直線束の評価核は次数$-2,-1$の直線束を商に持つフィルトレーションを許す。
この三次数の束でフィルターされた束にも同じことが成立する。
実際、拡大の切断列は$H^1=0$で完全、商と部分束が大域生成されるので中項も大域生成され、
三つの評価核も完全列をなす。従って$G_0=L_i$から始め、$K_{n+1}$に次の次数4の束を
掛けると再び次数2,3の直線束でフィルターされる。

上の評価列を$D_{i+n+1,j}$でひねり、大域切断を取る。尾部の長さが1以上なら
$K_{n+1}D_{i+n+1,j}$のフィルトレーションの各商は正次数で$H^1=0$、
長さ0なら評価は大域切断上恒等である。従って始点単純の線型分解が得られる。
具体的に第一射は$H^0(L_i)\otimes e_{i+1}S\to e_iS$、その核の位置$j>i$での
成分は$H^0(K_1D_{i+1,j})$である。以後も同じ議論で、この核を次の位置の
$H^0(G_1)$が生成し、各段の核と生成元が同様に続く。逆順でも成立する。
特に$S$は一次生成二次関係である。

偶奇二頂点に折り返すと、$W$の基底$a,b$はそれぞれ両頂点からの長さ2の切断の和として中心元になる。
$a$は正則で、$b$も$a$を割った後で正則である。実際、$bs=at$なら基点自由性から
$s$は$a$の零因子を重複度込みで含み、$s/a$は所定の尾部の直線束の大域切断となる。
次数が2未満ならその束の次数は負である。
$H=S/(a,b)$はKoszulである。一般にKoszul代数から次数2の中心正則元$z$を割るとき、
環変更Torスペクトル系列の二行は商のTorとその内部次数2の移動である。
商で最初に非対角成分が出るホモロジー次数$p$を取れば、第零行からの唯一あり得る微分の先は
次数$p-2$の第一行にあり、最小性から該当成分は零である。非対角成分が元の代数のTorに残って矛盾する。
この証明は次数零の代数が二頂点の半単純代数でも同じである。
長さ$n>0$の$S$の次元は$4n$なので、二つの次数2の正則元を割った各始点の次元列は$(1,4,6,4,1)$となる。

随伴完全対合も直接得られる。任意の$D=D_{ij}$、$0\le j-i\le4$について基点自由鉛筆の完全列
\[
 0\to DN^{-2}\to (DN^{-1})^{\oplus2}\to D\to0
\]
を使う。$Q(D)=H^0(D)/W H^0(DN^{-1})$と書けば、連結写像とSerre双対により
$Q(D)\simeq Q(N^2 D^{-1})^*$を得る。これは切断の積の後に
$H^0(N^2)/W H^0(N)\simeq H^1(\mathcal O_C)$へのtraceを取る対合である。
従って$H$は有限Frobeniusで、相補次数の対合により二次双対の左右のKoszul単純分解は
四移動して双対になる。大域次元4、左右AS条件と標準Hilbert級数が得られる。
その関係は表示の$R$と反転である。

$\ell$を上記traceとする。標準テンソルは四切断の積に$\ell$を施したものなので、
反対切断は$H^0(L^2)\times H^0(M^2)\to k$の積対合を通る。
次数4の完全線型系による楕円四次曲線は二二次式の完全交叉で、対称二乗は全二次切断に全射である。
この事実は二次式が少なくとも二本あり、非退化既約四次曲線を含む独立二式は共通平面を
持てず、その完全交叉が次数4を使い切ることからも従う。
$L\not\simeq M$なら$L^2N^{-1}$と$M^2N^{-1}$は非自明な次数零の束である。
先の鉛筆完全列を$D=L^2$にも適用すると、八次元の両切断空間がそのまま完全に対合する。
従って階数と核は主張通りである。$L=M$なら切断環は通常の楕円四次曲線の環で、
二本の定義式に$a,b$の二次式持ち上げを加えた対称完全交叉になる。
\end{proof}

\begin{theorem}[曲線データの復元と被覆]\label{thm:elliptic-pencil-reconstruction}
$L\not\simeq M$の上記全対象では、添字保持同型の必要十分条件は曲線同型が順序付きの
$(L,M,W)$を対応させることである。特に最小周期1の条件は$(M,L,W)$との曲線同型である。
このチャートは定理\ref{thm:reversible-component}の44次元成分の稠密開部分を覆う。
\end{theorem}
\begin{proof}
まず添字保持同型に現れる反対二位置の基底変換を揃えられることを示す。
第三切断の像は対称二乗内の八次元部分である。四次元空間の非スカラーな$g$に対し、
対称形式$A$で$A(gv,w)=A(v,gw)$を満たす空間の次元は7以下である。
実際、$v,gv$が独立な$v$を固定すると、$w\mapsto (gv)w-v(gw)$は対称二乗への階数3の写像である。
$v$を法として見ると核は$kv$に限られる。従って対称二乗の八次元部分を対称二乗に移す
二基底変換はスカラーを除いて同一である。これを両側に適用すれば、標準テンソルの同型は
$\PGL(E)\times\PGL(F)$で比較でき、両側の二次核から二つの埋込み楕円曲線$C_E,C_F$を得る。

両曲線の間の同定もテンソルから一意に決まる。$C_E\times C_F$上の束
$P=\mathcal O_{C_E}(2)\boxtimes\mathcal O_{C_F}(2)$を取る。
テンソルはその大域切断上の汎関数であり、実際の同定のグラフへ制限した切断に$\ell$を施す。
相異なる二つの同定のグラフ$G,G'$があったとする。$P(-G-G')$は豊富である。
実際、$C_F\to\mathbb P^1$の次数2の射で同定$\phi$の前後の点が等しい像を持つ条件を取ると、
$G_\phi$と対合を合成したグラフとの和が双次数$(2,2)$の外積束の因子となる。
各グラフは差の射のファイバーとしてnefである。従って$P(-G-G')$は次数$(4,4)$の
豊富な外積束に二つのnefなグラフ因子を加えたものとなる。
曲面の標準束は自明なのでKodaira消滅より両グラフの合併への制限は全射である。
二つのグラフの線型包の交わりは、その交点スキームの線型包に含まれる。
実際、合併上の切断は交点で一致する二切断の対であり、両制限は全射である
（各グラフ上の束は次数16で、下記の交点を引いても正次数）。

異なる二同定の一致点は高々四つで、全て横断的である。差を自己同型と見て固定点があればそれを原点に取る。
原点を保つ非自明な楕円自己同型は有限位数$n\ge2$で微分は1ではない。
巡回群による商へのHurwitz公式から固定点数は$2n/(n-1)$以下である。
従って両線型包の交点のテンソルは高々四つの点評価の和であり、反対階数8に達しない。
これでグラフの一意性を得る。グラフ上の混合乗法を$\mu$とすれば
標準テンソルの隣接切断の核を$\mu$で送ったものが$W$である。これは$H$の二次成分の表示と
その完全積対合による。逆に曲線同型は全切断の積と関係を移すので十分である。

楕円曲線の原点を平行移動して$L=\mathcal O(4o)$と置け、残る原点の選択は有限である。
$M$の選択は一次元、鉛筆の選択は$\Gr(2,8)$の非空開集合で12次元である。
曲線の一次元を加えて14次元となる。これらのデータは既約な族をなし、$L=M$にも特殊化できる。
例えば原点付きWeierstrass族上で$M=\mathcal O(3o+p)$と置き、相対Grassmann束を取ればよい。
標数零で$L$を保つ曲線自己同型群は有限であり、上記復元から射影基底を加えた第一窓の像は44次元である。
ここで反転を繰り返す対象では、第一窓の基底変換はそのまま全関係列の基底変換に延び、
第一窓を固定する射影変換も上の復元により有限群をなす。

例の四つの基点自由な二次式から一般の二式を選ぶと、Bertiniの定理で滑らかな楕円四次曲線が得られる
（完全交叉は連結で随伴により種数1）。残る二式はそこで基点自由な鉛筆となる。
従って$L=M$で二十点の第一窓はこの族の像に含まれる。全てのGramテンソルは両因子対称で
隣接階数6なので、その像は前節の二重交代化の小行列式スキーム$Z$に入る。
例を通る局所成分は唯一で次元44であるから、上の既約像の閉包はその成分である。
構成可能な稠密像は成分の稠密開部分を含み、そこにある全ての幾何点が前定理による全次数認証を持つ。
\end{proof}

\paragraph{有理数上の照合。}
\path{computations/verify_elliptic_reversible_pencil.py}では、八点の評価行列$X$のうち後四点を
一平面に取り、$Y$の列を$X^{\mathsf t}$の核の基底に取る。八点を通る二二次式の交叉は滑らかで、
その超平面束を$L$、八点の和を$P_8$、$M=\mathcal O(P_8)L^{-1}$とすると、$Y$は$M$による八点の評価を表す。
実際、$L$の評価の余核の双対は、因子$P_8$の完全列とSerre双対により$H^0(M)$であり、
各行の自明化で留数の非零定数を吸収できる。$N=\mathcal O(P_8)$の切断の評価像は
成分和が零の超平面である。$P_8$を零とする切断と、評価が
$(1,2,3,4,5,6,7,-28)$となる切断で基点自由鉛筆を作ると、traceはこの八係数の逆数を重みに持つ。
これは第二切断と任意の$N$の切断との積を消すことから分かる。両側の四次判別式から得る不変量の一致、
反転窓の零点、二つの切断階数6と8を有理数上で検算する。

\begin{corollary}[周期の判定と一般の直線束に対する一周期層]\label{cor:elliptic-pencil-period}
原点を固定し、次数4の束$L,M$の点和を$l,m\in C$とする。
$L\ne M$の場合、一周期となる必要十分条件は
\[
 u\in\operatorname{Aut}(C,o),\quad (1+u)(l-m)=0,\quad 4a=l-u m
\]
を満たす$\phi(z)=uz+a$の少なくとも一つが$W$を保つことである。
$j\ne0,1728$かつ$2(l-m)\ne0$なら、固定した$(C,L,M)$の鉛筆の12次元空間で
一周期部分の次元は6である。
\end{corollary}
\begin{proof}
$\phi^*$は次数4の点和$z$に$u^{-1}(z-4a)$として作用するので前定理で判定式を得る。
表示した一般条件では$u=-1$のみで、$a$の選択は十六個である。
この対合による商を$\pi:C\to\mathbb P^1$とすると、
$N=L\otimes\phi^*L$はノルムの引き戻し$\pi^*\mathcal O(4)$である。
$\pi_*\mathcal O_C=\mathcal O\oplus\mathcal O(-2)$より切断の固有空間次元は5と3である。
二次元不変部分空間の正負次元が$(2,0)$なら$\Gr(2,5)$の基点自由な非空開部分で次元6である。
$(1,1)$なら$\mathbb P^4\times\mathbb P^2$の開部分で、一般の正切断は負切断の有限個の零点を避ける。
$(0,2)$なら分岐上の共通零点を持つ。有限個の対合の和を取っても次元は6である。
\end{proof}

\begin{proposition}[正次元の点対応を持つ楕円鉛筆]\label{prop:elliptic-pencil-dihedral}
位数4の点$t\in C$と$T=\langle t\rangle$を取り、$f:C\to(C/T)/\{\pm1\}=\mathbb P^1$を取る。
$L=\mathcal O(3o+[t])$, $M=\mathcal O(3o+[-t])$とし、
$W=f^*H^0(\mathcal O(1))$とする。この標準AS対象は一周期、第三階数8である。
第一窓の点対応の被約支持は、正規化が$C$である一つの曲線と四つの孤立被約点からなる。
\end{proposition}
\begin{proof}
一般の$f$のファイバーは$p+T$と$-p+T$の和で、八点の点和は零であるから
$f^*\mathcal O(1)\simeq LM$である。$L,M$は異なり、反転がこれらを交換して鉛筆を保つ。
零点対とは$f$のファイバー因子の順序付き分割$D+D'$で$D\in|L|,D'\in|M|$となるものにほかならない。
これは第一関係の直交空間が$\mu^{-1}W$であることによる。

$p+T$から$k$点を取ると$D$の点和は$(2k-4)p$に$T$の元を加えたものである。
一般点で可能なのは$k=2$のみである。二点ずつの添字を$I,J\subset\mathbb Z/4$とすると条件は
$\sum I+\sum J=1\pmod4$である。得られる八組は
$(I,J)=(\{0,1\},\{1,3\})$から$(I,J)\mapsto(I+1,J-1)$と交換により
安定化群なしで一軌道となる。従ってこの曲線は$C$から次数1で被覆され、他の支配曲線はない。
各ファイバーの分割は有限個なので垂直な曲線もない。

分岐外で$k=1,3$なら$2p\in T$となって分岐外に反する。$k=4$の分割は
$D=p+T$、$4p=-t$となり、$T$の平行移動を除いて四つある。$k=0$は$p$を反転した同じ分割を与える。
これらは分岐外にある。単純な八点を局所的に順序付けると、分割の枝は局所的なファイバーパラメータで
表され、$|L|$に入る式$4p=-t$の微分は非零である。従って四点はいずれも孤立被約である。
最後に$2p=ht\in T$ならファイバーは$p+T$の各点を二回ずつ含む。
$D$の多重度が$(2,2,0,0)$または$(1,1,1,1)$の置換なら添字和が偶数で、
$4p=2ht$を加えても点和は$t$にならない。残る$(2,1,1,0)$型では重複点を両組に一回ずつ、
二つの単純点を別々の組に入れて変形できる。第二組の添字を$h$だけ戻して$-p+T$で書けば、
点和の条件がそのまま$\sum I+\sum J=1$になるので上の曲線の極限に属する。
分岐点での埋め込み成分の有無はこの被約支持の記述には用いない。
\end{proof}

\begin{theorem}[節点・尖点の整な境界]\label{thm:elliptic-pencil-integral-boundary}
$C$を節点または尖点を持つ整な算術種数1のGorenstein曲線とし、$L,M$は可逆層とする。
同じ次数4の二束と基点自由鉛筆の構成では、全ファイバーの標準AS正則性、
$L\not\simeq M$での第三階数8、および曲線同型による添字保持同型の必要十分条件が保たれる。
$L\not\simeq M$の節点層と尖点層の同型類の族はそれぞれ13次元、12次元であり、
上の滑らかな楕円チャートの特殊化に入る。
\end{theorem}
\begin{proof}
双対化層は自明である（次数零で非零切断を持つ）。長さ$h$の部分スキーム$Z$について
$\mathcal I_Z^*$は次数$h$の階数1無捩れ層である。実際、
$0\to\mathcal I_Z\to\mathcal O\to\mathcal O_Z\to0$を双対化すると、
局所Gorenstein双対性により最後の$\mathcal E xt^1(\mathcal O_Z,\mathcal O)$の長さが$h$である。
$\deg D>h$なら$D\to\mathcal I_Z^*$の非零写像は次数により存在しない。
従ってSerre双対で$H^1(D\mathcal I_Z)=0$となり、
次数2以上の束は大域生成、次数3以上では長さ2の部分スキームを分離する。
前定理の評価核のフィルトレーションでは、束の次数を下げる点を滑らかな点に選べば全証明が通る。
鉛筆の列、二中心元の正則性、完全対合の証明もそのまま通る。
非自明な次数零の束は切断を持たず、次数4の埋込みは同じ次数の議論で二二次式完全交叉なので、
反対階数と両側の曲線を再び復元できる。

同定の一意性を正規化$\nu:\mathbb P^1\to C$上で示す。各側の八次元二次切断空間から、
導手イデアルを掛けた七次元部分空間を取る。正規化では、節点の両原像または尖点の二重原像で
消える切断全体であり、固定因子を除けば$H^0(\mathbb P^1,\mathcal O(6))$である。
これは節点の二枝での値の貼り合わせ、尖点での一次係数の貼り合わせからも従う。
八次元完全対合を両側で余次元1の部分へ制限すると階数は6以上である。
この制限テンソルは、各同定が与える正規化のグラフ上の$(6,6)$切断を通って定まる。
実際、正規化上のグラフで消える切断は固定因子を戻した後でも元の整な曲線上で零である。
異なる射影直線の同定二つのグラフは長さ2で交わる。
合併のイデアルでひねると$\mathcal O(4,4)$となり$H^1=0$だから、
両グラフの共通線型包にあるテンソルの行列階数は高々2である。
非被約な長さ2の交点でも、積の対合はその長さ2の座標環を通る。従って同定は一つしかない。
正規化で一致する射は整な曲線上でも一致し、$W$は前と同じ隣接核から戻る。

抽象節点曲線は一種で、滑らかな部分は乗法群である。原点を保つ群は反転を含む有限群で、
平行移動で$L$を規格化した後は$M$の一次元と鉛筆の12次元が残る。
尖点の場合、滑らかな部分は加法群で、さらに原点を保つ乗法群が作用する。
次数4の点和を$L$では0、$M$では1に規格化できるので鉛筆の12次元だけが残る。
これらの自己同型の記述は、正規化上で節点の二原像を保つ群、尖点の一原像を保つアフィン群から得られる。
特殊化については$|L|$の二二次式を一般の二式の方へ変形すれば滑らかな完全交叉へ移る。
相対超平面束で$L$を延長し、$M$は特異点を避ける一般の四点の因子で表示して、
その滑らかな四点をエタール局所的に持ち上げる。$H^1(N)=0$なので$W$も延長でき、基点自由性は開である。
隣接関係は階数一定の乗法の核の直交空間であり、これが求める特殊化を与える。
\end{proof}

\begin{proposition}[有理曲線の環状境界]\label{prop:elliptic-pencil-polygons}
$C$が二、三または四本の滑らかな有理曲線を節点で環状につないだ曲線
（二本なら異なる二点でつなぐ）とする。$L,M$が各成分で同じ正の多重次数$d$を持ち、
$\sum d_i=4$とする。基点自由な全鉛筆について先と同じ構成は標準AS正則で、
$L\not\simeq M$なら第三階数8となり、両側の二次核はこの環状四次曲線を切り出す。
これらは滑らかな族の特殊化に入る。$L\not\simeq M$の場合には曲線データによる同型判定も保たれ、固定多重次数での有理商の次元は、成分数$n=2,3,4$に対して$14-n$である。
\end{proposition}
\begin{proof}
環状曲線では、正規化の各成分の切断を二つの端点で貼り合わせる。双対化層は自明である
（各成分に、両端だけで極を持ち留数が逆となる対数微分を置く）。
非負の多重次数で総次数が正の束は$H^1=0$となる。双対の切断が負次数成分で消えれば
次数零成分の列にもその消滅が伝わるためである。総次数が2以上なら大域生成される。
実際、正次数成分が二つ以上あれば各成分上で二つの端点値を独立に指定でき、
次数零の鎖では値を定数倍で引き継ぐ。正次数成分が一つだけなら、
その両端値に一つの比例条件を課した次数2以上の多項式系になり、基点を持たない。

評価核の証明では多重次数が$0\le b\le d$の束を取る。総次数3以上の場合は
正次数成分から滑らかな点を引いて非負の総次数2まで下げる。
最後の鉛筆の逆束に次の$L_i$を掛けると、非負の多重次数$d-b'$、総次数2となる。
他の評価核の商$\mathcal O(-p)$に掛けたものも総次数3で同じ範囲に残る。
これで線型分解の帰納と、尾部をひねった核の$H^1$消滅が成立する。
鉛筆の一般の第一切断$a$はどの成分でも恒等的に零にならず、第二切断は$a$の零点で単元である。
従って正則列と双対性の証明も通る。非自明な多重次数零の束には大域切断がない。

正の多重次数$d$の完全線型系は各成分を有理正規曲線に写す。
座標は各節点の値と、各成分上で両端で消える$d_i-1$個の内部座標である。
各成分は両端と自分の内部座標が張る空間に入り、異なる成分は指定の節点でだけ横断的に交わる。
二成分の場合も、二つの二次曲線を含む両平面の共通直線、または三次の有理正規曲線の割線は、
非線型成分と両端でのみ交わる。この表示で次数3の部分曲線が一平面に入ることはない。
独立な二本の消滅二次式は次元から存在し、もし共通平面因子があればその外の成分が
残る二平面の共通直線に入るので次数3以上が共通平面に残ってしまう。
従って二式は完全交叉を定め、次数と算術種数からこの曲線に等しい。
この完全交叉と$M$の一般の滑らかな四点を動かせば、前定理と同様に平滑化できる。

同定の復元を示す。一般の$a\in W$の零点は成分$i$に$r_i=2d_i$個ずつあり、
全て単純で節点を避ける。traceはこの点評価の非零重み付き和である。
実際、鉛筆列の連結写像はもう一方の切断$b$を選んで$a$の零点で$1/(ab)$の留数を取る写像であり、
$a$の単純性、$b$の非消滅、双対化層の生成微分から各重みは非零である。
$2L$を意味する束$L^2$の切断双対で、成分$i$の評価の線型包を$P_i$とする。
$M^2$側の同様の空間を$Q_i$とし、実際のグラフで対応するよう番号を付ける。
どちらも次元$r_i+1$である。節点評価と成分内部の$r_i-1$座標による基底では、
各$P_i,Q_i$は対応する座標部分空間となる。traceテンソルは$T=\sum T_i$で、
$T_i\in P_i\otimes Q_i$は上の$r_i$個の評価から作る階数$r_i$のテンソルである。

片方で成分$i$だけを台とする切断は端点二つで消え、固定因子を除けば
$H^0(\mathcal O(r_i-2))$である。その縮約で得る像は$Q_i$内で次元$r_i-1$を持つ。
別の対応成分を仮定すれば、この像は二成分の線型包の共通部分に入る。
$d_i\ge2$では共通部分は高々二次元なので矛盾する。
$d_i=1$で三成分以上なら、異なる成分の共通部分は高々一節点の直線である。
しかし縮約の像は二つの非節点での非零重み付き評価の和であり、
二次Veronese曲線の異なる三点の独立性から節点の直線には入らない。
二成分で$d_i=1$なら他方は$d_j=3$であり、次数でも区別される。従って成分対応は一意である。

三成分以上の場合は成分対応で節点対応も定まる。二つのグラフ表示で$T_i$が異なるとしても、
その差は他の$P_j\otimes Q_j$の和に入る。座標基底で共通するのは両端に対応する
二本の節点評価テンソルだけであり、これらは双方のグラフの線型包に入る。
従って$T_i$自身が二つの成分グラフの線型包に入る。
異なる同定なら二グラフの交点はその両端二点であり、共通線型包は二つの節点評価の包に等しい
（合併のイデアルでひねると$\mathcal O(r_i-2,r_i-2)$となる）。
$r_i>2$なら階数に反し、$r_i=2$でも、二つの非節点の評価が張る割線は節点二つの割線と異なる。

二成分の場合には少なくとも一方で$d_i\ge2$となる。片側だけをその成分支持に制限した
階数$r_i-1\ge3$のテンソルは正規化上の$(r_i-2,r_i)$グラフを通る。
二同定に共通ならランクは2以下となる（合併でひねると$\mathcal O(r_i-4,r_i-2)$で$H^1=0$）。
従ってこちらの同定と両節点の対応が定まる。
同じグラフ上の二つの$T_i$の差は、両側で両節点の評価が張る二次元空間のテンソル積に入る。
この四次元空間のうちグラフの線型包に入るのは、同じ節点どうしの二項だけである。
実際、両節点を$0,\infty$と取れば、異なる両端の単項式対の積は、
両端で消える内部単項式の対の積としても書ける。後者での値は零なので交差する二係数は零である。
ゆえに残る成分の$T_j$は二つの節点項を除いて定まり、先の割線の議論でその同定も一意となる。
$W$を戻せば同型の必要十分性が従う。

最後に、各束の同じ多重次数での貼り合わせパラメータは一つの乗法群である。
各有理成分の両端を保つ群の連結部分は$n$次元トーラスで、一般の鉛筆の安定化群は有限となる。
三成分以上なら一般の鉛筆は異なる節点を異なる値に送るため、基底の射影変換は有限個である。
これは節点値を独立に指定できることから従う。二成分なら少なくとも一成分の射は次数4以上で、
一般には三つ以上の分岐値を持つ（その成分への切断の制限は全射）。
基底の射影変換を固定した後で各曲線上の有限射の持ち上げも有限個である。
従って有理商の次元は$1+1+12-n$となる。
\end{proof}

\paragraph{異なる多重次数の境界での第三切断。}
環状曲線上の$L,M$が相異なる正の多重次数を持つ場合も、鉛筆のtraceテンソルの
反対切断階数だけは全て決まる。両端の対称二乗は各々の全二次切断へ全射なので、
鉛筆完全列による完全対合からその階数は
\[
 8-2h^0(LM^{-1})
\]
である。実際、$W\otimes H^0(LM^{-1})$からの積の核は
$H^0(M^{-2})=0$である。二成分で多重次数が$(3,1),(1,3)$（逆順も可）なら、
差の束の切断は一成分上の次数2の多項式で両端で消えるものからなり、階数は6である。
他の異なる正の多重次数では差は$+1,-1$と零だけである。負次数成分で切断が消え、
両側から零次数の鎖を経て正次数成分の両端で消えるため$h^0=0$、階数は8である。

\begin{corollary}[正の多重次数の必要十分条件]\label{cor:elliptic-polygon-degree-balance}
同じ環状曲線上の任意の正の多重次数$d,e$、$\sum d_i=\sum e_i=4$を許し、
第一窓を定理\ref{thm:elliptic-pencil-all-as}の表示で定める。鉛筆が基点自由なら、
反転二窓の代数が標準Koszul AS正則となる必要十分条件は$d=e$である。
異なる場合は全鉛筆で指数成長する。
\end{corollary}
\begin{proof}
十分性は命題\ref{prop:elliptic-pencil-polygons}である。
相補的な空でない二つの連続部分鎖$I,J$を取る。$I$だけを台に持つ$L$の切断は境界二点で消え、
次元$d(I)-1$を持つ。実際、正次数の各成分では両端値を独立に指定でき、内部節点で貼り合わせると
鎖上の切断次元は$d(I)+1$となり、両端値の零条件が二つである。
同様に$J$に支えを持つ$M$の切断の次元は$e(J)-1$である。
この二空間の積は順逆とも零である。従って元代数の矢をそれぞれの双対へ射影すると、
積の核の直交空間である全二次関係の像が消える。
この経路代数への全射により、長さ$2m$の次元は少なくとも
$((d(I)-1)(e(J)-1))^m$となる。

二成分で$d\ne e$なら、成分の順序と両束の役割を選んで一方の束は一成分に次数3、
もう一方は残る成分に次数2以上を持つ。三成分で異なる場合は、$J$を$M$の次数2の成分、
$I$を残り二成分とすれば$d(I)=3$である。四成分では正の多重次数は一種類しかない。
従って異なる全場合で指数成長し、標準Hilbert級数に反する。
\end{proof}

% !TeX root = ../papers/nc-p3.tex
\section{一つの多重点を持つ被約四次曲線の境界}
\label{sec:additive-pencil-boundary}
二つの接する枝、平面通常三重点、四枝の楕円多重点を同じ計算で扱う。
$C$の正規化を$m=2,3,4$本の射影直線とし、各枝の原点の局所座標を$z_i$とする。
それらの点を一つに合わせた局所環は、値が一致し、一次係数の和が零となる
正規化上の関数全体とする。枝のパラメータを変えれば二枝の場合は接する二枝の一次係数の一致となる。
直線束の貼り合わせでは、適切な自明化で、共通値$c$と一次係数$v_i$への条件は
$\sum v_i=\lambda c$となる。各枝での定数倍を除けば、この一つの加法パラメータが残る。
二つの束の積ではパラメータは加算される。

\begin{theorem}[正則性と被約境界の被覆]\label{thm:additive-degree-balance}
各成分上で正の次数を持ち総次数4の可逆層$L,M$と、$LM$の基点自由鉛筆を用いる。
定理\ref{thm:elliptic-pencil-all-as}の反転二窓が標準Koszul AS正則となる必要十分条件は
多重次数が一致することである。一致して$L\not\simeq M$なら反対階数は8で、
両側の二次核はそれぞれの完全線型系が定める$C$を切り出す。
添字保持同型の必要十分条件は、曲線同型が順序付きの二束と鉛筆を対応させることである。
このとき固定した多重次数での有理商の次元は$13-m$である。
\end{theorem}
\begin{proof}
導手は各枝の二重原点である。双対化層は自明である。実際各枝の$dz_i/z_i^2$を組にすると、
正則関数を掛けた留数和が零となる。逆に任意の双対化微分は導手条件で高々二重極となり、
各枝の二重極係数は等しく、単純極係数の和は零である。この組は上の微分の正則関数倍で尽くされる。
非負多重次数で総次数が正なら逆束の切断は零なので$H^1=0$である。
総次数が2以上なら元の束は大域生成される。共通値を非零に指定し、
正次数の枝の一次係数を調整して特異点で生成できる。
次数2以上の枝の他の点では導手で消える多項式で生成できる。次数1の枝ではもう一つの正次数の枝があれば
一次係数を互いに調整できる。正次数の枝が一つだけならその次数は2以上である。
次数零の枝は共通値で生成される。従って多重次数が一致する場合、
命題\ref{prop:elliptic-pencil-polygons}と同じ評価核のフィルトレーション、中心正則列、完全対合の
証明により全次数でAS正則性が従う。異なる場合は任意の真部分集合$I$上だけに支えを持つ
$L$の切断の次元が$d(I)-1$である。各枝で定数項を零にし、一次係数の和を零にするためである。
系\ref{cor:elliptic-polygon-degree-balance}と同じ相補的部分集合で指数成長の商を得る。

$|L|$は各枝を有理正規曲線に写す。線型包の座標は共通値、和に一式を課した一次係数、
枝ごとに独立な二次以上の係数である。したがって枝は原点以外で交わらず、局所環も復元される。
詳しくは、一次係数の超平面は埋込みの座標で実現する。二次以上を一枝だけで実現するには、
次数2以上の枝には二次の座標がある。$m\ge3$では指定枝と異なる一枝のみに
一次係数を持つ二関数を、異なる相手枝を選んで掛ければ指定枝だけの二次初項を得る。
$m=2$で一方が直線でも、他方の二次座標を共通接線座標の平方から引けばよい。
その後は一次座標を掛けて全ての高次初項を得る。完備化で導手全体を復元できるため閉埋込みとなる。
なお二枝の線型包は共通接線を共有するが、共通接線は非線型な有理正規曲線と二重原点でのみ交わる。
三枝以上の二つの線型包の交わりは共通点だけである。
この座標表示で平面に含まれる次数3の部分曲線はない。例えば三枝では二次曲線の内部座標と
独立な二つの接方向があり、四枝なら任意の三接方向は独立である。
従って前節と同じ次数の議論で二二次式完全交叉となる。
多重次数零の非自明な束では定数の組が一次貼り合わせを満たさず切断を持たない。
これで階数8と二次核が従う。完全交叉を平滑化して二束と鉛筆を延ばす議論も前節の通りである。

グラフの復元を詳しく示す。$r_i=2d_i$とし、鉛筆の一般の切断の零点を取ると、
前節と同様にテンソルは各枝の$r_i$個の異なる非特異点での非零重み付き評価の和である。
一枝だけに支えを持つ$L^2$の切断は、その枝で二重原点を引いた全切断である。
そのテンソル縮約の階数は$r_i-1$、像は対応する枝の評価の包に入る。
$d_i\ge2$ではこの階数は異なる枝の線型包の共通部分の次元より大きい。
$m\ge3,d_i=1$では像の生成元は二次Veronese上の異なる二点評価の非零重み付き和で、
他枝との共通点である原点評価の倍数にはならない。二枝で小次数の枝があれば、
大次数の枝の対応からその相手も定まる。

各正規化の原点を無限遠に置き、一方の候補グラフで両側のアフィン座標を$t$に揃える。
他の候補の枝写像を$t\mapsto at+b_i$と書く。傾き$a$は全枝で共通である。
これは二つの同定の比が、関数の一次係数の和零の条件を保つからである。
テンソルの$i$成分の表示は、$\deg u,\deg v\le r_i$に対して一方は
$\ell_i(u(t)v(t))$、他方は$\ell'_i(u(t)v(at+b_i))$となる（自明化の定数は汎関数に含める）。
$d_i\ge2$なら、片側の成分支持に制限した階数$r_i-1$のテンソルは、両候補の
$(r_i-2,r_i)$グラフの線型包に入る。異なる二グラフは長さ2で交わり、
合併のイデアルでひねると$\mathcal O(r_i-4,r_i-2)$で$H^1=0$になる。
よって二つの包の交わりの行列階数は高々2であり、その枝で同定は一致する。

$r_i=2$なら、源または先で次数0の多項式を使うとその切断は当該枝だけに支えを持つ。
相手側はこの枝への制限が全射なので、$m_j=\ell_i(t^j)$と置けば
$\ell_i'(t^j)=m_j$ ($0\le j\le2$)、
\[
 (a-1)m_1+b_i m_0=0,\qquad
 (a^2-1)m_2+2ab_i m_1+b_i^2 m_0=0.
\]
相異なる二点の重みは非零だから$m_0m_2-m_1^2\ne0$である。
大次数の枝があれば$a=1$となり、上の二式から$b_i=0$となる。

残るのは四直線の場合である。テンソルの二つの表示の差を枝ごとに取ると、
源か先の次数0の座標との組では差が零であり、各差は両側の原点の一次ジェットの包に入る。
全枝の和は零である。原点評価の直線で両側を割ると、残るのは接方向のテンソル
$h_i\otimes k_i$の係数だけである。三次元の四接方向には全係数非零の一つの関係だけがあり、
対応する四つの$h_i\otimes k_i$は独立である。実際、最初の三方向を両側の基底とすると、
第四項だけが基底の非対角積を持つ。従って各差で一次対一次の係数も零であり、
$m_2=a m_2+b_i m_1$を得る。最初の一次式と判別式の非消滅から$a=1,b_i=0$である。
全枝で一致すれば被約な曲線上の射は同じであり、鉛筆は隣接核から復元できる。

各多重次数での束の選択は一次元の加法群、鉛筆は12次元である。
曲線自己同型の連結部分は無限遠を保つ枝ごとのアフィン変換で、傾きが共通となるため次元$m+1$である。
一般の鉛筆の安定化群は有限である。各枝での二切断は一般に選べるので、次数4以上の枝があれば
その分岐値が三つ以上あり、基底の自己同型が有限に限られる。四直線のときも
異なる枝の二次射の分岐値の和集合は一般に三点以上となる。実際に二枝への切断の制限は
共通値を合わせれば一次係数の差を残る枝で調整でき、二つの二次射は無限遠での値だけを
共有して一般に選べる。基底変換を固定した後は各有限射の持ち上げは有限個である。
従って差は$14-(m+1)=13-m$となる。
\end{proof}

\begin{corollary}[被約二二次式完全交叉の曲線型]\label{cor:reduced-quartic-conductor}
$\mathbb P^3$内の被約二二次式完全交叉を曲線入力とし、各成分上で二束が正次数である鉛筆構成では、
上記と前節の曲線型で尽くされる。
\end{corollary}
\begin{proof}
$C$は連結で算術種数1、双対化層は自明である。正規化の成分の種数、成分と特異点の接続グラフの
第一Betti数、各特異点の$\delta_p-b_p+1$（$b_p$は枝数）の和は1である。
最後の数は、共通値だけを課す正規化の部分環からの余長であり非負である。
種数1の成分は導手因子の次数が$2-2g=0$なので他枝と交われず、滑らかな場合に尽きる。
ここで導手因子の次数は有限正規化の双対性
$\nu_*\omega_{\widetilde C}=\mathcal Hom_C(\nu_*\mathcal O_{\widetilde C},\omega_C)$から従う。

局所寄与が全て零なら、特異点の完備環は共通値だけを課す環である。
枝数$b$でその導手の正規化側の余長は$b$、$\delta=b-1$である。
Gorenstein曲線で導手のこの余長は$2\delta$である
（正規化完全列の双対を取り、長さ$\delta$の商の局所双対も同じ長さとなる）。
従って特異点は節点だけで、各有理成分上の導手の次数2から環状配置になる。自己節点の一成分も含む。
残りの場合は一つの特異点の局所寄与が1で、接続グラフは木となる。他の節点は分離節点になるが、
その一方の部分曲線は算術種数0であり、そこへの双対化層の制限次数は$-2+1=-1$で矛盾する。
従って他の特異点はなく、各有理成分は一枝だけでこの点と交わる。
各枝の導手次数は2なので、環は共通値に加えて一次係数の超平面を指定する。
この超平面の関係の全係数が非零でなければ、一枝の導手が一次となってしまう。
各枝の座標を変えて上の和零の表示を得る。一枝なら尖点であり、枝数は次数4以下である。
\end{proof}

% !TeX root = ../papers/nc-p3.tex
\section{二重二次曲線の鉛筆と滑らかな反対核からの逆判定}
\label{sec:ribbon-graph-recognition}
\begin{theorem}[二重二次曲線上の全次数認証と同型]\label{thm:conic-ribbon-as}
二二次式完全交叉$C$の被約台が滑らかな二次曲線$P\simeq\mathbb P^1$で、一般点での重複度を2とする。
$P$へ次数2で制限される二つの可逆層$L,M$と、$N=LM$の基点自由鉛筆$W$を取ると、
反転二窓は標準Koszul AS正則となる。$L\not\simeq M$なら反対階数8で、
$(C,L,M,W)$の曲線同型が添字保持同型の必要十分条件である。この層の有理商は9次元である。
$L\simeq M$なら対称完全交叉チャートに入る。
\end{theorem}
\begin{proof}
完全交叉はCohen--Macaulayなので点に支えを持つイデアルを含まない。
冪零根$J$は一般点で平方が零だから全体で$J^2=0$、$P$上で階数1無捩れである。
$\chi(\mathcal O_C)=0$から$J\simeq\mathcal O_P(-2)$となる。
滑らかなアフィン座標の持ち上げを局所的に選ぶと、貼り合わせの差は$T_P\otimes J$の導分である。
$H^1(T_P\otimes J)=H^1(\mathcal O_P)=0$なので構造層は分裂し、
$\mathcal O_C=\mathcal O_P\oplus J$と取れる。単元の完全列から、固定した制限次数での
可逆層の選択は$H^1(J)=k$の捩れとなる。非自明な制限次数零の束は、
$1$の持ち上げを妨げるこの類が非零なので大域切断を持たない。

二束の交互切断環のKoszul性を示す。束$G$の$P$への制限$G_0$が
$\mathcal O(1),\mathcal O(2)$の直和なら、$0\to G_0(-2)\to G\to G_0\to0$から
$H^1(G)=0$で$H^0(G)\to H^0(G_0)$は全射である。従って中山の補題で$G$は大域生成される。
評価核の制限は、$H^0(G_0(-2))$の次元個の自明束と
$\deg G_0$個の$\mathcal O(-1)$の直和になる。実際、制限した評価は$G_0$の評価を経由し、
各$\mathcal O(j)$の評価核は$\mathcal O(-1)^j$である。
次の$L_i$を掛ければ再び要求した制限型となる。尾部でさらに正の次数を掛けても$H^1=0$であり、
定理\ref{thm:elliptic-pencil-all-as}の線型分解がそのまま成立する。
鉛筆の一般切断は$P$の一般点で非零で正則元となる。基点自由性で第二切断は第一切断の
零因子上で単元となるので、二元の正則性とKoszul性の降下も通る。双対化層は自明であり、
同じSerre対合で有限Frobenius性を得る。

$|L|$の四つの切断として$P$上の二次単項式の持ち上げ$x_0,x_1,x_2$と
$H^0(JL)=H^0(\mathcal O_P)$の生成元$n$を取る。像の二次式は$n^2$と
$x_0x_2-x_1^2-n\ell(x)$である。後者は$H^0(JL^2)$を$n x_i$が張るため存在する。
被約台を埋め込み、$n$が各局所自明化で冪零根を生成するため閉埋込みになり、
二式の完全交叉は同じHilbert多項式でこの像に等しい。$M$側も同じである。
従って鉛筆対合で反対階数8および両側の二次核が従う。
同じ束のときは二定義式と二鉛筆切断の持ち上げが四二次式の正則列になる。

二つの埋込み二重曲線の間の同定を復元する。一般の$a\in W$の被約制限は単純な四零点を持つ。
その零因子上の環は四つの双対数環の積である。$\ell$はこの環の単元倍を除く留数対合から得られ、
各局所socleで非零である。実際、第二切断$b$はその上で単元で、随伴により$a$のCartier因子の
双対化加群は$N$の制限である。$1/b$で自明化した留数は完全対合を与える。
片側を$H^0(JL^2)=H^0(\mathcal O_P(2))$に制限すると他側は$H^0(\mathcal O_P(4))$を経由し、
四点での非零重み付き評価により対合の階数は3である。異なる二つの被約グラフの
$(2,4)$での共通線型包の階数は高々2なので、被約同定は一意である。

被約同定を固定する。第二の候補との差は、$J$でのスカラー$c$と、
$\mathcal O_P$から$J$への導分$v\in H^0(T_P\otimes J)=k$を持つ自己同型である。
二候補の乗法汎関数を比較する際は、値束を被約制限上で同一視する。
片側だけが冪零な積は$H^0(\mathcal O_P(6))$を張り、その汎関数は上の階数3により非零である。
源だけが冪零な積から両候補はこの同じ汎関数を持ち、先だけが冪零な積から$c=1$となる。
さらに被約切断の対の和$\sum p_i\otimes q_i$で$\sum p_iq_i=0$のものを取る。
持ち上げた積は冪零となるので、片側だけが冪零な積を加えて第一グラフ上の積を零にできる。
第二グラフ上には$v\sum p_i\,dq_i$だけが残る。局所枠の変更から来る零階項は
$\sum p_i q_i$を因子に持つため消える。$\mathcal O(4)$の切断の交代対だけでも
$\sum p_i\,dq_i$は$H^0(\Omega^1_P(8))=H^0(\mathcal O_P(6))$全体を張る。
アフィン座標で$(j-i)t^{i+j-1}$ ($0\le i<j\le4$)を取ればよい。
$v\ne0$なら先ほどの非零汎関数がこの全空間を消して矛盾する。よって曲線の同定は冪零方向まで一意で、
鉛筆も隣接核から戻る。

自己同型群は$\PGL_2$の持ち上げ、導分、冪零根のスカラーからなる五次元群である
（$J=\Omega^1_P$の自然作用で被約自己同型を持ち上げられる）。
二つの束のパラメータと鉛筆の次元は$1+1+12$である。一般の鉛筆の被約次数4の射は三つ以上の
分岐値を持つので、その被約安定化群は有限である。被約恒等の作用は次数$s$の束の
貼り合わせパラメータに$\lambda\mapsto c\lambda+s v$と作用する
（導分を$J=\Omega^1_P$への$d$で規格化し、二座標の遷移$t^s$の対数微分を取る）。
異なる二つの次数2のパラメータを固定すれば$c=1,v=0$となる。従って商次元は9である。
\end{proof}
\path{computations/verify_ribbon_trace.py}では二座標の遷移に一次項を含む束を使い、
異なる二束と等しい二束で両切断階数、および七次元の微分像を有理数上で検査できる。

\begin{theorem}[滑らかな反対核とグラフによる逆判定]\label{thm:smooth-apolar-graph-test}
$w\in\operatorname{Sym}^2 E\otimes\operatorname{Sym}^2 F$の隣接切断階数が6、反対階数が8とする。
両反対核の二次式が滑らかな四次完全交叉$X,Y$を切り出すとする。
$w$が楕円鉛筆の表示を持つ必要十分条件は、同定$\phi:X\to Y$のグラフへの
$(2,2)$切断の制限を通る汎関数であることである。この条件が成立すれば鉛筆の基点自由性も自動的で、
元の反転関係列は全次数で標準Koszul AS正則である。
各原点付き同定を一つ固定して平行移動を動かすとき、判定はパラメータ曲線上の次数64の
行列式線束の切断で行える。解があれば一つの被約点である。
\end{theorem}
\begin{proof}
$X,Y$上の二次切断を$A=\mathcal O_X(2),B=\mathcal O_Y(2)$と書く。
両核を割った$w$は$H^0(A)\otimes H^0(B)$上の階数8の汎関数である。
グラフ条件を仮定し、$X$上で$L=\mathcal O_X(1), M=\phi^*\mathcal O_Y(1), N=LM$とする。
$L=M$なら両者を同定したとき四切断の積は全対称で、反対と隣接の階数が等しくなるので$L\ne M$である。
$H^0(L)\otimes H^0(M)\to H^0(N)$は全射である。一般に次数4以上の二束の積の全射性は、
定理\ref{thm:elliptic-pencil-all-as}の評価核のフィルトレーションに他方の束を掛けた$H^1$消滅から従う。
グラフ上の汎関数を$\ell\in H^0(N^2)^*$とすれば、その$H^0(N)$上の積対合の核$W$は二次元である。
もし固定因子$D>0$があれば、$W$からその因子を除いた基点自由鉛筆は次数$8-\deg D\ge2$を持つ。
鉛筆の完全列で$W\otimes H^0(N)\to H^0(N^2(-D))$の核は$H^0(\mathcal O_X(D))$である。
次元を数えると積は全射となるので、$\ell$は長さ$\deg D\le6$の$D$への制限を通る。
反対対合の階数8と矛盾する。従って$W$は基点自由で、$\ell$は鉛筆の一意なtraceのスカラー倍である。
先の全次数定理を適用すればよい。逆と同定の一意性は定理\ref{thm:elliptic-pencil-reconstruction}による。

原点付きの同定を用い、$p_1,p_2:X\times X\to X$を二射影、$m$を加法とする。
$\mathcal B=p_{2*}(p_1^*A\otimes m^*B)$は階数16の束で、高次順像は零である。
各平行移動グラフへの積の制限は全射なので、$\mathcal B^*$は自明な64次元束の部分束である。
面のRiemann--Rochより$\deg\det\mathcal B=\chi(p_1^*A\otimes m^*B)=64$である。
実際、二因子の自己交叉は零で、$p_1,m$のファイバーは一点で交わる。
十六次外積に$w$を外積した全成分の同時消滅が求める条件である。いずれも$\det\mathcal B$の切断だから、
一つの恒等的に非零の成分の高々64個の零点で残りを検査すれば足りる。
異なる平行移動の線型包は前節の合併への全射により交わらない（射影線型包の意味）。
さらにグラフ$G$の法方向の一次微分は、その上で消える切断から
$H^0(G,(A\boxtimes B)|_G)$全体へ全射である。法束は自明であり、
$A\boxtimes B$を$-2G$でひねった束の$H^1$は前節と同じ豊富性で零だからである。
平行移動方向はこの法束を生成するので、非零の$w$に対して一次微分が全て消えることはない。
従って候補パラメータの閉スキームも被約である。
\end{proof}

\begin{corollary}[滑らかな反対核を保つ特殊化]\label{cor:smooth-apolar-specialization}
滑らかな楕円鉛筆から得た射影Gramテンソルの閉包内で、隣接階数6、反対階数8、
両反対核の滑らかな完全交叉性が保たれる任意の点は、実際に基点自由楕円鉛筆から得られる。
\end{corollary}
\begin{proof}
稠密な鉛筆表示から指定点への特殊化を曲線で取り、正規化し必要なら離散付値環を有限拡大する。
一定階数の反対核と二式の横断性から、両方の曲線はその上の滑らかな固有族となる。
一般ファイバーの同定は特殊ファイバーまで同定として延びる。
実際、滑らかな曲線族の正則な全空間の間の有理写像の不定点を点のブローアップで解消すると、
例外有理曲線は滑らかな種数1の特殊ファイバーの中で点に写る。
従って元の正規な面へ降下する（グラフの射影が一点ファイバーを持つ固有双有理射となる）。
逆写像も同じである。延長したグラフへの二次切断の積は両ファイバーで全射で、核は部分束である。
その核で$w$が消える条件が特殊点にも残るので、前定理を適用できる。
\end{proof}

\begin{theorem}[一つの単純零点からの楕円被覆認識]\label{thm:simple-zero-smooth-elliptic}
定理\ref{thm:smooth-apolar-graph-test}の切断階数と滑らかさのもとで、第一窓の六式に一つでも
単純零点対（六次元接空間でJacobi階数6）があれば、グラフ条件が成立する。
従ってこの場合は所属成分をあらかじめ仮定せずに全次数のAS性と楕円データの復元が得られる。
\end{theorem}
\begin{proof}
単純零点では六次元Gram空間$H$への行・列評価の像が相補三平面であり、
二重交代化の消滅により直交するので両方とも非退化である。
従って第\ref{sec:marked-opposite-rank}節の無跡対称三行列$A,B$による表示を取れる。
三列$A_i$が従属なら反対切断の最初の四行の間に関係ができ、反対核が
指定文字を因子とする可約二次式を含んでしまう。非退化な滑らかな楕円四次曲線を含む
可約二次曲面は存在しない。$B_i$も同様に独立である。

交換子行列$D_A$の階数が1以下の場合を除く。三次元の無跡対称行列空間系の中で
可換な独立二平面は、その中の正則行列（最小多項式の次数3）の無跡中心化代数に一致する。
実際、非正則な無跡対称行列の射影軌跡は$vv^{\mathsf t}$のスカラー行列を法とした像である
（ある固有値の固有空間が二次元以上）。これは二次Veronese曲面を階数3の点から射影した
埋込みであり、直線を含まない。正則行列の中心化代数はその多項式なので主張を得る。
階数0は三次元可換系を要するので不可能である。
階数1では三次元の系内に異なる二つの可換二平面があり、その共通直線は非正則行列を与える。
これは三次元で全ての二ベクトルが分解可能であることと上の中心化代数による。
共通直線の行列を$vv^{\mathsf t}$からスカラーを引いて表すと、三次元の系の全てが$v$を固有ベクトルに持つ。
従って$I,A_i$の直交補は$vu$ ($u\perp v$)という二次式の二次元空間である。
反対核の二式を指定文字の座標が零の平面に制限すると、この直交補に入る
（第\ref{sec:marked-opposite-rank}節の十行行列の最初の四行）。
従ってその平面内の直線$v=0$全体が楕円四次曲線に入ることになり矛盾する。

残る階数2のデータは可逆交換子を持つデータの閉包に属することを示す。
独立な無跡三列が張る三次元空間を固定せずに動かす。可換二平面の族は既約三次元である。
正則無跡対称行列の射影四次元開集合からその中心化二平面を取る写像が全射で、
各ファイバーは射影直線の開部分だからである。これを含む三次元空間の選択は射影二次元であり、
順序付き基底の選択は九次元である。階数2では可換二平面がただ一つとなるので、
この階数2の三列の族は既約14次元である。空でないことは対角な独立二行列と一般の第三行列で確認できる。

$A\mapsto D_A$の階数2の値には、左からの$\mathrm{SO}_3$と右からの$\mathrm{GL}_3$
（行列の同時直交変換と列変換）が作用する。右作用は余因子行列を通じるが全射である。
したがって値の軌道は像平面の法線が非等方か等方かで尽くされ、それぞれ次元8と7である。
非等方法線が実現することは、実有理数上で対角二行列と一般の実対称第三行列を取れば従う。
よって非等方軌道上のファイバー次元は6である。等方側は既約14次元の真閉部分に入るため、
その軌道が現れる場合にもファイバー次元は高々6である。
$D_B=-D_A^{\mathsf t}$により、階数2の対$(A,B)$の族の次元は高々$8+6+6=20$である。
一方、独立三列の対を動かす30次元空間ではこの等式は九式であり、各成分の次元は21以上である。
階数2以上の開部分で比較すれば、階数2だけに含まれる成分はない。

可逆交換子の全表示は命題\ref{prop:marked-net-parameters}のパラメータで覆われる。
二符号は基底変換で移り合い、射影基底を入れた像の閉包は既約である。
実際、有理な基底空間上の平方根被覆が分裂しても、二枚の像は連結な基底変換群で飽和すれば同じになる。
その一般点は楕円鉛筆の一般点である。これは同命題の支配性と
定理\ref{thm:elliptic-pencil-reconstruction}の稠密開被覆による（一般の第一窓からはGram直線が一意）。
従ってこの証明で得た階数2の極限も、基底付き射影テンソルとして楕円表示の閉包に入る。
系\ref{cor:smooth-apolar-specialization}により結論が従う。
\end{proof}

\begin{proposition}[片側の小行列式による単純零点の検査]\label{prop:smooth-apolar-defect}
上の階数と両反対核の滑らかさを仮定し、$\beta:E^*\times F^*\to H$を六次元Gram表示とする。
$B_x=\beta(x,-)$、$Q_x=B_x^{\mathsf t}h B_x$（$h$はGram計量）と書く。
第一窓の零点スキームは$\mathbb P(E^*)$内の$B_x$の四次小行列式の零点$Z$に同型である。
その各閉点で、唯一の零点対の接空間次元は$3-\operatorname{rank}Q_x\le2$である。
従って$Z$上の一点でも$Q_x$の三次小行列式が非零なら、楕円鉛筆表示とAS正則性が従う。
$Z$が有限なら、単純点の個数を$a$、接空間次元1,2の点数を$b,c$として
$a+2b+4c\le20$である。特に台に11点以上ある有限窓はこの判定で認証される。
\end{proposition}
\begin{proof}
$x\ne0$を固定する。反対核に可約二次式がないので、$x E^*$の四次元空間は
反対切断で単射に写る。その像を$F^*$上の二次形式と見ると、全てが$\ker B_x$を根基に含む。
実際、$f\in\ker B_x$なら$w(xe,ft)=\langle\beta(x,f),\beta(e,t)\rangle=0$である。
従って$\binom{\operatorname{rank}B_x+1}{2}\ge4$で階数は3以上となる。両側で同様である。
$Z$では階数3の小行列を可逆にすると、局所的な行列の行・列消去で単位三行と
残る一列三行となる。後者の三成分が四次小行列式のイデアルを生成し、
零点対のもう一方の射影座標は核直線で決まる。これはスキーム上での主張である。

零点対$(x,y)$では、三平面$U=\operatorname{im}B_x$と$V=\beta(-,y)(E^*)$は
Gram恒等式と$\beta(x,y)=0$により直交し、$V=U^\perp$である。
Jacobi像は$U+V$なので、接空間次元は$\dim(U\cap U^\perp)=3-\operatorname{rank}Q_x$となる。
反対核には$x^2$も入らず、$Q_x=w(x^2,-)$は非零である。
従って単純点の検出には$Z$を$Q_x$の三次小行列式の非消滅集合に制限すればよい。
同値な係数検査として、$B_x$の最大小行列式の斉次イデアルの根基に
$Q_x$の三次小行列式が全て属するかどうかを調べられる。

有限の場合は六つの$(1,1)$因子の固有交叉で、全長は$\binom63=20$である。
接空間次元が$k$の点では完備局所環で独立な線型部を消去すると、
$k$変数の正則局所環の極大イデアルの平方内の$k$元による零次元完全交叉となる。
$k=1$なら長さは2以上、$k=2$なら次数0,1で$1+2$次元、次数2にも
三つの単項式に対し関係が高々二本なので一次元以上残る。これで不等式を得る。
\end{proof}

% !TeX root = ../papers/nc-p3.tex
\section{非被約な四次完全交叉上の全鉛筆}
\label{sec:nonreduced-quartic-pencils}
\begin{theorem}\label{thm:nonreduced-quartic-pencils}
$C\subset\mathbb P^3$を非被約な二二次式完全交叉、$L,M$を各既約成分で正次数、
全次数4の可逆層、$W\subset H^0(C,LM)$を基点自由鉛筆とする。
反転二窓は全て標準Koszul AS正則である。$L=M$（同型を略記）では対称完全交叉となり、
$L\ne M$では反対階数8で、添字保持同型は$(C,L,M,W)$の順序付き曲線同型と同値である。
特に後者では反対核から被約台とその冪零構造も決まる。
\end{theorem}

\begin{proof}
まず埋込みを保つ乗法群で$H^1(\mathcal O_C)$に非自明に作用するものを作る。
二式は共通因子を持たず、一次元Cohen--Macaulay性により重複する既約成分が存在する。
この成分は次数から直線か滑らかな平面二次曲線である。後者のイデアルを$(h,f)$と書くと、
二式のうち一方は$h$の倍数にできる。二次曲線の一般点で交叉が単純にならないため
その倍数は$h^2$である（他方は$f$を非零係数で含む）。
他方の式の$h$一次項は非退化三変数二次式の変数の平行移動で除ける。

重複直線のイデアルを$(x,y)$とする。二式の$(x,y)^2$を法とした係数の二行二列行列は、
直線上で恒等的に行列式零である。階数1行列の直線はSegre曲面の母線なので、
共通左核または共通右核を持つ（射影像が一点以下の場合も含む）。
左核があれば鉛筆内に$x,y$の二次式があり、右核があれば線型部分を$x\ell_i(z,t)$にできる。
$\ell_i$が従属なら再び二変数の式を得る。独立なら$z,t$に$x,y$の倍数を加えて
$xz+a y^2,xt+b y^2$となり、$y^2$を消す線型結合が二平面を与える。
従っていずれも鉛筆に重平面または二平面がある。

重平面$h^2$を取ると、残る式の平面への制限の階数に応じて、滑らかな二次曲線のほかは
$(h^2,xy+\kappa hz)$または$(h^2,x^2+\kappa hz)$である
（$\kappa=0,1$、不要な変数を補って四変数とする）。
相異なる二平面$xy$を取る場合、その交線上に残る式が恒等零でなければ、
非被約性から片方の平面切断は二重直線である。そこでその制限を$z^2$として、
残りを$x$と一次式の積で書く。$z$に$x$の倍数を加え、$xy$を除くと、
第二式は$z^2+xt,z^2+x^2,z^2$のいずれかとなる。
交線全体が含まれれば、第二式は$x(\ell+a x)+y(m+b y)$と書ける。
交線上の$\ell,m$が独立なら$xz+yt$、非零比例なら$z(x+y)$にできる。
一方だけ非零なら共通因子を避けるため$xz+y^2$にできる。
両方零なら二変数正則列であり、二変数の平方の円錐曲線と鉛筆の二交点から
$(x^2,y^2)$にできる（接線なら共通一次因子が生じる）。これらの表示の重複は許す。

二式の重みを$e_1,e_2$、座標の重みを$w_i$とすると、随伴公式から
$H^1(\mathcal O_C)$の指標は符号を除いて$\sum w_i-e_1-e_2$である。
未使用変数があればそれだけに重み1を置く。滑らかな二次曲線の重複では$h$だけに1を置く。
残る$(h^2,xy+hz)$では$(h,x,y,z)$に$(0,1,1,2)$を、
$(xy,z^2+xt),(xy,xz+yt)$では$(x,y,z,t)$にそれぞれ$(0,0,1,2),(0,0,1,1)$を置く。
全て非零指標となる。

この表示から被約台は滑らかな有理曲線、二成分が一点で横断する木、三直線の鎖、
または三本の独立な方向を持つ同一点の三直線である。三直線の最後の場合も
正規化からの降下条件は値の一致だけで、全ての場合に被約台の$\operatorname{Pic}^0$は零である。
各成分の次数に一般点の長さを掛けて足すと4なので、正の多重次数は埋込み束
$L_0=\mathcal O_C(1)$のものに一意に決まる。実際、重複二次曲線は次数2、
四重直線は次数1、二重直線と単純二次曲線なら次数$(1,2)$で、
他の複数直線の場合は重複度の和が4なので全て次数1である。
単元列と冪零イデアルの対数から、多重次数零の可逆層全体は加法群である。
その接空間は$H^1(\mathcal O_C)=k$なので$\operatorname{Pic}^0(C)=\mathbb G_a$となる。

重複成分$P\simeq\mathbb P^1$に対し
$I=\mathcal Hom_C(\mathcal O_P,\mathcal O_C)\subset\mathcal O_C$を取る。
有限閉埋込みの双対性と$\omega_C=\mathcal O_C$により$I=\omega_P$で、
$H^1(I)\to H^1(\mathcal O_C)$は定数の制限の双対だから同型である。
$I$は一般点でArtin局所環のsocleにあり、$P$への像は零である。
従って$I^2=0$である。$\mathcal O_C\to\mathcal O_P\xrightarrow{d}\omega_P=I$の合成は
平方零の導分$\delta$で、$g_s=1+s\delta$は自己同型を与える。
$1+I$からの写像$H^1(I)\to\operatorname{Pic}^0(C)$は微分が同型の加法群の準同型であり、
標数零だから同型である。遷移関数$f$の変化は$1+s\,d\log(f|_P)$なので、
$g_s^*$は$P$上の次数$d$の束を、そのPicard座標で$sd$だけ移す（座標の定数倍を許す）。
従って任意の正または負の冪$L_0^n$の多重次数の束は全てその自己同型像である。
また$H^0(\mathcal O_C/I)=k$は上の$H^1$の同型から従う。多重次数零の非自明な束は
非零の$1+I$コサイクルで貼られ、商上の定数1は持ち上がらず、$H^0(I)=0$なので切断は零である。

第一束を$L=L_0$に揃え、第二束を$M_s=g_s^*L$とする。
上で構成した乗法群は$L$を固定し、$\operatorname{Pic}^0(C)$に非零指標で作用するので、
$s\ne0$の順序付き束の対は一軌道である。
交互切断環（偶奇二頂点に折り畳む）の各長さ$n>0$の束は、
上の平行移動により$L^n$と同じコホモロジーを持つ。
埋込み完全交叉の分解から$h^0=4n,h^1=0$で、負の冪には切断がない。
従って切断空間はパラメータ直線上で束となり基底変換と可換である。
$s=0$の切断環は二次完全交叉の斉次座標環を二頂点に引き戻したものでKoszulである
（二次正則列で多項式環を順に割る）。内部次数を一つ固定した相対正規化bar複体は
有限長のベクトル束の複体である。非対角Torの消滅はこの次数ごとに原点の開近傍へ移る。
各近傍が非零軌道に会うので、各次数で全パラメータ上の消滅を得る。
特に任意のこの多重次数の二束の切断の積は全射である。

交互切断環から$W$の二つの中心二次元で割る。一般の$a\in W$は全既約成分の一般点で非零、
ゆえに正則元である。第二の生成元は$a$の零因子上で単元である。
従って切断の等式$bu=av$から$u$は$a$で割れ、負次数部分の切断も零なので、
二元は正則列になる。定理\ref{thm:elliptic-pencil-all-as}の鉛筆完全列とSerre双対性により、
商はHilbert列$1,4,6,4,1$のKoszul Frobenius代数で、その二次双対が求める標準AS正則代数である。
同じ束なら二つの定義式と鉛筆が四二次式の完全交叉を与える。
異なる束なら$L/M$とその逆の切断が零で、同じ鉛筆列は$H^0(L^2)$と$H^0(M^2)$に
完全対合を与える。両束は$L_0$の自己同型像だから二次核は両埋込みの定義式そのものである。

最後に両埋込みの間の同定をスキーム上で一意に戻す。
階数8の両反対核には同時に四つの基底を二つの射影基底に還元する既出の議論が適用できる
（定理\ref{thm:elliptic-pencil-reconstruction}の証明）。
第一の候補グラフで両曲線を$C$と同定し、第二候補を自己同型$g$のグラフとする。
$s,t\in H^0(L)$の行列式$s(x)t(y)-t(x)s(y)$に、
任意の$H^0(L)\boxtimes H^0(M^2/L)$の切断を掛ける。積は両側の二次切断のテンソルで、
第一グラフで零だから与えられたテンソル汎関数で消える。
第二グラフでは$D=L\,g^*L$, $N'=L\,g^*M$として、
乗じた切断の積が$H^0(N'^2/D)$を張る。実際、二因子$L,g^*(M^2/L)$の
多重次数はともに一意な正の型なので先ほどの積の全射性が使える。
第二候補の鉛筆traceにおいて$H^0(D)$と$H^0(N'^2/D)$は完全に対合する。
これは$D/N'=g^*(L/M)$とその逆の切断が零だからである。
よって$s\,g^*t-t\,g^*s=0$が切断として成り立つ。$|L|$は閉埋込みなので、
切断の比を局所座標に使えば$g$は冪零方向を含め恒等である。
鉛筆は隣接核を積で写した像として戻り、逆に順序付きデータの同型は全ての積を保つ。
\end{proof}

% !TeX root = ../papers/nc-p3.tex
\section{非横断的な印付き零点からの楕円認識}
\label{sec:marked-gram-smoothing}
階数6の隣接Gram表示$\beta$とその非退化計量を用いる。
零点対を両基底の先頭に取り、$u_j=\beta_{0j}$、$v_i=\beta_{i0}$、
$T_{ij}=\beta_{ij}$ ($1\le i,j\le3$)と書く。
ここで$\beta_{00}=0$、$U=\langle u_j\rangle$、$V=\langle v_i\rangle$である。

\begin{lemma}[単純な印の三行列系の閉包]\label{lem:independent-marked-closure}
$U,V$が相補的で、両側で指定文字の一次倍四次元上に反対切断が単射なら、
その射影Gramテンソルは楕円鉛筆の閉包に属する。
\end{lemma}
\begin{proof}
直交分解後の第\ref{sec:marked-opposite-rank}節の二つの無跡三列は、
この単射性からいずれも独立である。九つの方程式$D_B=-D_A^{\mathsf t}$に対して
どの成分も次元21以上であり、階数2の解は定理\ref{thm:simple-zero-smooth-elliptic}の証明により
次元20以下である。階数1の解の次元も小さいことを確認する。
同じ証明の中心化代数の議論により、$A$の三次元空間はある共通固有ベクトル$v$を持つ。
無跡対称三行行列に対する$Xv\in kv$は二つの独立な線型条件なので、この三次元空間は$v$で決まる。
任意の対称行列で$Xv$を任意に指定でき、スカラーを引いても$V/kv$での値は変わらないからである。
二つの行列の交換子は$v$を殺し、歪対称行列と三次元ベクトルの同定で$D_A$の像は$kv$となる。
$v$が等方な場合も同じである。固定した$v$で基底の選択は九次元で、
余因子による列変換は像直線が固定された全ての非零な階数1行列（三次元）に推移的に作用する。
従って指定した階数1の値に対する三列のファイバーは六次元である。
階数1の値は高々五次元なので、対$(A,B)$は高々17次元である。
階数0は独立な可換三列を要し不可能である。
従って各成分の一般点で$D_A$は可逆で、全解は命題\ref{prop:marked-net-parameters}の
可逆チャートの閉包に入る。このチャートを基底付きで飽和した閉包と楕円鉛筆の閉包の一致は
定理\ref{thm:simple-zero-smooth-elliptic}の証明の最後による。
\end{proof}

$0\le i<k\le3,\ 0\le j<l\le3$のGram方程式を
$G_{ik,jl}=\langle\beta_{ij},\beta_{kl}\rangle-\langle\beta_{il},\beta_{kj}\rangle$とする。
$i=j=0$の九式は$U\perp V$を表す。他の27式を$T$の54座標で微分した行列を$J_T$とする。

\begin{theorem}[非横断点を通る相対Jacobi判定]\label{thm:marked-gram-smoothing}
隣接階数6、反対階数8、両側の反対核が滑らかな四次完全交叉を定めるとする。
ある零点対で$\operatorname{rank}J_T=27$なら、テンソルは基点自由楕円鉛筆の表示を持ち、
全次数で標準Koszul AS正則である。特に$R=U\cap V$の中に
\[
 C_r=(\langle r,T_{ij}\rangle)_{1\le i,j\le3},\qquad \det C_r\ne0
\]
を満たす$r$があれば、この判定が通る。
\end{theorem}
\begin{proof}
命題\ref{prop:smooth-apolar-defect}から$U,V$は三平面で$V=U^\perp$である。
$U$をGrassmann多様体上で動かし、$U,U^\perp$に基底を選ぶ空間を台とする。
これは滑らかで既約であり、非退化三平面が稠密である。
台の上で$T$を自由に動かす54次元アフィン束に残りの27式を課す。
指定の階数なら相対Jacobi判定により解空間から台への射はその点で滑らかである。
ゆえにこの解は$U$が非退化な解の閉包にある。
両側とも指定文字の一次倍へ制限した反対切断が単射となる開条件はもとの点で成立する
（滑らかな反対核には可約二次式がない）。同じ開部分内で近傍を取れる。
補題\ref{lem:independent-marked-closure}により楕円閉包への所属が特殊点にも残り、
系\ref{cor:smooth-apolar-specialization}から目的の表示が従う。

最後の十分条件を示す。27式のうち一方の指定文字だけを含む18式は$T$で線型である。
それらの微分が独立なことを確かめる。行関係の第一組は三階テンソルの係数として、
第一添字が$i$、第二添字が$j$、最終因子が$H$の元となり、
最終因子は$U$内で、$u_j$による同定で第二・第三添字について交代的である。
もう一組では最終因子が$V$内で、$v_i$による同定で第一・第三添字について交代的である。
二つの像の共通元は最終因子が$R$に入り、交代性から第二・第一因子も
各々対応する$R$の逆像に入る。従って三因子を$R$で同定した共通元は完全交代的となる。
$\dim R\le2$なので零である。各組だけの行は基底の独立性により独立である。

$\dot T_{ij}=h_{ij}r$とすればこの18式の微分は消える。残り九式への微分は、
二次小行列式の写像を$C_r$で微分して行列$(h_{ij})$に適用したものである。
三行可逆行列の余因子写像の微分は可逆である。恒等行列では
$h\mapsto\operatorname{tr}(h)I-h^{\mathsf t}$であり、他の点には基底変換で移せる。
従って残る九式にも全射である。
\end{proof}

\path{computations/verify_smooth_apolar_defect.py}は$y^2=x^3-x$上で無限遠点を$O$、
$P=(0,0)$として$L=\mathcal O(4O), M=\mathcal O(3O+P)$を取る。
鉛筆を$s$と$x^2y+y/x+11$で張り、印は$(s,1)$とする。
$s=1,x-2,y+1$で接空間次元は順に2,1,0となり、全てで$J_T$の階数は27である。
中央例の根基直線では$C_r$の階数が2なので、27式の検査は可逆な$C_r$を探すだけよりも広い。

\begin{corollary}[接空間二次元の印を持つ全候補]\label{cor:defect-two-decision}
同じ滑らかな反対核の仮定で、接空間次元2の零点対を一つ持つとする。
この印の根基$R$からの三行行列の鉛筆に階数1の元があれば標準ASではなく、
それ以外では楕円鉛筆表示が存在して標準Koszul AS正則である。
接空間次元1の印でも、根基直線の行列が階数1なら同じ排除が成立する。
\end{corollary}
\begin{proof}
$r\ne0$は$H$におけるGram像全体を消さない。対応する六式の線型結合は、
印付きの第一行と第一列が零で残る三行三列が$C_r$である。
階数1なら純テンソル関係があり、次の反転窓ではその転置が関係になる。
定理\ref{thm:no-consecutive-pure-relations}により標準ASと両立しない。
$R$が二次元で階数1の元がない場合、可逆な元があれば前定理を使える。
残るのは全ての非零元が階数2の行列鉛筆である。

このような三行鉛筆の型を確認する。$\mathbb P^1$上の写像
$\mathcal O^3\to\mathcal O(1)^3$の核を$\mathcal O(-e)$、余核を$\mathcal O(1+f)$とすれば、
両端で階数一定なので$e,f\ge0$かつ$e+f=2$である。
$e=0$なら共通核を除くと$\mathcal O(-1)^2\hookrightarrow\mathcal O^3$で余核は$\mathcal O(2)$となる。
商を与える三二次式は核に定数切断がないので独立で、$s^2,st,t^2$の基底に揃えられる。
その二本の一次関係を取れば、鉛筆は下の左行列に同値である。$f=0$は転置である。
$e=f=1$なら右・左の核直線を示す二つの一次ベクトルの成分はそれぞれ一次式全体を張る。
両方を$(-t,s,0)$に揃えると、最初の二列と最初の二行がそれぞれ$s,t$倍となり、
第三列から他の二列の倍数を引き、定数を正規化して右の形となる。
\[
 \begin{pmatrix}s&0&0\\t&s&0\\0&t&0\end{pmatrix},\qquad
 \begin{pmatrix}0&0&s\\0&0&t\\s&t&0\end{pmatrix}.
\]
各係数行列を$C_1,C_2$とする。根基内の変分だけで得られる九式の微分は、
二次小行列式写像の$C_1,C_2$での二つの微分の像の和である。
左の形とその転置では階数9、右では階数8で、右で不足する成分は左上二行二列の小行列式である。
これは二行・二列を独立にとり
$C_{ij}h_{kl}+h_{ij}C_{kl}-C_{il}h_{kj}-h_{il}C_{kj}$の係数を並べた行列の行消去で得る
（\path{computations/verify_smooth_apolar_defect.py}）。左の形ではこれで十分である。

右の形の不足する一式を埋める。補空間の添字空間を$I,J$と書くとこの形は
$C_\alpha=a_\alpha\otimes d+c\otimes b_\alpha$ ($\alpha=1,2$)で、
$a_1,a_2,c$および$b_1,b_2,d$は各双対空間の基底である。
$R$の基底$r_1,r_2$の逆像を軸の同型$J\to U,I\to V$でそれぞれ$j_\alpha,i_\alpha$とする。
線型Gram式から
$C_1(-,j_2)=C_2(-,j_1)$、$C_1(i_2,-)=C_2(i_1,-)$となる。
独立性により$J_R=\langle j_1,j_2\rangle=\ker d$、
$I_R=\langle i_1,i_2\rangle=\ker c$となり、
$a_1(i_2)=a_2(i_1),\ b_1(j_2)=b_2(j_1)$でもある。
不足した一式は$I_R\times J_R$の二行二列に対応する。
$j_3,i_3$を補い、対応する$u_3,v_3$の自己対合を1に規格化する。
$P=\langle u_3,T\rangle, Q=\langle v_3,T\rangle$とすると線型式は
\[
 P(i,j_\alpha)=C_\alpha(i,j_3),\qquad Q(i_\alpha,j)=C_\alpha(i_3,j)
\]
を与える。特にこの小矩形で$P$は列$a_\alpha|_{I_R}d(j_3)$を持ち可逆である。
ここで$\dot C_\alpha=a_\alpha\otimes d$とする。
先の二つの対称条件を満たし、$\dot P,\dot Q$を表示式に従って決めれば
線型18式を保つ変分に持ち上がる。実際、各$\dot T_{ij}$に対して指定するのは
独立な$r_1,r_2,u_3,v_3$との対合だけで、$\dot P$の第三列・$\dot Q$の第三行は自由である。

小矩形で$C_\alpha,\dot C_\alpha$はともに零だから、$T,\dot T$の各値は
$R^\perp=U+V$に入る。この矩形のGram式は$\det P+\det Q=0$であり、
構成した変分はその上で$\dot P=P,\dot Q=0$となる。その微分は$2\det P\ne0$である。
従って27式全体で階数27となり、前定理が適用できる。
\end{proof}

\begin{lemma}[根基直線上の階数2の印の平滑化]\label{lem:defect-one-rank-two}
前定理の仮定で$\dim R=1$かつ$C_r$が階数2なら、テンソルは楕円鉛筆の閉包に入る。
\end{lemma}
\begin{proof}
$r$を根基の生成元とし、両軸の$r$に対応する添字を$i_0\in I,j_0\in J$と書く
（指定した零点の文字を除いた三次元空間内の添字である）。
$U/kr,V/kr$の直交標準基底を持ち上げ、$r$に対する双曲対のもう一方の元をこれら四元と直交させる。
$T$の九個の値の座標を、この基底との対合として$C,H,P_1,P_2,Q_1,Q_2$の六行列で表す。
線型18式を解くと、$C,H$は自由である。$J=kj_0\oplus J'$とし、
二次元の$J'$を$U/kr$の上の基底で同定すると、各$i$に対して
\[
 P(i,j_0)=C(i,-)|_{J'},\qquad P(i,j')=S_i j'\quad (j'\in J')
\]
である。右辺の共変ベクトルは直交計量でベクトルと同一視し、
$S_i$は任意の対称二行行列である。$Q$側も$i,j$を交換した同形の九パラメータである。
よって自由な座標は$9+9+9+9=36$個となる。
残り九式は$C$での余因子微分を$H$に適用したものと、
各$P_\alpha,Q_\alpha$の二次小行列式との和である。

軸を固定し$C$の階数が2以上の解空間の次元が高々27であることを示す。
可逆なら$H$が一意で、$C,P,Q$の自由度が27である。
階数2なら$C$の選択は八次元で、余因子微分は階数6・核三次元である。
実際$C=\operatorname{diag}(1,1,0)$にすれば直接分かる。
左・右の核の生成元を$i_*,j_*$とし、二つの商を$\bar I=I/ki_*,\bar J=J/kj_*$とする。
この微分の余核の三座標は、九式のうち行対$(i,i_*)$、列対$(j,j_*)$の値を
$\bar I,\bar J$上の二行二列として読み、$kC$を法として取るものである。
$H$で解けるためにはこの三座標が全て零でなければならない。
従って条件を課す前の上界$8+9+9+3=29$を少なくとも二つ下げればよい。

$j_*=j_0$（直線の一致）なら、階数2の右核が指定直線となること自体が余次元2である。
以下$j_*,j_0$は独立とする。$S_{i_*}$がスカラー行列なら九個の$P$の自由座標に
独立な二条件がかかる。残る場合は$C,Q,S_{i_*}$を固定する。
$j_*$の$J'$成分を$k$と書き、それと独立な$l\in J'$を選ぶ。
$j_0,l$の像は$\bar J$の基底である。$i_*$を補う二つの$i$について$S_i$を独立に変える。
先の二行二列のうち$j_0$の列は変わらない。
実際$P(i,j_0)$は固定され、$P(i_*,j_0)=0$である。
$l$の列の$i$行の変化は
\[
  \langle \dot S_i l,S_{i_*} k\rangle-\langle\dot S_i k,S_{i_*} l\rangle
\]
である。$k,l$は基底で、全ての対称二行行列と可換な行列はスカラーに限られるので、
これは$\dot S_i$の非零の線型形式である。二行を独立に指定でき、
$C(-,j_0)\ne0$なので得られる二次元空間は$kC$を法としても単射に写る。
固定した他の座標のもとで条件は$S_i$についてアフィン線型である。従って少なくとも二条件がかかる。
これで全ての場合に固定軸の上界27が従う。

二つの軸を動かす台は先のGrassmann多様体の基底束である。
根基直線を持つ三平面はその台で余次元1以上である（非退化条件の行列式が恒等零ではない）。
根基が二次元以上の閉部分を除き、根基が直線の閉境界上で$C$の階数が1以下となる
閉条件も除いてよい。指定点は残る。54変数に27式を課した全解空間の成分の次元は
台の次元に27を加えた値以上である。一方、先の上界により残った閉境界だけでは
この次元に達しない。従って指定点は相補三平面を持つ解の閉包にある。
両軸での単射性の開条件を保ち、補題\ref{lem:independent-marked-closure}を用いればよい。
\end{proof}

\begin{theorem}[滑らかな反対四次曲線を持つ反転四窓の分類]\label{thm:smooth-apolar-complete}
両対称テンソルの隣接切断が階数6、反対切断が階数8で、
両反対核が滑らかな四次完全交叉を定めるとする。付随する反転四窓が標準Koszul AS正則となる
必要十分条件は、第一窓の六次元関係空間が非零の純テンソルを含まないことである。
この条件は全てのAS点が基点自由楕円鉛筆から得られることとも対応し、各点で
二つの曲線の同定と順序付きの二束・鉛筆が復元される。
\end{theorem}
\begin{proof}
純テンソルのある場合はその転置と連続するから既出の成長障害を使う。
純テンソルがない場合、六つの$(1,1)$式には少なくとも一つ零点対がある。
実際、六次元Segre多様体を射影空間で六つの超平面と交わらせている。
印の接空間次元は命題\ref{prop:smooth-apolar-defect}により$0,1,2$のいずれかである。
単純点では定理\ref{thm:simple-zero-smooth-elliptic}を使う。
次元1なら根基の行列は非零で階数1ではない。階数3は定理\ref{thm:marked-gram-smoothing}、
階数2は補題\ref{lem:defect-one-rank-two}と系\ref{cor:smooth-apolar-specialization}による。
次元2は系\ref{cor:defect-two-decision}で尽くされる。いずれも滑らかな曲線と基点自由鉛筆へ戻るので、
全次数の正則性とデータの一意性は定理\ref{thm:elliptic-pencil-reconstruction}による。
\end{proof}

\begin{corollary}[反転窓の反対階数10の排除]\label{cor:bisymmetric-high-rank-exclusion}
隣接階数6の両対称Gramテンソルで、両反対核がいずれも非零の一次式の積を含まないとする。
付随する反転窓が標準ASならテンソルは楕円鉛筆の閉包にあり、反対階数は8以下である。
特に反対階数10の反転窓は全て排除される。階数9の標準AS反転窓があれば、
少なくとも一方の一次元反対核は階数2以下の二次式で生成される。
\end{corollary}
\begin{proof}
指定文字の一次倍上の反対切断の単射性と、平方が核に入らないことは
両反対核についての仮定から従い、階数8や四次曲線の滑らかさを必要としない。
従って命題\ref{prop:smooth-apolar-defect}の三平面と接空間次元の議論が通る。
標準ASなら根基の階数1の元は純テンソル障害により存在しない。
単純点は補題\ref{lem:independent-marked-closure}、それ以外の印は前定理の
二つの平滑化証明を、楕円鉛筆の閉包を得る段階まで適用する。
閉包では反対切断の九次小行列式が消える。階数10では両反対核が零、
階数9ではそれぞれ一次元なので残る二主張を得る。
\end{proof}

% !TeX root = ../papers/nc-p3.tex
\section{整な特異四次曲線を保つ極限}
\label{sec:integral-apolar-limits}
\begin{theorem}\label{thm:integral-apolar-alternatives}
楕円鉛筆の特殊化で、片側の定義二次式の極限が整な節点または尖点四次曲線$X$を定め、
特殊点が標準ASとする。可逆な二束を持つ定理\ref{thm:elliptic-pencil-integral-boundary}の表示か、
次の正規化上の表示のいずれかとなる。後者の必要なデータは$\nu:P=\mathbb P^1\to X$、
$A=H^0(X,\mathcal O_X(1))\subset H^0(P,\mathcal O(4))$、$B=H^0(P,\mathcal O(3))$と
$H^0(P,\mathcal O(7))$内の二平面$W$である。第一窓は積を$W$で割った写像の核の直交補で定まる。
後者の反対階数は8で、他側の反対核はねじれ三次曲線とその割線（節点）または接線（尖点）の合併を切る。
従って整な四次曲線を保つどの特殊化にも反対階数7の標準AS点はない。
後者は候補表示であり、定理\ref{thm:all-conductor-pencils-excluded}で全鉛筆を排除する。
\end{theorem}
\begin{proof}
特殊ファイバーまで延ばした完全交叉曲線族を$\mathcal X$、第一束を$L$とする。
付値環を必要なら完備化・有限拡大し、特殊点が非特異な切断$o$を選ぶ。
一般点上の第二束は$L_K(-o)(p)$と書ける。次数1の束は唯一の有効因子を持つからである。
固有性で$p$を切断に延ばし、そのイデアルを$J$とする。
$M=L(-o)\otimes\mathcal Hom(J,\mathcal O_{\mathcal X})$を用いる。
この双対は平坦で、ファイバー上で切断のイデアルの双対になる。
実際、双対列で$\mathcal O$の余核は$\mathcal Ext^1(\mathcal O_p,\mathcal O)$で、
切断上の可逆加群である。これは相対Gorenstein曲線における切断の双対性から従い、
付値環の剰余体への基底変換と可換である。局所的にも、切断を$(u,v)$、
曲線式を$u a+v b$で書けば、イデアルの双対は恒等写像と$(u,v)\mapsto(b,-a)$で生成される。
後者の類が切断上の自由な余核を生成し、剰余体上でも同じ表示を得る。

$p$が非特異点に落ちれば$M_0$は可逆である。特異点に落ちればその極大イデアルは導手であり、
逆イデアルは正規化環である（節点では二枝の値、尖点では一次係数の条件を外す）。
従って$M_0=\nu_*\mathcal O_P(3)$で、$L_0M_0=\nu_*\mathcal O_P(7)$となる。
いずれの場合も$h^0(M_0)=4,h^0(L_0M_0)=8$で$H^1=0$である。
可逆の場合の切断の積は既出の評価核により全射である。非可逆の場合は導手因子
$\Delta\subset P$（相異なる二点または二重点）で消える四次式全体が$A$に入る。
それに全三次式を掛ければ$\Delta$で消える七次式を全て得る。
$\Delta$上は単元となる$A$の元と三次式の積で任意の値を得るので、ここでも全射である。

命題\ref{prop:one-smooth-quartic-specialization}の格子の最大指数の議論を適用する。
低い指数の列に独立な$b,b'$があれば、その積$A b+A b'$が六次元以上であることを示す。
正規化上で共通因子を除いた二式の次数を$d\ge1$とする。積の核は
次数$4-d$の形式$h$で、互いに素な残りの両式を掛けるといずれも$A$に入るものからなる。
$d\ge3$なら次元は2以下である。$d=2$なら条件が恒等零になるには両式が$\Delta$で消えなければならない。
実際、$h$は$\Delta$上で任意の値を取り、$A$の貼り合わせ超平面は節点で両端の係数が非零、
尖点で微分の係数が非零である。どちらの局所長さ2の環にも、その超平面に入る非零イデアルはない。
互いに素なので少なくとも一条件がかかる。$d=1$なら二つの一次式が$\Delta$上の
全切断を張り、同じ理由で$h$自身が$\Delta$で消える。三次式のこの空間は二次元である。
よって主張が成立し、最大指数の列の評価の階数が2以下になる。

低い列が$s$一本なら、可逆な場合はその長さ4の零スキーム内の任意の長さ3の部分で消える
非零$x\in A$を取ればよい。長さ4のArtin環には長さ1のイデアルがあり、
長さ3への評価は四次元の$A$上で核を持つ。非可逆な場合は$s$が三次式で、
$s H^0(\mathcal O_P(1))$は四次式の超平面$A$と非零の$x$で交わる。
この$x$と全ての高い列の積は$s H^0(\mathcal O_P(4))/(sA)$の一次元に入る。
いずれも$x$の評価が低い一列を含めて階数2以下となり、反転窓の成長障害に反する。
従って最大・最小指数は等しく、特殊窓は実際に$A\otimes H^0(M_0)\to H^0(L_0M_0)/W$の表示を持つ。
可逆な場合、基点は純テンソルを与えるので元の整な鉛筆定理が適用できる。

非可逆な場合に反対切断を決める。$P$に別の直線$T$を$\Delta$で貼り合わせた種数1の曲線$Y$を使う。
二点で貼る場合は二成分の環、二重点の場合は接する二枝で、双対化層はいずれも自明である。
$T$を一点に縮めて$X$を得る射の引き戻しを$L'$とし、$M'$は$P$上で次数3、$T$上で次数1として貼る。
$\Delta$上の同定と$M'$の貼り合わせは任意に一つ選ぶ。
$T$上の一次切断は$\Delta$上の値で一意に決まるから、
$H^0(L')=A,H^0(M')=B,H^0(L'M')=H^0(\mathcal O_P(7))$で積は所与のものとなる。
両束は大域生成されるので、$W$に$Y$上の基点があればその評価が純テンソル関係になる。
従ってASの仮定のもとで$W$は$Y$上で基点自由である。

$|M'|$は$P$をねじれ三次曲線に、$T$を$\Delta$の張る直線に写す。
異なる二点の割線または二重点の接線と三次曲線の交叉はちょうど$\Delta$である。
その合併は二二次式完全交叉である。実際、三次曲線の三本の二次式の直線への制限は
$\Delta$で消える一次元空間を張り（標準単項式座標で確かめられる）、二本が残る。
可約二次式はねじれ三次曲線を含まないので共通因子はなく、合併が次数4を使い切る。
ゆえに両束の対称二乗は各々八次元の全二次切断に全射である。
$L'/M'$およびその逆は多重次数$(1,-1)$と$(-1,1)$で、次数1の成分だけでは
長さ2の貼り合わせに沿って消えられないため切断を持たない。

鉛筆完全列とSerre双対で切断の積のtraceを四線型形式にすると、
隣接対合は階数6で所与の六式を像とし、反対対合は階数8で両埋込みの二次核を持つ。
これは鉛筆完全列の計算であり、ここでは切断環のKoszul性を用いない。
特殊四窓が標準ASなら四連続の二次関係からの最上交叉は一次元なので、
この非零テンソルがその標準テンソルを張る。これで階数と反対核が従う。
\end{proof}

この非可逆な場合でもAS点の同型から曲線データを戻せる。整な側では特異点で消える二次切断、
他側では直線上で消える二次切断に制限する。両者は$P$上で固定因子を除けば
次数6と4の全切断であり、次元は7と5である。完全な八次元反対対合をここに制限すると階数は4以上である。
積は$T$上で零なので、正規化とねじれ三次曲線を同定するグラフへの制限を通る。
異なる二つの射影直線同型のグラフは長さ2で交わり、合併のイデアルで$(6,4)$をひねると
$(4,2)$で$H^1=0$となる。従って二つの線型包の共通部分の行列階数は高々2である。
これで同定が一意に決まり、$W$も隣接核から戻る。逆にこの正規化上の順序付きデータの同型は六式を保つ。

\paragraph{切断環の三次障害。}
非可逆模型の$Y$上の交互切断環は二次生成関係だけでは表示できない。
節点では$A=\langle1+t^4,t,t^2,t^3\rangle$、尖点では
$A=\langle1,t^2,t^3,t^4\rangle$に揃えられる。前者では二枝の座標のスカラーを、
後者では$t/(1+ct)$の四次形式への作用を用いて貼り合わせ係数を規格化する。
$B=\langle1,t,t^2,t^3\rangle$として、$B,A,B$の二次乗法核から作る三次商は両方で13次元、
実際の切断は12次元である。三次の積は全射である。主成分では$AB$が全七次式で
さらに$B$を掛けると全十次式となる。主成分上で消えて$T$上で$\Delta$に沿って消える
一つの二次式は、ねじれ三次曲線上で消え直線上では消えない二次式を取り、
$T$で非零の$A$の切断を掛ければ得られる。
$A,B,A$の二次関係からの三次商は12次元となる。
\path{computations/verify_conductor_pencil.py}はそれぞれの二次乗法核を64行の行列に入れてこの次元を
有理数上で検証し、各型の基点自由な一鉛筆で隣接・反対階数6,8と両三次交叉の次元4を検査する。
曲線の交互切断環の非二次性から、二次の鉛筆を割った関係列の非AS性は導けない。

\begin{proposition}[非可逆模型の全鉛筆における有限双対]\label{prop:conductor-pencil-frobenius}
上の正規化表示で、幾何的な二成分曲線$Y$上で鉛筆$W$が基点自由とする。
全てのこのような$W$で、二次関係列の二次双対は次元列$1,4,6,4,1$の有限Frobenius代数である。
第三切断階数は8で、順序付き正規化データによる窓の同型判定も成立する。
\end{proposition}
\begin{proof}
まず二次の積だけで定めた交互切断環の二次表示を考える。
$A,B,A$の三次商は上の計算で12次元で、実際の切断への写像は全射である。
実際$AB$は全七次式であり、$A$と七次式の積も導手部分と$\Delta$上の値を分ければ全十一
次式を張る。この位置では$T$への制限は次数1で、主成分での値から決まる。
従って三次商と切断が一致し、$W$を両辺から割れば$W A$の八次元が消えて四次元になる。
$W A$の次元は、基点自由鉛筆の列で$A$に対応する束へ積を取り、積の核が
$H^0((M')^{-1})=0$となるためである。$B$についても同様に$\dim W B=8$である。

$B,A,B$の三次商から切断への全射の核は一次元である。その余分な方向は、
切断の積と合わせて、次の線型汎関数で検出できる。
\[
 b\otimes a\otimes b'\longmapsto
 a|_T\,\det(b|_T,b'|_T).
\]
ここで$L'$は$T$上で自明であり、$M'|_T=\mathcal O(1)$の二切断の行列式を取る。
この式は左の二次積の後にも右の二次積の後にも定義できるので、二次関係を消す。
しかも外側の二文字の交換で符号が変わり、通常の切断の積では得られない非零の汎関数である。
$w\in W$を左右の二次積へ互いに反転する形で持ち上げれば、
$w b-b w$のこの値は$2\det(w|_T,b|_T)$である。
基点自由性により$W$の$T$への制限は二次元であり、行列式を非零にできる。
従って$W$を両辺から割る際には、切断上で$W B$が消えるのに加えてこの一次元も消える。
結果は$13-8-1=4$次元となる。

上のtraceテンソルはどの三次交叉でもその全体を一文字縮約で張る。
縮約が四次元となるのは、文字が三次元以下で足りるならその側の対称二乗を通る反対階数が
6以下となるためである。四次交叉の元は、最初の三文字または最後の三文字をまとめると、
traceテンソルに第一位置の自己準同型を施したものでも、第四位置の自己準同型を施したものでもある。
後者は第一・第三位置で対称なので、前者の第一位置の作用も八次元の対称形式で自己随伴になる。
定理\ref{thm:elliptic-pencil-reconstruction}の次元7の上界の証明は非可逆な自己準同型にも通る。
従って第一位置の作用はスカラーで、四次交叉は一次元である。
五次交叉はこの四次テンソルを左・右に一文字延ばした二空間の共通部分である。
非零の共通元は第一位置だけを切ると、左が階数4、右が階数1となりえない。
従って五次以上は零である。二次双対はこれでtraceから得る完全対合の商と全次数で一致する。
両反対核とグラフの復元にはこのtraceとその階数だけを使ったので、順序付き同型判定も従う。
\end{proof}

純テンソル関係が無い条件もこの模型では幾何的に読める。
$H^0(\mathcal O_P(7))^*$内の汎関数が$A\otimes B$で階数1なら、
ある超平面$A'\subset A$について$A' B$を消す。
$A'$が$P$上で基点自由なら一般の二切断が互いに素な四次式で、
三次式との積は八次元全体を張るので矛盾する。従って$A'$は点での評価核である。
特異点の原像以外では、$|A|$はそこで埋込みで他の点と分離するので共通因子はその一点だけである。
一般二切断から共通因子を除いた互いに素な三次式と$B$の積は全六次式を張り、
元の汎関数はその点の評価に限られる。特異点の原像では$A'$は$\Delta$で消える全四次式なので、
汎関数は$\Delta$への制限を通る二次元空間で尽くされる。
この空間の全ての非零元は第一因子が特異点での評価の倍数で、実際に階数1である。
ゆえに射影化した階数1の軌跡は七次有理正規曲線と$\Delta$の張る直線の合併で、
$W^\perp$がこれを避けることは$Y$上の基点自由性と同値である。
また一つの文字の評価核が二次元以上なら、その文字と相手側の切断の積の中に$W$全体が入る。
主成分上では一文字が必ず零点を持つので、これも基点自由性に反する。

\begin{corollary}[七次で排除される二つの鉛筆]\label{cor:conductor-degree-seven}
節点型と尖点型の双方で、$W=\langle1+2t+4t^2+7t^6+3t^7,\ t+4t^4+t^7\rangle$
は基点自由で有限双対がFrobeniusであるが、標準ASではない。
\end{corollary}
\begin{proof}
両七次式の終結式は非零で、例外直線上でも二切断は独立である。
上記コードは、七次式の単項式座標における積の係数行列に$W^\perp$の行基底を掛け、
その六行と転置順の六行を交互の関係として使う。
既存の\path{computations/two_window_ext.py}の経路商で、次数$d-2$の基底に
最後の二文字の六関係を掛けた全ての像を、次数$d-1$の基底に一文字掛けた空間から消去する。
有理数上の両始点の次元は六次まで$1,4,10,20,35,56,84$であるが、
七次では$A$始点で120、$B$始点で121となる。標準値$\binom{10}{3}=120$に反する。
\end{proof}

節点型ではさらに$W_u=\langle1,t^7+u t^3\rangle$の全ての$u\in k$で標準AS性を排除できる。
主成分と例外直線で基点がない。上記コードは$1=0,t^7=-u t^3$で七次式の積を商に写し、
有理関数体$\mathbb Q(u)$上で$B$始点の七次次元121を得る。特殊化で一定階数の六関係が定める
各経路商の次元は上半連続なので、全ての有限パラメータで七次元数は121以上である。

% !TeX root = ../papers/nc-p3.tex
\section{二変数形式の積への商と非可逆境界の排除}
\label{sec:binary-product-obstruction}
$P_j=H^0(\mathbb P^1,\mathcal O(j))$と置く。
\begin{lemma}[二変数形式の積の双対関係]\label{lem:binary-diagonal-paths}
一次生成空間を$P_d^*,P_e^*$と交互に置き、二次関係を乗法
$P_d\otimes P_e\to P_{d+e}$の双対の像とその転置に取った経路代数を考える。
長さ$n\ge2$の次数列を$d_1,\ldots,d_n$、端の位置の$v_i$を1、
内部の$v_i$を2とすると、経路空間の次元は
\[
 \prod_{i=1}^n\max(0,d_i-v_i+1)
\]
である。特に一方の次数が2、他方が3なら指数成長する。
\end{lemma}
\begin{proof}
自由経路空間の双対を、各位置に別の斉次座標$(s_i,t_i)$を用いた多重斉次形式と同一視する。
位置$i,i+1$に関係を挿入した空間の直交補は、この二変数組を同一視したとき零となる形式である。
対角の斉次イデアルは$s_i t_{i+1}-t_i s_{i+1}$で生成される。
全位置の条件を同時に課すと、多項式環の相異なる既約元であるこれらの式の積で割り切れる。
逆に積の任意の倍数は条件を満たす。残る多重次数が$d_i-v_i$なので式が従う。
\end{proof}

\begin{theorem}[全ての導手鉛筆の指数成長]\label{thm:all-conductor-pencils-excluded}
定理\ref{thm:integral-apolar-alternatives}の非可逆模型では、節点・尖点のどちらでも
任意の基点自由鉛筆の反転二窓は指数成長する。従って標準ASではない。
片側の極限二次式が整な四次完全交叉を切る標準AS特殊化は、
その同じ曲線上の二つの可逆層と基点自由鉛筆から得られる。
\end{theorem}
\begin{proof}
導手因子$\Delta$を切る二次形式を$h$とする。$A$に含まれる三次元空間$hP_2$に双対で射影し、
もう一方は$B=P_3$のままとする。六関係はもともと$P_7/W$上の汎関数を積で引き戻した像なので、
射影後は$P_2\otimes P_3\to P_5$の双対像に含まれる。
従って補題の代数への次数付き全射がある。実際、$W$は例外直線の全一次切断に同型に制限され、
$hP_5$はその上で零なので、$hP_5\to P_7/W$は六次元同士の同型であり、射影後の関係は双対像全体である。
$B$始点の長さ$2m+1$で商の次元は$9\cdot 2^{m-1}$ ($m\ge1$)。
$m=8$で既に$1152>\binom{20}{3}=1140$である。滑らかな場合を含む最後の主張は
命題\ref{prop:one-smooth-quartic-specialization}と定理\ref{thm:integral-apolar-alternatives}による。
この一様な排除は、全ての鉛筆で七次が最初の障害かどうかを必要としない。
\end{proof}

\begin{corollary}[片側が整な階数8の判定]\label{cor:one-integral-apolar-recognition}
隣接階数6・反対階数8の両対称テンソルについて、一方の反対核が整な四次完全交叉を切るとする。
反転四窓が標準ASとなることと、その整曲線上の相異なる次数4の可逆層の対と基点自由鉛筆による
表示を持つことは同値である。このとき両側の反対核は同型な曲線を切り、順序付きデータを復元する。
\end{corollary}
\begin{proof}
標準ASなら定理\ref{thm:all-bisymmetric-as-closure}により滑らかな楕円鉛筆から特殊化する。
反対階数が保たれるので指定側の鉛筆は極限二次式と一致する。整な特異完全交叉四次曲線は
既出の種数公式により節点または尖点を持つ。前定理で非可逆な場合を排除でき、
可逆な場合の十分性・復元は整曲線の鉛筆定理による。
\end{proof}

\begin{proposition}[零次数の成分を許した多重次数判定]\label{prop:semipositive-pencil-balance}
有理曲線の環状多角形、または各枝が射影直線で唯一の特異点が楕円多重点となる曲線について、
$L,M$の多重次数$d,e$を非負とし、各総次数を4とする。
基点自由な$LM$の鉛筆の反転四窓が標準ASとなる必要十分条件は$d=e$である。
\end{proposition}
\begin{proof}
両束とも次数零の成分は切断を取る前に縮められる。環状の鎖の定数を消去すると
両端の値に元の貼り合わせのスカラーの積を課す。一成分だけが残れば二点を結ぶ節点となる。
多重点では零次数枝の定数を共通値とし一次係数を零とすれば、残りの枝に同じ一次の和の条件を得る。
一枝だけなら尖点である。束の積でもこの消去と貼り合わせ条件は可換である。
従って$d=e$なら既出の正次数の場合に帰着する。以下、全成分で$d_i+e_i>0$とする。成分数は高々8である。

相補的な非空の連続部分鎖$I,J$について、$I$だけに支えを持つ$L$の切断の次元は
$\max(d(I)-1,0)$である。鎖の部分の総次数が正なら、零次数の連続部分の定数を消去した後に、
二つの端点の値を独立に選べるためである。総次数零なら端点で消える定数は零である。
楕円多重点では任意の真部分集合で同じ式が成立する。共通値を零とし、
次数が正の枝があれば一次係数に独立な一式を課すからである。
これらの部分に支えを持つ相補的な切断の積は零で、各々の双対への射影で全二次関係が消える。
ゆえに二次元数の積が2以上なら指数成長となる。$d,e$の役割を交換した条件も使える。

この検査に残る$d\ne e$は、成分と束の順序を除き
$((4,0),(3,1))$と$((3,1,0),(3,0,1))$だけである。
有限列挙の方法は、各成分への四単位の配分を二回選び、$d_i+e_i>0$を課し、
全ての巡回区間$J$で和$d(J),e(J)$が同じ区分
$\{0,1\},\{2\},\{3,4\}$に属するものを残すことである。
区分が異なれば上の積の少なくとも一方が2以上となる。
\path{computations/verify_binary_product_obstruction.py}は成分数2から8を厳密列挙する。
多重点では枝に任意の巡回順を付けてよい。

二成分の場合、主成分での第一束の切断のうち、二つの貼り合わせ点
（接する場合はその長さ2の導手）で消えるものは$hP_2$で、他成分で恒等零である。
第二束は主成分へ全三次式として制限できる。従って積は$P_2,P_3$の積を通る。
三成分の場合には両束とも主成分上で全三次式を取れる。環状なら主成分の二端のうち、
第一束が次数零の枝と交わる点において消える三次式を第一束に取る。これは一次式倍の$P_2$である。
第一束の切断はその零次数枝全体でも消える。残りの枝の上では第二束が定数なので、
積は主成分での積の端点の値と、もう一端の値零から一意に決まる。
三重点なら三次式を共通点で消す。第一束は零次数枝で消え、残りの補助枝では第二束が定数である。
そこでの積の値は零で、一次係数は主成分上の積の一次係数から一意に決まる
（一次係数の和の条件による）。いずれも選んだ切断の積は主成分の$P_2\otimes P_3\to P_5$を通る。
補題\ref{lem:binary-diagonal-paths}の商により、残った二型も指数成長する。
\end{proof}

% !TeX root = ../papers/nc-p3.tex
\section{環状四次曲線を保つ反転窓の特殊化}
\label{sec:polygon-apolar-limits}
\begin{lemma}[特殊ファイバーの積と極限の六式]\label{lem:mixed-rank-loss}
楕円鉛筆の付値環上の特殊化で、一般ファイバーの二束の積の切断空間は八次元とする。
第一・第二座標格子を$A,B$、隣接の極限六式に双対な商写像を$\beta$とする。
両射影写像を解消した曲線族の特殊ファイバーの被約成分上へ積を制限した写像を
$\Phi:A_0\otimes B_0\to\bigoplus_i H^0((LM)|_{P_i})$とし、その階数を$r$とする。
$\Phi$で零になる任意の部分空間の$\beta$による像の次元は$8-r$以下である。
$r=8$なら$\beta$は$\Phi$の像を二平面で割る写像である。
\end{lemma}
\begin{proof}
分数体上で全積が張る八次元空間の中に$A\otimes B$の像の格子$D'$を取り、
さらに一般鉛筆で割った空間での像の格子$H'$を取る。ともに自由で階数は8,6である。
$A\otimes B\to D'$の核の元は一般ファイバーで零なので、両写像を延ばした整な全空間で
切断として零であり、特殊ファイバーの被約成分上でも零である。
従って$D'_0$から$\operatorname{im}\Phi$への全射がある。
$D'_0\to H'_0$も全射で、六式の一定階数の飽和極限に対応する。結論が従う。
\end{proof}

\begin{lemma}[二つの被約モデルの成分]\label{lem:nodal-core-resolution}
一般点で同定された二つの四次完全交叉曲線族について、一方の特殊ファイバーを節点だけを持つ
有理曲線の多角形（整な節点曲線も含む）、他方も被約とする。
両射影を持つ正則な共通解消の特殊ファイバーの被約支持は、有理曲線の一つの環とそれに付く木からなる。
第一射影の正次数成分は環上にある。他方も多角形ならその正次数成分も全てこの環上にあり、
楕円多重点を唯一の特異点とする有理枝の合併なら環全体が一点に写る。
\end{lemma}
\begin{proof}
節点の局所平滑化の完備方程式は$uv=t^a$である。実際、特殊ファイバーの二次項が非退化なので
平行移動と次数ごとの座標変換で$uv$と底の関数にでき、底の単元を吸収できる。
解消で二枝の間に有理二次曲線の鎖を挿入する（$u,t$のイデアルのブローアップの一方のチャートは
$u'v=t^{a-1}$、他方は正則）。必要なら先に有限基底変換し、自己節点も二成分間の交点にできる。
残る第二射影の有理写像は正則曲面上の点のブローアップで解消する。
交点でのブローアップは辺の細分、平滑点では末端の付加である。
各極限曲線族は完全交叉でCohen--Macaulay、一般ファイバーと特殊ファイバーの一般点で正則なので
全空間は正規である。共通解消からの固有双有理射は連結ファイバーを持つ。
特殊ファイバーの因子の押し出しは値域の特殊ファイバー因子である
（主因子$\operatorname{div}(t)$の押し出し）。被約性により、各像成分を支配する成分はただ一つで、
その重複度と写像次数はいずれも1である。対応する正規化間の写像は同型である。

第二像が多角形なら、各節点の二つの原像は連結ファイバーにより、直接の交点または
定値成分の鎖でつながる。各正次数成分の二端を順にたどると単一の環を作る。
異なる節点のファイバーは交わらず、整な自己節点の場合も一つの鎖が両端を結ぶので同じである。
第二像が多重点型で環上に正次数成分が一つだけあれば、環の残りを通じてその異なる二点が同じ像を持ち、
その枝の埋込み性に反する。二つ以上なら相隣る像は特異点で交わるため、各正次数成分の
環上の二端がその枝上の同一の原点に写り、やはり矛盾する。
整な尖点の場合も正規化写像の閉点での単射性で同じである。
\end{proof}

\begin{theorem}[二つの環状モデルの次数一致]\label{thm:two-polygon-specialization}
楕円鉛筆の特殊化で両側の極限二次式が被約な節点多角形を切るとする。
特殊窓が標準ASなら、二つの曲線は同じ多角形として同定でき、
その上の等しい正の多重次数の二束と基点自由鉛筆の表示を持つ。反対階数は6または8となる。
\end{theorem}
\begin{proof}
前補題の環に制限する。共に定値の木は特殊切断の値を制限しない。
それぞれの正次数成分の次数と束の貼り合わせを引き戻した環上の束を$L,M$、多重次数を$d,e$とする。
全切断は各四次元で、元のそれぞれの四座標がその全体を張る。
実際、多角形の零次数鎖の定数を消去すると元の束の切断に戻る。指定した座標は
被約完全交叉を埋め込む全一次切断であり、成分上への引き戻しは単射である。
共通の零次数成分を除くと$LM$の切断は八次元である。
$d=e$なら既出の正次数での積の全射性により補題\ref{lem:mixed-rank-loss}で$r=8$となる。
従って普通の二次元鉛筆の表示となる。基点があれば純テンソル関係を生じて標準ASと両立しない。
二曲線は同じ鎖を縮めた多角形であり、正則性・反対階数・復元は多角形の鉛筆定理による。

$d\ne e$とする。環上の貼り合わせ係数を共に1に正規化できる。
各成分の両端を固定する座標の倍率は二つの貼り合わせスカラーに指標$d,e$で作用するからである
（単項式の二端での係数を取る）。$d,e$は総次数が同じで異なるので有理数上独立で、
代数閉体上その二指標の写像は二次元トーラスへ全射である。
従って切断の積の階数下界は各多重次数ごとに整数行列で検査できる。

詳細な有限検査を示す。各成分で斉次単項式の係数を取り、各節点の両端係数を同じ変数にする。
零次数成分の両節点の変数も同じにする。内部係数は個別の変数で、
各束につき独立な四変数となる。各成分で二多項式を掛け、各係数を16個の双線型単項式で表す。
得た整数行列の二元体での階数$h$は$r$の下界である。
相補鎖への制限で積が零になる二切断空間の次元を$u,v$とする。
$\min(u,v)\ge1$かつ$\max(u,v)\ge 8-h+2$なら、補題\ref{lem:mixed-rank-loss}で
小さい側の一文字を固定すると大きい側にはなお二次元以上の評価核が残る。
転置の次窓と合わせると関係が消える一次元・二次元の交互商で指数成長する。

\path{computations/verify_polygon_multiplication.py}は四単位ずつの全配分を列挙し、
共通巡回移動・反転・二束交換で重複を除く。異なる配分174通りのうち、$h=7$の24通りは全て
上の次元不等式を満たす。$h=8$の150通りでは148通りが満たし、残る二型は
命題\ref{prop:semipositive-pencil-balance}の二変数の商で排除した二型である。
この二型では実際の六式も$\Phi$の像を経由するのでその排除を適用できる。
従って$d\ne e$は標準ASにならない。
\end{proof}

\begin{proposition}[多角形と楕円多重点を組み合わせた二側の排除]\label{prop:polygon-additive-exclusion}
楕円鉛筆の特殊化で、片側の極限二次式が有理多角形を、他側の極限二次式が
被約な楕円多重点型の四次曲線を定めれば、特殊窓は標準ASでない。
\end{proposition}
\begin{proof}
補題\ref{lem:nodal-core-resolution}により多角形の正次数は環に、多重点曲線の正次数は木に載る。
環は第二像の特異点に写る場合をまず扱う。第二像の各枝がどの木に属するかに応じ、
その木の付け根における第一像の点を$x_i$とする。第二枝では第一束の切断は$x_i$での値に等しい。
各線束の局所枠の定数倍は以下の次元に影響しない。
第二束の特異点で消える切断の空間$B'$は三次元である。環上は第二束の値が同一なので、
全積の環への制限は四次元である。$A_0 B'$の積は環で零であり、各第二枝では
$b_i$に$a(x_i)$を掛けた値である。
$m$枝の二次以上の係数は$4-m$個独立であり、一次係数は全係数非零の一式の核である。
$a$が全ての$x_i$で非零となるよう取れば、この三次元は積の像に入るので$r\ge7$である。
$x_i$が全て同一なら、そこを消す第一束の三次元空間と$B'$の積は全成分で零である
（両射影が定値の鎖でも零になることは各端で評価すればよい）。
実際の積の補正は一次元以下だから、一文字を固定して二次元の評価核が残る。

異なる$x_i$があれば、二つの第一切断の値の比が全枝で同じでないようにできる。
一次係数の超平面に非零の値を対角に掛けた二つの像は異なるので、合わせて全$m$座標を張る。
従って$r=8$である。次数2以上の第二枝があれば、その枝だけに支えを持つ切断と、
$x_i$で消える第一切断三次元の積が零である。全枝が次数1なら、二枝のみに支えを持つ
非零切断と、二つの付け根で消える第一切断二次元以上の積が零である。
どの場合も実際の六式で矩形の零点を生じ、指数成長する。

環の第二像が滑らかな点$y$に落ちる場合も含める。
特異点への各枝を含む木が二本以上なら、それらの特異点の原像を結ぶ道は環を通るため、
その点のファイバーの連結性に反する。ただ一つの木に全ての第二正次数成分があれば、
第一束の値はそれら全てで同一の$x$での評価となる。
$y$で消える第二束の三次元空間と$x$で消える第一束の三次元空間の積は零である。
実際、正次数のない他の木は$y$に写り、唯一の木は$x$に写る。
環上の切断とこの木の切断から再び$r\ge4+3$が従い、補正があっても排除される。
\end{proof}

\begin{corollary}\label{cor:reduced-polygon-window-test}
隣接階数6、反対階数8の両対称テンソルで、両反対核が被約四次完全交叉を切り、
一方は整な節点曲線または可約な節点多角形とする。反転窓が標準ASとなる必要十分条件は、
関係六平面が純テンソルを含まず、両曲線の同型のうち二つの埋込み束の多重次数を一致させるものの
グラフへの制限をテンソル汎関数が経由することである。そのとき順序付き鉛筆データを復元できる。
\end{corollary}
\begin{proof}
必要性は両モデルの被約曲線型を分け、上述の二側の定理と、他側が滑らかな場合には
系\ref{cor:one-smooth-apolar-as}を用いればよい。十分性ではグラフ上で$N=LM$として、
等しい正の多重次数だから$H^0(L)H^0(M)=H^0(N)$である。
与えられた汎関数は$H^0(N)$上の対称対合を誘導し、階数6なので核は二次元である。
この核を$W$とすれば第一窓は$H^0(N)/W$への積で定まり、基点は純テンソル関係を与えるため存在しない。
整曲線または多角形の全次数定理と復元定理が適用できる。
\end{proof}

% !TeX root = ../papers/nc-p3.tex
\section{両側が楕円多重点となる被約極限}
\label{sec:reduced-additive-limits}
\begin{theorem}[被約二側の特殊化]\label{thm:all-reduced-apolar-specializations}
滑らかな楕円鉛筆の特殊化で、両方の極限二次式の鉛筆が被約な四次完全交叉を切るとする。
特殊反転窓が標準ASなら、両曲線は同じ被約曲線として同定でき、二つの埋込み束は等しい
正の多重次数を持ち、六関係はその上の基点自由鉛筆から得られる。従って反対階数は6または8である。
\end{theorem}
\begin{proof}
整な側を持つ場合と節点多角形の場合は前二節による。残る両曲線$X,Y$は唯一の楕円多重点
$x_*,y_*$を持ち、いずれも二本以上の滑らかな有理枝からなる。各枝の次数を$d_i,e_j$、
枝数を$p,q$とする。総次数は各々4で、全て正である。埋込み束の切断空間を$A,B$とする。
各枝への制限は完全線型系である。他枝で一次係数の和を調整できるためである。
特異点で消え、指定一枝だけに支えを持つ切断はその原点で二重に消えるもの全体である。

一般点の同定のグラフを二曲線族の積内で閉包し、特殊ファイバーの被約支持を$G$とする。
両射影は正規な全空間への固有双有理射なのでファイバーは連結である。
被約な特殊曲線をサイクルとして押し出す議論は前節と同じで、各枝を支配する成分はただ一つで
その次数は1である（押し出すのはグラフの特殊ファイバーの重複度付きサイクルである）。
従って$G$の成分は、二つの枝の同型のグラフ、水平枝$X_i\times\{y_i\}$、
垂直枝$\{x_j\}\times Y_j$である。前者を対応枝と呼ぶ。
Segre埋込みで$G$の点評価が張るベクトル空間の次元$r$は積$AB$の被約成分への
制限の階数であり、補題\ref{lem:mixed-rank-loss}により$r\le8$である。

以下で連結な成分の合併のSegre線型包を使う。その両射影のベクトル空間次元が$a,b$なら、
この包のベクトル空間次元は$a+b-1$以上である。実際、一成分の二射影の切断は
射影直線の多項式で、その積の空間は次元の和から1を引いた値以上である
（互いに異なる先頭次数の基底をとり、先頭次数の和集合を見る）。
既存の連結な部分にそれと交わる一成分を加え、両側の増分を$c,d$とする。
$c>0$なら既存部分で消える左切断の新成分への像は$c$次元で、
右切断の制限は$d$に交点での非零評価を足した$d+1$次元以上である。
積は既存部分で消え、増分は$c+d$以上となる。$c=0,d>0$では左右を交換する。
成分を順に加えれば主張を得る。特に$7\le r\le8$である。
$\Phi$を$G$上の積とすれば、$\Phi_x$の階数が1以下の非零文字$x$は許されない。
実際の評価は$r=7$でも階数を高々1しか増やせず、二次元の評価核が指数成長を強制する。
$r=8$なら$\Phi_x$の階数は各非零文字で3以上必要である。

\paragraph{二つの特異点の対がグラフの外にある場合。}
$(x_*,y_*)\notin G$とする。$x_*$のファイバーは正次数の垂直枝を含めず連結なので一点$y_0\ne y_*$、
同様に$y_*$上は一点$x_0\ne x_*$である。全水平枝は$y_0$に、全垂直枝は$x_0$に載る。
対応枝は$x_0$の枝と$y_0$の枝の間の高々一つである。
対応枝がなければ$x_0,y_0$で消える三次元ずつの部分の積は零で、補正後も指数成長する。
一つあればその二次数を$d,e$、対応射影直線上で$x_*,y_*$に対応する二つの異なる点を$P,Q$とする。
対応枝上の積は全$P_{d+e}$を張り、それと独立に水平枝のみ、垂直枝のみへの
$3-d,3-e$次元の像を持つ（対応枝への制限核と、相手での単元の積）。ここまでで七次元である。

さらに対応枝上の二切断を$f,g$と書くと、他の水平枝の一次係数の重み付き和から
$f'(P)g(P)$を、他の垂直枝から$f(Q)g'(Q)$を読み取れる。これは$x_*,y_*$での
一次の貼り合わせ式によるもので、自明化による零階項は対応枝の積の値である。
この二座標は上記の$3-d,3-e$次元の像では消える。対応枝の乗法核に制限すると、
二座標は対角の括弧で割った$(d-1,e-1)$形式の$(P,P),(Q,Q)$での評価に比例する。
非零の定数倍は各点の座標の選択による。従って$d+e\ge3$なら二条件は独立で$r\ge9$となり矛盾する。
$d=e=1$なら$r=8$であるが、対応枝への二つの制限核は各二次元で相互の積が零となる。
これも標準ASを排除する。

\paragraph{二つの特異点の対がグラフ上にある場合。}
各対応枝は根（それぞれの特異点の原像）を対応させる。
実際、$x_*$上に対応枝の根以外の$Y_j$の点があれば、その枝には別の垂直成分がないので、
この点は$y_*$と同じファイバー内で孤立してしまう。
また$y_i\ne y_*$の水平枝は、$x_*$のファイバーの連結性により
$x_*$上の垂直枝の根以外の点に付く。逆も同じである。
従って対応枝以外は中心$(x_*,y_*)$を通る水平または垂直な枝と、
その根以外に付く反対向きの末端枝だけである。付属枝を持たない中心枝自身も末端に数える。

必要な線型包の計算を与える。$e,f$を中心での両評価ベクトルとし、
各々一般の第一・第二切断で値が1となるように点評価を規格化する。
その切断の核を$U,V$とし、各枝で中心からの差を張る部分空間を$U_i,V_j$とする。
その次元は$d_i,e_j$である。$U_i$の直和から$U$への写像の唯一の核を$(h_i)_i$とし、
$V_j$では$(k_j)_j$とする。$h_i,k_j$は各々零でない根の接方向である。
これは各枝の値を一致させた上で一次係数に一式だけを課す表示による。
正規化に用いる切断は全ての付け根で非零、また各対応枝上で両側の零因子が交わらないように選べる。

点評価の包から$U\otimes V$の成分を捨て、中心の直線も法として$U\oplus V$に射影する。
各対応枝の中心を法とした$d_i+e_j$次元の包は$U_i\oplus V_j$へ同型になる。
実際、二つの切断を規格化切断で割った有理関数の空間は、極因子が互いに素なので定数だけを共有する。
従って$U_i\oplus0$のベクトル$h_i$をこの包に持ち上げたときの
$U_i\otimes V_j$成分$D_{ij}$が一意に決まる。対応枝の根での一次微分は混合成分が零で、
投影が$(h_i,c_{ij}k_j)$となる（$c_{ij}\ne0$）。従って$0\oplus k_j$の持ち上げの混合成分は
$-c_{ij}^{-1}D_{ij}$である。
$D_{ij}$は非零である。さらに$d_i,e_j$の少なくとも一つが2以上なら、
$D_{ij}$は$h_i\otimes k_j$の倍数ではない。
もしその倍数なら$h_i\otimes(f+z k_j)$という純テンソルが対応枝の線型包に入る。
積$P_d\otimes P_e\to P_{d+e}$の双対像の純テンソルは一つの点での二つの評価の積だけである。
実際、単項式座標の行列は$(a_{i+j})$で、階数1なら相隣る二行二列の行列式と全二次小行列式が消える。
非零な分解$u_i v_j$に対し$u_{i+1}v_j=u_i v_{j+1}$なので、
二列のシフトが非零なら両列は同じ公比の等比数列、端のシフトが零なら両方が同じ端の評価である。
これで純テンソルの主張が従う。次数2以上の有理正規曲線は根の接線と根以外で交わらない。
$h_i$は根の評価とは異なり、$f$だけを取ることも対応枝の同一パラメータに反する。
$d_i=e_j=1$なら$D_{ij}$は非零の$h_i\otimes k_j$倍で、対応する$f+z k_j$は
その直線上の根以外の点を表す。

水平枝の$U_i$方向の増分は$u\otimes(f+v_i)$であり、$f+v_i$はその枝の第二像の点である。
垂直枝の$V_j$方向は$(e+u_j)\otimes v$で同様である。従って全点評価の包は
中心と各枝のこれらの増分で張られる。付け根の評価が増分の基点として必要な場合も、
それは中心枝の線型包内にある。$U\oplus V$への投影は全射で、直和の添字空間での関係は
先の$(h_i)_i,(k_j)_j$の二つである。ゆえに$r=7+\dim\langle H,K\rangle$で、ここで
\[
 H=\sum_{(i,j)\ {\rm 対応}}D_{ij}+\sum_{i\ {\rm 水平}}h_i\otimes v_i,\qquad
 K=-\sum_{(i,j)\ {\rm 対応}}c_{ij}^{-1}D_{ij}+\sum_{j\ {\rm 垂直}}u_j\otimes k_j .
\]

対応枝がなければ$H=0$は全$v_i$が等しいことと同値である。
$h_i$の関係が全係数1の関係だけだからである。$K=0$も同様で、
$r=7$なら全水平枝が同じ点に、全垂直枝も同じ点に載る。両側でその点において消える
三次元ずつの切断の積が零となり、一次元の補正を許しても排除される。
対応枝があれば$r=8$となることも確認する。$d_i\ge2$または$e_j\ge2$の対応枝を除いた
連結な残りの包は、両射影の次元下界により少なくとも$9-d_i-e_j$次元である。
その枝だけで二重に消える側の切断（根で二重に消え、他枝上では零）と相手の全切断の積は、
残りで零となり、その枝で$P_{d_i+e_j-2}$の次元$d_i+e_j-1$を与える。
従って$r\ge8$である。全対応枝が二次数$(1,1)$なら、$H$は$\sum_i h_i\otimes v'_i$の形で、
対応枝の$f+v'_i$は対応する第二直線の根以外の点である。水平枝の第二像は根または
別の垂直枝上の点なのでこれと一致しない。他の対応枝の根以外の点も一致しない。
$p\ge2$より全$v'_i$は一致せず$H\ne0$である。この場合にも$r=8$である。

従って残りでは、実際の評価は$\Phi$を経由する。末端の水平枝の次数が2以上なら、
その枝だけに支えを持つ非零切断での評価は第二像の一点での値だけを通り、階数は1以下である。
また末端の水平枝が二つあれば、その二枝だけに支えを持つ非零切断がある
（根の値を零、一次係数の重み付き和を零とする）。その評価は高々二点の値を通る。
末端枝上の切断は、他の垂直枝が水平枝内の根以外に付いていないため、そのような垂直枝上に
追加の支持を持つことはない。従って末端水平枝は高々一本で次数1である。末端垂直枝も同様である。
付属枝を持つ中心の水平枝には高々一本の次数1の垂直枝が付く。この水平枝の次数が2以上なら、
その枝だけに支えを持つ第一切断を選べる。積の値は付属直線上の第二切断の全二座標で決まり、
水平枝上で使う定数もその直線の根での値である。評価が階数2以下となり矛盾する。
従って全ての非対応枝は次数1で、非対応の各群は一本の中心枝または一本ずつの水平・垂直枝の組である。

対応枝が一本もなければ各群が一つの末端を持つため高々二群、各群高々二枝である。
$p,q\ge2$なら$p=q=2$で各側二本の次数1の和が4に達せず矛盾する。
少なくとも一つ対応枝があるとする。$H,K$に残る非対応群の非零項は、根以外の付け根を持つ
水平・垂直枝の組についての非零の$h_i\otimes k_j$倍で、$H,K$の一方にだけ現れる。
これは両枝が直線のためである。これらの$(i,j)$と対応枝の$(i,j)$は添字を共有しない部分的な対応である。
このような対応の各項$D_{ij}$または$h_i\otimes k_j$は独立である。
実際、$D_{ij}$に$kh_i\otimes k_j$の外の成分があれば、左または右で
根の全接方向の包を法とした成分が非零になる。$U_i$と根の全接方向の包の交わりは$kh_i$であり、
それ以外の$U_i$の方向はその枝に独立だから、その係数は関係中で零である。右側も同様である。
残った$h_i\otimes k_j$では、左右どちらかが全添字の真部分ならその側の方向は独立である。
全単射かつ$p=q\ge3$なら最初の$p-1$方向を基底に取ると最後の積だけが非対角成分を持つ。
$p=q=2$で両項とも単純な接方向の積になるには双方の全枝が次数1でなければならず、
各側の総次数4に反する。独立性が示された。

$r\le8$のため$H,K$は従属しなければならない。対応項があるので両方非零であり、
片方にしか現れない付属枝の項は存在できない。従って非対応枝があれば全て中心を通る。
さらに対応枝全てで$c_{ij}$は同じスカラーである。
これにより、全枝上の関数の値を一致させ、一次係数の重み付き和を一式で課す
一つの楕円多重点曲線$Z$を構成できる。詳しくは各第一枝で$h_i$を、各第二枝で
$k_j$を定める接パラメータを使い、対応枝では同型の微分が共通の$c_{ij}$である。
第二側の全接パラメータを同じ定数で取り直せば、一方だけの枝も含めて同じ和零の条件に揃う。
$h_i,k_j$の規格化が構造層のこの条件を与えるのは、規格化切断で割った局所正則関数の
一次係数の唯一の関係を取ったためである。対応枝は同一とし、水平または垂直枝を一方で
縮める二射を得る。根での関数の条件が合うのでこれらはスキームの射である。

この$Z$上の埋込み束の引き戻しの切断は元の四座標全体で、
積束の切断は八次元である。後者は正規化での共通値と一次係数の一式から計算できる。
$\Phi$が階数8なのでこの積の切断を全て張り、実際の六式はその二次元鉛筆の表示となる。
基点は純テンソルを与えるから存在しない。命題\ref{prop:semipositive-pencil-balance}により
二つの多重次数は一致し、非対応枝は存在できない。従って$X,Y$は同じ$Z$であり、
多重点の全次数定理と同型復元が適用できる。
\end{proof}

\begin{corollary}[階数8の被約反対核に対する判定]\label{cor:all-reduced-apolar-recognition}
隣接階数6・反対階数8の両対称テンソルの二つの反対核が被約四次完全交叉を定めるとする。
標準ASとなる必要十分条件は、第一窓に純テンソル関係がなく、両曲線の同型のうち
二つの埋込み束の多重次数を一致させるもののグラフへの制限をテンソル汎関数が経由することである。
このとき同定と鉛筆が復元される。反対階数7の標準AS特殊化では、二つの極限鉛筆を
ともに被約四次完全交叉に保つことはできない。
\end{corollary}
\begin{proof}
全両対称AS窓の閉包定理と前定理で必要性を得る。十分性では同定した曲線上で
二つの束の積は全切断を張るので、階数6の隣接対合の核は二次元である。
この鉛筆の基点は純テンソル関係を与えるため存在しない。
各被約曲線型での全次数定理と復元定理を適用する。最後は前定理の階数の結論である。
\end{proof}

% !TeX root = ../papers/nc-p3.tex
\section{重平面を持つ反対核の全次数判定}
\label{sec:square-apolar-ore}
$E^*,F^*$を二つの点座標空間とし、隣接階数6の両対称テンソルを
六次元非退化対称空間への全射$\beta:E^*\otimes F^*\to\mathcal H$で表す。
$\beta_{xf}=\beta(x,f)$と書く。両対称条件は
$\langle\beta_{xf},\beta_{yg}\rangle=\langle\beta_{xg},\beta_{yf}\rangle$である。
$K_E,K_F$を反対切断の左右の対称二次核とする。

\begin{lemma}[因子の相互対応]\label{lem:reciprocal-apolar-factors}
両側の全ての非零文字で$\beta$の評価階数が3以上とする。
$xx'\in K_E$を満たす射影因子の順序付き対のスキームと、
$ff'\in K_F$に対する同様のスキームは因子交換と可換な同型を持つ。
対応は$f=\ker\beta(x,-)$、$f'=\ker\beta(x',-)$である。
相異なる因子の対が標準AS窓にあれば、$U_x=\operatorname{im}\beta(x,-)$と
$U_{x'}$の交わりは高々一次元であり、一次元ならその生成元の関係行列は階数2である。
\end{lemma}
\begin{proof}
$xx'\in K_E$は両対称性により$U_x\perp U_{x'}$と同値である。
両方の次元が3以上なので互いに直交補である。
$V_f=\operatorname{im}\beta(-,f)$も$U_x$と直交するため$V_f=U_{x'}$、
同じく$V_{f'}=U_x$となる。従って$ff'\in K_F$で、逆は同じ構成である。
族上でも単元である三次小行列式の近傍を取れば、他側の三本の独立な像の直交補は
階数3の部分束で、評価像はそれに等しい。核直線の形成は基底変換と可換である。
交わりの元と$\beta$を対合して得る関係は$x,x'$両行、$f,f'$両列で消える。
従って二行二列行列の空間に入り、$\beta$の全射性からこの写像は単射である。
その中の二次元部分には行列式零の非零元がある。階数1の関係は反転窓の次の位置で
転置され、定理\ref{thm:no-consecutive-pure-relations}に反する。
評価階数2以下も一文字と二文字への交互商で全関係が消えるためASではない。
\end{proof}

$x_0^2\in K_E$とする。$U_{x_0}$は等方なので階数3以下であり、
標準ASの候補では相手の核$f_0$での評価とともに階数3でなければならない。
この二つの階数だけを以後の正規形の仮定とする。
両軸の像は同じ極大等方三平面$D$である。双曲分解$\mathcal H=D\oplus D^*$を取り、
両文字空間の三次元の補空間を軸評価によって$D$と同定すると
\begin{equation}\label{eq:double-plane-beta}
 \beta((r,x),(s,y))=(r y+s x+h(x,y),\ C(x,y,-)).
\end{equation}
軸を一つ含む両対称式により$C_{aij}$は三添字で対称であり、
全射性により三つの極形式$C_a=(C_{aij})$は独立である。
残りの条件は$h_{ij}=h(e_i,e_j)$、$C_{ij}=C(e_i,e_j,-)$として
\begin{equation}\label{eq:double-plane-flat}
 h_{ij}\cdot C_{kl}+C_{ij}\cdot h_{kl}
 =h_{il}\cdot C_{kj}+C_{il}\cdot h_{kj}
\end{equation}
で尽くされる。$I=\{S=S^{\mathsf t}:\sum S_{ij}C_{ij}=0\}$は
三次元二次系である。ここで$S$は$D$上の座標を係数とする対称テンソルで、
二次式としては$D^*$上で評価する。

\begin{theorem}[重平面での正則性と反対階数]\label{thm:double-plane-ore}
式\eqref{eq:double-plane-beta}の反転窓が標準Koszul AS正則である必要十分条件は、
$I$の三式が射影平面上に共通零点を持たないことである。
正則な場合の反対階数は6または8であり、階数8の両核は
軸同定の適切な補空間において同じ重平面付き四次完全交叉鉛筆になる。
\end{theorem}
\begin{proof}
$I$の共通射影零点$v$は$v v^{\mathsf t}\in\langle C_a\rangle$と同値で、
六関係の中に純テンソルを与える。以下$I$は正則列とする。
三文字$X_i$と三関係$r_a=\sum C_{aij}X_iX_j$の代数を$B$と書く。
二次双対は$\operatorname{Sym}(D)/(I)$で、三つの二次式の可換完全交叉である。
補題\ref{lem:ternary-clifford}のTate解消と完全対合の議論から$B$はKoszul AS三次元である。

第四文字を$Z$、$h_i=\sum h^i_{kl}X_kX_l$とし、$h_i^{\rm op}$は積の順を逆にした式とする。
六関係は$r_a$と$ZX_i=-X_iZ-h_i$で、奇位置では$h_i^{\rm op}$を使う。
$Zr_a$を右に移動すると$r_a Z$のほかに
$\sum C_{aij}(X_i h_j^{\rm op}-h_i X_j)$が残る。
$X_iX_kX_l$の係数は式\eqref{eq:double-plane-flat}から
\[
 C_{ia}\cdot h_{lk}-h_{ik}\cdot C_{la}
 = C_{ik}\cdot h_{la}-h_{ia}\cdot C_{lk}
\]
で、右辺は$r_bX_l,X_i r_b$の線型結合の係数である。転置した$h$でも同じ式が成立する。
従って$Z$の交換則は偶奇二頂点で折り畳んだ$B$の上のOre拡大になる。
実際、自己同型は両頂点を交換し各矢印を他頂点の負の矢印に移す。
補正項は頂点では零、矢印では出発位置に従って$-h_i$または$-h_i^{\rm op}$であり、
上の計算が歪導分の全ての定義関係への降下を与える。
Ore基底は各位置で$B$の基底に右から$Z^m$を掛けたものとなる。
$Z$の個数で濾過すれば随伴次数は符号自己同型による歪多項式拡大である。
定理\ref{thm:algebraic-nine}の証明と同じ内部次数ごとのbar濾過と
二次双対の完全対合の持ち上げから、元の窓も標準14641型Koszul AS正則となる。

次に$h=h_s+h_a$を対称・交代部分に分ける。
式\eqref{eq:double-plane-flat}は転置で保たれ各部分で成立する。
体積形式を選び$h^a_{ij,\mathrm{alt}}=\sum_b\epsilon_{ijb}P_{ba}$と書く。
左右補空間の取り直しで$h_{ij}$に$u_i e_j+v_j e_i$を足せる。
その交代部分は$P$の歪対称部分を任意に変えるので$P=P^{\mathsf t}$にできる。
順序付き補添字$(23),(31),(12)$で式\eqref{eq:double-plane-flat}の左辺と右辺の差を
二小行列式の行列に並べると、交代部分の$(u,v)$成分は
$\sum_a\epsilon_{uva}\sum_{ij}P_{ij}C_{aij}$である。
従って$P\in I$になる。対称部分だけを残したテンソル$w_s$は
四添字で完全対称であり、三つの極形式の独立性から隣接・反対階数がともに6である。

反対行列では軸平方の行列成分が消える。軸と一文字の積からなる三行と、
三文字の対称二乗の六列との長方形は極形式で、転置位置も同様である。
よって反対階数は6に$I\times I$に残る対合の階数を足したものである。
$w_s$の寄与はその上で零である。残る交代部分は$S,T$を対称行列として
\[
 2\sum_{ijb a k l}\epsilon_{ijb}P_{ba}S_{ik}T_{jl}C_{akl}.
\]
これは三次式$\det(Pz,Sz,Tz)$に$C$を縮約したものの二倍であり、
$P,S,T$に交代的である。$I$の基底を入れると行列式は
その三二次式のJacobi行列式の八分の一で、完全交叉のsocleを生成する。
実際、二次式のKoszul解消から文字のKoszul解消への比較写像は一次係数行列から作られ、
最上次数でこの行列式となる。三次元多項式環上で有限長加群の
$\operatorname{Ext}^3$は完全反変双対なので、剰余環から$k$への全射はこの最上次数で
単射を与える。$C$は$I D$を消し、非零なので一次元のsocleに非零である。
従って$P=0$なら反対階数6、$P\ne0$なら反対階数8で、残る核直線は$kP$である。

$P\ne0$では、$w_s$については$I$の各元が軸と一次式の積を加えることで
軸平方を法として一意な核元に持ち上がる。
交代部分は軸の行列成分で零であり、上の行列式は一つの引数も$P$なら恒等的に零である。
従って両核は同じ$\langle x_0^2,P+x_0\ell\rangle$となる。
$P\ne0$は軸を含まない三文字の二次式なので二式は互いに素である。
\end{proof}

\begin{corollary}[階数8の順序付き同型判定]\label{cor:double-plane-isom}
重平面チャートの階数8の標準窓の間で、指定した軸を保つ添字保持同型の存在は
次の係数比較と同値である。$a,b,c\in k^\times$、可逆な三行行列$A,B,G$、
二つの一次式$u,v$と歪対称写像$J:D^*\to D$が存在して
\[
\begin{aligned}
 G&=bA=aB, & C'(Ax,By,-)&=cG^{-\mathsf t}C(x,y,-),\\
 h'(Ax,By)&=G\big(h(x,y)+JC(x,y,-)\big)-u(x)By-v(y)Ax .
\end{aligned}
\]
軸を指定しないときは標的鉛筆内の全ての重平面（高々二個）とその対応軸を試せばよい。
\end{corollary}
\begin{proof}
反対階数8では定理\ref{thm:elliptic-pencil-reconstruction}の七次元上界により、
反対二位置の基底変換は射影的に一致する。Gram因子分解の全射性・非退化性により
テンソル直線を保つことはGram標的の相似変換が存在することと同値である。
基底変換を$(r,x)\mapsto (ar+u(x),Ax)$、$(s,y)\mapsto(bs+v(y),By)$と書くと、
軸での値から$G=bA=aB$で、相似変換は
$(d,\xi)\mapsto(G(d+J\xi),cG^{-\mathsf t}\xi)$となる。
残る値の比較が表示式であり、逆に全ての関係列を移す。
独立な二平方を張る鉛筆にはその二つしか平方がなく、右軸は
補題\ref{lem:reciprocal-apolar-factors}で決まる。
\end{proof}

\begin{corollary}[階数7に残る可約因子]\label{cor:rank-seven-distinct-factors}
両対称の標準AS窓で反対階数が7ならば、両反対核の中に相異なる二平面の和があり、
いずれの核にも重平面はない。各因子対の二軸の交わりは高々一次元である。
\end{corollary}
\begin{proof}
楕円閉包で両反対鉛筆の極限を取る。系\ref{cor:all-reduced-apolar-recognition}から、
少なくとも片側で二式の交叉は被約四次完全交叉ではない。共通因子を持つなら各式は可約で、
非被約完全交叉なら定理\ref{thm:nonreduced-quartic-pencils}の平面対の議論で可約な元がある。
その元は反対核に入る。前定理で平方の場合は消え、相手側と次元上界は
補題\ref{lem:reciprocal-apolar-factors}による。
\end{proof}

% !TeX root = ../papers/nc-p3.tex
\section{異なる平面因子の一意性と横断する二軸}
\label{sec:distinct-factor-rigidity}
記号は第\ref{sec:square-apolar-ore}節を引き継ぐ。

\begin{theorem}[各因子の相手の一意性]\label{thm:unique-linear-factor-partner}
両対称の標準AS反転窓では全ての$0\ne x\in E^*$に対して
$\dim\{w:xw\in K_E\}\le1$であり、右側でも同じである。
従って反対階数8では両反対鉛筆が四次完全交叉を定める。
反対階数を問わず、添字保持同型の反対二位置での基底変換は射影的に一致する。
\end{theorem}
\begin{proof}
$X=\{w:xw\in K_E\}$とし、$f=\ker\beta(x,-)$を取る（$X\ne0$の場合）。
補題\ref{lem:reciprocal-apolar-factors}は$\mathbb P(X)$から
$\mathbb P(Y)$、$Y=\{v:fv\in K_F\}$への同型を与える。
$U=U_x$、$V=U^\perp$と置くと、全ての$0\ne w\in X,0\ne v\in Y$で
$U_w=V$、$V_v=U$である。$\dim X\ge2$なら$x$とは異なる因子を選べるので
$R=U\cap V$は高々一次元である。$\beta(X,Y)\subset R$で、
$w$の$Y$上の核は丁度その対応直線であるから$\dim X=\dim Y=2$、$\dim R=1$となる。
一方、各$0\ne v\in Y$で$V_v=U$なので
$E^*/X\to U/R$は同型である。これは二次元空間間の二行行列の線型鉛筆で、
行列式が射影直線上で消える点が存在し矛盾する。
階数8の独立な二二次式に共通因子があれば、その因子には二次元の相手がある。

次に、反対切断の像の全ての対称形式に関して自己共役な$g$は
$(gv)w-v(gw)\in K_E$を全ての$v,w$で満たす。
固有値$\lambda$の一つを取ると$\ker(g-\lambda)\cdot\operatorname{im}(g-\lambda)\subset K_E$
である。固有ベクトルの相手が高々一次元なので$\operatorname{rank}(g-\lambda)\le1$、
階数1なら像の生成元に三次元の相手があり矛盾する。従って$g$はスカラーである。
反対二位置で別の基底変換をして再び対称な像を得る場合、それらの相対変換は
全ての像形式に関し自己共役となる。もう一側にも同じ議論を適用する。
\end{proof}

\begin{corollary}[スキームとしての因子曲面の切断]
$\theta_x:w\mapsto xw\bmod K_E$の階数3以下の小行列式スキーム上で、
因子交換は$[x]\mapsto[\ker\theta_x]$という正則な対合になる。
左右のこのスキームは同型で、対合とも可換である。
\end{corollary}
\begin{proof}
階数2以下の点はないので三次小行列式の一つを可逆にすると、残る小行列式は
核直線が存在するためのスキーム論的方程式である。従ってその上の核直線のグラフが
順序付き因子スキームであり、補題\ref{lem:reciprocal-apolar-factors}を適用できる。
\end{proof}

\begin{theorem}[横断する二因子での階数の飛び]\label{thm:transverse-product-even-rank}
相異なる因子の対$x_0 x_1\in K_E$と対応する$f_0 f_1\in K_F$を取り、
四評価の階数を3とする。$U=U_{x_0}=V_{f_1}$、$V=U_{x_1}=V_{f_0}$が
相補的なら、反対階数は6または8である。特に両対称の階数7の標準AS候補では
\emph{全ての}因子対が一次元で交わり、その根基は可逆な二行二列関係を与える。
\end{theorem}
\begin{proof}
$U,V$は直交し各計量は非退化である。$e=\beta(x_0,f_1)\ne0$、
$d=\beta(x_1,f_0)\ne0$とする。全双線型写像は二つの積の和として
\[
 \beta(x,f)=(a(x)*_U b(f),\ m(x)*_V n(f))
\]
と書ける。ここで$a=\beta(-,f_1), b=\beta(x_0,-)$は$U$全体への写像であり、
それぞれ$kx_1,kf_0$を核とする。直交射影した$\beta$はこの二商を経由するので、
ある積$*_U$が一意に定まり、$e$は両側の単位である。同様に
$m=\beta(-,f_0),n=\beta(x_1,-)$と$V$の単位$d$を得る。
両対称式に$x_0$または$f_1$を代入すると、三形式
$\langle u*v,w\rangle_U$が全対称であることが分かる。
従って積は可換で計量は不変である。$V$でも同様で、結合性は仮定しない。

$D_U=U/ke$とし
\[
 T(p,q,r,s)=\langle p*q,r*s\rangle_U,\quad
 F(p,q;r,s)=T(p,r,q,s)-T(p,s,q,r)
\]
と書く。$F$は各対で交代し単位を入れると消えるため、
$D_U$の面積形式$\omega$を選べば$F=\kappa_U\omega(p,q)\omega(r,s)$である。
同じく$V/kd$で$\kappa_V$を取る。
$X=E^*/\langle x_0,x_1\rangle$、$Y=F^*/\langle f_0,f_1\rangle$からの誘導同型を
$A:X\to D_U,B:Y\to D_U,M:X\to D_V,N:Y\to D_V$と書く。
面積座標で両対称式は
\[
 \kappa_U\det A\det B+\kappa_V\det M\det N=0
\]
の一式となる。

反対行列から零行$x_0x_1$と零列$f_0f_1$を除く。
$x_0$の一次倍の三行と$f_1$の一次倍の三列は$U$の計量を与え、
$x_1$と$f_0$の三行三列は$V$の計量を与える。反対の組のブロックは零である。
この六行六列を消去した残差は$\operatorname{Sym}^2 X\times\operatorname{Sym}^2 Y$上にある。
一因子$U$の寄与は
\[
 \tfrac12\{T(p,r,q,s)+T(p,s,q,r)\}-T(p,q,r,s)
\]
の$(A,A;B,B)$による引き戻しである。実際、六行の消去は各積$p*q$を
単位を持つ三行から計量を介して差し引くことである。
同式は$-\tfrac12\{F(p,r;q,s)+F(p,s;q,r)\}$であり、
平方$(p,p;r,r)$で$-\kappa_U\omega(p,r)^2$となる。
二次元空間の対称二乗上の非退化対合$J$を
$J(p^2,r^2)=\omega(p,r)^2$で定めれば、残差はこの$J$の二つの引き戻しの線型結合である。

$R(g)=\operatorname{Sym}^2(g)/\det g$は$J$を保つ。二つの寄与の係数は
前の行列式の一式により逆符号なので、どちらも零なら残差は零である。
さもなくば残差の階数は$R(A^{-1}B)-R(M^{-1}N)$の階数である
（$X,Y,D_U,D_V$の基底を取って比較する）。
$R(g)$は面積形式による同定で無跡二行行列への共役作用になる。
二行非スカラー行列の中心化代数はその行列と恒等行列が張り、無跡部分は一次元である。
従って二つの$R$の差の階数は射影変換が等しければ0、異なれば2である。
六つのピボットを戻して結論を得る。最後は系\ref{cor:rank-seven-distinct-factors}による。
\end{proof}

\begin{lemma}[根基直線上の三次形式]\label{lem:product-radical-cubic}
同じ四軸で$U\cap V=kr$とし、根基の関係
$t(\bar x,\bar f)=\langle r,\beta(x,f)\rangle$は$X\times Y$上非退化とする。
$U$から定まる同型$J_U:X\to Y$を、$a=\beta(-,f_1)$と$b=\beta(x_0,-)$の
$U/k\beta(x_0,f_1)$への誘導同型により$\bar b^{-1}\bar a$と置く。
$V$からも同様に$J_V=\bar n^{-1}\bar m$を作る。
$J_U,J_V$が射影的に異なれば反対階数は8以上である。
射影的に等しければ階数は6または8で、6の場合には完全対称なテンソルに同値である。
\end{lemma}
\begin{proof}
$U$の元$u,b_1,b_2$に対して
$S_U(u,b_1,b_2)=\langle u,\beta(x,f)\rangle$を$a(x)=b_1,b(f)=b_2$で定める。
二つの持ち上げの曖昧さの像は$V=U^\perp$に入るため良定義である。
$x_0,f_1$を両対称式に代入すれば、第一と第三、第一と第二の引数交換を得る。
従って$S_U$は完全対称で、単位$e=\beta(x_0,f_1)$を入れると$U$の計量となる。
$S_U(r,-,-)$は$U/ke$上に降り、その$(\bar a,\bar b)$による引き戻しが$t$である。
$V$でも同様である。

前定理の六行六列のブロックは二つの退化計量で、各々階数2である。
四ピボットを消去する。$U$の根基行は$x_0 x_r$、$a(x_r)=r$で、
$\operatorname{Sym}^2Y$に対する残りの行は$S_U(r,\bar b(-),\bar b(-))$、
すなわち$t(J_U^{-1}-, -)$となる。根基列は$t(-,J_U-)$である。
$V$の根基行・列では$J_V$を使う。従って残る二行は$J_U,J_V$が射影的に
異なるとき独立で、二列も同時に独立である。根基行と根基列の交叉ブロックは零だから
全階数は$4+2+2$以上である。

射影的に等しければ$g:E^*\to F^*$を選び、$b g=c_U a,n g=c_V m$とできる
（$c_U,c_V\ne0$）。実際、商$X\to Y$を$c_UJ_U=c_VJ_V$と取り、
$x_0,x_1$を$c_U f_1,c_V f_0$へ送る。商の持ち上げは二つの係数を独立に直せ、
$b,n$の各単位方向のずれをちょうど消す。
この同定で$\beta(x,gy)$と$\beta(y,gx)$の差は$S_U,S_V$の対称性により
$U+V$と直交し、従って$kr$に入る。どちらかの文字が軸にあれば差は零である。
ゆえに交代部分は$r\eta$、$\eta$は二次元商$X$上の交代形式である。
$t_g(x,y)=t(x,gy)$もそこで対称かつ非退化である。

対称部分$\beta_s$だけのGramテンソル$w_s$は完全対称である。実際、
$\langle r,r\rangle=0$なので元のテンソルと同時転置との平均がそのGramテンソルであり、
両対称式を保つ。$\beta_s$の像は軸から$U+V$を含み、
$t_g\ne0$なので全六次元を張る。従って$w_s$の反対階数は6である。
残る交代寄与は$\operatorname{Sym}^2 X$同士にのみ入り、平方での値は
$2\eta(p,q)t_g(p,q)$である。四つの計量ピボットの後の根基行と根基列は
それぞれ$t_g$の非零倍一つに縮約される。従ってさらに二ピボットを取った後、
変化するのは$\ker(t_g:\operatorname{Sym}^2 X\to k)$の二次元上の交代対合だけである。
具体的にはこの核の各元を四つの軸ピボットと根基行の一つを使って持ち上げられ、
交代寄与は持ち上げで変わらない。$w_s$の残差は階数6から零である。
$t_g$を単位行列とする基底では、二つの無跡元$e_1^2-e_2^2,e_1e_2$の対合が
$2\eta(e_1,e_2)$になる。従って$\eta=0$で階数6、そうでなければ8である。
\end{proof}

\begin{theorem}[両対称AS窓の階数二分法]\label{thm:bisymmetric-six-or-eight}
標準14641型の両対称AS反転窓の反対階数は6または8であり、6の場合は
完全対称な窓に同値である。従ってこの類の低階数部分は対称完全交叉の分類に帰着し、
残る階数8の両核はともに四次完全交叉を定める。
\end{theorem}
\begin{proof}
楕円閉包定理で階数は8以下である。階数7なら系\ref{cor:rank-seven-distinct-factors}から
相異なる因子があり、交わりの次元は0または1である。定理\ref{thm:transverse-product-even-rank}と
補題\ref{lem:product-radical-cubic}でいずれも不可能である。
階数6以下なら$\mathbb P(K_E)$は次元3以上で、$\mathbb P^9$内の
次元6の階数2以下の対称二次式の多様体と交わる。
重平面なら定理\ref{thm:double-plane-ore}の$P=0$の場合で完全対称となる。
相異なる因子なら前二命題により階数は6以上であり、根基直線の場合の等号は既に示した。

横断する場合には$\kappa_U=0$なら$U$の可換積が結合的となる。
実際、$F=0$と計量の不変性・非退化性から二通りの三重積は一致する。
有限次元可換単位代数の指標$\chi:U\to k$を取ると、$\chi$を$V$上で零に延ばして
$\beta$と合成した関係は$\chi(a(x))\chi(b(f))$という非零純テンソルで、AS性に反する。
従って$\kappa_U\kappa_V\ne0$で、階数6なら前定理の計算から
$B^{-1}A$と$N^{-1}M$は射影的に一致する。前補題と同じ持ち上げで
$b g=c_U a,n g=c_V m$と取れる。両方の積は可換なので$\beta(x,gy)$は対称となり、
Gramテンソルは完全対称になる。最後の完全交叉性は定理\ref{thm:unique-linear-factor-partner}による。
\end{proof}

% !TeX root = ../papers/nc-p3.tex
\section{横断する因子対の自由積表示と正則性}
\label{sec:transverse-factor-amalgam}
\begin{theorem}[横断する二平面での完全判定]\label{thm:transverse-factor-amalgam}
定理\ref{thm:transverse-product-even-rank}の四評価階数3・横断性の下で、
標準Koszul AS正則となる必要十分条件は$\kappa_U\ne0$である。
これは$U$の可換単位積の非結合性と同値で、$V$で検査しても同じである。
この条件は二つの反対鉛筆の被約性を問わない。
\end{theorem}
\begin{proof}
$\kappa_U=0$なら前節の指標の関係が純テンソルになる。
以後両方の$\kappa$が非零とする。
一つの積の三つの対称関係$\lambda(p*q)$ ($\lambda\in U^*$)が
階数1の形式を含まないことを示す。もし含めば単位$e$での縮約は
非零の$\lambda$だから、規格化して$\lambda(p*q)=\lambda(p)\lambda(q)$となる。
これは全ての結合子を消す。$L_p$を左乗法とすると
$\langle [L_p,L_q]r,s\rangle$は前節の面積形式の積に符号を掛けたものである。
$\bar p,\bar q$が独立ならこの交換子の階数は2、核は$ke$である。
$[L_p,L_q]r$も指標で消えるから$\lambda$の計量双対は$ke$に入る。
ところが$\langle e,p*q\rangle$は非退化な計量そのもので矛盾する。
従って三つの対称関係を課す三生成代数$B_U$は三変数二次完全交叉の
Koszul双対で、AS三次元となる。$B_V$も同じである。

$B_U$の三文字のうち二文字空間$W$を単位を消す線型形式全体とする。
三関係の線型結合で$W\otimes W$に入るものは単位評価で消えるので零である。
従って全対称積の像は$\operatorname{Sym}^2 W$の像と同じ三次元である。
これらの像は中心的である。実際、$\operatorname{Sym}^2 W$の基底を
独立な三つの次数2の中心変数に置き、全ての反交換子をこれらの対応する線型結合とする
Clifford代数へ$B_U$は全射する。Cliffordの順序付き平方自由基底から次数級数は
$(1+t)^3/(1-t^2)^3=(1-t)^{-3}$となり、$B_U$の級数と等しいため同型である。

二文字の部分代数$D$の三次関係は各語の最初と最後の文字の交換である。
実際、二つの平方は中心的で、二文字の反交換子も中心的である。
この交換だけで、各語の偶位置・奇位置を別々に整列できる。長さ$2m,2m+1$の次元は
それぞれ$(m+1)^2,(m+1)(m+2)$となる（交換関係は各軌道の語を同一視するだけである）。
Clifford基底で見た二文字部分も三つの中心変数上階数4で同じ級数なので、
上の三次関係で尽くされる。$B_U$は$D$上左右に自由で基底は$1,Z_U$である。
残り一文字$Z_U$をClifford基底で先頭または末尾に置けばよい。

第一・第二位置での二文字空間はそれぞれ$X^*,Y^*$である。
$U$表示と$V$表示の二文字座標は違っても、三次の両端を交換する関係空間は
両端に同じ、中央に独立な基底変換で保たれる。従って両側の$D$は同じ添字付き代数である。
六関係は$B_U,B_V$の各三関係の合併だから全代数は$A=B_U *_D B_V$となる。
各位置で$D$の右係数を持つ、$Z_U,Z_V$が交互に現れる語が基底である。
これを確かめるには、同じ側の二文字を掛けてその側の自由基底で書き、
間の$D$係数を右へ送る。重なりは同じ側の三文字、または同じ側の二文字と係数に限り、
その側の結合性と自由基底の一意性で解消する。反対向きでも同じ議論が通る。
従って$A$は両$B$上でも左右に自由である。

各頂点の自明加群を一緒に$k_*$と書く。誘導加群には短完全列
\[
 0\longrightarrow A\otimes_D k_*\longrightarrow
 (A\otimes_{B_U} k_*)\oplus(A\otimes_{B_V} k_*)\longrightarrow k_*\longrightarrow0
\]
がある。最初の写像は二つの商写像の差である。非空交互語はその末尾と異なる側の商でのみ
基底として残り、空語では対角の差が増大の核を張る。
$D$の自明加群の最小解消の階数は$1,2,2,1$、内部次数は$0,1,3,4$である。
実際、長さ$n$の語は二つの位置類の文字の別々の対称冪である。
先頭一文字の乗法の核はその位置類の交代二文字と残る対称冪で表され、
三次の二関係を掛けて全て得られる。その乗法の核は第二の位置類でも交代二文字を
取ったもので、一つの四次関係の自由倍となる。これは二次元空間での乗法
$W\otimes\operatorname{Sym}^j W\to\operatorname{Sym}^{j+1}W$の核が
$\bigwedge^2W\otimes\operatorname{Sym}^{j-1}W$であることを二回使ったものである。

自由性から誘導加群に対する$\operatorname{Tor}^A(k_*,-)$は
それぞれ$D,B_U,B_V$でのTorとなる。長完全列で$D$から両$B$への写像は
次数0と1では対角に単射、次数2と3では内部次数の相違で零である。
従って$A$の自明加群のTorは各位置で対角次数に集中し、階数は$1,4,6,4,1$で尽くされる。
元のGramテンソルは全ての一文字切断で階数4である。
実際、$\beta(x,-)=0$なら各単位での評価により$a(x)=m(x)=0$となり、
両者の核は別の直線である。右側も同様である。
隣接階数6と上のKoszul性から定理\ref{thm:frobenius-tensor-reconstruction}の
三つの条件を満たす。双対の完全対合と左右Koszul解消から標準AS正則性を得る。
\end{proof}

\begin{corollary}[印のある横断チャートの同型]
正則な二表示の間で順序付きの二因子を保つ添字保持同型は、
両文字空間の線型同型と直交二因子の相似変換の組で、
二つの表示$(a(x)*_U b(f),m(x)*_V n(f))$を対応させるものの存在と同値である。
二つの相似比は共通とする。印を忘れるときは標的の全ての横断する順序付き因子対を試す。
\end{corollary}
\begin{proof}
定理\ref{thm:unique-linear-factor-partner}で二位置ずつの変換が揃い、
全射Gram表示の同値から共通相似比の変換を得る。逆は関係列への代入で従う。
\end{proof}

\subsection{根基が直線の場合の共通な重み付き多項式環}
\begin{theorem}[根基直線の完全判定]\label{thm:radical-factor-amalgam}
補題\ref{lem:product-radical-cubic}の四評価階数3と$U\cap V=kr$を仮定する。
標準Koszul AS正則となる必要十分条件は、二文字関係$t$が非退化であることである。
\end{theorem}
\begin{proof}
六次元Gram空間の各元との対合を関係とみなす。
$t$は非零の二行二列関係であり、$U,V$の各三次元関係は
$S_U,S_V$の極形式そのものである。軸評価は各三文字空間に全射だから、
これらの階数1の関係は全窓にも純テンソルを与える。これで必要性を得る。
逆に$t$が非退化ならどちらの極形式系にも階数1の形式はない。
実際、$S_U(-,e,r)=0$なので階数1の対称形式があれば、その一次因子は$e$か$r$を消す。
$S_U(l,e,-)=0$なら$l\in kr$で、$S_U(r,-,-)$は階数2である。
$S_U(l,r,-)=0$なら$t$の非退化性から$l\in ke$で、$S_U(e,-,-)$も階数2である。
$e,r$が比例する場合も通る。$V$側も同様である。
三関係を各々に課した三生成代数を$B_U,B_V$とする。
全射Gram表示の双対の単射性により三関係は各々独立である。
前定理と同じ二次完全交叉の双対性により、どちらも標準のKoszul AS三次元である。

まず一方の三文字を単位$e$に関して分け、$W=\operatorname{ann}(e)$を二文字とする。
$S_U(l,e,-)$は$l$の退化計量による像なので、
二文字だけによる関係は$l\in kr$に対応するただ一式$t$である。
従って$\operatorname{Sym}^2 W$の像は二次元であり、全対称積の像は三次元である。
前定理と同じClifford表示を使う。すなわち全対称積の像の基底を
三つの独立な中心変数としてClifford代数を作ると、全射とHilbert級数の一致で同型である。
最初の二つの中心変数は二文字から取れる。このため二文字部分は非退化な一関係$t$だけで
表示される代数$D$になる。実際、対応する二文字のClifford基底は二中心変数上階数4で、
$1/(1-t_0)^2$の級数を与える（$t_0$は級数変数）。
非退化な二文字の一関係代数も各位置で基底を取り直せば二文字多項式環の級数となる。
第三の中心二次式$Q_U$を任意に補えば、$B_U$は$D[Q_U]$上左右に
$1,Z_U$を基底として自由である。これは二文字を先に、または後に並べたClifford基底による。
$V$側でも同じである。$D$の表示は第一・第二位置で$t,t^{\mathsf t}$であり、実際に同じ部分代数である。

ここで二側の第三の中心式の取り方が重要になる。$\langle s,r\rangle=1$なる
$s\in\mathcal H$を選ぶ。$U$上の対称双線型形式$q_U$を
$q_U(e,-)=\langle s,-\rangle|_U$として延長する。$q_U(e,r)=1$だが、
全ての極形式と二文字の二次式は$(e,r)$で消えるので、$q_U$の類を
第三の中心変数$Q_U$に取れる。$V$上でも同じように$q_V,Q_V$を選ぶ。
$q_U,q_V$を$(a,b),(m,n)$で引き戻すと、六つ目の関係は
\[
 \langle s,\beta(x,f)\rangle=q_U(a(x),b(f))+q_V(m(x),n(f))+b_0(\bar x,\bar f).
\]
実際、いずれかが軸のとき第一・第二項はそれぞれ指定した値になり、差は二文字だけの式となる。
反転位置では全てが転置される。

$D$については両位置を二次元空間に同定して$t$を非零交代形式とすることができる。
次位置の転置は符号だけ違うので、この同定で$D$は添字付き二変数多項式環である。
$b_0$と次位置での転置はこの可換環で同じ二次式を与える。
$Q_U$とその転置は三文字の対称な中心式を偶奇の表示で書いたもので、
二文字に対して$Q_{U,p}w_{p+2}=w_p Q_{U,p+1}$と交換する。
$V$でも同じである。従って$D$に次数2の中心変数$Q$を加えた共通部分
$D'=D[Q]$を、一方で$Q_U$、他方で$-Q_V-b_0$へ送って埋め込める。
後者は変数の平行移動で$D[Q_V]$と同じ部分代数となる。
位置を揃えた$D'$は重み$(1,1,2)$の多項式環である。
六関係のうち五つが両三関係の合併（$t$を共有）で、残りは
この$Q$の二表示の同一視であるから、$A=B_U*_{D'}B_V$である。

前定理の交互語基底と誘導加群の短完全列を$D'$に適用する。
$D'$のTorは内部次数$1,1,2$の三元上の外積代数である。
$\operatorname{Tor}_1$で二つの一次元方向は両側へ対角に単射、二次の生成元は零へ行く。
$\operatorname{Tor}_2$の内部次数2の一元は二文字の定義関係$t$に対応し、
各$B$の独立な二次関係の一つへ送られるため単射である。
これは二つの一次文字を含めたbar複体の二文字サイクルをそのまま写せば分かる。
内部次数3の二元と$\operatorname{Tor}_3$の内部次数4の一元は内部次数の差で零へ行く。
従って長完全列は再び$A$の対角Torだけを残し、その階数は$1,4,6,4,1$である。
各一文字切断の単射性は$a,m$と$b,n$の核直線の相違から従い、
Gramテンソルの完全対合により標準AS正則性を得る。
\end{proof}

\begin{corollary}[両対称反転窓の有限な分類手順]\label{cor:bisymmetric-effective-recognition}
隣接階数6の両対称テンソルの標準AS性と、正則なものの添字保持同型を有限な代数的条件で判定できる。
\end{corollary}
\begin{proof}
反対階数6では完全対称化可能性と対称完全交叉の条件を検査する。
階数6と8以外は定理\ref{thm:bisymmetric-six-or-eight}で排除される。
階数8では両鉛筆が四次完全交叉であることを確認する。
両側が被約なら系\ref{cor:all-reduced-apolar-recognition}のグラフ判定を使う。
その同型候補は射影的にも有限型の方程式で扱える。具体的には抽象グラフを
$\mathbb P^3\times\mathbb P^3$のHilbertスキームで探し、
両射影の同型性、二つの次数配分の一致および切断の制限条件を課す。
グラフの$(1,1)$に関するHilbert多項式はどれも$8n$である。

一側が非被約ならその鉛筆には階数2以下の元がある。平方の場合は
軸評価を検査し定理\ref{thm:double-plane-ore}を使う。
相異なる因子なら両側の四軸の評価階数3を検査し、
軸像の交わりが二次元以上のものは補題\ref{lem:reciprocal-apolar-factors}の
二行二列の純関係で排除する。残る二場合は本節の二定理で決まる。
同型は定理\ref{thm:unique-linear-factor-partner}から両文字空間での二つの基底変換と
六次元Gram空間の相似変換による係数同値と同値である。
必要なら印を指定した標的の因子を全て動かす。従って各条件は有限変数の方程式・不等式、
または指定された固定多項式のHilbertスキームで決まる。
この判定の対象は両対称表示を持つ反転窓である。
\end{proof}

% !TeX root = ../papers/nc-p3.tex
\section{反対位置の一組だけを対称化した四窓}
\label{sec:one-opposite-symmetric-evaluations}
四階テンソルを$E\otimes F\otimes E\otimes G$に置き、二つの$E$位置の交換で不変とする。
$F,G$の同定は仮定しない。以下では点座標を$E^*,F^*,G^*$と書く。
隣接階数6の因子分解を双対な六次元空間への全射$\beta:E^*\otimes F^*\to H$、
$\gamma:E^*\otimes G^*\to H^*$とすれば、テンソル対称性は
\[
 \langle\beta(x,f),\gamma(y,g)\rangle
 =\langle\beta(y,f),\gamma(x,g)\rangle
\]
である。関係の列は$R_\beta,R_\beta^{\rm op},R_\gamma,R_\gamma^{\rm op}$である。
ここで$R_\beta$は$H^*$との全対合であり、$R_\gamma$も同様である。

\begin{theorem}[四つの評価の下界]\label{thm:single-opposite-evaluation-ranks}
この四窓が標準AS正則なら、二つの写像をどちら側の非零文字で評価しても階数は3以上である。
\end{theorem}
\begin{proof}
各関係空間は次または前の転置窓と隣り合うので、
定理\ref{thm:no-consecutive-pure-relations}から純テンソルを含まない。
$d=\operatorname{rank}\beta(x,-)\le2$とする。$f\in\ker\beta(x,-)$では
$\beta(E^*,f)$が$\gamma(x,G^*)$を消す。従って後者は$R_\beta$内で
四行$d$列に支えられた関係空間に単射で写る。
その階数が4なら、射影三平面が四行二列の階数1のSegre多様体（次元4）に
$\mathbb P^7$内で交わる。$d<2$でも零列を補って同じである。
よって$\gamma(x,-)$の階数は3以下である。両$E$位置を$x$の一次元商、
$F,G$位置をそれぞれの核の双対へ制限すれば全関係が消え、
長さ$4m$の次元は$2^m$以上になる。$\gamma$側も同様である。

次に$a=\operatorname{rank}\beta(-,f)\le2$として
$S=\ker\beta(-,f)$、$s=4-a$と置く。両$E$位置を$S^*$、$F$位置を$f$の
一次元商へ制限する。この文字に触れる関係は全て消える。
$\gamma(S,G^*)$は$\beta(E^*,f)^\perp$に入り、そこは$6-a=s+2$次元である。
従って残る$S^*,G,S^*$の三文字ブロックの次元は少なくとも
$4s^2-2s(s+2)$で、$s\ge3$なら6以上である。

$s=2$では四次元を標的とする二行列$A,B:G^*\to k^4$で制限された$\gamma$を表す。
三文字ブロックの双対は$h_{ij}\in G^*$ ($i,j=0,1$)のうち
\[
 Ah_{0j}+Bh_{1j}=0,\qquad Ah_{i0}+Bh_{i1}=0
\]
を満たす空間である。$A$が可逆なら$h_{11}$を任意に選び、
$h_{01}=h_{10}=-A^{-1}B h_{11}$、$h_{00}=(A^{-1}B)^2 h_{11}$とできる。
一般の二行列でも核次元の上半連続性で4以上である。
一文字$f$でブロックを分離すると前後を結ぶ関係はないため、
$m$個のブロックはそのテンソル積となり指数成長する。
$\gamma(-,g)$については二つの奇位置を取り替えればよい。
\end{proof}

\begin{corollary}[同じ点スキームと恒等四移動]\label{cor:single-opposite-common-points}
同じ標準AS仮定の下で、
$\mathbb P(E^*)$上の$\beta(x,-)$と$\gamma(x,-)$の四次小行列式は
同じ退化スキームを定める。四辺の点対応は全てこのスキームのグラフで、
点列の四移動は恒等である。台の有限性・被約性は要らない。
\end{corollary}
\begin{proof}
$\beta(x,-)$の退化スキーム上では階数が丁度3であり、核直線$f$が定まる。
恒等式により$\gamma(x,G^*)$は$\beta(E^*,f)^\perp$に含まれる。
後者の直交条件に独立な三本があり、前者にも独立な三本がある。
局所的に単元となる小行列式を取れば、$\gamma$の評価像は階数3の部分束と一致し、
四次小行列式はスキーム上で消える。逆の包含も同じである。
各辺の零点スキームは定階数の核直線のグラフであり、反対側の評価でも階数3のため
両射影はそれぞれの退化スキームへの同型である。
従って$x$から出る点列はスキーム上で一意な$x,f,x,g,x,\ldots$となる。
\end{proof}

\begin{proposition}[片側対称の相手窓の決定]\label{prop:single-opposite-prolongation}
第一窓$R_\beta$から
\[
 P(\beta)=\{L:E^*\to H^*:\quad
 \langle\beta(x,f),L(y)\rangle=\langle\beta(y,f),L(x)\rangle\ \forall x,y,f\}
\]
を作る。片側対称な標準AS延長があれば$G^*=P(\beta)$で、
$\gamma$は普遍評価$\gamma(x,L)=L(x)$である。
逆の操作でも$P(\gamma)=F^*$である。各段は二十四元二十四式の線型演算で、
両端交代な類似の連立式の核は零でなければならない。
\end{proposition}
\begin{proof}
$R_\beta\otimes E$を$\operatorname{Hom}(E^*,H^*)$と同定すると、
$R_\beta\otimes E$と$E\otimes R_\beta^{\rm op}$の交叉は
両端交換に安定で、正固有部分が$P(\beta)$である。
全部の交叉は標準AS性から四次元、$\gamma(-,g)$の四元は一文字切断の全階数から
独立であり、対称式で正固有部分に入る。これで主張と逆向きの結果を得る。
\end{proof}

構成した相手が第二の反対対称性も持つかを調べるには、
$T:F^*\to P(\beta)$と対称写像$D:H\to H^*$について
$T(f)(x)=D\beta(x,f)$を線型に解き、両方が可逆な解を探せばよい。
実際、これが成り立つと前節のGram表示であり、
逆に$F^*$と$G^*$を同定した後で両対称なら二文字をまとめたGram行列が対称となる。
全射$\beta$の核はその行列の両核であるため相手の写像は$\beta$を経由し、
誘導される六次元の写像は可逆対称となる。次の未解決域は、この線型空間に
可逆解がなくても両方の普遍評価が四評価下界とAS条件を満たす場合である。

\begin{proposition}[有限な共通点スキームの固有な二乗像]\label{prop:single-opposite-intrinsic-square}
共通点スキーム$X\subset\mathbb P(E^*)$が有限な場合、その飽和イデアルの
最小解消の最後の一次行列から射$X\to\mathbb P^9$が決まり、
これは両方の奇位置への射の二次Veronese像を同時に与える。
このスキーム像を含むVeronese三次元多様体が一意なら、片側対称の標準AS延長は
両対称チャートに入る。
\end{proposition}
\begin{proof}
$6\times4$評価行列の最大小行列式イデアルは期待余次元3である。
Eagon--Northcott解消は次数$4,5,6$で階数$15,24,10$を持ち、
完全Cohen--Macaulay商なのでイデアルは飽和している。
両評価行列は系\ref{cor:single-opposite-common-points}で同じイデアルを解消する。
次数零の鎖同型は各項の定数基底変換となり、恒等写像の二つの持ち上げの差の
零ホモトピーは、次の項のシフトが厳密に大きいため零である。
従って最後の行列とその核の射影像は基底を除いて固有である。
体積直線を除くと、この行列を$x$で評価した写像は
$Q\mapsto\beta(x,-) Q$である。ここで$Q\in\operatorname{Sym}^2(F^*)$を
$F\to F^*$とみなす。これは最後の対称二乗から一文字を縮約する解消の微分である。
$X$上で評価核が直線$kf$なので、この行列の核は$kf^2$である。
階数一定の局所基底で同じことが言えるためスキーム上の射を得る。
$G$での表示も同じ十次元空間の射を与える。
Veroneseが一意ならその二つの表示は射影三空間の同型で対応し、$X$上で$f$と$g$が同定される。
すると二辺の有限グラフが一致する。有限な六双一次式完全交叉の$(1,1)$イデアルは
その六式で張られるので（定理\ref{thm:finite-square-chart}）、両奇窓の関係空間も一致する。
今度は奇位置を両端とする三次交叉も、命題\ref{prop:finite-flip-obstructions}により
端交換の正固有部分に入る。従ってテンソルはこの同定で両対称である。
\end{proof}

% !TeX root = ../papers/nc-p3.tex
\section{異なる二つのVeronese表示を持つ階数10の正則窓}
\label{sec:twenty-point-companion}
$u,v\in k^3$に対し、標準体積形式と内積で定まる外積を$u\times v$、
$D(u,v)_i=u_jv_k+u_kv_j$（$(i,j,k)$は巡回順）とする。
$\epsilon=\pm1$について六次元への写像を
\[
 B_\epsilon((r,u),(s,v))
 =\big(rv+\epsilon u\times v,\;su+\epsilon D(u,v)\big)
\]
と置く。前節の$\beta=B_+$、$\gamma=B_-$に通常の内積を使う。

\begin{theorem}\label{thm:twenty-point-distinct-companion}
この四窓は標準14641型Koszul AS正則で反対階数10、最小周期2である。
共通点スキームは被約な二十点で、その固有な二乗像は相異なる二つの
Veronese三次元多様体に入る。同じ第一窓からの片側対称な相手は一意である。
\end{theorem}

\paragraph{点と階数。}
$B_+,B_-$を六行十六列行列とする。$w=B_+^{\mathsf t}B_-$は二つの偶位置の交換で
不変である。これは外積の内積の公式を代入しても、四添字の係数比較でも従う。
両隣接階数は6、反対切断の階数は10である。
偶文字$(r,x,y,z)$で評価した$B_+$の四次小行列式を調べる。
$r=1$の辞書式消去は
\[
 x^2-4y^2z^2-y^2-z^2,\quad 2xy^2+2xz^2+x,\quad xyz,\quad
 2y^3+y,\quad 2z^3+z
\]
が生成するイデアルを与える。台は原点と、一座標が零で残り二つが
独立に$\pm\sqrt{-1/2}$である十二点である。
$r=0$では$u\times v=0$から$v$は$u$の非零倍であり、
$s u_i+2u_ju_k=0$に規格化できる。座標軸の三点と、
三座標が符号を除き等しい四点を得る。各点で奇側の核は直線である
（$r=1,u\ne0$では$I+[u]_\times$の核が直線で$s$も決まる）。
従って六本の双一次式は有限完全交叉で、その長さ
$[h_1^3h_2^3](h_1+h_2)^6=20$と相異なる点数が等しく、全て被約である。
$J=\operatorname{diag}(-1,1,1,1)$とすると
$B_-=-B_+(J\otimes J)$で点集合も$J$に安定である。
第三・第四位置を$J$で同定すれば、反復する二窓は
$R_0=\operatorname{row}(B_+)$、$R_1=\operatorname{row}(B_+^{\rm op}(1\otimes J))$となる。
ここで$\mathrm{op}$は二文字の順序交換である。

\begin{lemma}[正規二次式を持つ二周期窓の有限認証]\label{lem:normal-parameter-six-test}
二周期の四文字六関係の窓で、各位置に四つの二次式$q_{a,p}$があり、
$q_{a,p}v=\sigma_a(v)q_{a,p+1}$が全ての文字で成立するとする。
$\sigma_a$は二周期の次数付き自己同型で全ての$q_b$を固定する。
次の条件で全次数の標準Koszul AS正則性を認証できる。
\begin{enumerate}
\item 各位置で四式を割った次元は次数0から5で$1,4,6,4,1,0$、
元の次元は次数6まで$\binom{n+3}{3}$である。
\item 三次・四次Koszul交叉の次元は4と1で、その最上テンソルは
定理\ref{thm:frobenius-tensor-reconstruction}の一文字・隣接階数を持ち巡回で対応する。
\item 四次で止めた右Koszul列$C_i=W_i\otimes A$で、
各$i=0,\ldots,4$の位置$i$、内部次数$i+2$におけるホモロジーが零である。
\end{enumerate}
\end{lemma}
\begin{proof}
始点ごとに商の斉次基底を語の線型結合$b$に持ち上げ、抽象的な重み2の
四変数多項式環$P$上でこれらを基底とする自由加群$L$を取る。
正規性で四式を右端へ移せ、互いに可換なので$b p(q)$への写像$L\to A$は全射である。
商が長さ5以上で零になることは一次生成性による。
$H_L(t)=(1+t)^4/(1-t^2)^4$なので次数6まではこの全射が同型である。
各文字の左乗法を、$b$での積をこの同型で一意に持ち上げて右$P$線型に延長する。
各二次関係の作用は$b$で評価すれば次数が高々6で零なので、行列は全ての次数で関係を満たす。
空語に作用させると$A\to L$を得る。次数2以下の一致と$P$線型性より、
$q$の任意の単項式は空語を対応する$P$係数倍へ送る。さらに各$b$も空語を
その基底へ送るため合成$L\to A\to L$は恒等である。これで全次数の自由性を得る。
二周期性により、二文字を付けた始点・終点の移動はここで同一視できる。

次に各$q_a$について$C_i$上に内部次数2の右加群鎖写像
$m_i(w\otimes b)=\sigma_a(w)\otimes q_{a,p+i} b$を入れる。
微分が末尾一文字を掛ける定義と正規性から鎖写像である。
これは係数も含めて$\sigma_a$を施した後の右$q_a$倍に等しい。
生成元上で低い$i$から順に鎖ホモトピーを作る。
既に$h_{i-1}$を選んだとき$(m_i-h_{i-1}d_i)(w\otimes1)$は
内部次数$i+2$のサイクルである。仮定3から一つ逆像を取り、
$h_i$を右線型に延ばす。$i=4$ではそのサイクルが零である。
従って右$q_a$倍は全ホモロジーを消す。
各$C_i$は有限自由$P$加群で、全ホモロジーは原点に支えられる。
長さ4のこの複体の正ホモロジーは零である。実際、原点で局所化し、
最上項から帰納すると既に完全な尾部の余核は位置$i$で深さ$i$以上になる
（自由加群の深さ4と短完全列の深さの補題）。$i>0$ではそこに有限長部分加群がない。
原点の外では初めからホモロジーが零である。
次数零のホモロジーは始点単純であり、線型解消を得る。
Koszul双対の完全対合は仮定2とこの解消から従い、左右AS性と大域次元4を与える。
\end{proof}

\begin{proof}[定理\ref{thm:twenty-point-distinct-companion}の全次数証明]
二周期座標で$ii$を位置$p$での同名二文字の積とする。
四つの$q_a$を$00,11,22,33$の係数で与える。$q_0$は両位置で$(0,1,1,1)$、
$q_1$は偶位置で$(-1,0,2,-2)$、奇位置で$(-1,0,-2,2)$とし、
$q_2,q_3$は後三座標を巡回させる。両位置でこれらは四平方の基底である。
$\sigma_0=1$、$a=1,2,3$では$\sigma_a$は$0,a$の文字を固定し残り二文字を負にする。
六関係が標数2の二次元ベクトル空間の文字次数で斉次なのでこれらは自己同型である。
三次の経路商で$ q_{a,p}v=\sigma_a(v)q_{a,p+1}$を係数ごとに検算する。

\path{computations/verify_twenty_point_companion.py}は有理数体で以下を計算する。
長さ$n$の商は、一文字と既知の長さ$n-1$の商の積から、
六関係と長さ$n-2$の商の積の列空間を引くことで作る。
四平方を追加した場合と元の場合に補題の次元列を得る。
各始点で$W_3,W_4$は次数ごとの二つの部分空間の交叉から直接計算して次元4と1になる。
$w$の巡回と基底変換は$W_4$の非零元を与え、各一文字切断の階数は4である
（軸での評価、または係数行列）。$d_i$を右係数次数2で計算した階数は
$i=1,2,3,4$の順に$20,45,36,10$、右係数次数1では$10,20,15,4$である。
隣り合う二階数の和は各内部次数$i+2$で中間項の次元と等しく、補題が適用できる。

階数10なので定理\ref{thm:bisymmetric-six-or-eight}より両対称表示を持たない。
周期1なら偶位置の対称な交換を一位置移して奇位置の対称な交換を得る。
二組は異なる因子に作用するので同時に対称化でき、矛盾する。
二周期表示と合わせて最小周期は2である。
二つのVeroneseが一致すれば命題\ref{prop:single-opposite-intrinsic-square}の証明の通り
有限グラフと二つの奇窓が一致し両対称になる。従ってこの二つは相異なる。
相手窓の一意性は命題\ref{prop:single-opposite-prolongation}である。
\end{proof}

% !TeX root = ../papers/nc-p3.tex
\section{三軸表示の全係数と正則性の境界}
\label{sec:cross-product-parameters}
$L=\operatorname{diag}(a_1,a_2,a_3)$とし、前節の二写像を
\[
 B_\epsilon((r,u),(s,v))
   =(rv+\epsilon(Lu)\times v,\;su+\epsilon D(u,Lv))
\]
に広げる。係数について
$\Delta=a_1^2a_2^2+a_2^2a_3^2+a_3^2a_1^2$、$p_i=a_j^2+a_k^2$（巡回順）と書く。
偶位置の原点を$[1:0:0:0]$と呼ぶ。

\begin{theorem}[三軸族の全次数判定]\label{thm:cross-product-parameters}
標準Koszul AS正則である必要十分条件は$\Delta\ne0$である。
三係数が非零な正則点では反対階数10、最小周期2で、ちょうど二係数が非零なら
両対称チャートの反対階数8の正則点である。三係数が非零な正則点のうち、
共通点スキームが有限である必要十分条件は$p_1p_2p_3\ne0$であり、
このときスキームは被約な二十点である。
\end{theorem}

\begin{proof}
係数を不定元にした行列で$B_+^{\mathsf t} B_-$の偶位置対称性が成り立つ。
反対行列を四つの文字次数で分けると、次数零の四行四列は行列式$-4\Delta$、
残る三つの二行四列は定数の階数2の行列の$a_i$倍である。
以下の三次恒等式とともに\path{computations/verify_cross_product_parameters.py}で
係数ごとに再現できる。

まず$a_1a_2a_3\ne0$とする。二周期表示で、位置$p$からの積を文字で表し、
$e=(-1)^p$と書くと六関係は
\begin{equation}\label{eq:axial-reduction}
 0i=-a_j jk+e a_k kj,\qquad i0=-a_k jk-e a_j kj .
\end{equation}
ここで$0$は一文字である。三次の次元は全トーラス上で20である。
実際、$000$以外は零番文字を消して三文字の語にできる。文字次数$i\ne0$にある七語には
$j0k,k0j$の二通りの計算から二関係を得る。それぞれ
\[
 a_i(jij-kik)+e(a_j jji-a_k ikk),\quad
 e a_i(kik-jij)+a_k kki-a_j ijj
\]
で、二つは独立である。次数零には六順列と$000$があり、$i0i$の両側の計算は
$-a_k jki-e a_j kji+a_j ijk+e a_k ikj$を与える。
異なる二つの$i$では台が異なるので少なくとも二式は独立である。これが全体の上界20を与える。
下界は巡回Gramテンソルの四つの独立な一文字縮約が両位置の三次交叉に入ることによる。
一文字切断の全階数は、二つの軸における評価で各四文字の一次独立性から分かる。

四平方の係数で$q_{0,p}=(0,a_1,a_2,a_3)$とし、$q_{i,p}$の
$0,i,j,k$成分を順に
\[
 (-a_i,0,e a_k(a_i^2+a_j^2),-e a_j(a_i^2+a_k^2))
\]
と取る。前節と同じ四指標で$q_{h,p}v=\sigma_h(v)q_{h,p+1}$となる。
上記コードでは有理関数体で文字次数ごとの三次関係行列にこれらの式を追加し、
全てで階数11を保つことを調べる。三次商の次元はトーラス全体で一定と示したので、
各ファイバーの関係像は部分束であり、上の多項式恒等式は全ファイバーへ延びる。
四係数行列の行列式は$4\Delta^2$である。

四平方を割ったときの次元の上界は、三係数が非零なら常に$1,4,6,4,1,0$である。
二文字以上の語から式\eqref{eq:axial-reduction}で零番文字を除く。
$0ii$と$ii0$の計算から三文字全ての異なる並べ替えは互いに非零定数倍である。
$j0k$の計算から$jij=kik$である。従って三次は四次元以下である。
四次で三文字全てが現れ、既に平方がない語は、相異なる三文字の連続部分を並べ替えて
残る文字と隣り合わせにすれば零になる。残る交互語$ijij$は
$i j i=k j k$または$j i j=k i k$で第一または第二の添字を取り替えられ、
六つとも等しい。五次では四次の代表の末尾を第五文字と同じにして零とできる。

一般ファイバーの元代数の各次数の次元は標準値であることを先に確認する。
$(1,1,1)$がKoszulであり、二次関係と三次交叉はその近傍で束である。
$n$を固定し次数$n-1$まで平坦な近傍で、
$W_3\otimes A_{n-3}\to R\otimes A_{n-2}\to V\otimes A_{n-1}$を見る。
二写像の合成は零で基点では中間項で完全だから、階数の下半連続性で階数は近傍で一定となる。
最後の余核が$A_n$であり次数別帰納ができる。従ってどの特殊点でも各次元は
有理関数体での標準値以上である。

$\Delta\ne0$なら上の四正規元は四平方を張るので、
先の商の上界と正規性から各始点で多項式環上の高々$1,4,6,4,1$個の生成元があれば足りる。
この自由加群のHilbert級数は$(1-t)^{-4}$以下で、先の下界と一致する。
従って全次数の次元が標準で、両側とも四正規元の多項式環上で自由である。

ここからKoszul性も従う理由を記す。二頂点で折り畳み、各$q_i$は二頂点での和を指す。
四階テンソルの一文字・二文字の隣接切断は全階数である。
$k,V,R,W_3,kw$を係数空間とする右Koszul列の最初の二微分は表示によって完全
（$V\otimes A$の項での完全性）であり、最後の二微分は左の最初の二微分の
$A$双対になる。実際、$w$の$(1,3),(2,2),(3,1)$切断による完全対合で
一文字を移す両包含写像が互いに転置となる。
左の始点単純から$A$への$\operatorname{Hom},\operatorname{Ext}^1$はともに零である。
前者は$q_0$倍が単射だからである。後者について左増大イデアルからの左加群写像$f$を取る。
可換な二つの正規元$q_0,q_1$について$q_0 f(q_1)=q_1 f(q_0)$であり、
左多項式環上の自由性と互いに素な変数であることから$f(q_i)=q_i b$ ($i=0,1$)となる。
任意の正次数元$z$で$q_0$は中心的なので$q_0f(z)=z f(q_0)=q_0zb$、
従って$f$は右$b$倍に延びる。頂点別には冪等元で切ればよい。
これで右Koszul列の最上二項の正ホモロジーが消える。残る中央のホモロジーは
既に得たHilbert級数とEuler等式から零である。線型解消と完全対合から標準AS性が従う。

一方$\Delta=0$なら、$v_i=1/a_i$で$(s,v)=(1,v)$を取ると
$v\cdot v=0$、$Lv=(1,1,1)$となる。上側の値$rv+(Lu)\times v$は
$v^\perp$に入る。これが零のとき$Lu\in v^\perp$である。
実際、非零等方$v$では$w\times v\in kv$と$w\cdot v=0$が同値である。
従って$\sum u_i=0$となり、下側の値$(\sum u_i)(1,1,1)$も零である。
この奇文字の評価階数は2以下であり、定理\ref{thm:single-opposite-evaluation-ranks}で排除される。

ちょうど二係数が非零なら、例えば$a_i=0$として、$T$を$i$座標だけを負にする行列と置く。
$B_-((r,u),(-s,Tv))=\operatorname{diag}(T,-I)B_+((r,u),(s,v))$で両対称表示になる。
$0i$は偶側の反対核の相異なる因子で、二つの評価像は上下の相補三平面である。
対応する奇因子は$0,i$で、四評価とも階数3である。
上の三次元単位積の左乗法には非零な二つの軸回転$[a_j e_j]_\times,[a_k e_k]_\times$が現れ、
これらは可換でないので積は非結合的である。定理\ref{thm:transverse-factor-amalgam}を適用する。
階数は先のブロックから8である。一係数以下しか非零でないと、零係数の軸文字で
片側の評価が階数2以下となる。これで全パラメータを尽くした。

最後に三係数が非零な正則点の零点を調べる。$r=0$では$v$が$Lu$に比例し、
規格化後に$s u_i+p_i u_j u_k=0$となる。全$p_i\ne0$なら三座標軸と
$u_i^2$が$p_i$に比例する四点を得る。$p_i=0$なら$u_i=s=0$の直線を含む。
$r=1$では原点以外で$Lu\cdot v=v\cdot v=0$となる。
$u_1u_2u_3\ne0$と仮定して$w_i=a_i v_i/u_i$と置くと下側の式から
$w_i=c(2u_i^2-\sum u_j^2)$となる（$c\ne0$、さもなくば$v=0$）。
従って$Lu\cdot v=0$から$2\sum u_i^4=(\sum u_i^2)^2$となる。
差を四一次式に分解すれば符号変更後に一座標が他の二つの和である。
代入により$(a_i v_i)^2$が三つとも同じ非零値となり、
$v\cdot v=0$は$\Delta=0$を要して矛盾する。
例えば$u_k=0$で残りが非零のとき、上側から$v_i=-a_j u_j v_k$、
$v_j=a_i u_i v_k$で、下側と合わせると$u_i^2=u_j^2=-1/p_k$を得る。
一座標だけ非零の解はない。よって全$p_i\ne0$なら前節と同様に
$1+12+7=20$点で各奇核は直線であり、完全交叉の長さから被約である。
階数10のときの最小周期は前節と同じである。
\end{proof}

\begin{proposition}[原点を指定した同型と像の次元]\label{prop:cross-product-marked-isom}
$a_1a_2a_3\Delta\ne0$で原点を指定した二表示が添字保持同型である必要十分条件は、
三つの平方の順序を忘れた射影組$[a_1^2:a_2^2:a_3^2]$が一致することである。
原点を忘れた標準AS同型類の像も二次元である。
\end{proposition}
\begin{proof}
偶側の反対像は対称二乗の全体なので、その二つの基底変換は射影的に一致する。
固定した偶原点での奇核も両方一意である。三つの点空間の基底変換を
$(r,u)\mapsto(Ar+h(u),P u)$、$(s,v)\mapsto(Bs+k(v),Qv)$、
$(t,z)\mapsto(Ct+\ell(z),Tz)$と書く。
二つの全射の対による隣接分解の一意性から、$B_+$の標的の変換は両軸で
$AQ,BP$のブロックとなり、$B_-$の標的ではその逆転置の共通スカラー倍となる。
従って$P^{\mathsf t}P$はスカラーである。
上側の残る乗法は$h(u)I+[L'Pu]_\times=A Q[Lu]_\times Q^{-1}$であり、
跡から$h=0$となる。$L$が可逆なので$Q$は全ての歪対称行列を共役で歪対称へ送る。
$Q^{\mathsf t}Q$は従って三つの座標回転と可換で、スカラーに限られる。
$P=pO_V,Q=qO_U$とすると、$L'$は$O_U L O_V^{-1}$の非零定数倍であり、
$L^{\mathsf t}L$の固有値の射影組の一致が必要である。

逆に共通の置換で座標を並べ直せば、$L'=h_0 L S$と書ける
（$h_0\ne0$、$S$は対角符号行列）。共通置換は外積に置換符号を掛け、
$D$を保つため、$r$の符号で補正できる。
対角符号と倍率については偶変換を$(r,u)\mapsto(h_0\det(S) r,u)$、
両奇変換を$(s,v)\mapsto(h_0s,Sv)$と取る。
二つの六次元値が共通のブロック変換
$\operatorname{diag}(h_0\det(S) S,h_0 I)$で移り、計量の相似比は$h_0^2$である。
これで十分性を得る。一般の点スキームは二十点なので、印を忘れる写像のファイバーは
高々二十通りの原点の選択で尽くされる。印付きの二次元像は次元を保つ。
\end{proof}

\begin{theorem}[階数10の全特殊値で印を忘れた同型]\label{thm:cross-product-unmarked-isom}
$a_1a_2a_3\Delta\ne0$なら、原点を指定しない添字保持同型の条件も
三つの平方の順序を忘れた射影組の一致である。
\end{theorem}
\begin{proof}
まず全$p_i\ne0$のとき、偶位置の配置だけから不変量を回収する。
四つの座標点を頂点とし、残る十六点は各座標面に四点ずつ載る。
面$h$上の四点の座標平方は一つの射影行であり、非零座標に全ての符号を付けた点で尽くされる。
第零列の符号を変える対角基底変更（二乗前は$\sqrt{-1}$倍）により、四行は
対称な零対角行列$H$となり、$H_{0i}=p_i$、$H_{ij}=1$ ($1\le i<j\le3$)である。

三点以上を通る直線の本数は、各頂点で6、残る各点で3である。実際、一つの面の四符号点は
どの三つも独立で、その六対の結線には座標のうち一つの頂点が載る。
二頂点を結ぶ辺にさらに点はない。同じ面にある二点の直線がその面の外の点を通ることもない。
そこで三つの相異なる面の符号点だけを調べればよい。
頂点の置換後に残る面を無限遠面に取っても、面の規格化は同じ形に戻せる。
実際、対角基底変更と各行の規格化で$H_{ij}$を$t_i t_j H_{ij}$に変え、
第零以外の三頂点間の値を全て1にできる。新たな三値$p'_i$は
四頂点の三通りの対分割の辺積$H_{0i}H_{jk}$の並べ替えに共通の定数を掛けたものである。
従って$2\sum_{i<j}p'_ip'_j-\sum_i(p'_i)^2\ne0$は保たれる
（元の値は$4\Delta$）。アフィン座標で面$i$の点は第$i$座標だけが零で、
残り二座標は絶対符号を除いて共通の$m_i$、$m_i^2=-1/p'_i$となる。
これら三点が共線なら、アフィン依存の三係数$k_i$は和が零で、
三座標の式から$k_i m_i$は符号を除いて同じ非零値となる。
従って$\pm\sqrt{p'_1}\pm\sqrt{p'_2}\pm\sqrt{p'_3}=0$となり、上の非零条件に反する。

この本数で四頂点は配置から決まり、配置の同型は単項行列である。
同じ頂点順序で二つの対称行列表示が行・列のスカラー倍で対応すると、
対称性と全辺の非零性から倍率は共通定数を除き$t_i t_j$の形となる。
従って三つの対分割の辺積は、順不同の射影組として固有である。
これは$[p_1:p_2:p_3]$であり、$2a_i^2=p_j+p_k-p_i$により所望の不変量を得る。

次に$p_3=0,p_1p_2\ne0$とする。台は一本の直線$r=u_3=0$と、
孤立した十点（原点、第三軸、第一・第二座標面に各四点）からなる。
孤立点だけの三点直線は、原点と第三軸で各四本、他の各点では二本である。
実際、各四符号点の属する面には原点と第三軸があり、二点の結線のうち
これら二頂点を通るものが該当する。面の異なる孤立点三つにこのほかの結線はない。
従って二頂点が特定でき、これらと残る各点で張る二つの面を唯一の直線と交わらせて
他の二頂点も順序を除き復元できる。二つの四点面の平方行について、
第零・第三列の交叉比$H_{10}H_{23}/(H_{13}H_{20})=p_1/p_2$は
面または二頂点の交換による逆数を除き固有である。これで$[0:p_1:p_2]$を得る。
二つの$p_i$が零なら上の方程式から台に二本の直線があり、平方の組は
順序と共通倍率を除き$(1,1,-1)$しかない。各場合の十分性は印付きの同型による。
\end{proof}

\begin{corollary}[階数8の座標境界の同型]
三軸族の全正則点で、印なし同型は三係数の平方の順序を忘れた射影組の一致と同値である。
\end{corollary}
\begin{proof}
階数の相違で二つの部分は交わらない。残る$a_1=0$、$a_2=b,a_3=c$の場合を
先の両対称座標で比較する。両反対核は$\langle 01,23\rangle$であり、
四直線の交叉の四特異点が座標頂点を指定する。従って同型の各側の行列は単項行列で、
置換は対分割$01|23$を保つ。両側の平方同士の係数行列は
\[
 M=\begin{pmatrix}0&-1&1&1\\-1&0&-b^2&-c^2\\
 -1&b^2&0&-b^2-c^2\\-1&c^2&-b^2-c^2&0\end{pmatrix}.
\]
$b^2+c^2\ne0$なら零配置から両置換は一致する。
三通りの対分割$ij|kl$で$M_{ij}M_{ji}M_{kl}M_{lk}$を取ると
$[(b^2+c^2)^2:c^4:b^4]$が得られ、最初の分割が指定される。
$u=b^2/(b^2+c^2),v=c^2/(b^2+c^2)$は$u+v=1$、
$u-v=u^2-v^2$であるため、射影比$[b^2:c^2]$が交換を除き決まる。
$b^2+c^2=0$は零成分の数で区別されて一軌道である。
印付き十分性の符号・置換による構成は零係数を含めても成立する。
\end{proof}

% !TeX root = ../papers/nc-p3.tex
\section{四元群の一文字次数が欠ける片側対称窓}
\label{sec:klein-companion-windows}
$K=(\mathbb Z/2)^2$を文字次数群とし、片側対称な二写像
$\beta:E^*\otimes F^*\to H,\ \gamma:E^*\otimes G^*\to H^*$が斉次とする。
三つの点空間は各次数一次元の正則表現とする。
テンソル直線の指標は必要なら$G$の文字次数を平行移動して吸収する。
隣接階数は6である。以下で被覆を与えるのは、$H$で少なくとも一つの文字次数が欠ける場合である。
単純な固定点とは、四元群で固定される共通点でその辺の零点が横断的なものをいう。

\begin{theorem}\label{thm:klein-missing-character}
上の欠ける次数を持つ全入力で、標準AS正則性を有限に判定できる。
座標固定点の一つが単純な標準AS窓は、両対称チャートまたは前節の三軸族に属する。
特に欠ける次数を初めから仮定せずとも、単純な固定点のあるこの斉次類の
非両対称AS窓は全て、三係数の非零な正則三軸族で尽くされる。
全ての座標固定点が非横断的な場合には、後述する二平面ずつのPBW表示または
両対称チャートにより判定できる。
\end{theorem}

\paragraph{文字ブロック。}
定理\ref{thm:single-opposite-evaluation-ranks}の座標文字での四評価下界と
両関係空間に純テンソルがないことは必要である。以後この有限条件を満たすとする。
欠ける次数を平行移動して零とし、
$b_{ab}=\beta(e_a,f_b), c_{ab}=\gamma(e_a,g_b)$と書く。
$b_{aa}=c_{aa}=0$、$a\ne b$では全て非零である。
両対称性を課すのは偶側だけで、式は
$\langle b_{ab},c_{cd}\rangle=\langle b_{cb},c_{ad}\rangle$である。
従って$c_{ba}$は$b_{ab}$を消し、その非零な次数空間は二次元以上である。
総次元6より三つの非零次数は各二次元である。
一色の四ベクトル$b_{ab},b_{ba},b_{cd},b_{dc}$のうち三つが比例すれば、
その直線を消す共変ベクトルが一成分だけの純関係を与えるので、三つは比例しない。
$c_{ab}$の直線は$b_{ba}$の零化直線で尽くされる。

欠ける次数を仮定せずに単純な固定点から始めた場合もこの表示に入る。
偶座標点の唯一の奇核は固定された座標点である。二つの軸評価の像は、
それぞれ同じ一つの次数を除く三つの次数に直線を持つ。横断性から各次数で
この二直線は独立である。従って残る一つの次数は零である。

\paragraph{単純な頂点。}
$0$を単純な座標点とし、$U=\beta(e_0,F^*)$と$V=\beta(E^*,f_0)$を取ると
$H=U\oplus V$である。$\gamma(e_0,G^*)=V^\perp$、
$\gamma(E^*,g_0)=U^\perp$である。三つの軸評価で$u\in V,y\in U,z\in U^*$を
補座標に取ると、最後の軸は$V\to V^*$の可逆対称写像となる
（$f_0,g_0$での偶対称性）。異なる文字次数が直交するのでその対角を1に揃えられる。
両写像は
\[
 \beta=(r y+A_u y,\;s u+B_y u),\qquad
 \gamma=(r z+A_u^{\mathsf t}z,\;t u+C_z u)
\]
となる。転置の式は$e_0$を入れて得られ、$B_y,C_z$は$g_0,f_0$を入れて対称となる。
斉次性から、巡回順$i,j,k$で
$A_i e_j=\alpha_i e_k, A_i e_k=\delta_i e_j, A_i e_i=0$であり、
$B_i=b_i(E_{jk}+E_{kj}), C_i=c_i(E_{jk}+E_{kj})$である
（以下の$b_i$は一つの係数で、先の$b_{ab}$とは区別する）。
残る偶対称式は
\begin{equation}\label{eq:klein-flat-coefficients}
 b_i c_j=-\delta_i\delta_j,\qquad b_j c_i=-\alpha_i\alpha_j
 \qquad((i,j,k)\text{ 巡回})
\end{equation}
で尽くされる。これは$\beta(e_i,f_i)$と$\gamma(e_j,g_j)$の対合を、
偶文字を交換した値と比較し、奇文字も交換した比較を加えたものである。
その他の比較は文字次数で消えるか同じ式になる。

全係数が非零なら、三式を掛けて$\prod\alpha_i/\prod\delta_i=\epsilon=\pm1$を得る。
$U$の三座標を対角に変えると$\delta_i=\epsilon\alpha_i$にできる。
実際、各辺で対角倍率の平方の比を指定する三条件の積が1であり、代数閉体上で平方根が取れる。
変更後の六式から$b_i=d\alpha_i,c_i=-d^{-1}\alpha_i$となる。
$y,z$の相反する共通倍率で$d$を吸収する。
$\epsilon=-1$はまさに前節の三軸族である。$\epsilon=1$なら$A_i$は全て対称で、
$\gamma((r,u),(-s,y))=\operatorname{diag}(I,-I)\beta((r,u),(s,y))$となるため
両対称判定に帰着する。

係数が零になる場合も残さない。例えば$b_i=0$なら奇座標$f_i$への評価の
上側だけで階数3が必要なので$\delta_j\alpha_k\ne0$である。
式\eqref{eq:klein-flat-coefficients}から$\delta_i=\alpha_i=0$となる。
今度は$e_i$での両評価が下側だけで階数3となるため
$b_jb_kc_jc_k\ne0$、従って$c_i=0$である。
$c_i=0$から始めても転置で同じである。$b,c$が全て非零なら六式から他も全て非零になる。
残る二つの$A$は文字の三頂点上で二辺だけを支えるため、両方を対角変更で同時に対称にできる。
六式で$c_j/b_j=c_k/b_k$となるので、適切な奇単位の倍率と
$H$の下側計量の倍率を選ぶと再び両対称表示を得る。
以上で単純な頂点に関する被覆が従う。
\begin{remark}
全係数非零の$\epsilon=1$側も常に標準ASである。$L=\operatorname{diag}(\alpha_i)$として
上側の積は$D(Lu,y)$である。両核は
$u_i^2-u_j^2+(\alpha_j^2-\alpha_i^2)r^2$の張る同じ鉛筆である
（反対切断の平方以外の三ブロックは可逆二行交代行列、平方ブロックは階数2）。
三平方が相異なるときは対角完全交叉が滑らかである。対角鉛筆の特異な各元は
単独の座標頂点を頂点に持ち、その点はもう一式で消えないためである。
純関係もない。上側・下側の非零な二係数ベクトルをともに使えば、
零の$00$成分を挟む二行二列で階数2以上である。片側だけなら非零の
$D$の極行列が非零対角倍率を除いて含まれ、その非零成分の二行二列小行列式がある。
定理\ref{thm:smooth-apolar-complete}を適用する。
二平方が等しければ符号変更で$\alpha_i=\alpha_j$と取れる
（同時に三文字の対角符号と軸の倍率を変えればよい）。
因子$e_i+e_j,e_i-e_j$には逆順の同じ奇因子が対応する。
各三平面像は$e_i\pm e_j$の上側と下側および
$(e_k,e_k)$を張る（上下の非零な共通倍率を許す）。
よって交わりは直線であり、純関係がないので共通根基の二文字関係は可逆である。
定理\ref{thm:radical-factor-amalgam}からこの場合も正則となる。
\end{remark}

\paragraph{横断する頂点がない場合。}
四頂点の辺$ab$を$b_{ab},b_{ba}$が比例するとき選ぶ。頂点での二軸の交わりの次元は
そこに出入りする選択辺の本数である。従って各頂点に選択辺があれば、
互いに素な選択二辺があるか、選択辺が三辺の星をなすかのいずれかである。

\begin{lemma}[相補的な二つの選択辺]\label{lem:klein-matching-pbw}
相補的な二辺が選ばれた場合、座標評価下界と純関係の不在の下で窓は標準Koszul AS正則である。
\end{lemma}
\begin{proof}
二辺を$01,23$に取り、文字を$\{0,1\}$と$\{2,3\}$の二組に分ける。
この色の関係は各組に一本ずつの二項交換関係となる。
実際、二組がそれぞれ同一直線に写り、二直線は独立で全係数が非零である。
相手窓でも零化直線の表示により同じである。
残る四色成分の関係（他の二色で各二式）の、後の組を左、先の組を右に持つ
四単項式への射影は可逆である。例えば$b_{20},b_{31}$の独立性は
$\langle b_{20},c_{13}\rangle\ne0$と同値である。偶位置交換で
これは$\langle b_{10},c_{23}\rangle$となり、独立な二辺の直線の一つと
他方の零化共変ベクトルの対合だから非零である。$b_{30},b_{21}$、逆向きの二組、
$c$側も同じ理由による。従って四つの混合逆順を正順へ全て交換できる。
組内も二項式で整列する順序では、二次関係は全ての六逆順を先頭語とする。
次数3の正規候補語は20個である。一方テンソルの一文字切断は全階数である
（異なる次数の四評価が非零で、六次元への相手も全射）。
その四つの縮約は三次交叉に入るので三次商の次元は20以上となる。
従って三次で全ての重なりが解消され、全次数で二次Gr\"obner基底である。
このKoszul性と標準Hilbert級数、Gramテンソルの切断の完全対合から標準AS性を得る。
\end{proof}

最後に星なら、三色とも一方の辺が選択されている。相異なる奇添字$b,d$について
$m_{bd}=w_{abcd}$と置く。ここで$a,c$は残る二文字で、偶交換により順序によらない。
色$i$の二辺を$ab,cd$とすると、$b_{ab},b_{ba}$から$c_{cd},c_{dc}$への対合の
二行二列行列は、どちらかの辺が選択されていれば階数1である。
四成分は全て非零である。実際、選択辺の直線に残り一文字の値が入ることは
三つの比例による純関係障害で排除される。逆向きの二行二列行列も同様である。
これらは$m_{ij}$のうち相異なる四添字の二行二列式の全六本にあたる。
従って$m_{bd}=u_b v_d$ ($b\ne d$)と分離できる。
実際、行・列を頂点とする二部グラフ$K_{4,4}$から対角を除くと立方体の辺となり、
六つの小行列式は各面一周の倍率が1という条件である。六面の境界は全ての閉路を
整数係数で生成するので、任意の道で行・列の倍率を決められる。
奇文字$d$の一側を$u_d/v_d$倍すれば$m_{bd}$は$b,d$で対称である。
このほかの非零成分は偶・奇の各平方の対合だけなので、両対称表示を得る。
これで定理\ref{thm:klein-missing-character}の残りも証明された。

\begin{proposition}[三軸族の外の反対階数10]\label{ex:klein-matched-rank-ten}
交代行列$m$の上三角成分を辞書順に$(2,1,1,3,3,5)$とし、
偶平方・奇平方の対合を
\[
 T=(w_{abab})_{a,b}=
 \begin{pmatrix}0&0&1&1\\0&0&-35/2&-65/2\\
 2/15&21&0&0\\-2/15&-39&0&0\end{pmatrix}
\]
とする。四添字が相異なるときは$w_{abcd}=m_{bd}$、ほかは零とする。
三つの隣接色ブロックは各々階数2である。実際、辺を$ab=0h,cd$とし
$cd$を昇順に取ると、順序$ab,ba,cd,dc$でブロックは
$\left(\begin{smallmatrix}U&P\\-P^{\mathsf t}&V\end{smallmatrix}\right)$であり、
$P=m_{(b,a),(d,c)}$である。$h=1$なら$U=V=0$、$P$は全成分非零の階数1である。
$h=2,3$では$U$が可逆で$V=-P^{\mathsf t}U^{-1}P$となる。
各方向の関係行列で六逆順への射影が可逆であり、組内二関係は二項式である。
補題\ref{lem:klein-matching-pbw}の整列証明から全次数で標準ASとなる。
反対切断では平方ブロックの行列式が36で、残る六つの交代成分が独立だから階数10である。
各組の座標直線上では相手の同じ直線との間に非退化な二項式が一本だけ残る。
従って偶点スキームはねじれの位置にある二直線を含む。
三軸族の反対階数10では、全ての曲線成分が一つの平面内にある
（節\ref{sec:cross-product-parameters}）。この例は印を忘れても三軸族に入らない。
\end{proposition}

\begin{proposition}[次数次元が$(3,1,1,1)$の場合]\label{prop:klein-three-singletons}
四次数が全て現れ、その一つが三次元ならば、この斉次類の標準AS窓は
全対称表示を持つ。従ってこの型の任意の入力には全対称な完全交叉の判定が使える。
\end{proposition}
\begin{proof}
三次元の次数を零に移す。相異なる$a,b$について$b_{ab}=0$と仮定する。
残りの二つの$c$各々に$d=a+b+c$を取ると、偶対称式は
$0=\langle b_{cb},c_{ad}\rangle$となる。第$b$列の評価階数が3以上なので
両方の$b_{cb}$は非零である。この次数空間は一次元なので両方の$c_{ad}$が零となり、
$\gamma$の第$a$行の階数に反する。転置でも同じであり、全ての非対角成分は非零である。
各一文字空間の基底と双対基底を選ぶ。固定した$d$に対し残る$a,b,c$を巡回させて
$c_{ad}/c_{cd}=b_{ab}/b_{cb}$を掛けると、
$b_{ab}b_{bc}b_{ca}=b_{ba}b_{cb}b_{ac}$となる。
従って辺の比$b_{ab}/b_{ba}=u_a/u_b$と書ける。奇軸$f_b$を$u_b$倍すると
$\beta$は二文字に対して対称となる。零次数の対角成分はこの交換で動かない。
偶対称性に加えて隣接交換も得たので、系\ref{cor:two-positive-as-recognition}と
系\ref{cor:partial-symmetric-as-complete}を適用できる。
\end{proof}

% !TeX root = ../papers/nc-p3.tex
\section{二つの二次元次数と四次元の片側対称族}
\label{sec:klein-two-doublets}
前節の残る次数次元$(2,2,1,1)$を扱う。次数$0,1$を二次元とし、
$A=\{0,1\}, B=\{2,3\}$、$\bar i=i+1$と書く（加法は四元群内）。
$b_{ij},c_{ij}$は前節と同じである。異なる組を結ぶときは一次元の
値を双対基底のスカラーで表す。

\begin{theorem}\label{thm:klein-two-doublets-generic}
この型には、幾何学的一般点が反対階数10の標準Koszul AS正則窓となる既約な族がある。
添字保持同型に関する像の次元は4である。一般の偶点スキームは被約な二十点で、
どの三点も共線でない。従ってこの族は、点の印も群作用の印も忘れて三軸族と区別できる。
\end{theorem}

\paragraph{二つの行列の共通非対角成分。}
必要な座標評価階数3以上を課す。同じ組の$i,j$で$b_{ij}=0$なら、
他方の組の両方の$k$について$b_{kj}\ne0$である。
$l=i+j+k$として偶交換を使うと、一次元での対合から$c_{il}=0$を二つ得て矛盾する。
$c$側も同じである。異なる組の$i,j$には
$b_{ij}c_{\bar i,\bar j}=b_{\bar i,j}c_{i,\bar j}$が成り立つ。
二つの列ベクトルはいずれも非零なので、$g_{\bar j}$の倍率を変えれば
$c_{i,\bar j}=b_{ij}$とできる。個別の零成分もこれで含まれる。
以下、薄い次数の二方向を行列$U=(b_{i,2+k})$, $V=(b_{2+k,i})_{k,i}$で表す。
$U$の行順・列順は$i,k=0,1$である。

隣接切断の厚い次数$h=0,1$の行列を
$M_h(i,k)=w_{i,i+h,k,k+h}$とする。偶交換は$M_0$と$M_1^{\mathsf t}$の
対角以外の成分が一致することを与える。同じ組の異なる$i,k$では偶交換そのものであり、
異なる組では双方が同じ薄い二係数の積になる。従って二行列は
\begin{equation}\label{eq:doublet-completions}
 M_0=\begin{pmatrix}S&X\\Y&T\end{pmatrix},\qquad
 M_1^{\mathsf t}=\begin{pmatrix}S'&X\\Y&T'\end{pmatrix},\qquad
 X_{ik}=V_{ki}U_{i,\bar k},\quad Y_{ki}=U_{ik}V_{k,\bar i}.
\end{equation}
局所添字は$0,1$である。$S,S'$と$T,T'$はそれぞれ非対角成分が等しい。
逆にこれと薄いブロックの積表示で全ての偶交換式を尽くす。
実際、二つの偶添字が同じなら恒等式であり、同じ組の相異なる二つなら厚い側の
行列の比較または薄いスカラー積の交換となる。異なる組なら厚い側の成分が薄い側の積に決まる。

$\det Y=-\det X$である。$\det X\ne0$なら階数2条件は
$S,S'$が可逆で$T=YS^{-1}X$, $T'=Y(S')^{-1}X$であることと同値である。
$S'_{00}=i,S'_{11}=j,t=\det S'$を未知とすると、
$Y\operatorname{adj}(S')X-tT$の二つの非対角成分は$i,j,t$の一次式である。
$i,j$に関して解ける開集合では$i,j$は$t$の一次式となり、
行列式条件は高々二次で、既知根$\det S$を持つ。真に二次で残る根も非零なら、
もう一つの解は有理式で一意に得られる。二根が等しい場合もこの有理式は正則であり得る。
$S'=S$なら$g_{\bar j}$と$f_j$の同定で奇交換が成立する。
実際、偶交換も用いると奇交換は厚い二行列が全成分で一致する条件となる。

\paragraph{正則な合流点。}
$U=\left(\begin{smallmatrix}1&2\\2&1\end{smallmatrix}\right)$,
$V=\left(\begin{smallmatrix}3&5\\5&3\end{smallmatrix}\right)$,
$S=\left(\begin{smallmatrix}2&3\\44/87&7\end{smallmatrix}\right)$とする。
二本の一次式は独立で二次式には重根$\det S$がある。
この点の両対称な反対切断は階数8で、両側の核の鉛筆は対角である。
たとえば偶側の二行は平方順に
\[
 (25029,606,-3620,0),\quad (29189,3296,0,-5430)
\]
であり、奇側では$(254616,42891,-392770,0)$と
$(346834,272484,0,-589155)$である。どの二列の小行列式も非零なので両完全交叉は滑らかである。

さらに$U,V$が全成分非零で$\det X\ne0$のとき、第一窓の関係空間は純テンソルを含まない。
関係を四行四列にし、薄い二関係の係数を$z_2,z_3$とすると、
組をまたぐ二つの二行列の行列式は
$U_{00}U_{11}z_2^2-U_{01}U_{10}z_3^2$と
$V_{00}V_{11}z_2^2-V_{01}V_{10}z_3^2$である。
係数行列は$\det X\ne0$によって可逆で、純テンソルなら$z_2=z_3=0$となる。
組内二ブロックの片方は零でなければならないが、
式\eqref{eq:doublet-completions}の両方の列空間が各組へ同型に射影するので他方も零となる。
従って合流点に定理\ref{thm:smooth-apolar-complete}を適用できる。

$U,V,S$の十二係数から上の第二根で作る族は既約で、この合流点を含む。
隣接階数は6、一文字切断は合流点の近くで全階数である。三次・四次の交叉には
テンソルとその縮約から下界があり、上半連続性により合流点の近くで階数は一定となる。
五次交叉は零である開条件を保つ。この有限Frobenius双対を使い、
節\ref{sec:reversible-component}の内部次数別のbar複体の消滅により、
第二根の幾何学的一般点で全次数のKoszul性と左右AS性を得る。
これだけで個々の有理閉点の全次数正則性を判定したことにはならない。

\paragraph{二十点、直線と同型類の次元。}
次の有理入力で反対階数は10である。
\[
 U=\begin{pmatrix}1&2\\3&5\end{pmatrix},\quad
 V=\begin{pmatrix}2&3\\5&7\end{pmatrix},\quad
 S=\begin{pmatrix}2&1\\5&3\end{pmatrix},\quad
 S'=\begin{pmatrix}123/196&1\\5&1335/176\end{pmatrix}.
\]
素数12619で第一評価行列の十五個の最大小行列式を用いる。最初の非零座標を1とした
四つの層で辞書式消去を行うと、幾何点が全て素体上にあり、相異なる二十点を得る。
全ての三点組で三行三列の小行列式の少なくとも一つが非零である。
\path{computations/verify_klein_two_doublets.py}が、上の二入力とこの消去とを検算する。
有限で期待余次元3なのでEagon--Northcott解消から長さは20であり、特殊ファイバーも被約である。
射影性から有限性はこの点の近くで開き、同じ解消と長さにより有限平坦、さらに有限
エタールとなる。その相対三重積のうち相異なる三点の部分は開かつ閉であり、
共線条件の像は閉である。よって標数零の幾何学的一般点にも三点直線はない。
三軸族では有限な場合、各座標面の四符号点の結線が座標頂点を通る。

最後に座標変換を数える。固定座標で薄い二ブロックのうち一方は
二つの四係数ベクトルの外積であり、他方も同じ二ベクトルから決まる。
従って固定テンソルの表示で残る自由度は、二つの薄い次数の値を互いに逆の倍率で
変える一次元だけである。$S$はテンソル成分から直接読める。
射影化も除くと、第二根の射影テンソルの像は$12-1-1=10$次元である。
印付き文字の基底変更は対角で、反対階数10から偶二位置の変換は射影的に一致する。
実際、その反対像は対称二乗全体で、二つの変換による全平方の像が再び対称でなければならない。
奇二位置も$g_{\bar j},f_j$について射影的に同じ倍率になる。
これは組をまたぐ二つの奇添字では、規格化したテンソルが奇交換で等しく、
その成分が全て非零だからである。逆に偶・奇でそれぞれ独立に同じ対角変換を施してよく、
有効な射影対角群は六次元である。

一般の正則点では両奇点スキームも偶側と同型で有限被約である。
三つの射影表示で二次式は点全体で消え得ない（最大小行列式の最小次数は4）。
従ってこれらの有限集合の射影安定群は有限である。実際、各点を固定する非スカラー行列の
固有直線全体は二つの超平面の和集合に入る。よってテンソルを保つ射影変換群も有限で、
四元群作用の選び直しは有限通りである。各文字が重複しないので、作用を固定した
基底の選び直しは上で数えた対角変換と有限置換で尽きる。
定理\ref{thm:frobenius-tensor-reconstruction}によって添字保持同型を取った像は
$10-6=4$次元となる。

% !TeX root = ../papers/nc-p3.tex
\section{四次元族の固有な四元群・周期と同型判定}
\label{sec:klein-generic-isomorphisms}
前節の第二解の族で、三つの点配置を$X_E,X_F,X_G$と書く。

\begin{theorem}\label{thm:klein-generic-period-four}
幾何学的一般点で各配置の射影安定群は指定された四元群に等しく、三配置は互いに
射影同値でない。この標準AS窓の最小周期は4である。
さらに明示的な四つの有理不変量を高々十六通りの文字置換で比較すれば、
この条件を持つ開集合上の正則窓の添字保持同型が判定できる。
\end{theorem}
\begin{proof}
前節の素数12619の入力で、偶二十点の両評価行列は階数3である。
核の余因子から各々の奇配置を計算する。
各配置には三点直線がなく、四点以上の平面が48枚あり、それぞれ丁度四点を持つ。
このうち八平面を通す点が四つ、十平面を通す点が十六ある。
前者の四点は線型に独立で、その四つの座標面から外れる点もある。
従って配置間の射影同型は、四点の置換と、固定した第五点の行き先とで一意に決まる。
これらを全て作用させて二十点の一致を調べると、$(E,F,G)$の順の
射影同型の個数行列は$\operatorname{diag}(4,4,4)$である。
全点が素体上で第五点までが射影枠なので、これは代数閉包上でも同じ個数である。
\path{computations/verify_klein_generic_isomorphisms.py}がこの有限計算を行う。

この比較は一般点での上界を与える。有限エタールな順序付け被覆で二十点に番号を付け、
各置換$\pi$について$D x_i\wedge y_{\pi(i)}=0$を
$[D]\in\mathbb P(\operatorname{Mat}_4)$に課す。解の像は射影性から閉である。
特殊点で階数$1,2,3$の行列の核が含む点はそれぞれ高々$4,2,1$個で、
像空間が含む点は高々$1,2,4$個である。この和は20未満なので特異解はない。
番号は両奇配置にも核で運べる。実際、特殊点で全評価核は直線で相異なる奇点を与え、
偶側の有限エタール被覆上の階数3性は近傍で保たれる。両行列から得る偶点スキームの
一致は一般正則点で前述の通りであり、ある開集合へ延びる。より直接には
第一窓の各零点対で$\gamma(x,-)$は$\beta(-,f)^\perp$に入る。
特殊点で$\beta(-,f)$も階数3であることは、薄い次数が全成分非零で交差行列が
可逆なこの例では直接評価でも確認できる。従って近傍で核の点列が定まる。
有限個の置換を比較することで、一般点では射影同型の個数は上で数えた値以下となる。
指定された対角四元群は各配置を保つため各対角成分は丁度4である。
周期2なら一方の奇配置を他方に移す線型同型が必要で、周期1でも同様である。
四窓表示による周期4から最小性を得る。

従って窓の同型は偶配置の唯一の四元群を保つ。奇配置の変換も、偶配置の作用と
そのグラフ上で整合するため群上の同じ自己同型を誘導する。各配置の点ごとの
射影固定群は自明である。四文字の置換はある$\phi\in\operatorname{Aut}(K)$で
$ i\mapsto\phi(i)+a,\ j\mapsto\phi(j)+b,\ l\mapsto\phi(l)+b$
となる。実際、射影表現の持ち上げの差は共通の一文字の平行移動である。
偶二位置の変換は反対階数10から一致し、テンソルが全次数零なので奇二位置の移動量は
等しい。厚い次数の組を保つ条件は$\phi(1)=1$, $a+b\in\{0,1\}$である。
よって候補は$2\cdot4\cdot2=16$通りとなる。

文字を固定した後は、$U,V$が全成分非零で$S_{00}\ne0$の範囲を取る。
薄い成分の変換は$e_i f_j t^{\pm1}b_{ij}$であり、次数2に正、次数3に負を取る。
$e_i,f_j$は偶・奇座標の倍率で、$t$は二つの薄い次数の逆倍率である。
$\rho_U=U_{00}U_{11}/(U_{01}U_{10})$と置くと$t^4=1/\rho_U$で
$U$の行列積比が1になる。二つの行列の行・列の倍率を独立に選び、
$U$の成分を全て1、$V$の右下以外を1にすると、その右下は
$z=\rho_V/\rho_U$となる。この倍率が$S$の第1行と第0行に与える比と
第1列と第0列に与える比は、それぞれ
\[
 L=\frac{U_{01}V_{00}}{U_{11}V_{01}},\qquad
 R=\frac{U_{00}V_{01}}{U_{10}V_{00}}.
\]
実際、行比は$e_1 f_1/(e_0 f_0)$、列比は$e_1 f_0/(e_0 f_1)$である。
二組の間の残る共通倍率で$S_{00}=1$にできる。従って正規形の四係数は
\[
 \left(z,\ R\frac{S_{01}}{S_{00}},\ L\frac{S_{10}}{S_{00}},
               \ LR\frac{S_{11}}{S_{00}}\right).
\]
どの四乗根を選んでもこの表示は変わらない。逆に四係数は規格化した$U,V,S$を与え、
第二解は有理式で決まる。以上の分母と第二解の構成が十六置換の全てで有効な
開集合を取れば、置換して四係数を比較する必要十分条件となる（各置換は族の
一般点を保つ。厚い行列の各組内ブロックは一般に全成分非零で可逆である）。
正則性の主張は前節の幾何学的一般点、同型の判定は条件を満たす全正則入力に関する。
\end{proof}

% !TeX root = ../papers/nc-p3.tex
\section{二重次数の全境界と第二解の閉包}
\label{sec:klein-doublet-boundary}
\begin{theorem}\label{thm:klein-doublet-complete-cover}
次数次元$(2,2,1,1)$の片側対称な標準AS窓は、両対称チャートまたは
定理\ref{thm:klein-two-doublets-generic}の第二解の閉包に属する。
式\eqref{eq:doublet-completions}の$\det X=0$の全入力では両対称チャートの
判定で正則性が決まる。$\det X\ne0$では薄い次数の成分が零の場合を含め、
相異なる全ての行列補完が同じ既約族の閉包に入る。
後者の個々の特殊化での正則性は別に判定する必要がある。
\end{theorem}

\paragraph{交差ブロックの階数が落ちる場合。}
まず薄い次数の全成分が非零なら、前節の倍率変更で$U,V$の全成分を1にできる。
実際、$\det X=0$は二つの積比が等しい条件である。
$X,Y$は全成分1の行列$J$になる。$h=0,1$に対し、$M_h$の列空間の
$A,B$への射影階数を$u_{Ah},u_{Bh}$、行空間について$v_{Ah},v_{Bh}$とする。
各階数は1または2で、階数1なら像は$(1,1)$の直線である。
$u_{Ah}=v_{Bh}=2$は同時に起きない。階数2の因子分解を$A$行、$B$列に
制限すると可逆な二行列の積になり、交差ブロックの階数1に反する。
$u_{Bh},v_{Ah}$についても同様である。

$u_{A0}=u_{A1}=1$なら、$x=e_0+e_1$で偶評価すると、厚い二次数からの
奇$A$上の二式は共に$(1,1)$の倍数であり、薄い次数からの奇$B$上の
二式も比例する。評価の総階数は2以下である。他の組と相手窓についても同様である。
従って$u_{A0},u_{A1}$のうち一方だけが2であり、$v_{B0},v_{B1}$では
他方だけが2である。もう一つの組合せ$u_B,v_A$もそうなる。
もし$u_A,u_B$の階数2の色が同じなら、色1で両組の列射影がともに1か、
両組の行射影がともに1である。該当する窓の$01,10$の係数は同じで、
$23,32$も同じになる。色0はもとから対角、薄い次数の係数は全て1なので、
この窓は隣接対称である。二つの正交換から標準ASでは完全対称となる。

色が反対なら、たとえば色0で$A$の行・列射影がともに2、$B$がともに1である。
色1では逆になる。従って$T,S'$が$J$の定数倍であり、$S,S'$および
$T,T'$の非対角成分の一致から$S,T'$も対称行列である。組や色を逆にしても同じである。
$M_0,M_1$と薄い二行列は対称で、テンソルには二文字ずつの巡回交換がある。
これと偶交換を合成すると奇交換が得られる。

次に薄い成分に零を持つ場合、各行・列で零は高々一つである。$U_{00}=0$としてよい
（両組や組内のラベルを変え、必要なら薄い二色を交換する）。
$\det X=0$から$V_{00}V_{11}=0$である。
同じ薄い色の四成分中三つが零なら、残る一成分が純関係を与えるか隣接階数が落ちる。
従って両行列に丁度一つずつ、同じ色に零がある。$V_{11}=0$なら
$\gamma$の二行列で列を反転した零点同士が転置位置にあるため、
窓を取り替え、組内を表示し直せばよい。従って
\[
 U=\begin{pmatrix}0&u\\v&w\end{pmatrix},\quad
 V=\begin{pmatrix}0&p\\q&s\end{pmatrix},\quad
 X=\begin{pmatrix}0&0\\pw&sv\end{pmatrix},\quad
 Y=\begin{pmatrix}0&0\\us&wq\end{pmatrix}
\]
として六係数は非零である。$M_1^{\mathsf t}$の両組への列射影はともに2である。
もし1なら、交差ブロックの列方向が$(0,1)$なので一行全部が零になり、
厚い次数の座標評価成分が零になる。ゆえに$M_1$の両組の列射影はともに1で、
その二方向は$(us,wq),(pw,sv)$である。
奇座標を$f_{B1}/f_{A0}=q/u$, $f_{B0}/f_{A1}=p/v$,
$f_{B1}/f_{A1}=s/w$の倍率比で変える（ここだけ$f$は倍率）。
これは薄い二行列を転置位置で対称にし、同時に
$f_{A1}/f_{A0}=wq/(us)$, $f_{B1}/f_{B0}=sv/(pw)$により厚い色1でも
隣接対称にする。従ってこの場合の正則入力も完全対称である。

\paragraph{可逆な交差ブロックと全ての補完。}
$h=\det S'/\det S\ne0$とする。$h\ne1$では
$Z=(S'-hS)/(1-h)$と置くと、$S=Z+\operatorname{diag}(v,w)$,
$S'=Z+h\operatorname{diag}(v,w)$である。二行二列の余因子行列は線型なので、
下側の非対角成分の一致は
$Y\operatorname{adj}(Z)X=\operatorname{diag}(s,t)$と同値となる。
従って$Z$は$s,t$から一意に決まり、$\det Z=st/(\det X\det Y)$である。
二つの行列式の比を取る残りの条件はただ一つ
\[
             h v w=\det Z
\]
となる。この表示は可逆性を課す開部分で逆写像も持つ。基底$U,V$の八係数と
$h$について$h\det X\det Y$は単元であり、この方程式は$vw$と$st$を
結ぶ既約な二項式である（$v$について原始な一次多項式）。
空でない開条件を取り除いても既約であり、一般の薄い全成分非零の第二解が稠密である。

$h=1,S'\ne S$では$S'-S=\operatorname{diag}(m,n)$で、
$Y\operatorname{adj}(S'-S)X$の非対角が零になる。
$U=\left(\begin{smallmatrix}u&v\\w&d\end{smallmatrix}\right)$,
$V=\left(\begin{smallmatrix}p&q\\r&s\end{smallmatrix}\right)$と書くと、
$(m,n)$にかかる二行列の行列式は
$(psvw-rqud)(psvw+rqud)$となる。第一因子が$\det X\ne0$だから
第二が零で、八係数は全て非零である。前節の規格化で$U=J$,
$V=\left(\begin{smallmatrix}1&1\\1&z\end{smallmatrix}\right)$とすれば$z=-1$である。
$S=\left(\begin{smallmatrix}a&b\\c&d\end{smallmatrix}\right)$と書き直す。
このとき非対角条件は$m=n$で、行列式を保つ非自明解は
$(a',d')=(-d,-a)$で尽きる。

これも同じ閉包に入ることを直接示せる。$z\ne\pm1$で$k=(zb+c)/(z+1)$とすれば
$Y\operatorname{adj}(Z)X$を対角にする中心行列$Z$は
$Z_{00}=Z_{11}=k$, $Z_{01}=b,Z_{10}=c$である。従って第二解は
\[
 a'=k+\frac{k^2-bc}{d-k},\qquad d'=k+\frac{k^2-bc}{a-k}.
\]
$b\ne c$なら$z\to-1$で$k$は極を持ち、二値は$-d,-a$に延びる。
$b=c$の同じ解は$z=-1$で$c$を$b$に近づけて得る。
これで定理\ref{thm:klein-doublet-complete-cover}が従う。

\begin{remark}
最小周期4の主枝で正規パラメータを探すなら、次数にも制約がある。
各位置にある正則な正規元の族$q_i\in A_{i,i+n}$が
$A_{ij}q_j=q_i A_{i+n,j+n}$を全ての$i,j$で満たせば、
$a q_j=q_i\sigma(a)$で一意に定まる$\sigma$は乗法的な線型全単射である。
従って$n$は周期であり4の倍数になる。特にこの枝で先の二周期窓の
正則二次パラメータによる証明をそのまま適用することはできない。
\end{remark}

\begin{theorem}[有理閉点での最小周期4の正則性]\label{thm:klein-rational-period-four}
節\ref{sec:klein-two-doublets}で二十点の消去に使った有理入力自身が、
最小周期4、反対階数10の標準Koszul AS正則窓である。
四頂点で折り畳んだ環は四つの中心四次式の多項式環上で有限自由であり、
各始点での自由階数は256である。
\end{theorem}
\begin{proof}
\path{computations/verify_klein_period_four_center.py}で次の有限証明書を再生成する。
各内部次数・文字次数で左端に一文字を付け、先頭二文字の関係に既知の商基底を
右から掛けて掃き出す。次数5までの元の次元は各始点で$1,4,10,20,35,56$である。
各位置の四次式$q_p$の係数に$ q_p v_i=v_i q_{p+1}$を課す。添字$p$は4を法とする。
有理数体上の同次解空間の次元は、四つの文字次数順に$(3,2,0,0)$である。
計算コードに指定した有理核基底について、次数零の側から係数$(1,3,5),(2,7,11)$の二式、
次数1の側から基底二式を取る。各四式の係数の分母を払い原始整数化する。
これらは折り畳み環の中心に入る。

四式を関係に追加し、各始点で二次関係と四次関係の両方に右から商の基底語を掛けて
同じ次数別掃き出しを行う。素数1009で得る商の次元は次数0から13まで
\[
 1,4,10,20,31,40,44,40,31,20,10,4,1,0.
\]
従って標数零での各次元にはこれが上界となり、一次生成性で以後は零である。
この上界の母関数は$(1+t+t^2+t^3)^4$である。中心性により商の基底を持ち上げると、
重み4の四変数多項式環上の生成系を得る。その自由加群による上界は$(1-t)^{-4}$となる。
一方、当該有理点は正則な幾何学的一般点を持つ第二解の族に属し、各次数で関係を
挿入する行列の余核次元は特殊化で減少しない。従って全次数で上下界が等しく、
四式は代数的に独立で、この生成系は自由基底である。

四階テンソルの一文字と隣接の切断はそれぞれ全階数4と6である。
四頂点で折り畳んだ右Koszul列の最初の二微分は表示から完全で、最後の二微分は
テンソルの完全対合により左の表示の双対である。
多項式環上で自由なので、始点単純から環への$\operatorname{Hom}$と
$\operatorname{Ext}^1$は零となる。実際、左増大イデアル上の写像$f$に対し、
中心変数のうち二つを$u,v$として$u f(v)=v f(u)$を使えば、
互いに素な変数であることから$f(u)=u b$となり、$u f(z)=z f(u)$で
$f(z)=zb$が従う。よって双対側も上の二項で完全で、中央はHilbert級数の
Euler等式から完全になる。左右を逆にしても同様である。
これで全次数のKoszul性と左右AS性を得る。構成した有理点で三配置の非同値性は
節\ref{sec:klein-generic-isomorphisms}の同じ素数12619での比較から直接従い、
反対階数と最小周期も主張の値になる。
\end{proof}

% !TeX root = ../papers/nc-p3.tex
\section{四次中心の持ち上げと閉点の有限判定}
\label{sec:klein-quartic-center-lift}
第二補完の既約族のうち、上側の中心行列の対角を1に規格化できる開集合を取る。
節\ref{sec:klein-doublet-boundary}の記号で
\[
 z=\frac{1-c}{b-1},\quad Z=\begin{pmatrix}1&b\\c&1\end{pmatrix},
 \quad S=Z+\operatorname{diag}(v,w),\quad
 S'=Z+\frac{1-bc}{vw}\operatorname{diag}(v,w)
\]
である。$U$は全成分1で、$V$は右下の$z$以外が1である。
分母と二つの厚いブロックの行列式は非零とする。

\begin{theorem}\label{thm:klein-quartic-center-slice}
$v=4,w=6$の二次元断面では一般の四次中心の文字次数別次元は$(3,2,0,0)$である。
$(b,c)=(2,3)$を含む単一のザリスキー開集合上、全ての閉点が標準Koszul左右AS正則で、
各始点の代数は四つの中心四次式の多項式環上で階数256の自由加群となる。
\end{theorem}
\begin{proof}
厚い二次数について、各六関係の基底を先の組への射影が単位行列となるよう選ぶ。
$X=\left(\begin{smallmatrix}1&1\\1&z\end{smallmatrix}\right)$,
$Y=\left(\begin{smallmatrix}1&1\\z&1\end{smallmatrix}\right)$とすると、後の組への
二射影は第一窓で$YS^{-1},(S'^{-1}X)^{\mathsf t}$、相手窓で
$(S^{-1}X)^{\mathsf t},Y(S')^{-1}$である。これは
式\eqref{eq:doublet-completions}の列空間・行空間を取ったものである。
各始点の次数4と5の商基底を作り、内部次数4の元$q_p$に
$q_p x_i=x_i q_{p+1}$を課す。\path{computations/verify_klein_center_lift.py}は
$\mathbb Q(b,c)$上で関係を挿入して厳密に掃き出す。文字次数0では36列の階数33、
文字次数1では40列の階数38である。上記の有理点では全四次数の核次元が
$(3,2,0,0)$で、各始点の内部次数5までの商次元も標準値である。
第二補完の全ての特殊化で標準Hilbert関数が下界なので、次数5までの商はこの点の
近傍で次数別・文字別にベクトル束である。従って他の二文字次数の核は一般に零である。
同じ近傍を縮めれば、五つの中心元の空間は次数別の部分束である。

この有理点で文字次数0から二式、文字次数1から二式を取った商の次元は
定理\ref{thm:klein-rational-period-four}と同じ素数1009の計算で
$(1+t+t^2+t^3)^4$の係数以下で、次数13で零になる。
定階数の核束を用いて四式を局所切断に持ち上げる。商の次数13までの次元の上界は
有限個の行列の階数下界なので共通の開近傍で保たれる。中心性による多項式環上の
自由加群からの全射と標準Hilbert関数の下界を比べれば全次数で一致する。
同定理の双対完全性の証明は同じ四頂点でそのまま適用できる。
\end{proof}

\begin{proposition}\label{prop:klein-central-strata}
第二補完の閉包の任意の入力で、隣接階数6と一文字階数4を仮定する。
四つの中心四次式で上記の自由加群表示を得る点の集合は、有限個の階数計算で
判定できる構成可能集合であり、その全点で標準AS正則となる。
\end{proposition}
\begin{proof}
内部次数5までの商の階数を層別化し、各定階数小行列の上で四次の中心方程式の
階数を計算する。各核の基底の一次結合を四つ、係数を独立の不定元として作る。
四式を加えた商の次元を次数13まで計算し、各始点で上の母関数の係数以下かを判定する。
これは係数について開条件であり、無限体上で一組が存在することと有理関数体上で
成立することは同値である。基底や斉次な文字次数の選択を追加制約とする必要はない。
成立すれば前定理と同じ上下界が使える。逆に標準Hilbert級数を持つ中心四変数上の
自由表示があれば、その変数を割った母関数は必ず上の多項式である。
この条件が成立しない入力のAS性は、この判定だけからは従わない。
\end{proof}

% !TeX root = ../papers/nc-p3.tex
\section{四変数全体での中心四次式とPI恒等式}
\label{sec:klein-generic-central-parameters}
\begin{theorem}\label{thm:klein-generic-central-parameters}
四次元の第二補完族の幾何学的一般点で、四頂点環の四次中心の文字次数別次元は
$(3,2,0,0)$である。定理\ref{thm:klein-rational-period-four}の有理点を含む
単一の稠密開集合の全閉点で、中心四変数上の階数256の自由表示と標準Koszul左右AS性が成立する。
\end{theorem}
\begin{proof}
前節の二つのグラフ行列を$\mathbb Q(b,c,v,w)$上で用いる。
\path{computations/verify_klein_generic_quartics.py}は四変数を代数的に独立のまま保ち、
各次数の関係行に共通する分母を払い、行の多項式最大公約数を除いて掃き出す。
各始点で次数0から5の次元は$1,4,10,20,35,56$である。
次数4から5への中心方程式は文字次数0で36列・階数33、次数1で40列・階数38となる。
分母は有理関数体の零でない元だけで、階数は有限体や数値近似から推定していない。
残る二文字次数は前出の有理点で核が零であり、そこで次数5までの商次元も一般値なので
一般点でやはり零である。各文字次数の核次元が有理点でも一般値となるため、
その点の近くで五次までを局所自由化すれば中心元の核も部分束になる。
以前に選んだ四式をこの部分束の局所切断に延ばし、
命題\ref{prop:klein-central-strata}の有限個の商次元上界が保たれる開集合を取ればよい。
この議論は基底の規格化に依存せず、行列補完の既約閉包の定階数部分に適用できる。
\end{proof}

\begin{corollary}\label{cor:klein-pi-specializations}
第二補完の閉包の任意の標準ASファイバーで、四頂点環は次数512の標準多項式恒等式
$s_{512}=0$を満たす。従ってこの閉包には非PIの標準AS窓はない。
\end{corollary}
\begin{proof}
上の開集合で四つの中心四次式の多項式環を$P$、四頂点環を$B$とする。
自由基底の次数多項式$(1+t+t^2+t^3)^4$では、指数を4で割った各剰余類の係数和は64である。
中心四変数が位置を変えないため$e_i B e_j$は全て$P$上自由階数64、
$B e_j$は自由階数256である。四つの$B e_j$上の左作用を合わせると忠実なので、
$B$は四つの$256$次行列環の直積に埋め込まれる。次数$2n$の標準恒等式を$n=256$に適用する。実際、外部変数$u_i$に対し
$T=\sum_{i=1}^{2n}u_i A_i$と置くと、奇元の巡回符号で$\operatorname{tr}(T^{2j})=0$である。
偶部分は可換なのでNewton恒等式とCayley--Hamiltonから$(T^2)^n=0$となり、
$u_1\cdots u_{2n}$の係数が所望の交代多項式である。
族の任意の標準AS点で任意の有限な内部次数まで、経路商の次元は一般値に等しい。
従ってその次数範囲では関係挿入行列の像は局所的に部分束である。
$s_{512}$の変数に固定した斉次な語を代入して得る有限次数の切断は、
一般点で零なのでこの近傍でも零となる。多重線型性で任意の代入を得る。
\end{proof}

% !TeX root = ../papers/nc-p3.tex
\section{偶交換を外す楕円曲線と最大反対階数}
\label{sec:klein-elliptic-completion}
四位置の文字空間を独立な四元群の正則表現とし、テンソルの総文字次数を零とする。
偶交換を仮定しない。前節までの片側対称点の近くにある二つの隣接切断を
\[
 M_h(i,k)=w_{i,i+h,k,k+h},\qquad
 N_h(i,k)=w_{k,i+h,i,k+h}=M_{h+i+k}(k,i)
\]
と書く。$M_h,N_h$はいずれも$h=0,1$で階数2、$h=2,3$で階数1とする。

\begin{theorem}\label{thm:klein-elliptic-completion}
二つの隣接切断に上の文字次元を持つ窓には、反対階数16、最小周期4で
幾何学的一般点が標準Koszul左右AS正則となる既約な族がある。
印を忘れた添字保持同型類の像は五次元である。その標準形の最後の一変数は
一般の四つの係数に対して滑らかな種数1曲線を動く。
この族は、以下の非退化な表示の一般の基底係数の上にある全解を覆う。
\end{theorem}

\paragraph{二切断の同時規格化。}
二つの薄い色の全成分が非零とする。$M_2,N_2$は対角が等しい階数1行列なので、
偶二位置の相対的な対角基底変更で$M_2=N_2$とできる。この後の偶倍率は
両位置で射影的に同じである。$N_3$も$M_3$の対角共役であり、
各組$A=\{0,1\},B=\{2,3\}$内の非対角成分は
$N_2=M_2$から等しい。従って共役の対角倍率は$A$上で$\lambda$、$B$上で1と取れる。

さらに$M_h(i,k)=u_i^h v_k^h$と書けば、$M_2(i,\bar i)=M_3(\bar i,i)$から
$u_i^2v_{\bar i}^2=u_{\bar i}^3v_i^3$となる（上付きは色を表す）。
よって奇位置の各列の倍率で、両色について$v_k^h=u_k^{h+1}$にできる。
同じ奇列に現れる二つの条件の整合性がまさにこの等式である。
薄い係数は以前と同じ二行列$U,V$で表され、前と同じ倍率変更で
$U$を全成分1、$V$を右下$z$以外1とできる。

$A_0=\left(\begin{smallmatrix}a&b\\c&d\end{smallmatrix}\right)$,
$A_1=\left(\begin{smallmatrix}e&f\\g&h\end{smallmatrix}\right)$を$M_0,M_1^{\mathsf t}$の
左上とする。残る倍率で$A_0$の共通スカラーを除ける。
$X_l=\left(\begin{smallmatrix}1&1/l\\1/l&z\end{smallmatrix}\right)$,
$Y_l=\left(\begin{smallmatrix}1&l\\lz&1\end{smallmatrix}\right)$,
$F_l(A)=Y_l A^{-1}X_l$と置く。
$M_0$の交差ブロックは$X_\lambda,Y_\lambda$、
$M_1^{\mathsf t}$では$\lambda X_\lambda,Y_\lambda/\lambda$となる。
これは行と列が異なる組にあるとき、$N$の薄い側で
色3の成分だけが組の向きによって$\lambda^{\pm1}$倍されるためである。
従って両者の右下は$T=F_\lambda(A_0),T'=F_\lambda(A_1)$である。
$N_0,N_1^{\mathsf t}$の左上は$C=\left(\begin{smallmatrix}a&f\\g&d\end{smallmatrix}\right)$,
$D=\left(\begin{smallmatrix}e&b\\c&h\end{smallmatrix}\right)$となる。
これらの右下は$F_1(C),F_1(D)$でなければならず、同時にそれぞれ
$T$の対角と$T'$の非対角、$T'$の対角と$T$の非対角を組み合わせた行列である。
以後、使う各二行列の行列式と表示の分母が非零な部分で考える。

\paragraph{一つの判別式を持つ二次方程式。}
$x=\det C=ad-fg$, $y=\det D=eh-bc$とする。
$F_1(C)$の対角二式と$F_1(D)$の非対角二式は独立に解け、
\[
 f=f_0+f_1x,\quad g=g_0+g_1x,\qquad e=e_0+e_1y,\quad h=h_0+h_1y
\]
となる。具体的には
\[
\begin{split}
 f_0&=(a+d)/(z+1),\quad g_0=z f_0,\quad
 f_1=(T_{00}-T_{11})/(z^2-1),\quad g_1=-T_{00}-f_1,\\
 e_0&=h_0=(zb+c)/(z+1),\quad
 e_1=(zT_{01}-T_{10})/(z^2-1),\quad h_1=(T_{10}-e_1)/z .
\end{split}
\]
$P(x)=fg+x-ad=p x^2+m x+n$、$Q(y)=eh-y-bc=q y^2+k y+o$とする。
両者の判別式は等しく、$\delta=m^2-4pn=k^2-4qo$である。
二根を
\[
             2p x+m=s,\quad 2q y+k=-s,\qquad s^2=\delta
\]
と組ませると、残る四つのSchur補の等式も全て従う。
これらは六つの独立な係数$a,b,c,d,z,\lambda$の有理関数体上の恒等式である。
\path{computations/verify_klein_elliptic_completion.py}は、第一に二つの判別式を引き、
第二に$y=-(m+k)/(2q)-(p/q)x$を入れ、
残る等式の分母を払って$P(x)$で割った余りの定数項・一次項を照合する。
$\det A_1=x+y-ad+bc$なので二次拡大での行列の逆計算は不要である。

$\lambda^2\delta$は$\lambda$の高々四次の多項式である。実際、$f_1,g_1$は
$\lambda,1,\lambda^{-1}$の線型結合である。$(a,b,c,d,z)=(5,2,3,7,-2)$では
この多項式の7569倍は
\[
 31024\lambda^4-17712\lambda^3-22355\lambda^2+8856\lambda+7756
\]
で、四根は相異なる。一般の基底係数でも四つの単純分岐点があり、滑らかな射影モデル
の種数は1である。二重被覆は幾何学的に既約である。
もう一方の符号による組合せは一般には残る四式を満たさない。
例えば同じ五係数で$\lambda=2$とし、$2q y+k=2p x+m$を代入すると
対角の最初の等式の余りは非零になる（同コード）。$P$が一般に既約なので、
この組合せの解が一般の基底へ支配的に写ることはない。

\paragraph{正則性と同型類の次元。}
$\lambda=1$では$f=b,g=c,x=ad-bc$の有理切断があり、
上記の符号で$(e,h)$は片側対称族の第二補完の対角を与える。
実際、この切断における$P,Q$の符号関係は二つの有理関数の平方の一致である。
$(5,2,3,7,-2)$で$s=1$, $(e,h)=(1/6,-1/4)$となることで符号が決まる。
従って四次中心により正則と証明した点がこの既約族上にある。有限Frobenius性の
開条件と内部次数別のbar複体を用いた従来の特殊化論法により、幾何学的一般点も標準AS正則である。
反対階数16は下の例で達するため一般値である。
階数16では、どの二位置の正または負の交換が反対対にある場合も、その対の
像次元を10または6に抑えるので不可能である。隣接対の交換が可能なら残る二分割の
切断階数が等しくなるが、6と16なのでこれも不可能である。

片側対称切断の一般点で既知の被約二十点配置は、どの三点も共線でなく、
五点平面も持たない。後者は節\ref{sec:klein-generic-isomorphisms}の特殊点での
全平面の計算と有限エタール順序付け被覆から従う。
同節の射影行列の固有な比較を各辺に適用でき、両奇配置の非同値性が近傍で
保たれる。従って一般の最小周期は4である。また各点配置の射影安定群は有限なので、
四元群の印を忘れるファイバーは有限となる。上の規格化後は、
薄い次数の二つの積比から$z$が決まり、$N_3/M_3$の対角共役比から$\lambda$が決まる。
$M_2=N_2$が偶二位置の相対倍率を、$M_2=M_3^{\mathsf t}$が対応する奇二位置の
相対倍率を固定する。文字を固定した残る対角作用は前節までと同じで、
$A_0$の共通スカラーを除けば有限である。従って基底の四変数と楕円曲線の
一変数を合わせ、印なしの同型類の像は五次元となる。

\begin{proposition}\label{prop:klein-rational-full-opposite}
$(a,b,c,d,z,\lambda)=(1,1,1,2,-2,-2)$、$s=-27/2$から得る有理窓は、
反対階数16の標準Koszul左右AS正則窓であり、四つの中心四次式上で各始点が自由階数256となる。
\end{proposition}
\begin{proof}
$(e,f,g,h)=(-13/14,1/13,-14,-16/11)$となり、反対切断の行列式は
$-29506135761/19712$である。隣接切断はともに階数6、一文字切断は反対切断が
可逆なので全階数である。四辺の関係空間は$M,N$の列空間と行空間から読み、
奇始点では二文字を反転する。従って片側対称な場合の反復関係を仮定する必要はない。
実際の四窓について中心計算と商計算を再実行すると、文字次数順の中心次元は
$(3,2,0,0)$で、定理\ref{thm:klein-rational-period-four}と同じ選び方で
四式を割った商は全始点で次数13まで$(1+t+t^2+t^3)^4$の係数以下である。
幾何学的一般点の標準Hilbert関数が下界となるため、先と同じ自由加群と
Koszul列の双対の議論で全次数の主張が従う。
\end{proof}

\begin{corollary}\label{cor:klein-four-configurations}
上の反対階数16の有理点自身で、四位置の被約二十点配置は全て互いに射影非同値で、
それぞれの射影安定群は四元群である。従って最小周期は4であり、
両方の反対対で二次Veronese表示が相異なる。これらは楕円曲線族の一般点にも成立する。
\end{corollary}
\begin{proof}
素数2657で四辺それぞれの最大小行列式を四つの射影座標層で消去する。
全四位置で幾何点は相異なる素体上の二十点であり、核を四回取ると元の番号に戻る。
期待余次元の長さ20から各配置は被約である。全四配置で三点直線はなく、
四点平面48枚のうち八枚を通す四点が線型に独立である。
それらと第五点の像から写像を全列挙した個数行列は$4 I_4$となる。
\path{computations/verify_klein_full_rank_configurations.py}で検算する。
節\ref{sec:klein-generic-isomorphisms}と同じ固有な行列比較により、これが標数零での上界となる。
指定された群が全配置を保つため等号である。
定理\ref{thm:finite-square-chart}により反対位置の平方空間は単元倍で対応するが、
同じVeronese三次元多様体に載る二つの表示なら、その二十点配置は射影同値になってしまう。
\end{proof}

% !TeX root = ../papers/nc-p3.tex
\section{判別式の他方の符号と楕円曲線上の中心}
\label{sec:klein-second-pairings}
前節の記号を用い、$D_0=ad-bc$と書く。薄い成分は全て非零で、
$z\ne\pm1$、交差ブロックと四つの組内ブロックは可逆とする。

\begin{theorem}\label{thm:klein-other-pairings}
$\delta\ne0$または$pq\ne0$なら、全ての補完は楕円曲線族の根の組合せまたは
以下の四表示で尽くされる。後者は$x=y=D_0$であり、
\[
\begin{array}{c|c}
 \lambda & A_1 \\ \hline
 1&A_0\\
 -1&-A_0\\
 1/z&\left(\begin{smallmatrix}d&c/z\\zb&a\end{smallmatrix}\right)\\
 -1/z&-\left(\begin{smallmatrix}d&c/z\\zb&a\end{smallmatrix}\right)
\end{array}
\]
となる。四表示のAS性は全て、両対称窓の既存の有限判定で決められる。
ここでは$\delta\ne0$で二次式の最高次係数が消える場合も含め、
楕円曲線側の組合せとは二つの射影化した根のうち$s,-s$に対応するものをいう。
\end{theorem}
\begin{proof}
$\delta=0$, $pq\ne0$なら各々の根は一つしかない。以下$\delta\ne0$とする。
二つの二次式を各々二変数に斉次化すると、根はそれぞれ相異なる二点である。
有限な根での符号は$2p x+m$または$2q y+k$で決まる。たとえば$p=0$なら
有限根の符号は$m$、無限遠の根の符号は$-m$である。
この対応は判別式を可逆にした全体で正則である。実際、無限遠座標$u=1/x$では
符号を$-m-2nu$と書け、共通部分で一致する。

残る四式は、分母を払うと$(x,y)$で双一次になる。たとえば対角比較の二本は
\[
\begin{split}
 (b+c-f/\lambda-\lambda g)y
   -(e+h-b-c)(x-D_0)&=0,\\
 (z^2b+c-\lambda z^2 f-g/\lambda)y
   -\{z(e+h)-z^2 b-c\}(x-D_0)&=0 .
\end{split}
\]
非対角比較の二本は、第一項がそれぞれ
$(\lambda ze+h/\lambda-za-d)x$, $(e/\lambda+\lambda zd-a-zd)x$で、
そこから$(za+d-zf-g)(y-D_0)$, $(a+zd-zf-g)(y-D_0)$を引く。
各式の双斉次化は$s,-s$の組の上で恒等的に零である。これは前節の計算と、
判別式が単元である根の有限エタール被覆で、一般点での等式が保たれることによる。
従って他方の符号の有限対$(x,y)$も解なら、各式は$x$固定で異なる二つの
射影$y$に消え、$y$固定でも異なる二つの射影$x$に消える。

まず$x$を固定した対角二式の定数係数を取る。
$x\ne D_0$なら$2e_0=b+c$, $2ze_0=z^2b+c$から$b=c=0$である。
逆に$y\ne D_0$なら非対角二式の定数係数から
$za+d=a+zd=zf_0+g_0$、すなわち$a=d=0$である。
従って$b,c$のいずれかが非零なら$x=D_0$であり、
\[
 \begin{pmatrix}1&1\\z^2&1\end{pmatrix}\binom fg
 =\begin{pmatrix}1/\lambda&\lambda\\\lambda z^2&1/\lambda\end{pmatrix}\binom bc .
\]
これは$F_1(C)$の対角条件である。さらに上の対角二式が$y$について恒等的に
零なので、$(f,g)$と$(b,c)$を交換した式も成立する。
左の係数行列の逆と右の係数行列の積の固有値は
$(1/\lambda\pm\lambda z)/(1\pm z)$である。その一つの平方が1でなければならず、
$\lambda\in\{1,-1,1/z,-1/z\}$を得る。
$b=c=0$なら$(a,d)\ne(0,0)$で、$y=D_0$を使った非対角二式が同じ結論を与える。
今度の二行列は対角$z$、非対角1の行列と、対角$\lambda z$、非対角$1/\lambda$の行列である。

これら四倍率で表の$(e,f,g,h)$を入れると$x=y=D_0$で、元の四式を満たす。
この根で$2pD_0+m=2qD_0+k$となる。これは表示した一次式への代入でも、
両辺の平方が同じことと各倍率での非零な一例からも従う。
同じ符号のもう一つの組が有限であれば、その両座標はともに$D_0$と異なるため、
前の定数係数の結論が$A_0=0$を強いてしまう。これで全被覆を得る。

$\lambda=1,A_1=A_0$は両対称である。$\lambda=1/z$の表示でも
$r=(1,1,z,1)$, $t=(1,1/z,1,1)$として
\[
 r_i w_{ijkl}=r_{\bar k}w_{\bar k,j,\bar i,l},\qquad
 t_j w_{ijkl}=t_l w_{ilkj}.
\]
従って偶相手の文字を平行移動した後に対角倍率で両対称となる。
負の二倍率は四元群の二コサイクル
$\mu(i,j)=(-1)^{i_{\rm hi} j_{\rm lo}}$による捻りで上記二例に戻せる。
実際、$H=\operatorname{diag}(1,-1)$として
$(A_0,A_1,\lambda)\mapsto(HA_0H,-HA_1H,-\lambda)$で、
厚い色の$M_h(i,k)$には$(-1)^{j_{\rm lo}+k_{\rm lo}}$が掛かり薄い色は変わらない。
この符号と四文字の$\prod_{r<s}\mu(i_r,i_s)$の差は、総次数零の各成分で
$j_{\rm lo}+k_{\rm lo}+j_{\rm hi}j_{\rm lo}+l_{\rm hi}l_{\rm lo}$の符号であり、
一文字の基底変更で吸収できる。
コサイクル恒等式により捻った乗法は結合的で、任意の斉次テンソル積に
累積文字次数による同じ符号を付ければ左右のKoszul複体の微分が対応する。
双対の完全対合も次数対ごとの非零倍率しか変わらない。よって捻りとその逆は
標準Koszul左右AS性を保つ。\path{computations/verify_klein_exceptional_pairings.py}で
上の二交換と符号を係数ごとに検算する。
\end{proof}

\begin{proposition}\label{prop:klein-elliptic-center}
$(a,b,c,d,z)=(5,2,3,7,-2)$を固定した楕円曲線の関数体上で、
四次中心の文字次数別次元は$(3,2,0,0)$である。
$\lambda=1,s=1$を含む一つの開集合の全閉点が中心四変数上で
各始点自由階数256の標準AS窓となる。
\end{proposition}
\begin{proof}
前節の四次多項式を7569で割ったものを$J(\lambda)$と書き、
$\mathbb Q(\lambda)[r]/(r^2-J(\lambda))$で$s=r/\lambda$とする。
\path{computations/verify_klein_elliptic_center.py}は各係数を
$(P(\lambda)+r Q(\lambda))/D(\lambda)$の形で四つの実際の窓を順に掃き出す。
次数5までの次元は標準値、文字次数0の36列の中心方程式の核次元は3、
次数1の40列では2となる。$\lambda=1,s=1$は
定理\ref{thm:klein-quartic-center-slice}の認証した有理点で、残る二文字次数の核は零、
かつ四中心元の商は次数13で消える。定階数の核を近傍へ延ばす前と同じ論法を適用する。
\end{proof}

\begin{corollary}[楕円曲線族の一般の印なし同型]\label{cor:klein-elliptic-isom}
四配置の射影安定群が指定の四元群に等しく、規格化と単純な根の条件を満たす
楕円曲線族の開集合で、印なし添字保持同型は高々32通りの文字置換と係数比較で決まる。
\end{corollary}
\begin{proof}
辺がグラフで、各配置の点ごとの射影固定群は自明なので、四位置の群自己同型は同じである。
文字の置換は共通の$\phi\in\operatorname{Aut}(K)$に四つの平行移動量
$u,v,u',v'$を加えたものとなる。総次数と両方の厚い次数の組から
$u+v+u'+v'=0$, $\phi(1)=1$, $u+v,u'+v\in\{0,1\}$を得る。
従って$2\cdot4\cdot2\cdot2=32$通りである。各置換後に薄い係数を規格化し、
$A_0$の一つの非零成分で両方の$A_0,A_1$を割る。$z,\lambda$とこの二行列の一致が
必要十分である。実際、残る薄い次数の逆倍率の四乗根の選び方を変えても、
各組内での行・列の相対倍率は同一の符号に限られ、両行列は同じスカラー倍しか変化しない。
逆にこれらの値は全テンソルを再構成する。適切な基底成分が零なら別の非零成分で割ればよい。
\end{proof}

% !TeX root = ../papers/nc-p3.tex
\section{四元群を仮定しない十四次元族と二つの局所変形}
\label{sec:unrestricted-doublet}
$r$を節\ref{sec:klein-generic-isomorphisms}の反対階数10の有理点、
$t$を節\ref{sec:klein-elliptic-completion}の反対階数16の有理正則点とする。
$r$の正則性にも四次中心の証明書がある。ここでは両点で全ての文字次数の
変形を許す。自己同型は四位置の射影基底変換で測る。

\begin{theorem}\label{thm:unrestricted-doublet}
$r$を通る片側対称な窓には、添字保持同型を除いて十四次元の既約族がある。
幾何学的一般点は反対階数10、最小周期4の標準Koszul左右AS正則窓で、
その射影自己同型群は自明である。片側対称性まで外した軌道に垂直な接空間は
$r,t$ともに十五次元で、両点で全ての二次障害が消える。
$r$における片側対称族の外への一方向は、反対切断の六次元の核から余核への同型を誘導する。
\end{theorem}

\begin{proof}
$W=(w_{ij,kl})$、$W^\tau=(w_{kj,il})$とすると、二切断条件はこれらの階数が
共に6となる条件である。各切断で六次の可逆な小行列を選び、その逆を逆順の
行・列位置に埋めた行列を$U$と書く。左核・右核の基底行列を$L,R$とすれば、
Schur補行列の一次条件は$D(T)=L^{\mathsf t}TR=0$である。
二方向$S,T$の二次障害は$D(SUT+TUS)$を一次条件の像で割った類になる。
二切断を同時に扱うとき、二番目にはその切断自身の核と逆行列を使う。
これは可逆小行列の逆の二次までの展開から従う。標数は零である。

$r$では$W^\tau=W$なので、交換の対称・交代部分に分ければ、二切断の一次条件は
それぞれ単独の$D$になる。符号の固有空間をさらに四元群の文字で分け、
四つの位置の無限小基底変換を除いた結果を示す。各組は文字次数$0,1,2,3$の順である。
\[
\begin{array}{c|c|c}
 &\text{アフィン軌道方向}&\text{軌道を除いた接方向}\\\hline
r,\ \tau=+1 &(10,12,12,12)&(4,2,4,4)\\
r,\ \tau=-1 &(3,4,4,4)&(1,0,0,0)\\
t,\ \text{全方向}&(13,16,16,16)&(5,2,4,4)
\end{array}
\]
計算には$ i+j+k+l$が指定文字となるテンソル単位を用いる（加法は四元群内）。
$r$では偶二位置を交換した和または差を取り、対角項は一度だけ使う。
各$\mathfrak{gl}_4$の行列単位を一位置に作用させれば軌道方向を得る。
\path{computations/verify_unrestricted_doublet_deformations.py}は全てを有理数上で
掃き出し、軌道方向が実際に核に入ることと、この表の各階数を検算する。
対称部分の160変数から百個のSchur条件への微分は全射である。
よって固定した偶交換を持つ階数6軌跡は$r$で滑らかな60次元である。
反対像が対称二乗全体なので、二つの片側対称表示の間の偶二位置の基底変換は
射影的に一致する。実際、任意の平方の像が再び対称になるには二写像が各直線上で比例する。
従ってこの軌跡上の一般の軌道次元は46（アフィン倍率込み）で、
有理商の次元は$60-46=14$である。軌道の有限な射影安定群は以下からも従う。

三次、四次の縮約が作る有限Frobenius双対に対し、節\ref{sec:reversible-component}の
内部次数別の特殊化をこの既約族と正則点$r$に適用すると、幾何学的一般点で全次数の
正則性を得る。反対階数10は片側対称部分で開く。異なる二奇配置が互いに射影非同値で
あることは$r$での有限被約な二十点配置の比較により一般点に保たれるため、最小周期は4である。

同じ有限エタールな点の番号付けを$r$の近くで施す。一般点で窓を保つある置換が存在すれば、
各位置でその二十点を運ぶ行列の存在条件は射影的に閉で、$r$に特殊化する。
三点直線・五点平面がないため極限行列は可逆である
（節\ref{sec:klein-generic-isomorphisms}）。特殊点で各位置に五点の射影枠を選ぶと
候補行列は枠から一意かつ正則に復元され、窓を保つ等式は既約な近傍に延びる。
特殊点における窓の射影安定群は対角四元群である。実際、各配置の安定群は四元群で、
辺が点集合間のグラフであるため一位置の作用が全ての位置の作用を決める。
特殊値を$g$として、変動する行列によるテンソルの固定等式を微分すると、
任意の接ベクトル$v$について$gv-v$は軌道接空間に入る。番号付けはエタールなので
全接方向でこれを要求できる。しかし上表の対称行には全ての非自明文字が現れ、
$g$が商接空間を点ごとに固定すれば$g=1$である。有限エタールな番号付けの
異なる置換は特殊点でも異なるので、一般点の置換も恒等に限られる。

$r$の二次条件では対称同士と交代同士の積が対称側に入るので、上の全射で解ける。
交代方向の代表一つと十四の対称代表の混合積は、各文字の交代微分の像に入る。
$t$では十五代表の重複を許す全120組を二切断で同時に調べる。
どちらも上記コードで障害列を付け加えた階数増分が全て零となる。
基底変換自体は持ち上がるので軌道方向を除く計算で十分である。
最後に$r$で反対切断の左右核で唯一の交代代表を挟むと階数6を得る。
従ってその接線を持つ弧があれば反対行列式の最低次はちょうど6である。
\end{proof}

\paragraph{階数条件を保つ無限小変形の平坦性。}
二つの切断の階数が局所自由な意味で6である局所Artin環上の持ち上げで、
特殊ファイバーが標準AS窓なら、その四窓代数は各内部次数で自由となる。
実際、巡回したテンソル、末尾一文字と二文字を縮約した像、矢空間、基底環を
順に使い、各始点に階数$1,4,6,4,1$の係数を持つ線型複体を作る。
これらの像は単元小行列で直和因子となり、連続二縮約の積は二次関係に入るので
微分の合成は零である。内部次数$n$で、正のホモロジー位置の項に現れる
代数の次数は$n$未満である。帰納的にそれらは自由で、特殊ファイバーは
Koszul分解なので正の位置で完全である。左端から単元小行列を順に分離すると
基底環上でも正の位置で完全で、矢から掛け合わせた次数$n$商は自由になる。
より具体的に末尾の商は関係を挿入した
$R_i\otimes A_{i+2,i+n}\to V_i\otimes A_{i+1,i+n}$の余核であり、
二次表示の定義通り$A_{i,i+n}$である（$n=1$では関係項は零）。
左側も逆順で同じ証明となり、traceの相補次数対合は特殊化で完全なので完全である。
この議論は二次までの持ち上げにも適用される。十五の一次方向が全て同時に
任意次数まで持ち上がるかは、二次障害の消滅だけでは決まらない。

% !TeX root = ../papers/nc-p3.tex
\section{二十点が与える恒等式と十五次元の成分}
\label{sec:twenty-conormal-component}
前節の二つの有理点を$r,t$とする。全テンソルのアフィン空間の中で、
二つの隣接切断が共に階数6となる局所閉スキームを$Z$とする。
以下の恒等式は$Z$を定義する方程式の間に、周囲の環上で成立する。

\begin{proposition}\label{prop:twenty-conormal}
$W$とその第二切断$N=W^\tau$が階数6である点を考える。
第一切断の左六関係を消す$([a],[b])\in\mathbb P^3\times\mathbb P^3$が
横断的な二十点であるとする。各点で、$J_b c=c\otimes b$,
$J_a d=a\otimes d$と書く。$J_b^{\mathsf t}N$と$N J_a$が階数3で、
各々の左核・右核を$c,d$が生成するとする。二十個の
$(c\otimes b)\otimes(a\otimes d)$が線型独立なら、第一切断の
階数6軌跡上で第二切断の百個のSchur方程式は局所エタールに八十個で生成される。
さらに二方程式系を合わせた微分の階数が180なら、$Z$はその点で滑らかな76次元である。
\end{proposition}
\begin{proof}
第一切断の階数6軌跡は、六次小行列を可逆にすれば滑らかな156次元である。
六つの双一次式の横断的零点は射影的な完全交叉で、近傍で有限エタールに動く。
番号付け被覆上で二十対を切断として扱い、座標の局所代表を選ぶ。
第一切断の左核条件から、第二切断の階数に関係なく
\[
 J_b^{\mathsf t}N J_a=0
\]
が四行四列の行列として成立する。第二切断の可逆な六次小行列を左上に移して
\[
 N=P Q+E,\quad
 P=\binom{N_{11}}{N_{21}},\quad
 Q=(I_6\ \ N_{11}^{-1}N_{12}),\quad
 E=\begin{pmatrix}0&0\\0&B\end{pmatrix},\qquad
 B=N_{22}-N_{21}N_{11}^{-1}N_{12}
\]
と書く。$X=QJ_a$, $Y=J_b^{\mathsf t}P$は基点で階数3、$YX=0$なので
$\operatorname{im}X=\ker Y$となる。この三平面の補空間$H$を固定する。
近くで$k^4\xrightarrow{X} k^6/H$と$H\xrightarrow{Y}k^4$は階数3のままである。
前者の核を$d'$、後者の左核を$c'$で局所的に生成する。剰余の階数を仮定せずに
$Xd'\in H$かつ$(c')^{\mathsf t}YH=0$となる。従って
\[
 0=(c')^{\mathsf t}J_b^{\mathsf t}N J_a d'
   =(J_b c')^{\mathsf t} E(J_a d')
\]
が周囲の第一切断軌跡上の正確な恒等式になる。基点で$c',d'$は$c,d$である。
基点の第二切断の左右核を非主行・非主列の十座標へ射影する写像は、どちらも
六次主行列の可逆性から同型である。従って$B$の百係数の前に並ぶ二十行は
基点で独立である。対応する二十次小行列を可逆にすれば、二十個の$B$の成分が
残る八十個の線型結合となる。両切断を合わせた微分の階数が180なら、
この八十方程式の第一軌跡上の微分は全射である。Jacobi判定とエタール降下で結論を得る。
\end{proof}

\begin{theorem}\label{thm:full-fifteen-component}
$r,t$を含む$Z$の既約成分は両点で一意で、同一である。その添字保持同型を除いた
次元は15である。幾何学的一般点は標準Koszul左右AS正則、最小周期4、反対階数16で、
四位置の射影自己同型群は自明である。$r$の近くで片側対称窓の基底変換による像は
滑らかな因子をなし、その上の反対階数は10、外側では16である。
反対行列式はこの因子でちょうど六重に消える。
\end{theorem}
\begin{proof}
$r,t$の第一辺は六つの双一次式の被約二十点完全交叉である。
素数12619と2657で四配置を実際にたどり、二つの評価階数が各点で3、
二十行二百五十六列の上記行列が階数20であることを
\path{computations/verify_twenty_conormal_identities.py}で検算する。
第一辺の有限エタールな番号付けに沿って階数下界は標数零に延びる。
基点では$YX=0$なので両方の階数が3以上ならともに3であり、
核直線も特殊化と可換である。前節によれば各点の接空間は$61+15=76$次元で、
微分の階数は180である。命題\ref{prop:twenty-conormal}の仮定が全て成立する。

節\ref{sec:klein-elliptic-completion}の既約族の閉包は$\lambda=1$に
第二補完の一般点を含む。四基底の変換と全体倍率を許した閉包は既約であり、
第二補完の有理点$r$と楕円曲線族の有理点$t$をともに含む。
両点の局所環が正則なので一意な成分が一致する。一般の軌道は、四位置の
点配置が有限で射影枠を含むためアフィン次元61である。よって有理商は
$76-61=15$次元となる。ここでの滑らかさはテンソル軌跡の主張であり、
有限安定群を持つ点での粗商の滑らかさを前提としない。
正則点$t$から有限Frobenius条件と内部次数別の完全性を一般点に移す。
$t$の四配置の非同値性が一般点で保たれるので最小周期は4、反対階数も開条件で16となる。
$t$の射影安定群は四元群である。各置換を枠から復元して微分する前節の議論を、
今度は$Z$の滑らかな近傍で用いる。接空間の軌道商には全非自明文字が現れるため、
一般点での射影安定群は自明である。

最後に偶二位置を同定して対称となる条件を考える。基点$r$で反対切断の像は
対称二乗全体である。階数が10の近傍をこの対称な部分で取ると、同定は射影的に一意である。
これを動かす自由度は$\mathrm{PGL}_4$の15次元である。より具体的には、
対称二乗平面の基底変換軌道はGrassmann多様体の局所閉な斉次空間で、
その安定群は対角$\mathrm{PGL}_4$である。階数10の部分では像平面が正則に定まるので
同定も正則に回復できる。従って局所的に、この片側対称軌跡は滑らかな60次元の
固定同定軌跡に同定の15変数を加えた75次元の滑らかな部分多様体となる。
実際、反対階数が10以上の開集合で階数10は閉条件で、像平面の斉次軌道も
その各点の近くで閉にできる。$Z$は滑らかな76次元なので、$r$の近くでこの因子の
局所方程式を$u$とできる。反対切断の十次小行列でSchur補行列を取ると、
残る六行六列の各成分は因子上で零なので$u$で割れる。前節の横断方向の
核から余核への微分は可逆である。よって$u$で割った六行六列行列は
$r$の近傍で可逆となり、階数の二分と行列式の消滅次数が従う。
\end{proof}

% !TeX root = ../papers/nc-p3.tex
\section{二つの行列式二次式が消えるファイバー}
\label{sec:klein-zero-quadratics}
節\ref{sec:klein-elliptic-completion}の薄い成分が非零で、$z\ne\pm1$、
交差ブロックと組内のSchur主ブロックが可逆なチャートを用いる。
特に$z\lambda(1-\lambda^2z)D_0\ne0$である。

\begin{theorem}\label{thm:klein-all-pairings}
このチャートで二つの隣接階数が6となるテンソルは、五次元楕円曲線族の閉包と
定理\ref{thm:klein-other-pairings}の四つの表示で全て覆われる。
$P,Q$がともに零多項式のとき、共通倍率を$a=1$で固定すると八種類の基底条件があり、
それぞれのファイバーは単一の双一次方程式で決まる。
\end{theorem}
\begin{proof}
係数環は有理数上の多項式環をチャートの分母で局所化したものである。
$P(x)$と$Q(y)$はそれぞれ原始である。実際、定数項は順に
\[
 n=(a-zd)(za-d)/(z+1)^2,\qquad o=(b-c)(z^2b-c)/(z+1)^2
\]
で、各一次因子を零にしても対応する最高次係数は恒等的に零でない。
例えば$(a,b,c,d,z,\lambda)=(1,3,4,5,2,3)$から順に$d=za,a/z$、
$c=b,z^2b$と一箇所を置換すればよい。共通判別式は分数体で平方でない。
従ってGaussの補題により、$P(x)=0$と$Q(y)=0$の各全空間は整である。
両者の一般点から楕円曲線族への有理逆写像は
$2px+m=-(2qy+k)$で与えられる。

$\delta=0$で片方だけが零多項式の場合を考える。例えば$P\equiv0$なら、
解を持つ非零の$Q$は$q\ne0$の二重根を持つ。このとき
$y=-(2px+m+k)/(2q)$は$P(x)=0$上で正則であるため、任意の$x$に
付随する解は一般の族の特殊化である。逆側でも同様である。

$P\equiv Q\equiv0$なら上の定数項から$d\in\{za,a/z\}$,
$c\in\{b,z^2b\}$となる。$a=0$または$b=0$では$m=1$または$k=-1$となるので、
$a,b$は非零である。$a=1$とする。$f_0 g_0 e_0\ne0$であり、
$p=q=0$と$m=k=0$から、$f_1,g_1$の一方だけと$e_1,h_1$の一方だけが零である。
表の各対はこの零となる係数で、隣は$b$の値である。
\[
\begin{array}{cc|cc|cc}
 d&c& \multicolumn{2}{c|}{\text{第一の場合}}&\multicolumn{2}{c}{\text{第二の場合}}\\\hline
 z& b &(f_1,e_1)&1/\lambda &(g_1,h_1)&\lambda z\\
 z& z^2 b &(f_1,h_1)&\lambda &(g_1,e_1)&1/(\lambda z)\\
 1/z& b &(f_1,h_1)&1/(\lambda z)&(g_1,e_1)&\lambda\\
 1/z& z^2 b &(f_1,e_1)&\lambda/z &(g_1,h_1)&1/(\lambda z^2)
\end{array}
\]
検算の入力を明記する。$f_1=0$なら$T_{00}=T_{11}=1/f_0$、
$g_1=0$なら$(T_{00},T_{11})=(1,z^2)/g_0$である。
$e_1=0$なら$(T_{01},T_{10})=(1,z)/e_0$、
$h_1=0$ならこれが$(z,1)/e_0$となる。これらと
$T=Y_\lambda A_0^{-1}X_\lambda$を等置し、上記の非零因子で飽和させる。
十六の選択のうち八つでは単位イデアルを得る。他の八つでは表示した$b$を
強制する式がイデアルに入り、代入で全式が消える。
\path{computations/verify_klein_zero_quadratics.py}が有理係数のGröbner基底で検算する。

各行の代入後、$e,h,f,g$を前の一次式で表し、$F_\lambda(A_1)$に
要求される四式の分母を払う。零でない式は全て、係数環の単元倍を除いて同じ
$H(x,y)=\alpha xy+\beta x+\gamma y+\eta$であり、$\alpha$は単元である。
この主張も同コードが各式の係数を照合し、倍率の因子が全て可逆であることを検査する。
例えば表の左上の場合、単元倍を選ぶと
\[
 H=\lambda^2(\lambda^2 z-1)xy
   -\lambda^2(z-1)^2(x+y)+(z-1)^2(\lambda^2z-1)
\]
となる。ここでも$x,y,\det A_1$の非零条件を保つ。

一般に残る四式は$y$について線型である。その一つの$y$係数が基点で非零なら、
整な$P(x)=0$上の近傍でその式から$y$を解ける。一般点でその$y$は楕円曲線族の
ものであり、残る全ての等式は整性から近傍に延びる。よって当該点はその族の閉包に入る。
$x$係数が非零なら$Q(y)=0$から同様に従う。各固定ファイバーの$H$は$xy$係数が
非零なので既約な双一次曲線か、相異なる縦横二直線である。各成分の稠密開部では
少なくとも一方向に解け、交点も閉包に入る。この議論と
定理\ref{thm:klein-other-pairings}でチャート全体の被覆を得る。
ここで決めたのは階数軌跡であり、閉包上の各閉点の全次数正則性はなお別の判定を要する。
\end{proof}

% !TeX root = ../papers/nc-p3.tex
\section{一つの二十点配置への射影と平方根の衝突}
\label{sec:single-configuration-projection}
定理\ref{thm:full-fifteen-component}の十五次元の有理商を$M$とする。
第一位置の二十点の射影同値類を取る有理写像を$\pi$とする。

\begin{theorem}\label{thm:single-configuration-projection}
$\pi$は像へ生成的有限で、その関数体次数は偶数である。
第一関係六平面を固定座標で記録した像の次元は45、第一位置の点配置のみを
固定座標で記録した像の次元は30である。$t$および一般点で、四位置の各座標
四空間は自分の平方を固定したGrassmannスキームの孤立した被約点である。
$r$の偶位置ではこの局所ファイバーは長さ2であり、奇位置では被約点である。
\end{theorem}
\begin{proof}
第一切断の左核の基底を$L$とすると、第一関係空間を一定にする接方向は
$L^{\mathsf t}T=0$である。これと第二切断のSchur微分の共通核は$r,t$で
文字次数順に$(7,8,8,8)$次元となる。第三・第四位置の基底変更の軌道は
$\mathrm{GL}_4^2$の逆スカラー一組を除いて31次元で、常にこの共通核に入る。
従って一般にも31次元であり、第一辺の像の次元は$76-31=45$である。

次に任意の第一六平面の基底$r_1,\ldots,r_6$を動かす。点対$(a,b)$で
$E_a=(a^{\mathsf t}r_s)_s$は階数3で、第一点$a$を固定して第二点を動かせる
一次条件は$(a^{\mathsf t}\dot r_s b)_s\in\operatorname{im}E_a$である。
番号付き二十点がエタールに動くので、全二十対について課せば十分でもある。
各対につき三式を並べた60行96列の階数は両例で45である。下界は素数12619,2657で
得られ、上界は六関係の基底変更36と第二位置の射影変換15が核に入ることによる。
後者は先の点グラフと射影枠から独立に作用する。従ってGrassmann多様体の
接空間で第一配置を固定する方向は全て第二位置の基底変更であり、先の45次元像
にもこれらの軌道が含まれる。その第一配置への像は$45-15=30$次元、
第一位置の射影変換を除けば$30-15=15$で、$\pi$の生成的有限性を得る。
全階数の検査は\path{computations/verify_single_configuration_projection.py}にある。

奇二位置の交換は四窓を逆に読む操作で、左右AS条件を保存する。
第一位置の配置は逆向きのグラフの同じ端点となる。この操作は$M$を保つ。
実際、四元群の規格化前には$(M_h,N_h)$を$(N_h^{\mathsf t},M_h^{\mathsf t})$に
移すので、非零な薄いブロック、可逆な交差・組内ブロックを保つ。
薄い積からの再規格化は対角変換である。$z=\pm1$が一般に生じないことも、
$M'_2=N_2^{\mathsf t},\ M'_3=N_3^{\mathsf t}$の薄い係数の組内積比から従う。
より具体的には$\lambda=1$の片側対称な断面で、偶交換と奇交換の合成は
対置行列を転置し、薄い外積の二ベクトルを交換する。このとき薄い係数
$u_i^h v_k^h$の規格化前の$v_i^h=u_i^{h+1}$により両色が交換され、
$U,V$の各二列の順序が反転するので積比$z$は$1/z$に移る。
従って一般の像は定理\ref{thm:klein-all-pairings}の非退化チャートにある。
五次元の像は四つの例外表示（各々印なし次元は高々4）に入り得ず、
楕円曲線族の閉包に戻る。この族が$t$で滑らかな唯一の全成分に入るため
逆転はその成分も保つ。一般点で第二と第四の配置は射影非同値なので、
この対合は関数体に非自明に作用する。$\pi$は対合に不変であり、その次数は偶数となる。

最後に$D=k^{20}$で一位置の座標行列の列空間を$A$とする。
$B^2\subset A^2$となる$B\in\Gr(4,D)$を考え、$A$で独立な四行を固定する。
残り十六行の四係数がGrassmann座標になる。変分$S$の一次式は、全ての
$i\le j$について$A_i S_j+A_j S_i\bmod A^2$である。$t$の四位置と$r$の
二つの奇位置で階数64を得る。$r$の偶位置では還元後の階数は63で、
核の一方向$S$の二次項$(S_i S_j\bmod A^2)_{ij}$を付け加えると64となる。

$r$の偶位置の標数零での接次元が少なくとも1である理由も必要である。
二十点を番号付けした成分上、第一・第三位置の平方は単元倍を除いて等しい。
これは各辺の完全交叉のtraceを比較する定理\ref{thm:finite-square-chart}の証明で、
有限エタールな族上でもtraceとその比は局所切断となる。
必要な行ごとの平方根をエタール局所で選ぶ。$r$では第一・第三の点列が一致し、
行の自明化を合わせると平方を保つ単元は命題\ref{prop:finite-square-ratio}によりスカラーである。従って符号を合わせれば二つの平方根部分空間は一致する。
一般点では両配置が射影非同値なので二切断は相異なる。相対Grassmannスキームが
この点で不分岐なら対角がその近傍で開となり、二切断は一般点でも一致してしまう。
従って標数零でも接次元は1である。単元な63次小行列で核の生成元を
素数12619の剰余へ整合的に取れるため、上の二次障害も標数零で非零となる。
六十三変数を形式的に消去すると残りは一変数で、局所イデアルはその二乗で生成される。
これで重なった平方根の局所長さを得る。
\end{proof}

\begin{corollary}
十五次元成分の幾何学的一般点では、第一配置だけの射影安定群も自明である。
\end{corollary}
\begin{proof}
$t$で、上の二段の接計算は、$Z$の接空間を基底変更の方向で割った空間から
第一配置の接空間を射影変換の方向で割った空間への微分が単射であることを示す。
実際、第一配置を射影的に固定する接方向から第一位置の基底変換を引き、
次に第二位置の基底変換を引けば第一辺が一定になり、最後は第三・第四位置の軌道に入る。
一般の配置の安定群の元が与える置換を、有限エタールな番号付け上で$t$に特殊化する。
枠と特異行列の除外により、それは$t$の第一配置の四元群の元$g$となり、
枠からの行列は近傍で正則に動く。この$g$は$t$の四窓にも作用する。
従って固定等式を微分すると、全ての窓の商接方向$v$について
$\mathrm d\pi(gv-v)=0$が配置の商接空間で成立する（ここでは群軌道を引いた
接空間同士の写像を用いる）。単射性と全非自明文字が現れることから$g=1$である。
番号付け上で異なる置換は特殊点でも異なるため結論を得る。
\end{proof}

\paragraph{一つの配置からの有限候補。}
一般ファイバーの復元には、第一座標の四空間$A$から$Q=A^2\subset D$を作り、
$Q$と標準内積による$Q^\perp$に対して平方が入る四空間を求めればよい。
第一辺の完全交叉のtraceは各点で非零なので、第二座標の行単元を調整すれば
その平方は$Q^\perp$になる。第三座標は同様に$Q$の平方根である。
一般ファイバーの実際の解は上の孤立被約点に入るので、この二つの
Grassmannスキームの零次元成分を列挙すれば必ず候補に含まれる。
全階数の有理点では、四辺とも、第一空間自身を第三空間として再使用する候補は三次交叉で落ちる。実際、各六関係を列とする$R$について$R\otimes I_4$と$I_4\otimes R^{\mathrm{op}}$を並べると有理数上で階数48となり、交叉次元は必要な4に達しない。最大階数は一般点にも保たれる。残る各候補の隣接積の核から二窓を作り、系\ref{cor:two-window-extension}で
延長と同型を判定する。次節では生成的有限写像の正確な次数を決める。特殊配置での正則性判定はなお必要である。

従って復元次数$e$には、粗いが一様な上界
\[
 2\le e\le 2^{45}(2^{45}-1),\qquad 2\mid e
\]
がある。実際、一つの平方根スキームの孤立点は有限個なので、全てを一つの一般の
Grassmannアフィン座標に入れられる。六十四変数において平方条件は高々二次であり、
孤立解に対するBézout不等式からその数は$2^{64}$以下となる（正次元成分がある場合も
孤立点の局所重複度の和に同じ上界が成立する）。実際の四空間では平方を安定化する
対角単元はスカラーのみであるから、座標の符号倍は全体符号を除いた$2^{19}$個の
相異なる孤立解を与える。各符号軌道は同じ積の核を与える。よって実際の候補は
各側高々$2^{45}$通りで、第三位置では既知の第一空間の符号軌道を除ける。
二窓延長の一意性により積が上界となる。

% !TeX root = ../papers/nc-p3.tex
\section{三重点の四次鉛筆と復元次数二}
\label{sec:configuration-quartic-pencil}
$D=k^{20}$に成分ごとの積を入れ、$Q\subset D$を十次元とする。
座標評価を$p_i\in Q^*$とし、$H(Q)$を各$[p_i]$で三重以上に消える
$\mathbb P(Q^*)$の四次式の線型空間とする。以下、標数零を用い、
有限体での検算には標数$> 4$を用いる。

\begin{proposition}\label{prop:quartic-root-pencil}
$L\in\Gr(4,D)$が$L^2\subset Q$を満たすと、
$H(Q)$に自然な非零の射影元$[F_L]$がある。各$i$でその対称四重線型形式の切断
\[
 Q^*\longrightarrow \operatorname{Sym}^2 Q,\qquad
 v\longmapsto F_L(p_i,v,-,-)
\]
は階数6以下である。$L^2=Q$なら$F_L$の三重点軌跡は二次Veronese三次元多様体で、
$F_L$が一致する二つの全階数平方根は、対応する二十点で独立な行単元倍を除いて同じである
（各$p_i\ne0$とする）。
\end{proposition}
\begin{proof}
$\mu:\operatorname{Sym}^2L\to Q$を積として
$F_L(v)=\det(\mu^*v)$と取る。$\mu^*p_i$は$L$の第$i$行の平方である。
独立な四行の平方の和は非退化なので、$\mu$の全射性を仮定しなくても
この四次式は恒等的に零にならない。階数1の対称四行行列での最低次項は
核三平面上の対称行列の行列式である。その三次式が使う変数は六つであり、
零行に対応する階数0では三次項自身が零となる。制限写像によって変数の個数は増えない。
全階数なら三重以上の軌跡はちょうど階数1の射影行列の全体である。
この三次元多様体を保つ射影自己同型は元の$\mathbb P^3$の自己同型から誘導される。
従って二つの表示の各行は共通の基底変更の後に比例する。さらに双方の平方が$Q$で、
一方の三乗が$D$に同型なら、命題\ref{prop:finite-square-ratio}より
行単元の平方はスカラーで、違いは全体倍率と各行の符号に限られる。
\end{proof}

\begin{theorem}\label{thm:exact-two-configurations}
十五次元成分の幾何学的一般点と反対階数16の有理例$t$で、
第一配置からの添字保持同型類の復元次数はちょうど2である。
四辺が有限な標準AS延長は元の窓とその逆向きの窓で尽きる。
第一辺まで指定すると延長は一意である。反対階数10の有理例$r$の
第一配置でも同じ二つの向きで尽きる。
\end{theorem}
\begin{proof}
第一位置の行列の列空間を$A$とし、$Q=A^2$とする。
順に他の三位置を$B,C,E$と書く。辺上の点の番号を共有し、
定理\ref{thm:finite-square-chart}の完全交叉のtraceを使って
$C^2=Q$, $B^2=E^2=Q^\perp$とできる。後者では$D$の標準内積を使い、
行ごとの直線束の自明化を選ぶ。ここで扱う基点と一般点では四移動は恒等である。実際、二十点の番号付け上の戻り置換は局所一定で、両有理例で恒等と検算されている。従って$Q^\perp$の射影評価点は$B$の各行の
二次Veronese像に射影同値である。$t$の反対位置の点配置は互いに射影非同値
なので、両側で少なくとも二つの異なる行列式四次式がある。

計算を記す。基底の十変数の四次単項式は715個あり、各評価点で
$F(v+s p_i)$の$s^2$係数（$v$の二次式）を零と置く。
対称化により、これは三重以上の条件と同値である。二十点を合わせた
1100行715列の行列の核を取り、その基底$F_0,F_1$の各々から
最初の点の上記切断を作る。切断は55行10列であり、$F_0+uF_1$に対して
その七次小行列式の線型結合を二本取る。具体的には左右から七行と七列に圧縮した
行列式であり、各々次数7以下なので八値で正確に計算できる。
\path{computations/verify_configuration_quartic_pencil.py}の基底と圧縮を使った結果は次の通りである。
\[
\begin{array}{c|c|c|c|c}
 \text{基点}&\text{素数}&\text{平方の位置}&\dim H&\text{二本の最大公約式}\\\hline
 t&2657&A&2&u(u+277)\\
 t&2657&B&2&u(u+22)\\
 r&12619&A&2&u^2\\
 r&12619&B&2&u(u+2481)
\end{array}
\]
どれも無限遠で少なくとも一方が非零である。この計算が与えるのは、
可能な行列式を含む次数2の上界である。一般の鉛筆の全ての元が行列式になることは仮定しない。

$t$の近くで番号をエタール局所的に取り、二つの既知の行列式により
核次元は少なくとも2となる。上の小行列で同時に高々2なので鉛筆は局所自由である。
有限体から標数零の$t$へも同じ議論が使える。射影鉛筆上の二本の七次式の
共通零点は特殊ファイバーで長さ2なので、近傍で有限かつ長さは高々2となる。
これは射影性、有限射の直像に対する半連続性による。よって$t$と一般点で
可能な行列式は両側とも既知の二本だけである。これらは十変数全てを必要とする。
実際、行列式の三重軌跡は滑らかなVeronese像であり、非零な方向を変数から
除ければその射影点で四重になってしまう。従って$\mu$が全射でない平方根から
同じ式を得ることもない。命題\ref{prop:quartic-root-pencil}により、退化解も含め
Grassmannスキームの全幾何点が各側ちょうど二つの符号軌道で尽くされる。
既知の解は孤立被約で、一軌道は$2^{19}$点である。

こうして第二空間の候補は$B,E$、第三空間の候補は$A,C$に限られる。
$B$の後に$A$を取る候補は前節の三次交叉で排除される。$E$を先にした
逆向きの窓にも同じ検査を行うと$t$で交叉は零であり、一般点でも零である。
各組は連続二辺の積の核を決め、そのAS延長は系\ref{cor:two-window-extension}
により一意である。二つの向きは第二配置が異なるため同型でない。
特に十五次元成分の関数体は第一配置の像の関数体上二次であり、非自明な
自己同型は逆転である。ここで全四辺有限の他成分から同じ配置に来るAS窓も既に数えた。

$r$では奇側$B,E$が異なるので同じ証明が適用できる。偶側は標数12619で
上の共通零点が一つしかなく、既知の$A$の行列式がその点を占める。
$A$の三次評価は階数20であるから、上の議論により全幾何点は$A$の符号倍である。
前節の検算はこの剰余体でも各点の接次元1、二次障害非零を与えるので、
その局所長さは2で、全長さは$2^{20}$である。平方根Grassmannスキームは
固有なので有理点$r$の標数零ファイバーも有限で長さは高々$2^{20}$である。
前節で証明した$A$の長さ2の符号軌道が既にこれを尽くす。従って第三空間は
$A$だけとなる。二つの第二空間は実際に正則な両方向を与えるため結論を得る。
\end{proof}

% !TeX root = ../papers/nc-p3.tex
\section{二つのVeronese像の交わりと向きの分岐}
\label{sec:veronese-pencil-quotient}
$M$を十五次元成分の添字保持同型による有理商とする。$Q=A^2$の評価二十点を
$Y\subset\mathbb P(Q^*)$とし、それを含む二つのVeronese三次元多様体を
$V_A,V_C$と書く。以下の一般点に関する主張は標数零である。

\begin{theorem}\label{thm:veronese-pencil-quotient}
一般点では$V_A\cap V_C=Y$がスキームとして成立する。
奇位置の側でも同じであり、二つの二十点集合は互いの射影Gale変換である。
$M$から第一の四次鉛筆の射影同値類への関数体拡大は四元群をGalois群とする四次拡大である。
二鉛筆のうちどちらが第一かも忘れると、拡大の次数は8でGalois群は正方形の二面体群となる。
\end{theorem}
\begin{proof}
三次評価が同型なので$Y$を$V_A\simeq\mathbb P^3$で読むと、二十点の
イデアルの四次元数は$35-20=15$である。第一辺の六行四列の線型評価行列が期待余次元を持つため、
Eagon--Northcott分解によりイデアルはその四次最大小行列式で生成される。
$V_C$の二次小行列式を$V_A$へ制限すると、これらは二十点を消す四次式になる。
$t$の素数2657のファイバーでは、二十行十列の二次評価行列$X_A,X_C$について
$X_A T=\operatorname{diag}(s)X_C$を解く。$s$を決める百本の式は
$X_A$の左核を両辺に掛けた式であり、核は単元ベクトルの張る直線である。
得られた$T$を対称四行行列の全二次小行列式に入れ、四次の35単項式で展開すると
行の階数は15である。奇位置の側も15となる。
\path{computations/verify_configuration_intersections.py}が両方を検算する。
従って標数零の$t$および一般点でも全四次生成元を得る。これが交叉のスキーム構造まで決める。
一方$Q^\perp$の評価点は、$Q$の評価行列の関係を基底とした行列の二十個の対応行であり、
Gale変換の定義そのものである。各点の代表を変えても相手側の行倍率が変わるだけである。

一般鉛筆の中の行列式の集合は内在的である。既知の二本の共通三重点集合から上の
$Y$を回復できる。また鉛筆全体の共通三重点をまず取り、その任意の一点で
前節の偏極を課すこともできる。実際、二本の行列式では二階微分の成分が
対称四行行列の二次小行列式を張る。これは行列式の四重偏極で対角二方向または
非対角方向を選べば相補な二行二列式が得られることから従う（対称添字の
三通りの組合せの間には通常の直線関係がある）。従って集合としてはもちろん、
共通三重点のスキームも上記の交叉に一致する。

四位置の巡回移動を$T$、第一位置を固定した逆転を$J$とする。
逆転が$M$を保つことは前に示した。巡回移動は二つの隣接切断を行・列置換と
転置を除いて交換し、四元群の二つの薄い色を保つ。移動後に非退化チャートに
入ることは$t$で検算できる。実際、二つの切断の各厚い四行ブロックで四つの
組内・交差小行列式は全て非零で、薄い二行列の各成分も非零である。
薄い階数一行列の第0列を$u^2,u^3$とすると、規格化の$z$は
$u^2_2u^2_3u^3_0u^3_1/(u^3_2u^3_3u^2_0u^2_1)$であり、
移動後の第一切断で$-1/8$となる。列の選択や対角倍率はこの比に影響しない。
よって五次元の楕円曲線族の移動後の像は、非退化チャートの最大四次元の
例外表示には含まれず、同じ五次元族の閉包を支配する。$t$における全成分の一意性から
$T$は$M$を保つ。$T^4=J^2=1$, $JTJ=T^{-1}$であり、
四配置が一般に全て非同値であるため八作用は互いに異なる。

一つの鉛筆から二つの三重Veronese軌跡の一方を選ぶと、その上で交叉点を読むことにより
第一配置を射影同値を除いて復元できる。さらに向きを選べば前節により窓は決まる。
従って高々四通りであり、$1,T^2,J,T^2J$は全てその鉛筆を保つので丁度四通りである。
鉛筆は$Y$を決め、相手側はそのGale変換から回復できる。両鉛筆が射影同値なら、
対応するVeronese上で偶奇のいずれかの配置が射影同値になるので一般には相異なる。
従って順序なしの鉛筆対からのファイバーは八作用の一軌道で尽くされる。
標数零であるため、いずれの関数体拡大も表示した群の固定体拡大となる。
\end{proof}

\begin{corollary}\label{cor:orientation-ramification}
第一配置への二次拡大は、奇二位置が対称となる因子の一般点で分岐指数2を持つ。
その因子の一般配置における四辺有限のAS延長は一つである。
偶二位置が対称となる因子の一般点では逆転による二窓は異なる。
\end{corollary}
\begin{proof}
$r$を一位置移動して$r'=Tr$とする。前節の$r$の平方根スキームを
奇偶を交換して読むと、$r'$では第一側に二つの被約符号軌道があり、
第二側には一つの長さ2の軌道がある。従って第二空間は一意である。
第三位置で第一空間を再利用する候補は、前節の$r$における奇始点の三次交叉が
零である計算で排除される。

この計数は奇対称因子の一般点にも成立する。固有な平方根スキームのファイバーの
長さは$r'$の近傍で各側高々$2^{20}$である。第一側には各$2^{19}$点の相異なる
二つの軌道がある。第二側では二つの一般切断が因子上で一致するので、その根は
因子の一般点で不分岐ではない。これは二つの切断を番号と行平方根を選んで
エタール局所的に取れば、不分岐の場合の開対角に反するためである。従って
一つの符号軌道だけで既に全長さを使い切る。三次交叉零の条件は開であり、
実在する幾何学的一般AS窓だけが残る。

この因子は十四次元の片側対称族の巡回像であり、一般点では射影安定群が自明である。
四位置の点を番号付けし各位置の五点を枠に固定すれば、基底変換を除いた滑らかな
局所座標が得られる。番号の有限な選び直しは一般の安定群自明な点で自由なので、
商もそこで滑らかである。逆転$J$は因子を点ごとに固定し、全空間の一般点では
非自明である。因子に沿う座標の$J$平均と、横座標の反不変部分を取れる。
実際、滑らかな局所環上の有限位数作用は標数零で接空間の基底を固有ベクトルに
持ち上げられ、接空間上も恒等なら完備化上で恒等になってしまう。
従って横方向は符号表現で、局所不変環の横座標はその二乗となる。これは関数体の
因子付値についての指数を与える。射影像の正規化前の局所環を滑らかとは仮定していない。
偶対称因子の一般点では二奇配置の非同値性が保たれ、逆転は窓を固定しない。
\end{proof}

% !TeX root = ../papers/nc-p3.tex
\section{Euler導分による中心の持ち上げと十五次元成分の正則開集合}
\label{sec:quartic-euler-lifting}
$B$を四頂点で折り畳んだ経路長次数付き代数とし、$E(x)=n x$ ($x\in B_n$)を
Euler導分とする。以下の導分は各頂点冪等元を零に送る。従って内部次数$d$の
内導分の係数は$\bigoplus_i e_i B_d e_i$から取る。

\begin{proposition}\label{prop:euler-central-lifting}
体の標数を零、$d>0$とし、$B$の内部次数$d$の全導分が内導分と
$Z(B)_d E$の和であるとする。任意の局所完備Noether環上の次数別平坦な
斉次変形で、中心の次数$d$の非零因子は、同次数の中心元へ持ち上がる。
\end{proposition}
\begin{proof}
極大イデアルを$\mathfrak m$とし、$z\in Z(B)_d$の中心リフトが
$\mathfrak m^n$を法として得られたとする。その任意の次のリフト$y$を
頂点対角に選ぶ。$[y,-]$は$I=\mathfrak m^n/\mathfrak m^{n+1}$に値を持つ。
次数別平坦性と下の次数の中心性により、これはリフトの選択に関して
入力側で剰余体に降り、$B\to I\otimes B$の導分となる。
各$\ell\in I^*$で押し出した導分を$[v,-]+hE$と書く。
特に$z$のリフトとして$y$自身を使えば値は零である。
従って$0=dhz$であり、非零因子性から$h=0$である。
$I$の基底を取り、得られた内導分係数を$y$から引けば
$\mathfrak m^{n+1}$を法としても中心となる。完備化上では逆極限を取る。
ここで変形は頂点分解と次数を保ち、各斉次成分は有限階数とする。
\end{proof}

\begin{theorem}\label{thm:full-component-central-open}
十五次元成分の反対階数16の有理点$t$を含む一つの非空ザリスキー開集合で、
\emph{全ての幾何点}が標準Koszul左右AS正則となる。四次中心は五次元であり、
四つの中心四次式が生成する多項式環上で各始点が階数256の自由加群である。
\end{theorem}
\begin{proof}
$t$では既にこの自由表示と四次中心の文字次数別次元$(3,2,0,0)$を知っている。
次数4の導分は各矢を同じ端点の次数5成分に送り、六関係にはLeibniz則を課す。
文字次数$h$を固定すると各始点の四矢で合計56係数、四始点で224未知数である。
各関係の二文字の片方を順に五文字の商基底で置き換え、六文字商へ射影する。
\path{computations/verify_quartic_euler_derivations.py}が有理経路基底を
素数1009で還元し、$h=0,1,2,3$の順に階数$188,184,192,192$を検算する。
よって導分空間の次元の上界は$(36,40,32,32)$である。
他方、四次内導分の次元は$(33,38,32,32)$であり、各文字の四次中心の
Euler倍はこれと独立である。実際、中心四変数の一つに$hE$を作用させれば、
その非零因子と$4h$の積を得るが、内導分は中心を消す。各上界は従って等号である。

テンソル成分の$t$での完備局所環を取る。各Artin商上で隣接二切断の六次主小行列は
単元のままで、節\ref{sec:unrestricted-doublet}の次数別Koszul複体の分裂により
経路代数は次数別に自由である。命題\ref{prop:euler-central-lifting}を適用できる。
四次中心の基底は全て非零因子に取れる。最初の四つは認証した多項式変数とし、
残る元$z$にはその一変数$x$の一般定数倍を足す。多項式環の分数体上、
$x$の乗法は可逆なので$z+c x$の乗法の行列式は$c$について零多項式ではない。
元の代数も多項式環上で自由なため、適切な$c$で非零因子となる。
従って五つの独立な中心元が完備局所環上で同時に延びる。

次数4・5の経路空間を自由化した近傍で中心方程式はベクトル束の写像である。
完備化で五つの核切断は剰余でも独立なので、核次元が5未満となるために必要な
小行列式は完備環で全て零である。局所環から完備化への忠実平坦性により
それらは局所環でも零となる。一方$t$での階数を与える小行列式は単元である。
有限個の係数しか使わないため開近傍で四次中心は階数5の部分束となる。

最初の四変数をこの束の切断として取り、割った商の次数13までの上界が
$(1+t+t^2+t^3)^4$の係数となる有限個の開条件を課す。
$t$で全てが成立する。元の経路商の各次数には、この既約成分の正則な
幾何学的一般点からの標準Hilbert関数の下界がある。四中心式の多項式環からの
自由加群による全射とこの下界を比べ、定理\ref{thm:klein-rational-period-four}の
左右完全性の議論を適用すれば、同じ一つの開集合の全幾何点で結論を得る。
最後に開集合の基底変換による移動を合わせることで、印なしの主張となる。
\end{proof}

\begin{corollary}
この成分の閉包の標準AS点は全て多項式恒等式を持つ。
\end{corollary}
\begin{proof}
各終点の列は開集合で中心四変数上自由階数256である。
系\ref{cor:klein-pi-specializations}と同様に標準恒等式$s_{512}=0$を、
斉次な語に代入し、必要な有限個の経路次数の定階数性によって各標準AS点へ特殊化する。
\end{proof}

% !TeX root = ../papers/nc-p3.tex
\section{四次超曲面としての全中心とPI次数}
\label{sec:quartic-center-order}
前節の開集合で折り畳み環を$B$、多項式部分環を$P=k[x_1,x_2,x_3,x_4]$、
残る中心四次式を$z$とする。以下で中心の変数を通常の次数1に数えるときは、
元の経路長を4で割っている。

\begin{theorem}\label{thm:quartic-center-order}
$t$を含む単一の非空開集合の全幾何点で、全中心は
\[
 Z(B)=P[z]/\bigl(z^4+c_1z^3+c_2z^2+c_3z+c_4\bigr),
       \qquad c_i\in P_i
\]
の表示を持つ（左の$z$は環の元、右は表示変数）。その射影三次元多様体は
幾何学的に整かつ正規な四次超曲面である。$B$はこの中心上有限な素環で、
中心の分数体上で次数16の中心的単純環となる。各頂点の隅環の次数は4である。
特にこの開集合のPI次数は16であり、成分の閉包の全標準AS点が$s_{32}=0$を満たす。
\end{theorem}
\begin{proof}
各隅$e_i B e_j$は$P$上自由階数64で、全体は自由階数1024である。
まず有理点$t$で、既出の四変数の整数係数リフトを取り、文字次数0に残る
第三の独立な四次式を$z$とする。順序と原始整数係数への規格化は
\path{computations/verify_quartic_center_order.py}に指定する。
$B/(x_1,\ldots,x_4)$の第0頂点の隅で$z$の乗法の最小多項式は
素数1009の剰余体で$T^4$である。従って剰余の$1,z,z^2,z^3$は独立である
（隅の単位は$e_0$とする）。全折り畳み環でも同じ四元は独立で、
剰余基底へ延長して斉次リフトを取ると$B$の$P$自由基底になる。
特にこの四元の張る部分加群は直和因子である。

中心変数を$(1,2,3,4)$に指定すると、同じ素数で四頂点の隅の直和に
矢との可換式を課した核は次元4となる。第0隅での$z$の最小多項式は
\[
 T^4+421T^3+601T^2+542T+727\pmod {1009}
\]
で重根を持たず、特性多項式はその十六乗である。
同じ隅の64基底の左正則作用を$L_u$とすると、
$(\operatorname{tr}(L_uL_v))_{u,v}$の階数は64である。

これらの有限計算の特殊化を説明する。定数$\alpha_i$を指定した商で、
固定始点・終点の長さが$n$以下かつ$n$と法4で合同な語を扱う。
一文字を先頭に付けた四組の前次数基底に、$4\mid n$のときは単位も加える。
六関係と$x_i-\alpha_i$に残りの長さの基底語を右から掛けて掃き出す。
これは中央変数$y$を経路長4として$x_i=\alpha_i y$とした斉次商の次数$n$と同じである。
元が$P$上自由なので、この商は$k[y]$上自由であり、その次数$n$の次元は
$(1+t+t^2+t^3)^4$の該当する剰余類の次数$n$までの係数和である。
実際の検算は、全始点で$n\le16$までこの値と一致することを確かめる。
従ってこれらの有限範囲の主小行列は混標数局所環でも単元に取れ、
剰余体の計算は標数零での階数下界を与える。第0隅の乗法は
長さ12以下の基底に長さ4の切断を左から掛けて長さ16の基底へ移して求める。
正則作用のtraceは、同じ終点を持つ四つの安定した64次元空間上の
矢行列の積から求める。

$K=\operatorname{Frac}P$とする。中心は四隅の対角元が十六の矢と可換である
という$P$上の有限行列の核である。上の一点の階数から$K$上で核次元は高々4、
既知の四元から丁度4となる。自由直和因子の飽和性により
\[
 Z(B)=B\cap\operatorname{span}_K(1,z,z^2,z^3)
      =P\oplus Pz\oplus Pz^2\oplus Pz^3 .
\]
従って$z^4$の一意な表示が定理のモニック関係を与え、斉次性から係数次数も従う。
この段階では中心の既約性を使っていない。

この四次の平面切断を検査する。$(x_1,x_2,x_3,x_4)=(s,u,s+2u,3s-u)$とし、
$s=1$, $u=0,1,2,3,4$で上の有限乗法を計算する。
どの剰余でも第0隅における$z$の最小多項式は四次である。
よって各値で$1,z,z^2,z^3$に対する単元小行列を選べば、
$z^4$の一意な関係が標数零からその剰余の最小多項式に移る。
各係数は$u$で高々四次なので、Vandermonde行列で平面四次式全体が決まる。
全三射影座標層で偏導関数を検査すると、幾何学的な特異点がない。
従って$t$の標数零の平面切断も滑らかである。滑らかな射影平面曲線は
Bézoutにより幾何学的に整で、超曲面自身の可約性・重複成分もこれで排除される。
超曲面の特異軌跡が二次元以上ならこの平面に会うため、特異軌跡は高々曲線である。
超曲面はCohen--Macaulayなので余次元1の正則性から正規である。アフィン錐でも
同じ判定が通る。

上で用いた基底・有限階数条件と平面切断の滑らかさは$t$の一つの開近傍で保たれる。
より詳しくは、前節の四次中心の部分束から五切断を延ばし、自由表示の剰余基底を
取る。次数16までの係数比較と有限ファイバーのtraceをこの基底で作れば、全ての
主小行列と平面判別式は正則関数となる。従ってこの開集合の各幾何点で同じ議論が成立する。

最後に$K$上で$(1,2,3,4)$の検算から第0隅の正則trace対合は非退化である。
冪零根基に属する元は全ての積のtraceを消すので、第0隅は半単純となる。
$K$では$x_1$が可逆である。各頂点の$x_1$成分は長さ4の語の和で、どの語も
第0頂点を通る。従って$e_0$は充満冪等元であり、全代数も森田同値により半単純である。
中心は$K$上四次の体なので全体は中心的単純で、中心上の次元は$1024/4=256$、
各隅は$64/4=16$となる。元の環はその中央局所化に単射であり、非零な両側イデアルも
非零に局所化するので素環である。分裂体上では十六次行列環となり、標準恒等式
$s_{32}$が成立する。閉包の標準AS点へは系\ref{cor:klein-pi-specializations}と
同じ次数別定階数の論法を使う。
\end{proof}

\begin{corollary}\label{cor:quartic-division-corners}
同じ開集合で、中心の分数体$L$上の全商環は$M_4(\Delta)$であり、
$\Delta$は$L$上次数4の中心的可除代数である。特にどの始点・終点についても、
合成可能な二つの非零経路類の積は零にならない。
\end{corollary}
\begin{proof}
非零斉次中心元を全て可逆にした中心を$L_0[u,u^{-1}]$と書ける。
ここで$L_0$は射影中心の関数体、$u$には一つの長さ4の非零中心元を取る。
従って局所化した第0隅は$\Delta_0[u,u^{-1}]$で、
$\Delta_0$はその次数0部分である。$L=L_0(u)$に延長した隅が中心的単純なので、
$\Delta_0$も$L_0$上中心的単純で次元16である。

任意の有限生成次数付き右$B$加群の$L$上の次元は64の倍数である。
実際、頂点射影$e_iB$は全て$L$上の次元が64である。
$B$は中心上有限なのでNoether環であり、標準Koszul分解から次数付き大域次元は4である。
より具体的には、四単純の左線型分解から任意の次数付き加群の
$B/B_{\ge1}$へのTorは次数4より上で消える。有限生成加群の次数付き極小分解では
従って第5項は零で、全項は有限個の頂点射影の移動の和である。階数の交代和を取ればよい。

$\Delta_0$が可除でないなら、$0,1$と異なる冪等元$v\in\Delta_0$がある。
斉次中心分母を払えば$v=b/c$で、$b\in e_0Be_0$は斉次である。
右イデアル$bB$を$L$へ延ばすと、$e_0 B_L$の真の非零直和因子$v B_L$となる。
その次元は1以上63以下であり、先の整除性に反する。よって$\Delta_0$は可除である。
可除代数上の中心変数多項式環は整域なので、$L_0(u)$への延長$\Delta$も
有限次元の整域、従って可除である。同じ議論が各隅に適用でき、
四頂点冪等元は同値な原始冪等元をなす。各$e_i B e_j$を一行一列の
$\Delta$係数に埋め込めるので、最後の積の主張も従う。
\end{proof}

% !TeX root = ../papers/nc-p3.tex
\section{片側対称因子を横切る中心の二次・四次完全交叉}
\label{sec:symmetric-center-cover}
以下は偶二位置が対称となる因子$D$を扱う。基底付きテンソルの滑らかな成分を
$r$の近くで考え、中心の次数は経路長の四分の一で表示する。
$Z_j$をこの換算での次数$j$部分とする。

\begin{theorem}\label{thm:symmetric-center-cover}
$r$を含む成分の開近傍の全幾何点は標準Koszul左右AS正則である。
基底を局所的に選べば、$P=k[x_1,\ldots,x_4]$上の全中心は
$1,z,y,zy$を自由基底に持ち、$\deg z=1$, $\deg y=2$である。その表示は
\[
 z^2=a_2+a_1z+\tau y,\qquad
 y^2=b_4+b_3z+b_2y+b_1zy,\qquad a_i,b_i\in P_i .
\]
$\tau$は$D$の局所方程式である。因子上では$b_1=b_2=0$となるように
$y$を選べる。従って$D$上の射影中心は滑らかな三次元二次超曲面の二重被覆で、
分岐は四次式で切られ、経路の反転が被覆変換を誘導する。因子外では$y$を
消去して四次超曲面を得る。近傍の全幾何点で中心は幾何学的に整かつ正規、
折り畳み環のPI次数は16で、各隅の商可除代数の次数は4である。
\end{theorem}
\begin{proof}
$r$の四次中心の次元は5で、うち四切断は既出の多項式部分環を与える。
十五次元成分の一般点でも次元5であるから、経路長4,5の標準階数の開集合で
可換方程式の核は$r$の近くでも階数5の部分束である。四変数で割った次数13までの
判定をその切断に課せば、定理\ref{thm:full-component-central-open}の自由性と
左右完全性の証明がそのまま$r$を含む近傍に適用される。
一般点の四次超曲面表示から$\dim Z_2=15$を知るので、定階数の経路空間上の
可換方程式により$r$では少なくとも15である。

固定同定上、反転は長さ$4j$の語を逆順に読み、始点$i$を$3-i\pmod4$に送る。
二次関係の四つの表示からこれは反自己同型であり、中心では自己同型となる。
対称軌跡での同定の局所回復は定理\ref{thm:full-fifteen-component}による。
$K$を対角四元群の文字群とする。$r$で各矢と可換な頂点対角元を、
文字別に経路長9までの商基底で計算する。素数1009における
反転の正負固有空間の次元は
\[
\begin{array}{c|cccc}
 &0&1&2&3\\\hline
 Z_1&(3,0)&(2,0)&(0,0)&(0,0)\\
 Z_2&(8,1)&(6,0)&(0,0)&(0,0)
\end{array}
\]
である。商基底は右端から一文字ずつ延長して六関係の右倍を掃き出す。
各次数の経路次元は検算範囲で標準値と一致する
（各頂点で経路長$n$の次元は$\binom{n+3}{3}$）。
可換核の上界の合計が既知の下界と一致するため、局所整数環でも核は
直和因子で剰余への写像は全射である。$2$が可逆なので固有空間の射影も
特殊化と可換である。従って表は標数零で正しく、近傍の$Z_2$は階数15の束となる。
$D$上で反転不変な$Z_1$の四変数と第五切断$z$を取り、$Z_2$の
唯一の反不変直線の局所基底$y$を取る。この固有次元は因子上で局所一定である。
これらの切断を成分の近傍へ延ばす。

前節のフィルター付き有限商の基底と安定した矢作用で$z,y$を作用させる。
四変数の値を零とした第0隅で$z,y,zy$の乗法階数は剰余で順に32,2,1である。
特に次数の異なる$1,z,y,zy$の像は独立で、$P$自由な折り畳み環の
自由直和因子を張る。四変数の値$(1,2,3,4)$では剰余で可換核次元4、
第0隅の正則trace対合の階数64を得る。前節と同じ整数環からの階数比較と
分数体上の飽和性により、この四元が全中心の$P$自由基底になる。
斉次性により定理の二式を得る。これらが全関係を生成することは、
$P$上で$z,y$に比較用の重み2,3を付ければ従う。左辺の平方を右辺で
書き換えると重みが下がり、残る四単項式は独立である。
$D$上の固有符号から$\tau=b_1=b_2=0$が従う。

次に$\mu:\operatorname{Sym}^2 Z_1\to Z_2$を考える。
四変数の二次単項式十個と$x_i z$の四個は自由基底表示で独立である。
残る一列$z^2$の余核方向への係数が$\tau$なので、
$\det\mu$は単元倍の$\tau$である。$r$の唯一の横断商接方向で
その微分が非零であることを確認する。節\ref{sec:unrestricted-doublet}で
有理数上の軌道方向を除いて取った文字0の交代代表を$T$とする。
$p=1009$とし、$W+pT,W-pT,W^{\mathsf T}+pT^{\mathsf T},
W^{\mathsf T}-pT^{\mathsf T}$の六基底列を$\bmod\ p^2$で使う
（奇位置は二添字を交換）。経路長8までの商と長さ4の可換核は単元ピボットで自由である。
この核の文字0の三元と文字1の二元の全積を商に投影する。
文字0の九列には八つの単元ピボットがあり、残部は$p$で割れて$p^2$で
消えない。文字1の六列には六ピボットがある。
基点で階数14であることは既に標数零で分かった。これらの基底と掃き出しは
局所整数環上で単元のみを割る操作であり、基点で消える最大小行列の
$W+pT$での値は$p$とその微分の積に法$p^2$で等しい。
従って像を経路商に入れても非零の法線微分があり、$d\tau(T)\ne0$である。
滑らかな因子$D$上で$\tau=0$であったから、$\tau$はその局所方程式の単元倍となる。

$D$上では二次超曲面の座標環は$P\oplus Pz$、全中心はその上で$1,y$を
基底とする。分岐式は$g=b_4+b_3z$である。四変数を
$(s,t,s+2t,3s-t)$に制限する。$s=1$で$t=0,1,2,3,4$の五つの有限商の
正則作用から$z^2,y^2$を$1,z$で表示して補間すると、$r$の剰余では平面の
二次式のHesse行列が階数3となる。二次式と$g$に偏微分の二行二列小行列を
加え、全射影座標層で共通零点がないことを検算する。
従って標数零でも滑らかな平面二次曲線上で八個の分岐点は単純であり、
その被覆は滑らかで幾何学的に整な種数3の曲線である。
二次超曲面自身についても、$(1,i,j,k)$, $i+j+k\le2$の十点（非負整数）で
$z^2$の係数を補間すると五変数のHesse行列の階数は5であり、基底二次超曲面は滑らかである。
正確な行列計算は\path{computations/verify_symmetric_center.py}にある。

同じ二線型切断上の曲線の滑らかさと幾何学的整性は近傍で開く。
実際、二式で定まる重み付き射影族は四変数の射影空間上有限平坦であり、
指定した直線上への制限も有限平坦かつ固有である。重み2の頂点を通らないので
通常のJacobi判定を使える。全射影中心の各成分は三次元で、
特異軌跡が二次元以上なら有限射影で指定直線に会い、滑らかな曲線切断に反する。
アフィン錐も頂点以外は同じ正則性判定が通る。完全交叉ゆえに$S_2$、
従って正規で、次数0部分が体の正次数付き環には非自明な冪等元がないため整となる。
前節のtraceと斉次右イデアルの階数による議論が近傍全体で適用できる。
\end{proof}

% !TeX root = ../papers/nc-p3.tex
\section{中心の滑らかな部分と二十個の最小単純指標}
\label{sec:minimum-central-characters}
ここでは前二節の正規中心とPI次数16を持つ開集合を扱う。
体は代数閉、標数零とする。$B$は折り畳み環、$Z$は全中心、
$P$は四つの長さ4の中心変数の生成する環である。

\begin{proposition}\label{prop:smooth-azumaya}
$B$のAzumaya軌跡は$\operatorname{Spec}Z$の滑らかな部分に一致する。
$(u,v)\mapsto\operatorname{trd}(uv)$は$B\to\operatorname{Hom}_Z(B,Z)$
の同型を誘導する。ここで$\operatorname{trd}$は中心単純な一般商の被約traceである。
\end{proposition}
\begin{proof}
$B$は$P$上自由で、$\bar B=B/P_+B$は最高経路長12の斉次Frobenius代数である。
後者を確認する。正則中心列四元を一つずつ除くと、各頂点単純$S_i$に対し
$\operatorname{Hom}_{\bar B}(S_i,\bar B)$は内部次数の移動を除き
$\operatorname{Ext}^4_B(S_i,B)$に一致する。一元の場合は、非零因子の乗法による
二項分解から得る環変更のExtが一行に集中する。四段では変数の次数和16と
AS指数4の差により、各頂点に対応する左socleは経路長12に一つずつある。
このsocle上で各頂点成分に非零な長さ12の線型汎関数を選ぶと、その核は
非零左イデアルを含まない。実際、各単純socleは相異なる頂点型で重複しない。
従って積と汎関数から得る対合は完全である。

この汎関数を$P$自由な斉次基底で、経路長を12だけ下げる$P$線型汎関数へ
持ち上げる。積対合で定まる左$B$加群写像は斉次であり、$P_+$剰余で同型なので
次数付き中山と自由階数の一致から同型となる。一方$Z$は$P$上の四自由基底の
最後が長さ12で、その係数を取る写像が$Z\simeq\operatorname{Hom}_P(Z,P)$を
同じ次数移動付きで与える。実際、$P_+$剰余は四次モニック一元商または
重み付き二次・四次完全交叉で、最高次のsocleが一次元である。
随伴と二つの移動の相殺から、$\operatorname{Hom}_Z(B,Z)$と$B$は
次数移動なしに左$B$加群として同型である。

被約traceは$Z$に値を取る。各元は$Z$上整であり、分裂体での固有値の和は整、
かつ中心の分数体に入り、$Z$は正規だからである。次数作用との可換性により
trace対合の写像は次数0で、一般商で非退化である。先の左加群同型と合成すると
$B$の次数0の左加群自己準同型、すなわち$\sum_i c_i e_i$の右乗法になる。
一般商で可逆なので全$c_i\ne0$であり、元の写像も同型である。

$B$は有限な$Z$加群として極大Cohen--Macaulayである。
各極大点ではその下の$P$の四局所パラメータが正則列であり、局所化でもこの性質を保つ。
従って正則な中心局所環上ではAuslander--Buchsbaum公式により自由となる。
自由基底と被約trace双対基底を$u_j,v_j$とする。一般商での行列単位の計算から
$\sum_j u_j\otimes v_j$は両側作用を交換し、積の像は$16\cdot1$である。
両側の自由性によりこの等式は局所環上で成立する。16で割った元が分離元を与え、
中心が基底環に等しい有限射影的分離代数はAzumayaとなる。
逆にAzumaya点では$B$の局所化は中心上自由である。Koszul両側分解から
$B$の通常の大域次元は高々4であり、中央局所化後もこの上界を持つ。
ここで両側分解は$B\otimes_{B_0}W_j\otimes_{B_0}B$に左右の端文字の微分を置いたもの
で、各全次数で右因子次数により濾過すると片側Koszul完全性から対角$B$を解消する。
局所中心の剰余体を$k'$とし、$B\otimes_Z k'$の射影分解を中心へ制限する。
$k'$はこの自由な$k'$空間の直和因子なので中心上有限射影次元を持ち、
局所中心は正則である。
\end{proof}

\begin{theorem}\label{thm:twenty-minimal-characters}
中心錐の頂点外にある各単純右加群の頂点次元は$(d,d,d,d)$で、
$d\in\{1,2,3,4\}$である。十五次元成分の一般窓では、$d=1$の同型類は
次数によるスカラー作用を除くと丁度二十個であり、対応する射影中心上の点は
全て異なる特異点である。各点の一つの非零アフィン代表上には丁度二種類の
単純加群があり、次元4のものと、次元8または12のものからなる。
\end{theorem}
\begin{proof}
任意の頂点外の中心指標ではある$x\in P_4$が可逆である（添字4は経路長）。
実際$Z$は$P$上斉次有限である。$S$をこの指標の単純右加群とし、
$M_n=Se_{n\bmod4}$を$n\ge0$に置き、矢を$S$の作用で掛ける。
$x$は$M_n\to M_{n+4}$の同型を与え、$M$は次数付き有限生成である。
$B$はNoetherで、有限長の次数付き射影分解の各項は頂点射影の移動の有限直和である。
各頂点射影の全Hilbert級数は$(1-t)^{-4}$なので
\[
 H_M(t)=\frac{d_0+d_1t+d_2t^2+d_3t^3}{1-t^4}
       =\frac{h(t)}{(1-t)^4},\qquad d_i=\dim Se_i .
\]
非自明な四乗根三個で分子が消え、四次元数は等しい。PI次数16または
一般行列環の標準恒等式と単純像の全行列性から$4d\le16$である。

$d=1$なら各辺の四矢の値は同時に零ではない。逆に、四辺で非零な一次元空間を
結ぶ矢組は有向閉路をなすため単純である。関係式は各辺の双一次式そのものである。
第一辺は有限エタールな二十点、残りの三辺はグラフで、四移動は恒等である。
従って射影化した四組は丁度二十個で、スカラーを戻すと各組につき四辺の倍率は
頂点基底変換によりその積一つを残す。次数作用でこの積を任意の非零値へ移せる。
各単純は頂点外である。頂点指標の商は有限次元の正次数付き環であり、
正次数イデアルが冪零だから単純は四つの頂点単純に限られる。
一つの非零な$x$の値を固定すれば積が決まり、一つの射影組から一種類となる。

二十指標の分離を検算する。$t$の二つの切断の法線写像の共通核から、
節\ref{sec:unrestricted-doublet}の有理商接代表十五個の固定整数線型結合$T$を選ぶ。
$p=2657$では第一辺が完全分裂する。$W+pT$から二次関係を取り、法$p^2$の
経路長5までの商基底で四次の可換核の自由基底五元を求める。
各射影四組では各位置で一座標を1とし、残り十二座標の修正を二十四本の
辺方程式の微分で解く。全てで十二ピボットがあり、修正後の全式が消える。
第0隅の中心五元の各経路を四矢の積で評価し、単元座標で正規化すると、
剰余体では八値、法$p^2$では二十値が異なる。
\path{computations/verify_minimum_central_characters.py}が同じ接線と番号付けで再現する。

この分離は標数零の一般窓に移る。番号付き辺はエタールで、$t$の成分は滑らかなので、
標数零の接ベクトル$T$はこの被覆へ一意に持ち上がる。分裂素数での各番号は
$\mathbb Z_p$上へHenselで延びる。中心と経路の基底は検算の単元ピボットで取り、
一次微分の持ち上げも$\mathbb Z_p[\epsilon]/(\epsilon^2)$上で計算できる。
実際、余分な行の等式は$\mathbb Q_p[\epsilon]/(\epsilon^2)$上で
幾何学的な一次変形により成立し、単元による解の係数は整である。
$\epsilon\mapsto p\bmod p^2$による像が上の検算である。
よってどの二切断も近傍で恒等的には一致しない。番号付けは有限エタールであり、
対角を除いた対の指標一致軌跡の像を除けば一般窓で二十点全てを分離できる。
次数4の単純はAzumaya点の単純次元16より小さいため、各中心指標は前命題により
特異で、その射影点も特異となる（$x$を可逆にした中心は換算次数1の可逆元を持つ）。

最後に、一つの射影点上で$x$の値を1に規格化する。最小単純$S$の自己一次変形は
$B/(x-1)$上では全て自明である。実際、頂点基底を選んだ双対数上の変形でも
各頂点は自由階数1で、二十点のエタール性とグラフ性から各射影矢は固定である。
四倍率の積は$x=1$により固定で、残りは基底変換で消せる。
この三次元の中心を指標で完備化しても同じ自己拡大の消滅が成立する。
長さ2の拡大は中心極大イデアルの二乗で消え、完備化前後で同じだからである。
完備環上有限な代数は半完全である。もし単純が$S$しかなければ、その森田同値な
基本代数の根基$J$は$J/J^2=0$となる（これは単一単純の自己拡大の双対）。
中山により$J=0$で、中心が正次元であることに矛盾する。
ここで$x=1$近傍の中心は三次元の射影チャートである。例えば中心の$x$を
可逆にしてその四乗根をエタールに加えると、経路次数の作用で
$B/(x-1)$と可逆な一変数のテンソル積となる。中心、一般商の次数、有限性は
この自明化と両立する。

この完備中心は整域である（標数零の正規な多様体の局所環は優秀で正規性が完備化で保たれる）。
完備代数はこの環上torsion-freeで、一般商は次数16の中心単純環である。
各単純の次元を$n_j$とすると、半単純剰余の各行列環の原始対角冪等元は
直交した完備系へ持ち上がり、同じ行列環内では同値に持ち上がる。
一般分裂体で各原始元は正の行列階数$m_j$を持つので
$16=\sum_j n_jm_j$である。$n_j=4d_j$で割り直すと$4=\sum_j d_jm_j$。
$d=1$の単純は指標につき一つであり、ほかに少なくとも一つあるから、
残る可能性は$(d,m)=(1,1),(3,1)$または$(1,2),(2,1)$だけである。
\end{proof}

% !TeX root = ../papers/nc-p3.tex
\section{二十の中心指標での通常二重点と局所行列表示}
\label{sec:twenty-central-nodes}
\begin{theorem}\label{thm:twenty-central-nodes}
十五次元成分の一般窓の射影四次中心上で、前節の二十指標は全て通常二重点である。
その各点の単純加群の次元は4と12である。中心をその点で完備化し、
可逆な四次変数の値を1にして次数方向を除くと、中心と折り畳み環は
\[
 R\simeq k[[a,b,c,d]]/(ad-bc),\qquad
 B^{\wedge}\simeq\operatorname{End}_R(R^{\oplus4}\oplus I^{\oplus12})
\]
を満たす。$I$は一般二行二列行列の余核である（座標の選び方による）。
特にこの二十点はそれぞれ孤立している。
\end{theorem}
\begin{proof}
$A=e_0Be_0$とし、長さ4のループの固定した整数線型結合$u\in A$を取る。
その中心分数体上での四次被約特性多項式を$H(T)$とする。
中心は正規で、整な元の固有値の基本対称式も整であるから$H\in Z[T]$である。
四次中心の第五変数$z$の$P$上のモニック関係を$F$とする。$A$の$P$自由階数は64で、
中心四次拡大の一般分裂体では四つの四行行列環の直積になるので、補助変数$v$を入れて
\[
 \det\nolimits_{P}(T+vL_z-L_u)
   =\operatorname{Norm}_{Z/P}\bigl(H(T+vz)\bigr)^4,\qquad
 \det\nolimits_P(T-L_z)=F(T)^{16}.
\]
$L$は左正則作用である。これらは係数環上の多項式恒等式で、重根を持つ有限中心
ファイバーにも特殊化できる。特に中心根$\alpha$が$F$で重複度$e$なら、
$H_\alpha$の根$\lambda$の重複度$m$は、任意の定数$v$に対して
$L_u-\lambda+v(L_z-\alpha)$の零固有値に少なくとも$4em$を寄与する。
これはArtin中心のノルムが各剰余値を局所長だけ掛けた積であることによる。

この恒等式を有理点$t$で$p=2657$へ還元する手続きを明記する。
既定の四中心変数と$z$は局所整数環上に取り、四変数を零とした商で
規定の剰余基底と$1,z,z^2,z^3$の独立性を検算する。長さ13の剰余商は零である
（前節までのフィルター商の、零パラメータでの長さ13と長さ9の同一次元および
安定埋め込みによる。斉次剰余商の検算も下記コードに含める）。
従ってこれらの経路リフトから$P$自由加群の全射を得る。標数零で既知の
Hilbert関数に一致し、自由加群の部分加群である核はtorsion-freeなので、
局所整数環でも同型である。中心四元を含む斉次剰余基底へ延長でき、
$z^4$の関係も整係数となる。この整モデル$P[z]/F$は正規である。
実際、標数零では既知の正規環であり、特殊ファイバーも$(1,2,3,4)$で
$F$が四つの単根を持つので一般に被約である。従って垂直な高さ1素点では
$p$が極大イデアルを生成し、水平な高さ1素点でも正則である。
超曲面なので$S_2$であり、平坦性と標数零での整性から整域でもある。
従って$H$の係数はこの整モデルに入り、上の恒等式を剰余へ移せる。

二十の番号付き四組に四次中心元と$u$を評価し、$x_1=1$に規格化する。
全四組で$x_1$の剰余は非零である。$z,u$の値を$\alpha,\lambda$とすると、
第0隅の64次元正則作用で$L_z-\alpha$の零固有値の重複度は各々32となる。
$v=1,3,7$の三値で$L_u-\lambda+v(L_z-\alpha)$の零固有値重複度の
最小は各々8である。従って$e=2$であり、上の下界から$m\le1$となる。
$\lambda$は$H_\alpha$の根である。実際、$H(u)=0$を一次元の隅表現で評価すればよい。
従って$H_\alpha'(\lambda)\ne0$かつ$F''(\alpha)\ne0$である。
\path{computations/verify_central_nodes.py}は$u$の係数と以上の計算を再現する。

この二つの非消滅は各番号について標数零の$t$に持ち上がり、成分の一般窓へ移る。
ここで$H$を族でも正則に取れることを補足する。$t$の近くの滑らかなテンソル基底上で、
$P$自由表示、$Z=P[z]/F$、有限な経路の乗法表を前節までの単元小行列で取る。
全中心族は平坦で幾何学的に正規なファイバーを持ち、基底が滑らかなので正規である。
一般商の被約多項式の係数は再び整性からその環に入る。各ファイバーの一般の四変数の値で
$z$の四根は相異なり、上のノルム恒等式で$v$を不定として四根ごとの因子を分離すれば、
特殊化した$H$はその隅の被約多項式に一致する。実際、各因子の$T$の根は
$v$について傾き$-\alpha$を持ち、同一の傾きのモニック因子の積は一意である。
これで特異な中心指標にも適用できる。同様に、$H'_\alpha(\lambda)\ne0$の下で
最小単純に対応する原始冪等元の一般行列階数は1となる。詳しくは、
指標での被約traceは有限ファイバーの根基上で消える。根基元の高い冪は
中心極大イデアルと$A$の積に入り、そこでは全ての正冪のtraceが剰余で零となる。
Newton恒等式により根基元の全固有値は零である。従って誘導traceは各単純行列因子で
通常のtraceの定数倍となり、原始冪等元における値は持ち上げの一般行列階数である。
再びNewton恒等式から$H_\alpha$は各単純の特性多項式をこの階数だけ掛けた積となる。
$S e_0$は一次元で$u$の値が$\lambda$だから、その重複度1により前節の
$(1,2),(2,1)$の可能性が消え、残る単純の全次元は12となる。

次に$x_1=1$とし、$R$をこの射影中心の局所完備環とする。
点は特異であり、$F''(\alpha)\ne0$より四変数正則局所環中の超曲面の
重複度はちょうど2である。完備化した折り畳み環の二つの単純に対応する
原始冪等元を$p,q$と書く。ともに一般行列階数1なので一般商は分裂し、
$pB^{\wedge}p=R$となる。実際これは分数体の中で$R$上整な有限環である。
$B^{\wedge}$は$R$上極大Cohen--Macaulayなので反射的である。
この性質と高さ1でのAzumaya性は前節のエタールな次数自明化と完備化にも保たれる。
より詳しくは、完成写像の形式ファイバーが幾何学的に正則なので、完成した中心の
高さ1の点ではその下の点も正則である。よって左右乗法から
\[
 B^{\wedge}=\operatorname{End}_R(B^{\wedge}p)
\]
を得る。一般商では同型であり、高さ1のAzumaya環では階数1の$p$が充満で
森田同型となり、両辺の反射性から全体でも同型である。
$I=qB^{\wedge}p$と置けば原始冪等元の個数から
$B^{\wedge}p=R^{\oplus4}\oplus I^{\oplus12}$となる。
$I$は階数1の極大Cohen--Macaulay加群で、自由なら$p,q$が同値になり矛盾する。

$J=I^* I\subset R$を評価の像とする。基本代数
$\operatorname{End}_R(R\oplus I)$の$R$隅への射影は、$J$を割ると
$R/J$への代数全射になる。極大イデアルを$\mathfrak m$とすると$J=\mathfrak m$である。
実際$J$に単元があれば$I$は自由となる。一方$R/J$が非自明な余接空間を持てば、
その双対から剰余単純の非自明な自己一次拡大を得る。全射と森田同値で引き戻しても
非分裂であり、前節の$x_1=1$での最小単純の剛性に反する。中山より結論が従う。

最後に$R=S/(f)$、$S$を四変数の完備正則局所環とする。
$I$は$S$上深さ3なので最小自由表示は長さ1で、正方行列$M$による。
その行列式は単元倍の$f$である。実際、その素因子は加群の台から$f$だけであり、
$(f)$で局所化した長さは階数1ゆえ1である。$M$の成分は全て極大イデアルに入り、
$f$の重複度が2なのでサイズは高々2で、非自由性から丁度2である。
転置の核は$R$上で随伴行列の転置の像に一致する
（$S$に持ち上げ、積が$f$倍なら随伴行列を掛けて$f$を消去する）。
従って$I$の二生成元と双対との評価イデアル$J$は随伴行列、従って$M$の
四成分で生成される。$J=\mathfrak m$なのでこの四成分は四次元余接空間の基底をなし、
$S$の正則パラメータ系である。これを$a,b,c,d$とすれば主張の表示を得る。
\end{proof}

\begin{corollary}
一般の射影中心に残る特異点があれば、各点の単純加群は全て次元8で、
一種類または二種類である。二種類ならその完備環上の代数は
$\operatorname{End}_R(R^{\oplus8}\oplus J^{\oplus8})$であり、
$J$は非自由な階数1反射的加群である。
\end{corollary}
\begin{proof}
前節の完備化での階数和$4=\sum d_jm_j$を使う。
$d=4$があれば一種類だけで$m=1$である。全原始冪等元が同値で、
一般行列階数1の隅は正規中心に等しいため、完備化は十六次行列環である。
忠実平坦な降下によりAzumaya点となり、特異点に矛盾する。
$d=1$は既に二十点で尽き、$d=3$には$d=1$の相方が必要なので、
残る全ての$d_j$は2である。二種類なら両$m_j=1$であり、先の高さ1の
森田同型を反射性で延長する証明をそのまま適用できる。
\end{proof}

% !TeX root = ../papers/nc-p3.tex
\section{四中心切断を持つ窓に対する次数16の完全判定}
\label{sec:central-window-cutoff}
この判定は成分、反対階数、PBW性を仮定しない。四窓に六つずつの二次関係を与え、
有限双対の標準Frobenius検査を先に済ませる。
以下の次数は全て元の経路長である。
\begin{theorem}\label{thm:central-window-cutoff}
折り畳み環$B$に四つの中心四次式$x_1,\ldots,x_4$があり、
$Q=B/(x_1,\ldots,x_4)$が有限次元とする。このとき標準Koszul左右AS正則性は
次の有限条件と同値である。
\begin{enumerate}
\item 各始点の$Q$の次数$0$から13までの次元は$(1+t+t^2+t^3)^4$の係数に一致する。
\item 各始点の$B$の次数$n\le16$の次元は$\binom{n+3}{3}$であり、
左右の標準線型複体は内部次数16まで単純加群を解消する。
\end{enumerate}
特に$Q_{13}=0$だけで有限性の仮定を認証でき、以後の判定は経路長16を超えない。
適切な切断を持つ標準AS窓の近くで、四切断を含む有限Frobenius候補の
パラメータ空間上、これらの認証条件は開く。
\end{theorem}
\begin{proof}
$P=k[X_1,\dots,X_4]$, $\deg X_i=4$を$x_i$で作用させる。
$Q$が有限なら次数帰納で$B$は$P$上有限である。まず$B$を標準ASとする。
中央Koszul複体$K(X;B)$の正のホモロジーも含め全ホモロジーは$P_+$で消え、
有限次元の次数付き$B$加群である。有限次元次数付き加群は頂点単純の移動による
有限列を持つ。AS双対性と長完全列により、非零なそのような加群$H$について
$\operatorname{Ext}_B^p(H,B)$は$p=4$だけで非零となる。
各項が$B$自由な中央Koszul複体に$\operatorname{Hom}_B(-,B)$を施すと、
スペクトル系列の$\operatorname{Ext}^p_B(H_q,B)$は全次数$p+q$へ収束する。
$p=4$の一列しかなく、複体の長さは4なので、$q>0$の$H_q$は零である。
従って$B$の$P$上の高次Torは消え、極小次数付き自由分解から$B$は$P$自由である。
各始点の標準Hilbert級数に$(1-t^4)^4$を掛ければ第一条件を得る。第二条件は正則性から従う。

逆に二条件を仮定する。$Q$の各始点・終点の斉次基底を経路から選び、
その語$w$に対応する記号を基底とした自由$P$加群$F$を取る。最長は12である。
$Q_{13}=0$と一次生成性より全高次も零であり、評価$F\to B$は全射である。
第一条件から$F$の各始点級数は$(1-t)^{-4}$で、第二条件によりこの評価は
次数16まで同型となる。ここから全次数の同型が従うことを示す。
各矢$a$の左作用を、基底語への積$aw$の一意な座標表示で定義し、$P$線型に
$L_a:F\to F$へ延ばす。端点が合わない積は零とする。この操作は長さ13までで済む。
評価と矢は可換である。関係$r$の作用$r(L)$は各基底上で評価が零で、
長さは高々14なので$F$でも零である。全ての中心切断についても
$x_i(L)-X_i$は各基底上で長さが高々16、評価が零なので作用は零である。
従って$F$は$B$加群である。頂点単位の和に語を作用させれば、
対応する基底記号は次数12までの単射性により全て回収され、
さらに$X_i$の作用は$x_i$で回収される。これで$B\to F$が評価の逆写像となる。
特に標準Hilbert関数と自由性が全次数で確定した。

左右どちらかの頂点単純に対する線型複体を、増大写像付きで$C_\bullet$と書く。
各項は$P$自由で、係数空間の次数移動は0から4なので自由生成元の次数は高々16である。
同じ単純の$P$自由分解は四変数の通常のKoszul複体$D_\bullet$であり、
生成元次数は$0,4,8,12,16$である。増大写像の恒等を持ち上げて
$D_\bullet\to C_\bullet$を作れる。実際、各自由生成元の像を選ぶ際の閉元は
内部次数16以下にあり、仮定した範囲の完全性で境界として解ける。
比較写像の錐は次数16まで完全で、自由生成元に16を超えるものがない。
定数の単元成分を逐次縮約して極小化する。非零な項が残れば、内部次数が最小の
自由生成元は正次数係数の微分で零に移り、境界にもなれない。
この生成元の次数は16以下であり矛盾する。従って錐は可縮で、もとの複体が
全次数で単純を解消する。有限Frobenius性と合わせて求める標準AS性が従う。

最後に正則な基点で、次数16までの左右線型複体の単元小行列を開く。
次数別に帰納すれば、前次数の商空間が束であり、複体左端からの分裂完全性で
次の次数の商も標準階数の束となる。同じ範囲の中央Koszul複体も基点で
正の位置が完全であるから分裂する小行列を選べる。従って次数13までのその
余核の次元も所要値となる。有限個の開条件だけで定理を適用できる。
\end{proof}

\path{computations/verify_central_window_cutoff.py}の\texttt{certify}は有理係数で
この判定を実施する。四つの候補切断の中心性を長さ5で確認し、六関係と
四中心切断を挿入した語空間を長さ13まで構成する。$Q_{13}\ne0$で有限性が
別途未認証なら切断未確認を返す。切断を確認できた場合には、語の逐次商と
双対交差空間の一文字縮約から左右の複体の像階数を長さ16まで比較し、
正則・非正則のどちらも有限に判定する。
四切断を存在量化する場合、中心性とこれらの有限行列条件は可構成で、
標数零の代数閉体上でパラメータを消去できる。全標準AS点がこのような四切断を
持つことは必要とせず、その被覆は別の問題である。

\begin{corollary}
中心四次式を指定しない入力でも、$B$が四次中心全体の生成する部分環上有限である
という範囲では、同じ次数16まででAS正則性を決定できる。
\end{corollary}
\begin{proof}
長さ5の可換方程式で四次中心の基底を求め、独立な不定係数で四本の一般線型結合を取る。
標準ASなら$B$の成長次元は4で、有限性によりその中心部分環のKrull次元も4である。
この可換環は換算次数1で生成されるため、無限体上で一般の四線型式が
斉次Noether正規化を与える。従ってこの四本の剰余は有限次元であり、
定理の必要条件から長さ13で必ず零になる。よって有理関数体上でその次数が
非零なら、前記の有限性の仮定の下で非ASである。零なら定理の二条件を適用し、
体の忠実平坦な拡大に沿う完全性の降下で結論を得る。
この有限性には、全四次中心で割った商が長さ13で零という直接確認できる十分条件もある。
\end{proof}

% !TeX root = ../papers/nc-p3.tex
\section{片側有限な非中心パラメータによる全次数AS判定}
\label{sec:onesided-parameter-cutoff}
元の環の中で各パラメータが正規であることすら要しない有限判定を与える。
$B$ を四窓を折り畳んだ正次数付き環とし、二次双対の標準有限Frobenius検査を済ませる。
$P$ は四つの斉次生成元 $X_i$、次数 $d_i>0$ を持つ既知の重み付き反復Ore拡大とする。
各段の自己同型は次数を保ち、捩れ導分はその段の新しい文字の次数を持つ。
特に $X_jX_i=q_{ij}X_iX_j$ の量子多項式環を許し、$q_{ij}$ の位数を限定しない。
斉次経路の和 $x_i\in B_{d_i}$ が $P$ の定義交換式を満たすとする。
$P\to B$, $X_i\mapsto x_i$ は単位的で、最初から単射とは仮定しない。
$D=\sum d_i$, $L=D-4$、右 $B$ 加群 $Q=B/(\sum_i x_i B)$ と置く。

\begin{theorem}\label{thm:onesided-parameter-cutoff}
$Q$ が有限次元なら、標準14641型Koszul左右AS正則性は次と同値である。
\begin{enumerate}
\item 各終点で $Q$ の次数 $0,\ldots,L+1$ の次元が
$\prod_i(1+z+\cdots+z^{d_i-1})$ の係数に一致する。
\item 元の各経路次元が次数 $D$ まで標準値で、各左頂点単純の標準線型複体は
内部次数 $D$ まで完全である。
\end{enumerate}
第一条件が成立すれば $Q_{L+1}=0$ が有限性自体を認証するため、十分性には別の有限性証明を要しない。
成立時に $B$ は左 $P$ 加群として自由である。従ってこのような四パラメータを持つ枝への所属は、
$d_i$ と $P$ の交換式を固定すれば、全次数を調べず存在量化して判定できる。
\end{theorem}
\begin{proof}
$P$ は順序単項式基底を持ち、左・右Noetherである。後者はOre拡大の次数フィルトレーションと
可逆自己同型に対するHilbert基底定理を各段で使えばよい。自明な左右 $P$ 加群 $k$ は、
長さ4、生成次数が全て $D$ 以下の次数付き自由解消を持つ。
例えば前段の解消を新しい環へ誘導すると、新しい一文字だけの $k[X_j]$ の解消になる。
実際、これまでの正次数イデアルは自己同型と導分で保たれ、拡大は前段上で左右とも自由である。
商での $X_j$ 乗法は新しい環の加群写像であり、次数 $d_j$ だけ移した解消から比較写像を持ち上げ、
その錐を取れば $k$ の解消を得る。生成次数は従来のものと $d_j$ を足したものになる。

$Q=k\otimes_P B$ が有限次元なら、次数帰納で $B$ は左 $P$ 加群として有限生成である。
$B$ が標準ASなら $\operatorname{Tor}^P_q(k,B)$ は全て有限次元である。
これは左 $P$ 加群 $B$ の各段有限生成な極小解消を取って $k$ をテンソルすれば分かる。
同じTorを $k$ の右 $P$ 自由解消に $B$ をテンソルして取れば、
長さ4の右 $B$ 自由複体のホモロジーである。AS双対性により、非零な有限次元次数付き右加群 $H$ は
$\operatorname{Ext}^p_B(H,B)$ が $p=4$ だけで非零である。この複体への
hyper-Ext系列の $\operatorname{Ext}^p_B(\operatorname{Tor}^P_q(k,B),B)$ は
全次数 $p+q$ に収束する。残る一列に微分はなく到達次数は4以下だから $q>0$ のTorは零である。
極小解消に戻れば $B$ は左 $P$ 自由であり、終点冪等元ごとにも同じ結論を得る。
$P$ の級数 $\prod_i(1-z^{d_i})^{-1}$ で標準級数を割れば第一条件となる。

逆に $Q$ の各終点の基底を斉次経路 $v$ で持ち上げ、それらを記号とする自由左 $P$ 加群
$F$ を作る。$Q$ は頂点単位で生成される右 $B$ 加群なので $Q_{L+1}=0$ の後は全て零である。
従って評価 $F\to B$ は全射である。各終点で級数が標準であり、第二条件により次数 $D$ までは単射である。
矢 $a$ の右作用を各基底で $v a$ の一意な座標として定め、左 $P$ 線型に延ばす。
これには次数 $L+1$ までで足り、端点の異なる積は零に取る。二次関係は各基底で
次数 $L+2\le D$ の評価が零なので $F$ 上でも零である。従って $F$ は右 $B$ 加群にもなる。
頂点単位の基底記号の和を $u$ とすると、$u\cdot v$ は各基底記号そのものであり、
$u\cdot x_i=X_i u$ も次数 $d_i\le D$ の単射性から従う。
右作用が左 $P$ 線型であることを使うと、任意の $P$ の語 $p$ で
$u\cdot p(x)=p(X)u$ となる。従って $b\mapsto u\cdot b$ は評価の全次数の逆である。
この段階で $x_i$ が他の矢や経路とスカラー倍で交換することは使わない。

左頂点単純の標準複体を左 $P$ 加群として見れば、各項は終点別の $B$ の直和と次数移動である。
上で得た自由基底と四次までの移動から、その生成次数は $L+4=D$ 以下である。
$P$ に制限した頂点単純は自明な $k$ なので、最初に構成した自由解消からこの複体へ
増大写像を持ち上げられる。持上げに必要な完全性は内部次数 $D$ 以下だけである。
比較写像の錐の定数単元を消去し極小化する。非零な項が残れば、残る全自由生成元中で
最小内部次数のものは微分が零で境界に入らず、その次数は $D$ 以下となる。
これはその範囲の完全性に反する。よって左線型解消が全次数で完全である。
半単純頂点環のbar複体（またはKoszul双対性）によりKoszul性は左右共通で、
先の有限Frobenius条件から両側AS正則性が従う。
\end{proof}

\begin{remark}
$yx=-xy$, $wz=2zw$ で二組の間は可換とする四文字代数を取る。
$(x+y,y^2,z,w)$ は次数 $(1,2,1,1)$ の量子パラメータである。
$x+y$ の右倍のうち $x,y$ を掛けたものは $x^2-xy,xy+y^2$ を張るが、
$x(x+y)=x^2+xy$ はこの空間にない。従って最初のパラメータは正規元でない。
$y(x+y)=-(x+y)y+2y^2$ により左パラメータ基底は $1,y$ である。
また四周期の隅にある $z^4,w^4$ は無限位数の交換倍率を持つため、この隅はPIでなく
四次中心上有限の場合の判定では扱えない。ここでは経路長5までで正則性を認証できる。
\path{computations/verify_quantum_parameter_cutoff.py} はこの例で片側商、関係の倍率、有限複体を検査する。
\end{remark}

% !TeX root = ../papers/nc-p3.tex
\section{同時零二次式の境界におけるHilbert関数の超過}
\label{sec:zero-pairing-defect}
節\ref{sec:klein-zero-quadratics}の表の左上、すなわち
$a=1$, $d=z$, $c=b=1/\lambda$を取る。開いたチャートの条件を保つ。
\begin{proposition}
この境界の$x=z-1$の二変数部分では全ての窓が非標準ASである。
また$(z,\lambda)=(3,2)$に固定した双一次曲線の全ての窓も非標準ASである。
この曲線では$x=4$の第0始点は経路長5まで標準次元を持つが、長さ6で92となる。
\end{proposition}
\begin{proof}
$H=0$を解くと、両行列は
\[
 A_0=\begin{pmatrix}1&1/\lambda\\1/\lambda&z\end{pmatrix},\qquad
 A_1=\begin{pmatrix}1/\lambda&1\\ z-x&1/\lambda+\lambda y\end{pmatrix},\qquad
 y=\frac{(z-1)^2(\lambda^2 x-(\lambda^2z-1))}
 {\lambda^2((\lambda^2z-1)x-(z-1)^2)}
\]
となる。$x=z-1$では分母の括弧が$(z-1)z(\lambda^2-1)$になる。
$\lambda^2=1$のときはこの制限上で分子も消えるため、
この二変数部分とは上式の有理グラフの閉包を指す。
$(3,2)$では分子と分母に共通根はなく、指定チャートの曲線全体がこのグラフに入る。
それぞれの関数体上で各始点の二次関係から語商を計算する。
前次数の基底に末尾の四矢を掛け、二つ前の基底に末尾の六関係を掛けた行を消去する。
これは前の位置までの関係を先に割り、最後の位置に入る関係だけを追加した商である。
関係基底の六列は一つの許容標本で独立な列を選ぶ。既知の両切断階数6により
関数体上で全関係を張る。計算は有理関数の多項式最大公約数による正規化と
厳密な掃き出しである。$\mathbb Q(z,\lambda)$上で$x=z-1$とすると
長さ4の四始点の次元は$(35,36,36,35)$、
$\mathbb Q(x)$上で$(z,\lambda)=(3,2)$とすると長さ6では$(92,84,85,84)$である。
後者の$x=4$で第0始点の長さ0から6までの次元は$(1,4,10,20,35,56,92)$である。
\path{computations/verify_zero_pairing_defect.py}がこれらを検算する。
階数6を保つ部分で関係空間は部分束なので、基底の取り替えを許せば、
経路次元は任意の特殊化で減らない。従って標準値35または84の超過は
それぞれの全閉包のチャート内の各点へ引き継がれる。
\end{proof}

% !TeX root = ../papers/nc-p3.tex
\section{片方の行列式二次式が恒等的に消える全境界}
\label{sec:single-zero-pairings}
\begin{theorem}\label{thm:single-zero-pairings}
定理\ref{thm:klein-all-pairings}の可逆性を保つチャートでは、
$P,Q$の少なくとも一方が零多項式なら標準AS正則点はない。
片方だけが零となる点は
$\lambda\in\{1,-1,1/z,-1/z\}$に限られ、その始点0の
長さ7の経路次元は少なくとも128である。標準値は120である。
\end{theorem}
\begin{proof}
$P\equiv0$の場合には定数項から$d=za$または$d=a/z$である。
$a=0$なら$d=f_0=g_0=0$で一次係数$m=1$ゆえ矛盾する。
従って$a=1$に規格化できる。$Q\equiv0$の場合には
$c=b$または$c=z^2b$、同様に$k=-1$となる場合を除いて$b=1$とできる。
それぞれ次の二つの指定値を使う。$T=Y_\lambda A_0^{-1} X_\lambda$とする。
\[
\begin{array}{c|c|c}
 &\text{第一の場合}&\text{第二の場合}\\\hline
 P\equiv0:\ (T_{00},T_{11}) &(1,1)/f_0&(1,z^2)/g_0\\
 Q\equiv0:\ (T_{01},T_{10}) &(1,z)/e_0&(z,1)/e_0
\end{array}
\]
これは零にする最高次係数と一次係数を解いたもので、分母は非零である。
$D=ad-bc$とし、取り出した二成分の分子を$N_1,N_2$、
表の指定値を$u_1,u_2$と書く。
$E=N_1-u_1D$, $G=N_2-u_2D$を同時に零とする。

まず四つの例外倍率を除く。$u_2E-u_1G$は未定二変数
$(b,c)$または$(a,d)$の一次式で、後ろの変数の係数は単元である。
それを消去して$E$を取ると、残る変数の二つの一次式の積に単元を掛けた式になる。
例えば$P\equiv0,d=z,f_1=0$のとき、$c$を解いた後の式は
\[
 \frac{(b-\lambda)(\lambda b-1)(\lambda z-1)(\lambda z+1)}
 {\lambda(\lambda-1)(\lambda+1)}
\]
である。どの根でも他方の行列式二次式の三係数が全て消える。
\path{computations/verify_single_zero_pairings.py}は二種類の零多項式、
二つの比、表の二選択の各場合に消去を行い、
除数と二次式の最高次係数が局所化した環の単元であること、および両根での
三係数の消滅を検算する。従ってこれらは前節で全て排除済みである。

例外倍率を一つ固定すると、$u_2E-u_1G$は二変数の一方だけの一次式になる。
その係数の分子・分母の素因子は$z(z^2-1)$の因子に限られ、
この変数を解くと$E,G$はともに恒等的に零になる。残る変数を$u$とする。
例えば上の$P$の例で$\lambda=1$なら$b=1$、$u=c$である。
$Q$が非零で解を持つなら、判別式が零なので$Q$は真の二次式で、
$y=-k/(2q)$である。非零定数では解がなく、一次式なら判別式は非零だからである。
$x=v$は自由であり、四つの残りのSchur等式は恒等式になる。
$Q\equiv0$なら$x=-m/(2p)$, $y=v$として同じ議論を使う。
これで片方だけ零の全許容点は、二種類・二つの比・二選択・四倍率の
計32表示の各々で$(z,u,v)$のアフィン空間の開部分に入る。
この表示では重根の分母が零になる点を単に連続性で埋めず、
もう一式も零になる場合を前節、零にならない場合を上の次数比較で処理した。

同コードで各表示の$\mathbb Q(z,u,v)$上の四窓を作り、始点0で
長さ7までの語商を計算する。各表示で長さ0から6の次元は順に
$1,4,10,20,35,56,84$である。長さ7では128または134となる。
六関係の基底選択に使う点は$(z,u,v)=(3,5,7)$であり、
関数体上の階数計算はこの点の数値を用いない。
既約な表示上で特殊化した次元はこの値以上であるから、
各点で$\binom{10}{3}=120$を超える。
\end{proof}

% !TeX root = ../papers/nc-p3.tex
\section{真の二重根を与えるKummer四次曲面}
\label{sec:kummer-ramification}
$z\lambda(z^2-1)(\lambda^2z-1)(\lambda^2-1)(\lambda^2z^2-1)\ne0$とする。
この節の射影空間は$A_0$の係数$[a:b:c:d]$の空間である。
$D_0=ad-bc$と書き、節\ref{sec:klein-elliptic-completion}の判別式から
\[
 F=\lambda^2(z^2-1)^2D_0^2\,\delta
\]
を作る。これは全射影空間上の斉次四次式である。
\begin{theorem}\label{thm:kummer-ramification}
上の倍率条件を満たす全ての幾何点$(z,\lambda)$で、
$F=0$は幾何学的に整な四次曲面であり、特異点は相異なる十六の通常二重点で尽きる。
この曲面は六点で分岐する種数2曲線のJacobianの符号商と双有理である。
$P,Q$がともに真の二次式で重根を持つ許容窓は
$x=-m/(2p)$, $y=-k/(2q)$から一意に復元される。
この曲面を$(z,\lambda)$とともに動かした既約な族には、全幾何点が標準
Koszul左右AS正則で反対階数16となる非空開集合がある。
これは$(z,\lambda)=(3,2)$の曲面上の開集合も含む。開集合を縮めると
最小周期は4であり、四つの二十点配置それぞれの射影安定群は指定した四元群である。
\end{theorem}
\begin{proof}
$F=Aa^2+Ba+C$と展開し、$[1:0:0:0]$から$[b:c:d]$の平面に射影する。
$A,B,C$は順に二、三、四次である。次の六直線を取る。
\[
 b-c,\quad b-\lambda d,\quad -c+\lambda z d,\quad z^2 b-c,\quad
 \lambda c-d,\quad \lambda z b-d .
\]
$B^2-4AC$はこれらの積に$16(\lambda^2-1)(z-1)^2(\lambda^2z^2-1)$を
掛けたものである。$e=z\lambda+1/\lambda$とすると六直線の係数点は
双対平面の対称行列
\[
 N=\begin{pmatrix}
 2&1+z^2&e\\
 1+z^2&2z^2&ze\\
 e&ze&2z
 \end{pmatrix}
\]
の二次曲線にある。行列式は
$2z(\lambda^2-1)(z-1)^2(\lambda^2z^2-1)/\lambda^2$で非零、
かつ$A$のHesse行列は$-2\lambda^2\operatorname{adj}(N)$である。
従ってこれらは$A=0$への接線である。
各係数点の零座標と残り二座標の比を比べると六点は相異なり、滑らかな
双対二次曲線上の三点は共線でない。従って六直線の二本ずつの交点は十五個で、
いずれも$A=0$の外にある。以上の係数恒等式は
\path{computations/verify_kummer_discriminant.py}で検算する。

$A\ne0$では平方完成により単純分岐直線上で滑らか、二線交点上に各一個の
通常二重点を持つ。六直線の積は平方でないから関数体は平面の二次拡大である。
$A=0$と$B=0$の共通点は六接点に限られ、そこでは
$d(B^2-4AC)=-4C\,dA\ne0$なので$C\ne0$となる。
よって有限な$a$を持つ追加の特異点も、共通因子$\gcd(A,B,C)$もない。
射影中心$[1:0:0:0]$の接錐は非退化な$A$であり、第十六の通常二重点となる。
超曲面の整性は原始な二次式と関数体の既約性から従う。

代数閉体上で滑らかな二次曲線を$\mathbb P^1$と同定し、
六つの接点を分岐点とした二重被覆$C_2\to\mathbb P^1$を取る。
射影平面を$\operatorname{Sym}^2\mathbb P^1$として、二次曲線を二重点の軌跡とすれば、
指定点を含む次数2因子の直線がその接線である。
$C_2:v^2=f(t)$の二点について$v_1v_2$の平方は、両点が指定分岐点を含む
六つの直線式の積に定数と有理関数の平方を掛けたものである。
従って平面の上のこの二次拡大は、$\operatorname{Sym}^2 C_2$を同時の
超楕円反転で割った関数体に等しい。一般の次数2の線束はRiemann--Rochにより
有効因子を一意に持つ。標準因子を引くAbel写像で$\operatorname{Sym}^2 C_2$を
Jacobianと双有理に同定すると、同時反転は符号反転となる。

具体的に$(z,\lambda)=(3,2)$とし、二行列を
\[
 A_0=\begin{pmatrix}36/17&0\\9&4\end{pmatrix},\qquad
 A_1=\begin{pmatrix}36/19&36/37\\36/5&-12\end{pmatrix}
\]
とする。$p,q,x,y,D_0,\det A_1$は全て非零で$\delta=0$である。
両隣接切断は階数6、反対切断は階数16である。
四次中心は文字次数順に$(3,2,0,0)$次元で、既定の四中心切断で割った商は
素数1009において全始点・経路長13まで$(1+t+t^2+t^3)^4$の係数を与える。
既に正則な一般点を持つ五次元の二次補完族上にあり、次数ごとの標準Hilbert関数の
下界とこの上界から、定理\ref{thm:klein-rational-period-four}の
自由表示および左右Koszul列の議論により全次数で正則となる。

開集合についても理由を述べる。$pqD_0\ne0$における曲面族は、
上の幾何学的に整な四次の一般ファイバーを持ち、四次式に基底だけの因子がなく
整である。標準次元を持つ経路長4,5の近傍で、周囲の十五次元成分の一般点の
四次中心五次元は特殊化先でも下界となる。上の標本も五次元なので、
近傍で可換方程式は定階数となり、四切断を延ばせる。
次数13までの商次元の上界と反対階数16を保つ開条件を課せば、
同じ上下界と自由表示の証明がこの一つの開集合の全幾何点に適用できる。
また素数34381でこの有理標本の四位置の二十点は完全分裂し、四回の核の移動は恒等、
配置間の射影写像の個数行列は$4I_4$である。各配置の三点は非共線で五点は非共面であり、
系\ref{cor:klein-four-configurations}の射影枠の特殊化を適用できる。
四元群はもとから作用するので、開近傍で四つの安定群が確定し、位置の非同値性から
最小周期4が従う。
\end{proof}

% !TeX root = ../papers/nc-p3.tex
\section{Kummer曲面の極写像と十六の二重接触平面}
\label{sec:kummer-polar-completion}
前節の通常倍率を保ち、$A_0,A_1$を成分順$(00,01,10,11)$で列ベクトルと読む。
$u=z(\lambda^2+1)$, $v=\lambda(z+1)$, $w=\lambda^2z^2+1$として
\[
 T=\begin{pmatrix}
 u&v&zv&w\\v&\lambda^2+1&w&v\\zv&w&zu&zv\\w&v&zv&u
 \end{pmatrix}.
\]
\begin{theorem}\label{thm:kummer-polar-completion}
$F=0$の有理写像として$[A_1]=[T\nabla F(A_0)]$である。
この写像は対合で、窓の二位置移動を文字ラベルと対角倍率の変更まで実現する。
従って$T$は曲面の射影双対を再び同じ曲面に同定する。
十六の二重接触平面の各切断は滑らかな二次曲線を二重に数えたものである。
これらの上の許容窓は全て非標準ASである。
\end{theorem}
\begin{proof}
$pq\ne0$で$x=-m/(2p),y=-k/(2q)$を代入すると、完成行列の四成分は
\[
 \left(\frac{1+e_0(h_1-e_1)}{2h_1}, \frac{f_0g_1-g_0f_1-1}{2g_1}, \frac{g_0f_1-f_0g_1-1}{2f_1}, \frac{1+e_0(e_1-h_1)}{2e_1}\right)
\]
である。$\det T=(z^2-1)^2(\lambda^2z-1)^4$。
\path{computations/verify_kummer_polar.py}は各成分の分母を払い、
上の列と$T\nabla F$の六つの二次小行列が全て
$\mathbb Q[a,b,c,d,z,\lambda]$で$F$で割り切れることを検算する。
従って滑らかな点で両方の表示が定まる限り射影値は等しい。

$A_0,A_1$を交換すると$C,D$も交換され、両者のSchur補を対角と非対角に
分けた四つの条件がそのまま保たれる。テンソルについて具体的に、
元の成分$w_{k,l\mathbin\oplus1,i,j\mathbin\oplus1}$へ
$q^{I+J-K-L}$を掛ける。ここで$q^2=\lambda$、大文字は各添字が
$\{2,3\}$に属すとき1、他で0である。薄い色では指数が零で二つの薄い係数の積が
元の規格化に戻る。厚い色では$I=J,K=L$で、左上・右下は不変、
右上は$1/\lambda$、左下は$\lambda$倍となる。
従って得られるテンソルは同じ$(z,\lambda)$で$A_0,A_1$を交換したものに等しい。
この関係は二重根以外にも成立する。これは二位置移動を伴う関係であり、
移動を忘れた添字保持同型を自動的には与えない。

同コードで$F(T\nabla F)$も$F$で割り切れることを確認する。
どの許容倍率でもGauss像は二次元である。実際、通常二重点での接平面の極限は、
その非退化な二次接錐の各母線での接平面を全て含む。これは節点を一度吹き上げ、
例外二次曲線で勾配の先頭項を読むことで分かる。一節点につき双対平面内に
滑らかな二次曲線が一つあり、異なる節点ではその平面が異なる。
既約なGauss像の閉包はこれら二曲線を含むため次元は二である。
よって$T$を施した像は既約な$F=0$に稠密であり、両行列とも$pqD_0\ne0$の
一般点で交換できる。重根による完成が一意なので、再度の交換で元の行列に戻る。
この同定は射影双対性の双対定理と$T$の対称性からも従う。

十六平面を列挙する。四平面は
$b=c,\ c=z^2 b,\ d=za,\ a=zd$であり、
四平面は$f_1,g_1,e_1,h_1$の分子を零とするものである。残る八平面は
\[
 b=\lambda d,\quad d=\lambda z b,\quad c=\lambda z d,\quad d=\lambda c,
 \qquad
 a=\lambda z b,\quad b=\lambda a,\quad a=\lambda c,\quad c=\lambda z a .
\]
各法線を$L$と書く。同コードはチャートで単元の係数を一つ選んでその座標を消去し、
$F$の制限が単元倍の二次式の平方で、二次式のHesse行列式が単元であることを
因数分解で調べる。それぞれの二次曲線上の滑らかな点の接平面は$L$である。
$TL$は元の曲面の節点である。八つの交差成分平面では、$TL$を
$a=1$に規格化した値は定理\ref{thm:klein-all-pairings}の表の八つの$A_0$と一致する。
$d\in\{z,1/z\}$、$c/b\in\{1,z^2\}$、各組に二つの相異なる$b$があり、
$\lambda^2z\ne1$によりいずれも可逆である。他の八つの$TL$は四座標頂点と
\[
 (1,\lambda,1/\lambda,1),\quad(1,1/(\lambda z),\lambda z,1),\quad
 (1,\lambda,\lambda z,\lambda^2z),\quad
 (1,1/(\lambda z),1/\lambda,1/(\lambda^2z))
\]
で、全て$D_0=0$である。これらも通常倍率で相異なる。
これで前節の全十六節点と十六平面が対をなす。
ほかに平面切断が二重の滑らかな二次曲線となる平面はない。
実際、Gauss対合がその曲線を一点に縮めれば、逆写像はその点で
正則に定まらず、$T$が可逆なのでそこは節点である。

最初の四平面では$n=0$または$o=0$であり、真の二重根なら$x=0$または$y=0$に
なって可逆チャートを外れる。次の四平面では$pq=0$で、重判別式と有限な解を
仮定すると対応する二次式は恒等的に零になる。これらは前二節で処理した。
最後の八平面では$A_1$を新たな基底とすれば$P\equiv Q\equiv0$である。
実際、平面接触により一般点の$[A_1]$は上の特定の$[TL]$であり、
$pq\ne0$で有限な重根表示は正則なので節点への特殊化にも同じ値が保たれる。
従って交換した窓を定理\ref{thm:all-zero-pairing-defect}で除外できる。
二位置移動と基底変更は正則性を保つから元の窓も非正則である。
\end{proof}

% !TeX root = ../papers/nc-p3.tex
\section{符号付き判別被覆上の対根と節点の小解消}
\label{sec:paired-root-resolution}
通常倍率を一組固定し、重み$(1,1,1,1,2)$の射影空間に
$X:\ S^2=F(a,b,c,d)$を取る。$D_0\ne0$では
$s=S/(\lambda(z^2-1)D_0)$として$s^2=\delta$である。
共通倍率を局所的に一つの非零成分で固定し、根$x,y$を各々射影直線に完備化する。
規格化を変えると根は共通倍率の二乗で変換されるので、これらの射影直線束は貼り合う。
\begin{theorem}\label{thm:paired-root-small-resolution}
$X$の特異点は十六の通常三次元二重点である。$D_0\ne0$の部分で、
符号$s,-s$に対応する二つの根のグラフ閉包は射影的小解消を与える。
八つの節点を除けばこれは同型で、各例外ファイバーは正規束
$\mathcal O_{\mathbb P^1}(-1)^{\oplus2}$を持つ一本の射影直線である。
その第二行列$[A_1]$は対応する接触二次曲線の全体を一対一にパラメータ付ける。
被覆の符号反転は各節点で二つの小解消を交換する。
\end{theorem}
\begin{proof}
重み付き射影空間の頂点は$X$に属さない。分岐面の滑らかな点では$X$も滑らかで、
面の節点の非退化二次項に$S^2$を加えると三次元通常二重点となる。
既約な分岐四次の二重被覆で、特異点は孤立するから$X$は整かつ正規である。

$[U:V]$を$x$の斉次根とし、
\[
 M(s)=\begin{pmatrix}2p&m-s\\m+s&2n\end{pmatrix},
 \qquad M(s)\binom UV=0
\]
を取る。$\det M=s^2-\delta$で、$V\ne0$なら一行目が$2px+m=s$、
二行を合わせて$px^2+mx+n=0$となる。非零行列では核直線は一意に正則に定まる。
無限遠根や最高次係数の零点でも同じである。零行列の点は$P\equiv0$であり、
通常倍率では定理\ref{thm:klein-all-pairings}と前の片側零の排除により八節点で尽きる。
そこで$p,m,n$の微分は独立である。実際、$\delta=m^2-4pn$の二次項が非退化で、
この三つの一次項が従属ならその階数は三に達しない。
従ってこの三つはエタール局所座標であり、$s$を加えた四座標から
$M$の四成分への変換は線型同型になる。

階数一以下の二行二列行列の核直線を指定した空間は、射影直線上で各行が
指定直線を消す共変ベクトルの対、すなわち$\mathcal O(-1)^{\oplus2}$の全空間である。
ゆえにこの射影的な核の表示は滑らかで、小解消となる。
$Q$について$-s$を使っても同様である。両根の被約グラフ閉包から第一の小解消への
射影は固有かつ双有理なので全射である。節点上では全ての点が既知の
$H=xy+Bx+Cy+D=0$の双斉次化を満たす。$BC-D$は八表示の全てで単元である。
実際、分子分母の非定数因子は
$z,\lambda,z\pm1,\lambda^2z-1,\lambda\pm1,\lambda z\pm1$
に限られる。よってその曲線は射影直線の自己同型のグラフであり、どの第一根にも
第二根は高々一つである。射影は準有限かつ固有なので有限、正規な第一の小解消への
有限双有理射なので同型となる。第二根への射影も同様である。

\path{computations/verify_paired_root_resolution.py}は各節点で二組のJacobian行列式と
$BC-D$が単元であることも検算する。節点の列を$N$とすると、
$H$から得た$A_1(x)$は$F(A_1)=0$と$(T^{-1}N)^t A_1=0$を満たす。
四成分の共通分母を払った二次以下の四式の$1,x,x^2$係数行列は、
通常倍率で単元となる三次小行列を持つ。よって無限遠にも基点がなく、
射影化した表示は次数二の完全線型系である。その像が上記の接触二次曲線である。
この完備化のうち元の可逆窓に戻る点は全て長さ6の障害を持つ。

最後に$M(-s)=M(s)^t$である。従って符号反転は核直線と左核直線を交換する。
固定した右核を持つ階数一行列をスカラー倍で零へ近づけると、
左核は任意に選べる。よって符号反転は同じ小解消の例外直線を越える正則な自己写像にはならず、
二つの小解消の間の通常のフロップを与える。
\end{proof}

% !TeX root = ../papers/nc-p3.tex
\section{重根の二位置交換の固定点とEnriques商}
\label{sec:kummer-exchange-quotient}
\begin{theorem}\label{thm:kummer-exchange-enriques}
通常倍率の各幾何点で、規格化した二行列の交換は判別Kummer曲面の最小解消上の
固定点のない対合となる。従ってこの交換による解消曲面の商はEnriques曲面である。
これは印と倍率を固定した表示の二位置交換に関する商である。
\end{theorem}
\begin{proof}
通常倍率の連結な基底を$B$とする。曲面族の相対勾配の基点は互いに交わらない
十六の節点切断である。節点切断に沿って相対Hesse行列が非退化なので、アフィン
座標の一次偏微分は切断の余接空間の基底を与え、Nakayamaの補題で切断のイデアルを生成する。
従って極写像のグラフ閉包$Y$はこれら十六切断の同時吹上げである。
各例外曲線は滑らかな接錐二次曲線である。吹上げの各アフィン座標では、
二次方程式の非特異な射影接錐により$Y\to B$は例外上でも滑らかである。
これは固有な曲面族で、各ファイバーは通常二重点の最小解消である。
極写像が対合なので、二つの射影因子の交換は$Y$の正則な自己同型$\tau$となる。

$N=[u]$においてグラフに固定点があれば
\[
 F(u)=0,\qquad u^t\operatorname{adj}(T)u=0,\qquad
 u\wedge T\nabla F(u)=0 .
\]
滑らかな点ではEuler公式$u^t\nabla F=4F$を使う。
節点では極写像の極限は$u$を通る平面なので、グラフ上の第二座標は
$u^tT^{-1}v=0$を満たす。従って必要条件は節点にも及ぶ。
$(z,\lambda)=(3,2)$の整数係数でこの斉次閉集合を素数1009へ還元する。
先頭の非零座標を1とする四つの部分で、最初の三つではGröbner基底が$(1)$、
最後では二次式が非零定数である。
\path{computations/verify_kummer_exchange_quotient.py}がこれらを検算する。
射影閉集合の固有性により標数零のこの有理ファイバーにも幾何的な解はない。

固定点スキーム$Y^\tau$は$B$上固有かつ滑らかである。後者を詳しく見ると、
平方零のイデアルを持つ試験環への固定点の持ち上げについて、$Y/B$の滑らかさで
まず一つの持ち上げを取れる。持ち上げ全体は接ベクトルを係数とするアフィン空間で、
対合の作用はアフィン線型である。持ち上げとその対合像の中点を取れば固定した
持ち上げを得る。これは形式的滑らかさを示し、有限表示なので滑らかさとなる。
従って像は開かつ閉であり、連結基底の一点で空だから全体で空である。

解消四次曲面はK3曲面である。実際、四次超曲面の標準束は自明で通常二重点は
クレパントな有理特異点であり、四次の完全列から$H^1(\mathcal O)=0$である。
自由な対合による商は滑らかで射影的、被覆は次数2のエタール射である。
従って商でも$H^1(\mathcal O)=0$、$\chi(\mathcal O)=1$で、標準束の二乗は
ノルムにより自明である。Serre双対性から標準束自身は切断を持たず、自明でない。
これはEnriques曲面の条件を満たす。他の文字置換も許した添字保持同型の判定には、
別途四位置の作用を考慮する必要がある。
\end{proof}

% !TeX root = ../papers/nc-p3.tex
\section{四つの例外倍率の重根上の全次数正則性}
\label{sec:exceptional-ramification}
\begin{theorem}\label{thm:exceptional-ramification-as}
$z\ne0,\pm1$, $D_0\ne0$として$\lambda=\pm1,\pm1/z$を取る。
各判別四次は滑らかな二次曲面を二重に数えたものである。この二次曲面上の
可逆な窓が標準Koszul左右AS正則となる必要十分条件は、$P,Q$がともに真の二次式
であることである。その場合の完成は一意である。
\end{theorem}
\begin{proof}
$G=(a+d)(zb+c)-(z+1)(ad+bc)$と置く。正の二倍率$1,1/z$での$F$は
それぞれ$(z-1)^2G^2,(z-1)^2G^2/z^2$で、負の倍率では$b,c$の符号を変える。
$G$のHesse行列式は$(z^2-1)^2$である。真の重根の完成は
定理\ref{thm:klein-other-pairings}の表による。$pq=0$の有限解では、
その側の線型係数も判別式の零から消え、定数項も零になるので、
定理\ref{thm:all-zero-pairing-defect}, \ref{thm:single-zero-pairings}で除かれる。

まず$\lambda=1$とし、$A=A_0=A_1$とする。実際には$G=0$より広く、
次の八差が全て非零な全ての$A$で標準正則性を示す。
\[
 a-b,\quad a-c,\quad a-zb,\quad za-c,\quad
 b-d,\quad c-d,\quad zb-d,\quad c-zd .
\]
$G=0$では$x=y=D_0$なので真の重根なら$n=pD_0^2$と$o=qD_0^2$は非零である。
各差を零にすると、$G$の制限は$(a-zd)(za-d)(b-c)(z^2b-c)$の制限を割る。
例えば$a=b$では$G=(b-c)(zb-d)$である。従ってこの八差は全て非零となる。

両対称表示で反対切断を文字次数に分ける。四行二列、二行四列の二ブロックと、
行列式$z-1,1-z$の二行二列ブロック二つを得る。
長方形ブロックを両方とも四行二列に向けて行対を$01,02,03,12,13,23$の順に取る。
最初と最後の小行列は単元倍で順に$(a-d,zb-c)$と$(zb-c,a-d)$である。
中の四つはすべて上の八差のうち二つの積を$D_0$で割ったものとなる。
従って反対階数は8で、左右の二次核は四つの平方だけからなる二次鉛筆である。
行対の小行列は核二平面の相補小行列に対応する。
$a-d,zb-c$のいずれも零でなければ、各鉛筆の四つの列方向は相異なり、両曲線は滑らかである。
片方だけが零なら、両曲線とも横断する二平面二次曲線の環である。
実際、一対の列が比例するのでその二列を消す一式は残る二変数の異なる二直線因子に分かれ、
他方の式は各平面に非退化な対角二次曲線を切る。二対とも比例すれば、
独立な二組それぞれの二変数二次式が分解し、両曲線とも四直線の環となる。
これらの曲線の埋込み次数はそれぞれ$(2,2)$または$(1,1,1,1)$である。

六次元の関係空間に純テンソルがないことを確かめる。
文字を低い二文字と高い二文字に分けると、関係行列の右上と左下は
\[
 \begin{pmatrix}r&t\\t&r\end{pmatrix},\qquad
 \begin{pmatrix}r&t\\t&zr\end{pmatrix}.
\]
両行列式が零なら$z\ne1$より$r=t=0$である。残るブロック対角行列が階数一なら
対角ブロックの片方が零であるが、各厚い文字次数の関係二平面は低い二座標にも
高い二座標にも同型に射影するので、全体が零である。
従って滑らかな場合には定理\ref{thm:smooth-apolar-complete}が適用できる。

多角形の場合に必要な両側評価の階数を示す。零点対の低い二座標を$u,s$、
高い二座標を$v,t$とする。厚い四式は
\[
 \binom{u_0s_0}{u_1s_1}=L\binom{v_0t_0}{v_1t_1},\qquad
 \binom{u_0s_1}{u_1s_0}=K\binom{v_0t_1}{v_1t_0},\qquad
 L=-A^{-t}Y^t,\quad K=-A^{-1}X
\]
で、$X=\left(\begin{smallmatrix}1&1\\1&z\end{smallmatrix}\right)$,
$Y=\left(\begin{smallmatrix}1&1\\z&1\end{smallmatrix}\right)$である。
全八成分が非零で、$\det K=-\det L=(z-1)/D_0\ne0$である。
$(u,v)$を固定した評価の核が二次元以上と仮定する。
$u$または$v$が零なら厚い式が対応する二座標を消し、薄い式がさらに一式を課すので不可能である。
どちらも非零なら、核から$s$への射影も$t$への射影も単射である。
従って$t$は自由に動き、$s$も二次元を張る。上の四式の成分は全て非零なので、
$v$の一座標が零なら、例えば$v_1=0$で
$(u_0s_0)(u_1s_1)-(u_0s_1)(u_1s_0)$の右辺の$t_0^2$係数が非零となり矛盾する。
同様に全ての$u_i,v_i$が非零となる。
$t_0,t_1$の係数を比較し、$u_1/u_0$を二通りに求めると
$K_{11}K_{00}=L_{00}L_{11}$かつ$K_{10}K_{01}=L_{01}L_{10}$となる。
これは行列式の符号に反する。右評価では$K$の行と列を同時に交換するだけなので同様である。

滑らかな楕円鉛筆からこの窓への特殊化を取れる。実際、八差を非零に保った係数空間は
既約で、一般点の二鉛筆は滑らかであり、上でその正則表示を得た。
多角形への特殊化で、定理\ref{thm:two-polygon-specialization}の証明の共通の環を取る。
多重次数が異なる174型中172型では、同証明が実際の関係の片側評価に
二次元以上の核を与える。残る二型はどちらも次数3の成分を必要とし、
ここで得た埋込み次数に存在しない。従って両束の多重次数は等しく、同証明により
八次元の積の上の二次元鉛筆が得られる。基点は純テンソルを与えるため存在しない。
多角形鉛筆の全次数定理で正則となる。

$\lambda=1/z$の表の窓は$A'=\left(\begin{smallmatrix}zb&a\\zd&c\end{smallmatrix}\right)$を持つ
$\lambda=1$の窓と同型である。具体的に第三・第四位置の文字を各々1だけ平行移動し、
成分に$(1,1,1/z,1)_k(z,1,1,1)_l$を掛ける。$D'_0=-zD_0$,
$G(A')=zG(A)$で、$no$も$z^2$倍となる。負の二倍率は既証の可逆な
コサイクル捻りで正の二倍率へ戻り、$b,c$がともに反転する。
\path{computations/verify_exceptional_ramification.py}が四つの判別四次、上の小行列、
八差への制限、正倍率の基底変更の全256成分を検算する。
\end{proof}

% !TeX root = ../papers/nc-p3.tex
\section{例外交換表示の非正則な八直線}
\label{sec:exceptional-exchange-sheets}
\begin{lemma}\label{lem:bisymmetric-reducible-quick-test}
隣接階数6の両対称反転窓で、いずれかの反対核が二つの一次式の積を含むとする。
標準Koszul左右AS正則となる必要十分条件は、関係空間に純テンソルがなく、
両側の全ての非零評価が階数3以上であることである。
\end{lemma}
\begin{proof}
必要性は反転窓の二つの成長障害による。平方因子の場合、軸と相手の評価階数は
3であり、定理\ref{thm:double-plane-ore}の三変数完全交叉の条件に反すれば
極形式系に階数1が現れ、全窓に純テンソルを与える。
相異なる二因子なら補題\ref{lem:reciprocal-apolar-factors}で四軸の階数は3となる。
二つの直交三平面の交わりが二次元以上なら、二行二列関係の二平面に
行列式零の非零元があり、これも純テンソルとなる。
横断する場合の$\kappa_U=0$、および根基直線の場合の$\det t=0$も、
それぞれ純テンソルを与える。従って定理\ref{thm:transverse-factor-amalgam},
\ref{thm:radical-factor-amalgam}の十分条件が成立する。
\end{proof}

\begin{theorem}\label{thm:exceptional-exchange-eight-lines}
$\lambda=1,A_1=A_0$の全交換表示で、$zD_0(z^2-1)\ne0$とする。
非標準ASとなる$[a:b:c:d]$は下の八直線の和集合との交わりに等しい。
\[
\begin{array}{c|c|c|c}
 &\text{非零座標の位置}&\text{直線の式}& r^2\\\hline
\text{左}&02&a=b,\ c=b&b\\
 &03&a=zb,\ c=z^2b&b\\
 &12&c=b,\ d=b&b\\
 &13&d=zb,\ c=z^2b&b\\\hline
\text{右}&02&a=b,\ d=zb&b\\
 &03&a=zb,\ d=b&b\\
 &12&c=za,\ d=za&za\\
 &13&a=zd,\ c=zd&zd
\end{array}
\]
四つの例外交換表示には同じ判定が前節の変換で適用される。各表示で八直線は
全て重根二次曲面上にあり、特にその表示の単根窓は全て正則である。
\end{theorem}
\begin{proof}
純テンソルがないことは前節の角ブロックの証明で係数差の非零を必要としない。
反対切断の二つの長方形ブロックの各行（転置した側では各列）は零でない。
実際、前節の四行二列表示の第一側は各行に$1$または$z$があり、第二側でも同様である。
いずれかで二行が比例すれば、その二文字の平方の非自明な関係式を生じ、
二つの異なる一次因子を持つ。補題\ref{lem:bisymmetric-reducible-quick-test}が使える。
全ての行対が独立なら両対角鉛筆は滑らかで定理\ref{thm:smooth-apolar-complete}による。
従ってあとは評価核が二次元以上になる条件を決定すればよい。

左評価を前節の$u,v,s,t,L,K$で書く。一方の半分$u,v$が零の場合の排除、
および両半分が非零の場合の核から$s,t$への射影の単射性は二行列の可逆性だけで成立する。
$u_0u_1v_0v_1\ne0$とする。$c_0=u_1/u_0$, $v'=v_1/v_0$と置くと
係数比較から
\[
 c_0 L_{00}=K_{11}v',\quad c_0L_{01}v'=K_{10},\quad
 c_0K_{01}v'=L_{10},\quad c_0K_{00}=L_{11}v' .
\]
これらを掛け合わせて整理すれば、行列の成分が零でも前節と同じ行列式の矛盾になる。
$v_0v_1=0$で$u_0u_1\ne0$なら、例えば$v_1=0$で第一と第三の積式の
$t_0$係数、第ニと第四の積式の$t_0$係数を比較すると
$L_{00}=L_{10}=0$になる。$v_0=0$も同様である。
逆に$u$の一座標が零なら厚い式の対応する二行目または一行目により、
$v$の一座標が零である。右評価は$K$の行と列を同時に交換して同じ結論を与える。
従って評価核の次元超過は、低い一座標を1、高い一座標を$r\ne0$とする
表の八場合に限られる。

この座標形では薄い二式を使い相手の低い二座標を高い二座標から一意に解ける。
残る厚い四式が恒等的に消える条件を取る。例えば左02では
$r^2(a-b)=r^2(a-c)=0$と
$D_0=r^2(d-c)=r^2(d-b)$を得る。$D_0\ne0$なので表の通りとなる。
\path{computations/verify_exceptional_exchange_sheets.py}は八場合全てで、
$r^2$だけを因子とする係数から独立な二つの直線条件を取り出し、
残りが$D_0\ne0$の下で表の$r^2$を必要十分とすることを検算する。
右辺は該当直線の可逆部分で非零であり、幾何点では平方根が二つある。
各直線が前節の八差の平面のうち異なる二枚に含まれ、$G$がその上で消えることも
表に代入すれば分かる。$z\ne0,\pm1$ではこれらの直線は相異なる。
\end{proof}

% !TeX root = ../papers/nc-p3.tex
\section{両対称反転窓の二つの経路長による必要十分判定}
\label{sec:bisymmetric-finite-path-test}
$R\subset V\otimes W$を隣接階数6の両対称Gramテンソルの関係六平面とし、
反転窓の二次代数を$C=C(R,R^t,R,R^t)$とする。以下の判定は
反対階数、中心、反対鉛筆の曲線型の事前判定を必要としない。
$P_d$を二変数$d$次形式の空間と書く。
\begin{theorem}\label{thm:bisymmetric-two-length-test}
次は同値である。
\begin{enumerate}
\item $C$は標準14641型Koszul左右AS正則である。
\item いずれの始点でも長さ3の経路次元が20以下、長さ21の経路次元が2024以下である。
\item $R$は非零の純テンソルを含まず、どちらの非零片側評価も階数3以上であり、
次に定める二変数形式への商がどちらの向きにも存在しない。
\end{enumerate}
ここで二変数形式への商とは、矢空間の全射$V\to P_2^*$, $W\to P_3^*$で
$R$の像が乗法$P_2\otimes P_3\to P_5$の双対像に入るものをいう。
最初の二つの線型代数条件（純関係なし、評価階数3以上）を直接確認した場合は、
経路長17で両始点とも1140以下であることだけで全次数の正則性を判定できる。
特に両対称の隣接階数6の空間で標準正則窓はZariski開集合をなす。
\end{theorem}
\begin{proof}
まず第三条件の十分性を示す。反対核に可約な元があれば
補題\ref{lem:bisymmetric-reducible-quick-test}による。以下、どちらの核にも
そのような元がないとする。定理\ref{thm:all-bisymmetric-as-closure}の証明で、
楕円鉛筆閉包に入るために使う条件は純関係の不在と両軸での評価階数だけであった。
従ってこの窓を、一般点で相異なる二束を持つ滑らかな楕円鉛筆から特殊化できる。
二側の定義二次式の二平面をGrassmann多様体で極限に取る。それぞれ実際の
反対核に含まれるので、共通の一次因子も階数2以下の元もない。
従って両方とも四次完全交叉であり、非被約な鉛筆には可約な二次式があるという
定理\ref{thm:nonreduced-quartic-pencils}の議論から被約である。
反対階数が特殊点で一定であると仮定する必要はない。

被約極限について、既出の特殊化証明がどの障害で場合を除くかを確認する。
滑らかな側があれば、評価階数により切断格子が揃い、純関係がないので基点自由な
鉛筆が延びる。整な節点・尖点の場合にも定理\ref{thm:integral-apolar-alternatives}
の格子のずれは実際の評価階数2以下を与える。格子が揃った後、可逆束なら
基点自由な整曲線の表示になる。非可逆なら導手で消える三切断と他側の全四切断の
積が$P_2\otimes P_3\to P_5$を経由し、禁止された商を与える。

両側が節点多角形の場合、定理\ref{thm:two-polygon-specialization}の証明で
多重次数の異なる172型は、極限の積の階数の損失を含めても実際の評価核が
二次元以上になる。残る二型には命題\ref{prop:semipositive-pencil-balance}の
二変数形式の商がある。多角形と楕円多重点を組み合わせる場合は
命題\ref{prop:polygon-additive-exclusion}の証明が、実際の評価核に二次元を与える。
最後に両側が多重点の場合、定理\ref{thm:all-reduced-apolar-specializations}の
証明では、中心対のない場合は積の階数が8を超えるか実際の評価階数が2以下となる。
中心対がある場合も、積の階数7の不一致および次数の大きい末端は同じ評価障害で落ちる。
残る共通多重点模型では積が八次元を張り、命題\ref{prop:semipositive-pencil-balance}
によって、次数不一致は評価障害か二変数形式の商のいずれかを与える。
これらは上記各証明で正則性の必要条件として使われた具体的な排除であり、
ここでは第三条件が直接これらを禁じている。残った等次数の模型は積の二次元鉛筆を
持ち、基点は純関係を与えるため存在しない。各被約曲線の全次数定理が適用できる。

第二条件から第三条件を導く。非零評価の核が二次元以上なら、一文字とその核の
二文字へ交互に射影して全関係を消せる。二文字始点の長さ21では
$2^{11}=2048>2024$となる。従って両評価は階数3以上である。
特に各一文字での最上テンソルの縮約は単射で、三次交叉に四次元の像を持つ。
長さ3の次元は$64-2\cdot24+\dim W_{i,3}$なので、この像は交叉に等しい。
もし$a\otimes b\in R$が純関係なら、転置により$a\otimes b\otimes a$が交叉に入る。
従ってこれはある非零文字による最上テンソルの縮約である。
その先頭・末尾の一次元支持は、その文字で隣接関係を評価した像に等しい
（定理\ref{thm:no-consecutive-pure-relations}の縮約の議論）。
評価階数が1となり矛盾する。

二変数形式の商があれば反転先でも転置により同じ商が続く。
補題\ref{lem:binary-diagonal-paths}の経路商の$P_3^*$始点で長さ$2m+1$の次元は
$9\cdot2^{m-1}$である。長さ21では4608となり第二条件に反する。
長さ17でも$1152>1140$なので追加の判定も従う。
第一条件なら標準Hilbert級数から第二条件を得る。最後に経路関係の挿入行列の
階数下界は、六次元関係束の各基底チャートで開条件である。
\end{proof}

二変数形式の商は曲線の同型を探索せずに検査できる。三行四列の独立行列$U$を取り、
関係の各係数行列$A$について$UAY$がHankel行列となる可逆な四行四列行列$Y$を探す。
六つの反対角線が一定となることは、相隣る二行のずれから各関係につき六つの
線型方程式を与える。固定した$U$に対し合計36行16列の核を取り、その線型系で
行列式が恒等零かどうかで可逆解を判定できる。$U$については三次小行列が一つ
非零となる四つのチャートで解の存在を調べればよく、左右を交換した検査も行う。
\path{computations/verify_bisymmetric_finite_test.py}は節点・尖点の二標本で
座標を取り替え、指定した三平面上のこの核から可逆な解をスカラーまで復元する。

% !TeX root = ../papers/nc-p3.tex
\section{階数二曲線からの二変数商の検出と四辺での整合性}
\label{sec:binary-window-transfer}
$R\subset V\otimes W$を純テンソルを含まない六次元関係空間とし、
$V,W$は四次元とする。直線$D\subset V$で第一矢を割った像を
$H_D\subset(V/D)\otimes W$とする。
\begin{proposition}\label{prop:binary-quotient-rank-two-curve}
この向きで$D$を核とする二変数形式の商があるための必要十分条件は、
$\mathbb P(H_D)$内の階数一行列の軌跡に、両射影がそれぞれ
$\mathbb P(V/D),\mathbb P(W)$を線型に張る整曲線が存在することである。
このとき階数一軌跡はスキームとして有理正規五次曲線そのものであり、
二射影は滑らかな二次曲線およびねじれ三次曲線への同型である。
その逆像は$\mathbb P(R)$の階数二行列の有理正規五次曲線で、
全ての列像が$D$を含む。この曲線から$D$と二つの射影モデルが一意に戻る。
\end{proposition}
\begin{proof}
射影の核は$D\otimes W$に入り、その非零元は純テンソルなので、
$R\to H_D$は同型である。二変数形式の乗法の双対像は
$(h_{i+j})_{0\le i\le2,0\le j\le3}$の六平面である。
階数一の小行列は有理正規五次曲線のイデアルを生成する。実際、二次部分は
隣接する指数を一つずらす十本の標準二項式の空間と一致する。
こうしたずらしで同次数・同指数和の単項式を互いに移せるので、標準二項式が
パラメータ表示の核全体を生成する。
\path{computations/verify_binary_window_transfer.py}で両二次部分も比較する。

逆に指定した整曲線の関数体で、各射影の非零な座標で割り、
二組の座標関数の空間を$A,B$とする。次元は3と4で、$AB$は線型包から六次元以下である。
一般に曲線の関数体の有限次元部分空間は、滑らかな点での付値が相異なる基底を取れる。
二つの付値集合の和を取ると$\dim AB\ge\dim A+\dim B-1$で、ここでは等号である。
一般の点で消える$A$の二平面を$a\langle1,t\rangle$と書くと、$t$は非定数で、
その$B$との積は一点で消えるから五次元以下である。よって$\dim(B+tB)=5$。
$d$次元の$U$が$\dim(U+tU)=d+1$を満たせば、
$U_j=\{u:t^i u\in U,\ 0\le i\le j\}$に対し$U_{j+1}$は
$u\mapsto t^ju\bmod U_1$の核なので、各段で失う次元は高々1である。
$0\ne b\in U_{d-1}$から$U=b\langle1,t,\ldots,t^{d-1}\rangle$を得る。
従って$B$はこの形である。また$A+tA+t^2A+t^3A$は六次元で、
非定数$t$の乗法が保つ非零有限次元空間は存在しないので、各段で丁度一次元ずつ増える。
よって$A$も同じ$t$の三項の形である。曲線は二つのべき表示のグラフとなり、
その包は六平面全体である。これで必要十分性が従う。
持ち上げた行列の階数は高々2で純テンソルはないので丁度2である。
列像の射影は非退化な二次曲線を動くから、全ての列像の共通部分は丁度$D$となる。
具体的には階数二行列の曲線上で一つの共通直線を見つけ、その商の両射影の
線型包を調べれば、同じ判定を行える。
\end{proof}

\begin{corollary}\label{cor:four-window-binary-transfer}
反転を仮定しない四周期の窓で、交互の二位置で矢空間を一次元の核で割る。
各辺の向きを三次元側から見るとき前命題の条件が成立し、各頂点で左右の辺の
与える射影曲線が同一とする。このとき四窓に二変数形式の経路代数への
添字付き全射があり、長さ17には少なくとも1152次元の始点がある。
この整合性は、固定した四周期の核と埋込み曲線を用いた全射の存在と同値である。
\end{corollary}
\begin{proof}
各辺の階数一曲線の二射影は同型なので、隣接する頂点曲線間の同型を一つ与える。
一つの頂点で$\mathbb P^1$の座標を選び、この同型をたどって全ての整数位置へ
座標を運ぶ。二次または三次の完全線型系によって各座標変更は矢の線型同型に
スカラーを除き持ち上がる。スカラーの選択は関係空間を変えない。
全ての辺の関係が二変数の乗法の双対像に写るので添字付き商を得る。
逆は各辺のHankel六平面の一意な階数一曲線を読めばよい。
一周して得る$\operatorname{PGL}_2$の元は一般に恒等でなくてよい。
添字を一周ごとに区別して座標を選べるためである。四周期の座標選択まで要求する場合に限り、
射影的な一周の変換が恒等であることが必要十分となる。
\end{proof}

% !TeX root = ../papers/nc-p3.tex
\section{一辺の二変数商の余次元と残余六次曲線}
\label{sec:binary-incidence-sextic}
\begin{theorem}\label{thm:binary-incidence-codimension}
純関係を含まない$\operatorname{Gr}(6,16)$の開部分で、第一矢を三次元に割る
二変数形式の商の軌跡の閉包は既約で余次元13である。その一般点で核直線は一意である。
関係五平面内の階数二以下の軌跡は、被約支持として次数5の有理正規曲線と
次数15の曲線の和で、後者の正規化は滑らかな平面六次曲線となる。
\end{theorem}
\begin{proof}
核直線の選択は三次元、その商におけるHankel六平面の選択は
$\operatorname{PGL}_3\times\operatorname{PGL}_4$の軌道である。
階数一曲線の二射影が完全線型系で同じ$\mathbb P^1$を埋め込むため、
安定群は対角的な$\operatorname{PGL}_2$であり、軌道次元は$8+15-3=20$である。
分裂を一つ選ぶと持ち上げは$L:k^6\to k^4$で与えられ、24次元である。
従って印付き全空間は既約で次元47となる。一般の持ち上げは純関係を含まない。
実際、Hankel行列$H=(h_{i+j})$の階数一点は一次元しか動かず、
付加する行$Lh$がその行直線に入ることは各点で余次元3である。

階数二以下の四行行列を求める。$B_j(q)$を二次式$q=(q_0,q_1,q_2)$の係数を
一列ずつずらした$(j+1)\times(j+3)$行列とする。
$H$の階数が2なら左核の$[q]\in\mathbb P^2$は一意で$B_3(q)h=0$である。
$H$の二次元の行空間は$\ker B_1(q)$であり、両空間の次元と漸化式から等号となる。
従って付加した行もこの空間に入るための条件は
\[
 M_L(q)h=0,\qquad M_L(q)=\binom{B_3(q)}{B_1(q)L}.
\]
二つの$B_j$は常に全行階数である。$\det M_L$は平面六次式で、
\path{computations/verify_binary_incidence_sextic.py}の整数標本では
素数1009で三つの偏微分に共通射影零点がない。固有性により標数零の一般の
$L$でも滑らかで、従って幾何学的に整である。この場合六次上の$M_L$は常に
階数5なので、核直線が正則に$[h]$を与える。その一般点で$H$は階数2である。
実際、階数一のパラメータ$t$に対し$B_3(q)h(t)=0$は$q(t)=0$となり、
この旗$(t,q)$は二次元である。さらに$B_1(q)Lh(t)=0$は各旗で$L$に
独立な二条件を課すので、一般の$L$では旗の解は高々有限である。
よって$q$の一意性により残余曲線はこの六次と双有理であり、種数は10である。
$H$が階数1ならば付加後も常に階数二以下で、これが五次曲線を与える。

両成分の一般点の重複度は1である。$H$の階数2の開部分では可逆な小行列を用いて
その左核をスキームとして消去でき、上の滑らかな六次の核のグラフとなる。
五次曲線の一般点は二変数の座標で$h=(1,0,\ldots,0)$としてよい。
付加行の末尾三成分を$(u,v,w)$とすると、階数二以下の一次条件は
$(\dot h_2,\dot h_3,\dot h_4)$と$(\dot h_3,\dot h_4,\dot h_5)$がともに
$(u,v,w)$の倍数になることである。一般の$L$では$uw-v^2\ne0$なので、
四つの変分は全て零となり、五次の接方向だけが残る。
四行四列の階数二以下の射影多様体の次数は20である。例えば期待余次元4での
退化軌跡の公式は$ c_2^2-c_1c_3$で、$c_i$は$(1-H)^{-4}$の係数、
従って$100-80=20$である。ここで五平面との交叉は上記の通り期待次元1である。
よって残る次数は15となる。種数10の成分は有理曲線でなく、
前命題の検出する五次曲線は唯一である。その列像から核を復元できるので、
印を忘れても次元47である。Grassmann多様体の次元60との差が主張を与える。
\end{proof}

% !TeX root = ../papers/nc-p3.tex
\section{片側対称な第二補完の階数二曲線の次数}
\label{sec:companion-rank-two-sections}
\begin{proposition}\label{prop:companion-no-binary-quintic}
$(z,\lambda)=(-2,1)$とし、第一行列を
$\left(\begin{smallmatrix}5&2\\3&7\end{smallmatrix}\right)$、
第二行列を$\left(\begin{smallmatrix}1/6&2\\3&-1/4\end{smallmatrix}\right)$
とする。第一関係六平面の階数二以下の軌跡は、素数1009の幾何ファイバーで
次数4と16の二つの整成分を重複度1で持つ。標数零のこの六平面を含む
Grassmann多様体の開集合では、階数二以下の軌跡は次数5の有効一次元
サイクルを含まない。従ってこの辺では両方向とも二変数形式の商が存在しない。
\end{proposition}
\begin{proof}
$\mathbb P(R)=\mathbb P^5$内で関係行列の十六の三次小行列を取る。
$\mathbb F_{1009}$上の三つの超平面切断を
\path{computations/verify_companion_rank_two_sections.py}で計算する。
各切断の$\mathbb P^4$で先頭座標を1とし、残る四変数について
次数5以下の126単項式を並べる。十六の三次式に二次以下の15単項式を掛けた
係数行列の階数は106で、次数5の列は全て主列となる。
従って無限遠超平面に共通零点はなく、全ての高次単項式も四次以下に還元できる。
四行四列の階数二以下の多様体の次数は20で、
$\mathbb P^4$との交叉は固有かつ期待次元0だから切断の長さは20である
（同じ退化軌跡公式をこの四平面に適用できる）。
よって残る20列は座標環の基底である。一座標の乗法の特性多項式は
いずれも平方因子を持たず、その既約因子の次数は順に
\[
 (1,3,16),\qquad(2,2,4,4,8),\qquad(1,4,15)
\]
である。従って各切断は幾何学的に被約で、これらは実際のFrobenius巡回長である。

曲線は期待余次元で、その成分は全て一次元である。これは階数二以下の
行列多様体がCohen--Macaulayで、ここでの線型切断が正則列となることによる。
各成分は超平面に会うため被約切断から一般の重複度は1である。また被約切断の
各点は一つの成分にしか属さない。実際、その点で一つの超平面式を割った
局所環が体なので、曲線の極大イデアルは一元で生成され、正則点である。
Frobeniusが相異なる$r$成分（一つの次数を$e$とする）を推移的に動かすなら、
それらの切断に属する巡回長は$r$の倍数であり、和は$re$である。
どの超平面にも同じ成分の軌道型が適用される。
上の三分割の各々を組に分け、各組の最大公約数の約数$r$を試すと、
共通して残る$(r,e)$のリストは$((1,20))$と$((1,4),(1,16))$だけである。
コードはこの整数の列挙も行う。

薄い二つの関係係数を零にした$\mathbb P^3$では行列がブロック対角となる。
低い二行二列と高い二行二列の各ブロックへの射影は、厚い各色が二座標へ同型に
写るため線型同型である。両行列式が切る二つの滑らかな二次曲面は相異なる。
一致すれば階数二軌跡に曲面が含まれ、既に検算した有限切断に反する。
よってその完全交叉は次数4の有効曲線で、一成分だけの可能性は消える。
特にこの幾何ファイバーはいかなる次数5の有効一次元サイクルも含まない。

固定次数の有効一次元サイクルのChow多様体は固有であり、サイクルの台が
指定した線型空間と階数二以下の閉集合に含まれることは閉条件である。
例えば次数5のサイクルを有限個の成分に分解し、各成分のChow形式で
線型空間と小行列の消滅を課してもよい。従ってサイクルが存在する
六平面の軌跡は閉で、上記の良い還元から有理標本とその開近傍で存在しない。
純関係のない開部分を取ると命題\ref{prop:binary-quotient-rank-two-curve}を
両向きに使える。標本で純関係がないことは、前と同じ二つの角ブロックの行列式と
厚い部分の両射影の単射性から従う。これは非正則性の一つの障害を除く判定である。
\end{proof}

% !TeX root = ../papers/nc-p3.tex
\section{階数十六の四窓から二文字ブロックを回復する}
\label{sec:rank-two-block-recovery}
\begin{theorem}\label{thm:rank-two-block-recovery}
定理\ref{thm:kummer-ramification}の有理重根点と、
定理\ref{thm:klein-elliptic-completion}の$(z,\lambda)=(-2,-2)$の有理単根点を考える。
それぞれを含む印付き族の開近傍で、四辺の各関係五平面の階数二以下の曲線は
次数4と16の二つの幾何学的に整な成分を持つ。
各辺だけから両端の四文字空間の二つの二次元ブロックを、共通の交換を除き一意に回復できる。
第一辺からはさらに$z+1/z$を回復できる。これらは添字保持同型の必要不変量である。
\end{theorem}
\begin{proof}
\path{computations/verify_rank_two_block_recovery.py}は各標本の二切断と
それらの転置について、前節の有限超平面切断を同一素数1009で取り直す。
八つの曲線それぞれで残る成分軌道は$20$か$4+16$である。
厚い二色の四平面上で対角二ブロックの行列式が切る四次曲線があるため、全て後者となる。
奇位置で矢の順を交換しても行列の転置なので成分次数は変わらない。

開近傍への移行を詳しく述べる。期待次元の線型切断と厚い二行列式の
固有交叉を保つ開部分では、その次数20と4の交叉サイクルは族をなす。
これは固定次数の固有交叉のChow形式、または対応するCohen--Macaulayな
定次元の切断の平坦性による。幾何学的に二つ以上の有効サイクルの和に分解することは
次数ごとのChow多様体の加法の固有性から閉条件である。
まず四次の分解を除く。次に全サイクルが指定四次と、和の次数が16となる
二つの正次数サイクルに分解する閉条件を全て除く。
良い還元の二つの整成分はどの条件にも属さないので、標数零の二標本も残る。
これで所要の幾何学的整成分と重複度が開近傍で確定する。

次数4成分は厚い四次元関係空間$D$の射影内の二次式完全交叉であり、
一次式はそのイデアルにないのでその線型包が$\mathbb P(D)$となる。
辺の行列を第二矢空間の双対から第一矢空間への写像とみなす。
$SM=MT$が全ての$M\in D$について成立する自己準同型の対$(S,T)$を求める。
四つの行列についての64本の線型式は32未知数に対し、八辺全ての有理標本で階数30である。
二つの既知ブロックの同時射影が解なので、階数30を保つ開集合で解代数は$k\oplus k$である。
順序を付けない二つの原始べき等元はこの代数だけで決まる。それらの像および
双対像により両矢の二ブロックを取り戻せる。二変数形式の商に必要な次数5の
階数二曲線もどの辺にもない。

第一辺で、回復した二べき等元を用いて対角ブロックを除くと薄い二色の鉛筆を得る。
その二つの角ブロックの行列式は$r^2-t^2,zr^2-t^2$である。
従って鉛筆上の二つの零点対は相異なる四点で、二対の中を同時に交換する
射影直線の自己同型は一意である。現在の表示では$t/r\mapsto-t/r$であり、
その二固定点が薄い二色の方向を与える。二固定点を零と無限遠に取り、
別々の零点対から一つずつ選ぶと座標比の平方は$z$または$1/z$である。
座標の倍率、各対内の選択、二つの固定点の順序、二ブロックの交換によっても
$z+1/z$は変わらない。この計算は四元群の印を入力として必要としない。
\end{proof}

\begin{corollary}\label{cor:three-character-frames}
上の二標本を含む開集合を縮めれば、どの辺にも、各二文字ブロックを二本の
文字直線に分ける候補は三通りで足りる。四位置を通じても、互いに整合する
順序なし文字直線の分解は高々三通りである。
\end{corollary}
\begin{proof}
$D$の低・高ブロックの行列式の対称行列を$O,P$とする。いずれも非退化である。
コードで八辺の全てについて$O^{-1}P$の固有値が相異なり、トレースが非零であることを
有理数上で確認する。これらは開条件であり、両行列式の単元倍にも依存しない。
二つの二次曲面をそれぞれ保つ射影変換は$O^{-1}P$をスカラー倍まで保つ。
トレースからその倍率は1である。よって四固有直線を個々に固定し、
$O$について直交する固有基底で各係数は共通倍率を除き符号である。
両二次曲面は二行二列の階数一行列のSegre曲面であり、両方の母線族を保つには
符号の積が正でなければならない。実際、直交四次元空間で一座標だけの符号反転は
極大等方二平面をそれと直線で交わる平面に写し、二族を交換する。
従って残る群は同時符号を法とした偶符号の四元群である。

非自明三元の各々の二つの固有二平面では両形式とも非退化である。
ゆえに各Segre曲面の固定点は四点で、両方の射影直線因子で非自明な対合として作用する。
どちらの曲面の左・右因子でも、この四元群は忠実である。非自明な二対合の固定点対が
交わることはない。実際、一点を固定する射影直線の有限群は、その点を無限遠にして
有限軌道の平均を原点とすれば乗法群に入り、巡回群である。
もとの文字の下位ビットの符号は、低・高の各ブロックの各矢で非自明で、
厚い関係に作用し二つの行列式と各母線族を保つ。従ってこの三元の一つであり、
四つの二文字空間内の固定点対が求める文字直線を含む候補を与える。
隣の辺では一つの共通ブロックの固定点対を指定すれば候補は高々一つなので、
一周しても数は増えない。これは候補の上界であり、薄い係数の照合でさらに減る場合を含む。
\end{proof}

% !TeX root = ../papers/nc-p3.tex
\section{薄い二色による文字直線の一意性と印なし同型}
\label{sec:unique-character-frames}
\begin{theorem}\label{thm:unique-character-frames}
前節の重根・単根の両標本を含む開集合で、三通りに絞った四位置の文字直線の
分解は一意である。この開集合の標準正則窓同士の添字保持同型は、32個の置換と
非零係数の積だけで判定できる。二十点配置の安定群の検査を必要としない。
\end{theorem}
\begin{proof}
第一辺の低ブロックの列像の射影直線を使う。厚い四次成分の低ブロックを
$(1,t)^t(1,s)$と置き、高ブロックの行列式を$s$の二次式として判別式を取ると、
射影直線への次数2の射影の分岐式
$B(t)=b_4t^4+b_2t^2+b_0$を得る。奇数冪がないことは真の文字符号の作用からも従う。
二標本で四根は相異なり、定数項と最高次係数の比は順に
\[
 3195075625/15968016,\qquad 9604/184041 .
\]
従って$b_0^2\ne b_4^2$を保つ開集合を取れる。

真の文字の対合はこの直線で$t\mapsto-t$である。残る二元のいずれかを
別の文字分解の符号と仮定し、その行列を$S$とする。前節の四元群で可換で、
真の対合と固定点を共有しないため、$S$は反対角行列である。
新しい文字分解の符号も、薄い鉛筆の二つの文字方向をそれぞれ保つ必要がある。
これらは二つの零点対を同時に交換する対合の固定点として既に一意であり、
右上ブロックでは$I$と交換行列$J$である。高い側の相手の変換を$U$と書けば、
$SIU^{-1}$は$I$の倍数、$SJU^{-1}$は$J$の倍数である。
従って$SJS^{-1}$は$J$の倍数で、その係数は$\pm1$である。
反対角な$S$についてこれは$t\mapsto 1/t$または$-1/t$を与える。
ところが$S$は四次成分を保つので分岐因子も保つ。
そのためには$t^4B(1/t)$が$B(t)$の倍数となり、$b_0^2=b_4^2$が必要で矛盾する。
よって第一辺の一つの文字対しか残らず、系\ref{cor:three-character-frames}で
全四位置への延長も一意である。

文字直線を回復した後、窓同型は各位置で単項行列になる。ブロックを保つ置換は
$ i\mapsto i+u_j+\epsilon_j i_{\rm hi}\,1$の形である。
ここで和は四元群の和、$1$は下位ビットの元である。
テンソルの非零語は総文字零である。高ビットの全ての偶パターンで少なくとも
一つの非零語がある。全て低い場合と全て高い場合は可逆な主ブロック、
連続する同じ組から別組へ移る場合は非零な交差ブロック、他は薄い係数を使う。
四位置の置換後にも総文字零なので、全て低い語から$\sum u_j=0$、
二つずつ高い語から$\epsilon_j$は全位置に共通と分かる。
二つの隣接切断で厚い色の組を保つため$u_0,u_1,u_2$の高ビットが一致し、
和零から$u_3$も同じになる。従って候補数は$2\cdot4\cdot2\cdot2=32$であり、
係数の零を含む標本でも同じ上界である。

各置換後、支持が一致することを確かめる。非零語を行、位置と文字を列とする
整数行列の各行に、その語の四つの文字で1を置く。この行列を$M$、係数比を
$c_\alpha$とすれば、対角基底変換で一致する条件は
\[
 \prod_\alpha c_\alpha^{n_\alpha}=1\qquad(n\in\ker_{\mathbb Z}M^t)
\]
である。Smith標準形で整数核の基底を取り、残る非零対角成分では代数閉体で
根を取れば十分性も従う。全体スカラーは一位置の倍率で吸収できる。
標準正則窓では最上テンソルがスカラーまで決まり、その四切断で関係列が戻る。
定理\ref{thm:canonical-four}により標準四周期を選ぶこの比較が添字保持同型の必要十分判定となる。
\path{computations/verify_unique_character_frames.py}は上の分岐式と32置換を有理数上で検査する。

両標本自身で合格する自己置換は恒等一つである。対角倍率の重み行列の
非零Smith因子は両標本とも1が11個、2が2個である。
従って標準四窓表示の位置ごとのスカラーを除く射影安定群は、
既知の同時文字符号の四元群に一致する。
\end{proof}

% !TeX root = ../papers/nc-p3.tex
\section{同時文字交換が残る二直線と二つの印を持つ正則窓}
\label{sec:flip-character-locus}
$S_z\left(\begin{smallmatrix}a&b\\c&d\end{smallmatrix}\right)
=\left(\begin{smallmatrix}d&c/z\\zb&a\end{smallmatrix}\right)$とする。
\begin{theorem}\label{thm:simultaneous-character-flip}
規格化した可逆チャートで、四位置全てで$i\mapsto i+1$を誘導する単項自己同型が
存在する必要十分条件は、ある同じ$\sigma=\pm1$について
$S_z A_0=\sigma A_0$, $S_z A_1=\sigma A_1$となることである。
第一行列の条件$d=\sigma a$, $c=\sigma zb$だけで第二行列の条件も従う。
自己同型の平方は各位置で、射影的には低ブロック上の1と高ブロック上の$\sigma$である。
従って二ブロックと三候補を復元する前節までの条件の下では、
$\sigma=+1$の直線を除いて文字分解は一意である。この直線では実際に二つの異なる
印を持つ反対階数16の標準正則窓が存在する。
\end{theorem}
\begin{proof}
$u^2=z$とし、同時文字交換の後、四位置の成分ごとに順に
\[
 (u,u,1,z),\quad(1/z,1,1/u,1/u),\quad
 (u,u,1,z),\quad(1,1/z,1/u,1/u)
\]
を掛ける。薄い成分は元に戻る。厚い上左ブロックは二行列とも$S_z$で変わり、
二つの交差ブロックは不変である。右下はSchur補で一意に定まるため、
正規形は$(S_zA_0,S_zA_1,z,\lambda)$となる。
対角変換でこの二行列をさらに同じ定数だけ掛けられる。
逆に規格化を保つ対角変換の残りを確認する。薄い二色では低群から高群への
$2\times2$が全て非零である。二色で右・左薄ベクトルに許される逆倍率を
$t_2,t_3$とすると、低群の二行と高群の二列でたすき掛けに比較して
$t_3/t_2=\epsilon=\pm1$である。高群の二行・低群の二列でも、薄い元係数を
左右から消せば同じ比となる。従って四位置とも各組内の比は$\epsilon$である。
位置0の低・高の代表倍率を$A,B$、位置2を$C,D$とすると、薄い式で位置1,3の
倍率も決まり、厚い交差ブロックを保つ条件は$AD=BC$となる。
二つの上左ブロックには同じ$(A/B)^2$が掛かる。
特に対角自己同型なら$A/B=\pm1$で、位置ごとのスカラーを除けば同時文字符号だけである。
$S_z^2=1$なので、二行列を同時倍率で戻せるのは定理の固有条件に限られる。
第一固有条件を補完式に代入すると$g_0=zf_0$, $h_0=e_0$で、
負の固有値では両定数が零になる。また$g_1=\sigma z f_1$, $h_1=\sigma e_1$であり、
任意の根に対して$g=\sigma zf$, $h=\sigma e$が従う。

$u^2=\sigma z$と取り直した同じ四倍率と文字交換は、固有条件上の窓を固定する。
各二文字でこの単項変換の平方は二つの倍率の積であり、高・低ブロックの
その比は全位置で$\sigma$である。対角文字符号を掛けても積は両ブロックに
同一の符号だけ掛かり、この比は変わらない。
別の文字分解があれば、その下位ビット符号は三候補の非自明な別元である。
各二文字の射影直線で元の符号と可換で固定点を共有しないため、元の二文字を交換する。
四位置への伝播で全てのブロックに同じ結論が出る。これは射影位数2である必要があり、
上の必要条件は正の直線に絞られる。

具体的に$(z,\lambda)=(3,2)$で
\[
 A_0=\begin{pmatrix}0&1\\3&0\end{pmatrix},\qquad
 A_1=\begin{pmatrix}2&8/13\\24/13&2\end{pmatrix}
\]
とする。$x=-192/169,y=1$で二つの符号は$-1,1$、従って既約な対根族上にある。
三切断の階数は$6,16,6$で、四次中心次元は$(3,2,0,0)$である。
四中心切断による商は素数1009で全始点・次数13まで$(1+t+t^2+t^3)^4$の
次元を持つ。定理\ref{thm:kummer-ramification}と同じ上下界と自由性の証明により
この有理窓は標準Koszul左右AS正則である。

同時文字交換の各倍率に下位ビット符号を掛け、低・高の各二文字でその二固有直線を
新しい文字に取る。各位置で固有値の組は低・高で等しく、総テンソルは不変なので、
新しい文字も総次数零となる。これらの直線はいずれも元の座標直線と異なる。
二つの切断で厚い各ブロックの四つの二行二列主・交差ブロックが全て可逆で、
薄い二色の全成分が非零であることを、$\sqrt3,\sqrt{-1}$を含む良い還元1009で確認する。
両切断の全階数6とこれらの下界から各色の階数は$2,2,1,1$となる。
$z$に対応する薄い二行列の零点対は元と同じ交差比を与えるので、新表示の
$z$も$3$または$1/3$である。\path{computations/verify_flip_character_locus.py}
は倍率、二直線の係数、四中心切断および新固有基底の非零条件を検算する。
\end{proof}

% !TeX root = ../papers/nc-p3.tex
\section{文字交換の固定直線上の楕円曲線と正則性}
\label{sec:fixed-line-elliptic}
通常倍率、すなわち
$z\lambda(z^2-1)(\lambda^2z-1)(\lambda^2-1)(\lambda^2z^2-1)\ne0$
を取る。前節の固有直線を $L_\sigma$ ($\sigma=\pm1$) とし、その有限座標と補完を
\[
 A=\begin{pmatrix}a&1\\\sigma z&\sigma a\end{pmatrix},
 \qquad B=\begin{pmatrix}e&f\\\sigma zf&\sigma e\end{pmatrix}
\]
と書く。以下の点と開集合は標数零の幾何点について述べる。

\begin{theorem}\label{thm:fixed-line-elliptic}
各 $L_\sigma$ 上の符号付き対根の滑らかな完備化は楕円曲線である。
二曲線の $j$ 不変量は等しく、
\[
 j=256\frac{(1-k+k^2)^3}{k^2(1-k)^2},\qquad
 k=\frac{(\lambda^2-1)(\lambda^2z^2-1)}{(\lambda^2z-1)^2}.
\]
各直線はKummer四次面と横断的に四点で交わり、八点全てで両行列が可逆な真の重根を得る。
$(z,\lambda)=(3,2)$を含む倍率平面の非空開集合では、
八点全てが反対階数16の標準Koszul左右AS正則窓となる。
各始点は四つの中心四次式の多項式環上で階数256の自由加群である。
この開集合の各倍率で、各楕円曲線上の可逆窓は有限個の候補を除いて同じ性質を持つ。
\end{theorem}
\begin{proof}
$P(x)=px^2+mx+n$, $Q(y)=qy^2+hy+o$とし、
$F_\sigma=\lambda^2(z^2-1)^2(a^2-z)^2(m^2-4pn)$と置く。
$a$の四次式の最高次係数は $c=\lambda^2(z^2-1)^2$、定数項は $z^2c$ で、
その判別式は
\[
 2^{12}\lambda^4 z^6(\lambda^2-1)^2(\lambda^2 z^2-1)^2
 (\lambda^2z-1)^4(z-1)^4(z-\sigma)^{12}.
\]
従って四根は相異なる有限非零点であり、制限四次の単根なので面の滑らかな横断点である。
二直線は共通の射影点を持たない。二重被覆 $s^2=F_\sigma$ はRiemann--Hurwitzにより種数1、
無限遠に有理点を持つ。四次式の係数を降冪順に $A',B',C',D',E'$ と書き、
$I=12A'E'-3B'D'+(C')^2$を用いると $j=256I^3/\mathrm{disc}(F_\sigma)$。
$k=u/v$を上記の分子・分母で表すと
$I=16z^2(z-\sigma)^4(u^2-uv+v^2)$ となり、表示を得る。

$F_\sigma$と$a^2-z$の終結式は $16z^4(z-\sigma)^4(\lambda^2z-1)^4$。
重根公式 $x=-m/(2p),\ y=-h/(2q)$を補完式に入れた
$p,q,x,y,e,f,e^2-zf^2$ の各分子・分母についても、既約因子と
$F_\sigma$との終結式の因子は通常倍率で逆にした因子に限られる。
\path{computations/verify_fixed_line_elliptic_curves.py} がこれらの多項式恒等式を検算する。
ゆえに重根で元の可逆表示が定まる。

倍率 $(3,2)$ で原始化した二式は順に
\[
 16a^4-156a^3+555a^2-468a+144,\qquad 4a^4+339a^2+36.
\]
素数307ではそれぞれ $102,153,298,301$ と $80,119,188,227$ で単根に分裂する。
\path{computations/verify_fixed_branch_regular.py} は各剰余点の二行列から
四つの関係空間を作る。隣接階数はともに6、反対階数は16、
全始点で長さ4,5の次元は35,56、四次中心次元は文字順に $(3,2,0,0)$ である。
中心の最初の二文字次数から各二式を選ぶ。商の次元は全八点・四始点のいずれでも、
次数0から13の順に
\[
 1,4,10,20,31,40,44,40,31,20,10,4,1,0
\]
となる。

単根なので各剰余点は307進整数環に一意に持ち上がり、四根を両式とも尽くす。
各持上げで経路次元の上界は剰余体での値、下界は五次元対根族の正則一般点での値である。
従って長さ4,5の経路空間はこの離散付値環上で自由であり、中心方程式を同じ基底で比較できる。
一般点の中心からの下界5と剰余体の上界5が等しいので方程式は定階数で、
選んだ四式は中心に持ち上がる。商次元の上界と元の全次数の下界を合わせ、
定理\ref{thm:klein-rational-period-four}の自由基底と左右Koszul列の証明を適用する。
体拡大での完全性を降下でき、中心四式を選べる条件も核空間内の非空開条件である。
従って標数零の幾何点で主張が成り立つ。

長さ4,5の経路空間、中心方程式、商の長さ13までの階数条件は同じ議論で
各標本の開近傍に延びる。二つの四次式の根の軌跡は通常倍率上で有限エタールである。
八近傍の外の閉集合をこの有限射で倍率平面に写し、その閉像を除けば
$(3,2)$を残して八点を同時に扱える。固定倍率の滑らかな各被覆曲線にも
この正則な近傍が開集合として交わり、その補集合は有限である。
\end{proof}

\begin{proposition}\label{prop:fixed-line-defects}
各通常倍率で次の十二表示は全て真の相異なる対根の族に属し、長さ4で標準性を失う。
一行に同じ $a$ での二つの補完を記す。
\[
\begin{array}{cc|cc}
 \sigma&a&(e,f)&(e,f)\\ \hline
 +&\lambda z&(1,\lambda)&(z,\lambda z)\\
 +&1/\lambda&(z,1/\lambda)&(1,1/(\lambda z))\\
 -&\lambda z&(-1,\lambda)&(z,-\lambda z)\\
 -&-\lambda z&(1,\lambda)&(-z,-\lambda z)\\
 -&1/\lambda&(-z,1/\lambda)&(1,-1/(\lambda z))\\
 -&-1/\lambda&(z,1/\lambda)&(-1,-1/(\lambda z))
\end{array}
\]
\end{proposition}
\begin{proof}
$x=\sigma(a^2-zf^2)$, $y=\sigma(e^2-z)$で $P(x)=Q(y)=0$ かつ
$2px+m=-(2qy+h)\ne0$ となる。両行列式と $p,q,x,y$ も通常倍率では消えない。
前記の曲線の検算コードは十二表示について、関数体 $\mathbb Q(z,\lambda)$ 上で
四辺の六関係の全ての両側挿入を長さ4まで掃き出す。
正の四表示では全始点で次元36、負の八表示では始点を並べ替えて
$35,35,36,36$ である。通常倍率の特殊化でも関係は両切断で定階数6であり、
商次元は減少しない。標準窓の次元35に反する始点が常に残る。
\end{proof}

% !TeX root = ../papers/nc-p3.tex
\section{固定楕円曲線上の二位置移動と四次障害の商}
\label{sec:fixed-exchange-isogeny}
前節の楕円曲線を $E_\sigma:S^2=F_\sigma(a)$ とし、
$s=S/(\lambda(z^2-1)(a^2-z))$ を $2px+m$ に対応させる。
\begin{theorem}\label{thm:fixed-exchange-isogeny}
通常倍率で二行列を交換する二位置移動は $E_\sigma$ 上の固定点のない対合である。
原点を選べば位数2の点による平行移動であり、商写像は次数2の同種写像となる。
次の行列を $M=\left(\begin{smallmatrix}A&B\\C&-A\end{smallmatrix}\right)$、
$d=\det M$, $D=Ca-A$, $b=(Aa+B)/D$ と書く。
\[
 M_+=\begin{pmatrix}2\lambda z&-z(\lambda^2z+1)\\
 \lambda^2z+1&-2\lambda z\end{pmatrix},
 \qquad M_-=\begin{pmatrix}0&-z\\1&0\end{pmatrix}.
\]
交換の式は $(a,S)\mapsto (b,dS/D^2)$ である。
$C\ne0$では商の関数体を次の四次モデルで表示できる。
\[
 U=a+b,\quad V=S\frac{dU}{da},\quad
 V^2=\left(U^2-4(AU+B)/C\right)H(U),\quad
 H(a+b)D^2=C^2F_\sigma(a).
\]
$H$は二次式である。$C=0$の場合にも商写像の次数と平行移動の記述は保たれる。
\end{theorem}
\begin{proof}
固定直線上では $g_1=\sigma z f_1$, $h_1=\sigma e_1$ で、負の直線では定数項が消える。
$2px+m=s$ と $2qy+h=-s$ を使うと、第二行列の非対角・対角成分は
\[
 f=(s-1)/(2\sigma z f_1),\qquad e=(1-s)/(2\sigma e_1).
\]
従って $e/f=-zf_1/e_1=b$ は符号に依存しない。交換して両行列を $f$ で割ると、
新しい第一行列の座標は $b$、第二の非対角成分は $1/f$ になる。
定理\ref{thm:kummer-polar-completion}のテンソル恒等式によりこれは二位置移動を実現する。

$M^2=-dI$で $d$ は正負の順に $z(\lambda^2z-1)^2,z$ となり、
$F_\sigma(b)D^4=d^2F_\sigma(a)$ である。符号も決定するため、
$f=f_0+\sigma f_1(a^2-zf^2)$ を用いる。これを法として
\[
 (b^2-z)D^2(f+2\sigma z f_1(b))
 =d(a^2-z) f(1+2\sigma z f_1(a)f)
\]
となるので、実際の行列交換は上記の正符号の持上げである。
\path{computations/verify_fixed_exchange_isogeny.py} は両符号で恒等式を検算する。
滑らかな射影曲線の双有理写像なので全域に延びる。$db/da=d/D^2$ で正則微分
$da/S$ を保つ。楕円曲線の自己同型は平行移動と原点を固定する自己同型の積であり、
標数零では後者の微分への作用は忠実である。基底上の変換が非自明な対合なので
非零の位数2の平行移動となる。

$C\ne0$で $a^2-Ua+(AU+B)/C=0$ となる。$F_\sigma/(a-A/C)^2$ は交換不変で、
$U$の無限遠にのみ高々二位の極を持つから二次式 $H(U)$ である。
$a$の無限遠での最高次係数は非零なので次数は正確に2である。
微分公式 $(a-A/C)dU/da=2a-U$ から $V^2$ の式が出る。
$V$は交換不変で、$S$は $V$ と $a$ から復元できる。従って生成した体は指数2である。
\end{proof}

\begin{corollary}
四始点のいずれかで長さ4の標準次元を失う条件と、標準Koszul左右AS正則である条件は、
交換が窓に作用する共通定義域でこの商に降りる。一般倍率では前節の十二の非正則表示は
二つの商曲線上の合計六点に写る。
\end{corollary}
\begin{proof}
交換は始点を二つ移し、各文字空間の基底を変える。従っていずれかの始点での次元増加も、
全始点に課す左右正則性も保つ。十二表示の表は
$(a,e,f)\mapsto(e/f,a/f,1/f)$で閉じる。一般倍率で表の点は相異なり、対合は自由である。
この商は障害判定のためであって、添字保持同型を与えるものとは限らない。
\end{proof}

具体的な階数検査も小さくできる。第 $i$ 辺の二次商を $Q_i$、関係を $R_i$ と書くと、
長さ4の道空間は
\[
 V_i\otimes R_{i+1}\otimes V_{i+3}\longrightarrow Q_i\otimes Q_{i+2}
\]
の余核である。外側二対で先に関係を割り、残る中央の挿入を写したからである。
四文字次数で標準窓の長さ4の次元は順に $(9,10,8,8)$ である。
実際、群環で $v$ を四元群の正則表現、$r=(2,2,1,1)$ を二次関係の文字とすれば、
Koszul列の三次・四次項はそれぞれ $v$ と自明文字 $1$ である。
長さ4の係数は $v^4-3v^2r+r^2+2v^2-1$ となる。上記の余域は
$(26,26,24,24)$ なので、標準値に対応する階数は $(17,16,16,16)$ である。

\path{computations/verify_fixed_quartic_ranks.py} は倍率 $(3,2)$、素体
$\mathbb F_{101}$ で $a$ が有限、$pqxy\det A\det B\ne0$ の有理点を全て調べる。
正の曲線の101点のうち97点、負の曲線の104点のうち96点で全始点が標準次元を持つ。
他は全て前節の表の還元で、増加する文字は常に自明文字である。
この有限体の列挙は、高次拡大上の点や標数零での障害因子の全体を決定しない。

% !TeX root = ../papers/nc-p3.tex
\section{固定楕円曲線の四次障害の幾何点と反対行列式}
\label{sec:fixed-geometric-quartic-locus}
正の固定直線で $A=\left(\begin{smallmatrix}a&b\\zb&a\end{smallmatrix}\right)$,
$B=\left(\begin{smallmatrix}e&f\\zf&e\end{smallmatrix}\right)$とする。座標は共通倍率まで取る。
\begin{theorem}\label{thm:opposite-fixed-dets}
$r^2=z$ とし、$r$ の二通りの選択と $\epsilon=\pm1$ ごとに
\[
\begin{aligned}
 a&=r\{\lambda r(r+1)+\epsilon(r-1)\},&
 b&=r+1+\epsilon\lambda r(r-1),\\
 e&=r\{r+1-\epsilon\lambda r(r-1)\},&
 f&=\lambda r(r+1)-\epsilon(r-1)
\end{aligned}
\]
と置く。全ての通常倍率でこれらは可逆な対根窓で、$\det B=-\det A$、反対階数16を持つ。
全始点で長さ4の次元は36以上、長さ5は68以上で、一般倍率ではどちらも等号である。
\end{theorem}
\begin{proof}
$x=a^2-zf^2=4\epsilon\lambda r^3(z-1)=-y$,
$\det A=4r^3(\lambda^2z-1)=-\det B$を直接得る。
$P(x)=Q(y)=0$、符号の一致と補完式を代入で確認できる。
一般点では真の二次式の対根が窓を与え、全ての通常倍率で四つの行列式
$x,y,\det A,\det B$は単元であるから、切断のSchur関係も特殊化できる。
反対切断の文字別四行四列ブロックは全て行列式
$-(\lambda-1)^2(\lambda z-1)^2/\lambda^2$を持つ。
\path{computations/verify_opposite_fixed_determinants.py} はこれらを検算し、
$\mathbb Q(r,\lambda)$上で四始点の経路次元 $1,4,10,20,36,68$ を
関係の全挿入から計算する。各二次関係空間は定階数なので特殊化で商次元は減らない。
\end{proof}

\begin{theorem}\label{thm:fixed-geometric-quartic}
倍率 $(z,\lambda)=(3,2)$ の両固定曲線の全ての許容幾何点で、
いずれかの始点の長さ4が標準次元35を超える軌跡はちょうど十六点である。
正の曲線では命題\ref{prop:fixed-line-defects}の四点と上の四点、
負の曲線では同命題の八点で尽きる。二位置移動による商では各曲線に四点ずつとなる。
\end{theorem}
\begin{proof}
まず $b=1$、$pq\ne0$ とする。前節の $M$ の分子・分母を $N,D$ と書く。
$e=Nf/D$で、曲線の式は分母を払った
$g=\sigma z f_1 f^2+f-f_0-\sigma f_1 a^2=0$である。
正の側で $N^2-zD^2=-z(\lambda^2z-1)^2(a^2-z)$。
従って $\det B=-\det A$ は $f=D/(r(\lambda^2 z-1))$ と書ける。
これを $g$ に代入し分子を因数分解すると、$D$ と上の二表示の
$b_\epsilon a-a_\epsilon$ の積の定数倍である。
ここで $r$ の両根を取ると、新しい四点の $a$ 座標は
\[
 a^2-78a+69=0,\qquad 23a^2-78a+3=0.
\]
四根は相異なり、既知の $a=6,1/2$ からも離れる。各根に対応する新しい点は一つずつである。

除外の上界を得るため、前節の外側二対を割った挿入行列を使う。
四位置の各文字 $h$ で切断の最初の $r_h=2,2,1,1$ 本の列を取る。
関係行列の右端の $r_h$ 列が定数可逆なので、これを単位行列にして
行を $(L\mid I)$ と書ける。商の座標は核基底 $\left(\begin{smallmatrix}I\\-L\end{smallmatrix}\right)$から得られる。
挿入行列の各成分は $a,e,f$ の多項式となる。
\path{computations/verify_fixed_geometric_quartic_locus.py} は全四始点・四文字次数で
目標階数 $17,16,16,16$ の小行列式を取る。素数307での掃き出しは行と列を選ぶためだけに用い、
選んだ全行列式は有理係数で再計算する。

各小行列式に $e=Nf/D$ を入れて分母を払った多項式を $W_j$ とする。
$\operatorname{Res}_f(g,W_j)$ の最大公約数の平方自由部分から、
$a^2-3=0$, $pqD=0$ および既知の十二表示の $a$ 座標を取り除いたものを $u(a)$ とする。
既知の座標の上では既に二つの相異なる根が尽くされている。
$u$ が定数なら残余は空である。それ以外の場合は
$(u,g,W_1,\ldots,W_m)\subset\mathbb Q[f,a]$ のGröbner基底を求める。
$xy\det B$ の分母を払った式（正の側ではさらに $\det A+\det B$ を掛ける）の
平方の正規形は零になる。従って残った全幾何点は非許容点または上で同定した点である。
目標階数は正則一般点からの特殊化で上界でもあるため、これで残りの点の等号が従う。

最後に $b=0$ または $pqD=0$ のうち $\det A\ne0$ の点を直接調べる。
正では $[a:b]=[13:4],[12:13],[1:0]$ で各一点が許容され、
四始点とも長さ4,5の次元は35,56である。負では $[0:1],[1:0]$ に許容点はない。
検算は各 $P,Q$ の全有限根から符号と $xy\det B\ne0$ で選別する。
以上で正負それぞれ八点となり、前節の自由な対合はその集合を保つ。
\end{proof}

% !TeX root = ../papers/nc-p3.tex
\section{全通常倍率での四次障害の決定}
\label{sec:uniform-fixed-quartic}
\begin{theorem}\label{thm:uniform-fixed-quartic}
標数零の任意の通常倍率で、二本の固定直線の許容対根窓において、
どこかの始点の長さ4の次元が35を超える必要十分条件は次のいずれかである。
\begin{enumerate}
\item 命題\ref{prop:fixed-line-defects}の十二表示の一つである。
\item 正の直線上で $\det A+\det B=0$ となる。
この条件は定理\ref{thm:opposite-fixed-dets}の四表示で尽きる。
\end{enumerate}
正の曲線では常に八点、負の曲線では $\lambda^2z\ne-1$ なら八点、
$\lambda^2z=-1$ なら四点である。従って二位置移動の商では合計八点または六点である。
これは各倍率の幾何点の台に関する主張である。
\end{theorem}
\begin{proof}
$[a:b]$を第一行列の対角・非対角の比とし、
$(N,D)^t=M_\sigma(a,b)^t$ を定理\ref{thm:fixed-exchange-isogeny}の行列で定める。
許容点で $[e:f]=[N:D]$ なので、ただ一つの $u\ne0$ に対し $(e,f)=(Nu,Du)$ となる。
この比の式は稠密開集合から延びる。両行列は可逆で $M_\sigma$ も可逆なので両辺に基点はない。
$\Delta=a^2-zb^2$ と書く。$d=\det M_\sigma$ とすると
$N^2-zD^2=-d\Delta$ であり、$\det B/\det A=-d u^2$ である。
補完の二つの式は
\[
 \frac{\lambda(z+\sigma)\Delta}{D}
 \bigl(Du-f_0-\sigma f_1(a^2-zD^2u^2)\bigr)
 =
 \frac{\lambda(z+\sigma)\Delta}{N}
 \bigl(Nu-e_0-\sigma e_1(N^2u^2-zb^2)\bigr)
 =G(a,b,u)
\]
と共通の多項式になる。$N,D$は同時に消えないため、片方が零でも $G=0$ が必要である。
正で $u=1/(r(\lambda^2z-1))$ を代入すると $G$ は前節の
二つの一次式 $b_\epsilon a-a_\epsilon b$ の積の単元倍となる。
$r$の両根を取れば、無限遠を含めて反対行列式の軌跡が決定する。

次に $b=1$ での第一座標の候補を、倍率と同時に消去する。
前節の各始点・各文字次数の目標階数に対応する小行列式を一つ選ぶ。
$e=Nu,f=Du$を入れ、通常倍率で単元の分母を払うと
$T\in\mathbb Z[a,u,z,\lambda]$になる。
\path{computations/verify_uniform_fixed_quartic_support.py} は全32行列について
分数なしの消去で $T$ を計算し、各既約因子 $W$ について
$u$ を含めば $\operatorname{Res}_u(G,W)$、含まなければ $W$ を因数分解する。
いずれも非零で、倍率のみの因子は全て通常倍率の単元である。残りは次の表の積の因子に限る。
\[
\begin{array}{c|l}
 - & a\Delta(a^2-\lambda^2z^2)(\lambda^2a^2-1)\\
 + & a\Delta ND(a-\lambda z)(\lambda a-1) H_+(a)H_-(a)
\end{array}
\]
ここで $H_\epsilon=\operatorname{Res}_r(r^2-z,b_\epsilon a-a_\epsilon)$。
行列式の計算は四変数の整数多項式環内で行うので、表の制限は全ての通常倍率に
同時に適用できる。負の直線では $a=0$ なら $D=0$ より $x=0$、
$b=0$ なら $N=0$ より $y=0$ である。残る四つの第一座標の上の二根は
前の表の非正則表示で尽くされる。

正では $a=\lambda z$ または $1/\lambda$ のファイバーは同様に既知である。
$a=0$ または $b=0$ 上の許容点での追加検算を述べる。$a=0,b=1$ なら
$x=-zf^2\ne0$より $f=-1/(zf_1)$ かつ $e=Nf/D$。
$a=1,b=0$ なら $y=e^2\ne0$ より $e=1/e_1$ かつ $f=De/N$ である。
必要な分母が消えるとこれらの式を導く元の方程式に許容解はない。
$ND=0$ は二位置移動でこのいずれかに帰着する。
最後に $H_\epsilon=0$ で $abND\ne0$ とする。
第一行列を前節の $A_\epsilon$ に規格化でき、$f_1,e_1$ はともに非零である。
一つの根は既知の $B_\epsilon$。他の根は二次式の根の和から
\[
 f'=-1/(zf_1)-f_\epsilon,\qquad e'=1/e_1-e_\epsilon
\]
となる。

\path{computations/verify_uniform_fixed_quartic_residuals.py} は
二座標端とこの二つの相方に対して、全64個の目標階数の小行列式を
$\mathbb Q[r,\lambda]$で検算する。分子の各非定数因子は通常倍率で逆にした因子、
または当該表示の $xy\det A\det B(\det A+\det B)$ の分子・分母の因子に限られる。
よって分母が定義された許容点では、反対行列式の場合を除いて目標階数に達する。
必要性が従い、十分性は前の二命題の経路次元の下界による。

点の個数は $t=\sqrt d\,u$ によって確認できる。正では $\sqrt d=r(\lambda^2z-1)$、
負では $\sqrt d=r$ を取る。従来の表の $t$ はいずれも
$\{r,-r,1/r,-1/r\}$ を走り、通常倍率ではこれらは相異なる。
正では各一つであり、新しい四点は $t=\pm1$ の上に二つずつある。
同じ $t$ の二つの第一座標の行列式は
$-2r(r^2-1)(\lambda^2 r^2-1)$ で非零である。
負では各 $t$ の二つの $a$ の差の平方が
$(\lambda^2z+1)^2/\lambda^2$である。同じ $a,t$ は同じ窓を指定する。
このため指定の倍率だけ負の四組が合流する。自由な二位置移動で点の個数は半減する。
\end{proof}

% !TeX root = ../papers/nc-p3.tex
\section{固定曲線上の長さ5の挿入行列}
\label{sec:fixed-quintic-locus}
\begin{proposition}\label{prop:fixed-quintic-slice}
倍率 $(3,2)$ において両固定曲線上の許容窓の長さ4の次元が全始点で35ならば、
長さ5の次元は全始点で56である。従ってこの倍率で長さ5に初めて現れる障害はない。
\end{proposition}
\begin{proof}
始点 $i$ からの長さ5の道について、外側の二対で関係を割った後の空間は
$Q_i\otimes V_{i+2}\otimes Q_{i+3}$ である。左と右に残る二次関係の挿入は
\[
 (V_i\otimes R_{i+1}\otimes Q_{i+3})\oplus
 (Q_i\otimes R_{i+2}\otimes V_{i+4})
 \longrightarrow Q_i\otimes V_{i+2}\otimes Q_{i+3}.
\]
この余核は長さ5の道空間そのものである。文字ごとの余域次元は100、定義域次元は120となる。
実際に外側二対の文字を選ぶと中央の一文字が決まり、余域の次元は $10^2$。
片方の挿入では関係と対の文字を選べば残る一文字は一意で、次元は $6\cdot10$。
正則一般点では同時文字交換が奇数長の下位ビットを交換する。
さらにKoszul列の群環表示を長さ5まで展開すると、長さ5の文字は四元群の正則表現の14倍である。
実際、奇数次数の各項は正則表現 $v$ を因子に持ち、総次元は56である。
よって目標階数は全て86で、各文字の商次元14は特殊化での下界になる。

\path{computations/verify_fixed_quintic_locus.py} は正負と始点0,1と各文字次数について、
上の行列の二つの $86\times86$ 小行列式を有理数体上で計算する。
始点2,3は二位置移動で戻せる。$b=1$ とし、前節の $(e,f)=(Nu,Du)$ を使う。
各小行列式の既約因子を $G$ との終結式で第一座標に写し、両小行列式に
共通する第一座標の既約因子ごとに曲線上の同時零点を計算する。
$\det A=0$ と既知の全根が不適格な第一座標を除き、
正の側の反対行列式と 可逆表示の非許容条件は曲線上で飽和する。
$u$ の無限遠や $ND=0$ および $a=0$ の上の許容点は、
前の端点計算の長さ5と二位置移動で検査済みである。

負の側の残余は空である。正の側では各始点の文字0,1の検算において、
飽和後に第一座標の既約八次式の根が二組残る。いずれも消去イデアルは
$(P(a),u-U(a))$ の形で、$P$ は既約、$U$ は有理係数である。
各数体について素数1009で $P$ の素体根を一つ取り、$U$ の分母が単元で
実際の挿入行列の階数が86であることを確認する。
従って小行列式の一つはその数体で非零であり、全ての共役点で目標階数に達する。
この還元は既約因子ごとの判定で、素体の点の列挙だけから他の幾何点を除外していない。
\end{proof}

% !TeX root = ../papers/nc-p3.tex
\section{固定曲線の倍率を動かしたときの長さ5}
\label{sec:varying-fixed-quintics}
\begin{theorem}\label{thm:varying-fixed-quintics}
負の固定曲線では全通常倍率で、正の固定曲線では $(3,2)$ を含む倍率平面の
非空Zariski開集合で、長さ4を全始点で通過した全ての許容窓の長さ5の次元は全始点で56である。
\end{theorem}
\begin{proof}
同時文字交換の自己同型は各矢の次数を $j\mapsto j\mathbin\oplus1$ に送る。
従って奇数長の道で次数 $h$ と $h\mathbin\oplus1$ の次元が等しい。
二位置移動も併用すると、正負ごとに始点0,1、文字0,2を検算すればよい。
前節の行列から各一つの階数86の小行列式を取り、
$e=Nu,f=Du$ を入れて通常倍率で単元の分母を払う。
\path{computations/verify_varying_fixed_quintics.py} はこの行列式を
$\mathbb Z[a,u,z,\lambda]$ で因数分解する。具体的には非零成分を掃き出す際、
ピボット $p$ と消去成分 $x$ の最大公約数を $g$ とし、
行を $(p/g)R-(x/g)R_{\rm pivot}$ に置き換える。
各行の内容を割り、取り除いた因子と $p/g$ の逆倍率を記録する。
行・列の交換の符号と有理定数を除いた行列式の全因子を復元できる。

各因子を $G$ との $u$ に関する終結式で第一座標へ写す。
負では四つの検算の全てで、非定数因子は定理\ref{thm:uniform-fixed-quartic}
の表と通常倍率で逆にした因子に限られる。全根が四次で失格する座標と
非許容端点を除けば、従って全通常倍率で目標階数に達する。

正の側では、同定理の第一座標の表以外にも因子 $P(a,z,\lambda)$ が出る。
これらについて検算は個々に次の三条件を確認する。
\begin{enumerate}
\item 倍率だけに依存する追加因子はない。
\item $(3,2)$への特殊化は $a$ の次数を保ち、平方自由である。
\item 特殊化の全根で、$a\Delta ND$ と $u$ の二次式 $G$ の最高次係数は非零である。
既知の四次障害の第一座標はなく、$G$ と $xy\det B$ の終結式も消えない。
\end{enumerate}
これにより $(3,2)$ の近傍で $P=0$ は倍率平面に有限で、
その上の曲線 $G=0$ も有限になる。必要に応じて基底を縮め、上記の非零条件を保つ。
この有限族の中で目標階数に達しない閉集合の像は閉じる。
命題\ref{prop:fixed-quintic-slice}によりその像は $(3,2)$ を含まない。

表の中の第一座標についても注意する。正の古い二座標の根は四次で失格する。
座標端とその交換相手には、許容性の下で一意な有理表示がある。
二つの $H_\epsilon$ の根では、反対行列式の既知の点の相方を
定理\ref{thm:uniform-fixed-quartic}の式で比較する。$(3,2)$で
これらの四つの第一座標は相異なり $abND\ne0$ である。
相方では $xy\det B(\det A+\det B)\ne0$ である
（定理の四つの残余表示に $r=\pm\sqrt3,\lambda=2$ を代入すればよい）。
従って全ての相方と許容端点で、前節の長さ5の階数が開近傍に延びる。
$r^2=z$の有限被覆から悪い閉集合の像を除けば、この近傍は倍率平面上に取れる。
こうして残る全ての第一座標を同時に扱える。
\end{proof}

% !TeX root = ../papers/nc-p3.tex
\section{両固定曲線で四次と五次の障害点が一致すること}
\label{sec:uniform-positive-quintics}
\begin{theorem}\label{thm:uniform-fixed-quintics}
標数零の全ての通常倍率において、両固定曲線上の許容窓について次の三条件は同値である。
\begin{enumerate}
\item ある始点で長さ5の道空間の次元が56を超える。
\item ある始点で長さ4の道空間の次元が35を超える。
\item 命題\ref{prop:fixed-line-defects}の十二表示のいずれかであるか、
正の曲線上で $\det A+\det B=0$ である。
\end{enumerate}
従って各倍率での幾何点の個数は正で八点、負で八点
（$\lambda^2z=-1$ のとき四点）である。その補集合では全始点で長さ5の次元が56になる。
\end{theorem}
\begin{proof}
四次の判定は定理\ref{thm:uniform-fixed-quartic}による。
十二表示の道空間を $\mathbb Q(z,\lambda)$ 上で直接計算すると、
長さ5の四始点の次元を並べ替えた値は、正で $(68,68,68,68)$、
負で $(56,60,60,68)$ である。
\path{computations/verify_fixed_line_elliptic_curves.py} は各表示でこの値を検算する。
関係の挿入行列の階数は特殊化で増えないため、全通常倍率で少なくとも一始点は68以上になる。
新しい正の四表示にも定理\ref{thm:opposite-fixed-dets}の同じ下界がある。

逆向きは負では定理\ref{thm:varying-fixed-quintics}で示した。正の曲線で $b=1$ とし、
定理\ref{thm:uniform-fixed-quartic}の表の第一座標と通常倍率から除いた因子を逆にする。
$(e,f)=(Nu,Du)$ を代入して得る二次式を $G=g_2u^2+g_1u+g_0$ と書く。
この局所化上で $g_2$ は単元で、$v=g_2u$ とすれば
\[
 v^2=-g_1v-g_0g_2,\qquad
 \operatorname{Nm}(X+Yv)=X^2-g_1XY+g_0g_2Y^2 .
\]
命題\ref{prop:fixed-quintic-slice}の挿入行列に対し、始点0,1と文字0,2の各組で
86行86列を選ぶ。行と列の順序は
\path{research/fixed_quintic_pivot_certificate.json} に収録した。
分母を払った成分は $u$ について高々二次で、$g_2^2$ 倍すれば上式で $X+Yv$ に帰着する。
\path{computations/verify_uniform_positive_quintics.py} はこの形で各行列を掃き出す。
各ピボットのノルムを許した因子で割り尽くすと、344回全てで非零定数が残る。
従ってピボットはこの二次代数で可逆である。行の消去には共役とノルムを用い、
行から取り除く内容も全て単元であることを検査する。
よって四つの小行列式はいずれもこの開集合で消えない。
二位置移動と奇数長の同時文字交換により、残る始点・文字も目標階数86に達する。

表に残る第一座標のうち、四次を通過しうる許容点は同定理により、
二位置移動を除けば座標端の二表示と $H_\epsilon$ の既知の根の相方に尽きる。
$r^2=z$ として、これら四表示では係数を $r,\lambda$ の有理式で書ける。
\path{computations/verify_uniform_quintic_companions.py} は四表示の各四行列で
86次小行列式を求める。通常倍率の単元、定義域の分母、および
$xy\det A\det B(\det A+\det B)$ の因子を除き、その分子の各既約因子を $P$ とする。
別の小行列式の残余因子を $Q$ とすると、共通零点では少なくとも一つの $Q$ が消え、
その $r$ は $\operatorname{Res}_\lambda(P,Q)$ の根になる。
$Q$ が $r$ のみの式なら直接 $Q$ を使う。各終結式は非零である。
別の小行列式ごとに得た一変数既約因子の集合を交差すると、許した因子の外では全て空になる。
小行列を選ぶために使う素数307に依存せず、この因数分解と終結式の等式は整数上で成立する。
従って許容な相方と端点にも新しい五次障害はない。各文字の余核次元の下界14と合わせて結論を得る。
\end{proof}

% !TeX root = ../papers/nc-p3.tex
\section{正の固定曲線の端点・相方・八次閉点の全次数認証}
\label{sec:fixed-residual-regularity}
\begin{proposition}\label{prop:fixed-residual-regularity}
倍率 $(z,\lambda)=(3,2)$ の正の固定曲線上で、次の全22幾何点は反対階数16の
標準Koszul左右AS正則窓を与える。
\begin{enumerate}
\item $a=0$ と $b=0$ の二つの許容点。
\item $H_\epsilon=0$ の四つの第一座標の上で、$\det A+\det B=0$ を満たす点の相方。
\item 命題\ref{prop:fixed-quintic-slice}の小行列式の消去に残る二つの八次閉点の全共役。
\end{enumerate}
各始点の道空間の直和は四つの中心四次式の多項式環上で階数256の自由加群となる。
最初の二項の六点は、$(3,2)$ を含む倍率平面の共通の非空開集合上で同じ性質を保つ。
\end{proposition}
\begin{proof}
定理\ref{thm:uniform-fixed-quartic}の端点と相方の有理式を $r^2=3$ で読む。
相方の第一座標は二つの相異なる二次式の根であり、二つの端点から離れている。
193で $r$ は $14,179$ に単純分解する。この二値で相方の式を、端点では一方の値で式を還元する。
全ての分母と $xy\det A\det B(\det A+\det B)$ は剰余体で非零である。

八次閉点は同命題の検算が返す $(P(a),u-U(a))$ を重複なく取ると二つであり、
第一座標の既約八次式 $P$ も相異なる。各 $P$ の一つの単根を1009で取り、
$(a,U(a))$ の還元は $(60,14)$ と $(33,982)$ となる。
$P$ と $U$ の係数の分母はこの素数で可逆である。
Henselの補題で単根は持ち上がり、$\mathbb Q[a]/(P)$ の一つの付値で窓を還元できる。
証明すべき性質は体拡大と共役で保たれるので、各八次式につき一つの還元で足りる。

\path{computations/verify_fixed_residual_regularity.py} はこの八つの剰余窓で四つの関係空間を
作り、隣接階数6、反対階数16、四始点の長さ4,5の次元35,56と
四次中心の文字次元 $(3,2,0,0)$ を検算する。最初の二文字から各二式を選んだ商の次元は
全32組の窓と始点で、次数0から13まで順に
\[
 1,4,10,20,31,40,44,40,31,20,10,4,1,0
\]
となる。定理\ref{thm:fixed-line-elliptic}の持上げの議論を適用する。
すなわち経路と中心の次元は族の一般点から来る下界に等しく、
四中心は付値環上に持ち上がる。剰余商の次元上界と道空間の全次数の下界から
自由基底と左右Koszul完全性が従う。端点と相方の表示では分母の非消滅、経路と中心の階数、
商の有限性の証明に用いた小行列式がそれぞれ開近傍で保たれる。
$r^2=z$ の有限被覆で六点の近傍の補集合の像を除けば、所要の倍率の開集合を得る。
\end{proof}

% !TeX root = ../papers/nc-p3.tex
\section{固定曲線の反対階数と階数14の楕円曲線}
\label{sec:fixed-opposite-rank-drop}
\begin{theorem}\label{thm:fixed-opposite-rank-drop}
標数零の全通常倍率で、負の固定曲線の許容窓の反対切断は可逆である。
正の固定曲線では $A,B$ の第一行をそれぞれ $(a,b),(e,f)$ とすると、反対階数は
\[
 16-\boldsymbol{1}_{a^2+e^2=2zbf}-\boldsymbol{1}_{z(b^2+f^2)=2ae}.
\]
階数14,15の全窓は全始点で長さ4,5の標準次元を持つ。
階数14の窓が存在する必要十分条件は
\[
             (\lambda+1)^2(\lambda z+1)^2=16\lambda^2 z .
\]
この曲線上の各通常倍率で階数14の幾何点はちょうど四つある。
標数零の代数閉体上で、倍率も動かしたこの軌跡は二つの成分を持ち、
それぞれの滑らかな完備化は $Y^2=2(t-1)(t+1)(t+3)$ と同型で、$j=1728$ である。
\end{theorem}
\begin{proof}
反対切断を四文字のブロックに分ける。$b=1$、$\Delta=a^2-z$、
$(e,f)=u(N,D)$ と書く。文字2,3の行列式は両符号とも
$-(\lambda-1)^2(\lambda z-1)^2/\lambda^2$ である。
負の二つの残りの分子は、通常倍率の単元を除き
$(a^2-z^2u^2)^2,(a^2u^2-1)^2$ である。
曲線の二次式 $G$ と各因子の終結式を取ると、通常倍率の単元と
$a,\Delta$ の因子だけが残る。負の座標端は非許容である。

正の文字0では最後の行を $z$ で割り、文字1では二列目を $z$ で割る。
二つのブロックはともに
\[
 \begin{pmatrix}m&n&1&1\\ n&m&1&1\\1&1&C&D'\\1&1&D'&C\end{pmatrix}
\]
となる。文字0で $(m,n)=(a,e)$、文字1で $(m,n)=(b,f)$ である。
反対称部分の二つの係数は $m-n,\ k(m-n)$ で、$b=1$ では
$k=-1/(\lambda u z\Delta),-1/(\lambda u\Delta)$ である。
$m-n$ と $G$ の終結式の因子は定理\ref{thm:uniform-fixed-quartic}の
座標端とその交換像、$\Delta$、通常倍率の単元に限られる。
端点の公式では $a-e,b-f$ の分子は通常倍率の単元であり、行列交換でも非消滅を保つ。
従って反対称部分は可逆である。対称部分の二次の行列には非零の成分2があり、
その行列式は順に $-k(a^2+e^2-2zbf)$ と $-k(b^2+f^2-2ae/z)$ になる。
斉次表示で $\Delta=a^2-zb^2$ に置き換えれば、これらの恒等式は座標端にも延びる。
古い四次障害の正の四表示に二つの条件式を代入すると、いずれも通常倍率の単元である。
反対行列式の四つの障害点も定理\ref{thm:opposite-fixed-dets}で反対階数16である。
従って定理\ref{thm:uniform-fixed-quintics}は階数14,15の全点に標準次元を与える。
以上のブロックと終結式は \path{computations/verify_varying_opposite_rank.py} が検算する。

両条件が成り立つのは $(a+e)^2=z(b+f)^2$ かつ
$(a-e)^2=-z(b-f)^2$ のときである。$a+e=b+f=0$ は
$B=-A$ となり、固定点のない二位置移動に反する。従って $r^2=z$ の符号を
$a+e=r(b+f)$ により選べる。$i^2=-1$ のどちらかを取って
$a-e=ir(b-f)$ とし、$X=a/r$, $w=\lambda r$ と置く。この二式から
\[
 f=F=(1+i)X-ib,\qquad e=rE=r\big((1-i)b+iX\big).
\]
交換の比 $[e:f]=[N:D]$ は次の二次形式を与える。
\[
 K=E\big((w^2+1)X-2wb\big)-F\big(2wX-(w^2+1)b\big).
\]
$T=rw^2+wr^2-4rw+w+r$ とする。二つの補完条件
$f=f_0+f_1(a^2-zf^2)$、$e=e_0+e_1(e^2-zb^2)$ は、この代入後に共通分母
$w(r^2+1)(X+b)$ を持ち、分子は二次形式となる。分母が零になる許容点はない。
\path{computations/verify_fixed_rank_fourteen.py} は、両分子と $K$ の二変数斉次終結式が
$(w+1)^2T^2$ と、それぞれ
\[
  -2i\big(w^2+(-2+2i)w+1\big),\qquad
    2ir^2\big(w^2+(-2-2i)w+1\big)
\]
との積であることを検算する。括弧内の差は $4iw\ne0$ なので $T=0$ が必要である。
逆に $T=0$ では各分子と $K$ の係数行の全二次小行列式が消える。
$K$ の判別式は $16w(w-1)^2\ne0$ であり、二つの根は全補完条件を満たす。
実際、正の固定直線では $g_1=zf_1,h_1=e_1$ で残りの線型式も従い、
交換の比が $zf_1f+e_1e=0$、従って二つの二次式の導関数値の和が零となる対根を選ぶ。

許容性も全ての倍率で確認できる。$K$ と $X^2-b^2,E^2-F^2$ の終結式はともに
$-4(w^2-1)^2$、$X^2-F^2,E^2-b^2$ との終結式はともに
$-4(w-1)^2(w^2-6w+1)$ である。$T=r(w^2-6w+1)+w(r+1)^2$ なので、
最後の因子が零なら $r=-1$ となり通常性に反する。これらは斉次終結式なので無限遠の根も含む。
$T=0$ は $(\lambda+1)(\lambda z+1)=4\lambda r$ と同値であり、
倍率から $r$ が一意に定まる。二つの $i$ に共通する点は
$a=e,b=f$ となるため存在しない。これで各倍率の四点を得る。

倍率の曲線は
\[
 r=\frac{-2(t+3)}{t^2-1},\qquad \lambda=-\frac{(t-1)^2}{(t+3)^2}
\]
で有理に表示できる。逆は $t=(2r^2\lambda+r^2-5r+2)/(r(r-1))$ である。
各 $i$ で二次式 $K$ の根を添加することは、その判別式から $\sqrt w$ を添加することと等しい。
$w=2(t-1)/((t+1)(t+3))$ は平方でなく、$Y=(t+1)(t+3)\sqrt w$ が主張の滑らかな三次式を満たす。
$t+1$ を変数に取ると右辺は $2((t+1)^3-4(t+1))$ であり、$j=1728$ が従う。
\end{proof}

\begin{proposition}\label{prop:fixed-rank-drop-regularity}
倍率 $(3,2)$ で正の曲線の反対階数が16を下回るのは八幾何点で、全て階数15の標準Koszul左右AS正則窓である。
$(36,-1/4)$ では正の曲線の階数低下点は四幾何点で、全て階数14の標準Koszul左右AS正則窓である。
前者は倍率平面の非空開近傍で、後者は上の倍率曲線上の非空開近傍で成立する。
従って階数14の各楕円成分で全次数の正則性が未確定な許容点は高々有限個である。
\end{proposition}
\begin{proof}
正の二つの行列式を $G$ と消去し、許容な座標端を別に検算する。
$(3,2)$ の残余第一座標の多項式は文字0,1の順に
\[
 932a^4-8229a^3+18228a^2-5265a-576,\quad
 64a^4+1755a^3-18228a^2+24687a-8388 .
\]
$(36,-1/4)$ では両者とも $769a^4+3744a^3-10296a^2+134784a+996624$ となる。
各四次式は既約で、その根の上の該当点は $u=U(a)$ と一意に定まる。
後者では二つのブロックが同じ $U$ を与える。
\path{computations/verify_fixed_opposite_rank_drop.py} はこれらの等式と二座標端の可逆性を確かめ、
各閉点で101を上に持つ単純な剰余根を一つ取る。
その $(a,u)$ は順に $(40,4),(33,8),(15,71)$、反対階数は $15,15,14$ である。
消えるブロックが与える上界と剰余階数が一致する。
全ての剰余点で四次中心の文字次元は $(3,2,0,0)$、四式で割った各始点の商は
命題\ref{prop:fixed-residual-regularity}と同じ次数13までの次元を持つ。
同じ中心の持上げと自由基底の証明で、全共役の全次数の正則性が従う。

倍率を動かした終結式の残余は両文字で各四次式一つで、いずれも重複度1である。
$(3,2)$ への特殊化は次数を保ち、二式は分離的かつ互いに素である。
従って開近傍で有限エタールな八根を持つ。各根の上で共通の $u$ を与える一次の
部分終結式の最高次係数と、許容性および座標端の行列式はこの近傍を縮めれば零にならない。
$(36,-1/4)$ では倍率曲線上の四根が前定理で有限エタールであり、各四次式が全根を占める。
両場合とも有限被覆で正則性の証明に使った小行列式の非消滅が開であるから結論を得る。
\end{proof}

% !TeX root = ../papers/nc-p3.tex
\section{階数14の二位置移動による降下}
\label{sec:rank-fourteen-descent}
\begin{proposition}\label{prop:rank-fourteen-descent}
定理\ref{thm:fixed-opposite-rank-drop}の階数14の二成分は、それぞれ通常倍率の全域で
二位置移動により有理な $t$ 直線の開集合に降りる。
各成分の中心の階数と四中心による有限商の検査は、一変数の局所化多項式環上の行列で実行できる。
\end{proposition}
\begin{proof}
$i^2=-1$ を固定し、前定理の $w=\lambda r$ に対し $s^2=w$ と置く。
第一行列の第一行を $(rX,b)$ とする。全域で使える表示は
\[
 X(s)=i\big(s^2-(1+i)s-1\big),\qquad b(s)=s^2+(1-i)s-1 .
\]
この対は $K(X,b)=0$ を恒等的に満たす。$X/i-b=-2s$ で、$s=0$ でも同時には消えない。
$(E,F)=((1-i)b+iX,(1+i)X-ib)$ を代入すると
$(E(s),F(s))=(-X(-s),-b(-s))$ となる。従って $B(s)=-A(-s)$ である。
二位置移動は二つの行列を交換するため、共通スカラーを除き $s\mapsto-s$ と一致する。
この二次拡大の固定体は $\mathbb Q(i)(t)$ であり、各成分の二点のAS性は等しい。

通常倍率を与える $t$ の開部分の座標環を $R$ とする。$r=0,\infty$ は非通常なので
この開部分はアフィン直線に含まれ、$R$ は $\mathbb Q(i)[t]$ の局所化である。
$w$ は単元で、$S=R[s]/(s^2-w)$ は階数2の有限エタール代数となる。
上の表示の行列式は、単元 $r^2$ を除き
\[
 X^2-b^2=-2(s^2-1)(s-i)^2,\qquad E^2-F^2=-2(s^2-1)(s+i)^2
\]
である。いずれの因子も通常倍率で消えず、関係空間を表す元のテンソルを $S$ 上で構成できる。
従って $b=0$ を別の座標で張り替える必要もない。$c+ds$ による乗算は $R$ 基底 $(1,s)$ では
\[
 \begin{pmatrix}c&wd\\d&c\end{pmatrix}
\]
である。これで経路空間の関係の各挿入、中心の交換方程式、選んだ中心式の挿入を
$R$ 上の行列に移せる。道空間の商を作るときも、表示から核や余核を取ればよい。
$R$ は単項イデアル整域なので、標準次元が既知の長さ4,5の道空間は
$R$ 加群として有限自由である。階数条件は小行列式の最大公約数で判定できる。
幾何ファイバーで制限した行列は二つの $s$ に対応する行列の直和になり、
全始点での検査は二位置移動で同値になる。
\end{proof}

\begin{proposition}\label{prop:rank-fourteen-residue}
階数14のこの表示を標数零の離散付値環上で取り、剰余体を $\mathbb F_{101}$ とする。
その還元が通常な $\mathbb F_{101}$ 有理点ならば、対応する分数体上の窓は
標準Koszul左右AS正則である。該当する剰余点は $i$ の二値で各96点ある。
\end{proposition}
\begin{proof}
101では $i=10,91$ である。\path{computations/verify_rank_fourteen_descent.py} は
全ての $t$ と $s^2=w(t)$ の全根を列挙する。$t=\infty$ では $r=0$ となり通常性に反する。
許容な192点の全てで隣接階数6、反対階数14、経路長4,5の次元35,56と四次中心の文字次元
$(3,2,0,0)$ を検算する。最初の二文字から二式ずつ選ぶ。
四つの剰余点では文字0の二つの線型結合を取り直すと、
全192点・全四始点で商次元が命題\ref{prop:fixed-residual-regularity}の表に一致する。
同命題の付値環上の持上げにより、各剰余点に還元される窓で全次数の結論を得る。
\end{proof}

% !TeX root = ../papers/nc-p3.tex
\section{階数14の全成分上の四次中心}
\label{sec:rank-fourteen-center}
\begin{theorem}\label{thm:rank-fourteen-center}
標数零で、二本の階数14成分の全ての通常点で四次中心の文字次元は
$(3,2,0,0)$ である。有理な倍率座標 $t$ の上で四つの既約閉点
（次数 $8,8,20,34$）を除けば、関数体の消去で得た五つの中心元が各ファイバーの基底を与える。
この四閉点の上の全280幾何点は標準Koszul左右AS正則である。
\end{theorem}
\begin{proof}
$K=\mathbb Q(t)(i,Y)$、$i^2=-1,\ Y^2=C=2(t-1)(t+1)(t+3)$ と置く。
前節の大域表示により、二位置移動の二根と $i$ の二値を同じ有理係数の計算で扱える。
\path{computations/rank_fourteen_function_field.py} は $K$ の元を
$(A+Bi+DY+EiY)/q$、$A,B,D,E,q\in\mathbb Z[t]$ と表示し、共通因子を割る。
分子の $\mathbb Q(t)$ へのノルムは
\[
 \{A^2-B^2-C(D^2-E^2)\}^2+4(AB-CDE)^2 .
\]
この計算では数値補間を用いない。

始点・長さ・文字ごとに、直前の商基底の左に一文字を付け、
始端の二次関係に全ての残る語を付けて行消去する。この規則は全ての関係の挿入を含む。
自由列の語への射影を長さ5まで記録する。
\path{computations/verify_rank_fourteen_center.py} はこの既約分母を全て因数分解する。
通常倍率の定義で逆にした因子を除くと、残る極の第一座標の式はただ一つである。
\[
 t^8+8t^7+60t^6+248t^5+518t^4+568t^3+828t^2+1096t+769 .
\]
その外では各射影は定義され、関係を零にし、選んだ語を単位基底に送る。
四次・五次の道空間の次元は全点で既知なので、この全射はファイバーごとに同型となる。
三次以下でも同じ結論が成り立つ（固定曲線で標準の三次交叉を持つ）。
従って掃き出しの途中のピボットが消えても、最終射影が正則な点は取り除く必要がない。

各文字 $h$ の中心条件を、四始点の四次の語 $v_p$ に対する
$v_p x_j-x_jv_{p+1}=0$ の五次座標で記す。各行列は224行である。
独立な行を順に選び、以前のピボット列を零にした後に新たな列を一つ選ぶ。
選んだ行・列の順序では、ピボットの積はこの小行列式そのものである。
その積を体内で簡約してからノルムの分子を因数分解すると次の表を得る。
最終列には各小行列式について通常倍率の単元を除いた相異なる既約因子の次数を記す。
\[
\begin{array}{c|r|r|l}
 h&\text{列数}&\text{階数}&\text{因子の次数}\\ \hline
 0&36&33&8,8,20\\
 1&40&38&8,34\\
 2&32&32&8\\
 3&32&32&8
\end{array}
\]
表の和集合は上記の極を含む四因子である。各行列の一般階数は特殊化での上界となり、
選んだ小行列式が可逆な各点では等号である。余った列を自由変数に取ると、Cramerの公式で
この開集合上に中心の基底が得られる。

残る四因子については、各既約式の単純な素体根を一つずつ取る。
定数項193と769の八次式について $(p,t)=(157,141),(101,89)$、
二十次式と三十四次式については $(101,45),(101,85)$ を用いる。
全て通常な還元で、$i^2=-1$ と $s^2=w(t)$ はともに相異なる二根を持つ。
各四組で同じ中心次元を得る。また四つの中心四次式で割った商の次元は、
全始点で命題\ref{prop:fixed-residual-regularity}の次数0から13の表に一致する。
単根を持ち上げ、四組の拡大を全て比較したので、既約な $t$ の全共役が検査される。
同命題の定階数と中心持上げによって、これらの全点で全次数の正則性を得る。
\end{proof}

% !TeX root = ../papers/nc-p3.tex
\section{階数14の四中心商と六次の道空間}
\label{sec:rank-fourteen-central-cut}
\begin{theorem}\label{thm:rank-fourteen-central-cut}
標数零の階数14の二成分上で、各通常幾何点において、文字 $0,1$ から二本ずつ
四次中心元を選び、全四始点で商の長さ $0$ から $6$ の次元を
$1,4,10,20,31,40,44$ にできる。特に割る前の長さ6の道空間はどの始点でも次元84である。
以下の計算で生じる $t$ 直線上の22既約閉点の上では、全4768幾何点で
標準Koszul左右AS正則性が成立する。
\end{theorem}
\begin{proof}
前定理の中心行列の自由列を左から数え、文字0の第一・第三自由変数と
文字1の二自由変数に対応する四つの核ベクトルを取る。
\path{computations/verify_rank_fourteen_central_cut.py} は始点 $p=0,1$ ごとに
語基底の右に一文字を付ける。長さ $n$、文字 $h$ の列は、末尾 $b$ ごとに
長さ $n-1$、文字 $h+b$ の残存語を並べる（加法は四元群のもの）。
長さ $n-2$ の残存語の後ろに各二次関係を付け、最後から二番目の文字までを
直前の射影で展開して一行を得る。長さ4では選んだ中心式も行に加える。
これで中心式の全挿入を扱える。実際、四次中心元の各位置の成分は直前の矢と
交換でき、途中で始まる四次関係は先頭へ移してから右に矢を付けて得られる。

$(\ell,t,i,Y)=(101,2,10,38)$ の剰余点で、短い行から順に独立な行を選ぶ。
選んだ行について、ピボット列・行列式・自由列への射影は前定理の関数体で改めて計算する。
以前の行の関係を使い、残った非零行の最左列を順に消せば、
ピボットの積は選択小行列式に符号を掛けたものである。
この行列式が可逆なら、Cramerの公式で新しい自由列は全ての語を張る。
各段で一部の独立行だけを使っても、得られる語数は商次元の上界となる。
各始点で長さ $0$ から $6$ の上界は定理の列になる。
関係・中心式の分母、経路基底の極、および選択小行列式のノルムの分子と分母から
非通常因子を除くと、相異なる有理既約因子の次数は
\[
\begin{gathered}
8,8,8,14,16,16,20,22,24,24,34,34,\\
52,52,76,100,106,108,112,114,116,128 .
\end{gathered}
\]
この集合の外では、二位置移動と $Y$ の符号交換により残る二始点でも
同じ上界を与える切断が存在する。各始点でこの上界を満たす選択は
$\operatorname{Gr}(2,Z_{4,0})$ の非空開集合を含む（文字1では空間全体を取る）。
Grassmann多様体は既約なので、四つの開集合の共通部分から一組を取れる。

22因子については、それぞれの単純な剰余根を検算コードに指定した。
各根で $i$ と $s$ の全四組、全四始点について四中心商の長さ13までを再計算し、
命題\ref{prop:fixed-residual-regularity}の列を得る。従って持上げと共役により、
次数の和1192の四倍の幾何点を全次数で認証した。

最後に、割る前の次元には一般正則点から $\binom{n+3}{3}$ の下界がある。
各始点の商の斉次基底を持ち上げれば、それらは次数4の四変数多項式環上で
道空間を生成する。長さ6までは、上の商次元の上界を
$(1-u^4)^{-4}$ に掛けると各係数がこの下界に一致する。
従ってこの範囲では商と道空間の双方で等号が成立する。
\end{proof}

% !TeX root = ../papers/nc-p3.tex
\section{階数14の全幾何剰余点と円分倍率の正則性}
\label{sec:rank-fourteen-geometric-reduction}
\begin{theorem}\label{thm:rank-fourteen-geometric-reduction}
標数零の階数14の窓が離散付値環上の通常表示を持ち、その還元が標数101の通常点ならば、
分数体上で標準Koszul左右AS正則である。剰余座標 $t,s$ の拡大次数に制限はない。
特に、有理既約な原始整数多項式 $P(t)$ の101での還元が通常な幾何根を一つ持てば、
$P$ の全共役の上の全ての窓が正則である。この根には単根性を仮定しない。
\end{theorem}
\begin{proof}
$k=\mathbb F_{101}$、$R=k[t]$、$C=2(t-1)(t+1)(t+3)$、
$S=R[Y]/(Y^2-C)$ と置く。$i=10,91$ のそれぞれに対し、前節の四式を
分母が可逆な点で使う。二次関係と四次中心式は一行ずつ共通分母を払って $S$ 上に記す。
以下の計算は全ての $t$ を一変数のまま行う。

始点 $p$ を固定する。長さ $n$、文字 $h$ の語に $S$ の値を与え、
関係を零にする線型汎関数を考える。$R$ へ係数を制限すると、空語の値の空間は
階数2で、$Y$ の作用は $M_0=\left(\begin{smallmatrix}0&C\\1&0\end{smallmatrix}\right)$
である。最後の文字 $b$ を指定した制限写像を $J_{n,h}^{b}$ と書く。
長さ $n-1$ の四つの空間の直和の元を $(v_b)_b$ とする。
末尾に置く文字 $g$ の二次関係 $\sum c_{ab}x_ax_b$ に対し、条件は
\[
 \sum_{a,b} c_{ab}(M_{n-2,h+g})\,J_{n-1,h+b}^{a}\,v_b=0 .
\]
ここで $c=u+vY$ に対する $c(M)$ は $uI+vM$ である。
長さ4では各中心式にも同様に三つの制限写像の積を用いる。
後続の長さでは二次関係だけを加える。この条件だけでも元の商の双対を含み、
分数体で四式が中心なら前節の交換則により全条件となる。

得られた行列を $A$ とする。$R$ は単項イデアル整域なので、
多項式上の核の基底を並べた $J$ が取れる。その四つの行ブロックが次の制限写像となる。
$Y$ 作用を直和した行列 $M'$ に対し $JM=M'J$ を解き、全成分で確認する。
$A$ の列加群の基底を並べた行列を $B$、その行加群の基底を並べた正方行列を $D$ とする。
$r=\operatorname{rank}A$ とすれば、$A$ の $r$ 次小行列式のイデアルは $(\det D)$ である。
実際、生成列と基底列は互いに $R$ 線型結合で書け、行についても同様なので
Cauchy--Binet公式を両方向に適用できる。したがって $\det D$ が単元の局所環では
核と特殊化が交換する。一変数の被約Gröbner基底は先頭成分が相異なるため、
ここで必要な各部分加群の自由基底を与える。

\path{computations/rank_fourteen_dual.py} は始点 $0,1$ と各 $i$ で
長さ2から13の全四文字についてこの計算を行う。各長さの核の総階数は
命題\ref{prop:fixed-residual-regularity}の商次元の二倍である。
$C$ は $k(t)$ で平方ではなく、一般核は $k(t)(Y)$ 上の空間である。
通常点では二根が相異なる。よって定階数の開集合で $Y$ の二つの固有空間は
それぞれ一般階数の半分を持つ（有限エタール被覆上で各直和因子の階数は局所一定）。
二位置移動で全四始点が扱え、共通の四式の選択には前節と同じ開集合の交叉を使える。

$\det D$ の既約因子と四式および関係の分母を合わせ、その全根も調べる。
\path{computations/verify_rank_fourteen_geometric_reduction.py} は
各通常既約因子の次数を $d$ として $\mathbb F_{101^d}$ に根を取り、
必要なら二次拡大して $Y$ を取る。Frobenius共役と二位置移動により、
各 $i$ で一つの $t,Y$ と全四始点を検査すれば足りる。各点で経路長4,5の次元
35,56、中心の文字次元 $(3,2,0,0)$ を直接求め直し、その中心から選んだ
四式の長さ13までの商次元が所定の列に一致することを確かめる。

分母も行列式も消えない還元では有理式で指定した中心元が分数体上で使え、
還元で得た商次元は分数体での上界を与える。例外因子での直接検査では
定階数の経路・中心を経由して剰余切断を持ち上げる。
いずれも一般正則点からの全次数の下界を合わせれば
命題\ref{prop:fixed-residual-regularity}の議論が適用される。
最後に $P$ を射影直線上で原始化して閉包を取れば、101でのその幾何根は
正規化上の付値に支配される。そこで数体を必要なだけ拡大すれば $i,Y$ の全ての選択ができ、
同じ結論が共役で全ての根へ移る。
\end{proof}

\begin{corollary}\label{cor:rank-fourteen-cyclotomic}
階数14の二成分で、$t$ が1の冪根である全ての通常窓は、標数零で
標準Koszul左右AS正則である。
\end{corollary}
\begin{proof}
$t$ の位数を $101^a m$、$101\nmid m$ とする。101における共役の還元は
全ての原始 $m$ 乗根を尽くす。非通常な有限 $t$ は101で相異なる13点であり、
そのうちに現れる乗法位数は $1,2,25,51,100,3400,10200$ である。
この集合の $m>2$ は全て $\varphi(m)>13$ を満たす。従って $m>2$ なら
原始根の一つが通常であり、定理が適用される。

$m=1,2$ で $a=0$ のとき $t=\pm1$ は非通常である。$a>0$ のときは
素数109を使う。$109$ の101を法とする位数は100なので、原始 $101^a m$ 乗根の
109での剰余次数は100の倍数となる。$i=33$ として上の多項式核と行列式を
二始点で計算すると、分母と非通常因子を含む全例外因子の次数は100未満である。
\path{computations/verify_rank_fourteen_cyclotomic.py} が両素数での因数分解と位数を再現する。
$\mathbb Q(t)$ は奇数導手の円分体で2で不分岐だから $\mathbb Q(i)$ と交わるのは
$\mathbb Q$ だけである。従って $i$ のどちらの選択にもこの還元を使え、
$Y$ の両根は上と同じ核計算で扱われる。
\end{proof}

% !TeX root = ../papers/nc-p3.tex
\section{階数14の斉次表示における全境界テンソル}
\label{sec:rank-fourteen-boundary}
\begin{theorem}\label{thm:rank-fourteen-boundary}
命題\ref{prop:rank-fourteen-descent}の斉次表示を固定し、二本の楕円曲線上で
四階テンソルの射影極限を取る。通常域の外の48幾何パラメータのうち、
次の三式の根の上の32点は標準Koszul左右AS正則で反対階数14である。
\[
 t^2-2t-7,\qquad t^2+6t+1,\qquad t^4+2t^2+24t+37 .
\]
残る16点の極限テンソルは標準14641型を与えない。
また通常な標数零の窓が非正則なら、その全ての101進還元で倍率座標は
$\infty,1,-1,-3,19,80$ のいずれかである。
\end{theorem}
\begin{proof}
薄い成分は元の表示のままにし、厚い二成分を
$\left(\begin{smallmatrix}A&X\\Y_0&Y_0A^{-1}X\end{smallmatrix}\right)$ と
$\left(\begin{smallmatrix}B&\lambda X\\Y_0/\lambda&Y_0B^{-1}X\end{smallmatrix}\right)^t$
で作る。ここで $X=\left(\begin{smallmatrix}1&1/\lambda\\1/\lambda&z\end{smallmatrix}\right)$、
$Y_0=\left(\begin{smallmatrix}1&\lambda\\\lambda z&1\end{smallmatrix}\right)$ である。
\path{computations/verify_rank_fourteen_boundary.py} はこれを楕円関数体で簡約する。
全成分の分母の既約因子は $t-1,t+1,t+3$ だけである。特に $A$ または $B$ が
特異になっても、積全体が正則ならその値から関係を回復できる。

通常倍率の非消滅条件を $t$ で因数分解すると、三つの極と無限遠以外では
定理の三式と $t^2+2t+5$ が尽くす。定理の三式は101で順に
\[
 t^2+99t+94,\quad t^2+6t+1,\quad
 (t^2+81t+40)(t^2+20t+59)
\]
となり、各二次式は既約である。各根で、必要なら剰余体を二次拡大し、
$i,Y$ の全四通りについて生のテンソルの隣接四切断と反対切断の階数を求める。
結果は $6,6,6,6$ と14で、各一文字切断は異なる文字に異なる列支持を持ち、階数4である。
四つの隣接切断から直接得た関係で、経路次元35,56、中心次元 $(3,2,0,0)$、
四中心商の長さ13までの列を全始点で確認した。

標数零でもテンソルの切断階数には一般階数による上界とこれらの剰余階数による
下界がある。従って関係空間は近傍でベクトル束をなし、一般正則点から全次数の経路次元下界を得る。
長さ4,5の余核が定階数となり、中心の下界と剰余体での上界も一致するので、
前定理の中心持上げと左右Koszul完全性の証明が適用される。
同じ議論はこの還元を持つ全ての通常な標数零の窓にも使える。
$t^2+2t+5$ の101での根は19と80であるから、前定理との合併で最後の主張を得る。

$t=-1+2i$ とその共役では $s=\pm1$ を代入し、$\mathbb Q(i)$ 上で計算すると
隣接四切断と反対切断は全て階数4である。最後の四基点は二次被覆の分岐点である。
有限の $a$ では $t-a=u^2$、無限遠では $t=u^{-2}$ と置く。
有理部分と $Y$ 倍の部分は位数の偶奇が異なる。全係数の最小位数を抜くと、
同じ位数の $Y$ の先頭係数は共通スカラーなので射影化で除ける。
次表を $\mathbb Q(i)$ 上の掃き出しで得る。
\[
\begin{array}{c|c|c}
 t & \text{四つの隣接階数} & \text{反対階数}\\\hline
 -3 & 2,1,2,1&2\\
 -1 & 2,2,2,2&1\\
  1 & 6,6,6,6&6\\
 \infty&2,2,2,2&1
\end{array}
\]
$t=1$ では三文字の関係挿入の余核次元が始点順に $28,28,22,22$ となり、
全て標準値20を超える。他の三基点と $t^2+2t+5$ の根では既に二次関係数が不足する。
非分岐の十個の $t$ では $i,Y$ の四組があり、分岐四点では各 $i$ に一つの極限となる。
これは指定した斉次表示を射影テンソル空間に延長したときの全境界を数えている。
\end{proof}

% !TeX root = ../papers/nc-p3.tex
\section{生成元の重みによる四つの分岐剰余類の正則性}
\label{sec:rank-fourteen-ramified-disks}
\begin{theorem}\label{thm:rank-fourteen-ramified-disks}
標数零の通常な階数14の窓は、$t$ の101進還元の一つが $19,80$ 以外なら
標準Koszul左右AS正則である。無限遠への還元も含む。従って未判定の有理閉点の
原始既約式 $P(t)$ の最高次係数は101で単元であり、
\[
 \overline{P}(t)=c(t-19)^a(t-80)^b,\qquad c\ne0,\quad a+b=\deg P
\]
が必要である。
\end{theorem}
\begin{proof}
前定理の四つの分岐基点 $d=-3,-1,1,\infty$ を扱う。
$101$ 進付値環を有限拡大して $t-d=w^{18}$、無限遠なら $t=w^{-18}$ とする。
四位置の各文字 $a$ に対し、テンソルのその一文字に作用する対角変換を
$w^{-e_{p,a}}$ と取る。次表の各欄は文字順の四つの重みである。
\[
\begin{array}{c|r|r|r|r}
d &p=0&p=1&p=2&p=3\\\hline
1&(-9,-9,1,-17)&(4,-14,12,12)&(2,2,-6,-24)&(-15,3,-7,-7)\\
-1&(-7,-7,3,-15)&(-3,-21,5,5)&(-3,-3,7,-11)&(-23,-5,3,3)\\
-3&(10,10,-7,11)&(-12,6,-13,-13)&(-1,-1,0,18)&(3,-15,2,2)\\
\infty&(8,8,-9,9)&(-12,6,-13,-13)&(9,9,-8,10)&(13,-5,-6,-6)
\end{array}
\]
分岐一様化変数 $u=w^9$ での各非零成分の位数 $v$ と先頭係数を前節と同じ式で取る。
\path{computations/verify_rank_fourteen_ramified_disks.py} は全成分で
$9v-\sum_p e_{p,a_p}\ge0$ を確認する。分母の先頭係数は101で単元なので、
この不等式は $w$ の正の付値の大きさによらず変換後の成分の整性を与える。
等号は各基点で32成分に限られ、いずれも $Y$ を含まない項が先頭である。
従って $Y$ の符号にも $w$ の選択にも依存せず、表から剰余テンソルを得る。

$i=10,91$ の各値で、変換後の剰余テンソルの一文字切断と隣接切断は
全て階数4と6を持つ。四始点で経路次元35,56、中心の文字次元 $(3,2,0,0)$、
四中心商の長さ13までの標準次元を確認する。分数体上では可逆な対角変換であり、
文字を保つため中心次元も保存する。一般正則点からの下界と還元上界による
定階数および中心持上げを適用すれば、元の分数体上の窓で全次数の主張を得る。
これと前定理で残る還元は19と80だけである。原始化した $P$ の斉次化で
最高次係数が消えれば無限遠に根を持ち、それも認証されるため、最後の必要条件が従う。
\end{proof}

\begin{proposition}
残る二剰余類では、上の文字基底の対角変換だけでは隣接階数6の還元を作れない。
六つの二次関係の平坦極限は長さ3,4,5で各始点の次元 $20,38,72$ を与える。
\end{proposition}
\begin{proof}
$t=-1+2i$ とその共役で $s=\pm1$ を代入する。四階テンソルの全64個の
文字許容成分はどれも101進単元で、他の成分は曲線上で恒等的に零である。
前節によりこの境界テンソルの隣接階数は4である。
同じ還元を持つ窓の各非零成分と対応する境界成分との比は付値環で1に還元される。
どの位置別対角変換を両者に施してもこの比は変わらず、共通の最小位数を払った後も
還元は一致する。定数境界テンソルの対角変換では階数は4以下にしかならない。

一方、各文字の右端ブロックを単位行列にした一般の関係式は、この二点の上で
全係数が正則である。そのまま特殊化した六行は独立で、関係六平面の極限を与える。
同じ検算コードで $\mathbb Q(i)$ 上、各文字の両符号と両共役について
これらの関係の全挿入を消去すると、主張の次元を得る。
\end{proof}

% !TeX root = ../papers/nc-p3.tex
\section{反対階数14の二楕円成分の全通常点でのAS正則性}
\label{sec:rank-fourteen-all-ordinary}
\begin{theorem}\label{thm:rank-fourteen-all-ordinary}
標数零で定理\ref{thm:fixed-opposite-rank-drop}の二つの階数14成分に属する
全ての通常窓は、標準Koszul左右AS正則である。各始点の道空間は四次の中心元四本が
生成する多項式環上で階数256の自由加群である。特に、同定理の倍率方程式を満たす
各通常倍率の上の四つの階数14窓は全て正則である。
\end{theorem}
\begin{proof}
定理\ref{thm:rank-fourteen-ramified-disks}により $t$ が101で19または80に還元される
場合を扱えばよい。残る二値は $t^2+2t+5$ の相異なる根である。
$J^2=-1$、$\alpha=-1+2J$ と置き、$J$ の二つの101進埋込みを取る。
係数環を $D=\mathbb Z_{(101)}[J]$、局所変数を $u=t-\alpha$ とする。
成分の定数 $i$ には $\pm J$ を、$s$ の境界値には $\pm1$ を独立に取る。
$Y^2=2(t-1)(t+1)(t+3)$ の各解は、定数項
$Y_0=\pm(\alpha+1)(\alpha+3)$ から $D[[u]]$ に一意に持ち上がる。
実際 $2Y_0$ は単元である。通常な持上げでは $u$ の付値は正で、$u\ne0$ である。

定理\ref{thm:rank-fourteen-central-cut}で使った二次関係と四中心式を取る。
分母を払う際は中心元の全四位置で共通の倍率を取る。得られる多項式切断は中心である。
実際、長さ5の道空間は通常域で定階数であり、一般点で零の交換子は各点でも零である。
これらの切断を $t=\alpha+u$ に代入する。各始点・長さ・文字で前節までと同じ
右への一文字追加を行い、列に対応する語を定数項が単元のピボットで消す。
両複素埋込みを同時に扱うため、定数項のガウスノルムが101で単元かを検査する。

ピボットが見つからないときが境界の階数低下に対応する。その際は残る全行の
\emph{定数項が $\mathbb Q(J)$ 上で零}であることを直接確認し、各行を $u$ で一度割る。
これは $D[[u]][1/u]$ 上で有効な関係式である。各段の自由列が目標数まで減ったら
後退代入し、全列から残存語への射影を作る。一部の行を捨てても残存語は商を張る。
中心式の挿入は長さ4で行い、その後は矢との交換で全ての挿入を含む。

\path{computations/verify_rank_fourteen_all_ordinary.py} はこれを有限の厳密計算で行う。
初めに係数を $u^{27}$ を法として求める。単元の逆は同じ精度で計算でき、
$u$ による除算一回だけがその行を一桁短くする。ある文字ブロックで $k$ 回の
一斉除算を行うと、後退代入した全列も少なくとも元の精度から $k$ を引いた精度を持つ。
次の長さでは四ブロックの下限の最小値を使う。各除算時の零判定と単元判定には
定数項が完全に決まる精度が残る。このため、打切り展開から無限級数への外挿を要しない。
全16組（$i$、$s$ の符号と四始点）で長さ13まで正の精度が残り、商次元の上界は
\[
 1,4,10,20,31,40,44,40,31,20,10,4,1,0
\]
となる。各証明書の行操作は完備付値環で $u$ にその値を代入でき、必要な $u$ の逆だけを
分数体で取る。従って還元が19または80の通常窓にもこの上界が成り立つ。

未切断の道空間の各次数には一般正則点のHilbert次元が下界を与える。
四中心商から持ち上げた自由加群の上界がこの下界に一致し、
定理\ref{thm:klein-rational-period-four}の左右Koszul複体の証明が適用できる。
数体上の幾何点は101上の付値に埋め込み、体拡大による完全性を降下する。
超越的な $t$ はすでに正則な幾何学的一般点からの体拡大である。
\end{proof}

% !TeX root = ../papers/nc-p3.tex
\section{反対階数15の二つの四重被覆とその分岐曲線}
\label{sec:rank-fifteen-covers}
\begin{theorem}\label{thm:rank-fifteen-covers}
正の固定曲線の通常倍率平面上で
\[
 k=\frac{(\lambda+1)(\lambda z+1)}{4\lambda z},\qquad
 m=-\frac{(\lambda-1)(\lambda z-1)}{4\lambda z},\qquad
 P(n)=n^4+(1-zk^2)n^2+zm^2
\]
と置く。反対階数が低下する全窓は、斉次座標 $[a:b:e:f]$ で次の二表示に尽きる。
各表示で $P(n)=0$ とし、
\[
 (a,b,e,f)=(1+n,k+m/n,1-n,k-m/n)
\]
を取り、第二表示ではこれを $(zb,a,zf,e)$ に移す。
二表示はそれぞれ有限平坦な四重被覆を与え、幾何学的一般モノドロミーは位数8の二面体群である。
二位置移動は各被覆で $n\mapsto-n$ である。
\begin{enumerate}
\item $zk^2=1$ では二表示が一致し、四点全てが前定理の正則な階数14窓である。
\item $\delta=(1-zk^2)^2-4zm^2=0$ では合計四つの相異なる階数15窓があり、各根の重複度は2である。
\item その他の通常倍率では合計八つの相異なる階数15窓がある。
\end{enumerate}
個数は固定した座標での幾何パラメータ数である。
\end{theorem}
\begin{proof}
$C=\lambda^2z+1$ と書く。前の交換行列による $(e,f)=u(N,D)$ において
$af+be=C u(a^2-zb^2)$、$ae+zbf=2\lambda z u(a^2-zb^2)$ である。
補完の二つの二次式の残差はともに
$G=\lambda(z+1)u(a^2-zb^2)-ab-ef$ の倍数で、分子の倍率はそれぞれ $D,N$、
共通分母は $\lambda(z+1)(a^2-zb^2)$ となる。
$i$ 番目の二次式の選択が対根に対応することは $zf_1 f+e_1 e=0$ から従う。
従って許容座標では補完と交換の二条件が次の二次式に一致する。
\[
 2\lambda z(af+be)=C(ae+zbf),\qquad
 2\lambda z(ab+ef)=\lambda(z+1)(ae+zbf).
\]
逆にも第一式から $[e:f]=[N:D]$ が従い、第二式で $G=0$ となる。$C=0$ の場合も
交換行列は可逆であり、この議論は $C$ を分母に使わない。

最初の階数低下条件 $a^2+e^2=2zbf$ では $2(ae+zbf)=(a+e)^2$ である。
$a+e=0$ なら二式の和差から $(a-e)(b-f)=0$ となり、一方の行列が零になるか
$a^2=zb^2$ となるので除かれる。$a+e=1$ と規格化すると $b+f=k$、
$(a-e)(b-f)=m$ であり、$n=a-e$ として残条件が $P(n)=0$ になる。
もう一つの階数低下条件は上記の座標交換で移る。この交換は二つの補完式も保つ。
得られた全根で $n\ne0$ であり、$a^2-zb^2,e^2-zf^2,a^2-zf^2,e^2-zb^2$ は非零である。
\path{computations/verify_rank_fifteen_covers.py} は分母を払ったこれらの式と $P$ の
終結式が全て通常倍率の単元であることを両表示で検算する。このため座標端でも許容根を失わない。
第一表示で $a+e=1$ とした他方の階数低下条件の左辺 $z(b^2+f^2)-2ae$ は $zk^2-1$ である。
根数は $n^2$ の二次式の判別式と $zm^2\ne0$ から従う。

最後に $L=\mathbb C(z,\lambda)$ 上、$zm^2$ は平方でない。
$\delta$ も $z\delta$ も平方でない。実際 $\lambda=2$ で $\delta$ の既約分数の
分子は正次数の平方自由式で $z=0$ に根を持たず、その分母は平方である。
$n^2$ の二根を $v_1,v_2$ とすると、その二次分解体 $L(\sqrt\delta)$ においても
$v_1v_2=zm^2$ は非平方である。共役な $v_1,v_2$ は従って両方非平方で、
平方類も独立である。分解体の次数は8であり、二組の符号と二根の交換から二面体群を得る。
\end{proof}

\begin{proposition}\label{prop:rank-fifteen-ramification-curve}
各表示の分岐上の被約な窓の曲線は幾何学的に既約で、その滑らかな完備化は
$y^2=-t^4+14t^2-1$、$j=35152/9$ の種数1曲線である。
両曲線上で一般窓は標準Koszul左右AS正則である。
\end{proposition}
\begin{proof}
分岐上で $1-zk^2=2rm$ により $r^2=z$ の符号を一意に取り、
$w=\lambda r$、$K=rk$ と置く。$rk,rm$ は
$r+1/r$ と $w+1/w$ の和・差の四分の一である。従って
\[
 r+1/r=2-(K-1)^2,\qquad
 w+1/w=(K+1)^2-2,\qquad n^2=(K^2-1)/2.
\]
$K=\pm1$ は $m=0$ を与えて除かれる。
$A=(2r-2+(K-1)^2)/(K-1)$、$B=(2w+2-(K+1)^2)/(K+1)$ と置くと
$A^2=(K-3)(K+1)$、$B^2=(K-1)(K+3)$ となる。
$U=A/(K+1),V=n/(K+1)$ は $4V^2-U^2=1$ を満たす。
$t=2V+U$ を用いて
\[
 U=(t-t^{-1})/2,\quad V=(t+t^{-1})/4,\quad
 K=4/(1-U^2)-1,\quad y=2tB/n
\]
とすれば所定の四次式が従う。これらの式は両方向の有理逆を与える。
四根は相異なり、二元四次式の不変量から記載した $j$ を得る。

$K=2$ の八つの $(r,w,n)$ は139で全て単純分解する。二表示を合わせた16窓について、
上記コードはテンソルの隣接階数6、反対階数15、経路次元35,56、中心の文字次元
$(3,2,0,0)$ を調べ、全四始点の四中心商で長さ13まで標準次元を得る。
正の固定族の一般点からの経路・中心次元の下界と命題\ref{prop:fixed-residual-regularity}
の持上げにより、この16幾何点の全次数の正則性が従う。
定階数の中心から四式を近傍に延ばせるため、各既約分岐曲線でも非空開集合が正則となる。
\end{proof}

% !TeX root = ../papers/nc-p3.tex
\section{階数15の分岐曲線全体での中心と長さ六の商}
\label{sec:rank-fifteen-branch-center}
\begin{theorem}\label{thm:rank-fifteen-branch-center}
標数零で、前節の二本の被約分岐曲線の全ての通常点で四次中心の文字次元は
$(3,2,0,0)$ である。文字0と1から二式ずつ選ぶと、四始点で同時に、
長さ0から6の中心商の次元を $1,4,10,20,31,40,44$ にできる。
以下の消去で得られる有限集合は第一表示で2720点、第二表示で2616点あり、
この全5336幾何点は標準Koszul左右AS正則である。
\end{theorem}
\begin{proof}
$F=-t^4+14t^2-1$、$y^2=F$ とする。前節の有理逆を四つの関係空間に代入する。
$\mathbb Q(t,y)$ で分子 $A+Cy$、分母 $q$ を整数係数多項式で持ち、
文字ごとに長さ5まで道を左に延ばす。各消去では元の行の順に独立な行を選び、
ピボットの積でその行とピボット列の小行列式を求める。
この小行列式のノルムの分子、分母、および最終射影の分母の既約因子を記録する。
これで、選んだ小行列式が単元の点では同じ語が基底になる。
全四始点の次数0から5の次元は $1,4,10,20,35,56$ である。
得られた射影で四位置の交換子を計算し、各文字の中心行列にも同じ小行列式の検査を行う。
核の次元は順に $3,2,0,0$ であり、中心の一般基底が特殊化できる範囲が決まる。

両文字から四切断を取り、その係数と元の関係の分母を全て加える。
$t=4,y=13$ の101での還元で短い行から独立なものを選び、
実際の小行列式は標数零の関数体で計算する。
中心商では右に一文字延ばし、二つの連続する始点で長さ6まで
上記の次元を得る。使った小行列式のノルムと分母を再び加える。
これらの行で選んだ語が商を張る。一般正則点での次元と一致するので、
この範囲で特殊化した次元にも等しい。
二位置移動は $(t,y)\mapsto(-1/t,y/t^2)$ である。
実際この変換は $U,r,w,K$ を保ち、$V,n$ の符号を変える。
従って因子集合にこの変換の像を加えれば四始点全てが扱える。
必要なら各始点で有効な中心四式の開集合を同じ選択空間で交叉させる。

通常条件を外れる因子は $t$、$t\pm1$、$t^2+1$、
$t^2\pm2t-1$、$t^2\pm4t+1$、$t^4+10t^2+1$、
$t^8+20t^6-26t^4+20t^2+1$ および無限遠である。
これは $z,\lambda$ と通常条件の各因子の極・零点をノルムで求めて得られる。
$y$ の符号で $w$ は逆数に移り、$\lambda$ は $1/(z\lambda)$ に移るので
通常条件は二枚で一致する。極では $r$ が零または極になり、
無限遠では $K=-1$ となる。
この境界を除いた小行列式の例外集合は次の通りである。
\[
\begin{array}{c|r|r|r}
 \text{表示}&\text{第一座標の既約因子数}&\text{次数の和}&\text{幾何点数}\\ \hline
 1&49&1360&2720\\
 2&47&1308&2616
\end{array}
\]
\path{computations/verify_rank_fifteen_branch_center.py} は因子を関数体から
再生成し、各因子の単純な素体根と、その上の $y$ の相異なる両根を調べる。
各還元で長さ4,5の道、四次中心を求め直し、中心四式の商を全四始点で長さ13まで
計算する。次元は命題\ref{prop:fixed-residual-regularity}の列となる。
各既約因子の単根から得た付値で両根を検査したので、共役および体拡大で全点に戻せる。
同命題の中心持上げと全次数の下界により、表の全点は所要の正則性を持つ。
\end{proof}

\begin{corollary}\label{cor:rank-fifteen-near-collision}
通常倍率平面に $zk^2=1$ の軌跡全体を含む開集合があり、
その上の二つの四重被覆の全窓は標準Koszul左右AS正則である。
\end{corollary}
\begin{proof}
この軌跡上の全窓では定理\ref{thm:rank-fourteen-center}により
道と四次中心が一般次元を持ち、前節までの階数14の結果により中心商に標準次元を
与える四式がある。定階数の近傍に中心四式を延ばせば、長さ13までの
商の次元上界が開近傍で保たれる。道の全次数の下界と合わせて正則性が従う。
各被覆上でこれらの近傍の和を取り、閉補集合の倍率平面への像を除く。
二被覆は有限であり、指定の軌跡の上の全点を近傍が含むので、所要の開集合を得る。
\end{proof}

% !TeX root = ../papers/nc-p3.tex
\section{階数15の分岐曲線の境界円板と有理第一座標}
\label{sec:rank-fifteen-boundary-disks}
\begin{theorem}\label{thm:rank-fifteen-boundary-disks}
二つの階数15分岐曲線上で、標数零の通常点を取る。101上のある付値で
$t$ の射影的還元が前節の非通常境界上にあるか、$\mathbb P^1(\mathbb F_{101})$
に属するとする。この点は標準Koszul左右AS正則で、各始点の道空間の直和は
四変数の四次中心多項式環上で階数256の自由加群である。
特に $t\in\mathbb Q$ の全通常点が含まれる。
\end{theorem}
\begin{proof}
まず分岐曲線上の四次中心の一般切断を使う。各二次関係では行ごとに、
中心元では四位置にわたり共通の分母を払う。通常点の長さ5の経路が定階数なので、
得られた切断は全通常点で中心である。無限遠では $x=1/t$、$y'=x^2y$ とし、
同じ単位で各行または一つの中心元の全四位置の極を払う。
この座標でも四次式は $F(x)=-x^4+14x^2-1$ である。

還元の境界を三通りに調べる。以下の根を含む二次体では共役な根を一度に扱う。
\begin{enumerate}
\item $t=0,1,-1,\infty$、$t^2+1=0$、$t^2\pm4t+1=0$ の各剰余円板で、
展開の中心を標数零で $\alpha$ とする。$F(\alpha)\ne0$ のとき
$u=t-\alpha$（無限遠では $u=x$）とし、$y$ の二根をHenselで持ち上げる。
線形因子では二根は共役である。$\alpha^2=-1$ では $y_0=\pm4\alpha$
を別々に取る。$F(\alpha)=0$ のときは $u=y$ として
$F(t)=u^2$ を $t=\alpha+\cdots$ について解く。全ての微分の分母は101で単元である。

定理\ref{thm:rank-fourteen-all-ordinary}の行消去をこの局所変数に適用する。
定数体は $\mathbb Q(\sqrt d)$、$d=-1,2,3$ のいずれかで、単元の検査には
二次ノルムを使う。単元ピボットがなければ残る全係数の定数項の両座標が
有理数として零であることを検査してから $u$ で割る。
通常な持上げでは $u\ne0$ である。精度は $t=\pm1$ で6、
$t=0$ で16、無限遠と $t^2+1=0$ で17、
二つの分岐二次因子で14および15から始める。
各段で四文字の最大の除算回数だけ精度を減らしても長さ13で正の精度が残る。

\item $t^2\pm2t-1=0$ の二因子では $\alpha=\pm1+\sqrt2$ と置き、
$u=t-\alpha$ とする。$\alpha=1+\sqrt2$ では $y_0=\pm2(\alpha+1)$、
$\alpha=-1+\sqrt2$ では $y_0=\pm2(1-\alpha)$ である。
$u$ の平方根を添加し、位置 $p$ の第 $j$ 矢を $u^{w_{pj}/2}$ 倍する。
各二次関係を、その語の係数の付値の最小値で規格化する。
例えば第一因子 $\alpha=1+\sqrt2$ の第一表示、$y_0=-2(\alpha+1)$ では
整数重みの四行を
\[
 (0,0,-1,-5),\quad(0,-4,-1,-1),\quad
 (0,0,-1,-5),\quad(0,4,3,3)
\]
と取れる。他の七組の重みも下記の検算に収めた。
$u^{12}$ を法とする二次体上の厳密な展開で各行の最低項を取り、
打切り後の項の重みが最小値より大きいことを検査する。
得られた最低項を $\mathbb F_{101^2}$ に還元すると、
各隣接関係の階数は6、経路次元は長さ4,5で35,56、四次中心は文字別に
$(3,2,0,0)$ である。この中心から取った四式で全四始点の長さ13まで
所定の商次元を得る。重み付き基底変換は付値体の有限拡大で可逆であり、
中心の持上げはこの新しい格子に適用できる。

\item 残る四次と八次の境界因子は101で六つの相異なる二次因子に分かれる。
それぞれの根と、その上の $y$ の両方に多項式切断を直接代入し、
所定の四式の商を長さ13まで検算する。Frobenius共役は同じ結果を与える。
\end{enumerate}
\path{computations/verify_rank_fifteen_boundary_disks.py} が以上を再現する。
第一と第三の場合では二つの連続する始点を調べ、
前節の二位置移動で残る二始点を得る。境界の和はこの移動で保たれる。
同一の中心四式は、各始点で有効な選択の非空開集合を交叉して取れる。

通常な有理剰余も全て調べられる。有限な非通常有理剰余は $0,1,-1,10,-10$ の五つである。
残り96個は $t\mapsto-1/t$ で48対となる。各対の代表において、
$F(t)$ が平方ならその両根、非平方なら二次体上の一根を試す。後者の他根はFrobenius共役である。
\path{computations/verify_rank_fifteen_rational_residues.py} は両表示・全四始点で
経路、中心および新しく選んだ四式の商を検算する。
いずれも長さ13で商が消える。一般点からのHilbert次元の下界と
命題\ref{prop:fixed-residual-regularity}の議論で全次数の結論を得る。
\end{proof}

% !TeX root = ../papers/nc-p3.tex
\section{階数15の分岐上の全点正則性と種数五の曲線族}
\label{sec:rank-fifteen-all-ramification}
\begin{theorem}\label{thm:rank-fifteen-all-ramification}
標数零で二つの階数15分岐曲線の全ての通常窓は標準Koszul左右AS正則である。
各始点で道空間の直和は四つの中心四次式の多項式環上で階数256の自由加群となる。
通常倍率の開集合を、$zk^2=1$ と $\delta=0$ の両軌跡を全て含み、
その上の両四重被覆の全窓が同じ正則性を持つように取れる。
\end{theorem}
\begin{proof}
前節により、101で通常な次数2以上の剰余における長さ13の商だけが残る。
$S=\mathbb F_{101}[t]$、$B=S[y]/(y^2-F)$ とする。中心は前々節の四切断を
四位置の共通分母で多項式化する。各表示と始点0,1について、
\path{computations/verify_rank_fifteen_all_ramification.py} は
定理\ref{thm:rank-fourteen-geometric-reduction}の双対経路の計算を行う。
$B$ を $S$ 上の階数2の加群として、各段で適用する関係の行列 $A$、
$\ker A$ の基底と、この基底での $y$ の作用を求める。
独立な像の列からなる行列の行加群を取り、その正方基底行列の行列式を使う。
この行列式が単元となる局所環では像が直和因子になり、核が剰余体への
基底変換と交換する。従って前段からの矢の写像も含めて同じ計算が特殊化する。

四始点のうち二つと二表示について、長さ2から13までの全48文字ブロックをそれぞれ調べた。
一般次元は標準中心商の各次数の次元の二倍である。$F(t)\ne0$ の剰余点では
二つの $y$ の固有部分を考える。一般には $y^2-F$ が既約で両部分の次元が等しく、
核が基底変換と交換する領域ではエタールな二次被覆上で各部分が同じ階数を保つ。
例外因子を次の表に示す。中央二列は各始点の行列式の相異なる既約因子数、
最終列は両始点の和集合から境界と一次因子を除いた後の個数である。
\[
\begin{array}{c|rr|r}
 \text{表示}&\text{始点0}&\text{始点1}&\text{残る因子}\\ \hline
 1&29&29&8\\
 2&26&26&7
\end{array}
\]
各残余因子の根で四次中心を取り直す。$y$ が剰余体で分解するときは両根、
非分解なら二次拡大上の一根で、全四始点の商次元が標準列に一致する。
後者は相方と共役である。二位置移動の像で例外になる点も、この検査結果を
移すことで扱える。従ってどの通常な還元でも上界が持ち上がる。
前節の境界円板と合わせ、代数的な $t$ は数体上の付値を取れば全て尽くされる。
超越的な $t$ は一般点の体拡大である。中心の持上げと自由性は前節と同じである。
最後に定理\ref{thm:rank-fifteen-branch-center}で中心次元が全点で分かっている。
系\ref{cor:rank-fifteen-near-collision}の有限射の議論を分岐の全窓にも適用する。
\end{proof}

\begin{proposition}\label{prop:rank-fifteen-pencil}
二つの階数低下被覆の各々において、$\xi=zm^2/n^4$ は通常域で可逆な関数である。
$\xi=-1$ が階数14、$\xi=1$ が被約分岐曲線である。
他の各 $\xi$ のファイバーは幾何学的に既約で、滑らかな完備化の種数は
$\xi^2+\xi+1=0$ で3、その他で5となる。
二位置移動は各完備曲線上で不動点を持たず、商の種数はそれぞれ2と3である。
二つの種数3の値では両表示の一般窓が標準Koszul左右AS正則である。
\end{proposition}
\begin{proof}
幾何学的な値 $\xi=c^2\ne-1$ を固定し、$r=cn^2/m$、
$w=\lambda r$、$K=rk$ と置く。すると $r^2=z$ で、
\[
 n^2=(K^2-1)/(1+c^2),\quad
 (r+r^{-1})/2=K+cn^2,\quad (w+w^{-1})/2=K-cn^2.
\]
逆にこれらの式から $z,\lambda,n$ を復元でき、元の双二次式を満たす。
$P(n)=0$ から $zk^2=1$ は $\xi=-1$ と同値で、判別式は
$n^4(1-\xi)^2$ である。$s=c+c^{-1}$ とすると $K$ の射影直線上の関数体は
次の三つの平方根で生成される。
\[
 \sqrt{K^2-1},\qquad \sqrt{(K+s)^2-1},\qquad \sqrt{(K-s)^2-1}.
\]
実際 $(K\pm cn^2)^2-1= n^2 c^2 ((K\pm s)^2-1)/(1+c^2)$ である。
$r,w$ の二つのトレースは $K$ の非定数式で、恒等的に等しくも符号違いにもならない。
このため通常条件は各曲線の有限集合だけを除く。

三つの二次式の根対は $\{-1,1\}$ とその $\pm s$ 平行移動である。
$s\ne0$ で互いに重なるのは $s=\pm1,\pm2$ のときに限る。
これらの場合も根対を頂点間の辺と見ると閉路がなく、どの非空な積にも奇数位数の零点が残る。
従って拡大の次数は8である。分岐は各幾何根で指数2、無限遠で不分岐となる。
相異なる根の数は一般に6、$s=\pm1$ で5、$s=\pm2$ で4であり、
Riemann--Hurwitzから種数は5,3,1である。$s=\pm1$ が
$\xi^2+\xi+1=0$、$s=\pm2$ が $\xi=1$ に当たる。
二位置移動は $n$ のみを反転し $r,w$ を保つので、三平方根を同時に反転する。
三つの根対に共通点がないため、この元はどの惰性群にも属さない。

$s=1$、$K=3$ で $c^2-c+1=0$ の両根を取る。
139において二根と各 $n,r,w$ の二択は全て単純に分解する。
\path{computations/verify_rank_fifteen_pencil.py} は両表示の合計32窓で
通常条件、反対階数15、中心の文字次元と全四始点の長さ13までの商を確認する。
単純根の持上げと中心の開条件から、各種数3曲線の一般点も所要の正則性を持つ。
\end{proof}

\begin{corollary}\label{cor:rank-fifteen-octic}
両被覆の二つの種数3の特殊値は、幾何学的にはすべて同じ超楕円曲線
\[
 Y^2=t^8+14t^4+1
\]
を持つ。二位置移動は $(t,Y)\mapsto(-1/t,-Y/t^4)$ で、その商の種数2曲線は
$Z^2=(X^2+4)(X^4+4X^2+16)$ である。
\end{corollary}
\begin{proof}
$c^2-c+1=0$、$j^2=c$ を取り $H=jn$ とする。
$H^2=K^2-1$、$A^2=K(K+2)$、$B^2=K(K-2)$ と置くと、
$r=K+H^2+HA$、$w=K-H^2+HB$ が前命題の逆表示になる。
$U=A/K,V=B/K$ は $U^2+V^2=2$ を満たす。$t=(V-1)/(U-1)$ により
\[
 U=\frac{t^2-2t-1}{t^2+1},\quad V=\frac{1-2t-t^2}{t^2+1},\quad
 D=2t(t^2-1),\quad K=-\frac{(t^2+1)^2}{D}.
\]
$Y=DH$ で $D^2(K^2-1)=t^8+14t^4+1$ となり、各式は関数体で逆に解ける。
二位置移動は $H,A,B$ の全符号を反転するので $t'=(V+1)/(U+1)=-1/t$ である。
$D(t')=D(t)/t^4$ から $Y'$ が従う。$X=t-1/t$ と $Z=(t+1/t)Y/t^2$
は不変であり、上記の式を満たす。$t$ を添加すれば $Y$ が戻るので、これらが
不変体を生成する。両モデルの右辺は平方自由である。
\end{proof}

\begin{corollary}\label{cor:rank-fifteen-pencil-boundary}
$c^2\ne-1$ を固定した完備曲線から通常窓を得るために取り除く幾何点の個数は、
$s^2=1$ で44、$s^2\in\{4,-4,9/4\}$ で48、それ以外で56である。
\end{corollary}
\begin{proof}
$R=K+(K^2-1)/s$、$W=K-(K^2-1)/s$ は $r,w$ の半トレースである。
$K$ が有限なら $r,w$ は単元である。$r^4=1$ は $R=0,\pm1$、
$w^2=1$ は $W=\pm1$ と同値となる。残る $w=\pm r,\pm r^{-1}$ は
$W=\pm R$ と同値で $K=0,\pm1$ を与える。
従って境界は、前命題の全分岐値、$K=\infty,0$、$K^2+sK-1=0$
の上の全点である。無限遠では半トレースに極がある。
最後の二次式は $s^2=-4$ で重根を持つ。$\pm1,-s\pm1$ では零にならず、
$s\pm1$ における二つの値の積は $s^2(4s^2-9)$ である。$K=0$ が既に分岐値に
含まれるのは $s^2=1$ の場合である。各分岐値で4点、他の各値で8点を数えればよい。
\end{proof}

% !TeX root = ../papers/nc-p3.tex
\section{階数15曲線族の無限遠八点の重み付き極限}
\label{sec:rank-fifteen-pencil-infinity}
\begin{theorem}\label{thm:rank-fifteen-pencil-infinity}
階数15の各表示で、一般の水準 $\xi=c^2$ の $K=\infty$ 上の八点は
全て、矢の基底を取り直すと標準Koszul左右AS正則な極限を持つ。
これは二つの種数3の特殊水準でも全八点で成り立つ。さらにこの二水準では、
109上のある付値で $K$ が極を持つ任意の通常な標数零窓が正則であり、
各始点で四次中心の多項式環上の階数256の自由加群となる。
\end{theorem}
\begin{proof}
$u=1/K$、$L^2=1+c^2$ とし、$un$ の定数項が $1/L$ となるように $L$ の符号を選ぶ。
$v=L+c$ により $L=(v+v^{-1})/2$、$c=(v-v^{-1})/2$ となる。
$S=2c/L^2$ と置く。前節の二次式から
$r,w$ の各根の先頭項は、それぞれ
\[
 r\sim S u^{-2}\text{ または }u^2/S,\qquad
 w\sim -S u^{-2}\text{ または }-u^2/S
\]
となる。これらは指数を除けば単元の級数で、逆数は小さい方の根を与える。
例えば $u^2r$ について大きな根の方を解くと定数項 $S$ は単根である。
$n$ の二根も $un$ について単根となる。$L$ の二符号と $r,w$ の二選択で八点を尽くす。

$r\sim R u^d$、$w\sim W u^e$ と書く。第一表示の $a,b,e',f$（$e'$ は窓の第三座標）の
先頭項は順に
\[
 L^{-1}u^{-1},\quad (1+c/L)R^{-1}u^{-1-d},\quad
 -L^{-1}u^{-1},\quad(1-c/L)R^{-1}u^{-1-d}.
\]
ここで $b=(K+cn)/r$、$f=(K-cn)/r$ を用いた。$z=r^2,\lambda=w/r$ と
第二表示の交換にも同じ先頭項の積を使う。定理\ref{thm:rank-fifteen-boundary-disks}で
与えた整数重みによって第 $p$ 位置の第 $j$ 矢を $u^{w_{pj}/2}$ 倍する。
重みの表では $r$ が大きい場合に \texttt{pole-}、小さい場合に \texttt{pole+} を使い、
$w$ が大きい場合に符号 $+1$、小さい場合に $-1$ を使う。
各関係を行ごとに最低次数で割って特殊化する。

\path{computations/verify_rank_fifteen_pencil_infinity.py} は関係の各係数を既約分数にし、
分子・分母それぞれの最低次数の全単項式を $\mathbb Q(v)$ で加える。
その和の非消滅に加え、特殊化に使う全ての先頭係数が剰余でも単元であることを確認する。
$101$ で $v=2$ を使うと二表示と $r,w$ の四選択すべてで、隣接階数6、経路長4,5の
次元35,56、中心文字次元 $(3,2,0,0)$ および全四始点の長さ13までの標準商を得る。
これらの条件は $v$ の非空開集合に延びる。$v\mapsto c^2$ の一般ファイバーの
全てを含むように、有限個の例外値の像を除くことができる。

種数3の二水準では $c=j^2$、$j^4-j^2+1=0$ と取る。
109でこの式は四つの単根に分かれ、$v=j^2+j$ が必要な $L=j$ の両符号を含む。
全四根、二表示、四つの $r,w$ の選択で同じ剰余検算を行った。
一般窓の道次元の下界は重み付き特殊化でも保たれる。道と中心が剰余で一般次元を持つので
中心を持ち上げ、四中心商から全次数の自由性と正則性を得る。
これは重み付き極限そのものにも、その近くで $u$ の付値が正の通常窓にも適用できる。
\end{proof}

% !TeX root = ../papers/nc-p3.tex
\section{種数3の全通常点の四次中心と有限剰余}
\label{sec:rank-fifteen-octic-center}
\begin{theorem}\label{thm:rank-fifteen-octic-center}
階数15の二表示で、$\xi^2+\xi+1=0$ の両水準の全通常点が文字別四次中心次元
$(3,2,0,0)$ を持つ。以下の関数体消去で残る7104幾何点は
標準Koszul左右AS正則である。
\end{theorem}
\begin{proof}
八次式を $F=t^8+14t^4+1$ とする。$i^2=-1$ として
$E=\mathbb Q(i)(t,Y)$、$Y^2=F$ を取り、さらに $d^2=3$ とする。
$j=(d+i)/2$ は $j^4-j^2+1=0$ を満たす。
$\mathbb Q(i,d)=\mathbb Q(j)$ なので $E(d)$ における一つの消去で
全ての定数共役を扱える。$F$ は平方自由で、定数体を拡げても非平方である。
\path{computations/rank_fifteen_octic_field.py} は各係数を $A+dB$、$A,B\in E$ として
持ち、$d$ についてのノルム $A^2-3B^2$ で逆を取る。八次モデルの $K,H,r,w$ を使うと
$(a,b,e,f)=(1+H/j,(K+jH)/r,1-H/j,(K-jH)/r)$ となる。

\path{computations/verify_rank_fifteen_octic_center.py} は長さ5までの道を左に延ばす。
109の $(t,Y,i,d)=(2,46,33,49)$ で独立な行・ピボット列を選び、同じ消去を
標数零の上記の体で行う。各段で残る語の数は既知の経路次元と一致するので、
選んだ行だけで必要な全行の空間を張る。最終射影の係数 $A,B$ の分母と関係式の
分母を記録する。定理\ref{thm:fixed-opposite-rank-drop}で長さ4,5の経路次元は
各通常点で一定である。このため射影が正則な点では、恒等的に関係を殺し
自由語をそのまま返す射影を特殊化して実際の経路基底を得る。

各文字の交換子行列の独立な行を同様に取り、得られる小行列式のノルムを計算する。
具体的にはピボットごとに $E(d)/E$ と $E/\mathbb Q(t)$ のノルムを順に取り、
分子・分母を交叉約分して掛ける。得られる階数での核次元は $3,2,0,0$ である。
一般点での同じ下界は、定次元の経路上で中心の核を特殊化すれば各点でも成り立つ。
射影の極と小行列式の例外因子を合わせ、前節までに求めた非通常因子を除くと、
各表示につき17個の相異なる第一座標の既約因子が残る。双方でその次数の多重集合は
\[
 4,4,8,8,16,16,16,16,48,64,80,80,80,80,96,136,136
\]
である。

各既約因子が単根を持ち、$j$ とその $t$ 上の $Y$ が全て単純分解する素数を取る。
各素体で全四根 $j$ と両根 $Y$ において道、中心および四式の中心商を取り直すと、
全四始点で長さ13まで標準列を得る。単根の持上げおよび共役から、対応する通常点の
全次数の正則性が従う。一つの $\xi$ に $j$ の選択を固定すれば $t,Y$ の表示は
通常域で逆を持つ。従って二表示、二水準、次数和888、二枚の $Y$ を数えて7104点となる。
\end{proof}

\begin{proposition}\label{prop:rank-fifteen-octic-residues}
この四曲線の通常な標数零窓で、109での射影的な $t$ の還元が
$\mathbb P^1(\mathbb F_{109})$ に属し、
\[
 \bar t\notin\{5,22,33,37,53,56,72,76,87,104\}
\]
となるものは標準Koszul左右AS正則である。また $K^2+K-1$ が109で零に
還元される全通常窓も、$t$ の剰余次数によらず同じ正則性を持つ。
\end{proposition}
\begin{proof}
八次モデルの境界の第一座標を109で分解する。$K=\infty$ に対応する
$\bar t=0,1,-1,\infty$ は定理\ref{thm:rank-fifteen-pencil-infinity}で認証済みである。
残る有限の非通常有理剰余は命題に記した十個で、通常な96個は
$t\mapsto-1/t$ で48対になる。\path{computations/verify_rank_fifteen_octic_residues.py}
は各対の一代表で両表示を検算する。$Y$ の二根が素体にあれば両方、
非分解ならFrobenius共役の一方を取る。$j$ は符号を除いて $8,41$ で十分である。
実際 $j\mapsto-j$ は $H,r,w$ を保ち $n$ を反転し、窓の二位置移動を与える。
各代表では全四始点の中心商を長さ13まで確認した。

最後の境界は $K=10,98$ に分かれる。このとき $r=\pm i$、$w=2K+v$、
$v^2=4K^2-1$、$H^2=K^2-1$ である。$z=-1$ だが、使う関係の表示は正則である。
実際その分母に現れる $z,\lambda,z-1,\lambda^2 z-1$ は全て単元であり、
剰余関係は付値環上の通常窓の二次関係を特殊化して得られる。
$H,v$ はいずれも109で非平方の元の平方根で、Frobeniusは両符号を同時に反転する。
$K,r,w$ を固定して $H$ だけを反転すれば $n$ が反転して二位置移動になる。
従って $r$ の二値、それぞれの $j$ の符号代表、両表示について一組の $H,v$ を
検査すれば全選択を尽くす。ここでも隣接階数6、経路次元35,56、中心次元および
長さ13までの商の標準列が成立する。中心が持ち上がるので、通常な各付値体の
窓では全次数の自由性と正則性が従う。
\end{proof}

% !TeX root = ../papers/nc-p3.tex
\section{八次モデルの有限境界における二次関係の極限}
\label{sec:rank-fifteen-octic-boundary}
\begin{theorem}\label{thm:rank-fifteen-octic-boundary}
$\xi^2+\xi+1=0$ の二水準と階数15の両表示について、標数零の通常窓で
$109$ における $K$ の剰余が $0,2,-2$ のいずれかとなるものは、
標準Koszul左右AS正則である。従って前節の有理 $t$ 剰余のうち、
未認証のものは $K=\pm1$ 上の
$5,22,37,53,56,72,87,104$ に限られる。
\end{theorem}
\begin{proof}
$F=t^8+14t^4+1$ とし、$Y^2=F$、$D=2t(t^2-1)$、$H=Y/D$ とおく。
$K=0$ では $t=\pm i$、$H=\pm i$ である。$K=k=\pm2$ では
$t=-k/2\pm\sqrt2$、$H=\pm\sqrt3$ である。これらの式は109での
全幾何剰余も尽くす。$Y$ は各点で単元なので、各標数零の境界点
$(\alpha,Y_0)$ の近傍では $u=t-\alpha$ が相対局所座標となる。
$\eta$ を原始24乗根として $i=\eta^6$、$j=\eta^2$、
$\sqrt2=\eta^3+\eta^{-3}$、$\sqrt3=\eta^2+\eta^{-2}$ を選ぶ。

\path{computations/verify_rank_fifteen_octic_boundary.py} は各位置・各文字の
関係行列から全最大小行列式を取り、一組のPl\"ucker座標全体の分母を一度に払う。
各成分は $\mathbb Z[i,\sqrt3][t,Y]$ の元となる。
$Y(0)=Y_0$ として $Y(u)^2=F(\alpha+u)$ から係数を順次求めると、
$2Y_0$ が単元であるため全係数が109の上で整である。各組において
全成分に共通な $u$ の最大冪を除く。実際の位数は各組で10未満であり、
これより低い係数の消滅は数体内で厳密に判定する。
最初のベクトルを109で還元すると、各組に非零な座標がある。
それをピボットに選んだGrassmann座標の行列は局所環上の直和因子を与え、
通常な持上げでは元の関係空間を与える。従ってこの行列の剰余は、
指定した円板のすべての持上げで共通である。

計算では $j$ の剰余 $8,41$ を取り、符号は二位置移動で尽くす。
$t\mapsto-1/t$、$Y\mapsto-Y/t^4$ の窓移動は $H$ を反転する。
$k=0$ では両 $t$ を取り、各 $H$ の一符号を調べる。
$k=\pm2$ では109上のFrobeniusが $\sqrt2$ を反転し $\sqrt3$ を保つので、
一つの $t,H$ からFrobeniusと窓移動で全四点を得る。
検算は各剰余の四位置で二次関係の階数6、経路次元35,56、
四次中心の文字別次元 $(3,2,0,0)$ を確認し、四つの中心元による商で
長さ13まで標準列を得る。通常持上げの長さ4,5の経路次元と
中心次元は定理\ref{thm:fixed-opposite-rank-drop}、\ref{thm:rank-fifteen-octic-center}
で同じである。従ってこれらの経路と中心の核は付値環上で基底変換と可換で、
中心切断を持ち上げられる。標準四切断判定を適用して結論を得る。
\end{proof}

\begin{proposition}
同じ通常表示の二次関係空間を基底を変えずに $K=\pm1$ へ延長すると、
その極限関係が生成する二次代数の長さ3,4,5の次元は、両水準の両表示と
全始点でそれぞれ $20,38,72$ となる。
\end{proposition}
\begin{proof}
各点は $Y=0$ にあり、$F'(\alpha)\ne0$ なので $u=Y$ を取る。
$U=A/K$、$V=B/K$ とすると
$\alpha=(V-1)/(U-1)$ で、$K=1$ では
$(U,V)=(\pm\sqrt3,\pm i)$、$K=-1$ では
$(U,V)=(\pm i,\pm\sqrt3)$ である。各Pl\"ucker座標を
$P(t)+YQ(t)$ と書く。$t-\alpha=u^2/F'(\alpha)+O(u^4)$ より、
$P,Q$ の最初の非零Taylor係数の次数をそれぞれ $a,b$ とすれば、
二項の位数は $2a,2b+1$ となる。零多項式は無限大の位数で扱う。
小さい方の係数を $F'(\alpha)^a$ または $F'(\alpha)^b$ で割ると
先頭係数を得る。\path{computations/verify_rank_fifteen_octic_path_jump.py}
は各組の最低位数の座標を取り、対応する関係行列で数体上の道を左に延ばす。
両符号の全組合せで長さ3,4,5の自由語を数えると主張の値になる。
他の $j$ は定数体の共役で従う。従って通常な一般ファイバーの
四次経路次元35を保つには、これら二つの境界で別途高次関係の飽和が要る。
\end{proof}

% !TeX root = ../papers/nc-p3.tex
\section{種数5の水準の平面四次商と有限非分岐境界}
\label{sec:rank-fifteen-pencil-unramified}
\begin{proposition}\label{prop:rank-fifteen-pencil-plane}
$s=c+c^{-1}$、$s^2\notin\{0,1,4\}$ とする。
命題\ref{prop:rank-fifteen-pencil}の種数5曲線は、$[K:T:H:A:B]$ を座標とする
$\mathbb P^4$ の三つの二次式
\[
 H^2=K^2-T^2,\qquad A^2=(K+sT)^2-T^2,\qquad B^2=(K-sT)^2-T^2
\]
の滑らかな完全交叉である。二位置移動による商は $[H:A:B]$ を座標とする
滑らかな平面四次曲線で、その式は $D=A^2+B^2-2H^2$ と置いて
\[
 s^2(A^2-B^2)^2=4D(D+2s^2H^2)
\]
となる。
\end{proposition}
\begin{proof}
$[K:T]$ の各点上で三つの根を取るので射は有限であり、
三二次式は余次元3の完全交叉を与える。分岐値が相異なるので、
$T\ne0$ で零になる平方根は高々一つで、その平方の $K/T$ における零点は単純である。
$T=0$ では三根が全て非零である。従ってJacobi行列の階数は常に3で、
命題\ref{prop:rank-fifteen-pencil}の既約な関数体を持つ滑らかな完備曲線を得る。
随伴公式から埋込みの超平面束は標準束である。

$H,A,B$ は同時に零にならず、射影は二位置移動による三根の同時反転で不変である。
$D=2s^2T^2$、$A^2-B^2=4sKT$ を代入すると四次式を得る。
一般の像点で $T^2,KT$ が決まるため、残る二点は同時反転で移り合う。
射影の超平面束の次数は8であり、像の次数は4となる。従って表示した四次式が像の方程式である。
正規化は不動点のない商なので種数3で、これは平面四次の算術種数に等しい。
よって像は滑らかである。
\end{proof}

\begin{theorem}\label{thm:rank-fifteen-pencil-unramified}
一般の種数5水準の両表示で、$K=0$ と $K^2+sK-1=0$ 上の有限非分岐境界24点には、
通常部分の全ての窓が標準Koszul左右AS正則となるZariski近傍がある。
具体的に $c$ が $3/4$ に還元される標数零の通常窓で、
101において $K=0$ に還元されるもの、および197において $K^2+sK-1=0$ に
還元されるものは、各始点で四次中心の多項式環上の階数256の自由加群である。
\end{theorem}
\begin{proof}
$T=1$ とし、$L^2=1+c^2$、$R=K+(K^2-1)/s$、$W=K-(K^2-1)/s$ を取る。
$n=H/L$、$r=R\pm\sqrt{R^2-1}$、$w=W\pm\sqrt{W^2-1}$ で窓を表す。
$z=r^2$、$\lambda=w/r$ である。$R=0$ が上記の二次条件に当たる。
一般の $c$ では指定した三つの $K$ が相異なり、各々で三つの平方根は全て非零である。
関係の表示に必要な $z,\lambda,z-1,w^2-1$ は単元となる。
従ってこのチャートは境界そのものへ延びる。

$c=3/4$、$L=5/4$ と置く。$s=25/12$ で、101上の $K=0$ と
197上の二根 $K=35,45$ はいずれも上記の平方根を全て単純に分解する。
\path{computations/verify_rank_fifteen_pencil_unramified.py} は各根上の八通りすべてを取る。
各表示で計24点について隣接階数6、経路長4,5の次元35,56、
中心の文字次元 $(3,2,0,0)$、全四始点の四中心商の長さ13までの標準列を確認する。
$L$ の符号変更は $H$ の二択に含まれる。
命題\ref{prop:fixed-residual-regularity}の一般窓の経路・中心次元の下界と
剰余の上界が一致するので、中心を持ち上げ、四中心商の上界から
通常な持上げの全次数の自由性と正則性が従う。

各境界切断は $c$ の非空開集合の上で有限エタールである。
単純分解する上記の全点を検査したため、標数零でどの成分にも成功点がある。
道と中心の階数が一定の近傍で、四式の各始点の商の上界を与える小行列式の非消滅は開条件である。
有限な境界被覆から失敗集合の像を除けば、一般水準の24点を同時に含む所要の開近傍を得る。
\end{proof}

% !TeX root = ../papers/nc-p3.tex
\section{種数5の四つの分岐値における関係格子}
\label{sec:rank-fifteen-pencil-ramified-ends}
\begin{theorem}\label{thm:rank-fifteen-pencil-ramified-ends}
一般の種数5水準の両表示で、$K=\pm s\pm1$ 上の16境界点のそれぞれに、
通常窓が標準Koszul左右AS正則になる開近傍がある。さらに $c$ が101で
$3/4$ に還元され、$K$ がこの四値のいずれかに還元される全ての通常な標数零窓は
同じ正則性を持ち、各始点で四次中心の多項式環上階数256の自由加群となる。
一般水準の56境界点のうち、これまでの近傍判定で残るのは $K=\pm1$ 上の八点である。
\end{theorem}
\begin{proof}
$\delta,\epsilon\in\{1,-1\}$ とし、
$K_0=-\delta s-\epsilon$ と置く。$x$ は $\delta=1$ で $r$、
$\delta=-1$ で $w$ とし、その半トレースを $E=K+\delta(K^2-1)/s$ とする。
$E(K_0)=\epsilon$、$x=\epsilon$ である。他方の根を $q$ とすると境界値では
\[
 q^2-2(2K_0-\epsilon)q+1=0,\qquad n_0=\pm(c+\delta\epsilon)/c .
\]
前式の二根は相異なり、$E'(K_0)\ne0$、$n_0\ne0$ である。
これらは指定する101上の点でも単元条件を満たす。
$u=x-\epsilon$ と取ると $(x+x^{-1})/2=\epsilon+O(u^2)$ であり、
$c$ を係数に保って $K=K_0+O(u^2)$、$n=n_0+O(u^2)$、
他方の傾きも $q+O(u^2)$ と展開する。従って一次項の計算では $x$ だけを動かせばよい。

各位置・文字の関係行列から全最大小行列式を取り、共通分母と共通因子を除いて
$z,\lambda,a,b,e,f$ の多項式の組を得る。ここに各表示の窓を代入する。
\path{computations/verify_rank_fifteen_pencil_ramified_ends.py} は $c=3/4$ を
101で特殊化し、各 $\delta,\epsilon$ と両表示について $n_0,q$ の全四択を取る。
各組の定数項に単元があればその組を使う。定数項が全て零に還元される場合は、
$\mathbb Q(c)[q]$ で上記のモニック二次式を法として各分子の剰余が
\emph{恒等的に零}になること、分母が同二次式と互いに素であることを確かめる。
そのため $u$ で一度割ることが可能である。用いた窓の式の分母は $c,r,w$ の積の約数で、
これらは剰余でも単元である。割った後の定数項は $x$ による一次微分であり、
実際に各組で少なくとも一座標が単元となる。これにより、$c$ も変動する近傍で
元の関係を表すGrassmann座標が得られる。

各組の単元座標から行基底を復元すると、各剰余で二次関係の階数6、
経路次元35,56、中心の文字次元 $(3,2,0,0)$、全四始点の長さ13までの
四中心商の標準列を得る。前定理と同じ持上げにより、通常窓の近傍正則性が従う。
境界の有限被覆上で $c=3/4$ の全選択を認証したので、一般水準の全16点に適用できる。
無限遠八点と有限非分岐24点を合わせて48点となる。
\end{proof}

\begin{proposition}
残る $K=k=\pm1$ の八点では、各分岐と表示を固定すると、二次関係の
Grassmann極限は $c$ に依存しない。その二次代数の経路次元は
全始点で長さ3,4,5の順に $20,38,72$ である。
\end{proposition}
\begin{proof}
局所座標に $n$ を取る。$K^2=1+(1+c^2)n^2$、$R=K+cn^2$、$W=K-cn^2$ より
\[
 K=k+O(n^2),\quad r=k+\sigma(c+k)n+O(n^2),\quad
 w=k+\tau(c-k)n+O(n^2),\qquad \sigma,\tau\in\{1,-1\}.
\]
実際 $(r+r^{-1})/2=R$、$(w+w^{-1})/2=W$ の二次項を比べればよい。
両表示で $(z,\lambda,a,b,e,f)$ の境界値は全て1である。
上の検証コードは同じ多項式座標をこの点で評価し、組の定数項が全て零ならば、
表示の一次係数に沿う方向微分を取る。各組は一次までに非零となる。
非零ピボットで割った全係数を $\mathbb Q(c)$ で約分すると、全16組で分子・
分母の $c$ 次数はともに零である。この有理数の関係行列から長さ5まで道を延ばすと
上記の次元を得る。通常窓の四次・五次経路を復元するには、ここでも高次関係の飽和が要る。
\end{proof}

% !TeX root = ../papers/nc-p3.tex
\section{階数15の種数3四曲線上の全通常点}
\label{sec:rank-fifteen-octic-all}
\begin{theorem}\label{thm:rank-fifteen-octic-all}
標数零とする。$\xi^2+\xi+1=0$ の両水準で、階数15の両表示の
\emph{全通常点}は標準Koszul左右AS正則である。各始点で道空間は
四つの四次中心元が生成する多項式環上の階数256の自由加群になる。
\end{theorem}
\begin{proof}
まず109で $t$ が八次曲線の通常域に還元される場合を扱う。
$k=\mathbb F_{109}$、$F=t^8+14t^4+1$ とし、$k(t)[Y]/(Y^2-F)$ 上で
$j=8,41$ と両表示を取る。$F$ は109でも平方自由である。
\path{computations/rank_fifteen_octic_modular.py} は分母をモニックにした
$(a+bY)/d$ で消去する。道を長さ5まで延ばし、各射影の係数の極と実際の
階数小行列式の例外を記録する。四次中心の交換子にも同じ判定を行う。
有理関数の非零判定を一変数で行うには $a^2-Fb^2$ と $d$ を記録すればよい。
得られる道次元は $1,4,10,20,35,56$、中心次元は $(3,2,0,0)$ である。

\path{computations/verify_rank_fifteen_octic_ordinary.py} は文字1の中心二本を固定し、
文字0の三本の核基底のうち、第一・第三の対、または係数 $(1,3,5),(2,7,11)$ の対を使う。
始点0,1のそれぞれで商の道に一文字を加え、末尾の二次関係と長さ4での
中心関係を順次課す。分解する有理剰余で独立な行を選び、行列式は上の関数体で取る。
各段の選択小行列式はピボットの積である。前段までの関係・中心式・射影が
正則でこの小行列式が単元なら、選択列での逆をCramerの公式で取れる。
従って各始点で商次元の標準上界
\[
 1,4,10,20,31,40,44,40,31,20,10,4,1,0
\]
が長さ13まで得られる。二つの選択の例外多項式は、分母とノルム分子の
各平方自由部分の最小公倍数でよく、両方の例外に入る因子のみを残す。

始点0,1の残余の和集合を取り、既知の非通常因子と一次因子を除く。
$(\text{表示},j)=(0,8),(0,41),(1,8),(1,41)$ の順に
それぞれ $23,17,28,21$ 個の相異なる既約因子を得る。各因子の根と $Y$ の両シートについて
剰余体上で関係、道、中心、四中心商を再計算する。非分解の場合は
Frobeniusで交換される二点を一組として扱える。全てで全四始点の上界を得る。
一次の通常剰余は命題\ref{prop:rank-fifteen-octic-residues}で認証済みである。
$t\mapsto-1/t$ の窓移動で残る二始点へ、$j\mapsto-j$ の窓移動で残る定数根へ移る。
始点ごとに四式の選択が異なっても、良い四式の条件は各始点で同じ係数空間の
非空開条件であり、代数閉体上で四始点に共通な選択がある。
定理\ref{thm:rank-fifteen-octic-center}の中心次元と道次元が剰余でも保たれるので、
中心が持ち上がり、通常な持上げには全次数の自由性と正則性が従う。

次に $K=\pm1$ に還元される場合を扱う。$D=\mathbb Z_{(109)}[i,\sqrt3]$、
$j=(\sqrt3+i)/2$ とし、局所変数に $u=Y$ を取る。境界点の
$\alpha=(V-1)/(U-1)$ と $U,V$ の全符号は節\ref{sec:rank-fifteen-octic-boundary}の通りである。
$F'(\alpha)$ は $D$ の単元であり、$F(t)=u^2$ は $t(0)=\alpha$ から
$D[[u]]$ に解ける。\path{computations/rank_fifteen_octic_integral.py} は
標数零の八次曲線の関数体で、長さ4,5の経路基底と四次中心の核を正確に求める。
文字0の第一・第三と文字1の二本を選び、各中心元の全四位置で
一つの共通分母を払う。交換子は通常域の定階数の経路空間で一般点において零なので、
これらは分母を払った後も通常域全体で中心となる。

\path{computations/verify_rank_fifteen_octic_trace_disks.py} は
$D[[u]]/(u^{24})$ で二次関係と四中心式を展開する。
$1,i,\sqrt3,i\sqrt3$ の四成分を有理数係数の打切り多項式で保持し、
単元判定は定数項を109上の全四埋込み $(i,\sqrt3)=(\pm33,\pm49)$ で行う。
商の各文字ブロックで単元ピボットを順に取り、必要な階数に達する前に
単元が尽きた場合、残る全係数の定数項が数体上で厳密に零であることを確かめて
$u$ で割る。通常な持上げでは $u\ne0$ なので得られる式は関係として使える。
各文字で割った回数の最大値だけ次段の精度を下げ、各段の自由語は商を張る。
二位置移動が $(U,V)$ の両符号を同時に反転するため、第一の符号を固定して
第二の両符号と全四始点を取ればよい。全32組で長さ13に達しても
精度が正に残り、上記の標準上界を得る。この計算は二次極限で増加した経路を
高次関係の飽和で取り除く。

以上で109における通常剰余と $K=\pm1$ の全円板を尽くした。残りの非通常剰余は
定理\ref{thm:rank-fifteen-pencil-infinity}、命題\ref{prop:rank-fifteen-octic-residues}、
定理\ref{thm:rank-fifteen-octic-boundary}で認証済みである。完備曲線の各代数点は
109上の付値で還元でき、超越点は幾何学的一般点からの体拡大である。
通常な窓の一般Hilbert次元の下界に四中心商による上界を合わせ、
定理\ref{thm:rank-fourteen-all-ordinary}の自由性と左右Koszul複体の議論を適用する。
\end{proof}

% !TeX root = ../papers/nc-p3.tex
\section{種数5の有理水準における全境界円板}
\label{sec:rank-fifteen-pencil-trace-disks}
\begin{theorem}\label{thm:rank-fifteen-pencil-trace-disks}
$c=3/4$、$s=25/12$、$\xi=9/16$ の種数5水準の両表示を取る。
101で曲線の56境界点のいずれかに還元される標数零の通常窓は、
標準Koszul左右AS正則で、各始点で四中心の多項式環上の階数256の自由加群となる。
\end{theorem}
\begin{proof}
$K=k=\pm1$ 上では $u=n$ を使い、$D=\mathbb Z_{(101)}$、
$L^2=1+c^2$ とする。$K=k\sqrt{1+L^2u^2}$ とし、
$S^2=L^2+2cK+c^2u^2$、$T^2=L^2-2cK+c^2u^2$ について
$S(0)=c+k$、$T(0)=c-k$ の級数を取る。これらの初期値は単元である。
前節の関係から $r=K+cu^2+\sigma uS$、$w=K-cu^2+\tau uT$ である。
$K_0=\pm1$、$(\sigma,\tau)\in\{1,-1\}^2$ で全八点を尽くし、
二位置移動は $u$ の符号と二つの分岐符号を反転する。

\path{computations/verify_rank_fifteen_pencil_trace_disks.py} は、
二次関係の各Pl\"ucker座標を $D[[u]]/(u^{32})$ で評価する。
定数項に単元がない組は数体上での零を確認して $u$ で一度割り、
単元座標で行基底を取る。次に\emph{中心切断の前}に長さ5までの道を左に延ばす。
各位置で長さ1,2,3,4,5の文字次元を順に
$(1,1,1,1),(2,2,3,3),(5,5,5,5),(9,10,8,8),(14,14,14,14)$ として
付値環上で掃き出す。必要な階数に達するまで、単元ピボットがない段でだけ、
残る全行の定数項が有理数上で零であることを確認して $u$ で割る。
得られた語は通常持上げで全経路を張り、その総数は既知の長さ4,5の次元に等しいので
各文字の実際の基底を与える。

この基底で文字0,1の交換子方程式を作る。同様の飽和によって自由変数が
それぞれ3,2個になるまで消去する。これらは正の固定族の一般中心からの下界でもあるので、
残る方程式は分数体上、選んだ行の線型結合となる。単元ピボットを使う
後退代入は局所環の級数で行え、各通常持上げで同じ階数と核の自由変数を与える。
文字0で $(1,3,5),(2,7,11)$、文字1で二本の基底を選ぶと、
この手順で実際に持ち上がる四中心式が得られる。四中心商を右に延ばし、
定理\ref{thm:rank-fifteen-octic-all}と同じ上界を長さ13まで検証する。
$u$ で割った回数の最大値を各段の精度から引く。全32組で最後にも
正の精度が残るため、どの零判定にも十分な桁数がある。

残る有限境界では $K=0$ と $K=\pm s\pm1$ はこの同じ101で既に認証されている。
$K^2+sK-1=0$ では $r=\pm i$、$w=2K\pm\sqrt{4K^2-1}$、
$H=\pm\sqrt{K^2-1}$ を使う。上記コードは $\mathbb F_{101^4}$ に両根 $K$ と
三つの平方根の全選択を実現し、両表示で全四始点の四中心商を確認する。
最後に無限遠の八点では $L=\pm5/4$ に対応する $c+L=2,-1/2$ の両方で
定理\ref{thm:rank-fifteen-pencil-infinity}の重み付き関係を101で検算する。
六つの分岐値と他の有限境界値はこの標数でも所定の相異なる値を持つ。
以上と全次数の下界・中心商の上界から主張が従う。
\end{proof}

% !TeX root = ../papers/nc-p3.tex
\section{非正則な水平曲線の種数3水準への特殊化}
\label{sec:rank-fifteen-horizontal-specialization}
\begin{theorem}\label{thm:rank-fifteen-horizontal-specialization}
各階数15表示について、命題\ref{prop:rank-fifteen-pencil-plane}の三二次式を
$c$ の相対射影曲線とし、窓の表示のため $L^2=1+c^2$ を選ぶ。
この曲面の既約曲線 $B$ が $c$ の直線を支配し、その幾何学的一般窓が通常かつ
標準Koszul左右AS正則でないと仮定する。$c_0^2-c_0+1=0$ の各水準で
$B$ の射影閉包が取り得る特殊化は、$K=\pm1$ 上の八つの滑らかな点に限られる。特に各 $c_0$ 上で閉包はこの八点のいずれかを通る。
\end{theorem}
\begin{proof}
$c_0$ で $s=1$、$1+c_0^2\ne0$ である。$B$ の閉包の正規化上で
$c_0$ の上に付値を取る。曲線族は $c_0$ の近傍で射影的であるから、
この特殊化は存在する。必要なら有限拡大して幾何点と $L$ の選択を取る。

特殊点で $K$ が通常なトレースなら、三二次式の曲線は滑らかで
八次モデルの通常窓を与える。定理\ref{thm:rank-fifteen-octic-center}により
経路長4,5と四次中心の次元は周囲の通常部分の一般次元に一致する。
定理\ref{thm:rank-fifteen-octic-all}の証明はこの特殊点に中心四式の有限商の
証明書を与える。道・中心の定階数性により、付値環の一般窓にこの証明書を持ち上げられる。
従って $B$ の仮定に反する。

通常でない場合も $K=\infty$ なら定理\ref{thm:rank-fifteen-pencil-infinity}の
重みを使う。先頭係数は $c,L$ の有理式で、同定理の検算により $c_0$ の上で単元である。
従って $c$ も変動する付値環で同じ重みと剰余関係が使える。
$K^2+K-1=0$ の各有限点では、命題\ref{prop:rank-fifteen-octic-residues}が
実際の関係を評価した $r,w,z-1,w^2-1$ が単元のチャートが使える。
そこでも剰余の道と中心の次元および四中心商の証明書が得られており、
これらの単元は曲面上の特殊化でも保たれる。

$K=\pm2$ の八点は定理\ref{thm:rank-fifteen-pencil-ramified-ends}の
$\delta=\epsilon$ の二分岐に当たる。$q$ の二根は単純で
$E'(K_0)=-3$、$n_0\ne0$ なので、同じ相対座標 $u=x-\epsilon$ が使える。
\path{computations/verify_rank_fifteen_horizontal_specialization.py} は109における
$c^2-c+1$ の二根、二つの分岐、両表示および各四シートで関係の各文字ブロックを調べる。
定数項が全て零なら、固定した水準での零に加えて
$\mathbb Q(c)[q]/(q^2-2(2K_0-\epsilon)q+1)$ における恒等的消滅を確認する。
従って曲面でも $u$ で割ることができる。一度以内の除算後に単元ピボットがあり、
得られた剰余上で道・中心が所要の次元を持ち、全四始点で商の長さ13までの上界を満たす。
このため $c_0$ 上のこれらの点を中心とする付値でも一般窓が正則になる。

さらに $K=0$ では $H=\pm i$、$A=B=0$ である。特殊ファイバーの
$A^2=K(K+2)$、$B^2=K(K-2)$ の接錐は相異なる二直線で、ここには二つの節点がある。
一方、\emph{曲面上}では $u=r+1,v=w-1$ が正則局所座標になる。実際
\[
 R=(r+r^{-1})/2,\quad W=(w+w^{-1})/2,\quad
 K=(R+W)/2,\quad s=2(K^2-1)/(R-W).
\]
中心で $R=-1,W=1$ なので分母は単元である。
$c+c^{-1}=s$ と $n^2=(K^2-1)/(1+c^2)$ はいずれも境界値において単根で、
$n_0=\pm c_0$ となる。$H=Ln$ も単元であり、$A=s(r-R)/H$、$B=s(w-W)/H$ で戻せる。
$R,W$ はともに一次微分が零なので、$K,c,n$ の二方向の一次微分も零である。

上記の検算は二つの節点と両表示について各Pl\"ucker座標の組を評価する。
単元がない組に対し、$u=0$ の分岐 $(\delta,\epsilon)=(1,-1)$、
または $v=0$ の分岐 $(-1,1)$ で全座標が恒等的に零となることを、先の二次式を法として
確かめる。この二つの因子は正則座標そのものなので、対応する $u$ または $v$ で
割れる。先頭値はそれぞれ $r$ または $w$ での偏微分である。両定数根 $c_0$ と両表示で、
各節点につき二組を $u$、別の二組を $v$ で割ると全組が単元座標を持つ。
得られる関係で道・中心・四中心商の証明書を再び確認した。
この関係格子は独立な二変数を持つので、節点を通る任意の弧の通常一般窓に使える。

最後に $K=\pm1$ の上には $H=0$ で他の二根が非零となる点が各四つあり、いずれも滑らかである。
以上で他の全ての特殊化を除いた。
\end{proof}

% !TeX root = ../papers/nc-p3.tex
\section{トレース端での相対四次中心}
\label{sec:rank-fifteen-relative-trace-centers}
\begin{proposition}\label{prop:rank-fifteen-relative-trace-centers}
$c_0^2-c_0+1=0$ とし、前定理で残った $K=\pm1$ 上の各点を取る。
いずれの表示でも、この点に特殊化する付値環の通常な一般窓には、
長さ4,5でそれぞれ35語と56語からなる経路基底が取れ、
四次中心の文字0,1の部分にそれぞれ3本と2本の局所基底が取れる。
この中心から四式を選ぶと、全四始点で商の長さ4,5の文字別次元の上界は
それぞれ $(7,8,8,8)$ と $(10,10,10,10)$ である。
同じ式と語は、水準 $c$ も動く弧に適用できる。
\end{proposition}
\begin{proof}
$D=\mathbb Q[c]_{(c^2-c+1)}$ とし、$u=n$ を用いる。
$c\pm1$ および $1+c^2$ は $D$ の単元である。従って
定理\ref{thm:rank-fifteen-pencil-trace-disks}の $K,S,T$ の三つの平方根を
所定の初期値から $D[[u]]$ で取れる。二次関係を同じ式で展開し、
長さ5までの経路の射影、次に文字0,1の中心交換子の核を計算する。
\path{computations/verify_rank_fifteen_relative_trace_centers.py} は
各係数を $c$ の整数係数多項式の商として約分し、
ピボットの定数項の分子と分母がともに $c^2-c+1$ と互いに素であることを判定する。
単元が見つからず $u$ で割るときは、各定数項の分子が
$\mathbb Q[c]$ で零多項式であることを確認する。
したがって二変数の局所環でも同じ除算が成り立つ。
有限個の除算と単元の逆を取る操作は局所環の完備化で行え、
各弧の完備化にも写る。通常な一般窓で $u\ne0$ のため、
除算した行を関係として使える。

検算は $u^{11}$ で始める。二次関係で一桁を使い、長さ4,5の経路でそれぞれ
最大二回ずつ除算する。残る一回の中心交換子の除算後にも五桁が残る。
経路次元の下界と一般中心次元の下界により、この基底で得た射影と核は実際の
通常一般窓で所要の次元を持つ。文字0の三つの自由座標で
$(1,3,5),(2,7,11)$、文字1で二つの単位ベクトルを用い、四式による商を
右に延ばす。長さ4,5でそれぞれ最大二回ずつ除算し、
$K$ の両符号、第二分岐の両符号、両表示、全四始点の32組について
正の精度を保って主張の上界を得る。第一分岐の負符号は二位置移動で尽くされる。
\end{proof}

% !TeX root = ../papers/nc-p3.tex
\section{階数15の一般水準の全通常点}
\label{sec:rank-fifteen-generic-levels}
\begin{theorem}\label{thm:rank-fifteen-generic-levels}
標数零で、階数15の両表示に共通な有限集合
$E\subset\overline{\mathbb Q}$ が存在し、$\xi\notin E$ の種数5水準では
\emph{全ての通常窓}が標準Koszul左右AS正則である。
各始点は四つの中心四次式の多項式環上で階数256の自由加群である。
特に、通常で非正則な一般窓を持ち、水準を支配する既約曲線はない。
\end{theorem}
\begin{proof}
まず、定理\ref{thm:rank-fifteen-horizontal-specialization}で残った
$c_0^2-c_0+1=0$ の上の $K=\pm1$ の八点を扱う。
$t=n/(c^2-1)$ とする。分母はこれらの点の単元なので、
曲線だけでなく相対曲面でも $n$ と $t$ は同じ主イデアルを生成する。
命題\ref{prop:rank-fifteen-relative-trace-centers}の関係・中心を
$t^{25}$ を法として再構成する。$K(0)=1$ とし、$c$ の係数環では
$c^2-c+1$ と $c^2+c+1$ の\emph{両方}において単元となる分母だけを許す。
係数関数の定数項は $\mathbb Q(c)$ で正確に扱う。

\path{computations/verify_rank_fifteen_generic_levels.py} は級数ごとに
$\mathbb Z[c]$ の一つの共通分母を取り、$t$ の打切り多項式として和と積を取る。
逆元を取るのは定数項の分子が上の二つの既約多項式のどちらでも割れないときである。
単元が尽きて掃き出し中の行を $t$ で割るときは、残る全ての定数項の分子が
$\mathbb Z[c]$ の零多項式であることを調べる。このため任意の通常な一般窓を持つ弧で
同じ行が使える。平方根の初期値、関係格子の構成、長さ5までの未切断の道と
可換条件の核を取る手順は前命題の通りである。
文字0の二式は自由座標 $(1,3,5),(2,7,11)$、文字1の二式は座標基底から取る。
各始点で商に一文字を加え、末尾の二次関係と長さ4での中心式を課す。
各長さで四つの文字ブロックの最大除算回数だけ精度を下げる。
第二分岐の両符号、両表示、全四始点で長さ13まで正の精度を保ち、次元の上界は
\[
 1,4,10,20,31,40,44,40,31,20,10,4,1,0
\]
となる。

$K(0)=-1$ については $(c,K,r,w)\mapsto(-c,-K,-r,-w)$ を使う。
$n$ は同じで、二つの平方根の初期値 $c+K(0),c-K(0)$ はともに反転する。
$z=r^2,\lambda=w/r,(K\pm cn)/r$ は変わらず、第一・第二表示の窓はそれぞれ
まったく同じである。$-c_0$ は $c^2+c+1$ の根なので同時単元判定が適用できる。
二位置移動は第一分岐のもう一方の符号を尽くす。
よって八点に特殊化するどの通常な一般窓にも四中心商の証明書が得られる。
各次数で一般窓から来るHilbert次元の下界にこの上界を合わせると
四中心上の自由性と階数256が従う。通常な切断の完全対合を用いた
定理\ref{thm:klein-rational-period-four}の左右Koszul複体の議論が適用できる。

最後に四中心商の証明書を持つ通常窓の集合は構成可能である。実際、長さ5までの
経路基底と交換子行列の階数で有限層化できる。四式の係数は有限次元の核に属し、
四始点・長さ13までの各上界は有限行列の階数条件である。
よって係数の選択を含む構成可能集合を射影すればよい。
水準を支配する曲線の一般点には、この証明の八点の場合と
定理\ref{thm:rank-fifteen-horizontal-specialization}の証明の他の全ての場合により、
実際にこの有限証明書がある。従って幾何学的一般ファイバーのどの閉点も
証明書を持つ。もしその補集合と一般ファイバーとの交わりが非空なら、
構成可能性によりそこに閉点があるので矛盾する。補集合の $c$ 直線への像も
構成可能で一般点を含まないため有限である。この議論を
$\overline{\mathbb Q}$ 上で行い、$L^2=1+c^2$ の選択を有限射で降ろし、
両表示の例外の $\xi=c^2$ における像の和集合を取ればよい。
基礎体が代数閉でないときも、中心の係数空間における良い四式の非空開集合は
無限な基礎体上に点を持ち、左右の完全性は体拡大から降下する。
\end{proof}

% !TeX root = ../papers/nc-p3.tex
\section{階数15の二曲面における例外の零次元性}
\label{sec:rank-fifteen-vertical-curves}
\begin{theorem}\label{thm:rank-fifteen-vertical-curves}
標数零で、定理\ref{thm:rank-fifteen-covers}の二表示の各々において、どの種数5水準でも
少なくとも一つの無限遠点で、通常な一般窓に四中心商の証明書を持ち上げられる。
両曲面の通常な階数15窓のうち、この証明書を持たない幾何点の集合は有限で、
全て $\overline{\mathbb Q}$ 上に定義される。従って標準Koszul左右AS正則性が
失敗し得るのもこの有限集合の点に限られる。
\end{theorem}
\begin{proof}
定理\ref{thm:rank-fifteen-pencil-infinity}の $v=L+c$ を用い、
$r,w$ がともに小さい根である無限遠点を取る。$y=v^2$ なら
\[
 \xi=\frac{(y-1)^2}{4y}.
\]
$y$ と $y^{-1}$ は同じ水準を与え、$v$ と $-1/v$ は同じ $c$ の上で $L$ の符号を交換する。
どちらか一つの端で証明書が得られれば、既約な水準曲線の一般窓に適用できる。

まず無限遠での関係を $y$ の単項式だけに簡約する。
\path{computations/verify_rank_fifteen_vertical_curves.py} は重み付き極限で必要な全単元の
因子を取り、$2,v,v-1,v+1,v^2+1$ 以外に零や極を与えるものがないことを確かめる。
これらは種数5の各水準で全て非零である。各関係行で最初の語 $(a,b)$ の係数で割る。
位置 $p$ の別の語 $(i,j)$ に対して、行列 $M$ の一行を
\[
 e_{p,a}+e_{p+1,b}-e_{p,i}-e_{p+1,j}
\]
とする。位置は4を法とし、$e_{p,i}$ は16本の矢に付けた座標ベクトルである。
係数の符号を外し、$v,v-1,v+1,v^2+1,2$ の指数を並べた行列を $Y$ とする。
$M$ の独立な行上で $MG=Y$ を解くと、全行の $Y-MG$ では第一列以外が零で、
第一列も偶数であることを両表示で検算できる。
従って $G$ の有理数指数で矢をスカラー倍し、各行を再規格化すれば、
全係数は $\pm y^j$ になる。必要な根は幾何点で選べる。この基底変更は
四位置で周期的なので四次中心と全四始点の商を保つ。

この表示に対し $\mathbb Q(y)$ で長さ5までの経路基底、文字0,1の四次中心、
前節の四中心商を長さ13まで計算する。係数は整数係数の一変数多項式の既約分数とし、
実際に逆にした全ての分子の既約因子と初期分母の因子 $y$ を記録する。
記録した因子が消えなければ、全ての除算と階数上界が同じ点で成り立つ。
両表示の因子数は順に994,923である。$f(0)\ne0$ なら相反多項式
$y^{\deg f}f(1/y)$ を原始かつ最高係数が正となるよう規格化する。
両方の因子が記録にあるものだけを残すと17,33因子になる。
$y=0,\pm1$、$\xi=1$、$\xi^2+\xi+1=0$ を除き、相反で組にすると11,19軌道が残る。
これらの全軌道の代表について、101,103,107のいずれかで単根を指定し、
その根で経路と四次中心の次元および四始点の長さ13までの商を再計算する。
既約多項式の一単根での成功は標数零に持ち上がり、共役な全根に従う。
相反側の水準も同じである。これで全種数5水準で少なくとも一つの端の成功を得る。

証明書を持たない通常窓の集合は前定理の証明で定めた構成可能集合である。
水準を支配する曲線はその中に稠密な部分を持たず、上の議論によりどの種数5水準の
一般点もこの集合に入らない。他の通常な階数15水準は
定理\ref{thm:rank-fifteen-octic-all}と節\ref{sec:rank-fifteen-all-ramification}で認証済みである。
曲面上のこの集合の閉包は従って零次元であり、有限型性により有限である。
全ての方程式と階数条件は $\overline{\mathbb Q}$ 上で定義されるので最後の主張も従う。
\end{proof}

% !TeX root = ../papers/nc-p3.tex
\section{階数15の全通常窓の二次双対}
\label{sec:rank-fifteen-quadratic-dual}
\begin{theorem}\label{thm:rank-fifteen-quadratic-dual}
定理\ref{thm:rank-fifteen-covers}の両表示の全通常な階数15窓で、各始点の二次双対は
次数次元 $(1,4,6,4,1)$ の完全対合を持ち、次数5以上は零である。
特に前節に残る孤立点でも有限Frobenius条件は成立し、標準Koszul左右AS正則性は
元の代数のKoszul性と同値である。
\end{theorem}
\begin{proof}
位置 $p$ の二次関係を $R_p$、長さ $d$ の全ての隣接位置に $R$ を置いた空間の
共通部分を $K_{p,d}$ とする。これが二次双対の長さ $d$ の双対である。
文字の和を $\oplus$ と書く。全文字 $h$ の三次テンソルを、最初と最後の文字を
行 $i$、列 $k$ とする四行四列の表に書く。中央の文字は $i\oplus h\oplus k$ で一意である。
$K_{p,3}$ では列 $k$ が $R_p[h\oplus k]$ に、行 $i$ が
$R_{p+1}[h\oplus i]$ に入る。行の座標順は中央の文字と最後の文字の変換で並べ替える。

$J=h\oplus\{2,3\}$ に属する二行と二列のそれぞれは、全成分が非零の既知ベクトルの
スカラー倍である。$J\times J$ での四つの一致式により、この四本の倍率は
一つの倍率から決まる。残る行の関係空間は二次元で、対応する文字 $g$ は0または1である。
$J$ は $\{0,1\}$ または $\{2,3\}$ であり、$g\oplus J=J$ である。
厚い二次関係を $(L\mid I_2)$ と書くと $L$ も可逆なので、
$J$ 内の二成分を指定すれば残る二成分まで決まる。従って三次交叉は文字ごとに
高々一次元である。\path{computations/verify_rank_fifteen_quadratic_dual.py} は
四位置の $\det L$ の分子・分母が通常な単元の積であり、細い関係の全成分も単元である
ことを斉次座標から検算する。四階テンソル $W_p$ の一文字縮約が
各文字に非零の元を与えるので、これらがちょうど $K_{p,3}$ を張る。

$K_{p,4}$ の元は両端からの三次縮約により
\[
 (A\otimes1\otimes1\otimes1)W_p
   =(1\otimes1\otimes1\otimes B)W_p
\]
と書ける。反対切断を第一・第三文字を行、第二・第四文字を列とする行列 $M$ にすれば、
$(A\otimes1)M=M(1\otimes B^{\mathsf T})$ となる。
階数15では $M$ の核が一本であり、その生成元は二文字間の可逆な四行四列行列である。
実際、定理\ref{thm:fixed-opposite-rank-drop}のブロック分解で核を持つ文字は0または1で、
対反対称部分は可逆である。単元で行または列を規格化した対称ブロックは、
対称座標 $(t,t,s,s)$ 上で対角外成分がともに2の二行二列行列になる。
従って核の非零ベクトルは $t,s$ ともに非零である。これを二文字の行列に戻すと
$j$ 行 $j\oplus g$ 列に一つずつ非零成分があり、確かに可逆である。
規格化したブロックは対称なので左核にも同じ結論が成り立ち、
位置を巡回してもこの性質は保たれる。

$1\otimes B^{\mathsf T}$ はこの核直線を保つ。可逆な二文字行列に片側から作用して
スカラー倍となるため $B$ 自身がスカラーである。よって $K_{p,4}=kW_p$ である。
五次の元は $W_p\otimes v=w\otimes W_{p+1}$ と書けるが、非零なら第一文字での
$W_p$ の階数が1になり、階数4に反する。これで五次以後も零となる。
$W_p$ の一文字切断と隣接二文字切断の階数4,6により、最高次数への積は次数
$(1,3),(2,2),(3,1)$ で完全対合を与える。Koszulならこの二次双対から左右の
長さ4の自己双対な解消が得られ、Hilbert級数は $(1-t)^{-4}$ となるので主張が従う。
\end{proof}

\begin{corollary}\label{cor:rank-fifteen-sixth-test}
これらの窓のHilbert次元が最初に標準値を超え得るのは長さ6である。
各始点で文字0から3の順に、末尾の二次関係の挿入行列の大きさは
$56\times54,56\times54,56\times51,56\times51$ である。
長さ6の次元が84を超える条件は、このいずれかがそれぞれの目標階数
$36,36,34,34$ を下回ることと同値である。
\end{corollary}
\begin{proof}
長さ4,5は定理\ref{thm:fixed-opposite-rank-drop}により全通常点で一定であり、
各文字では一般点の下界に等しく $(9,10,8,8),(14,14,14,14)$ である。
長さ3は上の交叉から各文字で $16-6-6+1=5$ となる。
長さ6の商は $A_4\otimes R\longrightarrow A_5\otimes V$ の余核である
（各因子は対応する位置で取る）。出力は文字ごとに56次元、入力は関係の文字次元
$(2,2,1,1)$ と $A_4$ の畳み込みで上記の大きさを持つ。
正則な一般窓では直前の二項の次元がそれぞれ20と
$\dim A_2[h]=(2,2,3,3)_h$ である。従って完全性から表示した四階数を得る。
一定階数の経路基底の上で行列の階数は特殊化により増えず、目標時の余核次元の合計は
$20+20+22+22=84$ である。上記コードは各位置でこの挿入行列を直接生成する。
$\xi=9/16$、101上の $K=2$ では $H^2=3$ が非平方である。
二次剰余体で全八シートと両表示の行列を実際に構成し、全ての目標階数を確認した。
\end{proof}

% !TeX root = ../papers/nc-p3.tex
\section{階数15の添字保持同型と二つの独立な不変量}
\label{sec:rank-fifteen-adjoint-invariants}
\begin{theorem}\label{thm:rank-fifteen-adjoint-invariants}
階数15の両表示の各々は、添字保持同型で割っても二次元の標準Koszul左右AS正則窓を与える。
具体的には通常窓から二つの有理不変量 $J_2,J_3$ が得られ、各表示で代数的に独立である。
この二値への写像は一般に有限対一である。
\end{theorem}

最初に同型の必要十分条件を記す。$M:V_1^*\otimes V_3^*\to V_0\otimes V_2$ を
反対切断とし、非零の右核と左核の生成元を $N,D$ とする。前定理によりいずれも
非退化な二文字形式である。$D$ で $V_2$ と $V_0^*$ を、$N$ で $V_3^*$ と $V_1$
を同一視すると、$M$ の核は恒等行列の直線、像はトレース零の行列空間となる。
従ってスカラーを除く線型同型
\[
 T:\mathfrak{sl}(V_1)\longrightarrow\mathfrak{sl}(V_0)
\]
を得る。第一位置を固定した二窓が添字保持同型である必要十分条件は、対応する二つの $T$ が
$T'=a\operatorname{Ad}(g_0)T\operatorname{Ad}(g_1)^{-1}$ を満たすこと
（$a\ne0$）である。実際、添字保持同型は最初の四文字空間に作用し、
前定理の $K_{0,4}=kW$ を保つ。核形式を正規化すると、二つの反対側の基底変更は
双対基底変更とスカラーに限られるのでこの式が必要となる。逆にこの式から $M$ の
スカラー付き基底同値を復元できる。対応する四階テンソルの巡回する隣接切断が
四つの関係空間を決めるため、同じ基底変更を周期的に繰り返せば添字保持同型を得る。

一つの四行四列基底変更だけで調べるには、$X,Y\in\mathfrak{sl}(V_1)$ に
\[
 b_T(X,Y)=\tfrac14\operatorname{tr}(TX\,TY),\qquad
 m_T(X,Y)=T^{-1}\bigl(TX\,TY-b_T(X,Y)I\bigr)
\]
を付す。$U=\operatorname{Ad}(g_1)$ と $a\ne0$ が存在して
\[
 b_{T'}(UX,UY)=a^2 b_T(X,Y),\qquad
 m_{T'}(UX,UY)=a\,U m_T(X,Y)
\]
が全基底対で成り立つことが同値条件である。必要性は代入で従う。逆に
$I\mapsto I,\ TX\mapsto T'(UX)/a$ は二式により積を保つ全行列代数の同型である。
行列代数の単位的自己同型は内部自己同型なので $g_0$ が得られる。
この条件は $\det(g_1)$ と $a$ を逆にした有限多項式系として検査できる。

\begin{proof}[定理の証明]
両トレース零空間の非退化形式 $\operatorname{tr}(XY)$ で随伴を取り、
$S=T^*T$ とする。基底同値で $S$ はスカラー倍と共役しか受けないので、
$\operatorname{tr}S\ne0$ では
\[
 J_m=\frac{\operatorname{tr}(S^m)}{(\operatorname{tr}S)^m},\qquad m=2,3
\]
が添字保持同型不変である。チャートで核直線の生成元は小行列式の比で取れるから、
これらは有理関数である。

$x=K+H$ と $c$ を相対曲面の座標として、$101$ で $x=4,c=3/4,L=5/4$、
$A=10,B=6$ のシートを取る。三平方根の初期値が非零であるため、
$x,c$ はこの点でエタール座標である。
\path{computations/verify_rank_fifteen_adjoint_invariants.py} は
$\mathbb Z/(101^2)$ 上で三根を一回のNewton操作で持ち上げる。
反対切断で単元の15次小行列式を選んで両核を計算し、元の行列との積が零であることを
確認する。基底には非対角行列単位と $E_{ii}-E_{44}$ の三本を取り、そのトレース形式
の行列を $G$ とすれば、$S$ の行列は $G^{-1}T^{\mathsf T}GT$ である。
始点での $(J_2,J_3)$ は両表示とも $(1418,4898)\bmod 101^2$ である。
$x$ のみ、または $c$ のみを101だけ動かし、同じシートで差を101で割ると、偏微分の二列は
$(22,88)^{\mathsf T},(23,66)^{\mathsf T}$ となる。分母は全て単元である。
行列式は $34\bmod101$ で、標数零でも二つの微分の外積は恒等的に零ではない。

従って両表示の各々で不変量写像の像の次元は2である。始域も既約な曲面なので、
一般ファイバーは零次元で、有限型性から有限である。
定理\ref{thm:rank-fifteen-vertical-curves}の有限集合を除いても像の次元は変わらない。
\end{proof}

% !TeX root = ../papers/nc-p3.tex
\section{階数15の二表示の反転と六次階数条件の縮約}
\label{sec:rank-fifteen-word-symmetries}
\begin{theorem}\label{thm:rank-fifteen-word-symmetries}
標数零の幾何点で同じ $(z,\lambda,n)$ の二表示を取る。
第二表示の道代数は第一表示の道代数の反対代数に添字 $p\mapsto-p$ で同型となる。
標準Koszul左右AS正則性の失敗候補はこの二表示で一致する。
長さ6の経路の階数低下を決めるには、第一表示の始点 $0,1$、文字 $0,2$ の
四つの行列だけを調べればよい。目標階数は文字0で36、文字2で34である。
\end{theorem}
\begin{proof}
$W$ を第一表示の四階テンソル、$W'$ を第二表示のものとする。$r^2=z,d^2=\lambda$
とし、四位置で成分順を $0,1,2,3$ とする。単元からなる次の対角係数を取る。
\[
 (d^{-1},(rd)^{-1},r^{-1},1),\quad
 (d^{-1},r/d,r^{-1},1),\quad
 (d,d/r,r^{-1},1),\quad(zd,rd,r,1).
\]
これらを $D_p$ と書くと
\[
 W'_{ijkt}=D_0(i)D_1(j)D_2(k)D_3(t)W_{t,k,j\oplus1,i\oplus1}.
\]
\path{computations/verify_rank_fifteen_word_symmetries.py} は
$(a,b,e,f)\mapsto(zb,a,zf,e)$ を元のGram表示に代入し、この恒等式を分母を払って検算する。
テンソルの各隣接切断の像が関係空間なので、一文字ごとに基底を変え語を反転すると、
第二表示の始点 $p$、長さ $m$ の商が第一表示の始点 $-p-m$ の商に移る。
文字は語の位置のうち $0,1\bmod4$ に属する個数の偶奇だけ $1$ と排他的論理和を取る。
逆写像も同じ構成で得られ、左と右の条件が交換される。

次に $\pi=(2\;3)$ を文字群の線型置換とする。同じコードは
$W_{ijkt}$ が $W_{\pi(k),\pi(t),\pi(i),\pi(j)}$ に位置ごとの対角係数を掛けたもの
であることを検算する。従って始点の二位置移動は同じパラメータでの同型で、
文字には $\pi$ が作用する。最後に正の固定窓で
$(e,f)^{\mathsf T}=u\,\mathrm{exchange}(a,b)^{\mathsf T}$ と書く
（定理\ref{thm:fixed-opposite-rank-drop}の交換行列）。
通常域ではこの行列が可逆で $u$ と $zb^2-a^2$ は単元である。
$q^2=u(zb^2-a^2)$ と $\sqrt{-1}$ を取ると、各位置で文字に順に $2,3,2,3$ を加える変換も
対角係数を掛けて $W$ を保つ。コードは二つの変換の全対角係数を与え、
全256成分について恒等式を確認する。積に現れる根の偶数乗は元の単元で評価している。
長さ6では $2,3$ が各三回現れ、全文字に $1$ が加わる。従って四始点を二つに、
四文字を二つに減らせる。経路商そのものが各局所基底で同型なので、
余核の次元を検出するFittingイデアルも対応する。単元の根を添加した後の等式は
忠実平坦な基底変換から降下する。
\end{proof}

\begin{proposition}\label{prop:rank-fifteen-oriented-cubic}
両表示の幾何学的一般窓の最小周期はちょうど2である。
同一パラメータの第一・第二表示は、一般には添字保持同型でない。
\end{proposition}
\begin{proof}
前節の $S=T^*T$ は入力側のトレース零行列空間に作用する。
$d(X,Y,Z)=\operatorname{tr}((XY+YX)Z)$ とし、トレース形式で双対な基底 $E_i,E^i$ を取る。
\[
 C=\frac{\sum_{i,j,k}d(E_i,E_j,E_k)\,d(SE^i,SE^j,SE^k)}
                {(\operatorname{tr}S)^3}.
\]
完全縮約なので基底によらず、入力の共役で $d$ は保たれる。
$T$ のスカラー倍による因子も相殺するので、これは始点を固定した有理不変量である。
前節の $x=4,c=3/4$ の101上の点で元の反対切断と一位置巡回した反対切断から
それぞれ $T$ を作り直す。上記コードによる $C$ の剰余は第一表示で $(10,44)$、
第二表示で $(44,10)$ となる。核の規格化と分母はいずれも単元で、この点は標数零に持ち上がる。
従って各表示で一位置移動との同型は一般点で不可能であり、両表示も同一パラメータでは
区別される。二位置移動との同型は前定理で構成した。
\end{proof}

% !TeX root = ../papers/nc-p3.tex
\section{階数15の全座標端での六次経路}
\label{sec:rank-fifteen-coordinate-ends}
\begin{theorem}\label{thm:rank-fifteen-coordinate-ends}
正の固定曲線上の通常な階数15窓で、$a,b,e,f$ の少なくとも一つが零ならば、
全始点の長さ6の経路次元は84で、四次中心の文字次元は $(3,2,0,0)$ である。
\end{theorem}
\begin{proof}
定理\ref{thm:rank-fifteen-word-symmetries}の反転は座標端の和集合を保つので
第一表示を調べる。定理\ref{thm:rank-fifteen-covers}の $P(n)$ と標準座標を用いる。
$a=0$ または $e=0$ は $n=\pm1$ と同値である。二位置移動で $n=1$ とすれば
$z(k^2-m^2)=2$、すなわち $(z+1)(\lambda^2z+1)=8\lambda z$ となる。
$t=z$ と置くと、窓を与える楕円曲線と座標は
\[
 Y^2=-t^3+14t^2-t,\qquad \lambda=\frac{4t+Y}{t(t+1)},\qquad
 (a,b,e,f)=\left(2,\frac{t+1}{2t},0,\frac4{t+1}\right).
\]
一方 $b=0$ または $f=0$ のとき $m\ne0$ なので $k\ne0$ で、
$n=\pm m/k$ を $P(n)=0$ に入れると $m^2=-k^2$ を得る。
逆も成立する。二位置移動と $i^2=-1$ の根の選択により $m=ik$, $n=i$ とする。
すると $f=0$、$\lambda^2z+1=-i\lambda(z+1)$ で、今度は $t=\lambda/i$ により
\[
 z=\frac{1-t}{t(t+1)},\qquad \lambda=it,\qquad
 (a,b,e,f)=(1+i,2k,1-i,0).
\]
どちらの表示にも通常点で分母の零点はない。第二の場合 $t=-1$ は倍率式に反する。

各曲線の関数体で経路基底を長さ5まで構成し、
前節の四挿入行列の階数と、四つの文字ごとに全位置で矢と可換になる四次の元の空間を調べる。\path{computations/verify_rank_fifteen_coordinate_ends.py}
は全除算で逆元を取った分子のノルムを $\mathbb Q[t]$ 上で因数分解する。
被約表示の分母の新しい零点は、初期座標の分母または既に逆にした元の零点からしか生じない。
同じ長さの疎な行を同じ順で選び、主元の列を最小と最大に変えた二回の計算で
各々挿入の階数 $36,34,36,34$ と、中心の四つの核次元 $3,2,0,0$ を得る。表の数は相異なる既約因子の個数である。
\[
\begin{array}{c|rr|r|r}
 &\text{最小列}&\text{最大列}&\text{共通}&\text{通常域の残余}\\\hline
 n=1 &56&50&17&11\\
 m=ik,\ n=i &74&56&13&7
\end{array}
\]
楕円表示では $t=0,\pm1$ と $t^2-14t+1,t^2-6t+1,t^2+10t+1$ の零点は
非通常または表示外である。二根の交換は $\lambda\mapsto1/(t\lambda)$ であり、
$\lambda^2-1$ と $(t\lambda)^2-1$ の零点を交換する。
有理表示では $t=0,\pm1$ と $t^2+1,t^2\pm2t-1$ の零点を同様に除く。
階数14と15が同じ $t$ の上に来る場合は両方を計算に含めた。

残る各既約因子に対しては101,103,109,157のいずれかの素体上の単根を一つ指定し、
$Y$ または $i$ の二つの相異なる値の双方で経路基底・四階数・中心の核次元を再計算する。
単根での良い剰余は標数零の根に持ち上がる。二つの値を調べたので、根の体への
二次拡大が分解する場合も全ての幾何点を覆う。他の根へは共役で移る。
以上で全通常域において目標階数が従う。系\ref{cor:rank-fifteen-sixth-test}と
反転・二位置移動から結論を得る。
\end{proof}
\begin{corollary}
表の最後の列の18因子上では、対象の全ての通常窓が標準Koszul左右AS正則である。
\end{corollary}
\begin{proof}
上記コードは同じ全剰余点において四次中心の文字次元 $(3,2,0,0)$ と
四中心商の全四始点の長さ13までの次元も検算する。経路・中心の次元の下界と
命題\ref{prop:fixed-residual-regularity}の中心の持上げにより、
これらの閉点では全次数の正則性が従う。表の通常域の残余以外に対する
本節の一様な結論は六次の経路と四次中心に関するものである。
\end{proof}

% !TeX root = ../papers/nc-p3.tex
\section{倍率反転と階数15の座標端の一般同型類}
\label{sec:rank-fifteen-endpoint-orbits}
\begin{theorem}\label{thm:rank-fifteen-endpoint-orbits}
標数零の通常な階数15窓を考える。第一表示を $(z,\lambda,n)$ で取ると、
第二表示の $(z,1/(z\lambda),n)$ と添字保持同型である。
両表示が倍率を動かして与える添字保持同型類の集合は一致する。
また第一表示内で $\sigma(z,\lambda,n)=(1/z,1/\lambda,n)$ は添字保持同型を与える。
\end{theorem}
\begin{proof}
$k,m$ は $\lambda\mapsto1/(z\lambda)$ で不変であり、$\sigma$ ではともに $z$ 倍となる。
定理\ref{thm:rank-fifteen-covers}の $P(n)$ は両変換で同一である。
前者では第一表示の枠は同じで、対象の第二表示の枠は $(zb,a,zf,e)$ となる。
後者では対象の第一表示の枠が $(a,zb,e,zf)$ となる。
$r^2=z,d^2=\lambda$ とする。前者の対象テンソルは、元のテンソルの
$(i,j,k\oplus1,t\oplus1)$ 成分に、各位置で次の対角係数を掛けて得られる。
\[
 (rd,rd,1,1),\quad(rd,rd,1,1),\quad
 ((rd)^{-1},(rd)^{-1},z^{-1},1),\quad(r/d,(rd)^{-1},1,1).
\]
後者では四文字に同じ置換 $(2\;3)$ を施し、位置ごとに
\[
 (d/r,d/r,z^{-1},z^{-1}),\quad(rd,d/r,1,1),\quad
 (r/d,r/d,1,1),\quad((rd)^{-1},r/d,1,1)
\]
を掛ける。\path{computations/verify_rank_fifteen_word_symmetries.py} は
枠の四成分を独立変数として両恒等式を検算する。四位置で繰り返す矢の基底変更が
すべての隣接関係を対応させる。いずれも位置の順序と始点を保つ。
\end{proof}

\begin{theorem}\label{thm:endpoint-generic-fibers}
前節の第一表示の楕円曲線 $n=1$ 内で、一般点と添字保持同型なパラメータは
元の点とその $\sigma$ 像のちょうど二点である。
有理曲線 $m=ik,n=i$ の各成分内でも同様で、その二点は $t,-1/t$ である。
楕円曲線と有理曲線の添字保持同型類の像は異なる曲線である。
\end{theorem}
\begin{proof}
$T$ を節\ref{sec:rank-fifteen-adjoint-invariants}の写像とし、スカラー行列を零に送って
全行列空間に拡張する。枠の規格化 $a+e=2$ において反対切断の左右の核行列は
それぞれ $\operatorname{diag}(1,1,-1,-1/z)$ と $\operatorname{diag}(1,1,-1,-1)$ であり、
切断に掛けて確認できる。
$T$ をトレース形式で $H\in\operatorname{End}(V_0)\otimes\operatorname{End}(V_1)$
と見る。$P=(4,0,3,1,2)$、$Q=(2,4,3,0,1)$ を $0,\ldots,4$ での像の列とし、
\[
 D=\sum_{a,b\in\{0,1,2,3\}^5}
       \prod_{j=0}^4 H_{a_j,a_{P(j)},b_j,b_{Q(j)}},\qquad
 L=\frac{D^2}{(\operatorname{tr} T^*T)^5}
\]
と置く。各行列因子の行・列は置換の閉路に沿って縮約されるので、二つの独立共役を保ち、
$T$ のスカラー倍も $L$ では相殺する。従って $L$ は添字保持不変量である。

\path{computations/verify_rank_fifteen_endpoint_orbits.py} は各曲線の関数体上で
$J_2,J_3$ を計算する。楕円曲線では二式とも $t$ だけの関数になる。$t=2$ の不変量と
一致するための二式の分子のモニック最大公約式は、楕円・有理の順に
\[
 (t-2)(t-1/2)(t^2-110t/9+1),\qquad
 (t-2)(t-3)(t+1/2)(t+1/3).
\]
無限遠は最高次係数から除かれる。楕円上の $t=2$ の値を有理曲線上で探すと
最大公約式は1で無限遠にも解がなく、両像の相違が従う。

楕円上の $t=2,Y^2=46$ で $Y\equiv47\pmod{103}$ の剰余を取る。
上の四根は103で $2,52,25,33$ となり、それぞれ二つの $Y$ も相異なる。
$L$ は初めの二根上で99、残りの二根上で71となる。
命題\ref{prop:rank-fifteen-oriented-cubic}の $C$ は初めの二根の各々で65と63になり、
各一方だけが指定点と一致する。有理曲線では $i\equiv10\pmod{101}$ とし、
$t=2,-1/2$ で $C=93$、$t=3,-1/3$ で $C=98$ を得る。共役な $i$ の成分にも
同じ議論を適用できる。全分母は選んだ剰余点で単元である。

両最大公約式は平方自由であり、楕円の二次被覆もこの各点で不分岐である。
滑らかな完備曲線から有理不変量の値の射影直線への写像を取れば、
このファイバーは既に $J_2,J_3$ の段階で被約である。
$C,L$ を加えて残る長さは特殊点で高々2であるから、有限写像の一般ファイバーの
長さも高々2である。前定理の $\sigma$ は一般点を動かし、実際に二点を同型にする。
この比較は各正規化成分内のもので、有理曲線の二つの符号成分間の比較は残る。
\end{proof}

% !TeX root = ../papers/nc-p3.tex
\section{楕円座標端の全通常点での正則性}
\label{sec:rank-fifteen-coordinate-regularity}
\begin{theorem}\label{thm:rank-fifteen-coordinate-regularity}
標数零で $n=\pm1$ の両表示の全通常な階数15窓は標準Koszul左右AS正則である。
各始点の道空間は四本の四次中心元の生成する多項式環上で階数256の自由加群になる。
\end{theorem}
\begin{proof}
二位置移動と倍率反転により第一表示の $n=1$ を取る。
定理\ref{thm:rank-fifteen-coordinate-ends}の楕円曲線上では四次までの経路、
長さ5の経路、四次中心は既に一定次元である。
その関数体 $\mathbb Q(t,Y)$ 上で、中心の文字0の三自由座標の第一と第三、
文字1の二座標の両方を選ぶ。これら四本で割った商について、定理
\ref{thm:rank-fourteen-all-ordinary}の右への矢の追加と関係挿入を行う。
四本は全位置に共通の中心切断である。

\path{computations/verify_rank_fifteen_coordinate_regularity.py} は逆元を取る順を
疎な行と次数の低い係数を優先するように選ぶ。各係数を $(A+BY)/D$ の形で既約化し、
逆元を取る全ての分子のノルム $A^2-(-t^3+14t^2-t)B^2$ の既約因子を記録する。
初期座標での除算も含む。始点0,1の各々で長さ0から13までの商次元の上界は
\[
 1,4,10,20,31,40,44,40,31,20,10,4,1,0
\]
となる。各段で消去に用いた行の関係式だけでも、残した語が商を張る。
二位置移動の同型で他の二始点にも良い四本が存在する。四つの始点で共通の四本は、
中心の同じ係数空間で四つの非空開集合の交叉を取って得られる。

全除算に現れる相異なる既約因子は464個である。前定理の同時反転は楕円上で
$(t,Y)\mapsto(1/t,-Y/t^2)$ となる。ノルムは両方の $Y$ を覆うので、
$t^{\deg g}g(1/t)$ も記録にある因子だけを残せばよい。
$t=0$ を除くこの共通部分は69因子で、非通常因子を除き相反な因子を組にすると55軌道になる。
コードは各代表の単根と $Y$ の二つの相異なる剰余値を指定し、その両方で
経路・中心の次元と全四始点の長さ13までの商を検証する。
前節の座標端の場合と同じく、単根での成功と共役により各軌道全体を覆う。
一般点の道次元の下界とここでの中心商の上界から自由加群の上下界が一致する。
定理\ref{thm:rank-fourteen-all-ordinary}の自由性と左右の解消の議論で結論が従う。
\end{proof}

\begin{proposition}\label{prop:endpoint-moduli-normalizations}
座標端から得られる一般添字保持同型類を与える三つの像曲線は相異なる。
その滑らかな射影正規化は $j=35152/9$ の楕円曲線と二本の射影直線である。
\end{proposition}
\begin{proof}
定理\ref{thm:endpoint-generic-fibers}を与える不変量写像の各成分上の一般次数は2であり、
その関数体は同時反転による不変部分体に一致する。楕円上では
$s=t+1/t$、$v=(Y/t)(t-1/t)$ が同時反転で不変である。
\[
 v^2=(14-s)(s^2-4).
\]
$t^2-st+1=0$ と $v$ から元の関数体を二次で復元できるので、これが求める商である。
三根 $-2,2,14$ のLegendreパラメータは4で、記載の $j$ を得る。
有理成分では $s=t-1/t$ が不変部分体を生成する。

楕円と有理の像の相違は前定理で得た。有理曲線の二つの符号成分の間では、
$i$ 成分の $t=2$ の $J_2,J_3$ を $-i$ 成分で取る二式の最大公約式は
$(t+2)(t+3)(t-1/2)(t-1/3)$ で、無限遠にも解はない。
\path{computations/verify_rank_fifteen_endpoint_orbits.py} は
$109$ 上で $i=33$ を選び、元の点で $L=46$、相手の四候補全てで $L=63$ を得る。
各点は通常な階数15で全分母は単元である。従って元の不変量の値は相手の完備曲線の
像に属さず、両像も異なる。
\end{proof}

% !TeX root = ../papers/nc-p3.tex
\section{階数15の座標端で同型の重複が残る有限水準}
\label{sec:rank-fifteen-coordinate-collisions}
\begin{theorem}\label{thm:rank-fifteen-coordinate-collisions}
前節の同時反転で割った三曲線を $E,R_+,R_-$ とし、各点の水準を $U=zk^2$ とする。
二つの相異なる通常な階数15の点が添字保持同型を与えるなら、
表の該当欄の多項式のいずれかが零になる。
二つの有理成分は互いに共役である。最後の行では添字で水準を区別する。
\[
\begin{array}{c|l}
\text{比較する曲線} & \text{水準の必要条件を与える因子} \\\hline
 E,E & U+3,\ U-3,\ U-9,\ U^2-74U+169\\
 R_+,R_+\ \text{又は}\ R_-,R_- & U+3,\ U+1,\ U^2-42U+9\\
 R_+,R_- & U+3,\ U-9\\
 E,R_\pm &
 U_E^3+13U_E^2+259U_E-529,\quad U_R^3-147U_R^2-61U_R-49
\end{array}
\]
最後の行では二つの三次式がともに消えなければならない。
従ってこの表の外では、三つの座標端を合わせても、前節の商パラメータが添字保持同型類を
一意に定める。表の零点における同型の存在はこの必要条件からは従わない。
\end{theorem}
\begin{proof}
$E$ の座標を $s=t+1/t$, $v^2=(14-s)(s^2-4)$ とし、有理曲線では $s=t-1/t$ を使う。
$R_+$ では $i^2=-1$ を一つ固定し、$R_-$ では $i$ を $-i$ に替える。水準は
\[
 U_E=\frac{(s+10)^2}{16(s+2)},\qquad
 U_{R_+}=\frac{i(s^2+4)}{8s}.
\]
通常域ではこれらの分母は零にならない。
節\ref{sec:rank-fifteen-endpoint-orbits}の四不変量 $J_2,J_3,C,L$ を各曲線の関数体で求める。
\path{computations/verify_rank_fifteen_coordinate_collisions.py} は五重縮約の
各文字の和が零になる条件を、置換 $Q$ の二閉路に沿って解く。
こうして全成分の積を整数多項式で足し、最後に共通分母を戻す。
結果を $t^2-st\pm1$ で簡約すると、$J_2,J_3,L$ は全ての曲線上で $s$ の有理関数であり、
$E$ 上の $C$ は $v$ の一次式になる。その $v$ の係数は
\[
 \frac{-48(s+2)(s^2-28s+4)}{(s^2+68s+196)^3}.
\]
従って通常域で $s$ が一致して $v$ が異なる二点は、トレースが非零なら $U_E=3$ でしか
同型になり得ない。

二点の $s$ 座標を $s,u$ とする。各不変量の両曲線上の表示を既約分数 $A/B,A'/B'$ とし、
等式の分子 $A(s)B'(u)-B(s)A'(u)$ を取る。
$J_2,J_3,L$ を常に使い、双方が有理曲線なら $C$ も使う。同一成分の比較では
各行を $s-u$ で割る。得た行の全ての対について $u$ を消去し、非零な終結式の
最大公約式の平方自由部分を取る。非通常点、階数14の $U=1$ と
$\operatorname{tr}S=0$ の因子を除いた後、水準式で $s$ を消去する。
コードは五方向の比較で、表の因子のうちこの消去から来るものを最高係数1で検算する。
共通零点は各非零終結式を消すので、この操作は必要条件を失わない。

残る $\operatorname{tr}S=0$ は三成分とも $U=-3$ に一致する。
ここで $\operatorname{tr}(S^2)$ は非零であり、
$\operatorname{tr}(S^3)^2/\operatorname{tr}(S^2)^3$ の値は $E$ で $64/375$、
両有理成分で $-64/81$ となる。これは $J_3^2/J_2^3$ を約分して $U=-3$ の方程式で
剰余を取ると得られる。従ってこの水準でも楕円と有理の間の同型は除かれる。
\end{proof}

% !TeX root = ../papers/nc-p3.tex
\section{階数15の全座標端での標準AS正則性}
\label{sec:rank-fifteen-all-coordinate-ends}
\begin{theorem}\label{thm:rank-fifteen-all-coordinate-ends}
標数零の通常な階数15窓で枠のいずれかの座標が零ならば、両表示とも
標準Koszul左右AS正則である。各始点の道空間は四本の四次中心元が生成する
多項式環上で階数256の自由加群になる。従って定理\ref{thm:rank-fifteen-vertical-curves}
の有限な正則性失敗候補は、$abef\ne0$ の内部に限られる。
\end{theorem}
\begin{proof}
定理\ref{thm:rank-fifteen-coordinate-ends}で全座標端を楕円成分と
有理成分に還元し、定理\ref{thm:rank-fifteen-coordinate-regularity}で楕円成分を認証した。
残る第一表示の $m=ik,n=i$ で $\mathbb Q(i)(t)$ を係数体とする。
四次中心と経路の一定次元は全通常点で既知である。楕円の場合と同じ自由座標の四中心切断で割り、
始点0,1から右へ矢を追加する。
\path{computations/verify_rank_fifteen_coordinate_regularity.py} は、
両始点の各文字で長さ13まで前定理と同じ商次元上界を与える主元を取る。
逆にした分子 $A+Bi$ の各ノルム $A^2+B^2$ の既約因子は合わせて593個である。

同時反転 $t\mapsto-1/t$ により、$g(0)\ne0$ で
$t^{\deg g}g(-1/t)$ の規格化も記録にある因子は26個となる。
非通常点・表示外を除き、負の相反変換で組にすると残りは12軌道で、
代表因子の次数は順に
\[
 1,2,2,2,4,4,4,4,4,6,8,12
\]
となる。各代表について101,109,173,193のいずれかの単根を取る。
両方の $i$ の剰余値で、経路、四次中心および四始点の四中心商を独立に検算する。
$t$ の根の体に既に $i$ が含まれる場合にもこの二検査が全点を覆う。
よって各通常窓に、四始点で同時に使える中心商の証明書を得る。
中心多項式環に関する自由性と左右のKoszul複体の完全性は楕円成分と同じ議論で従う。
二位置移動、倍率反転と有理成分の共役を適用して全座標端に戻す。
\end{proof}

% !TeX root = ../papers/nc-p3.tex
\section{階数15のPauli基底と一般窓の有限同型判定}
\label{sec:rank-fifteen-pauli-comparison}
\begin{theorem}\label{thm:rank-fifteen-pauli-comparison}
第一表示の通常な階数15窓で、節\ref{sec:rank-fifteen-endpoint-orbits}の
$H\in M_4\otimes M_4$ は、両行列空間のPauli基底を対応させる十五項の和になる。
自然な十六次元空間を $H$ と単位元から復元できる非空開集合上では、
二窓の添字保持同型は48個のラベル置換について積の係数を比較すれば決定できる。
一般窓の独立共役に関する射影安定化群は位数16である。
\end{theorem}
\begin{proof}
$r^2=z,d^2=\lambda$ とし、$X_j$ を四文字への $j$ の平行移動とする。
$D(v_0,v_1,v_2,v_3)$ は対角行列を表す。出力と入力の順に
\[
 \begin{array}{c|cc}
     & A & B\\\hline
 0&D(r,r,1,z)X_1&D(1,1/r,d,rd)X_2\\
 1&D(1/r,r,1,1)X_1&D(1,r,d,rd)X_3
 \end{array}
\]
を取る。全四位置の $X_1$ は
\path{computations/verify_rank_fifteen_word_symmetries.py} の対角係数を伴って元のテンソルを保つ。
$B$ は定理\ref{thm:rank-fifteen-word-symmetries}の交互平行移動である。
実際 $a+e=2$ と階数低下式から $ae+zbf=2$ であり、交換式から
$q^2=u(zb^2-a^2)=-1/(\lambda z)$ となる。$q=i/(rd)$ を取れば表の係数を得る。
二つの位置のそれぞれで $A,B$ は可換で各平方はスカラーである。
また文字保存から任意の対角符号
$Z_b(j)=(-1)^{b\cdot j}$ の両側同時作用も使える。ここで $a,b,j\in\mathbb F_2^2$、
整数表示では下位ビットを第一座標とする。

各行列空間で $A,B$ の積に $Z_b$ を掛けると、物理的な文字移動 $a$ と位相 $b$ で指定される
十六本の基底直線を得る。同じ $a$ では四つの文字指標が独立であり、異なる $a$ の
台は交わらない。$A,B,Z$ の共役作用はこの空間を相異なる十六指標に分ける。
それぞれの共役の組は $H$ をスカラー倍まで保つ。入力のスカラー行列が $T$ の全核で、
出力のトレース零空間が全像であるから、そのスカラー指標は自明である。
実際、指標が非自明なら出力側の自明指標に対応する非自明な入力直線も核に入る。

$\pi$ を $1\mapsto1,\ 2\mapsto3$ の線型置換とし、$\beta=\pi^{-\mathsf T}$ とする。
対角符号から両側の移動は同じ $a$ で、$A,B$ の指標から位相はそれぞれ $b,\beta b$ となる。
従って $H$ はこの順で対応するPauli直線のテンソル積に分解する。
ラベル $x=(a,b)\in\mathbb F_2^4$ の成分を $H_x$ とし、便宜上
$H_0=I\otimes I$ と置く。$H$ 自身は非零ラベルの和で、十五成分は階数15により全て非零である。

二つの積を通常の積と $(A\otimes B)\circ(C\otimes D)=AC\otimes DB$ とする。
十六本の空間はどちらでも閉じ、各積のラベルは和になる。
$I,H$ を含み両積で閉じる最小の線型空間が十六次元なら、それがこの空間を内在的に決定する。
その可分テンソルの直線は基底の十六本だけである。実際、両Pauli基底で係数行列は対角である。
従って添字保持同型はこれらを置換し、通常積の保存から置換 $F$ は線型である。
左右の行列因子での交換子の符号も保つ。ラベル $(a,b),(c,e)$ の符号指数は
それぞれ $a\cdot e+c\cdot b$ と $a\cdot\beta e+c\cdot\beta b$ であり、
両方を保つ $F$ は48個である。コードで四基底ベクトルの像を列挙する。

$H_xH_y=c_{xy}H_{x+y}$、$e_x=\boldsymbol1_{x\ne0}$ とする。
二窓の同型の必要十分条件は、この48置換のいずれかとスカラー $h\ne0$ に対して
\[
 c'_{Fx,Fy}=h^{e_x+e_y-e_{x+y}}c_{xy}
\]
が全ラベルで成立することである。相異なる非零の一組から $h$ は一意に決まる。
十分性のため、各行列因子で四つのPauli生成元の平方を1に規格化し、移行先の
四生成元にも同じ規格化をする。四平方と六つの交換子が同じであり、
順序付き単項式は双方で $M_4$ の基底になる。生成元の対応は行列代数の同型を与え、内部共役で実現する。
二つの共役を適用した後で残る倍率は、上式により全ラベル群の指標である。
第一因子の交換子形式は非退化なので、一つのPauli元の共役がこの指標を実現する。
これで $H,H'$ のスカラーと独立共役による対応が得られ、節\ref{sec:rank-fifteen-adjoint-invariants}
の同値条件が適用される。

\path{computations/verify_rank_fifteen_pauli_comparison.py} は $101$ 上の
$(z,\lambda,n)=(4,9,46)$ で両積を閉じると十六次元になり、自己比較に成功する $F$ は
恒等写像だけになることを確かめる。単根なのでこの非消滅検査は標数零に持ち上がり、
不可約曲面の非空開集合で成り立つ。ラベルを固定する安定化元は各因子で指標を与え、
十五項を共通スカラー倍するには二指標の積が自明でなければならない。
一方の指標で他方が決まり、非退化性とPauli基底の生成性からちょうど十六組である。
\end{proof}
\begin{corollary}\label{cor:rank-fifteen-interior-101}
剰余標数101の付値を持つ標数零の体で、倍率が $(4,9)$ に還元されるとする。
この倍率の通常な階数15窓は両表示とも全て標準Koszul左右AS正則である。
\end{corollary}
\begin{proof}
上記コードは $P(n)$ の剰余根が $46,47,54,55$ の四単根で、全枠座標が非零であること、
全ての根で四中心商が全四始点の長さ13まで所定の次元を持つことを確かめる。
モニックな $P$ の定数項は単元なので全ての根は整かつ単元であり、
剰余はこの四つに限られる。経路・中心の次元を同時に確認したので、
命題\ref{prop:fixed-residual-regularity}の中心の持上げから主張が従う。
\end{proof}

% !TeX root = ../papers/nc-p3.tex
\section{階数15の同一倍率上の添字保持同型類}
\label{sec:rank-fifteen-root-exchange}
\begin{theorem}\label{thm:rank-fifteen-root-exchange}
標数零の幾何点で、正の固定曲線の通常倍率 $(z,\lambda)$ を固定する。
$zk^2\ne1$ とする。分岐式 $\delta=(1-zk^2)^2-4zm^2$ が非零なら、
二表示の八窓は添字保持同型でちょうど二類に分かれ、各類の最小周期は2である。
$\delta=0$ なら二表示の四窓は一類になる。ここで数えるのは同じ倍率の上の窓である。
\end{theorem}
\begin{proof}
$r^2=z$, $s^2=r$, $d^2=\lambda$, $i^2=-1$ と取る。
$n$ を定理\ref{thm:rank-fifteen-covers}の四次式の一根とすると、
$N=rm/n$ により全根は $n,-n,N,-N$ である。
$\lambda'=1/(z\lambda)$ としても $k,m$ は変わらない。
$H,H'$ を第一表示の $(z,\lambda,n),(z,\lambda',N)$ の規格化テンソルとする。
前節の出力・入力の行列 $A_0,A_1$ を用いて
\[
 G=D(1,1,(ds)^{-1},(ds)^{-1}),\quad
 G_0=G(I-iA_0/r),\qquad G_1=G(I+iA_1)
\]
と置く。$A_0^2=zI$, $A_1^2=I$ なので両行列は可逆である。すると
\[
 H'(G_0\otimes G_1)=(G_0\otimes G_1)H .
\]
\path{computations/verify_rank_fifteen_root_exchange.py} は元のGramブロックで各成分を展開する。
四変数 $s,d,n,i$ の有理式で224成分の差が恒等的に零で、残る32成分の分子は
$i^2+1$ で割り切れる。分母には $i$ を使わず、二つのGram逆行列は
定理\ref{thm:rank-fifteen-covers}により両枠で通常点に極を持たない。
この共役は始点を保つ。定理\ref{thm:rank-fifteen-endpoint-orbits}と合成すると
同じ倍率で第一表示の $n$ と第二表示の $N$ が同型になる。
また $n\mapsto-n$ は二位置移動であり、定理\ref{thm:rank-fifteen-word-symmetries}で
同じ窓はこの移動と添字保持同型である。従って高々二類で、$N=\pm n$ となる分岐では一類である。

二類を区別するため、$B(S)$ を命題\ref{prop:rank-fifteen-oriented-cubic}の $C$ の
分子とし、$S_1=T^*T$, $S_0=TT^*$ と置く。階数15ではトレース零空間上の $T$ は可逆なので
\[
 J=\frac{(B(S_1)-B(S_0))^5}{\det(S_1|_{\mathfrak{sl}_4})}
\]
は全通常窓で定義される。$T$ のスカラー倍に関する分子・分母の次数はともに30で、
両方とも独立な基底共役で不変である。語の反転と一位置移動はいずれも二つの
行列空間を交換する。実際、各位置をその反対位置と核形式で対にし、
$H$ の行・列の組を $(0,2),(1,3)$ とすれば、反転後の組は $(3,1),(2,0)$、
一位置移動後は $(1,3),(2,0)$ となる。対の順序を逆にすると、
核形式での基底同一視を除いて行列を転置する。転置もトレース形式と
$d(X,Y,Z)=\operatorname{tr}((XY+YX)Z)$ を保つので、いずれの操作でも二つの $B$ は交換する。
$S_0$ と $S_1$ は共役なので分母は等しく、$J$ の符号は反転する。

最後に $B(S_1)-B(S_0)$ の零点を調べる。上記コードで $H$ の四文字別の $T$ を
共通分母倍して多項式行列にし、トレース随伴を取り、二つの三次縮約を引く。
$P(n)$ の分子で割る前に最高係数の単元の冪を掛ける。剰余は $n$ に関して二次であり、
$n,z,\lambda$ の共通単項式を除いて $P$ の分子との終結式を取ると、非零定数を除いて
\[
 \begin{aligned}
 &\operatorname{num}(\delta)^2 (z-1)^4 z^{28}
 (z\lambda^2-1)^{52}\\
 &\hspace{20mm}{}\cdot\lambda^{56}(\lambda-1)^{60}(z\lambda-1)^{60}
 \end{aligned}
\]
となる。分岐因子以外は通常倍率で単元であり、$n$ も単元である。
従って非分岐点では $J\ne0$ で、同じ $n$ の反対表示も一位置移動も元の窓と区別される。
二位置移動の同型と合わせて最小周期も従う。
\end{proof}

\begin{corollary}
定理\ref{thm:rank-fifteen-coordinate-collisions}の楕円成分の $U=3$ では、
同じ $s$ の二点が同型になり、二つの $s$ は相異なる類を与える。
各有理成分の $U=-1$ ではその成分の二点が一類になる。
\end{corollary}
\begin{proof}
楕円成分の $n=1$ では $P(1)=0$ から $N^2=U-2$、有理成分の $m=ik,n=i$ では
$N^2=U$ である。指定水準では $N=\pm n$ なので、第一表示内で
$\lambda\mapsto1/(z\lambda)$ が同型を与える。楕円上では $t=z$ を固定して
$Y$ の符号を交換するので $s=t+1/t$ を保ち $v$ の符号を交換する。
$s^2-28s+4=0$ の異なる二根の間の同型は前掲定理の $s\ne u$ の消去で除かれる。
有理成分では $t=\lambda/i$ が $(t+1)/(t-1)$ に移り、$s=t-1/t$ は $4/s$ に移る。
これは $U=-1$ の水準式の異なる二根を交換する。共役成分でも同じ計算が成り立つ。
\end{proof}

% !TeX root = ../papers/nc-p3.tex
\section{階数15の分岐曲線上の最小周期}
\label{sec:rank-fifteen-branch-period}
\begin{theorem}\label{thm:rank-fifteen-branch-period}
正の固定曲線の通常倍率上で、反対階数15の全ての窓の最小周期は2である。
特に分岐上では各倍率に一つずつ現れる標準Koszul左右AS正則な添字保持同型類が、
いずれも厳密に周期2を持つ。
\end{theorem}
\begin{proof}
非分岐の場合は定理\ref{thm:rank-fifteen-root-exchange}による。
分岐上では同じ倍率の窓は全て添字保持同型なので、第一表示の
$Y^2=-t^4+14t^2-1$ の上で示せばよい。$H$ を四文字対に作用する十六次の行列として
$H_{i,j,k,u}$ の順で成分を書く。一位置移動したテンソルは核形式による独立共役を除き
$H^\rho_{i,j,k,u}=H_{u,i,j,k}$ である。実際、移動後の出力は旧入力で、
その入力は旧出力の対の順序を逆にしたものである。
節\ref{sec:rank-fifteen-adjoint-invariants}の同型条件により、周期1なら両行列は独立共役とスカラー倍で移り合う。

$M_j=\operatorname{tr}(H^j)$ とし、節\ref{sec:rank-fifteen-endpoint-orbits}と
同じ五重縮約 $D$ を $H$ に、$E$ を $H^\rho$ に対して取る。
\path{computations/verify_rank_fifteen_branch_period.py} は楕円曲線の関数体で
$j=2,3$ の双方について $\operatorname{tr}((H^\rho)^j)=M_j$ を確認する。
具体的に、文字差 $h$ のブロックを $T_h[i,u]=H_{i,u\oplus h,i\oplus h,u}$ と置くと、
移動後の成分は $T_h[u,i\oplus h]$ である。行列 $H$ は二文字の和 $c$ を保ち、
$c$ での成分は $T_{x\oplus y}[x,y\oplus c]$ となる。よって四つの四次行列の
二乗・三乗のトレースの和で $M_2,M_3$ を計算できる。
周期1なら、あるスカラー $a\ne0$ について $E=a^5D$, $M_j=a^j M_j$ が同時に成立する。
従って次の二式がともに消える必要がある。
\[
 M_2(D^2-E^2),\qquad M_3(D^3-E^3).
\]
これらのノルムの分子を整数多項式として取り、最大公約式の相異なる既約因子を求めると
\[
 t,\quad t-1,\quad t+1,\quad t^2+1,\quad
 t^2-4t+1,\quad t^2+4t+1,\quad t^4+10t^2+1
\]
となる。どの通常点でも両方の共役点でGram分母は単元である。実際、$Y$ の符号変更は
$w\mapsto1/w$、すなわち $\lambda\mapsto1/(z\lambda)$ で通常条件を保つ。
従って上記のノルムの分子は元の二式の共通零点を落とさない。

各因子を定理\ref{thm:rank-fifteen-covers}に続く分岐の表示で調べる。
$U=(t-t^{-1})/2$, $V=(t+t^{-1})/4$ とすると $t=2V+U$ と $1/t=2V-U$ は
通常点で有限であり $t=0,\infty$ は除かれる。
$t=\pm1$ は $K=3$, $r=-1$, $z=1$ を与え、$t^2=-1$ なら $n=(K+1)V=0$ となる。
$t^2\pm4t+1=0$ では $U^2=3$ より $K=-3$, $w=1$ となり $z\lambda^2=1$ である。
最後の四次因子では $U^2=-3$, $K=0$ であり、$r+1/r=1$ と $w+1/w=-1$ から
$w=-r$ または $w=-1/r$ を得る。従って $\lambda=-1$ または $z\lambda=-1$ である。
全て通常域の外にあるので、一位置移動との同型は不可能である。
二位置移動の同型と分岐上の全点正則性は既に証明されている。
\end{proof}

% !TeX root = ../papers/nc-p3.tex
\section{小標数の全有理通常倍率における正則性の持上げ}
\label{sec:rank-fifteen-small-residues}
\begin{theorem}\label{thm:rank-fifteen-small-residues}
$p=7,11,13,17,19$ のいずれかとし、標数零の幾何的な倍率 $(z,\lambda)$ がある付値で
$\mathbb F_p$ の通常倍率に還元されるとする。正の固定曲線のこの倍率の上にある
反対階数15の全窓は標準Koszul左右AS正則である。
各始点の経路の直和は四つの四次中心元の多項式環上で階数256の自由加群になる。
\end{theorem}
\begin{proof}
定理\ref{thm:rank-fifteen-root-exchange}と反対代数の対応により、
一つの倍率のある一根・一表示で四中心商の証明書があれば、両表示の全根が同じ証明書を持つ。
添字保持同型は四期の中心と商を運び、語の反転も四始点を並べ替えて商の次元を保つ。
倍率反転による二つの変換も使うと、剰余倍率では
\[
 (z,\lambda),\quad(z,1/(z\lambda)),\quad(1/z,1/\lambda),\quad(1/z,z\lambda)
\]
の組から一代表を選べばよい。後の二点は $z^2\ne1$ により前の二点と異なり、
$z\lambda^2\ne1$ により各対の中でも異なる。

\path{computations/verify_rank_fifteen_small_residues.py} は
$z\lambda(z^2-1)(\lambda^2-1)(z\lambda^2-1)(z^2\lambda^2-1)\ne0$ の全剰余対を
列挙し、この四点軌道に分ける。代表ごとに $P(n)$ の最小次数のモニック既約因子
$g$ を選び、剰余体 $\mathbb F_p[n]/(g)$ の根で第一表示を作る。
表の後三列は $\deg g$ による代表数である。
\[
\begin{array}{c|r|rrr}
 p&\text{通常倍率数}&1&2&4\\\hline
 7&4&0&0&1\\
 11&40&3&4&3\\
 13&72&6&6&6\\
 17&156&14&15&10\\
 19&208&12&20&20
\end{array}
\]
全120代表で始点 $0,1$、文字 $0,2$ の六次挿入行列は順に階数 $36,34$ を持つ。
続いて長さ4,5の経路次元 $35,56$、文字別の四次中心次元 $(3,2,0,0)$ を確認し、
中心から四式を選ぶ。全四始点でその商の長さ0から13の次元は
$1,4,10,20,31,40,44,40,31,20,10,4,1,0$ となる。

任意に指定した標数零の持上げで、モニックな $P$ の全根は付値環の整閉包に属し、
定数項は単元である。代数閉体まで延長して一次因子の積を還元すると、
選んだ剰余根に落ちる根が少なくとも一つある。重根でもこの議論は使える。
通常域のGram分母は各検査枠で可逆で、経路と中心の次元は一般点からの下界に等しい。
命題\ref{prop:fixed-residual-regularity}の持上げの証明により、
その標数零の根に全次数の自由性と左右Koszul完全性が従う。
最後に標数零で根交換・反転を施せばよい。このため剰余の四次式の分解型や
剰余での階数の低下を制限する必要はない。
\end{proof}

% !TeX root = ../papers/nc-p3.tex
\section{階数15の倍率復元と一般添字保持同型類の次数}
\label{sec:rank-fifteen-parameter-orbits}
\begin{theorem}\label{thm:rank-fifteen-parameter-orbits}
正の固定曲線の第一表示の通常な階数15窓を取る。
節\ref{sec:rank-fifteen-pauli-comparison}の二つの積で生成する空間が十六次元なら、
倍率を任意に動かして同じ添字保持同型類を与える第一表示の全パラメータを、
以下の192個の射影線型変換と有理式で列挙できる。
幾何学的一般点では相異なる $(z,\lambda,n)$ はちょうど24個であり、
倍率 $(z,\lambda)$ はちょうど12組である。
第二表示を含めると同じ類を与える表示付きパラメータは48個になる。
\end{theorem}
\begin{proof}
$s^4=z$, $d^2=\lambda$, $t=sd$, $r=s^2$, $w=t^2$, $K=rk$, $N=rm/n$ とする。
二つの行列因子で、節\ref{sec:rank-fifteen-pauli-comparison}の平行移動生成元を
各平方が1になるよう規格化する。出力の $A,B$ を $r,\sqrt d$ で、入力のものを
$1,s\sqrt d$ でそれぞれ割る。各対は可換で正則な四文字平行移動として作用するから、
基底の一文字を選びその四つの像を取れば、対角基底変更で通常の平行移動になる。
符号行列はこの変更で変わらない。よってテンソルを
\[
 H=\sum_{x=(h,j)\ne0}\gamma_x P_x,\qquad
 P_x=X_h Z_j\otimes X_h Z_{\beta j}
\]
と書ける。ラベルの整数表示は $x=h+4j$、$\beta$ の像の列は $(0,3,2,1)$ である。
次の八量を定める。
\[
 \begin{gathered}
 a=1,\quad b=(K-1)/2,\quad L=n/2,\quad M=N/2,\\
 p=s+s^{-1},\quad q=s-s^{-1},\quad u=t+t^{-1},\quad v=t-t^{-1},\\
 A=(u-p)/4,\quad B=(u+p)/4,\quad C=(v-q)/4,\quad D=(v+q)/4 .
 \end{gathered}
\]
$\gamma_0=0$ を補うと全係数は次の順である。
\[
 (0,b,A,A,L,M,-C,D,\ a,a+b,B,B,L,M,-D,C).
\]
\path{computations/verify_rank_fifteen_parameter_orbits.py} は元のGram表示から
各単項行列へのトレース射影を求める。規格化前の行列 $P$ なら射影係数は
$\operatorname{tr}(HP)/(16P^2)$ である（$P^2$ はスカラー）。
$B$ を含む対ではさらに $sd$ を掛けて上記の規格化に戻す。全成分で差の分子が
四次式の分子で割り切れ、Gram逆行列は通常域で定義される。

共通スカラーも許したこれらの係数は
\[
 \begin{gathered}
 L^2+M^2=b(a+b),\qquad A^2+B^2=a(a+b),\qquad C^2+D^2=ab,\\
 LM=-AB=-CD
 \end{gathered}
\]
を満たす。逆に $a\ne0$ として $a=1$ に戻せば
\[
 s=B-A+D-C,\qquad t=B+A+D+C,\qquad n=2L
\]
と復元できる。実際 $s^{-1}=B-A-D+C$, $t^{-1}=B+A-D-C$ は上の二次式から従う。
$r=s^2,w=t^2$ とすれば $(r+1/r+w+1/w)/4=1+2b$,
$(r+1/r-w-1/w)/4=4LM$ である。従って $z=s^4$, $\lambda=(t/s)^2$ において
$K=1+2b$, $N=2M$ であり、第一式から $P(n)=0$ となる。
通常な階数15の候補だけを残せば、復元の両方向が成り立つ。

二つの局所交換子形式を保つ48個のラベル置換を $F$ とする。
$P_xP_y=(-1)^{e(x,y)}P_{x+y}$ の符号指数は $x=(h,j),y=(h',j')$ に対し
$e(x,y)=(j+\beta j)\cdot h'$ である。各 $F$ に対し、係数配列の変換は
$\gamma'_x=i^{v(x)}\gamma_{Fx}$（$i^2=-1$）と共通スカラーで記述できる。
指数は $P$ の平方と積で決まり、四基底ラベルで各二択が残る。二解の比はラベル群の
符号指標である。逆に各局所形式を保つ $F$ は規格化した四生成元の対応により
各行列因子の内部共役に持ち上がる。積の符号条件を満たす別の解は一方の因子の
Pauli行列の共役で得られる。従ってこれで全 $48\cdot16$ 変換を得る。
二積の生成次元が16なら、同型は全てこの形になる。

上記コードは八量を独立変数として、係数配列の線型関係を保つ変換を列挙する。
各 $F$ につき四つ、合計192個が残り、全てで五つの二次式がそれらの生成イデアルに入る。
これらは射影係数の表示を保つ群をなす。各 $F$ で他の変換も特殊点で有効になり得るかを確認する。
一つの保存変換を施した後の違いは符号指標だけである。
係数 $\gamma_2=\gamma_3$ と $\gamma_4=\gamma_{12}$ はともに非零なので、
指標は $1,8$ の両ラベルで値 $+1$ を取る必要がある。この条件を満たす指標は四つだけで、
既に列挙したものに尽きる。これで任意の生成次元16の点に対し、
192像から復元した通常点が必要十分なリストになる。生成次元を仮定しなくても、
このリストに入る点は添字保持同型である。

最後に $(z,\lambda,n)=(4,9,46)\bmod101$ で生成次元16は既知である。
$r=2,d=3$ とし、剰余二次拡大で $s^2=2$ を解く。コードで192変換を施すと、
全てが通常かつ非分岐で、相異なる $(s,t,n)$ が192個、$(z,\lambda,n)$ が24個、
倍率が12組となる。元の $n$ は単根である。等式と分母の非消滅検査により、
少なくとも24個の異なるパラメータが得られる条件は標数零の非空開集合で成り立つ。
逆に同型な各 $(z,\lambda,n)$ には $(s,d)$ の八通りの選択があり、
いずれも上の必要条件から192像に入る。従って上限は24である。
各倍率の第一表示で一般の一類を与える根は正負の一対であるため倍率数は12となる。
倍率反転が第一・第二表示を全単射で対応させるので、最後の個数も従う。
\end{proof}

% !TeX root = ../papers/nc-p3.tex
\section{同型変換で得られる内部の三曲線の全点正則性}
\label{sec:rank-fifteen-regular-orbits}
\begin{theorem}\label{thm:rank-fifteen-regular-orbits}
$U=zk^2$, $V=zm^2$ とする。正の固定曲線の通常な階数15窓で
次の三因子のいずれかが零なら、両表示の全根は標準Koszul左右AS正則である。
各始点の経路の直和は四つの四次中心元の多項式環上で階数256の自由加群になる。
\[
\begin{aligned}
 F_1&=(U^2+14U+16V-7)^2-16U(3U-1)^2,\\
 F_2&=(U^2-2U+16V+9)^2-16U(U-3)^2,\\
 F_3&=(V-4(U+1))^2-64U .
\end{aligned}
\]
三式は $(U,V)$ 平面上で相異なる既約曲線を定める。
従来の座標端を含む倍率 $U-V=2$, $U+V=0$ および八次曲線の倍率 $V=(U-1)^2$ と合わせ、
前節の群の作用でこれらの条件に移る全ての通常窓が認証される。
\end{theorem}
\begin{proof}
前節の $a=1$, $b=(K-1)/2$, $c=a+b$ と五つの二次式を使う。
第一表示で $n^2=1$ は $a,e$ の一方が零になる元の枠であり、
$n^2=-1$ は $P(n)=0$ より $m^2=-k^2$ となり $b,f$ の一方が零になる元の枠である。
ここでいう枠の四文字は前節の係数配列の八量と区別する。
従って座標端の条件は係数配列で $16\gamma_4^4=\gamma_8^4$ と書ける。
192変換における $(F8,F4)$ は、第一ラベルが $8,1,9$ の順に、第二ラベルが
それぞれ $(4,5,12,13),(2,3,10,11),(6,7,14,15)$ の十二対に尽きる。
符号に用いる四乗根は四乗で消える。初めの四対では $n^2,N^2$ のいずれかが $\pm1$ であり、
和 $U-1$、積 $V$ を用いると元の二つの座標端水準になる。

次の四対では $4A^2,4B^2$ のいずれかが $\pm b^2$ である。この二根の和は $4c$、
積は $V$ なので条件は $b^4\pm4cb^2+V=0$ となる。
最後の四対では $b,c$ を交換した $c^4\pm4bc^2+V=0$ が得られる。
$K\mapsto-K$ の共役な二式ずつを掛けると $F_1,F_2$ になる。
八次曲線は $\xi=N^2/n^2$, $\xi^2+\xi+1=0$、すなわち
$(L^2+M^2)^2=L^2M^2$ である。各変換の $4,5$ の符号指数は同じ偶奇なので、
他の二対の条件も $(A^2+B^2)^2=A^2B^2$ と $(C^2+D^2)^2=C^2D^2$ になる。
従って $V=16c^2$ または $16b^2$ であり、その積の式が $F_3$ である。
各 $F_j$ は $V$ に関する二次式で、その判別式は非零平方の $U$ 倍であるから既約である。
例えば $U=4$ では順に $V=23/16,-25/16,36$ が各曲線上にあり、
これらでは両四重被覆の全根が座標端の外にある。

この条件で選んだ変換先が通常なら、座標端の全点正則性または
定理\ref{thm:rank-fifteen-octic-all}を適用できる。変換の十分性には二積の生成次元を仮定しない。
変換先が非通常のときも、前節の係数配列は定義され、元のテンソルに独立共役とスカラー倍で対応する。
以下、このテンソルから巡回する四つの隣接像を取る。$K,H$ を変換先の $rk,rm$ に相当する
$(r+1/r+w+1/w)/4,(r+1/r-w-1/w)/4$ とする。
復元式で $s,t$ は単元である。十五係数が非零なので $H\ne0$, $K\ne\pm1$ である。
通常条件の失敗には $K=0$ または $K+H,K-H$ の一方が $0,1,-1$ となることが必要である。
実際 $r+1/r=2(K+H)$, $w+1/w=2(K-H)$ で、$w=\pm r^{\pm1}$ のときは $KH=0$ となる。
これを各曲線に代入すると、検査すべき値は次に限られる。
\[
\begin{array}{c|l|l}
 \text{変換先} & H^2 & K\text{ の候補}\\\hline
 n^2=1 & K^2-2 & 0,\ 3/2,\ -3/2\\
 n^2=-1 & -K^2 & 4K^4+1=0\\
 \xi^2+\xi+1=0 & (K^2-1)^2 & 0,\ 2,\ -2,\ K^2\pm K-1=0
\end{array}
\]

\path{computations/verify_rank_fifteen_regular_orbits.py} はこれら全ての還元を
$\mathbb F_{17^{16}}$ で列挙する。最後の行では
$n^2=(K^2-1)(1\pm\sqrt{-3})/2$ である。各 $K,H,n$ に対し
$r+1/r,w+1/w$ から $r,w$ を求め、$s^2=r,t^2=w$ の全根を取る。
全て単数係数のモニックな方程式で整に選べ、逆元の分母 $2,s,t,n$ はこの還元で単元である。
符号変換で同じ剰余軌道に入るものは一代表にまとめる。
係数 $\gamma$ から直接
\[
 W_{ijku}=\boldsymbol1_{i\oplus k=j\oplus u}
  \sum_{v=0}^3 \gamma_{(i\oplus k)+4v}(-1)^{v\cdot k+\beta v\cdot u}
\]
を作り、各始点で隣接二文字像の基底を求める。全代表で各隣接像は六次元、
長さ4,5の経路は $35,56$ 次元、四次中心の文字次元は $(3,2,0,0)$ で、
四つの中心四次式による商は四始点とも長さ13まで定理\ref{thm:rank-fifteen-small-residues}の
次元列を与える。この計算は非通常倍率でのGram逆行列を使わない。

元の通常窓の経路と四次中心には一般点から来る次元の下界があり、
変換後の四文字テンソルも基底変更でその下界を保つ。上の剰余で隣接像の階数が
落ちないので四つの関係空間が持ち上がり、経路と中心の次元の一致から中心切断も持ち上がる。
よって前定理と同じ自由性・Koszul性の判定が適用される。
根交換と反対代数の対応によりこの認証は元の倍率の全根・両表示に及ぶ。
\end{proof}

% !TeX root = ../papers/nc-p3.tex
\section{階数15の有理な同型商と正則性の二変数への降下}
\label{sec:rank-fifteen-rational-quotient}
\begin{theorem}\label{thm:rank-fifteen-rational-quotient}
第一表示の通常な階数15窓で $U=zk^2$, $V=zm^2$ とする。次の量を取る。
\[
\begin{aligned}
 x&=\frac{V(U+3)}{(U-1)^2}, &y&=\frac{V^3}{(U-1)^4},\\
 p&=\frac{Vzkm(U-9)}{(U-1)^3},&
 q&=\frac{(2n^2+1-U)(z-1)(z\lambda^2-1)}{4z\lambda(U-1)^2}.
\end{aligned}
\]
これらは $U=-3$ も含む全通常点で定義される。
二積の生成次元16の開集合で、二窓が添字保持同型である必要十分条件は
$(x,p,q)$ の一致である。一般点では $x$ と $T=q+2p/(3\sqrt3)$ の一致だけでよい。
従ってこの階数層の添字保持同型類の有理関数体は $k(x,T)$ である。
生成次元の仮定を外しても $(x,p,q)$ の一致は同型の十分条件である。
第二表示には倍率反転で対応する第一表示の式を使う。
\end{theorem}
\begin{proof}
$k$ は標数零の代数閉体とする。前々節の規格化テンソルのパラメータ空間を
$U\ne1$, $V\ne0$ まで広げたものを $X$ とする。具体的には $s,t$ が単元、
$r=s^2,w=t^2$, $K=(r+1/r+w+1/w)/4$, $H=(r+1/r-w-1/w)/4$、
$n^4+(1-K^2)n^2+H^2=0$ であり、$U=K^2$, $V=H^2$ である。
$n$ も単元で、十五係数は非零になる。前々節の192変換は $X$ にも作用する。
実際、変換後も全係数は非零で、八量の復元式が $s,t$ の逆元と $n$ を与え、
$b(a+b)\ne0$ と $LM\ne0$ により $U\ne1$, $V\ne0$ を保つ。
$X$ は既約である。$n\pm H/n$ は $K^2-1\pm2H$ の平方根で、互いから $n$ を復元できる。
$s=2$ で二式を $t$ の既約分数にすると、二つの分子は互いに素な平方自由八次式で、
分母は平方である。従って $k(s,t)$ 上で二つの平方類が独立で、四次式は既約である。

八量の斉次表示で $c=a+b$ として
\[
\begin{gathered}
 P_2=2(a^2+b^2+c^2),\quad P_3=4abc,\quad D_3=(b+c)((b+c)^2-9a^2),\\
 h=4LM,\qquad E=16(L^2-M^2)(B^2-A^2)(D^2-C^2).
\end{gathered}
\]
$a=1$ に戻すと $P_2=U+3$, $P_3=U-1$, $h=H$ である。$E$ は
$(2n^2+1-U)(r-1/r)(w-1/w)/4$ になる。
\path{computations/verify_rank_fifteen_rational_quotient.py} は192変換で
$P_2,E$ が不変、$P_3,D_3$ の符号変化の積が $h$ の符号変化であることを五つの二次式の上で検算する。
従って $x=h^2P_2/P_3^2$, $y=h^6/P_3^4$, $p=h^3D_3/P_3^3$, $q=E/P_3^2$ は不変である。
直接展開すると
\[
 p^2=x^3-27y,\qquad q^2=(1-x)^2-4y
\]
となる。

これらの式で定義される $y\ne0$ の曲面の座標環を $R$ とする。
$y$ を消去し $T_{\pm}=q\pm2p/(3\sqrt3)$ と書けば
\[
 T_+T_- = f(x)=(1-x)^2-4x^3/27=-(x-3)^2(4x-3)/27.
\]
これは既約な超曲面で特異点は高々 $(x,T_+,T_-)=(3,0,0)$ なので、Serreの判定から正規である。
$T_+\ne0$ なら $T_-=f(x)/T_+$ から $p,q$ と $y$ が復元できる。

$X$ は $(x,y)$ 平面の $y\ne0$ の上で有限である。実際
\[
 (U+3)^3=(x^3/y)(U-1)^2,\qquad V^3=y(U-1)^4 .
\]
第一式はモニックで $U=1$ での値は64であるから $U-1$ の逆も整である。
$V$ とその逆、次いで $K,H,s,t,n$ と必要な逆元はモニックな方程式で整になる。
一般の $(x,y)$ に対し $U$ は三通り、各々で $V=x(U-1)^2/(U+3)$ である。
符号付きの $(K,H)$ は四通り、$s,t$ は各四通り、$n$ は四通りで、
これらは全て相異なる $X$ の点になる。従って関数体の次数は768である。

$X$ 上の $t\mapsto1/t$ は $p$ を保ち $q$ を反転し、$(s,t)\mapsto(t,s)$ は
$p$ を反転し $q$ を保つ。いずれも $n,x,y$ を固定するので
$k(x,y,p,q)$ の $k(x,y)$ 上の次数は4である。よって上の整域 $R$ が $k[X]$ に入り、
その分数体が192群の固定体になる（作用の忠実性は前々節の192個の相異なる像から従う）。
$k[X]$ は $R$ 上整で、$R$ は正規だから $k[X]^{G}=R$ である。
有限群の不変式は代数閉体上で相異なる有限軌道を分離する。従って例外ファイバーでも
三値の一致は規格化テンソルの192軌道の一致と同値である。前々節の添字保持同型判定と
$T_-=f(x)/T_+$ を合わせて主張を得る。
\end{proof}

\begin{corollary}\label{cor:rank-fifteen-regularity-quotient}
生成次元を問わず、同じ $(x,y)$ を持つ通常な階数15窓の標準Koszul左右AS正則性は一致する。
この二変数の一般ファイバーにある四つの添字保持同型類は、一つの窓、その一位置移動、
およびそれぞれの反対代数で尽くされる。
\end{corollary}
\begin{proof}
$t\mapsto1/t$ は $\lambda\mapsto1/(z\lambda)$ によって同じ根の反対表示を与え、
通常倍率を保つ。次に $s,t$ の交換では八量の $M,A,C$ の符号だけが変わる。
ラベル $(h,j)$ におけるその符号は
$(-1)^{h\cdot j}(-1)^{h\cdot 2}$ である。第一因子の $X_hZ_j$ の転置が初めの符号を、
対角符号行列の共役が後の符号を与える。第一因子の転置は四文字の位置 $0,2$ を交換するので、
四位置の巡回順序を反転する。同じ根の反対表示の語反転とは一位置移動だけ異なる。
この操作は規格化テンソルで直接成立する。

同じ $(x,y)$ での $(p,q)$ は独立な符号を除けば一致する。必要ならまず $t$ の逆転を、
次に $s,t$ の交換を施すと、目標とは $X$ で同じ192軌道に入る。途中の倍率が非通常でも、
最後の交換と192変換はテンソルの転置と独立共役として実行できる。
従って二つの通常窓の関係空間は巡回移動、語反転と基底変更で対応し、左右正則性が一致する。
一般には両符号がいずれも二択で、定理の必要条件により四類は相異なる。
\end{proof}

\begin{corollary}
正則性を検査する $(x,y)$ 平面では、次の四曲線に落ちる全ての通常窓が
標準Koszul左右AS正則である。
\[
\begin{aligned}
 &4y=(x-1)^2,\qquad y=(x-4)^2,\\
 &x^4-x^3+32x^2y-160xy+256y^2+125y=0,\\
 &x^4+x^3+32x^2y+256y^2-27y=0 .
\end{aligned}
\]
\end{corollary}
\begin{proof}
第一式では $q=0$ で、通常条件から $2n^2+1-U=0$、従って分岐上である。
第二式の左辺から右辺を引いて引き戻すと、単元を除いて
$((U-1)^2-V)F_3$ になる。後の二式の左辺は同様に
$(U-V-2)F_1$, $(U+V)F_2$ となる。
コードはこれらの因数分解を分母を払って検算する。分岐の全点正則性と
定理\ref{thm:rank-fifteen-regular-orbits}を適用する。
\end{proof}

% !TeX root = ../papers/nc-p3.tex
\section{階数15の全通常点を分離する三つの添字保持不変量}
\label{sec:rank-fifteen-global-invariants}
\begin{theorem}\label{thm:rank-fifteen-global-invariants}
正の固定曲線の任意の二つの通常な階数15窓が添字保持同型である必要十分条件は、
前節の $(x,p,q)$ が一致することである。二積の生成次元に条件は要らない。
特にこの一判定は分岐、座標端の例外水準、および異なる倍率間の比較を含む。
\end{theorem}
\begin{proof}
十分性は定理\ref{thm:rank-fifteen-rational-quotient}による。必要性のため
$H\in M_4\otimes M_4$ を任意の基底で規格化した反対切断とする。
$T$ をこれに対応する線型写像、$S_1=T^*T$, $S_0=TT^*$ とし、
$\Gamma$ を第一行列因子の転置とする。$D(H)$ は
節\ref{sec:rank-fifteen-endpoint-orbits}の二置換による五重縮約、
$B(S_j)$ は命題\ref{prop:rank-fifteen-oriented-cubic}の対称三次形式の縮約分子である。
次の量は全て元のテンソルから計算できる。
\[
 A_2=\frac{\operatorname{tr}H^2}{16},\qquad
 A_3=\frac{\operatorname{tr}H^3}{96},\qquad
 A_4=\frac{A_2^2-\operatorname{tr}(S_1^2)/64}{3}.
\]
すると前節の座標は次の式を持つ。
\[
\begin{aligned}
 x&=A_4 A_2/A_3^2, &y&=A_4^3/A_3^4,\\
 p&=-\frac{A_4\bigl(D(H)-D(\Gamma H)\bigr)}{128 A_3^3},
 &q&=\frac{B(S_1)-B(S_0)}{3072 A_3^2}.
\end{aligned}
\]
独立共役で全ての縮約は不変である。転置を取った因子の共役は逆転置行列の共役に移るので
$D(\Gamma H)$ も不変である。各分数の分子と分母の $H$ に関する次数は等しい。
規格化した係数表示では $A_3=U-1$ であり、階数15の全通常点で単元である。

恒等式を検算する。前節の斉次八量で $A_2=P_2$, $A_3=P_3$, $A_4=h^2$ となる。
実際、$P_j$ をここではラベル $j$ のPauli基底テンソルとすると、
$P_j^2=(-1)^{e(j,j)}I$ である。
$\mu_j=(-1)^{e(j,j)}\gamma_j^2$ と置けば、$S_1,S_0$ は対応する各局所Pauli基底で固有値
$16\mu_j$ を持つ。これは $T$ が一つのPauli行列を他方の行列へ、そのトレース双対による
係数を掛けて送ることから従う。よって $\operatorname{tr}H^2/16=\sum_j\mu_j$、
$\operatorname{tr}S_1^2/256=\sum_j\mu_j^2$ である。$H^3$ のトレースには
$j,k,j+k$ の組だけが寄与し、両因子の積の符号を掛ける。

$B(S_1)$ では入力側で可換な非零三ラベル $j,k,j+k$ の順序対 $(j,k)$ ごとに、
三つの固有値の積が一回寄与する。反可換な対の寄与は零である。
実際、局所Pauli基底 $E_j$ のトレース双対は $E_j/(4E_j^2)$ であり、可換な場合は
$d(E_j,E_k,E_{j+k})$ とその双対での値の積が1になる。
$B(S_0)$ では出力側の可換性を使う。この差の和を取ると $3072E$ となる。
五重縮約では入力の二閉路によりラベル列が $(a,c,b,a+b,c)$ となる。
第一因子の転置で符号が変わる項に限り、両因子で各閉路の積の符号を掛ける。
三閉路のトレース係数 $4^3$ と、差を取る因子2を戻すと
$D(H)-D(\Gamma H)=-128hD_3$ を得る。
\path{computations/verify_rank_fifteen_global_invariants.py} はこれらの多項式恒等式を
五つの二次式で簡約して検算する。さらに101上の四つの枠ではGram表示から
$H,T$ を再構成し、閉路の直接縮約とも比較する。
これで四つの分数が前節の値と一致し、全点での必要性が従う。
\end{proof}

\begin{corollary}
通常な階数15窓が語反転で得る反対窓と添字保持同型となるのは分岐上に限り、
一位置移動した反対窓と添字保持同型となるのは $U=9$ のときに限る。
二つの軌跡に共通の通常点はない。
\end{corollary}
\begin{proof}
前節の二操作で反転される座標はそれぞれ $q,p$ であり、
規格化テンソルでの群軌道の一致が十分である。上の縮約は、変換先の倍率が非通常でも
そのテンソルに直接適用できる。通常域では $q=0$ は $2n^2+1-U=0$ すなわち分岐で、
$p=0$ は $U=9$ である（$zkm\ne0$）。
両方が成り立てば $K=\pm3$, $H=\pm4$ なので $K+H,K-H$ の一方が $\pm1$ となる。
すると $z=1$ または $z\lambda^2=1$ となり通常条件に反する。
\end{proof}

% !TeX root = ../papers/nc-p3.tex
\section{正則性平面の小標数の全有理点}
\label{sec:rank-fifteen-quotient-residues}
\begin{theorem}\label{thm:rank-fifteen-quotient-residues}
剰余標数が $3,5,7$ のいずれかの付値を考える。標数零の通常な階数15窓の
正則性座標 $x,y$ が整で、$y$ は単元、両者の剰余が素体に属するとする。
この窓は標準Koszul左右AS正則であり、四つの中心四次式の多項式環上で
各始点の経路の直和が階数256の自由加群となる。
倍率や四次根自体の剰余には素体条件を課さない。
\end{theorem}
\begin{proof}
$v=V/(U-1)^2$ とすると正則性平面からの一つの持上げは
\[
 v(x-4v)^2=y,\qquad U=x/v-3,\qquad V=y/v^2
\]
により計算できる。$2,y$ は単元なので、三次式の最高係数と定数項は単元である。
従って $v,x-4v,U-1,V$ は全て単元である。次に $K^2=U,H^2=V$、
$r+1/r=2(K+H),w+1/w=2(K-H)$、$s^2=r,t^2=w$ と
$n^4+(1-U)n^2+V=0$ の根を取る。$s,t,n$ とその逆は整である。
これらは順にモニックな方程式なので、代数閉体の付値を延長すれば指定した剰余根に
落ちる根が各段で少なくとも一つある。三次式や四次式の還元が重根を持つ場合も同じである。

\path{computations/verify_rank_fifteen_quotient_residues.py} は各剰余標数の
全ての $(\bar x,\bar y)\in\mathbb F_p\times\mathbb F_p^\times$ を調べ、
上の方程式を有限拡大体で順に解く。検証した点は順に6,20,42個である。必要になった拡大次数は2,4,6,12で、
各次数の証明書の個数は標数3で $(1,3,1,1)$、標数5で $(2,10,6,2)$、
標数7で $(5,21,12,4)$ となる。
選んだ一つの $(s,t,n)$ から規格化テンソルを作り直し、その隣接像で
四始点の関係を生成する。隣接階数6、長さ4,5の経路次元 $35,56$、
中心の文字次元 $(3,2,0,0)$ を確認し、四つの中心式による商の長さ0から13の
次元を全始点で検査する。例えば $p=7,(\bar x,\bar y)=(0,2)$ では
文字零の中心の三自由座標上で $(1,0,1),(0,1,1)$ を使い所定の次元列を得る。

同じ標数零の $(x,y)$ の元の通常窓と、選んだ形式的な持上げのテンソルは
系\ref{cor:rank-fifteen-regularity-quotient}の証明により、基底変更・巡回移動・反転で対応する。
従って持上げ先の隣接像には経路と中心の所定の次元の下界がある。
剰余の階数と四中心商の証明書が定理\ref{thm:rank-fifteen-regular-orbits}と
同じ方法で持ち上がり、元の窓の全次数の主張が従う。
\end{proof}

% !TeX root = ../papers/nc-p3.tex
\section{階数15の通常倍率の像を補う四つの正則窓}
\label{sec:rank-fifteen-formal-boundary}
\begin{theorem}\label{thm:rank-fifteen-formal-boundary}
$y\ne0$, $p^2=x^3-27y$, $q^2=(1-x)^2-4y$ の規格化テンソルの商曲面を考える。
通常倍率の階数15窓の像は、この曲面から次の四点を除いたものにちょうど等しい。
全て $q=0$ 上にあり、$y=(1-x)^2/4$ である。
\begin{enumerate}
\item $x=3,3/4$ の各点（$p=0$）。
\item $U=3\pm2\sqrt2$, $x=(U+3)/4$,
$p=-U^2(U-9)/(U-1)^3$ の各点。
\end{enumerate}
四点の規格化テンソルも標準14641型のKoszul左右AS正則窓を与え、
各始点で四中心四次式の多項式環上自由階数256となる。
最初の二点の最小周期は1、残り二点の最小周期は2である。
\end{theorem}
\begin{proof}
$K=rk$, $H=rm$, $r=s^2$, $w=t^2$ とする。
形式的表示では $U=K^2\ne1$, $V=H^2\ne0$ を常に満たす。
通常条件はこの上で
\[
 K\ne0,\qquad K+H\notin\{0,1,-1\},\qquad K-H\notin\{1,-1\}
\]
に等しい。実際 $2(K+H)=r+1/r$, $2(K-H)=w+1/w$ である。
$w=-r,-1/r$ は $K=0$ に一致し、$w=r,1/r$ は $H=0$ に一致する。
また $z=-1$ は $K+H=0$ を与え、$z=1$ と $z\lambda^2=1$ は各和差が $\pm1$ となる条件である。

$q\ne0$ なら $K\pm H=\pm1$ は $q$ の積表示から起こらない。
$p\ne0$ なら三次式 $v(x-4v)^2-y$ の判別式は $256yp^2\ne0$ である。
その三根のうち $3v^2-xv+y\ne0$ となるものを取ると $U=x/v-3$ と $V=y/v^2$ は異なる。
$p\ne0$ により $K\ne0$ で、符号のどちらでも $K+H\ne0$ となる。
$KH$ の符号で所定の $p$ を、$w$ の逆転で所定の $q$ を選べる。
$p=0$, $q\ne0$ なら $x\ne0$ で $U=9$, $V=16x/3$ を選べばよい。
$K=3$ として $H$ の符号で $K+H\ne0$ にでき、同様に所定の $q$ になる。

$q=0$ で通常なら $2n^2+1-U=0$ が必要である。よって
$V=(U-1)^2/4$, $U=4x-3$, $H=\pm(U-1)/2$ である。
このとき $((K+H)^2-1)((K-H)^2-1)=(U-9)(U-1)^3/16$ なので、
$U=9$ と $K=0$ は除かれる。それ以外では $p$ が $KH$ を一意に指定する。
最後の例外 $K+H=0$ は $U=V$ と $KH=-U$ に相当し、
$(U-1)^2=4U$ を解いて記載した二点を得る。全ての除外条件を避けたときは
二次式から $r,w,n$ を復元できるので十分性も従う。

四点では形式的表示 $K=0,3,1+\sqrt2,1-\sqrt2$ を順に用いる。
初めの二つで $H=(U-1)/2$、後の二つで $H=-K$ とし、$n^2=(U-1)/2$ を取る。
\path{computations/verify_rank_fifteen_formal_boundary.py} は17の上で二平方根の符号と
$r,w,s,t,n$ の根を選び、四テンソルを有限体で直接作る。選んだ還元は逐次モニックな
方程式の根として標数零へ持ち上がり、上記の四点を尽くす。
各巡回隣接像の階数は6で、その直交補空間を二次関係とする双対の経路次元は
全始点で長さ0から5まで $(1,4,6,4,1,0)$ となる。
元の経路次元と四次中心の文字次元を確認し、四中心商も全始点で長さ13まで所定の次元列を得る。

この場合の持上げの下界は形式的パラメータ空間 $X$ の既約性から得る。
その稠密な通常窓で隣接階数6なので、剰余で階数6なら標数零の点とその近傍で関係は定階数である。
経路と双対経路の次元は正則一般点からの下界に等しく、四次中心も経路基底を取って同様に持ち上がる。
反対階数15は十五係数の非消滅から従い、各一文字切断の階数も4である
（一文字の台が三次元なら反対階数は高々12）。四次双対の双対は元のテンソル直線なので、
一文字と隣接二文字の切断階数から双対の全次数の完全対合を得る。
四中心商の上界と元の経路の全次数の下界を合わせた自由性・Koszul完全性の議論
（定理\ref{thm:klein-rational-period-four}）が適用できる。
また四次テンソルを関係から回収できるので、三つの縮約不変量はこの四窓でも添字保持不変であり、
通常倍率の窓とも互いにも混同されない。

形式的な $t$ の逆転は語反転である。実際Pauliテンソルの二因子を転置して交換すると、
ラベル $(a,b)$ は $(a,\beta b)$ に移り、符号は $(-1)^{e((a,b),(a,b))}$ になる。
コードはその係数配列が $C,D$ を $-D,-C$ に替えた配列と一致することを確認する。
$s,t$ の交換は既に第一因子の転置と共役である。従って一位置移動は
$(p,q)\mapsto(-p,-q)$ に対応し、二位置移動は元の192軌道に戻る。
軌道の分離と上の縮約不変性から記載した最小周期が従う。
\end{proof}

\begin{corollary}
上の形式的な階数15の族では $q=0$ の曲線の全点が正則である。
その正規化は $j\ne\pm1$ の直線で、
$x=(j^2+3)/4$, $p=j(j^2-9)/8$、$j=\pm3$ が同じ節点に移る。
\end{corollary}
\begin{proof}
通常点では $q=0$ は正則な分岐であり、残る四点は定理で認証した。
$p^2=(x-3)^2(4x-3)/4$ に $j=2p/(x-3)$ を代入すれば表示を得る。
\end{proof}

% !TeX root = ../papers/nc-p3.tex
\section{剰余標数3における正則性直線の全幾何点}
\label{sec:rank-fifteen-residue-line}
\begin{theorem}\label{thm:rank-fifteen-residue-line}
剰余標数3の付値を持つ標数零の体上で、通常な階数15窓の正則性座標が
$\bar x=0$、$y$ は単元を満たすとする。$\bar y$ は任意でよい。
この窓は標準Koszul左右AS正則で、各始点の経路の直和は
四つの中心四次式の多項式環上で階数256の自由加群である
（基礎体は代数閉体に拡大して考える）。
\end{theorem}
\begin{proof}
剰余体の代数閉包で $i^2=-1$ とする。$\bar y=H^6$、
$s^2+s^{-2}=2H$ と選ぶと、$s(s^4+1)\ne0$ である。
$t=is$ により $K=0$、$H=(s^2+s^{-2})/2$ となる。
四次式の根とその相方を
\[
 n=\tfrac12\bigl(s+s^{-1}-i(s-s^{-1})\bigr),\qquad
 N=\tfrac12\bigl(s+s^{-1}+i(s-s^{-1})\bigr)
\]
と取れる。実際 $nN=H$、$n^2+N^2=2=-1=K^2-1$ である。
ゆえに $U=0$, $V=H^2$ とすれば $(x,y)=(0,H^6)$、
$v=H^2$ は剰余三次式 $v(x-4v)^2=y$ の根となる。
こうして直線上の全剰余点に対し、係数を一変数のLaurent多項式で表示できる。

\path{computations/verify_rank_fifteen_residue_line.py} はこの表示を
$\mathbb F_9[s]$ 上で厳密に計算する。$P=s^2+1$, $Q=s^2-1$ とすると
$4s$ 倍した八係数 $(a,b,L,M,A,B,C,D)$ は
\[
 (4s,-2s,P-iQ,P+iQ,iQ-P,iQ+P,iP-Q,iP+Q).
\]
巡回する各隣接像の全生成行をここから取る。
分数体上の経路次元は長さ0から5で $(1,4,10,20,35,56)$、
四次中心の文字別次元は $(3,2,0,0)$ である。
この消去で逆にした主元の分子の全因子を記録する。
文字零の三自由座標における $(1,3,5),(2,7,11)$ と文字一の二基底を
四つの中心式に選び、各始点で分母を払う。
記録した主元が単元なら、これらは真の中心式を始点ごとに単元倍したものである。

次に各始点で、始点の四中心式を課した経路双対を長さ13まで計算する。
$\mathbb F_9[s]$ を基底 $1,i$ の $\mathbb F_3[s]$ 加群として扱い、
各段で終端の二次関係と、長さ4では中心式も課す。
得られる多項式行列の列像の基底を $C$ とする。$C$ の行加群の基底を並べた
正方行列の行列式は、単数を除いて元の行列の最大非零小行列式の最大公約元である。
これは二つの基底変更がそれぞれの生成加群を保つこととCauchy--Binet公式から従う。
その根の外では核が特殊化と可換なので、核と末尾文字ごとの射影を次段へ渡せる。
全始点で、係数拡大の次数2で割った商次元列は
\[
 (1,4,10,20,31,40,44,40,31,20,10,4,1,0)
\]
である。主元因子と全段の行列式の因子を合わせ、$\mathbb F_9$ 上で分解する。

$s\mapsto us,\ u/s$（$u^4=1$）は $H^6$ を保つ。
従って同じ $\bar y$ の認証には一つの代表で足り、因子の全軌道が残った場合だけ
直接調べればよい。$s(s^4+1)=0$ を除くと完全に残る既約因子の軌道は
次数1が一つ、次数2が五つである。各代表の一根で、全四始点の中心次元と
四中心商の所定の次元列を有限体で再検算した。他の共役根は
$\mathbb F_9$ 上のFrobeniusで移る。超越的な剰余値には主元や行列式の零点はない。
最後に定理\ref{thm:rank-fifteen-quotient-residues}のモニックな逐次持上げを
選んだ剰余代表に適用する。同じ標数零の $(x,y)$ を持つ通常窓への移送と、
中心式の持上げ、全次数の自由性・Koszul性の判定により結論を得る。
\end{proof}

% !TeX root = ../papers/nc-p3.tex
\section{反対階数12の量子極限と正則性平面の極}
\label{sec:rank-twelve-quantum-limit}
\begin{theorem}\label{thm:rank-twelve-quantum-limit}
Pauli八係数の五つの斉次二次式の射影解のうち $a=b=0$ となるものは32点で、
全て同一の添字保持同型類の四窓を与える。その反対階数は12であり、
四周期の基底変更で次の量子多項式環の窓になる。
\[
 X_jX_k=q_{jk}X_kX_j\quad(j<k),\qquad
 (q_{01},q_{02},q_{03},q_{12},q_{13},q_{23})=(1,1,-i,-i,1,-1),
 \qquad i^2=-1 .
\]
従ってこの境界類は標準Koszul左右AS正則で、最小周期は1である。
\end{theorem}
\begin{proof}
$L^2+M^2=A^2+B^2=C^2+D^2=0$、$LM=-AB=-CD$ なので
$L=0$ なら全係数が零になる。$L=1$ とすると
$M=\pm i$, $A^4=C^4=1$, $B=-M/A$, $D=-M/C$。
\path{computations/verify_rank_twelve_quantum_limit.py} は192変換による
$(L,M,A,B,C,D)=(1,-i,i,1,i,1)$ の射影軌道がこの32点と一致することを検算する。
各変換は四位置の基底変更を与える。非零なPauli係数は十二個であり、
両側のPauli行列がそれぞれ基底をなすので反対切断の階数は12である。

代表の各文字の基底で対 $(0,1),(2,3)$ に
$\left(\begin{smallmatrix}1&1\\1&-1\end{smallmatrix}\right)$ を施し、
奇数位置では文字を $j\mapsto j\oplus3$ で付け替える。四次テンソルの
非零な24成分は共通倍率を除き $\pm1\pm i$ である。各巡回隣接像は
六つの二項式を与え、順序を定理と同じに取ると係数は
\[
\begin{array}{c|rrrrrr}
0&1&1&-i&-i&1&-1\\
1&1&-1&i&i&-1&-1\\
2&-1&-1&-i&-i&-1&1\\
3&-1&1&i&i&1&1
\end{array}
\]
となる。各位置 $p$ の対角基底倍率を $d_{p,j}$ とする。
隣接比 $d_{p+1,j}/d_{p,j}$ を順に
$(1,1,1,1)$、$(1,1,-1,-1)$、$(1,-1,-1,1)$、$(1,-1,1,-1)$ とすると、
全行が最初の行になり、各文字で四つの比の積は1である。
これらの計算は $2$ が単元の任意の剰余体でも成立する。
順序単項式は基底をなし、四乗 $X_j^4$ による商には各指数が4未満の単項式が残る。
各単項式 $X^e$ に沿う四周期の中心切断は
$\prod_j q_{jk}^{e_j}$ が $k$ によらないとき一次元、それ以外は零である
（逆順の係数は逆数とする）。共通値の四乗は1なので一周の整合も成立する。
次数4では四つの純四乗と $e=(1,1,1,1)$ の計五本を得る。
\end{proof}

\begin{theorem}\label{thm:rank-fifteen-pole-cone}
標数零の通常な階数15窓の正則性座標 $x,y$ に付値 $\nu$ を取り、
剰余標数は零または奇素数とする。$\nu(0)=+\infty$ として
\[
                 \nu(y)<\min\{0,2\nu(x)\}
\]
ならばこの窓は標準Koszul左右AS正則で、各始点で四中心四次式の
多項式環上自由階数256となる。
\end{theorem}
\begin{proof}
$D=x^3/y$ として $(U+3)^3=D(U-1)^2$ を解く。
$D$ が整なら全根は整で $U-1$ は単元である（$4$ が単元）。
$V^3=y(U-1)^4$ より $V$ は極を持つ。$D$ が極なら $d^2=1/D$ として
$U=1+dE$ と置く。方程式 $(4+dE)^3=E^2$ の剰余根 $E=\pm8$ は単根であり、
付値を延長して単元根を選べる。このとき
$3\nu(V)=\nu(y)-2\nu(D)=3\nu(y)-6\nu(x)<0$ である。
いずれも $U$ は整、$V$ は極である。$x\ne0$ では
$V=x(U-1)^2/(U+3)$ を取り、$x=0$ では $U=-3$ として上の三乗根を取る。

$K^2=U$, $H^2=V$ とし、$r+1/r=2(K+H)$、$w+1/w=2(K-H)$
で極を持つ方の根を取る。$s^2=r$, $t^2=w$ の符号を選んで
$\overline{t/s}=i$、$\overline{H/s^2}=1/2$ となる。
$n/s$ の整方程式の還元は $e^4=-1/4$ なので、単根 $e=(1+i)/2$ を指定できる。
八係数全体を $4/s$ 倍すると $a,b$ の剰余は零、
$(L,M,A,B,C,D)$ の剰余は
$(2e,1/e,i-1,i+1,i-1,i+1)$ である。
$1+i$ で割ると前定理の代表を得る。特に $K=2$ を固定して
$H$ を無限遠へ動かせば、通常な階数15からこの境界類への退化が実現する。

剰余の六次元の隣接像は持ち上がり、長さ4,5の経路次元は $35,56$、
四次中心の全次元は5で、いずれも一般点の下界に等しい。
従って剰余の四つの純四乗を中心切断に持ち上げられる。
剰余の四中心商の基底は上で与えたので、長さ13で消え、長さ0から12で
$(1+t+t^2+t^3)^4$ の係数を与える。一般点からの経路次元の下界と合わせて
全次数の自由性とKoszul性を得る。選んだ $(U,V)$ の表示と元の通常窓は
同じ $(x,y)$ を持つので、基底変更・巡回移動・反転により元の主張が従う。
\end{proof}

\begin{corollary}
$p=0$ で $y$ が代数的なら、未認証点の候補は、$y$ が全ての奇素数上で整であり、
かつ3の上の全ての素点で正の付値を持つ場合に限られる。
\end{corollary}
\begin{proof}
$x^3=27y$ なので奇素数上で $y$ が極なら前定理の厳密不等式を満たす。
3の上で $y$ が単元なら $\bar x=0$ で定理\ref{thm:rank-fifteen-residue-line}が適用できる。
\end{proof}

% !TeX root = ../papers/nc-p3.tex
\section{Pauli射影閉包の境界での正則性の完全判定}
\label{sec:pauli-projective-boundary}
\begin{theorem}\label{thm:pauli-projective-boundary}
標数零の代数閉体上で、規格化テンソルの五つの二次式が定める射影曲面を考える。
$c=a+b$ とする。開集合 $abcLM\ne0$ の境界において、反対階数と
その四次テンソルの巡回隣接像が与える窓の標準14641型正則性は次の通りである。
\[
\begin{array}{c|c|c|l}
 & LM &\text{反対階数}&\text{判定}\\\hline
 abc=0,\ (a,b)\ne(0,0)&\ne0&14&\text{全点正則}\\
 a=b=0 &\ne0&12&\text{一つの量子同型類}\\
 abc\ne0&0&9&\text{長さ4の次元が36以上}\\
 abc=0,\ (a,b)\ne(0,0)&0&4&\text{隣接階数4}
\end{array}
\]
正則な両層では四中心四次式の多項式環上で各始点の経路の直和が自由階数256になる。
\end{theorem}
\begin{proof}
反対階数は非零なPauli係数の個数に等しい。$LM=-AB=-CD$、
三つの平方和はそれぞれ $bc,ac,ab$ である。$a=b=LM=0$ は射影点を与えない。
従って支持を数えれば表の四層で尽きる。192変換は $a,b,c$ の各々を
単元倍を除いて新しい $a$ に取れる。三つのうちちょうど一つが零のとき、
零のものを新しい $a$ とすることも、非零のものを新しい $a$ とすることもできる。

階数14では $a=1$ と取ると $b=0,-1$、従って $K=1,-1$ である。
復元式 $s=B-A+D-C$, $t=B+A+D+C$, $n=2L$ とその逆元の式は
$U=1$ でも五つの二次式だけから従う。$H=4LM\ne0$、$n^4=-H^2$ である。
$K+H\ne0,\pm1$ かつ $K-H\ne\pm1$ なら
定理\ref{thm:rank-fifteen-covers}の通常な階数14窓である。
Pauli係数の恒等式はこの通常域でも成り立つ。実際Gram表示の変換と
四次根による係数式はいずれもここで分母が単元である。
正則性と五次元の四次中心は定理\ref{thm:rank-fourteen-all-ordinary}と
定理\ref{thm:rank-fourteen-center}による。

残るのは $H^2=1,4$ に限られる。
\path{computations/verify_pauli_projective_boundary.py} は101上で $K,H$ の両符号、
$r,w,s,t,n$ の全根から形式的なテンソルを直接作る。
根の選択は順に256通りと128通りで、規格化した192軌道では二代表と一代表になる。
全代表で四始点の双対次元 $(1,4,6,4,1,0)$、経路長4,5の次元 $35,56$、
中心の次元 $(3,2,0,0)$、四中心商の長さ13までの標準列を確認した。
標数零の各根はこれらのいずれかに還元される。
隣接階数6と一文字階数4（反対階数14より従う）から双対の完全対合を得る。
関係と中心の定階数による持上げ、一般点からの経路次元の下界と四中心商による上界から
全次数の主張が従う。階数12は定理\ref{thm:rank-twelve-quantum-limit}である。

階数9では $a=1$ とし、$P=(u+u^{-1})/2$, $Q=(u-u^{-1})/2$ と置く。
$b=Q^2,\ c=P^2$ で、$(L,M),(A,B),(C,D)$ は各対の一方だけが非零である。
その三つの非零成分は順に $\pm PQ,P,Q$ と取れる。実際後の二成分を
$P,Q$ として $P^2-Q^2=1$ なので $u=P+Q$ と復元できる。
非零成分の位置と初めの符号による十六表示の各々で、有理関数体上の長さ4の
全ての関係挿入を文字別に行消去した。始点0で得る商次元は
$(10,10,8,8)$ で、合計36である。特殊点では元のテンソルの生成行を用いる
全挿入行列の階数は増えない。従って全点で標準次元35を超える。
階数4では $a=0,b=1$ とすれば $L,M$ の一方のみが $\pm1$ である。
この四テンソルの隣接階数は四始点とも4である。

これらの境界層はいずれも高々一次元である。五つの斉次方程式の
$\mathbb P^7$ における各既約成分は次元2以上なので、境界に含まれる別成分はない。
従ってこれは稠密な形式的開集合の射影閉包の境界全体の判定である。
\end{proof}

\begin{corollary}\label{cor:pauli-rank-fourteen-valuation-wall}
剰余標数が零または101で $\nu(y)=2\nu(x)<0$ なら、標数零の通常な階数15窓は
標準Koszul左右AS正則で、各始点で四次中心式四本の上に自由階数256となる。
\end{corollary}
\begin{proof}
前節の $D=x^3/y$ は極で、$U$ が1に還元される根を取れる。
$V=(U+3)^2 y/x^2$ より $\bar V=16\,\overline{y/x^2}\ne0$ である。
順にモニック方程式の根を取れば、表示は階数14の形式的なテンソルに還元される。
剰余標数零では定理の証明から、各剰余窓の四次中心は五次元で
四中心商の標準列が得られる。101の $\bar V=1,4$ でも上の直接計算が適用できる。

101の通常な階数14還元では定理\ref{thm:rank-fourteen-geometric-reduction}の
多項式核計算が四中心商を与える。ここでは階数15方向に中心が持ち上がるよう、
剰余の全四次中心の次元も検査する。
\path{computations/verify_rank_fourteen_wall_centers.py} は既存の階数14の
関数体上で、長さ5までの各経路行列と四つの中心方程式の主元の積を取る。
これは選択した行と列の実際の小行列式である。ノルムの分子と分母、
関係の分母を101で分解すると、最高次数12の既約因子33個が残る。
その外ではCramerの公式により経路基底と中心の五自由座標が特殊化できる。
各因子の通常な根でも両方の $i$ について全始点で経路・中心・四中心商を再計算した。
ゆえに全ての還元で中心が定階数となり、四切断が階数15の持上げに延びる。
前節と同じ自由性とKoszul性の判定が適用できる。
\end{proof}

% !TeX root = ../papers/nc-p3.tex
\section{三次曲線 $p=0$ の全点正則性}
\label{sec:rank-fifteen-cubic-zero-disk}
\begin{theorem}\label{thm:rank-fifteen-cubic-zero-disk}
標数零の形式的な階数15曲面で、$p=0$、すなわち $x^3=27y$ 上の
全ての窓は標準Koszul左右AS正則である。各始点で経路の直和は
四中心四次式の多項式環上自由階数256となる。
この曲線の正規化は $j\ne\pm1$ の直線で、
\[
 x=3(1-j^2)/4,\qquad y=(1-j^2)^3/64,\qquad q=j(j^2+3)/4 .
\]
$x=3,3/4$ の二つの点は最小周期1、残る全点は最小周期2である。
\end{theorem}

\begin{lemma}\label{lem:rank-fifteen-cubic-zero-disk}
剰余標数3の離散付値を持つ標数零の通常な窓で、
ある形式的表示が $0<\nu(H)$、$\nu(K-3)\ge25\nu(H)$ を満たすとする。
この窓は定理の正則性と自由性を持つ。
\end{lemma}
\begin{proof}
まず $K=3$、局所変数を $H$ とし、
$D=\mathbb Z_{(3)}[\sqrt2]$ 上の冪級数で考える。
$3$ はこの二次拡大で惰性を持つ。$P=s+s^{-1}$、$Q=s-s^{-1}$、
$P'=t+t^{-1}$、$Q'=t-t^{-1}$ の解を
\[
 P=\sqrt{8+2H},\quad Q=\sqrt{4+2H},\quad
 P'=\sqrt{8-2H},\quad Q'=\sqrt{4-2H},
\]
とし、四次根を $n=\sqrt{(8+\sqrt{64-4H^2})/2}$
と取る。定数項を順に $2\sqrt2,2,2\sqrt2,2,2\sqrt2$ とし、
内側の平方根は8に取る。いずれも単根である。
$N=H/n$ とし、八係数の規格化式から四つの巡回隣接像を得る。
$H=0$ は階数9のため、二次関係をそのまま特殊化すると
長さ4の余分な経路が現れる。

\path{computations/verify_rank_fifteen_cubic_zero_disk.py} は係数を
$D[H]/H^{25}$ で出発し、単元判定は定数項の二つの有理座標を3で還元して行う。
各経路行列で単元主元を消去する。目標階数に達する前に単元が尽きたら、
残った全ての係数の定数項が $\mathbb Q(\sqrt2)$ で厳密に零と確認し、
その全行を $H$ で割る。$H\ne0$ の窓ではこの行も関係である。
既知の経路次元の下界に達するまで消去するので、自由列の語は
実際の分数体上の経路基底となり、各語の座標を求められる。
長さ1から5で、$H$ による除算の回数の全始点・全文字での最大は
順に $0,0,0,2,2$ となる。従って得た経路基底での四次中心方程式に
使える精度は少なくとも21である。

各文字で中心次元 $3,2,0,0$ の下界に達するまで単元主元だけで消去できる。
よって得た自由座標から実際の中心切断を選べる。
文字零は三自由列の後二本、文字一は両自由列を取る。
各始点でこの四式を長さ4の始端に課し、以後も末尾の二次関係を挿入する。
長さ1から13で、各文字のうち最大の精度損失は全四始点とも
\[
                 0,0,0,2,4,2,2,2,2,2,2,2,0
\]
であり、最後にも一桁残る。各段で選択した消去は必要な単元主元と厳密な
定数項の零だけを使う。切断商の次元上界は $(1+t+t^2+t^3)^4$ の係数となり、
長さ13で消える。中心性による自由加群の上界、未切断経路の一般点からの下界、
定理\ref{thm:klein-rational-period-four}の左右Koszul複体の議論により主張を得る。

一般の $K=3+H^{25}T$ では $T$ は整である。
同じ定数根を使えば四次根と平方根は全て単根で、係数は
$D[T][[H]]$ で上の係数と $H^{25}$ を法として一致する。
各除算後も計算した桁数の合同は保たれる。従って証明で使った定数項は全て
$T$ に依存せず、完備付値環で特殊化して同じ上界を得る。
同じ $(x,y)$ における他の表示には基底変更・移動・反転で移せる。
\end{proof}

\begin{proof}[定理の証明]
$y$ が代数的なら数体の3の上の付値を選ぶ。$\nu(y)<0$ なら
$x^3=27y$ より定理\ref{thm:rank-fifteen-pole-cone}、
$\nu(y)=0$ なら $\bar x=0$ より定理\ref{thm:rank-fifteen-residue-line}を適用できる。
$\nu(y)>0$ では $K=3$, $H^2=16x/3$ とすると所定の $(x,y)$ を与え、
$\nu(H)=\nu(y)/6>0$ である。上の補題がこの最後の領域を尽くす。
超越的な $y$ には $\mathbb Q(y)$ の3進Gauss付値を使えば $y$ は単元である。
通常表示から外れる二点は定理\ref{thm:rank-fifteen-formal-boundary}で認証済みである。
最後に $q^2=(1-x/3)^2(1-4x/3)$ を $j=q/(1-x/3)$ で正規化すれば上式を得る。
$j^2=-3$ の二点は同じ節点に移り、$j=0$ はもう一つの $q=0$ の点に移る。
一位置移動は $(p,q)\mapsto(-p,-q)$ なので、完全な縮約不変量から最小周期も従う。
\end{proof}

\begin{corollary}
この3進の十分条件は正則性平面の二座標からも検査できる。
$\alpha=\nu(3)$, $\beta=\nu(y)>0$, $\delta=\nu(x^3/y-27)$ とする。
$\delta>6\alpha$ かつ $3\delta\ge12\alpha+25\beta$ ならば、
同じ正則性と自由性が成立する。
\end{corollary}
\begin{proof}
$\delta=+\infty$ は定理による。それ以外で $e^2=x^3/y-27$ と取る。
$D=(U+3)^3/(U-1)^2$ から $D-27=U(U-9)^2/(U-1)^2$ である。
$K=3+(e/9)z$ と置く。$\nu(e)>3\alpha$ なので方程式
$(K/3)((K+3)/3)z=K^2-1$ の還元は $2z=-1$ で単根である。
この解で $U=K^2$ とすれば所定の $D$ を与え、$U-1$ は単元である。
$V=x(U-1)^2/(U+3)$、$H^2=V$ より $\nu(H)=\beta/6$、
$\nu(K-3)=\delta/2-2\alpha$ なので補題の不等式が成立する。
\end{proof}

% !TeX root = ../papers/nc-p3.tex
\section{正則性座標 $x=0$ の全点とその添字保持分類}
\label{sec:rank-fifteen-x-zero}
\begin{theorem}\label{thm:rank-fifteen-x-zero}
標数零の階数15のPauli曲面で、$x=0$ の全点は標準Koszul左右AS正則で、
各始点で四中心四次式の多項式環上自由階数256となる。
この軌跡の添字保持同型類は $j\in\mathbb G_m\setminus\{1,-1\}$ で重複なく表され、
\[
 p=\frac{3\sqrt3}{4}(j-j^{-1}),\qquad
 q=\frac{j+j^{-1}}2,\qquad
 y=-\frac{(j-j^{-1})^2}{16}
\]
である。全ての最小周期は2である。$y$ を固定したとき $y=1/4$ で二類、
それ以外の $y\ne0$ で四類となる。
\end{theorem}
\begin{proof}
$x=0$ では $U=-3$、$H^6=256y$ と取れる。
3の上で $y$ が単元なら定理\ref{thm:rank-fifteen-residue-line}、
極を持てば定理\ref{thm:rank-fifteen-pole-cone}が適用される。
残る場合は $\nu(H)>0$ である。$i^2=-1,d^2=3$ とし、$K=id$ を取る。
$D=\mathbb Z_{(3)}[i,d]$ の $d$ 上の剰余体は $\mathbb F_9$ である。
$D[[H]]$ 上の四つの平方根
\[
 P,P'=\sqrt{2(K+1+H)},\sqrt{2(K+1-H)},\qquad
 Q,Q'=\sqrt{2(K-1+H)},\sqrt{2(K-1-H)}
\]
の定数項をそれぞれ $d+i,d+i,-(1+K),-(1+K)$ とする。
$(d+i)^2=2(K+1)$、$(1+K)^2=2(K-1)$ なので全て単根である。
$n=\sqrt{(-4-\sqrt{16-4H^2})/2}$ は内側の定数項が4、外側が $-2i$ の根とする。
これで $n^4+4n^2+H^2=0$ が成立する。

\path{computations/verify_rank_fifteen_x_zero.py} は $1,i,d,id$ の四座標を有理係数で持つ。
定数項が単元である条件は、全座標が3整で、初めの二座標の少なくとも一つが
3で非零になることである。$H$ で割る際は四座標とも定数項が
有理数として零であることを確認する。
前節と共通の経路・中心・商の行列を精度25から計算した。
長さ4,5の経路での精度損失はそれぞれ2、四次中心方程式では零となる。
文字零の後二自由列と文字一の両自由列を選び、四中心商を各始点で長さ13まで
計算すると、各長さの最大損失は前節と同列で最後の精度は1である。
従って中心の持上げと商次元の上界から全次数の正則性・自由性が従う。
より一般に $\nu(K-id)\ge25\nu(H)>0$ の窓にも同じ計算が適用できる。
実際 $K=id+H^{25}T$ としても四つの平方根と四次根の選んだ解は
精度25まで一致し、行操作に必要な全ての定数項が保たれる。

代数的な $y$ には数体の3の上の付値を取り、超越的な $y$ には
3進Gauss付値を取れば、以上の三場合で全点を尽くす。
商の方程式より $q^2-4p^2/27=1$ である。
$j=q+2p/(3\sqrt3)$ は単元で、その逆元は $q-2p/(3\sqrt3)$ である。
$y\ne0$ は $j\ne\pm1$ に等しい。定理\ref{thm:rank-fifteen-formal-boundary}の
四例外は全て $x\ne0$ なので、ここでは全点が通常表示を持つ。
$(x,p,q)$ は完全な添字保持同型不変量で、一位置移動は $j\mapsto-j$、
二位置移動は同型になる。$p$ は常に非零で、$q=0$ は $y=1/4$ に限られるから
最小周期と各ファイバーの個数も従う。
\end{proof}

\begin{corollary}
剰余標数3で $\alpha=\nu(3)$, $\beta=\nu(y)>0$,
$\delta=\nu(x^3/y)$ とする。$\delta>3\alpha$ かつ
$2\delta\ge3\alpha+25\beta$ ならば、通常な階数15窓は同じ正則性と自由性を持つ。
\end{corollary}
\begin{proof}
$x=0$ は定理に含まれる。$x\ne0$ では
$(U+3)^3=(x^3/y)(U-1)^2$ の根は整で、$U-1$ は単元である。
$\nu(U+3)=\delta/3>\alpha$ より $U/(-3)$ は1に還元される。
$K^2=U$ の符号を $K/(id)$ が1に還元されるように取れば
$\nu(K-id)=\delta/3-\alpha/2$ となる。
また $\nu(H)=\beta/6$ なので証明中の精度25の条件を満たす。
\end{proof}

% !TeX root = ../papers/nc-p3.tex
\section{可変の跡が3で消える円板における第六経路と四次中心}
\label{sec:rank-fifteen-variable-zero-disk}
\begin{theorem}\label{thm:rank-fifteen-variable-zero-disk}
標数零、剰余標数3の付値を $\nu$ とする。通常な階数15窓の座標が
$0<\nu(y)<3\nu(x)$ を満たすとき、各始点の長さ6の経路の文字別次元は
$(20,20,22,22)$、総次元は84である（文字は以下の規格化で並べる）。
巡回四次中心の文字別次元は $(3,2,0,0)$ であり、各始点で独立な四本を共通に選べる。
\end{theorem}
\begin{proof}
$(U+3)^3=(x^3/y)(U-1)^2$ は整係数のモニック方程式で、還元は $U^3=0$ である。
$U-1$ は単元であり、$K^2=U$, $H^6=y(U-1)^4$ とすると $K,H$ とも正の付値を持つ（$K=0$ も許す）。
$x\ne0$ では $H^2=x(U-1)^2/(U+3)$ を満たす符号・根を選ぶ。
$x=0$ では $U=-3$ として任意の六乗根を取れる。基礎体を拡大し、$i^2=-1$ を剰余体に取る。
$z^4-2Kz^2+1=0$ の単根 $z_*=1+i$ を持ち上げると、
\[
 K=(z^2+z^{-2})/2,\qquad P_0=z+z^{-1},\quad Q_0=z-z^{-1},\quad n_0=(z^2-z^{-2})/2 .
\]
$z,P_0,Q_0,n_0$ はいずれも単元である。そこで計算中は $z$ を不定元とし、
$R=\mathbb Z[z]_{(3,z^2+z+2)}$ とする。$R[[H]]$ において
$P_0^2\pm2H$, $Q_0^2\pm2H$ の平方根を定数項 $P_0,Q_0$ で、
\[
 n=\sqrt{\bigl(n_0^2+\sqrt{n_0^4-4H^2}\bigr)/2}
\]
を内側の定数項 $n_0^2$、外側の定数項 $n_0$ で一意に選べる。
$n_0^2=K^2-1$ なので、$n,H/n$ は積 $H$、平方和 $K^2-1$ を持つ。
これらが通常表示の平方根を与え、経路と中心は正則性座標で比較できる。

\path{computations/verify_rank_fifteen_variable_zero_disk.py} は
$\mathbb Q(z)[H]/(H^7)$ で平方根と関係を構成する。各主元については分母と分子を
$z_*$ で還元し、ともに非零であることを検査する。$H$ による除算では、
捨てる定数項の分子が $\mathbb Z[z]$ の零多項式であることを確かめる。
従ってこれらの行操作は $\bar z=z_*$ を満たすどの持上げにも適用できる。
長さ1から5の各段の最大精度損失は $(0,0,0,2,2)$ であり、全始点で経路基底を得た後に
精度3が残る。この基底で四次中心方程式を文字ごとに消去すると、単元主元だけで
残る自由列数が順に $3,2,0,0$ に達する。文字零の後二自由列と文字一の両自由列に
対応する中心式を各始点で評価すると、単元主元だけで商の長さ4の文字別次元は
$(7,8,8,8)$ に達する。四本が作る行はこれ以上の階数を持ち得ない。長さ6では四始点と四文字の各挿入行列の
幅が56で、最大2回の $H$ 除算で主張の余次元に達する。$H\ne0$ なので除算は
特殊化先の分数体で有効である。経路の各文字の次元には一般点からの下界があり、
長さ4と5の経路次元が一定の部分では中心方程式の核にも同様の下界がある。
以上の消去の上界と一致するので結論を得る。
\end{proof}

\begin{corollary}
$D\ne0$ が3の上の正付値を持つとき、$x^3=Dy$ 上の標数零の全通常窓で
各始点の第六経路の次元は84、四次中心の文字別次元は $(3,2,0,0)$ である。
\end{corollary}
\begin{proof}
係数体の付値を点の座標体まで延長する。$\nu(y)>0$ なら上の定理が適用される。
$\nu(y)=0$ なら $\nu(x)>0$ なので定理\ref{thm:rank-fifteen-residue-line}の経路・中心の計算を使う。
$\nu(y)<0$ なら $3\nu(x)=\nu(D)+\nu(y)$ より $\nu(y)<2\nu(x)$ であり、
定理\ref{thm:rank-fifteen-pole-cone}の経路と中心の計算を使う。
\end{proof}
正付値領域全体での全次数の正則性を決めるには、四次中心式を課した商の高次経路が残る。

% !TeX root = ../papers/nc-p3.tex
\section{可変跡の円板での中心四次式と一次経路の積}
\label{sec:rank-fifteen-first-center-products}
\begin{proposition}
定理\ref{thm:rank-fifteen-variable-zero-disk}の領域
$0<\nu(y)<3\nu(x)$ で、同定理で選んだ四中心式を $c_1,\ldots,c_4$ とする。
各始点において、各 $c_j$ に終端から出る四つの一次経路を掛けて得る
十六本は一次独立である。従ってこの四中心商の長さ5の文字別次元は
$(10,10,10,10)$、総次元は40である。
\end{proposition}
\begin{proof}
前節と同じ局所環 $R$ と変数 $H$ を用いる。精度9から長さ5までの未切断経路を
同じ順序で消去すると、長さ4と5で二桁ずつを使い、精度5が残る。
中心方程式の主元は単元なので、前節と同じ自由座標で四中心式を取れる。
各文字 $h$ に対し、文字 $g$ の中心式と末尾の文字 $g\oplus h$ の積を一行とする。
各単語を未切断経路の長さ5の基底で表すと、始点と $h$ の各組に対し
四行十四列の行列が得られる。

\path{computations/verify_rank_fifteen_first_center_products.py}
はこの十六行列を $R$ 上の同じ単元判定で消去する。
各行列とも、階数4に達するまでに必要な $H$ 除算は四回であり、精度が一桁残る。
各除算で捨てる定数項は $\mathbb Q(z)$ で恒等的に零であることを検査する。
従って $H$ が正付値で非零の全ての窓で、十六本の積が文字ごとに四本ずつ独立となる。
中心性により一次経路を先頭に置いた積も同じ空間に属し、
四中心式が生成する両側イデアルの長さ5成分をこれで尽くす。
\end{proof}

% !TeX root = ../papers/nc-p3.tex
\section{可変跡の円板での六次中心商の二文字の確定}
\label{sec:rank-fifteen-saturated-center-products}
\begin{proposition}
定理\ref{thm:rank-fifteen-variable-zero-disk}の領域で前節の四中心式による商を考える。
その長さ6の文字別次元を $d_h$ とすると、全始点で
\[
            d_0=d_1=12,\qquad 10\le d_2,d_3\le12 .
\]
従って六次中心商の総次元は44以上48以下である。
\end{proposition}
\begin{proof}
$I$ を四中心式の生成するイデアルとする。
前節の各四行十四列の行列を精度9から計算し、単元による行操作と
$H$ による飽和で得た階数4の行空間を、主列が単位行列になる基底で表す。
残り一桁の精度でその基底の全定数項が得られる。飽和の前後で $H\ne0$ に
おける行空間は同じなので、この基底は長さ5の $I$ の元からなる。
これを四つの末尾文字に沿って直和し、長さ6の未切断経路空間へ射影する。
この射影は前節までと同じ第六挿入行列で計算する。
五桁のうち二桁の損失で求まるので、積の基底・射影とも定数項が使える。
四式は中心的で経路は一次生成なので $I_6=I_5 A_1$ であり、
各文字について得る十六行は $I$ の六次成分を張る。

\path{computations/verify_rank_fifteen_saturated_center_products.py}
は、これらの定数項を $z=1+i$、標数3に還元して消去する。
四始点すべてで、十六行の階数は文字順に $(8,8,10,10)$ である。
対応する小行列式は単元なので、これは円板全体での $I$ の次数6成分の次元の下界を与える。
一方、中心性により各始点の六次成分は、中心式と長さ2の経路の積で生成される。
長さ2の文字別次元が $(2,2,3,3)$、中心式は文字0と文字1が各二本なので、
生成元数は文字ごとに $(8,8,12,12)$ である。これらを未切断経路の次元
$(20,20,22,22)$ から引けば主張を得る。

さらに文字 $h=2,3$ では、十個の単元主元を用いて十六行二十二列の行列を
十次の単位行列と六行十二列の残りの行列 $B_h$ に分けられる。
したがって $d_h=12-\operatorname{rank} B_h$、かつ $\operatorname{rank} B_h\le2$ である。
この文字で目標次元10に達する必要十分条件は $B_h$ の二次小行列式の少なくとも一つが
その点で非零となることである。
\end{proof}

% !TeX root = ../papers/nc-p3.tex
\section{ガウス有理数の跡から得る全点正則な種数二の鉛筆切片}
\label{sec:rank-fifteen-gaussian-pencil}
\begin{theorem}\label{thm:rank-fifteen-gaussian-pencil}
標数零の代数閉体上の階数15曲面で $D_0=59319/10000$ とおく。
$x^3=D_0y$ の全点は標準Koszul左右AS正則で、各始点で四中心四次式の
多項式環上自由階数256である。添字保持同型類の曲線は種数2の滑らかな完備曲線
\[
                 q^2=(1-s^2)^2-4s^6/D_0
\]
から $s=0,\infty$ 上の計四点を除いたもので、全点の最小周期は2である。
\end{theorem}
\begin{proof}
$i^2=-1$、$z=1+i$ とし、定理\ref{thm:rank-fifteen-variable-zero-disk}の表示で
$K=3i/4$, $n_0=5i/4$, $P_0=(3+i)/2$, $Q_0=(1+3i)/2$ を得る。
$U=K^2=-9/16$ であり、$(U+3)^3/(U-1)^2=D_0$ である。
$H^2=x(U-1)^2/(U+3)$ と取ると $H^6=y(U-1)^4$ を満たす。
$\mathbb Z_{(3)}[i]$ では $2,z,P_0,Q_0,n_0$ は単元で、
前出の平方根はこれらの定数項から一意に持ち上がる。

3の上の付値を取る。$\nu(0)=+\infty$ とする。
$\nu(y)>0$ では $\nu(H)=\nu(y)/6>0$ である。
\path{computations/verify_rank_fifteen_gaussian_pencil.py} は
$\mathbb Q(i)[H]/(H^{25})$ で四中心商を計算する。
単元判定には定数項のノルムを3で還元し、$H$ 除算には実部・虚部の両定数項の
有理数としての消滅を検査する。経路と四次中心で四桁を使った後、
長さ1から13までの各段の最大損失は四始点とも
$(0,0,0,2,4,2,2,2,2,2,2,2,0)$ である。従って四中心商は
$(1+t+t^2+t^3)^4$ の係数を上界とし、長さ13で消える。
定理\ref{thm:rank-fifteen-cubic-zero-disk}と同じ上下界および左右Koszul複体の
議論で全次数の結論を得る。
$\nu(y)=0$ なら $\nu(D_0)>0$ より $\nu(x)>0$ で、剰余直線の
定理\ref{thm:rank-fifteen-residue-line}を適用する。$\nu(y)<0$ なら
$\nu(y)<2\nu(x)$ なので定理\ref{thm:rank-fifteen-pole-cone}を使う。
代数的な $y$ には数体の素点、超越的な $y$ にはGauss付値を選べる。

一般に $D\ne0,27$ を固定し $a^2=1-27/D$ を一つ選ぶと、
商の方程式から $s=p/(ax)$, $x=s^2$, $p=as^3$, $y=s^6/D$ となる。
この逆写像により $s\ne0$ の二重被覆が添字保持パラメータを与える。
$(1-x)^2-4x^3/D$ の判別式は $16(D-27)/D^2$、定数項は1なので、六次式には単根が六個ある。
二重被覆は無限遠の二点で不分岐であり、種数は2である。
通常表示の外の四点の $D$ は $27$ または $38\pm6\sqrt2$ なので $D_0$ 上では全て通常である。
縮約不変量 $(x,p,q)$ が添字保持同型を分離する。
一位置移動は $(s,q)\mapsto(-s,-q)$、二位置移動は恒等であり $s\ne0$ なので周期も従う。
\end{proof}

\begin{corollary}
ある有限集合 $E\subset\overline{\mathbb Q}$ が存在し、$D\ne0,27$ かつ $D\notin E$ なら、
$x^3=Dy$ の全点についても、定理の全次数の正則性・自由性および最小周期2の
添字保持分類が成り立つ。曲線は上式の $D_0$ を $D$ に替えたものである。
\end{corollary}
\begin{proof}
定理\ref{thm:rank-fifteen-vertical-curves}の両表示に残る有限集合を
$D=x^3/y$ で写し、その像を $E$ とする。通常表示の外の四点も
定理\ref{thm:rank-fifteen-formal-boundary}により正則である。
これで各切片の全点を覆い、上で求めた逆写像と移動の式を適用できる。
\end{proof}

\begin{corollary}
上記の切片で一位置移動も同一視した商は、
楕円曲線 $Y^2=X(X-1)^2-4/D_0$ から無限遠点および $X=0$ 上の二点を除いた曲線である。
\end{corollary}
\begin{proof}
$X=s^{-2}$、$Y=q/s^3$ が移動の不変体を生成する。
一般の $D\ne0,27$ で、この三次曲線と $(s,q)\mapsto(-s,q)$ による
もう一方の商 $q^2=(1-x)^2-4x^3/D$ はともに滑らかである。
三次式の係数から計算した $j$ 不変量は順に
$16D^2/(D-27)$、$D(D-24)^3/(D-27)$ となる。
\end{proof}

\begin{corollary}
剰余標数3で $\alpha=\nu(3)$, $\beta=\nu(y)>0$, $\delta=\nu(x^3/y-D_0)$
とする。$\delta>6\alpha$ かつ $6(\delta-5\alpha)\ge25\beta$ ならば、
通常な階数15窓は同じ全次数の正則性と自由性を持つ。
\end{corollary}
\begin{proof}
$D=x^3/y$、$S_0=i/4$ とする。$K=3S$ と置くと
\[
 \frac{D-27}{3^6}=G(S),\qquad G(S)=\left(\frac{S(S^2-1)}{9S^2-1}\right)^2 .
\]
$\bar S_0=i$ では分母は単元で $\bar G'(i)\ne0$ である。
従って付値を延長して $S_0$ と同じ剰余の解を取れば
$\nu(S-S_0)=\delta-6\alpha$ となる。$\delta=+\infty$ は定理に含まれる。
$K-3i/4$ の付値は $\delta-5\alpha$ で、$K^2-1$ が単元なので $\nu(H)=\beta/6$ である。
仮定より $K=3i/4+H^{25}T$ と整な $T$ で書ける。
四次根と各平方根の定数根は単根なので $H^{25}$ を法として定理の計算と一致し、
用いた全ての主元と各除算で必要な合同が保たれる。
\end{proof}

% !TeX root = ../papers/nc-p3.tex
\section{可変跡円板の六次中心商の完全な文字別次元}
\label{sec:rank-fifteen-sixth-center-minors}
\begin{theorem}\label{thm:rank-fifteen-sixth-center-minors}
標数零、剰余標数3の通常な階数15窓で $0<\nu(y)<3\nu(x)$ とする。
定理\ref{thm:rank-fifteen-variable-zero-disk}の四中心式による商の次数6の次元は
全始点で44、文字別には $(12,12,10,10)$ である。
\end{theorem}
\begin{proof}
第\ref{sec:rank-fifteen-saturated-center-products}節の六行十二列の残余行列を
文字 $h=2,3$ について調べる。
$R=\mathbb Z[z]_{(3,z^2+z+2)}$ 上で $H^{11}$ を法として出発する。
長さ5の未切断経路と四次中心を構成すると七桁が残り、
各始点の五次イデアル基底にはさらに四桁、第六挿入の射影には二桁を使う。
両者を合成した十六行二十二列の行列は三桁まで正しい。

\path{computations/verify_rank_fifteen_sixth_center_minors.py}
は級数の係数に一つの共通分母を用いる。分母の単元性は
$z^2+z+2$ による剰余を取って二係数を3で還元して検査する。
十個の単元主元を消した後の全成分について、$H^0,H^1$ 係数の分子が
$\mathbb Z[z]$ で厳密に零になる。
従って残余行列は $B_h=H^2 C_h$ と書ける。$C_h$ の定数項を
$z=1+i$、標数3に還元し、行と列を二つずつ選んで行列式を計算すると、
全四始点で同じ二列 $(4,11)$ を使える。元の十六行の番号を零から数えると、
文字2で二行 $(0,5)$、文字3で $(1,4)$ を取った小行列式が単元となる。
対応する $B_h$ の二次小行列式は $H^4$ と単元の積である。
$H\ne0$ で残余階数は少なくとも2となり、二次経路と四中心式の積の個数から
上界も2である。文字0,1と合わせて主張が従う。
\end{proof}
\begin{corollary}
$D\ne0$ が3の上の正付値を持つとする。$x^3=Dy$ 上の全通常窓で、
四つの中心四次式を選んで各始点の商の長さ6の次元を44にできる。
\end{corollary}
\begin{proof}
$\nu(y)>0$ は本定理による。$\nu(y)=0$ では剰余直線、
$\nu(y)<0$ では極の定理の自由性の証明で四式を選べる。
その商のHilbert級数は $(1+t+t^2+t^3)^4$ で、六次係数は44である。
\end{proof}

% !TeX root = ../papers/nc-p3.tex
\section{六次中心商の証明を整な傾きへ降ろす}
\label{sec:rank-fifteen-integral-slope}
\begin{theorem}\label{thm:rank-fifteen-integral-slope}
標数零の通常な階数15窓に付値を取り、その剰余標数を零または奇素数とする。
$\nu(y)>0$ かつ $D=x^3/y$ が整であるとする。
このとき各始点の未切断六次経路は84次元で、
四つの中心四次式を選んだ商の六次成分は44次元である。
この主張は剰余傾きと剰余体の拡大次数に制限を課さない。
\end{theorem}
\begin{proof}
局所計算の $K=(z^2+z^{-2})/2$, $U=K^2$ に対して傾きは
\[
 \Phi(z)=\frac{(z^8+14z^4+1)^3}{4z^4(z^4-1)^4}=\frac{P(z)}{Q(z)} .
\]
前節の十一桁の消去で単元として使用した定数項の分母と分子、および
選んだ二次小行列式の先頭係数を記録する。文字0,1の未切断六次射影と
飽和五次イデアルの積の階数8の主元も記録する。
\path{computations/verify_rank_fifteen_integral_slope.py} は整数係数の各多項式の
内容が2の冪であることを検査し、原始因子を重複なく合わせた多項式 $F$ を作る。
その次数は 205 である。記録した主元での消去と $H$ 除算前の厳密な零は
$\mathbb Z[1/2,z,1/F][[H]]$ で同じ証明を与える。従って $F(\bar z)\ne0$
となる表示が一つあれば十分である。

$i^2=-1$ とし、$z\mapsto iz$ と $z\mapsto(z+1)/(z-1)$ のいずれでも
$\Phi$ が変わらないことは式を代入すれば分かる。有限の傾きでは
$z(z^4-1)=0$ の根は現れないので、まずその因子を除く。
$f$ と $f(-z)$ の最大公約元を取れば、零点が $z,-z$ の双方に残る。
モニック化したこの多項式は偶多項式なので、$iz$ の代入は偶数次係数の符号だけで扱える。
さらにその代入および $(z+1)/(z-1)$ の斉次代入との最大公約元を取る。
次数が下がる限りこれを反復する。全ての変換先が悪い点は各段で残るため、
最終の多項式が単元なら必ず良い表示がある。

有理数上では最終の最大公約元は1である。整数多項式 $a,b$ の
最大公約元 $d$ ごとに $a,b$ の先頭係数、内容と
$\operatorname{Res}(a/d,b/d)$ の素因子を検査する。
これらが単元となる素数では剰余でも同じ消去が有効である。
残る奇素数はコードに列挙し、各々の素体の多項式環で最初から最大公約元を取り直す。
残った根は全て、以下の六つの $\bar D$ に対する $P-\bar D Q$ の根で尽きることを
反復する最大公約元で検証する。この検査は剰余体の代数閉包でも有効である。

最初の主元で残った $(p,\bar D)$ と、その傾きの追加表示を与える
$\bar z$ の最小多項式 $m$ は順に
\[
\begin{array}{c|cccccc}
p&5&7&13&19&29&37\\
\bar D&0&0&0&17&3&0\\
m&z^2+z+2&z^2+z+6&z-2&z-2&z-2&z-2
\end{array}
\]
である。各 $m$ の既約性、$m\mid P-\bar D Q$ と $m\nmid Q$ を検査する。
19以外の五つの剰余体で同時に可逆な主元を取り、19ではその表示で可逆な主元を取って、有理関数体上の十一桁の計算をやり直す。
中心方程式の自由列をこの主元に従って選び直し、文字零の後二自由列と文字一の両自由列を使う。
全十六個の六次射影で目標次元に達し、五次中心イデアルの積の最終消去で
必要な除算回数はどの始点でも文字順に $(0,0,2,2)$ である。

所定の有限な $D$ に対し $P-DQ$ はモニックで定数項1である。
付値を代数閉包に延長すれば全根は整な単元であり、積の還元から
任意の指定した剰余根が持ち上がる。奇標数で $z^4=1$ と $P=0$ は両立しない。
よって $U-1$ は単元で、$H^6=y(U-1)^4$ は正付値を与える。
$x\ne0$ なら $H^2=x(U-1)^2/(U+3)$ と取る。
$x=0$ は定理\ref{thm:rank-fifteen-x-zero}でも認証済みである。
前節の経路・中心の下界と選んだ四式の上界を合わせ、同じ正則性座標の窓へ移せば結論を得る。
\end{proof}

% !TeX root = ../papers/nc-p3.tex
\section{左右の積の照合による混標数環帯の全次数認証}
\label{sec:rank-fifteen-overlap-annulus}
\begin{theorem}\label{thm:rank-fifteen-overlap-annulus}
標数零の通常な階数15窓に剰余標数3の付値を取り、
$0<\nu(y)<3\nu(x)$ かつ $5\nu(y)\le6\nu(3)$ とする。
この窓は標準Koszul左右AS正則で、各始点で四中心四次式の多項式環上
自由階数256である。
\end{theorem}
\begin{proof}
第\ref{sec:rank-fifteen-variable-zero-disk}節の表示で $\bar z=1+i$、
$\nu(H)=\nu(y)/6$ と取る。$R=\mathbb Z[z]_{(3,z^2+z+2)}[[H]]$ とし、普遍表示の分数体上の商の中で、
各長さと文字について語が生成する $R$ 格子を取る。
定理\ref{thm:rank-fifteen-sixth-center-minors}までの計算では、未切断経路の長さ5まで、
および四中心商の長さ6まで、必要な $H$ の冪で割った後に全主元が単元になる。
よって各格子を目標次元個の語で生成でき、分数体上の下界と合わせて自由基底となる。
$H$ 除算後の単元判定は有限精度で済むため、同じ消去を形式級数で持ち上げられる。
各格子は任意の小さい $H\ne0$ への特殊化後も、実際の語空間に全射する。

まず $R$ を3で還元し、十五桁を使って計算する。
四次中心の自由座標は標数零と3の双方で最後の始点の同じ語基底の同じ列で固定する。
未切断経路の剰余消去から得た生成列は、対応する還元格子を生成する。
実際、関係を課した後も、還元格子が自由で $H$ 捩れを持たないため、
計算に使う $H$ 除算後の行もその格子で関係を与える。
列数と階数が等しければこれらの列は基底であり、左右の積はその基底での還元に一致する。
中心方程式は単元主元で解け、前記の自由座標が一致するので、選んだ四式も一致する。
同じ議論は四式を課した格子に対して長さ6まで成り立つ。

四式を課した商を $B$ とする。まず $n\le5$ では、未切断経路の基底で
中心式と長さ $n-4$ の末尾との積を消去し、両側からの生成元の積の表を作る。
長さ5で四桁を使い、七桁残る。次に始点 $p$、文字 $h$ ごとに
末尾別の $B_{p,n-1}$ の四つの生成系と、先頭別の $B_{p+1,n-1}$ の
四つの生成系を直和に並べる。長さ $n-2$ の中間の生成元 $w$ と
先頭末尾の文字 $a,b$ に対して、結合則
\[
                       (aw)b-a(wb)=0
\]
の行を入れる。両項の座標は前段の左右の積表から得られる。
この直和は実際の商へ全射し、各行はその核に入る。消去後の自由列を
生成系に取り、元の直和の各列の像を次の両積表に用いる。
各 $H$ 除算は有効桁を一つ消費する。従って途中で余分な自由列や未使用の関係があっても
得られる次元は上界となる。

\path{computations/verify_rank_fifteen_overlap_annulus.py} は各 $H$ 係数を
$\mathbb F_3(z)$ の有理関数として計算する。除算前の定数係数はこの体で
厳密に零であり、単元主元の分子・分母は $z^2+z+2$ で還元して非零である。
長さ6から13で各段の最大精度損失は全始点で
\[
                             2,2,0,2,0,0,0,0
\]
となり、長さ6の後に五桁、最後に一桁が残る。ここで仮定は $3\in H^5\mathcal O$ と同値である。
長さ6までの五桁の表は上の自由格子の積表の還元なので、
$R\to\mathcal O/(H^5)$ により所定の窓の積表に移せる。
以後の除算はこの五桁だけで完了する。座標が $H$ を法として零なら、
実際の関係を分数体で $H$ で割って一桁少ない合同を得る。各始点における商の次元上界の列は
$(1+t+t^2+t^3)^4$ の係数列で、次数13で零になる。
四中心式の抽象的多項式環上の自由加群からの全射と一般点からの下界が
全次数で一致するので自由性が従う。定理\ref{thm:rank-fifteen-cubic-zero-disk}と
同じ左右Koszul複体の議論で正則性を得る。
\end{proof}

\begin{corollary}
数体 $L$ 内の $x,y$ が通常な窓を与え、$D=x^3/y$ が
3の上の全ての素点で正付値を持つとする。
\[
  5\operatorname{ord}_3\!\left(N_{L/\mathbb Q}(y)\right)\le6[L:\mathbb Q]
\]
ならば、その窓にも定理の全次数の結論が成り立つ。
\end{corollary}
\begin{proof}
各素点 $w\mid3$ で $\nu_w(3)=1$ と正規化する。
$\operatorname{ord}_3(N(y))/[L:\mathbb Q]$ は $\nu_w(y)$ の
重み $[L_w:\mathbb Q_3]/[L:\mathbb Q]$ の平均である。
よって $\nu_w(y)\le6/5$ となる $w$ がある。
この値が正なら $\nu_w(D)>0$ より本定理が適用できる。
零なら $\nu_w(x)>0$ なので定理\ref{thm:rank-fifteen-residue-line}、
負なら $\nu_w(y)<2\nu_w(x)$ なので定理\ref{thm:rank-fifteen-pole-cone}を使う。
\end{proof}

% !TeX root = ../papers/nc-p3.tex
\section{正則性座標の水平切片と添字保持同型曲線}
\label{sec:rank-fifteen-horizontal-fibers}
\begin{theorem}\label{thm:rank-fifteen-horizontal-fibers}
標数零の代数閉体で $y=c\ne0$ を固定する。階数15のPauli商曲面の添字保持同型類の
水平切片は、$t\ne0$ 上の曲線
\[
       V^2=2t\bigl((-t^2+2t-4c)^3-216ct^3\bigr)
\]
である。$c\ne1,1/64$ では滑らかな完備超楕円曲線の二つの
Weierstrass点を除いたもので、種数は3、全点の最小周期は2である。
一位置移動の商は種数2の曲線 $W^2=(x^3-27c)((1-x)^2-4c)$ から
無限遠点を除いたものである。
ある有限集合 $E\subset\overline{\mathbb Q}^{\times}$ の外では
水平切片の全点が標準Koszul左右AS正則で、中心四次式上の自由階数は各始点で256である。
\end{theorem}
\begin{proof}
定理\ref{thm:rank-fifteen-rational-quotient}、第\ref{sec:rank-fifteen-global-invariants}節および
定理\ref{thm:rank-fifteen-formal-boundary}より、$p^2=x^3-27c$, $q^2=(1-x)^2-4c$
の各点が一つの添字保持同型類を与え、一位置移動は $(p,q)\mapsto(-p,-q)$ である。
$t=1-x+q$ と取れば $(1-x+q)(1-x-q)=4c$ より
\[
 x=1-(t+4c/t)/2,\qquad q=(t-4c/t)/2,\qquad V=4t^2p .
\]
従って表示はアフィン曲線の同型である。一位置移動は
$(t,V)\mapsto(4c/t,-16c^2V/t^4)$ となる。
$0,\infty$ はいずれも二重被覆の分岐点であり、この二点だけを除く。
移動不変環は $k[x,pq]$ である。実際、$k[x]$ 上の自由基底
$1,p,q,pq$ のうち $1,pq$ が固定されるので $W=pq$ と取れる。

二つの被開平式の終結式は $(c-1)^2(1-64c)$ である。
\path{computations/verify_rank_fifteen_horizontal_fibers.py} は $t$ 表示の六次因子の
判別式が非零定数と $c^{10}(c-1)^2(64c-1)$ の積であることを検査する。
六次因子の定数項も非零なので、通常の切片の二重被覆は八点で分岐し種数3、
移動の商は五次式の五根と無限遠で分岐し種数2となる。
同時に $p=q=0$ となる点が無いので、最小周期は2である。
定理\ref{thm:rank-fifteen-vertical-curves}の有限例外集合の $y$ 座標を $E$ とすれば
全次数の主張が従い、通常表示外の四点は定理\ref{thm:rank-fifteen-formal-boundary}で覆われる。
\end{proof}
\begin{corollary}
$c=1,1/64$ の水平切片にはそれぞれ一つの通常二重点があり、
切片の正規化の種数は2、移動の商の正規化の種数は1である。
各二重点が表す窓は最小周期1の標準Koszul左右AS正則窓であり、
その切片の他の全点は最小周期2である。
\end{corollary}
\begin{proof}
二重点の $(x,t)$ は順に $(3,-2),(3/4,1/4)$ である。
七次式にある唯一の二重根を $b$ とすると、$(t-b)^2$ を取り除いた五次式は平方自由で、
$b$ で非零である。よって二つの異なる接線があり、正規化の二点が一つの同型類に重なる。
商の五次式も唯一の二重因子を除くと平方自由三次式となる。
二重点は $p=q=0$ の二つの既知の正則窓である。
\end{proof}
特に $c=3$ は滑らかな種数3切片である。3の上で $\nu(3)=1$ とすれば、
この曲線の $\nu(x)>1/3$ を満たす全点は定理\ref{thm:rank-fifteen-overlap-annulus}により
全次数で認証される。この明示的な範囲には、例えば $x\in3\mathbb Z_3$ の全ての持上げが入る。

% !TeX root = ../papers/nc-p3.tex
\section{三つの次数剰余類からの全次数の正則性}
\label{sec:rank-fifteen-three-residue-classes}
\begin{lemma}\label{lem:rank-fifteen-three-residue-classes}
標数零の四頂点折り畳み環 $A$ が、通常な階数15窓と同じ完全対合を持つ
標準二次表示と左右の $14641$ 複体を持つとする。
各始点の Hilbert 関数に $h_n=\binom{n+3}{3}$ の下界があり、
四中心四次式が定める $S=k[u_1,u_2,u_3,u_4]$ 上で $A$ は有限生成とする。
$M_r=\bigoplus_{n\equiv r\ (4)} A_n$ と置く。
$r=0,2,3$ で $M_r$ が $S$ 上自由で各始点の Hilbert 関数が $h_n$ なら、
残りの $M_1$ も自由で、全始点で全次数の Hilbert 関数が $h_n$ に等しい。
さらに $A$ は標準Koszul左右AS正則である。
\end{lemma}
\begin{proof}
全頂点を直和した右Koszul列の次数 $n\equiv2\pmod4$ を考える。
表示から $H_0=H_1=0$ である。上の二微分は完全対合で左の表示の双対になる。
左増大イデアル $J=A_{\ge1}$ から $A$ への次数 $d=n-4$ の左線形写像 $f$ を取る。
$u=u_1,v=u_2$ として $u f(v)=v f(u)$、$f(u),f(v)\in M_2$ を得る。
$M_2$ の自由性より $f(u)=ub$ と書ける。各矢 $a$ では $f(a)-ab\in M_3$ に対し
$u(f(a)-ab)=0$ なので、$M_3$ の自由性より $f(a)=ab$ である。
矢は左イデアルとして $J$ を生成するので $f$ は延長する。従って
$H_3=0$ であり、$H_4=0$ も $M_2$ 上で $u$ が単射であることから従う。
この列のEuler数は、各項のほかの次数が $2,0,3,2$ に属するので
\[
        \dim H_2=-4\bigl(\dim A_{n-1}-4h_{n-1}\bigr)\le0 .
\]
従って $H_2=0$、全頂点で下界が等号であり、$n-1\equiv1$ の全次数を得る。

次数 $2\bmod4$ の完全列を $S$ 加群として集める。次数シフトを省略すれば
\[
 0\longrightarrow F_4\longrightarrow F_3\longrightarrow F_2
 \longrightarrow M_1^{\oplus4}\longrightarrow F_0\longrightarrow0
\]
で、全ての $F_i$ は自由である。最後の全射は分裂するので $M_1$ の
射影次元は高々2である。四変数多項式環の斉次極大イデアルで深さは少なくとも2となる。
よって $A$ に同時に正則な二つの正次数の中心元 $u,v\in S$ を選べる。
上の延長論法で今度は $v$ が $A/uA$ に単射、$u$ が $A$ に単射であることを使うと、
全次数で $H_3=H_4=0$ を得る。左右を替えても同様である。
確定したHilbert関数のEuler等式から全次数で $H_2=0$ となり、
テンソルの双対性により左右Koszul・AS性が従う。

最後に次数 $1\bmod4$ の完全列を集めると、次数シフトを除き
$0\to M_1\to F'_3\to F'_2\to F'_1\to M_1\to0$ で中の三項は自由である。
射影次元が高々2の加群の第三シジジーは射影的なので $M_1$ は射影的であり、
次数付き自由基底を持つ。始点ごとの直和因子も自由である。
\end{proof}
\begin{corollary}
補題では不足する剰余類を任意の $r\in\mathbb Z/4$ に取っても同じ結論が成り立つ。
\end{corollary}
\begin{proof}
$n\equiv r+1$ では上の延長論法に $M_{r+1},M_{r+2}$ を使え、
Euler数で唯一未知の項は $-4\dim A_{n-1}$ である。これで同じHilbert関数を得る。
$r\ne3$ ならこの剰余類に始点単純は現れず、補題と同様に $M_r$ の射影次元は高々2である。
$r=3$ の場合は次数 $1\bmod4$ の列を使う。
$M_1,M_2$ の自由性より $H_3=H_4=0$、Euler等式から $H_2=0$ である。
未知項は中央の $M_3^{\oplus6}$ だけで、末端の自由加群への全射が分裂し、
その核への全射も分裂するので $M_3$ の射影次元は高々1である。
従って深さによる全複体の完全性が同様に従う。
$r\ne0$ では $M_r$ は自身の第三シジジーである。
$r=0$ では $K=\ker(M_0\to A_0)$ の第三シジジーとなる。
$A_0$ は $S$ 上の射影次元4であり $M_0$ の射影次元は高々2なので
$K$ の射影次元は高々3である。いずれも $M_r$ は射影的である。
\end{proof}
\begin{theorem}\label{thm:rank-fifteen-full-trace-disk}
標数零、剰余標数3の通常な階数15窓で $0<\nu(y)<3\nu(x)$ とする。
全始点で四つの中心四次式上自由階数256であり、標準Koszul左右AS正則である。
\end{theorem}
\begin{proof}
第\ref{sec:rank-fifteen-variable-zero-disk}節の $R$ と実際の四中心式で計算する。
\path{computations/verify_rank_fifteen_overlap_bound.py} は $\mathbb Q(z)$ 上の
十五桁の $H$ 級数から未切断経路と四中心式に四桁を使う。
各始点で四中心商の長さ1から7までの最大損失は $(0,0,0,2,4,2,2)$ で、
七次で文字別 $(10,10,10,10)$ の生成列と一桁の右積表が残る。
$H$ を落とす際には全係数の有理関数としての零を検査する。

右積表を $(3,z^2+z+2,H)$ で還元し、自由列で選んだ語の先頭に文字を付け、
その語を左から右へ評価して左積表も作る。第\ref{sec:rank-fifteen-overlap-annulus}節の結合則の行を使い、
長さ8から13で得る剰余体上の余次元は全始点で
\[
\begin{array}{c|cccccc}
 n&8&9&10&11&12&13\\ \hline
 \dim&31&24&10&4&1&0 .
\end{array}
\]
各段で剰余独立な行だけを持ち上げれば、その主元小行列式は単元である。
実際の商へ全射する自由な生成系を二通りの積の直和から取り、
持ち上げた独立行だけで割ると、得る両積表は計算した剰余表を持つ。
よって余次元はすべて所定の窓での上界である。
次数13の消滅と一次生成性から $A$ は $S$ 上有限である。
次数 $0,2,3\bmod4$ ではこれらの上界が $(1+t+t^2+t^3)^4$ の全係数に一致する。
各剰余類で列を持ち上げた自由 $S$ 加群の上界はちょうど $(1-t)^{-4}$ の
その剰余類部分で、各始点の一般点からの下界と等しい。従って三つの自由性が得られる。
補題\ref{lem:rank-fifteen-three-residue-classes}から全次数の結論が従う。
この議論により、実際の中心商の次数9の次元も各始点で20に決まる。
\end{proof}
\begin{corollary}\label{cor:rank-fifteen-open-integral-slopes}
ある $C(T)\in\mathbb Z[T]$ で $3\nmid C(0)$ となるものが存在する。
剰余標数が零または奇素数で、$\nu(y)>0$, $D=x^3/y$ が整、
$C(\bar D)\ne0$ なら上の全次数の結論が成立する。特に剰余標数3で
残り得る整な傾きの剰余は有限個である。
\end{corollary}
\begin{proof}
有理関数体上の全 $H$ 除算で捨てた係数は恒等的に零である。
七次までの有限個の逆元と、その後の剰余独立行を持ち上げたときの
主元の定数項は、分母と分子を合わせたある $F(z)\in\mathbb Z[z]$ を
逆にすれば全て単元である。ここに $2$ と平方根の初期分母も入れる。
$F$ は $(3,z^2+z+2)$ の外に取れる。従って $F(\bar z)\ne0$ の表示では
同じ行と列が三つの剰余類の自由性を与える。
第\ref{sec:rank-fifteen-integral-slope}節の $P,Q$ に対し、
$F(z)^{24}$ をモニックな $P(z)-TQ(z)$ で割った余りは
$\mathbb Z[T]$ 係数である。$T=0$ として3で還元すると、
$P(1+i)=0$, $F(1+i)\ne0$ なので余りは零でない。
定数項が3の単元となる余りの係数を一つ取り $C(T)$ とする。
$C(\bar D)\ne0$ なら次数24の $P-\bar DQ$ は $F^{24}$ を割らず、
従って少なくとも一つの幾何根で $F$ が消えない。その根の持上げと
$H$ の選択は同節による。
\end{proof}
\begin{corollary}\label{cor:rank-fifteen-positive-slope}
$D\ne0$ が3の上のある付値で正付値を持つとする。$x^3=Dy$ の全窓は
標準Koszul左右AS正則で、各始点で中心四次式上自由階数256である。
$D\ne27$ なら添字保持同型類は
$q^2=(1-s^2)^2-4s^6/D$ の種数2曲線から $s=0,\infty$ 上の四点を除いたもので、
全点の最小周期は2である。代数的な $D\ne0$ では、傾きの仮定は原始整数最小多項式の定数項が3の倍数であることに等しい。
\end{corollary}
\begin{proof}
付値を各点の座標体に延長する。$\nu(y)>0$ は本定理で尽くされ、
$\nu(y)=0$ は $\nu(x)>0$ なので定理\ref{thm:rank-fifteen-residue-line}、
$\nu(y)<0$ は $\nu(y)<2\nu(x)$ なので定理\ref{thm:rank-fifteen-pole-cone}を使う。
通常表示外の四点には定理\ref{thm:rank-fifteen-formal-boundary}が適用できる。
$D\ne27$ の分類と周期は第\ref{sec:rank-fifteen-gaussian-pencil}節の逆写像による。
$D=27$ は定理\ref{thm:rank-fifteen-cubic-zero-disk}の節点切片である。
最後の同値は、全ての3進素点で $\nu(D)\le0$ となることが $1/D$ の3整性に等しいことによる。
原始最小多項式の逆数多項式をモニック化した係数が全て3整である条件は、元の定数項が3の単元であることである。
\end{proof}

% !TeX root = ../papers/nc-p3.tex
\section{剰余傾き全体への混標数環帯の降下}
\label{sec:rank-fifteen-integral-annulus}
\begin{theorem}\label{thm:rank-fifteen-integral-annulus}
標数零、剰余標数3の通常な階数15窓で $\nu(y)>0$, $D=x^3/y$ は整とする。
$5\nu(y)\le6\nu(3)$ ならば、各始点で四中心四次式の多項式環上自由階数256であり、
標準Koszul左右AS正則である。剰余傾き $\bar D$ は任意でよい。
\end{theorem}
\begin{proof}
第\ref{sec:rank-fifteen-integral-slope}節の次数205の多項式を $F$ とする。
第\ref{sec:rank-fifteen-overlap-annulus}節の $\mathbb F_3(z)$ 上の十五桁の計算で使った
単元の定数係数の分子・分母に現れる全既約因子の積を $G$ とする。
\path{computations/verify_rank_fifteen_integral_annulus.py} は各逆元と
各主元を記録し、標数3の平方自由分解によって三乗因子に由来する零点も残す。
$\deg G=9835$ となる。$\bar F G$ の零点のうち有限な傾きの点について、
$z\mapsto iz,(z+1)/(z-1)$ による最大公約元の反復を行う。
モニック化した最終多項式は次数144であり、$P/Q$ を同節の傾きの式として
\[
 \prod_g Q^{\deg g}g(P/Q),\qquad
 g(T)=T+1,\quad T^2+2T+2,\quad T^3+2T^2+2T+2
\]
のモニック化と一致する。この三つは既約である。従ってこれらの六つの幾何根を除けば、
$\bar F G$ が零にならない表示があり、所定の $D$ に持ち上げられる。
$\bar F$ は六次までの標数零の自由格子、$G$ は十三次までの剰余計算を保証し、
第\ref{sec:rank-fifteen-overlap-annulus}節の五桁からの再開が適用できる。

三つの $g$ の順に、追加する $\bar z$ の既約多項式を
\[
 z^4+z^3+z^2+z+1,\qquad z^4+2z^3+2,\qquad z^3+2z+2
\]
と取る。これらは全て $\bar F$ と互いに素である。
三つの表示で同時に単元となる主元で剰余の経路基底を取ると、
最後の始点の四次の語基底がもとの計算と変わり得る。
そこで六次の標数零証明で使った語の像を現在の基底で求め、
その基底変更が単元主元だけでできることを検査する。
中心方程式の最後の始点ではこれらの語を列に使い、文字 $0,1$ の
最後の $3,2$ 個の自由座標を固定して単元主元で解く。
この四中心式は標数零のものの還元である。その商の左右の積を長さ13まで
再計算すると、各段の損失も上界も第\ref{sec:rank-fifteen-overlap-annulus}節に一致する。

各追加表示は $Q$ と互いに素である。対応する $g(P/Q)=0$ と
表示の既約性を検査しているので、Frobenius共役を取れば $g$ の全ての幾何根を覆う。
$P-DQ$ の根を持ち上げ $\nu(H)=\nu(y)/6$ として同じ全次数判定を適用する。
\end{proof}

\begin{theorem}\label{thm:rank-fifteen-unit-residues}
剰余標数3の付値で $(\bar x,\bar y)\in(\mathbb F_3^\times)^2$ となる
通常な標数零の全窓も標準Koszul左右AS正則で、各始点の自由階数は256である。
特に正則性座標が有理数で $x$ が3整、$y$ が3の単元なら全て認証される。
\end{theorem}
\begin{proof}
\path{computations/verify_rank_fifteen_unit_residues.py} は四対
$(1,1),(1,2),(2,1),(2,2)$ についてそれぞれ次数 $6,4,4,12$ の
有限体で $K,H,n,P_+,P_-,Q_+,Q_-$ を与える。$U=K^2,V=H^2$ として
両座標式、$n^4+(1-U)n^2+V=0$, $P_\pm^2=2(K+1\pm H)$,
$Q_\pm^2=2(K-1\pm H)$ を検査する。各表示で未切断経路を長さ5まで計算し、
四次中心の文字別次元が $(3,2,0,0)$ となることを確かめる。
四つの中心四次式を課した後の全四始点の次元は長さ13の消滅まで
$(1+t+t^2+t^3)^4$ の係数となる。

指定された $x,y$ について $D=x^3/y$ とし、$(U+3)^3-D(U-1)^2=0$ の
指定した剰余根を持ち上げる。$U+3$, $U-1$, $H$, $n$ は単元である。
$V=x(U-1)^2/(U+3)$ として残りのモニックな根を順に選ぶ。
代数閉包を取れば、根が重複する場合も指定された剰余根の上に根が存在する。
未切断経路の上下界と中心の核次元の下界により、単元主元で中心式を持ち上げられる。
四中心商の上界と同じ自由加群の比較で全次数の結論を得る。
有理座標で $\bar x=0$ の場合は定理\ref{thm:rank-fifteen-residue-line}による。
\end{proof}

\begin{corollary}
$x,y\in\mathbb Q$ で $x$ が3整とする。この範囲で未認証の点があるなら
$x\ne0$, $\nu_3(y)>0$, $\nu_3(y)\ge3\nu_3(x)$ を同時に満たす。
\end{corollary}
\begin{proof}
$x=0$ は定理\ref{thm:rank-fifteen-x-zero}、$\nu_3(y)<0$ は極の定理、
$\nu_3(y)=0$ は上の単元剰余の定理である。
残る正付値の場合も $\nu_3(y)<3\nu_3(x)$ なら定理\ref{thm:rank-fifteen-full-trace-disk}で尽くされる。
\end{proof}

% !TeX root = ../papers/nc-p3.tex
\section{七次までの主元次数の比較と整な傾き}
\label{sec:rank-fifteen-smith-slopes}
\begin{lemma}\label{lem:rank-fifteen-smith-transfer}
局所整域 $B$ の分数体を $K$、剰余体を $k$ とし、行列の係数を
$B[[H]]$ に取る。$K[[H]]$ および $k[[H]]$ 上で、必要な本数までの
Smith標準形の $H$ 指数の列が等しいとする。
剰余体での単元主元を持ち上げて消去すると、各段の $H$ 除算を $B[[H]]$ 上で
厳密に実行でき、同じ本数の列を消去できる。
\end{lemma}
\begin{proof}
指数零の剰余主元を一つずつ使う。各主元は $B[[H]]$ の単元なので、残りの行列の
一般体上の指数列からも先頭の零が一つずつ消える。
目標本数に達する前に指数零を使い切ったなら、残りの一般体上の全成分は $H$ の倍数である。
$B\hookrightarrow K$ により全定数項は $B$ でも零である。全体を $H$ で割り、
以後の指数を一つずつ下げて続ける。
途中までの列に対する主張にも、最後に使用する指数を越える桁は要らない。
\end{proof}
$B=\mathbb Z[z]_{(3,m)}$、$k=\mathbb F_3[z]/m$ としてこの補題を使う。
比較用の一般体は常に $\mathbb Q(z)$ である。
基準の $m=z^2+z+2$ では第\ref{sec:rank-fifteen-three-residue-classes}節の
有理関数上の消去が、単元主元が尽きるたびに全定数項の恒等的な零を確認している。
従って基準剰余体で再計算した各主元の除算回数は一般体の指数列に等しい。
単に各行列の最大損失を比較するだけでは足りず、各指数の重複度を含めて照合する。

他の $m$ で未切断経路から順に比較する。前段で残る語は一般体上では
同じ飽和した語の格子の基底を与える。語を替えても $K[[H]]$ 上の基底変更は可逆である。
関係の末尾あるいは先頭に掛ける語の変更も可逆な行変換なので、次段の一般の指数列は不変である。
四中心の選択だけは最後の始点の四次の語を基準表示に揃え、
同じ自由座標を指定する。これで中心商の七次までの比較にも同じ議論が通る。
各段の準同型を剰余に移し、長さ8以後の残余の次元上界が第\ref{sec:rank-fifteen-three-residue-classes}節に
一致すれば、同節の三剰余類の補題で全次数の結論が得られる。

\begin{theorem}\label{thm:rank-fifteen-integral-depth}
標数零、剰余標数3の通常な階数15窓で $\nu(y)>0$ かつ $D=x^3/y$ が整とする。
全ての剰余傾きと全ての深さで、各始点で四つの中心四次式上自由階数256となり、
標準Koszul左右AS正則である。
\end{theorem}
\begin{proof}
\path{computations/verify_rank_fifteen_integral_depth.py} は、まず $\mathbb F_3(z)$ で
十五桁を取り、第\ref{sec:rank-fifteen-three-residue-classes}節の七次までの消去を行う。
主元の定数項は基準剰余体で単元とし、除算前の零は $\mathbb F_3(z)$ で確かめる。
剰余体で選んだ指数の各重複度は、同じ先行格子に対する一般体のものと一致する。
この消去の分母・主元の先頭分子が消えない他の表示では、同じ行列操作が同じ指数を与える。
従って補題\ref{lem:rank-fifteen-smith-transfer}により七次までの消去を混標数の局所環で行える。
七次の右積表の定数項を $\mathbb F_3(z)$ に取り、左右の積の照合では
九次に各文字六列を残す。八次から十三次まで $H$ 除算をせず、
各段の剰余上界は第\ref{sec:rank-fifteen-three-residue-classes}節と同じとなる。

この全計算で使用した単元の定数係数の分子・分母の平方自由な積の次数は12017である。
第\ref{sec:rank-fifteen-integral-slope}節の変換による最大公約元の反復で
残った次数96の多項式は、モニック化すると
\[
               (P-Q)(P+Q)(P^2+2PQ+2Q^2)
\]
のモニック化に等しい。$P/Q$ は同節の傾きである。
このため追加すべき剰余傾きは $1,-1$ および $T^2+2T+2$ の共役な二根に尽きる。

\path{computations/verify_rank_fifteen_smith_slopes.py} は、この順に
\[
 z^6+2z^4+2z^3+z^2+2,\qquad z^4+z^3+z^2+z+1,\qquad z^4+2z^3+2
\]
の根を取る。各多項式の既約性と所定の傾きを確認し、
対応する有限体で基準表示と同じすべての指数列を得る。
比較するのは未切断経路80組、中心方程式4組、最後の始点の基底と自由座標の照合4組、
七次までの中心商112組の合計200組である。
いずれの表示でもその後の八次から十三次の上界も一致する。
剰余表示を指定した $D$ へ持ち上げて $\nu(H)=\nu(y)/6>0$ とし、
同じ三剰余類の自由性判定を適用すれば全点で結論を得る。
\end{proof}
\begin{corollary}\label{cor:rank-fifteen-rational-integral-slopes}
$x,y\in\mathbb Q$ の通常な階数15窓で $D=x^3/y$ が3整なら、
各始点で自由階数256の標準Koszul左右AS正則窓である。
\end{corollary}
\begin{proof}
$x=0$ は既知なので $D\ne0$ とする。$\nu_3(D)>0$ は
系\ref{cor:rank-fifteen-positive-slope}で全点を覆う。
$D$ が単元なら $3\nu_3(x)=\nu_3(y)$ である。
$\nu_3(y)>0$ の場合は $\bar D=1,-1$ の両方が上の計算に含まれる。
負なら $\nu_3(y)<2\nu_3(x)$ により極の定理を使い、
零なら両座標が素体単元なので定理\ref{thm:rank-fifteen-unit-residues}が適用できる。
\end{proof}

\begin{theorem}\label{thm:rank-fifteen-pole-slope-annulus}
剰余標数3の通常な標数零窓で $\bar x=-1$、$0<2\nu(y)\le\nu(3)$ とする。
この窓は標準Koszul左右AS正則で、各始点で四中心四次式上自由階数256である。
この領域では $x^3/y$ は極を持つ。例えば $y^2=3$ の水平切片の $\bar x=-1$ の全点が含まれる。
\end{theorem}
\begin{proof}
$d^2=y/x^3$ とし、$(4+de)^3=e^2$ の単純剰余根 $\pm8$ を持ち上げる。
$U=1+de$, $r^2=U-1$, $T^2=x/(U+3)$ とすれば
\[
    U=1+r^2,\quad V=r^4 T^2,\quad x=(4+r^2)T^2,\quad y=r^4 T^6.
\]
$T$ は単元で $\bar T^2=-1$、$\nu(r)=\nu(y)/4$ である。
$A=(1-s^2)/(1+s^2)$, $B=2s/(1+s^2)$、$T=AB$ と選ぶ。
$T^2+1$ の分子は標数3で
$(s^4+2s^3+s+1)(s^4+s^3+2s+1)$ となる。
第一の既約因子の剰余根を指定し、$T$ の符号を選ぶ。
$T(1+s^2)^2-2s(1-s^2)$ の先頭と定数は単元で、指定根は持ち上がる。
$A^2+B^2=1$ より $K=\sqrt{1+r^2}$ に定数項1を選ぶと、
$n=rA,H/n=rB,H=r^2T$ は必要な平方和と積を持つ。

$P_\pm=\sqrt{2(K+1)\pm2Tr^2}$ の定数項はともに2である。
残りの根を $Q_\pm=r\sqrt{2/(K+1)\pm2T}$ とし、
根号内の平方根の定数項は順に $A+B,B-A$ を使う。
これらの定数項は単元である。
\path{computations/verify_rank_fifteen_pole_slope_annulus.py} は
まず $\mathbb Q(s)$ 上で十桁を使い、未切断経路の長さ1から5まで、
中心方程式、五次までの中心商の左右の積を計算する。
全ての除算前の定数係数は恒等的に零であり、全ての主元は上の四次因子で単元である。
従って第\ref{sec:rank-fifteen-overlap-annulus}節と同じく、
語の格子はこの範囲で自由であり、その階数は一般の下界に一致する。
次に $\mathbb F_3(s)$ 上で十七桁を使う。中心方程式の自由座標を含む
始点ごとに、四次の語基底と指定した座標が標数零のものと一致することも照合する。
これで同じ中心式を課した格子の積表を還元できる。
未切断経路の各長さの最大損失は $(0,1,1,2,2)$、
中心方程式はさらに一桁、五次までの中心商はさらに二桁を使い、八桁が残る。
仮定は $3\in r^8\mathcal O$ と同値なので、還元された五次の積表から
八桁で再開できる。長さ6から13で $(2,1,1,1,1,0,1,0)$ 桁を使い、最後に一桁残る。
商の次元上界は全始点で $(1+t+t^2+t^3)^4$ の係数となり、十三次で消える。
標数零の経路と中心の下界を合わせる同じ自由性・Koszul性の議論で結論を得る。
\end{proof}

% !TeX root = ../papers/nc-p3.tex
\section{極をもつ傾きでの八重被覆と六次の積表}
\label{sec:rank-fifteen-pole-fibers}
\begin{theorem}\label{thm:rank-fifteen-pole-fibers}
剰余標数3の標数零の通常な階数15窓で $\nu(x)=0$, $\bar x\ne1$ とする。
$0<3\nu(y)\le2\nu(3)$ なら、各始点で四つの中心四次式上自由階数256の
標準Koszul左右AS正則窓である。

\end{theorem}
\begin{proof}
定理\ref{thm:rank-fifteen-pole-slope-annulus}の表示で
$\nu(r)=\nu(y)/4$、$x=(4+r^2)(AB)^2$ である。
$s$ から剰余座標 $X=\bar x$ への写像は
\[
  X=N(s)/D(s),\qquad N=16s^2(1-s^2)^2,\qquad D=(1+s^2)^4.
\]
互いに素な二式による次数8の写像で、$s\mapsto-s,1/s,(1+s)/(1-s)$ が生成する
八つの相異なる分数線形変換はこれを保つ。実際、第三の変換の二乗は $-1/s$ であり、
$(A,B)$ の符号変更と交換がこれらの表示を与える。従って一般のファイバーは一軌道で、
$X\ne0,1,\infty$ では $A,B,A+B,A-B$ がいずれも非零で八表示は相異なる。
各表示は符号を整えれば所定の窓へ持ち上がる。実際 $T=AB$ として
定理\ref{thm:rank-fifteen-pole-slope-annulus}の $T$ の平方根の符号を選び、
$T(1+s^2)^2-2s(1-s^2)=0$ の指定した剰余根を持ち上げればよい。

\path{computations/verify_rank_fifteen_pole_fibers.py} は、同じ十七桁の
$\mathbb F_3(s)$ 上の計算で、単元として用いる主元の定数係数の分子と分母を記録する。
平方自由な積の次数は121である。$X=0,1,\infty$ 上の因子を除き、
多項式と上の三変換による引き戻しの最大公約元を安定するまで取る。
モニック化した結果は次数48であり、
\[
  D^6 g_1(N/D)g_2(N/D),\qquad
  g_1(X)=X^3+2X^2+2X+2,\quad g_2(X)=X^3+X^2+2
\]
のモニック化に等しい。従ってこれらの根を除く各ファイバーに良い表示がある。
二つの三次式の根ではそれぞれ $s^3+s^2+s+2$,
$s^6+2s^4+2s^3+s^2+2$ の根を選ぶ。対応する二つの剰余体で同時に
単元となる主元を選んでも、除算前の消滅は $\mathbb F_3(s)$ 上で成り立つ。
どちらの表示でも中心商の長さ6から13での最大損失は $(2,1,1,1,1,0,1,0)$ である。

一方、基準表示で $\mathbb Q(s)$ 上の十二桁を使い、
未切断経路、四中心、五次の積表に続いて六次の全十六行列を消去する。
全ての $r$ 除算で捨てる係数はこの有理関数体で厳密に零である。
主元の先頭の分子・分母を3で還元した平方自由な積は正確に $s^9-s$ であり、
$ND(N-D)$ を割る。従って $X\ne0,1,\infty$ のどの表示でも
六次までの同じ消去を混標数の形式級数環で行える。
六次の最大損失は二桁で、一桁を残して文字別 $(12,12,10,10)$ の基底を得る。
中心の自由座標とそれを指定する四次語基底は、標数3の基準表示と追加表示の双方で一致するため、六次までの自由な語の格子の左右の積を
標数3に還元し、同じ商を得る。十七桁の剰余計算では六次の後に六桁が残り、
その後の最大損失の和は五である。従って $3\in r^6\mathcal O$ なら
六次から再開して十三次まで進める。十三次の消滅まで全次元上界が $(1+t+t^2+t^3)^4$ の係数に一致するので、
第\ref{sec:rank-fifteen-overlap-annulus}節の自由性と左右の解消の議論が適用できる。
\end{proof}

% !TeX root = ../papers/nc-p3.tex
\section{分岐方向の二つの平方根と接線の剰余座標}
\label{sec:rank-fifteen-ramified-pole}
\begin{lemma}\label{lem:rank-fifteen-ramified-pole}
剰余標数3で $\bar x=1$, $\nu(y)>0$, $2\nu(x-1)\ge\nu(y)$ とする。
$J=(x-1)^2/(4y)$ は整である。付値を延長すると、$\nu(H)=\nu(y)/4$ と
単元 $z$ による表示があり、$\bar J=(z^4+1)^2/(4z^4)\bmod\mathfrak m$ となる。
$\bar J\ne1$ ならこの式のどの幾何学的剰余根も指定できる。
\end{lemma}
\begin{proof}
前節と同じ $U=1+r^2$, $K=\sqrt{1+r^2}$, $x=(4+r^2)T^2$,
$y=r^4T^6$ において $K\equiv1$, $2T\equiv1$ を選ぶ。
$h=r/(K+1)$, $H=h/\sqrt2$ とし、$v^2=1-2T$, $u^2=1+2T$ と置く。
恒等式 $1-2T=(1+r^2/4-x)/((1+r^2/4)(1+2T))$ と仮定から
$L=v/h$ は整である。第二の小さい平方根の中身は
\[
  2/(K+1)-2T=v^2-h^2.
\]
$C^2=L^2-1$ と取ると $(L+C)(L-C)=1$ なので $z=L+C$ は単元である。
$L=(z+z^{-1})/2$, $C=(z-z^{-1})/2$ によりこの平方根を $hC$ と書ける。
一方 $A=(u-v)/2$, $B=(u+v)/2$ とすれば $A^2+B^2=1$, $AB=T$ である。

$D=(1-2H^2)^{-1}$, $E=\sqrt{1-H^2 L^2}$, $F=\sqrt{1-H^2(L^2+1)}$ の
平方根の定数項を1とする。$K=(1+2H^2)D$, $r^2=8H^2D^2$ となり、
経路の係数に入る量は
\[
 \begin{gathered}
  rA=2HD(E-HL),\quad rB=2HD(E+HL),\quad T=(1-2H^2L^2)/2,\\
  P_\pm=\sqrt{4D\pm 2Tr^2},\qquad Q_+=4HDF,\quad Q_-=4H^2DC .
 \end{gathered}
\]
$P_\pm$ の定数項は2である。従って係数は $2z$ を逆にした環上で展開できる。
$x=(1+r^2/4)(1-h^2 L^2)^2$ と $h^4/(4y)\equiv1$ より
$(x-1)/h^2\equiv1-2L^2=-(z^4+1)/(2z^2)$ であり、所要の剰余式を得る。

$r^2$ を与える二つの単純な解は第\ref{sec:rank-fifteen-smith-slopes}節の
$(4+de)^3=e^2$ の $e\equiv\pm8$ に対応し、それらの $h^2$ の先頭は互いに逆符号である。
従って $(x-1)/h^2$ の剰余の二つの符号を選べる。
一方、$v,C$ の符号により $z,-z,1/z,-1/z$ を選べる。
残る符号の交換は $z\mapsto iz$, $i^2=-1$ に対応する。
$\bar J=0$ では $z^4=-1$ の四根は既に最初の四表示で覆われる。
$\bar J\ne0,1$ では八表示は相異なる。
\end{proof}

\begin{theorem}\label{thm:rank-fifteen-ramified-pole}
補題の条件の下、通常な階数15窓で $\bar J\ne1$,
$0<3\nu(y)\le2\nu(3)$ とする。この窓は標準Koszul左右AS正則であり、
各始点で四中心四次式上自由階数256である。特に $x=1$ の全ての通常点で
$0<3\nu(y)\le2\nu(3)$ ならこの結論が成り立つ。
\end{theorem}
\begin{proof}
\path{computations/verify_rank_fifteen_ramified_pole.py} はまず
$\mathbb F_3(z)$ 上で $H$ の十七桁を使い、$m=z^2+z+2$ で単元の主元を選ぶ。
未切断経路の長さ1から5では最大 $(0,2,2,2,2)$ 桁を使い、中心式で一桁、
五次の積表で一桁を使う。長さ6から13での最大損失は $(1,0,1,1,2,0,1,0)$
であり、十三次で一桁を残して商が消える。

逆元に用いた定数係数の分子・分母の平方自由な積は次数897である。
$z=0,\infty$ を外して、$z\mapsto-z,iz,1/z$ の引き戻しで共通の零点を取る。
$-z$ の処理後は偶多項式なので $iz$ の代入にも $\mathbb F_3$ 上の計算で足りる。
得られる次数100の平方自由式は、$z^4-1$ と次の五つの多項式の
$J=(z^4+1)^2/(4z^4)$ によるファイバーの積である（分母と定数倍を除く）。
同じ列に、所定のファイバーで使う $z$ の既約式を記す。
\[
\begin{array}{c|l}
 J+1 & z^4+z^2+2\\
 J^4+J^3+J^2+1 & z^4+z+2\\
 J^2+1 & z^4+z^3+z^2+z+1\\
 J^2+J+2 & z^4+2z^3+z^2+1\\
 J^3+J^2+2 & z^6+2z^5+2z^3+2z+1
\end{array}
\]
$z^4-1$ は $J=1$ のファイバーの被約部分である。
他の五つの代表で同時に単元となる主元を選んで同じ有理関数上の消去を行う。
未切断経路80組、中心4組、五次までの中心商96組、左右の積の照合128組で、
全308組の主元次数の列は基準表示と一致する。四次中心の自由座標を
指定する語基底と列番号も一致する。これで $\bar J\ne1$ の全ての幾何点に剰余表示がある。

標数零では基準表示で $\mathbb Q(z)$ 上の十二桁を使って長さ7まで進み、
各 $H$ 除算の前に定数係数の恒等的な消滅を検査する。
上記308組のうち最初の212組で標数3と同じ全指数を得る。
補題\ref{lem:rank-fifteen-smith-transfer}により、他の良い表示でもこの範囲の
語の格子と積表は混標数で持ち上がる。実際、二つの積の照合行列で前段の
語基底を替える操作は、左右の生成系と中間語の空間にそれぞれ可逆な基底変更を施す。
中心座標を揃えてあるので同一の商でこの帰納を行える。
七次の剰余積表に六桁が残り、仮定は $3\in H^6\mathcal O$ と同値である。
その六桁から再開した八次から十三次の次元上界は
$(1+t+t^2+t^3)^4$ の全係数に一致するので、前節と同じ自由性と左右の解消の判定が適用できる。
\end{proof}

\begin{theorem}\label{thm:rank-fifteen-collision-tube}
同じ通常窓で $\nu(y)>0$, $\bar x=1$ とする。
$\nu(J-1)\ge9\nu(y)$ なら、深さの上限なしに標準Koszul左右AS正則であり、
各始点で四中心四次式上自由階数256である。$J=1$ も含む。
\end{theorem}
\begin{proof}
ここでは $J$ が単元なので $2\nu(x-1)=\nu(y)$ である。
$r^2$ の二表示のうち $L\equiv0$ となる方を使い、$z\equiv i$, $i^2=-1$ と選ぶ。
$a=H^2$, $l=L^2$, $t=1-2al$, $d=2a/(1-2a)^2$ と置くと
$x=(1+d)t^2$, $2\sqrt y=dt^3$ である。実際の二因子は
\[
\begin{aligned}
(1-2a)^2(x-1+2\sqrt y)&=-4a(l-1)(al-1)(4a^2l-1),\\
(1-2a)^2(x-1-2\sqrt y)&=4al(al+a-1)(4a^2l-4a+1).
\end{aligned}
\]
従って $J-1$ は $L^2(L^2-1)$ の単元倍である。
$2zL=(z-i)(z+i)$ より $\nu(z-i)=\nu(J-1)/2\ge18\nu(H)$。
$J=1$ のときには $z=i$ となる。

補題の式で同じ $H$ を固定して $z$ を $i$ にすると、各経路係数の差は
$H^2(z-i)$ の倍数である。実際、$rA,rB,Q_-$ で最初に $z$ に依存する項は
$H^2 L$ または $H^2 C$ に単元を掛けたもので、残りの単元平方根の差は
根の和で割って求められる。従って窓の係数は $H^{20}$ まで一致する。
コードは $\mathbb Q(i)$ 上でこの二十桁を計算する。全定数項の消滅は厳密であり、
逆元のノルムを3で還元すると非零である。未切断の各長さで $(0,2,2,3,2)$、
中心で一桁、五次の中心商で一桁を使う。長さ6から13での最大損失は
$(1,0,1,1,2,0,3,0)$ で、十三次の後に一桁が残る。
四次の中心商も一桁を使うが、五次の中心商は未切断経路の基底から直接作る。
この有限精度の主元・除算を実際の窓へ移し、中心の一般の下界と合わせて
中心式を持ち上げる。商の次元上界が全て一致するので全次数の結論を得る。
\end{proof}

% !TeX root = ../papers/nc-p3.tex
\section{単元座標での極の全深度判定}
\label{sec:rank-fifteen-unit-pole-depth}
\begin{theorem}\label{thm:rank-fifteen-unit-pole-depth}
標数零の通常な階数15窓で、剰余標数3の付値に対して $\nu(x)=0$, $\nu(y)>0$ とする。
次のいずれかの場合には、付値の深さによらず標準Koszul左右AS正則であり、
各始点で四中心四次式上自由階数256である。
\begin{enumerate}
\item $\bar x\ne1$.
\item $\bar x=1$, $2\nu(x-1)\ge\nu(y)$ で $J=(x-1)^2/(4y)$ の剰余が1と異なる。
\end{enumerate}
\end{theorem}
\begin{proof}
\path{computations/verify_rank_fifteen_unit_pole_depth.py} は各表示に十七桁を取り、
$\mathbb Q(s)$ または $\mathbb Q(z)$ 上で、実際に使う全てのパラメータ除算の前に
定数項が恒等的に零であることを確認する。係数環はそれぞれ
$\mathbb Z[s]_{(3,s^4+2s^3+s+1)}$、$\mathbb Z[z]_{(3,z^2+z+2)}$ である。
主元の定数項の分子・分母の3での還元を記録する。

第一の場合、未切断経路と中心方程式、五次の中心商に九桁を使い、
長さ6から13の最大損失は $(2,1,1,1,1,0,1,0)$ である。
従って形式級数環上の語の格子の両積表を最後まで持ち上げられ、十三次で商は零となる。
使用した定数係数の平方自由な積は、同じ深さの $\mathbb F_3(s)$ 上の
計算で得た次数121の多項式に一致する。第\ref{sec:rank-fifteen-pole-fibers}節の
八重被覆で残る二つの三次式のファイバーは追加表示で覆う。

この比較では、基準と追加表示の全308行列のSmith指数の重複度を照合する。
基準の標数零の計算では単元主元が尽きたときに分数体でも定数項が零なので、
基準の剰余体での指数列が分数体上の指数列に一致する。
補題\ref{lem:rank-fifteen-smith-transfer}を順に適用すると、
追加表示でも同じ本数の主元が混標数の格子上で得られる。
中心四次式を指定する最後の始点の語座標は同じであり、
左右からの積と中間語の基底変更は照合行列の可逆な行・列変換になる。
語の格子の自由性は一般点の次元下界による。従って基準と追加表示のいずれでも、
正付値のパラメータへの特殊化後に同じ上界を与える。

第二の場合、基準表示で十二次までの292行列を標数零で消去する。
五次までに十桁、続く長さ6から12では $(1,0,1,1,2,0,1)$ 桁を使い、
一桁を残す。全ての指数列と中心座標が剰余のものに一致する。
ここでの主元の平方自由な積は、第\ref{sec:rank-fifteen-ramified-pole}節で
十三次までを含めて得た次数897の多項式と等しい。
同節の五つの追加表示でこれらの指数列を移すと、十二次までの同じ自由格子を得る。
十三次は各表示でパラメータ除算を使わず単元主元だけで全列を消す。
従って十二次で特殊化した積表の剰余から十三次の消滅を持ち上げられる。
両方の場合に $(1+t+t^2+t^3)^4$ の全係数と商の上界が一致し、
第\ref{sec:rank-fifteen-overlap-annulus}節の左右の解消と自由加群の議論が適用できる。
\end{proof}
\begin{corollary}
$x,y\in\mathbb Q$, $y\ne0$ が通常な階数15窓を与えるとする。
$x$ が3進単元で $x\equiv-1\pmod3$ なら、全ての $y$ で上の正則性と自由性が成り立つ。
$x\equiv1\pmod3$ でも $0<\nu_3(y)<2\nu_3(x-1)$ の全点が含まれる。
特に $\nu_3(x)=0$, $\nu_3(y)=1$ の全有理点が含まれる。
\end{corollary}
\begin{proof}
負付値の $y$ には極の定理、単元 $y$ には素体単元の四点の定理が適用できる。
正付値では第一の場合か、厳密な不等式により $\bar J=0$ となる第二の場合を使う。
\end{proof}

% !TeX root = ../papers/nc-p3.tex
\section{外側の極における二つの小さい変数}
\label{sec:rank-fifteen-outer-pole}
\begin{lemma}\label{lem:rank-fifteen-outer-pole}
剰余標数3で $0<2\nu(x-1)<\nu(y)$ とする。付値を延長すると、
$\nu(w)=\nu(x-1)/2$, $\nu(u)=(\nu(y)-2\nu(x-1))/4$ の表示が存在する。
$H=uw$, $D=(1-2H^2)^{-1}$, $T=1-2w^2$ とすると
\[
             x=(1+2H^2D^2)T^2,\qquad y=H^4D^4 T^6.
\]
経路係数は $(u,w)=(0,0)$ における $\mathbb Z_{(3)}$ 上の整な形式級数で表せる。
\end{lemma}
\begin{proof}
補題\ref{lem:rank-fifteen-ramified-pole}と同じ $r,K,h,H$ を選ぶと
$\nu(H)=\nu(y)/4$ である。同補題の $1-2AB$ の恒等式では
$\nu(r^2)>\nu(x-1)$ なので $w^2=(1-2AB)/2$ は
$\nu(w^2)=\nu(x-1)$ を満たす。$u=H/w$ は正付値である。
そこで同補題の $HL=w$, $HC=w\sqrt{1-u^2}$ を選び、
$E=\sqrt{1-w^2}$, $F=\sqrt{1-w^2-H^2}$ を代入する。
各平方根の定数項は1または2であり、2は単元なので全係数は上記の環にある。
この代入で $x,y$ の二式も得る。
\end{proof}

\begin{theorem}\label{thm:rank-fifteen-outer-pole}
標数零の通常な階数15窓で、剰余標数3において $0<4\nu(x-1)=\nu(y)$ とする。
この窓は標準Koszul左右AS正則で、各始点で四中心四次式上自由階数256である。
\end{theorem}
\begin{proof}
前補題で $\nu(u)=\nu(w)$ となるので、$s=w^2$, $m=u/w$ とすれば $H=ms$ で
$m$ は単元である。$X=(1+2m^2s^2D^2)T^2$, $Y=m^4s^4D^4T^6$ と置く。
$B=(X-1)/s$ は整で $B(0,m)=-4$ であり、
$\Lambda=Y/(X-1)^4=m^4D^4T^6/B^4$ は $s=0$ で $m^4/256$ に等しい。
$\partial_s X(0,m)=-4$, $\partial_m X(0,m)=0$,
$\partial_m\Lambda(0,m)=m^3/64$ より、どの非零な剰余根でもヤコビ行列は可逆である。
従って $\bar m^4=256\bar\Lambda$ の四根のうち任意のものを指定して
所定の $x,y$ に持ち上げることができる。

\path{computations/verify_rank_fifteen_outer_pole.py} は $\mathbb Q(m)$ 上の
三十二桁を用いる。$m-1$ と3の上で単元の主元を選ぶと、未切断経路での最大損失は
$(0,3,3,4,4)$、四次中心の各文字では $(0,2,2,2)$、四次・五次の積表では
$(0,3)$ である。長さ6から13の $(3,1,2,1,3,0,2,0)$ 桁を使って最後の一桁で
十三次の消滅を得る。各 $w$ 除算前の定数項は有理関数として零であり、
使用した定数項の分子・分母の剰余の平方自由な積は
\[
                    (m^9-m)/(m-1)
\]
である。従って $\bar m\notin\mathbb F_9$ または $\bar m=1$ ならこの表示が使える。
$\mathbb F_9^\times$ の四乗は $\{1,-1\}$ である。四乗が1なら $\bar m=1$ を選べる。
四乗が $-1$ なら $m^2+m+2$ の根を選ぶ。この既約式と3の上で単元の主元を取り直し、
同じ三十二桁の標数零計算を十三次まで行うと、用いた全ての主元の分子・分母は
剰余でこの既約式と互いに素である。全ての剰余値で中心商の上界と一般の下界が一致し、
第\ref{sec:rank-fifteen-overlap-annulus}節の判定が適用できる。
\end{proof}

% !TeX root = ../papers/nc-p3.tex
\section{二変数のSmith指数と外側の極の全域}
\label{sec:rank-fifteen-outer-pole-fan}
\begin{lemma}\label{lem:rank-fifteen-two-parameter-smith}
$R$ を $\mathbb Z_{(3)}[S,T]_{(3,S,T)}$ のHensel化とし、$M$ を $R$ 上の行列とする。
高さ一の素点 $(S),(T)$ で必要な本数までのSmith指数をそれぞれ昇順に
$a_j,b_j$ とする。剰余の対角線 $S=T=t$ 上の $\mathbb F_3[[t]]$ で
同じ本数の指数が $c_j=a_j+b_j$ なら、対角線の主元を用いて
$R$ 上で単項式 $S^{a_j}T^{b_j}$ を対角成分として順に分離できる。
\end{lemma}
\begin{proof}
$R$ は正則局所環で、$S,T$ は非同伴な素元である。全成分の両素点での付値が
$a_1,b_1$ 以上なので、$S^{a_1}T^{b_1}$ は全成分を割る。
この単項式で割った行列は対角線で単元主元を持ち、その主元は $R$ でも単元である。
これを消去すると、両素点での次の指数は $a_j-a_1,b_j-b_1$、
対角線では $c_j-c_1$ であり、同じ等式が続く。帰納し、分離の際に除いた
共通単項式を掛け戻せばよい。特に対角線で選んだ行・列と同じ語基底が使える。
\end{proof}

\begin{theorem}\label{thm:rank-fifteen-outer-pole-fan}
通常な標数零の階数15窓で、剰余標数3の付値について
$0<2\nu(x-1)<\nu(y)$ とする。この外側の領域は全て標準Koszul左右AS正則で、
各始点で四中心四次式上自由階数256である。
\end{theorem}
\begin{proof}
補題\ref{lem:rank-fifteen-outer-pole}の $(w,u)$ はともに正付値である。
$(\nu(w),\nu(u))$ の第一象限を、順に
\[
          (1,0),\quad(1,1),\quad(1,2),\quad(0,1)
\]
の隣り合う二辺で分ける。最初と最後は付値零の境界で、内部の二つの半直線だけを
別に扱う。$(1,1)$ は定理\ref{thm:rank-fifteen-outer-pole}による。

$(1,2)$ では $u=mw^2$, $s=w^2$ として $H^2=m^2s^3$ である。
$B=(x-1)/s$ の剰余は $-4$、$y/(x-1)^6$ の剰余は $m^4/4096$ であり、
前節と同じ単元ヤコビ行列によって所定の $x,y$ に持ち上がる。
$\mathbb Q(m)$ 上で四十八桁を使うと、未切断経路は $(0,4,4,6,6)$、
四次中心は $(0,3,3,3)$、五次の積は五桁を要する。長さ6から13では
$(5,2,3,2,4,0,3,0)$ 桁を使う。実際に用いた全ての主元の定数分子・分母の
剰余の平方自由な積は $m$ であり、従ってこの半直線の全ての単元に適用できる。

隣り合う二辺を $(a,b),(c,d)$ と書くと $ad-bc=1$ なので、開いた錐の内部では
$w=S^a T^c$, $u=S^b T^d$ と一意に書け、$S=w^du^{-c},\ T=w^{-b}u^a$ はともに正付値である。
全平方根は定数項を指定した単純根なので、経路係数は補題の $R$ に入る。
$(1,0)$ に沿う一般指数は $w$ を消す接方向のものである。実際
$HL=w$, $H=uw$, $z=(1+\sqrt{1-u^2})/u$ によって $z$ は境界の一般元となる。
$(0,1)$ に沿う一般指数は $u$ を消す非分岐極のもので、
$A=(\sqrt{1-w^2}-w)/\sqrt2$, $B=(\sqrt{1-w^2}+w)/\sqrt2$ は非定数の円錐曲線の表示である。
それぞれ一様化元は既存表示の $H$ または $r$ の単元倍であり、
第\ref{sec:rank-fifteen-unit-pole-depth}節の両指数列が使える。
内部の二辺ではそれぞれ $u=mw,\,u=mw^2$ における一般の $m$ を使う。
その一般指数は $m=1$ の剰余でのものと一致する。実際それを基準に選んだ各消去は、
単元主元が尽きた際の定数項の消滅を $\mathbb Q(m)$ 上で厳密に確認している。

\path{computations/verify_rank_fifteen_outer_pole_fan.py} は各錐の対角線、すなわち
$(w,u)=(t^2,t),(t^2,t^3),(t,t^3)$ でそれぞれ48、79、64桁の標数零計算を行う。
各錐の未切断経路80組、四次中心4組、五次までの中心商96組、長さ6から13の左右の積の照合128組で、
対角線の指数は両辺の昇順の指数の項別の和に一致する。比較は各行列の全主元に対して行う。

未切断経路から順に補題を適用し、対角線の語基底で自由格子と積の表示を作る。
各辺の離散付値環上では、既知の基底とこの基底が同じ経路格子を張るので、
基底変更は可逆である。中心の自由座標を指定する最後の始点の全四次語基底と列番号も
両辺と対角線で一致する。従って同じ四本の中心を課して帰納できる。
左右の積の照合では、中間の語基底の変更も行と列の可逆変換に過ぎない。
十三次まで各段の一般の次元下界に達する本数を消去でき、
十三次で商は零である。所定の窓で $S,T$ はともに非零なので各単項式での除算が使え、
従来の自由加群と左右の解消の判定を適用できる。
\end{proof}

% !TeX root = ../papers/nc-p3.tex
\section{内側衝突の実数係数の表示}
\label{sec:rank-fifteen-inner-collision}
\begin{theorem}\label{thm:rank-fifteen-inner-collision}
剰余標数3の付値で $\bar x=1$, $\nu(y)>0$, $J=(x-1)^2/(4y)$, $\bar J=1$ とする。
通常な標数零の階数15窓で
\[
       \nu(J-1)\ge\nu(y) \qquad\text{または}\qquad
       2\nu(J-1)=\nu(y)
\]
なら、標準Koszul左右AS正則で、各始点で四中心四次式上自由階数256である。
第一の場合は $J=1$ を含む。
\end{theorem}
\begin{proof}
補題\ref{lem:rank-fifteen-ramified-pole}の二つの $h^2$ の選択で
$(x-1)/h^2$ の剰余は符号を替えられる。ここでその剰余を $-1$ とすれば
$1-2L^2\equiv-1$、すなわち $L^2\equiv1$ である。
$L\equiv1$ を選び、$C^2=L^2-1$ として $z=L+C$ と置けば $z\equiv1$ となる。
定理\ref{thm:rank-fifteen-collision-tube}の二因子の恒等式より
$J-1$ は $L^2 C^2$ の単元倍である。
$C=(z-1)(z+1)/(2z)$ より $\delta=z-1$ は
$2\nu(\delta)=\nu(J-1)$ を満たす。$C$ の符号を替えると $z\mapsto1/z$ となり、
$\delta\mapsto-\delta/(1+\delta)$ である。$\nu(y)=4\nu(H)$ は両方の表示で共通である。

$\nu(\delta)\ge2\nu(H)$ では $z=1+mH^2$ と書け、$m$ は整である。
\path{computations/verify_rank_fifteen_inner_collision.py} は
$\mathbb Q(m)$ 上の十六桁で全ての $H$ 除算前の定数項の消滅を確認する。
未切断経路の最大損失は $(0,2,2,3,2)$、四次中心の各文字で $(1,0,0,0)$、
四次と五次の積表はともに損失零である。
長さ6から13では $(2,0,1,0,1,0,1,0)$ 桁を使い一桁を残す。
主元の定数分子・分母の剰余の平方自由な積は $m-1$ である。
$\bar m=1$ なら $z\mapsto1/z$ により $\bar m=-1$ を選べるので全ての整な $m$ を覆う。

$\nu(\delta)=\nu(H)$ では $z=1+mH$ と書き、$m$ は単元である。
十五桁の基準表示では四次中心も全て単元主元で取れ、他の損失は上記と同じである。
主元の平方自由な積の次数は20である。その零点集合と $m\mapsto-m$ による像の共通部分は、
$m=0$ と次の三式の根およびそれらの符号反転で尽くされる。
\[
           m^2+m+2,\qquad m^2+1,\qquad m^3+2m^2+1 .
\]
この三つの剰余拡大で同時に単元の主元を使う別の $\mathbb Q(m)$ 上の計算で、
同じ十五桁から十三次の消滅までを検証する。符号の選択により全ての単元剰余に適用できる。
両方の場合に各始点の中心商の上界と一般の下界が一致し、
第\ref{sec:rank-fifteen-overlap-annulus}節の全次数判定が従う。
\end{proof}

\begin{corollary}
$x,y\in\mathbb Q$, $\nu_3(x)=0$, $y\ne0$ の通常な階数15窓を考える。
$x\not\equiv1\pmod9$ なら全ての $y$ で上記の正則性と自由性が成り立つ。
$x$ によらず、$\nu_3(y)=2$ または $\nu_3(y)$ が奇数のときにも成り立つ。
\end{corollary}
\begin{proof}
$\nu_3(y)\le0$ は極の定理と素体単元の定理で決まる。
$\nu_3(y)>0$ のとき、第\ref{sec:rank-fifteen-unit-pole-depth}節と
定理\ref{thm:rank-fifteen-outer-pole-fan}により、残り得るのは
$\bar x=1$, $\bar J=1$, $\nu_3(y)=2\nu_3(x-1)$ の場合に限る。
従って奇数の付値は含まれない。$\nu_3(y)=2$ なら
$\nu_3(J-1)$ は1または2以上であり、本節の定理の二条件で尽くされる。
$\bar x=1$ かつ $x\not\equiv1\pmod9$ の場合も $\nu_3(x-1)=1$ なので同じ結論になる。
\end{proof}

% !TeX root = ../papers/nc-p3.tex
\section{内側衝突の残りの錐と単元座標の極}
\label{sec:rank-fifteen-inner-collision-fan}
\begin{theorem}\label{thm:rank-fifteen-inner-collision-fan}
標数零の通常な階数15窓が剰余標数3において $\nu(x)=0$, $\nu(y)>0$ を満たせば、
標準Koszul左右AS正則で、各始点で四中心四次式上自由階数256である。
\end{theorem}
\begin{proof}
第\ref{sec:rank-fifteen-unit-pole-depth}節と第\ref{sec:rank-fifteen-outer-pole-fan}節により、
$\bar x=1$, $\bar J=1$ だけを扱えばよい。前節の $z=1+\delta$ の表示を使い、
$(\nu(H),\nu(\delta))$ の順に傾きを書く。
$(1,2)$ の半直線とそれより深い衝突、および $(1,1)$ の半直線は前節で決まった。
残る範囲は境界 $(1,0)$ から $(2,1),(1,1),(1,2)$ までの三つの開いた錐と
$(2,1)$ の半直線に尽きる。

$(2,1)$ では $H=m\delta^2$、$m$ は単元である。$z=1+\delta$ のまま
$\mathbb Q(m)$ 上で三十桁を使うと、長さ1から5の損失は $(0,4,4,5,4)$、
四次中心は $(0,1,1,1)$、五次の積は一桁を使う。
長さ6から13では $(3,0,2,1,2,0,2,0)$ 桁を使い、
主元の定数分子・分母の剰余の平方自由な積は $m(m+1)$ となる。
$m+1$ の上で単元の主元に取り直した同じ標数零計算ではその積は $m(m-1)$ となる。
二表示で全単元を覆う。

三つの錐で両辺の原始ベクトルの行列式を取るといずれも1であり、外側と同じ単項式の
座標変換が使える。対角線は順に $(3,1),(3,2),(2,3)$ である。
\path{computations/verify_rank_fifteen_inner_collision_fan.py} は
これらの対角線でそれぞれ46、44、30桁を用いて全308組の指数と中心座標を求める。
$(1,0)$ の辺では $z$ が一般の接線座標、$(2,1)$ では $H/\delta^2$ が一般元、
$(1,1)$ と $(1,2)$ では $\delta/H$ と $\delta/H^2$ がそれぞれ一般元である。

今回の錐では、同じ文字の二つの中心式を選ぶ際の座標が辺によって変わる。
例えば最後の錐の対角線では、文字0の三座標のうち一つは始点2、残る二つは始点3にある。
一方 $\delta=mH$ のもとの表示では三座標とも始点3にある。
このため、各錐で対角線上の中心座標を使う。
座標を指定する各四次語基底を辺の既知の経路格子に写し、単元主元で
正方行列が可逆であることを確認する。この基底で中心方程式を作り直し、
指定列を自由座標として残すように他の列から主元を取る。
その後の四中心商はこの基底で指定した同じ二本ずつの中心式から計算する。
境界 $(1,0)$ は最初の対角線の座標で、辺 $(1,1)$ は最後の錐に入る際に
その対角線の座標でも、標数零の有理関数体上で検証する。

これにより各錐の両辺と対角線で選んだ中心は同じ規格化を持つ。
各行列について昇順に並べた両辺の指数の和が対角線の全指数と一致する。
未切断経路から四次中心、五次までの中心商、左右の積表の順に
補題\ref{lem:rank-fifteen-two-parameter-smith}を適用すれば、
所定の正付値の両座標で十三次の消滅までの上界が得られる。
従来の全次数判定から結論が従う。
\end{proof}

\begin{corollary}
$x,y\in\mathbb Q$, $\nu_3(x)=0$, $y\ne0$ の通常な階数15窓は、
全て標準Koszul左右AS正則で、各始点で四中心四次式上自由階数256である。
\end{corollary}
\begin{proof}
$\nu_3(y)>0$ は本節、$\nu_3(y)=0$ は素体単元の定理、
$\nu_3(y)<0$ は $\nu_3(y)<\min(0,2\nu_3(x))$ の極の定理である。
\end{proof}

\begin{corollary}\label{cor:rank-fifteen-unit-vertical-conics}
3の上のある付値で $\bar a=1$ とする。標数零の代数閉体上の縦切片 $x=a$ では、
$y\ne0$ の全ての水準が標準Koszul左右AS正則である。
$f=(1-a)^2-4a^3/27$ と置くと、添字保持同型類は
\[
      T\in\mathbb G_m,\qquad
      y=\frac{a^3}{27}-\frac{(T-f/T)^2}{16}\ne0
\]
で重複なく表される。非零水準のうち $a^3/27$ または $(1-a)^2/4$ に等しいものは二類、
他は四類であり、全類の最小周期は2である。特に $a\in\mathbb Q$, $a\equiv1\pmod3$
なら $y$ は任意の非零な体の元でよい。
\end{corollary}
\begin{proof}
付値を係数体へ延長する。正付値の $y$ は本節の定理、
単元は定理\ref{thm:rank-fifteen-vertical-all-levels}、負付値は極の定理が覆う。
$\bar a=1$ なので $a\ne3,3/4$、従って $f\ne0$ である。
定理\ref{thm:rank-fifteen-formal-boundary}の残り二点の横座標は
$4a^2-12a+7=0$ を満たし、$\bar a=1$ とは両立しない。
従ってこの縦切片の全ての点は通常表示をもつ。$p^2=a^3-27y$,
$q^2=(1-a)^2-4y$ と全点での三不変量の分離を用いると、
$T=q+2p/(3\sqrt3)$, $f/T=q-2p/(3\sqrt3)$ から上の式を得る。
$f\ne0$ により二つの分岐水準は異なる。$p,q$ の一方の平方が零なら二類、
両方非零なら四類となる。通常窓の最小周期は2である。
\end{proof}

% !TeX root = ../papers/nc-p3.tex
\section{非単元座標の極の円錐曲線表示と深い領域}
\label{sec:rank-fifteen-nonunit-pole}
\begin{lemma}
剰余標数3で $0<3\nu(x)<\nu(y)$ とする。付値を延長すると
\[
   x=(4+r^2)T^2,\qquad y=r^4T^6,\qquad
   T=\frac{2s(1-s^2)}{(1+s^2)^2}
\]
において $\nu(s)=\nu(x)/2>0$, $\nu(r)=(\nu(y)-3\nu(x))/4>0$ の表示が取れる。
この原点で経路係数は $\mathbb Z_{(3)}$ 上の $(r,s)$ の整な形式級数である。
\end{lemma}
\begin{proof}
$d^2=y/x^3$ は正付値なので、第\ref{sec:rank-fifteen-smith-slopes}節の
$(4+de)^3=e^2$ の単純根から $r^2=de$ を取る。
$T^2=x/(4+r^2)$ は正付値で、上の $s\mapsto T$ は原点で微分2の単元を持つ。
従ってHensel化で小さい $s$ を選べる。$A=(1-s^2)/(1+s^2)$,
$B=2s/(1+s^2)$ は $A\equiv1$, $B\equiv0$ である。
$K=\sqrt{1+r^2}\equiv1$、$P_\pm$ の定数項を2、
$Q_\pm/r$ の定数項を $1,-1$ と選べば、全ての平方根で根の和が単元となる。
これで補題の整な表示と付値式を得る。
\end{proof}

\begin{theorem}\label{thm:rank-fifteen-nonunit-pole}
剰余標数3の標数零の通常な階数15窓で $0<5\nu(x)\le\nu(y)$ とする。
この窓は標準Koszul左右AS正則で、各始点で四中心四次式上自由階数256である。
\end{theorem}
\begin{proof}
まず等号の場合には $r=ms$ と書け、$m$ は単元である。$v=s^2$, $L=4(1-v)^2/(1+v)^4$ とすれば
$T^2=vL$ であり、$x=v(4+m^2v)L$、
$\Lambda=y/x^5=m^4/((4+m^2v)^5L^2)$ となる。
$v=0$ で $x_v=16$, $x_m=0$, $\Lambda=m^4/16384$ なので、
$\bar m^4=16384\bar\Lambda$ のどの根でもヤコビ行列は可逆である。

\path{computations/verify_rank_fifteen_nonunit_pole.py} は $\mathbb Q(m)$ 上で
三十一桁を使い、未切断経路で $(0,1,1,4,4)$、四次中心で $(0,1,1,1)$、
五次の積表で六桁を消費する。長さ6から13では $(4,3,1,3,1,0,1,0)$ 桁で
最後の一桁を残して十三次の消滅を得る。除算前の定数項は有理関数として零である。
使用した主元の定数分子・分母の剰余の平方自由な積は正確に $m$ である。
従って $m$ が単元である全ての表示でこの計算が使える。
次に $\nu(y)>5\nu(x)$ では $\nu(r)>\nu(s)>0$ なので $r=ST$, $s=T$ と置けば
$S,T$ はともに正付値である。$S=0$ の辺の一般指数は第\ref{sec:rank-fifteen-unit-pole-depth}節の非分岐極、
$T=0$ の辺は今回の $r=ms$ の有理関数体上の指数である。対角線 $r=t^2,s=t$ では
四十七桁を使い、全308組の各行列で昇順の指数が辺の指数の項別和となる。
両辺と対角線で中心座標を指定する語基底と自由列も一致する。
補題\ref{lem:rank-fifteen-two-parameter-smith}によりこの開いた錐の全点で
十三次までの消去が使える。どちらの場合も中心商の上界が一般の下界に達し、全次数の判定が従う。
\end{proof}

% !TeX root = ../papers/nc-p3.tex
\section{非単元座標の残る帯と3整な有理点}
\label{sec:rank-fifteen-nonunit-pole-strip}
\begin{theorem}\label{thm:rank-fifteen-nonunit-pole-strip}
剰余標数3の標数零の通常な階数15窓で $\nu(x)>0$ なら、任意の $y\ne0$ について
標準Koszul左右AS正則であり、各始点で四中心四次式上自由階数256である。
\end{theorem}
\begin{proof}
$x=0$ は既知である。$\nu(y)<0$ は極の定理、$\nu(y)=0$ は
定理\ref{thm:rank-fifteen-residue-line}に含まれる。$\nu(y)>0$ で
$\nu(y)\le3\nu(x)$ なら定理\ref{thm:rank-fifteen-integral-depth}、
$\nu(y)\ge5\nu(x)$ なら前節が適用できる。

残る $3\nu(x)<\nu(y)<5\nu(x)$ は円錐曲線の座標で $0<\nu(r)<\nu(s)$ である。
$r=S$, $s=ST$ と置くと $S,T$ は正付値である。$S=0$ の辺は前節の $r=ms$ で、
一般元 $m$ は $T^{-1}$ である。$T=0$ の辺は $s=0$ で $r$ は一般元となる。
この第二の辺の関数体を有限拡大して
\[
 v=\frac{2z}{1+z^2},\qquad r=\frac{2v}{1-v^2},\qquad
 K=\frac{1+v^2}{1-v^2},\qquad E=\frac{1+z^2}{1-z^2}
\]
とすれば $K^2=1+r^2$, $2(K+1)=4E^2$, $2/(K+1)=E^{-2}$ である。
$z=0$ では $r\sim4z$ であり、$s=0$ では円錐曲線の $A,B$ は $1,0$ である。
従って前節と同じ平方根の枝を $s$ に関する展開で選べる。
具体的には $P_\pm$ の定数項は $2E$、$Q_\pm/r$ の定数項は $E^{-1},-E^{-1}$ である。

\path{computations/verify_rank_fifteen_nonunit_pole_strip.py} は
$\mathbb Q(z)$ 上で $s$ の十五桁を取り、$3,z^3+2z^2+1$ の上で単元となる主元を選ぶ。
未切断経路は $(0,0,0,2,2)$、四次中心は全て零、四次の積は零、五次の積は四桁を使う。
長さ6から13の最大損失は $(2,2,0,2,0,0,0,0)$ である。
全ての $s$ 除算の前に定数係数の消滅を有理関数体上で確認する。
これは高さ一の辺の一般指数を求める計算なので、$z$ の全剰余値への特殊化を必要としない。

対角線 $r=t,s=t^2$ を標数零で四十五桁まで展開する。未切断経路は
$(0,1,1,6,6)$、四次中心は $(0,1,1,1)$、五次の積は十桁を使い、
長さ6から13では $(6,5,1,5,1,0,1,0)$ 桁を使う。
両辺と対角線で中心の自由座標とそれを指定する語基底が一致し、全308組の行列で
両辺の指数の項別の和が対角線の指数になる。従って
補題\ref{lem:rank-fifteen-two-parameter-smith}により、残る帯の全点で
十三次までの中心商の上界と全次数の正則性・自由性が得られる。
\end{proof}

\begin{corollary}
$x,y\in\mathbb Q$, $\nu_3(x)\ge0$, $y\ne0$ の通常な階数15窓は全て
標準Koszul左右AS正則で、各始点で四中心四次式上自由階数256である。
\end{corollary}
\begin{proof}
正付値は本節、単元は定理\ref{thm:rank-fifteen-inner-collision-fan}とその系による。
\end{proof}

\begin{corollary}
3の上のある付値で $\bar a=0$ となる縦切片 $x=a$ は、非零の全幾何水準で分類できる。
$a\ne3,3/4$ では系\ref{cor:rank-fifteen-unit-vertical-conics}と同じ $T$ による表示、
分岐水準での二類とその他の四類、最小周期2という結論が成り立つ。
$a=3,3/4$ では規格化テンソルの縦切片は $T_+T_-=0$ の二直線から $y=0$ を除いたもので、
水準 $y=a^3/27$ にただ一つの最小周期1の点があり、他の非零水準は最小周期2の四類である。
\end{corollary}
\begin{proof}
本節の定理が任意の $y$ への付値の延長で適用できる。
通常表示外の四点のうち $4a^2-12a+7=0$ の二点は $\bar a=0$ に入らない。
残り二点とその周期は定理\ref{thm:rank-fifteen-formal-boundary}で決まっている。
$f(a)\ne0$ では既述の円錐曲線の計算を使う。
$f(a)=0$ では $p=q=0$ が二直線の交点であり、それ以外の非零水準では
$p,q$ がいずれも非零で四つの符号が相異なる点を与える。
\end{proof}

\begin{lemma}\label{lem:rank-fifteen-negative-unit-cover}
次の単元剰余 $\bar x=-1$ において、符号を選んだ倍率の関数体は、
幾何学的剰余体上の $\bar y$ 直線の次数384のGalois被覆を与える。
正規化した被覆曲線の種数は97で、分岐値は $0,1,\infty$ の三つである。
\end{lemma}
\begin{proof}
標数3で $v=1/K$, $i^2=-1$, $H=i(K-v)$ と取れば
$\bar y=-(v-v^3)^2$ である。$n\pm H/n$, $P_\pm$, $Q_\pm$ の
六根は $v$ 直線上の平方類としてそれぞれ
\[
 (1-v^2)(1\pm2iv),\qquad
 v(1+v)\{1\pm i(1-v)\},\qquad
 v(1-v)\{1\pm i(1+v)\}
\]
を開平して得る。初めの二根の和の半分が $n$ となり、差の半分との積は $H$ である。
六つの三次式はそれぞれに固有な零点を $\mathbb F_9\setminus\mathbb F_3$ に一つずつ持つ。
従って平方類は独立で、$v$ 直線への次数は64である。奇数付値の点の和集合は
$\mathbb P^1(\mathbb F_9)$ の十点で、幾何学的な各慣性指数は2である。
Riemann--Hurwitz公式は $2g-2=-128+10\cdot32=192$ を与える。

六つの変換 $v\mapsto\epsilon v+b$ ($\epsilon=\pm1$, $b\in\mathbb F_3$) は
$\bar y$ を保つ。検証コードは各平方類の奇数付値を十点の指示ベクトルで表し、
各変換後の六ベクトルを加えても $\mathbb F_2$ 階数が6であることを確認する。
幾何学的な有理関数体ではこの付値の偶奇が平方類を決めるので、
六変換が全て開平拡大に持ち上がる。$v$ から $\bar y$ への写像は有限非零の
$\bar y$ 上で非分岐である。九つの有限分岐点のうち $v\in\mathbb F_3$ は
$\bar y=0$ に、残る六点は $\bar y=1$ に写り、無限遠も分岐する。
\end{proof}

% !TeX root = ../papers/nc-p3.tex
\section{種数97の被覆の低次数閉点と数体上の水準}
\label{sec:rank-fifteen-negative-unit-levels}
\begin{theorem}\label{thm:rank-fifteen-negative-unit-levels}
剰余標数3の付値で $\bar x=-1$ とする。$y$ が単元で $\bar y$ の
$\mathbb F_3$ 上の次数が高々4なら、標数零の通常な階数15窓は
標準Koszul左右AS正則で、各始点は四中心四次式上自由階数256である。
\end{theorem}
\begin{proof}
一次は定理\ref{thm:rank-fifteen-unit-residues}で得た。
二次、三次、四次の各閉点の次数を $d$ とし、$\bar y$ の根を
$\mathbb F_{3^d}$ に取る。$c^2=-\bar y$ と $i^2=-1$ は
$\mathbb F_{3^{2d}}$ で解ける。$v-v^3=c$ の三根は互いに
$\mathbb F_3$ の定数だけ異なるので、$3^{2d}$ 乗のFrobeniusは
これらに平行移動として作用する。従って $v$ の全根は
$\mathbb F_{3^{6d}}$ にあり、前節の六つの平方根は全て
$\mathbb F_{3^{12d}}$ に取れる。

\path{computations/verify_rank_fifteen_negative_unit_levels.py} は素体上のモニックな
既約式を列挙する。各次数の個数は順に $3,8,18$ である。それぞれの一根に対して
$K=1/v$, $H=i(K-v)$ とし、二つの平方根 $\sqrt{K^2-1\pm2H}$ の和の半分を $n$ に取る。
両座標式と $n,P_\pm,Q_\pm$ の方程式を照合した上で四始点の経路次元を五次まで、
中心の文字次元を四次で求める。四本の中心式を課して十三次まで消去すると、
各始点の次元が $(1+t+t^2+t^3)^4$ の係数となる。
これらの計算はFrobenius共役にも成り立つ。指定した剰余表示をモニックな式で
持ち上げる定理\ref{thm:rank-fifteen-unit-residues}の証明が適用できる。
\end{proof}

\begin{corollary}
$x\in\mathbb Q$ が3整とし、数体 $L$ が3の上に剰余次数4以下の素点を少なくとも一つ持つとする。
このとき任意の $y\in L^\times$ に対する通常な階数15窓は全て上の正則性と自由性を持つ。
特に $[\mathbb Q(y):\mathbb Q]\le4$ なら常に成り立つ。
\end{corollary}
\begin{proof}
その素点の付値を窓の係数体に延長する。$\bar x=0,1$ の全水準は既に認証した。
残る $\bar x=-1$ では $\nu(y)>0$ は定理\ref{thm:rank-fifteen-inner-collision-fan}、
$\nu(y)<0$ は極の定理、$\nu(y)=0$ は上の定理が覆う。
\end{proof}

\begin{lemma}\label{lem:rank-fifteen-negative-x-coordinates}
$\nu(x)<0$, $\nu(y)\ge2\nu(x)$ とする。付値を延長し $i^2=-1$ とすると、
円錐曲線の一つの極 $s=i+w$ の近くで $r=gw$ と書け、
\[
    \nu(w)=-\nu(x)/4>0,\qquad \nu(g)=(\nu(y)-2\nu(x))/4\ge0 .
\]
$w=0$ の辺の一般点で $K=1$, $n=-ig$, $H/n=g$ であり、
$P_\pm^2=4\mp2ig^2$, $(Q_\pm/g)^2=\mp2i$ である。
$\nu(g)>0$ の場合は、全倍率の係数を $(w,g)$ の原点で整に選べる。
\end{lemma}
\begin{proof}
$\nu(y)>3\nu(x)$ なので、$d^2=y/x^3$ を用いる既述の単純根から
$x=(4+r^2)T^2$, $y=r^4T^6$、$\nu(r)>0$ の表示を取る。
$\nu(T)=\nu(x)/2<0$ である。$(1+s^2)^2=2s(1-s^2)/T$ はモニックで、
剰余式は $(s-i)^2(s+i)^2$ となる。$i$ に還元される根では右辺の分子は単元なので
$2\nu(w)=-\nu(T)$ であり、$g=r/w$ の式も従う。
$L=2i+w$ とすると
\[
   n=g(1-s^2)/L,\qquad H/n=2gs/L,\qquad K^2=1+g^2 w^2 .
\]
$K\equiv1$ を選ぶ。$w=0$ での値を代入し、$H=n(H/n)$ と
$2(K-1)=2g^2w^2/(K+1)$ を使うと所定の根号を得る。
原点では $P_\pm^2$ と $(Q_\pm/g)^2$ は全て単元である。
平方根の定数項を有限拡大上で選べばHensel化で各平方根が得られる。
\end{proof}

\begin{corollary}
この辺で単元の傾き $\lambda=\overline{y/x^2}$ に対する符号付きの表示は、
幾何学的剰余体上で種数1、次数16のGalois被覆を持つ。分岐値は $0,1,\infty$ である。
\end{corollary}
\begin{proof}
$y/x^2=H^2/(4+r^2)^2$ より $w=0$ で $\lambda=-g^4/16$ となる。
$g$ 直線へは $4-2ig^2$ と $4+2ig^2$ を開平すればよく、
それぞれが相異なる二つの単根を持ち、互いに根を共有しないため次数は4である。
無限遠の極は偶数位数で、幾何学的な四つの有限分岐点の慣性指数は2である。
従ってRiemann--Hurwitz公式は $2g_{\rm curve}-2=-8+4\cdot2=0$ を与える。
$g\mapsto ig$ が二つの被開平式を交換するため、$\lambda$ 直線上の四変換は全て持ち上がる。
四つの有限分岐点では $g^4=-4$、従って標数3で $\lambda=1/4=1$ であり、
四乗写像は零と無限遠にも分岐する。

計算には既存の楕円関数体 $W^2=t-t^3$ がそのまま使える。実際
$A=(1-t^2)/(1+t^2)$, $B=2t/(1+t^2)$ とし、
\[
 P_+=2(A-B),\quad P_-=2(A+B),\quad
 g=2(1-i)W/(1+t^2),\quad \lambda=4(t-t^3)^2/(1+t^2)^4
\]
と置く。二乗すれば辺の式に一致し、一般点で $A,B$ は $P_\pm$ から、
$t=B/(1+A)$ と $W$ は $A,B,g$ から戻るので双有理である。
$(Q_+/g,Q_-/g)$ の定数には $(1-i,1+i)$ が取れる。
\end{proof}

\begin{theorem}\label{thm:rank-fifteen-negative-x-prime-slopes}
$\nu(x)<0$, $\nu(y)=2\nu(x)$ で $\lambda=\overline{y/x^2}=\pm1$ なら、
通常な階数15窓は標準Koszul左右AS正則で、各始点は四中心四次式上自由階数256である。
従ってこの境界の有理座標は全て正則である。
\end{theorem}
\begin{proof}
所定の $x,y$ に対し、$g^4=-16\lambda$ の指定した剰余根を選べることを確認する。
$z^4=1/x$, $w=zu$, $s=i+w$, $L=2i+w$ とおく。
前出の $H$ を用い、次の二式を $u,g$ について解く。
\[
  1=\frac{u^4L^4}{(4+g^2w^2)4s^2(1-s^2)^2},\qquad
  y/x^2=H^2/(4+g^2w^2)^2 .
\]
$z=0$ では右辺は $-u^4/4$, $-g^4/16$ で、各未知数は単元、
ヤコビ行列の行列式も単元である。Hensel化で所定の剰余根が持ち上がる。
二つの傾きで $\mathbb F_{3^8}$ に $g,P_\pm$ を取り、$K=1$, $n=-ig$, $H/n=g$ から
四つのテンソルを作る。コードは両方で経路長5までの一般次元、四次中心の
文字次元 $(3,2,0,0)$、および全四始点で十三次までの標準中心商を直接検証する。
$P_\pm$ の一方が零に還元される傾きでも、モニックな平方根は所定の剰余根の上に持ち上がる。
四次中心の持上げと中心商の上界から、既述の全次数判定を適用できる。
\end{proof}

% !TeX root = ../papers/nc-p3.tex
\section{負付値座標の境界の全傾きと対角線の近傍}
\label{sec:rank-fifteen-negative-x-boundary}
\begin{theorem}\label{thm:rank-fifteen-negative-x-boundary}
剰余標数3の付値で $\nu(x)<0$, $\nu(y)=2\nu(x)$ を満たす通常な
階数15窓は、剰余傾き $\lambda=\overline{y/x^2}$ の全水準で
標準Koszul左右AS正則であり、各始点は四中心四次式上自由階数256である。
\end{theorem}
\begin{proof}
前節の楕円表示 $W^2=t-t^3$ を使う。\path{computations/verify_rank_fifteen_negative_x_boundary.py}
は $\mathbb F_3(t)$ 上の基底 $1,i,W,iW$ でテンソルを構成する。
四次中心から文字0の二つと文字1の二つを選び、全四始点の十三次までの商を計算する。
逆にした元の分子のノルムと入力分母を除けば、同じ単元主元による計算が特殊化できる。
この平方自由な主元多項式の次数は 118054 である。

$N=4(t-t^3)^2$, $D=(1+t^2)^4$ とすれば $\lambda=N/D$ である。
$0,1,\infty$ の逆像を除いた後、$t\mapsto-t,\ 1/t,\ (1+t)/(1-t)$ の
引き戻しとの最大公約元を安定するまで取る。この八次写像を保つ八変換は
非分岐ファイバー上で推移的である。安定した多項式は次数 840 であり、
有限剰余体で各既約因子の $N/D$ の最小多項式を計算する。
その $N/D$ による引き戻しが、対応する既約因子の積に一致することも確認する。
残った閉点の次数 $1,2,3,4,5,6,8,12$ に対する個数はそれぞれ
$1,2,6,4,2,6,1,1$ である。全23閉点について各一根で平方根を取り、
経路・四次中心と全四始点の中心商を十三次まで再計算する。
この検証はFrobenius共役に及ぶ。分岐水準 $\lambda=1$ は前節で扱った。
各傾きで選んだ $g$ は前節の単元ヤコビ行列で所定の $x,y$ に持ち上がる。
楕円表示で選んだ平方根の符号を取り、中心商の上界を持ち上げると結論が得られる。
\end{proof}

\begin{proposition}\label{prop:rank-fifteen-negative-x-diagonal}
同じ負付値座標の $w,g$ 表示で $\nu(w)>0$ とする。
$\nu(g-w)\ge21\nu(w)$ なら上と同じ正則性と自由性を持つ。
\end{proposition}
\begin{proof}
$w=g=t$ として $\mathbb Q(i)[[t]]$ の二十一桁で計算する。
$A'=(1-s^2)/(2i+w)$, $B'=2s/(2i+w)$, $s=i+w$ とすれば
$n=gA'$, $H/n=gB'$ である。$K$ の定数項は1、$P_\pm$ は2、
$Q_\pm/g$ は $1\mp i$ として平方根を取る。
3の上で単元となる主元を用い、$t$ で割るときは実部・虚部の定数係数がともに
標数零で消えることを確認する。未切断経路の損失は $(0,1,1,3,2)$、
四次中心は $(0,1,2,2)$、五次の積は二桁を使う。
四中心商の六次から十三次の損失は $(4,1,1,1,1,0,1,0)$ となる。
$g=w+w^{21}u$ で $u$ が整なら、単純根の選択により経路係数は
対角線のものと $(w^{21})$ を法として一致する。全消去を通じて一桁が残るため、
この摂動でも中心の持上げと標準商の上界が得られる。
\end{proof}

$g=0$ の一般の辺では $r=gw$ が $g$ の単元倍で、$s=i+w$ は一般の円錐曲線座標である。
従って既知の非分岐極の指数を使える。$w=0$ の辺では上の中心商計算が示すように
一般点で未切断経路の全指数は零である。一方 $w=g$ の未切断四次の16行列のうち
8行列で、指数列は $g=0$ の辺と異なる。このため、この二辺だけで
二変数のSmith指数の補題を適用するための項別等式は成り立たない。
細分の候補 $g=mw$ で $\nu(m)=0$ なら $\nu(y)=\nu(x)$、
$\overline{y/x}=\bar m^4/4$ である。実際 $xw^4\equiv-4$,
$H/w^2\equiv-i\bar m^2$ と前節の $y/x^2$ の式を用いればよい。

% !TeX root = ../papers/nc-p3.tex
\section{負付値の二つの曲線近傍と境界指数の比較}
\label{sec:rank-fifteen-negative-x-arcs}
\begin{proposition}\label{prop:rank-fifteen-negative-x-arcs}
負付値座標 $w,g$ がともに正付値であるとする。通常な階数15窓が
\[
   \nu(g-w^2)\ge37\nu(w)\quad\text{または}\quad
   \nu(w-g^2)\ge21\nu(g)
\]
を満たすなら、標準Koszul左右AS正則であり、各始点は四つの中心四次式上で自由階数256である。
初めの近傍は $\nu(y)=0$ に、後の近傍は $2\nu(y)=3\nu(x)<0$ に属する。
\end{proposition}
\begin{proof}
前節の整な倍率表示を $(w,g)=(t,t^2)$ と $(t^2,t)$ に制限する。
平方根の定数項は前節と同じである。
\path{computations/verify_rank_fifteen_negative_x_arcs.py} は
$\mathbb Q(i)$ で展開し、3で還元したノルムが非零の主元を選ぶ。
前者の未切断経路の損失は $(0,2,2,4,4)$、中心は $(0,2,4,4)$、五次の積は四桁を使う。
六次から十三次は $(6,2,2,2,2,0,2,0)$ であり、三十七桁の最後の一桁が残る。
後者の各損失は 順に $(0,1,1,3,2)$、$(0,1,2,2)$、二桁、$(4,1,1,1,1,0,1,0)$ である。
両曲線で全四始点の中心商は十三次で消え、単純根と消去を上記の桁数で比較すると
前節と同様に近傍へ結論が伸びる。$\nu(x)=-4\nu(w)$,
$\nu(y)=4\nu(g)-8\nu(w)$ から水準の付値式も従う。
\end{proof}

単元 $m$ を残す二つの半直線 $g=mw^2$, $w=mg^2$ 上では、それぞれ
$\bar y=-\bar m^4$, $\overline{y^2/x^3}=-1/(64\bar m^4)$ となる。
実際 $xw^4\equiv-4$、$H/g^2\equiv-i$ を $y=x^2H^2/(4+r^2)^2$ に入れればよい。
$(t,t^2)$ の全308組の指数列を非分岐の辺と前節の対角線の項別和と比べると、
一致は254組であり、未切断四次でも16行列のうち4行列が異なる。
$(t^2,t)$ と前節の対角線では 307組が一致する。唯一異なるのは文字1の四次中心で、三つの指数が零から1に増える。最大指数だけを比べるとこの差は見えない。
ここで比較した対角線は倍率 $m=1$ の計算である。一般の半直線で
$\mathbb Q(i)(m)$ 上の指数と例外倍率を求める必要が残る。

% !TeX root = ../papers/nc-p3.tex
\section{円錐曲線の極の平方座標と負付値の閉領域}
\label{sec:rank-fifteen-negative-x-fold}
\begin{theorem}\label{thm:rank-fifteen-negative-x-fold}
剰余標数3の付値で $10\nu(y)\le19\nu(x)<0$ となる標数零の通常な階数15窓は
標準Koszul左右AS正則であり、各始点で四中心四次式上自由階数256である。
\end{theorem}
\begin{proof}
$\nu(y)\le2\nu(x)$ は定理\ref{thm:rank-fifteen-negative-x-boundary}と既知の極領域による。
残る範囲では $s=i+w$, $r=gw$ の表示で $w,g$ はいずれも正付値である。
$B=2s/(1+s^2)$, $A=(1-s^2)/(1+s^2)$ として
\[
 h=1/B,\qquad z=h^2,\qquad u=rB,\qquad v=A/B
\]
と置く。$u/g$ と $h/w$ は単元であり、$v^2=z-1$, $v\equiv-i$ である。
倍率は $n=vu$, $H/n=u$, $K^2=1+zu^2$, $H=vu^2$ を満たす。
$AB=v/z$ より正則性座標は
\[
 x=(4+zu^2)(z-1)/z^2,\qquad y=u^4(z-1)^3/z^4
\]
となる。従って $h$ の符号を使わず、変数 $z,u$ だけで倍率を展開できる。

$K\equiv1$ とし、$P_\pm^2=2(K+1\pm vu^2)$ の根の定数項を2、
$(Q_\pm/u)^2=2z/(K+1)\pm2v$ の根の定数項を $1\mp i$ とする。
全て単純根である。$z=0$ の値と比べると $v$ と $Q_\pm/u$ の差は $(z)$ に、
$K$ と $P_\pm$ の差は $(zu^2)$ に入る。従って四次テンソルを与える各係数の差は
$(zu)$ に入る。これは整な局所環のイデアルの包含である。

\path{computations/verify_rank_fifteen_negative_x_fold.py} は $z=0$ で
$\mathbb Q(i)$ 上の $u$ の二十一桁を使う。未切断経路の損失は $(0,1,1,3,2)$、
四次中心は $(0,1,2,2)$、五次の積は二桁、六次から十三次は $(4,1,1,1,1,0,1,0)$
であり、還元可能な単元主元だけで最後の一桁が残る。
全308行列の指数と中心の自由語座標は前節の $w=g^2$ のものと一致する。
仮定は $\nu(w)\ge10\nu(g)$、従って $\nu(z)\ge20\nu(u)$ に等しい。
このときテンソルの差は $(u^{21})$ に入り、同じ桁数の消去をそのまま持ち上げられる。
経路、四次中心、十三次での中心商の消滅から既述の全次数の判定が従う。
\end{proof}

% !TeX root = ../papers/nc-p3.tex
\section{平方座標の三つの境界による負付値縦切片の完全判定}
\label{sec:rank-fifteen-negative-x-fan}
\begin{theorem}\label{thm:rank-fifteen-negative-x-fan}
標数零の通常な階数15窓で、3の上のある付値に対して $\nu(x)<0$ なら、
任意の $y\ne0$ で標準Koszul左右AS正則である。各始点は四中心四次式上で自由階数256である。
\end{theorem}
\begin{proof}
$\nu(y)\le2\nu(x)$ は既知なので、前節の $z=1/B^2$, $u=rB$ がともに正付値の場合を考える。
$(\nu(u),\nu(z))$ 平面で $z=mu^d$、$d=1,2,3$ の三つの境界を用いる。
$d=3$ の表示は $m$ を整とすると、$\nu(z)\ge3\nu(u)$ の全範囲を含む。

\path{computations/verify_rank_fifteen_negative_x_fan.py} は、まず $m$ が一般元のとき
$\mathbb Q(i)(m)[[u]]$ で消去する。係数を実部と虚部に分け、それぞれ
$\mathbb Z[m]$ 上の共通分母で保持する。主元は3と $m$ の奇数次既約剰余式の上で単元に取る。
この剰余体では $-1$ は非平方なので、実部・虚部の少なくとも一方が単元ならよい。
逆元を取るときに使うノルム（実または純虚のときはその非零な座標）の
定数係数の分子・分母を記録し、$u$ 除算では両定数係数の標数零での消滅を検査する。
前節の根の定数項を使うため、初期分母は2の冪しか必要としない。

$d=1$ では十九桁を用い、未切断経路、四次中心、五次の積、六次から十三次の順の損失は
$(0,1,1,2,2)$、$(0,1,2,2)$、二桁、$(3,1,1,1,1,0,1,0)$ である。
$d=2,3$ では共に二十一桁で、これらは
$(0,1,1,3,2)$、$(0,1,2,2)$、二桁、$(4,1,1,1,1,0,1,0)$ となる。
主元多項式の剰余の平方自由な積は順に $m(m^4+1)$、$m(m+1)$、$1$ である。
第一の辺では $m^4+1=(m^2+m+2)(m^2+2m+2)$ の四根を、
第二の辺では $m=-1$ を検査する。第三の辺は $m=0$ を含め例外を持たない。
記録した多項式が消えない剰余では同じ計算が使える。
例外剰余ではその根を含む有限体で計算し、中心の自由列とそれを指定する語基底、
全308行列でのSmith指数列を重複度まで一般の辺と照合する。
補題\ref{lem:rank-fifteen-smith-transfer}を $m$ に関する局所環と一様化変数 $u$ に適用すれば、
この比較で一致した剰余の上の全ての持上げでも消去できる。
$d=1,2$ では $m=0$ は開いた隣接領域に属し、境界で検査するのは単元だけである。

残る三錐の辺の原始ベクトルは $(1,0),(1,1),(1,2),(1,3)$ で、
隣り合う二辺の行列式は1である。$(1,0)$ の辺では $z$ は一般元であり、
$h^2=z$ を添加すれば $r=hu$ が $u$ の単元倍である。
第\ref{sec:rank-fifteen-unit-pole-depth}節の一般円錐曲線で $v=A/B$,
$(Q_+/u,Q_-/u)|_{u=0}=(1+v,1-v)$ となり、同じ符号の規格化を与える。
この既知の辺と三つの一般境界を順に使う。

各錐の対角線 $(u,z)=(t^2,t),(t^2,t^3),(t^2,t^5)$ で、剰余体 $\mathbb F_9$ 上それぞれ
35、39、41桁を用いる。いずれも必要な中心の自由列とそれを指定する語基底が両辺で一致し、
全308行列について、昇順の指数の項別和が対角線の全指数に一致する。
補題\ref{lem:rank-fifteen-two-parameter-smith}は $i$ を添加した3上の非分岐局所環でも同じ証明で成り立つ。
各錐の単項式座標で未切断経路、四次中心、左右の積と中心商へ順次適用する。
十三次での消滅と全始点の上界から、従来の全次数判定で結論が従う。
\end{proof}

\begin{corollary}
$\nu(a)<0$ となる3の上の付値がある縦切片 $x=a$ では、全非零水準の添字保持同型類が
系\ref{cor:rank-fifteen-unit-vertical-conics}の $T$ の式で分類される。
分岐水準 $a^3/27,(1-a)^2/4$ で各二類、その他で各四類となり、最小周期は全て2である。
さらに任意の有理数 $x$ と非零の代数的数 $y$ で、$\mathbb Q(y)$ が
3上に剰余次数4以下の素点を持てば、全ての通常窓が上の正則性と自由性を持つ。
特に全ての有理座標対 $(x,y)$ が認証される。
\end{corollary}
\begin{proof}
通常表示外の四点の横座標 $3,3/4$ および $4a^2-12a+7=0$ の二根は全て3整である。
$(1-a)^2-4a^3/27$ の零点も $3,3/4$ なので、定理と既知の円錐曲線表示が縦切片の全点で使える。
数体については、$x$ が3整なら第\ref{sec:rank-fifteen-negative-unit-levels}節の系、
負付値なら上の定理による。
\end{proof}

% !TeX root = ../papers/nc-p3.tex
\section{一般単元剰余の二つの四次方程式と二次剰余の全点}
\label{sec:rank-fifteen-unit-quartic-towers}
\begin{theorem}\label{thm:rank-fifteen-unit-quartic-towers}
$k$ を標数3の代数閉体とし、$X=\bar x\ne0,1$ を固定する。
$c^2=X$, $b=1/c$, $d=b^3-b$ と置く。符号付き倍率の関数体の一表示は
\[
 D=v^3-v,\quad \bar y=X^3D^2,\quad
 R_{a,e}^2=f_{a,e}=(v-a-e b)\!\prod_{j\in\mathbb F_3\setminus\{a\}}(v-j)
 \qquad(a\in\mathbb F_3,\ e=\pm1)
\]
で与えられる。この被覆は $\bar y$ 上次数384、種数97であり、分岐値は $0,(1-X)^2,\infty$ である。
さらに $\eta_e=(\prod_a R_{a,e})/D$, $T=\eta_-+\eta_+$ とすると $k(T)$ は
$\bar y$ 上次数8の有理中間体である。その上の全被覆は次の二四次式の同時分解体である。
\[
 Z^4-2e b Z^2-8D\eta_e Z+b^2-8e bD,\qquad
 \eta_--\eta_+=2d/T,\qquad D=(T^2+4d^2/T^2)/4 .
\]
同時分解体の $k(T)$ 上の次数は48で、各四次式の分解体は次数12、Galois群は $A_4$ である。
\end{theorem}
\begin{proof}
$K=1/v$, $H=c(K-v)$ なら $U=K^2$, $V=H^2$ から二つの正則性座標が得られる。
$n\pm H/n$, $P_\pm$, $Q_\pm$ の六つの二乗に $v^2$ を掛けると、順に
$(1-v^2)(1\pm2cv)$、$2v(1+v)(1\pm c(1-v))$、$2v(1-v)(1\pm c(1+v))$ となる。
定数倍と符号の付け替えで上の六つのモニック三次式になり、定数は $k$ で開平できる。
初めの二根の和の半分から $n$ も復元できる。

$b\notin\mathbb F_3$ なので $\mathbb F_3$, $\mathbb F_3+b$, $\mathbb F_3-b$ は交わらない。
各三次式には他の五式が消えない単根 $a+e b$ があり、六平方類は独立である。
$v$ への64重被覆の分岐はこの九点と無限遠で、各慣性指数は2であるから
$2g-2=-128+10\cdot32=192$ となる。$v\mapsto\epsilon v+j$ は
$\bar y$ を保ち、$f_{a,e}$ を定数倍の置換に移すので、六変換とも持ち上がる。
三つの有限剰余類のうち元の素体は $D=0$、残る二類は $D=\pm d$ に写る。
$X^3 d^2=(1-X)^2$ であり、$v\mapsto D$ は有限点で非分岐である。

$f_{a,e}=D(1-e b/(v-a))$ より $\prod_a f_{a,e}=D^2(D-e d)$ である。
従って $\eta_e^2=D-e d$、$\eta_-^2-\eta_+^2=2d\ne0$ となり、
両根は $T$ の有理式で戻る。$D-d$, $D+d$ は異なる単根を持つので
$k(T)/k(D)$ は次数4であり、$\bar y$ への次数は8である。
$\sum_a f_{a,e}=e b$, $\sum_{a<a'} f_{a,e}f_{a',e}=2e bD$ は三次式の展開から従う。
三根 $R_{a,e}$ の符号を偶数個だけ反転した四つの和を取ると、
積 $D\eta_e$ が固定される。四つの和を根とするモニック四次式の展開が定理の式である。
四根のうち二つの和の半分から各 $R_{a,e}$ を戻せる。さらに
$f_{0,e}-f_{1,e}=-e b(2v-1)$ から $v$ が戻るので、同時分解体は全関数体である。
$k(T,v)/k(T)$ は、次数3の $k(v)/k(D)$ と次数4の $k(T)/k(D)$ が線型無関連なので次数3である。
この上では各三根の偶数符号変更が独立に作用し、合わせて $4\cdot4$ 個ある。
$v\mapsto v+1$ は二組の添字を同時に3巡回させ、各 $\eta_e$ を固定する。
従って全次数は48で、一方の四根に作用する像は $V_4\rtimes C_3=A_4$ である。
\end{proof}

\begin{corollary}\label{cor:rank-fifteen-unit-elliptic-quotient}
一方の四次式の根 $Z$ を選び、$\eta=\eta_e$ とする。点を一つ指定した関数体
$k(\eta,Z)$ は種数2であり、両符号とも同じ楕円曲線
\[
                 Y^2=(l+1)(l+1+Xl^2)
\]
の二重被覆になる。具体的には $s=e b$, $a=Z^2/s$,
$l=(Z^2+\eta Z)/s$, $Y=2a-l-1$ である。
\end{corollary}
\begin{proof}
標数3では四次式は $(Z^2-s)^2+D(\eta Z+s)=0$、
$D=\eta^2+s^3-s$ である。$A=Z^2$, $B=\eta Z$ を入れて $A$ 倍すると
$A(A-s)^2+(B^2+(s^3-s)A)(B+s)=0$ となる。
$A=sa$, $B=s(l-a)$ と変換すると $a^2-(l+1)a=Xl^2(l+1)$ に等しく、
平方完成して上の式を得る。二次因子の判別式は $1-X\ne0$、
$l=-1$ での値は $X\ne0$ なので三次式は平方自由である。
$Z^2=sa$ を添加すると $\eta=s(l-a)/Z$ も戻る。
楕円曲線上で $a=0$ は $(l,a)=(0,0),(-1,0)$ に限られる。
前者では $a-(l+1)$ が単元で $a$ の零位数は2である。
後者では $Y$ が一様化元で $l+1$ の位数が2なので、$a=(Y+l+1)/2$ は単純零である。
無限遠では $l,Y$ の極の位数は2,3で、$a$ は三位の極を持つ。
二重被覆の分岐点はこの単純零と無限遠の二点であり、種数は2となる。
さらに $t=l+1$, $W=Y/c$ で $W^2=t^3+(1+1/X)t^2+t$ となる。
標数3のHasse不変量は $t^2$ の係数なので、この楕円曲線が超特異となるのは
$X=-1$ のときに限る。$j=(X+1)^6/(X^4(1-X))$ である。
特に $X=-1$ で $l=i t-1$ と置くと $Y^2=-i(t-t^3)$ となり、
定数倍で既出の $W^2=t-t^3$ の関数体をそのまま利用できる。
\end{proof}

\begin{proposition}\label{prop:rank-fifteen-genus-thirteen-window}
上の六根を含む体を $E$ とする。四窓テンソルは $k(v)$ 上次数16、種数13の体
\[
 F=k(v)\big(\sqrt{(v-a)^2-b^2}\ (a\in\mathbb F_3),\ \eta_+\big)
\]
に係数を持つ表示を持ち、その表示でもKlein次数に関して斉次である。
\end{proposition}
\begin{proof}
同じ $a$ の両根を同時に反転する操作を対反転と呼ぶ。
$S$ は $1,-1$ の二対、$P$ は $0,-1$ の二対を反転し、位数4の群を生成する。
$R_{a,+}R_{a,-}$ は $D/(v-a)$ と上の平方根の積である。
三つの対積の符号指標と、三つの正号根の積の符号指標は独立である。
従ってこれらで生成した体の次数は16で、これが固定体である。
三つの二次式の単根は $a\pm b$ の六点で、$\eta_+^2=D-d$ の三単根もそのうちにある。
さらに無限遠で後者が奇数位数の極を持つ。他に分岐はなく各慣性次数は2なので、
$2g-2=-32+7\cdot8=24$ である。

添字 $t\in\{0,1,2,3\}$ の上下ビットを $t_H,t_L$ と書く。
正規形の倍率 $\gamma_{h,j}$ で $S$ の作用は $(-1)^{h_H}$ 倍、
$P$ の作用は $(-1)^{j_L}$ 倍である。これは $(L,M),(A,B),(C,D)$ が順に
$0,1,-1$ の各対に線型であることから従う。
四つの位置すべてで低ビットをWalsh基底 $e_{2r}+(-1)^q e_{2r+1}$ に変える。
第1、第3位置ではさらに添字 $t=2r+q$ の係数に
$\phi_t=R_{1,+}^{r}R_{0,+}^{q}$ を掛ける。この変換は関数体上可逆である。
$\beta=(0,3,2,1)$ と置くと、変換後の係数は
$a\mathbin\oplus c=b\mathbin\oplus d=h$ のとき
\[
 4\phi_a\phi_c\sum_{k=0}^1\ \,\sum_{j_L=h_L}
 (-1)^{(a_L+b_L)k+j_Hc_H+\beta(j)_Hd_H}\,
 \gamma_{2h_H+k,j}
\]
であり、それ以外は零となる。実際、旧低ビット $c_L,d_L$ で符号を足すと、
$\beta(j)_L=j_L$ により、二つの直交条件はいずれも新しい低ビット差を $j_L$ と等しくする。
この式での $S,P$ の符号を $\phi_a\phi_c$ が打ち消すので係数は $F$ に入る。
添字付きの基底変換なので、この表示に四次中心も移せば各始点の商の次元が保たれる。
\end{proof}

\begin{theorem}\label{thm:rank-fifteen-quadratic-residue-pairs}
数体 $L$ が3上に剰余次数2以下の素点を持つとする。任意の $x\in L$, $y\in L^\times$ で、
通常な階数15窓は標準Koszul左右AS正則で、各始点が四中心四次式上自由階数256である。
特に、$x,y$ が任意の有限個の二次数体の合成体に属すればよい。
\end{theorem}
\begin{proof}
素点の付値を窓の係数体に延長する。負付値の $x$ は前節、正付値または零の $x$ は
第\ref{sec:rank-fifteen-nonunit-pole-strip}節で決まる。単元の $x$ で $y$ が非単元なら
既知の正付値領域または極の定理による。単元同士で $\bar x=\pm1$ の場合も既に認証済みである。

残る $(X,Y)$ は $X\in\mathbb F_9\setminus\mathbb F_3$, $Y\in\mathbb F_9^\times$ であり、
同時Frobeniusで24軌道である。\path{computations/verify_rank_fifteen_unit_quartic_towers.py} は
各軌道で $c^2=X$, $(v^3-v)^2=Y/X^3$ を解き、倍率の方程式と二座標を検査する。
平方根の後のArtin--Schreier方程式は次数3以下で分解するため、$c,v$ は
$\mathbb F_{3^{12}}$ に取れ、残る平方根は $\mathbb F_{3^{24}}$ に取れる。
各表示で未切断経路を五次まで、文字別四次中心を計算し、四本の中心式を課して
全四始点で十三次まで標準商を確認する。単元座標でのモニックな持上げと
全次数の比較は定理\ref{thm:rank-fifteen-unit-residues}の証明と同じである。
二次数体の合成体では分解群の元の位数は高々2であり、その巡回な剰余体Galois群も位数2以下となる。
\end{proof}

% !TeX root = ../papers/nc-p3.tex
\section{種数2の四窓表示と全ての有理縦切片}
\label{sec:rank-fifteen-genus-two-unit-line}

\begin{proposition}\label{prop:rank-fifteen-diagonal-genus-two}
標数3の代数閉体 $k$ 上で $X\ne0,1$ を固定し、前節の $b=1/c$, $D=v^3-v$,
$d=b^3-b$ を使う。幾何学的一般点で四窓は種数2の体
$k(v,\sqrt{D^2-d^2})$ に係数を持つKlein斉次な表示と添字保持同型である。
この表示は $D(D^2-d^2)\ne0$ の有限点で使える。
\end{proposition}
\begin{proof}
前節の四位置のWalsh変換を使い、対角倍率はまだ掛けない。
$a=1$, $b_0=(K-1)/2$ とし、六根の組を
$(p_0,m_0)=(L+M,L-M)$, $(p_1,m_1)=(B-A,B+A)$,
$(p_2,m_2)=(D_0-C,D_0+C)$ とする。ここで $D_0$ は正規形の最後の倍率である。
共通定数を除く非零成分は36個で、各成分 $u_w$ は符号を除き、
$a,b_0,a+b_0$ と六根のうちの一つである。各語 $w$ の四位置の矢を示す
$0,1$ 行列を $M$ とし、辞書順に最初の独立な13行を $I$ とする。
$M_I$ の最初の独立な13列の小行列式は $-4$ で、$M=N M_I$ の係数 $N$ は全て整数である。
従って代数閉包上の対角基底変更で $I$ の13係数を同時に1とし、他の係数を
$u_w\prod_{r\in I}u_r^{-N_{wr}}$ とできる。実際、対応するトーラス準同型は
階数13で全射である。この操作は四位置で周期的である。

\path{computations/verify_rank_fifteen_genus_two_unit_line.py} はこの行列を四窓の式から生成する。
規格化した各単項式では六根の指数が全て偶数か全て奇数になる。
例えば語 $0231$ の係数は $-m_0m_1p_2/(m_2p_0p_1)$ である。
従って必要なのは六根の二乗とその積 $W$ のみである。
前節の三次式の積より $W^2$ の平方類は $D^2-d^2$ である。
これは相異なる六単根を持つので、二重被覆の種数は2である。
除算に現れる元は六根および $a,b_0,a+b_0$ に限る。
後二者の零では $K=\pm1$、従って $D=0$ であり、主張した開集合で変換は可逆である。
\end{proof}

\begin{theorem}\label{thm:rank-fifteen-all-negative-unit-residues}
剰余標数3の付値で $\bar x=-1$ なら、任意の非零の $y$ で通常な標数零の階数15窓は
標準Koszul左右AS正則で、各始点は四中心四次式上自由階数256である。
\end{theorem}
\begin{proof}
正付値・負付値の $y$ は既知なので単元を扱う。$i^2=-1$, $K=1/v$,
$H=i(K-v)$ とすると $Y=\bar y=-D^2$ である。上の六根の二乗は順に
\[
 (K^2-1\pm2H)/4,\qquad (K+1\pm H)/2,\qquad (K-1\pm H)/2 .
\]
$Z^2=D^2+1$ と置けば $W=Z D^2/v^6$ となり、二乗して六式の積と照合できる。
従って $\mathbb F_3(v)$ 上 $1,i,Z,iZ$ の四座標で演算できる。
符号に応じた $Z$ の選択は全て含まれる。

上のコードで五次まで経路基底を作り、変換後の文字別の四次中心次元が
$(5,0,0,0)$ になることを関数体上で求める。文字0の自由列の昇順の五本のうち
後の四本を課す。全始点で十三次まで中心商の次元は $(1+t+t^2+t^3)^4$ の係数に一致する。
各逆元の分子の $\mathbb F_3(v)$ へのノルムの零点を記録する。
入力・経路・中心を含めた各始点の平方自由な除外多項式の次数は
$1049,2866,1642,3847$、最小公倍式 $F$ の次数は8801である。

$Y=0,1$ を除くと各水準の六表示は $v\mapsto\pm v+j$ ($j\in\mathbb F_3$) の一軌道である。
$F$ からこの二水準の因子を除き、六変換による引き戻しの最大公約元を取ると次数156となる。
各既約因子で $Y$ の最小多項式 $g$ を求め、同じ $g$ の因子の積が
$g(-D^2)$ のモニック化に等しいことを確認する。全表示が除かれる八軌道の次数は
$2,2,2,3,3,4,4,6$ である。四次までと $Y=1$ は既知で、残る六次式は
\[
 Y^6+2Y^5+Y^4+2Y^3+Y+1 .
\]
この根に対し $\mathbb F_{3^{12}}$ で元の六根の表示を取り、四始点で長さ13まで直接検算する。
Frobenius共役でその六根を覆う。超越的な $Y$ はどの除外多項式にも当たらない。

選んだ剰余表示は定理\ref{thm:rank-fifteen-unit-residues}のモニックな方程式で持ち上がる。
上の基底変更の可逆行列も単元行列式のまま持ち上げる。
経路の一般下界と四次中心の全次元の下界5により、剰余の四つの中心式が持ち上がる。
十三次消滅を含む中心商の上界と既出の自由性判定から全次数の結論が従う。
\end{proof}

\begin{corollary}
標数零の代数閉体で、任意の有理数 $x$ と任意の非零の幾何水準 $y$ において、
通常な階数15窓は全て上の正則性と自由性を持つ。
さらに3の上のある付値で $\bar a=-1$ となる縦切片 $x=a$ では全非零水準の
添字保持同型類は系\ref{cor:rank-fifteen-unit-vertical-conics}の $T$ の式で一意に表される。
$a^3/27,(1-a)^2/4$ は各二類、その他は各四類で、最小周期は全て2である。
\end{corollary}
\begin{proof}
有理数 $x$ は3で負付値を持つか、剰余が $0,1,-1$ である。付値を係数体に延長し、
既知の三領域と本節の定理を適用する。$\bar a=-1$ では $a=3,3/4$ および
$4a^2-12a+7=0$ の二根はいずれも生じない。従って通常表示が全点を覆い、
二分岐水準は異なる。円錐曲線の表示と最小周期は前掲の系と同じ計算で得られる。
\end{proof}

% !TeX root = ../papers/nc-p3.tex
\section{二次剰余の全縦切片と二変数単元表示の有理化}
\label{sec:rank-fifteen-quadratic-verticals}
\begin{theorem}\label{thm:rank-fifteen-all-quadratic-verticals}
剰余標数3の付値で $\bar x\in\mathbb F_9$ なら、全ての非零の幾何水準 $y$ について、
通常な標数零の階数15窓は標準Koszul左右AS正則で、各始点が四つの中心四次式上自由階数256である。
特に3の上に剰余次数2以下の素点を持つ数体に $x$ が属すれば、任意の非零の $y$ にこの結論が成り立つ。
\end{theorem}
\begin{proof}
負付値の $x$ と素体に落ちる $x$ は既知である。残る六つの $X=\bar x$ は
$i^2=-1$ として $X=s+i$ ($s=0,1,-1$) とそのFrobenius共役である。
正付値の $y$ は単元横座標の内側の定理、負付値は極の定理で覆われる。
以下 $Y=\bar y$ は非零である。$c^2=X$, $b=1/c$, $d=b^3-b$, $D=v^3-v$ とし、
$Z^2=D^2-d^2$ と置く。前節の六根の積には
$W=(-ic)^3ZD^2/v^6$ が使える。実際、前節の六平方の積がこの二乗と一致する。
$Y=X^3D^2$ とし、素体上の写像で例外を扱うため $B=-D^2$ を用いる。
$B\in\mathbb F_9$ なら $Y\in\mathbb F_9$ なので単元同士の全二次剰余の定理を適用できる。

\path{computations/verify_rank_fifteen_quadratic_verticals.py} は36成分の表示で、
三つの $X$ とも前節と同じ未切断経路、文字0に五次元の四次中心、同じ自由列の後四本の切断を計算する。
四始点の次元は十三次の消滅まで標準商に一致する。
下表の四数は、入力と中心を含めた始点別の除外多項式の次数である。
\[
\begin{array}{c|c|r|r}
 X & \text{四始点} & \deg\mathrm{lcm} & \text{六表示で安定な部分の次数}\\ \hline
 i &(1205,1997,1898,2147)&6604&126\\
 1+i &(3240,3964,3281,2317)&12491&54\\
 -1+i &(3321,3604,3290,2843)&12756&48
\end{array}
\]
第一行では $c=1-i$, $d^2=-1$ なので前節の種数2の関数体を再利用し、
$B=0,1$ を除く。第二、三行では $c\in\mathbb F_{81}$ を使い、
各逆元のノルムを $\mathbb F_{81}(v)$ の多項式として記録する。
$B\in\mathbb F_9$ を除いた後、六変換 $v\mapsto\pm v+j$ の引き戻しの最大公約元を取る。
これと係数の三乗による四つの共役全体の最小公倍式は、素体上で次数72,48である。
$B$ の最小多項式ごとに完全ファイバーが得られることは前節と同じ引き戻しで照合する。

第一行で残る閉点の次数は $2,2,3,4,4,6$、第二、三行で素体上に拡大した除外集合は
それぞれ四次の三軌道と二軌道である。各素体多項式の根を九乗軌道に分け、
$Y=-X^3B$ ごとに元の七つの倍率を復元する。必要な体の素体上の次数は
第一行で $8,12$、第二行で $8,16$、第三行で8で足りる。各九乗軌道で四始点の十三次消滅までを
検査する。三乗は $X$ の他方の共役の全水準も与える。
安定除外集合の外ではいずれかの $v$ の全主元が単元になる。
四中心の持上げと自由加群比較は前節による。最後に $x$ が与えられた数体で負付値なら
第\ref{sec:rank-fifteen-negative-x-fan}節を、整なら上の結果を係数体への延長付値で使う。
\end{proof}

\begin{proposition}\label{prop:rank-fifteen-rational-unit-surface}
$X$ も動かすとき、$X\ne0,1$ かつ $Y\ne0,(1-X)^2$ の種数2表示は、
二変数の有理関数体による純非分離次数3の表示を持つ。
\end{proposition}
\begin{proof}
$S=v-b$, $T=v+b$ とすると、$F(z)=z^3-z$ の加法性から
$Z^2=F(S)F(T)$ である。この曲面は楕円曲線 $u^2=F(S)$ とそのコピーの積を
同時符号反転で割ったものに双有理である。$w=Z/F(S)$ の三乗根を $u$ とし、
$q=T-u^2S$ と取る。$F(T)=u^6F(S)$ を解けば
\[
 S=(q-q^3)/(u^6-u^2),\qquad T=q+u^2 S,\qquad Z=u^3 F(S)
\]
であり、$b=(T-S)/2$, $v=(S+T)/2$、従って四窓の全係数も戻る。
$F(S)F(T)=0$ は分岐水準である。$u\ne0$ で $u^4=1$ なら
$F(T)=\pm F(S)$ となり、$d=0$ または $D=0$ であって仮定に反する。
従ってこれらの式は指定した全幾何点で定義できる。
$S=0$ の上の一般分岐点で $F(S)$ の位数は2、$Z$ の位数は1である。
$w$ の位数は $-1$ で三乗ではないので、得た関数体の拡大次数はちょうど3である。
残る分岐水準では $F(S)=0$ または $F(T)=0$ となり、別の成分の規格化が必要である。
\end{proof}

% !TeX root = ../papers/nc-p3.tex
\section{素体の二つの放物線に沿う全剰余近傍}
\label{sec:rank-fifteen-prime-parabolic-pencils}
\begin{theorem}\label{thm:rank-fifteen-prime-parabolic-pencils}
剰余標数3の付値で $X=\bar x\ne0,1$ とし、$Y=\bar y$ は非零とする。
$Y/(1-X)^2\in\{1,-1\}$ なら全ての通常な標数零の階数15窓が
標準Koszul左右AS正則で、各始点が四中心四次式上自由階数256である。
$X$ の剰余次数に制限はない。
\end{theorem}
\begin{proof}
$Y=(1-X)^2$ では $c^2=X$, $v=1/c$ を取る。$K=c$, $H=K^2-1$ となり、
六根のうち $p_0^2=(K^2-1+2H)/4=0$ である。他の五根は $c\ne0,\pm1$ により非零である。
第\ref{sec:rank-fifteen-genus-two-unit-line}節の36成分から $p_0$ の四成分を除き、
残る32語に同じ対角変換の規格化を施す。
\path{computations/verify_rank_fifteen_prime_parabolic_pencils.py} と共通の規格化コードは、
指数行列の独立な13行の小行列式が絶対値4で、全ての規格化係数で五根の指数が偶数であることを確かめる。
従って $\mathbb F_3(c)$ で四窓を扱える。変換後の四次中心は文字0に五次元あり、
後ろの四自由列を選んだ四中心商が四始点で十三次まで標準次元に一致する。
始点別の除外多項式の次数は $7,3,5,5$、最小公倍式は
\[
 c(c-1)(c+1)(c^2+1)(c^2+c+2)
\]
である。$c,-c$ の両方が零点となるのは最初の四因子に限られる。
仮定のもとで残る $X=-1$ は前に認証した。

一般に $R=Y/(1-X)^2$ を固定し、$R\ne0,1$ とする。
$r^2=R$, $a^3=r$, $b=1/c$, $d=b^3-b$ とし、$v^3-v=r d$ の符号を選ぶ。
$q=v-ab$ と置くと
\[
 b=(q-q^3)/(r-a),\qquad v=q+ab,\qquad Z=\sqrt{R-1}\,d .
\]
$r-a\ne0$ であるから、この一変数表示は全ての $X\ne0,1$ に使える。
$Z^2=(v^3-v)^2-d^2$ であり、36成分の正規化に現れる六根の積も $q$ の有理関数になる。

$R=-1$ では $r=i$, $a=-i$, $Z=d$ とできる。$1/X=b^2=-(q^3-q)^2$ なので、
各有限非零水準の六表示は $q\mapsto\pm q+j$ ($j\in\mathbb F_3$) で尽くされる。
$\mathbb F_9(q)$ で同じ中心と四中心商を計算し、逆元を素体へのノルムで記録すると
四始点の除外次数は $1174,1285,1375,1074$、最小公倍式の次数は4438である。
$1/X=0,1$ を除き六変換で安定な部分は次数126で、完全ファイバーとなる
$1/X$ の閉点の次数は $1,2,2,2,3,3,4,4$ となる。
二次までは前節で認証済みである。二つの三次と二つの四次では各々
$\mathbb F_{3^{12}},\mathbb F_{3^8}$ で元の倍率を取り、四始点の中心商を十三次まで検算する。
$Y=-(1-X)^2$ は素体上の式なので、各多項式につき一根でFrobenius軌道が尽くされる。
両水準で単元主元による持上げと全次数判定は第\ref{sec:rank-fifteen-genus-two-unit-line}節による。
\end{proof}

\begin{corollary}
標数零の代数閉体で $\lambda\in\mathbb Q^\times$ とする。
放物線 $y=\lambda(1-x)^2$ の非零水準の通常窓は全て上の正則性と自由性を持つ。
\end{corollary}
\begin{proof}
3進付値を係数体に延長する。負付値の $x$ と $\bar x=0,1$ は全水準で認証済みである。
残る $x$ では $1-x$ が単元である。$\lambda$ が非単元なら $y$ も非単元で、
既知の内側または極の定理を使う。単元なら $\bar\lambda=\pm1$ なので本定理を使う。
\end{proof}

% !TeX root = ../papers/nc-p3.tex
\section{二次剰余の全放物線と経路行列の有理関数体への降下}
\label{sec:rank-fifteen-quadratic-parabolic-pencils}
\begin{theorem}
剰余標数3の付値で $X=\bar x\ne0,1$, $Y=\bar y\ne0$ とする。
$R=Y/(1-X)^2\in\mathbb F_9^\times$ なら、全ての通常な標数零の階数15窓は標準Koszul左右AS正則で、
各始点が四中心四次式上自由階数256である。
\end{theorem}
\begin{proof}
残る倍率の三乗共役の代表は $R=s+i$, $s=0,1,-1$, $i^2=-1$ である。
$r^2=R$, $a^3=r$, $\delta^2=R-1$ の全定数は $\mathbb F_{81}$ に取れる。
前節の $b=(q-q^3)/(r-a)$, $v=q+ab$, $Z=\delta(b^3-b)$ の表示で、
$B=-(q^3-q)^2=\kappa/X$, $\kappa=-(r-a)^2\in\mathbb F_9^\times$ である。
実際 $r^9=\pm r$ で $a^9$ にも同じ符号が付く。従って $B\in\mathbb F_9$ の有限非零水準は既知である。

\path{computations/verify_rank_fifteen_quadratic_parabolic_pencils.py} は
$\mathbb F_{81}(q)$ 上で五次までの経路と四次中心を計算する。
文字0の五本の自由列のうち第1,2,4,5列の中心を課すと、各始点の商が十三次まで標準次元を持つ。
下表で第3列は始点別の除外多項式の最小公倍式の次数、第4列は
$B\in\mathbb F_9$ を除いて $q\mapsto\pm q+j$ の六変換で最大公約元を取った次数、
第5列はその係数の三乗共役全体の和集合の次数である。
\[
\begin{array}{c|c|r|r|r|c}
R&\text{四始点の除外次数}&\text{最小公倍式}&\text{六変換}&\text{共役和}&\deg_B\\\hline
i &(1284,766,971,1142)&3831&66&84&4,4,6\\
1+i &(1476,743,1755,1723)&5416&60&96&4,4,8\\
-1+i &(1426,1136,1640,1844)&5768&60&72&4,4,4
\end{array}
\]
最終列は共役和を完全に覆う $B$ の素体上の既約因子の次数である。
各因子の根を九乗軌道に分け、$X=\kappa/B$, $Y=R(1-X)^2$ で元の七倍率を取り直す。
全軌道で四始点の十三次消滅までを直接確認する。最大の八次因子には
$\mathbb F_{3^{32}}$ を使う。六表示による被覆、係数の共役と標数零への持上げは
第\ref{sec:rank-fifteen-quadratic-verticals}節の証明と同じである。
\end{proof}

\begin{corollary}
3上に剰余次数2以下の素点を持つ数体 $L$ に対し、任意の $\lambda\in L^\times$ で
$y=\lambda(1-x)^2$ 上の非零幾何水準の通常窓は全て上の正則性と自由性を持つ。
さらに、通常な未認証点が残るなら $x,\lambda=y/(1-x)^2$ はともに代数的数であり、
各々のモニック最小多項式は3整で、定数項は3進単元、3での還元は $T^9-T$ と互いに素である。
\end{corollary}
\begin{proof}
前節の有理倍率の系と同じく、$\lambda$ が単元のときだけ今回の定理が必要である。
もし $x$ が超越的なら $3x$ に超越的剰余を与えるGauss付値を取り $\nu(x)<0$ とできる。
$\lambda$ が超越的なら同じく $\nu(\lambda)<0$ とでき、$\bar x\ne0,1$ の単元域で
$y$ は負付値となる。いずれも係数体への延長で全点を覆う。
代数的な $x,\lambda$ のいずれかが3上のある付値で非単元でも既知領域に入る。
全共役が単元ならモニック最小多項式は3整で定数項は単元である。
その還元に $\mathbb F_9$ の根があれば、代数閉包でモニック因子を一次式まで分解して還元すると、
その剰余を持つ共役がある。対応する付値で二次剰余の縦切片または本節の定理を使える。
\end{proof}

\begin{proposition}\label{prop:rank-fifteen-word-parity-descent}
36成分の正規化で、$F$ を六根の二乗と基底の三係数を含む体、
$E=F(W)$ を六根の積の添加とする。$W\ne0$ とする。任意の長さと始点について
文字別経路の関係行列、文字0の四次中心方程式、およびそこから選ぶ四中心商の関係行列は
語の基底変更と行のスカラー倍で $F$ 上に書ける。
特に二変数の通常表示では平方根の積を添加せず $F=k(K,H)$ でこれらの階数を計算できる。
\end{proposition}
\begin{proof}
$\omega(g,h)=g_L h_H\in\{0,1\}$ とし、位置 $p$ からの語 $w$ に
\[
 \ell_p(w)=\sum_j[\,p+j\equiv2\pmod4,\ w_j=3\,]+\sum_{j<k}\omega(w_j,w_k)\pmod2
\]
を付ける。添字 $j$ は0から始める。コードは全36係数の六根の共通指数の偶奇が
対応する $\ell_0(w)$ に等しいことを四つの巡回表示でも直接確認する。
連結 $UV$ の偶奇は両語の偶奇の和に $\omega(|U|,|V|)$ を加えたもので、$|U|$ はKlein次数である。

語基底を $W^{\ell_p(w)}w$ に変える。二次関係の固定した行は四語テンソルの残りの二文字が固定で、
変化する二文字のKlein次数は一定である。前後に固定した語を付けた任意の関係行でも、
各係数の $W$ 指数の偶奇はこの基底で行全体に共通になる。行をその偶奇の $W$ 冪で割ればよい。
従って各経路空間はこの基底で $F$ 上の商として計算できる。
四次の文字0の中心候補を同じ規格化語の線型結合とすると、矢を左右から掛けたときの
交差項 $\omega(0,h),\omega(h,0)$ はともに零である。
四位置周期性により矢の寄与も一致するので、交換子の行列も同じ操作で $F$ に降りる。
$F$ 上で解いた核は $E$ への拡大で全核を与える。得た文字0の切断の前後に語を
付けた行にも同じ議論が適用できる。$W^2\in F$ であることから結論が従う。
この議論は各長さでの線型代数を有理化する。積を再帰的に計算する場合、
次数 $g,h$ の積から $W^{\omega(g,h)}$ を除くと、括弧を変えたときの倍率は
\[
 (W^2)^{(\omega(h,k)+\omega(g,h+k)-\omega(g,h)-\omega(g+h,k))/2}.
\]
分子は偶数で、一つの次数が0なら零である。
\end{proof}

% !TeX root = ../papers/nc-p3.tex
\section{二変数の非分岐単元域全体の経路と四次中心}
\label{sec:rank-fifteen-bivariate-word-paths}
\begin{theorem}\label{thm:rank-fifteen-bivariate-word-paths}
標数3で $K,H$ を取り、
\[
 KH(K^2-1)\big((K^2-1)^2-4H^2\big)
 \big((K+1)^2-H^2\big)\big((K-1)^2-H^2\big)\ne0
\]
とする。36成分に規格化した窓の全始点で、長さ1から5までのKlein文字別経路次元は
\[
 (1,1,1,1),\quad (2,2,2,4),\quad (5,5,5,5),\quad
 (11,8,8,8),\quad(14,14,14,14)
\]
である。四次中心の文字別次元は $(5,0,0,0)$ であり、文字0の中心の同じ自由列が
開集合全体で基底を与える。この範囲には任意の幾何単元 $X\ne1$ と
$Y\ne0,(1-X)^2$ の表示が含まれる。
\end{theorem}
\begin{proof}
$F=\mathbb F_3(K,H)$ の二変数有理関数として独立に計算する。
$P=W^2$ を六根の二乗の積とし、前節の偶奇式を $\ell_p$ と書く。
\path{computations/verify_rank_fifteen_bivariate_word_paths.py} は全ての係数を
$\mathbb F_3[K,H]$ の互いに素な分子とモニック分母で持つ。

規格化した語 $W^{\ell_p(w)}w$ を、各長さ、始点、文字の商の基底に射影する。
先頭の矢を $a$ とし、末尾の語 $t$ の射影に現れる基底語を $s$ とする。
既存の $s$ の係数に掛ける補正は
\[
 P^{(\ell_p(at)-\ell_{p+1}(t)+\ell_{p+1}(s)-\ell_p(as))/2}
\]
である。$s,t$ の文字が等しいので分子は偶数である。
二次関係に固定した末尾を付ける際も、各項の $W$ 指数から行の共通偶奇を引いて
$P$ の整数冪を使う。これにより全関係行を構成し、主元行列を後退代入して
次の長さの語基底と射影を求める。各始点で表の次元が得られる。

文字0の四次中心方程式では未知数が $4\cdot11=44$ 個である。
各語の左右に矢を掛け、上の五次への射影で比較する。行共通の $W$ 冪を除くと
$F$ 上で階数39となる。ここまでに逆にした主元の分子の既約因子は、
$H,K-1,K+1$ と六根の二乗のモニック化の九つだけである。
他の三文字では交換子の左右で偶奇が違う場合がある。
各未知数を $u+Wv$ として $1,W$ 成分に分けると、各文字で $F$ 上64列となる。
上の九因子と $K$ に関して単元となる主元を優先して選べば、三つの行列とも階数64である。
従ってどの段階にも定理の開集合の内部で消える分母はない。
$W^2=P$ が特殊化で分解しても、$P\ne0$ なので二つの符号を合わせた行列が全列階数なら
両符号の核は零である。元の関係と全ての行操作が特殊化できるので、零行も保存される。

幾何単元表示で $K=1/v$, $H=c(K-v)$ とすると、$Y\ne0$ から $v\ne0,\pm1$、
従って $K,H,K\pm1$ は非零である。残る因子は六根の消失に対応し、これは分岐水準に限られる。
\end{proof}

% !TeX root = ../papers/nc-p3.tex
\section{単元曲面全体での五つの中心の独立性と六次中心商}
\label{sec:rank-fifteen-bivariate-sixth-cut}
\begin{theorem}
前節の開集合の各表示において、任意の独立な四つの四次中心を課した商は、
四始点の四次と五次でそれぞれ31,40次元である。
さらに任意の幾何単元対 $X\ne1$, $Y\ne0$ に対し、六次が全始点で44次元となる
四中心と表示が取れる。$Y\ne(1-X)^2$ かつ $X^2+X+2\ne0$ なら、前節の
文字0の五自由列のうち第1,2,4,5列を、同じ表示の四始点で使える。
\end{theorem}
\begin{proof}
$F=\mathbb F_3(K,H)$ とし、前節の五本の四次中心を同じ規格化語基底で書く。
\path{computations/verify_rank_fifteen_bivariate_sixth_cut.py} は、各始点でこれらを四次経路へ
評価する行列、および各矢を右に掛けて五次経路へ評価する行列を作る。
全部で $4(1+4)=20$ 個あり、五行の階数を調べればよい。
先に開集合の十因子しか持たない主元で消去する。残る行列の最大小行列式の分子から十因子を除き、
そのイデアルに1が入るかをGröbner基底で検査する。必要なときは開集合の積 $C$ に対して
$1-zC$ も加えて局所化する。どの行列も全点で階数5である。
従って任意の四次元部分空間が各始点で四本の独立な関係を与える。
五次では四つの矢の積が別々の文字にあり、各文字で四本が独立である。
中心性により左からの積もこの関係に含まれ、$35-4=31$, $56-16=40$ を得る。

次に定理の四本を固定し、先頭語から一矢を末尾に足す形で中心商の語基底を構成する。
二次関係と四次での中心切断に、前節と同じ $P=W^2$ の偶数指数補正を使う。
例えば末尾の二次関係の係数を $W^e f$、その語を $ab$、前置する語を $u$ とする。
$ua$ の射影の基底語を $s$ としたとき、列 $sb$ の $W$ 指数は
$e-\ell_p(ua)+\ell_p(s)-\ell_p(sb)$ である。行の共通偶奇を除いて $F$ 上で消去する。
四始点の六次文字別次元はいずれも $(12,12,12,8)$ となる。

主元が消える軌道を二変数で正確に制限する。$K=1/v$, $H=c(1-v^2)/v$,
$X=c^2$, $B=(v^3-v)^2$, $Y=X^3B$ と置く。四始点の除外因子を引き戻し、
もとの開集合ですでに単元な因子を除いて最小公倍式 $G(v,c)$ を取る。
五次まで、六次までの $(v,c)$ 次数はそれぞれ $(7,7),(55,44)$ である。
$\mathbb F_3[c,v]$ は $\mathbb F_3[X,B]$ 上で $c^i v^j$ ($i<2$, $j<6$) を基底に持つ。
実際 $c^2=X$, $v^6=B+2v^4-v^2$ の二つのモニックな式で順に割れる。
$XB\ne0$ では十二表示は相異なる。従って全表示が $G=0$ となる必要十分条件は
この基底での全係数が同時に零となることである。

非零係数の対から $B$ を消去した終結式の共通因子で、必要な $X$ の条件を得る。
検証コードで選ぶ対のモニック最大公約元は、五次まででは1、六次まででは
\[
                     X^{136}(X-1)^9(X^2+X+2)
\]
である。これは二変数曲面上の孤立した不良点も検出する必要条件である。
仮定の下で残る二次因子の根は第\ref{sec:rank-fifteen-quadratic-verticals}節で
各水準に十三次までの証明書がある。分岐 $Y=(1-X)^2$ も
第\ref{sec:rank-fifteen-prime-parabolic-pencils}節で全水準を認証した。以上で六次までの主張が従う。
\end{proof}

% !TeX root = ../papers/nc-p3.tex
\section{六次の例外ファイバーと九次で必要になる中心の選び直し}
\label{sec:rank-fifteen-fixed-center-obstruction}
\begin{proposition}
標数3の非分岐幾何単元対では、前節の第1,2,4,5列が四始点とも
六次まで所定の次元を与える表示を常に選べる。
しかしこの同じ四列では、非分岐域のある表示で九次の商の次元が
全始点で23となる。その表示を保ったまま第4列を第3列に替えると、
四始点とも次数0から13までの商の次元は
\[
1,4,10,20,31,40,44,40,31,20,10,4,1,0
\]
となる。また同じ表示に、九次まで所定の次元を持ち、十次で初めて増大する別の四中心もある。
\end{proposition}
\begin{proof}
前節の六次の終結式で残る横座標を $X^2+X+2=0$ とする。
$G$ の十二基底係数を $\mathbb F_9[B]$ に特殊化したモニック最大公約元は
$B-(X+1)$ である。前節の検証コードは同じ $G$ からこの多項式を計算する。
従って指定四列の一般主元が全十二表示で同時に消えうるのはこの共役な点対だけである。
ここでいう消失は、四始点の主元の積の消失を表す。

$a^4+2a^3+2=0$ とした $\mathbb F_{81}$ で、次の一表示を取る。
\[
 X=a^3+a^2+2,\qquad c=a^2+2,\qquad v=a^3+a^2+a.
\]
実際に $c^2=X$, $(v^3-v)^2=X+1$ であり、$K=1/v$, $H=c(K-v)$
で六根の二乗の積は非零である。共役点はFrobeniusで得る。
\path{computations/verify_rank_fifteen_fixed_center_obstruction.py} は普遍的な
五自由列の語係数をこの $K,H$ に特殊化し、体上で主元を選び直して
前節の末尾語の再帰を実行する。第1,2,4,5列では四始点とも八次まで上の次元に一致する。
九次の文字別次元は、始点0,2で $(6,6,6,5)$、始点1,3で $(5,6,6,6)$ である。
これで残っていた点対の六次と九次の増大が同時に分かる。
第1,2,3,5列で同じ再帰を十三次まで行うと表示した列になる。
さらに五列を $z_1,\ldots,z_5$ と書き、$z_1+az_3,z_2+a^2z_3,z_4+a^3z_3,z_5+a^4z_3$ を取る。
この四中心は独立である。全始点で九次まで期待値に一致し、十次から十三次の各次元は16となる。
特に九次までの一致だけでは、選んだ四中心が所定の商を与えるかは決まらない。
\end{proof}

% !TeX root = ../papers/nc-p3.tex
\section{全幾何単元対での七次四中心商と十三個の残余軌道}
\label{sec:rank-fifteen-bivariate-seventh-cut}
\begin{theorem}
標数3の任意の幾何単元対 $X\ne1$, $Y\ne0$ において、四始点すべてで
四次から七次の次元が順に $31,40,44,40$ となる共通の四次中心四本と表示が存在する。
下表の十三個のFrobenius軌道では、四中心商を十三次まで認証できる。
従ってこれらの剰余を持つ通常な標数零の階数15窓は標準Koszul左右AS正則で、
各始点は四つの中心四次式上で自由階数256である。
\end{theorem}
\begin{proof}
分岐水準と $X\in\mathbb F_9$ は既知である。
残りでは第\ref{sec:rank-fifteen-bivariate-sixth-cut}節の四列を使い、
語を右に延ばして七次まで進む。各始点で七次の文字別次元は $(10,10,10,10)$ となる。
四始点の主元分子の除外因子の積を同節と同じ最小公倍式として引き戻すと、
$G(v,c)$ の次数は $(328,280)$ である。$X=c^2$, $B=(v^3-v)^2=Y/X^3$ 上の
十二基底係数のモニック最大公約元は $D=X^3B-X+1$ となる。

まず $D\ne0$ とする。十二係数を $D$ で割り、開集合で単元となる
$X,X-1,B,X^3B-(1-X)^2$ の因子も除く。非零係数を $c$ の指数、次いで $v$ の指数の
昇順に並べ、最初の六係数を順に二本ずつ組にして $B$ を消去する。三終結式から
$X,X-1$ 因子を除いた最大公約元の次数は70である。その各既約因子の根で十二係数を
再度 $B$ の多項式として比較する。$[\mathbb F_3(X):\mathbb F_3]>2$ で
共通根を持つのは下表の十一組であり、全て $B\in\mathbb F_3(X)$ である。

次に $D=0$ とする。始点0,2の消去はそのまま使い、始点1,3で別の主元を選ぶ。
各有理関数行をスカラー倍の原始多項式行に直し、台が最小の行のうち原始係数の次数が
最小の位置を取る。同順位なら列順、行順を使う。行の内容と分母をスカラーで保持し、
記録する主元にはこのスカラーを戻す。従って有理関数行列の通常の消去をそのまま
特殊化できる領域が得られる。二通り目の除外積は次数 $(391,347)$ である。
十二係数に $B=(X-1)/X^3$ を代入し、各係数の分母を払った最大公約元は
\[
 X^{175}(X-1)^{39}(X+1)(X^2+X+2)(X^3+2X+2)(X^3+X^2+X+2)^2 .
\]
次数が二を超える二因子は下表の最後の列に対応する。
\[
\begin{array}{c|l|l}
 X\text{ の最小多項式}&B\ (D\ne0\text{ 側の候補})&B\ (D=0)\\ \hline
 X^3+2X+2 & X^2+X+1 & X^2+2X+1\\
 X^3+X^2+2 & 1,\ X^2+2X+2 & \\
 X^3+X^2+X+2 & 1,\ X & 2X^2+2X+1\\
 X^3+2X^2+2X+2 & 2X+2,\ X^2+2X+2,\ 2X^2+2X+1 & \\
 X^4+X^2+2X+1 & X^3+2X+1 & \\
 X^4+X^3+2X+1 & 2X^3+X^2+X & \\
 X^6+X^2+2X+1 & X^5+X^3+X+1 &
\end{array}
\]
\path{computations/verify_rank_fifteen_bivariate_seventh_cut.py} はこの各組を
有限体上の一つの表示で確かめ、Frobeniusで共役な全点を覆う。
実際に平方根を取り、元の四窓から五次までの経路、四次中心、四中心商の
十三次までの次元を計算する。必要な拡大次数はそれぞれ横座標の次数の二倍で足りる。
よって主元が全表示で失敗しうる残余にも同じ始点別の七次次元が存在する。
表の点では第\ref{sec:rank-fifteen-quadratic-verticals}節の四中心の持上げと
自由加群比較を適用できる。
\end{proof}

% !TeX root = ../papers/nc-p3.tex
\section{標数零の全通常階数15窓の七次までの中心上の自由性}
\label{sec:rank-fifteen-truncated-central-freeness}
\begin{theorem}
標数零の通常な階数15窓のすべてで、四始点に共通な四つの四次中心 $z_1,\ldots,z_4$ を選べる。
各始点 $p$ の長さ $0\le n\le7$ の経路空間 $A_{p,n}$ は次元 $h_n=\binom{n+3}{3}$ を持つ。
中心商 $Q=A/(z_1,\ldots,z_4)$ のこれらの次数での次元は
\[
             (q_0,\ldots,q_7)=(1,4,10,20,31,40,44,40)
\]
である。各 $Q_{p,n}$ の任意の基底を経路に持ち上げると、これらに中心単項式を
掛けたものが七次まで各 $A_{p,n}$ の基底となる。特に六次は84次元、七次は120次元である。
\end{theorem}
\begin{proof}
すでに四中心上の自由性が全次数で知られた窓には結論がある。残りは
第\ref{sec:rank-fifteen-quadratic-parabolic-pencils}節の系より、3上の付値で
$X\ne0,1$, $Y\ne0,(1-X)^2$ の単元域に入る。代数閉体上で始める。
前節の四中心を持つ剰余表示を取り、第\ref{sec:rank-fifteen-integral-annulus}節の
モニックな座標方程式でその表示を持ち上げる。五次までの未切断経路は、
一般点の下界と剰余での同じ次元によって単元主元の語基底を持つ。
この基底で四次中心の方程式を作ると剰余で核は五次元であり、
命題\ref{prop:fixed-residual-regularity}で用いた正則一般窓の四次中心から来る核次元の下界も五である。従って主元小行列式が単元のまま
自由列を選び、指定された四つの中心を同時に持ち上げられる。
各有限長さの中心商は語と関係の有限行列の余核だから、分数体で
$\dim Q_{p,n}\le q_n$ ($0\le n\le7$) となる。

形式的な次数4の変数を持つ多項式環を $S=k[T_1,\ldots,T_4]$ とし、$T_i$ を
$z_i$ として作用させる。$S_{4j}$ の次元は $s_j=\binom{j+3}{3}$ である。
任意の中心商の基底の持上げは各始点の経路を $S$ 上で生成する。実際、商における
像を差し引いた後は中心倍の和であり、残る経路の長さは4だけ短い。中心は次数4で
周期の始点に戻り、右に掛けるとこの生成を始点ごとに使える。
定理\ref{thm:rank-fifteen-generic-levels}の正則一般窓への比較による未切断経路の
下界は、各通常窓で $\dim A_{p,n}\ge h_n$ である。一方、
\[
 h_n\le \dim A_{p,n}\le\sum_{0\le4j\le n}s_j\dim Q_{p,n-4j}
 \le\sum_{0\le4j\le n}s_j q_{n-4j}=h_n
 \qquad(0\le n\le7).
\]
最後は $(1+t+t^2+t^3)^4/(1-t^4)^4=(1-t)^{-4}$ の係数比較である。
全て等号なので商の次元も確定し、生成する単項式は線型独立である。
体拡大を使った場合、元の体上の中心空間は核の体拡大と可換である。
四式の選択空間で七次までの商次元の上界は有限個の開条件をなし、非空である。
標数零の体は無限なので元の体でも同時に選べる。
\end{proof}
\begin{corollary}
この四中心を固定すれば、八次から十一次の商には四始点のそれぞれで下界
$31,20,10,4$ がある。これらの次数のいずれかで対応する上界が得られれば、
その次数で商次元、未切断経路次元、中心単項式基底も同時に確定する。
\end{corollary}
\begin{proof}
$n=8,9,10,11$ なら $n-4,n-8$ はすでに定理の範囲にある。
上の生成による不等式から
$\dim Q_{p,n}\ge h_n-4q_{n-4}-10q_{n-8}$ となる。右辺が掲げた四数である。
商次元がこの下界に達すれば、再び生成写像の両辺の次元が等しくなる。
\end{proof}

% !TeX root = ../papers/nc-p3.tex
\section{符号付きの六表示による全通常窓の八次判定}
\label{sec:rank-fifteen-oriented-eighth-cut}
\begin{theorem}
標数零のすべての通常な階数15窓において、四始点に共通な四次中心四本を選び、
八次まで中心単項式基底を取れる。各始点の八次で中心商は31次元、未切断経路は165次元である。
\end{theorem}
\begin{proof}
前節で残った非分岐単元対で、始点を保つために
定理\ref{thm:rank-fifteen-global-invariants}の座標 $(x,p,q)$ を使う。
$U=K^2$, $V=H^2$ と書くと
\[
 p=\frac{H^3K(U-9)}{(U-1)^3},\qquad p^2=x^3-27y,\qquad
 q^2=(1-x)^2-4y.
\]
剰余で $K=1/v$, $H=c(K-v)$, $X=c^2$ とすると $\bar p=c^3$ である。
そこで目標の $p$ に対して符号を含めて $c$ を固定する。$B=Y/X^3$ に対する
六つの $v$ のいずれでも、対応する $U$ と $K,H$ を既出のモニックな座標方程式で
持ち上げられる。この表示の $p'$ は $p'^2=p^2$ かつ $\bar p'=\bar p\ne0$ なので $p'=p$ である。
残る $q$ の符号は $K,H$ を保って選べる。実際 $n^2+(H/n)^2=U-1$ であり、
$n$ と $H/n$ の交換は $q$ の分子の因子 $2n^2+1-U$ の符号を反転する。
$q$ の他の因子はこの交換で変わらない。六根の符号は語の偶奇降下でどちらも含まれる。
従ってこの六表示のどれで得た始点0の証明書も、指定した窓へ\emph{添字を保って}
持ち上げられる。全十二表示を一緒に消去したときは、この最後の結論は直接には得られない。

第\ref{sec:rank-fifteen-bivariate-sixth-cut}節の第1,2,4,5列で始点0の右語を延ばすと、
八次の文字別次元は $(7,8,8,8)$ となる。
\path{computations/verify_rank_fifteen_oriented_eighth_cut.py} は原始多項式行に有理スカラー
$A/D$ を保持してこの消去を行う。原始行の非零係数 $b$ の有理主元としての分子・分母の次数和は
\[
 \deg A+\deg D+\deg b-2\deg\gcd(b,D)
\]
なので、第\ref{sec:rank-fifteen-bivariate-word-paths}節と同順位の主元を選べる。
除外多項式 $G(v,c)$ の次数は $(167,146)$ である。今度は $c$ を固定し、
$v^6=B+2v^4-v^2$ だけを使い、$G$ を $1,v,\ldots,v^5$ に関する六係数に分解する。
$B\ne0$ なので六表示は相異なり、全表示で除外される条件は六係数の同時消失である。
六係数の最大公約元は1である。$c,c-1,c+1,B,c^6B-(1-c^2)^2$ の単元因子を除き、
係数を低次側から二本ずつ組にして $B$ を消去する。三終結式から $c,c\pm1$ の因子を
除いた最大公約元の次数は33である。$c\in\mathbb F_9$ は既知なので、
より高次の既約因子の剰余体で六係数の $B$ に関する最大公約元を求める。
残る $(c,B)$ は六つの三次軌道と二つの四次軌道であり、個々の $B$ は
$\mathbb F_3(c)$ に属する。これらは元の四窓の有限体計算で四始点とも十三次まで確認する。
これでどの指定窓でも、少なくとも始点0において、八次商の上界31を実現できる。

最後に四中心の選択を共通化する。始点0での主張を一位置移動した窓にも適用する。
同じ $(x,y)$ 上の移動・反転は定理\ref{thm:rank-fifteen-rational-quotient}の符号変更で
表され、規格化テンソル上でも作用する。上で両方の $c$ と $q$ の符号を扱ったので、
添字を戻すと元の窓の各始点で少なくとも一つの四中心が取れる。
固定した窓の四次中心空間 $Z_4$ で、ある始点の八次商次元が31以下である条件は
$\operatorname{Gr}(4,Z_4)$ の開条件である。実際、普遍部分空間の局所基底を四式として、
関係行列の小行列式で書ける。各始点に対応する四つの非空開集合と、
前節の七次までの非空開集合はこの既約Grassmann多様体で同時に交わる。
標数零の基礎体上にも点がある。前節の下界31と生成写像の次元比較によって
八次の両次元と基底の主張が従う。
\end{proof}

\begin{corollary}
この四中心では十二次の商にも下界1がある。実際、中心単項式の生成により
$\dim Q_{p,12}\ge\binom{15}{3}-4\cdot31-10\cdot31-20\cdot1=1$ である。
従って九次から十二次までの目標上界 $20,10,4,1$ は、各次数で別々に証明してよい。
各上界の証明書が異なる四中心を用いる場合も、同じGrassmann多様体の開条件を交差させればよい。
\end{corollary}

% !TeX root = ../papers/nc-p3.tex
\section{九次で固定四列の全表示が失敗する剰余点}
\label{sec:rank-fifteen-ninth-orbit-obstruction}
\begin{proposition}
第\ref{sec:rank-fifteen-fixed-center-obstruction}節の $\mathbb F_{81}$ の点で
$X=a^3+a^2+2$, $B=X+1$, $c_0=a^2+2$, $v_0=a^3+a^2+a$ とする。
$c=\pm c_0$ とした全十二表示のいずれでも、第1,2,4,5列による九次商は四始点とも23次元である。
一方、$v=v_0$ で第4列を第3列に替えた四中心では九次が20次元となり、十三次で消える。
\end{proposition}
\begin{proof}
$v$ の全ての選択は $v=\pm v_0+j$ ($j\in\mathbb F_3$) である。
いずれも $(v^3-v)^2=B$ を満たし、$B\ne0$ なので相異なる。
\path{computations/verify_rank_fifteen_ninth_orbit_obstruction.py} は各 $K=1/v$, $H=c(K-v)$ で
同じ普遍的な四列を特殊化し、有限体上で経路の関係行列を初めから消去する。
両方の $c$ の符号で各表示の八次までは所定の次元である。九次の文字別次元は
始点0,2で $(6,6,6,5)$、始点1,3で $(5,6,6,6)$ となる。従ってこのファイバー全体で、表示や消去主元の
選び直しだけで、指定四列の八次証明書を九次へ延長することはできない。
交換後の四列の主張は第\ref{sec:rank-fifteen-fixed-center-obstruction}節の直接計算による。
\end{proof}

% !TeX root = ../papers/nc-p3.tex
\section{三中心商の乗法鉛筆と例外超平面}
\label{sec:three-center-pencil}

\begin{lemma}\label{lem:three-center-pencil}
四周期経路代数で三つの中心四次式を割った商を $C$ とする。
始点 $p$ の長さ $n$ の成分を $C_{p,n}$ と書く。次を仮定する。
\begin{enumerate}
\item $0\le n\le13$ で $\dim C_{p,n}$ は $(1-t^4)^3/(1-t)^4$ の次数 $n$ の係数である。
\item $0\le n<9$ で、$C_{p,n}$ の元のうち右から全ての矢を掛けて零になるものは零だけである。
\end{enumerate}
残る中心四次式の線型結合 $z$ について、$(C/zC)_{p,n}$ の次元が
次数 $0$ から $13$ で $(1+t+t^2+t^3)^4$ の係数に一致することと、
$z_p:C_{p,9}\longrightarrow C_{p,13}$ が同型であることは同値である。
\end{lemma}
\begin{proof}
両空間の次元は64であり、余核が十三次の商である。
核に入る長さ $m<9$ の非零元があれば、第二仮定を反復して右に矢を掛け、
長さ9の非零元を得る。この元も左乗法 $z_p$ の核に入り矛盾する。
従って各 $m\le9$ で単射であり、余核の次元は四つ前の係数との差になる。
逆は十三次の余核の消滅による。
\end{proof}

\begin{proposition}\label{prop:f81-center-pencil}
第\ref{sec:rank-fifteen-fixed-center-obstruction}節の $\mathbb F_{81}$ の表示
$(c,v)=(a^2+2,a^3+a^2+a)$ で、中心 $z_1,z_2,z_5$ を固定する。
第四式を $z_4+s z_3$ とする射影直線全体で、四始点の商が十三次まで所定の次元を持つための
必要十分条件は
\[
       s\notin\{0,-1,-(2a^3+a+2)\}.
\]
これは幾何点について成立し、$s=\infty$ は第四式 $z_3$ を表す。
\end{proposition}
\begin{proof}
\path{computations/verify_rank_fifteen_center_pencil.py} は三式だけで語を簡約する。
四始点のそれぞれで長さ0からの次元は
$1,4,10,20,32,44,54,60,63,64,64,64,64,64$ となる。
各文字次数で四つの右乗法の行列を横に連結し、長さ0から8で行が独立であることを確かめる。
長さ9と13の各文字空間は16次元である。三中心商の同一の語基底において
$z_4,z_3$ の左乗法行列を $M,N$ とすると、全十六ブロックで
\[
 \frac{\det(M+sN)}{\det N}=s^8(s+1)^4(s+2a^3+a+2)^4 .
\]
実際、行列式は高々十六次なので相異なる十七点での有限体消去と補間で全係数が決まり、
最高次係数も非零である。補題\ref{lem:three-center-pencil}を適用すればよい。
三例外で四式を直接課しても、長さ9から13で $s=0$ は $(23,22,22,22,22)$、
他の二値は $(20,16,16,16,16)$ となる。一方、五中心を全て課した商の長さ0から8の
次元は四始点とも $(1,4,10,20,30,36,34,20,0)$ である。従ってこの表示では
元の代数の各始点は五中心上で長さ7以下の155本の経路で生成される。実際、
商基底の持上げを引けば正次数の中心倍が残り、長さに関する帰納法が使える。
\end{proof}

\begin{proposition}\label{prop:five-center-bad-hypersurface}
標数零で五次元の中心四次式空間 $V$ を指定し、折り畳み環 $A$ が $k[V]$ 上有限で
標準Koszul左右AS正則であるとする。四平面 $U\subset V$ を双対射影空間
$\mathbb P(V^*)$ の点で表す。$A/(U)$ が無限次元となる点の集合は、
一つの頂点の隅から得られる次数64の斉次式の零点集合である。
特に有限性を実現する超平面を探す際、任意に固定したアフィン座標で
\[
 (i_1,i_2,i_3,i_4),\qquad i_j\in\mathbb Z_{\ge0},\quad \sum i_j\le64
\]
の有限リストが全てのこのような窓に共通して使える。
\end{proposition}
\begin{proof}
$k[V]$ は換算次数1で生成され、有限性と標準Hilbert関数から次元4である。
一般の四元 $x_1,\ldots,x_4$ を斉次正規化に取り、$z$ で $V$ の基底に延ばす。
$P=k[x_1,\ldots,x_4]$ とする。定理\ref{thm:central-window-cutoff}の必要性の証明により
$A$ は $P$ 自由であり、隅 $D=e_0 A e_0$ の商基底も斉次自由基底に持ち上がる。
経路長 $0,4,8,12$ の基底数は $1,31,31,1$ である。$D$ 上の $z$ の乗法を
この64基底で $M(x)$ と書くと $F(x,u)=\det(uI-M(x))$ は $u$ についてモニックである。
入力と出力の基底次数の差が置換積で相殺するため、換算次数で全次数64の斉次式となる。

$k[X_1,\ldots,X_4,u]$ から $A$ への中心作用を考える。閉点 $(b,d)$ の上での
$D$ のファイバーは $dI-M(b)$ の余核であり、その台は $V(F)$ である。
原点外ではこの台は $A$ の台とも一致する。実際、非零ファイバーで
ある中心四次式が非零スカラーとなる。各長さ4の閉路は頂点0を通るので、
このファイバーでは $1$ が $e_0$ の生成する両側イデアルに入る。
従ってファイバーが非零ならその第0隅も非零である。
超平面 $U$ を割ったときの台はこの台と $U$ の切り出す直線との交わりである
（有限加群に対する中山の補題）。斉次性により、その商が有限次元であることと
直線の射影点で $F$ が消えないことが同値となる。
この記述は線型座標変換で保たれる。

どのアフィンチャートでも $F$ の非斉次化は非零で次数は高々64である。
次数64以下の多項式に対する基底
$\prod_{j=1}^4\binom{t_j}{m_j}$（$\sum m_j\le64$）を取る。
指定した格子での評価行列は全次数順で対角成分1の三角行列である。
よって少なくとも一つの値で $F\ne0$ となる。
\end{proof}

この有限リストの各四平面について、十三次の消滅をまず一始点だけで検査できる。
そこを通過するまで高々三文字なので、その始点で消えれば全始点の長さ16以上が消える。
このとき中心商は有限であり、定理\ref{thm:central-window-cutoff}の二条件でAS性が決まる。
指定した五中心上の有限性が既知なら、全リストで一始点の十三次が非零のときも
命題の対偶と同定理の必要条件により非ASである。五中心上の有限性そのものは、
全五式で割った商の一始点で一つの次数が零になれば従う。

\begin{proposition}\label{prop:five-center-generic-eight}
標数3の二変数単元チャートの始点0で、規格化された固定四列 $(1,2,4,5)$ の
九次商は、全ての幾何点で少なくとも23次元である。
これに対して五列を全て課すと、関数体上の始点0の商次元は次数0から8で
\[
              1,4,10,20,30,36,34,20,0
\]
となる。
\end{proposition}
\begin{proof}
\path{computations/verify_rank_fifteen_full_center_cut.py} は
第\ref{sec:rank-fifteen-bivariate-sixth-cut}節の語を末尾へ延ばす再帰を
$\mathbb F_3(K,H)$ 上で実行する。四列の九次の文字別次元は $(6,6,6,5)$ である。
代入可能性には、最後まで同じ消去主元を取れることを仮定しない。
無切断の五次までの語基底と四次中心の規格化にはもとのチャートの単元だけが必要である。
長さ9の全語に二次関係と規格化四列を挿入した表示に戻ると、
余核の次元はどの特殊化でも一般点以上である。平方根による語の基底変換は
単元の体拡大上で実行でき、次元を変えない。この議論で得る下界は標数3に関するものである。
五列の場合、六次から八次までの四文字の次元は
$(10,10,10,4),(5,5,5,5),(0,0,0,0)$ である。
\end{proof}

\begin{theorem}\label{thm:ordinary-fifteen-finite-over-quartics}
正の固定曲線の標数零の全通常階数15窓で、折り畳み環 $A$ は四次中心の生成する部分環上で有限である。
四次中心全体で割った商は各始点で長さ13において零となる。
この族では次の手順が、非正則の場合も含め、各入力の標準Koszul左右AS正則性を決定する。
四次中心の基底を長さ5の交換子で求め、その四本の一般線型結合で割る。
係数を独立な不定元として、四始点・長さ0から13で商の次元が
$(1+t+t^2+t^3)^4$ の係数以下かを検査する。この条件は必要十分である。
\end{theorem}

\begin{proof}
四つの中心四次式上ですでに自由性が証明された窓には有限性がある。
第\ref{sec:rank-fifteen-quadratic-parabolic-pencils}節の帰着と
定理\ref{thm:rank-fifteen-all-quadratic-verticals},
\ref{thm:rank-fifteen-prime-parabolic-pencils}により、残るのは3上の付値で
$X=\bar x\notin\mathbb F_9$, $Y=\bar y\ne0,(1-X)^2$ の単元対である。
$X=c^2$, $B=Y/X^3=(v^3-v)^2$ とする。
前命題の五中心の八次までの主元を集め、開集合の単元を除いた
$G(v,c)$ の二次数は $(177,140)$ である。
全十二表示で $G=0$ となる必要十分条件は、自由基底
$c^i v^j$ ($0\le i<2$, $0\le j<6$) における十二係数が全て零となることである。
ここでは $XB\ne0$ により各ファイバーが被約である。

検証コードは十二係数の最大公約元が1であることを確認する。
第\ref{sec:rank-fifteen-bivariate-seventh-cut}節の手順で各係数から
$X, X-1,B,X^3B-(1-X)^2$ の因子を除き、各 $i$ で $(j,j+1)$
($j=0,2,4$) の組から $B$ を消去する。$X, X-1$ の因子を除いた
六つの終結式のモニック最大公約元は
\[
 (X+1)^{10}(X^2+1)^2 (X^3+2X+2)(X^5+2X^4+2X+2)
\]
である。後の二既約因子の剰余体で、除去前の十二係数を $B$ の多項式として
最大公約元を取ると、その全幾何根はそれぞれ
$B=X^2+2X+1$ と $B=2X^3+X^2$ の一点ずつである。
これらの二軌道では元の四窓を組み立て、四中心商の標準次元を
四始点・長さ13まで直接照合する。使用した体の絶対拡大次数は順に6と5である。
従ってここでは混標数への既出の四切断の持上げを使える。

その他の点では少なくとも一表示で $G\ne0$ なので、五中心商は始点0の長さ8で零である。
第\ref{sec:rank-fifteen-truncated-central-freeness}節と同じくモニックな座標方程式で
この表示を持ち上げる。長さ5までの経路基底が持ち上がり、中心方程式は剰余で核次元5、
標数零でも一般窓からの下界5を持つため、五つの自由列を全て持ち上げられる。
語の有限表示の余核次元は分数体で増えない。他の始点でも長さ11の経路は
始点0からの長さ8の区間を含むので、五中心商は有限次元である。
中心倍を除いて長さに関して帰納すれば、元の経路代数が五中心上で有限となる。
ここで目標と同じ $(x,y)$ を持つ表示の選択は十分である。
定理\ref{thm:rank-fifteen-rational-quotient}の証明の巡回移動、語反転、基底変更は
四次中心上での有限性を保つ。この理由で、固定した始点の商次元をそのまま運ぶときに
必要だった符号の制約を課さずに十二表示を使える。
四中心上の自由性を使った場合には四中心商の長さ13が零であり、
残りの場合は上の五中心商の長さ11が零なので、全四次中心で割った商はどちらも長さ13で零になる。
体拡大後の有限性も降下する。実際、四次中心全体は有限行列の核で体拡大と可換であり、
その商の高次消滅は忠実平坦な拡大から降りる。

正則な入力では有限性と斉次Noether正規化により、四つの一般四次中心が
有限次元商を与え、定理\ref{thm:central-window-cutoff}の必要条件から所定の係数に一致する。
逆に、正の固定曲線のこの族は第二補完の閉包の通常窓であり、当該上界が成立すれば
命題\ref{prop:klein-central-strata}の自由性・正則性の十分条件が適用できる。
これで両方向が従う。添字保持同型は定理\ref{thm:rank-fifteen-global-invariants}の
$(x,p,q)$ で判定できる。この有限手順から、全通常点で正則性の答が一律に肯定となることまでは結論しない。
\end{proof}

% !TeX root = ../papers/nc-p3.tex
\section{階数15の三次曲線束の全点と添字保持同型曲線}
\label{sec:rank-fifteen-cubic-pencils}
\begin{theorem}\label{thm:rank-fifteen-cubic-pencils}
正の固定曲線の通常な階数15窓の座標を $(x,y)$ とする。
3の上の付値で $X=\bar x\ne0,1$、$Y=\bar y\ne0$ が定義され、
$B=Y/X^3\in\mathbb F_9^\times$ なら、窓は標準Koszul左右AS正則である。
各始点は共通の四つの中心四次式の生成する多項式環上で自由階数256である。
従って $\lambda\ne0$ が3の上に剰余次数2以下の素点を持つ数体に属するとき、
曲線 $y=\lambda x^3$ の通常な全点でこの結論が成立する。$\lambda$ は任意の非零有理数でもよい。
\end{theorem}
\begin{proof}
$X\in\mathbb F_9$ と $Y/(1-X)^2\in\mathbb F_9$ の場合は既知なので除く。
$c^2=X$, $(v^3-v)^2=B$, $K=1/v$, $H=c(K-v)$ とする。
固定した $B$ では第\ref{sec:rank-fifteen-bivariate-word-paths}節の語の偶奇降下により、
各 $v$ の表示を純粋な一変数有理関数体 $\mathbb F_3(v)(c)$ で計算できる。
平方根の積を含む二次函数体は不要である。
五本の中心語の第1,2,3,5列を用いる。任意の幾何 $c$ を扱うため、消去で
逆数を用いた全ての分子と入力分母を平方自由な多項式 $D_v(c)$ に記録する。
各表示の始点0の四中心商は長さ1から13まで
\[
             4,10,20,31,40,44,40,31,20,10,4,1,0
\]
となる。これは $D_v(c)\ne0$ の各点で有効である。

計算の被覆を明記する。$i^2=-1$ として五つの代表 $B$ を取り、
非素体の三値の三乗共役も含める。二つの表示は $B=-1$ で $v,v+1$、
それ以外で $v,-v$ とする。さらに $\mathbb F_3(B)$ を固定する係数のFrobenius共役を全て使う。
$c$ の符号は変えない。全表示の除外多項式の最大公約元から
$X\in\mathbb F_9$ または $Y/(1-X)^2\in\mathbb F_9$ の根を除く。
後者には $N=Bc^6$, $D=(1-c^2)^2$ の $N^9D-ND^9$ を用いる。
非素体なら残りとその三乗共役の最小公倍式を取り、素体上の多項式にする。
\[
\begin{array}{c|rr|r|l}
 B &\deg D_v&\deg D_{v'}&\text{残余次数}&
   \text{既約因子の次数（括弧内は本数）}\\ \hline
 1&3418&3824&18&3(6)\\
 -1&3460&3963&38&4(8),\,6(1)\\
 -1+i&4058&4058&60&4(11),\,8(2)\\
 i&3056&3075&0&\\
 1+i&4097&3741&40&4(10)
\end{array}
\]
有限例外の各既約因子について、分解体で全根 $c$ を
$\mathbb F_3(B)$ のFrobenius軌道に分ける。
各代表で $X=c^2$, $Y=BX^3$ を戻し、必要な平方根が全て分解するまで体を延長して
元の四窓を組み立てる。既知の場合を除く全代表で、四次中心と未切断の五次までの道を再計算し、
四始点とも四中心商が十三次の消滅まで標準次元であることを確認する。
\path{computations/verify_rank_fifteen_cubic_pencils.py} は入力から中心式も再生成する。
五行の主元、最大公約元、全例外の直接計算は有限体での厳密演算である。

$c$ を固定した表示の変更は第\ref{sec:rank-fifteen-oriented-eighth-cut}節の
$\bar p=c^3$ と $q$ の符号選択により、指定した添字保持窓に持ち上がる。
同じ $(x,y)$ の各移動にもこの証明を適用すると、始点ごとに十三次までの次元上界を
与える四中心が取れる。中心四平面のGrassmann多様体上でこれら四つの非空開集合を交差させ、
共通の四中心を得る。命題\ref{prop:klein-central-strata}により、上界は全次数の自由性と
左右AS正則性を与える。数体の主張では所定の付値を係数体へ延長し、$x$ の負付値、
$X=0,1$、および $x$ が単元で $y$ が非単元の場合を既証の全水準の定理で除く。
残る単元域では $B=\bar\lambda$ である。
\end{proof}

\begin{corollary}
上の三次曲線束の添字保持同型不変量の曲線は、$\lambda\ne0,1/27$ のとき種数2の
滑らかな完備曲線の開部分である。通常な全点は最小周期2を持つ。
一位置移動で割った曲線の種数は1である。$\lambda=1/27$ では被約曲線の正規化は有理で、
$(x,p,q)=(3,0,0),(3/4,0,0)$ が周期1の点である。
\end{corollary}
\begin{proof}
定理\ref{thm:rank-fifteen-global-invariants}の座標では $p^2=x^3-27y$、
$q^2=(1-x)^2-4y$ であり、一位置移動は $(p,q)$ の両符号を反転する。
$a=1-27\lambda\ne0$ とすると、$t=p/x$ で $x=t^2/a$, $p=t^3/a$ と復元できる。
$q^2=(1-t^2/a)^2-4\lambda t^6/a^3$ の右辺は平方自由である。
実際、$x$ に関する三次式の判別式は $16\lambda(1-27\lambda)$ で、定数項は1である。
六単根を持つ二重被覆の種数は2となる。移動の不変量 $x,tq$ は
$(tq)^2=ax((1-x)^2-4\lambda x^3)$ の滑らかな種数1曲線を与える。
通常な領域では $x\ne0$ なので $p\ne0$ である。特殊値では $p^2=0$ となり、
被約曲線は $q^2=-(x-3)^2(4x-3)/27$ である。
$q/(x-3)$ が正規化の有理パラメータとなり、$q=0$ の二点とその他で移動の周期が分かれる。
\end{proof}

% !TeX root = ../papers/nc-p3.tex
\section{一文字次数の移動による全通常階数15窓の九次自由性}
\label{sec:rank-fifteen-ninth-character}
\begin{theorem}\label{thm:rank-fifteen-ninth-character}
標数零の正の固定曲線の全通常階数15窓で、四始点に共通の四つの中心四次式を選び、
九次まで中心単項式基底を取れる。各始点の九次経路は220次元、中心商は20次元である。
\end{theorem}
\begin{proof}
第\ref{sec:rank-fifteen-ninth-orbit-obstruction}節で固定四列の九次が23次元となった
障害を、別の四列の二変数消去なしに処理する。残った標数3の単元域で、正規化前の
36成分のWalsh表示を用いる。そのKlein文字次数を $0,1,2,3$ とする。
全四位置で文字に同じ $a$ を加える単項行列による自己同型が $a=1,2$ で存在する。
実際、第\ref{sec:rank-fifteen-genus-two-unit-line}節のWalsh和で
四文字を $(i,j,k,t)$、$h=i\oplus k=j\oplus t$ とする。
$1$ の加算では $i\oplus j$ と $k,t$ の上位ビットが保たれ、各項が同じである。
$2$ の加算では和の添字 $u$ と $\beta u$ の上位ビットの和が符号に加わる。
$\beta=(0,3,2,1)$ なのでこれは $u$ の下位ビット、すなわち $h$ の下位ビットである。
位置0,2の各文字にその下位ビットの符号を掛ければ元に戻る。
四周期対角基底変更による正規化後も、文字の置換は同じである。
長さ4では合計文字は変わらないので文字0の四次中心空間 $V$ は保たれ、
奇数長では合計文字を任意に平行移動できる。

まず固定した $c^2=X$ の六表示を調べる。固定四列 $(1,2,4,5)$ の
長さ8までの始点0の計算に続き、長さ9では文字3の関係行列だけを消去する。
この文字の商次元は5である。残る三文字の上界はここでは用いない。
\path{computations/verify_rank_fifteen_ninth_character.py} はこの計算で逆数を取った
主元を集める。既知のチャートの単元を除いた $G(v,c)$ の二次数は $(181,160)$ である。
$B=Y/X^3$、$(v^3-v)^2=B$ の六基底における係数の最大公約元は1である。
第\ref{sec:rank-fifteen-oriented-eighth-cut}節の符号を保つ終結式上界の次数は67である。
$c$ の次数2以下は既知で、それ以外に残る $(c,B)$ のFrobenius軌道は
三次が八つ、四次が三つである。各軌道で元の窓に戻して十三次まで全四始点の
標準商を検算する。従ってその十一軌道では既出の全次数の持上げを使える。

その他では、固定した $c$ の少なくとも一表示で、固定四平面 $U\subset V$ に対し
$\dim (A/(U))_{0,9,3}\le5$ である。まず\emph{同じ剰余窓の中で}四平面を動かす。
上の自己同型によりどの文字 $h$ でも $\dim(A/(U_h))_{0,9,h}\le5$ の四平面が存在する。
各不等式は $\operatorname{Gr}(4,V)$ 上で開く。中心が文字0なので選び直した全ての
四平面でも商に文字次数があり、各文字の関係挿入行列の小行列式が使える。
四つの非空開集合を交差させれば、同一の四平面で九次の全商が20次元以下となる。
この操作は固定四列のまま九次が20になることを必要としない。

五つの中心の基底が持ち上がるので、得た剰余四平面の係数も持ち上げられる。
$c$ を指定した六表示のため、$\bar p=c^3$ により目的の添字保持窓で始点0の上界を得る。
他の始点にも移動した窓で同じ議論を行い、標数零の全四次中心のGrassmann多様体で
四始点の開条件と八次までの自由性の開条件を交差させる。
中心倍による生成と未切断経路の一般点からの下界から
\[
  220=\binom{12}{3}\ \le\ \dim A_{p,9}\
     \le\ \dim Q_{p,9}+4\cdot40+10\cdot4\ \le\ 220 .
\]
よって全て等号であり、九次でも中心商基底の持上げと中心単項式の積が基底になる。
\end{proof}

% !TeX root = ../papers/nc-p3.tex
\section{五中心商の乗法表から全ての四中心商を求める}
\label{sec:rank-fifteen-center-border}
\begin{proposition}\label{prop:five-center-border}
四周期の二次経路代数 $A$ と五つの中心四次式 $z_1,\ldots,z_5$ を取り、
固定した始点で $M=A/(z_1,\ldots,z_5)$ の長さ8が零とする。
$[b_1:\cdots:b_5]\in\mathbb P^4$ に垂直な四式で割った始点加群を $Q_b$ とする。
$M$ の長さ7以下の語基底 $W$ を用い、次の一つの多項式行列の計算を全ての $b$ に使える。
\begin{enumerate}
\item $F=\bigoplus_{w\in W}k[Z]w$, $\deg Z=4$ と置く。
長さ8までの語消去で各矢の右作用を $F$ 上に持ち上げる。
除算するのは五中心商の語消去の主元だけで、$b$ による除算はない。
\item 持ち上げた二次関係を全基底語に作用させ、その残余と、最初の五中心関係の残余を取る。
$Z$ と全矢の右作用でそれらの線型包を閉じたものを $N$ とすれば $Q_b=F/N$ である。
従って各次数への伝播は、長さ7以下の語に対する固定した乗法表だけを使う。
\item 五中心商の各文字の次数0から7の次元が
\[
 (1,0,0,0),\ (1,1,1,1),\ (2,2,2,4),\ (5,5,5,5),\
 (6,8,8,8),\ (9,9,9,9),\ (10,10,10,4),\ (5,5,5,5)
\]
とする。四中心商の長さ9から13の文字別の目標上界を順に
$(5,5,5,5),(2,2,2,4),(1,1,1,1),(1,0,0,0),(0,0,0,0)$
とすると、各条件は $b$ の斉次小行列式の非消滅の和であり、その次数上界は順に
\[
 (6,6,6,6),\ (12,12,12,8),\ (14,14,14,14),\ (13,16,16,16),\
 (21,21,21,21)
\]
となる。ここでは五中心が文字0の場合を扱う。
\end{enumerate}
\end{proposition}
\begin{proof}
$Q_b$ は中心変数 $Z$ を加え $z_j=b_j Z$ を課した代数の始点行である。
実際、ある $b_i\ne0$ なので $Z$ を $z_i/b_i$ で消去できる。
$Z$ を割った商が $M$ であり、次数帰納で $F\to Q_b$ が全射である。

乗法表を具体化する。長さ $n$ でまだ決まらない列は、長さ $n-1$ の五中心商の
基底語に一文字を付けたものだけである。その他は $Z F_{n-4}$ に入り、既出の表を使う。
長さ $n-2$ の基底語に末尾二次関係を付けた行、および $n=4$ では
五行 $z_j-b_j Z$ を作る。新しい列だけへの射影は $M$ の語の関係行列そのものである。
これは $Z=0$ とすれば、各段で前段までの関係と最後の二次関係を全て課した商となることによる。
中央関係の右倍は前段の乗法の簡約に含まれ、中央性から途中へ挿入した関係も含まれる。
主元列を非主元列と $Z F_{n-4}$ の和で表し、残る列を長さ $n$ の基底語とする。
$M_8=0$ なら長さ8では全ての新しい列が消える。$Z$ と可換に延ばして全基底語の
右作用が定まった。全ての式は実際の $Q_b$ で成り立つ。

新しい列を消去した後の残余も関係として保持する。$F/N$ 上では全基底語で
二次関係が零なので、矢の作用は $A$ の右作用を与える。空語に矢を掛ければ全非主元語を回収できる。
$Z$ は作用と可換で、$z_j$ は $A$ の中心だから、空語で確認した
$z_j=b_j Z$ は全ての $Z^d w$ にも成立する。従って $Q_b\to F/N$ が得られ、
基底とその中心倍で両者の自然な写像は互いに逆である。
長さ8までに残余がなく、上の次元の場合には、最初の関係は長さ7の基底語に
二次関係を掛けた九次の残余だけで尽きる。これ以後はこれらの右倍と $Z$ 倍を取ればよい。

長さ $n$ の一行で、列 $Z^j w$ の係数は $b$ について次数 $j$ の斉次式となる。
実際、中央関係の右辺で $Z$ と一緒に $b_i$ が一つ加わり、元の二次関係は $b$ を含まない。
既出の行の $Z$ 倍は行全体で次数が一つ下がる。同次な全右倍をそのまま並べるか、
$Z$ 倍の代わりに全ての $b_iZ$ 倍を並べれば射影点での線型包は変わらない。
必要な階数を $r$ とすれば、$r$ 列の $j$ の和の最大値が小行列式次数の上界になる。
例えば十三次の各文字は $Z^2 M_5$ に九列、$Z^3 M_1$ に一列あり、
全階数の小行列式は $2\cdot9+3=21$ 次以下である。各次数で同じ計算を行うと表示した表を得る。
従って各文字・次数である良い $b$ が存在するかの検査には、四変数・全次数21の
補間格子で足りる。アフィン座標に相異なる22標本を用いるNewton補間により、
正標数でも必要なら有限体拡大で同じ結論が従う。
\end{proof}

\begin{proposition}\label{prop:ordinary-bivariate-border-cut}
標数3の非分岐単元表示の関数体 $\mathbb F_3(K,H)$ で、始点0において
五本の規格化中心の第2,3,4,5列による商は、長さ1から13まで
\[
               4,10,20,31,40,44,40,31,20,10,4,1,0
\]
の次元を持つ。この証明書の除外因子は元の $(K,H)$ 平面上で171個である。
\end{proposition}
\begin{proof}
前命題の二次関係は語の偶奇による二次拡大なしに持ち上げられる。
$J$ を六平方の積、$\alpha^2=J$ とし、次の矢 $a$ について
\[
 \delta(n,h,a)=\bigl([p+n=2\bmod4,\ a=3]+\beta(h,a)\bigr)\bmod2
\]
と置く。これは文字 $h$ の長さ $n$ の語に $a$ を付けたときの
$\ell_p$ の増分の偶奇であり、$n$ を4減らしても同じである。
そこで $a$ の右作用を $\alpha^\delta$ 倍して用いる。この正規化で
語 $w$ の係数の基底変換は $\alpha^{D(w)}$、ただし $D(w)$ は順次の $\delta$ の和である。
中心式で必要な追加倍率 $\alpha^{\ell_p(w)-D(w)}$ の指数は偶数である。
二次関係の各係数の $\alpha$ の指数から連続二回の $\delta$ を引いても、同一行では同じ偶奇となる。
従って行全体の共通倍率を除き、全ての新旧列を $\mathbb F_3(K,H)$ 上で扱える。

$Z$ 倍で短い語に戻る列を保存したまま、五中心商の新しい語の列だけを掃き出す。
初期の長さ8まで余分な関係はなく、九次には各文字の自由十列に独立な五つの残余が現れる。
以後は直前の残余の原行を選んで一文字付け、四つ前の残余の中心倍を併合する。
九次以後の文字別商は前命題の目標値にちょうど一致する。
\path{computations/verify_rank_fifteen_bivariate_border.py} は二変数のままこれを検算する。
主元の分子、最初の経路と中心の単元を含めて171既約因子が現れる。
各因子の零点を除けばこの始点での上界は特殊化できる。
さらに \path{computations/verify_rank_fifteen_center_border.py} は
前出の $\mathbb F_{81}$ の点の七つの中心方向と四始点で、全長について
元の四中心商から始める直接消去と照合する。
この関数体計算に残る例外点での正則性は、関数体の次元からは決まらない。
\end{proof}

% !TeX root = ../papers/nc-p3.tex
\section{正の固定曲線の通常階数15窓の全点分類}
\label{sec:rank-fifteen-all-ordinary}
\begin{theorem}\label{thm:rank-fifteen-all-ordinary}
標数零で定理\ref{thm:rank-fifteen-covers}の正の固定曲線から来る二表示の
全ての通常階数15窓は、標準14641型quadratic Koszul左右AS正則である。
四始点で共通の四つの中心四次式を選び、各始点の経路をその多項式環上の
自由階数256の加群にできる。従ってこの枝に通常な非正則点は存在しない。
定理\ref{thm:rank-fifteen-formal-boundary}の四点を加えた規格化窓の添字保持同型類は
\[
        y\ne0,\qquad p^2=x^3-27y,\qquad q^2=(1-x)^2-4y
\]
の各点に一対一に対応する。一位置移動は $(p,q)\mapsto(-p,-q)$ で、
$x=3,3/4$ で $p=q=0$ の二点だけが周期1、その他は全て周期2である。
\end{theorem}
\begin{proof}[九次以後の四つの上界]
第\ref{sec:three-center-pencil}節までの帰着と
定理\ref{thm:rank-fifteen-ninth-character}により、剰余標数3の
$X\notin\mathbb F_9$, $Y\ne0,(1-X)^2$ の単元対で、長さ10から13の
始点0の上界 $10,4,1,0$ を示せばよい。$c^2=X$ の符号を指定し、
$B=Y/X^3$、$(v^3-v)^2=B$ の六表示を使う。

命題\ref{prop:ordinary-bivariate-border-cut}と同じく第2,3,4,5列を課し、
第1列を変数 $Z$ に送る。初期の五中心商の主元は通常通り記録する。
長さ9以後では、右倍に用いるのは独立行を選ぶ前の元の行である。
よって後の次数の上界を証明するとき、それ以前に独立と判定した残余の
独立性を特殊化先で仮定する必要はない。前段の選択が従属に変わっても、
選ばれた各行とその右倍は依然として真の関係である。

現在の文字の段でも、選んだ原行と主元列でできる一つの小行列式だけを残す。
その値は、行と列の順序の符号を除き、有理係数としての主元の積である。
実際、選ばれた行だけを同じ順で消去する行加算はその小行列式を変えない。
原始多項式行に付随する有理倍率を含めて主元を掛け、約分後の分子を用いる。
特殊化でこの分子が単元なら、初期表が定義される限り、途中の個々の消去を再現できなくても
必要な原行の独立性が保証される。十一、十三次はWalsh文字の平行移動と
定理\ref{thm:rank-fifteen-ninth-character}の四平面の開条件により、
文字0だけでよい。偶数次は四文字の小行列式を同じ表示で全て使う。

各次数に必要な因子だけを $(v,c)$ へ引き戻した平方自由積を $D$ とする。
$D$ を $1,v,\ldots,v^5$ に展開した六係数の共通因子は、四行とも次の同じ多項式 $G(c,B)$
である。$u=cv$ として
\[
  F(c,v)=u^6-(c^2+1)u^4+(c^4-c^3+c)u^2+c^3-c^2-c+1,\qquad
  G(c,(v^3-v)^2)=\prod_{j\in\mathbb F_3}F(c,v+j).
\]
右辺の六係数は定数係数以外が零で、$G$ の二次数は $(18,3)$ である。
$G$ を除いた六係数に符号を固定した終結式上界を適用すると次の表を得る。
\[
\begin{array}{c|c|c|c|c}
 \text{長さ}&\text{検査する文字}&(\deg_v D,\deg_c D)&
  c\text{ の終結式上界の次数}&\text{残余軌道数}\\ \hline
 10&0,1,2,3&(120,114)&6&2\\
 11&0&(84,80)&0&0\\
 12&0,1,2,3&(171,170)&9&3\\
 13&0&(88,86)&3&1
\end{array}
\]
全四行の例外を合わせると三軌道だけである。
$c^3+2c^2+c+1=0$ で $B=1+c+2c^2$、
$c^3+2c^2+2c+2=0$ で $B=1$ または $2+2c+c^2$ となる。
各軌道は元の四窓に戻して全四始点の長さ13までを直接検査する。
\path{computations/verify_rank_fifteen_all_ordinary.py} は、二変数の語関係から
四行の小行列式・六係数・共通曲線・三軌道の検査を再現する。
\end{proof}

\begin{proof}[共通例外曲線の全点]
$G(c,B)=0$ なら、所定の $c$ を変えず、六つの $v$ のうち
$F(c,v)=0$ を満たすものが取れる。これは上の積の恒等式に従う。
$F$ は $\mathbb F_3[c,v]$ で既約であり、$c\ne0$ で $u=cv$ を使うと表示した
モニック六次式になる。$\mathbb F_3(c)[u]/(F)$ 上でまず $v=u/c$ の表示で第1列を落とし、
次に $v'=u/c+1$ の表示で第4列を落とす。いずれも始点0の四中心商は長さ13の消滅まで標準次元を持つ。

除外点の処理では、除算のノルムの平方自由因子を $c$ の一変数多項式として記録する。
具体的には $t=u^2$ として元を $(a(t)+u b(t))/d(c)$、
$\deg_t a,\deg_t b\le2$ と書く。逆元は $a^2-t b^2$ の三次拡大での乗法行列の余因子と
行列式で計算できる。二表示の記録の次数は6986と17641である。その最大公約元から
$c^{18}-c^2$ の根を除くと、残りは平方自由な次数237で、既約因子の数は
次数 $4,5,6,7,8$ の順に $5,11,15,8,2$ である。
各因子の一根 $c$ で $F(c,v)$ を有限体上分解し、全ての $B=(v^3-v)^2$ を
$\mathbb F_3(c)$ のFrobenius軌道に分ける。
$B\in\mathbb F_9$ は定理\ref{thm:rank-fifteen-cubic-pencils}、
$Y/(1-X)^2\in\mathbb F_9$ は既出の放物線束の定理で覆われる。
残る100軌道では、所定の $c$ の六表示と規格化五列の四中心から選び、
最初の二次関係に戻って十三次までの上界を検算する。これらの計算は
\path{computations/verify_rank_fifteen_signed_curve.py} にあり、
除算する六次拡大の算術は \path{computations/rank_fifteen_signed_curve_field.py} にある。

これで、各次数で、共通曲線の上も外も含めて、始点0で所定の上界を持つ表示が得られた。
必要な開集合の交叉は奇数次の四文字については一つの剰余表示の中で行う。
五つの中心を持ち上げ、指定した符号 $\bar p=c^3$ で目的の標数零の窓へ運ぶ。
四つの次数で同じ表示を選ぶ必要はない。各移動後の窓にも適用すると、元の窓の
全四次中心のGrassmann多様体上で各始点・各次数について非空開集合が得られる。
九次までに得た開集合も合わせて交差させ、同一の四中心で長さ13まで全ての上界が成立する。
命題\ref{prop:klein-central-strata}の全次数の自由性と左右Koszul完全性、有限双対の
完全対合から所望の正則性を得る。
最後に定理\ref{thm:rank-fifteen-global-invariants}と
\ref{thm:rank-fifteen-formal-boundary}の同型条件・通常表示の像と境界四点を適用する。
同型不変量の曲面は
$(2p+3\sqrt3q)(2p-3\sqrt3q)=(x-3)^2(4x-3)$
と書ける既約有理曲面の開部分で、一位置移動の固定点は述べた二点だけである。
\end{proof}

% !TeX root = ../papers/nc-p3.tex
\section{Walsh基底を持つ四窓の三つの切断と階数12,14,15の認識}
\label{sec:walsh-symmetric-windows}
$E=(\mathbb Z/2)^2=\{0,1,2,3\}$ の加法を $+$、元の低位、高位のビットを
$x_0,x_1$ と書く。四位置に $E$ でラベルされた基底を選び、四次テンソル $W$ が
$i+j+k+t=0$ に台を持つとする。全ラベルへの1の加算は係数を保ち、
2の加算は係数に $(-1)^{(i+k)_0}$ を掛けると仮定する。
以下「Walsh基底を持つ」とはこのような四基底の存在をいう。基底は同型のデータに含めない。
\begin{theorem}\label{thm:walsh-symmetric-windows}
二つの隣接切断がともに階数6のWalsh基底を持つ四次テンソルを取る。
各位置の二次関係として巡回隣接切断の像を用いる。
\begin{enumerate}
\item 反対階数15の全候補は、定理\ref{thm:rank-fifteen-all-ordinary}の境界四点を含む
規格化族にちょうど一致する。従って全候補が標準Koszul左右AS正則で、
指定基底を忘れた同型類と周期は同定理の有理曲面で尽くされる。
\item 反対階数12または14なら、隣接する二切断で非零の文字ブロックがそれぞれ三つあることが
標準Koszul左右AS正則性の必要十分条件である。階数12で正則なら
定理\ref{thm:rank-twelve-quantum-limit}の唯一の量子窓である。
階数14の全正則窓は二本の通常楕円表示とそのWalsh表示の端で覆われる。
\item この対称性を持つ標準窓には反対階数11,13は生じない。
\end{enumerate}
\end{theorem}
\begin{proof}
$b=i+j$, $e=j+k$, $d=b+e$ と置く。対称性から
\[
 W_{i,i+b,i+d,i+e}=(-1)^{i_1d_0}m_{be}
\]
と書ける。$01|23$ の $b$ ブロックは $i,k$ を行列添字とする。
$d$ ごとの符号付き移動 $P_d$ は $(P_d)_{ik}=(-1)^{i_1d_0}[k=i+d]$ である。
$P_1^2=P_2^2=1$, $P_1P_2=-P_2P_1$ から生成代数は $M_2$ となり、
四次元表現はその標準表現の二重である。$s_d=(-1)^{d_0+d_1}$ とすれば
二次元表現の行列式は $\sum_ds_d m_{b,b+d}^2$ である。実際 $P_1,P_2,P_3$ は
対ごとに反可換で、二乗は順に $1,1,-1$ である。$12|30$ の $e$ ブロックを $j,t$ で書くと
符号は $(-1)^{j_1 d_0}(-1)^{(d+e)_1d_0}$ となる。後の因子は各係数の平方を保つので、
列に対して同じ行列式条件を得る。他方 $02|13$ の $d$ ブロックは行の符号を除き
通常の巡回行列で、四固有値は $b\mapsto m_{b,b+d}$ の四つの文字和である。

両隣接切断とも各ブロックが偶数階数なので、最低一つずつ欠ける。
二つ以上欠ける側があれば、全辺で欠けている色を選んだ語は、関係を挿入した
どの語の台にも現れない。各辺の色と最初の文字を独立に選べるので、長さ $n$ で
少なくとも $4\cdot2^{\lfloor(n-1)/2\rfloor}$ 個が独立となり標準成長に反する。
また一方に非零の色が二つ以下なら、各 $d$ の対角列の台は高々二点である。
二点以下の文字和は非零の個数が偶数となるので、階数15では成長を仮定せず三色ずつが必要である。

四位置を $(0,B,B+E',E')$ ずらすと $b,e$ を任意にずらせる。
余分な符号 $(-1)^{i_1(B+E')_0}$ は最初の基底で除く。
両方の欠色を3にしたとき $m$ は左上 $3\times3$ に台を持つ。
全非零ブロックで階数2となる条件は $N_{be}=s_bs_em_{be}^2$ の全行和・列和の消滅である。
特に $N_{01}-N_{10}=N_{12}-N_{21}=N_{20}-N_{02}$ である。
この共通値が零なら $d\ne0$ の各対の係数は符号を除いて一致し、そのブロックは高々階数2、
従って反対階数は高々10である。非零ならこれら三ブロックは可逆で、
$\rho=12+\operatorname{rank} W_{02|13,0}$ となる。
最後の四つの文字和は台が高々三点なので一つだけが非零とはならない
（逆Fourier変換を取れば台が四点になる）。これで第三項も従う。

$\rho=12,14,15$ なら対角の文字和の一つが零である。$m_{be}$ に $b$ の文字を掛ける操作は
位置0,1に同じ対角文字を掛ける基底変換である。零を自明文字に移し、三つの対角成分を
$a+b,-b,-a$ と書く。残りは
\[
 m=\begin{pmatrix}a+b&p&n\\ m'&-b&l\\ -q&r&-a\end{pmatrix},\qquad
 r^2+l^2=2ab,\quad p^2=b^2+r^2,\quad (m')^2=b^2+l^2,\quad
 q^2=a^2+r^2,\quad n^2=a^2+l^2 .
\]
右の五式は全行和・列和と同値であり、$d\ne0$ の各ブロックが可逆なのは $r^2-l^2\ne0$ と同値である。
定理\ref{thm:rank-fifteen-parameter-orbits}のPauli表示の八係数は
\[
 (a,b,L,M,A,B,C,D)=
 \left(a,b,\tfrac{p+m'}2,\tfrac{p-m'}2,\tfrac{n-q}2,\tfrac{n+q}2,
                 \tfrac{l-r}2,\tfrac{l+r}2\right).
\]
四位置のWalsh変換でこれがまさに上の $W$ になることは各添字での文字和で得られる。
\path{computations/verify_walsh_symmetric_windows.py} は成分の恒等式と五式を独立変数で検算する。

階数15なら $ab(a+b)\ne0$ である。$a=1$ に規格化し
$s=q+r$, $t=n+l$, $N'=p+m'$, $H=r^2-l^2$, $K=1+2b$ とする。
$s^{-1}=q-r$, $t^{-1}=n-l$, $H/N'=p-m'$ なので三元は全て単元である。
$r^2+l^2=2b$ より
\[
 K=\tfrac14(s^2+s^{-2}+t^2+t^{-2}),\quad
 H=\tfrac14(s^2+s^{-2}-t^2-t^{-2}),\quad
 (N')^4+(1-K^2)(N')^2+H^2=0 .
\]
$K^2\ne1$, $H\ne0$ で、これが定理\ref{thm:rank-fifteen-rational-quotient}の
規格化表示そのものである。逆は同じWalsh変換から従う。
階数12なら $a=b=0$ であり、八係数の五式に定理\ref{thm:rank-twelve-quantum-limit}を適用できる。

階数14では三つの対角値のちょうど一つが零である。八係数の表示を保つ192変換は
ラベル $1,8,9$ に対応する $b,a,a+b$ のいずれも $b$ の位置へ運べる。
よって $b=0$, $a=1$ として上の同じ復元式を使い、$K=1$, $H\ne0$ を得る。
$H\notin\{1,-1,2,-2\}$ なら倍率は通常であり
（$K+H\ne0,\pm1$, $K-H\ne\pm1$）、定理\ref{thm:rank-fourteen-all-ordinary}による。
残りの二水準 $H^2=1,4$ では平方根の組 $(s,t,N')$ はそれぞれ128,64個である。
三元の逆の式があるので、異なる組は規格化係数として相異なる。

97の上の還元で全ての根を取り、$b=0$ を保つ表示変換の軌道を調べる。
各水準はそれぞれ64個ずつの二軌道、一軌道に分かれる。$H^2=1$ の各軌道は
通常な表示を含む。$H^2=4$ の一代表では全始点で双対経路は $(1,4,6,4,1,0)$、
元の経路は五次まで標準次元、四次中心は五次元で、
四中心商は長さ13の消滅まで所定の次元を持つ。前掲コードは全根とこれらの関係を直接計算する。
この有限集合での軌道計算は標数零にも成立する。実際、$a=1,b=0$ で変換した
$H^2=16(LM)^2$ は $LM=-AB=-CD$ により不変であり、
規格化の分母は四乗根の単元である。$s,t$ は $s^2+s^{-2}=2(1+H)$,
$t^2+t^{-2}=2(1-H)$ の根、$(N')^4=-H^2$ である。
$H=\pm2$ で重なる相反根は既に標数零で重なり、97で他の合流はない。
従って全パラメータ集合の還元は全単射で群作用と可換である。

最後の一代表を持ち上げる際、上の三つの根の方程式は $H$ 上モニックで定数項が単元なので
特殊点は通常な階数14窓の閉包にある（各根の環は有限自由、そのテンソル積も有限自由）。
通常点からの下界と剰余計算の上界が道空間の五次まで、双対経路、全四次中心で一致する。
中心の核は持ち上がり、四中心商の上界と全次数の経路次元の下界から
定理\ref{thm:rank-fourteen-all-ordinary}と同じ自由性・Koszul完全性・左右完全対合が従う。
\end{proof}

% !TeX root = ../papers/nc-p3.tex
\section{Walsh階数14窓の添字保持同型曲線と唯一の周期1点}
\label{sec:walsh-rank-fourteen-orbits}
\begin{theorem}\label{thm:walsh-rank-fourteen-orbits}
定理\ref{thm:walsh-symmetric-windows}の正則な階数14窓の添字保持同型類は
$\mathbb A^1\setminus\{-1,1\}$ の点と一対一に対応する。
$a=1,b=0$ の表示で $V=H^2$、
$E_6=16(L^2-M^2)(B^2-A^2)(D^2-C^2)$ と置けば座標は
$z=E_6/(4V)$、$V=4(1-z^2)$ である。一位置移動は $z\mapsto-z$ となり、
$z=0$ が唯一の最小周期1点、それ以外の全点の最小周期は2である。
\end{theorem}
\begin{proof}
$P_{h,u}$ は $(i,k)$ 成分が $[i=k+h](-1)^{u\cdot k}$ である四次のPauli行列とする。
反対切断を $\sum\gamma_{h+4u}P_{h,u}\otimes P_{h,\beta u}$ と書く。
$\beta$ はラベル順に $(0,3,2,1)$ である。二つの核はともに
$I,P_{1,0}$ で張られる。実際、この二ラベルだけが消え、Pauli基底のトレース双対は
各行列自身のスカラー倍である。各核内には階数2の直線が $I\pm P_{1,0}$ の
二つだけあり、いずれも二次元の像と相手側の二次元空間を結ぶ同型をスカラー倍まで定める。
これで各二頂点に結び付いた二次元直和分解が定まる。低位ビットをFourier変換すると、
二直線の代表は二直和因子それぞれの単位行列となる。

その分解で各行列因子の対角ブロック同士にテンソルを射影する。
異なる低位符号の組で射影は零、同じ低位符号の組では核の同型を通して
$\operatorname{End}(k^2)^{\otimes2}$ の元
\[
 2\big(A X\otimes X+ Z\otimes Z+B XZ\otimes XZ\big),\qquad
 X=\begin{pmatrix}0&1\\1&0\end{pmatrix},\quad Z=\operatorname{diag}(1,-1)
\]
を得る。従って偶位置と奇位置の二分解の対応も窓から決まる。
核の同型のスカラーと二つの基底変更は、上のテンソルに共通倍率と独立な共役で作用する。
各側のトレース形式で随伴との積を取ると、トレース零空間上の固有値の射影比は
両側で $[A^2:1:B^2]$ の順不同の組となる。
$A^2+B^2=1$, $AB=-H/4\ne0$ なので、他の二つの和に等しい値はただ一つである。
これを1に規格化すると残り二値の積が $V/16$ となる。よって $V$ は指定基底を忘れた不変量である。
さらに固有値1は単純であり、対応する無跡行列直線は $kZ$ である。
二固有線を取れば、四つの全頂点で各二空間がさらに二直線へ自然に分かれる。
$A^2=B^2$ のときも $kZ$ の固有値は単純なままである。

$H^2=V$ を固定し、$R^2=H/2$, $Q^2=1+R^2$, $N^2=1-R^2$, $i^2=-1$ を取る。
前節の三行列表示で
\[
 r=R,\quad l=iR,\quad q=Q,\quad n=N,\quad m'=iR,\quad p=\pm R
\]
の二代表を使う。この二代表と $b=0$ を保つ表示変換だけで、固定した $V$ の全ての窓を覆う。
実際、一般の $V$ では前節の復元で128組あり、各代表の64変換は相異なる。
$P_2=2(a^2+b^2+(a+b)^2)$ と $E_6$ は192変換で不変
（定理\ref{thm:rank-fifteen-rational-quotient}の斉次恒等式）なので、再規格化 $a'=\pm1$
も $E_6$ を変えない。二代表の $E_6$ は符号の異なる $4VQN$ である。
この被覆は特殊値でも成立する。固定した $R,Q,N$ に対する平方根の全組は、
$H$ を反転する場合の $R\mapsto iR$, $(Q,N)\mapsto(N,Q)$ も含めて
$R,Q,N$ の符号・四乗根倍で列挙される。一般値での被覆はこれらの線型式の恒等式なので
全点に特殊化できる。$V=4$ では前節により一軌道である。

残るのは二代表間の指定なし同型と一位置移動である。上の四直線の標識で、
各基底を低位ビット（分解の符号）と高位ビット（$Z$ の固有直線）で順序付ける。
同型の置換部分は $(P,Q',P,Q')$ の形で、$P,Q'$ は低位符号を同時に反転するかともに保ち、
各低位符号内の高位二直線を独立に交換できる。二つの核の階数2直線を単位行列にしたことから、
位置0,2の置換が一致し、位置1,3でも一致する。従って候補は32個で尽きる。
巡回したテンソルでも核の転置と二因子の交換によって同じ四直線を得る。

この基底での成分は、共通倍率を除き $1,R,Q,N$ の一つの四乗根倍である。
具体的に文字を $i',j,k,t$、$d=(i'+k)_0$, $g=(i'+k)_1$ とすれば、
$i'+k=j+t$ の成分は
\[
 \frac12\sum_{h,b'=0,1}\gamma_{h+2g+4(d+2b')}
   (-1)^{(i'+j)_0h+b'k_1+(b'+d)t_1}
\]
となる（ラベルビットの加法は標数2）。$QN\ne0$ では台は固定である。
台の16列の重み行列を作り、各語の行に対応する四位置の基底変数を記録する。
その整数左核の基底に沿った比の単項式が全て1であることが、対角基底変更の必要十分条件である。
Smith標準形と上の成分式で得る個数を示す。
\[
\begin{array}{c|rrr}
 (+R)\text{ 表示との比較先}&\text{台を保つ}&\text{全比が恒等的に1}&
 \text{定数比だけで不一致}\\ \hline
 -R&8&0&8\\
 (+R)\text{ の一位置移動}&8&0&8\\
 (-R)\text{ の一位置移動}&8&4&4\\
 (+R)\text{ の二位置移動}&8&4&4
\end{array}
\]
\path{computations/verify_walsh_rank_fourteen_orbits.py} は各比を
$i^eR^jQ^kN^l$ として計算する。最後の列では $(j,k,l)=(0,0,0)$,
$e\not\equiv0\bmod4$ の等式が要求されるので、どの非零特殊値でも不一致である。
$Q=0$ または $N=0$ では対応する成分を除いても、符号を替えた代表の一位置移動への
恒等的な対角同値が残る。両代表はそこで同型なので周期1である。
$E_6^2=16V^2(1-V/4)$ から所定の座標を得る。各値を実際に平方根で表示でき、
二代表の区別と移動が上表で確定したので主張が従う。
\end{proof}

% !TeX root = ../papers/nc-p3.tex
\section{Walsh表示の既約三次元族と指定を忘れた一般商}
\label{sec:walsh-presentation-space}
\begin{theorem}\label{thm:walsh-presentation-space}
三色ずつを持つWalsh窓の欠色をともに3にして射影化する。
その表示空間は既約三次元で、幾何学的一般点は反対階数16、最小周期2の
標準Koszul左右AS正則窓である。全ての正則なWalsh窓は最小周期が2を割る。
次の指定なし一般商も次元3である。$N_{be}=(-1)^{(b+e)_0+(b+e)_1}m_{be}^2$ とし
\[
 \textstyle \sum_b N_{be}=\sum_eN_{be}=0,\qquad
 P^2=\prod_{b\ne e}N_{be},\qquad (N,P)\sim(\lambda N,\lambda^3 P).
\]
この二重被覆を行列添字 $0,1,2$ の同時置換 $S_3$ で割ったものが、
一般の添字保持同型類を与える。移動は $(N,P)\mapsto (N^{\mathsf t},-P)$ である。
正則な階数12,14,15窓はこの幾何学的一般点からの特殊化であり、
いずれも全次数で平坦な付値環上の族で実現できる。
\end{theorem}
\begin{proof}
行和・列和が零の $3\times3$ 行列の射影空間は $\mathbb P^3$ である。
九つの成分は相異なる非比例な線型式である。第一成分を基準にした他の八成分の比は
平方類として独立である。各成分が消える超平面の付値で当該指数の偶奇を読めばよい。
射影座標ごとの平方写像は有限平坦であり、逆像の一般ファイバーは次数256の体である。
平坦性によりその閉包が全逆像であり、三色ずつの開条件を取っても既約である。
座標を三つの対角成分と $N_{01}-N_{10}$ に取れる。後者が非零なら
反対階数は $12$ に対角三成分のFourier変換の階数を足したものである。
対角の平方が独立な一般点では四つの文字和は非零である。

この既約族は前二節の正則点を含む。各一文字切断が全階数で隣接階数が6の近傍で、
テンソルとその縮約から双対経路の長さ3,4の次元には $4,1$ の下界がある。
正則点での値と上半連続性から近傍で各々一定であり、長さ5は零となる。
そこで得た有限平坦なFrobenius双対の内部次数別にbar複体を取ると、
定理\ref{thm:reversible-component}の議論で幾何学的一般点のKoszul性と左右AS性が従う。
この段階では全次数について単一の開集合を主張する必要はない。
一般点と任意の上記正則点で各次数のHilbert関数が等しいので、正則点を中心とし
関数体を分数体とする離散付値環（例えばその点の吹上げ・正規化の例外素因子）では
各次数の有限生成経路加群が自由である。

全ラベルへの平行移動 $x$（偶位置に $(-1)^{x_1i_0}$ を補う）と全位置での
文字倍が位数16の射影群 $G$ を与える。位数は一位置の作用だけでも16である。
一般の階数15点の射影安定化群は定理\ref{thm:rank-fifteen-pauli-comparison}により
ちょうど $G$ である。この際、反対核の非退化な二行列で残りの二基底変更を
射影的に決めるので、独立共役群での安定化群は四位置でのものに等しい。
幾何学的一般点でも安定化群はちょうど $G$ であることを説明する。

まず複素数上で四基底を正規直交とする。各一位置のGram行列はラベルの異なる
二切片の台が交わらないため対角で、平行移動が係数の絶対値を保つのでスカラーである。
従って全ての表示は $\mathrm{SL}_4^4$ のノルム最小ベクトルである。実際各因子の
跡零Hermite行列による一径数変換で、ノルム平方は実指数関数の非負和で、
その微分は原点で零である。射影安定化群が有限なら全安定化元は射影ユニタリである。
極分解の正値部分に沿うノルムの最小値が二点で一致すれば一定であり、
正値部分の一径数群がベクトルを固定するからである（最小ノルムの等しい軌道上で
スカラー倍率の絶対値は1）。安定化Lie環の零性は上の階数15点の近傍で開である。
そこへ余分な安定化元を持つ点が解析的に集積するなら、コンパクトな射影ユニタリ群で
部分列は階数15点の安定化元へ収束する。その $G$ の元を除き、単位元へ収束する
非自明な安定化元を得る。単位元近傍の対数をその長さで割り部分列を取れば
極限の非零なLie環元が階数15点の直線を固定し矛盾する。
余分な安定化元の存在は有限型の軌道方程式から構成可能であるため、
この解析的近傍の存在でザリスキー一般点でも排除できる。全条件は代数的なので標数零へ移る。

従って一般の二表示の同型は固定した $G$ を正規化する。$G$ の元を $(x,\chi)$ とすると、
偶頂点と奇頂点での行列の交換子形式の差は
$x_1y_0+y_1x_0$ である。その根基 $x=0$ は通常の対角文字群で、
四本の同時固有線が各頂点で定まる。従って一般の同型は単項行列である。
全九成分が非零の台を保つ置換を数える。四点の全置換は $E$ のアフィン変換で、
$0,1,2$ の全ての $b,e$ について四文字の和が零となる条件から、
四位置の線型部分は同じ $A\in\mathrm{GL}_2(\mathbb F_2)$ である。
欠色3を保つと位置0,2で同じ平行移動、位置1,3でも同じ平行移動となる。
行と列の色には同じ置換 $b\mapsto Ab+3+A3$ が働き、六置換全てが出る。

この置換の符号を補うには双線型形式
$T(x,y)=x_1y_0+(Ax)_1(Ay)_0$ を使う。これは対称であり、
$q(x+y)/(q(x)q(y))=(-1)^{T(x,y)}$ となる四乗根値の関数 $q$ がある。
$q(0)=1$ とし二つの基底で $q(e_j)^2=(-1)^{T(e_j,e_j)}$ を選べばよい。
$i,k=i+d$ の基底倍率で符号を補うと、係数に残るのは $q(d)$ またはその逆と $d$ の文字だけである。
$q(d)^2=s_d/s_{Ad}$ なので $N$ はそのまま置換される。
非対角六成分では各 $d\ne0$ が二回出て $\prod_{d\ne0}q(d)^2=1$ となり、$P=\prod_{b\ne e}m_{be}$ も保つ。

位置別の対角変換の効果を決める。ある $(b,e)$ の四語で基底倍率の積を取ると
各位置の全倍率を一度ずつ使うため、全係数比の四乗は一定である。
共通倍率を除いて平方の比は符号となる。三つの非零の行成分の和が零で、
符号変更後も零なら、三符号は同一である。列にも同じ議論を適用して
平方の共通倍率を除き九係数の符号だけが残る。
36語の各四位置で用いた基底を記録する整数重み行列のSmith因子は
1が十一個、2が二個である。その整数左核の基底を各色対の符号へ制限すると、
非零となる条件は六つの非対角符号の積が1という一式だけである。
従って $N,P$ は共通倍率を除いて必要十分であり、六置換を加えれば一般の指定なし同型が決まる。
\path{computations/verify_walsh_presentation_space.py} は重み行列と符号補正を再現する。

一位置移動は $m_{be}\mapsto(-1)^{b_1(b+e)_0}m_{eb}$ であり、
六つの非対角成分で符号の積は $-1$ である。一般の $N$ はその転置の同時置換・定数倍ではない
（対角三値を一般に取れば置換と倍率を固定し、$N_{01}-N_{10}\ne0$ とする）。
二位置移動の符号 $(-1)^{d_1d_0}$ は、全位置の文字 $i$ を
$(-1)^{i_0i_1}$ 倍すると消える。実際 $b,e\ne3$ より $b_0 b_1=e_0 e_1=0$ である。
この基底変換は各係数が零であっても有効であり、全正則点で周期は2を割る。
\end{proof}

\begin{proposition}\label{prop:walsh-dual-square-cut}
三色ずつを持つ任意のWalsh候補で、その二次双対は、元の窓に各位置の
四つの欠色語 $x_{p,i}x_{p+1,i+3}$ を零とする関係を追加した商と、
矢の対角基底変更で一致する。
\end{proposition}
\begin{proof}
$D=\operatorname{diag}(1,-1,-1,1)$ と置く。隣接ブロックを
$M=\sum z_dP_d$ とすれば $M^{\mathsf t}DM=MDM^{\mathsf t}=(\sum s_d z_d^2)D$ である。
各非零ブロックは階数2で右辺は零なので、関係平面 $R$ に対し $R^\perp=DR$ となる。
行ブロック、列ブロック、転置ブロックとも同じ恒等式で処理できる。
欠色では $R=0$ のため双対関係は四語全体である。
$D$ は偶位置の矢だけの文字倍で全切断に同時に実現する。
右の文字に掛かる場合の余分な $s_b$ は色ごとのスカラーである。
\end{proof}

\begin{theorem}\label{thm:walsh-dual-frobenius}
三色ずつを持つ全てのWalsh候補の二次双対 $B$ は各始点でHilbert関数
$(1,4,6,4,1)$ のFrobenius窓である。特にこの枝の標準AS性に残る条件はKoszul性だけである。
$B$ で各矢 $x_{p,i}$ の右零化イデアルは $x_{p+1,i+3}$ の生成する右イデアルに等しく、
左側でも前位置の $i+3$ が零化イデアルを生成する。
\end{theorem}
\begin{proof}
二次双対の三次の双対空間は元の隣接関係を左右から挿入した空間の交叉である。
三文字の和 $h$ を固定すると、語は隣接二色 $b,e$ で一意に定まる。
交叉の元は $b,e\ne3$ に台を持つ。各行では全四成分が所定の関係平面に入り、
さらに欠色の一成分が零となる。この平面は座標超平面に含まれない
（非退化三次元対角形式の全等方部分空間の次元は高々1である）。
従ってその切断は直線で、実際に最後の矢を $h$ に固定した $W$ の縮約が各行でこの直線を張る。
列でも同様である。この既知の縮約の非零成分上で各行倍率と各列倍率が一致する。
台の二部グラフは各行・各列に少なくとも二辺を持つ。これは三非零色と
$N$ の和の消滅から従う。各側が三頂点なのでグラフは連結であり、全倍率は共通となる。
巡回したテンソルでも符号と転置だけが変わり、各始点で三次は各色一次元となる。

四文字の和を $h$ とすれば、任意の四次双対の双対テンソルは
\[
        f(i)W_{i+h,j,k,t}=g(t)W_{i,j,k,t+h}
\]
の両表示を持つ。$h\ne0$ で $i+j=3$ とすれば右辺は零である。
左辺の行 $3+h$ は非零色なので、任意の $i$ で $k,t$ を選んで係数を非零にでき、
$f(i)=0$ となる。$h=0$ では $t=i+e$ で各 $e=0,1,2$ を選んで
$f(i)=g(i+e)$ を得るので $f,g$ は共通定数である。従って四次は $W$ の直線に双対である。
五次に全色 $h$ を固定すると、最初の四文字に $W$ を持つ表示は最後の文字が $h$ に固定され、
後の四文字に $W$ を持つ表示は最初の文字が $h$ に固定される。
$W$ は一つの最初の文字に台を持たないため交叉は零である。
全ての完全対合は、$W$ の単文字切断の階数4、隣接切断の階数6と上記の次元から従う。

命題\ref{prop:walsh-dual-square-cut}の表示で $x_{p,i}x_{p+1,i+3}=0$ である。
他の三文字との二次積はそれぞれ異なる色で非零である。実際関係平面は全等方であり、
一つの座標語は非等方なのでその平面に入らない。よって各単文字の左倍、右倍は
次数0,1でそれぞれ階数1,3である。四次の完全対合を使えば次数2,3の階数はそれぞれ
3,1となる。ここで相手側の写像は四文字テンソルを巡回した位置における単文字の乗法である。
隣り合う欠色の文字の乗法の合成は零であり、全次数で一方の像と他方の核の次元が一致する。
これで左右の零化イデアルの記述を得る。
\end{proof}

% !TeX root = ../papers/nc-p3.tex
\section{Walsh三次元族の全点に対するAS正則性の必要十分条件}
\label{sec:walsh-normal-quadrics}
\begin{theorem}\label{thm:walsh-normal-quadrics}
標数零の代数閉体でWalsh基底を持つ隣接階数6の四窓を取る。
標準14641型quadratic Koszul左右AS正則であるための必要十分条件は、両切断の
非零色が三つずつあり、欠色を3にした符号付き平方行列 $N$ の階数が2であることである。
従って反対階数16を含む全ての特殊値で、この判定に高次の経路計算は要らない。
正則点では四つの正規二次式が欠色の語を張り、その二乗は共通の四中心四次式を与える。
各始点でそれらの多項式環上自由階数256となる。
\end{theorem}
\begin{proof}[階数1の排除]
欠色が二つある側は定理\ref{thm:walsh-symmetric-windows}の成長障害で除かれる。
$N$ は行和・列和が零なので、残る可能性は階数1,2だけである。
階数1で三行三列が非零なら全九成分が非零で、係数の平方は行倍率と列倍率の積になる。
平方根を取ると $m_{be}=u_b v_e$ の九成分の符号変更となり、
$u_0^2=u_1^2+u_2^2$, $v_0^2=v_1^2+v_2^2$ である。
重み行列の符号判定により $m=(u_bv_e)$ とその $(0,1)$ 成分だけを反転したものの二表示で覆われる。
両円錐曲線を $(1+t^2,1-t^2,2t)$ で有理表示し、二独立変数の有理関数体上で
四次に元の六関係を各位置で挿入する。余核の文字次元は両表示で $(12,8,8,8)$、
合計36である。\path{computations/verify_walsh_normal_quadrics.py} が有理関数体の
正確な階数を再現する。両既約族の全特殊化でも四次元数は36以上であり、標準値35に反する。
\end{proof}

\begin{proof}[欠色四語を一文字の左右で移す写像]
以後 $\operatorname{rank}N=2$ とする。$y_{p,i}=x_{p,i}x_{p+1,i+3}$ と書く。
任意の文字 $a$ に対し
\[
 y_{p,i}x_{p+2,a}=\sum_j T_{p,a}[j,i]\,x_{p,a}y_{p+1,j+3}
 \tag{*}\label{eq:walsh-square-transport}
\]
となる可逆行列 $T_{p,a}$ が一意に定まることを示す。
$h=a+3$ の三文字を $(h+e,h+b+e,h+b)$ と書く。元の三次経路の双対は
$4\times4$ の成分 $X_{be}$ のうち各非欠色行・列が対応する関係直交平面に入るものである。
その次元は5である。実際各側の関係挿入が6次元、交叉が
定理\ref{thm:walsh-dual-frobenius}で1次元、挿入前は16次元である。
両側の欠色成分が零である部分は一次元であり、$W$ の最後の文字を固定した縮約に
命題\ref{prop:walsh-dual-square-cut}の対角文字倍を施した $C$ が張る。

第3行全体を零にしたとき、各非欠色列は $X_{be}=c_e C_{be}$ に制限される。
各非欠色行を第3列まで延長する条件は $\sum_e N_{be}c_e=0$ と同値である。
実際その関係直交平面は $s_e$ を対角とする形式の全等方平面で、
第3列を忘れた射影は単射である。像は $C_{b,*}$ との重み付き直交条件で与えられる二平面である。
$C_{be}^2=m_{be}^2$（巡回すると転置の可能性もある）なので所定の式となる。
階数2では $c_e$ は全て同じで、延長は一意である。この議論は個別の成分では割らない。
従って第3行への制限は全射で核が $kC$ である。第3列への制限も同じ核を持つ全射である。
これを双対化すれば式\eqref{eq:walsh-square-transport}の両端は同じ四空間の二基底である。

全等方性から、各 $X$ で非欠色の行ごとに $\sum_e s_e X_{be}^2=0$、
非欠色の列ごとに $\sum_b s_b X_{be}^2=0$ である。
$\sum_{be}s_bs_e X_{be}^2$ を二通りに取れば二つの境界の二次形式は一致する。
境界の順序はいずれも文字に $a+3$ を加えたものなので、共通符号を除き
$D=\operatorname{diag}(1,-1,-1,1)$ を得る。$D^{-1}=D$ により $T_{p,a}$ もその等長写像である。
$T_p=T_{p,0}$ は両側に共通な語 $x_{p,0}x_{p+1,3}x_{p+2,0}$ を固定するので、
第0基底ベクトルを固定する。行列式は $-1$ である。これは階数2の既約開集合上の
局所正則な関数で二値しか取らないから、一点で決まる。
前掲コードは $m$ の三行を $(0,1,i),(i,0,i),(-1,1,0)$ とした点で
$T_0=\operatorname{diag}(1,1,-1,1)$、全始点で行列式 $-1$ を得る。
\end{proof}

\begin{proof}[四つの正規二次式]
四語の係数ベクトルを分裂四元数環 $M_2(k)$ の基底 $e_0=1,e_1,e_2,e_3=e_1e_2$ で表す。
$e_a e_i=(-1)^{a_1 i_0}e_{a+i}$ とし、余因子共役を $z\mapsto\bar z$ とすれば
行列式ノルムは上の $D$ である。$T_p$ は1を固定する向き反転等長写像なので
$T_p(z)=q_p\bar z q_p^{-1}$ と書ける。実際、余因子共役を合成したものは
跡零三空間の特殊直交群の元であり、$\mathrm{PGL}_2\simeq\mathrm{SO}_3$ の随伴作用で尽くされる。
全頂点のラベルを $a$ だけずらすWalsh作用を式\eqref{eq:walsh-square-transport}に施すと
$T_{p,a}=L_{e_a}T_p L_{e_a}^{-1}$ となる。左右の基底を $i,j$ とした符号比は
$(-1)^{a_1(i+j)_0}$ である。ここで $L$ は左乗法を表す。

次の位置との整合を確認する。三語 $(i,j,k)$ を逆順にして
$(-1)^{b_0 e_1}$ を掛けると、位置 $p$ の三次から位置 $p+1$ の三次への線型同型となる。
確かに辺色 $c$ の関係空間は文字を $c$ ずらして第1文字 $t$ に $(-1)^{c_1t_0}$ を
掛ける操作で保たれる。二位置移動後の各関係は両文字への $J(i)=(-1)^{i_0i_1}$ 倍で得られる。
$h=i+j+k$ を固定すると $b_0e_1$ は、それぞれ $e_1 k_0$ と $e_1h_0$ だけ、
$b_1j_0+i_0i_1+j_0j_1$ と $b_0h_1$ だけ異なる。
これは各辺を逆にする際の必要な符号と、辺色だけの無害な倍率である。
境界の二基底を逆順にすれば
\[
 T_{p+1}=\theta_{\rm hi}\,T_p^{-1}\theta_{\rm lo},
 \qquad q_{p+1}=e_1q_p^{-1}e_2
\]
と選べる。ここで $\theta_{\rm hi}$ と $\theta_{\rm lo}$ は該当ビットの文字倍、
すなわち $e_1,e_2$ の各共役である。二回適用すると $q_{p+2}=e_3q_pe_3$、
四回で元に戻る。

係数ベクトル $e_c q_p$ の定める二次式を $u_{p,c}$ とする。式\eqref{eq:walsh-square-transport}
で次の位置の実際の係数順に戻す移動 $j\mapsto j+3$ は $z\mapsto e_1 z e_2$ である。
四元数の積則から
\[
 u_{p,c}x_{p+2,a}=\kappa_c\det(q_p)\,\phi_c(a)x_{p,a}u_{p+1,c},
 \quad
 \phi_c(a)=(-1)^{[a\ne0]+a_0c_1+a_1c_0},\quad
 \kappa_c=(-1)^{[c\ne0]+c_1}.
\]
実際 $q_p\overline{e_c q_p}=\det(q_p)\bar e_c$ を使い、
$e_a\bar e_c e_a/\det(e_a)=\phi_c(a)\bar e_c$ とすればよい。
$\det(q_{p+1})=\det(q_p)^{-1}$ より四位置で定数倍率の積は1である。
$u$ を位置ごとに規格化して係数を $\phi_c(a)$ のみにできる。
各 $q_p$ が可逆なので、四式は各位置の欠色空間の基底である。

折り畳んだ環で $u_c=\sum_p u_{p,c}$ と取る。矢を二位置ずらし
$\phi_c$ を掛ける対合で交換するため $u_c^2$ は中心である。
また $u_c$ と $u_d$ は非零スカラーで位置ごとに交換する。
実際 $\phi_c(i)\phi_c(i+3)$ は $s_i$ に一定符号を掛けたもので、
係数ベクトルへの $s_i$ 倍は $e_3$ の共役である。
従って $e_d q_p$ の直線は $e_d q_{p+2}$ の直線に移る。
\end{proof}

\begin{proof}[全次数の自由性とKoszul性]
四つの $u$ を全て割った商は定理\ref{thm:walsh-dual-frobenius}の
$(1+t)^4$ の窓である。各始点でその基底語を持ち上げる。
次数帰納と上の交換式により、順序付けた $u_c$ の単項式と持上げ語で全経路を張る。
各 $u_c^2$ を中心変数に移せば、中心四変数上の生成系の母関数の上界は
$(1+t)^4(1+t^2)^4=(1+t+t^2+t^3)^4$ となる。
各開始位置で $u$ の個数に応じて語の開始位置が移るが、各位置の商基底の次数数は同じである。
中心変数を含めた上界は $(1-t)^{-4}$ である。
一方この既約族の幾何学的一般点は正則なので、各次数の経路次元は全特殊化で
その標準値以上である。全次数で上下界が一致し、形式的中心変数上の自由加群からの
全射は同型である。特に中心四変数は独立である。
有限Frobenius双対と標準Hilbert級数に加え、中心多項式環上自由であるから、
定理\ref{thm:klein-rational-period-four}の左右線型複体の証明が適用できる。
両側で中心の最初の二変数が正則なので単純から環への $\operatorname{Hom},\operatorname{Ext}^1$
が消え、表示の完全性とその双対で両端を、Euler等式で中央を得る。
従って全次数でKoszul左右AS正則となる。
\end{proof}

% !TeX root = ../papers/nc-p3.tex
\section{十五次元成分内のWalsh三次元族とその外の周期層}
\label{sec:walsh-full-component}
\begin{theorem}\label{thm:walsh-full-component}
Walsh三次元族の全体は、添字保持同型を除いた次元が15の
定理\ref{thm:full-fifteen-component}の既約成分に属する。
一般のWalsh点で隣接階数スキームは滑らかである。指定なしの成分内で、
この点を通る周期が2を割る軌跡の被約局所枝は四つあり、次元は $9,5,5,5$ である。
各枝の幾何学的一般点は標準Koszul左右AS正則で最小周期2を持ち、
Walsh族の外にある。残る一般の変形は最小周期4である。
\end{theorem}
\begin{proof}[二つの表示族の交点の滑らかさ]
$a=1,b=3,r=4$, $l^2=-10$, $p=5$, $(m')^2=-1$, $q^2=17$, $n^2=-9$ とし、
53の上で $(l,m',q,n)=(19,23,21,16)$ となる根を選ぶ。
階数15の復元式では $(s,t,N')=(25,35,28)$ に還元される。通常な正の固定曲線の第一表示で
倍率 $(z,\lambda)=(15,38)$、枠 $(a',b',e',f')=(29,1,26,17)$ となる。
定理\ref{thm:klein-elliptic-completion}の二次式 $P,Q$ の先頭係数は剰余で
$4,36$、選んだ根での導関数は $7,-7$、二枠の行列式は $31,51$ である。
従ってこの標数零の窓は同定理の既約な二重被覆の対根表示に実際に含まれる。
分母が単元で、根を選ぶモニックな二重被覆の一般点をこの点へ特殊化できる。
従って既知の十五次元成分上のWalsh点である。

この点の第一辺は53の上で二十個の横断交点を持つ。四座標点と、
$(1,3,10,2)$ に四ラベルの平行移動・文字倍を施した十六点を使う。
各点で次の点の四座標は一次元核で決まり、六双一次式のJacobi階数は6である。
Segre多様体の次数は20なので、この二十個の孤立単純交点が全交叉を尽くす。
実際、線型切断を順に取ったときの全既約成分の次数の和は増えず、
正次元成分があればその次数も加算される。各点での二つの部分評価は階数3で、
得た核直線を用いる定理\ref{prop:twenty-conormal}の二十行は独立である。
\path{computations/verify_walsh_full_component.py} はこの計算と以下の接計算を再現する。
有限エタールな二十点と単元小行列によってこれらの条件は標数零に持ち上がる。

二切断の接方程式は、それぞれ摂動の両側に十次元の左右核を掛けて得る。
位数16の $G$ の文字を $(h,c)$ と書く。$h$ は語の総色、$c$ は同時移動に対する文字で、
$g=(x,z)$ は符号 $(-1)^{h\cdot z+c\cdot x}$ で作用する。
53での接空間の文字次元は $h=0$ の四文字で各4、$h\ne0$ の十二文字で各5である。
従って接方程式の階数は180以上で、二十恒等式により180以下でもある。
命題\ref{prop:twenty-conormal}で非射影化したスキームは滑らかな76次元となる。
文字次元の上界の合計が76なので個々の値も標数零で一致する。
滑らかな点では成分が一意であるから、この点を通る既約なWalsh族の全体が同じ成分に入る。
その一般点でも滑らかさと文字次元は保たれる。

アフィン軌道の接空間では自明文字は倍率の一つ、各非自明文字は四位置の行列から一つずつである。
実際各位置のPauli共役表現は行列空間で十六文字を一度ずつ持ち、
一般Walsh点の射影安定化群は有限である。
従って横断する十五次元空間は、自明表現三つと
$h\ne0,c\in(\mathbb Z/2)^2$ の十二の文字表現に分かれる。
\end{proof}

\begin{proof}[二位置移動の固定部分]
一般のWalsh点は幾何学的に安定で射影安定化群は $G$ である。
実際、正規直交基底に関するノルム最小性と安定化群の有限性は
定理\ref{thm:walsh-presentation-space}で確認した。二位置移動に四つの
対角変換 $J(i)=(-1)^{i_0 i_1}$ を合成して、表示点を固定する対合 $\tau$ を得る。
$\tau$ は四基底変更群を正規化する。安定軌道のエタール切片定理を、
二位置移動も加えた有限拡大の群に適用する
（D.~Luna, \emph{Slices étales}, Mémoires de la S.M.F. 33, 81--105, 1973）。
滑らかな横断切片の形式座標は有限安定化群で平均して線型化できる。
元の指定なし商の局所模型は上の十五次元空間の $G$ 商で、
その二位置移動には $G\rtimes\langle\tau\rangle$ の作用が誘導される。

$\tau$ の共役は $G$ の $(x,z)$ を $(x,z+x^\flat)$ に移す。
ここで $x^\flat=(x_1,x_0)$ であり、これは $J(i+x)/J(i)$ の線型部分である。
従って $h\ne0$ では $\tau$ が $(h,c)$ と $(h,c+h^\flat)$ の一次元方向を交換する。
各対の基底を選べば係数1の交換にできる。
三つの自明方向はWalsh族の接方向で、族の表示全体を $\tau$ が固定するので作用は恒等である。

有限商での移動不動点は、切片上の各 $g\tau$ の固定線型部分空間の像の合併である。
$g=(x,z)$ とし、各対の二座標を $A,B$ と書く。
$x\cdot h^\flat=0$ なら $B=(-1)^{h\cdot z+c\cdot x}A$、
そうでなければ両方零となる。$x=0$ なら六対の全部が残り、
各 $h$ で二対の符号が同じである。これら四つの九次元空間は $G$ の共役で一つの枝になる。
$x\ne0$ なら残るのは $h=x$ の二対だけであり、両対の符号は互いに逆になる。
二符号を同時に反転する二つの五次元空間は $G$ で同一視される。
三つの $x\ne0$ は異なる文字の支持を持つので互いに同一視されない。
また逆の二符号を持つ五次元空間は、同じ二符号を持つ九次元空間のいずれにも含まれない。
これで被約局所枝が尽きる。有限群の線型空間の軌道から来る各枝は
形式的にも既約である（その有限正規化の該当点は一軌道である）。

安定商の分離性により一位置移動の不動点は閉で、一般Walsh点はそこに属さない。
各局所枝の幾何学的一般点の全次数の正則性は、正則な特殊点から
有限Frobenius双対の内部次数別bar複体を持ち上げて得る。
各枝はWalsh三次元より大きく、成分の一般点には二位置移動も残らない。
\end{proof}

% !TeX root = ../papers/nc-p3.tex
\section{Walsh商の正規化と特殊な指定なし同型の有限復元}
\label{sec:walsh-spectral-fibers}
\begin{theorem}\label{thm:walsh-spectral-fibers}
定理\ref{thm:walsh-presentation-space}の平方行列の二重被覆の $S_3$ 商を $Y$ とする。
$Y$ は、Walsh表示の射影不変式商における閉像の正規化である。
標準正則像の逆像はちょうど $\operatorname{rank} N=2$ である。
全特殊点を含め、一つの窓の三つの切断の特異値から、対応し得る正規化点を有限リストに復元できる。
このリストに対する軌道閉包の線型比較で、指定を忘れた同型類の全ファイバーと最小周期が決まる。
同じ手順で任意の四次テンソルについて正則Walsh表示の有無も判定できる。
\end{theorem}
\begin{proof}[全ファイバーの有限性]
まず複素数体を使う。Walsh基底を正規直交としたテンソルは全一位置のGram行列が
スカラーで、ノルム最小ベクトルである。従って全ての非零表示が半安定であり、
同じ射影商の点を与えることとユニタリな四基底変更・共通倍率で移ることが同値である
（Kempf--Nessの最小ベクトル定理）。特に三切断の特異値を共通規格化して比較できる。
$\sum|m_{be}|^2=1$、すなわちテンソルのノルムを2に取る。

二つの隣接切断の各非零色のブロックは二次元行列表現の二重である。
四つの符号付き置換はユニタリな有限群の作用なので、直和分解もユニタリに取れる。
従って各ブロックには同じ正の特異値が二つあり、その平方は行または列の二乗和の2倍である。
零の行・列も許して三つの行和 $R_i=\sum_j|m_{ij}|^2$ と三つの列和 $C_i$ を
得る。添字順は有限通り全て取る。反対切断は各色でユニタリな対角符号と巡回行列の積である。
四つの単独枠と、三対の二重枠に十六個の特異値平方を割り当てる。単独枠は
\[
 L_\chi=\left|\sum_{i=0}^2(-1)^{i\cdot\chi}m_{ii}\right|^2
\]
であり、対 $ij=01,12,20$ の二つの値は
$L_{ij}^\pm=|m_{ij}\pm m_{ji}|^2$ でそれぞれ二度現れる。重複値、零も重複度を保って割り当てる。

$S_{ij}=(L_{ij}^++L_{ij}^-)/2$ とすると
\[
 |m_{ii}|^2=\tfrac12(R_i+C_i-\sum_{j\ne i}S_{ij}),\qquad
 \Re(m_{ii}\overline{m_{jj}})=\tfrac18\sum_\chi(-1)^{(i+j)\cdot\chi}L_\chi
 \quad(i\ne j).
\]
よって三つの複素対角成分を平面実ベクトルとしたGram行列が定まる。
その実現は非負で階数2以下のときに限り、共通複素位相を除けば共役を含め高々二通りである。
$d_i=m_{ii}^2$ とし、巡回順の対 $ij$ に対し $U_{ij}=d_k-d_i-d_j$（$k$ は残る添字）、
$J=N_{01}-N_{10}$ と置く。行列の和の条件から
\[
 N_{ij}=\tfrac12(U_{ij}+J),\qquad N_{ji}=\tfrac12(U_{ij}-J).
\]
三対の積 $L_{ij}^+ L_{ij}^-$ は全て $|J|^2$ である。この一致を確認する。
$S_{ij}=(|U_{ij}+J|+|U_{ij}-J|)/2$ なので
\[
 \Re(U_{ij}\bar J)^2=S_{ij}^2\big(|U_{ij}|^2+|J|^2-S_{ij}^2\big).
\]
$U_{ij}\ne0$ が一つでもあれば、その値と $|J|$ が決まっているため
$J$ は高々四点である。二乗で生じた余分な候補は実際の特異値比較で捨てる。
全て零なら $d_0=d_1=d_2=0$ で、自由な共通位相を使って $J=|J|$ に取れる。
こうして候補 $N$ は必ず有限個となる。各 $N$ では $P^2=\prod_{i\ne j}N_{ij}$ の根を取り、
$S_3$ と共通倍率を除く。前節の符号判定は零成分がある場合にも使える。
非対角成分が零なら未使用の符号で積を調整でき、全平方根表示が同値となる。
\path{computations/verify_walsh_spectral_fibers.py} は逆公式と重複固有値の割当てを実装し、
零対角と実でない $J$ の両方を検算する。
\end{proof}

\begin{proof}[正規化と最終比較]
九平方根の射影表示の座標環で許される対角符号を割ると、不変単項式は
全指数が偶数のものと、六非対角の指数だけが奇数のものに尽きる。
有限群の不変部分を取る操作は標数零で完全であり、
平方の五つの線型関係を課した後も被覆は $P^2=\prod_{i\ne j}N_{ij}$ である。
この二重被覆は相異なる六平面に沿う平方自由な分岐を持ち、正規である
（超曲面であり、余次元1では正則）。その有限 $S_3$ 商も正規である。
全非零表示で半安定なので、射影商への写像は $Y$ 上の射として定まる。
先の有限復元により全幾何ファイバーが有限で、射影的準有限射ゆえ有限である。
一般点では定理\ref{thm:walsh-presentation-space}の同型判定で一対一なので標数零で双有理であり、
正規化である。階数2なら三行三列は全て非零であり、正則像の記述は
定理\ref{thm:walsh-normal-quadrics}から従う。これはWalsh表示を持つ像の中での主張である。

有限リストの真のファイバーを選ぶため、任意の二テンソル $U,V$ に対して使える
基底未知数を消去した線型比較を記す。$d$ 次の対称冪で $U^d$ から始め、
四位置の行列単位をLie代数として反復して作用させ、線型包が増えなくなるまで取る。
得る空間 $L_d(U)$ は連結群 $\mathrm{GL}_4^4$ の軌道の $d$ 乗の線型包である。
標数零であるからLie代数不変部分空間はこの連結群で不変となる。
その零化空間は射影軌道閉包の $d$ 次の斉次方程式に等しい。
従って低次数から二空間を比較して違えば直ちに除ける。
終了を保証するには両軌道閉包の斉次イデアルをパラメータ表示から消去して比較してよく、
線型演算のみを使う場合は一様な有限上界 $d=4^{64}$ も使える。

実際、全行列による軌道写像は64変数で次数4である。写像のグラフに64個の一般線型条件を
課すとBézoutの次数上界は $4^{64}$ であり、線型射影した閉像の次数もこれ以下である。
可逆行列は稠密なのでその射影閉像が所望の軌道閉包である。
次数 $D$ の射影既約多様体の外の点を分離するには高々 $D$ 次の方程式で足りる。
次元を $r$ とすると、その点を通る一般の余次元 $r+1$ の線型空間は多様体と交わらず、
その中の超平面から $\mathbb P^{r+1}$ へ射影すれば、像の超曲面の次数は高々 $D$ で
対象点の像を含まない。方程式を引き戻せばよい。次数を上げるにはその点で非零の線型式を掛ける。
ゆえに上界で二つの $L_d$ が等しければ閉包が等しい。
各軌道はその既約閉包に稠密な開集合なので、閉包が等しいことと同型は同値である。
この最悪上界は実際の計算量の小ささを主張するものではない。

入力が任意のテンソルなら、まず四位置のGram行列がスカラーとなる代表を求める。
複素行列の実・虚部について、この条件と行列の可逆性は実代数的であり、
実閉体上の消去で解の有無と一つの代表を決定できる。代表が存在しなければWalsh軌道ではない。
存在すれば上の三つのスペクトルから有限個の $N,P$ を作り、実際の表示テンソルと
$L_d$ または軌道イデアルを比較する。隣接切断の必要な重複度も最初に照合する。
標数零の他の代数閉体上でも、入力の有限生成係数体を複素数へ埋め込んで比較できる。
同型の有無は有限型方程式の幾何学的可解性であり、選んだ埋込みや共役の選択に依存しない。
正則な分岐内の最小周期は、実際のファイバーが
$(N^{\mathsf t},-P)$ を含めば1、それ以外は2である。
\end{proof}

% !TeX root = ../papers/nc-p3.tex
\section{Walsh族を通る二つの大域的な半周期成分}
\label{sec:half-period-components}
\begin{theorem}\label{thm:half-period-components}
十五次元成分の安定商で、一般Walsh族を含む二位置移動の被約不動軌跡の既約成分は
九次元と五次元の二つである。前者の一般点の射影安定化群は自明、後者では位数2である。
後者の三つの局所枝はWalsh族に沿って全対称群 $S_3$ で置換される。
その正規化は一般Walsh点の上に三点を持ち、各枝はそこで滑らかである。
前者の正規化は一般Walsh点の上に一点を持つが、その近くに正規化の二点が同定される
五次元の部分が三つある。
\end{theorem}
\begin{proof}
前節の線型切片を使う。$x=0$ の九次元空間の一つを $L$ とし、六対の交換符号を全て正に取る。
この空間を保つ群は $D=\{(0,z)\}$ である。
各 $h\ne0$ の二つの独立座標に、同じ $(-1)^{z\cdot h}$ が働く。
全座標が非零なら、二位置移動を含まない安定化群は自明である。
$L$ の軌道の有限商の正規化は $L/D$ となる。これは有限な双有理写像で、有限線型群の商は正規である。
$G$ のうち $L$ を保たない元が $L$ の点を再び $L$ に送るには、その点の
非自明な文字の支持がただ一つの $h$ に含まれる必要がある。
実際、移動部分 $y\ne0$ は $y\cdot h^\flat\ne0$ の対の符号比を反転する。
$h=y$ の二座標を $u,v$ とすると、$\{(y,z)\}$ は同時符号のほかに相対符号を反転する。
よって $uv\ne0$ では $L/D$ の二点が像で同一となる。各 $h$ につきこの部分は
三自明方向と二座標の五次元である。六座標が零のWalsh点上には正規化の一点だけがある。

$x\ne0$ の局所枝では、$h=x$ の二対だけを使い一方の符号を逆にする。
両対が非零なら実際の安定化群は $\{1,(0,x^\flat)\}$ である。
持ち上げた二位置移動 $g\tau$ の二乗はその非自明元で、持上げの射影位数は4である。
固定した $x$ の一方の線型空間を保つ群は移動成分が $0,x$ の元である。
二座標への像は二符号の独立な変更全体なので、その商は三自明座標と $u^2,v^2$ の滑らかな空間になる。
他方の共役空間との交わりは三自明座標のみである。二つの線型空間の関数環は
その交叉上で一致する関数の対であり、交叉に働く共役作用は恒等であるから、
両者を合わせて $G$ で割ってもこの一つの滑らかな商を得る。
異なる三つの $x$ の支持は互いに異なり、交叉はWalsh方向である。

大域的には、欠色を固定した既約Walsh表示上で一定の $g\tau$ を固定する。
一般の表示点の近くでは滑らかな周囲空間への有限位数自己同型の不動点は滑らかである。
従って一般表示を含むこの固定点の既約成分は一つで、その基底変更飽和の閉包は既約である。
四基底を動かした場合も新しい局所成分は増えない。切片での有限個の固定空間に、
軌道内の $g\tau$ 不変方向を加えれば表示の固定スキームの接空間であり、
飽和して商に移せば対応する線型切片の全次元が出る。
Walsh表示を保つ基底変換は定理\ref{thm:walsh-presentation-space}で
移動ラベルの全 $\mathrm{GL}_2(\mathbb F_2)$ を実現する。
そのような変換 $n$ は $\tau$ を $g'\tau$ に共役し、
$g'$ は一般Walsh表示を固定するので $G$ に属する。
共役後も位数2であることから $g'$ の移動部分は零である。
ゆえに三つの非零 $x$ は全置換で移る。各 $x$ の二つの符号も既に $G$ で共役なので、
五次元の全ての局所枝の大域的な飽和は一致する。九次元側も同様に一成分であり、
二つの次元と局所非包含から両者は不動軌跡の別の既約成分である。
三つの標識は一般Walsh表示上で一定で、その指定を忘れる有限被覆の変換が
全六置換を与えるので、五次元側の正規化の三点の一般モノドロミーは $S_3$ である。
\end{proof}

\begin{proposition}\label{prop:half-period-pencils}
五次元成分の一般点では、二位置移動の位数4の持上げを取ると、各隣接切断は
八次の対称行列の階数4部分と八次の交代行列の階数2部分に分かれる。
交代部分の二つの対角ブロックの像は $2\times2$ 行列の二平面で、各々に
階数1の線が相異なって二本ある。これらは全頂点の四直線の基底枠を指定なしで有限通りに決める。
九次元成分の一般点では二位置移動の持上げは位数2で、両隣接切断を対称行列にできる。
\end{proposition}
\begin{proof}
二位置移動で対向する空間を同定する。五次元枝で持上げの二乗は二対二の固有空間を持つ
二つの射影対合である。偶、奇の符号を $p_i,q_j$ とすると、第一切断 $M$ は
$M_{ij,kl}=p_i q_j M_{kl,ij}$、$p_i q_j=p_k q_l$ に台を持つ形にできる。
第二切断は $N_{kj,il}=p_i M_{ij,kl}$ とすれば、同じ性質を $p_k q_j$ について持つ。
この符号の型は基点の実際の持上げを以下で確認すれば近傍でも一定である。

Walsh点で持上げを $g\tau$、$z=0$ とし、偶・奇四次行列を $A,B$ とする。
文字 $i$ の列をそれぞれ $i+x$ へ運び、$J(i)(-1)^{x_1 i_0},J(i)$ を掛ける。
第一切断を $M=W(A\otimes B)^{\mathsf t}$ に取れば
$A^2$ の文字は $x^\flat$、$B^2$ の文字も共通倍率 $J(x)$ を除いて $x^\flat$ である。
よって反対称側に来る非欠色 $b$ は各非零 $x$ で一つだけであり、切断の階数は順に4,2である。
階数の和は6なので、この二下界は近傍で等号のまま保たれる。
二つの負符号ブロックの内部の四次交代行列について、
\path{computations/verify_half_period_components.py} は先の53の点で
各像が二次元、射影鉛筆の行列式に二単根があることを三つの持上げ全てで確かめる。
単根の二行列の像と核はそれぞれ独立な二線である。
負符号空間は $E_+\otimes F_-$ と $E_-\otimes F_+$ の直和なので、これで全四空間の二分解を
一度ずつさらに二分できる。一般点では安定化群が位数2で、二つの持上げは互いに逆であるため
これらを両方取っても有限個の枠しか生じない。九次元の $x=0$ の場合は
$A,B$ が対角の符号行列で、双方の二乗とも1である。
同じ同定で全体が対称となり、族の対応する固定スキームでこの形が保たれる。
\end{proof}

% !TeX root = ../papers/nc-p3.tex
\section{位数4の半周期成分の二倍率と二つの二次式}
\label{sec:half-period-two-quadrics}
\begin{theorem}\label{thm:half-period-two-quadrics}
定理\ref{thm:half-period-components}の五次元成分は、有限な表示の選択を許せば、
二つの倍率上の $(2,2)$ 完全交叉三次元多様体の族で稠密開部分が覆われる。
倍率は $h^2-k^2=j^2-l^2=1$ とし $t=h^2-j^2\ne0$, $L=j^2+k^2$ とする。
四次元空間の基底で、以下の順に対角値を取る。
\[
 (d_i)=t^{-1}(jh,jh,-lk,-lk),\qquad
 (\lambda_i)=(jh,-jh,lk,-lk).
\]
平面 $U$ のPlücker座標 $p_{im}$（$1\le i<m\le4$）に対する方程式は
\[
 p_{12}p_{34}-p_{13}p_{24}+p_{14}p_{23}=0,\qquad
 \sum_{i<m}d_i d_m(\lambda_i+\lambda_m-L)p_{im}^2=0 .
\]
一般倍率で交叉は滑らかである。$\sum d_i d_m p_{im}^2\ne0$ 等の以下の開条件で、
平面の直交射影から元の四次テンソルを復元できる。
\end{theorem}
\begin{proof}[二つの交代部分の規格化]
前節の位数4の持上げを選び、$E_\pm,F_\pm$ の各二空間の基底を取る。
$E_+\otimes F_-$ と $E_-\otimes F_+$ 上の交代鉛筆の二つの階数1直線を
それぞれ $00,11$ に移す。第一切断の交代対角ブロックは
$00\wedge11$ の非零倍 $\alpha,\beta$ となる。
交代非対角ブロックは二平面間の $2\times2$ 行列であり、全交代部分が階数2である条件は
Pfaffianの消滅、すなわちその行列式が $\alpha\beta$ となることである。
第二切断 $N_{kj,il}=p_i M_{ij,kl}$ の二鉛筆は、第一切断の対応するブロックの
部分転置に符号を掛けたものなので、その像は $01,10$ の平面である。

四添字 $a=(a_1,a_2)\in\{0,1\}^2$ に $C a=(1-a_1,1-a_2)$ と $Ba=(a_2,a_1)$ を作用させ、
同じ記号を置換行列にも用いる。第一・第二の対称な八次ブロックの非対角 $4\times4$ 行列を
$H,J$ とする。鉛筆の台により $H$ の行 $a$ の非零位置は $CBa$、$J$ のそれは $Ba$ である。
非零値を $h_a,j_a$ と書く。第一の交代非対角行列は $(j_{uv})$、
第二は $(-h_{u,1-v})$ となる。第二の交代対角ブロックの係数は $-\alpha,\beta$ なので
\[
 h_{00}h_{11}-h_{01}h_{10}=j_{00}j_{11}-j_{01}j_{10}=\alpha\beta .
\]
逆に二つの交代部分はこのデータで全て復元できる。

$E_+,F_+,E_-,F_-$ の基底倍率を順に $x_i,y_i,u_i,v_i$ とすると、
$h_{ab}$ には $x_a v_{1-a}y_bu_{1-b}$、$j_{ab}$ には $x_a v_a y_bu_b$ が掛かる。
従って $\alpha=\beta=1$ の後でも二つの二行行列をそれぞれ独立に
行列式1の対角行・列倍率で規格化できる。実際 $x_0v_0,x_0v_1,y_0u_0,y_0u_1$ は
自由に選べ、相手添字の値は各々の逆となる。
全成分非零で $h_a=(h,k,k,h)$, $j_a=(j,l,l,j)$ にできる。
残る連続倍率は、作用しない倍率を除き二つの対称対角ブロックを互いに逆に拡大するものだけである。

これらの第一対角ブロックを $S,T$、$\Gamma$ を $2\times2$ 因子の第一位置での部分転置と書く。
$K=\operatorname{diag}(1,-1,-1,1)C$ とし、添字を $0,1,2,3$ の順にすれば
$G^\Gamma=G+\delta(G)K$、$\delta(G)=G_{12}-G_{03}$ が全対称行列で成立する。
両対称部分の階数4条件は
\[
 T=H^{\mathsf t}S^{-1}H,\qquad
 T^\Gamma=-J^{\mathsf t}(S^\Gamma)^{-1}J .
\]
両逆行列と $\delta(S)$ が非零の開部分を扱い、最後の倍率で $\delta(S)=1$ とする。
\end{proof}

\begin{proof}[対合の固有空間からの復元]
$\delta(T)=t=h^2-j^2=k^2-l^2$ と取る。この条件を課した族が五次元成分を支配することは
下で次元と交点から確認する。四次元空間とその双対の直和で、以下の写像を順に施す。
\[
 (x,y)\mapsto(H^{-1}y,H^{\mathsf t}x),\quad
 (x,y)\mapsto(x,y+tKx),\quad
 (x,y)\mapsto(J^{-\mathsf t}y,-Jx),\quad
 (x,y)\mapsto(x,y-Kx).
\]
合成 $F$ が $S$ のグラフを保てば二つのSchur条件が成立する。
$D=t^{-1}\operatorname{diag}(jh,-lk,-lk,jh)$、
$A=\operatorname{diag}(h/j,k/l,k/l,h/j)$、$R=-DAC$ と置く。
$F$ は反シンプレクティック対合で、その $\pm1$ 固有空間は $R\pm D$ のグラフである。
具体的に上側の二ブロックは $J^{-\mathsf t}H^{\mathsf t}$, $tJ^{-\mathsf t}KH^{-1}$、
下側はそれらへの左 $-K$ 倍と右下への $-JH^{-1}$ の和である。
$R\pm D$ のグラフに掛けて主張が従う。

正固有グラフに載せる水平平面を $U$ とする。負固有グラフに載せる平面を
$U^{\perp_D}$ とすれば両者の和はLagrange部分空間である。
$D|_U$ が非退化ならこれはグラフとなり、その対称行列は
\[
 S=R+D(2P_U-I),\qquad P_U=U(U^{\mathsf t}DU)^{-1}U^{\mathsf t}D .
\]
ここでは右辺の $U$ は平面の任意の列基底である。
$K$ によるトレースで $\delta(S)=1$ を書くと、
$\operatorname{tr}(D_{jh} C P_U)=L$ となる。ただし
$D_{jh}=\operatorname{diag}(jh,lk,lk,jh)$ である。
各相補添字対の和・差を $\sqrt2$ で割った基底で $D_{jh}C$ と $D$ は同時対角となり、
固有値は表記の $\lambda_i,d_i$ である。Cauchy--Binetの公式から
$\operatorname{tr}(D_{jh}CP_U)$ は
$\sum_{i<m}(\lambda_i+\lambda_m)d_i d_m p_{im}^2/\sum_{i<m}d_i d_m p_{im}^2$
であり、これが第二の二次式を与える。

$\delta(T)=t$ からは独立な方程式は増えない。実際、最初の写像は二つの固有グラフを
$t(R\mp D)$ のグラフに移し、正側の水平平面は $U'=H^{-1}(R+D)U$ となる。
よって $\delta(T)=t$ は $\operatorname{tr}(D_{jh}CP_{U'})=-L$ に等しい。
$U'$ の各固有座標の二乗の倍率は共通倍率を除いて
\[
 (z_i)=\big((j-h)^2,(j+h)^2,(l-k)^2,(l+k)^2\big)
\]
となる。置換 $CB$ は各固有線を保ち、そこでの追加符号は二乗で消える。
各 $i<m$ で
\[
 (\lambda_i+\lambda_m+L)z_i z_m=-t^2(\lambda_i+\lambda_m-L).
\]
同じ相補対なら $z_i z_m=t^2$、異なる対なら
$2L=(h^2+j^2)+(k^2+l^2)$ と $(j\pm h)(j\mp h)=(l\pm k)(l\mp k)=-t$ から従う。
従って両方のトレース条件はPlücker座標で同じ二次式となる。

二つの二次形式の鉛筆では、相補な二つのPlücker添字ごとに $2\times2$ に分かれる。
六つの判別根の平方は共通非零定数を除き
$L^2$, $L^2-(jh+lk)^2$, $L^2-(jh-lk)^2$ である。
一般にはこれらが相異なる非零値なので六根が単純であり、完全交叉は滑らかである。
滑らかな三次元完全交叉は連結ゆえ幾何学的に既約である。
\end{proof}

\begin{proof}[成分の被覆]
Walsh係数を行順に $a,b,c,d,e,f,g,v,w$ と書く。
持上げ $x=1$ で二鉛筆は既に対角線を持ち、規格化前の $H,J$ の重みは
$(c,-f,-f,c),(-v,g,-g,v)$ となる。$S$ の対角は $(b,d,d,b)$、
相補対の非対角は $a,e$、$T$ の対応値は $-b,d,d,-b$ と $a,-e$ となる。
従って $\delta(S)\delta(T)=a^2-e^2=c^2+v^2=f^2+g^2$ は符号平方の行・列和からの恒等式であり、
上で課した条件 $st=h^2-j^2$ はこのWalsh族上で成立する。
規格化前でも $F$ のグラフへの制限のトレースは零である。
上左ブロックの対角が零で、上右対角ブロックの相補二成分の符号が反対、
$S$ の該当対角値が等しいためである。従って固有次元は2ずつであり、
上の平面表示に入る。

53の共通点で $h^2,k^2,j^2,l^2$ は規格化後にそれぞれ $44,43,37,36$ であり、
各逆と六根の分離条件は全て単元である。
\path{computations/verify_half_period_two_quadrics.py} はこの点の固有次元と
一般倍率のグラフ・二次式の恒等式を検算する。
これで表示空間の一般ファイバーには必要な開条件を満たす点があり、総次元は5である。
二鉛筆の二根、残る四乗根倍率、持上げとその逆を除き基底は規格化されている。
$\delta(S)=1$ により有効に残る連続倍率もなく、$S$ から平面 $U$ は正固有グラフとの交叉で回収できる。
一般Walsh点の安定化群は有限なので、その近くでは持上げを忘れる繊維も有限である。
従ってこの既約表示の像は五次元で、一般Walsh点を通る位数4持上げの成分に属する。
定理\ref{thm:half-period-components}の一意性で、その大域的五次元成分を支配する。
\end{proof}

% !TeX root = ../papers/nc-p3.tex
\section{良い還元による対称性を仮定しない閉点の全次数認証}
\label{sec:unrestricted-good-reduction}
\begin{proposition}\label{prop:unrestricted-good-reduction}
離散付値環 $V$ 上の四つの自由な階数4の空間のテンソル $W$ を取る。
分数体上の二つの隣接切断が階数6で、幾何学的剰余体上の窓が
標準14641型quadratic Koszul左右AS正則なら、幾何学的分数体上でも同じ性質が成立する。
各始点、各次数の元の経路も二次双対の経路も $V$ 上自由である。
特に剰余標数が奇数で、剰余窓に三色ずつのWalsh表示があり符号平方行列が階数2なら適用できる。
\end{proposition}
\begin{proof}
特殊点で各隣接切断の階数は6である。従って単元な六次小行列があり、
分数体で七次小行列が全て零なので、関係像は直和因子で剰余と可換である。
四位置の直和を取った二次双対を $B$ とする。各経路加群は有限表示である。
三次にはテンソルの一文字縮約から双対次元の下界4があり、四次にはテンソル自身から
下界1がある。実際、一文字切断の非零四次小行列は特殊点から持ち上がり、
各縮約は左右の関係挿入の交叉に属する。四次でも三位置の交叉に $W$ が入る。
特殊点で両下界が等号なので分数体でも等号である。五次は剰余で零ゆえ中山の補題で零であり、
以後も生成次数から零である。各次数で剰余次元と分数体次元が等しいことから
$B$ は次数ごとに自由で、最高次への左右の完全対合は単元行列式により持ち上がる。

半単純な頂点環の相対版として、$V^4$ に関する $B$ の被約bar複体を取る。
内部次数 $n$ の項は正次数の経路加群のテンソル積の有限直和であり、有限自由である。
剰余でKoszul性によりホモロジー次数 $i\ne n$ で完全である。
各微分で剰余階数の小行列式を分数体へ移すと、隣り合う階数の和は項の次元以上、
複体条件から以下であり、完全性が従う。この議論を各 $n$ に適用するので、
一様な有限検査次数を必要としない。従って幾何学的分数体上の $B$ とその二次双対はKoszulで、
Frobenius対合から左右AS正則、逆Hilbert行列から標準成長となる。
元の経路加群でも特殊・一般の次元が一致するため自由である。

最後に奇標数での入力を確認する。定理\ref{thm:walsh-presentation-space}の
九つの相異なる超平面の付値による平方類独立性と有限平坦な平方写像は奇標数でそのまま成立する。
三色ずつの双対の完全対合は定理\ref{thm:walsh-dual-frobenius}の
非退化対角形式と連結な台の証明で全点に成立する。
この既約表示族は定理\ref{thm:rank-twelve-quantum-limit}の量子点を含む。
そこに用いた基底変換は $2$ の冪と四乗根の単元だけで行え、六つの量子関係の順序単項式は
奇標数でも基底となる。従って同じ内部次数別bar複体の階数比較で、
この標数の幾何学的一般点は正則である。各次数の元の経路次元は任意の特殊点で標準値以上となる。

定理\ref{thm:walsh-normal-quadrics}の階数2での境界間移送は全等方平面の交叉次元から定まり、
個別の非零成分による除算を要しない。直交移送の行列式の符号も奇標数で保たれる。
実際、向きを検算した量子点の係数 $(0,1,i),(i,0,i),(-1,1,0)$ を
$\mathbb Z[i,1/2]$ 上に取る。そこで符号平方行列の二次小行列に単元があり、
どの剰余でも左右の境界写像が同じ核を持つ全射である。
従って移送は各奇素点で正則に定まり、分数体での行列式 $-1$ に特殊化は可換である。
各幾何学的奇標数の階数2開集合は既約で、この符号が一定となる。
分裂四元数の乗法、$\mathrm{PGL}_2\simeq\mathrm{SO}_3$、四正規二次式の交換則も
$2$ が可逆なら成立する。その中心平方による同定理の上界は一般点からの下界と一致する。
中心多項式環上の自由性と左右線型複体の証明により全次数の正則性を得る。
\end{proof}

\begin{theorem}\label{thm:regular-53-adic-ball}
定理\ref{thm:walsh-full-component}の標数53のテンソルを $w$ とする。
二隣接階数6のアフィンスキーム $Z$ で $w$ に還元される
$\mathbb Z_{53}$ 点は、76個の自由なテンソル係数の剰余類の直積で表示される。
その全ての分数体窓は標準Koszul左右AS正則である。この球には
$\overline{\mathbb Q}$ 上で定義される正則窓が十五次元成分の中でZariski稠密にある。
最小周期4、反対階数16、自明な射影安定化群を同時に課しても稠密である。
\end{theorem}
\begin{proof}
命題\ref{prop:twenty-conormal}の二十個の線型恒等式は
六次単元小行列のSchur分解と有限エタールな交点の番号付けによる。
横断的な六双一次式の全交点は剰余で二十単純点なので、射影性とJacobi判定で
第一切断の階数6軌跡の相対近傍でも有限エタールである。
従って同命題の消去は $\mathbb Z_{(53)}$ 上でも成立する。
剰余で独立な180個の微分を選び、対応する180係数について単元小行列を取る。
これらの方程式の零点は相対的に滑らかで、選ばなかった76係数への射影はエタールである。
実際、全方程式は局所エタールに180個で生成され、独立な180式は余接空間で
この生成系と可逆な係数行列で結ばれる。従って局所環で同じイデアルを生成する。
Henselの補題で各自由係数の指定にただ一つの持上げがある。

具体的に剰余窓の $m$ の三行は $(4,5,16),(23,-3,19),(-21,4,-1)$ である。
\path{computations/verify_unrestricted_good_reduction.py} は符号平方の階数2と行・列和、
二十点、180独立なSchur微分を確認する。さらに76個の係数に剰余二桁を指定し、
$53^3$ を法として独立な180式でNewton補正した結果が全200式を満たすことを確認する。
全ての精度と全ての次数に対する結論は上のエタール性と命題\ref{prop:unrestricted-good-reduction}による。

自由係数にその剰余類の有理整数を取れば、エタール射の零次元ファイバーなので
得る係数は $\mathbb Q$ 上代数的である。この整数格子は標数零の76次元空間に稠密である
（一変数ずつ多項式の零点を数えればよい）。局所環は相対滑らかで、その完備化は
$\mathbb Z_{53}[[z_1,\ldots,z_{76}]]$ である。剰余点を通る成分は一意で、
標数零の上記Walsh持上げを含む。この局所成分は幾何学的にも一意である。
実際係数体の任意の有限拡大に対しその付値環上で剰余点を固定すれば、
同じ完備局所環の記述がその付値環の係数で成立する。
よって格子の持上げの閉包は次元から定理\ref{thm:full-fifteen-component}の成分全体である。
同定理の一般の三条件は構成可能で稠密開集合を含み、その開への制限と四基底の商で主張を得る。
\end{proof}

% !TeX root = ../papers/nc-p3.tex
\section{九次元半周期成分の四直線枠・三重みと二次の中心切断}
\label{sec:half-period-frame}
\begin{theorem}\label{thm:half-period-frame}
九次元半周期成分の稠密開部分には次の表示がある。四空間 $E$ に非退化対称形式 $q$ と
その二つの直線族の順序を取り、$E^*$ の順序付き四直線枠の代表を $u_i$ とする。
$q=\sum_{i=1}^4 u_i\otimes v_i$ で $v_i$ を定め、$p_i=u_i\otimes v_i$ と置く。
$\sum_i\omega_i=1$ の非零重みを取り、三平面
\[
 I=\{\textstyle\sum_i c_i p_i:\sum_i\omega_i c_i=0\}
\]
を作る。$L_+,L_-\subset\bigwedge^2 E^*$ を、$q$ の各直線族のPlücker座標上で
消える交代形式の三空間とする。二周期窓の関係は
\[
 R_0=I\oplus L_+,\qquad R_1^{\mathrm{op}}=I\oplus L_- .
\]
一般の窓から $[q]$、直線族、指定なしの四組 $(p_i,\omega_i)$ を共通射影基底変更まで回収できる。
二周期での $q$ の極化は中央の二次切断である。
八つの固定点は三つの二次式による $\mathbb P^3$ の被約完全交叉である。
五次元半周期成分の一般点には、可逆な二位置移動と交換する非零の二次切断は存在しない。
\end{theorem}
\begin{proof}[点加群と四空間の同定]
定理\ref{thm:half-period-components}の九次元枝の一般点で位数2の持上げを選んで
二周期化する。点の座標に作用する対称な第一行列を $M$ とし、
$N_{ij,kl}=M_{kj,il}$ を第二の逆順行列とする。第一辺の二十点 $(a,b)$ に
次の辺をたどって偶位置の点 $c$ を対応させる。半周期で同定すれば $a\mapsto c$ は対合となる。
実際四位置をたどった同定は有限四辺の四移動の恒等性による。
\path{computations/verify_half_period_frame.py} は53のWalsh点で
位数2持上げに固定八点・二巡回六組、三つの位数4持上げに二巡回二組・四巡回四組を得る。
六関係の有限エタールな交点を番号付けすれば置換は局所定数であり、
これらの巡回型は対応する標数零の各成分の一般点へ移る。
八つの固定点の両射影は各々射影枠を含む。五次元側では二つの持上げは互いに逆で、双方とも固定点がない。
移動と交換する切断の各点での二次評価が非零なら、三文字での交換式に
全ての一文字を代入して、二位置先の点と移動で運んだ始点とが一致する。
実際、両端の一次汎関数は非零で比例する。この比較は交換時の位置別スカラーを許しても成立する。
二次評価は六関係による十次元商上で単射である（53の評価階数10は開条件）。
従って五次元側でそのような切断は零である。

九次元側で $H=M-N$ を $\bigwedge^2 E\to(\bigwedge^2 F)^*$ の写像と読む。
二つの偶・奇点空間をまだ $E,F$ と区別している。$H$ は両添字対で交代的であり、
53の点でこの写像は同型である。固定点 $(a,b)$ では
$M(a\otimes b)=N(a\otimes b)=0$ なので、$H(a\wedge E)$ は
$b\wedge F$ の零化空間である。体積形式で双対を同定するとこれは $b\wedge F$ 自身となる。
従って $H$ によるPlücker二次形式の引戻しは八つの $a\wedge E$ の上で消える。
既に射影枠の五点でこの形式は定数倍に一意である。実際、最初の四点を座標点にすると
相補な二添字の積三つだけが残り、第五点を $(1,1,1,1)$ にして
$(1,1,1,1)\wedge z$ を代入すると係数比はPlücker式のものになる。
よって $H$ はKlein二次超曲面の射影同型で一方の平面族を保つ。
その平面族のパラメータ $\mathbb P(E)$ と $\mathbb P(F)$ の同型は射影線型で、
各固定 $a$ を対応する $b$ に送る。これで両空間を同定すると、
各 $a\wedge E$ を保つ写像は二平面の交線も保ち、全分解可能線を固定するのでスカラーである。
元の $H_{ij,kl}$ は従って体積テンソルの非零倍である。

\end{proof}
\begin{proof}[八点と二つの直線族]
$S=\mathrm{Sym}^2E$, $T=\bigwedge^2 E$ として $E\otimes E=S\oplus T$ に分ける。
$H$ の核は $S$、$T$ への制限は非退化である。
$r=\operatorname{rank}(M|_S)$ とすると、対称性で
$\dim(\operatorname{im} M\cap T^*)=6-r$ である。
よって $\operatorname{im}M+\operatorname{im}N=\operatorname{im}M+T^*$ は次元 $6+r$、
両像の交叉は次元 $6-r$ で、そこに $M(S)=N(S)$ が含まれる。従って $r\le3$。
対角上の第一辺は $r$ 本以下の二次式で定まる。その全ての点は $N$ も消すから
ちょうど固定八点である。有限性から $r=3$、Bézoutにより三つの二次式の被約完全交叉である。
一般点で対称対角ブロックの階数も3である（53で $S$ 上の階数3）。
共通像 $I=M(S)$ は二つの像の交叉で、$S$ への制限が単射となる。

$S$ の三次元非退化商でSchur消去すれば
$M=X+D_+$, $N=X+D_-$ と書ける。$X$ は対称階数3で像は $I$、
$D_+,D_-$ は $T$ 上の対称階数3で、$D_+-D_-=H$ である。
$H^{-1}D_+$ と $-H^{-1}D_-$ は和が恒等で階数の和が6であり、互いに相補な自己共役射影である。
従って両像はPlücker対合で非退化な相補直交三空間となる。
各三空間が $T$ で消す三平面とKlein二次超曲面の交叉は滑らかな平面円錐曲線である。
これは $\mathbb P^3$ の滑らかな二次曲面の一方の直線族を与え、直交する他平面は反対の直線族を与える。
実際、円錐曲線上の異なる二点の極対合は非零なので対応する直線は交わらない。
三つの互いにねじれた直線を選べば一つの滑らかな二次曲面の同じ直線族に入り、
この三直線を全て横切る直線は直交平面の円錐曲線そのものである。二度適用すると元の族を得る。
この二次曲面を $q=0$、両像を $L_+,L_-$ と書く。一般点では $q$ は固定八点の全てで非零である。
これは53で構成した対角係数 $(37,15,47,1)$ の形式で確かめられる。
\end{proof}

\begin{proof}[四直線枠の回収と逆構成]
$q$ を行列で表し、$G_{ij,kl}=q_{ik}q_{jl}-q_{il}q_{jk}$ とする。
$H=h\epsilon$ を固定体積で表す。二つの直線族の順序に合わせて
$c=\pm h/\sqrt{\det q}$ の符号を選ぶと
\[
 D_+=(H+cG)/2,\qquad D_-=(-H+cG)/2 .
\]
これは $H^{-1}cG$ の $T$ 上の二固有空間が二つの直線族の三空間となることによる。
例えば $q$ を単位行列にすれば、対合は相補二添字とその並べ替え符号を交換するHodge作用である。
$\Gamma$ を $i,k$ の部分転置とする。$Y=X+(c/2)q\otimes q$ では $Y^\Gamma=Y$ となる。
実際 $G^\Gamma-G=(q\otimes q)^\Gamma-q\otimes q$ を $M^\Gamma=N$ に代入すればよい。
一般に $Y$ の階数は4である。$q$ が固定点で非零なので $q\notin I$ である。

$Y$ の像の基底となる四つの双一次式の左変数を $a$ で評価した $4\times4$ 行列を $B(a)$ とし、
$Y$ はこれらの形式の積に非退化対称係数行列 $D$ を付けて書く。
部分転置不変性は $B(a)^{\mathsf t}D B(b)=B(b)^{\mathsf t}D B(a)$ である。
$B(a_0)$ が可逆なら $C(a)=B(a_0)^{-1}B(a)$ は
$B(a_0)^{\mathsf t}D B(a_0)$ について自己共役である。同じ式から全ての $C(a)$ は可換である。
53の検算では $a_0=(1,1,1,1)$ で行列式3、
$a=(7,4,2,9)$ で $C(a)$ の固有値は $-3,-4,-15,-52$ の四単根である。
従って一般には同時固有四直線があり、非退化対合で直交する。
相手側の四形式の基底も変えると $B(a)$ は対角となり、像の基底は四つの純双一次式になる。
その空間に非退化な $q$ が入るため両側の四因子がそれぞれ基底であり、
$q=\sum p_i$ で純形式の倍率を一意に定められる。純形式の直線はこの四本だけである。
$Y=(c/2)\sum_i p_i\otimes p_i/\omega_i$ と書くと全 $\omega_i\ne0$ で、
$X=Y-(c/2)q\otimes q$ の階数3は行列式補題から $\sum_i\omega_i=1$ と同値である。
その像が所定の $I$ である。

逆に任意のこの枠と非零重みから上式で $X,D_+,D_-$ を作れば $M^\Gamma=N$ で、
$I$ の交代三空間との交叉が零なら両隣接階数は6である。
識別に使った位数2の持上げは九次元側の一般点では射影的に一意である。
八点の両射影の同定、二つの交代三空間、$q$、$I$ を順に回収し、
$I+kq$ の四純直線と $I$ の方程式から四線と重みが回収される。
したがってこの開集合では同時添字置換と直線族の順序を保つ射影同型が同型判定を与える。
$q$ を固定して四線枠に $15-3=12$、その射影直交群に6、重みに3の自由度があり、
既知の成分次元9と一致する。最後に $R_0,R_1^{\mathrm{op}}$ の共通平面から、
偶・奇双方で $p_i=\omega_i q$（奇では転置）が成り立つ。
$(u_i v_i)u_i=u_i(v_i u_i)$ とその逆により、$q$ は全 $u_i,v_i$ と二周期で交換する。
\end{proof}

% !TeX root = ../papers/nc-p3.tex
\section{中心商を固定した三重み全体の正則性}
\label{sec:half-period-weight-fibers}
前節の枠で $K=\langle p_1,p_2,p_3,p_4\rangle$ と置く。
重みは $\omega_i\ne0$, $\sum\omega_i=1$ とし、$I\cap L_+=I\cap L_-=0$ である範囲を取る。
これは全ての隣接階数が6となる条件である。一文字切断は全階数である。
実際各 $L_\pm$ が可逆な交代形式を含むため、各辺の二方向の台はともに全四空間となる。
折り畳んだ二頂点代数を $A_\omega$、中心二次切断を $q$ と書く。
\begin{theorem}\label{thm:half-period-weight-fibers}
共通の中心商 $B$ は関係 $K+L_+$、$(K+L_-)^{\mathrm{op}}$ を持つ。
次の二条件は同値である。
\begin{enumerate}
\item $B$ は各始点で $H_B=(1-t^2)/(1-t)^4$ を持ち、左右の次数付きsocleが零である。
\item 全ての上記重みで $A_\omega$ は標準Koszul左右AS正則で $q$ は非零因子である。
\end{enumerate}
重みの範囲が空でないとする。この範囲のただ一つの重みで後者の正則性と非零因子性が
成立すれば両条件が従う。従ってこの判定には枠の六自由度だけが残る。
\end{theorem}
\begin{proof}
前節の積表示から $q$ は中心であり、$I+kq=K$ なので商は重みに依存しない。
全枠・重みの表示空間は既約である。実際、順序付き直線枠と基底空間の可逆変換は既約であり、
滑らかな形式と直線族の順序の対も固定した一対の $\mathrm{GL}_4$ 軌道である。
重みの超平面と上記の階数開条件を加えても既約である。
この族には正則点がある。定理\ref{thm:half-period-frame}の53の窓の枠と等重みの
特性零持上げが一例であり、下記の有理例で直接確かめる。
各内部次数における双対bar複体の議論を適用すれば、
この表示空間の幾何学的一般点で標準Koszul性と標準Hilbert関数が成立する
（命題\ref{prop:unrestricted-good-reduction}の縮約下界と同じ議論を局所の小行列上で行う）。
元の各経路次元の上半連続性により、全特殊点で標準Hilbert関数は下界となる。

第一条件の下で、$q$ による乗法と商への全射から、係数ごとに
$H_{A_\omega}\le H_B/(1-t^2)=(1-t)^{-4}$ である。下界と一致し、
乗法の核が全次数で零である。左右に同じ議論が使える。
中央の非零因子 $q$ の長さ1の自由解消により
$\operatorname{Ext}_{A_\omega}^1(S,A_\omega)
\simeq\operatorname{Hom}_B(S,B)(2)=0$、また $\operatorname{Hom}_{A_\omega}(S,A_\omega)=0$ となる。
ここで $S$ は各頂点単純である。前者は $A_\omega\to B$ の環変更スペクトル列で
$\operatorname{Ext}_{A_\omega}^j(B,A_\omega)$ が $j=1$ の $B(2)$ だけとなることから従う。
左右で同じである。

四次テンソルの縮約を項に持つ長さ4の線型列を取る。
両隣接階数6と一文字階数4から中間係数空間の次元は $1,4,6,4,1$ で、
反対側の表示の最初の二微分を双対化したものが最後の二微分である。
実際、相補縮約間の対合は同じテンソルの切断の左右像の対合であり、
三次から二次への含有は左の二次から一文字への含有に転置する。
関係による表示で最初の二項は完全、上の $\operatorname{Hom}$ と $\operatorname{Ext}^1$ の
消滅で最後の二項も完全である。残る中央では $H_{A_\omega}=(1-t)^{-4}$ の
Euler等式によりホモロジーが零となる。左右に適用して標準Koszul左右AS正則性を得る。
逆に一つの正則重みで $q$ が非零因子なら、乗法の短完全列が所定の商の級数を与え、
同じ環変更同型とAS消滅で商の両socleが零となる。
\end{proof}

\begin{corollary}\label{cor:half-period-rational-frame}
次の有理枠に対し、代数閉体の標数零で非零の全重み $\sum\omega_i=1$ が
標準Koszul左右AS正則な窓を与える。$q$ の対角と各 $u_i$ の行は
\[
 (37,15,47,9/(37\cdot15\cdot47)),\qquad
 \begin{pmatrix}
 1&-5&25&-26\\1&-5&-25&26\\1&5&25&26\\1&5&-25&-26
 \end{pmatrix}.
\]
直線族の順序は標準体積に対して $\sqrt{\det q}=3$ の符号で取る。
同じ剰余枠・直線族順序を持つ任意の非退化な53進持上げでも、
その幾何学的分数体の全ての非零重みに対して成立する。
\end{corollary}
\begin{proof}
$\det q=9$ で $K$ から対角二次式への制限は剰余で階数4である。
従って $K\cap \bigwedge^2E^*=0$ で全重みが二つの階数開条件を満たす。
\path{computations/verify_half_period_weight_fibers.py} は等重みで $h=45$, $c=15$ とした
第一テンソルの53での還元が、前節で基底を揃えたWalsh窓と一致することを直接確かめる。
命題\ref{prop:unrestricted-good-reduction}により等重みの持上げは全次数で正則である。

剰余で $q$ の非零因子性も確認する。Walshの四正規二次式のうち添字3のものは
交換符号 $\phi_3(i)=(-1)^{i_0i_1}=J(i)$ を持つため、この位数2持上げと交換する。
その平方は四中心変数の一つなので、左右とも非零因子である。
各位置で動く十二点の二次評価の階数は9で、両側とも十次元経路商からの核は直線である。
この正規式も、構成した $q$ もこの核に入り、いずれも非零である
（$q$ は固定八点で非零）。従って $q$ は各位置でこの正規式の非零倍となる。
等重みの整持上げの各経路加群は自由であり、剰余で単射な乗法の小行列を持ち上げれば
一般側でも全次数で $q$ は左右非零因子である。定理\ref{thm:half-period-weight-fibers}により
重みへの付値条件は消え、分母または分子が53で消える重みも含む。
\end{proof}

% !TeX root = ../papers/nc-p3.tex
\section{二次曲面との交点による枠のKoszulフィルトレーション}
\label{sec:half-period-semisimple}
\begin{theorem}\label{thm:half-period-semisimple}
前節の枠で $e_i$ を $u_i$ の双対基底とし、$G=(q(e_i,e_j))$ とする。
$G,G^{-1}$ の対角項が全て非零とする。
このとき標数零の代数閉体で、全ての非零重み $\sum\omega_i=1$ に対し
窓は標準Koszul左右AS正則で $q$ は非零因子である。
これは九次元成分の稠密開部分の全閉点を含む有限な十分条件である。
\end{theorem}
\begin{proof}[八次元半単純環での二次双対]
$D=k^4\oplus M_2(k)$ に積を入れ、$\ell(t,Z)=\sum_i t_i-\operatorname{tr}(Z)/2$、
$H=\ker\ell$ とする。$q(x,x)=\det X(x)$ となる線型同型 $X:E\to M_2(k)$ を選ぶ。
偶・奇の一次空間を $D$ に次の四空間 $V,W$ として入れる。
\[
 x\mapsto ((u_i(x)),X(x)),\qquad y\mapsto ((v_i(y)),\operatorname{adj}X(y)).
\]
直線族の順序の交換が必要なら転置を施す。$\operatorname{adj}$ は二次行列では線型である。
両空間ともスカラー四成分への射影と行列成分への射影は同型である。
極化恒等式で $\ell(VW)=\ell(WV)=0$ であり、非退化な跡対合 $\ell(ab)$ に関して $W=V^\perp$ である。
二周期でこれらを交互に掛けて生成する次数付き代数を $C$ とする（次数ごとに $D$ の別のコピーを使う）。
二次積のスカラー成分は四つの $p_i$、行列成分の跡零部分は $L_+$ または $L_-$ の三形式である。
実際、$X\operatorname{adj}Y$ の跡零部分は交代的で、階数一の二行列に共通の核があれば零である。
従ってその三形式は一方の直線族を消す三空間を張る。逆順積では共通の像を使い他方の族となる。
$\sum c_i p_i$ の対角を $e_j$ で評価すると $c_jG_{jj}$ となるので $K$ は交代空間に交わらない。
よって二次積の像は各位置でちょうど $H$、核は中心商 $B$ の関係空間の直交である。

$V,W$ はともに $D$ の単元を稠密に含む。$H$ の右安定化単元は
スカラー単元だけである。実際 $H$ の直交直線は $1$ で張られ、跡対合は非退化だからである。
従って $H$ に比例しない二単元を掛けた二超平面は異なり、三次積の像は $D$ になる。
以後も単元を掛ければ $D$ である。従って各始点で $C_n$ は次数順に $k,V$ または $W,H,D,D,\ldots$ となる。
\end{proof}

\begin{proof}[線型イデアルの構成]
始点を固定してその一次空間を $V_0$、次を $V_1$ とする（$V,W$ のどちらから始めてもよい）。
$D$ の右イデアルを
\[
 D_{S,L}=\{(t,Z):t_i=0\ (i\notin S),\ \operatorname{im}Z\subset L\}
 \quad(S\subset\{1,2,3,4\},\ L\subset k^2)
\]
と書き、これと正次数の $C$ との次数別交叉を $J_{S,L}$ とする。次数零は零とする。
$V_p$ のスカラーが $S$ に台を持つ部分は $|S|$ 次元である。
四添字を任意に二対に分け、そのいずれかを $S$ とする。対応する行列二空間の射影直線は行列式二次曲面に含まれない。
偶・奇の行列式は対角基底で $G$ または $G^{-1}$ の形式（共通倍率を除く）で、
各端の座標行列は可逆である。直線は少なくとも一つ階数1の根を持ち、その台はちょうど $S$ である。二重根も許す。以下では行列成分を大文字で、対応する一次元を小文字で表す。

根 $x$ の行列を $X$ とする。右零化空間は $D_{S^c,\ker X}$ で、
次位置の相補割線には像が $\ker X$ の根 $y$ がちょうど一つある。
実際、相補二空間は行列跡で直交する。第一の根 $X$ と同じ二空間内の可逆な $X'$ を基底とする。
像を $\ker X$ に制限した行列二空間では $\operatorname{tr}(XY)=0$ は自動で、
$\operatorname{tr}(X'Y)=0$ は非零の一方程式なので解は直線である。
従って $x V_1$ は次元3で $H\cap D_{S,\operatorname{im}X}$ に等しく、
$xH=xD=D_{S,\operatorname{im}X}$ は次元4である。後者は $\ker(x\cdot)$ に
$S^c$ 上の単位座標があり $\ell$ が非零だからである。一次でも同じ二空間に可逆な行列があるので
$J_{S,\operatorname{im}X}=xC$、その右零化イデアルは $yC$ となる。この議論は全根と両始点に適用する。

次にこの根 $x$ と同じ二空間の可逆行列を持つ $x'$ を加える。$xC$ に対する $x'$ の右コロン
$\{b:x'b\in xC\}$ は次位置の $J_{\mathrm{all},L'}$ である。
ただし $L'=X'^{-1}\operatorname{im}X$ である。正次数の積では $xC$ は上記の交叉を全て取り、$x'$ のスカラー台は $S$ 内であるからである。
このコロンの一次元数は2、二次は5、以後は6である。
従って $(x,x')C=J_{S,k^2}$ は次数1,2,3以後で次元 $2,5,6$ となる。
実際 $(x,x')C/xC$ はコロンによる商の一次ずらしであり、各次数で交叉の次元に達する。

ここで現れた $L=L'$ の $J_{\mathrm{all},L}$ も一次元を一つずつ加えて作れる。
この $L$ への $X,X'$ の制限は比例するため、像が $L$ に入る行列について二つの跡直交式は一つの非零式となる。従って $S^c$ 上で $\operatorname{im}Y=L$ の根 $y$ がある。それを取り、同じ行列像条件の
二空間から、全スカラーが非零な独立な $y'$ を選ぶ。
各スカラー座標汎関数は行列側で可逆な対向座標行列との跡対合である
（$V,W$ の直交と対向座標の非等方性を使う）。よってこの二空間全体で消える座標はない。
$yC$ に対する $y'$ の右コロンは $J_{S^c,k^2}$ となる。
行列の像条件は自動で、スカラーは $S^c$ の外で零となる必要があるからである。
同じ商の次元計算で $(y,y')C=J_{\mathrm{all},L}$ となる。

スカラー台が三点のときは必ず分割の一対を含む。その対から一座標ずつ増やして全体まで進む。
$J_{S,k^2}$ に座標 $j\notin S$ の一次元（スカラーが $j$ にだけ台を持つ元）を加えたとき、
右コロンは $J_{\mathrm{all}\setminus\{j\},k^2}$ である。
次元を各次数で一つずつ増やすので、得るイデアルはちょうど $J_{S\cup\{j\},k^2}$ である。
こうして零、根の主イデアル、二点・三点・全点のスカラー台のイデアル、上記の行列像直線のイデアルが
Koszulフィルトレーションをなす。すなわち各非零イデアルは同じ集合の小さいイデアルに
一次元を一つ加えたもので、そのコロンも集合に属する。
この条件が頂点単純の線型解消を与えることは、各商の短完全列のTor長完全列で従う。
ホモロジー次数に関する帰納の内側で一次生成元数に関して帰納すれば、
コロンの商は一つ低いホモロジー次数しか必要としない。
全正次数イデアルによる商もこの集合にあるから $C$ はKoszulであり、特に二次表示である。
\path{computations/verify_half_period_semisimple.py} は二つの始点の根・主イデアル・
二生成イデアルの各乗法階数を単根・二重根の両方で検算する。
\end{proof}

\begin{proof}[全重みへの適用]
$B^!=C$ が二次関係から従う。両位置の二次空間 $H$ には同じ元
$z=(\omega_1,\ldots,\omega_4,I_2)$ が入る。$D$ の中心単元であり、二周期で $C$ の中心非零因子である。
$B$ の七関係のうち $A_\omega$ に残るものは、$L_\pm$ と、$p_i$ の係数を $\omega$ に対して
直交させた $I$ である。従って直交空間を取れば $A_\omega^!=C/(z)$ である。
中心正則二次式によるこの商もKoszulである。実際 $C\to C/(z)$ のTor環変更スペクトル列は二行しかなく、
下段のホモロジー次数 $p$ のTorから出る可能な微分は上段の $p-2$ のTorの内部次数2ずらしへのものだけである。
その核は極限項で $C$ のKoszul性により内部次数 $p$ に集中する。像も $p$ に関する帰納で同じ次数に集中する。

商のHilbert関数は $1,4,6,4,1$ となり、四次 $D/zH$ の汎関数を $a\mapsto\ell(z^{-1}a)$ とする。
二次 $H/kz$ 上の積対合の根基は零である。$H$ の直交は $k1$ なので割る前の根基が $kz$ だからである。
一次と反対始点の三次 $D/zV_1$ でも $V_0^\perp=V_1$ から完全対合を得る。
従って $A_\omega$ は標準Koszul左右AS正則である。$B$ もKoszulで、
$C$ の級数を符号反転して逆に取れば $B$ は前節の所定の級数を持つ。短完全列の次元比較で $q$ は非零因子である。
具体例 $G=I_4$ でも全ての非零条件が成立し、開条件全体が枠空間に稠密である。
\end{proof}

% !TeX root = ../papers/nc-p3.tex
\section{可除代数の特殊化と十五次元成分の量子族との全交叉}
\label{sec:quantum-component-intersection}
\begin{theorem}\label{thm:quantum-component-intersection}
定理\ref{thm:full-fifteen-component}の成分の閉包に属する標準AS窓のうち、
量子多項式代数のindexed窓は添字保持同型を除いて一つだけである。
それは定理\ref{thm:rank-twelve-quantum-limit}の反対階数12の量子窓である。
より一般に、この閉包の標準AS点でFrobenius標準四周期の一つの隅環が
整域なら、その商可除代数の中心上の次数はちょうど4である。
\end{theorem}
\begin{proof}[同じ値群の可除代数の剰余]
まず付値の事実を示す。中心 $F$ 上次数 $d$ の可除代数 $D$ が離散付値を持ち、
値1の中心元 $t$ があるとする。剰余標数を零とする。このとき剰余可除代数も
その中心上の次数が $d$ である。

実際、付値完備化を取ると $D\otimes_F\widehat F\to\widehat D$ は単射である。
左辺は中心的単純環で写像は単位を保つ。像は完備体上有限次元のノルム空間なので完備かつ閉であり、
稠密な $D$ を含むため全射である。従って完備後も次数と中心は指定通りである。
剰余 $\bar D$ の $\bar F$ 上一次独立な元の持上げは、係数の最小付値を除いて還元することにより
$\widehat F$ 上でも独立である。従って剰余基底は有限で、その持上げは逆に全体を張る。
実際、各元の先頭係数を基底で消し、$t$ で割って反復すれば完備性で各係数が収束する。
よって $[\bar D:\bar F]=d^2$ である。

$\alpha\in Z(\bar D)$ は $\bar F$ 上有限分離である。最小多項式を中心の付値環上の
モニックな $f$ に持ち上げる。$\alpha$ の任意の持上げからNewton法を行うと、
各段階はその元の生成する可換部分体の中で行え、導関数の剰余値が非零なので根 $a$ を得る。
この根は非可換な付値環全体でも $\alpha$ の上にただ一つである。
差 $t^m b$ を持つ二根を $t^{m+1}$ を法として比較すると、
$\alpha$ の中心性で一次項は $f'(\alpha)\bar b$ となり矛盾する。
任意の単元による $a$ の共役は同じ剰余の根であり、$a$ は全単元と交換する。
全非零元は $t$ の冪と単元の積なので $a\in\widehat F$ である。従って $Z(\bar D)=\bar F$ を得る。

成分の一般点を分数体側、指定した標準点を特殊側に持つ離散付値環を取る。
例えばその点を吹き上げて正規化した例外素因子で取り、必要なら剰余体を拡張して考える。
基礎体は代数閉体なので、整域性はこの係数拡大で保たれる。
一般点と特殊点で各経路次元は同じ標準値であり、関係は単元な隣接小行列で定まる。
従って各斉次経路加群は自由で、隅の整モデル $A$ の剰余は所望の隅環である。
各元の $t$ 整除回数は有限個の斉次成分で読める。
剰余が整域なのでこの順序は非零元の積で加法的であり、
一般の隅のOre商可除代数に $v(ab^{-1})=v(a)-v(b)$ で延びる。
一般の隅が次数4の可除代数となることは系\ref{cor:quartic-division-corners}による。
値0の分数の分子・分母から共通の $t$ 冪を除けば、剰余は特殊隅の二元の分数となる。
逆に特殊隅の各非零元はこの剰余可除代数内で可逆であり、剰余は特殊隅の商可除代数に一致する。
標準恒等式は係数拡大と降下に可換であるから、剰余体を拡げて得た次数は元の点でも同じである。
これで一般の主張を得た。
\end{proof}

\begin{proof}[量子窓の標準四周期の隅と四次中心]
三角積 $T_{0ij}$ が生成する群の位数を $N$ と書く（無限も許す）。
与えられた量子窓を対角twist $q'_{ij}=q_{ij}d_i/d_j$ で
$\prod_jq'_{ij}=1$ にできる。実際 $d_i^4=(\prod_jq_{ij})^{-1}$ を取り、
その積は四乗根なので一つの $d_i$ の四乗根倍で積を1にする。
この表示では四変数を一度ずつ掛けた体積単項式 $m$ は中心である。
双対量子外積の最高次汎関数について一文字を残りの三文字の後ろへ移したときの係数は
$-\prod_jq'_{ij}$ またはその逆で、全てに共通の符号である。
従って通常の四位置反復が射影基底においてFrobenius標準四周期を与える。

隅は四次Veroneseで、その商可除代数は $\sum_i b_i\equiv0\pmod4$ の
整数指数のLaurent単項式で生成される。実際 $m$ の冪を掛ければ全指数を非負にできる。
この格子は中心な $(1,1,1,1)$ の直線と
$L=\{b:\sum_i b_i=0\}$ の直和である。$L$ の基底 $e_j-e_0$ の交換子は
三角積（または同時にその逆）であり、twistの選択に依らない。
非ねじれ交換子があればその二単項式の量子トーラスはPIにならない。
例えば中心上有限次元の可除代数で一方の共役に対する他方の冪が
無限個の異なる固有値を持つことはできない。
有限の場合、$L$ 上の交代指標の根基の指数は $N^2$ である。
原始 $N$ 乗根に対する三つの指数と $N$ の最大公約数は1であり、
整数の三次交代行列のSmith標準形は非零二因子がともに三係数の最大公約数となる。
従って法 $N$ の像の位数は $N^2$ である（$N=1$ は直接）。
交換子の根基の単項式がLaurent環の中心を張り、その中心分数体上で各剰余類の単項式は基底をなす。
よって隅のPI次数は $N$ である。上の特殊化定理から $N=4$ が必要となる。

さらに成分の四次中心は一般点で五次元なので、標準AS特殊点で少なくとも五次元である。
実際、長さ4,5の経路を自由化すれば矢との交換式の核次元は上半連続である。
量子窓で長さ4の単項式 $x^b$ が中心切断に寄与する必要十分条件は
$b-(1,1,1,1)\in\operatorname{rad}(L)$ である。これは四文字と交換する際の倍率が
文字に依らないことに等しい。その共通倍率は四文字の積 $m$ が中心なので四乗が1であり、
位置別の定数倍を選べば切断が一周する。各多重次数で寄与はちょうど一つである。

三角積を $i^a,i^b,i^c$ と書く。定理\ref{thm:quantum-indexed-all}の全置換で割ると、
位数4の候補と上式による中心次元は以下で尽きる。
\[
\begin{array}{c|ccccc}
 (a,b,c)&(0,0,1)&(0,1,2)&(1,1,1)&(1,1,2)&(1,3,1)\\\hline
 \text{軌道長}&12&24&6&12&2\\
 \text{四次中心次元}&3&2&3&1&5
\end{array}
\]
\path{computations/verify_quantum_component_intersection.py} は原始56組、四添字の全置換と
全次数4単項式を列挙する。残る最後の軌道は定理\ref{thm:rank-twelve-quantum-limit}の
量子窓の三角積であり、定理\ref{thm:walsh-full-component}で実際にこの成分に入る。
従って必要条件がここでは十分である。
\end{proof}

% !TeX root = ../papers/nc-p3.tex
\section{量子境界の五中心局所方程式と対角Ore枝との全交叉}
\label{sec:quantum-ore-incidence}
\begin{theorem}\label{thm:quantum-ore-incidence}
定理\ref{thm:quantum-component-intersection}の量子点の近くで、二隣接階数6のテンソルのうち
四次中心が五次元以上となる局所スキームは滑らかで、ちょうど十五次元成分に一致する。
さらに、基底の一径数部分群による射影極限がこの量子点である任意の隣接階数6テンソルは、
この成分に属し、それ自身が全次数で標準Koszul左右AS正則である。
したがって量子テンソルへのこの形の退化を持つ全ての窓で、成分所属は極限の量子同型類だけで決まる。
対角量子Ore表示 $A(q,f)$ が標準AS正則であるとする。
そのindexed類がこの成分の閉包にある必要十分条件は、背景 $A(q,0)$ が同定理の唯一の量子類であることである。
また定理\ref{thm:quantum-tangent-finite}の二つの九係数表示では $\zeta^2=-1$ の全係数がこの閉包に属し、
$\zeta=-1$ の表示は係数の特殊値を含めて全て外にある。有限四辺はここでは仮定しない。
\end{theorem}
\begin{proof}[階数条件の局所成分]
量子点で経路長4,5の経路加群は隣接階数スキームの近傍で局所自由である。
実際節\ref{sec:unrestricted-doublet}の局所Artin環上の平坦性を完備化の各商に適用し、
有限表示加群の完備化と忠実平坦性を用いればよい。
中心の交換行列の核が五次元以上となる小行列式スキームを $Z_5$ とする。
量子点で核はちょうど五次元であり、$Z_5$ 上の一次変形ではその基底が持ち上がる。

二隣接切断のSchur微分の全階数はこの点で少なくとも177である。
\path{computations/verify_quantum_ore_incidence.py} は四乗根を法97で22として、
量子交代テンソルの両切断の左右の十次元核で摂動を挟み、二百行の階数177を得る。
計算では三辺を1にしたゲージを使う。巡回標準化のための四基底変更はこの階数を変えない。
よってアフィン接空間次元は高々79である。三つの三角積の量子変形に制限すると、
$Z_5$ の一次条件の核は零である。実際体積単項式を中心にする前節のゲージは無限小でも選べる。
四つの純四乗の持上げを多重次数ごとに射影すると同じ純単項式の非零倍に限られる。
従って一階の交代指標の変化は全 $4e_i-\mathbf1$ と $L$ の各元との組で零になる。
四ベクトルは $L\otimes\mathbb Q$ を張るので三方向は全て消える。
従って $\dim T Z_5\le76$ である。他方既約な十五次元成分のアフィン閉包は次元76で、
一般点の五中心から $Z_5$ に含まれ量子点を通る。よって局所次元も接次元も76で、
局所環は正則、成分は一意である。
\end{proof}

\begin{proof}[量子点へ流入する全ての窓]
量子点を固定する連結射影安定化群は四頂点に共通な対角トーラスである。
実際、四三角積は全て1と異なるので四頂点の点配置は六本の座標直線で、その四三重点は内在的である。
連結群は各頂点で対角に作用する。六つの量子関係の両項の比を保つには、
位置の異なる対角行列は射影的に同じでなければならない。
このトーラスは三つの量子パラメータを変えたテンソルの直線も全て固定する。
アフィン空間ではその直線上の同じ倍率を除いて作用させる。従って
$T Z/T Z_5$ は自明指標だけからなる。実際、上の三つの量子接方向が商で独立で、
$\dim T Z\le79$ と $\dim T Z_5=76$ から商の全体を張る。

任意の一径数作用で量子点 $w$ を固定するものを取る。
$Z_5$ の完備局所イデアルはこの作用の重み零の三元で生成されることを確認する。
$Z$ の最大イデアルの余接空間の固有基底を選ぶ。$Z_5$ への制限の核の三ベクトルは重み零なので、
それらを $Z_5$ のイデアルから重み零で持ち上げ、残りを固有元で持ち上げる。
各有限jetでトーラス表現の重み分解が完全であるから、これらの選択は完備化で両立してできる。
全持上げは完備局所環を位相的に生成する。最後の76元は滑らかな $Z_5$ の正則形式座標なので、
最初の三元で割った後、76変数形式冪級数環からの全射とその $Z_5$ への写像の合成は同型である。
従って三元だけでイデアルを生成する。これは $Z$ の他の局所成分を排除する主張ではない。

いま $W$ の軌道弧 $t\mapsto\lambda(t)W$ の極限を $w$ とする。射影極限では
最小重みの全体倍率も除いて、非零なアフィン極限を取る。作用はこの極限を固定し、
余接商への作用は上の連結安定化群を通じて自明である。
三つの重み零生成元の $k[[t]]$ への引戻しは $t\mapsto st$ に不変だから定数で、定数項は零である。
よってこの弧は $Z_5$ の一意な局所成分に入り、四基底変更で不変な閉包には $W$ 自身が入る。
さらに特殊点が標準AS正則なので、付値環上の命題\ref{prop:unrestricted-good-reduction}で
分数体側も全次数で正則であり、軌道上の $W$ も正則となる。
逆に他の量子点へこの形で退化する窓が十五次元成分にあれば、基底変更不変な閉包にその量子点も属し、前節に反する。
\end{proof}

\begin{proof}[九係数表示の中心]
純三文字の巡回量子係数を $\zeta$、$\zeta^2=-1$ とする。混合式を
$X_n x_a-x_a X_{n+1}=f_{n,a}$ と書く。各項は $a$ を含まない $x_b^s x_c^{2-s}$ であり、
係数の位置倍率は $\zeta^{n(s+1)}$ である。この表示の厳密四移動はFrobenius標準である。
これを確認するため、双対の量子外積への重み付き随伴次数とその最高次基底でtraceを規格化する。
PBW基底で各次数の双対と対合は全係数空間上で自由、完全であるから、
局所Nakayama行列は係数に正則であり、対角基底倍率に共変である。
$X$ だけの倍率により各 $C_{a,s}$ は重み1にできる。従って基点と異なり得るのは
$X$ と純文字の間の非対角成分だけで、そこでも係数は一次でなければならない。
純文字の三つの独立な倍率で比べると、$C_{a,s}$ の純部分の指標は
$e_a-se_b-(2-s)e_c$（全指標の符号は作用の向きによる）であり、
一文字と $X$ の間に必要な $\pm e_j$ に等しくない。従って補正はない。
背景では全行積が1なので標準四移動は通常の反復である。

純三文字の量子Laurent環に一つの可換文字 $Y$ を加えた窓に、
$X_n=Y+\sum_a g_{n,a}$ と置いて埋め込む。$g_{n,a}$ は三つのLaurent単項式
$x^{\alpha_s}$、$\alpha_s=-e_a+s e_b+(2-s)e_c$ の線型結合である。
$x_j x^{\alpha_s}=\chi_j x^{\alpha_s}x_j$ とすれば
$\chi_b=\chi_c=\zeta^{-s-1}$、$\chi_a=\zeta^{2s-2}$ である。
従って位置倍率を $\chi_b^{-n}$ とすると $b,c$ で混合差は零となり、
$a$ では $1-\chi_a/\chi_b\ne0$ であるから所定の係数を作れる。
PBW順序で $X$ についてモニックな置換なのでこの表現は忠実である。

一つの極だけを残した $U_{n,a}=Y+g_{n,a}$ と
$P_{n,a}=U_{n,a}U_{n+1,a}U_{n+2,a}U_{n+3,a}$ を作る。$P_{n,a}$ はその一極窓の中心切断である。
$b,c$ とは $U$ 自身が交換し、$a$ を四回移した誤差は各 $s$ で
$\sum_{r=0}^3\zeta^{r(s+1)}=0$ を因子に持つ。実際、誤差の純二文字は $b,c$ だけなので
残る三つの $U$ を全て同じ側へ移せる。また $U$ とは四周期で交換する。
Laurent化すると $P_{n,a}$ は $Y$ とも交換するので $P_{n,a}=P_{n+1,a}$、
各純文字と通常の意味で可換である。以後 $P_a$ と書く。
$Y^{4-r}$ の係数に現れる純指数は $a$ 方向で $-r$、他の二方向で非負である。
三文字のLaurent中心の指数は三成分が法4で等しいので、$1\le r\le4$ では
他の二指数がともに $4-r$ 以上となる。

全三極に対し
\[
             z=Y^4+\sum_{a=0}^2(P_a-Y^4)
\]
はLaurent窓の中心切断であり、実際に元の窓に属する。
固定した $a$ 以外の純文字を可逆にすると、$U_{n,a}=X_n-\sum_{b\ne a}g_{n,b}$ はこの局所表示に入る。
$P_a$ も入る。$b\ne a$ の $P_b-Y^4$ の $Y^j$ の係数は $x_a^j$ で割れる。
$Y=U_{n,a}-g_{n,a}$ を代入すると $g_{n,a}$ の $x_a$ の極は高々一位である。
厳密には $Y^j$ の代わりに、順序付き $U$ 積の最高 $Y$ 項を順に消す。
長さ $d$ の $U$ 積の $Y^h$ 係数の付値は $h-d$ 以上なので、残る $Y^h$ の係数は
常に $x_a$ で少なくとも $h$ 回割れ、全ての $U$ 係数が整である。
これは各 $a$ について成立する。各局所表示の $X$ 順序基底は同じなので、
その係数の交叉は元の量子多項式環であり $z$ が整となる。
同様に $x_0x_1x_2Y$ では三極が一度ずつ消える。
これら二式と純四乗三式は $Y$ 次数と純多重次数により独立な五つの中心四次切断である。
前掲コードは三極を同時に持つ例で、四位置の $X$ 順序基底への書戻しも検算する。
\end{proof}

\begin{proof}[枝全体への適用]
九係数を同時に $t$ 倍する族は $t\ne0$ で $X$ の倍率変更により一定同型類で、
$t=0$ で量子背景になる。原始四乗根の背景は前節の唯一の類である。
全族に五中心があるため、量子点の $Z_5$ の一意な局所成分の結果が族全体に適用される。
一方 $\zeta=-1$ で閉包に入る点があれば同じ基底変更による特殊化でその量子背景も入るが、
その三角積の群は位数2であり前節に反する。

任意の正則な対角量子Ore表示でも、補題\ref{lem:quantum-ore-background-rank}の縮尺族により
背景の所属が必要である。逆に背景が指定量子類なら、第四文字との三つの交換係数を対角twistで1にする。
残る三辺は三角積であり、純三文字を付け替えて巡回順に同じ原始四乗根となる。
混合補正は位置依存の純二次式となり、式\eqref{eq:ternary-shift-derivation}を満たす。
定理\ref{thm:quantum-tangent-finite}の多重次数分解では、非負の $e_a$ 次数の補正は
$\mathbf1,\chi$ が独立なので位置別の $X$ への一次式の加算で消える。
残る次数は上記九係数だけで、その位相も同じ式で定まる。
この除去計算は辺の有限性を使わない。従って標準化したindexed窓は原始四乗根の九係数表示に入り、
十分性も従う。
\end{proof}

% !TeX root = ../papers/nc-p3.tex
\section{量子境界に退化する全窓の十八座標・四極大軌跡}
\label{sec:quantum-attraction-fiber}
\begin{theorem}\label{thm:quantum-attraction-fiber}
前節の階数12の量子類 $w$ を軌道閉包に含む窓の添字保持同型類は、次の十八座標で尽くされる。
四成分の和が零の指標格子で
\[
 \mathcal W=S_4(2,0,-1,-1)\ \cup\ S_4(1,1,-1,-1)
\]
とする。$z=(z_u)_{u\in\mathcal W}$ の非零台がある開半空間
$\langle u,\lambda\rangle>0$ に入るもの（空台を含む）を取り、
$T=(k^\times)^4/k^\times$ の $z_u\mapsto t^u z_u$ と四添字の偶置換で割る。
この対応は横断切片と固有線の基底を選んだ上での同型類の全単射である。
全点が標準Koszul左右AS正則である。極大な台は偶置換まで四つあり、
同型を除いた一般次元は $8,6,6,6$、これらの全台上で安定化群は自明、最小周期は4である。
\end{theorem}
\begin{proof}[横断する十八文字]
$q_{0j}=1$, $(q_{12},q_{13},q_{23})=(i,-i,i)$ と取る。
行積は全て1で、射影テンソルの巡回は量子体積を固定する。
基点は射影的にpolystableである。複素数上の正規直交基底では24個の量子交代係数が
全て絶対値1で、各一位置の四切片の台は互いに交わらずノルムが等しい。
従って定理\ref{thm:walsh-presentation-space}と同じ最小ノルム判定が使える。
四三角積による同型判定から、四座標線の置換で基点を保つものは $A_4$ である
（前節の軌道長2）。位置0での行列が与えられると、他位置の行列は六関係から射影的に一意である。
各許された置換は位置0の純置換行列に取れ、必要な位置別倍率は三角積判定の対角twistで得る。
純置換の積も位置0で決まるので、全射影安定化群は $H=T\rtimes A_4$ である。

文字符号をテンソル語の四文字の個数から $(1,1,1,1)$ を引いたものとする。
\path{computations/verify_quantum_attraction_fiber.py} は前節の接方程式をこの次数で分ける。
文字の型が $(3,-1,-1,-1),(2,0,-1,-1),(1,1,-1,-1),(1,0,0,-1),0$ の順に、
一文字あたりのアフィン接次元の上界は $0,1,1,4,13$ である。
合計79は前節の滑らかな76次元と三量子方向の和に等しく、個々に等号である。
$Z_5$ では零文字だけが3減る。四基底軌道と全体倍率の接では
$(1,0,0,-1)$ の各文字に四位置から四方向があり、零文字に十方向がある。
実際、射影安定化Lie環は対角三次元である。従って横断表現が定理の十八線となる。
十二本の軌道で $A_4$ は単純推移的であり、基底を置換して働く形にできる。
六本の軌道では $(1,1,-1,-1)$ の安定化元 $(01)(23)$ を調べる。
位置 $n$ で文字 $j$ を $i^{n d_j}$ 倍して置換し、$d=(0,0,3,1)$ と取ると
体積とこの変形線の双方の倍率が $-1$ である。従って射影接への倍率は1であり、
この六本にも純置換基底が取れる。計算は法97で行える。固有値は符号または四乗根で、
単根が剰余でも別々なので、得る指標は標数零で同じである。
\end{proof}

\begin{proof}[零ファイバーの全軌道]
閉軌道のLunaのエタール切片定理を、指数の巡回移動も加えた有限拡大の群で適用する。
射影空間では基点で非零の不変斉次式のアフィン開集合を使う。
窓の空間の切片はここで滑らかな $Z_5$ に取り、その切片から上の横断接表現へも
安定化群同変なエタール写像を取る。例えば、安定化群不変な局所アフィン座標で
軌道と横断方向の不変補空間を取り、横断接への射影を用いる。
いずれも基点の周囲でstrongly étale、すなわち不変商の射と元の射との正方形を
ファイバー積にできる（Lunaのエタール切片定理と基本補題）。
二つの商の基点ファイバーへの制限は従って同型である。
特にこれは完備局所環だけでなく、線型表現の零ファイバーとテンソル側の基点ファイバーの
全幾何軌道の対応を与える。飽和開集合は基点ファイバーの全体を含む。
巡回を加えても基点の元の四基底軌道は同じで、安定化群の元の部分群に戻して
添字保持同型だけを取れる。なお周囲の階数スキームでも $w$ を軌道閉包に含む点は、
閉軌道への一径数判定と前節の流入の定理により $Z_5$ に入る。

有限拡大 $T\rtimes A_4$ の線型零ファイバーは $T$ のものと同じである。
トーラスの対角作用では、原点へ退化する条件は非零台が開半空間に入ることである。
実際、その条件を満たす有理余指標を整数倍すれば直接退化でき、
満たさなければ台上で非零の正次数不変単項式がある。
その有限軌道の値を区別する不変量も取れる。成分表現を射影商でさらに商にした一点を
同型類空間とするのではなく、この零ファイバー内の個々の軌道を数えている。
\end{proof}

\begin{proof}[全特殊値の同型、周期と退化]
$\mathcal W$ を辞書式に並べて $0,\ldots,17$ と番号を付ける。
共通の台 $S$ の二点のトーラス同型は、$\mathbb Z^S\to X^*(T)$ の整数核の各基底元 $m$ で
$\prod_{j\in S} z_j^{m_j}$ が一致することと同値である。これは代数閉体上のトーラスの像の方程式であり、
飽和していない指標像でも成立する。先に偶置換全てを試せば台の異なる場合も含む。
一位置巡回は上記順の各座標に $i^{b_j}$ を掛け、向きを一方に取ると
\[
 (b_j)=(3,2,1,1,3,2,2,3,1,0,0,0,0,0,0,0,0,0).
\]
これは各一元接核で語を巡回して体積の符号を除く前掲コードの計算による。
従って $r=1,2,4$ の順に $(i^{rb_j}z_j)$ と $z$ の上の同型検査を行えば全台と全特殊値の最小周期が決まる。
退化も、偶置換後にある余指標が残す台 $S'$ で零、$S\setminus S'$ で正であり、
残す座標が上のトーラス同型条件を満たすことと同値である。
これは対角トーラス軌道閉包の面による分解で、各面上の非零座標は一軌道となる。

極大台の列挙では第四余指標座標を零に取って三空間の中心超平面配置を使う。
各室の符号は、二枚の独立な壁の交線の半直線、その壁内の隣接方向、壁外の方向の順に
無限小変位を取って尽くせる。実際、各室の閉錐は尖っており、その稜と稜を含む面を一組選べる。
壁上の台は周囲の室の台に含まれる。全168室の正台のうち包含極大なのは十一座標が12個、
九座標が36個で、偶置換を除いた代表は
\[
\begin{gathered}
 \{0,1,2,3,4,5,6,7,8,9,10\},\qquad
 \{0,1,2,3,4,5,7,9,12\},\\
 \{0,1,2,3,4,5,9,10,12\},\qquad
 \{0,1,2,3,4,6,8,10,13\}.
\end{gathered}
\]
各代表は十二個の偶置換全てで別の台となり、指標行列の三Smith因子は全て1である。
従って台全体が非零なら安定化群は自明である。各整数核への $b$ の値と4の最大公約数は1で、
台を保つ置換もないので全台上の周期は4である。台の大きさから3を引くと表示した次元を得る。
\end{proof}

% !TeX root = ../papers/nc-p3.tex
\section{量子四面の単純値・単一重複値を持つ全退化ファイバー}
\label{sec:quantum-face-resonance}
$q_{ab}$ は $x_bx_a=q_{ab}x_ax_b$ の規約とし、向き付き四面の値を
\[
 (h_0,h_1,h_2,h_3)=(T_{123},T_{023}^{-1},T_{013},T_{012}^{-1})
\]
とする。積は1で、奇置換は添字の置換に全値の逆数化を伴う。
\begin{theorem}\label{thm:quantum-face-resonance}
量子点を軌道閉包に持つ窓の全添字保持同型類について、次が成立する。
\begin{enumerate}
\item 四つの $h$ が相異なれば、この量子点自身だけである。
\item ちょうど一対 $h_k=h_l=c\ne1$ が等しく、残る $h_i,h_j$ は互いにも $c$ にも異なるとする。
このとき背景の六関係のうち $kl$ の一関係だけを変える、下記の三係数表示が全類を尽くす。
同型は各非零台と、三係数全てが非零なら $C_0C_2/C_1^2$ で決まる。
$c=-1$ の場合に限り台は $s\leftrightarrow2-s$ でも同一視する。
零係数を含む全ての類で最小周期は1、標準Koszul左右AS正則である。
\end{enumerate}
異なる背景間の重複は定理\ref{thm:quantum-indexed-all}の背景同型だけを考慮すればよい。
\end{theorem}

\begin{proof}[四面の値で決まる横断指標]
各行積を1にした量子座標を用いる。四位置のある量子窓に対する一次変形で、
関係を順序二次単項式へ書き直す。関係二文字の多重次数からの差を $\alpha$ とする。
$\sum\alpha_j=0$ で、負成分は $-1$ だけ、負の個数は高々二つである。
負の添字集合を $N$ とする。量子Laurent環で $x_j x^\alpha=\chi_j x^\alpha x_j$ と書く。
位置方向を $\eta^n$ ($\eta^4=1$) に分ける。$a,b$ 関係の一次補正を
$\eta^n c_{ab}x^\alpha x_bx_a$ の順で書けば、$c_{ab}$ は $N\subset\{a,b\}$ でのみ許される。
三文字重なりの線型条件は、$v_j=1-\eta\chi_j$ として $v\wedge c=0$ である。
実際、三文字を逆順から整列する二通りの差で、補正対の外の文字 $j$ を右に置いた項の係数は1、
左に置いた項は位置が一つずれて $x^\alpha$ を越えるため $\eta\chi_j$ になる。
背景の交換係数は共通の順序語で相殺する。
一次の文字変更も同じ計算で $c\mapsto c+v\wedge b$、$b_a$ は $N\subset\{a\}$ のときだけ許される。
この線型条件は少なくとも必要である。四窓の階数を保つ量子点の一次変形は各経路で自由なので
三文字で非零の重なりを生じない。関係の変化を基底変更で消せれば、
四次双対縮約直線の一意性からテンソル側にも別の射影変化は残らない。

負添字の外積を先に固定すると、残る条件と同値は相補添字上の $v$ との外積複体になる。
一成分でも非零ならその双対座標による内部積で可縮である。
従って $\alpha\ne0$ の横断指標が出るには $j\notin N$ の $\chi_j$ が全て同じでなければならない。
$N=\{b\}$ では $\alpha=e_a-e_b$ で、これは $a,b$ を含む二つの三角積が
ともに1となる条件である。定理の二場合では面値1が高々一つなので起こらない。
$N=\{k,l\}$ では $\alpha_s=se_i+(2-s)e_j-e_k-e_l$、$s=0,1,2$ である。
$\chi_i/\chi_j$ は $T_{ijk}T_{ijl}$ またはその逆なので、$\chi_i=\chi_j$ は $h_k=h_l$ に等しい。
このとき $\prod\chi_a=1$ と $\prod\chi_a^{\alpha_{s,a}}=1$ により $\chi_i^4=1$ である。
許されるモードは $\eta_s=\chi_i^{-1}$ のみで、各重みは高々一次元となる。

この三方向は全て同時に実現する。位置 $n$ の $x_lx_k-q_{kl}x_kx_l$ に
$\sum_{s=0}^2 C_s\eta_s^n x_i^s x_j^{2-s}$ を加える（固定順序 $i,j$ を使う）。
新旧各文字を通る全三文字重なりを取る。補正が現れると $k,l$ がともに消えるので
同じ書換えを再使用しない。従って上の線型条件だけで差は零である。
$k,l$ の個数を先に、次に反転数を下げて停止するため、全係数で順序語が基底になる。
量子窓への平坦退化またはこの二次PBW表示と量子外積への双対の随伴次数により標準正則である。
$\eta_s$ の比は連続する $s$ で一定で、四周期の各係数倍率は三指標への共通トーラス作用で実現できる。
従って一位置移動は各係数で同型になる。
\end{proof}

\begin{proof}[候補の尽くし方と同型]
量子テンソルの全24個の置換係数は非零である。この台の対角基底変更に関する重みの重心は
各頂点の四座標で均等なので、その重心は重み凸包の相対内部にある。
よって複素数上で四つの正対角行列により各一文字切断の四行のノルムを同じにできる。
具体的には跡零対数倍率上のノルム平方は有限個の正係数指数関数の和で、
有効な部分では凸包の内部条件から無限遠で発散し最小値を持つ。
異なる行の台は交わらないので全Gram行列がスカラーとなり、量子点はpolystableである。
代数的な閉軌道性は任意の標数零代数閉体に移る。
面値1が高々一つの範囲では四座標点が点スキームから内在的に復元され、
連結安定化群は共通対角トーラスである。異なる閉軌道は同じ商ファイバーを持たない。

三係数族のFrobenius標準四移動は普通の四反復である。
実際、純側 $i,j$ に重み1、$k,l$ に重み0を付けると、係数は全て重み2にできる。
PBW最高次の基底で規格化したNakayama行列は係数に正則で、各成分の重みの絶対値は高々1である。
従って基点から変化できない。最高次テンソルを量子体積の一つの係数で規格化すれば同様に、
全語の重みは0から4、体積語は重み2であるから補正は三係数の一次までである。
ゆえにテンソル族は基点を通るアフィン三平面となる。三方向は負添字が二つで
文字変更による境界がなく、異なる指標であるから軌道接と交わらない。

前節と同じLuna切片を用いるが、今回は周囲の滑らかさを仮定しない。
量子点の有限置換部分もこの三方向の空間を保つ。実際、反転した順序語の並べ替えでは
純 $i,j$ に補正がなく、$kl$ 以外の関係も変わらない。従って安定化群不変な軌道接の補空間にこの平面を含めて、
アフィン線型横断面と階数スキームの交叉で切片を選べる。射影倍率は体積係数の不変な局所規格化で除く。
切片の完備局所環の位相的生成元は、余接空間のトーラス固有基底を持ち上げて取れる。
上の接計算により重み零のもの以外は、単一重複の場合の三平面の座標だけで足りる。
その三座標は平面上で $C_s$ と一致する線型射影を選べる。重み零の生成元は
この平面に制限すると定数零である。純側二文字だけに同じ正の余重みを与えれば三方向とも正だからである。

切片の基点ファイバーの任意の点 $y$ は一径数部分群で原点へ退化できる。
三座標を $C_s(y)$ とし、同じ三値の平面上の点を取る。
前者が原点へ流れる余指標ではその非零座標は全て正なので、後者も同じ余指標で流れる。
二つの弧では重み零生成元が定数零、他の三座標が同じ単項式となるため、
全完備局所関数が一致する。従って二つの弧は一致し、元の点も同じである。
ここで切片の飽和エタール近傍は基点ファイバー全体を含む。横断平面上の全点も原点への
一径数退化を持つため同じ近傍に入る。四面値が相異なる場合は三座標もなく、同じ議論で基点だけが残る。

三指標の整数関係は $\alpha_0+\alpha_2=2\alpha_1$ の整数倍だけであり、
二つ以下では独立である。従ってトーラス商の同型条件は所定の台と単項式比である。
単一重複の場合、偶置換の安定化元は恒等だけである。奇置換では重複値が自分の逆となる必要があり、
$c\ne1$ より $c=-1$ となる。残る二値は相異なる逆数なので可能な置換は $(ij)$ だけである。
これを位置0の純置換とし、位置1で追加倍率 $d_i,d_j$ を使うと
$d_i/d_j=q_{ij}^2$ である。混合単項式の並べ替えの倍率 $q_{ij}$ と合わせれば、
単項式比に掛かる係数 $d_i/(d_jq_{ij}^2)$ は1となる。従ってこの場合も比の値は保たれる。
\path{computations/verify_quantum_face_resonance.py} は面の符号と置換条件を十二乗根格子でも検算する。
\end{proof}

% !TeX root = ../papers/nc-p3.tex
\section{三面の一致する量子点と九係数の中心一次拡大}
\label{sec:quantum-triple-face}
四つの向き付き面値のうちちょうど三つが一致するとする。残る面の添字を0に取る。
共通値を $Q$ とすれば $(h_0,h_1,h_2,h_3)=(Q^{-3},Q,Q,Q)$ で $Q^4\ne1$ である。
向きを反転すると $Q$ と $Q^{-1}$ が交換されるので各背景について一方を固定する。
文字を $t,x_1,x_2,x_3$ とし、以下 $(a,b,c)$ は $(1,2,3)$ の巡回置換を走る。
\begin{theorem}\label{thm:quantum-triple-face}
この量子点を軌道閉包に持つ全窓の表示は
\[
 [t,x_a]=0,\qquad x_bx_c-Qx_cx_b=A_a x_a^2+B_a tx_a+D_a t^2,
 \qquad A_1A_2A_3=0
\]
で尽くされる。全て標準Koszul左右AS正則で最小周期は1である。
$C_{a,0}=A_a,C_{a,1}=B_a,C_{a,2}=D_a$ に指標
$u_{a,s}=se_0+(2-s)e_a-e_b-e_c$ を与える。
二点が添字保持同型となる必要十分条件は、三添字の巡回置換後に非零台 $S$ が同じで、
$\mathbb Z^S\to\mathbb Z^4,\ m\mapsto\sum m_{a,s}u_{a,s}$ の整数核の各基底元で
$\prod C_{a,s}^{m_{a,s}}$ が一致することである。極大な三つの八次元座標平面は巡回で置換され、
指定を忘れた一般同型類の次元は5である。ある点の退化は、非零台のうち余指標の値が
零の項を残し正の項を消す操作と、上記の同型だけで得られる。
ここでの退化は同じ背景のファイバー内で取る。
\end{theorem}
\begin{proof}
$t$ を中心とする量子表示は定理\ref{thm:quantum-indexed-all}の三角積から得る。
各行積は1で、面値1は高々一つである。前節の接計算で、非零横断重みは
共通値の三面から二面を選んだ九つだけで、各重複度は高々1である。
全ての相補対に $t$ が入るため位置モードは1である。

テンソル積記号を一部省略して $L_a=t x_a-x_a t$, $R_a=x_bx_c-Qx_cx_b$,
$F_a=A_ax_a^2+B_a(t x_a+x_a t)/2+D_at^2$ を二文字空間に取る。四テンソル
\[
 W(C)=\sum_a \bigl(L_a\otimes(R_a-F_a)+(R_a-F_a)\otimes L_a\bigr)
\]
は一位置回転で符号を変える。背景部分は順列語の係数を巡回して分かる。
補正部分は各 $s$ で四語の符号付き巡回和の定数倍となる。
$L_a,R_a-F_a$ の六本は独立である。実際、相異なる非中心文字の積への射影で
後の三本の係数が消え、前の三本は別々の文字対にある。
従って両隣接切断は全係数で階数6、関係は表示の六本である。
九つの補正テンソルは非零で相異なる指標を持ち、負の添字が二つで
軌道接に入らない。従ってこれらは全非零横断接方向の基底である。

安定化群は共通対角トーラスと純巡回置換で生成される。
奇置換なら $Q=Q^{-1}$ を要するので起こらない。この群は上のアフィン九平面を保つ。
この平面を含む線型横断面を前節と同様に選ぶ。各座標への線型射影は平面上で $C_{a,s}$ と
一致するようトーラス同変に選べる。切片の完備局所環はこれらと重み零の生成元で位相的に生成される。
零ファイバーの点が余指標で背景へ流れるなら、九座標の非零台はその正半空間に含まれる。
同じ九値の平面上の点も同じ余指標で流れるので、両弧上で重み零生成元は背景での定数値に等しく、
九座標は同じ単項式となる。従って弧は一致する。飽和エタール切片による
全軌道の同定は定理\ref{thm:quantum-attraction-fiber}の証明と同じである。

$u_{1,0}+u_{2,0}+u_{3,0}=0$ なので三つの $A_a$ が同時に非零なら正半空間に入らない。
逆に消える平方係数の添字 $a$ に余重み $-2$、他の非中心二文字に1を取り、
$t$ に十分大きな正整数を与えれば、$a$ 以外の平方と六つの $B,D$ が全て正になる。
従って零ファイバーはちょうど三平面の和である。各点の一径数退化に
命題\ref{prop:unrestricted-good-reduction}を適用すれば全次数の正則性が従う。
表示は全位置で同一である。残る同型・面による退化は対角トーラス軌道の指標格子で決まる。
八座標の台では $u_{a,2}-u_{a,1}=e_0-e_a$ が三つとも取れるため
指標像は全格子であり、台を保つ非自明な巡回もない。一般次元は $8-3=5$ となる。
\end{proof}

\begin{proposition}\label{prop:triple-face-all-regularity}
上記の九係数表示で平方積 $P=A_1A_2A_3\ne0$ としても、二つの等式
$Q^3=-1,\ P=-1$ が同時に成立する場合を除き、全て標準Koszul左右AS正則である。
二等式が成立する場合は $B,D$ によらず指数成長する。
\end{proposition}
\begin{proof}
$t$ の余重みのみを1に取ると $B,D$ が消え、平方項は全て残る。
極限代数は $S[t]$ である。$S$ は対角文字変更によって三平方係数を
共通値 $r$ ($r^3=P$) に揃えた三生成Sklyanin代数である。実際、文字倍率
$d_a$ による $A_a$ の倍率は $d_a^3/(d_1d_2d_3)$ またはその逆で、積1の全組を取る。
通常の巡回関係 $p x_bx_c+q x_cx_b+s x_a^2$ のパラメータは $(p,q,s)=(1,-Q,-r)$ となる。

三生成Sklyanin代数のKoszul性とHilbert級数の定理
（Iyudu--Shkarin, \textit{Three dimensional Sklyanin algebras and Gr\"obner bases}）を用いる。
三係数が非零で三乗が全て等しい場合を除くと、$S$ はKoszulで級数 $(1-z)^{-3}$ を持つ。
左右AS性もここで直接読める。双対の次数は $1,3,3,1$ で、三つの関係 $r_a$ に対して
$\sum_a x_a\otimes r_a=\sum_a r_a\otimes x_a$ は非零な三次交叉の生成元である。
従って一次と二次の両向きの積対合は完全である。$S[t]$ はKoszulで双対は
$S^!$ と一生成外積代数の符号付きテンソル積となる。上の余指標と
命題\ref{prop:unrestricted-good-reduction}で $B,D$ を任意に戻せる。

例外では $A/(t)=S$ の成長を直接計算できる。三係数の三乗が等しいとき、
$x_j$ を $\epsilon^j$ 倍する対角自己同型（$\epsilon^3=1$）による次数twistは
関係の三係数を共通倍率を除き $(\epsilon^{-1}p,\epsilon q,s)$ に変える。
適切な $\epsilon$ で最初の二係数を揃え、先の対角文字変更で第三も揃える。
$(1,1,1)$ の関係空間は $X_i=\sum_{j=1}^3\omega^{ij}x_j$ ($i=0,1,2$、$\omega$ は原始三乗根)
の三平方が張る。実際、$b+c\equiv2a\pmod3$ で各平方は関係のFourier和である。
同じ文字の連続しない全語が基底となり、次数 $n\ge1$ の次元は $3\cdot2^{n-1}$。
twistは次数次元を保つので元の代数にも指数成長する商がある。
\end{proof}

% !TeX root = ../papers/nc-p3.tex
\section{全量子背景の三つの余指標表と有限重み再帰}
\label{sec:quantum-universal-chambers}
前二節で除いた面重複、根1、可換点も含め、軌道閉包に量子点を持つ範囲全体で、
候補の生成・正則性・添字保持同型を扱う。以下の三表では任意の体要素を自由欄に入れ、
線型方程式の右辺だけを検査する。退化先の量子点は定理\ref{thm:quantum-indexed-all}の
三パラメータと有限置換で指定する。

\begin{theorem}\label{thm:quantum-universal-chambers}
任意の量子背景 $w$ を固定する。以下の手順で受理する窓が、軌道閉包に $w$ を持つ
全幾何軌道の代表を尽くす。全て標準Koszul左右AS正則である。文字の番号付けを全て許して、
必要な余指標と正の整数重みの上限は次の三行だけである。
\[
\begin{array}{c|c|c}
 \lambda & \text{正の多重次数の個数}&\text{最大重み}\\ \hline
 (0,1,4,6)&17&13\\
 (0,3,5,9)&17&19\\
 (0,6,8,9)&20&13
\end{array}
\]
表内・表間の同型は、下記の横断面上で一つの四次行列の安定化群作用により判定できる。
非零台や安定化群の一般性を仮定しない。
\end{theorem}
\begin{proof}[安定化群と横断面]
同型と正則性に必要な基底は射影的に取る。まず全行積が1となる量子表示を使う。
さらに二文字 $a,b$ を含む三角積が全て1なら、その文字間の交換係数を1とし、他文字との交換係数も
二文字で一致するようにする。実際、定理\ref{thm:quantum-indexed-all}の証明によれば、
面値1が二つなら一つの二文字組、四つなら可換表示にできる。二文字を揃えた後、行積を
$P_j$ として対角twist倍率の四乗を $P_j^{-1}$ に取る。二文字組では同じ根を選び、
倍率の全積は組の外の一文字の四乗根の選択で1にできる。面値1が高々一つでも同じ行積調整を使う。
これで四位置に同じ対角トーラス $T$ が作用し、背景の量子体積は一位置回転で符号を変える。

四基底群を $G=(\mathrm{PGL}_4)^4$、射影安定化群を $H$ とする。
点対応の最初の射影変換で $H$ は忠実に表される。面値1が0または1個なら座標点とその右対応が内在的である。
2個なら二文字の全射影直線上でも右対応が同じ点であり、
外の二座標点とそれらから各内側座標点への直線上での右対応が残る対角倍率の比を決める。
4個なら右対応は全射影空間で恒等である。従って位置0の行列は、
同一交換係数の組の内部での可逆ブロック、各一文字の倍率、三角積を保つ置換の積で尽きる。
これらは全て持ち上がる。ブロックは四位置で共通であり、置換には量子同型判定の位置別倍率を付ける。
位置0を固定した持上げは他位置でも射影的に一意である。
$T$ は $H^\circ$ の極大トーラスである。

射影基点直線の $H$ 指標を除いてテンソル表現を正規化する
（線型表現は原空間と基点直線の双対との積に誘導される）。
$H$ は簡約的なので $H$ 不変線型形式 $f(w)=1$ と、
$f=1$ の中で $G$ の射影軌道接に横断する $H$ 不変補空間 $N$ を取れる。
巡回も同時に扱うときは一位置回転を加えた有限拡大の安定化群で選ぶ。
これらの選択は有限線型代数で可能である。例えば $f$ と軌道接への恒等な射影について、
既知ブロック群のLie環と有限置換、必要なら回転との可換性を線型条件として解く。
軌道接は四位置で各行列単位を $w$ に作用させ、射影方向へ写せば求まる。

二隣接階数が6の局所閉スキームを $Z$ とする。量子点のpolystabilityは
定理\ref{thm:quantum-face-resonance}の台の重心の証明が全 $q$ に適用できる。
Lunaの定理を $w$ で非零の不変斉次式のアフィン開集合に適用する。
$S=(w+N)\cap Z$ の基点ファイバーから $G$ 側の全基点ファイバーへ、
$H$ 軌道と $G$ 軌道の全単射を得る。詳しくは、周囲の射影空間における
$G\mathbin{\times}^H(w+N)$ の局所写像が基点でエタールであり、$Z$ で基底変換して切片とする。
強エタールな飽和開集合に縮めると商との正方形がファイバー積であり、閉点の商ファイバー全体を含む。
従ってこの主張は形式的な弧だけに限定されない。
逆に横断面で $w$ に流れる全点は、使った各 $H$ 不変開集合の中に入る。
Hilbert--Mumford判定と極大トーラスの共役性から、$H$ の作用後に
$n\in N$ が $T$ のある余指標で正の成分のみを持つことが必要十分である。
\end{proof}

\begin{proof}[三表の全被覆]
語 $v$ の指標は $\beta(v)=(\#_0v,\ldots,\#_3v)-(1,1,1,1)$ である。
非零な34通りの超平面を、定数ベクトルを除いた余指標の三空間で使う。
壁上の正台は隣の室の正台に含まれるため包含極大な正台のみを残す。
列挙は整数演算で行える。各法線の第四座標を除き、独立な順序付き二法線 $a,b$ ごとに
$r=a\times b,\ e=a\times r$ と置く。各符号 $i,j,k\in\{\pm1\}$ に対して、
別の法線 $v$ と $(ir,je,ka)$ との三内積のうち最初の非零値が正かを記録する。
これは一つの半直線、その上の二次元面、その外側の室を通る微小変位であり全室を尽くす。
実際、法線が三空間を張るので各全次元室の閉錐には稜があり、その稜を含む面を選べる。
重複を除くと336室、包含極大は17文字の48個と20文字の24個となる。
\path{computations/verify_quantum_weight_recursion.py} はその72個が表の三余指標の全置換に
一致することを確認する。余指標を極大台の代表へ換えても元の全非零成分は正のままである。
\end{proof}

\begin{proof}[重み順の復元と判定]
選んだ室の正の $N$ の座標だけを使う。各切断 $M$ の行列を順序対で表示する。
行・列とも $i<j$ の六対 $P$ を主部分に取る。量子背景では各行は相補対の列にちょうど一つの
非零体積係数を持ち、$M_{PP}$ は可逆である。この行列式の指標は零である。
実際、六対の文字総数は各添字3で、行と列を足して六つの体積指標を引くと零となる。
正の変位だけを許すと行列式は従って背景値のままである。残る十対を $E$ と置く。
求める方程式は各切断の百成分
\[
 B(M)=M_{EE}-M_{EP}M_{PP}^{-1}M_{PE}=0
\]
である。分母が非零定数なのでこの表の上では多項式となる。
行 $u$、列 $v$ の成分の指標は語 $uv$ の指標である。
これは六次主小行列を縁取る七次行列式をその主行列式で割る式からも分かる。

二切断を合わせた $B$ の基点での微分を $L$ とする。各多重次数 $\beta$ ごとに
$N_\beta=K_\beta\oplus D_\beta$, $K_\beta=\ker L_\beta$ と線型分解する。
整数重み $d=1,\ldots,\max\langle\beta,\lambda\rangle$ を順に処理する。
各 $\beta$ の自由欄に $z_\beta\in K_\beta$ を指定し、既に復元した
$n_{<d}$ を $B(w+n_{<d})$ に代入する。その $\beta$ 成分の負を $b_\beta$ とする。
$b_\beta\notin\operatorname{im} L_\beta$ なら棄却し、そうでなければ
\[
 n_\beta=z_\beta+(L_\beta|_{D_\beta})^{-1}b_\beta
\]
を加える。$N_\beta=0$ のときにも右辺の整合性を検査する。
多項式の非線型項に重み $d$ 以上の入力があれば、もう一つの正の入力を要するため
出力の重みは $d$ より大きい。従ってこの処理は重み $d$ の全式を厳密に解く。
非正重みの式は背景の値であり、最後の重みの後に方程式は残らない。
帰納的に、受理する各欄の値がただ一つのテンソル $w+n$ を与え、表内の全解がこの形となる。
両階数は常に6で、背景への一径数退化に命題\ref{prop:unrestricted-good-reduction}を適用できる。

異なる背景は定理\ref{thm:quantum-indexed-all}によって先に比較する。背景を同じ表示に揃えた後、
二つの復元値が同型となる条件は $n'=hn$ ($h\in H$) である。
有限置換を試し、残る四次のブロック行列に逆行列式用の一変数を付けて
多項式イデアルが単位イデアルかを消去すれば、非可換な安定化群でも判定が終わる。
周期は同じ検査で一位置・二位置の巡回を比較して決まる。
同じ商ファイバー内での軌道閉包も切片に移せる。ブロック群の各成分で
$y=hn$ のグラフから群変数を消去して得るイデアルを $n'$ に代入すればよい。
座標環の消去定理によりこの操作が像のZariski閉包を与える。トーラスだけの場合は
既出の整数核と台の面の検査に簡約する。

検算では三面一致の有理量子点で九方向のうち同じ余指標に正の全方向を同時に復元した。
可換点では、接空間の自由欄で実際に棄却する入力と、非線型な強制補正を一回
要して受理する入力の両方を有理数上で照合した。
\end{proof}

% !TeX root = ../papers/nc-p3.tex
\section{二対の非単位面値と三次の同時変形障害}
\label{sec:quantum-double-pair}
面値を $(c,c,d,d)$, $c\ne d$, $c,d\ne1$ とする。$\gamma=cd\in\{1,-1\}$ である。
両対を区別して高い側の文字を $x,y$、他方を $z,w$ と呼び、後の対では純平方の文字を $z$ に取る。
この選び方を全て許す。$a^2=d$ の一つを固定すれば、量子背景の三角積判定により
\[
 yx=xy,\quad zx=a xz,\quad zy=a^{-1}yz,\quad
 wx=a^{-1}xw,\quad wy=a yw,\quad wz=\gamma zw
\]
を使える。
\begin{theorem}\label{thm:quantum-double-pair}
この量子点へ退化する全窓は、下記の通常の次数付きOre表示で尽きる。
$yx=xy+D z^2$ とし、$zx,zy$ の二関係は背景の式を用いる。
$w$ による拡大は $w v=\sigma(v)w+\delta(v)$,
$(\sigma(x),\sigma(y),\sigma(z))=(a^{-1}x,ay,\gamma z)$ である。
\[
 \delta(z)=Ax^2+Bxy+Cy^2,\qquad
 \delta(x)=u xz+v yz,\quad \delta(y)=s xz+t yz,
\]
\[
 s=\frac{DA(1+\gamma a^2)}{a-a^{-1}},\qquad
 v=\frac{DC(1+\gamma a^{-2})}{a-a^{-1}} .
\]
$\gamma=1$ では条件は $BD=0$、$u=t=0$ とする。
$\gamma=-1$ では $A,B,C,D$ は全て自由で、$u=DB/(2a)$, $t=-DB/(2a)$ とする。
全類で最小周期は1であり、特殊値も標準Koszul左右AS正則である。
\end{theorem}
\begin{proof}[四欄の全被覆]
非零な横断接指標は前の面値計算により二つの三つ組だけであり、各々一次元以下である。
その三つ組は $x,y$ の二次単項式の次数から $e_z+e_w$ を引いたもの、
および $z,w$ の二次単項式の次数から $e_x+e_y$ を引いたものである。
それぞれを一組だけ使用した全三方向は一関係の単独書換えで実現し、各次元は1である。
実際、$wz$ の補正だけなら $\delta(z)$ は $x,y$ の任意の二次式でよい。
$yx$ の補正だけなら二対を交換した表示を使える。
両対の余重みの平均を比較し、$x,y$ 側を高く取ると、低い側の三つ組で正となり得るのは
余重みが最大の文字の純平方だけである。平均が等しい場合にも同じ包含が成立する。
従って流入点の横断接座標の非零台は $A,B,C,D$ の四欄に入る。

この場合、余重みは改めて $(1,1,2,-2)$ にできる。実際、切片の接空間への
トーラス同変線型射影 $\pi$ を接空間上の恒等と取ると、完備局所環はその座標で位相的に生成される。
元の流入弧で、重み零座標は背景値に等しく、他の非零座標は四欄だけである。
各テンソル座標をこの弧上の単項式で表すと、元の余重みが正なので各係数への寄与は有限で、
その多重次数は使った四欄の多重次数の非負和に限られる。
四欄はいずれも新しい余重みで2であり、従って新しい余指標でも同じ点が流入する。
同じ四値を持つ切片点は元の完備座標での弧の一致により一意である。

背景を通常の次数付き表示から標準四窓表示へ移す位置別対角変換は固定してよい。
この同定で、各欄の一次変形は上の通常関係の対応する係数で測る。
一つずつの補正には文字変更による一次境界がない（負添字が二つ）ので、
四係数と四接座標の対応は独立な非零倍率だけである。
通常のOre表示から得る族を切片上に戻したときも、この一次対応がそのまま
余重み2の四座標を与えることを確認しておく。
各次数で順序語が基底の族では関係と双対の最高次直線から標準四窓を基点の近くで正則に取れる。
最高次の倍率は体積の背景係数、四移動の基底は双対積で規格化でき、対角文字作用と可換する。
周囲の軌道接のLie補空間を対角同変に選ぶと、
その形式指数による基底変更とアフィン横断面の積から周囲への写像は一次で同型である。
従って形式逆を用いて切片点と基底変更を同変に取れる。
入力が全て余重み2なので余重み2の出力には線型項しか入らない。
所定の四値で得る二つの流入弧がこの形式同型で対応すれば、その一般点は同型である。
定数の二窓に戻した同型方程式は $k((u))$ で解を持つため、代数閉体 $k$ でも解を持つ。
以上により全四値のOre族を構成できた場合、または必要な障害を除いて構成できた場合に、
それ以上の窓は残らない。
\end{proof}

\begin{proof}[重なりと障害の必要性]
三文字部分は $z$ を先に置いた反復Ore拡大である。例えば $zx=a xz$ の量子平面へ
$y$ を加え、$yx=xy+D z^2$, $yz=a zy$ を用いれば基底ができる。
$\sigma$ はこの三関係を保つ。
途中では $\delta(z)=Ax^2+Bxy+Cy^2+H z^2$ を許す。
$\delta$ を順に $zx-a xz$, $zy-a^{-1}yz$, $yx-xy-Dz^2$ に適用して順序語の係数を取ると
\begin{align*}
 DB+H(a^2-1)+a(\gamma-1)u&=0,& DC(1+a^{-2})+(\gamma/a-a)v&=0,\\
 -DA(1+a^2)+(\gamma a-a^{-1})s&=0,&-DBa^{-2}+H(a^{-2}-1)+(\gamma-1)t/a&=0,\\
 (a-a^{-1})s&=DA(1+\gamma a^2),&(a-a^{-1})v&=DC(1+\gamma a^{-2}),\\
 (a-a^{-1})(t+u)&=DB(1+\gamma),&D\{a(t+u)-(1+\gamma)H\}&=0 .
\end{align*}
$\gamma=-1$ では定理の値と $H=0$ で全式が消える。
$\gamma=1$ では $H=-DB/(a^2-1)$ と $t+u=2DB/(a-a^{-1})$ を最後に代入すると
$2D^2B(a^2+1)/(a^2-1)=0$ が必要である。$a^2=1$ は非単位条件、
$a^2=-1$ はこの符号の場合の $c\ne d$ に反する。従って $DB=0$ で、
この条件下では定理の $H=u=t=0$ が全式を満たす。

$\gamma=1$ において最後の障害がOre表示の仮定を外しても必要であることを示す。
この場合普通表示の行積は1であり、背景体積は符号付き巡回不変である。
一関係だけの各一次変形も位置モード1である。実際、二つの相補文字の内部の交換係数は1で、
他の二文字との交換係数の積も1であるから、前節の $1-\eta\chi_j$ で $\eta=1$ となる。
巡回同変な切片と接射影を使えば、巡回した点も同じ四座標を持つ。
流入弧の一意性によりその点自体が巡回不変で、六関係を全位置で同じ表示に取れる。
順序を $x,y,z,w$ とし、逆順対を左辺とする六関係を行消去で規格化する。
必要なPlücker座標は背景で非零であり、多重次数を持つので流入点でも非零である。
各補正は四欄の重み2の非負和を持つ。二文字の余重みを並べれば、
$yx$ に加え得るのは $z^2$、$zx,zy$ には何もなく、$wx,wy$ には $xz,yz$、
$wz$ には $x^2,xy,y^2,z^2$ だけである。従ってその六関係は $H$ を許した上記の一般式に尽きる。
三文字の順序語は独立である。量子極限への各経路は命題\ref{prop:unrestricted-good-reduction}
により自由で、特殊点の順序語基底が持ち上がるためである。
従って $w$ と各三文字関係の重なりから八式を必ず満たす。
これで必要性を得た。逆に定理の係数では実際に捩れ導分を定めたのでOre拡大の
順序語が全次数で基底である。随伴次数を量子多項式環にする重み（または反復する
斉次Ore拡大の線型解消）から標準Koszul左右AS正則性が従い、表示は一位置周期である。
\end{proof}

\begin{corollary}
全特殊値では、二つの三つ組の非零台と、三項が揃った三つ組での
「両平方係数の積を混合係数の二乗で割った値」が全トーラス不変量である。
さらに $\gamma=1$ なら二組を交換する四巡回とその冪、
$\gamma=-1$ なら両組内の同時反転で台を同一視する。
同じファイバー内の退化はこれらの置換と台の露出面への制限で尽くされる。
$\gamma=1$ の極大台は指定を忘れると二型で同型類次元は1と0、
$\gamma=-1$ でも二型でどちらも1である。
\end{corollary}
\begin{proof}
正の側の三指標は端の和が中央の二倍という関係だけを持ち、他方の一純平方は
その有理線型包の外にある。従って整数核の不変量は述べた一比だけである。
面の向き付き置換条件から偶置換は恒等と $(xy)(zw)$ だけで、
奇置換が許されるのは $d=c^{-1}$ のときだけである。そのときの $(xzyw)$ は上記の
交換係数も保つ。各置換による左辺の並べ替えの共通倍率は比から消え、
混合単項式の並べ替えの倍率は $1$ または $-1$ で二乗が1である。
より高い余重みの項は四欄の値に寄与しない。切片ではこれらの座標と全余重み零生成元で
流入弧が決まるので、軌道閉包への面の条件もトーラスと同じである。
正符号では三項だけの台と、二平方および反対側の一平方の台が極大で、
後者の三指標は独立である。負符号では三項と反対側の一平方を持つ台が極大で、
三項をどちら側に置くかは同一視されない。
\end{proof}

% !TeX root = ../papers/nc-p3.tex
\section{正則正規な一次直線列を持つ窓と楕円三次商}
\label{sec:regular-normal-linear}
量子点への退化とは別の切り口として、一次の正規元を用いる。
正則正規な直線列とは、各 $i$ の非零元 $t_i\in C_{i,i+1}$ が、全 $i\le j$ で
\[
 C_{ij}t_j=t_i C_{i+1,j+1}
\]
を満たし、左右の乗法が単射であるような直線 $kt_i$ の列とする。周期は仮定しない。

\begin{theorem}\label{thm:regular-normal-linear}
この直線列を持つ標準14641型の添字代数は全て最小周期1である。
直線列付きの全同型類は、三生成三関係の標準Koszul左右AS正則代数
$S=T(V)/(R)$ と二写像 $\alpha:R\to V$, $\beta:R\to k$ から次のように与えられる。
$0\ne w\in(R\otimes V)\cap(V\otimes R)$ を取り、
\[
 h=(\alpha\otimes1-1\otimes\alpha)w,\qquad
 b=(\beta\otimes1-1\otimes\beta)w.
\]
採用する条件は $h\in R$, $\alpha(h)=-b$, $\beta(h)=0$ の三つである。
中心文字 $t$ を加え、$r=t\alpha(r)+t^2\beta(r)$ ($r\in R$) を課して得る代数を $A$ とする。
対象は $\check A$ と $kt$ である。二組の直線列付き同型は、
$(M\otimes M)R=R'$ となる可逆写像 $M:V\to V'$ と $e\in V^*$ が存在し、
全 $r\in R$, $r'=(M\otimes M)r$ について
\begin{align*}
 M\alpha(r)&=\alpha'(r')+(e\otimes M+M\otimes e)(r),\\
 \beta(r)+e(\alpha(r))&=\beta'(r')+(e\otimes e)(r)
\end{align*}
となることと同値である。
\end{theorem}
\begin{proof}
$t_i$ の代表を選ぶ。正規性と単射性により $a t_j=t_i\phi(a)$ で定まる写像
$\phi:C_{ij}\to C_{i+1,j+1}$ は全単射である。積 $ab$ にこの等式を二回適用すれば
$\phi(ab)=\phi(a)\phi(b)$、また $\phi(t_i)=t_{i+1}$ である。
従って標準四移動との可換性を先に仮定せず、$\phi$ の反復で通常の二次代数 $A$ の
反復表示が得られる。$t$ は中心で正則となる。
$S=A/(t)$ は三生成二次代数で、中心元の短完全列からHilbert級数 $(1-z)^{-3}$ を持つ。
Koszul性は変換環のTorスペクトル系列で従う。$\operatorname{Tor}^A_q(k,S)$ は
$q=0,1$ にだけあり、それぞれ $k,k(-1)$ である。もし $\operatorname{Tor}^S_p(k,k)$
に最初の非線型次数があれば、下段のその成分に入る微分はなく、出る可能性のある
$d_2$ は $p-2$ 列の線型な上段に入るのでそこにも消えない。$A$ のKoszul性に反する。
二次双対の級数は $(1+z)^3$ となり、大域次元は3である。
中心正則元に対する変換環Ext系列
（$\operatorname{Ext}^q_A(S,A)$ は $q=1$ の $S(1)$ だけ）を両側で使えばAS双対性も降りる。

$V$ を $A_1$ に持ち上げる。$A_2$ の次元から、中心三関係以外は $r$ の三つの持上げで尽き、
右辺には $tV$ と $kt^2$ の係数だけが現れる。$A$ は $U=A/(t-1)$ の
語長フィルトレーションのRees代数で、$\operatorname{gr} U=S$ である。
実際 $a\in A_n$ を $a t^{-n}$ として $A[t^{-1}]_0$ に入れると、単射な次数片の
増加和が $U$ となり、その逐次商が $S_n$ である。逆向きには
$U=T(V)/(r-\alpha(r)-\beta(r))$ がPBW、すなわち $S\to\operatorname{gr} U$ が同型なら、
上の斉次表示はそのRees代数である。各 $n$ で $S_n$ の語基底とその低次の $t$ 倍が
斉次表示を張り、Rees代数で独立だからである。
Koszul基底に対するBraverman--Gaitsgoryの二次PBW定理によれば、PBW性は
三次重なり上の表示した三条件に同値である。符号は $r-\alpha(r)-\beta(r)$ の規約による。
交叉は $S$ の三次双対成分の双対で一次元なので $w$ 一つで足りる。
ここで要求するPBW性は、三生成商が非可換でも成り立つ随伴次数の等式である。

PBWの場合は再び $t$ が正則である。上記Tor系列を逆に使うと、$S$ がKoszulなので
両段は内部次数と全ホモロジー次数が等しく、$A$ もKoszulである。
Hilbert級数とExt系列から両側AS、大域次元4、標準型を得る。
最後に、直線を保つ添字同型では頂点別スカラーを選んで全ての $t_i$ の像を $t'_i$ とできる。
中心性と正則性からこの同型は $\phi$ と可換し、位置によらない次数付き同型となる。
一次の持上げは $x\mapsto Mx+e(x)t'$ である。関係に代入し、$t'V'$ と $kt'^2$ の
二つの独立な係数を比較すると定理の式となる。逆もこの代入と可逆性による。
\end{proof}

\begin{theorem}\label{thm:elliptic-central-line}
上記の三生成商の二文字点対応が、滑らかな平面三次曲線 $E$ 上の自己同型
$\sigma$ のグラフであるとする。このとき正則正規一次直線列は一意であり、
定理\ref{thm:regular-normal-linear}の判定は直線の指定を忘れた添字保持同型にも必要十分である。
二つの商の間で試すべき射影行列は、二曲線の変曲点から有限個に絞れる。
\end{theorem}
\begin{proof}
$L=\mathcal O_E(1)$ とする。$\sigma^*L\simeq L$ なら、$\sigma$ は平面の射影変換に延びる。
このとき各二文字関係をその平面グラフ上で評価すると、$E$ 上で消える平面二次式である。
三次曲線のイデアルには非零二次式がないので、全てが平面グラフ上で零になる。
点対応の第一射影が $E$ という仮定に反する。

もう一つの直線列の生成元を商へ写して $v_i\in S_1$ とする。
$S_1 v_{i+1}=v_i S_1$ である。一方が零なら、もう一方も零である。
実際グラフ上での一次評価は単射で、$L$ と $\sigma^*L$ の全切断は基点を持たない。
全て非零なら、同じ等式を $E$ 上で評価して、
二つの積の線型系の共通零因子は $\sigma^*\operatorname{div}(v_{i+1})$ と
$\operatorname{div}(v_i)$ になる。各残り因子の線型系に基点がないため
零因子は重複度を込めて等しい。従って $\sigma^*L\simeq L$ となり、上と矛盾する。
よって全て零であり、元の直線列は各次数で $kt$ そのものである。
商の次数付き同型は滑らかな三次曲線間の射影同型を誘導する。
九つの変曲点から射影枠を選び、その像の順序付き候補を試して $R$ の像も検査すればよい。
実際、変曲点中に非共線な三点があり、その三辺上には合計で高々六点しかない。
\end{proof}

\begin{corollary}\label{cor:elliptic-central-linear-slice}
任意の滑らかな楕円三次商 $S$ を固定する。$e\mapsto D(e)$,
$D(e)(r)=(e\otimes1+1\otimes e)r$ の像に対して
$\operatorname{Hom}(R,V)$ の線型補空間を取り、$\alpha$ をこの補空間に限ることができる。
そこで $h(\alpha)\in R$ を課した各 $\alpha$ に対し、
\[
 \big(b(\beta),\beta(h(\alpha))\big)=\big(-\alpha(h(\alpha)),0\big)
\]
という線型系の全解が、採用する定数欄 $\beta$ を与える。異なる解の同型は
前定理の有限個の射影行列と重み $(1,2)$ の一斉倍率で尽くされる。
従って原点以外では、重み付き射影空間内の重み1、2、3の方程式と有限群の作用で全特殊値を扱える。
\end{corollary}
\begin{proof}
$D(e)=0$, $e\ne0$ と仮定する。全 $r$ で $r(e,e)=0$ なので $([e],[e])$ は
二文字点対応上にある。平面の対角に制限した式 $r(x,x)$ は全て $e$ で一次部分が零だから、
対角とのスキーム交叉のこの点の接次元は2である。滑らかな曲線のグラフの閉部分スキームで
接次元が1を超えることはできない。従って $D$ は単射であり、前定理の持上げ変更で
各 $\alpha$ をただ一回だけ補空間に入れられる。残る二条件は $\alpha$ を固定すると
まさに示した $\beta$ の線型系となる。
\end{proof}

\begin{corollary}\label{cor:sklyanin-central-classification}
滑らかな三次曲線を点対応に持つ通常の三生成Sklyanin商を固定し、
巡回関係を $r_a=p x_bx_c+q x_cx_b+s x_a^2$ と書く（$s\ne0$）。
この商に属する全窓の同型類は一点と重み付き射影空間
$\mathbb P(1,1,1,2,2,2)$ の有限群商である。ここでは正則正規一次直線列を通る商を指定する。
その直線列は一意なので、得る類の一般次元5は何も指定しない添字同型類の次元でもある。
全特殊値が正則であり、全て最小周期1、いずれも同じ $\check{S[t]}$ に平坦特殊化する。
\end{corollary}
\begin{proof}
$w=\sum_a x_a\otimes r_a=\sum_a r_a\otimes x_a$ を取れる。
$h=\sum_a(\alpha(r_a)\otimes x_a-x_a\otimes\alpha(r_a))$ は交代テンソルである。
$s\ne0$ より $R$ の基底は三つの純平方係数で区別でき、$R\cap\bigwedge^2 V=0$ である。
従ってPBW条件は $h=0$ に等しい。すなわち $\alpha(r_a)=\sum_j m_{aj}x_j$ の
$3\times3$ 係数行列が対称であればよく、三つの $\beta(r_a)$ は全て自由である。
一次の持上げの変更 $e$ は対角三係数にそれぞれ $2s e(x_a)$ を加減する。
従ってちょうど一回の変更で全対角を消せ、非対角三係数と定数三係数が残る。
文字を一斉にスカラー倍するとこれらの重みは1と2である（倍率の逆を選んでも同じ商となる）。
$\operatorname{Aut}_{\mathrm{gr}}(S)$ のスカラーを除いた群は前定理により有限である。
この群の作用は各射影行列の持上げを定理の二式に代入し、再び対角を消せば得られ、
持上げの違いは既に除いた一斉倍率だけである。原点以外の倍率軌道は上記重み付き
射影空間の点に対応する。原点への特殊化は $(\alpha,\beta)\mapsto
(u\alpha,u^2\beta)$ で与えられる。各次数で商の基底の持上げの $t$ 倍が族の基底となるため平坦である。
\end{proof}

このSklyanin商では全次数正則性を三次PBW定理によらず確かめることもできる。
$L_a=t x_a-x_a t$,
$P_a=r_a-(t\alpha(r_a)+\alpha(r_a)t)/2-\beta(r_a)t^2$ とする。
$\sum(L_a\otimes P_a+P_a\otimes L_a)$ は一位置回転で符号を変える。
基底部分では $w$ の巡回性を使う。$m$ が対称であるため、線型補正は
$v\in V$ に対する $(tv-vt)\otimes(tv+vt)+(tv+vt)\otimes(tv-vt)$ の線型結合であり、
定数補正も同様に符号付き巡回和となる。六関係は $R$ への射影と三中心関係から独立なので
両隣接階数は6である。$t$ だけに正の余重みを与えて全体倍率を除くと $\check{S[t]}$ に流れる。
命題\ref{prop:unrestricted-good-reduction}が任意の $m,\beta$ に適用できる。

% !TeX root = ../papers/nc-p3.tex
\section{正規一次直線列の四窓からの有限認識と整三次商}
\label{sec:normal-line-recognition}
前節の分類枝に属するかどうかを、直線列の指定も正則性の事前認証も使わずに調べる。
まず三生成商に必要な古典的な幾何的正則性判定を、本稿の関係空間の規約で示す。

\begin{lemma}\label{lem:ternary-finite-graph-test}
$S=T(V)/(R)$, $\dim V=\dim R=3$ とする。標準三次元Koszul左右AS正則性の必要十分条件は次である。
\begin{enumerate}
\item $(R\otimes V)\cap(V\otimes R)$ に、左右の一文字縮約がともに $R$ を張る元 $w$ がある。
\item $R$ の三つの双一次式の左右評価行列は射影平面の各点で階数2以上であり、
両行列式はともに恒等零か、同じ三次曲線をスキームとして定める。
\end{enumerate}
\end{lemma}
\begin{proof}
必要性は三次双対のFrobenius対合と補題\ref{lem:ternary-point-graph-growth}である。
十分性を示す。二次小行列が単元の開集合で核直線を取ると、三式の共通零点スキームは
両側の行列式軌跡上のグラフであり、両射影は同型である。平面全体ならその自己同型は
射影線型で、関係空間は平面グラフの双一次式全体となる。これは可換多項式環の次数twistである。

三次曲線 $E$（非被約でもよい）の場合、グラフ移動を $\sigma$、$L_j=(\sigma^j)^*\mathcal O_E(1)$、
$L_{[i,j)}=\bigotimes_{i\le h<j}L_h$ と置く。グラフは $\mathbb P^2\times\mathbb P^2$ で
余次元3の三式による完全交叉である。Koszul層解消を $(1,1)$ で捩ると、
$\mathcal O(-1,-1),\mathcal O(-2,-2)$ の全コホモロジーが消えることから
二文字の評価は $H^0(L_{[0,2)})$ へ全射で核はちょうど $R$ である。
非空な積 $L_{[i,j)}$ では $H^1=0$ である。実際 $L_i$ では平面三次式の完全列による。
新しい因子の基点自由な切断から全随伴点を避けるものを選ぶと、それによる乗法は単射で
余核の台が零次元なのでこの消滅は伝わる。同じ完全列と各 $L_j$ の次数3により
$h^0(L_{[0,n)})=3n$ ($n>0$) となる。

$D=\bigoplus_{n\ge0}H^0(L_{[0,n)})$ に $\sigma$ を用いた切断の積を入れる。
定理\ref{thm:maximal-tangent-complete}に用いた滑らかな場合と同じく、生成元・関係・
$w$ の複体を $E$ 上で評価すると $n\ge3$ で
\[
 0\to L_{[3,n)}\to R\otimes L_{[2,n)}
 \to V\otimes L_{[1,n)}\to L_{[0,n)}\to0
\]
が完全となる。合成は関係と三次交叉で零で、幾何学的各点の階数は $1,2,1$ である。
階数の小行列によって複体は局所的に分裂するので非被約点でも成立する。
大域切断では $n=3$ における $H^1(\mathcal O_E)=k$ だけが
$V\otimes H^0(L_{[1,3)})$ の位置のホモロジーを与える。他の正の次数ではここは完全で、
$R\otimes H^0(L_{[2,n)})$ の核は常に $w$ の像である。よって $D$ は一次生成である。
右単純加群の極小解消における最初の二微分は $V,R$ からなる。第三の微分の像には
$w$ が含まれるので、$D$ をテンソルした今の計算により
$\operatorname{Tor}_1^S(k,D)=k(-3)$, $\operatorname{Tor}_2^S(k,D)=0$ である。
従って左加群 $D$ の極小解消は $0\to S(-3)\to S\to D\to0$ となり、
左端は三次元 $g$ の右乗法で単射である。$H_S(z)=(1-z)^{-3}$ が従う。
$D$ に左socleはない。基点自由な $L_0$ の非零因子切断を掛ける写像は、
$\sigma$ による引戻しを合わせても単射である。従って $S$ の左socle元は $a g$ と書け、
$g$ による単射性から $a$ も左socleに属する。次数降下で零を得る。反対側でも同じである。
$w$ の端の縮約が全階数なので通常の長さ3のKoszul複体の左端は単射、
生成元と関係で右側は完全である。Hilbert級数のEuler等式から残る中間項も完全、
双対の相補積は両端の縮約から完全である。
\end{proof}

\begin{theorem}\label{thm:finite-normal-line-recognition}
標準四テンソル表示の二隣接切断が階数6とする。以下の有限条件は、
この窓が標準Koszul左右AS正則で正則正規一次直線列を持つことと同値である。
関係を $R_i\subset V_i\otimes V_{i+1}$ と書き、$i$ は法4とする。
\begin{enumerate}
\item 各 $V_i$ に非零 $t_i$ を取り、
$(V_i\otimes V_{i+1})/R_i$ 内の $t_i V_{i+1}$ と $V_i t_{i+1}$ が
同じ四次元空間になることを課す。
\item $x t_{i+1}=t_i g_i(x)$ で一意に定まる可逆行列について
$(g_i\otimes g_{i+1})R_i=R_{i+1}$ を四位置で検査する。
\item $g_0$ で最初の関係を一つの四次元空間の二乗に移し、$t=t_0$ で割る。
三生成商の関係次元3、補題\ref{lem:ternary-finite-graph-test}の二条件、
および持上げの $\alpha,\beta$ に対する定理\ref{thm:regular-normal-linear}の三条件を検査する。
\end{enumerate}
全ての正規直線列はこの四直線の解として現れる。従って二つの認証済み対象の
直線を指定しない同型は、一方の印を一つ固定し、他方の四直線の解と
定理\ref{thm:regular-normal-linear}の二式を同時に解いて判定できる。
\end{theorem}
\begin{proof}
四直線が得られれば $g_i$ は一移動の同型を与え、$g_i(t_i)=t_{i+1}$ である。
これに沿って全位置の基底を運べば、最初の六関係の反復で $t$ が中心になる。
三つの独立な中心関係以外の持上げを取って前節の定理を使うと十分性が従う。
判定は有限な多項式の等式・非消滅だけでよい。例えば射影点での階数1以下の不在は
各アフィンチャート上の小行列式のイデアルで検査できる。

必要性について、任意の印は前節で得る中心正則元 $t$ の形にすることができるが、
四個で十分とするにはFrobenius標準四移動がこの印を保つことを示す必要がある。
$A\to S=A/(t)$ の自然な二次双対写像 $S^!\to A^!$ は単射である。
実際、前節の変換環Tor系列では両段が線型で全微分が零であるから下段への辺準同型
$\operatorname{Tor}^A_i(k,k)\to\operatorname{Tor}^S_i(k,k)$ は全射である。
双対を取りKoszul双対との自然な同定を使えばよい。特に $S^!_3$ の生成元には
非零像 $u\in A^!_3$ がある。$t^\perp=S^!_1$ は左右とも $u$ を消し、
$A^!$ のFrobenius完全対合から両零化空間は正確に $t^\perp$ である。
$A^!$ のNakayama行列はこの超平面を保つ。実際、
$\operatorname{tr}(v u)=\operatorname{tr}(u\nu(v))$ を使えばよい。
よって逆双対で得る標準四移動は $kt$ を保つ。
代表元の倍率は自由なので厳密四周期の線として記録できる。
これで多くの印を持つ場合にも全印が有限な四直線検査から漏れない。
\end{proof}

\begin{proposition}\label{prop:integral-singular-central}
整な節点・尖点三次曲線のグラフを点対応に持つ商でも、前節の一次直線列は一意である。
持上げ変更 $e\mapsto D(e)$ は単射で、補正の重み1、2、3の方程式と有限射影同型検査がそのまま使える。
節点商の射影自己同型候補は高々六つ、尖点商では恒等のみである。
\end{proposition}
\begin{proof}
滑らかな場合の零因子の証明には超平面切断が有効Cartier因子であることだけで足り、
整三次曲線ではこれが成立する。$\sigma^*L\not\simeq L$ も同じく二次式の制限から従う。
$D(e)=0$ なら対角との交叉の接次元が2となるので $e$ は特異点で、
しかもグラフの接空間が対角内にある。従って $d\sigma_e=1$ である。

節点の場合、正規化で二枝を $0,\infty$ に取ると自己同型は $z\mapsto a z$ または $a/z$。
後者は枝を交換し、前者の二接線倍率は $a,a^{-1}$ であるから恒等微分を持つのは恒等だけである。
射影変換に延びるものは、節点の表示 $[z:z^2:1+z^3]$ の三次切断空間を保つので
$a^3=1$ の高々六つである。尖点では唯一の滑らかな変曲点を $u=0$ に取って
正規化の切断空間を $\langle1,u,u^3\rangle$ と書ける。自己同型は $u\mapsto a u+b$。
これが切断空間を保つ条件は $b=0$ である。従って $\sigma$ は $b\ne0$ を持ち、
これと可換する射影倍率 $u\mapsto d u$ は $d=1$ だけである。
$s=1/u$ として尖点の極大イデアルをその平方で割ると基底は $s^2,s^3$ である。
引戻しで $s^2\mapsto a^{-2}s^2-2ba^{-3}s^3$, $s^3\mapsto a^{-3}s^3$ となり、
恒等微分を持つには $a=1,b=0$ が必要である。両曲線で恒等移動は偏極を保つため除外され、
$D(e)=0$ から矛盾を得た。
\end{proof}

% !TeX root = ../papers/nc-p3.tex
\section{滑らかな二次曲面のグラフが続く窓の全分類}
\label{sec:smooth-quadric-graphs}
四辺の各射影の非空なファイバーがすべて零次元であり、少なくとも一辺の二文字点対応が
スキームとして滑らかな二次曲面間の同型のグラフであるとする。
他の辺の射影像の一致や全位置を通る点列は仮定しない。
\begin{theorem}\label{thm:smooth-quadric-graphs}
この枝の標準14641型の添字保持同型類は $\mathbb A^1\setminus\{2\}$ である。
四文字 $x,y,z,w$ を、二平面 $\langle x,y\rangle,\langle z,w\rangle$ に分ける。
異なる平面の二文字は交換し、ほかの二関係は位置 $i$ で
\[
 xy+yx+zw+wz=0,\qquad [z,w]=(-1)^i m[x,y].
\]
パラメータは $m\ne0,\ m^2\ne1$ で、上のアフィン座標は $m^2+m^{-2}$ である。
全点の最小周期は1である。これらの窓は正則正規一次直線列を持たず、
どの量子多項式点も軌道閉包に持たない。
\end{theorem}
\begin{proof}[二次核の正規性]
まず全辺のグラフがつながることを示す。各射影のファイバーは線型なので、
零次元条件は左右評価の全点での階数3以上に等しい。辺上の $(p,r)$ では
命題\ref{prop:three-finite-edges}の一般不等式 $b_i(p)+f_{i+1}(r)\le6$ により両方3である。
系\ref{cor:finite-edge-point-sequences}の局所証明で必要なのはこの階数だけである。
すなわち基点で単元な三次小行列を分裂させると、次辺評価の四次小行列がスキーム上で消え、
その核直線への射が定まる。逆向きの核直線も一意であるから共有文字への両像はスキームとして一致する。
これを反復すれば四辺とも滑らかな二次曲面のグラフである。
二次曲面の同型は二直線族を保つか交換するため、超平面束 $\mathcal O(1,1)$ を保つ。
従って点対応に沿って全整数位置の基底を運び、一つの $Q\subset\mathbb P(V^*)$ の
対角グラフにできる。$q\in\operatorname{Sym}^2 V$ を非退化なその方程式の対称持上げとする。
グラフを消す二次式の空間は $K=\bigwedge^2 V+kq$ である。各六関係はこの中の超平面で、
一次元像の生成元 $s_i$ を選ぶと
\[
 [p,r]_i=\Omega_i(p,r)s_i,\qquad q_i=c_i s_i
\]
と書ける。$\Omega_i=0$ なら辺は平面三空間全体の対角である。
階数2なら $\Omega_i$ を $V^*$ の可分解二ベクトルと見て、その二平面の射影直線と
$Q$ の交点 $[p]$ を取る。$(p,p)$ の接条件の交代部分は
$p\wedge(\dot p_2-\dot p_1)\in k\Omega_i$ で、差には射影接方向が一つ残る。
二つの変位を合わせた次元4から残る $q$ の一式を引いても接次元は3以上であり、
滑らかな対角グラフに反する。よって全ての $\Omega_i$ は非退化である。
逆に非退化なら $p\wedge r=t\Omega_i$ に $p$ を外積すると $p\wedge\Omega_i\ne0$ から
局所的に $t=0$ となり、残る条件は対角と $q(p,p)=0$ そのものである。

全 $s_i$ による商は可換二次曲面環 $D=k[V]/(q)$ の反復である。
$\ell_\Omega:\bigwedge^3 V\to V$ を
$\ell_\Omega(u,v,h)=\Omega(u,v)h+\Omega(v,h)u+\Omega(h,u)v$ とする。
三文字の交代和から
$s_i\ell_{\Omega_i}(u,v,h)=\ell_{\Omega_{i+1}}(u,v,h)s_{i+1}$ である。
両 $\ell$ は可逆なので $s_i V=V s_{i+1}$ である。これを全経路に繰り返せば
商への核はそれぞれ片側の $s$ の倍で尽くされる。標準Hilbert関数と
$\dim D_n=(n+1)^2$ の差は $\dim A_{i+2,i+n}$ に等しいから、左右の $s$ 乗法は全て単射である。
従って $s_i v=g_i(v)s_{i+1}$ は二移動の可逆な乗法移送を与え、さらに $s$ の
イデアルを保つ（$s_i s_{i+2}$ の式で両端の正則性を使う）。
商 $D$ の点対応は $Q$ の対角だから $g_i$ は全位置で射影的に同じで、
その射影変換は $Q$ を保つ。$s_i$ の倍率を順に直して $g_i=g$ と取る。
\end{proof}

\begin{proof}[二直交平面の復元]
$q$ の対称行列を $J:V^*\to V$ とし、$gJg^t=\mu J$ と書く。
体積を用いて $\ell_{\Omega_i}$ を可逆交代行列 $M_i:V^*\to V$ と見る。
共通の符号規約の下で $M_i$ は $\operatorname{Pf}(\Omega_i)\Omega_i^{-1}$ の
一定の非零倍である。三文字の交代恒等式から $M_{i+1}=g M_i$、
ゆえに $g^t\Omega_i=\Omega_i g$ となる。$s=\operatorname{Pf}(M_{i+1})/\operatorname{Pf}(M_i)$ とすると
$\Omega_{i+1}=s\Omega_i g^{-1}$。各一般化固有空間は交代形式で互いに直交し非退化なので
各固有値の重複度は偶数である。この比 $s$ は各固有値をその半重複度だけ掛けた値である。
実際 $\operatorname{Pf}((T I-g)M_i)/\operatorname{Pf}(M_i)$ は
$\det(T I-g)$ のモニック平方根で、その定数項がこの比である。
従って $s$ は位置によらない。

$qv-vq$ を二つの二文字交換子に展開すると
\[
 g(c_i I-J\Omega_i)=c_{i+1} I+J\Omega_{i+1}.
\]
ここで $\Omega_i$ の行列要素は $\Omega_i(e_j,v)$ の規約で取る。
相似式と自己随伴性から $gJ\Omega_i=(\mu/s)J\Omega_{i+1}$ なので
\[
 c_i g-c_{i+1} I=(1+\mu/s)J\Omega_{i+1}.
\]
$g$ がスカラーなら左辺はスカラー、右辺は非零な反対称形式に由来する非スカラーであり不可能である
（このとき $\mu=s$）。$\mu\ne -s$ なら $J\Omega_{i+1}$ は $g$ と可換する。
$J\Omega_i=s^{-1}J\Omega_{i+1}g$ も可換し、$gJ\Omega_i=\mu J\Omega_i g^{-1}$ から
$g^2=\mu I$ である。非スカラーなので二固有値は $b,-b$ で重複度は2ずつ、
$s=-b^2=-\mu$ となって矛盾する。よって $\mu=-s$ で全 $c_i=0$ である。

固有値が $b$ 一つなら、対称形式の相似性から $\mu=b^2=s$ が必要であり不可能である。
二固有値 $b,c$ の場合 $s=bc$ であり、$\mu=-bc$ では両一般化固有空間は
対称形式でも互いに直交する。それぞれで非退化なので $b^2=c^2=\mu$、従って $c=-b$。
二次元の非退化交代平面上の自己随伴写像はスカラーであり、Jordan部分はない。
$g$ の全体倍率を直して二固有値を $1,-1$ と取る。
両平面で $q$ を双曲型に取れば $q=xy+yx+zw+wz$、二平面間の交換子は零である。
$\Omega_{i+1}=s\Omega_i g^{-1}$ の二平面上の倍率の比が $-1$ なので、
超平面の式は定理の表示で尽くされる。
\end{proof}

\begin{proof}[双対Clifford環と除外される二値]
$D^!$ を通常次数で書いた環を $C$ とする。双対文字を $X,Y,Z,W$ とすると
四文字は平方零、二平面間は反交換、$XY+YX=ZW+WZ=h$ は中心である。
$C$ は $k[h]$ 上で十六の平方自由順序語を基底とする。実際各二文字対を
$YX=h-XY$ 等で並べ、平方と混合する重なり $YXX,YYX$ は零、二平面は符号付きで交換する。
これで $h$ を独立変数とした表示の基底が得られ、$h$ を消すと元の二次表示になる。
$U=XY-YX$, $T=ZW-WZ$ と置くと $U^2=T^2=h^2$, $UT=TU$ である。
$A^!$ は各位置で $C$ の一関係 $a_i=U+(-1)^i mT$ を加えたものとなる。
実際この関係の $K$ 上の汎関数は $q$ で零、二交換子の上で比 $(-1)^i m$ である。
$U$ は $X,Y$ と反交換し $Z,W$ と交換、$T$ はその逆なので $a_i$ は次位置の
$a_{i+1}$ へ符号付きで移り、$a_i a_{i+1}=(1-m^2)h^2$ である。

$m^2=1$ なら $a_{i+1}h^n\ne0$ が全て $a_i$ で消える。
$C_j$ は $j\ge3$ で常に次元8なので、正規性により商の高次偶数片は無限に残る。
標準型ではない。$m^2\ne1$ なら $a_i$ は左右正則で、商の各始点級数は $(1+z)^4$ である。
可換多項式環の正則二次式商 $D$ はKoszulであり、その双対 $C$ もKoszulである。
$C$ のこの正則正規二次式商もKoszulである。実際、変換環Tor系列の二段は
商のTorの下段と内部次数2を移した上段である。最初の非線型下段成分の出る微分は
二列前の線型な上段に入るので非線型次数を消せず、元のKoszul性に反する。
位置別の符号移送は頂点や矢の正則な捩りだけを加え、次数比較は同じである。
同じ議論を通常多項式環から $D$ にも適用できる。

$F=C/(h^2)$ は最高次数6のFrobenius環である。socleは $h$ に消されるので $hF$ に入り、
$C/(h)$ が四文字の外積環であることから左も右も $hXYZW$ の一本だけである。
位置 $i$ での商 $C/(a_i)$ の次数片を、$a_{i+1}$ の左乗法で $F$ の次数を2上げた部分へ入れる。
これは $a_i a_{i+1}=(1-m^2)h^2$ と単射性から核が正確に $a_i C$ である。
$F$ の積と最高traceを使えば、商の各非零な次数 $j$ の元に、右から掛けて非零となる
次数 $4-j$ の元がある。$a$ の正規性によりこの積は相手側でも商を通る。
相補次元が等しいので全積対合が完全である。よって $A^!$ はKoszulかつ標準Frobenius、
元の窓は全て標準Koszul左右AS正則である。
\end{proof}

\begin{proof}[同型と二つの補集合への所属]
点グラフから正規化後の全位置の同型は射影的に同じで $Q$ を保つ。
復元した $g$ の二固有平面を保つか交換する。保つものは二平面で同じ相似倍率を持ち、
交代面積に掛かる二つの行列式の比は $\pm1$ である。両符号が実現し、平面交換は比を逆にする。
ゆえに同型条件は $m'\in\{m,-m,m^{-1},-m^{-1}\}$ である。
この作用の不変環は $k[m^2+m^{-2}]$、除外する値は2である。一位置移動は $m\mapsto-m$ で同型である。
正規一次列があれば $D$ 上で連続する二つの超平面切断の零因子が等しい。
$Q$ は整で張るので全位置で同じ文字直線 $kv$ になる。$v V=V v$ の対応は、
可換整域 $D$ へ移すと各文字を固定するから、元の $A$ でも $[v,V]=0$ を要する。
$\Omega_i$ の非退化性に反する。

最後にこの窓の射影軌道はpolystableである。上で得た最高次四テンソルは、各語について
$F$ の $hXYZW$ 係数を $a_{i+1}$ とその四文字の積から取って表示できる。
$(X,Y,Z,W)$ を $(tX,t^{-1}Y,uZ,u^{-1}W)$ に動かすトーラスで不変なので、
異なる一文字切片の台は交わらない。語の台は $(XY)(ZW)$ の同時入替えで保たれる。
また一位置巡回と $Z,W$ 入替えの合成でも保たれる。実際、$F$ の最高traceに関する
一文字の巡回は符号 $-1$ を掛ける。平方自由基底で、相手の次数5が $h$ 倍の三文字であることから
外積の符号となる。さらに $a_{i+1}$ を一文字越すと符号を除き $a_i$ に変わり、
$Z,W$ 入替えがこれを元に戻す。この有限置換群で各語を平均すると、
各位置で第一平面の文字の重みは二文字とも同じで、他平面も同様、各平面の総重みは位置によらない。
台には $XYXY$ と $ZWZW$ がともにある（係数はそれぞれ $m$ および1の非零倍）。
従って全語に正の重みを与えた平均にこの二軌道の重みを調節し、全周辺重みを均等にできる。
相対内部の重心による対角ノルム最小化と切片の直交性
（定理\ref{thm:quantum-face-resonance}の証明）から全四基底群でpolystableとなる。
量子点の射影支持は座標線型空間の和であり、この滑らかな二次曲面とは異なる。
異なるpolystable軌道は互いの閉包に入らない。
\end{proof}

% !TeX root = ../papers/nc-p3.tex
\section{滑らかな二次曲面の各閉軌道への全退化}
\label{sec:quadric-attraction-fibers}
\begin{theorem}\label{thm:quadric-attraction-fibers}
前節の二次曲面窓を $W(m)$ とする。その軌道閉包への逆像、すなわち
$W(m)$ を軌道閉包に含む全窓の添字保持同型類は、
$m^2\ne-1$ では五つ、$m^2=-1$ では三つである。全て標準Koszul左右AS正則である。
次の表示は背景の点対応を保つという仮定を置かない。
$b$ を $X,Y,Z,W$ の一文字とし、他の平面の二文字を固定順に $c,d$ として
\[
 E_b=bcbd-bdbc-cbdb+dbcb
\]
と書く。四文字語はここでは四テンソルである。$W(m)$ を前節のClifford最高traceで規格化する。
全代表は $W(m)+A E_b+B E_{b'}$、$b\in\{X,Y\}$, $b'\in\{Z,W\}$、
$A,B\in\{0,1\}$ から得られる。残る同型は非零欄の台に対する両平面の同時反転、
および $m^2=-1$ の場合の四巡回である。
一般の場合、空台と二つの単独欄の同型類の最小周期は1、二つの両欄の類は2である。
例外 $m^2=-1$ では全三類が最小周期1である。同じファイバー内の全退化は、
これらの置換の後の台の包含である。
\end{theorem}
\begin{proof}[全パラメータでの横断接]
各文字に $(1,0),(-1,0),(0,1),(0,-1)$ を与える。
前節の軌道安定化群の連結部分はこの二次元トーラスである。実際、二つの非退化二平面を
保つ共通相似倍率の直交変換は、射影倍率を除けばこのトーラスと有限個の入替えとなる。
二平面を交換するものは $m^2=-1$ の場合にだけあり、片方の文字順を逆にするので
四文字上の四巡回となる。基点と各 $E_b$ はこの有限群とトーラスで張る部分空間内に移る。

\path{computations/verify_quadric_attraction.py} に接次元上界と持上げの有理式検算を置く。
ここで有限体での階数を全閉点へ外挿しないため、計算する小行列を記す。
各切断で $4i+j$ 順の六行六列 $P=(1,2,3,4,6,7)$ の行列式は
$-4m^2(m-1)^2(m+1)^2$ であり、許容パラメータ全域で単元である。
各切断の背景行列を $M$, $I=P^c$, $J=M_{PP}^{-1}$ とし、
$L$ を $I$ 行選択から $M_{IP}J$ 倍の $P$ 行選択を引いた行列、
$R$ を $I$ 列選択から $P$ 列選択の $JM_{PI}$ 倍を引いた行列とする。
テンソル変位の切断を $Y$ と書くと一次式は $LYR$ である。

二切断の一次式を語の指標で分け、さらに四位置に全行列単位を掛けた軌道接を引くと、
$(\pm2,0),(0,\pm2)$ には各高々一次元、零指標には高々二次元、
残りの全非零指標には何もない。具体的には各ブロックで $m=2$ での簡約階段形から
独立な列を選び、その列の転置から独立な行を選ぶ。一次式と軌道接の双方について、
選んだ小行列式を $\mathbb Q(m)$ で因数分解すると全て $m,m-1,m+1$ だけの
整数冪と非零定数の積である。従って上界は全パラメータに有効である。
四つの $E_b$ は各々一次式を消し、対応する軌道接に一列として加えた単元小行列により
横断四方向を実現する。
\end{proof}

\begin{proof}[四平面と全軌道]
余指標で基点へ流れるには四指標の各対から高々一方を選べる。
この二方向の非零台はそれぞれの文字を正にする $(\pm1,\pm1)$ で実現できる。
これらの表で変位の余重みは2、全方程式の最大余重みは4である。
実際、上の主行列式は基点で非零なので指標零、縁付き行列式の指標は余分な行と列の
四文字の指標である。重み4のSchur補正は $-LY_{*,P}JY_{P,*}R$ であり、
$Y$ を $E_b$ 一つ、および異なる平面から一つずつ選んだ二つの和にすると零になる。
前掲コードは二つの純平方項と偏極した混合項の全成分を記号的に消去する。
従って四つの表示平面は追加の非線型項を要さず、両隣接階数は正確に6である。

安定化群不変な線型横断面にこれら四方向を含めてLunaの切片を作る。
定理\ref{thm:quantum-face-resonance}の証明と同様に、完備局所環は四つの線型射影と
零指標の生成元で位相的に生成される。切片で基点へ流れる任意の点は、同じ余指標の弧上で
零指標が背景値、四射影が表示平面の点と同じ単項式であるから、その平面の点と一致する。
飽和開での強エタール性は閉軌道への全ファイバーを覆う。
各平面からの一径数退化に命題\ref{prop:unrestricted-good-reduction}を適用すれば全次数で正則である。

二平面から一つずつ選んだ指標は独立なので、代数閉体上では非零係数を全て1にできる。
$m^2\ne-1$ の場合には同時反転だけで割る。台の数は空台が1、単独が2、両方が2となる。
$m^2=-1$ の四巡回は単独四台および両方四台のそれぞれで推移的である。
一位置巡回で背景を前の表示に戻す変換は一平面内の入替えである（関係の $m$ が $-m$ に変わる）。
従って一つの指標成分だけを逆転する。これは一般の場合の両方の台だけを交換する。
形式的にもこの作用を確かめるには巡回を加えた有限拡大で切片と接射影を選べばよい。
四線の非零欄に掛かる倍率はトーラスで吸収できる。一般の場合の両台は二移動で戻り、
例外では一移動ですべて元の類に戻る。台の包含による退化は、各表示平面の二つの
独立な指標のトーラス軌道閉包から従う。
\end{proof}

% !TeX root = ../papers/nc-p3.tex
\section{滑らかな二次曲面商を持つ全窓と二次元射影像の次数}
\label{sec:smooth-quadric-quotients}
滑らかな二次曲面のグラフが点対応の一成分でよく、余分な点、直線、非被約構造を許す。
不変な仮定は、全整数位置に、滑らかな二次曲面の斉次座標環 $D$ の反復への
一次で同型な全射が存在することである。グラフの移動は超平面束を保つので
互換なグラフを含む場合は位置別基底を運べばこの形となる。
\begin{theorem}\label{thm:smooth-quadric-quotients}
この商を持つ標準14641型の全添字保持同型類は以下で尽き、全て最小周期1である。
\begin{enumerate}
\item 二つの可逆二次行列 $P,Q$ が同じ特性多項式を持つ場合。各々の共役、
両者の交換と共通スカラーだけを除く。相異なる二根なら根の比
$u\in k^\times\setminus\{1\}$ を $u\sim u^{-1}$ で割ったもの。
重根なら両方スカラー、片方だけJordan型、両方Jordan型の三点である。
\item $P=\operatorname{diag}(1,-1)$, $Q=\operatorname{diag}(r,-r)$,
$r\in k^\times$ とし、次位置で $r$ を $r^{-1}$ に替えて交互に繰り返す場合。
$r\sim-r\sim r^{-1}$ だけを除く。この枝の $r^2=1$ は第一の $u=-1$ と同じ一点であり、
それ以外では前節の二次曲面だけのグラフとなる。
\end{enumerate}
表示からの復元では、下記の七次元空間中の各 $(P_i,Q_i)$ を二次双対の追加関係に取る。
両方スカラーの場合は可換多項式窓であり、それ以外では二次曲面は点対応の唯一の二次元既約成分として復元できる。
この商を持つ有限Frobenius候補には非Koszul例はない。
\end{theorem}
\begin{proof}[四次対合による可逆性の強制]
$D^!$ の通常のClifford表示 $C$ と中心 $h$ は前節のものを使う。
$h=1$ での評価は各斉次片で単射である。実際、$k[h]$ 上の十六基底では
各語の係数は同じ次数の $h$ の冪の定数倍しか取り得ない。
$V^*$ の二次形式を $2\times2$ 行列の行列式とする。一次片のブロック表示は
\[
 v(U)=\begin{pmatrix}0&U\\ U^\#&0\end{pmatrix},\qquad
 U^\#=\operatorname{adj}(U),\qquad U\in M_2(k).
\]
この行列の平方は $\det(U)$ 倍の単位なのでClifford関係を満たす。
二つの行列単位の積で偶ブロックを取り、さらに一次を掛ければ全行列環が生成される。
具体的に $v(U)v(V)$ の対角ブロックは $(UV^\#,U^\#V)$ で跡が等しく、その線型包は
\[
 H=\{(P,Q)\in M_2(k)\oplus M_2(k):\operatorname{tr}P=\operatorname{tr}Q\}.
\]
行列単位で確認するか、$C_2,C_3,C_4$ の次元 $7,8,8$ と四次での体積を使ってもよい。
$C_3$ は非対角の全二ブロック、$C_4$ は対角の全二ブロックとなる。
$K$ を曲面のグラフを消す七つの関係の空間とすれば、窓の二次双対は
$C$ の反復で各位置の一つの非零 $a_i=(P_i,Q_i)\in H$ を消したものである。

四次の非零traceを $C_4$ に引き戻し、
$\operatorname{str}(b_i\,\cdot)$ と表す。ここで $b_i$ は二ブロック行列、
$\operatorname{str}(M,N)=\operatorname{tr}M-\operatorname{tr}N$ である。
$H^\perp=kI$ なので、$H\times H$ 上の四次積対合の左核は
$\{a\in H:b_i a\in kI\}$ である。標準Frobenius性からこの核は正確に $ka_i$、
右核は $ka_{i+2}$ である。$b_i$ が可逆なら $a_i$ も可逆である。
$b_i$ が特異なら $b_i a\in kI$ は $b_i a=0$ に限る。二ブロックの階数の和が
2以下なら全偶空間内の核は4次元以上で、超平面 $H$ 内でも3次元以上となり得ない。
従って階数は1と2である。その核内の $a_i$ は一ブロックだけに台を持ち、
その像は直線、跡は零の非零行列である。特に、どの $b_i$ も両ブロックが非零である。

しかしそのような一ブロック $a_i$ は前の四次traceと矛盾する。
$C_1a_i C_1$ は、$U,V$ が任意の二行二列行列であることから、
反対側の一ブロック全体を張る。任意の非零行列 $P$ で $M_2 P M_2=M_2$ だからである。
前のtraceはこの全空間を消すので $b_{i-1}$ の一ブロックが零になってしまう。
従って全 $b_i$ と全 $a_i$ は可逆である。この議論では各点への評価階数を仮定していない。
\end{proof}

\begin{proof}[二ブロックの移送と全特殊値]
三次関係イデアルの次元は $8-4=4$ であり、可逆性により $a_i C_1$ と $C_1 a_{i+1}$ は
ともに四次元で等しい。従って任意の $U$ で $a_i v(U)a_{i+1}^{-1}$ の
下ブロックが上ブロックの余因子となる。式は
\[
 Q_i U^\# P_{i+1}^{-1}
 =(Q_{i+1}^{-1})^\# U^\# P_i^\# .
\]
全 $U^\#$ に対する左右乗法の一致から $[P_{i+1}]=[P_i]$,
$[Q_{i+1}]=[Q_i]$ となる。$a_i$ のスカラーを除いて
$a_i=(P,\lambda_i Q)$ と書くと
\[
 \lambda_i\lambda_{i+1}=\det P/\det Q,\qquad
 \operatorname{tr}P=\lambda_i\operatorname{tr}Q.
\]
跡が非零なら倍率は一定で、下ブロックを取り直せば二行列の跡も行列式も等しい。
ともに跡零なら可逆だから異なる二根の対角形に取れ、表示の $r,r^{-1}$ が得られる。
初めから一定で跡零の場合はちょうど $r^2=1$ である。

逆にこのリストの各解では、移送式が成立し $a_i$ は各斉次片で正則な正規関係となる。
可逆性は $h=1$ の行列と各斉次片の単射性から読める。さらに両ブロックの逆行列は跡が等しい。
跡零の場合は両方零、それ以外の場合は行列式も等しいからである。
逆対を $C_2$ の元 $b$ へ持ち上げると $a_i b=b a_i=h^2$ である。
従って前節で用いた $C/(h^2)$ のFrobenius対合への埋込みが同様に働き、
商双対は級数 $(1+z)^4$ のKoszulかつFrobeniusである。
同じ次数の正規二次関係によるTor二段の議論はそのまま使える。これで全次数の正則性を得る。

任意の非スカラー表示では $\Omega_i\ne0$ である。階数4なら辺はちょうど曲面のグラフ。
階数2なら $p\wedge y=t\Omega_i$ の $t\ne0$ の解は、$\Omega_i$ の台である同じ直線に
両点が乗る場合に限る。そこでは $q(p,y)=t c_i$ の双一次式が残る。
この式が直線の積全体で消えるには $c_i=0$ かつ台が全等方でなければならない。
この場合 $a_i$ は直交する二つの等方文字のClifford積の非零倍となり平方零なので、可逆性と矛盾する。
従って余分な部分の次元は高々1で、最大の二次元成分から $Q$ とそのグラフが一意に戻る。
第二列の $r^2\ne1$ が曲面だけであることは前節の $U+mT$ での両ブロックの固有値が
それぞれ $\pm(1+m),\pm(1-m)$ であることによる。
同型は各位置で同じ曲面の射影変換である。
連結部分は $U\mapsto G U H^{-1}$（$G,H\in\mathrm{SL}_2$）で両ブロックに独立な共役を誘導し、
$U\mapsto U^\#$ がブロックを交換する。従ってリストにある同一視でちょうど尽きる。
共通スカラーだけが正規化されていないことに注意する。両方スカラーのときは追加関係が $h$ で
元の窓は可換多項式表示である。一移動は一定列なら恒等、跡零交互列ならブロック交換まで恒等である。
\end{proof}

\begin{proposition}\label{prop:surface-degree-bound}
四辺の各射影のファイバーが空または一点なら、前節と同様に全辺の共有射影像は
スキームとして一致し、点対応は相互の同型のグラフとなる。
一辺の第一射影像が平面三空間全体でなければ、その二次元部分の超曲面因子の総次数は
重複込みで高々2である。特に滑らかな二次曲面が一つの既約成分として現れる標準窓は、
余分な埋込み成分を許しても上のリストで尽くされる。
\end{proposition}
\begin{proof}
共有像の一致には双対の標準縮約を用いる。四次元文字に作用する六関係の評価行列を
$M(p)$ ($6\times4$) とし、一般階数4とする。その四次小行列は四次式である。
共通四次因子があれば四平面 $\operatorname{im}M(p)$ のPlücker点が一定となる。
よって六関係に定数の線型依存があり矛盾する。
共通三次因子があれば、それを除いたPlücker座標は一次式で、像の閉包は
$\operatorname{Gr}(4,6)$ 内の射影線型空間となる。その四平面全てが固定五平面内か、
全てが固定三平面を含む。実際二点を結ぶ直線がGrassmann多様体に入るには平面対が
三次元で交わる必要があり、二平面の共通三平面を含まない第三平面はその和五平面に入る。
両型を同時に外へ広げると三次元交叉条件に反する。固定五平面の場合は再び六関係が従属である。

固定三平面で割ると $3\times4$ の一次式行列の全評価が階数1以下である。
階数一行列の線型空間では像直線が共通か行直線が共通である。前の場合は既に除外した。
後の場合、三つの独立な関係が $\ell_j(p)v(y)$ と同じ右因子を持つ。
平面 $v(y)=0$ 上、残る三関係が左文字に作用する $3\times4$ 行列は
各点で階数3を持たなければならない。そうであれば
$\mathcal O^4\twoheadrightarrow\mathcal O(1)^3$ の核は $\mathcal O(-3)$ である。
平面上で全Chern類を掛けると $(1-3H)(1+H)^3=1-6H^2\ne1$ で矛盾する。
従って三次因子もない。
滑らかな二次曲面の成分を次の位置に送った像は、同型な二次元の非平面的整曲面で、
上の次数上界から再び滑らかな二次曲面である。これを繰り返せば全位置で商が得られる。
\end{proof}

% !TeX root = ../papers/nc-p3.tex
\section{整二次錐の全商と平面成分からの正規一次直線}
\label{sec:quadric-cone-quotients}
$D$ を階数3の二次式による四変数多項式環の商とする。
\begin{theorem}\label{thm:quadric-cone-quotients}
$\check D$ への一次で同型な全射を持つ標準14641型には、次の明示的な全同型類がある。
$\epsilon^2=0$ とし、三つの欄は二次双対への追加関係 $a_i$ を表す。
\begin{enumerate}
\item $a_i=I+\epsilon N$。$N$ は零、非零冪零Jordan行列、$\operatorname{diag}(1,-1)$ の三種。
\item $a_i=P$ は非スカラー可逆二次行列。共通倍率と共役を除き、
異なる根の比 $u\in k^\times\setminus\{1\}$ ($u\sim u^{-1}$) または非自明Jordan行列を取る。
\item $P=\operatorname{diag}(1,-1)$, $a_i=(1+(-1)^i\epsilon)P$ の一種。
\end{enumerate}
全て最小周期1であり、これら三欄の重複はない。第一欄の $N=0$ は可換多項式窓である。
最後の一種だけが余分な部分を持たず二次錐そのものを点対応とする。
整二次錐商を持つ有限標準Frobenius候補では、このリストに入ることもKoszul性も自動である。
\end{theorem}
\begin{proof}[退化Clifford環]
$T=k[\epsilon]/(\epsilon^2)$ とし $M_2(T)$ に $j^2=1$, $jBj=\bar B$
（$\epsilon$ の符号を変える）で奇元 $j$ を加える。
一次文字の像を $(U+\epsilon uI)j$, $U\in\mathfrak{sl}_2(k),u\in k$ とする。
平方は $U^2=-\det(U)I$ なので、これは階数3の形式のClifford表示である。
中心二次式 $h$ を独立にした前節の $k[h]$ 自由基底は退化した形式でも有効である。
例えば対角化して平方を $c_i h$、異なる二文字を反交換させれば順序語が基底となる。
$h=1$ では二文字積は $UV+\epsilon(uV-vU)$ となり、その線型包は
$H=\{P+\epsilon Q:\operatorname{tr} Q=0\}$ の全七次元である。
三文字は $M_2(T)j$、四文字は $M_2(T)$ の全空間を張る。これは無跡行列の積が
$M_2(k)$ を張ることと可逆な無跡文字を用いて分かる。
従って各斉次片のこの表示は単射である（次数4以後は $h$ による移送を使う）。

$M_2(T)$ のFrobenius汎関数を $\operatorname{tr}$ の $\epsilon$ 係数に取ると $H^\perp=kI$ となる。
四次traceに対応する行列を $b_i$ とする。前節と同じ二次対合から
$\{a\in H:b_i a\in kI\}=ka_i$ である。$b_i$ が非可逆ならその全行列空間での左乗法核の
$k$ 次元は2以下である。$T$ 上の行列は単元行列の左右乗法でSmith形にでき、
核の次元は $T^2$ 上の核の次元の二倍である。従って非可逆な場合のSmith形は
$\operatorname{diag}(1,\epsilon)$ だけである。その核にある $a_i$ は
$\epsilon N$ で $N\ne0$ は跡零かつ階数1、特に各 $b_i$ の剰余は非零である。
しかし $C_1a_i C_1$ は $\epsilon M_2(k)$ の全体を張る。実際、$N$ を行列単位 $E_{12}$ に
共役すれば $U E_{12} V$ で無跡な $U,V$ の第一列と第二行を任意に取れる。
前位置のtraceがこの空間を消すなら $b_{i-1}$ の剰余は零となり矛盾する。
従って全 $b_i,a_i$ が可逆である。
\end{proof}

\begin{proof}[移送、同型、正則性]
三次核次元4より再び $a_i C_1=C_1 a_{i+1}$ である。
$a_i=P_i+\epsilon Q_i$ と書く。全無跡 $U$ で $P_i U P_{i+1}^{-1}$ が無跡なので
$P_i$ は射影的に同じであり、$P_i=P$ とする。$D_i=Q_iP^{-1}$ と置くと
$a_i(U+\epsilon uI)\bar a_{i+1}^{-1}$ の $\epsilon$ 係数の条件は、
全無跡 $S=P U P^{-1}$ で $D_i S+S D_{i+1}$ がスカラーであることに等しい。
$D_i=d_i I+N_i$ と跡と無跡に分ける。無跡二行列の反交換子はスカラーなので、
無跡部分で $(d_i+d_{i+1})S+[N_i-N_{i+1},S]/2=0$ である。
無跡空間上の写像の跡から $d_{i+1}=-d_i$、中心化条件から $N_{i+1}=N_i$ となる。
よって $D_i=N+(-1)^i dI$, $\operatorname{tr}(NP)=0$, $d\operatorname{tr}P=0$ で尽くされる。

二次錐の射影自己同型は、全体倍率を除いて $(U,u)\mapsto(GUG^{-1},c u+f(U))$、
$c\ne0$, $G\in\mathrm{PGL}_2$ である。これは三次元の非退化二次形式の射影直交群と
一次元の根基を使えば従う。切断グラフの同型が抽象的に与えられても超平面束は保たれる。
実際、二次錐の極小解消 $\mathbb F_2$ で例外切断を $E$、母線を $F$ とすると、
引戻し超平面類は $E+2F$ であり、自己同型は唯一の負切断と母線を保つ。
$u$ の平行移動 $f(U)=\operatorname{tr}(K U)$ ($\operatorname{tr}K=0$) は
行列環の $I+\epsilon K$ による共役で実現する。奇文字では
$\{K,U\}=\operatorname{tr}(KU)I$、偶文字では $Q_i\mapsto Q_i+[K,P]$ となる。
従って $N$ は $(1-\operatorname{Ad}P)\mathfrak{sl}_2$ の加算まででよい。
非スカラー $P$ ではこの像は二次元で、$\operatorname{tr}(NP)=0$ の平面に等しい。
$P$ がスカラーなら $d=0$、残る $N$ は共役と非零倍率で割る。
これでリストの全同型を得る。交互な $d$ の符号は $c=-1$ で逆転できる。

逆行列の $\epsilon$ 係数にも跡がない。実際 $\operatorname{tr}(D_iP^{-1})=
(\operatorname{tr}P\operatorname{tr}D_i-\operatorname{tr}(D_iP))/\det P=0$。
従って各逆行列は二次の補元に持ち上がり、元との積は $h^2$ となる。
前節の正規二次式商のKoszul性と $C/(h^2)$ でのFrobenius対合を同様に使える。
$h$ を消すと四文字の外積環である点も変わらない。

可換表示以外では、前節と同じ $p\wedge y$ の議論で唯一の二次元既約成分が
二次錐である。交代部分が階数2でその台の直線が曲面に含まれなければ、
その直線の一般点に別の点対応が出る。含まれれば錐の滑らかな点で、前節の接計算から
グラフより大きな接空間が出る。交代部分が階数4のときだけグラフがそのまま得られる。
Clifford表示では交代二形式を、三つの無跡文字間の交換子と、根基文字とそれらとの交換子に分ける。
前者は $P$ の無跡部分 $P_0$、後者は $Q_i$ とそれぞれ無跡行列空間で同一視される。
両者を外積して非零となる条件は $\operatorname{tr}(P_0 Q_i)\ne0$ である。
実際、三無跡文字の交換子は外積二乗と無跡空間の同型で、
その体積対合は $\operatorname{tr}([U,V]W)$ の非零倍である。
リストの最後に限りこの跡が非零である。
\end{proof}

\begin{theorem}\label{thm:plane-surface-normal}
標準窓の全辺の射影ファイバーが空か一点で、一辺の射影像に平面成分があるとする。
この平面から正則正規一次直線列が復元され、中央化して得る三生成商は
三変数多項式環の次数付きtwistである。従って零次元射影ファイバーで二次元成分を持つ全標準窓は、
滑らかな二次曲面商、整二次錐商、または正規一次直線列の平面商のいずれかに入る。
\end{theorem}
\begin{proof}
前節の共有像の同型と次数上界から、平面の次の像も平面である。平面と同型な
埋込み曲面は次数2の整二次曲面になれない。射影像が平面三空間全体なら
自己同型は線型で、六関係はそのグラフの全関係となり、可換多項式窓に帰着する。
以下、平面の列に沿って座標を運び、商を三変数の可換多項式環 $S$ の反復に取る。
一次核を $kt_i$ とする。

元の標準の単純加群の右線型解消に左から作用されるこの商をテンソルする。
始点を固定して得る複体は $S$ 上の長さ4の有限自由複体で、
平面上の各点では順に各微分の階数が $1,3,3,1$ である。
中央二写像は隣接関係を片方の点で縮約したものと相補関係の転置であり、
グラフ上で階数3、両端は点の非零評価である。従って射影平面上では完全で、
各ホモロジーは無関係イデアルに台を持つ有限長 $S$ 加群である。
$\operatorname{Ext}^p_S(-,S)$ は非零な有限長加群上 $p=3$ にだけ残る。
複体のhyper-Extの全次数 $p+q$ は4以下なので $q>1$ のホモロジーは零である。
第零は $k$、Euler級数は $(1-z)^4/(1-z)^3=1-z$ だから第一は $k(-1)$ である。
この計算は各単純に対する $\operatorname{Tor}^A(k_i,S)$ を内部次数ごとに与える。
各終点を固定した $S$ の左極小解消を見ると、核はただ一つの一次元 $t$ の右乗法による
自由加群であり、その乗法は単射である。反対側でも同じ計算を行うと、
商への核は両方からその一次元だけで生成され、左右に正則となる。
定理\ref{thm:regular-normal-linear}による一移動で中央化すると、商の点対応は
依然平面のグラフである。その移動は射影線型なので、三生成三関係の商は多項式環の次数付きtwistである。
最後の被覆は真の曲面因子の総次数が高々2であることから従う。
\end{proof}

% !TeX root = ../papers/nc-p3.tex
\section{平面商の全Jordan型と平面を忘れた同型}
\label{sec:plane-quotient-normal-forms}
正則正規一次直線 $kt$ を中央化し、三生成商が平面グラフとなる枝を扱う。
$H\in\mathrm{GL}_3$ を射影共役まで取り、巡回添字 $(i,j,k)$ に対し
$r_i=x_j H(x_k)-x_k H(x_j)$ を商の関係とする。
四生成関係は $t$ の中心性と $r_i=t\sum_j M_{ji}x_j+t^2 b_i$ である。
以下の表で与える全係数が標準Koszul左右AS正則、最小周期1である。
\begin{theorem}\label{thm:plane-quotient-all-marked}
一つの平面を指定した添字保持同型類は、以下の各型で尽きる。
\begin{enumerate}
\item $H$ が非スカラーで対角化でき、固有値を $h_1,h_2,h_3$ とする。
$h_i^3=\det H$ を満たす添字だけに $M_{ii}=a_i,b_i$ を許し、他の係数を零とする。
該当添字数は0、1、3である。一つなら二係数の非零台だけを記録する。
三つなら比が原始三乗根の固有値になり、関係左辺を同じ巡回交換係数で規格化する。
この場合は六係数の台とトーラス
$a_i\mapsto d_i a_i/(d_jd_k),\ b_i\mapsto b_i/(d_jd_k)$ の整数核単項式を記録し、
循環置換も除く。
\item $H=J_2(1)\oplus[\lambda]$, $\lambda\ne1$ のとき、
$\lambda\ne-1$ では補正は全て零である。$\lambda=-1$ では $M_{33},b_3$ の二係数の台だけを記録する。
\item $H=J_3(1)$ のとき
\[
 M=\begin{pmatrix}0&0&0\\2a/3&-a/3&0\\a&0&0\end{pmatrix},
 \qquad b=(b_0,0,a^2/9)^t .
\]
$(a,b_0)$ の代表は $(1,0),(0,1),(0,0)$ の三つである。
\item $H=J_2(1)\oplus[1]$ では
\[
 M=\begin{pmatrix}0&-a/2&0\\a&0&-2c\\b&c&d\end{pmatrix},\qquad
 (b_0,b_1,b_2)^t\quad\text{を定数欄とする。}
\]
標準代表を次に示す。最初の二行の定数は全て零であり、その後の行では $b_1=0$。
\[
\begin{array}{c|c|l}
 \text{条件}&(a,b,c,d)&(b_0,b_2)\\ \hline
 c\ne0 &(0,-1/2,1,0)&(0,0)\\
 c=0,\ a\ne0 &(1,0,0,0)&(0,0)\\
 a=c=0,\ d\ne0 &(0,0,0,1)&(0,0)\text{ または }\mathbb P(3,2)\\
 a=c=d=0,\ b\ne0 &(0,1,0,0)&(0,v)\ (v\in k),\quad(1,1)\\
 M=0 &(0,0,0,0)&(0,0),(1,0),(0,1)
\end{array}
\]
\item $H$ がスカラーなら、三次元Lie環 $\mathfrak g$ とそのスカラー二コサイクル類の表を使う。
アーベル、Heisenberg、$\mathfrak{aff}_1\oplus k$ ではそれぞれ零類と一つの非零類である。
二次元アーベル導来環に残る一文字が可逆行列 $P$ で作用する型は $[P]$ の共役類で分類し、
跡零のときに限り零類と非零類に分かれる。$\mathfrak{sl}_2$ は零類だけである。
\end{enumerate}
\end{theorem}
\begin{proof}[PBW条件の座標化]
$R_i$ を $r_i$ の二文字係数行列、$D=\det H$ と置く。
三次交叉は交代体積に $1\otimes H\otimes H^2$ を作用させたもので、
右端の文字は $H^2x_i$、左端の文字は $DH^{-1}x_i$ である。
従って定理\ref{thm:regular-normal-linear}の条件は
\[
 F=(M(H^2)^t-DH^{-1}M^t)H^{-t},\quad F+F^t=0,\quad
 C=(F_{23},F_{31},F_{12})^t,\qquad
 MC+(H^2-DH^{-1})b=0,\quad b^t C=0
\]
である。持上げ変更 $e$ は $M$ の第 $i$ 列から $(R_i+R_i^t)e$ を引き、
定数列に $M^te-(e^tR_i e)_i$ を加える。
$H$ と可換する基底変換 $G$ は持上げ変更の前に
$(M,b)\mapsto(GMG^t/\det G,Gb/\det G)$ と作用する。
これは $r_i$ の変更行列が $\det(G)G^{-t}$ であるためである。
以上は各Jordan行列についても分母の消滅を伴わない有限計算である。

非スカラー対角の場合、$F+F^t=0$ の非対角式は
$(h_j-h_k)M_{ij}+(h_i-h_k)M_{ji}=0$ である。一対ごとに二係数の差の一方は非零で、
残る一方向は対応する持上げ変更で消せる。対角式は $(h_i^3-D)M_{ii}=0$ となる。
このとき $C=0$ なので定数にも同じ共鳴条件が掛かる。
固有値が重複し共鳴が一文字だけのとき、その二係数は外の文字の倍率と
二文字ブロックの行列式で独立に変えられる。相異なる固有値の場合にも同じく二倍率が独立である。
全三文字が共鳴する非スカラーの場合は固有値が三乗根の一軌道になり、固定した対角表示を
射影的に保つ有限置換は三巡回だけである。

$J_2(1)\oplus[\lambda]$ では、中央化群は第一ブロックの上三角Toeplitz行列と
第三文字の倍率である。先の線型式を消去すると $\lambda\ne\pm1$ では全 $M$ を消せ、
$H^2-DH^{-1}$ も可逆である。$\lambda=-1$ ではただ一つの $M_{33}$ と $b_3$ が残り、
ブロック倍率の二乗と第三倍率が台上で推移的に働く。
$J_3(1)$ では持上げ変更後の式は表の通りで、$G=pI+q(H-I)+r(H-I)^2$ の作用で
$a\mapsto a/p$ となる。$a\ne0$ なら $b_0$ に対する $r$ の一次係数は
$-8a^2/(9p^3)$ で非零、$a=0$ なら $b_0\mapsto b_0/p^2$ である。
\end{proof}

\begin{proof}[二重Jordanの残余作用とLie型]
$H=J_2(1)\oplus[1]$ で $F$ の線型式と持上げ変更を使うと定理の $M$ となる。
定数についての式は
\[
 3ad+4bc+2c^2=0,\qquad b_1=a(b-c)/6,\qquad
 -3a b_2+2(b-c)b_1-6c b_0=0 .
\]
未使用の第二持上げ変更は定数に $(a,0,-2c)^t$ の任意の倍を加える。
$G$ は行を $(p,q,r),(0,p,0),(0,z,t)$ と書ける（$pt\ne0$）。
$M$ を再び表の形に戻すと
\begin{align*}
 c'&=c/p,& a'&=(ap-2cr)/(pt),& d'&=(dt-cz)/p^2,\\
 b'&=(3apz+4bpt-6crz+6drt)/(4p^2t).
\end{align*}
$c\ne0$ なら $c'=1$, $a'=d'=0$ とでき、方程式が $b'=-1/2$ を与える。
$a\ne0=c$ なら $d=0$ で $a'=1,b'=0$ とできる。いずれも残った定数は
未使用の変更で全て零になる。$a=c=0$ なら $d\ne0$ で $d'=1,b'=0$ とでき、
この表示の残余作用は $(b_0,b_2)\mapsto (b_0/p^3,b_2/p^2)$ である。
$d=0,b\ne0$ では $b=1$ に直した後は $(b_0,b_2)\mapsto
((b_0+r(b_2-1))/t,b_2)$。$M=0$ の場合は
$(b_0,b_2)\mapsto(b_0/(pt)+r b_2/(p^2 t),b_2/p^2)$ となる。
これで全ての零値を含め表の代表が得られる。
\path{computations/verify_plane_quotient_jordan.py} は線型消去、残余二次式、倍率を記号的に照合する。

スカラーの場合、$\alpha:\bigwedge^2 V\to V$ はJacobi恒等式を満たす括弧、
定数は $d\beta=0$ を満たし、持上げ変更はコバウンダリの加算である。
導来環の次元0では交代形式の階数が0か2。次元1では
$[x,y]=z$（中心導来）または $[h,x]=x$（非中心導来）に基底を取れる。
前者の $H^2$ は $x^*\wedge z^*,y^*\wedge z^*$、後者は
$h^*\wedge z^*$ で張られ、自己同型は非零類で推移的に働く。
次元2では可解Lie環の導来環が冪零であることから二次元アーベルであり、
残る文字の作用 $P$ は全射である。コサイクル式は
$\operatorname{tr}P\cdot\beta(x,y)=0$、残る二成分はコバウンダリである。
跡零なら $P$ の中央化群の行列式で非零類を同定できる。
導来次元3では可解でなく、Levi分解と次元から $\mathfrak{sl}_2$ である。
括弧が $\bigwedge^2V\to V$ の同型なので全ての定数はコバウンダリとなる。
\end{proof}

\begin{theorem}\label{thm:plane-quotient-unmarked}
上の全表について平面を指定しない同型類は、可換多項式窓を除いて、
以下で得る高々二つの指定付き標準形の順序なしリストで決まる。
最初の平面 $t=0$ に加え、四つの可換変数での三つの二次式
\[
 f_i(x,t)=r_i(x,x)-t\sum_j M_{ji}x_j-t^2b_i
\]
の共通一次因子 $v\notin kt$ があれば、その平面を第二の印に取る。
各印について定理\ref{thm:finite-normal-line-recognition}の $g_i$ で中央化し、
表の該当する代表を記録する。可換多項式窓は一つの同型類である。
これらは平面商の全特殊値に通用し、全窓の各辺射影ファイバーは空または一点である。
\end{theorem}
\begin{proof}
$t=0$ の点で中心関係は次の $t'=0$ を強制し、次点は三生成商の射影自己同型で一意である。
$t\ne0$ では中心関係により次点は同じ点になり、残る条件は $f_i=0$ である。
右からも同じなので全射影に正次元ファイバーはない。
外に開部分を持つ別の平面は従って共通一次因子に等しく、反対に各その因子は
対角グラフを与える。全 $f_i=0$ なら三生成関係は交代関係で両補正も零である。
それ以外では命題\ref{prop:surface-degree-bound}から平面は合計高々二つである。
全位置で同じ表示を用いたとき最初の平面はそのグラフで移り、第二平面は対角で移る。
定理\ref{thm:plane-surface-normal}から両者とも正規一次直線列を与える。
同型は平面を送るのでリストは不変であり、一つの印を合わせて標準形が一致すれば
定理\ref{thm:regular-normal-linear}が裸の窓の同型も与える。
\end{proof}

% !TeX root = ../papers/nc-p3.tex
\section{一次元の点対応の次数と偏極}
\label{sec:point-graph-curves}
有限双対が標準Frobenius型で、四辺の左右射影ファイバーが空または一点、
一辺の像の次元が1とする。命題\ref{prop:surface-degree-bound}で全辺を共通の
スキーム $X$ に同定する。位置 $i$ の超平面束の引戻しを $L_i$ とする。
\begin{theorem}\label{thm:point-graph-curve-degree}
$X$ の一次元部分の総次数は、どちらの射影でも重複込みで高々7である。
任意の既約成分上で $\deg L_i$ は位置に依存しない。また $\operatorname{Pic}(X)$ で
\[
 L_i L_{i+2}^3=L_{i+1}^3 L_{i+3}.
\]
特に滑らかな楕円四次成分は存在すれば唯一で、その上で
$K=L_{i+2}L_i L_{i+1}^{-2}$ は $i$ に依存せず $K^{12}=\mathcal O$ である。
\end{theorem}
\begin{proof}
$B_j=\prod_{a=i+j}^{i+3} L_a$ と置き、$B_4=\mathcal O$ とする。
四文字のKoszul候補複体を点列上で評価すると、その各項は右端から順に
$B_j$ の $1,4,6,4,1$ 個の直和である。両端の微分は文字の点評価とその双対で階数1、
中央二微分は隣接関係の点評価と相補位置のその転置で階数3である。
複体条件と各ファイバーの完全性から局所分裂完全となる。
従って行列式の交代積は自明であり、表示した束の式を得る。
次数の三階差が零で、各成分の位置ごとの次数は二次以下の多項式である。
標準四移動は有限個の成分を置換し、各四位置の次数の集合は有限であるから、この多項式は定数である。

一辺の六関係が張る空間を $R$ とし、一般の超平面 $U\subset R$ を取る。
$U$ の五式の共通零点 $C\subset\mathbb P^3\times\mathbb P^3$ は完全交叉曲線である。
実際、元の基点集合の次元は1以下なので、一般の線型結合を順に選ぶと、
それ以前の交叉の次元2以上の成分を含むことはない。両射影は $C$ 上でも閉埋込みになる。
六式の片側評価が階数4の点で一式を除いても階数3以上である。階数3の点では、
階数を落とす $U$ は評価核（三次元）を含む超平面であり、これらは $\mathbb P^2$ をなす。
点自体は一次元以下の集合を動くので、両側についてこの閉集合の像は
$\mathbb P(R^*)$ を覆わない。残る $U$ の決定式軌跡で単元な三次小行列を使うと
次点の三座標比が正則に決まり、残りの式が射影像のスキームを定める。
随伴公式と交叉環より $\omega_C=\mathcal O_C(1,1)$、両次数10、算術種数11である。

$\mathcal O_X$ から零次元台の最大部分加群を除いた商の定める純曲線を $Y$ とする。
$C$ 内でそのイデアルの零化イデアルが定める残余曲線を $Z$ とする。
Gorenstein曲線上の双対性で $I_{Y/C}\otimes\omega_C=\omega_Z$ である。
実際、$\mathcal O_Y$ は極大Cohen--Macaulayで、その $\mathcal O_C$ 双対は
$\operatorname{ann} I_{Y/C}$、二重双対は元に戻る。イデアルの完全列を双対化すれば
残余商も純で、$\operatorname{ann} I_{Z/C}=I_{Y/C}$ を得る。
双対化層は各成分の一般点で商と同じ長さを持つため $\deg Y+\deg Z=10$ である。
第六式は $I_{Y/C}(1,1)$ の非零切断を与える。五式のKoszul解消を $(1,1)$ でひねると、
左の負の対角束は両射影空間のKünneth公式で次数5以下にコホモロジーを持たず、
$C$ を消す双一次式は $U$ だけだからである。Serre双対で $h^1(\mathcal O_Z)>0$ を得る。

射影三空間の純曲線で次数2以下なら $h^1(\mathcal O)=0$ である。
被約なら直線、平面二次曲線、または交わらない二直線である。
一本の直線 $T$ 上の二重構造では $I_T^2$ は一般点で零になり、純性により曲線上で零となる。
よって $I_{T/Z}$ は $T$ 上の階数1の無捩れ層で、$\mathcal O_T(-1)^2$ の商である。
これは次数 $\ge-1$ の直線束なので、構造層の完全列から同じ消滅が従う。
従って $\deg Z\ge3$ であり $\deg Y\le7$ である。

楕円四次成分は二本存在できず、標準四移動がその成分を保つ。
成分上で $D_i=L_{i+1}L_i^{-1}$ とすると $D_{i+1}=K D_i$。
四移動の引戻しが $\operatorname{Pic}^0$ に誘導する自己同型を $u$ として
$uK=K,\ uD_i=K^4D_i$（乗法記法）である。
$u=1$ なら $K^4=1$。他の場合 $u$ の位数は $2,3,4,6$ で、
$u-1$ の核は各々 $2,3,2,1$ で消える。これは各円分多項式を $1$ で評価すれば従う。
\end{proof}

\begin{corollary}\label{cor:point-graph-excess}
一辺のスキームが滑らかな連結曲線とこれに交わらない有限スキーム $F$ の直和であれば、
曲線の片側次数 $d$ と種数 $g$ に対して $\operatorname{length}F=18-4d+2g$ である。
従って楕円四次曲線なら残りの長さは4である。
\end{corollary}
\begin{proof}
六つの双一次式の最高Chern数は $\binom63=20$。
曲線に沿って六切断はイデアルを生成し、法束から $\mathcal O(1,1)^6$ への単射の余核は直線束である。
超過交叉公式で曲線の寄与はその次数となる。接束完全列より法束の次数は
$8d+2g-2$、余核の次数は $4d-2g+2$ である。
孤立部は滑らかな六次元空間で六式の完全交叉なので各局所長が寄与する。
\end{proof}

% !TeX root = ../papers/nc-p3.tex
\section{楕円四次点成分の二次核と全次数の判定}
\label{sec:elliptic-point-quartic}
前節の有限双対Frobenius窓のうち、点グラフの一次元部分に滑らかな楕円四次曲線 $E$ を持つものを取る。
平面内の滑らかな四次曲線は種数3なので $E$ は非退化である。従って
$V_i=H^0(E,L_i)$ は完全線型系である。$D_{ij}=H^0(E,L_i\cdots L_{j-1})$ ($i\le j$)
を積で結び、$A$ を元の窓の二次代数とする。定理\ref{thm:elliptic-pencil-all-as}の
評価核のフィルトレーションは各 $L_i$ の次数4だけを使うため、ここでも $D$ は一次生成かつ二次関係である。
$A\to D$ の核を $I$、$W_i=I_{i,i+2}$ と書く。$\dim W_i=2$ である。

\begin{theorem}\label{thm:elliptic-quartic-as-test}
$A$ が標準Koszul左右AS正則となる必要十分条件は、全四位置で
\[
 \ker(W_i\otimes W_{i+2}\longrightarrow A_{i,i+4})
\]
が階数2のテンソルを含むことである。そのとき核はその一本の直線であり、
$W$ の生成する部分代数は偶奇それぞれで次数を倍にした二変数多項式型の
添字付き代数である。$A$ はその上で左右自由である。
\end{theorem}
\begin{proof}[切断環を係数とするホモロジー]
始点 $i$ の右Koszul候補複体に $D$ をテンソルし終点 $i+n$ を取る。
これは前節の分裂局所完全な束複体に共通の直線束を掛けて切断を取ったもので、
項 $j$ の偏極は $L_{i+j}\cdots L_{i+n-1}$（逆順の区間は逆束）である。
$n\ge5$ なら全項の $H^1$ が零なので完全である。$n=4$ では $H^1$ は項4の
$H^1(\mathcal O_E)=k$ だけである。二行のhypercohomology系列は全体が零へ収束し、
この一項からの二階微分が切断列の項2のホモロジーへの同型となる。
$n=3$ で $H^1$ を持つ二項の間の射は同型である。その双対は末尾文字の完全線型系
$V_{i+3}\to H^0(L_{i+3})$ そのものであり、三次の係数との同定はFrobenius積対合による。
$n=2$ では六関係の $V_i\otimes V_{i+1}$ への包含と全切断乗法を使い、
項1のみに二次元が残る。$n=1,0$ は生成元の同型と次数零の単純を与える。
従って得られた正ホモロジーは、項1の次数2に二個、項2の次数4に一個である。
反対向きにも同じ議論が通る。

元の単純加群の極小自由解消は最初の二微分までこの候補複体と一致する。
三次の四つの重なりは独立な最小syzygyであり、極小解消の項3にこれらを含められる。
従って実際の $\operatorname{Tor}_1^A(k,D)$ は上の項1に等しく、
$\operatorname{Tor}_2^A(k,D)$ は上の項2の商である。左右を交換すれば、
各 $e_iD$ の右表示で $I_{i,*}=W_i A$、その関係加群は次数4の高々一元で生成される。
左表示でも $I_{*,j}=A W_{j-2}$ である。$A$ が標準なら、候補複体は最初から解消なので、
$D$ の極小左右解消はともに長さ2、生成次数は $0,2,4$、個数は $1,2,1$ である。
\end{proof}

\begin{proof}[十分性]
仮定のテンソル $r_i$ は $W_i\otimes A_{i+2,i+4}$ においても非零であり、
上の唯一のsyzygyを生成する。矢 $W_i$ と関係 $kr_i$ で定まる抽象代数を $P$ とする。
偶奇ごとに一位置の基底を決め、両方向に次の基底を選べば、
非退化な関係をいつも $x y-y x$ にできる。従って $P$ の長さ $2h$ 成分の次元は $h+1$ である。
$D$ の各成分の基底を $A$ に持ち上げ、左から $P$ の単項式を掛けると $A$ を張る。
これは $I_{i,*}=W_i A$ から次数帰納で従う。右からも張る。
従って各成分は $H_D(z)/(1-z^2)^2=(1-z)^{-4}$ の係数以下の次元を持つ。

逆向きの不等式には終点 $j$ を固定し、$a_n=\dim A_{j-n,j}$ とする。
上記の右表示から $a_n-2a_{n-2}+a_{n-4}=\dim D_{j-n,j}+k_n$、
ここで $k_n\ge0$ は最後の自由加群からの写像の核次元である。
$(1-z^2)^{-2}$ の係数が正または零なので両不等式は等号、全 $k_n=0$ である。
張るために使った基底は左右とも独立で、$P$ は $A$ に埋まり、$A$ は左右 $P$ 自由である。
$e_i D$ は $P$ の始点単純を $A$ に誘導した加群である。自由性によりその解消も誘導でき、
随伴と $P$ の二変数Koszul解消から $\underline{\operatorname{Ext}}_A^q(e_iD,A)=0$ ($q<2$) を得る。
ここで係数の $A$ は右 $P$ 自由でもある。環変更により $A$ の右単純から $A$ への
次数0,1のExtも消える。左側でも同様である。
標準候補複体の先頭二写像は生成・二次表示で完全であり、末尾二箇所も
反対側の先頭二写像の双対であるためこのExt消滅で完全となる。
残る中央は既に求めたHilbert関数とEuler式で完全である。
Frobenius対合による左右解消の双対性がAS条件を与える。
\end{proof}

\begin{proof}[必要性]
標準の場合を考える。唯一の楕円四次成分は標準四移動で保たれるので、
$A\to D$ は四頂点の正次数付き代数の全射に折り返せる。以下の双対 $(-)'$ は局所有限な次数付き双対である。
$D$ の長さ2の極小右解消を双対化して得る
$\underline{\operatorname{Ext}}_A^2(D,A)$ が左から $I$ で零化されることを示す。

まず有限次元右加群 $F$ には自然な左右双対
$\underline{\operatorname{Ext}}_A^4(F,A)=F'$（頂点・次数の標準移動とNakayamaねじれを除く）があり、
他のExtは零である。この同定の作用も確認しておく。$T_j$ を次数 $j$ のKoszul重なりとし、
対角解消 $A\otimes_{A_0}T_j\otimes_{A_0}A$ の微分を最初と最後の文字の積の符号付き和とする。
これは片端の次数でフィルターして単純のKoszul解消を使えば対角を解消する。
左から $F$ をテンソルした右自由解消の $A$ 双対の項は
$A\otimes_{A_0}T_j^*\otimes_{A_0} F'$ である。
$T_j^*$ をFrobenius積対合で相補位置の $T_{4-j}$ と同定すると、両端の微分は
逆順の二つの文字の積に移る。一周する端でだけtraceの巡回によるNakayama移動が入る。
従ってこれは $F'$ の左対角解消を逆に読んだものである。符号の次数自己同型と標準四移動は $I$ を保つ。

$F_N=A/A_{\ge N}$ とし、右加群の次数付きtorsion関手を
$\Gamma_A(M)=\varinjlim_N\underline{\operatorname{Hom}}_A(F_N,M)$ と書く。
$F_N$ の左作用を前合成すると結果は右加群になる。上の自然な同定と
$\varinjlim F_N'=A'$ により、$R\Gamma_A(A)$ は次数4だけに、標準のねじれ・移動を持つ
双加群 $A'$ を持つ。右 $D$ の有限自由右 $A$ 解消にこの関手を施せば、各自由項の導来torsionは
その $A$ 双対の次数付き双対にこの同じ移動を施したものである。
自由項間の射は左からの行列なので上の双加群同定の自然性で微分も対応する。
特に $\underline{\operatorname{Ext}}_A^2(D,A)$ は $R^2\Gamma_A(D)$ の次数付き双対のねじれ・移動である。

$R\Gamma_A(D)$ の右作用が $D$ を通ることは、ここではnoetherian仮定なしに確認できる。
右 $D$ 入射加群 $J$ に対して環変更で
$\underline{\operatorname{Ext}}_A^q(F_N,J)
=\underline{\operatorname{Hom}}_D(\operatorname{Tor}_q^A(F_N,D),J)$ となる
（次数移動した入射も使う）。$D$ の左 $A$ 解消の生成次数は $0,2,4$ であるから、
$q>0$ のTorは内部次数 $[N,N+3]$ にしかない。下限は $N$ 未満で元の解消のまま、
上限は $F_N$ と最も高い生成次数から従う。よって $N'\ge N+4$ からの遷移は零であり
$J$ は $\Gamma_A$ 非輪状である。$D$ 入射解消で計算すれば
$\Gamma_A(J)=\Gamma_D(J)$ は右 $D$ 加群であり、比較は右 $A$ 作用を保つ。
従って要求した $I$ による零化が従う。

最後に各始点の $D$ の右解消の唯一の最終列を $(u_i,v_i)^t$ とする。
$u_i,v_i\in A_{i+2,i+4}$ であり、その双対の余核は左巡回加群である。
その最小生成元に $W_{i+2}\subset I$ を左から作用させると零になる。
この内部次数で双対微分の像は $u_i,v_i$ のスカラー線型包に等しい。
よって $W_{i+2}=\langle u_i,v_i\rangle$、二係数は独立で、syzygyは $W_i\otimes W_{i+2}$ 内の階数2である。
\end{proof}

\begin{corollary}\label{cor:elliptic-point-finite-atlas}
この枝の添字保持同型類は次の有限判定付きパラメータと有限群同一視で尽くされる。
楕円曲線の各 $j$ について原点付きWeierstrass代表 $(E,o)$ を固定し、
$L_0=\mathcal O(4o)$、$L_1\in\operatorname{Pic}^4(E)$、$K\in E[12]$、
$L_2=K L_1^2/L_0$ とする。$a=0,1$ で切断乗法の八次元核から
$R_a\in\operatorname{Gr}(6,8)$ を一つずつ取る。
\begin{enumerate}
\item 二窓の三次重なり $U$ の次元が4であることを課し、$V_3=U^*$ と包含テンソルから
定理\ref{thm:frobenius-tensor-reconstruction}の条件と $R_a(w)=R_a$ ($a=0,1$) を検査する。
\item 四つの点対応の左右評価で階数2以下の点がないことと、最初の像が一次元であることを検査する。
各 $W_i$ は $E$ 上に制限して零になる二文字式を $R_i$ で割って得る。
復元した二辺への $E$ の埋込みは核直線の射で運ぶ。
\item 定理\ref{thm:elliptic-quartic-as-test}の四つの乗法核を検査する。
\end{enumerate}
生き残るパラメータは同じ $j$ で $E[4]\rtimes\operatorname{Aut}(E,o)$ の引戻しにより
二つの束と二つの関係が全て対応する時かつその時に限り同型となる。
最小周期は同じ比較を一移動、二移動した二窓に施して決まる。
\end{corollary}
\begin{proof}
任意の次数4の偏極は原点の平行移動で指定の $L_0$ にでき、残る平行移動は $E[4]$ である。
切断乗法の核の条件により最初の二辺は指定の $E$ のグラフを含む。
二番目の検査の下では全辺の射影像が同型で、他の成分や埋込み点を許しても
定理\ref{thm:point-graph-curve-degree}で楕円四次成分は一意となる。
三番目の検査は必要十分で、全ては有限次元の線型代数と射影消去で行える。
同型は唯一の曲線を各位置で保ち、グラフにより各位置の曲線同型は共通である。
完全線型系なのでその引戻しが文字の写像をスカラー倍まで決める。
逆に二窓を対応させる写像は系\ref{cor:two-window-extension}により全体に延びる。
移動した二窓は最初の偏極を再び $\mathcal O(4o)$ に直して同じ比較を行う。
\end{proof}

% !TeX root = ../papers/nc-p3.tex
\section{特異・可約な種数一の四次部分曲線と印の移動}
\label{sec:genus-one-point-marks}
各辺の左右の射影ファイバーが空か一点、像 $X$ が一次元である有限Frobenius窓を扱う。
$X$ の純四次部分曲線 $Y$ がGorensteinで、$H^0(\mathcal O_Y)=k$、
$\omega_Y=\mathcal O_Y$ を満たすときこれを印として使う。重複成分も許す。

\begin{lemma}\label{lem:genus-one-quartic-mark}
この印は各位置で二つの二次曲面の完全交叉となる。逆に一位置で型 $(2,2)$ の完全交叉部分曲線なら
全ての位置で上の印となる。相異なる印が同じ既約成分を共有することを排除しない。
印上の $L_i$ の積の全切断環 $D(Y)$ は一次生成で、
長さ $n>0$ の各成分の次元は $4n$ である。
\end{lemma}
\begin{proof}
各成分の次数は定理\ref{thm:point-graph-curve-degree}で一定であり、局所長を掛けて足せば
全位置で総次数4となる。真の純部分曲線 $Z\subsetneq Y$ の双対化層は
$\operatorname{Hom}_Y(\mathcal O_Z,\mathcal O_Y)$ なので大域切断は零である。
従って $h^1(\mathcal O_Z)=0$。多重次数零の可逆層 $M$ については
$H^0(M)\ne0$ なら $M=\mathcal O_Y$ であることも示す。
非零切断の像はある純部分曲線 $Z$ の構造層である（可逆層内に有限台の部分はない）。
$Z$ が真なら $M|_Z$ は自明である。実際、曲線の各被約成分で次数を取る写像の核は連結であり、
$h^1(\mathcal O_Z)=0$ ならこの核は零である。具体的には被約曲線の正規化の単元の完全列で、
各特異点の貼合せの群はトーラスと単冪群の拡大、滑らかな各成分の次数零Picard群は連結である。
冪零部分を戻すときも $1+\mathcal N^j$ の逐次商は加法群である。
これらが次数零部分の連結性を与え、その接空間は $H^1(\mathcal O_Z)$ である。
すると $\operatorname{Hom}_Y(\mathcal O_Z,M)=H^0(\omega_Z\otimes M|_Z)=0$ に反する。
$Z=Y$ なら切断の有限余核はEuler標数差が零なので消える。

一般の超平面切断は非零因子で非空な零因子を持つ。
従って正の長さの偏極積の逆束は $\mathcal O_Y$ に入り大域切断が零である。
双対性とRiemann--Rochにより長さ $n>0$ の積は $H^1=0,h^0=4n$ となる。
平面内の純四次曲線は算術種数3なので各埋込みは非退化で、完全線型系になる。
二次式は少なくとも二本ある。独立な二本が $\ell a,\ell b$ と共通平面因子を持つとする。
$\mathcal O_Y$ のイデアル $(\ell)$ は直線 $a=b=0$ 上の $\mathcal O(-1)$ の商であり、
非零で、その台の純性から同じ $\mathcal O(-1)$ に等しい。
従って $Y\cap\{\ell=0\}$ は次数3、Euler標数0である。平面内の純三次部分のEuler標数は0
なので孤立・埋込み点はなく、真の純部分曲線に $h^1=1$ を持つことになり矛盾する。
ゆえに二式は完全交叉であり、次数一致と純性で $Y$ と等しい。逆向きは随伴と完全交叉の連結性
（構造層のKoszul解消で $h^0=1$）である。

$L_i$ の一般の二切断は基点自由である。鉛筆の完全列を次の束 $L_{i+1}$ でひねると
切断乗法の余核は $H^1(L_{i+1}/L_i)$ の商である。多重次数零の先の議論と双対性により
$L_i\ne L_{i+1}$ なら零である。一致する場合は二次完全交叉の射影正規性を使う。
長さ3以上では同じ鉛筆列に正の尾部がさらに掛かる。
障害群の双対は多重次数零の束に正の偏極積の逆を掛けた束の切断である。
一般の正の切断を掛ければ多重次数零の束に単射し、後者の切断は零または定数なので障害は消える。
従って全切断乗法が全射となる。
\end{proof}

\begin{theorem}\label{thm:finite-orbit-quartic-marks}
四移動が印 $Y$ に持つ軌道 $\mathcal M$ が有限であるとする。
$W_i(Z)=\ker(A_{i,i+2}\to D(Z)_{i,i+2})$ は印 $Z$ ごとの二次元空間である。
次の条件が全次数の標準Koszul左右AS正則性と同値である。
軌道上に一つの置換 $\pi$ があり、各 $Z\in\mathcal M$ と各位置で
\[
 \ker\bigl(W_i(Z)\otimes W_{i+2}(\pi Z)\longrightarrow A_{i,i+4}\bigr)
\]
が階数2のテンソルを含む。位置は $4|\mathcal M|$ を法として調べればよい。
特に被約な四次完全交叉部分曲線を一つでも含む枝と、
重複を許す四次完全交叉が一次元部分の全体を占める枝にこの有限判定が適用される。
\end{theorem}
\begin{proof}
一つの固定した印で定理\ref{thm:elliptic-quartic-as-test}の束複体を評価する。
局所分裂、全ての正の偏極積に対する $H^1$ 消滅、$\omega_Z=\mathcal O_Z$ と
文字の完全線型系だけが使われている。従って同じ証明で左右とも第一Torは次数2に二つ、
第二Torは次数4に一つの単純の商である。標準の場合はこれで長さ2の左右解消を得る。
$4|\mathcal M|$ 頂点に折り返せば $D(Z)$ はこの代数の商である。
同定 $R^2\Gamma_A(D(Z))'\simeq\underline{\operatorname{Ext}}^2_A(D(Z),A)$ に付随するねじれは
標準四移動またはその逆（および次数スカラー）であり、全ての始点・印に共通である。
従って最後の右syzygyの二係数は今度は四移動した印の二次核を張る。
実際、ねじれたイデアルの左零化と二次元性から従う。
軌道でこのねじれを置換 $\pi$ として使えば必要性を得る。

逆に任意の $i,Z$ を固定する。偶奇二列の印 $Z_j$ を
$Z_i=Z,\ Z_{j+2}=\pi Z_j$ で選ぶ。他方の偶奇列の初期値は任意でよい。
必要なら周期の公倍数に展開する。矢 $W_j(Z_j)$ と仮定の非退化関係による抽象部分表示 $P$
は二変数多項式型である。始点 $j$ の矢で割ると $e_jD(Z_j)$、終点 $j+2$ の矢で割ると
$D(Z_j)e_{j+2}$ となる。どちらも同じHilbert関数を持つ。従って基底の持上げと
左からの $P$ の語で $A$ が張られ、逆向きにも張られる。
始点 $j$ の二次核の関係は四次に高々一元で生成されるので、仮定の関係がこの始点の全関係を生成する。
定理\ref{thm:elliptic-quartic-as-test}と同じ固定終点のEuler不等式より
$H_A=(1-z)^{-4}$、両側の張る集合が基底となる。特に $e_i D(Z)$ は $P$ 単純の誘導で
$\underline{\operatorname{Ext}}_A^q(e_iD(Z),A)=0$ ($q<2$) を得る。
同じ印 $Z$ を固定したまま各始点 $i$ でこの構成を行えることが重要である。
全始点の消滅を合わせれば、この一つの商 $A\to D(Z)$ の環変更で単純からの消滅を得る。
反対側では所定の終点の二つ手前の印を指定する。Frobenius双対とEuler式で正則性が従う。

被約な印は有限個の既約成分の和集合に限られる。
また一次元部分の全部を占める純な印は、有限台の部分を除いた商として一意である。
したがって最後の二場合はいずれも有限軌道である。
\end{proof}

\paragraph{印を指定しない場合。}
被約な印は射影像の被約既約成分の和をすべて調べ、純四次かつ二つの二次式の完全交叉になるものを残せば
有限に列挙できる。候補の同型は、二連続窓と位置 $0,1,2$ の偏極および選んだ印上の
曲線同型で比較できる。具体的には印上で三位置とも完全線型系なので、偏極を保つ曲線同型の引戻しが
二窓の関係を移せば、系\ref{cor:two-window-extension}で十分である。一つ目で固定した印を
二つ目の有限リストのいずれかに移すものがあることが、印を忘れた同型の必要十分条件である。
移動した候補とも同じ比較を行えば最小周期を判定できる。

% !TeX root = ../papers/nc-p3.tex
\section{三文字の可換部分と次数七の点グラフ}
\label{sec:septic-point-curves}
\begin{proposition}\label{prop:commuting-three-point-graph}
有限双対が標準Frobenius型で、各辺の左右の射影ファイバーが空または一点とする。
三次元空間 $H$ と各位置での単射 $H\hookrightarrow V_i$ があり、
全辺で $\bigwedge^2 H\subset R_i$ なら、元代数は標準Koszul左右AS正則である。
可換三変数多項式環の反復 $P$ はこれらの文字で埋め込まれ、元代数は左右 $P$ 自由である。
三文字の直交点を次の位置の三文字の直交点へ運ぶという条件は不要である。
\end{proposition}
\begin{proof}
位置 $i$ の三文字を消す点 $p$ で関係を評価すると、三つの可換関係が零になるから階数は高々3であり、
ファイバー条件でちょうど3である。従って $p$ は点グラフ上にあり、前後に一意の点列を持つ。
$p$ から出る右点加群を $M$ とする。$M$ と各終点単純の左Koszul候補複体とのテンソル積は、
長さ $n$ で右から $\min(n,4)+1$ 項だけを持つ。全項があれば点列上の階数は $1,3,3,1$ で完全である。
従って正の第一、第二ホモロジーはそれぞれ内部次数1,2のみに三次元ずつある。
極小解消の初めの二写像は候補と一致し、項3は全ての三次重なりを含むので、
真の第一Torはこの第一ホモロジーに等しく、第二Torは第二ホモロジーの商である。
よって $M=e_iA/H_iA$ であり、$H_i\otimes e_{i+1}A\to e_iA$ の関係は次数2に高々三個で生成される。
実際 $\bigwedge^2H$ の三つはそこに独立にあるので、ちょうどこれらが生成する。

抽象多項式代数 $P$ のKoszul複体から、直前の完全性は
$\operatorname{Tor}^P_1(k_i,A)=0$ を意味する。各終点の有界下次数付き加群の極小表示を取れば
左 $P$ 自由である。始点がどこでも商 $M$ は各長さで一次元なので、各終点で自由生成元は
各長さに一個あり、$\dim A_{i,i+n}=\binom{n+3}{3}$ を得る。次数零の持上げも自由生成元であり
$P$ 自身が埋まる。逆向きの点列で右自由性も従う。

右単純から $A$ へのHomは、非零な一文字による右 $P$ 乗法が単射なので零である。
正次数根基 $e_iA_+$ から $A$ への任意の斉次右加群写像を $e_iP_+$ に制限する。
$A$ が右 $P$ 自由で $P$ 上の単純の第一Extが $P$ に対して零なので、
その制限はある元の左乗法である。その左乗法を引いた写像は $H_i A$ を消し、
点加群の正次数尾部を経由する。この尾部から $A$ への加群写像は零である。
実際、各点加群の斉次元は次の $H$ の少なくとも二次元の部分で消えるが、
自由な右 $P$ 加群の非零な元を消す非零一次式はない。
よって第一Extも零である。反対側でも同じ二消滅を得る。候補Koszul複体の一端の完全性は
生成・二次表示、他端はFrobenius双対とこの二消滅から従い、中央はHilbert関数による。
\end{proof}

\begin{theorem}\label{thm:septic-point-graph}
定理\ref{thm:point-graph-curve-degree}の上限7に達する場合、点グラフ $X$ は純曲線で、
算術種数は5である。特に $X$ が滑らかな整曲線なら、全ての有限Frobenius候補が全次数で正則である。
その添字保持同型類は次の有限条件を満たす順序付きデータ $(X,p,q)$ の曲線同型類で尽きる。
\begin{enumerate}
\item $X$ は種数5の滑らかな射影曲線で $g^1_3$ を持たず、$p-q$ はちょうど位数4のPicard類を持つ。
\item $L=\omega_X(-p),\ N=\omega_X(-q)$ の完全線型系は閉埋込みで、
$R=\ker(H^0(L)\otimes H^0(N)\to H^0(LN))$ を取る。
$R,R^{\rm op}$ の交互窓に有限双対Frobenius検査を施し、全辺で射影ファイバーが空か一点、
一辺の像が一次元であることを課す。
\end{enumerate}
最小周期は $p,q$ を交換する曲線自己同型があれば1、それ以外は2である。
この記述は有限条件の解の存在やその次元を仮定しない。
\end{theorem}
\begin{proof}[残余三次曲線]
前節の五式の完全交叉 $C$ において、純曲線部分 $Y$ の残余 $Z$ は次数3で
$h^1(\mathcal O_Z)>0$、両側の射影三空間に埋まる。そのような純三次曲線は平面三次曲線である。
実際、一般の直線への線型射影 $f$ は有限平坦であり、跡で単位を分裂させると
$f_*\mathcal O_Z=\mathcal O\oplus\mathcal O(a)\oplus\mathcal O(b)$、$a\ge b$, $b\le-2$ となる。
有限平坦性は純性により直線上無捩れであることから従う。残りの二つの空間座標は、
射影直線の局所座標で割るとこの代数を生成する。それらの跡零成分は
$\mathcal O(-1)$ からの写像であり、$\mathcal O(b)$ 成分は零なので $a\ge-1$。
$\mathcal O(a)$ の平方が $\mathcal O(b)$ に非零成分を持たなければ二成分だけの部分代数に留まるため
$2a\le b$ である。従って $(a,b)=(-1,-2)$、残り二座標の跡零成分は比例する。
これが空間の一次式を与え、純性から平面内の三次因子である。

随伴により $\omega_Z=\mathcal O_Z$ かつ $h^0(\mathcal O_Z)=1$。
第六式は $I_{Y/C}(1,1)=\omega_Z$ の生成元である。従って六式は $Y$ をスキームとして切り出し、
$X=Y$ である。残余の両側次数は3なので
$\chi(I_{Y/C})=\chi(\mathcal O_Z(-1,-1))=-6$ であり、$\chi(\mathcal O_X)=-10+6=-4$ を得る。
\end{proof}

\begin{proof}[滑らかな場合の偏極と関係]
次数7の滑らかな種数5の空間曲線は平面にも二次曲面にも含まれない。
平面では平面種数公式に反する。滑らかな二次曲面上なら二次数 $(a,7-a)$ に対して
$(a-1)(6-a)=5$ に整数解がない。二次錐の解消 $\mathbb F_2$ 上の狭義変換は
元の滑らかな曲線に同型である（射は双有理かつ有限）。負切断 $s$ とファイバー $f$ で
その類は $as+7f$、随伴による種数は $-a^2+7a-6$ でやはり5にならない。
各 $L_i$ は独立な四切断を持ち、Riemann--Rochで次数1の $\omega_X/L_i$ が有効だから
一意な点 $p_i$ で $L_i=\omega_X(-p_i)$、$h^0(L_i)=4$ となる。
次数3以下の鉛筆があれば、その動く部分の $p_i$ を通る因子に点を加えて
$p_i+D$、$\deg D=2,\ h^0(p_i+D)\ge2$ とできる。
Riemann--Rochで $h^0(L_i(-D))\ge3$ となり埋込みに反する。

偏極の三階差が零なので $(1-T)^3(1+T)=1-2T+2T^3-T^4$ を使えば
$p_i+2p_{i+3}\sim 2p_{i+1}+p_{i+4}$ である。
三次の相異なる有効因子が線型同値なら鉛筆になるので、両辺は因子として等しい。
重複点の比較により $p_{i+3}=p_{i+1}$、すなわち二位置で戻る。
全て同じ点なら空間曲線に二次式がないため六関係が全て可換となり像が三次元になる。
従って相異なる二点 $p,q$ が交互に現れ、三階差から $4p\sim4q$ となる。
次数2以下の鉛筆もないので位数はちょうど4である。

$U=H^0(\omega_X)$ を五次元の標準座標空間とする。
標準像を消す二次式の空間の次元は3である。実際 $\dim\operatorname{Sym}^2U=15$、
$h^0(\omega_X^2)=12$ で下限3を得る。点 $p$ での微分は三次元の余法空間に入り、
微分の消える二次式は $\operatorname{Sym}^2 H^0(\omega_X(-p))$ に属するから
上の空間二次式の非存在で零である。相異なる二つの超平面
$V_p,V_q\subset U$ の積の対称二乗での像は、$p,q$ で消える13次元の空間全体となる。
共通の三文字 $H=V_p\cap V_q$ と両側の補文字で単項式を書けば、補文字の二つの平方だけが欠ける。
上の標準像の三式は全てここに入り、曲線上の積の階数は10である。
従って六関係は全乗法核であり、その中に $\bigwedge^2H$ が入る。
全窓は記載した反転列となり、命題\ref{prop:commuting-three-point-graph}が正則性を与える。

逆に表の入力が有限検査を通れば、最初の点対応は指定の曲線を含む。
前節の次数上限と本節の純性から点グラフがその曲線に等しい
（純な曲線への有限台のイデアルは零）。同型は点グラフと各埋込みの偏極を保ち、
二点を順序付きで復元する。逆にこの二点を対応させる曲線同型は全ての切断乗法核を保つ。
一移動は二点の順序を交換し、二移動は元の表示を返す。
\end{proof}

% !TeX root = ../papers/nc-p3.tex
\section{滑らかな七次点曲線の全係数と五次元の同型類}
\label{sec:septic-bielliptic-forms}
\begin{theorem}\label{thm:septic-bielliptic-forms}
前節の滑らかな七次枝の有限条件は次のように解ける。$R=r^2,S=s^2,H=h^2+\lambda rs$ と書き、
射影四空間の座標 $r,s,h,u,v$ における三式を
\[
 ru=F=f_S S+f_H H,\qquad sv=G=g_R R+g_H H,\qquad
 v^2-u^2=4D=4(d_R R+d_S S+d_H H)
\]
とする。 $\lambda$ を含む八係数の許容範囲は三式が滑らかな完全交叉曲線を定める開集合全体である。
二点は $r=s=h=0,\ v=\pm u$ であり、前節の順序付き二射影から代数を作る。
同型の全同一視は
\[
 (r,s,h,u,v)\longmapsto (a r,b s,c h,e u,e v),\qquad
 (r,s,h,u,v)\longmapsto(s,r,h,v,u)
\]
から誘導される係数変更で尽くされる。後者は $D$ の符号も変える。
座標の共通倍率を除く三次元トーラスの重み作用なので、零係数を含め台とその重みの
整数核単項式を比較し、最後に交換を試せばよい。
この開集合は非空で同型類の次元は5である。最小周期1の必要十分条件は $\lambda$ だけを $-\lambda$ にした係数が同じ軌道に入ることであり、一般点の最小周期は2である。
全曲線が $h$ の符号反転で滑らかな楕円曲線の二重被覆になる。
\end{theorem}

\begin{proof}[標準二次式の制約]
前節の順序付き曲線から始める。三つの共通文字の空間を $J$ とし、
第一の補文字を $x$、第二を $y$ として標準空間を $J\oplus kx\oplus ky$ と書く。
標準像の二次式の基底を
\[
 Q_j=xB_j+C_j,\qquad B_j=a_j y+b_j,\qquad C_j=y c_j+d_j
 \quad(b_j,c_j\in J,\ d_j\in\operatorname{Sym}^2 J)
\]
とする。両補文字の軸が二つの指定点である。
$B_j$ は $x$ 軸での勾配なので独立である（前節の余法空間への単射）。
全 $a_j$ が零なら二辺の三文字直交点の対が三つの可換関係と全 $Q_j$ を消す。
その対は相異なる二つの曲線点に由来するのでグラフ性に反する。
従って $a_1=1,a_2=a_3=0$ とできる。$b_2,b_3$ と $c_2,c_3$ はそれぞれ独立である。
曲線の標準像はこの三式の完全交叉である。実際 $J$ 座標が非零の開集合では
二つの射影像での $J$ 座標を同一に規格化でき、三式の解がそのまま点グラフに入る。
外に残る直線上の解は二指定点だけであり、完全交叉の次数8と標準像の次数が一致する。

第一三次重なりを外二文字の反転で分解する。反対称部分では、三つの $Q_j$ に掛かる
外文字 $f_j$ は $\sum Q_j f_j=0$ を満たす。完全交叉には一次係数のsyzygyがないので
全て零である。残りの重なりは三つの可換関係だけを両端に持ち
$\bigwedge^3J$ で尽きる。従って対称部分は三次元である。
その元を最初の二文字の六関係の右倍で表し、$Q_j$ の係数を $f_j=A_jx+w_j$、
$A_j\in k,w_j\in J$ とする。対称部分から $(f_j)$ への射は単射である。
核なら上の $\bigwedge^3J$ に入り、これは反対称だからである。

対称性を中央文字 $y$ で縮約し、外側の $x$ と $J$ の成分を比較すると
$\sum a_j w_j=\sum A_j c_j$、また $\sum c_j\wedge w_j=0$ となる。
中央を $J^*$ で縮約した同じ比較から、$J\otimes J$ の対称商では
$\sum b_j w_j=\sum A_j d_j$ となる。可換関係の $x$ 倍はこの対称商で消える。
従って $\sum B_j w_j=\sum A_j C_j$ であり、
$\sum A_j Q_j$ は $(B_1,B_2,B_3)$ に入る。この三一次式の共通直線は $x$ 軸の点での接線である。
接線は完全交叉曲線に含まれず、従って $(A_j)$ の像は高々二次元である。
三次元の対称部分には $A_j=0$ の非零元がある。上の式から $w_1=0$、
$(w_2,w_3)=t(b_3,-b_2)$、$t\ne0$ となる。従って
$c_2\wedge b_3-c_3\wedge b_2=0$。$b_2,b_3$ を基底に延ばして係数を比べると、
$c_2,c_3$ は同じ二平面にあり、$b_j\mapsto c_j$ の跡は零である。
この可逆写像は対角化でき、二固有値は互いに逆符号である。

二固有文字を $r,s$ に取り、$y$ を再尺度化する。第一式の $x,y$ の一次交差項を
それぞれ $y,x$ の $J$ による持上げ変更で吸収すれば、三標準式は
\[
 xy=D,\qquad xr-r y=F,\qquad xs+s y=G
\]
となる。ここではまだ $D,F,G$ は $r,s,h$ の任意の二次式である。
非可換な二窓ではこの順の式と逆順の式を交互に置き、$r,s,h$ は交換する。
次数3だけの書換えを考える。$r<s<h<x$（次位置では $x$ を $y$ に替える）とし、
各可換逆順対と補文字に続く $r,s$, 補文字対を解消する。
次数辞書式順序は各書換えで下がる。先頭単項式の三文字重複は
$hsr,xsr,xyr,xys,xyx$ の五個である。最初の二個の差は零、
残る三個の差は非零の定数倍を除き
$xF-Fx,xG-Gx,xD-Dx$ の書換え結果である。これは $yr=rx-F,\ ys=-sx+G$ を代入すれば従う。
二次先頭関係の左右倍の行消去では各重複に一つの差が残るから、
三次重なりの次元は5からこれらの差の階数を引いた数である。

$T\in\operatorname{Sym}^2J$ に対する $xT-Tx$ の書換えが零となる核は
$\langle r^2,s^2\rangle$ である。二つの平方は直接交換する。他の四単項式
$rs,rh,sh,h^2$ の結果にはそれぞれ $-2rsx,\ ryh,\ -syh,\ xh^2$ が現れ、
その各語は他の三つの結果には現れない。従って $D,F,G$ のこの核を法とした像は一本の直線を張る。
共通の代表二次式で $h^2$ の係数が零なら、標準空間の $h$ 軸の点が曲線上にあり、
三式のそこでの勾配は高々一方向となって滑らかさに反する。
平方完成で $h$ の $r,s$ による持上げを変更し、代表の $rh,sh$ 項を消す。
残る $rs$ 項を保持して共通代表を $H=h^2+\lambda rs$ とする。
$u=x-y, v=x+y$ に移り、$u$ に $r$ の倍、$v$ に $s$ の倍を加えて
$F$ の $R$ 係数、$G$ の $S$ 係数を零にする。この操作で $D$ に加わる項も $R,S,H$ の線型結合である。
\end{proof}

\begin{proof}[滑らかな全係数での正則性]
表示された滑らかな完全交叉は連結で、随伴で超平面束が標準束、種数が5である。
長さ3の部分スキームを含む直線があれば三二次式はその直線で恒等的に消え、整曲線性に反する。
よって二指定点からの射影は接方向を含めて閉埋込みになる。
三つの可換式と三つの混合式が定める点対応もこの曲線に等しい。
両側の $J$ 座標が非零の部分は既に見た。片側だけ零のところではもう一方の $J$ 座標の
一つを1にし、零となる側の補文字を1にする。可換式により小さい側の三座標は
他方の三座標の $t$ 倍である。三混合式は補文字軸の点での標準三式の
厳密変換（代入して $t$ を一つ除いた式）で、$t=0$ での式は独立な三勾配である。
従って例外ファイバーはその接線方向の一点で、その三つのアフィン方向変数に関する
Jacobianは可逆となる。ここでも滑らかな曲線で余分な構造はない。
両側で $J=0$ は $xy$ の式に反する。

以下 $x_i$ は偶数位置の $x$ と奇数位置の $y$ を表す。関係は
$x_i r=r x_{i+1}+(-1)^i F,\ x_i s=-s x_{i+1}+G,\ x_i x_{i+1}=D$ と書ける。
二移動の意味で $R,S$ は中心的である。$r,s,h$ は交換し、$x_i r^2, x_i s^2$
を二度書き換えると補正が相殺する。$x_i x_{i+1}r, x_i x_{i+1}s, x_i x_{i+1} x_{i+2}$
の両括弧付けは $F,G,D$ のそれぞれと $x$ の交換を与える。
滑らかさから $d_H,f_H,g_H$ の少なくとも一つは非零で、$H$ も中心的である。
抽象的な次数2の中央多項式環 $B=k[R,S,H]$ 上で、$r^2=R,s^2=S,h^2=H-\lambda rs$
と上の三つの $x_i$ の式を課すと元代数に戻る。$B$ 上の基底は
$r^a s^b$ ($0\le a,b\le1$) に $h,x$ の交互語を続けたものである。
実際、$r,s$ を左に移し、平方と補文字対を消す。$r,s,h$ の内部では全て交換し、
平方の値もそこで交換する。補文字を含む重複は $xrr,xss,xsr,xxr,xxs,xxx$ で、
最初の二つは $F,G$ と $r,s$ の交換、三番目は $r,s$ の交換と
$F$ の交代符号、最後の三つは $F,G,D$ の中心性により一致する。
次数辞書式で減少するためこれで基底が従い、Hilbert級数は
$(1+z)^2(1+z)/((1-z)(1-z^2)^3)=(1-z)^{-4}$ である。

二次双対を $T$ とする。先の三次計算は全四位置で $\dim T_3=4$ を与える。
ここで下付きは各始点からの長さとする。また $\dim T_4\ge1$ である。
実際、単純の極小解消の初めは線型で生成・関係数が4,6、
三次の最小syzygyは4個である。四次のEuler式は
$\dim\operatorname{Tor}_{4,4}-\dim\operatorname{Tor}_{3,4}=1$ を与え、
対角Torの双対は四次の重なりである。

$T$ の積を具体的に調べるため、元代数の奇数位置だけ $s\mapsto-s,\ y\mapsto-X$ とし、
偶数では $x=X$ とする。双対の文字を小文字 $r,s,h,z$（最後が $X$ の双対）で再び書く。
以下の関係は六式との係数対合で読める。元の $J$ の三可換式は
$r,s$ と $s,h$ について対称、$r,h$ について交代となる。残る三式の
$X$ を含む項は $X^2,\ Xr+rX,\ Xs+sX$ であり、その右辺は
$r^2,s^2,h^2\pm\lambda rs$ の線型結合である。最後の $rs$ は $r,s$ の対称関係を法として
交代表示に取れ、符号だけが位置で交互となる。従って双対で
\begin{itemize}
\item $zh=hz=0$、$h$ は $r$ と符号を替えて、$s$ とはそのまま交換する。
$z$ は $r,s$ と交換する。$rs-sr$ は位置別スカラー倍の $h^2$ である。
\item $r^2,s^2,h^2$ はそれぞれ $z\langle z,r,s\rangle$ で表せる。
これら三つの表示は位置に依存しない。従って $T/(h)$ は
$k[r,s,z]$ の三つの二次式による商の反復である。
\end{itemize}
この三式に共通の射影零点があれば、$h=0$ を付け加えた点のテンソル平方が
元の六関係に入る。奇数位置の変更を戻すと標準曲線上の二つの非零切断の積が
恒等的に零となり、整性に反する。従って $T/(h)$ は可換完全交叉で各始点の次元が
$(1,3,3,1)$、最上次の積対合は完全である。

$h^3=0$ かつ $hr^2=hs^2=0$ なので、$(h)$ は
$h^e r^a s^b$ ($e=1,2$、$0\le a,b\le1$) で張られる。
$h$ を動かすときに $r,s$ の交換から生じる項も $h^3=0$ で消える。
よって $T$ の次数別上限は $(1,4,6,4,1)$、これまでの下限と一致し、各語が独立である。
特に $h^2rs\ne0$ である。$(h)$ の一次層、二次層の任意の非零元は
$h,r,s$ の相補語との積で最上次に届く。具体的には一次層ではまず $h$ を掛け、
残る $r,s$ の平方零の二文字で完全対合を取る。
$z$ は $(h)$ を消すが、$h^2=z\ell$ と書けば $z(\ell rs)=h^2rs\ne0$。
従って $T/(h)$ のsocleの持上げに対する $z$ の乗法は非零である。
$(h)$ の外の任意の斉次元も、まず完全交叉商の対合でsocleの持上げに移し、
$z$ を掛ければ最上次に届く。左右で同じ議論が使え、全相補積が完全となる。
これで有限Frobenius検査が通り、前節の三可換文字の命題で全次数AS正則性が従う。
\end{proof}

\begin{proof}[同型、次元と実在]
標準座標で二点の順序を保つ変換は $J$ を保ち、補文字をそれぞれ
$ax+\ell,by+m$ ($\ell,m\in J$) に移す。二点を通る割線を含む二次式の鉛筆から、
前出の跡零写像の二固有直線 $kr,ks$ が無順序で決まる。
交換する場合は表示した有限交換を先に施せる。各直線を保つときは
鉛筆の $x,y$ の係数比から $a=b$ である。$xy$ の式の変換が三式の包に入るには
具体的に
$m=m_r r+m_s s,\ \ell=-m_r r+m_s s$ となる。
これは $xm+y\ell$ の補文字を含む部分が $r(x-y),s(x+y)$ の包に入るためである。
従って $u,v$ の許される平行移動はそれぞれ $r,s$ の倍に限られ、
正規化 $f_R=g_S=0$ がこの自由度を消す。残る $h$ の持上げには $rh,sh$ の係数が生じるが、
少なくとも一つの $H$ 係数が非零なのでこれも零である。従って表示した倍率と交換だけが残る。
$h$ の符号反転の固定点は $h=0$ にある。滑らかな曲線の固定点で微分は $-1$、
他の座標比は不変なので $h$ はそこで単純に消える。固定点数は標準次数8、Hurwitz公式で商の種数は1である。

$u$ を保ち $v,s$ の符号を変えれば二点を交換して $\lambda$ の符号だけを変える。
従って述べた周期判定を得る。一般の台では
$f_S^2d_R/d_S^2,\ g_R^2d_S/d_R^2$ は倍率不変で、有限交換で互いの負に移る。
一般にはこの二数は互いの負でなく、交換を伴う同一視はない。
$d_R,d_S,f_S,g_R,d_H$ を保つ倍率は $a=b=e,\ c^2=e^2$ で $\lambda$ を変えない。
従って共通倍率を除いた一般安定化群は有限で、$\lambda\ne0$ の一般点は最小周期2である。

非空性は正規化前の $a=2$ の例
\[
 xy=r^2+2s^2+3h^2,\quad xr-2ry=2r^2+5s^2+h^2,\quad
 xs+2sy=3r^2-s^2+2h^2
\]
で確認できる。$y$ を再尺度化して平方完成と二平行移動を施せば定理の表示になる。
$h^2$ を消した二次式二本の射影三空間での交叉と、そこでの $h=0$ の分岐因子は滑らかである。
付随コードは標数101で両Jacobianの特異イデアルが全射影チャートで単位を含むことを確かめる。
次元定理で二交叉の次元はそれぞれ1,0であり、射影性によって標数零でも滑らかである。
この二重被覆が滑らかな完全交叉を与える。従って八次元係数空間の非空開集合が得られ、
一般軌道の三次元を引いて同型類の次元は5である。
\end{proof}

% !TeX root = ../papers/nc-p3.tex
\section{曲線の二点縮約の階数低下因子}
\label{sec:opposite-curve-divisor}
四辺の左右射影ファイバーが空か一点で、標準有限Frobenius候補の点グラフが
滑らかな整曲線成分 $C$ を持つとする。$C$ が全位置で射影三空間を張り、
標準四移動がこの成分を保つとする。点列に沿う埋込み $j_i:C\to\mathbb P(V_i^*)$
で $j_{i+4}=j_i\circ\sigma$ となるよう $\sigma$ を定める。
$L_i=j_i^*\mathcal O(1)$、$K=L_i L_{i+2}/L_{i+1}^2$ は前出の一定の多重差である。
\begin{theorem}\label{thm:opposite-curve-divisor}
$C\times C$ の点 $(p,q)$ で、位置0と2の二点 $j_0(p),j_2(q)$ により $w$ を縮約する。
階数が1以下になる軌跡は有効Cartier因子 $D$ で、その類は
\[
 \mathcal O(D)=\mathcal O(\Delta+\Gamma_{\sigma^{-1}})
       \otimes\operatorname{pr}_1^*K^{-1}\otimes\operatorname{pr}_2^*K^{-1}.
\]
ここで $\Delta$ は対角、$\Gamma$ はグラフである。$D$ は両側に有限平坦で次数2となる。
\end{theorem}
\begin{proof}
$H_1(p),H_3(p)$ を $j_1(p),j_3(\sigma^{-1}p)$ がそれぞれ消す超平面とし、
これらを第一曲線上の部分束と見る。$S=\operatorname{pr}_1^*L_0\otimes
\operatorname{pr}_2^*L_2$ とする。二点縮約は隣接点での消滅により
$H_3^*\to H_1\otimes S$ を与える（第一因子からの引戻しを略記する）。
その右核には $j_3(q)$、左核には $j_1(q)$ の制限が入る。
それぞれの三座標の共通零点は $q=\sigma^{-1}p$, $q=p$ のグラフで、スキームとして重複度1である。
実際、座標は一つの埋込みに対する点の座標差を局所的に生成し、
滑らかな曲線の対角イデアルは可逆である。従ってこの共通因子を除けば
右・左の核にそれぞれ直和局所因子となる直線束
\[
 A=\operatorname{pr}_2^*L_3^{-1}(\Gamma_{\sigma^{-1}}),\qquad
 B=\operatorname{pr}_2^*L_1^{-1}(\Delta)
\]
が入る。局所的な因子による除算後も縮約が零なのは曲面が整だからである。
そこで二点行列は $(H_3^*/A)\to B^\perp\otimes S$ という二階数束の間の写像となる。

固定した $p$ の上でこの写像が恒等的に階数1以下になることはない。
仮にそうなら、$j_2(C)$ を三行三列の行列空間に線型に写した像は
零行列を除いて階数1である。両端の射影像はいずれも射影平面を張る。
これは $j_2(C)$ が射影三空間を張り、$w(j_0(p))$ の端支持がともに3であるためである。
関数体上で行列を列と行の積に書くと、両側の三成分が張る関数空間は各三次元である。
その積は少なくとも五次元となる。実際、各空間で一つの滑らかな点の位数が相異なる基底を取り、
二つの三元整数集合の和集合（対ごとの和の集合）には五つ以上の位数が現れる。
ところが九成分は元の四座標の線型結合なので四次元以下であり矛盾する。
$q$ を固定しても同様である。

よって二行二列の行列式は零でなく、階数低下を切る。
$\det H_1=L_1^{-1},\ \det H_3=(\sigma^{-1})^*L_3^{-1}$ より行列式の束は
\[
 S^2\otimes \det H_1\otimes\det H_3\otimes A\otimes B.
\]
$L_{-1}=(\sigma^{-1})^*L_3$ と束の三階差を用いれば主張の類となる。
どのファイバーも因子に含まれず、制限次数は2である。従って固有準有限射として有限、
純曲線から滑らかな曲線へ台の落ちる成分がないので平坦である。
\end{proof}

\begin{corollary}\label{cor:curve-divisor-monodromy}
次のモノドロミー制限が成り立つ。
\begin{enumerate}
\item 種数 $g\ge2$ で $\sigma=1$、または $\sigma\ne1$ で固定点が $2g-2$ 未満なら
$D=\Delta+\Gamma_{\sigma^{-1}}$ かつ $K=1$ である。
\item $g\ge2$ で $\sigma$ が対合なら、$K\ne1$ が可能なのは商が楕円曲線のときだけである。
この場合も $K^2=1$ であり、商写像 $f:C\to E$ に対して $D$ は
$f(q)-f(p)=e$、$e\in E[2]$ の一ファイバーとなる。
\item $C$ が楕円曲線で $\sigma$ が平行移動なら $K^2=1$、
$D$ は重複を許す二つの平行移動グラフの和である。
一般の楕円四次成分でも常に $K^6=1$ である。
\end{enumerate}
\end{corollary}
\begin{proof}
対角とグラフの自己交叉は $2-2g$ である。相異なる場合の交叉は固定点の個数である。
非恒等の有限位数自己同型の固定点は標数零で単純である。
交叉数が負なら該当成分が $D$ に入るので、一方を引き、残りの自己交叉も負であることを使う。
$\sigma=1$ でも二度引けば同じである。類を一方のファイバーに制限すると $K=1$ を得る。

対合の商の種数を $b$ とすると固定点数は $2g+2-4b$ である。
$b\ge2$ なら第一項に帰着する。$b=0$ なら二重被覆の束を $H=f^*\mathcal O(1)$ として、
二グラフの和の束は $H\boxtimes H$ である。$D$ は縦横成分を持たないので
$h^0(H/K)\ge2$ が必要である。種数2以上では次数2の鉛筆は唯一である。
実際、二つあれば両射の像は $\mathbb P^1\times\mathbb P^1$ の $(2,2)$ 曲線に双有理で
種数が1以下になるか、同じ二重被覆を通る。従って $K=1$ である。

$b=1$ なら差の射 $d(p,q)=f(q)-f(p)$ の零ファイバーが二グラフの和である。
一般ファイバーは二つの二次関数体の合成に対応して整となる。
実際、分岐因子は非空で有限であり、その一般の平行移動とは一致しない。
Stein分解で全ファイバーが連結となり、全てのファイバーは一般点で分離的な
二重被覆の積として被約である（Cartier因子の純性を使う）。
$D$ はファイバーと数値同値なので各成分が $d$ で点に写る。
各ファイバーの成分の交叉行列の核はファイバー全体の重複度列が張る。
実際、最大係数の成分でファイバーの定数倍を引いて交叉数零を使えば、
隣接成分も同じ係数となり連結性で伝わる。
ゆえに $D$ は丸ごとの一ファイバーである。両座標ファイバーへの制限から
$K$ は $E$ 上の互いに逆な二つの次数零束のどちらの引戻しでもある。
$f^*$ は単射である。非自明な束が自明になるならノルムで二乗が自明となり、
対応する非自明なエタール二重被覆を $f$ が経由するが、分岐に反する。
従ってこの二束は同じ二捩れである。

楕円曲線の平行移動の場合は、二グラフはいずれも $q-p$ のファイバーであり、
数値同値を用いれば $D$ も二ファイバーの和となる。対角に沿って次数零の束の差を制限すると
$K^2=1$ である。他の楕円自己同型では定理\ref{thm:point-graph-curve-degree}の
$uK=K$ により位数は2または3を割る。従って楕円四次ではどの場合も6で消える。
\end{proof}

\begin{corollary}\label{cor:septic-canonical-monodromy}
定理\ref{thm:septic-bielliptic-forms}の全ての滑らかな七次窓で、標準四移動が
点曲線に誘導する自己同型は $h\mapsto-h$ である。
商を $f:C\to E$、順序付き二点の像を $P,Q$ とすれば
\[
 D=\{(a,b):f(b)-f(a)=P-Q\},\qquad P-Q\in E[2]\setminus\{0\}.
\]
\end{corollary}
\begin{proof}
偏極が二つの標準射影を交互に繰り返すので、二階差は $\mathcal O(2p-2q)$ またはその逆である。
これは非自明な二捩れだから、前系で標準四移動は恒等になれない。
八係数の一般点では、順序付き二点を保つ曲線自己同型は恒等と $h$ の符号変更だけである。
実際、定理\ref{thm:septic-bielliptic-forms}の二つの倍率不変量が有限交換を排除し、
残る倍率は共通倍率と $h$ の符号だけである。標準四移動は両偏極を保ち二点を個別に保つので、
一般点で表示の対合に一致する。この一致は全滑らかな係数に及ぶ。
実際、一定のFrobenius次元を持つこの開集合上でtraceを局所的に取り、
完全積対合の逆行列から四移動の射影作用が正則に定まる。traceの取り替えは射影作用を変えない。
滑らかな完全交叉族は平坦で全空間に一般ファイバーが稠密であり、二射の一致は閉条件である。
両点は $h=0$ の分岐点なので $K=f^*\mathcal O(P-Q)$、$f^*$ の単射性と前系から因子の式を得る。
\end{proof}

% !TeX root = ../papers/nc-p3.tex
\section{非超楕円な種数三の滑らかな六次点グラフの排除}
\label{sec:nonhyperelliptic-sextic-exclusion}
\begin{theorem}\label{thm:nonhyperelliptic-sextic-exclusion}
有限Frobenius候補で一辺の点グラフが滑らかな非超楕円曲線 $C$（種数3、次数6）そのもので、
四辺の各射影ファイバーが空か一点となるものは存在しない。
\end{theorem}
\begin{proof}[点曲線の三次式]
各位置で次数6の偏極 $L$ は $h^0=4$ なので完全線型系となる。
$T=L/\omega_C$ は次数2で非有効である。有効ならその長さ2の因子が $L$ の切断に
一条件しか課さず、閉埋込みに反する。$\omega_C/T$ もRiemann--Rochで非有効である。
この空間曲線は二次曲面に入らない。滑らかな二次曲面なら二次数は $(2,4)$ または逆順となり
次数2の鉛筆ができる。二次錐の解消での類は $as+6f$、随伴種数は $-a^2+6a-5$。
種数3なら $a=2,4$ のうち負切断との交叉 $6-2a\ge0$ により $a=2$ となり同様である。

反対位置を変数にした前節の行列式三次式 $S_p$ が各 $p$ で非零であることを示す。
全関係空間が可分解元を含まないのは、非退化な整曲線上の非零切断の積が非零であるためである。
従って三次縮約の中支持は3以上（前述の特異な二元行列束の議論）である。
中支持4では定理\ref{thm:finite-cubic-separators}のUFD上の証明をそのまま使える。
中支持3で行列式が恒等零なら、定理\ref{thm:ternary-point-subalgebras}による
三文字の点グラフが平面であり、ある辺で共通三文字の可換関係を含む。
その三文字の零点を両側で $a,b\in C$ とする。両側の三座標の比が共通だから、
基点を除いた偏極 $L_i(-a)=L_{i+1}(-b)=N$ は同型である。
$a=b$ なら元の二偏極も一致し、曲線を含む二次曲面がないため六関係は交換子全体となって矛盾する。
$a\ne b$ なら次数7の $N(a+b)$ は非常に豊富で、五次元の全切断を持つ。
共通三文字とそれぞれの補文字 $x,y$ をこの標準でない射影四空間の座標とする。
残る三関係はそこで三つの二次式となる。共通三文字が非零の開部分ではこの三式の解が
そのまま二辺座標の対に復元されるため、零点は与えた曲線に限る。
外の割線に $xy$ の係数を持つ式がなければ $(j_i(a),j_{i+1}(b))$ が点対応に入ってしまう。
従って割線にも二指定点しかない。三式は完全交叉で次数8なのに、
台が次数7の整曲線だけとなり矛盾する。

反対位置の $L=L_2$ に戻る。$I_C$ は四つの三次式、三つの一次係数syzygyを持つ。
実際、二次制限は次元10の同型である。$L$ の基点自由鉛筆の完全列と
$n\ge3$ での $H^1(L^{n-2})=0$ から全高次の制限も全射になる。
従って曲線は算術的Cohen--Macaulayであり、$I_C$ は3正則である
（$H^1(I_C(2))=H^2(I_C(1))=H^3(I_C)=0$）。
Hilbert--BurchとHilbert関数により、四つの生成元は線型な $3\times4$ 行列 $A(z)$ の最大小行列である。
$A$ の曲線上の階数は常に2である。階数1以下では全小行列の微分が零となって滑らかさに反する。
\end{proof}

\begin{proof}[各長さ二の階数低下因子が三割線に乗ること]
定理\ref{thm:opposite-curve-divisor}の $D_p$ は長さ2であり、そのスキーム全体が $S_p$ の特異点に入る。
実際、三次行列式の微分は二次小行列の一次結合で、それら全てが $D_p$ 上で零である。
四つの符号付き最大小行列の列を $m$ として、$S_p=\ell m$、$\ell\ne0$ と書く。
曲線上で $dm$ の像は $\ker A$ に等しい。$Am=0$ を微分すれば包含が従い、
両階数が2であることは滑らかな曲線のイデアルを四式が生成することによる。
従って長さ2の $D_p$ 上では階数2の束 $\ker A$ が定数三平面 $W=\ker\ell$ に入る。
$W\to H^0(D_p,(W\otimes\mathcal O)/\ker A)$ の像は高々二次元である。
非零な定数核 $y$ を取れば、一次式の列 $A(z)y$ は $D_p$ とその線型包の直線で零になる。
この直線では $A$ の四つの最大小行列が一本の三次式の倍になり、その式は恒等零ではない
（$C$ は直線を含まない）。よって直線と $C$ の交叉は $D_p$ を含む長さ3の因子である。

$p$ を動かすとこの直線は正則に動く。長さ2の因子への切断評価が非常に豊富性で全射だからである。
曲線との交叉は有限で全ファイバー長が3、滑らかな基曲線上で平坦となる。
曲面 $C\times C$ 上のこの族は純な因子で、そこから $D$ を除いた残りは次数1のグラフ $r(p)$ である。
長さ3の交叉因子 $D_p+r(p)$ に対して $L(-D_p-r(p))$ は直線を含む平面から二切断を持つ。
Riemann--Rochにより
\[
 D_p+r(p)\ \sim\ T+s(p)
\]
となる一意な点 $s(p)$ がある。これは次数1のAbel埋込みを逆に読んだ射である。
どの点 $z$ でも $T+z$（加法的束記法）はただ一つの有効因子を持つ。
実際、$h^0(\omega_C/T)=0$ である。従って $s$ が定数なら $D$ の第二射影は有限台となり不可能である。
$r$ が定数 $r_0$ なら、全 $s(p)$ について $T+s(p)-r_0$ が有効である。
次数2のSerre双対により $s(p)$ は $(\omega_C/T)+r_0$ の唯一の有効因子に入ってしまう。
よって $r,s$ はともに非定数で、Hurwitzにより自己同型である。

もし $r(p)=s(p)$ なら $D_p\sim T$ となって非有効性に反する。
従って $g=rs^{-1}$ は固定点を持たない。$\tau=\sigma^{-1}$ として前節の因子の類を使い、
$p$ を動かすとJacobianに作用する準同型に次の恒等式を得る。
\[
 (1-g_*)s_* =1+\tau_* .
\]
\end{proof}

\begin{proof}[Fermat四次曲線への限定]
有限個のデータを複素数へ埋めて議論できる。標準平面上で固定点を持たない $g$ は、
相異なる三固有方向を持つ。重複なら固定直線が曲線と交わるからである。
固有座標 $u,v,w$ の純四乗の係数はいずれも非零であり、従って三固有値の四乗が等しい。
共通倍率と順序により $g$ は $(1,\mathrm i,-\mathrm i)$ に規格化できる。
これはJacobianの接空間上の実際の作用でもある。一般に平面四次式 $f$ に
$f(Mz)=c f(z)$ となる行列 $M$ が作用すると、標準微分の留数表示から接空間上の行列は
$(\det M/c)M$ である。ここで $c=\det M=1$ である。
同じ事実により、実際の接空間行列 $N$ については $f(Nz)=\det(N) f(z)$ となる。

有限自己同型群の作用を不変Hermite内積でユニタリとし、上のJacobian等式に
その随伴を右から掛けると $\tau+\tau^{-1}=-(g+g^{-1})$ である（接空間上の記号）。
従って $ku$ では $\tau=-1$、残る二平面では固有値が $\pm\mathrm i$ である。
両方の符号が出ると $\det\tau=-1$ となり、$u^4$ の非零係数に反する。
ゆえに残る二平面では一つのスカラー $\eta=\pm\mathrm i$ である。
$g$ 不変の四次単項式は $u^4,v^4,w^4,u^2vw,v^2w^2$ であり、
$\tau$ の行列式は1なので $u^2vw$ は現れない。
さらに $(1-g)s=1+\tau$ から、$s$ の $v,w$ の行は対角で、
対角成分の対は $(\mathrm i,1)$ または $(1,-\mathrm i)$ である。
ユニタリ性で $u$ 方向も保つ。純四乗の係数は保たれるが $v^2w^2$ の符号は変わるのでこれも消える。
従って $C$ はFermat平面四次曲線で、$\sigma$ は一座標を原始四乗根倍する射影変換である。
\end{proof}

\begin{proof}[偏極の全列挙と三次重なりによる排除]
$K=L_{i+2}L_i/L_{i+1}^2$ とする。$\sigma$ の四固定点は $u=0$ の標準超平面切断 $B_0,\ldots,B_3$ である。
$K\ne1$ なら $D$ は対角を含まない。含めば残りが負の自己交叉のグラフと数値同値なので
そのグラフに等しく、$K=1$ となる。$D\cdot\Delta=0$ ゆえ両者は交わらない。
対角への制限と $\mathcal O(\Delta)|_\Delta=\omega_C^{-1}$、
$\mathcal O(\Gamma_{\sigma^{-1}})|_\Delta=\mathcal O(\sum B_j)=\omega_C$ から $K^2=1$ を得る。
$K=1$ でも同じである。

$\sigma$ に不変な $\operatorname{Pic}^0$ の群は $(\mathbb Z/4)^2$ である。
実際、接空間の固有値から $1-\sigma^*$ の核の位数は
$|2(1-\eta)^2|^2=16$。各 $4B_j$ は $\omega_C$ と同値であり、
$A=B_1-B_0,B=B_2-B_0$ はこの核で独立な四捩れである。独立性が破れれば
$2A,2B,2(A-B)$ のいずれかが零となって次数2の鉛筆ができる。
$\sum B_j\sim4B_0$ より $B_3-B_0=-A-B$ である。
一階差 $d_i=L_{i+1}/L_i$ は $\sigma^*d_i=d_i K^4=d_i$ を満たすので $d_i^4=1$、
従って $\sigma^*L_i=L_i d_i^4 K^6=L_i$ となる。

不変な次数2の束を $2B_0+mA+nB$ と書く。不変かつ有効なら、
有効因子が唯一であり、その台は四固定点に限られる（それ以外では軌道の長さが4）。
従って有効な十個は $B_j+B_k$、残る非有効なものの $(m,n)$ は
\[
 (1,2),(1,3),(2,1),(2,3),(3,1),(3,2).
\]
各 $L_i/\omega_C$ はこのリストの点であり、位置 $i$ でのラベルは
$e+i d+\binom{i}{2}k$、$2k=0$ となる。隣接ラベルが等しいと二次曲面の非存在で可換となる。
これを除くと順序付き四列は30個である。

\path{computations/verify_nonhyperelliptic_sextic.py}で全30列の三次重なりを厳密に除外する。
モデル $u^4=t^4-1$ とし、$B_j$ の基点を $1,\mathrm i,-1,-\mathrm i$ の順に取る。
$L=\sum a_j B_j$ に $a=(6-m-n,m,n,0)$ を使う。全切断の基底は
\[
 \frac{u^e t^h}{\prod_j(t-t_j)^{b_j}},\qquad
 b_j=\lfloor(a_j+e)/4\rfloor,\quad 0\le e<4,\quad 0\le h\le\sum b_j-e .
\]
分岐点での位数と無限遠での極の上限、巡回群の固有分解から従う。
各隣接対の16積を共通分母にすると、標数101、$\mathrm i=10$ で階数は10である。
従って標数零でも全六関係が積の核であり、良い素点での特殊化は計算した六平面になる。
三次交叉の次元は特殊化で下がらない。各始点で次元0の四列が6個、次元1が16個、次元2が8個であり、
いずれも標準有限Frobeniusに必要な次元4に達しない。以上で全候補を除外した。
\end{proof}

% !TeX root = ../papers/nc-p3.tex
\section{滑らかな六次曲線と孤立残余の全排除}
\label{sec:smooth-sextic-exclusion}
\begin{theorem}\label{thm:smooth-sextic-exclusion}
標準有限Frobenius候補で四辺の左右の射影ファイバーが空か一点なら、
点グラフが滑らかな整六次曲線 $C$ とそれに交わらない零次元スキームの直和になることはない。
\end{theorem}
\begin{proof}[種数の限定]
各位置の空間次数は6で一定である。平面六次曲線なら種数は10であり、
非退化な空間六次曲線のCastelnuovo上界4により全位置で平面に入る。
両平面の積の中で全関係の基点は次元1以下なので、一般の三線型結合は三本の
双一次式による完全交叉曲線を定める。その両側次数は3であり、六次曲線を含めない。
従って全位置で非退化である。系\ref{cor:point-graph-excess}で孤立残余の長さは $2g-6$、
Castelnuovoにより $g\le4$ である。
$g=4$ なら $L_i$ の四つの独立切断とRiemann--Rochにより次数零の
$\omega_C/L_i$ が自明であり、全位置が同じ標準埋込みになる。
$h^0(\omega_C^2)=9$ なのでこれを含む二次曲面がある。
独立な二式は存在しない。完全交叉なら次数4に過ぎず、共通平面因子なら
整な非退化曲線が平面または直線に入ってしまう。
従って曲線上の双一次関係は全てこの二次曲面の対角上でも零であり、点グラフの次元に反する。
残る $g=3$ の場合は孤立部の長さが零で、非超楕円の場合は前節で排除した。
\end{proof}

\begin{proof}[超楕円な場合の分裂型]
二重被覆を $f:C\to\mathbb P^1$、$J=f^*\mathcal O(1)$、
$W=H^0(\mathbb P^1,\mathcal O(1))$ とする。$L_i$ は次数6の完全な四切断の偏極で、
$\omega_C=J^2$、$L_i/J^2$ は次数2で切断を持たない（持てば二点を分離できない）。
$f_*L_i$ は階数2、次数2であり、$f_*L_i(-2)$（直線上でのひねり）は切断を持たない。
従って分裂型は $\mathcal O(1)^2$ である。
$P_i=L_i/J$ は基点自由な次数4の鉛筆で、$U_i=H^0(P_i)$ とすると
$f_*P_i=U_i\otimes\mathcal O$、$V_i=W\otimes U_i$ となる。
各埋込みは二鉛筆の積で滑らかなSegre二次曲面に入る。
もし $P_i=P_{i+1}$ なら二次曲面の対角も点グラフに入り矛盾する。
ここでも六次の非退化整曲線を含む二次曲面は一つだけである。

各二文字で積の束写像の核を $K_i$ と書く。
\[
 0\longrightarrow K_i\longrightarrow
 U_i\otimes U_{i+1}\otimes\mathcal O\longrightarrow f_*(P_i P_{i+1})\longrightarrow0 .
\]
両束 $P_i$ は有限平坦二重被覆の各ファイバー上でその全切断空間を $U_i$ から受け取るため、
局所Artin環上の可逆加群を自明化した通常の乗法により最後の射は全射である。
従って $K_i$ は階数2、次数 $-4$ であり、$H^0(K_i)=0$ である。
最後の消滅は $P_i$ の鉛筆列を $P_{i+1}$ でひねたときの次数零の核束
$P_{i+1}/P_i$ が非自明なためである。
直和分解の可能性は $(-1,-3),(-2,-2)$ だけであり、いずれも $H^1(K_i(2))=0$。
よって $V_i\otimes V_{i+1}$ から曲線の全二次切断への積は全射、六関係は全ての核である。
このうち四本は $\bigwedge^2 W\otimes U_i\otimes U_{i+1}$、残り二本は
$H^0(K_i(2))$ で与えられる。

分裂型 $(-1,-3)$ なら最後の二式は各基点 $z\in\mathbb P^1$ で一形式の倍にしかならない。
両側の第一鉛筆を $z$ に固定した $\mathbb P(U_i^*)\times\mathbb P(U_{i+1}^*)$ では
四つの交代式は全て零で、残りが正次元の零点を持つ。
実際の点グラフはこの基点上で有限なので矛盾する。従って全て $K_i=\mathcal O(-2)^2$ である。

三文字の全積核の束を調べる。二つの部分束 $K_i\otimes U_{i+2}$ と $U_i\otimes K_{i+1}$ を
階数8の自明束に入れ、その交叉を $B_i$ とする。和の余核は $f_*(P_iP_{i+1}P_{i+2})$ である。
実際、長さ2の各ファイバーの二次元環 $D$ で、三つの可逆加群を自明化すると
$D^{\otimes3}$ を隣接二箇所の乗法核で割ったものは $D$ である。
単位を挿入することでどの三テンソルも全積と二つの単位に等しくなるためである。
従って和は定数階数、$B_i$ は階数2の部分束で、
$\deg B_i=-8-8+(12-4)=-8$ となる。
$B_i(3)$ の次数は $-2$ で、$K_i(3)\otimes U_{i+2}=\mathcal O(1)^4$ に入るため、
二つの直和因子の次数は高々1である。よって $h^0(B_i(3))\le2$。

四つの交代式を各辺で割った三文字空間は
$\operatorname{Sym}^3 W\otimes U_i\otimes U_{i+1}\otimes U_{i+2}$ で次元32である。
$W$ の隣り合う二文字を自由に交換し、三つの $U$ の係数は独立に保持するためである。
残る左と右の二次関係の挿入はそれぞれ $H^0(K_i(3))\otimes U_{i+2}$,
$U_i\otimes H^0(K_{i+1}(3))$ の八次元空間を与える。
$K_i(2)$ は自明束なのでここで $W$ との積は全切断を張る。
両者の交叉は $H^0(B_i(3))$ である。従って三文字商の次元は高々
$32-8-8+2=18$、有限Frobenius候補に必要な20に反する。
\end{proof}

% !TeX root = ../papers/nc-p3.tex
\section{種数二の滑らかな五次点グラフの偏極}
\label{sec:genus-two-quintic-ruling}
\begin{theorem}\label{thm:genus-two-quintic-ruling}
四辺の左右射影ファイバーが空か一点となる標準有限Frobenius候補で、
点グラフの一次元部分が滑らかな整五次曲線 $C$ で、孤立残余と交わらないとする。
種数は1または2である。種数2では孤立残余の長さは2で、標準四移動は超楕円対合となる。
$H=\omega_C$ とし、添字を固定すると、偏極は二つの相異なるWeierstrass点 $a,b$ に対する
$H^2(a),H^2(b)$ を交互に繰り返す。各埋込みを含む二次曲面は錐である。
\end{theorem}
\begin{proof}
平面五次なら種数6であり全位置で平面となるが、前節の平面積の完全交叉次数3に反する。
非退化ならCastelnuovo上界が2、系\ref{cor:point-graph-excess}の孤立残余長は $2g-2$ である。
以下 $g=2$ とする。次数5の $L_i$ は完全な四切断を持ち、$h^0(L_i^2)=9$ だから曲線は二次曲面 $Q_i$ に入る。
非退化整曲線の次数が4を超えるのでこの曲面は一意かつ既約である。
$L_i=L_{i+1}$ なら全関係がその二次曲面の対角で零となるため、隣接偏極は相異なる。

$Q_2$ が滑らかなら曲線の型は $(3,2)$ または逆で、二つの射影の次数は2と3である。
錐ならその解消 $\mathbb F_2$ の負切断を $s$、ファイバーを $f$ とする。
曲線の固有変換は元の滑らかな曲線と同型である（頂点イデアルの曲線への制限は可逆）。
その類は $ms+5f$、種数は $-m^2+5m-4$ で $5-2m\ge0$ なので $m=2$ となる。
$s$ は一点 $t$ で切り、超平面 $s+2f$ の制限は $L_2=H^2(t)$ である。
どちらの曲面でも次数2の鉛筆は種数2曲線唯一の超楕円鉛筆である。

一般の $p\in C$ を固定する。反対二点縮約を三行三列に書いた行列の三次行列式 $S_p$ は $C$ で零になる。
その $Q_2$ への制限が非零であることを示す。恒等零なら曲面の一般点で行列の階数は2である
（曲線の一般点ですでに2）。二次小行列からなる随伴行列の各成分を
滑らかな二次曲面または錐の解消のCox環に書く。この環は多項式環である。
階数1ゆえ、共通因子を除いた列と行の外積と一つの同次共通因子に分解できる。
実際、一列の最大公約元を除き、他の列をその有理函数倍に書くと分母は単元となる。
得た列と行は各々同じ類の三切断からなり、それぞれ固定因子を持たない。

これら二写像は $C$ 上で、$j_1(C),j_3(C)$ をそれぞれ $j_1(p),j_3(\sigma^{-1}p)$ から射影したものを与える。
核直線は一般点で一意であり、これは定理\ref{thm:opposite-curve-divisor}の左右核の記述による。
両射影は平面を張り、一般の $p$ では像は円錐曲線でもない。実際、射影束は次数4で、
円錐曲線なら唯一の二重被覆から $H^2$ を引き戻した系、すなわち $L_1(-p)=H^2$
（他方も同様）である。これは高々一点に限られる。

滑らかな二次曲面では、片方の次数が零の切断列は一つの射影を経由する。
その次数は小行列の二次数 $(2,2)$ により高々2で、像が平面を張れば円錐曲線になる。
従って列と行は共に両次数が正で、双方の型は $(1,1)$ である。
錐の解消でも、固定因子のない系の類 $u s+v f$ は $v\ge2u$ を満たす。
$u=0$ の場合は射が超楕円射影を経由し、その像の次数は元の射影束の次数4から高々2である。
両方で $u\ge1$ なので、小行列の類 $2s+4f$ から双方とも $s+2f$ になる
（共通因子の有効類の二係数も非負）。
従っていずれの曲面でも各射影は三つの超平面切断で与えられる。
特に曲線への制限で次数5から4へ下がり、ある点 $x$ で $L_2(-x)=L_1(-p)$ である。
これは一般の $p$ に成り立たない。実際 $L_2/L_1$ は非自明で、この束を $x-p$ と書く方法は有限である。
一つの表示 $x_0-p_0$ と比較すれば $x+p_0\sim x_0+p$、
二因子が異なるならどちらも超楕円ファイバーである。この矛盾で $S_p|_{Q_2}\ne0$ を得る。

曲面上でこの三次切断から $C$ を除くと、滑らかな場合は次数3で $C$ を切る一本の母線、
錐の解消では $s+f$ の有効因子が残る。後者は $s$ と一本のファイバーの和である。
どちらの場合も残余の $C$ への制限 $E_p$ は次数3で、束は $P_2=L_2/H$ となる。
$D_p$ を前出の定理の長さ2の階数低下因子とする。$S_p$ の微分は二次小行列の線型結合なので
$D_p$ で零である。一般の $p$ では $D_p$ は錐の頂点を避ける（両射影は有限）。
滑らかな曲面で局所的に $S_p=c e$、$c=0$ が滑らかな $C$ を切るため、
制限微分は $e\,dc$ であり $D_p\subset E_p$ がスキームとして従う。
従って残余は一点 $r(p)$ であり、束として
\[
 \mathcal O(r(p))=P_2\otimes K\otimes\mathcal O(-p-\tau p),\qquad \tau=\sigma^{-1}.
\]
右辺は全ての $p$ で次数1のAbel像に入る（その像は閉）。ゆえに $r:C\to C$ は射である。

$r$ が非定数ならHurwitzにより自己同型で、Jacobian上で $r_*=-1-\tau_*$ となる。
不変Hermite内積で両者がユニタリなので $\tau_*+\tau_*^{-1}=-1$、従って $r_*=\tau_*^2$ である
（有限個の係数を複素数に埋めればよい）。自己同型のJacobian作用は忠実なので
$\tau$ は位数3で $r=\tau^2$ となる。どの位置でも同じJacobian恒等式を適用でき、
右辺の準同型が非零なので残余写像は全て同じ $\tau^2$ になる。
上式を各位置で比較すると全ての $P_i$ が一致する。これは隣接偏極が異なることに反する。

ゆえに $r$ は定数で $\tau_*=-1$、$\tau$ は超楕円対合である。
系\ref{cor:curve-divisor-monodromy}により $K=1$。上の式より $P_i=H(t_i)$、
$L_i=H^2(t_i)$ が全位置で従う。滑らかな二次曲面の場合に $P_i$ は基点自由な第三次鉛筆だが、
$H(t_i)$ は $t_i$ を基点に持つので全位置が錐である。
一階差 $d=L_{i+1}/L_i$ は $K=1$ により一定、$\sigma^*d=d$ かつ $\sigma^*d=d^{-1}$ なので $d^2=1$。
従って $L_{i+2}=L_i$ で、$L_{i+4}=\sigma^*L_i=L_i$ より各 $t_i$ はWeierstrass点となる。
\end{proof}

\begin{theorem}\label{thm:genus-two-quintic-exclusion}
前定理の種数2の場合は存在しない。従って滑らかな整五次曲線と交わらない有限残余からなる
点グラフの候補は、楕円五次曲線そのものに限られる。
\end{theorem}
\begin{proof}
超楕円座標 $x$ の零と極を $a,b$ の下に取ると
$y^2=x F(x)$、$F$ は四次で零点は互いに異なり $F(0)\ne0$ となる。
$H=\mathcal O(2b)$ とすると $L_0=L_2=H^2(a)$ の基底は $1,x,x^2,y/x$、
$L_1=L_3=H^2(b)$ の基底は $1,x,x^2,y$ である。これらを添字 $0,1,2,3$ で表す。
第四の文字を奇文字と呼ぶ。偶数位置で $a$ は第四座標点、$b$ は第三座標点である
（$b$ で $x^2$ の極は位数4、$y/x$ の極は位数3）。
標準四移動は $y\mapsto-y$ なので、四番目の辺での位置0の座標同一視は
奇文字だけの符号変更である。偏極は完全で曲線は非退化なので射影同一視はこれで定まる。

四次Frobeniusテンソルを $w$ とする。他の二添字を固定した係数は
隣接関係に属するので、隣接二切断の積で評価して零になる。巡回する最後の辺では上の符号変更を施す。
$y$ を一個含む関数だけを取れば、偶数から奇数位置への辺で
\[
 w_{30}=w_{23}=0,\qquad w_{3,j+1}=-w_{j,3}\quad(j=0,1)
\]
が成り立つ（他の位置を略記する）。逆の型の辺では転置した式を使う。
位置3から0の辺でのみ交換式の負号が正号になる。これらは $y$ の奇部分での
$x^{-1},1,x,x^2$ の独立性だけを用い、$F$ の係数を選ばない。

奇文字が一個か三個ならその四添字の係数は零である。
一個の場合は奇文字を右に一つずつ、他の三文字と12回交換する。
交換不能なら端の零式により係数が零である。全交換が可能なら、他の各文字との交換は
3ステップおきに4回起き、数値添字の増減は辺の偶奇に従い打ち消し合う。
三個の場合は一つだけの通常文字を右へ12回交換し、毎回交互にその値が増減する。
どちらも全添字が元へ戻り、最後の辺を3回、他の辺を計9回通るので符号は負となる。
従って標数零では消滅する。
特に $w_{3,j,2,k}=0$ が全 $j,k$ で成り立つ。$j,k$ が通常文字なら奇文字は一個、
$j=3$ なら位置1から2で転置した端の零式、$k=3$ なら位置2から3で $w_{23}=0$ を使う。
これらの係数消滅は\path{computations/verify_genus_two_quintic_ruling.py}でも有理数上の行空間で検算する。

よって $j_0(a)$ での三次縮約は、中位置で $j_2(b)$ が消す三平面 $U_2$ に支持される。
左端は $j_1(a)$ が消す三平面 $U_1$ に支持され、右端支持は次元3である。
実際、隣接切断の階数6と左右の点射影ファイバーの一意性により、隣接関係の評価の像は三平面全体である。
従って末尾の縮約は $R_1$ に入る $U_1\otimes U_2$ の三次元部分空間を与える。
$U_1,U_2$ はそれぞれ $L_1(-a),L_2(-b)$ の全切断である。
$2a\sim2b$ より両束は共通の次数4の束 $N$ となる。この平面写像の像は平面を張り、二次曲線に入らない。
そうなら唯一の二重被覆から $N=H^2$ となるが $a\ne b$ に反する。
従って $U_1\otimes U_2$ の全双一次関係はこの同一視のもとで三つの交換子であり、$R_1$ は全てを含む。

$N(a+b)=L_1(b)=L_2(a)$ は次数6の非常に豊富な束で、五つの全切断を持つ。
二つの四切断空間の共通部分は上の三切断であり、和は全切断となる。
対応する射影四空間でこれら三座標を $h$、二つの補座標を $u,v$ と書く。
三つの交換子以外の三関係は三つの二次式になる。
$h\ne0$ では点 $[h:u:v]$ と共通座標が比例かつ非零な二点の組とは一対一に正則に復元される。
したがってこの開集合上で三式の共通零点の一次元の台は、次数6の $C$ に限られる。
$h=0$ の割線では全二次式は $uv$ の倍数で、どれかの係数が非零である。
全て零なら $(j_1(a),j_2(b))$ が点対応に入り、$a\ne b$ かつ $j_1(a)$ の次点の一意性に反する。
従って割線も成分ではない。三式は余次元3の完全交叉を定め、次数8である。
その一次元成分の台は次数6の曲線一つなので次数は6の倍数となるはずで矛盾する。
\end{proof}

% !TeX root = ../papers/nc-p3.tex
\section{楕円五次点グラフの三割線と射影中心}
\label{sec:elliptic-quintic-trisecants}
以下、四辺で左右の射影ファイバーが空か一点の標準有限Frobenius候補で、
点グラフが滑らかな楕円五次曲線 $E$ の場合を扱う。$j_i:E\to\mathbb P(V_i^*)$、
$V_i\subset H^0(L_i)$ は非常に豊富な四次元部分系である。
完全な五切断系からの射影中心 $o_i$ は、接線を含む割線多様体の外にある。

\begin{proposition}\label{prop:quintic-mixed-kernel}
$R_i$ は常に $V_i\otimes V_{i+1}\to H^0(L_iL_{i+1})$ の全核である。
\end{proposition}
\begin{proof}
$L=L_i,M=L_{i+1}$ とする。積の像を消す非零な汎関数 $\ell$ があれば、
完全切断系上の双線型形式 $\ell(st)$ の階数は高々2である。
実際、二つの超平面の積上で零なので行と列を一本ずつ除けば零行列になる。
Serre双対でこの汎関数に対応する拡大を
$0\to\mathcal O_E\to F\to LM\to0$ とする。$L^{-1}$ でひねった接続射は上の双線型形式であり、
$F L^{-1}$ は次数0、階数2で三つ以上の独立な切断を持つ。
その切断の作る部分束が一般点で階数2なら、独立な二切断の外積は次数零で非零、
ゆえにどこでも消えず、この束は自明で切断は二次元しかない。
よって切断は一つの飽和直線束に入り、その次数は少なくとも3である。
元に戻すと $F$ には次数8以上の直線束が入り、$LM$ への像は $LM(-Z)$、$\ell(Z)\le2$
（ここで長さを表す）となる。引戻し拡大は分裂するので、汎関数は $H^0(LM(-Z))$ を消す。
両 $V_i,V_{i+1}$ は $Z$ 上の全切断を与え、その積は $LM|_Z$ の全切断を張るから矛盾する。
従って積は全射で、核次元 $16-10=6$ により主張を得る。
\end{proof}

\begin{lemma}\label{lem:elliptic-quintic-trisecants}
楕円五次埋込み $j:E\hookrightarrow\mathbb P^3$ で、各 $t\in\operatorname{Pic}^3 E$ に
その類の三割線因子 $T_t$ がスキームとして一意に定まる。
$Z=\{(t,q):q\in T_t\}$ の両射影は有限平坦で、次数は順に3と2である。
\end{lemma}
\begin{proof}
$N=t$、$M=L/N$ と書く。非零切断 $s\in H^0(N)$ の因子が直線に入る条件は
$sH^0(M)\subset V$ である。$H^0(N)\to H^0(M)^*$ を射影中心 $o$ と乗法で定めると階数は2である。
階数が落ちれば、ある非零 $m\in H^0(M)$ で $o$ は $mH^0(N)$ を消し、
長さ2の $(m)_0$ の張る割線上に入るからである。従って核はただ一本である。
次数4の因子が空間直線に乗ることもない。$L$ からそれを引けば次数1で、二つの独立平面を含めない。

$E_t\times E_q$ 上で $N$ を正規化Poincaré束に第二因子の次数3の束を掛けて取り、
$M=\operatorname{pr}_q^*L/N$ とする。$N,M$ の第一射影での順像は階数3,2で次数は共に $-1$。
実際、高次順像が零で、Poincaré因子 $\Delta-E\times\{0\}-\{0\}\times E$ の自乗は $-2$、
両座標ファイバーとの交叉は零なので曲面Riemann--Rochで各Euler数が $-1$ になる。
上の全射の核は次数 $-2$ の束 $B$ で、$Z$ の束は $N\otimes\operatorname{pr}_t^* B^{-1}$、
両側のファイバー次数は3と2となる。どの $q$ にも無限個の三割線は通らない。
実際、$j(q)$ からの平面射影は次数4の基点自由系であり、双有理なら長さ2の分離に失敗する
因子は有限である（各特異像の有限逆像で、重複点は曲線上で一意）。
二重被覆なら像は二次曲線で、そのファイバーの長さ2の因子は全て同じ類を持つため、
$q$ を加えた三割線因子は既に示した一意性に反する。従って両射影は有限で、因子の純性から平坦である。
\end{proof}

\begin{lemma}\label{lem:quintic-cubic-polar}
$V^*$ 上に非退化交代形式 $\omega$ があり、$j(E)$ を切る五次元の三次式系を
$P=(\bigwedge^2 V^*)_0$（$\omega$ による縮約の核）の双対と同一視できる。
曲線上の点 $z$ で三次式が特異となる条件は、対応する汎関数が
$F_z=z\wedge z^\perp$ を消すことである（$z$ は同次座標の代表）。
さらに長さ2の因子 $B$ の張る直線が等方的なら、その直線は三割線である。
\end{lemma}
\begin{proof}
原点を選べば $L=\mathcal O(5O)$ とできる。Weierstrass式 $y^2=x^3+A x+B_0$ に対して
完全系の座標を $(a,b,c,d,e)=(1,x,y,x^2,xy)$ とする。交代五行列の初めの二行の右三成分を
$s=(a,b,c),t=(b,d,e)$、左上の二行二列を零、右下の交代三行列を
ベクトル $(-B_0 a,-d-Aa,c)$ の外積行列とする。
最大Pfaffianは符号を除いて $s,t$ の二次小行列三つと、このベクトルとの内積二つである。
これらの共通零点はスキームとして完全五次曲線である。
$a\ne0$ ではWeierstrass式そのものになり、$a=0$ では $[0:0:0:0:1]$ のみが残る。
$e=1$ で $b=cd,a=bc$ を消去すると、残りは $c-d^2-Aad-B_0 ab=0$ の
滑らかな一方程式となる（他の二式はその倍数）。

中心 $o$ での行列の階数は4である。その核を最後の基底ベクトルに取り、行列を
\[
 \begin{pmatrix} uJ+D(z)&z\\-z^t&0\end{pmatrix}
\]
と書く。最後の列の四切断 $z$ は $o$ からの射影の座標を与える。
独立でなければ $o$ を通る直線全体で同じ列が零となり、そこではPfaffianの条件は一本の二次式だけである。
これは $o$ で非零で、長さ2の割線を与えて矛盾する。$J$ は可逆交代行列である。
初めの四つの符号付きPfaffianは $uJ' z+b(z)$、$J'$ は可逆交代行列
（$J$ のPfaffian余因子）、$b$ は二次の列である。Pfaffian列は元の行列の核に入るので $z^tb(z)=0$。
$v=J'^{-1}b$、$\omega(z,w)=z^tJ'w$ とすると $z\wedge v(z)$ は $P$ に値を取る。
この列の成分は射影曲線をスキームとして切る。実際 $z_i\ne0$ の局所で
$v(z)=-uz$ の $u$ を消去でき、最後のPfaffianも行列の核の恒等式により従う。

各等方二平面 $W$ 上で $v$ の $V^*/W=W^*$ への像は $z$ での評価が零の二次共変ベクトルである。
二変数の一次syzygyにより、これは $z$ に垂直な一次共変ベクトルと一次式 $h_W$ の積である。
$\mathbb P(W)$ が $E$ と交わらなければ $h_W$ の零点で $z\wedge v(z)$ は $W$ を表す。
このような $W$ は一般的である（各 $q\in E$ を通る等方直線は一次元族）。
ゆえに三次写像の像は $\mathbb P(P)$ の滑らかな三次元のPlücker二次曲面に稠密で、五成分は独立である。
全三次系の次元も5である。$j(E)$ は二次曲面に入らない。滑らかな曲面と錐の解消の次数・種数式は
前節のとおりで、次数5では種数1にならない。従って二次の制限は次元10の同型で、
$V$ の基点自由鉛筆の列と $H^1(L)=0$ から三次の制限も全射である。
$z\in E$ では $v(z)$ は $z$ の倍数で、微分像は $F_z$ に入る。
上のスキーム生成と滑らかさにより微分像は二次元、従って $F_z$ 全体となる。
この定数階数の同一視は曲線上の束の同一視である。
最後に $W$ が長さ2の $B$ を含む等方平面なら、$h_W$ は $B$ で零となり恒等零である。
そこで $z\wedge v$ は $W$ 上で一本の三次式に限られ、曲線との交叉は長さ3となる。
\end{proof}

\begin{theorem}\label{thm:elliptic-quintic-centers}
楕円五次点グラフの有限Frobenius候補では $j(E)=1728$（ここで $j$ は楕円曲線の不変量）であり、
標準四移動は平行移動である。各位置で、完全五切断系からの射影中心は、$E$ の四次の単数
$\beta=\pm\mathrm i$ のいずれかについて $\ker(2-\beta)$ の固有直線である。
従って各 $L_i$ を固定すれば中心は高々十点に限られる。
反対二点の階数低下因子の各一般ファイバーは、二点の差が $\ker(1+\mathrm i)$ の
非零元となる被約な対である。
\end{theorem}
\begin{proof}[行列式の三次式が一般に非零であること]
位置0の $p$ を固定し、位置2を変数とした三行三列の行列式を $S_p$ とする。
一般の $p$ で恒等零と仮定する。随伴行列は一般階数1である
（定理\ref{thm:opposite-curve-divisor}により曲線上で元の階数が一般に2）。
四変数多項式環でその二次成分を共通因子と列・行の外積に分解する。
両因子は曲線上で隣接埋込みを一点から射影した平面像を与えるので次数が正で、共に一次となる。
特に $j_1$ を $p$ から射影した平面系と $j_2$ の三つの一次式が同じ写像を与える。
後者には基因子が一点 $x$ だけあり、$L_2(-x)=L_1(-p)$ となるので $x=p+d$、
$d\in E$ は固定である。両平面系は基因子を除けば同じ三切断空間である。

$j_1$ の一般の三割線因子 $T_t$ を取る。補題\ref{lem:elliptic-quintic-trisecants}の有限性により、
その各台 $p$ は上の一般性を満たす。また一般の $T_t$ は少なくとも二つの台を持つ。
全て $3p$ なら標数零でその台は $t$ の有理写像となり、次数9の三倍写像に有理切断を与えてしまう。
$p$ からの平面系では $T_t-p$ が一点へ写る。切断の言葉では、この長さ2の因子で零となる
平面系の部分が二次元以上である。従って $T_t-p+(p+d)$ は $j_2$ で三割線因子である。
$p$ によらず同じ因子類なので、一意性によりこの因子は $p$ によらない。
二つの異なる台で比較すると $d=0$ となる。ゆえに $L_1=L_2$ で、
一般の $p$ で両切断超平面のうち $p$ で消える部分も等しい。
この部分は $p$ とともに動いて各超平面を張るので両系も等しい。
曲線を含む二次曲面がないため、全関係は射影三空間の対角上で零となって矛盾する。

\smallskip
\noindent\textit{階数低下から三割線へ。}
$D_p$ は定理\ref{thm:opposite-curve-divisor}の長さ2の因子とする。
行列の二次小行列が $D_p$ 上で消えるので、$S_p$ はこのスキーム上で特異である。
補題\ref{lem:quintic-cubic-polar}により、これは対応する $P$ 上の汎関数が $F_z$ を $D_p$ 全体で消す条件である。
$D_p$ の張る直線が一般に非等方的だとする。このとき条件は四つ独立で、
汎関数は両点を結ぶ二ベクトルの原始部分に比例する
（外積による非退化な対合で $P$ と $P^*$ を同一視する）。
実際、二点が相異なれば $F_z,F_w$ の共通元は $z\wedge w$ の倍に限られ、非等方なので共通元は零。
二重点なら接方向を $h$ として、$\omega(z,h)\ne0$ である。
$F_z$ の微分を $z\wedge V^*$ で割ると $z\wedge u\mapsto h\wedge u\pmod z$ となり、
$u\in z^\perp/\langle z\rangle$ 上で単射である。この場合の汎関数は $(z\wedge h)_0$ である。
両点の場合の直交恒等式から差を取って対角で制限すればこれが従う。

$p\mapsto [S_p]$ が使う束の次数は高々5である。実際、その係数は三行列式の束
$L_0^3/(L_1L_{-1})=L_0/K$ の切断である。
一方 $p\mapsto \langle j_2(D_p)\rangle$ のPlücker束は
$\det\pi_*(\operatorname{pr}_2^*L_2|_D)$ である（$\pi:D\to E$ は第一射影）。
$D^2=2n$、$n=\Delta\cdot\Gamma_{\sigma^{-1}}$ とおくと、この束の次数は
曲面上の因子列のRiemann--Rochで $10-n$ となる。
$n$ は平行移動なら0、他は $\deg(1-\sigma_*)\le4$ である。
二ベクトルの原始部分への射影中心は階数4の二ベクトルで、Plücker像にないため次数は落ちない。
従ってこの写像の次数は6以上となり、$[S_p]$ と一般に等しいことに反する。
ゆえに一般の $D_p$ は等方直線にあり、補題により三割線に乗る。

\smallskip
\noindent\textit{三割線の入射曲線。}
三割線が切る長さ3の因子の残余一点を $r(p)$、その三次因子類を $t(p)$ とすれば
$p\mapsto (t(p),r(p))$ は $Z$ に値を取る射である。
実際、直線が正則に動き、長さ3の有限交叉は一般に平坦で余因子を取れる。得た有理射を滑らかな基曲線上で延長する。
$t$ は非定数（固定した三割線は有限個の点しか含まない）で、$r$ も補題の有限性で非定数である。
よって $Z$ に、正規化が楕円曲線となる成分がある。
層の重複度も数えれば $Z$ の二つの射影次数は3と2で、各成分は両側に支配的である。
一成分だけなら重複度は $\gcd(3,2)=1$。その正規化からの両射を $f,g$ とすると、
$f$ の三つの逆像の $g$ による像の和が三次因子類に等しい。原点を選び準同型部分を取れば
$f_*=3g_*$（三つの逆像は核元だけ異なる）。これは $3=9\deg g$ を与えて矛盾する。
従って二成分で、どちらも被約かつ第二射影で同型となる（有限双有理射の像の基曲線は正規）。
第一射影の次数は1と2である。原点選択による平行移動定数を略すと、第一成分は
$r=\beta t$、第二成分は $t=\gamma r$ となる。ここで $\beta$ は単数、$\deg\gamma=2$。
三点の和が $t$ なので $1=\beta+\widehat\gamma$ である。二点の和の準同型部分が双対同種写像になるのは
$\widehat\gamma\gamma=[2]$ から従う。
$|1-\beta|^2=2$ より $\beta=\pm\mathrm i$、$j(E)=1728$ である。

第二成分の点を $r(p)$ として除くなら、$D_p$ には同じ成分のもう一点と、第一成分の一点が残る。
これらの $p$ に関する次数はそれぞれ $\deg r$ と $2\deg r$ で、和が第二射影次数2になりえない。
従って除くのは第一成分の点のみで、$D_p$ は $\gamma$ の一ファイバーを与え、
その二点の差は $\ker\gamma=\ker(1+\mathrm i)$ の非零元である。
$D\to E$ の第二射影次数から $\deg(t)=2$、また $t_*-r_*=1+\sigma^{-1}_*$ である。
よって $\deg(1+\sigma^{-1}_*)=\deg(1-\beta)\deg(t)=4$ となり、$\sigma$ の準同型部分は恒等である。

\smallskip
\noindent\textit{射影中心。}
$k$ を $\ker\gamma$ の非零元とすると全三割線因子は
$q+(q+k)+(\beta\gamma q+c)$ と書ける。$\gamma=1-\bar\beta$ なので $\beta\gamma=\beta-1$。
$G=\ker(2-\beta)$ は位数5であり、その各元による平行移動は上の三点を同時に平行移動する。
従ってこの三割線族を保つ。$G$ は次数5の $L$ を保つので完全系の射影四空間に作用する。
各因子の完全系での包は平面であり、これらの平面の共通部分は中心 $o$ 一点である。
他にも共通点があれば射影後に全三割線が同じ点を通る。しかし一般の $q$ を通る相異なる三割線は二本あり、
固定点と $j(q)$ を結ぶ直線の一意性に反する。従って $G$ は $o$ を固定する。
巡回群の生成元を五乗が恒等となるよう線型化すると、$L$ は $E/G$ の次数1の束から引き戻される。
五指標の各固有切断空間はその束の平坦ひねりの切断で一次元なので固有直線は五本である。
位置を回せば全ての中心で同じ結論を得る。
\end{proof}

\begin{corollary}\label{cor:quintic-finite-elliptic-parameters}
$E:y^2=x^3-x$ と固定する。添字を保ち、共通のパラメータ平行移動と各位置の線型座標変更を許すと、
楕円五次窓の全候補は有限個の一パラメータ族の次のリストに含まれる。
四切断系を
\[
 j_\epsilon(x,y)=[1+\epsilon x^2:x:y:xy],\qquad \epsilon=\mathrm i,-\mathrm i
\]
として、位置 $i$ の写像は $p\mapsto j_{\epsilon_i}(p+t_i)$ と取れる。
$t_0=0,t_1=t,t_2=2t+d,t_3=3t+3d+e$、ここで $t\in E,\ d\in E[10],\ e\in E[5]$。
四移動 $p\mapsto p+s$ では $t_4=s$ と選べて $s=4t+6d+f$、$f\in E[5]$ である。
各辺の関係は命題\ref{prop:quintic-mixed-kernel}の全核から復元する。
ここでの有限リストは必要条件で、正則性や全てのパラメータ間の同型は別に決める必要がある。
\end{corollary}
\begin{proof}
$L=\mathcal O(5O)$ の完全座標で $o_\epsilon=[-\epsilon:0:0:1:0]$ とおく。
これは割線外である。実際、$a,b,d$ の部分における有限二点の線型結合で $b=0$ とし、
$c=e=0$ とするには、異なる $x$ の場合は両点の $y=0$ が必要である。
$x=0,1,-1$ の相異なる二つの評価の組からは上の $d/a=\epsilon$ は得られない。
同じ $x$ なら $b=0,a\ne0$ より $x=0$ で $d=0$。一方が $O$ の場合も同様に $b=0$ なら $d=0$。
有限点での接線上では、$o_\epsilon$ は曲線上にないので非零な接方向係数を要する。
$e=xy$ を微分し $c=e=0$ を課すと $y\,dx=0$、ゆえに分岐点である。
そこでは $db=dd=0$ なので $b=0$ は $x=0$ を強制する。$O$ の接線では $a=0$ である。

$k=(0,0)$、$[\epsilon](x,y)=(-x,\epsilon y)$ とする。
$j_\epsilon$ で $q,q+k,([\epsilon]-1)q$ は一般に共線である。実際、座標は
\[
 q+k=(-1/x,y/x^2),\quad
 ([\epsilon]-1)q=\left(\frac{\epsilon(x^2-1)}{2x},\frac{(1-\epsilon)y(x^2+1)}{4x^2}\right)
\]
となり、三行四列の座標の全三次小行列が零である。
三点の和の準同型部分は $1+\epsilon$ で次数2なのでこの族は全ての三割線因子を与える。
前定理の共通平面の議論で $o_\epsilon$ は $\ker(2-\epsilon)$ の固有直線である。
完全系の $E[5]$ 平行移動はこの五つの固有直線を置換し、その作用は推移的である。
実際、$E[5]$ が一つの直線を固定すればその直線に恒等となるよう表現を線型化でき、
$L$ が次数25の商を通じて降下することになり次数5に反する。
位数25の群が五点に固定点なく作用するので軌道は一つである。
一般の $L_i$ は $\mathcal O(5O)$ の平行移動であり、これで二つの系と $t_i$ の表示が従う。

各 $L_i$ の次数5のAbel座標は $-5t_i$ である。平行移動の場合の
系\ref{cor:curve-divisor-monodromy}より二階差 $K$ は二捩れで、添字によらない。
従って $d=t_2-2t_1+t_0$ は十捩れ、次の二階差と $d$ との差は五捩れである。
四番目を含む同じ漸化式を用いれば表示の通りとなる。
\end{proof}

\begin{corollary}\label{cor:quintic-same-center-type}
前系の四つの符号は全て同じで、どの隣接差 $t_{i+1}-t_i$ も五捩れではない。
従って四位置の埋込みは同じ位数5の平行移動群を保つ。
\end{corollary}
\begin{proof}
二つの隣接埋込みが同じ長さ3の三割線因子を持つことはない。
両割線の積 $\mathbb P^1\times\mathbb P^1$ 上で六関係はこの長さ3のグラフに消える。
双一次式の制限はそこへ全射である。実際、どちらの二切断も長さ2を分離する。
三相異点なら二つの一次式の積で任意の二点だけを零にできる。
一つの二重点と単点なら二重点で同時に消える二因子の積で単点を分け、
単点で零かつ二重点で単元の因子と他方の二切断で二重点を分ける。
三重点なら単元と一次の一様化元を両側で選び、積の位数0,1,2を得る。
従って制限された六式は高々一次元を張り、割線の積上に余分な曲線を与える。

各 $j_\epsilon(p+t_i)$ の三割線は $q,q+k,(\epsilon-1)q+c_i$ である。
二符号が違えば、その二つの第三点を一致させる方程式は非零な同種写像の一ファイバーで、必ず解がある。
特殊な $q$ でも因子の一致は重複度つきで成り立つので、これは前段に反する。
従って符号は一定である。

隣接差が五捩れなら両偏極は同じ $L$ である。
中心が同じなら二次曲面がないため $\mathbb P^3$ の対角が点グラフに入る。
中心が異なるなら二中心は同じ巡回群の固有点である。
それらを結ぶ直線が完全系の $E$ に交われば位数5の平行移動の全軌道を含むが、
完全五次系では長さ3の因子は平面を張る。従って二つの四切断空間の共通三切断は基点自由である。
六関係は共通三切断の三交換子を含む。残る三式を完全系の射影四空間に移せば、
共通三座標が非零の開部分で零点は曲線のみで、補座標だけの割線も成分にならない。
割線全体が入れば両側の三座標が共に零になる点対を与え、基点自由性と点グラフの仮定に反する。
従ってこの三二次式の完全交叉の次数8が唯一の曲線成分の次数5の倍数となり矛盾する。
\end{proof}

% !TeX root = ../papers/nc-p3.tex
\section{楕円五次窓の二捩れへの限定}
\label{sec:elliptic-quintic-cubic-overlaps}
\begin{theorem}\label{thm:elliptic-quintic-torsion}
前節の仮定で、各添字に独立な線型基底変更と点曲線の同時再パラメータ化を許すと、
楕円五次窓の候補は $j_{\mathrm i},j_{-\mathrm i}$ に対する各一つだけである。
具体的には $a=(1,0), b=(-1,0)$ とし、隣接辺に
\[
 R(a),R(b),R(a),R(b),\ldots,\qquad
 R(v)=\ker\bigl(V\otimes V\to H^0(L\,T_v^*L),\quad f\otimes g\mapsto f\,T_v^*g\bigr)
\]
を置く。ここで $V$ はそれぞれの明示楕円五次系、$T_v$ は平行移動である。
逆順の二捩れ列は同じ添字保持同型類を与える。両符号が実現するなら二類は相異なる。
\end{theorem}
\begin{proof}[三次重なりの指標分解]
複素共役で片方の符号を扱えば足りる。以下 $\epsilon=\mathrm i$、$H=\ker(2-\mathrm i)$ とする。
隣接二辺は $R(t),R(t+d)$、$d\in E[10]$ で、各辺の差は $H$ を法として選べる。
実際、頂点の平行移動に $H$ の元を独立に加えても、埋込みの線型座標変更で吸収できる。
二辺が有限Frobenius窓に属するなら
\[
 (R(t)\otimes V)\cap(V\otimes R(t+d))
\]
は四次元である。その四つの $H$ 指標は全て異なる。実際、全四位置の関係を保つ $H$ の作用は
一次元の四次交叉直線も保つ。この直線を第四因子で縮約して、第四文字の双対（共通指標ひねり）と
三次交叉との同型を得る。$V$ 自身の四指標は前節により相異なる。

\smallskip
\noindent\textit{良い還元における全幾何点の検算。}
曲線と十捩れを標数181、$\mathrm i=19$ に還元する。十捩れは有限エタールで、
全100点がこの素体上有理である。$H$ を法として二階差 $d$ は20通り。
\path{computations/verify_elliptic_quintic_overlaps.py}が行う厳密検算の内容を示す。
指標を $1,2,3,4\pmod5$ とするよう線型化し、その基底での点座標を $z(P)$ とする。
$h$ ブロックの未知係数を三指標の和が $h$ の語に限る（12または13列）。
$q$ は相異なる三つの $H$ 剰余類の代表を動く。左側は最初の二位置を
$z(q),z(q+t)$ で評価し、右側は後ろ二位置を $z(q),z(q+t+d)$ で評価する。
各々残る四つの添字で係数を取り、計24行とする。
これらの行の共通核は当該指標の全三次交叉である。
実際、対の評価で得る曲線切断は、残る一添字を固定すれば一定指標を持つ。
次数10の束は巡回五重商で次数2の束に降下し、一指標の切断が三剰余類で消えれば恒等零である。

一般階数の証明には各ブロックで $t=(6,36)$ において全列階数を持つ一小行列を選ぶ。
この最大次小行列は $t$ の関数であり、各行の極はその行の変数点が $O$ となる一点で高々5次である。
有限な極の $x$ 座標 $x_j$ と使用行数 $m_j$ に対して
$\prod (x-x_j)^{5m_j}$ を掛け、行列式を $A(x)+y B(x)$ に直す。
無限遠の極上限は $10\sum m_j+5m_O\le130$。よってそのRiemann--Roch基底
$x^i,x^i y$ に対し、正則な有理点での全列階数の補間行列で厳密に復元できる。
各ブロックのノルム $N_h=A^2-(x^3-x)B^2$ を作る。清算に用いた極でも、
射影的な小行列が零ならば清算後の値は零である（各行の極上限で既に清算している）。
無限遠は別に直接調べる。

四指標が落ちる点では任意の相異なる $h,j$ に対して $N_hN_j=0$。
全十積の最大公約元を各 $d$ で因数分解すると、全ての因子は一次で、
全ての根の上で曲線の $y$ も素体にある。従ってここでの候補点の直接階数検査は
代数閉包上の全候補を覆う。残る点を $(d,t)$ の $H$ 剰余類で示すと
\[
 (0,0),\qquad (k,a),\qquad(k,b),\qquad k=(0,0).
\]
三点とも、指標0は単射で、他の四指標の核は各一次元である。

\smallskip
\noindent\textit{標数零への移行。}
素点181での離散付値環を必要に応じ有限拡大し、各捩れ点はそのエタール持上げを用いる。
対象の階数条件は $E$ 上の閉条件である。曲線・捩れ・線型系が代数的に定義されているため、
特定の捩れ点以外に標数零の解があるなら、代数閉体 $\overline{\mathbb Q}$ 上にもある。
固有性でその解を付値環上に延長できる。特殊化で上の三剰余類のいずれかに入るので、
各類において捩れ点と同じところへ特殊化する別解を排除すればよい。

上の三点で、評価用 $q$ を極を避けて取り直し、局所変数 $z$ で評価行列を二次まで展開する。
不変微分に双対な $\partial x=2y,\ \partial y=3x^2-1$ による
$1,\partial,\partial^2/2$ の係数は、標数181でもこの位数まで正則局所座標を与える。
単元で消去し、列余階数が1なら、残った一列の各係数の最小消滅次数が最大次小行列の最小次数になる。
五指標での結果は順に
\[
 (0,2,2,2,2),\qquad (0,1,1,1,1),\qquad (0,1,1,1,1).
\]
$H$ による平行移動に関しては埋込みの基底変更で同じ局所次数となる。

$(k,a),(k,b)$ では標数零でも三次交叉が四次元である。
実際 $u=\pm1$ に対する二捩れ移動の座標は
$x'=u+(3u^2-1)/(x-u),\ y'=-(3u^2-1)y/(x-u)^2$ である。
二系の16積を $(x-u)^3$ で清算し $y^2=x^3-x$ を使って核を取ると、二辺の挿入行列の階数は
両順とも $\mathbb Q(\mathrm i)$ 上で44（同コード末尾）となる。
従って四指標の核は捩れ点の持上げで実際に残る。上の単純零の小行列は、
捩れ切断からの局所変数 $z$ で、完備局所環において $z$ と単元の積である。

$(0,0)$ では二辺とも同じ $R(t)$ で、$t=0$ で交換子六本となる。
各非零指標の一次元核は任意の一次変形で少なくとも一次まで持ち上がる。
実際、変形関係を $ab-ba-zP_{ab}+O(z^2)$ と取り、
三文字の交代syzygyを左右からこれで持ち上げた際の一次の差は
$\sum_{\rm cyc}(P_{ab}c-cP_{ab})$ である。この差は可換三次商で零なので
零次関係の左・右挿入で消せる。指標射影で持上げを同じ指標に保てる。
従って当該最大次小行列は捩れ切断で $z^2$ により整除される。
ここでの関係束は乗法が全射で局所自由であり、この議論は付値環上でも行える
（または標数零での定数・一次係数の消滅と付値環の単射を使う）。
特殊ファイバーでちょうど二次の小行列を選べば、完備局所環で $z^2$ と単元の積である。
以上、どの場合も同じ剰余点への他の持上げを許さない。評価点と指標基底は滑らかな曲線、
エタール捩れおよび相異なる五乗根を用いて持ち上げればよい。

隣接差0の剰余類は前節により排除済みである。従って各位置の差は $a$ と $b$ を $H$ を法として
交互に繰り返す。各埋込みの座標変更で剰余類の代表を順に正規化でき、
全辺で表示の列を得る。$[\mathrm i]$ は $a,b$ を交換し、どちらの $j_\epsilon$ にも
$[1+\epsilon x^2:-x:\mathrm i y:-\mathrm i xy]$ と線型に作用するので逆順は同型である。
最後に、三割線入射のうち第一射影次数1の成分が決める単数 $\beta$ は
曲線同型で共役される。楕円曲線の準同型環は可換で、平行移動は準同型部分を変えないため、
$\mathrm i,-\mathrm i$ は区別される。
\end{proof}

% !TeX root = ../papers/nc-p3.tex
\section{滑らかな整五次点グラフの二同型類}
\label{sec:elliptic-quintic-regularity}
\begin{theorem}\label{thm:elliptic-quintic-regularity}
滑らかな整五次曲線とそれに交わらない孤立残余を点グラフとし、四辺の左右射影ファイバーが
空か一点となる標準14641型の添字保持同型類は、前定理の二つで尽きる。
どちらも孤立残余は空で、点グラフは被約な楕円五次曲線である。
両者とも最小の裸周期は1で、標準四移動は点曲線上で恒等である。
\end{theorem}
\begin{proof}[隣接点グラフのスキーム]
$L,M$ を一辺の二偏極、$V,W$ をその切断系とする。隣接差 $a,b$ は五捩れではないので $L\ne M$。
二つの曲線には共通の三割線因子がない。前節の三つ組のうち差が $k$ の部分対は、
重複を含めても一意であり、同じ部分対に対する第三点は $\ker(2-\epsilon)$ 外の平行移動で変わる。

$H^0(LM)^*$ の非零な $\lambda$ が $V\otimes W$ 上で階数1を与えるとする。
完全系での双線型形式は階数3以下で、命題\ref{prop:quintic-mixed-kernel}の拡大を用いると
$F L^{-1}$ は二つ以上の独立切断を持つ。行列式 $M/L$ が非自明な次数零の束なので、
これらの切断は全て一本の飽和直線束に入り、その次数は2以上となる。
従って $\lambda$ は長さ3以下の因子 $Z$ 上の汎関数である。$Z$ を最小に取れば、
その上の乗法対合は非退化となる。実際、核イデアルが非零ならそのイデアルで割ったより短い部分を通る。
長さ2までは両空間の評価が全射で、その制限対合の階数は長さに等しい。
長さ3なら両像は次元2以上であり、一方でも全射なら階数2以上である。
両像が次元2なら共通三割線ができてしまう。よって $Z$ は一点であり、
点対は曲線のグラフ上にある。逆に曲線上の全点は関係を消すので、これで台が決まる。

曲線上の $p$ で $V(-p), W(-p)$ と書いて、それぞれ $p$ で消える三切断を基点を除いて見る。
両者は $L(-p),M(-p)$ の基点自由な超平面である。その積は $H^0(LM(-2p))$ を張る。
実際、この積を消す汎関数の全双線型階数は2以下で、次数4の二偏極について上と同じ拡大を使えば
二つ以上の切断と行列式 $M/L\ne1$ により、汎関数は長さ2以下の部分に支持される。
基点自由性で長さ1は不可能。長さ2で一方の評価が全射なら他方の単元切断との積で全切断を得る。
両方とも全射でなければこの長さ2の因子に $p$ を加えたものが共通三割線となる。
階数1行列でのSegre錐の接条件は、左右の消滅超平面の積上で微分が零となることである。
従ってグラフのアフィン錐の接空間は $H^0(LM(-2p))$ の零化空間で次元2、
射影接空間は次元1である。全閉点が滑らかとなりスキームとして曲線のグラフに等しい。

\smallskip
\noindent\textit{有限Frobenius性と中心切断。}
以下は $\epsilon=\mathrm i$ で示す。共役により逆符号も従う。
$D=\operatorname{diag}(1,-1,\mathrm i,-\mathrm i)$ は曲線上で $[\mathrm i]$ に対応し、
$(D\otimes D)R(a)=R(b)$、$(D^2\otimes D^2)R(a)=R(a)$ である。
位置 $i$ の基底を $D^i$ で変えれば全辺の関係は $U=(1\otimes D)R(a)$ に一致する。
よって通常二次代数 $T(V)/(U)$ の添字展開として裸周期1を得る。

\path{computations/verify_elliptic_quintic_as.py}は $\mathbb Q(\mathrm i)$ 上で次を確認する。
三文字、四文字の隣接関係の交叉次元は4,1。四次テンソルの三つの切断の階数は4,6,4で、
循環置換はこのテンソルを固定する。これは双対の全積対合の完全性を与える。
五次交叉は零である。両端支持が4の四次テンソルに対し、その右に一文字を付けた空間と
左に一文字を付けた空間が非零に交われば、最初の一文字の支持が同時に4と1になってしまう。
これで標準有限Frobenius性が従い、循環恒等式から標準四移動も射影的に恒等である。

通常代数の四次部分で、全一次文字 $x$ に対し $u x=\lambda x u$ となる空間を
$\lambda=1,-1,\mathrm i,-\mathrm i$ の順に求めると、次元は $2,1,1,1$。
計算は単項式を長さ5まで逐次商にし、五次でこの線型方程式を解く。
各解 $u_j$（$1\le j\le5$、基底選択は同コード）に対し、頂点 $m$ で
$\lambda_j^m u_j$ とする族は折り畳み環で中心となる。
これらを $j^g$（$0\le g\le3$）で線型結合した四つを $z_g$ とする。

\smallskip
\noindent\textit{全次数の認証。}
これらの式は181で良い還元を持つ。標数181において $z_g$ による商は四始点全てで
次数0から13まで
\[
 (1,4,10,20,31,40,44,40,31,20,10,4,1,0)
\]
となる。これは同コードが各長さで、直前までの商に最後の二次関係と四つの四次中心関係を挿入して求める。

標数181での通常の道空間と左右線型複体は
\path{computations/verify_elliptic_quintic_koszul.py}で別に計算する。
交代二捩れの元の基底に戻し、位数5の指標で各行列を分ける。始点0について、長さ $n-2$ の道の
末尾に六関係を挿入し、長さ $n-1$ の道と最後の文字の積を商にする。
すべての $n\le16$ で次元は $h_n=\binom{n+3}{3}$。
線型複体の右端二箇所の完全性は一次生成・二次表示から従う。左端の微分については、
四次テンソルの一文字縮約が独立なので、接頭部の道への作用の核は四つの末尾文字の共通零化空間である。
各次数 $n\le16$ の該当する道でこの共通零化空間は零である。
中央の三次交叉から二次交叉への微分は、三次テンソルを一文字と二次関係に分けて、
その一文字を接頭部に掛けた行列である。像階数は $4h_{n-3}-h_{n-4}$
（負次数は零）となる。従ってEuler式により有限範囲で完全である。
始点を一つ進めることは $[\mathrm i]$ による基底変更であり、両隣接差は二捩れゆえ
$R(a),R(b)$ はそれぞれ転置で不変である。逆順の交代窓も同じなので反対側と全始点にこの計算が及ぶ。
還元後の三つの切断の階数と四次交叉の次元も同コードで確認する。

定理\ref{thm:central-window-cutoff}の十分条件が標数181で成立する。その証明は体の標数を用いない。
標数零への適用では、四中心式の存在を特殊化から推測する必要はない。上で標数零の式を先に解いている。
同定理の開性の証明をこの離散付値環上で用いる。具体的には、次数16まで、有限Frobenius枠の一定階数と
線型複体の単元小行列から、次の道空間が自由であることと完全性を次数帰納で持ち上げる。
中心Koszul複体は特殊ファイバーで正の位置が完全であり、その同じ階数小行列により
中心商の必要な次元も次数13まで持ち上がる。よって同定理を標数零で適用でき、
二候補はともに全次数でKoszul左右AS正則である。先の点グラフの照合と二類の非同型性を併せれば主張を得る。
\end{proof}

% !TeX root = ../papers/nc-p3.tex
\section{滑らかな有理四次点グラフの排除}
\label{sec:rational-quartic-exclusion}
\begin{theorem}\label{thm:rational-quartic-exclusion}
標準有限Frobenius候補で四辺の左右射影ファイバーが空か一点なら、
点グラフが滑らかな整有理四次曲線とそれに交わらない有限スキームの和となることはない。
\end{theorem}
\begin{proof}[二項形式と三割線鉛筆]
$W=H^0(\mathbb P^1,\mathcal O(1))$、$S_n=\operatorname{Sym}^n W$ と書く。
曲線の四次埋込みは非退化である（滑らかな平面四次の種数は3）。
各位置の四切断 $V_i\subset S_4$ の定める中心を $o_i\in\mathbb P(S_4^*)$ とする。
基底 $s,t$ に対して $c_l=o_i(s^{4-l}t^l)$、$C_i=(c_{u+v})_{0\le u,v\le2}$ とする。
このHankel行列は可逆である。特異なら $o_i$ はある二次式 $q$ に対して $qS_2$ を消し、
長さ2の因子の張る割線に乗ってしまう。三次鉛筆
\[
 U_i=\{f\in S_3:o_i(fW)=0\}
\]
は二次元である。$S_3\to W^*$ の階数が2に達しなければ $C_i$ の行も従属する。
共通一次因子 $l$ があれば二つの独立二次式が $C_i$ に関して $lW$ に直交してしまうので
この鉛筆は基点自由である。従って $W\otimes U_i\to V_i$ は同型となる
（互いに素な二つの三次式に一次のsyzygyはない）。
$0\ne f\in U_i$ の零因子は空間三割線に入る。実際、$fW$ の二平面がそこを通る。

隣接する鉛筆は $U_i\cap U_j=0$（以下 $j=i+1,k=i+2$）を満たす。
共通元があれば同じ三割線因子があり、系\ref{cor:quintic-same-center-type}の長さ3の制限と同様、
二割線の積には関係式が高々一本しか残らず、余分な曲線が生じる。
この相補性により $V_i\otimes U_j\to S_7$ は同型である。
鉛筆の基底 $u,v$ との積の核は $(-gv,gu)$、$g\in W$ に限られ、
$gU_j\subset V_i$ は $g=0$ を強制する。実際、$U_j\to W^*$、$f\mapsto o_i(f\,{-})$ は同型である。
さらに $W$ を掛けて $V_i\otimes V_j\to S_8$ は全射で、全核 $R'_{ij}$ は七次元となる。
実際の六関係はその中の超平面 $R_{ij}=\ker\lambda$ である。

\smallskip
\noindent\textit{七関係の二次syzygy束。}
$M_i=\ker(V_i\otimes\mathcal O\to\mathcal O(4))$ とする。
$W\otimes U_i$ の二段の評価列で
$0\to\mathcal O(-1)^2\to M_i\to\mathcal O(-2)\to0$ を得る。
拡大は分裂する。点 $z$ で後ろの文字を評価すると
\[
 0\longrightarrow B\longrightarrow R'_{ij}\otimes\mathcal O
   \longrightarrow M_i(4)\longrightarrow0.
\]
実際、$R_{ij}$ の評価だけで階数3である（$j_j(z)$ の前点が一点）。
$B$ は階数4、次数 $-8$ で $H^0(B(1))=0$ である。
最後の消滅は、$sA+tB'=0$、$A,B'\in R'_{ij}$ なら二つのベクトルが $(tP,-sP)$ で
$P\in V_i\otimes U_j$、かつ $P$ の積が零となることから従う。
よって $B=F\otimes\mathcal O(-2)$、$\dim F=4$。
その包含を二次成分の七行四列行列 $C(z)$ と書く。
行空間 $U\subset F^*\otimes S_2$ は七次元である。定数行の従属は
$H^0(M_i(4)^*)$ に入るが、この空間は零である。
$Q(z)=\lambda C(z)\in F^*\otimes S_2$ はどの $z$ でも零にならない。
もし零なら $B_z\subset R_{ij}$ となり、$R_{ij}$ の評価階数が2に落ちる。

三文字で最初に $R'_{ij}\otimes S_4$ を取り、後ろ二切断の積を零にすると
上の列で $H^0(B(4))=F\otimes S_2$ を得る。最後の切断を $V_k$ に限ると
\[
 K'=(R'_{ij}\otimes V_k)\cap(V_i\otimes R'_{jk})
   =\{v\in F\otimes S_2:\ o_k(Cv)=0\}.
\]
$C_k$ が非退化なので七つの行が課す条件は独立、$\dim K'=12-7=5$ である。
その中に実際の四次元三次交叉があり、$\lambda$ で最初の二位置を縮約した像は高々一次元となる。
すなわち $v\mapsto Qv\in V_k$（二次式の積）の $K'$ での像は高々一本の四次式 $h$ が張る。

\smallskip
\noindent\textit{長さ二の因子上の階数条件。}
$S_2$ 上の点評価列を $e_z=(s(z)^2,s(z)t(z),t(z)^2)^t$ とし、
$q_z=C_k^{-1}e_z\in S_2$ と置く（局所で座標を自明化する）。
$h$ が非零ならその零スキーム上で、像が零なら全直線上で、
\[
                       Q(z)\otimes q_z\in U
\]
が成り立つ。これは $K'$ の零化空間が $(1\otimes C_k)U$ であり、
$Qv$ の点評価が $Q(z)\otimes e_z$ による縮約であることによる。
以下この包含を満たす局所基環上で考える。$Q(z)$ の成分の一つは単元なので、
ある原始的な（剰余体で非零な）行 $\mu$ に対し、独立変数 $x$ について
$\mu C(x)=q_z(x)Q(z)$ となる。

$q_z(x)$ の切る相対次数2の因子を $D$ とする。$D$ 上で $\mu$ は $B$ を消すため、
$M_i(4)|_D$ の双対切断 $\eta$ に降りる。$\eta$ の $R'_{ij}$ の各定数関係での値は定数である。
局所基環上で $D$ の代数は基底 $1,u$ を取り、$u$ の係数を取る汎関数は定数を消し
完全な乗法対合を与える（単位二次式で割った一変数環の表示を取れる）。
これを $\eta$ に合成した原始的な汎関数は $R'_{ij}\to H^0(D,M_i(4)|_D)$ の像を消す。
原始性は、$\mu$ が剰余体で非零でありこの対合が完全なことによる。
この余核は
\[
                  S_6/(V_i\cdot P),\qquad
 P=\ker\bigl(S_2\xrightarrow{\,o_j(q_z\,{-})\,}\mathcal O\bigr)
\]
に等しい。実際、$R'_{ij}\to V_i\otimes V_j\to S_8$ の全射列を $D$ 上の評価列に写す。
中央の評価は全射で、その核は $V_i\otimes q_zP$、右端の核は $q_z S_6$。
蛇の補題で $q_z$ を除けば表示を得る。従って $V_i\otimes P\to S_6$ の七次小行列は全て零になる。

この最後の階数低下の方程式を $P\subset S_2$ の変化する射影平面上で求める。
$P$ が基点自由なら、基底 $u,v$ の積の核は $(-gv,gu)$、$g\in S_2$ で
$o_i(gu)=o_i(gv)=0$ を満たす一次元である。従って積の階数は7。
共通一次因子 $l$ があれば $P=lW$ で、積は $lS_5$ 全体となる。
実際 $W V_i$ は次元 $8-\dim U_i=6$ である。ゆえに台は、$P$ の消滅共変ベクトルが
点評価となる滑らかな円錐曲線である。小行列はこの曲線を被約に切ることも確認できる。
$l(p)=0$、$w(p)\ne0$ とし、基底を $lw,l^2$ から $lw,l^2+\delta w^2$ に動かす。
$g\in U_i$ を $g(p)\ne0$ に選べば、零次のsyzygy $(-gl,gw)$ の一次の像は
$gw^3$ で、余核 $S_6/lS_5$ に非零である。階数6の単元小行列で消去すると、
ある七次小行列の横微分が非零となる。全小行列が円錐曲線で零なので、
各点でこの曲線の被約局所方程式を生成する。

\smallskip
\noindent\textit{Hankel行列の恒等式。}
上の共変ベクトルは $C_jq_z$ であるから、その第一成分と第三成分の積から第二成分の二乗を引いた四次式
\[
 b(z)=(C_jC_k^{-1}e_z)_0(C_jC_k^{-1}e_z)_2-(C_jC_k^{-1}e_z)_1^2
\]
は $h$ の零スキーム上で零である（像が零なら恒等零）。
一方、$A=C_j,B_0=C_k$ とし両者の五つのHankel係数を $a_m,b_m$ と書くと
\[
 (\det B_0)o_k(b)=
 (A\operatorname{adj}(B_0)A)_{02}-(A\operatorname{adj}(B_0)A)_{11}
 =
 \det\!\begin{pmatrix}
 a_0&a_1&a_2&a_3\\a_1&a_2&a_3&a_4\\
 b_0&b_1&b_2&b_3\\b_1&b_2&b_3&b_4
 \end{pmatrix}.
\]
第一の等式は $o_k(e_z e_z^t)=B_0$、第二は展開である。
右辺は $U_j\cap U_k=0$ により非零。従って $b(z)$ は非零で $V_k$ に属さない。
これは像が零の場合に直ちに矛盾し、非零の場合にも同次数の $b$ が $h$ の零因子を含むため
$h\in V_k$ のスカラー倍にならねばならず矛盾する。
\end{proof}

\begin{corollary}
点ファイバーを持つ滑らかな整四次曲線と孤立残余の枝では、
非退化な曲線なら楕円四次しか残らない。平面四次なら種数3なので全位置で平面に入り、
平面の積における三つの一般双一次式の完全交叉の片側次数3に反する。
従って全枝で曲線は楕円四次、孤立部の長さは4である。
\end{corollary}

% !TeX root = ../papers/nc-p3.tex
\section{滑らかな平面三次点グラフと正規一次列の同値}
\label{sec:plane-cubic-point-graphs}
\begin{theorem}\label{thm:plane-cubic-point-graphs}
四辺の左右射影ファイバーが空か一点となる有限Frobenius窓で、点グラフが滑らかな平面三次曲線
$E$ とそれに交わらない有限スキームの和になるものは、次の全ての窓で尽きる。
次数3の偏極 $L$ を持つ $E$ と $e\in E\setminus E[3]$ に対し、平行移動 $T_e$ の
三生成Sklyanin商を取り、系\ref{cor:sklyanin-central-classification}の全ての中心一次PBW持上げを使う。
各窓は全次数で標準Koszul左右AS正則、裸周期1である。有限部は長さ8の三二次式の完全交叉で、
非被約でもよい。固定商ごとの添字保持同型類は同系の分裂点と
$\mathbb P(1,1,1,2,2,2)$ の有限群商であり、商を動かす同型は偏極曲線と平行移動の同型で判定する。
\end{theorem}
\begin{proof}[混合関係の両射影が全階数であること]
滑らかな種数1の三次埋込みは各位置で平面に入り、各平面を切る文字を $t_i$ とする。
その上の文字は次数3の偏極 $L_i$ の完全系である。隣接平面の積に制限した六関係の像は三次元。
実際、次元が3未満なら余分な曲面があり、次元は3を超えない。
後者は次数3の二完全系の積が全射であることによる。隣接偏極が等しければ平面対角を含むので等しくない。
異なる場合、前の偏極の基点自由鉛筆と後の偏極の積は鉛筆列と
$H^1(L_{i+1}/L_i)=0$ より全射となる。平面対積での三式の零スキームはこの三次グラフに等しい。
実際、曲面は存在せず、三式の完全交叉の両側次数は3で、既にこの曲線が全次数を占めている。

制限で零となる三関係を、平面上の座標 $X,Y$ を持ち上げて列で書けば
\[
          t_i B_iY-t_{i+1} A_iX+c_i t_it_{i+1}=0
\]
となる。$A_i,B_i$ は三行三列である。$A_i$ の核に射影点 $P$ があると仮定する。
$P$ が平面三次上なら、その次点との対で三混合式の微分の階数が1以下となり、
残る三式の微分を加えても六次元の周囲空間で階数4以下である。曲線はそこで滑らかなため矛盾する。
従って $P$ は平面三次の外にある。そこで混合三式は後ろの座標に関して恒等零、
残る三式は後ろの平面だけで階数3となる。従ってただ一つの次点が存在し、その点は平面の外にある。
点グラフの両射影の整合性から $P$ には前点も存在する。その前点も平面外でなければならない。
前の辺の混合式に代入すると $B_{i-1}P=0$ を得る。

この二つの階数欠損は許されない。$A_i$ の左核により $t_i\otimes v\in R_i$ なる非零な
純関係が生じ、$B_{i-1}$ の左核により $u\otimes t_i\in R_{i-1}$ が生じる。
従って三次交叉は純テンソル $u\otimes t_i\otimes v$ を含む。
しかし有限Frobenius対合により、これもある非零な第四文字評価による四次テンソルの縮約である。
各点評価で端の支持は少なくとも3である。隣接切断の像が全六関係であり、
その評価の階数が2以下ならば対応する射影ファイバーが直線を含むからである。
この矛盾で $A_i$ は可逆、逆向きで $B_i$ も可逆となる。

\smallskip
\noindent\textit{三次の単射性から一周期と全次数へ。}
混合式より $t_i V_{i+1}=V_i t_{i+1}$ は二次商の中で四次元となる。
実際、混合式は純な片側の関係を含まず、制限が零の関係は混合三式で尽きる。
よって $x t_{i+1}=t_i g_i(x)$ を与える可逆な $g_i:V_i\to V_{i+1}$ があり、
$g_i(t_i)=t_{i+1}$ である。両側から生成するイデアルは片側からも生成できる。
$B=A/(t_i)$ の三文字商は従って次元が少なくとも $20-10=10$ である。

$M_i=\ker(H^0(L_i)\otimes\mathcal O_E\to L_i)$ とおく。これは階数2、次数 $-3$ の安定束である。
次数 $-1$ 以上の直線部分束があれば、その自明三次元束への写像の成分は高々一つの切断の倍となる。
次数 $-1$ ならその共通零点で部分束でなく、次数0なら切断の一次独立性に反する。
$B$ の三次交叉は
\[
 H^0(M_i\otimes M_{i+2}\otimes L_{i+1})
 =\operatorname{Hom}(M_i^*,M_{i+2}\otimes L_{i+1})
\]
である。実際、三テンソルの中央の文字を曲線点に沿って評価すれば、
両隣接積が零という条件はちょうど両評価核束への包含になる（中央が完全切断系）。
右辺の二束は同じ傾きの安定束で、非零射は同型、次元は高々1である。
従って $\dim B_{i,i+3}=27-18+\dim H^0(M_i\otimes M_{i+2}\otimes L_{i+1})\le10$。
等号となり、$t$ の左右乗法は二次から三次へ単射である。
非零な商の三次交叉テンソルは両端の支持が3である。各点で上の束写像は同型なので、
その値の端支持は $M_i$ と $M_{i+2}$ のその点の二平面全体となり、これらは点を動かせば三空間を張る。

二文字関係の後ろに $t$ を付けて二回移動すると、今得た単射性で
$(g_i\otimes g_{i+1})R_i=R_{i+1}$ が従う。この一移動で基底を運べば窓は同じ通常関係を繰り返し、
$t$ は中心になる。商 $B$ は三次交叉の全支持と平面三次グラフから
補題\ref{lem:ternary-finite-graph-test}により三次元Koszul左右AS正則である。
商の関係 $r$ の持上げを $r=t\alpha(r)+t^2\beta(r)$ と書く。三次syzygyの二表示を引くと
定理\ref{thm:regular-normal-linear}の記法で $t(h+t b)=0$。
二次から三次への単射性で $h+t b=0$、商に移して $h\in R_B$ となる。
$h$ の持上げを代入し、$tV$ の四次元性を使えば $\alpha(h)=-b,\ \beta(h)=0$ となる。
従って同定理のPBW三条件が成立し、正則性は全次数に及ぶ。

\smallskip
\noindent\textit{商の自己同型と逆構成。}
中心表示の平面三次グラフの移動 $\rho$ は固定点を持たない。
固定点では中心三関係の二つの横方向に関する列が比例し、先と同じ接空間の矛盾が生じる。
楕円曲線で固定点を持たない自己同型は平行移動 $T_e$ であり、隣接偏極が異なるので $3e\ne0$ である。
逆にこの条件の $T_e$ に対し、完全切断系の三つの積関係はスキームとして滑らかな三次グラフを切る。
台については、積の汎関数の階数が1なら、命題\ref{prop:quintic-mixed-kernel}の拡大で
次数零の二階数束に二つ以上の独立切断がある。その行列式は $T_e^*L/L\ne1$ なので
二切断は同じ次数2以上の直線部分束に入り、汎関数は長さ1以下の因子上に支持される。
接空間については、$p$ で消える二平面の積が前後の偏極の積から $2p$ を引いた全切断を張る。
二つの次数2の完全鉛筆は相異なるため、鉛筆列で四次元から四次元への積が同型となるからである。
従ってこのグラフは被約である。

この三関係は通常の三生成Sklyanin形となる。実際、$E[3]$ は両偏極とグラフを保ち、
$H^0(L)$ 上で三指標の対角作用とその三つの巡回置換を与える。
これは位数3部分群での降下で各固有切断が一次元となり、他の三捩れがそれらを推移的に置換することによる。
全群が固有直線を保てば線型化により次数9の商に次数3の束が降下してしまう。
従って関係空間でも一つの対角指標の関係とその二つの巡回置換が基底である。
実際、文字の指標を一つ送れば二文字の指標は二つ送られ、三つが推移的になる。
三式は $A yz+B zy+C x^2$ とその巡回置換であり、通常の循環三次テンソルが全支持のsyzygyを与える。
よってこの商に系\ref{cor:sklyanin-central-classification}を適用できる。

全てのPBW持上げで、中心関係のため $t\ne0$ なら両点は等しく、
不等な平面三次の各点対の近くでは逆に両側で $t=0$ になる。
一方だけで $t=0$ の点対は存在しない。従って対角上の残余は射影三空間中の三つの二次式で、
無限遠では固定点がないため空、全体が有限である（完備な正次元部分は超平面に交わる）。
完全交叉で長さ8となる。これで全持上げの点グラフの仮定も確認した。
\end{proof}

\begin{corollary}
系\ref{cor:sklyanin-central-classification}の係数 $p,q,s$ と対称な一次補正行列 $M$、
定数列 $b$ を使うと、長さ8の孤立部はアフィン三空間における
\[
 \Phi(X)=\tfrac{s}{3}(X_1^3+X_2^3+X_3^3)+(p+q)X_1X_2X_3
           -\tfrac12 X^tMX-b^t X
\]
の臨界スキームである。一般の $b$ で八点は全て単純で、各特殊値の重複はこの勾配イデアルで決まる。
\end{corollary}
\begin{proof}
$t=1$ の対角へ三関係を制限すれば勾配に一致する。定数 $b$ を動かす勾配写像は
固有である。実際、そのグラフを射影三空間とパラメータ空間の積で閉じても
固定点のない最高次三式のため無限遠点がない。全ファイバーは長さ8で、有限かつ支配的、
標数零なので一般点でエタールである。
\end{proof}

% !TeX root = ../papers/nc-p3.tex
\section{空間有理正規三次の六点残余と四移動}
\label{sec:twisted-cubic-residual}
四辺で左右の点射影ファイバーが空か一点となる標準有限Frobenius窓で、
点グラフが空間有理正規三次曲線 $C$ と交わらない零次元部分 $F$ の和となる場合を扱う。
各文字空間は同じ $\mathbb P^1$ の完全な二項三次形式の空間として同一視できる。
\begin{theorem}\label{thm:twisted-cubic-residual}
辺の六関係を二項変数対で表し、対角の一次交代因子で割ると、
二次形式二組の九次元積空間内の六平面 $T$ を得る。
$\Pi=\mathbb P(T^\perp)$ の行列座標を $B=(B_{ij})_{0\le i,j\le2}$ とする。
$F$ はこの平面での次の二式の完全交叉とスキームとして同型である。
\[
 \det B=0,\qquad
 Q=B_{00}B_{22}-(B_{01}+B_{10})(B_{12}+B_{21})
       + B_{11}(B_{02}+B_{20}+B_{11})=0 .
\]
この長さ6の交叉の全ての点で $B$ は階数2である。
\end{theorem}
\begin{proof}
曲線対の基点を一致させた座標で、六関係を対角上で零となる二次数 $(3,3)$ 形式とする。
対角因子による除算後の全九形式の双対写像は、二点 $X,Y\in\mathbb P^3$ に対し
\[
 B_{ij}=X_iY_{j+1}-X_{i+1}Y_j,\qquad 0\le i,j\le2
\]
を与える。その基点は曲線の対角である。実際、二つの三行二列Hankel行列
$H(X)=(X_i,X_{i+1})$ と $H(Y)$ の間に非退化交代行列を挟んだ積が $B$ である。
両方が階数2なら $B$ も階数2で、一方だけが階数2でも零にならない。
両方階数1なら二点は曲線上で、同じ点でだけ零となる。
基点は被約でもある。曲線の点を $X=Y=[1:0:0:0]$ とすると、
一次部は $Y_1-X_1,Y_2,Y_3,-X_2,-X_3$ を含み余次元5の法方向を全て切る。
二項変数の同時基底変更で全点に運べる。

写像の像の閉包は二式の完全交叉である。行列積による行列式消滅と、
$X\wedge Y$ のPlücker関係による $Q$ の消滅は表示に代入しても確かめられる。
対称成分の六座標は $X\wedge Y$ の六座標を可逆線型変換したものである
（$B_{ii}$ は隣接小行列、$B_{ij}+B_{ji}$ で残りを順に得る）。
階数2行列の左右核を二項二次式として射影平面二つに写す。
各核式の零因子は長さ2であり、その因子の張る三次曲線の割線上を $X,Y$ が動くと、
同じ両核を持つ行列の射影三空間内で写像が得られる。
一般の両因子の台 $p_1,p_2$ と $q_1,q_2$ が全て異なれば、割線の座標で
写像はSegre写像の四成分にそれぞれ $[p_i,q_j]$ を掛けるだけである（角括弧は二次元の行列式）。
従ってこのファイバーでは既約な二次曲面をちょうど一回被覆する。

余分な六次元成分が二核の特殊部分に映ることもない。階数1以下の行列は射影次元4である。
階数2では核の各ファイバーの次元は3以下なので、基が余次元2以上なら足りない。
基の因子上で $Q$ が恒等零になるものはない。実際、相異なる単純根で台が交わらなければ既述の二次曲面である。
片方に二重根があっても両因子が交わらないなら、長さ2の各座標環のテンソル積上で
角括弧は単元であり、同じSegre座標の可逆変換となる。ここで $Q$ 自身がその単元倍の式であることは、
二つの割線が交わらないため、$X\wedge Y$ による四交差座標の線型写像が単射で、
Plücker式のその部分空間への制限が通常の二行列式になることからも従う。
単純根が一つだけ一致する場合には、両核で定めた四座標での $Q$ は非零の二座標の積である。
具体的には評価ベクトル $v_2(p_i)v_2(q_j)^t$ による座標を $b_{ij}$ とすると、
$Q$ は $[p_1,q_2][p_2,q_1]b_{11}b_{22}-[p_1,q_1][p_2,q_2]b_{12}b_{21}$ の非零倍である。
比例係数は両対の内部の行列式だけからなる単元である（$Q$ の式に四つの評価ベクトルを代入する）。
従って二式の零点は既約、一般点で被約で、完全交叉の純性からスキームとしてこの像閉包になる。

基点の曲線を吹き上げて写像を延長する。例外因子上の像は、基点での九式の法微分の像を射影化したものである。
$\Pi$ はその像に交わらない。交われば六関係の微分が一つの非零法方向で全て零になり、
滑らかで孤立部と交わらない曲線の接空間が二次元以上になってしまう。
ゆえに $\Pi$ 上の交叉は全て実際の曲線外の点対から来る。その逆像は $F$ なので交叉は有限、次数6である。
これらの点対はどちらの座標も曲線外であるため $B$ は階数2である。
二核に対応する二割線が同じなら、その割線の積上で $B$ の像は高々直線となる。
実際、同じ二根なら角括弧による乗法の階数は2以下である
（二重根では双対数二組の一様化元の差の乗法が階数2）。よって曲線外の点対を通るファイバーが正次元となり矛盾する。
二割線は異なり、$X\ne Y$ である。$B$ が決める $X\wedge Y$ の直線を各割線と交わらせれば $X,Y$ を復元できる。
この直線が割線の一方に等しければ、他方の点が両割線の交点になる。
有理正規三次で異なる二割線が交わるのは因子の共通台の曲線点だけである（長さ4まで線型独立）。
よって交点構成は各点の近くで定数階数の線型代数により正則である。
被約な像閉包のこの開部分でこれはもとの写像の逆であり、平面切断にも同型が残る。
\end{proof}

\begin{theorem}\label{thm:twisted-cubic-monodromy}
この三次曲線上の標準四移動 $\sigma$ は恒等または対合である。
窓テンソルは四つの隣接対角因子を掛けた
二次元文字四つのテンソル $f$（係数16）で与えられる。
対合の場合は二固定点を基底に取ると $f$ は偶数個または奇数個の第二座標を持つ八成分だけに支持される。
\end{theorem}
\begin{proof}
各点写像を $j_i$ とし、$j_{i+4}=j_i\circ\sigma$ の規約で曲線座標をそろえる。
四次テンソル $w$ を四組の二項変数で各次数3の形式として見る。
任意の隣接対をそのグラフの点で評価すると残り二位置の全係数が零になる。
従って $w$ は四つの既約な対角因子の積で割り切れ、商 $f$ は各次数1である。
最後の対角にだけ $\sigma$ または座標同一視の逆向きの $\sigma^{-1}$ が入る。
標準四移動は四次交叉直線を保つので $f$ は四位置同時の $\sigma$ の固有ベクトルとなる。

非自明なユニポテントなら、$\sigma$ のべきの閉包は連結加法群である。
長さ6の非空な有限部の台の各点はこの群で固定されねばならない。
しかし有理正規三次の文字表現にはこの群の射影固定点は曲線上の一点しかない。

半単純の場合は各点の座標 $X_0,\ldots,X_3$ の重みを共通倍率を除いて $0,1,2,3$ とする。
位数が5以上または無限なら、点空間の曲線外の固定点は第二・第三座標点だけとなる。
また $B_{ij}$ の重みは共通倍率を除いて $i+j$ である。重み0と4の純行列は階数1、
重み1と3の純行列はともに $Q=\det=0$ を満たす（それぞれ単一の斜め列だけに支持）。
従って $\Pi$ に重み $\ne2$ の固有ベクトルがあれば平面切断に固定点が生じる。
前定理の復元により、それは曲線外の固定点の対でなければならず、
かつ二つの二次核は異なる。よって第二・第三座標点を互いに結ぶ二通りだけであり、その行列の重みは2である。
従って $\Pi$ は重み2の平面全体に等しい。しかしこの平面の $B_{11}=0$ の直線は二式で消え、有限性に反する。
この議論は窓が対角トーラス全体を保つ場合にも適用できる。

位数4と3を調べる。アフィン変数を $z_0,\ldots,z_3$ とし、順に隣接因子を
$z_1-z_0,z_2-z_1,z_3-z_2,u z_0-z_3$ とする（$u$ は原始根）。
もし $f$ が整数重み一つに支持されれば窓を対角トーラスが保つため既に除外した。
位数4なら従って $f$ は定数項と四変数の積だけを持つ。
隣接切断 $01|23$ の行列は、その両側の対角因子を除いた
$f(z_2-z_1)(u z_0-z_3)$ の係数行列（各側九次元）で階数が決まる。
位置0,1で指数対 $(0,2),(2,0)$ の行はともに零なので、
$\Pi$ に $B_{02},B_{20}$ だけを動かす直線が入り、再び有限性に反する。

位数3なら二固定点を交換して、定数項1と、三次単項式
$\prod_{j\ne m}z_j$ の係数 $A,B,C,D$ だけを仮定できる。
上の九行九列を行の全指数を法3として分けると、三ブロックは列順序を辞書順として
\[
 \begin{pmatrix}0&-1&0\\ C&D-Au&0\\0&(C-B)u&Du\end{pmatrix},\quad
 \begin{pmatrix}1&0&-A\\0&u&-B\\-Cu&-Du&0\end{pmatrix},\quad
 \begin{pmatrix}0&A&0\\-u&B-C&Au-D\\0&0&Bu\end{pmatrix}.
\]
総階数は $4+\#\{A,B,C,D\text{ の非零係数}\}+\mathbf1_{ACu+BD\ne0}$ である。
標準の階数6になるには非零係数はちょうど二つで、それらは反対位置ではない。
すると、係数が零の隣接二位置では、$f$ の全三次項がその二変数の両方を含む。
この二位置を一方にした切断でも、定数項から各変数の指数は高々1、
三次項からは両指数が少なくとも1となる。従って $(0,2),(2,0)$ の行が零で同じ矛盾となる。
これで残るのは恒等または位数2だけで、対合の八成分の記述はその固有指標から従う。
\end{proof}

% !TeX root = ../papers/nc-p3.tex
\section{四移動の固定イデアル、余接層と滑らかな点曲線の全被覆}
\label{sec:point-graph-fixed-differentials}
四辺の左右点ファイバーが空または一点の標準有限Frobenius窓を取る。
点グラフを $X$ とし、埋込みを $j_i$、偏極を $L_i$、
$j_{i+4}=j_i\sigma$ となる四移動を $\sigma$ とする。
定理\ref{thm:point-graph-curve-degree}の一定の二階差を $K=L_{i-1}L_{i+1}/L_i^2$ と書く。
次の式は曲線が全位置で空間を張ることを要求しない。

\begin{theorem}\label{thm:fixed-differential}
一次元の成分で一般点が被約なものの正規化を $\nu:C\to X$ とする。
$M$ を $\nu^*\Omega_X$ の無捩れ商、$T$ をその捩れ部分の局所長の因子とする。
$M=\omega_C(-B)$ で、$B$ は正規化写像の微分の共通零因子である。
\begin{enumerate}
\item $\sigma^{-1}\nu\ne\nu$ なら、二写像の一致スキームは $C$ 上の有効因子 $Z$ で
\[
                 \mathcal O_C(Z)=\nu^*K^2\otimes M(T).
\]
別の成分に移る場合にも、二写像の像が交わる点での一致を数える。
\item $\sigma^{-1}\nu=\nu$ なら
\[
                   \nu^*K^2\otimes M^2(T)=\mathcal O_C .
\]
\item $X$ が被約な局所完全交叉曲線である開閉部分 $Y$ 上では、四移動の固定イデアルは
$\omega_Y\otimes K^2|_Y$ の一切断の零イデアルである。固定成分がなければ固定スキームはCartier因子となる。
\end{enumerate}
\end{theorem}
\begin{proof}
一辺の周囲を $P=\mathbb P^3\times\mathbb P^3$、$H=L_0L_1$ とする。
以下正規化への引戻しを略記する。六関係のJacobi写像
\[
           R_0\otimes H^{-1}\longrightarrow\Omega_P|_C
                          \longrightarrow\nu^*\Omega_X\longrightarrow0
\]
の一般階数は5である。最初の写像の核を $S\otimes H^{-1}$ と書くと、
$S$ は $R_0\otimes\mathcal O_C$ の飽和直線部分束である。
$\det\Omega_P|_C=H^{-4}$ と完全列の行列式から
$S=H^{-1}\otimes M(T)$ を得る。捩れ層の行列式は局所長の因子の束である。
標数零なので $\nu^*\Omega_X\to\Omega_C$ は一般点で同型となり、像は $M=\omega_C(-B)$ である。

四次テンソル $w\in R_0\otimes R_2$ は両六次元空間の完全対合を与える。
後ろ二位置を $j_2(p),j_3(\sigma^{-1}p)$ で評価すると、
$R_0\otimes L_2 L_{-1}$ の切断 $s$ を得る。ここで位置3と位置 $-1$ の文字同定を用いた。
左二位置の第一を $j_0(p)$ に固定すると最後から最初への辺により全ての値が零、
第二を $j_1(p)$ に固定しても位置1から2への辺により全て零である。
よって $s$ は六式の微分の核 $S\otimes L_2L_{-1}$ に入る。
またその六成分の零イデアルは後ろの点対がその辺のグラフに入る条件そのものである。
グラフの両射影が閉埋込みなので、これは $p=\sigma^{-1}p$ の一致イデアルをスキームとして引き戻したものである。
$S$ は局所直和因子であり、$L_{-1}L_2/H=K^2$ は偏極の三階差の式から従う。
切断が一般点で非零なら第一項を得る。

成分上で $\sigma$ が恒等の場合、後ろ二点を一旦独立に $j_2(p),j_3(q)$ とし、
$q=p$ に沿う $q$ の一次微分を取る。零切断の微分なのでこれは自明化によらない
$R_0\otimes L_2 L_{-1}\otimes\omega_C$ の切断である。
第二文字を $j_1(p)$ に固定すれば微分後も零である。
第一文字を $j_0(q)$ に固定した消滅を微分すると、第一文字を微分する項は後ろ二位置の隣接消滅で零となる。
従って再び $S$ に値を持つ。その零因子は正確に $B$ である。
実際、後ろの六式が $C\times C$ に引き戻すイデアルは $\nu(p)=\nu(q)$ のイデアルで、
局所的に $X$ のアフィン座標の二組の差で生成される。
対角イデアルをその二乗で割ると、六式の一次係数はこれら座標の微分と同じ部分加群 $M\subset\omega_C$ を生成する。
これが第二項を与える。

最後に $Y$ 上では $I/I^2$ は階数5の束で、六式が全射
$R_0\otimes H^{-1}\to I/I^2$ を与える。
余法束から周囲の余接束への写像は各成分の一般点で単射であり、$Y$ が被約なので単射である。
従って微分の核はこの全射の核に一致し、$H$ でひねって得る $S_Y$ は局所直和因子である。
随伴公式より $S_Y=H^{-1}\otimes\omega_Y$ で、上の $s$ の零スキームの同定がそのまま使える。
\end{proof}

\begin{corollary}\label{cor:smooth-point-curves-covered}
点グラフが空でない滑らかな射影曲線と、これと交わらない有限スキームとの和なら、
曲線は連結で、正則窓の全候補は次の四行で尽きる。
\[
\begin{array}{c|c|c|l}
\text{次数}&\text{種数}&\text{有限部の長さ}&\text{正則性と添字保持同型の判定}\\ \hline
3&1&8&\text{定理}\ \ref{thm:plane-cubic-point-graphs}\\
4&1&4&\text{定理}\ \ref{thm:elliptic-quartic-as-test},\ \text{系}\ \ref{cor:elliptic-point-finite-atlas}\\
5&1&0&\text{定理}\ \ref{thm:elliptic-quintic-regularity}\\
7&5&0&\text{定理}\ \ref{thm:septic-bielliptic-forms}
\end{array}
\]
特に有理正規三次に対する前節の残余平面と階数条件は同時には実現されない。
\end{corollary}
\begin{proof}
滑らかな開閉成分では $T=B=0$ で、$K$ の次数は零である。
成分を点ごとに保てば第二項から種数1となる。
そうでなければ第一項から固定因子の長さは $2g-2$ であり、ここでも有理曲線は現れない。
各成分の次数は従って少なくとも3で、総次数上限7により二成分があれば次数は $(3,3)$ または $(3,4)$。
三次は種数1、四次の種数は高々3である。系\ref{cor:point-graph-excess}の超過交叉の証明を
各開閉成分に適用すると、三次の寄与は12、残る成分の寄与も少なくとも12となり、総数20を超える。
連結な曲線は既出の各次数の結果で足りる。次数4の平面曲線は隣接平面の積の点グラフの次数上限3で除かれ、
次数6は定理\ref{thm:smooth-sextic-exclusion}で除かれる。
\end{proof}

\begin{corollary}\label{cor:nodal-fixed-count}
被約な節点曲線が点グラフの開閉部分をなすとする。
正規化成分の種数を $g$、その上の節点の枝の本数を $n$ とすると、
四移動がその成分上で恒等なら $n=4-4g$、
そうでなければ正規化上の固定イデアルの長さは $2g-2+n$ である。
\end{corollary}
\begin{proof}
節点 $xy=0$ の $x$ 枝で余接層を引き戻すと自由な $dx$ と $x\,dy=0$ の直和となる。
従って $B=0$、$T$ は全節点の逆像を一回ずつ数える。前定理の次数を取ればよい。
\end{proof}

% !TeX root = ../papers/nc-p3.tex
\section{被約な低次数の点曲線と三角形の全持上げ}
\label{sec:reduced-cubic-point-curves}
前節と同じ点ファイバー条件のもとで、点グラフが純な被約曲線 $Y$ と交わらない有限部 $F$ の和とする。

\begin{theorem}\label{thm:reduced-cubic-point-curves}
$1\le\deg Y\le3$ では、滑らかな平面三次の既出の枝以外に、三角形の次の枝だけがある。
$k^*\setminus\{1\}$ の元 $\lambda$ に対して $q^3=\lambda$ を一つ選び、添字を法3として
\[
              r_i=x_{i+1}x_{i+2}-q x_i^2
\]
を三生成商の関係とする（向きは頂点座標の選択による）。
各 $q$ について中心文字 $t$ を加え、
\[
       r_i=t\sum_j m_{ij}x_j+t^2 b_i,\qquad m_{ij}=m_{ji},\quad m_{ii}=0
\]
とする全六係数が採用される。いずれも裸周期1、全次数で標準Koszul左右AS正則であり、$F$ の長さは8である。
異なる $\lambda$ の間に添字保持同型はない。固定 $\lambda$ の同型類は分裂点と
$\mathbb P(1,1,1,2,2,2)/(\mathbb Z/3)^2$ である。群作用は文字の巡回置換と
$\operatorname{diag}(1,\zeta,\zeta^2)$（$\zeta$ は原始三乗根）の関係への代入で与える。
この枝の全点は滑らかな平面三次の枝からの平坦特殊化である。
\end{theorem}
\begin{proof}[曲線の被覆]
線と滑らかな二次曲線の開閉成分は系\ref{cor:smooth-point-curves-covered}の証明で除かれる。
節点を高々一本しか持たない滑らかな有理成分（正規化上で枝を数える）も
系\ref{cor:nodal-fixed-count}の証明をその成分に適用すれば、次数が足りない。
従って二直線、平面を異にする二次曲線と直線、三直線の鎖などの有理な端を持つ配置は現れない。
整な非平面三次は有理正規三次である（正規化の次数3の束の独立四切断は種数0の完全系を与える）。
三直線が一点で交わり平面を張らない場合の局所式は $xy=xz=yz=0$ であり、
各枝上の余接層は自由部分と長さ2の捩れを持つ。成分上の恒等は定理\ref{thm:fixed-differential}の第二項に反し、
それ以外では第一項の固定長は零だが唯一の三重点は固定される。これで非平面の場合も尽きる。

被約な特異平面三次では $\omega_Y=\mathcal O_Y$ であり、全成分の正規化は有理である。
局所完全交叉ゆえ各正規化上で $M(T)=\nu^*\omega_Y$ である。実際、余法束列を正規化に引き戻すと
余法束の射は一般点で単射であり、自由な始域上で単射となるから、行列式は随伴公式の束を返す。
点ごとに恒等となる成分があれば固定微分の第二式で $\deg M=0$ を得るが $M\subset\omega_{\mathbb P^1}$ で矛盾する。
従って固定スキームは次数零の有効因子、従って空である。唯一の特異点を持つ節点・尖点の整三次、
接する二次曲線と直線、共点の平面三直線は除かれる。横断する二次曲線と直線では次数の異なる各成分を保つが、
射影直線の自己同型は固定点を持つ。残る三角形では四移動は三頂点と三辺を巡回置換する。
三直線が相異なる三節点で会えば全位置で同一平面に入る。

\smallskip
\noindent\textit{平面を消す文字の正規性。}
三角形の頂点を点列に沿って標識し、各辺の両端の文字を対応する座標とする。
隣接平面に制限した双一次関係は三次元以下である。各頂点で対角係数が消え、
二頂点間の各辺ではその二つの交差係数の間に一つの比例式があるためである。
二式以下なら平面の積で曲面が残るので、制限はこの三次元全体である。基底は
$x_ax'_b-c_{ab}x_bx'_a$ ($a<b$, $c_{ab}\ne0$) となる。
左右行列式は三座標の積の非零倍となる（恒等零なら平面全体が点射影像に入る）。
各曲線点で階数2ゆえ、平面の積での零スキームは三角形のグラフである。

平面を消す文字を $t_i$ とする。制限で消える三関係を定理\ref{thm:plane-cubic-point-graphs}の
$A_i,B_i$ の表示で書く。$A_i$ の核に平面内の点 $P$ があるとする。
$P$ が三角形の滑らかな点なら、三混合式の微分は次の点との対で階数1以下となり、
全六式の微分階数が5に達しない。核に頂点があれば、標準四移動が平面も混合関係も保つため
三頂点が全て核に入る。すると $A_i=0$ となり滑らかな点でも同じ矛盾を得る。
ここで核は $t_i=0$ にある前点で全混合式が後ろの任意の文字に対して消える条件なので内在的である。
$P$ が三角形外なら次点と前点は共に平面外、かつ $B_{i-1}P=0$。
前掲定理の純な三次交叉の矛盾によりこれも不可能である。逆向きも同様で両行列は可逆となる。

従って $t$ の両側イデアルは片側から生成され、三文字商の三次元数は $20-10=10$ 以上である。
先ほどの三つの交換関係の三次交叉の次元は高々1である。
実際、両方の隣接対で文字の交換は非零倍率を与え、隣接文字が一致する語の係数は零である。
文字が重複する全語で零、相異なる三文字の六語の係数は一つから全て決まる。
よって次元はちょうど1、三次元数は10、三次交叉は左右で全支持となる。
$t$ の二次から三次への左右単射性を得たので、前掲定理と同じ二回の移動で通常の中心表示に戻る。
補題\ref{lem:ternary-finite-graph-test}と定理\ref{thm:regular-normal-linear}の三次PBW条件により
三文字商と四文字の持上げは全次数で正則である。

\smallskip
\noindent\textit{中心表示での移動と同型。}
中心表示の三文字商の移動 $\rho$ は三角形上で固定点を持たない。
これには節点での接空間を仮定しない。曲線のグラフは平面の積で三式の完全交叉であり、
周囲の二平面方程式を加えて局所完全交叉である。従って六式の超過交叉の寄与は
$\deg(\omega_Y^{-1}(2,2))=12$ である（余法束の行列式または系\ref{cor:point-graph-excess}の法束計算）。
よって孤立部の長さは8である。中心関係によりこれらは全てアフィン対角 $t=1$ 上にあり、
そこで三つの二次式の孤立共通零スキームになる。対角の射影閉包で三式が無限遠にも零点を持てば
アフィン孤立部の長さは8未満となる。実際、一般の二つの線型結合は共通曲面を持たず
（三式の共通曲面は有限なアフィン部にも無限遠の三角形にも入らない）、
次数4の完全交叉曲線を与える。共通曲線成分があればその正次数分を除いて第三式と交わる
残りのサイクルの交数は8未満であり、各アフィン孤立長を含む。共通曲線がなければ
無限遠点の正の交数を8から引くことになる。よって $\rho$ は固定点を持たず、三辺を巡回する。

頂点の順と座標を選び、$\rho$ が頂点を一つ送る規約にそろえる。
辺の座標比の倍率から関係は $x_{i+1}x_{i+2}-q_i x_i^2$ となる
（頂点の巡回の向きは表示に合わせて取る）。$q_i$ は全て非零、積 $\lambda$ が1なら
行列式が恒等零である。対角座標変更で $q_i$ は $d_{i+1}d_{i+2}/d_i^2$ の倍率を受け、
その像は積1の全二次元トーラスである。従って三つとも $q^3=\lambda$ の同じ $q$ にできる。
この積は同型不変である。平面の同型は三角形を保つ単項行列であり、
標準の巡回を標準の巡回に運ぶ頂点置換は巡回置換だけである。同じ $q$ を保つ対角変換は
共通倍率を除き $\operatorname{diag}(1,\zeta,\zeta^2)$ のべきで尽きる。

逆に $q^3\ne1$ では三式の左右行列式は三角形を簡約に切り、各点で階数2である。
巡回syzygy $\sum x_i r_i=\sum r_i x_i$ は両端で全支持なので三生成商は正則である。
系\ref{cor:sklyanin-central-classification}の係数計算は、純平方係数 $-q$ が非零で
巡回syzygyを持つ今の三式にも適用できる。対称行列と任意の定数欄を得て、
文字への $t$ の加算で対角を一意に零にできる。固定点のない平面グラフの近くでは
中心関係が両位置の $t$ を消し、対角の三式は無限遠に零点を持たず長さ8の完全交叉となる。
三角形の張る平面は内在的なので、線を指定しない全ての同型についても上の有限群と
重み $(1,2)$ の共通倍率で尽きる。

最後に $r_i$ に $u x_{i+2}x_{i+1}$ を加える。三次の行列式は全体倍率を除き
$(1+u^3-q^3)x_0x_1x_2+uq(x_0^3+x_1^3+x_2^3)$ である（$u$ の符号規約による全体倍率を許す）。
一般の $u$ で滑らかなHesse三次となり、対角制限も一般に固定点を持たない
（存在すれば $(1+u)^3=q^3$）。三次syzygyと同じ六補正係数は有効で、
一般点と $u=0$ の全次数のHilbert関数が一定である。各次数の有限表示加群は
$k[u]_{(u)}$ 上で自由なので、指定した全点への平坦特殊化を得る。
\end{proof}

% !TeX root = ../papers/nc-p3.tex
\section{被約な特異四次点曲線の被覆と多角形の係数判定}
\label{sec:reduced-quartic-point-curves}
\begin{theorem}\label{thm:reduced-quartic-point-curves}
四辺の左右が点ファイバーで、点グラフが純な被約四次曲線 $Y$ と交わらない有限部の和である
標準有限Frobenius窓を取る。滑らかな楕円四次のほかに残り得る $Y$ は、
二つの節点で交わる二つの平面二次曲線、または四辺四頂点の空間四角形である。
いずれも $(2,2)$ 完全交叉で、有限部は長さ4である。多角形では四移動は固定点のない回転を誘導し
（成分上の座標倍率を許す）、二階差は $K^2=1$ を満たす。
\end{theorem}
\begin{proof}
非連結なら全連結成分が次数3以下となる。合計が4なので次数1または2の成分を含む。
滑らかな開閉有理成分、または節点が一本しかない有理成分は固定微分の次数式で除かれる。
複数の成分を四移動が入れ替える場合にもこの式は成立する。

一位置で平面四次なら算術種数3である。一方、連結被約な非退化空間四次の算術種数は高々1である。
実際、$H^1(I_Y)=0$ なので一般平面切断の相異なる四点は平面を張る。
そのHilbert関数は次数 $0,1,2,\ldots$ で $1,3,4,\ldots$ である（四点の各々を他の三点から二直線で分離できる）。
非零因子となる平面式による斉次座標環の商はこの切断の座標環に全射するから、
四次曲線のHilbert多項式は $4n$ 以上となる。ゆえに一位置で平面なら全位置で平面であり、
隣接平面の積で一般の三つの双一次関係を取ると完全交叉曲線になる。共通基点が一次元以下だからである。
その各射影次数は3で四次曲線を含めない。

整で特異な非退化四次なら算術種数1、正規化は有理、全 $\delta$ は1となる。
正規化への余長1の局所環は二枝の共通値を課す節点か、一枝の $k+t^2k[[t]]$ の尖点である。
従ってGorensteinであり、双対化束は次数零の非零切断を持つので自明である。
前節の平面三次と同じ固定微分の証明で、唯一の特異点が固定されることに矛盾する。

平面三次部分曲線と外の直線の和も除ける。交点はただ一点で、平面方程式が直線上で単純に消える。
交点が三次の滑らかな点なら直線は節点一本の端となる。特異点なら局所座標で
平面を $z=0$、直線を $x=y=0$、三次の局所式を $f(x,y)$ として
和集合の式は $zx,zy,f$ である。直線上の余接層には長さ2の捩れがあり、
$B=0$、他の特異点はこの直線に乗らない。直線を点ごとに保つことも、
定理\ref{thm:fixed-differential}の第一項の零次の固定因子がこの交点を含むこともできない。
実際ここは全曲線で唯一接次元3の点なので固定される。

残りの成分次数を列挙する。
\begin{itemize}
\item $(3,1)$ の非平面三次は有理正規三次であり、直線との交わりは長さ2以下である
（この三次は二次式で切られる）。長さ1なら端の節点である。長さ2なら異なる二点での横断か一つの接点である。
二つの滑らかな有理成分を長さ2で貼った曲線の双対化束は自明である。
実際二点の場合は両側で留数が逆になる対数微分が生成し、一点の場合は通常の二枝接点となる。
後者は共通の一次jetを持つ二つの形式冪級数を貼った平面特異点で、導手の逆像は各枝上で次数2、
算術種数1である。各成分上で双対化束の次数は零で、非零の大域切断は
片方で恒等零なら接点で他方の定数も零となるため生成元となる。
四移動は以下このような自明な双対化束の有理成分の和では固定点を持てない
（定理\ref{thm:fixed-differential}、$\deg M<0$、Cartier固定因子の総次数零）。
ところが次数の異なる二成分を保つ自己同型は有理成分上で固定点を持つ。
\item $(2,2)$ では二平面の共通直線上に二曲線の交点があり、長さは高々2。
空、単純一交点、重複一交点は既述の議論で除ける。二単純交点なら記載した二角形である。
\item $(2,1,1)$ では三次部分曲線が平面なら既に除いた。従って各直線は二次曲線と高々一点で単純に交わる。
三つの対交点が相異なれば三角形の双対グラフで、双対化束は自明だが次数2の成分を保たねばならない。
対交点が不足すれば節点一本以下の端がある。三枝が同一点で会う場合は、どの二枝も接方向が異なる。
この三重点での各枝の余接捩れは長さ2である。三方向が独立なら三座標軸の和に形式同型である
（各枝をパラメータ表示し、その三つの表示のベクトル和による形式座標変更を用いる）。
三方向が一平面に入ればその平面への射影の三枝は通常の平面三重点で、
高さ座標はこの三枝上で二次以上から始まるため平面三重点の局所環に属する。
実際一枝にだけ二次以上の任意の級数を載せるには他の二枝の局所式の積を掛ければよい。
平面三枝の積の勾配は各枝でちょうど二次の共通因子を持つ。
いずれも次数式で恒等は除かれ、非恒等なら唯一の三重点が固定されることに反する。
\item 四直線では三本が同一平面に入る場合を既に除いた。四本共点なら方向は三点共線のない平面四点で、
二次式二本の完全交叉である。この錐の双対化束は自明、唯一の特異点が固定されてしまう。
三本だけが共点なら残る直線はその三本の高々一本としか会えず、端になってしまう。
三重点がなければ単純な交点グラフは三角形を含まない（相異なる三交点を作る三直線は同一平面に入る）。
全頂点の価数が2以上であることから四角形だけが残る。
\end{itemize}
多角形は節点の各枝に留数 $+1,-1$ の微分を順に貼って $\omega_Y=\mathcal O_Y$ を得る。
補題\ref{lem:genus-one-quartic-mark}から $(2,2)$ 完全交叉である。グラフの埋込みも局所完全交叉である
（片側で余次元2の完全交叉、他の三アフィン座標は第一側の正則関数の持上げで消去する）。
超過交叉公式で曲線の寄与は16となる。四移動に固定点はなく、向きを反す多角形自己同型は
成分か節点を固定するから、成分置換は非自明な回転となる。固定切断が消えないことから $K^2=1$ を得る。
\end{proof}

\begin{corollary}\label{cor:quartic-cycle-atlas}
二角形と四角形のそれぞれで、全係数に対するAS性、添字保持同型および最小周期は有限に判定できる。
$m=2,4$、$d=4/m$ とし、$m$ 本の射影直線の $\infty$ と次の $0$ を貼ってモデルを固定する。
多重次数 $(d,\ldots,d)$、貼合せ倍率1の束を $L$ とする。
$D\in k^*$、$K\in\{1,-1\}\subset\operatorname{Pic}^0(Y)=k^*$ に対し
$L_0=L,L_1=LD,L_2=LD^2K$ を取る。最初の二辺の切断乗法核（各次元8）から六平面を一つずつ選ぶ。
系\ref{cor:elliptic-point-finite-atlas}の二窓復元・点ファイバー検査に加え、第一グラフが
$Y$ と交わらない有限部の和となることを課す。これはグラフのイデアルを $I_Y$ で飽和して残余を取り、
その次元、非交叉とイデアルの交わりの一致を射影上で検査できる。
最後に定理\ref{thm:finite-orbit-quartic-marks}の四位置の二次元二次核の積を検査する。

同型比較は多角形の回転・反転と座標倍率 $z_a\mapsto s_a z_a$、
$(\prod_a s_a)^d=1$ の引戻しで二束・二関係を対応させることで必要十分である。
多角形の二種類の間に同型はない。移動後に最初の束を $L$ に戻して同じ比較を行えば
最小周期 $1,2,4$ を区別できる。
\end{corollary}
\begin{proof}
全切断系と乗法核の次元は補題\ref{lem:genus-one-quartic-mark}による。
各成分で束を標準自明化すると、貼合せの唯一の不変量は倍率の積である。
座標倍率を作用させると次数 $d$ の座標の両端の値の比が $s_a^d$ 倍となるので、
任意の第一束を $L$ に移せ、その安定群が記載のものになる。反転は $D$ を逆数にする。
曲線は点グラフの唯一の純曲線全体なので同型がこれを保ち、全切断系のため文字写像は引戻しで決まる。
逆は二窓の延長定理である。正則性は印の軌道が一点の有限判定を用いた。

連続な安定群の商も全ての零係数パターンで単項式比較に直せる。
連結トーラス $\prod s_a=1$ は各切断空間に対角的に作用する。
節点値を一つだけ立てる隣接二成分の端単項式と、各成分内部の単項式を基底に取ればよい。
貼合せを保つ各成分の倍率補正は、一成分で1と決め、次に $s_a^d$ を順に掛けて得る指標である。
両Grassmann座標を周囲の文字テンソル基底で書き、各非零Plücker支持で最初の座標を1とする。
残った非零座標の指標差の整数行列をSmith形にすれば、独立な指標の座標は全て1にでき、
残る単項式座標の一致が連結トーラス同型と同値である。代数閉体なので非零の整数乗写像は全射である。
最後に $\prod s_a\in\mu_d$ と多角形の有限置換を比較すればよい。
\end{proof}

% !TeX root = ../papers/nc-p3.tex
\section{整な特異五次点曲線の排除}
\label{sec:singular-integral-quintic-exclusion}
\begin{theorem}\label{thm:singular-integral-quintic-exclusion}
四辺の左右が点ファイバーで、点グラフが整な被約五次曲線と交わらない有限部の和となる標準AS窓は、
定理\ref{thm:elliptic-quintic-regularity}の滑らかな楕円五次の二類で尽きる。
曲線の滑らかさもGorenstein性も仮定しない。
\end{theorem}
\begin{proof}[二次曲面とGorensteinの場合]
曲線を $C$ とする。平面五次なら算術種数6で、非退化整五次のCastelnuovo上界2により全位置で平面となり、
隣接平面の積での次数上限3に反する。算術種数0では滑らかであり、
種数1の特異整曲線は唯一の節点か尖点を持ち、定理\ref{thm:reduced-quartic-point-curves}の固定イデアルの証明で次数によらず除かれる。
以後 $p_a(C)=2$ とする。次数5の $L_i$ は完全な四切断を持ち、$h^0(L_i^2)=9$ である。
実際、整曲線上の階数1無捩れ双対化層の次数は2であり、これに $L_i^{-n}$ を掛けた負次数の層には非零切断がない。
各位置で $C$ を含む既約二次曲面 $Q_i$ がただ一つあり、$L_i\ne L_{i+1}$ である
（偏極が等しければ全双一次式がその曲面の対角で消える）。
滑らかな二次曲面なら曲線の型は $(2,3)$ またはその逆である。
二次錐なら解消上の狭義変換の型は $m s+5f$、$m=1,2$ である。ここで $s$ は負切断である。
$m=2$ なら $s$ と長さ1で交わり、像は頂点で滑らかである。実際錐への写像の局所座標は
$(z,zu,zu^2)$ で、$s$ の方程式 $z$ が枝の一様化元になる。従って狭義変換は元の曲線と同型である。

$m=1$ では曲線を二項形式で
\[
 [Fs^2:Fst:Ft^2:G],\qquad \deg F=3,\quad\deg G=5,\quad (F,G)=1
\]
と書ける（母線への射影は同型で、頂点に写る因子が $F=0$ である）。
頂点以外は滑らかで、頂点の重複度は3、接次元は3である。後者は辺の三アフィン座標の線型結合が
極大イデアルの二乗に入れば、その二項二次係数が $F$ の全零因子上で消えてしまうからである。
この点はGorensteinではない。余次元2のGorenstein曲線イデアルは完全交叉である
（極小Hilbert--Burch解消の末尾階数が型に等しい）。接次元3での二つの方程式はいずれも二次以上から始まり、
一般平面切断での局所交数は4以上になるからである。以上によりGorensteinの場合には $m=1$ は混在しない。

Gorensteinとして $H=\omega_C$ と置く。滑らかな二次曲面の次数2の母線束も
$m=2$ の錐の母線束も $H$ に一致する。両者は次数2で二切断を持つため、
Riemann--Rochと次数零の可逆層の切断条件から従う。これが唯一の次数2の鉛筆である。
一般の滑らかな $p$ に対して、定理\ref{thm:genus-two-quintic-ruling}で使った反対縮約の
三次行列式 $S_p$ と階数低下因子 $D_p$ を取る。この場合にも $D_p$ は $C$ の滑らかな部分に台を持つ長さ2の
Cartier因子で、$\mathcal O_C(D_p)=K^{-1}(p+\sigma^{-1}p)$ である。理由を明記する。
定理\ref{thm:opposite-curve-divisor}の二核ベクトルの共通因子は、$p$ を一般に固定すれば
相異なる位置でそれぞれ滑らかな点 $p,\sigma^{-1}p$ の可逆イデアルである。
残る二階数束の行列式とその束の計算は $C$ 上でそのまま成立する。
特異点 $q$ を固定しても、位置 $q$ での端支持はともに3、反対位置の曲線は空間を張る。
従って $p$ が動くとき行列の一般階数が1以下となることはない
（正規化の滑らかな点の位数を使う同定理の二つの三次元関数空間の積の下界5）。
よって一般の $p$ では $D_p$ が全特異点と錐の頂点を避ける。

$S_p$ の $Q_2$ への制限も非零である。前掲の五次の定理と同様、
恒等零なら余因子行列を曲面または錐の解消のCox環で列と行に分ける。
各々の $C$ 上の写像は $L_1(-p)$ と $L_3(-\sigma^{-1}p)$ による平面射影である。
これらは一般の $p$ で二次曲線には入らない。入れば次数2の鉛筆の一意性から
該当する次数4の束が $H^2$ となり、$p$ は高々一つである。
従って滑らかな二次曲面での両因子の型は $(1,1)$、錐の解消では $s+2f$ である
（片側の母線次数が零なら二次までの一変数表示となり、または次数2の射を経由する）。
よって前者の射影を $L_2$ の三切断でも表すことになる。
この三切断は基点自由な $L_1(-p)$ の三切断の共通倍率であり、その倍率は
$C$ 自身で正則な線束の射である。基点自由な三切断の一つが単元の開集合で比を取ればよい。
従って $L_2(-x)=L_1(-p)$ となる長さ1のCartier因子 $x$ があり、これは滑らかな点である。
この等式が動く $p$ で成り立てば、二つの表示を比較した相異なる二次因子はどちらも唯一の $|H|$ のファイバーになる。
一方の点を固定すると相手も決まるので不可能である（$L_1\ne L_2$）。

従って三次曲面の切断から $C$ を除いて得る残余は前掲定理の通り一本の母線、または解消上の $s+f$ である。
$C$ への制限は $P_2=L_2/H$ の次数3の有効因子 $E_p$ となる。
実際滑らかな二次曲面で残る母線の束は $L_2/H$、錐でも $s+f=(s+2f)-f$ である。
二次小行列が生成する階数低下のイデアルには $S_p$ の微分が入る。
$D_p$ は曲面と $C$ の滑らかな部分にあるため、局所的に $S_p=c e$ として $E_p$ は $D_p$ を含む。
差は滑らかな一点 $r_p$（長さ1の有効Cartier因子）で、
\[
                p+\tau p+r_p\in |P_2 K|,\qquad \tau=\sigma^{-1}.
\]
これは次元1の線型系である。次数3なのでRiemann--Rochで切断はちょうど二つとなる。

この鉛筆を正規化に引き戻して基点を除いた非定数写像を $h$ とする。
一般の $p$ はただ一つのファイバーに入り、その同じファイバーに $\tau p$ も入る。
$\tau=1$ なら一般点で重複度2になり標数零の分離性に反する。
従って $\tau$ の位数は $\deg h\le3$ を割り2か3である。
位数3なら基点はなく、$r_p=\tau^2p$ である。同じ議論を各位置の一般の $p$ に適用すれば
三点の因子が共通なので全ての $P_i K$ が同じになり、$L_i\ne L_{i+1}$ に反する。
位数2なら写像は次数2で、正規化上の固定部は長さ1である。
この基点は $C$ の滑らかな点にある。実際、一般の $p$ で $p,\tau p$ は滑らかであり、
その二点を除いた有効Cartier因子 $r_p$ も滑らかである。従って $C$ 上の鉛筆の固定部は
滑らかな一点で、これを除いた $C\to\mathbb P^1$ は基点自由な次数2の有限平坦射である。
$\tau$ はその被覆対合である。しかし二重被覆の固定因子の長さは6である。
実際、trace分解は $\mathcal O\oplus N^{-1}$ で、$\chi(\mathcal O_C)=-1$ より $\deg N=3$、
固定因子の局所式は反不変生成元 $y=0$、$y^2$ は $N^2$ の切断である。
他方、定理\ref{thm:fixed-differential}の第三項により固定因子は次数 $2p_a-2=2$ となり矛盾する。
以上は特異な二重被覆でも固定イデアルの長さをそのまま数えている。

\smallskip
\noindent\textit{三重の頂点を持つ場合。}
非Gorensteinの場合は全位置で $m=1$ の錐表示になる。唯一の特異点は四移動が固定し、
その接次元は3である。従って命題\ref{prop:ore-mark-recognition}が適用できる。
水平三文字を全位置で同一視して、辺は一定の三文字商関係 $S$ と
$X_i h-h X_{i+1}=d_i(h)$ である。AS窓なら三文字部分も標準ASである。
頂点からの水平射影は全位置で正規化を滑らかな平面二次曲線に同型に写す。

三文字の点グラフが平面全体の非恒等射影変換なら、定理\ref{thm:plane-ore-unmarked}の表示で
点グラフは曲面を含む（零補正は $X=0$ の平面、固定直線上の非零補正は同定理の二次曲面）。
平面上で恒等なら、任意の持ち上がる一般の水平点で前後の高さに同じ量を加えられ、
二次曲線全体が持ち上がる以上やはり曲面を得る。
残る三文字グラフは平面三次上の自己同型である。水平二次曲線を含むので二次曲線と直線の和であり、
二次曲線は唯一だから全位置で同じである。その上の自己同型を $g$ とする。
$g=1$ も高さを同時に動かせるため除かれる。

一般の正規化点 $p$ で二次形式の座標列を $v(p)$ とすると、混合式は高さ $a,b$ に対して
\[
               a\,v(gp)-b\,v(p)=d_i(v(p),v(gp))
\]
となる（座標の持上げを固定する）。$a$ は本来水平束 $\mathcal O_{\mathbb P^1}(2)$ の有理切断である。
Cramerの公式により、その極因子の次数は $v(p)\wedge v(gp)$ の三成分の共通零因子の次数以下である。
局所的に各座標束を自明化して単元となる二次小行列を共通因子除去後に使えばよい。
二次埋込みは閉埋込みなので、この共通零イデアルは $g$ の固定イデアルである。
非恒等な射影直線自己同型の固定長は重複込みで2である。ところが五次の表示で高さは $G/F$、
その極因子の次数は3で、水平切断の加算でも変化しない。これで最後の可能性も消える。
\end{proof}

% !TeX root = ../papers/nc-p3.tex
\section{特異点を含む整七次点曲線の全係数}
\label{sec:integral-septic-boundary}
\begin{proposition}\label{prop:automatic-spatial-tangent-mark}
四辺の左右が点ファイバーの標準有限Frobenius候補では、共通点グラフの接次元3の全ての点は
標準四移動の固定点であり、それぞれ命題\ref{prop:ore-mark-recognition}のOre表示を与える。
従ってこのような点が一つでもあれば、全次数AS性は三生成部分の有限判定に帰着する。
\end{proposition}
\begin{proof}
任意の接次元3の点からグラフに沿って全位置に点 $p_i$ を取り、初めは周期性を仮定しない。
$H_i=\ker p_i$ とし、定理\ref{thm:maximal-tangent-ore}の証明と同じく個別の評価階数3と結合階数3から、
三つの水平関係 $S_i$ と、残りの文字を右へ移す三つの混合関係を得る。
水平三次交叉の次元を $b$ とすれば全三次元数が $19+b$ 以下なので $b\ge1$。
第一を $p_i$ で縮約して零になる三次交叉は、高さ4の完全対合と一意な前点から
$w(p_{i-1})$ の張る直線である（標準四窓の文字同定を使用）。これが水平三次交叉である。
末端支持は直前の点列から $H_{i-2}$ 全体であり、一方水平性から $H_{i+2}$ に入る。
従って両者は一致する。全位置で同じ議論が通るので四周期であり、Ore認識命題の条件が全て成立する。
\end{proof}

\begin{theorem}\label{thm:integral-septic-boundary}
点ファイバーを持ち点曲線が整な被約七次となる全標準AS窓は、定理\ref{thm:septic-bielliptic-forms}の
八係数のうち、表示した三式が\emph{整な完全交叉曲線}を定めるもの全部である。
全係数で同じ倍率と交換の同型判定、$\lambda\mapsto-\lambda$ による最小周期判定が成り立つ。
この全層は既約で同型類は五次元、全ての特異曲線の窓は滑らかな七次窓から平坦に特殊化する。
\end{theorem}
\begin{proof}[非Gorenstein点の除外]
純性と算術種数5は定理\ref{thm:septic-point-graph}による。どの埋込みにも二次式がない。
滑らかな二次曲面の二次数では種数5にならない。二次錐でも、次数7の整曲線のWeil類は
母線 $A$ と超平面 $H$ により $A+3H$ である。実際解消の負切断を法とした類群は母線が生成し、$H=2A$。
従って反射的な曲線イデアルは $I_{A/Q}(-3)$ と同型で、
$\chi(\mathcal O_C(n))=(n+1)^2-(n-2)(n-3)=7n-5$ となり算術種数6である。
平面七次の種数も5に等しくない。

$h^0(L_i)=4$ である。これ以上なら完全系の五切断による非退化七次の射影四空間模型が得られ、
埋込み四切断を含めるので算術種数5のままであり、Castelnuovo上界3に反する。
双対性により $L_i\hookrightarrow\omega_C$ はスカラーを除き一意で、商は長さ1である。
従って非Gorenstein点があればただ一点 $P$ で、全ての商がそこに支持される。
$P$ の埋込み次元は3である（接次元2以下の被約曲線特異点は平面超曲面または滑らか）。
その極大イデアルを $\mathfrak m$ とすると $\mathfrak m\omega_C\subset L_i$ である。
$\omega_C$ が局所巡回でなく、$\omega_C/L_i$ の長さが1、$L_i$ が自由階数1なので、
$\mathfrak m\omega_C=\mathfrak m L_i$ でなければならない。
実際 $\mathfrak mL_i\subset\mathfrak m\omega_C\subset L_i$ の上の包含が等号なら
$\omega_C/\mathfrak m\omega_C$ が一次元になってしまう。

これで $P$ で消える三文字は全ての位置で $H^0(\mathfrak m_P\omega_C)$ という同じ空間になり、
積を曲線の関数体で比較できる。この水平三文字の間に可換二次式はない（七次を消す二次式がない）。
従ってOre表示の三水平関係は三つの交換子で尽きる。命題\ref{prop:automatic-spatial-tangent-mark}で
混合関係の一次係数により水平空間を順に同定すると、一定の三生成基底の点グラフは平面の射影変換である。
非恒等の場合は定理\ref{thm:plane-ore-unmarked}により二次曲面か平面が四文字の点グラフに入る。
恒等なら一般の持ち上がる水平点で二位置の高さを同時に変えられ、水平射影は曲線を描くので
やはり全グラフの次元が2以上となる。これで非Gorenstein点は除かれる。

\smallskip
\noindent\textit{二つの滑らかな射影中心。}
以下 $\omega_C$ は可逆なので $L_i=\omega_C(-p_i)$、$p_i$ は滑らかな点である
（長さ1のCartier因子は正則点）。滑らかな点だけに支持される相異なる二つの次数3以下の
有効Cartier因子が線型同値になることはない。共通部分を除けば基点自由な次数3以下の鉛筆になる。
その $p_i$ を通るファイバーに必要なら点を加えれば $p_i+D$、$\deg D=2$、
$h^0(p_i+D)\ge2$ となる。残余 $D$ が特異点に乗ってもCartier因子である。
Riemann--Rochにより $h^0(L_i(-D))\ge3$ となり、四切断による閉埋込みが長さ2を分離することに反する。
従って偏極の三階差から得る $p_i+2p_{i+3}\sim2p_{i+1}+p_{i+4}$ は因子の等式である。
定理\ref{thm:septic-point-graph}と同じく、二次式の非存在により異なる二点 $p,q$ が交互に現れ、
$p-q$ の位数はちょうど4である。

$\omega_C$ の五つの切断は閉埋込みを与える。$p$ の外では四切断への射影が既に埋込みであり、
$p$ で消えない第五切断と単純に消える $L_0$ 由来の切断が $p$ の接方向も分離する。
標準二次式は $h^0(\omega_C^2)=12$ により少なくとも三つであり、$p$ の三次元余法空間への微分が
単射なのでちょうど三つである。単射性は $p$ を中心とする射影像に二次式がないことと同値である。
これ以降の六積関係の計算、共通三文字の三交換子と順序付き二点からの同型復元は
定理\ref{thm:septic-point-graph}の証明が適用できる。

\smallskip
\noindent\textit{係数の必要性と逆構成。}
定理\ref{thm:septic-bielliptic-forms}の三標準式は整な完全交叉である。
共通三文字 $J$ が非零の開部分はグラフから復元でき、$J=0$ の直線には二つの指定点しかないからである。
二点が滑らかで接線が曲線に含まれないことから、同定理の三次重なり計算で
$xy=D,\ xr-ry=F,\ xs+sy=G$ を得る。ここで $r,s,h$ は共通三文字である。
$D,F,G$ の $\langle r^2,s^2\rangle$ を法とした方向は一本である。
共通代表の $h^2$ 係数が零だとすると $h$ 軸の点で六関係を左評価した像は
$\langle r,s\rangle$ に入る。実際三交換子の評価もここに入り、残る三式には $hx,hy,h^2$ がない。
従って右点ファイバーが少なくとも直線となり矛盾する。
平方完成と二つの補文字の変更により滑らかな場合と同じ八係数表示になる。

逆にこの表示が整な完全交叉を定めれば、両指定点 $r=s=h=0$, $v=\pm u$ は滑らかである
（三式のそこでの微分は $dr,ds,dv\mp du$ の非零倍）。長さ3のスキームを含む直線は三式で
恒等的に消えて曲線成分となるので存在しない。従って各指定点からの完全射影は閉埋込みである
（長さ2の分離が失敗すれば、指定点と合わせた長さ3が直線に乗る）。
$J\ne0$ での点対の復元は座標比の正則な逆を持つため、特異点でもスキーム同型である。
一方だけで $J=0$ の箇所は滑らかな指定点での厳密変換とその三つの独立な微分だけを用いた
前掲定理の計算で同型になる。両側で $J=0$ は $xy$ の式で不可能である。

全次数の計算にも新たな例外はない。$F,G,D$ の $H=h^2+\lambda rs$ 係数の少なくとも一本は非零である
（全て零なら三式は $h$ と独立で、整な非退化曲線にならない）。
従って中央三二次式上の単項式基底と三次重なりの計算が成立する。
双対の可換三二次式に共通射影零点がない理由は整曲線上で二つの非零切断の積が零にならないことであり、
これも特異点を許す。定理\ref{thm:septic-bielliptic-forms}の有限Frobenius計算と
命題\ref{prop:commuting-three-point-graph}が全次数の正則性を与える。
二つの射影中心と二つの固有文字の復元も変わらず、従って同型と周期の倍率・交換判定が全域で成立する。
射影完全交叉の平坦な開集合で幾何学的整性は開であり、滑らかな係数はその中で稠密である。
任意の許容係数を一般の直線で滑らかな係数につなぐと、各次数のHilbert関数の一致から
代数族もその点の離散付値環上で平坦となる。同型類の次元は滑らかな開集合の5である。
\end{proof}

\paragraph{特異な許容係数の例。}
$\lambda=0$, $F=S+H$, $G=R+H$, $4D=3R-3S+8H$ とする。
$r=1$, $s=x\ne0$, $z=h^2$ と置けば $u=x^2+z$, $v=(1+z)/x$、残る式は
\[
        (1-x^2)z^2+(2-2x^4-8x^2)z+(1-x^2)^3=0
\]
である。判別式は $16x^2(1+x^2)(3x^2-1)$ で平方でない。
$x=1,z=0$ では $z$ の係数が非零で解の位数は3、従ってこの二次関数体で $z$ も平方でない。
$h^2=z$ の全曲線はこの開集合で一般点が被約かつ既約であり、原始多項式なので縦の曲線成分もない。
$r=0$ または $s=0$ に三式の解は有限個しかなく、完全交叉の純性により曲線は整である。
$x=1,h=0$ は尖点であり、この窓も全次数で正則となる。

% !TeX root = ../papers/nc-p3.tex
\section{整六次点曲線の特異点と正規化の制限}
\label{sec:singular-sextic-normalization}
\begin{theorem}\label{thm:singular-sextic-normalization}
点ファイバーを持つ標準AS窓の点グラフが整な被約六次曲線 $C$ と交わらない有限部の和なら、
有限部は空で、$C$ は抽象曲線として整な特異平面四次である。
正規化の種数は0または1に限る。標準四移動が恒等でなければ、標準平面上の座標を取ると
$[u:v:w]\mapsto[\eta u:v:w]$、$\eta$ の位数は2または4であり、$K^2=1$ である。
位数4の特異模型は $u^4+F_4(v,w)=0$ で、二項四次の根の重複型は $(2,1,1)$ または $(3,1)$ である。
\end{theorem}
\begin{proof}[三次元接点の高さと極]
平面六次は隣接平面積の次数で除かれる。接次元3の点があれば
命題\ref{prop:automatic-spatial-tangent-mark}のOre表示を取る。
その水平射影の像は整な非退化平面曲線 $D$ である。三生成基底が平面グラフなら
定理\ref{thm:integral-septic-boundary}と同じく曲面が入る。残る基底は三次曲線のグラフで、
二次曲線なら定理\ref{thm:singular-integral-quintic-exclusion}の二次の極上限により
持ち上がる曲線の次数は高々4となる。実際、二つの高さは一般点で水平点から一意に戻るため
水平点曲線の関数体は $C$ の関数体である。従って $D$ は整三次で、基底の移動を $g$ とする。

$g$ の固定点 $x$ では $v(x)\wedge d_i(v(x),v(gx))\ne0$ である。零なら混合式が解を持ち、
両方の高さを同時に変える直線が点グラフの成分になるからである。
一般点で $a\,v(gx)-b\,v(x)=d_i(v(x),v(gx))$ を満たすので、
一つが単元となる成分で割ると、固定点の局所イデアルは $1/a$ で生成される（局所自明化を用いる）。
その正規化への引戻しは正確に高さの極因子となる。従って元の曲線の次数は
$3$ とこの極次数の和である。高さと水平三座標で基点のない射を与え、一般点では逆もあるからである。

滑らかな三次の基底なら定理\ref{thm:all-elliptic-ore-windows}を使える。
固定点のない移動では極はなく、固定点がある位数3,4の基底では補正が零に同値で、
各固定点に直線が生じる。残る反転は四つの単純固定点を持ち次数は7である。
節点三次の自己同型で二枝を保つ非恒等なものは正規化で $z\mapsto cz$、
節点での固定イデアルは二枝の極大イデアルで主でない。枝を交換すれば $z\mapsto c/z$ となり、
枝での固定イデアルの引戻しは各長さ1、その外に単純固定点二つで極次数は4である。
尖点三次の滑らかな部分で移動は $z\mapsto az+b$ と書ける。
三生成ASの偏極条件で $g^*L/L$ は非自明かつ $g^*$ 不変な次数零類である
（補題\ref{lem:ternary-finite-graph-test}の三次関係の行列式、または補題\ref{lem:reduced-ternary-fourth-shift}の束計算）。
$\operatorname{Pic}^0=\mathbb G_a$ への作用は $a$ または $a^{-1}$ 倍である
（尖点での正規化の一次jetの単元商）。よって $a=1$, $b\ne0$ である。
尖点のパラメータ $t=1/z$ に対し $t\mapsto t/(1+bt)$ で、
$t^2,t^3$ の差の位数は3,4となる。$k[[t^2,t^3]]$ に位数1の元がなく、この固定イデアルも主でない。
以上で六次では接次元3の点はない。

\smallskip
\noindent\textit{標準平面と固定直線。}
従って全特異点は局所平面、グラフは局所完全交叉である。
超過交叉の長さ $2p_a-6$ とCastelnuovo上界から $p_a=3,4$ であり、
種数4なら全埋込みが標準埋込みで同じ二次曲面の対角を含むという既出の矛盾になる。
種数3では有限部は空である。標準束は基点自由である。実際任意の点 $x$ について
$\mathcal Hom(\mathfrak m_x,\mathcal O_C)$ の大域切断は定数のみである。
滑らかな点で単純極一つを持つ函数は次数1の射を与えてしまい、特異点では
$(\mathcal O:\mathfrak m_x)=\operatorname{End}(\mathfrak m_x)$ である
（像に単元があれば極大イデアルが主になる）。後者は有限な整な部分正規化の函数層である。
双対性から基点自由性を得る。三切断による次数4の標準写像は二次曲線への有限平坦二重被覆か、
算術種数が等しいため整平面四次への同型である。二重被覆では定理\ref{thm:smooth-sextic-exclusion}の
鉛筆分裂の証明がそのまま成立する。$L_i/\omega_C$ の有効な長さ2のCartier因子は埋込みに反し、
有限平坦な二重被覆のtrace分解から $\omega_C=J^2$ である。
以後は直線上の束の分裂と各長さ2の環上の乗法しか用いず、三次の次元上限18に達する。

標準平面で $\tau=\sigma^{-1}$ が非恒等なら、固定切断は $C$ 上で長さ4のCartier因子を切る。
射影変換が固定直線を持たなければ、平面での固定スキームの長さは3で足りない
（接束の切断の孤立零の最高Chern数、またはJordan形での小行列式）。
固定直線を持つ非半単純な変換は一つの推移 $u\mapsto u+v$ なので、
不変な四次曲線の式は $v,w$ だけになり整でない。
半単純なら固定直線の四交点で全長が尽き、中心点は曲線の外である。
$u^4$ と $u$ を含まない部分がともに非零だから倍率の四乗は1である。
固定因子は標準直線切断なので定理\ref{thm:fixed-differential}より $K^2=1$。
位数4の式で全根の多重度が偶数なら可約で、全て単純なら滑らかである。

\smallskip
\noindent\textit{正規化の種数2の排除。}
この場合はただ一つの節点か尖点がある。定理\ref{thm:fixed-differential}の恒等の場合の式で、
$\deg M\ge 2-1=1$ かつ $\deg M(T)=\deg\omega_C=4$ なので $\tau\ne1$ である。
節点の二枝を保つ有限位数の非恒等写像は両枝で接倍率が1でなく、
固定イデアルが主にならない。交換する場合は接作用の固有値に1が必要なので、
その平方の枝接倍率が1となり、有限位数性から $\tau^2=1$ である。
尖点に位数4の作用があれば、正規化の一様化元を固有元にして、
$t^2,t^3$ の差がともに非零倍率となるため主でない。従ってどちらも対合である。
節点では固定長2と二つの滑らかな固定点、尖点では固定長3と一つの滑らかな固定点を得る。
正規化 $N$ 上の対合はいずれも二つの固定点を持ち、Hurwitzにより商の種数は1である。

一般の滑らかな $p$ で定理\ref{thm:nonhyperelliptic-sextic-exclusion}の
長さ2の反対階数低下因子 $D_p$ と三割線を取る。ここで要る特異曲線への移行を説明する。
どの六次埋込みにも二次式はない。滑らかな二次曲面上の種数3の型は $(2,4)$ またはその逆で二重被覆になる。
錐では次数6の曲線のWeil類は三超平面の和なので算術種数4である。
$T=L_2/\omega_C$ と $\omega_C/T$ は非有効な次数2の束である。
曲線の二次式がなく、全二次切断が十次元なのでHilbert--Burchの四三次小行列による記述は同じである。
線型三行四列は曲線上で常に階数2で、階数1なら接次元が3になってしまう。
行列式三次 $S_p$ の非消滅について、中支持4の証明は次数や滑らかさによらず、
中支持3のときに平面グラフから得る二点 $a,b$ は一般の $p$ の点列上にあり滑らかである。
同じ共通三文字を用いる次数7の束は非常に豊富である（長さ2のイデアルでひねった後の
$H^1$ の双対は負次数の階数1無捩れ層の切断となる）。従って三二次式の次数8との矛盾も同じである。
$D_p$ が一般に特異点を避け、$\mathcal O(D_p)=\mathcal O(p+\tau p)/K$ となる理由は
定理\ref{thm:singular-integral-quintic-exclusion}の二核束の議論による。

三割線との長さ3の交叉のうち $D_p$ を除いた長さ1の残余は一般には滑らかな可動点 $r(p)$ である。
もしある固定点 $z$ を通る直線ばかりなら、$z$ からの射影の一般ファイバーが
長さ2の $D_p$ を含む。射影平面を張る像に対し射影束の次数は $6-\operatorname{mult}_z C\le5$ なので、
分離性から二次曲線への二重被覆しかない。そのとき曲線が二次錐に入ってしまう。
$D_p$ の一般台は可動であるから上のファイバーは一般に $z$ 以外の滑らかな部分で数えてよい。
この議論で特異点が残余に常駐することも除かれる。
従って一般の $p$ で交叉は有効Cartier因子となる。たとえ余点が $D_p$ に重なっても
滑らかな点では長さの加算が因子の加算となる。直線を含む平面とRiemann--Rochから
$D_p+r(p)\sim T+s(p)$、$s(p)$ は有効な次数1の束が定める滑らかな点である。
一意な残余と一意な有効切断により両点は関数体上定義され、$N$ の自己写像に延びる。
$|T+s|$ は任意の滑らかな $s$ で一点なので $s(p)$ も非定数である。
種数2ではHurwitzにより $r,s$ はいずれも自己同型である。

$g=rs^{-1}$ とすると、$N$ のJacobianに
\[
                  (1-g_*)s_*=1+\tau_*
\]
が成立する。これは一般 $p$ での $p+\tau p+r(p)-s(p)\sim T K$ を正規化へ引き戻したもの。
$T K$ は非有効である。そうでなければ非超楕円性からその有効因子 $Z$ は唯一であり、
一般の $s$ では $|T K+s|=\{Z+s\}$ となる（双対な次数2の束の切断も高々一つ）。
これは可動な相異なる二点 $p,\tau p$ を含めない。従って $s$ は恒等でも $\tau$ でもない。
そうであれば上の次数3の因子から $s(p)$ を引いて $T K$ が有効になるためである。

標数零の有限個のデータを複素数へ埋め、有限自己同型群に不変なHermite内積を取る。
Jacobian等式に随伴を右から掛けると $g+g^{-1}=-2\tau$ である。
$\tau$ の $+1,-1$ 固有空間が各一次元のため、$g=-\tau$、$s$ は
この順に $1,\xi$ で作用する。有限位数の非恒等写像 $s,\tau s$ の固定点数は
Lefschetz数 $-2\operatorname{Re}\xi,2\operatorname{Re}\xi$ に等しい
（両方とも全固定点が単純）。従って両者とも固定点がなく、$\xi=\pm\mathrm i$。
種数2のJacobian作用は忠実なので $s^2=\tau$ である。
巡回四重被覆 $N\to N/\langle s\rangle$ の唯一の分岐軌道は $\tau$ の二固定点で、
分岐指数2である。これは不可能である。実際Kummer表示の四乗される根号の中の有理函数は
その一箇所でのみ位数が $2\bmod4$、他はすべて $0\bmod4$ となり主因子の次数が零にならない。
滑らかな種数3の場合は既出の全排除を用いれば、正規化種数は0,1のみである。
\end{proof}

% !TeX root = ../papers/nc-p3.tex
\section{整六次点曲線の全排除と整曲線層の分類}
\label{sec:integral-sextic-exclusion}
\begin{theorem}\label{thm:integral-sextic-exclusion}
標準Koszul左右AS正則窓で四辺の左右が点ファイバーなら、整な被約六次曲線と
これと交わらない零次元スキームが点グラフになることはない。
\end{theorem}
\begin{proof}[位数4と巡回被覆の固有束]
前節により残るのは正規化 $N$ の種数が0または1の整な特異平面四次 $C$ である。
位数4のときは $C\to\mathbb P^1$、$u^4=F_4(v,w)$ の有限平坦巡回被覆を $f$ とする。
$d_i=L_{i+1}/L_i$ は $\sigma^*d_i=d_iK^4=d_i$ を満たし、$d_i^4=1$ である。
実際四共役の積はノルムの引戻しで、次数零なので直線上で自明である。
従って $\sigma^*L_i=L_i d_i^4K^6=L_i$、各非有効次数2束 $T=L_i/\omega_C$ は不変である。
生成元の同型をスカラーで調整して四乗を1にし、$T$ を線型化する。
$f_*T$ の四つの固有層は各々直線束である（一般点で各指標が一回ずつ現れ、直線上で無捩れ）。
$H^0(T)=H^1(T)=0$ なのでこれらは全て $\mathcal O(-1)$ である。
$u$ の乗法は隣接固有束を基底の $\mathcal O(1)$ でひねる四つの一次式となり、積は $F_4$ である。
根の上のファイバーでは $T$ が可逆ゆえ $k[u]/u^4$ の自由階数1加群で、$u$ の作用の階数は3である。
従って同じ根で二つの一次式は消えない。$F_4$ は四単純根を持ち $C$ が滑らかになってしまう。

\smallskip
\noindent\textit{三割線からの二つの自己写像の次数。}
$\tau=\sigma^{-1}$ は恒等または対合とする。前節の三割線の構成を使う。
正規化の種数2を用いたのは、二つの非定数自己写像が得られた\emph{後}のHurwitzの箇所だけである。
従って今も $r,s:N\to N$ は非定数で、一般の滑らかな $p$ に対し
\[
 D_p+r(p)\in |T+s(p)|,\qquad p+\tau p+r(p)\in |U+s(p)|,\qquad
 U=T K .
\]
ここで $D_p$ は長さ2の可動因子である。$U$ は非有効である。
有効なら $U$ と $\omega_C/U$ はそれぞれ高々一つの切断を持ち、一般の $s(p)$ で
$|U+s(p)|$ は固定因子と一点 $s(p)$ の和だけとなる。これは可動な二点 $p,\tau p$ を含めず、
$\tau=1$ のときも一般の $p$ を二回含めない。

非有効な次数2束 $U$ について、一般の滑らかな $x$ に対し $|U+x|$ は一点で、次数3の因子を定める。
これを $N\times N$ に閉じた対応を重複込みで $B_U$ とする。両射影次数は3であり、
転置は $U'=\omega_C/U$ から得る対応である。実際、滑らかな点対では
$H^0(U(x-y))\ne0$ と $H^0(U'(y-x))\ne0$ がSerre双対で同値である。
因子の重複も一致する。一般の点対を動かす局所環上で $R\Gamma(U(x-y))$ は
$H^0(U(x))\to U(x)|_y$ の二直線束の評価で表せ、その行列式が因子を切る。
相対Serre双対は転置側をその双対二項複体と同定するので消滅位数も同じである。
どの特異点 $z$ も一般の $|U+x|$ に入らない。
$\mathfrak m_z U$ は $h^0=0,\chi=-1$ で、そのSerre双対にはただ一つの独立切断があり、
一般点 $x$ で消えない。従って $h^0(\mathfrak m_z U(x))=0$ である。$U'$ 側でも同じで、
両方向で三つの点が一般の滑らかな部分で数えられる。

$m=\deg s$ とすれば、一般の $y$ が $|U+s(p)|$ に入る $p$ を両表示から数えて
$3m=2+\deg r$ を得る。$B_U$ は $(s(p),p)$ と $(s(p),\tau p)$ の像を含む。
異なる像なら $2m\le3$、同じで $\tau=1$ なら重複度が2以上なのでやはり $2m\le3$ である。
同じで $\tau\ne1$ なら $s\tau=s$ で $m$ は偶数かつ $m\le3$、従って $m=2$ である。
この場合一般の $|U+x|$ は $s^{-1}(x)$ と残る一点 $f(x)$ の和となり、
$r=fs$ と $\deg r=4$ から $f:N\to N$ は次数2である。
任意の正則双対化微分 $\eta\in H^0(C,\omega_C)$ を $N$ 上の有理微分と見ると
\[
                            f^*\eta+s_*\eta=\eta .
\]
これは動くCartier因子のAbel写像の微分である。滑らかな点で点が動くときの微分と
$\eta$ の対はその点での微分評価であり（局所パラメータの単純極の留数）、
複数の動点なら和、有限被覆の一ファイバーなら各逆枝の和すなわちtraceを取る。
標準平面の対合は非自明なため反不変な非零 $\eta$ があり、$s\tau=s$ より $s_*\eta=0$ である。

ところが種数0,1の滑らかな曲線の次数2の自己写像は非零の有理微分をそのまま保てない。
実際、各極の位数から1を引いた非負部分の合計は引戻しで2倍になる。
従って極があれば単純極だけで、その台は逆像で不変である。種数1の自己写像はエタールだからこれは不可能で、
正則微分に対する倍率も絶対値 $\sqrt2$ である（複素一様化、有限生成体を複素数に埋めればよい）。
射影直線では留数の和から極が少なくとも二つあり、Hurwitzの全分岐次数2から逆像不変な台は高々二点である。
二点を $0,\infty$ とすれば写像は $cz^{\pm2}$、微分は $dz/z$ の倍で、倍率は $\pm2$ である。
従って $m=2$ は消え、$r,s$ はともに自己同型である。

\smallskip
\noindent\textit{対合の反不変微分。}
$A=s^{-1},\ g=rs^{-1}$ とすると $B_U$ は $A,\tau A,g$ の三グラフの和である。
$g=1$ なら $D_p\sim T$ に反する。任意の $\eta\in H^0(\omega_C)$ に上と同じ微分を取ると
\[
             A^*(1+\tau^*)\eta+g^*\eta=\eta .
\]
標準平面の対合の座標で $x=v/w$ と置けば、関数体は
$F=u^4+a(x)u^2+b(x)=0$ で表せる。標準随伴公式による微分 $dx/F_u$ と
$x\,dx/F_u$ はともに反不変である。これは特異点でも平面の双対化束の切断である。
従って $g$ は両方を保ち、比 $x$ と $F_u$ を保つ。
しかし $F_u=2u(2u^2+a)$ と $F_u^2=4u^2(a^2-4b)$ より $x,F_u$ は全関数体を生成する。
判別式 $a^2-4b$ は整性から恒等零でない。よって $g=1$ となり矛盾する。

\smallskip
\noindent\textit{恒等四移動と正規化の分岐。}
$\tau=1$ では定理\ref{thm:fixed-differential}で $\deg M+4=0$、すなわち
正規化微分の零因子 $B$ の次数は $2g(N)+2$ である。
一枝の最小正位数を $e$ とすると $B$ の係数は $e-1$ で、全枝での和は
$\delta=3-g(N)$ を超えない。実際、正規化局所環でその枝の一様化元の第1から第 $e-1$ 乗の係数は
元の局所環からは得られず、別々の枝でこれらは独立な余長を与える。
従って $N=\mathbb P^1$、$\deg B=2$ である。
$B_U$ とその転置の三グラフ表示から、全双対化微分について
\[
      (2A^*+g^*)\eta=\eta,\qquad
      (2(A^{-1})^*+(g^{-1})^*)\eta=\eta .
\]
$A\ne1$ である。$s=1$ なら $2p+r(p)\sim U+p$ から $U$ が有効になる。

$B$ の台の点 $a$ に零点を持つ射影直線座標を $t$ とすると、$dt/t^2$ は $C$ 上の正則双対化微分である。
実際この微分は $a$ 以外で正則で、$C$ 上の任意の局所函数を掛けても
その枝の一次係数が零なので留数が零である。ここで有限正規化の留数表示
（各局所函数との積の枝留数の和が零であること）を用いた。
$e\ge3$ なら $dt/t^3$ にも同じことが成り立つ。
第一の等式の右辺に一つしか極がないため、$A,g$ は両方 $a$ を固定する。
異なる二点に極が出れば消し合えず、両方が同じ別の点に出れば $a$ の極が作れない。
$a$ での接倍率の逆数を $c,d$ とし、二つの等式の二重極係数を比較すると
$2c+d=1,\ 2/c+1/d=1$、よって $c=1,d=-1$ である。
$B$ の台が二点なら $A$ は異なる二点を固定しかつ接倍率1なので恒等になってしまう。
一点なら $e=3$ で、三重極を使うと $2c^2+d^2=1$ を要し矛盾する。これで全ての六次を除いた。
\end{proof}

\begin{corollary}\label{cor:integral-point-curves-covered}
点ファイバーで、一次元部分が整な被約曲線かつ有限部と交わらない標準AS窓の全体は、
滑らかな平面三次、滑らかな楕円四次、二類の滑らかな楕円五次、
および特異点を許す八係数の整七次の四行で尽きる。各行の正則性と添字保持同型の判定は
系\ref{cor:smooth-point-curves-covered}と定理\ref{thm:integral-septic-boundary}の通りである。
四行の次数は相異なるので行を越える同型はない。
\end{corollary}
\begin{proof}
次数上限は7であり、次数4以下、5、6、7をそれぞれ
定理\ref{thm:reduced-cubic-point-curves}、定理\ref{thm:reduced-quartic-point-curves}、
定理\ref{thm:singular-integral-quintic-exclusion}と本節および前々節の七次の定理で合わせる。
\end{proof}

% !TeX root = ../papers/nc-p3.tex
\section{埋込み点を含む点曲線の余長と高次数の純性}
\label{sec:pure-curve-finite-defect}
点ファイバーを持つ標準有限Frobenius窓で共通点グラフ $X$ が一次元とする。
$\mathcal E\subset\mathcal O_X$ を有限台の最大イデアル、$Y$ をその商の純曲線とし、
$d$ を片側次数、$g=p_a(Y)$ とする。孤立したArtin環も $\mathcal E$ の長さに含む。
\begin{theorem}\label{thm:pure-curve-defect}
重複・可約成分や埋込み点を許して
\[
                    \ell\mathcal E=18-4d+2g .
\]
$d=7$ では $g=5,\ \mathcal E=0$、$d=6$ では $g=3,\ \mathcal E=0$ である。
六次の各辺で一般の五つの関係を取ると、五式内の残余曲線は両射影で型 $(2,2)$ の四次完全交叉となる。
$d=5$ なら $(g,\ell\mathcal E)=(1,0),(2,2)$ のいずれかであり、
$d=4$ なら $g=-1,0,1$ に限られる。
\end{theorem}
\begin{proof}
定理\ref{thm:point-graph-curve-degree}の一般五式の完全交叉 $W$ と、$Y$ の残余 $Z$ を取る。
$\omega_W=\mathcal O_W(1,1)$、$I_{Y/W}(1,1)=\omega_Z$ である。
第六式の切断 $\mathcal O_Z\to\omega_Z$ は各一般点で全射である。
実際そこで $X$ と $Y$ は等しく、$W$ 上の $I_X$ は第六式が生成する。
$\omega_Z$ と $\mathcal O_Z$ は各一般点で長さが同じなので、この射はそこでは同型である。
$Z$ は純ゆえ全体で単射となる。その余核を $(-1,-1)$ でひねったものは
$I_{Y/W}/I_{X/W}=\mathcal E$ である。従って長さは $2p_a(Z)-2$。
$\chi(\mathcal O_W)=-10$、$Z$ の両側次数 $10-d$ と上のイデアルのEuler標数から
$p_a(Z)=g+10-2d$、主張の式を得る。

純な次数4の非平面空間曲線 $V$ は、非被約でも算術種数1以下であり、
種数1なら二つ以上の独立な二次式があることを示す。一般平面断面の四点スキーム $A$ への
斉次座標環の商の全射を使う。長さ $n$ の平面零次元スキームのHilbert関数は
$n$ に達するまで各次数で少なくとも1増加する。一般一次非零因子で割った標準次数代数の
Hilbert関数がその一階差で、一度零なら以後も零になるためである。
元の座標環を一般平面式で割った環の次数1は3、次数2は3以上、次数3以降は4以上である。
もし次数2が3なら元の二次式は三本以上ある。しかし純四次を含む独立二次式が共通因子を持たなければ
型 $(2,2)$ で既に全次数を占め、三本目はない。いずれの二本も共通因子を持つ場合は
三次元の二次式空間に共通の一次因子がある。一本を固定してその高々二つの一次因子で割れる
部分空間の和集合を使えばよい。三本が $\ell a_j$ なら $\ell\mathcal O_V$ は三独立一次式の
共通点に支持され、純性により零となり平面に入ってしまう。
よって次数2は4以上、総和から $g(V)\le1$、等号なら元の次数2の次元は8である。
特に算術種数3の純四次は平面に入り、この性質は同じ次数の他の埋込みにも及ぶ。

$d=6$ の四次残余は先の単射で種数1以上である。平面なら両側で平面四次だが、
五式の基点は曲線 $W$ であり、両平面の積で一般の三式が切る完全交叉の両次数は3である。
従って平面でなく種数1、$\mathcal E=0$、$\omega_Z=\mathcal O_Z$ である。
二つの二次式が共通一次因子 $\ell$ を持てば、$(\ell)$ はその残りの二一次式の共通直線 $T$ 上の
$\mathcal O_T(-1)$ の商で、非平面性と純性からこれに等しい。
この単射は $\mathcal Hom_Z(\mathcal O_T,\omega_Z)=\omega_T=\mathcal O_T(-2)$ を通るので不可能である。
従って両側で完全交叉となる。七次の主張は定理\ref{thm:septic-point-graph}による。

五次の場合は $g\ge1$ である。$Y$ に純な平面四次部分曲線は含まれない。
その部分曲線の算術種数3と先の上界、成分ごとの偏極次数一定性から両位置で平面四次となり、
平面の積での切断に反する。$Y$ 自身が平面五次となる場合も除く。
実際、非平面なら上と同じ平面断面の議論で、一般平面式で割った環の次数 $0,1,2,3,4,\ldots$ の
次元が少なくとも $1,3,3,4,5,\ldots$ なので算術種数は5以下で、平面五次の種数6にならない。
従って平面なら両位置で平面でありやはり不可能となる。

非平面五次で $g\ge3$ と仮定する。二次式がなければ上の次数2の下界が6となり、
総和から $g\le2$ で矛盾する。既約な二次曲面に入れば滑らかな場合は有効二次数の和5から種数は高々2、
錐なら純な因子の類は母線と二超平面の和で、前節の七次と同じHilbert多項式計算から種数2となる。
ここで純な余次元一のイデアルは深さの完全列により正規曲面上で反射的である。
残る二次式を $\ell m$ とし、二因子の順を選ぶと $(\ell)$ の次数 $e$ は2以下となる。
実際 $(\ell)$ は $\mathcal O_{Y\cap(m=0)}(-1)$ の商であり、$(\ell),(m)$ の次数の和は5以下である。
非平面性から $e=1,2$ で、$(\ell)=\mathcal O_{D}(-1)$ の $D$ は平面 $m=0$ の純な次数 $e$ の因子である。
反対の商 $Y\cap(\ell=0)$ の平面内の純部分は次数 $5-e$ の因子で、他の有限余長を $b\ge0$ とする。
完全列から
\[
             g=p_a(D)+p_a(\text{平面次数 }5-e)+e-1-b .
\]
$e=2$ なら右辺は2以下、$e=1$ で $g\ge3$ なら純な平面四次部分曲線を持ってしまう。
これで五次を得る。四次の純部分自体も同様に平面ではなく、上界1と余長の非負性から結論する。
\end{proof}

\begin{proposition}\label{prop:embedded-fixed-length}
$Y$ が被約な局所完全交叉曲線で、四移動がどの成分上でも恒等でないとする。
$Z_{\rm fix}$ を $Y$ 上の固定スキームとすると
\[
        \ell Z_{\rm fix}=2p_a(Y)-2+\ell (\mathcal E/\mathcal E^2)|_Y .
\]
右の有限層は孤立部では零である。この式は固定スキームがCartierでなくても成立する。
\end{proposition}
\begin{proof}
$P=\mathbb P^3\times\mathbb P^3$、$H=L_0L_1$、
$J=I_Y/(I_X+I_Y^2)=\mathcal E/\mathcal E^2$ と置く。
$R_0 H^{-1}\to I_Y/I_Y^2$ は階数6から5への写像で、余核が $J$ である。
その核を $S H^{-1}$ と書く。余法束は被約曲線の一般点で周囲の微分に単射なので
$S$ は微分の飽和核でもある。定理\ref{thm:fixed-differential}の反対評価切断は
$S L_2 L_{-1}$ に入り、その係数は $Y$ 上の固定イデアル $\mathfrak f$ を生成する。
実際後ろの辺を $Y\times Y$ に制限したファイバー積は $Y$ のグラフそのものである。
一般点でこの切断が核を張るので、飽和核は
$S=(L_2 L_{-1})^{-1}\mathcal Hom(\mathfrak f,\mathcal O_Y)$ である。
局所的には切断の有理倍率で六係数全てを正則にするものを取ったことになる。
Gorenstein曲線上の零次元層の局所双対により、$\mathcal Hom(\mathfrak f,\mathcal O_Y)$ は
$\mathcal O_Y$ を余長 $\ell Z_{\rm fix}$ で含む。余法束の行列式は $H^{-4}\omega_Y^{-1}$。
上の四項完全列の次数（無捩れ層ではEuler標数による次数）を取り、
$\deg L_2 L_{-1}=\deg H$ を代入すれば表示式となる。
\end{proof}

\begin{corollary}\label{cor:integral-hull-high-degrees}
曲線に重なる有限構造も許し、純部分 $Y$ が整な被約曲線であるとする。
標準AS窓で $d\ge5$ なら $X=Y$ で、楕円五次の二類と八係数の整七次だけが現れる。
$d=4$ なら $Y$ は種数1の $(2,2)$ 完全交叉で、節点・尖点や埋込み点を許しても
定理\ref{thm:finite-orbit-quartic-marks}が全次数の必要十分判定を与える。
\end{corollary}
\begin{proof}
六次・七次の純性と整曲線の分類を使う。五次では算術種数1なら純性が従う。
種数2のGorenstein曲線では定理\ref{thm:singular-integral-quintic-exclusion}の
二次曲面の残余母線から得る三次鉛筆の議論を適用する。
一般の滑らかな $p$ を取る二核束、$Q_2$ への非零制限、$D_p+r_p$ の計算は
有限台の追加構造を避け、全て $Y$ 上の切断と完全線型系のみによる。
従って四移動は二重被覆の対合で固定長は6である。前命題では
$\ell J=6-2=4$ を必要とするが、$J$ は長さ2の $\mathcal E$ の商なので不可能である。
非Gorensteinなら同定理の唯一の錐頂点は $X$ 上でも接次元3である
（$Y$ で既に3、$X$ は三空間に埋まる）。自動Ore表示の水平射影と極次数上限は
他の孤立点・埋込み点があっても同じである。曲面はなく、可動な水平二次曲線の上の高さは
一般点で一意なので非Gorenstein五次も消える。

整四次の算術種数は0または1である。種数0なら $Y=\mathbb P^1$ で、
$\nu^*\Omega_X$ の捩れ長は常に $\ell J$ に一致する。
実際 $Y$ が滑らかなら余法列は局所分裂し、微分の余核は $\Omega_Y$ に加えて $J$ を持つ。
非恒等な射影直線自己同型の固定長は2、恒等の場合は定理\ref{thm:fixed-differential}の第二項から、
どちらでも $\ell J=4$ となる。ところが $\ell\mathcal E=2$。
残る整な種数1の曲線は滑らかまたは唯一の節点・尖点を持ち双対化束は自明である。
ゆえに記載の四次印の定理を適用できる。
\end{proof}

% !TeX root = ../papers/nc-p3.tex
\section{種数一の被約連結成分と五角形の非存在}
\label{sec:genus-one-polygon-exclusion}
\begin{theorem}\label{thm:reduced-genus-one-components}
点ファイバーを持つ標準有限Frobenius窓の共通点グラフで、純な被約曲線が開閉な連結成分 $C$ をなすとする。
$p_a(C)=0$ は不可能である。$p_a(C)=1$ なら滑らかな楕円曲線か、滑らかな有理成分を
二つずつの節点で輪につないだ多角形である。後者では四移動が $C$ を保つときは固定点のない回転を誘導する。
$C$ が点グラフの唯一の曲線成分をなす、という仮定はここでいう連結成分の分類には不要である。
\end{theorem}
\begin{proof}
正規化の成分と特異点を頂点とし、各特異点への各枝を辺とする連結二部多重グラフを取る。
正規化の完全列から
\[
 p_a(C)=\sum g(\widetilde C_j)+b_1(\text{グラフ})+
            \sum_p(\delta_p-r_p+1)
\]
である。$r_p$ は枝数で、最後の和の各項は非負である。
実際、完備局所環は正規化の各枝の定数項だけを一致させた環 $A'$ に入り、
この項は $A'$ までの余長となる。余長零の特異点は $r$ 本の独立な座標軸の和であり、
埋込み次元上限から $r=2,3$ である。以下、$C$ が他の連結成分に移れば一致点を持たないことにも注意する。

木の末端にある有理正規化成分がただ一つの上記の座標軸特異点で会う場合を考える。
この枝では $B=0$ で余接捩れは $r-1$ である。局所式 $x_i x_j=0$ を一枝上で微分すればよい。
定理\ref{thm:fixed-differential}の恒等の場合は $r-1=4$ が必要であり、非恒等の場合に一致因子の長さは
$r-3$ となる。$r=2$ では負、$r=3$ では命題\ref{prop:automatic-spatial-tangent-mark}により
その点が四移動で固定されるため長さ零に反する。滑らかな開閉有理成分も同じ式で除ける。
従って種数零の木は消える。種数一で正規化に楕円成分があれば、木に他の成分があると
この楕円成分以外の末端を持つので、滑らかな楕円だけとなる。
$b_1=1$ の場合は全て座標軸特異点で、同じ末端排除により全頂点の価数は2以上、Euler式から全て2である。
従って節点の輪となる。これが四移動で保たれるとき各有理成分を点ごとに固定することはできず、
一致因子の長さは各成分で零である。向きの反転は節点か成分を保ち、後者でも射影直線上に固定点を持つ。
同様に一成分の自己節点の輪も、唯一の節点が固定されるため保たれる場合は生じない。

残る場合は木にただ一つ、$\delta_p-r_p+1=1$ の特異点 $p$ がある。
$p$ から遠い方へ辿れば、$p$ に直接付く成分の外には先に除いた末端しか作れない。
従って全成分は $p$ の一枝ずつに対応する有理成分で、他の特異点はない。
この場合 $C$ が保たれるなら $p$ は固定される。保たれない可能性については次数上限も使う。
$C$ とその像が相異なれば両方の次数は同じで合計7以下なので $\deg C\le3$ である。
次数2以下の純被約連結曲線の種数は零なので、双方とも次数3に限られる。
定理\ref{thm:pure-curve-defect}を使う。
すなわち次数6なら全純部分の種数3に対し二つの開閉楕円種数成分の和の算術種数は1、
次数7なら残る純直線を加えた種数は0で所要の5に達しない。
従って種数一の連結成分は保たれる（同じ議論は節点の輪にも適用できる）。

この唯一の特異点の環を明記する。$A\subset A'$ の余長が1なら $\mathfrak m_A A'\subset A$、
定数項が零なので $\mathfrak m_A$ は $A'$ のイデアルである。
$\mathfrak m_{A'}/\mathfrak m_A$ は一次元の加群なので $\mathfrak m_{A'}^2\subset\mathfrak m_A$。
従って定数を貼ったうえに一階係数に一つの線型条件を課すだけである。
係数の零となる枝は残りと値だけを貼る滑らかな軸である。この軸上の余接捩れは残りの環の接次元に等しく、高々2である。
実際、残りの最小座標 $y_j$ と軸の $t$ に対し追加の関係は $t y_j=0$ であり、残りの関係には一次項がない。
よってこの枝で恒等なら必要な長さ4を得られず、非恒等でも $p$ の固定により少なくとも3の捩れが必要である。

全係数が非零なら各枝の変数の倍率を変えて一階係数の和を零とできる。
$r=1,2$ は尖点と二枝接点であり、形式表示はそれぞれ $y^2=x^3$、$y(y-x^2)=0$ である。
$r\ge3$ ならこの環は $r-1$ 次元空間中の $r$ 本の直線の錐で、方向は射影枠をなす。
実際、和零超平面の成分ごとの積は各座標の二次項を張る
（相異なる $i,j,k$ に対し $(e_i-e_j)(e_i-e_k)=e_i$）。
以後全ての高次項を生成でき、接次元から $r\le4$ である。
$r=3$ は平面三重点、$r=4$ は平面内の枠四点を切る二つの二次式の錐である。
尖点では $(\deg B,\deg T)=(1,3)$、他の三型の各枝では $(0,2)$ である。
接点と三重点では平面方程式の勾配の共通位数を取る。四枝では各方向で二つの勾配は独立で、
枝でともに一位となる。いずれも $\deg M+\deg T=0$、$\deg M<0$。
四移動が枝上で恒等なら $2\deg M+\deg T=0$ に反し、それ以外では一致長零に固定点 $p$ が反する。
\end{proof}

\begin{theorem}\label{thm:pentagon-node-rank}
点ファイバーを持つ標準有限Frobenius窓の点グラフは、開閉な空間五角形を含まない。
特に被約な種数一の純部分を持つ五次点グラフの標準AS窓は滑らかな楕円五次の二類だけである。
\end{theorem}
\begin{proof}
五頂点を隣接順に番号付けし、全位置でグラフ対応により同じ番号を使う。
各位置で任意の四頂点は独立である。四点は巡回的な連続四点であり、
同一平面に入ればその両端の非隣接二辺が交わってしまう。
従って各位置で頂点を同じ射影枠 $v_0,\ldots,v_4$ に取り、代表を $\sum v_j=0$ とできる。
各直線上の隣接埋込み間の写像は
$[s v_j+t v_{j+1}]\mapsto[s v_j+q_j t v_{j+1}]$、$q_j\ne0$ である。
五倍率の一致や添字方向の周期を仮定しない。

この辺の関係 $B$ に対し $N_{jk}=B(v_j,v_k)$ は行和・列和が零で、
$N_{jj}=0,\ N_{j+1,j}=-q_j N_{j,j+1}$ を満たす。
逆にこれらの行列が六次元をなす。隣接五成分 $N_{j,j+1}$ を任意に取ると、
非隣接の十成分は行と列を頂点とする二部十角形の辺に対応する
（行 $j$ は列 $j\pm2$ に接する）。各頂点での和を指定する条件は一つだけ従属し、
行和総計と列和総計の一致で可解である。従って六関係はこの空間全体である。

$U_i$ をこの空間で行 $i$ が零の部分とする。片側点ファイバー条件から次元3である。
$U_0$ の列 $j$ での評価階数は $j=0$ で高々1、他の四点でちょうど2である。
計算を示す。$a,b,c=N_{12},N_{23},N_{34}$ を自由変数とする。行0が零なら
残りの非隣接成分は和条件から一意である（十角形を行0で切った木で和を順に解く）。
各列のうち行 $1,2,3$ を取れば十分で、その $(a,b,c)$ からの係数行列は順に
\[
\begin{gathered}
\begin{pmatrix}0&0&0\\q_1-1&0&1-q_3\\1-q_1&0&q_3-1\end{pmatrix},\quad
\begin{pmatrix}0&0&0\\-q_1&0&0\\q_1-1&q_2&-q_3\end{pmatrix},\quad
\begin{pmatrix}1&0&0\\0&0&0\\0&-q_2&0\end{pmatrix},\\
\begin{pmatrix}0&-1&q_3\\0&1&0\\0&0&0\end{pmatrix},\quad
\begin{pmatrix}-1&1&-q_3\\1&-1&q_3-1\\0&0&1\end{pmatrix}.
\end{gathered}
\]
最後の四行列には順に $-q_1q_2,-q_2,-q_3,1$ の二次小行列がある。
転置した辺では倍率を逆数にすれば同じ結論となる。

四次テンソル $w$ の位置0で節点 $p$ を評価する。位置1と2の関係を用いると
残る三次テンソルは $U_p\otimes V_3$ に入り、$V_3^*$ の縮約像は $U_p$ 全体である。
実際 $R_3$ と $R_1$ の完全対合のもとでこれは位置0での $R_3$ の評価の転置に当たり階数3である。
同様に位置2と3の関係を転置して使えば、位置3の前点 $j=\sigma^{-1}p$ での列が零になる
三次元部分全部を与える（位置3の消滅超平面は最後から最初への辺で決まる）。
従って位置2でも節点 $z$ を評価した $w(p,-,z,-)$ の階数は、
第一の表示では $z=p$ のときだけ1以下、第二では $z=j$ のときだけ1以下である。
よって全節点で $\sigma^{-1}p=p$ を要し、前定理の固定点を持たない回転に反する。
次数上限7から五角形自体は四移動で保たれる。

最後に五次で純部分が被約な種数一なら定理\ref{thm:pure-curve-defect}より有限部分は零である。
非連結なら次数2以下の被約な連結成分があり、種数零で除外される。
多角形の回転の位数 $k>1$ は成分の全軌道の長さであり、各軌道の成分次数は等しい。
よって $k\mid 5$、成分は五本の直線に限られ、今排除した。整の場合は
定理\ref{thm:elliptic-quintic-regularity}の二類で尽きる。
\end{proof}

% !TeX root = ../papers/nc-p3.tex
\section{直線成分の反対縮約による全次数AS判定}
\label{sec:line-component-as-test}
標準有限双対Frobenius窓で各辺の左右のファイバーが空か一点、共通グラフ $X$ が一次元とする。
$X$ の被約成分に直線があれば、定理\ref{thm:point-graph-curve-degree}により全位置で直線である。
重複した台の直線、埋込み点、他の任意の曲線成分を許す。
各位置の文字空間を $V_i$、直線成分 $Z$ の消滅二平面を $U_i(Z)\subset V_i$ と書く。

\begin{theorem}\label{thm:line-component-as-test}
任意の一本の直線成分の四移動 $\sigma$ による軌道 $\mathcal L$ を選ぶ。
位置 $i+2,i+3$ の二文字で四次テンソルを $Z\times Z$ に制限し、
直線の対角方程式 $s_0t_1-s_1t_0$ を除いた二文字関係を $r_i(Z)$ とする。このとき
\[
 0\ne r_i(Z)\in R_i\cap\big(U_i(\sigma Z)\otimes U_{i+1}(Z)\big).
\]
全次数で標準Koszul左右AS正則であるための必要十分条件は、
$i=0,1,2,3$ と全 $Z\in\mathcal L$ について $r_i(Z)$ が階数2を持つことである。
特に各位置で空間を張る既約曲線成分と一本の直線成分が共存すれば、
全ての有限Frobenius候補はAS正則である。
\end{theorem}
\begin{proof}[関係と低次の解消]
$Z$ をどの位置でも同じ $\mathbb P^1$ でパラメータ化する。最後の二文字の各関係の制限は
対角上で消える $(1,1)$ 形式で、その像は一次元である。零なら対応が直線の積を含み点ファイバーに反する。
$w$ の六対六対合から $r_i(Z)\ne0$ である。第三文字を $s$ で評価し、第二文字も $s$ で評価すれば
隣接関係により零となる。$t\ne s$ を選べば $r_i(Z)$ の第二文字は $U_{i+1}(Z)$ に入る。
同様に第四文字を $t$、最初を $\sigma t$ と評価すれば、最後から最初への関係で第一の包含も得る。

$D(Z)$ を直線の切断環とする。全位置で $\mathcal O(1)$ の二切断が得られるので、
$A\to D(Z)$ は全射で、$D(Z)$ は二変数多項式環の反復である。
定理\ref{thm:elliptic-quartic-as-test}の束複体を $\mathbb P^1$ に制限する。
内部次数 $n$ で項 $j$ の束は $\mathcal O(n-j)$ の $1,4,6,4,1$ 個の直和である。
点評価の階数は $1,3,3,1$ で局所分裂完全である。$n\ge3$ では全項の $H^1$ が零、
$n=2$ では項4に $H^1(\mathcal O(-2))=k$ が残り、二行の超コホモロジーの微分で
切断列の項2にちょうど一つのホモロジーを与える。$n=1$ では切断列自体が四文字から二切断への
全射なので項1に二つ、$n=0$ は項0のみである。反対側も同じである。
単純の極小解消の初めの二微分は候補と同一、三次の重なりも第三項に入るので、
実際の第一Torはこの計算と一致し、第二Torは計算した一元の商となる。
従って $D(Z)$ の始点 $i$ での右イデアルは $U_i(Z)A$ で、
その関係は次数2で高々一つ、左イデアルも各終点で同様である。
既に非零な $r_i(\sigma^{-1} Z)$ がその関係を与えるので必ず一つである。

\smallskip
\noindent\textit{必要性。}
標準ASなら候補は解消であり、直線商の左右極小解消は次数 $0,1,2$、個数 $1,2,1$ で終わる。
もし最後の二係数が比例するなら、右解消を双対化した第二Extの最小生成元に対する
左側の一次消滅空間は高々一次元である。一方、軌道長の四倍の周期に折り返せば、
定理\ref{thm:elliptic-quartic-as-test}の局所双対計算がそのまま適用され、
この第二Extは標準四移動を除いて $R^2\Gamma_A(D(Z))'$ に一致する。
その作用はひねられた直線商を通るので一次消滅空間は少なくとも二次元である。
実際ここで使う直線商の左右解消は有限自由であり、$A/A_{\ge N}$ との正Torは
内部次数 $[N,N+1]$ に限られる。従って同定の証明で用いた入射加群の非輪状性も同様に成立する。
これで階数1は排除される。

\smallskip
\noindent\textit{十分性。}
全整数で $Z_{i+1}=\sigma^{-1}Z_i$ とし、二文字平面 $U_i(Z_i)$ と
非退化関係 $r_i(Z_{i+1})$ で抽象的な部分表示 $P$ を作る。各位置で順に基底を取れば
関係が交換子となり $H_P=(1-t)^{-2}$ である。
$U_i(Z_i)A$ による右商は $D(Z_i)$ の始点部分、$A U_{j-1}(Z_{j-1})$ による左商は
$D(Z_{j-1})$ の終点部分で、いずれも長さ $n$ で次元 $n+1$ である。
各商の基底を持ち上げて $P$ の語を前または後ろに掛ければ $A$ を張る。
従って各成分の次元は $(1-t)^{-4}$ の係数以下である。一方、右表示の唯一の関係で始めた列のEuler式は、
固定終点で $a_n-2a_{n-1}+a_{n-2}=n+1+e_n$（$e_n\ge0$ は末項の核次元）である。
$(1-t)^{-2}$ を掛ければ上界で等号、全 $e_n=0$ である。左右の張る集合は基底で、
$A$ は $P$ 上左右自由、始点 $i$ の直線商加群は $P$ 単純の誘導となる。
任意の所定の $i,Z$ に対し $Z_i=Z$ としてこの構成を行えるから、同じ $D(Z)$ の全始点で
$\underline{\operatorname{Ext}}^j_A(e_iD(Z),A)=0$ ($j<2$) を得る。終点側も同じである。
直線商を通る環変更で左右の単純から $A$ への最初の二Extも消える。
生成と二次表示で候補複体の始めは完全であり、上の消滅とFrobenius双対で末尾も完全、
中央は既に求めたHilbert関数で完全である。従って全次数のKoszul性と左右AS性が従う。

最後に $r_i(Z)=a\otimes b$ ならグラフの各既約成分で最初の線型形式 $a$ または二番目の
線型形式 $b$ が恒等零となる。両位置で空間を張る同じ既約成分があれば不可能である。
\end{proof}

\begin{corollary}\label{cor:line-component-atlas}
直線を含む一次元点グラフの全AS窓は、次の有限パラメータと同型判定で被覆される。
最初の三位置の指定直線を同じ座標直線にして、その上の隣接写像を恒等とする。
二つの辺ごとに、直線上の切断乗法核である十三次元空間の六平面を選ぶ。
二窓からの有限Frobenius復元、四辺の点ファイバー、一次元の像と上の軌道の階数検査を課す。
指定を忘れた同型は、第二の窓の有限個の直線成分のどれかを指定して、
三位置の線型変換によって二つの関係を同時に運べることと同値である。
これらの変換は当該直線を対応させ、そこで同じ $\mathrm{PGL}_2$ を誘導するものに限れば十分である。
同じ比較を一移動・二移動に用いれば最小周期を得る。
\end{corollary}
\begin{proof}
全窓で一つの初期直線は点グラフによって一意に運ばれる。成分ごとの次数一定性で
それは各位置の直線で、全次数は高々7だから軌道探索は有限である。
逆に最初の二辺に指定直線のグラフが入り、復元・ファイバー・次元検査を通れば、
共通像の既約成分としてその直線が入る。二窓の延長は系\ref{cor:two-window-extension}による。
座標変更の存在はこの三つの放物部分群と共通の二次行列を用いた有限な多項式条件で判定できる。
\end{proof}

% !TeX root = ../papers/nc-p3.tex
\section{平面二次曲線成分の一次・二次生成元による全次数判定}
\label{sec:conic-component-as-test}
同じ点ファイバーの仮定で、全位置で平面二次曲線になる純な部分スキーム $Z\subset X$ を印とする。
四移動によるその軌道 $\mathcal C$ が有限とする。特に滑らかな二次曲線成分はこの条件を満たす。
$D(Z)$ は全切断商、$I(Z)$ はその核である。平面式を $x_i(Z)\in V_i$ とし、スカラー倍は任意とする。

\begin{theorem}\label{thm:conic-component-as-test}
各 $Z\in\mathcal C$ と位置 $i$ で
$y_i(Z)\in I_{i,i+2}(Z)\setminus x_i(Z)V_{i+1}$ を選ぶ。
標準Koszul左右AS正則性は次の有限条件と同値であり、どの持上げ $y_i$ を使ってもよい。
軌道の置換 $\pi$ が存在し、全 $i,Z$ について $Z'=\pi Z$ と書くと
\[
 x_i(Z) a + y_i(Z) b=0\quad\text{in } A_{i,i+3},\qquad
 0\ne b\in k x_{i+2}(Z'),\quad
 a\in I_{i+1,i+3}(Z')\setminus V_{i+1}b
\]
を満たす係数がある。検査する位置は $4|\mathcal C|$ 個で十分である。
\end{theorem}
\begin{proof}
二次曲線の各偏極はその平面の $\mathcal O(1)$ である。滑らかなら $\mathbb P^1$ の $\mathcal O(2)$、
二直線なら各直線上次数1で木の貼合せ、二重直線なら直線上の冪零部分 $\mathcal O(-1)$ の
$H^1$ が零なので、いずれも全位置で同じ束 $L$ となる。
平面での一つの二次式の解消から
$\omega_Z=L^{-1}$、$h^0(L^n)=2n+1$ ($n\ge0$) で、切断環は一次生成である。
$D(Z)$ はこの環の反復となる。$I_{i,i+2}$ の次元は5である。

直線の場合と同じ局所分裂束複体を用いる。$n\ge4$ なら全項の $H^1$ は零。
$n=3$ では項4に $H^1(L^{-1})=k$ があり、項2の切断ホモロジーを一元与える。
$n=2$ では項4から項3への $H^1$ の射が三次元から四次元への単射である。
その双対は四文字から三つの平面切断への全射である。従って項1の切断ホモロジーだけが一次元残る。
$n=1$ では平面式一本が項1に残る。$n=0$ は単純のみである。
よって右の第一Torは次数1,2に一元ずつ、第二Torは次数3の一元の商であり、反対側も同じ。
特に $x_i V_{i+1}$ は四次元で、右イデアルは $x_i,y_i$ で生成され、関係は次数3に高々一つ。
左イデアルも終点に応じて同様である。

標準ASならこれが長さ2の極小解消となる。最終列を $(a,b)^t$ として双対化する。
定理\ref{thm:elliptic-quartic-as-test}の局所双対は軌道を折り返した同じ証明で適用でき、
第二Extの最小生成元は左から四移動で運ばれた印 $Z'$ のイデアルで零化される。
正Torの遷移消滅には $[N,N+2]$ の内部次数範囲を使えばよい。
従って次数1において $x_{i+2}(Z')$ は $b$ の倍数、次数2では
$V_{i+1}b+k a$ が五次元の $I_{i+1,i+3}(Z')$ を含む。以上から表示の全条件を得る。
置換はねじれまたはその逆が全軌道に誘導するものを取る。
$y_i$ を非零倍と $x_i c$ の加算で取り替えても、$a$ の変化は倍率と $V_{i+1}b$ の元だけである。

逆に有限条件を仮定する。$b\ne0$ より表示の関係は上で得た高々一つの最小関係を実現する。
固定終点で $h_n=\dim A_{j-n,j}$、$d_n=2n+1$ ($n\ge0$、負なら零) とおく。
右表示のEuler式は
\[
 h_n-h_{n-1}-h_{n-2}+h_{n-3}=d_n+e_n,\qquad e_n\ge0 ,
\]
ここで $e_n$ は末項の核次元である。一方 $m_n(Z)=\dim(A_{i,j}/x_i(Z)A_{i+1,j})$、$i=j-n$ とする。
この商から $D(Z)_{i,j}$ への核は $y_i(Z)A_{i+2,j}$ の像であり、
表示の関係により $y_i(Z)b A_{i+3,j}$ は消える。従って
$m_n(Z)\le d_n+m_{n-2}(\pi Z)$ である。$h_n\le h_{n-1}+m_n(Z)$ と合わせると上界は
\[
 \frac{H_D}{(1-t)(1-t^2)}=(1-t)^{-4}.
\]
先のEuler式に同じ非負係数の逆多項式 $1/((1-t)(1-t^2))$ を掛ければ下界もこれに等しい。
よって $e_n=0$、各不等式で等号、$x_i(Z)$ の左乗法は単射である。
反対側でも同じ上界の議論が通る。実際、終点 $i+3$ で $a,b$ は $Z'$ の左イデアルの二次・一次の
最小生成元であり、$x_i(Z) a=-y_i(Z)b$ を使い、置換を逆向きに辿ればよい。
既知の $H_A$ はこの上界と等しいから、全ての $x_i(Z)$ の右乗法も単射である。

右の長さ2の自由解消を双対化すると、最初の写像は $x_i$ の右乗法の単射性から単射で、
最後の余核は $a,b$ が生成する左イデアルによる商、すなわち $D(Z')$ の終点部分である。
中央のホモロジーも先のEuler等式（内部次数を末項に合わせる）で消える。
従って右切断商加群 $e_iD(Z)$ からの最初の二Extは零、反対側も同じ。
既出の環変更・Frobenius双対・Euler式で標準Koszul左右AS正則性が従う。
\end{proof}

\begin{corollary}\label{cor:conic-component-atlas}
滑らかな二次曲線成分を含む全窓は有限条件付き二窓アトラスを持つ。
同じ三つの座標平面に $[s^2:st:t^2:0]$ を取り、二辺で対角の二次曲線を消す
次元11の双一次式空間から六平面を選ぶ。
系\ref{cor:line-component-atlas}の二窓復元・点ファイバー・次元の検査と本定理を用いる。
指定を忘れた同型は第二窓の各二次曲線成分を順に標準化し、最初の三位置で
曲線上同じ $\mathrm{PGL}_2$ を誘導する平面安定化群の三変換により
二辺が同時に移るかで判定できる。一移動・二移動と同じ比較をすれば最小周期が決まる。
\end{corollary}
\begin{proof}
二平面の二次曲線切断の積は $H^0(\mathcal O_{\mathbb P^1}(4))$ 全体を張る。
六平面の周囲の次元は従って $16-5=11$ である。全ての曲線成分は位置ごとの次数を保ち、
二次曲線上の写像はその完全二次線型系により平面の射影変換に延びる。
残りは二窓の延長と有限個の成分の列挙による。
\end{proof}

\begin{corollary}\label{cor:full-quintic-point-window-atlas}
点ファイバーを持つ五次点グラフの全標準AS窓は、非被約構造と埋込み点を含め、
次の互いに素な三行で被覆される。
\[
\begin{array}{c|l}
\text{被約台の成分}&\text{全次数AS性と添字保持同型}\\ \hline
\text{直線あり}&\text{系}\ \ref{cor:line-component-atlas}\\
\text{直線なし、二次曲線あり}&\text{系}\ \ref{cor:conic-component-atlas}\\
\text{両方なし}&\text{滑らかな楕円五次の二類、定理}\ \ref{thm:elliptic-quintic-regularity}
\end{array}
\]
最初の二行では純部分の次数5を課す。これは射影像の一次のHilbert多項式の最高係数で検査できる。
種数二の可約五次での長さ二の有限構造も個別の接型を指定せずこの判定に入る。
さらに任意の次数について、直線・二次曲線も整な純部分もない残りの曲線台は、
三次二成分、三次と四次の二成分、または一本の整三次上の二重構造だけである。
\end{corollary}
\begin{proof}
各成分の一般長と次数の積の和が片側次数である。次数1の整曲線は直線、次数2は滑らかな平面二次曲線である。
これらがなければ二成分以上または一般長2以上には次数6以上を要し、
総次数上限7の下では記載の三通りに尽きる。一成分で一般長1なら純部分の冪零イデアルは
有限台なので零であり、系\ref{cor:integral-hull-high-degrees}を適用できる。
最初の二行では指定を忘れる際、同一窓の複数の印も同型判定で同一視される。
\end{proof}

% !TeX root = ../papers/nc-p3.tex
\section{点グラフ内の任意の平面三次印からの中心一次表示}
\label{sec:plane-cubic-mark-normality}
\begin{theorem}\label{thm:plane-cubic-mark-normality}
標準有限Frobenius窓で左右が点ファイバー、共通グラフ $X$ が一次元とする。
一位置の像が平面内に純な三次部分曲線 $C$ を含めば、全位置のその平面を消す一次式は正則正規列になる。
$C$ の重複成分、その上の埋込み点、外の余分な曲線や有限構造を許す。
全窓は定理\ref{thm:regular-normal-linear}の三生成三次元AS商の中心一次PBW持上げとなり、
全次数で標準Koszul左右AS正則、最小周期1である。
逆に三生成商の点対応が平面三次のグラフとなる同定理の全PBW持上げで、
四文字の各辺が点ファイバーかつ一次元ならこの枝に入る。
\end{theorem}
\begin{proof}
$C$ は純で $p_a=1$ であり、全位置で次数3であるから
定理\ref{thm:septic-point-graph}の残余三次の証明により各位置で平面三次になる。
その平面方程式を $t_i$、平面上の偏極を $L_i$ とする。
$H^0(\mathcal O_C)=k$、$\omega_C=\mathcal O_C$、全切断が埋込みを与える。
隣接する三切断系の積の像は六次元である。同じ束なら平面に二次式がないことから従う。
相異なるなら基点自由な鉛筆の列で $H^1(L_{i+1}/L_i)=0$ を使う。
この消滅は補題\ref{lem:genus-one-quartic-mark}と同じく、
多重次数零の束に切断があれば自明であることによる。実際、真の純部分曲線は平面で次数2以下の因子で
$H^1(\mathcal O)=0$ であり、同補題の部分曲線への双対性の証明をそのまま使える。
偏極が等しければ三次曲線に二次式がないため平面対角全体が入ってしまう。
六関係の隣接平面積への制限は次元3以下で、二式以下なら曲面が残るためちょうど3である。
この三式は平面積で完全交叉、既にこれに入る $C$ の両側次数が3なので完全交叉自体が $C$ である。

定理\ref{thm:plane-cubic-point-graphs}の記法で、平面上で零になる三関係を
$t_i B_i Y-t_{i+1} A_i X+c_i t_it_{i+1}$ と書く。
$A_i$ の核に射影点 $p$ があれば $B_{i-1}$ も特異となることを示す。
$p$ が三次曲線外なら制限三式の平面上での評価が可逆であり、次点も前点も平面外となる。
実際、混合式は次点に何も課さず、残りの三式は独立なので唯一の次点がある。
前点の存在は両側グラフの同定による。そこで一つ前の混合式で $B_{i-1}p=0$ となる。

$p\in C$ の場合は、曲線の滑らかさや、$X$ がそこで $C$ と一致することを要しない。
局所環 $\mathcal O_{X,p}$ で $C$ のイデアルを $J$ とし、余法方向の最小生成元空間
$T=J/\mathfrak m_pJ$ を全位置で同一視する。局所自明化した $t_i,t_{i+1}$ は $T$ を張り、
この二元の線型関係はちょうど三混合式の二列 $B_i j_{i+1}(p),-A_i j_i(p)$ の行空間である。
実際、平面積上の三式の持上げと二つの平面式は、周囲の滑らかな六次元空間で
$C$ のイデアルの長さ5の正則列である。極小生成元空間でまず持上げ三元を除き、
残る二元に混合三式の定数係数だけが作用する。これは $J/\mathfrak m_pJ$ の表示そのものである。
$A_i j_i(p)=0$ ゆえ、他方の列が非零なら $T$ は一次元で $t_i=0$、
他方も零なら $T$ は二次元である。前の辺の二元 $t_{i-1},t_i$ も同じ $T$ を張る。
前者の場合 $t_{i-1}\ne0$ なので混合式が $B_{i-1}j_i(p)=0$ を強制する。
後者なら二元は独立で、やはりこの列は零である。
以上どちらでも $B_{i-1}$ は特異となる。
左右核から $u\otimes t_i\in R_{i-1}$、$t_i\otimes v\in R_i$ なる非零純関係が生じ、
三次交叉に純テンソルが入る。しかし四次対合と点ファイバー条件から
どの非零三次交叉も両端支持が少なくとも3である。矛盾により全 $A_i$、逆向きに全 $B_i$ は可逆である。

商三関係 $S_i$ の非零元は階数1を持たない。もし $a\otimes b$ があれば $C$ の各一般点で積が零となる。
各平面の非零線型式は高々一つの直線成分だけを含み、その一般点の
$k(C_{\rm red})[u]/(u^m)$ で消せる重複は一層だけである。平面三次の局所式が
その直線の方程式の $m$ 乗と単元の積だからである（$m=1$ ではその一層全部）。
成分ごとの次数は次位置でも等しく、二式で消せる成分次数の合計は重複込みで高々2、三次全体は消えない。

商三次交叉の非零テンソルの第一支持が二次元以下なら、それを消す非零文字評価 $p$ を選ぶ。
末尾の全縮約は $S_i$ の $p$ での評価核に入るが、この核は点ファイバー条件から高々一次元である。
よって末尾支持が一次元となり、先頭縮約で $S_{i+1}$ に非零な階数1が生じる。従って両端支持は3である。
このようなテンソルは三文字の一端と三次元関係空間の間で同型を与えるので、
交叉は高々一次元である（代数閉体上で可逆行列だけの射影直線はない）。
他方、混合行列の可逆性から $t_i V_{i+1}=V_i t_{i+1}$ は四次元で、
三文字商の三次元数は $20-10=10$ 以上である。従って交叉はちょうど一元、商三次元数は10であり、
$t$ の二次から三次への左右乗法は単射である。一次の可逆な移送を関係に二回用いれば
この単射性で二次関係空間そのものを次へ運ぶ。従って中心表示への同定、
補題\ref{lem:ternary-finite-graph-test}による商の正則性、三次PBW条件からの全次数正則性は
定理\ref{thm:plane-cubic-point-graphs}と同じである。
逆向きは中心PBW表示を平面対に制限すればよい。
\end{proof}

\paragraph{指定を忘れた同型。}
定理\ref{thm:regular-normal-linear}の有限係数式でPBW補正と指定平面の同型を判定する。
平面印の総体は各位置での平面をGrassmann変数に取り、二つの連続する平面積上で
三式の完全交叉三次グラフとその共有中間像の一致を課して記述できる。
指定を忘れた比較には、第二窓で第一側の印に対応する印を許す。
厳密には二窓延長の系\ref{cor:two-window-extension}により、両関係を運ぶ三位置の座標変換を検査すれば十分で、
第一の印が移った後の平面はその変換で定まる。従って二重直線や三重直線を含む場合に印の個数の有限性を仮定しない。

% !TeX root = ../papers/nc-p3.tex
\section{有理正規三次上の六次二重曲線の全排除}
\label{sec:double-twisted-cubic-exclusion}
\begin{theorem}\label{thm:double-twisted-cubic-exclusion}
一次元の点ファイバーグラフを持つ標準有限Frobenius窓において、純部分が
空間有理正規三次の上の重複度2の曲線となることはない。
\end{theorem}
\begin{proof}[二重曲線の抽象型]
定理\ref{thm:pure-curve-defect}より点グラフ $X$ 自体が純で $p_a=3$ である。
被約台を $C=\mathbb P^1$ とし、そのイデアルを $N$ とする。
$N^2$ は一般点で零、純性から全体で零である。$N$ は滑らかな $C$ 上で階数1無捩れ、
Euler標数から $N=\mathcal O_C(-4)$ である。このような二重曲線（リボン）の抽象型は二つしかない。
二つのアフィン座標直線では各パラメータを持ち上げて構造層の商写像を環として分裂できる。
重なりでの分裂の差は $N$ に値を持つ導分で、変更を除いた全体は
$H^1(T_C\otimes N)=H^1(\mathcal O(-2))=k$、$N$ の倍率を許すと零と非零の二型である。

分裂型では引戻し $J=f^*\mathcal O(1)$ を持つ有限平坦二重被覆 $f:X\to C$ がある。
その相対双対は $\operatorname{Hom}(\mathcal O\oplus\mathcal O(-4),\mathcal O(-2))$ なので
$\omega_X=J^2$ である。各 $L_i$ の制限は $\mathcal O_C(3)$ で、$N L_i|_C=\mathcal O(-1)$ の
コホモロジー消滅から全四切断になる。$L_i/J^2$ が非零切断を持てば、冪零部分の次数は $-3$ なので
一般点で非零因子であり、長さ2のCartier零因子を持つ。その上で $L_i$ の三切断が零になり埋込みに反する。
ゆえに $f_*L_i=\mathcal O(1)^2$。$P_i=L_i/J$ の二切断は各長さ2ファイバーを張る。
これで定理\ref{thm:smooth-sextic-exclusion}の超楕円分裂型の証明を適用できる。
同じ $P_i=P_{i+1}$ ならば二鉛筆の積の滑らかな二次曲面の対角が入る。
実際 $X$ を含む二次曲面はこの場合一意で、独立な二式は次数6を切れず、
共通平面因子を持てば非平面な台が直線に入ってしまう。
$P_i\ne P_{i+1}$ の場合の必要な次数零切断の消滅も、冪零部分の負次数と
$C$ 上の定数を用いれば従う。他は有限平坦な各長さ2の環の乗法だけを使うので同じ三次元数上限18を得る。

非分裂型は平面内の二重円錐曲線 $(xz-y^2)^2=0$ である。
実際 $x=1$ の座標を $t=y/x,e=z/x-t^2$、$z=1$ の座標を $u=y/z,e'=x/z-u^2$ とすると
\[
                 e^2=0,\qquad u=t^{-1}-e t^{-3},\quad e'=e t^{-4}.
\]
この分裂の差は $e\partial_t$ の自明化で $t^{-1}$ に対応し、上の $H^1(\mathcal O(-2))$ の生成元である。
随伴によりこの平面埋込みが完全標準系であり、全自己同型は標準平面の二次曲線安定群 $\mathrm{PGL}_2$ である。

\smallskip
\noindent\textit{四移動と全切断系。}
$1+N$ の単元完全列より $\operatorname{Pic}^0 X=H^1(N)$ は三次元ベクトル群である。
$H^1(\mathcal O(-4))$ への $\mathrm{PGL}_2$ 作用は二項二次形式の三次元表現である。
実際、Serre双対で $H^0(\mathcal O(2))$ の双対、$N$ の線型化は
標準平面二次曲線の余法束によるものなので、半単純群に指標の違いもない。
被約次数3の不変な参照束として、$\omega_X^3$ の一意な平方根を取れる
（被約上の平方根を持ち上げてベクトル群で2を割る）。これとの差を $l_i\in H^1(N)$ と書く。
偏極三階差の消滅から $l_i$ は $i$ の二次以下の多項式であり、$l_{i+4}=\sigma^*l_i$。
$\sigma$ が半単純なら各固有成分でこの等式を見て多項式は定数である。
二次曲面に入らないことは下記の座標でも示すので、この場合全隣接関係が交換子となり点グラフは全空間になる。
恒等もここに含める。
非自明ユニポテントなら同じ一径数加法群内の四乗根を $\tau$ として
$l_i=(\tau^*)^i l_0$ となる。実際右辺も三次元表現の冪零指数から二次以下の多項式で、
$0,4,8$ で一致する。従って全四切断系は一つの系と連続する $\tau$ の移動で与えられる。

各束 $L$ の遷移函数は二つのアフィン座標で
$t^3(1+e(a/t+(b+3/2)/t^2+c/t^3))$ と書ける。
三係数は $H^1(\mathcal O(-4))$ のLaurent基底であり、$3/2$ は基準の $\omega_X^{3/2}$ による。
$t^j$ の一意な切断持上げは $t^j+e g_j(t)$ で
\[
 (g_0,g_1,g_2,g_3)=(0,a,at+b+1/2,at^2+(b+3/2)t+c).
\]
これは $u^{3-j}$ を先の遷移式で書き、$e$ の係数の非負べきだけを残せば得られる。
三次曲線の三つのHankel二次式に代入した $e$ 係数は、斉次化して
$sU,tU+sV,tV$、$U=(b+1/2)s-at$、$V=cs+(1/2-b)t$ である。
従って埋込みであるためには
\[
                          1/4-b^2+ac\ne0
\]
が必要である。さもなければ共通零点の法方向が三次曲線の接方向に入る。
この不等式で三つの二次係数は独立で、二重曲線を含む二次曲面はない
（三係数の行列式が同じ $1/4-b^2+ac$ である）。

二つの隣接偏極が異なる場合にも、六関係は全切断乗法核である。
被約台での七次元の二文字像に加え、冪零部分の $H^0(\mathcal O(2))$ も全て得られる。
二つの束の上記表示の係数差を $A,B,C'$ とする。交換子の $e$ 係数には
\[
 A,\quad At+B,\quad At^2+Bt+C',\quad Bt,\quad Bt^2+C't,\quad C't^2
\]
があり、三係数の少なくとも一つが非零ならこれらは二次以下の多項式を張る。
従って像次元10、核次元6である。同じ偏極の場合にも先の二次係数で同じ結論となる。
従って $\tau$ で系を運べば全辺は同一の通常関係であり、$j_i=j_0\tau^i$ と選べる。
$\tau^4=\sigma$ なので、この通常表示での位置4から位置0へのFrobenius同定はスカラーである。
実際その射影変換は両側の $j_4=j_0\sigma$ の同じ座標を同一視し、非退化な被約三次を点ごとに固定する。
従って高さ4のテンソルは通常の四位置巡回の固有ベクトルである。固有値を先に指定する必要はない。

\smallskip
\noindent\textit{巡回テンソルの全係数排除。}
標準平面の座標を取り $\tau(t,e)=(t+1,e)$ とできる。四次テンソルを
四つの変数 $t_j$ ($j$ は法4) で各次数3以下の多項式 $w$ と見る。
全辺で $t_{j+1}=t_j+1$ の三次曲線のグラフ上に消えるため
\[
                  w=\prod_j(t_{j+1}-t_j-1)\, f(t_0,t_1,t_2,t_3).
\]
$0\ne f$ は各次数1以下であり、前の積が巡回不変なので $f$ も巡回固有である。
最初の辺がリボン全体を消すための一次の必要条件は
\[
       [(D_0+D_1)w]_{t_0=t,\ t_1=t+1}=0,\qquad D(t^k)=g_k(t).
\]
これは他の二位置を被約台に制限しても全四文字を張るので係数恒等式となる。

使用した消去を再現できる形で示す。$\{0,1\}^4$ の巡回軌道を辞書順に取り、
固有値の候補 $z=1,-1,\mathrm i,-\mathrm i$ に対し長さ $h$ が $z^h=1$ を満たす軌道のみ採用する。
その単項式を順に $z^j$ 倍して足した列を使う。各列に四つの括弧を掛け、位置0,1の指数
$k$ を $m=1,2,3$ だけ下げる係数を順に $2a,2b+2k-3,2c$ とする（$k<m$ は零）。
$(t+1)^d$ を展開して $t,t_2,t_3$ の係数を並べれば上式の二倍である。
軌道係数の最初の非零項を1にした各射影チャートで、この行列の消滅イデアルを
$\mathbb Q(\mathrm i)$ 上で計算した。$z=-1,\mathrm i$ の全 $4,3$ チャート、および
$z=1$ の最初以外の五チャートは単位イデアルである。共役で $-\mathrm i$ も同じである。
$z=1$ の最初のチャートのイデアルは $a^2,(b+1/2)^2$ を含む。
この計算は\path{computations/verify_double_twisted_cubic.py}が全行と各Gröbner基底を構成し、
二多項式の剰余と単位イデアルを検証する。最後の場合も $a=0,b=-1/2$ となり
埋込みの不等式に反する。これで全ての二重構造を排除した。
\end{proof}

% !TeX root = ../papers/nc-p3.tex
\section{空間三次と四次の七次曲線の排除および七次点グラフの被覆}
\label{sec:reducible-septic-coverage}
\begin{theorem}\label{thm:spatial-cubic-quartic-septic}
点ファイバーを持つ標準AS窓の七次点グラフは、空間有理正規三次 $C$ と整な四次 $D$ の和にはならない。
\end{theorem}
\begin{proof}
七次の純性により点グラフは被約な $Y=C\cup D$ そのものであり $p_a(Y)=5$ である。
平面四次は他の全位置でも算術種数3により平面に入り、隣接平面積の三双一次式の次数上限3に反する。
そこで $D$ は非平面、$h=p_a(D)=0$ または1である。$T=C\cap D$ は長さ $6-h$ で、
$C$ 上では有効Cartier因子である。$h=0$ なら $D$ は滑らかな有理曲線、
$h=1$ なら滑らかな楕円曲線または一点が節点・尖点の種数一曲線で、いずれもGorensteinである。

まず $Y$ に接次元3の点 $v$ がある場合を除く。二つの滑らかな枝の和の接次元は高々2である。
実際、二つのパラメータを交点スキーム $k[[u]]/(u^m)$ 上で一致させると、両局所環のファイバー積は
$k[[u,y]]/(y(y-u^m))$ となる。従って $v$ は $C$ と特異な $D$ の二重特異点との交点に限る。
命題\ref{prop:automatic-spatial-tangent-mark}のOre表示を使う。この点からの水平射影で
$C$ の像は平面二次曲線、$D$ も同様である。後者は正規化上で四次の偏極から二重特異点の
基点次数2を引くことと、非平面性から従う。両者とも各位置でこの二次曲線に双有理に写る。
三生成基底が平面全体の非恒等グラフなら定理\ref{thm:plane-ore-unmarked}により点曲面が入る。
恒等なら持上げの高さを同時に動かせるため同じ矛盾になる。残る三生成基底の点曲線は平面三次であり、
二つの既約な二次曲線を含めず、両水平像は同じである。
そこで一般の水平点で $v(gp)$ と $v(p)$ が独立なら、混合式
$a v(gp)-b v(p)=d_i(v(p),v(gp))$ が両高さを一意に決め、異なる二成分が同じ持上げを占めてしまう。
比例する場合は、持上げのある各点で高さが一次元動き曲面となる。これで空間接点はない。

従って $Y$ は局所平面で、$T$ は $D$ 上でもCartierである。
実際、両成分の共通点では滑らかな局所曲面の中の二つの因子として交わる。
随伴または留数により $\omega_Y|_C=\omega_C(T)$、$\omega_Y|_D=\omega_D(T)$ である。
$L_i$ の多重次数は $(3,4)$ で、$\omega_Y/L_i$ の多重次数は $(1-h,h)$ となる。
Riemann--Rochと四つの非退化な文字から $h^0(\omega_Y/L_i)\ge1$。
次数零側への制限は単射である。そこで零になる切断は、他方の次数1の束の切断として長さ $6-h$ で消えねばならないからである。
次数零の整曲線上の非零切断は生成元となるので $h^0=1$。非零切断は他方でも恒等零でなく、
$T$ と交わらない長さ1のCartier零因子を持つ。従って滑らかな一点 $p_i$ により
$L_i=\omega_Y(-p_i)$ であり、全ての $p_i$ は同じ次数1側成分にある。四文字は完全系となる。

次数1側にのみ台を持ち $T$ を避ける等次数 $e\le4$ の二つの有効因子が線型同値なら、両者は一致する。
その比の有理関数は他方の成分で非零定数であり、前者でも $T$ に沿ってその定数に等しい。
非定数であればこの差は極の総次数 $e$ 以下に対し長さ $6-h>4$ の零点を持ってしまう。
この比較は特異な整四次上でも、$T$ がCartierであり極が滑らかな点に限られるので同じである。
偏極三階差から $p_i+2p_{i+3}\sim2p_{i+1}+p_{i+4}$、従って因子として等しく
$p_{i+3}=p_{i+1}$。三階差に戻すと $4p_i\sim4p_{i+1}$ なのでこれも等しい。
全偏極が等しく、全切断系で座標を揃えれば六関係は $Y$ の対角上で消える。
$Y$ を含む二次式は高々一つである。独立な二式は次数7を切れず、共通平面因子があれば
非平面な $C$ が残る二一次式の直線に入ってしまう。従って対角上に少なくとも二次曲面が残り矛盾する。
\end{proof}

\begin{corollary}\label{cor:all-septic-point-windows}
七次の全標準AS窓は、直線成分がある場合、直線がなく二次曲線成分がある場合、そのいずれもなく
平面三次成分がある場合、および整な七次八係数族で被覆される。
前の三行はそれぞれ系\ref{cor:line-component-atlas}、系\ref{cor:conic-component-atlas}、
定理\ref{thm:plane-cubic-mark-normality}の必要十分な有限条件に純次数7を課す。
行の所属は添字保持同型不変で、印の指定を忘れた比較と最小周期は各行の判定で得られる。
六次で直線・二次曲線・平面三次のどの成分も持たない点グラフが存在すれば、
相異なる二本の有理正規三次の被約な和で、その交わりは長さ4である。
\end{corollary}
\begin{proof}
系\ref{cor:full-quintic-point-window-atlas}の次数分割を使う。
七次で残る $3+4$ の空間三次の場合は前定理で、整七次は定理\ref{thm:integral-septic-boundary}である。
六次は純で、整な場合は定理\ref{thm:integral-sextic-exclusion}、二重三次は前節で除かれた。
二成分の算術種数3から交わりの長さ4が従う。
\end{proof}

% !TeX root = ../papers/nc-p3.tex
\section{二本の有理正規三次からなる六次点グラフの全排除}
\label{sec:binary-sextic-exclusion}
\begin{theorem}\label{thm:binary-sextic-exclusion}
四辺が点ファイバーの標準AS窓の共通点グラフは、二本の有理正規三次 $C_1,C_2$ の和 $Y$ にはならない。
四点が重なる場合も含める。
\end{theorem}
\begin{proof}[曲線の切断系と可動な三割線]
純性と種数公式より $T=C_1\cap C_2$ は長さ4である。$T$ は両射影直線上のCartier因子であり、
$Y$ はその貼合せである。各共通点の形式環は二つの滑らかな枝を同じ長さ $m$ の点で貼った
$k[[x,y]]/(y(y-x^m))$ で、$Y$ はGorenstein、$\omega_Y|_{C_j}=\mathcal O(2)$。
各 $L_i$ の両成分上の三次切断は $T$ 上の四切断に同型に制限される。
従って四文字は $L_i$ の完全切断系であり、$H^1(L_i)=0$。

$T$ の貼合せが両直線間の射影同型に延びれば、$Y$ は直線への有限平坦二重被覆となる。
局所的には同じ直線の函数二つを $T$ 上で一致させる代数であり、trace零部分は $\mathcal O(-4)$ である。
被覆鉛筆を $J$ として $\omega_Y=J^2$。$L_i/J^2$ の切断は片側で零ならもう一方の
次数1の切断も $T$ 上で消えて零になる。非零なら長さ2のCartier零因子を持ち、
$L_i$ の三切断がその上で消え、埋込みに反する。従って $f_*L_i=\mathcal O(1)^2$。
多重次数零の束の非零切断は連結な被約曲線上で生成元なので、
定理\ref{thm:smooth-sextic-exclusion}の有限平坦二重被覆の計算で三次元数は高々18となり不可能である。
この計算で二鉛筆が隣接して等しいときは、Segre二次曲面の対角が入る。
実際二曲線の和を含む二次式が二本あれば、独立二式の次数4または共通平面因子と非平面な台の矛盾になる。

以後貼合せが射影同型に延びないとする。各位置の $Y$ は二次曲面に入らない。
二平面では非平面な成分を含めず、錐では純六次因子の反射的イデアルは
$\mathcal O_Q(-3)$ で算術種数4になる。滑らかな二次曲面で各三次は $(1,2)$ または $(2,1)$。
両型が異なれば算術種数4、同じなら一方の射影が各成分で同型になり、上の貼合せが延びてしまう。
標準系は各成分の二次完全系に同型に制限される。実際、片方で零の標準微分は他方で正則な微分になるからである。
二つの標準像は異なる滑らかな平面二次曲線となる。同じなら両成分で同型な標準射影が貼合せを延ばす。
相異なるなら像の平面四次と $Y$ の算術種数が一致するので有限双有理な標準写像は同型である。

$\tau=\sigma^{-1}$ と書く。位置0で一般の滑らかな $p$ を固定する。
反対二点縮約の階数1以下の因子 $D_p$ は $Y$ 上で
\[
                   \mathcal O(D_p)=K^{-1}(p+\tau p)
\tag{*}
\]
である。定理\ref{thm:opposite-curve-divisor}の二つの核直線の共通零は、
固定した $p$ の場合には滑らかな二点 $p,\tau p$（同一点でもよい）の可逆イデアルである。
各成分で一般階数は2である。階数1の外積なら、両端の三つずつの函数の積は少なくとも五つを張るが、
変数側は各成分が張る四文字だけだからである。逆に $q$ を固定した場合にもこの議論が通るので、
一般の $p$ では全ての交点を避けて因子が得られる。
二成分上で二点とも動かす際、対角や $\tau$ のグラフがある成分の積ではその因子を核直線から除き、
異なる成分の像の座標が一致する有限個の点は除いて行列式を取る。
各射影直線の積は滑らかなので、二次小行列と核の外積の比は余次元2を越えて延びる
（極は余次元1で検出される）。従って階数低下の因子の二次数は、
当該成分積に存在する対角と $\tau$ のグラフの二次数の和である。
どの縦横ファイバーも因子として入らない。実際、交点を除いたその一般点では
既述の三次元函数空間の積の下界が使える。従って一方の射影に次数1なら他方にも次数1で射影同型のグラフとなる。

$D_p$ は位置2で $Y$ の三割線に乗ることを示す。
反対縮約の $3\times3$ 行列式 $S_p$ は一般の $p$ で非零である。
どの $R_i$ にも階数1関係はない（一本の三次だけでも全位置で空間を張る）。
従って縮約の中支持は3または4であり、中支持4の場合は定理\ref{thm:finite-cubic-separators}の
UFDでの余因子行列の証明を使える。中支持3で三次式が恒等零なら、
定理\ref{thm:ternary-point-subalgebras}の三支持の三文字部分代数は平面グラフを持つ
（階数1縮約がなく、その中央の行列式が恒等零）。従って第一と中央の三文字での平面射影が
$Y$ 全体で共通になる。その第一の中心は $j_1(p)$ で、中央の中心を $b$ とする。
$p$ の属する三次の射影像は二次曲線だから、$b$ も同じ成分に属する。
外からの三次の射影が非退化な二次曲線になることは次数から不可能である。
同じ比較を他方で行えば $b\notin T$ でもある。
全 $Y$ で基点を除いた偏極は従って $L_1(-p)=L_2(-b)$ である。
二つの一般の $p$ で比較すると、その差は一成分上の二次因子の線型同値となる。
他成分で定数の有理函数が長さ4の $T$ でその値を取り、極は二次以下なので定数である。
従って因子として等しく、$p$ が動くことから $b=p$、$L_1=L_2$ となる。
二次曲面がないので全関係が交換子になり矛盾する。

$Y$ は四つの三次最大小行列で切られる。実際二次式がなく
$h^0(L^2)=10$（両側の六次切断は長さ4へ全射）、
$H^1(L^n)=0$ ($n>0$) なので、基点自由鉛筆の完全列から全高次の制限も全射である。
従って定理\ref{thm:nonhyperelliptic-sextic-exclusion}のHilbert--Burchの証明を使える。
一般の $D_p$ は滑らかな部分にある。$S_p$ はこの三次系に入り $D_p$ で特異なので、
同定理の $3\times4$ 線型行列での核の議論により $D_p$ の張る直線上で
四つの最大小行列は一本の三次式の倍数となる。$Y$ 自体は直線を含まず、交わりはちょうど長さ3である。

\smallskip
\noindent\textit{四移動の位数。}
$\tau$ が一成分を点ごとに固定すれば、標準平面の滑らかな二次曲線を点ごとに固定するので全体で恒等である。
非恒等の場合は定理\ref{thm:fixed-differential}により固定スキームは長さ4のCartier因子である。
従って標準平面の射影変換は固定直線を持つ（孤立固定スキームの長さは3）。
非半単純な固定直線型は推移変換であり、不変な四次式は二変数だけになって二次曲線二本の積にならない。
半単純の場合、固定直線の切断が既に長さ4であるため中心は曲線外で、固定因子はこの標準切断であり $K^2=1$。
固定直線を $u=0$ とすれば四次式に $u^4$ と $u$ を含まない部分が共にあり、倍率の位数は2または4である。
位数4は定理\ref{thm:integral-sextic-exclusion}の巡回固有束の議論で除ける。
詳しくは標準式が $u^4=F_4(v,w)$ となり、$d_i=L_{i+1}/L_i$ は不変でノルムにより $d_i^4=1$、
従って $\tau^*L_i=L_i$。ノルムはここでも有限平坦な巡回四重射の一般軌道上の積で、
Cartier因子のノルム・引戻しの式は滑らかな点を台とする因子で確認できる
（十分正の束を用い特異点を避けた二因子の差で任意の線束を表せる）。
$P=L_2/\omega_Y$ は多重次数 $(1,1)$ で $H^0(P)=H^1(P)=0$。
非零切断は他成分で恒等零なら長さ4で消えるので零、さもなければ長さ2のCartier零因子で埋込みに反する。
よって不変な $P$ を線型化した順像の四固有束は $\mathcal O(-1)$ であり、これは一般軌道が自由推移的なので可約でも同じである。
$u$ の四つの一次乗法の積が $F_4$ で、根の上で $k[u]/u^4$ の自由階数1加群になるためどの二因子も根を共有しない。
すると全根が単純で標準曲線は滑らかになり矛盾する。以後 $\tau$ は恒等または対合である。

\smallskip
\noindent\textit{成分を保つ場合。}
一般の $p\in C_1$ に対して $(*)$ の多重次数は $(2,0)$。
三割線の残る点 $r$ は必ず $C_2\setminus T$ にある。有理正規三次 $C_1$ と直線は長さ3では交われず、
共通点が三番目となることもない。従って $D_p+r$ は滑らかな台上のCartier因子である。
$L_2(-D_p-r)$ には直線を含む平面から二切断があるのでRiemann--Rochにより
$D_p+r\sim P+s$ となる滑らかな $s\in C_1$ が得られる。実際残余の束は多重次数 $(1,0)$ で、
非零切断は $C_2$ で恒等零になれず生成元となり、$C_1$ 上の一点だけで消える。
$P$ は非有効で双対残余も非有効だから $|P+s|$ は一元で、$D_p$ が動く以上 $s$ は非定数である。
$(*)$ より $p+\tau p+r\sim PK+s$ でもある。

多重次数 $(1,1)$ の束 $P'$ について、$|P'+s|$ の因子の $C_2$ 上の一点がどう動くかを記す。
両正規化上の一次切断の $T$ での積
$H^0(\mathcal O_{C_1}(1))\otimes H^0(\mathcal O_{C_2}(1))\to H^0(\mathcal O_T(1,1))$ は同型である。
核があれば $(1,1)$ 形式が $T$ で消えるが、階数2なら貼合せが射影同型に延び、
階数1なら二因子で高々長さ2しか消えないからである。
$C_1$ の二次切断の $T$ 上への像は超平面で、その零化汎関数 $\lambda$ は各局所socleで非零である。
実際長さ3部分では二次切断の評価が同型である。従って $\lambda$ は $T$ 上で完全積対合を与える。
各束の自明化を選べば、$P'$ の貼合せは長さ4の環の単元 $v$ で、全体倍率を除いて決まる。
$s$ で消える一次式 $l_s$ と $C_2$ の一次切断 $m$ に対し、$m$ が $|P'+s|$ の切断の制限となる条件は
$\lambda(v l_s m)=0$ である。$C_1$ の二次切断にその値が持ち上がる条件そのものである。
上の二つの同型から $v\mapsto (l,m\mapsto\lambda(vlm))$ は四次元空間の線型同型である。
$P'=P$ では行列は可逆である。階数が落ちれば $s$ に関係なく条件を満たす $m\ne0$ があり、
$vm$ は $C_1$ の一次切断の評価像に入る。実際その像は両方の $l$ との積を $\lambda$ が消す二次元空間である。
これは非有効性に反する。従って $r$ は $s$ の射影同型像である。
$P'=PK$ も同じ非定数の対応を全ての一般の $s$ で与えるので、その行列は比例し、$P=PK$、$K=1$ である。
$(*)$ の二次因子は $D_p=p+\tau p$ として一致する。他方の成分で定数となる有理函数の極次数2と
$T$ の長さ4を再び比較すればよい。同じ結論が両成分で成立する。

$\tau$ が非自明な対合なら両成分上で非自明であり、
$C_1$ の $p,\tau p$ を結ぶ割線の和は滑らかな二次曲面となる。
実際、座標を $t\mapsto-t$ とすれば割線は $[1:0:t^2:0],[0:1:0:t^2]$ の結ぶ直線である。
各割線が $C_2$ にも会い、その交点が動く（一定なら一本の母線を除く全てにも含まれてしまう）。
よって $Y$ が二次曲面に入り矛盾する。

$\tau=1$ ならば両三次の各接線は相手に交わる。三次の接線面を二項三次の
$[u^2v]$ ($u,v$ は二項一次式) で表示すると、対角が元の三次で、
その外では重根 $u$ と残りの根 $v$ は一意である。
従って $C_2$ は $u\mapsto[u^2 h(u)]$ と書ける。実際、$C_2$ から両射影直線への二つの射の次数を
$e,f$ とすると $2e+f=3$ である。接点は動くので $e\ge1$、ゆえに $e=f=1$ で $h$ は射影同型となる。
$h=1$ なら同じ三次であり、$h$ の固有点だけが二曲線の共通点になる。
共通点が一つなら $\operatorname{Pic}^0Y$ は長さ4の局所環の単元を定数倍で割った群なので
標数零で捩れを持たない。既に $K=1$、四移動が恒等なので $L_{i+1}/L_i$ の四乗は1、全偏極が等しく不可能である。
共通点が二つなら $h$ は半単純である。
$C_2=H C_1$ とすると $H$ は対角行列で、その四成分は等差数列
$1,1+d,1+2d,1+3d$ に規格化できる。これは $u^2h(u)$ の二項展開による
（$\operatorname{Sym}^3$ の単項式基底を両側で揃える）。全て非零であり $d\ne0$。
逆に $C_1$ も $C_2$ の接線面に入るため $H^{-1} C_1$ は $C_1$ の接線面に入る。
対角行列 $H^{-1}$ は二つの共通端点を別々に固定するので、重根のパラメータ変更も残りの根の対応もその二点を別々に固定する。
従って $v_j=(1+jd)^{-1}$ は他の等差数列に等比数列を掛けたものである。
ある $b\ne0$ で $v_{j+2}-2b v_{j+1}+b^2v_j=0$ ($j=0,1$) が必要だが、
分母を除き $b$ を消去した終結式は $8d^6(1+d)(1+2d)$ で、矛盾する。

\smallskip
\noindent\textit{成分を交換する対合。}
標準平面の同次座標で $\tau:[u:v:w]\mapsto[-u:v:w]$ とし、一方の二次曲線を
$u^2+u l(v,w)+q(v,w)=0$ とする。もう一方は $u$ の符号を替えた式であり、$l\ne0$。
$T$ は各二次曲線上で $T_0+T_1$、ここで $T_0=(u=0)$、$T_1=(l=0)$ はともに次数2である。
二因子は相異なる。実際同じ切断因子を与える平面線型式なら二次曲線上で比例する。
$T$ を $C_1$ 上で見る埋込みを $x$、$C_2$ 側の埋込みを $\tau^{-1}$ で $C_1$ に運んだものを $y$ とする。
$x=y$ のイデアルは $T$ 上で $(u)$ で、その零化イデアルは $T_1$ のイデアル $(l)$ である。
中心外で射影座標の差を取れば $(u)$ となり、$C_1$ 上で $T$ は $ul$ のCartier因子だからである。
$T_1$ の二点が異なれば $y$ は $x$ の二点を交換する。
二重点なら $u=0=l$ の点で接線 $l=0$ が中心を通り、その接方向に対合は倍率 $-1$ を誘導する。
この場合 $y$ は $x$ に対し $T_1$ の一階jetの符号を替える。

前述の二次数より $p\in C_1$ に対して $D_p=A(p)+B(p)$ で、$A:C_1\to C_1$、$B:C_1\to C_2$ は射影同型である。
$(*)$ の線束同値を $T$ 上の貼合せで比較すると、
\[
         A^{-1}x+y=x+B^{-1}\tau y\qquad\text{in }\operatorname{Sym}^2(C_1)(T).
\tag{**}
\]
詳細には $[x,A(p)][y,p]$ と $[x,p][\tau y,B(p)]$ は、$p$ によらない貼合せ単元と
$p$ によるスカラーを除き一致する（角括弧は二項一次の行列式で、$C_2$ 側は $\tau$ で座標を運ぶ）。
両者は $p$ の二次式で、どの $p$ においても二つの一次因子の積は長さ4の $T$ 上で恒等零にならない。
従って $p$ によるスカラーは両側で極も零も持たず定数となる。
$T$ 上の $p$ の二つの根を因子として比較すれば $(**)$ であり、非被約な $T$ でも有効である。
$F=A^{-1},G=B^{-1}\tau$ と置く。$T_0$ では $(**)$ の同一の一次因子を消して $F=G$ となる。
局所アフィン座標で二次式に $p=x,y$ を代入すれば $(Fx-x)(y-x)=0$、
$(Gy-y)(y-x)=0$。従って $F,G$ は $T_1$ を各々点ごとにスキームとして固定する
（$y$ による $T_1$ の像も同じ因子である）。
$T_0\cup T_1$ の長さは3以上なので $F=G$ は射影同型として成立し、$B=\tau A$ となる。

ゆえに $z\in C_1$ と $\tau z\in C_2$ を結ぶ一般の直線は三割線である。三番目の点が乗る成分を
$C_1$ と呼び直せる（対合なのでこの直線族は同じである）。その点を $g(z)$ とすれば $g$ は射影同型である。
実際、$C_1$ の三つのHankel二次式を $C_2$ に制限し、長さ4の共通零因子を除くと
基点自由な二次式系が得られる。これが $C_2$ の点を通る $C_1$ の唯一の割線の
二次因子へ送る写像である。Hankel行列の核の二項二次式を取れば、相異なる根の二評価ベクトルを結ぶ直線が戻り、
重根では接線が戻る。$z$ と $g(z)$ の非順序対の系の次数が従って2であるから $1+\deg g=2$ となる。
この議論は三番目の点が $z$ に重なる場合も含む。

この二次因子族 $z+g(z)$ の割線に $\tau z$ が乗る包含は全基底で閉条件として成立する。
$T$ 上で $\tau z$ 自身が $C_1$ にあるので、$z=x$ とすれば $y$ は因子 $x+g(x)$ に入る。
割線上の $C_1$ のスキームは普遍長さ2因子そのものである。実際、二項二次式を $Q$ とすれば
その割線の二平面式は $Qs,Qt$ であり、$\mathbb P^1$ の $s,t$ は基点を持たない。
従って $(y-x)(y-gx)=0$、$T_1$ 上で $gx=y$ である。
$T_1$ が二単点なら $g$ は二点を交換し、二重点なら一点を固定して接倍率が $-1$ である。
どちらも $g$ は非自明な対合となる。
その割線は先に示した滑らかな二次曲面を掃き、$C_2$ もそこに入るので矛盾する。
\end{proof}

\begin{corollary}\label{cor:point-curve-high-coverage}
点ファイバーの一次元点グラフで純部分の次数が4以上なら、全標準AS窓は既出の有限アトラスで被覆される。
六次では直線、二次曲線、平面三次のいずれかの成分が必要である。
全次数を通じ、点グラフ一次元の枝でこれらの印、四次完全交叉印、整五次・七次の既出分類の
どれにも入らないとすれば、純部分は空間有理正規三次で有限余長6を持つ。
\end{corollary}
\begin{proof}
六次は系\ref{cor:all-septic-point-windows}と本定理による。
四次で直線・二次曲線がなければ純部分は整であり、系\ref{cor:integral-hull-high-degrees}を適用する。
五次・七次は既出の全被覆、三次の残余の有限長は定理\ref{thm:pure-curve-defect}である。
\end{proof}

% !TeX root = ../papers/nc-p3.tex
\section{有理正規三次上の基点を持つ階数二関係}
\label{sec:twisted-cubic-line-modules}
各辺の両射影のファイバーが空か一点である標準有限Frobenius窓で、
点列の中に全位置で完全な三次系により埋め込まれる同じ $C=\mathbb P^1$ があるとする。
曲線上の点の移動に沿って座標を取り、$V_i=S^3$ と書く（$S^m=H^0(C,\mathcal O(m))$）。
$H_i(p)$ は点 $p$ の消滅超平面である。曲線外の有限構造には条件を付けない。
\begin{theorem}\label{thm:twisted-cubic-line-resolution}
$r\in R_i$ が階数2で、その左支持 $U\subset V_i$ の二つの三次式が共通零点を持つとする。
右支持を $W$ とする。右加群 $L_i(U)=e_i A/U A$ の長さ $n$ の次元は全て $n+1$ であり、
$U\otimes e_{i+1}A\to e_iA$ の関係は $r$ 一個で生成される。
左加群についても両端を反転した結論を得る。
さらにこの窓が標準Koszul左右AS正則であることと、全始点・長さで
$\dim A_{i,i+n}=\binom{n+3}{3}$ となることは同値である。
\end{theorem}
\begin{proof}[点加群による線形表示]
$r$ の両支持の鉛筆は基因子を除くと同じ次数の射を与える。
実際、二つの独立な左係数 $u,v$ と右係数 $s,t$ で $r=u\otimes t-v\otimes s$ と書くと、
対角上で $ut=vs$ である。従って $u=dF,v=dG,s=gF,t=gG$ と書け、
$F,G$ は互いに素で $\deg d=\deg g$ である。特に左右の基点の有無は一致する。
非零の階数1関係は積が整曲線上で消えないため存在しない。

共通零点 $p$ を一つ取り、$q\in\mathbb P(U^\perp)$ を $p$ と異なる点とする。
$H_i(p)=U\oplus kv$ として $q(v)=1$ とする。
右点加群 $M_i(p)=e_iA/H_i(p)A$ のこの超平面の関係は
$K=R_i\cap(H_i(p)\otimes V_{i+1})$ の三個で生成される。
これは命題\ref{prop:commuting-three-point-graph}の点列での階数 $1,3,3,1$ の計算である。
$p$ の前点で四テンソルを縮約した三テンソルを $h$ とすれば、その末尾縮約が $K$ を張る。
$h$ の先頭を $q$ で縮約した $r'$ は零でない。
実際、前点から出る辺の六関係上の対評価が零なら、点ファイバーの一意性で $q=p$ となる。
相補二位置との完全対によりその対評価はこの縮約で読める。
$K$ での先頭 $q$ 評価は $r$ を消すので、その像 $U'$ は次元2以下であり、
$r'\in R_{i+1}$ の左支持である。階数1の関係がないので次元は2、核は $kr$ となる。
$h$ の最終支持は四移動で同一視した前点の消滅超平面に入る（巡回した隣接グラフでの消滅）。
従って $r'$ の右支持も曲線上に基点を持ち、先の分解により $U'$ にも基点がある。

$va\in UA$ は $H_i(p)$ の関係を意味し、その $v$ 成分全体は上の三生成と $q$ 評価により
ちょうど $U'A$ である。従って次数を各始点から測ると
\[
 0\longrightarrow L_{i+1}(U')(-1)\xrightarrow{\ v\ }
 L_i(U)\longrightarrow M_i(p)\longrightarrow0 .
\]
全ての後続加群でこの操作を反復し、長さ0から帰納すれば主張の次元を得る。

$K$ の基底を $r,s_1,s_2$ とし、後二者の $v$ 成分を $u_1',u_2'$ とする。
$h=r\otimes c_0+s_1\otimes c_1+s_2\otimes c_2$ と書けば
$r'=u_1'\otimes c_1+u_2'\otimes c_2$ であり、
$h\in H_i(p)\otimes R_{i+1}$ により $r c_0+s_1 c_1+s_2 c_2=0$ が $H_i(p)\otimes e_{i+1}A$ で成り立つ。
$U$ の任意の関係を $r a_0+s_1 a_1+s_2 a_2$ と表示すると
$u_1'a_1+u_2'a_2=0$ である。一つ短い長さでの帰納仮定により
$(a_1,a_2)=(c_1 b,c_2 b)$ だから、元の関係は $r(a_0-c_0 b)$ である。
開始次数2では $q$ 評価の核 $kr$ そのものとなる。左側でも逆向きの点列で同じ計算ができる。
\end{proof}

\begin{proof}[Hilbert関数からの十分性]
どの位置でも任意の $p\in C$ を含む二次有効因子 $D$ で
$U=H^0(\mathcal O(3)(-D))\subset H_i(p)$ を左支持とする階数2関係が存在する。
実際、$R_i=[z,t]T_i$ と二次の六平面 $T_i\subset S^2\otimes S^2$ を書き、
$D$ の式を $d$ とすると、求める条件はある $g\ne0$ に対する $d\otimes g\in T_i$ である。
$T_i^\perp$ の三行を $d$ で評価した三次行列式は、
$p$ を根に持つ $d$ の射影直線で必ず零点を持つ。左加群用の逆向きでも同じである。

所定のHilbert関数を仮定する。上で得た三項線形表示のEuler式から、
左端の $e_{i+2}A\to U\otimes e_{i+1}A$ も単射となる。
左表示にも適用できる。右表示の双対は、階数2の係数対の基底を合わせると
同じ $r$ の右支持 $W$ から作る左表示の末尾の商を除いた複体である。
従って各右 $L_i(U)$ から自由加群への次数0,1のExtは消える（各表現可能加群を標的とする）。
任意の点加群は先に構成したこの種の $L$ の商なので自由加群へのHomは零であり、その尾部も点加群である。
上の短完全列を正次数に制限し、$L_{i+1}(U')$ と点加群の尾部の両方でHomが消えることから
$L_i(U)_{\ge1}$ のHomも消える。

右単純の第一Extの消滅を見るため、$e_iA_+$ から自由加群への任意の写像を $UA$ に制限する。
$L_i(U)$ の第一Ext消滅でこれは $e_iA$ 上に延びる。差は $L_i(U)_{\ge1}$ を通るので零である。
単純からのHomも、$L_i(U)$ の商であるため零である。左側でも同じで、
候補Koszul複体の生成・二次表示の端とそのFrobenius双対の端が完全になり、中央はEuler式による。
逆向きは標準Koszul解消のEuler式である。
\end{proof}

\begin{corollary}\label{cor:twisted-cubic-compatible-pencils}
全添字で二次元鉛筆 $U_i\subset V_i$ があり、
$R_i\cap(U_i\otimes U_{i+1})$ が基点を持つ階数2関係 $r_i$ を含めば、この窓は標準正則である。
特に四移動の任意の指定した有限べきで閉じるこのような列は有限方程式による十分条件を与える。
\end{corollary}
\begin{proof}
二文字・一非退化関係の添字付き代数 $P$ は、基底を両方向に順に合わせれば二変数多項式環の反復である。
定理の関係生成を $P$ 単純のKoszul複体に適用すると $\operatorname{Tor}_1^P(k_i,A)=0$。
各終点で有界下加群の極小表示を取ると $A$ は左 $P$ 自由である。自由生成元の個数は
$e_iA/U_iA$ の次元から長さ $n$ で $n+1$ なので、Hilbert級数は $(1-t)^{-4}$ となる。
\end{proof}

\begin{proposition}\label{prop:twisted-cubic-secant-cubic}
$T_i^\perp$ の $d$ 評価の行列式が恒等零でなければ、その平面三次式は必ず直線因子を持つ。
$g$ 側の行列式についても同じである。
\end{proposition}
\begin{proof}
$d$ の根を順序付きで $p,q$ とする。$q\ne p$ では $q$ の曲線上の像を
上記の $q$ 評価に使い、$q=p$ では接線上の微分評価を使う。
このとき $H_i(p)$ 内の評価核はいずれも $d S^1$ である。従って求める三次式の消滅は、
三テンソルの先頭をこの評価で縮約した行列の階数が2以下であることに等しい。
四テンソルを四辺の角括弧と多重一次式 $f$ の積に書く。
その順序を $p,q,z,t$ とすると、$p,q$ での縮約から $[p,q]$ と両端の
一次式（残った二因子に作用する単射）を除いたものは $[z,t] f(p,q,z,t)$ である。
継ぎ目の四移動は外側の一次式か $p$ の座標同一視に繰り入れる。
$q=p$ では $q$ 側の微分が消えた括弧にだけ作用し、同じ行列を与える。

$f(p,q,z,t)=a+bz+ct+d'tz$ と書くアフィン座標では、残った三行三列の行列式は
$(c-b)(ad'-bc)$ である。ここで $c-b$ は $p,q$ について二次数 $(1,1)$ の切断である。
これが恒等零なら元の三次式も恒等零である。そうでなければ、この切断の零因子の
各既約成分の対称平方平面での像は直線または二次曲線である。
実際、順序を忘れる有限写像による超平面束の引戻しは $\mathcal O(1,1)$ で、
その各成分上の次数は1または2である。元の三次式はこの像で消える。
二次曲線を含む場合も三次式の残余は直線である。両側の交換で他方の結論を得る。
\end{proof}

% !TeX root = ../papers/nc-p3.tex
\section{二項四テンソルの整な種数一対応と整合する割線鉛筆}
\label{sec:elliptic-binary-cubic-pencils}
前節の有理正規三次 $C$ が点グラフの一次元の台全体で、残りの構造は有限とする。
各埋込点の埋込み次元は2以下とする（3の場合の判定は定理\ref{thm:maximal-tangent-ore}）。
四移動の曲線上の座標変換を $S$ と書き、四テンソルから四つの隣接角括弧を除いた多重一次式を $f$ とする。
\begin{theorem}\label{thm:elliptic-binary-cubic-pencils}
隣り合う二変数を固定し、残り二変数の $2\times2$ 行列式を取る。
このうち一つが整な $(2,2)$ 曲線なら、この窓は標準Koszul左右AS正則である。
節点型・尖点型も含み、各位置で高々四個の割線鉛筆の集合を選んで、隣り合う集合を関係で全射的に結ぶことができる。
\end{theorem}

\begin{proof}[順序付き二点の曲線]
四移動の継ぎ目を最後に置き、必要なら $S$ を逆に取る。
無限個の位置に各三次系を $C$ の座標で自明化する。
全ての隣接二座標で $f$ の評価行列が零になることはない。
対角外では元の四テンソルもその二点で消え、一意の点ファイバーに反する。
対角上の $(p,p)$ なら第二座標だけ曲線に沿って微分しても、括弧と $f$ が共に消えるので零となる。
これは第二位置の曲線の接線全体が第一位置の $p$ のファイバーに入ることを意味し、やはり不可能である。
$E$ を仮定の整因子とするとこの上で行列は階数1である。
左核、右核から、どの隣接二座標にも前と次の座標がそれぞれ一意に付く。
実際、隣接三点の条件は三変数を入れた四テンソルの商 $f$ が残る一変数で恒等零となることで、
隣の窓では同じ条件を巡回して読む。より正確には、任意の隣接行列式の零点スキームから、階数1行列の核を足して得る三点の
スキームは、隣の行列式の零点スキームからも逆側の核を足したものである。
これは単元成分で行列を掃き出せば階数1条件が残り一つの行列式となることによる。
従って隣の行列式が恒等零になることもなく、全て同じ整な算術種数1曲線となる。
一つの $E$ 上に全ての二重被覆 $p_j:E\to C$ を取れ、
$(p_{j-1},p_j)$ は対応する $(2,2)$ 曲線への同型である。
$D_j(x)=p_{j-1}(x)+p_j(x)\in\operatorname{Sym}^2C$ とする。
辺 $j$ には全ての $x\in E$ について次の割線関係がある。
\[
            D_j(x)\ \longmapsto\ S D_{j-2}(x).
\]
ここで因子 $d$ の鉛筆は $d S^1$ である。実際、四テンソルの先頭二変数に
$p_{j-2},p_{j-1}$ を入れると、残る二変数について $f$ は一次式の積で、
その根は $p_j,S p_{j-3}$ である。外側の二つの括弧が $p_{j-1},S p_{j-2}$ を加え、
中央の括弧が二鉛筆の階数2関係を与える。外側で評価した二点が一致する有限箇所にも
割線関係の閉性で延びる。四位置の移動で $D_{j+4}(E)=S D_j(E)$ となる。

\end{proof}
\begin{proof}[四次平面像の排除]
各辺の $\Pi=\mathbb P(T^\perp)$ は可逆な三行三列行列を含む。
実際、定理\ref{thm:twisted-cubic-residual}の基点の吹上げは完全交叉
$\{\det B=Q(B)=0\}$ に全射である。この像の計算は点グラフと無関係である。
$\Pi$ 上の吹上げの像で、例外因子以外から来るものは曲線外の有限個の点対から来る。
例外因子では $C$ の法微分を射影化しており、六式が消える法方向の次元は
点グラフの接空間の次元から1を引いたものである。従って有限個の埋込点の上で各高々一つの射影点しかない。
ゆえに $\Pi\cap\{\det B=Q=0\}$ は有限である。$\det B$ が $\Pi$ 上で恒等零なら
二次式だけで平面に無限個の零点を持つので矛盾する。

さらに $\Pi$ は片側の因子が二次Veronese曲線上にある階数1行列を含まない。
実際、基点写像 $B(X,Y)$ で $X$ に曲線上の $p$ を入れると、アフィン表示では
$B=u(p)(Y_{k+1}-pY_k)_k^t$、$u(p)=(1,p,p^2)^t$ を得る。
三差分は全射であるから、該当する非零行列が $\Pi$ にあれば $Y$ は $p$ 以外のファイバー点になる。
右側および無限遠点も同様である。

基底 $B_1,B_2,B_3$ を使い $M(d)$ の行を $d^t B_k$ とする。割線関係は $M(d)g=0$ である。
この三次行列式は恒等零にならず、階数1以下の軌跡は有限である。以下は $E$ の整性によらない。
可逆な $B_1$ を取れば、一般階数1以下は各 $d^t B_k$ が全 $d$ で $d^t B_1$ に比例することになり不可能である。
一般階数2で行列式が恒等零なら、随伴行列は多項式環の一意分解により $h(d)u(d)v(d)^t$ と書ける。
ここで $u,v$ は成分の最大公約数1の斉次ベクトルである（有理関数上で階数1に分解し分母を払う）。
定数左核は基底の独立性に反し、定数右核は可逆な $B_1$ に反する。
従って $u,v$ は一次である。$u$ が可逆な線型写像なら全 $d$ で $d\otimes u(d)\in T$ となり、
六次元の $T$ は対称二乗の片側を変換したものに等しい。直交する $\Pi$ に可逆行列がなく矛盾する。
そうでなければ、原始性により $u$ は階数2で、基底を変えて $u=(-d_2,d_1,0)^t$ とできる。
$M u=0$ より各行の先頭二成分は $(\lambda_k d_1,\lambda_k d_2)$ となる。
$\lambda$ の核から来る行列は一列しか持たず、$\Pi$ は階数1行列の直線を含む。
動く側の因子の直線は必ずVeronese二次曲線に交わるので、これも不可能である。
階数1以下の曲線があれば三次行列式の重複因子なので直線でなければならない。
その直線で $d^t B_k$ は $d^t B_1$ に比例し、$B_1$ の可逆性により各比例係数は直線上で一定である。
再び二つの独立な一次結合の行列はその左側二平面で消え、禁止された階数1行列の直線を得る。

$D_j(E)$ は二次曲線か四次曲線である。順序を忘れる射が次数2以下で
$D_j^*\mathcal O(1)$ の次数が4だからである。非零の三次行列式に入るため前者となり、
$D_j:E\to D_j(E)$ は次数2である。一般点で $M$ の階数が2なので $S D_{j-2}$ は有理的に $D_j$ を経由する。
両者が次数2なので被覆の対合 $s_j,s_{j-2}$ は等しい。
$s_j$ は $p_{j-1},p_j$ を交換する。順序付き二点の写像が埋込みだからである。
\end{proof}

\begin{proof}[楕円曲線の加法と鉛筆の連結]
滑らかな部分 $G$ に原点を取って群とする。滑らかな場合は楕円曲線、
節点なら $\mathbb G_m$、尖点なら $\mathbb G_a$ である。全て加法記号で書き、
$G\simeq\operatorname{Pic}^0 E$ は $x\mapsto\mathcal O(x-0)$ で同一視する。
この同一視は次数1の可逆層にRiemann--Rochで一意の零点があり、長さ1のCartier点は滑らかであることからも従う。
一つの二重被覆を $p$ とし、一般のファイバーの二点の和が一定なので、
原点を平行移動してその対合を $x\mapsto-x$ とする。$G$ 上の2倍は全射である。
$s_j$ は反射 $x\mapsto a_j-x$ である。実際、$G$ の自己同型は平行移動と原点を保つ群自己同型の合成で、
位数2ならその線型部は $1$ か $-1$ である。位数2の平行移動なら
$p_{j-1}$ の被覆対合と可換で、$p_j=p_{j-1}s_j$ は同じ対合を持ち、二点の埋込みに反する。
よって全 $j$ で $p_j(x)=p(x+c_j)$ と書ける。
交換対合の一致から $c_{j-1}+c_j=c_{j-3}+c_{j-2}$ が成り立つ。
また $L_j=p_j^*\mathcal O_C(1)$ とすると $L_jL_{j+3}=L_{j+1}L_{j+2}$ である。
実際、中二点で評価した階数1行列は $\mathbb P^1\times\mathbb P^1$ のSegre曲面への射を与え、
基点のないその四切断の束は $L_{j+1}L_{j+2}$、両射影で得る核の束は $L_j,L_{j+3}$ である。
$\operatorname{Pic}^2 E$ の加法座標で $L_j$ は $-2c_j$ なので、
$\delta_j=c_j-c_{j-1}$ とすれば $2\delta_{j+2}=2\delta_j$ となる。
さらに二周期の和 $c_{j-1}+c_j$ から、$c_{j+2}-c_j=t_j$ は $t_{j+1}=-t_j$ を満たし
$\delta_{j+2}-\delta_j=2t_j$ である。従って $4t_j=0$、$\delta_{j+4}=\delta_j$ である。

$S$ は $G$ に $v\mapsto \alpha v+b$ と持ち上がる。実際、四移動で隣接順序対の曲線を同型に写せる。
第一射影はそれぞれ $p$ を平行移動したものなので、これら平行移動で補正すれば $p$ に関する持上げとなる。
被覆対合との可換性から $2b=0$ である。
ここで $\alpha$ は原点を保つ自己同型であり、特殊な $j$ 不変量も許す。
二次曲線 $D_j(E)$ は順序を付けると $(p(v),p(v+\delta_j))$ の像で、これが決める差は符号を除き一意である。
実際、第一座標の一般持上げは $v,-v$ のみで、第二座標の一致で差を決定できる。
四移動の像と $\delta_{j+4}=\delta_j$ により $\alpha\delta_j=\pm\delta_j$ である。

$2d_j=\delta_j$ とし $m_j=c_{j-1}+d_j$ と取る。
$P(y,d)=p(y-d)+p(y+d)$ と書くと $D_j(x)=P(x+m_j,d_j)$ であり
\[
 P(y,d)=P(d,y),\qquad P(y+e,d)=P(y,d+e)\quad(e\in G[2]).
\]
各引数は独立に符号を変えられる。$2(m_j-m_{j-2})=0$ なので辺の対応は
$P(y,d_j)$ を $P(\alpha y+e_j,\alpha d_{j-2})$ に運ぶ（$e_j\in G[2]$）。
$\alpha d_j$ と $\pm d_j$ の差、$\alpha d_{j-2}$ と $\pm d_{j+2}$ の差はいずれも二捩れである。
従って有限集合
\[
 \{\,P(d_{j+1}+e,d_j):e\in G[2]\,\}
\]
から次の同じ集合に割線関係で全射する。実際、二引数を交換して二捩れを第一引数に集めれば
$P(d_{j+2}+e',d_{j+1})$ となり、$e$ を変えれば $e'$ が二捩れ全体を走る。
この列を両方向へ延ばし、系\ref{cor:twisted-cubic-compatible-pencils}を適用する。
\end{proof}

\begin{corollary}\label{cor:binary-vertical-ternary}
同じ窓で隣接行列式が恒等零、または一つでも垂直・水平な成分を持つ場合も標準正則である。
従ってこの三次曲線の残余でここまでの試験の外に残る二項対応は、
二本の $(1,1)$ グラフの和またはその二重化だけである。
\end{corollary}
\begin{proof}
一変数をその成分の値に固定する。次の変数を動かすと残り二位置の行列は全て階数1である。
この線型鉛筆には零行列もなく、Segre曲面の直線だから、どちらかの因子で核が一定である。
隣接位置の一定核なら先に排除した隣接零行列になるため、反対位置で一定核となる。
従って四テンソルも相対する一対の曲線点で消える。
一方の点での三テンソルの中支持は高々3である。両端支持は点ファイバー条件で3、
各 $R_i$ に階数1関係はない。定理\ref{thm:finite-frobenius-low-opposite-as}の末段の鉛筆の議論により
中支持は3であり、定理\ref{thm:ternary-point-subalgebras}の三方向の縮約にも階数1はない。
これで全次数正則性を得る。残る因子の主張は二次数を加えるだけである。
\end{proof}

% !TeX root = ../papers/nc-p3.tex
\section{二重メビウスグラフと二次因子による有限鎖判定}
\label{sec:doubled-binary-pencils}
前二節と同じ三次曲線の残余で、埋込み次元が2以下とする。
前節の $E$ は隣接位置での $f$ の行列式を取ったスキームとする。
\begin{lemma}\label{lem:cubic-secant-isomorphism}
各辺の割線因子の対応 $M(d)g=0$ は、二つの平面三次スキームの同型のグラフである。
さらに $E$ の任意の成分・無限小構造を含め、$D_j(x)\mapsto S D_{j-2}(x)$ はこのグラフに入る。
\end{lemma}
\begin{proof}
三次式が非零であることは前節で証明した。
ある $d$ で $M(d)$ の核が二次元以上なら、$d S^1$ に支持される独立な二つの関係
$[z,t]d(z)g(t)$ ができ、定理\ref{thm:twisted-cubic-line-resolution}の二次での唯一性に反する。
逆向きも同じである。よって零点上では単元の二次小行列式による消去で射影がスキームとして同型になる。
$E$ 自身の隣の位置への移動も前節の単元による核の構成でスキーム同型である。
順序付き二点の全平面積で、四テンソルの二点縮約を先頭の括弧で割った二形式は
関係空間に値を取る。これは括弧が非零の開集合で成り立つ多項式恒等式だからである。
この恒等式を $E$ に制限し、階数1の行列を局所分解すれば所定の二因子を得る。
先頭の括弧が $E$ 上で零因子となる場合もこれでよい。
\end{proof}

\begin{theorem}\label{thm:doubled-binary-pencils}
$E$ が $(1,1)$ グラフの二重化なら、窓は標準Koszul左右AS正則である。
被約な二グラフの場合も、次のいずれかで同じ結論となる。
\begin{enumerate}
\item 一つの成分を全位置へ運んだとき、その上で各隣接座標の変換が恒等または対合となる。
\item 全ての位置で二つのグラフが互いの逆変換のグラフとなる。
\end{enumerate}
第一条件は四移動が成分を交換しても八つの連続位置で調べられる。
\end{theorem}
\begin{proof}
成分を一本取り、$h_j:\mathbb P^1\to C$ をその上の各位置の射影同型とする。
二重グラフの場合、$\phi=h_jh_{j-1}^{-1}$ が恒等でも対合でもなければ、
順序を忘れる射の下で一般点の近くで二重の二次曲線が現れる。
実際、グラフと逆グラフは異なり、順序を忘れる射は対角外でエタール、
被約グラフの像は次数2で一般点上で一対一である。
従ってこの二重グラフ上で消える平面多項式は二次曲線の一般点で少なくとも二重に消える。
補題の非零の三次式では不可能である。よって二重化は第一条件に入る。

第一条件で $a_j=h_{j-1}^{-1}h_j$ は恒等または対合である。
$x\mapsto D_j(x)$ は $a_j$ が恒等なら対角二次曲線への同型、
非自明なら直線への二重被覆で対合は $a_j$ である。
辺のグラフが同型なので各 $x$ について $D_j(x)$ と $D_{j-2}(x)$ は互いを決め、$a_j=a_{j-2}$ である。
二つの変換を $a,b$ と書く。これらに共通固定点 $x$ があれば、全 $h_j(x)$ は同じ $p$ で、
$D_j(x)=2p$ は位置によらない。共通固定点がなければ、両者は非自明で異なり、
$ab$ の固定点は相異なる二点である。一点なら両変換が $ab$ を逆に共役するのでそこを固定するからである。
また共通固定点の非存在により $a,b$ はともにこの二点を交換する。
実際、一方が二点を個別に固定すれば、積も固定するため他方も個別に固定する。
この二点集合を $Z$ とすると $h_j(Z)=h_0(Z)$ で、$x\in Z$ では $D_j(x)$ が常に同じ二次因子である。

第二条件では二つの異なるグラフは $\phi,\phi^{-1}$ で、$\phi$ は対合でない。
両者が交わる点の座標は $\phi^2$ の固定点であり、いずれも $\phi$ 自身の固定点である。
二成分の共通点を一つ選んで全位置へ運ぶと、各辺でその両座標が等しいので $D_j(x)=2p$ はやはり一定である。

以上のいずれでも、ある一定因子 $D$ について全 $j$ の辺に $D\mapsto SD$ がある。
四移動により各整数 $m$ について $S^mD\mapsto S^{m+1}D$ も全辺にある。
これで両方向の鉛筆列ができ、系\ref{cor:twisted-cubic-compatible-pencils}による。
\end{proof}

\begin{proposition}\label{prop:decide-compatible-secants}
この三次曲線の残余で、基因子の次数2の鉛筆による両方向の整合列が存在するかどうかは、
四窓の三次式、核同型および四移動に対する有限な閉集合演算で決定できる。
この試験に通る窓は全次数で正則である。
\end{proposition}
\begin{proof}
被約な射影平面の四つのコピー上で各窓の出発三次曲線を $B_i$ とする。
四辺目の終点を四移動で最初の位置に戻す。辺の同型を $f_i:B_i\to B'_i$ と書く。
$Z_i=B_i$ から始め、四位置を同時に
\[
 Z_i\ \longleftarrow\ Z_i\cap f_i^{-1}(Z_{i+1}\cap B'_i)
                       \cap f_{i-1}(Z_{i-1})
\]
で更新する（最後の項は位置 $i$ の平面に取る）。閉埋込みの像と逆像なので全て閉で、
根基イデアルの計算と等号判定で反復できる。閉集合の降鎖条件により有限回で安定する。
任意の両方向列は常に残る。逆に安定後の各点には残った前点・次点があり、
辺の同型によりこれを両方向へ延ばせる。四移動をほどけばもとの整合列である。
\end{proof}

% !TeX root = ../papers/nc-p3.tex
\section{二つのメビウス成分と一次元点グラフの全被覆}
\label{sec:split-binary-pencil-charts}
\begin{theorem}\label{thm:split-binary-pencil-charts}
前節の三次残余で $E$ が二つの相異なる $(1,1)$ グラフの和なら、全ての窓は標準Koszul左右AS正則である。
\end{theorem}
\begin{proof}[二つの節点での座標]
反対二位置のある二点で $f$ が零行列なら四テンソルもそこで消え、系\ref{cor:binary-vertical-ternary}の
三支持の議論で済む。以下、隣接・反対のどの二位置にも零行列がないとする。
二成分の交点は長さ2である。まず二単点の場合を扱う。
どの位置でもその二点の像は相異なる。四変数の各々でこの二点に対応する座標を $0,\infty$ とする。
ここで四変数のラベルは一つの特異点での核を用いて合わせる。実際、最初の二変数の
行列式の特異点では、その両核を残りの二変数に入れた四組が $f$ の全ての一階偏微分を消す。
これは階数1行列の行列式の微分が左右核での評価に比例するためである。
この四組から任意の隣接二座標を取っても同じ特異点の移動となる。
従って二つの四組は全て0と全て無限遠になり、$f$ は二変数ずつの積六項だけを持つ。
係数が一つでも零ならその相補二位置の $(0,0)$ で零行列となる。
各座標の非零定数倍で相補係数を等しくできる（三対の比の文字は独立で、代数閉体上そのトーラス写像は全射）。
ベクトル座標を $u,v,z,t$ とすると
\[
 F=a(uv+zt)+b(uz+vt)+c(ut+vz),\qquad abc\ne0 .
\]
ここでの四つの座標は独立に選んだもので、点グラフ $C$ に沿った同一視とは一般に異なる。
$u,v$ を同じ $C$ に戻すと二点は $P_0 u,P_1 v$ となる。
$z,t$ も同様に $P_2 z,P_3 t$ と書く。$\Theta=P_0^{-1}P_1,\ \eta=P_2^{-1}P_3$ と置く。

$u,v$ の二成分は $v=\lambda u,\lambda^{-1}u$ で、
$bc(\lambda+\lambda^{-1})=a^2-b^2-c^2$, $\lambda^2\ne1$ である。
$m=-(c+b\lambda)/a,\ n=-(b+c\lambda)/a$ とすると $mn=\lambda$ で、
第一成分での残る二根は $z=mu,t=nu$、第二成分では $m,n$ を逆数にする。
特に $m^2\ne1$ である。同じ $u,z=mu$ で二つの $v$ が核になるなら反対二位置で零行列になるからである。
補題\ref{lem:cubic-secant-isomorphism}の辺で、第一成分は
\[
 \{P_0 x,P_1(\lambda x)\}\ \longmapsto \{P_2(x/m),P_3(\lambda x/m)\},
\]
第二成分は $\lambda,m$ を逆数にした対応を与える。実際、$t=nx,u=x$ を固定したとき
$v$ 全体で消える根は $z=x/m$ である。これは $F$ に代入して従う。
四移動の継ぎ目があれば補題と同じく終点の二座標を一緒に移すので相対変換 $\eta$ には影響しない。
\end{proof}

\begin{proof}[混合した変換の制限]
$\Lambda(x)=\lambda x$ とし、二つの変換 $\Theta\Lambda,\Theta\Lambda^{-1}$ を調べる。
どちらかが恒等なら $\Theta$ は対角である。ともに非自明な対合なら二つの行列の跡が消えて
$\Theta$ は反対角である。どちらも対合でも恒等でもなければ、順序を忘れた両像は二次曲線で
同じ三次式に入るため一致する。順序を戻す二つのグラフは互いの逆であり
$(\Theta\Lambda^{-1})^2=\Lambda^{-2}$ が射影的に成り立つ。
非自明な倍率の平方根はその二固定点を保ち、反対角行列の平方は恒等だから $\Theta$ は対角である。

片方だけが非自明な対合で他方が恒等でない場合、第一成分を対合としてよい。
$\Theta=\bigl(\begin{smallmatrix}A&B\\C&-\lambda A\end{smallmatrix}\bigr)$ と書くと $A\ne0$。
第一成分上の割線因子は次数2で直線を被覆し、被覆対合は $\Theta\Lambda$ である。
辺が同型なので、$M(x)=mx$ とすると終点の対合は
$\eta\Lambda=M^{-1}\Theta\Lambda M$、ゆえに $\eta=M^{-1}\Theta M$ となる。

命題\ref{prop:twisted-cubic-secant-cubic}の二因子のうち、残る二根の一致を検出する双一次式を $H$ と書く。
その零点は、第二成分の逆グラフを含む。実際、順序を忘れた三次式は第一成分の直線と第二成分の二次曲線を含み、
後者の引戻しは元のグラフと逆グラフの和である。逆グラフは初めの二成分と異なるので $H$ に入る。
双一次式の次数により両者の式は比例する。残る二根の一致は $z=\eta t$ であるから、
$F$ の四係数にこの交代縮約を適用し、共通の非零倍を除くと
\[
 H=m Cuv-nAu-mAv-B/m .
\]
逆グラフは $u=(\Theta\Lambda^{-1}\Theta)v$ であり、その式を $G$ とすると
\[
 G=AC(\lambda^{-1}-\lambda)uv+
 (BC/\lambda+A^2\lambda^2)u-(A^2/\lambda+BC)v-AB(\lambda^{-1}-\lambda).
\]
$B,C$ が共に非零なら $uv$ と定数の係数比で $m^2=1$ となり不可能である。
$B\ne0,C=0$ なら $u,v$ の係数比で $m^2=-1/\lambda^2$、
$v$ と定数の比で $m^2(1-\lambda^2)=1$ となり矛盾する。
$C\ne0,B=0$ は四座標を一緒に逆数にして同じ場合に帰着する。
混合型でも $\Theta$ は対角である。

従って全ての隣接位置で二つの節点の像の非順序対は同じである。これを $0,\infty$ と正規化する。
全ての位置の割線三次曲線は、$J:x\mapsto-x$ の点対 $\{x,-x\}$ の直線 $L$ を含む。
実際、$\Theta$ が対角なら元の二成分は倍率変換のグラフである。
一方が $-1$ の倍率ならすでに $L$ がある。そうでなければ二次像は一致し、二つの非自明な倍率は互いの逆となる。
各成分から割線因子への次数は1で、終点でも同じだから $\eta$ は対角である
（反対角なら両成分が直線への二重被覆になる）。
$\Theta$ が反対角なら元の二グラフはともに $x\mapsto k/x$ 型の対合で、二つの相異なる直線を与える。
各成分での次数2により終点でも二つの相異なる対合が必要で $\eta$ は反対角である。

いずれの場合も残る縮約 $H$ は共通の物理座標 $p,q$ において、
斉次化すれば $h_1p+h_2q$ の形である。対角の場合は $F$ の $z,t$ の一次係数を使う。
反対角の場合は $uv$ と定数を使い、第二物理座標が $v$ の逆数の定数倍である。
原点と無限遠の同時交換でもこの形は保たれる。$H$ が恒等零なら割線三次式も恒等零となり不可能である。
一方の係数だけ非零なら、その二本の縦横の零線は順序交換で保たれず、その交換像は
元の二グラフにも入らないので三次式の対称性に反する。
対角で逆グラフどうしの場合、二次曲線の残余直線の像は $H$ から来なければならない。
反対角の場合も $H$ の倍率グラフは元の二つの対合と異なる。
よってこの倍率グラフは直線に写らなければならず、倍率は $-1$ である
（恒等なら対角二次曲線、その他の非対合も二次曲線に写る）。これで $L$ の包含を得た。

四位置の節点座標での同時 $J$ は六項の $F$ を保ち、物理座標でも同じ $J$ なので角括弧も保つ（射影的に）。
四移動も二節点の像の対を保つので $J$ と可換である。
従って割線三次の同型は $J$ と可換で、固定軌跡の正次元成分 $L$ は次の $L$ に写る。
実際、対称平方平面の固定点集合は $L$ と点 $\{0,\infty\}$ である。
これで両方向の割線列ができ、全次数正則性を得る。
\end{proof}

\begin{proof}[一つの二重交点]
二成分が一点で接するとする。各成分を第一射影でパラメータ化し、接点を無限遠に取る。
各位置で第一成分の前向きの座標を $x$ 自身に取ると、第二成分は接点の一次ジェットを共有するので、
四つの座標は $x,x+e,x+f,x+g$ となる。$e\ne0,f-e\ne0$ は隣り合うグラフの相異性による。
$z=v$ と $z=v+f-e$ 上で、$u,t$ の行列の二根は前後の座標である。
従って四テンソルの商 $F$ の制限は非零定数倍の $(u-v)(t-v)$ と $(u-v+e)(t-v+e-g)$ である。
倍率が定数なのは斉次化した同じ次数の切断であり、零行列もないからである。
$z$ で一次補間し、$v^3,v^2$ の係数を消すと二定数は等しく $g=f+e$ で、
\[
 F=(u-v)(t-v)+\frac{z-v}{f-e}\bigl(-f(u-v)+e(t-v)-ef\bigr)
\]
となる。第一成分の二根は $v=u$ で $z=u,t=u+f$ であり、
$t=u+f,u$ の組から来る逆側の根も $z'=u+f$ である。
従って第一成分の割線の移動は、始点の組 $(x,x)$ から終点の組 $(x+f,x+f)$ を使う。

物理座標の相対変換を同じく $\Theta,\eta$ とすると、始点の二グラフは
$\Theta$ と $\Theta T_e$（$T_e(x)=x+e$）である。
いずれかが恒等なら $\Theta$ は無限遠を固定する。ともに対合なら二つの跡の差が
$e\Theta_{21}$ なのでやはり固定する。ともに対合でも恒等でもなければ二次曲線の一致により
互いに逆で、$\Theta^2=T_{-e}$ だから唯一の固定点の無限遠を保つ。
混合なら対合の成分を第一に選ぶ。無限遠が固定されないと仮定し、四座標を一緒に平行移動して
$\Theta(x)=k/x,\ k\ne0$ とできる。第一成分の対合の移動から
$\eta=T_f\Theta T_{-f}$ である。残余縮約 $H$ は先と同じく第二成分の逆グラフを含まねばならない。
このグラフの式は $u=(\Theta T_e\Theta)v$、すなわち $euv+ku-kv=0$ である。
その上の $(u,v)=(0,0)$ では $F$ の二根は $(z,t)=(0,f)$ だが $\eta(f)=\infty$ で一致しない。
$H$ の消滅に反する。従って全辺で無限遠の像が一致し、接点を運ぶと物理座標が一定の $p$ になる。
全位置で $D_j=2p$ を使えば前節同様に両方向の割線列を得る。
\end{proof}

\begin{corollary}\label{cor:all-one-dimensional-point-graphs}
両側点ファイバーで一次元の点グラフを持つ枝では、全標準窓が有限の正則性・添字保持同型判定で被覆される。
純部分が空間有理正規三次となる最後の枝では、有限Frobenius検査と点グラフ検査を通った候補は、
埋込み次元2以下なら全て正則である。埋込み次元3では定理\ref{thm:maximal-tangent-ore}の判定を用いる。
三次曲線が唯一の一次元の台であることを課し、四つの完全三次系で座標をそろえると、この枝の
パラメータは四移動 $S$ と多重一次形式の射影類 $[f]$ である。許容域は四括弧の積と $f$ から作る
四テンソルについて上記の有限検査を課したものである。
同型条件は同時 $\mathrm{PGL}_2$ 作用による $S$ の共役と $[f]$ の対応であり、
周期は一・二位置回転について同じ同型条件を試し、いずれも失敗すれば4となる。
\end{corollary}
\begin{proof}
系\ref{cor:point-curve-high-coverage}、定理\ref{thm:elliptic-binary-cubic-pencils}、
系\ref{cor:binary-vertical-ternary}、定理\ref{thm:doubled-binary-pencils}と本定理を合わせる。
同型は唯一の一次元の被約曲線を保つ。最初の三辺の点列がそこで対角であるため
四位置の基底変更はこの曲線上で同じ射影変換を誘導する。完全三次系なので各周囲空間の射影変換もこれで決まる。
第四辺で $S$ が対応し、四テンソルの交叉直線から $[f]$ が対応する。逆も各関係の復元により従う。
標準四周期構造を忘れた同型と周期については定理\ref{thm:canonical-four}と系\ref{cor:period-orbits}を用いる。
\end{proof}

% !TeX root = ../papers/nc-p3.tex
\section{縦切片 $x=1$ の円分体表示と有理水準の完全判定}
\label{sec:rank-fifteen-vertical-one}
\begin{theorem}\label{thm:rank-fifteen-vertical-one-tube}
標数零の通常な階数15窓で、剰余標数3の付値が
$\nu(y)>0$, $4\nu(x-1)\ge17\nu(y)$ を満たすとする。
窓は標準Koszul左右AS正則で、各始点で四中心四次式上自由階数256である。
特に $x=1$ では全ての正付値 $y$ に適用できる。
\end{theorem}
\begin{proof}
前節の表示で $\bar J=0$ であり、$z$ の剰余は原始八乗根である。
$\zeta^4=-1$ として $\zeta$ の上の素点を $\bar\zeta^2+\bar\zeta+2=0$
で指定し、補題\ref{lem:rank-fifteen-ramified-pole}により $z\equiv\zeta$ と選ぶ。
同じ $H$ に対して $x=1$ になる表示を次のように作る。
\[
 D=(1-2H^2)^{-1},\qquad G=\sqrt{1+2H^2D^2},\qquad
 L_0^2=\frac{D^2}{G(G+1)},\qquad C_0^2=L_0^2-1 .
\]
$G$ の定数項は1、$L_0,C_0$ の定数項は
$(\zeta+\zeta^{-1})/2,(\zeta-\zeta^{-1})/2$ とする。
これらはすべて単元で、$z_0=L_0+C_0$ は単元である。
$G^2-1=2H^2D^2$ より $G(1-2H^2L_0^2)=1$ であり、所定の縦切片となる。

\path{computations/verify_rank_fifteen_vertical_one.py} は、
$\mathbb Q[Z]/(Z^4+1)$ の係数で十七桁を取り、全ての $H$ 除算の定数項が厳密に零で
あることを検査する。単元判定には $c\equiv1\pmod3$, $c^2=-2$ を持ち上げ、
$Z^4+1=(Z^2+cZ-1)(Z^2-cZ-1)$ の第一因子を使う。
係数の共通分母の3部分が $3^e$ なら、分子を第一因子で $3^{e+1}$ まで還元する。
剰余の二係数が $3^e$ の倍数であることを確かめ、その商の少なくとも一方が
3の単元となる主元を使う。これで分母に現れる他の素点での極も扱える。
未切断経路、中心、五次の中心商に合計十桁を使う。長さ6から13の
損失列は $(1,0,1,1,2,0,1,0)$ で、一桁を残して十三次で消える。

元の $x$ については $x=G^2(1-2H^2L^2)^2$ である。
二乗の差を取ると、$1-2H^2L^2$ と $1-2H^2L_0^2$ の和、
および $L+L_0$ はいずれも単元である。
また $2(L-L_0)=(z-z_0)(1-1/(zz_0))$ の後の因子も単元なので
\[
                  \nu(x-1)=2\nu(H)+\nu(z-z_0)
\]
となる。$x=1$ なら同じ枝で $z=z_0$ である。
定理\ref{thm:rank-fifteen-collision-tube}と同じく、経路係数の変化は
$H^2(z-z_0)$ の倍数である。仮定と $\nu(y)=4\nu(H)$ から
実際の窓と円分体の表示は $H^{17}$ を法として一致する。
主元を持ち上げ、中心商の上界と一般の下界を合わせると全次数の結論を得る。
\end{proof}

\begin{corollary}\label{cor:rank-fifteen-vertical-one-rational}
標数零の代数閉体で $x=1$, $y\ne0$ とする。$y$ が代数的で原始整数最小多項式
$f$ の斉次化が $\mathbb P^1(\mathbb F_3)$ に根を持つなら、その水準の全窓が
定理の正則性と自由性を持つ。特に有理水準 $y\in\mathbb Q^\times$ の全てで成り立つ。
縦切片の添字保持同型類の座標は $t\in\mathbb G_m\setminus\{1,-1\}$ で、上記の各水準では
\[
                    y=-\frac{(t-t^{-1})^2}{108}
\]
の解がそのまま正則窓の添字保持同型類を与える。$y=1/27$ なら二類、それ以外の
上記水準では四類であり、全類の最小周期は2である。
\end{corollary}
\begin{proof}
ある3進素点で $y$ が正付値なら本定理、負付値なら
$\nu(y)<2\nu(x)$ の極の定理を使う。全ての共役根が単元の場合には、
原始式の先頭と定数はいずれも3の単元であり、斉次根は $0,\infty$ 以外である。
従って仮定から剰余根 $1$ または $-1$ があり、その上の
共役根を選べる。定理\ref{thm:rank-fifteen-unit-residues}が適用できる。
一次式の斉次化は必ず素体上に根を持つ。

この切片で $p^2=1-27y$, $q^2=-4y$ より $p^2-27q^2/4=1$ である。
$t=p+3\sqrt3\,q/2$ と取れば $p=(t+t^{-1})/2$,
$q=(t-t^{-1})/(3\sqrt3)$ となる。$q\ne0$ なので通常表示外の四点は含まれず、
定理\ref{thm:rank-fifteen-rational-quotient}と第\ref{sec:rank-fifteen-global-invariants}節の
全点分離より、相異なる解は相異なる添字保持同型類を与える。
一位置移動は $t\mapsto-t$ で固定点がなく、$p=0$ の水準だけで解は二つになる。
\end{proof}

% !TeX root = ../papers/nc-p3.tex
\section{座標1の単元剰余を表す楕円曲線と三次までの全水準}
\label{sec:rank-fifteen-vertical-unit-levels}
\begin{lemma}\label{lem:rank-fifteen-vertical-elliptic-residues}
標数3の $\bar x=1$, $\bar y\ne0$ の剰余表示は
\[
 \bar y=(v-v^3)^2,\qquad
 k(v)\bigl(\sqrt v,\sqrt{v-1},\sqrt{v+1}\bigr)
\]
から得られる。ただし $k$ は代数閉体とする。この関数体の滑らかな完備曲線は
種数1で、$\bar y$ 直線への写像は次数48のGalois被覆である。
\end{lemma}
\begin{proof}
$K=1/v$, $U=K^2$, $H=K-v$, $V=H^2$ とすれば $V(U+3)/(U-1)^2=1$,
$V^3/(U-1)^4=(v-v^3)^2$ となる。$\bar y\ne0$ なら $v\ne0,\pm1$ である。
$i^2=-1$, $w^2=v$ とし、$A+B=\sqrt{1+v}$, $A-B=\sqrt{1-v}$ で $A,B$ を取る。
$A^2+B^2=1$, $2AB=v$ より $E=A^2-B^2$ は $E^2=1-v^2$ を満たす。
$n=EA/v$ とすると $H/n=-EB/v$,
$n^2+(H/n)^2=K^2-1$ となる。他の四根は順に
\[
 P_+=(v+1)/w,\quad P_-=i(A+B),\quad Q_+=(v-1)/w,\quad Q_-=A-B
\]
と取れる。いずれも $v\ne0,\pm1$ で使える。

$v,v-1,v+1$ の平方類は三つの相異なる単純零点での付値により独立である。
従って $v$ 直線への被覆は次数8で、$0,1,-1,\infty$ にのみ分岐する。
各点では単元は完備局所体で平方なので慣性指数は2である。
Riemann--Hurwitz公式から $2g-2=-16+4\cdot4=0$ となる。
一方 $v\mapsto\epsilon v+b$ ($\epsilon=\pm1$, $b\in\mathbb F_3$) は
$(v-v^3)^2$ を保つ六変換である。三つの被開平式を符号を除き置換するため
いずれも次数8の被覆へ持ち上がり、全次数48だけの自己同型を与える。
具体的には $A=(1-t^2)/(1+t^2)$, $B=2t/(1+t^2)$ と置けるので、
$W^2=t-t^3$ の表示では $w=W/(1+t^2)$ となる。
\end{proof}

\begin{theorem}\label{thm:rank-fifteen-vertical-unit-levels}
標数零の通常な階数15窓で、3の上のある付値で $\bar x=1$ とし、
$\bar y\ne0$ の $\mathbb F_3$ 上の次数が高々3とする。この窓は標準Koszul左右AS正則で、
各始点で四中心四次式上自由階数256である。
\end{theorem}
\begin{proof}
次数1は定理\ref{thm:rank-fifteen-unit-residues}による。
\path{computations/verify_rank_fifteen_vertical_unit_levels.py} は素体上の
モニックな二次・三次式を列挙し、既約な三つの二次式と八つの三次式の
それぞれについて有限体表示を検証する。必要な拡大次数は $6,8,12,18$ である。
指定した $\bar y$ の最小多項式、$v=1/K$, $H=K-v$ の二式と全ての平方根を照合する。
各表示で未切断経路の長さ0から5の次元は $(1,4,10,20,35,56)$、
四次中心の文字別次元は $(3,2,0,0)$ である。文字0の三つの自由座標の後二つと、
文字1の二つの自由座標を選び、四始点全てで長さ13までの中心商の上界を得る。

所定の $x,y$ へはまず $(U+3)^3=(x^3/y)(U-1)^2$ の指定した根を持ち上げる。
$U-1$ は単元であり、剰余で微分は非零である。
$V=x(U-1)^2/(U+3)$ も単元で、残る根はモニックな式から順に持ち上がる。
単元主元、一般の経路と中心の下界、中心商の上界の一致は
定理\ref{thm:rank-fifteen-unit-residues}と同じ全次数の証明を与える。
有限体の共役によりこれら十一式の全ての根に適用できる。
\end{proof}
\begin{corollary}
$x=1$, $y\ne0$ で $[\mathbb Q(y):\mathbb Q]\le3$ とする。
全ての添字保持同型類は標準Koszul左右AS正則で、
系\ref{cor:rank-fifteen-vertical-one-rational}の $t$ の式による四類
（$y=1/27$ のみ二類）であり、最小周期はすべて2である。
より一般に $y$ の原始整数最小多項式の斉次化が、3で還元して次数3以下の
閉点を零点に持てば同じ結論が成り立つ。
\end{corollary}
\begin{proof}
3での正付値と負付値は前節で尽くされた。全共役が単元ならば、
所定の剰余閉点の上の共役根を選び、本定理を適用する。次数が高々3なら、
還元した原始斉次式の零点のいずれも次数が高々3である。
\end{proof}

% !TeX root = ../papers/nc-p3.tex
\section{階数15の縦切片 $x=1$ の完全分類}
\label{sec:rank-fifteen-vertical-all-levels}
\begin{theorem}\label{thm:rank-fifteen-vertical-all-levels}
剰余標数3の付値で $\bar x=1$, $\nu(y)=0$ となる標数零の通常な階数15窓は、
全て標準Koszul左右AS正則で、各始点で四中心四次式上自由階数256である。
\end{theorem}
\begin{proof}
第\ref{sec:rank-fifteen-vertical-unit-levels}節の楕円曲線 $W^2=t-t^3$ 上で、
$i^2=-1$ とし、係数を $1,i,W,iW$ の四座標で保持する。
\path{computations/verify_rank_fifteen_vertical_all_levels.py} は
$\mathbb F_3(t,i,W)$ 上で未切断経路を五次まで構成し、文字0と1の中心方程式を解く。
前節と同じ四本の中心を課して四始点の商を十三次まで消去する。
各文字で必要な本数の主元を選び、逆元を取るたびに分子の $\mathbb F_3(t)$ への
ノルムを記録する。入力の逆元から帰納すると、これらのノルムの零点を除けば
全ての除算が可能である。得られた平方自由な積 $F(t)$ の次数は 204451 である。

$N=t(1-t^2)$, $D=(1+t^2)^2$ と置けば、$v=N/D$,
$\bar y=P/Q$ において $P=(ND^2-N^3)^2$, $Q=D^6$ である。
この次数24の写像で $P,Q$ は互いに素で、$P'Q-PQ'$ の零点は $PQ=0$ 上に限る。
従って有限の非零な $\bar y$ の各ファイバーは相異なる24点を持つ。
$t\mapsto-t,1/t,(1+t)/(1-t)$ のいずれでも $P/Q$ は変わらない。
$F$ から $PQ=0$ の点を除き、これら三変換の引き戻しとの最大公約元を反復すると、
安定した式の次数は 784 である。全ての表示が $F=0$ に入るファイバーはこの式に含まれる。
残る各既約因子で $P/Q$ の最小多項式 $g$ を求め、同じ $g$ に対応する因子の積が
$Q^{\deg g}g(P/Q)$ のモニック化と等しいかを厳密に比較する。

全表示が除かれる剰余軌道は、$Y+1$, $Y-1$, $Y^3+Y^2+2Y+1$,
$Y^6+Y^3+2Y^2+2Y+1$ の四つに限る。一次と三次は前節の証明で認証済みである。
六次では $\mathbb F_{3^6}$ の表示で最小多項式と七つの倍率座標を直接検査する。
未切断経路の長さ0から5、中心の文字別次元 $(3,2,0,0)$、
四中心商の四始点の長さ13までの次元を取り直す。有限体の共役で六根が尽くされる。

従って各剰余値では少なくとも一表示で全ての主元が取れる。
未切断経路・中心の一般の下界、中心商の上界と持ち上げは
定理\ref{thm:rank-fifteen-unit-residues}の通りであり、剰余値が超越的なときも
非零な $F$ の零点を外れるため同じ判定が成り立つ。
\end{proof}

\begin{corollary}
標数零の代数閉体において、通常な階数15の縦切片 $x=1$ の全点が上記の正則性と自由性をもつ。
添字保持同型類は $t\in\mathbb G_m\setminus\{1,-1\}$ で一意に表され、
水準 $y=-(t-t^{-1})^2/108$ ごとに $y=1/27$ では二類、それ以外の $y\ne0$ では四類である。
全類の最小周期は2である。
\end{corollary}
\begin{proof}
係数体に3進付値を延長する。$\nu(y)>0$ なら
定理\ref{thm:rank-fifteen-unit-pole-depth}で $J=0$、$\nu(y)<0$ なら極の定理、
$\nu(y)=0$ なら本節の定理を適用する。同型座標と一位置移動
$t\mapsto-t$ は系\ref{cor:rank-fifteen-vertical-one-rational}による。
\end{proof}

% !TeX root = ../papers/nc-p3.tex
\section{一つの交代証明書の剛性}
\label{sec:single-negative-rigidity}
ここでは単一の交代証明書を持つ残余も決定する。
行列因子化に用いるClifford代数の議論は線型代数として証明し、AS対象にClifford表示を仮定しない。

\begin{lemma}[非退化階数2行列空間の圧縮形]\label{lem:rank-two-square-compression}
一次式$4\times4$行列$M(z)$の一般階数が2で、全ての像の和が$k^4$、全ての核の交叉が零なら、
定数の行列基底変更によって、その非零成分は一行または一列にしか存在しない形になる。
\end{lemma}
\begin{proof}
補題\ref{lem:no-rank-two-four-by-six}と同じ階数2基点のブロック分解で、
$D=0$, $CB=0$を得る。ここでは$B,C$はともに$2\times2$である。
$C$の一般階数0または2、共通像直線の場合は、それぞれ共通像または共通核の仮定に反する。
従って$C=c(z)\lambda$、$B=u b(z)$、$ku=\ker\lambda$となる。
Schur補空間の次の係数から$CAB=0$となり、$c,b$が非零なので$\lambda A u=0$である。
$u=e_1$, $\lambda=e_2^*$とすれば、行列の可能な成分は第一行と第二列に限られる。
\end{proof}

\begin{lemma}[Klein二次式の$4\times4$行列因子化]\label{lem:klein-four-matrix-factorization}
$z=(z_{01},z_{02},z_{03},z_{12},z_{13},z_{23})$、
$q=z_{01}z_{23}-z_{02}z_{13}+z_{03}z_{12}$とする。
一次式$4\times4$行列$M(z)$が一般階数4で、$q=0$上の一般階数が2以下なら、
定数の可逆な左右変換によって標準交代行列$X(z)$またはそのPfaff随伴$X^\#(z)$になる。
ここで$X_{ij}=z_{ij}$ ($i<j$)、$X_{ji}=-z_{ij}$であり、
$XX^\#=X^\#X=qI_4$と符号を定める。
\end{lemma}
\begin{proof}
三次小行列式は全て既約二次式$q$で割れる。
従って$\operatorname{adj}M$は$q$の倍であり、$\det M$とその全偏微分も$q$で割れる。
標数零で$q$は非退化なので、$\det M=cq^2$、$c\ne0$となる。
$N=\operatorname{adj}M/(cq)$は一次式行列で、$MN=NM=qI_4$である。
従って
\[
 \Gamma(z)=\begin{pmatrix}0&M(z)\\N(z)&0\end{pmatrix},\qquad \Gamma(z)^2=q(z)I_8
\]
は六次元非退化二次空間のClifford代数の八次元表現を与える。
このClifford代数は$\operatorname{Mat}_8(k)$である。
実際、双曲基底の生成元を三変数外積空間$\bigwedge^*k^3$上の外積・縮約として作用させると、
それらは全ての行列単位を生成する。具体的には各変数の数作用素とその補を掛けて各外積基底への射影を作り、
外積・縮約で任意の二基底間の行列単位を得る。一方、Clifford関係で整列した単項式は高々$2^6=64$個なので、
この全射は64次元行列代数への同型である。

従ってこの八次元表現と$(X,X^\#)$から作る標準八次元表現は同型である。
次数符号行列$J=\operatorname{diag}(I_4,-I_4)$は全生成元と反可換である。
表現同型で移した$J$と標準$J$の積は全行列と可換なのでスカラーであり、
両者の平方が1であることからそのスカラーは$\pm1$となる。
従って表現同型は二つの四次元次数部分を保つか交換し、
$M$は左右変換で$X$または$X^\#$となる。
\end{proof}

$E_{cd}$を後二因子の行列単位とし、$ij=e_i\wedge e_j$と略記する。
二つの圧縮テンソルを
\[
 \begin{aligned}
 T_4&=\sum_{j=1}^3 A_j\otimes(E_{0j}+E_{j0})+03\otimes E_{00},
       &&(A_1,A_2,A_3)=(01,02,12),\\
 T_5&=\sum_{j=1}^3 A_j\otimes E_{0j}+\sum_{j=1}^3 B_j\otimes E_{j0},
       &&(B_1,B_2,B_3)=(01,03,13)
 \end{aligned}
\]
と置く。最初の二因子は既に同一視されている。

\begin{theorem}[部分交代四テンソルの三正規形]\label{thm:partially-alternating-three-forms}
$w\in\bigwedge^2V\otimes W\otimes U$の全一因子階数を4とし、
$02|13$と$12|03$の切断階数を6とする。
共通の$\operatorname{GL}(V)$と後二因子の独立基底変更により、$w$は
交代体積テンソル、$T_4$、$T_5$のいずれかに移る。
$01|23$階数はそれぞれ6,4,5である。
\end{theorem}
\begin{proof}
$02|13$の階数分解で作る$P(x),Q(x)$は
$P(x)Q(y)+P(y)Q(x)=0$を満たす。
補題\ref{lem:no-rank-two-four-by-six}を$P,Q^{\mathsf T}$へ適用すると、
系\ref{cor:single-negative-ranks-three}と同じ理由で一般階数は両方3となる。

$P(x)$の左核と$Q(x)$の右核の射影点を$c(x),d(x)$とする。
どちらも定数なら対応する一因子階数が4未満となるので、非定数である。
$P(x)Q(y)=-P(y)Q(x)$は$c(x),c(y)$を左核、$d(x),d(y)$を右核に持つ。
一般の$x,y$で各対は独立なので、その階数は2以下である。
従って$w$の最初の二因子を二ベクトル$z=x\wedge y$で縮約した行列$M(z)$は、
Klein二次曲面$q(z)=0$上で階数2以下となる。

$M$の一般階数3は不可能である。
その場合$\operatorname{adj}M=q L(z)$で、$L$は非零の一般階数1の一次式行列となる。
階数1行列空間の共通像・共通核の二分と$ML=LM=0$から、
$M$に共通右核または共通左核が生じ、一因子階数4に反する。
一般階数1以下も同じ二分から不可能である。

一般階数4なら前補題を適用する。
$X(z)$が与えるテンソルは$w_{abcd}=\delta_{ac}\delta_{bd}-\delta_{ad}\delta_{bc}$で、
$02|13$作用は行列$H$を$\operatorname{tr}(H)I-H^{\mathsf T}$へ送る。
その核は零なので階数16となり、仮定に反する。
$X^\#$が与えるものは交代体積テンソルの非零スカラー倍である。

残る一般階数2では、圧縮形の補題により
\[
 w=\sum_{j=1}^3 A_j\otimes E_{0j}+\sum_{j=1}^3 B_j\otimes E_{j0}+C\otimes E_{00}
\]
と書ける。$A_j,B_j,C$は四次元$V$の交代行列である。
後二因子の一因子階数4から、$A_1,A_2,A_3$と$B_1,B_2,B_3$はそれぞれ独立である。
$02|13$行列はブロック$\left(\begin{smallmatrix}C&A\\B&0\end{smallmatrix}\right)$を持つ。
独立な三交代行列の像の和は三次元以上なので、$\rk A,\rk B\geq3$である。
全階数6から両方3となり、ある三空間$H_A,H_B\subset V$について
\[
 \langle A_j\rangle=\bigwedge^2H_A,\qquad
 \langle B_j\rangle=\bigwedge^2H_B
\]
となる。また階数6の等号から$C(\ker B)\subset\operatorname{im}A=H_A$となる。

$H_A\ne H_B$なら、その零化共変量を$\alpha,\beta$とした条件は$\alpha C\beta=0$であり、
これは$C\in\bigwedge^2H_A+\bigwedge^2H_B$と同値である。
後二因子のshearで$C$を消し、二三空間と各基底を標準化すれば$T_5$となる。
$H_A=H_B=H$なら、$C\in\bigwedge^2H$は最初の一因子階数4に反する。
$C$の$\bigwedge^2H$成分をshearで消せば、残りは$e_3\wedge v$、$0\ne v\in H$である。
$H$内の基底と残るスカラーを変えて$T_4$を得る。
両正規形の第一切断階数は、それぞれ張る二形式の次元4と5である。
\end{proof}

\begin{corollary}[単一負符号証明書の認識]\label{cor:one-negative-frobenius-recognition}
有限Frobenius候補の三つの$S_j^-(w)$のうち一つにでも可逆元があれば、
その四窓は可換多項式のindexed類である。
この認識には元代数のKoszul性を前提としない。
従って既知の完全対称枝と可換多項式枝の外の未分類AS残余では、
可逆な正符号証明書は高々一つであり、可逆な負符号証明書は一つも存在しない。
\end{corollary}
\begin{proof}
隣接証明書の場合は補題\ref{lem:adjacent-witness-rank-six}により三切断とも階数6である。
対角証明書の場合はそれ以外の二切断が元の隣接切断なので階数6である。
従って証明書の二因子を最初の二因子へ置いて前定理を適用できる。
隣接二因子を交換した場合、$T_4,T_5$はその切断階数4,5から既に除かれる。
対角の二因子の場合は巡回始点を選んで元の順序へ戻す。
このとき両圧縮形で$R_0\supset H_A\otimes ke_0$、$R_1\supset ke_0\otimes H_A$である。
実際、$d\ne0$の$A_j\otimes E_{0j}$部分だけを縮約すると、各々この三次元空間を得る。
従って$W_{0,3}$は$H_A\otimes ke_0\otimes H_A$の独立な九つの単項式を含み、
有限Frobenius条件の次元4に反する。
従って残るのは交代体積テンソルだけであり、可換多項式のindexed類である。
正確な交叉次元は両圧縮形で$(10,16),(8,16),(10,16),(8,16)$となる。
有理数の関係・全核基底と四次までの検算は
\path{computations/verify_negative_witness_rigidity.py}に置く。
同じコードは$XX^\#=qI$、標準Clifford像の次元64、恒等二形式型の第三階数16も検算する。
\end{proof}

\section{基点を持つ四二次式系の支配的チャート}
\label{sec:product-net-boundary}
ここでは最初の三二次式も動かし、基点を持つ四次元二次式空間の一般点を覆う。
この被覆はGrassmann多様体上の接写像で証明し、有限個の標本から正則性を推測しない。

$p=(0,0,0,1)$で消える二次式の九次元空間を$T_p$、その四次元部分空間の多様体を
$Y=\operatorname{Gr}(4,T_p)$とする。$\dim Y=20$である。
次の八係数族を取る：
\begin{align*}
 f_0&=e_0^2-e_0e_3-\alpha e_0e_1,&
 f_1&=e_1^2-e_1e_3-\beta e_1e_2,\\
 f_2&=e_2^2-e_2e_3-\gamma e_2e_0,&
 q&=e_0e_3+u e_1e_3+v e_2e_3+a e_0e_1+b e_0e_2+c e_1e_2.
\end{align*}
$H\subset\operatorname{GL}(V^*)$を、$p$で消える一次式の三次元空間を保つ部分群とする。
これは13次元の滑らかで既約な群であり、共通スカラーは$Y$上自明に作用する。

\begin{theorem}[八係数積型netチャートの支配性]\label{thm:product-net-dominance}
射$H\times\mathbb A^8\to Y$、$(g,\theta)\mapsto g\langle f_0,f_1,f_2,q\rangle$は支配的である。
さらに最初の三式が相異なる八点完全交叉で、第四式が$p$だけで消える部分へ制限しても支配的である。
これは二次式系の一般点の被覆であり、indexed AS同型類の次元やAS正則性の認定ではない。
\end{theorem}
\begin{proof}
四式は最初の三平方係数と$q$の非零係数から独立であり、全係数で$Y$への射が定まる。
係数値
\[
 (\alpha,\beta,\gamma,u,v,a,b,c)=(2,3,5,2,3,5,7,11)
\]
と$g=I$で接写像を計算する。
$T_p/L$の五次元商へ各四式の変分を射影すれば、
接空間$\operatorname{Hom}(L,T_p/L)$の20座標が得られる。
$H$の13個の行列単位変分と八係数の変分を並べた行列は$20\times21$であり、
列$0,1,\ldots,11,13,14,\ldots,20$による20次小行列式は
\[
 -\frac{696}{78125}\ne0
\]
となる。単項式順と商基底の選択、微分行列の全構成は
\path{computations/verify_product_net_dominance.py}に置く。
従って接写像は全射であり、滑らかな既約な$Y$への射は支配的である。

この点で最初の三式の八点は、$e_3=1$として
\[
\begin{gathered}
(0,0,0),\ (1,0,0),\ (0,1,0),\ (0,0,1),\\
(3,1,0),\ (1,0,6),\ (0,4,1),\ (-9/29,-19/29,-16/29)
\end{gathered}
\]
である。無限遠には、三変数の積の連鎖から
$\alpha\beta\gamma=1$でなければ非零解はなく、この点では$30\ne1$である。
上の八点は相異なり、従って三二次式の次数8の完全交叉を尽くす。
$q$の他の七点での値は
\[
1,\ 2,\ 3,\ 20,\ 61,\ 55,\ 2452/841
\]
と全て非零である。
この条件はパラメータの開条件なので、制限した射も同じ全射接写像を持ち、支配的である。
\end{proof}

\begin{proposition}[支配的チャートから全基点域への障害の移送]
\label{prop:product-net-universal-obstruction-transfer}
上記八変数の一般点における三次双対フレーム・テンソルの五次$d_2$階数を$r$とする。
$r<84$が厳密に証明されれば、基点を持つ$L$から作る三因子対称テンソルは
標準型Koszul AS点を一つも含まない。
より正確には、その有限Frobenius・四次正確部分で始点0の$H_2$次元は全て$84-r$以上である。
ここで$L$に基点が一つだけであることや、相異なる八点netを持つことを仮定する必要はない。
\end{proposition}
\begin{proof}
まず基点を$p$に固定する。
$Y$上の可換二次商$B$について$\dim B_3\geq4$であり、等号となる部分$Y_3$は空でない開部分である。
実際、次数3の関係挿入は16本であり、Boolean標本で等号を達成する。
$Y_3$上の三次双対は階数4のベクトル束であり、そのフレーム束$\mathcal F$は既約である。
$\mathcal F$から三因子対称テンソルへの写像は正則である。
その有限Frobenius・四次正確部分$\mathcal U$は、二対二階数の上界6から開であり、既知標本を含む。

前定理の支配的射を$Y_3$へ制限し、さらに戻り因子の$\operatorname{GL}_4$基底変更を許せば、
$\mathcal F$への支配的な射が得られる。
八係数族で選んだ三次座標正規化が定まる部分も空でない開部分である。
共通の三因子変換と戻り因子の基底変更は、五次微分の階数を保存する。
従って八変数一般点の階数$r$は$\mathcal U$の一般階数でもある。
全ての$\mathcal U$上で$\operatorname{rk}d_2\leq r$、$\operatorname{rk}d_3\leq36$となり、
$\dim H_2\geq120-r-36=84-r$を得る。

最後に、有限Frobeniusな三因子対称テンソルでは、対称な三連続因子の三次交叉が
$B_3^*$に一致し、次元は4である。従って必ず上の$Y_3$のフレームから来る。
基点が他の場所なら共通の射影基底変換で$p$へ移せる。
AS点は有限Frobenius・四次正確部分に入るので、$r<84$なら全て排除される。
\end{proof}

この条件付き移行は一般点の精密な階数検算にも使えるが、基点枝全体のAS排除は
定理\ref{thm:partial-symmetric-basepoint-exclusion}で、一般点計算を介さずに得られた。

\section{基点での一次部分を使う三変数三次縮約}
\label{sec:ternary-cubic-reduction}
八係数チャートの$f_i$を$x_i h=g_i(x)$と書き、$h=e_3$とする。
ここで
\[
 g_0=x_0^2-\alpha x_0x_1,\qquad
 g_1=x_1^2-\beta x_1x_2,\qquad
 g_2=x_2^2-\gamma x_2x_0.
\]
第四式を最初の三式で変更すれば、純粋な三変数二次式
\[
 Q=q+f_0+u f_1+v f_2
 =g_0+u g_1+v g_2+a x_0x_1+b x_0x_2+c x_1x_2
\]
になる。

\begin{proposition}[七本の三変数三次条件]\label{prop:ternary-cubic-reduction}
$B_3^*$は、基点$p=(0,0,0,1)$のcube$p^{\otimes3}$の一次元空間と、次の七つの三変数三次式を消す
対称三テンソルの空間との直和として表示できる：
\[
 Qx_i\ (i=0,1,2),\qquad
 g_i x_j-g_j x_i\ (i<j),\qquad
 \mu(Q),
\]
ここで$\mu(x_a x_b)=g_a x_b$（$a\leq b$）とする。
純三次座標を$c_{abc}$とすると、残りの座標は
\[
 c_{ijh}=c(g_i x_j),\qquad c_{ihh}=c(\mu(g_i)),\qquad c_{hhh}=0
\]
で復元される。$i,j$の選択の違いは表示した交差条件で一致する。
一般点では七条件は独立であり、三つの座標$012,111,222$による正規化で三次元基底を取れる。
これは二次式係数のチャートの縮約であり、indexed同型類の次元を主張しない。
\end{proposition}
\begin{proof}
元の条件は$f_i$または$q$に任意の一次式を掛けた16本の消滅条件である。
$f_i x_j$から六つの$x_i x_j h$座標を純三次座標で消去し、
同じ座標の二通りの消去の整合性が三つの$g_i x_j-g_j x_i$となる。
$f_i h$は三つの$x_i h^2$座標を消去する。
残る$Qx_i$と$Qh$の条件は、それぞれ表示の三条件と$\mu(Q)$である。
$h^3$座標には全ての式が作用しないので、基点cubeが独立に残る。
従ってこの操作は必要十分であり、九座標の消去と七条件を与える。

七条件と三つの正規化座標からなる$10\times10$行列は、
$\alpha=\beta=\gamma=a=b=c=0$, $u=v=1$で可逆である。
この点では$x_i^3=x_i h^2$なので、既知の$012,133,233$正規化フレームに一致する。
従って一般点でも可逆であり、七条件は独立となる。
\path{computations/ternary_cubic_frame.py}はこの10座標系のみを解き、
復元後には元の全16条件も検算する。
\end{proof}

戻り辺Torの$Q_2$には、六つの$ x_i h-g_i$、$h x_i-g_i$、
三つの三変数交換子、および$Q$を基底として取れる。
これは元の六交換子と四二次式の基底変更であり、$h$への縮約の左右を分離する。
\texttt{verify_trivial_extension_bitor.py --frame ternary --point-quotient}は、
この表示と補題\ref{lem:one-return-point-quotient}を組み合わせる。
二変数の全零合成・階数は再検算済みであり、八変数の最終階数はなお未確定である。

\section{7変数の代数的族と一般点の最小周期2}\label{sec:algebraic-family}
ここでは形式的な持上げから離れ、全パラメータで定義された代数的族を構成する。
基底体は引き続き標数零とする。基底を
\[
 S=\operatorname{Spec}k[u,v,r,s,a,b,c]\simeq\mathbb A^7
\]
とし、位置$i$で$\varepsilon=(-1)^i$と置く。積は隣接する一次生成空間の積を表す。
\begin{equation}\label{eq:seven-family}
\begin{aligned}
 YW+WY&=0,& ZW+WZ&=0,& ZY+YZ&=W^2,\\
 XW+WX&=uY^2+vZ^2+\varepsilon a(YZ-\tfrac12W^2),\\
 XY+YX&=W^2+Z^2+rY^2+\varepsilon(aWY+bWZ),\\
 XZ+ZX&=sY^2+\varepsilon cWY.
\end{aligned}
\end{equation}
この二次$\ZZ$-代数を$\mathcal C$と書く。
点$p_0=(0,1,0,0,0,0,0)$では$\mathcal C_{p_0}=C_0$である。
$u=r=s=a=b=0,\ v=1$に制限すると、既存の2周期族$A(c)$を得る。

\begin{theorem}[全ファイバーで正則な代数的族]\label{thm:algebraic-seven}
式\eqref{eq:seven-family}の$\mathcal C$は$S$上成分ごとに自由であり、
各$\mathcal C_{i,i+n}$の基底は
\[
 \{W^{h}Y^{j}Z^{\ell}X^{m}:h+j+\ell+m=n\}
\]
である。任意の幾何学的ファイバーは行Hilbert級数$(1-t)^{-4}$の
標準14641型quadratic Koszul AS正則$\ZZ$-代数である。
全ファイバーが2周期的である。ここには一般性条件や分母の非零条件を課さない。
\end{theorem}
\begin{proof}
$W<Y<Z<X$とし、逆順の二次語を順序付き語へ書き換える。
補助重みを$(1,2,2,4)$とし、重みが同じなら辞書式順序で比較する。
各規則の$-x_ax_b$以外の項は左辺より真に重みが小さく、
$-x_ax_b$も辞書式順序を小さくする。
従って全パラメータで還元は停止する。
重なりは$ZYW,XYW,XZW,XZY$の四つで、始点の偶奇を含めて8個である。
全ての差は$k[u,v,r,s,a,b,c]$上で恒等的に零となる。
再現コード\path{computations/verify_algebraic_family.py}は、特殊化や次数打切りを行わず
8個を多項式として検査する。Diamond lemma\cite[Theorem 1.2]{Bergman}から主張の自由基底を得る。

この恒等式には次の構造的説明もある。
\[
 B=k\langle W,Y,Z\rangle/(YW+WY,ZW+WZ,ZY+YZ-W^2),\quad T=YZ-\tfrac12W^2
\]
と置く。$W^2,Y^2,Z^2$は中心的で、$T$は$W$と可換、$Y,Z$と反可換である。
通常の次数1の導分$d$を
\[
 d(W)=aT,\qquad d(Y)=aWY+bWZ,\qquad d(Z)=cWY
\]
で定められる。最初の二関係の像は零であり、第三関係については
\[
 \{d(Z),Y\}+\{Z,d(Y)\}=2aWT=\{d(W),W\}
\]
となる。ここで$\{x,y\}=xy+yx$である。
一方、$\sigma_B(W,Y,Z)=(-W,-Y,-Z)$に対し、
\[
 p_W=uY^2+vZ^2,\quad p_Y=W^2+Z^2+rY^2,\quad p_Z=sY^2
\]
は、$\delta_p(fh)=\delta_p(f)h+\sigma_B(f)\delta_p(h)$を満たす$\sigma_B$-導分を定める。
各$p_j$が中心的なので関係の像は交換子の和として零になる。
$g^2=1$、$gW=-Wg,gY=-Yg,gZ=-Zg$を追加した環に
\[
 \sigma(g)=-g,\quad\sigma(W,Y,Z)=(-W,-Y,-Z),\quad
 \delta(g)=0,\quad\delta(j)=p_j+g d(j)
\]
を定めると、$\delta$は$\sigma$-導分である。
Ore拡大$Xh=\sigma(h)X+\delta(h)$の二つの冪等元$(1\pm g)/2$に沿った
展開が式\eqref{eq:seven-family}である。これも全パラメータでの基底を説明する。

AS正則性をHilbert級数だけから結論しない。
二次Gr\"obner基底からKoszul性が従う\cite[Theorem 5.3]{Priddy}。
そのbar複体のフィルトレーションによる証明は、各経路の始点・終点を固定すれば$\ZZ$-代数にも適用できる。
補助重みの随伴次数代数は、全ての相異なる生成元が反可換する四変数の量子多項式代数である。
二次双対$E$には逆重みのフィルトレーションを入れる。
二次関係の直交補空間の初期空間を取ることで、その随伴次数代数は
\[
 k[\xi_W,\xi_Y,\xi_Z,\xi_X]/(\xi_W^2,\xi_Y^2,\xi_Z^2,\xi_X^2)
\]
の対応する$\ZZ$-版からの商となる。
Koszul性とHilbert級数の逆数公式により$\dim E_{i,i+n}=\binom4n$なので、この全射は同型である。
随伴次数代数の相補次数積は完全である。
任意の非零な先頭項に相補単項式を掛ければ最高次数で非零になるから、
元の$E$の相補次数積も非退化であり、同じ次元より完全である。
従って局所Frobenius性が全ファイバーで成り立つ。
左右のKoszul複体をこの完全対合で双対化すると、単純加群に対するExtは
次数4だけに一次元存在し、左右AS条件と14641型解消を得る。
Koszul双加群解消の長さも4であり、大域次元は4となる。
係数が偶奇だけに依存するので厳密2移動は同型である。
\end{proof}

\begin{theorem}[同型類の7次元性と三つの交代方向]\label{thm:seven-image}
この族のindexed $\ZZ$-代数同型類の像は7次元である。
正確には、稠密開集合上でパラメータの同型類への写像は一般に有限対一であり、
この意味で有効パラメータは7個ある。
$p_0$における接方向は偶4方向と交代3方向に分かれ、いずれも同型で消えない。
交代3方向の前節の座標は
\begin{equation}\label{eq:seven-odd-plane}
 (o_0,o_1,o_2,o_3,o_4,o_5)=(0,-b,-a,-b,0,c).
\end{equation}
これは$Q=0$の最大等方3平面であり、本族はその全方向を代数的に同時に積分している。
\end{theorem}
\begin{proof}
位置0の一つの関係空間だけを取り、
$\phi:S\to\Gr(6,k^4\otimes k^4)$と
$G=\PGL_4\times\PGL_4$の作用を考える。
$p_0$での関係座標における群作用の微分は$[N_0\ N_1]$であり、階数は30である。
式\eqref{eq:seven-family}を七つのパラメータで微分した$60\times7$行列を$J$とすると
\[
 \rk[N_0\ N_1]=30,\qquad \rk[N_0\ N_1\ J]=37.
\]
同梱コードはこれらを有理数上で厳密に検証する。
従って滑らかな37次元の源$G\times S$からGrassmannianへの軌道飽和写像は
$p_0$で階数37を持ち、その像の閉包は37次元である。
一般の$G$軌道は30次元なので、一般軌道を分離する有理商の次元は7である。
別の言い方では、一般のパラメータに対し、その関係空間と$G$-同値なパラメータは有限個である。
裸のindexed同型は必ず位置0の関係空間の$G$-同値を誘導するので、
同型による同一視はこれ以上の正次元ファイバーを生まない。
全域モジュライstackの存在や全周期構造の分類は、この次元評価に不要である。

偶方向について$\rk[D_1(1)\ J_{u,v,r,s}]=19=15+4$、
交代方向について$\rk[D_1(-1)\ J_{a,b,c}]=18=15+3$となる。
商基底へ移すと式\eqref{eq:seven-odd-plane}を得る。
代入すれば$Q=0$であり、階数6の非退化二次形式の等方空間として次元3は最大である。
積分可能性は単なる接空間の主張ではなく、前定理の全$S$上の族によって与えられる。
\end{proof}

\subsection{隣接する点対応の射影配置による非1周期性}
関係を双線型化し、始点の座標を$[w:y:z:x]\in\mathbb P^3$とする。
次の$6\times4$行列が標点の座標に掛かる。
\begin{equation}\label{eq:seven-point-matrix}
 M_\varepsilon(w,y,z,x)=
 \begin{pmatrix}
 y&w&0&0\\z&0&w&0\\-w&z&y&0\\
 x+\varepsilon aw/2&-uy&-vz-\varepsilon ay&w\\
 -w&x-ry-\varepsilon aw&-z-\varepsilon bw&y\\
 0&-sy-\varepsilon cw&x&z
 \end{pmatrix}.
\end{equation}
$E_\varepsilon\subset\mathbb P^3$をその$4\times4$小行列式の零点スキームとする。
これは位置の偶奇に対応する点対応の左射影である。
1周期同型が存在すれば$E_+$と$E_-$は射影線型同型でなければならない。

\begin{theorem}[一般のファイバーは最小周期2]\label{thm:seven-period-two}
$S$の非空なZariski開集合上で、$\mathcal C$の最小周期は2である。
従って7次元の同型類の像の一般点は、どのgraded代数にも由来しない。
これは特定のgraded表示からの変形枝だけを排除する主張ではない。
\end{theorem}
\begin{proof}
$e_X=[0:0:0:1]$以外の点を$e_X$から$\mathbb P^2_{w,y,z}$へ射影する。
行列\eqref{eq:seven-point-matrix}の最初の三行の行列式は$-w(w^2+2yz)$である。
従って射影像は直線$L=(w=0)$と非特異二次曲線$D=(w^2+2yz=0)$の上にある。

$L$上で$[w:y:z]=[0:1:\ell]$とすると、標点の前三座標は$[0:1:-\ell]$であり、
必要十分条件は
\begin{equation}\label{eq:seven-line}
 v\ell^2+\varepsilon a\ell-u=0.
\end{equation}
$uv(a^2+4uv)\ne0$なら相異なる二点を得る。
残る座標は$x+x'=r-\ell^2,\ x'-x=s/\ell$から一意に決まる。

$D\setminus L$上では$[w:y:z]=[1:t:-1/(2t)]$で$t\ne0$と書け、
標点は$[1:-t:1/(2t)]$へ射影される。
最後の三関係から$x,x'$を消去すると
\begin{equation}\label{eq:seven-sextic}
 F_\varepsilon(t)=8st^6+8\varepsilon ct^5+4rt^4+4\varepsilon at^3
                   -4t^2-2\varepsilon bt+1=0
\end{equation}
を得る。$s\ne0$かつ$F_+$が重根を持たなければ、二次曲線上に相異なる六点がある。
それぞれの点で$x,x'$は一意に決まる。
この八点は、対応する一変数方程式の単根性から$E_\varepsilon$の被約点である。
一方$e_X$では最後の三行を用いて標点の三座標を局所的に消去でき、
残る方程式は二次以上である。従ってその接空間は三次元である。
上の条件では支持が有限なので$e_X$は非被約点であり、唯一の非被約支持点となる。

さらに、射影された八点のうち六点を通る二次曲線が$D$だけであるという条件を、
両符号について課す。これは単根の有限\'etale被覆上で有限個の行列式の非零として書ける開条件である。
この条件下で$E_+\simeq E_-$という射影線型同型は$e_X$を固定し、
商平面の射影変換を誘導する。その変換は$D$を保ち、残る二点を結ぶ$L$も保つ。
$D$上の座標$t$では、従って$t\mapsto\rho t$または$t\mapsto\rho/t$である。

$F_-(\rho t)$の根が$F_+(t)$の根に一致する場合、定数項は両方1なので両多項式は等しい。
$t^2$係数から$\rho^2=1$、$a\ne0$なら$t^3$係数から$\rho=-1$である。
反転型の場合は$t^6F_-(\rho/t)=F_+(t)/(8s)$となる。
$t^4,t^2$係数を比較して
\[
 -4\rho^2=\frac{r}{2s},\qquad 4r\rho^4=-\frac1{2s}
\]
を得るので、$r^3+8s=0$が必要である。これを非零と仮定すれば反転型はない。
残る$t\mapsto-t$は平面上で$[w:y:z]\mapsto[w:-y:-z]$であり、$L$上では恒等である。
しかし$a\ne0$なら式\eqref{eq:seven-line}の正負の二点集合は異なる。
従ってこの変換も両配置を対応させず、$E_+$と$E_-$は射影線型非同型である。

以上の開条件が非空であることを、次の有理数点で確認する。
\begin{equation}\label{eq:seven-rational-point}
 (u,v,r,s,a,b,c)=
 (13/3,-13/6,-1/4,1/2,13/2,13/5,-13/5).
\end{equation}
$F_+$の根は$1/5,1/2,1,5,-1,-1/2$、直線上の正符号の根は$1,2$、負符号では$-1,-2$である。
$r^3+8s=255/64$、$a^2+4uv=169/36$となる。
射影平面上の八点から六点を選ぶ28通りについて、二次単項式の評価行列を有理数上で計算すると、
行列式が零になるのは二次曲線上の六点を取った場合だけである。
その場合の階数は5であり、二次曲線は一意である。両符号とも同じ検査に成功する。
具体的な点、行列、根の検査は同梱コードに実装し、証明書に記録する。
従って全条件を満たす非空開集合が存在し、その上では1周期性が排除される。
全ファイバーが2周期的なので、そこでの最小周期は2である。
\end{proof}

\begin{remark}[何を決定し、何をまだ決定していないか]
本節は、形式級数の存在ではなく全$\mathbb A^7$上の多項式関係、全閉ファイバーの正則性、
一般点の非1周期性、同型類の像の次元7を与える。
ただしこの像が全モジュライの既約成分であるとは主張しない。
Clifford点の全九つの偶方向を含む族でもなく、残る交代方向の代数的積分、
全成分の次元、全1周期構造の分類は未解決である。
特に本節の7次元を、前節の14次元候補と取り違えてはならない。
\end{remark}



% !TeX root = ../papers/nc-p3.tex
\section{楕円Clifford部分代数を動かす九変数族}\label{sec:elliptic-ore}
七変数族の三生成部分代数を二方向へ動かす。新しいパラメータを$t,h$とし、
\[
 S_9=\Spec k[u,v,r,s,a,b,c,t,h,(1+th)^{-1}]
\]
を基底とする。以下の$A,B,C$は二次式の略記であり、代数の名前ではない。
\begin{equation}\label{eq:nine-quadratics}
 A=YZ-\tfrac12W^2,\qquad B=WY-\tfrac t2Z^2,\qquad C=WZ-\tfrac h2Y^2.
\end{equation}
位置$i$で$\varepsilon=(-1)^i$とし、隣接する生成空間の積に
\begin{equation}\label{eq:nine-family}
\begin{aligned}
 YW+WY&=tZ^2,& ZW+WZ&=hY^2,& ZY+YZ&=W^2,\\
 XW+WX&=uY^2+vZ^2+\varepsilon(aA-bhB-ctC),\\
 XY+YX&=W^2+Z^2+rY^2+\varepsilon(aB+bC),\\
 XZ+ZX&=sY^2+\varepsilon cB
\end{aligned}
\end{equation}
を課す。これを$\mathcal C^{(9)}$と書く。$t=h=0$への制限は
式\eqref{eq:seven-family}と\emph{表示そのものが一致}する。

\begin{theorem}[九変数の代数的拡張]\label{thm:algebraic-nine}
$\mathcal C^{(9)}$は$S_9$上成分ごとに自由で、各次数$n$の基底は
$W^{i_0}Y^{i_1}Z^{i_2}X^{i_3}$（$i_0+i_1+i_2+i_3=n$）である。
全幾何学的ファイバーは標準14641型quadratic Koszul AS正則$\ZZ$-代数である。
同型類の像は9次元であり、稠密開集合上の最小周期は2である。
従って、どのgraded代数にも由来しない対象の9次元の代数的族を得る。
この像の全モジュライ内での既約成分性は主張しない。
\end{theorem}
証明を正則性、次元、非1周期性に分ける。
後二者を接空間の次元だけから結論することはしない。

\subsection{Clifford基底と全ファイバーのAS正則性}
\begin{lemma}\label{lem:ternary-clifford}
$R=k[t,h,(1+th)^{-1}]$上の
\[
 B_{t,h}=R\langle W,Y,Z\rangle/(YW+WY-tZ^2,ZW+WZ-hY^2,ZY+YZ-W^2)
\]
では三つの平方が中心的である。
$B_{t,h}$は$R[W^2,Y^2,Z^2]$上、
$W^{e_0}Y^{e_1}Z^{e_2}$（$e_j\in\{0,1\}$）を基底とする自由加群である。
全ファイバーは三次元Koszul AS正則代数である。
\end{lemma}
\begin{proof}
例えば$[W,Y^2]=-t[Y,Z^2]$、$[Y,Z^2]=-[Z,W^2]$、
$[Z,W^2]=-h[W,Y^2]$なので$(1+th)[W,Y^2]=0$である。
巡回して全ての平方の中心性を得る。
独立な中心変数$U,V,T$を導入し、$W^2=U,Y^2=V,Z^2=T$と
$YW+WY=tT,ZW+WZ=hV,ZY+YZ=U$を課した通常のClifford代数を考える。
平方を消し、逆順対を交換すると長さが減るか反転数が減り、
Clifford関係の全重なりは解消する。従って平方自由単項式が基底である。
平方の中心性により、$U,V,T$をそれぞれ$W^2,Y^2,Z^2$に送る逆写像があり、
このClifford代数と$B_{t,h}$は同型である。
従って$R$上でも全順序付き単項式を基底に持つ。

各ファイバーの二次双対は可換な
\begin{equation}\label{eq:ternary-dual}
 E=k[\xi,\eta,\zeta]/
 (\xi^2+\eta\zeta,\ \eta^2+h\xi\zeta,\ \zeta^2+t\xi\eta)
\end{equation}
である。三二次式の共通零点が非零なら、ある座標が零の場合は三座標全てが零となって矛盾し、
三座標が非零の場合は三式を掛けて$1+th=0$を得る。
従って三二次式は正則列で、$E$はHilbert級数$(1+q)^3$のFrobenius代数である。

この二次完全交叉のKoszul性は、次のTate複体でも確認できる。
$S=k[\xi,\eta,\zeta]$、$f_j=\sum_i x_i\ell_{ji}$（$\ell_{ji}$は一次式）と書き、
$E\otimes\Lambda(e_1,e_2,e_3)\otimes\Gamma(z_1,z_2,z_3)$に
$d(e_i)=x_i$、$d(z_j)=\sum_i\ell_{ji}e_i$を入れる。
$e_i,z_j$の（ホモロジー次数、内部次数）は$(1,1),(2,2)$である。
$E\otimes_S K_S$のホモロジーは正則列のKoszul解消との比較により
三つの$f_j$の類が生成する外積代数であり、$z_j$を加えた複体はそれらを殺して$k$を解消する。
全微分の係数は一次なので解消は線型である。
従って$E$とその二次双対$B_{t,h}$はKoszulであり、$E$の完全対合から
$B_{t,h}$の左右AS条件と長さ3の線型解消が従う。
これはTateの標準的な構成であり、新しいKoszul性定理を主張するものではない。
\end{proof}

\begin{lemma}[導分の三つの独立な係数]\label{lem:elliptic-derivation}
式\eqref{eq:nine-quadratics}を$B_{t,h}$の元とすると、
\begin{equation}\label{eq:elliptic-derivation}
 d(W)=aA-bhB-ctC,\qquad d(Y)=aB+bC,\qquad d(Z)=cB
\end{equation}
は通常の次数1の導分を定める。
また各生成元を任意の中心二次式に送る写像は、
$\sigma_B(W,Y,Z)=(-W,-Y,-Z)$に関する次数1の$\sigma_B$-導分を定める。
\end{lemma}
\begin{proof}
三つの定義関係にLeibniz則を適用し、Clifford基底へ還元すれば全て零になる。
この計算は$k[t,h,a,b,c,U,V,T]$上の多項式恒等式であり、特殊化を要しない。
さらに
\[
 \Omega=WYZ-\tfrac t2Z^3+\tfrac h2Y^3-\tfrac12W^3
\]
と置くと$d(W^2)=2a\Omega,d(Y^2)=-2b\Omega,d(Z^2)=2c\Omega$である。
これらの恒等式と三関係の検査を
\path{computations/verify_elliptic_family.py}で独立に再現する。
中心二次式の場合は、例えば$\delta_p(YW+WY-tZ^2)$が
中心元との交換子の和となるため零である。他の二関係も同様である。
\end{proof}

\begin{proof}[定理\ref{thm:algebraic-nine}の自由性・AS正則性]
$B_{t,h}$に$g^2=1,\ gj=-jg$（$j=W,Y,Z$）を追加した環で
\[
 \sigma(g)=-g,\quad \sigma(j)=-j,\qquad
 \delta(g)=0,\quad \delta(j)=p_j+g d(j)
\]
と置く。$p_W=uY^2+vZ^2,p_Y=W^2+Z^2+rY^2,p_Z=sY^2$である。
補題\ref{lem:elliptic-derivation}により$\delta$は$\sigma$-導分である。
実際、$B_{t,h}$の二次関係には中心二次式の$\sigma_B$-導分と$g d$が作用し、
$gj+jg$の像も$-gp_j-d(j)+p_jg+d(j)=0$となる。
Ore拡大$Xb=\sigma(b)X+\delta(b)$は左係数環上自由である。
$Xg=-gX$なので二冪等元$(1\pm g)/2$に沿って展開すると
式\eqref{eq:nine-family}を得る。補題\ref{lem:ternary-clifford}から所望の自由基底を得る。

補助重みを$W,Y,Z$に1、$X$に2とする。
随伴次数$\ZZ$-代数は$B_{t,h}[X;\sigma_B]$に付随する1周期的代数であり、
$B_{t,h}$のKoszul解消と一変数多項式環の解消の捩れたテンソル積から、
Koszul性と標準14641型の左右AS双対性を持つ。
各内部次数で有限な重みフィルトレーションをbar複体に入れると、
随伴次数のTorが対角上にしかないため元の代数のTorも対角上に集中する。
従って元の代数はKoszulである。
二次双対には逆フィルトレーションを入れる。
随伴次数の双対からの全射は、Hilbert級数の逆数公式で両辺の次数別次元が
$(1,4,6,4,1)$と一致するため同型である。
随伴次数のFrobenius完全対合は元の双対へ持ち上がる。
左右Koszul複体の双対性から左右AS条件を得て、Koszul双加群解消から大域次元4を得る。
全ての議論は$1+th\ne0$の各ファイバーで成立する。
\end{proof}

\begin{remark}
$t,h$が共に非零の場合、上の自由基底の存在を
六つの関係式が固定単項式順序に関する二次Gr\"obner基底であるという主張に読み替えない。
ここでは中心平方上のClifford基底とOre基底を使った。
$1+th=0$を無断で含めることもできない。
例えば$(t,h)=(1,-1)$では三生成代数の次数3の関係行列の階数が17から15へ下がる。
\end{remark}

\subsection{同型類の九次元性}
\begin{proof}[定理\ref{thm:algebraic-nine}の次元主張]
基点$(u,v,r,s,a,b,c,t,h)=(0,1,0,0,0,0,0,0,0)$での代数は$C_0$である。
一つの関係空間に対する$\PGL_4\times\PGL_4$の微分を$[N_0\ N_1]$、
九パラメータの微分を$J_9$とすると、厳密な有理数計算で
\[
 \rk[N_0\ N_1]=30,\qquad \rk[N_0\ N_1\ J_9]=39
\]
を得る。また偶六方向では階数$21=15+6$、交代三方向では$18=15+3$となる。
従って39次元の滑らかな源$(\PGL_4\times\PGL_4)\times S_9$からの
軌道飽和写像の像は39次元で、一般軌道は30次元である。
定理\ref{thm:seven-image}と同じ有理商の議論から、関係空間の同型類は9次元である。
indexed同型は必ずこの関係空間の同型を誘導するので、
九パラメータの一般ファイバーで正次元の同一視は生じない。
これは実際の九変数族の像の次元であり、全変形接空間の次元からの推測ではない。
\end{proof}

\subsection{滑らかな楕円曲線上の12点による非1周期性}
以下の議論は、七変数族の非1周期点があることから九変数族の一般点を推測するものではない。
九変数族の楕円領域で独立に非1周期性を証明する。
三生成部分の双線型関係の行列は
\[
 K(p)=\begin{pmatrix}y&w&-tz\\z&-hy&w\\-w&z&y\end{pmatrix},\qquad p=[w:y:z].
\]
その行列式は
\begin{equation}\label{eq:elliptic-cubic}
 -F_{t,h}(w,y,z),\qquad
 F_{t,h}=w^3+hy^3+tz^3+(2-th)wyz
\end{equation}
である。$th(1+th)(th-8)\ne0$なら$E=(F_{t,h}=0)$は滑らかな平面三次曲線である。
実際、Hesse型三次式の特異点条件を直接計算すると
\[
 (2-th)^3+27th=-(1+th)^2(th-8)
\]
となる。$K(p)$は$E$上階数2であり、その核を$q=\tau(p)$と書く。
双線型関係が対称なので$\tau$は正則な対合である。
固定点は$2wy=tz^2,2wz=hy^2,2yz=w^2$を満たすが、
滑らかな領域では非零解がない。従って$\tau$は非零2等分点による平行移動である。
$L=\mathcal O_E(1)$は次数3なので$\tau^*L\not\simeq L$であり、
$\tau$自体は平面の射影変換から来ない。

まず$p_W,p_Y,p_Z$を三つの平方の\emph{任意の}線型結合に広げ、
その九係数を$P=(P_{ij})$とする。導分は式\eqref{eq:elliptic-derivation}のままとする。
この拡大族も前と同じ証明でAS正則である。
$q=\tau(p)$に対し
\[
 U=(p_0q_0,p_1q_1,p_2q_2)^{\mathsf T},\qquad \omega=p\times q
\]
と置く。第四座標$x,x'$が存在するための条件は
\begin{equation}\label{eq:elliptic-divisor}
 S_\varepsilon(p)=\omega^{\mathsf T}\bigl(PU+\varepsilon d(p,q)\bigr)=0.
\end{equation}
$p,q$は常に一次独立なので、条件が成立すれば$x,x'$は一意に決まる。
ここで$S_\varepsilon$は$M=(L\otimes\tau^*L)^2$、次数12の直線束の切断である。
三二次式$A,B,C$は双線型化後、それぞれ$\omega_0/2,\omega_2/2,-\omega_1/2$となるので
\begin{equation}\label{eq:even-divisor-section}
 \omega^{\mathsf T}d(p,q)=\tfrac12\bigl(
 a(\omega_0^2+\omega_1\omega_2)
 -b(h\omega_0\omega_2+\omega_1^2)
 +c(t\omega_0\omega_1+\omega_2^2)\bigr).
\end{equation}
$\tau$によって$U$は不変、$\omega$は符号を変える。
従って$D_\varepsilon=(S_\varepsilon=0)$は$D_-=\tau(D_+)$を満たす。

\begin{lemma}[使える切断は反不変六次元と不変三次元]\label{lem:elliptic-sections}
$P,a,b,c$を動かした$S_+$の空間$\mathcal W\subset H^0(E,M)$は9次元である。
その反不変部分は$H^0(E,M)^-$全体の6次元、不変部分は3次元である。
$\mathcal W$は基点自由で、一般の切断の零因子は相異なる12点である。
\end{lemma}
\begin{proof}
$\pi:E\to E'=E/\langle\tau\rangle$とし、
$\pi_*\mathcal O_E=\mathcal O_{E'}\oplus\eta$と書く。$\eta$は非自明な2-torsionである。
$L\otimes\tau^*L$の因子交換による線型化を降下させた次数3の直線束を$N$とすると、
$U_i$は$H^0(N)$、$\omega_i$は$H^0(N\otimes\eta)$の基底になる。
この独立性を確認するには、$L$と$\tau^*L$が非同型な次数3の直線束であることを使う。
基点自由pencilの完全列と$H^0(\tau^*L\otimes L^{-1})=H^1(\tau^*L\otimes L^{-1})=0$より、
$H^0(L)\otimes H^0(\tau^*L)\to H^0(L\otimes\tau^*L)$は全射で、核は3次元である。
核は既存の三つの対称な双線型関係に一致する。その関係空間に非零の対角式も
反対称式もないので、$U_i$と$\omega_i$はそれぞれ独立である。

次数3の$N$に基点自由pencilを取り、その完全列を$N\otimes\eta$で捩ると
$H^0(\eta)=H^1(\eta)=0$より
\[
 H^0(N)\otimes H^0(N\otimes\eta)\longrightarrow H^0(N^2\otimes\eta)
\]
は全射である。従って$\omega_iU_j$は反不変部分の6次元全体を張る。
一方$\omega_i$による$E'$の写像は次数3の埋込みである。
その像は平面三次曲線で、非零二次式は恒等的に零にならない。
式\eqref{eq:even-divisor-section}の三二次式は平方係数を見れば独立なので、不変部分は3次元となる。
反不変部分だけで、$E'$上の次数6の完全線型系の引戻しになっている。
従って基点自由で、$\pi$が\'etaleであることから相異なる12点を切る切断が存在する。
単根性は開条件なので一般の$\mathcal W$の切断も同じ性質を持つ。
\end{proof}

\begin{lemma}[一般の12点配置は対合後の配置と射影同型でない]\label{lem:elliptic-nonperiod}
$\mathcal W$の一般の切断について、$D_+$を$\tau(D_+)$へ送る平面の射影変換は存在しない。
\end{lemma}
\begin{proof}
12点は相異なるので、B\'ezoutの定理によりそれらを含む三次曲線は$E$だけである。
従って候補は有限群$\Aut(E,L)$に限られる。
任意の$\gamma$に対して$f=\tau^{-1}\gamma$と置く。
$\tau$は$L$を保たないため$f\ne\mathrm{id}$である。
また$M\simeq L^4$であり、$\gamma,\tau$はともに$M$を保つ。
$\mathcal W$全体で$f(D)=D$となるためには、$f^*$が$\mathcal W$上スカラーでなければならない。

反不変六次元の完全線型系が与える写像は、$\pi$と$E'$の次数6の埋込みの合成である。
従って$f^*$がこの空間上スカラーなら$\pi f=\pi$であり、$f=\mathrm{id}$または$\tau$である。
前者は除外済みで、後者は$\gamma=\mathrm{id}$のときだけである。
しかし$\tau$は反不変部分に$-1$、非零な不変三次元部分に$+1$として作用するため、
$\mathcal W$全体ではスカラーでない。
従って各$\gamma$について、$f(D)=D$を満たす切断は固有空間の有限和集合という真の閉部分集合に入る。
有限個の$\gamma$を全て除けば所望の非空開集合を得る。
\end{proof}

\begin{proof}[定理\ref{thm:algebraic-nine}の一般の最小周期2]
拡大族の一般点で、$e_X=[0:0:0:1]$以外の左点スキームは$D_\varepsilon$の12点の
一意な持上げであり、全て被約である。
$e_X$の近傍では最後の三関係を使って標点座標を消去でき、残る三式の先頭二次項は
$YW+WY-tZ^2,ZW+WZ-hY^2,ZY+YZ-W^2$を同一点で評価したものである。
それらが射影共通零点を持てば$\tau$の固定点を与えるため、三式は正則列である。
従って$e_X$は長さ8、接空間次元3の局所完全交叉であり、唯一の非被約支持点となる。
左点スキームの射影同型は$e_X$を固定し、商平面で$D_+$と$D_-$を対応させる。
補題\ref{lem:elliptic-nonperiod}はこれを一般点で排除する。
よって拡大族は一般に1周期的ではない。

拡大族から式\eqref{eq:nine-family}の正規形へ戻る。
$X$に$\ell_WW+\ell_YY+\ell_ZZ$を足すと中心係数は
\[
\begin{array}{c|ccc}
 &W^2&Y^2&Z^2\\\hline
 p_W&2\ell_W&h\ell_Z&t\ell_Y\\
 p_Y&\ell_Z&2\ell_Y&t\ell_W\\
 p_Z&\ell_Y&h\ell_W&2\ell_Z
\end{array}
\]
だけ変化し、交代項は変わらない。
従って$p_W$の$W^2$係数と$p_Z$の$W^2,Z^2$係数を一意に零にできる。
残る$p_Y$の$W^2,Z^2$係数を$A_0,B_0$とし、$A_0B_0\ne0$の開集合を取る。
$W'=\lambda W,Y'=\mu Y,Z'=\nu Z,X'=\rho X$、$\mu\nu=\lambda^2$とすると
それらは$\rho A_0/\nu$、$\rho\mu B_0/\nu^2$へ移る。
$\nu=1,\mu=A_0/B_0,\lambda^2=A_0/B_0,\rho=1/A_0$で両方を1にできる。
この変換は三生成関係の形と導分の三次元空間を保つ。
すなわち
\[
 t'=\frac{\mu\lambda}{\nu^2}t,\quad h'=\frac{\nu\lambda}{\mu^2}h,
 \quad a'=\frac\rho\lambda a,\quad b'=\frac{\rho\mu}{\lambda\nu}b,
 \quad c'=\frac{\rho\nu}{\lambda\mu}c
\]
であり、$t'h'=th$である。
従って拡大族の一般点は九変数の正規形と同型であり、非1周期性も保たれる。

開性については、滑らかな楕円領域を必要に応じて縮め、有限\'etale被覆で
$\Aut(E,L)$を一定化する。上の単根条件と各有限自己同型に関する非固有ベクトル条件は開であり、
被覆から降下する。これを正規形へ制限すれば$S_9$の非空開集合を得る。
全ファイバーが厳密2周期的であることと合わせて、そこでの最小周期は2である。
\end{proof}

\begin{corollary}[1周期的対象への代数的特殊化]\label{cor:nine-to-period-one}
九変数族の任意の対象は、全ファイバーが標準14641型quadratic Koszul AS正則である
$\mathbb A^1$上の成分ごとに自由な族によって、1周期的対象とつながる。
同じことは七変数族についても成立する。
\end{corollary}
\begin{proof}
$p=(u,v,r,s,a,b,c,t,h)\in S_9(k)$を固定し、
\[
 \mathcal D_\lambda
 =\mathcal C^{(9)}(u,v,r,s,\lambda a,\lambda b,\lambda c,t,h)
 \qquad(\lambda\in\mathbb A^1)
\]
と置く。$t,h$を変えないため$1+th\ne0$は保たれ、
定理\ref{thm:algebraic-nine}の基底変換により自由性と全ファイバーのAS正則性を得る。
$\mathcal D_1=\mathcal C^{(9)}_p$であり、
$\lambda=0$では$(-1)^i$を含む全ての項が消える。
従って関係式は位置$i$に依存せず、$\mathcal D_0$は厳密1周期である。
$t=h=0$に制限すれば七変数族の場合を得る。
\end{proof}

\begin{remark}[非1周期性と変形連結性]
一般点の最小周期が2であることは、AS正則な連結平坦族の中で1周期的対象へ到達できないことを意味しない。
上の系は本族では到達できることを示す。
また、1周期的な特殊ファイバーを持つことと、1周期的な一般ファイバーを持つ族の特殊化になることは別問題である。
本族の最大性が証明されても、1周期的locusと交わらない変形連結成分の存在は従わない。
\end{remark}

\begin{remark}[確定範囲]
九変数族は同型類の次元9を持ち、七変数族を特殊化として含む。
三生成部分の点曲線は、$t=h=0$の直線と二次曲線の和から、一般の$t,h$で滑らかな楕円曲線へ変わる。
従って追加の二変数は、同じ七変数族の冗長な表示ではない。
ただし全モジュライの既約成分性、Clifford点の14次元局所模型、残る三つの交代方向の同時積分、
全1周期構造および全族内外同型の分類はなお未解決である。
本節の代数的構成と点配置の証明の先行研究上の新規性は未確認である。
\end{remark}

% !TeX root = ../papers/nc-p3.tex
\section{Ore型構成の内部完全性と楕円同型不変量}\label{sec:ore-complete}
\subsection{次数一のOre導分の完全な分類}\label{subsec:ore-derivations}
固定された$B_{t,h}$の上でのOre型構成がどの程度一般的かを、
具体的な導分空間の分類によって明らかにする。ここで$R=k[t,h,(1+th)^{-1}]$、
$\sigma_B|_{(B_{t,h})_1}=-1$とする。
$\operatorname{Der}_1(B)$は内部次数1の通常の導分、
$\operatorname{Der}_{\sigma_B,1}(B)$は内部次数1の$\sigma_B$-導分の$R$-加群とする。

\begin{theorem}[固定Clifford部分のOre導分の全分類]\label{thm:ore-derivations}
$R$上、次の自由加群としての等式が成立する。
\begin{align}
 \operatorname{Der}_1(B_{t,h})
 &=R\operatorname{ad}_W\oplus R\operatorname{ad}_Y\oplus R\operatorname{ad}_Z
   \oplus Rd_a\oplus Rd_b\oplus Rd_c,\label{eq:ordinary-derivations}\\
 \operatorname{Der}_{\sigma_B,1}(B_{t,h})
 &=\operatorname{Hom}_R\bigl((B_{t,h})_1,
       RW^2\oplus RY^2\oplus RZ^2\bigr)\simeq R^9.
       \label{eq:sign-derivations}
\end{align}
ここで$d_a,d_b,d_c$は式\eqref{eq:elliptic-derivation}の$a,b,c$の係数である。
従って、$Xj+jX=p_j+(-1)^i d_j$で表される二周期Ore型拡大で、
$p$を次数1の$\sigma_B$-導分、$d$を通常の次数1の導分とするものは、
固定した$B_{t,h}$上では$9+6=15$個の線型係数を持つ。
$X_i\mapsto X_i+\ell+(-1)^i m$（$\ell,m\in(B_{t,h})_1$）による
六つの独立な冗長性を除くと、残る係数は$6+3=9$個である。
この完全性は\emph{指定したClifford部分代数とOre形式を保持する構成}に限る。
一般の四次元AS正則$\mathbb Z$-代数全体の分類を意味しない。
\end{theorem}
\begin{proof}
$(B_{t,h})_2$の基底を$(W^2,WY,WZ,Y^2,YZ,Z^2)$とする。
補題\ref{lem:ternary-clifford}の中心平方上の表示によって、
三次成分には
\[
 W(W^2,Y^2,Z^2),\quad Y(W^2,Y^2,Z^2),\quad
 Z(W^2,Y^2,Z^2),\quad WYZ
\]
の10個の基底がある。生成元3個の各像を六つの基底で書くと18未知数となる。
$YW+WY-tZ^2,ZW+WZ-hY^2,ZY+YZ-W^2$への
Leibniz則または$\sigma_B$-Leibniz則の適用は、
これらの三次基底上に30行18列の行列$M_{\rm ord},M_\sigma$を与える。

通常導分の行列には、$t,h$に依存しない行列式$-2048$を持つ
12次正方小行列がある。よって$\operatorname{rank}M_{\rm ord}\ge12$。
一方、三つの内導分と式\eqref{eq:elliptic-derivation}の三つの導分は核に属し、
それらの18行6列の係数行列には、行列式$4$の6次正方小行列がある。
従って全ファイバーで核の次元は少なくとも6であり、上の12次小行列とあわせて
厳密に6となる。二つの定数小行列式が単元であるため、$R$上でも
六つの導分が核の自由基底となる。

$\sigma_B$-導分については、三つの生成元の像を任意の中心二次式とすれば
三関係への適用が零となるので、$R^9$が核に含まれる。
一方$M_\sigma$には、二つの9次正方小行列があり、その行列式は
\[
 4(th-2)(th+4)^2,\qquad 3ht^3(th-8)^2
\]
である。これらは同時に零にならないため、全ファイバーで階数は9以上である。
中心二次式を取る九つの標準導分は直和因子であり、
二つの小行列が覆う開集合上で$M_\sigma$の核と一致する。
ゆえに$R$上でも式\eqref{eq:sign-derivations}を得る。

最後に$\ell,m\in(B_{t,h})_1$として生成元を変えると、
関係の右辺への変化は
\[
 p_j\longmapsto p_j+\ell j+j\ell,\qquad
 d_j\longmapsto d_j+mj-jm
\]
である。前者は三次元の中心二次式値の$\sigma_B$-内導分、
後者は三次元の通常の内導分であり、互いに独立である。
実際$\ell=\alpha W+\beta Y+\gamma Z$とすれば、
$p_W$の$W^2$係数、$p_Z$の$W^2,Z^2$係数への寄与は
$2\alpha,\beta,2\gamma$なので、前者は全ファイバーで自由に作用する。
後者の独立性は上の六導分の自由基底性から従う。
全行列、小行列の行・列番号、厳密な多項式恒等式は
\path{computations/verify_ore_derivations.py}で再現できる。
\end{proof}

\begin{corollary}[九変数正規形のOre型での十分性]\label{cor:ore-slice}
$B_{t,h}$と符号自己同型$\sigma_B$を固定した上の二周期Ore型構成を考える。
$X_i\mapsto X_i+\ell+(-1)^im$という変更を許せば、
$d$は一意に$d_a,d_b,d_c$の線型結合へ移せる。
また$p$は$p_W$の$W^2$係数と$p_Z$の$W^2,Z^2$係数を零にした六係数の形へ移せる。

さらに$B_{t,h}$の三関係の形を保つ対角変換
$W\mapsto\lambda W,Y\mapsto\mu Y,Z\mapsto\nu Z,X\mapsto\rho X$
（$\lambda^2=\mu\nu$）を許し、残る$p_Y$の$W^2,Z^2$係数が非零な開集合に限ると、
両係数を1にできる。従ってこの\emph{印付きOre型構成}の稠密開集合は、
有限個の対角正規化の選択を除いて式\eqref{eq:nine-family}の九変数族で表される。
従って同じClifford部分のOre形式を固定したまま、さらに自由な
次数一の導分係数を付け加えても、一般の同型類の次元は増えない。
\end{corollary}
\begin{proof}
定理\ref{thm:ore-derivations}の六つの基底変更によって所要の切片を取る。
前者三個は通常の内導分を除き、後者三個は中心二次式の三係数を除く。
$\lambda^2=\mu\nu$を満たす対角変換は、全体スカラーを除いて二次元のトーラスをなす。
$p_Y$の$W^2,Z^2$係数を$A_0,B_0$とすると、変換後は
$\rho A_0/\nu,\rho\mu B_0/\nu^2$である。
従って$A_0B_0\ne0$なら$\rho/\nu=A_0^{-1}$、
$\mu/\nu=A_0/B_0$と取って両者を1にできる。
$\lambda^2=\mu\nu$に残る符号の有限選択だけは同型類の次元に影響しない。
残る六つの中心係数は$(u,v,r,s)$および規格化された二係数、
通常導分の外部三係数は$(a,b,c)$へ整理され、
式\eqref{eq:nine-family}と一致する。$t,h$も動かした全パラメータ数は
$2+15-3-3-2=9$である。
これは印付きを忘れた同型の完全分類ではない。
九変数族の\emph{裸の}indexed同型類の像の次元9は、前節の独立な階数証明による。
\end{proof}

\subsection{内在的な楕円$j$-不変量による同型の必要条件}
\begin{proposition}[九変数族の同型に対する明示的障害]\label{prop:elliptic-j-obstruction}
九変数族のうち、$E_{t,h}$が滑らかで、唯一の非被約左点が$e_X$、
その点からの射影で残りの左点が$E_{t,h}$上の相異なる12点となる開集合を考える。
この開集合の二つのファイバーが\emph{添字保持}$\ZZ$-代数として同型なら、
$x=th$、$x'=t'h'$に対して
\begin{equation}\label{eq:elliptic-j}
 J(x)=J(x'),\qquad
 J(x)=\frac{(x-2)^3(x^3-6x^2+228x-8)^3}
                 {x(x-8)^3(x+1)^6}
\end{equation}
が必要である。$J:\mathbb P^1_x\to\mathbb P^1_j$は次数12の非定数有理写像である。
従って同型類を一つ固定したとき、その同型類に属し得る一般の$x=th$の値は
高々12個である。ただし$j$の一致は同型の\emph{十分条件ではない}。
\end{proposition}
\begin{proof}
添字保持同型は左点スキームの射影同型を誘導する。
その唯一の非被約点$e_X$は内在的であるから射影同型で保たれる。
従ってそこからの射影が作る平面上の12点集合も射影同型で対応する。
それらを通る平面三次曲線は$E_{t,h}$ただ一つである。
実際、滑らかな既約三次曲線と異なる三次曲線は、共通成分がなければ
B\'ezoutの定理により9点を超えて相異なる点で交わらない。
よって二つのファイバーが同型なら$E_{t,h}$と$E_{t',h'}$が同型となり、
それらの$j$-不変量は一致する。

標数零の代数閉体上で$w,y,z$をそれぞれ適切に三乗根で規格化すると、
式\eqref{eq:elliptic-cubic}はHesse標準形
$W^3+Y^3+Z^3-3\kappa WYZ=0$になり、
\[
 \kappa^3=-\frac{(2-th)^3}{27th}.
\]
Hesse三次曲線の古典的公式
$j=27\kappa^3(\kappa^3+8)^3/(\kappa^3-1)^3$
（\cite[\S7.3, (42)]{KLV}）へ代入して式\eqref{eq:elliptic-j}を得る。
分子の次数12、分母の次数10で、両者に共通因子はない。
従って有理写像の次数は12であり、各幾何学的ファイバーに含まれる
$x$の値は高々12個である。
多項式恒等式と次数の検証は\path{computations/verify_ore_derivations.py}にある。
\end{proof}

\section{楕円2-isogenyの内在性と四対一のモジュライ座標}\label{sec:isogeny}
九変数族の楕円曲線$E_{t,h}$に現れる非零2等分点による平行移動を
$\tau:E\to E$と書き、$L=\mathcal O_E(1)$とする。
本節では定理\ref{thm:algebraic-nine}の開集合をさらに必要に応じて縮め、
$E$が滑らかで、各位置の左点スキームの非被約支持点が唯一の$e_X$であり、
$e_X$から射影した残りの支持が$E$上の相異なる12点である場合を考える。
これらは証明済みの非空な開条件である。

\begin{lemma}[Clifford部分と2等分点の内在的復元]\label{lem:intrinsic-elliptic-tau}
上の開集合の二つのファイバーが添字保持$\ZZ$-代数として同型なら、
両者の偏極付き楕円曲線と非零2等分点による平行移動の組
\[
 (E,L,\tau),\qquad (E',L',\tau')
\]
は同型である。具体的には、各整数位置に誘導される平面の射影同型は
楕円曲線上では同一の$\alpha:E\xrightarrow\sim E'$となり、
$\alpha^*L'\simeq L$、$\alpha\tau=\tau'\alpha$を満たす。
従って$E/\langle\tau\rangle$の同型類もindexed $\ZZ$-代数の不変量である。
\end{lemma}
\begin{proof}
左点スキームの唯一の非被約点$e_X$は、次数一空間の射影空間内で内在的に定まる。
その消滅超平面$H_i\subset C_{i,i+1}$は、表示では$W,Y,Z$が張る三次元空間であり、
任意の添字保持同型で保存される。これらの$H_i$から生成される
部分$\ZZ$-代数は、全位置で同じ三生成Clifford代数$B_{t,h}$に由来する。
従って任意のindexed同型は、この部分代数のindexed同型を誘導する。
その次数一写像を双対化すると、各位置で偏極を保つ射影同型
$\alpha_i:(E,L)\xrightarrow\sim(E',L')$が得られる。

三生成部分の双線型関係の点対応は$\tau$のグラフであるから、
$\alpha_{i+1}\tau=\tau'\alpha_i$が成立する。
$\beta_i=\alpha_i\tau\alpha_i^{-1}$は$E'$上の非零2等分点の平行移動である。
従って$\alpha_{i+1}\alpha_i^{-1}=\tau'\beta_i^{-1}$は、
高々2位数の点による平行移動である。
一方この自己同型は$L'$を保つ。
楕円曲線上の次数3の直線束を保つ平行移動の点は3-torsionに限るので、
2-torsionでもあるこの点は零である。
ゆえに$\alpha_{i+1}=\alpha_i$であり、同時に
$\alpha_i\tau=\tau'\alpha_i$を得る。
\end{proof}

\begin{theorem}[商曲線の方程式と強化された同型障害]\label{thm:explicit-two-isogeny}
$x=th$、$x\ne0,-1,8$とする。非零2等分点$\tau$で割った商楕円曲線
$E'=E/\langle\tau\rangle$には次の平面三次表示がある。
\begin{equation}\label{eq:isogenous-cubic}
 th\,A^3+t\,B^3+h\,C^3-(th+4)ABC=0.
\end{equation}
$j(E)$と$j(E')$はそれぞれ
\begin{align}
 J(x)&=\frac{(x-2)^3(x^3-6x^2+228x-8)^3}
                    {x(x-8)^3(x+1)^6},\label{eq:source-j}\\
 J'(x)&=\frac{(x+4)^3(x^3+228x^2+48x+64)^3}
                    {x^2(x-8)^6(x+1)^3}.\label{eq:isogenous-j}
\end{align}
さらに
\begin{equation}\label{eq:isogeny-modular-u}
 U(x)=\frac{4096(x+1)^3}{x(x-8)^3}
\end{equation}
と置けば
\begin{equation}\label{eq:2-isogeny-modular-forms}
 J(x)=\frac{(U(x)+256)^3}{U(x)^2},\qquad
 J'(x)=\frac{(U(x)+16)^3}{U(x)}.
\end{equation}
有理写像$U:\mathbb P^1_x\to\mathbb P^1_U$の次数は4である。
また$k(J(U),J'(U))=k(U)$なので、一般の二つの九変数族のファイバーが
indexed同型なら$U(th)=U(t'h')$が必要である。
従って、同型類を固定すると、一般の$x=th$に可能な値は高々\emph{4個}である。
これは命題\ref{prop:elliptic-j-obstruction}の12個という上界を改善する。
$U$の一致は同型の十分条件ではない。
\end{theorem}
\begin{proof}
式\eqref{eq:elliptic-cubic}の$K(p)$について$p=(w,y,z)$と置き、
第一行と第二行の外積を計算すると
\[
 q=(w^2-thyz,-wy-tz^2,-wz-hy^2)
\]
となる。直接計算で$K(p)q=(0,0,-F_{t,h}(p))^{\mathsf T}$だから、
$E$の稠密開集合で$q$は$\tau(p)$を表す。
このとき$p\times q$の三座標は
\begin{equation}\label{eq:isogeny-cross-product}
 \begin{aligned}
 A&=tz^3-hy^3,\\
 B&=2zw^2+hy^2w-thyz^2,\\
 C&=-2yw^2-tz^2w+thy^2z.
 \end{aligned}
\end{equation}
これらは$q$の共通のスカラーを除けば$\omega=p\times\tau(p)$であり、
$\tau$で符号が変わるため射影的には不変である。
より直接的には、$\omega$の三座標は
$L\otimes\tau^*L$（次数6）の切断であり、同時に零となるのは
$p=\tau(p)$のときに限る。$\tau$は固定点を持たないから基点自由である。
三座標が線型独立なことは式\eqref{eq:isogeny-cross-product}から分かる：
線型結合が$E$上零ならその三次多項式は$F_{t,h}$の定数倍だが、
$w^3$係数を比較すると倍数が零となり、$w^2y,w^2z$および
$y^3,z^3$の係数比較で三つの係数も零となる。
従って像は直線ではない。
この写像は$\tau$で不変だから次数は偶数であり、
像の次数と写像次数の積は6となる。
像が直線でないことから写像次数は2、像の次数は3である。
これは補題\ref{lem:elliptic-sections}の商の次数3線型系の記述とも一致する。

多項式として次の恒等式が成り立つ。
\[
\begin{split}
 &th A^3+t B^3+h C^3-(th+4)ABC\\
 &\quad=F_{t,h}(w,y,z)\,(tz^3-hy^3)
       (h^2ty^3+ht^2z^3-6htwyz+8w^3).
\end{split}
\]
従って像は式\eqref{eq:isogenous-cubic}に含まれる。
この三次曲線は$x\ne0,-1,8$で滑らかであり、先の商の平面三次曲線と一致する。

両曲線はHesse標準形に変形でき、そのパラメータの三乗は
\[
 \kappa^3=-\frac{(2-x)^3}{27x},\qquad
 (\kappa')^3=\frac{(x+4)^3}{27x^2}
\]
である。Hesseの古典的な$j$公式\cite[\S7.3, (42)]{KLV}へ代入すると、
式\eqref{eq:source-j}と\eqref{eq:isogenous-j}が得られる。
式\eqref{eq:2-isogeny-modular-forms}は有理関数の恒等式であり、
\path{computations/verify_isogeny_invariant.py}で全分子を多項式として検査する。
$U(x)$の分子・分母は互いに素で、それぞれ次数3と4なので写像の次数は4である。

次に$J_0(U)=(U+256)^3/U^2$、$J_1(U)=(U+16)^3/U$と置く。
$[k(U):k(J_0)]=3$である。
もし$J_1\in k(J_0)$なら、両者の有理写像の次数が3なので
$J_1$は$J_0$の一次分数変換でなければならない。
しかし$J_0$の極は$U=0$で2位、$U=\infty$で1位であり、
$J_1$の極は同じ二点で逆に1位と2位である。
この極の次数の違いは一次分数変換では実現できない。
ゆえに$J_1\notin k(J_0)$であり、次数3が素数だから
$k(J_0,J_1)=k(U)$となる。
従って一般の$U$は順序付き対$(J_0,J_1)$で一意に決まる。
補題\ref{lem:intrinsic-elliptic-tau}によりindexed同型はこの順序付き対を保つから、
一般点では$U(th)$も保存される。
次数4の有理写像$U(x)$のファイバーには高々4点しかない。
\end{proof}

\begin{corollary}[四対一被覆の分岐と四面体群]\label{cor:isogeny-a4}
有理関数\eqref{eq:isogeny-modular-u}の分岐値は$0,-64,\infty$であり、
分岐指数はそれぞれ$(3,1),(2,2),(3,1)$である。
特に$k(x)/k(U)$のGalois閉包の群は四面体群$A_4$であり、
この次数4の被覆はGalois被覆ではない。
これはパラメータの代数的被覆についての主張であって、
一般ファイバーの自己同型群を$A_4$と同定するものではない。
\end{corollary}
\begin{proof}
微分は
\[
 \frac{dU}{dx}=-\frac{4096(x+1)^2(x^2+20x-8)}{x^2(x-8)^4}.
\]
$U^{-1}(0)=3[-1]+[\infty]$、$U^{-1}(\infty)=3[8]+[0]$であり、
$x=-10\pm6\sqrt3$の各点は分岐指数2で$U=-64$へ写る。
分岐寄与は$2+2+1+1=6=2\cdot4-2$なので、他の分岐はない。
分岐置換は3-cycleとdouble transpositionで、すべて偶置換である。
被覆は連結なのでmonodromy群は$A_4$の推移的部分群であり、
3-cycleを含むから$V_4$ではなく$A_4$そのものである。
\end{proof}

\begin{remark}
$U(x)$が4対1であることを、異なる四つの$x$が
必ず同型な九変数族の代数を与えるという結論に読み替えてはならない。
ここで得たのは楕円曲線と指定2-isogenyから来る\emph{必要な}代数的不変量である。
十二点の因子と第四生成元の持上げのデータはまだ比較する必要がある。
\end{remark}

\section{Clifford--Ore表示の外に出る代数的な二周期変形}\label{sec:nonore}
固定された三生成部分代数のOre拡大とは異なる変形を得るため、
三生成元の関係自体に第四生成元$X$を含む項を加える。
ここでは記号$k_0,v_0,b_0,s_0,a_0$をパラメータとし、
基底体$k$との混同を避ける。$S_5=\mathbb A^5_{k_0,v_0,b_0,s_0,a_0}$とする。
位置$i$で$\varepsilon=(-1)^i$と置き、以下の関係を課す。
\begin{equation}\label{eq:nonore-five}
\begin{aligned}
 YW+WY&=\varepsilon a_0ZX,& ZW+WZ&=0,& ZY+YZ&=k_0W^2,\\
 XW+WX&=v_0Z^2,& XY+YX&=b_0W^2+s_0Z^2,& XZ+ZX&=0.
\end{aligned}
\end{equation}
$k_0=v_0=b_0=s_0=1,a_0=0$でClifford基点$C_0$になる。
$a_0$の変化は$YW+WY$に$X$を含ませるので、
\emph{選んだ}$W,Y,Z$部分代数によるOre構成には属さない。

\begin{theorem}[全五変数で正則な非Ore型チャート]\label{thm:nonore-five}
式\eqref{eq:nonore-five}は$S_5$上、全次数で自由な二周期$\ZZ$-代数を定める。
各次数$n$で$W^iY^jZ^\ell X^m$（$i+j+\ell+m=n$）が基底となり、
全ての幾何学的ファイバーは標準14641型のquadratic Koszul AS正則代数である。
添字保持同型による同型類の像の次元は正確に\emph{2}である。
Clifford基点での一次変形は独立な偶方向1個と交代方向1個を含み、
後者は前節の九変数Ore族の接像には入らない。
一般ファイバーの最小周期が2であるかどうかは本定理では主張しない。
\end{theorem}
\begin{proof}
$W<Y<Z<X$とし、補助重みを
\[
 \operatorname{wt}(W,Y,Z,X)=(3,5,2,2)
\]
と定め、正の重みと辞書式順序で語を比較する。
式\eqref{eq:nonore-five}の六つの規則では、反順序の二次語を
$-1$倍の順序付き語へ換える項は同じ重みで辞書式順序が小さく、
残る全ての項は\emph{真に小さい重み}を持つ。
例えば$YW$の重みは8だが$ZX$は4、$ZY$は7だが$W^2$は6、
$XW$は5だが$Z^2$は4、$XY$は7だが$W^2$は6、$Z^2$は4である。
従って還元は全パラメータ上で停止する。

三次重なりは$ZYW$、$XYW$、$XZW$、$XZY$の四つであり、
始点の偶奇を含めて全部で8個である。
$\mathbb Q[k_0,v_0,b_0,s_0,a_0]$上で、二つの還元順序の差は
\emph{恒等的に零}になる。
再現コード\path{computations/verify_nonore_family.py}は8個すべてを
記号多項式として検査する。Diamond lemmaより全基底上の自由性が従う。
重み付き随伴次数代数は、四つの生成元が相異なるとき反可換する
標準四変数量子多項式代数である。
二次Gr\"obner性からKoszul性が従い、二次双対のHilbert級数は
$(1+t)^4$となる。双対への逆フィルトレーションを使うと、
その随伴次数代数は量子外積代数からの商であり、次数別次元が等しいため同型となる。
量子外積の相補次数積は完全であるから元の双対も局所Frobeniusである。
左右Koszul解消を双対化して標準14641型の左右AS条件と大域次元4を得る。
係数は位置の偶奇だけに依存するので厳密2周期的である。

次元を計算する。$S_5$には四生成元の対角変更
$W\mapsto\alpha W,Y\mapsto\beta Y,Z\mapsto\gamma Z,X\mapsto\delta X$
が作用し、係数は
\[
 \begin{aligned}
 k_0&\mapsto k_0\alpha^2/(\beta\gamma),&
 v_0&\mapsto v_0\gamma^2/(\alpha\delta),\\
 b_0&\mapsto b_0\alpha^2/(\beta\delta),&
 s_0&\mapsto s_0\gamma^2/(\beta\delta),&
 a_0&\mapsto a_0\gamma\delta/(\alpha\beta)
 \end{aligned}
\]
と変化する。共通スカラーを除いた有効なトーラス作用の階数は3なので、
その一般軌道は3次元である。
従って同型類の像の次元は高々$5-3=2$である。
このトーラスの有理不変量として、例えば
\[
 I_1=\frac{b_0s_0}{k_0^2v_0^2},\qquad
 I_2=\frac{a_0b_0^4}{k_0^5v_0^3}
\]
を取れる。これらが既に\emph{全ての}indexed同型類を分類するとは主張しない。

逆向きの次元不等式はClifford基点での厳密線型代数から従う。
一つの関係空間上の$\PGL_4\times\PGL_4$軌道微分は階数30、
五パラメータを加えた微分は階数32である。
その偶方向は軌道と合わせて階数16、交代方向も階数16である。
軌道飽和像の次元は少なくとも32で、一般軌道は30次元だから
同型類の像の次元は少なくとも2である。
以上により正確に2となる。

さらに$a_0$の接方向は$YW$の書換え則へ$(-1)^iZX$を加えるだけで、
先に固定した交代商基底$o_0$そのものである。
九変数Ore族の交代接像は
$(o_0,o_1,o_2,o_3,o_4,o_5)=(0,-b,-a,-b,0,c)$なので、
$o_0$はその像に含まれない。
これはClifford基点での\emph{実際の新しい代数的変形方向}を証明する。
ただし他のOre分解を用いた表示や、全モジュライ成分への所属については未決定である。
\end{proof}

\begin{remark}
式\eqref{eq:nonore-five}では一般の$k_0,v_0,b_0,s_0$が零でもAS正則性は成立する。
これは九変数楕円族の$1+th\ne0$という制限とは別の全パラメータ主張である。
一方、二周期的表示を持つことから最小周期2とは結論できない。
一般の最小周期や、九変数族の別のチャートとの大域的な交わりは未解決である。
\end{remark}

\subsection{非Ore型族の一般点は最小周期2である}\label{subsec:nonore-period}
式\eqref{eq:nonore-five}について，この小節だけ
$(K,V,B,S,A)=(k_0,v_0,b_0,s_0,a_0)$と略記する．
点対応の双線型化を用いて，接方向だけでなく実際の
非1周期的な代数を得る．

\begin{theorem}[非Ore型五変数族の一般最小周期2]\label{thm:nonore-generic-two}
式\eqref{eq:nonore-five}の基底$\mathbb A^5$には非空なZariski開集合
$\mathcal U$があり，その上で最小正周期はちょうど$2$である．
したがって\emph{どのgraded algebraにも由来しない}二次元の
同型類の代数的族が得られる．さらにこの族の一般点の左点スキームは
一次元の支持を持つため，十二個の単純点と一個の非被約点を持つ
九変数楕円Clifford--Ore族の一般点とは非同型である．
ただし二つの族の閉包が異なる既約成分であるとは主張しない．
\end{theorem}
\begin{proof}
位置$i$で$\varepsilon=(-1)^i$とする．左点スキーム$E_\varepsilon$は，
始点座標$[w:y:z:x]\in\mathbb P^3$に対し，次の$6\times4$行列の
階数が$3$以下となる零点スキームである．
\begin{equation}\label{eq:nonore-point-matrix}
 M_\varepsilon=\begin{pmatrix}
 y&w&0&-\varepsilon Az\\
 z&0&w&0\\
 -Kw&z&y&0\\
 x&0&-Vz&w\\
 -Bw&x&-Sz&y\\
 0&0&x&z
 \end{pmatrix}.
\end{equation}
$AKVBS\ne0$とし，次の五次式は重根を持たないと仮定する．
\begin{equation}\label{eq:nonore-quintic}
 Q_\varepsilon(z)=\varepsilon AV^2z^5+4Sz^2-4B.
\end{equation}
全$4\times4$小行列式を調べると，$E_\varepsilon$の支持は一本の直線
\[
 \ell=\{w=z=0\}
\]
および，$Q_\varepsilon(z)=0$を満たす相異なる五つの点
\begin{equation}\label{eq:nonore-five-points}
 p_\varepsilon(z)=
 \left[1:-\frac K{2z}-\frac{\varepsilon AV}4z^3:
              z:-\frac V2z^2\right]
\end{equation}
に正確に一致する．必要性を得るには，例えば行列の行番号を$0$から数え，
$(1,2,3,5)$行の小行列式
$-z^2(Vz^2+2wx)$，$(0,1,2,5)$行の
$z(\varepsilon Axz^2-Kw^3-2wyz)$，$(0,1,4,5)$行の
$z(\varepsilon Ax^2z-Bw^3+Swz^2)$を用いればよい．
$z=0$なら$(0,1,2,3)$行の小行列式から$Kw^4=0$，従って$w=0$となる．
$z\ne0$では第一式により$w\ne0$なので$w=1$とでき，
後二式が\eqref{eq:nonore-five-points}と\eqref{eq:nonore-quintic}を与える．
逆に式\eqref{eq:nonore-five-points}を代入した全15小行列式は
$Q_\varepsilon$で割り切れる．$w=z=0$上でも全15小行列式は零になる．
これらの恒等式は\path{computations/verify_nonore_point_geometry.py}で，
パラメータを特殊化せず厳密に確認する．
従って$\ell$は唯一の正次元既約支持成分であり，
任意の$E_+\simeq E_-$の射影線型同型は$\ell$を固定する．

$Q_-( -z)=Q_+(z)$なので，$E_-$の五つの孤立点の
直線$\ell$からの射影座標は，$E_+$の五つの根の負である．
ここで一般の$Q_+$の根の集合には$\operatorname{PGL}_2$の非自明な
自己同型がないことを確認する．非空性は
\[
 K=V=A=1,\qquad B=-\tfrac14,\quad S=\tfrac14
\]
という有理数点で示せる．この点では$Q_+(z)=z^5+z^2+1$である．
この多項式は$\mathbb F_2$上既約で，$\mathbb F_5$上では
\[
 (z^2-z+2)(z^3+z^2-z-2)
\]
と，次数$(2,3)$の相異なる既約因子に分かれる．従ってガロア群は
推移的な5次置換群で，5-cycleと$(2,3)$型の置換を含み，
特にtranspositionを含む．推移的な5次置換群は素数次数ゆえ原始的であり，
原始置換群にtranspositionがあれば，その共役の作る辺グラフは連結で，
対応するtranspositionが$S_5$を生成する．従ってガロア群は$S_5$である．

五根の$\operatorname{PGL}_2$安定化群$\Gamma$は，五根への作用として
$S_5$へ忠実に埋まる．ガロア共役で$\Gamma$は保たれるから，
その像は$S_5$の正規部分群である．非自明なら$A_5$を含むが，
その五点作用で一点を固定する部分群は$A_4$を含む．
これは有限な$\operatorname{PGL}_2$部分群の一点安定化群が，
標数零では巡回群でなければならないことに矛盾する．
よって$\Gamma=1$である．異なる五点に対する非自明な置換が
$\operatorname{PGL}_2$で実現する条件は，各置換ごとの閉条件である．
五点の有限\'etale被覆から降下すれば，安定化群が自明な領域は
非空なZariski開集合となる．以下この開集合上で議論する．

$\ell$を固定する射影変換は商平面の射影直線
$[w:z]$にM\"obius変換を誘導する．
五根の集合をその負の集合へ送るM\"obius変換は，
上の安定化群の自明性から$z\mapsto-z$ただ一つである．
従ってもし$E_+\simeq E_-$なら，
$w,z$の二座標上では$[w:z]\mapsto[w:-z]$であり，
その$y$座標は五つの根上で
$1,z,z^2,y_+(z)$の線型結合として$y_-(-z)$に一致しなければならない．

ところが$Q_+(z)=0$より
\[
 z^{-1}=\frac{AV^2z^4+4Sz}{4B}
\]
である．従って$k[z]/(Q_+)$の部分空間
$\langle1,z,z^2\rangle$で割ると，
\[
 \begin{aligned}
 y_+(z)&\equiv-\frac{AV}{4}z^3-\frac{KAV^2}{8B}z^4,\\
 y_-(-z)&\equiv-\frac{AV}{4}z^3+\frac{KAV^2}{8B}z^4.
 \end{aligned}
\]
この二つの係数ベクトルの行列式は
\[
 -\frac{KA^2V^3}{16B}\ne0.
\]
すなわち後者は$1,z,z^2,y_+(z)$の線型包に入らない．
相異なる五根での関数の一致は$k[z]/(Q_+)$での一致と同値なので，
上で必要とされた射影変換は存在しない．
従って$E_+$と$E_-$は射影同型でない．

1周期同型があれば，隣接する位置の左点スキームが射影同型になるから，
これは1周期性を排除する．全ファイバーには厳密2周期があるので，
最小正周期は2である．同型類の像が二次元であることは
定理\ref{thm:nonore-five}ですでに証明した．
最後に，一般の九変数楕円Clifford--Ore族の左点スキームは
非被約一点と十二個の単純点という零次元支持を持つのに対し，
本族の左点スキームは直線$\ell$を含む．
同型は点スキームを保つので，両一般点は非同型である．
\end{proof}

\begin{remark}
一般の二周期性および九変数楕円族の一般点との非同型性は上で証明した．
別の三生成部分を選ぶと本族はOre表示を持ち，七変数族への閉包包含も
定理\ref{thm:nonore-in-seven-closure}で証明される．
両族が異なる既約成分をなすという主張はしない．
\end{remark}

\section{別のOre表示と七変数族への境界退化}\label{sec:ore-boundary}
「非Ore型」という呼称は，最初に選んだ三生成部分
$\langle W,Y,Z\rangle$に関するものであった．
別の三生成部分を選べば，この族はOre表示を持つ．
さらにその表示により，既存の七変数族への閉包関係を完全に決定できる．

\begin{theorem}[非Ore型二次元族は七変数族の閉包に含まれる]
\label{thm:nonore-in-seven-closure}
式\eqref{eq:nonore-five}の全パラメータから得られる幾何学的同型類の像は，
式\eqref{eq:seven-family}の七変数族の同型類の像のZariski閉包に含まれる．
正確には，一つの関係空間を記録する
$\Gr(6,k^4\otimes k^4)$で，七変数族の
$\PGL_4\times\PGL_4$軌道飽和の閉包に，五変数境界族の関係空間が全て属する．

実際，$v_0\ne0$上に，パラメータ$d$を持つ次の代数的平坦族が存在する．
式\eqref{eq:nonore-five}の最初の関係を
\begin{equation}\label{eq:ore-boundary-family}
 YW+WY=(-1)^i a_0ZX+d\left(Z^2+\frac{X^2}{v_0^2}\right)
\end{equation}
に置き換え，残りの五関係は変えない．
$d=0$は元の五変数境界族であり，$d\ne0$の各ファイバーは
七変数族に添字保持同型である．
\end{theorem}
\begin{proof}
$K=k_0,V=v_0,B=b_0,S=s_0,A=a_0$と書く．
$V\ne0$のもとで新しい生成元
\[
 \widehat W=Z,\qquad \widehat Y=W,\qquad
 \widehat Z=X/V,\qquad\widehat X=Y
\]
を取ると，全$d$において関係は
\begin{equation}\label{eq:ore-boundary-primed}
\begin{aligned}
 \widehat Y\widehat W+\widehat W\widehat Y&=0, &
 \widehat Z\widehat W+\widehat W\widehat Z&=0, &
 \widehat Z\widehat Y+\widehat Y\widehat Z&=\widehat W^2,\\
 \widehat X\widehat W+\widehat W\widehat X&=K\widehat Y^2,\\
 \widehat X\widehat Y+\widehat Y\widehat X
  &=d(\widehat W^2+\widehat Z^2)+(-1)^i AV\widehat W\widehat Z,\\
 \widehat X\widehat Z+\widehat Z\widehat X
  &=\frac BV\widehat Y^2+\frac SV\widehat W^2
\end{aligned}
\end{equation}
となる．上の三関係はClifford型三生成代数$B_{0,0}$を与え，
下の三関係は符号自己同型に関するOre拡大である．
通常導分の交代部分を$\partial$と書くと
$\partial(\widehat W)=\partial(\widehat Z)=0$，
$\partial(\widehat Y)=AV\widehat W\widehat Z$であり，
偶部分は三つの生成元を中心二次式に送る符号捩れ導分である．
既に証明したClifford--Oreの自由基底とフィルター付きKoszul--Frobenius性を適用できる．
また補助重み$(1,2,2,4)$のモニック二次規則を使うと，
全$d$で書換えが停止し，始点二通りの全八つの三次重なりが
$k[K,V,V^{-1},B,S,A,d]$上の恒等式として解消する．
従ってこの族は全次数で自由で，全ファイバーが標準14641型AS正則である．

$d\ne0$ではさらに
\[
 \widetilde W=\widehat W,\quad
 \widetilde Y=\widehat Y,\quad
 \widetilde Z=\widehat Z,\quad
 \widetilde X=\frac{\widehat X}{d}-\frac{S}{Vd}\widehat Y
\]
を取る．関係\eqref{eq:ore-boundary-primed}に代入すると，
七変数族\eqref{eq:seven-family}の正規形を得る．
具体的なパラメータ対応は
\begin{equation}\label{eq:boundary-to-seven}
 (u,v,r,s,a,b,c)
 =\left(\frac Kd,0,-\frac{2S}{Vd},
          \frac B{Vd},0,\frac{AV}{d},0\right).
\end{equation}
いずれも隣接する生成空間の可逆な基底変更であり，
添字の平行移動を同型として使っていない．

よって$V\ne0$を固定して$d$を零へ近づけると，
$d\ne0$のファイバーは七変数族と同型で，特殊ファイバーは五変数境界族である．
この操作は関係空間のGrassmannianの代数的な射で実現される．
$V=0$のファイバーも，非Ore型五変数族の$V\ne0$部分から
その関係空間を特殊化すれば得られる．
従って閉包への包含は全パラメータで成立する．
式と三次重なりの厳密検算は
\path{computations/verify_ore_boundary.py}に収めた．
\end{proof}

\begin{theorem}[七変数族の軌道飽和閉包はClifford点で特異]
\label{thm:seven-closure-singular}
$\Phi_7:\mathbb A^7\to\Gr(6,k^4\otimes k^4)$を，
七変数族の位置零の関係空間を記録する射とし，
$G=\PGL_4\times\PGL_4$と置く．
$Z_7=\overline{G\cdot\Phi_7(\mathbb A^7)}$は既約な37次元部分多様体である．
Clifford基点$z_0=\Phi_7(0,1,0,0,0,0,0)$において
\[
 \dim T_{z_0}Z_7\ge39>37=\dim Z_7.
\]
特に$Z_7$は$z_0$で特異である．
これは有限表示の軌道飽和閉包についての主張であり，
全AS正則モジュライの既約成分や粗モジュライの特異性を断定しない．
\end{theorem}
\begin{proof}
前節の厳密計算では，$G\times\mathbb A^7$からGrassmannianへの
微分の階数は$z_0$で37，$G$軌道の階数は30である．
源$G\times\mathbb A^7$は既約37次元なので，
$Z_7$は既約かつ次元37である．

Clifford基点の元の$W,Y,Z,X$座標で，次の二つの関係の接方向を考える．
\[
 \eta:\quad YW+WY\mapsto Z^2+X^2,\qquad
 \xi:\quad YW+WY\mapsto ZX.
\]
第一は\eqref{eq:ore-boundary-family}で
$K=V=B=S=1,A=0,d=q$とした$q=0$での導関数である．
第二は$K=V=B=S=1,A=q,d=q^2$とした導関数である．
$q\ne0$ではいずれも定理\ref{thm:nonore-in-seven-closure}によって
七変数族の$G$軌道に属するため，
$\eta,\xi\in T_{z_0}Z_7$である．

群作用の微分$[N_0\ N_1]$，七変数族のパラメータ微分$J_7$に対し，
$\mathbb Q$上の厳密階数計算は
\[
 \begin{aligned}
 \operatorname{rank}[N_0\ N_1]&=30,\\
 \operatorname{rank}[N_0\ N_1\ J_7]&=37,\\
 \operatorname{rank}[N_0\ N_1\ J_7\ \eta]&=38,\\
 \operatorname{rank}[N_0\ N_1\ J_7\ \xi]&=38,\\
 \operatorname{rank}[N_0\ N_1\ J_7\ \eta\ \xi]&=39.
 \end{aligned}
\]
これは\path{computations/verify_ore_boundary.py}が関係行列から直接計算する．
従って$T_{z_0}Z_7$には39次元の部分空間があり，
37次元の既約多様体のこの閉点は正則でない．
\end{proof}

\begin{remark}
以前の「九変数Ore族の選んだパラメータ切片では$o_0$方向が出ない」
という接空間計算は正しいが，
その事実から「$o_0$方向はOre族の閉包の外」という結論は出ない．
実際，基底変更が$d^{-1}$の極を持つ上の退化が別の接方向を作る．
また非Ore型五変数族の一般点と楕円Ore九変数族の一般点は
点スキームの次元により非同型であるが，
それと閉包への包含は矛盾しない．
\textbf{二つの族の像を，独立した既約成分として数えるべきではない．}
これらの像を含む全モジュライの既約成分の次元と最大性はなお未解決である．
\end{remark}

% !TeX root = ../papers/nc-p3.tex
\section{九変数族の軌道飽和閉包とClifford点の特異性}\label{sec:nine-closure}
九変数の楕円Clifford--Ore族の関係空間を位置$0$で取り、
\[
 \Phi_9:S_9\longrightarrow\Gr(6,k^4\otimes k^4),\qquad
 G=\PGL_4\times\PGL_4
\]
とする。ここで$S_9=\operatorname{Spec}k[u,v,r,s,a,b,c,t,h,(1+th)^{-1}]$である。
この射の像の$G$-軌道飽和のZariski閉包を
$Z_9=\overline{G\cdot\Phi_9(S_9)}$とし、七変数族について同様に$Z_7$を定める。
これは\emph{位置0の有限関係表示の空間}における閉包であり、
裸の$\mathbb Z$-代数の粗モジュライ空間であるとは主張しない。

\begin{theorem}[九変数族の閉包のClifford特異性]\label{thm:nine-closure-singular}
標数$0$の代数閉体上で
\[
 Z_7\subset Z_9,\qquad
 \dim Z_7=37,\qquad \dim Z_9=39
\]
が成立する。特に$Z_7$は$Z_9$の余次元$2$の真の閉部分多様体である。
Clifford基点$z_0$において
\[
 \boxed{\dim T_{z_0}Z_9\ge41>39=\dim Z_9}
\]
だから、$Z_9$は$z_0$で特異である。
さらに、$G$の基底変更による接方向を除くと、
$z_0$における\emph{実際の代数的一変数変形から得られる方向の線型包}は
少なくとも$11$次元であり、偶方向$7$個、交代方向$4$個を含む。
ただし、これら全方向を同時に積分する滑らかな$11$次元族は存在するとは主張しない。
\end{theorem}
\begin{proof}
七変数族は九変数族の$t=h=0$という閉切片だから$Z_7\subset Z_9$である。
$S_9$は既約な9次元の滑らかなスキームであり、$G$の次元は30である。
Clifford点で、$G\times S_9$からGrassmannianへの微分の階数は39である
（定理\ref{thm:algebraic-nine}の次元計算）。
従って像の閉包は39次元で既約である。
同様に七変数族では次元37となる（定理\ref{thm:seven-image}）。

前節の代数的退化\eqref{eq:ore-boundary-family}から得られる
二つの接方向を$\eta,\xi$とする。元の$W,Y,Z,X$の順序で、
$\eta$は$YW+WY$の右辺に$Z^2+X^2$を加える偶方向、
$\xi$は$YW+WY$に$(-1)^iZX$を加える交代方向である。
定理\ref{thm:nonore-in-seven-closure}と
\ref{thm:seven-closure-singular}により、両者は$T_{z_0}Z_7$、
従って$T_{z_0}Z_9$にも属する。

$N_0,N_1$を基底変更の$60\times16$行列、$J_9$を九変数の
関係係数の$60\times9$ Jacobianとする。Clifford点での階数は
\begin{equation}\label{eq:nine-closure-ranks}
\begin{aligned}
 \operatorname{rk}[N_0\ N_1]&=30,\\
 \operatorname{rk}[N_0\ N_1\ J_9]&=39,\\
 \operatorname{rk}[N_0\ N_1\ J_9\ \eta]&=40,\\
 \operatorname{rk}[N_0\ N_1\ J_9\ \xi]&=40,\\
 \operatorname{rk}[N_0\ N_1\ J_9\ \eta\ \xi]&=41.
\end{aligned}
\end{equation}
すべて有理数上の厳密階数であり、
\path{computations/verify_nine_closure_tangent.py}が標準四生成の関係行列から
再生成して検査する。従って$\dim T_{z_0}Z_9\ge41$を得て、特異性が従う。

最後に偶・交代のFourierモードを分ける。
九変数族の偶パラメータ$u,v,r,s,t,h$の微分$J_9^{+}$と
交代パラメータ$a,b,c$の微分$J_9^{-}$について
\[
 \begin{array}{c|ccc}
 &\operatorname{rk}D_1&\operatorname{rk}[D_1\ J_9]&
 \operatorname{rk}[D_1\ J_9\ \text{境界接方向}]\\\hline
 \text{偶}&15&21&22\\
 \text{交代}&15&18&19
 \end{array}
\]
である。従って基底変更で割った線型包には偶方向7個、交代方向4個が存在する。
九変数族の九方向は滑らかな代数的基底上で与えられ、追加の$\eta,\xi$は
前節の明示的な一変数退化で与えられるので、これらは全て\emph{個別には}
実際の代数的変形の接方向である。
\end{proof}

\begin{proposition}[11接方向の同時積分を妨げる二次障害]
\label{prop:nine-closure-obstruction}
上の$11$個の変形方向の線型包$L$をClifford基点の
基点付きindexed変形の$15$次元接空間に写す。
交代方向の座標を定理\ref{thm:clifford-quadratic}の
$(o_0,o_1,o_2,o_3,o_4,o_5)$とすると、$L$の交代部分は
\[
 (o_0,o_1,o_2,o_3,o_4,o_5)=(\lambda,-b,-a,-b,0,c)
\]
であり、二次障害は
\[
 \left.Q\right|_L=\lambda c.
\]
従って、$L$を接像とする\emph{滑らかな}基底上の基点付き平坦AS正則族は存在しない。
これは特異な基底上の同時積分や全モジュライの成分次元を排除するものではない。
\end{proposition}
\begin{proof}
七・九変数族の交代基底は式\eqref{eq:seven-odd-plane}で与えられ、
追加の$\xi$は$o_0$である。
従って式\eqref{eq:clifford-Q}へ代入すると
\[
 Q=o_0o_5+o_1o_2-o_2o_3+o_3o_4=\lambda c.
\]
接写像が$L$へ同型となる滑らかな局所基底があれば、その完備局所環で
$\lambda,c$は独立な一次座標となる。
変形の二次持上げの必要条件$Q=0$は一次座標の二次部分に
$\lambda c=0$を要求するが、滑らかな基底の形式冪級数環では
この積は零でない。矛盾である。
\end{proof}

\begin{remark}
$Z_9$の特異性は，九次元族とは異なる成分を構成したことを意味しない。
また，$Z_9$の接空間が大きいことをもって非1周期的な成分の
同型類次元が$11$以上であるとは結論しない。
本節の二つの追加方向のうち、$\xi$と九変数族の$c$方向は
二次で混合障害を持つ。閉包の大きな接空間は、滑らかな高次元族の存在とは
両立しない場合があることを明示する。
\end{remark}


\section{残る分類問題と検証範囲}
標準四窓への抽象的な還元は得られたが、一般四窓のKoszul/AS正則locusの
明示的な必要十分条件、正則な代数的族による全被覆、全閉特殊値、
相対族の大域的下降、全モジュライの既約成分は未解決である。
対称閉包の残余は定理\ref{thm:symmetric-closure-as-exact}で排除した。
残る構造分類は、その閉包外の低階数領域、特に$\rho=6$の未分類部分と、
$\rho\ge7$の高階数領域を含む。対称閉包での認識を全分類と同一視しない。

九変数族については、一般点の全変形と既存族の接像を比較し、
像が成分を尽くすかを決める必要がある。
楕円曲線・2等分点・12点因子と第四座標の持上げから、
式\eqref{eq:isogeny-modular-u}を超える十分な同型条件を求めることも課題である。
五変数境界族の一般最小周期2と七変数族への閉包包含は既に証明したので、
残るのは全AS正則locusにおける最大性と、閉包の外にある代数的変形である。
Clifford点では全1周期構造と高次障害が未決定であり、
特定のgraded表示の9次元接像を1周期的部分全体と同一視しない。

有限線型代数、全二次障害、PBW重なり、導分恒等式、
点配置の小行列式、2-isogenyの明示式と分岐、境界の基底変更には
有理数・多項式による検算コードを付す。
各計算は対応する有限恒等式を支え、一般非1周期性、全ファイバーAS正則性、
閉包包含はそれぞれ本文の構造的な証明に依存する。
\path{research/validation.json}に今回の論理監査、計算、組版、保存の範囲を分けて記録する。
全定理の形式化、第三者査読、新規性の確定、
外部原稿\cite{as4period}の全計算の再実行は行っていない。

\begin{thebibliography}{99}
\bibitem{VVRW}
M.~Vancliff, K.~Van Rompay and L.~Willaert,
\emph{Some quantum $\mathbb P^3$s with finitely many points},
Communications in Algebra \textbf{26} (1998), 1193--1208.
\url{https://bpb-us-e1.wpmucdn.com/websites.uta.edu/dist/0/5123/files/2022/10/finpts.pdf}.
使用箇所：Introduction, \S3（graded Clifford代数のOre変形）。
\bibitem{SV}
D.~R.~Stephenson and M.~Vancliff,
\emph{Constructing Clifford quantum $\mathbb P^3$s with finitely many points},
Journal of Algebra \textbf{312} (2007), 86--110.
\href{https://doi.org/10.1016/j.jalgebra.2007.02.015}{doi:10.1016/j.jalgebra.2007.02.015}.
本稿の全特殊化との同定は未確認。
\bibitem{BI}
V.~Bhatoy and C.~Ingalls,
\emph{Discrete Invariants of Koszul Artin-Schelter Regular Algebras of Dimension four},
\href{https://arxiv.org/abs/2602.13178v1}{arXiv:2602.13178v1} (2026).
使用箇所：Introduction, \S3, \S6。graded側の比較文献であり、
本稿のindexed成分数の根拠としては使わない。

\bibitem{BP} A.~I.~Bondal and A.~E.~Polishchuk,
\emph{Homological properties of associative algebras: the method of helices},
Russian Acad. Sci. Izv. Math. \textbf{42} (1994), no.~2, 219--260.
\href{https://doi.org/10.1070/IM1994v042n02ABEH001536}{doi:10.1070/IM1994v042n02ABEH001536}.
\bibitem{MoriNyman} I.~Mori and A.~Nyman,
\emph{A categorical characterization of quantum projective $\mathbb Z$-spaces},
Applied Categorical Structures \textbf{33} (2025).
\href{https://doi.org/10.1007/s10485-025-09806-2}{doi:10.1007/s10485-025-09806-2};
\href{https://arxiv.org/abs/2307.15253}{arXiv:2307.15253}.
\bibitem{Zhang} J.~J.~Zhang,
\emph{Twisted graded algebras and equivalences of graded categories},
Proc. Lond. Math. Soc. \textbf{72} (1996), 281--311.
\bibitem{Sierra} S.~J.~Sierra,
\emph{G-algebras, twistings, and equivalences of graded categories},
Algebras and Representation Theory \textbf{14} (2011), no.~2, 377--390.
\href{https://doi.org/10.1007/s10468-009-9193-y}{doi:10.1007/s10468-009-9193-y};
\href{https://arxiv.org/abs/math/0608791}{arXiv:math/0608791}.
\bibitem{VdB} M.~Van den Bergh, \emph{Non-commutative quadrics},
\href{https://arxiv.org/abs/0807.3753v3}{arXiv:0807.3753v3} (2010).
使用箇所：Proposition 3.3, Theorem 5.5.10, \S5.6.
\bibitem{BILN} V.~Bhatoy, C.~Ingalls, F.~LaRoche and R.~Nookala,
\emph{Some components of the moduli space of Koszul Artin-Schelter regular algebras of dimension four},
\href{https://arxiv.org/abs/2511.08390v1}{arXiv:2511.08390v1} (2025).
使用箇所：\S3.1, Table 2. Proposition 5.1の一般的な接空間判定は本稿では採用しない。
\bibitem{KLV}
V.~Knibbeler, S.~Lombardo and A.~P.~Veselov,
\emph{Automorphic Lie algebras and modular forms},
International Mathematics Research Notices \textbf{2023}, no.~6, 5209--5262.
\href{https://doi.org/10.1093/imrn/rnab376}{doi:10.1093/imrn/rnab376}.
使用箇所：\S7.3、式(42)のHesse標準形の$j$-不変量。
\bibitem{as4period} AS4周期性研究原稿,
\emph{四次元二次AS正則$\mathbb Z$-代数の周期性：最小周期1,2,4と非1周期的な平坦族},
2026年11月12日・第2版（私的原稿、独立査読・新規性未確認）.
\url{https://github.com/uedakazushi/as4/blob/main/research/z_periodicity/papers/as4_z_periodicity_ja.tex}.
確認済み本文のGit blob SHAは\path{research/BASELINE.md}を参照。
\bibitem{Bergman}
G.~M.~Bergman, \emph{The diamond lemma for ring theory},
Adv. Math. \textbf{29} (1978), 178--218.
\url{https://math.berkeley.edu/~gbergman/papers/updates/diamond.html}.
\bibitem{Priddy}
S.~B.~Priddy, \emph{Koszul resolutions}, Trans. Amer. Math. Soc.
\textbf{152} (1970), 39--60.
\url{https://math.mit.edu/~hrm/palestine/priddy-koszul-resolutions.pdf}.
\bibitem{BS2006} G.~R.~Belitskii and V.~V.~Sergeichuk,
\emph{Congruence of multilinear forms}, Linear Algebra Appl. \textbf{418} (2006), 751--762.
\href{https://doi.org/10.1016/j.laa.2006.03.003}{doi:10.1016/j.laa.2006.03.003};
\href{https://arxiv.org/abs/0710.0834}{arXiv:0710.0834}.
使用箇所：Theorem 2(a), Lemma 4(a)（複素数体上の主張）。
\bibitem{DDLpencils} F.~De Ter\'an, F.~M.~Dopico and J.~M.~Landsberg,
\emph{An explicit description of the irreducible components of the set of matrix pencils with bounded normal rank},
著者原稿（2017年1月27日版）,
\url{https://gauss.uc3m.es/fdopico/papers/laa2016-2.pdf}.
使用箇所：Theorem 1（Kronecker標準形）と零最小指数の説明。
\bibitem{FIstability} M.~Fedorchuk and A.~Isaev,
\emph{Stability of associated forms},
\href{https://arxiv.org/abs/1703.00438v5}{arXiv:1703.00438v5} (2018).
使用箇所：Theorem 1.1（入力のpolystabilityを仮定）。
\bibitem{FIassociated} M.~Fedorchuk and A.~Isaev,
\emph{Associated form morphism},
\href{https://arxiv.org/abs/1807.02082v2}{arXiv:1807.02082v2} (2018).
使用箇所：\S2（対角群のassociated-form射）。
\bibitem{HIOrep}
M.~Herschend, O.~Iyama and S.~Oppermann, $n$-representation infinite algebras,
\emph{Advances in Mathematics} \textbf{252} (2014), 292--342.
\href{https://arxiv.org/abs/1205.1272}{arXiv:1205.1272}.
\bibitem{GrantIyama}
J.~Grant and O.~Iyama, Higher preprojective algebras, Koszul algebras, and superpotentials,
\href{https://arxiv.org/abs/1902.07878}{arXiv:1902.07878}.
\bibitem{GrantZigzag}
J.~Grant, Higher zigzag algebras,
\emph{Documenta Mathematica} \textbf{24} (2019), 749--814.
\href{https://doi.org/10.4171/DM/693}{doi:10.4171/DM/693}.
\bibitem{PSVQuartics}
D.~Plaumann, B.~Sturmfels, C.~Vinzant, Quartic curves and their bitangents,
\emph{Journal of Symbolic Computation} \textbf{46} (2011), 712--733.
\href{https://doi.org/10.1016/j.jsc.2011.01.007}{doi:10.1016/j.jsc.2011.01.007}.
\end{thebibliography}
\end{document}
